% ---------------------------------------------------------------------------
% Preamble for the generated book. Typeset with XeLaTeX.
% Design brief: match the interior conventions of O'Reilly technical books --
% 7 x 9.1875in trim, serif body, humanist-sans headings, tracked-out chapter
% labels, ruled admonitions with margin icons, hairline booktabs tables and
% tinted syntax-highlighted code listings.
% ---------------------------------------------------------------------------
\documentclass[twoside,openright]{book}

\usepackage[fontsize=10.5pt]{fontsize}
\usepackage[
  paperwidth=7in, paperheight=9.1875in,
  inner=1.0in, outer=0.72in, top=0.8in, bottom=0.9in,
  headsep=13pt, footskip=26pt
]{geometry}

\usepackage{fontspec}
\usepackage[table]{xcolor}
\usepackage{graphicx}
\usepackage{booktabs}
\usepackage{array}
\usepackage{tabularx}
\usepackage{xltabular}
\usepackage{enumitem}
\usepackage{fancyhdr}
\usepackage{titlesec}
\usepackage{titletoc}
\usepackage{fvextra}
\usepackage[breakable,skins]{tcolorbox}
\usepackage{fontawesome5}
\usepackage[normalem]{ulem}   % normalem: keep \emph italic, only add \sout
\usepackage{needspace}
\usepackage{emptypage}
\usepackage[protrusion=true,expansion=false]{microtype}
\usepackage[hyphens]{url}

% --- Palette -------------------------------------------------------------
\definecolor{ink}{HTML}{16191C}
\definecolor{headgray}{HTML}{55606A}
\definecolor{rulegray}{HTML}{AEB6BD}
\definecolor{hairline}{HTML}{CDD3D8}
\definecolor{numbergray}{HTML}{BFC6CC}
\definecolor{codebg}{HTML}{F6F7F8}
\definecolor{codeframe}{HTML}{DFE3E6}
\definecolor{tablehdr}{HTML}{EDEFF1}
\definecolor{linkblue}{HTML}{1C5C99}
\definecolor{noteblue}{HTML}{2A6496}
\definecolor{tipgreen}{HTML}{2F6B45}
\definecolor{warnred}{HTML}{A03328}

% --- Fonts ---------------------------------------------------------------
% Ligatures=TeX (i.e. Mapping=tex-text) is applied to the text fonts only.
% It must never reach the monospace font: XeTeX applies font mappings even to
% verbatim material, which would turn straight quotes in code into curly ones.
\setmainfont{Erewhon}[
  Extension      = .otf,
  Ligatures      = TeX,
  UprightFont    = *-Regular,
  ItalicFont     = *-Italic,
  BoldFont       = *-Bold,
  BoldItalicFont = *-BoldItalic ]
\setsansfont{SourceSansPro}[
  Extension      = .otf,
  Ligatures      = TeX,
  UprightFont    = *-Regular,
  ItalicFont     = *-RegularIt,
  BoldFont       = *-Semibold,
  BoldItalicFont = *-SemiboldIt ]
\setmonofont{SourceCodePro}[
  Extension      = .otf,
  Ligatures      = NoCommon,
  UprightFont    = *-Regular,
  ItalicFont     = *-RegularIt,
  BoldFont       = *-Semibold,
  BoldItalicFont = *-SemiboldIt,
  Scale          = 0.94 ]
% Tracked-out sans for chapter labels and running heads.
\newfontfamily\labelfont{SourceSansPro}[
  Extension = .otf, Ligatures = TeX, UprightFont = *-Semibold, LetterSpace = 7.0 ]
\newfontfamily\numeralfont{SourceSansPro}[
  Extension = .otf, Ligatures = TeX, UprightFont = *-Light ]

% A handful of CJK glyphs appear in the localization examples; map them to a
% font that covers them. newunicodechar substitutions also work inside
% verbatim material, which is where these occur.
\usepackage{newunicodechar}
\newfontfamily\cjkfont{Noto Sans Mono CJK JP}[Scale=MatchLowercase]
\newunicodechar{年}{{\cjkfont 年}}
\newunicodechar{月}{{\cjkfont 月}}
\newunicodechar{日}{{\cjkfont 日}}
\newunicodechar{￥}{{\cjkfont ￥}}

% Arrows/comparators used in prose, tables, and ASCII diagrams that the text
% and mono fonts don't carry glyphs for.
\newunicodechar{→}{\ensuremath{\rightarrow}}
\newunicodechar{≥}{\ensuremath{\geq}}

\linespread{1.07}
\setlength{\parindent}{0pt}
\setlength{\parskip}{6.5pt plus 2pt minus 1pt}
\setlength{\emergencystretch}{2em}
\widowpenalty=10000
\clubpenalty=10000
\brokenpenalty=10000
\raggedbottom
\frenchspacing

% --- Inline code ---------------------------------------------------------
% Straight quotes, no ligatures, and break opportunities inside long
% identifiers so they never run into the margin.
\newcommand{\code}[1]{\texttt{#1}}
\newcommand{\headcode}[1]{\texttt{#1}}

% --- Headings ------------------------------------------------------------
\setcounter{secnumdepth}{0}   % O'Reilly interiors leave sections unnumbered
\setcounter{tocdepth}{1}

\titleformat{\chapter}[display]
  {\sffamily\raggedright}
  {\labelfont\fontsize{9}{9}\selectfont\color{headgray}%
   \MakeUppercase{\chaptertitlename}\ \MakeUppercase{\thechapter}}
  {10pt}
  {\color{ink}\sffamily\bfseries\fontsize{25}{29}\selectfont}
  [{\vspace{9pt}{\color{rulegray}\titlerule[1pt]}}]
\titlespacing*{\chapter}{0pt}{16pt}{30pt}

\titleformat{\section}
  {\sffamily\bfseries\fontsize{14.5}{17}\selectfont\color{ink}\raggedright}{}{0pt}{}
\titlespacing*{\section}{0pt}{19pt plus 5pt minus 3pt}{6pt}

\titleformat{\subsection}
  {\sffamily\bfseries\fontsize{11.8}{14}\selectfont\color{ink}\raggedright}{}{0pt}{}
\titlespacing*{\subsection}{0pt}{15pt plus 4pt minus 2pt}{4pt}

\titleformat{\subsubsection}
  {\sffamily\bfseries\itshape\fontsize{10.5}{13}\selectfont\color{ink}\raggedright}{}{0pt}{}
\titlespacing*{\subsubsection}{0pt}{12pt plus 3pt minus 2pt}{3pt}

% --- Running heads -------------------------------------------------------
\newcommand{\runhead}[1]{{\sffamily\fontsize{8.5}{10}\selectfont\color{headgray}#1}}
\newcommand{\folio}{{\sffamily\bfseries\fontsize{9}{11}\selectfont\color{ink}\thepage}}
\renewcommand{\chaptermark}[1]{\markboth{#1}{#1}}
\renewcommand{\sectionmark}[1]{\markright{#1}}
\pagestyle{fancy}
\fancyhf{}
\fancyhead[LE]{\folio\hspace{1em}\runhead{\leftmark}}
\fancyhead[RO]{\runhead{\rightmark}\hspace{1em}\folio}
\renewcommand{\headrulewidth}{0.4pt}
\renewcommand{\headrule}{{\color{hairline}%
  \hrule height \headrulewidth width\headwidth}}
\fancypagestyle{plain}{%
  \fancyhf{}\renewcommand{\headrulewidth}{0pt}%
  \fancyfoot[LE,RO]{\folio}}

% --- Lists ---------------------------------------------------------------
\setlist{topsep=5pt, itemsep=3pt, parsep=2pt, leftmargin=1.35em}
\setlist[itemize]{label=\raisebox{0.12ex}{\scriptsize\color{headgray}$\bullet$}}
\setlist[itemize,2]{label=\raisebox{0.12ex}{\scriptsize\color{headgray}$-$}}
\setlist[enumerate]{font=\sffamily\bfseries}

% --- Code listings -------------------------------------------------------
\newcommand{\codesize}{\fontsize{8.5}{10.4}\selectfont}
% xleftmargin reserves a gutter inside the listing so that wrapped lines and
% the continuation arrow stay within the tinted box instead of hanging into
% the page margin.
\fvset{
  fontsize=\codesize,
  breaklines=true,
  breakanywhere=true,
  breakautoindent=false,
  breakindent=0pt,
  breaksymbolleft=\hbox{\color{numbergray}\tiny\raisebox{0.15ex}{$\hookrightarrow$}},
  breaksymbolsepleft=3pt,
  breaksymbolindentleft=0pt,
  xleftmargin=11pt
}
\newtcolorbox{codebox}{
  breakable, enhanced jigsaw, sharp corners,
  colback=codebg, colframe=codeframe, boxrule=0.4pt,
  left=1pt, right=5pt, top=5pt, bottom=5pt,
  before skip=9pt plus 2pt, after skip=11pt plus 2pt,
  pad at break*=3pt, toprule at break=0pt, bottomrule at break=0pt
}

% --- Admonitions ---------------------------------------------------------
% Rules above and below, icon hanging in the left indent; fully breakable
% across pages, unlike a boxed environment.
\newlength{\admonindent}
\setlength{\admonindent}{30pt}
\newenvironment{admon}[2]{%
  \par\addvspace{11pt}%
  \noindent{\color{rulegray}\rule{\linewidth}{0.4pt}}\par\nobreak\vspace{4pt}%
  \begingroup
  \fontsize{9.6}{12.4}\selectfont
  \setlength{\leftskip}{\admonindent}%
  \noindent\hspace*{-\admonindent}%
  \makebox[\admonindent][l]{\raisebox{-1.2pt}{\color{#1}\large #2}}%
  \ignorespaces
}{%
  \par\endgroup\vspace{4pt}%
  \noindent{\color{rulegray}\rule{\linewidth}{0.4pt}}\par\addvspace{11pt}%
}
\newenvironment{notebox}{\begin{admon}{noteblue}{\faInfoCircle}}{\end{admon}}
\newenvironment{tipbox}{\begin{admon}{tipgreen}{\faLightbulb}}{\end{admon}}
\newenvironment{warnbox}{\begin{admon}{warnred}{\faExclamationTriangle}}{\end{admon}}

% --- Sidebar (for the "show the review" solution blocks) -----------------
\newenvironment{sidebar}[1]{%
  \par\addvspace{13pt}%
  \noindent{\color{headgray}\rule{\linewidth}{1pt}}\par\nobreak\vspace{5pt}%
  \noindent{\labelfont\fontsize{8.6}{10}\selectfont\color{headgray}\MakeUppercase{#1}}%
  \par\nobreak\vspace{3pt}%
  \begingroup\fontsize{9.8}{12.6}\selectfont
}{%
  \par\endgroup\vspace{5pt}%
  \noindent{\color{headgray}\rule{\linewidth}{1pt}}\par\addvspace{13pt}%
}

% --- Tables --------------------------------------------------------------
\newcolumntype{L}[1]{>{\hsize=#1\hsize\raggedright\arraybackslash}X}
\newcolumntype{C}[1]{>{\hsize=#1\hsize\centering\arraybackslash}X}
\newcolumntype{R}[1]{>{\hsize=#1\hsize\raggedleft\arraybackslash}X}
\newcommand{\thd}[1]{{\sffamily\bfseries\color{ink}#1}}
\renewcommand{\arraystretch}{1.22}
\setlength{\tabcolsep}{5pt}
\setlength{\LTpre}{10pt}
\setlength{\LTpost}{12pt}
\arrayrulecolor{ink}
\newcommand{\tablefont}{\fontsize{9}{11.4}\selectfont}
\newcommand{\tablefontsmall}{\fontsize{8.2}{10.4}\selectfont}

% --- Thematic break ------------------------------------------------------
\newcommand{\ornament}{\par\addvspace{8pt}%
  \centerline{\color{rulegray}\sffamily\footnotesize $*$\quad$*$\quad$*$}%
  \par\addvspace{8pt}}

% --- Title page helpers --------------------------------------------------
\newcommand{\booktitleblock}[3]{%
  \thispagestyle{empty}%
  \vspace*{1.15in}%
  {\raggedright
   {\sffamily\bfseries\fontsize{34}{38}\selectfont\color{ink}#1\par}
   \vspace{12pt}
   {\color{rulegray}\rule{\linewidth}{1.2pt}}\par
   \vspace{12pt}
   {\sffamily\fontsize{14}{18}\selectfont\color{headgray}#2\par}
   \vfill
   {\labelfont\fontsize{10}{12}\selectfont\color{ink}\MakeUppercase{#3}\par}
   \vspace{0.35in}}%
  \clearpage}


% --- Pygments style definitions ---

\makeatletter
\def\PY@reset{\let\PY@it=\relax \let\PY@bf=\relax%
    \let\PY@ul=\relax \let\PY@tc=\relax%
    \let\PY@bc=\relax \let\PY@ff=\relax}
\def\PY@tok#1{\csname PY@tok@#1\endcsname}
\def\PY@toks#1+{\ifx\relax#1\empty\else%
    \PY@tok{#1}\expandafter\PY@toks\fi}
\def\PY@do#1{\PY@bc{\PY@tc{\PY@ul{%
    \PY@it{\PY@bf{\PY@ff{#1}}}}}}}
\def\PY#1#2{\PY@reset\PY@toks#1+\relax+\PY@do{#2}}

\@namedef{PY@tok@}{\def\PY@tc##1{\textcolor[rgb]{0.09,0.10,0.11}{##1}}}
\@namedef{PY@tok@c}{\let\PY@it=\textit\def\PY@tc##1{\textcolor[rgb]{0.36,0.40,0.44}{##1}}}
\@namedef{PY@tok@cp}{\def\PY@tc##1{\textcolor[rgb]{0.58,0.22,0.00}{##1}}}
\@namedef{PY@tok@k}{\let\PY@bf=\textbf\def\PY@tc##1{\textcolor[rgb]{0.66,0.13,0.23}{##1}}}
\@namedef{PY@tok@kt}{\let\PY@bf=\textbf\def\PY@tc##1{\textcolor[rgb]{0.66,0.13,0.23}{##1}}}
\@namedef{PY@tok@kc}{\let\PY@bf=\textbf\def\PY@tc##1{\textcolor[rgb]{0.66,0.13,0.23}{##1}}}
\@namedef{PY@tok@o}{\def\PY@tc##1{\textcolor[rgb]{0.09,0.10,0.11}{##1}}}
\@namedef{PY@tok@ow}{\let\PY@bf=\textbf\def\PY@tc##1{\textcolor[rgb]{0.66,0.13,0.23}{##1}}}
\@namedef{PY@tok@p}{\def\PY@tc##1{\textcolor[rgb]{0.09,0.10,0.11}{##1}}}
\@namedef{PY@tok@n}{\def\PY@tc##1{\textcolor[rgb]{0.09,0.10,0.11}{##1}}}
\@namedef{PY@tok@nb}{\def\PY@tc##1{\textcolor[rgb]{0.04,0.31,0.55}{##1}}}
\@namedef{PY@tok@bp}{\let\PY@it=\textit\def\PY@tc##1{\textcolor[rgb]{0.04,0.31,0.55}{##1}}}
\@namedef{PY@tok@nc}{\def\PY@tc##1{\textcolor[rgb]{0.54,0.29,0.00}{##1}}}
\@namedef{PY@tok@nn}{\def\PY@tc##1{\textcolor[rgb]{0.54,0.29,0.00}{##1}}}
\@namedef{PY@tok@ne}{\def\PY@tc##1{\textcolor[rgb]{0.54,0.29,0.00}{##1}}}
\@namedef{PY@tok@nf}{\def\PY@tc##1{\textcolor[rgb]{0.42,0.20,0.66}{##1}}}
\@namedef{PY@tok@nd}{\def\PY@tc##1{\textcolor[rgb]{0.54,0.29,0.00}{##1}}}
\@namedef{PY@tok@na}{\def\PY@tc##1{\textcolor[rgb]{0.04,0.31,0.55}{##1}}}
\@namedef{PY@tok@nt}{\def\PY@tc##1{\textcolor[rgb]{0.07,0.35,0.17}{##1}}}
\@namedef{PY@tok@no}{\def\PY@tc##1{\textcolor[rgb]{0.04,0.31,0.55}{##1}}}
\@namedef{PY@tok@nv}{\def\PY@tc##1{\textcolor[rgb]{0.09,0.10,0.11}{##1}}}
\@namedef{PY@tok@l}{\def\PY@tc##1{\textcolor[rgb]{0.04,0.31,0.55}{##1}}}
\@namedef{PY@tok@s}{\def\PY@tc##1{\textcolor[rgb]{0.04,0.23,0.44}{##1}}}
\@namedef{PY@tok@se}{\let\PY@bf=\textbf\def\PY@tc##1{\textcolor[rgb]{0.04,0.23,0.44}{##1}}}
\@namedef{PY@tok@sd}{\let\PY@it=\textit\def\PY@tc##1{\textcolor[rgb]{0.36,0.40,0.44}{##1}}}
\@namedef{PY@tok@m}{\def\PY@tc##1{\textcolor[rgb]{0.04,0.31,0.55}{##1}}}
\@namedef{PY@tok@g}{\def\PY@tc##1{\textcolor[rgb]{0.09,0.10,0.11}{##1}}}
\@namedef{PY@tok@gp}{\def\PY@tc##1{\textcolor[rgb]{0.36,0.40,0.44}{##1}}}
\@namedef{PY@tok@go}{\def\PY@tc##1{\textcolor[rgb]{0.09,0.10,0.11}{##1}}}
\@namedef{PY@tok@ge}{\let\PY@it=\textit\def\PY@tc##1{\textcolor[rgb]{0.09,0.10,0.11}{##1}}}
\@namedef{PY@tok@gs}{\let\PY@bf=\textbf\def\PY@tc##1{\textcolor[rgb]{0.09,0.10,0.11}{##1}}}
\@namedef{PY@tok@err}{\def\PY@tc##1{\textcolor[rgb]{0.66,0.13,0.23}{##1}}}
\@namedef{PY@tok@w}{\def\PY@tc##1{\textcolor[rgb]{0.09,0.10,0.11}{##1}}}
\@namedef{PY@tok@esc}{\def\PY@tc##1{\textcolor[rgb]{0.09,0.10,0.11}{##1}}}
\@namedef{PY@tok@x}{\def\PY@tc##1{\textcolor[rgb]{0.09,0.10,0.11}{##1}}}
\@namedef{PY@tok@kd}{\let\PY@bf=\textbf\def\PY@tc##1{\textcolor[rgb]{0.66,0.13,0.23}{##1}}}
\@namedef{PY@tok@kn}{\let\PY@bf=\textbf\def\PY@tc##1{\textcolor[rgb]{0.66,0.13,0.23}{##1}}}
\@namedef{PY@tok@kp}{\let\PY@bf=\textbf\def\PY@tc##1{\textcolor[rgb]{0.66,0.13,0.23}{##1}}}
\@namedef{PY@tok@kr}{\let\PY@bf=\textbf\def\PY@tc##1{\textcolor[rgb]{0.66,0.13,0.23}{##1}}}
\@namedef{PY@tok@ni}{\def\PY@tc##1{\textcolor[rgb]{0.09,0.10,0.11}{##1}}}
\@namedef{PY@tok@fm}{\def\PY@tc##1{\textcolor[rgb]{0.42,0.20,0.66}{##1}}}
\@namedef{PY@tok@py}{\def\PY@tc##1{\textcolor[rgb]{0.09,0.10,0.11}{##1}}}
\@namedef{PY@tok@nl}{\def\PY@tc##1{\textcolor[rgb]{0.09,0.10,0.11}{##1}}}
\@namedef{PY@tok@nx}{\def\PY@tc##1{\textcolor[rgb]{0.09,0.10,0.11}{##1}}}
\@namedef{PY@tok@vc}{\def\PY@tc##1{\textcolor[rgb]{0.09,0.10,0.11}{##1}}}
\@namedef{PY@tok@vg}{\def\PY@tc##1{\textcolor[rgb]{0.09,0.10,0.11}{##1}}}
\@namedef{PY@tok@vi}{\def\PY@tc##1{\textcolor[rgb]{0.09,0.10,0.11}{##1}}}
\@namedef{PY@tok@vm}{\def\PY@tc##1{\textcolor[rgb]{0.09,0.10,0.11}{##1}}}
\@namedef{PY@tok@ld}{\def\PY@tc##1{\textcolor[rgb]{0.04,0.31,0.55}{##1}}}
\@namedef{PY@tok@sa}{\def\PY@tc##1{\textcolor[rgb]{0.04,0.23,0.44}{##1}}}
\@namedef{PY@tok@sb}{\def\PY@tc##1{\textcolor[rgb]{0.04,0.23,0.44}{##1}}}
\@namedef{PY@tok@sc}{\def\PY@tc##1{\textcolor[rgb]{0.04,0.23,0.44}{##1}}}
\@namedef{PY@tok@dl}{\def\PY@tc##1{\textcolor[rgb]{0.04,0.23,0.44}{##1}}}
\@namedef{PY@tok@s2}{\def\PY@tc##1{\textcolor[rgb]{0.04,0.23,0.44}{##1}}}
\@namedef{PY@tok@sh}{\def\PY@tc##1{\textcolor[rgb]{0.04,0.23,0.44}{##1}}}
\@namedef{PY@tok@si}{\def\PY@tc##1{\textcolor[rgb]{0.04,0.23,0.44}{##1}}}
\@namedef{PY@tok@sx}{\def\PY@tc##1{\textcolor[rgb]{0.04,0.23,0.44}{##1}}}
\@namedef{PY@tok@sr}{\def\PY@tc##1{\textcolor[rgb]{0.04,0.23,0.44}{##1}}}
\@namedef{PY@tok@s1}{\def\PY@tc##1{\textcolor[rgb]{0.04,0.23,0.44}{##1}}}
\@namedef{PY@tok@ss}{\def\PY@tc##1{\textcolor[rgb]{0.04,0.23,0.44}{##1}}}
\@namedef{PY@tok@mb}{\def\PY@tc##1{\textcolor[rgb]{0.04,0.31,0.55}{##1}}}
\@namedef{PY@tok@mf}{\def\PY@tc##1{\textcolor[rgb]{0.04,0.31,0.55}{##1}}}
\@namedef{PY@tok@mh}{\def\PY@tc##1{\textcolor[rgb]{0.04,0.31,0.55}{##1}}}
\@namedef{PY@tok@mi}{\def\PY@tc##1{\textcolor[rgb]{0.04,0.31,0.55}{##1}}}
\@namedef{PY@tok@il}{\def\PY@tc##1{\textcolor[rgb]{0.04,0.31,0.55}{##1}}}
\@namedef{PY@tok@mo}{\def\PY@tc##1{\textcolor[rgb]{0.04,0.31,0.55}{##1}}}
\@namedef{PY@tok@pm}{\def\PY@tc##1{\textcolor[rgb]{0.09,0.10,0.11}{##1}}}
\@namedef{PY@tok@ch}{\let\PY@it=\textit\def\PY@tc##1{\textcolor[rgb]{0.36,0.40,0.44}{##1}}}
\@namedef{PY@tok@cm}{\let\PY@it=\textit\def\PY@tc##1{\textcolor[rgb]{0.36,0.40,0.44}{##1}}}
\@namedef{PY@tok@cpf}{\let\PY@it=\textit\def\PY@tc##1{\textcolor[rgb]{0.36,0.40,0.44}{##1}}}
\@namedef{PY@tok@c1}{\let\PY@it=\textit\def\PY@tc##1{\textcolor[rgb]{0.36,0.40,0.44}{##1}}}
\@namedef{PY@tok@cs}{\let\PY@it=\textit\def\PY@tc##1{\textcolor[rgb]{0.36,0.40,0.44}{##1}}}
\@namedef{PY@tok@gd}{\def\PY@tc##1{\textcolor[rgb]{0.09,0.10,0.11}{##1}}}
\@namedef{PY@tok@gr}{\def\PY@tc##1{\textcolor[rgb]{0.09,0.10,0.11}{##1}}}
\@namedef{PY@tok@gh}{\def\PY@tc##1{\textcolor[rgb]{0.09,0.10,0.11}{##1}}}
\@namedef{PY@tok@gi}{\def\PY@tc##1{\textcolor[rgb]{0.09,0.10,0.11}{##1}}}
\@namedef{PY@tok@gu}{\def\PY@tc##1{\textcolor[rgb]{0.09,0.10,0.11}{##1}}}
\@namedef{PY@tok@ges}{\def\PY@tc##1{\textcolor[rgb]{0.09,0.10,0.11}{##1}}}
\@namedef{PY@tok@gt}{\def\PY@tc##1{\textcolor[rgb]{0.09,0.10,0.11}{##1}}}

\def\PYZbs{\char`\\}
\def\PYZus{\char`\_}
\def\PYZob{\char`\{}
\def\PYZcb{\char`\}}
\def\PYZca{\char`\^}
\def\PYZam{\char`\&}
\def\PYZlt{\char`\<}
\def\PYZgt{\char`\>}
\def\PYZsh{\char`\#}
\def\PYZpc{\char`\%}
\def\PYZdl{\char`\$}
\def\PYZhy{\char`\-}
\def\PYZsq{\char`\'}
\def\PYZdq{\char`\"}
\def\PYZti{\char`\~}
% for compatibility with earlier versions
\def\PYZat{@}
\def\PYZlb{[}
\def\PYZrb{]}
\makeatother


\usepackage[unicode,bookmarksnumbered,bookmarksdepth=2,colorlinks=true,linkcolor=linkblue,urlcolor=linkblue,citecolor=linkblue,pdfborder={0 0 0}]{hyperref}
\hypersetup{pdftitle={Core Java Preparation},pdfauthor={Mahdi Sadeghi},pdfsubject={A Code-Review Interview Companion for Java 21+},pdfcreator={XeLaTeX}}
\begin{document}
\frontmatter
\booktitleblock{Core Java Preparation}{A Code-Review Interview Companion for Java 21+}{Mahdi Sadeghi}
\thispagestyle{empty}
\vspace*{\fill}
{\raggedright\fontsize{9}{12.5}\selectfont\color{headgray}
\textbf{Core Java Preparation}\par\vspace{3pt}
by Mahdi Sadeghi.\par\vspace{9pt}
Typeset from the Markdown sources in \code{chapters/} on 08 August 2026: 22 chapters, roughly 157,778 words. Rebuild with \code{python3 tools/build\_pdf.py}.\par\vspace{9pt}
Body text is set in Erewhon, headings in Source Sans Pro, and code in Source Code Pro. Code listings are highlighted with Pygments and the book is typeset with XeLaTeX.\par\vspace{9pt}
Examples target Java 21 and later. Code in this book is intended for study and may be used freely.\par}
\clearpage
\pdfbookmark[0]{Contents}{toc}
\tableofcontents
\cleardoublepage

\chapter*{Preface}
\addcontentsline{toc}{chapter}{Preface}
\markboth{Preface}{Preface}

Each topic in this book is a self-contained chapter. Every chapter explains its
subtopics in plain language with runnable Java examples, and ends with a
\emph{Common Code-Review Interview Pitfalls} section that collects the mistakes
reviewers actually catch.

The material is written for someone who already knows Java but wants to review
it sharply --- the kind of review you do before a code-review-style interview.
Examples target Java~21 and later, with notes where older or newer releases
differ.

\section{Conventions Used in This Book}

The following typographic conventions are used:

\begin{description}[leftmargin=0pt, style=nextline, itemsep=5pt]
  \item[\emph{Italic}] Indicates emphasis and new terms where they are first
    defined.
  \item[\code{Constant width}] Used for program listings, as well as within
    paragraphs to refer to program elements such as variable and method names,
    data types, statements, keywords, class names, and command-line tools.
\end{description}

Code listings are set on a tinted background and syntax-highlighted. Where a
single line of code or shell output is too long for the measure, it wraps and
the continuation is marked with a $\hookrightarrow$ symbol in the left margin.

\begin{notebox}
This element signifies a general note.
\end{notebox}

\begin{tipbox}
This element signifies a tip or suggestion.
\end{tipbox}

\begin{warnbox}
This element indicates a warning or a caution.
\end{warnbox}

Sections marked off by rules and a small heading --- like \textsc{the review}
blocks in the final chapter --- hold worked answers. Try the exercise before
reading them.

\mainmatter
\setcounter{chapter}{0}
\chapter{Java Fundamentals}
\label{chap:1}
\label{h:1:1-java-fundamentals}
Every code review starts with the basics: how data is stored, how variables behave, and how control flow is structured. Get these wrong and everything built on top of them — collections, streams, concurrency — inherits the bug. This chapter is a fast refresher on the building blocks of Java, written for someone who already knows the language but wants to review it sharply before a code-review-style interview. Examples target Java 21+, with notes where older or newer versions differ.

\section[Primitive Data Types]{Primitive Data Types}
\label{h:1:primitive-data-types}
A \textbf{primitive type} is a basic building block for data — not an object, just a raw value stored directly in memory (on the stack, or inline inside an object). Java has eight of them. They’re fast and predictable because there’s no object overhead: no header, no pointer indirection, no garbage collection tracking.

{\tablefont
\begin{xltabular}{\linewidth}{@{}L{0.462}L{0.538}L{2.385}L{0.615}@{}}
\toprule
\rowcolor{tablehdr}
\thd{Type} & \thd{Size} & \thd{Range} & \thd{Default value} \\
\midrule
\endhead
\code{byte} & 8 bits & -128 to 127 & \code{0} \\
\code{short} & 16 bits & -32,768 to 32,767 & \code{0} \\
\code{int} & 32 bits & -2,147,483,648 to 2,147,483,647 & \code{0} \\
\code{long} & 64 bits & \textasciitilde{}-9.2×10\textasciicircum{}18 to 9.2×10\textasciicircum{}18 & \code{0L} \\
\code{float} & 32 bits & \textasciitilde{}±3.4×10\textasciicircum{}38 (7 sig. digits) & \code{0.\allowbreak{}0f} \\
\code{double} & 64 bits & \textasciitilde{}±1.8×10\textasciicircum{}308 (15 sig. digits) & \code{0.\allowbreak{}0d} \\
\code{char} & 16 bits & 0 to 65,535 (unsigned, UTF-16 code unit) & \code{\textquotesingle{}\textbackslash{}u0000\textquotesingle{}} \\
\code{boolean} & JVM-dependent (not specified) & \code{true} / \code{false} & \code{false} \\
\bottomrule
\end{xltabular}
}

\begin{codebox}
\begin{Verbatim}[commandchars=\\\{\}]
\PY{k+kd}{public}\PY{+w}{ }\PY{k+kd}{class} \PY{n+nc}{PrimitiveBasics}\PY{+w}{ }\PY{p}{\PYZob{}}
\PY{+w}{    }\PY{k+kd}{public}\PY{+w}{ }\PY{k+kd}{static}\PY{+w}{ }\PY{k+kt}{void}\PY{+w}{ }\PY{n+nf}{main}\PY{p}{(}\PY{n}{String}\PY{o}{[}\PY{o}{]}\PY{+w}{ }\PY{n}{args}\PY{p}{)}\PY{+w}{ }\PY{p}{\PYZob{}}
\PY{+w}{        }\PY{k+kt}{byte}\PY{+w}{ }\PY{n}{smallCount}\PY{+w}{ }\PY{o}{=}\PY{+w}{ }\PY{l+m+mi}{100}\PY{p}{;}
\PY{+w}{        }\PY{k+kt}{int}\PY{+w}{ }\PY{n}{userId}\PY{+w}{ }\PY{o}{=}\PY{+w}{ }\PY{l+m+mi}{4\PYZus{}500\PYZus{}231}\PY{p}{;}\PY{+w}{          }\PY{c+c1}{// underscores improve readability}
\PY{+w}{        }\PY{k+kt}{long}\PY{+w}{ }\PY{n}{totalBytes}\PY{+w}{ }\PY{o}{=}\PY{+w}{ }\PY{l+m+mi}{9\PYZus{}000\PYZus{}000\PYZus{}000L}\PY{p}{;}\PY{+w}{ }\PY{c+c1}{// \PYZsq{}L\PYZsq{} suffix required beyond int range}
\PY{+w}{        }\PY{k+kt}{double}\PY{+w}{ }\PY{n}{price}\PY{+w}{ }\PY{o}{=}\PY{+w}{ }\PY{l+m+mf}{19.99}\PY{p}{;}
\PY{+w}{        }\PY{k+kt}{char}\PY{+w}{ }\PY{n}{grade}\PY{+w}{ }\PY{o}{=}\PY{+w}{ }\PY{l+s+sc}{\PYZsq{}A\PYZsq{}}\PY{p}{;}
\PY{+w}{        }\PY{k+kt}{boolean}\PY{+w}{ }\PY{n}{isActive}\PY{+w}{ }\PY{o}{=}\PY{+w}{ }\PY{k+kc}{true}\PY{p}{;}

\PY{+w}{        }\PY{n}{System}\PY{p}{.}\PY{n+na}{out}\PY{p}{.}\PY{n+na}{println}\PY{p}{(}\PY{n}{userId}\PY{p}{)}\PY{p}{;}\PY{+w}{      }\PY{c+c1}{// prints: 4500231}
\PY{+w}{        }\PY{n}{System}\PY{p}{.}\PY{n+na}{out}\PY{p}{.}\PY{n+na}{println}\PY{p}{(}\PY{p}{(}\PY{k+kt}{int}\PY{p}{)}\PY{+w}{ }\PY{n}{grade}\PY{p}{)}\PY{p}{;}\PY{+w}{ }\PY{c+c1}{// prints: 65 (ASCII code for \PYZsq{}A\PYZsq{})}
\PY{+w}{    }\PY{p}{\PYZcb{}}
\PY{p}{\PYZcb{}}
\end{Verbatim}
\end{codebox}

\textbf{Gotcha — integer overflow is silent.} Primitives wrap around instead of throwing an exception:

\begin{codebox}
\begin{Verbatim}[commandchars=\\\{\}]
\PY{k+kt}{int}\PY{+w}{ }\PY{n}{max}\PY{+w}{ }\PY{o}{=}\PY{+w}{ }\PY{n}{Integer}\PY{p}{.}\PY{n+na}{MAX\PYZus{}VALUE}\PY{p}{;}
\PY{n}{System}\PY{p}{.}\PY{n+na}{out}\PY{p}{.}\PY{n+na}{println}\PY{p}{(}\PY{n}{max}\PY{+w}{ }\PY{o}{+}\PY{+w}{ }\PY{l+m+mi}{1}\PY{p}{)}\PY{p}{;}\PY{+w}{ }\PY{c+c1}{// prints: \PYZhy{}2147483648 (wraps around, no error!)}
\end{Verbatim}
\end{codebox}

\textbf{Gotcha — floating-point is not exact.} \code{float} and \code{double} use binary fractions, so they can’t represent every decimal value exactly. This matters a lot for money-related code.

\begin{codebox}
\begin{Verbatim}[commandchars=\\\{\}]
\PY{k+kt}{double}\PY{+w}{ }\PY{n}{a}\PY{+w}{ }\PY{o}{=}\PY{+w}{ }\PY{l+m+mf}{0.1}\PY{p}{;}
\PY{k+kt}{double}\PY{+w}{ }\PY{n}{b}\PY{+w}{ }\PY{o}{=}\PY{+w}{ }\PY{l+m+mf}{0.2}\PY{p}{;}
\PY{n}{System}\PY{p}{.}\PY{n+na}{out}\PY{p}{.}\PY{n+na}{println}\PY{p}{(}\PY{n}{a}\PY{+w}{ }\PY{o}{+}\PY{+w}{ }\PY{n}{b}\PY{p}{)}\PY{p}{;}\PY{+w}{          }\PY{c+c1}{// prints: 0.30000000000000004}
\PY{n}{System}\PY{p}{.}\PY{n+na}{out}\PY{p}{.}\PY{n+na}{println}\PY{p}{(}\PY{n}{a}\PY{+w}{ }\PY{o}{+}\PY{+w}{ }\PY{n}{b}\PY{+w}{ }\PY{o}{=}\PY{o}{=}\PY{+w}{ }\PY{l+m+mf}{0.3}\PY{p}{)}\PY{p}{;}\PY{+w}{   }\PY{c+c1}{// prints: false}

\PY{c+c1}{// Use BigDecimal for money or exact decimal math}
\PY{n}{java}\PY{p}{.}\PY{n+na}{math}\PY{p}{.}\PY{n+na}{BigDecimal}\PY{+w}{ }\PY{n}{x}\PY{+w}{ }\PY{o}{=}\PY{+w}{ }\PY{k}{new}\PY{+w}{ }\PY{n}{java}\PY{p}{.}\PY{n+na}{math}\PY{p}{.}\PY{n+na}{BigDecimal}\PY{p}{(}\PY{l+s}{\PYZdq{}}\PY{l+s}{0.1}\PY{l+s}{\PYZdq{}}\PY{p}{)}\PY{p}{;}
\PY{n}{java}\PY{p}{.}\PY{n+na}{math}\PY{p}{.}\PY{n+na}{BigDecimal}\PY{+w}{ }\PY{n}{y}\PY{+w}{ }\PY{o}{=}\PY{+w}{ }\PY{k}{new}\PY{+w}{ }\PY{n}{java}\PY{p}{.}\PY{n+na}{math}\PY{p}{.}\PY{n+na}{BigDecimal}\PY{p}{(}\PY{l+s}{\PYZdq{}}\PY{l+s}{0.2}\PY{l+s}{\PYZdq{}}\PY{p}{)}\PY{p}{;}
\PY{n}{System}\PY{p}{.}\PY{n+na}{out}\PY{p}{.}\PY{n+na}{println}\PY{p}{(}\PY{n}{x}\PY{p}{.}\PY{n+na}{add}\PY{p}{(}\PY{n}{y}\PY{p}{)}\PY{p}{)}\PY{p}{;}\PY{+w}{       }\PY{c+c1}{// prints: 0.3}
\end{Verbatim}
\end{codebox}

\section[Variables, Scope, and Lifetime]{Variables, Scope, and Lifetime}
\label{h:1:variables-scope-and-lifetime}
A \textbf{variable} is a named storage location. \textbf{Scope} is \emph{where in the code} a variable’s name is visible. \textbf{Lifetime} is \emph{how long} the variable’s storage actually exists in memory. These two often line up, but not always — that’s where bugs hide.

\begin{codebox}
\begin{Verbatim}[commandchars=\\\{\}]
\PY{k+kd}{public}\PY{+w}{ }\PY{k+kd}{class} \PY{n+nc}{ScopeExample}\PY{+w}{ }\PY{p}{\PYZob{}}
\PY{+w}{    }\PY{k+kd}{static}\PY{+w}{ }\PY{k+kt}{int}\PY{+w}{ }\PY{n}{instanceCounter}\PY{+w}{ }\PY{o}{=}\PY{+w}{ }\PY{l+m+mi}{0}\PY{p}{;}\PY{+w}{ }\PY{c+c1}{// class (static) scope: lives for the whole program}

\PY{+w}{    }\PY{k+kd}{public}\PY{+w}{ }\PY{k+kd}{static}\PY{+w}{ }\PY{k+kt}{void}\PY{+w}{ }\PY{n+nf}{main}\PY{p}{(}\PY{n}{String}\PY{o}{[}\PY{o}{]}\PY{+w}{ }\PY{n}{args}\PY{p}{)}\PY{+w}{ }\PY{p}{\PYZob{}}
\PY{+w}{        }\PY{k+kt}{int}\PY{+w}{ }\PY{n}{outer}\PY{+w}{ }\PY{o}{=}\PY{+w}{ }\PY{l+m+mi}{10}\PY{p}{;}\PY{+w}{ }\PY{c+c1}{// method scope: lives for the duration of main()}

\PY{+w}{        }\PY{k}{if}\PY{+w}{ }\PY{p}{(}\PY{n}{outer}\PY{+w}{ }\PY{o}{\PYZgt{}}\PY{+w}{ }\PY{l+m+mi}{5}\PY{p}{)}\PY{+w}{ }\PY{p}{\PYZob{}}
\PY{+w}{            }\PY{k+kt}{int}\PY{+w}{ }\PY{n}{inner}\PY{+w}{ }\PY{o}{=}\PY{+w}{ }\PY{l+m+mi}{20}\PY{p}{;}\PY{+w}{ }\PY{c+c1}{// block scope: only visible inside this if\PYZhy{}block}
\PY{+w}{            }\PY{n}{System}\PY{p}{.}\PY{n+na}{out}\PY{p}{.}\PY{n+na}{println}\PY{p}{(}\PY{n}{outer}\PY{+w}{ }\PY{o}{+}\PY{+w}{ }\PY{n}{inner}\PY{p}{)}\PY{p}{;}\PY{+w}{ }\PY{c+c1}{// prints: 30}
\PY{+w}{        }\PY{p}{\PYZcb{}}

\PY{+w}{        }\PY{c+c1}{// System.out.println(inner); // COMPILE ERROR: inner is out of scope here}

\PY{+w}{        }\PY{k}{for}\PY{+w}{ }\PY{p}{(}\PY{k+kt}{int}\PY{+w}{ }\PY{n}{i}\PY{+w}{ }\PY{o}{=}\PY{+w}{ }\PY{l+m+mi}{0}\PY{p}{;}\PY{+w}{ }\PY{n}{i}\PY{+w}{ }\PY{o}{\PYZlt{}}\PY{+w}{ }\PY{l+m+mi}{3}\PY{p}{;}\PY{+w}{ }\PY{n}{i}\PY{o}{+}\PY{o}{+}\PY{p}{)}\PY{+w}{ }\PY{p}{\PYZob{}}
\PY{+w}{            }\PY{k+kt}{int}\PY{+w}{ }\PY{n}{loopLocal}\PY{+w}{ }\PY{o}{=}\PY{+w}{ }\PY{n}{i}\PY{+w}{ }\PY{o}{*}\PY{+w}{ }\PY{n}{i}\PY{p}{;}\PY{+w}{ }\PY{c+c1}{// re\PYZhy{}created fresh on every iteration}
\PY{+w}{            }\PY{n}{System}\PY{p}{.}\PY{n+na}{out}\PY{p}{.}\PY{n+na}{println}\PY{p}{(}\PY{n}{loopLocal}\PY{p}{)}\PY{p}{;}\PY{+w}{ }\PY{c+c1}{// prints: 0, then 1, then 4}
\PY{+w}{        }\PY{p}{\PYZcb{}}
\PY{+w}{    }\PY{p}{\PYZcb{}}
\PY{p}{\PYZcb{}}
\end{Verbatim}
\end{codebox}

Local variables live on the \textbf{stack} (or are inlined by the JVM) and disappear when their block ends. Instance fields live as long as their object does, tracked by the \textbf{heap} and garbage collector. Static fields live as long as the class is loaded — effectively the whole program run.

\textbf{Gotcha — shadowing.} A variable declared in an inner scope can hide (shadow) one with the same name in an outer scope, which is legal for local variables but easy to misread:

\begin{codebox}
\begin{Verbatim}[commandchars=\\\{\}]
\PY{k+kd}{public}\PY{+w}{ }\PY{k+kd}{class} \PY{n+nc}{Shadowing}\PY{+w}{ }\PY{p}{\PYZob{}}
\PY{+w}{    }\PY{k+kd}{static}\PY{+w}{ }\PY{k+kt}{int}\PY{+w}{ }\PY{n}{value}\PY{+w}{ }\PY{o}{=}\PY{+w}{ }\PY{l+m+mi}{100}\PY{p}{;}

\PY{+w}{    }\PY{k+kd}{public}\PY{+w}{ }\PY{k+kd}{static}\PY{+w}{ }\PY{k+kt}{void}\PY{+w}{ }\PY{n+nf}{main}\PY{p}{(}\PY{n}{String}\PY{o}{[}\PY{o}{]}\PY{+w}{ }\PY{n}{args}\PY{p}{)}\PY{+w}{ }\PY{p}{\PYZob{}}
\PY{+w}{        }\PY{k+kt}{int}\PY{+w}{ }\PY{n}{value}\PY{+w}{ }\PY{o}{=}\PY{+w}{ }\PY{l+m+mi}{5}\PY{p}{;}\PY{+w}{ }\PY{c+c1}{// shadows the static field within main()}
\PY{+w}{        }\PY{n}{System}\PY{p}{.}\PY{n+na}{out}\PY{p}{.}\PY{n+na}{println}\PY{p}{(}\PY{n}{value}\PY{p}{)}\PY{p}{;}\PY{+w}{       }\PY{c+c1}{// prints: 5}
\PY{+w}{        }\PY{n}{System}\PY{p}{.}\PY{n+na}{out}\PY{p}{.}\PY{n+na}{println}\PY{p}{(}\PY{n}{Shadowing}\PY{p}{.}\PY{n+na}{value}\PY{p}{)}\PY{p}{;}\PY{+w}{ }\PY{c+c1}{// prints: 100 (explicit access to the field)}
\PY{+w}{    }\PY{p}{\PYZcb{}}
\PY{p}{\PYZcb{}}
\end{Verbatim}
\end{codebox}

\textbf{Gotcha — variables must be definitely assigned before use.} Java won’t let you read a local variable that might not have been initialized on every path:

\begin{codebox}
\begin{Verbatim}[commandchars=\\\{\}]
\PY{k+kt}{int}\PY{+w}{ }\PY{n}{result}\PY{p}{;}
\PY{k}{if}\PY{+w}{ }\PY{p}{(}\PY{n}{Math}\PY{p}{.}\PY{n+na}{random}\PY{p}{(}\PY{p}{)}\PY{+w}{ }\PY{o}{\PYZgt{}}\PY{+w}{ }\PY{l+m+mf}{0.5}\PY{p}{)}\PY{+w}{ }\PY{p}{\PYZob{}}
\PY{+w}{    }\PY{n}{result}\PY{+w}{ }\PY{o}{=}\PY{+w}{ }\PY{l+m+mi}{1}\PY{p}{;}
\PY{p}{\PYZcb{}}
\PY{c+c1}{// System.out.println(result); // COMPILE ERROR: result might not have been initialized}
\end{Verbatim}
\end{codebox}

\section[Operators and Expressions]{Operators and Expressions}
\label{h:1:operators-and-expressions}
An \textbf{operator} combines values into an \textbf{expression} that produces a new value. Java groups them into arithmetic, relational, logical, bitwise, assignment, and ternary categories. Precedence and evaluation order matter for review because subtle mistakes often hide in dense one-liners.

\begin{codebox}
\begin{Verbatim}[commandchars=\\\{\}]
\PY{k+kd}{public}\PY{+w}{ }\PY{k+kd}{class} \PY{n+nc}{Operators}\PY{+w}{ }\PY{p}{\PYZob{}}
\PY{+w}{    }\PY{k+kd}{public}\PY{+w}{ }\PY{k+kd}{static}\PY{+w}{ }\PY{k+kt}{void}\PY{+w}{ }\PY{n+nf}{main}\PY{p}{(}\PY{n}{String}\PY{o}{[}\PY{o}{]}\PY{+w}{ }\PY{n}{args}\PY{p}{)}\PY{+w}{ }\PY{p}{\PYZob{}}
\PY{+w}{        }\PY{k+kt}{int}\PY{+w}{ }\PY{n}{a}\PY{+w}{ }\PY{o}{=}\PY{+w}{ }\PY{l+m+mi}{7}\PY{p}{,}\PY{+w}{ }\PY{n}{b}\PY{+w}{ }\PY{o}{=}\PY{+w}{ }\PY{l+m+mi}{2}\PY{p}{;}

\PY{+w}{        }\PY{n}{System}\PY{p}{.}\PY{n+na}{out}\PY{p}{.}\PY{n+na}{println}\PY{p}{(}\PY{n}{a}\PY{+w}{ }\PY{o}{/}\PY{+w}{ }\PY{n}{b}\PY{p}{)}\PY{p}{;}\PY{+w}{   }\PY{c+c1}{// prints: 3  (integer division truncates)}
\PY{+w}{        }\PY{n}{System}\PY{p}{.}\PY{n+na}{out}\PY{p}{.}\PY{n+na}{println}\PY{p}{(}\PY{n}{a}\PY{+w}{ }\PY{o}{\PYZpc{}}\PY{+w}{ }\PY{n}{b}\PY{p}{)}\PY{p}{;}\PY{+w}{   }\PY{c+c1}{// prints: 1  (remainder)}
\PY{+w}{        }\PY{n}{System}\PY{p}{.}\PY{n+na}{out}\PY{p}{.}\PY{n+na}{println}\PY{p}{(}\PY{p}{(}\PY{k+kt}{double}\PY{p}{)}\PY{+w}{ }\PY{n}{a}\PY{+w}{ }\PY{o}{/}\PY{+w}{ }\PY{n}{b}\PY{p}{)}\PY{p}{;}\PY{+w}{ }\PY{c+c1}{// prints: 3.5 (cast forces floating\PYZhy{}point division)}

\PY{+w}{        }\PY{k+kt}{boolean}\PY{+w}{ }\PY{n}{eligible}\PY{+w}{ }\PY{o}{=}\PY{+w}{ }\PY{p}{(}\PY{n}{a}\PY{+w}{ }\PY{o}{\PYZgt{}}\PY{+w}{ }\PY{l+m+mi}{5}\PY{p}{)}\PY{+w}{ }\PY{o}{\PYZam{}}\PY{o}{\PYZam{}}\PY{+w}{ }\PY{p}{(}\PY{n}{b}\PY{+w}{ }\PY{o}{\PYZlt{}}\PY{+w}{ }\PY{l+m+mi}{10}\PY{p}{)}\PY{p}{;}\PY{+w}{ }\PY{c+c1}{// logical AND}
\PY{+w}{        }\PY{n}{System}\PY{p}{.}\PY{n+na}{out}\PY{p}{.}\PY{n+na}{println}\PY{p}{(}\PY{n}{eligible}\PY{p}{)}\PY{p}{;}\PY{+w}{ }\PY{c+c1}{// prints: true}

\PY{+w}{        }\PY{k+kt}{int}\PY{+w}{ }\PY{n}{shifted}\PY{+w}{ }\PY{o}{=}\PY{+w}{ }\PY{l+m+mi}{1}\PY{+w}{ }\PY{o}{\PYZlt{}}\PY{o}{\PYZlt{}}\PY{+w}{ }\PY{l+m+mi}{4}\PY{p}{;}\PY{+w}{ }\PY{c+c1}{// bitwise left shift: 1 * 2\PYZca{}4}
\PY{+w}{        }\PY{n}{System}\PY{p}{.}\PY{n+na}{out}\PY{p}{.}\PY{n+na}{println}\PY{p}{(}\PY{n}{shifted}\PY{p}{)}\PY{p}{;}\PY{+w}{ }\PY{c+c1}{// prints: 16}

\PY{+w}{        }\PY{n}{String}\PY{+w}{ }\PY{n}{label}\PY{+w}{ }\PY{o}{=}\PY{+w}{ }\PY{p}{(}\PY{n}{a}\PY{+w}{ }\PY{o}{\PYZgt{}}\PY{+w}{ }\PY{n}{b}\PY{p}{)}\PY{+w}{ }\PY{o}{?}\PY{+w}{ }\PY{l+s}{\PYZdq{}}\PY{l+s}{a wins}\PY{l+s}{\PYZdq{}}\PY{+w}{ }\PY{p}{:}\PY{+w}{ }\PY{l+s}{\PYZdq{}}\PY{l+s}{b wins}\PY{l+s}{\PYZdq{}}\PY{p}{;}\PY{+w}{ }\PY{c+c1}{// ternary operator}
\PY{+w}{        }\PY{n}{System}\PY{p}{.}\PY{n+na}{out}\PY{p}{.}\PY{n+na}{println}\PY{p}{(}\PY{n}{label}\PY{p}{)}\PY{p}{;}\PY{+w}{ }\PY{c+c1}{// prints: a wins}
\PY{+w}{    }\PY{p}{\PYZcb{}}
\PY{p}{\PYZcb{}}
\end{Verbatim}
\end{codebox}

\textbf{Short-circuit evaluation.} \code{\&\&} and \code{\textbar{}\textbar{}} stop evaluating as soon as the result is known. \code{\&} and \code{\textbar{}} (single character, also usable on booleans) always evaluate both sides. This distinction matters when the right-hand side has side effects or could throw:

\begin{codebox}
\begin{Verbatim}[commandchars=\\\{\}]
\PY{n}{String}\PY{+w}{ }\PY{n}{input}\PY{+w}{ }\PY{o}{=}\PY{+w}{ }\PY{k+kc}{null}\PY{p}{;}
\PY{k}{if}\PY{+w}{ }\PY{p}{(}\PY{n}{input}\PY{+w}{ }\PY{o}{!}\PY{o}{=}\PY{+w}{ }\PY{k+kc}{null}\PY{+w}{ }\PY{o}{\PYZam{}}\PY{o}{\PYZam{}}\PY{+w}{ }\PY{n}{input}\PY{p}{.}\PY{n+na}{length}\PY{p}{(}\PY{p}{)}\PY{+w}{ }\PY{o}{\PYZgt{}}\PY{+w}{ }\PY{l+m+mi}{0}\PY{p}{)}\PY{+w}{ }\PY{p}{\PYZob{}}\PY{+w}{ }\PY{c+c1}{// safe: length() never called on null}
\PY{+w}{    }\PY{n}{System}\PY{p}{.}\PY{n+na}{out}\PY{p}{.}\PY{n+na}{println}\PY{p}{(}\PY{l+s}{\PYZdq{}}\PY{l+s}{non\PYZhy{}empty}\PY{l+s}{\PYZdq{}}\PY{p}{)}\PY{p}{;}
\PY{p}{\PYZcb{}}

\PY{c+c1}{// if (input != null \PYZam{} input.length() \PYZgt{} 0) // would throw NullPointerException!}
\end{Verbatim}
\end{codebox}

\textbf{Gotcha — integer division truncates toward zero.}

\begin{codebox}
\begin{Verbatim}[commandchars=\\\{\}]
\PY{n}{System}\PY{p}{.}\PY{n+na}{out}\PY{p}{.}\PY{n+na}{println}\PY{p}{(}\PY{o}{\PYZhy{}}\PY{l+m+mi}{7}\PY{+w}{ }\PY{o}{/}\PY{+w}{ }\PY{l+m+mi}{2}\PY{p}{)}\PY{p}{;}\PY{+w}{  }\PY{c+c1}{// prints: \PYZhy{}3 (not \PYZhy{}4; truncates toward zero)}
\PY{n}{System}\PY{p}{.}\PY{n+na}{out}\PY{p}{.}\PY{n+na}{println}\PY{p}{(}\PY{o}{\PYZhy{}}\PY{l+m+mi}{7}\PY{+w}{ }\PY{o}{\PYZpc{}}\PY{+w}{ }\PY{l+m+mi}{2}\PY{p}{)}\PY{p}{;}\PY{+w}{  }\PY{c+c1}{// prints: \PYZhy{}1 (sign follows the dividend)}
\end{Verbatim}
\end{codebox}

\textbf{Gotcha — mixing \code{++}/\code{--} inside larger expressions is legal but hard to read.}

\begin{codebox}
\begin{Verbatim}[commandchars=\\\{\}]
\PY{k+kt}{int}\PY{+w}{ }\PY{n}{i}\PY{+w}{ }\PY{o}{=}\PY{+w}{ }\PY{l+m+mi}{5}\PY{p}{;}
\PY{k+kt}{int}\PY{+w}{ }\PY{n}{j}\PY{+w}{ }\PY{o}{=}\PY{+w}{ }\PY{n}{i}\PY{o}{+}\PY{o}{+}\PY{+w}{ }\PY{o}{+}\PY{+w}{ }\PY{o}{+}\PY{o}{+}\PY{n}{i}\PY{p}{;}\PY{+w}{ }\PY{c+c1}{// i++ uses 5, then i becomes 6; ++i makes it 7 and uses 7}
\PY{n}{System}\PY{p}{.}\PY{n+na}{out}\PY{p}{.}\PY{n+na}{println}\PY{p}{(}\PY{n}{j}\PY{p}{)}\PY{p}{;}\PY{+w}{ }\PY{c+c1}{// prints: 12}
\PY{n}{System}\PY{p}{.}\PY{n+na}{out}\PY{p}{.}\PY{n+na}{println}\PY{p}{(}\PY{n}{i}\PY{p}{)}\PY{p}{;}\PY{+w}{ }\PY{c+c1}{// prints: 7}
\end{Verbatim}
\end{codebox}

Avoid writing code like this in production — it compiles fine but reviewers should flag it for clarity.

\section[Control Flow Statements]{Control Flow Statements}
\label{h:1:control-flow-statements}
\textbf{Control flow} statements decide which code runs, and how many times. Java’s core set is \code{if}/\code{else}, \code{switch}, \code{for}, \code{while}, \code{do-\allowbreak{}while}, plus \code{break}, \code{continue}, and (since Java 14) the \textbf{switch expression}.

\begin{codebox}
\begin{Verbatim}[commandchars=\\\{\}]
\PY{k+kd}{public}\PY{+w}{ }\PY{k+kd}{class} \PY{n+nc}{ControlFlow}\PY{+w}{ }\PY{p}{\PYZob{}}
\PY{+w}{    }\PY{k+kd}{public}\PY{+w}{ }\PY{k+kd}{static}\PY{+w}{ }\PY{k+kt}{void}\PY{+w}{ }\PY{n+nf}{main}\PY{p}{(}\PY{n}{String}\PY{o}{[}\PY{o}{]}\PY{+w}{ }\PY{n}{args}\PY{p}{)}\PY{+w}{ }\PY{p}{\PYZob{}}
\PY{+w}{        }\PY{k+kt}{int}\PY{+w}{ }\PY{n}{score}\PY{+w}{ }\PY{o}{=}\PY{+w}{ }\PY{l+m+mi}{82}\PY{p}{;}

\PY{+w}{        }\PY{c+c1}{// Classic if/else chain}
\PY{+w}{        }\PY{n}{String}\PY{+w}{ }\PY{n}{grade}\PY{p}{;}
\PY{+w}{        }\PY{k}{if}\PY{+w}{ }\PY{p}{(}\PY{n}{score}\PY{+w}{ }\PY{o}{\PYZgt{}}\PY{o}{=}\PY{+w}{ }\PY{l+m+mi}{90}\PY{p}{)}\PY{+w}{ }\PY{p}{\PYZob{}}
\PY{+w}{            }\PY{n}{grade}\PY{+w}{ }\PY{o}{=}\PY{+w}{ }\PY{l+s}{\PYZdq{}}\PY{l+s}{A}\PY{l+s}{\PYZdq{}}\PY{p}{;}
\PY{+w}{        }\PY{p}{\PYZcb{}}\PY{+w}{ }\PY{k}{else}\PY{+w}{ }\PY{k}{if}\PY{+w}{ }\PY{p}{(}\PY{n}{score}\PY{+w}{ }\PY{o}{\PYZgt{}}\PY{o}{=}\PY{+w}{ }\PY{l+m+mi}{80}\PY{p}{)}\PY{+w}{ }\PY{p}{\PYZob{}}
\PY{+w}{            }\PY{n}{grade}\PY{+w}{ }\PY{o}{=}\PY{+w}{ }\PY{l+s}{\PYZdq{}}\PY{l+s}{B}\PY{l+s}{\PYZdq{}}\PY{p}{;}
\PY{+w}{        }\PY{p}{\PYZcb{}}\PY{+w}{ }\PY{k}{else}\PY{+w}{ }\PY{p}{\PYZob{}}
\PY{+w}{            }\PY{n}{grade}\PY{+w}{ }\PY{o}{=}\PY{+w}{ }\PY{l+s}{\PYZdq{}}\PY{l+s}{C}\PY{l+s}{\PYZdq{}}\PY{p}{;}
\PY{+w}{        }\PY{p}{\PYZcb{}}
\PY{+w}{        }\PY{n}{System}\PY{p}{.}\PY{n+na}{out}\PY{p}{.}\PY{n+na}{println}\PY{p}{(}\PY{n}{grade}\PY{p}{)}\PY{p}{;}\PY{+w}{ }\PY{c+c1}{// prints: B}

\PY{+w}{        }\PY{c+c1}{// Traditional switch statement (fall\PYZhy{}through by default)}
\PY{+w}{        }\PY{k}{switch}\PY{+w}{ }\PY{p}{(}\PY{n}{grade}\PY{p}{)}\PY{+w}{ }\PY{p}{\PYZob{}}
\PY{+w}{            }\PY{k}{case}\PY{+w}{ }\PY{l+s}{\PYZdq{}}\PY{l+s}{A}\PY{l+s}{\PYZdq{}}\PY{p}{:}
\PY{+w}{                }\PY{n}{System}\PY{p}{.}\PY{n+na}{out}\PY{p}{.}\PY{n+na}{println}\PY{p}{(}\PY{l+s}{\PYZdq{}}\PY{l+s}{Excellent}\PY{l+s}{\PYZdq{}}\PY{p}{)}\PY{p}{;}
\PY{+w}{                }\PY{k}{break}\PY{p}{;}\PY{+w}{ }\PY{c+c1}{// without break, execution falls into the next case!}
\PY{+w}{            }\PY{k}{case}\PY{+w}{ }\PY{l+s}{\PYZdq{}}\PY{l+s}{B}\PY{l+s}{\PYZdq{}}\PY{p}{:}
\PY{+w}{                }\PY{n}{System}\PY{p}{.}\PY{n+na}{out}\PY{p}{.}\PY{n+na}{println}\PY{p}{(}\PY{l+s}{\PYZdq{}}\PY{l+s}{Good}\PY{l+s}{\PYZdq{}}\PY{p}{)}\PY{p}{;}
\PY{+w}{                }\PY{k}{break}\PY{p}{;}
\PY{+w}{            }\PY{k}{default}\PY{p}{:}
\PY{+w}{                }\PY{n}{System}\PY{p}{.}\PY{n+na}{out}\PY{p}{.}\PY{n+na}{println}\PY{p}{(}\PY{l+s}{\PYZdq{}}\PY{l+s}{Keep trying}\PY{l+s}{\PYZdq{}}\PY{p}{)}\PY{p}{;}
\PY{+w}{        }\PY{p}{\PYZcb{}}

\PY{+w}{        }\PY{c+c1}{// Modern switch expression (Java 14+): no fall\PYZhy{}through, returns a value}
\PY{+w}{        }\PY{n}{String}\PY{+w}{ }\PY{n}{message}\PY{+w}{ }\PY{o}{=}\PY{+w}{ }\PY{k}{switch}\PY{+w}{ }\PY{p}{(}\PY{n}{grade}\PY{p}{)}\PY{+w}{ }\PY{p}{\PYZob{}}
\PY{+w}{            }\PY{k}{case}\PY{+w}{ }\PY{l+s}{\PYZdq{}}\PY{l+s}{A}\PY{l+s}{\PYZdq{}}\PY{+w}{ }\PY{o}{\PYZhy{}}\PY{o}{\PYZgt{}}\PY{+w}{ }\PY{l+s}{\PYZdq{}}\PY{l+s}{Excellent}\PY{l+s}{\PYZdq{}}\PY{p}{;}
\PY{+w}{            }\PY{k}{case}\PY{+w}{ }\PY{l+s}{\PYZdq{}}\PY{l+s}{B}\PY{l+s}{\PYZdq{}}\PY{+w}{ }\PY{o}{\PYZhy{}}\PY{o}{\PYZgt{}}\PY{+w}{ }\PY{l+s}{\PYZdq{}}\PY{l+s}{Good}\PY{l+s}{\PYZdq{}}\PY{p}{;}
\PY{+w}{            }\PY{k}{default}\PY{+w}{ }\PY{o}{\PYZhy{}}\PY{o}{\PYZgt{}}\PY{+w}{ }\PY{l+s}{\PYZdq{}}\PY{l+s}{Keep trying}\PY{l+s}{\PYZdq{}}\PY{p}{;}
\PY{+w}{        }\PY{p}{\PYZcb{}}\PY{p}{;}
\PY{+w}{        }\PY{n}{System}\PY{p}{.}\PY{n+na}{out}\PY{p}{.}\PY{n+na}{println}\PY{p}{(}\PY{n}{message}\PY{p}{)}\PY{p}{;}\PY{+w}{ }\PY{c+c1}{// prints: Good}

\PY{+w}{        }\PY{c+c1}{// Loops}
\PY{+w}{        }\PY{k}{for}\PY{+w}{ }\PY{p}{(}\PY{k+kt}{int}\PY{+w}{ }\PY{n}{i}\PY{+w}{ }\PY{o}{=}\PY{+w}{ }\PY{l+m+mi}{0}\PY{p}{;}\PY{+w}{ }\PY{n}{i}\PY{+w}{ }\PY{o}{\PYZlt{}}\PY{+w}{ }\PY{l+m+mi}{3}\PY{p}{;}\PY{+w}{ }\PY{n}{i}\PY{o}{+}\PY{o}{+}\PY{p}{)}\PY{+w}{ }\PY{p}{\PYZob{}}
\PY{+w}{            }\PY{n}{System}\PY{p}{.}\PY{n+na}{out}\PY{p}{.}\PY{n+na}{print}\PY{p}{(}\PY{n}{i}\PY{+w}{ }\PY{o}{+}\PY{+w}{ }\PY{l+s}{\PYZdq{}}\PY{l+s}{ }\PY{l+s}{\PYZdq{}}\PY{p}{)}\PY{p}{;}\PY{+w}{ }\PY{c+c1}{// prints: 0 1 2}
\PY{+w}{        }\PY{p}{\PYZcb{}}
\PY{+w}{        }\PY{n}{System}\PY{p}{.}\PY{n+na}{out}\PY{p}{.}\PY{n+na}{println}\PY{p}{(}\PY{p}{)}\PY{p}{;}

\PY{+w}{        }\PY{k+kt}{int}\PY{+w}{ }\PY{n}{n}\PY{+w}{ }\PY{o}{=}\PY{+w}{ }\PY{l+m+mi}{3}\PY{p}{;}
\PY{+w}{        }\PY{k}{while}\PY{+w}{ }\PY{p}{(}\PY{n}{n}\PY{+w}{ }\PY{o}{\PYZgt{}}\PY{+w}{ }\PY{l+m+mi}{0}\PY{p}{)}\PY{+w}{ }\PY{p}{\PYZob{}}
\PY{+w}{            }\PY{n}{System}\PY{p}{.}\PY{n+na}{out}\PY{p}{.}\PY{n+na}{print}\PY{p}{(}\PY{n}{n}\PY{+w}{ }\PY{o}{+}\PY{+w}{ }\PY{l+s}{\PYZdq{}}\PY{l+s}{ }\PY{l+s}{\PYZdq{}}\PY{p}{)}\PY{p}{;}\PY{+w}{ }\PY{c+c1}{// prints: 3 2 1}
\PY{+w}{            }\PY{n}{n}\PY{o}{\PYZhy{}}\PY{o}{\PYZhy{}}\PY{p}{;}
\PY{+w}{        }\PY{p}{\PYZcb{}}
\PY{+w}{        }\PY{n}{System}\PY{p}{.}\PY{n+na}{out}\PY{p}{.}\PY{n+na}{println}\PY{p}{(}\PY{p}{)}\PY{p}{;}

\PY{+w}{        }\PY{c+c1}{// do\PYZhy{}while always runs the body at least once}
\PY{+w}{        }\PY{k+kt}{int}\PY{+w}{ }\PY{n}{attempts}\PY{+w}{ }\PY{o}{=}\PY{+w}{ }\PY{l+m+mi}{0}\PY{p}{;}
\PY{+w}{        }\PY{k}{do}\PY{+w}{ }\PY{p}{\PYZob{}}
\PY{+w}{            }\PY{n}{attempts}\PY{o}{+}\PY{o}{+}\PY{p}{;}
\PY{+w}{        }\PY{p}{\PYZcb{}}\PY{+w}{ }\PY{k}{while}\PY{+w}{ }\PY{p}{(}\PY{n}{attempts}\PY{+w}{ }\PY{o}{\PYZlt{}}\PY{+w}{ }\PY{l+m+mi}{0}\PY{p}{)}\PY{p}{;}\PY{+w}{ }\PY{c+c1}{// condition false immediately, but body ran once}
\PY{+w}{        }\PY{n}{System}\PY{p}{.}\PY{n+na}{out}\PY{p}{.}\PY{n+na}{println}\PY{p}{(}\PY{n}{attempts}\PY{p}{)}\PY{p}{;}\PY{+w}{ }\PY{c+c1}{// prints: 1}
\PY{+w}{    }\PY{p}{\PYZcb{}}
\PY{p}{\PYZcb{}}
\end{Verbatim}
\end{codebox}

\textbf{Gotcha — switch fall-through.} Forgetting \code{break} in a classic \code{switch} statement lets execution continue into the next case:

\begin{codebox}
\begin{Verbatim}[commandchars=\\\{\}]
\PY{k+kt}{int}\PY{+w}{ }\PY{n}{day}\PY{+w}{ }\PY{o}{=}\PY{+w}{ }\PY{l+m+mi}{6}\PY{p}{;}
\PY{k}{switch}\PY{+w}{ }\PY{p}{(}\PY{n}{day}\PY{p}{)}\PY{+w}{ }\PY{p}{\PYZob{}}
\PY{+w}{    }\PY{k}{case}\PY{+w}{ }\PY{l+m+mi}{6}\PY{p}{:}
\PY{+w}{    }\PY{k}{case}\PY{+w}{ }\PY{l+m+mi}{7}\PY{p}{:}
\PY{+w}{        }\PY{n}{System}\PY{p}{.}\PY{n+na}{out}\PY{p}{.}\PY{n+na}{println}\PY{p}{(}\PY{l+s}{\PYZdq{}}\PY{l+s}{Weekend}\PY{l+s}{\PYZdq{}}\PY{p}{)}\PY{p}{;}
\PY{+w}{        }\PY{k}{break}\PY{p}{;}
\PY{+w}{    }\PY{k}{default}\PY{p}{:}
\PY{+w}{        }\PY{n}{System}\PY{p}{.}\PY{n+na}{out}\PY{p}{.}\PY{n+na}{println}\PY{p}{(}\PY{l+s}{\PYZdq{}}\PY{l+s}{Weekday}\PY{l+s}{\PYZdq{}}\PY{p}{)}\PY{p}{;}
\PY{p}{\PYZcb{}}
\PY{c+c1}{// prints: Weekend (cases 6 and 7 intentionally share the same code — this is fine)}
\end{Verbatim}
\end{codebox}

Prefer the arrow-style \code{switch} expression in new code: it removes fall-through entirely and forces every path to return a value (the compiler checks exhaustiveness for enums and sealed types).

\textbf{Gotcha — labeled break/continue.} Rarely used, but useful for breaking out of nested loops:

\begin{codebox}
\begin{Verbatim}[commandchars=\\\{\}]
\PY{n+nl}{outer}\PY{p}{:}
\PY{k}{for}\PY{+w}{ }\PY{p}{(}\PY{k+kt}{int}\PY{+w}{ }\PY{n}{i}\PY{+w}{ }\PY{o}{=}\PY{+w}{ }\PY{l+m+mi}{0}\PY{p}{;}\PY{+w}{ }\PY{n}{i}\PY{+w}{ }\PY{o}{\PYZlt{}}\PY{+w}{ }\PY{l+m+mi}{3}\PY{p}{;}\PY{+w}{ }\PY{n}{i}\PY{o}{+}\PY{o}{+}\PY{p}{)}\PY{+w}{ }\PY{p}{\PYZob{}}
\PY{+w}{    }\PY{k}{for}\PY{+w}{ }\PY{p}{(}\PY{k+kt}{int}\PY{+w}{ }\PY{n}{j}\PY{+w}{ }\PY{o}{=}\PY{+w}{ }\PY{l+m+mi}{0}\PY{p}{;}\PY{+w}{ }\PY{n}{j}\PY{+w}{ }\PY{o}{\PYZlt{}}\PY{+w}{ }\PY{l+m+mi}{3}\PY{p}{;}\PY{+w}{ }\PY{n}{j}\PY{o}{+}\PY{o}{+}\PY{p}{)}\PY{+w}{ }\PY{p}{\PYZob{}}
\PY{+w}{        }\PY{k}{if}\PY{+w}{ }\PY{p}{(}\PY{n}{j}\PY{+w}{ }\PY{o}{=}\PY{o}{=}\PY{+w}{ }\PY{l+m+mi}{1}\PY{p}{)}\PY{+w}{ }\PY{p}{\PYZob{}}
\PY{+w}{            }\PY{k}{continue}\PY{+w}{ }\PY{n}{outer}\PY{p}{;}\PY{+w}{ }\PY{c+c1}{// skips to the next i, not just the next j}
\PY{+w}{        }\PY{p}{\PYZcb{}}
\PY{+w}{        }\PY{n}{System}\PY{p}{.}\PY{n+na}{out}\PY{p}{.}\PY{n+na}{println}\PY{p}{(}\PY{n}{i}\PY{+w}{ }\PY{o}{+}\PY{+w}{ }\PY{l+s}{\PYZdq{}}\PY{l+s}{,}\PY{l+s}{\PYZdq{}}\PY{+w}{ }\PY{o}{+}\PY{+w}{ }\PY{n}{j}\PY{p}{)}\PY{p}{;}
\PY{+w}{    }\PY{p}{\PYZcb{}}
\PY{p}{\PYZcb{}}
\PY{c+c1}{// prints: 0,0  1,0  2,0}
\end{Verbatim}
\end{codebox}

\section[Methods and Parameter Passing]{Methods and Parameter Passing}
\label{h:1:methods-and-parameter-passing}
A \textbf{method} is a named, reusable block of code. Java is famous for one rule that trips up almost every interview candidate: \textbf{Java is always pass-by-value.} What gets copied differs for primitives versus objects, and that difference explains most of the confusion.

\begin{itemize}
\item 
For \textbf{primitives}, the value itself is copied. Changes inside the method never affect the caller’s variable.

\item 
For \textbf{objects}, the \emph{reference} (the “address” pointing to the object) is copied — not the object itself. The method can use that reference to mutate the object’s internal state, but it cannot make the caller’s variable point to a different object.

\end{itemize}

\begin{codebox}
\begin{Verbatim}[commandchars=\\\{\}]
\PY{k+kd}{public}\PY{+w}{ }\PY{k+kd}{class} \PY{n+nc}{ParameterPassing}\PY{+w}{ }\PY{p}{\PYZob{}}

\PY{+w}{    }\PY{k+kd}{static}\PY{+w}{ }\PY{k+kt}{void}\PY{+w}{ }\PY{n+nf}{incrementPrimitive}\PY{p}{(}\PY{k+kt}{int}\PY{+w}{ }\PY{n}{value}\PY{p}{)}\PY{+w}{ }\PY{p}{\PYZob{}}
\PY{+w}{        }\PY{n}{value}\PY{+w}{ }\PY{o}{=}\PY{+w}{ }\PY{n}{value}\PY{+w}{ }\PY{o}{+}\PY{+w}{ }\PY{l+m+mi}{1}\PY{p}{;}\PY{+w}{ }\PY{c+c1}{// only modifies the local copy}
\PY{+w}{    }\PY{p}{\PYZcb{}}

\PY{+w}{    }\PY{k+kd}{static}\PY{+w}{ }\PY{k+kt}{void}\PY{+w}{ }\PY{n+nf}{appendToList}\PY{p}{(}\PY{n}{java}\PY{p}{.}\PY{n+na}{util}\PY{p}{.}\PY{n+na}{List}\PY{o}{\PYZlt{}}\PY{n}{String}\PY{o}{\PYZgt{}}\PY{+w}{ }\PY{n}{list}\PY{p}{)}\PY{+w}{ }\PY{p}{\PYZob{}}
\PY{+w}{        }\PY{n}{list}\PY{p}{.}\PY{n+na}{add}\PY{p}{(}\PY{l+s}{\PYZdq{}}\PY{l+s}{added inside method}\PY{l+s}{\PYZdq{}}\PY{p}{)}\PY{p}{;}\PY{+w}{ }\PY{c+c1}{// mutates the SAME object the caller sees}
\PY{+w}{    }\PY{p}{\PYZcb{}}

\PY{+w}{    }\PY{k+kd}{static}\PY{+w}{ }\PY{k+kt}{void}\PY{+w}{ }\PY{n+nf}{reassignReference}\PY{p}{(}\PY{n}{java}\PY{p}{.}\PY{n+na}{util}\PY{p}{.}\PY{n+na}{List}\PY{o}{\PYZlt{}}\PY{n}{String}\PY{o}{\PYZgt{}}\PY{+w}{ }\PY{n}{list}\PY{p}{)}\PY{+w}{ }\PY{p}{\PYZob{}}
\PY{+w}{        }\PY{n}{list}\PY{+w}{ }\PY{o}{=}\PY{+w}{ }\PY{k}{new}\PY{+w}{ }\PY{n}{java}\PY{p}{.}\PY{n+na}{util}\PY{p}{.}\PY{n+na}{ArrayList}\PY{o}{\PYZlt{}}\PY{o}{\PYZgt{}}\PY{p}{(}\PY{p}{)}\PY{p}{;}\PY{+w}{ }\PY{c+c1}{// rebinds the LOCAL copy of the reference only}
\PY{+w}{        }\PY{n}{list}\PY{p}{.}\PY{n+na}{add}\PY{p}{(}\PY{l+s}{\PYZdq{}}\PY{l+s}{this is invisible to the caller}\PY{l+s}{\PYZdq{}}\PY{p}{)}\PY{p}{;}
\PY{+w}{    }\PY{p}{\PYZcb{}}

\PY{+w}{    }\PY{k+kd}{public}\PY{+w}{ }\PY{k+kd}{static}\PY{+w}{ }\PY{k+kt}{void}\PY{+w}{ }\PY{n+nf}{main}\PY{p}{(}\PY{n}{String}\PY{o}{[}\PY{o}{]}\PY{+w}{ }\PY{n}{args}\PY{p}{)}\PY{+w}{ }\PY{p}{\PYZob{}}
\PY{+w}{        }\PY{k+kt}{int}\PY{+w}{ }\PY{n}{number}\PY{+w}{ }\PY{o}{=}\PY{+w}{ }\PY{l+m+mi}{5}\PY{p}{;}
\PY{+w}{        }\PY{n}{incrementPrimitive}\PY{p}{(}\PY{n}{number}\PY{p}{)}\PY{p}{;}
\PY{+w}{        }\PY{n}{System}\PY{p}{.}\PY{n+na}{out}\PY{p}{.}\PY{n+na}{println}\PY{p}{(}\PY{n}{number}\PY{p}{)}\PY{p}{;}\PY{+w}{ }\PY{c+c1}{// prints: 5 (unchanged)}

\PY{+w}{        }\PY{n}{java}\PY{p}{.}\PY{n+na}{util}\PY{p}{.}\PY{n+na}{List}\PY{o}{\PYZlt{}}\PY{n}{String}\PY{o}{\PYZgt{}}\PY{+w}{ }\PY{n}{names}\PY{+w}{ }\PY{o}{=}\PY{+w}{ }\PY{k}{new}\PY{+w}{ }\PY{n}{java}\PY{p}{.}\PY{n+na}{util}\PY{p}{.}\PY{n+na}{ArrayList}\PY{o}{\PYZlt{}}\PY{o}{\PYZgt{}}\PY{p}{(}\PY{p}{)}\PY{p}{;}
\PY{+w}{        }\PY{n}{names}\PY{p}{.}\PY{n+na}{add}\PY{p}{(}\PY{l+s}{\PYZdq{}}\PY{l+s}{Alice}\PY{l+s}{\PYZdq{}}\PY{p}{)}\PY{p}{;}
\PY{+w}{        }\PY{n}{appendToList}\PY{p}{(}\PY{n}{names}\PY{p}{)}\PY{p}{;}
\PY{+w}{        }\PY{n}{System}\PY{p}{.}\PY{n+na}{out}\PY{p}{.}\PY{n+na}{println}\PY{p}{(}\PY{n}{names}\PY{p}{)}\PY{p}{;}\PY{+w}{ }\PY{c+c1}{// prints: [Alice, added inside method]}

\PY{+w}{        }\PY{n}{reassignReference}\PY{p}{(}\PY{n}{names}\PY{p}{)}\PY{p}{;}
\PY{+w}{        }\PY{n}{System}\PY{p}{.}\PY{n+na}{out}\PY{p}{.}\PY{n+na}{println}\PY{p}{(}\PY{n}{names}\PY{p}{)}\PY{p}{;}\PY{+w}{ }\PY{c+c1}{// prints: [Alice, added inside method] (unchanged!)}
\PY{+w}{    }\PY{p}{\PYZcb{}}
\PY{p}{\PYZcb{}}
\end{Verbatim}
\end{codebox}

Think of a reference like a piece of paper with a house address written on it. Passing it to a method hands over a \emph{photocopy} of that paper. The method can walk to the house and rearrange the furniture (mutate the object) — everyone with a copy of the address sees the new furniture. But if the method writes a \emph{different} address on its photocopy, that only affects its own slip of paper; your original paper still points to the old house.

\textbf{Gotcha — “pass by reference” is a myth in Java.} There is no way to make a method change which object a caller’s variable points to, without returning a new value:

\begin{codebox}
\begin{Verbatim}[commandchars=\\\{\}]
\PY{k+kd}{static}\PY{+w}{ }\PY{k+kt}{void}\PY{+w}{ }\PY{n+nf}{swap}\PY{p}{(}\PY{n}{String}\PY{+w}{ }\PY{n}{a}\PY{p}{,}\PY{+w}{ }\PY{n}{String}\PY{+w}{ }\PY{n}{b}\PY{p}{)}\PY{+w}{ }\PY{p}{\PYZob{}}
\PY{+w}{    }\PY{n}{String}\PY{+w}{ }\PY{n}{temp}\PY{+w}{ }\PY{o}{=}\PY{+w}{ }\PY{n}{a}\PY{p}{;}
\PY{+w}{    }\PY{n}{a}\PY{+w}{ }\PY{o}{=}\PY{+w}{ }\PY{n}{b}\PY{p}{;}
\PY{+w}{    }\PY{n}{b}\PY{+w}{ }\PY{o}{=}\PY{+w}{ }\PY{n}{temp}\PY{p}{;}\PY{+w}{ }\PY{c+c1}{// only swaps the local copies}
\PY{p}{\PYZcb{}}

\PY{n}{String}\PY{+w}{ }\PY{n}{x}\PY{+w}{ }\PY{o}{=}\PY{+w}{ }\PY{l+s}{\PYZdq{}}\PY{l+s}{first}\PY{l+s}{\PYZdq{}}\PY{p}{;}
\PY{n}{String}\PY{+w}{ }\PY{n}{y}\PY{+w}{ }\PY{o}{=}\PY{+w}{ }\PY{l+s}{\PYZdq{}}\PY{l+s}{second}\PY{l+s}{\PYZdq{}}\PY{p}{;}
\PY{n}{swap}\PY{p}{(}\PY{n}{x}\PY{p}{,}\PY{+w}{ }\PY{n}{y}\PY{p}{)}\PY{p}{;}
\PY{n}{System}\PY{p}{.}\PY{n+na}{out}\PY{p}{.}\PY{n+na}{println}\PY{p}{(}\PY{n}{x}\PY{+w}{ }\PY{o}{+}\PY{+w}{ }\PY{l+s}{\PYZdq{}}\PY{l+s}{ }\PY{l+s}{\PYZdq{}}\PY{+w}{ }\PY{o}{+}\PY{+w}{ }\PY{n}{y}\PY{p}{)}\PY{p}{;}\PY{+w}{ }\PY{c+c1}{// prints: first second (NOT swapped)}
\end{Verbatim}
\end{codebox}

\textbf{Method overloading} lets multiple methods share a name if their parameter lists differ. It’s resolved at compile time based on the declared (static) types of the arguments.

\begin{codebox}
\begin{Verbatim}[commandchars=\\\{\}]
\PY{k+kd}{static}\PY{+w}{ }\PY{k+kt}{void}\PY{+w}{ }\PY{n+nf}{log}\PY{p}{(}\PY{k+kt}{int}\PY{+w}{ }\PY{n}{code}\PY{p}{)}\PY{+w}{ }\PY{p}{\PYZob{}}
\PY{+w}{    }\PY{n}{System}\PY{p}{.}\PY{n+na}{out}\PY{p}{.}\PY{n+na}{println}\PY{p}{(}\PY{l+s}{\PYZdq{}}\PY{l+s}{int overload: }\PY{l+s}{\PYZdq{}}\PY{+w}{ }\PY{o}{+}\PY{+w}{ }\PY{n}{code}\PY{p}{)}\PY{p}{;}
\PY{p}{\PYZcb{}}

\PY{k+kd}{static}\PY{+w}{ }\PY{k+kt}{void}\PY{+w}{ }\PY{n+nf}{log}\PY{p}{(}\PY{n}{String}\PY{+w}{ }\PY{n}{message}\PY{p}{)}\PY{+w}{ }\PY{p}{\PYZob{}}
\PY{+w}{    }\PY{n}{System}\PY{p}{.}\PY{n+na}{out}\PY{p}{.}\PY{n+na}{println}\PY{p}{(}\PY{l+s}{\PYZdq{}}\PY{l+s}{String overload: }\PY{l+s}{\PYZdq{}}\PY{+w}{ }\PY{o}{+}\PY{+w}{ }\PY{n}{message}\PY{p}{)}\PY{p}{;}
\PY{p}{\PYZcb{}}

\PY{n}{log}\PY{p}{(}\PY{l+m+mi}{404}\PY{p}{)}\PY{p}{;}\PY{+w}{      }\PY{c+c1}{// prints: int overload: 404}
\PY{n}{log}\PY{p}{(}\PY{l+s}{\PYZdq{}}\PY{l+s}{Error}\PY{l+s}{\PYZdq{}}\PY{p}{)}\PY{p}{;}\PY{+w}{  }\PY{c+c1}{// prints: String overload: Error}
\end{Verbatim}
\end{codebox}

\section[Arrays]{Arrays}
\label{h:1:arrays}
An \textbf{array} is a fixed-size, ordered container of elements of the same type, stored contiguously in memory. Once created, its length can never change — that’s the key trade-off against dynamic structures like \code{ArrayList}.

\begin{codebox}
\begin{Verbatim}[commandchars=\\\{\}]
\PY{k+kd}{public}\PY{+w}{ }\PY{k+kd}{class} \PY{n+nc}{ArrayBasics}\PY{+w}{ }\PY{p}{\PYZob{}}
\PY{+w}{    }\PY{k+kd}{public}\PY{+w}{ }\PY{k+kd}{static}\PY{+w}{ }\PY{k+kt}{void}\PY{+w}{ }\PY{n+nf}{main}\PY{p}{(}\PY{n}{String}\PY{o}{[}\PY{o}{]}\PY{+w}{ }\PY{n}{args}\PY{p}{)}\PY{+w}{ }\PY{p}{\PYZob{}}
\PY{+w}{        }\PY{k+kt}{int}\PY{o}{[}\PY{o}{]}\PY{+w}{ }\PY{n}{scores}\PY{+w}{ }\PY{o}{=}\PY{+w}{ }\PY{k}{new}\PY{+w}{ }\PY{k+kt}{int}\PY{o}{[}\PY{l+m+mi}{3}\PY{o}{]}\PY{p}{;}\PY{+w}{      }\PY{c+c1}{// all elements default to 0}
\PY{+w}{        }\PY{n}{scores}\PY{o}{[}\PY{l+m+mi}{0}\PY{o}{]}\PY{+w}{ }\PY{o}{=}\PY{+w}{ }\PY{l+m+mi}{90}\PY{p}{;}
\PY{+w}{        }\PY{n}{scores}\PY{o}{[}\PY{l+m+mi}{1}\PY{o}{]}\PY{+w}{ }\PY{o}{=}\PY{+w}{ }\PY{l+m+mi}{75}\PY{p}{;}
\PY{+w}{        }\PY{n}{scores}\PY{o}{[}\PY{l+m+mi}{2}\PY{o}{]}\PY{+w}{ }\PY{o}{=}\PY{+w}{ }\PY{l+m+mi}{88}\PY{p}{;}

\PY{+w}{        }\PY{k+kt}{int}\PY{o}{[}\PY{o}{]}\PY{+w}{ }\PY{n}{literalScores}\PY{+w}{ }\PY{o}{=}\PY{+w}{ }\PY{p}{\PYZob{}}\PY{l+m+mi}{90}\PY{p}{,}\PY{+w}{ }\PY{l+m+mi}{75}\PY{p}{,}\PY{+w}{ }\PY{l+m+mi}{88}\PY{p}{\PYZcb{}}\PY{p}{;}\PY{+w}{ }\PY{c+c1}{// shorthand array literal}

\PY{+w}{        }\PY{n}{System}\PY{p}{.}\PY{n+na}{out}\PY{p}{.}\PY{n+na}{println}\PY{p}{(}\PY{n}{scores}\PY{p}{.}\PY{n+na}{length}\PY{p}{)}\PY{p}{;}\PY{+w}{ }\PY{c+c1}{// prints: 3}
\PY{+w}{        }\PY{n}{System}\PY{p}{.}\PY{n+na}{out}\PY{p}{.}\PY{n+na}{println}\PY{p}{(}\PY{n}{java}\PY{p}{.}\PY{n+na}{util}\PY{p}{.}\PY{n+na}{Arrays}\PY{p}{.}\PY{n+na}{toString}\PY{p}{(}\PY{n}{scores}\PY{p}{)}\PY{p}{)}\PY{p}{;}\PY{+w}{ }\PY{c+c1}{// prints: [90, 75, 88]}

\PY{+w}{        }\PY{c+c1}{// 2D array: an array of arrays}
\PY{+w}{        }\PY{k+kt}{int}\PY{o}{[}\PY{o}{]}\PY{o}{[}\PY{o}{]}\PY{+w}{ }\PY{n}{grid}\PY{+w}{ }\PY{o}{=}\PY{+w}{ }\PY{p}{\PYZob{}}
\PY{+w}{            }\PY{p}{\PYZob{}}\PY{l+m+mi}{1}\PY{p}{,}\PY{+w}{ }\PY{l+m+mi}{2}\PY{p}{,}\PY{+w}{ }\PY{l+m+mi}{3}\PY{p}{\PYZcb{}}\PY{p}{,}
\PY{+w}{            }\PY{p}{\PYZob{}}\PY{l+m+mi}{4}\PY{p}{,}\PY{+w}{ }\PY{l+m+mi}{5}\PY{p}{,}\PY{+w}{ }\PY{l+m+mi}{6}\PY{p}{\PYZcb{}}
\PY{+w}{        }\PY{p}{\PYZcb{}}\PY{p}{;}
\PY{+w}{        }\PY{n}{System}\PY{p}{.}\PY{n+na}{out}\PY{p}{.}\PY{n+na}{println}\PY{p}{(}\PY{n}{grid}\PY{o}{[}\PY{l+m+mi}{1}\PY{o}{]}\PY{o}{[}\PY{l+m+mi}{2}\PY{o}{]}\PY{p}{)}\PY{p}{;}\PY{+w}{ }\PY{c+c1}{// prints: 6}

\PY{+w}{        }\PY{c+c1}{// Iterating}
\PY{+w}{        }\PY{k}{for}\PY{+w}{ }\PY{p}{(}\PY{k+kt}{int}\PY{+w}{ }\PY{n}{score}\PY{+w}{ }\PY{p}{:}\PY{+w}{ }\PY{n}{literalScores}\PY{p}{)}\PY{+w}{ }\PY{p}{\PYZob{}}
\PY{+w}{            }\PY{n}{System}\PY{p}{.}\PY{n+na}{out}\PY{p}{.}\PY{n+na}{print}\PY{p}{(}\PY{n}{score}\PY{+w}{ }\PY{o}{+}\PY{+w}{ }\PY{l+s}{\PYZdq{}}\PY{l+s}{ }\PY{l+s}{\PYZdq{}}\PY{p}{)}\PY{p}{;}\PY{+w}{ }\PY{c+c1}{// prints: 90 75 88}
\PY{+w}{        }\PY{p}{\PYZcb{}}
\PY{+w}{        }\PY{n}{System}\PY{p}{.}\PY{n+na}{out}\PY{p}{.}\PY{n+na}{println}\PY{p}{(}\PY{p}{)}\PY{p}{;}
\PY{+w}{    }\PY{p}{\PYZcb{}}
\PY{p}{\PYZcb{}}
\end{Verbatim}
\end{codebox}

\textbf{Gotcha — arrays of objects default to \code{null}, not empty objects.}

\begin{codebox}
\begin{Verbatim}[commandchars=\\\{\}]
\PY{n}{String}\PY{o}{[}\PY{o}{]}\PY{+w}{ }\PY{n}{names}\PY{+w}{ }\PY{o}{=}\PY{+w}{ }\PY{k}{new}\PY{+w}{ }\PY{n}{String}\PY{o}{[}\PY{l+m+mi}{3}\PY{o}{]}\PY{p}{;}
\PY{n}{System}\PY{p}{.}\PY{n+na}{out}\PY{p}{.}\PY{n+na}{println}\PY{p}{(}\PY{n}{names}\PY{o}{[}\PY{l+m+mi}{0}\PY{o}{]}\PY{p}{)}\PY{p}{;}\PY{+w}{ }\PY{c+c1}{// prints: null}
\PY{c+c1}{// names[0].length(); // throws NullPointerException}
\end{Verbatim}
\end{codebox}

\textbf{Gotcha — \code{Array\allowbreak{}Index\allowbreak{}Out\allowbreak{}Of\allowbreak{}Bounds\allowbreak{}Exception} is a runtime error, not a compile-time one.}

\begin{codebox}
\begin{Verbatim}[commandchars=\\\{\}]
\PY{k+kt}{int}\PY{o}{[}\PY{o}{]}\PY{+w}{ }\PY{n}{values}\PY{+w}{ }\PY{o}{=}\PY{+w}{ }\PY{p}{\PYZob{}}\PY{l+m+mi}{1}\PY{p}{,}\PY{+w}{ }\PY{l+m+mi}{2}\PY{p}{,}\PY{+w}{ }\PY{l+m+mi}{3}\PY{p}{\PYZcb{}}\PY{p}{;}
\PY{c+c1}{// System.out.println(values[3]); // throws: ArrayIndexOutOfBoundsException: Index 3 out of bounds for length 3}
\end{Verbatim}
\end{codebox}

\textbf{Gotcha — arrays are covariant, which can cause a runtime exception.} \code{String[]} is a subtype of \code{Object[]}, so you can assign one to the other — but the JVM still remembers the real element type and enforces it at runtime:

\begin{codebox}
\begin{Verbatim}[commandchars=\\\{\}]
\PY{n}{Object}\PY{o}{[}\PY{o}{]}\PY{+w}{ }\PY{n}{objects}\PY{+w}{ }\PY{o}{=}\PY{+w}{ }\PY{k}{new}\PY{+w}{ }\PY{n}{String}\PY{o}{[}\PY{l+m+mi}{3}\PY{o}{]}\PY{p}{;}\PY{+w}{ }\PY{c+c1}{// legal: array covariance}
\PY{c+c1}{// objects[0] = Integer.valueOf(42); // throws: ArrayStoreException at runtime}
\end{Verbatim}
\end{codebox}

\textbf{Gotcha — \code{==} compares array references, not contents; \code{equals()} on arrays does too.} Use \code{Arrays.\allowbreak{}equals()} for content comparison:

\begin{codebox}
\begin{Verbatim}[commandchars=\\\{\}]
\PY{k+kt}{int}\PY{o}{[}\PY{o}{]}\PY{+w}{ }\PY{n}{a}\PY{+w}{ }\PY{o}{=}\PY{+w}{ }\PY{p}{\PYZob{}}\PY{l+m+mi}{1}\PY{p}{,}\PY{+w}{ }\PY{l+m+mi}{2}\PY{p}{,}\PY{+w}{ }\PY{l+m+mi}{3}\PY{p}{\PYZcb{}}\PY{p}{;}
\PY{k+kt}{int}\PY{o}{[}\PY{o}{]}\PY{+w}{ }\PY{n}{b}\PY{+w}{ }\PY{o}{=}\PY{+w}{ }\PY{p}{\PYZob{}}\PY{l+m+mi}{1}\PY{p}{,}\PY{+w}{ }\PY{l+m+mi}{2}\PY{p}{,}\PY{+w}{ }\PY{l+m+mi}{3}\PY{p}{\PYZcb{}}\PY{p}{;}
\PY{n}{System}\PY{p}{.}\PY{n+na}{out}\PY{p}{.}\PY{n+na}{println}\PY{p}{(}\PY{n}{a}\PY{+w}{ }\PY{o}{=}\PY{o}{=}\PY{+w}{ }\PY{n}{b}\PY{p}{)}\PY{p}{;}\PY{+w}{                       }\PY{c+c1}{// prints: false}
\PY{n}{System}\PY{p}{.}\PY{n+na}{out}\PY{p}{.}\PY{n+na}{println}\PY{p}{(}\PY{n}{a}\PY{p}{.}\PY{n+na}{equals}\PY{p}{(}\PY{n}{b}\PY{p}{)}\PY{p}{)}\PY{p}{;}\PY{+w}{                  }\PY{c+c1}{// prints: false (Object.equals, reference check)}
\PY{n}{System}\PY{p}{.}\PY{n+na}{out}\PY{p}{.}\PY{n+na}{println}\PY{p}{(}\PY{n}{java}\PY{p}{.}\PY{n+na}{util}\PY{p}{.}\PY{n+na}{Arrays}\PY{p}{.}\PY{n+na}{equals}\PY{p}{(}\PY{n}{a}\PY{p}{,}\PY{+w}{ }\PY{n}{b}\PY{p}{)}\PY{p}{)}\PY{p}{;}\PY{+w}{ }\PY{c+c1}{// prints: true}
\end{Verbatim}
\end{codebox}

\section[Strings, StringBuilder, and StringBuffer]{Strings, StringBuilder, and StringBuffer}
\label{h:1:strings-stringbuilder-and-stringbuffer}
A \code{String} in Java is \textbf{immutable} — once created, its contents can never change. Every “modification” (\code{concat}, \code{substring}, \code{replace}, \code{+}) actually creates a brand-new \code{String} object. This is why heavy string manipulation should avoid \code{String} and use \code{StringBuilder} or \code{StringBuffer} instead — mutable, resizable character containers.

\begin{codebox}
\begin{Verbatim}[commandchars=\\\{\}]
\PY{k+kd}{public}\PY{+w}{ }\PY{k+kd}{class} \PY{n+nc}{StringBasics}\PY{+w}{ }\PY{p}{\PYZob{}}
\PY{+w}{    }\PY{k+kd}{public}\PY{+w}{ }\PY{k+kd}{static}\PY{+w}{ }\PY{k+kt}{void}\PY{+w}{ }\PY{n+nf}{main}\PY{p}{(}\PY{n}{String}\PY{o}{[}\PY{o}{]}\PY{+w}{ }\PY{n}{args}\PY{p}{)}\PY{+w}{ }\PY{p}{\PYZob{}}
\PY{+w}{        }\PY{n}{String}\PY{+w}{ }\PY{n}{greeting}\PY{+w}{ }\PY{o}{=}\PY{+w}{ }\PY{l+s}{\PYZdq{}}\PY{l+s}{Hello}\PY{l+s}{\PYZdq{}}\PY{p}{;}
\PY{+w}{        }\PY{n}{String}\PY{+w}{ }\PY{n}{upper}\PY{+w}{ }\PY{o}{=}\PY{+w}{ }\PY{n}{greeting}\PY{p}{.}\PY{n+na}{toUpperCase}\PY{p}{(}\PY{p}{)}\PY{p}{;}\PY{+w}{ }\PY{c+c1}{// returns a NEW String}
\PY{+w}{        }\PY{n}{System}\PY{p}{.}\PY{n+na}{out}\PY{p}{.}\PY{n+na}{println}\PY{p}{(}\PY{n}{greeting}\PY{p}{)}\PY{p}{;}\PY{+w}{ }\PY{c+c1}{// prints: Hello (unchanged)}
\PY{+w}{        }\PY{n}{System}\PY{p}{.}\PY{n+na}{out}\PY{p}{.}\PY{n+na}{println}\PY{p}{(}\PY{n}{upper}\PY{p}{)}\PY{p}{;}\PY{+w}{    }\PY{c+c1}{// prints: HELLO}

\PY{+w}{        }\PY{c+c1}{// StringBuilder: mutable, efficient for repeated concatenation}
\PY{+w}{        }\PY{n}{StringBuilder}\PY{+w}{ }\PY{n}{sb}\PY{+w}{ }\PY{o}{=}\PY{+w}{ }\PY{k}{new}\PY{+w}{ }\PY{n}{StringBuilder}\PY{p}{(}\PY{p}{)}\PY{p}{;}
\PY{+w}{        }\PY{k}{for}\PY{+w}{ }\PY{p}{(}\PY{k+kt}{int}\PY{+w}{ }\PY{n}{i}\PY{+w}{ }\PY{o}{=}\PY{+w}{ }\PY{l+m+mi}{1}\PY{p}{;}\PY{+w}{ }\PY{n}{i}\PY{+w}{ }\PY{o}{\PYZlt{}}\PY{o}{=}\PY{+w}{ }\PY{l+m+mi}{3}\PY{p}{;}\PY{+w}{ }\PY{n}{i}\PY{o}{+}\PY{o}{+}\PY{p}{)}\PY{+w}{ }\PY{p}{\PYZob{}}
\PY{+w}{            }\PY{n}{sb}\PY{p}{.}\PY{n+na}{append}\PY{p}{(}\PY{l+s}{\PYZdq{}}\PY{l+s}{item}\PY{l+s}{\PYZdq{}}\PY{p}{)}\PY{p}{.}\PY{n+na}{append}\PY{p}{(}\PY{n}{i}\PY{p}{)}\PY{p}{.}\PY{n+na}{append}\PY{p}{(}\PY{l+s}{\PYZdq{}}\PY{l+s}{,}\PY{l+s}{\PYZdq{}}\PY{p}{)}\PY{p}{;}
\PY{+w}{        }\PY{p}{\PYZcb{}}
\PY{+w}{        }\PY{n}{sb}\PY{p}{.}\PY{n+na}{deleteCharAt}\PY{p}{(}\PY{n}{sb}\PY{p}{.}\PY{n+na}{length}\PY{p}{(}\PY{p}{)}\PY{+w}{ }\PY{o}{\PYZhy{}}\PY{+w}{ }\PY{l+m+mi}{1}\PY{p}{)}\PY{p}{;}\PY{+w}{ }\PY{c+c1}{// remove trailing comma}
\PY{+w}{        }\PY{n}{System}\PY{p}{.}\PY{n+na}{out}\PY{p}{.}\PY{n+na}{println}\PY{p}{(}\PY{n}{sb}\PY{p}{.}\PY{n+na}{toString}\PY{p}{(}\PY{p}{)}\PY{p}{)}\PY{p}{;}\PY{+w}{ }\PY{c+c1}{// prints: item1,item2,item3}
\PY{+w}{    }\PY{p}{\PYZcb{}}
\PY{p}{\PYZcb{}}
\end{Verbatim}
\end{codebox}

\textbf{The string pool.} String literals are interned (cached and reused) in a special pool. Two literals with the same content share the same object; objects created with \code{new~\allowbreak{}String(...)} do not.

\begin{codebox}
\begin{Verbatim}[commandchars=\\\{\}]
\PY{n}{String}\PY{+w}{ }\PY{n}{a}\PY{+w}{ }\PY{o}{=}\PY{+w}{ }\PY{l+s}{\PYZdq{}}\PY{l+s}{cache}\PY{l+s}{\PYZdq{}}\PY{p}{;}
\PY{n}{String}\PY{+w}{ }\PY{n}{b}\PY{+w}{ }\PY{o}{=}\PY{+w}{ }\PY{l+s}{\PYZdq{}}\PY{l+s}{cache}\PY{l+s}{\PYZdq{}}\PY{p}{;}
\PY{n}{System}\PY{p}{.}\PY{n+na}{out}\PY{p}{.}\PY{n+na}{println}\PY{p}{(}\PY{n}{a}\PY{+w}{ }\PY{o}{=}\PY{o}{=}\PY{+w}{ }\PY{n}{b}\PY{p}{)}\PY{p}{;}\PY{+w}{ }\PY{c+c1}{// prints: true (same pooled literal)}

\PY{n}{String}\PY{+w}{ }\PY{n}{c}\PY{+w}{ }\PY{o}{=}\PY{+w}{ }\PY{k}{new}\PY{+w}{ }\PY{n}{String}\PY{p}{(}\PY{l+s}{\PYZdq{}}\PY{l+s}{cache}\PY{l+s}{\PYZdq{}}\PY{p}{)}\PY{p}{;}
\PY{n}{System}\PY{p}{.}\PY{n+na}{out}\PY{p}{.}\PY{n+na}{println}\PY{p}{(}\PY{n}{a}\PY{+w}{ }\PY{o}{=}\PY{o}{=}\PY{+w}{ }\PY{n}{c}\PY{p}{)}\PY{p}{;}\PY{+w}{          }\PY{c+c1}{// prints: false (different object on the heap)}
\PY{n}{System}\PY{p}{.}\PY{n+na}{out}\PY{p}{.}\PY{n+na}{println}\PY{p}{(}\PY{n}{a}\PY{p}{.}\PY{n+na}{equals}\PY{p}{(}\PY{n}{c}\PY{p}{)}\PY{p}{)}\PY{p}{;}\PY{+w}{     }\PY{c+c1}{// prints: true (same content)}
\end{Verbatim}
\end{codebox}

\textbf{Rule of thumb:} always compare strings with \code{.\allowbreak{}equals()}, never \code{==}, unless you specifically know you’re comparing interned constants.

{\tablefont
\begin{xltabular}{\linewidth}{@{}L{0.446}L{1.592}L{0.481}L{1.481}@{}}
\toprule
\rowcolor{tablehdr}
\thd{Feature} & \thd{\code{String}} & \thd{\code{StringBuilder}} & \thd{\code{StringBuffer}} \\
\midrule
\endhead
Mutable? & No & Yes & Yes \\
Thread-safe? & Yes (immutable objects are inherently safe) & No & Yes (synchronized methods) \\
Performance & Slow for repeated edits (new object each time) & Fast & Slower than \code{StringBuilder} due to locking \\
Typical use & Fixed or rarely-changed text & Single-threaded string building & Shared mutable buffer across threads (rare today) \\
Introduced & Java 1.0 & Java 1.5 & Java 1.0 \\
\bottomrule
\end{xltabular}
}

In modern code, \code{StringBuffer} is mostly legacy. If you don’t need cross-thread safety (and usually you don’t, since string building is typically local to one method), use \code{StringBuilder}.

\textbf{Gotcha — \code{+} in a loop is quietly expensive.} The compiler optimizes a single chained \code{+} expression into a \code{StringBuilder} automatically, but each iteration of a loop using \code{+=} on a \code{String} creates a brand-new object every time:

\begin{codebox}
\begin{Verbatim}[commandchars=\\\{\}]
\PY{n}{String}\PY{+w}{ }\PY{n}{result}\PY{+w}{ }\PY{o}{=}\PY{+w}{ }\PY{l+s}{\PYZdq{}}\PY{l+s}{\PYZdq{}}\PY{p}{;}
\PY{k}{for}\PY{+w}{ }\PY{p}{(}\PY{k+kt}{int}\PY{+w}{ }\PY{n}{i}\PY{+w}{ }\PY{o}{=}\PY{+w}{ }\PY{l+m+mi}{0}\PY{p}{;}\PY{+w}{ }\PY{n}{i}\PY{+w}{ }\PY{o}{\PYZlt{}}\PY{+w}{ }\PY{l+m+mi}{5}\PY{p}{;}\PY{+w}{ }\PY{n}{i}\PY{o}{+}\PY{o}{+}\PY{p}{)}\PY{+w}{ }\PY{p}{\PYZob{}}
\PY{+w}{    }\PY{n}{result}\PY{+w}{ }\PY{o}{+}\PY{o}{=}\PY{+w}{ }\PY{n}{i}\PY{p}{;}\PY{+w}{ }\PY{c+c1}{// creates a new String object on every iteration — O(n\PYZca{}2) overall}
\PY{p}{\PYZcb{}}
\PY{n}{System}\PY{p}{.}\PY{n+na}{out}\PY{p}{.}\PY{n+na}{println}\PY{p}{(}\PY{n}{result}\PY{p}{)}\PY{p}{;}\PY{+w}{ }\PY{c+c1}{// prints: 01234}
\PY{c+c1}{// Better: use a StringBuilder and call .append(i) in the loop}
\end{Verbatim}
\end{codebox}

\textbf{Gotcha — text blocks (Java 15+)} make multi-line strings much cleaner, but indentation is significant:

\begin{codebox}
\begin{Verbatim}[commandchars=\\\{\}]
\PY{n}{String}\PY{+w}{ }\PY{n}{json}\PY{+w}{ }\PY{o}{=}\PY{+w}{ }\PY{l+s}{\PYZdq{}\PYZdq{}\PYZdq{}}
\PY{l+s}{    \PYZob{}}
\PY{l+s}{      }\PY{l+s}{\PYZdq{}}\PY{l+s}{name}\PY{l+s}{\PYZdq{}}\PY{l+s}{: }\PY{l+s}{\PYZdq{}}\PY{l+s}{Ada}\PY{l+s}{\PYZdq{}}
\PY{l+s}{    \PYZcb{}}
\PY{l+s}{    }\PY{l+s}{\PYZdq{}\PYZdq{}\PYZdq{}}\PY{p}{;}
\PY{n}{System}\PY{p}{.}\PY{n+na}{out}\PY{p}{.}\PY{n+na}{println}\PY{p}{(}\PY{n}{json}\PY{p}{.}\PY{n+na}{strip}\PY{p}{(}\PY{p}{)}\PY{p}{)}\PY{p}{;}
\PY{c+c1}{// prints:}
\PY{c+c1}{// \PYZob{}}
\PY{c+c1}{//   \PYZdq{}name\PYZdq{}: \PYZdq{}Ada\PYZdq{}}
\PY{c+c1}{// \PYZcb{}}
\end{Verbatim}
\end{codebox}

\section[Autoboxing and Unboxing]{Autoboxing and Unboxing}
\label{h:1:autoboxing-and-unboxing}
\textbf{Autoboxing} is the automatic conversion of a primitive value into its corresponding \textbf{wrapper object} (e.g., \code{int} → \code{Integer}). \textbf{Unboxing} is the reverse. The compiler inserts these conversions for you, which is convenient — but it hides real object creation and real \code{NullPointerException} risk.

{\tablefont
\begin{xltabular}{\linewidth}{@{}L{0.875}L{1.125}@{}}
\toprule
\rowcolor{tablehdr}
\thd{Primitive} & \thd{Wrapper class} \\
\midrule
\endhead
\code{byte} & \code{Byte} \\
\code{short} & \code{Short} \\
\code{int} & \code{Integer} \\
\code{long} & \code{Long} \\
\code{float} & \code{Float} \\
\code{double} & \code{Double} \\
\code{char} & \code{Character} \\
\code{boolean} & \code{Boolean} \\
\bottomrule
\end{xltabular}
}

\begin{codebox}
\begin{Verbatim}[commandchars=\\\{\}]
\PY{k+kd}{public}\PY{+w}{ }\PY{k+kd}{class} \PY{n+nc}{BoxingBasics}\PY{+w}{ }\PY{p}{\PYZob{}}
\PY{+w}{    }\PY{k+kd}{public}\PY{+w}{ }\PY{k+kd}{static}\PY{+w}{ }\PY{k+kt}{void}\PY{+w}{ }\PY{n+nf}{main}\PY{p}{(}\PY{n}{String}\PY{o}{[}\PY{o}{]}\PY{+w}{ }\PY{n}{args}\PY{p}{)}\PY{+w}{ }\PY{p}{\PYZob{}}
\PY{+w}{        }\PY{k+kt}{int}\PY{+w}{ }\PY{n}{primitiveValue}\PY{+w}{ }\PY{o}{=}\PY{+w}{ }\PY{l+m+mi}{42}\PY{p}{;}
\PY{+w}{        }\PY{n}{Integer}\PY{+w}{ }\PY{n}{boxed}\PY{+w}{ }\PY{o}{=}\PY{+w}{ }\PY{n}{primitiveValue}\PY{p}{;}\PY{+w}{     }\PY{c+c1}{// autoboxing: int \PYZhy{}\PYZgt{} Integer}
\PY{+w}{        }\PY{k+kt}{int}\PY{+w}{ }\PY{n}{unboxed}\PY{+w}{ }\PY{o}{=}\PY{+w}{ }\PY{n}{boxed}\PY{p}{;}\PY{+w}{                }\PY{c+c1}{// unboxing: Integer \PYZhy{}\PYZgt{} int}

\PY{+w}{        }\PY{n}{java}\PY{p}{.}\PY{n+na}{util}\PY{p}{.}\PY{n+na}{List}\PY{o}{\PYZlt{}}\PY{n}{Integer}\PY{o}{\PYZgt{}}\PY{+w}{ }\PY{n}{numbers}\PY{+w}{ }\PY{o}{=}\PY{+w}{ }\PY{k}{new}\PY{+w}{ }\PY{n}{java}\PY{p}{.}\PY{n+na}{util}\PY{p}{.}\PY{n+na}{ArrayList}\PY{o}{\PYZlt{}}\PY{o}{\PYZgt{}}\PY{p}{(}\PY{p}{)}\PY{p}{;}
\PY{+w}{        }\PY{n}{numbers}\PY{p}{.}\PY{n+na}{add}\PY{p}{(}\PY{l+m+mi}{10}\PY{p}{)}\PY{p}{;}\PY{+w}{ }\PY{c+c1}{// autoboxed to Integer before being stored}
\PY{+w}{        }\PY{k+kt}{int}\PY{+w}{ }\PY{n}{first}\PY{+w}{ }\PY{o}{=}\PY{+w}{ }\PY{n}{numbers}\PY{p}{.}\PY{n+na}{get}\PY{p}{(}\PY{l+m+mi}{0}\PY{p}{)}\PY{p}{;}\PY{+w}{ }\PY{c+c1}{// unboxed back to int}
\PY{+w}{        }\PY{n}{System}\PY{p}{.}\PY{n+na}{out}\PY{p}{.}\PY{n+na}{println}\PY{p}{(}\PY{n}{first}\PY{p}{)}\PY{p}{;}\PY{+w}{ }\PY{c+c1}{// prints: 10}
\PY{+w}{    }\PY{p}{\PYZcb{}}
\PY{p}{\PYZcb{}}
\end{Verbatim}
\end{codebox}

\textbf{Gotcha — unboxing \code{null} throws \code{NullPointerException}.} This is one of the most common review findings in real codebases:

\begin{codebox}
\begin{Verbatim}[commandchars=\\\{\}]
\PY{n}{Integer}\PY{+w}{ }\PY{n}{count}\PY{+w}{ }\PY{o}{=}\PY{+w}{ }\PY{k+kc}{null}\PY{p}{;}
\PY{c+c1}{// int total = count + 1; // throws NullPointerException during unboxing!}

\PY{n}{Map}\PY{o}{\PYZlt{}}\PY{n}{String}\PY{p}{,}\PY{+w}{ }\PY{n}{Integer}\PY{o}{\PYZgt{}}\PY{+w}{ }\PY{n}{inventory}\PY{+w}{ }\PY{o}{=}\PY{+w}{ }\PY{k}{new}\PY{+w}{ }\PY{n}{HashMap}\PY{o}{\PYZlt{}}\PY{o}{\PYZgt{}}\PY{p}{(}\PY{p}{)}\PY{p}{;}
\PY{c+c1}{// int qty = inventory.get(\PYZdq{}widget\PYZdq{}); // NPE if \PYZdq{}widget\PYZdq{} isn\PYZsq{}t in the map (get returns null)}
\PY{k+kt}{int}\PY{+w}{ }\PY{n}{qty}\PY{+w}{ }\PY{o}{=}\PY{+w}{ }\PY{n}{inventory}\PY{p}{.}\PY{n+na}{getOrDefault}\PY{p}{(}\PY{l+s}{\PYZdq{}}\PY{l+s}{widget}\PY{l+s}{\PYZdq{}}\PY{p}{,}\PY{+w}{ }\PY{l+m+mi}{0}\PY{p}{)}\PY{p}{;}\PY{+w}{ }\PY{c+c1}{// safe}
\end{Verbatim}
\end{codebox}

\textbf{Gotcha — the Integer cache means \code{==} sometimes “works” and sometimes doesn’t.} The JVM caches boxed \code{Integer} values from -128 to 127. Comparing cached values with \code{==} accidentally passes; comparing larger values fails:

\begin{codebox}
\begin{Verbatim}[commandchars=\\\{\}]
\PY{n}{Integer}\PY{+w}{ }\PY{n}{a}\PY{+w}{ }\PY{o}{=}\PY{+w}{ }\PY{l+m+mi}{100}\PY{p}{;}
\PY{n}{Integer}\PY{+w}{ }\PY{n}{b}\PY{+w}{ }\PY{o}{=}\PY{+w}{ }\PY{l+m+mi}{100}\PY{p}{;}
\PY{n}{System}\PY{p}{.}\PY{n+na}{out}\PY{p}{.}\PY{n+na}{println}\PY{p}{(}\PY{n}{a}\PY{+w}{ }\PY{o}{=}\PY{o}{=}\PY{+w}{ }\PY{n}{b}\PY{p}{)}\PY{p}{;}\PY{+w}{ }\PY{c+c1}{// prints: true (both pulled from the internal cache)}

\PY{n}{Integer}\PY{+w}{ }\PY{n}{c}\PY{+w}{ }\PY{o}{=}\PY{+w}{ }\PY{l+m+mi}{200}\PY{p}{;}
\PY{n}{Integer}\PY{+w}{ }\PY{n}{d}\PY{+w}{ }\PY{o}{=}\PY{+w}{ }\PY{l+m+mi}{200}\PY{p}{;}
\PY{n}{System}\PY{p}{.}\PY{n+na}{out}\PY{p}{.}\PY{n+na}{println}\PY{p}{(}\PY{n}{c}\PY{+w}{ }\PY{o}{=}\PY{o}{=}\PY{+w}{ }\PY{n}{d}\PY{p}{)}\PY{p}{;}\PY{+w}{ }\PY{c+c1}{// prints: false (outside cache range, different objects)}

\PY{n}{System}\PY{p}{.}\PY{n+na}{out}\PY{p}{.}\PY{n+na}{println}\PY{p}{(}\PY{n}{c}\PY{p}{.}\PY{n+na}{equals}\PY{p}{(}\PY{n}{d}\PY{p}{)}\PY{p}{)}\PY{p}{;}\PY{+w}{ }\PY{c+c1}{// prints: true (always correct — compare wrapper objects with equals)}
\end{Verbatim}
\end{codebox}

\textbf{Gotcha — mixing primitives and wrappers in arithmetic can silently unbox and throw.} Also, autoboxing in tight loops (e.g., \code{Long~\allowbreak{}sum~\allowbreak{}=\allowbreak{}~\allowbreak{}0L;\allowbreak{}~\allowbreak{}sum~\allowbreak{}+=\allowbreak{}~\allowbreak{}i;} inside a million-iteration loop) creates a huge number of short-lived wrapper objects, hurting performance. Prefer primitive accumulator variables in hot loops.

\section[Varargs]{Varargs}
\label{h:1:varargs}
\textbf{Varargs} (variable-length arguments) let a method accept zero or more arguments of a type without the caller having to build an array explicitly. Internally, the compiler packages the arguments into an array. The syntax is \code{Type...\allowbreak{}~\allowbreak{}name}, and it must be the \emph{last} parameter in the method signature.

\begin{codebox}
\begin{Verbatim}[commandchars=\\\{\}]
\PY{k+kd}{public}\PY{+w}{ }\PY{k+kd}{class} \PY{n+nc}{VarargsBasics}\PY{+w}{ }\PY{p}{\PYZob{}}

\PY{+w}{    }\PY{k+kd}{static}\PY{+w}{ }\PY{k+kt}{int}\PY{+w}{ }\PY{n+nf}{sum}\PY{p}{(}\PY{k+kt}{int}\PY{p}{.}\PY{p}{.}\PY{p}{.}\PY{+w}{ }\PY{n}{numbers}\PY{p}{)}\PY{+w}{ }\PY{p}{\PYZob{}}
\PY{+w}{        }\PY{k+kt}{int}\PY{+w}{ }\PY{n}{total}\PY{+w}{ }\PY{o}{=}\PY{+w}{ }\PY{l+m+mi}{0}\PY{p}{;}
\PY{+w}{        }\PY{k}{for}\PY{+w}{ }\PY{p}{(}\PY{k+kt}{int}\PY{+w}{ }\PY{n}{n}\PY{+w}{ }\PY{p}{:}\PY{+w}{ }\PY{n}{numbers}\PY{p}{)}\PY{+w}{ }\PY{p}{\PYZob{}}
\PY{+w}{            }\PY{n}{total}\PY{+w}{ }\PY{o}{+}\PY{o}{=}\PY{+w}{ }\PY{n}{n}\PY{p}{;}
\PY{+w}{        }\PY{p}{\PYZcb{}}
\PY{+w}{        }\PY{k}{return}\PY{+w}{ }\PY{n}{total}\PY{p}{;}
\PY{+w}{    }\PY{p}{\PYZcb{}}

\PY{+w}{    }\PY{k+kd}{public}\PY{+w}{ }\PY{k+kd}{static}\PY{+w}{ }\PY{k+kt}{void}\PY{+w}{ }\PY{n+nf}{main}\PY{p}{(}\PY{n}{String}\PY{o}{[}\PY{o}{]}\PY{+w}{ }\PY{n}{args}\PY{p}{)}\PY{+w}{ }\PY{p}{\PYZob{}}
\PY{+w}{        }\PY{n}{System}\PY{p}{.}\PY{n+na}{out}\PY{p}{.}\PY{n+na}{println}\PY{p}{(}\PY{n}{sum}\PY{p}{(}\PY{p}{)}\PY{p}{)}\PY{p}{;}\PY{+w}{           }\PY{c+c1}{// prints: 0  (zero arguments is allowed)}
\PY{+w}{        }\PY{n}{System}\PY{p}{.}\PY{n+na}{out}\PY{p}{.}\PY{n+na}{println}\PY{p}{(}\PY{n}{sum}\PY{p}{(}\PY{l+m+mi}{5}\PY{p}{)}\PY{p}{)}\PY{p}{;}\PY{+w}{           }\PY{c+c1}{// prints: 5}
\PY{+w}{        }\PY{n}{System}\PY{p}{.}\PY{n+na}{out}\PY{p}{.}\PY{n+na}{println}\PY{p}{(}\PY{n}{sum}\PY{p}{(}\PY{l+m+mi}{1}\PY{p}{,}\PY{+w}{ }\PY{l+m+mi}{2}\PY{p}{,}\PY{+w}{ }\PY{l+m+mi}{3}\PY{p}{,}\PY{+w}{ }\PY{l+m+mi}{4}\PY{p}{)}\PY{p}{)}\PY{p}{;}\PY{+w}{   }\PY{c+c1}{// prints: 10}

\PY{+w}{        }\PY{k+kt}{int}\PY{o}{[}\PY{o}{]}\PY{+w}{ }\PY{n}{existingArray}\PY{+w}{ }\PY{o}{=}\PY{+w}{ }\PY{p}{\PYZob{}}\PY{l+m+mi}{10}\PY{p}{,}\PY{+w}{ }\PY{l+m+mi}{20}\PY{p}{,}\PY{+w}{ }\PY{l+m+mi}{30}\PY{p}{\PYZcb{}}\PY{p}{;}
\PY{+w}{        }\PY{n}{System}\PY{p}{.}\PY{n+na}{out}\PY{p}{.}\PY{n+na}{println}\PY{p}{(}\PY{n}{sum}\PY{p}{(}\PY{n}{existingArray}\PY{p}{)}\PY{p}{)}\PY{p}{;}\PY{+w}{ }\PY{c+c1}{// prints: 60 (an array can be passed directly)}
\PY{+w}{    }\PY{p}{\PYZcb{}}
\PY{p}{\PYZcb{}}
\end{Verbatim}
\end{codebox}

You’ve already used varargs if you’ve called \code{String.\allowbreak{}format(...)}, \code{List.\allowbreak{}of(...)}, or \code{System.\allowbreak{}out.\allowbreak{}printf(...)}.

\textbf{Gotcha — varargs and overload resolution can be ambiguous or surprising.} A specific-arity overload is always preferred over the varargs version when both match:

\begin{codebox}
\begin{Verbatim}[commandchars=\\\{\}]
\PY{k+kd}{static}\PY{+w}{ }\PY{k+kt}{void}\PY{+w}{ }\PY{n+nf}{report}\PY{p}{(}\PY{k+kt}{int}\PY{+w}{ }\PY{n}{a}\PY{p}{,}\PY{+w}{ }\PY{k+kt}{int}\PY{+w}{ }\PY{n}{b}\PY{p}{)}\PY{+w}{ }\PY{p}{\PYZob{}}
\PY{+w}{    }\PY{n}{System}\PY{p}{.}\PY{n+na}{out}\PY{p}{.}\PY{n+na}{println}\PY{p}{(}\PY{l+s}{\PYZdq{}}\PY{l+s}{specific overload}\PY{l+s}{\PYZdq{}}\PY{p}{)}\PY{p}{;}
\PY{p}{\PYZcb{}}

\PY{k+kd}{static}\PY{+w}{ }\PY{k+kt}{void}\PY{+w}{ }\PY{n+nf}{report}\PY{p}{(}\PY{k+kt}{int}\PY{p}{.}\PY{p}{.}\PY{p}{.}\PY{+w}{ }\PY{n}{values}\PY{p}{)}\PY{+w}{ }\PY{p}{\PYZob{}}
\PY{+w}{    }\PY{n}{System}\PY{p}{.}\PY{n+na}{out}\PY{p}{.}\PY{n+na}{println}\PY{p}{(}\PY{l+s}{\PYZdq{}}\PY{l+s}{varargs overload}\PY{l+s}{\PYZdq{}}\PY{p}{)}\PY{p}{;}
\PY{p}{\PYZcb{}}

\PY{n}{report}\PY{p}{(}\PY{l+m+mi}{1}\PY{p}{,}\PY{+w}{ }\PY{l+m+mi}{2}\PY{p}{)}\PY{p}{;}\PY{+w}{ }\PY{c+c1}{// prints: specific overload (exact match wins over varargs)}
\PY{n}{report}\PY{p}{(}\PY{l+m+mi}{1}\PY{p}{,}\PY{+w}{ }\PY{l+m+mi}{2}\PY{p}{,}\PY{+w}{ }\PY{l+m+mi}{3}\PY{p}{)}\PY{p}{;}\PY{+w}{ }\PY{c+c1}{// prints: varargs overload (no exact match exists)}
\end{Verbatim}
\end{codebox}

\textbf{Gotcha — mixing varargs with generics risks heap pollution.} The compiler warns about “unchecked generics array creation” because a varargs array of a generic type can be assigned to \code{Object[]} and have incompatible elements inserted, breaking type safety at runtime. Annotate such methods with \code{@\allowbreak{}SafeVarargs} only if you’ve verified the method doesn’t misuse the array.

\begin{codebox}
\begin{Verbatim}[commandchars=\\\{\}]
\PY{n+nd}{@SafeVarargs}
\PY{k+kd}{static}\PY{+w}{ }\PY{o}{\PYZlt{}}\PY{n}{T}\PY{o}{\PYZgt{}}\PY{+w}{ }\PY{n}{java}\PY{p}{.}\PY{n+na}{util}\PY{p}{.}\PY{n+na}{List}\PY{o}{\PYZlt{}}\PY{n}{T}\PY{o}{\PYZgt{}}\PY{+w}{ }\PY{n+nf}{listOf}\PY{p}{(}\PY{n}{T}\PY{p}{.}\PY{p}{.}\PY{p}{.}\PY{+w}{ }\PY{n}{items}\PY{p}{)}\PY{+w}{ }\PY{p}{\PYZob{}}
\PY{+w}{    }\PY{k}{return}\PY{+w}{ }\PY{n}{java}\PY{p}{.}\PY{n+na}{util}\PY{p}{.}\PY{n+na}{Arrays}\PY{p}{.}\PY{n+na}{asList}\PY{p}{(}\PY{n}{items}\PY{p}{)}\PY{p}{;}
\PY{p}{\PYZcb{}}
\end{Verbatim}
\end{codebox}

\textbf{Gotcha — an ambiguous null argument.} Passing \code{null} directly to a varargs method is ambiguous to a human reader (and sometimes to the compiler, if multiple overloads exist) — it’s treated as a \code{null} \emph{array}, not a single \code{null} element:

\begin{codebox}
\begin{Verbatim}[commandchars=\\\{\}]
\PY{k+kd}{static}\PY{+w}{ }\PY{k+kt}{void}\PY{+w}{ }\PY{n+nf}{printAll}\PY{p}{(}\PY{n}{String}\PY{p}{.}\PY{p}{.}\PY{p}{.}\PY{+w}{ }\PY{n}{items}\PY{p}{)}\PY{+w}{ }\PY{p}{\PYZob{}}
\PY{+w}{    }\PY{n}{System}\PY{p}{.}\PY{n+na}{out}\PY{p}{.}\PY{n+na}{println}\PY{p}{(}\PY{n}{items}\PY{+w}{ }\PY{o}{=}\PY{o}{=}\PY{+w}{ }\PY{k+kc}{null}\PY{+w}{ }\PY{o}{?}\PY{+w}{ }\PY{l+s}{\PYZdq{}}\PY{l+s}{null array}\PY{l+s}{\PYZdq{}}\PY{+w}{ }\PY{p}{:}\PY{+w}{ }\PY{n}{items}\PY{p}{.}\PY{n+na}{length}\PY{+w}{ }\PY{o}{+}\PY{+w}{ }\PY{l+s}{\PYZdq{}}\PY{l+s}{ items}\PY{l+s}{\PYZdq{}}\PY{p}{)}\PY{p}{;}
\PY{p}{\PYZcb{}}

\PY{n}{printAll}\PY{p}{(}\PY{k+kc}{null}\PY{p}{)}\PY{p}{;}\PY{+w}{ }\PY{c+c1}{// prints: null array (the whole array reference is null)}
\PY{n}{printAll}\PY{p}{(}\PY{p}{(}\PY{n}{String}\PY{p}{)}\PY{+w}{ }\PY{k+kc}{null}\PY{p}{)}\PY{p}{;}\PY{+w}{ }\PY{c+c1}{// prints: 1 items (a single\PYZhy{}element array containing null)}
\end{Verbatim}
\end{codebox}

\section[Local Variable Type Inference (var)]{Local Variable Type Inference (var)}
\label{h:1:local-variable-type-inference-var}
Since Java 10, the \code{var} keyword lets the compiler infer a local variable’s type from the right-hand side of its initialization. It is \textbf{not} dynamic typing — the type is fixed at compile time, just written implicitly. \code{var} only works for local variables with an initializer; it cannot be used for fields, method parameters, or return types.

\begin{codebox}
\begin{Verbatim}[commandchars=\\\{\}]
\PY{k+kd}{public}\PY{+w}{ }\PY{k+kd}{class} \PY{n+nc}{VarBasics}\PY{+w}{ }\PY{p}{\PYZob{}}
\PY{+w}{    }\PY{k+kd}{public}\PY{+w}{ }\PY{k+kd}{static}\PY{+w}{ }\PY{k+kt}{void}\PY{+w}{ }\PY{n+nf}{main}\PY{p}{(}\PY{n}{String}\PY{o}{[}\PY{o}{]}\PY{+w}{ }\PY{n}{args}\PY{p}{)}\PY{+w}{ }\PY{p}{\PYZob{}}
\PY{+w}{        }\PY{k+kd}{var}\PY{+w}{ }\PY{n}{name}\PY{+w}{ }\PY{o}{=}\PY{+w}{ }\PY{l+s}{\PYZdq{}}\PY{l+s}{Alice}\PY{l+s}{\PYZdq{}}\PY{p}{;}\PY{+w}{                 }\PY{c+c1}{// inferred as String}
\PY{+w}{        }\PY{k+kd}{var}\PY{+w}{ }\PY{n}{count}\PY{+w}{ }\PY{o}{=}\PY{+w}{ }\PY{l+m+mi}{10}\PY{p}{;}\PY{+w}{                      }\PY{c+c1}{// inferred as int}
\PY{+w}{        }\PY{k+kd}{var}\PY{+w}{ }\PY{n}{prices}\PY{+w}{ }\PY{o}{=}\PY{+w}{ }\PY{k}{new}\PY{+w}{ }\PY{n}{java}\PY{p}{.}\PY{n+na}{util}\PY{p}{.}\PY{n+na}{ArrayList}\PY{o}{\PYZlt{}}\PY{n}{Double}\PY{o}{\PYZgt{}}\PY{p}{(}\PY{p}{)}\PY{p}{;}\PY{+w}{ }\PY{c+c1}{// inferred as ArrayList\PYZlt{}Double\PYZgt{}}
\PY{+w}{        }\PY{n}{prices}\PY{p}{.}\PY{n+na}{add}\PY{p}{(}\PY{l+m+mf}{19.99}\PY{p}{)}\PY{p}{;}

\PY{+w}{        }\PY{c+c1}{// name = 42; // COMPILE ERROR: name\PYZsq{}s type is fixed as String, not dynamic}

\PY{+w}{        }\PY{k}{for}\PY{+w}{ }\PY{p}{(}\PY{k+kd}{var}\PY{+w}{ }\PY{n}{price}\PY{+w}{ }\PY{p}{:}\PY{+w}{ }\PY{n}{prices}\PY{p}{)}\PY{+w}{ }\PY{p}{\PYZob{}}\PY{+w}{ }\PY{c+c1}{// var also works in for\PYZhy{}each loops}
\PY{+w}{            }\PY{n}{System}\PY{p}{.}\PY{n+na}{out}\PY{p}{.}\PY{n+na}{println}\PY{p}{(}\PY{n}{price}\PY{p}{)}\PY{p}{;}\PY{+w}{ }\PY{c+c1}{// prints: 19.99}
\PY{+w}{        }\PY{p}{\PYZcb{}}
\PY{+w}{    }\PY{p}{\PYZcb{}}
\PY{p}{\PYZcb{}}
\end{Verbatim}
\end{codebox}

\textbf{Gotcha — \code{var} requires an initializer, and can’t infer from \code{null} alone.}

\begin{codebox}
\begin{Verbatim}[commandchars=\\\{\}]
\PY{c+c1}{// var x; // COMPILE ERROR: cannot infer type without an initializer}
\PY{c+c1}{// var y = null; // COMPILE ERROR: cannot infer type from null}
\PY{k+kd}{var}\PY{+w}{ }\PY{n}{z}\PY{+w}{ }\PY{o}{=}\PY{+w}{ }\PY{p}{(}\PY{n}{String}\PY{p}{)}\PY{+w}{ }\PY{k+kc}{null}\PY{p}{;}\PY{+w}{ }\PY{c+c1}{// legal: explicit cast tells the compiler the type}
\end{Verbatim}
\end{codebox}

\textbf{Gotcha — \code{var} can obscure the actual type in a review, hurting readability.} This compiles but is a common review comment:

\begin{codebox}
\begin{Verbatim}[commandchars=\\\{\}]
\PY{k+kd}{var}\PY{+w}{ }\PY{n}{result}\PY{+w}{ }\PY{o}{=}\PY{+w}{ }\PY{n}{service}\PY{p}{.}\PY{n+na}{process}\PY{p}{(}\PY{n}{request}\PY{p}{)}\PY{p}{;}\PY{+w}{ }\PY{c+c1}{// what type is \PYZdq{}result\PYZdq{}? Not obvious from this line alone.}
\end{Verbatim}
\end{codebox}

\code{var} is best used when the type is already obvious from the right-hand side (e.g., \code{var~\allowbreak{}list~\allowbreak{}=\allowbreak{}~\allowbreak{}new~\allowbreak{}Array\allowbreak{}List\textless{}\allowbreak{}String\textgreater{}();}) and best avoided when it hides meaningful information (e.g., the return type of a method call the reader doesn’t already know).

\textbf{Gotcha — \code{var} with diamond operator loses generic type information if there’s no explicit type argument.}

\begin{codebox}
\begin{Verbatim}[commandchars=\\\{\}]
\PY{k+kd}{var}\PY{+w}{ }\PY{n}{list}\PY{+w}{ }\PY{o}{=}\PY{+w}{ }\PY{k}{new}\PY{+w}{ }\PY{n}{ArrayList}\PY{o}{\PYZlt{}}\PY{o}{\PYZgt{}}\PY{p}{(}\PY{p}{)}\PY{p}{;}\PY{+w}{ }\PY{c+c1}{// inferred as ArrayList\PYZlt{}Object\PYZgt{}, NOT the type you might expect}
\PY{n}{list}\PY{p}{.}\PY{n+na}{add}\PY{p}{(}\PY{l+s}{\PYZdq{}}\PY{l+s}{text}\PY{l+s}{\PYZdq{}}\PY{p}{)}\PY{p}{;}
\PY{n}{list}\PY{p}{.}\PY{n+na}{add}\PY{p}{(}\PY{l+m+mi}{42}\PY{p}{)}\PY{p}{;}\PY{+w}{ }\PY{c+c1}{// compiles fine — the list is untyped (Object), losing compile\PYZhy{}time safety}
\end{Verbatim}
\end{codebox}

\section[Common Code-Review Interview Pitfalls]{Common Code-Review Interview Pitfalls}
\label{h:1:common-code-review-interview-pitfalls}
\begin{enumerate}
\item 
\textbf{Using \code{==} to compare \code{String} or wrapper objects instead of \code{.\allowbreak{}equals()}.} Why it matters: \code{==} compares references (or cached instances), not logical content, and fails unpredictably outside the small-integer cache or literal pool.

\begin{codebox}
\begin{Verbatim}[commandchars=\\\{\}]
\PY{c+c1}{// Before}
\PY{k}{if}\PY{+w}{ }\PY{p}{(}\PY{n}{userInput}\PY{+w}{ }\PY{o}{=}\PY{o}{=}\PY{+w}{ }\PY{l+s}{\PYZdq{}}\PY{l+s}{admin}\PY{l+s}{\PYZdq{}}\PY{p}{)}\PY{+w}{ }\PY{p}{\PYZob{}}\PY{+w}{ }\PY{p}{.}\PY{p}{.}\PY{p}{.}\PY{+w}{ }\PY{p}{\PYZcb{}}
\PY{c+c1}{// After}
\PY{k}{if}\PY{+w}{ }\PY{p}{(}\PY{l+s}{\PYZdq{}}\PY{l+s}{admin}\PY{l+s}{\PYZdq{}}\PY{p}{.}\PY{n+na}{equals}\PY{p}{(}\PY{n}{userInput}\PY{p}{)}\PY{p}{)}\PY{+w}{ }\PY{p}{\PYZob{}}\PY{+w}{ }\PY{p}{.}\PY{p}{.}\PY{p}{.}\PY{+w}{ }\PY{p}{\PYZcb{}}
\end{Verbatim}
\end{codebox}

\item 
\textbf{Unboxing a possibly-null wrapper type.} Why it matters: it throws \code{NullPointerException} at the unboxing point, often far from where the \code{null} originated, making the bug hard to trace.

\begin{codebox}
\begin{Verbatim}[commandchars=\\\{\}]
\PY{c+c1}{// Before}
\PY{k+kt}{int}\PY{+w}{ }\PY{n}{total}\PY{+w}{ }\PY{o}{=}\PY{+w}{ }\PY{n}{getCachedCount}\PY{p}{(}\PY{n}{id}\PY{p}{)}\PY{+w}{ }\PY{o}{+}\PY{+w}{ }\PY{l+m+mi}{1}\PY{p}{;}\PY{+w}{ }\PY{c+c1}{// getCachedCount may return null}
\PY{c+c1}{// After}
\PY{k+kt}{int}\PY{+w}{ }\PY{n}{total}\PY{+w}{ }\PY{o}{=}\PY{+w}{ }\PY{n}{Optional}\PY{p}{.}\PY{n+na}{ofNullable}\PY{p}{(}\PY{n}{getCachedCount}\PY{p}{(}\PY{n}{id}\PY{p}{)}\PY{p}{)}\PY{p}{.}\PY{n+na}{orElse}\PY{p}{(}\PY{l+m+mi}{0}\PY{p}{)}\PY{+w}{ }\PY{o}{+}\PY{+w}{ }\PY{l+m+mi}{1}\PY{p}{;}
\end{Verbatim}
\end{codebox}

\item 
\textbf{Building strings with \code{+=} inside a loop.} Why it matters: each iteration allocates a new \code{String}, turning an O(n) loop into O(n²) work.

\begin{codebox}
\begin{Verbatim}[commandchars=\\\{\}]
\PY{c+c1}{// Before}
\PY{n}{String}\PY{+w}{ }\PY{n}{csv}\PY{+w}{ }\PY{o}{=}\PY{+w}{ }\PY{l+s}{\PYZdq{}}\PY{l+s}{\PYZdq{}}\PY{p}{;}
\PY{k}{for}\PY{+w}{ }\PY{p}{(}\PY{n}{String}\PY{+w}{ }\PY{n}{s}\PY{+w}{ }\PY{p}{:}\PY{+w}{ }\PY{n}{items}\PY{p}{)}\PY{+w}{ }\PY{n}{csv}\PY{+w}{ }\PY{o}{+}\PY{o}{=}\PY{+w}{ }\PY{n}{s}\PY{+w}{ }\PY{o}{+}\PY{+w}{ }\PY{l+s}{\PYZdq{}}\PY{l+s}{,}\PY{l+s}{\PYZdq{}}\PY{p}{;}
\PY{c+c1}{// After}
\PY{n}{StringBuilder}\PY{+w}{ }\PY{n}{csv}\PY{+w}{ }\PY{o}{=}\PY{+w}{ }\PY{k}{new}\PY{+w}{ }\PY{n}{StringBuilder}\PY{p}{(}\PY{p}{)}\PY{p}{;}
\PY{k}{for}\PY{+w}{ }\PY{p}{(}\PY{n}{String}\PY{+w}{ }\PY{n}{s}\PY{+w}{ }\PY{p}{:}\PY{+w}{ }\PY{n}{items}\PY{p}{)}\PY{+w}{ }\PY{n}{csv}\PY{p}{.}\PY{n+na}{append}\PY{p}{(}\PY{n}{s}\PY{p}{)}\PY{p}{.}\PY{n+na}{append}\PY{p}{(}\PY{l+s}{\PYZdq{}}\PY{l+s}{,}\PY{l+s}{\PYZdq{}}\PY{p}{)}\PY{p}{;}
\end{Verbatim}
\end{codebox}

\item 
\textbf{Forgetting \code{break} in a classic \code{switch} statement.} Why it matters: execution silently falls through into the next case, running unintended code.

\begin{codebox}
\begin{Verbatim}[commandchars=\\\{\}]
\PY{c+c1}{// Before}
\PY{k}{switch}\PY{+w}{ }\PY{p}{(}\PY{n}{status}\PY{p}{)}\PY{+w}{ }\PY{p}{\PYZob{}}
\PY{+w}{    }\PY{k}{case}\PY{+w}{ }\PY{n}{ACTIVE}\PY{p}{:}\PY{+w}{ }\PY{n}{enable}\PY{p}{(}\PY{p}{)}\PY{p}{;}\PY{+w}{ }\PY{c+c1}{// falls through to PENDING\PYZsq{}s code too!}
\PY{+w}{    }\PY{k}{case}\PY{+w}{ }\PY{n}{PENDING}\PY{p}{:}\PY{+w}{ }\PY{n}{notify}\PY{p}{(}\PY{p}{)}\PY{p}{;}
\PY{p}{\PYZcb{}}
\PY{c+c1}{// After: add break, or better, use a switch expression with \PYZhy{}\PYZgt{} arrows}
\end{Verbatim}
\end{codebox}

\item 
\textbf{Assuming a method can reassign the caller’s object reference.} Why it matters: Java is pass-by-value for references, so reassigning a parameter inside a method never affects the caller’s variable — this is a frequent source of “why didn’t my fix take effect” bugs.

\begin{codebox}
\begin{Verbatim}[commandchars=\\\{\}]
\PY{c+c1}{// Before (does nothing useful)}
\PY{k+kd}{static}\PY{+w}{ }\PY{k+kt}{void}\PY{+w}{ }\PY{n+nf}{reset}\PY{p}{(}\PY{n}{List}\PY{o}{\PYZlt{}}\PY{n}{String}\PY{o}{\PYZgt{}}\PY{+w}{ }\PY{n}{list}\PY{p}{)}\PY{+w}{ }\PY{p}{\PYZob{}}\PY{+w}{ }\PY{n}{list}\PY{+w}{ }\PY{o}{=}\PY{+w}{ }\PY{k}{new}\PY{+w}{ }\PY{n}{ArrayList}\PY{o}{\PYZlt{}}\PY{o}{\PYZgt{}}\PY{p}{(}\PY{p}{)}\PY{p}{;}\PY{+w}{ }\PY{p}{\PYZcb{}}
\PY{c+c1}{// After}
\PY{k+kd}{static}\PY{+w}{ }\PY{n}{List}\PY{o}{\PYZlt{}}\PY{n}{String}\PY{o}{\PYZgt{}}\PY{+w}{ }\PY{n+nf}{reset}\PY{p}{(}\PY{n}{List}\PY{o}{\PYZlt{}}\PY{n}{String}\PY{o}{\PYZgt{}}\PY{+w}{ }\PY{n}{list}\PY{p}{)}\PY{+w}{ }\PY{p}{\PYZob{}}\PY{+w}{ }\PY{k}{return}\PY{+w}{ }\PY{k}{new}\PY{+w}{ }\PY{n}{ArrayList}\PY{o}{\PYZlt{}}\PY{o}{\PYZgt{}}\PY{p}{(}\PY{p}{)}\PY{p}{;}\PY{+w}{ }\PY{p}{\PYZcb{}}
\end{Verbatim}
\end{codebox}

\item 
\textbf{Relying on floating-point (\code{double}/\code{float}) for monetary or exact decimal calculations.} Why it matters: binary floating-point can’t represent many decimal fractions exactly, causing rounding errors that compound over many operations.

\begin{codebox}
\begin{Verbatim}[commandchars=\\\{\}]
\PY{c+c1}{// Before}
\PY{k+kt}{double}\PY{+w}{ }\PY{n}{total}\PY{+w}{ }\PY{o}{=}\PY{+w}{ }\PY{l+m+mf}{0.1}\PY{+w}{ }\PY{o}{+}\PY{+w}{ }\PY{l+m+mf}{0.2}\PY{p}{;}\PY{+w}{ }\PY{c+c1}{// 0.30000000000000004}
\PY{c+c1}{// After}
\PY{n}{BigDecimal}\PY{+w}{ }\PY{n}{total}\PY{+w}{ }\PY{o}{=}\PY{+w}{ }\PY{k}{new}\PY{+w}{ }\PY{n}{BigDecimal}\PY{p}{(}\PY{l+s}{\PYZdq{}}\PY{l+s}{0.1}\PY{l+s}{\PYZdq{}}\PY{p}{)}\PY{p}{.}\PY{n+na}{add}\PY{p}{(}\PY{k}{new}\PY{+w}{ }\PY{n}{BigDecimal}\PY{p}{(}\PY{l+s}{\PYZdq{}}\PY{l+s}{0.2}\PY{l+s}{\PYZdq{}}\PY{p}{)}\PY{p}{)}\PY{p}{;}\PY{+w}{ }\PY{c+c1}{// 0.3}
\end{Verbatim}
\end{codebox}

\item 
\textbf{Ignoring silent integer overflow.} Why it matters: arithmetic that exceeds a type’s range wraps around without any warning or exception, producing corrupted values that pass silently into further logic.

\begin{codebox}
\begin{Verbatim}[commandchars=\\\{\}]
\PY{c+c1}{// Before}
\PY{k+kt}{int}\PY{+w}{ }\PY{n}{total}\PY{+w}{ }\PY{o}{=}\PY{+w}{ }\PY{n}{bigValue1}\PY{+w}{ }\PY{o}{*}\PY{+w}{ }\PY{n}{bigValue2}\PY{p}{;}\PY{+w}{ }\PY{c+c1}{// may silently overflow}
\PY{c+c1}{// After}
\PY{k+kt}{long}\PY{+w}{ }\PY{n}{total}\PY{+w}{ }\PY{o}{=}\PY{+w}{ }\PY{n}{Math}\PY{p}{.}\PY{n+na}{multiplyExact}\PY{p}{(}\PY{p}{(}\PY{k+kt}{long}\PY{p}{)}\PY{+w}{ }\PY{n}{bigValue1}\PY{p}{,}\PY{+w}{ }\PY{n}{bigValue2}\PY{p}{)}\PY{p}{;}\PY{+w}{ }\PY{c+c1}{// throws on overflow}
\end{Verbatim}
\end{codebox}

\item 
\textbf{Using \code{var} where the inferred type isn’t obvious to the reader.} Why it matters: it trades a small typing convenience for reduced readability, especially with method calls whose return type isn’t well known.

\begin{codebox}
\begin{Verbatim}[commandchars=\\\{\}]
\PY{c+c1}{// Before}
\PY{k+kd}{var}\PY{+w}{ }\PY{n}{result}\PY{+w}{ }\PY{o}{=}\PY{+w}{ }\PY{n}{repository}\PY{p}{.}\PY{n+na}{fetch}\PY{p}{(}\PY{n}{id}\PY{p}{)}\PY{p}{;}\PY{+w}{ }\PY{c+c1}{// unclear type at a glance}
\PY{c+c1}{// After}
\PY{n}{CustomerRecord}\PY{+w}{ }\PY{n}{result}\PY{+w}{ }\PY{o}{=}\PY{+w}{ }\PY{n}{repository}\PY{p}{.}\PY{n+na}{fetch}\PY{p}{(}\PY{n}{id}\PY{p}{)}\PY{p}{;}
\end{Verbatim}
\end{codebox}

\item 
\textbf{Comparing or checking array contents with \code{==} or \code{.\allowbreak{}equals()}.} Why it matters: arrays don’t override \code{equals()}, so both operators check reference identity, not element-by-element equality — a classic false-negative bug.

\begin{codebox}
\begin{Verbatim}[commandchars=\\\{\}]
\PY{c+c1}{// Before}
\PY{k}{if}\PY{+w}{ }\PY{p}{(}\PY{n}{expectedBytes}\PY{p}{.}\PY{n+na}{equals}\PY{p}{(}\PY{n}{actualBytes}\PY{p}{)}\PY{p}{)}\PY{+w}{ }\PY{p}{\PYZob{}}\PY{+w}{ }\PY{p}{.}\PY{p}{.}\PY{p}{.}\PY{+w}{ }\PY{p}{\PYZcb{}}\PY{+w}{ }\PY{c+c1}{// always compares references}
\PY{c+c1}{// After}
\PY{k}{if}\PY{+w}{ }\PY{p}{(}\PY{n}{Arrays}\PY{p}{.}\PY{n+na}{equals}\PY{p}{(}\PY{n}{expectedBytes}\PY{p}{,}\PY{+w}{ }\PY{n}{actualBytes}\PY{p}{)}\PY{p}{)}\PY{+w}{ }\PY{p}{\PYZob{}}\PY{+w}{ }\PY{p}{.}\PY{p}{.}\PY{p}{.}\PY{+w}{ }\PY{p}{\PYZcb{}}
\end{Verbatim}
\end{codebox}

\item 
\textbf{Passing a single \code{null} to a varargs parameter and expecting one null element.} Why it matters: an untyped \code{null} is interpreted as “no array at all,” not “an array containing one null,” which can cause a \code{NullPointerException} inside the method instead of the intended behavior.

\begin{codebox}
\begin{Verbatim}[commandchars=\\\{\}]
\PY{c+c1}{// Before}
\PY{n}{printAll}\PY{p}{(}\PY{k+kc}{null}\PY{p}{)}\PY{p}{;}\PY{+w}{ }\PY{c+c1}{// passes a null array reference}
\PY{c+c1}{// After}
\PY{n}{printAll}\PY{p}{(}\PY{p}{(}\PY{n}{String}\PY{p}{)}\PY{+w}{ }\PY{k+kc}{null}\PY{p}{)}\PY{p}{;}\PY{+w}{ }\PY{c+c1}{// passes a one\PYZhy{}element array: \PYZob{} null \PYZcb{}}
\end{Verbatim}
\end{codebox}

\item 
\textbf{Not checking for \code{Array\allowbreak{}Index\allowbreak{}Out\allowbreak{}Of\allowbreak{}Bounds\allowbreak{}Exception} risk on user-controlled indices.} Why it matters: array bounds are only checked at runtime, and code review should catch any index math derived from external input before it reaches the array access.

\begin{codebox}
\begin{Verbatim}[commandchars=\\\{\}]
\PY{c+c1}{// Before}
\PY{n}{process}\PY{p}{(}\PY{n}{items}\PY{o}{[}\PY{n}{requestedIndex}\PY{o}{]}\PY{p}{)}\PY{p}{;}\PY{+w}{ }\PY{c+c1}{// requestedIndex comes straight from a request param}
\PY{c+c1}{// After}
\PY{k}{if}\PY{+w}{ }\PY{p}{(}\PY{n}{requestedIndex}\PY{+w}{ }\PY{o}{\PYZgt{}}\PY{o}{=}\PY{+w}{ }\PY{l+m+mi}{0}\PY{+w}{ }\PY{o}{\PYZam{}}\PY{o}{\PYZam{}}\PY{+w}{ }\PY{n}{requestedIndex}\PY{+w}{ }\PY{o}{\PYZlt{}}\PY{+w}{ }\PY{n}{items}\PY{p}{.}\PY{n+na}{length}\PY{p}{)}\PY{+w}{ }\PY{n}{process}\PY{p}{(}\PY{n}{items}\PY{o}{[}\PY{n}{requestedIndex}\PY{o}{]}\PY{p}{)}\PY{p}{;}
\end{Verbatim}
\end{codebox}

\item 
\textbf{Using \code{StringBuffer} by default instead of \code{StringBuilder}.} Why it matters: \code{StringBuffer}’s synchronization overhead is wasted when the builder never leaves a single thread — which is the vast majority of real-world usage.

\begin{codebox}
\begin{Verbatim}[commandchars=\\\{\}]
\PY{c+c1}{// Before}
\PY{n}{StringBuffer}\PY{+w}{ }\PY{n}{sb}\PY{+w}{ }\PY{o}{=}\PY{+w}{ }\PY{k}{new}\PY{+w}{ }\PY{n}{StringBuffer}\PY{p}{(}\PY{p}{)}\PY{p}{;}
\PY{c+c1}{// After}
\PY{n}{StringBuilder}\PY{+w}{ }\PY{n}{sb}\PY{+w}{ }\PY{o}{=}\PY{+w}{ }\PY{k}{new}\PY{+w}{ }\PY{n}{StringBuilder}\PY{p}{(}\PY{p}{)}\PY{p}{;}\PY{+w}{ }\PY{c+c1}{// same API, no locking cost}
\end{Verbatim}
\end{codebox}

\item 
\textbf{Autoboxing wrapper types inside performance-critical loops.} Why it matters: each boxing operation allocates a short-lived object, adding GC pressure that’s easy to miss because the code looks like plain arithmetic.

\begin{codebox}
\begin{Verbatim}[commandchars=\\\{\}]
\PY{c+c1}{// Before}
\PY{n}{Long}\PY{+w}{ }\PY{n}{sum}\PY{+w}{ }\PY{o}{=}\PY{+w}{ }\PY{l+m+mi}{0}\PY{n}{L}\PY{p}{;}
\PY{k}{for}\PY{+w}{ }\PY{p}{(}\PY{k+kt}{int}\PY{+w}{ }\PY{n}{i}\PY{+w}{ }\PY{o}{=}\PY{+w}{ }\PY{l+m+mi}{0}\PY{p}{;}\PY{+w}{ }\PY{n}{i}\PY{+w}{ }\PY{o}{\PYZlt{}}\PY{+w}{ }\PY{l+m+mi}{1\PYZus{}000\PYZus{}000}\PY{p}{;}\PY{+w}{ }\PY{n}{i}\PY{o}{+}\PY{o}{+}\PY{p}{)}\PY{+w}{ }\PY{n}{sum}\PY{+w}{ }\PY{o}{+}\PY{o}{=}\PY{+w}{ }\PY{n}{i}\PY{p}{;}\PY{+w}{ }\PY{c+c1}{// boxes/unboxes every iteration}
\PY{c+c1}{// After}
\PY{k+kt}{long}\PY{+w}{ }\PY{n}{sum}\PY{+w}{ }\PY{o}{=}\PY{+w}{ }\PY{l+m+mi}{0}\PY{n}{L}\PY{p}{;}
\PY{k}{for}\PY{+w}{ }\PY{p}{(}\PY{k+kt}{int}\PY{+w}{ }\PY{n}{i}\PY{+w}{ }\PY{o}{=}\PY{+w}{ }\PY{l+m+mi}{0}\PY{p}{;}\PY{+w}{ }\PY{n}{i}\PY{+w}{ }\PY{o}{\PYZlt{}}\PY{+w}{ }\PY{l+m+mi}{1\PYZus{}000\PYZus{}000}\PY{p}{;}\PY{+w}{ }\PY{n}{i}\PY{o}{+}\PY{o}{+}\PY{p}{)}\PY{+w}{ }\PY{n}{sum}\PY{+w}{ }\PY{o}{+}\PY{o}{=}\PY{+w}{ }\PY{n}{i}\PY{p}{;}\PY{+w}{ }\PY{c+c1}{// stays primitive}
\end{Verbatim}
\end{codebox}

\item 
\textbf{Relying on operator precedence in dense expressions without parentheses.} Why it matters: mixing \code{\&\&}/\code{\textbar{}\textbar{}}, bitwise operators, or chained \code{++}/\code{--} without parentheses makes intent ambiguous to a reviewer, even when the compiler’s answer is well-defined.

\begin{codebox}
\begin{Verbatim}[commandchars=\\\{\}]
\PY{c+c1}{// Before}
\PY{k}{if}\PY{+w}{ }\PY{p}{(}\PY{n}{a}\PY{+w}{ }\PY{o}{|}\PY{o}{|}\PY{+w}{ }\PY{n}{b}\PY{+w}{ }\PY{o}{\PYZam{}}\PY{o}{\PYZam{}}\PY{+w}{ }\PY{n}{c}\PY{p}{)}\PY{+w}{ }\PY{p}{\PYZob{}}\PY{+w}{ }\PY{p}{.}\PY{p}{.}\PY{p}{.}\PY{+w}{ }\PY{p}{\PYZcb{}}\PY{+w}{ }\PY{c+c1}{// precedence is correct but not obvious to a reader}
\PY{c+c1}{// After}
\PY{k}{if}\PY{+w}{ }\PY{p}{(}\PY{n}{a}\PY{+w}{ }\PY{o}{|}\PY{o}{|}\PY{+w}{ }\PY{p}{(}\PY{n}{b}\PY{+w}{ }\PY{o}{\PYZam{}}\PY{o}{\PYZam{}}\PY{+w}{ }\PY{n}{c}\PY{p}{)}\PY{p}{)}\PY{+w}{ }\PY{p}{\PYZob{}}\PY{+w}{ }\PY{p}{.}\PY{p}{.}\PY{p}{.}\PY{+w}{ }\PY{p}{\PYZcb{}}
\end{Verbatim}
\end{codebox}

\item 
\textbf{Declaring \code{var} with the diamond operator and no type witness.} Why it matters: \code{var~\allowbreak{}list~\allowbreak{}=\allowbreak{}~\allowbreak{}new~\allowbreak{}Array\allowbreak{}List\textless{}\textgreater{}();} infers \code{ArrayList\textless{}\allowbreak{}Object\textgreater{}}, silently discarding the compile-time type safety generics are supposed to provide.

\begin{codebox}
\begin{Verbatim}[commandchars=\\\{\}]
\PY{c+c1}{// Before}
\PY{k+kd}{var}\PY{+w}{ }\PY{n}{names}\PY{+w}{ }\PY{o}{=}\PY{+w}{ }\PY{k}{new}\PY{+w}{ }\PY{n}{ArrayList}\PY{o}{\PYZlt{}}\PY{o}{\PYZgt{}}\PY{p}{(}\PY{p}{)}\PY{p}{;}\PY{+w}{ }\PY{c+c1}{// ArrayList\PYZlt{}Object\PYZgt{} — accepts anything}
\PY{c+c1}{// After}
\PY{k+kd}{var}\PY{+w}{ }\PY{n}{names}\PY{+w}{ }\PY{o}{=}\PY{+w}{ }\PY{k}{new}\PY{+w}{ }\PY{n}{ArrayList}\PY{o}{\PYZlt{}}\PY{n}{String}\PY{o}{\PYZgt{}}\PY{p}{(}\PY{p}{)}\PY{p}{;}\PY{+w}{ }\PY{c+c1}{// ArrayList\PYZlt{}String\PYZgt{} — type\PYZhy{}safe}
\end{Verbatim}
\end{codebox}

\end{enumerate}


\setcounter{chapter}{1}
\chapter{Object-Oriented Programming}
\label{chap:2}
\label{h:2:2-object-oriented-programming}
Object-Oriented Programming (OOP) is a way of organizing code around “objects” instead of just functions and data. Java is built on OOP from the ground up, so almost every code review question about design quality traces back to these basics. This chapter walks through the core OOP building blocks in Java 21+, with small, realistic examples and the “gotchas” that trip people up in interviews and in real pull requests.

\section[Classes and Objects]{Classes and Objects}
\label{h:2:classes-and-objects}
A \textbf{class} is a blueprint. It describes what data (fields) and behavior (methods) something has, but it is not itself a “thing” you can use directly. An \textbf{object} is an actual instance created from that blueprint, living in memory, with its own copy of the fields (unless they are \code{static}, see \hyperref[h:2:static-members]{Static Members}).

Think of \code{Account} as the blueprint for a bank account, and each customer’s actual account as an object built from that blueprint.

\begin{codebox}
\begin{Verbatim}[commandchars=\\\{\}]
\PY{k+kd}{public}\PY{+w}{ }\PY{k+kd}{class} \PY{n+nc}{Account}\PY{+w}{ }\PY{p}{\PYZob{}}
\PY{+w}{    }\PY{k+kd}{private}\PY{+w}{ }\PY{n}{String}\PY{+w}{ }\PY{n}{ownerName}\PY{p}{;}
\PY{+w}{    }\PY{k+kd}{private}\PY{+w}{ }\PY{k+kt}{double}\PY{+w}{ }\PY{n}{balance}\PY{p}{;}

\PY{+w}{    }\PY{k+kd}{public}\PY{+w}{ }\PY{n+nf}{Account}\PY{p}{(}\PY{n}{String}\PY{+w}{ }\PY{n}{ownerName}\PY{p}{,}\PY{+w}{ }\PY{k+kt}{double}\PY{+w}{ }\PY{n}{balance}\PY{p}{)}\PY{+w}{ }\PY{p}{\PYZob{}}
\PY{+w}{        }\PY{k}{this}\PY{p}{.}\PY{n+na}{ownerName}\PY{+w}{ }\PY{o}{=}\PY{+w}{ }\PY{n}{ownerName}\PY{p}{;}
\PY{+w}{        }\PY{k}{this}\PY{p}{.}\PY{n+na}{balance}\PY{+w}{ }\PY{o}{=}\PY{+w}{ }\PY{n}{balance}\PY{p}{;}
\PY{+w}{    }\PY{p}{\PYZcb{}}

\PY{+w}{    }\PY{k+kd}{public}\PY{+w}{ }\PY{k+kt}{void}\PY{+w}{ }\PY{n+nf}{deposit}\PY{p}{(}\PY{k+kt}{double}\PY{+w}{ }\PY{n}{amount}\PY{p}{)}\PY{+w}{ }\PY{p}{\PYZob{}}
\PY{+w}{        }\PY{n}{balance}\PY{+w}{ }\PY{o}{+}\PY{o}{=}\PY{+w}{ }\PY{n}{amount}\PY{p}{;}
\PY{+w}{    }\PY{p}{\PYZcb{}}

\PY{+w}{    }\PY{k+kd}{public}\PY{+w}{ }\PY{k+kt}{double}\PY{+w}{ }\PY{n+nf}{getBalance}\PY{p}{(}\PY{p}{)}\PY{+w}{ }\PY{p}{\PYZob{}}
\PY{+w}{        }\PY{k}{return}\PY{+w}{ }\PY{n}{balance}\PY{p}{;}
\PY{+w}{    }\PY{p}{\PYZcb{}}
\PY{p}{\PYZcb{}}

\PY{k+kd}{public}\PY{+w}{ }\PY{k+kd}{class} \PY{n+nc}{Bank}\PY{+w}{ }\PY{p}{\PYZob{}}
\PY{+w}{    }\PY{k+kd}{public}\PY{+w}{ }\PY{k+kd}{static}\PY{+w}{ }\PY{k+kt}{void}\PY{+w}{ }\PY{n+nf}{main}\PY{p}{(}\PY{n}{String}\PY{o}{[}\PY{o}{]}\PY{+w}{ }\PY{n}{args}\PY{p}{)}\PY{+w}{ }\PY{p}{\PYZob{}}
\PY{+w}{        }\PY{n}{Account}\PY{+w}{ }\PY{n}{aliceAccount}\PY{+w}{ }\PY{o}{=}\PY{+w}{ }\PY{k}{new}\PY{+w}{ }\PY{n}{Account}\PY{p}{(}\PY{l+s}{\PYZdq{}}\PY{l+s}{Alice}\PY{l+s}{\PYZdq{}}\PY{p}{,}\PY{+w}{ }\PY{l+m+mf}{100.0}\PY{p}{)}\PY{p}{;}
\PY{+w}{        }\PY{n}{Account}\PY{+w}{ }\PY{n}{bobAccount}\PY{+w}{ }\PY{o}{=}\PY{+w}{ }\PY{k}{new}\PY{+w}{ }\PY{n}{Account}\PY{p}{(}\PY{l+s}{\PYZdq{}}\PY{l+s}{Bob}\PY{l+s}{\PYZdq{}}\PY{p}{,}\PY{+w}{ }\PY{l+m+mf}{50.0}\PY{p}{)}\PY{p}{;}

\PY{+w}{        }\PY{n}{aliceAccount}\PY{p}{.}\PY{n+na}{deposit}\PY{p}{(}\PY{l+m+mf}{25.0}\PY{p}{)}\PY{p}{;}

\PY{+w}{        }\PY{n}{System}\PY{p}{.}\PY{n+na}{out}\PY{p}{.}\PY{n+na}{println}\PY{p}{(}\PY{n}{aliceAccount}\PY{p}{.}\PY{n+na}{getBalance}\PY{p}{(}\PY{p}{)}\PY{p}{)}\PY{p}{;}\PY{+w}{ }\PY{c+c1}{// 125.0}
\PY{+w}{        }\PY{n}{System}\PY{p}{.}\PY{n+na}{out}\PY{p}{.}\PY{n+na}{println}\PY{p}{(}\PY{n}{bobAccount}\PY{p}{.}\PY{n+na}{getBalance}\PY{p}{(}\PY{p}{)}\PY{p}{)}\PY{p}{;}\PY{+w}{   }\PY{c+c1}{// 50.0 (independent object)}
\PY{+w}{    }\PY{p}{\PYZcb{}}
\PY{p}{\PYZcb{}}
\end{Verbatim}
\end{codebox}

Each \code{new~\allowbreak{}Account(...)} call creates a separate object with its own \code{balance}. Changing \code{aliceAccount} never affects \code{bobAccount}. That independence is the whole point of objects: state is bundled with the behavior that operates on it.

\textbf{Gotcha: object identity vs. object equality.} Two different objects can have the same data but are still not “the same object” in memory.

\begin{codebox}
\begin{Verbatim}[commandchars=\\\{\}]
\PY{n}{Account}\PY{+w}{ }\PY{n}{a1}\PY{+w}{ }\PY{o}{=}\PY{+w}{ }\PY{k}{new}\PY{+w}{ }\PY{n}{Account}\PY{p}{(}\PY{l+s}{\PYZdq{}}\PY{l+s}{Alice}\PY{l+s}{\PYZdq{}}\PY{p}{,}\PY{+w}{ }\PY{l+m+mf}{100.0}\PY{p}{)}\PY{p}{;}
\PY{n}{Account}\PY{+w}{ }\PY{n}{a2}\PY{+w}{ }\PY{o}{=}\PY{+w}{ }\PY{k}{new}\PY{+w}{ }\PY{n}{Account}\PY{p}{(}\PY{l+s}{\PYZdq{}}\PY{l+s}{Alice}\PY{l+s}{\PYZdq{}}\PY{p}{,}\PY{+w}{ }\PY{l+m+mf}{100.0}\PY{p}{)}\PY{p}{;}

\PY{n}{System}\PY{p}{.}\PY{n+na}{out}\PY{p}{.}\PY{n+na}{println}\PY{p}{(}\PY{n}{a1}\PY{+w}{ }\PY{o}{=}\PY{o}{=}\PY{+w}{ }\PY{n}{a2}\PY{p}{)}\PY{p}{;}\PY{+w}{       }\PY{c+c1}{// false: different objects in memory}
\PY{n}{System}\PY{p}{.}\PY{n+na}{out}\PY{p}{.}\PY{n+na}{println}\PY{p}{(}\PY{n}{a1}\PY{p}{.}\PY{n+na}{equals}\PY{p}{(}\PY{n}{a2}\PY{p}{)}\PY{p}{)}\PY{p}{;}\PY{+w}{  }\PY{c+c1}{// true only if equals() is overridden;}
\PY{+w}{                                     }\PY{c+c1}{// otherwise also false (default Object.equals uses ==)}
\end{Verbatim}
\end{codebox}

\code{==} compares references (memory addresses), not content. Unless a class overrides \code{equals()}, \code{.\allowbreak{}equals()} behaves exactly like \code{==}. This is covered more in the “Object Class \& Common APIs” chapter, but it starts here with basic object identity.

\section[Constructors]{Constructors}
\label{h:2:constructors}
A \textbf{constructor} is a special method that runs when an object is created with \code{new}. It has the same name as the class, no return type (not even \code{void}), and its job is to put the object into a valid starting state.

\begin{codebox}
\begin{Verbatim}[commandchars=\\\{\}]
\PY{k+kd}{public}\PY{+w}{ }\PY{k+kd}{class} \PY{n+nc}{Order}\PY{+w}{ }\PY{p}{\PYZob{}}
\PY{+w}{    }\PY{k+kd}{private}\PY{+w}{ }\PY{k+kd}{final}\PY{+w}{ }\PY{n}{String}\PY{+w}{ }\PY{n}{orderId}\PY{p}{;}
\PY{+w}{    }\PY{k+kd}{private}\PY{+w}{ }\PY{k+kd}{final}\PY{+w}{ }\PY{n}{List}\PY{o}{\PYZlt{}}\PY{n}{String}\PY{o}{\PYZgt{}}\PY{+w}{ }\PY{n}{items}\PY{p}{;}
\PY{+w}{    }\PY{k+kd}{private}\PY{+w}{ }\PY{k+kt}{double}\PY{+w}{ }\PY{n}{total}\PY{p}{;}

\PY{+w}{    }\PY{c+c1}{// No\PYZhy{}args constructor}
\PY{+w}{    }\PY{k+kd}{public}\PY{+w}{ }\PY{n+nf}{Order}\PY{p}{(}\PY{p}{)}\PY{+w}{ }\PY{p}{\PYZob{}}
\PY{+w}{        }\PY{k}{this}\PY{p}{(}\PY{l+s}{\PYZdq{}}\PY{l+s}{ORD\PYZhy{}0000}\PY{l+s}{\PYZdq{}}\PY{p}{)}\PY{p}{;}\PY{+w}{ }\PY{c+c1}{// calls the other constructor below}
\PY{+w}{    }\PY{p}{\PYZcb{}}

\PY{+w}{    }\PY{c+c1}{// Constructor with a parameter}
\PY{+w}{    }\PY{k+kd}{public}\PY{+w}{ }\PY{n+nf}{Order}\PY{p}{(}\PY{n}{String}\PY{+w}{ }\PY{n}{orderId}\PY{p}{)}\PY{+w}{ }\PY{p}{\PYZob{}}
\PY{+w}{        }\PY{k}{this}\PY{p}{.}\PY{n+na}{orderId}\PY{+w}{ }\PY{o}{=}\PY{+w}{ }\PY{n}{orderId}\PY{p}{;}
\PY{+w}{        }\PY{k}{this}\PY{p}{.}\PY{n+na}{items}\PY{+w}{ }\PY{o}{=}\PY{+w}{ }\PY{k}{new}\PY{+w}{ }\PY{n}{ArrayList}\PY{o}{\PYZlt{}}\PY{o}{\PYZgt{}}\PY{p}{(}\PY{p}{)}\PY{p}{;}
\PY{+w}{        }\PY{k}{this}\PY{p}{.}\PY{n+na}{total}\PY{+w}{ }\PY{o}{=}\PY{+w}{ }\PY{l+m+mf}{0.0}\PY{p}{;}
\PY{+w}{    }\PY{p}{\PYZcb{}}
\PY{p}{\PYZcb{}}
\end{Verbatim}
\end{codebox}

If you don’t write any constructor, Java gives you a free \textbf{default constructor} with no arguments. The moment you write any constructor yourself, that free one disappears.

\begin{codebox}
\begin{Verbatim}[commandchars=\\\{\}]
\PY{k+kd}{public}\PY{+w}{ }\PY{k+kd}{class} \PY{n+nc}{Shape}\PY{+w}{ }\PY{p}{\PYZob{}}
\PY{+w}{    }\PY{k+kd}{private}\PY{+w}{ }\PY{n}{String}\PY{+w}{ }\PY{n}{color}\PY{p}{;}
\PY{+w}{    }\PY{c+c1}{// No constructor written here.}
\PY{p}{\PYZcb{}}

\PY{c+c1}{// Elsewhere:}
\PY{n}{Shape}\PY{+w}{ }\PY{n}{s}\PY{+w}{ }\PY{o}{=}\PY{+w}{ }\PY{k}{new}\PY{+w}{ }\PY{n}{Shape}\PY{p}{(}\PY{p}{)}\PY{p}{;}\PY{+w}{ }\PY{c+c1}{// Works: compiler\PYZhy{}generated default constructor}
\end{Verbatim}
\end{codebox}

\begin{codebox}
\begin{Verbatim}[commandchars=\\\{\}]
\PY{k+kd}{public}\PY{+w}{ }\PY{k+kd}{class} \PY{n+nc}{Shape}\PY{+w}{ }\PY{p}{\PYZob{}}
\PY{+w}{    }\PY{k+kd}{private}\PY{+w}{ }\PY{n}{String}\PY{+w}{ }\PY{n}{color}\PY{p}{;}

\PY{+w}{    }\PY{k+kd}{public}\PY{+w}{ }\PY{n+nf}{Shape}\PY{p}{(}\PY{n}{String}\PY{+w}{ }\PY{n}{color}\PY{p}{)}\PY{+w}{ }\PY{p}{\PYZob{}}
\PY{+w}{        }\PY{k}{this}\PY{p}{.}\PY{n+na}{color}\PY{+w}{ }\PY{o}{=}\PY{+w}{ }\PY{n}{color}\PY{p}{;}
\PY{+w}{    }\PY{p}{\PYZcb{}}
\PY{+w}{    }\PY{c+c1}{// No no\PYZhy{}args constructor written, and the default one is gone now.}
\PY{p}{\PYZcb{}}

\PY{c+c1}{// Elsewhere:}
\PY{n}{Shape}\PY{+w}{ }\PY{n}{s}\PY{+w}{ }\PY{o}{=}\PY{+w}{ }\PY{k}{new}\PY{+w}{ }\PY{n}{Shape}\PY{p}{(}\PY{p}{)}\PY{p}{;}\PY{+w}{ }\PY{c+c1}{// Compile error: no matching constructor}
\end{Verbatim}
\end{codebox}

\textbf{Constructor chaining} means one constructor calls another using \code{this(...)}, so shared setup logic lives in one place. A constructor can also call a parent class constructor using \code{super(...)}. If you don’t call \code{super(...)} explicitly, Java inserts an implicit no-argument \code{super()} call as the very first line.

\begin{codebox}
\begin{Verbatim}[commandchars=\\\{\}]
\PY{k+kd}{public}\PY{+w}{ }\PY{k+kd}{class} \PY{n+nc}{Vehicle}\PY{+w}{ }\PY{p}{\PYZob{}}
\PY{+w}{    }\PY{k+kd}{protected}\PY{+w}{ }\PY{n}{String}\PY{+w}{ }\PY{n}{plateNumber}\PY{p}{;}

\PY{+w}{    }\PY{k+kd}{public}\PY{+w}{ }\PY{n+nf}{Vehicle}\PY{p}{(}\PY{n}{String}\PY{+w}{ }\PY{n}{plateNumber}\PY{p}{)}\PY{+w}{ }\PY{p}{\PYZob{}}
\PY{+w}{        }\PY{k}{this}\PY{p}{.}\PY{n+na}{plateNumber}\PY{+w}{ }\PY{o}{=}\PY{+w}{ }\PY{n}{plateNumber}\PY{p}{;}
\PY{+w}{        }\PY{n}{System}\PY{p}{.}\PY{n+na}{out}\PY{p}{.}\PY{n+na}{println}\PY{p}{(}\PY{l+s}{\PYZdq{}}\PY{l+s}{Vehicle created: }\PY{l+s}{\PYZdq{}}\PY{+w}{ }\PY{o}{+}\PY{+w}{ }\PY{n}{plateNumber}\PY{p}{)}\PY{p}{;}
\PY{+w}{    }\PY{p}{\PYZcb{}}
\PY{p}{\PYZcb{}}

\PY{k+kd}{public}\PY{+w}{ }\PY{k+kd}{class} \PY{n+nc}{Car}\PY{+w}{ }\PY{k+kd}{extends}\PY{+w}{ }\PY{n}{Vehicle}\PY{+w}{ }\PY{p}{\PYZob{}}
\PY{+w}{    }\PY{k+kd}{private}\PY{+w}{ }\PY{k+kt}{int}\PY{+w}{ }\PY{n}{doors}\PY{p}{;}

\PY{+w}{    }\PY{k+kd}{public}\PY{+w}{ }\PY{n+nf}{Car}\PY{p}{(}\PY{n}{String}\PY{+w}{ }\PY{n}{plateNumber}\PY{p}{,}\PY{+w}{ }\PY{k+kt}{int}\PY{+w}{ }\PY{n}{doors}\PY{p}{)}\PY{+w}{ }\PY{p}{\PYZob{}}
\PY{+w}{        }\PY{k+kd}{super}\PY{p}{(}\PY{n}{plateNumber}\PY{p}{)}\PY{p}{;}\PY{+w}{ }\PY{c+c1}{// must be the first statement}
\PY{+w}{        }\PY{k}{this}\PY{p}{.}\PY{n+na}{doors}\PY{+w}{ }\PY{o}{=}\PY{+w}{ }\PY{n}{doors}\PY{p}{;}
\PY{+w}{        }\PY{n}{System}\PY{p}{.}\PY{n+na}{out}\PY{p}{.}\PY{n+na}{println}\PY{p}{(}\PY{l+s}{\PYZdq{}}\PY{l+s}{Car created with }\PY{l+s}{\PYZdq{}}\PY{+w}{ }\PY{o}{+}\PY{+w}{ }\PY{n}{doors}\PY{+w}{ }\PY{o}{+}\PY{+w}{ }\PY{l+s}{\PYZdq{}}\PY{l+s}{ doors}\PY{l+s}{\PYZdq{}}\PY{p}{)}\PY{p}{;}
\PY{+w}{    }\PY{p}{\PYZcb{}}
\PY{p}{\PYZcb{}}

\PY{c+c1}{// new Car(\PYZdq{}XYZ\PYZhy{}123\PYZdq{}, 4);}
\PY{c+c1}{// Output:}
\PY{c+c1}{// Vehicle created: XYZ\PYZhy{}123}
\PY{c+c1}{// Car created with 4 doors}
\end{Verbatim}
\end{codebox}

\textbf{Gotcha: calling an overridable method from a constructor.} If a constructor calls a method that a subclass overrides, the subclass version can run \emph{before} the subclass has finished initializing its own fields.

\begin{codebox}
\begin{Verbatim}[commandchars=\\\{\}]
\PY{k+kd}{public}\PY{+w}{ }\PY{k+kd}{class} \PY{n+nc}{Shape}\PY{+w}{ }\PY{p}{\PYZob{}}
\PY{+w}{    }\PY{k+kd}{public}\PY{+w}{ }\PY{n+nf}{Shape}\PY{p}{(}\PY{p}{)}\PY{+w}{ }\PY{p}{\PYZob{}}
\PY{+w}{        }\PY{n}{draw}\PY{p}{(}\PY{p}{)}\PY{p}{;}\PY{+w}{ }\PY{c+c1}{// calling an overridable method during construction}
\PY{+w}{    }\PY{p}{\PYZcb{}}

\PY{+w}{    }\PY{k+kd}{public}\PY{+w}{ }\PY{k+kt}{void}\PY{+w}{ }\PY{n+nf}{draw}\PY{p}{(}\PY{p}{)}\PY{+w}{ }\PY{p}{\PYZob{}}
\PY{+w}{        }\PY{n}{System}\PY{p}{.}\PY{n+na}{out}\PY{p}{.}\PY{n+na}{println}\PY{p}{(}\PY{l+s}{\PYZdq{}}\PY{l+s}{Drawing a generic shape}\PY{l+s}{\PYZdq{}}\PY{p}{)}\PY{p}{;}
\PY{+w}{    }\PY{p}{\PYZcb{}}
\PY{p}{\PYZcb{}}

\PY{k+kd}{public}\PY{+w}{ }\PY{k+kd}{class} \PY{n+nc}{Circle}\PY{+w}{ }\PY{k+kd}{extends}\PY{+w}{ }\PY{n}{Shape}\PY{+w}{ }\PY{p}{\PYZob{}}
\PY{+w}{    }\PY{k+kd}{private}\PY{+w}{ }\PY{k+kt}{int}\PY{+w}{ }\PY{n}{radius}\PY{+w}{ }\PY{o}{=}\PY{+w}{ }\PY{l+m+mi}{5}\PY{p}{;}

\PY{+w}{    }\PY{n+nd}{@Override}
\PY{+w}{    }\PY{k+kd}{public}\PY{+w}{ }\PY{k+kt}{void}\PY{+w}{ }\PY{n+nf}{draw}\PY{p}{(}\PY{p}{)}\PY{+w}{ }\PY{p}{\PYZob{}}
\PY{+w}{        }\PY{c+c1}{// radius may not be assigned yet when this runs from Shape\PYZsq{}s constructor!}
\PY{+w}{        }\PY{n}{System}\PY{p}{.}\PY{n+na}{out}\PY{p}{.}\PY{n+na}{println}\PY{p}{(}\PY{l+s}{\PYZdq{}}\PY{l+s}{Drawing circle with radius }\PY{l+s}{\PYZdq{}}\PY{+w}{ }\PY{o}{+}\PY{+w}{ }\PY{n}{radius}\PY{p}{)}\PY{p}{;}
\PY{+w}{    }\PY{p}{\PYZcb{}}
\PY{p}{\PYZcb{}}

\PY{c+c1}{// new Circle();}
\PY{c+c1}{// Output: Drawing circle with radius 0}
\PY{c+c1}{// (Shape\PYZsq{}s constructor runs first and calls draw(), but Circle\PYZsq{}s field}
\PY{c+c1}{//  initializer \PYZdq{}radius = 5\PYZdq{} hasn\PYZsq{}t executed yet.)}
\end{Verbatim}
\end{codebox}

The fix: avoid calling overridable (non-\code{private}, non-\code{final}, non-\code{static}) methods from a constructor, or mark the method \code{final} if it must be called during construction.

\section[Encapsulation]{Encapsulation}
\label{h:2:encapsulation}
\textbf{Encapsulation} means hiding the internal state of an object and only exposing it through controlled methods. Fields are usually \code{private}, and public \textbf{getters/setters} (or no setters at all, for immutable data) control access. This protects invariants — rules that must always be true, like “balance can never go negative.”

\begin{codebox}
\begin{Verbatim}[commandchars=\\\{\}]
\PY{k+kd}{public}\PY{+w}{ }\PY{k+kd}{class} \PY{n+nc}{Account}\PY{+w}{ }\PY{p}{\PYZob{}}
\PY{+w}{    }\PY{k+kd}{private}\PY{+w}{ }\PY{k+kt}{double}\PY{+w}{ }\PY{n}{balance}\PY{p}{;}

\PY{+w}{    }\PY{k+kd}{public}\PY{+w}{ }\PY{n+nf}{Account}\PY{p}{(}\PY{k+kt}{double}\PY{+w}{ }\PY{n}{initialBalance}\PY{p}{)}\PY{+w}{ }\PY{p}{\PYZob{}}
\PY{+w}{        }\PY{k}{if}\PY{+w}{ }\PY{p}{(}\PY{n}{initialBalance}\PY{+w}{ }\PY{o}{\PYZlt{}}\PY{+w}{ }\PY{l+m+mi}{0}\PY{p}{)}\PY{+w}{ }\PY{p}{\PYZob{}}
\PY{+w}{            }\PY{k}{throw}\PY{+w}{ }\PY{k}{new}\PY{+w}{ }\PY{n}{IllegalArgumentException}\PY{p}{(}\PY{l+s}{\PYZdq{}}\PY{l+s}{Initial balance cannot be negative}\PY{l+s}{\PYZdq{}}\PY{p}{)}\PY{p}{;}
\PY{+w}{        }\PY{p}{\PYZcb{}}
\PY{+w}{        }\PY{k}{this}\PY{p}{.}\PY{n+na}{balance}\PY{+w}{ }\PY{o}{=}\PY{+w}{ }\PY{n}{initialBalance}\PY{p}{;}
\PY{+w}{    }\PY{p}{\PYZcb{}}

\PY{+w}{    }\PY{k+kd}{public}\PY{+w}{ }\PY{k+kt}{double}\PY{+w}{ }\PY{n+nf}{getBalance}\PY{p}{(}\PY{p}{)}\PY{+w}{ }\PY{p}{\PYZob{}}
\PY{+w}{        }\PY{k}{return}\PY{+w}{ }\PY{n}{balance}\PY{p}{;}
\PY{+w}{    }\PY{p}{\PYZcb{}}

\PY{+w}{    }\PY{k+kd}{public}\PY{+w}{ }\PY{k+kt}{void}\PY{+w}{ }\PY{n+nf}{withdraw}\PY{p}{(}\PY{k+kt}{double}\PY{+w}{ }\PY{n}{amount}\PY{p}{)}\PY{+w}{ }\PY{p}{\PYZob{}}
\PY{+w}{        }\PY{k}{if}\PY{+w}{ }\PY{p}{(}\PY{n}{amount}\PY{+w}{ }\PY{o}{\PYZgt{}}\PY{+w}{ }\PY{n}{balance}\PY{p}{)}\PY{+w}{ }\PY{p}{\PYZob{}}
\PY{+w}{            }\PY{k}{throw}\PY{+w}{ }\PY{k}{new}\PY{+w}{ }\PY{n}{IllegalStateException}\PY{p}{(}\PY{l+s}{\PYZdq{}}\PY{l+s}{Insufficient funds}\PY{l+s}{\PYZdq{}}\PY{p}{)}\PY{p}{;}
\PY{+w}{        }\PY{p}{\PYZcb{}}
\PY{+w}{        }\PY{n}{balance}\PY{+w}{ }\PY{o}{\PYZhy{}}\PY{o}{=}\PY{+w}{ }\PY{n}{amount}\PY{p}{;}
\PY{+w}{    }\PY{p}{\PYZcb{}}
\PY{p}{\PYZcb{}}
\end{Verbatim}
\end{codebox}

Without encapsulation, any code could do \code{account.\allowbreak{}balance~\allowbreak{}=\allowbreak{}~\allowbreak{}-\allowbreak{}500;} directly and break the invariant. With a \code{private} field, the only way to change \code{balance} is through \code{withdraw()}, which enforces the rule.

\textbf{Gotcha: a getter that leaks a mutable field breaks encapsulation even though the field is \code{private}.}

\begin{codebox}
\begin{Verbatim}[commandchars=\\\{\}]
\PY{k+kd}{public}\PY{+w}{ }\PY{k+kd}{class} \PY{n+nc}{Order}\PY{+w}{ }\PY{p}{\PYZob{}}
\PY{+w}{    }\PY{k+kd}{private}\PY{+w}{ }\PY{k+kd}{final}\PY{+w}{ }\PY{n}{List}\PY{o}{\PYZlt{}}\PY{n}{String}\PY{o}{\PYZgt{}}\PY{+w}{ }\PY{n}{items}\PY{+w}{ }\PY{o}{=}\PY{+w}{ }\PY{k}{new}\PY{+w}{ }\PY{n}{ArrayList}\PY{o}{\PYZlt{}}\PY{o}{\PYZgt{}}\PY{p}{(}\PY{p}{)}\PY{p}{;}

\PY{+w}{    }\PY{k+kd}{public}\PY{+w}{ }\PY{n}{List}\PY{o}{\PYZlt{}}\PY{n}{String}\PY{o}{\PYZgt{}}\PY{+w}{ }\PY{n+nf}{getItems}\PY{p}{(}\PY{p}{)}\PY{+w}{ }\PY{p}{\PYZob{}}
\PY{+w}{        }\PY{k}{return}\PY{+w}{ }\PY{n}{items}\PY{p}{;}\PY{+w}{ }\PY{c+c1}{// returns the real internal list!}
\PY{+w}{    }\PY{p}{\PYZcb{}}
\PY{p}{\PYZcb{}}

\PY{n}{Order}\PY{+w}{ }\PY{n}{order}\PY{+w}{ }\PY{o}{=}\PY{+w}{ }\PY{k}{new}\PY{+w}{ }\PY{n}{Order}\PY{p}{(}\PY{p}{)}\PY{p}{;}
\PY{n}{order}\PY{p}{.}\PY{n+na}{getItems}\PY{p}{(}\PY{p}{)}\PY{p}{.}\PY{n+na}{add}\PY{p}{(}\PY{l+s}{\PYZdq{}}\PY{l+s}{Hacked item}\PY{l+s}{\PYZdq{}}\PY{p}{)}\PY{p}{;}\PY{+w}{ }\PY{c+c1}{// mutates internal state from outside}
\end{Verbatim}
\end{codebox}

The fix is a \textbf{defensive copy}, or returning an unmodifiable view:

\begin{codebox}
\begin{Verbatim}[commandchars=\\\{\}]
\PY{k+kd}{public}\PY{+w}{ }\PY{n}{List}\PY{o}{\PYZlt{}}\PY{n}{String}\PY{o}{\PYZgt{}}\PY{+w}{ }\PY{n+nf}{getItems}\PY{p}{(}\PY{p}{)}\PY{+w}{ }\PY{p}{\PYZob{}}
\PY{+w}{    }\PY{k}{return}\PY{+w}{ }\PY{n}{List}\PY{p}{.}\PY{n+na}{copyOf}\PY{p}{(}\PY{n}{items}\PY{p}{)}\PY{p}{;}\PY{+w}{ }\PY{c+c1}{// caller gets a snapshot, can\PYZsq{}t mutate internals}
\PY{p}{\PYZcb{}}
\end{Verbatim}
\end{codebox}

\section[Inheritance]{Inheritance}
\label{h:2:inheritance}
\textbf{Inheritance} lets one class (the subclass) reuse and extend the fields and methods of another class (the superclass), using the \code{extends} keyword. It models an “is-a” relationship: a \code{Car} \textbf{is a} \code{Vehicle}.

\begin{codebox}
\begin{Verbatim}[commandchars=\\\{\}]
\PY{k+kd}{public}\PY{+w}{ }\PY{k+kd}{class} \PY{n+nc}{Vehicle}\PY{+w}{ }\PY{p}{\PYZob{}}
\PY{+w}{    }\PY{k+kd}{protected}\PY{+w}{ }\PY{n}{String}\PY{+w}{ }\PY{n}{plateNumber}\PY{p}{;}
\PY{+w}{    }\PY{k+kd}{protected}\PY{+w}{ }\PY{k+kt}{int}\PY{+w}{ }\PY{n}{speed}\PY{p}{;}

\PY{+w}{    }\PY{k+kd}{public}\PY{+w}{ }\PY{n+nf}{Vehicle}\PY{p}{(}\PY{n}{String}\PY{+w}{ }\PY{n}{plateNumber}\PY{p}{)}\PY{+w}{ }\PY{p}{\PYZob{}}
\PY{+w}{        }\PY{k}{this}\PY{p}{.}\PY{n+na}{plateNumber}\PY{+w}{ }\PY{o}{=}\PY{+w}{ }\PY{n}{plateNumber}\PY{p}{;}
\PY{+w}{    }\PY{p}{\PYZcb{}}

\PY{+w}{    }\PY{k+kd}{public}\PY{+w}{ }\PY{k+kt}{void}\PY{+w}{ }\PY{n+nf}{accelerate}\PY{p}{(}\PY{k+kt}{int}\PY{+w}{ }\PY{n}{amount}\PY{p}{)}\PY{+w}{ }\PY{p}{\PYZob{}}
\PY{+w}{        }\PY{n}{speed}\PY{+w}{ }\PY{o}{+}\PY{o}{=}\PY{+w}{ }\PY{n}{amount}\PY{p}{;}
\PY{+w}{    }\PY{p}{\PYZcb{}}

\PY{+w}{    }\PY{k+kd}{public}\PY{+w}{ }\PY{n}{String}\PY{+w}{ }\PY{n+nf}{describe}\PY{p}{(}\PY{p}{)}\PY{+w}{ }\PY{p}{\PYZob{}}
\PY{+w}{        }\PY{k}{return}\PY{+w}{ }\PY{l+s}{\PYZdq{}}\PY{l+s}{Vehicle }\PY{l+s}{\PYZdq{}}\PY{+w}{ }\PY{o}{+}\PY{+w}{ }\PY{n}{plateNumber}\PY{+w}{ }\PY{o}{+}\PY{+w}{ }\PY{l+s}{\PYZdq{}}\PY{l+s}{ at speed }\PY{l+s}{\PYZdq{}}\PY{+w}{ }\PY{o}{+}\PY{+w}{ }\PY{n}{speed}\PY{p}{;}
\PY{+w}{    }\PY{p}{\PYZcb{}}
\PY{p}{\PYZcb{}}

\PY{k+kd}{public}\PY{+w}{ }\PY{k+kd}{class} \PY{n+nc}{Car}\PY{+w}{ }\PY{k+kd}{extends}\PY{+w}{ }\PY{n}{Vehicle}\PY{+w}{ }\PY{p}{\PYZob{}}
\PY{+w}{    }\PY{k+kd}{private}\PY{+w}{ }\PY{k+kt}{int}\PY{+w}{ }\PY{n}{doors}\PY{p}{;}

\PY{+w}{    }\PY{k+kd}{public}\PY{+w}{ }\PY{n+nf}{Car}\PY{p}{(}\PY{n}{String}\PY{+w}{ }\PY{n}{plateNumber}\PY{p}{,}\PY{+w}{ }\PY{k+kt}{int}\PY{+w}{ }\PY{n}{doors}\PY{p}{)}\PY{+w}{ }\PY{p}{\PYZob{}}
\PY{+w}{        }\PY{k+kd}{super}\PY{p}{(}\PY{n}{plateNumber}\PY{p}{)}\PY{p}{;}
\PY{+w}{        }\PY{k}{this}\PY{p}{.}\PY{n+na}{doors}\PY{+w}{ }\PY{o}{=}\PY{+w}{ }\PY{n}{doors}\PY{p}{;}
\PY{+w}{    }\PY{p}{\PYZcb{}}

\PY{+w}{    }\PY{c+c1}{// Inherits accelerate() and describe() for free}
\PY{p}{\PYZcb{}}

\PY{n}{Car}\PY{+w}{ }\PY{n}{car}\PY{+w}{ }\PY{o}{=}\PY{+w}{ }\PY{k}{new}\PY{+w}{ }\PY{n}{Car}\PY{p}{(}\PY{l+s}{\PYZdq{}}\PY{l+s}{XYZ\PYZhy{}123}\PY{l+s}{\PYZdq{}}\PY{p}{,}\PY{+w}{ }\PY{l+m+mi}{4}\PY{p}{)}\PY{p}{;}
\PY{n}{car}\PY{p}{.}\PY{n+na}{accelerate}\PY{p}{(}\PY{l+m+mi}{30}\PY{p}{)}\PY{p}{;}
\PY{n}{System}\PY{p}{.}\PY{n+na}{out}\PY{p}{.}\PY{n+na}{println}\PY{p}{(}\PY{n}{car}\PY{p}{.}\PY{n+na}{describe}\PY{p}{(}\PY{p}{)}\PY{p}{)}\PY{p}{;}\PY{+w}{ }\PY{c+c1}{// Vehicle XYZ\PYZhy{}123 at speed 30}
\end{Verbatim}
\end{codebox}

Java only allows \textbf{single inheritance} for classes: a class can \code{extends} only one other class. This avoids the “diamond problem” (ambiguity when two parent classes define the same method). Java does allow a class to implement multiple interfaces, since interfaces (mostly) don’t carry conflicting state.

\begin{codebox}
\begin{Verbatim}[commandchars=\\\{\}]
\PY{k+kd}{public}\PY{+w}{ }\PY{k+kd}{class} \PY{n+nc}{Car}\PY{+w}{ }\PY{k+kd}{extends}\PY{+w}{ }\PY{n}{Vehicle}\PY{+w}{ }\PY{k+kd}{implements}\PY{+w}{ }\PY{n}{Insurable}\PY{p}{,}\PY{+w}{ }\PY{n}{Sellable}\PY{+w}{ }\PY{p}{\PYZob{}}
\PY{+w}{    }\PY{c+c1}{// one superclass, multiple interfaces — this is fine}
\PY{p}{\PYZcb{}}
\end{Verbatim}
\end{codebox}

\textbf{Gotcha: field hiding.} Unlike methods, fields in Java are not polymorphic. If a subclass declares a field with the same name as a superclass field, both fields exist — which one you see depends on the \emph{declared type} of the reference, not the actual object type.

\begin{codebox}
\begin{Verbatim}[commandchars=\\\{\}]
\PY{k+kd}{public}\PY{+w}{ }\PY{k+kd}{class} \PY{n+nc}{Vehicle}\PY{+w}{ }\PY{p}{\PYZob{}}
\PY{+w}{    }\PY{k+kd}{protected}\PY{+w}{ }\PY{k+kt}{int}\PY{+w}{ }\PY{n}{speed}\PY{+w}{ }\PY{o}{=}\PY{+w}{ }\PY{l+m+mi}{10}\PY{p}{;}
\PY{p}{\PYZcb{}}

\PY{k+kd}{public}\PY{+w}{ }\PY{k+kd}{class} \PY{n+nc}{Car}\PY{+w}{ }\PY{k+kd}{extends}\PY{+w}{ }\PY{n}{Vehicle}\PY{+w}{ }\PY{p}{\PYZob{}}
\PY{+w}{    }\PY{k+kd}{protected}\PY{+w}{ }\PY{k+kt}{int}\PY{+w}{ }\PY{n}{speed}\PY{+w}{ }\PY{o}{=}\PY{+w}{ }\PY{l+m+mi}{20}\PY{p}{;}\PY{+w}{ }\PY{c+c1}{// hides Vehicle\PYZsq{}s speed, does NOT override it}
\PY{p}{\PYZcb{}}

\PY{n}{Vehicle}\PY{+w}{ }\PY{n}{v}\PY{+w}{ }\PY{o}{=}\PY{+w}{ }\PY{k}{new}\PY{+w}{ }\PY{n}{Car}\PY{p}{(}\PY{p}{)}\PY{p}{;}
\PY{n}{System}\PY{p}{.}\PY{n+na}{out}\PY{p}{.}\PY{n+na}{println}\PY{p}{(}\PY{n}{v}\PY{p}{.}\PY{n+na}{speed}\PY{p}{)}\PY{p}{;}\PY{+w}{ }\PY{c+c1}{// 10 (uses Vehicle\PYZsq{}s field, based on reference type!)}

\PY{n}{Car}\PY{+w}{ }\PY{n}{c}\PY{+w}{ }\PY{o}{=}\PY{+w}{ }\PY{k}{new}\PY{+w}{ }\PY{n}{Car}\PY{p}{(}\PY{p}{)}\PY{p}{;}
\PY{n}{System}\PY{p}{.}\PY{n+na}{out}\PY{p}{.}\PY{n+na}{println}\PY{p}{(}\PY{n}{c}\PY{p}{.}\PY{n+na}{speed}\PY{p}{)}\PY{p}{;}\PY{+w}{ }\PY{c+c1}{// 20}
\end{Verbatim}
\end{codebox}

This is a classic code-review red flag: shadowing fields across a class hierarchy is confusing and error-prone. Prefer accessing state through methods, not fields, across inheritance boundaries.

\section[Polymorphism]{Polymorphism}
\label{h:2:polymorphism}
\textbf{Polymorphism} (“many forms”) means a single reference type can point to different actual object types, and calling a method on it runs the correct version for the \emph{actual} object — decided at runtime. This is also called \textbf{dynamic dispatch} or \textbf{runtime polymorphism}.

\begin{codebox}
\begin{Verbatim}[commandchars=\\\{\}]
\PY{k+kd}{public}\PY{+w}{ }\PY{k+kd}{abstract}\PY{+w}{ }\PY{k+kd}{class} \PY{n+nc}{Shape}\PY{+w}{ }\PY{p}{\PYZob{}}
\PY{+w}{    }\PY{k+kd}{public}\PY{+w}{ }\PY{k+kd}{abstract}\PY{+w}{ }\PY{k+kt}{double}\PY{+w}{ }\PY{n+nf}{area}\PY{p}{(}\PY{p}{)}\PY{p}{;}
\PY{p}{\PYZcb{}}

\PY{k+kd}{public}\PY{+w}{ }\PY{k+kd}{class} \PY{n+nc}{Circle}\PY{+w}{ }\PY{k+kd}{extends}\PY{+w}{ }\PY{n}{Shape}\PY{+w}{ }\PY{p}{\PYZob{}}
\PY{+w}{    }\PY{k+kd}{private}\PY{+w}{ }\PY{k+kd}{final}\PY{+w}{ }\PY{k+kt}{double}\PY{+w}{ }\PY{n}{radius}\PY{p}{;}
\PY{+w}{    }\PY{k+kd}{public}\PY{+w}{ }\PY{n+nf}{Circle}\PY{p}{(}\PY{k+kt}{double}\PY{+w}{ }\PY{n}{radius}\PY{p}{)}\PY{+w}{ }\PY{p}{\PYZob{}}\PY{+w}{ }\PY{k}{this}\PY{p}{.}\PY{n+na}{radius}\PY{+w}{ }\PY{o}{=}\PY{+w}{ }\PY{n}{radius}\PY{p}{;}\PY{+w}{ }\PY{p}{\PYZcb{}}
\PY{+w}{    }\PY{n+nd}{@Override}
\PY{+w}{    }\PY{k+kd}{public}\PY{+w}{ }\PY{k+kt}{double}\PY{+w}{ }\PY{n+nf}{area}\PY{p}{(}\PY{p}{)}\PY{+w}{ }\PY{p}{\PYZob{}}\PY{+w}{ }\PY{k}{return}\PY{+w}{ }\PY{n}{Math}\PY{p}{.}\PY{n+na}{PI}\PY{+w}{ }\PY{o}{*}\PY{+w}{ }\PY{n}{radius}\PY{+w}{ }\PY{o}{*}\PY{+w}{ }\PY{n}{radius}\PY{p}{;}\PY{+w}{ }\PY{p}{\PYZcb{}}
\PY{p}{\PYZcb{}}

\PY{k+kd}{public}\PY{+w}{ }\PY{k+kd}{class} \PY{n+nc}{Rectangle}\PY{+w}{ }\PY{k+kd}{extends}\PY{+w}{ }\PY{n}{Shape}\PY{+w}{ }\PY{p}{\PYZob{}}
\PY{+w}{    }\PY{k+kd}{private}\PY{+w}{ }\PY{k+kd}{final}\PY{+w}{ }\PY{k+kt}{double}\PY{+w}{ }\PY{n}{width}\PY{p}{,}\PY{+w}{ }\PY{n}{height}\PY{p}{;}
\PY{+w}{    }\PY{k+kd}{public}\PY{+w}{ }\PY{n+nf}{Rectangle}\PY{p}{(}\PY{k+kt}{double}\PY{+w}{ }\PY{n}{width}\PY{p}{,}\PY{+w}{ }\PY{k+kt}{double}\PY{+w}{ }\PY{n}{height}\PY{p}{)}\PY{+w}{ }\PY{p}{\PYZob{}}
\PY{+w}{        }\PY{k}{this}\PY{p}{.}\PY{n+na}{width}\PY{+w}{ }\PY{o}{=}\PY{+w}{ }\PY{n}{width}\PY{p}{;}
\PY{+w}{        }\PY{k}{this}\PY{p}{.}\PY{n+na}{height}\PY{+w}{ }\PY{o}{=}\PY{+w}{ }\PY{n}{height}\PY{p}{;}
\PY{+w}{    }\PY{p}{\PYZcb{}}
\PY{+w}{    }\PY{n+nd}{@Override}
\PY{+w}{    }\PY{k+kd}{public}\PY{+w}{ }\PY{k+kt}{double}\PY{+w}{ }\PY{n+nf}{area}\PY{p}{(}\PY{p}{)}\PY{+w}{ }\PY{p}{\PYZob{}}\PY{+w}{ }\PY{k}{return}\PY{+w}{ }\PY{n}{width}\PY{+w}{ }\PY{o}{*}\PY{+w}{ }\PY{n}{height}\PY{p}{;}\PY{+w}{ }\PY{p}{\PYZcb{}}
\PY{p}{\PYZcb{}}

\PY{k+kd}{public}\PY{+w}{ }\PY{k+kd}{class} \PY{n+nc}{AreaCalculator}\PY{+w}{ }\PY{p}{\PYZob{}}
\PY{+w}{    }\PY{k+kd}{public}\PY{+w}{ }\PY{k+kd}{static}\PY{+w}{ }\PY{k+kt}{void}\PY{+w}{ }\PY{n+nf}{main}\PY{p}{(}\PY{n}{String}\PY{o}{[}\PY{o}{]}\PY{+w}{ }\PY{n}{args}\PY{p}{)}\PY{+w}{ }\PY{p}{\PYZob{}}
\PY{+w}{        }\PY{n}{List}\PY{o}{\PYZlt{}}\PY{n}{Shape}\PY{o}{\PYZgt{}}\PY{+w}{ }\PY{n}{shapes}\PY{+w}{ }\PY{o}{=}\PY{+w}{ }\PY{n}{List}\PY{p}{.}\PY{n+na}{of}\PY{p}{(}\PY{k}{new}\PY{+w}{ }\PY{n}{Circle}\PY{p}{(}\PY{l+m+mi}{2}\PY{p}{)}\PY{p}{,}\PY{+w}{ }\PY{k}{new}\PY{+w}{ }\PY{n}{Rectangle}\PY{p}{(}\PY{l+m+mi}{3}\PY{p}{,}\PY{+w}{ }\PY{l+m+mi}{4}\PY{p}{)}\PY{p}{)}\PY{p}{;}
\PY{+w}{        }\PY{k}{for}\PY{+w}{ }\PY{p}{(}\PY{n}{Shape}\PY{+w}{ }\PY{n}{shape}\PY{+w}{ }\PY{p}{:}\PY{+w}{ }\PY{n}{shapes}\PY{p}{)}\PY{+w}{ }\PY{p}{\PYZob{}}
\PY{+w}{            }\PY{n}{System}\PY{p}{.}\PY{n+na}{out}\PY{p}{.}\PY{n+na}{println}\PY{p}{(}\PY{n}{shape}\PY{p}{.}\PY{n+na}{area}\PY{p}{(}\PY{p}{)}\PY{p}{)}\PY{p}{;}\PY{+w}{ }\PY{c+c1}{// correct area() runs for each real type}
\PY{+w}{        }\PY{p}{\PYZcb{}}
\PY{+w}{        }\PY{c+c1}{// Output:}
\PY{+w}{        }\PY{c+c1}{// 12.566370614359172}
\PY{+w}{        }\PY{c+c1}{// 12.0}
\PY{+w}{    }\PY{p}{\PYZcb{}}
\PY{p}{\PYZcb{}}
\end{Verbatim}
\end{codebox}

The loop variable is declared as \code{Shape}, but each call to \code{area()} runs the actual subclass’s implementation. The calling code doesn’t need \code{if~\allowbreak{}(\allowbreak{}shape~\allowbreak{}instanceof~\allowbreak{}Circle)\allowbreak{}~\allowbreak{}...} checks — that’s the benefit of polymorphism.

There are two flavors people mention in interviews:

{\tablefont
\begin{xltabular}{\linewidth}{@{}L{0.962}L{0.577}L{1.769}L{0.692}@{}}
\toprule
\rowcolor{tablehdr}
\thd{Type} & \thd{Also called} & \thd{Resolved when} & \thd{Example} \\
\midrule
\endhead
Compile-time polymorphism & Static binding & At compile time, by looking at parameter types & Method overloading \\
Runtime polymorphism & Dynamic binding & At runtime, by looking at the actual object & Method overriding \\
\bottomrule
\end{xltabular}
}

\textbf{Gotcha: \code{static} methods are not polymorphic.} They are resolved by the \emph{reference type}, not the actual object — this is called \textbf{method hiding}, not overriding.

\begin{codebox}
\begin{Verbatim}[commandchars=\\\{\}]
\PY{k+kd}{public}\PY{+w}{ }\PY{k+kd}{class} \PY{n+nc}{Vehicle}\PY{+w}{ }\PY{p}{\PYZob{}}
\PY{+w}{    }\PY{k+kd}{public}\PY{+w}{ }\PY{k+kd}{static}\PY{+w}{ }\PY{n}{String}\PY{+w}{ }\PY{n+nf}{category}\PY{p}{(}\PY{p}{)}\PY{+w}{ }\PY{p}{\PYZob{}}\PY{+w}{ }\PY{k}{return}\PY{+w}{ }\PY{l+s}{\PYZdq{}}\PY{l+s}{Generic Vehicle}\PY{l+s}{\PYZdq{}}\PY{p}{;}\PY{+w}{ }\PY{p}{\PYZcb{}}
\PY{p}{\PYZcb{}}

\PY{k+kd}{public}\PY{+w}{ }\PY{k+kd}{class} \PY{n+nc}{Car}\PY{+w}{ }\PY{k+kd}{extends}\PY{+w}{ }\PY{n}{Vehicle}\PY{+w}{ }\PY{p}{\PYZob{}}
\PY{+w}{    }\PY{k+kd}{public}\PY{+w}{ }\PY{k+kd}{static}\PY{+w}{ }\PY{n}{String}\PY{+w}{ }\PY{n+nf}{category}\PY{p}{(}\PY{p}{)}\PY{+w}{ }\PY{p}{\PYZob{}}\PY{+w}{ }\PY{k}{return}\PY{+w}{ }\PY{l+s}{\PYZdq{}}\PY{l+s}{Car}\PY{l+s}{\PYZdq{}}\PY{p}{;}\PY{+w}{ }\PY{p}{\PYZcb{}}\PY{+w}{ }\PY{c+c1}{// hides, not overrides}
\PY{p}{\PYZcb{}}

\PY{n}{Vehicle}\PY{+w}{ }\PY{n}{v}\PY{+w}{ }\PY{o}{=}\PY{+w}{ }\PY{k}{new}\PY{+w}{ }\PY{n}{Car}\PY{p}{(}\PY{p}{)}\PY{p}{;}
\PY{n}{System}\PY{p}{.}\PY{n+na}{out}\PY{p}{.}\PY{n+na}{println}\PY{p}{(}\PY{n}{v}\PY{p}{.}\PY{n+na}{category}\PY{p}{(}\PY{p}{)}\PY{p}{)}\PY{p}{;}\PY{+w}{ }\PY{c+c1}{// \PYZdq{}Generic Vehicle\PYZdq{} — decided by reference type!}
\end{Verbatim}
\end{codebox}

\section[Abstraction]{Abstraction}
\label{h:2:abstraction}
\textbf{Abstraction} means exposing \emph{what} something does while hiding \emph{how} it does it. In Java this is done with \code{abstract} classes and interfaces. An \code{abstract} class can have some implemented methods and some unimplemented (\code{abstract}) ones; it cannot be instantiated directly.

\begin{codebox}
\begin{Verbatim}[commandchars=\\\{\}]
\PY{k+kd}{public}\PY{+w}{ }\PY{k+kd}{abstract}\PY{+w}{ }\PY{k+kd}{class} \PY{n+nc}{PaymentMethod}\PY{+w}{ }\PY{p}{\PYZob{}}
\PY{+w}{    }\PY{c+c1}{// Abstract method: no body, subclasses must implement it}
\PY{+w}{    }\PY{k+kd}{public}\PY{+w}{ }\PY{k+kd}{abstract}\PY{+w}{ }\PY{k+kt}{boolean}\PY{+w}{ }\PY{n+nf}{authorize}\PY{p}{(}\PY{k+kt}{double}\PY{+w}{ }\PY{n}{amount}\PY{p}{)}\PY{p}{;}

\PY{+w}{    }\PY{c+c1}{// Concrete method: shared logic for all payment methods}
\PY{+w}{    }\PY{k+kd}{public}\PY{+w}{ }\PY{k+kd}{final}\PY{+w}{ }\PY{k+kt}{void}\PY{+w}{ }\PY{n+nf}{pay}\PY{p}{(}\PY{k+kt}{double}\PY{+w}{ }\PY{n}{amount}\PY{p}{)}\PY{+w}{ }\PY{p}{\PYZob{}}
\PY{+w}{        }\PY{k}{if}\PY{+w}{ }\PY{p}{(}\PY{n}{authorize}\PY{p}{(}\PY{n}{amount}\PY{p}{)}\PY{p}{)}\PY{+w}{ }\PY{p}{\PYZob{}}
\PY{+w}{            }\PY{n}{System}\PY{p}{.}\PY{n+na}{out}\PY{p}{.}\PY{n+na}{println}\PY{p}{(}\PY{l+s}{\PYZdq{}}\PY{l+s}{Paid }\PY{l+s}{\PYZdq{}}\PY{+w}{ }\PY{o}{+}\PY{+w}{ }\PY{n}{amount}\PY{p}{)}\PY{p}{;}
\PY{+w}{        }\PY{p}{\PYZcb{}}\PY{+w}{ }\PY{k}{else}\PY{+w}{ }\PY{p}{\PYZob{}}
\PY{+w}{            }\PY{n}{System}\PY{p}{.}\PY{n+na}{out}\PY{p}{.}\PY{n+na}{println}\PY{p}{(}\PY{l+s}{\PYZdq{}}\PY{l+s}{Payment declined}\PY{l+s}{\PYZdq{}}\PY{p}{)}\PY{p}{;}
\PY{+w}{        }\PY{p}{\PYZcb{}}
\PY{+w}{    }\PY{p}{\PYZcb{}}
\PY{p}{\PYZcb{}}

\PY{k+kd}{public}\PY{+w}{ }\PY{k+kd}{class} \PY{n+nc}{CreditCardPayment}\PY{+w}{ }\PY{k+kd}{extends}\PY{+w}{ }\PY{n}{PaymentMethod}\PY{+w}{ }\PY{p}{\PYZob{}}
\PY{+w}{    }\PY{n+nd}{@Override}
\PY{+w}{    }\PY{k+kd}{public}\PY{+w}{ }\PY{k+kt}{boolean}\PY{+w}{ }\PY{n+nf}{authorize}\PY{p}{(}\PY{k+kt}{double}\PY{+w}{ }\PY{n}{amount}\PY{p}{)}\PY{+w}{ }\PY{p}{\PYZob{}}
\PY{+w}{        }\PY{k}{return}\PY{+w}{ }\PY{n}{amount}\PY{+w}{ }\PY{o}{\PYZlt{}}\PY{o}{=}\PY{+w}{ }\PY{l+m+mi}{5000}\PY{p}{;}\PY{+w}{ }\PY{c+c1}{// simplified rule}
\PY{+w}{    }\PY{p}{\PYZcb{}}
\PY{p}{\PYZcb{}}

\PY{n}{PaymentMethod}\PY{+w}{ }\PY{n}{payment}\PY{+w}{ }\PY{o}{=}\PY{+w}{ }\PY{k}{new}\PY{+w}{ }\PY{n}{CreditCardPayment}\PY{p}{(}\PY{p}{)}\PY{p}{;}
\PY{n}{payment}\PY{p}{.}\PY{n+na}{pay}\PY{p}{(}\PY{l+m+mi}{1200}\PY{p}{)}\PY{p}{;}\PY{+w}{ }\PY{c+c1}{// Paid 1200.0}
\end{Verbatim}
\end{codebox}

\begin{codebox}
\begin{Verbatim}[commandchars=\\\{\}]
\PY{n}{PaymentMethod}\PY{+w}{ }\PY{n}{payment}\PY{+w}{ }\PY{o}{=}\PY{+w}{ }\PY{k}{new}\PY{+w}{ }\PY{n}{PaymentMethod}\PY{p}{(}\PY{p}{)}\PY{p}{;}\PY{+w}{ }\PY{c+c1}{// Compile error: PaymentMethod is abstract}
\end{Verbatim}
\end{codebox}

Interfaces (covered fully in the “Modern Java Language Features” chapter) are a purer form of abstraction: historically 100\% abstract, though modern Java interfaces can also have \code{default}, \code{static}, and \code{private} methods.

\begin{codebox}
\begin{Verbatim}[commandchars=\\\{\}]
\PY{k+kd}{public}\PY{+w}{ }\PY{k+kd}{interface} \PY{n+nc}{Discountable}\PY{+w}{ }\PY{p}{\PYZob{}}
\PY{+w}{    }\PY{k+kt}{double}\PY{+w}{ }\PY{n+nf}{applyDiscount}\PY{p}{(}\PY{k+kt}{double}\PY{+w}{ }\PY{n}{price}\PY{p}{)}\PY{p}{;}
\PY{p}{\PYZcb{}}

\PY{k+kd}{public}\PY{+w}{ }\PY{k+kd}{class} \PY{n+nc}{Order}\PY{+w}{ }\PY{k+kd}{implements}\PY{+w}{ }\PY{n}{Discountable}\PY{+w}{ }\PY{p}{\PYZob{}}
\PY{+w}{    }\PY{n+nd}{@Override}
\PY{+w}{    }\PY{k+kd}{public}\PY{+w}{ }\PY{k+kt}{double}\PY{+w}{ }\PY{n+nf}{applyDiscount}\PY{p}{(}\PY{k+kt}{double}\PY{+w}{ }\PY{n}{price}\PY{p}{)}\PY{+w}{ }\PY{p}{\PYZob{}}
\PY{+w}{        }\PY{k}{return}\PY{+w}{ }\PY{n}{price}\PY{+w}{ }\PY{o}{*}\PY{+w}{ }\PY{l+m+mf}{0.9}\PY{p}{;}\PY{+w}{ }\PY{c+c1}{// 10\PYZpc{} off}
\PY{+w}{    }\PY{p}{\PYZcb{}}
\PY{p}{\PYZcb{}}
\end{Verbatim}
\end{codebox}

\textbf{Rule of thumb for code review:} use an abstract class when subclasses share state or common implementation; use an interface when you’re describing a capability that unrelated classes might implement (\code{Comparable}, \code{Discountable}, \code{Serializable}).

\section[Method Overloading]{Method Overloading}
\label{h:2:method-overloading}
\textbf{Method overloading} means defining multiple methods with the \emph{same name} but \emph{different parameter lists} (different type, number, or order of parameters) in the same class. It is resolved at \textbf{compile time} based on the arguments’ declared types — this is why it’s called compile-time (static) polymorphism.

\begin{codebox}
\begin{Verbatim}[commandchars=\\\{\}]
\PY{k+kd}{public}\PY{+w}{ }\PY{k+kd}{class} \PY{n+nc}{Invoice}\PY{+w}{ }\PY{p}{\PYZob{}}
\PY{+w}{    }\PY{k+kd}{public}\PY{+w}{ }\PY{k+kt}{double}\PY{+w}{ }\PY{n+nf}{calculateTotal}\PY{p}{(}\PY{k+kt}{double}\PY{+w}{ }\PY{n}{price}\PY{p}{,}\PY{+w}{ }\PY{k+kt}{int}\PY{+w}{ }\PY{n}{quantity}\PY{p}{)}\PY{+w}{ }\PY{p}{\PYZob{}}
\PY{+w}{        }\PY{k}{return}\PY{+w}{ }\PY{n}{price}\PY{+w}{ }\PY{o}{*}\PY{+w}{ }\PY{n}{quantity}\PY{p}{;}
\PY{+w}{    }\PY{p}{\PYZcb{}}

\PY{+w}{    }\PY{k+kd}{public}\PY{+w}{ }\PY{k+kt}{double}\PY{+w}{ }\PY{n+nf}{calculateTotal}\PY{p}{(}\PY{k+kt}{double}\PY{+w}{ }\PY{n}{price}\PY{p}{,}\PY{+w}{ }\PY{k+kt}{int}\PY{+w}{ }\PY{n}{quantity}\PY{p}{,}\PY{+w}{ }\PY{k+kt}{double}\PY{+w}{ }\PY{n}{discount}\PY{p}{)}\PY{+w}{ }\PY{p}{\PYZob{}}
\PY{+w}{        }\PY{k}{return}\PY{+w}{ }\PY{n}{price}\PY{+w}{ }\PY{o}{*}\PY{+w}{ }\PY{n}{quantity}\PY{+w}{ }\PY{o}{*}\PY{+w}{ }\PY{p}{(}\PY{l+m+mi}{1}\PY{+w}{ }\PY{o}{\PYZhy{}}\PY{+w}{ }\PY{n}{discount}\PY{p}{)}\PY{p}{;}
\PY{+w}{    }\PY{p}{\PYZcb{}}

\PY{+w}{    }\PY{k+kd}{public}\PY{+w}{ }\PY{k+kt}{double}\PY{+w}{ }\PY{n+nf}{calculateTotal}\PY{p}{(}\PY{k+kt}{double}\PY{+w}{ }\PY{n}{price}\PY{p}{,}\PY{+w}{ }\PY{k+kt}{int}\PY{+w}{ }\PY{n}{quantity}\PY{p}{,}\PY{+w}{ }\PY{k+kt}{double}\PY{+w}{ }\PY{n}{discount}\PY{p}{,}\PY{+w}{ }\PY{k+kt}{double}\PY{+w}{ }\PY{n}{taxRate}\PY{p}{)}\PY{+w}{ }\PY{p}{\PYZob{}}
\PY{+w}{        }\PY{k}{return}\PY{+w}{ }\PY{n}{price}\PY{+w}{ }\PY{o}{*}\PY{+w}{ }\PY{n}{quantity}\PY{+w}{ }\PY{o}{*}\PY{+w}{ }\PY{p}{(}\PY{l+m+mi}{1}\PY{+w}{ }\PY{o}{\PYZhy{}}\PY{+w}{ }\PY{n}{discount}\PY{p}{)}\PY{+w}{ }\PY{o}{*}\PY{+w}{ }\PY{p}{(}\PY{l+m+mi}{1}\PY{+w}{ }\PY{o}{+}\PY{+w}{ }\PY{n}{taxRate}\PY{p}{)}\PY{p}{;}
\PY{+w}{    }\PY{p}{\PYZcb{}}
\PY{p}{\PYZcb{}}

\PY{n}{Invoice}\PY{+w}{ }\PY{n}{invoice}\PY{+w}{ }\PY{o}{=}\PY{+w}{ }\PY{k}{new}\PY{+w}{ }\PY{n}{Invoice}\PY{p}{(}\PY{p}{)}\PY{p}{;}
\PY{n}{System}\PY{p}{.}\PY{n+na}{out}\PY{p}{.}\PY{n+na}{println}\PY{p}{(}\PY{n}{invoice}\PY{p}{.}\PY{n+na}{calculateTotal}\PY{p}{(}\PY{l+m+mf}{10.0}\PY{p}{,}\PY{+w}{ }\PY{l+m+mi}{3}\PY{p}{)}\PY{p}{)}\PY{p}{;}\PY{+w}{             }\PY{c+c1}{// 30.0}
\PY{n}{System}\PY{p}{.}\PY{n+na}{out}\PY{p}{.}\PY{n+na}{println}\PY{p}{(}\PY{n}{invoice}\PY{p}{.}\PY{n+na}{calculateTotal}\PY{p}{(}\PY{l+m+mf}{10.0}\PY{p}{,}\PY{+w}{ }\PY{l+m+mi}{3}\PY{p}{,}\PY{+w}{ }\PY{l+m+mf}{0.1}\PY{p}{)}\PY{p}{)}\PY{p}{;}\PY{+w}{        }\PY{c+c1}{// 27.0}
\end{Verbatim}
\end{codebox}

You cannot overload by return type alone — the parameter list must differ.

\begin{codebox}
\begin{Verbatim}[commandchars=\\\{\}]
\PY{k+kd}{public}\PY{+w}{ }\PY{k+kt}{int}\PY{+w}{ }\PY{n+nf}{process}\PY{p}{(}\PY{n}{String}\PY{+w}{ }\PY{n}{data}\PY{p}{)}\PY{+w}{ }\PY{p}{\PYZob{}}\PY{+w}{ }\PY{k}{return}\PY{+w}{ }\PY{l+m+mi}{1}\PY{p}{;}\PY{+w}{ }\PY{p}{\PYZcb{}}
\PY{k+kd}{public}\PY{+w}{ }\PY{k+kt}{double}\PY{+w}{ }\PY{n+nf}{process}\PY{p}{(}\PY{n}{String}\PY{+w}{ }\PY{n}{data}\PY{p}{)}\PY{+w}{ }\PY{p}{\PYZob{}}\PY{+w}{ }\PY{k}{return}\PY{+w}{ }\PY{l+m+mf}{1.0}\PY{p}{;}\PY{+w}{ }\PY{p}{\PYZcb{}}\PY{+w}{ }\PY{c+c1}{// Compile error: same signature}
\end{Verbatim}
\end{codebox}

\textbf{Gotcha: overload resolution with autoboxing and varargs.} Java picks the \emph{most specific} match, and it prefers, in order: (1) exact match with widening of primitives, (2) autoboxing/unboxing, (3) varargs. This can surprise people.

\begin{codebox}
\begin{Verbatim}[commandchars=\\\{\}]
\PY{k+kd}{public}\PY{+w}{ }\PY{k+kd}{class} \PY{n+nc}{Logger}\PY{+w}{ }\PY{p}{\PYZob{}}
\PY{+w}{    }\PY{k+kd}{public}\PY{+w}{ }\PY{k+kt}{void}\PY{+w}{ }\PY{n+nf}{log}\PY{p}{(}\PY{k+kt}{int}\PY{+w}{ }\PY{n}{code}\PY{p}{)}\PY{+w}{ }\PY{p}{\PYZob{}}
\PY{+w}{        }\PY{n}{System}\PY{p}{.}\PY{n+na}{out}\PY{p}{.}\PY{n+na}{println}\PY{p}{(}\PY{l+s}{\PYZdq{}}\PY{l+s}{int version: }\PY{l+s}{\PYZdq{}}\PY{+w}{ }\PY{o}{+}\PY{+w}{ }\PY{n}{code}\PY{p}{)}\PY{p}{;}
\PY{+w}{    }\PY{p}{\PYZcb{}}

\PY{+w}{    }\PY{k+kd}{public}\PY{+w}{ }\PY{k+kt}{void}\PY{+w}{ }\PY{n+nf}{log}\PY{p}{(}\PY{n}{Integer}\PY{+w}{ }\PY{n}{code}\PY{p}{)}\PY{+w}{ }\PY{p}{\PYZob{}}
\PY{+w}{        }\PY{n}{System}\PY{p}{.}\PY{n+na}{out}\PY{p}{.}\PY{n+na}{println}\PY{p}{(}\PY{l+s}{\PYZdq{}}\PY{l+s}{Integer version: }\PY{l+s}{\PYZdq{}}\PY{+w}{ }\PY{o}{+}\PY{+w}{ }\PY{n}{code}\PY{p}{)}\PY{p}{;}
\PY{+w}{    }\PY{p}{\PYZcb{}}

\PY{+w}{    }\PY{k+kd}{public}\PY{+w}{ }\PY{k+kt}{void}\PY{+w}{ }\PY{n+nf}{log}\PY{p}{(}\PY{n}{Object}\PY{p}{.}\PY{p}{.}\PY{p}{.}\PY{+w}{ }\PY{n}{codes}\PY{p}{)}\PY{+w}{ }\PY{p}{\PYZob{}}
\PY{+w}{        }\PY{n}{System}\PY{p}{.}\PY{n+na}{out}\PY{p}{.}\PY{n+na}{println}\PY{p}{(}\PY{l+s}{\PYZdq{}}\PY{l+s}{varargs version}\PY{l+s}{\PYZdq{}}\PY{p}{)}\PY{p}{;}
\PY{+w}{    }\PY{p}{\PYZcb{}}
\PY{p}{\PYZcb{}}

\PY{n}{Logger}\PY{+w}{ }\PY{n}{logger}\PY{+w}{ }\PY{o}{=}\PY{+w}{ }\PY{k}{new}\PY{+w}{ }\PY{n}{Logger}\PY{p}{(}\PY{p}{)}\PY{p}{;}
\PY{n}{logger}\PY{p}{.}\PY{n+na}{log}\PY{p}{(}\PY{l+m+mi}{5}\PY{p}{)}\PY{p}{;}\PY{+w}{ }\PY{c+c1}{// \PYZdq{}int version: 5\PYZdq{}}
\PY{c+c1}{// Java prefers the exact primitive match over autoboxing to Integer,}
\PY{c+c1}{// and prefers both of those over varargs.}
\end{Verbatim}
\end{codebox}

\begin{codebox}
\begin{Verbatim}[commandchars=\\\{\}]
\PY{k+kd}{public}\PY{+w}{ }\PY{k+kd}{class} \PY{n+nc}{Calculator}\PY{+w}{ }\PY{p}{\PYZob{}}
\PY{+w}{    }\PY{k+kd}{public}\PY{+w}{ }\PY{k+kt}{void}\PY{+w}{ }\PY{n+nf}{print}\PY{p}{(}\PY{k+kt}{long}\PY{+w}{ }\PY{n}{value}\PY{p}{)}\PY{+w}{   }\PY{p}{\PYZob{}}\PY{+w}{ }\PY{n}{System}\PY{p}{.}\PY{n+na}{out}\PY{p}{.}\PY{n+na}{println}\PY{p}{(}\PY{l+s}{\PYZdq{}}\PY{l+s}{long: }\PY{l+s}{\PYZdq{}}\PY{+w}{ }\PY{o}{+}\PY{+w}{ }\PY{n}{value}\PY{p}{)}\PY{p}{;}\PY{+w}{ }\PY{p}{\PYZcb{}}
\PY{+w}{    }\PY{k+kd}{public}\PY{+w}{ }\PY{k+kt}{void}\PY{+w}{ }\PY{n+nf}{print}\PY{p}{(}\PY{n}{Integer}\PY{+w}{ }\PY{n}{value}\PY{p}{)}\PY{+w}{ }\PY{p}{\PYZob{}}\PY{+w}{ }\PY{n}{System}\PY{p}{.}\PY{n+na}{out}\PY{p}{.}\PY{n+na}{println}\PY{p}{(}\PY{l+s}{\PYZdq{}}\PY{l+s}{Integer: }\PY{l+s}{\PYZdq{}}\PY{+w}{ }\PY{o}{+}\PY{+w}{ }\PY{n}{value}\PY{p}{)}\PY{p}{;}\PY{+w}{ }\PY{p}{\PYZcb{}}

\PY{+w}{    }\PY{k+kd}{public}\PY{+w}{ }\PY{k+kd}{static}\PY{+w}{ }\PY{k+kt}{void}\PY{+w}{ }\PY{n+nf}{main}\PY{p}{(}\PY{n}{String}\PY{o}{[}\PY{o}{]}\PY{+w}{ }\PY{n}{args}\PY{p}{)}\PY{+w}{ }\PY{p}{\PYZob{}}
\PY{+w}{        }\PY{k+kt}{int}\PY{+w}{ }\PY{n}{x}\PY{+w}{ }\PY{o}{=}\PY{+w}{ }\PY{l+m+mi}{5}\PY{p}{;}
\PY{+w}{        }\PY{k}{new}\PY{+w}{ }\PY{n}{Calculator}\PY{p}{(}\PY{p}{)}\PY{p}{.}\PY{n+na}{print}\PY{p}{(}\PY{n}{x}\PY{p}{)}\PY{p}{;}
\PY{+w}{        }\PY{c+c1}{// Output: \PYZdq{}long: 5\PYZdq{}}
\PY{+w}{        }\PY{c+c1}{// Widening (int \PYZhy{}\PYZgt{} long) is preferred over autoboxing (int \PYZhy{}\PYZgt{} Integer)!}
\PY{+w}{    }\PY{p}{\PYZcb{}}
\PY{p}{\PYZcb{}}
\end{Verbatim}
\end{codebox}

This is a frequent code-review discussion point: overloads that mix primitives, wrapper types, and varargs are a maintenance hazard because the “obviously correct” overload isn’t always the one that’s chosen.

\section[Method Overriding]{Method Overriding}
\label{h:2:method-overriding}
\textbf{Method overriding} means a subclass provides its own implementation of a method already defined in its superclass, with the \emph{same signature}. It is resolved at \textbf{runtime} based on the actual object type — runtime (dynamic) polymorphism. Always use the \code{@\allowbreak{}Override} annotation; it makes the compiler check that you’re actually overriding something.

\begin{codebox}
\begin{Verbatim}[commandchars=\\\{\}]
\PY{k+kd}{public}\PY{+w}{ }\PY{k+kd}{class} \PY{n+nc}{Shape}\PY{+w}{ }\PY{p}{\PYZob{}}
\PY{+w}{    }\PY{k+kd}{public}\PY{+w}{ }\PY{k+kt}{double}\PY{+w}{ }\PY{n+nf}{area}\PY{p}{(}\PY{p}{)}\PY{+w}{ }\PY{p}{\PYZob{}}
\PY{+w}{        }\PY{k}{return}\PY{+w}{ }\PY{l+m+mf}{0.0}\PY{p}{;}
\PY{+w}{    }\PY{p}{\PYZcb{}}

\PY{+w}{    }\PY{n+nd}{@Override}
\PY{+w}{    }\PY{k+kd}{public}\PY{+w}{ }\PY{n}{String}\PY{+w}{ }\PY{n+nf}{toString}\PY{p}{(}\PY{p}{)}\PY{+w}{ }\PY{p}{\PYZob{}}
\PY{+w}{        }\PY{k}{return}\PY{+w}{ }\PY{l+s}{\PYZdq{}}\PY{l+s}{Shape with area }\PY{l+s}{\PYZdq{}}\PY{+w}{ }\PY{o}{+}\PY{+w}{ }\PY{n}{area}\PY{p}{(}\PY{p}{)}\PY{p}{;}
\PY{+w}{    }\PY{p}{\PYZcb{}}
\PY{p}{\PYZcb{}}

\PY{k+kd}{public}\PY{+w}{ }\PY{k+kd}{class} \PY{n+nc}{Square}\PY{+w}{ }\PY{k+kd}{extends}\PY{+w}{ }\PY{n}{Shape}\PY{+w}{ }\PY{p}{\PYZob{}}
\PY{+w}{    }\PY{k+kd}{private}\PY{+w}{ }\PY{k+kd}{final}\PY{+w}{ }\PY{k+kt}{double}\PY{+w}{ }\PY{n}{side}\PY{p}{;}
\PY{+w}{    }\PY{k+kd}{public}\PY{+w}{ }\PY{n+nf}{Square}\PY{p}{(}\PY{k+kt}{double}\PY{+w}{ }\PY{n}{side}\PY{p}{)}\PY{+w}{ }\PY{p}{\PYZob{}}\PY{+w}{ }\PY{k}{this}\PY{p}{.}\PY{n+na}{side}\PY{+w}{ }\PY{o}{=}\PY{+w}{ }\PY{n}{side}\PY{p}{;}\PY{+w}{ }\PY{p}{\PYZcb{}}

\PY{+w}{    }\PY{n+nd}{@Override}
\PY{+w}{    }\PY{k+kd}{public}\PY{+w}{ }\PY{k+kt}{double}\PY{+w}{ }\PY{n+nf}{area}\PY{p}{(}\PY{p}{)}\PY{+w}{ }\PY{p}{\PYZob{}}
\PY{+w}{        }\PY{k}{return}\PY{+w}{ }\PY{n}{side}\PY{+w}{ }\PY{o}{*}\PY{+w}{ }\PY{n}{side}\PY{p}{;}
\PY{+w}{    }\PY{p}{\PYZcb{}}
\PY{p}{\PYZcb{}}

\PY{n}{Shape}\PY{+w}{ }\PY{n}{shape}\PY{+w}{ }\PY{o}{=}\PY{+w}{ }\PY{k}{new}\PY{+w}{ }\PY{n}{Square}\PY{p}{(}\PY{l+m+mi}{4}\PY{p}{)}\PY{p}{;}
\PY{n}{System}\PY{p}{.}\PY{n+na}{out}\PY{p}{.}\PY{n+na}{println}\PY{p}{(}\PY{n}{shape}\PY{p}{.}\PY{n+na}{area}\PY{p}{(}\PY{p}{)}\PY{p}{)}\PY{p}{;}\PY{+w}{     }\PY{c+c1}{// 16.0}
\PY{n}{System}\PY{p}{.}\PY{n+na}{out}\PY{p}{.}\PY{n+na}{println}\PY{p}{(}\PY{n}{shape}\PY{p}{)}\PY{p}{;}\PY{+w}{            }\PY{c+c1}{// Shape with area 16.0 (toString overridden indirectly via area())}
\end{Verbatim}
\end{codebox}

Overriding rules the compiler enforces:

\begin{itemize}
\item 
Same method name and parameter list.

\item 
Return type must be the same or a \textbf{covariant} (narrower) subtype.

\item 
Access modifier can stay the same or become \emph{more} visible, never less visible.

\item 
Cannot override a method marked \code{final} or \code{static} (that becomes hiding, not overriding — see \hyperref[h:2:static-members]{Static Members}).

\end{itemize}

{\tablefont
\begin{xltabular}{\linewidth}{@{}L{0.652}L{1.261}L{1.087}@{}}
\toprule
\rowcolor{tablehdr}
\thd{} & \thd{Method Overloading} & \thd{Method Overriding} \\
\midrule
\endhead
Same method name? & Yes & Yes \\
Parameter list & Must differ & Must be identical \\
Return type & Can differ freely & Must be same or covariant \\
Resolved & Compile time (static binding) & Runtime (dynamic binding) \\
Relationship & Same class (or same class hierarchy) & Superclass / subclass \\
Access modifier & Can differ freely & Cannot be more restrictive \\
\bottomrule
\end{xltabular}
}

\textbf{Gotcha: narrowing the access modifier when overriding is a compile error.}

\begin{codebox}
\begin{Verbatim}[commandchars=\\\{\}]
\PY{k+kd}{public}\PY{+w}{ }\PY{k+kd}{class} \PY{n+nc}{Vehicle}\PY{+w}{ }\PY{p}{\PYZob{}}
\PY{+w}{    }\PY{k+kd}{public}\PY{+w}{ }\PY{k+kt}{void}\PY{+w}{ }\PY{n+nf}{start}\PY{p}{(}\PY{p}{)}\PY{+w}{ }\PY{p}{\PYZob{}}
\PY{+w}{        }\PY{n}{System}\PY{p}{.}\PY{n+na}{out}\PY{p}{.}\PY{n+na}{println}\PY{p}{(}\PY{l+s}{\PYZdq{}}\PY{l+s}{Vehicle starting}\PY{l+s}{\PYZdq{}}\PY{p}{)}\PY{p}{;}
\PY{+w}{    }\PY{p}{\PYZcb{}}
\PY{p}{\PYZcb{}}

\PY{k+kd}{public}\PY{+w}{ }\PY{k+kd}{class} \PY{n+nc}{Car}\PY{+w}{ }\PY{k+kd}{extends}\PY{+w}{ }\PY{n}{Vehicle}\PY{+w}{ }\PY{p}{\PYZob{}}
\PY{+w}{    }\PY{n+nd}{@Override}
\PY{+w}{    }\PY{k+kd}{protected}\PY{+w}{ }\PY{k+kt}{void}\PY{+w}{ }\PY{n+nf}{start}\PY{p}{(}\PY{p}{)}\PY{+w}{ }\PY{p}{\PYZob{}}\PY{+w}{ }\PY{c+c1}{// Compile error: cannot reduce visibility from public to protected}
\PY{+w}{        }\PY{n}{System}\PY{p}{.}\PY{n+na}{out}\PY{p}{.}\PY{n+na}{println}\PY{p}{(}\PY{l+s}{\PYZdq{}}\PY{l+s}{Car starting}\PY{l+s}{\PYZdq{}}\PY{p}{)}\PY{p}{;}
\PY{+w}{    }\PY{p}{\PYZcb{}}
\PY{p}{\PYZcb{}}
\end{Verbatim}
\end{codebox}

\textbf{Gotcha: \code{@\allowbreak{}Override} catches signature mistakes early.}

\begin{codebox}
\begin{Verbatim}[commandchars=\\\{\}]
\PY{k+kd}{public}\PY{+w}{ }\PY{k+kd}{class} \PY{n+nc}{Shape}\PY{+w}{ }\PY{p}{\PYZob{}}
\PY{+w}{    }\PY{k+kd}{public}\PY{+w}{ }\PY{k+kt}{double}\PY{+w}{ }\PY{n+nf}{area}\PY{p}{(}\PY{p}{)}\PY{+w}{ }\PY{p}{\PYZob{}}\PY{+w}{ }\PY{k}{return}\PY{+w}{ }\PY{l+m+mf}{0.0}\PY{p}{;}\PY{+w}{ }\PY{p}{\PYZcb{}}
\PY{p}{\PYZcb{}}

\PY{k+kd}{public}\PY{+w}{ }\PY{k+kd}{class} \PY{n+nc}{Circle}\PY{+w}{ }\PY{k+kd}{extends}\PY{+w}{ }\PY{n}{Shape}\PY{+w}{ }\PY{p}{\PYZob{}}
\PY{+w}{    }\PY{c+c1}{// Typo: overload instead of override, silently creates a NEW method}
\PY{+w}{    }\PY{k+kd}{public}\PY{+w}{ }\PY{k+kt}{double}\PY{+w}{ }\PY{n+nf}{area}\PY{p}{(}\PY{k+kt}{double}\PY{+w}{ }\PY{n}{scale}\PY{p}{)}\PY{+w}{ }\PY{p}{\PYZob{}}
\PY{+w}{        }\PY{k}{return}\PY{+w}{ }\PY{l+m+mf}{3.14}\PY{+w}{ }\PY{o}{*}\PY{+w}{ }\PY{n}{scale}\PY{p}{;}
\PY{+w}{    }\PY{p}{\PYZcb{}}
\PY{p}{\PYZcb{}}
\PY{c+c1}{// Without @Override, this compiles but never overrides Shape.area().}
\PY{c+c1}{// Adding @Override to area(double scale) would immediately fail to compile,}
\PY{c+c1}{// revealing the mistake.}
\end{Verbatim}
\end{codebox}

\section[Access Modifiers]{Access Modifiers}
\label{h:2:access-modifiers}
Access modifiers control which code can see a class, field, method, or constructor. Java has four levels:

{\tablefont
\begin{xltabular}{\linewidth}{@{}L{1.703}L{0.549}L{0.659}L{1.538}L{0.549}@{}}
\toprule
\rowcolor{tablehdr}
\thd{Modifier} & \thd{Same class} & \thd{Same package} & \thd{Subclass (different package)} & \thd{Everywhere} \\
\midrule
\endhead
\code{private} & Yes & No & No & No \\
(no modifier / package-private) & Yes & Yes & No & No \\
\code{protected} & Yes & Yes & Yes & No \\
\code{public} & Yes & Yes & Yes & Yes \\
\bottomrule
\end{xltabular}
}

\begin{codebox}
\begin{Verbatim}[commandchars=\\\{\}]
\PY{k+kn}{package}\PY{+w}{ }\PY{n+nn}{com.bank.account}\PY{p}{;}

\PY{k+kd}{public}\PY{+w}{ }\PY{k+kd}{class} \PY{n+nc}{Account}\PY{+w}{ }\PY{p}{\PYZob{}}
\PY{+w}{    }\PY{k+kd}{private}\PY{+w}{ }\PY{k+kt}{double}\PY{+w}{ }\PY{n}{balance}\PY{p}{;}\PY{+w}{         }\PY{c+c1}{// only inside Account itself}
\PY{+w}{    }\PY{n}{String}\PY{+w}{ }\PY{n}{accountType}\PY{p}{;}\PY{+w}{             }\PY{c+c1}{// package\PYZhy{}private: visible to classes in com.bank.account}
\PY{+w}{    }\PY{k+kd}{protected}\PY{+w}{ }\PY{n}{String}\PY{+w}{ }\PY{n}{branchCode}\PY{p}{;}\PY{+w}{    }\PY{c+c1}{// visible in package + subclasses anywhere}
\PY{+w}{    }\PY{k+kd}{public}\PY{+w}{ }\PY{n}{String}\PY{+w}{ }\PY{n}{ownerName}\PY{p}{;}\PY{+w}{        }\PY{c+c1}{// visible everywhere}
\PY{p}{\PYZcb{}}
\end{Verbatim}
\end{codebox}

\begin{codebox}
\begin{Verbatim}[commandchars=\\\{\}]
\PY{k+kn}{package}\PY{+w}{ }\PY{n+nn}{com.bank.account}\PY{p}{;}

\PY{k+kd}{public}\PY{+w}{ }\PY{k+kd}{class} \PY{n+nc}{SavingsAccount}\PY{+w}{ }\PY{k+kd}{extends}\PY{+w}{ }\PY{n}{Account}\PY{+w}{ }\PY{p}{\PYZob{}}
\PY{+w}{    }\PY{k+kd}{public}\PY{+w}{ }\PY{k+kt}{void}\PY{+w}{ }\PY{n+nf}{printBranch}\PY{p}{(}\PY{p}{)}\PY{+w}{ }\PY{p}{\PYZob{}}
\PY{+w}{        }\PY{n}{System}\PY{p}{.}\PY{n+na}{out}\PY{p}{.}\PY{n+na}{println}\PY{p}{(}\PY{n}{branchCode}\PY{p}{)}\PY{p}{;}\PY{+w}{ }\PY{c+c1}{// OK: protected, and this is a subclass}
\PY{+w}{        }\PY{c+c1}{// System.out.println(balance); // Compile error: balance is private to Account}
\PY{+w}{    }\PY{p}{\PYZcb{}}
\PY{p}{\PYZcb{}}
\end{Verbatim}
\end{codebox}

\begin{codebox}
\begin{Verbatim}[commandchars=\\\{\}]
\PY{k+kn}{package}\PY{+w}{ }\PY{n+nn}{com.bank.reporting}\PY{p}{;}

\PY{k+kd}{public}\PY{+w}{ }\PY{k+kd}{class} \PY{n+nc}{ReportGenerator}\PY{+w}{ }\PY{p}{\PYZob{}}
\PY{+w}{    }\PY{k+kd}{public}\PY{+w}{ }\PY{k+kt}{void}\PY{+w}{ }\PY{n+nf}{printOwner}\PY{p}{(}\PY{n}{Account}\PY{+w}{ }\PY{n}{account}\PY{p}{)}\PY{+w}{ }\PY{p}{\PYZob{}}
\PY{+w}{        }\PY{n}{System}\PY{p}{.}\PY{n+na}{out}\PY{p}{.}\PY{n+na}{println}\PY{p}{(}\PY{n}{account}\PY{p}{.}\PY{n+na}{ownerName}\PY{p}{)}\PY{p}{;}\PY{+w}{ }\PY{c+c1}{// OK: public}
\PY{+w}{        }\PY{c+c1}{// System.out.println(account.branchCode); // Compile error: different package, not a subclass}
\PY{+w}{    }\PY{p}{\PYZcb{}}
\PY{p}{\PYZcb{}}
\end{Verbatim}
\end{codebox}

\textbf{Gotcha: \code{protected} in a different package is only accessible through inheritance, not through an arbitrary reference.}

\begin{codebox}
\begin{Verbatim}[commandchars=\\\{\}]
\PY{k+kn}{package}\PY{+w}{ }\PY{n+nn}{com.bank.reporting}\PY{p}{;}

\PY{k+kn}{import}\PY{+w}{ }\PY{n+nn}{com.bank.account.Account}\PY{p}{;}

\PY{k+kd}{public}\PY{+w}{ }\PY{k+kd}{class} \PY{n+nc}{AuditTool}\PY{+w}{ }\PY{k+kd}{extends}\PY{+w}{ }\PY{n}{Account}\PY{+w}{ }\PY{p}{\PYZob{}}
\PY{+w}{    }\PY{k+kd}{public}\PY{+w}{ }\PY{k+kt}{void}\PY{+w}{ }\PY{n+nf}{check}\PY{p}{(}\PY{n}{Account}\PY{+w}{ }\PY{n}{other}\PY{p}{)}\PY{+w}{ }\PY{p}{\PYZob{}}
\PY{+w}{        }\PY{k}{this}\PY{p}{.}\PY{n+na}{branchCode}\PY{+w}{ }\PY{o}{=}\PY{+w}{ }\PY{l+s}{\PYZdq{}}\PY{l+s}{B01}\PY{l+s}{\PYZdq{}}\PY{p}{;}\PY{+w}{        }\PY{c+c1}{// OK: accessing through \PYZdq{}this\PYZdq{} (the subclass instance)}
\PY{+w}{        }\PY{c+c1}{// other.branchCode = \PYZdq{}B02\PYZdq{};    // Compile error: accessing through another Account reference}
\PY{+w}{    }\PY{p}{\PYZcb{}}
\PY{p}{\PYZcb{}}
\end{Verbatim}
\end{codebox}

The rule of thumb code reviewers apply: \textbf{use the most restrictive modifier that still works.} Start with \code{private}, widen only when there’s a real need. Overly public fields and methods make future refactoring risky because you don’t know who depends on them.

\section[Static Members]{Static Members}
\label{h:2:static-members}
\code{static} members belong to the \textbf{class itself}, not to any particular object. There is exactly one copy, shared by all instances.

\begin{codebox}
\begin{Verbatim}[commandchars=\\\{\}]
\PY{k+kd}{public}\PY{+w}{ }\PY{k+kd}{class} \PY{n+nc}{Account}\PY{+w}{ }\PY{p}{\PYZob{}}
\PY{+w}{    }\PY{k+kd}{private}\PY{+w}{ }\PY{k+kd}{static}\PY{+w}{ }\PY{k+kt}{int}\PY{+w}{ }\PY{n}{totalAccountsCreated}\PY{+w}{ }\PY{o}{=}\PY{+w}{ }\PY{l+m+mi}{0}\PY{p}{;}\PY{+w}{ }\PY{c+c1}{// shared across all Account objects}
\PY{+w}{    }\PY{k+kd}{private}\PY{+w}{ }\PY{k+kd}{final}\PY{+w}{ }\PY{n}{String}\PY{+w}{ }\PY{n}{ownerName}\PY{p}{;}

\PY{+w}{    }\PY{k+kd}{public}\PY{+w}{ }\PY{n+nf}{Account}\PY{p}{(}\PY{n}{String}\PY{+w}{ }\PY{n}{ownerName}\PY{p}{)}\PY{+w}{ }\PY{p}{\PYZob{}}
\PY{+w}{        }\PY{k}{this}\PY{p}{.}\PY{n+na}{ownerName}\PY{+w}{ }\PY{o}{=}\PY{+w}{ }\PY{n}{ownerName}\PY{p}{;}
\PY{+w}{        }\PY{n}{totalAccountsCreated}\PY{o}{+}\PY{o}{+}\PY{p}{;}
\PY{+w}{    }\PY{p}{\PYZcb{}}

\PY{+w}{    }\PY{k+kd}{public}\PY{+w}{ }\PY{k+kd}{static}\PY{+w}{ }\PY{k+kt}{int}\PY{+w}{ }\PY{n+nf}{getTotalAccountsCreated}\PY{p}{(}\PY{p}{)}\PY{+w}{ }\PY{p}{\PYZob{}}
\PY{+w}{        }\PY{k}{return}\PY{+w}{ }\PY{n}{totalAccountsCreated}\PY{p}{;}
\PY{+w}{    }\PY{p}{\PYZcb{}}
\PY{p}{\PYZcb{}}

\PY{k}{new}\PY{+w}{ }\PY{n}{Account}\PY{p}{(}\PY{l+s}{\PYZdq{}}\PY{l+s}{Alice}\PY{l+s}{\PYZdq{}}\PY{p}{)}\PY{p}{;}
\PY{k}{new}\PY{+w}{ }\PY{n}{Account}\PY{p}{(}\PY{l+s}{\PYZdq{}}\PY{l+s}{Bob}\PY{l+s}{\PYZdq{}}\PY{p}{)}\PY{p}{;}
\PY{k}{new}\PY{+w}{ }\PY{n}{Account}\PY{p}{(}\PY{l+s}{\PYZdq{}}\PY{l+s}{Carol}\PY{l+s}{\PYZdq{}}\PY{p}{)}\PY{p}{;}

\PY{n}{System}\PY{p}{.}\PY{n+na}{out}\PY{p}{.}\PY{n+na}{println}\PY{p}{(}\PY{n}{Account}\PY{p}{.}\PY{n+na}{getTotalAccountsCreated}\PY{p}{(}\PY{p}{)}\PY{p}{)}\PY{p}{;}\PY{+w}{ }\PY{c+c1}{// 3}
\end{Verbatim}
\end{codebox}

Common uses of \code{static}:

\begin{itemize}
\item 
\textbf{Static fields}: shared/global state per class, like counters or constants.

\item 
\textbf{Static methods}: utility logic that doesn’t need object state, e.g. \code{Math.\allowbreak{}sqrt(\allowbreak{}x)}.

\item 
\textbf{Static nested classes}: see \hyperref[h:2:nested-inner-local-and-anonymous-classes]{Nested, Inner, Local, and Anonymous Classes}.

\item 
\textbf{Static initializer blocks}: run once when the class is first loaded.

\end{itemize}

\begin{codebox}
\begin{Verbatim}[commandchars=\\\{\}]
\PY{k+kd}{public}\PY{+w}{ }\PY{k+kd}{class} \PY{n+nc}{Configuration}\PY{+w}{ }\PY{p}{\PYZob{}}
\PY{+w}{    }\PY{k+kd}{public}\PY{+w}{ }\PY{k+kd}{static}\PY{+w}{ }\PY{k+kd}{final}\PY{+w}{ }\PY{n}{String}\PY{+w}{ }\PY{n}{ENVIRONMENT}\PY{p}{;}

\PY{+w}{    }\PY{k+kd}{static}\PY{+w}{ }\PY{p}{\PYZob{}}
\PY{+w}{        }\PY{c+c1}{// Runs once, when the class is loaded — good for expensive one\PYZhy{}time setup}
\PY{+w}{        }\PY{n}{ENVIRONMENT}\PY{+w}{ }\PY{o}{=}\PY{+w}{ }\PY{n}{System}\PY{p}{.}\PY{n+na}{getenv}\PY{p}{(}\PY{p}{)}\PY{p}{.}\PY{n+na}{getOrDefault}\PY{p}{(}\PY{l+s}{\PYZdq{}}\PY{l+s}{APP\PYZus{}ENV}\PY{l+s}{\PYZdq{}}\PY{p}{,}\PY{+w}{ }\PY{l+s}{\PYZdq{}}\PY{l+s}{development}\PY{l+s}{\PYZdq{}}\PY{p}{)}\PY{p}{;}
\PY{+w}{        }\PY{n}{System}\PY{p}{.}\PY{n+na}{out}\PY{p}{.}\PY{n+na}{println}\PY{p}{(}\PY{l+s}{\PYZdq{}}\PY{l+s}{Configuration loaded for: }\PY{l+s}{\PYZdq{}}\PY{+w}{ }\PY{o}{+}\PY{+w}{ }\PY{n}{ENVIRONMENT}\PY{p}{)}\PY{p}{;}
\PY{+w}{    }\PY{p}{\PYZcb{}}
\PY{p}{\PYZcb{}}
\end{Verbatim}
\end{codebox}

A \code{static} method cannot directly access instance (non-static) fields or methods, because there’s no specific object to work on.

\begin{codebox}
\begin{Verbatim}[commandchars=\\\{\}]
\PY{k+kd}{public}\PY{+w}{ }\PY{k+kd}{class} \PY{n+nc}{Account}\PY{+w}{ }\PY{p}{\PYZob{}}
\PY{+w}{    }\PY{k+kd}{private}\PY{+w}{ }\PY{k+kt}{double}\PY{+w}{ }\PY{n}{balance}\PY{+w}{ }\PY{o}{=}\PY{+w}{ }\PY{l+m+mi}{100}\PY{p}{;}

\PY{+w}{    }\PY{k+kd}{public}\PY{+w}{ }\PY{k+kd}{static}\PY{+w}{ }\PY{k+kt}{void}\PY{+w}{ }\PY{n+nf}{printBalance}\PY{p}{(}\PY{p}{)}\PY{+w}{ }\PY{p}{\PYZob{}}
\PY{+w}{        }\PY{c+c1}{// System.out.println(balance); // Compile error: no instance to read balance from}
\PY{+w}{    }\PY{p}{\PYZcb{}}
\PY{p}{\PYZcb{}}
\end{Verbatim}
\end{codebox}

\textbf{Gotcha: static methods are inherited but hidden, not overridden} (also shown earlier in \hyperref[h:2:polymorphism]{Polymorphism}). Calling a static method through an instance reference is legal but misleading — code reviewers should flag it.

\begin{codebox}
\begin{Verbatim}[commandchars=\\\{\}]
\PY{k+kd}{public}\PY{+w}{ }\PY{k+kd}{class} \PY{n+nc}{ReportUtil}\PY{+w}{ }\PY{p}{\PYZob{}}
\PY{+w}{    }\PY{k+kd}{public}\PY{+w}{ }\PY{k+kd}{static}\PY{+w}{ }\PY{k+kt}{void}\PY{+w}{ }\PY{n+nf}{printHeader}\PY{p}{(}\PY{p}{)}\PY{+w}{ }\PY{p}{\PYZob{}}
\PY{+w}{        }\PY{n}{System}\PY{p}{.}\PY{n+na}{out}\PY{p}{.}\PY{n+na}{println}\PY{p}{(}\PY{l+s}{\PYZdq{}}\PY{l+s}{=== Report ===}\PY{l+s}{\PYZdq{}}\PY{p}{)}\PY{p}{;}
\PY{+w}{    }\PY{p}{\PYZcb{}}
\PY{p}{\PYZcb{}}

\PY{n}{ReportUtil}\PY{+w}{ }\PY{n}{util}\PY{+w}{ }\PY{o}{=}\PY{+w}{ }\PY{k}{new}\PY{+w}{ }\PY{n}{ReportUtil}\PY{p}{(}\PY{p}{)}\PY{p}{;}
\PY{n}{util}\PY{p}{.}\PY{n+na}{printHeader}\PY{p}{(}\PY{p}{)}\PY{p}{;}\PY{+w}{ }\PY{c+c1}{// Works, but should be called as ReportUtil.printHeader()}
\end{Verbatim}
\end{codebox}

\textbf{Gotcha: mutable static state is a shared, hidden dependency} — dangerous in multi-threaded code and hard to test because it persists between tests.

\begin{codebox}
\begin{Verbatim}[commandchars=\\\{\}]
\PY{k+kd}{public}\PY{+w}{ }\PY{k+kd}{class} \PY{n+nc}{IdGenerator}\PY{+w}{ }\PY{p}{\PYZob{}}
\PY{+w}{    }\PY{k+kd}{private}\PY{+w}{ }\PY{k+kd}{static}\PY{+w}{ }\PY{k+kt}{int}\PY{+w}{ }\PY{n}{nextId}\PY{+w}{ }\PY{o}{=}\PY{+w}{ }\PY{l+m+mi}{1}\PY{p}{;}\PY{+w}{ }\PY{c+c1}{// shared mutable state!}

\PY{+w}{    }\PY{k+kd}{public}\PY{+w}{ }\PY{k+kd}{static}\PY{+w}{ }\PY{k+kt}{int}\PY{+w}{ }\PY{n+nf}{next}\PY{p}{(}\PY{p}{)}\PY{+w}{ }\PY{p}{\PYZob{}}
\PY{+w}{        }\PY{k}{return}\PY{+w}{ }\PY{n}{nextId}\PY{o}{+}\PY{o}{+}\PY{p}{;}\PY{+w}{ }\PY{c+c1}{// not thread\PYZhy{}safe: race condition under concurrent access}
\PY{+w}{    }\PY{p}{\PYZcb{}}
\PY{p}{\PYZcb{}}
\end{Verbatim}
\end{codebox}

\section[Final Keyword]{Final Keyword}
\label{h:2:final-keyword}
\code{final} means “cannot be changed again,” but what exactly it locks down depends on where it’s applied:

{\tablefont
\begin{xltabular}{\linewidth}{@{}L{0.548}L{1.452}@{}}
\toprule
\rowcolor{tablehdr}
\thd{Applied to} & \thd{Meaning} \\
\midrule
\endhead
\code{final} variable/field & Value (or object reference) can only be assigned once \\
\code{final} method & Cannot be overridden by subclasses \\
\code{final} class & Cannot be extended (no subclasses at all) \\
\bottomrule
\end{xltabular}
}

\begin{codebox}
\begin{Verbatim}[commandchars=\\\{\}]
\PY{k+kd}{public}\PY{+w}{ }\PY{k+kd}{final}\PY{+w}{ }\PY{k+kd}{class} \PY{n+nc}{ImmutablePoint}\PY{+w}{ }\PY{p}{\PYZob{}}\PY{+w}{ }\PY{c+c1}{// cannot be subclassed}
\PY{+w}{    }\PY{k+kd}{private}\PY{+w}{ }\PY{k+kd}{final}\PY{+w}{ }\PY{k+kt}{int}\PY{+w}{ }\PY{n}{x}\PY{p}{;}\PY{+w}{ }\PY{c+c1}{// must be assigned exactly once (constructor or declaration)}
\PY{+w}{    }\PY{k+kd}{private}\PY{+w}{ }\PY{k+kd}{final}\PY{+w}{ }\PY{k+kt}{int}\PY{+w}{ }\PY{n}{y}\PY{p}{;}

\PY{+w}{    }\PY{k+kd}{public}\PY{+w}{ }\PY{n+nf}{ImmutablePoint}\PY{p}{(}\PY{k+kt}{int}\PY{+w}{ }\PY{n}{x}\PY{p}{,}\PY{+w}{ }\PY{k+kt}{int}\PY{+w}{ }\PY{n}{y}\PY{p}{)}\PY{+w}{ }\PY{p}{\PYZob{}}
\PY{+w}{        }\PY{k}{this}\PY{p}{.}\PY{n+na}{x}\PY{+w}{ }\PY{o}{=}\PY{+w}{ }\PY{n}{x}\PY{p}{;}
\PY{+w}{        }\PY{k}{this}\PY{p}{.}\PY{n+na}{y}\PY{+w}{ }\PY{o}{=}\PY{+w}{ }\PY{n}{y}\PY{p}{;}
\PY{+w}{    }\PY{p}{\PYZcb{}}

\PY{+w}{    }\PY{k+kd}{public}\PY{+w}{ }\PY{k+kd}{final}\PY{+w}{ }\PY{k+kt}{int}\PY{+w}{ }\PY{n+nf}{getX}\PY{p}{(}\PY{p}{)}\PY{+w}{ }\PY{p}{\PYZob{}}\PY{+w}{ }\PY{c+c1}{// cannot be overridden, even though the class already forbids subclassing}
\PY{+w}{        }\PY{k}{return}\PY{+w}{ }\PY{n}{x}\PY{p}{;}
\PY{+w}{    }\PY{p}{\PYZcb{}}
\PY{p}{\PYZcb{}}
\end{Verbatim}
\end{codebox}

\begin{codebox}
\begin{Verbatim}[commandchars=\\\{\}]
\PY{k+kd}{public}\PY{+w}{ }\PY{k+kd}{class} \PY{n+nc}{Vehicle}\PY{+w}{ }\PY{p}{\PYZob{}}
\PY{+w}{    }\PY{k+kd}{protected}\PY{+w}{ }\PY{k+kd}{final}\PY{+w}{ }\PY{k+kt}{int}\PY{+w}{ }\PY{n}{maxSpeed}\PY{p}{;}

\PY{+w}{    }\PY{k+kd}{public}\PY{+w}{ }\PY{n+nf}{Vehicle}\PY{p}{(}\PY{k+kt}{int}\PY{+w}{ }\PY{n}{maxSpeed}\PY{p}{)}\PY{+w}{ }\PY{p}{\PYZob{}}
\PY{+w}{        }\PY{k}{this}\PY{p}{.}\PY{n+na}{maxSpeed}\PY{+w}{ }\PY{o}{=}\PY{+w}{ }\PY{n}{maxSpeed}\PY{p}{;}\PY{+w}{ }\PY{c+c1}{// OK: first (and only) assignment}
\PY{+w}{    }\PY{p}{\PYZcb{}}

\PY{+w}{    }\PY{k+kd}{public}\PY{+w}{ }\PY{k+kt}{void}\PY{+w}{ }\PY{n+nf}{boost}\PY{p}{(}\PY{p}{)}\PY{+w}{ }\PY{p}{\PYZob{}}
\PY{+w}{        }\PY{c+c1}{// this.maxSpeed = maxSpeed + 10; // Compile error: final field can\PYZsq{}t be reassigned}
\PY{+w}{    }\PY{p}{\PYZcb{}}
\PY{p}{\PYZcb{}}
\end{Verbatim}
\end{codebox}

\begin{codebox}
\begin{Verbatim}[commandchars=\\\{\}]
\PY{k+kd}{public}\PY{+w}{ }\PY{k+kd}{class} \PY{n+nc}{Vehicle}\PY{+w}{ }\PY{p}{\PYZob{}}
\PY{+w}{    }\PY{k+kd}{public}\PY{+w}{ }\PY{k+kd}{final}\PY{+w}{ }\PY{k+kt}{void}\PY{+w}{ }\PY{n+nf}{identify}\PY{p}{(}\PY{p}{)}\PY{+w}{ }\PY{p}{\PYZob{}}
\PY{+w}{        }\PY{n}{System}\PY{p}{.}\PY{n+na}{out}\PY{p}{.}\PY{n+na}{println}\PY{p}{(}\PY{l+s}{\PYZdq{}}\PY{l+s}{I am a vehicle}\PY{l+s}{\PYZdq{}}\PY{p}{)}\PY{p}{;}
\PY{+w}{    }\PY{p}{\PYZcb{}}
\PY{p}{\PYZcb{}}

\PY{k+kd}{public}\PY{+w}{ }\PY{k+kd}{class} \PY{n+nc}{Car}\PY{+w}{ }\PY{k+kd}{extends}\PY{+w}{ }\PY{n}{Vehicle}\PY{+w}{ }\PY{p}{\PYZob{}}
\PY{+w}{    }\PY{n+nd}{@Override}
\PY{+w}{    }\PY{k+kd}{public}\PY{+w}{ }\PY{k+kt}{void}\PY{+w}{ }\PY{n+nf}{identify}\PY{p}{(}\PY{p}{)}\PY{+w}{ }\PY{p}{\PYZob{}}\PY{+w}{ }\PY{c+c1}{// Compile error: cannot override a final method}
\PY{+w}{        }\PY{n}{System}\PY{p}{.}\PY{n+na}{out}\PY{p}{.}\PY{n+na}{println}\PY{p}{(}\PY{l+s}{\PYZdq{}}\PY{l+s}{I am a car}\PY{l+s}{\PYZdq{}}\PY{p}{)}\PY{p}{;}
\PY{+w}{    }\PY{p}{\PYZcb{}}
\PY{p}{\PYZcb{}}
\end{Verbatim}
\end{codebox}

\textbf{Gotcha: \code{final} on a reference only locks the reference, not the object’s contents.} A \code{final} field can still point to a mutable object whose internal state changes.

\begin{codebox}
\begin{Verbatim}[commandchars=\\\{\}]
\PY{k+kd}{public}\PY{+w}{ }\PY{k+kd}{class} \PY{n+nc}{Order}\PY{+w}{ }\PY{p}{\PYZob{}}
\PY{+w}{    }\PY{k+kd}{private}\PY{+w}{ }\PY{k+kd}{final}\PY{+w}{ }\PY{n}{List}\PY{o}{\PYZlt{}}\PY{n}{String}\PY{o}{\PYZgt{}}\PY{+w}{ }\PY{n}{items}\PY{+w}{ }\PY{o}{=}\PY{+w}{ }\PY{k}{new}\PY{+w}{ }\PY{n}{ArrayList}\PY{o}{\PYZlt{}}\PY{o}{\PYZgt{}}\PY{p}{(}\PY{p}{)}\PY{p}{;}

\PY{+w}{    }\PY{k+kd}{public}\PY{+w}{ }\PY{k+kt}{void}\PY{+w}{ }\PY{n+nf}{addItem}\PY{p}{(}\PY{n}{String}\PY{+w}{ }\PY{n}{item}\PY{p}{)}\PY{+w}{ }\PY{p}{\PYZob{}}
\PY{+w}{        }\PY{n}{items}\PY{p}{.}\PY{n+na}{add}\PY{p}{(}\PY{n}{item}\PY{p}{)}\PY{p}{;}\PY{+w}{ }\PY{c+c1}{// legal! We\PYZsq{}re mutating the List, not reassigning \PYZdq{}items\PYZdq{}}
\PY{+w}{    }\PY{p}{\PYZcb{}}

\PY{+w}{    }\PY{k+kd}{public}\PY{+w}{ }\PY{k+kt}{void}\PY{+w}{ }\PY{n+nf}{replaceList}\PY{p}{(}\PY{p}{)}\PY{+w}{ }\PY{p}{\PYZob{}}
\PY{+w}{        }\PY{c+c1}{// items = new ArrayList\PYZlt{}\PYZgt{}(); // Compile error: can\PYZsq{}t reassign a final field}
\PY{+w}{    }\PY{p}{\PYZcb{}}
\PY{p}{\PYZcb{}}
\end{Verbatim}
\end{codebox}

This is a common interview trap: \code{final} does \textbf{not} mean immutable. True immutability requires the class itself to prevent internal mutation too (see the “Immutability” topic in the “Object Class \& Common APIs” chapter).

\section[This and Super]{This and Super}
\label{h:2:this-and-super}
\code{this} refers to the \textbf{current object} — the one whose method or constructor is currently executing. \code{super} refers to the \textbf{immediate superclass}, used to call its constructor or access members it defines.

\begin{codebox}
\begin{Verbatim}[commandchars=\\\{\}]
\PY{k+kd}{public}\PY{+w}{ }\PY{k+kd}{class} \PY{n+nc}{Vehicle}\PY{+w}{ }\PY{p}{\PYZob{}}
\PY{+w}{    }\PY{k+kd}{protected}\PY{+w}{ }\PY{n}{String}\PY{+w}{ }\PY{n}{plateNumber}\PY{p}{;}

\PY{+w}{    }\PY{k+kd}{public}\PY{+w}{ }\PY{n+nf}{Vehicle}\PY{p}{(}\PY{n}{String}\PY{+w}{ }\PY{n}{plateNumber}\PY{p}{)}\PY{+w}{ }\PY{p}{\PYZob{}}
\PY{+w}{        }\PY{k}{this}\PY{p}{.}\PY{n+na}{plateNumber}\PY{+w}{ }\PY{o}{=}\PY{+w}{ }\PY{n}{plateNumber}\PY{p}{;}
\PY{+w}{    }\PY{p}{\PYZcb{}}

\PY{+w}{    }\PY{k+kd}{public}\PY{+w}{ }\PY{n}{String}\PY{+w}{ }\PY{n+nf}{describe}\PY{p}{(}\PY{p}{)}\PY{+w}{ }\PY{p}{\PYZob{}}
\PY{+w}{        }\PY{k}{return}\PY{+w}{ }\PY{l+s}{\PYZdq{}}\PY{l+s}{Vehicle }\PY{l+s}{\PYZdq{}}\PY{+w}{ }\PY{o}{+}\PY{+w}{ }\PY{n}{plateNumber}\PY{p}{;}
\PY{+w}{    }\PY{p}{\PYZcb{}}
\PY{p}{\PYZcb{}}

\PY{k+kd}{public}\PY{+w}{ }\PY{k+kd}{class} \PY{n+nc}{Car}\PY{+w}{ }\PY{k+kd}{extends}\PY{+w}{ }\PY{n}{Vehicle}\PY{+w}{ }\PY{p}{\PYZob{}}
\PY{+w}{    }\PY{k+kd}{private}\PY{+w}{ }\PY{n}{String}\PY{+w}{ }\PY{n}{plateNumber}\PY{p}{;}\PY{+w}{ }\PY{c+c1}{// shadows Vehicle\PYZsq{}s field — usually a bad idea, shown here for illustration}

\PY{+w}{    }\PY{k+kd}{public}\PY{+w}{ }\PY{n+nf}{Car}\PY{p}{(}\PY{n}{String}\PY{+w}{ }\PY{n}{plateNumber}\PY{p}{)}\PY{+w}{ }\PY{p}{\PYZob{}}
\PY{+w}{        }\PY{k+kd}{super}\PY{p}{(}\PY{n}{plateNumber}\PY{p}{)}\PY{p}{;}\PY{+w}{       }\PY{c+c1}{// calls Vehicle\PYZsq{}s constructor}
\PY{+w}{        }\PY{k}{this}\PY{p}{.}\PY{n+na}{plateNumber}\PY{+w}{ }\PY{o}{=}\PY{+w}{ }\PY{n}{plateNumber}\PY{p}{;}\PY{+w}{ }\PY{c+c1}{// \PYZdq{}this\PYZdq{} disambiguates the field from the parameter}
\PY{+w}{    }\PY{p}{\PYZcb{}}

\PY{+w}{    }\PY{n+nd}{@Override}
\PY{+w}{    }\PY{k+kd}{public}\PY{+w}{ }\PY{n}{String}\PY{+w}{ }\PY{n+nf}{describe}\PY{p}{(}\PY{p}{)}\PY{+w}{ }\PY{p}{\PYZob{}}
\PY{+w}{        }\PY{k}{return}\PY{+w}{ }\PY{k+kd}{super}\PY{p}{.}\PY{n+na}{describe}\PY{p}{(}\PY{p}{)}\PY{+w}{ }\PY{o}{+}\PY{+w}{ }\PY{l+s}{\PYZdq{}}\PY{l+s}{ (Car version, plate=}\PY{l+s}{\PYZdq{}}\PY{+w}{ }\PY{o}{+}\PY{+w}{ }\PY{k}{this}\PY{p}{.}\PY{n+na}{plateNumber}\PY{+w}{ }\PY{o}{+}\PY{+w}{ }\PY{l+s}{\PYZdq{}}\PY{l+s}{)}\PY{l+s}{\PYZdq{}}\PY{p}{;}
\PY{+w}{    }\PY{p}{\PYZcb{}}
\PY{p}{\PYZcb{}}

\PY{n}{Car}\PY{+w}{ }\PY{n}{car}\PY{+w}{ }\PY{o}{=}\PY{+w}{ }\PY{k}{new}\PY{+w}{ }\PY{n}{Car}\PY{p}{(}\PY{l+s}{\PYZdq{}}\PY{l+s}{ABC\PYZhy{}999}\PY{l+s}{\PYZdq{}}\PY{p}{)}\PY{p}{;}
\PY{n}{System}\PY{p}{.}\PY{n+na}{out}\PY{p}{.}\PY{n+na}{println}\PY{p}{(}\PY{n}{car}\PY{p}{.}\PY{n+na}{describe}\PY{p}{(}\PY{p}{)}\PY{p}{)}\PY{p}{;}
\PY{c+c1}{// Output: Vehicle ABC\PYZhy{}999 (Car version, plate=ABC\PYZhy{}999)}
\end{Verbatim}
\end{codebox}

Common uses of \code{this}:

\begin{itemize}
\item 
Disambiguate a field from a constructor/method parameter with the same name (\code{this.\allowbreak{}balance~\allowbreak{}=\allowbreak{}~\allowbreak{}balance;}).

\item 
Call another constructor in the same class (\code{this(...)} — constructor chaining, see \hyperref[h:2:constructors]{Constructors}).

\item 
Pass the current object as an argument (\code{someMethod(\allowbreak{}this)}).

\item 
Return the current object for method chaining (\code{return~\allowbreak{}this;}).

\end{itemize}

\begin{codebox}
\begin{Verbatim}[commandchars=\\\{\}]
\PY{k+kd}{public}\PY{+w}{ }\PY{k+kd}{class} \PY{n+nc}{InvoiceBuilder}\PY{+w}{ }\PY{p}{\PYZob{}}
\PY{+w}{    }\PY{k+kd}{private}\PY{+w}{ }\PY{k+kt}{double}\PY{+w}{ }\PY{n}{amount}\PY{p}{;}
\PY{+w}{    }\PY{k+kd}{private}\PY{+w}{ }\PY{n}{String}\PY{+w}{ }\PY{n}{currency}\PY{p}{;}

\PY{+w}{    }\PY{k+kd}{public}\PY{+w}{ }\PY{n}{InvoiceBuilder}\PY{+w}{ }\PY{n+nf}{withAmount}\PY{p}{(}\PY{k+kt}{double}\PY{+w}{ }\PY{n}{amount}\PY{p}{)}\PY{+w}{ }\PY{p}{\PYZob{}}
\PY{+w}{        }\PY{k}{this}\PY{p}{.}\PY{n+na}{amount}\PY{+w}{ }\PY{o}{=}\PY{+w}{ }\PY{n}{amount}\PY{p}{;}
\PY{+w}{        }\PY{k}{return}\PY{+w}{ }\PY{k}{this}\PY{p}{;}\PY{+w}{ }\PY{c+c1}{// enables chaining}
\PY{+w}{    }\PY{p}{\PYZcb{}}

\PY{+w}{    }\PY{k+kd}{public}\PY{+w}{ }\PY{n}{InvoiceBuilder}\PY{+w}{ }\PY{n+nf}{withCurrency}\PY{p}{(}\PY{n}{String}\PY{+w}{ }\PY{n}{currency}\PY{p}{)}\PY{+w}{ }\PY{p}{\PYZob{}}
\PY{+w}{        }\PY{k}{this}\PY{p}{.}\PY{n+na}{currency}\PY{+w}{ }\PY{o}{=}\PY{+w}{ }\PY{n}{currency}\PY{p}{;}
\PY{+w}{        }\PY{k}{return}\PY{+w}{ }\PY{k}{this}\PY{p}{;}
\PY{+w}{    }\PY{p}{\PYZcb{}}
\PY{p}{\PYZcb{}}

\PY{n}{InvoiceBuilder}\PY{+w}{ }\PY{n}{builder}\PY{+w}{ }\PY{o}{=}\PY{+w}{ }\PY{k}{new}\PY{+w}{ }\PY{n}{InvoiceBuilder}\PY{p}{(}\PY{p}{)}
\PY{+w}{    }\PY{p}{.}\PY{n+na}{withAmount}\PY{p}{(}\PY{l+m+mf}{99.99}\PY{p}{)}
\PY{+w}{    }\PY{p}{.}\PY{n+na}{withCurrency}\PY{p}{(}\PY{l+s}{\PYZdq{}}\PY{l+s}{EUR}\PY{l+s}{\PYZdq{}}\PY{p}{)}\PY{p}{;}\PY{+w}{ }\PY{c+c1}{// fluent chaining thanks to \PYZdq{}return this;\PYZdq{}}
\end{Verbatim}
\end{codebox}

Common uses of \code{super}:

\begin{itemize}
\item 
Call the parent constructor (\code{super(...)} — must be the first statement in a constructor).

\item 
Call the parent’s version of an overridden method (\code{super.\allowbreak{}describe()}).

\item 
Access a parent field that’s shadowed by a subclass field.

\end{itemize}

\textbf{Gotcha: \code{super} only reaches one level up.} There’s no \code{super.\allowbreak{}super.\allowbreak{}method()} in Java — if \code{Car~\allowbreak{}extends~\allowbreak{}Vehicle~\allowbreak{}extends~\allowbreak{}Machine}, \code{Car} can call \code{super.\allowbreak{}describe()} to reach \code{Vehicle}’s version, but not \code{Machine}’s directly (unless \code{Vehicle} doesn’t override it, in which case \code{super.\allowbreak{}describe()} naturally falls through to \code{Machine}’s implementation).

\section[Nested, Inner, Local, and Anonymous Classes]{Nested, Inner, Local, and Anonymous Classes}
\label{h:2:nested-inner-local-and-anonymous-classes}
Java allows classes defined inside other classes or even inside methods. They’re grouped into four kinds:

{\tablefont
\begin{xltabular}{\linewidth}{@{}L{0.445}L{1.485}L{0.962}L{1.108}@{}}
\toprule
\rowcolor{tablehdr}
\thd{Kind} & \thd{Defined} & \thd{Needs outer instance?} & \thd{Typical use} \\
\midrule
\endhead
Static nested class & Inside a class, marked \code{static} & No & Grouping a helper class that doesn’t need outer state \\
Inner class & Inside a class, not \code{static} & Yes (holds implicit reference to outer object) & Tight coupling with an outer object’s state \\
Local class & Inside a method body & Yes, if the method is an instance method & One-off helper used only within a method \\
Anonymous class & Inline, no name, usually implementing an interface or extending a class & Depends on context & Quick one-time implementation, e.g. a callback \\
\bottomrule
\end{xltabular}
}

\textbf{Static nested class} — behaves like a top-level class, just namespaced inside another for organization. Does not hold a reference to an outer instance.

\begin{codebox}
\begin{Verbatim}[commandchars=\\\{\}]
\PY{k+kd}{public}\PY{+w}{ }\PY{k+kd}{class} \PY{n+nc}{Order}\PY{+w}{ }\PY{p}{\PYZob{}}
\PY{+w}{    }\PY{k+kd}{private}\PY{+w}{ }\PY{k+kd}{final}\PY{+w}{ }\PY{n}{List}\PY{o}{\PYZlt{}}\PY{n}{LineItem}\PY{o}{\PYZgt{}}\PY{+w}{ }\PY{n}{lineItems}\PY{+w}{ }\PY{o}{=}\PY{+w}{ }\PY{k}{new}\PY{+w}{ }\PY{n}{ArrayList}\PY{o}{\PYZlt{}}\PY{o}{\PYZgt{}}\PY{p}{(}\PY{p}{)}\PY{p}{;}

\PY{+w}{    }\PY{k+kd}{public}\PY{+w}{ }\PY{k+kd}{static}\PY{+w}{ }\PY{k+kd}{class} \PY{n+nc}{LineItem}\PY{+w}{ }\PY{p}{\PYZob{}}\PY{+w}{ }\PY{c+c1}{// static nested class}
\PY{+w}{        }\PY{k+kd}{private}\PY{+w}{ }\PY{k+kd}{final}\PY{+w}{ }\PY{n}{String}\PY{+w}{ }\PY{n}{productName}\PY{p}{;}
\PY{+w}{        }\PY{k+kd}{private}\PY{+w}{ }\PY{k+kd}{final}\PY{+w}{ }\PY{k+kt}{int}\PY{+w}{ }\PY{n}{quantity}\PY{p}{;}

\PY{+w}{        }\PY{k+kd}{public}\PY{+w}{ }\PY{n+nf}{LineItem}\PY{p}{(}\PY{n}{String}\PY{+w}{ }\PY{n}{productName}\PY{p}{,}\PY{+w}{ }\PY{k+kt}{int}\PY{+w}{ }\PY{n}{quantity}\PY{p}{)}\PY{+w}{ }\PY{p}{\PYZob{}}
\PY{+w}{            }\PY{k}{this}\PY{p}{.}\PY{n+na}{productName}\PY{+w}{ }\PY{o}{=}\PY{+w}{ }\PY{n}{productName}\PY{p}{;}
\PY{+w}{            }\PY{k}{this}\PY{p}{.}\PY{n+na}{quantity}\PY{+w}{ }\PY{o}{=}\PY{+w}{ }\PY{n}{quantity}\PY{p}{;}
\PY{+w}{        }\PY{p}{\PYZcb{}}
\PY{+w}{    }\PY{p}{\PYZcb{}}

\PY{+w}{    }\PY{k+kd}{public}\PY{+w}{ }\PY{k+kt}{void}\PY{+w}{ }\PY{n+nf}{addLineItem}\PY{p}{(}\PY{n}{String}\PY{+w}{ }\PY{n}{productName}\PY{p}{,}\PY{+w}{ }\PY{k+kt}{int}\PY{+w}{ }\PY{n}{quantity}\PY{p}{)}\PY{+w}{ }\PY{p}{\PYZob{}}
\PY{+w}{        }\PY{n}{lineItems}\PY{p}{.}\PY{n+na}{add}\PY{p}{(}\PY{k}{new}\PY{+w}{ }\PY{n}{LineItem}\PY{p}{(}\PY{n}{productName}\PY{p}{,}\PY{+w}{ }\PY{n}{quantity}\PY{p}{)}\PY{p}{)}\PY{p}{;}
\PY{+w}{    }\PY{p}{\PYZcb{}}
\PY{p}{\PYZcb{}}

\PY{c+c1}{// Created without needing an Order instance:}
\PY{n}{Order}\PY{p}{.}\PY{n+na}{LineItem}\PY{+w}{ }\PY{n}{item}\PY{+w}{ }\PY{o}{=}\PY{+w}{ }\PY{k}{new}\PY{+w}{ }\PY{n}{Order}\PY{p}{.}\PY{n+na}{LineItem}\PY{p}{(}\PY{l+s}{\PYZdq{}}\PY{l+s}{Widget}\PY{l+s}{\PYZdq{}}\PY{p}{,}\PY{+w}{ }\PY{l+m+mi}{3}\PY{p}{)}\PY{p}{;}
\end{Verbatim}
\end{codebox}

\textbf{Inner class (non-static)} — tied to a specific outer object; it can access the outer object’s fields directly, even \code{private} ones.

\begin{codebox}
\begin{Verbatim}[commandchars=\\\{\}]
\PY{k+kd}{public}\PY{+w}{ }\PY{k+kd}{class} \PY{n+nc}{Order}\PY{+w}{ }\PY{p}{\PYZob{}}
\PY{+w}{    }\PY{k+kd}{private}\PY{+w}{ }\PY{k+kt}{double}\PY{+w}{ }\PY{n}{total}\PY{p}{;}
\PY{+w}{    }\PY{k+kd}{private}\PY{+w}{ }\PY{k+kd}{final}\PY{+w}{ }\PY{n}{List}\PY{o}{\PYZlt{}}\PY{n}{String}\PY{o}{\PYZgt{}}\PY{+w}{ }\PY{n}{discountLog}\PY{+w}{ }\PY{o}{=}\PY{+w}{ }\PY{k}{new}\PY{+w}{ }\PY{n}{ArrayList}\PY{o}{\PYZlt{}}\PY{o}{\PYZgt{}}\PY{p}{(}\PY{p}{)}\PY{p}{;}

\PY{+w}{    }\PY{k+kd}{public}\PY{+w}{ }\PY{k+kd}{class} \PY{n+nc}{DiscountApplier}\PY{+w}{ }\PY{p}{\PYZob{}}\PY{+w}{ }\PY{c+c1}{// inner (non\PYZhy{}static) class}
\PY{+w}{        }\PY{k+kd}{public}\PY{+w}{ }\PY{k+kt}{void}\PY{+w}{ }\PY{n+nf}{apply}\PY{p}{(}\PY{k+kt}{double}\PY{+w}{ }\PY{n}{percentage}\PY{p}{)}\PY{+w}{ }\PY{p}{\PYZob{}}
\PY{+w}{            }\PY{n}{total}\PY{+w}{ }\PY{o}{=}\PY{+w}{ }\PY{n}{total}\PY{+w}{ }\PY{o}{*}\PY{+w}{ }\PY{p}{(}\PY{l+m+mi}{1}\PY{+w}{ }\PY{o}{\PYZhy{}}\PY{+w}{ }\PY{n}{percentage}\PY{p}{)}\PY{p}{;}\PY{+w}{ }\PY{c+c1}{// accesses outer field directly}
\PY{+w}{            }\PY{n}{discountLog}\PY{p}{.}\PY{n+na}{add}\PY{p}{(}\PY{l+s}{\PYZdq{}}\PY{l+s}{Applied }\PY{l+s}{\PYZdq{}}\PY{+w}{ }\PY{o}{+}\PY{+w}{ }\PY{p}{(}\PY{n}{percentage}\PY{+w}{ }\PY{o}{*}\PY{+w}{ }\PY{l+m+mi}{100}\PY{p}{)}\PY{+w}{ }\PY{o}{+}\PY{+w}{ }\PY{l+s}{\PYZdq{}}\PY{l+s}{\PYZpc{} discount}\PY{l+s}{\PYZdq{}}\PY{p}{)}\PY{p}{;}
\PY{+w}{        }\PY{p}{\PYZcb{}}
\PY{+w}{    }\PY{p}{\PYZcb{}}
\PY{p}{\PYZcb{}}

\PY{n}{Order}\PY{+w}{ }\PY{n}{order}\PY{+w}{ }\PY{o}{=}\PY{+w}{ }\PY{k}{new}\PY{+w}{ }\PY{n}{Order}\PY{p}{(}\PY{p}{)}\PY{p}{;}
\PY{n}{Order}\PY{p}{.}\PY{n+na}{DiscountApplier}\PY{+w}{ }\PY{n}{applier}\PY{+w}{ }\PY{o}{=}\PY{+w}{ }\PY{n}{order}\PY{p}{.}\PY{n+na}{new}\PY{+w}{ }\PY{n+nf}{DiscountApplier}\PY{p}{(}\PY{p}{)}\PY{p}{;}\PY{+w}{ }\PY{c+c1}{// needs an outer instance}
\PY{n}{applier}\PY{p}{.}\PY{n+na}{apply}\PY{p}{(}\PY{l+m+mf}{0.1}\PY{p}{)}\PY{p}{;}
\end{Verbatim}
\end{codebox}

\textbf{Local class} — declared inside a method, visible only within that method. Useful when a helper type is needed just for one algorithm and shouldn’t leak outside it.

\begin{codebox}
\begin{Verbatim}[commandchars=\\\{\}]
\PY{k+kd}{public}\PY{+w}{ }\PY{k+kd}{class} \PY{n+nc}{ReportGenerator}\PY{+w}{ }\PY{p}{\PYZob{}}
\PY{+w}{    }\PY{k+kd}{public}\PY{+w}{ }\PY{n}{List}\PY{o}{\PYZlt{}}\PY{n}{String}\PY{o}{\PYZgt{}}\PY{+w}{ }\PY{n+nf}{generateSummaries}\PY{p}{(}\PY{n}{List}\PY{o}{\PYZlt{}}\PY{n}{Double}\PY{o}{\PYZgt{}}\PY{+w}{ }\PY{n}{amounts}\PY{p}{)}\PY{+w}{ }\PY{p}{\PYZob{}}
\PY{+w}{        }\PY{k+kd}{class} \PY{n+nc}{Formatter}\PY{+w}{ }\PY{p}{\PYZob{}}\PY{+w}{ }\PY{c+c1}{// local class, scoped to this method}
\PY{+w}{            }\PY{n}{String}\PY{+w}{ }\PY{n+nf}{format}\PY{p}{(}\PY{k+kt}{double}\PY{+w}{ }\PY{n}{amount}\PY{p}{)}\PY{+w}{ }\PY{p}{\PYZob{}}
\PY{+w}{                }\PY{k}{return}\PY{+w}{ }\PY{n}{String}\PY{p}{.}\PY{n+na}{format}\PY{p}{(}\PY{l+s}{\PYZdq{}}\PY{l+s}{\PYZdl{}\PYZpc{}.2f}\PY{l+s}{\PYZdq{}}\PY{p}{,}\PY{+w}{ }\PY{n}{amount}\PY{p}{)}\PY{p}{;}
\PY{+w}{            }\PY{p}{\PYZcb{}}
\PY{+w}{        }\PY{p}{\PYZcb{}}

\PY{+w}{        }\PY{n}{Formatter}\PY{+w}{ }\PY{n}{formatter}\PY{+w}{ }\PY{o}{=}\PY{+w}{ }\PY{k}{new}\PY{+w}{ }\PY{n}{Formatter}\PY{p}{(}\PY{p}{)}\PY{p}{;}
\PY{+w}{        }\PY{n}{List}\PY{o}{\PYZlt{}}\PY{n}{String}\PY{o}{\PYZgt{}}\PY{+w}{ }\PY{n}{summaries}\PY{+w}{ }\PY{o}{=}\PY{+w}{ }\PY{k}{new}\PY{+w}{ }\PY{n}{ArrayList}\PY{o}{\PYZlt{}}\PY{o}{\PYZgt{}}\PY{p}{(}\PY{p}{)}\PY{p}{;}
\PY{+w}{        }\PY{k}{for}\PY{+w}{ }\PY{p}{(}\PY{k+kt}{double}\PY{+w}{ }\PY{n}{amount}\PY{+w}{ }\PY{p}{:}\PY{+w}{ }\PY{n}{amounts}\PY{p}{)}\PY{+w}{ }\PY{p}{\PYZob{}}
\PY{+w}{            }\PY{n}{summaries}\PY{p}{.}\PY{n+na}{add}\PY{p}{(}\PY{n}{formatter}\PY{p}{.}\PY{n+na}{format}\PY{p}{(}\PY{n}{amount}\PY{p}{)}\PY{p}{)}\PY{p}{;}
\PY{+w}{        }\PY{p}{\PYZcb{}}
\PY{+w}{        }\PY{k}{return}\PY{+w}{ }\PY{n}{summaries}\PY{p}{;}
\PY{+w}{    }\PY{p}{\PYZcb{}}
\PY{p}{\PYZcb{}}
\end{Verbatim}
\end{codebox}

\textbf{Anonymous class} — a class with no name, defined and instantiated in a single expression. Common before lambdas existed, and still used for interfaces with more than one method, or when you need to keep extra state.

\begin{codebox}
\begin{Verbatim}[commandchars=\\\{\}]
\PY{k+kd}{public}\PY{+w}{ }\PY{k+kd}{interface} \PY{n+nc}{DiscountRule}\PY{+w}{ }\PY{p}{\PYZob{}}
\PY{+w}{    }\PY{k+kt}{double}\PY{+w}{ }\PY{n+nf}{apply}\PY{p}{(}\PY{k+kt}{double}\PY{+w}{ }\PY{n}{price}\PY{p}{)}\PY{p}{;}
\PY{p}{\PYZcb{}}

\PY{k+kd}{public}\PY{+w}{ }\PY{k+kd}{class} \PY{n+nc}{CheckoutService}\PY{+w}{ }\PY{p}{\PYZob{}}
\PY{+w}{    }\PY{k+kd}{public}\PY{+w}{ }\PY{k+kt}{double}\PY{+w}{ }\PY{n+nf}{checkout}\PY{p}{(}\PY{k+kt}{double}\PY{+w}{ }\PY{n}{price}\PY{p}{,}\PY{+w}{ }\PY{n}{DiscountRule}\PY{+w}{ }\PY{n}{rule}\PY{p}{)}\PY{+w}{ }\PY{p}{\PYZob{}}
\PY{+w}{        }\PY{k}{return}\PY{+w}{ }\PY{n}{rule}\PY{p}{.}\PY{n+na}{apply}\PY{p}{(}\PY{n}{price}\PY{p}{)}\PY{p}{;}
\PY{+w}{    }\PY{p}{\PYZcb{}}

\PY{+w}{    }\PY{k+kd}{public}\PY{+w}{ }\PY{k+kt}{void}\PY{+w}{ }\PY{n+nf}{run}\PY{p}{(}\PY{p}{)}\PY{+w}{ }\PY{p}{\PYZob{}}
\PY{+w}{        }\PY{k+kt}{double}\PY{+w}{ }\PY{n}{finalPrice}\PY{+w}{ }\PY{o}{=}\PY{+w}{ }\PY{n}{checkout}\PY{p}{(}\PY{l+m+mf}{100.0}\PY{p}{,}\PY{+w}{ }\PY{k}{new}\PY{+w}{ }\PY{n}{DiscountRule}\PY{p}{(}\PY{p}{)}\PY{+w}{ }\PY{p}{\PYZob{}}\PY{+w}{ }\PY{c+c1}{// anonymous class}
\PY{+w}{            }\PY{n+nd}{@Override}
\PY{+w}{            }\PY{k+kd}{public}\PY{+w}{ }\PY{k+kt}{double}\PY{+w}{ }\PY{n+nf}{apply}\PY{p}{(}\PY{k+kt}{double}\PY{+w}{ }\PY{n}{price}\PY{p}{)}\PY{+w}{ }\PY{p}{\PYZob{}}
\PY{+w}{                }\PY{k}{return}\PY{+w}{ }\PY{n}{price}\PY{+w}{ }\PY{o}{*}\PY{+w}{ }\PY{l+m+mf}{0.85}\PY{p}{;}\PY{+w}{ }\PY{c+c1}{// 15\PYZpc{} off}
\PY{+w}{            }\PY{p}{\PYZcb{}}
\PY{+w}{        }\PY{p}{\PYZcb{}}\PY{p}{)}\PY{p}{;}
\PY{+w}{        }\PY{n}{System}\PY{p}{.}\PY{n+na}{out}\PY{p}{.}\PY{n+na}{println}\PY{p}{(}\PY{n}{finalPrice}\PY{p}{)}\PY{p}{;}\PY{+w}{ }\PY{c+c1}{// 85.0}
\PY{+w}{    }\PY{p}{\PYZcb{}}
\PY{p}{\PYZcb{}}
\end{Verbatim}
\end{codebox}

For a single-method (functional) interface like \code{DiscountRule}, a \textbf{lambda expression} is usually preferred over an anonymous class in modern Java — it’s shorter and equally clear:

\begin{codebox}
\begin{Verbatim}[commandchars=\\\{\}]
\PY{k+kt}{double}\PY{+w}{ }\PY{n}{finalPrice}\PY{+w}{ }\PY{o}{=}\PY{+w}{ }\PY{n}{checkout}\PY{p}{(}\PY{l+m+mf}{100.0}\PY{p}{,}\PY{+w}{ }\PY{n}{price}\PY{+w}{ }\PY{o}{\PYZhy{}}\PY{o}{\PYZgt{}}\PY{+w}{ }\PY{n}{price}\PY{+w}{ }\PY{o}{*}\PY{+w}{ }\PY{l+m+mf}{0.85}\PY{p}{)}\PY{p}{;}
\end{Verbatim}
\end{codebox}

\textbf{Gotcha: a non-static inner class secretly holds a reference to its outer instance}, which can cause \textbf{memory leaks} — the outer object can’t be garbage collected while the inner object is still reachable (e.g., stored in a long-lived collection or passed to another thread).

\begin{codebox}
\begin{Verbatim}[commandchars=\\\{\}]
\PY{k+kd}{public}\PY{+w}{ }\PY{k+kd}{class} \PY{n+nc}{Order}\PY{+w}{ }\PY{p}{\PYZob{}}
\PY{+w}{    }\PY{k+kd}{private}\PY{+w}{ }\PY{k+kd}{final}\PY{+w}{ }\PY{k+kt}{byte}\PY{o}{[}\PY{o}{]}\PY{+w}{ }\PY{n}{largePayload}\PY{+w}{ }\PY{o}{=}\PY{+w}{ }\PY{k}{new}\PY{+w}{ }\PY{k+kt}{byte}\PY{o}{[}\PY{l+m+mi}{10\PYZus{}000\PYZus{}000}\PY{o}{]}\PY{p}{;}\PY{+w}{ }\PY{c+c1}{// 10 MB}

\PY{+w}{    }\PY{k+kd}{public}\PY{+w}{ }\PY{k+kd}{class} \PY{n+nc}{Listener}\PY{+w}{ }\PY{p}{\PYZob{}}\PY{+w}{ }\PY{c+c1}{// inner class — implicitly holds a reference to the enclosing Order}
\PY{+w}{        }\PY{k+kt}{void}\PY{+w}{ }\PY{n+nf}{onEvent}\PY{p}{(}\PY{p}{)}\PY{+w}{ }\PY{p}{\PYZob{}}
\PY{+w}{            }\PY{n}{System}\PY{p}{.}\PY{n+na}{out}\PY{p}{.}\PY{n+na}{println}\PY{p}{(}\PY{l+s}{\PYZdq{}}\PY{l+s}{Event for order with payload size }\PY{l+s}{\PYZdq{}}\PY{+w}{ }\PY{o}{+}\PY{+w}{ }\PY{n}{largePayload}\PY{p}{.}\PY{n+na}{length}\PY{p}{)}\PY{p}{;}
\PY{+w}{        }\PY{p}{\PYZcb{}}
\PY{+w}{    }\PY{p}{\PYZcb{}}
\PY{p}{\PYZcb{}}

\PY{c+c1}{// If a Listener instance is stored somewhere long\PYZhy{}lived (e.g., a static list}
\PY{c+c1}{// or an event bus), the entire enclosing Order — including its 10 MB payload —}
\PY{c+c1}{// stays reachable and cannot be garbage collected, even if nobody needs the Order anymore.}
\end{Verbatim}
\end{codebox}

The fix: prefer a \code{static} nested class (optionally taking any needed outer data explicitly as a constructor parameter) when you don’t truly need access to the outer instance’s state.

\begin{codebox}
\begin{Verbatim}[commandchars=\\\{\}]
\PY{k+kd}{public}\PY{+w}{ }\PY{k+kd}{class} \PY{n+nc}{Order}\PY{+w}{ }\PY{p}{\PYZob{}}
\PY{+w}{    }\PY{k+kd}{private}\PY{+w}{ }\PY{k+kd}{final}\PY{+w}{ }\PY{k+kt}{byte}\PY{o}{[}\PY{o}{]}\PY{+w}{ }\PY{n}{largePayload}\PY{+w}{ }\PY{o}{=}\PY{+w}{ }\PY{k}{new}\PY{+w}{ }\PY{k+kt}{byte}\PY{o}{[}\PY{l+m+mi}{10\PYZus{}000\PYZus{}000}\PY{o}{]}\PY{p}{;}

\PY{+w}{    }\PY{k+kd}{public}\PY{+w}{ }\PY{k+kd}{static}\PY{+w}{ }\PY{k+kd}{class} \PY{n+nc}{Listener}\PY{+w}{ }\PY{p}{\PYZob{}}\PY{+w}{ }\PY{c+c1}{// static: no hidden reference to Order}
\PY{+w}{        }\PY{k+kt}{void}\PY{+w}{ }\PY{n+nf}{onEvent}\PY{p}{(}\PY{p}{)}\PY{+w}{ }\PY{p}{\PYZob{}}
\PY{+w}{            }\PY{n}{System}\PY{p}{.}\PY{n+na}{out}\PY{p}{.}\PY{n+na}{println}\PY{p}{(}\PY{l+s}{\PYZdq{}}\PY{l+s}{Event received}\PY{l+s}{\PYZdq{}}\PY{p}{)}\PY{p}{;}
\PY{+w}{        }\PY{p}{\PYZcb{}}
\PY{+w}{    }\PY{p}{\PYZcb{}}
\PY{p}{\PYZcb{}}
\end{Verbatim}
\end{codebox}

\textbf{Gotcha: local and anonymous classes can only capture \code{effectively~\allowbreak{}final} local variables} — variables that are never reassigned after initialization.

\begin{codebox}
\begin{Verbatim}[commandchars=\\\{\}]
\PY{k+kd}{public}\PY{+w}{ }\PY{n}{List}\PY{o}{\PYZlt{}}\PY{n}{Runnable}\PY{o}{\PYZgt{}}\PY{+w}{ }\PY{n+nf}{createTasks}\PY{p}{(}\PY{p}{)}\PY{+w}{ }\PY{p}{\PYZob{}}
\PY{+w}{    }\PY{n}{List}\PY{o}{\PYZlt{}}\PY{n}{Runnable}\PY{o}{\PYZgt{}}\PY{+w}{ }\PY{n}{tasks}\PY{+w}{ }\PY{o}{=}\PY{+w}{ }\PY{k}{new}\PY{+w}{ }\PY{n}{ArrayList}\PY{o}{\PYZlt{}}\PY{o}{\PYZgt{}}\PY{p}{(}\PY{p}{)}\PY{p}{;}
\PY{+w}{    }\PY{k}{for}\PY{+w}{ }\PY{p}{(}\PY{k+kt}{int}\PY{+w}{ }\PY{n}{i}\PY{+w}{ }\PY{o}{=}\PY{+w}{ }\PY{l+m+mi}{0}\PY{p}{;}\PY{+w}{ }\PY{n}{i}\PY{+w}{ }\PY{o}{\PYZlt{}}\PY{+w}{ }\PY{l+m+mi}{3}\PY{p}{;}\PY{+w}{ }\PY{n}{i}\PY{o}{+}\PY{o}{+}\PY{p}{)}\PY{+w}{ }\PY{p}{\PYZob{}}
\PY{+w}{        }\PY{k+kt}{int}\PY{+w}{ }\PY{n}{taskNumber}\PY{+w}{ }\PY{o}{=}\PY{+w}{ }\PY{n}{i}\PY{p}{;}\PY{+w}{ }\PY{c+c1}{// effectively final: assigned once per loop iteration}
\PY{+w}{        }\PY{n}{tasks}\PY{p}{.}\PY{n+na}{add}\PY{p}{(}\PY{k}{new}\PY{+w}{ }\PY{n}{Runnable}\PY{p}{(}\PY{p}{)}\PY{+w}{ }\PY{p}{\PYZob{}}
\PY{+w}{            }\PY{n+nd}{@Override}
\PY{+w}{            }\PY{k+kd}{public}\PY{+w}{ }\PY{k+kt}{void}\PY{+w}{ }\PY{n+nf}{run}\PY{p}{(}\PY{p}{)}\PY{+w}{ }\PY{p}{\PYZob{}}
\PY{+w}{                }\PY{n}{System}\PY{p}{.}\PY{n+na}{out}\PY{p}{.}\PY{n+na}{println}\PY{p}{(}\PY{l+s}{\PYZdq{}}\PY{l+s}{Task }\PY{l+s}{\PYZdq{}}\PY{+w}{ }\PY{o}{+}\PY{+w}{ }\PY{n}{taskNumber}\PY{p}{)}\PY{p}{;}\PY{+w}{ }\PY{c+c1}{// OK: captures effectively final variable}
\PY{+w}{            }\PY{p}{\PYZcb{}}
\PY{+w}{        }\PY{p}{\PYZcb{}}\PY{p}{)}\PY{p}{;}
\PY{+w}{    }\PY{p}{\PYZcb{}}
\PY{+w}{    }\PY{k}{return}\PY{+w}{ }\PY{n}{tasks}\PY{p}{;}
\PY{p}{\PYZcb{}}
\end{Verbatim}
\end{codebox}

\begin{codebox}
\begin{Verbatim}[commandchars=\\\{\}]
\PY{k+kd}{public}\PY{+w}{ }\PY{n}{Runnable}\PY{+w}{ }\PY{n+nf}{createTask}\PY{p}{(}\PY{p}{)}\PY{+w}{ }\PY{p}{\PYZob{}}
\PY{+w}{    }\PY{k+kt}{int}\PY{+w}{ }\PY{n}{counter}\PY{+w}{ }\PY{o}{=}\PY{+w}{ }\PY{l+m+mi}{0}\PY{p}{;}
\PY{+w}{    }\PY{n}{Runnable}\PY{+w}{ }\PY{n}{task}\PY{+w}{ }\PY{o}{=}\PY{+w}{ }\PY{p}{(}\PY{p}{)}\PY{+w}{ }\PY{o}{\PYZhy{}}\PY{o}{\PYZgt{}}\PY{+w}{ }\PY{n}{System}\PY{p}{.}\PY{n+na}{out}\PY{p}{.}\PY{n+na}{println}\PY{p}{(}\PY{n}{counter}\PY{p}{)}\PY{p}{;}
\PY{+w}{    }\PY{n}{counter}\PY{+w}{ }\PY{o}{=}\PY{+w}{ }\PY{l+m+mi}{5}\PY{p}{;}\PY{+w}{ }\PY{c+c1}{// Compile error: \PYZdq{}counter\PYZdq{} is captured by the lambda, so it must stay effectively final}
\PY{+w}{    }\PY{k}{return}\PY{+w}{ }\PY{n}{task}\PY{p}{;}
\PY{p}{\PYZcb{}}
\end{Verbatim}
\end{codebox}

\section[Common Code-Review Interview Pitfalls]{Common Code-Review Interview Pitfalls}
\label{h:2:common-code-review-interview-pitfalls}
\begin{enumerate}
\item 
\textbf{Public mutable fields instead of encapsulated access.} Why it matters: callers can corrupt an object’s invariants directly, bypassing any validation.

\begin{codebox}
\begin{Verbatim}[commandchars=\\\{\}]
\PY{c+c1}{// Before}
\PY{k+kd}{public}\PY{+w}{ }\PY{k+kd}{class} \PY{n+nc}{Account}\PY{+w}{ }\PY{p}{\PYZob{}}\PY{+w}{ }\PY{k+kd}{public}\PY{+w}{ }\PY{k+kt}{double}\PY{+w}{ }\PY{n}{balance}\PY{p}{;}\PY{+w}{ }\PY{p}{\PYZcb{}}
\PY{c+c1}{// After}
\PY{k+kd}{public}\PY{+w}{ }\PY{k+kd}{class} \PY{n+nc}{Account}\PY{+w}{ }\PY{p}{\PYZob{}}
\PY{+w}{    }\PY{k+kd}{private}\PY{+w}{ }\PY{k+kt}{double}\PY{+w}{ }\PY{n}{balance}\PY{p}{;}
\PY{+w}{    }\PY{k+kd}{public}\PY{+w}{ }\PY{k+kt}{void}\PY{+w}{ }\PY{n+nf}{withdraw}\PY{p}{(}\PY{k+kt}{double}\PY{+w}{ }\PY{n}{amount}\PY{p}{)}\PY{+w}{ }\PY{p}{\PYZob{}}\PY{+w}{ }\PY{c+cm}{/* validate, then mutate */}\PY{+w}{ }\PY{p}{\PYZcb{}}
\PY{p}{\PYZcb{}}
\end{Verbatim}
\end{codebox}

\item 
\textbf{Getters that return direct references to mutable internal collections/objects.} Why it matters: callers can mutate “private” state from outside, silently breaking encapsulation.

\begin{codebox}
\begin{Verbatim}[commandchars=\\\{\}]
\PY{c+c1}{// Before}
\PY{k+kd}{public}\PY{+w}{ }\PY{n}{List}\PY{o}{\PYZlt{}}\PY{n}{String}\PY{o}{\PYZgt{}}\PY{+w}{ }\PY{n+nf}{getItems}\PY{p}{(}\PY{p}{)}\PY{+w}{ }\PY{p}{\PYZob{}}\PY{+w}{ }\PY{k}{return}\PY{+w}{ }\PY{n}{items}\PY{p}{;}\PY{+w}{ }\PY{p}{\PYZcb{}}
\PY{c+c1}{// After}
\PY{k+kd}{public}\PY{+w}{ }\PY{n}{List}\PY{o}{\PYZlt{}}\PY{n}{String}\PY{o}{\PYZgt{}}\PY{+w}{ }\PY{n+nf}{getItems}\PY{p}{(}\PY{p}{)}\PY{+w}{ }\PY{p}{\PYZob{}}\PY{+w}{ }\PY{k}{return}\PY{+w}{ }\PY{n}{List}\PY{p}{.}\PY{n+na}{copyOf}\PY{p}{(}\PY{n}{items}\PY{p}{)}\PY{p}{;}\PY{+w}{ }\PY{p}{\PYZcb{}}
\end{Verbatim}
\end{codebox}

\item 
\textbf{Calling an overridable instance method from a constructor.} Why it matters: the subclass’s fields may not be initialized yet when the overridden method runs, producing wrong values (e.g., \code{0} or \code{null}).

\begin{codebox}
\begin{Verbatim}[commandchars=\\\{\}]
\PY{c+c1}{// Before: Shape() calls draw(), Circle overrides draw() using a field set after super() runs}
\PY{c+c1}{// After: mark draw() final, or move the call out of the constructor into an init method}
\end{Verbatim}
\end{codebox}

\item 
\textbf{Confusing method overloading with overriding, especially with a missing \code{@\allowbreak{}Override}.} Why it matters: a typo in parameters silently creates a new overload instead of overriding, and the bug hides until runtime.

\begin{codebox}
\begin{Verbatim}[commandchars=\\\{\}]
\PY{c+c1}{// Before (no annotation, silent overload)}
\PY{k+kd}{public}\PY{+w}{ }\PY{k+kt}{double}\PY{+w}{ }\PY{n+nf}{area}\PY{p}{(}\PY{k+kt}{double}\PY{+w}{ }\PY{n}{scale}\PY{p}{)}\PY{+w}{ }\PY{p}{\PYZob{}}\PY{+w}{ }\PY{p}{.}\PY{p}{.}\PY{p}{.}\PY{+w}{ }\PY{p}{\PYZcb{}}
\PY{c+c1}{// After (compiler enforces correctness)}
\PY{n+nd}{@Override}
\PY{k+kd}{public}\PY{+w}{ }\PY{k+kt}{double}\PY{+w}{ }\PY{n+nf}{area}\PY{p}{(}\PY{p}{)}\PY{+w}{ }\PY{p}{\PYZob{}}\PY{+w}{ }\PY{p}{.}\PY{p}{.}\PY{p}{.}\PY{+w}{ }\PY{p}{\PYZcb{}}
\end{Verbatim}
\end{codebox}

\item 
\textbf{Field hiding across a class hierarchy.} Why it matters: which field you read depends on the \emph{reference type}, not the actual object — a frequent source of confusing bugs.

\begin{codebox}
\begin{Verbatim}[commandchars=\\\{\}]
\PY{c+c1}{// Before: Car declares its own \PYZdq{}speed\PYZdq{} field that shadows Vehicle.speed}
\PY{c+c1}{// After: don\PYZsq{}t redeclare inherited fields; add new fields with distinct names, or use methods}
\end{Verbatim}
\end{codebox}

\item 
\textbf{Treating \code{static} methods as polymorphic.} Why it matters: \code{static} methods are hidden, not overridden — calling them through a superclass-typed reference always uses the reference’s declared type, not the runtime type.

\begin{codebox}
\begin{Verbatim}[commandchars=\\\{\}]
\PY{n}{Vehicle}\PY{+w}{ }\PY{n}{v}\PY{+w}{ }\PY{o}{=}\PY{+w}{ }\PY{k}{new}\PY{+w}{ }\PY{n}{Car}\PY{p}{(}\PY{p}{)}\PY{p}{;}
\PY{n}{v}\PY{p}{.}\PY{n+na}{category}\PY{p}{(}\PY{p}{)}\PY{p}{;}\PY{+w}{ }\PY{c+c1}{// resolves using Vehicle\PYZsq{}s static method, surprising reviewers expecting Car\PYZsq{}s}
\end{Verbatim}
\end{codebox}

\item 
\textbf{Assuming \code{final} means immutable.} Why it matters: \code{final} only prevents reassigning the reference; the referenced object (e.g., a \code{List}) can still be mutated internally.

\begin{codebox}
\begin{Verbatim}[commandchars=\\\{\}]
\PY{k+kd}{private}\PY{+w}{ }\PY{k+kd}{final}\PY{+w}{ }\PY{n}{List}\PY{o}{\PYZlt{}}\PY{n}{String}\PY{o}{\PYZgt{}}\PY{+w}{ }\PY{n}{items}\PY{+w}{ }\PY{o}{=}\PY{+w}{ }\PY{k}{new}\PY{+w}{ }\PY{n}{ArrayList}\PY{o}{\PYZlt{}}\PY{o}{\PYZgt{}}\PY{p}{(}\PY{p}{)}\PY{p}{;}
\PY{n}{items}\PY{p}{.}\PY{n+na}{add}\PY{p}{(}\PY{l+s}{\PYZdq{}}\PY{l+s}{x}\PY{l+s}{\PYZdq{}}\PY{p}{)}\PY{p}{;}\PY{+w}{ }\PY{c+c1}{// still allowed, even though \PYZdq{}items\PYZdq{} is final}
\end{Verbatim}
\end{codebox}

\item 
\textbf{Overly broad access modifiers (\code{public} by default).} Why it matters: wide visibility becomes a long-term maintenance liability — every public member is part of your API contract and hard to change later.

\begin{codebox}
\begin{Verbatim}[commandchars=\\\{\}]
\PY{c+c1}{// Before}
\PY{k+kd}{public}\PY{+w}{ }\PY{k+kt}{double}\PY{+w}{ }\PY{n}{balance}\PY{p}{;}
\PY{c+c1}{// After}
\PY{k+kd}{private}\PY{+w}{ }\PY{k+kt}{double}\PY{+w}{ }\PY{n}{balance}\PY{p}{;}
\PY{k+kd}{public}\PY{+w}{ }\PY{k+kt}{double}\PY{+w}{ }\PY{n+nf}{getBalance}\PY{p}{(}\PY{p}{)}\PY{+w}{ }\PY{p}{\PYZob{}}\PY{+w}{ }\PY{k}{return}\PY{+w}{ }\PY{n}{balance}\PY{p}{;}\PY{+w}{ }\PY{p}{\PYZcb{}}
\end{Verbatim}
\end{codebox}

\item 
\textbf{Non-static inner classes causing memory leaks.} Why it matters: an inner class instance implicitly holds a reference to its enclosing object, keeping it alive longer than expected if stored somewhere long-lived.

\begin{codebox}
\begin{Verbatim}[commandchars=\\\{\}]
\PY{c+c1}{// Before}
\PY{k+kd}{public}\PY{+w}{ }\PY{k+kd}{class} \PY{n+nc}{Listener}\PY{+w}{ }\PY{p}{\PYZob{}}\PY{+w}{ }\PY{c+cm}{/* implicit outer reference */}\PY{+w}{ }\PY{p}{\PYZcb{}}
\PY{c+c1}{// After}
\PY{k+kd}{public}\PY{+w}{ }\PY{k+kd}{static}\PY{+w}{ }\PY{k+kd}{class} \PY{n+nc}{Listener}\PY{+w}{ }\PY{p}{\PYZob{}}\PY{+w}{ }\PY{c+cm}{/* no implicit outer reference */}\PY{+w}{ }\PY{p}{\PYZcb{}}
\end{Verbatim}
\end{codebox}

\item 
\textbf{Overload resolution ambiguity with autoboxing and varargs.} Why it matters: adding a new overload can silently change which overload existing calls resolve to, without any compile error.

\begin{codebox}
\begin{Verbatim}[commandchars=\\\{\}]
\PY{k+kt}{void}\PY{+w}{ }\PY{n+nf}{log}\PY{p}{(}\PY{k+kt}{int}\PY{+w}{ }\PY{n}{code}\PY{p}{)}\PY{+w}{ }\PY{p}{\PYZob{}}\PY{+w}{ }\PY{p}{.}\PY{p}{.}\PY{p}{.}\PY{+w}{ }\PY{p}{\PYZcb{}}
\PY{k+kt}{void}\PY{+w}{ }\PY{n+nf}{log}\PY{p}{(}\PY{n}{Integer}\PY{+w}{ }\PY{n}{code}\PY{p}{)}\PY{+w}{ }\PY{p}{\PYZob{}}\PY{+w}{ }\PY{p}{.}\PY{p}{.}\PY{p}{.}\PY{+w}{ }\PY{p}{\PYZcb{}}
\PY{k+kt}{void}\PY{+w}{ }\PY{n+nf}{log}\PY{p}{(}\PY{n}{Object}\PY{p}{.}\PY{p}{.}\PY{p}{.}\PY{+w}{ }\PY{n}{codes}\PY{p}{)}\PY{+w}{ }\PY{p}{\PYZob{}}\PY{+w}{ }\PY{p}{.}\PY{p}{.}\PY{p}{.}\PY{+w}{ }\PY{p}{\PYZcb{}}
\PY{c+c1}{// log(5) always prefers int over Integer over varargs — surprising to reviewers unfamiliar with the resolution order}
\end{Verbatim}
\end{codebox}

\item 
\textbf{Narrowing visibility when overriding (or trying to).} Why it matters: it’s a compile error, but the intent (restricting subclass behavior) usually signals a deeper design problem — maybe the method shouldn’t be overridable at all.

\begin{codebox}
\begin{Verbatim}[commandchars=\\\{\}]
\PY{c+c1}{// Before: attempting to override public start() as protected — fails to compile}
\PY{c+c1}{// After: keep the same or wider access, or reconsider whether the base method should be public}
\end{Verbatim}
\end{codebox}

\item 
\textbf{Mutable \code{static} fields used as shared/global state.} Why it matters: they’re a hidden dependency between unrelated pieces of code, cause race conditions in multi-threaded contexts, and leak state between unit tests.

\begin{codebox}
\begin{Verbatim}[commandchars=\\\{\}]
\PY{c+c1}{// Before}
\PY{k+kd}{private}\PY{+w}{ }\PY{k+kd}{static}\PY{+w}{ }\PY{k+kt}{int}\PY{+w}{ }\PY{n}{nextId}\PY{+w}{ }\PY{o}{=}\PY{+w}{ }\PY{l+m+mi}{1}\PY{p}{;}\PY{+w}{ }\PY{c+c1}{// shared, unsynchronized}
\PY{c+c1}{// After}
\PY{k+kd}{private}\PY{+w}{ }\PY{k+kd}{static}\PY{+w}{ }\PY{k+kd}{final}\PY{+w}{ }\PY{n}{AtomicInteger}\PY{+w}{ }\PY{n}{nextId}\PY{+w}{ }\PY{o}{=}\PY{+w}{ }\PY{k}{new}\PY{+w}{ }\PY{n}{AtomicInteger}\PY{p}{(}\PY{l+m+mi}{1}\PY{p}{)}\PY{p}{;}
\end{Verbatim}
\end{codebox}

\item 
\textbf{No-args constructor silently disappearing once another constructor is added.} Why it matters: code (or a framework needing a default constructor, e.g. some serializers) that relied on \code{new~\allowbreak{}Foo()} breaks at compile or runtime.

\begin{codebox}
\begin{Verbatim}[commandchars=\\\{\}]
\PY{c+c1}{// Before: only Shape(String color) exists}
\PY{c+c1}{// After: explicitly add Shape() if a no\PYZhy{}args constructor is still needed}
\end{Verbatim}
\end{codebox}

\item 
\textbf{Using an abstract class where a simple interface would do (or vice versa).} Why it matters: abstract classes force single inheritance and couple unrelated types; interfaces should be preferred for pure capability contracts, abstract classes for shared implementation/state.

\begin{codebox}
\begin{Verbatim}[commandchars=\\\{\}]
\PY{c+c1}{// Before: abstract class Discountable with no shared state or implementation}
\PY{c+c1}{// After: interface Discountable \PYZob{} double applyDiscount(double price); \PYZcb{}}
\end{Verbatim}
\end{codebox}

\item 
\textbf{Anonymous classes used where a lambda would be clearer.} Why it matters: for single-method functional interfaces, an anonymous class is verbose boilerplate compared to a lambda, and reviewers should flag it as an easy simplification.

\begin{codebox}
\begin{Verbatim}[commandchars=\\\{\}]
\PY{c+c1}{// Before}
\PY{k}{new}\PY{+w}{ }\PY{n}{DiscountRule}\PY{p}{(}\PY{p}{)}\PY{+w}{ }\PY{p}{\PYZob{}}
\PY{+w}{    }\PY{n+nd}{@Override}
\PY{+w}{    }\PY{k+kd}{public}\PY{+w}{ }\PY{k+kt}{double}\PY{+w}{ }\PY{n+nf}{apply}\PY{p}{(}\PY{k+kt}{double}\PY{+w}{ }\PY{n}{price}\PY{p}{)}\PY{+w}{ }\PY{p}{\PYZob{}}\PY{+w}{ }\PY{k}{return}\PY{+w}{ }\PY{n}{price}\PY{+w}{ }\PY{o}{*}\PY{+w}{ }\PY{l+m+mf}{0.85}\PY{p}{;}\PY{+w}{ }\PY{p}{\PYZcb{}}
\PY{p}{\PYZcb{}}\PY{p}{;}
\PY{c+c1}{// After}
\PY{p}{(}\PY{k+kt}{double}\PY{+w}{ }\PY{n}{price}\PY{p}{)}\PY{+w}{ }\PY{o}{\PYZhy{}}\PY{o}{\PYZgt{}}\PY{+w}{ }\PY{n}{price}\PY{+w}{ }\PY{o}{*}\PY{+w}{ }\PY{l+m+mf}{0.85}\PY{p}{;}
\end{Verbatim}
\end{codebox}

\end{enumerate}


\setcounter{chapter}{2}
\chapter{Modern Java Language Features}
\label{chap:3}
\label{h:3:3-modern-java-language-features}
Java has changed a lot since Java 8. Newer versions added features that make code shorter, safer, and easier to read. This chapter covers the features you are most likely to see in a modern codebase or be asked about in a code-review interview. Each section explains what the feature is, which JDK version it shipped in, and shows a small, realistic example.

\begin{itemize}
\item 
\hyperref[h:3:enums]{Enums}

\item 
\hyperref[h:3:records]{Records}

\item 
\hyperref[h:3:sealed-classes-and-interfaces]{Sealed Classes and Interfaces}

\item 
\hyperref[h:3:interfaces-default-static-private-methods]{Interfaces (default, static, private methods)}

\item 
\hyperref[h:3:abstract-classes]{Abstract Classes}

\item 
\hyperref[h:3:pattern-matching-instanceof-switch]{Pattern Matching (instanceof, switch)}

\item 
\hyperref[h:3:primitive-types-in-patterns]{Primitive Types in Patterns}

\item 
\hyperref[h:3:switch-expressions]{Switch Expressions}

\item 
\hyperref[h:3:text-blocks]{Text Blocks}

\item 
\hyperref[h:3:string-templates-preview]{String Templates (Preview)}

\item 
\hyperref[h:3:unnamed-variables-and-patterns]{Unnamed Variables and Patterns}

\item 
\hyperref[h:3:value-based-classes]{Value-Based Classes}

\item 
\hyperref[h:3:preview-and-incubator-features]{Preview and Incubator Features}

\item 
\hyperref[h:3:deprecated-and-removed-features-across-java-versions]{Deprecated and Removed Features Across Java Versions}

\item 
\hyperref[h:3:common-code-review-interview-pitfalls]{Common Code-Review Interview Pitfalls}

\end{itemize}

\section[Enums]{Enums}
\label{h:3:enums}
An \textbf{enum} (short for “enumeration”) is a type with a fixed set of constant values. Use an enum when a variable can only be one of a small, known list of options, like days of the week or order statuses.

Enums in Java are full classes. They can have fields, constructors, methods, and they can even implement interfaces. This is more powerful than enums in many other languages.

\begin{codebox}
\begin{Verbatim}[commandchars=\\\{\}]
\PY{k+kd}{public}\PY{+w}{ }\PY{k+kd}{enum}\PY{+w}{ }\PY{n}{OrderStatus}\PY{+w}{ }\PY{p}{\PYZob{}}
\PY{+w}{    }\PY{n}{NEW}\PY{p}{,}
\PY{+w}{    }\PY{n}{PAID}\PY{p}{,}
\PY{+w}{    }\PY{n}{SHIPPED}\PY{p}{,}
\PY{+w}{    }\PY{n}{DELIVERED}\PY{p}{,}
\PY{+w}{    }\PY{n}{CANCELLED}
\PY{p}{\PYZcb{}}
\end{Verbatim}
\end{codebox}

Enums can carry data and behavior:

\begin{codebox}
\begin{Verbatim}[commandchars=\\\{\}]
\PY{k+kd}{public}\PY{+w}{ }\PY{k+kd}{enum}\PY{+w}{ }\PY{n}{Planet}\PY{+w}{ }\PY{p}{\PYZob{}}
\PY{+w}{    }\PY{n}{MERCURY}\PY{p}{(}\PY{l+m+mf}{3.303e+23}\PY{p}{,}\PY{+w}{ }\PY{l+m+mf}{2.4397e6}\PY{p}{)}\PY{p}{,}
\PY{+w}{    }\PY{n}{VENUS}\PY{p}{(}\PY{l+m+mf}{4.869e+24}\PY{p}{,}\PY{+w}{ }\PY{l+m+mf}{6.0518e6}\PY{p}{)}\PY{p}{,}
\PY{+w}{    }\PY{n}{EARTH}\PY{p}{(}\PY{l+m+mf}{5.976e+24}\PY{p}{,}\PY{+w}{ }\PY{l+m+mf}{6.37814e6}\PY{p}{)}\PY{p}{;}

\PY{+w}{    }\PY{k+kd}{private}\PY{+w}{ }\PY{k+kd}{final}\PY{+w}{ }\PY{k+kt}{double}\PY{+w}{ }\PY{n}{mass}\PY{p}{;}\PY{+w}{   }\PY{c+c1}{// in kilograms}
\PY{+w}{    }\PY{k+kd}{private}\PY{+w}{ }\PY{k+kd}{final}\PY{+w}{ }\PY{k+kt}{double}\PY{+w}{ }\PY{n}{radius}\PY{p}{;}\PY{+w}{ }\PY{c+c1}{// in meters}

\PY{+w}{    }\PY{n}{Planet}\PY{p}{(}\PY{k+kt}{double}\PY{+w}{ }\PY{n}{mass}\PY{p}{,}\PY{+w}{ }\PY{k+kt}{double}\PY{+w}{ }\PY{n}{radius}\PY{p}{)}\PY{+w}{ }\PY{p}{\PYZob{}}
\PY{+w}{        }\PY{k}{this}\PY{p}{.}\PY{n+na}{mass}\PY{+w}{ }\PY{o}{=}\PY{+w}{ }\PY{n}{mass}\PY{p}{;}
\PY{+w}{        }\PY{k}{this}\PY{p}{.}\PY{n+na}{radius}\PY{+w}{ }\PY{o}{=}\PY{+w}{ }\PY{n}{radius}\PY{p}{;}
\PY{+w}{    }\PY{p}{\PYZcb{}}

\PY{+w}{    }\PY{k+kt}{double}\PY{+w}{ }\PY{n+nf}{surfaceGravity}\PY{p}{(}\PY{p}{)}\PY{+w}{ }\PY{p}{\PYZob{}}
\PY{+w}{        }\PY{k+kd}{final}\PY{+w}{ }\PY{k+kt}{double}\PY{+w}{ }\PY{n}{G}\PY{+w}{ }\PY{o}{=}\PY{+w}{ }\PY{l+m+mf}{6.67300E\PYZhy{}11}\PY{p}{;}
\PY{+w}{        }\PY{k}{return}\PY{+w}{ }\PY{n}{G}\PY{+w}{ }\PY{o}{*}\PY{+w}{ }\PY{n}{mass}\PY{+w}{ }\PY{o}{/}\PY{+w}{ }\PY{p}{(}\PY{n}{radius}\PY{+w}{ }\PY{o}{*}\PY{+w}{ }\PY{n}{radius}\PY{p}{)}\PY{p}{;}
\PY{+w}{    }\PY{p}{\PYZcb{}}
\PY{p}{\PYZcb{}}

\PY{c+c1}{// Usage:}
\PY{k+kt}{double}\PY{+w}{ }\PY{n}{g}\PY{+w}{ }\PY{o}{=}\PY{+w}{ }\PY{n}{Planet}\PY{p}{.}\PY{n+na}{EARTH}\PY{p}{.}\PY{n+na}{surfaceGravity}\PY{p}{(}\PY{p}{)}\PY{p}{;}\PY{+w}{ }\PY{c+c1}{// \PYZti{}9.8}
\end{Verbatim}
\end{codebox}

Enums can also have abstract methods, where each constant provides its own implementation:

\begin{codebox}
\begin{Verbatim}[commandchars=\\\{\}]
\PY{k+kd}{public}\PY{+w}{ }\PY{k+kd}{enum}\PY{+w}{ }\PY{n}{Operation}\PY{+w}{ }\PY{p}{\PYZob{}}
\PY{+w}{    }\PY{n}{PLUS}\PY{+w}{ }\PY{p}{\PYZob{}}
\PY{+w}{        }\PY{k+kd}{public}\PY{+w}{ }\PY{k+kt}{int}\PY{+w}{ }\PY{n+nf}{apply}\PY{p}{(}\PY{k+kt}{int}\PY{+w}{ }\PY{n}{a}\PY{p}{,}\PY{+w}{ }\PY{k+kt}{int}\PY{+w}{ }\PY{n}{b}\PY{p}{)}\PY{+w}{ }\PY{p}{\PYZob{}}\PY{+w}{ }\PY{k}{return}\PY{+w}{ }\PY{n}{a}\PY{+w}{ }\PY{o}{+}\PY{+w}{ }\PY{n}{b}\PY{p}{;}\PY{+w}{ }\PY{p}{\PYZcb{}}
\PY{+w}{    }\PY{p}{\PYZcb{}}\PY{p}{,}
\PY{+w}{    }\PY{n}{MINUS}\PY{+w}{ }\PY{p}{\PYZob{}}
\PY{+w}{        }\PY{k+kd}{public}\PY{+w}{ }\PY{k+kt}{int}\PY{+w}{ }\PY{n+nf}{apply}\PY{p}{(}\PY{k+kt}{int}\PY{+w}{ }\PY{n}{a}\PY{p}{,}\PY{+w}{ }\PY{k+kt}{int}\PY{+w}{ }\PY{n}{b}\PY{p}{)}\PY{+w}{ }\PY{p}{\PYZob{}}\PY{+w}{ }\PY{k}{return}\PY{+w}{ }\PY{n}{a}\PY{+w}{ }\PY{o}{\PYZhy{}}\PY{+w}{ }\PY{n}{b}\PY{p}{;}\PY{+w}{ }\PY{p}{\PYZcb{}}
\PY{+w}{    }\PY{p}{\PYZcb{}}\PY{p}{;}

\PY{+w}{    }\PY{k+kd}{public}\PY{+w}{ }\PY{k+kd}{abstract}\PY{+w}{ }\PY{k+kt}{int}\PY{+w}{ }\PY{n+nf}{apply}\PY{p}{(}\PY{k+kt}{int}\PY{+w}{ }\PY{n}{a}\PY{p}{,}\PY{+w}{ }\PY{k+kt}{int}\PY{+w}{ }\PY{n}{b}\PY{p}{)}\PY{p}{;}
\PY{p}{\PYZcb{}}

\PY{c+c1}{// Usage:}
\PY{k+kt}{int}\PY{+w}{ }\PY{n}{result}\PY{+w}{ }\PY{o}{=}\PY{+w}{ }\PY{n}{Operation}\PY{p}{.}\PY{n+na}{PLUS}\PY{p}{.}\PY{n+na}{apply}\PY{p}{(}\PY{l+m+mi}{2}\PY{p}{,}\PY{+w}{ }\PY{l+m+mi}{3}\PY{p}{)}\PY{p}{;}\PY{+w}{ }\PY{c+c1}{// 5}
\end{Verbatim}
\end{codebox}

Key facts every reviewer should know:

{\tablefont
\begin{xltabular}{\linewidth}{@{}L{0.427}L{1.573}@{}}
\toprule
\rowcolor{tablehdr}
\thd{Fact} & \thd{Detail} \\
\midrule
\endhead
Since & Java 5 \\
Base class & Implicitly extends \code{java.\allowbreak{}lang.\allowbreak{}Enum} (cannot extend anything else) \\
Interfaces & Can implement one or more interfaces \\
Comparison & Always safe to compare with \code{==} (constants are singletons) \\
Switch & Enums work great in \code{switch} statements and expressions \\
Serialization & Enum serialization is handled specially and is safe by default \\
\code{values()} / \code{valueOf()} & Auto-generated static methods to list constants or parse from a string \\
\bottomrule
\end{xltabular}
}

\begin{codebox}
\begin{Verbatim}[commandchars=\\\{\}]
\PY{k}{for}\PY{+w}{ }\PY{p}{(}\PY{n}{OrderStatus}\PY{+w}{ }\PY{n}{s}\PY{+w}{ }\PY{p}{:}\PY{+w}{ }\PY{n}{OrderStatus}\PY{p}{.}\PY{n+na}{values}\PY{p}{(}\PY{p}{)}\PY{p}{)}\PY{+w}{ }\PY{p}{\PYZob{}}
\PY{+w}{    }\PY{n}{System}\PY{p}{.}\PY{n+na}{out}\PY{p}{.}\PY{n+na}{println}\PY{p}{(}\PY{n}{s}\PY{p}{)}\PY{p}{;}\PY{+w}{ }\PY{c+c1}{// NEW, PAID, SHIPPED, DELIVERED, CANCELLED}
\PY{p}{\PYZcb{}}

\PY{n}{OrderStatus}\PY{+w}{ }\PY{n}{s}\PY{+w}{ }\PY{o}{=}\PY{+w}{ }\PY{n}{OrderStatus}\PY{p}{.}\PY{n+na}{valueOf}\PY{p}{(}\PY{l+s}{\PYZdq{}}\PY{l+s}{PAID}\PY{l+s}{\PYZdq{}}\PY{p}{)}\PY{p}{;}\PY{+w}{ }\PY{c+c1}{// throws IllegalArgumentException if not found}
\end{Verbatim}
\end{codebox}

\section[Records]{Records}
\label{h:3:records}
A \textbf{record} is a compact way to declare an immutable data class. Before records, developers had to hand-write (or generate) a constructor, getters, \code{equals()}, \code{hashCode()}, and \code{toString()} for simple “data holder” classes. Records generate all of that automatically.

Records were previewed in Java 14/15 and finalized as a standard feature in \textbf{Java 16}.

\begin{codebox}
\begin{Verbatim}[commandchars=\\\{\}]
\PY{k+kd}{public}\PY{+w}{ }\PY{k+kd}{record} \PY{n+nc}{Point}\PY{p}{(}\PY{k+kt}{int}\PY{+w}{ }\PY{n}{x}\PY{p}{,}\PY{+w}{ }\PY{k+kt}{int}\PY{+w}{ }\PY{n}{y}\PY{p}{)}\PY{+w}{ }\PY{p}{\PYZob{}}\PY{+w}{ }\PY{p}{\PYZcb{}}

\PY{c+c1}{// The compiler generates:}
\PY{c+c1}{// \PYZhy{} a constructor Point(int x, int y)}
\PY{c+c1}{// \PYZhy{} accessors x() and y() (note: no \PYZdq{}get\PYZdq{} prefix)}
\PY{c+c1}{// \PYZhy{} equals(), hashCode(), toString()}

\PY{n}{Point}\PY{+w}{ }\PY{n}{p1}\PY{+w}{ }\PY{o}{=}\PY{+w}{ }\PY{k}{new}\PY{+w}{ }\PY{n}{Point}\PY{p}{(}\PY{l+m+mi}{1}\PY{p}{,}\PY{+w}{ }\PY{l+m+mi}{2}\PY{p}{)}\PY{p}{;}
\PY{n}{Point}\PY{+w}{ }\PY{n}{p2}\PY{+w}{ }\PY{o}{=}\PY{+w}{ }\PY{k}{new}\PY{+w}{ }\PY{n}{Point}\PY{p}{(}\PY{l+m+mi}{1}\PY{p}{,}\PY{+w}{ }\PY{l+m+mi}{2}\PY{p}{)}\PY{p}{;}

\PY{n}{System}\PY{p}{.}\PY{n+na}{out}\PY{p}{.}\PY{n+na}{println}\PY{p}{(}\PY{n}{p1}\PY{p}{)}\PY{p}{;}\PY{+w}{          }\PY{c+c1}{// Point[x=1, y=2]}
\PY{n}{System}\PY{p}{.}\PY{n+na}{out}\PY{p}{.}\PY{n+na}{println}\PY{p}{(}\PY{n}{p1}\PY{p}{.}\PY{n+na}{x}\PY{p}{(}\PY{p}{)}\PY{p}{)}\PY{p}{;}\PY{+w}{      }\PY{c+c1}{// 1}
\PY{n}{System}\PY{p}{.}\PY{n+na}{out}\PY{p}{.}\PY{n+na}{println}\PY{p}{(}\PY{n}{p1}\PY{p}{.}\PY{n+na}{equals}\PY{p}{(}\PY{n}{p2}\PY{p}{)}\PY{p}{)}\PY{p}{;}\PY{+w}{ }\PY{c+c1}{// true}
\end{Verbatim}
\end{codebox}

You can add validation with a \textbf{compact constructor}, plus extra methods and static factories:

\begin{codebox}
\begin{Verbatim}[commandchars=\\\{\}]
\PY{k+kd}{public}\PY{+w}{ }\PY{k+kd}{record} \PY{n+nc}{Range}\PY{p}{(}\PY{k+kt}{int}\PY{+w}{ }\PY{n}{min}\PY{p}{,}\PY{+w}{ }\PY{k+kt}{int}\PY{+w}{ }\PY{n}{max}\PY{p}{)}\PY{+w}{ }\PY{p}{\PYZob{}}

\PY{+w}{    }\PY{c+c1}{// Compact constructor: runs before fields are assigned.}
\PY{+w}{    }\PY{k+kd}{public}\PY{+w}{ }\PY{n}{Range}\PY{+w}{ }\PY{p}{\PYZob{}}
\PY{+w}{        }\PY{k}{if}\PY{+w}{ }\PY{p}{(}\PY{n}{min}\PY{+w}{ }\PY{o}{\PYZgt{}}\PY{+w}{ }\PY{n}{max}\PY{p}{)}\PY{+w}{ }\PY{p}{\PYZob{}}
\PY{+w}{            }\PY{k}{throw}\PY{+w}{ }\PY{k}{new}\PY{+w}{ }\PY{n}{IllegalArgumentException}\PY{p}{(}\PY{l+s}{\PYZdq{}}\PY{l+s}{min must be \PYZlt{}= max}\PY{l+s}{\PYZdq{}}\PY{p}{)}\PY{p}{;}
\PY{+w}{        }\PY{p}{\PYZcb{}}
\PY{+w}{    }\PY{p}{\PYZcb{}}

\PY{+w}{    }\PY{k+kd}{public}\PY{+w}{ }\PY{k+kt}{int}\PY{+w}{ }\PY{n+nf}{length}\PY{p}{(}\PY{p}{)}\PY{+w}{ }\PY{p}{\PYZob{}}
\PY{+w}{        }\PY{k}{return}\PY{+w}{ }\PY{n}{max}\PY{+w}{ }\PY{o}{\PYZhy{}}\PY{+w}{ }\PY{n}{min}\PY{p}{;}
\PY{+w}{    }\PY{p}{\PYZcb{}}

\PY{+w}{    }\PY{k+kd}{public}\PY{+w}{ }\PY{k+kd}{static}\PY{+w}{ }\PY{n}{Range}\PY{+w}{ }\PY{n+nf}{of}\PY{p}{(}\PY{k+kt}{int}\PY{+w}{ }\PY{n}{min}\PY{p}{,}\PY{+w}{ }\PY{k+kt}{int}\PY{+w}{ }\PY{n}{max}\PY{p}{)}\PY{+w}{ }\PY{p}{\PYZob{}}
\PY{+w}{        }\PY{k}{return}\PY{+w}{ }\PY{k}{new}\PY{+w}{ }\PY{n}{Range}\PY{p}{(}\PY{n}{min}\PY{p}{,}\PY{+w}{ }\PY{n}{max}\PY{p}{)}\PY{p}{;}
\PY{+w}{    }\PY{p}{\PYZcb{}}
\PY{p}{\PYZcb{}}

\PY{n}{Range}\PY{+w}{ }\PY{n}{r}\PY{+w}{ }\PY{o}{=}\PY{+w}{ }\PY{n}{Range}\PY{p}{.}\PY{n+na}{of}\PY{p}{(}\PY{l+m+mi}{1}\PY{p}{,}\PY{+w}{ }\PY{l+m+mi}{10}\PY{p}{)}\PY{p}{;}
\PY{n}{System}\PY{p}{.}\PY{n+na}{out}\PY{p}{.}\PY{n+na}{println}\PY{p}{(}\PY{n}{r}\PY{p}{.}\PY{n+na}{length}\PY{p}{(}\PY{p}{)}\PY{p}{)}\PY{p}{;}\PY{+w}{ }\PY{c+c1}{// 9}
\end{Verbatim}
\end{codebox}

Records vs. regular classes:

{\tablefont
\begin{xltabular}{\linewidth}{@{}L{0.655}L{1.418}L{0.927}@{}}
\toprule
\rowcolor{tablehdr}
\thd{Aspect} & \thd{Record} & \thd{Regular class} \\
\midrule
\endhead
Fields & All fields are \code{private~\allowbreak{}final} automatically & You choose mutability \\
Accessors & Auto-generated (\code{x()}, not \code{getX()}) & You write them, often \code{getX()} \\
\code{equals}/\code{hashCode}/\code{toString} & Auto-generated from all fields & You write them (or use IDE/Lombok) \\
Inheritance & Cannot extend another class (implicitly extends \code{Record}) & Can extend any class \\
Can implement interfaces & Yes & Yes \\
Mutability & Always immutable (fields cannot change) & Can be mutable or immutable \\
Best use case & Simple immutable data carriers (DTOs, value objects) & Entities with behavior, mutable state, or inheritance \\
\bottomrule
\end{xltabular}
}

Records can implement interfaces, which is useful with sealed types (see next section):

\begin{codebox}
\begin{Verbatim}[commandchars=\\\{\}]
\PY{k+kd}{public}\PY{+w}{ }\PY{k+kd}{interface} \PY{n+nc}{Shape}\PY{+w}{ }\PY{p}{\PYZob{}}
\PY{+w}{    }\PY{k+kt}{double}\PY{+w}{ }\PY{n+nf}{area}\PY{p}{(}\PY{p}{)}\PY{p}{;}
\PY{p}{\PYZcb{}}

\PY{k+kd}{public}\PY{+w}{ }\PY{k+kd}{record} \PY{n+nc}{Circle}\PY{p}{(}\PY{k+kt}{double}\PY{+w}{ }\PY{n}{radius}\PY{p}{)}\PY{+w}{ }\PY{k+kd}{implements}\PY{+w}{ }\PY{n}{Shape}\PY{+w}{ }\PY{p}{\PYZob{}}
\PY{+w}{    }\PY{k+kd}{public}\PY{+w}{ }\PY{k+kt}{double}\PY{+w}{ }\PY{n+nf}{area}\PY{p}{(}\PY{p}{)}\PY{+w}{ }\PY{p}{\PYZob{}}
\PY{+w}{        }\PY{k}{return}\PY{+w}{ }\PY{n}{Math}\PY{p}{.}\PY{n+na}{PI}\PY{+w}{ }\PY{o}{*}\PY{+w}{ }\PY{n}{radius}\PY{+w}{ }\PY{o}{*}\PY{+w}{ }\PY{n}{radius}\PY{p}{;}
\PY{+w}{    }\PY{p}{\PYZcb{}}
\PY{p}{\PYZcb{}}
\end{Verbatim}
\end{codebox}

Records can also be \textbf{generic}:

\begin{codebox}
\begin{Verbatim}[commandchars=\\\{\}]
\PY{k+kd}{public}\PY{+w}{ }\PY{k+kd}{record} \PY{n+nc}{Pair}\PY{o}{\PYZlt{}}\PY{n}{A}\PY{p}{,}\PY{+w}{ }\PY{n}{B}\PY{o}{\PYZgt{}}\PY{p}{(}\PY{n}{A}\PY{+w}{ }\PY{n}{first}\PY{p}{,}\PY{+w}{ }\PY{n}{B}\PY{+w}{ }\PY{n}{second}\PY{p}{)}\PY{+w}{ }\PY{p}{\PYZob{}}\PY{+w}{ }\PY{p}{\PYZcb{}}

\PY{n}{Pair}\PY{o}{\PYZlt{}}\PY{n}{String}\PY{p}{,}\PY{+w}{ }\PY{n}{Integer}\PY{o}{\PYZgt{}}\PY{+w}{ }\PY{n}{pair}\PY{+w}{ }\PY{o}{=}\PY{+w}{ }\PY{k}{new}\PY{+w}{ }\PY{n}{Pair}\PY{o}{\PYZlt{}}\PY{o}{\PYZgt{}}\PY{p}{(}\PY{l+s}{\PYZdq{}}\PY{l+s}{age}\PY{l+s}{\PYZdq{}}\PY{p}{,}\PY{+w}{ }\PY{l+m+mi}{30}\PY{p}{)}\PY{p}{;}
\end{Verbatim}
\end{codebox}

\section[Sealed Classes and Interfaces]{Sealed Classes and Interfaces}
\label{h:3:sealed-classes-and-interfaces}
A \textbf{sealed} class or interface restricts which other classes or interfaces may extend or implement it. This gives you controlled inheritance: you (the author) decide the exact, closed list of subtypes. Sealed types were finalized in \textbf{Java 17}.

Why this matters: with a sealed hierarchy, the compiler (and tools like \code{switch}) knows every possible subtype. This enables \textbf{exhaustiveness checking} — the compiler can tell you if you forgot to handle a case.

\begin{codebox}
\begin{Verbatim}[commandchars=\\\{\}]
\PY{k+kd}{public}\PY{+w}{ }\PY{k+kd}{sealed}\PY{+w}{ }\PY{k+kd}{interface} \PY{n+nc}{Shape}\PY{+w}{ }\PY{n}{permits}\PY{+w}{ }\PY{n}{Circle}\PY{p}{,}\PY{+w}{ }\PY{n}{Square}\PY{p}{,}\PY{+w}{ }\PY{n}{Rectangle}\PY{+w}{ }\PY{p}{\PYZob{}}\PY{+w}{ }\PY{p}{\PYZcb{}}

\PY{k+kd}{public}\PY{+w}{ }\PY{k+kd}{record} \PY{n+nc}{Circle}\PY{p}{(}\PY{k+kt}{double}\PY{+w}{ }\PY{n}{radius}\PY{p}{)}\PY{+w}{ }\PY{k+kd}{implements}\PY{+w}{ }\PY{n}{Shape}\PY{+w}{ }\PY{p}{\PYZob{}}\PY{+w}{ }\PY{p}{\PYZcb{}}
\PY{k+kd}{public}\PY{+w}{ }\PY{k+kd}{record} \PY{n+nc}{Square}\PY{p}{(}\PY{k+kt}{double}\PY{+w}{ }\PY{n}{side}\PY{p}{)}\PY{+w}{ }\PY{k+kd}{implements}\PY{+w}{ }\PY{n}{Shape}\PY{+w}{ }\PY{p}{\PYZob{}}\PY{+w}{ }\PY{p}{\PYZcb{}}
\PY{k+kd}{public}\PY{+w}{ }\PY{k+kd}{record} \PY{n+nc}{Rectangle}\PY{p}{(}\PY{k+kt}{double}\PY{+w}{ }\PY{n}{width}\PY{p}{,}\PY{+w}{ }\PY{k+kt}{double}\PY{+w}{ }\PY{n}{height}\PY{p}{)}\PY{+w}{ }\PY{k+kd}{implements}\PY{+w}{ }\PY{n}{Shape}\PY{+w}{ }\PY{p}{\PYZob{}}\PY{+w}{ }\PY{p}{\PYZcb{}}
\end{Verbatim}
\end{codebox}

If the permitted subtypes are in the same file, you can drop the \code{permits} clause — the compiler figures it out:

\begin{codebox}
\begin{Verbatim}[commandchars=\\\{\}]
\PY{k+kd}{public}\PY{+w}{ }\PY{k+kd}{sealed}\PY{+w}{ }\PY{k+kd}{interface} \PY{n+nc}{Payment}\PY{+w}{ }\PY{p}{\PYZob{}}
\PY{+w}{    }\PY{k+kd}{record} \PY{n+nc}{CreditCard}\PY{p}{(}\PY{n}{String}\PY{+w}{ }\PY{n}{number}\PY{p}{)}\PY{+w}{ }\PY{k+kd}{implements}\PY{+w}{ }\PY{n}{Payment}\PY{+w}{ }\PY{p}{\PYZob{}}\PY{+w}{ }\PY{p}{\PYZcb{}}
\PY{+w}{    }\PY{k+kd}{record} \PY{n+nc}{BankTransfer}\PY{p}{(}\PY{n}{String}\PY{+w}{ }\PY{n}{iban}\PY{p}{)}\PY{+w}{ }\PY{k+kd}{implements}\PY{+w}{ }\PY{n}{Payment}\PY{+w}{ }\PY{p}{\PYZob{}}\PY{+w}{ }\PY{p}{\PYZcb{}}
\PY{+w}{    }\PY{k+kd}{record} \PY{n+nc}{Cash}\PY{p}{(}\PY{p}{)}\PY{+w}{ }\PY{k+kd}{implements}\PY{+w}{ }\PY{n}{Payment}\PY{+w}{ }\PY{p}{\PYZob{}}\PY{+w}{ }\PY{p}{\PYZcb{}}
\PY{p}{\PYZcb{}}
\end{Verbatim}
\end{codebox}

A permitted subclass must be declared as one of: \code{final}, \code{sealed} (with its own permitted subtypes), or \code{non-\allowbreak{}sealed} (reopens the hierarchy for further extension).

\begin{codebox}
\begin{Verbatim}[commandchars=\\\{\}]
\PY{k+kd}{public}\PY{+w}{ }\PY{k+kd}{sealed}\PY{+w}{ }\PY{k+kd}{class} \PY{n+nc}{Vehicle}\PY{+w}{ }\PY{n}{permits}\PY{+w}{ }\PY{n}{Car}\PY{p}{,}\PY{+w}{ }\PY{n}{Truck}\PY{+w}{ }\PY{p}{\PYZob{}}\PY{+w}{ }\PY{p}{\PYZcb{}}

\PY{k+kd}{public}\PY{+w}{ }\PY{k+kd}{final}\PY{+w}{ }\PY{k+kd}{class} \PY{n+nc}{Car}\PY{+w}{ }\PY{k+kd}{extends}\PY{+w}{ }\PY{n}{Vehicle}\PY{+w}{ }\PY{p}{\PYZob{}}\PY{+w}{ }\PY{p}{\PYZcb{}}\PY{+w}{          }\PY{c+c1}{// no further subclassing}
\PY{k+kd}{public}\PY{+w}{ }\PY{n}{non}\PY{o}{\PYZhy{}}\PY{k+kd}{sealed}\PY{+w}{ }\PY{k+kd}{class} \PY{n+nc}{Truck}\PY{+w}{ }\PY{k+kd}{extends}\PY{+w}{ }\PY{n}{Vehicle}\PY{+w}{ }\PY{p}{\PYZob{}}\PY{+w}{ }\PY{p}{\PYZcb{}}\PY{+w}{   }\PY{c+c1}{// anyone can extend Truck now}
\end{Verbatim}
\end{codebox}

Sealed vs \code{final} vs open (normal) inheritance:

{\tablefont
\begin{xltabular}{\linewidth}{@{}L{0.429}L{1.370}L{1.201}@{}}
\toprule
\rowcolor{tablehdr}
\thd{Modifier} & \thd{Who can extend it} & \thd{Typical use} \\
\midrule
\endhead
(none, open) & Anyone & Flexible libraries, frameworks \\
\code{final} & No one & Utility classes, immutable value types \\
\code{sealed} & Only the classes listed in \code{permits} (or same file/package) & Closed hierarchies you want exhaustive \code{switch} over \\
\code{non-\allowbreak{}sealed} & Anyone (used to reopen a branch of a sealed hierarchy) & Escape hatch inside a sealed hierarchy \\
\bottomrule
\end{xltabular}
}

Sealed types shine with pattern matching in \code{switch} (covered below):

\begin{codebox}
\begin{Verbatim}[commandchars=\\\{\}]
\PY{k+kd}{static}\PY{+w}{ }\PY{k+kt}{double}\PY{+w}{ }\PY{n+nf}{area}\PY{p}{(}\PY{n}{Shape}\PY{+w}{ }\PY{n}{shape}\PY{p}{)}\PY{+w}{ }\PY{p}{\PYZob{}}
\PY{+w}{    }\PY{k}{return}\PY{+w}{ }\PY{k}{switch}\PY{+w}{ }\PY{p}{(}\PY{n}{shape}\PY{p}{)}\PY{+w}{ }\PY{p}{\PYZob{}}
\PY{+w}{        }\PY{k}{case}\PY{+w}{ }\PY{n}{Circle}\PY{+w}{ }\PY{n}{c}\PY{+w}{ }\PY{o}{\PYZhy{}}\PY{o}{\PYZgt{}}\PY{+w}{ }\PY{n}{Math}\PY{p}{.}\PY{n+na}{PI}\PY{+w}{ }\PY{o}{*}\PY{+w}{ }\PY{n}{c}\PY{p}{.}\PY{n+na}{radius}\PY{p}{(}\PY{p}{)}\PY{+w}{ }\PY{o}{*}\PY{+w}{ }\PY{n}{c}\PY{p}{.}\PY{n+na}{radius}\PY{p}{(}\PY{p}{)}\PY{p}{;}
\PY{+w}{        }\PY{k}{case}\PY{+w}{ }\PY{n}{Square}\PY{+w}{ }\PY{n}{s}\PY{+w}{ }\PY{o}{\PYZhy{}}\PY{o}{\PYZgt{}}\PY{+w}{ }\PY{n}{s}\PY{p}{.}\PY{n+na}{side}\PY{p}{(}\PY{p}{)}\PY{+w}{ }\PY{o}{*}\PY{+w}{ }\PY{n}{s}\PY{p}{.}\PY{n+na}{side}\PY{p}{(}\PY{p}{)}\PY{p}{;}
\PY{+w}{        }\PY{k}{case}\PY{+w}{ }\PY{n}{Rectangle}\PY{+w}{ }\PY{n}{r}\PY{+w}{ }\PY{o}{\PYZhy{}}\PY{o}{\PYZgt{}}\PY{+w}{ }\PY{n}{r}\PY{p}{.}\PY{n+na}{width}\PY{p}{(}\PY{p}{)}\PY{+w}{ }\PY{o}{*}\PY{+w}{ }\PY{n}{r}\PY{p}{.}\PY{n+na}{height}\PY{p}{(}\PY{p}{)}\PY{p}{;}
\PY{+w}{        }\PY{c+c1}{// no \PYZdq{}default\PYZdq{} needed: compiler knows these are ALL the cases}
\PY{+w}{    }\PY{p}{\PYZcb{}}\PY{p}{;}
\PY{p}{\PYZcb{}}
\end{Verbatim}
\end{codebox}

\section[Interfaces (default, static, private methods)]{Interfaces (default, static, private methods)}
\label{h:3:interfaces-default-static-private-methods}
Classic interfaces (before Java 8) could only declare abstract methods — no bodies allowed. Modern interfaces can now include real code.

{\tablefont
\begin{xltabular}{\linewidth}{@{}L{0.549}L{0.426}L{0.426}L{1.234}L{2.366}@{}}
\toprule
\rowcolor{tablehdr}
\thd{Method type} & \thd{Since} & \thd{Has a body?} & \thd{Can be overridden?} & \thd{Purpose} \\
\midrule
\endhead
\code{abstract} (plain) & Java 1 & No & Yes (must be implemented) & The contract \\
\code{default} & Java 8 & Yes & Yes & Provide a default implementation so existing implementers don’t break \\
\code{static} & Java 8 & Yes & No (belongs to the interface itself) & Utility/helper methods related to the interface \\
\code{private} & Java 9 & Yes & No & Share code between default/static methods without exposing it \\
\bottomrule
\end{xltabular}
}

\begin{codebox}
\begin{Verbatim}[commandchars=\\\{\}]
\PY{k+kd}{public}\PY{+w}{ }\PY{k+kd}{interface} \PY{n+nc}{Vehicle}\PY{+w}{ }\PY{p}{\PYZob{}}

\PY{+w}{    }\PY{c+c1}{// Abstract method: every implementer must define this.}
\PY{+w}{    }\PY{k+kt}{void}\PY{+w}{ }\PY{n+nf}{drive}\PY{p}{(}\PY{p}{)}\PY{p}{;}

\PY{+w}{    }\PY{c+c1}{// Default method: implementers get this for free, but may override it.}
\PY{+w}{    }\PY{k}{default}\PY{+w}{ }\PY{k+kt}{void}\PY{+w}{ }\PY{n+nf}{honk}\PY{p}{(}\PY{p}{)}\PY{+w}{ }\PY{p}{\PYZob{}}
\PY{+w}{        }\PY{n}{System}\PY{p}{.}\PY{n+na}{out}\PY{p}{.}\PY{n+na}{println}\PY{p}{(}\PY{l+s}{\PYZdq{}}\PY{l+s}{Beep beep!}\PY{l+s}{\PYZdq{}}\PY{p}{)}\PY{p}{;}
\PY{+w}{    }\PY{p}{\PYZcb{}}

\PY{+w}{    }\PY{c+c1}{// Static method: called on the interface itself, e.g. Vehicle.create(...)}
\PY{+w}{    }\PY{k+kd}{static}\PY{+w}{ }\PY{n}{Vehicle}\PY{+w}{ }\PY{n+nf}{create}\PY{p}{(}\PY{n}{String}\PY{+w}{ }\PY{n}{type}\PY{p}{)}\PY{+w}{ }\PY{p}{\PYZob{}}
\PY{+w}{        }\PY{k}{return}\PY{+w}{ }\PY{k}{switch}\PY{+w}{ }\PY{p}{(}\PY{n}{type}\PY{p}{)}\PY{+w}{ }\PY{p}{\PYZob{}}
\PY{+w}{            }\PY{k}{case}\PY{+w}{ }\PY{l+s}{\PYZdq{}}\PY{l+s}{car}\PY{l+s}{\PYZdq{}}\PY{+w}{ }\PY{o}{\PYZhy{}}\PY{o}{\PYZgt{}}\PY{+w}{ }\PY{k}{new}\PY{+w}{ }\PY{n}{Car}\PY{p}{(}\PY{p}{)}\PY{p}{;}
\PY{+w}{            }\PY{k}{default}\PY{+w}{ }\PY{o}{\PYZhy{}}\PY{o}{\PYZgt{}}\PY{+w}{ }\PY{k}{throw}\PY{+w}{ }\PY{k}{new}\PY{+w}{ }\PY{n}{IllegalArgumentException}\PY{p}{(}\PY{l+s}{\PYZdq{}}\PY{l+s}{Unknown type: }\PY{l+s}{\PYZdq{}}\PY{+w}{ }\PY{o}{+}\PY{+w}{ }\PY{n}{type}\PY{p}{)}\PY{p}{;}
\PY{+w}{        }\PY{p}{\PYZcb{}}\PY{p}{;}
\PY{+w}{    }\PY{p}{\PYZcb{}}

\PY{+w}{    }\PY{c+c1}{// Private method (Java 9+): helper reused by default methods below.}
\PY{+w}{    }\PY{k+kd}{private}\PY{+w}{ }\PY{k+kt}{void}\PY{+w}{ }\PY{n+nf}{log}\PY{p}{(}\PY{n}{String}\PY{+w}{ }\PY{n}{action}\PY{p}{)}\PY{+w}{ }\PY{p}{\PYZob{}}
\PY{+w}{        }\PY{n}{System}\PY{p}{.}\PY{n+na}{out}\PY{p}{.}\PY{n+na}{println}\PY{p}{(}\PY{l+s}{\PYZdq{}}\PY{l+s}{[Vehicle] }\PY{l+s}{\PYZdq{}}\PY{+w}{ }\PY{o}{+}\PY{+w}{ }\PY{n}{action}\PY{p}{)}\PY{p}{;}
\PY{+w}{    }\PY{p}{\PYZcb{}}

\PY{+w}{    }\PY{k}{default}\PY{+w}{ }\PY{k+kt}{void}\PY{+w}{ }\PY{n+nf}{startTrip}\PY{p}{(}\PY{p}{)}\PY{+w}{ }\PY{p}{\PYZob{}}
\PY{+w}{        }\PY{n}{log}\PY{p}{(}\PY{l+s}{\PYZdq{}}\PY{l+s}{starting trip}\PY{l+s}{\PYZdq{}}\PY{p}{)}\PY{p}{;}
\PY{+w}{        }\PY{n}{drive}\PY{p}{(}\PY{p}{)}\PY{p}{;}
\PY{+w}{    }\PY{p}{\PYZcb{}}
\PY{p}{\PYZcb{}}

\PY{k+kd}{class} \PY{n+nc}{Car}\PY{+w}{ }\PY{k+kd}{implements}\PY{+w}{ }\PY{n}{Vehicle}\PY{+w}{ }\PY{p}{\PYZob{}}
\PY{+w}{    }\PY{k+kd}{public}\PY{+w}{ }\PY{k+kt}{void}\PY{+w}{ }\PY{n+nf}{drive}\PY{p}{(}\PY{p}{)}\PY{+w}{ }\PY{p}{\PYZob{}}
\PY{+w}{        }\PY{n}{System}\PY{p}{.}\PY{n+na}{out}\PY{p}{.}\PY{n+na}{println}\PY{p}{(}\PY{l+s}{\PYZdq{}}\PY{l+s}{Car is driving}\PY{l+s}{\PYZdq{}}\PY{p}{)}\PY{p}{;}
\PY{+w}{    }\PY{p}{\PYZcb{}}
\PY{p}{\PYZcb{}}
\end{Verbatim}
\end{codebox}

Why this matters for code review: \code{default} methods let library authors add new methods to an interface \textbf{without breaking existing implementations} (this is exactly why \code{Collection} and \code{Map} gained many \code{default} methods over the years, like \code{forEach} and \code{getOrDefault}). Watch out, though — if a class implements two interfaces with the same \code{default} method signature, the class must override it to resolve the conflict, or the code will not compile.

\begin{codebox}
\begin{Verbatim}[commandchars=\\\{\}]
\PY{k+kd}{interface} \PY{n+nc}{A}\PY{+w}{ }\PY{p}{\PYZob{}}\PY{+w}{ }\PY{k}{default}\PY{+w}{ }\PY{n}{String}\PY{+w}{ }\PY{n+nf}{name}\PY{p}{(}\PY{p}{)}\PY{+w}{ }\PY{p}{\PYZob{}}\PY{+w}{ }\PY{k}{return}\PY{+w}{ }\PY{l+s}{\PYZdq{}}\PY{l+s}{A}\PY{l+s}{\PYZdq{}}\PY{p}{;}\PY{+w}{ }\PY{p}{\PYZcb{}}\PY{+w}{ }\PY{p}{\PYZcb{}}
\PY{k+kd}{interface} \PY{n+nc}{B}\PY{+w}{ }\PY{p}{\PYZob{}}\PY{+w}{ }\PY{k}{default}\PY{+w}{ }\PY{n}{String}\PY{+w}{ }\PY{n+nf}{name}\PY{p}{(}\PY{p}{)}\PY{+w}{ }\PY{p}{\PYZob{}}\PY{+w}{ }\PY{k}{return}\PY{+w}{ }\PY{l+s}{\PYZdq{}}\PY{l+s}{B}\PY{l+s}{\PYZdq{}}\PY{p}{;}\PY{+w}{ }\PY{p}{\PYZcb{}}\PY{+w}{ }\PY{p}{\PYZcb{}}

\PY{k+kd}{class} \PY{n+nc}{C}\PY{+w}{ }\PY{k+kd}{implements}\PY{+w}{ }\PY{n}{A}\PY{p}{,}\PY{+w}{ }\PY{n}{B}\PY{+w}{ }\PY{p}{\PYZob{}}
\PY{+w}{    }\PY{c+c1}{// Must override: ambiguous otherwise.}
\PY{+w}{    }\PY{k+kd}{public}\PY{+w}{ }\PY{n}{String}\PY{+w}{ }\PY{n+nf}{name}\PY{p}{(}\PY{p}{)}\PY{+w}{ }\PY{p}{\PYZob{}}
\PY{+w}{        }\PY{k}{return}\PY{+w}{ }\PY{n}{A}\PY{p}{.}\PY{n+na}{super}\PY{p}{.}\PY{n+na}{name}\PY{p}{(}\PY{p}{)}\PY{+w}{ }\PY{o}{+}\PY{+w}{ }\PY{n}{B}\PY{p}{.}\PY{n+na}{super}\PY{p}{.}\PY{n+na}{name}\PY{p}{(}\PY{p}{)}\PY{p}{;}\PY{+w}{ }\PY{c+c1}{// \PYZdq{}AB\PYZdq{}}
\PY{+w}{    }\PY{p}{\PYZcb{}}
\PY{p}{\PYZcb{}}
\end{Verbatim}
\end{codebox}

\section[Abstract Classes]{Abstract Classes}
\label{h:3:abstract-classes}
An \textbf{abstract class} cannot be instantiated directly. It exists to be extended. It can mix fully implemented methods with abstract ones that subclasses must fill in. Abstract classes have existed since Java 1, but it’s worth revisiting how they compare to modern interfaces since both can now hold real code.

\begin{codebox}
\begin{Verbatim}[commandchars=\\\{\}]
\PY{k+kd}{public}\PY{+w}{ }\PY{k+kd}{abstract}\PY{+w}{ }\PY{k+kd}{class} \PY{n+nc}{Employee}\PY{+w}{ }\PY{p}{\PYZob{}}

\PY{+w}{    }\PY{k+kd}{protected}\PY{+w}{ }\PY{k+kd}{final}\PY{+w}{ }\PY{n}{String}\PY{+w}{ }\PY{n}{name}\PY{p}{;}

\PY{+w}{    }\PY{k+kd}{protected}\PY{+w}{ }\PY{n+nf}{Employee}\PY{p}{(}\PY{n}{String}\PY{+w}{ }\PY{n}{name}\PY{p}{)}\PY{+w}{ }\PY{p}{\PYZob{}}
\PY{+w}{        }\PY{k}{this}\PY{p}{.}\PY{n+na}{name}\PY{+w}{ }\PY{o}{=}\PY{+w}{ }\PY{n}{name}\PY{p}{;}
\PY{+w}{    }\PY{p}{\PYZcb{}}

\PY{+w}{    }\PY{c+c1}{// Concrete method: shared by all subclasses.}
\PY{+w}{    }\PY{k+kd}{public}\PY{+w}{ }\PY{n}{String}\PY{+w}{ }\PY{n+nf}{greet}\PY{p}{(}\PY{p}{)}\PY{+w}{ }\PY{p}{\PYZob{}}
\PY{+w}{        }\PY{k}{return}\PY{+w}{ }\PY{l+s}{\PYZdq{}}\PY{l+s}{Hello, my name is }\PY{l+s}{\PYZdq{}}\PY{+w}{ }\PY{o}{+}\PY{+w}{ }\PY{n}{name}\PY{p}{;}
\PY{+w}{    }\PY{p}{\PYZcb{}}

\PY{+w}{    }\PY{c+c1}{// Abstract method: each subclass must define its own pay calculation.}
\PY{+w}{    }\PY{k+kd}{public}\PY{+w}{ }\PY{k+kd}{abstract}\PY{+w}{ }\PY{k+kt}{double}\PY{+w}{ }\PY{n+nf}{calculatePay}\PY{p}{(}\PY{p}{)}\PY{p}{;}
\PY{p}{\PYZcb{}}

\PY{k+kd}{public}\PY{+w}{ }\PY{k+kd}{class} \PY{n+nc}{Manager}\PY{+w}{ }\PY{k+kd}{extends}\PY{+w}{ }\PY{n}{Employee}\PY{+w}{ }\PY{p}{\PYZob{}}
\PY{+w}{    }\PY{k+kd}{private}\PY{+w}{ }\PY{k+kd}{final}\PY{+w}{ }\PY{k+kt}{double}\PY{+w}{ }\PY{n}{baseSalary}\PY{p}{;}
\PY{+w}{    }\PY{k+kd}{private}\PY{+w}{ }\PY{k+kd}{final}\PY{+w}{ }\PY{k+kt}{double}\PY{+w}{ }\PY{n}{bonus}\PY{p}{;}

\PY{+w}{    }\PY{k+kd}{public}\PY{+w}{ }\PY{n+nf}{Manager}\PY{p}{(}\PY{n}{String}\PY{+w}{ }\PY{n}{name}\PY{p}{,}\PY{+w}{ }\PY{k+kt}{double}\PY{+w}{ }\PY{n}{baseSalary}\PY{p}{,}\PY{+w}{ }\PY{k+kt}{double}\PY{+w}{ }\PY{n}{bonus}\PY{p}{)}\PY{+w}{ }\PY{p}{\PYZob{}}
\PY{+w}{        }\PY{k+kd}{super}\PY{p}{(}\PY{n}{name}\PY{p}{)}\PY{p}{;}
\PY{+w}{        }\PY{k}{this}\PY{p}{.}\PY{n+na}{baseSalary}\PY{+w}{ }\PY{o}{=}\PY{+w}{ }\PY{n}{baseSalary}\PY{p}{;}
\PY{+w}{        }\PY{k}{this}\PY{p}{.}\PY{n+na}{bonus}\PY{+w}{ }\PY{o}{=}\PY{+w}{ }\PY{n}{bonus}\PY{p}{;}
\PY{+w}{    }\PY{p}{\PYZcb{}}

\PY{+w}{    }\PY{n+nd}{@Override}
\PY{+w}{    }\PY{k+kd}{public}\PY{+w}{ }\PY{k+kt}{double}\PY{+w}{ }\PY{n+nf}{calculatePay}\PY{p}{(}\PY{p}{)}\PY{+w}{ }\PY{p}{\PYZob{}}
\PY{+w}{        }\PY{k}{return}\PY{+w}{ }\PY{n}{baseSalary}\PY{+w}{ }\PY{o}{+}\PY{+w}{ }\PY{n}{bonus}\PY{p}{;}
\PY{+w}{    }\PY{p}{\PYZcb{}}
\PY{p}{\PYZcb{}}
\end{Verbatim}
\end{codebox}

Abstract class vs. interface — the modern comparison:

{\tablefont
\begin{xltabular}{\linewidth}{@{}L{0.480}L{1.056}L{1.464}@{}}
\toprule
\rowcolor{tablehdr}
\thd{Aspect} & \thd{Abstract class} & \thd{Interface} \\
\midrule
\endhead
Instance fields & Yes, any visibility & No instance fields (only \code{static~\allowbreak{}final} constants) \\
Constructors & Yes & No \\
Multiple inheritance & A class can extend only \textbf{one} abstract class & A class can implement \textbf{many} interfaces \\
Method bodies & Any method can have a body & \code{default}, \code{static}, \code{private} can have bodies; plain methods cannot \\
Access modifiers on methods & \code{public}, \code{protected}, \code{private} all allowed & Methods are implicitly \code{public} (except \code{private} helper methods) \\
State/encapsulation & Good fit — designed to hold and manage state & Poor fit — designed to describe capability/contract \\
When to use & “Is-a” relationship with shared state and shared code & “Can-do” capability, especially across unrelated classes \\
\bottomrule
\end{xltabular}
}

Rule of thumb for reviews: if two classes need to \textbf{share fields or a constructor}, lean toward an abstract class. If you’re just describing a \textbf{capability} (e.g. \code{Comparable}, \code{Runnable}, \code{Serializable}), use an interface.

\section[Pattern Matching (instanceof, switch)]{Pattern Matching (instanceof, switch)}
\label{h:3:pattern-matching-instanceof-switch}
\textbf{Pattern matching} lets you check the shape or type of a value and extract its parts in one step, instead of checking the type and then manually casting.

\subsection[instanceof pattern matching (Java 16, finalized)]{\code{instanceof} pattern matching (Java 16, finalized)}
\label{h:3:instanceof-pattern-matching-java-16-finalized}
Before:

\begin{codebox}
\begin{Verbatim}[commandchars=\\\{\}]
\PY{k}{if}\PY{+w}{ }\PY{p}{(}\PY{n}{obj}\PY{+w}{ }\PY{k}{instanceof}\PY{+w}{ }\PY{n}{String}\PY{p}{)}\PY{+w}{ }\PY{p}{\PYZob{}}
\PY{+w}{    }\PY{n}{String}\PY{+w}{ }\PY{n}{s}\PY{+w}{ }\PY{o}{=}\PY{+w}{ }\PY{p}{(}\PY{n}{String}\PY{p}{)}\PY{+w}{ }\PY{n}{obj}\PY{p}{;}\PY{+w}{ }\PY{c+c1}{// manual, error\PYZhy{}prone cast}
\PY{+w}{    }\PY{n}{System}\PY{p}{.}\PY{n+na}{out}\PY{p}{.}\PY{n+na}{println}\PY{p}{(}\PY{n}{s}\PY{p}{.}\PY{n+na}{length}\PY{p}{(}\PY{p}{)}\PY{p}{)}\PY{p}{;}
\PY{p}{\PYZcb{}}
\end{Verbatim}
\end{codebox}

After — the cast is automatic, and the variable \code{s} is only in scope where it makes sense:

\begin{codebox}
\begin{Verbatim}[commandchars=\\\{\}]
\PY{k}{if}\PY{+w}{ }\PY{p}{(}\PY{n}{obj}\PY{+w}{ }\PY{k}{instanceof}\PY{+w}{ }\PY{n}{String}\PY{+w}{ }\PY{n}{s}\PY{p}{)}\PY{+w}{ }\PY{p}{\PYZob{}}
\PY{+w}{    }\PY{n}{System}\PY{p}{.}\PY{n+na}{out}\PY{p}{.}\PY{n+na}{println}\PY{p}{(}\PY{n}{s}\PY{p}{.}\PY{n+na}{length}\PY{p}{(}\PY{p}{)}\PY{p}{)}\PY{p}{;}\PY{+w}{ }\PY{c+c1}{// \PYZdq{}s\PYZdq{} is already a String here}
\PY{p}{\PYZcb{}}
\end{Verbatim}
\end{codebox}

You can combine it with extra conditions:

\begin{codebox}
\begin{Verbatim}[commandchars=\\\{\}]
\PY{k}{if}\PY{+w}{ }\PY{p}{(}\PY{n}{obj}\PY{+w}{ }\PY{k}{instanceof}\PY{+w}{ }\PY{n}{String}\PY{+w}{ }\PY{n}{s}\PY{+w}{ }\PY{o}{\PYZam{}}\PY{o}{\PYZam{}}\PY{+w}{ }\PY{o}{!}\PY{n}{s}\PY{p}{.}\PY{n+na}{isBlank}\PY{p}{(}\PY{p}{)}\PY{p}{)}\PY{+w}{ }\PY{p}{\PYZob{}}
\PY{+w}{    }\PY{n}{System}\PY{p}{.}\PY{n+na}{out}\PY{p}{.}\PY{n+na}{println}\PY{p}{(}\PY{l+s}{\PYZdq{}}\PY{l+s}{Non\PYZhy{}blank string: }\PY{l+s}{\PYZdq{}}\PY{+w}{ }\PY{o}{+}\PY{+w}{ }\PY{n}{s}\PY{p}{)}\PY{p}{;}
\PY{p}{\PYZcb{}}
\end{Verbatim}
\end{codebox}

\subsection[switch pattern matching (Java 21, finalized)]{\code{switch} pattern matching (Java 21, finalized)}
\label{h:3:switch-pattern-matching-java-21-finalized}
Pattern matching for \code{switch} lets you match on type (and, with records, on structure) directly in a \code{switch}.

\begin{codebox}
\begin{Verbatim}[commandchars=\\\{\}]
\PY{k+kd}{sealed}\PY{+w}{ }\PY{k+kd}{interface} \PY{n+nc}{Shape}\PY{+w}{ }\PY{n}{permits}\PY{+w}{ }\PY{n}{Circle}\PY{p}{,}\PY{+w}{ }\PY{n}{Square}\PY{+w}{ }\PY{p}{\PYZob{}}\PY{+w}{ }\PY{p}{\PYZcb{}}
\PY{k+kd}{record} \PY{n+nc}{Circle}\PY{p}{(}\PY{k+kt}{double}\PY{+w}{ }\PY{n}{radius}\PY{p}{)}\PY{+w}{ }\PY{k+kd}{implements}\PY{+w}{ }\PY{n}{Shape}\PY{+w}{ }\PY{p}{\PYZob{}}\PY{+w}{ }\PY{p}{\PYZcb{}}
\PY{k+kd}{record} \PY{n+nc}{Square}\PY{p}{(}\PY{k+kt}{double}\PY{+w}{ }\PY{n}{side}\PY{p}{)}\PY{+w}{ }\PY{k+kd}{implements}\PY{+w}{ }\PY{n}{Shape}\PY{+w}{ }\PY{p}{\PYZob{}}\PY{+w}{ }\PY{p}{\PYZcb{}}

\PY{k+kd}{static}\PY{+w}{ }\PY{n}{String}\PY{+w}{ }\PY{n+nf}{describe}\PY{p}{(}\PY{n}{Shape}\PY{+w}{ }\PY{n}{shape}\PY{p}{)}\PY{+w}{ }\PY{p}{\PYZob{}}
\PY{+w}{    }\PY{k}{return}\PY{+w}{ }\PY{k}{switch}\PY{+w}{ }\PY{p}{(}\PY{n}{shape}\PY{p}{)}\PY{+w}{ }\PY{p}{\PYZob{}}
\PY{+w}{        }\PY{k}{case}\PY{+w}{ }\PY{n}{Circle}\PY{+w}{ }\PY{n}{c}\PY{+w}{ }\PY{n}{when}\PY{+w}{ }\PY{n}{c}\PY{p}{.}\PY{n+na}{radius}\PY{p}{(}\PY{p}{)}\PY{+w}{ }\PY{o}{\PYZgt{}}\PY{+w}{ }\PY{l+m+mi}{100}\PY{+w}{ }\PY{o}{\PYZhy{}}\PY{o}{\PYZgt{}}\PY{+w}{ }\PY{l+s}{\PYZdq{}}\PY{l+s}{huge circle}\PY{l+s}{\PYZdq{}}\PY{p}{;}\PY{+w}{   }\PY{c+c1}{// guarded pattern}
\PY{+w}{        }\PY{k}{case}\PY{+w}{ }\PY{n}{Circle}\PY{+w}{ }\PY{n}{c}\PY{+w}{                        }\PY{o}{\PYZhy{}}\PY{o}{\PYZgt{}}\PY{+w}{ }\PY{l+s}{\PYZdq{}}\PY{l+s}{circle r=}\PY{l+s}{\PYZdq{}}\PY{+w}{ }\PY{o}{+}\PY{+w}{ }\PY{n}{c}\PY{p}{.}\PY{n+na}{radius}\PY{p}{(}\PY{p}{)}\PY{p}{;}
\PY{+w}{        }\PY{k}{case}\PY{+w}{ }\PY{n}{Square}\PY{+w}{ }\PY{n}{s}\PY{+w}{                        }\PY{o}{\PYZhy{}}\PY{o}{\PYZgt{}}\PY{+w}{ }\PY{l+s}{\PYZdq{}}\PY{l+s}{square side=}\PY{l+s}{\PYZdq{}}\PY{+w}{ }\PY{o}{+}\PY{+w}{ }\PY{n}{s}\PY{p}{.}\PY{n+na}{side}\PY{p}{(}\PY{p}{)}\PY{p}{;}
\PY{+w}{    }\PY{p}{\PYZcb{}}\PY{p}{;}
\PY{p}{\PYZcb{}}
\end{Verbatim}
\end{codebox}

\textbf{Record patterns} (Java 21) let you destructure a record right in the pattern, pulling out its components:

\begin{codebox}
\begin{Verbatim}[commandchars=\\\{\}]
\PY{k+kd}{record} \PY{n+nc}{Point}\PY{p}{(}\PY{k+kt}{int}\PY{+w}{ }\PY{n}{x}\PY{p}{,}\PY{+w}{ }\PY{k+kt}{int}\PY{+w}{ }\PY{n}{y}\PY{p}{)}\PY{+w}{ }\PY{p}{\PYZob{}}\PY{+w}{ }\PY{p}{\PYZcb{}}
\PY{k+kd}{record} \PY{n+nc}{Line}\PY{p}{(}\PY{n}{Point}\PY{+w}{ }\PY{n}{start}\PY{p}{,}\PY{+w}{ }\PY{n}{Point}\PY{+w}{ }\PY{n}{end}\PY{p}{)}\PY{+w}{ }\PY{p}{\PYZob{}}\PY{+w}{ }\PY{p}{\PYZcb{}}

\PY{k+kd}{static}\PY{+w}{ }\PY{n}{String}\PY{+w}{ }\PY{n+nf}{describe}\PY{p}{(}\PY{n}{Object}\PY{+w}{ }\PY{n}{obj}\PY{p}{)}\PY{+w}{ }\PY{p}{\PYZob{}}
\PY{+w}{    }\PY{k}{return}\PY{+w}{ }\PY{k}{switch}\PY{+w}{ }\PY{p}{(}\PY{n}{obj}\PY{p}{)}\PY{+w}{ }\PY{p}{\PYZob{}}
\PY{+w}{        }\PY{c+c1}{// Nested destructuring: pull x1,y1,x2,y2 straight out of the Line.}
\PY{+w}{        }\PY{k}{case}\PY{+w}{ }\PY{n}{Line}\PY{p}{(}\PY{n}{Point}\PY{p}{(}\PY{k+kd}{var}\PY{+w}{ }\PY{n}{x1}\PY{p}{,}\PY{+w}{ }\PY{k+kd}{var}\PY{+w}{ }\PY{n}{y1}\PY{p}{)}\PY{p}{,}\PY{+w}{ }\PY{n}{Point}\PY{p}{(}\PY{k+kd}{var}\PY{+w}{ }\PY{n}{x2}\PY{p}{,}\PY{+w}{ }\PY{k+kd}{var}\PY{+w}{ }\PY{n}{y2}\PY{p}{)}\PY{p}{)}\PY{+w}{ }\PY{o}{\PYZhy{}}\PY{o}{\PYZgt{}}
\PY{+w}{            }\PY{l+s}{\PYZdq{}}\PY{l+s}{Line from (\PYZpc{}d,\PYZpc{}d) to (\PYZpc{}d,\PYZpc{}d)}\PY{l+s}{\PYZdq{}}\PY{p}{.}\PY{n+na}{formatted}\PY{p}{(}\PY{n}{x1}\PY{p}{,}\PY{+w}{ }\PY{n}{y1}\PY{p}{,}\PY{+w}{ }\PY{n}{x2}\PY{p}{,}\PY{+w}{ }\PY{n}{y2}\PY{p}{)}\PY{p}{;}
\PY{+w}{        }\PY{k}{case}\PY{+w}{ }\PY{n}{Point}\PY{p}{(}\PY{k+kd}{var}\PY{+w}{ }\PY{n}{x}\PY{p}{,}\PY{+w}{ }\PY{k+kd}{var}\PY{+w}{ }\PY{n}{y}\PY{p}{)}\PY{+w}{ }\PY{o}{\PYZhy{}}\PY{o}{\PYZgt{}}\PY{+w}{ }\PY{l+s}{\PYZdq{}}\PY{l+s}{Point(\PYZpc{}d,\PYZpc{}d)}\PY{l+s}{\PYZdq{}}\PY{p}{.}\PY{n+na}{formatted}\PY{p}{(}\PY{n}{x}\PY{p}{,}\PY{+w}{ }\PY{n}{y}\PY{p}{)}\PY{p}{;}
\PY{+w}{        }\PY{k}{default}\PY{+w}{ }\PY{o}{\PYZhy{}}\PY{o}{\PYZgt{}}\PY{+w}{ }\PY{l+s}{\PYZdq{}}\PY{l+s}{unknown}\PY{l+s}{\PYZdq{}}\PY{p}{;}
\PY{+w}{    }\PY{p}{\PYZcb{}}\PY{p}{;}
\PY{p}{\PYZcb{}}
\end{Verbatim}
\end{codebox}

Timeline summary:

{\tablefont
\begin{xltabular}{\linewidth}{@{}L{1.958}L{0.542}L{0.500}@{}}
\toprule
\rowcolor{tablehdr}
\thd{Feature} & \thd{Preview since} & \thd{Finalized in} \\
\midrule
\endhead
\code{instanceof} pattern matching & Java 14 & Java 16 \\
\code{switch} pattern matching (type patterns, guards) & Java 17 & Java 21 \\
Record patterns (destructuring) & Java 19 & Java 21 \\
\bottomrule
\end{xltabular}
}

\section[Primitive Types in Patterns]{Primitive Types in Patterns}
\label{h:3:primitive-types-in-patterns}
Historically, patterns (in \code{instanceof} and \code{switch}) only worked with reference types (objects), not primitives like \code{int} or \code{double}. \textbf{Primitive type patterns} extend pattern matching to primitives, including inside record patterns.

This feature is a \textbf{preview feature starting in JDK 23} (JEP 455, "Primitive Types in Patterns, \code{instanceof}, and \code{switch\textquotedbl{}).\allowbreak{}~\allowbreak{}It~\allowbreak{}is~\allowbreak{}still~\allowbreak{}evolving~\allowbreak{}—~\allowbreak{}as~\allowbreak{}of~\allowbreak{}Java~\allowbreak{}25~\allowbreak{}it~\allowbreak{}remains~\allowbreak{}a~\allowbreak{}preview~\allowbreak{}feature~\allowbreak{}(\allowbreak{}re-\allowbreak{}previewed~\allowbreak{}as~\allowbreak{}JEP~\allowbreak{}507).\allowbreak{}~\allowbreak{}You~\allowbreak{}must~\allowbreak{}compile~\allowbreak{}and~\allowbreak{}run~\allowbreak{}with~}–enable-preview` to use it.

\begin{codebox}
\begin{Verbatim}[commandchars=\\\{\}]
\PY{c+c1}{// Requires: javac \PYZhy{}\PYZhy{}release 23 \PYZhy{}\PYZhy{}enable\PYZhy{}preview Main.java}
\PY{c+c1}{//           java \PYZhy{}\PYZhy{}enable\PYZhy{}preview Main}

\PY{n}{Object}\PY{+w}{ }\PY{n}{obj}\PY{+w}{ }\PY{o}{=}\PY{+w}{ }\PY{l+m+mi}{42}\PY{p}{;}

\PY{c+c1}{// instanceof with a primitive pattern (preview):}
\PY{k}{if}\PY{+w}{ }\PY{p}{(}\PY{n}{obj}\PY{+w}{ }\PY{k}{instanceof}\PY{+w}{ }\PY{k+kt}{int}\PY{+w}{ }\PY{n}{i}\PY{p}{)}\PY{+w}{ }\PY{p}{\PYZob{}}
\PY{+w}{    }\PY{n}{System}\PY{p}{.}\PY{n+na}{out}\PY{p}{.}\PY{n+na}{println}\PY{p}{(}\PY{l+s}{\PYZdq{}}\PY{l+s}{It\PYZsq{}s an int: }\PY{l+s}{\PYZdq{}}\PY{+w}{ }\PY{o}{+}\PY{+w}{ }\PY{n}{i}\PY{p}{)}\PY{p}{;}
\PY{p}{\PYZcb{}}
\end{Verbatim}
\end{codebox}

It also removes an old rough edge: \code{switch} on primitives couldn’t easily be combined with ranges or exact numeric matching against boxed types. With primitive patterns, you get natural, type-safe matching:

\begin{codebox}
\begin{Verbatim}[commandchars=\\\{\}]
\PY{k+kd}{static}\PY{+w}{ }\PY{n}{String}\PY{+w}{ }\PY{n+nf}{classify}\PY{p}{(}\PY{n}{Object}\PY{+w}{ }\PY{n}{o}\PY{p}{)}\PY{+w}{ }\PY{p}{\PYZob{}}
\PY{+w}{    }\PY{k}{return}\PY{+w}{ }\PY{k}{switch}\PY{+w}{ }\PY{p}{(}\PY{n}{o}\PY{p}{)}\PY{+w}{ }\PY{p}{\PYZob{}}
\PY{+w}{        }\PY{k}{case}\PY{+w}{ }\PY{k+kt}{int}\PY{+w}{ }\PY{n}{i}\PY{+w}{ }\PY{n}{when}\PY{+w}{ }\PY{n}{i}\PY{+w}{ }\PY{o}{\PYZlt{}}\PY{+w}{ }\PY{l+m+mi}{0}\PY{+w}{    }\PY{o}{\PYZhy{}}\PY{o}{\PYZgt{}}\PY{+w}{ }\PY{l+s}{\PYZdq{}}\PY{l+s}{negative int}\PY{l+s}{\PYZdq{}}\PY{p}{;}
\PY{+w}{        }\PY{k}{case}\PY{+w}{ }\PY{k+kt}{int}\PY{+w}{ }\PY{n}{i}\PY{+w}{               }\PY{o}{\PYZhy{}}\PY{o}{\PYZgt{}}\PY{+w}{ }\PY{l+s}{\PYZdq{}}\PY{l+s}{non\PYZhy{}negative int: }\PY{l+s}{\PYZdq{}}\PY{+w}{ }\PY{o}{+}\PY{+w}{ }\PY{n}{i}\PY{p}{;}
\PY{+w}{        }\PY{k}{case}\PY{+w}{ }\PY{k+kt}{long}\PY{+w}{ }\PY{n}{l}\PY{+w}{               }\PY{o}{\PYZhy{}}\PY{o}{\PYZgt{}}\PY{+w}{ }\PY{l+s}{\PYZdq{}}\PY{l+s}{a long: }\PY{l+s}{\PYZdq{}}\PY{+w}{ }\PY{o}{+}\PY{+w}{ }\PY{n}{l}\PY{p}{;}
\PY{+w}{        }\PY{k}{case}\PY{+w}{ }\PY{k+kt}{double}\PY{+w}{ }\PY{n}{d}\PY{+w}{             }\PY{o}{\PYZhy{}}\PY{o}{\PYZgt{}}\PY{+w}{ }\PY{l+s}{\PYZdq{}}\PY{l+s}{a double: }\PY{l+s}{\PYZdq{}}\PY{+w}{ }\PY{o}{+}\PY{+w}{ }\PY{n}{d}\PY{p}{;}
\PY{+w}{        }\PY{k}{default}\PY{+w}{                   }\PY{o}{\PYZhy{}}\PY{o}{\PYZgt{}}\PY{+w}{ }\PY{l+s}{\PYZdq{}}\PY{l+s}{something else}\PY{l+s}{\PYZdq{}}\PY{p}{;}
\PY{+w}{    }\PY{p}{\PYZcb{}}\PY{p}{;}
\PY{p}{\PYZcb{}}
\end{Verbatim}
\end{codebox}

Because this is preview, do not rely on it in production code yet; the syntax could still change before final release. Always check the JDK release notes before using it on a real project.

\begin{codebox}
\begin{Verbatim}[commandchars=\\\{\}]
\PY{c+c1}{// Compiling without \PYZhy{}\PYZhy{}enable\PYZhy{}preview gives an error like:}
\PY{c+c1}{// error: primitive patterns in switch and instanceof are a preview feature}
\PY{c+c1}{//        and are disabled by default.}
\end{Verbatim}
\end{codebox}

\section[Switch Expressions]{Switch Expressions}
\label{h:3:switch-expressions}
A traditional \code{switch} is a \textbf{statement}: it doesn’t produce a value, and it needs \code{break} to avoid falling through to the next case. A \textbf{switch expression} (finalized in \textbf{Java 14}) produces a value directly, uses \code{-\textgreater{}} (arrow) syntax, and has no fall-through.

\begin{codebox}
\begin{Verbatim}[commandchars=\\\{\}]
\PY{c+c1}{// Old style: switch STATEMENT}
\PY{k+kt}{int}\PY{+w}{ }\PY{n}{day}\PY{+w}{ }\PY{o}{=}\PY{+w}{ }\PY{l+m+mi}{3}\PY{p}{;}
\PY{n}{String}\PY{+w}{ }\PY{n}{name}\PY{p}{;}
\PY{k}{switch}\PY{+w}{ }\PY{p}{(}\PY{n}{day}\PY{p}{)}\PY{+w}{ }\PY{p}{\PYZob{}}
\PY{+w}{    }\PY{k}{case}\PY{+w}{ }\PY{l+m+mi}{1}\PY{p}{:}
\PY{+w}{        }\PY{n}{name}\PY{+w}{ }\PY{o}{=}\PY{+w}{ }\PY{l+s}{\PYZdq{}}\PY{l+s}{Monday}\PY{l+s}{\PYZdq{}}\PY{p}{;}
\PY{+w}{        }\PY{k}{break}\PY{p}{;}
\PY{+w}{    }\PY{k}{case}\PY{+w}{ }\PY{l+m+mi}{2}\PY{p}{:}
\PY{+w}{        }\PY{n}{name}\PY{+w}{ }\PY{o}{=}\PY{+w}{ }\PY{l+s}{\PYZdq{}}\PY{l+s}{Tuesday}\PY{l+s}{\PYZdq{}}\PY{p}{;}
\PY{+w}{        }\PY{k}{break}\PY{p}{;}
\PY{+w}{    }\PY{k}{default}\PY{p}{:}
\PY{+w}{        }\PY{n}{name}\PY{+w}{ }\PY{o}{=}\PY{+w}{ }\PY{l+s}{\PYZdq{}}\PY{l+s}{Unknown}\PY{l+s}{\PYZdq{}}\PY{p}{;}
\PY{+w}{        }\PY{k}{break}\PY{p}{;}
\PY{p}{\PYZcb{}}
\end{Verbatim}
\end{codebox}

\begin{codebox}
\begin{Verbatim}[commandchars=\\\{\}]
\PY{c+c1}{// New style: switch EXPRESSION}
\PY{k+kt}{int}\PY{+w}{ }\PY{n}{day}\PY{+w}{ }\PY{o}{=}\PY{+w}{ }\PY{l+m+mi}{3}\PY{p}{;}
\PY{n}{String}\PY{+w}{ }\PY{n}{name}\PY{+w}{ }\PY{o}{=}\PY{+w}{ }\PY{k}{switch}\PY{+w}{ }\PY{p}{(}\PY{n}{day}\PY{p}{)}\PY{+w}{ }\PY{p}{\PYZob{}}
\PY{+w}{    }\PY{k}{case}\PY{+w}{ }\PY{l+m+mi}{1}\PY{+w}{ }\PY{o}{\PYZhy{}}\PY{o}{\PYZgt{}}\PY{+w}{ }\PY{l+s}{\PYZdq{}}\PY{l+s}{Monday}\PY{l+s}{\PYZdq{}}\PY{p}{;}
\PY{+w}{    }\PY{k}{case}\PY{+w}{ }\PY{l+m+mi}{2}\PY{+w}{ }\PY{o}{\PYZhy{}}\PY{o}{\PYZgt{}}\PY{+w}{ }\PY{l+s}{\PYZdq{}}\PY{l+s}{Tuesday}\PY{l+s}{\PYZdq{}}\PY{p}{;}
\PY{+w}{    }\PY{k}{default}\PY{+w}{ }\PY{o}{\PYZhy{}}\PY{o}{\PYZgt{}}\PY{+w}{ }\PY{l+s}{\PYZdq{}}\PY{l+s}{Unknown}\PY{l+s}{\PYZdq{}}\PY{p}{;}
\PY{p}{\PYZcb{}}\PY{p}{;}
\end{Verbatim}
\end{codebox}

Multiple case labels can share one arrow:

\begin{codebox}
\begin{Verbatim}[commandchars=\\\{\}]
\PY{n}{String}\PY{+w}{ }\PY{n}{category}\PY{+w}{ }\PY{o}{=}\PY{+w}{ }\PY{k}{switch}\PY{+w}{ }\PY{p}{(}\PY{n}{day}\PY{p}{)}\PY{+w}{ }\PY{p}{\PYZob{}}
\PY{+w}{    }\PY{k}{case}\PY{+w}{ }\PY{l+m+mi}{1}\PY{p}{,}\PY{+w}{ }\PY{l+m+mi}{2}\PY{p}{,}\PY{+w}{ }\PY{l+m+mi}{3}\PY{p}{,}\PY{+w}{ }\PY{l+m+mi}{4}\PY{p}{,}\PY{+w}{ }\PY{l+m+mi}{5}\PY{+w}{ }\PY{o}{\PYZhy{}}\PY{o}{\PYZgt{}}\PY{+w}{ }\PY{l+s}{\PYZdq{}}\PY{l+s}{Weekday}\PY{l+s}{\PYZdq{}}\PY{p}{;}
\PY{+w}{    }\PY{k}{case}\PY{+w}{ }\PY{l+m+mi}{6}\PY{p}{,}\PY{+w}{ }\PY{l+m+mi}{7}\PY{+w}{ }\PY{o}{\PYZhy{}}\PY{o}{\PYZgt{}}\PY{+w}{ }\PY{l+s}{\PYZdq{}}\PY{l+s}{Weekend}\PY{l+s}{\PYZdq{}}\PY{p}{;}
\PY{+w}{    }\PY{k}{default}\PY{+w}{ }\PY{o}{\PYZhy{}}\PY{o}{\PYZgt{}}\PY{+w}{ }\PY{l+s}{\PYZdq{}}\PY{l+s}{Invalid}\PY{l+s}{\PYZdq{}}\PY{p}{;}
\PY{p}{\PYZcb{}}\PY{p}{;}
\end{Verbatim}
\end{codebox}

When a case needs more than one statement, use a block and \code{yield} to return the value:

\begin{codebox}
\begin{Verbatim}[commandchars=\\\{\}]
\PY{k+kt}{int}\PY{+w}{ }\PY{n}{score}\PY{+w}{ }\PY{o}{=}\PY{+w}{ }\PY{l+m+mi}{85}\PY{p}{;}
\PY{n}{String}\PY{+w}{ }\PY{n}{grade}\PY{+w}{ }\PY{o}{=}\PY{+w}{ }\PY{k}{switch}\PY{+w}{ }\PY{p}{(}\PY{n}{score}\PY{+w}{ }\PY{o}{/}\PY{+w}{ }\PY{l+m+mi}{10}\PY{p}{)}\PY{+w}{ }\PY{p}{\PYZob{}}
\PY{+w}{    }\PY{k}{case}\PY{+w}{ }\PY{l+m+mi}{10}\PY{p}{,}\PY{+w}{ }\PY{l+m+mi}{9}\PY{+w}{ }\PY{o}{\PYZhy{}}\PY{o}{\PYZgt{}}\PY{+w}{ }\PY{l+s}{\PYZdq{}}\PY{l+s}{A}\PY{l+s}{\PYZdq{}}\PY{p}{;}
\PY{+w}{    }\PY{k}{case}\PY{+w}{ }\PY{l+m+mi}{8}\PY{+w}{ }\PY{o}{\PYZhy{}}\PY{o}{\PYZgt{}}\PY{+w}{ }\PY{p}{\PYZob{}}
\PY{+w}{        }\PY{n}{System}\PY{p}{.}\PY{n+na}{out}\PY{p}{.}\PY{n+na}{println}\PY{p}{(}\PY{l+s}{\PYZdq{}}\PY{l+s}{Good job!}\PY{l+s}{\PYZdq{}}\PY{p}{)}\PY{p}{;}\PY{+w}{ }\PY{c+c1}{// side effect allowed}
\PY{+w}{        }\PY{k+kd}{yield}\PY{+w}{ }\PY{l+s}{\PYZdq{}}\PY{l+s}{B}\PY{l+s}{\PYZdq{}}\PY{p}{;}
\PY{+w}{    }\PY{p}{\PYZcb{}}
\PY{+w}{    }\PY{k}{default}\PY{+w}{ }\PY{o}{\PYZhy{}}\PY{o}{\PYZgt{}}\PY{+w}{ }\PY{l+s}{\PYZdq{}}\PY{l+s}{C or below}\PY{l+s}{\PYZdq{}}\PY{p}{;}
\PY{p}{\PYZcb{}}\PY{p}{;}
\end{Verbatim}
\end{codebox}

Switch statement vs. switch expression:

{\tablefont
\begin{xltabular}{\linewidth}{@{}L{0.676}L{0.761}L{1.563}@{}}
\toprule
\rowcolor{tablehdr}
\thd{Aspect} & \thd{\code{switch} statement} & \thd{\code{switch} expression} \\
\midrule
\endhead
Produces a value & No & Yes \\
Syntax & \code{case~\allowbreak{}X:} with \code{break} & \code{case~\allowbreak{}X~\allowbreak{}-\textgreater{}} (arrow) or block with \code{yield} \\
Fall-through & Yes, unless you add \code{break} & No fall-through \\
Exhaustiveness check (enums/sealed types) & Not enforced & Compiler can enforce it — missing cases won’t compile \\
Since & Java 1 & Java 14 \\
\bottomrule
\end{xltabular}
}

Exhaustiveness with \code{enum} or \code{sealed} types is a big win for reviewers: the compiler forces you to handle every case (or add a \code{default}), catching bugs when a new enum constant is added later but a \code{switch} elsewhere wasn’t updated.

\begin{codebox}
\begin{Verbatim}[commandchars=\\\{\}]
\PY{k+kd}{enum}\PY{+w}{ }\PY{n}{TrafficLight}\PY{+w}{ }\PY{p}{\PYZob{}}\PY{+w}{ }\PY{n}{RED}\PY{p}{,}\PY{+w}{ }\PY{n}{YELLOW}\PY{p}{,}\PY{+w}{ }\PY{n}{GREEN}\PY{+w}{ }\PY{p}{\PYZcb{}}

\PY{k+kd}{static}\PY{+w}{ }\PY{n}{String}\PY{+w}{ }\PY{n+nf}{action}\PY{p}{(}\PY{n}{TrafficLight}\PY{+w}{ }\PY{n}{light}\PY{p}{)}\PY{+w}{ }\PY{p}{\PYZob{}}
\PY{+w}{    }\PY{c+c1}{// No default needed: compiler verifies all 3 enum values are covered.}
\PY{+w}{    }\PY{k}{return}\PY{+w}{ }\PY{k}{switch}\PY{+w}{ }\PY{p}{(}\PY{n}{light}\PY{p}{)}\PY{+w}{ }\PY{p}{\PYZob{}}
\PY{+w}{        }\PY{k}{case}\PY{+w}{ }\PY{n}{RED}\PY{+w}{ }\PY{o}{\PYZhy{}}\PY{o}{\PYZgt{}}\PY{+w}{ }\PY{l+s}{\PYZdq{}}\PY{l+s}{Stop}\PY{l+s}{\PYZdq{}}\PY{p}{;}
\PY{+w}{        }\PY{k}{case}\PY{+w}{ }\PY{n}{YELLOW}\PY{+w}{ }\PY{o}{\PYZhy{}}\PY{o}{\PYZgt{}}\PY{+w}{ }\PY{l+s}{\PYZdq{}}\PY{l+s}{Slow down}\PY{l+s}{\PYZdq{}}\PY{p}{;}
\PY{+w}{        }\PY{k}{case}\PY{+w}{ }\PY{n}{GREEN}\PY{+w}{ }\PY{o}{\PYZhy{}}\PY{o}{\PYZgt{}}\PY{+w}{ }\PY{l+s}{\PYZdq{}}\PY{l+s}{Go}\PY{l+s}{\PYZdq{}}\PY{p}{;}
\PY{+w}{    }\PY{p}{\PYZcb{}}\PY{p}{;}
\PY{p}{\PYZcb{}}
\end{Verbatim}
\end{codebox}

\section[Text Blocks]{Text Blocks}
\label{h:3:text-blocks}
A \textbf{text block} is a multi-line string literal. It removes the need for \code{\textbackslash{}n} and string concatenation when writing multi-line text like JSON, SQL, or HTML. Text blocks were finalized in \textbf{Java 15}.

A text block starts and ends with \code{\textquotedbl{}\textquotedbl{}\textquotedbl{}} (three double quotes), and the opening \code{\textquotedbl{}\textquotedbl{}\textquotedbl{}} must be followed by a line break.

\begin{codebox}
\begin{Verbatim}[commandchars=\\\{\}]
\PY{n}{String}\PY{+w}{ }\PY{n}{json}\PY{+w}{ }\PY{o}{=}\PY{+w}{ }\PY{l+s}{\PYZdq{}\PYZdq{}\PYZdq{}}
\PY{l+s}{        \PYZob{}}
\PY{l+s}{          }\PY{l+s}{\PYZdq{}}\PY{l+s}{name}\PY{l+s}{\PYZdq{}}\PY{l+s}{: }\PY{l+s}{\PYZdq{}}\PY{l+s}{Ada}\PY{l+s}{\PYZdq{}}\PY{l+s}{,}
\PY{l+s}{          }\PY{l+s}{\PYZdq{}}\PY{l+s}{role}\PY{l+s}{\PYZdq{}}\PY{l+s}{: }\PY{l+s}{\PYZdq{}}\PY{l+s}{Engineer}\PY{l+s}{\PYZdq{}}
\PY{l+s}{        \PYZcb{}}
\PY{l+s}{        }\PY{l+s}{\PYZdq{}\PYZdq{}\PYZdq{}}\PY{p}{;}

\PY{n}{System}\PY{p}{.}\PY{n+na}{out}\PY{p}{.}\PY{n+na}{println}\PY{p}{(}\PY{n}{json}\PY{p}{)}\PY{p}{;}
\PY{c+c1}{// Output:}
\PY{c+c1}{// \PYZob{}}
\PY{c+c1}{//   \PYZdq{}name\PYZdq{}: \PYZdq{}Ada\PYZdq{},}
\PY{c+c1}{//   \PYZdq{}role\PYZdq{}: \PYZdq{}Engineer\PYZdq{}}
\PY{c+c1}{// \PYZcb{}}
\end{Verbatim}
\end{codebox}

Before text blocks, the same JSON needed manual escaping:

\begin{codebox}
\begin{Verbatim}[commandchars=\\\{\}]
\PY{n}{String}\PY{+w}{ }\PY{n}{jsonOld}\PY{+w}{ }\PY{o}{=}\PY{+w}{ }\PY{l+s}{\PYZdq{}}\PY{l+s}{\PYZob{}}\PY{l+s}{\PYZbs{}}\PY{l+s}{n}\PY{l+s}{\PYZdq{}}\PY{+w}{ }\PY{o}{+}
\PY{+w}{                 }\PY{l+s}{\PYZdq{}}\PY{l+s}{  }\PY{l+s}{\PYZbs{}\PYZdq{}}\PY{l+s}{name}\PY{l+s}{\PYZbs{}\PYZdq{}}\PY{l+s}{: }\PY{l+s}{\PYZbs{}\PYZdq{}}\PY{l+s}{Ada}\PY{l+s}{\PYZbs{}\PYZdq{}}\PY{l+s}{,}\PY{l+s}{\PYZbs{}}\PY{l+s}{n}\PY{l+s}{\PYZdq{}}\PY{+w}{ }\PY{o}{+}
\PY{+w}{                 }\PY{l+s}{\PYZdq{}}\PY{l+s}{  }\PY{l+s}{\PYZbs{}\PYZdq{}}\PY{l+s}{role}\PY{l+s}{\PYZbs{}\PYZdq{}}\PY{l+s}{: }\PY{l+s}{\PYZbs{}\PYZdq{}}\PY{l+s}{Engineer}\PY{l+s}{\PYZbs{}\PYZdq{}}\PY{l+s}{\PYZbs{}}\PY{l+s}{n}\PY{l+s}{\PYZdq{}}\PY{+w}{ }\PY{o}{+}
\PY{+w}{                 }\PY{l+s}{\PYZdq{}}\PY{l+s}{\PYZcb{}}\PY{l+s}{\PYZbs{}}\PY{l+s}{n}\PY{l+s}{\PYZdq{}}\PY{p}{;}
\end{Verbatim}
\end{codebox}

Indentation rules: the compiler looks at the \textbf{least-indented line} (including the closing \code{\textquotedbl{}\textquotedbl{}\textquotedbl{}}) and strips that much leading whitespace from every line. This means the position of the closing \code{\textquotedbl{}\textquotedbl{}\textquotedbl{}} controls the indentation of the output.

\begin{codebox}
\begin{Verbatim}[commandchars=\\\{\}]
\PY{n}{String}\PY{+w}{ }\PY{n}{sql}\PY{+w}{ }\PY{o}{=}\PY{+w}{ }\PY{l+s}{\PYZdq{}\PYZdq{}\PYZdq{}}
\PY{l+s}{    SELECT id, name}
\PY{l+s}{    FROM employees}
\PY{l+s}{    WHERE department = \PYZsq{}Engineering\PYZsq{}}
\PY{l+s}{    }\PY{l+s}{\PYZdq{}\PYZdq{}\PYZdq{}}\PY{p}{;}
\end{Verbatim}
\end{codebox}

You can still use escape sequences, and a few special ones help control whitespace:

\begin{codebox}
\begin{Verbatim}[commandchars=\\\{\}]
\PY{n}{String}\PY{+w}{ }\PY{n}{noTrailingNewline}\PY{+w}{ }\PY{o}{=}\PY{+w}{ }\PY{l+s}{\PYZdq{}\PYZdq{}\PYZdq{}}
\PY{l+s}{        Line one}
\PY{l+s}{        Line two}\PY{l+s}{\PYZbs{}}
\PY{l+s}{        }\PY{l+s}{\PYZdq{}\PYZdq{}\PYZdq{}}\PY{p}{;}\PY{+w}{ }\PY{c+c1}{// \PYZdq{}\PYZbs{}\PYZdq{} at end of line joins it with the next line (no newline inserted)}

\PY{n}{String}\PY{+w}{ }\PY{n}{withTrailingSpaces}\PY{+w}{ }\PY{o}{=}\PY{+w}{ }\PY{l+s}{\PYZdq{}\PYZdq{}\PYZdq{}}
\PY{l+s}{        Trailing spaces matter here.   }\PY{l+s}{\PYZbs{}}\PY{l+s}{s}
\PY{l+s}{        }\PY{l+s}{\PYZdq{}\PYZdq{}\PYZdq{}}\PY{p}{;}\PY{+w}{ }\PY{c+c1}{// \PYZbs{}s forces a space that would otherwise be trimmed}
\end{Verbatim}
\end{codebox}

You can also combine text blocks with \code{String.\allowbreak{}format} or \code{formatted}:

\begin{codebox}
\begin{Verbatim}[commandchars=\\\{\}]
\PY{n}{String}\PY{+w}{ }\PY{n}{template}\PY{+w}{ }\PY{o}{=}\PY{+w}{ }\PY{l+s}{\PYZdq{}\PYZdq{}\PYZdq{}}
\PY{l+s}{        Hello, \PYZpc{}s!}
\PY{l+s}{        You have \PYZpc{}d new messages.}
\PY{l+s}{        }\PY{l+s}{\PYZdq{}\PYZdq{}\PYZdq{}}\PY{p}{;}
\PY{n}{String}\PY{+w}{ }\PY{n}{message}\PY{+w}{ }\PY{o}{=}\PY{+w}{ }\PY{n}{template}\PY{p}{.}\PY{n+na}{formatted}\PY{p}{(}\PY{l+s}{\PYZdq{}}\PY{l+s}{Ada}\PY{l+s}{\PYZdq{}}\PY{p}{,}\PY{+w}{ }\PY{l+m+mi}{5}\PY{p}{)}\PY{p}{;}
\end{Verbatim}
\end{codebox}

\section[String Templates (Preview)]{String Templates (Preview)}
\label{h:3:string-templates-preview}
\textbf{String templates} were a proposed feature to embed expressions directly inside string literals, similar to JavaScript template literals or Kotlin string interpolation. They used the syntax \code{STR.\allowbreak{}\textquotedbl{}Hello~\allowbreak{}\textbackslash{}\{\allowbreak{}name\}\allowbreak{}\textquotedbl{}}.

\begin{codebox}
\begin{Verbatim}[commandchars=\\\{\}]
\PY{c+c1}{// Example of the PREVIEWED (now withdrawn) syntax — do not use in real code:}
\PY{c+c1}{// String name = \PYZdq{}Ada\PYZdq{};}
\PY{c+c1}{// int count = 5;}
\PY{c+c1}{// String message = STR.\PYZdq{}Hello \PYZbs{}\PYZob{}name\PYZcb{}, you have \PYZbs{}\PYZob{}count\PYZcb{} messages.\PYZdq{};}
\end{Verbatim}
\end{codebox}

\textbf{Important — be accurate about status:} String templates were previewed in \textbf{Java 21} (JEP 430) and again, with refinements, in \textbf{Java 22} (JEP 459). However, based on developer feedback, the OpenJDK team decided the design needed more work, and the feature was \textbf{withdrawn} — it did not ship as a preview in Java 23, and there is no finalized version as of Java 25. If you see \code{STR.\allowbreak{}\textquotedbl{}...\allowbreak{}\textquotedbl{}} syntax in an article or an older tutorial, know that it will not compile on any released JDK.

Until a template mechanism (re-)ships, the standard ways to build formatted strings remain:

\begin{codebox}
\begin{Verbatim}[commandchars=\\\{\}]
\PY{n}{String}\PY{+w}{ }\PY{n}{name}\PY{+w}{ }\PY{o}{=}\PY{+w}{ }\PY{l+s}{\PYZdq{}}\PY{l+s}{Ada}\PY{l+s}{\PYZdq{}}\PY{p}{;}
\PY{k+kt}{int}\PY{+w}{ }\PY{n}{count}\PY{+w}{ }\PY{o}{=}\PY{+w}{ }\PY{l+m+mi}{5}\PY{p}{;}

\PY{c+c1}{// String concatenation}
\PY{n}{String}\PY{+w}{ }\PY{n}{a}\PY{+w}{ }\PY{o}{=}\PY{+w}{ }\PY{l+s}{\PYZdq{}}\PY{l+s}{Hello, }\PY{l+s}{\PYZdq{}}\PY{+w}{ }\PY{o}{+}\PY{+w}{ }\PY{n}{name}\PY{+w}{ }\PY{o}{+}\PY{+w}{ }\PY{l+s}{\PYZdq{}}\PY{l+s}{, you have }\PY{l+s}{\PYZdq{}}\PY{+w}{ }\PY{o}{+}\PY{+w}{ }\PY{n}{count}\PY{+w}{ }\PY{o}{+}\PY{+w}{ }\PY{l+s}{\PYZdq{}}\PY{l+s}{ messages.}\PY{l+s}{\PYZdq{}}\PY{p}{;}

\PY{c+c1}{// String.format / formatted (classic, since Java 5 / Java 15)}
\PY{n}{String}\PY{+w}{ }\PY{n}{b}\PY{+w}{ }\PY{o}{=}\PY{+w}{ }\PY{l+s}{\PYZdq{}}\PY{l+s}{Hello, \PYZpc{}s, you have \PYZpc{}d messages.}\PY{l+s}{\PYZdq{}}\PY{p}{.}\PY{n+na}{formatted}\PY{p}{(}\PY{n}{name}\PY{p}{,}\PY{+w}{ }\PY{n}{count}\PY{p}{)}\PY{p}{;}

\PY{c+c1}{// Text block + formatted (Java 15+)}
\PY{n}{String}\PY{+w}{ }\PY{n}{c}\PY{+w}{ }\PY{o}{=}\PY{+w}{ }\PY{l+s}{\PYZdq{}\PYZdq{}\PYZdq{}}
\PY{l+s}{        Hello, \PYZpc{}s, you have \PYZpc{}d messages.}
\PY{l+s}{        }\PY{l+s}{\PYZdq{}\PYZdq{}\PYZdq{}}\PY{p}{.}\PY{n+na}{formatted}\PY{p}{(}\PY{n}{name}\PY{p}{,}\PY{+w}{ }\PY{n}{count}\PY{p}{)}\PY{p}{;}
\end{Verbatim}
\end{codebox}

Code-review takeaway: if a candidate claims string templates are “a stable Java feature,” that is a \textbf{red flag} — the honest answer is “previewed twice, then withdrawn; not currently available.”

\section[Unnamed Variables and Patterns]{Unnamed Variables and Patterns}
\label{h:3:unnamed-variables-and-patterns}
Sometimes you must declare a variable (for a lambda parameter, a \code{catch} block, or a pattern component) that you never actually use. The \textbf{unnamed variable}, written as a single underscore \code{\_}, lets you say “I don’t need this” explicitly. Finalized in \textbf{Java 22} (JEP 456).

\begin{codebox}
\begin{Verbatim}[commandchars=\\\{\}]
\PY{n}{map}\PY{p}{.}\PY{n+na}{forEach}\PY{p}{(}\PY{p}{(}\PY{n}{key}\PY{p}{,}\PY{+w}{ }\PY{n}{\PYZus{}}\PY{p}{)}\PY{+w}{ }\PY{o}{\PYZhy{}}\PY{o}{\PYZgt{}}\PY{+w}{ }\PY{n}{System}\PY{p}{.}\PY{n+na}{out}\PY{p}{.}\PY{n+na}{println}\PY{p}{(}\PY{n}{key}\PY{p}{)}\PY{p}{)}\PY{p}{;}\PY{+w}{ }\PY{c+c1}{// value is unused}
\end{Verbatim}
\end{codebox}

\begin{codebox}
\begin{Verbatim}[commandchars=\\\{\}]
\PY{k}{try}\PY{+w}{ }\PY{p}{\PYZob{}}
\PY{+w}{    }\PY{n}{riskyOperation}\PY{p}{(}\PY{p}{)}\PY{p}{;}
\PY{p}{\PYZcb{}}\PY{+w}{ }\PY{k}{catch}\PY{+w}{ }\PY{p}{(}\PY{n}{IOException}\PY{+w}{ }\PY{n}{\PYZus{}}\PY{p}{)}\PY{+w}{ }\PY{p}{\PYZob{}}
\PY{+w}{    }\PY{c+c1}{// We don\PYZsq{}t care about the exception details, just that it failed.}
\PY{+w}{    }\PY{n}{System}\PY{p}{.}\PY{n+na}{out}\PY{p}{.}\PY{n+na}{println}\PY{p}{(}\PY{l+s}{\PYZdq{}}\PY{l+s}{Operation failed}\PY{l+s}{\PYZdq{}}\PY{p}{)}\PY{p}{;}
\PY{p}{\PYZcb{}}
\end{Verbatim}
\end{codebox}

Unnamed patterns work well with record deconstruction when you only care about some components:

\begin{codebox}
\begin{Verbatim}[commandchars=\\\{\}]
\PY{k+kd}{record} \PY{n+nc}{Point3D}\PY{p}{(}\PY{k+kt}{int}\PY{+w}{ }\PY{n}{x}\PY{p}{,}\PY{+w}{ }\PY{k+kt}{int}\PY{+w}{ }\PY{n}{y}\PY{p}{,}\PY{+w}{ }\PY{k+kt}{int}\PY{+w}{ }\PY{n}{z}\PY{p}{)}\PY{+w}{ }\PY{p}{\PYZob{}}\PY{+w}{ }\PY{p}{\PYZcb{}}

\PY{k+kd}{static}\PY{+w}{ }\PY{n}{String}\PY{+w}{ }\PY{n+nf}{flatten}\PY{p}{(}\PY{n}{Point3D}\PY{+w}{ }\PY{n}{p}\PY{p}{)}\PY{+w}{ }\PY{p}{\PYZob{}}
\PY{+w}{    }\PY{c+c1}{// We only care about x and y; z is discarded.}
\PY{+w}{    }\PY{k}{if}\PY{+w}{ }\PY{p}{(}\PY{n}{p}\PY{+w}{ }\PY{k}{instanceof}\PY{+w}{ }\PY{n+nf}{Point3D}\PY{p}{(}\PY{k+kd}{var}\PY{+w}{ }\PY{n}{x}\PY{p}{,}\PY{+w}{ }\PY{k+kd}{var}\PY{+w}{ }\PY{n}{y}\PY{p}{,}\PY{+w}{ }\PY{k+kd}{var}\PY{+w}{ }\PY{n}{\PYZus{}}\PY{p}{)}\PY{p}{)}\PY{+w}{ }\PY{p}{\PYZob{}}
\PY{+w}{        }\PY{k}{return}\PY{+w}{ }\PY{l+s}{\PYZdq{}}\PY{l+s}{(}\PY{l+s}{\PYZdq{}}\PY{+w}{ }\PY{o}{+}\PY{+w}{ }\PY{n}{x}\PY{+w}{ }\PY{o}{+}\PY{+w}{ }\PY{l+s}{\PYZdq{}}\PY{l+s}{, }\PY{l+s}{\PYZdq{}}\PY{+w}{ }\PY{o}{+}\PY{+w}{ }\PY{n}{y}\PY{+w}{ }\PY{o}{+}\PY{+w}{ }\PY{l+s}{\PYZdq{}}\PY{l+s}{)}\PY{l+s}{\PYZdq{}}\PY{p}{;}
\PY{+w}{    }\PY{p}{\PYZcb{}}
\PY{+w}{    }\PY{k}{return}\PY{+w}{ }\PY{l+s}{\PYZdq{}}\PY{l+s}{invalid}\PY{l+s}{\PYZdq{}}\PY{p}{;}
\PY{p}{\PYZcb{}}
\end{Verbatim}
\end{codebox}

You can also use \code{\_} for a loop variable you never read:

\begin{codebox}
\begin{Verbatim}[commandchars=\\\{\}]
\PY{k+kt}{int}\PY{+w}{ }\PY{n}{count}\PY{+w}{ }\PY{o}{=}\PY{+w}{ }\PY{l+m+mi}{0}\PY{p}{;}
\PY{k}{for}\PY{+w}{ }\PY{p}{(}\PY{k+kd}{var}\PY{+w}{ }\PY{n}{\PYZus{}}\PY{+w}{ }\PY{p}{:}\PY{+w}{ }\PY{n}{someList}\PY{p}{)}\PY{+w}{ }\PY{p}{\PYZob{}}
\PY{+w}{    }\PY{n}{count}\PY{o}{+}\PY{o}{+}\PY{p}{;}\PY{+w}{ }\PY{c+c1}{// we only care how many elements there are}
\PY{p}{\PYZcb{}}
\end{Verbatim}
\end{codebox}

Why this matters in review: unlike an unused variable named \code{ignored} or \code{e}, the unnamed variable \code{\_} is a language-level signal — some static analysis tools and the compiler itself understand “this is intentionally unused,” reducing false-positive “unused variable” warnings and making intent explicit to the next reader.

\section[Value-Based Classes]{Value-Based Classes}
\label{h:3:value-based-classes}
A \textbf{value-based class} is a class whose instances are treated as pure values, not as objects with identity. \code{Integer}, \code{LocalDate}, \code{Optional}, and (since Java 16) \code{record} types are all examples. The JDK documentation defines conventions for such classes; they’re a design pattern/contract, not a language keyword.

Rules for a value-based class (per the JDK’s \code{java.\allowbreak{}lang.\allowbreak{}doc-\allowbreak{}files/\allowbreak{}Value\allowbreak{}Based.\allowbreak{}html} conventions):

\begin{enumerate}
\item 
Fields are \code{final}, and the object is immutable.

\item 
The class is \code{final} (or effectively behaves that way) — no subclassing that adds identity-sensitive behavior.

\item 
Instances are considered interchangeable if \code{equals()} says they’re equal — you should never rely on \code{==} or synchronize on them.

\item 
The class does not expose an identity-sensitive API (no exposed lock, no \code{==}-based contracts).

\item 
Instances are typically created via static factory methods, not always via visible constructors (e.g. \code{Optional.\allowbreak{}of(...)}, \code{LocalDate.\allowbreak{}of(...)}).

\end{enumerate}

\begin{codebox}
\begin{Verbatim}[commandchars=\\\{\}]
\PY{c+c1}{// Integer is value\PYZhy{}based: NEVER rely on identity (\PYZdq{}==\PYZdq{}) for comparison.}
\PY{n}{Integer}\PY{+w}{ }\PY{n}{a}\PY{+w}{ }\PY{o}{=}\PY{+w}{ }\PY{l+m+mi}{100}\PY{p}{;}
\PY{n}{Integer}\PY{+w}{ }\PY{n}{b}\PY{+w}{ }\PY{o}{=}\PY{+w}{ }\PY{l+m+mi}{100}\PY{p}{;}
\PY{n}{System}\PY{p}{.}\PY{n+na}{out}\PY{p}{.}\PY{n+na}{println}\PY{p}{(}\PY{n}{a}\PY{+w}{ }\PY{o}{=}\PY{o}{=}\PY{+w}{ }\PY{n}{b}\PY{p}{)}\PY{p}{;}\PY{+w}{      }\PY{c+c1}{// true (small int cache, but DO NOT rely on this)}

\PY{n}{Integer}\PY{+w}{ }\PY{n}{c}\PY{+w}{ }\PY{o}{=}\PY{+w}{ }\PY{l+m+mi}{200}\PY{p}{;}
\PY{n}{Integer}\PY{+w}{ }\PY{n}{d}\PY{+w}{ }\PY{o}{=}\PY{+w}{ }\PY{l+m+mi}{200}\PY{p}{;}
\PY{n}{System}\PY{p}{.}\PY{n+na}{out}\PY{p}{.}\PY{n+na}{println}\PY{p}{(}\PY{n}{c}\PY{+w}{ }\PY{o}{=}\PY{o}{=}\PY{+w}{ }\PY{n}{d}\PY{p}{)}\PY{p}{;}\PY{+w}{      }\PY{c+c1}{// false! Outside the cached range (\PYZhy{}128..127)}
\PY{n}{System}\PY{p}{.}\PY{n+na}{out}\PY{p}{.}\PY{n+na}{println}\PY{p}{(}\PY{n}{c}\PY{p}{.}\PY{n+na}{equals}\PY{p}{(}\PY{n}{d}\PY{p}{)}\PY{p}{)}\PY{p}{;}\PY{+w}{ }\PY{c+c1}{// true — this is the correct way to compare}
\end{Verbatim}
\end{codebox}

\begin{codebox}
\begin{Verbatim}[commandchars=\\\{\}]
\PY{c+c1}{// Records are value\PYZhy{}based by nature.}
\PY{k+kd}{public}\PY{+w}{ }\PY{k+kd}{record} \PY{n+nc}{Money}\PY{p}{(}\PY{n}{String}\PY{+w}{ }\PY{n}{currency}\PY{p}{,}\PY{+w}{ }\PY{k+kt}{long}\PY{+w}{ }\PY{n}{cents}\PY{p}{)}\PY{+w}{ }\PY{p}{\PYZob{}}\PY{+w}{ }\PY{p}{\PYZcb{}}

\PY{n}{Money}\PY{+w}{ }\PY{n}{price1}\PY{+w}{ }\PY{o}{=}\PY{+w}{ }\PY{k}{new}\PY{+w}{ }\PY{n}{Money}\PY{p}{(}\PY{l+s}{\PYZdq{}}\PY{l+s}{EUR}\PY{l+s}{\PYZdq{}}\PY{p}{,}\PY{+w}{ }\PY{l+m+mi}{1999}\PY{p}{)}\PY{p}{;}
\PY{n}{Money}\PY{+w}{ }\PY{n}{price2}\PY{+w}{ }\PY{o}{=}\PY{+w}{ }\PY{k}{new}\PY{+w}{ }\PY{n}{Money}\PY{p}{(}\PY{l+s}{\PYZdq{}}\PY{l+s}{EUR}\PY{l+s}{\PYZdq{}}\PY{p}{,}\PY{+w}{ }\PY{l+m+mi}{1999}\PY{p}{)}\PY{p}{;}

\PY{n}{System}\PY{p}{.}\PY{n+na}{out}\PY{p}{.}\PY{n+na}{println}\PY{p}{(}\PY{n}{price1}\PY{+w}{ }\PY{o}{=}\PY{o}{=}\PY{+w}{ }\PY{n}{price2}\PY{p}{)}\PY{p}{;}\PY{+w}{      }\PY{c+c1}{// false — different instances}
\PY{n}{System}\PY{p}{.}\PY{n+na}{out}\PY{p}{.}\PY{n+na}{println}\PY{p}{(}\PY{n}{price1}\PY{p}{.}\PY{n+na}{equals}\PY{p}{(}\PY{n}{price2}\PY{p}{)}\PY{p}{)}\PY{p}{;}\PY{+w}{ }\PY{c+c1}{// true — same value}
\end{Verbatim}
\end{codebox}

Because value-based instances may be interchangeable, the JDK explicitly warns against \textbf{synchronizing on them}:

\begin{codebox}
\begin{Verbatim}[commandchars=\\\{\}]
\PY{n}{Integer}\PY{+w}{ }\PY{n}{lock}\PY{+w}{ }\PY{o}{=}\PY{+w}{ }\PY{l+m+mi}{42}\PY{p}{;}
\PY{k+kd}{synchronized}\PY{+w}{ }\PY{p}{(}\PY{n}{lock}\PY{p}{)}\PY{+w}{ }\PY{p}{\PYZob{}}\PY{+w}{ }\PY{c+c1}{// BAD: value\PYZhy{}based classes may be cached/shared/pooled}
\PY{+w}{    }\PY{c+c1}{// ...}
\PY{p}{\PYZcb{}}
\end{Verbatim}
\end{codebox}

Use a dedicated \code{Object} or \code{ReentrantLock} for locking instead:

\begin{codebox}
\begin{Verbatim}[commandchars=\\\{\}]
\PY{k+kd}{private}\PY{+w}{ }\PY{k+kd}{final}\PY{+w}{ }\PY{n}{Object}\PY{+w}{ }\PY{n}{lock}\PY{+w}{ }\PY{o}{=}\PY{+w}{ }\PY{k}{new}\PY{+w}{ }\PY{n}{Object}\PY{p}{(}\PY{p}{)}\PY{p}{;}

\PY{k+kd}{synchronized}\PY{+w}{ }\PY{p}{(}\PY{n}{lock}\PY{p}{)}\PY{+w}{ }\PY{p}{\PYZob{}}
\PY{+w}{    }\PY{c+c1}{// safe: this object has real identity and is never shared}
\PY{p}{\PYZcb{}}
\end{Verbatim}
\end{codebox}

\section[Preview and Incubator Features]{Preview and Incubator Features}
\label{h:3:preview-and-incubator-features}
The JDK ships experimental features in two ways, so developers can try them before they become permanent:

{\tablefontsmall
\begin{xltabular}{\linewidth}{@{}L{0.424}L{1.422}L{1.154}@{}}
\toprule
\rowcolor{tablehdr}
\thd{Mechanism} & \thd{What it means} & \thd{How to enable} \\
\midrule
\endhead
\textbf{Preview feature} & A complete, fully-specified language or API feature, but not yet guaranteed stable. Might change or be withdrawn based on feedback. & Compile and run with \code{--\allowbreak{}enable-\allowbreak{}preview} (plus \code{--\allowbreak{}release~\allowbreak{}\textless{}\allowbreak{}N\textgreater{}} when compiling) \\
\textbf{Incubator module} & An API delivered in a \code{jdk.\allowbreak{}incubator.*} module, for APIs that need real-world testing before joining the standard \code{java.*} API. & Add \code{--\allowbreak{}add-\allowbreak{}modules~\allowbreak{}jdk.\allowbreak{}incubator.\textless{}\allowbreak{}name\textgreater{}} \\
\bottomrule
\end{xltabular}
}

Preview features must be explicitly enabled on \textbf{both} compilation and execution, and the produced \code{.\allowbreak{}class} files are marked so they cannot run on a JDK that isn’t the exact same preview-enabled version.

\begin{codebox}
\begin{Verbatim}[commandchars=\\\{\}]
\PY{c+c1}{\PYZsh{} Compiling code that uses a preview feature (e.g. primitive patterns, JDK 23):}
javac\PY{+w}{ }\PYZhy{}\PYZhy{}release\PY{+w}{ }\PY{l+m}{23}\PY{+w}{ }\PYZhy{}\PYZhy{}enable\PYZhy{}preview\PY{+w}{ }Main.java

\PY{c+c1}{\PYZsh{} Running it:}
java\PY{+w}{ }\PYZhy{}\PYZhy{}enable\PYZhy{}preview\PY{+w}{ }Main
\end{Verbatim}
\end{codebox}

Examples of features that went through the preview process (with outcomes):

{\tablefont
\begin{xltabular}{\linewidth}{@{}L{1.191}L{0.618}L{1.191}@{}}
\toprule
\rowcolor{tablehdr}
\thd{Feature} & \thd{Preview JDKs} & \thd{Outcome} \\
\midrule
\endhead
Records & 14, 15 & Finalized in 16 \\
Sealed classes & 15, 16 & Finalized in 17 \\
Pattern matching for \code{switch} & 17, 18, 19, 20 & Finalized in 21 \\
Record patterns & 19, 20 & Finalized in 21 \\
Virtual threads & 19, 20 & Finalized in 21 \\
Structured concurrency & 21, 22, 23, 24 & Still preview as of Java 25 \\
String templates & 21, 22 & \textbf{Withdrawn} — not shipped \\
Unnamed variables \& patterns & 21 & Finalized in 22 \\
Primitive types in patterns & 23 (JEP 455), 24 (JEP 488) & Still preview as of Java 25 (re-previewed as JEP 507) \\
\bottomrule
\end{xltabular}
}

Incubator module examples:

\begin{codebox}
\begin{Verbatim}[commandchars=\\\{\}]
\PY{c+c1}{// Vector API (jdk.incubator.vector) — SIMD\PYZhy{}style computation, still incubating}
\PY{c+c1}{// across many releases (first incubated in Java 16, still incubating in Java 25).}
\PY{k+kn}{import}\PY{+w}{ }\PY{n+nn}{jdk.incubator.vector.*}\PY{p}{;}

\PY{c+c1}{// Foreign Function \PYZam{} Memory API started as incubator/preview (Java 17\PYZhy{}21)}
\PY{c+c1}{// and was finalized as a standard feature in Java 22.}
\end{Verbatim}
\end{codebox}

Code-review guidance: preview and incubator features are \textbf{not recommended for production code} unless the team has explicitly decided to accept the risk (API changes, or the feature being withdrawn entirely, as happened with string templates). Flag any use of \code{--\allowbreak{}enable-\allowbreak{}preview} in a build file as something that needs a deliberate, documented decision.

\section[Deprecated and Removed Features Across Java Versions]{Deprecated and Removed Features Across Java Versions}
\label{h:3:deprecated-and-removed-features-across-java-versions}
Java evolves by deprecating old APIs before removing them, usually giving developers multiple releases of warning. \textbf{Deprecated} means “still works, but don’t use it — it may be removed.” \textbf{Removed} means “it’s actually gone; the code won’t compile or run.”

\begin{codebox}
\begin{Verbatim}[commandchars=\\\{\}]
\PY{n+nd}{@Deprecated}\PY{p}{(}\PY{n}{since}\PY{+w}{ }\PY{o}{=}\PY{+w}{ }\PY{l+s}{\PYZdq{}}\PY{l+s}{9}\PY{l+s}{\PYZdq{}}\PY{p}{,}\PY{+w}{ }\PY{n}{forRemoval}\PY{+w}{ }\PY{o}{=}\PY{+w}{ }\PY{k+kc}{true}\PY{p}{)}
\PY{k+kd}{public}\PY{+w}{ }\PY{k+kt}{void}\PY{+w}{ }\PY{n+nf}{oldMethod}\PY{p}{(}\PY{p}{)}\PY{+w}{ }\PY{p}{\PYZob{}}
\PY{+w}{    }\PY{c+c1}{// ...}
\PY{p}{\PYZcb{}}
\end{Verbatim}
\end{codebox}

\begin{codebox}
\begin{Verbatim}[commandchars=\\\{\}]
\PY{c+c1}{// Calling a deprecated method produces a compiler warning:}
\PY{n}{obj}\PY{p}{.}\PY{n+na}{oldMethod}\PY{p}{(}\PY{p}{)}\PY{p}{;}
\PY{c+c1}{// warning: [deprecation] oldMethod() in Foo has been deprecated and marked for removal}
\end{Verbatim}
\end{codebox}

Notable deprecations and removals:

{\tablefontsmall
\begin{xltabular}{\linewidth}{@{}L{1.053}L{0.443}L{0.950}L{1.554}@{}}
\toprule
\rowcolor{tablehdr}
\thd{Feature} & \thd{Status} & \thd{Since / Removed in} & \thd{Why} \\
\midrule
\endhead
\code{Object.\allowbreak{}finalize()} & Deprecated for removal & Deprecated in Java 9, deprecated-for-removal in Java 18 & Unpredictable timing, GC overhead, security issues. Use \code{try-\allowbreak{}with-\allowbreak{}resources} or \code{Cleaner} instead. \\
Security Manager (\code{SecurityManager}, \code{System.\allowbreak{}set\allowbreak{}Security\allowbreak{}Manager}) & Removed & Deprecated for removal in Java 17, removed in Java 24 & Rarely used correctly in practice, and complicated the JVM. No replacement — sandboxing should happen outside the JVM (containers, OS-level controls). \\
Applets (\code{java.\allowbreak{}applet}) & Removed & Deprecated in Java 9, removed in Java 17 & Browsers dropped plugin support (NPAPI); applets became unusable. \\
\code{Thread.\allowbreak{}stop()}, \code{Thread.\allowbreak{}suspend()}, \code{Thread.\allowbreak{}resume()} & Deprecated for removal & Deprecated since Java 1.2, still present but strongly discouraged & Inherently unsafe — can leave shared state in a corrupted, half-updated condition. Use cooperative cancellation (\code{volatile} flags, \code{interrupt()}) instead. \\
\code{new~\allowbreak{}Integer(...)}, \code{new~\allowbreak{}Long(...)}, \code{new~\allowbreak{}Boolean(...)}, etc. (boxed-type constructors) & Deprecated for removal & Deprecated in Java 9 & Wastes memory by creating unnecessary objects; \code{valueOf(...)} (or autoboxing) reuses cached instances for common values. \\
Nashorn JavaScript engine (\code{javax.\allowbreak{}script}, \code{jjs}) & Removed & Deprecated in Java 11, removed in Java 15 & Fell behind modern JavaScript (ES6+); GraalVM’s JS engine is the modern replacement. \\
\code{Runtime.\allowbreak{}exec(\allowbreak{}String)} (single-string overload usage patterns) & Still present, but risky usage patterns discouraged & Ongoing guidance, not removed & Splitting a single command string on spaces is fragile and can be an injection risk; prefer the \code{String[]} / \code{ProcessBuilder} form with explicit arguments. \\
Java Web Start / \code{javaws} & Removed & Removed in Java 11 (was part of Oracle JDK) & Superseded by other deployment approaches; browser plugin model became obsolete. \\
CORBA and Java EE modules (\code{java.\allowbreak{}corba}, \code{java.\allowbreak{}xml.\allowbreak{}ws}, \code{java.\allowbreak{}xml.\allowbreak{}bind}, etc.) & Removed & Deprecated in Java 9, removed in Java 11 & Moved out of the JDK; available as separate Maven/Gradle dependencies if still needed. \\
\code{finalization} mechanism as a whole (the \code{Finalizable} protocol) & Disabled by default, planned for removal & Disabled by default via \code{--\allowbreak{}finalization=\allowbreak{}disabled} option available since Java 18, deprecated for removal since Java 18 & Same reasons as \code{finalize()} above — the whole mechanism is being phased out. \\
Biased Locking (JVM internal optimization) & Removed & Deprecated in Java 15, removed in Java 18 & Modern GC and JIT improvements made the added complexity not worth it. \\
\code{com.\allowbreak{}sun.\allowbreak{}security.\allowbreak{}auth} classes / old JAAS internals & Deprecated & Various & Prefer supported \code{javax.\allowbreak{}security.\allowbreak{}auth} public APIs. \\
\bottomrule
\end{xltabular}
}

\begin{codebox}
\begin{Verbatim}[commandchars=\\\{\}]
\PY{c+c1}{// BAD (deprecated boxed constructor):}
\PY{n}{Integer}\PY{+w}{ }\PY{n}{wasteful}\PY{+w}{ }\PY{o}{=}\PY{+w}{ }\PY{k}{new}\PY{+w}{ }\PY{n}{Integer}\PY{p}{(}\PY{l+m+mi}{42}\PY{p}{)}\PY{p}{;}\PY{+w}{ }\PY{c+c1}{// creates a brand\PYZhy{}new object every time}

\PY{c+c1}{// GOOD (uses the internal cache for common values):}
\PY{n}{Integer}\PY{+w}{ }\PY{n}{efficient}\PY{+w}{ }\PY{o}{=}\PY{+w}{ }\PY{n}{Integer}\PY{p}{.}\PY{n+na}{valueOf}\PY{p}{(}\PY{l+m+mi}{42}\PY{p}{)}\PY{p}{;}
\PY{n}{Integer}\PY{+w}{ }\PY{n}{autoboxed}\PY{+w}{ }\PY{o}{=}\PY{+w}{ }\PY{l+m+mi}{42}\PY{p}{;}\PY{+w}{ }\PY{c+c1}{// compiler calls Integer.valueOf() automatically}
\end{Verbatim}
\end{codebox}

\begin{codebox}
\begin{Verbatim}[commandchars=\\\{\}]
\PY{c+c1}{// BAD (Thread.stop is unsafe — can corrupt shared state mid\PYZhy{}update):}
\PY{c+c1}{// worker.stop();}

\PY{c+c1}{// GOOD (cooperative cancellation):}
\PY{k+kd}{class} \PY{n+nc}{Worker}\PY{+w}{ }\PY{k+kd}{implements}\PY{+w}{ }\PY{n}{Runnable}\PY{+w}{ }\PY{p}{\PYZob{}}
\PY{+w}{    }\PY{k+kd}{private}\PY{+w}{ }\PY{k+kd}{volatile}\PY{+w}{ }\PY{k+kt}{boolean}\PY{+w}{ }\PY{n}{running}\PY{+w}{ }\PY{o}{=}\PY{+w}{ }\PY{k+kc}{true}\PY{p}{;}

\PY{+w}{    }\PY{k+kd}{public}\PY{+w}{ }\PY{k+kt}{void}\PY{+w}{ }\PY{n+nf}{run}\PY{p}{(}\PY{p}{)}\PY{+w}{ }\PY{p}{\PYZob{}}
\PY{+w}{        }\PY{k}{while}\PY{+w}{ }\PY{p}{(}\PY{n}{running}\PY{p}{)}\PY{+w}{ }\PY{p}{\PYZob{}}
\PY{+w}{            }\PY{c+c1}{// do work}
\PY{+w}{        }\PY{p}{\PYZcb{}}
\PY{+w}{    }\PY{p}{\PYZcb{}}

\PY{+w}{    }\PY{k+kd}{public}\PY{+w}{ }\PY{k+kt}{void}\PY{+w}{ }\PY{n+nf}{shutdown}\PY{p}{(}\PY{p}{)}\PY{+w}{ }\PY{p}{\PYZob{}}
\PY{+w}{        }\PY{n}{running}\PY{+w}{ }\PY{o}{=}\PY{+w}{ }\PY{k+kc}{false}\PY{p}{;}\PY{+w}{ }\PY{c+c1}{// worker checks this flag and exits cleanly}
\PY{+w}{    }\PY{p}{\PYZcb{}}
\PY{p}{\PYZcb{}}
\end{Verbatim}
\end{codebox}

Code review tip: when you see \code{@\allowbreak{}Deprecated} in your own project’s code (not a library), check the Javadoc for \code{forRemoval~\allowbreak{}=\allowbreak{}~\allowbreak{}true} and a \code{since} version. That tells you how urgent the migration is — \code{forRemoval~\allowbreak{}=\allowbreak{}~\allowbreak{}true} means it is scheduled to actually disappear.

\section[Common Code-Review Interview Pitfalls]{Common Code-Review Interview Pitfalls}
\label{h:3:common-code-review-interview-pitfalls}
\begin{enumerate}
\item 
\textbf{Using \code{==} to compare boxed types or records instead of \code{equals()}.}
Why it matters: \code{Integer}, \code{Long}, and record instances are value-based; identity comparison is misleading and can pass by accident for small cached values, then fail in production for larger ones.

\begin{codebox}
\begin{Verbatim}[commandchars=\\\{\}]
\PY{c+c1}{// Before (bug hiding in plain sight)}
\PY{n}{Integer}\PY{+w}{ }\PY{n}{a}\PY{+w}{ }\PY{o}{=}\PY{+w}{ }\PY{l+m+mi}{1000}\PY{p}{,}\PY{+w}{ }\PY{n}{b}\PY{+w}{ }\PY{o}{=}\PY{+w}{ }\PY{l+m+mi}{1000}\PY{p}{;}
\PY{k}{if}\PY{+w}{ }\PY{p}{(}\PY{n}{a}\PY{+w}{ }\PY{o}{=}\PY{o}{=}\PY{+w}{ }\PY{n}{b}\PY{p}{)}\PY{+w}{ }\PY{p}{\PYZob{}}\PY{+w}{ }\PY{c+cm}{/* false! outside Integer cache range */}\PY{+w}{ }\PY{p}{\PYZcb{}}

\PY{c+c1}{// After}
\PY{k}{if}\PY{+w}{ }\PY{p}{(}\PY{n}{a}\PY{p}{.}\PY{n+na}{equals}\PY{p}{(}\PY{n}{b}\PY{p}{)}\PY{p}{)}\PY{+w}{ }\PY{p}{\PYZob{}}\PY{+w}{ }\PY{c+cm}{/* correctly true */}\PY{+w}{ }\PY{p}{\PYZcb{}}
\end{Verbatim}
\end{codebox}

\item 
\textbf{Forgetting \code{permits} doesn’t limit exhaustiveness — missing a \code{default} on a non-sealed \code{switch}.}
Why it matters: only sealed hierarchies and enums give you compiler-enforced exhaustiveness. A \code{switch} over a plain interface or class still needs a \code{default}, or it won’t compile as an expression.

\begin{codebox}
\begin{Verbatim}[commandchars=\\\{\}]
\PY{c+c1}{// Before: compile error, no default and Shape is not sealed}
\PY{c+c1}{// String s = switch (shape) \PYZob{} case Circle c \PYZhy{}\PYZgt{} \PYZdq{}circle\PYZdq{}; \PYZcb{};}

\PY{c+c1}{// After: either seal Shape, or add a default}
\PY{n}{String}\PY{+w}{ }\PY{n}{s}\PY{+w}{ }\PY{o}{=}\PY{+w}{ }\PY{k}{switch}\PY{+w}{ }\PY{p}{(}\PY{n}{shape}\PY{p}{)}\PY{+w}{ }\PY{p}{\PYZob{}}
\PY{+w}{    }\PY{k}{case}\PY{+w}{ }\PY{n}{Circle}\PY{+w}{ }\PY{n}{c}\PY{+w}{ }\PY{o}{\PYZhy{}}\PY{o}{\PYZgt{}}\PY{+w}{ }\PY{l+s}{\PYZdq{}}\PY{l+s}{circle}\PY{l+s}{\PYZdq{}}\PY{p}{;}
\PY{+w}{    }\PY{k}{default}\PY{+w}{ }\PY{o}{\PYZhy{}}\PY{o}{\PYZgt{}}\PY{+w}{ }\PY{l+s}{\PYZdq{}}\PY{l+s}{other}\PY{l+s}{\PYZdq{}}\PY{p}{;}
\PY{p}{\PYZcb{}}\PY{p}{;}
\end{Verbatim}
\end{codebox}

\item 
\textbf{Assuming records are always the right replacement for a class.}
Why it matters: records cannot extend another class, and all fields are immutable. If a candidate needs inheritance or mutable state, a record is the wrong tool.

\begin{codebox}
\begin{Verbatim}[commandchars=\\\{\}]
\PY{c+c1}{// Before: trying to add mutable state to a record — won\PYZsq{}t compile}
\PY{c+c1}{// record Counter(int count) \PYZob{} void increment() \PYZob{} count++; \PYZcb{} \PYZcb{} // ERROR}

\PY{c+c1}{// After: use a regular class for mutable state}
\PY{k+kd}{class} \PY{n+nc}{Counter}\PY{+w}{ }\PY{p}{\PYZob{}}
\PY{+w}{    }\PY{k+kd}{private}\PY{+w}{ }\PY{k+kt}{int}\PY{+w}{ }\PY{n}{count}\PY{p}{;}
\PY{+w}{    }\PY{k+kt}{void}\PY{+w}{ }\PY{n+nf}{increment}\PY{p}{(}\PY{p}{)}\PY{+w}{ }\PY{p}{\PYZob{}}\PY{+w}{ }\PY{n}{count}\PY{o}{+}\PY{o}{+}\PY{p}{;}\PY{+w}{ }\PY{p}{\PYZcb{}}
\PY{p}{\PYZcb{}}
\end{Verbatim}
\end{codebox}

\item 
\textbf{Believing string templates (\code{STR.\allowbreak{}\textquotedbl{}...\allowbreak{}\textquotedbl{}}) are usable today.}
Why it matters: this feature was previewed in Java 21/22 and then withdrawn. Code using \code{STR.\allowbreak{}\textquotedbl{}...\allowbreak{}\textquotedbl{}} will not compile on any released JDK, including 21 through 25.

\begin{codebox}
\begin{Verbatim}[commandchars=\\\{\}]
\PY{c+c1}{// Before (does not compile on any released JDK):}
\PY{c+c1}{// String msg = STR.\PYZdq{}Hello \PYZbs{}\PYZob{}name\PYZcb{}\PYZdq{};}

\PY{c+c1}{// After (works today):}
\PY{n}{String}\PY{+w}{ }\PY{n}{msg}\PY{+w}{ }\PY{o}{=}\PY{+w}{ }\PY{l+s}{\PYZdq{}}\PY{l+s}{Hello \PYZpc{}s}\PY{l+s}{\PYZdq{}}\PY{p}{.}\PY{n+na}{formatted}\PY{p}{(}\PY{n}{name}\PY{p}{)}\PY{p}{;}
\end{Verbatim}
\end{codebox}

\item 
\textbf{Shipping preview-feature code without \code{--\allowbreak{}enable-\allowbreak{}preview} documented in the build.}
Why it matters: preview class files are tagged and will fail to run on a JDK that isn’t the exact matching preview-enabled version — this can break CI/CD or production silently.

\begin{codebox}
\begin{Verbatim}[commandchars=\\\{\}]
\PY{c+c1}{// Before: build script has no mention of preview, but code uses primitive patterns}
\PY{c+c1}{// (compiles locally with a preview\PYZhy{}enabled IDE, fails in CI)}

\PY{c+c1}{// After: explicit and documented}
\PY{c+c1}{// javac \PYZhy{}\PYZhy{}release 23 \PYZhy{}\PYZhy{}enable\PYZhy{}preview Main.java}
\PY{c+c1}{// java \PYZhy{}\PYZhy{}enable\PYZhy{}preview Main}
\end{Verbatim}
\end{codebox}

\item 
\textbf{Synchronizing on a value-based class instance (e.g. \code{Integer}, \code{Long}, boxed values).}
Why it matters: value-based instances may be cached or otherwise shared; two “different” locks might secretly be the same object, or a future JDK might change caching behavior, causing subtle deadlocks or lost synchronization.

\begin{codebox}
\begin{Verbatim}[commandchars=\\\{\}]
\PY{c+c1}{// Before}
\PY{n}{Integer}\PY{+w}{ }\PY{n}{lock}\PY{+w}{ }\PY{o}{=}\PY{+w}{ }\PY{l+m+mi}{1}\PY{p}{;}
\PY{k+kd}{synchronized}\PY{+w}{ }\PY{p}{(}\PY{n}{lock}\PY{p}{)}\PY{+w}{ }\PY{p}{\PYZob{}}\PY{+w}{ }\PY{c+cm}{/* risky: cached instance could be shared elsewhere */}\PY{+w}{ }\PY{p}{\PYZcb{}}

\PY{c+c1}{// After}
\PY{k+kd}{private}\PY{+w}{ }\PY{k+kd}{final}\PY{+w}{ }\PY{n}{Object}\PY{+w}{ }\PY{n}{lock}\PY{+w}{ }\PY{o}{=}\PY{+w}{ }\PY{k}{new}\PY{+w}{ }\PY{n}{Object}\PY{p}{(}\PY{p}{)}\PY{p}{;}
\PY{k+kd}{synchronized}\PY{+w}{ }\PY{p}{(}\PY{n}{lock}\PY{p}{)}\PY{+w}{ }\PY{p}{\PYZob{}}\PY{+w}{ }\PY{c+cm}{/* safe: unique identity, never reused */}\PY{+w}{ }\PY{p}{\PYZcb{}}
\end{Verbatim}
\end{codebox}

\item 
\textbf{Using \code{Thread.\allowbreak{}stop()} or relying on \code{finalize()} for cleanup.}
Why it matters: both are deprecated for removal and unsafe/unreliable. Reviewers should flag any use immediately.

\begin{codebox}
\begin{Verbatim}[commandchars=\\\{\}]
\PY{c+c1}{// Before}
\PY{c+c1}{// worker.stop(); // can corrupt shared state}
\PY{c+c1}{// protected void finalize() \PYZob{} resource.close(); \PYZcb{} // timing not guaranteed}

\PY{c+c1}{// After}
\PY{n}{worker}\PY{p}{.}\PY{n+na}{requestShutdown}\PY{p}{(}\PY{p}{)}\PY{p}{;}\PY{+w}{ }\PY{c+c1}{// cooperative flag}
\PY{k}{try}\PY{+w}{ }\PY{p}{(}\PY{k+kd}{var}\PY{+w}{ }\PY{n}{resource}\PY{+w}{ }\PY{o}{=}\PY{+w}{ }\PY{n}{openResource}\PY{p}{(}\PY{p}{)}\PY{p}{)}\PY{+w}{ }\PY{p}{\PYZob{}}\PY{+w}{ }\PY{c+cm}{/* auto\PYZhy{}closed */}\PY{+w}{ }\PY{p}{\PYZcb{}}
\end{Verbatim}
\end{codebox}

\item 
\textbf{Treating \code{default} interface methods as free multiple inheritance without checking for conflicts.}
Why it matters: implementing two interfaces with clashing \code{default} methods forces the implementing class to resolve the conflict explicitly, or the code won’t compile — this is easy to overlook when interfaces evolve independently.

\begin{codebox}
\begin{Verbatim}[commandchars=\\\{\}]
\PY{c+c1}{// Before: compile error if both A and B declare default String name()}
\PY{c+c1}{// class C implements A, B \PYZob{} \PYZcb{}}

\PY{c+c1}{// After: explicitly resolve}
\PY{k+kd}{class} \PY{n+nc}{C}\PY{+w}{ }\PY{k+kd}{implements}\PY{+w}{ }\PY{n}{A}\PY{p}{,}\PY{+w}{ }\PY{n}{B}\PY{+w}{ }\PY{p}{\PYZob{}}
\PY{+w}{    }\PY{k+kd}{public}\PY{+w}{ }\PY{n}{String}\PY{+w}{ }\PY{n+nf}{name}\PY{p}{(}\PY{p}{)}\PY{+w}{ }\PY{p}{\PYZob{}}\PY{+w}{ }\PY{k}{return}\PY{+w}{ }\PY{n}{A}\PY{p}{.}\PY{n+na}{super}\PY{p}{.}\PY{n+na}{name}\PY{p}{(}\PY{p}{)}\PY{p}{;}\PY{+w}{ }\PY{p}{\PYZcb{}}
\PY{p}{\PYZcb{}}
\end{Verbatim}
\end{codebox}

\item 
\textbf{Forgetting that text block indentation depends on the closing \code{\textquotedbl{}\textquotedbl{}\textquotedbl{}} position.}
Why it matters: moving the closing delimiter changes the output’s leading whitespace — a common source of “why does my JSON/SQL have extra spaces” bugs.

\begin{codebox}
\begin{Verbatim}[commandchars=\\\{\}]
\PY{c+c1}{// Before: closing \PYZdq{}\PYZdq{}\PYZdq{} at column 0 adds no stripping, keeping unwanted indentation}
\PY{n}{String}\PY{+w}{ }\PY{n}{s}\PY{+w}{ }\PY{o}{=}\PY{+w}{ }\PY{l+s}{\PYZdq{}\PYZdq{}\PYZdq{}}
\PY{l+s}{        line1}
\PY{l+s}{\PYZdq{}\PYZdq{}\PYZdq{}}\PY{p}{;}

\PY{c+c1}{// After: align closing delimiter with the intended left margin}
\PY{n}{String}\PY{+w}{ }\PY{n}{s}\PY{+w}{ }\PY{o}{=}\PY{+w}{ }\PY{l+s}{\PYZdq{}\PYZdq{}\PYZdq{}}
\PY{l+s}{        line1}
\PY{l+s}{        }\PY{l+s}{\PYZdq{}\PYZdq{}\PYZdq{}}\PY{p}{;}
\end{Verbatim}
\end{codebox}

\item 
\textbf{Using an unnamed variable \code{\_} and then trying to read it.}
Why it matters: \code{\_} is not a normal variable — the compiler forbids referencing it after declaration. Reusing multiple \code{\_} in the same scope is fine (that’s the point), but reading it is not.

\begin{codebox}
\begin{Verbatim}[commandchars=\\\{\}]
\PY{c+c1}{// Before: compile error}
\PY{c+c1}{// map.forEach((key, \PYZus{}) \PYZhy{}\PYZgt{} System.out.println(\PYZus{})); // ERROR: cannot use \PYZus{} as a value}

\PY{c+c1}{// After}
\PY{n}{map}\PY{p}{.}\PY{n+na}{forEach}\PY{p}{(}\PY{p}{(}\PY{n}{key}\PY{p}{,}\PY{+w}{ }\PY{n}{\PYZus{}}\PY{p}{)}\PY{+w}{ }\PY{o}{\PYZhy{}}\PY{o}{\PYZgt{}}\PY{+w}{ }\PY{n}{System}\PY{p}{.}\PY{n+na}{out}\PY{p}{.}\PY{n+na}{println}\PY{p}{(}\PY{n}{key}\PY{p}{)}\PY{p}{)}\PY{p}{;}
\end{Verbatim}
\end{codebox}

\item 
\textbf{Assuming an enum’s \code{ordinal()} is a stable identifier for persistence.}
Why it matters: \code{ordinal()} reflects declaration order. Reordering or inserting a new constant silently shifts every ordinal after it, corrupting any stored data (database rows, serialized files) that relied on the old numbering.

\begin{codebox}
\begin{Verbatim}[commandchars=\\\{\}]
\PY{c+c1}{// Before: storing ordinal() as a DB value}
\PY{c+c1}{// int code = status.ordinal(); // fragile: breaks if enum order changes}

\PY{c+c1}{// After: store an explicit, stable code}
\PY{k+kd}{enum}\PY{+w}{ }\PY{n}{OrderStatus}\PY{+w}{ }\PY{p}{\PYZob{}}
\PY{+w}{    }\PY{n}{NEW}\PY{p}{(}\PY{l+m+mi}{1}\PY{p}{)}\PY{p}{,}\PY{+w}{ }\PY{n}{PAID}\PY{p}{(}\PY{l+m+mi}{2}\PY{p}{)}\PY{p}{,}\PY{+w}{ }\PY{n}{SHIPPED}\PY{p}{(}\PY{l+m+mi}{3}\PY{p}{)}\PY{p}{;}
\PY{+w}{    }\PY{k+kd}{final}\PY{+w}{ }\PY{k+kt}{int}\PY{+w}{ }\PY{n}{code}\PY{p}{;}
\PY{+w}{    }\PY{n}{OrderStatus}\PY{p}{(}\PY{k+kt}{int}\PY{+w}{ }\PY{n}{code}\PY{p}{)}\PY{+w}{ }\PY{p}{\PYZob{}}\PY{+w}{ }\PY{k}{this}\PY{p}{.}\PY{n+na}{code}\PY{+w}{ }\PY{o}{=}\PY{+w}{ }\PY{n}{code}\PY{p}{;}\PY{+w}{ }\PY{p}{\PYZcb{}}
\PY{p}{\PYZcb{}}
\end{Verbatim}
\end{codebox}

\item 
\textbf{Marking a class \code{sealed} but forgetting \code{permits} subclasses must each declare \code{final}, \code{sealed}, or \code{non-\allowbreak{}sealed}.}
Why it matters: this is a compile error waiting to happen for anyone extending a sealed type for the first time, and reviewers should know it’s mandatory, not optional style.

\begin{codebox}
\begin{Verbatim}[commandchars=\\\{\}]
\PY{c+c1}{// Before: compile error — Car doesn\PYZsq{}t declare its inheritance modifier}
\PY{c+c1}{// sealed class Vehicle permits Car \PYZob{} \PYZcb{}}
\PY{c+c1}{// class Car extends Vehicle \PYZob{} \PYZcb{} // ERROR}

\PY{c+c1}{// After}
\PY{k+kd}{sealed}\PY{+w}{ }\PY{k+kd}{class} \PY{n+nc}{Vehicle}\PY{+w}{ }\PY{n}{permits}\PY{+w}{ }\PY{n}{Car}\PY{+w}{ }\PY{p}{\PYZob{}}\PY{+w}{ }\PY{p}{\PYZcb{}}
\PY{k+kd}{final}\PY{+w}{ }\PY{k+kd}{class} \PY{n+nc}{Car}\PY{+w}{ }\PY{k+kd}{extends}\PY{+w}{ }\PY{n}{Vehicle}\PY{+w}{ }\PY{p}{\PYZob{}}\PY{+w}{ }\PY{p}{\PYZcb{}}
\end{Verbatim}
\end{codebox}

\item 
\textbf{Confusing “deprecated” with “removed” and assuming old code will “just still work.”}
Why it matters: some features (Security Manager, Applets, Nashorn, CORBA modules) are fully removed, not just discouraged. Code or dependencies relying on them will fail to build or run on modern JDKs, which matters a lot when reviewing an upgrade PR.

\begin{codebox}
\begin{Verbatim}[commandchars=\\\{\}]
\PY{c+c1}{// Before (Java 8\PYZhy{}era code, will NOT run on Java 24+):}
\PY{c+c1}{// System.setSecurityManager(new SecurityManager()); // REMOVED in Java 24}

\PY{c+c1}{// After: use OS/container\PYZhy{}level sandboxing instead}
\PY{c+c1}{// (no direct JVM replacement — redesign the isolation strategy)}
\end{Verbatim}
\end{codebox}

\item 
\textbf{Using \code{new~\allowbreak{}Integer(...)} / \code{new~\allowbreak{}Boolean(...)} style boxed constructors in new code.}
Why it matters: these constructors are deprecated for removal and always allocate a new object, defeating the JVM’s integer cache and wasting memory for no benefit.

\begin{codebox}
\begin{Verbatim}[commandchars=\\\{\}]
\PY{c+c1}{// Before}
\PY{n}{Boolean}\PY{+w}{ }\PY{n}{flag}\PY{+w}{ }\PY{o}{=}\PY{+w}{ }\PY{k}{new}\PY{+w}{ }\PY{n}{Boolean}\PY{p}{(}\PY{k+kc}{true}\PY{p}{)}\PY{p}{;}\PY{+w}{ }\PY{c+c1}{// deprecated, wasteful}

\PY{c+c1}{// After}
\PY{n}{Boolean}\PY{+w}{ }\PY{n}{flag}\PY{+w}{ }\PY{o}{=}\PY{+w}{ }\PY{n}{Boolean}\PY{p}{.}\PY{n+na}{TRUE}\PY{p}{;}\PY{+w}{ }\PY{c+c1}{// or plain `true` with autoboxing}
\end{Verbatim}
\end{codebox}

\item 
\textbf{Overusing guarded patterns (\code{case~\allowbreak{}X~\allowbreak{}when~\allowbreak{}...}) instead of simple boolean logic, hurting readability.}
Why it matters: pattern matching in \code{switch} is powerful, but stacking many \code{when} clauses with complex conditions can make a \code{switch} harder to read than a plain \code{if/\allowbreak{}else} chain — reviewers should ask whether the pattern actually clarifies intent.

\begin{codebox}
\begin{Verbatim}[commandchars=\\\{\}]
\PY{c+c1}{// Before: hard to scan, conditions bury the intent}
\PY{n}{String}\PY{+w}{ }\PY{n}{result}\PY{+w}{ }\PY{o}{=}\PY{+w}{ }\PY{k}{switch}\PY{+w}{ }\PY{p}{(}\PY{n}{shape}\PY{p}{)}\PY{+w}{ }\PY{p}{\PYZob{}}
\PY{+w}{    }\PY{k}{case}\PY{+w}{ }\PY{n}{Circle}\PY{+w}{ }\PY{n}{c}\PY{+w}{ }\PY{n}{when}\PY{+w}{ }\PY{n}{c}\PY{p}{.}\PY{n+na}{radius}\PY{p}{(}\PY{p}{)}\PY{+w}{ }\PY{o}{\PYZgt{}}\PY{+w}{ }\PY{l+m+mi}{100}\PY{+w}{ }\PY{o}{\PYZam{}}\PY{o}{\PYZam{}}\PY{+w}{ }\PY{n}{c}\PY{p}{.}\PY{n+na}{radius}\PY{p}{(}\PY{p}{)}\PY{+w}{ }\PY{o}{\PYZlt{}}\PY{+w}{ }\PY{l+m+mi}{1000}\PY{+w}{ }\PY{o}{\PYZam{}}\PY{o}{\PYZam{}}\PY{+w}{ }\PY{n}{isValid}\PY{p}{(}\PY{n}{c}\PY{p}{)}\PY{+w}{ }\PY{o}{\PYZhy{}}\PY{o}{\PYZgt{}}\PY{+w}{ }\PY{l+s}{\PYZdq{}}\PY{l+s}{big circle}\PY{l+s}{\PYZdq{}}\PY{p}{;}
\PY{+w}{    }\PY{k}{default}\PY{+w}{ }\PY{o}{\PYZhy{}}\PY{o}{\PYZgt{}}\PY{+w}{ }\PY{l+s}{\PYZdq{}}\PY{l+s}{other}\PY{l+s}{\PYZdq{}}\PY{p}{;}
\PY{p}{\PYZcb{}}\PY{p}{;}

\PY{c+c1}{// After: extract the condition to a named method for clarity}
\PY{n}{String}\PY{+w}{ }\PY{n}{result}\PY{+w}{ }\PY{o}{=}\PY{+w}{ }\PY{k}{switch}\PY{+w}{ }\PY{p}{(}\PY{n}{shape}\PY{p}{)}\PY{+w}{ }\PY{p}{\PYZob{}}
\PY{+w}{    }\PY{k}{case}\PY{+w}{ }\PY{n}{Circle}\PY{+w}{ }\PY{n}{c}\PY{+w}{ }\PY{n}{when}\PY{+w}{ }\PY{n+nf}{isBigValidCircle}\PY{p}{(}\PY{n}{c}\PY{p}{)}\PY{+w}{ }\PY{o}{\PYZhy{}}\PY{o}{\PYZgt{}}\PY{+w}{ }\PY{l+s}{\PYZdq{}}\PY{l+s}{big circle}\PY{l+s}{\PYZdq{}}\PY{p}{;}
\PY{+w}{    }\PY{k}{default}\PY{+w}{ }\PY{o}{\PYZhy{}}\PY{o}{\PYZgt{}}\PY{+w}{ }\PY{l+s}{\PYZdq{}}\PY{l+s}{other}\PY{l+s}{\PYZdq{}}\PY{p}{;}
\PY{p}{\PYZcb{}}\PY{p}{;}
\end{Verbatim}
\end{codebox}

\end{enumerate}


\setcounter{chapter}{3}
\chapter{Generics and Type System}
\label{chap:4}
\label{h:4:4-generics-and-type-system}
Generics let you write code that works with different types while keeping type safety at compile time. They are one of the most interview-heavy topics in Java because the compiler does a lot of invisible work — and that invisible work (erasure) causes surprising rules that trip up even experienced developers. This chapter builds up from basic generic classes to the trickiest interview questions: erasure, bridge methods, wildcards, and why arrays and generics disagree about variance.

By the end, you should be able to explain \emph{why} a rule exists, not just recite it. That is what separates a strong code-review answer from a memorized one.

\begin{itemize}
\item 
\hyperref[h:4:generics]{Generics}

\begin{itemize}
\item 
\hyperref[h:4:why-generics-exist]{Why generics exist}

\item 
\hyperref[h:4:generic-classes]{Generic classes}

\item 
\hyperref[h:4:generic-methods]{Generic methods}

\item 
\hyperref[h:4:bounded-type-parameters]{Bounded type parameters}

\item 
\hyperref[h:4:multiple-bounds]{Multiple bounds}

\item 
\hyperref[h:4:recursive-generic-bounds-t-extends-comparablet]{Recursive generic bounds}

\item 
\hyperref[h:4:raw-types-and-unchecked-warnings]{Raw types and unchecked warnings}

\item 
\hyperref[h:4:generic-type-inference-and-the-diamond-operator]{Generic type inference and the diamond operator}

\item 
\hyperref[h:4:var-with-generics]{\code{var} with generics}

\item 
\hyperref[h:4:generics-and-optional--streams]{Generics and \code{Optional} / streams}

\end{itemize}

\item 
\hyperref[h:4:type-erasure]{Type Erasure}

\begin{itemize}
\item 
\hyperref[h:4:what-erasure-does-at-the-bytecode-level]{What erasure does at the bytecode level}

\item 
\hyperref[h:4:bridge-methods]{Bridge methods}

\item 
\hyperref[h:4:why-you-cannot-do-new-t]{Why you cannot do \code{new~\allowbreak{}T[]}}

\item 
\hyperref[h:4:why-you-cannot-do-instanceof-liststring]{Why you cannot do \code{instanceof~\allowbreak{}List\textless{}\allowbreak{}String\textgreater{}}}

\item 
\hyperref[h:4:why-you-cannot-have-static-generic-fields]{Why you cannot have static generic fields}

\item 
\hyperref[h:4:heap-pollution-and-safevarargs]{Heap pollution and \code{@\allowbreak{}SafeVarargs}}

\end{itemize}

\item 
\hyperref[h:4:wildcards-extends-super]{Wildcards (extends, super)}

\begin{itemize}
\item 
\hyperref[h:4:unbounded-wildcards-]{Unbounded wildcards \code{\textless{}?\textgreater{}}}

\item 
\hyperref[h:4:upper-bounded-wildcards--extends-t]{Upper-bounded wildcards \code{\textless{}?\allowbreak{}~\allowbreak{}extends~\allowbreak{}T\textgreater{}}}

\item 
\hyperref[h:4:lower-bounded-wildcards--super-t]{Lower-bounded wildcards \code{\textless{}?\allowbreak{}~\allowbreak{}super~\allowbreak{}T\textgreater{}}}

\item 
\hyperref[h:4:pecs-producer-extends-consumer-super]{PECS: Producer Extends, Consumer Super}

\item 
\hyperref[h:4:list-vs-listobject-vs-raw-list]{\code{List\textless{}?\textgreater{}} vs \code{List\textless{}\allowbreak{}Object\textgreater{}} vs raw \code{List}}

\end{itemize}

\item 
\hyperref[h:4:covariance-and-contravariance]{Covariance and Contravariance}

\begin{itemize}
\item 
\hyperref[h:4:arrays-are-covariant]{Arrays are covariant}

\item 
\hyperref[h:4:generics-are-invariant]{Generics are invariant}

\item 
\hyperref[h:4:arraystoreexception-demo]{\code{ArrayStoreException} demo}

\item 
\hyperref[h:4:method-return-type-covariance]{Method return type covariance}

\end{itemize}

\item 
\hyperref[h:4:common-code-review-interview-pitfalls]{Common Code-Review Interview Pitfalls}

\end{itemize}

\section[Generics]{Generics}
\label{h:4:generics}
\textbf{Generics} let a class, interface, or method work with a \emph{type parameter} — a placeholder for a real type that gets filled in later. Think of a type parameter like a function argument, except it stands for a type (like \code{String} or \code{Integer}) instead of a value.

\subsection[Why generics exist]{Why generics exist}
\label{h:4:why-generics-exist}
Before Java 5 (2004), collections held plain \code{Object} references. You had to cast everything, and mistakes only showed up at runtime.

\begin{codebox}
\begin{Verbatim}[commandchars=\\\{\}]
\PY{c+c1}{// Pre\PYZhy{}generics style (Java 1.4 and earlier) — do not write this today}
\PY{n}{List}\PY{+w}{ }\PY{n}{oldList}\PY{+w}{ }\PY{o}{=}\PY{+w}{ }\PY{k}{new}\PY{+w}{ }\PY{n}{ArrayList}\PY{p}{(}\PY{p}{)}\PY{p}{;}
\PY{n}{oldList}\PY{p}{.}\PY{n+na}{add}\PY{p}{(}\PY{l+s}{\PYZdq{}}\PY{l+s}{hello}\PY{l+s}{\PYZdq{}}\PY{p}{)}\PY{p}{;}
\PY{n}{oldList}\PY{p}{.}\PY{n+na}{add}\PY{p}{(}\PY{l+m+mi}{42}\PY{p}{)}\PY{p}{;}\PY{+w}{ }\PY{c+c1}{// no compile\PYZhy{}time check, mixed types allowed}

\PY{n}{String}\PY{+w}{ }\PY{n}{s}\PY{+w}{ }\PY{o}{=}\PY{+w}{ }\PY{p}{(}\PY{n}{String}\PY{p}{)}\PY{+w}{ }\PY{n}{oldList}\PY{p}{.}\PY{n+na}{get}\PY{p}{(}\PY{l+m+mi}{1}\PY{p}{)}\PY{p}{;}\PY{+w}{ }\PY{c+c1}{// compiles fine...}
\PY{c+c1}{// ...but throws java.lang.ClassCastException at runtime}
\end{Verbatim}
\end{codebox}

Generics move that error from runtime to compile time:

\begin{codebox}
\begin{Verbatim}[commandchars=\\\{\}]
\PY{n}{List}\PY{o}{\PYZlt{}}\PY{n}{String}\PY{o}{\PYZgt{}}\PY{+w}{ }\PY{n}{names}\PY{+w}{ }\PY{o}{=}\PY{+w}{ }\PY{k}{new}\PY{+w}{ }\PY{n}{ArrayList}\PY{o}{\PYZlt{}}\PY{o}{\PYZgt{}}\PY{p}{(}\PY{p}{)}\PY{p}{;}
\PY{n}{names}\PY{p}{.}\PY{n+na}{add}\PY{p}{(}\PY{l+s}{\PYZdq{}}\PY{l+s}{Ada}\PY{l+s}{\PYZdq{}}\PY{p}{)}\PY{p}{;}
\PY{c+c1}{// names.add(42); // compile error: incompatible types: int cannot be converted to String}
\PY{n}{String}\PY{+w}{ }\PY{n}{first}\PY{+w}{ }\PY{o}{=}\PY{+w}{ }\PY{n}{names}\PY{p}{.}\PY{n+na}{get}\PY{p}{(}\PY{l+m+mi}{0}\PY{p}{)}\PY{p}{;}\PY{+w}{ }\PY{c+c1}{// no cast needed}
\end{Verbatim}
\end{codebox}

That is the whole point of generics: \textbf{catch type mistakes earlier, and remove casts}.

\subsection[Generic classes]{Generic classes}
\label{h:4:generic-classes}
A generic class declares one or more type parameters in angle brackets after the class name. By convention, type parameters are single, uppercase letters: \code{T} (Type), \code{E} (Element), \code{K}/\code{V} (Key/Value), \code{R} (Result).

\begin{codebox}
\begin{Verbatim}[commandchars=\\\{\}]
\PY{k+kd}{public}\PY{+w}{ }\PY{k+kd}{class} \PY{n+nc}{Box}\PY{o}{\PYZlt{}}\PY{n}{T}\PY{o}{\PYZgt{}}\PY{+w}{ }\PY{p}{\PYZob{}}
\PY{+w}{    }\PY{k+kd}{private}\PY{+w}{ }\PY{n}{T}\PY{+w}{ }\PY{n}{content}\PY{p}{;}

\PY{+w}{    }\PY{k+kd}{public}\PY{+w}{ }\PY{k+kt}{void}\PY{+w}{ }\PY{n+nf}{set}\PY{p}{(}\PY{n}{T}\PY{+w}{ }\PY{n}{content}\PY{p}{)}\PY{+w}{ }\PY{p}{\PYZob{}}
\PY{+w}{        }\PY{k}{this}\PY{p}{.}\PY{n+na}{content}\PY{+w}{ }\PY{o}{=}\PY{+w}{ }\PY{n}{content}\PY{p}{;}
\PY{+w}{    }\PY{p}{\PYZcb{}}

\PY{+w}{    }\PY{k+kd}{public}\PY{+w}{ }\PY{n}{T}\PY{+w}{ }\PY{n+nf}{get}\PY{p}{(}\PY{p}{)}\PY{+w}{ }\PY{p}{\PYZob{}}
\PY{+w}{        }\PY{k}{return}\PY{+w}{ }\PY{n}{content}\PY{p}{;}
\PY{+w}{    }\PY{p}{\PYZcb{}}
\PY{p}{\PYZcb{}}
\end{Verbatim}
\end{codebox}

Usage:

\begin{codebox}
\begin{Verbatim}[commandchars=\\\{\}]
\PY{n}{Box}\PY{o}{\PYZlt{}}\PY{n}{String}\PY{o}{\PYZgt{}}\PY{+w}{ }\PY{n}{stringBox}\PY{+w}{ }\PY{o}{=}\PY{+w}{ }\PY{k}{new}\PY{+w}{ }\PY{n}{Box}\PY{o}{\PYZlt{}}\PY{o}{\PYZgt{}}\PY{p}{(}\PY{p}{)}\PY{p}{;}
\PY{n}{stringBox}\PY{p}{.}\PY{n+na}{set}\PY{p}{(}\PY{l+s}{\PYZdq{}}\PY{l+s}{hello}\PY{l+s}{\PYZdq{}}\PY{p}{)}\PY{p}{;}
\PY{n}{String}\PY{+w}{ }\PY{n}{value}\PY{+w}{ }\PY{o}{=}\PY{+w}{ }\PY{n}{stringBox}\PY{p}{.}\PY{n+na}{get}\PY{p}{(}\PY{p}{)}\PY{p}{;}\PY{+w}{ }\PY{c+c1}{// no cast}

\PY{n}{Box}\PY{o}{\PYZlt{}}\PY{n}{Integer}\PY{o}{\PYZgt{}}\PY{+w}{ }\PY{n}{intBox}\PY{+w}{ }\PY{o}{=}\PY{+w}{ }\PY{k}{new}\PY{+w}{ }\PY{n}{Box}\PY{o}{\PYZlt{}}\PY{o}{\PYZgt{}}\PY{p}{(}\PY{p}{)}\PY{p}{;}
\PY{n}{intBox}\PY{p}{.}\PY{n+na}{set}\PY{p}{(}\PY{l+m+mi}{5}\PY{p}{)}\PY{p}{;}
\PY{c+c1}{// intBox.set(\PYZdq{}oops\PYZdq{}); // compile error: incompatible types: String cannot be converted to Integer}
\end{Verbatim}
\end{codebox}

A class can have multiple type parameters:

\begin{codebox}
\begin{Verbatim}[commandchars=\\\{\}]
\PY{k+kd}{public}\PY{+w}{ }\PY{k+kd}{class} \PY{n+nc}{Pair}\PY{o}{\PYZlt{}}\PY{n}{K}\PY{p}{,}\PY{+w}{ }\PY{n}{V}\PY{o}{\PYZgt{}}\PY{+w}{ }\PY{p}{\PYZob{}}
\PY{+w}{    }\PY{k+kd}{private}\PY{+w}{ }\PY{k+kd}{final}\PY{+w}{ }\PY{n}{K}\PY{+w}{ }\PY{n}{key}\PY{p}{;}
\PY{+w}{    }\PY{k+kd}{private}\PY{+w}{ }\PY{k+kd}{final}\PY{+w}{ }\PY{n}{V}\PY{+w}{ }\PY{n}{value}\PY{p}{;}

\PY{+w}{    }\PY{k+kd}{public}\PY{+w}{ }\PY{n+nf}{Pair}\PY{p}{(}\PY{n}{K}\PY{+w}{ }\PY{n}{key}\PY{p}{,}\PY{+w}{ }\PY{n}{V}\PY{+w}{ }\PY{n}{value}\PY{p}{)}\PY{+w}{ }\PY{p}{\PYZob{}}
\PY{+w}{        }\PY{k}{this}\PY{p}{.}\PY{n+na}{key}\PY{+w}{ }\PY{o}{=}\PY{+w}{ }\PY{n}{key}\PY{p}{;}
\PY{+w}{        }\PY{k}{this}\PY{p}{.}\PY{n+na}{value}\PY{+w}{ }\PY{o}{=}\PY{+w}{ }\PY{n}{value}\PY{p}{;}
\PY{+w}{    }\PY{p}{\PYZcb{}}

\PY{+w}{    }\PY{k+kd}{public}\PY{+w}{ }\PY{n}{K}\PY{+w}{ }\PY{n+nf}{getKey}\PY{p}{(}\PY{p}{)}\PY{+w}{   }\PY{p}{\PYZob{}}\PY{+w}{ }\PY{k}{return}\PY{+w}{ }\PY{n}{key}\PY{p}{;}\PY{+w}{ }\PY{p}{\PYZcb{}}
\PY{+w}{    }\PY{k+kd}{public}\PY{+w}{ }\PY{n}{V}\PY{+w}{ }\PY{n+nf}{getValue}\PY{p}{(}\PY{p}{)}\PY{+w}{ }\PY{p}{\PYZob{}}\PY{+w}{ }\PY{k}{return}\PY{+w}{ }\PY{n}{value}\PY{p}{;}\PY{+w}{ }\PY{p}{\PYZcb{}}
\PY{p}{\PYZcb{}}

\PY{n}{Pair}\PY{o}{\PYZlt{}}\PY{n}{String}\PY{p}{,}\PY{+w}{ }\PY{n}{Integer}\PY{o}{\PYZgt{}}\PY{+w}{ }\PY{n}{ageOf}\PY{+w}{ }\PY{o}{=}\PY{+w}{ }\PY{k}{new}\PY{+w}{ }\PY{n}{Pair}\PY{o}{\PYZlt{}}\PY{o}{\PYZgt{}}\PY{p}{(}\PY{l+s}{\PYZdq{}}\PY{l+s}{Grace}\PY{l+s}{\PYZdq{}}\PY{p}{,}\PY{+w}{ }\PY{l+m+mi}{85}\PY{p}{)}\PY{p}{;}
\end{Verbatim}
\end{codebox}

\subsection[Generic methods]{Generic methods}
\label{h:4:generic-methods}
A method can introduce its own type parameter, independent of the class it lives in. The type parameter is declared right before the return type.

\begin{codebox}
\begin{Verbatim}[commandchars=\\\{\}]
\PY{k+kd}{public}\PY{+w}{ }\PY{k+kd}{class} \PY{n+nc}{Utils}\PY{+w}{ }\PY{p}{\PYZob{}}
\PY{+w}{    }\PY{c+c1}{// \PYZlt{}T\PYZgt{} here is the method\PYZsq{}s own type parameter}
\PY{+w}{    }\PY{k+kd}{public}\PY{+w}{ }\PY{k+kd}{static}\PY{+w}{ }\PY{o}{\PYZlt{}}\PY{n}{T}\PY{o}{\PYZgt{}}\PY{+w}{ }\PY{n}{T}\PY{+w}{ }\PY{n+nf}{firstNonNull}\PY{p}{(}\PY{n}{T}\PY{+w}{ }\PY{n}{a}\PY{p}{,}\PY{+w}{ }\PY{n}{T}\PY{+w}{ }\PY{n}{b}\PY{p}{)}\PY{+w}{ }\PY{p}{\PYZob{}}
\PY{+w}{        }\PY{k}{return}\PY{+w}{ }\PY{n}{a}\PY{+w}{ }\PY{o}{!}\PY{o}{=}\PY{+w}{ }\PY{k+kc}{null}\PY{+w}{ }\PY{o}{?}\PY{+w}{ }\PY{n}{a}\PY{+w}{ }\PY{p}{:}\PY{+w}{ }\PY{n}{b}\PY{p}{;}
\PY{+w}{    }\PY{p}{\PYZcb{}}

\PY{+w}{    }\PY{c+c1}{// Works with any two types that share a common supertype}
\PY{+w}{    }\PY{k+kd}{public}\PY{+w}{ }\PY{k+kd}{static}\PY{+w}{ }\PY{o}{\PYZlt{}}\PY{n}{T}\PY{o}{\PYZgt{}}\PY{+w}{ }\PY{n}{List}\PY{o}{\PYZlt{}}\PY{n}{T}\PY{o}{\PYZgt{}}\PY{+w}{ }\PY{n+nf}{listOf}\PY{p}{(}\PY{n}{T}\PY{+w}{ }\PY{n}{a}\PY{p}{,}\PY{+w}{ }\PY{n}{T}\PY{+w}{ }\PY{n}{b}\PY{p}{)}\PY{+w}{ }\PY{p}{\PYZob{}}
\PY{+w}{        }\PY{n}{List}\PY{o}{\PYZlt{}}\PY{n}{T}\PY{o}{\PYZgt{}}\PY{+w}{ }\PY{n}{list}\PY{+w}{ }\PY{o}{=}\PY{+w}{ }\PY{k}{new}\PY{+w}{ }\PY{n}{ArrayList}\PY{o}{\PYZlt{}}\PY{o}{\PYZgt{}}\PY{p}{(}\PY{p}{)}\PY{p}{;}
\PY{+w}{        }\PY{n}{list}\PY{p}{.}\PY{n+na}{add}\PY{p}{(}\PY{n}{a}\PY{p}{)}\PY{p}{;}
\PY{+w}{        }\PY{n}{list}\PY{p}{.}\PY{n+na}{add}\PY{p}{(}\PY{n}{b}\PY{p}{)}\PY{p}{;}
\PY{+w}{        }\PY{k}{return}\PY{+w}{ }\PY{n}{list}\PY{p}{;}
\PY{+w}{    }\PY{p}{\PYZcb{}}
\PY{p}{\PYZcb{}}

\PY{n}{String}\PY{+w}{ }\PY{n}{result}\PY{+w}{ }\PY{o}{=}\PY{+w}{ }\PY{n}{Utils}\PY{p}{.}\PY{n+na}{firstNonNull}\PY{p}{(}\PY{k+kc}{null}\PY{p}{,}\PY{+w}{ }\PY{l+s}{\PYZdq{}}\PY{l+s}{fallback}\PY{l+s}{\PYZdq{}}\PY{p}{)}\PY{p}{;}\PY{+w}{ }\PY{c+c1}{// T inferred as String}
\PY{n}{List}\PY{o}{\PYZlt{}}\PY{n}{Integer}\PY{o}{\PYZgt{}}\PY{+w}{ }\PY{n}{nums}\PY{+w}{ }\PY{o}{=}\PY{+w}{ }\PY{n}{Utils}\PY{p}{.}\PY{n+na}{listOf}\PY{p}{(}\PY{l+m+mi}{1}\PY{p}{,}\PY{+w}{ }\PY{l+m+mi}{2}\PY{p}{)}\PY{p}{;}\PY{+w}{               }\PY{c+c1}{// T inferred as Integer}
\end{Verbatim}
\end{codebox}

You rarely need to spell out the type parameter explicitly — the compiler infers it from the arguments. But you can be explicit if inference fails:

\begin{codebox}
\begin{Verbatim}[commandchars=\\\{\}]
\PY{n}{List}\PY{o}{\PYZlt{}}\PY{n}{String}\PY{o}{\PYZgt{}}\PY{+w}{ }\PY{n}{empty}\PY{+w}{ }\PY{o}{=}\PY{+w}{ }\PY{n}{Utils}\PY{p}{.}\PY{o}{\PYZlt{}}\PY{n}{String}\PY{o}{\PYZgt{}}\PY{n}{listOf}\PY{p}{(}\PY{k+kc}{null}\PY{p}{,}\PY{+w}{ }\PY{k+kc}{null}\PY{p}{)}\PY{p}{;}
\end{Verbatim}
\end{codebox}

\subsection[Bounded type parameters]{Bounded type parameters}
\label{h:4:bounded-type-parameters}
Sometimes you want to restrict what types are allowed. An \textbf{upper bound} (\code{extends}) says “T must be this type or a subtype of it.” This lets you call methods from the bound inside the generic code.

\begin{codebox}
\begin{Verbatim}[commandchars=\\\{\}]
\PY{c+c1}{// Without a bound, you cannot call numeric methods on T}
\PY{k+kd}{public}\PY{+w}{ }\PY{k+kd}{class} \PY{n+nc}{NumberBox}\PY{o}{\PYZlt{}}\PY{n}{T}\PY{o}{\PYZgt{}}\PY{+w}{ }\PY{p}{\PYZob{}}
\PY{+w}{    }\PY{k+kd}{private}\PY{+w}{ }\PY{n}{T}\PY{+w}{ }\PY{n}{value}\PY{p}{;}
\PY{+w}{    }\PY{c+c1}{// double asDouble() \PYZob{} return value.doubleValue(); \PYZcb{} // compile error: cannot find symbol}
\PY{p}{\PYZcb{}}

\PY{c+c1}{// With an upper bound, Number\PYZsq{}s methods become available}
\PY{k+kd}{public}\PY{+w}{ }\PY{k+kd}{class} \PY{n+nc}{BoundedNumberBox}\PY{o}{\PYZlt{}}\PY{n}{T}\PY{+w}{ }\PY{k+kd}{extends}\PY{+w}{ }\PY{n}{Number}\PY{o}{\PYZgt{}}\PY{+w}{ }\PY{p}{\PYZob{}}
\PY{+w}{    }\PY{k+kd}{private}\PY{+w}{ }\PY{n}{T}\PY{+w}{ }\PY{n}{value}\PY{p}{;}

\PY{+w}{    }\PY{k+kd}{public}\PY{+w}{ }\PY{n+nf}{BoundedNumberBox}\PY{p}{(}\PY{n}{T}\PY{+w}{ }\PY{n}{value}\PY{p}{)}\PY{+w}{ }\PY{p}{\PYZob{}}
\PY{+w}{        }\PY{k}{this}\PY{p}{.}\PY{n+na}{value}\PY{+w}{ }\PY{o}{=}\PY{+w}{ }\PY{n}{value}\PY{p}{;}
\PY{+w}{    }\PY{p}{\PYZcb{}}

\PY{+w}{    }\PY{k+kd}{public}\PY{+w}{ }\PY{k+kt}{double}\PY{+w}{ }\PY{n+nf}{asDouble}\PY{p}{(}\PY{p}{)}\PY{+w}{ }\PY{p}{\PYZob{}}
\PY{+w}{        }\PY{k}{return}\PY{+w}{ }\PY{n}{value}\PY{p}{.}\PY{n+na}{doubleValue}\PY{p}{(}\PY{p}{)}\PY{p}{;}\PY{+w}{ }\PY{c+c1}{// OK — every Number has doubleValue()}
\PY{+w}{    }\PY{p}{\PYZcb{}}
\PY{p}{\PYZcb{}}

\PY{n}{BoundedNumberBox}\PY{o}{\PYZlt{}}\PY{n}{Integer}\PY{o}{\PYZgt{}}\PY{+w}{ }\PY{n}{b1}\PY{+w}{ }\PY{o}{=}\PY{+w}{ }\PY{k}{new}\PY{+w}{ }\PY{n}{BoundedNumberBox}\PY{o}{\PYZlt{}}\PY{o}{\PYZgt{}}\PY{p}{(}\PY{l+m+mi}{5}\PY{p}{)}\PY{p}{;}
\PY{n}{BoundedNumberBox}\PY{o}{\PYZlt{}}\PY{n}{Double}\PY{o}{\PYZgt{}}\PY{+w}{ }\PY{n}{b2}\PY{+w}{ }\PY{o}{=}\PY{+w}{ }\PY{k}{new}\PY{+w}{ }\PY{n}{BoundedNumberBox}\PY{o}{\PYZlt{}}\PY{o}{\PYZgt{}}\PY{p}{(}\PY{l+m+mf}{5.5}\PY{p}{)}\PY{p}{;}
\PY{c+c1}{// BoundedNumberBox\PYZlt{}String\PYZgt{} b3 = new BoundedNumberBox\PYZlt{}\PYZgt{}(\PYZdq{}x\PYZdq{});}
\PY{c+c1}{// compile error: String is not a subtype of Number}
\end{Verbatim}
\end{codebox}

Note that \code{extends} here means “extends or implements” — it works for both classes and interfaces.

\subsection[Multiple bounds]{Multiple bounds}
\label{h:4:multiple-bounds}
A type parameter can have more than one bound: at most one class, plus any number of interfaces, joined with \code{\&}. The class (if present) must come first.

\begin{codebox}
\begin{Verbatim}[commandchars=\\\{\}]
\PY{k+kd}{interface} \PY{n+nc}{Named}\PY{+w}{ }\PY{p}{\PYZob{}}
\PY{+w}{    }\PY{n}{String}\PY{+w}{ }\PY{n+nf}{name}\PY{p}{(}\PY{p}{)}\PY{p}{;}
\PY{p}{\PYZcb{}}

\PY{c+c1}{// T must be a Number AND implement Named}
\PY{k+kd}{static}\PY{+w}{ }\PY{o}{\PYZlt{}}\PY{n}{T}\PY{+w}{ }\PY{k+kd}{extends}\PY{+w}{ }\PY{n}{Number}\PY{+w}{ }\PY{o}{\PYZam{}}\PY{+w}{ }\PY{n}{Named}\PY{o}{\PYZgt{}}\PY{+w}{ }\PY{k+kt}{void}\PY{+w}{ }\PY{n+nf}{printLabeled}\PY{p}{(}\PY{n}{T}\PY{+w}{ }\PY{n}{value}\PY{p}{)}\PY{+w}{ }\PY{p}{\PYZob{}}
\PY{+w}{    }\PY{n}{System}\PY{p}{.}\PY{n+na}{out}\PY{p}{.}\PY{n+na}{println}\PY{p}{(}\PY{n}{value}\PY{p}{.}\PY{n+na}{name}\PY{p}{(}\PY{p}{)}\PY{+w}{ }\PY{o}{+}\PY{+w}{ }\PY{l+s}{\PYZdq{}}\PY{l+s}{ = }\PY{l+s}{\PYZdq{}}\PY{+w}{ }\PY{o}{+}\PY{+w}{ }\PY{n}{value}\PY{p}{.}\PY{n+na}{doubleValue}\PY{p}{(}\PY{p}{)}\PY{p}{)}\PY{p}{;}
\PY{p}{\PYZcb{}}

\PY{c+c1}{// class Weird\PYZlt{}T extends Named \PYZam{} Number\PYZgt{} \PYZob{}\PYZcb{} // compile error: class types must come first}
\end{Verbatim}
\end{codebox}

\subsection[Recursive generic bounds (\textless{}T extends Comparable\textless{}T\textgreater{}\textgreater{})]{Recursive generic bounds (\code{\textless{}\allowbreak{}T~\allowbreak{}extends~\allowbreak{}Comparable\textless{}\allowbreak{}T\textgreater{}\textgreater{}})}
\label{h:4:recursive-generic-bounds-t-extends-comparablet}
This pattern shows up constantly in interviews: a type parameter bounded by a generic type that refers back to itself. It says “T must be comparable to other T’s.”

\begin{codebox}
\begin{Verbatim}[commandchars=\\\{\}]
\PY{k+kd}{public}\PY{+w}{ }\PY{k+kd}{static}\PY{+w}{ }\PY{o}{\PYZlt{}}\PY{n}{T}\PY{+w}{ }\PY{k+kd}{extends}\PY{+w}{ }\PY{n}{Comparable}\PY{o}{\PYZlt{}}\PY{n}{T}\PY{o}{\PYZgt{}}\PY{o}{\PYZgt{}}\PY{+w}{ }\PY{n}{T}\PY{+w}{ }\PY{n+nf}{max}\PY{p}{(}\PY{n}{List}\PY{o}{\PYZlt{}}\PY{n}{T}\PY{o}{\PYZgt{}}\PY{+w}{ }\PY{n}{list}\PY{p}{)}\PY{+w}{ }\PY{p}{\PYZob{}}
\PY{+w}{    }\PY{n}{T}\PY{+w}{ }\PY{n}{best}\PY{+w}{ }\PY{o}{=}\PY{+w}{ }\PY{n}{list}\PY{p}{.}\PY{n+na}{get}\PY{p}{(}\PY{l+m+mi}{0}\PY{p}{)}\PY{p}{;}
\PY{+w}{    }\PY{k}{for}\PY{+w}{ }\PY{p}{(}\PY{n}{T}\PY{+w}{ }\PY{n}{item}\PY{+w}{ }\PY{p}{:}\PY{+w}{ }\PY{n}{list}\PY{p}{)}\PY{+w}{ }\PY{p}{\PYZob{}}
\PY{+w}{        }\PY{k}{if}\PY{+w}{ }\PY{p}{(}\PY{n}{item}\PY{p}{.}\PY{n+na}{compareTo}\PY{p}{(}\PY{n}{best}\PY{p}{)}\PY{+w}{ }\PY{o}{\PYZgt{}}\PY{+w}{ }\PY{l+m+mi}{0}\PY{p}{)}\PY{+w}{ }\PY{p}{\PYZob{}}
\PY{+w}{            }\PY{n}{best}\PY{+w}{ }\PY{o}{=}\PY{+w}{ }\PY{n}{item}\PY{p}{;}
\PY{+w}{        }\PY{p}{\PYZcb{}}
\PY{+w}{    }\PY{p}{\PYZcb{}}
\PY{+w}{    }\PY{k}{return}\PY{+w}{ }\PY{n}{best}\PY{p}{;}
\PY{p}{\PYZcb{}}

\PY{n}{List}\PY{o}{\PYZlt{}}\PY{n}{Integer}\PY{o}{\PYZgt{}}\PY{+w}{ }\PY{n}{nums}\PY{+w}{ }\PY{o}{=}\PY{+w}{ }\PY{n}{List}\PY{p}{.}\PY{n+na}{of}\PY{p}{(}\PY{l+m+mi}{3}\PY{p}{,}\PY{+w}{ }\PY{l+m+mi}{7}\PY{p}{,}\PY{+w}{ }\PY{l+m+mi}{2}\PY{p}{)}\PY{p}{;}
\PY{n}{Integer}\PY{+w}{ }\PY{n}{top}\PY{+w}{ }\PY{o}{=}\PY{+w}{ }\PY{n}{max}\PY{p}{(}\PY{n}{nums}\PY{p}{)}\PY{p}{;}\PY{+w}{ }\PY{c+c1}{// works, Integer implements Comparable\PYZlt{}Integer\PYZgt{}}

\PY{k+kd}{class} \PY{n+nc}{Point}\PY{+w}{ }\PY{p}{\PYZob{}}\PY{+w}{ }\PY{k+kt}{int}\PY{+w}{ }\PY{n}{x}\PY{p}{,}\PY{+w}{ }\PY{n}{y}\PY{p}{;}\PY{+w}{ }\PY{p}{\PYZcb{}}\PY{+w}{ }\PY{c+c1}{// does NOT implement Comparable}
\PY{n}{List}\PY{o}{\PYZlt{}}\PY{n}{Point}\PY{o}{\PYZgt{}}\PY{+w}{ }\PY{n}{points}\PY{+w}{ }\PY{o}{=}\PY{+w}{ }\PY{n}{List}\PY{p}{.}\PY{n+na}{of}\PY{p}{(}\PY{k}{new}\PY{+w}{ }\PY{n}{Point}\PY{p}{(}\PY{p}{)}\PY{p}{)}\PY{p}{;}
\PY{c+c1}{// max(points); // compile error: Point is not a valid substitute for T extends Comparable\PYZlt{}T\PYZgt{}}
\end{Verbatim}
\end{codebox}

Why not just \code{\textless{}\allowbreak{}T~\allowbreak{}extends~\allowbreak{}Comparable\textgreater{}} (raw)? Because then \code{compareTo} would accept any \code{Object}, losing type safety. The recursive form \code{Comparable\textless{}\allowbreak{}T\textgreater{}} guarantees you can only compare \code{T} to other \code{T}. This exact signature is how \code{Collections.\allowbreak{}max} is declared in the JDK.

\subsection[Raw types and unchecked warnings]{Raw types and unchecked warnings}
\label{h:4:raw-types-and-unchecked-warnings}
A \textbf{raw type} is a generic class used without its type argument, e.g. \code{List} instead of \code{List\textless{}\allowbreak{}String\textgreater{}}. Raw types exist only for backward compatibility with pre-Java-5 code. Using them today throws away all compile-time type checking.

\begin{codebox}
\begin{Verbatim}[commandchars=\\\{\}]
\PY{n+nd}{@SuppressWarnings}\PY{p}{(}\PY{l+s}{\PYZdq{}}\PY{l+s}{unchecked}\PY{l+s}{\PYZdq{}}\PY{p}{)}
\PY{n}{List}\PY{+w}{ }\PY{n}{raw}\PY{+w}{ }\PY{o}{=}\PY{+w}{ }\PY{k}{new}\PY{+w}{ }\PY{n}{ArrayList}\PY{p}{(}\PY{p}{)}\PY{p}{;}\PY{+w}{     }\PY{c+c1}{// raw type — legal but discouraged}
\PY{n}{raw}\PY{p}{.}\PY{n+na}{add}\PY{p}{(}\PY{l+s}{\PYZdq{}}\PY{l+s}{a string}\PY{l+s}{\PYZdq{}}\PY{p}{)}\PY{p}{;}
\PY{n}{raw}\PY{p}{.}\PY{n+na}{add}\PY{p}{(}\PY{l+m+mi}{99}\PY{p}{)}\PY{p}{;}\PY{+w}{                    }\PY{c+c1}{// no error here, but it\PYZsq{}s a landmine}

\PY{n}{List}\PY{o}{\PYZlt{}}\PY{n}{String}\PY{o}{\PYZgt{}}\PY{+w}{ }\PY{n}{typed}\PY{+w}{ }\PY{o}{=}\PY{+w}{ }\PY{n}{raw}\PY{p}{;}\PY{+w}{       }\PY{c+c1}{// unchecked warning: unchecked conversion}
\PY{k}{for}\PY{+w}{ }\PY{p}{(}\PY{n}{String}\PY{+w}{ }\PY{n}{s}\PY{+w}{ }\PY{p}{:}\PY{+w}{ }\PY{n}{typed}\PY{p}{)}\PY{+w}{ }\PY{p}{\PYZob{}}\PY{+w}{        }\PY{c+c1}{// compiles, but...}
\PY{+w}{    }\PY{n}{System}\PY{p}{.}\PY{n+na}{out}\PY{p}{.}\PY{n+na}{println}\PY{p}{(}\PY{n}{s}\PY{p}{.}\PY{n+na}{length}\PY{p}{(}\PY{p}{)}\PY{p}{)}\PY{p}{;}
\PY{p}{\PYZcb{}}
\PY{c+c1}{// Throws java.lang.ClassCastException at the Integer 99,}
\PY{c+c1}{// because the compiler inserts a hidden (String) cast on read.}
\end{Verbatim}
\end{codebox}

The compiler warning here is \code{unchecked~\allowbreak{}conversion} — it means “I cannot verify this is actually a \code{List\textless{}\allowbreak{}String\textgreater{}}; trust me.” Never silence these warnings without understanding why they are safe (see \code{@\allowbreak{}SafeVarargs} below for a legitimate case).

\subsection[Generic type inference and the diamond operator]{Generic type inference and the diamond operator}
\label{h:4:generic-type-inference-and-the-diamond-operator}
Before Java 7, you had to repeat the type argument on both sides:

\begin{codebox}
\begin{Verbatim}[commandchars=\\\{\}]
\PY{n}{Map}\PY{o}{\PYZlt{}}\PY{n}{String}\PY{p}{,}\PY{+w}{ }\PY{n}{List}\PY{o}{\PYZlt{}}\PY{n}{Integer}\PY{o}{\PYZgt{}}\PY{o}{\PYZgt{}}\PY{+w}{ }\PY{n}{map}\PY{+w}{ }\PY{o}{=}\PY{+w}{ }\PY{k}{new}\PY{+w}{ }\PY{n}{HashMap}\PY{o}{\PYZlt{}}\PY{n}{String}\PY{p}{,}\PY{+w}{ }\PY{n}{List}\PY{o}{\PYZlt{}}\PY{n}{Integer}\PY{o}{\PYZgt{}}\PY{o}{\PYZgt{}}\PY{p}{(}\PY{p}{)}\PY{p}{;}\PY{+w}{ }\PY{c+c1}{// verbose, pre\PYZhy{}Java\PYZhy{}7}
\end{Verbatim}
\end{codebox}

The \textbf{diamond operator} \code{\textless{}\textgreater{}} (Java 7+) tells the compiler to infer the type argument from the left-hand side (target type):

\begin{codebox}
\begin{Verbatim}[commandchars=\\\{\}]
\PY{n}{Map}\PY{o}{\PYZlt{}}\PY{n}{String}\PY{p}{,}\PY{+w}{ }\PY{n}{List}\PY{o}{\PYZlt{}}\PY{n}{Integer}\PY{o}{\PYZgt{}}\PY{o}{\PYZgt{}}\PY{+w}{ }\PY{n}{map}\PY{+w}{ }\PY{o}{=}\PY{+w}{ }\PY{k}{new}\PY{+w}{ }\PY{n}{HashMap}\PY{o}{\PYZlt{}}\PY{o}{\PYZgt{}}\PY{p}{(}\PY{p}{)}\PY{p}{;}\PY{+w}{ }\PY{c+c1}{// compiler infers \PYZlt{}String, List\PYZlt{}Integer\PYZgt{}\PYZgt{}}
\end{Verbatim}
\end{codebox}

Inference also works across method calls and generic methods, and Java 10+ can combine it with \code{var} for local variables (see below). Inference has limits — it works left-to-right and top-down, so it cannot look \emph{forward} into how you use the result:

\begin{codebox}
\begin{Verbatim}[commandchars=\\\{\}]
\PY{c+c1}{// Compiler cannot infer T just from an empty list with no target type}
\PY{k+kd}{var}\PY{+w}{ }\PY{n}{list}\PY{+w}{ }\PY{o}{=}\PY{+w}{ }\PY{n}{List}\PY{p}{.}\PY{n+na}{of}\PY{p}{(}\PY{p}{)}\PY{p}{;}\PY{+w}{ }\PY{c+c1}{// inferred as List\PYZlt{}Object\PYZgt{}, probably not what you wanted}

\PY{n}{List}\PY{o}{\PYZlt{}}\PY{n}{String}\PY{o}{\PYZgt{}}\PY{+w}{ }\PY{n}{strings}\PY{+w}{ }\PY{o}{=}\PY{+w}{ }\PY{n}{List}\PY{p}{.}\PY{n+na}{of}\PY{p}{(}\PY{p}{)}\PY{p}{;}\PY{+w}{ }\PY{c+c1}{// OK: target type String is known here}
\end{Verbatim}
\end{codebox}

\subsection[var with generics]{\code{var} with generics}
\label{h:4:var-with-generics}
\code{var} (Java 10+) infers the \emph{declared} type from the initializer — it does not create a raw or wildcard type, and it does not weaken generics. \code{var} is not “dynamic typing”; the type is fixed at compile time, just not written out.

\begin{codebox}
\begin{Verbatim}[commandchars=\\\{\}]
\PY{k+kd}{var}\PY{+w}{ }\PY{n}{names}\PY{+w}{ }\PY{o}{=}\PY{+w}{ }\PY{k}{new}\PY{+w}{ }\PY{n}{ArrayList}\PY{o}{\PYZlt{}}\PY{n}{String}\PY{o}{\PYZgt{}}\PY{p}{(}\PY{p}{)}\PY{p}{;}\PY{+w}{ }\PY{c+c1}{// names has static type ArrayList\PYZlt{}String\PYZgt{}}
\PY{n}{names}\PY{p}{.}\PY{n+na}{add}\PY{p}{(}\PY{l+s}{\PYZdq{}}\PY{l+s}{Linus}\PY{l+s}{\PYZdq{}}\PY{p}{)}\PY{p}{;}
\PY{c+c1}{// names.add(42); // compile error: incompatible types: int cannot be converted to String}
\end{Verbatim}
\end{codebox}

Caution: \code{var} needs a type on the right to infer from. You cannot use \code{var} with the diamond alone:

\begin{codebox}
\begin{Verbatim}[commandchars=\\\{\}]
\PY{c+c1}{// var list = new ArrayList\PYZlt{}\PYZgt{}(); // compiles, but infers ArrayList\PYZlt{}Object\PYZgt{} — loses precision}
\PY{k+kd}{var}\PY{+w}{ }\PY{n}{list}\PY{+w}{ }\PY{o}{=}\PY{+w}{ }\PY{k}{new}\PY{+w}{ }\PY{n}{ArrayList}\PY{o}{\PYZlt{}}\PY{n}{String}\PY{o}{\PYZgt{}}\PY{p}{(}\PY{p}{)}\PY{p}{;}\PY{+w}{ }\PY{c+c1}{// be explicit on one side}
\end{Verbatim}
\end{codebox}

Rule of thumb for code review: if \code{var} combined with \code{\textless{}\textgreater{}} makes the element type unclear at the call site, prefer spelling out the generic type.

\subsection[Generics and Optional / streams]{Generics and \code{Optional} / streams}
\label{h:4:generics-and-optional--streams}
\code{Optional\textless{}\allowbreak{}T\textgreater{}} and the Streams API (\code{Stream\textless{}\allowbreak{}T\textgreater{}}) are generic types that lean heavily on inference and bounded wildcards internally. They are covered fully in Chapter 5 and Chapter 8, but a couple of generics-specific points belong here.

\begin{codebox}
\begin{Verbatim}[commandchars=\\\{\}]
\PY{n}{Optional}\PY{o}{\PYZlt{}}\PY{n}{String}\PY{o}{\PYZgt{}}\PY{+w}{ }\PY{n}{maybeName}\PY{+w}{ }\PY{o}{=}\PY{+w}{ }\PY{n}{Optional}\PY{p}{.}\PY{n+na}{of}\PY{p}{(}\PY{l+s}{\PYZdq{}}\PY{l+s}{Ada}\PY{l+s}{\PYZdq{}}\PY{p}{)}\PY{p}{;}

\PY{c+c1}{// map() is a generic method: \PYZlt{}U\PYZgt{} Optional\PYZlt{}U\PYZgt{} map(Function\PYZlt{}? super T, ? extends U\PYZgt{} mapper)}
\PY{n}{Optional}\PY{o}{\PYZlt{}}\PY{n}{Integer}\PY{o}{\PYZgt{}}\PY{+w}{ }\PY{n}{length}\PY{+w}{ }\PY{o}{=}\PY{+w}{ }\PY{n}{maybeName}\PY{p}{.}\PY{n+na}{map}\PY{p}{(}\PY{n}{String}\PY{p}{:}\PY{p}{:}\PY{n}{length}\PY{p}{)}\PY{p}{;}

\PY{n}{List}\PY{o}{\PYZlt{}}\PY{n}{String}\PY{o}{\PYZgt{}}\PY{+w}{ }\PY{n}{names}\PY{+w}{ }\PY{o}{=}\PY{+w}{ }\PY{n}{List}\PY{p}{.}\PY{n+na}{of}\PY{p}{(}\PY{l+s}{\PYZdq{}}\PY{l+s}{Ada}\PY{l+s}{\PYZdq{}}\PY{p}{,}\PY{+w}{ }\PY{l+s}{\PYZdq{}}\PY{l+s}{Grace}\PY{l+s}{\PYZdq{}}\PY{p}{,}\PY{+w}{ }\PY{l+s}{\PYZdq{}}\PY{l+s}{Linus}\PY{l+s}{\PYZdq{}}\PY{p}{)}\PY{p}{;}
\PY{n}{List}\PY{o}{\PYZlt{}}\PY{n}{Integer}\PY{o}{\PYZgt{}}\PY{+w}{ }\PY{n}{lengths}\PY{+w}{ }\PY{o}{=}\PY{+w}{ }\PY{n}{names}\PY{p}{.}\PY{n+na}{stream}\PY{p}{(}\PY{p}{)}
\PY{+w}{        }\PY{p}{.}\PY{n+na}{map}\PY{p}{(}\PY{n}{String}\PY{p}{:}\PY{p}{:}\PY{n}{length}\PY{p}{)}\PY{+w}{   }\PY{c+c1}{// Stream\PYZlt{}String\PYZgt{} \PYZhy{}\PYZgt{} Stream\PYZlt{}Integer\PYZgt{}, T inferred per stage}
\PY{+w}{        }\PY{p}{.}\PY{n+na}{toList}\PY{p}{(}\PY{p}{)}\PY{p}{;}\PY{+w}{              }\PY{c+c1}{// Java 16+ shorthand for collect(Collectors.toList())}
\end{Verbatim}
\end{codebox}

Notice \code{Function\textless{}?\allowbreak{}~\allowbreak{}super~\allowbreak{}T,\allowbreak{}~\allowbreak{}?\allowbreak{}~\allowbreak{}extends~\allowbreak{}U\textgreater{}} in \code{Optional.\allowbreak{}map}’s real signature — that is the PECS wildcard pattern explained later in this chapter. Streams use the same pattern throughout (\code{Collector\textless{}?\allowbreak{}~\allowbreak{}super~\allowbreak{}T,\allowbreak{}~\allowbreak{}A,\allowbreak{}~\allowbreak{}R\textgreater{}}, etc.) so that a pipeline can accept producers and consumers of \emph{related} types, not just exact matches.

\section[Type Erasure]{Type Erasure}
\label{h:4:type-erasure}
\textbf{Type erasure} is how the Java compiler implements generics: type parameters exist only in source code and are checked at compile time, then \emph{erased} (removed) before bytecode is generated. At runtime, there is no \code{List\textless{}\allowbreak{}String\textgreater{}} — there is only \code{List}, with a compiler-inserted cast wherever you read an element.

Erasure exists for one reason: \textbf{backward compatibility}. Java 5 needed generic collections to run on a JVM whose bytecode format did not change, and to interoperate with pre-generics \code{.\allowbreak{}class} files already in production. Erasure made that possible at the cost of some runtime type information.

\subsection[What erasure does at the bytecode level]{What erasure does at the bytecode level}
\label{h:4:what-erasure-does-at-the-bytecode-level}
Given this source:

\begin{codebox}
\begin{Verbatim}[commandchars=\\\{\}]
\PY{k+kd}{public}\PY{+w}{ }\PY{k+kd}{class} \PY{n+nc}{Box}\PY{o}{\PYZlt{}}\PY{n}{T}\PY{o}{\PYZgt{}}\PY{+w}{ }\PY{p}{\PYZob{}}
\PY{+w}{    }\PY{k+kd}{private}\PY{+w}{ }\PY{n}{T}\PY{+w}{ }\PY{n}{value}\PY{p}{;}

\PY{+w}{    }\PY{k+kd}{public}\PY{+w}{ }\PY{k+kt}{void}\PY{+w}{ }\PY{n+nf}{set}\PY{p}{(}\PY{n}{T}\PY{+w}{ }\PY{n}{value}\PY{p}{)}\PY{+w}{ }\PY{p}{\PYZob{}}
\PY{+w}{        }\PY{k}{this}\PY{p}{.}\PY{n+na}{value}\PY{+w}{ }\PY{o}{=}\PY{+w}{ }\PY{n}{value}\PY{p}{;}
\PY{+w}{    }\PY{p}{\PYZcb{}}

\PY{+w}{    }\PY{k+kd}{public}\PY{+w}{ }\PY{n}{T}\PY{+w}{ }\PY{n+nf}{get}\PY{p}{(}\PY{p}{)}\PY{+w}{ }\PY{p}{\PYZob{}}
\PY{+w}{        }\PY{k}{return}\PY{+w}{ }\PY{n}{value}\PY{p}{;}
\PY{+w}{    }\PY{p}{\PYZcb{}}
\PY{p}{\PYZcb{}}
\end{Verbatim}
\end{codebox}

The compiler erases \code{T} to its bound — \code{Object}, since \code{T} has no bound here — and produces bytecode roughly equivalent to:

\begin{codebox}
\begin{Verbatim}[commandchars=\\\{\}]
\PY{c+c1}{// What the erased bytecode is equivalent to (not real source you write)}
\PY{k+kd}{public}\PY{+w}{ }\PY{k+kd}{class} \PY{n+nc}{Box}\PY{+w}{ }\PY{p}{\PYZob{}}
\PY{+w}{    }\PY{k+kd}{private}\PY{+w}{ }\PY{n}{Object}\PY{+w}{ }\PY{n}{value}\PY{p}{;}

\PY{+w}{    }\PY{k+kd}{public}\PY{+w}{ }\PY{k+kt}{void}\PY{+w}{ }\PY{n+nf}{set}\PY{p}{(}\PY{n}{Object}\PY{+w}{ }\PY{n}{value}\PY{p}{)}\PY{+w}{ }\PY{p}{\PYZob{}}
\PY{+w}{        }\PY{k}{this}\PY{p}{.}\PY{n+na}{value}\PY{+w}{ }\PY{o}{=}\PY{+w}{ }\PY{n}{value}\PY{p}{;}
\PY{+w}{    }\PY{p}{\PYZcb{}}

\PY{+w}{    }\PY{k+kd}{public}\PY{+w}{ }\PY{n}{Object}\PY{+w}{ }\PY{n+nf}{get}\PY{p}{(}\PY{p}{)}\PY{+w}{ }\PY{p}{\PYZob{}}
\PY{+w}{        }\PY{k}{return}\PY{+w}{ }\PY{n}{value}\PY{p}{;}
\PY{+w}{    }\PY{p}{\PYZcb{}}
\PY{p}{\PYZcb{}}
\end{Verbatim}
\end{codebox}

Anywhere you call \code{box.\allowbreak{}get()} from code that expects \code{String}, the compiler inserts a cast:

\begin{codebox}
\begin{Verbatim}[commandchars=\\\{\}]
\PY{n}{Box}\PY{o}{\PYZlt{}}\PY{n}{String}\PY{o}{\PYZgt{}}\PY{+w}{ }\PY{n}{box}\PY{+w}{ }\PY{o}{=}\PY{+w}{ }\PY{k}{new}\PY{+w}{ }\PY{n}{Box}\PY{o}{\PYZlt{}}\PY{o}{\PYZgt{}}\PY{p}{(}\PY{p}{)}\PY{p}{;}
\PY{n}{box}\PY{p}{.}\PY{n+na}{set}\PY{p}{(}\PY{l+s}{\PYZdq{}}\PY{l+s}{hi}\PY{l+s}{\PYZdq{}}\PY{p}{)}\PY{p}{;}
\PY{n}{String}\PY{+w}{ }\PY{n}{s}\PY{+w}{ }\PY{o}{=}\PY{+w}{ }\PY{n}{box}\PY{p}{.}\PY{n+na}{get}\PY{p}{(}\PY{p}{)}\PY{p}{;}
\PY{c+c1}{// Compiled roughly as:}
\PY{c+c1}{// String s = (String) box.get();}
\end{Verbatim}
\end{codebox}

You can verify this yourself with \code{javap~\allowbreak{}-\allowbreak{}c}:

\begin{codebox}
\begin{Verbatim}[commandchars=\\\{\}]
javac\PY{+w}{ }Box.java
javap\PY{+w}{ }\PYZhy{}c\PY{+w}{ }Box.class
\PY{c+c1}{\PYZsh{} Look at the set(Ljava/lang/Object;)V and get()Ljava/lang/Object; signatures —}
\PY{c+c1}{\PYZsh{} there is no T anywhere in the compiled descriptor.}
\end{Verbatim}
\end{codebox}

If \code{T} is bounded, e.g. \code{\textless{}\allowbreak{}T~\allowbreak{}extends~\allowbreak{}Number\textgreater{}}, erasure uses the bound instead of \code{Object}:

\begin{codebox}
\begin{Verbatim}[commandchars=\\\{\}]
\PY{k+kd}{public}\PY{+w}{ }\PY{k+kd}{class} \PY{n+nc}{NumberBox}\PY{o}{\PYZlt{}}\PY{n}{T}\PY{+w}{ }\PY{k+kd}{extends}\PY{+w}{ }\PY{n}{Number}\PY{o}{\PYZgt{}}\PY{+w}{ }\PY{p}{\PYZob{}}
\PY{+w}{    }\PY{k+kd}{private}\PY{+w}{ }\PY{n}{T}\PY{+w}{ }\PY{n}{value}\PY{p}{;}
\PY{+w}{    }\PY{k+kd}{public}\PY{+w}{ }\PY{n}{T}\PY{+w}{ }\PY{n+nf}{get}\PY{p}{(}\PY{p}{)}\PY{+w}{ }\PY{p}{\PYZob{}}\PY{+w}{ }\PY{k}{return}\PY{+w}{ }\PY{n}{value}\PY{p}{;}\PY{+w}{ }\PY{p}{\PYZcb{}}
\PY{p}{\PYZcb{}}
\PY{c+c1}{// Erased: private Number value; public Number get() \PYZob{} return value; \PYZcb{}}
\end{Verbatim}
\end{codebox}

\subsection[Bridge methods]{Bridge methods}
\label{h:4:bridge-methods}
Erasure creates a subtle problem with overriding. Consider a generic interface and a class that implements it with a concrete type:

\begin{codebox}
\begin{Verbatim}[commandchars=\\\{\}]
\PY{k+kd}{interface} \PY{n+nc}{Container}\PY{o}{\PYZlt{}}\PY{n}{T}\PY{o}{\PYZgt{}}\PY{+w}{ }\PY{p}{\PYZob{}}
\PY{+w}{    }\PY{k+kt}{void}\PY{+w}{ }\PY{n+nf}{set}\PY{p}{(}\PY{n}{T}\PY{+w}{ }\PY{n}{value}\PY{p}{)}\PY{p}{;}
\PY{p}{\PYZcb{}}

\PY{k+kd}{class} \PY{n+nc}{StringContainer}\PY{+w}{ }\PY{k+kd}{implements}\PY{+w}{ }\PY{n}{Container}\PY{o}{\PYZlt{}}\PY{n}{String}\PY{o}{\PYZgt{}}\PY{+w}{ }\PY{p}{\PYZob{}}
\PY{+w}{    }\PY{n+nd}{@Override}
\PY{+w}{    }\PY{k+kd}{public}\PY{+w}{ }\PY{k+kt}{void}\PY{+w}{ }\PY{n+nf}{set}\PY{p}{(}\PY{n}{String}\PY{+w}{ }\PY{n}{value}\PY{p}{)}\PY{+w}{ }\PY{p}{\PYZob{}}
\PY{+w}{        }\PY{n}{System}\PY{p}{.}\PY{n+na}{out}\PY{p}{.}\PY{n+na}{println}\PY{p}{(}\PY{l+s}{\PYZdq{}}\PY{l+s}{Set: }\PY{l+s}{\PYZdq{}}\PY{+w}{ }\PY{o}{+}\PY{+w}{ }\PY{n}{value}\PY{p}{)}\PY{p}{;}
\PY{+w}{    }\PY{p}{\PYZcb{}}
\PY{p}{\PYZcb{}}
\end{Verbatim}
\end{codebox}

After erasure, \code{Container\textless{}\allowbreak{}T\textgreater{}.\allowbreak{}set} becomes \code{set(\allowbreak{}Object)}. But \code{String\allowbreak{}Container.\allowbreak{}set(\allowbreak{}String)} does not have the same erased signature as \code{set(\allowbreak{}Object)} — so it would not actually override the interface method at the bytecode level, breaking polymorphism.

The compiler fixes this by generating a hidden \textbf{bridge method}: a synthetic \code{set(\allowbreak{}Object)} method in \code{StringContainer} that casts and delegates to \code{set(\allowbreak{}String)}.

\begin{codebox}
\begin{Verbatim}[commandchars=\\\{\}]
\PY{c+c1}{// What the compiler generates behind the scenes (not something you write)}
\PY{k+kd}{class} \PY{n+nc}{StringContainer}\PY{+w}{ }\PY{k+kd}{implements}\PY{+w}{ }\PY{n}{Container}\PY{o}{\PYZlt{}}\PY{n}{String}\PY{o}{\PYZgt{}}\PY{+w}{ }\PY{p}{\PYZob{}}
\PY{+w}{    }\PY{k+kd}{public}\PY{+w}{ }\PY{k+kt}{void}\PY{+w}{ }\PY{n+nf}{set}\PY{p}{(}\PY{n}{String}\PY{+w}{ }\PY{n}{value}\PY{p}{)}\PY{+w}{ }\PY{p}{\PYZob{}}\PY{+w}{ }\PY{c+cm}{/* real body */}\PY{+w}{ }\PY{p}{\PYZcb{}}

\PY{+w}{    }\PY{c+c1}{// synthetic bridge method, added automatically}
\PY{+w}{    }\PY{k+kd}{public}\PY{+w}{ }\PY{k+kt}{void}\PY{+w}{ }\PY{n+nf}{set}\PY{p}{(}\PY{n}{Object}\PY{+w}{ }\PY{n}{value}\PY{p}{)}\PY{+w}{ }\PY{p}{\PYZob{}}
\PY{+w}{        }\PY{n}{set}\PY{p}{(}\PY{p}{(}\PY{n}{String}\PY{p}{)}\PY{+w}{ }\PY{n}{value}\PY{p}{)}\PY{p}{;}\PY{+w}{ }\PY{c+c1}{// cast, then delegate}
\PY{+w}{    }\PY{p}{\PYZcb{}}
\PY{p}{\PYZcb{}}
\end{Verbatim}
\end{codebox}

You can see this with reflection — \code{String\allowbreak{}Container.\allowbreak{}class.\allowbreak{}get\allowbreak{}Declared\allowbreak{}Methods()} lists two \code{set} methods, and \code{Method.\allowbreak{}isBridge()} is \code{true} for the synthetic one. This is a classic “gotcha” interview question: “why does my class have two \code{set} methods when I only wrote one?”

\subsection[Why you cannot do new T[]]{Why you cannot do \code{new~\allowbreak{}T[]}}
\label{h:4:why-you-cannot-do-new-t}
You cannot create an array of a generic type parameter directly, because array creation needs a real, runtime-known component type — and after erasure, \code{T} is gone.

\begin{codebox}
\begin{Verbatim}[commandchars=\\\{\}]
\PY{k+kd}{public}\PY{+w}{ }\PY{k+kd}{class} \PY{n+nc}{Stack}\PY{o}{\PYZlt{}}\PY{n}{T}\PY{o}{\PYZgt{}}\PY{+w}{ }\PY{p}{\PYZob{}}
\PY{+w}{    }\PY{c+c1}{// private T[] items = new T[10];}
\PY{+w}{    }\PY{c+c1}{// compile error: generic array creation}

\PY{+w}{    }\PY{n+nd}{@SuppressWarnings}\PY{p}{(}\PY{l+s}{\PYZdq{}}\PY{l+s}{unchecked}\PY{l+s}{\PYZdq{}}\PY{p}{)}
\PY{+w}{    }\PY{k+kd}{private}\PY{+w}{ }\PY{n}{T}\PY{o}{[}\PY{o}{]}\PY{+w}{ }\PY{n}{items}\PY{+w}{ }\PY{o}{=}\PY{+w}{ }\PY{p}{(}\PY{n}{T}\PY{o}{[}\PY{o}{]}\PY{p}{)}\PY{+w}{ }\PY{k}{new}\PY{+w}{ }\PY{n}{Object}\PY{o}{[}\PY{l+m+mi}{10}\PY{o}{]}\PY{p}{;}\PY{+w}{ }\PY{c+c1}{// common workaround, still an unchecked cast}
\PY{p}{\PYZcb{}}
\end{Verbatim}
\end{codebox}

The workaround (\code{(\allowbreak{}T[])\allowbreak{}~\allowbreak{}new~\allowbreak{}Object[\allowbreak{}10]}) compiles but is unsafe: the actual runtime array is \code{Object[]}, not \code{T[]}. If you ever leak that array reference out and someone assigns it to a more specific array variable, you get a \code{ClassCastException} later. The safer alternative, when you truly need a typed array, is to require a \code{Class\textless{}\allowbreak{}T\textgreater{}} token from the caller and use \code{java.\allowbreak{}util.\allowbreak{}reflect.\allowbreak{}Array.\allowbreak{}new\allowbreak{}Instance}:

\begin{codebox}
\begin{Verbatim}[commandchars=\\\{\}]
\PY{n+nd}{@SuppressWarnings}\PY{p}{(}\PY{l+s}{\PYZdq{}}\PY{l+s}{unchecked}\PY{l+s}{\PYZdq{}}\PY{p}{)}
\PY{k+kd}{public}\PY{+w}{ }\PY{k+kd}{static}\PY{+w}{ }\PY{o}{\PYZlt{}}\PY{n}{T}\PY{o}{\PYZgt{}}\PY{+w}{ }\PY{n}{T}\PY{o}{[}\PY{o}{]}\PY{+w}{ }\PY{n+nf}{newArray}\PY{p}{(}\PY{n}{Class}\PY{o}{\PYZlt{}}\PY{n}{T}\PY{o}{\PYZgt{}}\PY{+w}{ }\PY{n}{type}\PY{p}{,}\PY{+w}{ }\PY{k+kt}{int}\PY{+w}{ }\PY{n}{size}\PY{p}{)}\PY{+w}{ }\PY{p}{\PYZob{}}
\PY{+w}{    }\PY{k}{return}\PY{+w}{ }\PY{p}{(}\PY{n}{T}\PY{o}{[}\PY{o}{]}\PY{p}{)}\PY{+w}{ }\PY{n}{java}\PY{p}{.}\PY{n+na}{lang}\PY{p}{.}\PY{n+na}{reflect}\PY{p}{.}\PY{n+na}{Array}\PY{p}{.}\PY{n+na}{newInstance}\PY{p}{(}\PY{n}{type}\PY{p}{,}\PY{+w}{ }\PY{n}{size}\PY{p}{)}\PY{p}{;}
\PY{p}{\PYZcb{}}

\PY{n}{String}\PY{o}{[}\PY{o}{]}\PY{+w}{ }\PY{n}{strings}\PY{+w}{ }\PY{o}{=}\PY{+w}{ }\PY{n}{newArray}\PY{p}{(}\PY{n}{String}\PY{p}{.}\PY{n+na}{class}\PY{p}{,}\PY{+w}{ }\PY{l+m+mi}{5}\PY{p}{)}\PY{p}{;}\PY{+w}{ }\PY{c+c1}{// works, and is actually a String[] at runtime}
\end{Verbatim}
\end{codebox}

\subsection[Why you cannot do instanceof List\textless{}String\textgreater{}]{Why you cannot do \code{instanceof~\allowbreak{}List\textless{}\allowbreak{}String\textgreater{}}}
\label{h:4:why-you-cannot-do-instanceof-liststring}
\code{instanceof} is a runtime check, and after erasure there is no runtime distinction between \code{List\textless{}\allowbreak{}String\textgreater{}} and \code{List\textless{}\allowbreak{}Integer\textgreater{}} — both are just \code{List}.

\begin{codebox}
\begin{Verbatim}[commandchars=\\\{\}]
\PY{n}{List}\PY{o}{\PYZlt{}}\PY{n}{String}\PY{o}{\PYZgt{}}\PY{+w}{ }\PY{n}{list}\PY{+w}{ }\PY{o}{=}\PY{+w}{ }\PY{k}{new}\PY{+w}{ }\PY{n}{ArrayList}\PY{o}{\PYZlt{}}\PY{o}{\PYZgt{}}\PY{p}{(}\PY{p}{)}\PY{p}{;}

\PY{c+c1}{// if (list instanceof List\PYZlt{}String\PYZgt{}) \PYZob{} \PYZcb{} // compile error: illegal generic type for instanceof}

\PY{k}{if}\PY{+w}{ }\PY{p}{(}\PY{n}{list}\PY{+w}{ }\PY{k}{instanceof}\PY{+w}{ }\PY{n}{List}\PY{o}{\PYZlt{}}\PY{o}{?}\PY{o}{\PYZgt{}}\PY{p}{)}\PY{+w}{ }\PY{p}{\PYZob{}}\PY{+w}{ }\PY{c+c1}{// OK — unbounded wildcard is allowed}
\PY{+w}{    }\PY{n}{System}\PY{p}{.}\PY{n+na}{out}\PY{p}{.}\PY{n+na}{println}\PY{p}{(}\PY{l+s}{\PYZdq{}}\PY{l+s}{it\PYZsq{}s some kind of List}\PY{l+s}{\PYZdq{}}\PY{p}{)}\PY{p}{;}
\PY{p}{\PYZcb{}}

\PY{k}{if}\PY{+w}{ }\PY{p}{(}\PY{n}{list}\PY{+w}{ }\PY{k}{instanceof}\PY{+w}{ }\PY{n}{ArrayList}\PY{o}{\PYZlt{}}\PY{n}{String}\PY{o}{\PYZgt{}}\PY{+w}{ }\PY{n}{al}\PY{p}{)}\PY{+w}{ }\PY{p}{\PYZob{}}\PY{+w}{ }\PY{c+c1}{// still compile error, same reason}
\PY{+w}{    }\PY{c+c1}{// ...}
\PY{p}{\PYZcb{}}
\end{Verbatim}
\end{codebox}

Only the unbounded wildcard form (\code{List\textless{}?\textgreater{}}) or the raw form (\code{List}) is legal with \code{instanceof}, because those do not claim any specific type argument that would need a runtime check.

\subsection[Why you cannot have static generic fields]{Why you cannot have static generic fields}
\label{h:4:why-you-cannot-have-static-generic-fields}
A \code{static} field belongs to the class itself, and a class exists once at runtime, regardless of how many different type arguments (\code{Box\textless{}\allowbreak{}String\textgreater{}}, \code{Box\textless{}\allowbreak{}Integer\textgreater{}}, …) are used in source code. Since erasure removes \code{T} from the class entirely, there is nothing for a static field of type \code{T} to mean.

\begin{codebox}
\begin{Verbatim}[commandchars=\\\{\}]
\PY{k+kd}{public}\PY{+w}{ }\PY{k+kd}{class} \PY{n+nc}{Box}\PY{o}{\PYZlt{}}\PY{n}{T}\PY{o}{\PYZgt{}}\PY{+w}{ }\PY{p}{\PYZob{}}
\PY{+w}{    }\PY{c+c1}{// private static T defaultValue; // compile error: non\PYZhy{}static type variable T cannot be referenced from a static context}

\PY{+w}{    }\PY{k+kd}{private}\PY{+w}{ }\PY{k+kd}{static}\PY{+w}{ }\PY{n}{Object}\PY{+w}{ }\PY{n}{defaultValue}\PY{p}{;}\PY{+w}{ }\PY{c+c1}{// fine — just use Object, or a concrete type}
\PY{p}{\PYZcb{}}
\end{Verbatim}
\end{codebox}

Static generic \emph{methods} are fine (\code{static~\allowbreak{}\textless{}\allowbreak{}T\textgreater{}\allowbreak{}~\allowbreak{}T~\allowbreak{}identity(\allowbreak{}T~\allowbreak{}t)}), because the method itself introduces a fresh type parameter each time it is called — it does not depend on the class’s erased type parameter.

\subsection[Heap pollution and @SafeVarargs]{Heap pollution and \code{@\allowbreak{}SafeVarargs}}
\label{h:4:heap-pollution-and-safevarargs}
\textbf{Heap pollution} happens when a variable of a parameterized type refers to an object that is not actually of that parameterized type — usually caused by mixing raw types, unchecked casts, or varargs with generics.

\begin{codebox}
\begin{Verbatim}[commandchars=\\\{\}]
\PY{k+kd}{static}\PY{+w}{ }\PY{k+kt}{void}\PY{+w}{ }\PY{n+nf}{dangerous}\PY{p}{(}\PY{n}{List}\PY{o}{\PYZlt{}}\PY{n}{String}\PY{o}{\PYZgt{}}\PY{p}{.}\PY{p}{.}\PY{p}{.}\PY{+w}{ }\PY{n}{lists}\PY{p}{)}\PY{+w}{ }\PY{p}{\PYZob{}}\PY{+w}{ }\PY{c+c1}{// varargs of a generic type — allowed, but risky}
\PY{+w}{    }\PY{n}{Object}\PY{o}{[}\PY{o}{]}\PY{+w}{ }\PY{n}{array}\PY{+w}{ }\PY{o}{=}\PY{+w}{ }\PY{n}{lists}\PY{p}{;}\PY{+w}{       }\PY{c+c1}{// varargs is implemented as an array, arrays are covariant (see below)}
\PY{+w}{    }\PY{n}{array}\PY{o}{[}\PY{l+m+mi}{0}\PY{o}{]}\PY{+w}{ }\PY{o}{=}\PY{+w}{ }\PY{n}{List}\PY{p}{.}\PY{n+na}{of}\PY{p}{(}\PY{l+m+mi}{42}\PY{p}{)}\PY{p}{;}\PY{+w}{       }\PY{c+c1}{// legal at the array level: List\PYZlt{}Integer\PYZgt{} stored where List\PYZlt{}String\PYZgt{} expected}
\PY{+w}{    }\PY{n}{String}\PY{+w}{ }\PY{n}{s}\PY{+w}{ }\PY{o}{=}\PY{+w}{ }\PY{n}{lists}\PY{o}{[}\PY{l+m+mi}{0}\PY{o}{]}\PY{p}{.}\PY{n+na}{get}\PY{p}{(}\PY{l+m+mi}{0}\PY{p}{)}\PY{p}{;}\PY{+w}{   }\PY{c+c1}{// compiles fine...}
\PY{+w}{    }\PY{c+c1}{// ...throws ClassCastException at runtime: Integer cannot be cast to String}
\PY{p}{\PYZcb{}}
\end{Verbatim}
\end{codebox}

This compiles only with an \code{unchecked} warning at the declaration site, because \code{T...} varargs create a hidden array of a generic type — exactly the “arrays are covariant, generics are not” mismatch. If you are certain your varargs method never stores anything unsafe into that array (i.e., it only \emph{reads} from it), you can suppress the warning with \code{@\allowbreak{}SafeVarargs}:

\begin{codebox}
\begin{Verbatim}[commandchars=\\\{\}]
\PY{n+nd}{@SafeVarargs}\PY{+w}{ }\PY{c+c1}{// a promise to the compiler AND to readers: this method is safe}
\PY{k+kd}{static}\PY{+w}{ }\PY{o}{\PYZlt{}}\PY{n}{T}\PY{o}{\PYZgt{}}\PY{+w}{ }\PY{n}{List}\PY{o}{\PYZlt{}}\PY{n}{T}\PY{o}{\PYZgt{}}\PY{+w}{ }\PY{n+nf}{listOf}\PY{p}{(}\PY{n}{T}\PY{p}{.}\PY{p}{.}\PY{p}{.}\PY{+w}{ }\PY{n}{items}\PY{p}{)}\PY{+w}{ }\PY{p}{\PYZob{}}
\PY{+w}{    }\PY{k}{return}\PY{+w}{ }\PY{n}{Arrays}\PY{p}{.}\PY{n+na}{asList}\PY{p}{(}\PY{n}{items}\PY{p}{)}\PY{p}{;}\PY{+w}{ }\PY{c+c1}{// only reads from the array, never writes into it}
\PY{p}{\PYZcb{}}
\end{Verbatim}
\end{codebox}

\code{@\allowbreak{}SafeVarargs} is only allowed on \code{static} or \code{final} methods, \code{private} instance methods (Java 9+), and constructors — never on a regular overridable instance method, because a subclass could override it with unsafe behavior, breaking the promise.

\begin{codebox}
\begin{Verbatim}[commandchars=\\\{\}]
\PY{c+c1}{// @SafeVarargs}
\PY{c+c1}{// void notAllowed(List\PYZlt{}String\PYZgt{}... lists) \PYZob{} \PYZcb{} // compile error if the method is a non\PYZhy{}final, non\PYZhy{}private instance method}
\end{Verbatim}
\end{codebox}

\section[Wildcards (extends, super)]{Wildcards (extends, super)}
\label{h:4:wildcards-extends-super}
A \textbf{wildcard} (\code{?}) represents an unknown type argument. Wildcards exist to make generic APIs more flexible when you do not need to know or care about the exact type argument — only what you can do with it (read vs. write).

\subsection[Unbounded wildcards \textless{}?\textgreater{}]{Unbounded wildcards \code{\textless{}?\textgreater{}}}
\label{h:4:unbounded-wildcards-}
\code{\textless{}?\textgreater{}} means “a list of some type, I don’t know or care which.” Use it when your code only needs methods that do not depend on the type argument at all (e.g., \code{size()}, \code{clear()}).

\begin{codebox}
\begin{Verbatim}[commandchars=\\\{\}]
\PY{k+kd}{static}\PY{+w}{ }\PY{k+kt}{void}\PY{+w}{ }\PY{n+nf}{printSize}\PY{p}{(}\PY{n}{List}\PY{o}{\PYZlt{}}\PY{o}{?}\PY{o}{\PYZgt{}}\PY{+w}{ }\PY{n}{list}\PY{p}{)}\PY{+w}{ }\PY{p}{\PYZob{}}
\PY{+w}{    }\PY{n}{System}\PY{p}{.}\PY{n+na}{out}\PY{p}{.}\PY{n+na}{println}\PY{p}{(}\PY{l+s}{\PYZdq{}}\PY{l+s}{Size: }\PY{l+s}{\PYZdq{}}\PY{+w}{ }\PY{o}{+}\PY{+w}{ }\PY{n}{list}\PY{p}{.}\PY{n+na}{size}\PY{p}{(}\PY{p}{)}\PY{p}{)}\PY{p}{;}\PY{+w}{ }\PY{c+c1}{// fine, size() doesn\PYZsq{}t depend on T}
\PY{p}{\PYZcb{}}

\PY{k+kd}{static}\PY{+w}{ }\PY{k+kt}{void}\PY{+w}{ }\PY{n+nf}{printAll}\PY{p}{(}\PY{n}{List}\PY{o}{\PYZlt{}}\PY{o}{?}\PY{o}{\PYZgt{}}\PY{+w}{ }\PY{n}{list}\PY{p}{)}\PY{+w}{ }\PY{p}{\PYZob{}}
\PY{+w}{    }\PY{k}{for}\PY{+w}{ }\PY{p}{(}\PY{n}{Object}\PY{+w}{ }\PY{n}{o}\PY{+w}{ }\PY{p}{:}\PY{+w}{ }\PY{n}{list}\PY{p}{)}\PY{+w}{ }\PY{p}{\PYZob{}}\PY{+w}{   }\PY{c+c1}{// you can only read elements as Object}
\PY{+w}{        }\PY{n}{System}\PY{p}{.}\PY{n+na}{out}\PY{p}{.}\PY{n+na}{println}\PY{p}{(}\PY{n}{o}\PY{p}{)}\PY{p}{;}
\PY{+w}{    }\PY{p}{\PYZcb{}}
\PY{+w}{    }\PY{c+c1}{// list.add(\PYZdq{}x\PYZdq{}); // compile error: no way to prove String matches the unknown type}
\PY{p}{\PYZcb{}}
\end{Verbatim}
\end{codebox}

\subsection[Upper-bounded wildcards \textless{}? extends T\textgreater{}]{Upper-bounded wildcards \code{\textless{}?\allowbreak{}~\allowbreak{}extends~\allowbreak{}T\textgreater{}}}
\label{h:4:upper-bounded-wildcards--extends-t}
\code{\textless{}?\allowbreak{}~\allowbreak{}extends~\allowbreak{}Number\textgreater{}} means “a list of some unknown type that is \code{Number} or a subtype of \code{Number}.” You can safely \textbf{read} elements out as \code{Number}, but you cannot safely \textbf{add} to it (the compiler cannot know the exact subtype, so any insert could violate it).

\begin{codebox}
\begin{Verbatim}[commandchars=\\\{\}]
\PY{k+kd}{static}\PY{+w}{ }\PY{k+kt}{double}\PY{+w}{ }\PY{n+nf}{sumAll}\PY{p}{(}\PY{n}{List}\PY{o}{\PYZlt{}}\PY{o}{?}\PY{+w}{ }\PY{k+kd}{extends}\PY{+w}{ }\PY{n}{Number}\PY{o}{\PYZgt{}}\PY{+w}{ }\PY{n}{numbers}\PY{p}{)}\PY{+w}{ }\PY{p}{\PYZob{}}
\PY{+w}{    }\PY{k+kt}{double}\PY{+w}{ }\PY{n}{total}\PY{+w}{ }\PY{o}{=}\PY{+w}{ }\PY{l+m+mi}{0}\PY{p}{;}
\PY{+w}{    }\PY{k}{for}\PY{+w}{ }\PY{p}{(}\PY{n}{Number}\PY{+w}{ }\PY{n}{n}\PY{+w}{ }\PY{p}{:}\PY{+w}{ }\PY{n}{numbers}\PY{p}{)}\PY{+w}{ }\PY{p}{\PYZob{}}\PY{+w}{ }\PY{c+c1}{// reading is safe — every element IS\PYZhy{}A Number}
\PY{+w}{        }\PY{n}{total}\PY{+w}{ }\PY{o}{+}\PY{o}{=}\PY{+w}{ }\PY{n}{n}\PY{p}{.}\PY{n+na}{doubleValue}\PY{p}{(}\PY{p}{)}\PY{p}{;}
\PY{+w}{    }\PY{p}{\PYZcb{}}
\PY{+w}{    }\PY{c+c1}{// numbers.add(1);       // compile error: cannot add to a List\PYZlt{}? extends Number\PYZgt{}}
\PY{+w}{    }\PY{c+c1}{// numbers.add(1.0);     // compile error, same reason}
\PY{+w}{    }\PY{k}{return}\PY{+w}{ }\PY{n}{total}\PY{p}{;}
\PY{p}{\PYZcb{}}

\PY{n}{List}\PY{o}{\PYZlt{}}\PY{n}{Integer}\PY{o}{\PYZgt{}}\PY{+w}{ }\PY{n}{ints}\PY{+w}{ }\PY{o}{=}\PY{+w}{ }\PY{n}{List}\PY{p}{.}\PY{n+na}{of}\PY{p}{(}\PY{l+m+mi}{1}\PY{p}{,}\PY{+w}{ }\PY{l+m+mi}{2}\PY{p}{,}\PY{+w}{ }\PY{l+m+mi}{3}\PY{p}{)}\PY{p}{;}
\PY{n}{List}\PY{o}{\PYZlt{}}\PY{n}{Double}\PY{o}{\PYZgt{}}\PY{+w}{ }\PY{n}{doubles}\PY{+w}{ }\PY{o}{=}\PY{+w}{ }\PY{n}{List}\PY{p}{.}\PY{n+na}{of}\PY{p}{(}\PY{l+m+mf}{1.0}\PY{p}{,}\PY{+w}{ }\PY{l+m+mf}{2.0}\PY{p}{)}\PY{p}{;}
\PY{n}{sumAll}\PY{p}{(}\PY{n}{ints}\PY{p}{)}\PY{p}{;}\PY{+w}{    }\PY{c+c1}{// OK}
\PY{n}{sumAll}\PY{p}{(}\PY{n}{doubles}\PY{p}{)}\PY{p}{;}\PY{+w}{ }\PY{c+c1}{// OK — this flexibility is the whole point}
\end{Verbatim}
\end{codebox}

\subsection[Lower-bounded wildcards \textless{}? super T\textgreater{}]{Lower-bounded wildcards \code{\textless{}?\allowbreak{}~\allowbreak{}super~\allowbreak{}T\textgreater{}}}
\label{h:4:lower-bounded-wildcards--super-t}
\code{\textless{}?\allowbreak{}~\allowbreak{}super~\allowbreak{}Integer\textgreater{}} means “a list of some unknown type that is \code{Integer} or a supertype of \code{Integer}.” You can safely \textbf{write} \code{Integer}s into it, but reading only gives you \code{Object} (the compiler only knows the list holds \emph{at least} \code{Integer}-compatible items, so it cannot promise a more specific type back).

\begin{codebox}
\begin{Verbatim}[commandchars=\\\{\}]
\PY{k+kd}{static}\PY{+w}{ }\PY{k+kt}{void}\PY{+w}{ }\PY{n+nf}{addNumbers}\PY{p}{(}\PY{n}{List}\PY{o}{\PYZlt{}}\PY{o}{?}\PY{+w}{ }\PY{k+kd}{super}\PY{+w}{ }\PY{n}{Integer}\PY{o}{\PYZgt{}}\PY{+w}{ }\PY{n}{list}\PY{p}{)}\PY{+w}{ }\PY{p}{\PYZob{}}
\PY{+w}{    }\PY{n}{list}\PY{p}{.}\PY{n+na}{add}\PY{p}{(}\PY{l+m+mi}{1}\PY{p}{)}\PY{p}{;}\PY{+w}{   }\PY{c+c1}{// safe — any supertype of Integer can hold an Integer}
\PY{+w}{    }\PY{n}{list}\PY{p}{.}\PY{n+na}{add}\PY{p}{(}\PY{l+m+mi}{2}\PY{p}{)}\PY{p}{;}
\PY{+w}{    }\PY{c+c1}{// Integer first = list.get(0); // compile error: get() returns Object, not Integer}
\PY{+w}{    }\PY{n}{Object}\PY{+w}{ }\PY{n}{first}\PY{+w}{ }\PY{o}{=}\PY{+w}{ }\PY{n}{list}\PY{p}{.}\PY{n+na}{get}\PY{p}{(}\PY{l+m+mi}{0}\PY{p}{)}\PY{p}{;}\PY{+w}{     }\PY{c+c1}{// this is all you get back}
\PY{p}{\PYZcb{}}

\PY{n}{List}\PY{o}{\PYZlt{}}\PY{n}{Number}\PY{o}{\PYZgt{}}\PY{+w}{ }\PY{n}{numbers}\PY{+w}{ }\PY{o}{=}\PY{+w}{ }\PY{k}{new}\PY{+w}{ }\PY{n}{ArrayList}\PY{o}{\PYZlt{}}\PY{o}{\PYZgt{}}\PY{p}{(}\PY{p}{)}\PY{p}{;}
\PY{n}{List}\PY{o}{\PYZlt{}}\PY{n}{Object}\PY{o}{\PYZgt{}}\PY{+w}{ }\PY{n}{objects}\PY{+w}{ }\PY{o}{=}\PY{+w}{ }\PY{k}{new}\PY{+w}{ }\PY{n}{ArrayList}\PY{o}{\PYZlt{}}\PY{o}{\PYZgt{}}\PY{p}{(}\PY{p}{)}\PY{p}{;}
\PY{n}{addNumbers}\PY{p}{(}\PY{n}{numbers}\PY{p}{)}\PY{p}{;}\PY{+w}{ }\PY{c+c1}{// OK, Number is a supertype of Integer}
\PY{n}{addNumbers}\PY{p}{(}\PY{n}{objects}\PY{p}{)}\PY{p}{;}\PY{+w}{ }\PY{c+c1}{// OK, Object is a supertype of Integer}
\end{Verbatim}
\end{codebox}

\subsection[PECS: Producer Extends, Consumer Super]{PECS: Producer Extends, Consumer Super}
\label{h:4:pecs-producer-extends-consumer-super}
\textbf{PECS} is the mnemonic for choosing the right wildcard: if a parameterized type \emph{produces} (you read from it), use \code{extends}. If it \emph{consumes} (you write into it), use \code{super}. If it does both, use no wildcard at all (an exact type).

{\tablefont
\begin{xltabular}{\linewidth}{@{}L{0.659}L{1.224}L{0.659}L{1.459}@{}}
\toprule
\rowcolor{tablehdr}
\thd{Role} & \thd{Wildcard} & \thd{You can read} & \thd{You can write} \\
\midrule
\endhead
Producer only & \code{?\allowbreak{}~\allowbreak{}extends~\allowbreak{}T} & Yes, as \code{T} & No (compile error) \\
Consumer only & \code{?\allowbreak{}~\allowbreak{}super~\allowbreak{}T} & Only as \code{Object} & Yes, \code{T} and subtypes \\
Both & exact type \code{T} (no wildcard) & Yes, as \code{T} & Yes, as \code{T} \\
Neither (rare) & \code{?} (unbounded) & Only as \code{Object} & No (compile error, except \code{null}) \\
\bottomrule
\end{xltabular}
}

The JDK’s own \code{Collections.\allowbreak{}copy} is the textbook PECS example:

\begin{codebox}
\begin{Verbatim}[commandchars=\\\{\}]
\PY{c+c1}{// Real JDK signature:}
\PY{c+c1}{// public static \PYZlt{}T\PYZgt{} void copy(List\PYZlt{}? super T\PYZgt{} dest, List\PYZlt{}? extends T\PYZgt{} src)}
\PY{k+kd}{public}\PY{+w}{ }\PY{k+kd}{static}\PY{+w}{ }\PY{o}{\PYZlt{}}\PY{n}{T}\PY{o}{\PYZgt{}}\PY{+w}{ }\PY{k+kt}{void}\PY{+w}{ }\PY{n+nf}{copy}\PY{p}{(}\PY{n}{List}\PY{o}{\PYZlt{}}\PY{o}{?}\PY{+w}{ }\PY{k+kd}{super}\PY{+w}{ }\PY{n}{T}\PY{o}{\PYZgt{}}\PY{+w}{ }\PY{n}{dest}\PY{p}{,}\PY{+w}{ }\PY{n}{List}\PY{o}{\PYZlt{}}\PY{o}{?}\PY{+w}{ }\PY{k+kd}{extends}\PY{+w}{ }\PY{n}{T}\PY{o}{\PYZgt{}}\PY{+w}{ }\PY{n}{src}\PY{p}{)}\PY{+w}{ }\PY{p}{\PYZob{}}
\PY{+w}{    }\PY{k}{for}\PY{+w}{ }\PY{p}{(}\PY{k+kt}{int}\PY{+w}{ }\PY{n}{i}\PY{+w}{ }\PY{o}{=}\PY{+w}{ }\PY{l+m+mi}{0}\PY{p}{;}\PY{+w}{ }\PY{n}{i}\PY{+w}{ }\PY{o}{\PYZlt{}}\PY{+w}{ }\PY{n}{src}\PY{p}{.}\PY{n+na}{size}\PY{p}{(}\PY{p}{)}\PY{p}{;}\PY{+w}{ }\PY{n}{i}\PY{o}{+}\PY{o}{+}\PY{p}{)}\PY{+w}{ }\PY{p}{\PYZob{}}
\PY{+w}{        }\PY{n}{dest}\PY{p}{.}\PY{n+na}{set}\PY{p}{(}\PY{n}{i}\PY{p}{,}\PY{+w}{ }\PY{n}{src}\PY{p}{.}\PY{n+na}{get}\PY{p}{(}\PY{n}{i}\PY{p}{)}\PY{p}{)}\PY{p}{;}\PY{+w}{ }\PY{c+c1}{// read (produce) from src, write (consume) into dest}
\PY{+w}{    }\PY{p}{\PYZcb{}}
\PY{p}{\PYZcb{}}

\PY{n}{List}\PY{o}{\PYZlt{}}\PY{n}{Object}\PY{o}{\PYZgt{}}\PY{+w}{ }\PY{n}{dest}\PY{+w}{ }\PY{o}{=}\PY{+w}{ }\PY{k}{new}\PY{+w}{ }\PY{n}{ArrayList}\PY{o}{\PYZlt{}}\PY{o}{\PYZgt{}}\PY{p}{(}\PY{n}{List}\PY{p}{.}\PY{n+na}{of}\PY{p}{(}\PY{l+m+mi}{0}\PY{p}{,}\PY{+w}{ }\PY{l+m+mi}{0}\PY{p}{,}\PY{+w}{ }\PY{l+m+mi}{0}\PY{p}{)}\PY{p}{)}\PY{p}{;}
\PY{n}{List}\PY{o}{\PYZlt{}}\PY{n}{Integer}\PY{o}{\PYZgt{}}\PY{+w}{ }\PY{n}{src}\PY{+w}{ }\PY{o}{=}\PY{+w}{ }\PY{n}{List}\PY{p}{.}\PY{n+na}{of}\PY{p}{(}\PY{l+m+mi}{1}\PY{p}{,}\PY{+w}{ }\PY{l+m+mi}{2}\PY{p}{,}\PY{+w}{ }\PY{l+m+mi}{3}\PY{p}{)}\PY{p}{;}
\PY{n}{copy}\PY{p}{(}\PY{n}{dest}\PY{p}{,}\PY{+w}{ }\PY{n}{src}\PY{p}{)}\PY{p}{;}\PY{+w}{ }\PY{c+c1}{// src produces Integers, dest consumes them as Objects — PECS in action}
\end{Verbatim}
\end{codebox}

Without PECS, \code{copy} would need to be \code{copy(\allowbreak{}List\textless{}\allowbreak{}T\textgreater{}\allowbreak{}~\allowbreak{}dest,\allowbreak{}~\allowbreak{}List\textless{}\allowbreak{}T\textgreater{}\allowbreak{}~\allowbreak{}src)}, forcing both lists to hold the \emph{exact} same type — you could not copy \code{List\textless{}\allowbreak{}Integer\textgreater{}} into \code{List\textless{}\allowbreak{}Object\textgreater{}}, even though that is perfectly safe.

\subsection[List\textless{}?\textgreater{} vs List\textless{}Object\textgreater{} vs raw List]{\code{List\textless{}?\textgreater{}} vs \code{List\textless{}\allowbreak{}Object\textgreater{}} vs raw \code{List}}
\label{h:4:list-vs-listobject-vs-raw-list}
These three look similar but behave very differently. This is a favorite whiteboard question.

{\tablefont
\begin{xltabular}{\linewidth}{@{}L{0.492}L{1.475}L{0.820}L{0.820}L{1.393}@{}}
\toprule
\rowcolor{tablehdr}
\thd{Expression} & \thd{Meaning} & \thd{Can you \code{add(\allowbreak{}String)}?} & \thd{Can you \code{add(\allowbreak{}Object)}?} & \thd{Can you assign \code{List\textless{}\allowbreak{}String\textgreater{}} to it?} \\
\midrule
\endhead
\code{List\textless{}?\textgreater{}} & list of \emph{some} unknown, fixed type & No & No & Yes \\
\code{List\textless{}\allowbreak{}Object\textgreater{}} & list whose element type IS \code{Object} & Yes & Yes & No (compile error) \\
\code{List} (raw) & no compile-time type checking at all & Yes (unchecked) & Yes (unchecked) & Yes (with unchecked warning) \\
\bottomrule
\end{xltabular}
}

\begin{codebox}
\begin{Verbatim}[commandchars=\\\{\}]
\PY{n}{List}\PY{o}{\PYZlt{}}\PY{n}{String}\PY{o}{\PYZgt{}}\PY{+w}{ }\PY{n}{strings}\PY{+w}{ }\PY{o}{=}\PY{+w}{ }\PY{k}{new}\PY{+w}{ }\PY{n}{ArrayList}\PY{o}{\PYZlt{}}\PY{o}{\PYZgt{}}\PY{p}{(}\PY{n}{List}\PY{p}{.}\PY{n+na}{of}\PY{p}{(}\PY{l+s}{\PYZdq{}}\PY{l+s}{a}\PY{l+s}{\PYZdq{}}\PY{p}{,}\PY{+w}{ }\PY{l+s}{\PYZdq{}}\PY{l+s}{b}\PY{l+s}{\PYZdq{}}\PY{p}{)}\PY{p}{)}\PY{p}{;}

\PY{n}{List}\PY{o}{\PYZlt{}}\PY{o}{?}\PY{o}{\PYZgt{}}\PY{+w}{ }\PY{n}{wildcard}\PY{+w}{ }\PY{o}{=}\PY{+w}{ }\PY{n}{strings}\PY{p}{;}\PY{+w}{      }\PY{c+c1}{// OK, ? matches any type argument}
\PY{c+c1}{// wildcard.add(\PYZdq{}c\PYZdq{});             // compile error: unknown type, can\PYZsq{}t prove safety}

\PY{n}{List}\PY{o}{\PYZlt{}}\PY{n}{Object}\PY{o}{\PYZgt{}}\PY{+w}{ }\PY{n}{objects}\PY{+w}{ }\PY{o}{=}\PY{+w}{ }\PY{n}{strings}\PY{p}{;}\PY{+w}{  }\PY{c+c1}{// compile error: incompatible types}
\PY{c+c1}{// The compiler will not let this compile:}
\PY{c+c1}{// error: incompatible types: List\PYZlt{}String\PYZgt{} cannot be converted to List\PYZlt{}Object\PYZgt{}}

\PY{n}{List}\PY{+w}{ }\PY{n}{raw}\PY{+w}{ }\PY{o}{=}\PY{+w}{ }\PY{n}{strings}\PY{p}{;}\PY{+w}{              }\PY{c+c1}{// compiles, with an \PYZdq{}unchecked\PYZdq{} nature — no error, no warning here}
\PY{n}{raw}\PY{p}{.}\PY{n+na}{add}\PY{p}{(}\PY{l+m+mi}{42}\PY{p}{)}\PY{p}{;}\PY{+w}{                     }\PY{c+c1}{// compiles! but corrupts the underlying List\PYZlt{}String\PYZgt{} at runtime}
\PY{c+c1}{// strings.get(2); // throws ClassCastException — the \PYZdq{}42\PYZdq{} is really an Integer, cast to String fails}
\end{Verbatim}
\end{codebox}

That last example is exactly why raw types are dangerous: they let you silently break the type-safety promise that generics exist to enforce.

\section[Covariance and Contravariance]{Covariance and Contravariance}
\label{h:4:covariance-and-contravariance}
\textbf{Covariance} means “a subtype relationship is preserved”: if \code{Cat} is a subtype of \code{Animal}, then \code{Cat[]} is treated as a subtype of \code{Animal[]}. \textbf{Contravariance} is the reverse direction (used for function parameters, e.g. \code{?\allowbreak{}~\allowbreak{}super~\allowbreak{}T}). \textbf{Invariance} means no such relationship exists at all — \code{List\textless{}\allowbreak{}Cat\textgreater{}} and \code{List\textless{}\allowbreak{}Animal\textgreater{}} are completely unrelated types, even though \code{Cat} and \code{Animal} are related.

Java arrays are covariant. Java generics are invariant. This mismatch is one of the most important things to understand for a code-review interview, because it explains half of the wildcard rules above.

\subsection[Arrays are covariant]{Arrays are covariant}
\label{h:4:arrays-are-covariant}
\begin{codebox}
\begin{Verbatim}[commandchars=\\\{\}]
\PY{n}{Object}\PY{o}{[}\PY{o}{]}\PY{+w}{ }\PY{n}{objects}\PY{+w}{ }\PY{o}{=}\PY{+w}{ }\PY{k}{new}\PY{+w}{ }\PY{n}{String}\PY{o}{[}\PY{l+m+mi}{3}\PY{o}{]}\PY{p}{;}\PY{+w}{ }\PY{c+c1}{// legal: String[] IS\PYZhy{}A Object[]}
\PY{n}{objects}\PY{o}{[}\PY{l+m+mi}{0}\PY{o}{]}\PY{+w}{ }\PY{o}{=}\PY{+w}{ }\PY{l+s}{\PYZdq{}}\PY{l+s}{hello}\PY{l+s}{\PYZdq{}}\PY{p}{;}\PY{+w}{              }\PY{c+c1}{// fine}
\end{Verbatim}
\end{codebox}

This is convenient, but it means the \emph{compile-time} type of the array reference (\code{Object[]}) does not match its \emph{actual runtime} type (\code{String[]}). The JVM has to check every array write at runtime to protect against violations.

\subsection[Generics are invariant]{Generics are invariant}
\label{h:4:generics-are-invariant}
\begin{codebox}
\begin{Verbatim}[commandchars=\\\{\}]
\PY{c+c1}{// List\PYZlt{}Object\PYZgt{} objects = new ArrayList\PYZlt{}String\PYZgt{}();}
\PY{c+c1}{// compile error: incompatible types: ArrayList\PYZlt{}String\PYZgt{} cannot be converted to List\PYZlt{}Object\PYZgt{}}
\end{Verbatim}
\end{codebox}

Even though \code{String} IS-A \code{Object}, \code{List\textless{}\allowbreak{}String\textgreater{}} is NOT a \code{List\textless{}\allowbreak{}Object\textgreater{}}. This is invariance, and it is deliberate — generics were designed to catch the exact mistake that array covariance allows to slip through to runtime (see next section). If you need covariant-like behavior for generics, that is exactly what \code{?\allowbreak{}~\allowbreak{}extends~\allowbreak{}T} wildcards are for.

\subsection[ArrayStoreException demo]{\code{ArrayStoreException} demo}
\label{h:4:arraystoreexception-demo}
Because arrays check types at runtime, a covariant array write that violates the actual element type fails with \code{ArrayStoreException} — at runtime, not compile time.

\begin{codebox}
\begin{Verbatim}[commandchars=\\\{\}]
\PY{n}{Object}\PY{o}{[}\PY{o}{]}\PY{+w}{ }\PY{n}{objects}\PY{+w}{ }\PY{o}{=}\PY{+w}{ }\PY{k}{new}\PY{+w}{ }\PY{n}{String}\PY{o}{[}\PY{l+m+mi}{3}\PY{o}{]}\PY{p}{;}\PY{+w}{ }\PY{c+c1}{// compiles: array covariance}
\PY{n}{objects}\PY{o}{[}\PY{l+m+mi}{0}\PY{o}{]}\PY{+w}{ }\PY{o}{=}\PY{+w}{ }\PY{l+s}{\PYZdq{}}\PY{l+s}{safe}\PY{l+s}{\PYZdq{}}\PY{p}{;}\PY{+w}{               }\PY{c+c1}{// OK, actually a String}
\PY{n}{objects}\PY{o}{[}\PY{l+m+mi}{1}\PY{o}{]}\PY{+w}{ }\PY{o}{=}\PY{+w}{ }\PY{l+m+mi}{42}\PY{p}{;}\PY{+w}{                   }\PY{c+c1}{// compiles (Integer IS\PYZhy{}A Object)...}
\PY{c+c1}{// ...but throws java.lang.ArrayStoreException at runtime,}
\PY{c+c1}{// because the actual array is a String[] and 42 is not a String.}
\end{Verbatim}
\end{codebox}

Generics avoid this entire class of bug at compile time:

\begin{codebox}
\begin{Verbatim}[commandchars=\\\{\}]
\PY{n}{List}\PY{o}{\PYZlt{}}\PY{n}{String}\PY{o}{\PYZgt{}}\PY{+w}{ }\PY{n}{list}\PY{+w}{ }\PY{o}{=}\PY{+w}{ }\PY{k}{new}\PY{+w}{ }\PY{n}{ArrayList}\PY{o}{\PYZlt{}}\PY{o}{\PYZgt{}}\PY{p}{(}\PY{p}{)}\PY{p}{;}
\PY{n}{list}\PY{p}{.}\PY{n+na}{add}\PY{p}{(}\PY{l+s}{\PYZdq{}}\PY{l+s}{safe}\PY{l+s}{\PYZdq{}}\PY{p}{)}\PY{p}{;}
\PY{c+c1}{// list.add(42); // compile error: incompatible types — caught here, never reaches runtime}
\end{Verbatim}
\end{codebox}

This is the single clearest argument for “generics over arrays” in a code review: \textbf{prefer \code{List\textless{}\allowbreak{}T\textgreater{}} to \code{T[]}} in APIs whenever possible, because the failure mode moves from a runtime exception to a compile error.

{\tablefont
\begin{xltabular}{\linewidth}{@{}L{0.637}L{1.142}L{1.221}@{}}
\toprule
\rowcolor{tablehdr}
\thd{Aspect} & \thd{Arrays} & \thd{Generics} \\
\midrule
\endhead
Variance & Covariant (\code{String[]} IS-A \code{Object[]}) & Invariant (\code{List\textless{}\allowbreak{}String\textgreater{}} is NOT a \code{List\textless{}\allowbreak{}Object\textgreater{}}) \\
Type check timing & Runtime (per-element store check) & Compile time (erased after that) \\
Failure mode on mismatch & \code{ArrayStoreException} at runtime & Compile error, or none if avoided by design \\
Runtime type info & Preserved (\code{array.\allowbreak{}getClass()} knows \code{String[]}) & Erased (no \code{List\textless{}\allowbreak{}String\textgreater{}.\allowbreak{}class}) \\
\bottomrule
\end{xltabular}
}

\subsection[Method return type covariance]{Method return type covariance}
\label{h:4:method-return-type-covariance}
There is a second, unrelated meaning of “covariance” that shows up in overriding: an overriding method’s return type may be a \emph{subtype} of the overridden method’s return type. This is legal and has nothing to do with type erasure — it is about the class hierarchy, not generics.

\begin{codebox}
\begin{Verbatim}[commandchars=\\\{\}]
\PY{k+kd}{class} \PY{n+nc}{Animal}\PY{+w}{ }\PY{p}{\PYZob{}}
\PY{+w}{    }\PY{n}{Animal}\PY{+w}{ }\PY{n+nf}{reproduce}\PY{p}{(}\PY{p}{)}\PY{+w}{ }\PY{p}{\PYZob{}}\PY{+w}{ }\PY{k}{return}\PY{+w}{ }\PY{k}{new}\PY{+w}{ }\PY{n}{Animal}\PY{p}{(}\PY{p}{)}\PY{p}{;}\PY{+w}{ }\PY{p}{\PYZcb{}}
\PY{p}{\PYZcb{}}

\PY{k+kd}{class} \PY{n+nc}{Cat}\PY{+w}{ }\PY{k+kd}{extends}\PY{+w}{ }\PY{n}{Animal}\PY{+w}{ }\PY{p}{\PYZob{}}
\PY{+w}{    }\PY{n+nd}{@Override}
\PY{+w}{    }\PY{n}{Cat}\PY{+w}{ }\PY{n+nf}{reproduce}\PY{p}{(}\PY{p}{)}\PY{+w}{ }\PY{p}{\PYZob{}}\PY{+w}{ }\PY{k}{return}\PY{+w}{ }\PY{k}{new}\PY{+w}{ }\PY{n}{Cat}\PY{p}{(}\PY{p}{)}\PY{p}{;}\PY{+w}{ }\PY{p}{\PYZcb{}}\PY{+w}{ }\PY{c+c1}{// covariant return type — legal since Java 5}
\PY{p}{\PYZcb{}}
\end{Verbatim}
\end{codebox}

Interview tip: do not confuse this (covariant return types, a class-hierarchy feature) with generic type variance (\code{List\textless{}\allowbreak{}Cat\textgreater{}} vs \code{List\textless{}\allowbreak{}Animal\textgreater{}}, which is invariant unless you use wildcards). They are two different “covariance” concepts that happen to share a name.

\section[Common Code-Review Interview Pitfalls]{Common Code-Review Interview Pitfalls}
\label{h:4:common-code-review-interview-pitfalls}
\begin{enumerate}
\item 
\textbf{Using raw types “to make the compiler stop complaining.”}
Why it matters: raw types disable all generic type checking, turning compile-time errors into runtime \code{ClassCastException}s.

\begin{codebox}
\begin{Verbatim}[commandchars=\\\{\}]
\PY{c+c1}{// Before}
\PY{n}{List}\PY{+w}{ }\PY{n}{list}\PY{+w}{ }\PY{o}{=}\PY{+w}{ }\PY{k}{new}\PY{+w}{ }\PY{n}{ArrayList}\PY{p}{(}\PY{p}{)}\PY{p}{;}
\PY{n}{list}\PY{p}{.}\PY{n+na}{add}\PY{p}{(}\PY{l+s}{\PYZdq{}}\PY{l+s}{a}\PY{l+s}{\PYZdq{}}\PY{p}{)}\PY{p}{;}
\PY{n}{list}\PY{p}{.}\PY{n+na}{add}\PY{p}{(}\PY{l+m+mi}{1}\PY{p}{)}\PY{p}{;}\PY{+w}{ }\PY{c+c1}{// silently allowed}

\PY{c+c1}{// After}
\PY{n}{List}\PY{o}{\PYZlt{}}\PY{n}{String}\PY{o}{\PYZgt{}}\PY{+w}{ }\PY{n}{list}\PY{+w}{ }\PY{o}{=}\PY{+w}{ }\PY{k}{new}\PY{+w}{ }\PY{n}{ArrayList}\PY{o}{\PYZlt{}}\PY{o}{\PYZgt{}}\PY{p}{(}\PY{p}{)}\PY{p}{;}
\PY{n}{list}\PY{p}{.}\PY{n+na}{add}\PY{p}{(}\PY{l+s}{\PYZdq{}}\PY{l+s}{a}\PY{l+s}{\PYZdq{}}\PY{p}{)}\PY{p}{;}
\PY{c+c1}{// list.add(1); // compile error — caught immediately}
\end{Verbatim}
\end{codebox}

\item 
\textbf{Suppressing \code{unchecked} warnings without verifying safety.}
Why it matters: \code{@\allowbreak{}Suppress\allowbreak{}Warnings(\allowbreak{}\textquotedbl{}unchecked\textquotedbl{})} is a promise to the compiler that you manually verified type safety. Blind suppression hides real bugs.

\begin{codebox}
\begin{Verbatim}[commandchars=\\\{\}]
\PY{c+c1}{// Before — suppressed without justification}
\PY{n+nd}{@SuppressWarnings}\PY{p}{(}\PY{l+s}{\PYZdq{}}\PY{l+s}{unchecked}\PY{l+s}{\PYZdq{}}\PY{p}{)}
\PY{n}{List}\PY{o}{\PYZlt{}}\PY{n}{String}\PY{o}{\PYZgt{}}\PY{+w}{ }\PY{n}{list}\PY{+w}{ }\PY{o}{=}\PY{+w}{ }\PY{p}{(}\PY{n}{List}\PY{o}{\PYZlt{}}\PY{n}{String}\PY{o}{\PYZgt{}}\PY{p}{)}\PY{+w}{ }\PY{n}{someRawList}\PY{p}{;}

\PY{c+c1}{// After — scoped narrowly, with a comment explaining why it\PYZsq{}s safe}
\PY{n+nd}{@SuppressWarnings}\PY{p}{(}\PY{l+s}{\PYZdq{}}\PY{l+s}{unchecked}\PY{l+s}{\PYZdq{}}\PY{p}{)}\PY{+w}{ }\PY{c+c1}{// safe: someRawList is populated exclusively by addString() above}
\PY{n}{List}\PY{o}{\PYZlt{}}\PY{n}{String}\PY{o}{\PYZgt{}}\PY{+w}{ }\PY{n}{list}\PY{+w}{ }\PY{o}{=}\PY{+w}{ }\PY{p}{(}\PY{n}{List}\PY{o}{\PYZlt{}}\PY{n}{String}\PY{o}{\PYZgt{}}\PY{p}{)}\PY{+w}{ }\PY{n}{someRawList}\PY{p}{;}
\end{Verbatim}
\end{codebox}

\item 
\textbf{Designing an API with \code{T[]} instead of \code{List\textless{}\allowbreak{}T\textgreater{}}.}
Why it matters: arrays are covariant and only check types at runtime (\code{ArrayStoreException}); \code{List\textless{}\allowbreak{}T\textgreater{}} is invariant and checked at compile time.

\begin{codebox}
\begin{Verbatim}[commandchars=\\\{\}]
\PY{c+c1}{// Before}
\PY{k+kt}{void}\PY{+w}{ }\PY{n+nf}{process}\PY{p}{(}\PY{n}{Number}\PY{o}{[}\PY{o}{]}\PY{+w}{ }\PY{n}{values}\PY{p}{)}\PY{+w}{ }\PY{p}{\PYZob{}}\PY{+w}{ }\PY{n}{values}\PY{o}{[}\PY{l+m+mi}{0}\PY{o}{]}\PY{+w}{ }\PY{o}{=}\PY{+w}{ }\PY{l+m+mf}{3.14}\PY{p}{;}\PY{+w}{ }\PY{p}{\PYZcb{}}\PY{+w}{ }\PY{c+c1}{// ArrayStoreException risk if called with Integer[]}

\PY{c+c1}{// After}
\PY{k+kt}{void}\PY{+w}{ }\PY{n+nf}{process}\PY{p}{(}\PY{n}{List}\PY{o}{\PYZlt{}}\PY{o}{?}\PY{+w}{ }\PY{k+kd}{extends}\PY{+w}{ }\PY{n}{Number}\PY{o}{\PYZgt{}}\PY{+w}{ }\PY{n}{values}\PY{p}{)}\PY{+w}{ }\PY{p}{\PYZob{}}\PY{+w}{ }\PY{c+cm}{/* read\PYZhy{}only access is compile\PYZhy{}time safe */}\PY{+w}{ }\PY{p}{\PYZcb{}}
\end{Verbatim}
\end{codebox}

\item 
\textbf{Using \code{?\allowbreak{}~\allowbreak{}extends~\allowbreak{}T} on a parameter you need to write into.}
Why it matters: \code{extends} wildcards are producer-only; the compiler will reject \code{add()} calls, and “fixing” it by removing the wildcard silently loses flexibility.

\begin{codebox}
\begin{Verbatim}[commandchars=\\\{\}]
\PY{c+c1}{// Before}
\PY{k+kt}{void}\PY{+w}{ }\PY{n+nf}{fill}\PY{p}{(}\PY{n}{List}\PY{o}{\PYZlt{}}\PY{o}{?}\PY{+w}{ }\PY{k+kd}{extends}\PY{+w}{ }\PY{n}{Number}\PY{o}{\PYZgt{}}\PY{+w}{ }\PY{n}{list}\PY{p}{)}\PY{+w}{ }\PY{p}{\PYZob{}}
\PY{+w}{    }\PY{c+c1}{// list.add(1); // compile error}
\PY{p}{\PYZcb{}}

\PY{c+c1}{// After — use super if you need to write, or an exact type if you need both}
\PY{k+kt}{void}\PY{+w}{ }\PY{n+nf}{fill}\PY{p}{(}\PY{n}{List}\PY{o}{\PYZlt{}}\PY{o}{?}\PY{+w}{ }\PY{k+kd}{super}\PY{+w}{ }\PY{n}{Number}\PY{o}{\PYZgt{}}\PY{+w}{ }\PY{n}{list}\PY{p}{)}\PY{+w}{ }\PY{p}{\PYZob{}}
\PY{+w}{    }\PY{n}{list}\PY{p}{.}\PY{n+na}{add}\PY{p}{(}\PY{l+m+mi}{1}\PY{p}{)}\PY{p}{;}\PY{+w}{ }\PY{c+c1}{// OK}
\PY{p}{\PYZcb{}}
\end{Verbatim}
\end{codebox}

\item 
\textbf{Assuming \code{List\textless{}\allowbreak{}String\textgreater{}} and \code{List\textless{}\allowbreak{}Object\textgreater{}} are interchangeable.}
Why it matters: generics are invariant; this is one of the most common “why won’t this compile” confusions.

\begin{codebox}
\begin{Verbatim}[commandchars=\\\{\}]
\PY{c+c1}{// Before}
\PY{c+c1}{// List\PYZlt{}Object\PYZgt{} objs = new ArrayList\PYZlt{}String\PYZgt{}(); // compile error}

\PY{c+c1}{// After — use a wildcard if you only need read access}
\PY{n}{List}\PY{o}{\PYZlt{}}\PY{o}{?}\PY{+w}{ }\PY{k+kd}{extends}\PY{+w}{ }\PY{n}{Object}\PY{o}{\PYZgt{}}\PY{+w}{ }\PY{n}{objs}\PY{+w}{ }\PY{o}{=}\PY{+w}{ }\PY{k}{new}\PY{+w}{ }\PY{n}{ArrayList}\PY{o}{\PYZlt{}}\PY{n}{String}\PY{o}{\PYZgt{}}\PY{p}{(}\PY{p}{)}\PY{p}{;}\PY{+w}{ }\PY{c+c1}{// fine, though rarely useful since Object adds nothing}
\end{Verbatim}
\end{codebox}

\item 
\textbf{Trying \code{instanceof~\allowbreak{}List\textless{}\allowbreak{}String\textgreater{}} or catching a specific generic exception type per type argument.}
Why it matters: type erasure removes the type argument at runtime; the JVM cannot check it.

\begin{codebox}
\begin{Verbatim}[commandchars=\\\{\}]
\PY{c+c1}{// Before}
\PY{c+c1}{// if (obj instanceof List\PYZlt{}String\PYZgt{}) \PYZob{} \PYZcb{} // compile error}

\PY{c+c1}{// After}
\PY{k}{if}\PY{+w}{ }\PY{p}{(}\PY{n}{obj}\PY{+w}{ }\PY{k}{instanceof}\PY{+w}{ }\PY{n}{List}\PY{o}{\PYZlt{}}\PY{o}{?}\PY{o}{\PYZgt{}}\PY{+w}{ }\PY{n}{list}\PY{+w}{ }\PY{o}{\PYZam{}}\PY{o}{\PYZam{}}\PY{+w}{ }\PY{o}{!}\PY{n}{list}\PY{p}{.}\PY{n+na}{isEmpty}\PY{p}{(}\PY{p}{)}\PY{+w}{ }\PY{o}{\PYZam{}}\PY{o}{\PYZam{}}\PY{+w}{ }\PY{n}{list}\PY{p}{.}\PY{n+na}{get}\PY{p}{(}\PY{l+m+mi}{0}\PY{p}{)}\PY{+w}{ }\PY{k}{instanceof}\PY{+w}{ }\PY{n}{String}\PY{p}{)}\PY{+w}{ }\PY{p}{\PYZob{}}
\PY{+w}{    }\PY{c+c1}{// manually verify the first element\PYZsq{}s runtime type instead}
\PY{p}{\PYZcb{}}
\end{Verbatim}
\end{codebox}

\item 
\textbf{Forgetting that overloaded generic methods can collide after erasure.}
Why it matters: two overloads that differ only by type argument have the \emph{same} erased signature and will not compile.

\begin{codebox}
\begin{Verbatim}[commandchars=\\\{\}]
\PY{c+c1}{// Before}
\PY{c+c1}{// void process(List\PYZlt{}String\PYZgt{} list) \PYZob{} \PYZcb{}}
\PY{c+c1}{// void process(List\PYZlt{}Integer\PYZgt{} list) \PYZob{} \PYZcb{} // compile error: both erase to process(List)}

\PY{c+c1}{// After — differentiate by method name or take a wrapper/marker type}
\PY{k+kt}{void}\PY{+w}{ }\PY{n+nf}{processStrings}\PY{p}{(}\PY{n}{List}\PY{o}{\PYZlt{}}\PY{n}{String}\PY{o}{\PYZgt{}}\PY{+w}{ }\PY{n}{list}\PY{p}{)}\PY{+w}{ }\PY{p}{\PYZob{}}\PY{+w}{ }\PY{p}{\PYZcb{}}
\PY{k+kt}{void}\PY{+w}{ }\PY{n+nf}{processIntegers}\PY{p}{(}\PY{n}{List}\PY{o}{\PYZlt{}}\PY{n}{Integer}\PY{o}{\PYZgt{}}\PY{+w}{ }\PY{n}{list}\PY{p}{)}\PY{+w}{ }\PY{p}{\PYZob{}}\PY{+w}{ }\PY{p}{\PYZcb{}}
\end{Verbatim}
\end{codebox}

\item 
\textbf{Writing an unsafe generic varargs method without \code{@\allowbreak{}SafeVarargs}, or worse, adding \code{@\allowbreak{}SafeVarargs} to an unsafe one.}
Why it matters: \code{@\allowbreak{}SafeVarargs} is a promise; misusing it hides real heap pollution bugs, and omitting it on a genuinely safe method just produces noisy warnings for every caller.

\begin{codebox}
\begin{Verbatim}[commandchars=\\\{\}]
\PY{c+c1}{// Before — unsafe, and no annotation to flag the risk}
\PY{k+kd}{static}\PY{+w}{ }\PY{o}{\PYZlt{}}\PY{n}{T}\PY{o}{\PYZgt{}}\PY{+w}{ }\PY{k+kt}{void}\PY{+w}{ }\PY{n+nf}{addAll}\PY{p}{(}\PY{n}{List}\PY{o}{\PYZlt{}}\PY{n}{T}\PY{o}{\PYZgt{}}\PY{+w}{ }\PY{n}{list}\PY{p}{,}\PY{+w}{ }\PY{n}{T}\PY{p}{.}\PY{p}{.}\PY{p}{.}\PY{+w}{ }\PY{n}{items}\PY{p}{)}\PY{+w}{ }\PY{p}{\PYZob{}}
\PY{+w}{    }\PY{n}{Object}\PY{o}{[}\PY{o}{]}\PY{+w}{ }\PY{n}{array}\PY{+w}{ }\PY{o}{=}\PY{+w}{ }\PY{n}{items}\PY{p}{;}
\PY{+w}{    }\PY{n}{array}\PY{o}{[}\PY{l+m+mi}{0}\PY{o}{]}\PY{+w}{ }\PY{o}{=}\PY{+w}{ }\PY{l+s}{\PYZdq{}}\PY{l+s}{corrupt}\PY{l+s}{\PYZdq{}}\PY{p}{;}\PY{+w}{ }\PY{c+c1}{// writes an unrelated type into the array — heap pollution}
\PY{p}{\PYZcb{}}

\PY{c+c1}{// After — either make it actually safe (read\PYZhy{}only) and annotate, or take a List parameter instead}
\PY{n+nd}{@SafeVarargs}
\PY{k+kd}{static}\PY{+w}{ }\PY{o}{\PYZlt{}}\PY{n}{T}\PY{o}{\PYZgt{}}\PY{+w}{ }\PY{n}{List}\PY{o}{\PYZlt{}}\PY{n}{T}\PY{o}{\PYZgt{}}\PY{+w}{ }\PY{n+nf}{listOf}\PY{p}{(}\PY{n}{T}\PY{p}{.}\PY{p}{.}\PY{p}{.}\PY{+w}{ }\PY{n}{items}\PY{p}{)}\PY{+w}{ }\PY{p}{\PYZob{}}
\PY{+w}{    }\PY{k}{return}\PY{+w}{ }\PY{n}{List}\PY{p}{.}\PY{n+na}{of}\PY{p}{(}\PY{n}{items}\PY{p}{)}\PY{p}{;}\PY{+w}{ }\PY{c+c1}{// read\PYZhy{}only, never stores into items}
\PY{p}{\PYZcb{}}
\end{Verbatim}
\end{codebox}

\item 
\textbf{Expecting a static field or method to “remember” the type argument per instance.}
Why it matters: erasure means there is one class, one set of static members, shared across all parameterizations (\code{Box\textless{}\allowbreak{}String\textgreater{}}, \code{Box\textless{}\allowbreak{}Integer\textgreater{}}, …).

\begin{codebox}
\begin{Verbatim}[commandchars=\\\{\}]
\PY{c+c1}{// Before}
\PY{k+kd}{class} \PY{n+nc}{Box}\PY{o}{\PYZlt{}}\PY{n}{T}\PY{o}{\PYZgt{}}\PY{+w}{ }\PY{p}{\PYZob{}}
\PY{+w}{    }\PY{c+c1}{// static T lastCreated; // compile error}
\PY{p}{\PYZcb{}}

\PY{c+c1}{// After}
\PY{k+kd}{class} \PY{n+nc}{Box}\PY{o}{\PYZlt{}}\PY{n}{T}\PY{o}{\PYZgt{}}\PY{+w}{ }\PY{p}{\PYZob{}}
\PY{+w}{    }\PY{k+kd}{static}\PY{+w}{ }\PY{n}{Object}\PY{+w}{ }\PY{n}{lastCreated}\PY{p}{;}\PY{+w}{ }\PY{c+c1}{// shared across all Box\PYZlt{}...\PYZgt{} instances, by design}
\PY{p}{\PYZcb{}}
\end{Verbatim}
\end{codebox}

\item 
\textbf{Not recognizing bridge methods when reading decompiled/reflected code.}
Why it matters: reviewers who see two overloaded-looking methods with the same name after implementing a generic interface may think it’s a bug, when it’s a compiler-generated bridge method required for correct polymorphism.

\begin{codebox}
\begin{Verbatim}[commandchars=\\\{\}]
\PY{k+kd}{interface} \PY{n+nc}{Container}\PY{o}{\PYZlt{}}\PY{n}{T}\PY{o}{\PYZgt{}}\PY{+w}{ }\PY{p}{\PYZob{}}\PY{+w}{ }\PY{k+kt}{void}\PY{+w}{ }\PY{n+nf}{set}\PY{p}{(}\PY{n}{T}\PY{+w}{ }\PY{n}{value}\PY{p}{)}\PY{p}{;}\PY{+w}{ }\PY{p}{\PYZcb{}}
\PY{k+kd}{class} \PY{n+nc}{StringContainer}\PY{+w}{ }\PY{k+kd}{implements}\PY{+w}{ }\PY{n}{Container}\PY{o}{\PYZlt{}}\PY{n}{String}\PY{o}{\PYZgt{}}\PY{+w}{ }\PY{p}{\PYZob{}}
\PY{+w}{    }\PY{k+kd}{public}\PY{+w}{ }\PY{k+kt}{void}\PY{+w}{ }\PY{n+nf}{set}\PY{p}{(}\PY{n}{String}\PY{+w}{ }\PY{n}{value}\PY{p}{)}\PY{+w}{ }\PY{p}{\PYZob{}}\PY{+w}{ }\PY{c+cm}{/* ... */}\PY{+w}{ }\PY{p}{\PYZcb{}}
\PY{+w}{    }\PY{c+c1}{// compiler silently also generates: public void set(Object value) \PYZob{} set((String) value); \PYZcb{}}
\PY{p}{\PYZcb{}}
\end{Verbatim}
\end{codebox}

\item 
\textbf{Using a recursive bound incorrectly, e.g. \code{\textless{}\allowbreak{}T~\allowbreak{}extends~\allowbreak{}Comparable\textgreater{}} (raw) instead of \code{\textless{}\allowbreak{}T~\allowbreak{}extends~\allowbreak{}Comparable\textless{}\allowbreak{}T\textgreater{}\textgreater{}}.}
Why it matters: the raw bound loses type safety — \code{compareTo} accepts any \code{Object}, so mismatched comparisons compile and fail at runtime instead of compile time.

\begin{codebox}
\begin{Verbatim}[commandchars=\\\{\}]
\PY{c+c1}{// Before}
\PY{k+kd}{static}\PY{+w}{ }\PY{o}{\PYZlt{}}\PY{n}{T}\PY{+w}{ }\PY{k+kd}{extends}\PY{+w}{ }\PY{n}{Comparable}\PY{o}{\PYZgt{}}\PY{+w}{ }\PY{n}{T}\PY{+w}{ }\PY{n+nf}{max}\PY{p}{(}\PY{n}{List}\PY{o}{\PYZlt{}}\PY{n}{T}\PY{o}{\PYZgt{}}\PY{+w}{ }\PY{n}{list}\PY{p}{)}\PY{+w}{ }\PY{p}{\PYZob{}}\PY{+w}{ }\PY{c+cm}{/* ... */}\PY{+w}{ }\PY{p}{\PYZcb{}}\PY{+w}{ }\PY{c+c1}{// raw Comparable, weak guarantee}

\PY{c+c1}{// After}
\PY{k+kd}{static}\PY{+w}{ }\PY{o}{\PYZlt{}}\PY{n}{T}\PY{+w}{ }\PY{k+kd}{extends}\PY{+w}{ }\PY{n}{Comparable}\PY{o}{\PYZlt{}}\PY{n}{T}\PY{o}{\PYZgt{}}\PY{o}{\PYZgt{}}\PY{+w}{ }\PY{n}{T}\PY{+w}{ }\PY{n+nf}{max}\PY{p}{(}\PY{n}{List}\PY{o}{\PYZlt{}}\PY{n}{T}\PY{o}{\PYZgt{}}\PY{+w}{ }\PY{n}{list}\PY{p}{)}\PY{+w}{ }\PY{p}{\PYZob{}}\PY{+w}{ }\PY{c+cm}{/* ... */}\PY{+w}{ }\PY{p}{\PYZcb{}}\PY{+w}{ }\PY{c+c1}{// T can only compare to T}
\end{Verbatim}
\end{codebox}

\item 
\textbf{Confusing \code{var} with type erasure or loss of type safety.}
Why it matters: \code{var} infers a concrete compile-time type from the initializer; it is not related to generics’ runtime erasure and does not weaken checking.

\begin{codebox}
\begin{Verbatim}[commandchars=\\\{\}]
\PY{c+c1}{// Before — reviewer flags this as \PYZdq{}unsafe, dynamic typing\PYZdq{}}
\PY{k+kd}{var}\PY{+w}{ }\PY{n}{names}\PY{+w}{ }\PY{o}{=}\PY{+w}{ }\PY{k}{new}\PY{+w}{ }\PY{n}{ArrayList}\PY{o}{\PYZlt{}}\PY{n}{String}\PY{o}{\PYZgt{}}\PY{p}{(}\PY{p}{)}\PY{p}{;}

\PY{c+c1}{// After — clarify: names has static type ArrayList\PYZlt{}String\PYZgt{}, exactly as if written out}
\PY{n}{ArrayList}\PY{o}{\PYZlt{}}\PY{n}{String}\PY{o}{\PYZgt{}}\PY{+w}{ }\PY{n}{names}\PY{+w}{ }\PY{o}{=}\PY{+w}{ }\PY{k}{new}\PY{+w}{ }\PY{n}{ArrayList}\PY{o}{\PYZlt{}}\PY{o}{\PYZgt{}}\PY{p}{(}\PY{p}{)}\PY{p}{;}\PY{+w}{ }\PY{c+c1}{// equivalent compile\PYZhy{}time behavior}
\end{Verbatim}
\end{codebox}

\item 
\textbf{Relying on \code{List.\allowbreak{}of()} or an empty generic literal without a target type.}
Why it matters: without a target type, the compiler infers the most general type (often \code{Object}), which silently defeats type checking downstream.

\begin{codebox}
\begin{Verbatim}[commandchars=\\\{\}]
\PY{c+c1}{// Before}
\PY{k+kd}{var}\PY{+w}{ }\PY{n}{empty}\PY{+w}{ }\PY{o}{=}\PY{+w}{ }\PY{n}{List}\PY{p}{.}\PY{n+na}{of}\PY{p}{(}\PY{p}{)}\PY{p}{;}\PY{+w}{ }\PY{c+c1}{// inferred as List\PYZlt{}Object\PYZgt{}}

\PY{c+c1}{// After — give the compiler a target type}
\PY{n}{List}\PY{o}{\PYZlt{}}\PY{n}{String}\PY{o}{\PYZgt{}}\PY{+w}{ }\PY{n}{empty}\PY{+w}{ }\PY{o}{=}\PY{+w}{ }\PY{n}{List}\PY{p}{.}\PY{n+na}{of}\PY{p}{(}\PY{p}{)}\PY{p}{;}\PY{+w}{ }\PY{c+c1}{// inferred correctly as List\PYZlt{}String\PYZgt{}}
\end{Verbatim}
\end{codebox}

\item 
\textbf{Treating \code{PECS} as optional API polish instead of a correctness tool.}
Why it matters: skipping \code{extends}/\code{super} on library-style methods forces callers into exact-type matches, breaking otherwise-valid calls and causing unnecessary casts or \code{List.\allowbreak{}copyOf()} calls at call sites.

\begin{codebox}
\begin{Verbatim}[commandchars=\\\{\}]
\PY{c+c1}{// Before — overly strict, exact type required}
\PY{k+kd}{static}\PY{+w}{ }\PY{k+kt}{double}\PY{+w}{ }\PY{n+nf}{sum}\PY{p}{(}\PY{n}{List}\PY{o}{\PYZlt{}}\PY{n}{Double}\PY{o}{\PYZgt{}}\PY{+w}{ }\PY{n}{values}\PY{p}{)}\PY{+w}{ }\PY{p}{\PYZob{}}\PY{+w}{ }\PY{c+cm}{/* ... */}\PY{+w}{ }\PY{p}{\PYZcb{}}
\PY{c+c1}{// sum(List.of(1, 2, 3)); // compile error: List\PYZlt{}Integer\PYZgt{} is not List\PYZlt{}Double\PYZgt{}}

\PY{c+c1}{// After — PECS: this is a producer, so use extends}
\PY{k+kd}{static}\PY{+w}{ }\PY{k+kt}{double}\PY{+w}{ }\PY{n+nf}{sum}\PY{p}{(}\PY{n}{List}\PY{o}{\PYZlt{}}\PY{o}{?}\PY{+w}{ }\PY{k+kd}{extends}\PY{+w}{ }\PY{n}{Number}\PY{o}{\PYZgt{}}\PY{+w}{ }\PY{n}{values}\PY{p}{)}\PY{+w}{ }\PY{p}{\PYZob{}}\PY{+w}{ }\PY{c+cm}{/* ... */}\PY{+w}{ }\PY{p}{\PYZcb{}}
\PY{c+c1}{// sum(List.of(1, 2, 3)); // now compiles}
\end{Verbatim}
\end{codebox}

\item 
\textbf{Assuming array covariance is “safe” because it compiles.}
Why it matters: covariant array writes that violate the actual runtime type compile cleanly but throw \code{ArrayStoreException} in production — often only under rarely-hit code paths, making it a nasty review-time-invisible bug.

\begin{codebox}
\begin{Verbatim}[commandchars=\\\{\}]
\PY{c+c1}{// Before}
\PY{n}{Object}\PY{o}{[}\PY{o}{]}\PY{+w}{ }\PY{n}{items}\PY{+w}{ }\PY{o}{=}\PY{+w}{ }\PY{k}{new}\PY{+w}{ }\PY{n}{String}\PY{o}{[}\PY{l+m+mi}{3}\PY{o}{]}\PY{p}{;}
\PY{n}{items}\PY{o}{[}\PY{l+m+mi}{0}\PY{o}{]}\PY{+w}{ }\PY{o}{=}\PY{+w}{ }\PY{l+m+mi}{42}\PY{p}{;}\PY{+w}{ }\PY{c+c1}{// compiles, throws ArrayStoreException at runtime}

\PY{c+c1}{// After — use a generic collection so the mistake is caught at compile time}
\PY{n}{List}\PY{o}{\PYZlt{}}\PY{n}{String}\PY{o}{\PYZgt{}}\PY{+w}{ }\PY{n}{items}\PY{+w}{ }\PY{o}{=}\PY{+w}{ }\PY{k}{new}\PY{+w}{ }\PY{n}{ArrayList}\PY{o}{\PYZlt{}}\PY{o}{\PYZgt{}}\PY{p}{(}\PY{p}{)}\PY{p}{;}
\PY{c+c1}{// items.add(42); // compile error, caught immediately}
\end{Verbatim}
\end{codebox}

\end{enumerate}


\setcounter{chapter}{4}
\chapter{Object Class \& Common APIs}
\label{chap:5}
\label{h:5:5-object-class--common-apis}
Every class in Java, whether you write it or not, is a child of \code{java.\allowbreak{}lang.\allowbreak{}Object}. That means every object you create already has a few methods built in: \code{equals()}, \code{hashCode()}, \code{toString()}, \code{clone()}, and others. In code review, bugs in these “invisible” methods are some of the most common findings — a broken \code{equals()} can silently corrupt a \code{HashMap}, and a missing defensive copy can let external code mutate your “immutable” class. This chapter walks through these APIs in depth, plus the related tools (\code{Optional}, \code{Comparator}, immutability patterns) that reviewers check every day.

\section[equals(), hashCode(), and toString()]{equals(), hashCode(), and toString()}
\label{h:5:equals-hashcode-and-tostring}
\code{Object} gives every class three methods for free:

\begin{itemize}
\item 
\code{equals(\allowbreak{}Object~\allowbreak{}other)} — decides if two objects are “equal”.

\item 
\code{hashCode()} — returns an \code{int} “fingerprint” used by hash-based collections (\code{HashMap}, \code{HashSet}).

\item 
\code{toString()} — returns a human-readable \code{String} representation, used by logging, debuggers, and string concatenation.

\end{itemize}

The default implementations (inherited straight from \code{Object}) compare \textbf{identity} (are they the same object in memory?) and print the class name plus a memory-based hash, like \code{com.\allowbreak{}acme.\allowbreak{}User@\allowbreak{}1b6d3586}. That is almost never what you want for a data-carrying class.

\begin{codebox}
\begin{Verbatim}[commandchars=\\\{\}]
\PY{k+kd}{class} \PY{n+nc}{User}\PY{+w}{ }\PY{p}{\PYZob{}}
\PY{+w}{    }\PY{k+kd}{private}\PY{+w}{ }\PY{k+kd}{final}\PY{+w}{ }\PY{n}{String}\PY{+w}{ }\PY{n}{email}\PY{p}{;}
\PY{+w}{    }\PY{n}{User}\PY{p}{(}\PY{n}{String}\PY{+w}{ }\PY{n}{email}\PY{p}{)}\PY{+w}{ }\PY{p}{\PYZob{}}\PY{+w}{ }\PY{k}{this}\PY{p}{.}\PY{n+na}{email}\PY{+w}{ }\PY{o}{=}\PY{+w}{ }\PY{n}{email}\PY{p}{;}\PY{+w}{ }\PY{p}{\PYZcb{}}
\PY{p}{\PYZcb{}}

\PY{n}{User}\PY{+w}{ }\PY{n}{a}\PY{+w}{ }\PY{o}{=}\PY{+w}{ }\PY{k}{new}\PY{+w}{ }\PY{n}{User}\PY{p}{(}\PY{l+s}{\PYZdq{}}\PY{l+s}{a@acme.com}\PY{l+s}{\PYZdq{}}\PY{p}{)}\PY{p}{;}
\PY{n}{User}\PY{+w}{ }\PY{n}{b}\PY{+w}{ }\PY{o}{=}\PY{+w}{ }\PY{k}{new}\PY{+w}{ }\PY{n}{User}\PY{p}{(}\PY{l+s}{\PYZdq{}}\PY{l+s}{a@acme.com}\PY{l+s}{\PYZdq{}}\PY{p}{)}\PY{p}{;}

\PY{n}{System}\PY{p}{.}\PY{n+na}{out}\PY{p}{.}\PY{n+na}{println}\PY{p}{(}\PY{n}{a}\PY{p}{.}\PY{n+na}{equals}\PY{p}{(}\PY{n}{b}\PY{p}{)}\PY{p}{)}\PY{p}{;}\PY{+w}{ }\PY{c+c1}{// false! Different objects, same data.}
\PY{n}{System}\PY{p}{.}\PY{n+na}{out}\PY{p}{.}\PY{n+na}{println}\PY{p}{(}\PY{n}{a}\PY{p}{)}\PY{p}{;}\PY{+w}{           }\PY{c+c1}{// User@1b6d3586 — useless in logs.}
\end{Verbatim}
\end{codebox}

\subsection[The equals() contract]{The equals() contract}
\label{h:5:the-equals-contract}
If you override \code{equals()}, you must honor a strict contract. It is defined in the Javadoc of \code{Object.\allowbreak{}equals} and every reviewer should be able to recite it:

{\tablefont
\begin{xltabular}{\linewidth}{@{}L{0.409}L{1.591}@{}}
\toprule
\rowcolor{tablehdr}
\thd{Property} & \thd{Meaning} \\
\midrule
\endhead
Reflexive & \code{x.\allowbreak{}equals(\allowbreak{}x)} must be \code{true}. \\
Symmetric & If \code{x.\allowbreak{}equals(\allowbreak{}y)} is \code{true}, then \code{y.\allowbreak{}equals(\allowbreak{}x)} must also be \code{true}. \\
Transitive & If \code{x.\allowbreak{}equals(\allowbreak{}y)} and \code{y.\allowbreak{}equals(\allowbreak{}z)} are \code{true}, then \code{x.\allowbreak{}equals(\allowbreak{}z)} must be \code{true}. \\
Consistent & Calling \code{x.\allowbreak{}equals(\allowbreak{}y)} repeatedly returns the same result, as long as neither object changed. \\
Null comparison & \code{x.\allowbreak{}equals(\allowbreak{}null)} must return \code{false} — never throw a \code{NullPointerException}. \\
\bottomrule
\end{xltabular}
}

A typical, correct \code{equals()}/\code{hashCode()} pair:

\begin{codebox}
\begin{Verbatim}[commandchars=\\\{\}]
\PY{k+kn}{import}\PY{+w}{ }\PY{n+nn}{java.util.Objects}\PY{p}{;}

\PY{k+kd}{final}\PY{+w}{ }\PY{k+kd}{class} \PY{n+nc}{Point}\PY{+w}{ }\PY{p}{\PYZob{}}
\PY{+w}{    }\PY{k+kd}{private}\PY{+w}{ }\PY{k+kd}{final}\PY{+w}{ }\PY{k+kt}{int}\PY{+w}{ }\PY{n}{x}\PY{p}{;}
\PY{+w}{    }\PY{k+kd}{private}\PY{+w}{ }\PY{k+kd}{final}\PY{+w}{ }\PY{k+kt}{int}\PY{+w}{ }\PY{n}{y}\PY{p}{;}

\PY{+w}{    }\PY{n}{Point}\PY{p}{(}\PY{k+kt}{int}\PY{+w}{ }\PY{n}{x}\PY{p}{,}\PY{+w}{ }\PY{k+kt}{int}\PY{+w}{ }\PY{n}{y}\PY{p}{)}\PY{+w}{ }\PY{p}{\PYZob{}}
\PY{+w}{        }\PY{k}{this}\PY{p}{.}\PY{n+na}{x}\PY{+w}{ }\PY{o}{=}\PY{+w}{ }\PY{n}{x}\PY{p}{;}
\PY{+w}{        }\PY{k}{this}\PY{p}{.}\PY{n+na}{y}\PY{+w}{ }\PY{o}{=}\PY{+w}{ }\PY{n}{y}\PY{p}{;}
\PY{+w}{    }\PY{p}{\PYZcb{}}

\PY{+w}{    }\PY{n+nd}{@Override}
\PY{+w}{    }\PY{k+kd}{public}\PY{+w}{ }\PY{k+kt}{boolean}\PY{+w}{ }\PY{n+nf}{equals}\PY{p}{(}\PY{n}{Object}\PY{+w}{ }\PY{n}{o}\PY{p}{)}\PY{+w}{ }\PY{p}{\PYZob{}}
\PY{+w}{        }\PY{k}{if}\PY{+w}{ }\PY{p}{(}\PY{k}{this}\PY{+w}{ }\PY{o}{=}\PY{o}{=}\PY{+w}{ }\PY{n}{o}\PY{p}{)}\PY{+w}{ }\PY{k}{return}\PY{+w}{ }\PY{k+kc}{true}\PY{p}{;}\PY{+w}{                 }\PY{c+c1}{// reflexive shortcut}
\PY{+w}{        }\PY{k}{if}\PY{+w}{ }\PY{p}{(}\PY{n}{o}\PY{+w}{ }\PY{o}{=}\PY{o}{=}\PY{+w}{ }\PY{k+kc}{null}\PY{+w}{ }\PY{o}{|}\PY{o}{|}\PY{+w}{ }\PY{n}{getClass}\PY{p}{(}\PY{p}{)}\PY{+w}{ }\PY{o}{!}\PY{o}{=}\PY{+w}{ }\PY{n}{o}\PY{p}{.}\PY{n+na}{getClass}\PY{p}{(}\PY{p}{)}\PY{p}{)}\PY{+w}{ }\PY{k}{return}\PY{+w}{ }\PY{k+kc}{false}\PY{p}{;}\PY{+w}{ }\PY{c+c1}{// null\PYZhy{}safe + type check}
\PY{+w}{        }\PY{n}{Point}\PY{+w}{ }\PY{n}{other}\PY{+w}{ }\PY{o}{=}\PY{+w}{ }\PY{p}{(}\PY{n}{Point}\PY{p}{)}\PY{+w}{ }\PY{n}{o}\PY{p}{;}
\PY{+w}{        }\PY{k}{return}\PY{+w}{ }\PY{n}{x}\PY{+w}{ }\PY{o}{=}\PY{o}{=}\PY{+w}{ }\PY{n}{other}\PY{p}{.}\PY{n+na}{x}\PY{+w}{ }\PY{o}{\PYZam{}}\PY{o}{\PYZam{}}\PY{+w}{ }\PY{n}{y}\PY{+w}{ }\PY{o}{=}\PY{o}{=}\PY{+w}{ }\PY{n}{other}\PY{p}{.}\PY{n+na}{y}\PY{p}{;}
\PY{+w}{    }\PY{p}{\PYZcb{}}

\PY{+w}{    }\PY{n+nd}{@Override}
\PY{+w}{    }\PY{k+kd}{public}\PY{+w}{ }\PY{k+kt}{int}\PY{+w}{ }\PY{n+nf}{hashCode}\PY{p}{(}\PY{p}{)}\PY{+w}{ }\PY{p}{\PYZob{}}
\PY{+w}{        }\PY{k}{return}\PY{+w}{ }\PY{n}{Objects}\PY{p}{.}\PY{n+na}{hash}\PY{p}{(}\PY{n}{x}\PY{p}{,}\PY{+w}{ }\PY{n}{y}\PY{p}{)}\PY{p}{;}
\PY{+w}{    }\PY{p}{\PYZcb{}}

\PY{+w}{    }\PY{n+nd}{@Override}
\PY{+w}{    }\PY{k+kd}{public}\PY{+w}{ }\PY{n}{String}\PY{+w}{ }\PY{n+nf}{toString}\PY{p}{(}\PY{p}{)}\PY{+w}{ }\PY{p}{\PYZob{}}
\PY{+w}{        }\PY{k}{return}\PY{+w}{ }\PY{l+s}{\PYZdq{}}\PY{l+s}{Point\PYZob{}x=}\PY{l+s}{\PYZdq{}}\PY{+w}{ }\PY{o}{+}\PY{+w}{ }\PY{n}{x}\PY{+w}{ }\PY{o}{+}\PY{+w}{ }\PY{l+s}{\PYZdq{}}\PY{l+s}{, y=}\PY{l+s}{\PYZdq{}}\PY{+w}{ }\PY{o}{+}\PY{+w}{ }\PY{n}{y}\PY{+w}{ }\PY{o}{+}\PY{+w}{ }\PY{l+s}{\PYZdq{}}\PY{l+s}{\PYZcb{}}\PY{l+s}{\PYZdq{}}\PY{p}{;}
\PY{+w}{    }\PY{p}{\PYZcb{}}
\PY{p}{\PYZcb{}}
\end{Verbatim}
\end{codebox}

\subsection[The hashCode() contract, and why it must match equals()]{The hashCode() contract, and why it must match equals()}
\label{h:5:the-hashcode-contract-and-why-it-must-match-equals}
The rule is simple but easy to forget: \textbf{if two objects are equal according to \code{equals()}, they must have the same \code{hashCode()}.} The reverse is not required — unequal objects are \emph{allowed} to share a hash code (a “hash collision”), but equal objects can never disagree on their hash.

\begin{codebox}
\begin{Verbatim}[commandchars=\\\{\}]
\PY{c+c1}{// BROKEN: overrides equals() but not hashCode()}
\PY{k+kd}{class} \PY{n+nc}{Bad}\PY{+w}{ }\PY{p}{\PYZob{}}
\PY{+w}{    }\PY{k+kd}{private}\PY{+w}{ }\PY{k+kd}{final}\PY{+w}{ }\PY{n}{String}\PY{+w}{ }\PY{n}{id}\PY{p}{;}
\PY{+w}{    }\PY{n}{Bad}\PY{p}{(}\PY{n}{String}\PY{+w}{ }\PY{n}{id}\PY{p}{)}\PY{+w}{ }\PY{p}{\PYZob{}}\PY{+w}{ }\PY{k}{this}\PY{p}{.}\PY{n+na}{id}\PY{+w}{ }\PY{o}{=}\PY{+w}{ }\PY{n}{id}\PY{p}{;}\PY{+w}{ }\PY{p}{\PYZcb{}}

\PY{+w}{    }\PY{n+nd}{@Override}
\PY{+w}{    }\PY{k+kd}{public}\PY{+w}{ }\PY{k+kt}{boolean}\PY{+w}{ }\PY{n+nf}{equals}\PY{p}{(}\PY{n}{Object}\PY{+w}{ }\PY{n}{o}\PY{p}{)}\PY{+w}{ }\PY{p}{\PYZob{}}
\PY{+w}{        }\PY{k}{return}\PY{+w}{ }\PY{n}{o}\PY{+w}{ }\PY{k}{instanceof}\PY{+w}{ }\PY{n}{Bad}\PY{+w}{ }\PY{n}{b}\PY{+w}{ }\PY{o}{\PYZam{}}\PY{o}{\PYZam{}}\PY{+w}{ }\PY{n}{b}\PY{p}{.}\PY{n+na}{id}\PY{p}{.}\PY{n+na}{equals}\PY{p}{(}\PY{n}{id}\PY{p}{)}\PY{p}{;}
\PY{+w}{    }\PY{p}{\PYZcb{}}
\PY{+w}{    }\PY{c+c1}{// hashCode() NOT overridden — still identity\PYZhy{}based!}
\PY{p}{\PYZcb{}}

\PY{n}{Set}\PY{o}{\PYZlt{}}\PY{n}{Bad}\PY{o}{\PYZgt{}}\PY{+w}{ }\PY{n}{set}\PY{+w}{ }\PY{o}{=}\PY{+w}{ }\PY{k}{new}\PY{+w}{ }\PY{n}{HashSet}\PY{o}{\PYZlt{}}\PY{o}{\PYZgt{}}\PY{p}{(}\PY{p}{)}\PY{p}{;}
\PY{n}{set}\PY{p}{.}\PY{n+na}{add}\PY{p}{(}\PY{k}{new}\PY{+w}{ }\PY{n}{Bad}\PY{p}{(}\PY{l+s}{\PYZdq{}}\PY{l+s}{x}\PY{l+s}{\PYZdq{}}\PY{p}{)}\PY{p}{)}\PY{p}{;}
\PY{n}{System}\PY{p}{.}\PY{n+na}{out}\PY{p}{.}\PY{n+na}{println}\PY{p}{(}\PY{n}{set}\PY{p}{.}\PY{n+na}{contains}\PY{p}{(}\PY{k}{new}\PY{+w}{ }\PY{n}{Bad}\PY{p}{(}\PY{l+s}{\PYZdq{}}\PY{l+s}{x}\PY{l+s}{\PYZdq{}}\PY{p}{)}\PY{p}{)}\PY{p}{)}\PY{p}{;}\PY{+w}{ }\PY{c+c1}{// false! Different hash buckets.}
\end{Verbatim}
\end{codebox}

\code{HashSet} and \code{HashMap} first look at \code{hashCode()} to find the right “bucket,” then use \code{equals()} inside that bucket. If \code{hashCode()} is inconsistent with \code{equals()}, two “equal” objects can land in different buckets and the collection will never find one using the other. This is one of the most common code-review findings: \textbf{“you overrode equals() but not hashCode().”} Modern IDEs and static analyzers (and the compiler with certain lint flags) will warn about this.

\subsection[Mutable fields as HashMap keys — a classic bug]{Mutable fields as HashMap keys — a classic bug}
\label{h:5:mutable-fields-as-hashmap-keys--a-classic-bug}
Never use a mutable object as a \code{HashMap} key (or \code{HashSet} element) if its \code{hashCode()} depends on fields that can change after insertion.

\begin{codebox}
\begin{Verbatim}[commandchars=\\\{\}]
\PY{k+kd}{class} \PY{n+nc}{MutableKey}\PY{+w}{ }\PY{p}{\PYZob{}}
\PY{+w}{    }\PY{k+kd}{private}\PY{+w}{ }\PY{k+kt}{int}\PY{+w}{ }\PY{n}{value}\PY{p}{;}
\PY{+w}{    }\PY{n}{MutableKey}\PY{p}{(}\PY{k+kt}{int}\PY{+w}{ }\PY{n}{value}\PY{p}{)}\PY{+w}{ }\PY{p}{\PYZob{}}\PY{+w}{ }\PY{k}{this}\PY{p}{.}\PY{n+na}{value}\PY{+w}{ }\PY{o}{=}\PY{+w}{ }\PY{n}{value}\PY{p}{;}\PY{+w}{ }\PY{p}{\PYZcb{}}
\PY{+w}{    }\PY{k+kt}{void}\PY{+w}{ }\PY{n+nf}{setValue}\PY{p}{(}\PY{k+kt}{int}\PY{+w}{ }\PY{n}{value}\PY{p}{)}\PY{+w}{ }\PY{p}{\PYZob{}}\PY{+w}{ }\PY{k}{this}\PY{p}{.}\PY{n+na}{value}\PY{+w}{ }\PY{o}{=}\PY{+w}{ }\PY{n}{value}\PY{p}{;}\PY{+w}{ }\PY{p}{\PYZcb{}}

\PY{+w}{    }\PY{n+nd}{@Override}
\PY{+w}{    }\PY{k+kd}{public}\PY{+w}{ }\PY{k+kt}{boolean}\PY{+w}{ }\PY{n+nf}{equals}\PY{p}{(}\PY{n}{Object}\PY{+w}{ }\PY{n}{o}\PY{p}{)}\PY{+w}{ }\PY{p}{\PYZob{}}
\PY{+w}{        }\PY{k}{return}\PY{+w}{ }\PY{n}{o}\PY{+w}{ }\PY{k}{instanceof}\PY{+w}{ }\PY{n}{MutableKey}\PY{+w}{ }\PY{n}{k}\PY{+w}{ }\PY{o}{\PYZam{}}\PY{o}{\PYZam{}}\PY{+w}{ }\PY{n}{k}\PY{p}{.}\PY{n+na}{value}\PY{+w}{ }\PY{o}{=}\PY{o}{=}\PY{+w}{ }\PY{n}{value}\PY{p}{;}
\PY{+w}{    }\PY{p}{\PYZcb{}}
\PY{+w}{    }\PY{n+nd}{@Override}
\PY{+w}{    }\PY{k+kd}{public}\PY{+w}{ }\PY{k+kt}{int}\PY{+w}{ }\PY{n+nf}{hashCode}\PY{p}{(}\PY{p}{)}\PY{+w}{ }\PY{p}{\PYZob{}}
\PY{+w}{        }\PY{k}{return}\PY{+w}{ }\PY{n}{Integer}\PY{p}{.}\PY{n+na}{hashCode}\PY{p}{(}\PY{n}{value}\PY{p}{)}\PY{p}{;}
\PY{+w}{    }\PY{p}{\PYZcb{}}
\PY{p}{\PYZcb{}}

\PY{n}{Map}\PY{o}{\PYZlt{}}\PY{n}{MutableKey}\PY{p}{,}\PY{+w}{ }\PY{n}{String}\PY{o}{\PYZgt{}}\PY{+w}{ }\PY{n}{map}\PY{+w}{ }\PY{o}{=}\PY{+w}{ }\PY{k}{new}\PY{+w}{ }\PY{n}{HashMap}\PY{o}{\PYZlt{}}\PY{o}{\PYZgt{}}\PY{p}{(}\PY{p}{)}\PY{p}{;}
\PY{n}{MutableKey}\PY{+w}{ }\PY{n}{key}\PY{+w}{ }\PY{o}{=}\PY{+w}{ }\PY{k}{new}\PY{+w}{ }\PY{n}{MutableKey}\PY{p}{(}\PY{l+m+mi}{1}\PY{p}{)}\PY{p}{;}
\PY{n}{map}\PY{p}{.}\PY{n+na}{put}\PY{p}{(}\PY{n}{key}\PY{p}{,}\PY{+w}{ }\PY{l+s}{\PYZdq{}}\PY{l+s}{hello}\PY{l+s}{\PYZdq{}}\PY{p}{)}\PY{p}{;}

\PY{n}{key}\PY{p}{.}\PY{n+na}{setValue}\PY{p}{(}\PY{l+m+mi}{2}\PY{p}{)}\PY{p}{;}\PY{+w}{                 }\PY{c+c1}{// mutate the key AFTER insertion}
\PY{n}{System}\PY{p}{.}\PY{n+na}{out}\PY{p}{.}\PY{n+na}{println}\PY{p}{(}\PY{n}{map}\PY{p}{.}\PY{n+na}{get}\PY{p}{(}\PY{n}{key}\PY{p}{)}\PY{p}{)}\PY{p}{;}\PY{+w}{ }\PY{c+c1}{// null! It\PYZsq{}s stored in the bucket for hashCode(1),}
\PY{+w}{                                   }\PY{c+c1}{// but now key.hashCode() == 2.}
\end{Verbatim}
\end{codebox}

\textbf{Rule of thumb:} keys used in hash-based collections should be immutable, or at least never mutated while stored as a key. This is why \code{String}, \code{Integer}, records, and enums are popular map keys — they are immutable.

\subsection[getClass() vs instanceof — the symmetry trap with inheritance]{getClass() vs instanceof — the symmetry trap with inheritance}
\label{h:5:getclass-vs-instanceof--the-symmetry-trap-with-inheritance}
There are two common ways to type-check inside \code{equals()}, and they behave differently with subclassing:

{\tablefontsmall
\begin{xltabular}{\linewidth}{@{}L{0.481}L{0.852}L{1.667}@{}}
\toprule
\rowcolor{tablehdr}
\thd{Approach} & \thd{Behavior} & \thd{Risk} \\
\midrule
\endhead
\code{get\allowbreak{}Class()\allowbreak{}~\allowbreak{}!=\allowbreak{}~\allowbreak{}o.\allowbreak{}get\allowbreak{}Class()} & Only equal to objects of the \emph{exact same} class. & Safe from symmetry bugs, but two logically-equal instances of different subclasses are never equal. \\
\code{o~\allowbreak{}instanceof~\allowbreak{}MyClass} & Equal to any subclass instance too. & Can break \textbf{symmetry} if a subclass adds fields and also overrides \code{equals()}. \\
\bottomrule
\end{xltabular}
}

\begin{codebox}
\begin{Verbatim}[commandchars=\\\{\}]
\PY{k+kd}{class} \PY{n+nc}{Point}\PY{+w}{ }\PY{p}{\PYZob{}}
\PY{+w}{    }\PY{k+kt}{int}\PY{+w}{ }\PY{n}{x}\PY{p}{,}\PY{+w}{ }\PY{n}{y}\PY{p}{;}
\PY{+w}{    }\PY{n}{Point}\PY{p}{(}\PY{k+kt}{int}\PY{+w}{ }\PY{n}{x}\PY{p}{,}\PY{+w}{ }\PY{k+kt}{int}\PY{+w}{ }\PY{n}{y}\PY{p}{)}\PY{+w}{ }\PY{p}{\PYZob{}}\PY{+w}{ }\PY{k}{this}\PY{p}{.}\PY{n+na}{x}\PY{+w}{ }\PY{o}{=}\PY{+w}{ }\PY{n}{x}\PY{p}{;}\PY{+w}{ }\PY{k}{this}\PY{p}{.}\PY{n+na}{y}\PY{+w}{ }\PY{o}{=}\PY{+w}{ }\PY{n}{y}\PY{p}{;}\PY{+w}{ }\PY{p}{\PYZcb{}}

\PY{+w}{    }\PY{n+nd}{@Override}
\PY{+w}{    }\PY{k+kd}{public}\PY{+w}{ }\PY{k+kt}{boolean}\PY{+w}{ }\PY{n+nf}{equals}\PY{p}{(}\PY{n}{Object}\PY{+w}{ }\PY{n}{o}\PY{p}{)}\PY{+w}{ }\PY{p}{\PYZob{}}
\PY{+w}{        }\PY{k}{if}\PY{+w}{ }\PY{p}{(}\PY{o}{!}\PY{p}{(}\PY{n}{o}\PY{+w}{ }\PY{k}{instanceof}\PY{+w}{ }\PY{n}{Point}\PY{+w}{ }\PY{n}{p}\PY{p}{)}\PY{p}{)}\PY{+w}{ }\PY{k}{return}\PY{+w}{ }\PY{k+kc}{false}\PY{p}{;}\PY{+w}{ }\PY{c+c1}{// instanceof\PYZhy{}based}
\PY{+w}{        }\PY{k}{return}\PY{+w}{ }\PY{n}{x}\PY{+w}{ }\PY{o}{=}\PY{o}{=}\PY{+w}{ }\PY{n}{p}\PY{p}{.}\PY{n+na}{x}\PY{+w}{ }\PY{o}{\PYZam{}}\PY{o}{\PYZam{}}\PY{+w}{ }\PY{n}{y}\PY{+w}{ }\PY{o}{=}\PY{o}{=}\PY{+w}{ }\PY{n}{p}\PY{p}{.}\PY{n+na}{y}\PY{p}{;}
\PY{+w}{    }\PY{p}{\PYZcb{}}
\PY{p}{\PYZcb{}}

\PY{k+kd}{class} \PY{n+nc}{ColorPoint}\PY{+w}{ }\PY{k+kd}{extends}\PY{+w}{ }\PY{n}{Point}\PY{+w}{ }\PY{p}{\PYZob{}}
\PY{+w}{    }\PY{n}{String}\PY{+w}{ }\PY{n}{color}\PY{p}{;}
\PY{+w}{    }\PY{n}{ColorPoint}\PY{p}{(}\PY{k+kt}{int}\PY{+w}{ }\PY{n}{x}\PY{p}{,}\PY{+w}{ }\PY{k+kt}{int}\PY{+w}{ }\PY{n}{y}\PY{p}{,}\PY{+w}{ }\PY{n}{String}\PY{+w}{ }\PY{n}{color}\PY{p}{)}\PY{+w}{ }\PY{p}{\PYZob{}}
\PY{+w}{        }\PY{k+kd}{super}\PY{p}{(}\PY{n}{x}\PY{p}{,}\PY{+w}{ }\PY{n}{y}\PY{p}{)}\PY{p}{;}
\PY{+w}{        }\PY{k}{this}\PY{p}{.}\PY{n+na}{color}\PY{+w}{ }\PY{o}{=}\PY{+w}{ }\PY{n}{color}\PY{p}{;}
\PY{+w}{    }\PY{p}{\PYZcb{}}

\PY{+w}{    }\PY{n+nd}{@Override}
\PY{+w}{    }\PY{k+kd}{public}\PY{+w}{ }\PY{k+kt}{boolean}\PY{+w}{ }\PY{n+nf}{equals}\PY{p}{(}\PY{n}{Object}\PY{+w}{ }\PY{n}{o}\PY{p}{)}\PY{+w}{ }\PY{p}{\PYZob{}}
\PY{+w}{        }\PY{k}{if}\PY{+w}{ }\PY{p}{(}\PY{o}{!}\PY{p}{(}\PY{n}{o}\PY{+w}{ }\PY{k}{instanceof}\PY{+w}{ }\PY{n}{ColorPoint}\PY{+w}{ }\PY{n}{cp}\PY{p}{)}\PY{p}{)}\PY{+w}{ }\PY{k}{return}\PY{+w}{ }\PY{k+kc}{false}\PY{p}{;}
\PY{+w}{        }\PY{k}{return}\PY{+w}{ }\PY{k+kd}{super}\PY{p}{.}\PY{n+na}{equals}\PY{p}{(}\PY{n}{cp}\PY{p}{)}\PY{+w}{ }\PY{o}{\PYZam{}}\PY{o}{\PYZam{}}\PY{+w}{ }\PY{n}{color}\PY{p}{.}\PY{n+na}{equals}\PY{p}{(}\PY{n}{cp}\PY{p}{.}\PY{n+na}{color}\PY{p}{)}\PY{p}{;}
\PY{+w}{    }\PY{p}{\PYZcb{}}
\PY{p}{\PYZcb{}}

\PY{n}{Point}\PY{+w}{ }\PY{n}{p}\PY{+w}{ }\PY{o}{=}\PY{+w}{ }\PY{k}{new}\PY{+w}{ }\PY{n}{Point}\PY{p}{(}\PY{l+m+mi}{1}\PY{p}{,}\PY{+w}{ }\PY{l+m+mi}{2}\PY{p}{)}\PY{p}{;}
\PY{n}{ColorPoint}\PY{+w}{ }\PY{n}{cp}\PY{+w}{ }\PY{o}{=}\PY{+w}{ }\PY{k}{new}\PY{+w}{ }\PY{n}{ColorPoint}\PY{p}{(}\PY{l+m+mi}{1}\PY{p}{,}\PY{+w}{ }\PY{l+m+mi}{2}\PY{p}{,}\PY{+w}{ }\PY{l+s}{\PYZdq{}}\PY{l+s}{red}\PY{l+s}{\PYZdq{}}\PY{p}{)}\PY{p}{;}

\PY{n}{System}\PY{p}{.}\PY{n+na}{out}\PY{p}{.}\PY{n+na}{println}\PY{p}{(}\PY{n}{p}\PY{p}{.}\PY{n+na}{equals}\PY{p}{(}\PY{n}{cp}\PY{p}{)}\PY{p}{)}\PY{p}{;}\PY{+w}{  }\PY{c+c1}{// true  \PYZhy{}\PYZgt{} cp IS a Point with matching x,y}
\PY{n}{System}\PY{p}{.}\PY{n+na}{out}\PY{p}{.}\PY{n+na}{println}\PY{p}{(}\PY{n}{cp}\PY{p}{.}\PY{n+na}{equals}\PY{p}{(}\PY{n}{p}\PY{p}{)}\PY{p}{)}\PY{p}{;}\PY{+w}{  }\PY{c+c1}{// false \PYZhy{}\PYZgt{} p is not a ColorPoint}
\PY{c+c1}{// SYMMETRY BROKEN: p.equals(cp) != cp.equals(p)}
\end{Verbatim}
\end{codebox}

Effective Java’s advice (Joshua Bloch): \textbf{favor composition over inheritance for value classes}, or use \code{getClass()} equality when subclassing is possible, or make the class \code{final}. If a class is designed to be extended, consider giving it no \code{equals()}/\code{hashCode()} at all, or documenting clearly that subclasses must not add fields relevant to equality.

\begin{codebox}
\begin{Verbatim}[commandchars=\\\{\}]
\PY{c+c1}{// SAFER: getClass() check avoids the symmetry problem}
\PY{n+nd}{@Override}
\PY{k+kd}{public}\PY{+w}{ }\PY{k+kt}{boolean}\PY{+w}{ }\PY{n+nf}{equals}\PY{p}{(}\PY{n}{Object}\PY{+w}{ }\PY{n}{o}\PY{p}{)}\PY{+w}{ }\PY{p}{\PYZob{}}
\PY{+w}{    }\PY{k}{if}\PY{+w}{ }\PY{p}{(}\PY{k}{this}\PY{+w}{ }\PY{o}{=}\PY{o}{=}\PY{+w}{ }\PY{n}{o}\PY{p}{)}\PY{+w}{ }\PY{k}{return}\PY{+w}{ }\PY{k+kc}{true}\PY{p}{;}
\PY{+w}{    }\PY{k}{if}\PY{+w}{ }\PY{p}{(}\PY{n}{o}\PY{+w}{ }\PY{o}{=}\PY{o}{=}\PY{+w}{ }\PY{k+kc}{null}\PY{+w}{ }\PY{o}{|}\PY{o}{|}\PY{+w}{ }\PY{n}{getClass}\PY{p}{(}\PY{p}{)}\PY{+w}{ }\PY{o}{!}\PY{o}{=}\PY{+w}{ }\PY{n}{o}\PY{p}{.}\PY{n+na}{getClass}\PY{p}{(}\PY{p}{)}\PY{p}{)}\PY{+w}{ }\PY{k}{return}\PY{+w}{ }\PY{k+kc}{false}\PY{p}{;}
\PY{+w}{    }\PY{n}{Point}\PY{+w}{ }\PY{n}{p}\PY{+w}{ }\PY{o}{=}\PY{+w}{ }\PY{p}{(}\PY{n}{Point}\PY{p}{)}\PY{+w}{ }\PY{n}{o}\PY{p}{;}
\PY{+w}{    }\PY{k}{return}\PY{+w}{ }\PY{n}{x}\PY{+w}{ }\PY{o}{=}\PY{o}{=}\PY{+w}{ }\PY{n}{p}\PY{p}{.}\PY{n+na}{x}\PY{+w}{ }\PY{o}{\PYZam{}}\PY{o}{\PYZam{}}\PY{+w}{ }\PY{n}{y}\PY{+w}{ }\PY{o}{=}\PY{o}{=}\PY{+w}{ }\PY{n}{p}\PY{p}{.}\PY{n+na}{y}\PY{p}{;}
\PY{p}{\PYZcb{}}
\end{Verbatim}
\end{codebox}

\subsection[Objects.equals and Objects.hash — null-safe helpers]{Objects.equals and Objects.hash — null-safe helpers}
\label{h:5:objectsequals-and-objectshash--null-safe-helpers}
\code{java.\allowbreak{}util.\allowbreak{}Objects} provides small utilities that remove a lot of boilerplate and, importantly, handle \code{null} correctly:

\begin{codebox}
\begin{Verbatim}[commandchars=\\\{\}]
\PY{k+kn}{import}\PY{+w}{ }\PY{n+nn}{java.util.Objects}\PY{p}{;}

\PY{k+kd}{class} \PY{n+nc}{Person}\PY{+w}{ }\PY{p}{\PYZob{}}
\PY{+w}{    }\PY{k+kd}{private}\PY{+w}{ }\PY{k+kd}{final}\PY{+w}{ }\PY{n}{String}\PY{+w}{ }\PY{n}{name}\PY{p}{;}\PY{+w}{   }\PY{c+c1}{// could be null}
\PY{+w}{    }\PY{k+kd}{private}\PY{+w}{ }\PY{k+kd}{final}\PY{+w}{ }\PY{n}{Integer}\PY{+w}{ }\PY{n}{age}\PY{p}{;}\PY{+w}{   }\PY{c+c1}{// could be null}

\PY{+w}{    }\PY{n}{Person}\PY{p}{(}\PY{n}{String}\PY{+w}{ }\PY{n}{name}\PY{p}{,}\PY{+w}{ }\PY{n}{Integer}\PY{+w}{ }\PY{n}{age}\PY{p}{)}\PY{+w}{ }\PY{p}{\PYZob{}}
\PY{+w}{        }\PY{k}{this}\PY{p}{.}\PY{n+na}{name}\PY{+w}{ }\PY{o}{=}\PY{+w}{ }\PY{n}{name}\PY{p}{;}
\PY{+w}{        }\PY{k}{this}\PY{p}{.}\PY{n+na}{age}\PY{+w}{ }\PY{o}{=}\PY{+w}{ }\PY{n}{age}\PY{p}{;}
\PY{+w}{    }\PY{p}{\PYZcb{}}

\PY{+w}{    }\PY{n+nd}{@Override}
\PY{+w}{    }\PY{k+kd}{public}\PY{+w}{ }\PY{k+kt}{boolean}\PY{+w}{ }\PY{n+nf}{equals}\PY{p}{(}\PY{n}{Object}\PY{+w}{ }\PY{n}{o}\PY{p}{)}\PY{+w}{ }\PY{p}{\PYZob{}}
\PY{+w}{        }\PY{k}{if}\PY{+w}{ }\PY{p}{(}\PY{k}{this}\PY{+w}{ }\PY{o}{=}\PY{o}{=}\PY{+w}{ }\PY{n}{o}\PY{p}{)}\PY{+w}{ }\PY{k}{return}\PY{+w}{ }\PY{k+kc}{true}\PY{p}{;}
\PY{+w}{        }\PY{k}{if}\PY{+w}{ }\PY{p}{(}\PY{o}{!}\PY{p}{(}\PY{n}{o}\PY{+w}{ }\PY{k}{instanceof}\PY{+w}{ }\PY{n}{Person}\PY{+w}{ }\PY{n}{p}\PY{p}{)}\PY{p}{)}\PY{+w}{ }\PY{k}{return}\PY{+w}{ }\PY{k+kc}{false}\PY{p}{;}
\PY{+w}{        }\PY{k}{return}\PY{+w}{ }\PY{n}{Objects}\PY{p}{.}\PY{n+na}{equals}\PY{p}{(}\PY{n}{name}\PY{p}{,}\PY{+w}{ }\PY{n}{p}\PY{p}{.}\PY{n+na}{name}\PY{p}{)}\PY{+w}{ }\PY{o}{\PYZam{}}\PY{o}{\PYZam{}}\PY{+w}{ }\PY{n}{Objects}\PY{p}{.}\PY{n+na}{equals}\PY{p}{(}\PY{n}{age}\PY{p}{,}\PY{+w}{ }\PY{n}{p}\PY{p}{.}\PY{n+na}{age}\PY{p}{)}\PY{p}{;}
\PY{+w}{    }\PY{p}{\PYZcb{}}

\PY{+w}{    }\PY{n+nd}{@Override}
\PY{+w}{    }\PY{k+kd}{public}\PY{+w}{ }\PY{k+kt}{int}\PY{+w}{ }\PY{n+nf}{hashCode}\PY{p}{(}\PY{p}{)}\PY{+w}{ }\PY{p}{\PYZob{}}
\PY{+w}{        }\PY{k}{return}\PY{+w}{ }\PY{n}{Objects}\PY{p}{.}\PY{n+na}{hash}\PY{p}{(}\PY{n}{name}\PY{p}{,}\PY{+w}{ }\PY{n}{age}\PY{p}{)}\PY{p}{;}\PY{+w}{ }\PY{c+c1}{// null\PYZhy{}safe, handles any number of fields}
\PY{+w}{    }\PY{p}{\PYZcb{}}
\PY{p}{\PYZcb{}}
\end{Verbatim}
\end{codebox}

Without \code{Objects.\allowbreak{}equals}, comparing potentially-null fields with \code{name.\allowbreak{}equals(\allowbreak{}p.\allowbreak{}name)} risks a \code{NullPointerException}. \code{Objects.\allowbreak{}equals(\allowbreak{}a,\allowbreak{}~\allowbreak{}b)} returns \code{true} if both are \code{null}, \code{false} if only one is \code{null}, and delegates to \code{a.\allowbreak{}equals(\allowbreak{}b)} otherwise.

\subsection[Records give you equals/hashCode/toString for free]{Records give you equals/hashCode/toString for free}
\label{h:5:records-give-you-equalshashcodetostring-for-free}
Since Java 16 (standard), \code{record} types auto-generate \code{equals()}, \code{hashCode()}, and \code{toString()} based on all their components, and the generated code is correct by construction (using \code{getClass()} semantics, so it is symmetric).

\begin{codebox}
\begin{Verbatim}[commandchars=\\\{\}]
\PY{k+kd}{record} \PY{n+nc}{Point}\PY{p}{(}\PY{k+kt}{int}\PY{+w}{ }\PY{n}{x}\PY{p}{,}\PY{+w}{ }\PY{k+kt}{int}\PY{+w}{ }\PY{n}{y}\PY{p}{)}\PY{+w}{ }\PY{p}{\PYZob{}}\PY{p}{\PYZcb{}}

\PY{n}{Point}\PY{+w}{ }\PY{n}{a}\PY{+w}{ }\PY{o}{=}\PY{+w}{ }\PY{k}{new}\PY{+w}{ }\PY{n}{Point}\PY{p}{(}\PY{l+m+mi}{1}\PY{p}{,}\PY{+w}{ }\PY{l+m+mi}{2}\PY{p}{)}\PY{p}{;}
\PY{n}{Point}\PY{+w}{ }\PY{n}{b}\PY{+w}{ }\PY{o}{=}\PY{+w}{ }\PY{k}{new}\PY{+w}{ }\PY{n}{Point}\PY{p}{(}\PY{l+m+mi}{1}\PY{p}{,}\PY{+w}{ }\PY{l+m+mi}{2}\PY{p}{)}\PY{p}{;}

\PY{n}{System}\PY{p}{.}\PY{n+na}{out}\PY{p}{.}\PY{n+na}{println}\PY{p}{(}\PY{n}{a}\PY{p}{.}\PY{n+na}{equals}\PY{p}{(}\PY{n}{b}\PY{p}{)}\PY{p}{)}\PY{p}{;}\PY{+w}{ }\PY{c+c1}{// true}
\PY{n}{System}\PY{p}{.}\PY{n+na}{out}\PY{p}{.}\PY{n+na}{println}\PY{p}{(}\PY{n}{a}\PY{p}{)}\PY{p}{;}\PY{+w}{           }\PY{c+c1}{// Point[x=1, y=2]}
\PY{n}{System}\PY{p}{.}\PY{n+na}{out}\PY{p}{.}\PY{n+na}{println}\PY{p}{(}\PY{n}{a}\PY{p}{.}\PY{n+na}{hashCode}\PY{p}{(}\PY{p}{)}\PY{+w}{ }\PY{o}{=}\PY{o}{=}\PY{+w}{ }\PY{n}{b}\PY{p}{.}\PY{n+na}{hashCode}\PY{p}{(}\PY{p}{)}\PY{p}{)}\PY{p}{;}\PY{+w}{ }\PY{c+c1}{// true}
\end{Verbatim}
\end{codebox}

In a code review, if you see a hand-written data class with manually written \code{equals()}/\code{hashCode()}/\code{toString()}/getters and no extra behavior, a very reasonable comment is: \textbf{“could this be a record?”}

\subsection[IDE-generated equals/hashCode]{IDE-generated equals/hashCode}
\label{h:5:ide-generated-equalshashcode}
IntelliJ, Eclipse, and \code{lombok}’s \code{@\allowbreak{}EqualsAndHashCode} all generate correct, contract-following code — but they generate it based on a snapshot of the fields \emph{at generation time}. Common review issue: a field is added later and the developer forgets to regenerate \code{equals()}/\code{hashCode()}, so the new field is silently ignored in comparisons.

\begin{codebox}
\begin{Verbatim}[commandchars=\\\{\}]
\PY{c+c1}{// Generated when the class had only `id`. Later, `status` was added but equals() was}
\PY{c+c1}{// never regenerated:}
\PY{k+kd}{class} \PY{n+nc}{Order}\PY{+w}{ }\PY{p}{\PYZob{}}
\PY{+w}{    }\PY{k+kd}{private}\PY{+w}{ }\PY{k+kd}{final}\PY{+w}{ }\PY{n}{String}\PY{+w}{ }\PY{n}{id}\PY{p}{;}
\PY{+w}{    }\PY{k+kd}{private}\PY{+w}{ }\PY{n}{String}\PY{+w}{ }\PY{n}{status}\PY{p}{;}\PY{+w}{ }\PY{c+c1}{// added later, NOT included below}

\PY{+w}{    }\PY{n+nd}{@Override}
\PY{+w}{    }\PY{k+kd}{public}\PY{+w}{ }\PY{k+kt}{boolean}\PY{+w}{ }\PY{n+nf}{equals}\PY{p}{(}\PY{n}{Object}\PY{+w}{ }\PY{n}{o}\PY{p}{)}\PY{+w}{ }\PY{p}{\PYZob{}}
\PY{+w}{        }\PY{k}{if}\PY{+w}{ }\PY{p}{(}\PY{o}{!}\PY{p}{(}\PY{n}{o}\PY{+w}{ }\PY{k}{instanceof}\PY{+w}{ }\PY{n}{Order}\PY{+w}{ }\PY{n}{order}\PY{p}{)}\PY{p}{)}\PY{+w}{ }\PY{k}{return}\PY{+w}{ }\PY{k+kc}{false}\PY{p}{;}
\PY{+w}{        }\PY{k}{return}\PY{+w}{ }\PY{n}{Objects}\PY{p}{.}\PY{n+na}{equals}\PY{p}{(}\PY{n}{id}\PY{p}{,}\PY{+w}{ }\PY{n}{order}\PY{p}{.}\PY{n+na}{id}\PY{p}{)}\PY{p}{;}\PY{+w}{ }\PY{c+c1}{// status is silently ignored!}
\PY{+w}{    }\PY{p}{\PYZcb{}}
\PY{+w}{    }\PY{n+nd}{@Override}
\PY{+w}{    }\PY{k+kd}{public}\PY{+w}{ }\PY{k+kt}{int}\PY{+w}{ }\PY{n+nf}{hashCode}\PY{p}{(}\PY{p}{)}\PY{+w}{ }\PY{p}{\PYZob{}}
\PY{+w}{        }\PY{k}{return}\PY{+w}{ }\PY{n}{Objects}\PY{p}{.}\PY{n+na}{hash}\PY{p}{(}\PY{n}{id}\PY{p}{)}\PY{p}{;}
\PY{+w}{    }\PY{p}{\PYZcb{}}
\PY{p}{\PYZcb{}}
\end{Verbatim}
\end{codebox}

Whether that is a bug depends on intent — sometimes you \emph{want} equality by ID only (e.g., entity semantics). But it must be a deliberate choice, documented in a comment, not an oversight.

\section[Object Cloning]{Object Cloning}
\label{h:5:object-cloning}
\code{Object} declares a \code{protected} method called \code{clone()} that is supposed to create a copy of an object. In practice, it is widely considered one of Java’s worst-designed APIs, and most style guides (including Effective Java) recommend avoiding it.

\subsection[Why Cloneable is broken]{Why Cloneable is broken}
\label{h:5:why-cloneable-is-broken}
To use \code{clone()}, a class must implement the empty marker interface \code{Cloneable} — but \code{Cloneable} does not declare a \code{clone()} method itself. It just flips a hidden flag that \code{Object.\allowbreak{}clone()} checks at runtime. If you forget \code{implements~\allowbreak{}Cloneable}, calling \code{clone()} throws \code{Clone\allowbreak{}Not\allowbreak{}Supported\allowbreak{}Exception} — a checked exception, discovered only at runtime, not compile time.

\begin{codebox}
\begin{Verbatim}[commandchars=\\\{\}]
\PY{k+kd}{class} \PY{n+nc}{Box}\PY{+w}{ }\PY{k+kd}{implements}\PY{+w}{ }\PY{n}{Cloneable}\PY{+w}{ }\PY{p}{\PYZob{}}
\PY{+w}{    }\PY{k+kd}{private}\PY{+w}{ }\PY{k+kt}{int}\PY{+w}{ }\PY{n}{size}\PY{p}{;}
\PY{+w}{    }\PY{n}{Box}\PY{p}{(}\PY{k+kt}{int}\PY{+w}{ }\PY{n}{size}\PY{p}{)}\PY{+w}{ }\PY{p}{\PYZob{}}\PY{+w}{ }\PY{k}{this}\PY{p}{.}\PY{n+na}{size}\PY{+w}{ }\PY{o}{=}\PY{+w}{ }\PY{n}{size}\PY{p}{;}\PY{+w}{ }\PY{p}{\PYZcb{}}

\PY{+w}{    }\PY{n+nd}{@Override}
\PY{+w}{    }\PY{k+kd}{public}\PY{+w}{ }\PY{n}{Box}\PY{+w}{ }\PY{n+nf}{clone}\PY{p}{(}\PY{p}{)}\PY{+w}{ }\PY{p}{\PYZob{}}
\PY{+w}{        }\PY{k}{try}\PY{+w}{ }\PY{p}{\PYZob{}}
\PY{+w}{            }\PY{k}{return}\PY{+w}{ }\PY{p}{(}\PY{n}{Box}\PY{p}{)}\PY{+w}{ }\PY{k+kd}{super}\PY{p}{.}\PY{n+na}{clone}\PY{p}{(}\PY{p}{)}\PY{p}{;}\PY{+w}{ }\PY{c+c1}{// field\PYZhy{}by\PYZhy{}field shallow copy}
\PY{+w}{        }\PY{p}{\PYZcb{}}\PY{+w}{ }\PY{k}{catch}\PY{+w}{ }\PY{p}{(}\PY{n}{CloneNotSupportedException}\PY{+w}{ }\PY{n}{e}\PY{p}{)}\PY{+w}{ }\PY{p}{\PYZob{}}
\PY{+w}{            }\PY{k}{throw}\PY{+w}{ }\PY{k}{new}\PY{+w}{ }\PY{n}{AssertionError}\PY{p}{(}\PY{n}{e}\PY{p}{)}\PY{p}{;}\PY{+w}{ }\PY{c+c1}{// \PYZdq{}can\PYZsq{}t happen\PYZdq{} since we implement Cloneable}
\PY{+w}{        }\PY{p}{\PYZcb{}}
\PY{+w}{    }\PY{p}{\PYZcb{}}
\PY{p}{\PYZcb{}}
\end{Verbatim}
\end{codebox}

Problems with this design:

\begin{itemize}
\item 
\code{Cloneable}/\code{clone()} bypass constructors entirely — \code{super.\allowbreak{}clone()} does a raw memory copy, so invariants enforced in your constructor are never re-checked.

\item 
The contract is vague and unenforceable by the compiler (no \code{@\allowbreak{}Override}-style safety net for correctness).

\item 
Subclassing is fragile: if a subclass adds fields, it must remember to also override \code{clone()} correctly, and \code{final} fields cannot be “fixed up” after \code{super.\allowbreak{}clone()}.

\item 
Checked exception (\code{Clone\allowbreak{}Not\allowbreak{}Supported\allowbreak{}Exception}) makes calling code messy even though, in a correct implementation, it never actually happens.

\end{itemize}

\subsection[Shallow copy vs deep copy]{Shallow copy vs deep copy}
\label{h:5:shallow-copy-vs-deep-copy}
\code{Object.\allowbreak{}clone()} (and any naive copy) performs a \textbf{shallow copy}: primitive fields are duplicated, but reference fields (objects, arrays, collections) still point to the \emph{same} underlying object. A \textbf{deep copy} recursively copies referenced objects too, so the two copies share no mutable state.

\begin{codebox}
\begin{Verbatim}[commandchars=\\\{\}]
\PY{k+kd}{class} \PY{n+nc}{Team}\PY{+w}{ }\PY{k+kd}{implements}\PY{+w}{ }\PY{n}{Cloneable}\PY{+w}{ }\PY{p}{\PYZob{}}
\PY{+w}{    }\PY{k+kd}{private}\PY{+w}{ }\PY{n}{String}\PY{+w}{ }\PY{n}{name}\PY{p}{;}
\PY{+w}{    }\PY{k+kd}{private}\PY{+w}{ }\PY{n}{List}\PY{o}{\PYZlt{}}\PY{n}{String}\PY{o}{\PYZgt{}}\PY{+w}{ }\PY{n}{members}\PY{p}{;}\PY{+w}{ }\PY{c+c1}{// reference field}

\PY{+w}{    }\PY{n}{Team}\PY{p}{(}\PY{n}{String}\PY{+w}{ }\PY{n}{name}\PY{p}{,}\PY{+w}{ }\PY{n}{List}\PY{o}{\PYZlt{}}\PY{n}{String}\PY{o}{\PYZgt{}}\PY{+w}{ }\PY{n}{members}\PY{p}{)}\PY{+w}{ }\PY{p}{\PYZob{}}
\PY{+w}{        }\PY{k}{this}\PY{p}{.}\PY{n+na}{name}\PY{+w}{ }\PY{o}{=}\PY{+w}{ }\PY{n}{name}\PY{p}{;}
\PY{+w}{        }\PY{k}{this}\PY{p}{.}\PY{n+na}{members}\PY{+w}{ }\PY{o}{=}\PY{+w}{ }\PY{n}{members}\PY{p}{;}
\PY{+w}{    }\PY{p}{\PYZcb{}}

\PY{+w}{    }\PY{n+nd}{@Override}
\PY{+w}{    }\PY{k+kd}{public}\PY{+w}{ }\PY{n}{Team}\PY{+w}{ }\PY{n+nf}{clone}\PY{p}{(}\PY{p}{)}\PY{+w}{ }\PY{p}{\PYZob{}}
\PY{+w}{        }\PY{k}{try}\PY{+w}{ }\PY{p}{\PYZob{}}
\PY{+w}{            }\PY{n}{Team}\PY{+w}{ }\PY{n}{copy}\PY{+w}{ }\PY{o}{=}\PY{+w}{ }\PY{p}{(}\PY{n}{Team}\PY{p}{)}\PY{+w}{ }\PY{k+kd}{super}\PY{p}{.}\PY{n+na}{clone}\PY{p}{(}\PY{p}{)}\PY{p}{;}\PY{+w}{ }\PY{c+c1}{// shallow: `members` list is SHARED}
\PY{+w}{            }\PY{k}{return}\PY{+w}{ }\PY{n}{copy}\PY{p}{;}
\PY{+w}{        }\PY{p}{\PYZcb{}}\PY{+w}{ }\PY{k}{catch}\PY{+w}{ }\PY{p}{(}\PY{n}{CloneNotSupportedException}\PY{+w}{ }\PY{n}{e}\PY{p}{)}\PY{+w}{ }\PY{p}{\PYZob{}}
\PY{+w}{            }\PY{k}{throw}\PY{+w}{ }\PY{k}{new}\PY{+w}{ }\PY{n}{AssertionError}\PY{p}{(}\PY{n}{e}\PY{p}{)}\PY{p}{;}
\PY{+w}{        }\PY{p}{\PYZcb{}}
\PY{+w}{    }\PY{p}{\PYZcb{}}
\PY{p}{\PYZcb{}}

\PY{n}{List}\PY{o}{\PYZlt{}}\PY{n}{String}\PY{o}{\PYZgt{}}\PY{+w}{ }\PY{n}{members}\PY{+w}{ }\PY{o}{=}\PY{+w}{ }\PY{k}{new}\PY{+w}{ }\PY{n}{ArrayList}\PY{o}{\PYZlt{}}\PY{o}{\PYZgt{}}\PY{p}{(}\PY{n}{List}\PY{p}{.}\PY{n+na}{of}\PY{p}{(}\PY{l+s}{\PYZdq{}}\PY{l+s}{Alice}\PY{l+s}{\PYZdq{}}\PY{p}{,}\PY{+w}{ }\PY{l+s}{\PYZdq{}}\PY{l+s}{Bob}\PY{l+s}{\PYZdq{}}\PY{p}{)}\PY{p}{)}\PY{p}{;}
\PY{n}{Team}\PY{+w}{ }\PY{n}{original}\PY{+w}{ }\PY{o}{=}\PY{+w}{ }\PY{k}{new}\PY{+w}{ }\PY{n}{Team}\PY{p}{(}\PY{l+s}{\PYZdq{}}\PY{l+s}{Red}\PY{l+s}{\PYZdq{}}\PY{p}{,}\PY{+w}{ }\PY{n}{members}\PY{p}{)}\PY{p}{;}
\PY{n}{Team}\PY{+w}{ }\PY{n}{shallow}\PY{+w}{ }\PY{o}{=}\PY{+w}{ }\PY{n}{original}\PY{p}{.}\PY{n+na}{clone}\PY{p}{(}\PY{p}{)}\PY{p}{;}

\PY{n}{shallow}\PY{p}{.}\PY{n+na}{members}\PY{p}{.}\PY{n+na}{add}\PY{p}{(}\PY{l+s}{\PYZdq{}}\PY{l+s}{Eve}\PY{l+s}{\PYZdq{}}\PY{p}{)}\PY{p}{;}\PY{+w}{ }\PY{c+c1}{// mutates the list shared by BOTH teams!}
\PY{n}{System}\PY{p}{.}\PY{n+na}{out}\PY{p}{.}\PY{n+na}{println}\PY{p}{(}\PY{n}{original}\PY{p}{.}\PY{n+na}{members}\PY{p}{)}\PY{p}{;}\PY{+w}{ }\PY{c+c1}{// [Alice, Bob, Eve] — surprise mutation}
\end{Verbatim}
\end{codebox}

A correct deep copy must clone (or defensively copy) every mutable reference field:

\begin{codebox}
\begin{Verbatim}[commandchars=\\\{\}]
\PY{n+nd}{@Override}
\PY{k+kd}{public}\PY{+w}{ }\PY{n}{Team}\PY{+w}{ }\PY{n+nf}{clone}\PY{p}{(}\PY{p}{)}\PY{+w}{ }\PY{p}{\PYZob{}}
\PY{+w}{    }\PY{k}{try}\PY{+w}{ }\PY{p}{\PYZob{}}
\PY{+w}{        }\PY{n}{Team}\PY{+w}{ }\PY{n}{copy}\PY{+w}{ }\PY{o}{=}\PY{+w}{ }\PY{p}{(}\PY{n}{Team}\PY{p}{)}\PY{+w}{ }\PY{k+kd}{super}\PY{p}{.}\PY{n+na}{clone}\PY{p}{(}\PY{p}{)}\PY{p}{;}
\PY{+w}{        }\PY{n}{copy}\PY{p}{.}\PY{n+na}{members}\PY{+w}{ }\PY{o}{=}\PY{+w}{ }\PY{k}{new}\PY{+w}{ }\PY{n}{ArrayList}\PY{o}{\PYZlt{}}\PY{o}{\PYZgt{}}\PY{p}{(}\PY{k}{this}\PY{p}{.}\PY{n+na}{members}\PY{p}{)}\PY{p}{;}\PY{+w}{ }\PY{c+c1}{// deep\PYZhy{}ish copy of the list}
\PY{+w}{        }\PY{k}{return}\PY{+w}{ }\PY{n}{copy}\PY{p}{;}
\PY{+w}{    }\PY{p}{\PYZcb{}}\PY{+w}{ }\PY{k}{catch}\PY{+w}{ }\PY{p}{(}\PY{n}{CloneNotSupportedException}\PY{+w}{ }\PY{n}{e}\PY{p}{)}\PY{+w}{ }\PY{p}{\PYZob{}}
\PY{+w}{        }\PY{k}{throw}\PY{+w}{ }\PY{k}{new}\PY{+w}{ }\PY{n}{AssertionError}\PY{p}{(}\PY{n}{e}\PY{p}{)}\PY{p}{;}
\PY{+w}{    }\PY{p}{\PYZcb{}}
\PY{p}{\PYZcb{}}
\end{Verbatim}
\end{codebox}

\subsection[The better alternative: copy constructors and static factories]{The better alternative: copy constructors and static factories}
\label{h:5:the-better-alternative-copy-constructors-and-static-factories}
Most experienced reviewers push back on \code{clone()} in favor of a \textbf{copy constructor} or a \textbf{static factory method}. Both are simpler, do not throw checked exceptions, run through normal constructor logic (so invariants are validated), and work fine with \code{final} fields.

\begin{codebox}
\begin{Verbatim}[commandchars=\\\{\}]
\PY{k+kd}{final}\PY{+w}{ }\PY{k+kd}{class} \PY{n+nc}{Team}\PY{+w}{ }\PY{p}{\PYZob{}}
\PY{+w}{    }\PY{k+kd}{private}\PY{+w}{ }\PY{k+kd}{final}\PY{+w}{ }\PY{n}{String}\PY{+w}{ }\PY{n}{name}\PY{p}{;}
\PY{+w}{    }\PY{k+kd}{private}\PY{+w}{ }\PY{k+kd}{final}\PY{+w}{ }\PY{n}{List}\PY{o}{\PYZlt{}}\PY{n}{String}\PY{o}{\PYZgt{}}\PY{+w}{ }\PY{n}{members}\PY{p}{;}

\PY{+w}{    }\PY{n}{Team}\PY{p}{(}\PY{n}{String}\PY{+w}{ }\PY{n}{name}\PY{p}{,}\PY{+w}{ }\PY{n}{List}\PY{o}{\PYZlt{}}\PY{n}{String}\PY{o}{\PYZgt{}}\PY{+w}{ }\PY{n}{members}\PY{p}{)}\PY{+w}{ }\PY{p}{\PYZob{}}
\PY{+w}{        }\PY{k}{this}\PY{p}{.}\PY{n+na}{name}\PY{+w}{ }\PY{o}{=}\PY{+w}{ }\PY{n}{name}\PY{p}{;}
\PY{+w}{        }\PY{k}{this}\PY{p}{.}\PY{n+na}{members}\PY{+w}{ }\PY{o}{=}\PY{+w}{ }\PY{k}{new}\PY{+w}{ }\PY{n}{ArrayList}\PY{o}{\PYZlt{}}\PY{o}{\PYZgt{}}\PY{p}{(}\PY{n}{members}\PY{p}{)}\PY{p}{;}\PY{+w}{ }\PY{c+c1}{// defensive copy, see later section}
\PY{+w}{    }\PY{p}{\PYZcb{}}

\PY{+w}{    }\PY{c+c1}{// Copy constructor}
\PY{+w}{    }\PY{n}{Team}\PY{p}{(}\PY{n}{Team}\PY{+w}{ }\PY{n}{other}\PY{p}{)}\PY{+w}{ }\PY{p}{\PYZob{}}
\PY{+w}{        }\PY{k}{this}\PY{p}{(}\PY{n}{other}\PY{p}{.}\PY{n+na}{name}\PY{p}{,}\PY{+w}{ }\PY{n}{other}\PY{p}{.}\PY{n+na}{members}\PY{p}{)}\PY{p}{;}
\PY{+w}{    }\PY{p}{\PYZcb{}}

\PY{+w}{    }\PY{c+c1}{// Static factory alternative}
\PY{+w}{    }\PY{k+kd}{static}\PY{+w}{ }\PY{n}{Team}\PY{+w}{ }\PY{n+nf}{copyOf}\PY{p}{(}\PY{n}{Team}\PY{+w}{ }\PY{n}{other}\PY{p}{)}\PY{+w}{ }\PY{p}{\PYZob{}}
\PY{+w}{        }\PY{k}{return}\PY{+w}{ }\PY{k}{new}\PY{+w}{ }\PY{n}{Team}\PY{p}{(}\PY{n}{other}\PY{p}{.}\PY{n+na}{name}\PY{p}{,}\PY{+w}{ }\PY{n}{other}\PY{p}{.}\PY{n+na}{members}\PY{p}{)}\PY{p}{;}
\PY{+w}{    }\PY{p}{\PYZcb{}}
\PY{p}{\PYZcb{}}

\PY{n}{Team}\PY{+w}{ }\PY{n}{original}\PY{+w}{ }\PY{o}{=}\PY{+w}{ }\PY{k}{new}\PY{+w}{ }\PY{n}{Team}\PY{p}{(}\PY{l+s}{\PYZdq{}}\PY{l+s}{Red}\PY{l+s}{\PYZdq{}}\PY{p}{,}\PY{+w}{ }\PY{n}{List}\PY{p}{.}\PY{n+na}{of}\PY{p}{(}\PY{l+s}{\PYZdq{}}\PY{l+s}{Alice}\PY{l+s}{\PYZdq{}}\PY{p}{,}\PY{+w}{ }\PY{l+s}{\PYZdq{}}\PY{l+s}{Bob}\PY{l+s}{\PYZdq{}}\PY{p}{)}\PY{p}{)}\PY{p}{;}
\PY{n}{Team}\PY{+w}{ }\PY{n}{copy}\PY{+w}{ }\PY{o}{=}\PY{+w}{ }\PY{k}{new}\PY{+w}{ }\PY{n}{Team}\PY{p}{(}\PY{n}{original}\PY{p}{)}\PY{p}{;}\PY{+w}{ }\PY{c+c1}{// clean, no checked exceptions, no Cloneable dance}
\end{Verbatim}
\end{codebox}

{\tablefont
\begin{xltabular}{\linewidth}{@{}L{0.459}L{1.094}L{1.447}@{}}
\toprule
\rowcolor{tablehdr}
\thd{Approach} & \thd{Pros} & \thd{Cons} \\
\midrule
\endhead
\code{Object.\allowbreak{}clone()} / \code{Cloneable} & Built into the language & Bypasses constructors, checked exception, fragile with subclasses and \code{final} fields \\
Copy constructor & Simple, runs constructor logic, works with \code{final} fields & Must be written per class \\
Static factory (\code{copyOf}) & Same benefits, can return an interface type or cache instances & Same as above \\
\bottomrule
\end{xltabular}
}

\textbf{Interview takeaway:} if asked “how do you copy an object in Java,” the strong answer is “prefer a copy constructor or static factory; avoid \code{Cloneable} because it’s fundamentally broken.”

\section[Optional]{Optional}
\label{h:5:optional}
\code{java.\allowbreak{}util.\allowbreak{}Optional\textless{}\allowbreak{}T\textgreater{}} is a container that either holds a value or is empty. It was added in Java 8 specifically to make \textbf{absence of a value explicit in a method’s return type}, instead of relying on \code{null} and hoping callers remember to check.

\begin{codebox}
\begin{Verbatim}[commandchars=\\\{\}]
\PY{k+kn}{import}\PY{+w}{ }\PY{n+nn}{java.util.Optional}\PY{p}{;}

\PY{n}{Optional}\PY{o}{\PYZlt{}}\PY{n}{String}\PY{o}{\PYZgt{}}\PY{+w}{ }\PY{n+nf}{findUserEmail}\PY{p}{(}\PY{k+kt}{long}\PY{+w}{ }\PY{n}{userId}\PY{p}{)}\PY{+w}{ }\PY{p}{\PYZob{}}
\PY{+w}{    }\PY{n}{User}\PY{+w}{ }\PY{n}{user}\PY{+w}{ }\PY{o}{=}\PY{+w}{ }\PY{n}{repository}\PY{p}{.}\PY{n+na}{findById}\PY{p}{(}\PY{n}{userId}\PY{p}{)}\PY{p}{;}\PY{+w}{ }\PY{c+c1}{// may not exist}
\PY{+w}{    }\PY{k}{return}\PY{+w}{ }\PY{n}{user}\PY{+w}{ }\PY{o}{=}\PY{o}{=}\PY{+w}{ }\PY{k+kc}{null}\PY{+w}{ }\PY{o}{?}\PY{+w}{ }\PY{n}{Optional}\PY{p}{.}\PY{n+na}{empty}\PY{p}{(}\PY{p}{)}\PY{+w}{ }\PY{p}{:}\PY{+w}{ }\PY{n}{Optional}\PY{p}{.}\PY{n+na}{of}\PY{p}{(}\PY{n}{user}\PY{p}{.}\PY{n+na}{getEmail}\PY{p}{(}\PY{p}{)}\PY{p}{)}\PY{p}{;}
\PY{p}{\PYZcb{}}
\end{Verbatim}
\end{codebox}

Callers are nudged (though not forced) to handle the empty case instead of risking a \code{NullPointerException}:

\begin{codebox}
\begin{Verbatim}[commandchars=\\\{\}]
\PY{n}{Optional}\PY{o}{\PYZlt{}}\PY{n}{String}\PY{o}{\PYZgt{}}\PY{+w}{ }\PY{n}{email}\PY{+w}{ }\PY{o}{=}\PY{+w}{ }\PY{n}{findUserEmail}\PY{p}{(}\PY{l+m+mi}{42}\PY{p}{)}\PY{p}{;}

\PY{n}{String}\PY{+w}{ }\PY{n}{display}\PY{+w}{ }\PY{o}{=}\PY{+w}{ }\PY{n}{email}\PY{p}{.}\PY{n+na}{orElse}\PY{p}{(}\PY{l+s}{\PYZdq{}}\PY{l+s}{no email on file}\PY{l+s}{\PYZdq{}}\PY{p}{)}\PY{p}{;}
\end{Verbatim}
\end{codebox}

\subsection[map, flatMap, filter]{map, flatMap, filter}
\label{h:5:map-flatmap-filter}
\code{Optional} supports a small functional pipeline, similar to \code{Stream}:

\begin{codebox}
\begin{Verbatim}[commandchars=\\\{\}]
\PY{n}{Optional}\PY{o}{\PYZlt{}}\PY{n}{String}\PY{o}{\PYZgt{}}\PY{+w}{ }\PY{n}{email}\PY{+w}{ }\PY{o}{=}\PY{+w}{ }\PY{n}{findUserEmail}\PY{p}{(}\PY{l+m+mi}{42}\PY{p}{)}\PY{p}{;}

\PY{c+c1}{// map: transform the value if present, otherwise stay empty}
\PY{n}{Optional}\PY{o}{\PYZlt{}}\PY{n}{Integer}\PY{o}{\PYZgt{}}\PY{+w}{ }\PY{n}{domainLength}\PY{+w}{ }\PY{o}{=}\PY{+w}{ }\PY{n}{email}\PY{p}{.}\PY{n+na}{map}\PY{p}{(}\PY{n}{e}\PY{+w}{ }\PY{o}{\PYZhy{}}\PY{o}{\PYZgt{}}\PY{+w}{ }\PY{n}{e}\PY{p}{.}\PY{n+na}{split}\PY{p}{(}\PY{l+s}{\PYZdq{}}\PY{l+s}{@}\PY{l+s}{\PYZdq{}}\PY{p}{)}\PY{o}{[}\PY{l+m+mi}{1}\PY{o}{]}\PY{p}{.}\PY{n+na}{length}\PY{p}{(}\PY{p}{)}\PY{p}{)}\PY{p}{;}

\PY{c+c1}{// filter: keep the value only if it matches a predicate}
\PY{n}{Optional}\PY{o}{\PYZlt{}}\PY{n}{String}\PY{o}{\PYZgt{}}\PY{+w}{ }\PY{n}{corporateEmail}\PY{+w}{ }\PY{o}{=}\PY{+w}{ }\PY{n}{email}\PY{p}{.}\PY{n+na}{filter}\PY{p}{(}\PY{n}{e}\PY{+w}{ }\PY{o}{\PYZhy{}}\PY{o}{\PYZgt{}}\PY{+w}{ }\PY{n}{e}\PY{p}{.}\PY{n+na}{endsWith}\PY{p}{(}\PY{l+s}{\PYZdq{}}\PY{l+s}{@acme.com}\PY{l+s}{\PYZdq{}}\PY{p}{)}\PY{p}{)}\PY{p}{;}

\PY{c+c1}{// flatMap: like map, but the function itself returns an Optional (avoids Optional\PYZlt{}Optional\PYZlt{}T\PYZgt{}\PYZgt{})}
\PY{n}{Optional}\PY{o}{\PYZlt{}}\PY{n}{User}\PY{o}{\PYZgt{}}\PY{+w}{ }\PY{n+nf}{findUserByEmail}\PY{p}{(}\PY{n}{String}\PY{+w}{ }\PY{n}{email}\PY{p}{)}\PY{+w}{ }\PY{p}{\PYZob{}}\PY{+w}{ }\PY{c+cm}{/* ... */}\PY{+w}{ }\PY{k}{return}\PY{+w}{ }\PY{n}{Optional}\PY{p}{.}\PY{n+na}{empty}\PY{p}{(}\PY{p}{)}\PY{p}{;}\PY{+w}{ }\PY{p}{\PYZcb{}}

\PY{n}{Optional}\PY{o}{\PYZlt{}}\PY{n}{User}\PY{o}{\PYZgt{}}\PY{+w}{ }\PY{n}{user}\PY{+w}{ }\PY{o}{=}\PY{+w}{ }\PY{n}{email}\PY{p}{.}\PY{n+na}{flatMap}\PY{p}{(}\PY{k}{this}\PY{p}{:}\PY{p}{:}\PY{n}{findUserByEmail}\PY{p}{)}\PY{p}{;}
\end{Verbatim}
\end{codebox}

\begin{codebox}
\begin{Verbatim}[commandchars=\\\{\}]
\PY{c+c1}{// Without flatMap, map would produce a nested Optional\PYZlt{}Optional\PYZlt{}User\PYZgt{}\PYZgt{} — awkward:}
\PY{n}{Optional}\PY{o}{\PYZlt{}}\PY{n}{Optional}\PY{o}{\PYZlt{}}\PY{n}{User}\PY{o}{\PYZgt{}}\PY{o}{\PYZgt{}}\PY{+w}{ }\PY{n}{nested}\PY{+w}{ }\PY{o}{=}\PY{+w}{ }\PY{n}{email}\PY{p}{.}\PY{n+na}{map}\PY{p}{(}\PY{k}{this}\PY{p}{:}\PY{p}{:}\PY{n}{findUserByEmail}\PY{p}{)}\PY{p}{;}\PY{+w}{ }\PY{c+c1}{// wrong shape}
\end{Verbatim}
\end{codebox}

\subsection[orElse vs orElseGet vs orElseThrow]{orElse vs orElseGet vs orElseThrow}
\label{h:5:orelse-vs-orelseget-vs-orelsethrow}
These three “unwrap” an \code{Optional}, but they differ in when their argument runs — a frequent code-review gotcha:

{\tablefont
\begin{xltabular}{\linewidth}{@{}L{0.766}L{0.670}L{1.564}@{}}
\toprule
\rowcolor{tablehdr}
\thd{Method} & \thd{Argument} & \thd{When the argument runs} \\
\midrule
\endhead
\code{orElse(\allowbreak{}T~\allowbreak{}value)} & A plain value & \textbf{Always evaluated}, even if the \code{Optional} is present \\
\code{orElseGet(\allowbreak{}Supplier\textless{}\allowbreak{}T\textgreater{})} & A lazy supplier & Only called if the \code{Optional} is empty \\
\code{or\allowbreak{}Else\allowbreak{}Throw(\allowbreak{}Supplier\textless{}\allowbreak{}X\textgreater{})} & An exception supplier & Only called (and thrown) if the \code{Optional} is empty \\
\bottomrule
\end{xltabular}
}

\begin{codebox}
\begin{Verbatim}[commandchars=\\\{\}]
\PY{c+c1}{// BROKEN\PYZhy{}ish: orElse() always builds the fallback, even when not needed.}
\PY{n}{String}\PY{+w}{ }\PY{n}{email}\PY{+w}{ }\PY{o}{=}\PY{+w}{ }\PY{n}{findUserEmail}\PY{p}{(}\PY{l+m+mi}{42}\PY{p}{)}\PY{p}{.}\PY{n+na}{orElse}\PY{p}{(}\PY{n}{buildExpensiveDefault}\PY{p}{(}\PY{p}{)}\PY{p}{)}\PY{p}{;}\PY{+w}{ }\PY{c+c1}{// buildExpensiveDefault() ALWAYS runs}

\PY{c+c1}{// FIXED: orElseGet() only calls the supplier when the Optional is empty.}
\PY{n}{String}\PY{+w}{ }\PY{n}{email2}\PY{+w}{ }\PY{o}{=}\PY{+w}{ }\PY{n}{findUserEmail}\PY{p}{(}\PY{l+m+mi}{42}\PY{p}{)}\PY{p}{.}\PY{n+na}{orElseGet}\PY{p}{(}\PY{k}{this}\PY{p}{:}\PY{p}{:}\PY{n}{buildExpensiveDefault}\PY{p}{)}\PY{p}{;}
\end{Verbatim}
\end{codebox}

For a simple constant like \code{\textquotedbl{}unknown\textquotedbl{}}, \code{orElse} is fine and slightly clearer. The rule: \textbf{use \code{orElse} for cheap constants, \code{orElseGet} for anything expensive or side-effecting.}

\begin{codebox}
\begin{Verbatim}[commandchars=\\\{\}]
\PY{n}{User}\PY{+w}{ }\PY{n}{user}\PY{+w}{ }\PY{o}{=}\PY{+w}{ }\PY{n}{repository}\PY{p}{.}\PY{n+na}{findById}\PY{p}{(}\PY{l+m+mi}{42}\PY{p}{)}
\PY{+w}{        }\PY{p}{.}\PY{n+na}{orElseThrow}\PY{p}{(}\PY{p}{(}\PY{p}{)}\PY{+w}{ }\PY{o}{\PYZhy{}}\PY{o}{\PYZgt{}}\PY{+w}{ }\PY{k}{new}\PY{+w}{ }\PY{n}{NoSuchElementException}\PY{p}{(}\PY{l+s}{\PYZdq{}}\PY{l+s}{User 42 not found}\PY{l+s}{\PYZdq{}}\PY{p}{)}\PY{p}{)}\PY{p}{;}
\end{Verbatim}
\end{codebox}

\code{orElseThrow()} with no argument throws \code{NoSuchElementException}; the overload with a \code{Supplier\textless{}\allowbreak{}X\textgreater{}} lets you throw a domain-specific exception.

\subsection[ifPresent, ifPresentOrElse, and Optional in a switch/pattern context]{ifPresent, ifPresentOrElse, and Optional in a switch/pattern context}
\label{h:5:ifpresent-ifpresentorelse-and-optional-in-a-switchpattern-context}
\begin{codebox}
\begin{Verbatim}[commandchars=\\\{\}]
\PY{n}{findUserEmail}\PY{p}{(}\PY{l+m+mi}{42}\PY{p}{)}\PY{p}{.}\PY{n+na}{ifPresent}\PY{p}{(}\PY{n}{email}\PY{+w}{ }\PY{o}{\PYZhy{}}\PY{o}{\PYZgt{}}\PY{+w}{ }\PY{n}{System}\PY{p}{.}\PY{n+na}{out}\PY{p}{.}\PY{n+na}{println}\PY{p}{(}\PY{l+s}{\PYZdq{}}\PY{l+s}{Found: }\PY{l+s}{\PYZdq{}}\PY{+w}{ }\PY{o}{+}\PY{+w}{ }\PY{n}{email}\PY{p}{)}\PY{p}{)}\PY{p}{;}

\PY{n}{findUserEmail}\PY{p}{(}\PY{l+m+mi}{42}\PY{p}{)}\PY{p}{.}\PY{n+na}{ifPresentOrElse}\PY{p}{(}
\PY{+w}{        }\PY{n}{email}\PY{+w}{ }\PY{o}{\PYZhy{}}\PY{o}{\PYZgt{}}\PY{+w}{ }\PY{n}{System}\PY{p}{.}\PY{n+na}{out}\PY{p}{.}\PY{n+na}{println}\PY{p}{(}\PY{l+s}{\PYZdq{}}\PY{l+s}{Found: }\PY{l+s}{\PYZdq{}}\PY{+w}{ }\PY{o}{+}\PY{+w}{ }\PY{n}{email}\PY{p}{)}\PY{p}{,}
\PY{+w}{        }\PY{p}{(}\PY{p}{)}\PY{+w}{ }\PY{o}{\PYZhy{}}\PY{o}{\PYZgt{}}\PY{+w}{ }\PY{n}{System}\PY{p}{.}\PY{n+na}{out}\PY{p}{.}\PY{n+na}{println}\PY{p}{(}\PY{l+s}{\PYZdq{}}\PY{l+s}{No email found}\PY{l+s}{\PYZdq{}}\PY{p}{)}
\PY{p}{)}\PY{p}{;}
\end{Verbatim}
\end{codebox}

\subsection[What Optional is NOT for]{What Optional is NOT for}
\label{h:5:what-optional-is-not-for}
This is the part reviewers check most: \code{Optional} was designed for \textbf{return types}, not for fields, parameters, or collection elements.

\begin{codebox}
\begin{Verbatim}[commandchars=\\\{\}]
\PY{c+c1}{// BAD: Optional as a field — adds boxing overhead, complicates serialization,}
\PY{c+c1}{// and Optional itself is not Serializable.}
\PY{k+kd}{class} \PY{n+nc}{User}\PY{+w}{ }\PY{p}{\PYZob{}}
\PY{+w}{    }\PY{k+kd}{private}\PY{+w}{ }\PY{n}{Optional}\PY{o}{\PYZlt{}}\PY{n}{String}\PY{o}{\PYZgt{}}\PY{+w}{ }\PY{n}{nickname}\PY{p}{;}\PY{+w}{ }\PY{c+c1}{// avoid}
\PY{p}{\PYZcb{}}

\PY{c+c1}{// BAD: Optional as a method parameter — forces every caller to wrap values,}
\PY{c+c1}{// and null is still possible anyway (Optional can itself be null!).}
\PY{k+kt}{void}\PY{+w}{ }\PY{n+nf}{updateNickname}\PY{p}{(}\PY{n}{Optional}\PY{o}{\PYZlt{}}\PY{n}{String}\PY{o}{\PYZgt{}}\PY{+w}{ }\PY{n}{nickname}\PY{p}{)}\PY{+w}{ }\PY{p}{\PYZob{}}\PY{+w}{ }\PY{c+cm}{/* avoid */}\PY{+w}{ }\PY{p}{\PYZcb{}}

\PY{c+c1}{// BAD: Optional inside a collection — just omit the entry, or use an empty value.}
\PY{n}{List}\PY{o}{\PYZlt{}}\PY{n}{Optional}\PY{o}{\PYZlt{}}\PY{n}{String}\PY{o}{\PYZgt{}}\PY{o}{\PYZgt{}}\PY{+w}{ }\PY{n}{names}\PY{p}{;}\PY{+w}{ }\PY{c+c1}{// avoid}
\end{Verbatim}
\end{codebox}

Better alternatives:

\begin{codebox}
\begin{Verbatim}[commandchars=\\\{\}]
\PY{c+c1}{// GOOD: plain nullable field, document with @Nullable or just a comment.}
\PY{k+kd}{class} \PY{n+nc}{User}\PY{+w}{ }\PY{p}{\PYZob{}}
\PY{+w}{    }\PY{k+kd}{private}\PY{+w}{ }\PY{n}{String}\PY{+w}{ }\PY{n}{nickname}\PY{p}{;}\PY{+w}{ }\PY{c+c1}{// may be null}
\PY{p}{\PYZcb{}}

\PY{c+c1}{// GOOD: overload instead of Optional parameter.}
\PY{k+kt}{void}\PY{+w}{ }\PY{n+nf}{updateNickname}\PY{p}{(}\PY{n}{String}\PY{+w}{ }\PY{n}{nickname}\PY{p}{)}\PY{+w}{ }\PY{p}{\PYZob{}}\PY{+w}{ }\PY{c+cm}{/* ... */}\PY{+w}{ }\PY{p}{\PYZcb{}}

\PY{c+c1}{// GOOD: absence in a collection is just \PYZdq{}not in the list\PYZdq{} or an empty String.}
\PY{n}{List}\PY{o}{\PYZlt{}}\PY{n}{String}\PY{o}{\PYZgt{}}\PY{+w}{ }\PY{n}{names}\PY{p}{;}\PY{+w}{ }\PY{c+c1}{// simply don\PYZsq{}t add a null/absent entry}
\end{Verbatim}
\end{codebox}

\textbf{Interview one-liner:} “\code{Optional} is for values a method \emph{returns} that might not exist — never for fields, parameters, or collections. It’s a return-type tool, not a general null-replacement.”

\section[Comparator and Comparable]{Comparator and Comparable}
\label{h:5:comparator-and-comparable}
Java has two interfaces for ordering objects, and mixing them up is a common review comment.

{\tablefont
\begin{xltabular}{\linewidth}{@{}L{0.722}L{0.944}L{1.278}L{1.056}@{}}
\toprule
\rowcolor{tablehdr}
\thd{Interface} & \thd{Method} & \thd{Where the logic lives} & \thd{How many orderings?} \\
\midrule
\endhead
\code{Comparable\textless{}\allowbreak{}T\textgreater{}} & \code{compareTo(\allowbreak{}T~\allowbreak{}o)} & Inside the class itself & One “natural order” \\
\code{Comparator\textless{}\allowbreak{}T\textgreater{}} & \code{compare(\allowbreak{}T~\allowbreak{}a,\allowbreak{}~\allowbreak{}T~\allowbreak{}b)} & External, standalone & As many as you want \\
\bottomrule
\end{xltabular}
}

\subsection[The compareTo contract]{The compareTo contract}
\label{h:5:the-compareto-contract}
\code{compareTo} must return a negative number, zero, or a positive number, meaning “less than,” “equal to,” or “greater than” respectively. Like \code{equals()}, it has a formal contract:

\begin{itemize}
\item 
\textbf{Sign consistency:} \code{x.\allowbreak{}compareTo(\allowbreak{}y)} and \code{y.\allowbreak{}compareTo(\allowbreak{}x)} must have opposite signs (or both be zero).

\item 
\textbf{Transitivity:} if \code{x.\allowbreak{}compareTo(\allowbreak{}y)\allowbreak{}~\allowbreak{}\textgreater{}\allowbreak{}~\allowbreak{}0} and \code{y.\allowbreak{}compareTo(\allowbreak{}z)\allowbreak{}~\allowbreak{}\textgreater{}\allowbreak{}~\allowbreak{}0}, then \code{x.\allowbreak{}compareTo(\allowbreak{}z)\allowbreak{}~\allowbreak{}\textgreater{}\allowbreak{}~\allowbreak{}0}.

\item 
\textbf{Consistency with equals (strongly recommended, not strictly required):} \code{x.\allowbreak{}compareTo(\allowbreak{}y)\allowbreak{}~\allowbreak{}==\allowbreak{}~\allowbreak{}0} should imply \code{x.\allowbreak{}equals(\allowbreak{}y)}. If not, document it — \code{TreeSet}/\code{TreeMap} use \code{compareTo} for equality, \emph{not} \code{equals()}, so violating this causes confusing duplicate-removal behavior.

\end{itemize}

\begin{codebox}
\begin{Verbatim}[commandchars=\\\{\}]
\PY{k+kd}{final}\PY{+w}{ }\PY{k+kd}{class} \PY{n+nc}{Money}\PY{+w}{ }\PY{k+kd}{implements}\PY{+w}{ }\PY{n}{Comparable}\PY{o}{\PYZlt{}}\PY{n}{Money}\PY{o}{\PYZgt{}}\PY{+w}{ }\PY{p}{\PYZob{}}
\PY{+w}{    }\PY{k+kd}{private}\PY{+w}{ }\PY{k+kd}{final}\PY{+w}{ }\PY{k+kt}{long}\PY{+w}{ }\PY{n}{cents}\PY{p}{;}
\PY{+w}{    }\PY{n}{Money}\PY{p}{(}\PY{k+kt}{long}\PY{+w}{ }\PY{n}{cents}\PY{p}{)}\PY{+w}{ }\PY{p}{\PYZob{}}\PY{+w}{ }\PY{k}{this}\PY{p}{.}\PY{n+na}{cents}\PY{+w}{ }\PY{o}{=}\PY{+w}{ }\PY{n}{cents}\PY{p}{;}\PY{+w}{ }\PY{p}{\PYZcb{}}

\PY{+w}{    }\PY{n+nd}{@Override}
\PY{+w}{    }\PY{k+kd}{public}\PY{+w}{ }\PY{k+kt}{int}\PY{+w}{ }\PY{n+nf}{compareTo}\PY{p}{(}\PY{n}{Money}\PY{+w}{ }\PY{n}{other}\PY{p}{)}\PY{+w}{ }\PY{p}{\PYZob{}}
\PY{+w}{        }\PY{k}{return}\PY{+w}{ }\PY{n}{Long}\PY{p}{.}\PY{n+na}{compare}\PY{p}{(}\PY{k}{this}\PY{p}{.}\PY{n+na}{cents}\PY{p}{,}\PY{+w}{ }\PY{n}{other}\PY{p}{.}\PY{n+na}{cents}\PY{p}{)}\PY{p}{;}\PY{+w}{ }\PY{c+c1}{// safe, no overflow}
\PY{+w}{    }\PY{p}{\PYZcb{}}
\PY{p}{\PYZcb{}}
\end{Verbatim}
\end{codebox}

\subsection[The integer-subtraction overflow bug]{The integer-subtraction overflow bug}
\label{h:5:the-integer-subtraction-overflow-bug}
A very common — and very wrong — shortcut is subtracting two numbers to get the comparison result:

\begin{codebox}
\begin{Verbatim}[commandchars=\\\{\}]
\PY{c+c1}{// BROKEN: integer overflow can flip the sign!}
\PY{k+kd}{class} \PY{n+nc}{Score}\PY{+w}{ }\PY{k+kd}{implements}\PY{+w}{ }\PY{n}{Comparable}\PY{o}{\PYZlt{}}\PY{n}{Score}\PY{o}{\PYZgt{}}\PY{+w}{ }\PY{p}{\PYZob{}}
\PY{+w}{    }\PY{k+kd}{private}\PY{+w}{ }\PY{k+kd}{final}\PY{+w}{ }\PY{k+kt}{int}\PY{+w}{ }\PY{n}{value}\PY{p}{;}
\PY{+w}{    }\PY{n}{Score}\PY{p}{(}\PY{k+kt}{int}\PY{+w}{ }\PY{n}{value}\PY{p}{)}\PY{+w}{ }\PY{p}{\PYZob{}}\PY{+w}{ }\PY{k}{this}\PY{p}{.}\PY{n+na}{value}\PY{+w}{ }\PY{o}{=}\PY{+w}{ }\PY{n}{value}\PY{p}{;}\PY{+w}{ }\PY{p}{\PYZcb{}}

\PY{+w}{    }\PY{n+nd}{@Override}
\PY{+w}{    }\PY{k+kd}{public}\PY{+w}{ }\PY{k+kt}{int}\PY{+w}{ }\PY{n+nf}{compareTo}\PY{p}{(}\PY{n}{Score}\PY{+w}{ }\PY{n}{other}\PY{p}{)}\PY{+w}{ }\PY{p}{\PYZob{}}
\PY{+w}{        }\PY{k}{return}\PY{+w}{ }\PY{k}{this}\PY{p}{.}\PY{n+na}{value}\PY{+w}{ }\PY{o}{\PYZhy{}}\PY{+w}{ }\PY{n}{other}\PY{p}{.}\PY{n+na}{value}\PY{p}{;}\PY{+w}{ }\PY{c+c1}{// overflow risk}
\PY{+w}{    }\PY{p}{\PYZcb{}}
\PY{p}{\PYZcb{}}

\PY{n}{Score}\PY{+w}{ }\PY{n}{a}\PY{+w}{ }\PY{o}{=}\PY{+w}{ }\PY{k}{new}\PY{+w}{ }\PY{n}{Score}\PY{p}{(}\PY{n}{Integer}\PY{p}{.}\PY{n+na}{MAX\PYZus{}VALUE}\PY{p}{)}\PY{p}{;}\PY{+w}{   }\PY{c+c1}{// 2147483647}
\PY{n}{Score}\PY{+w}{ }\PY{n}{b}\PY{+w}{ }\PY{o}{=}\PY{+w}{ }\PY{k}{new}\PY{+w}{ }\PY{n}{Score}\PY{p}{(}\PY{o}{\PYZhy{}}\PY{l+m+mi}{1}\PY{p}{)}\PY{p}{;}
\PY{n}{System}\PY{p}{.}\PY{n+na}{out}\PY{p}{.}\PY{n+na}{println}\PY{p}{(}\PY{n}{a}\PY{p}{.}\PY{n+na}{compareTo}\PY{p}{(}\PY{n}{b}\PY{p}{)}\PY{p}{)}\PY{p}{;}\PY{+w}{ }\PY{c+c1}{// expected positive, but overflows to negative!}
\end{Verbatim}
\end{codebox}

Because \code{Integer.\allowbreak{}MAX\_\allowbreak{}VALUE~\allowbreak{}-\allowbreak{}~\allowbreak{}(-\allowbreak{}1)} overflows past \code{Integer.\allowbreak{}MAX\_\allowbreak{}VALUE} and wraps around to a negative number, \code{compareTo} reports \code{a~\allowbreak{}\textless{}\allowbreak{}~\allowbreak{}b} when actually \code{a~\allowbreak{}\textgreater{}\allowbreak{}~\allowbreak{}b}. This breaks sorting and binary search silently, often only on large or negative inputs — exactly the kind of bug that slips past small tests.

\begin{codebox}
\begin{Verbatim}[commandchars=\\\{\}]
\PY{c+c1}{// FIXED: use the boxed type\PYZsq{}s static compare method — never overflows.}
\PY{n+nd}{@Override}
\PY{k+kd}{public}\PY{+w}{ }\PY{k+kt}{int}\PY{+w}{ }\PY{n+nf}{compareTo}\PY{p}{(}\PY{n}{Score}\PY{+w}{ }\PY{n}{other}\PY{p}{)}\PY{+w}{ }\PY{p}{\PYZob{}}
\PY{+w}{    }\PY{k}{return}\PY{+w}{ }\PY{n}{Integer}\PY{p}{.}\PY{n+na}{compare}\PY{p}{(}\PY{k}{this}\PY{p}{.}\PY{n+na}{value}\PY{p}{,}\PY{+w}{ }\PY{n}{other}\PY{p}{.}\PY{n+na}{value}\PY{p}{)}\PY{p}{;}
\PY{p}{\PYZcb{}}
\end{Verbatim}
\end{codebox}

This applies to \code{Integer.\allowbreak{}compare}, \code{Long.\allowbreak{}compare}, \code{Double.\allowbreak{}compare}, etc. \textbf{Rule for reviewers: any \code{compareTo}/\code{compare} that uses \code{a~\allowbreak{}-\allowbreak{}~\allowbreak{}b} on \code{int}/\code{long} fields should be flagged.}

\subsection[Comparator.comparing, thenComparing, reversed]{Comparator.comparing, thenComparing, reversed}
\label{h:5:comparatorcomparing-thencomparing-reversed}
\code{Comparator} has static and default methods (since Java 8) that make building complex orderings declarative instead of hand-written \code{if}/\code{else} chains:

\begin{codebox}
\begin{Verbatim}[commandchars=\\\{\}]
\PY{k+kd}{record} \PY{n+nc}{Employee}\PY{p}{(}\PY{n}{String}\PY{+w}{ }\PY{n}{name}\PY{p}{,}\PY{+w}{ }\PY{n}{String}\PY{+w}{ }\PY{n}{department}\PY{p}{,}\PY{+w}{ }\PY{k+kt}{int}\PY{+w}{ }\PY{n}{salary}\PY{p}{)}\PY{+w}{ }\PY{p}{\PYZob{}}\PY{p}{\PYZcb{}}

\PY{n}{List}\PY{o}{\PYZlt{}}\PY{n}{Employee}\PY{o}{\PYZgt{}}\PY{+w}{ }\PY{n}{employees}\PY{+w}{ }\PY{o}{=}\PY{+w}{ }\PY{k}{new}\PY{+w}{ }\PY{n}{ArrayList}\PY{o}{\PYZlt{}}\PY{o}{\PYZgt{}}\PY{p}{(}\PY{n}{List}\PY{p}{.}\PY{n+na}{of}\PY{p}{(}
\PY{+w}{    }\PY{k}{new}\PY{+w}{ }\PY{n}{Employee}\PY{p}{(}\PY{l+s}{\PYZdq{}}\PY{l+s}{Alice}\PY{l+s}{\PYZdq{}}\PY{p}{,}\PY{+w}{ }\PY{l+s}{\PYZdq{}}\PY{l+s}{Engineering}\PY{l+s}{\PYZdq{}}\PY{p}{,}\PY{+w}{ }\PY{l+m+mi}{90000}\PY{p}{)}\PY{p}{,}
\PY{+w}{    }\PY{k}{new}\PY{+w}{ }\PY{n}{Employee}\PY{p}{(}\PY{l+s}{\PYZdq{}}\PY{l+s}{Bob}\PY{l+s}{\PYZdq{}}\PY{p}{,}\PY{+w}{   }\PY{l+s}{\PYZdq{}}\PY{l+s}{Engineering}\PY{l+s}{\PYZdq{}}\PY{p}{,}\PY{+w}{ }\PY{l+m+mi}{95000}\PY{p}{)}\PY{p}{,}
\PY{+w}{    }\PY{k}{new}\PY{+w}{ }\PY{n}{Employee}\PY{p}{(}\PY{l+s}{\PYZdq{}}\PY{l+s}{Carol}\PY{l+s}{\PYZdq{}}\PY{p}{,}\PY{+w}{ }\PY{l+s}{\PYZdq{}}\PY{l+s}{Sales}\PY{l+s}{\PYZdq{}}\PY{p}{,}\PY{+w}{       }\PY{l+m+mi}{80000}\PY{p}{)}
\PY{p}{)}\PY{p}{)}\PY{p}{;}

\PY{n}{Comparator}\PY{o}{\PYZlt{}}\PY{n}{Employee}\PY{o}{\PYZgt{}}\PY{+w}{ }\PY{n}{byDeptThenSalaryDesc}\PY{+w}{ }\PY{o}{=}
\PY{+w}{        }\PY{n}{Comparator}\PY{p}{.}\PY{n+na}{comparing}\PY{p}{(}\PY{n}{Employee}\PY{p}{:}\PY{p}{:}\PY{n}{department}\PY{p}{)}
\PY{+w}{                   }\PY{p}{.}\PY{n+na}{thenComparing}\PY{p}{(}\PY{n}{Comparator}\PY{p}{.}\PY{n+na}{comparing}\PY{p}{(}\PY{n}{Employee}\PY{p}{:}\PY{p}{:}\PY{n}{salary}\PY{p}{)}\PY{p}{.}\PY{n+na}{reversed}\PY{p}{(}\PY{p}{)}\PY{p}{)}\PY{p}{;}

\PY{n}{employees}\PY{p}{.}\PY{n+na}{sort}\PY{p}{(}\PY{n}{byDeptThenSalaryDesc}\PY{p}{)}\PY{p}{;}
\PY{c+c1}{// Engineering: Bob (95000), Alice (90000); then Sales: Carol (80000)}
\end{Verbatim}
\end{codebox}

Read it left to right: sort by department first; for ties, sort by salary, descending.

\begin{codebox}
\begin{Verbatim}[commandchars=\\\{\}]
\PY{c+c1}{// Manual equivalent — verbose and error\PYZhy{}prone by comparison:}
\PY{n}{Comparator}\PY{o}{\PYZlt{}}\PY{n}{Employee}\PY{o}{\PYZgt{}}\PY{+w}{ }\PY{n}{manual}\PY{+w}{ }\PY{o}{=}\PY{+w}{ }\PY{p}{(}\PY{n}{e1}\PY{p}{,}\PY{+w}{ }\PY{n}{e2}\PY{p}{)}\PY{+w}{ }\PY{o}{\PYZhy{}}\PY{o}{\PYZgt{}}\PY{+w}{ }\PY{p}{\PYZob{}}
\PY{+w}{    }\PY{k+kt}{int}\PY{+w}{ }\PY{n}{deptCompare}\PY{+w}{ }\PY{o}{=}\PY{+w}{ }\PY{n}{e1}\PY{p}{.}\PY{n+na}{department}\PY{p}{(}\PY{p}{)}\PY{p}{.}\PY{n+na}{compareTo}\PY{p}{(}\PY{n}{e2}\PY{p}{.}\PY{n+na}{department}\PY{p}{(}\PY{p}{)}\PY{p}{)}\PY{p}{;}
\PY{+w}{    }\PY{k}{if}\PY{+w}{ }\PY{p}{(}\PY{n}{deptCompare}\PY{+w}{ }\PY{o}{!}\PY{o}{=}\PY{+w}{ }\PY{l+m+mi}{0}\PY{p}{)}\PY{+w}{ }\PY{k}{return}\PY{+w}{ }\PY{n}{deptCompare}\PY{p}{;}
\PY{+w}{    }\PY{k}{return}\PY{+w}{ }\PY{n}{Integer}\PY{p}{.}\PY{n+na}{compare}\PY{p}{(}\PY{n}{e2}\PY{p}{.}\PY{n+na}{salary}\PY{p}{(}\PY{p}{)}\PY{p}{,}\PY{+w}{ }\PY{n}{e1}\PY{p}{.}\PY{n+na}{salary}\PY{p}{(}\PY{p}{)}\PY{p}{)}\PY{p}{;}\PY{+w}{ }\PY{c+c1}{// reversed manually}
\PY{p}{\PYZcb{}}\PY{p}{;}
\end{Verbatim}
\end{codebox}

\subsection[naturalOrder and nullsFirst / nullsLast]{naturalOrder and nullsFirst / nullsLast}
\label{h:5:naturalorder-and-nullsfirst--nullslast}
\begin{codebox}
\begin{Verbatim}[commandchars=\\\{\}]
\PY{n}{List}\PY{o}{\PYZlt{}}\PY{n}{String}\PY{o}{\PYZgt{}}\PY{+w}{ }\PY{n}{names}\PY{+w}{ }\PY{o}{=}\PY{+w}{ }\PY{k}{new}\PY{+w}{ }\PY{n}{ArrayList}\PY{o}{\PYZlt{}}\PY{o}{\PYZgt{}}\PY{p}{(}\PY{n}{List}\PY{p}{.}\PY{n+na}{of}\PY{p}{(}\PY{l+s}{\PYZdq{}}\PY{l+s}{Bob}\PY{l+s}{\PYZdq{}}\PY{p}{,}\PY{+w}{ }\PY{l+s}{\PYZdq{}}\PY{l+s}{alice}\PY{l+s}{\PYZdq{}}\PY{p}{,}\PY{+w}{ }\PY{l+s}{\PYZdq{}}\PY{l+s}{Carol}\PY{l+s}{\PYZdq{}}\PY{p}{)}\PY{p}{)}\PY{p}{;}
\PY{n}{names}\PY{p}{.}\PY{n+na}{sort}\PY{p}{(}\PY{n}{Comparator}\PY{p}{.}\PY{n+na}{naturalOrder}\PY{p}{(}\PY{p}{)}\PY{p}{)}\PY{p}{;}\PY{+w}{ }\PY{c+c1}{// uses String\PYZsq{}s own compareTo (case\PYZhy{}sensitive)}

\PY{c+c1}{// Null\PYZhy{}safe sorting: nulls can crash a naive comparator with NPE.}
\PY{n}{List}\PY{o}{\PYZlt{}}\PY{n}{String}\PY{o}{\PYZgt{}}\PY{+w}{ }\PY{n}{withNulls}\PY{+w}{ }\PY{o}{=}\PY{+w}{ }\PY{k}{new}\PY{+w}{ }\PY{n}{ArrayList}\PY{o}{\PYZlt{}}\PY{o}{\PYZgt{}}\PY{p}{(}\PY{n}{List}\PY{p}{.}\PY{n+na}{of}\PY{p}{(}\PY{l+s}{\PYZdq{}}\PY{l+s}{Bob}\PY{l+s}{\PYZdq{}}\PY{p}{,}\PY{+w}{ }\PY{k+kc}{null}\PY{p}{,}\PY{+w}{ }\PY{l+s}{\PYZdq{}}\PY{l+s}{Alice}\PY{l+s}{\PYZdq{}}\PY{p}{)}\PY{p}{)}\PY{p}{;}

\PY{c+c1}{// BROKEN: NullPointerException on the null element.}
\PY{c+c1}{// withNulls.sort(Comparator.naturalOrder());}

\PY{c+c1}{// FIXED: push nulls to the front (or use nullsLast to push them to the end).}
\PY{n}{withNulls}\PY{p}{.}\PY{n+na}{sort}\PY{p}{(}\PY{n}{Comparator}\PY{p}{.}\PY{n+na}{nullsFirst}\PY{p}{(}\PY{n}{Comparator}\PY{p}{.}\PY{n+na}{naturalOrder}\PY{p}{(}\PY{p}{)}\PY{p}{)}\PY{p}{)}\PY{p}{;}
\end{Verbatim}
\end{codebox}

\subsection[Comparable vs Comparator — when to use which]{Comparable vs Comparator — when to use which}
\label{h:5:comparable-vs-comparator--when-to-use-which}
\begin{codebox}
\begin{Verbatim}[commandchars=\\\{\}]
\PY{c+c1}{// Comparable: ONE natural ordering, baked into the class.}
\PY{k+kd}{class} \PY{n+nc}{Invoice}\PY{+w}{ }\PY{k+kd}{implements}\PY{+w}{ }\PY{n}{Comparable}\PY{o}{\PYZlt{}}\PY{n}{Invoice}\PY{o}{\PYZgt{}}\PY{+w}{ }\PY{p}{\PYZob{}}
\PY{+w}{    }\PY{k+kd}{private}\PY{+w}{ }\PY{k+kd}{final}\PY{+w}{ }\PY{n}{LocalDate}\PY{+w}{ }\PY{n}{dueDate}\PY{p}{;}
\PY{+w}{    }\PY{n+nd}{@Override}
\PY{+w}{    }\PY{k+kd}{public}\PY{+w}{ }\PY{k+kt}{int}\PY{+w}{ }\PY{n+nf}{compareTo}\PY{p}{(}\PY{n}{Invoice}\PY{+w}{ }\PY{n}{other}\PY{p}{)}\PY{+w}{ }\PY{p}{\PYZob{}}
\PY{+w}{        }\PY{k}{return}\PY{+w}{ }\PY{k}{this}\PY{p}{.}\PY{n+na}{dueDate}\PY{p}{.}\PY{n+na}{compareTo}\PY{p}{(}\PY{n}{other}\PY{p}{.}\PY{n+na}{dueDate}\PY{p}{)}\PY{p}{;}\PY{+w}{ }\PY{c+c1}{// \PYZdq{}natural\PYZdq{} order = due date}
\PY{+w}{    }\PY{p}{\PYZcb{}}
\PY{p}{\PYZcb{}}

\PY{c+c1}{// Comparator: as many EXTRA orderings as you need, defined externally.}
\PY{n}{Comparator}\PY{o}{\PYZlt{}}\PY{n}{Invoice}\PY{o}{\PYZgt{}}\PY{+w}{ }\PY{n}{byAmountDesc}\PY{+w}{ }\PY{o}{=}\PY{+w}{ }\PY{n}{Comparator}\PY{p}{.}\PY{n+na}{comparing}\PY{p}{(}\PY{n}{Invoice}\PY{p}{:}\PY{p}{:}\PY{n}{getAmount}\PY{p}{)}\PY{p}{.}\PY{n+na}{reversed}\PY{p}{(}\PY{p}{)}\PY{p}{;}
\PY{n}{Comparator}\PY{o}{\PYZlt{}}\PY{n}{Invoice}\PY{o}{\PYZgt{}}\PY{+w}{ }\PY{n}{byCustomerName}\PY{+w}{ }\PY{o}{=}\PY{+w}{ }\PY{n}{Comparator}\PY{p}{.}\PY{n+na}{comparing}\PY{p}{(}\PY{n}{Invoice}\PY{p}{:}\PY{p}{:}\PY{n}{getCustomerName}\PY{p}{)}\PY{p}{;}
\end{Verbatim}
\end{codebox}

\textbf{Interview one-liner:} “\code{Comparable} defines a type’s single natural order and lives inside the class; \code{Comparator} defines external, swappable orderings, and you can have as many as you like.”

\section[Immutability]{Immutability}
\label{h:5:immutability}
An \textbf{immutable} object is one whose observable state never changes after construction. Once built, it stays exactly as it is — no setters, no mutation, ever. Immutable objects are automatically thread-safe (no synchronization needed), easier to reason about, and safe to share and cache. \code{String}, \code{Integer}, \code{LocalDate}, and records (when built carefully) are all immutable or close to it.

\subsection[Recipe for a truly immutable class]{Recipe for a truly immutable class}
\label{h:5:recipe-for-a-truly-immutable-class}
\begin{enumerate}
\item 
Make the class \code{final} (or use only private constructors) so it cannot be subclassed to add mutable state or override behavior.

\item 
Make every field \code{private} and \code{final}.

\item 
Do not provide setters or any other mutator methods.

\item 
If a field is a mutable type (array, \code{List}, \code{Date}, mutable custom object), make a \textbf{defensive copy} on the way in (constructor) and on the way out (getter). Covered in detail in the next section.

\end{enumerate}

\begin{codebox}
\begin{Verbatim}[commandchars=\\\{\}]
\PY{k+kn}{import}\PY{+w}{ }\PY{n+nn}{java.util.List}\PY{p}{;}

\PY{k+kd}{final}\PY{+w}{ }\PY{k+kd}{class} \PY{n+nc}{ImmutableInvoice}\PY{+w}{ }\PY{p}{\PYZob{}}
\PY{+w}{    }\PY{k+kd}{private}\PY{+w}{ }\PY{k+kd}{final}\PY{+w}{ }\PY{n}{String}\PY{+w}{ }\PY{n}{id}\PY{p}{;}
\PY{+w}{    }\PY{k+kd}{private}\PY{+w}{ }\PY{k+kd}{final}\PY{+w}{ }\PY{n}{List}\PY{o}{\PYZlt{}}\PY{n}{String}\PY{o}{\PYZgt{}}\PY{+w}{ }\PY{n}{items}\PY{p}{;}
\PY{+w}{    }\PY{k+kd}{private}\PY{+w}{ }\PY{k+kd}{final}\PY{+w}{ }\PY{k+kt}{int}\PY{+w}{ }\PY{n}{totalCents}\PY{p}{;}

\PY{+w}{    }\PY{n}{ImmutableInvoice}\PY{p}{(}\PY{n}{String}\PY{+w}{ }\PY{n}{id}\PY{p}{,}\PY{+w}{ }\PY{n}{List}\PY{o}{\PYZlt{}}\PY{n}{String}\PY{o}{\PYZgt{}}\PY{+w}{ }\PY{n}{items}\PY{p}{,}\PY{+w}{ }\PY{k+kt}{int}\PY{+w}{ }\PY{n}{totalCents}\PY{p}{)}\PY{+w}{ }\PY{p}{\PYZob{}}
\PY{+w}{        }\PY{k}{this}\PY{p}{.}\PY{n+na}{id}\PY{+w}{ }\PY{o}{=}\PY{+w}{ }\PY{n}{id}\PY{p}{;}
\PY{+w}{        }\PY{k}{this}\PY{p}{.}\PY{n+na}{items}\PY{+w}{ }\PY{o}{=}\PY{+w}{ }\PY{n}{List}\PY{p}{.}\PY{n+na}{copyOf}\PY{p}{(}\PY{n}{items}\PY{p}{)}\PY{p}{;}\PY{+w}{ }\PY{c+c1}{// defensive copy + unmodifiable}
\PY{+w}{        }\PY{k}{this}\PY{p}{.}\PY{n+na}{totalCents}\PY{+w}{ }\PY{o}{=}\PY{+w}{ }\PY{n}{totalCents}\PY{p}{;}
\PY{+w}{    }\PY{p}{\PYZcb{}}

\PY{+w}{    }\PY{n}{String}\PY{+w}{ }\PY{n+nf}{getId}\PY{p}{(}\PY{p}{)}\PY{+w}{ }\PY{p}{\PYZob{}}\PY{+w}{ }\PY{k}{return}\PY{+w}{ }\PY{n}{id}\PY{p}{;}\PY{+w}{ }\PY{p}{\PYZcb{}}
\PY{+w}{    }\PY{k+kt}{int}\PY{+w}{ }\PY{n+nf}{getTotalCents}\PY{p}{(}\PY{p}{)}\PY{+w}{ }\PY{p}{\PYZob{}}\PY{+w}{ }\PY{k}{return}\PY{+w}{ }\PY{n}{totalCents}\PY{p}{;}\PY{+w}{ }\PY{p}{\PYZcb{}}
\PY{+w}{    }\PY{n}{List}\PY{o}{\PYZlt{}}\PY{n}{String}\PY{o}{\PYZgt{}}\PY{+w}{ }\PY{n+nf}{getItems}\PY{p}{(}\PY{p}{)}\PY{+w}{ }\PY{p}{\PYZob{}}\PY{+w}{ }\PY{k}{return}\PY{+w}{ }\PY{n}{items}\PY{p}{;}\PY{+w}{ }\PY{p}{\PYZcb{}}\PY{+w}{ }\PY{c+c1}{// safe: List.copyOf is already immutable}
\PY{p}{\PYZcb{}}
\end{Verbatim}
\end{codebox}

\begin{codebox}
\begin{Verbatim}[commandchars=\\\{\}]
\PY{c+c1}{// Usage: attempts to mutate the exposed list fail loudly, not silently.}
\PY{n}{ImmutableInvoice}\PY{+w}{ }\PY{n}{invoice}\PY{+w}{ }\PY{o}{=}\PY{+w}{ }\PY{k}{new}\PY{+w}{ }\PY{n}{ImmutableInvoice}\PY{p}{(}\PY{l+s}{\PYZdq{}}\PY{l+s}{INV\PYZhy{}1}\PY{l+s}{\PYZdq{}}\PY{p}{,}\PY{+w}{ }\PY{n}{List}\PY{p}{.}\PY{n+na}{of}\PY{p}{(}\PY{l+s}{\PYZdq{}}\PY{l+s}{Widget}\PY{l+s}{\PYZdq{}}\PY{p}{)}\PY{p}{,}\PY{+w}{ }\PY{l+m+mi}{999}\PY{p}{)}\PY{p}{;}
\PY{n}{invoice}\PY{p}{.}\PY{n+na}{getItems}\PY{p}{(}\PY{p}{)}\PY{p}{.}\PY{n+na}{add}\PY{p}{(}\PY{l+s}{\PYZdq{}}\PY{l+s}{Sneaky item}\PY{l+s}{\PYZdq{}}\PY{p}{)}\PY{p}{;}\PY{+w}{ }\PY{c+c1}{// throws UnsupportedOperationException}
\end{Verbatim}
\end{codebox}

\subsection[Records and immutability]{Records and immutability}
\label{h:5:records-and-immutability}
\code{record} gives you \code{private~\allowbreak{}final} fields and no setters automatically — but it does \textbf{not} automatically give you deep immutability. If a record component is a mutable type, you still need to defend it yourself:

\begin{codebox}
\begin{Verbatim}[commandchars=\\\{\}]
\PY{c+c1}{// Looks immutable, but ISN\PYZsq{}T fully:}
\PY{k+kd}{record} \PY{n+nc}{Invoice}\PY{p}{(}\PY{n}{String}\PY{+w}{ }\PY{n}{id}\PY{p}{,}\PY{+w}{ }\PY{n}{List}\PY{o}{\PYZlt{}}\PY{n}{String}\PY{o}{\PYZgt{}}\PY{+w}{ }\PY{n}{items}\PY{p}{)}\PY{+w}{ }\PY{p}{\PYZob{}}\PY{p}{\PYZcb{}}

\PY{n}{List}\PY{o}{\PYZlt{}}\PY{n}{String}\PY{o}{\PYZgt{}}\PY{+w}{ }\PY{n}{mutable}\PY{+w}{ }\PY{o}{=}\PY{+w}{ }\PY{k}{new}\PY{+w}{ }\PY{n}{ArrayList}\PY{o}{\PYZlt{}}\PY{o}{\PYZgt{}}\PY{p}{(}\PY{n}{List}\PY{p}{.}\PY{n+na}{of}\PY{p}{(}\PY{l+s}{\PYZdq{}}\PY{l+s}{Widget}\PY{l+s}{\PYZdq{}}\PY{p}{)}\PY{p}{)}\PY{p}{;}
\PY{n}{Invoice}\PY{+w}{ }\PY{n}{invoice}\PY{+w}{ }\PY{o}{=}\PY{+w}{ }\PY{k}{new}\PY{+w}{ }\PY{n}{Invoice}\PY{p}{(}\PY{l+s}{\PYZdq{}}\PY{l+s}{INV\PYZhy{}1}\PY{l+s}{\PYZdq{}}\PY{p}{,}\PY{+w}{ }\PY{n}{mutable}\PY{p}{)}\PY{p}{;}
\PY{n}{mutable}\PY{p}{.}\PY{n+na}{add}\PY{p}{(}\PY{l+s}{\PYZdq{}}\PY{l+s}{Sneaky item}\PY{l+s}{\PYZdq{}}\PY{p}{)}\PY{p}{;}\PY{+w}{ }\PY{c+c1}{// mutates the record\PYZsq{}s internal state from outside!}
\PY{n}{System}\PY{p}{.}\PY{n+na}{out}\PY{p}{.}\PY{n+na}{println}\PY{p}{(}\PY{n}{invoice}\PY{p}{.}\PY{n+na}{items}\PY{p}{(}\PY{p}{)}\PY{p}{)}\PY{p}{;}\PY{+w}{ }\PY{c+c1}{// [Widget, Sneaky item]}
\end{Verbatim}
\end{codebox}

\begin{codebox}
\begin{Verbatim}[commandchars=\\\{\}]
\PY{c+c1}{// FIXED: use a compact constructor to defensively copy.}
\PY{k+kd}{record} \PY{n+nc}{Invoice}\PY{p}{(}\PY{n}{String}\PY{+w}{ }\PY{n}{id}\PY{p}{,}\PY{+w}{ }\PY{n}{List}\PY{o}{\PYZlt{}}\PY{n}{String}\PY{o}{\PYZgt{}}\PY{+w}{ }\PY{n}{items}\PY{p}{)}\PY{+w}{ }\PY{p}{\PYZob{}}
\PY{+w}{    }\PY{n}{Invoice}\PY{p}{(}\PY{n}{String}\PY{+w}{ }\PY{n}{id}\PY{p}{,}\PY{+w}{ }\PY{n}{List}\PY{o}{\PYZlt{}}\PY{n}{String}\PY{o}{\PYZgt{}}\PY{+w}{ }\PY{n}{items}\PY{p}{)}\PY{+w}{ }\PY{p}{\PYZob{}}
\PY{+w}{        }\PY{k}{this}\PY{p}{.}\PY{n+na}{id}\PY{+w}{ }\PY{o}{=}\PY{+w}{ }\PY{n}{id}\PY{p}{;}
\PY{+w}{        }\PY{k}{this}\PY{p}{.}\PY{n+na}{items}\PY{+w}{ }\PY{o}{=}\PY{+w}{ }\PY{n}{List}\PY{p}{.}\PY{n+na}{copyOf}\PY{p}{(}\PY{n}{items}\PY{p}{)}\PY{p}{;}\PY{+w}{ }\PY{c+c1}{// defends against caller mutation}
\PY{+w}{    }\PY{p}{\PYZcb{}}
\PY{p}{\PYZcb{}}
\end{Verbatim}
\end{codebox}

\subsection[Immutability and performance / thread-safety trade-offs]{Immutability and performance / thread-safety trade-offs}
\label{h:5:immutability-and-performance--thread-safety-trade-offs}
{\tablefont
\begin{xltabular}{\linewidth}{@{}L{0.841}L{1.159}@{}}
\toprule
\rowcolor{tablehdr}
\thd{Benefit} & \thd{Cost} \\
\midrule
\endhead
Naturally thread-safe, no locks needed & Every “change” allocates a new object \\
Safe to cache, share, use as map keys & Can be wasteful for large objects mutated frequently in a loop \\
Simpler to reason about, fewer invariant bugs & Requires discipline (defensive copies) for every mutable field \\
\bottomrule
\end{xltabular}
}

For a value that changes frequently inside a hot loop (e.g., a running total), a mutable local variable or builder is usually fine — immutability shines for objects that are shared across threads or passed around as stable values (DTOs, config, keys, events).

\section[Defensive Copying]{Defensive Copying}
\label{h:5:defensive-copying}
\textbf{Defensive copying} means making an independent copy of a mutable object at a trust boundary, so that neither side can affect the other through shared references. Without it, “immutable” classes can be mutated indirectly by whoever handed you the object — or whoever you handed the object to.

\subsection[The two places you need it: constructor AND getter]{The two places you need it: constructor AND getter}
\label{h:5:the-two-places-you-need-it-constructor-and-getter}
A very common code-review finding is defending only \emph{one} side (usually the constructor) and forgetting the other (the getter).

\begin{codebox}
\begin{Verbatim}[commandchars=\\\{\}]
\PY{k+kn}{import}\PY{+w}{ }\PY{n+nn}{java.util.ArrayList}\PY{p}{;}
\PY{k+kn}{import}\PY{+w}{ }\PY{n+nn}{java.util.Date}\PY{p}{;}
\PY{k+kn}{import}\PY{+w}{ }\PY{n+nn}{java.util.List}\PY{p}{;}

\PY{c+c1}{// BROKEN: no defensive copies anywhere.}
\PY{k+kd}{final}\PY{+w}{ }\PY{k+kd}{class} \PY{n+nc}{Trip}\PY{+w}{ }\PY{p}{\PYZob{}}
\PY{+w}{    }\PY{k+kd}{private}\PY{+w}{ }\PY{k+kd}{final}\PY{+w}{ }\PY{n}{Date}\PY{+w}{ }\PY{n}{startDate}\PY{p}{;}\PY{+w}{   }\PY{c+c1}{// Date is mutable!}
\PY{+w}{    }\PY{k+kd}{private}\PY{+w}{ }\PY{k+kd}{final}\PY{+w}{ }\PY{n}{List}\PY{o}{\PYZlt{}}\PY{n}{String}\PY{o}{\PYZgt{}}\PY{+w}{ }\PY{n}{stops}\PY{p}{;}

\PY{+w}{    }\PY{n}{Trip}\PY{p}{(}\PY{n}{Date}\PY{+w}{ }\PY{n}{startDate}\PY{p}{,}\PY{+w}{ }\PY{n}{List}\PY{o}{\PYZlt{}}\PY{n}{String}\PY{o}{\PYZgt{}}\PY{+w}{ }\PY{n}{stops}\PY{p}{)}\PY{+w}{ }\PY{p}{\PYZob{}}
\PY{+w}{        }\PY{k}{this}\PY{p}{.}\PY{n+na}{startDate}\PY{+w}{ }\PY{o}{=}\PY{+w}{ }\PY{n}{startDate}\PY{p}{;}\PY{+w}{ }\PY{c+c1}{// caller\PYZsq{}s Date object is stored directly}
\PY{+w}{        }\PY{k}{this}\PY{p}{.}\PY{n+na}{stops}\PY{+w}{ }\PY{o}{=}\PY{+w}{ }\PY{n}{stops}\PY{p}{;}\PY{+w}{         }\PY{c+c1}{// caller\PYZsq{}s List is stored directly}
\PY{+w}{    }\PY{p}{\PYZcb{}}

\PY{+w}{    }\PY{n}{Date}\PY{+w}{ }\PY{n+nf}{getStartDate}\PY{p}{(}\PY{p}{)}\PY{+w}{ }\PY{p}{\PYZob{}}\PY{+w}{ }\PY{k}{return}\PY{+w}{ }\PY{n}{startDate}\PY{p}{;}\PY{+w}{ }\PY{p}{\PYZcb{}}\PY{+w}{ }\PY{c+c1}{// returns the SAME mutable Date}
\PY{+w}{    }\PY{n}{List}\PY{o}{\PYZlt{}}\PY{n}{String}\PY{o}{\PYZgt{}}\PY{+w}{ }\PY{n+nf}{getStops}\PY{p}{(}\PY{p}{)}\PY{+w}{ }\PY{p}{\PYZob{}}\PY{+w}{ }\PY{k}{return}\PY{+w}{ }\PY{n}{stops}\PY{p}{;}\PY{+w}{ }\PY{p}{\PYZcb{}}\PY{+w}{  }\PY{c+c1}{// returns the SAME mutable List}
\PY{p}{\PYZcb{}}

\PY{n}{Date}\PY{+w}{ }\PY{n}{d}\PY{+w}{ }\PY{o}{=}\PY{+w}{ }\PY{k}{new}\PY{+w}{ }\PY{n}{Date}\PY{p}{(}\PY{p}{)}\PY{p}{;}
\PY{n}{List}\PY{o}{\PYZlt{}}\PY{n}{String}\PY{o}{\PYZgt{}}\PY{+w}{ }\PY{n}{stops}\PY{+w}{ }\PY{o}{=}\PY{+w}{ }\PY{k}{new}\PY{+w}{ }\PY{n}{ArrayList}\PY{o}{\PYZlt{}}\PY{o}{\PYZgt{}}\PY{p}{(}\PY{n}{List}\PY{p}{.}\PY{n+na}{of}\PY{p}{(}\PY{l+s}{\PYZdq{}}\PY{l+s}{Paris}\PY{l+s}{\PYZdq{}}\PY{p}{)}\PY{p}{)}\PY{p}{;}
\PY{n}{Trip}\PY{+w}{ }\PY{n}{trip}\PY{+w}{ }\PY{o}{=}\PY{+w}{ }\PY{k}{new}\PY{+w}{ }\PY{n}{Trip}\PY{p}{(}\PY{n}{d}\PY{p}{,}\PY{+w}{ }\PY{n}{stops}\PY{p}{)}\PY{p}{;}

\PY{n}{d}\PY{p}{.}\PY{n+na}{setTime}\PY{p}{(}\PY{l+m+mi}{0}\PY{p}{)}\PY{p}{;}\PY{+w}{              }\PY{c+c1}{// mutates the Trip\PYZsq{}s startDate from OUTSIDE}
\PY{n}{stops}\PY{p}{.}\PY{n+na}{add}\PY{p}{(}\PY{l+s}{\PYZdq{}}\PY{l+s}{Berlin}\PY{l+s}{\PYZdq{}}\PY{p}{)}\PY{p}{;}\PY{+w}{       }\PY{c+c1}{// mutates the Trip\PYZsq{}s stops from OUTSIDE}

\PY{n}{trip}\PY{p}{.}\PY{n+na}{getStartDate}\PY{p}{(}\PY{p}{)}\PY{p}{.}\PY{n+na}{setTime}\PY{p}{(}\PY{l+m+mi}{123456789L}\PY{p}{)}\PY{p}{;}\PY{+w}{ }\PY{c+c1}{// mutates the Trip\PYZsq{}s internal Date via the GETTER}
\end{Verbatim}
\end{codebox}

\begin{codebox}
\begin{Verbatim}[commandchars=\\\{\}]
\PY{c+c1}{// FIXED: defend both directions.}
\PY{k+kd}{final}\PY{+w}{ }\PY{k+kd}{class} \PY{n+nc}{Trip}\PY{+w}{ }\PY{p}{\PYZob{}}
\PY{+w}{    }\PY{k+kd}{private}\PY{+w}{ }\PY{k+kd}{final}\PY{+w}{ }\PY{n}{Date}\PY{+w}{ }\PY{n}{startDate}\PY{p}{;}
\PY{+w}{    }\PY{k+kd}{private}\PY{+w}{ }\PY{k+kd}{final}\PY{+w}{ }\PY{n}{List}\PY{o}{\PYZlt{}}\PY{n}{String}\PY{o}{\PYZgt{}}\PY{+w}{ }\PY{n}{stops}\PY{p}{;}

\PY{+w}{    }\PY{n}{Trip}\PY{p}{(}\PY{n}{Date}\PY{+w}{ }\PY{n}{startDate}\PY{p}{,}\PY{+w}{ }\PY{n}{List}\PY{o}{\PYZlt{}}\PY{n}{String}\PY{o}{\PYZgt{}}\PY{+w}{ }\PY{n}{stops}\PY{p}{)}\PY{+w}{ }\PY{p}{\PYZob{}}
\PY{+w}{        }\PY{k}{this}\PY{p}{.}\PY{n+na}{startDate}\PY{+w}{ }\PY{o}{=}\PY{+w}{ }\PY{k}{new}\PY{+w}{ }\PY{n}{Date}\PY{p}{(}\PY{n}{startDate}\PY{p}{.}\PY{n+na}{getTime}\PY{p}{(}\PY{p}{)}\PY{p}{)}\PY{p}{;}\PY{+w}{ }\PY{c+c1}{// copy IN}
\PY{+w}{        }\PY{k}{this}\PY{p}{.}\PY{n+na}{stops}\PY{+w}{ }\PY{o}{=}\PY{+w}{ }\PY{k}{new}\PY{+w}{ }\PY{n}{ArrayList}\PY{o}{\PYZlt{}}\PY{o}{\PYZgt{}}\PY{p}{(}\PY{n}{stops}\PY{p}{)}\PY{p}{;}\PY{+w}{              }\PY{c+c1}{// copy IN}
\PY{+w}{    }\PY{p}{\PYZcb{}}

\PY{+w}{    }\PY{n}{Date}\PY{+w}{ }\PY{n+nf}{getStartDate}\PY{p}{(}\PY{p}{)}\PY{+w}{ }\PY{p}{\PYZob{}}
\PY{+w}{        }\PY{k}{return}\PY{+w}{ }\PY{k}{new}\PY{+w}{ }\PY{n}{Date}\PY{p}{(}\PY{n}{startDate}\PY{p}{.}\PY{n+na}{getTime}\PY{p}{(}\PY{p}{)}\PY{p}{)}\PY{p}{;}\PY{+w}{ }\PY{c+c1}{// copy OUT}
\PY{+w}{    }\PY{p}{\PYZcb{}}

\PY{+w}{    }\PY{n}{List}\PY{o}{\PYZlt{}}\PY{n}{String}\PY{o}{\PYZgt{}}\PY{+w}{ }\PY{n+nf}{getStops}\PY{p}{(}\PY{p}{)}\PY{+w}{ }\PY{p}{\PYZob{}}
\PY{+w}{        }\PY{k}{return}\PY{+w}{ }\PY{n}{List}\PY{p}{.}\PY{n+na}{copyOf}\PY{p}{(}\PY{n}{stops}\PY{p}{)}\PY{p}{;}\PY{+w}{ }\PY{c+c1}{// copy OUT, and unmodifiable too}
\PY{+w}{    }\PY{p}{\PYZcb{}}
\PY{p}{\PYZcb{}}
\end{Verbatim}
\end{codebox}

Now, mutating the caller’s original \code{Date}/\code{List} after construction — or mutating whatever the getter returns — has zero effect on the \code{Trip}’s internal state.

\subsection[List.copyOf vs unmodifiable views vs true immutable copies]{List.copyOf vs unmodifiable views vs true immutable copies}
\label{h:5:listcopyof-vs-unmodifiable-views-vs-true-immutable-copies}
These three look similar but behave very differently, and mixing them up is a frequent bug:

{\tablefont
\begin{xltabular}{\linewidth}{@{}L{0.739}L{1.413}L{0.848}@{}}
\toprule
\rowcolor{tablehdr}
\thd{Technique} & \thd{Independent from source?} & \thd{Can it be mutated directly?} \\
\midrule
\endhead
\code{Collections.\allowbreak{}unmodifiable\allowbreak{}List(\allowbreak{}list)} & No — it’s a \textbf{view}; changes to the original \code{list} still show through & No (throws on direct mutation attempts) \\
\code{List.\allowbreak{}copyOf(\allowbreak{}list)} & Yes — makes an actual independent copy & No (immutable) \\
\code{new~\allowbreak{}ArrayList\textless{}\textgreater{}(\allowbreak{}list)} & Yes — independent copy & Yes (still mutable) \\
\bottomrule
\end{xltabular}
}

\begin{codebox}
\begin{Verbatim}[commandchars=\\\{\}]
\PY{n}{List}\PY{o}{\PYZlt{}}\PY{n}{String}\PY{o}{\PYZgt{}}\PY{+w}{ }\PY{n}{source}\PY{+w}{ }\PY{o}{=}\PY{+w}{ }\PY{k}{new}\PY{+w}{ }\PY{n}{ArrayList}\PY{o}{\PYZlt{}}\PY{o}{\PYZgt{}}\PY{p}{(}\PY{n}{List}\PY{p}{.}\PY{n+na}{of}\PY{p}{(}\PY{l+s}{\PYZdq{}}\PY{l+s}{A}\PY{l+s}{\PYZdq{}}\PY{p}{,}\PY{+w}{ }\PY{l+s}{\PYZdq{}}\PY{l+s}{B}\PY{l+s}{\PYZdq{}}\PY{p}{)}\PY{p}{)}\PY{p}{;}

\PY{n}{List}\PY{o}{\PYZlt{}}\PY{n}{String}\PY{o}{\PYZgt{}}\PY{+w}{ }\PY{n}{view}\PY{+w}{ }\PY{o}{=}\PY{+w}{ }\PY{n}{Collections}\PY{p}{.}\PY{n+na}{unmodifiableList}\PY{p}{(}\PY{n}{source}\PY{p}{)}\PY{p}{;}
\PY{n}{source}\PY{p}{.}\PY{n+na}{add}\PY{p}{(}\PY{l+s}{\PYZdq{}}\PY{l+s}{C}\PY{l+s}{\PYZdq{}}\PY{p}{)}\PY{p}{;}
\PY{n}{System}\PY{p}{.}\PY{n+na}{out}\PY{p}{.}\PY{n+na}{println}\PY{p}{(}\PY{n}{view}\PY{p}{)}\PY{p}{;}\PY{+w}{ }\PY{c+c1}{// [A, B, C] — the view is NOT independent, it just blocks writes}

\PY{n}{List}\PY{o}{\PYZlt{}}\PY{n}{String}\PY{o}{\PYZgt{}}\PY{+w}{ }\PY{n}{copy}\PY{+w}{ }\PY{o}{=}\PY{+w}{ }\PY{n}{List}\PY{p}{.}\PY{n+na}{copyOf}\PY{p}{(}\PY{n}{source}\PY{p}{)}\PY{p}{;}
\PY{n}{source}\PY{p}{.}\PY{n+na}{add}\PY{p}{(}\PY{l+s}{\PYZdq{}}\PY{l+s}{D}\PY{l+s}{\PYZdq{}}\PY{p}{)}\PY{p}{;}
\PY{n}{System}\PY{p}{.}\PY{n+na}{out}\PY{p}{.}\PY{n+na}{println}\PY{p}{(}\PY{n}{copy}\PY{p}{)}\PY{p}{;}\PY{+w}{ }\PY{c+c1}{// [A, B, C] — the copy IS independent}
\end{Verbatim}
\end{codebox}

\begin{codebox}
\begin{Verbatim}[commandchars=\\\{\}]
\PY{c+c1}{// BROKEN: developer thinks unmodifiableList() protects against upstream mutation.}
\PY{k+kd}{class} \PY{n+nc}{Roster}\PY{+w}{ }\PY{p}{\PYZob{}}
\PY{+w}{    }\PY{k+kd}{private}\PY{+w}{ }\PY{k+kd}{final}\PY{+w}{ }\PY{n}{List}\PY{o}{\PYZlt{}}\PY{n}{String}\PY{o}{\PYZgt{}}\PY{+w}{ }\PY{n}{names}\PY{p}{;}
\PY{+w}{    }\PY{n}{Roster}\PY{p}{(}\PY{n}{List}\PY{o}{\PYZlt{}}\PY{n}{String}\PY{o}{\PYZgt{}}\PY{+w}{ }\PY{n}{names}\PY{p}{)}\PY{+w}{ }\PY{p}{\PYZob{}}
\PY{+w}{        }\PY{k}{this}\PY{p}{.}\PY{n+na}{names}\PY{+w}{ }\PY{o}{=}\PY{+w}{ }\PY{n}{Collections}\PY{p}{.}\PY{n+na}{unmodifiableList}\PY{p}{(}\PY{n}{names}\PY{p}{)}\PY{p}{;}\PY{+w}{ }\PY{c+c1}{// still a VIEW over the caller\PYZsq{}s list!}
\PY{+w}{    }\PY{p}{\PYZcb{}}
\PY{p}{\PYZcb{}}

\PY{n}{List}\PY{o}{\PYZlt{}}\PY{n}{String}\PY{o}{\PYZgt{}}\PY{+w}{ }\PY{n}{mutable}\PY{+w}{ }\PY{o}{=}\PY{+w}{ }\PY{k}{new}\PY{+w}{ }\PY{n}{ArrayList}\PY{o}{\PYZlt{}}\PY{o}{\PYZgt{}}\PY{p}{(}\PY{n}{List}\PY{p}{.}\PY{n+na}{of}\PY{p}{(}\PY{l+s}{\PYZdq{}}\PY{l+s}{Alice}\PY{l+s}{\PYZdq{}}\PY{p}{)}\PY{p}{)}\PY{p}{;}
\PY{n}{Roster}\PY{+w}{ }\PY{n}{roster}\PY{+w}{ }\PY{o}{=}\PY{+w}{ }\PY{k}{new}\PY{+w}{ }\PY{n}{Roster}\PY{p}{(}\PY{n}{mutable}\PY{p}{)}\PY{p}{;}
\PY{n}{mutable}\PY{p}{.}\PY{n+na}{add}\PY{p}{(}\PY{l+s}{\PYZdq{}}\PY{l+s}{Mallory}\PY{l+s}{\PYZdq{}}\PY{p}{)}\PY{p}{;}\PY{+w}{ }\PY{c+c1}{// Roster\PYZsq{}s \PYZdq{}immutable\PYZdq{} view now includes \PYZdq{}Mallory\PYZdq{} too!}
\end{Verbatim}
\end{codebox}

\begin{codebox}
\begin{Verbatim}[commandchars=\\\{\}]
\PY{c+c1}{// FIXED: copy first, THEN wrap (or just use List.copyOf, which does both).}
\PY{k+kd}{class} \PY{n+nc}{Roster}\PY{+w}{ }\PY{p}{\PYZob{}}
\PY{+w}{    }\PY{k+kd}{private}\PY{+w}{ }\PY{k+kd}{final}\PY{+w}{ }\PY{n}{List}\PY{o}{\PYZlt{}}\PY{n}{String}\PY{o}{\PYZgt{}}\PY{+w}{ }\PY{n}{names}\PY{p}{;}
\PY{+w}{    }\PY{n}{Roster}\PY{p}{(}\PY{n}{List}\PY{o}{\PYZlt{}}\PY{n}{String}\PY{o}{\PYZgt{}}\PY{+w}{ }\PY{n}{names}\PY{p}{)}\PY{+w}{ }\PY{p}{\PYZob{}}
\PY{+w}{        }\PY{k}{this}\PY{p}{.}\PY{n+na}{names}\PY{+w}{ }\PY{o}{=}\PY{+w}{ }\PY{n}{List}\PY{p}{.}\PY{n+na}{copyOf}\PY{p}{(}\PY{n}{names}\PY{p}{)}\PY{p}{;}\PY{+w}{ }\PY{c+c1}{// independent AND unmodifiable}
\PY{+w}{    }\PY{p}{\PYZcb{}}
\PY{p}{\PYZcb{}}
\end{Verbatim}
\end{codebox}

\textbf{Rule of thumb:} \code{unmodifiableXxx()} prevents \emph{your} code (or the caller of your getter) from mutating a collection through that particular reference, but it does nothing to stop the \emph{original owner} of the backing collection from mutating it. Use it for a getter when you already own an independent copy; use \code{copyOf} (or a manual copy) when you need true independence.

\subsection[Arrays are always a leak vector]{Arrays are always a leak vector}
\label{h:5:arrays-are-always-a-leak-vector}
Arrays are mutable in Java, and unlike \code{List}, there is no built-in “unmodifiable array” wrapper. Any method that returns or accepts an array is a defensive-copying hotspot.

\begin{codebox}
\begin{Verbatim}[commandchars=\\\{\}]
\PY{c+c1}{// BROKEN: exposes the internal array directly.}
\PY{k+kd}{class} \PY{n+nc}{PasswordPolicy}\PY{+w}{ }\PY{p}{\PYZob{}}
\PY{+w}{    }\PY{k+kd}{private}\PY{+w}{ }\PY{k+kd}{final}\PY{+w}{ }\PY{k+kt}{char}\PY{o}{[}\PY{o}{]}\PY{+w}{ }\PY{n}{allowedSymbols}\PY{p}{;}
\PY{+w}{    }\PY{n}{PasswordPolicy}\PY{p}{(}\PY{k+kt}{char}\PY{o}{[}\PY{o}{]}\PY{+w}{ }\PY{n}{allowedSymbols}\PY{p}{)}\PY{+w}{ }\PY{p}{\PYZob{}}
\PY{+w}{        }\PY{k}{this}\PY{p}{.}\PY{n+na}{allowedSymbols}\PY{+w}{ }\PY{o}{=}\PY{+w}{ }\PY{n}{allowedSymbols}\PY{p}{;}\PY{+w}{ }\PY{c+c1}{// stores caller\PYZsq{}s array directly}
\PY{+w}{    }\PY{p}{\PYZcb{}}
\PY{+w}{    }\PY{k+kt}{char}\PY{o}{[}\PY{o}{]}\PY{+w}{ }\PY{n+nf}{getAllowedSymbols}\PY{p}{(}\PY{p}{)}\PY{+w}{ }\PY{p}{\PYZob{}}\PY{+w}{ }\PY{k}{return}\PY{+w}{ }\PY{n}{allowedSymbols}\PY{p}{;}\PY{+w}{ }\PY{p}{\PYZcb{}}\PY{+w}{ }\PY{c+c1}{// returns internal array directly}
\PY{p}{\PYZcb{}}

\PY{k+kt}{char}\PY{o}{[}\PY{o}{]}\PY{+w}{ }\PY{n}{symbols}\PY{+w}{ }\PY{o}{=}\PY{+w}{ }\PY{p}{\PYZob{}}\PY{l+s+sc}{\PYZsq{}!\PYZsq{}}\PY{p}{,}\PY{+w}{ }\PY{l+s+sc}{\PYZsq{}@\PYZsq{}}\PY{p}{,}\PY{+w}{ }\PY{l+s+sc}{\PYZsq{}\PYZsh{}\PYZsq{}}\PY{p}{\PYZcb{}}\PY{p}{;}
\PY{n}{PasswordPolicy}\PY{+w}{ }\PY{n}{policy}\PY{+w}{ }\PY{o}{=}\PY{+w}{ }\PY{k}{new}\PY{+w}{ }\PY{n}{PasswordPolicy}\PY{p}{(}\PY{n}{symbols}\PY{p}{)}\PY{p}{;}
\PY{n}{policy}\PY{p}{.}\PY{n+na}{getAllowedSymbols}\PY{p}{(}\PY{p}{)}\PY{o}{[}\PY{l+m+mi}{0}\PY{o}{]}\PY{+w}{ }\PY{o}{=}\PY{+w}{ }\PY{l+s+sc}{\PYZsq{}X\PYZsq{}}\PY{p}{;}\PY{+w}{ }\PY{c+c1}{// mutates internal state via the returned array!}
\PY{n}{symbols}\PY{o}{[}\PY{l+m+mi}{1}\PY{o}{]}\PY{+w}{ }\PY{o}{=}\PY{+w}{ }\PY{l+s+sc}{\PYZsq{}Y\PYZsq{}}\PY{p}{;}\PY{+w}{                    }\PY{c+c1}{// mutates internal state via the original reference!}
\end{Verbatim}
\end{codebox}

\begin{codebox}
\begin{Verbatim}[commandchars=\\\{\}]
\PY{c+c1}{// FIXED: clone() on the way in AND on the way out.}
\PY{k+kd}{class} \PY{n+nc}{PasswordPolicy}\PY{+w}{ }\PY{p}{\PYZob{}}
\PY{+w}{    }\PY{k+kd}{private}\PY{+w}{ }\PY{k+kd}{final}\PY{+w}{ }\PY{k+kt}{char}\PY{o}{[}\PY{o}{]}\PY{+w}{ }\PY{n}{allowedSymbols}\PY{p}{;}
\PY{+w}{    }\PY{n}{PasswordPolicy}\PY{p}{(}\PY{k+kt}{char}\PY{o}{[}\PY{o}{]}\PY{+w}{ }\PY{n}{allowedSymbols}\PY{p}{)}\PY{+w}{ }\PY{p}{\PYZob{}}
\PY{+w}{        }\PY{k}{this}\PY{p}{.}\PY{n+na}{allowedSymbols}\PY{+w}{ }\PY{o}{=}\PY{+w}{ }\PY{n}{allowedSymbols}\PY{p}{.}\PY{n+na}{clone}\PY{p}{(}\PY{p}{)}\PY{p}{;}\PY{+w}{ }\PY{c+c1}{// array clone() is a fast shallow copy}
\PY{+w}{    }\PY{p}{\PYZcb{}}
\PY{+w}{    }\PY{k+kt}{char}\PY{o}{[}\PY{o}{]}\PY{+w}{ }\PY{n+nf}{getAllowedSymbols}\PY{p}{(}\PY{p}{)}\PY{+w}{ }\PY{p}{\PYZob{}}
\PY{+w}{        }\PY{k}{return}\PY{+w}{ }\PY{n}{allowedSymbols}\PY{p}{.}\PY{n+na}{clone}\PY{p}{(}\PY{p}{)}\PY{p}{;}
\PY{+w}{    }\PY{p}{\PYZcb{}}
\PY{p}{\PYZcb{}}
\end{Verbatim}
\end{codebox}

(Note: \code{array.\allowbreak{}clone()} is one of the few reasonable uses of \code{clone()} in Java — arrays have well-defined, non-broken clone semantics, unlike ordinary objects.)

\subsection[Dates are a classic leak vector too]{Dates are a classic leak vector too}
\label{h:5:dates-are-a-classic-leak-vector-too}
\code{java.\allowbreak{}util.\allowbreak{}Date} is mutable (it has \code{setTime()}), which is one of the reasons the legacy date API is considered a design mistake. The modern \code{java.\allowbreak{}time} types (\code{LocalDate}, \code{Instant}, \code{LocalDateTime}, etc.) are all immutable, so \textbf{switching to \code{java.\allowbreak{}time} eliminates this entire category of bug} — no defensive copy is needed because there is nothing to mutate.

\begin{codebox}
\begin{Verbatim}[commandchars=\\\{\}]
\PY{k+kn}{import}\PY{+w}{ }\PY{n+nn}{java.time.LocalDate}\PY{p}{;}

\PY{c+c1}{// GOOD: no defensive copying needed, because LocalDate is immutable.}
\PY{k+kd}{final}\PY{+w}{ }\PY{k+kd}{class} \PY{n+nc}{Membership}\PY{+w}{ }\PY{p}{\PYZob{}}
\PY{+w}{    }\PY{k+kd}{private}\PY{+w}{ }\PY{k+kd}{final}\PY{+w}{ }\PY{n}{LocalDate}\PY{+w}{ }\PY{n}{startDate}\PY{p}{;}
\PY{+w}{    }\PY{n}{Membership}\PY{p}{(}\PY{n}{LocalDate}\PY{+w}{ }\PY{n}{startDate}\PY{p}{)}\PY{+w}{ }\PY{p}{\PYZob{}}
\PY{+w}{        }\PY{k}{this}\PY{p}{.}\PY{n+na}{startDate}\PY{+w}{ }\PY{o}{=}\PY{+w}{ }\PY{n}{startDate}\PY{p}{;}\PY{+w}{ }\PY{c+c1}{// safe: LocalDate can\PYZsq{}t be mutated by anyone}
\PY{+w}{    }\PY{p}{\PYZcb{}}
\PY{+w}{    }\PY{n}{LocalDate}\PY{+w}{ }\PY{n+nf}{getStartDate}\PY{p}{(}\PY{p}{)}\PY{+w}{ }\PY{p}{\PYZob{}}
\PY{+w}{        }\PY{k}{return}\PY{+w}{ }\PY{n}{startDate}\PY{p}{;}\PY{+w}{ }\PY{c+c1}{// safe: caller can\PYZsq{}t mutate it either}
\PY{+w}{    }\PY{p}{\PYZcb{}}
\PY{p}{\PYZcb{}}
\end{Verbatim}
\end{codebox}

\textbf{Interview one-liner:} “Defensive copying is only necessary for mutable types. The best fix is often to avoid mutable types altogether — prefer \code{java.\allowbreak{}time} over \code{Date}, and \code{List.\allowbreak{}copyOf}/records over hand-rolled mutable collections.”

\section[Common Code-Review Interview Pitfalls]{Common Code-Review Interview Pitfalls}
\label{h:5:common-code-review-interview-pitfalls}
\begin{enumerate}
\item 
\textbf{Overriding \code{equals()} without \code{hashCode()}.}
Why it matters: breaks the equals/hashCode contract; the object silently “disappears” from \code{HashMap}/\code{HashSet} lookups.

\begin{codebox}
\begin{Verbatim}[commandchars=\\\{\}]
\PY{c+c1}{// Before: equals() overridden, hashCode() left default (identity\PYZhy{}based)}
\PY{n+nd}{@Override}\PY{+w}{ }\PY{k+kd}{public}\PY{+w}{ }\PY{k+kt}{boolean}\PY{+w}{ }\PY{n+nf}{equals}\PY{p}{(}\PY{n}{Object}\PY{+w}{ }\PY{n}{o}\PY{p}{)}\PY{+w}{ }\PY{p}{\PYZob{}}\PY{+w}{ }\PY{k}{return}\PY{+w}{ }\PY{n}{o}\PY{+w}{ }\PY{k}{instanceof}\PY{+w}{ }\PY{n}{Id}\PY{+w}{ }\PY{n}{i}\PY{+w}{ }\PY{o}{\PYZam{}}\PY{o}{\PYZam{}}\PY{+w}{ }\PY{n}{i}\PY{p}{.}\PY{n+na}{value}\PY{+w}{ }\PY{o}{=}\PY{o}{=}\PY{+w}{ }\PY{n}{value}\PY{p}{;}\PY{+w}{ }\PY{p}{\PYZcb{}}

\PY{c+c1}{// After: keep them in sync}
\PY{n+nd}{@Override}\PY{+w}{ }\PY{k+kd}{public}\PY{+w}{ }\PY{k+kt}{int}\PY{+w}{ }\PY{n+nf}{hashCode}\PY{p}{(}\PY{p}{)}\PY{+w}{ }\PY{p}{\PYZob{}}\PY{+w}{ }\PY{k}{return}\PY{+w}{ }\PY{n}{Integer}\PY{p}{.}\PY{n+na}{hashCode}\PY{p}{(}\PY{n}{value}\PY{p}{)}\PY{p}{;}\PY{+w}{ }\PY{p}{\PYZcb{}}
\end{Verbatim}
\end{codebox}

\item 
\textbf{Using a mutable field as a \code{HashMap}/\code{HashSet} key’s basis for \code{hashCode()}.}
Why it matters: mutating the key after insertion makes it unfindable — a silent data-loss bug.

\begin{codebox}
\begin{Verbatim}[commandchars=\\\{\}]
\PY{c+c1}{// Before: key.setStatus(\PYZdq{}DONE\PYZdq{}) after map.put(key, ...) \PYZhy{}\PYZgt{} map.get(key) returns null}
\PY{c+c1}{// After: use an immutable id (String/UUID/record) as the key instead of the mutable entity}
\end{Verbatim}
\end{codebox}

\item 
\textbf{\code{instanceof}-based \code{equals()} in a class hierarchy where subclasses add fields.}
Why it matters: breaks symmetry (\code{a.\allowbreak{}equals(\allowbreak{}b)\allowbreak{}~\allowbreak{}!=\allowbreak{}~\allowbreak{}b.\allowbreak{}equals(\allowbreak{}a)}), causing inconsistent behavior in collections.

\begin{codebox}
\begin{Verbatim}[commandchars=\\\{\}]
\PY{c+c1}{// Before: if (!(o instanceof Point p)) ... // ColorPoint subclass breaks symmetry}
\PY{c+c1}{// After: if (o == null || getClass() != o.getClass()) return false; // or make the class final}
\end{Verbatim}
\end{codebox}

\item 
\textbf{Integer subtraction inside \code{compareTo}/\code{compare}.}
Why it matters: overflows silently for large or negative values, corrupting sort order.

\begin{codebox}
\begin{Verbatim}[commandchars=\\\{\}]
\PY{c+c1}{// Before: return this.value \PYZhy{} other.value;}
\PY{c+c1}{// After:  return Integer.compare(this.value, other.value);}
\end{Verbatim}
\end{codebox}

\item 
\textbf{Using \code{Cloneable}/\code{Object.\allowbreak{}clone()} for copying.}
Why it matters: bypasses constructors (invariants unchecked), fragile with \code{final} fields and subclassing, throws a checked exception unnecessarily.

\begin{codebox}
\begin{Verbatim}[commandchars=\\\{\}]
\PY{c+c1}{// Before: class Team implements Cloneable \PYZob{} ... super.clone() ... \PYZcb{}}
\PY{c+c1}{// After:  Team(Team other) \PYZob{} this(other.name, other.members); \PYZcb{} // copy constructor}
\end{Verbatim}
\end{codebox}

\item 
\textbf{Shallow-copying an object that has mutable reference fields.}
Why it matters: the “copy” secretly shares mutable state with the original — mutating one affects the other.

\begin{codebox}
\begin{Verbatim}[commandchars=\\\{\}]
\PY{c+c1}{// Before: Team copy = (Team) super.clone(); // members list is SHARED}
\PY{c+c1}{// After:  copy.members = new ArrayList\PYZlt{}\PYZgt{}(this.members); // deep\PYZhy{}ish copy}
\end{Verbatim}
\end{codebox}

\item 
\textbf{\code{Optional} used as a field, parameter, or inside a collection.}
Why it matters: adds overhead, is not \code{Serializable}, and \code{Optional} itself can be \code{null} — defeating its purpose.

\begin{codebox}
\begin{Verbatim}[commandchars=\\\{\}]
\PY{c+c1}{// Before: void updateNickname(Optional\PYZlt{}String\PYZgt{} nickname)}
\PY{c+c1}{// After:  void updateNickname(String nickname) // nullable, documented}
\end{Verbatim}
\end{codebox}

\item 
\textbf{Calling \code{or\allowbreak{}Else(\allowbreak{}expensive\allowbreak{}Call())} instead of \code{orElseGet(...)}.}
Why it matters: \code{orElse}’s argument is always evaluated eagerly, even when the \code{Optional} is present — wasted work or unwanted side effects.

\begin{codebox}
\begin{Verbatim}[commandchars=\\\{\}]
\PY{c+c1}{// Before: opt.orElse(buildExpensiveDefault());}
\PY{c+c1}{// After:  opt.orElseGet(this::buildExpensiveDefault);}
\end{Verbatim}
\end{codebox}

\item 
\textbf{Exposing an internal mutable collection or array directly from a getter.}
Why it matters: callers can mutate “encapsulated” internal state from outside the class, breaking invariants silently.

\begin{codebox}
\begin{Verbatim}[commandchars=\\\{\}]
\PY{c+c1}{// Before: List\PYZlt{}String\PYZgt{} getItems() \PYZob{} return items; \PYZcb{}}
\PY{c+c1}{// After:  List\PYZlt{}String\PYZgt{} getItems() \PYZob{} return List.copyOf(items); \PYZcb{}}
\end{Verbatim}
\end{codebox}

\item 
\textbf{Storing a caller-provided mutable object directly in the constructor without copying.}
Why it matters: the caller can mutate your object’s “immutable” internal state after construction, from their original reference.

\begin{codebox}
\begin{Verbatim}[commandchars=\\\{\}]
\PY{c+c1}{// Before: this.startDate = startDate; // Date is mutable}
\PY{c+c1}{// After:  this.startDate = new Date(startDate.getTime()); // defensive copy in}
\end{Verbatim}
\end{codebox}

\item 
\textbf{Confusing \code{Collections.\allowbreak{}unmodifiable\allowbreak{}List()} (a view) with a true immutable copy.}
Why it matters: an unmodifiable \emph{view} still reflects changes made to the original backing collection — it does not provide independence.

\begin{codebox}
\begin{Verbatim}[commandchars=\\\{\}]
\PY{c+c1}{// Before: this.names = Collections.unmodifiableList(names); // still a view over caller\PYZsq{}s list}
\PY{c+c1}{// After:  this.names = List.copyOf(names); // independent AND unmodifiable}
\end{Verbatim}
\end{codebox}

\item 
\textbf{\code{compareTo} that is inconsistent with \code{equals()}, used with \code{TreeSet}/\code{TreeMap}.}
Why it matters: \code{TreeSet}/\code{TreeMap} use \code{compareTo} (not \code{equals()}) to decide duplicates, so “equal” elements by \code{equals()} can both end up in the set if \code{compareTo} disagrees — or vice versa, distinct elements silently vanish as “duplicates.”

\begin{codebox}
\begin{Verbatim}[commandchars=\\\{\}]
\PY{c+c1}{// Before: compareTo compares only `lastName`, but equals() compares full name \PYZhy{}\PYZgt{} surprises in TreeSet}
\PY{c+c1}{// After:  make compareTo and equals() agree, or clearly document the difference}
\end{Verbatim}
\end{codebox}

\item 
\textbf{Sorting a list that may contain \code{null} with \code{Comparator.\allowbreak{}natural\allowbreak{}Order()} directly.}
Why it matters: throws \code{NullPointerException} at runtime instead of handling nulls predictably.

\begin{codebox}
\begin{Verbatim}[commandchars=\\\{\}]
\PY{c+c1}{// Before: list.sort(Comparator.naturalOrder());}
\PY{c+c1}{// After:  list.sort(Comparator.nullsFirst(Comparator.naturalOrder()));}
\end{Verbatim}
\end{codebox}

\item 
\textbf{Forgetting to regenerate IDE/Lombok-generated \code{equals()}/\code{hashCode()} after adding a field.}
Why it matters: the new field is silently excluded from equality checks, causing subtly wrong deduplication or map behavior.

\begin{codebox}
\begin{Verbatim}[commandchars=\\\{\}]
\PY{c+c1}{// Before: `status` field added, but equals()/hashCode() still only reference `id`}
\PY{c+c1}{// After:  regenerate, or intentionally document ID\PYZhy{}only equality semantics}
\end{Verbatim}
\end{codebox}

\item 
\textbf{Hand-writing a data class (fields, constructor, getters, \code{equals}/\code{hashCode}/\code{toString}) that could be a \code{record}.}
Why it matters: more boilerplate to maintain, more surface area for the bugs above (missing \code{hashCode}, stale generated code); records get correct semantics for free.

\begin{codebox}
\begin{Verbatim}[commandchars=\\\{\}]
\PY{c+c1}{// Before: class Point \PYZob{} private final int x, y; /* \PYZti{}30 lines of boilerplate */ \PYZcb{}}
\PY{c+c1}{// After:  record Point(int x, int y) \PYZob{}\PYZcb{}}
\end{Verbatim}
\end{codebox}

\end{enumerate}


\setcounter{chapter}{5}
\chapter{Exception Handling}
\label{chap:6}
\label{h:6:6-exception-handling}
Things go wrong at runtime: files are missing, networks time out, arrays get the wrong index. Java uses \textbf{exceptions} to signal that something went wrong, separate from your normal return values. This chapter explains how exceptions work under the hood, how to handle them correctly, and the mistakes that reviewers flag most often in real code review. We target Java 21+ throughout.

\section[Exceptions]{Exceptions}
\label{h:6:exceptions}
An \textbf{exception} is an object that describes an unexpected event. When something goes wrong, Java \textbf{throws} an exception object. If nothing handles it, the program (or thread) stops and prints a \textbf{stack trace}. If code \textbf{catches} it, the program can recover or fail gracefully instead of crashing.

\begin{codebox}
\begin{Verbatim}[commandchars=\\\{\}]
\PY{k+kd}{public}\PY{+w}{ }\PY{k+kd}{class} \PY{n+nc}{Divider}\PY{+w}{ }\PY{p}{\PYZob{}}
\PY{+w}{    }\PY{k+kd}{public}\PY{+w}{ }\PY{k+kd}{static}\PY{+w}{ }\PY{k+kt}{int}\PY{+w}{ }\PY{n+nf}{divide}\PY{p}{(}\PY{k+kt}{int}\PY{+w}{ }\PY{n}{a}\PY{p}{,}\PY{+w}{ }\PY{k+kt}{int}\PY{+w}{ }\PY{n}{b}\PY{p}{)}\PY{+w}{ }\PY{p}{\PYZob{}}
\PY{+w}{        }\PY{k}{return}\PY{+w}{ }\PY{n}{a}\PY{+w}{ }\PY{o}{/}\PY{+w}{ }\PY{n}{b}\PY{p}{;}\PY{+w}{ }\PY{c+c1}{// throws ArithmeticException if b == 0}
\PY{+w}{    }\PY{p}{\PYZcb{}}

\PY{+w}{    }\PY{k+kd}{public}\PY{+w}{ }\PY{k+kd}{static}\PY{+w}{ }\PY{k+kt}{void}\PY{+w}{ }\PY{n+nf}{main}\PY{p}{(}\PY{n}{String}\PY{o}{[}\PY{o}{]}\PY{+w}{ }\PY{n}{args}\PY{p}{)}\PY{+w}{ }\PY{p}{\PYZob{}}
\PY{+w}{        }\PY{k}{try}\PY{+w}{ }\PY{p}{\PYZob{}}
\PY{+w}{            }\PY{n}{System}\PY{p}{.}\PY{n+na}{out}\PY{p}{.}\PY{n+na}{println}\PY{p}{(}\PY{n}{divide}\PY{p}{(}\PY{l+m+mi}{10}\PY{p}{,}\PY{+w}{ }\PY{l+m+mi}{0}\PY{p}{)}\PY{p}{)}\PY{p}{;}
\PY{+w}{        }\PY{p}{\PYZcb{}}\PY{+w}{ }\PY{k}{catch}\PY{+w}{ }\PY{p}{(}\PY{n}{ArithmeticException}\PY{+w}{ }\PY{n}{e}\PY{p}{)}\PY{+w}{ }\PY{p}{\PYZob{}}
\PY{+w}{            }\PY{n}{System}\PY{p}{.}\PY{n+na}{out}\PY{p}{.}\PY{n+na}{println}\PY{p}{(}\PY{l+s}{\PYZdq{}}\PY{l+s}{Cannot divide by zero: }\PY{l+s}{\PYZdq{}}\PY{+w}{ }\PY{o}{+}\PY{+w}{ }\PY{n}{e}\PY{p}{.}\PY{n+na}{getMessage}\PY{p}{(}\PY{p}{)}\PY{p}{)}\PY{p}{;}
\PY{+w}{        }\PY{p}{\PYZcb{}}
\PY{+w}{    }\PY{p}{\PYZcb{}}
\PY{p}{\PYZcb{}}
\end{Verbatim}
\end{codebox}

\subsection[The Throwable Hierarchy]{The Throwable Hierarchy}
\label{h:6:the-throwable-hierarchy}
Every exception in Java is a subclass of \code{Throwable}. There are three important branches: \code{Error}, \code{Exception}, and inside \code{Exception}, \code{RuntimeException}.

\begin{codebox}
\begin{Verbatim}
Throwable
├── Error                          (serious, usually unrecoverable JVM problems)
│   ├── OutOfMemoryError
│   ├── StackOverflowError
│   └── VirtualMachineError
│
└── Exception                      (things application code can reasonably handle)
    ├── IOException                (checked)
    ├── SQLException                (checked)
    ├── InterruptedException        (checked)
    │
    └── RuntimeException            (unchecked)
        ├── NullPointerException
        ├── IllegalArgumentException
        │     └── NumberFormatException
        ├── IllegalStateException
        ├── IndexOutOfBoundsException
        │     ├── ArrayIndexOutOfBoundsException
        │     └── StringIndexOutOfBoundsException
        ├── ClassCastException
        ├── ArithmeticException
        └── UnsupportedOperationException
\end{Verbatim}
\end{codebox}

\begin{itemize}
\item 
\code{Error} — things like \code{OutOfMemoryError} or \code{StackOverflowError}. These signal that the JVM itself is in trouble. Application code almost never catches these, because there is usually nothing useful you can do.

\item 
\code{Exception} — problems your application should anticipate and possibly recover from, such as a missing file.

\item 
\code{RuntimeException} — a subclass of \code{Exception} for programming errors, like calling a method on \code{null}. These are \textbf{unchecked} (explained below).

\end{itemize}

\code{Throwable} itself is rarely used directly. Catching \code{Throwable} is even broader than catching \code{Exception} — it also captures \code{Error}, which is almost never what you want.

\begin{codebox}
\begin{Verbatim}[commandchars=\\\{\}]
\PY{c+c1}{// Avoid: this also swallows OutOfMemoryError and StackOverflowError}
\PY{k}{try}\PY{+w}{ }\PY{p}{\PYZob{}}
\PY{+w}{    }\PY{n}{riskyOperation}\PY{p}{(}\PY{p}{)}\PY{p}{;}
\PY{p}{\PYZcb{}}\PY{+w}{ }\PY{k}{catch}\PY{+w}{ }\PY{p}{(}\PY{n}{Throwable}\PY{+w}{ }\PY{n}{t}\PY{p}{)}\PY{+w}{ }\PY{p}{\PYZob{}}
\PY{+w}{    }\PY{n}{log}\PY{p}{.}\PY{n+na}{error}\PY{p}{(}\PY{l+s}{\PYZdq{}}\PY{l+s}{Something went wrong}\PY{l+s}{\PYZdq{}}\PY{p}{,}\PY{+w}{ }\PY{n}{t}\PY{p}{)}\PY{p}{;}
\PY{p}{\PYZcb{}}
\end{Verbatim}
\end{codebox}

\subsection[Stack Traces and How to Read Them]{Stack Traces and How to Read Them}
\label{h:6:stack-traces-and-how-to-read-them}
A \textbf{stack trace} is a printed list of method calls that were active when the exception was created. It reads top to bottom: the top line is where the exception happened, and each line below is “called from here.”

\begin{codebox}
\begin{Verbatim}[commandchars=\\\{\}]
\PY{k+kd}{public}\PY{+w}{ }\PY{k+kd}{class} \PY{n+nc}{Chain}\PY{+w}{ }\PY{p}{\PYZob{}}
\PY{+w}{    }\PY{k+kd}{public}\PY{+w}{ }\PY{k+kd}{static}\PY{+w}{ }\PY{k+kt}{void}\PY{+w}{ }\PY{n+nf}{main}\PY{p}{(}\PY{n}{String}\PY{o}{[}\PY{o}{]}\PY{+w}{ }\PY{n}{args}\PY{p}{)}\PY{+w}{ }\PY{p}{\PYZob{}}
\PY{+w}{        }\PY{n}{step1}\PY{p}{(}\PY{p}{)}\PY{p}{;}
\PY{+w}{    }\PY{p}{\PYZcb{}}
\PY{+w}{    }\PY{k+kd}{static}\PY{+w}{ }\PY{k+kt}{void}\PY{+w}{ }\PY{n+nf}{step1}\PY{p}{(}\PY{p}{)}\PY{+w}{ }\PY{p}{\PYZob{}}\PY{+w}{ }\PY{n}{step2}\PY{p}{(}\PY{p}{)}\PY{p}{;}\PY{+w}{ }\PY{p}{\PYZcb{}}
\PY{+w}{    }\PY{k+kd}{static}\PY{+w}{ }\PY{k+kt}{void}\PY{+w}{ }\PY{n+nf}{step2}\PY{p}{(}\PY{p}{)}\PY{+w}{ }\PY{p}{\PYZob{}}\PY{+w}{ }\PY{n}{step3}\PY{p}{(}\PY{p}{)}\PY{p}{;}\PY{+w}{ }\PY{p}{\PYZcb{}}
\PY{+w}{    }\PY{k+kd}{static}\PY{+w}{ }\PY{k+kt}{void}\PY{+w}{ }\PY{n+nf}{step3}\PY{p}{(}\PY{p}{)}\PY{+w}{ }\PY{p}{\PYZob{}}\PY{+w}{ }\PY{k}{throw}\PY{+w}{ }\PY{k}{new}\PY{+w}{ }\PY{n}{IllegalStateException}\PY{p}{(}\PY{l+s}{\PYZdq{}}\PY{l+s}{boom}\PY{l+s}{\PYZdq{}}\PY{p}{)}\PY{p}{;}\PY{+w}{ }\PY{p}{\PYZcb{}}
\PY{p}{\PYZcb{}}
\end{Verbatim}
\end{codebox}

Output:

\begin{codebox}
\begin{Verbatim}
Exception in thread "main" java.lang.IllegalStateException: boom
	at Chain.step3(Chain.java:7)
	at Chain.step2(Chain.java:6)
	at Chain.step1(Chain.java:5)
	at Chain.main(Chain.java:2)
\end{Verbatim}
\end{codebox}

How to read it:

\begin{enumerate}
\item 
\textbf{First line} — exception class and message. This tells you \emph{what} went wrong.

\item 
\textbf{\code{at} lines} — the call stack at the moment of the throw, most recent call first. \code{Chain.\allowbreak{}step3(\allowbreak{}Chain.\allowbreak{}java:\allowbreak{}7)} means line 7 of \code{Chain.\allowbreak{}java}, inside method \code{step3}.

\item 
\textbf{Caused by} — if present, a chained exception (see below). Always read the \emph{innermost} “Caused by” first; that is usually the real root cause.

\item 
\textbf{\code{...\allowbreak{}~\allowbreak{}N~\allowbreak{}more}} — Java collapses stack frames that are identical to a previous trace, to save space, when printing a chained exception.

\end{enumerate}

\begin{codebox}
\begin{Verbatim}[commandchars=\\\{\}]
\PY{k}{try}\PY{+w}{ }\PY{p}{\PYZob{}}
\PY{+w}{    }\PY{n}{loadConfig}\PY{p}{(}\PY{p}{)}\PY{p}{;}
\PY{p}{\PYZcb{}}\PY{+w}{ }\PY{k}{catch}\PY{+w}{ }\PY{p}{(}\PY{n}{IOException}\PY{+w}{ }\PY{n}{e}\PY{p}{)}\PY{+w}{ }\PY{p}{\PYZob{}}
\PY{+w}{    }\PY{k}{throw}\PY{+w}{ }\PY{k}{new}\PY{+w}{ }\PY{n}{RuntimeException}\PY{p}{(}\PY{l+s}{\PYZdq{}}\PY{l+s}{Failed to start application}\PY{l+s}{\PYZdq{}}\PY{p}{,}\PY{+w}{ }\PY{n}{e}\PY{p}{)}\PY{p}{;}
\PY{p}{\PYZcb{}}
\end{Verbatim}
\end{codebox}

\begin{codebox}
\begin{Verbatim}
Exception in thread "main" java.lang.RuntimeException: Failed to start application
	at App.main(App.java:12)
Caused by: java.io.FileNotFoundException: config.yml (No such file or directory)
	at java.base/java.io.FileInputStream.open0(Native Method)
	... 5 more
\end{Verbatim}
\end{codebox}

\section[Checked vs Unchecked Exceptions]{Checked vs Unchecked Exceptions}
\label{h:6:checked-vs-unchecked-exceptions}
Java splits exceptions into two categories based on whether the compiler forces you to deal with them.

{\tablefont
\begin{xltabular}{\linewidth}{@{}L{0.541}L{1.270}L{1.189}@{}}
\toprule
\rowcolor{tablehdr}
\thd{} & \thd{Checked} & \thd{Unchecked} \\
\midrule
\endhead
Base class & \code{Exception} (not \code{RuntimeException}) & \code{RuntimeException} or \code{Error} \\
Compiler enforcement & Must be caught or declared with \code{throws} & No compiler enforcement \\
Typical cause & External conditions (file missing, network down) & Programming bugs (null access, bad argument) \\
Examples & \code{IOException}, \code{SQLException}, \code{InterruptedException} & \code{NullPointerException}, \code{Illegal\allowbreak{}Argument\allowbreak{}Exception}, \code{Index\allowbreak{}Out\allowbreak{}Of\allowbreak{}Bounds\allowbreak{}Exception} \\
Should caller recover? & Often yes — retry, fallback, ask user again & Usually no — fix the bug instead \\
\bottomrule
\end{xltabular}
}

\begin{codebox}
\begin{Verbatim}[commandchars=\\\{\}]
\PY{k+kn}{import}\PY{+w}{ }\PY{n+nn}{java.io.*}\PY{p}{;}

\PY{k+kd}{public}\PY{+w}{ }\PY{k+kd}{class} \PY{n+nc}{CheckedExample}\PY{+w}{ }\PY{p}{\PYZob{}}
\PY{+w}{    }\PY{c+c1}{// Checked exception: caller MUST catch it or declare \PYZdq{}throws IOException\PYZdq{}}
\PY{+w}{    }\PY{k+kd}{public}\PY{+w}{ }\PY{k+kd}{static}\PY{+w}{ }\PY{n}{String}\PY{+w}{ }\PY{n+nf}{readFirstLine}\PY{p}{(}\PY{n}{String}\PY{+w}{ }\PY{n}{path}\PY{p}{)}\PY{+w}{ }\PY{k+kd}{throws}\PY{+w}{ }\PY{n}{IOException}\PY{+w}{ }\PY{p}{\PYZob{}}
\PY{+w}{        }\PY{k}{try}\PY{+w}{ }\PY{p}{(}\PY{n}{BufferedReader}\PY{+w}{ }\PY{n}{reader}\PY{+w}{ }\PY{o}{=}\PY{+w}{ }\PY{k}{new}\PY{+w}{ }\PY{n}{BufferedReader}\PY{p}{(}\PY{k}{new}\PY{+w}{ }\PY{n}{FileReader}\PY{p}{(}\PY{n}{path}\PY{p}{)}\PY{p}{)}\PY{p}{)}\PY{+w}{ }\PY{p}{\PYZob{}}
\PY{+w}{            }\PY{k}{return}\PY{+w}{ }\PY{n}{reader}\PY{p}{.}\PY{n+na}{readLine}\PY{p}{(}\PY{p}{)}\PY{p}{;}
\PY{+w}{        }\PY{p}{\PYZcb{}}
\PY{+w}{    }\PY{p}{\PYZcb{}}
\PY{p}{\PYZcb{}}
\end{Verbatim}
\end{codebox}

\begin{codebox}
\begin{Verbatim}[commandchars=\\\{\}]
\PY{k+kd}{public}\PY{+w}{ }\PY{k+kd}{class} \PY{n+nc}{UncheckedExample}\PY{+w}{ }\PY{p}{\PYZob{}}
\PY{+w}{    }\PY{c+c1}{// Unchecked: compiler does not force the caller to do anything}
\PY{+w}{    }\PY{k+kd}{public}\PY{+w}{ }\PY{k+kd}{static}\PY{+w}{ }\PY{k+kt}{int}\PY{+w}{ }\PY{n+nf}{parsePositive}\PY{p}{(}\PY{n}{String}\PY{+w}{ }\PY{n}{value}\PY{p}{)}\PY{+w}{ }\PY{p}{\PYZob{}}
\PY{+w}{        }\PY{k+kt}{int}\PY{+w}{ }\PY{n}{n}\PY{+w}{ }\PY{o}{=}\PY{+w}{ }\PY{n}{Integer}\PY{p}{.}\PY{n+na}{parseInt}\PY{p}{(}\PY{n}{value}\PY{p}{)}\PY{p}{;}\PY{+w}{       }\PY{c+c1}{// may throw NumberFormatException}
\PY{+w}{        }\PY{k}{if}\PY{+w}{ }\PY{p}{(}\PY{n}{n}\PY{+w}{ }\PY{o}{\PYZlt{}}\PY{+w}{ }\PY{l+m+mi}{0}\PY{p}{)}\PY{+w}{ }\PY{p}{\PYZob{}}
\PY{+w}{            }\PY{k}{throw}\PY{+w}{ }\PY{k}{new}\PY{+w}{ }\PY{n}{IllegalArgumentException}\PY{p}{(}\PY{l+s}{\PYZdq{}}\PY{l+s}{Value must be positive: }\PY{l+s}{\PYZdq{}}\PY{+w}{ }\PY{o}{+}\PY{+w}{ }\PY{n}{n}\PY{p}{)}\PY{p}{;}
\PY{+w}{        }\PY{p}{\PYZcb{}}
\PY{+w}{        }\PY{k}{return}\PY{+w}{ }\PY{n}{n}\PY{p}{;}
\PY{+w}{    }\PY{p}{\PYZcb{}}
\PY{p}{\PYZcb{}}
\end{Verbatim}
\end{codebox}

\subsection[When to Use Which]{When to Use Which}
\label{h:6:when-to-use-which}
{\tablefont
\begin{xltabular}{\linewidth}{@{}L{1.560}L{0.440}@{}}
\toprule
\rowcolor{tablehdr}
\thd{Situation} & \thd{Use} \\
\midrule
\endhead
Caller can realistically recover (retry, prompt again, use a default) & Checked exception \\
The failure is a bug in the calling code (bad argument, wrong state) & Unchecked exception \\
Library/API boundary where you want to force callers to think about failure & Checked exception \\
Internal code where wrapping every call in \code{try/\allowbreak{}catch} would just add noise & Unchecked exception \\
\bottomrule
\end{xltabular}
}

Modern Java style (post Java 8, especially with streams and lambdas) leans toward \textbf{unchecked} exceptions, because checked exceptions do not compose well with functional interfaces (a \code{Function\textless{}\allowbreak{}T,\allowbreak{}~\allowbreak{}R\textgreater{}} cannot declare \code{throws~\allowbreak{}IOException}). This is why many modern libraries (e.g., \code{UncheckedIOException}) provide unchecked wrappers around checked ones.

\begin{codebox}
\begin{Verbatim}[commandchars=\\\{\}]
\PY{k+kn}{import}\PY{+w}{ }\PY{n+nn}{java.io.UncheckedIOException}\PY{p}{;}
\PY{k+kn}{import}\PY{+w}{ }\PY{n+nn}{java.io.IOException}\PY{p}{;}
\PY{k+kn}{import}\PY{+w}{ }\PY{n+nn}{java.nio.file.*}\PY{p}{;}
\PY{k+kn}{import}\PY{+w}{ }\PY{n+nn}{java.util.List}\PY{p}{;}
\PY{k+kn}{import}\PY{+w}{ }\PY{n+nn}{java.util.stream.Stream}\PY{p}{;}

\PY{k+kd}{public}\PY{+w}{ }\PY{k+kd}{class} \PY{n+nc}{StreamsAndCheckedExceptions}\PY{+w}{ }\PY{p}{\PYZob{}}
\PY{+w}{    }\PY{k+kd}{public}\PY{+w}{ }\PY{k+kd}{static}\PY{+w}{ }\PY{n}{List}\PY{o}{\PYZlt{}}\PY{n}{String}\PY{o}{\PYZgt{}}\PY{+w}{ }\PY{n+nf}{linesOf}\PY{p}{(}\PY{n}{Path}\PY{+w}{ }\PY{n}{path}\PY{p}{)}\PY{+w}{ }\PY{p}{\PYZob{}}
\PY{+w}{        }\PY{k}{try}\PY{+w}{ }\PY{p}{(}\PY{n}{Stream}\PY{o}{\PYZlt{}}\PY{n}{String}\PY{o}{\PYZgt{}}\PY{+w}{ }\PY{n}{lines}\PY{+w}{ }\PY{o}{=}\PY{+w}{ }\PY{n}{Files}\PY{p}{.}\PY{n+na}{lines}\PY{p}{(}\PY{n}{path}\PY{p}{)}\PY{p}{)}\PY{+w}{ }\PY{p}{\PYZob{}}
\PY{+w}{            }\PY{k}{return}\PY{+w}{ }\PY{n}{lines}\PY{p}{.}\PY{n+na}{toList}\PY{p}{(}\PY{p}{)}\PY{p}{;}
\PY{+w}{        }\PY{p}{\PYZcb{}}\PY{+w}{ }\PY{k}{catch}\PY{+w}{ }\PY{p}{(}\PY{n}{IOException}\PY{+w}{ }\PY{n}{e}\PY{p}{)}\PY{+w}{ }\PY{p}{\PYZob{}}
\PY{+w}{            }\PY{c+c1}{// Files.lines already throws unchecked UncheckedIOException lazily during iteration,}
\PY{+w}{            }\PY{c+c1}{// but the try\PYZhy{}with\PYZhy{}resources open call itself is still checked, so we wrap it too.}
\PY{+w}{            }\PY{k}{throw}\PY{+w}{ }\PY{k}{new}\PY{+w}{ }\PY{n}{UncheckedIOException}\PY{p}{(}\PY{n}{e}\PY{p}{)}\PY{p}{;}
\PY{+w}{        }\PY{p}{\PYZcb{}}
\PY{+w}{    }\PY{p}{\PYZcb{}}
\PY{p}{\PYZcb{}}
\end{Verbatim}
\end{codebox}

\section[Exception Handling Best Practices]{Exception Handling Best Practices}
\label{h:6:exception-handling-best-practices}
Handling exceptions well is more about discipline than syntax. This section covers the patterns reviewers look for.

\subsection[finally Semantics]{\code{finally} Semantics}
\label{h:6:finally-semantics}
Code in a \code{finally} block always runs, whether the \code{try} succeeded, threw, or returned. It is the right place for cleanup that must happen no matter what.

\begin{codebox}
\begin{Verbatim}[commandchars=\\\{\}]
\PY{k+kd}{public}\PY{+w}{ }\PY{k+kd}{class} \PY{n+nc}{FinallyDemo}\PY{+w}{ }\PY{p}{\PYZob{}}
\PY{+w}{    }\PY{k+kd}{static}\PY{+w}{ }\PY{k+kt}{int}\PY{+w}{ }\PY{n+nf}{compute}\PY{p}{(}\PY{p}{)}\PY{+w}{ }\PY{p}{\PYZob{}}
\PY{+w}{        }\PY{k}{try}\PY{+w}{ }\PY{p}{\PYZob{}}
\PY{+w}{            }\PY{n}{System}\PY{p}{.}\PY{n+na}{out}\PY{p}{.}\PY{n+na}{println}\PY{p}{(}\PY{l+s}{\PYZdq{}}\PY{l+s}{try}\PY{l+s}{\PYZdq{}}\PY{p}{)}\PY{p}{;}
\PY{+w}{            }\PY{k}{throw}\PY{+w}{ }\PY{k}{new}\PY{+w}{ }\PY{n}{RuntimeException}\PY{p}{(}\PY{l+s}{\PYZdq{}}\PY{l+s}{fail}\PY{l+s}{\PYZdq{}}\PY{p}{)}\PY{p}{;}
\PY{+w}{        }\PY{p}{\PYZcb{}}\PY{+w}{ }\PY{k}{finally}\PY{+w}{ }\PY{p}{\PYZob{}}
\PY{+w}{            }\PY{n}{System}\PY{p}{.}\PY{n+na}{out}\PY{p}{.}\PY{n+na}{println}\PY{p}{(}\PY{l+s}{\PYZdq{}}\PY{l+s}{finally always runs}\PY{l+s}{\PYZdq{}}\PY{p}{)}\PY{p}{;}
\PY{+w}{        }\PY{p}{\PYZcb{}}
\PY{+w}{    }\PY{p}{\PYZcb{}}
\PY{p}{\PYZcb{}}
\end{Verbatim}
\end{codebox}

\subsection[How return in finally Swallows Exceptions]{How \code{return} in \code{finally} Swallows Exceptions}
\label{h:6:how-return-in-finally-swallows-exceptions}
This is a classic trap: if \code{finally} itself contains a \code{return} (or a \code{throw}), it \textbf{overrides} whatever the \code{try} or \code{catch} block was about to do — including silently discarding an in-flight exception.

\begin{codebox}
\begin{Verbatim}[commandchars=\\\{\}]
\PY{c+c1}{// BROKEN: the exception disappears without a trace}
\PY{k+kd}{public}\PY{+w}{ }\PY{k+kd}{class} \PY{n+nc}{SwallowedException}\PY{+w}{ }\PY{p}{\PYZob{}}
\PY{+w}{    }\PY{k+kd}{static}\PY{+w}{ }\PY{k+kt}{int}\PY{+w}{ }\PY{n+nf}{risky}\PY{p}{(}\PY{p}{)}\PY{+w}{ }\PY{p}{\PYZob{}}
\PY{+w}{        }\PY{k}{try}\PY{+w}{ }\PY{p}{\PYZob{}}
\PY{+w}{            }\PY{k}{throw}\PY{+w}{ }\PY{k}{new}\PY{+w}{ }\PY{n}{IllegalStateException}\PY{p}{(}\PY{l+s}{\PYZdq{}}\PY{l+s}{real problem}\PY{l+s}{\PYZdq{}}\PY{p}{)}\PY{p}{;}
\PY{+w}{        }\PY{p}{\PYZcb{}}\PY{+w}{ }\PY{k}{finally}\PY{+w}{ }\PY{p}{\PYZob{}}
\PY{+w}{            }\PY{k}{return}\PY{+w}{ }\PY{l+m+mi}{42}\PY{p}{;}\PY{+w}{ }\PY{c+c1}{// this return wins; the exception above is discarded}
\PY{+w}{        }\PY{p}{\PYZcb{}}
\PY{+w}{    }\PY{p}{\PYZcb{}}

\PY{+w}{    }\PY{k+kd}{public}\PY{+w}{ }\PY{k+kd}{static}\PY{+w}{ }\PY{k+kt}{void}\PY{+w}{ }\PY{n+nf}{main}\PY{p}{(}\PY{n}{String}\PY{o}{[}\PY{o}{]}\PY{+w}{ }\PY{n}{args}\PY{p}{)}\PY{+w}{ }\PY{p}{\PYZob{}}
\PY{+w}{        }\PY{n}{System}\PY{p}{.}\PY{n+na}{out}\PY{p}{.}\PY{n+na}{println}\PY{p}{(}\PY{n}{risky}\PY{p}{(}\PY{p}{)}\PY{p}{)}\PY{p}{;}\PY{+w}{ }\PY{c+c1}{// prints 42, exception is gone forever}
\PY{+w}{    }\PY{p}{\PYZcb{}}
\PY{p}{\PYZcb{}}
\end{Verbatim}
\end{codebox}

\begin{codebox}
\begin{Verbatim}[commandchars=\\\{\}]
\PY{c+c1}{// FIXED: never put a return (or throw) inside finally}
\PY{k+kd}{public}\PY{+w}{ }\PY{k+kd}{class} \PY{n+nc}{NoSwallowedException}\PY{+w}{ }\PY{p}{\PYZob{}}
\PY{+w}{    }\PY{k+kd}{static}\PY{+w}{ }\PY{k+kt}{int}\PY{+w}{ }\PY{n+nf}{risky}\PY{p}{(}\PY{p}{)}\PY{+w}{ }\PY{p}{\PYZob{}}
\PY{+w}{        }\PY{k}{try}\PY{+w}{ }\PY{p}{\PYZob{}}
\PY{+w}{            }\PY{k}{throw}\PY{+w}{ }\PY{k}{new}\PY{+w}{ }\PY{n}{IllegalStateException}\PY{p}{(}\PY{l+s}{\PYZdq{}}\PY{l+s}{real problem}\PY{l+s}{\PYZdq{}}\PY{p}{)}\PY{p}{;}
\PY{+w}{        }\PY{p}{\PYZcb{}}\PY{+w}{ }\PY{k}{finally}\PY{+w}{ }\PY{p}{\PYZob{}}
\PY{+w}{            }\PY{n}{System}\PY{p}{.}\PY{n+na}{out}\PY{p}{.}\PY{n+na}{println}\PY{p}{(}\PY{l+s}{\PYZdq{}}\PY{l+s}{cleanup only, no return here}\PY{l+s}{\PYZdq{}}\PY{p}{)}\PY{p}{;}
\PY{+w}{        }\PY{p}{\PYZcb{}}
\PY{+w}{    }\PY{p}{\PYZcb{}}
\PY{p}{\PYZcb{}}
\end{Verbatim}
\end{codebox}

Rule of thumb: \code{finally} blocks should only do cleanup (closing resources, releasing locks, logging). They should never \code{return}, \code{break}, \code{continue}, or \code{throw}.

\subsection[Multi-Catch]{Multi-Catch}
\label{h:6:multi-catch}
When several exception types need the same handling, \code{catch} can list them with \code{\textbar{}} instead of duplicating the block.

\begin{codebox}
\begin{Verbatim}[commandchars=\\\{\}]
\PY{k+kd}{public}\PY{+w}{ }\PY{k+kd}{class} \PY{n+nc}{MultiCatchDemo}\PY{+w}{ }\PY{p}{\PYZob{}}
\PY{+w}{    }\PY{k+kd}{static}\PY{+w}{ }\PY{k+kt}{void}\PY{+w}{ }\PY{n+nf}{parse}\PY{p}{(}\PY{n}{String}\PY{+w}{ }\PY{n}{input}\PY{p}{)}\PY{+w}{ }\PY{p}{\PYZob{}}
\PY{+w}{        }\PY{k}{try}\PY{+w}{ }\PY{p}{\PYZob{}}
\PY{+w}{            }\PY{k+kt}{int}\PY{+w}{ }\PY{n}{value}\PY{+w}{ }\PY{o}{=}\PY{+w}{ }\PY{n}{Integer}\PY{p}{.}\PY{n+na}{parseInt}\PY{p}{(}\PY{n}{input}\PY{p}{)}\PY{p}{;}
\PY{+w}{            }\PY{n}{System}\PY{p}{.}\PY{n+na}{out}\PY{p}{.}\PY{n+na}{println}\PY{p}{(}\PY{l+m+mi}{100}\PY{+w}{ }\PY{o}{/}\PY{+w}{ }\PY{n}{value}\PY{p}{)}\PY{p}{;}
\PY{+w}{        }\PY{p}{\PYZcb{}}\PY{+w}{ }\PY{k}{catch}\PY{+w}{ }\PY{p}{(}\PY{n}{NumberFormatException}\PY{+w}{ }\PY{o}{|}\PY{+w}{ }\PY{n}{ArithmeticException}\PY{+w}{ }\PY{n}{e}\PY{p}{)}\PY{+w}{ }\PY{p}{\PYZob{}}
\PY{+w}{            }\PY{n}{System}\PY{p}{.}\PY{n+na}{out}\PY{p}{.}\PY{n+na}{println}\PY{p}{(}\PY{l+s}{\PYZdq{}}\PY{l+s}{Invalid input or division error: }\PY{l+s}{\PYZdq{}}\PY{+w}{ }\PY{o}{+}\PY{+w}{ }\PY{n}{e}\PY{p}{.}\PY{n+na}{getMessage}\PY{p}{(}\PY{p}{)}\PY{p}{)}\PY{p}{;}
\PY{+w}{        }\PY{p}{\PYZcb{}}
\PY{+w}{    }\PY{p}{\PYZcb{}}
\PY{p}{\PYZcb{}}
\end{Verbatim}
\end{codebox}

Note: the caught variable (\code{e} here) is effectively \code{final}, and its static type is the least upper bound of the listed types — you cannot reassign it, and you cannot catch a type and its subtype together in the same multi-catch (the compiler will reject that as redundant).

\subsection[Rethrow with Precise Type Inference]{Rethrow with Precise Type Inference}
\label{h:6:rethrow-with-precise-type-inference}
Since Java 7, if a \code{catch} block only rethrows the caught exception (without reassigning it to a broader type), the compiler tracks the \textbf{precise} type that was actually thrown, not just the declared catch type. This lets a method declare \code{throws~\allowbreak{}IOException,\allowbreak{}~\allowbreak{}SQLException} instead of the broader \code{throws~\allowbreak{}Exception}.

\begin{codebox}
\begin{Verbatim}[commandchars=\\\{\}]
\PY{k+kd}{public}\PY{+w}{ }\PY{k+kd}{class} \PY{n+nc}{PreciseRethrow}\PY{+w}{ }\PY{p}{\PYZob{}}
\PY{+w}{    }\PY{k+kd}{static}\PY{+w}{ }\PY{k+kt}{void}\PY{+w}{ }\PY{n+nf}{doWork}\PY{p}{(}\PY{k+kt}{boolean}\PY{+w}{ }\PY{n}{useDb}\PY{p}{)}\PY{+w}{ }\PY{k+kd}{throws}\PY{+w}{ }\PY{n}{IOException}\PY{p}{,}\PY{+w}{ }\PY{n}{SQLException}\PY{+w}{ }\PY{p}{\PYZob{}}
\PY{+w}{        }\PY{k}{try}\PY{+w}{ }\PY{p}{\PYZob{}}
\PY{+w}{            }\PY{k}{if}\PY{+w}{ }\PY{p}{(}\PY{n}{useDb}\PY{p}{)}\PY{+w}{ }\PY{p}{\PYZob{}}
\PY{+w}{                }\PY{k}{throw}\PY{+w}{ }\PY{k}{new}\PY{+w}{ }\PY{n}{SQLException}\PY{p}{(}\PY{l+s}{\PYZdq{}}\PY{l+s}{db error}\PY{l+s}{\PYZdq{}}\PY{p}{)}\PY{p}{;}
\PY{+w}{            }\PY{p}{\PYZcb{}}\PY{+w}{ }\PY{k}{else}\PY{+w}{ }\PY{p}{\PYZob{}}
\PY{+w}{                }\PY{k}{throw}\PY{+w}{ }\PY{k}{new}\PY{+w}{ }\PY{n}{IOException}\PY{p}{(}\PY{l+s}{\PYZdq{}}\PY{l+s}{io error}\PY{l+s}{\PYZdq{}}\PY{p}{)}\PY{p}{;}
\PY{+w}{            }\PY{p}{\PYZcb{}}
\PY{+w}{        }\PY{p}{\PYZcb{}}\PY{+w}{ }\PY{k}{catch}\PY{+w}{ }\PY{p}{(}\PY{n}{Exception}\PY{+w}{ }\PY{n}{e}\PY{p}{)}\PY{+w}{ }\PY{p}{\PYZob{}}
\PY{+w}{            }\PY{c+c1}{// e is caught as Exception, but the compiler knows only}
\PY{+w}{            }\PY{c+c1}{// IOException or SQLException could actually reach here,}
\PY{+w}{            }\PY{c+c1}{// so this rethrow is allowed without \PYZdq{}throws Exception\PYZdq{}.}
\PY{+w}{            }\PY{k}{throw}\PY{+w}{ }\PY{n}{e}\PY{p}{;}
\PY{+w}{        }\PY{p}{\PYZcb{}}
\PY{+w}{    }\PY{p}{\PYZcb{}}
\PY{p}{\PYZcb{}}
\end{Verbatim}
\end{codebox}

If you reassign \code{e} inside the catch block, the compiler loses this precision and falls back to the declared catch type.

\subsection[Exception Chaining (Cause)]{Exception Chaining (Cause)}
\label{h:6:exception-chaining-cause}
When you catch a low-level exception and throw a higher-level one, always pass the original as the \textbf{cause}. This preserves the full stack trace instead of hiding the real reason.

\begin{codebox}
\begin{Verbatim}[commandchars=\\\{\}]
\PY{c+c1}{// BROKEN: original cause is lost}
\PY{k+kd}{public}\PY{+w}{ }\PY{k+kd}{class} \PY{n+nc}{LostCause}\PY{+w}{ }\PY{p}{\PYZob{}}
\PY{+w}{    }\PY{k+kd}{static}\PY{+w}{ }\PY{k+kt}{void}\PY{+w}{ }\PY{n+nf}{loadUser}\PY{p}{(}\PY{n}{String}\PY{+w}{ }\PY{n}{id}\PY{p}{)}\PY{+w}{ }\PY{p}{\PYZob{}}
\PY{+w}{        }\PY{k}{try}\PY{+w}{ }\PY{p}{\PYZob{}}
\PY{+w}{            }\PY{n}{fetchFromDatabase}\PY{p}{(}\PY{n}{id}\PY{p}{)}\PY{p}{;}
\PY{+w}{        }\PY{p}{\PYZcb{}}\PY{+w}{ }\PY{k}{catch}\PY{+w}{ }\PY{p}{(}\PY{n}{SQLException}\PY{+w}{ }\PY{n}{e}\PY{p}{)}\PY{+w}{ }\PY{p}{\PYZob{}}
\PY{+w}{            }\PY{k}{throw}\PY{+w}{ }\PY{k}{new}\PY{+w}{ }\PY{n}{RuntimeException}\PY{p}{(}\PY{l+s}{\PYZdq{}}\PY{l+s}{Could not load user }\PY{l+s}{\PYZdq{}}\PY{+w}{ }\PY{o}{+}\PY{+w}{ }\PY{n}{id}\PY{p}{)}\PY{p}{;}\PY{+w}{ }\PY{c+c1}{// cause dropped!}
\PY{+w}{        }\PY{p}{\PYZcb{}}
\PY{+w}{    }\PY{p}{\PYZcb{}}
\PY{+w}{    }\PY{k+kd}{static}\PY{+w}{ }\PY{k+kt}{void}\PY{+w}{ }\PY{n+nf}{fetchFromDatabase}\PY{p}{(}\PY{n}{String}\PY{+w}{ }\PY{n}{id}\PY{p}{)}\PY{+w}{ }\PY{k+kd}{throws}\PY{+w}{ }\PY{n}{SQLException}\PY{+w}{ }\PY{p}{\PYZob{}}
\PY{+w}{        }\PY{k}{throw}\PY{+w}{ }\PY{k}{new}\PY{+w}{ }\PY{n}{SQLException}\PY{p}{(}\PY{l+s}{\PYZdq{}}\PY{l+s}{connection refused}\PY{l+s}{\PYZdq{}}\PY{p}{)}\PY{p}{;}
\PY{+w}{    }\PY{p}{\PYZcb{}}
\PY{p}{\PYZcb{}}
\end{Verbatim}
\end{codebox}

\begin{codebox}
\begin{Verbatim}[commandchars=\\\{\}]
\PY{c+c1}{// FIXED: chain the cause}
\PY{k+kd}{public}\PY{+w}{ }\PY{k+kd}{class} \PY{n+nc}{PreservedCause}\PY{+w}{ }\PY{p}{\PYZob{}}
\PY{+w}{    }\PY{k+kd}{static}\PY{+w}{ }\PY{k+kt}{void}\PY{+w}{ }\PY{n+nf}{loadUser}\PY{p}{(}\PY{n}{String}\PY{+w}{ }\PY{n}{id}\PY{p}{)}\PY{+w}{ }\PY{p}{\PYZob{}}
\PY{+w}{        }\PY{k}{try}\PY{+w}{ }\PY{p}{\PYZob{}}
\PY{+w}{            }\PY{n}{fetchFromDatabase}\PY{p}{(}\PY{n}{id}\PY{p}{)}\PY{p}{;}
\PY{+w}{        }\PY{p}{\PYZcb{}}\PY{+w}{ }\PY{k}{catch}\PY{+w}{ }\PY{p}{(}\PY{n}{SQLException}\PY{+w}{ }\PY{n}{e}\PY{p}{)}\PY{+w}{ }\PY{p}{\PYZob{}}
\PY{+w}{            }\PY{k}{throw}\PY{+w}{ }\PY{k}{new}\PY{+w}{ }\PY{n}{RuntimeException}\PY{p}{(}\PY{l+s}{\PYZdq{}}\PY{l+s}{Could not load user }\PY{l+s}{\PYZdq{}}\PY{+w}{ }\PY{o}{+}\PY{+w}{ }\PY{n}{id}\PY{p}{,}\PY{+w}{ }\PY{n}{e}\PY{p}{)}\PY{p}{;}\PY{+w}{ }\PY{c+c1}{// cause preserved}
\PY{+w}{        }\PY{p}{\PYZcb{}}
\PY{+w}{    }\PY{p}{\PYZcb{}}
\PY{+w}{    }\PY{k+kd}{static}\PY{+w}{ }\PY{k+kt}{void}\PY{+w}{ }\PY{n+nf}{fetchFromDatabase}\PY{p}{(}\PY{n}{String}\PY{+w}{ }\PY{n}{id}\PY{p}{)}\PY{+w}{ }\PY{k+kd}{throws}\PY{+w}{ }\PY{n}{SQLException}\PY{+w}{ }\PY{p}{\PYZob{}}
\PY{+w}{        }\PY{k}{throw}\PY{+w}{ }\PY{k}{new}\PY{+w}{ }\PY{n}{SQLException}\PY{p}{(}\PY{l+s}{\PYZdq{}}\PY{l+s}{connection refused}\PY{l+s}{\PYZdq{}}\PY{p}{)}\PY{p}{;}
\PY{+w}{    }\PY{p}{\PYZcb{}}
\PY{p}{\PYZcb{}}
\end{Verbatim}
\end{codebox}

You can also inspect the cause chain programmatically with \code{getCause()}, and Java 9+ adds \code{Throwable.\allowbreak{}get\allowbreak{}Suppressed()} for suppressed exceptions (covered later).

\begin{codebox}
\begin{Verbatim}[commandchars=\\\{\}]
\PY{k}{try}\PY{+w}{ }\PY{p}{\PYZob{}}
\PY{+w}{    }\PY{n}{PreservedCause}\PY{p}{.}\PY{n+na}{loadUser}\PY{p}{(}\PY{l+s}{\PYZdq{}}\PY{l+s}{42}\PY{l+s}{\PYZdq{}}\PY{p}{)}\PY{p}{;}
\PY{p}{\PYZcb{}}\PY{+w}{ }\PY{k}{catch}\PY{+w}{ }\PY{p}{(}\PY{n}{RuntimeException}\PY{+w}{ }\PY{n}{e}\PY{p}{)}\PY{+w}{ }\PY{p}{\PYZob{}}
\PY{+w}{    }\PY{n}{Throwable}\PY{+w}{ }\PY{n}{cause}\PY{+w}{ }\PY{o}{=}\PY{+w}{ }\PY{n}{e}\PY{p}{.}\PY{n+na}{getCause}\PY{p}{(}\PY{p}{)}\PY{p}{;}
\PY{+w}{    }\PY{n}{System}\PY{p}{.}\PY{n+na}{out}\PY{p}{.}\PY{n+na}{println}\PY{p}{(}\PY{l+s}{\PYZdq{}}\PY{l+s}{Root cause: }\PY{l+s}{\PYZdq{}}\PY{+w}{ }\PY{o}{+}\PY{+w}{ }\PY{n}{cause}\PY{p}{)}\PY{p}{;}
\PY{p}{\PYZcb{}}
\end{Verbatim}
\end{codebox}

\subsection[Custom Exception Design]{Custom Exception Design}
\label{h:6:custom-exception-design}
Design custom exceptions when the standard library types don’t communicate enough domain meaning. Keep them simple, immutable, and always provide constructors that support chaining.

\begin{codebox}
\begin{Verbatim}[commandchars=\\\{\}]
\PY{k+kd}{public}\PY{+w}{ }\PY{k+kd}{class} \PY{n+nc}{InsufficientFundsException}\PY{+w}{ }\PY{k+kd}{extends}\PY{+w}{ }\PY{n}{Exception}\PY{+w}{ }\PY{p}{\PYZob{}}
\PY{+w}{    }\PY{k+kd}{private}\PY{+w}{ }\PY{k+kd}{final}\PY{+w}{ }\PY{n}{BigDecimal}\PY{+w}{ }\PY{n}{shortfall}\PY{p}{;}

\PY{+w}{    }\PY{k+kd}{public}\PY{+w}{ }\PY{n+nf}{InsufficientFundsException}\PY{p}{(}\PY{n}{String}\PY{+w}{ }\PY{n}{message}\PY{p}{,}\PY{+w}{ }\PY{n}{BigDecimal}\PY{+w}{ }\PY{n}{shortfall}\PY{p}{)}\PY{+w}{ }\PY{p}{\PYZob{}}
\PY{+w}{        }\PY{k+kd}{super}\PY{p}{(}\PY{n}{message}\PY{p}{)}\PY{p}{;}
\PY{+w}{        }\PY{k}{this}\PY{p}{.}\PY{n+na}{shortfall}\PY{+w}{ }\PY{o}{=}\PY{+w}{ }\PY{n}{shortfall}\PY{p}{;}
\PY{+w}{    }\PY{p}{\PYZcb{}}

\PY{+w}{    }\PY{k+kd}{public}\PY{+w}{ }\PY{n+nf}{InsufficientFundsException}\PY{p}{(}\PY{n}{String}\PY{+w}{ }\PY{n}{message}\PY{p}{,}\PY{+w}{ }\PY{n}{BigDecimal}\PY{+w}{ }\PY{n}{shortfall}\PY{p}{,}\PY{+w}{ }\PY{n}{Throwable}\PY{+w}{ }\PY{n}{cause}\PY{p}{)}\PY{+w}{ }\PY{p}{\PYZob{}}
\PY{+w}{        }\PY{k+kd}{super}\PY{p}{(}\PY{n}{message}\PY{p}{,}\PY{+w}{ }\PY{n}{cause}\PY{p}{)}\PY{p}{;}
\PY{+w}{        }\PY{k}{this}\PY{p}{.}\PY{n+na}{shortfall}\PY{+w}{ }\PY{o}{=}\PY{+w}{ }\PY{n}{shortfall}\PY{p}{;}
\PY{+w}{    }\PY{p}{\PYZcb{}}

\PY{+w}{    }\PY{k+kd}{public}\PY{+w}{ }\PY{n}{BigDecimal}\PY{+w}{ }\PY{n+nf}{getShortfall}\PY{p}{(}\PY{p}{)}\PY{+w}{ }\PY{p}{\PYZob{}}
\PY{+w}{        }\PY{k}{return}\PY{+w}{ }\PY{n}{shortfall}\PY{p}{;}
\PY{+w}{    }\PY{p}{\PYZcb{}}
\PY{p}{\PYZcb{}}
\end{Verbatim}
\end{codebox}

Guidelines for custom exceptions:

\begin{itemize}
\item 
Extend \code{Exception} for checked, recoverable business errors (e.g., \code{Insufficient\allowbreak{}Funds\allowbreak{}Exception}); extend \code{RuntimeException} for programming errors or when using functional/stream code.

\item 
Always provide a \code{(\allowbreak{}String~\allowbreak{}message,\allowbreak{}~\allowbreak{}Throwable~\allowbreak{}cause)} constructor for chaining.

\item 
Add extra fields (like \code{shortfall} above) only when calling code actually needs them programmatically — don’t just stuff everything into the message string.

\item 
Avoid deep exception hierarchies. One or two custom types per module is usually enough.

\item 
Don’t make exceptions mutable; they should describe a moment in time.

\end{itemize}

\subsection[NullPointerException Helpful Messages (JDK 14+)]{\code{NullPointerException} Helpful Messages (JDK 14+)}
\label{h:6:nullpointerexception-helpful-messages-jdk-14}
Since JDK 14 (stabilized as a default from JDK 15), the JVM can generate a \textbf{helpful NPE message} that names exactly which variable or method call was null, instead of just a bare stack trace line.

\begin{codebox}
\begin{Verbatim}[commandchars=\\\{\}]
\PY{k+kd}{public}\PY{+w}{ }\PY{k+kd}{class} \PY{n+nc}{HelpfulNpeDemo}\PY{+w}{ }\PY{p}{\PYZob{}}
\PY{+w}{    }\PY{k+kd}{record} \PY{n+nc}{Address}\PY{p}{(}\PY{n}{String}\PY{+w}{ }\PY{n}{city}\PY{p}{)}\PY{+w}{ }\PY{p}{\PYZob{}}\PY{p}{\PYZcb{}}
\PY{+w}{    }\PY{k+kd}{record} \PY{n+nc}{Person}\PY{p}{(}\PY{n}{Address}\PY{+w}{ }\PY{n}{address}\PY{p}{)}\PY{+w}{ }\PY{p}{\PYZob{}}\PY{p}{\PYZcb{}}

\PY{+w}{    }\PY{k+kd}{public}\PY{+w}{ }\PY{k+kd}{static}\PY{+w}{ }\PY{k+kt}{void}\PY{+w}{ }\PY{n+nf}{main}\PY{p}{(}\PY{n}{String}\PY{o}{[}\PY{o}{]}\PY{+w}{ }\PY{n}{args}\PY{p}{)}\PY{+w}{ }\PY{p}{\PYZob{}}
\PY{+w}{        }\PY{n}{Person}\PY{+w}{ }\PY{n}{person}\PY{+w}{ }\PY{o}{=}\PY{+w}{ }\PY{k}{new}\PY{+w}{ }\PY{n}{Person}\PY{p}{(}\PY{k+kc}{null}\PY{p}{)}\PY{p}{;}
\PY{+w}{        }\PY{n}{System}\PY{p}{.}\PY{n+na}{out}\PY{p}{.}\PY{n+na}{println}\PY{p}{(}\PY{n}{person}\PY{p}{.}\PY{n+na}{address}\PY{p}{(}\PY{p}{)}\PY{p}{.}\PY{n+na}{city}\PY{p}{(}\PY{p}{)}\PY{p}{)}\PY{p}{;}\PY{+w}{ }\PY{c+c1}{// NPE here}
\PY{+w}{    }\PY{p}{\PYZcb{}}
\PY{p}{\PYZcb{}}
\end{Verbatim}
\end{codebox}

Without the feature you’d just see \code{NullPointerException} with no detail. With helpful NPE messages enabled (default since Java 15), you get:

\begin{codebox}
\begin{Verbatim}
Exception in thread "main" java.lang.NullPointerException:
    Cannot invoke "HelpfulNpeDemo$Address.city()" because the return value of
    "HelpfulNpeDemo$Person.address()" is null
	at HelpfulNpeDemo.main(HelpfulNpeDemo.java:8)
\end{Verbatim}
\end{codebox}

This tells you precisely which call returned \code{null}, which is invaluable when a line has several chained calls. If you’re on an older JVM or it’s disabled, enable it with \code{-\allowbreak{}XX:+\allowbreak{}Show\allowbreak{}Code\allowbreak{}Details\allowbreak{}In\allowbreak{}Exception\allowbreak{}Messages} (unnecessary on modern JDKs, where it’s on by default).

\subsection[Suppressed Exceptions]{Suppressed Exceptions}
\label{h:6:suppressed-exceptions}
When a \code{try-\allowbreak{}with-\allowbreak{}resources} block throws an exception in the body \emph{and} closing a resource also throws, Java doesn’t discard either one. The body’s exception is thrown as the primary exception, and the close-time exception is attached to it as a \textbf{suppressed exception}.

\begin{codebox}
\begin{Verbatim}[commandchars=\\\{\}]
\PY{k+kd}{public}\PY{+w}{ }\PY{k+kd}{class} \PY{n+nc}{SuppressedDemo}\PY{+w}{ }\PY{p}{\PYZob{}}
\PY{+w}{    }\PY{k+kd}{static}\PY{+w}{ }\PY{k+kd}{class} \PY{n+nc}{Resource}\PY{+w}{ }\PY{k+kd}{implements}\PY{+w}{ }\PY{n}{AutoCloseable}\PY{+w}{ }\PY{p}{\PYZob{}}
\PY{+w}{        }\PY{k+kd}{public}\PY{+w}{ }\PY{k+kt}{void}\PY{+w}{ }\PY{n+nf}{use}\PY{p}{(}\PY{p}{)}\PY{+w}{ }\PY{p}{\PYZob{}}
\PY{+w}{            }\PY{k}{throw}\PY{+w}{ }\PY{k}{new}\PY{+w}{ }\PY{n}{RuntimeException}\PY{p}{(}\PY{l+s}{\PYZdq{}}\PY{l+s}{failure during use}\PY{l+s}{\PYZdq{}}\PY{p}{)}\PY{p}{;}
\PY{+w}{        }\PY{p}{\PYZcb{}}
\PY{+w}{        }\PY{n+nd}{@Override}
\PY{+w}{        }\PY{k+kd}{public}\PY{+w}{ }\PY{k+kt}{void}\PY{+w}{ }\PY{n+nf}{close}\PY{p}{(}\PY{p}{)}\PY{+w}{ }\PY{p}{\PYZob{}}
\PY{+w}{            }\PY{k}{throw}\PY{+w}{ }\PY{k}{new}\PY{+w}{ }\PY{n}{RuntimeException}\PY{p}{(}\PY{l+s}{\PYZdq{}}\PY{l+s}{failure during close}\PY{l+s}{\PYZdq{}}\PY{p}{)}\PY{p}{;}
\PY{+w}{        }\PY{p}{\PYZcb{}}
\PY{+w}{    }\PY{p}{\PYZcb{}}

\PY{+w}{    }\PY{k+kd}{public}\PY{+w}{ }\PY{k+kd}{static}\PY{+w}{ }\PY{k+kt}{void}\PY{+w}{ }\PY{n+nf}{main}\PY{p}{(}\PY{n}{String}\PY{o}{[}\PY{o}{]}\PY{+w}{ }\PY{n}{args}\PY{p}{)}\PY{+w}{ }\PY{p}{\PYZob{}}
\PY{+w}{        }\PY{k}{try}\PY{+w}{ }\PY{p}{(}\PY{n}{Resource}\PY{+w}{ }\PY{n}{r}\PY{+w}{ }\PY{o}{=}\PY{+w}{ }\PY{k}{new}\PY{+w}{ }\PY{n}{Resource}\PY{p}{(}\PY{p}{)}\PY{p}{)}\PY{+w}{ }\PY{p}{\PYZob{}}
\PY{+w}{            }\PY{n}{r}\PY{p}{.}\PY{n+na}{use}\PY{p}{(}\PY{p}{)}\PY{p}{;}
\PY{+w}{        }\PY{p}{\PYZcb{}}\PY{+w}{ }\PY{k}{catch}\PY{+w}{ }\PY{p}{(}\PY{n}{RuntimeException}\PY{+w}{ }\PY{n}{e}\PY{p}{)}\PY{+w}{ }\PY{p}{\PYZob{}}
\PY{+w}{            }\PY{n}{System}\PY{p}{.}\PY{n+na}{out}\PY{p}{.}\PY{n+na}{println}\PY{p}{(}\PY{l+s}{\PYZdq{}}\PY{l+s}{Primary: }\PY{l+s}{\PYZdq{}}\PY{+w}{ }\PY{o}{+}\PY{+w}{ }\PY{n}{e}\PY{p}{.}\PY{n+na}{getMessage}\PY{p}{(}\PY{p}{)}\PY{p}{)}\PY{p}{;}
\PY{+w}{            }\PY{k}{for}\PY{+w}{ }\PY{p}{(}\PY{n}{Throwable}\PY{+w}{ }\PY{n}{suppressed}\PY{+w}{ }\PY{p}{:}\PY{+w}{ }\PY{n}{e}\PY{p}{.}\PY{n+na}{getSuppressed}\PY{p}{(}\PY{p}{)}\PY{p}{)}\PY{+w}{ }\PY{p}{\PYZob{}}
\PY{+w}{                }\PY{n}{System}\PY{p}{.}\PY{n+na}{out}\PY{p}{.}\PY{n+na}{println}\PY{p}{(}\PY{l+s}{\PYZdq{}}\PY{l+s}{Suppressed: }\PY{l+s}{\PYZdq{}}\PY{+w}{ }\PY{o}{+}\PY{+w}{ }\PY{n}{suppressed}\PY{p}{.}\PY{n+na}{getMessage}\PY{p}{(}\PY{p}{)}\PY{p}{)}\PY{p}{;}
\PY{+w}{            }\PY{p}{\PYZcb{}}
\PY{+w}{        }\PY{p}{\PYZcb{}}
\PY{+w}{    }\PY{p}{\PYZcb{}}
\PY{p}{\PYZcb{}}
\end{Verbatim}
\end{codebox}

Output:

\begin{codebox}
\begin{Verbatim}
Primary: failure during use
Suppressed: failure during close
\end{Verbatim}
\end{codebox}

Without try-with-resources, if you manually close a resource inside \code{finally}, the close-time exception would simply \emph{replace} the original one — losing the real cause. Suppressed exceptions solve exactly that problem.

\section[Try-with-Resources]{Try-with-Resources}
\label{h:6:try-with-resources}
\code{try-\allowbreak{}with-\allowbreak{}resources} (introduced in Java 7) automatically closes resources for you, even if an exception is thrown, without needing a manual \code{finally} block.

\begin{codebox}
\begin{Verbatim}[commandchars=\\\{\}]
\PY{c+c1}{// BROKEN: manual close, easy to get wrong, and swallows exceptions from use()}
\PY{k+kd}{public}\PY{+w}{ }\PY{k+kd}{class} \PY{n+nc}{ManualClose}\PY{+w}{ }\PY{p}{\PYZob{}}
\PY{+w}{    }\PY{k+kd}{static}\PY{+w}{ }\PY{k+kt}{void}\PY{+w}{ }\PY{n+nf}{copy}\PY{p}{(}\PY{n}{String}\PY{+w}{ }\PY{n}{from}\PY{p}{,}\PY{+w}{ }\PY{n}{String}\PY{+w}{ }\PY{n}{to}\PY{p}{)}\PY{+w}{ }\PY{k+kd}{throws}\PY{+w}{ }\PY{n}{IOException}\PY{+w}{ }\PY{p}{\PYZob{}}
\PY{+w}{        }\PY{n}{FileInputStream}\PY{+w}{ }\PY{n}{in}\PY{+w}{ }\PY{o}{=}\PY{+w}{ }\PY{k+kc}{null}\PY{p}{;}
\PY{+w}{        }\PY{n}{FileOutputStream}\PY{+w}{ }\PY{n}{out}\PY{+w}{ }\PY{o}{=}\PY{+w}{ }\PY{k+kc}{null}\PY{p}{;}
\PY{+w}{        }\PY{k}{try}\PY{+w}{ }\PY{p}{\PYZob{}}
\PY{+w}{            }\PY{n}{in}\PY{+w}{ }\PY{o}{=}\PY{+w}{ }\PY{k}{new}\PY{+w}{ }\PY{n}{FileInputStream}\PY{p}{(}\PY{n}{from}\PY{p}{)}\PY{p}{;}
\PY{+w}{            }\PY{n}{out}\PY{+w}{ }\PY{o}{=}\PY{+w}{ }\PY{k}{new}\PY{+w}{ }\PY{n}{FileOutputStream}\PY{p}{(}\PY{n}{to}\PY{p}{)}\PY{p}{;}
\PY{+w}{            }\PY{n}{in}\PY{p}{.}\PY{n+na}{transferTo}\PY{p}{(}\PY{n}{out}\PY{p}{)}\PY{p}{;}
\PY{+w}{        }\PY{p}{\PYZcb{}}\PY{+w}{ }\PY{k}{finally}\PY{+w}{ }\PY{p}{\PYZob{}}
\PY{+w}{            }\PY{k}{if}\PY{+w}{ }\PY{p}{(}\PY{n}{out}\PY{+w}{ }\PY{o}{!}\PY{o}{=}\PY{+w}{ }\PY{k+kc}{null}\PY{p}{)}\PY{+w}{ }\PY{n}{out}\PY{p}{.}\PY{n+na}{close}\PY{p}{(}\PY{p}{)}\PY{p}{;}\PY{+w}{ }\PY{c+c1}{// if this throws, exception from in.close() below is lost}
\PY{+w}{            }\PY{k}{if}\PY{+w}{ }\PY{p}{(}\PY{n}{in}\PY{+w}{ }\PY{o}{!}\PY{o}{=}\PY{+w}{ }\PY{k+kc}{null}\PY{p}{)}\PY{+w}{ }\PY{n}{in}\PY{p}{.}\PY{n+na}{close}\PY{p}{(}\PY{p}{)}\PY{p}{;}
\PY{+w}{        }\PY{p}{\PYZcb{}}
\PY{+w}{    }\PY{p}{\PYZcb{}}
\PY{p}{\PYZcb{}}
\end{Verbatim}
\end{codebox}

\begin{codebox}
\begin{Verbatim}[commandchars=\\\{\}]
\PY{c+c1}{// FIXED: try\PYZhy{}with\PYZhy{}resources, both resources closed automatically, in reverse order}
\PY{k+kd}{public}\PY{+w}{ }\PY{k+kd}{class} \PY{n+nc}{TryWithResourcesDemo}\PY{+w}{ }\PY{p}{\PYZob{}}
\PY{+w}{    }\PY{k+kd}{static}\PY{+w}{ }\PY{k+kt}{void}\PY{+w}{ }\PY{n+nf}{copy}\PY{p}{(}\PY{n}{String}\PY{+w}{ }\PY{n}{from}\PY{p}{,}\PY{+w}{ }\PY{n}{String}\PY{+w}{ }\PY{n}{to}\PY{p}{)}\PY{+w}{ }\PY{k+kd}{throws}\PY{+w}{ }\PY{n}{IOException}\PY{+w}{ }\PY{p}{\PYZob{}}
\PY{+w}{        }\PY{k}{try}\PY{+w}{ }\PY{p}{(}\PY{n}{FileInputStream}\PY{+w}{ }\PY{n}{in}\PY{+w}{ }\PY{o}{=}\PY{+w}{ }\PY{k}{new}\PY{+w}{ }\PY{n}{FileInputStream}\PY{p}{(}\PY{n}{from}\PY{p}{)}\PY{p}{;}
\PY{+w}{             }\PY{n}{FileOutputStream}\PY{+w}{ }\PY{n}{out}\PY{+w}{ }\PY{o}{=}\PY{+w}{ }\PY{k}{new}\PY{+w}{ }\PY{n}{FileOutputStream}\PY{p}{(}\PY{n}{to}\PY{p}{)}\PY{p}{)}\PY{+w}{ }\PY{p}{\PYZob{}}
\PY{+w}{            }\PY{n}{in}\PY{p}{.}\PY{n+na}{transferTo}\PY{p}{(}\PY{n}{out}\PY{p}{)}\PY{p}{;}
\PY{+w}{        }\PY{p}{\PYZcb{}}
\PY{+w}{    }\PY{p}{\PYZcb{}}
\PY{p}{\PYZcb{}}
\end{Verbatim}
\end{codebox}

\subsection[Desugaring: What the Compiler Actually Generates]{Desugaring: What the Compiler Actually Generates}
\label{h:6:desugaring-what-the-compiler-actually-generates}
The compiler rewrites the try-with-resources block into ordinary \code{try/\allowbreak{}finally} code with null-checks and suppressed-exception wiring. Roughly, \code{copy} above desugars to something like this:

\begin{codebox}
\begin{Verbatim}[commandchars=\\\{\}]
\PY{k+kd}{static}\PY{+w}{ }\PY{k+kt}{void}\PY{+w}{ }\PY{n+nf}{copy}\PY{p}{(}\PY{n}{String}\PY{+w}{ }\PY{n}{from}\PY{p}{,}\PY{+w}{ }\PY{n}{String}\PY{+w}{ }\PY{n}{to}\PY{p}{)}\PY{+w}{ }\PY{k+kd}{throws}\PY{+w}{ }\PY{n}{IOException}\PY{+w}{ }\PY{p}{\PYZob{}}
\PY{+w}{    }\PY{n}{FileInputStream}\PY{+w}{ }\PY{n}{in}\PY{+w}{ }\PY{o}{=}\PY{+w}{ }\PY{k}{new}\PY{+w}{ }\PY{n}{FileInputStream}\PY{p}{(}\PY{n}{from}\PY{p}{)}\PY{p}{;}
\PY{+w}{    }\PY{k}{try}\PY{+w}{ }\PY{p}{\PYZob{}}
\PY{+w}{        }\PY{n}{FileOutputStream}\PY{+w}{ }\PY{n}{out}\PY{+w}{ }\PY{o}{=}\PY{+w}{ }\PY{k}{new}\PY{+w}{ }\PY{n}{FileOutputStream}\PY{p}{(}\PY{n}{to}\PY{p}{)}\PY{p}{;}
\PY{+w}{        }\PY{k}{try}\PY{+w}{ }\PY{p}{\PYZob{}}
\PY{+w}{            }\PY{n}{in}\PY{p}{.}\PY{n+na}{transferTo}\PY{p}{(}\PY{n}{out}\PY{p}{)}\PY{p}{;}
\PY{+w}{        }\PY{p}{\PYZcb{}}\PY{+w}{ }\PY{k}{finally}\PY{+w}{ }\PY{p}{\PYZob{}}
\PY{+w}{            }\PY{k}{if}\PY{+w}{ }\PY{p}{(}\PY{n}{out}\PY{+w}{ }\PY{o}{!}\PY{o}{=}\PY{+w}{ }\PY{k+kc}{null}\PY{p}{)}\PY{+w}{ }\PY{p}{\PYZob{}}
\PY{+w}{                }\PY{n}{out}\PY{p}{.}\PY{n+na}{close}\PY{p}{(}\PY{p}{)}\PY{p}{;}\PY{+w}{ }\PY{c+c1}{// any exception here becomes suppressed if \PYZsq{}try\PYZsq{} already threw}
\PY{+w}{            }\PY{p}{\PYZcb{}}
\PY{+w}{        }\PY{p}{\PYZcb{}}
\PY{+w}{    }\PY{p}{\PYZcb{}}\PY{+w}{ }\PY{k}{finally}\PY{+w}{ }\PY{p}{\PYZob{}}
\PY{+w}{        }\PY{k}{if}\PY{+w}{ }\PY{p}{(}\PY{n}{in}\PY{+w}{ }\PY{o}{!}\PY{o}{=}\PY{+w}{ }\PY{k+kc}{null}\PY{p}{)}\PY{+w}{ }\PY{p}{\PYZob{}}
\PY{+w}{            }\PY{n}{in}\PY{p}{.}\PY{n+na}{close}\PY{p}{(}\PY{p}{)}\PY{p}{;}
\PY{+w}{        }\PY{p}{\PYZcb{}}
\PY{+w}{    }\PY{p}{\PYZcb{}}
\PY{p}{\PYZcb{}}
\end{Verbatim}
\end{codebox}

Key detail: this is why any exception thrown while closing becomes a \textbf{suppressed} exception rather than replacing the original — the generated code catches the primary exception, tries to close, and calls \code{addSuppressed} on failure before rethrowing.

\subsection[Resource Close Order]{Resource Close Order}
\label{h:6:resource-close-order}
Resources declared in a try-with-resources header are closed in the \textbf{reverse} order of declaration — last declared, first closed. This mirrors how you’d unwind nested \code{finally} blocks by hand.

\begin{codebox}
\begin{Verbatim}[commandchars=\\\{\}]
\PY{k+kd}{public}\PY{+w}{ }\PY{k+kd}{class} \PY{n+nc}{CloseOrderDemo}\PY{+w}{ }\PY{p}{\PYZob{}}
\PY{+w}{    }\PY{k+kd}{record} \PY{n+nc}{Loud}\PY{p}{(}\PY{n}{String}\PY{+w}{ }\PY{n}{name}\PY{p}{)}\PY{+w}{ }\PY{k+kd}{implements}\PY{+w}{ }\PY{n}{AutoCloseable}\PY{+w}{ }\PY{p}{\PYZob{}}
\PY{+w}{        }\PY{n}{Loud}\PY{+w}{ }\PY{n+nf}{open}\PY{p}{(}\PY{p}{)}\PY{+w}{ }\PY{p}{\PYZob{}}\PY{+w}{ }\PY{n}{System}\PY{p}{.}\PY{n+na}{out}\PY{p}{.}\PY{n+na}{println}\PY{p}{(}\PY{l+s}{\PYZdq{}}\PY{l+s}{open }\PY{l+s}{\PYZdq{}}\PY{+w}{ }\PY{o}{+}\PY{+w}{ }\PY{n}{name}\PY{p}{)}\PY{p}{;}\PY{+w}{ }\PY{k}{return}\PY{+w}{ }\PY{k}{this}\PY{p}{;}\PY{+w}{ }\PY{p}{\PYZcb{}}
\PY{+w}{        }\PY{n+nd}{@Override}\PY{+w}{ }\PY{k+kd}{public}\PY{+w}{ }\PY{k+kt}{void}\PY{+w}{ }\PY{n+nf}{close}\PY{p}{(}\PY{p}{)}\PY{+w}{ }\PY{p}{\PYZob{}}\PY{+w}{ }\PY{n}{System}\PY{p}{.}\PY{n+na}{out}\PY{p}{.}\PY{n+na}{println}\PY{p}{(}\PY{l+s}{\PYZdq{}}\PY{l+s}{close }\PY{l+s}{\PYZdq{}}\PY{+w}{ }\PY{o}{+}\PY{+w}{ }\PY{n}{name}\PY{p}{)}\PY{p}{;}\PY{+w}{ }\PY{p}{\PYZcb{}}
\PY{+w}{    }\PY{p}{\PYZcb{}}

\PY{+w}{    }\PY{k+kd}{public}\PY{+w}{ }\PY{k+kd}{static}\PY{+w}{ }\PY{k+kt}{void}\PY{+w}{ }\PY{n+nf}{main}\PY{p}{(}\PY{n}{String}\PY{o}{[}\PY{o}{]}\PY{+w}{ }\PY{n}{args}\PY{p}{)}\PY{+w}{ }\PY{p}{\PYZob{}}
\PY{+w}{        }\PY{k}{try}\PY{+w}{ }\PY{p}{(}\PY{n}{Loud}\PY{+w}{ }\PY{n}{a}\PY{+w}{ }\PY{o}{=}\PY{+w}{ }\PY{k}{new}\PY{+w}{ }\PY{n}{Loud}\PY{p}{(}\PY{l+s}{\PYZdq{}}\PY{l+s}{A}\PY{l+s}{\PYZdq{}}\PY{p}{)}\PY{p}{.}\PY{n+na}{open}\PY{p}{(}\PY{p}{)}\PY{p}{;}
\PY{+w}{             }\PY{n}{Loud}\PY{+w}{ }\PY{n}{b}\PY{+w}{ }\PY{o}{=}\PY{+w}{ }\PY{k}{new}\PY{+w}{ }\PY{n}{Loud}\PY{p}{(}\PY{l+s}{\PYZdq{}}\PY{l+s}{B}\PY{l+s}{\PYZdq{}}\PY{p}{)}\PY{p}{.}\PY{n+na}{open}\PY{p}{(}\PY{p}{)}\PY{p}{;}
\PY{+w}{             }\PY{n}{Loud}\PY{+w}{ }\PY{n}{c}\PY{+w}{ }\PY{o}{=}\PY{+w}{ }\PY{k}{new}\PY{+w}{ }\PY{n}{Loud}\PY{p}{(}\PY{l+s}{\PYZdq{}}\PY{l+s}{C}\PY{l+s}{\PYZdq{}}\PY{p}{)}\PY{p}{.}\PY{n+na}{open}\PY{p}{(}\PY{p}{)}\PY{p}{)}\PY{+w}{ }\PY{p}{\PYZob{}}
\PY{+w}{            }\PY{n}{System}\PY{p}{.}\PY{n+na}{out}\PY{p}{.}\PY{n+na}{println}\PY{p}{(}\PY{l+s}{\PYZdq{}}\PY{l+s}{body}\PY{l+s}{\PYZdq{}}\PY{p}{)}\PY{p}{;}
\PY{+w}{        }\PY{p}{\PYZcb{}}
\PY{+w}{    }\PY{p}{\PYZcb{}}
\PY{p}{\PYZcb{}}
\end{Verbatim}
\end{codebox}

Output:

\begin{codebox}
\begin{Verbatim}
open A
open B
open C
body
close C
close B
close A
\end{Verbatim}
\end{codebox}

\subsection[Effectively-Final Resources (Java 9+)]{Effectively-Final Resources (Java 9+)}
\label{h:6:effectively-final-resources-java-9}
Since Java 9, you can use an already-declared effectively-final variable directly in the try-with-resources header, without redeclaring it.

\begin{codebox}
\begin{Verbatim}[commandchars=\\\{\}]
\PY{k+kd}{public}\PY{+w}{ }\PY{k+kd}{class} \PY{n+nc}{ExistingResource}\PY{+w}{ }\PY{p}{\PYZob{}}
\PY{+w}{    }\PY{k+kd}{static}\PY{+w}{ }\PY{k+kt}{void}\PY{+w}{ }\PY{n+nf}{process}\PY{p}{(}\PY{p}{)}\PY{+w}{ }\PY{k+kd}{throws}\PY{+w}{ }\PY{n}{IOException}\PY{+w}{ }\PY{p}{\PYZob{}}
\PY{+w}{        }\PY{n}{BufferedReader}\PY{+w}{ }\PY{n}{reader}\PY{+w}{ }\PY{o}{=}\PY{+w}{ }\PY{k}{new}\PY{+w}{ }\PY{n}{BufferedReader}\PY{p}{(}\PY{k}{new}\PY{+w}{ }\PY{n}{FileReader}\PY{p}{(}\PY{l+s}{\PYZdq{}}\PY{l+s}{data.txt}\PY{l+s}{\PYZdq{}}\PY{p}{)}\PY{p}{)}\PY{p}{;}
\PY{+w}{        }\PY{k}{try}\PY{+w}{ }\PY{p}{(}\PY{n}{reader}\PY{p}{)}\PY{+w}{ }\PY{p}{\PYZob{}}\PY{+w}{ }\PY{c+c1}{// just reference it, no re\PYZhy{}declaration needed}
\PY{+w}{            }\PY{n}{System}\PY{p}{.}\PY{n+na}{out}\PY{p}{.}\PY{n+na}{println}\PY{p}{(}\PY{n}{reader}\PY{p}{.}\PY{n+na}{readLine}\PY{p}{(}\PY{p}{)}\PY{p}{)}\PY{p}{;}
\PY{+w}{        }\PY{p}{\PYZcb{}}
\PY{+w}{    }\PY{p}{\PYZcb{}}
\PY{p}{\PYZcb{}}
\end{Verbatim}
\end{codebox}

\subsection[AutoCloseable vs Closeable]{\code{AutoCloseable} vs \code{Closeable}}
\label{h:6:autocloseable-vs-closeable}
{\tablefont
\begin{xltabular}{\linewidth}{@{}L{0.729}L{0.813}L{1.458}@{}}
\toprule
\rowcolor{tablehdr}
\thd{} & \thd{\code{AutoCloseable}} & \thd{\code{Closeable}} \\
\midrule
\endhead
Introduced & Java 7 & Java 5 (retrofitted to extend \code{AutoCloseable} in Java 7) \\
\code{close()} signature & \code{void~\allowbreak{}close()\allowbreak{}~\allowbreak{}throws~\allowbreak{}Exception} & \code{void~\allowbreak{}close()\allowbreak{}~\allowbreak{}throws~\allowbreak{}IOException} \\
Exception type & Any \code{Exception} & Only \code{IOException} (narrower) \\
Idempotent close required? & Not guaranteed by contract & Yes — calling \code{close()} more than once must be a no-op \\
Typical use & Any resource: locks, DB connections, custom types & I/O-specific types (streams, readers, writers) \\
Can be used in try-with-resources? & Yes & Yes (it extends \code{AutoCloseable}) \\
\bottomrule
\end{xltabular}
}

\begin{codebox}
\begin{Verbatim}[commandchars=\\\{\}]
\PY{k+kd}{public}\PY{+w}{ }\PY{k+kd}{class} \PY{n+nc}{LockResource}\PY{+w}{ }\PY{k+kd}{implements}\PY{+w}{ }\PY{n}{AutoCloseable}\PY{+w}{ }\PY{p}{\PYZob{}}
\PY{+w}{    }\PY{k+kd}{private}\PY{+w}{ }\PY{k+kd}{final}\PY{+w}{ }\PY{n}{java}\PY{p}{.}\PY{n+na}{util}\PY{p}{.}\PY{n+na}{concurrent}\PY{p}{.}\PY{n+na}{locks}\PY{p}{.}\PY{n+na}{Lock}\PY{+w}{ }\PY{n}{lock}\PY{p}{;}

\PY{+w}{    }\PY{k+kd}{public}\PY{+w}{ }\PY{n+nf}{LockResource}\PY{p}{(}\PY{n}{java}\PY{p}{.}\PY{n+na}{util}\PY{p}{.}\PY{n+na}{concurrent}\PY{p}{.}\PY{n+na}{locks}\PY{p}{.}\PY{n+na}{Lock}\PY{+w}{ }\PY{n}{lock}\PY{p}{)}\PY{+w}{ }\PY{p}{\PYZob{}}
\PY{+w}{        }\PY{k}{this}\PY{p}{.}\PY{n+na}{lock}\PY{+w}{ }\PY{o}{=}\PY{+w}{ }\PY{n}{lock}\PY{p}{;}
\PY{+w}{        }\PY{n}{lock}\PY{p}{.}\PY{n+na}{lock}\PY{p}{(}\PY{p}{)}\PY{p}{;}
\PY{+w}{    }\PY{p}{\PYZcb{}}

\PY{+w}{    }\PY{n+nd}{@Override}
\PY{+w}{    }\PY{k+kd}{public}\PY{+w}{ }\PY{k+kt}{void}\PY{+w}{ }\PY{n+nf}{close}\PY{p}{(}\PY{p}{)}\PY{+w}{ }\PY{p}{\PYZob{}}\PY{+w}{ }\PY{c+c1}{// narrows the throws clause to nothing, which is fine}
\PY{+w}{        }\PY{n}{lock}\PY{p}{.}\PY{n+na}{unlock}\PY{p}{(}\PY{p}{)}\PY{p}{;}
\PY{+w}{    }\PY{p}{\PYZcb{}}
\PY{p}{\PYZcb{}}
\end{Verbatim}
\end{codebox}

\begin{codebox}
\begin{Verbatim}[commandchars=\\\{\}]
\PY{k+kd}{public}\PY{+w}{ }\PY{k+kd}{static}\PY{+w}{ }\PY{k+kt}{void}\PY{+w}{ }\PY{n+nf}{withLock}\PY{p}{(}\PY{n}{java}\PY{p}{.}\PY{n+na}{util}\PY{p}{.}\PY{n+na}{concurrent}\PY{p}{.}\PY{n+na}{locks}\PY{p}{.}\PY{n+na}{Lock}\PY{+w}{ }\PY{n}{lock}\PY{p}{,}\PY{+w}{ }\PY{n}{Runnable}\PY{+w}{ }\PY{n}{action}\PY{p}{)}\PY{+w}{ }\PY{p}{\PYZob{}}
\PY{+w}{    }\PY{k}{try}\PY{+w}{ }\PY{p}{(}\PY{n}{LockResource}\PY{+w}{ }\PY{n}{guard}\PY{+w}{ }\PY{o}{=}\PY{+w}{ }\PY{k}{new}\PY{+w}{ }\PY{n}{LockResource}\PY{p}{(}\PY{n}{lock}\PY{p}{)}\PY{p}{)}\PY{+w}{ }\PY{p}{\PYZob{}}
\PY{+w}{        }\PY{n}{action}\PY{p}{.}\PY{n+na}{run}\PY{p}{(}\PY{p}{)}\PY{p}{;}
\PY{+w}{    }\PY{p}{\PYZcb{}}
\PY{p}{\PYZcb{}}
\end{Verbatim}
\end{codebox}

Prefer \code{Closeable} for I/O-flavored classes so callers get the narrower, more specific \code{IOException} instead of a generic \code{Exception}. Use \code{AutoCloseable} for everything else (locks, custom pooled resources, transactions).

\section[Assertions]{Assertions}
\label{h:6:assertions}
An \textbf{assertion} is a statement that checks something you believe must always be true. If it’s false, the JVM throws an \code{AssertionError}. Assertions are for catching \emph{your own bugs} during development and testing — they are not a validation mechanism for public APIs.

\begin{codebox}
\begin{Verbatim}[commandchars=\\\{\}]
\PY{k+kd}{public}\PY{+w}{ }\PY{k+kd}{class} \PY{n+nc}{AssertionDemo}\PY{+w}{ }\PY{p}{\PYZob{}}
\PY{+w}{    }\PY{k+kd}{static}\PY{+w}{ }\PY{k+kt}{int}\PY{+w}{ }\PY{n+nf}{half}\PY{p}{(}\PY{k+kt}{int}\PY{+w}{ }\PY{n}{even}\PY{p}{)}\PY{+w}{ }\PY{p}{\PYZob{}}
\PY{+w}{        }\PY{k}{assert}\PY{+w}{ }\PY{n}{even}\PY{+w}{ }\PY{o}{\PYZpc{}}\PY{+w}{ }\PY{l+m+mi}{2}\PY{+w}{ }\PY{o}{=}\PY{o}{=}\PY{+w}{ }\PY{l+m+mi}{0}\PY{+w}{ }\PY{p}{:}\PY{+w}{ }\PY{l+s}{\PYZdq{}}\PY{l+s}{Expected an even number but got }\PY{l+s}{\PYZdq{}}\PY{+w}{ }\PY{o}{+}\PY{+w}{ }\PY{n}{even}\PY{p}{;}
\PY{+w}{        }\PY{k}{return}\PY{+w}{ }\PY{n}{even}\PY{+w}{ }\PY{o}{/}\PY{+w}{ }\PY{l+m+mi}{2}\PY{p}{;}
\PY{+w}{    }\PY{p}{\PYZcb{}}
\PY{p}{\PYZcb{}}
\end{Verbatim}
\end{codebox}

\subsection[The -ea Flag]{The \code{-\allowbreak{}ea} Flag}
\label{h:6:the--ea-flag}
Assertions are \textbf{disabled by default} at runtime for performance reasons. You must explicitly enable them with \code{-\allowbreak{}ea} (or \code{-\allowbreak{}enableassertions}) when running the JVM:

\begin{codebox}
\begin{Verbatim}[commandchars=\\\{\}]
java\PY{+w}{ }\PYZhy{}ea\PY{+w}{ }com.example.AssertionDemo
\end{Verbatim}
\end{codebox}

Because assertions can be silently switched off, code must never depend on them running. If \code{-\allowbreak{}ea} is off, the \code{assert} line is skipped entirely — no side effect happens at all.

\subsection[Why Assertions Must Not Have Side Effects]{Why Assertions Must Not Have Side Effects}
\label{h:6:why-assertions-must-not-have-side-effects}
If an assertion is disabled, its expression is \textbf{never evaluated}. Any side effect inside it simply won’t happen in production, creating a bug that only shows up when assertions are off — which is the default.

\begin{codebox}
\begin{Verbatim}[commandchars=\\\{\}]
\PY{k+kn}{import}\PY{+w}{ }\PY{n+nn}{java.util.List}\PY{p}{;}

\PY{k+kd}{public}\PY{+w}{ }\PY{k+kd}{class} \PY{n+nc}{SideEffectAssertion}\PY{+w}{ }\PY{p}{\PYZob{}}
\PY{+w}{    }\PY{k+kd}{static}\PY{+w}{ }\PY{k+kt}{void}\PY{+w}{ }\PY{n+nf}{process}\PY{p}{(}\PY{n}{List}\PY{o}{\PYZlt{}}\PY{n}{String}\PY{o}{\PYZgt{}}\PY{+w}{ }\PY{n}{items}\PY{p}{)}\PY{+w}{ }\PY{p}{\PYZob{}}
\PY{+w}{        }\PY{c+c1}{// BROKEN: removeDuplicates() only runs when \PYZhy{}ea is enabled!}
\PY{+w}{        }\PY{k}{assert}\PY{+w}{ }\PY{n}{removeDuplicates}\PY{p}{(}\PY{n}{items}\PY{p}{)}\PY{p}{;}
\PY{+w}{        }\PY{n}{System}\PY{p}{.}\PY{n+na}{out}\PY{p}{.}\PY{n+na}{println}\PY{p}{(}\PY{n}{items}\PY{p}{)}\PY{p}{;}
\PY{+w}{    }\PY{p}{\PYZcb{}}

\PY{+w}{    }\PY{k+kd}{static}\PY{+w}{ }\PY{k+kt}{boolean}\PY{+w}{ }\PY{n+nf}{removeDuplicates}\PY{p}{(}\PY{n}{List}\PY{o}{\PYZlt{}}\PY{n}{String}\PY{o}{\PYZgt{}}\PY{+w}{ }\PY{n}{items}\PY{p}{)}\PY{+w}{ }\PY{p}{\PYZob{}}
\PY{+w}{        }\PY{c+c1}{// pretend this mutates \PYZsq{}items\PYZsq{} and returns true}
\PY{+w}{        }\PY{k}{return}\PY{+w}{ }\PY{k+kc}{true}\PY{p}{;}
\PY{+w}{    }\PY{p}{\PYZcb{}}
\PY{p}{\PYZcb{}}
\end{Verbatim}
\end{codebox}

\begin{codebox}
\begin{Verbatim}[commandchars=\\\{\}]
\PY{k+kn}{import}\PY{+w}{ }\PY{n+nn}{java.util.List}\PY{p}{;}

\PY{k+kd}{public}\PY{+w}{ }\PY{k+kd}{class} \PY{n+nc}{NoSideEffectAssertion}\PY{+w}{ }\PY{p}{\PYZob{}}
\PY{+w}{    }\PY{k+kd}{static}\PY{+w}{ }\PY{k+kt}{void}\PY{+w}{ }\PY{n+nf}{process}\PY{p}{(}\PY{n}{List}\PY{o}{\PYZlt{}}\PY{n}{String}\PY{o}{\PYZgt{}}\PY{+w}{ }\PY{n}{items}\PY{p}{)}\PY{+w}{ }\PY{p}{\PYZob{}}
\PY{+w}{        }\PY{n}{removeDuplicates}\PY{p}{(}\PY{n}{items}\PY{p}{)}\PY{p}{;}\PY{+w}{ }\PY{c+c1}{// always runs, regardless of \PYZhy{}ea}
\PY{+w}{        }\PY{k}{assert}\PY{+w}{ }\PY{o}{!}\PY{n}{hasDuplicates}\PY{p}{(}\PY{n}{items}\PY{p}{)}\PY{+w}{ }\PY{p}{:}\PY{+w}{ }\PY{l+s}{\PYZdq{}}\PY{l+s}{Duplicates remain after cleanup}\PY{l+s}{\PYZdq{}}\PY{p}{;}\PY{+w}{ }\PY{c+c1}{// pure check, no side effect}
\PY{+w}{        }\PY{n}{System}\PY{p}{.}\PY{n+na}{out}\PY{p}{.}\PY{n+na}{println}\PY{p}{(}\PY{n}{items}\PY{p}{)}\PY{p}{;}
\PY{+w}{    }\PY{p}{\PYZcb{}}

\PY{+w}{    }\PY{k+kd}{static}\PY{+w}{ }\PY{k+kt}{void}\PY{+w}{ }\PY{n+nf}{removeDuplicates}\PY{p}{(}\PY{n}{List}\PY{o}{\PYZlt{}}\PY{n}{String}\PY{o}{\PYZgt{}}\PY{+w}{ }\PY{n}{items}\PY{p}{)}\PY{+w}{ }\PY{p}{\PYZob{}}\PY{+w}{ }\PY{c+cm}{/* mutates items */}\PY{+w}{ }\PY{p}{\PYZcb{}}
\PY{+w}{    }\PY{k+kd}{static}\PY{+w}{ }\PY{k+kt}{boolean}\PY{+w}{ }\PY{n+nf}{hasDuplicates}\PY{p}{(}\PY{n}{List}\PY{o}{\PYZlt{}}\PY{n}{String}\PY{o}{\PYZgt{}}\PY{+w}{ }\PY{n}{items}\PY{p}{)}\PY{+w}{ }\PY{p}{\PYZob{}}\PY{+w}{ }\PY{k}{return}\PY{+w}{ }\PY{k+kc}{false}\PY{p}{;}\PY{+w}{ }\PY{p}{\PYZcb{}}
\PY{p}{\PYZcb{}}
\end{Verbatim}
\end{codebox}

\subsection[Why Assertions Must Not Validate Public API Arguments]{Why Assertions Must Not Validate Public API Arguments}
\label{h:6:why-assertions-must-not-validate-public-api-arguments}
Because assertions can be disabled, they are the wrong tool for checking arguments on a \textbf{public} method — anyone calling your API from outside might run with assertions off, and your validation would silently vanish, letting bad input through.

\begin{codebox}
\begin{Verbatim}[commandchars=\\\{\}]
\PY{c+c1}{// BROKEN: public API validation via assert — disappears when \PYZhy{}ea is off}
\PY{k+kd}{public}\PY{+w}{ }\PY{k+kd}{class} \PY{n+nc}{AccountBroken}\PY{+w}{ }\PY{p}{\PYZob{}}
\PY{+w}{    }\PY{k+kd}{private}\PY{+w}{ }\PY{k+kt}{double}\PY{+w}{ }\PY{n}{balance}\PY{p}{;}

\PY{+w}{    }\PY{k+kd}{public}\PY{+w}{ }\PY{k+kt}{void}\PY{+w}{ }\PY{n+nf}{withdraw}\PY{p}{(}\PY{k+kt}{double}\PY{+w}{ }\PY{n}{amount}\PY{p}{)}\PY{+w}{ }\PY{p}{\PYZob{}}
\PY{+w}{        }\PY{k}{assert}\PY{+w}{ }\PY{n}{amount}\PY{+w}{ }\PY{o}{\PYZgt{}}\PY{+w}{ }\PY{l+m+mi}{0}\PY{+w}{ }\PY{p}{:}\PY{+w}{ }\PY{l+s}{\PYZdq{}}\PY{l+s}{amount must be positive}\PY{l+s}{\PYZdq{}}\PY{p}{;}\PY{+w}{ }\PY{c+c1}{// NOT enforced by default!}
\PY{+w}{        }\PY{n}{balance}\PY{+w}{ }\PY{o}{\PYZhy{}}\PY{o}{=}\PY{+w}{ }\PY{n}{amount}\PY{p}{;}
\PY{+w}{    }\PY{p}{\PYZcb{}}
\PY{p}{\PYZcb{}}
\end{Verbatim}
\end{codebox}

\begin{codebox}
\begin{Verbatim}[commandchars=\\\{\}]
\PY{c+c1}{// FIXED: public API validation should throw a real, unconditional exception}
\PY{k+kd}{public}\PY{+w}{ }\PY{k+kd}{class} \PY{n+nc}{AccountFixed}\PY{+w}{ }\PY{p}{\PYZob{}}
\PY{+w}{    }\PY{k+kd}{private}\PY{+w}{ }\PY{k+kt}{double}\PY{+w}{ }\PY{n}{balance}\PY{p}{;}

\PY{+w}{    }\PY{k+kd}{public}\PY{+w}{ }\PY{k+kt}{void}\PY{+w}{ }\PY{n+nf}{withdraw}\PY{p}{(}\PY{k+kt}{double}\PY{+w}{ }\PY{n}{amount}\PY{p}{)}\PY{+w}{ }\PY{p}{\PYZob{}}
\PY{+w}{        }\PY{k}{if}\PY{+w}{ }\PY{p}{(}\PY{n}{amount}\PY{+w}{ }\PY{o}{\PYZlt{}}\PY{o}{=}\PY{+w}{ }\PY{l+m+mi}{0}\PY{p}{)}\PY{+w}{ }\PY{p}{\PYZob{}}
\PY{+w}{            }\PY{k}{throw}\PY{+w}{ }\PY{k}{new}\PY{+w}{ }\PY{n}{IllegalArgumentException}\PY{p}{(}\PY{l+s}{\PYZdq{}}\PY{l+s}{amount must be positive: }\PY{l+s}{\PYZdq{}}\PY{+w}{ }\PY{o}{+}\PY{+w}{ }\PY{n}{amount}\PY{p}{)}\PY{p}{;}
\PY{+w}{        }\PY{p}{\PYZcb{}}
\PY{+w}{        }\PY{n}{balance}\PY{+w}{ }\PY{o}{\PYZhy{}}\PY{o}{=}\PY{+w}{ }\PY{n}{amount}\PY{p}{;}
\PY{+w}{    }\PY{p}{\PYZcb{}}
\PY{p}{\PYZcb{}}
\end{Verbatim}
\end{codebox}

\subsection[Where Assertions Belong]{Where Assertions Belong}
\label{h:6:where-assertions-belong}
{\tablefont
\begin{xltabular}{\linewidth}{@{}L{1.108}L{0.892}@{}}
\toprule
\rowcolor{tablehdr}
\thd{Use case} & \thd{Right tool} \\
\midrule
\endhead
Checking a public method’s arguments & \code{if~\allowbreak{}(...)\allowbreak{}~\allowbreak{}throw~\allowbreak{}new~\allowbreak{}Illegal\allowbreak{}Argument\allowbreak{}Exception(...)} \\
Checking internal invariants only your own code can violate (e.g., “this private helper should never see a negative index here”) & \code{assert} \\
Documenting an assumption for future maintainers, checked only during testing & \code{assert} \\
Guarding against \code{null} from an external caller & Explicit \code{null} check + exception (e.g., \code{Objects.\allowbreak{}requireNonNull}) \\
Verifying test expectations & Use a test framework’s assertions (e.g., JUnit’s \code{assertEquals}), not the \code{assert} keyword \\
\bottomrule
\end{xltabular}
}

\begin{codebox}
\begin{Verbatim}[commandchars=\\\{\}]
\PY{k+kd}{public}\PY{+w}{ }\PY{k+kd}{class} \PY{n+nc}{InvariantExample}\PY{+w}{ }\PY{p}{\PYZob{}}
\PY{+w}{    }\PY{k+kd}{private}\PY{+w}{ }\PY{k+kt}{int}\PY{+w}{ }\PY{n+nf}{index}\PY{p}{(}\PY{k+kt}{int}\PY{+w}{ }\PY{n}{size}\PY{p}{,}\PY{+w}{ }\PY{k+kt}{int}\PY{+w}{ }\PY{n}{position}\PY{p}{)}\PY{+w}{ }\PY{p}{\PYZob{}}
\PY{+w}{        }\PY{k+kt}{int}\PY{+w}{ }\PY{n}{result}\PY{+w}{ }\PY{o}{=}\PY{+w}{ }\PY{n}{position}\PY{+w}{ }\PY{o}{\PYZpc{}}\PY{+w}{ }\PY{n}{size}\PY{p}{;}
\PY{+w}{        }\PY{c+c1}{// Internal invariant: our own arithmetic guarantees this is always \PYZgt{}= 0.}
\PY{+w}{        }\PY{c+c1}{// Safe to use assert here because it\PYZsq{}s checking OUR logic, not caller input.}
\PY{+w}{        }\PY{k}{assert}\PY{+w}{ }\PY{n}{result}\PY{+w}{ }\PY{o}{\PYZgt{}}\PY{o}{=}\PY{+w}{ }\PY{l+m+mi}{0}\PY{+w}{ }\PY{p}{:}\PY{+w}{ }\PY{l+s}{\PYZdq{}}\PY{l+s}{computed negative index: }\PY{l+s}{\PYZdq{}}\PY{+w}{ }\PY{o}{+}\PY{+w}{ }\PY{n}{result}\PY{p}{;}
\PY{+w}{        }\PY{k}{return}\PY{+w}{ }\PY{n}{result}\PY{p}{;}
\PY{+w}{    }\PY{p}{\PYZcb{}}
\PY{p}{\PYZcb{}}
\end{Verbatim}
\end{codebox}

\section[Common Code-Review Interview Pitfalls]{Common Code-Review Interview Pitfalls}
\label{h:6:common-code-review-interview-pitfalls}
\begin{enumerate}
\item 
\textbf{Catching \code{Exception} (or \code{Throwable}) and doing nothing (\code{catch~\allowbreak{}(\allowbreak{}Exception~\allowbreak{}e)\allowbreak{}~\allowbreak{}\{\}}).}
Why it matters: it hides every failure, including bugs, making the system silently wrong instead of loudly failing. It’s one of the fastest ways to fail a code review.

\begin{codebox}
\begin{Verbatim}[commandchars=\\\{\}]
\PY{c+c1}{// Before}
\PY{k}{try}\PY{+w}{ }\PY{p}{\PYZob{}}\PY{+w}{ }\PY{n}{save}\PY{p}{(}\PY{n}{record}\PY{p}{)}\PY{p}{;}\PY{+w}{ }\PY{p}{\PYZcb{}}\PY{+w}{ }\PY{k}{catch}\PY{+w}{ }\PY{p}{(}\PY{n}{Exception}\PY{+w}{ }\PY{n}{e}\PY{p}{)}\PY{+w}{ }\PY{p}{\PYZob{}}\PY{p}{\PYZcb{}}
\PY{c+c1}{// After}
\PY{k}{try}\PY{+w}{ }\PY{p}{\PYZob{}}\PY{+w}{ }\PY{n}{save}\PY{p}{(}\PY{n}{record}\PY{p}{)}\PY{p}{;}\PY{+w}{ }\PY{p}{\PYZcb{}}\PY{+w}{ }\PY{k}{catch}\PY{+w}{ }\PY{p}{(}\PY{n}{IOException}\PY{+w}{ }\PY{n}{e}\PY{p}{)}\PY{+w}{ }\PY{p}{\PYZob{}}\PY{+w}{ }\PY{k}{throw}\PY{+w}{ }\PY{k}{new}\PY{+w}{ }\PY{n}{UncheckedIOException}\PY{p}{(}\PY{l+s}{\PYZdq{}}\PY{l+s}{Failed to save record}\PY{l+s}{\PYZdq{}}\PY{p}{,}\PY{+w}{ }\PY{n}{e}\PY{p}{)}\PY{p}{;}\PY{+w}{ }\PY{p}{\PYZcb{}}
\end{Verbatim}
\end{codebox}

\item 
\textbf{\code{return} inside a \code{finally} block, silently discarding an exception.}
Why it matters: the caller thinks the operation succeeded and never learns about the real failure.

\begin{codebox}
\begin{Verbatim}[commandchars=\\\{\}]
\PY{c+c1}{// Before}
\PY{k}{finally}\PY{+w}{ }\PY{p}{\PYZob{}}\PY{+w}{ }\PY{k}{return}\PY{+w}{ }\PY{n}{result}\PY{p}{;}\PY{+w}{ }\PY{p}{\PYZcb{}}
\PY{c+c1}{// After}
\PY{k}{finally}\PY{+w}{ }\PY{p}{\PYZob{}}\PY{+w}{ }\PY{n}{cleanup}\PY{p}{(}\PY{p}{)}\PY{p}{;}\PY{+w}{ }\PY{p}{\PYZcb{}}\PY{+w}{ }\PY{c+c1}{// no return/throw in finally}
\end{Verbatim}
\end{codebox}

\item 
\textbf{Losing the cause when wrapping exceptions.}
Why it matters: without the cause, the stack trace stops at the wrapping point and the real root cause is unrecoverable in logs.

\begin{codebox}
\begin{Verbatim}[commandchars=\\\{\}]
\PY{c+c1}{// Before}
\PY{k}{throw}\PY{+w}{ }\PY{k}{new}\PY{+w}{ }\PY{n}{RuntimeException}\PY{p}{(}\PY{l+s}{\PYZdq{}}\PY{l+s}{failed}\PY{l+s}{\PYZdq{}}\PY{p}{)}\PY{p}{;}
\PY{c+c1}{// After}
\PY{k}{throw}\PY{+w}{ }\PY{k}{new}\PY{+w}{ }\PY{n}{RuntimeException}\PY{p}{(}\PY{l+s}{\PYZdq{}}\PY{l+s}{failed}\PY{l+s}{\PYZdq{}}\PY{p}{,}\PY{+w}{ }\PY{n}{e}\PY{p}{)}\PY{p}{;}
\end{Verbatim}
\end{codebox}

\item 
\textbf{Logging an exception and then rethrowing it (double reporting).}
Why it matters: the same error gets logged multiple times up the call stack, flooding logs and confusing on-call engineers about how many times it actually happened. Log at the boundary where you handle it, or rethrow — never both at every layer.

\begin{codebox}
\begin{Verbatim}[commandchars=\\\{\}]
\PY{c+c1}{// Before}
\PY{k}{catch}\PY{+w}{ }\PY{p}{(}\PY{n}{IOException}\PY{+w}{ }\PY{n}{e}\PY{p}{)}\PY{+w}{ }\PY{p}{\PYZob{}}\PY{+w}{ }\PY{n}{log}\PY{p}{.}\PY{n+na}{error}\PY{p}{(}\PY{l+s}{\PYZdq{}}\PY{l+s}{failed}\PY{l+s}{\PYZdq{}}\PY{p}{,}\PY{+w}{ }\PY{n}{e}\PY{p}{)}\PY{p}{;}\PY{+w}{ }\PY{k}{throw}\PY{+w}{ }\PY{n}{e}\PY{p}{;}\PY{+w}{ }\PY{p}{\PYZcb{}}
\PY{c+c1}{// After}
\PY{k}{catch}\PY{+w}{ }\PY{p}{(}\PY{n}{IOException}\PY{+w}{ }\PY{n}{e}\PY{p}{)}\PY{+w}{ }\PY{p}{\PYZob{}}\PY{+w}{ }\PY{k}{throw}\PY{+w}{ }\PY{k}{new}\PY{+w}{ }\PY{n}{ServiceException}\PY{p}{(}\PY{l+s}{\PYZdq{}}\PY{l+s}{Upload failed}\PY{l+s}{\PYZdq{}}\PY{p}{,}\PY{+w}{ }\PY{n}{e}\PY{p}{)}\PY{p}{;}\PY{+w}{ }\PY{p}{\PYZcb{}}\PY{+w}{ }\PY{c+c1}{// log once, at the top\PYZhy{}level handler}
\end{Verbatim}
\end{codebox}

\item 
\textbf{Catching \code{InterruptedException} and swallowing it without restoring the interrupt flag.}
Why it matters: \code{Thread.\allowbreak{}sleep}/\code{wait}/blocking calls clear the interrupt flag when they throw; if you don’t restore it, the thread’s owning framework (e.g., an executor) can no longer tell the thread was asked to stop, breaking graceful shutdown.

\begin{codebox}
\begin{Verbatim}[commandchars=\\\{\}]
\PY{c+c1}{// Before}
\PY{k}{try}\PY{+w}{ }\PY{p}{\PYZob{}}\PY{+w}{ }\PY{n}{Thread}\PY{p}{.}\PY{n+na}{sleep}\PY{p}{(}\PY{l+m+mi}{1000}\PY{p}{)}\PY{p}{;}\PY{+w}{ }\PY{p}{\PYZcb{}}\PY{+w}{ }\PY{k}{catch}\PY{+w}{ }\PY{p}{(}\PY{n}{InterruptedException}\PY{+w}{ }\PY{n}{e}\PY{p}{)}\PY{+w}{ }\PY{p}{\PYZob{}}\PY{+w}{ }\PY{c+cm}{/* ignored */}\PY{+w}{ }\PY{p}{\PYZcb{}}
\PY{c+c1}{// After}
\PY{k}{try}\PY{+w}{ }\PY{p}{\PYZob{}}\PY{+w}{ }\PY{n}{Thread}\PY{p}{.}\PY{n+na}{sleep}\PY{p}{(}\PY{l+m+mi}{1000}\PY{p}{)}\PY{p}{;}\PY{+w}{ }\PY{p}{\PYZcb{}}\PY{+w}{ }\PY{k}{catch}\PY{+w}{ }\PY{p}{(}\PY{n}{InterruptedException}\PY{+w}{ }\PY{n}{e}\PY{p}{)}\PY{+w}{ }\PY{p}{\PYZob{}}\PY{+w}{ }\PY{n}{Thread}\PY{p}{.}\PY{n+na}{currentThread}\PY{p}{(}\PY{p}{)}\PY{p}{.}\PY{n+na}{interrupt}\PY{p}{(}\PY{p}{)}\PY{p}{;}\PY{+w}{ }\PY{p}{\PYZcb{}}
\end{Verbatim}
\end{codebox}

\item 
\textbf{Using exceptions for normal control flow.}
Why it matters: exceptions capture a full stack trace and are expensive; using them to signal ordinary conditions (like “not found”) hurts performance and readability.

\begin{codebox}
\begin{Verbatim}[commandchars=\\\{\}]
\PY{c+c1}{// Before}
\PY{k}{try}\PY{+w}{ }\PY{p}{\PYZob{}}\PY{+w}{ }\PY{k}{return}\PY{+w}{ }\PY{n}{map}\PY{p}{.}\PY{n+na}{get}\PY{p}{(}\PY{n}{key}\PY{p}{)}\PY{p}{;}\PY{+w}{ }\PY{p}{\PYZcb{}}\PY{+w}{ }\PY{k}{catch}\PY{+w}{ }\PY{p}{(}\PY{n}{NullPointerException}\PY{+w}{ }\PY{n}{e}\PY{p}{)}\PY{+w}{ }\PY{p}{\PYZob{}}\PY{+w}{ }\PY{k}{return}\PY{+w}{ }\PY{n}{defaultValue}\PY{p}{;}\PY{+w}{ }\PY{p}{\PYZcb{}}
\PY{c+c1}{// After}
\PY{k}{return}\PY{+w}{ }\PY{n}{map}\PY{p}{.}\PY{n+na}{getOrDefault}\PY{p}{(}\PY{n}{key}\PY{p}{,}\PY{+w}{ }\PY{n}{defaultValue}\PY{p}{)}\PY{p}{;}
\end{Verbatim}
\end{codebox}

\item 
\textbf{Throwing generic exceptions (\code{throw~\allowbreak{}new~\allowbreak{}Exception(\allowbreak{}\textquotedbl{}error\textquotedbl{})} or \code{RuntimeException}) instead of a specific, meaningful type.}
Why it matters: callers can’t distinguish failure modes to handle them differently, and generic catches become tempting, widening the blast radius of \code{catch} blocks.

\begin{codebox}
\begin{Verbatim}[commandchars=\\\{\}]
\PY{c+c1}{// Before}
\PY{k}{throw}\PY{+w}{ }\PY{k}{new}\PY{+w}{ }\PY{n}{RuntimeException}\PY{p}{(}\PY{l+s}{\PYZdq{}}\PY{l+s}{user not found}\PY{l+s}{\PYZdq{}}\PY{p}{)}\PY{p}{;}
\PY{c+c1}{// After}
\PY{k}{throw}\PY{+w}{ }\PY{k}{new}\PY{+w}{ }\PY{n}{NoSuchElementException}\PY{p}{(}\PY{l+s}{\PYZdq{}}\PY{l+s}{No user with id }\PY{l+s}{\PYZdq{}}\PY{+w}{ }\PY{o}{+}\PY{+w}{ }\PY{n}{id}\PY{p}{)}\PY{p}{;}
\end{Verbatim}
\end{codebox}

\item 
\textbf{Validating public API arguments with \code{assert} instead of a real check.}
Why it matters: assertions are off by default in production (\code{-\allowbreak{}ea} not set), so the validation silently disappears, letting invalid state propagate deep into the system before it fails somewhere confusing.

\begin{codebox}
\begin{Verbatim}[commandchars=\\\{\}]
\PY{c+c1}{// Before}
\PY{k}{assert}\PY{+w}{ }\PY{n}{amount}\PY{+w}{ }\PY{o}{\PYZgt{}}\PY{+w}{ }\PY{l+m+mi}{0}\PY{p}{;}
\PY{c+c1}{// After}
\PY{k}{if}\PY{+w}{ }\PY{p}{(}\PY{n}{amount}\PY{+w}{ }\PY{o}{\PYZlt{}}\PY{o}{=}\PY{+w}{ }\PY{l+m+mi}{0}\PY{p}{)}\PY{+w}{ }\PY{k}{throw}\PY{+w}{ }\PY{k}{new}\PY{+w}{ }\PY{n}{IllegalArgumentException}\PY{p}{(}\PY{l+s}{\PYZdq{}}\PY{l+s}{amount must be positive: }\PY{l+s}{\PYZdq{}}\PY{+w}{ }\PY{o}{+}\PY{+w}{ }\PY{n}{amount}\PY{p}{)}\PY{p}{;}
\end{Verbatim}
\end{codebox}

\item 
\textbf{Putting side effects inside an \code{assert} expression.}
Why it matters: when assertions are disabled, the expression is never evaluated, so that side effect silently stops happening — a correctness bug that appears only in production.

\begin{codebox}
\begin{Verbatim}[commandchars=\\\{\}]
\PY{c+c1}{// Before}
\PY{k}{assert}\PY{+w}{ }\PY{n}{list}\PY{p}{.}\PY{n+na}{remove}\PY{p}{(}\PY{n}{item}\PY{p}{)}\PY{p}{;}
\PY{c+c1}{// After}
\PY{k+kt}{boolean}\PY{+w}{ }\PY{n}{removed}\PY{+w}{ }\PY{o}{=}\PY{+w}{ }\PY{n}{list}\PY{p}{.}\PY{n+na}{remove}\PY{p}{(}\PY{n}{item}\PY{p}{)}\PY{p}{;}
\PY{k}{assert}\PY{+w}{ }\PY{n}{removed}\PY{+w}{ }\PY{p}{:}\PY{+w}{ }\PY{l+s}{\PYZdq{}}\PY{l+s}{expected item to be present}\PY{l+s}{\PYZdq{}}\PY{p}{;}
\end{Verbatim}
\end{codebox}

\item 
\textbf{Not closing resources, or closing them manually instead of using try-with-resources.}
Why it matters: manual \code{finally}-based closing is verbose and easy to get subtly wrong (e.g., an exception from \code{close()} masking the real exception, or a missing null-check causing a \code{NullPointerException} during cleanup).

\begin{codebox}
\begin{Verbatim}[commandchars=\\\{\}]
\PY{c+c1}{// Before}
\PY{n}{FileInputStream}\PY{+w}{ }\PY{n}{in}\PY{+w}{ }\PY{o}{=}\PY{+w}{ }\PY{k}{new}\PY{+w}{ }\PY{n}{FileInputStream}\PY{p}{(}\PY{n}{path}\PY{p}{)}\PY{p}{;}
\PY{k}{try}\PY{+w}{ }\PY{p}{\PYZob{}}\PY{+w}{ }\PY{n}{use}\PY{p}{(}\PY{n}{in}\PY{p}{)}\PY{p}{;}\PY{+w}{ }\PY{p}{\PYZcb{}}\PY{+w}{ }\PY{k}{finally}\PY{+w}{ }\PY{p}{\PYZob{}}\PY{+w}{ }\PY{n}{in}\PY{p}{.}\PY{n+na}{close}\PY{p}{(}\PY{p}{)}\PY{p}{;}\PY{+w}{ }\PY{p}{\PYZcb{}}
\PY{c+c1}{// After}
\PY{k}{try}\PY{+w}{ }\PY{p}{(}\PY{n}{FileInputStream}\PY{+w}{ }\PY{n}{in}\PY{+w}{ }\PY{o}{=}\PY{+w}{ }\PY{k}{new}\PY{+w}{ }\PY{n}{FileInputStream}\PY{p}{(}\PY{n}{path}\PY{p}{)}\PY{p}{)}\PY{+w}{ }\PY{p}{\PYZob{}}\PY{+w}{ }\PY{n}{use}\PY{p}{(}\PY{n}{in}\PY{p}{)}\PY{p}{;}\PY{+w}{ }\PY{p}{\PYZcb{}}
\end{Verbatim}
\end{codebox}

\item 
\textbf{Throwing checked exceptions from lambdas passed to standard functional interfaces, causing awkward wrapping or compile errors.}
Why it matters: \code{Function}, \code{Consumer}, etc. don’t declare \code{throws}, so checked exceptions inside lambdas must be caught locally or wrapped — forgetting this leads to compile errors or ugly ad hoc wrapping scattered across the codebase.

\begin{codebox}
\begin{Verbatim}[commandchars=\\\{\}]
\PY{c+c1}{// Before (does not compile: readAllBytes throws IOException)}
\PY{n}{paths}\PY{p}{.}\PY{n+na}{stream}\PY{p}{(}\PY{p}{)}\PY{p}{.}\PY{n+na}{map}\PY{p}{(}\PY{n}{p}\PY{+w}{ }\PY{o}{\PYZhy{}}\PY{o}{\PYZgt{}}\PY{+w}{ }\PY{n}{Files}\PY{p}{.}\PY{n+na}{readAllBytes}\PY{p}{(}\PY{n}{p}\PY{p}{)}\PY{p}{)}\PY{p}{;}
\PY{c+c1}{// After}
\PY{n}{paths}\PY{p}{.}\PY{n+na}{stream}\PY{p}{(}\PY{p}{)}\PY{p}{.}\PY{n+na}{map}\PY{p}{(}\PY{n}{p}\PY{+w}{ }\PY{o}{\PYZhy{}}\PY{o}{\PYZgt{}}\PY{+w}{ }\PY{p}{\PYZob{}}
\PY{+w}{    }\PY{k}{try}\PY{+w}{ }\PY{p}{\PYZob{}}\PY{+w}{ }\PY{k}{return}\PY{+w}{ }\PY{n}{Files}\PY{p}{.}\PY{n+na}{readAllBytes}\PY{p}{(}\PY{n}{p}\PY{p}{)}\PY{p}{;}\PY{+w}{ }\PY{p}{\PYZcb{}}
\PY{+w}{    }\PY{k}{catch}\PY{+w}{ }\PY{p}{(}\PY{n}{IOException}\PY{+w}{ }\PY{n}{e}\PY{p}{)}\PY{+w}{ }\PY{p}{\PYZob{}}\PY{+w}{ }\PY{k}{throw}\PY{+w}{ }\PY{k}{new}\PY{+w}{ }\PY{n}{UncheckedIOException}\PY{p}{(}\PY{n}{e}\PY{p}{)}\PY{p}{;}\PY{+w}{ }\PY{p}{\PYZcb{}}
\PY{p}{\PYZcb{}}\PY{p}{)}\PY{p}{;}
\end{Verbatim}
\end{codebox}

\item 
\textbf{Catching a broad type just to rethrow a narrower, unrelated one, losing the original stack trace.}
Why it matters: creating a brand-new exception without chaining the original throws away the actual failure location, making the bug much harder to diagnose from logs alone.

\begin{codebox}
\begin{Verbatim}[commandchars=\\\{\}]
\PY{c+c1}{// Before}
\PY{k}{catch}\PY{+w}{ }\PY{p}{(}\PY{n}{Exception}\PY{+w}{ }\PY{n}{e}\PY{p}{)}\PY{+w}{ }\PY{p}{\PYZob{}}\PY{+w}{ }\PY{k}{throw}\PY{+w}{ }\PY{k}{new}\PY{+w}{ }\PY{n}{BusinessException}\PY{p}{(}\PY{l+s}{\PYZdq{}}\PY{l+s}{failed}\PY{l+s}{\PYZdq{}}\PY{p}{)}\PY{p}{;}\PY{+w}{ }\PY{p}{\PYZcb{}}
\PY{c+c1}{// After}
\PY{k}{catch}\PY{+w}{ }\PY{p}{(}\PY{n}{Exception}\PY{+w}{ }\PY{n}{e}\PY{p}{)}\PY{+w}{ }\PY{p}{\PYZob{}}\PY{+w}{ }\PY{k}{throw}\PY{+w}{ }\PY{k}{new}\PY{+w}{ }\PY{n}{BusinessException}\PY{p}{(}\PY{l+s}{\PYZdq{}}\PY{l+s}{failed}\PY{l+s}{\PYZdq{}}\PY{p}{,}\PY{+w}{ }\PY{n}{e}\PY{p}{)}\PY{p}{;}\PY{+w}{ }\PY{p}{\PYZcb{}}
\end{Verbatim}
\end{codebox}

\item 
\textbf{Declaring \code{throws~\allowbreak{}Exception} on a method signature instead of the specific checked exceptions it actually throws.}
Why it matters: it forces every caller to catch the overly broad \code{Exception}, defeating the whole purpose of checked exceptions — communicating precisely what can go wrong.

\begin{codebox}
\begin{Verbatim}[commandchars=\\\{\}]
\PY{c+c1}{// Before}
\PY{k+kt}{void}\PY{+w}{ }\PY{n+nf}{process}\PY{p}{(}\PY{p}{)}\PY{+w}{ }\PY{k+kd}{throws}\PY{+w}{ }\PY{n}{Exception}\PY{+w}{ }\PY{p}{\PYZob{}}\PY{+w}{ }\PY{p}{.}\PY{p}{.}\PY{p}{.}\PY{+w}{ }\PY{p}{\PYZcb{}}
\PY{c+c1}{// After}
\PY{k+kt}{void}\PY{+w}{ }\PY{n+nf}{process}\PY{p}{(}\PY{p}{)}\PY{+w}{ }\PY{k+kd}{throws}\PY{+w}{ }\PY{n}{IOException}\PY{p}{,}\PY{+w}{ }\PY{n}{SQLException}\PY{+w}{ }\PY{p}{\PYZob{}}\PY{+w}{ }\PY{p}{.}\PY{p}{.}\PY{p}{.}\PY{+w}{ }\PY{p}{\PYZcb{}}
\end{Verbatim}
\end{codebox}

\item 
\textbf{Comparing exception messages (\code{e.\allowbreak{}get\allowbreak{}Message().\allowbreak{}contains(...)}) to detect a specific failure instead of catching a specific type.}
Why it matters: messages are not part of the API contract and can change between JDK versions or library updates, silently breaking the check.

\begin{codebox}
\begin{Verbatim}[commandchars=\\\{\}]
\PY{c+c1}{// Before}
\PY{k}{catch}\PY{+w}{ }\PY{p}{(}\PY{n}{Exception}\PY{+w}{ }\PY{n}{e}\PY{p}{)}\PY{+w}{ }\PY{p}{\PYZob{}}\PY{+w}{ }\PY{k}{if}\PY{+w}{ }\PY{p}{(}\PY{n}{e}\PY{p}{.}\PY{n+na}{getMessage}\PY{p}{(}\PY{p}{)}\PY{p}{.}\PY{n+na}{contains}\PY{p}{(}\PY{l+s}{\PYZdq{}}\PY{l+s}{not found}\PY{l+s}{\PYZdq{}}\PY{p}{)}\PY{p}{)}\PY{+w}{ }\PY{p}{\PYZob{}}\PY{+w}{ }\PY{p}{.}\PY{p}{.}\PY{p}{.}\PY{+w}{ }\PY{p}{\PYZcb{}}\PY{+w}{ }\PY{p}{\PYZcb{}}
\PY{c+c1}{// After}
\PY{k}{catch}\PY{+w}{ }\PY{p}{(}\PY{n}{NoSuchElementException}\PY{+w}{ }\PY{n}{e}\PY{p}{)}\PY{+w}{ }\PY{p}{\PYZob{}}\PY{+w}{ }\PY{p}{.}\PY{p}{.}\PY{p}{.}\PY{+w}{ }\PY{p}{\PYZcb{}}
\end{Verbatim}
\end{codebox}

\item 
\textbf{Ignoring suppressed exceptions when debugging try-with-resources failures.}
Why it matters: if closing a resource also fails, that information is attached as a suppressed exception on the primary one; skipping \code{getSuppressed()} during triage can hide a second, independent root cause (e.g., a leaked connection on top of the original error).

\begin{codebox}
\begin{Verbatim}[commandchars=\\\{\}]
\PY{c+c1}{// Before}
\PY{k}{catch}\PY{+w}{ }\PY{p}{(}\PY{n}{Exception}\PY{+w}{ }\PY{n}{e}\PY{p}{)}\PY{+w}{ }\PY{p}{\PYZob{}}\PY{+w}{ }\PY{n}{log}\PY{p}{.}\PY{n+na}{error}\PY{p}{(}\PY{l+s}{\PYZdq{}}\PY{l+s}{Failed}\PY{l+s}{\PYZdq{}}\PY{p}{,}\PY{+w}{ }\PY{n}{e}\PY{p}{)}\PY{p}{;}\PY{+w}{ }\PY{p}{\PYZcb{}}\PY{+w}{ }\PY{c+c1}{// suppressed causes buried in nested trace}
\PY{c+c1}{// After}
\PY{k}{catch}\PY{+w}{ }\PY{p}{(}\PY{n}{Exception}\PY{+w}{ }\PY{n}{e}\PY{p}{)}\PY{+w}{ }\PY{p}{\PYZob{}}
\PY{+w}{    }\PY{n}{log}\PY{p}{.}\PY{n+na}{error}\PY{p}{(}\PY{l+s}{\PYZdq{}}\PY{l+s}{Failed}\PY{l+s}{\PYZdq{}}\PY{p}{,}\PY{+w}{ }\PY{n}{e}\PY{p}{)}\PY{p}{;}
\PY{+w}{    }\PY{k}{for}\PY{+w}{ }\PY{p}{(}\PY{n}{Throwable}\PY{+w}{ }\PY{n}{s}\PY{+w}{ }\PY{p}{:}\PY{+w}{ }\PY{n}{e}\PY{p}{.}\PY{n+na}{getSuppressed}\PY{p}{(}\PY{p}{)}\PY{p}{)}\PY{+w}{ }\PY{n}{log}\PY{p}{.}\PY{n+na}{error}\PY{p}{(}\PY{l+s}{\PYZdq{}}\PY{l+s}{Also failed while closing}\PY{l+s}{\PYZdq{}}\PY{p}{,}\PY{+w}{ }\PY{n}{s}\PY{p}{)}\PY{p}{;}
\PY{p}{\PYZcb{}}
\end{Verbatim}
\end{codebox}

\end{enumerate}


\setcounter{chapter}{6}
\chapter{Collections Framework}
\label{chap:7}
\label{h:7:7-collections-framework}
The Collections Framework is Java’s standard toolkit for storing and processing groups of objects. It gives you ready-made data structures (lists, sets, maps, queues) so you don’t write your own array-resizing or hash-table code. This chapter walks through the interfaces, the concrete classes that implement them, and the mistakes that come up most often in code review. The goal is that you can pick the right collection quickly and explain \emph{why} it is right.

\section[Collections Framework Overview]{Collections Framework Overview}
\label{h:7:collections-framework-overview}
A \textbf{collection} is an object that holds other objects, called \emph{elements}. The Collections Framework is a set of interfaces (contracts) and classes (implementations) that all work together. Almost every Java program uses it, so interviewers expect you to know the shape of it cold.

Here is the interface hierarchy you need in your head. \code{Iterable} is the root because “something you can loop over” is the most basic idea. \code{Map} is drawn separately because a map does not extend \code{Collection} — it holds key-value pairs, not single elements.

\begin{codebox}
\begin{Verbatim}
Iterable<E>
   |
   +-- Collection<E>
          |
          +-- List<E>            (ordered, index-based, duplicates allowed)
          |      +-- ArrayList
          |      +-- LinkedList
          |      +-- Vector (legacy)
          |
          +-- Set<E>              (no duplicates)
          |      +-- HashSet
          |      |      +-- LinkedHashSet
          |      +-- SortedSet<E>
          |             +-- NavigableSet<E>
          |                    +-- TreeSet
          |
          +-- Queue<E>            (head/tail processing)
                 +-- PriorityQueue
                 +-- Deque<E>      (double-ended queue)
                        +-- ArrayDeque
                        +-- LinkedList (also a List)

Map<K,V>                          (separate hierarchy, NOT a Collection)
   +-- HashMap
   |      +-- LinkedHashMap
   +-- SortedMap<K,V>
   |      +-- NavigableMap<K,V>
   |             +-- TreeMap
   +-- Hashtable (legacy)
   +-- IdentityHashMap
   +-- WeakHashMap
   +-- EnumMap
\end{Verbatim}
\end{codebox}

Java 21 added two small but handy interfaces:

\begin{itemize}
\item 
\textbf{\code{SequencedCollection\textless{}\allowbreak{}E\textgreater{}}} — for collections that have a well-defined \emph{first} and \emph{last} element (like \code{List} and \code{Deque}). It adds \code{getFirst()}, \code{getLast()}, \code{addFirst()}, \code{addLast()}, \code{removeFirst()}, \code{removeLast()}, and \code{reversed()}. \code{List} and \code{Deque} now implement it, so you get these methods on \code{ArrayList}, \code{LinkedList}, and \code{ArrayDeque} for free.

\item 
\textbf{\code{SequencedMap\textless{}\allowbreak{}K,\allowbreak{}V\textgreater{}}} — the map equivalent, with \code{firstEntry()}, \code{lastEntry()}, \code{putFirst()}, \code{putLast()}, and \code{reversed()}. \code{LinkedHashMap} and \code{TreeMap} implement it.

\end{itemize}

\begin{codebox}
\begin{Verbatim}[commandchars=\\\{\}]
\PY{n}{List}\PY{o}{\PYZlt{}}\PY{n}{String}\PY{o}{\PYZgt{}}\PY{+w}{ }\PY{n}{names}\PY{+w}{ }\PY{o}{=}\PY{+w}{ }\PY{k}{new}\PY{+w}{ }\PY{n}{ArrayList}\PY{o}{\PYZlt{}}\PY{o}{\PYZgt{}}\PY{p}{(}\PY{n}{List}\PY{p}{.}\PY{n+na}{of}\PY{p}{(}\PY{l+s}{\PYZdq{}}\PY{l+s}{Ann}\PY{l+s}{\PYZdq{}}\PY{p}{,}\PY{+w}{ }\PY{l+s}{\PYZdq{}}\PY{l+s}{Bob}\PY{l+s}{\PYZdq{}}\PY{p}{,}\PY{+w}{ }\PY{l+s}{\PYZdq{}}\PY{l+s}{Cara}\PY{l+s}{\PYZdq{}}\PY{p}{)}\PY{p}{)}\PY{p}{;}
\PY{n}{System}\PY{p}{.}\PY{n+na}{out}\PY{p}{.}\PY{n+na}{println}\PY{p}{(}\PY{n}{names}\PY{p}{.}\PY{n+na}{getFirst}\PY{p}{(}\PY{p}{)}\PY{p}{)}\PY{p}{;}\PY{+w}{ }\PY{c+c1}{// Ann}
\PY{n}{System}\PY{p}{.}\PY{n+na}{out}\PY{p}{.}\PY{n+na}{println}\PY{p}{(}\PY{n}{names}\PY{p}{.}\PY{n+na}{getLast}\PY{p}{(}\PY{p}{)}\PY{p}{)}\PY{p}{;}\PY{+w}{  }\PY{c+c1}{// Cara}
\PY{n}{List}\PY{o}{\PYZlt{}}\PY{n}{String}\PY{o}{\PYZgt{}}\PY{+w}{ }\PY{n}{reversed}\PY{+w}{ }\PY{o}{=}\PY{+w}{ }\PY{n}{names}\PY{p}{.}\PY{n+na}{reversed}\PY{p}{(}\PY{p}{)}\PY{p}{;}
\PY{n}{System}\PY{p}{.}\PY{n+na}{out}\PY{p}{.}\PY{n+na}{println}\PY{p}{(}\PY{n}{reversed}\PY{p}{)}\PY{p}{;}\PY{+w}{         }\PY{c+c1}{// [Cara, Bob, Ann]}
\end{Verbatim}
\end{codebox}

\textbf{Why the framework matters in a code review}: picking the wrong collection is one of the most common review findings — using a \code{List} where a \code{Set} was meant, using \code{ArrayList} where a \code{Deque} fits better, or forgetting that \code{HashMap} iteration order is unspecified. Knowing the contracts lets you spot these fast.

\section[Collection Interface]{Collection Interface}
\label{h:7:collection-interface}
\code{Collection\textless{}\allowbreak{}E\textgreater{}} is the root interface for anything that is a “bunch of elements” (excluding maps). It defines the operations every collection must support: \code{add}, \code{remove}, \code{contains}, \code{size}, \code{isEmpty}, \code{iterator}, \code{stream}, \code{forEach}, and bulk operations like \code{addAll}, \code{removeAll}, \code{retainAll}.

\begin{codebox}
\begin{Verbatim}[commandchars=\\\{\}]
\PY{n}{Collection}\PY{o}{\PYZlt{}}\PY{n}{String}\PY{o}{\PYZgt{}}\PY{+w}{ }\PY{n}{col}\PY{+w}{ }\PY{o}{=}\PY{+w}{ }\PY{k}{new}\PY{+w}{ }\PY{n}{ArrayList}\PY{o}{\PYZlt{}}\PY{o}{\PYZgt{}}\PY{p}{(}\PY{p}{)}\PY{p}{;}
\PY{n}{col}\PY{p}{.}\PY{n+na}{add}\PY{p}{(}\PY{l+s}{\PYZdq{}}\PY{l+s}{apple}\PY{l+s}{\PYZdq{}}\PY{p}{)}\PY{p}{;}
\PY{n}{col}\PY{p}{.}\PY{n+na}{add}\PY{p}{(}\PY{l+s}{\PYZdq{}}\PY{l+s}{banana}\PY{l+s}{\PYZdq{}}\PY{p}{)}\PY{p}{;}
\PY{n}{System}\PY{p}{.}\PY{n+na}{out}\PY{p}{.}\PY{n+na}{println}\PY{p}{(}\PY{n}{col}\PY{p}{.}\PY{n+na}{contains}\PY{p}{(}\PY{l+s}{\PYZdq{}}\PY{l+s}{apple}\PY{l+s}{\PYZdq{}}\PY{p}{)}\PY{p}{)}\PY{p}{;}\PY{+w}{ }\PY{c+c1}{// true}
\PY{n}{System}\PY{p}{.}\PY{n+na}{out}\PY{p}{.}\PY{n+na}{println}\PY{p}{(}\PY{n}{col}\PY{p}{.}\PY{n+na}{size}\PY{p}{(}\PY{p}{)}\PY{p}{)}\PY{p}{;}\PY{+w}{            }\PY{c+c1}{// 2}
\PY{n}{col}\PY{p}{.}\PY{n+na}{removeIf}\PY{p}{(}\PY{n}{s}\PY{+w}{ }\PY{o}{\PYZhy{}}\PY{o}{\PYZgt{}}\PY{+w}{ }\PY{n}{s}\PY{p}{.}\PY{n+na}{startsWith}\PY{p}{(}\PY{l+s}{\PYZdq{}}\PY{l+s}{b}\PY{l+s}{\PYZdq{}}\PY{p}{)}\PY{p}{)}\PY{p}{;}
\PY{n}{System}\PY{p}{.}\PY{n+na}{out}\PY{p}{.}\PY{n+na}{println}\PY{p}{(}\PY{n}{col}\PY{p}{)}\PY{p}{;}\PY{+w}{                   }\PY{c+c1}{// [apple]}
\end{Verbatim}
\end{codebox}

Key things every \code{Collection} implementation must get right:

\begin{itemize}
\item 
\textbf{\code{equals}/\code{hashCode} contract}: \code{Collection.\allowbreak{}equals} is defined precisely for \code{List} and \code{Set} (element-by-element), so if you put custom objects in a collection, your \code{equals}/\code{hashCode} must be consistent (see Chapter 5).

\item 
\textbf{Optional operations}: some collections (like the ones from \code{List.\allowbreak{}of()}) throw \code{Unsupported\allowbreak{}Operation\allowbreak{}Exception} for \code{add}/\code{remove}. The interface documents these as “optional operations.”

\item 
\textbf{Fail-fast iterators}: most \code{java.\allowbreak{}util} collections detect structural modification during iteration and throw \code{Concurrent\allowbreak{}Modification\allowbreak{}Exception} (covered in detail below).

\end{itemize}

\section[List]{List}
\label{h:7:list}
A \code{List\textless{}\allowbreak{}E\textgreater{}} is an \textbf{ordered} collection accessed by a zero-based index. It allows \textbf{duplicate} elements. Think of it as a resizable array with extra methods.

Key methods: \code{get(\allowbreak{}int)}, \code{set(\allowbreak{}int,\allowbreak{}~\allowbreak{}E)}, \code{add(\allowbreak{}int,\allowbreak{}~\allowbreak{}E)}, \code{remove(\allowbreak{}int)}, \code{indexOf(\allowbreak{}Object)}, \code{subList(\allowbreak{}from,\allowbreak{}~\allowbreak{}to)}, plus everything from \code{Collection} and (since Java 21) \code{SequencedCollection}.

\begin{codebox}
\begin{Verbatim}[commandchars=\\\{\}]
\PY{n}{List}\PY{o}{\PYZlt{}}\PY{n}{String}\PY{o}{\PYZgt{}}\PY{+w}{ }\PY{n}{fruits}\PY{+w}{ }\PY{o}{=}\PY{+w}{ }\PY{k}{new}\PY{+w}{ }\PY{n}{ArrayList}\PY{o}{\PYZlt{}}\PY{o}{\PYZgt{}}\PY{p}{(}\PY{p}{)}\PY{p}{;}
\PY{n}{fruits}\PY{p}{.}\PY{n+na}{add}\PY{p}{(}\PY{l+s}{\PYZdq{}}\PY{l+s}{apple}\PY{l+s}{\PYZdq{}}\PY{p}{)}\PY{p}{;}
\PY{n}{fruits}\PY{p}{.}\PY{n+na}{add}\PY{p}{(}\PY{l+s}{\PYZdq{}}\PY{l+s}{banana}\PY{l+s}{\PYZdq{}}\PY{p}{)}\PY{p}{;}
\PY{n}{fruits}\PY{p}{.}\PY{n+na}{add}\PY{p}{(}\PY{l+m+mi}{1}\PY{p}{,}\PY{+w}{ }\PY{l+s}{\PYZdq{}}\PY{l+s}{kiwi}\PY{l+s}{\PYZdq{}}\PY{p}{)}\PY{p}{;}\PY{+w}{       }\PY{c+c1}{// insert at index 1}
\PY{n}{System}\PY{p}{.}\PY{n+na}{out}\PY{p}{.}\PY{n+na}{println}\PY{p}{(}\PY{n}{fruits}\PY{p}{)}\PY{p}{;}\PY{+w}{  }\PY{c+c1}{// [apple, kiwi, banana]}
\PY{n}{System}\PY{p}{.}\PY{n+na}{out}\PY{p}{.}\PY{n+na}{println}\PY{p}{(}\PY{n}{fruits}\PY{p}{.}\PY{n+na}{get}\PY{p}{(}\PY{l+m+mi}{0}\PY{p}{)}\PY{p}{)}\PY{p}{;}\PY{+w}{ }\PY{c+c1}{// apple}
\end{Verbatim}
\end{codebox}

Main implementations: \code{ArrayList} (array-backed, default choice), \code{LinkedList} (doubly-linked list, rarely the right default), \code{Vector} (legacy, synchronized, avoid).

\code{List.\allowbreak{}equals} compares element order and values, so \code{[\allowbreak{}1,\allowbreak{}~\allowbreak{}2,\allowbreak{}~\allowbreak{}3]} and \code{[\allowbreak{}3,\allowbreak{}~\allowbreak{}2,\allowbreak{}~\allowbreak{}1]} are \textbf{not} equal even though they hold the same elements.

\section[Set]{Set}
\label{h:7:set}
A \code{Set\textless{}\allowbreak{}E\textgreater{}} is a collection with \textbf{no duplicate elements}, based on \code{equals}. It models mathematical sets. It has no index-based access.

Sub-interfaces:

\begin{itemize}
\item 
\code{SortedSet\textless{}\allowbreak{}E\textgreater{}} — elements kept in sorted order (by natural ordering or a \code{Comparator}).

\item 
\code{NavigableSet\textless{}\allowbreak{}E\textgreater{}} — adds navigation methods like \code{floor}, \code{ceiling}, \code{higher}, \code{lower}.

\end{itemize}

\begin{codebox}
\begin{Verbatim}[commandchars=\\\{\}]
\PY{n}{Set}\PY{o}{\PYZlt{}}\PY{n}{String}\PY{o}{\PYZgt{}}\PY{+w}{ }\PY{n}{tags}\PY{+w}{ }\PY{o}{=}\PY{+w}{ }\PY{k}{new}\PY{+w}{ }\PY{n}{HashSet}\PY{o}{\PYZlt{}}\PY{o}{\PYZgt{}}\PY{p}{(}\PY{p}{)}\PY{p}{;}
\PY{n}{tags}\PY{p}{.}\PY{n+na}{add}\PY{p}{(}\PY{l+s}{\PYZdq{}}\PY{l+s}{java}\PY{l+s}{\PYZdq{}}\PY{p}{)}\PY{p}{;}
\PY{n}{tags}\PY{p}{.}\PY{n+na}{add}\PY{p}{(}\PY{l+s}{\PYZdq{}}\PY{l+s}{java}\PY{l+s}{\PYZdq{}}\PY{p}{)}\PY{p}{;}\PY{+w}{ }\PY{c+c1}{// ignored, duplicate}
\PY{n}{tags}\PY{p}{.}\PY{n+na}{add}\PY{p}{(}\PY{l+s}{\PYZdq{}}\PY{l+s}{spring}\PY{l+s}{\PYZdq{}}\PY{p}{)}\PY{p}{;}
\PY{n}{System}\PY{p}{.}\PY{n+na}{out}\PY{p}{.}\PY{n+na}{println}\PY{p}{(}\PY{n}{tags}\PY{p}{.}\PY{n+na}{size}\PY{p}{(}\PY{p}{)}\PY{p}{)}\PY{p}{;}\PY{+w}{ }\PY{c+c1}{// 2}
\PY{n}{System}\PY{p}{.}\PY{n+na}{out}\PY{p}{.}\PY{n+na}{println}\PY{p}{(}\PY{n}{tags}\PY{p}{.}\PY{n+na}{contains}\PY{p}{(}\PY{l+s}{\PYZdq{}}\PY{l+s}{java}\PY{l+s}{\PYZdq{}}\PY{p}{)}\PY{p}{)}\PY{p}{;}\PY{+w}{ }\PY{c+c1}{// true}
\end{Verbatim}
\end{codebox}

Main implementations: \code{HashSet} (fast, no order), \code{LinkedHashSet} (insertion order), \code{TreeSet} (sorted order).

\section[Queue]{Queue}
\label{h:7:queue}
A \code{Queue\textless{}\allowbreak{}E\textgreater{}} models a line of elements processed in a specific order, typically \textbf{FIFO} (first-in-first-out). Core methods come in two flavors: ones that throw exceptions and ones that return a special value (\code{null} or \code{false}) on failure.

{\tablefont
\begin{xltabular}{\linewidth}{@{}L{0.587}L{1.043}L{1.370}@{}}
\toprule
\rowcolor{tablehdr}
\thd{Operation} & \thd{Throws exception} & \thd{Returns special value} \\
\midrule
\endhead
Insert & \code{add(\allowbreak{}e)} & \code{offer(\allowbreak{}e)} \\
Remove & \code{remove()} & \code{poll()} \\
Examine & \code{element()} & \code{peek()} \\
\bottomrule
\end{xltabular}
}

\begin{codebox}
\begin{Verbatim}[commandchars=\\\{\}]
\PY{n}{Queue}\PY{o}{\PYZlt{}}\PY{n}{Integer}\PY{o}{\PYZgt{}}\PY{+w}{ }\PY{n}{queue}\PY{+w}{ }\PY{o}{=}\PY{+w}{ }\PY{k}{new}\PY{+w}{ }\PY{n}{java}\PY{p}{.}\PY{n+na}{util}\PY{p}{.}\PY{n+na}{ArrayDeque}\PY{o}{\PYZlt{}}\PY{o}{\PYZgt{}}\PY{p}{(}\PY{p}{)}\PY{p}{;}
\PY{n}{queue}\PY{p}{.}\PY{n+na}{offer}\PY{p}{(}\PY{l+m+mi}{1}\PY{p}{)}\PY{p}{;}
\PY{n}{queue}\PY{p}{.}\PY{n+na}{offer}\PY{p}{(}\PY{l+m+mi}{2}\PY{p}{)}\PY{p}{;}
\PY{n}{queue}\PY{p}{.}\PY{n+na}{offer}\PY{p}{(}\PY{l+m+mi}{3}\PY{p}{)}\PY{p}{;}
\PY{n}{System}\PY{p}{.}\PY{n+na}{out}\PY{p}{.}\PY{n+na}{println}\PY{p}{(}\PY{n}{queue}\PY{p}{.}\PY{n+na}{poll}\PY{p}{(}\PY{p}{)}\PY{p}{)}\PY{p}{;}\PY{+w}{ }\PY{c+c1}{// 1 (FIFO order)}
\PY{n}{System}\PY{p}{.}\PY{n+na}{out}\PY{p}{.}\PY{n+na}{println}\PY{p}{(}\PY{n}{queue}\PY{p}{.}\PY{n+na}{peek}\PY{p}{(}\PY{p}{)}\PY{p}{)}\PY{p}{;}\PY{+w}{ }\PY{c+c1}{// 2, does not remove}
\end{Verbatim}
\end{codebox}

Main implementations: \code{PriorityQueue} (elements ordered by priority, not insertion order), \code{ArrayDeque} (fast general-purpose queue/stack), \code{LinkedList} (also implements \code{Queue}, rarely the best choice).

Prefer the \code{offer}/\code{poll}/\code{peek} family in production code — they signal “empty” or “full” with a return value instead of an exception, which is usually easier to handle in a loop.

\section[Deque]{Deque}
\label{h:7:deque}
A \code{Deque\textless{}\allowbreak{}E\textgreater{}} (“deck”, \textbf{d}ouble-\textbf{e}nded \textbf{queue}) supports insertion and removal at \textbf{both ends}. It can act as a \code{Queue} (FIFO), a \code{Stack} (LIFO), or both.

\begin{codebox}
\begin{Verbatim}[commandchars=\\\{\}]
\PY{n}{Deque}\PY{o}{\PYZlt{}}\PY{n}{Integer}\PY{o}{\PYZgt{}}\PY{+w}{ }\PY{n}{stack}\PY{+w}{ }\PY{o}{=}\PY{+w}{ }\PY{k}{new}\PY{+w}{ }\PY{n}{java}\PY{p}{.}\PY{n+na}{util}\PY{p}{.}\PY{n+na}{ArrayDeque}\PY{o}{\PYZlt{}}\PY{o}{\PYZgt{}}\PY{p}{(}\PY{p}{)}\PY{p}{;}
\PY{n}{stack}\PY{p}{.}\PY{n+na}{push}\PY{p}{(}\PY{l+m+mi}{1}\PY{p}{)}\PY{p}{;}\PY{+w}{   }\PY{c+c1}{// addFirst}
\PY{n}{stack}\PY{p}{.}\PY{n+na}{push}\PY{p}{(}\PY{l+m+mi}{2}\PY{p}{)}\PY{p}{;}
\PY{n}{stack}\PY{p}{.}\PY{n+na}{push}\PY{p}{(}\PY{l+m+mi}{3}\PY{p}{)}\PY{p}{;}
\PY{n}{System}\PY{p}{.}\PY{n+na}{out}\PY{p}{.}\PY{n+na}{println}\PY{p}{(}\PY{n}{stack}\PY{p}{.}\PY{n+na}{pop}\PY{p}{(}\PY{p}{)}\PY{p}{)}\PY{p}{;}\PY{+w}{ }\PY{c+c1}{// 3 (LIFO, like a Stack)}

\PY{n}{Deque}\PY{o}{\PYZlt{}}\PY{n}{Integer}\PY{o}{\PYZgt{}}\PY{+w}{ }\PY{n}{dq}\PY{+w}{ }\PY{o}{=}\PY{+w}{ }\PY{k}{new}\PY{+w}{ }\PY{n}{java}\PY{p}{.}\PY{n+na}{util}\PY{p}{.}\PY{n+na}{ArrayDeque}\PY{o}{\PYZlt{}}\PY{o}{\PYZgt{}}\PY{p}{(}\PY{p}{)}\PY{p}{;}
\PY{n}{dq}\PY{p}{.}\PY{n+na}{addFirst}\PY{p}{(}\PY{l+m+mi}{10}\PY{p}{)}\PY{p}{;}
\PY{n}{dq}\PY{p}{.}\PY{n+na}{addLast}\PY{p}{(}\PY{l+m+mi}{20}\PY{p}{)}\PY{p}{;}
\PY{n}{System}\PY{p}{.}\PY{n+na}{out}\PY{p}{.}\PY{n+na}{println}\PY{p}{(}\PY{n}{dq}\PY{p}{)}\PY{p}{;}\PY{+w}{ }\PY{c+c1}{// [10, 20]}
\PY{n}{System}\PY{p}{.}\PY{n+na}{out}\PY{p}{.}\PY{n+na}{println}\PY{p}{(}\PY{n}{dq}\PY{p}{.}\PY{n+na}{pollFirst}\PY{p}{(}\PY{p}{)}\PY{p}{)}\PY{p}{;}\PY{+w}{ }\PY{c+c1}{// 10}
\PY{n}{System}\PY{p}{.}\PY{n+na}{out}\PY{p}{.}\PY{n+na}{println}\PY{p}{(}\PY{n}{dq}\PY{p}{.}\PY{n+na}{pollLast}\PY{p}{(}\PY{p}{)}\PY{p}{)}\PY{p}{;}\PY{+w}{  }\PY{c+c1}{// 20}
\end{Verbatim}
\end{codebox}

Because \code{Deque} extends \code{SequencedCollection}, it also has \code{getFirst()}, \code{getLast()}, and \code{reversed()}.

\textbf{Deque replaces both \code{Stack} and \code{java.\allowbreak{}util.\allowbreak{}Stack}/legacy usage.} The official Javadoc recommends \code{ArrayDeque} over \code{java.\allowbreak{}util.\allowbreak{}Stack} (legacy, synchronized, extends \code{Vector}) for stack behavior, and over \code{LinkedList} for queue behavior, because \code{ArrayDeque} is faster and has no synchronization overhead.

\section[Map]{Map}
\label{h:7:map}
A \code{Map\textless{}\allowbreak{}K,\allowbreak{}V\textgreater{}} stores \textbf{key-value pairs}. Keys are unique; each key maps to exactly one value. \code{Map} does \textbf{not} extend \code{Collection} because its elements are pairs, not single objects — but it does expose collection views: \code{keySet()}, \code{values()}, and \code{entrySet()}.

\begin{codebox}
\begin{Verbatim}[commandchars=\\\{\}]
\PY{n}{Map}\PY{o}{\PYZlt{}}\PY{n}{String}\PY{p}{,}\PY{+w}{ }\PY{n}{Integer}\PY{o}{\PYZgt{}}\PY{+w}{ }\PY{n}{ages}\PY{+w}{ }\PY{o}{=}\PY{+w}{ }\PY{k}{new}\PY{+w}{ }\PY{n}{HashMap}\PY{o}{\PYZlt{}}\PY{o}{\PYZgt{}}\PY{p}{(}\PY{p}{)}\PY{p}{;}
\PY{n}{ages}\PY{p}{.}\PY{n+na}{put}\PY{p}{(}\PY{l+s}{\PYZdq{}}\PY{l+s}{Ann}\PY{l+s}{\PYZdq{}}\PY{p}{,}\PY{+w}{ }\PY{l+m+mi}{30}\PY{p}{)}\PY{p}{;}
\PY{n}{ages}\PY{p}{.}\PY{n+na}{put}\PY{p}{(}\PY{l+s}{\PYZdq{}}\PY{l+s}{Bob}\PY{l+s}{\PYZdq{}}\PY{p}{,}\PY{+w}{ }\PY{l+m+mi}{25}\PY{p}{)}\PY{p}{;}
\PY{n}{ages}\PY{p}{.}\PY{n+na}{put}\PY{p}{(}\PY{l+s}{\PYZdq{}}\PY{l+s}{Ann}\PY{l+s}{\PYZdq{}}\PY{p}{,}\PY{+w}{ }\PY{l+m+mi}{31}\PY{p}{)}\PY{p}{;}\PY{+w}{ }\PY{c+c1}{// overwrites, put returns old value (30)}

\PY{k}{for}\PY{+w}{ }\PY{p}{(}\PY{n}{Map}\PY{p}{.}\PY{n+na}{Entry}\PY{o}{\PYZlt{}}\PY{n}{String}\PY{p}{,}\PY{+w}{ }\PY{n}{Integer}\PY{o}{\PYZgt{}}\PY{+w}{ }\PY{n}{entry}\PY{+w}{ }\PY{p}{:}\PY{+w}{ }\PY{n}{ages}\PY{p}{.}\PY{n+na}{entrySet}\PY{p}{(}\PY{p}{)}\PY{p}{)}\PY{+w}{ }\PY{p}{\PYZob{}}
\PY{+w}{    }\PY{n}{System}\PY{p}{.}\PY{n+na}{out}\PY{p}{.}\PY{n+na}{println}\PY{p}{(}\PY{n}{entry}\PY{p}{.}\PY{n+na}{getKey}\PY{p}{(}\PY{p}{)}\PY{+w}{ }\PY{o}{+}\PY{+w}{ }\PY{l+s}{\PYZdq{}}\PY{l+s}{ \PYZhy{}\PYZgt{} }\PY{l+s}{\PYZdq{}}\PY{+w}{ }\PY{o}{+}\PY{+w}{ }\PY{n}{entry}\PY{p}{.}\PY{n+na}{getValue}\PY{p}{(}\PY{p}{)}\PY{p}{)}\PY{p}{;}
\PY{p}{\PYZcb{}}
\PY{c+c1}{// Ann \PYZhy{}\PYZgt{} 31}
\PY{c+c1}{// Bob \PYZhy{}\PYZgt{} 25 (order not guaranteed for plain HashMap)}
\end{Verbatim}
\end{codebox}

Useful modern default methods (Java 8+), all worth knowing for reviews:

\begin{codebox}
\begin{Verbatim}[commandchars=\\\{\}]
\PY{n}{Map}\PY{o}{\PYZlt{}}\PY{n}{String}\PY{p}{,}\PY{+w}{ }\PY{n}{Integer}\PY{o}{\PYZgt{}}\PY{+w}{ }\PY{n}{scores}\PY{+w}{ }\PY{o}{=}\PY{+w}{ }\PY{k}{new}\PY{+w}{ }\PY{n}{HashMap}\PY{o}{\PYZlt{}}\PY{o}{\PYZgt{}}\PY{p}{(}\PY{p}{)}\PY{p}{;}
\PY{n}{scores}\PY{p}{.}\PY{n+na}{putIfAbsent}\PY{p}{(}\PY{l+s}{\PYZdq{}}\PY{l+s}{Ann}\PY{l+s}{\PYZdq{}}\PY{p}{,}\PY{+w}{ }\PY{l+m+mi}{10}\PY{p}{)}\PY{p}{;}\PY{+w}{        }\PY{c+c1}{// inserts only if key missing}
\PY{n}{scores}\PY{p}{.}\PY{n+na}{merge}\PY{p}{(}\PY{l+s}{\PYZdq{}}\PY{l+s}{Ann}\PY{l+s}{\PYZdq{}}\PY{p}{,}\PY{+w}{ }\PY{l+m+mi}{5}\PY{p}{,}\PY{+w}{ }\PY{n}{Integer}\PY{p}{:}\PY{p}{:}\PY{n}{sum}\PY{p}{)}\PY{p}{;}\PY{+w}{ }\PY{c+c1}{// Ann \PYZhy{}\PYZgt{} 15}
\PY{n}{scores}\PY{p}{.}\PY{n+na}{computeIfAbsent}\PY{p}{(}\PY{l+s}{\PYZdq{}}\PY{l+s}{Bob}\PY{l+s}{\PYZdq{}}\PY{p}{,}\PY{+w}{ }\PY{n}{k}\PY{+w}{ }\PY{o}{\PYZhy{}}\PY{o}{\PYZgt{}}\PY{+w}{ }\PY{l+m+mi}{0}\PY{p}{)}\PY{p}{;}\PY{+w}{ }\PY{c+c1}{// Bob \PYZhy{}\PYZgt{} 0}
\PY{n}{scores}\PY{p}{.}\PY{n+na}{compute}\PY{p}{(}\PY{l+s}{\PYZdq{}}\PY{l+s}{Bob}\PY{l+s}{\PYZdq{}}\PY{p}{,}\PY{+w}{ }\PY{p}{(}\PY{n}{k}\PY{p}{,}\PY{+w}{ }\PY{n}{v}\PY{p}{)}\PY{+w}{ }\PY{o}{\PYZhy{}}\PY{o}{\PYZgt{}}\PY{+w}{ }\PY{n}{v}\PY{+w}{ }\PY{o}{+}\PY{+w}{ }\PY{l+m+mi}{1}\PY{p}{)}\PY{p}{;}\PY{+w}{ }\PY{c+c1}{// Bob \PYZhy{}\PYZgt{} 1}
\PY{n}{scores}\PY{p}{.}\PY{n+na}{getOrDefault}\PY{p}{(}\PY{l+s}{\PYZdq{}}\PY{l+s}{Cara}\PY{l+s}{\PYZdq{}}\PY{p}{,}\PY{+w}{ }\PY{o}{\PYZhy{}}\PY{l+m+mi}{1}\PY{p}{)}\PY{p}{;}\PY{+w}{       }\PY{c+c1}{// \PYZhy{}1, no exception}
\end{Verbatim}
\end{codebox}

Sub-interfaces:

\begin{itemize}
\item 
\code{SortedMap\textless{}\allowbreak{}K,\allowbreak{}V\textgreater{}} — keys kept sorted.

\item 
\code{NavigableMap\textless{}\allowbreak{}K,\allowbreak{}V\textgreater{}} — adds \code{floorKey}, \code{ceilingKey}, \code{firstEntry}, \code{lastEntry}, etc.

\item 
\code{SequencedMap\textless{}\allowbreak{}K,\allowbreak{}V\textgreater{}} (Java 21) — \code{putFirst}, \code{putLast}, \code{firstEntry}, \code{lastEntry}, \code{reversed()}. Implemented by \code{LinkedHashMap} and \code{TreeMap}.

\end{itemize}

Main implementations: \code{HashMap}, \code{LinkedHashMap}, \code{TreeMap}, plus special-purpose ones: \code{IdentityHashMap}, \code{WeakHashMap}, \code{EnumMap}.

\section[Master “Which Collection Should I Pick?” Table]{Master “Which Collection Should I Pick?” Table}
\label{h:7:master-which-collection-should-i-pick-table}
Use this as your mental cheat sheet in an interview. “Amortized” means the average cost over many operations; an occasional resize can cost more.

{\tablefont
\begin{xltabular}{\linewidth}{@{}L{1.334}L{0.446}L{0.821}L{1.437}L{0.446}L{0.821}L{1.539}L{0.821}L{1.334}@{}}
\toprule
\rowcolor{tablehdr}
\thd{Collection} & \thd{Get by index} & \thd{Get by key} & \thd{Add/remove at end} & \thd{Add/remove at front} & \thd{Contains} & \thd{Ordering} & \thd{Duplicates} & \thd{Null elements} \\
\midrule
\endhead
\code{ArrayList} & O(1) & — & O(1) amortized & O(n) & O(n) & insertion order & yes & yes \\
\code{LinkedList} & O(n) & — & O(1) & O(1) & O(n) & insertion order & yes & yes \\
\code{HashMap} & — & O(1) avg & — & — & O(1) avg (key) & unspecified & keys: no, values: yes & 1 null key, many null values \\
\code{LinkedHashMap} & — & O(1) avg & — & — & O(1) avg & insertion (or access) order & keys: no & same as HashMap \\
\code{TreeMap} & — & O(log n) & — & — & O(log n) & sorted by key & keys: no & no null keys \\
\code{HashSet} & — & — & O(1) avg & — & O(1) avg & unspecified & no & one null \\
\code{LinkedHashSet} & — & — & O(1) avg & — & O(1) avg & insertion order & no & one null \\
\code{TreeSet} & — & — & O(log n) & — & O(log n) & sorted & no & no null \\
\code{ArrayDeque} & — & — & O(1) & O(1) & O(n) & insertion order & yes & no null \\
\code{PriorityQueue} & — & — & O(log n) insert & O(log n) poll head & O(n) & priority order (only head is guaranteed smallest) & yes & no null \\
\bottomrule
\end{xltabular}
}

Quick decision guide:

\begin{itemize}
\item 
Need index access, mostly append at the end → \textbf{\code{ArrayList}}.

\item 
Need a stack or a queue, or fast add/remove at both ends → \textbf{\code{ArrayDeque}}.

\item 
Need uniqueness, don’t care about order → \textbf{\code{HashSet}}.

\item 
Need uniqueness and insertion order preserved → \textbf{\code{LinkedHashSet}}.

\item 
Need uniqueness and sorted order → \textbf{\code{TreeSet}}.

\item 
Need key-value lookups, don’t care about order → \textbf{\code{HashMap}}.

\item 
Need key-value lookups with predictable iteration order → \textbf{\code{LinkedHashMap}}.

\item 
Need key-value lookups sorted by key, or range queries → \textbf{\code{TreeMap}}.

\item 
Need “always process smallest/largest next” → \textbf{\code{PriorityQueue}}.

\item 
Multithreaded access → see \hyperref[h:7:concurrent-collections]{Concurrent Collections}.

\end{itemize}

\section[ArrayList]{ArrayList}
\label{h:7:arraylist}
\textbf{How it works internally}: \code{ArrayList} wraps a plain \code{Object[]} array. When the array is full and you add another element, it allocates a new array (roughly 1.5x the old capacity) and copies everything over (\code{Arrays.\allowbreak{}copyOf}). This is why appends are “amortized” O(1) — most calls are cheap, but occasionally one call pays for the copy.

\begin{codebox}
\begin{Verbatim}[commandchars=\\\{\}]
\PY{n}{List}\PY{o}{\PYZlt{}}\PY{n}{Integer}\PY{o}{\PYZgt{}}\PY{+w}{ }\PY{n}{list}\PY{+w}{ }\PY{o}{=}\PY{+w}{ }\PY{k}{new}\PY{+w}{ }\PY{n}{ArrayList}\PY{o}{\PYZlt{}}\PY{o}{\PYZgt{}}\PY{p}{(}\PY{p}{)}\PY{p}{;}\PY{+w}{       }\PY{c+c1}{// default capacity, grows as needed}
\PY{n}{List}\PY{o}{\PYZlt{}}\PY{n}{Integer}\PY{o}{\PYZgt{}}\PY{+w}{ }\PY{n}{sized}\PY{+w}{ }\PY{o}{=}\PY{+w}{ }\PY{k}{new}\PY{+w}{ }\PY{n}{ArrayList}\PY{o}{\PYZlt{}}\PY{o}{\PYZgt{}}\PY{p}{(}\PY{l+m+mi}{1000}\PY{p}{)}\PY{p}{;}\PY{+w}{   }\PY{c+c1}{// pre\PYZhy{}size to avoid early resizes}
\PY{k}{for}\PY{+w}{ }\PY{p}{(}\PY{k+kt}{int}\PY{+w}{ }\PY{n}{i}\PY{+w}{ }\PY{o}{=}\PY{+w}{ }\PY{l+m+mi}{0}\PY{p}{;}\PY{+w}{ }\PY{n}{i}\PY{+w}{ }\PY{o}{\PYZlt{}}\PY{+w}{ }\PY{l+m+mi}{5}\PY{p}{;}\PY{+w}{ }\PY{n}{i}\PY{o}{+}\PY{o}{+}\PY{p}{)}\PY{+w}{ }\PY{n}{list}\PY{p}{.}\PY{n+na}{add}\PY{p}{(}\PY{n}{i}\PY{p}{)}\PY{p}{;}
\PY{n}{list}\PY{p}{.}\PY{n+na}{remove}\PY{p}{(}\PY{n}{Integer}\PY{p}{.}\PY{n+na}{valueOf}\PY{p}{(}\PY{l+m+mi}{2}\PY{p}{)}\PY{p}{)}\PY{p}{;}\PY{+w}{  }\PY{c+c1}{// removes value 2, not index 2}
\PY{n}{list}\PY{p}{.}\PY{n+na}{remove}\PY{p}{(}\PY{l+m+mi}{2}\PY{p}{)}\PY{p}{;}\PY{+w}{                   }\PY{c+c1}{// removes index 2 (int overload)}
\PY{n}{System}\PY{p}{.}\PY{n+na}{out}\PY{p}{.}\PY{n+na}{println}\PY{p}{(}\PY{n}{list}\PY{p}{)}\PY{p}{;}\PY{+w}{         }\PY{c+c1}{// [0, 1, 4] after both removals}
\end{Verbatim}
\end{codebox}

\textbf{Big-O}:

\begin{itemize}
\item 
\code{get(\allowbreak{}index)} / \code{set(\allowbreak{}index,\allowbreak{}~\allowbreak{}e)}: O(1) — direct array access.

\item 
\code{add(\allowbreak{}e)} (append at end): O(1) amortized, O(n) worst case (on resize).

\item 
\code{add(\allowbreak{}index,\allowbreak{}~\allowbreak{}e)} / \code{remove(\allowbreak{}index)} in the middle or front: O(n) — must shift elements.

\item 
\code{contains} / \code{indexOf}: O(n) — linear scan.

\end{itemize}

\textbf{Ordering}: preserves insertion order (it’s index-based).

\textbf{Null handling}: allows multiple \code{null} elements.

\textbf{Thread safety}: not thread-safe. Concurrent modification from multiple threads can corrupt internal state or throw \code{Concurrent\allowbreak{}Modification\allowbreak{}Exception}.

\textbf{When to choose it}: default \code{List} choice. Use it when you mostly read by index or append at the end, and rarely insert/remove in the middle or at the front.

\section[LinkedList]{LinkedList}
\label{h:7:linkedlist}
\textbf{How it works internally}: a doubly-linked list of nodes; each node holds the element plus references to the previous and next node. It implements both \code{List} and \code{Deque}.

\begin{codebox}
\begin{Verbatim}[commandchars=\\\{\}]
\PY{n}{LinkedList}\PY{o}{\PYZlt{}}\PY{n}{String}\PY{o}{\PYZgt{}}\PY{+w}{ }\PY{n}{chain}\PY{+w}{ }\PY{o}{=}\PY{+w}{ }\PY{k}{new}\PY{+w}{ }\PY{n}{LinkedList}\PY{o}{\PYZlt{}}\PY{o}{\PYZgt{}}\PY{p}{(}\PY{p}{)}\PY{p}{;}
\PY{n}{chain}\PY{p}{.}\PY{n+na}{addLast}\PY{p}{(}\PY{l+s}{\PYZdq{}}\PY{l+s}{b}\PY{l+s}{\PYZdq{}}\PY{p}{)}\PY{p}{;}
\PY{n}{chain}\PY{p}{.}\PY{n+na}{addFirst}\PY{p}{(}\PY{l+s}{\PYZdq{}}\PY{l+s}{a}\PY{l+s}{\PYZdq{}}\PY{p}{)}\PY{p}{;}
\PY{n}{chain}\PY{p}{.}\PY{n+na}{addLast}\PY{p}{(}\PY{l+s}{\PYZdq{}}\PY{l+s}{c}\PY{l+s}{\PYZdq{}}\PY{p}{)}\PY{p}{;}
\PY{n}{System}\PY{p}{.}\PY{n+na}{out}\PY{p}{.}\PY{n+na}{println}\PY{p}{(}\PY{n}{chain}\PY{p}{)}\PY{p}{;}\PY{+w}{ }\PY{c+c1}{// [a, b, c]}
\PY{n}{chain}\PY{p}{.}\PY{n+na}{removeFirst}\PY{p}{(}\PY{p}{)}\PY{p}{;}
\PY{n}{System}\PY{p}{.}\PY{n+na}{out}\PY{p}{.}\PY{n+na}{println}\PY{p}{(}\PY{n}{chain}\PY{p}{)}\PY{p}{;}\PY{+w}{ }\PY{c+c1}{// [b, c]}
\end{Verbatim}
\end{codebox}

\textbf{Big-O}:

\begin{itemize}
\item 
\code{addFirst} / \code{addLast} / \code{removeFirst} / \code{removeLast}: O(1) — just pointer updates.

\item 
\code{get(\allowbreak{}index)}: O(n) — must walk the list from the nearer end.

\item 
\code{add(\allowbreak{}index,\allowbreak{}~\allowbreak{}e)} / \code{remove(\allowbreak{}index)}: O(n) to find the position, O(1) to splice.

\end{itemize}

\textbf{Ordering}: insertion order.

\textbf{Null handling}: allows \code{null} elements.

\textbf{Thread safety}: not thread-safe.

\textbf{When to choose it}: rarely the best default. Choose it only when you need frequent insert/remove at both ends \emph{and} also need \code{List} semantics (index access) in the same structure. For pure queue/stack behavior, prefer \code{ArrayDeque} — it’s faster and uses less memory (no per-node object overhead, no boxing of pointers).

\begin{codebox}
\begin{Verbatim}[commandchars=\\\{\}]
\PY{c+c1}{// Common mistake: using LinkedList as a general\PYZhy{}purpose List}
\PY{n}{List}\PY{o}{\PYZlt{}}\PY{n}{Integer}\PY{o}{\PYZgt{}}\PY{+w}{ }\PY{n}{bad}\PY{+w}{ }\PY{o}{=}\PY{+w}{ }\PY{k}{new}\PY{+w}{ }\PY{n}{LinkedList}\PY{o}{\PYZlt{}}\PY{o}{\PYZgt{}}\PY{p}{(}\PY{p}{)}\PY{p}{;}
\PY{k}{for}\PY{+w}{ }\PY{p}{(}\PY{k+kt}{int}\PY{+w}{ }\PY{n}{i}\PY{+w}{ }\PY{o}{=}\PY{+w}{ }\PY{l+m+mi}{0}\PY{p}{;}\PY{+w}{ }\PY{n}{i}\PY{+w}{ }\PY{o}{\PYZlt{}}\PY{+w}{ }\PY{l+m+mi}{100\PYZus{}000}\PY{p}{;}\PY{+w}{ }\PY{n}{i}\PY{o}{+}\PY{o}{+}\PY{p}{)}\PY{+w}{ }\PY{n}{bad}\PY{p}{.}\PY{n+na}{add}\PY{p}{(}\PY{n}{i}\PY{p}{)}\PY{p}{;}
\PY{c+c1}{// bad.get(50\PYZus{}000) is O(n) \PYZhy{}\PYZhy{} walks half the list every time!}
\end{Verbatim}
\end{codebox}

\section[HashMap]{HashMap}
\label{h:7:hashmap}
\code{HashMap} is the workhorse map implementation. Understanding its internals is one of the most frequently asked collection topics in interviews.

\textbf{How it works internally}:

\begin{enumerate}
\item 
\textbf{Buckets}: internally, \code{HashMap} holds an array called the \emph{table}. Each slot in that array is called a \textbf{bucket}. A bucket can hold zero, one, or many entries.

\item 
\textbf{Hashing}: when you call \code{put(\allowbreak{}key,\allowbreak{}~\allowbreak{}value)}, the map computes \code{key.\allowbreak{}hashCode()}, then applies an internal “spread” function (XORs the hash with its own upper bits) to reduce collisions, then uses \code{hash~\allowbreak{}\&\allowbreak{}~\allowbreak{}(\allowbreak{}table.\allowbreak{}length~\allowbreak{}-\allowbreak{}~\allowbreak{}1)} to pick a bucket index. This is why table size is always a power of two — \code{\&\allowbreak{}~\allowbreak{}(\allowbreak{}length~\allowbreak{}-\allowbreak{}~\allowbreak{}1)} is a fast way to do \code{\%~\allowbreak{}length} when \code{length} is a power of two.

\item 
\textbf{Collision handling}: if two keys land in the same bucket, they used to be stored as a \textbf{linked list} of entries in that bucket. You look up a key by finding its bucket, then walking the list and comparing with \code{equals()}.

\item 
\textbf{Treeification (JDK 8+)}: if a single bucket’s linked list grows to \textbf{8 or more} entries (and the table has at least 64 buckets), that bucket is converted into a small \textbf{red-black tree} instead of a linked list. This turns worst-case lookup within a bucket from O(n) into O(log n). This mostly protects against pathological hash collisions (e.g., poorly written \code{hashCode()} or adversarial input) — it is a safety net, not something you should rely on for normal use. If entries in a treeified bucket drop below 6, it converts back to a linked list.

\item 
\textbf{Resizing / load factor}: the map tracks a \textbf{load factor}, default \code{0.\allowbreak{}75}. When \code{size~\allowbreak{}\textgreater{}\allowbreak{}~\allowbreak{}capacity~\allowbreak{}*\allowbreak{}~\allowbreak{}load\allowbreak{}Factor}, the table \textbf{doubles} in size and every entry is rehashed into the new, larger table (a “resize” or “rehash”). Default initial capacity is 16, so the first resize happens after the 13th entry (16 * 0.75 = 12). Resizing is O(n) but happens rarely, so \code{put} is amortized O(1).

\end{enumerate}

\begin{codebox}
\begin{Verbatim}[commandchars=\\\{\}]
\PY{n}{Map}\PY{o}{\PYZlt{}}\PY{n}{String}\PY{p}{,}\PY{+w}{ }\PY{n}{Integer}\PY{o}{\PYZgt{}}\PY{+w}{ }\PY{n}{inventory}\PY{+w}{ }\PY{o}{=}\PY{+w}{ }\PY{k}{new}\PY{+w}{ }\PY{n}{HashMap}\PY{o}{\PYZlt{}}\PY{o}{\PYZgt{}}\PY{p}{(}\PY{p}{)}\PY{p}{;}\PY{+w}{ }\PY{c+c1}{// capacity 16, load factor 0.75}
\PY{n}{inventory}\PY{p}{.}\PY{n+na}{put}\PY{p}{(}\PY{l+s}{\PYZdq{}}\PY{l+s}{apples}\PY{l+s}{\PYZdq{}}\PY{p}{,}\PY{+w}{ }\PY{l+m+mi}{10}\PY{p}{)}\PY{p}{;}
\PY{n}{inventory}\PY{p}{.}\PY{n+na}{put}\PY{p}{(}\PY{l+s}{\PYZdq{}}\PY{l+s}{bananas}\PY{l+s}{\PYZdq{}}\PY{p}{,}\PY{+w}{ }\PY{l+m+mi}{20}\PY{p}{)}\PY{p}{;}
\PY{n}{inventory}\PY{p}{.}\PY{n+na}{put}\PY{p}{(}\PY{k+kc}{null}\PY{p}{,}\PY{+w}{ }\PY{l+m+mi}{0}\PY{p}{)}\PY{p}{;}\PY{+w}{          }\PY{c+c1}{// one null key is allowed}
\PY{n}{inventory}\PY{p}{.}\PY{n+na}{put}\PY{p}{(}\PY{l+s}{\PYZdq{}}\PY{l+s}{apples}\PY{l+s}{\PYZdq{}}\PY{p}{,}\PY{+w}{ }\PY{l+m+mi}{15}\PY{p}{)}\PY{p}{;}\PY{+w}{     }\PY{c+c1}{// overwrites, returns old value 10}

\PY{n}{System}\PY{p}{.}\PY{n+na}{out}\PY{p}{.}\PY{n+na}{println}\PY{p}{(}\PY{n}{inventory}\PY{p}{.}\PY{n+na}{get}\PY{p}{(}\PY{l+s}{\PYZdq{}}\PY{l+s}{apples}\PY{l+s}{\PYZdq{}}\PY{p}{)}\PY{p}{)}\PY{p}{;}\PY{+w}{ }\PY{c+c1}{// 15}
\PY{n}{System}\PY{p}{.}\PY{n+na}{out}\PY{p}{.}\PY{n+na}{println}\PY{p}{(}\PY{n}{inventory}\PY{p}{.}\PY{n+na}{get}\PY{p}{(}\PY{l+s}{\PYZdq{}}\PY{l+s}{kiwi}\PY{l+s}{\PYZdq{}}\PY{p}{)}\PY{p}{)}\PY{p}{;}\PY{+w}{   }\PY{c+c1}{// null: no such key}
\PY{n}{System}\PY{p}{.}\PY{n+na}{out}\PY{p}{.}\PY{n+na}{println}\PY{p}{(}\PY{n}{inventory}\PY{p}{.}\PY{n+na}{get}\PY{p}{(}\PY{k+kc}{null}\PY{p}{)}\PY{p}{)}\PY{p}{;}\PY{+w}{     }\PY{c+c1}{// 0}
\end{Verbatim}
\end{codebox}

\begin{codebox}
\begin{Verbatim}[commandchars=\\\{\}]
\PY{c+c1}{// Pre\PYZhy{}sizing avoids repeated resizes when you know the approximate size}
\PY{n}{Map}\PY{o}{\PYZlt{}}\PY{n}{String}\PY{p}{,}\PY{+w}{ }\PY{n}{String}\PY{o}{\PYZgt{}}\PY{+w}{ }\PY{n}{big}\PY{+w}{ }\PY{o}{=}\PY{+w}{ }\PY{k}{new}\PY{+w}{ }\PY{n}{HashMap}\PY{o}{\PYZlt{}}\PY{o}{\PYZgt{}}\PY{p}{(}\PY{l+m+mi}{200}\PY{p}{)}\PY{p}{;}\PY{+w}{ }\PY{c+c1}{// avoids \PYZti{}4 resize/rehash cycles}
\end{Verbatim}
\end{codebox}

\textbf{Big-O}: \code{get}/\code{put}/\code{remove}/\code{containsKey}: O(1) average, O(log n) worst case per bucket after treeification (was O(n) before JDK 8). This assumes a good \code{hashCode()} distribution; a bad \code{hashCode()} (e.g., always returning \code{0}) degrades performance because everything lands in one bucket.

\textbf{Ordering}: \textbf{unspecified and not guaranteed to be stable} across JDK versions or even across resizes of the same map. Never rely on \code{HashMap} iteration order.

\textbf{Null handling}: allows exactly \textbf{one} \code{null} key and any number of \code{null} values.

\textbf{Thread safety}: not thread-safe. Concurrent \code{put} calls from multiple threads can corrupt the internal structure (classic bug: infinite loop during resize under older JDKs, or lost updates). Use \code{ConcurrentHashMap} instead.

\textbf{When to choose it}: the default map. Use it whenever you need key-value lookups and don’t care about iteration order.

\begin{codebox}
\begin{Verbatim}[commandchars=\\\{\}]
\PY{c+c1}{// hashCode()/equals() contract matters for keys!}
\PY{k+kd}{class} \PY{n+nc}{BadKey}\PY{+w}{ }\PY{p}{\PYZob{}}
\PY{+w}{    }\PY{k+kt}{int}\PY{+w}{ }\PY{n}{id}\PY{p}{;}
\PY{+w}{    }\PY{n}{BadKey}\PY{p}{(}\PY{k+kt}{int}\PY{+w}{ }\PY{n}{id}\PY{p}{)}\PY{+w}{ }\PY{p}{\PYZob{}}\PY{+w}{ }\PY{k}{this}\PY{p}{.}\PY{n+na}{id}\PY{+w}{ }\PY{o}{=}\PY{+w}{ }\PY{n}{id}\PY{p}{;}\PY{+w}{ }\PY{p}{\PYZcb{}}
\PY{+w}{    }\PY{c+c1}{// no equals()/hashCode() override \PYZhy{}\PYZgt{} uses Object identity}
\PY{p}{\PYZcb{}}
\PY{n}{Map}\PY{o}{\PYZlt{}}\PY{n}{BadKey}\PY{p}{,}\PY{+w}{ }\PY{n}{String}\PY{o}{\PYZgt{}}\PY{+w}{ }\PY{n}{map}\PY{+w}{ }\PY{o}{=}\PY{+w}{ }\PY{k}{new}\PY{+w}{ }\PY{n}{HashMap}\PY{o}{\PYZlt{}}\PY{o}{\PYZgt{}}\PY{p}{(}\PY{p}{)}\PY{p}{;}
\PY{n}{map}\PY{p}{.}\PY{n+na}{put}\PY{p}{(}\PY{k}{new}\PY{+w}{ }\PY{n}{BadKey}\PY{p}{(}\PY{l+m+mi}{1}\PY{p}{)}\PY{p}{,}\PY{+w}{ }\PY{l+s}{\PYZdq{}}\PY{l+s}{a}\PY{l+s}{\PYZdq{}}\PY{p}{)}\PY{p}{;}
\PY{n}{System}\PY{p}{.}\PY{n+na}{out}\PY{p}{.}\PY{n+na}{println}\PY{p}{(}\PY{n}{map}\PY{p}{.}\PY{n+na}{get}\PY{p}{(}\PY{k}{new}\PY{+w}{ }\PY{n}{BadKey}\PY{p}{(}\PY{l+m+mi}{1}\PY{p}{)}\PY{p}{)}\PY{p}{)}\PY{p}{;}\PY{+w}{ }\PY{c+c1}{// null! different object, same \PYZdq{}logical\PYZdq{} key}
\end{Verbatim}
\end{codebox}

\section[LinkedHashMap]{LinkedHashMap}
\label{h:7:linkedhashmap}
\textbf{How it works internally}: extends \code{HashMap} and adds a doubly-linked list threading through all entries, to remember order. Every \code{Map.\allowbreak{}Entry} also has \code{before}/\code{after} pointers.

\begin{codebox}
\begin{Verbatim}[commandchars=\\\{\}]
\PY{n}{Map}\PY{o}{\PYZlt{}}\PY{n}{String}\PY{p}{,}\PY{+w}{ }\PY{n}{Integer}\PY{o}{\PYZgt{}}\PY{+w}{ }\PY{n}{lru}\PY{+w}{ }\PY{o}{=}\PY{+w}{ }\PY{k}{new}\PY{+w}{ }\PY{n}{LinkedHashMap}\PY{o}{\PYZlt{}}\PY{o}{\PYZgt{}}\PY{p}{(}\PY{l+m+mi}{16}\PY{p}{,}\PY{+w}{ }\PY{l+m+mf}{0.75f}\PY{p}{,}\PY{+w}{ }\PY{k+kc}{true}\PY{p}{)}\PY{p}{;}\PY{+w}{ }\PY{c+c1}{// accessOrder = true}
\PY{n}{lru}\PY{p}{.}\PY{n+na}{put}\PY{p}{(}\PY{l+s}{\PYZdq{}}\PY{l+s}{a}\PY{l+s}{\PYZdq{}}\PY{p}{,}\PY{+w}{ }\PY{l+m+mi}{1}\PY{p}{)}\PY{p}{;}
\PY{n}{lru}\PY{p}{.}\PY{n+na}{put}\PY{p}{(}\PY{l+s}{\PYZdq{}}\PY{l+s}{b}\PY{l+s}{\PYZdq{}}\PY{p}{,}\PY{+w}{ }\PY{l+m+mi}{2}\PY{p}{)}\PY{p}{;}
\PY{n}{lru}\PY{p}{.}\PY{n+na}{put}\PY{p}{(}\PY{l+s}{\PYZdq{}}\PY{l+s}{c}\PY{l+s}{\PYZdq{}}\PY{p}{,}\PY{+w}{ }\PY{l+m+mi}{3}\PY{p}{)}\PY{p}{;}
\PY{n}{lru}\PY{p}{.}\PY{n+na}{get}\PY{p}{(}\PY{l+s}{\PYZdq{}}\PY{l+s}{a}\PY{l+s}{\PYZdq{}}\PY{p}{)}\PY{p}{;}\PY{+w}{ }\PY{c+c1}{// touches \PYZdq{}a\PYZdq{}, moves it to the end in access\PYZhy{}order mode}
\PY{n}{System}\PY{p}{.}\PY{n+na}{out}\PY{p}{.}\PY{n+na}{println}\PY{p}{(}\PY{n}{lru}\PY{p}{.}\PY{n+na}{keySet}\PY{p}{(}\PY{p}{)}\PY{p}{)}\PY{p}{;}\PY{+w}{ }\PY{c+c1}{// [b, c, a]}
\end{Verbatim}
\end{codebox}

Because it tracks access order, \code{LinkedHashMap} is the standard building block for a simple \textbf{LRU cache}:

\begin{codebox}
\begin{Verbatim}[commandchars=\\\{\}]
\PY{k+kd}{class} \PY{n+nc}{LruCache}\PY{o}{\PYZlt{}}\PY{n}{K}\PY{p}{,}\PY{+w}{ }\PY{n}{V}\PY{o}{\PYZgt{}}\PY{+w}{ }\PY{k+kd}{extends}\PY{+w}{ }\PY{n}{LinkedHashMap}\PY{o}{\PYZlt{}}\PY{n}{K}\PY{p}{,}\PY{+w}{ }\PY{n}{V}\PY{o}{\PYZgt{}}\PY{+w}{ }\PY{p}{\PYZob{}}
\PY{+w}{    }\PY{k+kd}{private}\PY{+w}{ }\PY{k+kd}{final}\PY{+w}{ }\PY{k+kt}{int}\PY{+w}{ }\PY{n}{maxSize}\PY{p}{;}
\PY{+w}{    }\PY{n}{LruCache}\PY{p}{(}\PY{k+kt}{int}\PY{+w}{ }\PY{n}{maxSize}\PY{p}{)}\PY{+w}{ }\PY{p}{\PYZob{}}
\PY{+w}{        }\PY{k+kd}{super}\PY{p}{(}\PY{l+m+mi}{16}\PY{p}{,}\PY{+w}{ }\PY{l+m+mf}{0.75f}\PY{p}{,}\PY{+w}{ }\PY{k+kc}{true}\PY{p}{)}\PY{p}{;}\PY{+w}{ }\PY{c+c1}{// true = access\PYZhy{}order}
\PY{+w}{        }\PY{k}{this}\PY{p}{.}\PY{n+na}{maxSize}\PY{+w}{ }\PY{o}{=}\PY{+w}{ }\PY{n}{maxSize}\PY{p}{;}
\PY{+w}{    }\PY{p}{\PYZcb{}}
\PY{+w}{    }\PY{n+nd}{@Override}
\PY{+w}{    }\PY{k+kd}{protected}\PY{+w}{ }\PY{k+kt}{boolean}\PY{+w}{ }\PY{n+nf}{removeEldestEntry}\PY{p}{(}\PY{n}{Map}\PY{p}{.}\PY{n+na}{Entry}\PY{o}{\PYZlt{}}\PY{n}{K}\PY{p}{,}\PY{+w}{ }\PY{n}{V}\PY{o}{\PYZgt{}}\PY{+w}{ }\PY{n}{eldest}\PY{p}{)}\PY{+w}{ }\PY{p}{\PYZob{}}
\PY{+w}{        }\PY{k}{return}\PY{+w}{ }\PY{n}{size}\PY{p}{(}\PY{p}{)}\PY{+w}{ }\PY{o}{\PYZgt{}}\PY{+w}{ }\PY{n}{maxSize}\PY{p}{;}\PY{+w}{ }\PY{c+c1}{// evict the oldest entry once full}
\PY{+w}{    }\PY{p}{\PYZcb{}}
\PY{p}{\PYZcb{}}

\PY{n}{LruCache}\PY{o}{\PYZlt{}}\PY{n}{String}\PY{p}{,}\PY{+w}{ }\PY{n}{String}\PY{o}{\PYZgt{}}\PY{+w}{ }\PY{n}{cache}\PY{+w}{ }\PY{o}{=}\PY{+w}{ }\PY{k}{new}\PY{+w}{ }\PY{n}{LruCache}\PY{o}{\PYZlt{}}\PY{o}{\PYZgt{}}\PY{p}{(}\PY{l+m+mi}{2}\PY{p}{)}\PY{p}{;}
\PY{n}{cache}\PY{p}{.}\PY{n+na}{put}\PY{p}{(}\PY{l+s}{\PYZdq{}}\PY{l+s}{a}\PY{l+s}{\PYZdq{}}\PY{p}{,}\PY{+w}{ }\PY{l+s}{\PYZdq{}}\PY{l+s}{1}\PY{l+s}{\PYZdq{}}\PY{p}{)}\PY{p}{;}
\PY{n}{cache}\PY{p}{.}\PY{n+na}{put}\PY{p}{(}\PY{l+s}{\PYZdq{}}\PY{l+s}{b}\PY{l+s}{\PYZdq{}}\PY{p}{,}\PY{+w}{ }\PY{l+s}{\PYZdq{}}\PY{l+s}{2}\PY{l+s}{\PYZdq{}}\PY{p}{)}\PY{p}{;}
\PY{n}{cache}\PY{p}{.}\PY{n+na}{get}\PY{p}{(}\PY{l+s}{\PYZdq{}}\PY{l+s}{a}\PY{l+s}{\PYZdq{}}\PY{p}{)}\PY{p}{;}\PY{+w}{      }\PY{c+c1}{// \PYZdq{}a\PYZdq{} becomes most\PYZhy{}recently\PYZhy{}used}
\PY{n}{cache}\PY{p}{.}\PY{n+na}{put}\PY{p}{(}\PY{l+s}{\PYZdq{}}\PY{l+s}{c}\PY{l+s}{\PYZdq{}}\PY{p}{,}\PY{+w}{ }\PY{l+s}{\PYZdq{}}\PY{l+s}{3}\PY{l+s}{\PYZdq{}}\PY{p}{)}\PY{p}{;}\PY{+w}{ }\PY{c+c1}{// evicts \PYZdq{}b\PYZdq{} (least recently used)}
\PY{n}{System}\PY{p}{.}\PY{n+na}{out}\PY{p}{.}\PY{n+na}{println}\PY{p}{(}\PY{n}{cache}\PY{p}{.}\PY{n+na}{keySet}\PY{p}{(}\PY{p}{)}\PY{p}{)}\PY{p}{;}\PY{+w}{ }\PY{c+c1}{// [a, c]}
\end{Verbatim}
\end{codebox}

\textbf{Big-O}: same as \code{HashMap} — O(1) average for \code{get}/\code{put}/\code{remove} — plus a small constant overhead to maintain the linked list.

\textbf{Ordering}: insertion order by default, or access order if constructed with \code{accessOrder~\allowbreak{}=\allowbreak{}~\allowbreak{}true}.

\textbf{Null handling}: same as \code{HashMap} (one null key, many null values).

\textbf{Thread safety}: not thread-safe.

\textbf{When to choose it}: when you need predictable, repeatable iteration order (e.g., for deterministic tests, or to preserve the order fields were added for serialization), or to build an LRU cache.

\section[TreeMap]{TreeMap}
\label{h:7:treemap}
\textbf{How it works internally}: a \textbf{red-black tree} (a self-balancing binary search tree), keyed by \code{Comparable} order or a supplied \code{Comparator}. Every insert/lookup walks down the tree comparing keys.

\begin{codebox}
\begin{Verbatim}[commandchars=\\\{\}]
\PY{n}{TreeMap}\PY{o}{\PYZlt{}}\PY{n}{String}\PY{p}{,}\PY{+w}{ }\PY{n}{Integer}\PY{o}{\PYZgt{}}\PY{+w}{ }\PY{n}{sorted}\PY{+w}{ }\PY{o}{=}\PY{+w}{ }\PY{k}{new}\PY{+w}{ }\PY{n}{TreeMap}\PY{o}{\PYZlt{}}\PY{o}{\PYZgt{}}\PY{p}{(}\PY{p}{)}\PY{p}{;}
\PY{n}{sorted}\PY{p}{.}\PY{n+na}{put}\PY{p}{(}\PY{l+s}{\PYZdq{}}\PY{l+s}{banana}\PY{l+s}{\PYZdq{}}\PY{p}{,}\PY{+w}{ }\PY{l+m+mi}{2}\PY{p}{)}\PY{p}{;}
\PY{n}{sorted}\PY{p}{.}\PY{n+na}{put}\PY{p}{(}\PY{l+s}{\PYZdq{}}\PY{l+s}{apple}\PY{l+s}{\PYZdq{}}\PY{p}{,}\PY{+w}{ }\PY{l+m+mi}{1}\PY{p}{)}\PY{p}{;}
\PY{n}{sorted}\PY{p}{.}\PY{n+na}{put}\PY{p}{(}\PY{l+s}{\PYZdq{}}\PY{l+s}{cherry}\PY{l+s}{\PYZdq{}}\PY{p}{,}\PY{+w}{ }\PY{l+m+mi}{3}\PY{p}{)}\PY{p}{;}
\PY{n}{System}\PY{p}{.}\PY{n+na}{out}\PY{p}{.}\PY{n+na}{println}\PY{p}{(}\PY{n}{sorted}\PY{p}{)}\PY{p}{;}\PY{+w}{ }\PY{c+c1}{// \PYZob{}apple=1, banana=2, cherry=3\PYZcb{} \PYZhy{}\PYZhy{} sorted by key}

\PY{n}{System}\PY{p}{.}\PY{n+na}{out}\PY{p}{.}\PY{n+na}{println}\PY{p}{(}\PY{n}{sorted}\PY{p}{.}\PY{n+na}{firstKey}\PY{p}{(}\PY{p}{)}\PY{p}{)}\PY{p}{;}\PY{+w}{        }\PY{c+c1}{// apple}
\PY{n}{System}\PY{p}{.}\PY{n+na}{out}\PY{p}{.}\PY{n+na}{println}\PY{p}{(}\PY{n}{sorted}\PY{p}{.}\PY{n+na}{lastKey}\PY{p}{(}\PY{p}{)}\PY{p}{)}\PY{p}{;}\PY{+w}{         }\PY{c+c1}{// cherry}
\PY{n}{System}\PY{p}{.}\PY{n+na}{out}\PY{p}{.}\PY{n+na}{println}\PY{p}{(}\PY{n}{sorted}\PY{p}{.}\PY{n+na}{ceilingKey}\PY{p}{(}\PY{l+s}{\PYZdq{}}\PY{l+s}{b}\PY{l+s}{\PYZdq{}}\PY{p}{)}\PY{p}{)}\PY{p}{;}\PY{+w}{   }\PY{c+c1}{// banana (smallest key \PYZgt{}= \PYZdq{}b\PYZdq{})}
\PY{n}{System}\PY{p}{.}\PY{n+na}{out}\PY{p}{.}\PY{n+na}{println}\PY{p}{(}\PY{n}{sorted}\PY{p}{.}\PY{n+na}{floorKey}\PY{p}{(}\PY{l+s}{\PYZdq{}}\PY{l+s}{b}\PY{l+s}{\PYZdq{}}\PY{p}{)}\PY{p}{)}\PY{p}{;}\PY{+w}{     }\PY{c+c1}{// apple  (largest key \PYZlt{}= \PYZdq{}b\PYZdq{})}
\PY{n}{System}\PY{p}{.}\PY{n+na}{out}\PY{p}{.}\PY{n+na}{println}\PY{p}{(}\PY{n}{sorted}\PY{p}{.}\PY{n+na}{headMap}\PY{p}{(}\PY{l+s}{\PYZdq{}}\PY{l+s}{banana}\PY{l+s}{\PYZdq{}}\PY{p}{)}\PY{p}{)}\PY{p}{;}\PY{+w}{ }\PY{c+c1}{// \PYZob{}apple=1\PYZcb{} \PYZhy{}\PYZhy{} keys strictly less than \PYZdq{}banana\PYZdq{}}
\end{Verbatim}
\end{codebox}

\textbf{Big-O}: \code{get}/\code{put}/\code{remove}/\code{containsKey}: O(log n). Range operations (\code{headMap}, \code{tailMap}, \code{subMap}) are O(log n) to locate the boundary, and the returned view is backed by the original map.

\textbf{Ordering}: sorted by natural ordering of keys, or by the \code{Comparator} given at construction time.

\textbf{Null handling}: \textbf{no null keys} allowed (throws \code{NullPointerException} — comparisons need a real key). Null values are allowed if the comparator doesn’t touch values.

\textbf{Thread safety}: not thread-safe.

\textbf{When to choose it}: when you need keys sorted, or range queries (\code{floorKey}, \code{ceilingKey}, \code{subMap}), or the smallest/largest key on demand.

\begin{codebox}
\begin{Verbatim}[commandchars=\\\{\}]
\PY{c+c1}{// Custom key type needs a Comparator, or must implement Comparable}
\PY{n}{TreeMap}\PY{o}{\PYZlt{}}\PY{n}{String}\PY{p}{,}\PY{+w}{ }\PY{n}{Integer}\PY{o}{\PYZgt{}}\PY{+w}{ }\PY{n}{byLength}\PY{+w}{ }\PY{o}{=}\PY{+w}{ }\PY{k}{new}\PY{+w}{ }\PY{n}{TreeMap}\PY{o}{\PYZlt{}}\PY{o}{\PYZgt{}}\PY{p}{(}\PY{n}{Comparator}\PY{p}{.}\PY{n+na}{comparingInt}\PY{p}{(}\PY{n}{String}\PY{p}{:}\PY{p}{:}\PY{n}{length}\PY{p}{)}\PY{p}{)}\PY{p}{;}
\PY{n}{byLength}\PY{p}{.}\PY{n+na}{put}\PY{p}{(}\PY{l+s}{\PYZdq{}}\PY{l+s}{kiwi}\PY{l+s}{\PYZdq{}}\PY{p}{,}\PY{+w}{ }\PY{l+m+mi}{1}\PY{p}{)}\PY{p}{;}
\PY{n}{byLength}\PY{p}{.}\PY{n+na}{put}\PY{p}{(}\PY{l+s}{\PYZdq{}}\PY{l+s}{fig}\PY{l+s}{\PYZdq{}}\PY{p}{,}\PY{+w}{ }\PY{l+m+mi}{2}\PY{p}{)}\PY{p}{;}
\PY{n}{System}\PY{p}{.}\PY{n+na}{out}\PY{p}{.}\PY{n+na}{println}\PY{p}{(}\PY{n}{byLength}\PY{p}{)}\PY{p}{;}\PY{+w}{ }\PY{c+c1}{// \PYZob{}fig=2, kiwi=1\PYZcb{} \PYZhy{}\PYZhy{} sorted by string length}
\end{Verbatim}
\end{codebox}

\section[HashSet]{HashSet}
\label{h:7:hashset}
\textbf{How it works internally}: backed internally by a \code{HashMap\textless{}\allowbreak{}E,\allowbreak{}~\allowbreak{}Object\textgreater{}} where every element is stored as a key, and all values point to a shared dummy constant object. So everything said about \code{HashMap} buckets, hashing, treeification, and resizing applies here too.

\begin{codebox}
\begin{Verbatim}[commandchars=\\\{\}]
\PY{n}{Set}\PY{o}{\PYZlt{}}\PY{n}{String}\PY{o}{\PYZgt{}}\PY{+w}{ }\PY{n}{set}\PY{+w}{ }\PY{o}{=}\PY{+w}{ }\PY{k}{new}\PY{+w}{ }\PY{n}{HashSet}\PY{o}{\PYZlt{}}\PY{o}{\PYZgt{}}\PY{p}{(}\PY{p}{)}\PY{p}{;}
\PY{n}{set}\PY{p}{.}\PY{n+na}{add}\PY{p}{(}\PY{l+s}{\PYZdq{}}\PY{l+s}{x}\PY{l+s}{\PYZdq{}}\PY{p}{)}\PY{p}{;}
\PY{n}{set}\PY{p}{.}\PY{n+na}{add}\PY{p}{(}\PY{l+s}{\PYZdq{}}\PY{l+s}{y}\PY{l+s}{\PYZdq{}}\PY{p}{)}\PY{p}{;}
\PY{n}{set}\PY{p}{.}\PY{n+na}{add}\PY{p}{(}\PY{l+s}{\PYZdq{}}\PY{l+s}{x}\PY{l+s}{\PYZdq{}}\PY{p}{)}\PY{p}{;}\PY{+w}{ }\PY{c+c1}{// ignored \PYZhy{}\PYZhy{} already present}
\PY{n}{System}\PY{p}{.}\PY{n+na}{out}\PY{p}{.}\PY{n+na}{println}\PY{p}{(}\PY{n}{set}\PY{p}{.}\PY{n+na}{size}\PY{p}{(}\PY{p}{)}\PY{p}{)}\PY{p}{;}\PY{+w}{ }\PY{c+c1}{// 2}
\PY{n}{System}\PY{p}{.}\PY{n+na}{out}\PY{p}{.}\PY{n+na}{println}\PY{p}{(}\PY{n}{set}\PY{p}{)}\PY{p}{;}\PY{+w}{        }\PY{c+c1}{// order unspecified, e.g. [x, y] or [y, x]}
\end{Verbatim}
\end{codebox}

\textbf{Big-O}: \code{add}/\code{remove}/\code{contains}: O(1) average, same worst-case caveats as \code{HashMap}.

\textbf{Ordering}: unspecified.

\textbf{Null handling}: allows one \code{null} element.

\textbf{Thread safety}: not thread-safe.

\textbf{When to choose it}: default \code{Set} choice, when you need fast membership tests and don’t care about order.

\section[LinkedHashSet]{LinkedHashSet}
\label{h:7:linkedhashset}
\textbf{How it works internally}: backed by a \code{LinkedHashMap} internally, so it has the same linked-list-through-buckets trick to preserve insertion order.

\begin{codebox}
\begin{Verbatim}[commandchars=\\\{\}]
\PY{n}{Set}\PY{o}{\PYZlt{}}\PY{n}{String}\PY{o}{\PYZgt{}}\PY{+w}{ }\PY{n}{visited}\PY{+w}{ }\PY{o}{=}\PY{+w}{ }\PY{k}{new}\PY{+w}{ }\PY{n}{LinkedHashSet}\PY{o}{\PYZlt{}}\PY{o}{\PYZgt{}}\PY{p}{(}\PY{p}{)}\PY{p}{;}
\PY{n}{visited}\PY{p}{.}\PY{n+na}{add}\PY{p}{(}\PY{l+s}{\PYZdq{}}\PY{l+s}{home}\PY{l+s}{\PYZdq{}}\PY{p}{)}\PY{p}{;}
\PY{n}{visited}\PY{p}{.}\PY{n+na}{add}\PY{p}{(}\PY{l+s}{\PYZdq{}}\PY{l+s}{cart}\PY{l+s}{\PYZdq{}}\PY{p}{)}\PY{p}{;}
\PY{n}{visited}\PY{p}{.}\PY{n+na}{add}\PY{p}{(}\PY{l+s}{\PYZdq{}}\PY{l+s}{home}\PY{l+s}{\PYZdq{}}\PY{p}{)}\PY{p}{;}\PY{+w}{ }\PY{c+c1}{// ignored, but original position is kept}
\PY{n}{visited}\PY{p}{.}\PY{n+na}{add}\PY{p}{(}\PY{l+s}{\PYZdq{}}\PY{l+s}{checkout}\PY{l+s}{\PYZdq{}}\PY{p}{)}\PY{p}{;}
\PY{n}{System}\PY{p}{.}\PY{n+na}{out}\PY{p}{.}\PY{n+na}{println}\PY{p}{(}\PY{n}{visited}\PY{p}{)}\PY{p}{;}\PY{+w}{ }\PY{c+c1}{// [home, cart, checkout]}
\end{Verbatim}
\end{codebox}

\textbf{Big-O}: same as \code{HashSet} — O(1) average.

\textbf{Ordering}: insertion order, preserved.

\textbf{Null handling}: one \code{null} allowed.

\textbf{Thread safety}: not thread-safe.

\textbf{When to choose it}: when you want set semantics (no duplicates) but also need to remember the order elements were added — e.g., de-duplicating a list while preserving first-seen order.

\begin{codebox}
\begin{Verbatim}[commandchars=\\\{\}]
\PY{n}{List}\PY{o}{\PYZlt{}}\PY{n}{String}\PY{o}{\PYZgt{}}\PY{+w}{ }\PY{n}{withDupes}\PY{+w}{ }\PY{o}{=}\PY{+w}{ }\PY{n}{List}\PY{p}{.}\PY{n+na}{of}\PY{p}{(}\PY{l+s}{\PYZdq{}}\PY{l+s}{b}\PY{l+s}{\PYZdq{}}\PY{p}{,}\PY{+w}{ }\PY{l+s}{\PYZdq{}}\PY{l+s}{a}\PY{l+s}{\PYZdq{}}\PY{p}{,}\PY{+w}{ }\PY{l+s}{\PYZdq{}}\PY{l+s}{b}\PY{l+s}{\PYZdq{}}\PY{p}{,}\PY{+w}{ }\PY{l+s}{\PYZdq{}}\PY{l+s}{c}\PY{l+s}{\PYZdq{}}\PY{p}{,}\PY{+w}{ }\PY{l+s}{\PYZdq{}}\PY{l+s}{a}\PY{l+s}{\PYZdq{}}\PY{p}{)}\PY{p}{;}
\PY{n}{Set}\PY{o}{\PYZlt{}}\PY{n}{String}\PY{o}{\PYZgt{}}\PY{+w}{ }\PY{n}{deduped}\PY{+w}{ }\PY{o}{=}\PY{+w}{ }\PY{k}{new}\PY{+w}{ }\PY{n}{LinkedHashSet}\PY{o}{\PYZlt{}}\PY{o}{\PYZgt{}}\PY{p}{(}\PY{n}{withDupes}\PY{p}{)}\PY{p}{;}
\PY{n}{System}\PY{p}{.}\PY{n+na}{out}\PY{p}{.}\PY{n+na}{println}\PY{p}{(}\PY{n}{deduped}\PY{p}{)}\PY{p}{;}\PY{+w}{ }\PY{c+c1}{// [b, a, c] \PYZhy{}\PYZhy{} de\PYZhy{}duplicated, first\PYZhy{}seen order kept}
\end{Verbatim}
\end{codebox}

\section[TreeSet]{TreeSet}
\label{h:7:treeset}
\textbf{How it works internally}: backed by a \code{TreeMap\textless{}\allowbreak{}E,\allowbreak{}~\allowbreak{}Object\textgreater{}}, so it’s a red-black tree under the hood.

\begin{codebox}
\begin{Verbatim}[commandchars=\\\{\}]
\PY{n}{TreeSet}\PY{o}{\PYZlt{}}\PY{n}{Integer}\PY{o}{\PYZgt{}}\PY{+w}{ }\PY{n}{scores}\PY{+w}{ }\PY{o}{=}\PY{+w}{ }\PY{k}{new}\PY{+w}{ }\PY{n}{TreeSet}\PY{o}{\PYZlt{}}\PY{o}{\PYZgt{}}\PY{p}{(}\PY{n}{List}\PY{p}{.}\PY{n+na}{of}\PY{p}{(}\PY{l+m+mi}{42}\PY{p}{,}\PY{+w}{ }\PY{l+m+mi}{7}\PY{p}{,}\PY{+w}{ }\PY{l+m+mi}{15}\PY{p}{,}\PY{+w}{ }\PY{l+m+mi}{4}\PY{p}{)}\PY{p}{)}\PY{p}{;}
\PY{n}{System}\PY{p}{.}\PY{n+na}{out}\PY{p}{.}\PY{n+na}{println}\PY{p}{(}\PY{n}{scores}\PY{p}{)}\PY{p}{;}\PY{+w}{          }\PY{c+c1}{// [4, 7, 15, 42] \PYZhy{}\PYZhy{} sorted}
\PY{n}{System}\PY{p}{.}\PY{n+na}{out}\PY{p}{.}\PY{n+na}{println}\PY{p}{(}\PY{n}{scores}\PY{p}{.}\PY{n+na}{first}\PY{p}{(}\PY{p}{)}\PY{p}{)}\PY{p}{;}\PY{+w}{  }\PY{c+c1}{// 4}
\PY{n}{System}\PY{p}{.}\PY{n+na}{out}\PY{p}{.}\PY{n+na}{println}\PY{p}{(}\PY{n}{scores}\PY{p}{.}\PY{n+na}{last}\PY{p}{(}\PY{p}{)}\PY{p}{)}\PY{p}{;}\PY{+w}{   }\PY{c+c1}{// 42}
\PY{n}{System}\PY{p}{.}\PY{n+na}{out}\PY{p}{.}\PY{n+na}{println}\PY{p}{(}\PY{n}{scores}\PY{p}{.}\PY{n+na}{higher}\PY{p}{(}\PY{l+m+mi}{15}\PY{p}{)}\PY{p}{)}\PY{p}{;}\PY{+w}{ }\PY{c+c1}{// 42 (smallest element strictly greater than 15)}
\PY{n}{System}\PY{p}{.}\PY{n+na}{out}\PY{p}{.}\PY{n+na}{println}\PY{p}{(}\PY{n}{scores}\PY{p}{.}\PY{n+na}{lower}\PY{p}{(}\PY{l+m+mi}{15}\PY{p}{)}\PY{p}{)}\PY{p}{;}\PY{+w}{  }\PY{c+c1}{// 7  (largest element strictly less than 15)}
\end{Verbatim}
\end{codebox}

\textbf{Big-O}: \code{add}/\code{remove}/\code{contains}: O(log n).

\textbf{Ordering}: sorted, by natural order or supplied \code{Comparator}.

\textbf{Null handling}: \textbf{no nulls} (comparisons would throw \code{NullPointerException}).

\textbf{Thread safety}: not thread-safe.

\textbf{When to choose it}: when you need a de-duplicated, always-sorted collection, or navigation methods like \code{first()}, \code{last()}, \code{higher()}, \code{lower()}, \code{subSet()}.

\section[PriorityQueue]{PriorityQueue}
\label{h:7:priorityqueue}
\textbf{How it works internally}: a \textbf{binary heap} stored in an array. The heap property guarantees the smallest element (by natural order or comparator) is always at the root/head, so \code{peek()}/\code{poll()} are cheap. It is \textbf{not} a full sort — only the head is guaranteed to be the minimum; the rest of the internal array is \emph{not} in sorted order.

\begin{codebox}
\begin{Verbatim}[commandchars=\\\{\}]
\PY{n}{PriorityQueue}\PY{o}{\PYZlt{}}\PY{n}{Integer}\PY{o}{\PYZgt{}}\PY{+w}{ }\PY{n}{pq}\PY{+w}{ }\PY{o}{=}\PY{+w}{ }\PY{k}{new}\PY{+w}{ }\PY{n}{PriorityQueue}\PY{o}{\PYZlt{}}\PY{o}{\PYZgt{}}\PY{p}{(}\PY{p}{)}\PY{p}{;}
\PY{n}{pq}\PY{p}{.}\PY{n+na}{offer}\PY{p}{(}\PY{l+m+mi}{5}\PY{p}{)}\PY{p}{;}
\PY{n}{pq}\PY{p}{.}\PY{n+na}{offer}\PY{p}{(}\PY{l+m+mi}{1}\PY{p}{)}\PY{p}{;}
\PY{n}{pq}\PY{p}{.}\PY{n+na}{offer}\PY{p}{(}\PY{l+m+mi}{3}\PY{p}{)}\PY{p}{;}
\PY{n}{System}\PY{p}{.}\PY{n+na}{out}\PY{p}{.}\PY{n+na}{println}\PY{p}{(}\PY{n}{pq}\PY{p}{.}\PY{n+na}{poll}\PY{p}{(}\PY{p}{)}\PY{p}{)}\PY{p}{;}\PY{+w}{ }\PY{c+c1}{// 1 (smallest first)}
\PY{n}{System}\PY{p}{.}\PY{n+na}{out}\PY{p}{.}\PY{n+na}{println}\PY{p}{(}\PY{n}{pq}\PY{p}{.}\PY{n+na}{poll}\PY{p}{(}\PY{p}{)}\PY{p}{)}\PY{p}{;}\PY{+w}{ }\PY{c+c1}{// 3}
\PY{n}{System}\PY{p}{.}\PY{n+na}{out}\PY{p}{.}\PY{n+na}{println}\PY{p}{(}\PY{n}{pq}\PY{p}{.}\PY{n+na}{poll}\PY{p}{(}\PY{p}{)}\PY{p}{)}\PY{p}{;}\PY{+w}{ }\PY{c+c1}{// 5}

\PY{c+c1}{// Max\PYZhy{}heap using a reverse comparator}
\PY{n}{PriorityQueue}\PY{o}{\PYZlt{}}\PY{n}{Integer}\PY{o}{\PYZgt{}}\PY{+w}{ }\PY{n}{maxHeap}\PY{+w}{ }\PY{o}{=}\PY{+w}{ }\PY{k}{new}\PY{+w}{ }\PY{n}{PriorityQueue}\PY{o}{\PYZlt{}}\PY{o}{\PYZgt{}}\PY{p}{(}\PY{n}{Comparator}\PY{p}{.}\PY{n+na}{reverseOrder}\PY{p}{(}\PY{p}{)}\PY{p}{)}\PY{p}{;}
\PY{n}{maxHeap}\PY{p}{.}\PY{n+na}{offer}\PY{p}{(}\PY{l+m+mi}{5}\PY{p}{)}\PY{p}{;}
\PY{n}{maxHeap}\PY{p}{.}\PY{n+na}{offer}\PY{p}{(}\PY{l+m+mi}{1}\PY{p}{)}\PY{p}{;}
\PY{n}{maxHeap}\PY{p}{.}\PY{n+na}{offer}\PY{p}{(}\PY{l+m+mi}{3}\PY{p}{)}\PY{p}{;}
\PY{n}{System}\PY{p}{.}\PY{n+na}{out}\PY{p}{.}\PY{n+na}{println}\PY{p}{(}\PY{n}{maxHeap}\PY{p}{.}\PY{n+na}{poll}\PY{p}{(}\PY{p}{)}\PY{p}{)}\PY{p}{;}\PY{+w}{ }\PY{c+c1}{// 5 (largest first)}
\end{Verbatim}
\end{codebox}

\begin{codebox}
\begin{Verbatim}[commandchars=\\\{\}]
\PY{c+c1}{// Common use case: \PYZdq{}top\PYZhy{}k\PYZdq{} or scheduling by priority}
\PY{k+kd}{record} \PY{n+nc}{Task}\PY{p}{(}\PY{n}{String}\PY{+w}{ }\PY{n}{name}\PY{p}{,}\PY{+w}{ }\PY{k+kt}{int}\PY{+w}{ }\PY{n}{priority}\PY{p}{)}\PY{+w}{ }\PY{p}{\PYZob{}}\PY{p}{\PYZcb{}}
\PY{n}{PriorityQueue}\PY{o}{\PYZlt{}}\PY{n}{Task}\PY{o}{\PYZgt{}}\PY{+w}{ }\PY{n}{tasks}\PY{+w}{ }\PY{o}{=}\PY{+w}{ }\PY{k}{new}\PY{+w}{ }\PY{n}{PriorityQueue}\PY{o}{\PYZlt{}}\PY{o}{\PYZgt{}}\PY{p}{(}\PY{n}{Comparator}\PY{p}{.}\PY{n+na}{comparingInt}\PY{p}{(}\PY{n}{Task}\PY{p}{:}\PY{p}{:}\PY{n}{priority}\PY{p}{)}\PY{p}{)}\PY{p}{;}
\PY{n}{tasks}\PY{p}{.}\PY{n+na}{offer}\PY{p}{(}\PY{k}{new}\PY{+w}{ }\PY{n}{Task}\PY{p}{(}\PY{l+s}{\PYZdq{}}\PY{l+s}{cleanup}\PY{l+s}{\PYZdq{}}\PY{p}{,}\PY{+w}{ }\PY{l+m+mi}{3}\PY{p}{)}\PY{p}{)}\PY{p}{;}
\PY{n}{tasks}\PY{p}{.}\PY{n+na}{offer}\PY{p}{(}\PY{k}{new}\PY{+w}{ }\PY{n}{Task}\PY{p}{(}\PY{l+s}{\PYZdq{}}\PY{l+s}{urgent\PYZhy{}fix}\PY{l+s}{\PYZdq{}}\PY{p}{,}\PY{+w}{ }\PY{l+m+mi}{1}\PY{p}{)}\PY{p}{)}\PY{p}{;}
\PY{n}{tasks}\PY{p}{.}\PY{n+na}{offer}\PY{p}{(}\PY{k}{new}\PY{+w}{ }\PY{n}{Task}\PY{p}{(}\PY{l+s}{\PYZdq{}}\PY{l+s}{report}\PY{l+s}{\PYZdq{}}\PY{p}{,}\PY{+w}{ }\PY{l+m+mi}{2}\PY{p}{)}\PY{p}{)}\PY{p}{;}
\PY{n}{System}\PY{p}{.}\PY{n+na}{out}\PY{p}{.}\PY{n+na}{println}\PY{p}{(}\PY{n}{tasks}\PY{p}{.}\PY{n+na}{poll}\PY{p}{(}\PY{p}{)}\PY{p}{.}\PY{n+na}{name}\PY{p}{(}\PY{p}{)}\PY{p}{)}\PY{p}{;}\PY{+w}{ }\PY{c+c1}{// urgent\PYZhy{}fix (lowest number = highest priority)}
\end{Verbatim}
\end{codebox}

\textbf{Big-O}: \code{offer}/\code{add} (insert): O(log n). \code{poll}/\code{remove} head: O(log n). \code{peek}: O(1). \code{contains} (arbitrary element): O(n).

\textbf{Ordering}: only the head element order is guaranteed; iterating the queue directly does \textbf{not} yield sorted order (use repeated \code{poll()} for that).

\textbf{Null handling}: \textbf{no nulls} — throws \code{NullPointerException} (nulls can’t be compared).

\textbf{Thread safety}: not thread-safe. For a concurrent priority queue, use \code{PriorityBlockingQueue}.

\textbf{When to choose it}: task scheduling, “always process the most urgent/smallest item next,” or top-k problems (e.g., “k largest elements in a stream”).

\section[ArrayDeque]{ArrayDeque}
\label{h:7:arraydeque}
\textbf{How it works internally}: a resizable circular array (a ring buffer) with head and tail indices. Unlike \code{ArrayList}, it can grow efficiently at both ends because it wraps around the array instead of always shifting from index 0.

\begin{codebox}
\begin{Verbatim}[commandchars=\\\{\}]
\PY{n}{Deque}\PY{o}{\PYZlt{}}\PY{n}{String}\PY{o}{\PYZgt{}}\PY{+w}{ }\PY{n}{deque}\PY{+w}{ }\PY{o}{=}\PY{+w}{ }\PY{k}{new}\PY{+w}{ }\PY{n}{ArrayDeque}\PY{o}{\PYZlt{}}\PY{o}{\PYZgt{}}\PY{p}{(}\PY{p}{)}\PY{p}{;}
\PY{n}{deque}\PY{p}{.}\PY{n+na}{addFirst}\PY{p}{(}\PY{l+s}{\PYZdq{}}\PY{l+s}{b}\PY{l+s}{\PYZdq{}}\PY{p}{)}\PY{p}{;}
\PY{n}{deque}\PY{p}{.}\PY{n+na}{addFirst}\PY{p}{(}\PY{l+s}{\PYZdq{}}\PY{l+s}{a}\PY{l+s}{\PYZdq{}}\PY{p}{)}\PY{p}{;}
\PY{n}{deque}\PY{p}{.}\PY{n+na}{addLast}\PY{p}{(}\PY{l+s}{\PYZdq{}}\PY{l+s}{c}\PY{l+s}{\PYZdq{}}\PY{p}{)}\PY{p}{;}
\PY{n}{System}\PY{p}{.}\PY{n+na}{out}\PY{p}{.}\PY{n+na}{println}\PY{p}{(}\PY{n}{deque}\PY{p}{)}\PY{p}{;}\PY{+w}{ }\PY{c+c1}{// [a, b, c]}

\PY{c+c1}{// as a Stack (LIFO)}
\PY{n}{Deque}\PY{o}{\PYZlt{}}\PY{n}{Integer}\PY{o}{\PYZgt{}}\PY{+w}{ }\PY{n}{stack}\PY{+w}{ }\PY{o}{=}\PY{+w}{ }\PY{k}{new}\PY{+w}{ }\PY{n}{ArrayDeque}\PY{o}{\PYZlt{}}\PY{o}{\PYZgt{}}\PY{p}{(}\PY{p}{)}\PY{p}{;}
\PY{n}{stack}\PY{p}{.}\PY{n+na}{push}\PY{p}{(}\PY{l+m+mi}{1}\PY{p}{)}\PY{p}{;}\PY{+w}{ }\PY{n}{stack}\PY{p}{.}\PY{n+na}{push}\PY{p}{(}\PY{l+m+mi}{2}\PY{p}{)}\PY{p}{;}\PY{+w}{ }\PY{n}{stack}\PY{p}{.}\PY{n+na}{push}\PY{p}{(}\PY{l+m+mi}{3}\PY{p}{)}\PY{p}{;}
\PY{n}{System}\PY{p}{.}\PY{n+na}{out}\PY{p}{.}\PY{n+na}{println}\PY{p}{(}\PY{n}{stack}\PY{p}{.}\PY{n+na}{pop}\PY{p}{(}\PY{p}{)}\PY{p}{)}\PY{p}{;}\PY{+w}{ }\PY{c+c1}{// 3}

\PY{c+c1}{// as a Queue (FIFO)}
\PY{n}{Deque}\PY{o}{\PYZlt{}}\PY{n}{Integer}\PY{o}{\PYZgt{}}\PY{+w}{ }\PY{n}{queue}\PY{+w}{ }\PY{o}{=}\PY{+w}{ }\PY{k}{new}\PY{+w}{ }\PY{n}{ArrayDeque}\PY{o}{\PYZlt{}}\PY{o}{\PYZgt{}}\PY{p}{(}\PY{p}{)}\PY{p}{;}
\PY{n}{queue}\PY{p}{.}\PY{n+na}{offer}\PY{p}{(}\PY{l+m+mi}{1}\PY{p}{)}\PY{p}{;}\PY{+w}{ }\PY{n}{queue}\PY{p}{.}\PY{n+na}{offer}\PY{p}{(}\PY{l+m+mi}{2}\PY{p}{)}\PY{p}{;}\PY{+w}{ }\PY{n}{queue}\PY{p}{.}\PY{n+na}{offer}\PY{p}{(}\PY{l+m+mi}{3}\PY{p}{)}\PY{p}{;}
\PY{n}{System}\PY{p}{.}\PY{n+na}{out}\PY{p}{.}\PY{n+na}{println}\PY{p}{(}\PY{n}{queue}\PY{p}{.}\PY{n+na}{poll}\PY{p}{(}\PY{p}{)}\PY{p}{)}\PY{p}{;}\PY{+w}{ }\PY{c+c1}{// 1}
\end{Verbatim}
\end{codebox}

\textbf{Big-O}: \code{addFirst}/\code{addLast}/\code{removeFirst}/\code{removeLast}/\code{peekFirst}/\code{peekLast}: O(1) amortized. \code{get(\allowbreak{}index)}-style access is not offered (it’s not a \code{List}); \code{contains}: O(n).

\textbf{Ordering}: insertion order (head to tail).

\textbf{Null handling}: \textbf{no nulls} — \code{null} is used internally as a sentinel to mean “empty,” so \code{add(\allowbreak{}null)} throws \code{NullPointerException}.

\textbf{Thread safety}: not thread-safe. For concurrent use, see \code{ConcurrentLinkedDeque} or \code{LinkedBlockingDeque}.

\textbf{When to choose it}: the default choice for both stacks and queues. Faster than \code{Stack} (legacy, synchronized) and faster than \code{LinkedList} for the same job, with less memory overhead.

\section[IdentityHashMap]{IdentityHashMap}
\label{h:7:identityhashmap}
\textbf{How it works internally}: looks like \code{HashMap}, but uses \textbf{reference equality} (\code{==}) instead of \code{.\allowbreak{}equals()}, and \code{System.\allowbreak{}identity\allowbreak{}Hash\allowbreak{}Code()} instead of \code{.\allowbreak{}hashCode()}, to compare and hash keys. Two distinct objects that are “equal” by \code{.\allowbreak{}equals()} are treated as \textbf{different} keys if they are not the same object in memory.

\begin{codebox}
\begin{Verbatim}[commandchars=\\\{\}]
\PY{n}{Map}\PY{o}{\PYZlt{}}\PY{n}{String}\PY{p}{,}\PY{+w}{ }\PY{n}{String}\PY{o}{\PYZgt{}}\PY{+w}{ }\PY{n}{identityMap}\PY{+w}{ }\PY{o}{=}\PY{+w}{ }\PY{k}{new}\PY{+w}{ }\PY{n}{IdentityHashMap}\PY{o}{\PYZlt{}}\PY{o}{\PYZgt{}}\PY{p}{(}\PY{p}{)}\PY{p}{;}
\PY{n}{String}\PY{+w}{ }\PY{n}{a}\PY{+w}{ }\PY{o}{=}\PY{+w}{ }\PY{k}{new}\PY{+w}{ }\PY{n}{String}\PY{p}{(}\PY{l+s}{\PYZdq{}}\PY{l+s}{key}\PY{l+s}{\PYZdq{}}\PY{p}{)}\PY{p}{;}
\PY{n}{String}\PY{+w}{ }\PY{n}{b}\PY{+w}{ }\PY{o}{=}\PY{+w}{ }\PY{k}{new}\PY{+w}{ }\PY{n}{String}\PY{p}{(}\PY{l+s}{\PYZdq{}}\PY{l+s}{key}\PY{l+s}{\PYZdq{}}\PY{p}{)}\PY{p}{;}\PY{+w}{ }\PY{c+c1}{// different object, same content}
\PY{n}{identityMap}\PY{p}{.}\PY{n+na}{put}\PY{p}{(}\PY{n}{a}\PY{p}{,}\PY{+w}{ }\PY{l+s}{\PYZdq{}}\PY{l+s}{first}\PY{l+s}{\PYZdq{}}\PY{p}{)}\PY{p}{;}
\PY{n}{identityMap}\PY{p}{.}\PY{n+na}{put}\PY{p}{(}\PY{n}{b}\PY{p}{,}\PY{+w}{ }\PY{l+s}{\PYZdq{}}\PY{l+s}{second}\PY{l+s}{\PYZdq{}}\PY{p}{)}\PY{p}{;}\PY{+w}{ }\PY{c+c1}{// treated as a DIFFERENT key}
\PY{n}{System}\PY{p}{.}\PY{n+na}{out}\PY{p}{.}\PY{n+na}{println}\PY{p}{(}\PY{n}{identityMap}\PY{p}{.}\PY{n+na}{size}\PY{p}{(}\PY{p}{)}\PY{p}{)}\PY{p}{;}\PY{+w}{ }\PY{c+c1}{// 2, not 1!}
\PY{n}{System}\PY{p}{.}\PY{n+na}{out}\PY{p}{.}\PY{n+na}{println}\PY{p}{(}\PY{n}{identityMap}\PY{p}{.}\PY{n+na}{get}\PY{p}{(}\PY{n}{a}\PY{p}{)}\PY{p}{)}\PY{p}{;}\PY{+w}{ }\PY{c+c1}{// first}
\PY{n}{System}\PY{p}{.}\PY{n+na}{out}\PY{p}{.}\PY{n+na}{println}\PY{p}{(}\PY{n}{identityMap}\PY{p}{.}\PY{n+na}{get}\PY{p}{(}\PY{k}{new}\PY{+w}{ }\PY{n}{String}\PY{p}{(}\PY{l+s}{\PYZdq{}}\PY{l+s}{key}\PY{l+s}{\PYZdq{}}\PY{p}{)}\PY{p}{)}\PY{p}{)}\PY{p}{;}\PY{+w}{ }\PY{c+c1}{// null \PYZhy{}\PYZhy{} not the same reference}
\end{Verbatim}
\end{codebox}

\textbf{Big-O}: O(1) average for \code{get}/\code{put} — same style of hash table as \code{HashMap}, just with identity hashing.

\textbf{Ordering}: unspecified.

\textbf{Null handling}: allows a \code{null} key and null values.

\textbf{Thread safety}: not thread-safe.

\textbf{When to choose it}: rare. Useful for object-graph traversal algorithms (e.g., serialization, deep-copy, cycle detection) where you need to track “have I visited this exact object instance before,” regardless of its \code{.\allowbreak{}equals()} implementation. The JDK itself uses it in \code{ObjectOutputStream}.

\section[WeakHashMap]{WeakHashMap}
\label{h:7:weakhashmap}
\textbf{How it works internally}: stores keys as \textbf{weak references}. If nothing else in the program holds a strong reference to a key, the garbage collector is free to reclaim it. Once a key is collected, its entry is removed from the map automatically (lazily, on the next access or via an internal cleanup on the next operation).

\begin{codebox}
\begin{Verbatim}[commandchars=\\\{\}]
\PY{n}{Map}\PY{o}{\PYZlt{}}\PY{n}{Object}\PY{p}{,}\PY{+w}{ }\PY{n}{String}\PY{o}{\PYZgt{}}\PY{+w}{ }\PY{n}{cache}\PY{+w}{ }\PY{o}{=}\PY{+w}{ }\PY{k}{new}\PY{+w}{ }\PY{n}{WeakHashMap}\PY{o}{\PYZlt{}}\PY{o}{\PYZgt{}}\PY{p}{(}\PY{p}{)}\PY{p}{;}
\PY{n}{Object}\PY{+w}{ }\PY{n}{key}\PY{+w}{ }\PY{o}{=}\PY{+w}{ }\PY{k}{new}\PY{+w}{ }\PY{n}{Object}\PY{p}{(}\PY{p}{)}\PY{p}{;}
\PY{n}{cache}\PY{p}{.}\PY{n+na}{put}\PY{p}{(}\PY{n}{key}\PY{p}{,}\PY{+w}{ }\PY{l+s}{\PYZdq{}}\PY{l+s}{metadata}\PY{l+s}{\PYZdq{}}\PY{p}{)}\PY{p}{;}
\PY{n}{System}\PY{p}{.}\PY{n+na}{out}\PY{p}{.}\PY{n+na}{println}\PY{p}{(}\PY{n}{cache}\PY{p}{.}\PY{n+na}{size}\PY{p}{(}\PY{p}{)}\PY{p}{)}\PY{p}{;}\PY{+w}{ }\PY{c+c1}{// 1}

\PY{n}{key}\PY{+w}{ }\PY{o}{=}\PY{+w}{ }\PY{k+kc}{null}\PY{p}{;}\PY{+w}{          }\PY{c+c1}{// drop the only strong reference}
\PY{n}{System}\PY{p}{.}\PY{n+na}{gc}\PY{p}{(}\PY{p}{)}\PY{p}{;}\PY{+w}{          }\PY{c+c1}{// suggest GC (not guaranteed, but likely in this simple example)}
\PY{c+c1}{// after GC runs, the entry may be gone:}
\PY{c+c1}{// System.out.println(cache.size()); // often 0, though timing isn\PYZsq{}t guaranteed}
\end{Verbatim}
\end{codebox}

\textbf{Big-O}: O(1) average for \code{get}/\code{put}, same as \code{HashMap}.

\textbf{Ordering}: unspecified.

\textbf{Null handling}: allows a null key/value (though a null key is a bit unusual conceptually).

\textbf{Thread safety}: not thread-safe.

\textbf{When to choose it}: caches or metadata maps keyed by objects you don’t want to keep alive artificially — e.g., attaching extra data to objects managed elsewhere, such as classloader-scoped caches. Not a general-purpose cache eviction tool; for size/time-based eviction, prefer a real caching library (Caffeine, Guava) or \code{LinkedHashMap}-based LRU.

\section[EnumMap]{EnumMap}
\label{h:7:enummap}
\textbf{How it works internally}: a specialized \code{Map} implementation for \textbf{enum keys only}, backed internally by a plain array indexed by the enum constant’s \code{ordinal()}. No hashing needed at all.

\begin{codebox}
\begin{Verbatim}[commandchars=\\\{\}]
\PY{k+kd}{enum}\PY{+w}{ }\PY{n}{Day}\PY{+w}{ }\PY{p}{\PYZob{}}\PY{+w}{ }\PY{n}{MON}\PY{p}{,}\PY{+w}{ }\PY{n}{TUE}\PY{p}{,}\PY{+w}{ }\PY{n}{WED}\PY{p}{,}\PY{+w}{ }\PY{n}{THU}\PY{p}{,}\PY{+w}{ }\PY{n}{FRI}\PY{p}{,}\PY{+w}{ }\PY{n}{SAT}\PY{p}{,}\PY{+w}{ }\PY{n}{SUN}\PY{+w}{ }\PY{p}{\PYZcb{}}

\PY{n}{EnumMap}\PY{o}{\PYZlt{}}\PY{n}{Day}\PY{p}{,}\PY{+w}{ }\PY{n}{String}\PY{o}{\PYZgt{}}\PY{+w}{ }\PY{n}{schedule}\PY{+w}{ }\PY{o}{=}\PY{+w}{ }\PY{k}{new}\PY{+w}{ }\PY{n}{EnumMap}\PY{o}{\PYZlt{}}\PY{o}{\PYZgt{}}\PY{p}{(}\PY{n}{Day}\PY{p}{.}\PY{n+na}{class}\PY{p}{)}\PY{p}{;}
\PY{n}{schedule}\PY{p}{.}\PY{n+na}{put}\PY{p}{(}\PY{n}{Day}\PY{p}{.}\PY{n+na}{MON}\PY{p}{,}\PY{+w}{ }\PY{l+s}{\PYZdq{}}\PY{l+s}{Standup}\PY{l+s}{\PYZdq{}}\PY{p}{)}\PY{p}{;}
\PY{n}{schedule}\PY{p}{.}\PY{n+na}{put}\PY{p}{(}\PY{n}{Day}\PY{p}{.}\PY{n+na}{FRI}\PY{p}{,}\PY{+w}{ }\PY{l+s}{\PYZdq{}}\PY{l+s}{Retro}\PY{l+s}{\PYZdq{}}\PY{p}{)}\PY{p}{;}
\PY{n}{System}\PY{p}{.}\PY{n+na}{out}\PY{p}{.}\PY{n+na}{println}\PY{p}{(}\PY{n}{schedule}\PY{p}{)}\PY{p}{;}\PY{+w}{ }\PY{c+c1}{// \PYZob{}MON=Standup, FRI=Retro\PYZcb{} \PYZhy{}\PYZhy{} always in enum declaration order}
\end{Verbatim}
\end{codebox}

\textbf{Big-O}: \code{get}/\code{put}/\code{remove}: O(1) — direct array indexing by ordinal, no hash collisions possible.

\textbf{Ordering}: always iterates in the natural order the enum constants are declared (i.e., ordinal order), regardless of insertion order.

\textbf{Null handling}: no null keys (throws \code{NullPointerException}); null values allowed.

\textbf{Thread safety}: not thread-safe.

\textbf{When to choose it}: whenever your map key is an enum. It is faster and more memory-efficient than \code{HashMap\textless{}\allowbreak{}EnumType,\allowbreak{}~\allowbreak{}V\textgreater{}}, and its guaranteed ordering is a nice bonus.

\section[EnumSet]{EnumSet}
\label{h:7:enumset}
\textbf{How it works internally}: a specialized \code{Set} for enum values, internally represented as a \textbf{bitmask} (a \code{long}, or an array of \code{long}s for enums with more than 64 constants). Membership tests and unions/intersections become simple bitwise operations.

\begin{codebox}
\begin{Verbatim}[commandchars=\\\{\}]
\PY{k+kd}{enum}\PY{+w}{ }\PY{n}{Permission}\PY{+w}{ }\PY{p}{\PYZob{}}\PY{+w}{ }\PY{n}{READ}\PY{p}{,}\PY{+w}{ }\PY{n}{WRITE}\PY{p}{,}\PY{+w}{ }\PY{n}{EXECUTE}\PY{p}{,}\PY{+w}{ }\PY{n}{DELETE}\PY{+w}{ }\PY{p}{\PYZcb{}}

\PY{n}{EnumSet}\PY{o}{\PYZlt{}}\PY{n}{Permission}\PY{o}{\PYZgt{}}\PY{+w}{ }\PY{n}{readWrite}\PY{+w}{ }\PY{o}{=}\PY{+w}{ }\PY{n}{EnumSet}\PY{p}{.}\PY{n+na}{of}\PY{p}{(}\PY{n}{Permission}\PY{p}{.}\PY{n+na}{READ}\PY{p}{,}\PY{+w}{ }\PY{n}{Permission}\PY{p}{.}\PY{n+na}{WRITE}\PY{p}{)}\PY{p}{;}
\PY{n}{EnumSet}\PY{o}{\PYZlt{}}\PY{n}{Permission}\PY{o}{\PYZgt{}}\PY{+w}{ }\PY{n}{all}\PY{+w}{ }\PY{o}{=}\PY{+w}{ }\PY{n}{EnumSet}\PY{p}{.}\PY{n+na}{allOf}\PY{p}{(}\PY{n}{Permission}\PY{p}{.}\PY{n+na}{class}\PY{p}{)}\PY{p}{;}
\PY{n}{EnumSet}\PY{o}{\PYZlt{}}\PY{n}{Permission}\PY{o}{\PYZgt{}}\PY{+w}{ }\PY{n}{none}\PY{+w}{ }\PY{o}{=}\PY{+w}{ }\PY{n}{EnumSet}\PY{p}{.}\PY{n+na}{noneOf}\PY{p}{(}\PY{n}{Permission}\PY{p}{.}\PY{n+na}{class}\PY{p}{)}\PY{p}{;}
\PY{n}{EnumSet}\PY{o}{\PYZlt{}}\PY{n}{Permission}\PY{o}{\PYZgt{}}\PY{+w}{ }\PY{n}{everythingButDelete}\PY{+w}{ }\PY{o}{=}\PY{+w}{ }\PY{n}{EnumSet}\PY{p}{.}\PY{n+na}{complementOf}\PY{p}{(}\PY{n}{EnumSet}\PY{p}{.}\PY{n+na}{of}\PY{p}{(}\PY{n}{Permission}\PY{p}{.}\PY{n+na}{DELETE}\PY{p}{)}\PY{p}{)}\PY{p}{;}

\PY{n}{System}\PY{p}{.}\PY{n+na}{out}\PY{p}{.}\PY{n+na}{println}\PY{p}{(}\PY{n}{readWrite}\PY{p}{)}\PY{p}{;}\PY{+w}{            }\PY{c+c1}{// [READ, WRITE]}
\PY{n}{System}\PY{p}{.}\PY{n+na}{out}\PY{p}{.}\PY{n+na}{println}\PY{p}{(}\PY{n}{everythingButDelete}\PY{p}{)}\PY{p}{;}\PY{+w}{  }\PY{c+c1}{// [READ, WRITE, EXECUTE]}
\PY{n}{System}\PY{p}{.}\PY{n+na}{out}\PY{p}{.}\PY{n+na}{println}\PY{p}{(}\PY{n}{readWrite}\PY{p}{.}\PY{n+na}{contains}\PY{p}{(}\PY{n}{Permission}\PY{p}{.}\PY{n+na}{READ}\PY{p}{)}\PY{p}{)}\PY{p}{;}\PY{+w}{ }\PY{c+c1}{// true}
\end{Verbatim}
\end{codebox}

\textbf{Big-O}: \code{add}/\code{remove}/\code{contains}: O(1) — bitmask operations.

\textbf{Ordering}: natural order of enum declaration (ordinal order).

\textbf{Null handling}: no nulls (throws \code{NullPointerException}).

\textbf{Thread safety}: not thread-safe (though because it’s so compact, wrapping with \code{Collections.\allowbreak{}synchronized\allowbreak{}Set} is cheap if needed).

\textbf{When to choose it}: whenever you have a set of flags/options drawn from a fixed enum — e.g., permissions, feature flags, days of the week. Much faster and more compact than \code{HashSet\textless{}\allowbreak{}EnumType\textgreater{}}.

\section[Immutable Collections (List.of())]{Immutable Collections (List.of())}
\label{h:7:immutable-collections-listof}
Since Java 9, the JDK ships \textbf{factory methods} that create truly immutable collections directly, without needing a mutable builder first.

\begin{codebox}
\begin{Verbatim}[commandchars=\\\{\}]
\PY{n}{List}\PY{o}{\PYZlt{}}\PY{n}{String}\PY{o}{\PYZgt{}}\PY{+w}{ }\PY{n}{list}\PY{+w}{ }\PY{o}{=}\PY{+w}{ }\PY{n}{List}\PY{p}{.}\PY{n+na}{of}\PY{p}{(}\PY{l+s}{\PYZdq{}}\PY{l+s}{a}\PY{l+s}{\PYZdq{}}\PY{p}{,}\PY{+w}{ }\PY{l+s}{\PYZdq{}}\PY{l+s}{b}\PY{l+s}{\PYZdq{}}\PY{p}{,}\PY{+w}{ }\PY{l+s}{\PYZdq{}}\PY{l+s}{c}\PY{l+s}{\PYZdq{}}\PY{p}{)}\PY{p}{;}
\PY{n}{Set}\PY{o}{\PYZlt{}}\PY{n}{String}\PY{o}{\PYZgt{}}\PY{+w}{ }\PY{n}{set}\PY{+w}{ }\PY{o}{=}\PY{+w}{ }\PY{n}{Set}\PY{p}{.}\PY{n+na}{of}\PY{p}{(}\PY{l+s}{\PYZdq{}}\PY{l+s}{x}\PY{l+s}{\PYZdq{}}\PY{p}{,}\PY{+w}{ }\PY{l+s}{\PYZdq{}}\PY{l+s}{y}\PY{l+s}{\PYZdq{}}\PY{p}{,}\PY{+w}{ }\PY{l+s}{\PYZdq{}}\PY{l+s}{z}\PY{l+s}{\PYZdq{}}\PY{p}{)}\PY{p}{;}
\PY{n}{Map}\PY{o}{\PYZlt{}}\PY{n}{String}\PY{p}{,}\PY{+w}{ }\PY{n}{Integer}\PY{o}{\PYZgt{}}\PY{+w}{ }\PY{n}{map}\PY{+w}{ }\PY{o}{=}\PY{+w}{ }\PY{n}{Map}\PY{p}{.}\PY{n+na}{of}\PY{p}{(}\PY{l+s}{\PYZdq{}}\PY{l+s}{a}\PY{l+s}{\PYZdq{}}\PY{p}{,}\PY{+w}{ }\PY{l+m+mi}{1}\PY{p}{,}\PY{+w}{ }\PY{l+s}{\PYZdq{}}\PY{l+s}{b}\PY{l+s}{\PYZdq{}}\PY{p}{,}\PY{+w}{ }\PY{l+m+mi}{2}\PY{p}{)}\PY{p}{;}
\PY{n}{Map}\PY{o}{\PYZlt{}}\PY{n}{String}\PY{p}{,}\PY{+w}{ }\PY{n}{Integer}\PY{o}{\PYZgt{}}\PY{+w}{ }\PY{n}{bigMap}\PY{+w}{ }\PY{o}{=}\PY{+w}{ }\PY{n}{Map}\PY{p}{.}\PY{n+na}{ofEntries}\PY{p}{(}
\PY{+w}{    }\PY{n}{Map}\PY{p}{.}\PY{n+na}{entry}\PY{p}{(}\PY{l+s}{\PYZdq{}}\PY{l+s}{a}\PY{l+s}{\PYZdq{}}\PY{p}{,}\PY{+w}{ }\PY{l+m+mi}{1}\PY{p}{)}\PY{p}{,}
\PY{+w}{    }\PY{n}{Map}\PY{p}{.}\PY{n+na}{entry}\PY{p}{(}\PY{l+s}{\PYZdq{}}\PY{l+s}{b}\PY{l+s}{\PYZdq{}}\PY{p}{,}\PY{+w}{ }\PY{l+m+mi}{2}\PY{p}{)}\PY{p}{,}
\PY{+w}{    }\PY{n}{Map}\PY{p}{.}\PY{n+na}{entry}\PY{p}{(}\PY{l+s}{\PYZdq{}}\PY{l+s}{c}\PY{l+s}{\PYZdq{}}\PY{p}{,}\PY{+w}{ }\PY{l+m+mi}{3}\PY{p}{)}
\PY{p}{)}\PY{p}{;}
\end{Verbatim}
\end{codebox}

Key rules of these factories, all commonly tested in interviews:

\begin{enumerate}
\item 
\textbf{Null-hostile}: passing \code{null} as an element, key, or value throws \code{NullPointerException} immediately — even just constructing the list fails, it’s not a lazy check.

\begin{codebox}
\begin{Verbatim}[commandchars=\\\{\}]
\PY{n}{List}\PY{p}{.}\PY{n+na}{of}\PY{p}{(}\PY{l+s}{\PYZdq{}}\PY{l+s}{a}\PY{l+s}{\PYZdq{}}\PY{p}{,}\PY{+w}{ }\PY{k+kc}{null}\PY{p}{)}\PY{p}{;}\PY{+w}{ }\PY{c+c1}{// throws NullPointerException}
\end{Verbatim}
\end{codebox}

\item 
\textbf{No duplicates in \code{Set.\allowbreak{}of}/\code{Map.\allowbreak{}of}}: passing a duplicate element or key throws \code{Illegal\allowbreak{}Argument\allowbreak{}Exception} at creation time.

\begin{codebox}
\begin{Verbatim}[commandchars=\\\{\}]
\PY{n}{Set}\PY{p}{.}\PY{n+na}{of}\PY{p}{(}\PY{l+s}{\PYZdq{}}\PY{l+s}{a}\PY{l+s}{\PYZdq{}}\PY{p}{,}\PY{+w}{ }\PY{l+s}{\PYZdq{}}\PY{l+s}{a}\PY{l+s}{\PYZdq{}}\PY{p}{)}\PY{p}{;}\PY{+w}{ }\PY{c+c1}{// throws IllegalArgumentException: duplicate element}
\PY{n}{Map}\PY{p}{.}\PY{n+na}{of}\PY{p}{(}\PY{l+s}{\PYZdq{}}\PY{l+s}{a}\PY{l+s}{\PYZdq{}}\PY{p}{,}\PY{+w}{ }\PY{l+m+mi}{1}\PY{p}{,}\PY{+w}{ }\PY{l+s}{\PYZdq{}}\PY{l+s}{a}\PY{l+s}{\PYZdq{}}\PY{p}{,}\PY{+w}{ }\PY{l+m+mi}{2}\PY{p}{)}\PY{p}{;}\PY{+w}{ }\PY{c+c1}{// throws IllegalArgumentException: duplicate key}
\end{Verbatim}
\end{codebox}

\item 
\textbf{Truly unmodifiable}: any mutating call (\code{add}, \code{remove}, \code{set}, \code{put}, \code{clear}, …) throws \code{Unsupported\allowbreak{}Operation\allowbreak{}Exception}.

\begin{codebox}
\begin{Verbatim}[commandchars=\\\{\}]
\PY{n}{List}\PY{o}{\PYZlt{}}\PY{n}{String}\PY{o}{\PYZgt{}}\PY{+w}{ }\PY{n}{immutable}\PY{+w}{ }\PY{o}{=}\PY{+w}{ }\PY{n}{List}\PY{p}{.}\PY{n+na}{of}\PY{p}{(}\PY{l+s}{\PYZdq{}}\PY{l+s}{a}\PY{l+s}{\PYZdq{}}\PY{p}{,}\PY{+w}{ }\PY{l+s}{\PYZdq{}}\PY{l+s}{b}\PY{l+s}{\PYZdq{}}\PY{p}{)}\PY{p}{;}
\PY{n}{immutable}\PY{p}{.}\PY{n+na}{add}\PY{p}{(}\PY{l+s}{\PYZdq{}}\PY{l+s}{c}\PY{l+s}{\PYZdq{}}\PY{p}{)}\PY{p}{;}\PY{+w}{ }\PY{c+c1}{// throws UnsupportedOperationException}
\end{Verbatim}
\end{codebox}

\end{enumerate}

\textbf{\code{List.\allowbreak{}copyOf} / \code{Set.\allowbreak{}copyOf} / \code{Map.\allowbreak{}copyOf}}: produce an immutable \emph{copy} of an existing collection. If the source is already an immutable instance from these same factories, it may return the same instance instead of copying (an optimization, not something to rely on).

\begin{codebox}
\begin{Verbatim}[commandchars=\\\{\}]
\PY{n}{List}\PY{o}{\PYZlt{}}\PY{n}{String}\PY{o}{\PYZgt{}}\PY{+w}{ }\PY{n}{mutable}\PY{+w}{ }\PY{o}{=}\PY{+w}{ }\PY{k}{new}\PY{+w}{ }\PY{n}{ArrayList}\PY{o}{\PYZlt{}}\PY{o}{\PYZgt{}}\PY{p}{(}\PY{n}{List}\PY{p}{.}\PY{n+na}{of}\PY{p}{(}\PY{l+s}{\PYZdq{}}\PY{l+s}{a}\PY{l+s}{\PYZdq{}}\PY{p}{,}\PY{+w}{ }\PY{l+s}{\PYZdq{}}\PY{l+s}{b}\PY{l+s}{\PYZdq{}}\PY{p}{)}\PY{p}{)}\PY{p}{;}
\PY{n}{List}\PY{o}{\PYZlt{}}\PY{n}{String}\PY{o}{\PYZgt{}}\PY{+w}{ }\PY{n}{snapshot}\PY{+w}{ }\PY{o}{=}\PY{+w}{ }\PY{n}{List}\PY{p}{.}\PY{n+na}{copyOf}\PY{p}{(}\PY{n}{mutable}\PY{p}{)}\PY{p}{;}\PY{+w}{ }\PY{c+c1}{// independent, immutable snapshot}
\PY{n}{mutable}\PY{p}{.}\PY{n+na}{add}\PY{p}{(}\PY{l+s}{\PYZdq{}}\PY{l+s}{c}\PY{l+s}{\PYZdq{}}\PY{p}{)}\PY{p}{;}
\PY{n}{System}\PY{p}{.}\PY{n+na}{out}\PY{p}{.}\PY{n+na}{println}\PY{p}{(}\PY{n}{snapshot}\PY{p}{)}\PY{p}{;}\PY{+w}{ }\PY{c+c1}{// [a, b] \PYZhy{}\PYZhy{} unaffected by later changes to \PYZdq{}mutable\PYZdq{}}
\end{Verbatim}
\end{codebox}

\textbf{\code{Collectors.\allowbreak{}to\allowbreak{}Unmodifiable\allowbreak{}List/\allowbreak{}Set/\allowbreak{}Map}}: the Streams equivalent, producing an immutable result directly from a stream pipeline.

\begin{codebox}
\begin{Verbatim}[commandchars=\\\{\}]
\PY{n}{List}\PY{o}{\PYZlt{}}\PY{n}{String}\PY{o}{\PYZgt{}}\PY{+w}{ }\PY{n}{upper}\PY{+w}{ }\PY{o}{=}\PY{+w}{ }\PY{n}{java}\PY{p}{.}\PY{n+na}{util}\PY{p}{.}\PY{n+na}{stream}\PY{p}{.}\PY{n+na}{Stream}\PY{p}{.}\PY{n+na}{of}\PY{p}{(}\PY{l+s}{\PYZdq{}}\PY{l+s}{a}\PY{l+s}{\PYZdq{}}\PY{p}{,}\PY{+w}{ }\PY{l+s}{\PYZdq{}}\PY{l+s}{b}\PY{l+s}{\PYZdq{}}\PY{p}{,}\PY{+w}{ }\PY{l+s}{\PYZdq{}}\PY{l+s}{c}\PY{l+s}{\PYZdq{}}\PY{p}{)}
\PY{+w}{        }\PY{p}{.}\PY{n+na}{map}\PY{p}{(}\PY{n}{String}\PY{p}{:}\PY{p}{:}\PY{n}{toUpperCase}\PY{p}{)}
\PY{+w}{        }\PY{p}{.}\PY{n+na}{collect}\PY{p}{(}\PY{n}{java}\PY{p}{.}\PY{n+na}{util}\PY{p}{.}\PY{n+na}{stream}\PY{p}{.}\PY{n+na}{Collectors}\PY{p}{.}\PY{n+na}{toUnmodifiableList}\PY{p}{(}\PY{p}{)}\PY{p}{)}\PY{p}{;}
\PY{c+c1}{// upper.add(\PYZdq{}D\PYZdq{}) would throw UnsupportedOperationException}
\end{Verbatim}
\end{codebox}

\textbf{Immutable factories vs \code{Collections.\allowbreak{}unmodifiable\allowbreak{}List} — not the same thing}: \code{Collections.\allowbreak{}unmodifiable\allowbreak{}Xxx} wraps an existing mutable collection in a \textbf{read-only view}. The view itself can’t be modified directly, but if you keep a reference to the underlying mutable collection, changes there still show through the view.

\begin{codebox}
\begin{Verbatim}[commandchars=\\\{\}]
\PY{n}{List}\PY{o}{\PYZlt{}}\PY{n}{String}\PY{o}{\PYZgt{}}\PY{+w}{ }\PY{n}{backing}\PY{+w}{ }\PY{o}{=}\PY{+w}{ }\PY{k}{new}\PY{+w}{ }\PY{n}{ArrayList}\PY{o}{\PYZlt{}}\PY{o}{\PYZgt{}}\PY{p}{(}\PY{n}{List}\PY{p}{.}\PY{n+na}{of}\PY{p}{(}\PY{l+s}{\PYZdq{}}\PY{l+s}{a}\PY{l+s}{\PYZdq{}}\PY{p}{,}\PY{+w}{ }\PY{l+s}{\PYZdq{}}\PY{l+s}{b}\PY{l+s}{\PYZdq{}}\PY{p}{)}\PY{p}{)}\PY{p}{;}
\PY{n}{List}\PY{o}{\PYZlt{}}\PY{n}{String}\PY{o}{\PYZgt{}}\PY{+w}{ }\PY{n}{view}\PY{+w}{ }\PY{o}{=}\PY{+w}{ }\PY{n}{Collections}\PY{p}{.}\PY{n+na}{unmodifiableList}\PY{p}{(}\PY{n}{backing}\PY{p}{)}\PY{p}{;}

\PY{n}{view}\PY{p}{.}\PY{n+na}{add}\PY{p}{(}\PY{l+s}{\PYZdq{}}\PY{l+s}{c}\PY{l+s}{\PYZdq{}}\PY{p}{)}\PY{p}{;}\PY{+w}{     }\PY{c+c1}{// throws UnsupportedOperationException (the view is read\PYZhy{}only)}
\PY{n}{backing}\PY{p}{.}\PY{n+na}{add}\PY{p}{(}\PY{l+s}{\PYZdq{}}\PY{l+s}{c}\PY{l+s}{\PYZdq{}}\PY{p}{)}\PY{p}{;}\PY{+w}{  }\PY{c+c1}{// allowed! and it changes what \PYZdq{}view\PYZdq{} shows}
\PY{n}{System}\PY{p}{.}\PY{n+na}{out}\PY{p}{.}\PY{n+na}{println}\PY{p}{(}\PY{n}{view}\PY{p}{)}\PY{p}{;}\PY{+w}{ }\PY{c+c1}{// [a, b, c] \PYZhy{}\PYZhy{} the \PYZdq{}immutable\PYZdq{} view just changed}
\end{Verbatim}
\end{codebox}

{\tablefont
\begin{xltabular}{\linewidth}{@{}L{0.918}L{1.200}L{0.882}@{}}
\toprule
\rowcolor{tablehdr}
\thd{} & \thd{\code{Collections.\allowbreak{}unmodifiable\allowbreak{}List}} & \thd{\code{List.\allowbreak{}of()} / \code{List.\allowbreak{}copyOf()}} \\
\midrule
\endhead
Backed by & the original mutable list (a view) & its own internal storage \\
Changes to source visible? & yes & no (copy) \\
Allows null elements? & depends on the backing list & never (throws NPE) \\
Created before Java 9? & yes (Java 1.2) & no (Java 9+) \\
\bottomrule
\end{xltabular}
}

\section[Concurrent Collections]{Concurrent Collections}
\label{h:7:concurrent-collections}
Plain \code{java.\allowbreak{}util} collections (\code{ArrayList}, \code{HashMap}, etc.) are \textbf{not} thread-safe. \code{java.\allowbreak{}util.\allowbreak{}concurrent} provides collections designed for multithreaded use, with better scalability than simply synchronizing everything.

\subsection[Why Collections.synchronizedList is not enough]{Why \code{Collections.\allowbreak{}synchronized\allowbreak{}List} is not enough}
\label{h:7:why-collectionssynchronizedlist-is-not-enough}
\code{Collections.\allowbreak{}synchronized\allowbreak{}Xxx} wraps a collection so each \textbf{individual method call} is synchronized (guarded by a single lock). That protects against corruption from concurrent single calls, but it does \textbf{not} make compound actions (check-then-act, like “if absent, then add”) atomic.

\begin{codebox}
\begin{Verbatim}[commandchars=\\\{\}]
\PY{n}{List}\PY{o}{\PYZlt{}}\PY{n}{String}\PY{o}{\PYZgt{}}\PY{+w}{ }\PY{n}{list}\PY{+w}{ }\PY{o}{=}\PY{+w}{ }\PY{n}{Collections}\PY{p}{.}\PY{n+na}{synchronizedList}\PY{p}{(}\PY{k}{new}\PY{+w}{ }\PY{n}{ArrayList}\PY{o}{\PYZlt{}}\PY{o}{\PYZgt{}}\PY{p}{(}\PY{p}{)}\PY{p}{)}\PY{p}{;}

\PY{c+c1}{// BUG: this is two separate synchronized calls, not one atomic operation.}
\PY{c+c1}{// Between contains() and add(), another thread could add the same item.}
\PY{k}{if}\PY{+w}{ }\PY{p}{(}\PY{o}{!}\PY{n}{list}\PY{p}{.}\PY{n+na}{contains}\PY{p}{(}\PY{l+s}{\PYZdq{}}\PY{l+s}{item}\PY{l+s}{\PYZdq{}}\PY{p}{)}\PY{p}{)}\PY{+w}{ }\PY{p}{\PYZob{}}
\PY{+w}{    }\PY{n}{list}\PY{p}{.}\PY{n+na}{add}\PY{p}{(}\PY{l+s}{\PYZdq{}}\PY{l+s}{item}\PY{l+s}{\PYZdq{}}\PY{p}{)}\PY{p}{;}\PY{+w}{ }\PY{c+c1}{// race condition: might now add a duplicate}
\PY{p}{\PYZcb{}}

\PY{c+c1}{// Correct: manually synchronize the whole compound action on the list\PYZsq{}s own lock}
\PY{k+kd}{synchronized}\PY{+w}{ }\PY{p}{(}\PY{n}{list}\PY{p}{)}\PY{+w}{ }\PY{p}{\PYZob{}}
\PY{+w}{    }\PY{k}{if}\PY{+w}{ }\PY{p}{(}\PY{o}{!}\PY{n}{list}\PY{p}{.}\PY{n+na}{contains}\PY{p}{(}\PY{l+s}{\PYZdq{}}\PY{l+s}{item}\PY{l+s}{\PYZdq{}}\PY{p}{)}\PY{p}{)}\PY{+w}{ }\PY{p}{\PYZob{}}
\PY{+w}{        }\PY{n}{list}\PY{p}{.}\PY{n+na}{add}\PY{p}{(}\PY{l+s}{\PYZdq{}}\PY{l+s}{item}\PY{l+s}{\PYZdq{}}\PY{p}{)}\PY{p}{;}
\PY{+w}{    }\PY{p}{\PYZcb{}}
\PY{p}{\PYZcb{}}

\PY{c+c1}{// Also: even a single iteration needs manual synchronization to avoid}
\PY{c+c1}{// ConcurrentModificationException from another thread\PYZsq{}s concurrent write}
\PY{k+kd}{synchronized}\PY{+w}{ }\PY{p}{(}\PY{n}{list}\PY{p}{)}\PY{+w}{ }\PY{p}{\PYZob{}}
\PY{+w}{    }\PY{k}{for}\PY{+w}{ }\PY{p}{(}\PY{n}{String}\PY{+w}{ }\PY{n}{s}\PY{+w}{ }\PY{p}{:}\PY{+w}{ }\PY{n}{list}\PY{p}{)}\PY{+w}{ }\PY{p}{\PYZob{}}
\PY{+w}{        }\PY{n}{System}\PY{p}{.}\PY{n+na}{out}\PY{p}{.}\PY{n+na}{println}\PY{p}{(}\PY{n}{s}\PY{p}{)}\PY{p}{;}
\PY{+w}{    }\PY{p}{\PYZcb{}}
\PY{p}{\PYZcb{}}
\end{Verbatim}
\end{codebox}

This is exactly why \code{java.\allowbreak{}util.\allowbreak{}concurrent} collections exist: they provide atomic compound operations (like \code{putIfAbsent}) without you needing an external lock.

\subsection[ConcurrentHashMap]{ConcurrentHashMap}
\label{h:7:concurrenthashmap}
A highly-scalable map for concurrent access. Internally (JDK 8+) it does \textbf{not} lock the whole map; it uses fine-grained locking per bin (bucket) plus CAS (compare-and-swap) operations for many updates, so multiple threads can write to different parts of the map at the same time.

\begin{codebox}
\begin{Verbatim}[commandchars=\\\{\}]
\PY{n}{ConcurrentHashMap}\PY{o}{\PYZlt{}}\PY{n}{String}\PY{p}{,}\PY{+w}{ }\PY{n}{Integer}\PY{o}{\PYZgt{}}\PY{+w}{ }\PY{n}{counts}\PY{+w}{ }\PY{o}{=}\PY{+w}{ }\PY{k}{new}\PY{+w}{ }\PY{n}{ConcurrentHashMap}\PY{o}{\PYZlt{}}\PY{o}{\PYZgt{}}\PY{p}{(}\PY{p}{)}\PY{p}{;}
\PY{n}{counts}\PY{p}{.}\PY{n+na}{putIfAbsent}\PY{p}{(}\PY{l+s}{\PYZdq{}}\PY{l+s}{errors}\PY{l+s}{\PYZdq{}}\PY{p}{,}\PY{+w}{ }\PY{l+m+mi}{0}\PY{p}{)}\PY{p}{;}
\PY{n}{counts}\PY{p}{.}\PY{n+na}{compute}\PY{p}{(}\PY{l+s}{\PYZdq{}}\PY{l+s}{errors}\PY{l+s}{\PYZdq{}}\PY{p}{,}\PY{+w}{ }\PY{p}{(}\PY{n}{k}\PY{p}{,}\PY{+w}{ }\PY{n}{v}\PY{p}{)}\PY{+w}{ }\PY{o}{\PYZhy{}}\PY{o}{\PYZgt{}}\PY{+w}{ }\PY{n}{v}\PY{+w}{ }\PY{o}{+}\PY{+w}{ }\PY{l+m+mi}{1}\PY{p}{)}\PY{p}{;}\PY{+w}{ }\PY{c+c1}{// atomic read\PYZhy{}modify\PYZhy{}write}
\PY{n}{counts}\PY{p}{.}\PY{n+na}{merge}\PY{p}{(}\PY{l+s}{\PYZdq{}}\PY{l+s}{errors}\PY{l+s}{\PYZdq{}}\PY{p}{,}\PY{+w}{ }\PY{l+m+mi}{1}\PY{p}{,}\PY{+w}{ }\PY{n}{Integer}\PY{p}{:}\PY{p}{:}\PY{n}{sum}\PY{p}{)}\PY{p}{;}\PY{+w}{   }\PY{c+c1}{// also atomic}

\PY{c+c1}{// Safe iteration while other threads mutate the map \PYZhy{}\PYZhy{} no ConcurrentModificationException,}
\PY{c+c1}{// but the view may or may not reflect very recent concurrent updates (\PYZdq{}weakly consistent\PYZdq{})}
\PY{k}{for}\PY{+w}{ }\PY{p}{(}\PY{n}{String}\PY{+w}{ }\PY{n}{key}\PY{+w}{ }\PY{p}{:}\PY{+w}{ }\PY{n}{counts}\PY{p}{.}\PY{n+na}{keySet}\PY{p}{(}\PY{p}{)}\PY{p}{)}\PY{+w}{ }\PY{p}{\PYZob{}}
\PY{+w}{    }\PY{n}{System}\PY{p}{.}\PY{n+na}{out}\PY{p}{.}\PY{n+na}{println}\PY{p}{(}\PY{n}{key}\PY{+w}{ }\PY{o}{+}\PY{+w}{ }\PY{l+s}{\PYZdq{}}\PY{l+s}{ = }\PY{l+s}{\PYZdq{}}\PY{+w}{ }\PY{o}{+}\PY{+w}{ }\PY{n}{counts}\PY{p}{.}\PY{n+na}{get}\PY{p}{(}\PY{n}{key}\PY{p}{)}\PY{p}{)}\PY{p}{;}
\PY{p}{\PYZcb{}}
\end{Verbatim}
\end{codebox}

Notes:

\begin{itemize}
\item 
Does \textbf{not} allow \code{null} keys or values (unlike \code{HashMap}) — this is deliberate, because \code{null} would be ambiguous in a concurrent \code{get} (is it “not present” or “present with null value”?).

\item 
Iterators are \textbf{weakly consistent}: they never throw \code{Concurrent\allowbreak{}Modification\allowbreak{}Exception}, and they reflect the state of the map at some point during or after the iterator was created, but not necessarily every single concurrent update.

\item 
\code{size()} is an estimate under heavy concurrent modification, not a hard snapshot.

\end{itemize}

\subsection[CopyOnWriteArrayList]{CopyOnWriteArrayList}
\label{h:7:copyonwritearraylist}
A \code{List} where every mutating operation (\code{add}, \code{remove}, \code{set}) makes a \textbf{fresh copy of the entire underlying array}. Reads never block and never see partial writes — they just read a stable snapshot array.

\begin{codebox}
\begin{Verbatim}[commandchars=\\\{\}]
\PY{n}{CopyOnWriteArrayList}\PY{o}{\PYZlt{}}\PY{n}{String}\PY{o}{\PYZgt{}}\PY{+w}{ }\PY{n}{listeners}\PY{+w}{ }\PY{o}{=}\PY{+w}{ }\PY{k}{new}\PY{+w}{ }\PY{n}{CopyOnWriteArrayList}\PY{o}{\PYZlt{}}\PY{o}{\PYZgt{}}\PY{p}{(}\PY{p}{)}\PY{p}{;}
\PY{n}{listeners}\PY{p}{.}\PY{n+na}{add}\PY{p}{(}\PY{l+s}{\PYZdq{}}\PY{l+s}{listenerA}\PY{l+s}{\PYZdq{}}\PY{p}{)}\PY{p}{;}
\PY{n}{listeners}\PY{p}{.}\PY{n+na}{add}\PY{p}{(}\PY{l+s}{\PYZdq{}}\PY{l+s}{listenerB}\PY{l+s}{\PYZdq{}}\PY{p}{)}\PY{p}{;}

\PY{c+c1}{// Safe to iterate even if another thread adds/removes concurrently \PYZhy{}\PYZhy{}}
\PY{c+c1}{// the iterator works on the snapshot array taken when the loop started,}
\PY{c+c1}{// so no ConcurrentModificationException, ever.}
\PY{k}{for}\PY{+w}{ }\PY{p}{(}\PY{n}{String}\PY{+w}{ }\PY{n}{listener}\PY{+w}{ }\PY{p}{:}\PY{+w}{ }\PY{n}{listeners}\PY{p}{)}\PY{+w}{ }\PY{p}{\PYZob{}}
\PY{+w}{    }\PY{n}{System}\PY{p}{.}\PY{n+na}{out}\PY{p}{.}\PY{n+na}{println}\PY{p}{(}\PY{l+s}{\PYZdq{}}\PY{l+s}{Notifying }\PY{l+s}{\PYZdq{}}\PY{+w}{ }\PY{o}{+}\PY{+w}{ }\PY{n}{listener}\PY{p}{)}\PY{p}{;}
\PY{+w}{    }\PY{n}{listeners}\PY{p}{.}\PY{n+na}{remove}\PY{p}{(}\PY{n}{listener}\PY{p}{)}\PY{p}{;}\PY{+w}{ }\PY{c+c1}{// does NOT affect this ongoing iteration\PYZsq{}s snapshot}
\PY{p}{\PYZcb{}}
\end{Verbatim}
\end{codebox}

\textbf{When to use it}: read-heavy, write-rare scenarios, like a list of event listeners/observers. \textbf{Avoid} it for write-heavy workloads — every single write copies the whole array, which is O(n) per write and creates a lot of garbage.

\subsection[BlockingQueue family]{BlockingQueue family}
\label{h:7:blockingqueue-family}
\code{BlockingQueue\textless{}\allowbreak{}E\textgreater{}} extends \code{Queue} and adds methods that \textbf{block} the calling thread instead of failing, until space or an element becomes available. This is the standard building block for producer-consumer pipelines.

{\tablefont
\begin{xltabular}{\linewidth}{@{}L{0.925}L{1.037}L{1.037}@{}}
\toprule
\rowcolor{tablehdr}
\thd{Method} & \thd{On empty (poll side)} & \thd{On full (offer side)} \\
\midrule
\endhead
\code{add}/\code{remove} & throws exception & throws exception \\
\code{offer}/\code{poll} & returns \code{null}/\code{false} immediately & returns \code{false} immediately \\
\code{put}/\code{take} & \textbf{blocks} until available & \textbf{blocks} until space \\
\code{offer(\allowbreak{}e,\allowbreak{}~\allowbreak{}timeout)} / \code{poll(\allowbreak{}timeout)} & blocks up to a timeout, then gives up & blocks up to a timeout, then gives up \\
\bottomrule
\end{xltabular}
}

\begin{codebox}
\begin{Verbatim}[commandchars=\\\{\}]
\PY{n}{BlockingQueue}\PY{o}{\PYZlt{}}\PY{n}{String}\PY{o}{\PYZgt{}}\PY{+w}{ }\PY{n}{jobs}\PY{+w}{ }\PY{o}{=}\PY{+w}{ }\PY{k}{new}\PY{+w}{ }\PY{n}{java}\PY{p}{.}\PY{n+na}{util}\PY{p}{.}\PY{n+na}{concurrent}\PY{p}{.}\PY{n+na}{LinkedBlockingQueue}\PY{o}{\PYZlt{}}\PY{o}{\PYZgt{}}\PY{p}{(}\PY{l+m+mi}{100}\PY{p}{)}\PY{p}{;}\PY{+w}{ }\PY{c+c1}{// bounded capacity}

\PY{c+c1}{// Producer thread}
\PY{k}{new}\PY{+w}{ }\PY{n}{Thread}\PY{p}{(}\PY{p}{(}\PY{p}{)}\PY{+w}{ }\PY{o}{\PYZhy{}}\PY{o}{\PYZgt{}}\PY{+w}{ }\PY{p}{\PYZob{}}
\PY{+w}{    }\PY{k}{try}\PY{+w}{ }\PY{p}{\PYZob{}}
\PY{+w}{        }\PY{n}{jobs}\PY{p}{.}\PY{n+na}{put}\PY{p}{(}\PY{l+s}{\PYZdq{}}\PY{l+s}{job\PYZhy{}1}\PY{l+s}{\PYZdq{}}\PY{p}{)}\PY{p}{;}\PY{+w}{ }\PY{c+c1}{// blocks if the queue is full}
\PY{+w}{    }\PY{p}{\PYZcb{}}\PY{+w}{ }\PY{k}{catch}\PY{+w}{ }\PY{p}{(}\PY{n}{InterruptedException}\PY{+w}{ }\PY{n}{e}\PY{p}{)}\PY{+w}{ }\PY{p}{\PYZob{}}
\PY{+w}{        }\PY{n}{Thread}\PY{p}{.}\PY{n+na}{currentThread}\PY{p}{(}\PY{p}{)}\PY{p}{.}\PY{n+na}{interrupt}\PY{p}{(}\PY{p}{)}\PY{p}{;}
\PY{+w}{    }\PY{p}{\PYZcb{}}
\PY{p}{\PYZcb{}}\PY{p}{)}\PY{p}{.}\PY{n+na}{start}\PY{p}{(}\PY{p}{)}\PY{p}{;}

\PY{c+c1}{// Consumer thread}
\PY{k}{new}\PY{+w}{ }\PY{n}{Thread}\PY{p}{(}\PY{p}{(}\PY{p}{)}\PY{+w}{ }\PY{o}{\PYZhy{}}\PY{o}{\PYZgt{}}\PY{+w}{ }\PY{p}{\PYZob{}}
\PY{+w}{    }\PY{k}{try}\PY{+w}{ }\PY{p}{\PYZob{}}
\PY{+w}{        }\PY{n}{String}\PY{+w}{ }\PY{n}{job}\PY{+w}{ }\PY{o}{=}\PY{+w}{ }\PY{n}{jobs}\PY{p}{.}\PY{n+na}{take}\PY{p}{(}\PY{p}{)}\PY{p}{;}\PY{+w}{ }\PY{c+c1}{// blocks until a job is available}
\PY{+w}{        }\PY{n}{System}\PY{p}{.}\PY{n+na}{out}\PY{p}{.}\PY{n+na}{println}\PY{p}{(}\PY{l+s}{\PYZdq{}}\PY{l+s}{processing }\PY{l+s}{\PYZdq{}}\PY{+w}{ }\PY{o}{+}\PY{+w}{ }\PY{n}{job}\PY{p}{)}\PY{p}{;}
\PY{+w}{    }\PY{p}{\PYZcb{}}\PY{+w}{ }\PY{k}{catch}\PY{+w}{ }\PY{p}{(}\PY{n}{InterruptedException}\PY{+w}{ }\PY{n}{e}\PY{p}{)}\PY{+w}{ }\PY{p}{\PYZob{}}
\PY{+w}{        }\PY{n}{Thread}\PY{p}{.}\PY{n+na}{currentThread}\PY{p}{(}\PY{p}{)}\PY{p}{.}\PY{n+na}{interrupt}\PY{p}{(}\PY{p}{)}\PY{p}{;}
\PY{+w}{    }\PY{p}{\PYZcb{}}
\PY{p}{\PYZcb{}}\PY{p}{)}\PY{p}{.}\PY{n+na}{start}\PY{p}{(}\PY{p}{)}\PY{p}{;}
\end{Verbatim}
\end{codebox}

Common implementations:

\begin{itemize}
\item 
\textbf{\code{ArrayBlockingQueue}}: fixed-capacity, array-backed, FIFO.

\item 
\textbf{\code{LinkedBlockingQueue}}: optionally bounded, linked-node-backed, FIFO. Generally higher throughput than \code{ArrayBlockingQueue} under contention.

\item 
\textbf{\code{PriorityBlockingQueue}}: unbounded, orders elements by priority like \code{PriorityQueue}, but with blocking \code{take()}.

\item 
\textbf{\code{SynchronousQueue}}: zero capacity — a \code{put} must wait for a matching \code{take} (a direct handoff). Used inside some \code{ExecutorService} configurations.

\item 
\textbf{\code{DelayQueue}}: elements become available only after their delay expires — useful for scheduling/retry logic.

\end{itemize}

\subsection[ConcurrentLinkedQueue / ConcurrentLinkedDeque]{ConcurrentLinkedQueue / ConcurrentLinkedDeque}
\label{h:7:concurrentlinkedqueue--concurrentlinkeddeque}
Unbounded, \textbf{non-blocking}, thread-safe queues/deques built with CAS (lock-free) operations instead of locks. They implement plain \code{Queue}/\code{Deque}, not \code{BlockingQueue} — there is no \code{put}/\code{take} that waits; \code{offer}/\code{poll} return immediately.

\begin{codebox}
\begin{Verbatim}[commandchars=\\\{\}]
\PY{n}{ConcurrentLinkedQueue}\PY{o}{\PYZlt{}}\PY{n}{String}\PY{o}{\PYZgt{}}\PY{+w}{ }\PY{n}{queue}\PY{+w}{ }\PY{o}{=}\PY{+w}{ }\PY{k}{new}\PY{+w}{ }\PY{n}{ConcurrentLinkedQueue}\PY{o}{\PYZlt{}}\PY{o}{\PYZgt{}}\PY{p}{(}\PY{p}{)}\PY{p}{;}
\PY{n}{queue}\PY{p}{.}\PY{n+na}{offer}\PY{p}{(}\PY{l+s}{\PYZdq{}}\PY{l+s}{task1}\PY{l+s}{\PYZdq{}}\PY{p}{)}\PY{p}{;}
\PY{n}{queue}\PY{p}{.}\PY{n+na}{offer}\PY{p}{(}\PY{l+s}{\PYZdq{}}\PY{l+s}{task2}\PY{l+s}{\PYZdq{}}\PY{p}{)}\PY{p}{;}
\PY{n}{System}\PY{p}{.}\PY{n+na}{out}\PY{p}{.}\PY{n+na}{println}\PY{p}{(}\PY{n}{queue}\PY{p}{.}\PY{n+na}{poll}\PY{p}{(}\PY{p}{)}\PY{p}{)}\PY{p}{;}\PY{+w}{ }\PY{c+c1}{// task1, thread\PYZhy{}safe, no blocking, no external lock needed}
\end{Verbatim}
\end{codebox}

\textbf{When to use it}: high-throughput producer-consumer scenarios where you don’t need threads to block and wait — e.g., a metrics collector where producers just fire-and-forget and a consumer periodically drains what’s there.

\subsection[Quick summary table]{Quick summary table}
\label{h:7:quick-summary-table}
{\tablefont
\begin{xltabular}{\linewidth}{@{}L{1.067}L{0.426}L{0.587}L{1.920}@{}}
\toprule
\rowcolor{tablehdr}
\thd{Class} & \thd{Blocking?} & \thd{Bounded?} & \thd{Typical use} \\
\midrule
\endhead
\code{ConcurrentHashMap} & no & no & general-purpose concurrent map \\
\code{CopyOnWriteArrayList} & no & no & read-heavy list, e.g. listener lists \\
\code{ArrayBlockingQueue} & yes & yes (fixed) & producer-consumer with backpressure \\
\code{LinkedBlockingQueue} & yes & optional & producer-consumer, general purpose \\
\code{ConcurrentLinkedQueue} & no & no & lock-free, high-throughput, non-blocking \\
\code{SynchronousQueue} & yes & zero capacity & direct handoff between threads \\
\bottomrule
\end{xltabular}
}

\section[Common Code-Review Interview Pitfalls]{Common Code-Review Interview Pitfalls}
\label{h:7:common-code-review-interview-pitfalls}
\begin{enumerate}
\item 
\textbf{Using \code{HashMap}/\code{HashSet} iteration order as if it were guaranteed.}
Why it matters: order can change between JDK versions, or even between runs, and silently breaks logic or tests that depend on it.

\begin{codebox}
\begin{Verbatim}[commandchars=\\\{\}]
\PY{c+c1}{// Before: relies on unspecified order}
\PY{n}{Map}\PY{o}{\PYZlt{}}\PY{n}{String}\PY{p}{,}\PY{+w}{ }\PY{n}{Integer}\PY{o}{\PYZgt{}}\PY{+w}{ }\PY{n}{m}\PY{+w}{ }\PY{o}{=}\PY{+w}{ }\PY{k}{new}\PY{+w}{ }\PY{n}{HashMap}\PY{o}{\PYZlt{}}\PY{o}{\PYZgt{}}\PY{p}{(}\PY{p}{)}\PY{p}{;}
\PY{c+c1}{// ... assumes keys always print in insertion order}

\PY{c+c1}{// After: use LinkedHashMap if order must be predictable}
\PY{n}{Map}\PY{o}{\PYZlt{}}\PY{n}{String}\PY{p}{,}\PY{+w}{ }\PY{n}{Integer}\PY{o}{\PYZgt{}}\PY{+w}{ }\PY{n}{m}\PY{+w}{ }\PY{o}{=}\PY{+w}{ }\PY{k}{new}\PY{+w}{ }\PY{n}{LinkedHashMap}\PY{o}{\PYZlt{}}\PY{o}{\PYZgt{}}\PY{p}{(}\PY{p}{)}\PY{p}{;}
\end{Verbatim}
\end{codebox}

\item 
\textbf{Modifying a collection while iterating with a for-each loop.}
Why it matters: fail-fast iterators detect structural changes and throw \code{Concurrent\allowbreak{}Modification\allowbreak{}Exception} — this includes single-threaded code, not just concurrent access.

\begin{codebox}
\begin{Verbatim}[commandchars=\\\{\}]
\PY{c+c1}{// Before: throws ConcurrentModificationException}
\PY{k}{for}\PY{+w}{ }\PY{p}{(}\PY{n}{String}\PY{+w}{ }\PY{n}{s}\PY{+w}{ }\PY{p}{:}\PY{+w}{ }\PY{n}{list}\PY{p}{)}\PY{+w}{ }\PY{p}{\PYZob{}}
\PY{+w}{    }\PY{k}{if}\PY{+w}{ }\PY{p}{(}\PY{n}{s}\PY{p}{.}\PY{n+na}{isEmpty}\PY{p}{(}\PY{p}{)}\PY{p}{)}\PY{+w}{ }\PY{n}{list}\PY{p}{.}\PY{n+na}{remove}\PY{p}{(}\PY{n}{s}\PY{p}{)}\PY{p}{;}
\PY{p}{\PYZcb{}}

\PY{c+c1}{// After: use the iterator\PYZsq{}s own remove, or removeIf}
\PY{n}{list}\PY{p}{.}\PY{n+na}{removeIf}\PY{p}{(}\PY{n}{String}\PY{p}{:}\PY{p}{:}\PY{n}{isEmpty}\PY{p}{)}\PY{p}{;}
\PY{c+c1}{// or:}
\PY{n}{Iterator}\PY{o}{\PYZlt{}}\PY{n}{String}\PY{o}{\PYZgt{}}\PY{+w}{ }\PY{n}{it}\PY{+w}{ }\PY{o}{=}\PY{+w}{ }\PY{n}{list}\PY{p}{.}\PY{n+na}{iterator}\PY{p}{(}\PY{p}{)}\PY{p}{;}
\PY{k}{while}\PY{+w}{ }\PY{p}{(}\PY{n}{it}\PY{p}{.}\PY{n+na}{hasNext}\PY{p}{(}\PY{p}{)}\PY{p}{)}\PY{+w}{ }\PY{p}{\PYZob{}}
\PY{+w}{    }\PY{k}{if}\PY{+w}{ }\PY{p}{(}\PY{n}{it}\PY{p}{.}\PY{n+na}{next}\PY{p}{(}\PY{p}{)}\PY{p}{.}\PY{n+na}{isEmpty}\PY{p}{(}\PY{p}{)}\PY{p}{)}\PY{+w}{ }\PY{n}{it}\PY{p}{.}\PY{n+na}{remove}\PY{p}{(}\PY{p}{)}\PY{p}{;}
\PY{p}{\PYZcb{}}
\end{Verbatim}
\end{codebox}

\item 
\textbf{Storing mutable objects as \code{HashMap}/\code{HashSet} keys and then mutating them.}
Why it matters: the bucket a key lives in is chosen from its hash at insertion time; if the key’s hash changes afterward, the entry becomes unreachable (“lost” in the map).

\begin{codebox}
\begin{Verbatim}[commandchars=\\\{\}]
\PY{c+c1}{// Before: mutable key breaks lookup}
\PY{n}{List}\PY{o}{\PYZlt{}}\PY{n}{Integer}\PY{o}{\PYZgt{}}\PY{+w}{ }\PY{n}{key}\PY{+w}{ }\PY{o}{=}\PY{+w}{ }\PY{k}{new}\PY{+w}{ }\PY{n}{ArrayList}\PY{o}{\PYZlt{}}\PY{o}{\PYZgt{}}\PY{p}{(}\PY{n}{List}\PY{p}{.}\PY{n+na}{of}\PY{p}{(}\PY{l+m+mi}{1}\PY{p}{,}\PY{+w}{ }\PY{l+m+mi}{2}\PY{p}{)}\PY{p}{)}\PY{p}{;}
\PY{n}{Set}\PY{o}{\PYZlt{}}\PY{n}{List}\PY{o}{\PYZlt{}}\PY{n}{Integer}\PY{o}{\PYZgt{}}\PY{o}{\PYZgt{}}\PY{+w}{ }\PY{n}{set}\PY{+w}{ }\PY{o}{=}\PY{+w}{ }\PY{k}{new}\PY{+w}{ }\PY{n}{HashSet}\PY{o}{\PYZlt{}}\PY{o}{\PYZgt{}}\PY{p}{(}\PY{p}{)}\PY{p}{;}
\PY{n}{set}\PY{p}{.}\PY{n+na}{add}\PY{p}{(}\PY{n}{key}\PY{p}{)}\PY{p}{;}
\PY{n}{key}\PY{p}{.}\PY{n+na}{add}\PY{p}{(}\PY{l+m+mi}{3}\PY{p}{)}\PY{p}{;}\PY{+w}{              }\PY{c+c1}{// hashCode changes!}
\PY{n}{set}\PY{p}{.}\PY{n+na}{contains}\PY{p}{(}\PY{n}{key}\PY{p}{)}\PY{p}{;}\PY{+w}{       }\PY{c+c1}{// often false \PYZhy{}\PYZhy{} looked up in the wrong bucket}

\PY{c+c1}{// After: use immutable keys (e.g., List.copyOf, records, Strings)}
\PY{n}{List}\PY{o}{\PYZlt{}}\PY{n}{Integer}\PY{o}{\PYZgt{}}\PY{+w}{ }\PY{n}{key}\PY{+w}{ }\PY{o}{=}\PY{+w}{ }\PY{n}{List}\PY{p}{.}\PY{n+na}{of}\PY{p}{(}\PY{l+m+mi}{1}\PY{p}{,}\PY{+w}{ }\PY{l+m+mi}{2}\PY{p}{)}\PY{p}{;}
\end{Verbatim}
\end{codebox}

\item 
\textbf{Assuming \code{List.\allowbreak{}remove(\allowbreak{}int)} and \code{List.\allowbreak{}remove(\allowbreak{}Object)} do the same thing with Integer arguments.}
Why it matters: \code{list.\allowbreak{}remove(\allowbreak{}2)} on a \code{List\textless{}\allowbreak{}Integer\textgreater{}} calls the \code{int} overload (removes by index), not the \code{Object} overload (removes by value) — a classic autoboxing trap.

\begin{codebox}
\begin{Verbatim}[commandchars=\\\{\}]
\PY{n}{List}\PY{o}{\PYZlt{}}\PY{n}{Integer}\PY{o}{\PYZgt{}}\PY{+w}{ }\PY{n}{nums}\PY{+w}{ }\PY{o}{=}\PY{+w}{ }\PY{k}{new}\PY{+w}{ }\PY{n}{ArrayList}\PY{o}{\PYZlt{}}\PY{o}{\PYZgt{}}\PY{p}{(}\PY{n}{List}\PY{p}{.}\PY{n+na}{of}\PY{p}{(}\PY{l+m+mi}{10}\PY{p}{,}\PY{+w}{ }\PY{l+m+mi}{20}\PY{p}{,}\PY{+w}{ }\PY{l+m+mi}{30}\PY{p}{)}\PY{p}{)}\PY{p}{;}
\PY{n}{nums}\PY{p}{.}\PY{n+na}{remove}\PY{p}{(}\PY{l+m+mi}{2}\PY{p}{)}\PY{p}{;}\PY{+w}{                  }\PY{c+c1}{// removes INDEX 2 \PYZhy{}\PYZgt{} [10, 20]}
\PY{n}{nums}\PY{p}{.}\PY{n+na}{remove}\PY{p}{(}\PY{n}{Integer}\PY{p}{.}\PY{n+na}{valueOf}\PY{p}{(}\PY{l+m+mi}{20}\PY{p}{)}\PY{p}{)}\PY{p}{;}\PY{+w}{ }\PY{c+c1}{// removes VALUE 20 \PYZhy{}\PYZgt{} [10]}
\end{Verbatim}
\end{codebox}

\item 
\textbf{Not overriding \code{equals()}/\code{hashCode()} on custom keys/elements.}
Why it matters: without them, \code{HashMap}/\code{HashSet} fall back to identity comparison, so logically-equal objects are treated as distinct.

\begin{codebox}
\begin{Verbatim}[commandchars=\\\{\}]
\PY{c+c1}{// Before: no equals()/hashCode() \PYZhy{}\PYZgt{} identity semantics}
\PY{k+kd}{class} \PY{n+nc}{Point}\PY{+w}{ }\PY{p}{\PYZob{}}\PY{+w}{ }\PY{k+kt}{int}\PY{+w}{ }\PY{n}{x}\PY{p}{,}\PY{+w}{ }\PY{n}{y}\PY{p}{;}\PY{+w}{ }\PY{p}{\PYZcb{}}

\PY{c+c1}{// After: value semantics, works correctly as a map key / set element}
\PY{k+kd}{record} \PY{n+nc}{Point}\PY{p}{(}\PY{k+kt}{int}\PY{+w}{ }\PY{n}{x}\PY{p}{,}\PY{+w}{ }\PY{k+kt}{int}\PY{+w}{ }\PY{n}{y}\PY{p}{)}\PY{+w}{ }\PY{p}{\PYZob{}}\PY{p}{\PYZcb{}}
\end{Verbatim}
\end{codebox}

\item 
\textbf{Choosing \code{LinkedList} as a default \code{List} implementation.}
Why it matters: \code{get(\allowbreak{}index)} is O(n) on \code{LinkedList}, and it has higher memory overhead per element (node objects with two pointers each) than \code{ArrayList}’s flat array.

\begin{codebox}
\begin{Verbatim}[commandchars=\\\{\}]
\PY{c+c1}{// Before: slow random access}
\PY{n}{List}\PY{o}{\PYZlt{}}\PY{n}{String}\PY{o}{\PYZgt{}}\PY{+w}{ }\PY{n}{list}\PY{+w}{ }\PY{o}{=}\PY{+w}{ }\PY{k}{new}\PY{+w}{ }\PY{n}{LinkedList}\PY{o}{\PYZlt{}}\PY{o}{\PYZgt{}}\PY{p}{(}\PY{p}{)}\PY{p}{;}

\PY{c+c1}{// After: ArrayList is the right default unless you specifically need}
\PY{c+c1}{// O(1) insert/remove at both ends combined with List semantics}
\PY{n}{List}\PY{o}{\PYZlt{}}\PY{n}{String}\PY{o}{\PYZgt{}}\PY{+w}{ }\PY{n}{list}\PY{+w}{ }\PY{o}{=}\PY{+w}{ }\PY{k}{new}\PY{+w}{ }\PY{n}{ArrayList}\PY{o}{\PYZlt{}}\PY{o}{\PYZgt{}}\PY{p}{(}\PY{p}{)}\PY{p}{;}
\end{Verbatim}
\end{codebox}

\item 
\textbf{Using \code{Vector} or \code{Stack} in new code.}
Why it matters: both are legacy, synchronized classes with unnecessary locking overhead for single-threaded use, and \code{Stack} extends \code{Vector} (an odd, dated design).

\begin{codebox}
\begin{Verbatim}[commandchars=\\\{\}]
\PY{c+c1}{// Before: legacy, synchronized, slower}
\PY{n}{Stack}\PY{o}{\PYZlt{}}\PY{n}{Integer}\PY{o}{\PYZgt{}}\PY{+w}{ }\PY{n}{stack}\PY{+w}{ }\PY{o}{=}\PY{+w}{ }\PY{k}{new}\PY{+w}{ }\PY{n}{Stack}\PY{o}{\PYZlt{}}\PY{o}{\PYZgt{}}\PY{p}{(}\PY{p}{)}\PY{p}{;}

\PY{c+c1}{// After: modern, faster, unsynchronized}
\PY{n}{Deque}\PY{o}{\PYZlt{}}\PY{n}{Integer}\PY{o}{\PYZgt{}}\PY{+w}{ }\PY{n}{stack}\PY{+w}{ }\PY{o}{=}\PY{+w}{ }\PY{k}{new}\PY{+w}{ }\PY{n}{ArrayDeque}\PY{o}{\PYZlt{}}\PY{o}{\PYZgt{}}\PY{p}{(}\PY{p}{)}\PY{p}{;}
\end{Verbatim}
\end{codebox}

\item 
\textbf{Passing an immutable collection (\code{List.\allowbreak{}of(...)}) somewhere it will be mutated later.}
Why it matters: it throws \code{Unsupported\allowbreak{}Operation\allowbreak{}Exception} at runtime, often far from where the immutable list was created, making the bug hard to trace.

\begin{codebox}
\begin{Verbatim}[commandchars=\\\{\}]
\PY{c+c1}{// Before: crashes deep inside some other method}
\PY{n}{List}\PY{o}{\PYZlt{}}\PY{n}{String}\PY{o}{\PYZgt{}}\PY{+w}{ }\PY{n}{config}\PY{+w}{ }\PY{o}{=}\PY{+w}{ }\PY{n}{List}\PY{p}{.}\PY{n+na}{of}\PY{p}{(}\PY{l+s}{\PYZdq{}}\PY{l+s}{a}\PY{l+s}{\PYZdq{}}\PY{p}{,}\PY{+w}{ }\PY{l+s}{\PYZdq{}}\PY{l+s}{b}\PY{l+s}{\PYZdq{}}\PY{p}{)}\PY{p}{;}
\PY{n}{someLegacyCodeThatAddsDefaults}\PY{p}{(}\PY{n}{config}\PY{p}{)}\PY{p}{;}\PY{+w}{ }\PY{c+c1}{// throws UnsupportedOperationException}

\PY{c+c1}{// After: wrap in a mutable copy if the caller will mutate it}
\PY{n}{List}\PY{o}{\PYZlt{}}\PY{n}{String}\PY{o}{\PYZgt{}}\PY{+w}{ }\PY{n}{config}\PY{+w}{ }\PY{o}{=}\PY{+w}{ }\PY{k}{new}\PY{+w}{ }\PY{n}{ArrayList}\PY{o}{\PYZlt{}}\PY{o}{\PYZgt{}}\PY{p}{(}\PY{n}{List}\PY{p}{.}\PY{n+na}{of}\PY{p}{(}\PY{l+s}{\PYZdq{}}\PY{l+s}{a}\PY{l+s}{\PYZdq{}}\PY{p}{,}\PY{+w}{ }\PY{l+s}{\PYZdq{}}\PY{l+s}{b}\PY{l+s}{\PYZdq{}}\PY{p}{)}\PY{p}{)}\PY{p}{;}
\end{Verbatim}
\end{codebox}

\item 
\textbf{Assuming \code{Collections.\allowbreak{}synchronized\allowbreak{}List}/\code{synchronizedMap} makes compound operations atomic.}
Why it matters: each individual call is synchronized, but a sequence like check-then-add is still a race condition across threads.

\begin{codebox}
\begin{Verbatim}[commandchars=\\\{\}]
\PY{c+c1}{// Before: race condition between contains() and add()}
\PY{k}{if}\PY{+w}{ }\PY{p}{(}\PY{o}{!}\PY{n}{syncList}\PY{p}{.}\PY{n+na}{contains}\PY{p}{(}\PY{n}{x}\PY{p}{)}\PY{p}{)}\PY{+w}{ }\PY{n}{syncList}\PY{p}{.}\PY{n+na}{add}\PY{p}{(}\PY{n}{x}\PY{p}{)}\PY{p}{;}

\PY{c+c1}{// After: synchronize the whole compound action, or use a concurrent collection}
\PY{k+kd}{synchronized}\PY{+w}{ }\PY{p}{(}\PY{n}{syncList}\PY{p}{)}\PY{+w}{ }\PY{p}{\PYZob{}}
\PY{+w}{    }\PY{k}{if}\PY{+w}{ }\PY{p}{(}\PY{o}{!}\PY{n}{syncList}\PY{p}{.}\PY{n+na}{contains}\PY{p}{(}\PY{n}{x}\PY{p}{)}\PY{p}{)}\PY{+w}{ }\PY{n}{syncList}\PY{p}{.}\PY{n+na}{add}\PY{p}{(}\PY{n}{x}\PY{p}{)}\PY{p}{;}
\PY{p}{\PYZcb{}}
\end{Verbatim}
\end{codebox}

\item 
\textbf{Using \code{HashMap}/\code{HashSet} from multiple threads without synchronization.}
Why it matters: concurrent structural modification can corrupt the internal bucket structure, causing silent data loss, infinite loops (older JDKs), or \code{Concurrent\allowbreak{}Modification\allowbreak{}Exception}.

\begin{codebox}
\begin{Verbatim}[commandchars=\\\{\}]
\PY{c+c1}{// Before: not thread\PYZhy{}safe}
\PY{n}{Map}\PY{o}{\PYZlt{}}\PY{n}{String}\PY{p}{,}\PY{+w}{ }\PY{n}{Integer}\PY{o}{\PYZgt{}}\PY{+w}{ }\PY{n}{counters}\PY{+w}{ }\PY{o}{=}\PY{+w}{ }\PY{k}{new}\PY{+w}{ }\PY{n}{HashMap}\PY{o}{\PYZlt{}}\PY{o}{\PYZgt{}}\PY{p}{(}\PY{p}{)}\PY{p}{;}

\PY{c+c1}{// After: use a collection designed for concurrency}
\PY{n}{Map}\PY{o}{\PYZlt{}}\PY{n}{String}\PY{p}{,}\PY{+w}{ }\PY{n}{Integer}\PY{o}{\PYZgt{}}\PY{+w}{ }\PY{n}{counters}\PY{+w}{ }\PY{o}{=}\PY{+w}{ }\PY{k}{new}\PY{+w}{ }\PY{n}{ConcurrentHashMap}\PY{o}{\PYZlt{}}\PY{o}{\PYZgt{}}\PY{p}{(}\PY{p}{)}\PY{p}{;}
\end{Verbatim}
\end{codebox}

\item 
\textbf{Putting \code{null} into a \code{ConcurrentHashMap} or a \code{TreeMap}/\code{TreeSet}.}
Why it matters: \code{ConcurrentHashMap} rejects null keys/values outright (ambiguity with “absent”); \code{TreeMap}/\code{TreeSet} need to \emph{compare} keys, and comparing against \code{null} throws \code{NullPointerException}.

\begin{codebox}
\begin{Verbatim}[commandchars=\\\{\}]
\PY{c+c1}{// Before: throws NullPointerException at runtime}
\PY{k}{new}\PY{+w}{ }\PY{n}{ConcurrentHashMap}\PY{o}{\PYZlt{}}\PY{n}{String}\PY{p}{,}\PY{+w}{ }\PY{n}{String}\PY{o}{\PYZgt{}}\PY{p}{(}\PY{p}{)}\PY{p}{.}\PY{n+na}{put}\PY{p}{(}\PY{l+s}{\PYZdq{}}\PY{l+s}{key}\PY{l+s}{\PYZdq{}}\PY{p}{,}\PY{+w}{ }\PY{k+kc}{null}\PY{p}{)}\PY{p}{;}
\PY{k}{new}\PY{+w}{ }\PY{n}{TreeSet}\PY{o}{\PYZlt{}}\PY{n}{String}\PY{o}{\PYZgt{}}\PY{p}{(}\PY{p}{)}\PY{p}{.}\PY{n+na}{add}\PY{p}{(}\PY{k+kc}{null}\PY{p}{)}\PY{p}{;}

\PY{c+c1}{// After: use Optional, a sentinel value, or filter out nulls before inserting}
\end{Verbatim}
\end{codebox}

\item 
\textbf{Forgetting that \code{PriorityQueue} iteration order is not sorted order.}
Why it matters: only \code{peek()}/\code{poll()} guarantee the smallest element; iterating the queue (\code{for~\allowbreak{}(\allowbreak{}var~\allowbreak{}x~\allowbreak{}:\allowbreak{}~\allowbreak{}pq)}) or printing it shows heap-array order, which looks unsorted.

\begin{codebox}
\begin{Verbatim}[commandchars=\\\{\}]
\PY{n}{PriorityQueue}\PY{o}{\PYZlt{}}\PY{n}{Integer}\PY{o}{\PYZgt{}}\PY{+w}{ }\PY{n}{pq}\PY{+w}{ }\PY{o}{=}\PY{+w}{ }\PY{k}{new}\PY{+w}{ }\PY{n}{PriorityQueue}\PY{o}{\PYZlt{}}\PY{o}{\PYZgt{}}\PY{p}{(}\PY{n}{List}\PY{p}{.}\PY{n+na}{of}\PY{p}{(}\PY{l+m+mi}{5}\PY{p}{,}\PY{+w}{ }\PY{l+m+mi}{1}\PY{p}{,}\PY{+w}{ }\PY{l+m+mi}{3}\PY{p}{)}\PY{p}{)}\PY{p}{;}
\PY{n}{System}\PY{p}{.}\PY{n+na}{out}\PY{p}{.}\PY{n+na}{println}\PY{p}{(}\PY{n}{pq}\PY{p}{)}\PY{p}{;}\PY{+w}{ }\PY{c+c1}{// NOT guaranteed to print [1, 3, 5]}

\PY{c+c1}{// To get sorted output, drain with poll():}
\PY{n}{List}\PY{o}{\PYZlt{}}\PY{n}{Integer}\PY{o}{\PYZgt{}}\PY{+w}{ }\PY{n}{sorted}\PY{+w}{ }\PY{o}{=}\PY{+w}{ }\PY{k}{new}\PY{+w}{ }\PY{n}{ArrayList}\PY{o}{\PYZlt{}}\PY{o}{\PYZgt{}}\PY{p}{(}\PY{p}{)}\PY{p}{;}
\PY{k}{while}\PY{+w}{ }\PY{p}{(}\PY{o}{!}\PY{n}{pq}\PY{p}{.}\PY{n+na}{isEmpty}\PY{p}{(}\PY{p}{)}\PY{p}{)}\PY{+w}{ }\PY{n}{sorted}\PY{p}{.}\PY{n+na}{add}\PY{p}{(}\PY{n}{pq}\PY{p}{.}\PY{n+na}{poll}\PY{p}{(}\PY{p}{)}\PY{p}{)}\PY{p}{;}
\end{Verbatim}
\end{codebox}

\item 
\textbf{Using \code{Arrays.\allowbreak{}asList()} and expecting a fully mutable, independent list.}
Why it matters: it returns a fixed-size list backed by the array — \code{add}/\code{remove} throw \code{Unsupported\allowbreak{}Operation\allowbreak{}Exception}, and \code{set} writes through to the original array.

\begin{codebox}
\begin{Verbatim}[commandchars=\\\{\}]
\PY{c+c1}{// Before: surprising behavior}
\PY{n}{Integer}\PY{o}{[}\PY{o}{]}\PY{+w}{ }\PY{n}{arr}\PY{+w}{ }\PY{o}{=}\PY{+w}{ }\PY{p}{\PYZob{}}\PY{l+m+mi}{1}\PY{p}{,}\PY{+w}{ }\PY{l+m+mi}{2}\PY{p}{,}\PY{+w}{ }\PY{l+m+mi}{3}\PY{p}{\PYZcb{}}\PY{p}{;}
\PY{n}{List}\PY{o}{\PYZlt{}}\PY{n}{Integer}\PY{o}{\PYZgt{}}\PY{+w}{ }\PY{n}{list}\PY{+w}{ }\PY{o}{=}\PY{+w}{ }\PY{n}{Arrays}\PY{p}{.}\PY{n+na}{asList}\PY{p}{(}\PY{n}{arr}\PY{p}{)}\PY{p}{;}
\PY{n}{list}\PY{p}{.}\PY{n+na}{add}\PY{p}{(}\PY{l+m+mi}{4}\PY{p}{)}\PY{p}{;}\PY{+w}{ }\PY{c+c1}{// throws UnsupportedOperationException}

\PY{c+c1}{// After: wrap in a real, resizable, independent list}
\PY{n}{List}\PY{o}{\PYZlt{}}\PY{n}{Integer}\PY{o}{\PYZgt{}}\PY{+w}{ }\PY{n}{list}\PY{+w}{ }\PY{o}{=}\PY{+w}{ }\PY{k}{new}\PY{+w}{ }\PY{n}{ArrayList}\PY{o}{\PYZlt{}}\PY{o}{\PYZgt{}}\PY{p}{(}\PY{n}{Arrays}\PY{p}{.}\PY{n+na}{asList}\PY{p}{(}\PY{n}{arr}\PY{p}{)}\PY{p}{)}\PY{p}{;}
\end{Verbatim}
\end{codebox}

\item 
\textbf{Not pre-sizing collections when the final size is known.}
Why it matters: repeated resizing (\code{ArrayList} array copies, \code{HashMap} rehashing) wastes CPU and creates avoidable garbage under load.

\begin{codebox}
\begin{Verbatim}[commandchars=\\\{\}]
\PY{c+c1}{// Before: several resize/copy operations as the list grows}
\PY{n}{List}\PY{o}{\PYZlt{}}\PY{n}{String}\PY{o}{\PYZgt{}}\PY{+w}{ }\PY{n}{results}\PY{+w}{ }\PY{o}{=}\PY{+w}{ }\PY{k}{new}\PY{+w}{ }\PY{n}{ArrayList}\PY{o}{\PYZlt{}}\PY{o}{\PYZgt{}}\PY{p}{(}\PY{p}{)}\PY{p}{;}
\PY{k}{for}\PY{+w}{ }\PY{p}{(}\PY{k+kt}{int}\PY{+w}{ }\PY{n}{i}\PY{+w}{ }\PY{o}{=}\PY{+w}{ }\PY{l+m+mi}{0}\PY{p}{;}\PY{+w}{ }\PY{n}{i}\PY{+w}{ }\PY{o}{\PYZlt{}}\PY{+w}{ }\PY{l+m+mi}{100\PYZus{}000}\PY{p}{;}\PY{+w}{ }\PY{n}{i}\PY{o}{+}\PY{o}{+}\PY{p}{)}\PY{+w}{ }\PY{n}{results}\PY{p}{.}\PY{n+na}{add}\PY{p}{(}\PY{n}{fetch}\PY{p}{(}\PY{n}{i}\PY{p}{)}\PY{p}{)}\PY{p}{;}

\PY{c+c1}{// After: allocate once}
\PY{n}{List}\PY{o}{\PYZlt{}}\PY{n}{String}\PY{o}{\PYZgt{}}\PY{+w}{ }\PY{n}{results}\PY{+w}{ }\PY{o}{=}\PY{+w}{ }\PY{k}{new}\PY{+w}{ }\PY{n}{ArrayList}\PY{o}{\PYZlt{}}\PY{o}{\PYZgt{}}\PY{p}{(}\PY{l+m+mi}{100\PYZus{}000}\PY{p}{)}\PY{p}{;}
\end{Verbatim}
\end{codebox}

\item 
\textbf{Assuming \code{TreeMap}/\code{TreeSet} ordering matches \code{equals()}.}
Why it matters: \code{TreeSet}/\code{TreeMap} decide “duplicate” using the \code{Comparator}/\code{compareTo}, not \code{equals()}. Two elements that compare as \code{0} are treated as duplicates, even if \code{equals()} would say they differ.

\begin{codebox}
\begin{Verbatim}[commandchars=\\\{\}]
\PY{c+c1}{// Before: silently drops \PYZdq{}second\PYZdq{} because comparator says they\PYZsq{}re \PYZdq{}equal\PYZdq{}}
\PY{n}{TreeSet}\PY{o}{\PYZlt{}}\PY{n}{String}\PY{o}{\PYZgt{}}\PY{+w}{ }\PY{n}{set}\PY{+w}{ }\PY{o}{=}\PY{+w}{ }\PY{k}{new}\PY{+w}{ }\PY{n}{TreeSet}\PY{o}{\PYZlt{}}\PY{o}{\PYZgt{}}\PY{p}{(}\PY{n}{Comparator}\PY{p}{.}\PY{n+na}{comparingInt}\PY{p}{(}\PY{n}{String}\PY{p}{:}\PY{p}{:}\PY{n}{length}\PY{p}{)}\PY{p}{)}\PY{p}{;}
\PY{n}{set}\PY{p}{.}\PY{n+na}{add}\PY{p}{(}\PY{l+s}{\PYZdq{}}\PY{l+s}{cat}\PY{l+s}{\PYZdq{}}\PY{p}{)}\PY{p}{;}
\PY{n}{set}\PY{p}{.}\PY{n+na}{add}\PY{p}{(}\PY{l+s}{\PYZdq{}}\PY{l+s}{dog}\PY{l+s}{\PYZdq{}}\PY{p}{)}\PY{p}{;}\PY{+w}{ }\PY{c+c1}{// same length as \PYZdq{}cat\PYZdq{} \PYZhy{}\PYZgt{} compareTo returns 0 \PYZhy{}\PYZgt{} treated as duplicate, ignored!}
\PY{n}{System}\PY{p}{.}\PY{n+na}{out}\PY{p}{.}\PY{n+na}{println}\PY{p}{(}\PY{n}{set}\PY{p}{)}\PY{p}{;}\PY{+w}{ }\PY{c+c1}{// [cat]}
\end{Verbatim}
\end{codebox}

\item 
\textbf{Choosing \code{CopyOnWriteArrayList} for a write-heavy workload.}
Why it matters: every mutation copies the entire backing array — O(n) per write — which becomes a serious bottleneck as the list grows or writes become frequent.

\begin{codebox}
\begin{Verbatim}[commandchars=\\\{\}]
\PY{c+c1}{// Before: O(n) per add(), terrible for a write\PYZhy{}heavy queue\PYZhy{}like list}
\PY{n}{CopyOnWriteArrayList}\PY{o}{\PYZlt{}}\PY{n}{String}\PY{o}{\PYZgt{}}\PY{+w}{ }\PY{n}{events}\PY{+w}{ }\PY{o}{=}\PY{+w}{ }\PY{k}{new}\PY{+w}{ }\PY{n}{CopyOnWriteArrayList}\PY{o}{\PYZlt{}}\PY{o}{\PYZgt{}}\PY{p}{(}\PY{p}{)}\PY{p}{;}
\PY{k}{for}\PY{+w}{ }\PY{p}{(}\PY{k+kt}{int}\PY{+w}{ }\PY{n}{i}\PY{+w}{ }\PY{o}{=}\PY{+w}{ }\PY{l+m+mi}{0}\PY{p}{;}\PY{+w}{ }\PY{n}{i}\PY{+w}{ }\PY{o}{\PYZlt{}}\PY{+w}{ }\PY{l+m+mi}{1\PYZus{}000\PYZus{}000}\PY{p}{;}\PY{+w}{ }\PY{n}{i}\PY{o}{+}\PY{o}{+}\PY{p}{)}\PY{+w}{ }\PY{n}{events}\PY{p}{.}\PY{n+na}{add}\PY{p}{(}\PY{l+s}{\PYZdq{}}\PY{l+s}{event\PYZhy{}}\PY{l+s}{\PYZdq{}}\PY{+w}{ }\PY{o}{+}\PY{+w}{ }\PY{n}{i}\PY{p}{)}\PY{p}{;}

\PY{c+c1}{// After: use a concurrent queue built for frequent writes}
\PY{n}{ConcurrentLinkedQueue}\PY{o}{\PYZlt{}}\PY{n}{String}\PY{o}{\PYZgt{}}\PY{+w}{ }\PY{n}{events}\PY{+w}{ }\PY{o}{=}\PY{+w}{ }\PY{k}{new}\PY{+w}{ }\PY{n}{ConcurrentLinkedQueue}\PY{o}{\PYZlt{}}\PY{o}{\PYZgt{}}\PY{p}{(}\PY{p}{)}\PY{p}{;}
\end{Verbatim}
\end{codebox}

\item 
\textbf{Relying on \code{Map.\allowbreak{}get()} returning \code{null} to mean “absent” when \code{null} values are also stored.}
Why it matters: you can’t tell “key missing” apart from “key present with a null value” using \code{get()} alone, which leads to subtle bugs.

\begin{codebox}
\begin{Verbatim}[commandchars=\\\{\}]
\PY{c+c1}{// Before: ambiguous}
\PY{k}{if}\PY{+w}{ }\PY{p}{(}\PY{n}{map}\PY{p}{.}\PY{n+na}{get}\PY{p}{(}\PY{l+s}{\PYZdq{}}\PY{l+s}{key}\PY{l+s}{\PYZdq{}}\PY{p}{)}\PY{+w}{ }\PY{o}{=}\PY{o}{=}\PY{+w}{ }\PY{k+kc}{null}\PY{p}{)}\PY{+w}{ }\PY{p}{\PYZob{}}\PY{+w}{ }\PY{c+cm}{/* missing, or present\PYZhy{}with\PYZhy{}null? unclear */}\PY{+w}{ }\PY{p}{\PYZcb{}}

\PY{c+c1}{// After: use containsKey, or getOrDefault with a real sentinel, and avoid}
\PY{c+c1}{// storing null values in the first place when possible}
\PY{k}{if}\PY{+w}{ }\PY{p}{(}\PY{o}{!}\PY{n}{map}\PY{p}{.}\PY{n+na}{containsKey}\PY{p}{(}\PY{l+s}{\PYZdq{}}\PY{l+s}{key}\PY{l+s}{\PYZdq{}}\PY{p}{)}\PY{p}{)}\PY{+w}{ }\PY{p}{\PYZob{}}\PY{+w}{ }\PY{c+cm}{/* definitely missing */}\PY{+w}{ }\PY{p}{\PYZcb{}}
\end{Verbatim}
\end{codebox}

\item 
\textbf{Using a plain \code{Iterator.\allowbreak{}remove()} loop when \code{removeIf} says exactly what you mean.}
Why it matters: not a correctness bug, but a clarity/maintainability one that reviewers flag — \code{removeIf} is shorter, less error-prone (no risk of forgetting to call \code{it.\allowbreak{}remove()} and hitting \code{Concurrent\allowbreak{}Modification\allowbreak{}Exception}), and communicates intent directly.

\begin{codebox}
\begin{Verbatim}[commandchars=\\\{\}]
\PY{c+c1}{// Before: verbose, easy to get wrong}
\PY{n}{Iterator}\PY{o}{\PYZlt{}}\PY{n}{String}\PY{o}{\PYZgt{}}\PY{+w}{ }\PY{n}{it}\PY{+w}{ }\PY{o}{=}\PY{+w}{ }\PY{n}{names}\PY{p}{.}\PY{n+na}{iterator}\PY{p}{(}\PY{p}{)}\PY{p}{;}
\PY{k}{while}\PY{+w}{ }\PY{p}{(}\PY{n}{it}\PY{p}{.}\PY{n+na}{hasNext}\PY{p}{(}\PY{p}{)}\PY{p}{)}\PY{+w}{ }\PY{p}{\PYZob{}}
\PY{+w}{    }\PY{k}{if}\PY{+w}{ }\PY{p}{(}\PY{n}{it}\PY{p}{.}\PY{n+na}{next}\PY{p}{(}\PY{p}{)}\PY{p}{.}\PY{n+na}{isBlank}\PY{p}{(}\PY{p}{)}\PY{p}{)}\PY{+w}{ }\PY{n}{it}\PY{p}{.}\PY{n+na}{remove}\PY{p}{(}\PY{p}{)}\PY{p}{;}
\PY{p}{\PYZcb{}}

\PY{c+c1}{// After: same result, clearer intent}
\PY{n}{names}\PY{p}{.}\PY{n+na}{removeIf}\PY{p}{(}\PY{n}{String}\PY{p}{:}\PY{p}{:}\PY{n}{isBlank}\PY{p}{)}\PY{p}{;}
\end{Verbatim}
\end{codebox}

\end{enumerate}


\setcounter{chapter}{7}
\chapter{Functional Programming}
\label{chap:8}
\label{h:8:8-functional-programming}
Java added lambdas and the Streams API in Java 8, and it has kept building on them ever since. These features let you write code that describes \emph{what} you want to happen instead of \emph{how} to loop and mutate to get there. This style is called “functional programming” — you pass behavior around as values, instead of only passing data. This chapter is a fast, practical tour of the pieces you need to read and write this kind of code confidently, and to catch the subtle mistakes reviewers are expected to flag.

All examples target Java 21+.

\section[Functional Interfaces]{Functional Interfaces}
\label{h:8:functional-interfaces}
A \textbf{functional interface} is an interface with exactly one abstract method (often called a “SAM” — single abstract method). It can have any number of \code{default} or \code{static} methods too — those don’t count. Java uses functional interfaces as the “shape” that a lambda expression or method reference plugs into.

The \code{@\allowbreak{}FunctionalInterface} annotation is optional, but you should always add it to interfaces you design for this purpose. It doesn’t change runtime behavior — it just tells the compiler “this interface must have exactly one abstract method.” If someone later adds a second abstract method, the build fails immediately instead of breaking lambdas elsewhere in the codebase at a mysterious call site.

\begin{codebox}
\begin{Verbatim}[commandchars=\\\{\}]
\PY{n+nd}{@FunctionalInterface}
\PY{k+kd}{interface} \PY{n+nc}{Validator}\PY{o}{\PYZlt{}}\PY{n}{T}\PY{o}{\PYZgt{}}\PY{+w}{ }\PY{p}{\PYZob{}}
\PY{+w}{    }\PY{k+kt}{boolean}\PY{+w}{ }\PY{n+nf}{isValid}\PY{p}{(}\PY{n}{T}\PY{+w}{ }\PY{n}{value}\PY{p}{)}\PY{p}{;}

\PY{+w}{    }\PY{c+c1}{// default methods are allowed \PYZhy{} they don\PYZsq{}t break the \PYZdq{}single abstract method\PYZdq{} rule}
\PY{+w}{    }\PY{k}{default}\PY{+w}{ }\PY{n}{Validator}\PY{o}{\PYZlt{}}\PY{n}{T}\PY{o}{\PYZgt{}}\PY{+w}{ }\PY{n+nf}{negate}\PY{p}{(}\PY{p}{)}\PY{+w}{ }\PY{p}{\PYZob{}}
\PY{+w}{        }\PY{k}{return}\PY{+w}{ }\PY{n}{value}\PY{+w}{ }\PY{o}{\PYZhy{}}\PY{o}{\PYZgt{}}\PY{+w}{ }\PY{o}{!}\PY{n}{isValid}\PY{p}{(}\PY{n}{value}\PY{p}{)}\PY{p}{;}
\PY{+w}{    }\PY{p}{\PYZcb{}}
\PY{p}{\PYZcb{}}
\end{Verbatim}
\end{codebox}

If you try to add a second abstract method, the compiler rejects it:

\begin{codebox}
\begin{Verbatim}[commandchars=\\\{\}]
\PY{n+nd}{@FunctionalInterface}
\PY{k+kd}{interface} \PY{n+nc}{Broken}\PY{o}{\PYZlt{}}\PY{n}{T}\PY{o}{\PYZgt{}}\PY{+w}{ }\PY{p}{\PYZob{}}
\PY{+w}{    }\PY{k+kt}{boolean}\PY{+w}{ }\PY{n+nf}{isValid}\PY{p}{(}\PY{n}{T}\PY{+w}{ }\PY{n}{value}\PY{p}{)}\PY{p}{;}
\PY{+w}{    }\PY{k+kt}{boolean}\PY{+w}{ }\PY{n+nf}{isInvalid}\PY{p}{(}\PY{n}{T}\PY{+w}{ }\PY{n}{value}\PY{p}{)}\PY{p}{;}\PY{+w}{ }\PY{c+c1}{// compile error: Broken is not a functional interface}
\PY{p}{\PYZcb{}}
\end{Verbatim}
\end{codebox}

\subsection[The java.util.function catalogue]{The \code{java.\allowbreak{}util.\allowbreak{}function} catalogue}
\label{h:8:the-javautilfunction-catalogue}
Java ships a standard set of general-purpose functional interfaces in \code{java.\allowbreak{}util.\allowbreak{}function} so you rarely need to declare your own. Learning this table is worth it — reviewers expect you to reach for the right one instead of inventing a custom interface every time.

{\tablefont
\begin{xltabular}{\linewidth}{@{}L{0.740}L{0.431}L{0.740}L{2.089}@{}}
\toprule
\rowcolor{tablehdr}
\thd{Interface} & \thd{Abstract method} & \thd{Signature (conceptually)} & \thd{Typical use} \\
\midrule
\endhead
\code{Function\textless{}\allowbreak{}T,\allowbreak{}R\textgreater{}} & \code{apply} & \code{T~\allowbreak{}-\textgreater{}\allowbreak{}~\allowbreak{}R} & Transform a value \\
\code{BiFunction\textless{}\allowbreak{}T,\allowbreak{}U,\allowbreak{}R\textgreater{}} & \code{apply} & \code{(\allowbreak{}T,\allowbreak{}~\allowbreak{}U)\allowbreak{}~\allowbreak{}-\textgreater{}\allowbreak{}~\allowbreak{}R} & Combine two values into one \\
\code{Supplier\textless{}\allowbreak{}T\textgreater{}} & \code{get} & \code{()\allowbreak{}~\allowbreak{}-\textgreater{}\allowbreak{}~\allowbreak{}T} & Produce/lazily create a value \\
\code{Consumer\textless{}\allowbreak{}T\textgreater{}} & \code{accept} & \code{T~\allowbreak{}-\textgreater{}\allowbreak{}~\allowbreak{}void} & Do something with a value, no result \\
\code{BiConsumer\textless{}\allowbreak{}T,\allowbreak{}U\textgreater{}} & \code{accept} & \code{(\allowbreak{}T,\allowbreak{}~\allowbreak{}U)\allowbreak{}~\allowbreak{}-\textgreater{}\allowbreak{}~\allowbreak{}void} & Do something with two values \\
\code{Predicate\textless{}\allowbreak{}T\textgreater{}} & \code{test} & \code{T~\allowbreak{}-\textgreater{}\allowbreak{}~\allowbreak{}boolean} & Yes/no check on a value \\
\code{BiPredicate\textless{}\allowbreak{}T,\allowbreak{}U\textgreater{}} & \code{test} & \code{(\allowbreak{}T,\allowbreak{}~\allowbreak{}U)\allowbreak{}~\allowbreak{}-\textgreater{}\allowbreak{}~\allowbreak{}boolean} & Yes/no check on two values \\
\code{UnaryOperator\textless{}\allowbreak{}T\textgreater{}} & \code{apply} & \code{T~\allowbreak{}-\textgreater{}\allowbreak{}~\allowbreak{}T} & \code{Function\textless{}\allowbreak{}T,\allowbreak{}T\textgreater{}} specialization, same input/output type \\
\code{BinaryOperator\textless{}\allowbreak{}T\textgreater{}} & \code{apply} & \code{(\allowbreak{}T,\allowbreak{}~\allowbreak{}T)\allowbreak{}~\allowbreak{}-\textgreater{}\allowbreak{}~\allowbreak{}T} & \code{BiFunction\textless{}\allowbreak{}T,\allowbreak{}T,\allowbreak{}T\textgreater{}} specialization, used in \code{reduce} \\
\bottomrule
\end{xltabular}
}

Boxing \code{int}/\code{long}/\code{double} into \code{Integer}/\code{Long}/\code{Double} costs memory and CPU (see \hyperref[h:8:stream-performance]{Stream Performance}). To avoid it, the JDK also provides primitive specializations:

{\tablefont
\begin{xltabular}{\linewidth}{@{}L{1.457}L{0.543}@{}}
\toprule
\rowcolor{tablehdr}
\thd{Interface} & \thd{Purpose} \\
\midrule
\endhead
\code{IntFunction\textless{}\allowbreak{}R\textgreater{}}, \code{LongFunction\textless{}\allowbreak{}R\textgreater{}}, \code{DoubleFunction\textless{}\allowbreak{}R\textgreater{}} & \code{int/\allowbreak{}long/\allowbreak{}double~\allowbreak{}-\textgreater{}\allowbreak{}~\allowbreak{}R} \\
\code{ToIntFunction\textless{}\allowbreak{}T\textgreater{}}, \code{ToLongFunction\textless{}\allowbreak{}T\textgreater{}}, \code{ToDoubleFunction\textless{}\allowbreak{}T\textgreater{}} & \code{T~\allowbreak{}-\textgreater{}\allowbreak{}~\allowbreak{}int/\allowbreak{}long/\allowbreak{}double} \\
\code{IntUnaryOperator}, \code{LongUnaryOperator}, \code{DoubleUnaryOperator} & \code{int~\allowbreak{}-\textgreater{}\allowbreak{}~\allowbreak{}int}, etc. \\
\code{IntBinaryOperator}, \code{LongBinaryOperator}, \code{DoubleBinaryOperator} & \code{(\allowbreak{}int,\allowbreak{}~\allowbreak{}int)\allowbreak{}~\allowbreak{}-\textgreater{}\allowbreak{}~\allowbreak{}int}, etc. \\
\code{IntPredicate}, \code{LongPredicate}, \code{DoublePredicate} & \code{int~\allowbreak{}-\textgreater{}\allowbreak{}~\allowbreak{}boolean}, etc. \\
\code{IntConsumer}, \code{LongConsumer}, \code{DoubleConsumer} & \code{int~\allowbreak{}-\textgreater{}\allowbreak{}~\allowbreak{}void}, etc. \\
\code{IntSupplier}, \code{LongSupplier}, \code{DoubleSupplier} & \code{()\allowbreak{}~\allowbreak{}-\textgreater{}\allowbreak{}~\allowbreak{}int}, etc. \\
\code{ObjIntConsumer\textless{}\allowbreak{}T\textgreater{}}, \code{ObjLongConsumer\textless{}\allowbreak{}T\textgreater{}}, \code{ObjDoubleConsumer\textless{}\allowbreak{}T\textgreater{}} & \code{(\allowbreak{}T,\allowbreak{}~\allowbreak{}int)\allowbreak{}~\allowbreak{}-\textgreater{}\allowbreak{}~\allowbreak{}void}, etc. \\
\bottomrule
\end{xltabular}
}

\begin{codebox}
\begin{Verbatim}[commandchars=\\\{\}]
\PY{k+kn}{import}\PY{+w}{ }\PY{n+nn}{java.util.function.*}\PY{p}{;}

\PY{n}{Function}\PY{o}{\PYZlt{}}\PY{n}{String}\PY{p}{,}\PY{+w}{ }\PY{n}{Integer}\PY{o}{\PYZgt{}}\PY{+w}{ }\PY{n}{length}\PY{+w}{ }\PY{o}{=}\PY{+w}{ }\PY{n}{String}\PY{p}{:}\PY{p}{:}\PY{n}{length}\PY{p}{;}
\PY{n}{BiFunction}\PY{o}{\PYZlt{}}\PY{n}{Integer}\PY{p}{,}\PY{+w}{ }\PY{n}{Integer}\PY{p}{,}\PY{+w}{ }\PY{n}{Integer}\PY{o}{\PYZgt{}}\PY{+w}{ }\PY{n}{add}\PY{+w}{ }\PY{o}{=}\PY{+w}{ }\PY{n}{Integer}\PY{p}{:}\PY{p}{:}\PY{n}{sum}\PY{p}{;}
\PY{n}{Supplier}\PY{o}{\PYZlt{}}\PY{n}{List}\PY{o}{\PYZlt{}}\PY{n}{String}\PY{o}{\PYZgt{}}\PY{o}{\PYZgt{}}\PY{+w}{ }\PY{n}{newList}\PY{+w}{ }\PY{o}{=}\PY{+w}{ }\PY{n}{ArrayList}\PY{p}{:}\PY{p}{:}\PY{k}{new}\PY{p}{;}
\PY{n}{Consumer}\PY{o}{\PYZlt{}}\PY{n}{String}\PY{o}{\PYZgt{}}\PY{+w}{ }\PY{n}{print}\PY{+w}{ }\PY{o}{=}\PY{+w}{ }\PY{n}{System}\PY{p}{.}\PY{n+na}{out}\PY{p}{:}\PY{p}{:}\PY{n}{println}\PY{p}{;}
\PY{n}{Predicate}\PY{o}{\PYZlt{}}\PY{n}{String}\PY{o}{\PYZgt{}}\PY{+w}{ }\PY{n}{isBlank}\PY{+w}{ }\PY{o}{=}\PY{+w}{ }\PY{n}{String}\PY{p}{:}\PY{p}{:}\PY{n}{isBlank}\PY{p}{;}
\PY{n}{UnaryOperator}\PY{o}{\PYZlt{}}\PY{n}{String}\PY{o}{\PYZgt{}}\PY{+w}{ }\PY{n}{upper}\PY{+w}{ }\PY{o}{=}\PY{+w}{ }\PY{n}{String}\PY{p}{:}\PY{p}{:}\PY{n}{toUpperCase}\PY{p}{;}
\PY{n}{BinaryOperator}\PY{o}{\PYZlt{}}\PY{n}{Integer}\PY{o}{\PYZgt{}}\PY{+w}{ }\PY{n}{max}\PY{+w}{ }\PY{o}{=}\PY{+w}{ }\PY{n}{Integer}\PY{p}{:}\PY{p}{:}\PY{n}{max}\PY{p}{;}
\PY{n}{IntPredicate}\PY{+w}{ }\PY{n}{isEven}\PY{+w}{ }\PY{o}{=}\PY{+w}{ }\PY{n}{n}\PY{+w}{ }\PY{o}{\PYZhy{}}\PY{o}{\PYZgt{}}\PY{+w}{ }\PY{n}{n}\PY{+w}{ }\PY{o}{\PYZpc{}}\PY{+w}{ }\PY{l+m+mi}{2}\PY{+w}{ }\PY{o}{=}\PY{o}{=}\PY{+w}{ }\PY{l+m+mi}{0}\PY{p}{;}

\PY{n}{System}\PY{p}{.}\PY{n+na}{out}\PY{p}{.}\PY{n+na}{println}\PY{p}{(}\PY{n}{length}\PY{p}{.}\PY{n+na}{apply}\PY{p}{(}\PY{l+s}{\PYZdq{}}\PY{l+s}{hello}\PY{l+s}{\PYZdq{}}\PY{p}{)}\PY{p}{)}\PY{p}{;}\PY{+w}{   }\PY{c+c1}{// 5}
\PY{n}{System}\PY{p}{.}\PY{n+na}{out}\PY{p}{.}\PY{n+na}{println}\PY{p}{(}\PY{n}{add}\PY{p}{.}\PY{n+na}{apply}\PY{p}{(}\PY{l+m+mi}{2}\PY{p}{,}\PY{+w}{ }\PY{l+m+mi}{3}\PY{p}{)}\PY{p}{)}\PY{p}{;}\PY{+w}{         }\PY{c+c1}{// 5}
\PY{n}{System}\PY{p}{.}\PY{n+na}{out}\PY{p}{.}\PY{n+na}{println}\PY{p}{(}\PY{n}{isEven}\PY{p}{.}\PY{n+na}{test}\PY{p}{(}\PY{l+m+mi}{4}\PY{p}{)}\PY{p}{)}\PY{p}{;}\PY{+w}{          }\PY{c+c1}{// true}
\end{Verbatim}
\end{codebox}

\subsection[Writing your own]{Writing your own}
\label{h:8:writing-your-own}
Reach for a custom functional interface only when none of the standard ones fit — usually because you need a \emph{named}, self-documenting method, or you need to throw a checked exception.

\begin{codebox}
\begin{Verbatim}[commandchars=\\\{\}]
\PY{n+nd}{@FunctionalInterface}
\PY{k+kd}{interface} \PY{n+nc}{RowMapper}\PY{o}{\PYZlt{}}\PY{n}{T}\PY{o}{\PYZgt{}}\PY{+w}{ }\PY{p}{\PYZob{}}
\PY{+w}{    }\PY{n}{T}\PY{+w}{ }\PY{n+nf}{map}\PY{p}{(}\PY{n}{ResultSet}\PY{+w}{ }\PY{n}{row}\PY{p}{)}\PY{+w}{ }\PY{k+kd}{throws}\PY{+w}{ }\PY{n}{SQLException}\PY{p}{;}\PY{+w}{ }\PY{c+c1}{// checked exception allowed here}
\PY{p}{\PYZcb{}}

\PY{n}{RowMapper}\PY{o}{\PYZlt{}}\PY{n}{String}\PY{o}{\PYZgt{}}\PY{+w}{ }\PY{n}{nameMapper}\PY{+w}{ }\PY{o}{=}\PY{+w}{ }\PY{n}{row}\PY{+w}{ }\PY{o}{\PYZhy{}}\PY{o}{\PYZgt{}}\PY{+w}{ }\PY{n}{row}\PY{p}{.}\PY{n+na}{getString}\PY{p}{(}\PY{l+s}{\PYZdq{}}\PY{l+s}{name}\PY{l+s}{\PYZdq{}}\PY{p}{)}\PY{p}{;}
\end{Verbatim}
\end{codebox}

\section[Lambda Expressions]{Lambda Expressions}
\label{h:8:lambda-expressions}
A \textbf{lambda expression} is an anonymous, inline implementation of a functional interface’s single method. Syntax ranges from terse to explicit:

\begin{codebox}
\begin{Verbatim}[commandchars=\\\{\}]
\PY{n}{Runnable}\PY{+w}{ }\PY{n}{r1}\PY{+w}{ }\PY{o}{=}\PY{+w}{ }\PY{p}{(}\PY{p}{)}\PY{+w}{ }\PY{o}{\PYZhy{}}\PY{o}{\PYZgt{}}\PY{+w}{ }\PY{n}{System}\PY{p}{.}\PY{n+na}{out}\PY{p}{.}\PY{n+na}{println}\PY{p}{(}\PY{l+s}{\PYZdq{}}\PY{l+s}{run}\PY{l+s}{\PYZdq{}}\PY{p}{)}\PY{p}{;}\PY{+w}{                 }\PY{c+c1}{// no params}
\PY{n}{Function}\PY{o}{\PYZlt{}}\PY{n}{Integer}\PY{p}{,}\PY{+w}{ }\PY{n}{Integer}\PY{o}{\PYZgt{}}\PY{+w}{ }\PY{n}{square}\PY{+w}{ }\PY{o}{=}\PY{+w}{ }\PY{n}{x}\PY{+w}{ }\PY{o}{\PYZhy{}}\PY{o}{\PYZgt{}}\PY{+w}{ }\PY{n}{x}\PY{+w}{ }\PY{o}{*}\PY{+w}{ }\PY{n}{x}\PY{p}{;}\PY{+w}{                }\PY{c+c1}{// one param, inferred type}
\PY{n}{Function}\PY{o}{\PYZlt{}}\PY{n}{Integer}\PY{p}{,}\PY{+w}{ }\PY{n}{Integer}\PY{o}{\PYZgt{}}\PY{+w}{ }\PY{n}{square2}\PY{+w}{ }\PY{o}{=}\PY{+w}{ }\PY{p}{(}\PY{n}{Integer}\PY{+w}{ }\PY{n}{x}\PY{p}{)}\PY{+w}{ }\PY{o}{\PYZhy{}}\PY{o}{\PYZgt{}}\PY{+w}{ }\PY{n}{x}\PY{+w}{ }\PY{o}{*}\PY{+w}{ }\PY{n}{x}\PY{p}{;}\PY{+w}{      }\PY{c+c1}{// one param, explicit type}
\PY{n}{BiFunction}\PY{o}{\PYZlt{}}\PY{n}{Integer}\PY{p}{,}\PY{+w}{ }\PY{n}{Integer}\PY{p}{,}\PY{+w}{ }\PY{n}{Integer}\PY{o}{\PYZgt{}}\PY{+w}{ }\PY{n}{add}\PY{+w}{ }\PY{o}{=}\PY{+w}{ }\PY{p}{(}\PY{n}{a}\PY{p}{,}\PY{+w}{ }\PY{n}{b}\PY{p}{)}\PY{+w}{ }\PY{o}{\PYZhy{}}\PY{o}{\PYZgt{}}\PY{+w}{ }\PY{n}{a}\PY{+w}{ }\PY{o}{+}\PY{+w}{ }\PY{n}{b}\PY{p}{;}\PY{+w}{    }\PY{c+c1}{// two params, block\PYZhy{}less}
\PY{n}{BiFunction}\PY{o}{\PYZlt{}}\PY{n}{Integer}\PY{p}{,}\PY{+w}{ }\PY{n}{Integer}\PY{p}{,}\PY{+w}{ }\PY{n}{Integer}\PY{o}{\PYZgt{}}\PY{+w}{ }\PY{n}{addBlock}\PY{+w}{ }\PY{o}{=}\PY{+w}{ }\PY{p}{(}\PY{n}{a}\PY{p}{,}\PY{+w}{ }\PY{n}{b}\PY{p}{)}\PY{+w}{ }\PY{o}{\PYZhy{}}\PY{o}{\PYZgt{}}\PY{+w}{ }\PY{p}{\PYZob{}}\PY{+w}{    }\PY{c+c1}{// block body, needs return}
\PY{+w}{    }\PY{k+kt}{int}\PY{+w}{ }\PY{n}{sum}\PY{+w}{ }\PY{o}{=}\PY{+w}{ }\PY{n}{a}\PY{+w}{ }\PY{o}{+}\PY{+w}{ }\PY{n}{b}\PY{p}{;}
\PY{+w}{    }\PY{k}{return}\PY{+w}{ }\PY{n}{sum}\PY{p}{;}
\PY{p}{\PYZcb{}}\PY{p}{;}
\end{Verbatim}
\end{codebox}

Rules of thumb:

\begin{itemize}
\item 
Parentheses around a single inferred parameter are optional: \code{x~\allowbreak{}-\textgreater{}\allowbreak{}~\allowbreak{}x~\allowbreak{}*\allowbreak{}~\allowbreak{}x} and \code{(\allowbreak{}x)\allowbreak{}~\allowbreak{}-\textgreater{}\allowbreak{}~\allowbreak{}x~\allowbreak{}*\allowbreak{}~\allowbreak{}x} are the same.

\item 
If you type one parameter, type them all (\code{(\allowbreak{}a,\allowbreak{}~\allowbreak{}b)} or \code{(\allowbreak{}int~\allowbreak{}a,\allowbreak{}~\allowbreak{}int~\allowbreak{}b)}, not \code{(\allowbreak{}a,\allowbreak{}~\allowbreak{}int~\allowbreak{}b)}).

\item 
A block body (\code{\{\allowbreak{}~\allowbreak{}...\allowbreak{}~\allowbreak{}\}}) needs an explicit \code{return} (unless the return type is \code{void}); an expression body does not.

\end{itemize}

\subsection[Effectively-final capture]{Effectively-final capture}
\label{h:8:effectively-final-capture}
A lambda can read local variables from its enclosing scope, but only if those variables are \textbf{effectively final} — assigned exactly once, even without the \code{final} keyword. This is called “capturing” a variable.

\begin{codebox}
\begin{Verbatim}[commandchars=\\\{\}]
\PY{k+kt}{int}\PY{+w}{ }\PY{n}{factor}\PY{+w}{ }\PY{o}{=}\PY{+w}{ }\PY{l+m+mi}{10}\PY{p}{;}\PY{+w}{ }\PY{c+c1}{// effectively final: never reassigned}
\PY{n}{Function}\PY{o}{\PYZlt{}}\PY{n}{Integer}\PY{p}{,}\PY{+w}{ }\PY{n}{Integer}\PY{o}{\PYZgt{}}\PY{+w}{ }\PY{n}{scale}\PY{+w}{ }\PY{o}{=}\PY{+w}{ }\PY{n}{x}\PY{+w}{ }\PY{o}{\PYZhy{}}\PY{o}{\PYZgt{}}\PY{+w}{ }\PY{n}{x}\PY{+w}{ }\PY{o}{*}\PY{+w}{ }\PY{n}{factor}\PY{p}{;}\PY{+w}{ }\PY{c+c1}{// OK, captures factor}

\PY{k+kt}{int}\PY{+w}{ }\PY{n}{counter}\PY{+w}{ }\PY{o}{=}\PY{+w}{ }\PY{l+m+mi}{0}\PY{p}{;}
\PY{n}{Runnable}\PY{+w}{ }\PY{n}{bad}\PY{+w}{ }\PY{o}{=}\PY{+w}{ }\PY{p}{(}\PY{p}{)}\PY{+w}{ }\PY{o}{\PYZhy{}}\PY{o}{\PYZgt{}}\PY{+w}{ }\PY{p}{\PYZob{}}
\PY{+w}{    }\PY{c+c1}{// counter++; // compile error: counter is not effectively final}
\PY{p}{\PYZcb{}}\PY{p}{;}
\end{Verbatim}
\end{codebox}

Why the restriction exists: the lambda might run later, or on another thread. Java captures a \emph{snapshot} of the variable’s value (for primitives/references) rather than a live reference to the variable itself, so it forbids anything that could make that snapshot stale or ambiguous. If you need a mutable counter across calls, wrap it in a container:

\begin{codebox}
\begin{Verbatim}[commandchars=\\\{\}]
\PY{k+kt}{int}\PY{o}{[}\PY{o}{]}\PY{+w}{ }\PY{n}{counter}\PY{+w}{ }\PY{o}{=}\PY{+w}{ }\PY{p}{\PYZob{}}\PY{l+m+mi}{0}\PY{p}{\PYZcb{}}\PY{p}{;}
\PY{n}{Runnable}\PY{+w}{ }\PY{n}{increment}\PY{+w}{ }\PY{o}{=}\PY{+w}{ }\PY{p}{(}\PY{p}{)}\PY{+w}{ }\PY{o}{\PYZhy{}}\PY{o}{\PYZgt{}}\PY{+w}{ }\PY{n}{counter}\PY{o}{[}\PY{l+m+mi}{0}\PY{o}{]}\PY{o}{+}\PY{o}{+}\PY{p}{;}\PY{+w}{ }\PY{c+c1}{// OK: counter (the array reference) never changes}
\PY{n}{increment}\PY{p}{.}\PY{n+na}{run}\PY{p}{(}\PY{p}{)}\PY{p}{;}
\PY{n}{increment}\PY{p}{.}\PY{n+na}{run}\PY{p}{(}\PY{p}{)}\PY{p}{;}
\PY{n}{System}\PY{p}{.}\PY{n+na}{out}\PY{p}{.}\PY{n+na}{println}\PY{p}{(}\PY{n}{counter}\PY{o}{[}\PY{l+m+mi}{0}\PY{o}{]}\PY{p}{)}\PY{p}{;}\PY{+w}{ }\PY{c+c1}{// 2}
\end{Verbatim}
\end{codebox}

Fields (instance or static) are exempt from this rule — only \emph{local} variables and parameters must be effectively final.

\subsection[this semantics vs anonymous classes]{\code{this} semantics vs anonymous classes}
\label{h:8:this-semantics-vs-anonymous-classes}
A lambda does \textbf{not} introduce its own \code{this}. Inside a lambda, \code{this} refers to the enclosing instance — exactly as if the lambda’s body were pasted directly into the surrounding method. An anonymous inner class, by contrast, introduces a new \code{this} that refers to the anonymous class instance itself.

\begin{codebox}
\begin{Verbatim}[commandchars=\\\{\}]
\PY{k+kd}{class} \PY{n+nc}{Counter}\PY{+w}{ }\PY{p}{\PYZob{}}
\PY{+w}{    }\PY{k+kt}{int}\PY{+w}{ }\PY{n}{value}\PY{+w}{ }\PY{o}{=}\PY{+w}{ }\PY{l+m+mi}{0}\PY{p}{;}

\PY{+w}{    }\PY{n}{Runnable}\PY{+w}{ }\PY{n+nf}{lambdaVersion}\PY{p}{(}\PY{p}{)}\PY{+w}{ }\PY{p}{\PYZob{}}
\PY{+w}{        }\PY{k}{return}\PY{+w}{ }\PY{p}{(}\PY{p}{)}\PY{+w}{ }\PY{o}{\PYZhy{}}\PY{o}{\PYZgt{}}\PY{+w}{ }\PY{p}{\PYZob{}}
\PY{+w}{            }\PY{k}{this}\PY{p}{.}\PY{n+na}{value}\PY{o}{+}\PY{o}{+}\PY{p}{;}\PY{+w}{          }\PY{c+c1}{// \PYZsq{}this\PYZsq{} is the Counter instance}
\PY{+w}{            }\PY{n}{System}\PY{p}{.}\PY{n+na}{out}\PY{p}{.}\PY{n+na}{println}\PY{p}{(}\PY{k}{this}\PY{p}{.}\PY{n+na}{getClass}\PY{p}{(}\PY{p}{)}\PY{p}{)}\PY{p}{;}\PY{+w}{ }\PY{c+c1}{// class Counter}
\PY{+w}{        }\PY{p}{\PYZcb{}}\PY{p}{;}
\PY{+w}{    }\PY{p}{\PYZcb{}}

\PY{+w}{    }\PY{n}{Runnable}\PY{+w}{ }\PY{n+nf}{anonymousVersion}\PY{p}{(}\PY{p}{)}\PY{+w}{ }\PY{p}{\PYZob{}}
\PY{+w}{        }\PY{k}{return}\PY{+w}{ }\PY{k}{new}\PY{+w}{ }\PY{n}{Runnable}\PY{p}{(}\PY{p}{)}\PY{+w}{ }\PY{p}{\PYZob{}}
\PY{+w}{            }\PY{k+kt}{int}\PY{+w}{ }\PY{n}{value}\PY{+w}{ }\PY{o}{=}\PY{+w}{ }\PY{l+m+mi}{100}\PY{p}{;}\PY{+w}{ }\PY{c+c1}{// shadows Counter.value inside this class}
\PY{+w}{            }\PY{n+nd}{@Override}
\PY{+w}{            }\PY{k+kd}{public}\PY{+w}{ }\PY{k+kt}{void}\PY{+w}{ }\PY{n+nf}{run}\PY{p}{(}\PY{p}{)}\PY{+w}{ }\PY{p}{\PYZob{}}
\PY{+w}{                }\PY{k}{this}\PY{p}{.}\PY{n+na}{value}\PY{o}{+}\PY{o}{+}\PY{p}{;}\PY{+w}{      }\PY{c+c1}{// \PYZsq{}this\PYZsq{} is the anonymous Runnable}
\PY{+w}{                }\PY{n}{System}\PY{p}{.}\PY{n+na}{out}\PY{p}{.}\PY{n+na}{println}\PY{p}{(}\PY{k}{this}\PY{p}{.}\PY{n+na}{getClass}\PY{p}{(}\PY{p}{)}\PY{p}{)}\PY{p}{;}\PY{+w}{ }\PY{c+c1}{// class Counter\PYZdl{}1 (anonymous)}
\PY{+w}{            }\PY{p}{\PYZcb{}}
\PY{+w}{        }\PY{p}{\PYZcb{}}\PY{p}{;}
\PY{+w}{    }\PY{p}{\PYZcb{}}
\PY{p}{\PYZcb{}}
\end{Verbatim}
\end{codebox}

This matters in code review: if you convert an anonymous class to a lambda, double check any use of \code{this} — the meaning can change and silently start referring to a different object.

\subsection[No checked exceptions]{No checked exceptions}
\label{h:8:no-checked-exceptions}
None of the standard \code{java.\allowbreak{}util.\allowbreak{}function} interfaces declare \code{throws}. If the method you’re calling inside a lambda throws a checked exception, you must catch it, wrap it in an unchecked exception, or use a custom functional interface that declares \code{throws}.

\begin{codebox}
\begin{Verbatim}[commandchars=\\\{\}]
\PY{n}{List}\PY{o}{\PYZlt{}}\PY{n}{String}\PY{o}{\PYZgt{}}\PY{+w}{ }\PY{n}{paths}\PY{+w}{ }\PY{o}{=}\PY{+w}{ }\PY{n}{List}\PY{p}{.}\PY{n+na}{of}\PY{p}{(}\PY{l+s}{\PYZdq{}}\PY{l+s}{a.txt}\PY{l+s}{\PYZdq{}}\PY{p}{,}\PY{+w}{ }\PY{l+s}{\PYZdq{}}\PY{l+s}{b.txt}\PY{l+s}{\PYZdq{}}\PY{p}{)}\PY{p}{;}

\PY{c+c1}{// Function\PYZlt{}String, byte[]\PYZgt{} reader = Files::readAllBytes; // won\PYZsq{}t compile: readAllBytes throws IOException}

\PY{n}{Function}\PY{o}{\PYZlt{}}\PY{n}{String}\PY{p}{,}\PY{+w}{ }\PY{k+kt}{byte}\PY{o}{[}\PY{o}{]}\PY{o}{\PYZgt{}}\PY{+w}{ }\PY{n}{reader}\PY{+w}{ }\PY{o}{=}\PY{+w}{ }\PY{n}{path}\PY{+w}{ }\PY{o}{\PYZhy{}}\PY{o}{\PYZgt{}}\PY{+w}{ }\PY{p}{\PYZob{}}
\PY{+w}{    }\PY{k}{try}\PY{+w}{ }\PY{p}{\PYZob{}}
\PY{+w}{        }\PY{k}{return}\PY{+w}{ }\PY{n}{Files}\PY{p}{.}\PY{n+na}{readAllBytes}\PY{p}{(}\PY{n}{Path}\PY{p}{.}\PY{n+na}{of}\PY{p}{(}\PY{n}{path}\PY{p}{)}\PY{p}{)}\PY{p}{;}
\PY{+w}{    }\PY{p}{\PYZcb{}}\PY{+w}{ }\PY{k}{catch}\PY{+w}{ }\PY{p}{(}\PY{n}{IOException}\PY{+w}{ }\PY{n}{e}\PY{p}{)}\PY{+w}{ }\PY{p}{\PYZob{}}
\PY{+w}{        }\PY{k}{throw}\PY{+w}{ }\PY{k}{new}\PY{+w}{ }\PY{n}{UncheckedIOException}\PY{p}{(}\PY{n}{e}\PY{p}{)}\PY{p}{;}\PY{+w}{ }\PY{c+c1}{// wrap checked \PYZhy{}\PYZgt{} unchecked}
\PY{+w}{    }\PY{p}{\PYZcb{}}
\PY{p}{\PYZcb{}}\PY{p}{;}
\end{Verbatim}
\end{codebox}

\section[Method References]{Method References}
\label{h:8:method-references}
A \textbf{method reference} is shorthand for a lambda that just calls one existing method. It uses the \code{::} operator. It’s not a separate concept from lambdas — the compiler converts it to the same kind of functional-interface implementation. Use it whenever a lambda would do nothing but forward its argument(s) to an existing method.

{\tablefont
\begin{xltabular}{\linewidth}{@{}L{0.868}L{0.943}L{1.472}L{0.717}@{}}
\toprule
\rowcolor{tablehdr}
\thd{Kind} & \thd{Syntax} & \thd{Equivalent lambda} & \thd{Example} \\
\midrule
\endhead
Static method & \code{Class\allowbreak{}Name::\allowbreak{}static\allowbreak{}Method} & \code{(\allowbreak{}args)\allowbreak{}~\allowbreak{}-\textgreater{}\allowbreak{}~\allowbreak{}Class\allowbreak{}Name.\allowbreak{}static\allowbreak{}Method(\allowbreak{}args)} & \code{Integer::\allowbreak{}parseInt} \\
Bound instance method & \code{instance::\allowbreak{}method} & \code{(\allowbreak{}args)\allowbreak{}~\allowbreak{}-\textgreater{}\allowbreak{}~\allowbreak{}instance.\allowbreak{}method(\allowbreak{}args)} & \code{myList::\allowbreak{}add} \\
Unbound instance method & \code{Class\allowbreak{}Name::\allowbreak{}instance\allowbreak{}Method} & \code{(\allowbreak{}obj,\allowbreak{}~\allowbreak{}args)\allowbreak{}~\allowbreak{}-\textgreater{}\allowbreak{}~\allowbreak{}obj.\allowbreak{}instance\allowbreak{}Method(\allowbreak{}args)} & \code{String::\allowbreak{}toUpperCase} \\
Constructor & \code{ClassName::\allowbreak{}new} & \code{(\allowbreak{}args)\allowbreak{}~\allowbreak{}-\textgreater{}\allowbreak{}~\allowbreak{}new~\allowbreak{}Class\allowbreak{}Name(\allowbreak{}args)} & \code{ArrayList::\allowbreak{}new} \\
\bottomrule
\end{xltabular}
}

\begin{codebox}
\begin{Verbatim}[commandchars=\\\{\}]
\PY{c+c1}{// 1. Static method reference}
\PY{n}{Function}\PY{o}{\PYZlt{}}\PY{n}{String}\PY{p}{,}\PY{+w}{ }\PY{n}{Integer}\PY{o}{\PYZgt{}}\PY{+w}{ }\PY{n}{parse}\PY{+w}{ }\PY{o}{=}\PY{+w}{ }\PY{n}{Integer}\PY{p}{:}\PY{p}{:}\PY{n}{parseInt}\PY{p}{;}
\PY{n}{System}\PY{p}{.}\PY{n+na}{out}\PY{p}{.}\PY{n+na}{println}\PY{p}{(}\PY{n}{parse}\PY{p}{.}\PY{n+na}{apply}\PY{p}{(}\PY{l+s}{\PYZdq{}}\PY{l+s}{42}\PY{l+s}{\PYZdq{}}\PY{p}{)}\PY{p}{)}\PY{p}{;}\PY{+w}{ }\PY{c+c1}{// 42}

\PY{c+c1}{// 2. Bound instance method reference \PYZhy{} the receiver is a specific, already\PYZhy{}existing object}
\PY{n}{List}\PY{o}{\PYZlt{}}\PY{n}{String}\PY{o}{\PYZgt{}}\PY{+w}{ }\PY{n}{names}\PY{+w}{ }\PY{o}{=}\PY{+w}{ }\PY{k}{new}\PY{+w}{ }\PY{n}{ArrayList}\PY{o}{\PYZlt{}}\PY{o}{\PYZgt{}}\PY{p}{(}\PY{p}{)}\PY{p}{;}
\PY{n}{Consumer}\PY{o}{\PYZlt{}}\PY{n}{String}\PY{o}{\PYZgt{}}\PY{+w}{ }\PY{n}{addToNames}\PY{+w}{ }\PY{o}{=}\PY{+w}{ }\PY{n}{names}\PY{p}{:}\PY{p}{:}\PY{n}{add}\PY{p}{;}
\PY{n}{addToNames}\PY{p}{.}\PY{n+na}{accept}\PY{p}{(}\PY{l+s}{\PYZdq{}}\PY{l+s}{Ada}\PY{l+s}{\PYZdq{}}\PY{p}{)}\PY{p}{;}
\PY{n}{System}\PY{p}{.}\PY{n+na}{out}\PY{p}{.}\PY{n+na}{println}\PY{p}{(}\PY{n}{names}\PY{p}{)}\PY{p}{;}\PY{+w}{ }\PY{c+c1}{// [Ada]}

\PY{c+c1}{// 3. Unbound instance method reference \PYZhy{} the receiver becomes the first parameter}
\PY{n}{Function}\PY{o}{\PYZlt{}}\PY{n}{String}\PY{p}{,}\PY{+w}{ }\PY{n}{String}\PY{o}{\PYZgt{}}\PY{+w}{ }\PY{n}{toUpper}\PY{+w}{ }\PY{o}{=}\PY{+w}{ }\PY{n}{String}\PY{p}{:}\PY{p}{:}\PY{n}{toUpperCase}\PY{p}{;}
\PY{n}{System}\PY{p}{.}\PY{n+na}{out}\PY{p}{.}\PY{n+na}{println}\PY{p}{(}\PY{n}{toUpper}\PY{p}{.}\PY{n+na}{apply}\PY{p}{(}\PY{l+s}{\PYZdq{}}\PY{l+s}{hi}\PY{l+s}{\PYZdq{}}\PY{p}{)}\PY{p}{)}\PY{p}{;}\PY{+w}{ }\PY{c+c1}{// HI}

\PY{n}{BiFunction}\PY{o}{\PYZlt{}}\PY{n}{String}\PY{p}{,}\PY{+w}{ }\PY{n}{String}\PY{p}{,}\PY{+w}{ }\PY{n}{Boolean}\PY{o}{\PYZgt{}}\PY{+w}{ }\PY{n}{startsWith}\PY{+w}{ }\PY{o}{=}\PY{+w}{ }\PY{n}{String}\PY{p}{:}\PY{p}{:}\PY{n}{startsWith}\PY{p}{;}
\PY{n}{System}\PY{p}{.}\PY{n+na}{out}\PY{p}{.}\PY{n+na}{println}\PY{p}{(}\PY{n}{startsWith}\PY{p}{.}\PY{n+na}{apply}\PY{p}{(}\PY{l+s}{\PYZdq{}}\PY{l+s}{hello}\PY{l+s}{\PYZdq{}}\PY{p}{,}\PY{+w}{ }\PY{l+s}{\PYZdq{}}\PY{l+s}{he}\PY{l+s}{\PYZdq{}}\PY{p}{)}\PY{p}{)}\PY{p}{;}\PY{+w}{ }\PY{c+c1}{// true}

\PY{c+c1}{// 4. Constructor reference}
\PY{n}{Supplier}\PY{o}{\PYZlt{}}\PY{n}{List}\PY{o}{\PYZlt{}}\PY{n}{String}\PY{o}{\PYZgt{}}\PY{o}{\PYZgt{}}\PY{+w}{ }\PY{n}{listFactory}\PY{+w}{ }\PY{o}{=}\PY{+w}{ }\PY{n}{ArrayList}\PY{p}{:}\PY{p}{:}\PY{k}{new}\PY{p}{;}
\PY{n}{List}\PY{o}{\PYZlt{}}\PY{n}{String}\PY{o}{\PYZgt{}}\PY{+w}{ }\PY{n}{fresh}\PY{+w}{ }\PY{o}{=}\PY{+w}{ }\PY{n}{listFactory}\PY{p}{.}\PY{n+na}{get}\PY{p}{(}\PY{p}{)}\PY{p}{;}

\PY{n}{Function}\PY{o}{\PYZlt{}}\PY{n}{String}\PY{p}{,}\PY{+w}{ }\PY{n}{StringBuilder}\PY{o}{\PYZgt{}}\PY{+w}{ }\PY{n}{sbFactory}\PY{+w}{ }\PY{o}{=}\PY{+w}{ }\PY{n}{StringBuilder}\PY{p}{:}\PY{p}{:}\PY{k}{new}\PY{p}{;}
\PY{n}{System}\PY{p}{.}\PY{n+na}{out}\PY{p}{.}\PY{n+na}{println}\PY{p}{(}\PY{n}{sbFactory}\PY{p}{.}\PY{n+na}{apply}\PY{p}{(}\PY{l+s}{\PYZdq{}}\PY{l+s}{go}\PY{l+s}{\PYZdq{}}\PY{p}{)}\PY{p}{.}\PY{n+na}{append}\PY{p}{(}\PY{l+s}{\PYZdq{}}\PY{l+s}{!}\PY{l+s}{\PYZdq{}}\PY{p}{)}\PY{p}{)}\PY{p}{;}\PY{+w}{ }\PY{c+c1}{// go!}
\end{Verbatim}
\end{codebox}

A quick way to tell \#2 from \#3: if the part before \code{::} is a \emph{variable} (an existing object), it’s bound. If it’s a \emph{type name}, it’s unbound and the receiver becomes the method’s first parameter.

\section[Streams API]{Streams API}
\label{h:8:streams-api}
A \textbf{stream} is a sequence of elements that supports functional-style operations — \code{map}, \code{filter}, \code{reduce}, and so on — computed in a pipeline. A stream is \emph{not} a data structure; it doesn’t store elements. It’s a view over a source (a collection, an array, an I/O channel, a generator function) that computes results on demand.

\subsection[Creating streams]{Creating streams}
\label{h:8:creating-streams}
\begin{codebox}
\begin{Verbatim}[commandchars=\\\{\}]
\PY{n}{Stream}\PY{o}{\PYZlt{}}\PY{n}{String}\PY{o}{\PYZgt{}}\PY{+w}{ }\PY{n}{fromValues}\PY{+w}{ }\PY{o}{=}\PY{+w}{ }\PY{n}{Stream}\PY{p}{.}\PY{n+na}{of}\PY{p}{(}\PY{l+s}{\PYZdq{}}\PY{l+s}{a}\PY{l+s}{\PYZdq{}}\PY{p}{,}\PY{+w}{ }\PY{l+s}{\PYZdq{}}\PY{l+s}{b}\PY{l+s}{\PYZdq{}}\PY{p}{,}\PY{+w}{ }\PY{l+s}{\PYZdq{}}\PY{l+s}{c}\PY{l+s}{\PYZdq{}}\PY{p}{)}\PY{p}{;}
\PY{n}{Stream}\PY{o}{\PYZlt{}}\PY{n}{String}\PY{o}{\PYZgt{}}\PY{+w}{ }\PY{n}{fromList}\PY{+w}{ }\PY{o}{=}\PY{+w}{ }\PY{n}{List}\PY{p}{.}\PY{n+na}{of}\PY{p}{(}\PY{l+s}{\PYZdq{}}\PY{l+s}{a}\PY{l+s}{\PYZdq{}}\PY{p}{,}\PY{+w}{ }\PY{l+s}{\PYZdq{}}\PY{l+s}{b}\PY{l+s}{\PYZdq{}}\PY{p}{,}\PY{+w}{ }\PY{l+s}{\PYZdq{}}\PY{l+s}{c}\PY{l+s}{\PYZdq{}}\PY{p}{)}\PY{p}{.}\PY{n+na}{stream}\PY{p}{(}\PY{p}{)}\PY{p}{;}
\PY{n}{Stream}\PY{o}{\PYZlt{}}\PY{n}{String}\PY{o}{\PYZgt{}}\PY{+w}{ }\PY{n}{fromArray}\PY{+w}{ }\PY{o}{=}\PY{+w}{ }\PY{n}{Arrays}\PY{p}{.}\PY{n+na}{stream}\PY{p}{(}\PY{k}{new}\PY{+w}{ }\PY{n}{String}\PY{o}{[}\PY{o}{]}\PY{p}{\PYZob{}}\PY{l+s}{\PYZdq{}}\PY{l+s}{a}\PY{l+s}{\PYZdq{}}\PY{p}{,}\PY{+w}{ }\PY{l+s}{\PYZdq{}}\PY{l+s}{b}\PY{l+s}{\PYZdq{}}\PY{p}{,}\PY{+w}{ }\PY{l+s}{\PYZdq{}}\PY{l+s}{c}\PY{l+s}{\PYZdq{}}\PY{p}{\PYZcb{}}\PY{p}{)}\PY{p}{;}
\PY{n}{Stream}\PY{o}{\PYZlt{}}\PY{n}{String}\PY{o}{\PYZgt{}}\PY{+w}{ }\PY{n}{empty}\PY{+w}{ }\PY{o}{=}\PY{+w}{ }\PY{n}{Stream}\PY{p}{.}\PY{n+na}{empty}\PY{p}{(}\PY{p}{)}\PY{p}{;}
\PY{n}{IntStream}\PY{+w}{ }\PY{n}{fromRange}\PY{+w}{ }\PY{o}{=}\PY{+w}{ }\PY{n}{IntStream}\PY{p}{.}\PY{n+na}{range}\PY{p}{(}\PY{l+m+mi}{0}\PY{p}{,}\PY{+w}{ }\PY{l+m+mi}{5}\PY{p}{)}\PY{p}{;}\PY{+w}{        }\PY{c+c1}{// 0,1,2,3,4}
\PY{n}{IntStream}\PY{+w}{ }\PY{n}{fromRangeClosed}\PY{+w}{ }\PY{o}{=}\PY{+w}{ }\PY{n}{IntStream}\PY{p}{.}\PY{n+na}{rangeClosed}\PY{p}{(}\PY{l+m+mi}{1}\PY{p}{,}\PY{+w}{ }\PY{l+m+mi}{5}\PY{p}{)}\PY{p}{;}\PY{+w}{ }\PY{c+c1}{// 1,2,3,4,5}
\end{Verbatim}
\end{codebox}

\subsection[Intermediate vs terminal operations]{Intermediate vs terminal operations}
\label{h:8:intermediate-vs-terminal-operations}
\textbf{Intermediate operations} (\code{map}, \code{filter}, \code{sorted}, …) return a new stream and are \emph{lazy} — they don’t run anything by themselves. \textbf{Terminal operations} (\code{collect}, \code{forEach}, \code{reduce}, \code{count}, …) trigger the pipeline to actually execute and produce a result or side effect. A stream pipeline with no terminal operation does nothing at all.

\begin{codebox}
\begin{Verbatim}[commandchars=\\\{\}]
\PY{n}{Stream}\PY{o}{\PYZlt{}}\PY{n}{String}\PY{o}{\PYZgt{}}\PY{+w}{ }\PY{n}{pipeline}\PY{+w}{ }\PY{o}{=}\PY{+w}{ }\PY{n}{List}\PY{p}{.}\PY{n+na}{of}\PY{p}{(}\PY{l+s}{\PYZdq{}}\PY{l+s}{banana}\PY{l+s}{\PYZdq{}}\PY{p}{,}\PY{+w}{ }\PY{l+s}{\PYZdq{}}\PY{l+s}{apple}\PY{l+s}{\PYZdq{}}\PY{p}{,}\PY{+w}{ }\PY{l+s}{\PYZdq{}}\PY{l+s}{cherry}\PY{l+s}{\PYZdq{}}\PY{p}{)}\PY{p}{.}\PY{n+na}{stream}\PY{p}{(}\PY{p}{)}
\PY{+w}{    }\PY{p}{.}\PY{n+na}{filter}\PY{p}{(}\PY{n}{s}\PY{+w}{ }\PY{o}{\PYZhy{}}\PY{o}{\PYZgt{}}\PY{+w}{ }\PY{n}{s}\PY{p}{.}\PY{n+na}{length}\PY{p}{(}\PY{p}{)}\PY{+w}{ }\PY{o}{\PYZgt{}}\PY{+w}{ }\PY{l+m+mi}{5}\PY{p}{)}\PY{+w}{  }\PY{c+c1}{// nothing happens yet}
\PY{+w}{    }\PY{p}{.}\PY{n+na}{map}\PY{p}{(}\PY{n}{String}\PY{p}{:}\PY{p}{:}\PY{n}{toUpperCase}\PY{p}{)}\PY{p}{;}\PY{+w}{    }\PY{c+c1}{// still nothing happens}

\PY{n}{List}\PY{o}{\PYZlt{}}\PY{n}{String}\PY{o}{\PYZgt{}}\PY{+w}{ }\PY{n}{result}\PY{+w}{ }\PY{o}{=}\PY{+w}{ }\PY{n}{pipeline}\PY{p}{.}\PY{n+na}{collect}\PY{p}{(}\PY{n}{Collectors}\PY{p}{.}\PY{n+na}{toList}\PY{p}{(}\PY{p}{)}\PY{p}{)}\PY{p}{;}\PY{+w}{ }\PY{c+c1}{// NOW it runs}
\PY{n}{System}\PY{p}{.}\PY{n+na}{out}\PY{p}{.}\PY{n+na}{println}\PY{p}{(}\PY{n}{result}\PY{p}{)}\PY{p}{;}\PY{+w}{ }\PY{c+c1}{// [BANANA, CHERRY]}
\end{Verbatim}
\end{codebox}

\subsection[Laziness and short-circuiting]{Laziness and short-circuiting}
\label{h:8:laziness-and-short-circuiting}
Because intermediate ops are lazy, the pipeline processes one element at a time, end-to-end, rather than materializing an intermediate list at every stage. Combined with \textbf{short-circuiting} operations (\code{findFirst}, \code{anyMatch}, \code{limit}, …), this means a stream can stop early without touching the rest of the source.

\begin{codebox}
\begin{Verbatim}[commandchars=\\\{\}]
\PY{n}{Optional}\PY{o}{\PYZlt{}}\PY{n}{Integer}\PY{o}{\PYZgt{}}\PY{+w}{ }\PY{n}{firstBig}\PY{+w}{ }\PY{o}{=}\PY{+w}{ }\PY{n}{Stream}\PY{p}{.}\PY{n+na}{of}\PY{p}{(}\PY{l+m+mi}{1}\PY{p}{,}\PY{+w}{ }\PY{l+m+mi}{2}\PY{p}{,}\PY{+w}{ }\PY{l+m+mi}{3}\PY{p}{,}\PY{+w}{ }\PY{l+m+mi}{4}\PY{p}{,}\PY{+w}{ }\PY{l+m+mi}{5}\PY{p}{,}\PY{+w}{ }\PY{l+m+mi}{6}\PY{p}{)}
\PY{+w}{    }\PY{p}{.}\PY{n+na}{peek}\PY{p}{(}\PY{n}{n}\PY{+w}{ }\PY{o}{\PYZhy{}}\PY{o}{\PYZgt{}}\PY{+w}{ }\PY{n}{System}\PY{p}{.}\PY{n+na}{out}\PY{p}{.}\PY{n+na}{println}\PY{p}{(}\PY{l+s}{\PYZdq{}}\PY{l+s}{checking }\PY{l+s}{\PYZdq{}}\PY{+w}{ }\PY{o}{+}\PY{+w}{ }\PY{n}{n}\PY{p}{)}\PY{p}{)}
\PY{+w}{    }\PY{p}{.}\PY{n+na}{filter}\PY{p}{(}\PY{n}{n}\PY{+w}{ }\PY{o}{\PYZhy{}}\PY{o}{\PYZgt{}}\PY{+w}{ }\PY{n}{n}\PY{+w}{ }\PY{o}{\PYZgt{}}\PY{+w}{ }\PY{l+m+mi}{3}\PY{p}{)}
\PY{+w}{    }\PY{p}{.}\PY{n+na}{findFirst}\PY{p}{(}\PY{p}{)}\PY{p}{;}
\PY{c+c1}{// Output:}
\PY{c+c1}{// checking 1}
\PY{c+c1}{// checking 2}
\PY{c+c1}{// checking 3}
\PY{c+c1}{// checking 4}
\PY{c+c1}{// (stops here \PYZhy{} findFirst is short\PYZhy{}circuiting)}
\PY{n}{System}\PY{p}{.}\PY{n+na}{out}\PY{p}{.}\PY{n+na}{println}\PY{p}{(}\PY{n}{firstBig}\PY{p}{)}\PY{p}{;}\PY{+w}{ }\PY{c+c1}{// Optional[4]}
\end{Verbatim}
\end{codebox}

\subsection[map, filter, flatMap, mapMulti]{\code{map}, \code{filter}, \code{flatMap}, \code{mapMulti}}
\label{h:8:map-filter-flatmap-mapmulti}
\begin{itemize}
\item 
\code{map}: transform each element 1-to-1.

\item 
\code{filter}: keep elements matching a \code{Predicate}.

\item 
\code{flatMap}: transform each element into a \emph{stream}, then flatten all those streams into one (1-to-many, or 1-to-zero).

\item 
\code{mapMulti} (Java 16+): an imperative alternative to \code{flatMap} for cases where building an intermediate stream is wasteful — you push zero or more results into a consumer yourself.

\end{itemize}

\begin{codebox}
\begin{Verbatim}[commandchars=\\\{\}]
\PY{n}{List}\PY{o}{\PYZlt{}}\PY{n}{List}\PY{o}{\PYZlt{}}\PY{n}{Integer}\PY{o}{\PYZgt{}}\PY{o}{\PYZgt{}}\PY{+w}{ }\PY{n}{nested}\PY{+w}{ }\PY{o}{=}\PY{+w}{ }\PY{n}{List}\PY{p}{.}\PY{n+na}{of}\PY{p}{(}\PY{n}{List}\PY{p}{.}\PY{n+na}{of}\PY{p}{(}\PY{l+m+mi}{1}\PY{p}{,}\PY{+w}{ }\PY{l+m+mi}{2}\PY{p}{)}\PY{p}{,}\PY{+w}{ }\PY{n}{List}\PY{p}{.}\PY{n+na}{of}\PY{p}{(}\PY{l+m+mi}{3}\PY{p}{,}\PY{+w}{ }\PY{l+m+mi}{4}\PY{p}{)}\PY{p}{,}\PY{+w}{ }\PY{n}{List}\PY{p}{.}\PY{n+na}{of}\PY{p}{(}\PY{p}{)}\PY{p}{)}\PY{p}{;}

\PY{n}{List}\PY{o}{\PYZlt{}}\PY{n}{Integer}\PY{o}{\PYZgt{}}\PY{+w}{ }\PY{n}{flat}\PY{+w}{ }\PY{o}{=}\PY{+w}{ }\PY{n}{nested}\PY{p}{.}\PY{n+na}{stream}\PY{p}{(}\PY{p}{)}
\PY{+w}{    }\PY{p}{.}\PY{n+na}{flatMap}\PY{p}{(}\PY{n}{List}\PY{p}{:}\PY{p}{:}\PY{n}{stream}\PY{p}{)}
\PY{+w}{    }\PY{p}{.}\PY{n+na}{toList}\PY{p}{(}\PY{p}{)}\PY{p}{;}
\PY{n}{System}\PY{p}{.}\PY{n+na}{out}\PY{p}{.}\PY{n+na}{println}\PY{p}{(}\PY{n}{flat}\PY{p}{)}\PY{p}{;}\PY{+w}{ }\PY{c+c1}{// [1, 2, 3, 4]}

\PY{n}{List}\PY{o}{\PYZlt{}}\PY{n}{Integer}\PY{o}{\PYZgt{}}\PY{+w}{ }\PY{n}{viaMapMulti}\PY{+w}{ }\PY{o}{=}\PY{+w}{ }\PY{n}{nested}\PY{p}{.}\PY{n+na}{stream}\PY{p}{(}\PY{p}{)}
\PY{+w}{    }\PY{p}{.}\PY{o}{\PYZlt{}}\PY{n}{Integer}\PY{o}{\PYZgt{}}\PY{n}{mapMulti}\PY{p}{(}\PY{p}{(}\PY{n}{list}\PY{p}{,}\PY{+w}{ }\PY{n}{consumer}\PY{p}{)}\PY{+w}{ }\PY{o}{\PYZhy{}}\PY{o}{\PYZgt{}}\PY{+w}{ }\PY{n}{list}\PY{p}{.}\PY{n+na}{forEach}\PY{p}{(}\PY{n}{consumer}\PY{p}{)}\PY{p}{)}
\PY{+w}{    }\PY{p}{.}\PY{n+na}{toList}\PY{p}{(}\PY{p}{)}\PY{p}{;}
\PY{n}{System}\PY{p}{.}\PY{n+na}{out}\PY{p}{.}\PY{n+na}{println}\PY{p}{(}\PY{n}{viaMapMulti}\PY{p}{)}\PY{p}{;}\PY{+w}{ }\PY{c+c1}{// [1, 2, 3, 4]}
\end{Verbatim}
\end{codebox}

\subsection[sorted, distinct]{\code{sorted}, \code{distinct}}
\label{h:8:sorted-distinct}
\begin{codebox}
\begin{Verbatim}[commandchars=\\\{\}]
\PY{n}{List}\PY{o}{\PYZlt{}}\PY{n}{String}\PY{o}{\PYZgt{}}\PY{+w}{ }\PY{n}{words}\PY{+w}{ }\PY{o}{=}\PY{+w}{ }\PY{n}{List}\PY{p}{.}\PY{n+na}{of}\PY{p}{(}\PY{l+s}{\PYZdq{}}\PY{l+s}{pear}\PY{l+s}{\PYZdq{}}\PY{p}{,}\PY{+w}{ }\PY{l+s}{\PYZdq{}}\PY{l+s}{fig}\PY{l+s}{\PYZdq{}}\PY{p}{,}\PY{+w}{ }\PY{l+s}{\PYZdq{}}\PY{l+s}{apple}\PY{l+s}{\PYZdq{}}\PY{p}{,}\PY{+w}{ }\PY{l+s}{\PYZdq{}}\PY{l+s}{fig}\PY{l+s}{\PYZdq{}}\PY{p}{,}\PY{+w}{ }\PY{l+s}{\PYZdq{}}\PY{l+s}{date}\PY{l+s}{\PYZdq{}}\PY{p}{)}\PY{p}{;}

\PY{n}{List}\PY{o}{\PYZlt{}}\PY{n}{String}\PY{o}{\PYZgt{}}\PY{+w}{ }\PY{n}{sortedDistinct}\PY{+w}{ }\PY{o}{=}\PY{+w}{ }\PY{n}{words}\PY{p}{.}\PY{n+na}{stream}\PY{p}{(}\PY{p}{)}
\PY{+w}{    }\PY{p}{.}\PY{n+na}{distinct}\PY{p}{(}\PY{p}{)}\PY{+w}{                        }\PY{c+c1}{// relies on equals()/hashCode()}
\PY{+w}{    }\PY{p}{.}\PY{n+na}{sorted}\PY{p}{(}\PY{n}{Comparator}\PY{p}{.}\PY{n+na}{comparing}\PY{p}{(}\PY{n}{String}\PY{p}{:}\PY{p}{:}\PY{n}{length}\PY{p}{)}\PY{p}{.}\PY{n+na}{thenComparing}\PY{p}{(}\PY{n}{Comparator}\PY{p}{.}\PY{n+na}{naturalOrder}\PY{p}{(}\PY{p}{)}\PY{p}{)}\PY{p}{)}
\PY{+w}{    }\PY{p}{.}\PY{n+na}{toList}\PY{p}{(}\PY{p}{)}\PY{p}{;}
\PY{n}{System}\PY{p}{.}\PY{n+na}{out}\PY{p}{.}\PY{n+na}{println}\PY{p}{(}\PY{n}{sortedDistinct}\PY{p}{)}\PY{p}{;}\PY{+w}{ }\PY{c+c1}{// [fig, date, pear, apple]}
\end{Verbatim}
\end{codebox}

\code{sorted()} on a stream is a \emph{stateful} intermediate op — it must buffer the whole stream before it can emit anything, so it doesn’t play well with infinite streams unless combined with \code{limit} first.

\subsection[limit/skip, takeWhile/dropWhile]{\code{limit}/\code{skip}, \code{takeWhile}/\code{dropWhile}}
\label{h:8:limitskip-takewhiledropwhile}
\code{limit}/\code{skip} cut by \emph{position}. \code{takeWhile}/\code{dropWhile} (Java 9+) cut by a \emph{predicate}, and both stop evaluating the predicate as soon as it flips (they assume, but do not enforce, that the input is suitably ordered for the use case).

\begin{codebox}
\begin{Verbatim}[commandchars=\\\{\}]
\PY{n}{List}\PY{o}{\PYZlt{}}\PY{n}{Integer}\PY{o}{\PYZgt{}}\PY{+w}{ }\PY{n}{nums}\PY{+w}{ }\PY{o}{=}\PY{+w}{ }\PY{n}{List}\PY{p}{.}\PY{n+na}{of}\PY{p}{(}\PY{l+m+mi}{1}\PY{p}{,}\PY{+w}{ }\PY{l+m+mi}{2}\PY{p}{,}\PY{+w}{ }\PY{l+m+mi}{3}\PY{p}{,}\PY{+w}{ }\PY{l+m+mi}{10}\PY{p}{,}\PY{+w}{ }\PY{l+m+mi}{4}\PY{p}{,}\PY{+w}{ }\PY{l+m+mi}{5}\PY{p}{)}\PY{p}{;}

\PY{n}{System}\PY{p}{.}\PY{n+na}{out}\PY{p}{.}\PY{n+na}{println}\PY{p}{(}\PY{n}{nums}\PY{p}{.}\PY{n+na}{stream}\PY{p}{(}\PY{p}{)}\PY{p}{.}\PY{n+na}{skip}\PY{p}{(}\PY{l+m+mi}{1}\PY{p}{)}\PY{p}{.}\PY{n+na}{limit}\PY{p}{(}\PY{l+m+mi}{2}\PY{p}{)}\PY{p}{.}\PY{n+na}{toList}\PY{p}{(}\PY{p}{)}\PY{p}{)}\PY{p}{;}\PY{+w}{      }\PY{c+c1}{// [2, 3]}
\PY{n}{System}\PY{p}{.}\PY{n+na}{out}\PY{p}{.}\PY{n+na}{println}\PY{p}{(}\PY{n}{nums}\PY{p}{.}\PY{n+na}{stream}\PY{p}{(}\PY{p}{)}\PY{p}{.}\PY{n+na}{takeWhile}\PY{p}{(}\PY{n}{n}\PY{+w}{ }\PY{o}{\PYZhy{}}\PY{o}{\PYZgt{}}\PY{+w}{ }\PY{n}{n}\PY{+w}{ }\PY{o}{\PYZlt{}}\PY{+w}{ }\PY{l+m+mi}{5}\PY{p}{)}\PY{p}{.}\PY{n+na}{toList}\PY{p}{(}\PY{p}{)}\PY{p}{)}\PY{p}{;}\PY{+w}{ }\PY{c+c1}{// [1, 2, 3] \PYZhy{} stops at 10}
\PY{n}{System}\PY{p}{.}\PY{n+na}{out}\PY{p}{.}\PY{n+na}{println}\PY{p}{(}\PY{n}{nums}\PY{p}{.}\PY{n+na}{stream}\PY{p}{(}\PY{p}{)}\PY{p}{.}\PY{n+na}{dropWhile}\PY{p}{(}\PY{n}{n}\PY{+w}{ }\PY{o}{\PYZhy{}}\PY{o}{\PYZgt{}}\PY{+w}{ }\PY{n}{n}\PY{+w}{ }\PY{o}{\PYZlt{}}\PY{+w}{ }\PY{l+m+mi}{5}\PY{p}{)}\PY{p}{.}\PY{n+na}{toList}\PY{p}{(}\PY{p}{)}\PY{p}{)}\PY{p}{;}\PY{+w}{ }\PY{c+c1}{// [10, 4, 5] \PYZhy{} drops until first \PYZgt{}= 5, then keeps rest}
\end{Verbatim}
\end{codebox}

\subsection[peek]{\code{peek}}
\label{h:8:peek}
\code{peek} lets you look at each element as it flows through, without transforming it. It exists for debugging a pipeline — \textbf{not} for side effects that affect the program’s result. Two reasons:

\begin{enumerate}
\item 
If the stream is short-circuited (e.g., a later \code{findFirst}), \code{peek} may not run for every element.

\item 
Implementations are free to skip calling \code{peek} at all if they can prove the pipeline’s result doesn’t depend on it (an optimization the JDK explicitly reserves the right to make).

\end{enumerate}

\begin{codebox}
\begin{Verbatim}[commandchars=\\\{\}]
\PY{c+c1}{// Fine: debugging only, no logic depends on it}
\PY{k+kt}{long}\PY{+w}{ }\PY{n}{count}\PY{+w}{ }\PY{o}{=}\PY{+w}{ }\PY{n}{Stream}\PY{p}{.}\PY{n+na}{of}\PY{p}{(}\PY{l+s}{\PYZdq{}}\PY{l+s}{a}\PY{l+s}{\PYZdq{}}\PY{p}{,}\PY{+w}{ }\PY{l+s}{\PYZdq{}}\PY{l+s}{bb}\PY{l+s}{\PYZdq{}}\PY{p}{,}\PY{+w}{ }\PY{l+s}{\PYZdq{}}\PY{l+s}{ccc}\PY{l+s}{\PYZdq{}}\PY{p}{)}
\PY{+w}{    }\PY{p}{.}\PY{n+na}{peek}\PY{p}{(}\PY{n}{s}\PY{+w}{ }\PY{o}{\PYZhy{}}\PY{o}{\PYZgt{}}\PY{+w}{ }\PY{n}{System}\PY{p}{.}\PY{n+na}{out}\PY{p}{.}\PY{n+na}{println}\PY{p}{(}\PY{l+s}{\PYZdq{}}\PY{l+s}{seen: }\PY{l+s}{\PYZdq{}}\PY{+w}{ }\PY{o}{+}\PY{+w}{ }\PY{n}{s}\PY{p}{)}\PY{p}{)}
\PY{+w}{    }\PY{p}{.}\PY{n+na}{count}\PY{p}{(}\PY{p}{)}\PY{p}{;}\PY{+w}{ }\PY{c+c1}{// JDK may actually SKIP peek entirely here since count doesn\PYZsq{}t need the elements!}
\end{Verbatim}
\end{codebox}

\subsection[reduce]{\code{reduce}}
\label{h:8:reduce}
\code{reduce} folds a stream down to a single value using a combining function.

\begin{codebox}
\begin{Verbatim}[commandchars=\\\{\}]
\PY{n}{List}\PY{o}{\PYZlt{}}\PY{n}{Integer}\PY{o}{\PYZgt{}}\PY{+w}{ }\PY{n}{nums}\PY{+w}{ }\PY{o}{=}\PY{+w}{ }\PY{n}{List}\PY{p}{.}\PY{n+na}{of}\PY{p}{(}\PY{l+m+mi}{1}\PY{p}{,}\PY{+w}{ }\PY{l+m+mi}{2}\PY{p}{,}\PY{+w}{ }\PY{l+m+mi}{3}\PY{p}{,}\PY{+w}{ }\PY{l+m+mi}{4}\PY{p}{,}\PY{+w}{ }\PY{l+m+mi}{5}\PY{p}{)}\PY{p}{;}

\PY{k+kt}{int}\PY{+w}{ }\PY{n}{sum}\PY{+w}{ }\PY{o}{=}\PY{+w}{ }\PY{n}{nums}\PY{p}{.}\PY{n+na}{stream}\PY{p}{(}\PY{p}{)}\PY{p}{.}\PY{n+na}{reduce}\PY{p}{(}\PY{l+m+mi}{0}\PY{p}{,}\PY{+w}{ }\PY{n}{Integer}\PY{p}{:}\PY{p}{:}\PY{n}{sum}\PY{p}{)}\PY{p}{;}\PY{+w}{               }\PY{c+c1}{// identity + accumulator}
\PY{n}{System}\PY{p}{.}\PY{n+na}{out}\PY{p}{.}\PY{n+na}{println}\PY{p}{(}\PY{n}{sum}\PY{p}{)}\PY{p}{;}\PY{+w}{ }\PY{c+c1}{// 15}

\PY{n}{Optional}\PY{o}{\PYZlt{}}\PY{n}{Integer}\PY{o}{\PYZgt{}}\PY{+w}{ }\PY{n}{product}\PY{+w}{ }\PY{o}{=}\PY{+w}{ }\PY{n}{nums}\PY{p}{.}\PY{n+na}{stream}\PY{p}{(}\PY{p}{)}\PY{p}{.}\PY{n+na}{reduce}\PY{p}{(}\PY{p}{(}\PY{n}{a}\PY{p}{,}\PY{+w}{ }\PY{n}{b}\PY{p}{)}\PY{+w}{ }\PY{o}{\PYZhy{}}\PY{o}{\PYZgt{}}\PY{+w}{ }\PY{n}{a}\PY{+w}{ }\PY{o}{*}\PY{+w}{ }\PY{n}{b}\PY{p}{)}\PY{p}{;}\PY{+w}{ }\PY{c+c1}{// no identity \PYZhy{}\PYZgt{} Optional}
\PY{n}{System}\PY{p}{.}\PY{n+na}{out}\PY{p}{.}\PY{n+na}{println}\PY{p}{(}\PY{n}{product}\PY{p}{)}\PY{p}{;}\PY{+w}{ }\PY{c+c1}{// Optional[120]}

\PY{c+c1}{// Three\PYZhy{}arg form: identity, accumulator, combiner (combiner used only when running in parallel)}
\PY{k+kt}{int}\PY{+w}{ }\PY{n}{totalLength}\PY{+w}{ }\PY{o}{=}\PY{+w}{ }\PY{n}{Stream}\PY{p}{.}\PY{n+na}{of}\PY{p}{(}\PY{l+s}{\PYZdq{}}\PY{l+s}{a}\PY{l+s}{\PYZdq{}}\PY{p}{,}\PY{+w}{ }\PY{l+s}{\PYZdq{}}\PY{l+s}{bb}\PY{l+s}{\PYZdq{}}\PY{p}{,}\PY{+w}{ }\PY{l+s}{\PYZdq{}}\PY{l+s}{ccc}\PY{l+s}{\PYZdq{}}\PY{p}{)}
\PY{+w}{    }\PY{p}{.}\PY{n+na}{reduce}\PY{p}{(}\PY{l+m+mi}{0}\PY{p}{,}\PY{+w}{ }\PY{p}{(}\PY{n}{acc}\PY{p}{,}\PY{+w}{ }\PY{n}{s}\PY{p}{)}\PY{+w}{ }\PY{o}{\PYZhy{}}\PY{o}{\PYZgt{}}\PY{+w}{ }\PY{n}{acc}\PY{+w}{ }\PY{o}{+}\PY{+w}{ }\PY{n}{s}\PY{p}{.}\PY{n+na}{length}\PY{p}{(}\PY{p}{)}\PY{p}{,}\PY{+w}{ }\PY{n}{Integer}\PY{p}{:}\PY{p}{:}\PY{n}{sum}\PY{p}{)}\PY{p}{;}
\PY{n}{System}\PY{p}{.}\PY{n+na}{out}\PY{p}{.}\PY{n+na}{println}\PY{p}{(}\PY{n}{totalLength}\PY{p}{)}\PY{p}{;}\PY{+w}{ }\PY{c+c1}{// 6}
\end{Verbatim}
\end{codebox}

\subsection[findFirst vs findAny]{\code{findFirst} vs \code{findAny}}
\label{h:8:findfirst-vs-findany}
Both are short-circuiting and return \code{Optional\textless{}\allowbreak{}T\textgreater{}}. \code{findFirst} always returns the first matching element in encounter order. \code{findAny} may return \emph{any} matching element — it exists so that a parallel stream can return as soon as one thread finds a match, without waiting to determine which was truly “first.” On a sequential stream they usually behave the same, but only \code{findFirst} is guaranteed to.

\begin{codebox}
\begin{Verbatim}[commandchars=\\\{\}]
\PY{n}{List}\PY{o}{\PYZlt{}}\PY{n}{Integer}\PY{o}{\PYZgt{}}\PY{+w}{ }\PY{n}{nums}\PY{+w}{ }\PY{o}{=}\PY{+w}{ }\PY{n}{List}\PY{p}{.}\PY{n+na}{of}\PY{p}{(}\PY{l+m+mi}{1}\PY{p}{,}\PY{+w}{ }\PY{l+m+mi}{2}\PY{p}{,}\PY{+w}{ }\PY{l+m+mi}{3}\PY{p}{,}\PY{+w}{ }\PY{l+m+mi}{4}\PY{p}{,}\PY{+w}{ }\PY{l+m+mi}{5}\PY{p}{)}\PY{p}{;}

\PY{n}{Optional}\PY{o}{\PYZlt{}}\PY{n}{Integer}\PY{o}{\PYZgt{}}\PY{+w}{ }\PY{n}{first}\PY{+w}{ }\PY{o}{=}\PY{+w}{ }\PY{n}{nums}\PY{p}{.}\PY{n+na}{parallelStream}\PY{p}{(}\PY{p}{)}\PY{p}{.}\PY{n+na}{filter}\PY{p}{(}\PY{n}{n}\PY{+w}{ }\PY{o}{\PYZhy{}}\PY{o}{\PYZgt{}}\PY{+w}{ }\PY{n}{n}\PY{+w}{ }\PY{o}{\PYZpc{}}\PY{+w}{ }\PY{l+m+mi}{2}\PY{+w}{ }\PY{o}{=}\PY{o}{=}\PY{+w}{ }\PY{l+m+mi}{0}\PY{p}{)}\PY{p}{.}\PY{n+na}{findFirst}\PY{p}{(}\PY{p}{)}\PY{p}{;}\PY{+w}{ }\PY{c+c1}{// always 2}
\PY{n}{Optional}\PY{o}{\PYZlt{}}\PY{n}{Integer}\PY{o}{\PYZgt{}}\PY{+w}{ }\PY{n}{any}\PY{+w}{ }\PY{o}{=}\PY{+w}{ }\PY{n}{nums}\PY{p}{.}\PY{n+na}{parallelStream}\PY{p}{(}\PY{p}{)}\PY{p}{.}\PY{n+na}{filter}\PY{p}{(}\PY{n}{n}\PY{+w}{ }\PY{o}{\PYZhy{}}\PY{o}{\PYZgt{}}\PY{+w}{ }\PY{n}{n}\PY{+w}{ }\PY{o}{\PYZpc{}}\PY{+w}{ }\PY{l+m+mi}{2}\PY{+w}{ }\PY{o}{=}\PY{o}{=}\PY{+w}{ }\PY{l+m+mi}{0}\PY{p}{)}\PY{p}{.}\PY{n+na}{findAny}\PY{p}{(}\PY{p}{)}\PY{p}{;}\PY{+w}{     }\PY{c+c1}{// 2 or 4, whichever thread finished first}
\end{Verbatim}
\end{codebox}

\subsection[Primitive streams]{Primitive streams}
\label{h:8:primitive-streams}
\code{IntStream}, \code{LongStream}, \code{DoubleStream} avoid boxing overhead and add numeric operations (\code{sum}, \code{average}, \code{max}, \code{min}, \code{summaryStatistics}). Convert between object streams and primitive streams with \code{mapToInt}/\code{mapToObj}/\code{boxed}.

\begin{codebox}
\begin{Verbatim}[commandchars=\\\{\}]
\PY{k+kt}{int}\PY{o}{[}\PY{o}{]}\PY{+w}{ }\PY{n}{values}\PY{+w}{ }\PY{o}{=}\PY{+w}{ }\PY{p}{\PYZob{}}\PY{l+m+mi}{3}\PY{p}{,}\PY{+w}{ }\PY{l+m+mi}{1}\PY{p}{,}\PY{+w}{ }\PY{l+m+mi}{4}\PY{p}{,}\PY{+w}{ }\PY{l+m+mi}{1}\PY{p}{,}\PY{+w}{ }\PY{l+m+mi}{5}\PY{p}{,}\PY{+w}{ }\PY{l+m+mi}{9}\PY{p}{\PYZcb{}}\PY{p}{;}

\PY{n}{IntSummaryStatistics}\PY{+w}{ }\PY{n}{stats}\PY{+w}{ }\PY{o}{=}\PY{+w}{ }\PY{n}{IntStream}\PY{p}{.}\PY{n+na}{of}\PY{p}{(}\PY{n}{values}\PY{p}{)}\PY{p}{.}\PY{n+na}{summaryStatistics}\PY{p}{(}\PY{p}{)}\PY{p}{;}
\PY{n}{System}\PY{p}{.}\PY{n+na}{out}\PY{p}{.}\PY{n+na}{println}\PY{p}{(}\PY{n}{stats}\PY{p}{.}\PY{n+na}{getSum}\PY{p}{(}\PY{p}{)}\PY{p}{)}\PY{p}{;}\PY{+w}{     }\PY{c+c1}{// 23}
\PY{n}{System}\PY{p}{.}\PY{n+na}{out}\PY{p}{.}\PY{n+na}{println}\PY{p}{(}\PY{n}{stats}\PY{p}{.}\PY{n+na}{getAverage}\PY{p}{(}\PY{p}{)}\PY{p}{)}\PY{p}{;}\PY{+w}{ }\PY{c+c1}{// 3.8333...}
\PY{n}{System}\PY{p}{.}\PY{n+na}{out}\PY{p}{.}\PY{n+na}{println}\PY{p}{(}\PY{n}{stats}\PY{p}{.}\PY{n+na}{getMax}\PY{p}{(}\PY{p}{)}\PY{p}{)}\PY{p}{;}\PY{+w}{     }\PY{c+c1}{// 9}

\PY{n}{List}\PY{o}{\PYZlt{}}\PY{n}{Integer}\PY{o}{\PYZgt{}}\PY{+w}{ }\PY{n}{boxed}\PY{+w}{ }\PY{o}{=}\PY{+w}{ }\PY{n}{IntStream}\PY{p}{.}\PY{n+na}{of}\PY{p}{(}\PY{n}{values}\PY{p}{)}\PY{p}{.}\PY{n+na}{boxed}\PY{p}{(}\PY{p}{)}\PY{p}{.}\PY{n+na}{toList}\PY{p}{(}\PY{p}{)}\PY{p}{;}\PY{+w}{ }\PY{c+c1}{// [3, 1, 4, 1, 5, 9]}

\PY{n}{List}\PY{o}{\PYZlt{}}\PY{n}{String}\PY{o}{\PYZgt{}}\PY{+w}{ }\PY{n}{words}\PY{+w}{ }\PY{o}{=}\PY{+w}{ }\PY{n}{List}\PY{p}{.}\PY{n+na}{of}\PY{p}{(}\PY{l+s}{\PYZdq{}}\PY{l+s}{a}\PY{l+s}{\PYZdq{}}\PY{p}{,}\PY{+w}{ }\PY{l+s}{\PYZdq{}}\PY{l+s}{bb}\PY{l+s}{\PYZdq{}}\PY{p}{,}\PY{+w}{ }\PY{l+s}{\PYZdq{}}\PY{l+s}{ccc}\PY{l+s}{\PYZdq{}}\PY{p}{)}\PY{p}{;}
\PY{k+kt}{int}\PY{+w}{ }\PY{n}{totalChars}\PY{+w}{ }\PY{o}{=}\PY{+w}{ }\PY{n}{words}\PY{p}{.}\PY{n+na}{stream}\PY{p}{(}\PY{p}{)}\PY{p}{.}\PY{n+na}{mapToInt}\PY{p}{(}\PY{n}{String}\PY{p}{:}\PY{p}{:}\PY{n}{length}\PY{p}{)}\PY{p}{.}\PY{n+na}{sum}\PY{p}{(}\PY{p}{)}\PY{p}{;}\PY{+w}{ }\PY{c+c1}{// 6}
\end{Verbatim}
\end{codebox}

\subsection[Stream.iterate/Stream.generate]{\code{Stream.\allowbreak{}iterate}/\code{Stream.\allowbreak{}generate}}
\label{h:8:streamiteratestreamgenerate}
Both produce potentially \emph{infinite} streams — always pair them with \code{limit}, or with \code{takeWhile} on \code{iterate}, or you’ll hang the program.

\begin{codebox}
\begin{Verbatim}[commandchars=\\\{\}]
\PY{n}{List}\PY{o}{\PYZlt{}}\PY{n}{Integer}\PY{o}{\PYZgt{}}\PY{+w}{ }\PY{n}{powersOfTwo}\PY{+w}{ }\PY{o}{=}\PY{+w}{ }\PY{n}{Stream}\PY{p}{.}\PY{n+na}{iterate}\PY{p}{(}\PY{l+m+mi}{1}\PY{p}{,}\PY{+w}{ }\PY{n}{n}\PY{+w}{ }\PY{o}{\PYZhy{}}\PY{o}{\PYZgt{}}\PY{+w}{ }\PY{n}{n}\PY{+w}{ }\PY{o}{*}\PY{+w}{ }\PY{l+m+mi}{2}\PY{p}{)}
\PY{+w}{    }\PY{p}{.}\PY{n+na}{limit}\PY{p}{(}\PY{l+m+mi}{5}\PY{p}{)}
\PY{+w}{    }\PY{p}{.}\PY{n+na}{toList}\PY{p}{(}\PY{p}{)}\PY{p}{;}
\PY{n}{System}\PY{p}{.}\PY{n+na}{out}\PY{p}{.}\PY{n+na}{println}\PY{p}{(}\PY{n}{powersOfTwo}\PY{p}{)}\PY{p}{;}\PY{+w}{ }\PY{c+c1}{// [1, 2, 4, 8, 16]}

\PY{c+c1}{// Java 9+ three\PYZhy{}arg iterate: seed, hasNext predicate (built\PYZhy{}in short\PYZhy{}circuit), next function}
\PY{n}{List}\PY{o}{\PYZlt{}}\PY{n}{Integer}\PY{o}{\PYZgt{}}\PY{+w}{ }\PY{n}{below100}\PY{+w}{ }\PY{o}{=}\PY{+w}{ }\PY{n}{Stream}\PY{p}{.}\PY{n+na}{iterate}\PY{p}{(}\PY{l+m+mi}{1}\PY{p}{,}\PY{+w}{ }\PY{n}{n}\PY{+w}{ }\PY{o}{\PYZhy{}}\PY{o}{\PYZgt{}}\PY{+w}{ }\PY{n}{n}\PY{+w}{ }\PY{o}{\PYZlt{}}\PY{+w}{ }\PY{l+m+mi}{100}\PY{p}{,}\PY{+w}{ }\PY{n}{n}\PY{+w}{ }\PY{o}{\PYZhy{}}\PY{o}{\PYZgt{}}\PY{+w}{ }\PY{n}{n}\PY{+w}{ }\PY{o}{*}\PY{+w}{ }\PY{l+m+mi}{2}\PY{p}{)}\PY{p}{.}\PY{n+na}{toList}\PY{p}{(}\PY{p}{)}\PY{p}{;}
\PY{n}{System}\PY{p}{.}\PY{n+na}{out}\PY{p}{.}\PY{n+na}{println}\PY{p}{(}\PY{n}{below100}\PY{p}{)}\PY{p}{;}\PY{+w}{ }\PY{c+c1}{// [1, 2, 4, 8, 16, 32, 64]}

\PY{n}{Stream}\PY{o}{\PYZlt{}}\PY{n}{Double}\PY{o}{\PYZgt{}}\PY{+w}{ }\PY{n}{randoms}\PY{+w}{ }\PY{o}{=}\PY{+w}{ }\PY{n}{Stream}\PY{p}{.}\PY{n+na}{generate}\PY{p}{(}\PY{n}{Math}\PY{p}{:}\PY{p}{:}\PY{n}{random}\PY{p}{)}\PY{p}{.}\PY{n+na}{limit}\PY{p}{(}\PY{l+m+mi}{3}\PY{p}{)}\PY{p}{;}
\end{Verbatim}
\end{codebox}

\subsection[Optional interop]{\code{Optional} interop}
\label{h:8:optional-interop}
Streams and \code{Optional} are designed to work together. \code{findFirst}, \code{reduce} (no identity), \code{max}, \code{min} all return \code{Optional}. \code{Optional} itself has a \code{stream()} method (0 or 1 element), handy for flattening a stream of \code{Optional}s.

\begin{codebox}
\begin{Verbatim}[commandchars=\\\{\}]
\PY{n}{List}\PY{o}{\PYZlt{}}\PY{n}{Optional}\PY{o}{\PYZlt{}}\PY{n}{String}\PY{o}{\PYZgt{}}\PY{o}{\PYZgt{}}\PY{+w}{ }\PY{n}{maybeNames}\PY{+w}{ }\PY{o}{=}\PY{+w}{ }\PY{n}{List}\PY{p}{.}\PY{n+na}{of}\PY{p}{(}\PY{n}{Optional}\PY{p}{.}\PY{n+na}{of}\PY{p}{(}\PY{l+s}{\PYZdq{}}\PY{l+s}{Ada}\PY{l+s}{\PYZdq{}}\PY{p}{)}\PY{p}{,}\PY{+w}{ }\PY{n}{Optional}\PY{p}{.}\PY{n+na}{empty}\PY{p}{(}\PY{p}{)}\PY{p}{,}\PY{+w}{ }\PY{n}{Optional}\PY{p}{.}\PY{n+na}{of}\PY{p}{(}\PY{l+s}{\PYZdq{}}\PY{l+s}{Linus}\PY{l+s}{\PYZdq{}}\PY{p}{)}\PY{p}{)}\PY{p}{;}

\PY{n}{List}\PY{o}{\PYZlt{}}\PY{n}{String}\PY{o}{\PYZgt{}}\PY{+w}{ }\PY{n}{presentOnly}\PY{+w}{ }\PY{o}{=}\PY{+w}{ }\PY{n}{maybeNames}\PY{p}{.}\PY{n+na}{stream}\PY{p}{(}\PY{p}{)}
\PY{+w}{    }\PY{p}{.}\PY{n+na}{flatMap}\PY{p}{(}\PY{n}{Optional}\PY{p}{:}\PY{p}{:}\PY{n}{stream}\PY{p}{)}\PY{+w}{ }\PY{c+c1}{// Optional.stream() turns Optional\PYZlt{}T\PYZgt{} into Stream\PYZlt{}T\PYZgt{} of 0 or 1 elements}
\PY{+w}{    }\PY{p}{.}\PY{n+na}{toList}\PY{p}{(}\PY{p}{)}\PY{p}{;}
\PY{n}{System}\PY{p}{.}\PY{n+na}{out}\PY{p}{.}\PY{n+na}{println}\PY{p}{(}\PY{n}{presentOnly}\PY{p}{)}\PY{p}{;}\PY{+w}{ }\PY{c+c1}{// [Ada, Linus]}
\end{Verbatim}
\end{codebox}

\subsection[Streams are single-use]{Streams are single-use}
\label{h:8:streams-are-single-use}
A stream can be consumed by a terminal operation exactly once. Reusing a stream throws \code{IllegalStateException}. If you need to run the pipeline twice, build it from the source again.

\begin{codebox}
\begin{Verbatim}[commandchars=\\\{\}]
\PY{n}{Stream}\PY{o}{\PYZlt{}}\PY{n}{String}\PY{o}{\PYZgt{}}\PY{+w}{ }\PY{n}{stream}\PY{+w}{ }\PY{o}{=}\PY{+w}{ }\PY{n}{Stream}\PY{p}{.}\PY{n+na}{of}\PY{p}{(}\PY{l+s}{\PYZdq{}}\PY{l+s}{a}\PY{l+s}{\PYZdq{}}\PY{p}{,}\PY{+w}{ }\PY{l+s}{\PYZdq{}}\PY{l+s}{b}\PY{l+s}{\PYZdq{}}\PY{p}{,}\PY{+w}{ }\PY{l+s}{\PYZdq{}}\PY{l+s}{c}\PY{l+s}{\PYZdq{}}\PY{p}{)}\PY{p}{;}
\PY{n}{stream}\PY{p}{.}\PY{n+na}{forEach}\PY{p}{(}\PY{n}{System}\PY{p}{.}\PY{n+na}{out}\PY{p}{:}\PY{p}{:}\PY{n}{println}\PY{p}{)}\PY{p}{;}
\PY{c+c1}{// stream.count(); // throws IllegalStateException: stream has already been operated upon or closed}

\PY{c+c1}{// Correct: create a new stream from the source each time}
\PY{n}{List}\PY{o}{\PYZlt{}}\PY{n}{String}\PY{o}{\PYZgt{}}\PY{+w}{ }\PY{n}{source}\PY{+w}{ }\PY{o}{=}\PY{+w}{ }\PY{n}{List}\PY{p}{.}\PY{n+na}{of}\PY{p}{(}\PY{l+s}{\PYZdq{}}\PY{l+s}{a}\PY{l+s}{\PYZdq{}}\PY{p}{,}\PY{+w}{ }\PY{l+s}{\PYZdq{}}\PY{l+s}{b}\PY{l+s}{\PYZdq{}}\PY{p}{,}\PY{+w}{ }\PY{l+s}{\PYZdq{}}\PY{l+s}{c}\PY{l+s}{\PYZdq{}}\PY{p}{)}\PY{p}{;}
\PY{n}{source}\PY{p}{.}\PY{n+na}{stream}\PY{p}{(}\PY{p}{)}\PY{p}{.}\PY{n+na}{forEach}\PY{p}{(}\PY{n}{System}\PY{p}{.}\PY{n+na}{out}\PY{p}{:}\PY{p}{:}\PY{n}{println}\PY{p}{)}\PY{p}{;}
\PY{k+kt}{long}\PY{+w}{ }\PY{n}{count}\PY{+w}{ }\PY{o}{=}\PY{+w}{ }\PY{n}{source}\PY{p}{.}\PY{n+na}{stream}\PY{p}{(}\PY{p}{)}\PY{p}{.}\PY{n+na}{count}\PY{p}{(}\PY{p}{)}\PY{p}{;}
\end{Verbatim}
\end{codebox}

\section[Collectors]{Collectors}
\label{h:8:collectors}
A \textbf{collector} is a recipe for how to accumulate the elements of a stream into a result — a \code{List}, a \code{Map}, a single summary number, a \code{String}, or anything else. You use one via the terminal \code{collect(...)} operation. \code{Collectors} is the factory class of built-in recipes.

\subsection[toList/toSet/toMap]{\code{toList}/\code{toSet}/\code{toMap}}
\label{h:8:tolisttosettomap}
\begin{codebox}
\begin{Verbatim}[commandchars=\\\{\}]
\PY{n}{List}\PY{o}{\PYZlt{}}\PY{n}{String}\PY{o}{\PYZgt{}}\PY{+w}{ }\PY{n}{list}\PY{+w}{ }\PY{o}{=}\PY{+w}{ }\PY{n}{Stream}\PY{p}{.}\PY{n+na}{of}\PY{p}{(}\PY{l+s}{\PYZdq{}}\PY{l+s}{a}\PY{l+s}{\PYZdq{}}\PY{p}{,}\PY{+w}{ }\PY{l+s}{\PYZdq{}}\PY{l+s}{b}\PY{l+s}{\PYZdq{}}\PY{p}{,}\PY{+w}{ }\PY{l+s}{\PYZdq{}}\PY{l+s}{a}\PY{l+s}{\PYZdq{}}\PY{p}{)}\PY{p}{.}\PY{n+na}{collect}\PY{p}{(}\PY{n}{Collectors}\PY{p}{.}\PY{n+na}{toList}\PY{p}{(}\PY{p}{)}\PY{p}{)}\PY{p}{;}\PY{+w}{ }\PY{c+c1}{// mutable, allows duplicates}
\PY{n}{List}\PY{o}{\PYZlt{}}\PY{n}{String}\PY{o}{\PYZgt{}}\PY{+w}{ }\PY{n}{immutableList}\PY{+w}{ }\PY{o}{=}\PY{+w}{ }\PY{n}{Stream}\PY{p}{.}\PY{n+na}{of}\PY{p}{(}\PY{l+s}{\PYZdq{}}\PY{l+s}{a}\PY{l+s}{\PYZdq{}}\PY{p}{,}\PY{+w}{ }\PY{l+s}{\PYZdq{}}\PY{l+s}{b}\PY{l+s}{\PYZdq{}}\PY{p}{,}\PY{+w}{ }\PY{l+s}{\PYZdq{}}\PY{l+s}{a}\PY{l+s}{\PYZdq{}}\PY{p}{)}\PY{p}{.}\PY{n+na}{toList}\PY{p}{(}\PY{p}{)}\PY{p}{;}\PY{+w}{            }\PY{c+c1}{// Java 16+ shortcut, unmodifiable}

\PY{n}{Set}\PY{o}{\PYZlt{}}\PY{n}{String}\PY{o}{\PYZgt{}}\PY{+w}{ }\PY{n}{set}\PY{+w}{ }\PY{o}{=}\PY{+w}{ }\PY{n}{Stream}\PY{p}{.}\PY{n+na}{of}\PY{p}{(}\PY{l+s}{\PYZdq{}}\PY{l+s}{a}\PY{l+s}{\PYZdq{}}\PY{p}{,}\PY{+w}{ }\PY{l+s}{\PYZdq{}}\PY{l+s}{b}\PY{l+s}{\PYZdq{}}\PY{p}{,}\PY{+w}{ }\PY{l+s}{\PYZdq{}}\PY{l+s}{a}\PY{l+s}{\PYZdq{}}\PY{p}{)}\PY{p}{.}\PY{n+na}{collect}\PY{p}{(}\PY{n}{Collectors}\PY{p}{.}\PY{n+na}{toSet}\PY{p}{(}\PY{p}{)}\PY{p}{)}\PY{p}{;}\PY{+w}{ }\PY{c+c1}{// [a, b] \PYZhy{} duplicates dropped}

\PY{n}{Map}\PY{o}{\PYZlt{}}\PY{n}{String}\PY{p}{,}\PY{+w}{ }\PY{n}{Integer}\PY{o}{\PYZgt{}}\PY{+w}{ }\PY{n}{byLength}\PY{+w}{ }\PY{o}{=}\PY{+w}{ }\PY{n}{Stream}\PY{p}{.}\PY{n+na}{of}\PY{p}{(}\PY{l+s}{\PYZdq{}}\PY{l+s}{fig}\PY{l+s}{\PYZdq{}}\PY{p}{,}\PY{+w}{ }\PY{l+s}{\PYZdq{}}\PY{l+s}{date}\PY{l+s}{\PYZdq{}}\PY{p}{,}\PY{+w}{ }\PY{l+s}{\PYZdq{}}\PY{l+s}{kiwi}\PY{l+s}{\PYZdq{}}\PY{p}{)}
\PY{+w}{    }\PY{p}{.}\PY{n+na}{collect}\PY{p}{(}\PY{n}{Collectors}\PY{p}{.}\PY{n+na}{toMap}\PY{p}{(}\PY{n}{Function}\PY{p}{.}\PY{n+na}{identity}\PY{p}{(}\PY{p}{)}\PY{p}{,}\PY{+w}{ }\PY{n}{String}\PY{p}{:}\PY{p}{:}\PY{n}{length}\PY{p}{)}\PY{p}{)}\PY{p}{;}
\PY{n}{System}\PY{p}{.}\PY{n+na}{out}\PY{p}{.}\PY{n+na}{println}\PY{p}{(}\PY{n}{byLength}\PY{p}{)}\PY{p}{;}\PY{+w}{ }\PY{c+c1}{// \PYZob{}date=4, fig=3, kiwi=4\PYZcb{} (order not guaranteed)}
\end{Verbatim}
\end{codebox}

\code{Collectors.\allowbreak{}toMap} throws \code{IllegalStateException} if two elements produce the same key — because it has no idea how you’d like the conflict resolved. Fix it by supplying a merge function.

\begin{codebox}
\begin{Verbatim}[commandchars=\\\{\}]
\PY{n}{List}\PY{o}{\PYZlt{}}\PY{n}{String}\PY{o}{\PYZgt{}}\PY{+w}{ }\PY{n}{words}\PY{+w}{ }\PY{o}{=}\PY{+w}{ }\PY{n}{List}\PY{p}{.}\PY{n+na}{of}\PY{p}{(}\PY{l+s}{\PYZdq{}}\PY{l+s}{fig}\PY{l+s}{\PYZdq{}}\PY{p}{,}\PY{+w}{ }\PY{l+s}{\PYZdq{}}\PY{l+s}{kiwi}\PY{l+s}{\PYZdq{}}\PY{p}{,}\PY{+w}{ }\PY{l+s}{\PYZdq{}}\PY{l+s}{date}\PY{l+s}{\PYZdq{}}\PY{p}{,}\PY{+w}{ }\PY{l+s}{\PYZdq{}}\PY{l+s}{pear}\PY{l+s}{\PYZdq{}}\PY{p}{)}\PY{p}{;}\PY{+w}{ }\PY{c+c1}{// \PYZdq{}fig\PYZdq{} and \PYZdq{}pear\PYZdq{} both length 4? no \PYZhy{} kiwi/date/pear length 4}

\PY{c+c1}{// This throws: two words with the same length collide as keys}
\PY{c+c1}{// Map\PYZlt{}Integer, String\PYZgt{} byLen = words.stream()}
\PY{c+c1}{//     .collect(Collectors.toMap(String::length, w \PYZhy{}\PYZgt{} w)); // IllegalStateException: Duplicate key 4}

\PY{c+c1}{// Fix: provide a merge function to resolve collisions}
\PY{n}{Map}\PY{o}{\PYZlt{}}\PY{n}{Integer}\PY{p}{,}\PY{+w}{ }\PY{n}{String}\PY{o}{\PYZgt{}}\PY{+w}{ }\PY{n}{byLenFixed}\PY{+w}{ }\PY{o}{=}\PY{+w}{ }\PY{n}{words}\PY{p}{.}\PY{n+na}{stream}\PY{p}{(}\PY{p}{)}
\PY{+w}{    }\PY{p}{.}\PY{n+na}{collect}\PY{p}{(}\PY{n}{Collectors}\PY{p}{.}\PY{n+na}{toMap}\PY{p}{(}\PY{n}{String}\PY{p}{:}\PY{p}{:}\PY{n}{length}\PY{p}{,}\PY{+w}{ }\PY{n}{w}\PY{+w}{ }\PY{o}{\PYZhy{}}\PY{o}{\PYZgt{}}\PY{+w}{ }\PY{n}{w}\PY{p}{,}\PY{+w}{ }\PY{p}{(}\PY{n}{existing}\PY{p}{,}\PY{+w}{ }\PY{n}{incoming}\PY{p}{)}\PY{+w}{ }\PY{o}{\PYZhy{}}\PY{o}{\PYZgt{}}\PY{+w}{ }\PY{n}{existing}\PY{+w}{ }\PY{o}{+}\PY{+w}{ }\PY{l+s}{\PYZdq{}}\PY{l+s}{,}\PY{l+s}{\PYZdq{}}\PY{+w}{ }\PY{o}{+}\PY{+w}{ }\PY{n}{incoming}\PY{p}{)}\PY{p}{)}\PY{p}{;}
\PY{n}{System}\PY{p}{.}\PY{n+na}{out}\PY{p}{.}\PY{n+na}{println}\PY{p}{(}\PY{n}{byLenFixed}\PY{p}{)}\PY{p}{;}\PY{+w}{ }\PY{c+c1}{// \PYZob{}3=fig, 4=kiwi,date,pear\PYZcb{}}
\end{Verbatim}
\end{codebox}

\subsection[groupingBy with downstream collectors]{\code{groupingBy} with downstream collectors}
\label{h:8:groupingby-with-downstream-collectors}
\code{groupingBy} buckets elements by a classifier function into a \code{Map\textless{}\allowbreak{}K,\allowbreak{}~\allowbreak{}List\textless{}\allowbreak{}T\textgreater{}\textgreater{}} by default. Pass a \textbf{downstream collector} as a second argument to summarize each bucket instead of just listing it.

\begin{codebox}
\begin{Verbatim}[commandchars=\\\{\}]
\PY{k+kd}{record} \PY{n+nc}{Employee}\PY{p}{(}\PY{n}{String}\PY{+w}{ }\PY{n}{department}\PY{p}{,}\PY{+w}{ }\PY{n}{String}\PY{+w}{ }\PY{n}{name}\PY{p}{,}\PY{+w}{ }\PY{k+kt}{double}\PY{+w}{ }\PY{n}{salary}\PY{p}{)}\PY{+w}{ }\PY{p}{\PYZob{}}\PY{p}{\PYZcb{}}

\PY{n}{List}\PY{o}{\PYZlt{}}\PY{n}{Employee}\PY{o}{\PYZgt{}}\PY{+w}{ }\PY{n}{employees}\PY{+w}{ }\PY{o}{=}\PY{+w}{ }\PY{n}{List}\PY{p}{.}\PY{n+na}{of}\PY{p}{(}
\PY{+w}{    }\PY{k}{new}\PY{+w}{ }\PY{n}{Employee}\PY{p}{(}\PY{l+s}{\PYZdq{}}\PY{l+s}{Eng}\PY{l+s}{\PYZdq{}}\PY{p}{,}\PY{+w}{ }\PY{l+s}{\PYZdq{}}\PY{l+s}{Ada}\PY{l+s}{\PYZdq{}}\PY{p}{,}\PY{+w}{ }\PY{l+m+mi}{9000}\PY{p}{)}\PY{p}{,}
\PY{+w}{    }\PY{k}{new}\PY{+w}{ }\PY{n}{Employee}\PY{p}{(}\PY{l+s}{\PYZdq{}}\PY{l+s}{Eng}\PY{l+s}{\PYZdq{}}\PY{p}{,}\PY{+w}{ }\PY{l+s}{\PYZdq{}}\PY{l+s}{Linus}\PY{l+s}{\PYZdq{}}\PY{p}{,}\PY{+w}{ }\PY{l+m+mi}{9500}\PY{p}{)}\PY{p}{,}
\PY{+w}{    }\PY{k}{new}\PY{+w}{ }\PY{n}{Employee}\PY{p}{(}\PY{l+s}{\PYZdq{}}\PY{l+s}{Sales}\PY{l+s}{\PYZdq{}}\PY{p}{,}\PY{+w}{ }\PY{l+s}{\PYZdq{}}\PY{l+s}{Grace}\PY{l+s}{\PYZdq{}}\PY{p}{,}\PY{+w}{ }\PY{l+m+mi}{7000}\PY{p}{)}
\PY{p}{)}\PY{p}{;}

\PY{n}{Map}\PY{o}{\PYZlt{}}\PY{n}{String}\PY{p}{,}\PY{+w}{ }\PY{n}{List}\PY{o}{\PYZlt{}}\PY{n}{String}\PY{o}{\PYZgt{}}\PY{o}{\PYZgt{}}\PY{+w}{ }\PY{n}{namesByDept}\PY{+w}{ }\PY{o}{=}\PY{+w}{ }\PY{n}{employees}\PY{p}{.}\PY{n+na}{stream}\PY{p}{(}\PY{p}{)}
\PY{+w}{    }\PY{p}{.}\PY{n+na}{collect}\PY{p}{(}\PY{n}{Collectors}\PY{p}{.}\PY{n+na}{groupingBy}\PY{p}{(}\PY{n}{Employee}\PY{p}{:}\PY{p}{:}\PY{n}{department}\PY{p}{,}\PY{+w}{ }\PY{n}{Collectors}\PY{p}{.}\PY{n+na}{mapping}\PY{p}{(}\PY{n}{Employee}\PY{p}{:}\PY{p}{:}\PY{n}{name}\PY{p}{,}\PY{+w}{ }\PY{n}{Collectors}\PY{p}{.}\PY{n+na}{toList}\PY{p}{(}\PY{p}{)}\PY{p}{)}\PY{p}{)}\PY{p}{)}\PY{p}{;}
\PY{n}{System}\PY{p}{.}\PY{n+na}{out}\PY{p}{.}\PY{n+na}{println}\PY{p}{(}\PY{n}{namesByDept}\PY{p}{)}\PY{p}{;}\PY{+w}{ }\PY{c+c1}{// \PYZob{}Eng=[Ada, Linus], Sales=[Grace]\PYZcb{}}

\PY{n}{Map}\PY{o}{\PYZlt{}}\PY{n}{String}\PY{p}{,}\PY{+w}{ }\PY{n}{Double}\PY{o}{\PYZgt{}}\PY{+w}{ }\PY{n}{totalSalaryByDept}\PY{+w}{ }\PY{o}{=}\PY{+w}{ }\PY{n}{employees}\PY{p}{.}\PY{n+na}{stream}\PY{p}{(}\PY{p}{)}
\PY{+w}{    }\PY{p}{.}\PY{n+na}{collect}\PY{p}{(}\PY{n}{Collectors}\PY{p}{.}\PY{n+na}{groupingBy}\PY{p}{(}\PY{n}{Employee}\PY{p}{:}\PY{p}{:}\PY{n}{department}\PY{p}{,}\PY{+w}{ }\PY{n}{Collectors}\PY{p}{.}\PY{n+na}{summingDouble}\PY{p}{(}\PY{n}{Employee}\PY{p}{:}\PY{p}{:}\PY{n}{salary}\PY{p}{)}\PY{p}{)}\PY{p}{)}\PY{p}{;}
\PY{n}{System}\PY{p}{.}\PY{n+na}{out}\PY{p}{.}\PY{n+na}{println}\PY{p}{(}\PY{n}{totalSalaryByDept}\PY{p}{)}\PY{p}{;}\PY{+w}{ }\PY{c+c1}{// \PYZob{}Eng=18500.0, Sales=7000.0\PYZcb{}}
\end{Verbatim}
\end{codebox}

\subsection[partitioningBy]{\code{partitioningBy}}
\label{h:8:partitioningby}
A special case of grouping into exactly two buckets: \code{true} and \code{false}. Both keys always appear, even if a bucket is empty.

\begin{codebox}
\begin{Verbatim}[commandchars=\\\{\}]
\PY{n}{Map}\PY{o}{\PYZlt{}}\PY{n}{Boolean}\PY{p}{,}\PY{+w}{ }\PY{n}{List}\PY{o}{\PYZlt{}}\PY{n}{Integer}\PY{o}{\PYZgt{}}\PY{o}{\PYZgt{}}\PY{+w}{ }\PY{n}{evenOdd}\PY{+w}{ }\PY{o}{=}\PY{+w}{ }\PY{n}{IntStream}\PY{p}{.}\PY{n+na}{rangeClosed}\PY{p}{(}\PY{l+m+mi}{1}\PY{p}{,}\PY{+w}{ }\PY{l+m+mi}{6}\PY{p}{)}\PY{p}{.}\PY{n+na}{boxed}\PY{p}{(}\PY{p}{)}
\PY{+w}{    }\PY{p}{.}\PY{n+na}{collect}\PY{p}{(}\PY{n}{Collectors}\PY{p}{.}\PY{n+na}{partitioningBy}\PY{p}{(}\PY{n}{n}\PY{+w}{ }\PY{o}{\PYZhy{}}\PY{o}{\PYZgt{}}\PY{+w}{ }\PY{n}{n}\PY{+w}{ }\PY{o}{\PYZpc{}}\PY{+w}{ }\PY{l+m+mi}{2}\PY{+w}{ }\PY{o}{=}\PY{o}{=}\PY{+w}{ }\PY{l+m+mi}{0}\PY{p}{)}\PY{p}{)}\PY{p}{;}
\PY{n}{System}\PY{p}{.}\PY{n+na}{out}\PY{p}{.}\PY{n+na}{println}\PY{p}{(}\PY{n}{evenOdd}\PY{p}{)}\PY{p}{;}\PY{+w}{ }\PY{c+c1}{// \PYZob{}false=[1, 3, 5], true=[2, 4, 6]\PYZcb{}}
\end{Verbatim}
\end{codebox}

\subsection[joining, counting, summingInt, averagingDouble]{\code{joining}, \code{counting}, \code{summingInt}, \code{averagingDouble}}
\label{h:8:joining-counting-summingint-averagingdouble}
\begin{codebox}
\begin{Verbatim}[commandchars=\\\{\}]
\PY{n}{String}\PY{+w}{ }\PY{n}{csv}\PY{+w}{ }\PY{o}{=}\PY{+w}{ }\PY{n}{Stream}\PY{p}{.}\PY{n+na}{of}\PY{p}{(}\PY{l+s}{\PYZdq{}}\PY{l+s}{a}\PY{l+s}{\PYZdq{}}\PY{p}{,}\PY{+w}{ }\PY{l+s}{\PYZdq{}}\PY{l+s}{b}\PY{l+s}{\PYZdq{}}\PY{p}{,}\PY{+w}{ }\PY{l+s}{\PYZdq{}}\PY{l+s}{c}\PY{l+s}{\PYZdq{}}\PY{p}{)}\PY{p}{.}\PY{n+na}{collect}\PY{p}{(}\PY{n}{Collectors}\PY{p}{.}\PY{n+na}{joining}\PY{p}{(}\PY{l+s}{\PYZdq{}}\PY{l+s}{, }\PY{l+s}{\PYZdq{}}\PY{p}{,}\PY{+w}{ }\PY{l+s}{\PYZdq{}}\PY{l+s}{[}\PY{l+s}{\PYZdq{}}\PY{p}{,}\PY{+w}{ }\PY{l+s}{\PYZdq{}}\PY{l+s}{]}\PY{l+s}{\PYZdq{}}\PY{p}{)}\PY{p}{)}\PY{p}{;}
\PY{n}{System}\PY{p}{.}\PY{n+na}{out}\PY{p}{.}\PY{n+na}{println}\PY{p}{(}\PY{n}{csv}\PY{p}{)}\PY{p}{;}\PY{+w}{ }\PY{c+c1}{// [a, b, c]}

\PY{k+kt}{long}\PY{+w}{ }\PY{n}{count}\PY{+w}{ }\PY{o}{=}\PY{+w}{ }\PY{n}{Stream}\PY{p}{.}\PY{n+na}{of}\PY{p}{(}\PY{l+s}{\PYZdq{}}\PY{l+s}{a}\PY{l+s}{\PYZdq{}}\PY{p}{,}\PY{+w}{ }\PY{l+s}{\PYZdq{}}\PY{l+s}{b}\PY{l+s}{\PYZdq{}}\PY{p}{,}\PY{+w}{ }\PY{l+s}{\PYZdq{}}\PY{l+s}{c}\PY{l+s}{\PYZdq{}}\PY{p}{)}\PY{p}{.}\PY{n+na}{collect}\PY{p}{(}\PY{n}{Collectors}\PY{p}{.}\PY{n+na}{counting}\PY{p}{(}\PY{p}{)}\PY{p}{)}\PY{p}{;}\PY{+w}{       }\PY{c+c1}{// 3}
\PY{k+kt}{int}\PY{+w}{ }\PY{n}{totalLen}\PY{+w}{ }\PY{o}{=}\PY{+w}{ }\PY{n}{Stream}\PY{p}{.}\PY{n+na}{of}\PY{p}{(}\PY{l+s}{\PYZdq{}}\PY{l+s}{aa}\PY{l+s}{\PYZdq{}}\PY{p}{,}\PY{+w}{ }\PY{l+s}{\PYZdq{}}\PY{l+s}{b}\PY{l+s}{\PYZdq{}}\PY{p}{,}\PY{+w}{ }\PY{l+s}{\PYZdq{}}\PY{l+s}{ccc}\PY{l+s}{\PYZdq{}}\PY{p}{)}\PY{p}{.}\PY{n+na}{collect}\PY{p}{(}\PY{n}{Collectors}\PY{p}{.}\PY{n+na}{summingInt}\PY{p}{(}\PY{n}{String}\PY{p}{:}\PY{p}{:}\PY{n}{length}\PY{p}{)}\PY{p}{)}\PY{p}{;}\PY{+w}{ }\PY{c+c1}{// 6}
\PY{k+kt}{double}\PY{+w}{ }\PY{n}{avgLen}\PY{+w}{ }\PY{o}{=}\PY{+w}{ }\PY{n}{Stream}\PY{p}{.}\PY{n+na}{of}\PY{p}{(}\PY{l+s}{\PYZdq{}}\PY{l+s}{aa}\PY{l+s}{\PYZdq{}}\PY{p}{,}\PY{+w}{ }\PY{l+s}{\PYZdq{}}\PY{l+s}{b}\PY{l+s}{\PYZdq{}}\PY{p}{,}\PY{+w}{ }\PY{l+s}{\PYZdq{}}\PY{l+s}{ccc}\PY{l+s}{\PYZdq{}}\PY{p}{)}\PY{p}{.}\PY{n+na}{collect}\PY{p}{(}\PY{n}{Collectors}\PY{p}{.}\PY{n+na}{averagingDouble}\PY{p}{(}\PY{n}{String}\PY{p}{:}\PY{p}{:}\PY{n}{length}\PY{p}{)}\PY{p}{)}\PY{p}{;}\PY{+w}{ }\PY{c+c1}{// 2.0}
\end{Verbatim}
\end{codebox}

\subsection[mapping, flatMapping, filtering]{\code{mapping}, \code{flatMapping}, \code{filtering}}
\label{h:8:mapping-flatmapping-filtering}
These are \textbf{downstream collectors} — meant to be nested inside \code{groupingBy}/\code{partitioningBy} to reshape or pre-process each bucket before the final collector runs.

\begin{codebox}
\begin{Verbatim}[commandchars=\\\{\}]
\PY{k+kd}{record} \PY{n+nc}{Order}\PY{p}{(}\PY{n}{String}\PY{+w}{ }\PY{n}{customer}\PY{p}{,}\PY{+w}{ }\PY{n}{List}\PY{o}{\PYZlt{}}\PY{n}{String}\PY{o}{\PYZgt{}}\PY{+w}{ }\PY{n}{items}\PY{p}{)}\PY{+w}{ }\PY{p}{\PYZob{}}\PY{p}{\PYZcb{}}

\PY{n}{List}\PY{o}{\PYZlt{}}\PY{n}{Order}\PY{o}{\PYZgt{}}\PY{+w}{ }\PY{n}{orders}\PY{+w}{ }\PY{o}{=}\PY{+w}{ }\PY{n}{List}\PY{p}{.}\PY{n+na}{of}\PY{p}{(}
\PY{+w}{    }\PY{k}{new}\PY{+w}{ }\PY{n}{Order}\PY{p}{(}\PY{l+s}{\PYZdq{}}\PY{l+s}{Ada}\PY{l+s}{\PYZdq{}}\PY{p}{,}\PY{+w}{ }\PY{n}{List}\PY{p}{.}\PY{n+na}{of}\PY{p}{(}\PY{l+s}{\PYZdq{}}\PY{l+s}{pen}\PY{l+s}{\PYZdq{}}\PY{p}{,}\PY{+w}{ }\PY{l+s}{\PYZdq{}}\PY{l+s}{paper}\PY{l+s}{\PYZdq{}}\PY{p}{)}\PY{p}{)}\PY{p}{,}
\PY{+w}{    }\PY{k}{new}\PY{+w}{ }\PY{n}{Order}\PY{p}{(}\PY{l+s}{\PYZdq{}}\PY{l+s}{Ada}\PY{l+s}{\PYZdq{}}\PY{p}{,}\PY{+w}{ }\PY{n}{List}\PY{p}{.}\PY{n+na}{of}\PY{p}{(}\PY{l+s}{\PYZdq{}}\PY{l+s}{stapler}\PY{l+s}{\PYZdq{}}\PY{p}{)}\PY{p}{)}\PY{p}{,}
\PY{+w}{    }\PY{k}{new}\PY{+w}{ }\PY{n}{Order}\PY{p}{(}\PY{l+s}{\PYZdq{}}\PY{l+s}{Grace}\PY{l+s}{\PYZdq{}}\PY{p}{,}\PY{+w}{ }\PY{n}{List}\PY{p}{.}\PY{n+na}{of}\PY{p}{(}\PY{l+s}{\PYZdq{}}\PY{l+s}{pen}\PY{l+s}{\PYZdq{}}\PY{p}{)}\PY{p}{)}
\PY{p}{)}\PY{p}{;}

\PY{c+c1}{// flatMapping: flatten each customer\PYZsq{}s nested item lists into one flat set}
\PY{n}{Map}\PY{o}{\PYZlt{}}\PY{n}{String}\PY{p}{,}\PY{+w}{ }\PY{n}{Set}\PY{o}{\PYZlt{}}\PY{n}{String}\PY{o}{\PYZgt{}}\PY{o}{\PYZgt{}}\PY{+w}{ }\PY{n}{itemsByCustomer}\PY{+w}{ }\PY{o}{=}\PY{+w}{ }\PY{n}{orders}\PY{p}{.}\PY{n+na}{stream}\PY{p}{(}\PY{p}{)}
\PY{+w}{    }\PY{p}{.}\PY{n+na}{collect}\PY{p}{(}\PY{n}{Collectors}\PY{p}{.}\PY{n+na}{groupingBy}\PY{p}{(}\PY{n}{Order}\PY{p}{:}\PY{p}{:}\PY{n}{customer}\PY{p}{,}
\PY{+w}{             }\PY{n}{Collectors}\PY{p}{.}\PY{n+na}{flatMapping}\PY{p}{(}\PY{n}{o}\PY{+w}{ }\PY{o}{\PYZhy{}}\PY{o}{\PYZgt{}}\PY{+w}{ }\PY{n}{o}\PY{p}{.}\PY{n+na}{items}\PY{p}{(}\PY{p}{)}\PY{p}{.}\PY{n+na}{stream}\PY{p}{(}\PY{p}{)}\PY{p}{,}\PY{+w}{ }\PY{n}{Collectors}\PY{p}{.}\PY{n+na}{toSet}\PY{p}{(}\PY{p}{)}\PY{p}{)}\PY{p}{)}\PY{p}{)}\PY{p}{;}
\PY{n}{System}\PY{p}{.}\PY{n+na}{out}\PY{p}{.}\PY{n+na}{println}\PY{p}{(}\PY{n}{itemsByCustomer}\PY{p}{)}\PY{p}{;}\PY{+w}{ }\PY{c+c1}{// \PYZob{}Ada=[pen, paper, stapler], Grace=[pen]\PYZcb{}}

\PY{c+c1}{// filtering: keep only orders with more than one item, per customer, before collecting}
\PY{n}{Map}\PY{o}{\PYZlt{}}\PY{n}{String}\PY{p}{,}\PY{+w}{ }\PY{n}{List}\PY{o}{\PYZlt{}}\PY{n}{Order}\PY{o}{\PYZgt{}}\PY{o}{\PYZgt{}}\PY{+w}{ }\PY{n}{multiItemOrdersByCustomer}\PY{+w}{ }\PY{o}{=}\PY{+w}{ }\PY{n}{orders}\PY{p}{.}\PY{n+na}{stream}\PY{p}{(}\PY{p}{)}
\PY{+w}{    }\PY{p}{.}\PY{n+na}{collect}\PY{p}{(}\PY{n}{Collectors}\PY{p}{.}\PY{n+na}{groupingBy}\PY{p}{(}\PY{n}{Order}\PY{p}{:}\PY{p}{:}\PY{n}{customer}\PY{p}{,}
\PY{+w}{             }\PY{n}{Collectors}\PY{p}{.}\PY{n+na}{filtering}\PY{p}{(}\PY{n}{o}\PY{+w}{ }\PY{o}{\PYZhy{}}\PY{o}{\PYZgt{}}\PY{+w}{ }\PY{n}{o}\PY{p}{.}\PY{n+na}{items}\PY{p}{(}\PY{p}{)}\PY{p}{.}\PY{n+na}{size}\PY{p}{(}\PY{p}{)}\PY{+w}{ }\PY{o}{\PYZgt{}}\PY{+w}{ }\PY{l+m+mi}{1}\PY{p}{,}\PY{+w}{ }\PY{n}{Collectors}\PY{p}{.}\PY{n+na}{toList}\PY{p}{(}\PY{p}{)}\PY{p}{)}\PY{p}{)}\PY{p}{)}\PY{p}{;}
\PY{n}{System}\PY{p}{.}\PY{n+na}{out}\PY{p}{.}\PY{n+na}{println}\PY{p}{(}\PY{n}{multiItemOrdersByCustomer}\PY{p}{)}\PY{p}{;}\PY{+w}{ }\PY{c+c1}{// \PYZob{}Ada=[Order[customer=Ada, items=[pen, paper]]], Grace=[]\PYZcb{}}
\end{Verbatim}
\end{codebox}

Note the difference from \code{stream().\allowbreak{}filter(...)} \emph{before} \code{groupingBy}: filtering before grouping can make a whole group disappear; \code{Collectors.\allowbreak{}filtering} keeps the group key present (as shown above, \code{Grace=[]} still shows up).

\subsection[teeing]{\code{teeing}}
\label{h:8:teeing}
\code{teeing} (Java 12+) runs a stream through \emph{two} collectors simultaneously and merges their results with a \code{BiFunction}. Useful when you need two different summaries from a single pass.

\begin{codebox}
\begin{Verbatim}[commandchars=\\\{\}]
\PY{k+kd}{record} \PY{n+nc}{MinMax}\PY{p}{(}\PY{k+kt}{int}\PY{+w}{ }\PY{n}{min}\PY{p}{,}\PY{+w}{ }\PY{k+kt}{int}\PY{+w}{ }\PY{n}{max}\PY{p}{)}\PY{+w}{ }\PY{p}{\PYZob{}}\PY{p}{\PYZcb{}}

\PY{n}{MinMax}\PY{+w}{ }\PY{n}{result}\PY{+w}{ }\PY{o}{=}\PY{+w}{ }\PY{n}{Stream}\PY{p}{.}\PY{n+na}{of}\PY{p}{(}\PY{l+m+mi}{4}\PY{p}{,}\PY{+w}{ }\PY{l+m+mi}{1}\PY{p}{,}\PY{+w}{ }\PY{l+m+mi}{7}\PY{p}{,}\PY{+w}{ }\PY{l+m+mi}{3}\PY{p}{,}\PY{+w}{ }\PY{l+m+mi}{9}\PY{p}{,}\PY{+w}{ }\PY{l+m+mi}{2}\PY{p}{)}
\PY{+w}{    }\PY{p}{.}\PY{n+na}{collect}\PY{p}{(}\PY{n}{Collectors}\PY{p}{.}\PY{n+na}{teeing}\PY{p}{(}
\PY{+w}{        }\PY{n}{Collectors}\PY{p}{.}\PY{n+na}{minBy}\PY{p}{(}\PY{n}{Integer}\PY{p}{:}\PY{p}{:}\PY{n}{compareTo}\PY{p}{)}\PY{p}{,}
\PY{+w}{        }\PY{n}{Collectors}\PY{p}{.}\PY{n+na}{maxBy}\PY{p}{(}\PY{n}{Integer}\PY{p}{:}\PY{p}{:}\PY{n}{compareTo}\PY{p}{)}\PY{p}{,}
\PY{+w}{        }\PY{p}{(}\PY{n}{min}\PY{p}{,}\PY{+w}{ }\PY{n}{max}\PY{p}{)}\PY{+w}{ }\PY{o}{\PYZhy{}}\PY{o}{\PYZgt{}}\PY{+w}{ }\PY{k}{new}\PY{+w}{ }\PY{n}{MinMax}\PY{p}{(}\PY{n}{min}\PY{p}{.}\PY{n+na}{orElseThrow}\PY{p}{(}\PY{p}{)}\PY{p}{,}\PY{+w}{ }\PY{n}{max}\PY{p}{.}\PY{n+na}{orElseThrow}\PY{p}{(}\PY{p}{)}\PY{p}{)}
\PY{+w}{    }\PY{p}{)}\PY{p}{)}\PY{p}{;}
\PY{n}{System}\PY{p}{.}\PY{n+na}{out}\PY{p}{.}\PY{n+na}{println}\PY{p}{(}\PY{n}{result}\PY{p}{)}\PY{p}{;}\PY{+w}{ }\PY{c+c1}{// MinMax[min=1, max=9]}
\end{Verbatim}
\end{codebox}

\subsection[collectingAndThen]{\code{collectingAndThen}}
\label{h:8:collectingandthen}
Wraps a collector and applies a finishing transformation to its result — commonly used to make a mutable result immutable.

\begin{codebox}
\begin{Verbatim}[commandchars=\\\{\}]
\PY{n}{List}\PY{o}{\PYZlt{}}\PY{n}{String}\PY{o}{\PYZgt{}}\PY{+w}{ }\PY{n}{immutableUpper}\PY{+w}{ }\PY{o}{=}\PY{+w}{ }\PY{n}{Stream}\PY{p}{.}\PY{n+na}{of}\PY{p}{(}\PY{l+s}{\PYZdq{}}\PY{l+s}{b}\PY{l+s}{\PYZdq{}}\PY{p}{,}\PY{+w}{ }\PY{l+s}{\PYZdq{}}\PY{l+s}{a}\PY{l+s}{\PYZdq{}}\PY{p}{,}\PY{+w}{ }\PY{l+s}{\PYZdq{}}\PY{l+s}{c}\PY{l+s}{\PYZdq{}}\PY{p}{)}
\PY{+w}{    }\PY{p}{.}\PY{n+na}{collect}\PY{p}{(}\PY{n}{Collectors}\PY{p}{.}\PY{n+na}{collectingAndThen}\PY{p}{(}\PY{n}{Collectors}\PY{p}{.}\PY{n+na}{toList}\PY{p}{(}\PY{p}{)}\PY{p}{,}\PY{+w}{ }\PY{n}{Collections}\PY{p}{:}\PY{p}{:}\PY{n}{unmodifiableList}\PY{p}{)}\PY{p}{)}\PY{p}{;}
\PY{c+c1}{// immutableUpper.add(\PYZdq{}x\PYZdq{}); // throws UnsupportedOperationException}
\end{Verbatim}
\end{codebox}

\subsection[Writing a custom Collector]{Writing a custom Collector}
\label{h:8:writing-a-custom-collector}
\code{Collector\textless{}\allowbreak{}T,\allowbreak{}~\allowbreak{}A,\allowbreak{}~\allowbreak{}R\textgreater{}} needs four ingredients: a \textbf{supplier} (create the mutable accumulator), an \textbf{accumulator} (fold one element in), a \textbf{combiner} (merge two accumulators — used in parallel streams), and a \textbf{finisher} (convert the accumulator to the final result).

\begin{codebox}
\begin{Verbatim}[commandchars=\\\{\}]
\PY{n}{Collector}\PY{o}{\PYZlt{}}\PY{n}{String}\PY{p}{,}\PY{+w}{ }\PY{n}{StringBuilder}\PY{p}{,}\PY{+w}{ }\PY{n}{String}\PY{o}{\PYZgt{}}\PY{+w}{ }\PY{n}{customJoiner}\PY{+w}{ }\PY{o}{=}\PY{+w}{ }\PY{n}{Collector}\PY{p}{.}\PY{n+na}{of}\PY{p}{(}
\PY{+w}{    }\PY{n}{StringBuilder}\PY{p}{:}\PY{p}{:}\PY{k}{new}\PY{p}{,}\PY{+w}{                    }\PY{c+c1}{// supplier}
\PY{+w}{    }\PY{p}{(}\PY{n}{sb}\PY{p}{,}\PY{+w}{ }\PY{n}{s}\PY{p}{)}\PY{+w}{ }\PY{o}{\PYZhy{}}\PY{o}{\PYZgt{}}\PY{+w}{ }\PY{n}{sb}\PY{p}{.}\PY{n+na}{append}\PY{p}{(}\PY{n}{s}\PY{p}{)}\PY{p}{.}\PY{n+na}{append}\PY{p}{(}\PY{l+s+sc}{\PYZsq{}|\PYZsq{}}\PY{p}{)}\PY{p}{,}\PY{+w}{   }\PY{c+c1}{// accumulator}
\PY{+w}{    }\PY{p}{(}\PY{n}{sb1}\PY{p}{,}\PY{+w}{ }\PY{n}{sb2}\PY{p}{)}\PY{+w}{ }\PY{o}{\PYZhy{}}\PY{o}{\PYZgt{}}\PY{+w}{ }\PY{n}{sb1}\PY{p}{.}\PY{n+na}{append}\PY{p}{(}\PY{n}{sb2}\PY{p}{)}\PY{p}{,}\PY{+w}{         }\PY{c+c1}{// combiner (for parallel merging)}
\PY{+w}{    }\PY{n}{StringBuilder}\PY{p}{:}\PY{p}{:}\PY{n}{toString}\PY{+w}{                }\PY{c+c1}{// finisher}
\PY{p}{)}\PY{p}{;}

\PY{n}{String}\PY{+w}{ }\PY{n}{result}\PY{+w}{ }\PY{o}{=}\PY{+w}{ }\PY{n}{Stream}\PY{p}{.}\PY{n+na}{of}\PY{p}{(}\PY{l+s}{\PYZdq{}}\PY{l+s}{a}\PY{l+s}{\PYZdq{}}\PY{p}{,}\PY{+w}{ }\PY{l+s}{\PYZdq{}}\PY{l+s}{b}\PY{l+s}{\PYZdq{}}\PY{p}{,}\PY{+w}{ }\PY{l+s}{\PYZdq{}}\PY{l+s}{c}\PY{l+s}{\PYZdq{}}\PY{p}{)}\PY{p}{.}\PY{n+na}{collect}\PY{p}{(}\PY{n}{customJoiner}\PY{p}{)}\PY{p}{;}
\PY{n}{System}\PY{p}{.}\PY{n+na}{out}\PY{p}{.}\PY{n+na}{println}\PY{p}{(}\PY{n}{result}\PY{p}{)}\PY{p}{;}\PY{+w}{ }\PY{c+c1}{// a|b|c|}
\end{Verbatim}
\end{codebox}

For most cases, though, composing the built-in collectors (as shown above) is clearer and less error-prone than hand-rolling one.

\section[Spliterator]{Spliterator}
\label{h:8:spliterator}
A \textbf{\code{Spliterator}} (“splitable iterator”) is the low-level engine that powers both sequential and parallel stream traversal. You rarely write one yourself, but understanding it explains \emph{why} some sources parallelize well and others don’t.

\begin{itemize}
\item 
\code{try\allowbreak{}Advance(\allowbreak{}Consumer\textless{}?\allowbreak{}~\allowbreak{}super~\allowbreak{}T\textgreater{}\allowbreak{}~\allowbreak{}action)}: process the next element if there is one and return \code{true}; return \code{false} if the source is exhausted. This is how sequential traversal works — one element at a time, like \code{Iterator.\allowbreak{}next()} but pushing the value into a callback instead of returning it.

\item 
\code{trySplit()}: attempt to carve off a chunk of the remaining elements into a \emph{new} \code{Spliterator} that a different thread can process, returning \code{null} if the source can’t (or won’t) be split further. This is the method parallel streams call to divide work across the common fork/join pool.

\item 
\textbf{Characteristics}: a bitmask describing properties of the source — \code{ORDERED}, \code{SORTED}, \code{SIZED} (size known up front), \code{SUBSIZED} (splits also report exact size), \code{DISTINCT}, \code{NONNULL}, \code{IMMUTABLE}, \code{CONCURRENT}. Stream operations use these as hints to skip unnecessary work — e.g., \code{distinct()} on a stream whose spliterator reports \code{DISTINCT} is a no-op.

\end{itemize}

\begin{codebox}
\begin{Verbatim}[commandchars=\\\{\}]
\PY{n}{List}\PY{o}{\PYZlt{}}\PY{n}{Integer}\PY{o}{\PYZgt{}}\PY{+w}{ }\PY{n}{data}\PY{+w}{ }\PY{o}{=}\PY{+w}{ }\PY{n}{List}\PY{p}{.}\PY{n+na}{of}\PY{p}{(}\PY{l+m+mi}{1}\PY{p}{,}\PY{+w}{ }\PY{l+m+mi}{2}\PY{p}{,}\PY{+w}{ }\PY{l+m+mi}{3}\PY{p}{,}\PY{+w}{ }\PY{l+m+mi}{4}\PY{p}{,}\PY{+w}{ }\PY{l+m+mi}{5}\PY{p}{,}\PY{+w}{ }\PY{l+m+mi}{6}\PY{p}{,}\PY{+w}{ }\PY{l+m+mi}{7}\PY{p}{,}\PY{+w}{ }\PY{l+m+mi}{8}\PY{p}{)}\PY{p}{;}
\PY{n}{Spliterator}\PY{o}{\PYZlt{}}\PY{n}{Integer}\PY{o}{\PYZgt{}}\PY{+w}{ }\PY{n}{spliterator}\PY{+w}{ }\PY{o}{=}\PY{+w}{ }\PY{n}{data}\PY{p}{.}\PY{n+na}{spliterator}\PY{p}{(}\PY{p}{)}\PY{p}{;}

\PY{n}{System}\PY{p}{.}\PY{n+na}{out}\PY{p}{.}\PY{n+na}{println}\PY{p}{(}\PY{n}{spliterator}\PY{p}{.}\PY{n+na}{estimateSize}\PY{p}{(}\PY{p}{)}\PY{p}{)}\PY{p}{;}\PY{+w}{ }\PY{c+c1}{// 8}
\PY{n}{System}\PY{p}{.}\PY{n+na}{out}\PY{p}{.}\PY{n+na}{println}\PY{p}{(}\PY{n}{spliterator}\PY{p}{.}\PY{n+na}{characteristics}\PY{p}{(}\PY{p}{)}\PY{+w}{ }\PY{o}{\PYZam{}}\PY{+w}{ }\PY{n}{Spliterator}\PY{p}{.}\PY{n+na}{SIZED}\PY{p}{)}\PY{p}{;}\PY{+w}{ }\PY{c+c1}{// non\PYZhy{}zero: SIZED is set}

\PY{n}{Spliterator}\PY{o}{\PYZlt{}}\PY{n}{Integer}\PY{o}{\PYZgt{}}\PY{+w}{ }\PY{n}{firstHalf}\PY{+w}{ }\PY{o}{=}\PY{+w}{ }\PY{n}{spliterator}\PY{p}{.}\PY{n+na}{trySplit}\PY{p}{(}\PY{p}{)}\PY{p}{;}\PY{+w}{ }\PY{c+c1}{// spliterator now covers the second half}
\PY{n}{firstHalf}\PY{p}{.}\PY{n+na}{forEachRemaining}\PY{p}{(}\PY{n}{n}\PY{+w}{ }\PY{o}{\PYZhy{}}\PY{o}{\PYZgt{}}\PY{+w}{ }\PY{n}{System}\PY{p}{.}\PY{n+na}{out}\PY{p}{.}\PY{n+na}{print}\PY{p}{(}\PY{n}{n}\PY{+w}{ }\PY{o}{+}\PY{+w}{ }\PY{l+s}{\PYZdq{}}\PY{l+s}{ }\PY{l+s}{\PYZdq{}}\PY{p}{)}\PY{p}{)}\PY{p}{;}\PY{+w}{ }\PY{c+c1}{// 1 2 3 4}
\PY{n}{System}\PY{p}{.}\PY{n+na}{out}\PY{p}{.}\PY{n+na}{println}\PY{p}{(}\PY{p}{)}\PY{p}{;}
\PY{n}{spliterator}\PY{p}{.}\PY{n+na}{forEachRemaining}\PY{p}{(}\PY{n}{n}\PY{+w}{ }\PY{o}{\PYZhy{}}\PY{o}{\PYZgt{}}\PY{+w}{ }\PY{n}{System}\PY{p}{.}\PY{n+na}{out}\PY{p}{.}\PY{n+na}{print}\PY{p}{(}\PY{n}{n}\PY{+w}{ }\PY{o}{+}\PY{+w}{ }\PY{l+s}{\PYZdq{}}\PY{l+s}{ }\PY{l+s}{\PYZdq{}}\PY{p}{)}\PY{p}{)}\PY{p}{;}\PY{+w}{ }\PY{c+c1}{// 5 6 7 8}
\end{Verbatim}
\end{codebox}

\subsection[When you’d implement one]{When you’d implement one}
\label{h:8:when-youd-implement-one}
You’d write a custom \code{Spliterator} when you’re exposing a custom data source as a \code{Stream} — for example, streaming rows from a specialized data structure (a skip list, a memory-mapped file, a custom ring buffer) and you want that stream to parallelize efficiently. You’d implement \code{trySplit()} to divide the source in a way that respects its actual structure, and report accurate characteristics so downstream operations can optimize correctly. Getting \code{SIZED}/\code{SUBSIZED} wrong, for instance, can make parallel splitting produce badly unbalanced chunks.

\begin{codebox}
\begin{Verbatim}[commandchars=\\\{\}]
\PY{k+kd}{class} \PY{n+nc}{RangeSpliterator}\PY{+w}{ }\PY{k+kd}{implements}\PY{+w}{ }\PY{n}{Spliterator}\PY{o}{\PYZlt{}}\PY{n}{Integer}\PY{o}{\PYZgt{}}\PY{+w}{ }\PY{p}{\PYZob{}}
\PY{+w}{    }\PY{k+kd}{private}\PY{+w}{ }\PY{k+kt}{int}\PY{+w}{ }\PY{n}{current}\PY{p}{;}
\PY{+w}{    }\PY{k+kd}{private}\PY{+w}{ }\PY{k+kd}{final}\PY{+w}{ }\PY{k+kt}{int}\PY{+w}{ }\PY{n}{end}\PY{p}{;}

\PY{+w}{    }\PY{n}{RangeSpliterator}\PY{p}{(}\PY{k+kt}{int}\PY{+w}{ }\PY{n}{start}\PY{p}{,}\PY{+w}{ }\PY{k+kt}{int}\PY{+w}{ }\PY{n}{end}\PY{p}{)}\PY{+w}{ }\PY{p}{\PYZob{}}
\PY{+w}{        }\PY{k}{this}\PY{p}{.}\PY{n+na}{current}\PY{+w}{ }\PY{o}{=}\PY{+w}{ }\PY{n}{start}\PY{p}{;}
\PY{+w}{        }\PY{k}{this}\PY{p}{.}\PY{n+na}{end}\PY{+w}{ }\PY{o}{=}\PY{+w}{ }\PY{n}{end}\PY{p}{;}
\PY{+w}{    }\PY{p}{\PYZcb{}}

\PY{+w}{    }\PY{n+nd}{@Override}
\PY{+w}{    }\PY{k+kd}{public}\PY{+w}{ }\PY{k+kt}{boolean}\PY{+w}{ }\PY{n+nf}{tryAdvance}\PY{p}{(}\PY{n}{Consumer}\PY{o}{\PYZlt{}}\PY{o}{?}\PY{+w}{ }\PY{k+kd}{super}\PY{+w}{ }\PY{n}{Integer}\PY{o}{\PYZgt{}}\PY{+w}{ }\PY{n}{action}\PY{p}{)}\PY{+w}{ }\PY{p}{\PYZob{}}
\PY{+w}{        }\PY{k}{if}\PY{+w}{ }\PY{p}{(}\PY{n}{current}\PY{+w}{ }\PY{o}{\PYZgt{}}\PY{o}{=}\PY{+w}{ }\PY{n}{end}\PY{p}{)}\PY{+w}{ }\PY{k}{return}\PY{+w}{ }\PY{k+kc}{false}\PY{p}{;}
\PY{+w}{        }\PY{n}{action}\PY{p}{.}\PY{n+na}{accept}\PY{p}{(}\PY{n}{current}\PY{o}{+}\PY{o}{+}\PY{p}{)}\PY{p}{;}
\PY{+w}{        }\PY{k}{return}\PY{+w}{ }\PY{k+kc}{true}\PY{p}{;}
\PY{+w}{    }\PY{p}{\PYZcb{}}

\PY{+w}{    }\PY{n+nd}{@Override}
\PY{+w}{    }\PY{k+kd}{public}\PY{+w}{ }\PY{n}{Spliterator}\PY{o}{\PYZlt{}}\PY{n}{Integer}\PY{o}{\PYZgt{}}\PY{+w}{ }\PY{n+nf}{trySplit}\PY{p}{(}\PY{p}{)}\PY{+w}{ }\PY{p}{\PYZob{}}
\PY{+w}{        }\PY{k+kt}{int}\PY{+w}{ }\PY{n}{remaining}\PY{+w}{ }\PY{o}{=}\PY{+w}{ }\PY{n}{end}\PY{+w}{ }\PY{o}{\PYZhy{}}\PY{+w}{ }\PY{n}{current}\PY{p}{;}
\PY{+w}{        }\PY{k}{if}\PY{+w}{ }\PY{p}{(}\PY{n}{remaining}\PY{+w}{ }\PY{o}{\PYZlt{}}\PY{+w}{ }\PY{l+m+mi}{2}\PY{p}{)}\PY{+w}{ }\PY{k}{return}\PY{+w}{ }\PY{k+kc}{null}\PY{p}{;}\PY{+w}{ }\PY{c+c1}{// too small to split}
\PY{+w}{        }\PY{k+kt}{int}\PY{+w}{ }\PY{n}{mid}\PY{+w}{ }\PY{o}{=}\PY{+w}{ }\PY{n}{current}\PY{+w}{ }\PY{o}{+}\PY{+w}{ }\PY{n}{remaining}\PY{+w}{ }\PY{o}{/}\PY{+w}{ }\PY{l+m+mi}{2}\PY{p}{;}
\PY{+w}{        }\PY{n}{Spliterator}\PY{o}{\PYZlt{}}\PY{n}{Integer}\PY{o}{\PYZgt{}}\PY{+w}{ }\PY{n}{firstHalf}\PY{+w}{ }\PY{o}{=}\PY{+w}{ }\PY{k}{new}\PY{+w}{ }\PY{n}{RangeSpliterator}\PY{p}{(}\PY{n}{current}\PY{p}{,}\PY{+w}{ }\PY{n}{mid}\PY{p}{)}\PY{p}{;}
\PY{+w}{        }\PY{k}{this}\PY{p}{.}\PY{n+na}{current}\PY{+w}{ }\PY{o}{=}\PY{+w}{ }\PY{n}{mid}\PY{p}{;}\PY{+w}{ }\PY{c+c1}{// this spliterator keeps the second half}
\PY{+w}{        }\PY{k}{return}\PY{+w}{ }\PY{n}{firstHalf}\PY{p}{;}
\PY{+w}{    }\PY{p}{\PYZcb{}}

\PY{+w}{    }\PY{n+nd}{@Override}
\PY{+w}{    }\PY{k+kd}{public}\PY{+w}{ }\PY{k+kt}{long}\PY{+w}{ }\PY{n+nf}{estimateSize}\PY{p}{(}\PY{p}{)}\PY{+w}{ }\PY{p}{\PYZob{}}\PY{+w}{ }\PY{k}{return}\PY{+w}{ }\PY{n}{end}\PY{+w}{ }\PY{o}{\PYZhy{}}\PY{+w}{ }\PY{n}{current}\PY{p}{;}\PY{+w}{ }\PY{p}{\PYZcb{}}

\PY{+w}{    }\PY{n+nd}{@Override}
\PY{+w}{    }\PY{k+kd}{public}\PY{+w}{ }\PY{k+kt}{int}\PY{+w}{ }\PY{n+nf}{characteristics}\PY{p}{(}\PY{p}{)}\PY{+w}{ }\PY{p}{\PYZob{}}
\PY{+w}{        }\PY{k}{return}\PY{+w}{ }\PY{n}{ORDERED}\PY{+w}{ }\PY{o}{|}\PY{+w}{ }\PY{n}{SIZED}\PY{+w}{ }\PY{o}{|}\PY{+w}{ }\PY{n}{SUBSIZED}\PY{+w}{ }\PY{o}{|}\PY{+w}{ }\PY{n}{IMMUTABLE}\PY{+w}{ }\PY{o}{|}\PY{+w}{ }\PY{n}{NONNULL}\PY{p}{;}
\PY{+w}{    }\PY{p}{\PYZcb{}}
\PY{p}{\PYZcb{}}

\PY{n}{Stream}\PY{o}{\PYZlt{}}\PY{n}{Integer}\PY{o}{\PYZgt{}}\PY{+w}{ }\PY{n}{customStream}\PY{+w}{ }\PY{o}{=}\PY{+w}{ }\PY{n}{StreamSupport}\PY{p}{.}\PY{n+na}{stream}\PY{p}{(}\PY{k}{new}\PY{+w}{ }\PY{n}{RangeSpliterator}\PY{p}{(}\PY{l+m+mi}{0}\PY{p}{,}\PY{+w}{ }\PY{l+m+mi}{10}\PY{p}{)}\PY{p}{,}\PY{+w}{ }\PY{k+kc}{false}\PY{p}{)}\PY{p}{;}
\PY{n}{System}\PY{p}{.}\PY{n+na}{out}\PY{p}{.}\PY{n+na}{println}\PY{p}{(}\PY{n}{customStream}\PY{p}{.}\PY{n+na}{mapToInt}\PY{p}{(}\PY{n}{Integer}\PY{p}{:}\PY{p}{:}\PY{n}{intValue}\PY{p}{)}\PY{p}{.}\PY{n+na}{sum}\PY{p}{(}\PY{p}{)}\PY{p}{)}\PY{p}{;}\PY{+w}{ }\PY{c+c1}{// 45}
\end{Verbatim}
\end{codebox}

\section[Stream Performance]{Stream Performance}
\label{h:8:stream-performance}
Streams read nicely, but they are not free. Every stage adds a virtual call, and for simple operations on already-in-memory data, a classic \code{for} loop is usually faster and, more importantly, has fewer surprising costs.

\subsection[When streams are slower than loops]{When streams are slower than loops}
\label{h:8:when-streams-are-slower-than-loops}
For small collections or trivial operations, the overhead of building the pipeline (lambda allocation, boxing, virtual dispatch through each stage) can outweigh the benefit of readability. In hot paths — code that runs millions of times per second — measure before assuming a stream refactor is free.

\begin{codebox}
\begin{Verbatim}[commandchars=\\\{\}]
\PY{k+kt}{int}\PY{o}{[}\PY{o}{]}\PY{+w}{ }\PY{n}{data}\PY{+w}{ }\PY{o}{=}\PY{+w}{ }\PY{p}{\PYZob{}}\PY{l+m+mi}{1}\PY{p}{,}\PY{+w}{ }\PY{l+m+mi}{2}\PY{p}{,}\PY{+w}{ }\PY{l+m+mi}{3}\PY{p}{,}\PY{+w}{ }\PY{l+m+mi}{4}\PY{p}{,}\PY{+w}{ }\PY{l+m+mi}{5}\PY{p}{\PYZcb{}}\PY{p}{;}

\PY{c+c1}{// Loop: simple, no allocation beyond the primitive array}
\PY{k+kt}{int}\PY{+w}{ }\PY{n}{sumLoop}\PY{+w}{ }\PY{o}{=}\PY{+w}{ }\PY{l+m+mi}{0}\PY{p}{;}
\PY{k}{for}\PY{+w}{ }\PY{p}{(}\PY{k+kt}{int}\PY{+w}{ }\PY{n}{n}\PY{+w}{ }\PY{p}{:}\PY{+w}{ }\PY{n}{data}\PY{p}{)}\PY{+w}{ }\PY{n}{sumLoop}\PY{+w}{ }\PY{o}{+}\PY{o}{=}\PY{+w}{ }\PY{n}{n}\PY{p}{;}

\PY{c+c1}{// Stream: clean, but goes through IntStream machinery + a lambda for a trivial task}
\PY{k+kt}{int}\PY{+w}{ }\PY{n}{sumStream}\PY{+w}{ }\PY{o}{=}\PY{+w}{ }\PY{n}{IntStream}\PY{p}{.}\PY{n+na}{of}\PY{p}{(}\PY{n}{data}\PY{p}{)}\PY{p}{.}\PY{n+na}{sum}\PY{p}{(}\PY{p}{)}\PY{p}{;}
\end{Verbatim}
\end{codebox}

Both are fine here — the point is that for tiny, simple aggregations, a loop is not something to apologize for in review.

\subsection[Boxing costs]{Boxing costs}
\label{h:8:boxing-costs}
\code{Stream\textless{}\allowbreak{}Integer\textgreater{}} stores boxed \code{Integer} objects; every \code{int} you put into it gets wrapped in an object, and every arithmetic operation on it may need to unbox, compute, and re-box. \code{IntStream}/\code{LongStream}/\code{DoubleStream} avoid this entirely by keeping primitives unboxed through the pipeline.

\begin{codebox}
\begin{Verbatim}[commandchars=\\\{\}]
\PY{n}{List}\PY{o}{\PYZlt{}}\PY{n}{Integer}\PY{o}{\PYZgt{}}\PY{+w}{ }\PY{n}{numbers}\PY{+w}{ }\PY{o}{=}\PY{+w}{ }\PY{n}{List}\PY{p}{.}\PY{n+na}{of}\PY{p}{(}\PY{l+m+mi}{1}\PY{p}{,}\PY{+w}{ }\PY{l+m+mi}{2}\PY{p}{,}\PY{+w}{ }\PY{l+m+mi}{3}\PY{p}{,}\PY{+w}{ }\PY{l+m+mi}{4}\PY{p}{,}\PY{+w}{ }\PY{l+m+mi}{5}\PY{p}{)}\PY{p}{;}

\PY{c+c1}{// Boxes every element into Integer, unboxes for the sum}
\PY{k+kt}{int}\PY{+w}{ }\PY{n}{slow}\PY{+w}{ }\PY{o}{=}\PY{+w}{ }\PY{n}{numbers}\PY{p}{.}\PY{n+na}{stream}\PY{p}{(}\PY{p}{)}\PY{p}{.}\PY{n+na}{reduce}\PY{p}{(}\PY{l+m+mi}{0}\PY{p}{,}\PY{+w}{ }\PY{n}{Integer}\PY{p}{:}\PY{p}{:}\PY{n}{sum}\PY{p}{)}\PY{p}{;}

\PY{c+c1}{// No boxing at all \PYZhy{} stays as int the whole way}
\PY{k+kt}{int}\PY{+w}{ }\PY{n}{fast}\PY{+w}{ }\PY{o}{=}\PY{+w}{ }\PY{n}{numbers}\PY{p}{.}\PY{n+na}{stream}\PY{p}{(}\PY{p}{)}\PY{p}{.}\PY{n+na}{mapToInt}\PY{p}{(}\PY{n}{Integer}\PY{p}{:}\PY{p}{:}\PY{n}{intValue}\PY{p}{)}\PY{p}{.}\PY{n+na}{sum}\PY{p}{(}\PY{p}{)}\PY{p}{;}
\end{Verbatim}
\end{codebox}

\subsection[Splittability of sources]{Splittability of sources}
\label{h:8:splittability-of-sources}
Parallel streams rely on \code{trySplit()} dividing work into roughly equal, cheaply-computed chunks (see \hyperref[h:8:spliterator]{Spliterator}). Not all sources split well:

\begin{itemize}
\item 
\textbf{Good}: \code{ArrayList}, arrays, \code{IntStream.\allowbreak{}range} — random access means splitting a range in half is O(1) and produces balanced chunks.

\item 
\textbf{Bad}: \code{LinkedList}, \code{Iterator}-based sources, \code{Stream.\allowbreak{}iterate} — no random access, so splitting requires walking elements one at a time, which defeats the purpose, or can’t be split at all.

\end{itemize}

\begin{codebox}
\begin{Verbatim}[commandchars=\\\{\}]
\PY{n}{List}\PY{o}{\PYZlt{}}\PY{n}{Integer}\PY{o}{\PYZgt{}}\PY{+w}{ }\PY{n}{arrayList}\PY{+w}{ }\PY{o}{=}\PY{+w}{ }\PY{k}{new}\PY{+w}{ }\PY{n}{ArrayList}\PY{o}{\PYZlt{}}\PY{o}{\PYZgt{}}\PY{p}{(}\PY{n}{List}\PY{p}{.}\PY{n+na}{of}\PY{p}{(}\PY{l+m+mi}{1}\PY{p}{,}\PY{+w}{ }\PY{l+m+mi}{2}\PY{p}{,}\PY{+w}{ }\PY{l+m+mi}{3}\PY{p}{,}\PY{+w}{ }\PY{l+m+mi}{4}\PY{p}{,}\PY{+w}{ }\PY{l+m+mi}{5}\PY{p}{,}\PY{+w}{ }\PY{l+m+mi}{6}\PY{p}{,}\PY{+w}{ }\PY{l+m+mi}{7}\PY{p}{,}\PY{+w}{ }\PY{l+m+mi}{8}\PY{p}{)}\PY{p}{)}\PY{p}{;}
\PY{n}{List}\PY{o}{\PYZlt{}}\PY{n}{Integer}\PY{o}{\PYZgt{}}\PY{+w}{ }\PY{n}{linkedList}\PY{+w}{ }\PY{o}{=}\PY{+w}{ }\PY{k}{new}\PY{+w}{ }\PY{n}{LinkedList}\PY{o}{\PYZlt{}}\PY{o}{\PYZgt{}}\PY{p}{(}\PY{n}{arrayList}\PY{p}{)}\PY{p}{;}

\PY{n}{arrayList}\PY{p}{.}\PY{n+na}{parallelStream}\PY{p}{(}\PY{p}{)}\PY{p}{.}\PY{n+na}{forEach}\PY{p}{(}\PY{n}{n}\PY{+w}{ }\PY{o}{\PYZhy{}}\PY{o}{\PYZgt{}}\PY{+w}{ }\PY{p}{\PYZob{}}\PY{p}{\PYZcb{}}\PY{p}{)}\PY{p}{;}\PY{+w}{  }\PY{c+c1}{// splits into balanced chunks cheaply \PYZhy{} scales well}
\PY{n}{linkedList}\PY{p}{.}\PY{n+na}{parallelStream}\PY{p}{(}\PY{p}{)}\PY{p}{.}\PY{n+na}{forEach}\PY{p}{(}\PY{n}{n}\PY{+w}{ }\PY{o}{\PYZhy{}}\PY{o}{\PYZgt{}}\PY{+w}{ }\PY{p}{\PYZob{}}\PY{p}{\PYZcb{}}\PY{p}{)}\PY{p}{;}\PY{+w}{ }\PY{c+c1}{// trySplit degrades to linear traversal \PYZhy{} little/no benefit}
\end{Verbatim}
\end{codebox}

\subsection[The common ForkJoinPool and why blocking is dangerous]{The common \code{ForkJoinPool} and why blocking is dangerous}
\label{h:8:the-common-forkjoinpool-and-why-blocking-is-dangerous}
\code{parallelStream()} doesn’t spin up its own threads — by default it submits work to the JVM-wide \textbf{common \code{ForkJoinPool}}, sized to \code{Runtime.\allowbreak{}get\allowbreak{}Runtime().\allowbreak{}available\allowbreak{}Processors()\allowbreak{}~\allowbreak{}-\allowbreak{}~\allowbreak{}1} worker threads. That pool is shared by \emph{every} parallel stream in the whole application (and by anything else using \code{Fork\allowbreak{}Join\allowbreak{}Pool.\allowbreak{}common\allowbreak{}Pool()} directly, including \code{CompletableFuture}’s async methods).

If a task inside a parallel stream \textbf{blocks} — a network call, a JDBC query, a \code{Thread.\allowbreak{}sleep}, waiting on a lock — it occupies one of those few shared worker threads and can starve unrelated parallel work elsewhere in the app. Since the pool is small and shared, one badly-behaved parallel stream can quietly degrade performance for the whole application, not just itself.

\begin{codebox}
\begin{Verbatim}[commandchars=\\\{\}]
\PY{c+c1}{// Dangerous: blocking I/O inside a parallel stream ties up common pool threads}
\PY{n}{List}\PY{o}{\PYZlt{}}\PY{n}{String}\PY{o}{\PYZgt{}}\PY{+w}{ }\PY{n}{urls}\PY{+w}{ }\PY{o}{=}\PY{+w}{ }\PY{n}{List}\PY{p}{.}\PY{n+na}{of}\PY{p}{(}\PY{l+s}{\PYZdq{}}\PY{l+s}{https://a}\PY{l+s}{\PYZdq{}}\PY{p}{,}\PY{+w}{ }\PY{l+s}{\PYZdq{}}\PY{l+s}{https://b}\PY{l+s}{\PYZdq{}}\PY{p}{,}\PY{+w}{ }\PY{l+s}{\PYZdq{}}\PY{l+s}{https://c}\PY{l+s}{\PYZdq{}}\PY{p}{)}\PY{p}{;}
\PY{n}{List}\PY{o}{\PYZlt{}}\PY{n}{String}\PY{o}{\PYZgt{}}\PY{+w}{ }\PY{n}{bodies}\PY{+w}{ }\PY{o}{=}\PY{+w}{ }\PY{n}{urls}\PY{p}{.}\PY{n+na}{parallelStream}\PY{p}{(}\PY{p}{)}
\PY{+w}{    }\PY{p}{.}\PY{n+na}{map}\PY{p}{(}\PY{n}{url}\PY{+w}{ }\PY{o}{\PYZhy{}}\PY{o}{\PYZgt{}}\PY{+w}{ }\PY{n}{httpClient}\PY{p}{.}\PY{n+na}{get}\PY{p}{(}\PY{n}{url}\PY{p}{)}\PY{p}{)}\PY{+w}{ }\PY{c+c1}{// blocks a shared ForkJoinPool worker thread per call}
\PY{+w}{    }\PY{p}{.}\PY{n+na}{toList}\PY{p}{(}\PY{p}{)}\PY{p}{;}

\PY{c+c1}{// Safer: use a dedicated executor for blocking work, keep the common pool free}
\PY{n}{ExecutorService}\PY{+w}{ }\PY{n}{ioPool}\PY{+w}{ }\PY{o}{=}\PY{+w}{ }\PY{n}{Executors}\PY{p}{.}\PY{n+na}{newFixedThreadPool}\PY{p}{(}\PY{l+m+mi}{10}\PY{p}{)}\PY{p}{;}
\PY{n}{List}\PY{o}{\PYZlt{}}\PY{n}{CompletableFuture}\PY{o}{\PYZlt{}}\PY{n}{String}\PY{o}{\PYZgt{}}\PY{o}{\PYZgt{}}\PY{+w}{ }\PY{n}{futures}\PY{+w}{ }\PY{o}{=}\PY{+w}{ }\PY{n}{urls}\PY{p}{.}\PY{n+na}{stream}\PY{p}{(}\PY{p}{)}
\PY{+w}{    }\PY{p}{.}\PY{n+na}{map}\PY{p}{(}\PY{n}{url}\PY{+w}{ }\PY{o}{\PYZhy{}}\PY{o}{\PYZgt{}}\PY{+w}{ }\PY{n}{CompletableFuture}\PY{p}{.}\PY{n+na}{supplyAsync}\PY{p}{(}\PY{p}{(}\PY{p}{)}\PY{+w}{ }\PY{o}{\PYZhy{}}\PY{o}{\PYZgt{}}\PY{+w}{ }\PY{n}{httpClient}\PY{p}{.}\PY{n+na}{get}\PY{p}{(}\PY{n}{url}\PY{p}{)}\PY{p}{,}\PY{+w}{ }\PY{n}{ioPool}\PY{p}{)}\PY{p}{)}
\PY{+w}{    }\PY{p}{.}\PY{n+na}{toList}\PY{p}{(}\PY{p}{)}\PY{p}{;}
\PY{n}{List}\PY{o}{\PYZlt{}}\PY{n}{String}\PY{o}{\PYZgt{}}\PY{+w}{ }\PY{n}{results}\PY{+w}{ }\PY{o}{=}\PY{+w}{ }\PY{n}{futures}\PY{p}{.}\PY{n+na}{stream}\PY{p}{(}\PY{p}{)}\PY{p}{.}\PY{n+na}{map}\PY{p}{(}\PY{n}{CompletableFuture}\PY{p}{:}\PY{p}{:}\PY{n}{join}\PY{p}{)}\PY{p}{.}\PY{n+na}{toList}\PY{p}{(}\PY{p}{)}\PY{p}{;}
\end{Verbatim}
\end{codebox}

\subsection[Stateful lambdas and shared mutable state]{Stateful lambdas and shared mutable state}
\label{h:8:stateful-lambdas-and-shared-mutable-state}
Passing a lambda that mutates shared state into a stream operation — especially a parallel one — is a race condition waiting to happen. Streams give no guarantee about thread-safety of your own captured state; that’s on you.

\begin{codebox}
\begin{Verbatim}[commandchars=\\\{\}]
\PY{c+c1}{// BROKEN: ArrayList is not thread\PYZhy{}safe, and multiple threads add() concurrently}
\PY{n}{List}\PY{o}{\PYZlt{}}\PY{n}{Integer}\PY{o}{\PYZgt{}}\PY{+w}{ }\PY{n}{results}\PY{+w}{ }\PY{o}{=}\PY{+w}{ }\PY{k}{new}\PY{+w}{ }\PY{n}{ArrayList}\PY{o}{\PYZlt{}}\PY{o}{\PYZgt{}}\PY{p}{(}\PY{p}{)}\PY{p}{;}
\PY{n}{IntStream}\PY{p}{.}\PY{n+na}{range}\PY{p}{(}\PY{l+m+mi}{0}\PY{p}{,}\PY{+w}{ }\PY{l+m+mi}{10\PYZus{}000}\PY{p}{)}\PY{p}{.}\PY{n+na}{parallel}\PY{p}{(}\PY{p}{)}\PY{p}{.}\PY{n+na}{forEach}\PY{p}{(}\PY{n}{results}\PY{p}{:}\PY{p}{:}\PY{n}{add}\PY{p}{)}\PY{p}{;}\PY{+w}{ }\PY{c+c1}{// corrupts internal state, may throw or lose elements}

\PY{c+c1}{// FIXED: let the stream own the accumulation via a proper collector}
\PY{n}{List}\PY{o}{\PYZlt{}}\PY{n}{Integer}\PY{o}{\PYZgt{}}\PY{+w}{ }\PY{n}{results2}\PY{+w}{ }\PY{o}{=}\PY{+w}{ }\PY{n}{IntStream}\PY{p}{.}\PY{n+na}{range}\PY{p}{(}\PY{l+m+mi}{0}\PY{p}{,}\PY{+w}{ }\PY{l+m+mi}{10\PYZus{}000}\PY{p}{)}\PY{p}{.}\PY{n+na}{parallel}\PY{p}{(}\PY{p}{)}\PY{p}{.}\PY{n+na}{boxed}\PY{p}{(}\PY{p}{)}\PY{p}{.}\PY{n+na}{collect}\PY{p}{(}\PY{n}{Collectors}\PY{p}{.}\PY{n+na}{toList}\PY{p}{(}\PY{p}{)}\PY{p}{)}\PY{p}{;}
\end{Verbatim}
\end{codebox}

The rule: lambdas passed to stream operations should be \textbf{stateless} with respect to shared mutable data. Let \code{collect}/\code{reduce} do the accumulating — those are designed to combine partial results safely across threads.

\subsection[forEach vs forEachOrdered]{\code{forEach} vs \code{forEachOrdered}}
\label{h:8:foreach-vs-foreachordered}
On a parallel stream, \code{forEach} makes \textbf{no promise about order} — elements are processed as threads happen to finish, which is faster but unpredictable. \code{forEachOrdered} forces processing back into encounter order, which reintroduces most of the coordination overhead you were trying to avoid by going parallel.

\begin{codebox}
\begin{Verbatim}[commandchars=\\\{\}]
\PY{n}{List}\PY{o}{\PYZlt{}}\PY{n}{Integer}\PY{o}{\PYZgt{}}\PY{+w}{ }\PY{n}{nums}\PY{+w}{ }\PY{o}{=}\PY{+w}{ }\PY{n}{List}\PY{p}{.}\PY{n+na}{of}\PY{p}{(}\PY{l+m+mi}{1}\PY{p}{,}\PY{+w}{ }\PY{l+m+mi}{2}\PY{p}{,}\PY{+w}{ }\PY{l+m+mi}{3}\PY{p}{,}\PY{+w}{ }\PY{l+m+mi}{4}\PY{p}{,}\PY{+w}{ }\PY{l+m+mi}{5}\PY{p}{)}\PY{p}{;}

\PY{n}{nums}\PY{p}{.}\PY{n+na}{parallelStream}\PY{p}{(}\PY{p}{)}\PY{p}{.}\PY{n+na}{forEach}\PY{p}{(}\PY{n}{System}\PY{p}{.}\PY{n+na}{out}\PY{p}{:}\PY{p}{:}\PY{n}{println}\PY{p}{)}\PY{p}{;}\PY{+w}{        }\PY{c+c1}{// order NOT guaranteed, e.g. 3 1 4 2 5}
\PY{n}{nums}\PY{p}{.}\PY{n+na}{parallelStream}\PY{p}{(}\PY{p}{)}\PY{p}{.}\PY{n+na}{forEachOrdered}\PY{p}{(}\PY{n}{System}\PY{p}{.}\PY{n+na}{out}\PY{p}{:}\PY{p}{:}\PY{n}{println}\PY{p}{)}\PY{p}{;}\PY{+w}{ }\PY{c+c1}{// always 1 2 3 4 5, but loses most parallel speedup}
\end{Verbatim}
\end{codebox}

If you need parallel speed and ordered output, prefer collecting to a \code{List} (which preserves encounter order) over forcing \code{forEachOrdered}.

\section[Parallel Streams]{Parallel Streams}
\label{h:8:parallel-streams}
A \textbf{parallel stream} splits the source, processes chunks on multiple threads from the common \code{ForkJoinPool}, then merges the partial results. \code{stream()} and \code{parallelStream()} build the \emph{same} pipeline API — the parallel-ness is just a flag, toggled with \code{.\allowbreak{}parallel()} / \code{.\allowbreak{}sequential()} at any point in the chain (the \emph{last} call before the terminal operation wins).

\begin{codebox}
\begin{Verbatim}[commandchars=\\\{\}]
\PY{n}{List}\PY{o}{\PYZlt{}}\PY{n}{Integer}\PY{o}{\PYZgt{}}\PY{+w}{ }\PY{n}{nums}\PY{+w}{ }\PY{o}{=}\PY{+w}{ }\PY{n}{IntStream}\PY{p}{.}\PY{n+na}{rangeClosed}\PY{p}{(}\PY{l+m+mi}{1}\PY{p}{,}\PY{+w}{ }\PY{l+m+mi}{20}\PY{p}{)}\PY{p}{.}\PY{n+na}{boxed}\PY{p}{(}\PY{p}{)}\PY{p}{.}\PY{n+na}{toList}\PY{p}{(}\PY{p}{)}\PY{p}{;}

\PY{k+kt}{int}\PY{+w}{ }\PY{n}{sumSequential}\PY{+w}{ }\PY{o}{=}\PY{+w}{ }\PY{n}{nums}\PY{p}{.}\PY{n+na}{stream}\PY{p}{(}\PY{p}{)}\PY{p}{.}\PY{n+na}{mapToInt}\PY{p}{(}\PY{n}{Integer}\PY{p}{:}\PY{p}{:}\PY{n}{intValue}\PY{p}{)}\PY{p}{.}\PY{n+na}{sum}\PY{p}{(}\PY{p}{)}\PY{p}{;}
\PY{k+kt}{int}\PY{+w}{ }\PY{n}{sumParallel}\PY{+w}{ }\PY{o}{=}\PY{+w}{ }\PY{n}{nums}\PY{p}{.}\PY{n+na}{parallelStream}\PY{p}{(}\PY{p}{)}\PY{p}{.}\PY{n+na}{mapToInt}\PY{p}{(}\PY{n}{Integer}\PY{p}{:}\PY{p}{:}\PY{n}{intValue}\PY{p}{)}\PY{p}{.}\PY{n+na}{sum}\PY{p}{(}\PY{p}{)}\PY{p}{;}
\PY{n}{System}\PY{p}{.}\PY{n+na}{out}\PY{p}{.}\PY{n+na}{println}\PY{p}{(}\PY{n}{sumSequential}\PY{+w}{ }\PY{o}{=}\PY{o}{=}\PY{+w}{ }\PY{n}{sumParallel}\PY{p}{)}\PY{p}{;}\PY{+w}{ }\PY{c+c1}{// true \PYZhy{} same result, different execution}
\end{Verbatim}
\end{codebox}

\subsection[When parallel actually pays off — the “N × Q” rule of thumb]{When parallel actually pays off — the “N × Q” rule of thumb}
\label{h:8:when-parallel-actually-pays-off--the-n--q-rule-of-thumb}
Parallelism has fixed costs: splitting the source, scheduling tasks, merging results. It only pays off when the \emph{total work} is large enough to amortize those costs. A widely used heuristic (from Brian Goetz and the JDK’s own stream design notes) is:

\begin{notebox}
\textbf{N × Q \textgreater{} 10,000} (roughly), where \textbf{N} = number of elements and \textbf{Q} = cost of processing one element.

\end{notebox}

\begin{itemize}
\item 
Large \textbf{N}, cheap \textbf{Q} (e.g., summing a million \code{int}s): parallel can help, but boxing and splitting overhead may eat the gain — measure.

\item 
Small \textbf{N}, expensive \textbf{Q} (e.g., 50 items, each needing a slow computation): parallel can help a lot, since 50 splits handled by several threads still swamps the coordination cost.

\item 
Small \textbf{N}, cheap \textbf{Q} (e.g., summing 10 numbers): parallel almost always loses — pure overhead, no upside.

\item 
Large \textbf{N}, expensive \textbf{Q}: parallel usually wins clearly.

\end{itemize}

\begin{codebox}
\begin{Verbatim}[commandchars=\\\{\}]
\PY{c+c1}{// Small N, cheap Q: parallel overhead dominates \PYZhy{} don\PYZsq{}t bother}
\PY{k+kt}{int}\PY{+w}{ }\PY{n}{tinySum}\PY{+w}{ }\PY{o}{=}\PY{+w}{ }\PY{n}{IntStream}\PY{p}{.}\PY{n+na}{rangeClosed}\PY{p}{(}\PY{l+m+mi}{1}\PY{p}{,}\PY{+w}{ }\PY{l+m+mi}{10}\PY{p}{)}\PY{p}{.}\PY{n+na}{parallel}\PY{p}{(}\PY{p}{)}\PY{p}{.}\PY{n+na}{sum}\PY{p}{(}\PY{p}{)}\PY{p}{;}\PY{+w}{ }\PY{c+c1}{// slower than sequential in practice}

\PY{c+c1}{// Large N, expensive Q per element: parallel is likely to help}
\PY{n}{List}\PY{o}{\PYZlt{}}\PY{n}{Long}\PY{o}{\PYZgt{}}\PY{+w}{ }\PY{n}{primesNear}\PY{+w}{ }\PY{o}{=}\PY{+w}{ }\PY{n}{LongStream}\PY{p}{.}\PY{n+na}{rangeClosed}\PY{p}{(}\PY{l+m+mi}{2}\PY{p}{,}\PY{+w}{ }\PY{l+m+mi}{2\PYZus{}000\PYZus{}000}\PY{p}{)}
\PY{+w}{    }\PY{p}{.}\PY{n+na}{parallel}\PY{p}{(}\PY{p}{)}
\PY{+w}{    }\PY{p}{.}\PY{n+na}{filter}\PY{p}{(}\PY{n}{BigIntegerUtil}\PY{p}{:}\PY{p}{:}\PY{n}{isPrimeSlowCheck}\PY{p}{)}\PY{+w}{ }\PY{c+c1}{// pretend this is CPU\PYZhy{}heavy per element}
\PY{+w}{    }\PY{p}{.}\PY{n+na}{boxed}\PY{p}{(}\PY{p}{)}
\PY{+w}{    }\PY{p}{.}\PY{n+na}{toList}\PY{p}{(}\PY{p}{)}\PY{p}{;}
\end{Verbatim}
\end{codebox}

Always benchmark with something like JMH before trusting intuition — modern JITs and cache effects can make the “obvious” answer wrong, especially at small N. Never guess in a code review either way; ask for a benchmark or a clear justification when a PR reaches for \code{.\allowbreak{}parallel()}.

\section[Common Code-Review Interview Pitfalls]{Common Code-Review Interview Pitfalls}
\label{h:8:common-code-review-interview-pitfalls}
\begin{enumerate}
\item 
\textbf{Missing \code{@\allowbreak{}FunctionalInterface}, so a second abstract method sneaks in later.}
Why it matters: without the annotation, the compiler won’t stop a teammate from breaking every lambda that implements the interface.

\begin{codebox}
\begin{Verbatim}[commandchars=\\\{\}]
\PY{c+c1}{// Before}
\PY{k+kd}{interface} \PY{n+nc}{Handler}\PY{o}{\PYZlt{}}\PY{n}{T}\PY{o}{\PYZgt{}}\PY{+w}{ }\PY{p}{\PYZob{}}\PY{+w}{ }\PY{k+kt}{void}\PY{+w}{ }\PY{n+nf}{handle}\PY{p}{(}\PY{n}{T}\PY{+w}{ }\PY{n}{t}\PY{p}{)}\PY{p}{;}\PY{+w}{ }\PY{p}{\PYZcb{}}

\PY{c+c1}{// After}
\PY{n+nd}{@FunctionalInterface}
\PY{k+kd}{interface} \PY{n+nc}{Handler}\PY{o}{\PYZlt{}}\PY{n}{T}\PY{o}{\PYZgt{}}\PY{+w}{ }\PY{p}{\PYZob{}}\PY{+w}{ }\PY{k+kt}{void}\PY{+w}{ }\PY{n+nf}{handle}\PY{p}{(}\PY{n}{T}\PY{+w}{ }\PY{n}{t}\PY{p}{)}\PY{p}{;}\PY{+w}{ }\PY{p}{\PYZcb{}}
\end{Verbatim}
\end{codebox}

\item 
\textbf{Capturing a mutable loop variable or field expecting per-iteration semantics.}
Why it matters: lambdas capture variables, not values in the “live” sense for fields, so shared mutable state read later can be wrong or racy.

\begin{codebox}
\begin{Verbatim}[commandchars=\\\{\}]
\PY{c+c1}{// Before \PYZhy{} all callbacks share and read the mutable field at call time, likely all seeing the final value}
\PY{k}{for}\PY{+w}{ }\PY{p}{(}\PY{k+kt}{int}\PY{+w}{ }\PY{n}{i}\PY{+w}{ }\PY{o}{=}\PY{+w}{ }\PY{l+m+mi}{0}\PY{p}{;}\PY{+w}{ }\PY{n}{i}\PY{+w}{ }\PY{o}{\PYZlt{}}\PY{+w}{ }\PY{n}{tasks}\PY{p}{.}\PY{n+na}{size}\PY{p}{(}\PY{p}{)}\PY{p}{;}\PY{+w}{ }\PY{n}{i}\PY{o}{+}\PY{o}{+}\PY{p}{)}\PY{+w}{ }\PY{p}{\PYZob{}}
\PY{+w}{    }\PY{k}{this}\PY{p}{.}\PY{n+na}{currentIndex}\PY{+w}{ }\PY{o}{=}\PY{+w}{ }\PY{n}{i}\PY{p}{;}
\PY{+w}{    }\PY{n}{tasks}\PY{p}{.}\PY{n+na}{get}\PY{p}{(}\PY{n}{i}\PY{p}{)}\PY{p}{.}\PY{n+na}{setCallback}\PY{p}{(}\PY{p}{(}\PY{p}{)}\PY{+w}{ }\PY{o}{\PYZhy{}}\PY{o}{\PYZgt{}}\PY{+w}{ }\PY{n}{process}\PY{p}{(}\PY{k}{this}\PY{p}{.}\PY{n+na}{currentIndex}\PY{p}{)}\PY{p}{)}\PY{p}{;}
\PY{p}{\PYZcb{}}

\PY{c+c1}{// After \PYZhy{} capture an effectively\PYZhy{}final local per iteration}
\PY{k}{for}\PY{+w}{ }\PY{p}{(}\PY{k+kt}{int}\PY{+w}{ }\PY{n}{i}\PY{+w}{ }\PY{o}{=}\PY{+w}{ }\PY{l+m+mi}{0}\PY{p}{;}\PY{+w}{ }\PY{n}{i}\PY{+w}{ }\PY{o}{\PYZlt{}}\PY{+w}{ }\PY{n}{tasks}\PY{p}{.}\PY{n+na}{size}\PY{p}{(}\PY{p}{)}\PY{p}{;}\PY{+w}{ }\PY{n}{i}\PY{o}{+}\PY{o}{+}\PY{p}{)}\PY{+w}{ }\PY{p}{\PYZob{}}
\PY{+w}{    }\PY{k+kt}{int}\PY{+w}{ }\PY{n}{index}\PY{+w}{ }\PY{o}{=}\PY{+w}{ }\PY{n}{i}\PY{p}{;}\PY{+w}{ }\PY{c+c1}{// effectively final, one per iteration}
\PY{+w}{    }\PY{n}{tasks}\PY{p}{.}\PY{n+na}{get}\PY{p}{(}\PY{n}{i}\PY{p}{)}\PY{p}{.}\PY{n+na}{setCallback}\PY{p}{(}\PY{p}{(}\PY{p}{)}\PY{+w}{ }\PY{o}{\PYZhy{}}\PY{o}{\PYZgt{}}\PY{+w}{ }\PY{n}{process}\PY{p}{(}\PY{n}{index}\PY{p}{)}\PY{p}{)}\PY{p}{;}
\PY{p}{\PYZcb{}}
\end{Verbatim}
\end{codebox}

\item 
\textbf{Using \code{peek} to drive actual program logic instead of debugging.}
Why it matters: \code{peek} may be skipped by the JDK if the result doesn’t depend on it, and short-circuiting can stop it from ever running for all elements.

\begin{codebox}
\begin{Verbatim}[commandchars=\\\{\}]
\PY{c+c1}{// Before \PYZhy{} relies on peek to populate a list as a side effect}
\PY{n}{List}\PY{o}{\PYZlt{}}\PY{n}{String}\PY{o}{\PYZgt{}}\PY{+w}{ }\PY{n}{seen}\PY{+w}{ }\PY{o}{=}\PY{+w}{ }\PY{k}{new}\PY{+w}{ }\PY{n}{ArrayList}\PY{o}{\PYZlt{}}\PY{o}{\PYZgt{}}\PY{p}{(}\PY{p}{)}\PY{p}{;}
\PY{k+kt}{long}\PY{+w}{ }\PY{n}{count}\PY{+w}{ }\PY{o}{=}\PY{+w}{ }\PY{n}{names}\PY{p}{.}\PY{n+na}{stream}\PY{p}{(}\PY{p}{)}\PY{p}{.}\PY{n+na}{peek}\PY{p}{(}\PY{n}{seen}\PY{p}{:}\PY{p}{:}\PY{n}{add}\PY{p}{)}\PY{p}{.}\PY{n+na}{count}\PY{p}{(}\PY{p}{)}\PY{p}{;}\PY{+w}{ }\PY{c+c1}{// may not run peek at all!}

\PY{c+c1}{// After \PYZhy{} do the real work in map/forEach, not peek}
\PY{n}{List}\PY{o}{\PYZlt{}}\PY{n}{String}\PY{o}{\PYZgt{}}\PY{+w}{ }\PY{n}{seen2}\PY{+w}{ }\PY{o}{=}\PY{+w}{ }\PY{n}{names}\PY{p}{.}\PY{n+na}{stream}\PY{p}{(}\PY{p}{)}\PY{p}{.}\PY{n+na}{collect}\PY{p}{(}\PY{n}{Collectors}\PY{p}{.}\PY{n+na}{toList}\PY{p}{(}\PY{p}{)}\PY{p}{)}\PY{p}{;}
\end{Verbatim}
\end{codebox}

\item 
\textbf{Mutating shared state from inside a parallel stream lambda.}
Why it matters: non-thread-safe collections corrupt under concurrent \code{add}, causing silent data loss or exceptions.

\begin{codebox}
\begin{Verbatim}[commandchars=\\\{\}]
\PY{c+c1}{// Before}
\PY{n}{List}\PY{o}{\PYZlt{}}\PY{n}{Integer}\PY{o}{\PYZgt{}}\PY{+w}{ }\PY{n}{out}\PY{+w}{ }\PY{o}{=}\PY{+w}{ }\PY{k}{new}\PY{+w}{ }\PY{n}{ArrayList}\PY{o}{\PYZlt{}}\PY{o}{\PYZgt{}}\PY{p}{(}\PY{p}{)}\PY{p}{;}
\PY{n}{data}\PY{p}{.}\PY{n+na}{parallelStream}\PY{p}{(}\PY{p}{)}\PY{p}{.}\PY{n+na}{forEach}\PY{p}{(}\PY{n}{out}\PY{p}{:}\PY{p}{:}\PY{n}{add}\PY{p}{)}\PY{p}{;}\PY{+w}{ }\PY{c+c1}{// race condition}

\PY{c+c1}{// After}
\PY{n}{List}\PY{o}{\PYZlt{}}\PY{n}{Integer}\PY{o}{\PYZgt{}}\PY{+w}{ }\PY{n}{out2}\PY{+w}{ }\PY{o}{=}\PY{+w}{ }\PY{n}{data}\PY{p}{.}\PY{n+na}{parallelStream}\PY{p}{(}\PY{p}{)}\PY{p}{.}\PY{n+na}{collect}\PY{p}{(}\PY{n}{Collectors}\PY{p}{.}\PY{n+na}{toList}\PY{p}{(}\PY{p}{)}\PY{p}{)}\PY{p}{;}
\end{Verbatim}
\end{codebox}

\item 
\textbf{Reusing a stream after a terminal operation.}
Why it matters: streams are single-use; a second terminal call throws \code{IllegalStateException} at runtime, not compile time.

\begin{codebox}
\begin{Verbatim}[commandchars=\\\{\}]
\PY{c+c1}{// Before}
\PY{n}{Stream}\PY{o}{\PYZlt{}}\PY{n}{String}\PY{o}{\PYZgt{}}\PY{+w}{ }\PY{n}{s}\PY{+w}{ }\PY{o}{=}\PY{+w}{ }\PY{n}{items}\PY{p}{.}\PY{n+na}{stream}\PY{p}{(}\PY{p}{)}\PY{p}{;}
\PY{n}{s}\PY{p}{.}\PY{n+na}{forEach}\PY{p}{(}\PY{n}{System}\PY{p}{.}\PY{n+na}{out}\PY{p}{:}\PY{p}{:}\PY{n}{println}\PY{p}{)}\PY{p}{;}
\PY{k+kt}{long}\PY{+w}{ }\PY{n}{n}\PY{+w}{ }\PY{o}{=}\PY{+w}{ }\PY{n}{s}\PY{p}{.}\PY{n+na}{count}\PY{p}{(}\PY{p}{)}\PY{p}{;}\PY{+w}{ }\PY{c+c1}{// IllegalStateException}

\PY{c+c1}{// After}
\PY{k+kt}{long}\PY{+w}{ }\PY{n}{n2}\PY{+w}{ }\PY{o}{=}\PY{+w}{ }\PY{n}{items}\PY{p}{.}\PY{n+na}{stream}\PY{p}{(}\PY{p}{)}\PY{p}{.}\PY{n+na}{count}\PY{p}{(}\PY{p}{)}\PY{p}{;}\PY{+w}{ }\PY{c+c1}{// build a fresh stream from the source}
\end{Verbatim}
\end{codebox}

\item 
\textbf{Calling \code{.\allowbreak{}parallel()} on a small or cheap workload “for speed.”}
Why it matters: fork/join scheduling and merge overhead usually outweighs the gain unless N × Q is large; it can actually be slower and it steals threads from the shared pool.

\begin{codebox}
\begin{Verbatim}[commandchars=\\\{\}]
\PY{c+c1}{// Before}
\PY{k+kt}{int}\PY{+w}{ }\PY{n}{sum}\PY{+w}{ }\PY{o}{=}\PY{+w}{ }\PY{n}{IntStream}\PY{p}{.}\PY{n+na}{rangeClosed}\PY{p}{(}\PY{l+m+mi}{1}\PY{p}{,}\PY{+w}{ }\PY{l+m+mi}{20}\PY{p}{)}\PY{p}{.}\PY{n+na}{parallel}\PY{p}{(}\PY{p}{)}\PY{p}{.}\PY{n+na}{sum}\PY{p}{(}\PY{p}{)}\PY{p}{;}\PY{+w}{ }\PY{c+c1}{// pure overhead}

\PY{c+c1}{// After}
\PY{k+kt}{int}\PY{+w}{ }\PY{n}{sum2}\PY{+w}{ }\PY{o}{=}\PY{+w}{ }\PY{n}{IntStream}\PY{p}{.}\PY{n+na}{rangeClosed}\PY{p}{(}\PY{l+m+mi}{1}\PY{p}{,}\PY{+w}{ }\PY{l+m+mi}{20}\PY{p}{)}\PY{p}{.}\PY{n+na}{sum}\PY{p}{(}\PY{p}{)}\PY{p}{;}
\end{Verbatim}
\end{codebox}

\item 
\textbf{Blocking I/O inside a parallel stream.}
Why it matters: it occupies threads in the shared common \code{ForkJoinPool}, potentially starving unrelated parallel work elsewhere in the JVM.

\begin{codebox}
\begin{Verbatim}[commandchars=\\\{\}]
\PY{c+c1}{// Before}
\PY{n}{List}\PY{o}{\PYZlt{}}\PY{n}{String}\PY{o}{\PYZgt{}}\PY{+w}{ }\PY{n}{bodies}\PY{+w}{ }\PY{o}{=}\PY{+w}{ }\PY{n}{urls}\PY{p}{.}\PY{n+na}{parallelStream}\PY{p}{(}\PY{p}{)}\PY{p}{.}\PY{n+na}{map}\PY{p}{(}\PY{n}{httpClient}\PY{p}{:}\PY{p}{:}\PY{n}{get}\PY{p}{)}\PY{p}{.}\PY{n+na}{toList}\PY{p}{(}\PY{p}{)}\PY{p}{;}

\PY{c+c1}{// After}
\PY{n}{ExecutorService}\PY{+w}{ }\PY{n}{ioPool}\PY{+w}{ }\PY{o}{=}\PY{+w}{ }\PY{n}{Executors}\PY{p}{.}\PY{n+na}{newFixedThreadPool}\PY{p}{(}\PY{l+m+mi}{10}\PY{p}{)}\PY{p}{;}
\PY{n}{List}\PY{o}{\PYZlt{}}\PY{n}{String}\PY{o}{\PYZgt{}}\PY{+w}{ }\PY{n}{bodies2}\PY{+w}{ }\PY{o}{=}\PY{+w}{ }\PY{n}{urls}\PY{p}{.}\PY{n+na}{stream}\PY{p}{(}\PY{p}{)}
\PY{+w}{    }\PY{p}{.}\PY{n+na}{map}\PY{p}{(}\PY{n}{u}\PY{+w}{ }\PY{o}{\PYZhy{}}\PY{o}{\PYZgt{}}\PY{+w}{ }\PY{n}{CompletableFuture}\PY{p}{.}\PY{n+na}{supplyAsync}\PY{p}{(}\PY{p}{(}\PY{p}{)}\PY{+w}{ }\PY{o}{\PYZhy{}}\PY{o}{\PYZgt{}}\PY{+w}{ }\PY{n}{httpClient}\PY{p}{.}\PY{n+na}{get}\PY{p}{(}\PY{n}{u}\PY{p}{)}\PY{p}{,}\PY{+w}{ }\PY{n}{ioPool}\PY{p}{)}\PY{p}{)}
\PY{+w}{    }\PY{p}{.}\PY{n+na}{toList}\PY{p}{(}\PY{p}{)}\PY{p}{.}\PY{n+na}{stream}\PY{p}{(}\PY{p}{)}\PY{p}{.}\PY{n+na}{map}\PY{p}{(}\PY{n}{CompletableFuture}\PY{p}{:}\PY{p}{:}\PY{n}{join}\PY{p}{)}\PY{p}{.}\PY{n+na}{toList}\PY{p}{(}\PY{p}{)}\PY{p}{;}
\end{Verbatim}
\end{codebox}

\item 
\textbf{Assuming \code{Collectors.\allowbreak{}toMap} silently overwrites on duplicate keys.}
Why it matters: it throws \code{IllegalStateException} at runtime instead — a merge function is required if duplicates are possible.

\begin{codebox}
\begin{Verbatim}[commandchars=\\\{\}]
\PY{c+c1}{// Before}
\PY{n}{Map}\PY{o}{\PYZlt{}}\PY{n}{Integer}\PY{p}{,}\PY{+w}{ }\PY{n}{String}\PY{o}{\PYZgt{}}\PY{+w}{ }\PY{n}{byLen}\PY{+w}{ }\PY{o}{=}\PY{+w}{ }\PY{n}{words}\PY{p}{.}\PY{n+na}{stream}\PY{p}{(}\PY{p}{)}\PY{p}{.}\PY{n+na}{collect}\PY{p}{(}\PY{n}{Collectors}\PY{p}{.}\PY{n+na}{toMap}\PY{p}{(}\PY{n}{String}\PY{p}{:}\PY{p}{:}\PY{n}{length}\PY{p}{,}\PY{+w}{ }\PY{n}{w}\PY{+w}{ }\PY{o}{\PYZhy{}}\PY{o}{\PYZgt{}}\PY{+w}{ }\PY{n}{w}\PY{p}{)}\PY{p}{)}\PY{p}{;}
\PY{c+c1}{// throws on duplicate length}

\PY{c+c1}{// After}
\PY{n}{Map}\PY{o}{\PYZlt{}}\PY{n}{Integer}\PY{p}{,}\PY{+w}{ }\PY{n}{String}\PY{o}{\PYZgt{}}\PY{+w}{ }\PY{n}{byLen2}\PY{+w}{ }\PY{o}{=}\PY{+w}{ }\PY{n}{words}\PY{p}{.}\PY{n+na}{stream}\PY{p}{(}\PY{p}{)}
\PY{+w}{    }\PY{p}{.}\PY{n+na}{collect}\PY{p}{(}\PY{n}{Collectors}\PY{p}{.}\PY{n+na}{toMap}\PY{p}{(}\PY{n}{String}\PY{p}{:}\PY{p}{:}\PY{n}{length}\PY{p}{,}\PY{+w}{ }\PY{n}{w}\PY{+w}{ }\PY{o}{\PYZhy{}}\PY{o}{\PYZgt{}}\PY{+w}{ }\PY{n}{w}\PY{p}{,}\PY{+w}{ }\PY{p}{(}\PY{n}{a}\PY{p}{,}\PY{+w}{ }\PY{n}{b}\PY{p}{)}\PY{+w}{ }\PY{o}{\PYZhy{}}\PY{o}{\PYZgt{}}\PY{+w}{ }\PY{n}{a}\PY{p}{)}\PY{p}{)}\PY{p}{;}\PY{+w}{ }\PY{c+c1}{// keep first on conflict}
\end{Verbatim}
\end{codebox}

\item 
\textbf{Filtering before \code{groupingBy} when you meant to filter within each group.}
Why it matters: filtering first can make a whole key disappear from the resulting map, changing downstream logic that expects every key to be present.

\begin{codebox}
\begin{Verbatim}[commandchars=\\\{\}]
\PY{c+c1}{// Before \PYZhy{} Sales might vanish entirely if no order has \PYZgt{}1 item}
\PY{n}{Map}\PY{o}{\PYZlt{}}\PY{n}{String}\PY{p}{,}\PY{+w}{ }\PY{n}{List}\PY{o}{\PYZlt{}}\PY{n}{Order}\PY{o}{\PYZgt{}}\PY{o}{\PYZgt{}}\PY{+w}{ }\PY{n}{m}\PY{+w}{ }\PY{o}{=}\PY{+w}{ }\PY{n}{orders}\PY{p}{.}\PY{n+na}{stream}\PY{p}{(}\PY{p}{)}
\PY{+w}{    }\PY{p}{.}\PY{n+na}{filter}\PY{p}{(}\PY{n}{o}\PY{+w}{ }\PY{o}{\PYZhy{}}\PY{o}{\PYZgt{}}\PY{+w}{ }\PY{n}{o}\PY{p}{.}\PY{n+na}{items}\PY{p}{(}\PY{p}{)}\PY{p}{.}\PY{n+na}{size}\PY{p}{(}\PY{p}{)}\PY{+w}{ }\PY{o}{\PYZgt{}}\PY{+w}{ }\PY{l+m+mi}{1}\PY{p}{)}
\PY{+w}{    }\PY{p}{.}\PY{n+na}{collect}\PY{p}{(}\PY{n}{Collectors}\PY{p}{.}\PY{n+na}{groupingBy}\PY{p}{(}\PY{n}{Order}\PY{p}{:}\PY{p}{:}\PY{n}{customer}\PY{p}{)}\PY{p}{)}\PY{p}{;}

\PY{c+c1}{// After \PYZhy{} every customer key is still present, possibly with an empty list}
\PY{n}{Map}\PY{o}{\PYZlt{}}\PY{n}{String}\PY{p}{,}\PY{+w}{ }\PY{n}{List}\PY{o}{\PYZlt{}}\PY{n}{Order}\PY{o}{\PYZgt{}}\PY{o}{\PYZgt{}}\PY{+w}{ }\PY{n}{m2}\PY{+w}{ }\PY{o}{=}\PY{+w}{ }\PY{n}{orders}\PY{p}{.}\PY{n+na}{stream}\PY{p}{(}\PY{p}{)}
\PY{+w}{    }\PY{p}{.}\PY{n+na}{collect}\PY{p}{(}\PY{n}{Collectors}\PY{p}{.}\PY{n+na}{groupingBy}\PY{p}{(}\PY{n}{Order}\PY{p}{:}\PY{p}{:}\PY{n}{customer}\PY{p}{,}
\PY{+w}{        }\PY{n}{Collectors}\PY{p}{.}\PY{n+na}{filtering}\PY{p}{(}\PY{n}{o}\PY{+w}{ }\PY{o}{\PYZhy{}}\PY{o}{\PYZgt{}}\PY{+w}{ }\PY{n}{o}\PY{p}{.}\PY{n+na}{items}\PY{p}{(}\PY{p}{)}\PY{p}{.}\PY{n+na}{size}\PY{p}{(}\PY{p}{)}\PY{+w}{ }\PY{o}{\PYZgt{}}\PY{+w}{ }\PY{l+m+mi}{1}\PY{p}{,}\PY{+w}{ }\PY{n}{Collectors}\PY{p}{.}\PY{n+na}{toList}\PY{p}{(}\PY{p}{)}\PY{p}{)}\PY{p}{)}\PY{p}{)}\PY{p}{;}
\end{Verbatim}
\end{codebox}

\item 
\textbf{Wrapping a checked exception badly, or swallowing it, inside a lambda.}
Why it matters: standard functional interfaces don’t declare \code{throws}; a careless catch can hide real failures instead of surfacing them.

\begin{codebox}
\begin{Verbatim}[commandchars=\\\{\}]
\PY{c+c1}{// Before \PYZhy{} swallows the real error and returns garbage}
\PY{n}{Function}\PY{o}{\PYZlt{}}\PY{n}{String}\PY{p}{,}\PY{+w}{ }\PY{k+kt}{byte}\PY{o}{[}\PY{o}{]}\PY{o}{\PYZgt{}}\PY{+w}{ }\PY{n}{reader}\PY{+w}{ }\PY{o}{=}\PY{+w}{ }\PY{n}{path}\PY{+w}{ }\PY{o}{\PYZhy{}}\PY{o}{\PYZgt{}}\PY{+w}{ }\PY{p}{\PYZob{}}
\PY{+w}{    }\PY{k}{try}\PY{+w}{ }\PY{p}{\PYZob{}}\PY{+w}{ }\PY{k}{return}\PY{+w}{ }\PY{n}{Files}\PY{p}{.}\PY{n+na}{readAllBytes}\PY{p}{(}\PY{n}{Path}\PY{p}{.}\PY{n+na}{of}\PY{p}{(}\PY{n}{path}\PY{p}{)}\PY{p}{)}\PY{p}{;}\PY{+w}{ }\PY{p}{\PYZcb{}}
\PY{+w}{    }\PY{k}{catch}\PY{+w}{ }\PY{p}{(}\PY{n}{IOException}\PY{+w}{ }\PY{n}{e}\PY{p}{)}\PY{+w}{ }\PY{p}{\PYZob{}}\PY{+w}{ }\PY{k}{return}\PY{+w}{ }\PY{k}{new}\PY{+w}{ }\PY{k+kt}{byte}\PY{o}{[}\PY{l+m+mi}{0}\PY{o}{]}\PY{p}{;}\PY{+w}{ }\PY{p}{\PYZcb{}}
\PY{p}{\PYZcb{}}\PY{p}{;}

\PY{c+c1}{// After \PYZhy{} preserve the failure, just change its exception type}
\PY{n}{Function}\PY{o}{\PYZlt{}}\PY{n}{String}\PY{p}{,}\PY{+w}{ }\PY{k+kt}{byte}\PY{o}{[}\PY{o}{]}\PY{o}{\PYZgt{}}\PY{+w}{ }\PY{n}{reader2}\PY{+w}{ }\PY{o}{=}\PY{+w}{ }\PY{n}{path}\PY{+w}{ }\PY{o}{\PYZhy{}}\PY{o}{\PYZgt{}}\PY{+w}{ }\PY{p}{\PYZob{}}
\PY{+w}{    }\PY{k}{try}\PY{+w}{ }\PY{p}{\PYZob{}}\PY{+w}{ }\PY{k}{return}\PY{+w}{ }\PY{n}{Files}\PY{p}{.}\PY{n+na}{readAllBytes}\PY{p}{(}\PY{n}{Path}\PY{p}{.}\PY{n+na}{of}\PY{p}{(}\PY{n}{path}\PY{p}{)}\PY{p}{)}\PY{p}{;}\PY{+w}{ }\PY{p}{\PYZcb{}}
\PY{+w}{    }\PY{k}{catch}\PY{+w}{ }\PY{p}{(}\PY{n}{IOException}\PY{+w}{ }\PY{n}{e}\PY{p}{)}\PY{+w}{ }\PY{p}{\PYZob{}}\PY{+w}{ }\PY{k}{throw}\PY{+w}{ }\PY{k}{new}\PY{+w}{ }\PY{n}{UncheckedIOException}\PY{p}{(}\PY{n}{e}\PY{p}{)}\PY{p}{;}\PY{+w}{ }\PY{p}{\PYZcb{}}
\PY{p}{\PYZcb{}}\PY{p}{;}
\end{Verbatim}
\end{codebox}

\item 
\textbf{Confusing \code{findFirst} and \code{findAny} in code that depends on order.}
Why it matters: \code{findAny} gives no ordering guarantee, especially in parallel; using it where determinism matters introduces flaky behavior.

\begin{codebox}
\begin{Verbatim}[commandchars=\\\{\}]
\PY{c+c1}{// Before \PYZhy{} flaky in tests/parallel runs, expects the smallest even number}
\PY{n}{Optional}\PY{o}{\PYZlt{}}\PY{n}{Integer}\PY{o}{\PYZgt{}}\PY{+w}{ }\PY{n}{firstEven}\PY{+w}{ }\PY{o}{=}\PY{+w}{ }\PY{n}{nums}\PY{p}{.}\PY{n+na}{parallelStream}\PY{p}{(}\PY{p}{)}\PY{p}{.}\PY{n+na}{filter}\PY{p}{(}\PY{n}{n}\PY{+w}{ }\PY{o}{\PYZhy{}}\PY{o}{\PYZgt{}}\PY{+w}{ }\PY{n}{n}\PY{+w}{ }\PY{o}{\PYZpc{}}\PY{+w}{ }\PY{l+m+mi}{2}\PY{+w}{ }\PY{o}{=}\PY{o}{=}\PY{+w}{ }\PY{l+m+mi}{0}\PY{p}{)}\PY{p}{.}\PY{n+na}{findAny}\PY{p}{(}\PY{p}{)}\PY{p}{;}

\PY{c+c1}{// After}
\PY{n}{Optional}\PY{o}{\PYZlt{}}\PY{n}{Integer}\PY{o}{\PYZgt{}}\PY{+w}{ }\PY{n}{firstEven2}\PY{+w}{ }\PY{o}{=}\PY{+w}{ }\PY{n}{nums}\PY{p}{.}\PY{n+na}{parallelStream}\PY{p}{(}\PY{p}{)}\PY{p}{.}\PY{n+na}{filter}\PY{p}{(}\PY{n}{n}\PY{+w}{ }\PY{o}{\PYZhy{}}\PY{o}{\PYZgt{}}\PY{+w}{ }\PY{n}{n}\PY{+w}{ }\PY{o}{\PYZpc{}}\PY{+w}{ }\PY{l+m+mi}{2}\PY{+w}{ }\PY{o}{=}\PY{o}{=}\PY{+w}{ }\PY{l+m+mi}{0}\PY{p}{)}\PY{p}{.}\PY{n+na}{findFirst}\PY{p}{(}\PY{p}{)}\PY{p}{;}
\end{Verbatim}
\end{codebox}

\item 
\textbf{Using an infinite \code{Stream.\allowbreak{}iterate}/\code{Stream.\allowbreak{}generate} without a bound.}
Why it matters: without \code{limit} or a \code{takeWhile}/predicate form, the pipeline never terminates and hangs the thread.

\begin{codebox}
\begin{Verbatim}[commandchars=\\\{\}]
\PY{c+c1}{// Before}
\PY{n}{List}\PY{o}{\PYZlt{}}\PY{n}{Integer}\PY{o}{\PYZgt{}}\PY{+w}{ }\PY{n}{powers}\PY{+w}{ }\PY{o}{=}\PY{+w}{ }\PY{n}{Stream}\PY{p}{.}\PY{n+na}{iterate}\PY{p}{(}\PY{l+m+mi}{1}\PY{p}{,}\PY{+w}{ }\PY{n}{n}\PY{+w}{ }\PY{o}{\PYZhy{}}\PY{o}{\PYZgt{}}\PY{+w}{ }\PY{n}{n}\PY{+w}{ }\PY{o}{*}\PY{+w}{ }\PY{l+m+mi}{2}\PY{p}{)}\PY{p}{.}\PY{n+na}{toList}\PY{p}{(}\PY{p}{)}\PY{p}{;}\PY{+w}{ }\PY{c+c1}{// hangs forever}

\PY{c+c1}{// After}
\PY{n}{List}\PY{o}{\PYZlt{}}\PY{n}{Integer}\PY{o}{\PYZgt{}}\PY{+w}{ }\PY{n}{powers2}\PY{+w}{ }\PY{o}{=}\PY{+w}{ }\PY{n}{Stream}\PY{p}{.}\PY{n+na}{iterate}\PY{p}{(}\PY{l+m+mi}{1}\PY{p}{,}\PY{+w}{ }\PY{n}{n}\PY{+w}{ }\PY{o}{\PYZhy{}}\PY{o}{\PYZgt{}}\PY{+w}{ }\PY{n}{n}\PY{+w}{ }\PY{o}{*}\PY{+w}{ }\PY{l+m+mi}{2}\PY{p}{)}\PY{p}{.}\PY{n+na}{limit}\PY{p}{(}\PY{l+m+mi}{10}\PY{p}{)}\PY{p}{.}\PY{n+na}{toList}\PY{p}{(}\PY{p}{)}\PY{p}{;}
\end{Verbatim}
\end{codebox}

\item 
\textbf{Boxing unnecessarily in numeric pipelines.}
Why it matters: \code{Stream\textless{}\allowbreak{}Integer\textgreater{}} allocates and unboxes on every element; \code{IntStream} avoids it and is measurably faster on large data.

\begin{codebox}
\begin{Verbatim}[commandchars=\\\{\}]
\PY{c+c1}{// Before}
\PY{k+kt}{int}\PY{+w}{ }\PY{n}{total}\PY{+w}{ }\PY{o}{=}\PY{+w}{ }\PY{n}{list}\PY{p}{.}\PY{n+na}{stream}\PY{p}{(}\PY{p}{)}\PY{p}{.}\PY{n+na}{reduce}\PY{p}{(}\PY{l+m+mi}{0}\PY{p}{,}\PY{+w}{ }\PY{n}{Integer}\PY{p}{:}\PY{p}{:}\PY{n}{sum}\PY{p}{)}\PY{p}{;}

\PY{c+c1}{// After}
\PY{k+kt}{int}\PY{+w}{ }\PY{n}{total2}\PY{+w}{ }\PY{o}{=}\PY{+w}{ }\PY{n}{list}\PY{p}{.}\PY{n+na}{stream}\PY{p}{(}\PY{p}{)}\PY{p}{.}\PY{n+na}{mapToInt}\PY{p}{(}\PY{n}{Integer}\PY{p}{:}\PY{p}{:}\PY{n}{intValue}\PY{p}{)}\PY{p}{.}\PY{n+na}{sum}\PY{p}{(}\PY{p}{)}\PY{p}{;}
\end{Verbatim}
\end{codebox}

\item 
\textbf{Confusing anonymous-class \code{this} with lambda \code{this} after a refactor.}
Why it matters: converting an anonymous class to a lambda changes what \code{this} refers to, which can silently break field access or logging that relied on the anonymous class’s own identity.

\begin{codebox}
\begin{Verbatim}[commandchars=\\\{\}]
\PY{c+c1}{// Before \PYZhy{} \PYZsq{}this\PYZsq{} refers to the anonymous Runnable}
\PY{n}{Runnable}\PY{+w}{ }\PY{n}{r}\PY{+w}{ }\PY{o}{=}\PY{+w}{ }\PY{k}{new}\PY{+w}{ }\PY{n}{Runnable}\PY{p}{(}\PY{p}{)}\PY{+w}{ }\PY{p}{\PYZob{}}
\PY{+w}{    }\PY{k+kd}{public}\PY{+w}{ }\PY{k+kt}{void}\PY{+w}{ }\PY{n+nf}{run}\PY{p}{(}\PY{p}{)}\PY{+w}{ }\PY{p}{\PYZob{}}\PY{+w}{ }\PY{n}{log}\PY{p}{(}\PY{k}{this}\PY{p}{.}\PY{n+na}{getClass}\PY{p}{(}\PY{p}{)}\PY{p}{.}\PY{n+na}{getSimpleName}\PY{p}{(}\PY{p}{)}\PY{p}{)}\PY{p}{;}\PY{+w}{ }\PY{p}{\PYZcb{}}
\PY{p}{\PYZcb{}}\PY{p}{;}

\PY{c+c1}{// After (lambda) \PYZhy{} \PYZsq{}this\PYZsq{} now refers to the ENCLOSING instance, not the lambda}
\PY{n}{Runnable}\PY{+w}{ }\PY{n}{r2}\PY{+w}{ }\PY{o}{=}\PY{+w}{ }\PY{p}{(}\PY{p}{)}\PY{+w}{ }\PY{o}{\PYZhy{}}\PY{o}{\PYZgt{}}\PY{+w}{ }\PY{n}{log}\PY{p}{(}\PY{k}{this}\PY{p}{.}\PY{n+na}{getClass}\PY{p}{(}\PY{p}{)}\PY{p}{.}\PY{n+na}{getSimpleName}\PY{p}{(}\PY{p}{)}\PY{p}{)}\PY{p}{;}\PY{+w}{ }\PY{c+c1}{// meaning changed \PYZhy{} verify this is intended}
\end{Verbatim}
\end{codebox}

\item 
\textbf{Parallelizing a poorly-splittable source and expecting a speedup.}
Why it matters: \code{LinkedList}/iterator-backed sources can’t split efficiently, so \code{parallelStream()} on them often gives little to no benefit while still paying coordination overhead.

\begin{codebox}
\begin{Verbatim}[commandchars=\\\{\}]
\PY{c+c1}{// Before}
\PY{n}{LinkedList}\PY{o}{\PYZlt{}}\PY{n}{Integer}\PY{o}{\PYZgt{}}\PY{+w}{ }\PY{n}{data}\PY{+w}{ }\PY{o}{=}\PY{+w}{ }\PY{k}{new}\PY{+w}{ }\PY{n}{LinkedList}\PY{o}{\PYZlt{}}\PY{o}{\PYZgt{}}\PY{p}{(}\PY{n}{bigList}\PY{p}{)}\PY{p}{;}
\PY{n}{data}\PY{p}{.}\PY{n+na}{parallelStream}\PY{p}{(}\PY{p}{)}\PY{p}{.}\PY{n+na}{map}\PY{p}{(}\PY{k}{this}\PY{p}{:}\PY{p}{:}\PY{n}{heavyWork}\PY{p}{)}\PY{p}{.}\PY{n+na}{toList}\PY{p}{(}\PY{p}{)}\PY{p}{;}\PY{+w}{ }\PY{c+c1}{// splitting degrades to near\PYZhy{}linear walk}

\PY{c+c1}{// After \PYZhy{} use a splittable source}
\PY{n}{List}\PY{o}{\PYZlt{}}\PY{n}{Integer}\PY{o}{\PYZgt{}}\PY{+w}{ }\PY{n}{data2}\PY{+w}{ }\PY{o}{=}\PY{+w}{ }\PY{k}{new}\PY{+w}{ }\PY{n}{ArrayList}\PY{o}{\PYZlt{}}\PY{o}{\PYZgt{}}\PY{p}{(}\PY{n}{bigList}\PY{p}{)}\PY{p}{;}
\PY{n}{data2}\PY{p}{.}\PY{n+na}{parallelStream}\PY{p}{(}\PY{p}{)}\PY{p}{.}\PY{n+na}{map}\PY{p}{(}\PY{k}{this}\PY{p}{:}\PY{p}{:}\PY{n}{heavyWork}\PY{p}{)}\PY{p}{.}\PY{n+na}{toList}\PY{p}{(}\PY{p}{)}\PY{p}{;}
\end{Verbatim}
\end{codebox}

\item 
\textbf{Relying on \code{forEach} order in a parallel stream.}
Why it matters: parallel \code{forEach} makes no ordering guarantee; code that assumes sequential-looking output will be flaky under load.

\begin{codebox}
\begin{Verbatim}[commandchars=\\\{\}]
\PY{c+c1}{// Before \PYZhy{} assumes output prints 1..N in order}
\PY{n}{nums}\PY{p}{.}\PY{n+na}{parallelStream}\PY{p}{(}\PY{p}{)}\PY{p}{.}\PY{n+na}{forEach}\PY{p}{(}\PY{n}{System}\PY{p}{.}\PY{n+na}{out}\PY{p}{:}\PY{p}{:}\PY{n}{println}\PY{p}{)}\PY{p}{;}

\PY{c+c1}{// After \PYZhy{} if order matters, either go sequential or collect then iterate}
\PY{n}{List}\PY{o}{\PYZlt{}}\PY{n}{Integer}\PY{o}{\PYZgt{}}\PY{+w}{ }\PY{n}{ordered}\PY{+w}{ }\PY{o}{=}\PY{+w}{ }\PY{n}{nums}\PY{p}{.}\PY{n+na}{parallelStream}\PY{p}{(}\PY{p}{)}\PY{p}{.}\PY{n+na}{map}\PY{p}{(}\PY{n}{n}\PY{+w}{ }\PY{o}{\PYZhy{}}\PY{o}{\PYZgt{}}\PY{+w}{ }\PY{n}{n}\PY{+w}{ }\PY{o}{*}\PY{+w}{ }\PY{l+m+mi}{2}\PY{p}{)}\PY{p}{.}\PY{n+na}{toList}\PY{p}{(}\PY{p}{)}\PY{p}{;}\PY{+w}{ }\PY{c+c1}{// preserves encounter order}
\PY{n}{ordered}\PY{p}{.}\PY{n+na}{forEach}\PY{p}{(}\PY{n}{System}\PY{p}{.}\PY{n+na}{out}\PY{p}{:}\PY{p}{:}\PY{n}{println}\PY{p}{)}\PY{p}{;}
\end{Verbatim}
\end{codebox}

\item 
\textbf{Writing an overly clever one-liner stream pipeline that’s hard to review.}
Why it matters: dense chained pipelines with side effects buried in \code{map}/\code{peek} are hard to test and hide bugs; reviewers should push back on readability, not just correctness.

\begin{codebox}
\begin{Verbatim}[commandchars=\\\{\}]
\PY{c+c1}{// Before \PYZhy{} hard to review, mixes transformation with hidden side effects}
\PY{n}{List}\PY{o}{\PYZlt{}}\PY{n}{String}\PY{o}{\PYZgt{}}\PY{+w}{ }\PY{n}{r}\PY{+w}{ }\PY{o}{=}\PY{+w}{ }\PY{n}{data}\PY{p}{.}\PY{n+na}{stream}\PY{p}{(}\PY{p}{)}\PY{p}{.}\PY{n+na}{peek}\PY{p}{(}\PY{k}{this}\PY{p}{:}\PY{p}{:}\PY{n}{audit}\PY{p}{)}\PY{p}{.}\PY{n+na}{map}\PY{p}{(}\PY{k}{this}\PY{p}{:}\PY{p}{:}\PY{n}{transform}\PY{p}{)}\PY{p}{.}\PY{n+na}{filter}\PY{p}{(}\PY{n}{Objects}\PY{p}{:}\PY{p}{:}\PY{n}{nonNull}\PY{p}{)}
\PY{+w}{    }\PY{p}{.}\PY{n+na}{peek}\PY{p}{(}\PY{k}{this}\PY{p}{:}\PY{p}{:}\PY{n}{cache}\PY{p}{)}\PY{p}{.}\PY{n+na}{sorted}\PY{p}{(}\PY{p}{)}\PY{p}{.}\PY{n+na}{distinct}\PY{p}{(}\PY{p}{)}\PY{p}{.}\PY{n+na}{toList}\PY{p}{(}\PY{p}{)}\PY{p}{;}

\PY{c+c1}{// After \PYZhy{} explicit steps, side effects named and separated}
\PY{n}{List}\PY{o}{\PYZlt{}}\PY{n}{String}\PY{o}{\PYZgt{}}\PY{+w}{ }\PY{n}{transformed}\PY{+w}{ }\PY{o}{=}\PY{+w}{ }\PY{n}{data}\PY{p}{.}\PY{n+na}{stream}\PY{p}{(}\PY{p}{)}\PY{p}{.}\PY{n+na}{map}\PY{p}{(}\PY{k}{this}\PY{p}{:}\PY{p}{:}\PY{n}{transform}\PY{p}{)}\PY{p}{.}\PY{n+na}{filter}\PY{p}{(}\PY{n}{Objects}\PY{p}{:}\PY{p}{:}\PY{n}{nonNull}\PY{p}{)}\PY{p}{.}\PY{n+na}{toList}\PY{p}{(}\PY{p}{)}\PY{p}{;}
\PY{n}{transformed}\PY{p}{.}\PY{n+na}{forEach}\PY{p}{(}\PY{k}{this}\PY{p}{:}\PY{p}{:}\PY{n}{audit}\PY{p}{)}\PY{p}{;}
\PY{n}{List}\PY{o}{\PYZlt{}}\PY{n}{String}\PY{o}{\PYZgt{}}\PY{+w}{ }\PY{n}{r2}\PY{+w}{ }\PY{o}{=}\PY{+w}{ }\PY{n}{transformed}\PY{p}{.}\PY{n+na}{stream}\PY{p}{(}\PY{p}{)}\PY{p}{.}\PY{n+na}{sorted}\PY{p}{(}\PY{p}{)}\PY{p}{.}\PY{n+na}{distinct}\PY{p}{(}\PY{p}{)}\PY{p}{.}\PY{n+na}{toList}\PY{p}{(}\PY{p}{)}\PY{p}{;}
\end{Verbatim}
\end{codebox}

\item 
\textbf{Treating a custom \code{Collector}’s combiner as optional in a stream that might run in parallel.}
Why it matters: the combiner is only exercised under parallel execution; a broken or missing combiner passes all sequential tests but produces wrong results the moment someone adds \code{.\allowbreak{}parallel()}.

\begin{codebox}
\begin{Verbatim}[commandchars=\\\{\}]
\PY{c+c1}{// Before \PYZhy{} combiner just picks one side, silently dropping data when split}
\PY{n}{Collector}\PY{o}{\PYZlt{}}\PY{n}{String}\PY{p}{,}\PY{+w}{ }\PY{n}{StringBuilder}\PY{p}{,}\PY{+w}{ }\PY{n}{String}\PY{o}{\PYZgt{}}\PY{+w}{ }\PY{n}{broken}\PY{+w}{ }\PY{o}{=}\PY{+w}{ }\PY{n}{Collector}\PY{p}{.}\PY{n+na}{of}\PY{p}{(}
\PY{+w}{    }\PY{n}{StringBuilder}\PY{p}{:}\PY{p}{:}\PY{k}{new}\PY{p}{,}
\PY{+w}{    }\PY{p}{(}\PY{n}{sb}\PY{p}{,}\PY{+w}{ }\PY{n}{s}\PY{p}{)}\PY{+w}{ }\PY{o}{\PYZhy{}}\PY{o}{\PYZgt{}}\PY{+w}{ }\PY{n}{sb}\PY{p}{.}\PY{n+na}{append}\PY{p}{(}\PY{n}{s}\PY{p}{)}\PY{p}{,}
\PY{+w}{    }\PY{p}{(}\PY{n}{sb1}\PY{p}{,}\PY{+w}{ }\PY{n}{sb2}\PY{p}{)}\PY{+w}{ }\PY{o}{\PYZhy{}}\PY{o}{\PYZgt{}}\PY{+w}{ }\PY{n}{sb1}\PY{p}{,}\PY{+w}{      }\PY{c+c1}{// BUG: drops sb2\PYZsq{}s contents entirely on merge}
\PY{+w}{    }\PY{n}{StringBuilder}\PY{p}{:}\PY{p}{:}\PY{n}{toString}\PY{p}{)}\PY{p}{;}

\PY{c+c1}{// After \PYZhy{} combiner actually merges both partial results}
\PY{n}{Collector}\PY{o}{\PYZlt{}}\PY{n}{String}\PY{p}{,}\PY{+w}{ }\PY{n}{StringBuilder}\PY{p}{,}\PY{+w}{ }\PY{n}{String}\PY{o}{\PYZgt{}}\PY{+w}{ }\PY{n}{fixed}\PY{+w}{ }\PY{o}{=}\PY{+w}{ }\PY{n}{Collector}\PY{p}{.}\PY{n+na}{of}\PY{p}{(}
\PY{+w}{    }\PY{n}{StringBuilder}\PY{p}{:}\PY{p}{:}\PY{k}{new}\PY{p}{,}
\PY{+w}{    }\PY{p}{(}\PY{n}{sb}\PY{p}{,}\PY{+w}{ }\PY{n}{s}\PY{p}{)}\PY{+w}{ }\PY{o}{\PYZhy{}}\PY{o}{\PYZgt{}}\PY{+w}{ }\PY{n}{sb}\PY{p}{.}\PY{n+na}{append}\PY{p}{(}\PY{n}{s}\PY{p}{)}\PY{p}{,}
\PY{+w}{    }\PY{p}{(}\PY{n}{sb1}\PY{p}{,}\PY{+w}{ }\PY{n}{sb2}\PY{p}{)}\PY{+w}{ }\PY{o}{\PYZhy{}}\PY{o}{\PYZgt{}}\PY{+w}{ }\PY{n}{sb1}\PY{p}{.}\PY{n+na}{append}\PY{p}{(}\PY{n}{sb2}\PY{p}{)}\PY{p}{,}\PY{+w}{ }\PY{c+c1}{// correct merge}
\PY{+w}{    }\PY{n}{StringBuilder}\PY{p}{:}\PY{p}{:}\PY{n}{toString}\PY{p}{)}\PY{p}{;}
\end{Verbatim}
\end{codebox}

\end{enumerate}


\setcounter{chapter}{8}
\chapter{Date, Time \& Localization}
\label{chap:9}
\label{h:9:9-date-time--localization}
Dates, times, and text that must work in more than one country trip up even experienced Java developers. This chapter covers the modern \code{java.\allowbreak{}time} API, how to format and parse dates safely, and how to write code that behaves correctly for users in different languages, countries, and time zones. By the end you should be able to spot the classic review-time mistakes: mutable date bugs, wrong time zone assumptions, locale-sensitive string comparisons, and formatter thread-safety issues.

\section[Date and Time API (java.time)]{Date and Time API (java.time)}
\label{h:9:date-and-time-api-javatime}
\subsection[Why the old Date/Calendar/SimpleDateFormat were replaced]{Why the old \code{Date}/\code{Calendar}/\code{SimpleDateFormat} were replaced}
\label{h:9:why-the-old-datecalendarsimpledateformat-were-replaced}
Before Java 8, date and time handling used \code{java.\allowbreak{}util.\allowbreak{}Date}, \code{java.\allowbreak{}util.\allowbreak{}Calendar}, and \code{java.\allowbreak{}text.\allowbreak{}Simple\allowbreak{}Date\allowbreak{}Format}. These classes had deep design problems:

\begin{itemize}
\item 
\textbf{Mutable}: A \code{Date} object can be changed after creation. Passing one to another method risks it being silently mutated.

\item 
\textbf{Confusing month numbering}: Months in \code{Calendar} are zero-based (January is \code{0}), a constant source of off-by-one bugs.

\item 
\textbf{Not thread-safe}: \code{SimpleDateFormat} is not safe to share across threads, but people often make it a \code{static} field “for performance” and get corrupted output under load.

\item 
\textbf{Poor API design}: \code{Date} mixes up “a point in time” with “a calendar date,” which are conceptually different things.

\item 
\textbf{No clear time zone model}: It was easy to write code that silently used the JVM’s default time zone.

\end{itemize}

\code{java.\allowbreak{}time} (JSR-310, led by Stephen Colebourne, the author of Joda-Time) fixed all of this: every type is immutable, thread-safe, and clearly named for what it represents.

\begin{codebox}
\begin{Verbatim}[commandchars=\\\{\}]
\PY{k+kn}{import}\PY{+w}{ }\PY{n+nn}{java.util.Date}\PY{p}{;}
\PY{k+kn}{import}\PY{+w}{ }\PY{n+nn}{java.util.Calendar}\PY{p}{;}

\PY{k+kd}{public}\PY{+w}{ }\PY{k+kd}{class} \PY{n+nc}{LegacyProblems}\PY{+w}{ }\PY{p}{\PYZob{}}
\PY{+w}{    }\PY{k+kd}{public}\PY{+w}{ }\PY{k+kd}{static}\PY{+w}{ }\PY{k+kt}{void}\PY{+w}{ }\PY{n+nf}{main}\PY{p}{(}\PY{n}{String}\PY{o}{[}\PY{o}{]}\PY{+w}{ }\PY{n}{args}\PY{p}{)}\PY{+w}{ }\PY{p}{\PYZob{}}
\PY{+w}{        }\PY{n}{Date}\PY{+w}{ }\PY{n}{date}\PY{+w}{ }\PY{o}{=}\PY{+w}{ }\PY{k}{new}\PY{+w}{ }\PY{n}{Date}\PY{p}{(}\PY{p}{)}\PY{p}{;}
\PY{+w}{        }\PY{n}{Calendar}\PY{+w}{ }\PY{n}{cal}\PY{+w}{ }\PY{o}{=}\PY{+w}{ }\PY{n}{Calendar}\PY{p}{.}\PY{n+na}{getInstance}\PY{p}{(}\PY{p}{)}\PY{p}{;}
\PY{+w}{        }\PY{n}{cal}\PY{p}{.}\PY{n+na}{set}\PY{p}{(}\PY{l+m+mi}{2026}\PY{p}{,}\PY{+w}{ }\PY{n}{Calendar}\PY{p}{.}\PY{n+na}{JANUARY}\PY{p}{,}\PY{+w}{ }\PY{l+m+mi}{15}\PY{p}{)}\PY{p}{;}\PY{+w}{ }\PY{c+c1}{// JANUARY == 0, easy to get wrong}
\PY{+w}{        }\PY{n}{System}\PY{p}{.}\PY{n+na}{out}\PY{p}{.}\PY{n+na}{println}\PY{p}{(}\PY{n}{cal}\PY{p}{.}\PY{n+na}{getTime}\PY{p}{(}\PY{p}{)}\PY{p}{)}\PY{p}{;}
\PY{+w}{        }\PY{c+c1}{// Output (example): Thu Jan 15 00:00:00 UTC 2026}

\PY{+w}{        }\PY{c+c1}{// Mutation trap: passing \PYZdq{}date\PYZdq{} elsewhere lets other code change it.}
\PY{+w}{        }\PY{n}{Date}\PY{+w}{ }\PY{n}{copy}\PY{+w}{ }\PY{o}{=}\PY{+w}{ }\PY{n}{date}\PY{p}{;}
\PY{+w}{        }\PY{n}{copy}\PY{p}{.}\PY{n+na}{setTime}\PY{p}{(}\PY{l+m+mi}{0}\PY{n}{L}\PY{p}{)}\PY{p}{;}\PY{+w}{ }\PY{c+c1}{// also changes the original \PYZdq{}date\PYZdq{} reference!}
\PY{+w}{        }\PY{n}{System}\PY{p}{.}\PY{n+na}{out}\PY{p}{.}\PY{n+na}{println}\PY{p}{(}\PY{n}{date}\PY{p}{.}\PY{n+na}{getTime}\PY{p}{(}\PY{p}{)}\PY{p}{)}\PY{p}{;}
\PY{+w}{        }\PY{c+c1}{// Output: 0  \PYZlt{}\PYZhy{}\PYZhy{} surprising, \PYZdq{}date\PYZdq{} was mutated through \PYZdq{}copy\PYZdq{}}
\PY{+w}{    }\PY{p}{\PYZcb{}}
\PY{p}{\PYZcb{}}
\end{Verbatim}
\end{codebox}

{\tablefont
\begin{xltabular}{\linewidth}{@{}L{0.667}L{1.154}L{1.179}@{}}
\toprule
\rowcolor{tablehdr}
\thd{Legacy type} & \thd{Modern replacement} & \thd{Notes} \\
\midrule
\endhead
\code{java.\allowbreak{}util.\allowbreak{}Date} & \code{Instant} (point in time) or \code{LocalDateTime} (no zone) & \code{Date} conflated “instant” and “calendar” concepts \\
\code{java.\allowbreak{}util.\allowbreak{}Calendar} & \code{ZonedDateTime} & Zero-based months gone; clear zone handling \\
\code{java.\allowbreak{}util.\allowbreak{}Gregorian\allowbreak{}Calendar} & \code{ZonedDateTime} / \code{LocalDate} &  \\
\code{java.\allowbreak{}text.\allowbreak{}Simple\allowbreak{}Date\allowbreak{}Format} & \code{DateTimeFormatter} & Immutable and thread-safe \\
\code{java.\allowbreak{}sql.\allowbreak{}Date} / \code{Timestamp} & \code{LocalDate} / \code{Instant} (via JDBC 4.2+ \code{getObject}) & Modern JDBC drivers support \code{java.\allowbreak{}time} directly \\
\code{TimeZone} & \code{ZoneId} / \code{ZoneOffset} &  \\
\bottomrule
\end{xltabular}
}

\subsection[The core types]{The core types}
\label{h:9:the-core-types}
\code{java.\allowbreak{}time} splits “date and time” into precise, single-purpose types. Picking the right one is a common review topic.

{\tablefont
\begin{xltabular}{\linewidth}{@{}L{0.442}L{1.339}L{1.220}@{}}
\toprule
\rowcolor{tablehdr}
\thd{Type} & \thd{Represents} & \thd{Example use case} \\
\midrule
\endhead
\code{LocalDate} & Date only, no time, no zone & Birthdate, invoice date \\
\code{LocalTime} & Time only, no date, no zone & Store opening time (09:00) \\
\code{LocalDateTime} & Date + time, no zone & “Meeting at 2026-08-07T14:00” in an unspecified zone \\
\code{ZonedDateTime} & Date + time + full time zone rules (\code{ZoneId}) & Flight departure time in \code{Europe/\allowbreak{}Amsterdam} \\
\code{OffsetDateTime} & Date + time + fixed UTC offset (no DST rules) & Timestamps in APIs / logs that need an offset but not a full zone \\
\code{Instant} & A point on the machine timeline (nanoseconds since epoch) & Event timestamps, “now” for logging \\
\code{Duration} & Amount of time in seconds/nanoseconds & “Timeout after 30 seconds” \\
\code{Period} & Amount of time in years/months/days & “Subscription lasts 1 year” \\
\code{Year} & A calendar year, e.g. 2026 & Leap-year checks \\
\code{YearMonth} & Year + month, no day & Credit card expiry (MM/YYYY) \\
\code{MonthDay} & Month + day, no year & Recurring anniversary, without a year \\
\code{DayOfWeek} & Enum: MONDAY…SUNDAY & “Is this a business day?” \\
\code{ZoneId} & A geographical/political time zone, e.g. \code{Europe/\allowbreak{}Berlin} & Zone-aware scheduling \\
\code{ZoneOffset} & A fixed offset from UTC, e.g. \code{+\allowbreak{}02:\allowbreak{}00} & Offset math without DST rules \\
\bottomrule
\end{xltabular}
}

\begin{codebox}
\begin{Verbatim}[commandchars=\\\{\}]
\PY{k+kn}{import}\PY{+w}{ }\PY{n+nn}{java.time.*}\PY{p}{;}

\PY{k+kd}{public}\PY{+w}{ }\PY{k+kd}{class} \PY{n+nc}{CoreTypesDemo}\PY{+w}{ }\PY{p}{\PYZob{}}
\PY{+w}{    }\PY{k+kd}{public}\PY{+w}{ }\PY{k+kd}{static}\PY{+w}{ }\PY{k+kt}{void}\PY{+w}{ }\PY{n+nf}{main}\PY{p}{(}\PY{n}{String}\PY{o}{[}\PY{o}{]}\PY{+w}{ }\PY{n}{args}\PY{p}{)}\PY{+w}{ }\PY{p}{\PYZob{}}
\PY{+w}{        }\PY{n}{LocalDate}\PY{+w}{ }\PY{n}{date}\PY{+w}{ }\PY{o}{=}\PY{+w}{ }\PY{n}{LocalDate}\PY{p}{.}\PY{n+na}{of}\PY{p}{(}\PY{l+m+mi}{2026}\PY{p}{,}\PY{+w}{ }\PY{l+m+mi}{8}\PY{p}{,}\PY{+w}{ }\PY{l+m+mi}{7}\PY{p}{)}\PY{p}{;}
\PY{+w}{        }\PY{n}{LocalTime}\PY{+w}{ }\PY{n}{time}\PY{+w}{ }\PY{o}{=}\PY{+w}{ }\PY{n}{LocalTime}\PY{p}{.}\PY{n+na}{of}\PY{p}{(}\PY{l+m+mi}{14}\PY{p}{,}\PY{+w}{ }\PY{l+m+mi}{30}\PY{p}{)}\PY{p}{;}
\PY{+w}{        }\PY{n}{LocalDateTime}\PY{+w}{ }\PY{n}{dateTime}\PY{+w}{ }\PY{o}{=}\PY{+w}{ }\PY{n}{LocalDateTime}\PY{p}{.}\PY{n+na}{of}\PY{p}{(}\PY{n}{date}\PY{p}{,}\PY{+w}{ }\PY{n}{time}\PY{p}{)}\PY{p}{;}
\PY{+w}{        }\PY{n}{ZonedDateTime}\PY{+w}{ }\PY{n}{zoned}\PY{+w}{ }\PY{o}{=}\PY{+w}{ }\PY{n}{dateTime}\PY{p}{.}\PY{n+na}{atZone}\PY{p}{(}\PY{n}{ZoneId}\PY{p}{.}\PY{n+na}{of}\PY{p}{(}\PY{l+s}{\PYZdq{}}\PY{l+s}{Europe/Amsterdam}\PY{l+s}{\PYZdq{}}\PY{p}{)}\PY{p}{)}\PY{p}{;}
\PY{+w}{        }\PY{n}{OffsetDateTime}\PY{+w}{ }\PY{n}{offset}\PY{+w}{ }\PY{o}{=}\PY{+w}{ }\PY{n}{zoned}\PY{p}{.}\PY{n+na}{toOffsetDateTime}\PY{p}{(}\PY{p}{)}\PY{p}{;}
\PY{+w}{        }\PY{n}{Instant}\PY{+w}{ }\PY{n}{instant}\PY{+w}{ }\PY{o}{=}\PY{+w}{ }\PY{n}{zoned}\PY{p}{.}\PY{n+na}{toInstant}\PY{p}{(}\PY{p}{)}\PY{p}{;}

\PY{+w}{        }\PY{n}{System}\PY{p}{.}\PY{n+na}{out}\PY{p}{.}\PY{n+na}{println}\PY{p}{(}\PY{n}{date}\PY{p}{)}\PY{p}{;}\PY{+w}{      }\PY{c+c1}{// 2026\PYZhy{}08\PYZhy{}07}
\PY{+w}{        }\PY{n}{System}\PY{p}{.}\PY{n+na}{out}\PY{p}{.}\PY{n+na}{println}\PY{p}{(}\PY{n}{time}\PY{p}{)}\PY{p}{;}\PY{+w}{      }\PY{c+c1}{// 14:30}
\PY{+w}{        }\PY{n}{System}\PY{p}{.}\PY{n+na}{out}\PY{p}{.}\PY{n+na}{println}\PY{p}{(}\PY{n}{dateTime}\PY{p}{)}\PY{p}{;}\PY{+w}{  }\PY{c+c1}{// 2026\PYZhy{}08\PYZhy{}07T14:30}
\PY{+w}{        }\PY{n}{System}\PY{p}{.}\PY{n+na}{out}\PY{p}{.}\PY{n+na}{println}\PY{p}{(}\PY{n}{zoned}\PY{p}{)}\PY{p}{;}\PY{+w}{     }\PY{c+c1}{// 2026\PYZhy{}08\PYZhy{}07T14:30+02:00[Europe/Amsterdam]}
\PY{+w}{        }\PY{n}{System}\PY{p}{.}\PY{n+na}{out}\PY{p}{.}\PY{n+na}{println}\PY{p}{(}\PY{n}{offset}\PY{p}{)}\PY{p}{;}\PY{+w}{    }\PY{c+c1}{// 2026\PYZhy{}08\PYZhy{}07T14:30+02:00}
\PY{+w}{        }\PY{n}{System}\PY{p}{.}\PY{n+na}{out}\PY{p}{.}\PY{n+na}{println}\PY{p}{(}\PY{n}{instant}\PY{p}{)}\PY{p}{;}\PY{+w}{   }\PY{c+c1}{// 2026\PYZhy{}08\PYZhy{}07T12:30:00Z}

\PY{+w}{        }\PY{n}{Duration}\PY{+w}{ }\PY{n}{meetingLength}\PY{+w}{ }\PY{o}{=}\PY{+w}{ }\PY{n}{Duration}\PY{p}{.}\PY{n+na}{ofMinutes}\PY{p}{(}\PY{l+m+mi}{45}\PY{p}{)}\PY{p}{;}
\PY{+w}{        }\PY{n}{Period}\PY{+w}{ }\PY{n}{subscription}\PY{+w}{ }\PY{o}{=}\PY{+w}{ }\PY{n}{Period}\PY{p}{.}\PY{n+na}{ofYears}\PY{p}{(}\PY{l+m+mi}{1}\PY{p}{)}\PY{p}{;}
\PY{+w}{        }\PY{n}{System}\PY{p}{.}\PY{n+na}{out}\PY{p}{.}\PY{n+na}{println}\PY{p}{(}\PY{n}{meetingLength}\PY{p}{)}\PY{p}{;}\PY{+w}{ }\PY{c+c1}{// PT45M}
\PY{+w}{        }\PY{n}{System}\PY{p}{.}\PY{n+na}{out}\PY{p}{.}\PY{n+na}{println}\PY{p}{(}\PY{n}{subscription}\PY{p}{)}\PY{p}{;}\PY{+w}{  }\PY{c+c1}{// P1Y}

\PY{+w}{        }\PY{n}{YearMonth}\PY{+w}{ }\PY{n}{expiry}\PY{+w}{ }\PY{o}{=}\PY{+w}{ }\PY{n}{YearMonth}\PY{p}{.}\PY{n+na}{of}\PY{p}{(}\PY{l+m+mi}{2027}\PY{p}{,}\PY{+w}{ }\PY{l+m+mi}{12}\PY{p}{)}\PY{p}{;}
\PY{+w}{        }\PY{n}{MonthDay}\PY{+w}{ }\PY{n}{anniversary}\PY{+w}{ }\PY{o}{=}\PY{+w}{ }\PY{n}{MonthDay}\PY{p}{.}\PY{n+na}{of}\PY{p}{(}\PY{l+m+mi}{6}\PY{p}{,}\PY{+w}{ }\PY{l+m+mi}{15}\PY{p}{)}\PY{p}{;}
\PY{+w}{        }\PY{n}{DayOfWeek}\PY{+w}{ }\PY{n}{dow}\PY{+w}{ }\PY{o}{=}\PY{+w}{ }\PY{n}{date}\PY{p}{.}\PY{n+na}{getDayOfWeek}\PY{p}{(}\PY{p}{)}\PY{p}{;}
\PY{+w}{        }\PY{n}{System}\PY{p}{.}\PY{n+na}{out}\PY{p}{.}\PY{n+na}{println}\PY{p}{(}\PY{n}{expiry}\PY{p}{)}\PY{p}{;}\PY{+w}{      }\PY{c+c1}{// 2027\PYZhy{}12}
\PY{+w}{        }\PY{n}{System}\PY{p}{.}\PY{n+na}{out}\PY{p}{.}\PY{n+na}{println}\PY{p}{(}\PY{n}{anniversary}\PY{p}{)}\PY{p}{;}\PY{+w}{ }\PY{c+c1}{// \PYZhy{}\PYZhy{}06\PYZhy{}15}
\PY{+w}{        }\PY{n}{System}\PY{p}{.}\PY{n+na}{out}\PY{p}{.}\PY{n+na}{println}\PY{p}{(}\PY{n}{dow}\PY{p}{)}\PY{p}{;}\PY{+w}{         }\PY{c+c1}{// FRIDAY}
\PY{+w}{    }\PY{p}{\PYZcb{}}
\PY{p}{\PYZcb{}}
\end{Verbatim}
\end{codebox}

\subsection[Immutability and the fluent plus/minus/with API]{Immutability and the fluent \code{plus}/\code{minus}/\code{with} API}
\label{h:9:immutability-and-the-fluent-plusminuswith-api}
Every \code{java.\allowbreak{}time} type is immutable: methods like \code{plusDays}, \code{minusMonths}, and \code{withYear} return a \textbf{new} object instead of mutating the original. This removes an entire category of bugs where a shared date gets silently changed. It also means you must always use the return value — a very common code-review pitfall is calling \code{date.\allowbreak{}plusDays(\allowbreak{}1);} and discarding the result.

\begin{codebox}
\begin{Verbatim}[commandchars=\\\{\}]
\PY{k+kn}{import}\PY{+w}{ }\PY{n+nn}{java.time.LocalDate}\PY{p}{;}

\PY{k+kd}{public}\PY{+w}{ }\PY{k+kd}{class} \PY{n+nc}{ImmutabilityDemo}\PY{+w}{ }\PY{p}{\PYZob{}}
\PY{+w}{    }\PY{k+kd}{public}\PY{+w}{ }\PY{k+kd}{static}\PY{+w}{ }\PY{k+kt}{void}\PY{+w}{ }\PY{n+nf}{main}\PY{p}{(}\PY{n}{String}\PY{o}{[}\PY{o}{]}\PY{+w}{ }\PY{n}{args}\PY{p}{)}\PY{+w}{ }\PY{p}{\PYZob{}}
\PY{+w}{        }\PY{n}{LocalDate}\PY{+w}{ }\PY{n}{original}\PY{+w}{ }\PY{o}{=}\PY{+w}{ }\PY{n}{LocalDate}\PY{p}{.}\PY{n+na}{of}\PY{p}{(}\PY{l+m+mi}{2026}\PY{p}{,}\PY{+w}{ }\PY{l+m+mi}{8}\PY{p}{,}\PY{+w}{ }\PY{l+m+mi}{7}\PY{p}{)}\PY{p}{;}

\PY{+w}{        }\PY{n}{LocalDate}\PY{+w}{ }\PY{n}{nextWeek}\PY{+w}{ }\PY{o}{=}\PY{+w}{ }\PY{n}{original}\PY{p}{.}\PY{n+na}{plusWeeks}\PY{p}{(}\PY{l+m+mi}{1}\PY{p}{)}\PY{p}{;}
\PY{+w}{        }\PY{n}{LocalDate}\PY{+w}{ }\PY{n}{lastYear}\PY{+w}{ }\PY{o}{=}\PY{+w}{ }\PY{n}{original}\PY{p}{.}\PY{n+na}{minusYears}\PY{p}{(}\PY{l+m+mi}{1}\PY{p}{)}\PY{p}{;}
\PY{+w}{        }\PY{n}{LocalDate}\PY{+w}{ }\PY{n}{changedYear}\PY{+w}{ }\PY{o}{=}\PY{+w}{ }\PY{n}{original}\PY{p}{.}\PY{n+na}{withYear}\PY{p}{(}\PY{l+m+mi}{2030}\PY{p}{)}\PY{p}{;}

\PY{+w}{        }\PY{n}{System}\PY{p}{.}\PY{n+na}{out}\PY{p}{.}\PY{n+na}{println}\PY{p}{(}\PY{n}{original}\PY{p}{)}\PY{p}{;}\PY{+w}{    }\PY{c+c1}{// 2026\PYZhy{}08\PYZhy{}07  (unchanged!)}
\PY{+w}{        }\PY{n}{System}\PY{p}{.}\PY{n+na}{out}\PY{p}{.}\PY{n+na}{println}\PY{p}{(}\PY{n}{nextWeek}\PY{p}{)}\PY{p}{;}\PY{+w}{    }\PY{c+c1}{// 2026\PYZhy{}08\PYZhy{}14}
\PY{+w}{        }\PY{n}{System}\PY{p}{.}\PY{n+na}{out}\PY{p}{.}\PY{n+na}{println}\PY{p}{(}\PY{n}{lastYear}\PY{p}{)}\PY{p}{;}\PY{+w}{    }\PY{c+c1}{// 2025\PYZhy{}08\PYZhy{}07}
\PY{+w}{        }\PY{n}{System}\PY{p}{.}\PY{n+na}{out}\PY{p}{.}\PY{n+na}{println}\PY{p}{(}\PY{n}{changedYear}\PY{p}{)}\PY{p}{;}\PY{+w}{ }\PY{c+c1}{// 2030\PYZhy{}08\PYZhy{}07}

\PY{+w}{        }\PY{c+c1}{// Fluent chaining \PYZhy{} each call returns a new instance}
\PY{+w}{        }\PY{n}{LocalDate}\PY{+w}{ }\PY{n}{result}\PY{+w}{ }\PY{o}{=}\PY{+w}{ }\PY{n}{original}\PY{p}{.}\PY{n+na}{plusMonths}\PY{p}{(}\PY{l+m+mi}{2}\PY{p}{)}\PY{p}{.}\PY{n+na}{minusDays}\PY{p}{(}\PY{l+m+mi}{3}\PY{p}{)}\PY{p}{.}\PY{n+na}{withDayOfMonth}\PY{p}{(}\PY{l+m+mi}{1}\PY{p}{)}\PY{p}{;}
\PY{+w}{        }\PY{n}{System}\PY{p}{.}\PY{n+na}{out}\PY{p}{.}\PY{n+na}{println}\PY{p}{(}\PY{n}{result}\PY{p}{)}\PY{p}{;}\PY{+w}{ }\PY{c+c1}{// 2026\PYZhy{}10\PYZhy{}01}
\PY{+w}{    }\PY{p}{\PYZcb{}}
\PY{p}{\PYZcb{}}
\end{Verbatim}
\end{codebox}

\subsection[TemporalAdjusters]{\code{TemporalAdjusters}}
\label{h:9:temporaladjusters}
\code{TemporalAdjusters} provides ready-made rules for common “find this special date” queries, like “next Monday” or “last day of the month.” You can also write your own adjuster as a lambda.

\begin{codebox}
\begin{Verbatim}[commandchars=\\\{\}]
\PY{k+kn}{import}\PY{+w}{ }\PY{n+nn}{java.time.DayOfWeek}\PY{p}{;}
\PY{k+kn}{import}\PY{+w}{ }\PY{n+nn}{java.time.LocalDate}\PY{p}{;}
\PY{k+kn}{import}\PY{+w}{ }\PY{n+nn}{java.time.temporal.TemporalAdjusters}\PY{p}{;}

\PY{k+kd}{public}\PY{+w}{ }\PY{k+kd}{class} \PY{n+nc}{AdjustersDemo}\PY{+w}{ }\PY{p}{\PYZob{}}
\PY{+w}{    }\PY{k+kd}{public}\PY{+w}{ }\PY{k+kd}{static}\PY{+w}{ }\PY{k+kt}{void}\PY{+w}{ }\PY{n+nf}{main}\PY{p}{(}\PY{n}{String}\PY{o}{[}\PY{o}{]}\PY{+w}{ }\PY{n}{args}\PY{p}{)}\PY{+w}{ }\PY{p}{\PYZob{}}
\PY{+w}{        }\PY{n}{LocalDate}\PY{+w}{ }\PY{n}{date}\PY{+w}{ }\PY{o}{=}\PY{+w}{ }\PY{n}{LocalDate}\PY{p}{.}\PY{n+na}{of}\PY{p}{(}\PY{l+m+mi}{2026}\PY{p}{,}\PY{+w}{ }\PY{l+m+mi}{8}\PY{p}{,}\PY{+w}{ }\PY{l+m+mi}{7}\PY{p}{)}\PY{p}{;}\PY{+w}{ }\PY{c+c1}{// a Friday}

\PY{+w}{        }\PY{n}{LocalDate}\PY{+w}{ }\PY{n}{nextMonday}\PY{+w}{ }\PY{o}{=}\PY{+w}{ }\PY{n}{date}\PY{p}{.}\PY{n+na}{with}\PY{p}{(}\PY{n}{TemporalAdjusters}\PY{p}{.}\PY{n+na}{next}\PY{p}{(}\PY{n}{DayOfWeek}\PY{p}{.}\PY{n+na}{MONDAY}\PY{p}{)}\PY{p}{)}\PY{p}{;}
\PY{+w}{        }\PY{n}{LocalDate}\PY{+w}{ }\PY{n}{lastDayOfMonth}\PY{+w}{ }\PY{o}{=}\PY{+w}{ }\PY{n}{date}\PY{p}{.}\PY{n+na}{with}\PY{p}{(}\PY{n}{TemporalAdjusters}\PY{p}{.}\PY{n+na}{lastDayOfMonth}\PY{p}{(}\PY{p}{)}\PY{p}{)}\PY{p}{;}
\PY{+w}{        }\PY{n}{LocalDate}\PY{+w}{ }\PY{n}{firstDayOfNextMonth}\PY{+w}{ }\PY{o}{=}\PY{+w}{ }\PY{n}{date}\PY{p}{.}\PY{n+na}{with}\PY{p}{(}\PY{n}{TemporalAdjusters}\PY{p}{.}\PY{n+na}{firstDayOfNextMonth}\PY{p}{(}\PY{p}{)}\PY{p}{)}\PY{p}{;}

\PY{+w}{        }\PY{n}{System}\PY{p}{.}\PY{n+na}{out}\PY{p}{.}\PY{n+na}{println}\PY{p}{(}\PY{n}{nextMonday}\PY{p}{)}\PY{p}{;}\PY{+w}{          }\PY{c+c1}{// 2026\PYZhy{}08\PYZhy{}10}
\PY{+w}{        }\PY{n}{System}\PY{p}{.}\PY{n+na}{out}\PY{p}{.}\PY{n+na}{println}\PY{p}{(}\PY{n}{lastDayOfMonth}\PY{p}{)}\PY{p}{;}\PY{+w}{       }\PY{c+c1}{// 2026\PYZhy{}08\PYZhy{}31}
\PY{+w}{        }\PY{n}{System}\PY{p}{.}\PY{n+na}{out}\PY{p}{.}\PY{n+na}{println}\PY{p}{(}\PY{n}{firstDayOfNextMonth}\PY{p}{)}\PY{p}{;}\PY{+w}{  }\PY{c+c1}{// 2026\PYZhy{}09\PYZhy{}01}

\PY{+w}{        }\PY{c+c1}{// Custom adjuster: move to the next even day of the month}
\PY{+w}{        }\PY{n}{LocalDate}\PY{+w}{ }\PY{n}{nextEvenDay}\PY{+w}{ }\PY{o}{=}\PY{+w}{ }\PY{n}{date}\PY{p}{.}\PY{n+na}{with}\PY{p}{(}\PY{n}{temporal}\PY{+w}{ }\PY{o}{\PYZhy{}}\PY{o}{\PYZgt{}}\PY{+w}{ }\PY{p}{\PYZob{}}
\PY{+w}{            }\PY{n}{LocalDate}\PY{+w}{ }\PY{n}{d}\PY{+w}{ }\PY{o}{=}\PY{+w}{ }\PY{n}{LocalDate}\PY{p}{.}\PY{n+na}{from}\PY{p}{(}\PY{n}{temporal}\PY{p}{)}\PY{p}{;}
\PY{+w}{            }\PY{k}{do}\PY{+w}{ }\PY{p}{\PYZob{}}
\PY{+w}{                }\PY{n}{d}\PY{+w}{ }\PY{o}{=}\PY{+w}{ }\PY{n}{d}\PY{p}{.}\PY{n+na}{plusDays}\PY{p}{(}\PY{l+m+mi}{1}\PY{p}{)}\PY{p}{;}
\PY{+w}{            }\PY{p}{\PYZcb{}}\PY{+w}{ }\PY{k}{while}\PY{+w}{ }\PY{p}{(}\PY{n}{d}\PY{p}{.}\PY{n+na}{getDayOfMonth}\PY{p}{(}\PY{p}{)}\PY{+w}{ }\PY{o}{\PYZpc{}}\PY{+w}{ }\PY{l+m+mi}{2}\PY{+w}{ }\PY{o}{!}\PY{o}{=}\PY{+w}{ }\PY{l+m+mi}{0}\PY{p}{)}\PY{p}{;}
\PY{+w}{            }\PY{k}{return}\PY{+w}{ }\PY{n}{d}\PY{p}{;}
\PY{+w}{        }\PY{p}{\PYZcb{}}\PY{p}{)}\PY{p}{;}
\PY{+w}{        }\PY{n}{System}\PY{p}{.}\PY{n+na}{out}\PY{p}{.}\PY{n+na}{println}\PY{p}{(}\PY{n}{nextEvenDay}\PY{p}{)}\PY{p}{;}\PY{+w}{ }\PY{c+c1}{// 2026\PYZhy{}08\PYZhy{}08}
\PY{+w}{    }\PY{p}{\PYZcb{}}
\PY{p}{\PYZcb{}}
\end{Verbatim}
\end{codebox}

\subsection[ChronoUnit.between]{\code{ChronoUnit.\allowbreak{}between}}
\label{h:9:chronounitbetween}
\code{ChronoUnit} lets you measure the difference between two temporal objects in a specific unit (days, hours, years, etc.). It is often clearer than manually computing \code{Period} or \code{Duration} when you just need a single number.

\begin{codebox}
\begin{Verbatim}[commandchars=\\\{\}]
\PY{k+kn}{import}\PY{+w}{ }\PY{n+nn}{java.time.LocalDate}\PY{p}{;}
\PY{k+kn}{import}\PY{+w}{ }\PY{n+nn}{java.time.LocalDateTime}\PY{p}{;}
\PY{k+kn}{import}\PY{+w}{ }\PY{n+nn}{java.time.temporal.ChronoUnit}\PY{p}{;}

\PY{k+kd}{public}\PY{+w}{ }\PY{k+kd}{class} \PY{n+nc}{ChronoUnitDemo}\PY{+w}{ }\PY{p}{\PYZob{}}
\PY{+w}{    }\PY{k+kd}{public}\PY{+w}{ }\PY{k+kd}{static}\PY{+w}{ }\PY{k+kt}{void}\PY{+w}{ }\PY{n+nf}{main}\PY{p}{(}\PY{n}{String}\PY{o}{[}\PY{o}{]}\PY{+w}{ }\PY{n}{args}\PY{p}{)}\PY{+w}{ }\PY{p}{\PYZob{}}
\PY{+w}{        }\PY{n}{LocalDate}\PY{+w}{ }\PY{n}{start}\PY{+w}{ }\PY{o}{=}\PY{+w}{ }\PY{n}{LocalDate}\PY{p}{.}\PY{n+na}{of}\PY{p}{(}\PY{l+m+mi}{2026}\PY{p}{,}\PY{+w}{ }\PY{l+m+mi}{1}\PY{p}{,}\PY{+w}{ }\PY{l+m+mi}{1}\PY{p}{)}\PY{p}{;}
\PY{+w}{        }\PY{n}{LocalDate}\PY{+w}{ }\PY{n}{end}\PY{+w}{ }\PY{o}{=}\PY{+w}{ }\PY{n}{LocalDate}\PY{p}{.}\PY{n+na}{of}\PY{p}{(}\PY{l+m+mi}{2026}\PY{p}{,}\PY{+w}{ }\PY{l+m+mi}{8}\PY{p}{,}\PY{+w}{ }\PY{l+m+mi}{7}\PY{p}{)}\PY{p}{;}

\PY{+w}{        }\PY{k+kt}{long}\PY{+w}{ }\PY{n}{daysBetween}\PY{+w}{ }\PY{o}{=}\PY{+w}{ }\PY{n}{ChronoUnit}\PY{p}{.}\PY{n+na}{DAYS}\PY{p}{.}\PY{n+na}{between}\PY{p}{(}\PY{n}{start}\PY{p}{,}\PY{+w}{ }\PY{n}{end}\PY{p}{)}\PY{p}{;}
\PY{+w}{        }\PY{k+kt}{long}\PY{+w}{ }\PY{n}{monthsBetween}\PY{+w}{ }\PY{o}{=}\PY{+w}{ }\PY{n}{ChronoUnit}\PY{p}{.}\PY{n+na}{MONTHS}\PY{p}{.}\PY{n+na}{between}\PY{p}{(}\PY{n}{start}\PY{p}{,}\PY{+w}{ }\PY{n}{end}\PY{p}{)}\PY{p}{;}

\PY{+w}{        }\PY{n}{System}\PY{p}{.}\PY{n+na}{out}\PY{p}{.}\PY{n+na}{println}\PY{p}{(}\PY{n}{daysBetween}\PY{p}{)}\PY{p}{;}\PY{+w}{   }\PY{c+c1}{// 218}
\PY{+w}{        }\PY{n}{System}\PY{p}{.}\PY{n+na}{out}\PY{p}{.}\PY{n+na}{println}\PY{p}{(}\PY{n}{monthsBetween}\PY{p}{)}\PY{p}{;}\PY{+w}{ }\PY{c+c1}{// 7}

\PY{+w}{        }\PY{n}{LocalDateTime}\PY{+w}{ }\PY{n}{t1}\PY{+w}{ }\PY{o}{=}\PY{+w}{ }\PY{n}{LocalDateTime}\PY{p}{.}\PY{n+na}{of}\PY{p}{(}\PY{l+m+mi}{2026}\PY{p}{,}\PY{+w}{ }\PY{l+m+mi}{8}\PY{p}{,}\PY{+w}{ }\PY{l+m+mi}{7}\PY{p}{,}\PY{+w}{ }\PY{l+m+mi}{9}\PY{p}{,}\PY{+w}{ }\PY{l+m+mi}{0}\PY{p}{)}\PY{p}{;}
\PY{+w}{        }\PY{n}{LocalDateTime}\PY{+w}{ }\PY{n}{t2}\PY{+w}{ }\PY{o}{=}\PY{+w}{ }\PY{n}{LocalDateTime}\PY{p}{.}\PY{n+na}{of}\PY{p}{(}\PY{l+m+mi}{2026}\PY{p}{,}\PY{+w}{ }\PY{l+m+mi}{8}\PY{p}{,}\PY{+w}{ }\PY{l+m+mi}{7}\PY{p}{,}\PY{+w}{ }\PY{l+m+mi}{17}\PY{p}{,}\PY{+w}{ }\PY{l+m+mi}{30}\PY{p}{)}\PY{p}{;}
\PY{+w}{        }\PY{k+kt}{long}\PY{+w}{ }\PY{n}{minutesWorked}\PY{+w}{ }\PY{o}{=}\PY{+w}{ }\PY{n}{ChronoUnit}\PY{p}{.}\PY{n+na}{MINUTES}\PY{p}{.}\PY{n+na}{between}\PY{p}{(}\PY{n}{t1}\PY{p}{,}\PY{+w}{ }\PY{n}{t2}\PY{p}{)}\PY{p}{;}
\PY{+w}{        }\PY{n}{System}\PY{p}{.}\PY{n+na}{out}\PY{p}{.}\PY{n+na}{println}\PY{p}{(}\PY{n}{minutesWorked}\PY{p}{)}\PY{p}{;}\PY{+w}{ }\PY{c+c1}{// 510}
\PY{+w}{    }\PY{p}{\PYZcb{}}
\PY{p}{\PYZcb{}}
\end{Verbatim}
\end{codebox}

\subsection[Clock and testable time]{\code{Clock} and testable time}
\label{h:9:clock-and-testable-time}
Calling \code{LocalDate.\allowbreak{}now()} or \code{Instant.\allowbreak{}now()} directly inside business logic makes that logic hard to unit test — the result changes every time you run it. \code{Clock} is an abstraction over “the current time” that you can inject and replace with a fixed clock in tests.

\begin{codebox}
\begin{Verbatim}[commandchars=\\\{\}]
\PY{k+kn}{import}\PY{+w}{ }\PY{n+nn}{java.time.Clock}\PY{p}{;}
\PY{k+kn}{import}\PY{+w}{ }\PY{n+nn}{java.time.Instant}\PY{p}{;}
\PY{k+kn}{import}\PY{+w}{ }\PY{n+nn}{java.time.LocalDate}\PY{p}{;}
\PY{k+kn}{import}\PY{+w}{ }\PY{n+nn}{java.time.ZoneOffset}\PY{p}{;}

\PY{k+kd}{public}\PY{+w}{ }\PY{k+kd}{class} \PY{n+nc}{ClockDemo}\PY{+w}{ }\PY{p}{\PYZob{}}

\PY{+w}{    }\PY{c+c1}{// Production code depends on a Clock, not on LocalDate.now() directly}
\PY{+w}{    }\PY{k+kd}{static}\PY{+w}{ }\PY{k+kd}{class} \PY{n+nc}{InvoiceService}\PY{+w}{ }\PY{p}{\PYZob{}}
\PY{+w}{        }\PY{k+kd}{private}\PY{+w}{ }\PY{k+kd}{final}\PY{+w}{ }\PY{n}{Clock}\PY{+w}{ }\PY{n}{clock}\PY{p}{;}

\PY{+w}{        }\PY{n}{InvoiceService}\PY{p}{(}\PY{n}{Clock}\PY{+w}{ }\PY{n}{clock}\PY{p}{)}\PY{+w}{ }\PY{p}{\PYZob{}}
\PY{+w}{            }\PY{k}{this}\PY{p}{.}\PY{n+na}{clock}\PY{+w}{ }\PY{o}{=}\PY{+w}{ }\PY{n}{clock}\PY{p}{;}
\PY{+w}{        }\PY{p}{\PYZcb{}}

\PY{+w}{        }\PY{n}{LocalDate}\PY{+w}{ }\PY{n+nf}{today}\PY{p}{(}\PY{p}{)}\PY{+w}{ }\PY{p}{\PYZob{}}
\PY{+w}{            }\PY{k}{return}\PY{+w}{ }\PY{n}{LocalDate}\PY{p}{.}\PY{n+na}{now}\PY{p}{(}\PY{n}{clock}\PY{p}{)}\PY{p}{;}
\PY{+w}{        }\PY{p}{\PYZcb{}}
\PY{+w}{    }\PY{p}{\PYZcb{}}

\PY{+w}{    }\PY{k+kd}{public}\PY{+w}{ }\PY{k+kd}{static}\PY{+w}{ }\PY{k+kt}{void}\PY{+w}{ }\PY{n+nf}{main}\PY{p}{(}\PY{n}{String}\PY{o}{[}\PY{o}{]}\PY{+w}{ }\PY{n}{args}\PY{p}{)}\PY{+w}{ }\PY{p}{\PYZob{}}
\PY{+w}{        }\PY{c+c1}{// Real usage: system clock}
\PY{+w}{        }\PY{n}{InvoiceService}\PY{+w}{ }\PY{n}{realService}\PY{+w}{ }\PY{o}{=}\PY{+w}{ }\PY{k}{new}\PY{+w}{ }\PY{n}{InvoiceService}\PY{p}{(}\PY{n}{Clock}\PY{p}{.}\PY{n+na}{systemUTC}\PY{p}{(}\PY{p}{)}\PY{p}{)}\PY{p}{;}
\PY{+w}{        }\PY{n}{System}\PY{p}{.}\PY{n+na}{out}\PY{p}{.}\PY{n+na}{println}\PY{p}{(}\PY{n}{realService}\PY{p}{.}\PY{n+na}{today}\PY{p}{(}\PY{p}{)}\PY{p}{)}\PY{p}{;}\PY{+w}{ }\PY{c+c1}{// e.g. 2026\PYZhy{}08\PYZhy{}07}

\PY{+w}{        }\PY{c+c1}{// Test usage: fixed clock, deterministic and repeatable}
\PY{+w}{        }\PY{n}{Clock}\PY{+w}{ }\PY{n}{fixedClock}\PY{+w}{ }\PY{o}{=}\PY{+w}{ }\PY{n}{Clock}\PY{p}{.}\PY{n+na}{fixed}\PY{p}{(}
\PY{+w}{                }\PY{n}{Instant}\PY{p}{.}\PY{n+na}{parse}\PY{p}{(}\PY{l+s}{\PYZdq{}}\PY{l+s}{2026\PYZhy{}01\PYZhy{}01T00:00:00Z}\PY{l+s}{\PYZdq{}}\PY{p}{)}\PY{p}{,}\PY{+w}{ }\PY{n}{ZoneOffset}\PY{p}{.}\PY{n+na}{UTC}\PY{p}{)}\PY{p}{;}
\PY{+w}{        }\PY{n}{InvoiceService}\PY{+w}{ }\PY{n}{testService}\PY{+w}{ }\PY{o}{=}\PY{+w}{ }\PY{k}{new}\PY{+w}{ }\PY{n}{InvoiceService}\PY{p}{(}\PY{n}{fixedClock}\PY{p}{)}\PY{p}{;}
\PY{+w}{        }\PY{n}{System}\PY{p}{.}\PY{n+na}{out}\PY{p}{.}\PY{n+na}{println}\PY{p}{(}\PY{n}{testService}\PY{p}{.}\PY{n+na}{today}\PY{p}{(}\PY{p}{)}\PY{p}{)}\PY{p}{;}\PY{+w}{ }\PY{c+c1}{// 2026\PYZhy{}01\PYZhy{}01 (always, in tests)}
\PY{+w}{    }\PY{p}{\PYZcb{}}
\PY{p}{\PYZcb{}}
\end{Verbatim}
\end{codebox}

\subsection[DST transitions: gaps and overlaps]{DST transitions: gaps and overlaps}
\label{h:9:dst-transitions-gaps-and-overlaps}
Daylight Saving Time (DST) creates two tricky situations in local time:

\begin{itemize}
\item 
\textbf{Gap}: Clocks jump forward, and some local times never occur (e.g. 02:30 does not exist the night clocks spring forward).

\item 
\textbf{Overlap}: Clocks jump backward, and some local times occur twice (e.g. 02:30 happens twice the night clocks fall back).

\end{itemize}

\code{ZonedDateTime} resolves these automatically but silently — it is important to know the behavior instead of assuming “one hour is always one hour.”

\begin{codebox}
\begin{Verbatim}[commandchars=\\\{\}]
\PY{k+kn}{import}\PY{+w}{ }\PY{n+nn}{java.time.LocalDateTime}\PY{p}{;}
\PY{k+kn}{import}\PY{+w}{ }\PY{n+nn}{java.time.ZoneId}\PY{p}{;}
\PY{k+kn}{import}\PY{+w}{ }\PY{n+nn}{java.time.ZonedDateTime}\PY{p}{;}

\PY{k+kd}{public}\PY{+w}{ }\PY{k+kd}{class} \PY{n+nc}{DstDemo}\PY{+w}{ }\PY{p}{\PYZob{}}
\PY{+w}{    }\PY{k+kd}{public}\PY{+w}{ }\PY{k+kd}{static}\PY{+w}{ }\PY{k+kt}{void}\PY{+w}{ }\PY{n+nf}{main}\PY{p}{(}\PY{n}{String}\PY{o}{[}\PY{o}{]}\PY{+w}{ }\PY{n}{args}\PY{p}{)}\PY{+w}{ }\PY{p}{\PYZob{}}
\PY{+w}{        }\PY{n}{ZoneId}\PY{+w}{ }\PY{n}{amsterdam}\PY{+w}{ }\PY{o}{=}\PY{+w}{ }\PY{n}{ZoneId}\PY{p}{.}\PY{n+na}{of}\PY{p}{(}\PY{l+s}{\PYZdq{}}\PY{l+s}{Europe/Amsterdam}\PY{l+s}{\PYZdq{}}\PY{p}{)}\PY{p}{;}

\PY{+w}{        }\PY{c+c1}{// Gap: 2026\PYZhy{}03\PYZhy{}29 02:30 does not exist in Europe/Amsterdam (spring forward at 02:00 \PYZhy{}\PYZgt{} 03:00)}
\PY{+w}{        }\PY{n}{LocalDateTime}\PY{+w}{ }\PY{n}{gapTime}\PY{+w}{ }\PY{o}{=}\PY{+w}{ }\PY{n}{LocalDateTime}\PY{p}{.}\PY{n+na}{of}\PY{p}{(}\PY{l+m+mi}{2026}\PY{p}{,}\PY{+w}{ }\PY{l+m+mi}{3}\PY{p}{,}\PY{+w}{ }\PY{l+m+mi}{29}\PY{p}{,}\PY{+w}{ }\PY{l+m+mi}{2}\PY{p}{,}\PY{+w}{ }\PY{l+m+mi}{30}\PY{p}{)}\PY{p}{;}
\PY{+w}{        }\PY{n}{ZonedDateTime}\PY{+w}{ }\PY{n}{resolvedGap}\PY{+w}{ }\PY{o}{=}\PY{+w}{ }\PY{n}{gapTime}\PY{p}{.}\PY{n+na}{atZone}\PY{p}{(}\PY{n}{amsterdam}\PY{p}{)}\PY{p}{;}
\PY{+w}{        }\PY{n}{System}\PY{p}{.}\PY{n+na}{out}\PY{p}{.}\PY{n+na}{println}\PY{p}{(}\PY{n}{resolvedGap}\PY{p}{)}\PY{p}{;}
\PY{+w}{        }\PY{c+c1}{// Output: 2026\PYZhy{}03\PYZhy{}29T03:30+02:00[Europe/Amsterdam]  \PYZlt{}\PYZhy{}\PYZhy{} shifted forward, not what you typed!}

\PY{+w}{        }\PY{c+c1}{// Overlap: 2026\PYZhy{}10\PYZhy{}25 02:30 happens twice (fall back at 03:00 \PYZhy{}\PYZgt{} 02:00)}
\PY{+w}{        }\PY{n}{LocalDateTime}\PY{+w}{ }\PY{n}{overlapTime}\PY{+w}{ }\PY{o}{=}\PY{+w}{ }\PY{n}{LocalDateTime}\PY{p}{.}\PY{n+na}{of}\PY{p}{(}\PY{l+m+mi}{2026}\PY{p}{,}\PY{+w}{ }\PY{l+m+mi}{10}\PY{p}{,}\PY{+w}{ }\PY{l+m+mi}{25}\PY{p}{,}\PY{+w}{ }\PY{l+m+mi}{2}\PY{p}{,}\PY{+w}{ }\PY{l+m+mi}{30}\PY{p}{)}\PY{p}{;}
\PY{+w}{        }\PY{n}{ZonedDateTime}\PY{+w}{ }\PY{n}{firstOccurrence}\PY{+w}{ }\PY{o}{=}\PY{+w}{ }\PY{n}{overlapTime}\PY{p}{.}\PY{n+na}{atZone}\PY{p}{(}\PY{n}{amsterdam}\PY{p}{)}\PY{p}{;}
\PY{+w}{        }\PY{n}{System}\PY{p}{.}\PY{n+na}{out}\PY{p}{.}\PY{n+na}{println}\PY{p}{(}\PY{n}{firstOccurrence}\PY{p}{)}\PY{p}{;}
\PY{+w}{        }\PY{c+c1}{// Output: 2026\PYZhy{}10\PYZhy{}25T02:30+02:00[Europe/Amsterdam] (the earlier of the two offsets, by default)}

\PY{+w}{        }\PY{c+c1}{// Adding a \PYZdq{}day\PYZdq{} across a DST boundary: still 24h wall\PYZhy{}clock days, but may be 23 or 25 actual hours.}
\PY{+w}{        }\PY{n}{ZonedDateTime}\PY{+w}{ }\PY{n}{beforeDst}\PY{+w}{ }\PY{o}{=}\PY{+w}{ }\PY{n}{ZonedDateTime}\PY{p}{.}\PY{n+na}{of}\PY{p}{(}\PY{l+m+mi}{2026}\PY{p}{,}\PY{+w}{ }\PY{l+m+mi}{3}\PY{p}{,}\PY{+w}{ }\PY{l+m+mi}{28}\PY{p}{,}\PY{+w}{ }\PY{l+m+mi}{12}\PY{p}{,}\PY{+w}{ }\PY{l+m+mi}{0}\PY{p}{,}\PY{+w}{ }\PY{l+m+mi}{0}\PY{p}{,}\PY{+w}{ }\PY{l+m+mi}{0}\PY{p}{,}\PY{+w}{ }\PY{n}{amsterdam}\PY{p}{)}\PY{p}{;}
\PY{+w}{        }\PY{n}{ZonedDateTime}\PY{+w}{ }\PY{n}{afterOneDay}\PY{+w}{ }\PY{o}{=}\PY{+w}{ }\PY{n}{beforeDst}\PY{p}{.}\PY{n+na}{plusDays}\PY{p}{(}\PY{l+m+mi}{1}\PY{p}{)}\PY{p}{;}
\PY{+w}{        }\PY{n}{System}\PY{p}{.}\PY{n+na}{out}\PY{p}{.}\PY{n+na}{println}\PY{p}{(}\PY{n}{afterOneDay}\PY{p}{)}\PY{p}{;}\PY{+w}{ }\PY{c+c1}{// 2026\PYZhy{}03\PYZhy{}29T12:00+02:00[Europe/Amsterdam] \PYZhy{} only 23 real hours passed}
\PY{+w}{    }\PY{p}{\PYZcb{}}
\PY{p}{\PYZcb{}}
\end{Verbatim}
\end{codebox}

\subsection[Comparing instants vs local times]{Comparing instants vs local times}
\label{h:9:comparing-instants-vs-local-times}
\code{Instant} represents an absolute point on the universal timeline; two instants can always be safely compared regardless of where they came from. \code{LocalDateTime} has no time zone information at all, so comparing two \code{LocalDateTime} values only tells you which “wall clock reading” is earlier — it says nothing about which one happened first in reality if they came from different zones.

\begin{codebox}
\begin{Verbatim}[commandchars=\\\{\}]
\PY{k+kn}{import}\PY{+w}{ }\PY{n+nn}{java.time.*}\PY{p}{;}

\PY{k+kd}{public}\PY{+w}{ }\PY{k+kd}{class} \PY{n+nc}{ComparisonDemo}\PY{+w}{ }\PY{p}{\PYZob{}}
\PY{+w}{    }\PY{k+kd}{public}\PY{+w}{ }\PY{k+kd}{static}\PY{+w}{ }\PY{k+kt}{void}\PY{+w}{ }\PY{n+nf}{main}\PY{p}{(}\PY{n}{String}\PY{o}{[}\PY{o}{]}\PY{+w}{ }\PY{n}{args}\PY{p}{)}\PY{+w}{ }\PY{p}{\PYZob{}}
\PY{+w}{        }\PY{n}{ZonedDateTime}\PY{+w}{ }\PY{n}{tokyoMeeting}\PY{+w}{ }\PY{o}{=}
\PY{+w}{                }\PY{n}{ZonedDateTime}\PY{p}{.}\PY{n+na}{of}\PY{p}{(}\PY{l+m+mi}{2026}\PY{p}{,}\PY{+w}{ }\PY{l+m+mi}{8}\PY{p}{,}\PY{+w}{ }\PY{l+m+mi}{7}\PY{p}{,}\PY{+w}{ }\PY{l+m+mi}{22}\PY{p}{,}\PY{+w}{ }\PY{l+m+mi}{0}\PY{p}{,}\PY{+w}{ }\PY{l+m+mi}{0}\PY{p}{,}\PY{+w}{ }\PY{l+m+mi}{0}\PY{p}{,}\PY{+w}{ }\PY{n}{ZoneId}\PY{p}{.}\PY{n+na}{of}\PY{p}{(}\PY{l+s}{\PYZdq{}}\PY{l+s}{Asia/Tokyo}\PY{l+s}{\PYZdq{}}\PY{p}{)}\PY{p}{)}\PY{p}{;}
\PY{+w}{        }\PY{n}{ZonedDateTime}\PY{+w}{ }\PY{n}{amsterdamMeeting}\PY{+w}{ }\PY{o}{=}
\PY{+w}{                }\PY{n}{ZonedDateTime}\PY{p}{.}\PY{n+na}{of}\PY{p}{(}\PY{l+m+mi}{2026}\PY{p}{,}\PY{+w}{ }\PY{l+m+mi}{8}\PY{p}{,}\PY{+w}{ }\PY{l+m+mi}{7}\PY{p}{,}\PY{+w}{ }\PY{l+m+mi}{16}\PY{p}{,}\PY{+w}{ }\PY{l+m+mi}{0}\PY{p}{,}\PY{+w}{ }\PY{l+m+mi}{0}\PY{p}{,}\PY{+w}{ }\PY{l+m+mi}{0}\PY{p}{,}\PY{+w}{ }\PY{n}{ZoneId}\PY{p}{.}\PY{n+na}{of}\PY{p}{(}\PY{l+s}{\PYZdq{}}\PY{l+s}{Europe/Amsterdam}\PY{l+s}{\PYZdq{}}\PY{p}{)}\PY{p}{)}\PY{p}{;}

\PY{+w}{        }\PY{c+c1}{// Wrong\PYZhy{}ish: comparing LocalDateTime ignores the zone entirely}
\PY{+w}{        }\PY{k+kt}{boolean}\PY{+w}{ }\PY{n}{localLooksLater}\PY{+w}{ }\PY{o}{=}\PY{+w}{ }\PY{n}{tokyoMeeting}\PY{p}{.}\PY{n+na}{toLocalDateTime}\PY{p}{(}\PY{p}{)}
\PY{+w}{                }\PY{p}{.}\PY{n+na}{isAfter}\PY{p}{(}\PY{n}{amsterdamMeeting}\PY{p}{.}\PY{n+na}{toLocalDateTime}\PY{p}{(}\PY{p}{)}\PY{p}{)}\PY{p}{;}
\PY{+w}{        }\PY{n}{System}\PY{p}{.}\PY{n+na}{out}\PY{p}{.}\PY{n+na}{println}\PY{p}{(}\PY{n}{localLooksLater}\PY{p}{)}\PY{p}{;}\PY{+w}{ }\PY{c+c1}{// true (22:00 \PYZgt{} 16:00) \PYZhy{} misleading!}

\PY{+w}{        }\PY{c+c1}{// Correct: compare the actual instants}
\PY{+w}{        }\PY{k+kt}{boolean}\PY{+w}{ }\PY{n}{instantIsLater}\PY{+w}{ }\PY{o}{=}\PY{+w}{ }\PY{n}{tokyoMeeting}\PY{p}{.}\PY{n+na}{toInstant}\PY{p}{(}\PY{p}{)}\PY{p}{.}\PY{n+na}{isAfter}\PY{p}{(}\PY{n}{amsterdamMeeting}\PY{p}{.}\PY{n+na}{toInstant}\PY{p}{(}\PY{p}{)}\PY{p}{)}\PY{p}{;}
\PY{+w}{        }\PY{n}{System}\PY{p}{.}\PY{n+na}{out}\PY{p}{.}\PY{n+na}{println}\PY{p}{(}\PY{n}{instantIsLater}\PY{p}{)}\PY{p}{;}\PY{+w}{ }\PY{c+c1}{// false \PYZhy{} Tokyo 22:00 JST happens BEFORE Amsterdam 16:00 CEST}

\PY{+w}{        }\PY{n}{System}\PY{p}{.}\PY{n+na}{out}\PY{p}{.}\PY{n+na}{println}\PY{p}{(}\PY{n}{tokyoMeeting}\PY{p}{.}\PY{n+na}{toInstant}\PY{p}{(}\PY{p}{)}\PY{p}{)}\PY{p}{;}\PY{+w}{     }\PY{c+c1}{// 2026\PYZhy{}08\PYZhy{}07T13:00:00Z}
\PY{+w}{        }\PY{n}{System}\PY{p}{.}\PY{n+na}{out}\PY{p}{.}\PY{n+na}{println}\PY{p}{(}\PY{n}{amsterdamMeeting}\PY{p}{.}\PY{n+na}{toInstant}\PY{p}{(}\PY{p}{)}\PY{p}{)}\PY{p}{;}\PY{+w}{ }\PY{c+c1}{// 2026\PYZhy{}08\PYZhy{}07T14:00:00Z}
\PY{+w}{    }\PY{p}{\PYZcb{}}
\PY{p}{\PYZcb{}}
\end{Verbatim}
\end{codebox}

\subsection[Storing time in a database]{Storing time in a database}
\label{h:9:storing-time-in-a-database}
The safest rule: \textbf{store instants in UTC, and store the original time zone separately if you need to display local wall-clock time later.} A raw \code{LocalDateTime} alone loses the information needed to know “when” that really was in absolute terms; a \code{ZonedDateTime}’s zone rules can even change in the future (governments do change DST rules).

{\tablefont
\begin{xltabular}{\linewidth}{@{}L{0.994}L{1.006}@{}}
\toprule
\rowcolor{tablehdr}
\thd{What you need to store} & \thd{Recommended column type(s)} \\
\midrule
\endhead
“This event happened at an exact moment” & \code{TIMESTAMP~\allowbreak{}WITH~\allowbreak{}TIME~\allowbreak{}ZONE} (stored internally as UTC), mapped to \code{Instant} or \code{OffsetDateTime} \\
“This event happened at an exact moment, and I must show the original local time later” & UTC instant column + a separate \code{zone\_\allowbreak{}id} (String) column \\
“A calendar date with no time meaning” (birthday, deadline date) & \code{DATE}, mapped to \code{LocalDate} \\
“A recurring local time with no date” (store opens at 9am, always local) & \code{TIME}, mapped to \code{LocalTime} \\
\bottomrule
\end{xltabular}
}

\begin{codebox}
\begin{Verbatim}[commandchars=\\\{\}]
\PY{k+kn}{import}\PY{+w}{ }\PY{n+nn}{java.time.Instant}\PY{p}{;}
\PY{k+kn}{import}\PY{+w}{ }\PY{n+nn}{java.time.ZoneId}\PY{p}{;}
\PY{k+kn}{import}\PY{+w}{ }\PY{n+nn}{java.time.ZonedDateTime}\PY{p}{;}

\PY{k+kd}{public}\PY{+w}{ }\PY{k+kd}{class} \PY{n+nc}{DbStorageExample}\PY{+w}{ }\PY{p}{\PYZob{}}

\PY{+w}{    }\PY{c+c1}{// What you\PYZsq{}d persist for an appointment booked by a user in a given zone}
\PY{+w}{    }\PY{k+kd}{record} \PY{n+nc}{AppointmentRecord}\PY{p}{(}\PY{n}{Instant}\PY{+w}{ }\PY{n}{utcInstant}\PY{p}{,}\PY{+w}{ }\PY{n}{String}\PY{+w}{ }\PY{n}{zoneId}\PY{p}{)}\PY{+w}{ }\PY{p}{\PYZob{}}

\PY{+w}{        }\PY{n}{ZonedDateTime}\PY{+w}{ }\PY{n+nf}{toLocalDisplayTime}\PY{p}{(}\PY{p}{)}\PY{+w}{ }\PY{p}{\PYZob{}}
\PY{+w}{            }\PY{k}{return}\PY{+w}{ }\PY{n}{utcInstant}\PY{p}{.}\PY{n+na}{atZone}\PY{p}{(}\PY{n}{ZoneId}\PY{p}{.}\PY{n+na}{of}\PY{p}{(}\PY{n}{zoneId}\PY{p}{)}\PY{p}{)}\PY{p}{;}
\PY{+w}{        }\PY{p}{\PYZcb{}}
\PY{+w}{    }\PY{p}{\PYZcb{}}

\PY{+w}{    }\PY{k+kd}{public}\PY{+w}{ }\PY{k+kd}{static}\PY{+w}{ }\PY{k+kt}{void}\PY{+w}{ }\PY{n+nf}{main}\PY{p}{(}\PY{n}{String}\PY{o}{[}\PY{o}{]}\PY{+w}{ }\PY{n}{args}\PY{p}{)}\PY{+w}{ }\PY{p}{\PYZob{}}
\PY{+w}{        }\PY{n}{ZonedDateTime}\PY{+w}{ }\PY{n}{bookedAt}\PY{+w}{ }\PY{o}{=}
\PY{+w}{                }\PY{n}{ZonedDateTime}\PY{p}{.}\PY{n+na}{of}\PY{p}{(}\PY{l+m+mi}{2026}\PY{p}{,}\PY{+w}{ }\PY{l+m+mi}{8}\PY{p}{,}\PY{+w}{ }\PY{l+m+mi}{7}\PY{p}{,}\PY{+w}{ }\PY{l+m+mi}{15}\PY{p}{,}\PY{+w}{ }\PY{l+m+mi}{0}\PY{p}{,}\PY{+w}{ }\PY{l+m+mi}{0}\PY{p}{,}\PY{+w}{ }\PY{l+m+mi}{0}\PY{p}{,}\PY{+w}{ }\PY{n}{ZoneId}\PY{p}{.}\PY{n+na}{of}\PY{p}{(}\PY{l+s}{\PYZdq{}}\PY{l+s}{America/New\PYZus{}York}\PY{l+s}{\PYZdq{}}\PY{p}{)}\PY{p}{)}\PY{p}{;}

\PY{+w}{        }\PY{n}{AppointmentRecord}\PY{+w}{ }\PY{n}{record}\PY{+w}{ }\PY{o}{=}\PY{+w}{ }\PY{k}{new}\PY{+w}{ }\PY{n}{AppointmentRecord}\PY{p}{(}\PY{n}{bookedAt}\PY{p}{.}\PY{n+na}{toInstant}\PY{p}{(}\PY{p}{)}\PY{p}{,}\PY{+w}{ }\PY{l+s}{\PYZdq{}}\PY{l+s}{America/New\PYZus{}York}\PY{l+s}{\PYZdq{}}\PY{p}{)}\PY{p}{;}

\PY{+w}{        }\PY{n}{System}\PY{p}{.}\PY{n+na}{out}\PY{p}{.}\PY{n+na}{println}\PY{p}{(}\PY{n}{record}\PY{p}{.}\PY{n+na}{utcInstant}\PY{p}{(}\PY{p}{)}\PY{p}{)}\PY{p}{;}\PY{+w}{ }\PY{c+c1}{// 2026\PYZhy{}08\PYZhy{}07T19:00:00Z}
\PY{+w}{        }\PY{n}{System}\PY{p}{.}\PY{n+na}{out}\PY{p}{.}\PY{n+na}{println}\PY{p}{(}\PY{n}{record}\PY{p}{.}\PY{n+na}{toLocalDisplayTime}\PY{p}{(}\PY{p}{)}\PY{p}{)}\PY{p}{;}
\PY{+w}{        }\PY{c+c1}{// Output: 2026\PYZhy{}08\PYZhy{}07T15:00\PYZhy{}04:00[America/New\PYZus{}York]}
\PY{+w}{    }\PY{p}{\PYZcb{}}
\PY{p}{\PYZcb{}}
\end{Verbatim}
\end{codebox}

\section[Formatting and Parsing]{Formatting and Parsing}
\label{h:9:formatting-and-parsing}
\subsection[DateTimeFormatter predefined formats]{\code{DateTimeFormatter} predefined formats}
\label{h:9:datetimeformatter-predefined-formats}
\code{DateTimeFormatter} ships with several constants for standard formats, most notably the ISO family. These are good defaults for machine-to-machine communication (logs, APIs, file names).

\begin{codebox}
\begin{Verbatim}[commandchars=\\\{\}]
\PY{k+kn}{import}\PY{+w}{ }\PY{n+nn}{java.time.LocalDate}\PY{p}{;}
\PY{k+kn}{import}\PY{+w}{ }\PY{n+nn}{java.time.LocalDateTime}\PY{p}{;}
\PY{k+kn}{import}\PY{+w}{ }\PY{n+nn}{java.time.format.DateTimeFormatter}\PY{p}{;}

\PY{k+kd}{public}\PY{+w}{ }\PY{k+kd}{class} \PY{n+nc}{PredefinedFormatsDemo}\PY{+w}{ }\PY{p}{\PYZob{}}
\PY{+w}{    }\PY{k+kd}{public}\PY{+w}{ }\PY{k+kd}{static}\PY{+w}{ }\PY{k+kt}{void}\PY{+w}{ }\PY{n+nf}{main}\PY{p}{(}\PY{n}{String}\PY{o}{[}\PY{o}{]}\PY{+w}{ }\PY{n}{args}\PY{p}{)}\PY{+w}{ }\PY{p}{\PYZob{}}
\PY{+w}{        }\PY{n}{LocalDate}\PY{+w}{ }\PY{n}{date}\PY{+w}{ }\PY{o}{=}\PY{+w}{ }\PY{n}{LocalDate}\PY{p}{.}\PY{n+na}{of}\PY{p}{(}\PY{l+m+mi}{2026}\PY{p}{,}\PY{+w}{ }\PY{l+m+mi}{8}\PY{p}{,}\PY{+w}{ }\PY{l+m+mi}{7}\PY{p}{)}\PY{p}{;}
\PY{+w}{        }\PY{n}{LocalDateTime}\PY{+w}{ }\PY{n}{dateTime}\PY{+w}{ }\PY{o}{=}\PY{+w}{ }\PY{n}{LocalDateTime}\PY{p}{.}\PY{n+na}{of}\PY{p}{(}\PY{l+m+mi}{2026}\PY{p}{,}\PY{+w}{ }\PY{l+m+mi}{8}\PY{p}{,}\PY{+w}{ }\PY{l+m+mi}{7}\PY{p}{,}\PY{+w}{ }\PY{l+m+mi}{14}\PY{p}{,}\PY{+w}{ }\PY{l+m+mi}{30}\PY{p}{,}\PY{+w}{ }\PY{l+m+mi}{15}\PY{p}{)}\PY{p}{;}

\PY{+w}{        }\PY{n}{System}\PY{p}{.}\PY{n+na}{out}\PY{p}{.}\PY{n+na}{println}\PY{p}{(}\PY{n}{date}\PY{p}{.}\PY{n+na}{format}\PY{p}{(}\PY{n}{DateTimeFormatter}\PY{p}{.}\PY{n+na}{ISO\PYZus{}LOCAL\PYZus{}DATE}\PY{p}{)}\PY{p}{)}\PY{p}{;}
\PY{+w}{        }\PY{c+c1}{// Output: 2026\PYZhy{}08\PYZhy{}07}

\PY{+w}{        }\PY{n}{System}\PY{p}{.}\PY{n+na}{out}\PY{p}{.}\PY{n+na}{println}\PY{p}{(}\PY{n}{dateTime}\PY{p}{.}\PY{n+na}{format}\PY{p}{(}\PY{n}{DateTimeFormatter}\PY{p}{.}\PY{n+na}{ISO\PYZus{}LOCAL\PYZus{}DATE\PYZus{}TIME}\PY{p}{)}\PY{p}{)}\PY{p}{;}
\PY{+w}{        }\PY{c+c1}{// Output: 2026\PYZhy{}08\PYZhy{}07T14:30:15}

\PY{+w}{        }\PY{n}{System}\PY{p}{.}\PY{n+na}{out}\PY{p}{.}\PY{n+na}{println}\PY{p}{(}\PY{n}{dateTime}\PY{p}{.}\PY{n+na}{format}\PY{p}{(}\PY{n}{DateTimeFormatter}\PY{p}{.}\PY{n+na}{ISO\PYZus{}DATE\PYZus{}TIME}\PY{p}{)}\PY{p}{)}\PY{p}{;}
\PY{+w}{        }\PY{c+c1}{// Output: 2026\PYZhy{}08\PYZhy{}07T14:30:15 (no offset here because LocalDateTime has no zone)}

\PY{+w}{        }\PY{n}{System}\PY{p}{.}\PY{n+na}{out}\PY{p}{.}\PY{n+na}{println}\PY{p}{(}\PY{n}{dateTime}\PY{p}{.}\PY{n+na}{format}\PY{p}{(}\PY{n}{DateTimeFormatter}\PY{p}{.}\PY{n+na}{BASIC\PYZus{}ISO\PYZus{}DATE}\PY{p}{)}\PY{p}{)}\PY{p}{;}
\PY{+w}{        }\PY{c+c1}{// Output: 20260807}
\PY{+w}{    }\PY{p}{\PYZcb{}}
\PY{p}{\PYZcb{}}
\end{Verbatim}
\end{codebox}

\subsection[ofPattern and custom patterns]{\code{ofPattern} and custom patterns}
\label{h:9:ofpattern-and-custom-patterns}
\code{Date\allowbreak{}Time\allowbreak{}Formatter.\allowbreak{}of\allowbreak{}Pattern} lets you define your own layout using pattern letters (borrowed conceptually from \code{SimpleDateFormat}, but not identical — see the \code{YYYY} pitfall below).

\begin{codebox}
\begin{Verbatim}[commandchars=\\\{\}]
\PY{k+kn}{import}\PY{+w}{ }\PY{n+nn}{java.time.LocalDateTime}\PY{p}{;}
\PY{k+kn}{import}\PY{+w}{ }\PY{n+nn}{java.time.format.DateTimeFormatter}\PY{p}{;}
\PY{k+kn}{import}\PY{+w}{ }\PY{n+nn}{java.util.Locale}\PY{p}{;}

\PY{k+kd}{public}\PY{+w}{ }\PY{k+kd}{class} \PY{n+nc}{OfPatternDemo}\PY{+w}{ }\PY{p}{\PYZob{}}
\PY{+w}{    }\PY{k+kd}{public}\PY{+w}{ }\PY{k+kd}{static}\PY{+w}{ }\PY{k+kt}{void}\PY{+w}{ }\PY{n+nf}{main}\PY{p}{(}\PY{n}{String}\PY{o}{[}\PY{o}{]}\PY{+w}{ }\PY{n}{args}\PY{p}{)}\PY{+w}{ }\PY{p}{\PYZob{}}
\PY{+w}{        }\PY{n}{LocalDateTime}\PY{+w}{ }\PY{n}{dt}\PY{+w}{ }\PY{o}{=}\PY{+w}{ }\PY{n}{LocalDateTime}\PY{p}{.}\PY{n+na}{of}\PY{p}{(}\PY{l+m+mi}{2026}\PY{p}{,}\PY{+w}{ }\PY{l+m+mi}{8}\PY{p}{,}\PY{+w}{ }\PY{l+m+mi}{7}\PY{p}{,}\PY{+w}{ }\PY{l+m+mi}{14}\PY{p}{,}\PY{+w}{ }\PY{l+m+mi}{30}\PY{p}{,}\PY{+w}{ }\PY{l+m+mi}{15}\PY{p}{)}\PY{p}{;}

\PY{+w}{        }\PY{n}{DateTimeFormatter}\PY{+w}{ }\PY{n}{fmt1}\PY{+w}{ }\PY{o}{=}\PY{+w}{ }\PY{n}{DateTimeFormatter}\PY{p}{.}\PY{n+na}{ofPattern}\PY{p}{(}\PY{l+s}{\PYZdq{}}\PY{l+s}{dd/MM/yyyy HH:mm:ss}\PY{l+s}{\PYZdq{}}\PY{p}{)}\PY{p}{;}
\PY{+w}{        }\PY{n}{DateTimeFormatter}\PY{+w}{ }\PY{n}{fmt2}\PY{+w}{ }\PY{o}{=}\PY{+w}{ }\PY{n}{DateTimeFormatter}\PY{p}{.}\PY{n+na}{ofPattern}\PY{p}{(}\PY{l+s}{\PYZdq{}}\PY{l+s}{EEEE, MMMM d, yyyy}\PY{l+s}{\PYZdq{}}\PY{p}{,}\PY{+w}{ }\PY{n}{Locale}\PY{p}{.}\PY{n+na}{ENGLISH}\PY{p}{)}\PY{p}{;}

\PY{+w}{        }\PY{n}{System}\PY{p}{.}\PY{n+na}{out}\PY{p}{.}\PY{n+na}{println}\PY{p}{(}\PY{n}{dt}\PY{p}{.}\PY{n+na}{format}\PY{p}{(}\PY{n}{fmt1}\PY{p}{)}\PY{p}{)}\PY{p}{;}\PY{+w}{ }\PY{c+c1}{// 07/08/2026 14:30:15}
\PY{+w}{        }\PY{n}{System}\PY{p}{.}\PY{n+na}{out}\PY{p}{.}\PY{n+na}{println}\PY{p}{(}\PY{n}{dt}\PY{p}{.}\PY{n+na}{format}\PY{p}{(}\PY{n}{fmt2}\PY{p}{)}\PY{p}{)}\PY{p}{;}\PY{+w}{ }\PY{c+c1}{// Friday, August 7, 2026}
\PY{+w}{    }\PY{p}{\PYZcb{}}
\PY{p}{\PYZcb{}}
\end{Verbatim}
\end{codebox}

{\tablefont
\begin{xltabular}{\linewidth}{@{}L{0.462}L{1.731}L{0.808}@{}}
\toprule
\rowcolor{tablehdr}
\thd{Pattern letter} & \thd{Meaning} & \thd{Example} \\
\midrule
\endhead
\code{y} & Year of era & \code{2026} \\
\code{M} & Month of year & \code{08} or \code{Aug} (with \code{MMM}) \\
\code{d} & Day of month & \code{07} \\
\code{E} & Day of week & \code{Fri} or \code{Friday} (with \code{EEEE}) \\
\code{H} & Hour (0-23) & \code{14} \\
\code{h} & Hour (1-12), needs \code{a} for AM/PM & \code{02~\allowbreak{}PM} \\
\code{m} & Minute & \code{30} \\
\code{s} & Second & \code{15} \\
\code{S} & Fraction of second & \code{123} \\
\code{z} & Time zone name & \code{CEST} \\
\code{Z} & Zone offset & \code{+\allowbreak{}0200} \\
\code{X} & ISO-8601 offset & \code{+\allowbreak{}02:\allowbreak{}00} \\
\bottomrule
\end{xltabular}
}

\subsection[ofLocalizedDate and locale-aware styles]{\code{ofLocalizedDate} and locale-aware styles}
\label{h:9:oflocalizeddate-and-locale-aware-styles}
Instead of hand-rolling a pattern, \code{ofLocalizedDate}, \code{ofLocalizedTime}, and \code{ofLocalizedDateTime} produce formats that automatically match the conventions of a given \code{Locale} (order of day/month/year, separators, month names).

\begin{codebox}
\begin{Verbatim}[commandchars=\\\{\}]
\PY{k+kn}{import}\PY{+w}{ }\PY{n+nn}{java.time.LocalDate}\PY{p}{;}
\PY{k+kn}{import}\PY{+w}{ }\PY{n+nn}{java.time.format.DateTimeFormatter}\PY{p}{;}
\PY{k+kn}{import}\PY{+w}{ }\PY{n+nn}{java.time.format.FormatStyle}\PY{p}{;}
\PY{k+kn}{import}\PY{+w}{ }\PY{n+nn}{java.util.Locale}\PY{p}{;}

\PY{k+kd}{public}\PY{+w}{ }\PY{k+kd}{class} \PY{n+nc}{LocalizedFormatDemo}\PY{+w}{ }\PY{p}{\PYZob{}}
\PY{+w}{    }\PY{k+kd}{public}\PY{+w}{ }\PY{k+kd}{static}\PY{+w}{ }\PY{k+kt}{void}\PY{+w}{ }\PY{n+nf}{main}\PY{p}{(}\PY{n}{String}\PY{o}{[}\PY{o}{]}\PY{+w}{ }\PY{n}{args}\PY{p}{)}\PY{+w}{ }\PY{p}{\PYZob{}}
\PY{+w}{        }\PY{n}{LocalDate}\PY{+w}{ }\PY{n}{date}\PY{+w}{ }\PY{o}{=}\PY{+w}{ }\PY{n}{LocalDate}\PY{p}{.}\PY{n+na}{of}\PY{p}{(}\PY{l+m+mi}{2026}\PY{p}{,}\PY{+w}{ }\PY{l+m+mi}{8}\PY{p}{,}\PY{+w}{ }\PY{l+m+mi}{7}\PY{p}{)}\PY{p}{;}

\PY{+w}{        }\PY{n}{DateTimeFormatter}\PY{+w}{ }\PY{n}{us}\PY{+w}{ }\PY{o}{=}\PY{+w}{ }\PY{n}{DateTimeFormatter}\PY{p}{.}\PY{n+na}{ofLocalizedDate}\PY{p}{(}\PY{n}{FormatStyle}\PY{p}{.}\PY{n+na}{LONG}\PY{p}{)}\PY{p}{.}\PY{n+na}{withLocale}\PY{p}{(}\PY{n}{Locale}\PY{p}{.}\PY{n+na}{US}\PY{p}{)}\PY{p}{;}
\PY{+w}{        }\PY{n}{DateTimeFormatter}\PY{+w}{ }\PY{n}{de}\PY{+w}{ }\PY{o}{=}\PY{+w}{ }\PY{n}{DateTimeFormatter}\PY{p}{.}\PY{n+na}{ofLocalizedDate}\PY{p}{(}\PY{n}{FormatStyle}\PY{p}{.}\PY{n+na}{LONG}\PY{p}{)}\PY{p}{.}\PY{n+na}{withLocale}\PY{p}{(}\PY{n}{Locale}\PY{p}{.}\PY{n+na}{GERMANY}\PY{p}{)}\PY{p}{;}
\PY{+w}{        }\PY{n}{DateTimeFormatter}\PY{+w}{ }\PY{n}{jp}\PY{+w}{ }\PY{o}{=}\PY{+w}{ }\PY{n}{DateTimeFormatter}\PY{p}{.}\PY{n+na}{ofLocalizedDate}\PY{p}{(}\PY{n}{FormatStyle}\PY{p}{.}\PY{n+na}{LONG}\PY{p}{)}\PY{p}{.}\PY{n+na}{withLocale}\PY{p}{(}\PY{n}{Locale}\PY{p}{.}\PY{n+na}{JAPAN}\PY{p}{)}\PY{p}{;}

\PY{+w}{        }\PY{n}{System}\PY{p}{.}\PY{n+na}{out}\PY{p}{.}\PY{n+na}{println}\PY{p}{(}\PY{n}{date}\PY{p}{.}\PY{n+na}{format}\PY{p}{(}\PY{n}{us}\PY{p}{)}\PY{p}{)}\PY{p}{;}\PY{+w}{ }\PY{c+c1}{// August 7, 2026}
\PY{+w}{        }\PY{n}{System}\PY{p}{.}\PY{n+na}{out}\PY{p}{.}\PY{n+na}{println}\PY{p}{(}\PY{n}{date}\PY{p}{.}\PY{n+na}{format}\PY{p}{(}\PY{n}{de}\PY{p}{)}\PY{p}{)}\PY{p}{;}\PY{+w}{ }\PY{c+c1}{// 7. August 2026}
\PY{+w}{        }\PY{n}{System}\PY{p}{.}\PY{n+na}{out}\PY{p}{.}\PY{n+na}{println}\PY{p}{(}\PY{n}{date}\PY{p}{.}\PY{n+na}{format}\PY{p}{(}\PY{n}{jp}\PY{p}{)}\PY{p}{)}\PY{p}{;}\PY{+w}{ }\PY{c+c1}{// 2026年8月7日}
\PY{+w}{    }\PY{p}{\PYZcb{}}
\PY{p}{\PYZcb{}}
\end{Verbatim}
\end{codebox}

\subsection[DateTimeFormatterBuilder]{\code{Date\allowbreak{}Time\allowbreak{}Formatter\allowbreak{}Builder}}
\label{h:9:datetimeformatterbuilder}
For complex or conditional formats — optional sections, case-insensitive parsing, custom text lookups — \code{Date\allowbreak{}Time\allowbreak{}Formatter\allowbreak{}Builder} gives fine-grained control beyond a simple pattern string.

\begin{codebox}
\begin{Verbatim}[commandchars=\\\{\}]
\PY{k+kn}{import}\PY{+w}{ }\PY{n+nn}{java.time.LocalDate}\PY{p}{;}
\PY{k+kn}{import}\PY{+w}{ }\PY{n+nn}{java.time.format.DateTimeFormatter}\PY{p}{;}
\PY{k+kn}{import}\PY{+w}{ }\PY{n+nn}{java.time.format.DateTimeFormatterBuilder}\PY{p}{;}
\PY{k+kn}{import}\PY{+w}{ }\PY{n+nn}{java.time.temporal.ChronoField}\PY{p}{;}
\PY{k+kn}{import}\PY{+w}{ }\PY{n+nn}{java.util.Locale}\PY{p}{;}

\PY{k+kd}{public}\PY{+w}{ }\PY{k+kd}{class} \PY{n+nc}{FormatterBuilderDemo}\PY{+w}{ }\PY{p}{\PYZob{}}
\PY{+w}{    }\PY{k+kd}{public}\PY{+w}{ }\PY{k+kd}{static}\PY{+w}{ }\PY{k+kt}{void}\PY{+w}{ }\PY{n+nf}{main}\PY{p}{(}\PY{n}{String}\PY{o}{[}\PY{o}{]}\PY{+w}{ }\PY{n}{args}\PY{p}{)}\PY{+w}{ }\PY{p}{\PYZob{}}
\PY{+w}{        }\PY{n}{DateTimeFormatter}\PY{+w}{ }\PY{n}{formatter}\PY{+w}{ }\PY{o}{=}\PY{+w}{ }\PY{k}{new}\PY{+w}{ }\PY{n}{DateTimeFormatterBuilder}\PY{p}{(}\PY{p}{)}
\PY{+w}{                }\PY{p}{.}\PY{n+na}{appendValue}\PY{p}{(}\PY{n}{ChronoField}\PY{p}{.}\PY{n+na}{YEAR}\PY{p}{,}\PY{+w}{ }\PY{l+m+mi}{4}\PY{p}{)}
\PY{+w}{                }\PY{p}{.}\PY{n+na}{appendLiteral}\PY{p}{(}\PY{l+s+sc}{\PYZsq{}\PYZhy{}\PYZsq{}}\PY{p}{)}
\PY{+w}{                }\PY{p}{.}\PY{n+na}{appendValue}\PY{p}{(}\PY{n}{ChronoField}\PY{p}{.}\PY{n+na}{MONTH\PYZus{}OF\PYZus{}YEAR}\PY{p}{,}\PY{+w}{ }\PY{l+m+mi}{2}\PY{p}{)}
\PY{+w}{                }\PY{p}{.}\PY{n+na}{appendLiteral}\PY{p}{(}\PY{l+s+sc}{\PYZsq{}\PYZhy{}\PYZsq{}}\PY{p}{)}
\PY{+w}{                }\PY{p}{.}\PY{n+na}{appendValue}\PY{p}{(}\PY{n}{ChronoField}\PY{p}{.}\PY{n+na}{DAY\PYZus{}OF\PYZus{}MONTH}\PY{p}{,}\PY{+w}{ }\PY{l+m+mi}{2}\PY{p}{)}
\PY{+w}{                }\PY{p}{.}\PY{n+na}{optionalStart}\PY{p}{(}\PY{p}{)}
\PY{+w}{                }\PY{p}{.}\PY{n+na}{appendLiteral}\PY{p}{(}\PY{l+s}{\PYZdq{}}\PY{l+s}{ (}\PY{l+s}{\PYZdq{}}\PY{p}{)}
\PY{+w}{                }\PY{p}{.}\PY{n+na}{appendText}\PY{p}{(}\PY{n}{ChronoField}\PY{p}{.}\PY{n+na}{DAY\PYZus{}OF\PYZus{}WEEK}\PY{p}{)}
\PY{+w}{                }\PY{p}{.}\PY{n+na}{appendLiteral}\PY{p}{(}\PY{l+s+sc}{\PYZsq{})\PYZsq{}}\PY{p}{)}
\PY{+w}{                }\PY{p}{.}\PY{n+na}{optionalEnd}\PY{p}{(}\PY{p}{)}
\PY{+w}{                }\PY{p}{.}\PY{n+na}{toFormatter}\PY{p}{(}\PY{n}{Locale}\PY{p}{.}\PY{n+na}{ENGLISH}\PY{p}{)}\PY{p}{;}

\PY{+w}{        }\PY{n}{LocalDate}\PY{+w}{ }\PY{n}{date}\PY{+w}{ }\PY{o}{=}\PY{+w}{ }\PY{n}{LocalDate}\PY{p}{.}\PY{n+na}{of}\PY{p}{(}\PY{l+m+mi}{2026}\PY{p}{,}\PY{+w}{ }\PY{l+m+mi}{8}\PY{p}{,}\PY{+w}{ }\PY{l+m+mi}{7}\PY{p}{)}\PY{p}{;}
\PY{+w}{        }\PY{n}{System}\PY{p}{.}\PY{n+na}{out}\PY{p}{.}\PY{n+na}{println}\PY{p}{(}\PY{n}{date}\PY{p}{.}\PY{n+na}{format}\PY{p}{(}\PY{n}{formatter}\PY{p}{)}\PY{p}{)}\PY{p}{;}
\PY{+w}{        }\PY{c+c1}{// Output: 2026\PYZhy{}08\PYZhy{}07 (Friday)}
\PY{+w}{    }\PY{p}{\PYZcb{}}
\PY{p}{\PYZcb{}}
\end{Verbatim}
\end{codebox}

\subsection[Strict vs lenient ResolverStyle]{Strict vs lenient \code{ResolverStyle}}
\label{h:9:strict-vs-lenient-resolverstyle}
\code{ResolverStyle} controls how aggressively the parser accepts out-of-range or ambiguous field values. \code{STRICT} rejects anything invalid, \code{SMART} (the default) applies sensible corrections, and \code{LENIENT} allows arithmetic overflow into adjacent fields.

\begin{codebox}
\begin{Verbatim}[commandchars=\\\{\}]
\PY{k+kn}{import}\PY{+w}{ }\PY{n+nn}{java.time.LocalDate}\PY{p}{;}
\PY{k+kn}{import}\PY{+w}{ }\PY{n+nn}{java.time.format.DateTimeFormatter}\PY{p}{;}
\PY{k+kn}{import}\PY{+w}{ }\PY{n+nn}{java.time.format.ResolverStyle}\PY{p}{;}

\PY{k+kd}{public}\PY{+w}{ }\PY{k+kd}{class} \PY{n+nc}{ResolverStyleDemo}\PY{+w}{ }\PY{p}{\PYZob{}}
\PY{+w}{    }\PY{k+kd}{public}\PY{+w}{ }\PY{k+kd}{static}\PY{+w}{ }\PY{k+kt}{void}\PY{+w}{ }\PY{n+nf}{main}\PY{p}{(}\PY{n}{String}\PY{o}{[}\PY{o}{]}\PY{+w}{ }\PY{n}{args}\PY{p}{)}\PY{+w}{ }\PY{p}{\PYZob{}}
\PY{+w}{        }\PY{n}{DateTimeFormatter}\PY{+w}{ }\PY{n}{strict}\PY{+w}{ }\PY{o}{=}\PY{+w}{ }\PY{n}{DateTimeFormatter}\PY{p}{.}\PY{n+na}{ofPattern}\PY{p}{(}\PY{l+s}{\PYZdq{}}\PY{l+s}{yyyy\PYZhy{}MM\PYZhy{}dd}\PY{l+s}{\PYZdq{}}\PY{p}{)}
\PY{+w}{                }\PY{p}{.}\PY{n+na}{withResolverStyle}\PY{p}{(}\PY{n}{ResolverStyle}\PY{p}{.}\PY{n+na}{STRICT}\PY{p}{)}\PY{p}{;}
\PY{+w}{        }\PY{n}{DateTimeFormatter}\PY{+w}{ }\PY{n}{smart}\PY{+w}{ }\PY{o}{=}\PY{+w}{ }\PY{n}{DateTimeFormatter}\PY{p}{.}\PY{n+na}{ofPattern}\PY{p}{(}\PY{l+s}{\PYZdq{}}\PY{l+s}{yyyy\PYZhy{}MM\PYZhy{}dd}\PY{l+s}{\PYZdq{}}\PY{p}{)}
\PY{+w}{                }\PY{p}{.}\PY{n+na}{withResolverStyle}\PY{p}{(}\PY{n}{ResolverStyle}\PY{p}{.}\PY{n+na}{SMART}\PY{p}{)}\PY{p}{;}\PY{+w}{ }\PY{c+c1}{// default}
\PY{+w}{        }\PY{n}{DateTimeFormatter}\PY{+w}{ }\PY{n}{lenient}\PY{+w}{ }\PY{o}{=}\PY{+w}{ }\PY{n}{DateTimeFormatter}\PY{p}{.}\PY{n+na}{ofPattern}\PY{p}{(}\PY{l+s}{\PYZdq{}}\PY{l+s}{yyyy\PYZhy{}MM\PYZhy{}dd}\PY{l+s}{\PYZdq{}}\PY{p}{)}
\PY{+w}{                }\PY{p}{.}\PY{n+na}{withResolverStyle}\PY{p}{(}\PY{n}{ResolverStyle}\PY{p}{.}\PY{n+na}{LENIENT}\PY{p}{)}\PY{p}{;}

\PY{+w}{        }\PY{c+c1}{// \PYZdq{}2026\PYZhy{}02\PYZhy{}30\PYZdq{} \PYZhy{} February never has 30 days}
\PY{+w}{        }\PY{k}{try}\PY{+w}{ }\PY{p}{\PYZob{}}
\PY{+w}{            }\PY{n}{LocalDate}\PY{p}{.}\PY{n+na}{parse}\PY{p}{(}\PY{l+s}{\PYZdq{}}\PY{l+s}{2026\PYZhy{}02\PYZhy{}30}\PY{l+s}{\PYZdq{}}\PY{p}{,}\PY{+w}{ }\PY{n}{strict}\PY{p}{)}\PY{p}{;}
\PY{+w}{        }\PY{p}{\PYZcb{}}\PY{+w}{ }\PY{k}{catch}\PY{+w}{ }\PY{p}{(}\PY{n}{Exception}\PY{+w}{ }\PY{n}{e}\PY{p}{)}\PY{+w}{ }\PY{p}{\PYZob{}}
\PY{+w}{            }\PY{n}{System}\PY{p}{.}\PY{n+na}{out}\PY{p}{.}\PY{n+na}{println}\PY{p}{(}\PY{l+s}{\PYZdq{}}\PY{l+s}{STRICT fails: }\PY{l+s}{\PYZdq{}}\PY{+w}{ }\PY{o}{+}\PY{+w}{ }\PY{n}{e}\PY{p}{.}\PY{n+na}{getMessage}\PY{p}{(}\PY{p}{)}\PY{p}{)}\PY{p}{;}
\PY{+w}{            }\PY{c+c1}{// Output: STRICT fails: Invalid date \PYZsq{}FEBRUARY 30\PYZsq{}}
\PY{+w}{        }\PY{p}{\PYZcb{}}

\PY{+w}{        }\PY{k}{try}\PY{+w}{ }\PY{p}{\PYZob{}}
\PY{+w}{            }\PY{n}{LocalDate}\PY{p}{.}\PY{n+na}{parse}\PY{p}{(}\PY{l+s}{\PYZdq{}}\PY{l+s}{2026\PYZhy{}02\PYZhy{}30}\PY{l+s}{\PYZdq{}}\PY{p}{,}\PY{+w}{ }\PY{n}{smart}\PY{p}{)}\PY{p}{;}
\PY{+w}{        }\PY{p}{\PYZcb{}}\PY{+w}{ }\PY{k}{catch}\PY{+w}{ }\PY{p}{(}\PY{n}{Exception}\PY{+w}{ }\PY{n}{e}\PY{p}{)}\PY{+w}{ }\PY{p}{\PYZob{}}
\PY{+w}{            }\PY{n}{System}\PY{p}{.}\PY{n+na}{out}\PY{p}{.}\PY{n+na}{println}\PY{p}{(}\PY{l+s}{\PYZdq{}}\PY{l+s}{SMART fails: }\PY{l+s}{\PYZdq{}}\PY{+w}{ }\PY{o}{+}\PY{+w}{ }\PY{n}{e}\PY{p}{.}\PY{n+na}{getMessage}\PY{p}{(}\PY{p}{)}\PY{p}{)}\PY{p}{;}
\PY{+w}{            }\PY{c+c1}{// Output: SMART fails: Invalid date \PYZsq{}FEBRUARY 30\PYZsq{}}
\PY{+w}{        }\PY{p}{\PYZcb{}}

\PY{+w}{        }\PY{n}{LocalDate}\PY{+w}{ }\PY{n}{lenientResult}\PY{+w}{ }\PY{o}{=}\PY{+w}{ }\PY{n}{LocalDate}\PY{p}{.}\PY{n+na}{parse}\PY{p}{(}\PY{l+s}{\PYZdq{}}\PY{l+s}{2026\PYZhy{}02\PYZhy{}30}\PY{l+s}{\PYZdq{}}\PY{p}{,}\PY{+w}{ }\PY{n}{lenient}\PY{p}{)}\PY{p}{;}
\PY{+w}{        }\PY{n}{System}\PY{p}{.}\PY{n+na}{out}\PY{p}{.}\PY{n+na}{println}\PY{p}{(}\PY{n}{lenientResult}\PY{p}{)}\PY{p}{;}
\PY{+w}{        }\PY{c+c1}{// Output: 2026\PYZhy{}03\PYZhy{}02  (overflowed into March)}
\PY{+w}{    }\PY{p}{\PYZcb{}}
\PY{p}{\PYZcb{}}
\end{Verbatim}
\end{codebox}

\subsection[ISO-8601]{ISO-8601}
\label{h:9:iso-8601}
ISO-8601 is the international standard for representing dates and times as text (\code{2026-\allowbreak{}08-\allowbreak{}07T14:\allowbreak{}30:\allowbreak{}15+\allowbreak{}02:\allowbreak{}00}). It is unambiguous, sortable as a plain string, and the default choice for APIs, logs, and config files. \code{java.\allowbreak{}time}’s \code{toString()} methods and \code{Date\allowbreak{}Time\allowbreak{}Formatter.\allowbreak{}ISO\_*} constants follow it by default.

\begin{codebox}
\begin{Verbatim}[commandchars=\\\{\}]
\PY{k+kn}{import}\PY{+w}{ }\PY{n+nn}{java.time.OffsetDateTime}\PY{p}{;}
\PY{k+kn}{import}\PY{+w}{ }\PY{n+nn}{java.time.ZoneOffset}\PY{p}{;}

\PY{k+kd}{public}\PY{+w}{ }\PY{k+kd}{class} \PY{n+nc}{Iso8601Demo}\PY{+w}{ }\PY{p}{\PYZob{}}
\PY{+w}{    }\PY{k+kd}{public}\PY{+w}{ }\PY{k+kd}{static}\PY{+w}{ }\PY{k+kt}{void}\PY{+w}{ }\PY{n+nf}{main}\PY{p}{(}\PY{n}{String}\PY{o}{[}\PY{o}{]}\PY{+w}{ }\PY{n}{args}\PY{p}{)}\PY{+w}{ }\PY{p}{\PYZob{}}
\PY{+w}{        }\PY{n}{OffsetDateTime}\PY{+w}{ }\PY{n}{odt}\PY{+w}{ }\PY{o}{=}\PY{+w}{ }\PY{n}{OffsetDateTime}\PY{p}{.}\PY{n+na}{of}\PY{p}{(}\PY{l+m+mi}{2026}\PY{p}{,}\PY{+w}{ }\PY{l+m+mi}{8}\PY{p}{,}\PY{+w}{ }\PY{l+m+mi}{7}\PY{p}{,}\PY{+w}{ }\PY{l+m+mi}{14}\PY{p}{,}\PY{+w}{ }\PY{l+m+mi}{30}\PY{p}{,}\PY{+w}{ }\PY{l+m+mi}{15}\PY{p}{,}\PY{+w}{ }\PY{l+m+mi}{0}\PY{p}{,}\PY{+w}{ }\PY{n}{ZoneOffset}\PY{p}{.}\PY{n+na}{ofHours}\PY{p}{(}\PY{l+m+mi}{2}\PY{p}{)}\PY{p}{)}\PY{p}{;}

\PY{+w}{        }\PY{n}{System}\PY{p}{.}\PY{n+na}{out}\PY{p}{.}\PY{n+na}{println}\PY{p}{(}\PY{n}{odt}\PY{p}{)}\PY{p}{;}\PY{+w}{ }\PY{c+c1}{// 2026\PYZhy{}08\PYZhy{}07T14:30:15+02:00 (ISO\PYZhy{}8601, no formatter needed)}

\PY{+w}{        }\PY{n}{OffsetDateTime}\PY{+w}{ }\PY{n}{parsedBack}\PY{+w}{ }\PY{o}{=}\PY{+w}{ }\PY{n}{OffsetDateTime}\PY{p}{.}\PY{n+na}{parse}\PY{p}{(}\PY{l+s}{\PYZdq{}}\PY{l+s}{2026\PYZhy{}08\PYZhy{}07T14:30:15+02:00}\PY{l+s}{\PYZdq{}}\PY{p}{)}\PY{p}{;}
\PY{+w}{        }\PY{n}{System}\PY{p}{.}\PY{n+na}{out}\PY{p}{.}\PY{n+na}{println}\PY{p}{(}\PY{n}{parsedBack}\PY{p}{.}\PY{n+na}{equals}\PY{p}{(}\PY{n}{odt}\PY{p}{)}\PY{p}{)}\PY{p}{;}\PY{+w}{ }\PY{c+c1}{// true}
\PY{+w}{    }\PY{p}{\PYZcb{}}
\PY{p}{\PYZcb{}}
\end{Verbatim}
\end{codebox}

\subsection[The YYYY vs yyyy week-year bug]{The \code{YYYY} vs \code{yyyy} week-year bug}
\label{h:9:the-yyyy-vs-yyyy-week-year-bug}
This is a famous, easy-to-miss bug. Lowercase \code{yyyy} is the \textbf{calendar year}. Uppercase \code{YYYY} is the \textbf{week-based year} (ISO week-year, used together with week-of-year \code{w}). Near the start or end of a year, they can disagree, silently shifting dates by a year.

\begin{codebox}
\begin{Verbatim}[commandchars=\\\{\}]
\PY{k+kn}{import}\PY{+w}{ }\PY{n+nn}{java.time.LocalDate}\PY{p}{;}
\PY{k+kn}{import}\PY{+w}{ }\PY{n+nn}{java.time.format.DateTimeFormatter}\PY{p}{;}

\PY{k+kd}{public}\PY{+w}{ }\PY{k+kd}{class} \PY{n+nc}{WeekYearBugDemo}\PY{+w}{ }\PY{p}{\PYZob{}}
\PY{+w}{    }\PY{k+kd}{public}\PY{+w}{ }\PY{k+kd}{static}\PY{+w}{ }\PY{k+kt}{void}\PY{+w}{ }\PY{n+nf}{main}\PY{p}{(}\PY{n}{String}\PY{o}{[}\PY{o}{]}\PY{+w}{ }\PY{n}{args}\PY{p}{)}\PY{+w}{ }\PY{p}{\PYZob{}}
\PY{+w}{        }\PY{c+c1}{// 2027\PYZhy{}01\PYZhy{}01 is a Friday, and it belongs to ISO week 53 of year 2026}
\PY{+w}{        }\PY{n}{LocalDate}\PY{+w}{ }\PY{n}{date}\PY{+w}{ }\PY{o}{=}\PY{+w}{ }\PY{n}{LocalDate}\PY{p}{.}\PY{n+na}{of}\PY{p}{(}\PY{l+m+mi}{2027}\PY{p}{,}\PY{+w}{ }\PY{l+m+mi}{1}\PY{p}{,}\PY{+w}{ }\PY{l+m+mi}{1}\PY{p}{)}\PY{p}{;}

\PY{+w}{        }\PY{n}{DateTimeFormatter}\PY{+w}{ }\PY{n}{withLowercaseY}\PY{+w}{ }\PY{o}{=}\PY{+w}{ }\PY{n}{DateTimeFormatter}\PY{p}{.}\PY{n+na}{ofPattern}\PY{p}{(}\PY{l+s}{\PYZdq{}}\PY{l+s}{yyyy\PYZhy{}MM\PYZhy{}dd}\PY{l+s}{\PYZdq{}}\PY{p}{)}\PY{p}{;}
\PY{+w}{        }\PY{n}{DateTimeFormatter}\PY{+w}{ }\PY{n}{withUppercaseY}\PY{+w}{ }\PY{o}{=}\PY{+w}{ }\PY{n}{DateTimeFormatter}\PY{p}{.}\PY{n+na}{ofPattern}\PY{p}{(}\PY{l+s}{\PYZdq{}}\PY{l+s}{YYYY\PYZhy{}MM\PYZhy{}dd}\PY{l+s}{\PYZdq{}}\PY{p}{)}\PY{p}{;}

\PY{+w}{        }\PY{n}{System}\PY{p}{.}\PY{n+na}{out}\PY{p}{.}\PY{n+na}{println}\PY{p}{(}\PY{n}{date}\PY{p}{.}\PY{n+na}{format}\PY{p}{(}\PY{n}{withLowercaseY}\PY{p}{)}\PY{p}{)}\PY{p}{;}\PY{+w}{ }\PY{c+c1}{// 2027\PYZhy{}01\PYZhy{}01 (correct calendar date)}
\PY{+w}{        }\PY{n}{System}\PY{p}{.}\PY{n+na}{out}\PY{p}{.}\PY{n+na}{println}\PY{p}{(}\PY{n}{date}\PY{p}{.}\PY{n+na}{format}\PY{p}{(}\PY{n}{withUppercaseY}\PY{p}{)}\PY{p}{)}\PY{p}{;}\PY{+w}{ }\PY{c+c1}{// 2026\PYZhy{}01\PYZhy{}01 (!) week\PYZhy{}year, NOT what most people expect}

\PY{+w}{        }\PY{c+c1}{// Rule of thumb for code review: always use lowercase yyyy unless}
\PY{+w}{        }\PY{c+c1}{// you are deliberately building an ISO week\PYZhy{}date format (with \PYZsq{}w\PYZsq{}).}
\PY{+w}{    }\PY{p}{\PYZcb{}}
\PY{p}{\PYZcb{}}
\end{Verbatim}
\end{codebox}

\subsection[Parsing failures and DateTimeParseException]{Parsing failures and \code{DateTimeParseException}}
\label{h:9:parsing-failures-and-datetimeparseexception}
Any failed parse throws an unchecked \code{DateTimeParseException}. Because it’s unchecked, it is easy to forget to handle at the boundary where you accept external, untrusted date strings (user input, CSV files, third-party APIs).

\begin{codebox}
\begin{Verbatim}[commandchars=\\\{\}]
\PY{k+kn}{import}\PY{+w}{ }\PY{n+nn}{java.time.LocalDate}\PY{p}{;}
\PY{k+kn}{import}\PY{+w}{ }\PY{n+nn}{java.time.format.DateTimeParseException}\PY{p}{;}

\PY{k+kd}{public}\PY{+w}{ }\PY{k+kd}{class} \PY{n+nc}{ParseFailureDemo}\PY{+w}{ }\PY{p}{\PYZob{}}
\PY{+w}{    }\PY{k+kd}{public}\PY{+w}{ }\PY{k+kd}{static}\PY{+w}{ }\PY{k+kt}{void}\PY{+w}{ }\PY{n+nf}{main}\PY{p}{(}\PY{n}{String}\PY{o}{[}\PY{o}{]}\PY{+w}{ }\PY{n}{args}\PY{p}{)}\PY{+w}{ }\PY{p}{\PYZob{}}
\PY{+w}{        }\PY{n}{String}\PY{o}{[}\PY{o}{]}\PY{+w}{ }\PY{n}{inputs}\PY{+w}{ }\PY{o}{=}\PY{+w}{ }\PY{p}{\PYZob{}}\PY{l+s}{\PYZdq{}}\PY{l+s}{2026\PYZhy{}08\PYZhy{}07}\PY{l+s}{\PYZdq{}}\PY{p}{,}\PY{+w}{ }\PY{l+s}{\PYZdq{}}\PY{l+s}{07/08/2026}\PY{l+s}{\PYZdq{}}\PY{p}{,}\PY{+w}{ }\PY{l+s}{\PYZdq{}}\PY{l+s}{not\PYZhy{}a\PYZhy{}date}\PY{l+s}{\PYZdq{}}\PY{p}{\PYZcb{}}\PY{p}{;}

\PY{+w}{        }\PY{k}{for}\PY{+w}{ }\PY{p}{(}\PY{n}{String}\PY{+w}{ }\PY{n}{input}\PY{+w}{ }\PY{p}{:}\PY{+w}{ }\PY{n}{inputs}\PY{p}{)}\PY{+w}{ }\PY{p}{\PYZob{}}
\PY{+w}{            }\PY{k}{try}\PY{+w}{ }\PY{p}{\PYZob{}}
\PY{+w}{                }\PY{n}{LocalDate}\PY{+w}{ }\PY{n}{parsed}\PY{+w}{ }\PY{o}{=}\PY{+w}{ }\PY{n}{LocalDate}\PY{p}{.}\PY{n+na}{parse}\PY{p}{(}\PY{n}{input}\PY{p}{)}\PY{p}{;}\PY{+w}{ }\PY{c+c1}{// expects ISO\PYZus{}LOCAL\PYZus{}DATE by default}
\PY{+w}{                }\PY{n}{System}\PY{p}{.}\PY{n+na}{out}\PY{p}{.}\PY{n+na}{println}\PY{p}{(}\PY{l+s}{\PYZdq{}}\PY{l+s}{Parsed: }\PY{l+s}{\PYZdq{}}\PY{+w}{ }\PY{o}{+}\PY{+w}{ }\PY{n}{parsed}\PY{p}{)}\PY{p}{;}
\PY{+w}{            }\PY{p}{\PYZcb{}}\PY{+w}{ }\PY{k}{catch}\PY{+w}{ }\PY{p}{(}\PY{n}{DateTimeParseException}\PY{+w}{ }\PY{n}{e}\PY{p}{)}\PY{+w}{ }\PY{p}{\PYZob{}}
\PY{+w}{                }\PY{n}{System}\PY{p}{.}\PY{n+na}{out}\PY{p}{.}\PY{n+na}{println}\PY{p}{(}\PY{l+s}{\PYZdq{}}\PY{l+s}{Failed to parse \PYZsq{}}\PY{l+s}{\PYZdq{}}\PY{+w}{ }\PY{o}{+}\PY{+w}{ }\PY{n}{input}\PY{+w}{ }\PY{o}{+}\PY{+w}{ }\PY{l+s}{\PYZdq{}}\PY{l+s}{\PYZsq{}: }\PY{l+s}{\PYZdq{}}\PY{+w}{ }\PY{o}{+}\PY{+w}{ }\PY{n}{e}\PY{p}{.}\PY{n+na}{getMessage}\PY{p}{(}\PY{p}{)}\PY{p}{)}\PY{p}{;}
\PY{+w}{            }\PY{p}{\PYZcb{}}
\PY{+w}{        }\PY{p}{\PYZcb{}}
\PY{+w}{        }\PY{c+c1}{// Output:}
\PY{+w}{        }\PY{c+c1}{// Parsed: 2026\PYZhy{}08\PYZhy{}07}
\PY{+w}{        }\PY{c+c1}{// Failed to parse \PYZsq{}07/08/2026\PYZsq{}: Text \PYZsq{}07/08/2026\PYZsq{} could not be parsed at index 2}
\PY{+w}{        }\PY{c+c1}{// Failed to parse \PYZsq{}not\PYZhy{}a\PYZhy{}date\PYZsq{}: Text \PYZsq{}not\PYZhy{}a\PYZhy{}date\PYZsq{} could not be parsed at index 0}
\PY{+w}{    }\PY{p}{\PYZcb{}}
\PY{p}{\PYZcb{}}
\end{Verbatim}
\end{codebox}

\subsection[Thread safety: DateTimeFormatter vs SimpleDateFormat]{Thread safety: \code{DateTimeFormatter} vs \code{SimpleDateFormat}}
\label{h:9:thread-safety-datetimeformatter-vs-simpledateformat}
\code{SimpleDateFormat} is \textbf{not thread-safe} — internally it uses a mutable \code{Calendar} field, so two threads calling \code{format()} on the same instance concurrently can corrupt each other’s results. \code{DateTimeFormatter} is \textbf{immutable and thread-safe}, so it is safe (and recommended) to share a single instance as a \code{static~\allowbreak{}final} constant.

\begin{codebox}
\begin{Verbatim}[commandchars=\\\{\}]
\PY{k+kn}{import}\PY{+w}{ }\PY{n+nn}{java.text.SimpleDateFormat}\PY{p}{;}
\PY{k+kn}{import}\PY{+w}{ }\PY{n+nn}{java.time.format.DateTimeFormatter}\PY{p}{;}
\PY{k+kn}{import}\PY{+w}{ }\PY{n+nn}{java.util.Date}\PY{p}{;}
\PY{k+kn}{import}\PY{+w}{ }\PY{n+nn}{java.util.concurrent.ExecutorService}\PY{p}{;}
\PY{k+kn}{import}\PY{+w}{ }\PY{n+nn}{java.util.concurrent.Executors}\PY{p}{;}

\PY{k+kd}{public}\PY{+w}{ }\PY{k+kd}{class} \PY{n+nc}{ThreadSafetyDemo}\PY{+w}{ }\PY{p}{\PYZob{}}

\PY{+w}{    }\PY{c+c1}{// UNSAFE: sharing a SimpleDateFormat across threads can produce garbled output or exceptions}
\PY{+w}{    }\PY{k+kd}{private}\PY{+w}{ }\PY{k+kd}{static}\PY{+w}{ }\PY{k+kd}{final}\PY{+w}{ }\PY{n}{SimpleDateFormat}\PY{+w}{ }\PY{n}{UNSAFE\PYZus{}FORMAT}\PY{+w}{ }\PY{o}{=}\PY{+w}{ }\PY{k}{new}\PY{+w}{ }\PY{n}{SimpleDateFormat}\PY{p}{(}\PY{l+s}{\PYZdq{}}\PY{l+s}{yyyy\PYZhy{}MM\PYZhy{}dd HH:mm:ss}\PY{l+s}{\PYZdq{}}\PY{p}{)}\PY{p}{;}

\PY{+w}{    }\PY{c+c1}{// SAFE: DateTimeFormatter is immutable, sharing it is fine and idiomatic}
\PY{+w}{    }\PY{k+kd}{private}\PY{+w}{ }\PY{k+kd}{static}\PY{+w}{ }\PY{k+kd}{final}\PY{+w}{ }\PY{n}{DateTimeFormatter}\PY{+w}{ }\PY{n}{SAFE\PYZus{}FORMAT}\PY{+w}{ }\PY{o}{=}
\PY{+w}{            }\PY{n}{DateTimeFormatter}\PY{p}{.}\PY{n+na}{ofPattern}\PY{p}{(}\PY{l+s}{\PYZdq{}}\PY{l+s}{yyyy\PYZhy{}MM\PYZhy{}dd HH:mm:ss}\PY{l+s}{\PYZdq{}}\PY{p}{)}\PY{p}{;}

\PY{+w}{    }\PY{k+kd}{public}\PY{+w}{ }\PY{k+kd}{static}\PY{+w}{ }\PY{k+kt}{void}\PY{+w}{ }\PY{n+nf}{main}\PY{p}{(}\PY{n}{String}\PY{o}{[}\PY{o}{]}\PY{+w}{ }\PY{n}{args}\PY{p}{)}\PY{+w}{ }\PY{k+kd}{throws}\PY{+w}{ }\PY{n}{InterruptedException}\PY{+w}{ }\PY{p}{\PYZob{}}
\PY{+w}{        }\PY{n}{ExecutorService}\PY{+w}{ }\PY{n}{pool}\PY{+w}{ }\PY{o}{=}\PY{+w}{ }\PY{n}{Executors}\PY{p}{.}\PY{n+na}{newFixedThreadPool}\PY{p}{(}\PY{l+m+mi}{4}\PY{p}{)}\PY{p}{;}
\PY{+w}{        }\PY{k}{for}\PY{+w}{ }\PY{p}{(}\PY{k+kt}{int}\PY{+w}{ }\PY{n}{i}\PY{+w}{ }\PY{o}{=}\PY{+w}{ }\PY{l+m+mi}{0}\PY{p}{;}\PY{+w}{ }\PY{n}{i}\PY{+w}{ }\PY{o}{\PYZlt{}}\PY{+w}{ }\PY{l+m+mi}{4}\PY{p}{;}\PY{+w}{ }\PY{n}{i}\PY{o}{+}\PY{o}{+}\PY{p}{)}\PY{+w}{ }\PY{p}{\PYZob{}}
\PY{+w}{            }\PY{n}{pool}\PY{p}{.}\PY{n+na}{submit}\PY{p}{(}\PY{p}{(}\PY{p}{)}\PY{+w}{ }\PY{o}{\PYZhy{}}\PY{o}{\PYZgt{}}\PY{+w}{ }\PY{p}{\PYZob{}}
\PY{+w}{                }\PY{c+c1}{// Calling UNSAFE\PYZus{}FORMAT.format(new Date()) here concurrently}
\PY{+w}{                }\PY{c+c1}{// may throw NumberFormatException or return a corrupted string.}
\PY{+w}{                }\PY{n}{String}\PY{+w}{ }\PY{n}{safeResult}\PY{+w}{ }\PY{o}{=}\PY{+w}{ }\PY{n}{java}\PY{p}{.}\PY{n+na}{time}\PY{p}{.}\PY{n+na}{LocalDateTime}\PY{p}{.}\PY{n+na}{now}\PY{p}{(}\PY{p}{)}\PY{p}{.}\PY{n+na}{format}\PY{p}{(}\PY{n}{SAFE\PYZus{}FORMAT}\PY{p}{)}\PY{p}{;}
\PY{+w}{                }\PY{n}{System}\PY{p}{.}\PY{n+na}{out}\PY{p}{.}\PY{n+na}{println}\PY{p}{(}\PY{n}{Thread}\PY{p}{.}\PY{n+na}{currentThread}\PY{p}{(}\PY{p}{)}\PY{p}{.}\PY{n+na}{getName}\PY{p}{(}\PY{p}{)}\PY{+w}{ }\PY{o}{+}\PY{+w}{ }\PY{l+s}{\PYZdq{}}\PY{l+s}{: }\PY{l+s}{\PYZdq{}}\PY{+w}{ }\PY{o}{+}\PY{+w}{ }\PY{n}{safeResult}\PY{p}{)}\PY{p}{;}
\PY{+w}{            }\PY{p}{\PYZcb{}}\PY{p}{)}\PY{p}{;}
\PY{+w}{        }\PY{p}{\PYZcb{}}
\PY{+w}{        }\PY{n}{pool}\PY{p}{.}\PY{n+na}{shutdown}\PY{p}{(}\PY{p}{)}\PY{p}{;}
\PY{+w}{        }\PY{c+c1}{// Fix for legacy code: wrap SimpleDateFormat in a ThreadLocal, or just switch to DateTimeFormatter.}
\PY{+w}{    }\PY{p}{\PYZcb{}}
\PY{p}{\PYZcb{}}
\end{Verbatim}
\end{codebox}

\section[Internationalization (i18n)]{Internationalization (i18n)}
\label{h:9:internationalization-i18n}
Internationalization (“i18n” — 18 letters between “i” and “n”) is designing software so it \emph{can} be adapted to different languages and regions without code changes. Localization (“l10n”) is the actual act of adapting it for one specific locale (translating text, formatting dates, etc.).

\subsection[Locale]{\code{Locale}}
\label{h:9:locale}
A \code{Locale} identifies a language, optionally a country/region, and optionally a variant. It drives almost every locale-sensitive API in the JDK.

\begin{codebox}
\begin{Verbatim}[commandchars=\\\{\}]
\PY{k+kn}{import}\PY{+w}{ }\PY{n+nn}{java.util.Locale}\PY{p}{;}

\PY{k+kd}{public}\PY{+w}{ }\PY{k+kd}{class} \PY{n+nc}{LocaleDemo}\PY{+w}{ }\PY{p}{\PYZob{}}
\PY{+w}{    }\PY{k+kd}{public}\PY{+w}{ }\PY{k+kd}{static}\PY{+w}{ }\PY{k+kt}{void}\PY{+w}{ }\PY{n+nf}{main}\PY{p}{(}\PY{n}{String}\PY{o}{[}\PY{o}{]}\PY{+w}{ }\PY{n}{args}\PY{p}{)}\PY{+w}{ }\PY{p}{\PYZob{}}
\PY{+w}{        }\PY{n}{Locale}\PY{+w}{ }\PY{n}{usEnglish}\PY{+w}{ }\PY{o}{=}\PY{+w}{ }\PY{n}{Locale}\PY{p}{.}\PY{n+na}{US}\PY{p}{;}\PY{+w}{                       }\PY{c+c1}{// language=en, country=US}
\PY{+w}{        }\PY{n}{Locale}\PY{+w}{ }\PY{n}{german}\PY{+w}{ }\PY{o}{=}\PY{+w}{ }\PY{n}{Locale}\PY{p}{.}\PY{n+na}{GERMANY}\PY{p}{;}\PY{+w}{                     }\PY{c+c1}{// language=de, country=DE}
\PY{+w}{        }\PY{n}{Locale}\PY{+w}{ }\PY{n}{custom}\PY{+w}{ }\PY{o}{=}\PY{+w}{ }\PY{k}{new}\PY{+w}{ }\PY{n}{Locale}\PY{p}{(}\PY{l+s}{\PYZdq{}}\PY{l+s}{en}\PY{l+s}{\PYZdq{}}\PY{p}{,}\PY{+w}{ }\PY{l+s}{\PYZdq{}}\PY{l+s}{IN}\PY{l+s}{\PYZdq{}}\PY{p}{)}\PY{p}{;}\PY{+w}{              }\PY{c+c1}{// English as used in India}
\PY{+w}{        }\PY{n}{Locale}\PY{+w}{ }\PY{n}{fromTag}\PY{+w}{ }\PY{o}{=}\PY{+w}{ }\PY{n}{Locale}\PY{p}{.}\PY{n+na}{forLanguageTag}\PY{p}{(}\PY{l+s}{\PYZdq{}}\PY{l+s}{pt\PYZhy{}BR}\PY{l+s}{\PYZdq{}}\PY{p}{)}\PY{p}{;}\PY{+w}{      }\PY{c+c1}{// Brazilian Portuguese, BCP 47 tag}
\PY{+w}{        }\PY{n}{Locale}\PY{+w}{ }\PY{n}{root}\PY{+w}{ }\PY{o}{=}\PY{+w}{ }\PY{n}{Locale}\PY{p}{.}\PY{n+na}{ROOT}\PY{p}{;}\PY{+w}{                            }\PY{c+c1}{// no language/country \PYZhy{} a neutral, culture\PYZhy{}independent locale}

\PY{+w}{        }\PY{n}{System}\PY{p}{.}\PY{n+na}{out}\PY{p}{.}\PY{n+na}{println}\PY{p}{(}\PY{n}{usEnglish}\PY{p}{.}\PY{n+na}{getLanguage}\PY{p}{(}\PY{p}{)}\PY{+w}{ }\PY{o}{+}\PY{+w}{ }\PY{l+s}{\PYZdq{}}\PY{l+s}{\PYZhy{}}\PY{l+s}{\PYZdq{}}\PY{+w}{ }\PY{o}{+}\PY{+w}{ }\PY{n}{usEnglish}\PY{p}{.}\PY{n+na}{getCountry}\PY{p}{(}\PY{p}{)}\PY{p}{)}\PY{p}{;}\PY{+w}{ }\PY{c+c1}{// en\PYZhy{}US}
\PY{+w}{        }\PY{n}{System}\PY{p}{.}\PY{n+na}{out}\PY{p}{.}\PY{n+na}{println}\PY{p}{(}\PY{n}{german}\PY{p}{.}\PY{n+na}{getDisplayName}\PY{p}{(}\PY{n}{Locale}\PY{p}{.}\PY{n+na}{US}\PY{p}{)}\PY{p}{)}\PY{p}{;}\PY{+w}{                        }\PY{c+c1}{// German (Germany)}
\PY{+w}{        }\PY{n}{System}\PY{p}{.}\PY{n+na}{out}\PY{p}{.}\PY{n+na}{println}\PY{p}{(}\PY{n}{fromTag}\PY{p}{.}\PY{n+na}{getDisplayName}\PY{p}{(}\PY{p}{)}\PY{p}{)}\PY{p}{;}\PY{+w}{                                }\PY{c+c1}{// Portuguese (Brazil)}
\PY{+w}{        }\PY{n}{System}\PY{p}{.}\PY{n+na}{out}\PY{p}{.}\PY{n+na}{println}\PY{p}{(}\PY{n}{root}\PY{p}{)}\PY{p}{;}
\PY{+w}{        }\PY{c+c1}{// Output: (empty string) \PYZhy{} Locale.ROOT has no language, country, or variant}

\PY{+w}{        }\PY{c+c1}{// Locale.ROOT is the right choice for locale\PYZhy{}independent operations,}
\PY{+w}{        }\PY{c+c1}{// like normalizing internal identifiers or keys, not user\PYZhy{}facing text.}
\PY{+w}{        }\PY{n}{System}\PY{p}{.}\PY{n+na}{out}\PY{p}{.}\PY{n+na}{println}\PY{p}{(}\PY{l+s}{\PYZdq{}}\PY{l+s}{HELLO}\PY{l+s}{\PYZdq{}}\PY{p}{.}\PY{n+na}{toLowerCase}\PY{p}{(}\PY{n}{Locale}\PY{p}{.}\PY{n+na}{ROOT}\PY{p}{)}\PY{p}{)}\PY{p}{;}\PY{+w}{ }\PY{c+c1}{// hello}
\PY{+w}{    }\PY{p}{\PYZcb{}}
\PY{p}{\PYZcb{}}
\end{Verbatim}
\end{codebox}

\subsection[ResourceBundle and property files]{\code{ResourceBundle} and property files}
\label{h:9:resourcebundle-and-property-files}
\code{ResourceBundle} loads translated text from \code{.\allowbreak{}properties} files named by locale, such as \code{Messages\_\allowbreak{}en.\allowbreak{}properties}, \code{Messages\_\allowbreak{}de.\allowbreak{}properties}, and \code{Messages.\allowbreak{}properties} (the default fallback).

\begin{codebox}
\begin{Verbatim}
# Messages.properties (default/fallback)
greeting=Hello, {0}!

# Messages_de.properties
greeting=Hallo, {0}!

# Messages_fr.properties
greeting=Bonjour, {0} !
\end{Verbatim}
\end{codebox}

\begin{codebox}
\begin{Verbatim}[commandchars=\\\{\}]
\PY{k+kn}{import}\PY{+w}{ }\PY{n+nn}{java.text.MessageFormat}\PY{p}{;}
\PY{k+kn}{import}\PY{+w}{ }\PY{n+nn}{java.util.Locale}\PY{p}{;}
\PY{k+kn}{import}\PY{+w}{ }\PY{n+nn}{java.util.ResourceBundle}\PY{p}{;}

\PY{k+kd}{public}\PY{+w}{ }\PY{k+kd}{class} \PY{n+nc}{ResourceBundleDemo}\PY{+w}{ }\PY{p}{\PYZob{}}
\PY{+w}{    }\PY{k+kd}{public}\PY{+w}{ }\PY{k+kd}{static}\PY{+w}{ }\PY{k+kt}{void}\PY{+w}{ }\PY{n+nf}{main}\PY{p}{(}\PY{n}{String}\PY{o}{[}\PY{o}{]}\PY{+w}{ }\PY{n}{args}\PY{p}{)}\PY{+w}{ }\PY{p}{\PYZob{}}
\PY{+w}{        }\PY{n}{ResourceBundle}\PY{+w}{ }\PY{n}{enBundle}\PY{+w}{ }\PY{o}{=}\PY{+w}{ }\PY{n}{ResourceBundle}\PY{p}{.}\PY{n+na}{getBundle}\PY{p}{(}\PY{l+s}{\PYZdq{}}\PY{l+s}{Messages}\PY{l+s}{\PYZdq{}}\PY{p}{,}\PY{+w}{ }\PY{n}{Locale}\PY{p}{.}\PY{n+na}{ENGLISH}\PY{p}{)}\PY{p}{;}
\PY{+w}{        }\PY{n}{ResourceBundle}\PY{+w}{ }\PY{n}{deBundle}\PY{+w}{ }\PY{o}{=}\PY{+w}{ }\PY{n}{ResourceBundle}\PY{p}{.}\PY{n+na}{getBundle}\PY{p}{(}\PY{l+s}{\PYZdq{}}\PY{l+s}{Messages}\PY{l+s}{\PYZdq{}}\PY{p}{,}\PY{+w}{ }\PY{n}{Locale}\PY{p}{.}\PY{n+na}{GERMAN}\PY{p}{)}\PY{p}{;}

\PY{+w}{        }\PY{n}{String}\PY{+w}{ }\PY{n}{enGreeting}\PY{+w}{ }\PY{o}{=}\PY{+w}{ }\PY{n}{MessageFormat}\PY{p}{.}\PY{n+na}{format}\PY{p}{(}\PY{n}{enBundle}\PY{p}{.}\PY{n+na}{getString}\PY{p}{(}\PY{l+s}{\PYZdq{}}\PY{l+s}{greeting}\PY{l+s}{\PYZdq{}}\PY{p}{)}\PY{p}{,}\PY{+w}{ }\PY{l+s}{\PYZdq{}}\PY{l+s}{Alice}\PY{l+s}{\PYZdq{}}\PY{p}{)}\PY{p}{;}
\PY{+w}{        }\PY{n}{String}\PY{+w}{ }\PY{n}{deGreeting}\PY{+w}{ }\PY{o}{=}\PY{+w}{ }\PY{n}{MessageFormat}\PY{p}{.}\PY{n+na}{format}\PY{p}{(}\PY{n}{deBundle}\PY{p}{.}\PY{n+na}{getString}\PY{p}{(}\PY{l+s}{\PYZdq{}}\PY{l+s}{greeting}\PY{l+s}{\PYZdq{}}\PY{p}{)}\PY{p}{,}\PY{+w}{ }\PY{l+s}{\PYZdq{}}\PY{l+s}{Alice}\PY{l+s}{\PYZdq{}}\PY{p}{)}\PY{p}{;}

\PY{+w}{        }\PY{n}{System}\PY{p}{.}\PY{n+na}{out}\PY{p}{.}\PY{n+na}{println}\PY{p}{(}\PY{n}{enGreeting}\PY{p}{)}\PY{p}{;}\PY{+w}{ }\PY{c+c1}{// Hello, Alice!}
\PY{+w}{        }\PY{n}{System}\PY{p}{.}\PY{n+na}{out}\PY{p}{.}\PY{n+na}{println}\PY{p}{(}\PY{n}{deGreeting}\PY{p}{)}\PY{p}{;}\PY{+w}{ }\PY{c+c1}{// Hallo, Alice!}

\PY{+w}{        }\PY{c+c1}{// Requesting a locale with no matching file falls back to Messages.properties}
\PY{+w}{        }\PY{n}{ResourceBundle}\PY{+w}{ }\PY{n}{jaBundle}\PY{+w}{ }\PY{o}{=}\PY{+w}{ }\PY{n}{ResourceBundle}\PY{p}{.}\PY{n+na}{getBundle}\PY{p}{(}\PY{l+s}{\PYZdq{}}\PY{l+s}{Messages}\PY{l+s}{\PYZdq{}}\PY{p}{,}\PY{+w}{ }\PY{n}{Locale}\PY{p}{.}\PY{n+na}{JAPANESE}\PY{p}{)}\PY{p}{;}
\PY{+w}{        }\PY{n}{System}\PY{p}{.}\PY{n+na}{out}\PY{p}{.}\PY{n+na}{println}\PY{p}{(}\PY{n}{jaBundle}\PY{p}{.}\PY{n+na}{getString}\PY{p}{(}\PY{l+s}{\PYZdq{}}\PY{l+s}{greeting}\PY{l+s}{\PYZdq{}}\PY{p}{)}\PY{p}{)}\PY{p}{;}\PY{+w}{ }\PY{c+c1}{// Hello, \PYZob{}0\PYZcb{}! (fallback, untranslated placeholder)}
\PY{+w}{    }\PY{p}{\PYZcb{}}
\PY{p}{\PYZcb{}}
\end{Verbatim}
\end{codebox}

\subsection[MessageFormat: plurals and argument indexes]{\code{MessageFormat}: plurals and argument indexes}
\label{h:9:messageformat-plurals-and-argument-indexes}
\code{MessageFormat} supports positional arguments (\code{\{\allowbreak{}0\}}, \code{\{\allowbreak{}1\}}) that can be reordered per language, and a \code{choice}/\code{plural}-style syntax for picking text based on a number. This matters because different languages pluralize differently (some languages have more than two plural forms).

\begin{codebox}
\begin{Verbatim}[commandchars=\\\{\}]
\PY{k+kn}{import}\PY{+w}{ }\PY{n+nn}{java.text.MessageFormat}\PY{p}{;}
\PY{k+kn}{import}\PY{+w}{ }\PY{n+nn}{java.util.Locale}\PY{p}{;}

\PY{k+kd}{public}\PY{+w}{ }\PY{k+kd}{class} \PY{n+nc}{MessageFormatDemo}\PY{+w}{ }\PY{p}{\PYZob{}}
\PY{+w}{    }\PY{k+kd}{public}\PY{+w}{ }\PY{k+kd}{static}\PY{+w}{ }\PY{k+kt}{void}\PY{+w}{ }\PY{n+nf}{main}\PY{p}{(}\PY{n}{String}\PY{o}{[}\PY{o}{]}\PY{+w}{ }\PY{n}{args}\PY{p}{)}\PY{+w}{ }\PY{p}{\PYZob{}}
\PY{+w}{        }\PY{c+c1}{// Argument indexes let translators reorder words for grammar reasons}
\PY{+w}{        }\PY{n}{String}\PY{+w}{ }\PY{n}{pattern}\PY{+w}{ }\PY{o}{=}\PY{+w}{ }\PY{l+s}{\PYZdq{}}\PY{l+s}{\PYZob{}0\PYZcb{} sent \PYZob{}1\PYZcb{} a message.}\PY{l+s}{\PYZdq{}}\PY{p}{;}\PY{+w}{       }\PY{c+c1}{// English word order}
\PY{+w}{        }\PY{n}{String}\PY{+w}{ }\PY{n}{patternGerman}\PY{+w}{ }\PY{o}{=}\PY{+w}{ }\PY{l+s}{\PYZdq{}}\PY{l+s}{\PYZob{}1\PYZcb{} bekam eine Nachricht von \PYZob{}0\PYZcb{}.}\PY{l+s}{\PYZdq{}}\PY{p}{;}\PY{+w}{  }\PY{c+c1}{// reordered for German}

\PY{+w}{        }\PY{n}{System}\PY{p}{.}\PY{n+na}{out}\PY{p}{.}\PY{n+na}{println}\PY{p}{(}\PY{n}{MessageFormat}\PY{p}{.}\PY{n+na}{format}\PY{p}{(}\PY{n}{pattern}\PY{p}{,}\PY{+w}{ }\PY{l+s}{\PYZdq{}}\PY{l+s}{Alice}\PY{l+s}{\PYZdq{}}\PY{p}{,}\PY{+w}{ }\PY{l+s}{\PYZdq{}}\PY{l+s}{Bob}\PY{l+s}{\PYZdq{}}\PY{p}{)}\PY{p}{)}\PY{p}{;}
\PY{+w}{        }\PY{c+c1}{// Output: Alice sent Bob a message.}
\PY{+w}{        }\PY{n}{System}\PY{p}{.}\PY{n+na}{out}\PY{p}{.}\PY{n+na}{println}\PY{p}{(}\PY{n}{MessageFormat}\PY{p}{.}\PY{n+na}{format}\PY{p}{(}\PY{n}{patternGerman}\PY{p}{,}\PY{+w}{ }\PY{l+s}{\PYZdq{}}\PY{l+s}{Alice}\PY{l+s}{\PYZdq{}}\PY{p}{,}\PY{+w}{ }\PY{l+s}{\PYZdq{}}\PY{l+s}{Bob}\PY{l+s}{\PYZdq{}}\PY{p}{)}\PY{p}{)}\PY{p}{;}
\PY{+w}{        }\PY{c+c1}{// Output: Bob bekam eine Nachricht von Alice.}

\PY{+w}{        }\PY{c+c1}{// Simple pluralization using choice format}
\PY{+w}{        }\PY{n}{MessageFormat}\PY{+w}{ }\PY{n}{itemsFormat}\PY{+w}{ }\PY{o}{=}\PY{+w}{ }\PY{k}{new}\PY{+w}{ }\PY{n}{MessageFormat}\PY{p}{(}
\PY{+w}{                }\PY{l+s}{\PYZdq{}}\PY{l+s}{There \PYZob{}0,choice,0\PYZsh{}are no items|1\PYZsh{}is one item|1\PYZlt{}are \PYZob{}0,number,integer\PYZcb{} items\PYZcb{}.}\PY{l+s}{\PYZdq{}}\PY{p}{,}
\PY{+w}{                }\PY{n}{Locale}\PY{p}{.}\PY{n+na}{ENGLISH}\PY{p}{)}\PY{p}{;}

\PY{+w}{        }\PY{n}{System}\PY{p}{.}\PY{n+na}{out}\PY{p}{.}\PY{n+na}{println}\PY{p}{(}\PY{n}{itemsFormat}\PY{p}{.}\PY{n+na}{format}\PY{p}{(}\PY{k}{new}\PY{+w}{ }\PY{n}{Object}\PY{o}{[}\PY{o}{]}\PY{p}{\PYZob{}}\PY{l+m+mi}{0}\PY{p}{\PYZcb{}}\PY{p}{)}\PY{p}{)}\PY{p}{;}\PY{+w}{ }\PY{c+c1}{// There are no items.}
\PY{+w}{        }\PY{n}{System}\PY{p}{.}\PY{n+na}{out}\PY{p}{.}\PY{n+na}{println}\PY{p}{(}\PY{n}{itemsFormat}\PY{p}{.}\PY{n+na}{format}\PY{p}{(}\PY{k}{new}\PY{+w}{ }\PY{n}{Object}\PY{o}{[}\PY{o}{]}\PY{p}{\PYZob{}}\PY{l+m+mi}{1}\PY{p}{\PYZcb{}}\PY{p}{)}\PY{p}{)}\PY{p}{;}\PY{+w}{ }\PY{c+c1}{// There is one item.}
\PY{+w}{        }\PY{n}{System}\PY{p}{.}\PY{n+na}{out}\PY{p}{.}\PY{n+na}{println}\PY{p}{(}\PY{n}{itemsFormat}\PY{p}{.}\PY{n+na}{format}\PY{p}{(}\PY{k}{new}\PY{+w}{ }\PY{n}{Object}\PY{o}{[}\PY{o}{]}\PY{p}{\PYZob{}}\PY{l+m+mi}{5}\PY{p}{\PYZcb{}}\PY{p}{)}\PY{p}{)}\PY{p}{;}\PY{+w}{ }\PY{c+c1}{// There are 5 items.}
\PY{+w}{        }\PY{c+c1}{// Note: for serious plural\PYZhy{}rule handling (many languages have 3\PYZhy{}6 plural forms),}
\PY{+w}{        }\PY{c+c1}{// prefer a library that implements full CLDR plural rules (e.g. ICU4J).}
\PY{+w}{    }\PY{p}{\PYZcb{}}
\PY{p}{\PYZcb{}}
\end{Verbatim}
\end{codebox}

\subsection[NumberFormat and Currency]{\code{NumberFormat} and \code{Currency}}
\label{h:9:numberformat-and-currency}
Numbers and currency amounts are formatted very differently across locales — decimal separators, grouping separators, and currency symbol placement all vary.

\begin{codebox}
\begin{Verbatim}[commandchars=\\\{\}]
\PY{k+kn}{import}\PY{+w}{ }\PY{n+nn}{java.text.NumberFormat}\PY{p}{;}
\PY{k+kn}{import}\PY{+w}{ }\PY{n+nn}{java.util.Currency}\PY{p}{;}
\PY{k+kn}{import}\PY{+w}{ }\PY{n+nn}{java.util.Locale}\PY{p}{;}

\PY{k+kd}{public}\PY{+w}{ }\PY{k+kd}{class} \PY{n+nc}{NumberFormatDemo}\PY{+w}{ }\PY{p}{\PYZob{}}
\PY{+w}{    }\PY{k+kd}{public}\PY{+w}{ }\PY{k+kd}{static}\PY{+w}{ }\PY{k+kt}{void}\PY{+w}{ }\PY{n+nf}{main}\PY{p}{(}\PY{n}{String}\PY{o}{[}\PY{o}{]}\PY{+w}{ }\PY{n}{args}\PY{p}{)}\PY{+w}{ }\PY{p}{\PYZob{}}
\PY{+w}{        }\PY{k+kt}{double}\PY{+w}{ }\PY{n}{amount}\PY{+w}{ }\PY{o}{=}\PY{+w}{ }\PY{l+m+mf}{1234567.891}\PY{p}{;}

\PY{+w}{        }\PY{n}{NumberFormat}\PY{+w}{ }\PY{n}{us}\PY{+w}{ }\PY{o}{=}\PY{+w}{ }\PY{n}{NumberFormat}\PY{p}{.}\PY{n+na}{getNumberInstance}\PY{p}{(}\PY{n}{Locale}\PY{p}{.}\PY{n+na}{US}\PY{p}{)}\PY{p}{;}
\PY{+w}{        }\PY{n}{NumberFormat}\PY{+w}{ }\PY{n}{de}\PY{+w}{ }\PY{o}{=}\PY{+w}{ }\PY{n}{NumberFormat}\PY{p}{.}\PY{n+na}{getNumberInstance}\PY{p}{(}\PY{n}{Locale}\PY{p}{.}\PY{n+na}{GERMANY}\PY{p}{)}\PY{p}{;}

\PY{+w}{        }\PY{n}{System}\PY{p}{.}\PY{n+na}{out}\PY{p}{.}\PY{n+na}{println}\PY{p}{(}\PY{n}{us}\PY{p}{.}\PY{n+na}{format}\PY{p}{(}\PY{n}{amount}\PY{p}{)}\PY{p}{)}\PY{p}{;}\PY{+w}{ }\PY{c+c1}{// 1,234,567.891}
\PY{+w}{        }\PY{n}{System}\PY{p}{.}\PY{n+na}{out}\PY{p}{.}\PY{n+na}{println}\PY{p}{(}\PY{n}{de}\PY{p}{.}\PY{n+na}{format}\PY{p}{(}\PY{n}{amount}\PY{p}{)}\PY{p}{)}\PY{p}{;}\PY{+w}{ }\PY{c+c1}{// 1.234.567,891}

\PY{+w}{        }\PY{n}{NumberFormat}\PY{+w}{ }\PY{n}{usCurrency}\PY{+w}{ }\PY{o}{=}\PY{+w}{ }\PY{n}{NumberFormat}\PY{p}{.}\PY{n+na}{getCurrencyInstance}\PY{p}{(}\PY{n}{Locale}\PY{p}{.}\PY{n+na}{US}\PY{p}{)}\PY{p}{;}
\PY{+w}{        }\PY{n}{NumberFormat}\PY{+w}{ }\PY{n}{jpCurrency}\PY{+w}{ }\PY{o}{=}\PY{+w}{ }\PY{n}{NumberFormat}\PY{p}{.}\PY{n+na}{getCurrencyInstance}\PY{p}{(}\PY{n}{Locale}\PY{p}{.}\PY{n+na}{JAPAN}\PY{p}{)}\PY{p}{;}
\PY{+w}{        }\PY{n}{System}\PY{p}{.}\PY{n+na}{out}\PY{p}{.}\PY{n+na}{println}\PY{p}{(}\PY{n}{usCurrency}\PY{p}{.}\PY{n+na}{format}\PY{p}{(}\PY{l+m+mf}{1234.5}\PY{p}{)}\PY{p}{)}\PY{p}{;}\PY{+w}{ }\PY{c+c1}{// \PYZdl{}1,234.50}
\PY{+w}{        }\PY{n}{System}\PY{p}{.}\PY{n+na}{out}\PY{p}{.}\PY{n+na}{println}\PY{p}{(}\PY{n}{jpCurrency}\PY{p}{.}\PY{n+na}{format}\PY{p}{(}\PY{l+m+mf}{1234.5}\PY{p}{)}\PY{p}{)}\PY{p}{;}\PY{+w}{ }\PY{c+c1}{// ￥1,235  (Yen has no minor unit, gets rounded)}

\PY{+w}{        }\PY{n}{Currency}\PY{+w}{ }\PY{n}{eur}\PY{+w}{ }\PY{o}{=}\PY{+w}{ }\PY{n}{Currency}\PY{p}{.}\PY{n+na}{getInstance}\PY{p}{(}\PY{l+s}{\PYZdq{}}\PY{l+s}{EUR}\PY{l+s}{\PYZdq{}}\PY{p}{)}\PY{p}{;}
\PY{+w}{        }\PY{n}{System}\PY{p}{.}\PY{n+na}{out}\PY{p}{.}\PY{n+na}{println}\PY{p}{(}\PY{n}{eur}\PY{p}{.}\PY{n+na}{getSymbol}\PY{p}{(}\PY{n}{Locale}\PY{p}{.}\PY{n+na}{GERMANY}\PY{p}{)}\PY{p}{)}\PY{p}{;}\PY{+w}{ }\PY{c+c1}{// €}
\PY{+w}{        }\PY{n}{System}\PY{p}{.}\PY{n+na}{out}\PY{p}{.}\PY{n+na}{println}\PY{p}{(}\PY{n}{eur}\PY{p}{.}\PY{n+na}{getDefaultFractionDigits}\PY{p}{(}\PY{p}{)}\PY{p}{)}\PY{p}{;}\PY{+w}{  }\PY{c+c1}{// 2}
\PY{+w}{    }\PY{p}{\PYZcb{}}
\PY{p}{\PYZcb{}}
\end{Verbatim}
\end{codebox}

\subsection[Collation with Collator]{Collation with \code{Collator}}
\label{h:9:collation-with-collator}
Sorting text alphabetically is locale-dependent. A plain \code{String.\allowbreak{}compareTo} sorts by raw UTF-16 code unit values, which does not match human alphabetical order in most languages. \code{Collator} sorts the way a native speaker of a given locale would expect.

\begin{codebox}
\begin{Verbatim}[commandchars=\\\{\}]
\PY{k+kn}{import}\PY{+w}{ }\PY{n+nn}{java.text.Collator}\PY{p}{;}
\PY{k+kn}{import}\PY{+w}{ }\PY{n+nn}{java.util.Arrays}\PY{p}{;}
\PY{k+kn}{import}\PY{+w}{ }\PY{n+nn}{java.util.Locale}\PY{p}{;}

\PY{k+kd}{public}\PY{+w}{ }\PY{k+kd}{class} \PY{n+nc}{CollatorDemo}\PY{+w}{ }\PY{p}{\PYZob{}}
\PY{+w}{    }\PY{k+kd}{public}\PY{+w}{ }\PY{k+kd}{static}\PY{+w}{ }\PY{k+kt}{void}\PY{+w}{ }\PY{n+nf}{main}\PY{p}{(}\PY{n}{String}\PY{o}{[}\PY{o}{]}\PY{+w}{ }\PY{n}{args}\PY{p}{)}\PY{+w}{ }\PY{p}{\PYZob{}}
\PY{+w}{        }\PY{n}{String}\PY{o}{[}\PY{o}{]}\PY{+w}{ }\PY{n}{words}\PY{+w}{ }\PY{o}{=}\PY{+w}{ }\PY{p}{\PYZob{}}\PY{l+s}{\PYZdq{}}\PY{l+s}{Apfel}\PY{l+s}{\PYZdq{}}\PY{p}{,}\PY{+w}{ }\PY{l+s}{\PYZdq{}}\PY{l+s}{Äpfel}\PY{l+s}{\PYZdq{}}\PY{p}{,}\PY{+w}{ }\PY{l+s}{\PYZdq{}}\PY{l+s}{banane}\PY{l+s}{\PYZdq{}}\PY{p}{,}\PY{+w}{ }\PY{l+s}{\PYZdq{}}\PY{l+s}{Banane}\PY{l+s}{\PYZdq{}}\PY{p}{\PYZcb{}}\PY{p}{;}\PY{+w}{ }\PY{c+c1}{// Ä is U+00C4}

\PY{+w}{        }\PY{n}{String}\PY{o}{[}\PY{o}{]}\PY{+w}{ }\PY{n}{naturalSort}\PY{+w}{ }\PY{o}{=}\PY{+w}{ }\PY{n}{words}\PY{p}{.}\PY{n+na}{clone}\PY{p}{(}\PY{p}{)}\PY{p}{;}
\PY{+w}{        }\PY{n}{Arrays}\PY{p}{.}\PY{n+na}{sort}\PY{p}{(}\PY{n}{naturalSort}\PY{p}{)}\PY{p}{;}\PY{+w}{ }\PY{c+c1}{// plain code\PYZhy{}unit comparison}

\PY{+w}{        }\PY{n}{String}\PY{o}{[}\PY{o}{]}\PY{+w}{ }\PY{n}{germanSort}\PY{+w}{ }\PY{o}{=}\PY{+w}{ }\PY{n}{words}\PY{p}{.}\PY{n+na}{clone}\PY{p}{(}\PY{p}{)}\PY{p}{;}
\PY{+w}{        }\PY{n}{Collator}\PY{+w}{ }\PY{n}{germanCollator}\PY{+w}{ }\PY{o}{=}\PY{+w}{ }\PY{n}{Collator}\PY{p}{.}\PY{n+na}{getInstance}\PY{p}{(}\PY{n}{Locale}\PY{p}{.}\PY{n+na}{GERMANY}\PY{p}{)}\PY{p}{;}
\PY{+w}{        }\PY{n}{Arrays}\PY{p}{.}\PY{n+na}{sort}\PY{p}{(}\PY{n}{germanSort}\PY{p}{,}\PY{+w}{ }\PY{n}{germanCollator}\PY{p}{)}\PY{p}{;}

\PY{+w}{        }\PY{n}{System}\PY{p}{.}\PY{n+na}{out}\PY{p}{.}\PY{n+na}{println}\PY{p}{(}\PY{n}{Arrays}\PY{p}{.}\PY{n+na}{toString}\PY{p}{(}\PY{n}{naturalSort}\PY{p}{)}\PY{p}{)}\PY{p}{;}
\PY{+w}{        }\PY{c+c1}{// Output: [Apfel, Banane, banane, Äpfel]  \PYZlt{}\PYZhy{}\PYZhy{} uppercase/accents sort oddly by raw code points}
\PY{+w}{        }\PY{n}{System}\PY{p}{.}\PY{n+na}{out}\PY{p}{.}\PY{n+na}{println}\PY{p}{(}\PY{n}{Arrays}\PY{p}{.}\PY{n+na}{toString}\PY{p}{(}\PY{n}{germanSort}\PY{p}{)}\PY{p}{)}\PY{p}{;}
\PY{+w}{        }\PY{c+c1}{// Output: [Apfel, Äpfel, banane, Banane]  \PYZlt{}\PYZhy{}\PYZhy{} matches how a German speaker expects it sorted}
\PY{+w}{    }\PY{p}{\PYZcb{}}
\PY{p}{\PYZcb{}}
\end{Verbatim}
\end{codebox}

\subsection[String.format with a Locale]{\code{String.\allowbreak{}format} with a \code{Locale}}
\label{h:9:stringformat-with-a-locale}
\code{String.\allowbreak{}format} (and \code{Formatter}/\code{printf}) is locale-sensitive for things like decimal points and digit grouping. Always pass an explicit \code{Locale} in server-side or shared code — otherwise it silently uses the JVM’s default locale, which can differ between environments.

\begin{codebox}
\begin{Verbatim}[commandchars=\\\{\}]
\PY{k+kn}{import}\PY{+w}{ }\PY{n+nn}{java.util.Locale}\PY{p}{;}

\PY{k+kd}{public}\PY{+w}{ }\PY{k+kd}{class} \PY{n+nc}{StringFormatLocaleDemo}\PY{+w}{ }\PY{p}{\PYZob{}}
\PY{+w}{    }\PY{k+kd}{public}\PY{+w}{ }\PY{k+kd}{static}\PY{+w}{ }\PY{k+kt}{void}\PY{+w}{ }\PY{n+nf}{main}\PY{p}{(}\PY{n}{String}\PY{o}{[}\PY{o}{]}\PY{+w}{ }\PY{n}{args}\PY{p}{)}\PY{+w}{ }\PY{p}{\PYZob{}}
\PY{+w}{        }\PY{k+kt}{double}\PY{+w}{ }\PY{n}{price}\PY{+w}{ }\PY{o}{=}\PY{+w}{ }\PY{l+m+mf}{1234.5}\PY{p}{;}

\PY{+w}{        }\PY{n}{String}\PY{+w}{ }\PY{n}{defaultResult}\PY{+w}{ }\PY{o}{=}\PY{+w}{ }\PY{n}{String}\PY{p}{.}\PY{n+na}{format}\PY{p}{(}\PY{l+s}{\PYZdq{}}\PY{l+s}{\PYZpc{},.2f}\PY{l+s}{\PYZdq{}}\PY{p}{,}\PY{+w}{ }\PY{n}{price}\PY{p}{)}\PY{p}{;}\PY{+w}{      }\PY{c+c1}{// uses JVM default locale \PYZhy{} risky!}
\PY{+w}{        }\PY{n}{String}\PY{+w}{ }\PY{n}{usResult}\PY{+w}{ }\PY{o}{=}\PY{+w}{ }\PY{n}{String}\PY{p}{.}\PY{n+na}{format}\PY{p}{(}\PY{n}{Locale}\PY{p}{.}\PY{n+na}{US}\PY{p}{,}\PY{+w}{ }\PY{l+s}{\PYZdq{}}\PY{l+s}{\PYZpc{},.2f}\PY{l+s}{\PYZdq{}}\PY{p}{,}\PY{+w}{ }\PY{n}{price}\PY{p}{)}\PY{p}{;}
\PY{+w}{        }\PY{n}{String}\PY{+w}{ }\PY{n}{deResult}\PY{+w}{ }\PY{o}{=}\PY{+w}{ }\PY{n}{String}\PY{p}{.}\PY{n+na}{format}\PY{p}{(}\PY{n}{Locale}\PY{p}{.}\PY{n+na}{GERMANY}\PY{p}{,}\PY{+w}{ }\PY{l+s}{\PYZdq{}}\PY{l+s}{\PYZpc{},.2f}\PY{l+s}{\PYZdq{}}\PY{p}{,}\PY{+w}{ }\PY{n}{price}\PY{p}{)}\PY{p}{;}

\PY{+w}{        }\PY{n}{System}\PY{p}{.}\PY{n+na}{out}\PY{p}{.}\PY{n+na}{println}\PY{p}{(}\PY{n}{usResult}\PY{p}{)}\PY{p}{;}\PY{+w}{ }\PY{c+c1}{// 1,234.50}
\PY{+w}{        }\PY{n}{System}\PY{p}{.}\PY{n+na}{out}\PY{p}{.}\PY{n+na}{println}\PY{p}{(}\PY{n}{deResult}\PY{p}{)}\PY{p}{;}\PY{+w}{ }\PY{c+c1}{// 1.234,50}
\PY{+w}{        }\PY{c+c1}{// \PYZdq{}defaultResult\PYZdq{} could be either of the above depending on the machine\PYZsq{}s default locale \PYZhy{}}
\PY{+w}{        }\PY{c+c1}{// this is a classic source of \PYZdq{}works on my machine\PYZdq{} bugs in server\PYZhy{}side code.}
\PY{+w}{    }\PY{p}{\PYZcb{}}
\PY{p}{\PYZcb{}}
\end{Verbatim}
\end{codebox}

\subsection[toUpperCase() and the Turkish dotless-i problem]{\code{toUpperCase()} and the Turkish dotless-i problem}
\label{h:9:touppercase-and-the-turkish-dotless-i-problem}
Calling \code{toUpperCase()} or \code{toLowerCase()} without a \code{Locale} uses the JVM’s default locale. In Turkish (and Azerbaijani), the letter \code{i} uppercases to \code{İ} (dotted capital I), not the ASCII \code{I} most code expects — and lowercase \code{I} becomes \code{ı} (dotless). This breaks case-insensitive comparisons against constants like \code{\textquotedbl{}ID\textquotedbl{}} or \code{\textquotedbl{}HTTP\textquotedbl{}} when the default locale is Turkish.

\begin{codebox}
\begin{Verbatim}[commandchars=\\\{\}]
\PY{k+kn}{import}\PY{+w}{ }\PY{n+nn}{java.util.Locale}\PY{p}{;}

\PY{k+kd}{public}\PY{+w}{ }\PY{k+kd}{class} \PY{n+nc}{TurkishLocaleBugDemo}\PY{+w}{ }\PY{p}{\PYZob{}}
\PY{+w}{    }\PY{k+kd}{public}\PY{+w}{ }\PY{k+kd}{static}\PY{+w}{ }\PY{k+kt}{void}\PY{+w}{ }\PY{n+nf}{main}\PY{p}{(}\PY{n}{String}\PY{o}{[}\PY{o}{]}\PY{+w}{ }\PY{n}{args}\PY{p}{)}\PY{+w}{ }\PY{p}{\PYZob{}}
\PY{+w}{        }\PY{n}{String}\PY{+w}{ }\PY{n}{input}\PY{+w}{ }\PY{o}{=}\PY{+w}{ }\PY{l+s}{\PYZdq{}}\PY{l+s}{id}\PY{l+s}{\PYZdq{}}\PY{p}{;}

\PY{+w}{        }\PY{n}{String}\PY{+w}{ }\PY{n}{upperTurkish}\PY{+w}{ }\PY{o}{=}\PY{+w}{ }\PY{n}{input}\PY{p}{.}\PY{n+na}{toUpperCase}\PY{p}{(}\PY{k}{new}\PY{+w}{ }\PY{n}{Locale}\PY{p}{(}\PY{l+s}{\PYZdq{}}\PY{l+s}{tr}\PY{l+s}{\PYZdq{}}\PY{p}{,}\PY{+w}{ }\PY{l+s}{\PYZdq{}}\PY{l+s}{TR}\PY{l+s}{\PYZdq{}}\PY{p}{)}\PY{p}{)}\PY{p}{;}
\PY{+w}{        }\PY{n}{String}\PY{+w}{ }\PY{n}{upperRoot}\PY{+w}{ }\PY{o}{=}\PY{+w}{ }\PY{n}{input}\PY{p}{.}\PY{n+na}{toUpperCase}\PY{p}{(}\PY{n}{Locale}\PY{p}{.}\PY{n+na}{ROOT}\PY{p}{)}\PY{p}{;}

\PY{+w}{        }\PY{n}{System}\PY{p}{.}\PY{n+na}{out}\PY{p}{.}\PY{n+na}{println}\PY{p}{(}\PY{n}{upperTurkish}\PY{p}{)}\PY{p}{;}\PY{+w}{ }\PY{c+c1}{// İD  (dotted capital İ, not plain \PYZdq{}I\PYZdq{})}
\PY{+w}{        }\PY{n}{System}\PY{p}{.}\PY{n+na}{out}\PY{p}{.}\PY{n+na}{println}\PY{p}{(}\PY{n}{upperRoot}\PY{p}{)}\PY{p}{;}\PY{+w}{    }\PY{c+c1}{// ID  (plain ASCII, as expected)}

\PY{+w}{        }\PY{n}{System}\PY{p}{.}\PY{n+na}{out}\PY{p}{.}\PY{n+na}{println}\PY{p}{(}\PY{n}{upperTurkish}\PY{p}{.}\PY{n+na}{equals}\PY{p}{(}\PY{l+s}{\PYZdq{}}\PY{l+s}{ID}\PY{l+s}{\PYZdq{}}\PY{p}{)}\PY{p}{)}\PY{p}{;}\PY{+w}{ }\PY{c+c1}{// false \PYZhy{} a real bug if this runs on a Turkish\PYZhy{}locale server!}
\PY{+w}{        }\PY{n}{System}\PY{p}{.}\PY{n+na}{out}\PY{p}{.}\PY{n+na}{println}\PY{p}{(}\PY{n}{upperRoot}\PY{p}{.}\PY{n+na}{equals}\PY{p}{(}\PY{l+s}{\PYZdq{}}\PY{l+s}{ID}\PY{l+s}{\PYZdq{}}\PY{p}{)}\PY{p}{)}\PY{p}{;}\PY{+w}{     }\PY{c+c1}{// true}

\PY{+w}{        }\PY{c+c1}{// Rule for code review: for internal/programmatic comparisons (protocol keywords,}
\PY{+w}{        }\PY{c+c1}{// enum names, HTTP headers), always use toUpperCase(Locale.ROOT) / toLowerCase(Locale.ROOT),}
\PY{+w}{        }\PY{c+c1}{// or better, String.equalsIgnoreCase() when you just need a comparison.}
\PY{+w}{    }\PY{p}{\PYZcb{}}
\PY{p}{\PYZcb{}}
\end{Verbatim}
\end{codebox}

\subsection[Character encoding and UTF-8 as the default charset]{Character encoding and UTF-8 as the default charset}
\label{h:9:character-encoding-and-utf-8-as-the-default-charset}
Historically, \code{String.\allowbreak{}getBytes()}, \code{new~\allowbreak{}String(\allowbreak{}bytes)}, \code{FileReader}, and similar APIs used the \textbf{platform default charset} if none was specified — which could be UTF-8 on Linux but something else (like Windows-1252) on Windows, causing garbled text (“mojibake”) when code moved between environments. Since \textbf{JDK 18} (\href{https://openjdk.org/jeps/400}{JEP 400}), the JVM’s default charset is \textbf{always UTF-8}, regardless of OS or locale, unless explicitly overridden with the \code{file.\allowbreak{}encoding} system property.

\begin{codebox}
\begin{Verbatim}[commandchars=\\\{\}]
\PY{k+kn}{import}\PY{+w}{ }\PY{n+nn}{java.nio.charset.Charset}\PY{p}{;}
\PY{k+kn}{import}\PY{+w}{ }\PY{n+nn}{java.nio.charset.StandardCharsets}\PY{p}{;}

\PY{k+kd}{public}\PY{+w}{ }\PY{k+kd}{class} \PY{n+nc}{DefaultCharsetDemo}\PY{+w}{ }\PY{p}{\PYZob{}}
\PY{+w}{    }\PY{k+kd}{public}\PY{+w}{ }\PY{k+kd}{static}\PY{+w}{ }\PY{k+kt}{void}\PY{+w}{ }\PY{n+nf}{main}\PY{p}{(}\PY{n}{String}\PY{o}{[}\PY{o}{]}\PY{+w}{ }\PY{n}{args}\PY{p}{)}\PY{+w}{ }\PY{p}{\PYZob{}}
\PY{+w}{        }\PY{n}{System}\PY{p}{.}\PY{n+na}{out}\PY{p}{.}\PY{n+na}{println}\PY{p}{(}\PY{n}{Charset}\PY{p}{.}\PY{n+na}{defaultCharset}\PY{p}{(}\PY{p}{)}\PY{p}{)}\PY{p}{;}
\PY{+w}{        }\PY{c+c1}{// Output (JDK 18+): UTF\PYZhy{}8   (was platform\PYZhy{}dependent before JDK 18)}

\PY{+w}{        }\PY{n}{String}\PY{+w}{ }\PY{n}{text}\PY{+w}{ }\PY{o}{=}\PY{+w}{ }\PY{l+s}{\PYZdq{}}\PY{l+s}{café ☕}\PY{l+s}{\PYZdq{}}\PY{p}{;}
\PY{+w}{        }\PY{k+kt}{byte}\PY{o}{[}\PY{o}{]}\PY{+w}{ }\PY{n}{defaultBytes}\PY{+w}{ }\PY{o}{=}\PY{+w}{ }\PY{n}{text}\PY{p}{.}\PY{n+na}{getBytes}\PY{p}{(}\PY{p}{)}\PY{p}{;}\PY{+w}{                 }\PY{c+c1}{// uses default charset \PYZhy{} UTF\PYZhy{}8 on JDK 18+}
\PY{+w}{        }\PY{k+kt}{byte}\PY{o}{[}\PY{o}{]}\PY{+w}{ }\PY{n}{explicitBytes}\PY{+w}{ }\PY{o}{=}\PY{+w}{ }\PY{n}{text}\PY{p}{.}\PY{n+na}{getBytes}\PY{p}{(}\PY{n}{StandardCharsets}\PY{p}{.}\PY{n+na}{UTF\PYZus{}8}\PY{p}{)}\PY{p}{;}\PY{+w}{ }\PY{c+c1}{// always explicit, safest for review}

\PY{+w}{        }\PY{n}{System}\PY{p}{.}\PY{n+na}{out}\PY{p}{.}\PY{n+na}{println}\PY{p}{(}\PY{n}{defaultBytes}\PY{p}{.}\PY{n+na}{length}\PY{p}{)}\PY{p}{;}\PY{+w}{  }\PY{c+c1}{// 8 (matches UTF\PYZhy{}8 encoding)}
\PY{+w}{        }\PY{n}{System}\PY{p}{.}\PY{n+na}{out}\PY{p}{.}\PY{n+na}{println}\PY{p}{(}\PY{n}{explicitBytes}\PY{p}{.}\PY{n+na}{length}\PY{p}{)}\PY{p}{;}\PY{+w}{ }\PY{c+c1}{// 8}

\PY{+w}{        }\PY{c+c1}{// Best practice even on JDK 18+: always pass an explicit Charset for I/O boundaries}
\PY{+w}{        }\PY{c+c1}{// (files, network, serialization) so behavior never depends on JVM defaults or \PYZhy{}Dfile.encoding.}
\PY{+w}{    }\PY{p}{\PYZcb{}}
\PY{p}{\PYZcb{}}
\end{Verbatim}
\end{codebox}

\subsection[Normalizer for Unicode]{\code{Normalizer} for Unicode}
\label{h:9:normalizer-for-unicode}
The same visible character can sometimes be represented by different sequences of Unicode code points — for example, “é” can be one precomposed code point (\code{U+\allowbreak{}00E9}) or an “e” followed by a combining accent (\code{U+\allowbreak{}0065~\allowbreak{}U+\allowbreak{}0301}). These compare as unequal with plain \code{.\allowbreak{}equals()} even though they look identical. \code{java.\allowbreak{}text.\allowbreak{}Normalizer} converts text into a consistent form before comparing or storing it.

\begin{codebox}
\begin{Verbatim}[commandchars=\\\{\}]
\PY{k+kn}{import}\PY{+w}{ }\PY{n+nn}{java.text.Normalizer}\PY{p}{;}

\PY{k+kd}{public}\PY{+w}{ }\PY{k+kd}{class} \PY{n+nc}{NormalizerDemo}\PY{+w}{ }\PY{p}{\PYZob{}}
\PY{+w}{    }\PY{k+kd}{public}\PY{+w}{ }\PY{k+kd}{static}\PY{+w}{ }\PY{k+kt}{void}\PY{+w}{ }\PY{n+nf}{main}\PY{p}{(}\PY{n}{String}\PY{o}{[}\PY{o}{]}\PY{+w}{ }\PY{n}{args}\PY{p}{)}\PY{+w}{ }\PY{p}{\PYZob{}}
\PY{+w}{        }\PY{n}{String}\PY{+w}{ }\PY{n}{precomposed}\PY{+w}{ }\PY{o}{=}\PY{+w}{ }\PY{l+s}{\PYZdq{}}\PY{l+s}{café}\PY{l+s}{\PYZdq{}}\PY{p}{;}\PY{+w}{        }\PY{c+c1}{// \PYZdq{}café\PYZdq{} using single code point U+00E9}
\PY{+w}{        }\PY{n}{String}\PY{+w}{ }\PY{n}{decomposed}\PY{+w}{ }\PY{o}{=}\PY{+w}{ }\PY{l+s}{\PYZdq{}}\PY{l+s}{café}\PY{l+s}{\PYZdq{}}\PY{p}{;}\PY{+w}{        }\PY{c+c1}{// \PYZdq{}café\PYZdq{} using \PYZsq{}e\PYZsq{} + combining acute accent U+0301}

\PY{+w}{        }\PY{n}{System}\PY{p}{.}\PY{n+na}{out}\PY{p}{.}\PY{n+na}{println}\PY{p}{(}\PY{n}{precomposed}\PY{p}{.}\PY{n+na}{equals}\PY{p}{(}\PY{n}{decomposed}\PY{p}{)}\PY{p}{)}\PY{p}{;}\PY{+w}{ }\PY{c+c1}{// false \PYZhy{} look identical, but different bytes!}

\PY{+w}{        }\PY{n}{String}\PY{+w}{ }\PY{n}{normalizedA}\PY{+w}{ }\PY{o}{=}\PY{+w}{ }\PY{n}{Normalizer}\PY{p}{.}\PY{n+na}{normalize}\PY{p}{(}\PY{n}{precomposed}\PY{p}{,}\PY{+w}{ }\PY{n}{Normalizer}\PY{p}{.}\PY{n+na}{Form}\PY{p}{.}\PY{n+na}{NFC}\PY{p}{)}\PY{p}{;}
\PY{+w}{        }\PY{n}{String}\PY{+w}{ }\PY{n}{normalizedB}\PY{+w}{ }\PY{o}{=}\PY{+w}{ }\PY{n}{Normalizer}\PY{p}{.}\PY{n+na}{normalize}\PY{p}{(}\PY{n}{decomposed}\PY{p}{,}\PY{+w}{ }\PY{n}{Normalizer}\PY{p}{.}\PY{n+na}{Form}\PY{p}{.}\PY{n+na}{NFC}\PY{p}{)}\PY{p}{;}

\PY{+w}{        }\PY{n}{System}\PY{p}{.}\PY{n+na}{out}\PY{p}{.}\PY{n+na}{println}\PY{p}{(}\PY{n}{normalizedA}\PY{p}{.}\PY{n+na}{equals}\PY{p}{(}\PY{n}{normalizedB}\PY{p}{)}\PY{p}{)}\PY{p}{;}\PY{+w}{ }\PY{c+c1}{// true \PYZhy{} now safe to compare}

\PY{+w}{        }\PY{c+c1}{// NFC = canonical composition (preferred for storage/display)}
\PY{+w}{        }\PY{c+c1}{// NFD = canonical decomposition (useful for stripping accents, searching)}
\PY{+w}{        }\PY{n}{String}\PY{+w}{ }\PY{n}{stripped}\PY{+w}{ }\PY{o}{=}\PY{+w}{ }\PY{n}{Normalizer}\PY{p}{.}\PY{n+na}{normalize}\PY{p}{(}\PY{n}{precomposed}\PY{p}{,}\PY{+w}{ }\PY{n}{Normalizer}\PY{p}{.}\PY{n+na}{Form}\PY{p}{.}\PY{n+na}{NFD}\PY{p}{)}
\PY{+w}{                }\PY{p}{.}\PY{n+na}{replaceAll}\PY{p}{(}\PY{l+s}{\PYZdq{}}\PY{l+s}{\PYZbs{}\PYZbs{}}\PY{l+s}{p\PYZob{}M\PYZcb{}}\PY{l+s}{\PYZdq{}}\PY{p}{,}\PY{+w}{ }\PY{l+s}{\PYZdq{}}\PY{l+s}{\PYZdq{}}\PY{p}{)}\PY{p}{;}\PY{+w}{ }\PY{c+c1}{// remove combining marks}
\PY{+w}{        }\PY{n}{System}\PY{p}{.}\PY{n+na}{out}\PY{p}{.}\PY{n+na}{println}\PY{p}{(}\PY{n}{stripped}\PY{p}{)}\PY{p}{;}\PY{+w}{ }\PY{c+c1}{// cafe}
\PY{+w}{    }\PY{p}{\PYZcb{}}
\PY{p}{\PYZcb{}}
\end{Verbatim}
\end{codebox}

\subsection[Right-to-left text, code points, and grapheme clusters]{Right-to-left text, code points, and grapheme clusters}
\label{h:9:right-to-left-text-code-points-and-grapheme-clusters}
Some languages (Arabic, Hebrew) are written right-to-left (RTL). Java strings store text left-to-right in memory regardless of visual direction — rendering direction is a display concern (handled by UI toolkits, not the \code{String} class itself). A separate, very common bug source: \code{String.\allowbreak{}length()} counts UTF-16 \emph{code units}, not visible characters. Emoji, some CJK (Chinese/Japanese/Korean) characters, and combining accents can each take more than one code unit, and a single visible “character” (grapheme cluster) can be made of multiple code points.

\begin{codebox}
\begin{Verbatim}[commandchars=\\\{\}]
\PY{k+kn}{import}\PY{+w}{ }\PY{n+nn}{java.text.BreakIterator}\PY{p}{;}
\PY{k+kn}{import}\PY{+w}{ }\PY{n+nn}{java.util.Locale}\PY{p}{;}

\PY{k+kd}{public}\PY{+w}{ }\PY{k+kd}{class} \PY{n+nc}{TextLengthDemo}\PY{+w}{ }\PY{p}{\PYZob{}}
\PY{+w}{    }\PY{k+kd}{public}\PY{+w}{ }\PY{k+kd}{static}\PY{+w}{ }\PY{k+kt}{void}\PY{+w}{ }\PY{n+nf}{main}\PY{p}{(}\PY{n}{String}\PY{o}{[}\PY{o}{]}\PY{+w}{ }\PY{n}{args}\PY{p}{)}\PY{+w}{ }\PY{p}{\PYZob{}}
\PY{+w}{        }\PY{n}{String}\PY{+w}{ }\PY{n}{flagEmoji}\PY{+w}{ }\PY{o}{=}\PY{+w}{ }\PY{l+s}{\PYZdq{}}\PY{l+s}{\PYZbs{}}\PY{l+s}{uD83C}\PY{l+s}{\PYZbs{}}\PY{l+s}{uDDF3}\PY{l+s}{\PYZbs{}}\PY{l+s}{uD83C}\PY{l+s}{\PYZbs{}}\PY{l+s}{uDDF1}\PY{l+s}{\PYZdq{}}\PY{p}{;}\PY{+w}{ }\PY{c+c1}{// the Netherlands flag emoji, one visible glyph}
\PY{+w}{        }\PY{n}{String}\PY{+w}{ }\PY{n}{familyEmoji}\PY{+w}{ }\PY{o}{=}\PY{+w}{ }\PY{l+s}{\PYZdq{}}\PY{l+s}{\PYZbs{}}\PY{l+s}{uD83D}\PY{l+s}{\PYZbs{}}\PY{l+s}{uDC68}\PY{l+s}{\PYZbs{}}\PY{l+s}{u200D}\PY{l+s}{\PYZbs{}}\PY{l+s}{uD83D}\PY{l+s}{\PYZbs{}}\PY{l+s}{uDC69}\PY{l+s}{\PYZbs{}}\PY{l+s}{u200D}\PY{l+s}{\PYZbs{}}\PY{l+s}{uD83D}\PY{l+s}{\PYZbs{}}\PY{l+s}{uDC67}\PY{l+s}{\PYZdq{}}\PY{p}{;}\PY{+w}{ }\PY{c+c1}{// family emoji, one glyph}

\PY{+w}{        }\PY{n}{System}\PY{p}{.}\PY{n+na}{out}\PY{p}{.}\PY{n+na}{println}\PY{p}{(}\PY{n}{flagEmoji}\PY{p}{.}\PY{n+na}{length}\PY{p}{(}\PY{p}{)}\PY{p}{)}\PY{p}{;}\PY{+w}{              }\PY{c+c1}{// 4 (UTF\PYZhy{}16 code units, NOT 1!)}
\PY{+w}{        }\PY{n}{System}\PY{p}{.}\PY{n+na}{out}\PY{p}{.}\PY{n+na}{println}\PY{p}{(}\PY{n}{flagEmoji}\PY{p}{.}\PY{n+na}{codePointCount}\PY{p}{(}\PY{l+m+mi}{0}\PY{p}{,}\PY{+w}{ }\PY{n}{flagEmoji}\PY{p}{.}\PY{n+na}{length}\PY{p}{(}\PY{p}{)}\PY{p}{)}\PY{p}{)}\PY{p}{;}\PY{+w}{ }\PY{c+c1}{// 2 (code points, still not 1)}

\PY{+w}{        }\PY{c+c1}{// BreakIterator finds actual user\PYZhy{}perceived \PYZdq{}characters\PYZdq{} (grapheme clusters)}
\PY{+w}{        }\PY{n}{BreakIterator}\PY{+w}{ }\PY{n}{iterator}\PY{+w}{ }\PY{o}{=}\PY{+w}{ }\PY{n}{BreakIterator}\PY{p}{.}\PY{n+na}{getCharacterInstance}\PY{p}{(}\PY{n}{Locale}\PY{p}{.}\PY{n+na}{ROOT}\PY{p}{)}\PY{p}{;}
\PY{+w}{        }\PY{n}{iterator}\PY{p}{.}\PY{n+na}{setText}\PY{p}{(}\PY{n}{familyEmoji}\PY{p}{)}\PY{p}{;}
\PY{+w}{        }\PY{k+kt}{int}\PY{+w}{ }\PY{n}{graphemeCount}\PY{+w}{ }\PY{o}{=}\PY{+w}{ }\PY{l+m+mi}{0}\PY{p}{;}
\PY{+w}{        }\PY{k+kt}{int}\PY{+w}{ }\PY{n}{boundary}\PY{+w}{ }\PY{o}{=}\PY{+w}{ }\PY{n}{iterator}\PY{p}{.}\PY{n+na}{first}\PY{p}{(}\PY{p}{)}\PY{p}{;}
\PY{+w}{        }\PY{k}{while}\PY{+w}{ }\PY{p}{(}\PY{n}{boundary}\PY{+w}{ }\PY{o}{!}\PY{o}{=}\PY{+w}{ }\PY{n}{BreakIterator}\PY{p}{.}\PY{n+na}{DONE}\PY{p}{)}\PY{+w}{ }\PY{p}{\PYZob{}}
\PY{+w}{            }\PY{k+kt}{int}\PY{+w}{ }\PY{n}{next}\PY{+w}{ }\PY{o}{=}\PY{+w}{ }\PY{n}{iterator}\PY{p}{.}\PY{n+na}{next}\PY{p}{(}\PY{p}{)}\PY{p}{;}
\PY{+w}{            }\PY{k}{if}\PY{+w}{ }\PY{p}{(}\PY{n}{next}\PY{+w}{ }\PY{o}{!}\PY{o}{=}\PY{+w}{ }\PY{n}{BreakIterator}\PY{p}{.}\PY{n+na}{DONE}\PY{p}{)}\PY{+w}{ }\PY{n}{graphemeCount}\PY{o}{+}\PY{o}{+}\PY{p}{;}
\PY{+w}{            }\PY{n}{boundary}\PY{+w}{ }\PY{o}{=}\PY{+w}{ }\PY{n}{next}\PY{p}{;}
\PY{+w}{        }\PY{p}{\PYZcb{}}
\PY{+w}{        }\PY{n}{System}\PY{p}{.}\PY{n+na}{out}\PY{p}{.}\PY{n+na}{println}\PY{p}{(}\PY{n}{graphemeCount}\PY{p}{)}\PY{p}{;}\PY{+w}{ }\PY{c+c1}{// 1 (one visible \PYZdq{}character\PYZdq{} made of several code points)}

\PY{+w}{        }\PY{c+c1}{// Rule for code review: never assume String.length() == \PYZdq{}number of visible characters\PYZdq{}}
\PY{+w}{        }\PY{c+c1}{// when the input may contain emoji, combining marks, or characters outside the Basic}
\PY{+w}{        }\PY{c+c1}{// Multilingual Plane. Use codePointCount() or BreakIterator depending on what you truly need.}
\PY{+w}{    }\PY{p}{\PYZcb{}}
\PY{p}{\PYZcb{}}
\end{Verbatim}
\end{codebox}

\section[Localization (l10n)]{Localization (l10n)}
\label{h:9:localization-l10n}
\subsection[Organizing bundles and fallback rules]{Organizing bundles and fallback rules}
\label{h:9:organizing-bundles-and-fallback-rules}
A typical project has one base bundle plus one file per supported locale, sometimes split by variant (language, then language+country). \code{ResourceBundle} looks up a candidate locale chain from most to least specific, and falls back to the base bundle if nothing else matches.

\begin{codebox}
\begin{Verbatim}
Messages.properties          <- default / fallback (usually English)
Messages_en.properties       <- English (falls back to default automatically if same content)
Messages_en_GB.properties    <- British English overrides (e.g. "colour" vs "color")
Messages_de.properties       <- German
Messages_de_CH.properties    <- Swiss German overrides
Messages_fr.properties       <- French
\end{Verbatim}
\end{codebox}

\begin{codebox}
\begin{Verbatim}[commandchars=\\\{\}]
\PY{k+kn}{import}\PY{+w}{ }\PY{n+nn}{java.util.Locale}\PY{p}{;}
\PY{k+kn}{import}\PY{+w}{ }\PY{n+nn}{java.util.ResourceBundle}\PY{p}{;}

\PY{k+kd}{public}\PY{+w}{ }\PY{k+kd}{class} \PY{n+nc}{BundleFallbackDemo}\PY{+w}{ }\PY{p}{\PYZob{}}
\PY{+w}{    }\PY{k+kd}{public}\PY{+w}{ }\PY{k+kd}{static}\PY{+w}{ }\PY{k+kt}{void}\PY{+w}{ }\PY{n+nf}{main}\PY{p}{(}\PY{n}{String}\PY{o}{[}\PY{o}{]}\PY{+w}{ }\PY{n}{args}\PY{p}{)}\PY{+w}{ }\PY{p}{\PYZob{}}
\PY{+w}{        }\PY{c+c1}{// Requesting en\PYZus{}GB will look for: Messages\PYZus{}en\PYZus{}GB \PYZhy{}\PYZgt{} Messages\PYZus{}en \PYZhy{}\PYZgt{} Messages (base)}
\PY{+w}{        }\PY{n}{ResourceBundle}\PY{+w}{ }\PY{n}{bundle}\PY{+w}{ }\PY{o}{=}\PY{+w}{ }\PY{n}{ResourceBundle}\PY{p}{.}\PY{n+na}{getBundle}\PY{p}{(}\PY{l+s}{\PYZdq{}}\PY{l+s}{Messages}\PY{l+s}{\PYZdq{}}\PY{p}{,}\PY{+w}{ }\PY{k}{new}\PY{+w}{ }\PY{n}{Locale}\PY{p}{(}\PY{l+s}{\PYZdq{}}\PY{l+s}{en}\PY{l+s}{\PYZdq{}}\PY{p}{,}\PY{+w}{ }\PY{l+s}{\PYZdq{}}\PY{l+s}{GB}\PY{l+s}{\PYZdq{}}\PY{p}{)}\PY{p}{)}\PY{p}{;}
\PY{+w}{        }\PY{n}{System}\PY{p}{.}\PY{n+na}{out}\PY{p}{.}\PY{n+na}{println}\PY{p}{(}\PY{n}{bundle}\PY{p}{.}\PY{n+na}{getLocale}\PY{p}{(}\PY{p}{)}\PY{p}{)}\PY{p}{;}
\PY{+w}{        }\PY{c+c1}{// Output depends on which files actually exist; falls back gracefully instead of throwing.}

\PY{+w}{        }\PY{c+c1}{// Requesting a totally unsupported locale (e.g. Icelandic) still works,}
\PY{+w}{        }\PY{c+c1}{// it just falls all the way back to the base Messages.properties file.}
\PY{+w}{        }\PY{n}{ResourceBundle}\PY{+w}{ }\PY{n}{fallback}\PY{+w}{ }\PY{o}{=}\PY{+w}{ }\PY{n}{ResourceBundle}\PY{p}{.}\PY{n+na}{getBundle}\PY{p}{(}\PY{l+s}{\PYZdq{}}\PY{l+s}{Messages}\PY{l+s}{\PYZdq{}}\PY{p}{,}\PY{+w}{ }\PY{k}{new}\PY{+w}{ }\PY{n}{Locale}\PY{p}{(}\PY{l+s}{\PYZdq{}}\PY{l+s}{is}\PY{l+s}{\PYZdq{}}\PY{p}{,}\PY{+w}{ }\PY{l+s}{\PYZdq{}}\PY{l+s}{IS}\PY{l+s}{\PYZdq{}}\PY{p}{)}\PY{p}{)}\PY{p}{;}
\PY{+w}{        }\PY{n}{System}\PY{p}{.}\PY{n+na}{out}\PY{p}{.}\PY{n+na}{println}\PY{p}{(}\PY{n}{fallback}\PY{p}{.}\PY{n+na}{getBaseBundleName}\PY{p}{(}\PY{p}{)}\PY{p}{)}\PY{p}{;}\PY{+w}{ }\PY{c+c1}{// Messages}
\PY{+w}{    }\PY{p}{\PYZcb{}}
\PY{p}{\PYZcb{}}
\end{Verbatim}
\end{codebox}

{\tablefont
\begin{xltabular}{\linewidth}{@{}L{0.979}L{1.021}@{}}
\toprule
\rowcolor{tablehdr}
\thd{Requested locale} & \thd{Lookup order (most to least specific)} \\
\midrule
\endhead
\code{de\_\allowbreak{}CH} (German, Switzerland) & \code{Messages\_\allowbreak{}de\_\allowbreak{}CH} -\textgreater{} \code{Messages\_\allowbreak{}de} -\textgreater{} \code{Messages} \\
\code{en\_\allowbreak{}GB} (English, UK) & \code{Messages\_\allowbreak{}en\_\allowbreak{}GB} -\textgreater{} \code{Messages\_\allowbreak{}en} -\textgreater{} \code{Messages} \\
\code{is\_\allowbreak{}IS} (Icelandic, Iceland) - no matching files & \code{Messages\_\allowbreak{}is\_\allowbreak{}IS} -\textgreater{} \code{Messages\_\allowbreak{}is} -\textgreater{} \code{Messages} (base) \\
\bottomrule
\end{xltabular}
}

\subsection[Pseudo-localization]{Pseudo-localization}
\label{h:9:pseudo-localization}
Pseudo-localization is a testing technique: you generate a fake “translation” that expands text length, adds accented characters, and wraps strings with markers, without doing real translation. It quickly reveals UI bugs (truncated text, hardcoded string concatenation, untranslated strings) before real translators get involved.

\begin{codebox}
\begin{Verbatim}[commandchars=\\\{\}]
\PY{k+kn}{import}\PY{+w}{ }\PY{n+nn}{java.util.Locale}\PY{p}{;}
\PY{k+kn}{import}\PY{+w}{ }\PY{n+nn}{java.util.ResourceBundle}\PY{p}{;}

\PY{k+kd}{public}\PY{+w}{ }\PY{k+kd}{class} \PY{n+nc}{PseudoLocalizationDemo}\PY{+w}{ }\PY{p}{\PYZob{}}

\PY{+w}{    }\PY{c+c1}{// A crude pseudo\PYZhy{}localizer: wraps text and pads length by \PYZti{}40\PYZpc{}, which is}
\PY{+w}{    }\PY{c+c1}{// roughly how much German/French text tends to expand versus English.}
\PY{+w}{    }\PY{k+kd}{static}\PY{+w}{ }\PY{n}{String}\PY{+w}{ }\PY{n+nf}{pseudoLocalize}\PY{p}{(}\PY{n}{String}\PY{+w}{ }\PY{n}{text}\PY{p}{)}\PY{+w}{ }\PY{p}{\PYZob{}}
\PY{+w}{        }\PY{n}{String}\PY{+w}{ }\PY{n}{expanded}\PY{+w}{ }\PY{o}{=}\PY{+w}{ }\PY{n}{text}\PY{p}{.}\PY{n+na}{replace}\PY{p}{(}\PY{l+s}{\PYZdq{}}\PY{l+s}{a}\PY{l+s}{\PYZdq{}}\PY{p}{,}\PY{+w}{ }\PY{l+s}{\PYZdq{}}\PY{l+s}{àà}\PY{l+s}{\PYZdq{}}\PY{p}{)}\PY{p}{.}\PY{n+na}{replace}\PY{p}{(}\PY{l+s}{\PYZdq{}}\PY{l+s}{e}\PY{l+s}{\PYZdq{}}\PY{p}{,}\PY{+w}{ }\PY{l+s}{\PYZdq{}}\PY{l+s}{éé}\PY{l+s}{\PYZdq{}}\PY{p}{)}\PY{p}{.}\PY{n+na}{replace}\PY{p}{(}\PY{l+s}{\PYZdq{}}\PY{l+s}{o}\PY{l+s}{\PYZdq{}}\PY{p}{,}\PY{+w}{ }\PY{l+s}{\PYZdq{}}\PY{l+s}{òò}\PY{l+s}{\PYZdq{}}\PY{p}{)}\PY{p}{;}
\PY{+w}{        }\PY{k}{return}\PY{+w}{ }\PY{l+s}{\PYZdq{}}\PY{l+s}{[}\PY{l+s}{\PYZdq{}}\PY{+w}{ }\PY{o}{+}\PY{+w}{ }\PY{n}{expanded}\PY{+w}{ }\PY{o}{+}\PY{+w}{ }\PY{l+s}{\PYZdq{}}\PY{l+s}{]}\PY{l+s}{\PYZdq{}}\PY{p}{;}
\PY{+w}{    }\PY{p}{\PYZcb{}}

\PY{+w}{    }\PY{k+kd}{public}\PY{+w}{ }\PY{k+kd}{static}\PY{+w}{ }\PY{k+kt}{void}\PY{+w}{ }\PY{n+nf}{main}\PY{p}{(}\PY{n}{String}\PY{o}{[}\PY{o}{]}\PY{+w}{ }\PY{n}{args}\PY{p}{)}\PY{+w}{ }\PY{p}{\PYZob{}}
\PY{+w}{        }\PY{n}{ResourceBundle}\PY{+w}{ }\PY{n}{bundle}\PY{+w}{ }\PY{o}{=}\PY{+w}{ }\PY{n}{ResourceBundle}\PY{p}{.}\PY{n+na}{getBundle}\PY{p}{(}\PY{l+s}{\PYZdq{}}\PY{l+s}{Messages}\PY{l+s}{\PYZdq{}}\PY{p}{,}\PY{+w}{ }\PY{n}{Locale}\PY{p}{.}\PY{n+na}{ENGLISH}\PY{p}{)}\PY{p}{;}
\PY{+w}{        }\PY{n}{String}\PY{+w}{ }\PY{n}{original}\PY{+w}{ }\PY{o}{=}\PY{+w}{ }\PY{n}{bundle}\PY{p}{.}\PY{n+na}{getString}\PY{p}{(}\PY{l+s}{\PYZdq{}}\PY{l+s}{greeting}\PY{l+s}{\PYZdq{}}\PY{p}{)}\PY{p}{;}\PY{+w}{ }\PY{c+c1}{// \PYZdq{}Hello, \PYZob{}0\PYZcb{}!\PYZdq{}}

\PY{+w}{        }\PY{n}{System}\PY{p}{.}\PY{n+na}{out}\PY{p}{.}\PY{n+na}{println}\PY{p}{(}\PY{n}{pseudoLocalize}\PY{p}{(}\PY{n}{original}\PY{p}{)}\PY{p}{)}\PY{p}{;}
\PY{+w}{        }\PY{c+c1}{// Output: [Héééllòò, \PYZob{}0\PYZcb{}!]}

\PY{+w}{        }\PY{c+c1}{// If the UI truncates this or breaks layout, it will likely break for real}
\PY{+w}{        }\PY{c+c1}{// translations too (German and Finnish strings are often 30\PYZhy{}50\PYZpc{} longer than English).}
\PY{+w}{    }\PY{p}{\PYZcb{}}
\PY{p}{\PYZcb{}}
\end{Verbatim}
\end{codebox}

\subsection[Testing with -Duser.language]{Testing with \code{-\allowbreak{}Duser.\allowbreak{}language}}
\label{h:9:testing-with--duserlanguage}
You can override the JVM’s default locale from the command line without changing any code, which is very useful for manual and automated locale testing.

\begin{codebox}
\begin{Verbatim}[commandchars=\\\{\}]
\PY{k+kn}{import}\PY{+w}{ }\PY{n+nn}{java.util.Locale}\PY{p}{;}

\PY{k+kd}{public}\PY{+w}{ }\PY{k+kd}{class} \PY{n+nc}{DefaultLocaleDemo}\PY{+w}{ }\PY{p}{\PYZob{}}
\PY{+w}{    }\PY{k+kd}{public}\PY{+w}{ }\PY{k+kd}{static}\PY{+w}{ }\PY{k+kt}{void}\PY{+w}{ }\PY{n+nf}{main}\PY{p}{(}\PY{n}{String}\PY{o}{[}\PY{o}{]}\PY{+w}{ }\PY{n}{args}\PY{p}{)}\PY{+w}{ }\PY{p}{\PYZob{}}
\PY{+w}{        }\PY{n}{System}\PY{p}{.}\PY{n+na}{out}\PY{p}{.}\PY{n+na}{println}\PY{p}{(}\PY{l+s}{\PYZdq{}}\PY{l+s}{Default locale: }\PY{l+s}{\PYZdq{}}\PY{+w}{ }\PY{o}{+}\PY{+w}{ }\PY{n}{Locale}\PY{p}{.}\PY{n+na}{getDefault}\PY{p}{(}\PY{p}{)}\PY{p}{)}\PY{p}{;}
\PY{+w}{        }\PY{c+c1}{// Run with:   java \PYZhy{}Duser.language=de \PYZhy{}Duser.country=DE DefaultLocaleDemo}
\PY{+w}{        }\PY{c+c1}{// Output:     Default locale: de\PYZus{}DE}

\PY{+w}{        }\PY{c+c1}{// Run with:   java \PYZhy{}Duser.language=tr \PYZhy{}Duser.country=TR DefaultLocaleDemo}
\PY{+w}{        }\PY{c+c1}{// Output:     Default locale: tr\PYZus{}TR}
\PY{+w}{        }\PY{c+c1}{// (useful to specifically reproduce the Turkish dotless\PYZhy{}i bug shown earlier)}
\PY{+w}{    }\PY{p}{\PYZcb{}}
\PY{p}{\PYZcb{}}
\end{Verbatim}
\end{codebox}

\begin{codebox}
\begin{Verbatim}[commandchars=\\\{\}]
\PY{c+c1}{// A small JUnit\PYZhy{}style test that temporarily swaps the default locale to verify behavior.}
\PY{k+kn}{import}\PY{+w}{ }\PY{n+nn}{java.util.Locale}\PY{p}{;}
\PY{k+kn}{import}\PY{+w}{ }\PY{n+nn}{org.junit.jupiter.api.*}\PY{p}{;}

\PY{k+kd}{class} \PY{n+nc}{LocaleSensitiveTest}\PY{+w}{ }\PY{p}{\PYZob{}}
\PY{+w}{    }\PY{k+kd}{private}\PY{+w}{ }\PY{n}{Locale}\PY{+w}{ }\PY{n}{originalLocale}\PY{p}{;}

\PY{+w}{    }\PY{n+nd}{@BeforeEach}
\PY{+w}{    }\PY{k+kt}{void}\PY{+w}{ }\PY{n+nf}{saveLocale}\PY{p}{(}\PY{p}{)}\PY{+w}{ }\PY{p}{\PYZob{}}
\PY{+w}{        }\PY{n}{originalLocale}\PY{+w}{ }\PY{o}{=}\PY{+w}{ }\PY{n}{Locale}\PY{p}{.}\PY{n+na}{getDefault}\PY{p}{(}\PY{p}{)}\PY{p}{;}
\PY{+w}{    }\PY{p}{\PYZcb{}}

\PY{+w}{    }\PY{n+nd}{@AfterEach}
\PY{+w}{    }\PY{k+kt}{void}\PY{+w}{ }\PY{n+nf}{restoreLocale}\PY{p}{(}\PY{p}{)}\PY{+w}{ }\PY{p}{\PYZob{}}
\PY{+w}{        }\PY{n}{Locale}\PY{p}{.}\PY{n+na}{setDefault}\PY{p}{(}\PY{n}{originalLocale}\PY{p}{)}\PY{p}{;}\PY{+w}{ }\PY{c+c1}{// always restore \PYZhy{} global state affects other tests!}
\PY{+w}{    }\PY{p}{\PYZcb{}}

\PY{+w}{    }\PY{n+nd}{@Test}
\PY{+w}{    }\PY{k+kt}{void}\PY{+w}{ }\PY{n+nf}{formatsCorrectlyUnderGermanLocale}\PY{p}{(}\PY{p}{)}\PY{+w}{ }\PY{p}{\PYZob{}}
\PY{+w}{        }\PY{n}{Locale}\PY{p}{.}\PY{n+na}{setDefault}\PY{p}{(}\PY{k}{new}\PY{+w}{ }\PY{n}{Locale}\PY{p}{(}\PY{l+s}{\PYZdq{}}\PY{l+s}{de}\PY{l+s}{\PYZdq{}}\PY{p}{,}\PY{+w}{ }\PY{l+s}{\PYZdq{}}\PY{l+s}{DE}\PY{l+s}{\PYZdq{}}\PY{p}{)}\PY{p}{)}\PY{p}{;}
\PY{+w}{        }\PY{n}{String}\PY{+w}{ }\PY{n}{result}\PY{+w}{ }\PY{o}{=}\PY{+w}{ }\PY{n}{String}\PY{p}{.}\PY{n+na}{format}\PY{p}{(}\PY{l+s}{\PYZdq{}}\PY{l+s}{\PYZpc{},.2f}\PY{l+s}{\PYZdq{}}\PY{p}{,}\PY{+w}{ }\PY{l+m+mf}{1234.5}\PY{p}{)}\PY{p}{;}
\PY{+w}{        }\PY{n}{Assertions}\PY{p}{.}\PY{n+na}{assertEquals}\PY{p}{(}\PY{l+s}{\PYZdq{}}\PY{l+s}{1.234,50}\PY{l+s}{\PYZdq{}}\PY{p}{,}\PY{+w}{ }\PY{n}{result}\PY{p}{)}\PY{p}{;}
\PY{+w}{    }\PY{p}{\PYZcb{}}
\PY{p}{\PYZcb{}}
\end{Verbatim}
\end{codebox}

\subsection[CLDR as the default locale data provider]{CLDR as the default locale data provider}
\label{h:9:cldr-as-the-default-locale-data-provider}
CLDR (Unicode Common Locale Data Repository) is the industry-standard source of locale data — date/time patterns, number formats, currency symbols, plural rules, and more — maintained collaboratively (contributors include Google, Apple, Microsoft, IBM, and others). Since \textbf{JDK 9}, the JDK uses \textbf{CLDR} as the default locale data provider instead of its own older, less complete data (JRE-specific data, sometimes called “COMPAT”). This changed some formatting output compared to Java 8 — a real source of subtle behavior differences when upgrading old codebases.

\begin{codebox}
\begin{Verbatim}[commandchars=\\\{\}]
\PY{k+kn}{import}\PY{+w}{ }\PY{n+nn}{java.text.NumberFormat}\PY{p}{;}
\PY{k+kn}{import}\PY{+w}{ }\PY{n+nn}{java.util.Locale}\PY{p}{;}

\PY{k+kd}{public}\PY{+w}{ }\PY{k+kd}{class} \PY{n+nc}{CldrProviderDemo}\PY{+w}{ }\PY{p}{\PYZob{}}
\PY{+w}{    }\PY{k+kd}{public}\PY{+w}{ }\PY{k+kd}{static}\PY{+w}{ }\PY{k+kt}{void}\PY{+w}{ }\PY{n+nf}{main}\PY{p}{(}\PY{n}{String}\PY{o}{[}\PY{o}{]}\PY{+w}{ }\PY{n}{args}\PY{p}{)}\PY{+w}{ }\PY{p}{\PYZob{}}
\PY{+w}{        }\PY{n}{System}\PY{p}{.}\PY{n+na}{out}\PY{p}{.}\PY{n+na}{println}\PY{p}{(}\PY{n}{System}\PY{p}{.}\PY{n+na}{getProperty}\PY{p}{(}\PY{l+s}{\PYZdq{}}\PY{l+s}{java.locale.providers}\PY{l+s}{\PYZdq{}}\PY{p}{)}\PY{p}{)}\PY{p}{;}
\PY{+w}{        }\PY{c+c1}{// Output (JDK 9+): null by default (meaning CLDR is used implicitly), or \PYZdq{}CLDR\PYZdq{} if explicitly set}

\PY{+w}{        }\PY{n}{NumberFormat}\PY{+w}{ }\PY{n}{currency}\PY{+w}{ }\PY{o}{=}\PY{+w}{ }\PY{n}{NumberFormat}\PY{p}{.}\PY{n+na}{getCurrencyInstance}\PY{p}{(}\PY{n}{Locale}\PY{p}{.}\PY{n+na}{US}\PY{p}{)}\PY{p}{;}
\PY{+w}{        }\PY{n}{System}\PY{p}{.}\PY{n+na}{out}\PY{p}{.}\PY{n+na}{println}\PY{p}{(}\PY{n}{currency}\PY{p}{.}\PY{n+na}{format}\PY{p}{(}\PY{l+m+mf}{1234.5}\PY{p}{)}\PY{p}{)}\PY{p}{;}
\PY{+w}{        }\PY{c+c1}{// Output (CLDR, JDK 9+): \PYZdl{}1,234.50}
\PY{+w}{        }\PY{c+c1}{// On JDK 8 with the old COMPAT provider, some locales format currency slightly}
\PY{+w}{        }\PY{c+c1}{// differently (e.g. different symbol placement or spacing for certain locales).}

\PY{+w}{        }\PY{c+c1}{// You can force the old JDK 8 behavior (not recommended, mostly for migration debugging):}
\PY{+w}{        }\PY{c+c1}{// java \PYZhy{}Djava.locale.providers=COMPAT,CLDR CldrProviderDemo}
\PY{+w}{    }\PY{p}{\PYZcb{}}
\PY{p}{\PYZcb{}}
\end{Verbatim}
\end{codebox}

\section[Common Code-Review Interview Pitfalls]{Common Code-Review Interview Pitfalls}
\label{h:9:common-code-review-interview-pitfalls}
\begin{enumerate}
\item 
\textbf{Discarding the result of an immutable \code{plus}/\code{minus}/\code{with} call.}
Why it matters: \code{java.\allowbreak{}time} types are immutable — the call does nothing unless you capture the return value. This is the single most common \code{java.\allowbreak{}time} bug.

\begin{codebox}
\begin{Verbatim}[commandchars=\\\{\}]
\PY{c+c1}{// Before (bug): does nothing, \PYZdq{}deadline\PYZdq{} is unchanged}
\PY{n}{deadline}\PY{p}{.}\PY{n+na}{plusDays}\PY{p}{(}\PY{l+m+mi}{7}\PY{p}{)}\PY{p}{;}

\PY{c+c1}{// After (fixed): capture the new value}
\PY{n}{deadline}\PY{+w}{ }\PY{o}{=}\PY{+w}{ }\PY{n}{deadline}\PY{p}{.}\PY{n+na}{plusDays}\PY{p}{(}\PY{l+m+mi}{7}\PY{p}{)}\PY{p}{;}
\end{Verbatim}
\end{codebox}

\item 
\textbf{Sharing a mutable \code{SimpleDateFormat} as a \code{static} field across threads.}
Why it matters: \code{SimpleDateFormat} is not thread-safe; concurrent \code{format()}/\code{parse()} calls can corrupt results or throw exceptions.

\begin{codebox}
\begin{Verbatim}[commandchars=\\\{\}]
\PY{c+c1}{// Before (bug): shared mutable state}
\PY{k+kd}{static}\PY{+w}{ }\PY{k+kd}{final}\PY{+w}{ }\PY{n}{SimpleDateFormat}\PY{+w}{ }\PY{n}{FMT}\PY{+w}{ }\PY{o}{=}\PY{+w}{ }\PY{k}{new}\PY{+w}{ }\PY{n}{SimpleDateFormat}\PY{p}{(}\PY{l+s}{\PYZdq{}}\PY{l+s}{yyyy\PYZhy{}MM\PYZhy{}dd}\PY{l+s}{\PYZdq{}}\PY{p}{)}\PY{p}{;}

\PY{c+c1}{// After (fixed): DateTimeFormatter is immutable and thread\PYZhy{}safe}
\PY{k+kd}{static}\PY{+w}{ }\PY{k+kd}{final}\PY{+w}{ }\PY{n}{DateTimeFormatter}\PY{+w}{ }\PY{n}{FMT}\PY{+w}{ }\PY{o}{=}\PY{+w}{ }\PY{n}{DateTimeFormatter}\PY{p}{.}\PY{n+na}{ofPattern}\PY{p}{(}\PY{l+s}{\PYZdq{}}\PY{l+s}{yyyy\PYZhy{}MM\PYZhy{}dd}\PY{l+s}{\PYZdq{}}\PY{p}{)}\PY{p}{;}
\end{Verbatim}
\end{codebox}

\item 
\textbf{Storing \code{LocalDateTime} for events that need an absolute point in time.}
Why it matters: \code{LocalDateTime} has no zone, so “2026-08-07T14:00” is ambiguous across regions and can’t be safely compared or converted.

\begin{codebox}
\begin{Verbatim}[commandchars=\\\{\}]
\PY{c+c1}{// Before (bug): ambiguous, no zone info}
\PY{n}{LocalDateTime}\PY{+w}{ }\PY{n}{eventTime}\PY{+w}{ }\PY{o}{=}\PY{+w}{ }\PY{n}{LocalDateTime}\PY{p}{.}\PY{n+na}{now}\PY{p}{(}\PY{p}{)}\PY{p}{;}

\PY{c+c1}{// After (fixed): store an absolute instant, plus the zone if you need to redisplay locally}
\PY{n}{Instant}\PY{+w}{ }\PY{n}{eventInstant}\PY{+w}{ }\PY{o}{=}\PY{+w}{ }\PY{n}{Instant}\PY{p}{.}\PY{n+na}{now}\PY{p}{(}\PY{p}{)}\PY{p}{;}
\PY{n}{String}\PY{+w}{ }\PY{n}{zoneId}\PY{+w}{ }\PY{o}{=}\PY{+w}{ }\PY{l+s}{\PYZdq{}}\PY{l+s}{Europe/Amsterdam}\PY{l+s}{\PYZdq{}}\PY{p}{;}
\end{Verbatim}
\end{codebox}

\item 
\textbf{Using \code{YYYY} (week-year) instead of \code{yyyy} (calendar year) in a pattern.}
Why it matters: near year boundaries, the week-based year can differ from the calendar year, silently shifting dates by one year.

\begin{codebox}
\begin{Verbatim}[commandchars=\\\{\}]
\PY{c+c1}{// Before (bug): week\PYZhy{}year, wrong for calendar dates}
\PY{n}{DateTimeFormatter}\PY{p}{.}\PY{n+na}{ofPattern}\PY{p}{(}\PY{l+s}{\PYZdq{}}\PY{l+s}{YYYY\PYZhy{}MM\PYZhy{}dd}\PY{l+s}{\PYZdq{}}\PY{p}{)}\PY{p}{;}

\PY{c+c1}{// After (fixed): calendar year}
\PY{n}{DateTimeFormatter}\PY{p}{.}\PY{n+na}{ofPattern}\PY{p}{(}\PY{l+s}{\PYZdq{}}\PY{l+s}{yyyy\PYZhy{}MM\PYZhy{}dd}\PY{l+s}{\PYZdq{}}\PY{p}{)}\PY{p}{;}
\end{Verbatim}
\end{codebox}

\item 
\textbf{Calling \code{toUpperCase()}/\code{toLowerCase()} without a \code{Locale} for internal comparisons.}
Why it matters: the Turkish locale (and others) can transform \code{i}/\code{I} unexpectedly, breaking comparisons against ASCII constants.

\begin{codebox}
\begin{Verbatim}[commandchars=\\\{\}]
\PY{c+c1}{// Before (bug): depends on default locale}
\PY{k}{if}\PY{+w}{ }\PY{p}{(}\PY{n}{input}\PY{p}{.}\PY{n+na}{toUpperCase}\PY{p}{(}\PY{p}{)}\PY{p}{.}\PY{n+na}{equals}\PY{p}{(}\PY{l+s}{\PYZdq{}}\PY{l+s}{ID}\PY{l+s}{\PYZdq{}}\PY{p}{)}\PY{p}{)}\PY{+w}{ }\PY{p}{\PYZob{}}\PY{+w}{ }\PY{p}{.}\PY{p}{.}\PY{p}{.}\PY{+w}{ }\PY{p}{\PYZcb{}}

\PY{c+c1}{// After (fixed): locale\PYZhy{}independent}
\PY{k}{if}\PY{+w}{ }\PY{p}{(}\PY{n}{input}\PY{p}{.}\PY{n+na}{equalsIgnoreCase}\PY{p}{(}\PY{l+s}{\PYZdq{}}\PY{l+s}{ID}\PY{l+s}{\PYZdq{}}\PY{p}{)}\PY{p}{)}\PY{+w}{ }\PY{p}{\PYZob{}}\PY{+w}{ }\PY{p}{.}\PY{p}{.}\PY{p}{.}\PY{+w}{ }\PY{p}{\PYZcb{}}
\end{Verbatim}
\end{codebox}

\item 
\textbf{Calling \code{.\allowbreak{}now()} directly inside business logic instead of injecting a \code{Clock}.}
Why it matters: hardcoded “now” calls make unit tests flaky or impossible to make deterministic.

\begin{codebox}
\begin{Verbatim}[commandchars=\\\{\}]
\PY{c+c1}{// Before (hard to test)}
\PY{n}{LocalDate}\PY{+w}{ }\PY{n}{today}\PY{+w}{ }\PY{o}{=}\PY{+w}{ }\PY{n}{LocalDate}\PY{p}{.}\PY{n+na}{now}\PY{p}{(}\PY{p}{)}\PY{p}{;}

\PY{c+c1}{// After (testable)}
\PY{n}{LocalDate}\PY{+w}{ }\PY{n}{today}\PY{+w}{ }\PY{o}{=}\PY{+w}{ }\PY{n}{LocalDate}\PY{p}{.}\PY{n+na}{now}\PY{p}{(}\PY{n}{clock}\PY{p}{)}\PY{p}{;}\PY{+w}{ }\PY{c+c1}{// clock injected, Clock.fixed(...) in tests}
\end{Verbatim}
\end{codebox}

\item 
\textbf{Assuming “add one day” always adds exactly 24 hours.}
Why it matters: DST transitions mean a calendar day can be 23 or 25 real hours; use \code{Duration} for exact elapsed time, \code{Period}/\code{plusDays} for calendar-based reasoning.

\begin{codebox}
\begin{Verbatim}[commandchars=\\\{\}]
\PY{c+c1}{// Before (bug): assumes 24 real hours always pass}
\PY{n}{Instant}\PY{+w}{ }\PY{n}{next}\PY{+w}{ }\PY{o}{=}\PY{+w}{ }\PY{n}{instant}\PY{p}{.}\PY{n+na}{plus}\PY{p}{(}\PY{n}{Duration}\PY{p}{.}\PY{n+na}{ofDays}\PY{p}{(}\PY{l+m+mi}{1}\PY{p}{)}\PY{p}{)}\PY{p}{;}

\PY{c+c1}{// After (fixed): use calendar semantics when you mean \PYZdq{}the next calendar day\PYZdq{}}
\PY{n}{ZonedDateTime}\PY{+w}{ }\PY{n}{next}\PY{+w}{ }\PY{o}{=}\PY{+w}{ }\PY{n}{zonedDateTime}\PY{p}{.}\PY{n+na}{plusDays}\PY{p}{(}\PY{l+m+mi}{1}\PY{p}{)}\PY{p}{;}
\end{Verbatim}
\end{codebox}

\item 
\textbf{Comparing \code{LocalDateTime} values that originated from different time zones.}
Why it matters: \code{LocalDateTime} ignores zone offsets entirely, so “later on the clock” does not mean “later in reality.”

\begin{codebox}
\begin{Verbatim}[commandchars=\\\{\}]
\PY{c+c1}{// Before (bug): compares wall\PYZhy{}clock time, ignoring zone}
\PY{k+kt}{boolean}\PY{+w}{ }\PY{n}{isLater}\PY{+w}{ }\PY{o}{=}\PY{+w}{ }\PY{n}{a}\PY{p}{.}\PY{n+na}{toLocalDateTime}\PY{p}{(}\PY{p}{)}\PY{p}{.}\PY{n+na}{isAfter}\PY{p}{(}\PY{n}{b}\PY{p}{.}\PY{n+na}{toLocalDateTime}\PY{p}{(}\PY{p}{)}\PY{p}{)}\PY{p}{;}

\PY{c+c1}{// After (fixed): compare actual instants}
\PY{k+kt}{boolean}\PY{+w}{ }\PY{n}{isLater}\PY{+w}{ }\PY{o}{=}\PY{+w}{ }\PY{n}{a}\PY{p}{.}\PY{n+na}{toInstant}\PY{p}{(}\PY{p}{)}\PY{p}{.}\PY{n+na}{isAfter}\PY{p}{(}\PY{n}{b}\PY{p}{.}\PY{n+na}{toInstant}\PY{p}{(}\PY{p}{)}\PY{p}{)}\PY{p}{;}
\end{Verbatim}
\end{codebox}

\item 
\textbf{Letting \code{String.\allowbreak{}format}/\code{NumberFormat} use the JVM default locale in shared/server code.}
Why it matters: default locale can differ between machines and environments, causing inconsistent output (“works on my machine” bugs, wrong decimal separators sent to clients).

\begin{codebox}
\begin{Verbatim}[commandchars=\\\{\}]
\PY{c+c1}{// Before (bug): implicit, environment\PYZhy{}dependent}
\PY{n}{String}\PY{p}{.}\PY{n+na}{format}\PY{p}{(}\PY{l+s}{\PYZdq{}}\PY{l+s}{\PYZpc{},.2f}\PY{l+s}{\PYZdq{}}\PY{p}{,}\PY{+w}{ }\PY{n}{price}\PY{p}{)}\PY{p}{;}

\PY{c+c1}{// After (fixed): explicit locale}
\PY{n}{String}\PY{p}{.}\PY{n+na}{format}\PY{p}{(}\PY{n}{Locale}\PY{p}{.}\PY{n+na}{US}\PY{p}{,}\PY{+w}{ }\PY{l+s}{\PYZdq{}}\PY{l+s}{\PYZpc{},.2f}\PY{l+s}{\PYZdq{}}\PY{p}{,}\PY{+w}{ }\PY{n}{price}\PY{p}{)}\PY{p}{;}
\end{Verbatim}
\end{codebox}

\item 
\textbf{Not handling \code{DateTimeParseException} at boundaries that accept external date strings.}
Why it matters: it’s an unchecked exception; forgetting to catch it lets malformed user/API input crash a request instead of returning a clean validation error.

\begin{codebox}
\begin{Verbatim}[commandchars=\\\{\}]
\PY{c+c1}{// Before (bug): unhandled, crashes on bad input}
\PY{n}{LocalDate}\PY{+w}{ }\PY{n}{date}\PY{+w}{ }\PY{o}{=}\PY{+w}{ }\PY{n}{LocalDate}\PY{p}{.}\PY{n+na}{parse}\PY{p}{(}\PY{n}{userInput}\PY{p}{)}\PY{p}{;}

\PY{c+c1}{// After (fixed): handled explicitly}
\PY{k}{try}\PY{+w}{ }\PY{p}{\PYZob{}}
\PY{+w}{    }\PY{n}{LocalDate}\PY{+w}{ }\PY{n}{date}\PY{+w}{ }\PY{o}{=}\PY{+w}{ }\PY{n}{LocalDate}\PY{p}{.}\PY{n+na}{parse}\PY{p}{(}\PY{n}{userInput}\PY{p}{)}\PY{p}{;}
\PY{p}{\PYZcb{}}\PY{+w}{ }\PY{k}{catch}\PY{+w}{ }\PY{p}{(}\PY{n}{DateTimeParseException}\PY{+w}{ }\PY{n}{e}\PY{p}{)}\PY{+w}{ }\PY{p}{\PYZob{}}
\PY{+w}{    }\PY{k}{throw}\PY{+w}{ }\PY{k}{new}\PY{+w}{ }\PY{n}{InvalidRequestException}\PY{p}{(}\PY{l+s}{\PYZdq{}}\PY{l+s}{Invalid date: }\PY{l+s}{\PYZdq{}}\PY{+w}{ }\PY{o}{+}\PY{+w}{ }\PY{n}{userInput}\PY{p}{)}\PY{p}{;}
\PY{p}{\PYZcb{}}
\end{Verbatim}
\end{codebox}

\item 
\textbf{Using \code{String.\allowbreak{}length()} to validate or truncate text that may contain emoji or combining characters.}
Why it matters: \code{length()} counts UTF-16 code units, not visible characters, so limits can cut a string mid-character or reject valid short input.

\begin{codebox}
\begin{Verbatim}[commandchars=\\\{\}]
\PY{c+c1}{// Before (bug): may split a surrogate pair or combining sequence}
\PY{n}{String}\PY{+w}{ }\PY{n}{truncated}\PY{+w}{ }\PY{o}{=}\PY{+w}{ }\PY{n}{input}\PY{p}{.}\PY{n+na}{substring}\PY{p}{(}\PY{l+m+mi}{0}\PY{p}{,}\PY{+w}{ }\PY{n}{Math}\PY{p}{.}\PY{n+na}{min}\PY{p}{(}\PY{n}{input}\PY{p}{.}\PY{n+na}{length}\PY{p}{(}\PY{p}{)}\PY{p}{,}\PY{+w}{ }\PY{l+m+mi}{10}\PY{p}{)}\PY{p}{)}\PY{p}{;}

\PY{c+c1}{// After (better): use code point aware truncation, or a grapheme\PYZhy{}aware}
\PY{c+c1}{// BreakIterator for user\PYZhy{}facing character counts}
\PY{k+kt}{int}\PY{+w}{ }\PY{n}{end}\PY{+w}{ }\PY{o}{=}\PY{+w}{ }\PY{n}{input}\PY{p}{.}\PY{n+na}{offsetByCodePoints}\PY{p}{(}\PY{l+m+mi}{0}\PY{p}{,}\PY{+w}{ }\PY{n}{Math}\PY{p}{.}\PY{n+na}{min}\PY{p}{(}\PY{l+m+mi}{10}\PY{p}{,}\PY{+w}{ }\PY{n}{input}\PY{p}{.}\PY{n+na}{codePointCount}\PY{p}{(}\PY{l+m+mi}{0}\PY{p}{,}\PY{+w}{ }\PY{n}{input}\PY{p}{.}\PY{n+na}{length}\PY{p}{(}\PY{p}{)}\PY{p}{)}\PY{p}{)}\PY{p}{)}\PY{p}{;}
\PY{n}{String}\PY{+w}{ }\PY{n}{truncated}\PY{+w}{ }\PY{o}{=}\PY{+w}{ }\PY{n}{input}\PY{p}{.}\PY{n+na}{substring}\PY{p}{(}\PY{l+m+mi}{0}\PY{p}{,}\PY{+w}{ }\PY{n}{end}\PY{p}{)}\PY{p}{;}
\end{Verbatim}
\end{codebox}

\item 
\textbf{Comparing Unicode text with \code{.\allowbreak{}equals()} without normalizing first.}
Why it matters: visually identical strings can have different code point sequences (precomposed vs. decomposed accents), causing “duplicate” entries or failed lookups.

\begin{codebox}
\begin{Verbatim}[commandchars=\\\{\}]
\PY{c+c1}{// Before (bug): may treat identical\PYZhy{}looking text as different}
\PY{k}{if}\PY{+w}{ }\PY{p}{(}\PY{n}{nameFromUserA}\PY{p}{.}\PY{n+na}{equals}\PY{p}{(}\PY{n}{nameFromUserB}\PY{p}{)}\PY{p}{)}\PY{+w}{ }\PY{p}{\PYZob{}}\PY{+w}{ }\PY{p}{.}\PY{p}{.}\PY{p}{.}\PY{+w}{ }\PY{p}{\PYZcb{}}

\PY{c+c1}{// After (fixed): normalize first}
\PY{k}{if}\PY{+w}{ }\PY{p}{(}\PY{n}{Normalizer}\PY{p}{.}\PY{n+na}{normalize}\PY{p}{(}\PY{n}{nameFromUserA}\PY{p}{,}\PY{+w}{ }\PY{n}{Normalizer}\PY{p}{.}\PY{n+na}{Form}\PY{p}{.}\PY{n+na}{NFC}\PY{p}{)}
\PY{+w}{        }\PY{p}{.}\PY{n+na}{equals}\PY{p}{(}\PY{n}{Normalizer}\PY{p}{.}\PY{n+na}{normalize}\PY{p}{(}\PY{n}{nameFromUserB}\PY{p}{,}\PY{+w}{ }\PY{n}{Normalizer}\PY{p}{.}\PY{n+na}{Form}\PY{p}{.}\PY{n+na}{NFC}\PY{p}{)}\PY{p}{)}\PY{p}{)}\PY{+w}{ }\PY{p}{\PYZob{}}\PY{+w}{ }\PY{p}{.}\PY{p}{.}\PY{p}{.}\PY{+w}{ }\PY{p}{\PYZcb{}}
\end{Verbatim}
\end{codebox}

\item 
\textbf{Assuming the platform default charset is always UTF-8 on any JDK version.}
Why it matters: before JDK 18, the default charset was platform-dependent; code relying on the implicit default can misbehave when moved to another OS or an older JDK.

\begin{codebox}
\begin{Verbatim}[commandchars=\\\{\}]
\PY{c+c1}{// Before (risky on JDK \PYZlt{} 18 / non\PYZhy{}UTF\PYZhy{}8 platforms): implicit charset}
\PY{k+kt}{byte}\PY{o}{[}\PY{o}{]}\PY{+w}{ }\PY{n}{bytes}\PY{+w}{ }\PY{o}{=}\PY{+w}{ }\PY{n}{text}\PY{p}{.}\PY{n+na}{getBytes}\PY{p}{(}\PY{p}{)}\PY{p}{;}

\PY{c+c1}{// After (fixed): always explicit, safe on any JDK version}
\PY{k+kt}{byte}\PY{o}{[}\PY{o}{]}\PY{+w}{ }\PY{n}{bytes}\PY{+w}{ }\PY{o}{=}\PY{+w}{ }\PY{n}{text}\PY{p}{.}\PY{n+na}{getBytes}\PY{p}{(}\PY{n}{StandardCharsets}\PY{p}{.}\PY{n+na}{UTF\PYZus{}8}\PY{p}{)}\PY{p}{;}
\end{Verbatim}
\end{codebox}

\item 
\textbf{Hardcoding date/number formats instead of using locale-aware formatters for user-facing text.}
Why it matters: users expect dates and numbers formatted the way their own locale/region does, not a fixed pattern chosen by the developer.

\begin{codebox}
\begin{Verbatim}[commandchars=\\\{\}]
\PY{c+c1}{// Before (bug): fixed pattern regardless of user\PYZsq{}s locale}
\PY{n}{DateTimeFormatter}\PY{p}{.}\PY{n+na}{ofPattern}\PY{p}{(}\PY{l+s}{\PYZdq{}}\PY{l+s}{MM/dd/yyyy}\PY{l+s}{\PYZdq{}}\PY{p}{)}\PY{p}{;}

\PY{c+c1}{// After (fixed): adapts to the user\PYZsq{}s locale automatically}
\PY{n}{DateTimeFormatter}\PY{p}{.}\PY{n+na}{ofLocalizedDate}\PY{p}{(}\PY{n}{FormatStyle}\PY{p}{.}\PY{n+na}{MEDIUM}\PY{p}{)}\PY{p}{.}\PY{n+na}{withLocale}\PY{p}{(}\PY{n}{userLocale}\PY{p}{)}\PY{p}{;}
\end{Verbatim}
\end{codebox}

\item 
\textbf{Sorting user-facing text with \code{String.\allowbreak{}compareTo} instead of a locale-aware \code{Collator}.}
Why it matters: plain code-unit comparison does not match human alphabetical order for accented letters, case, or non-Latin scripts, producing confusing sort order in UIs.

\begin{codebox}
\begin{Verbatim}[commandchars=\\\{\}]
\PY{c+c1}{// Before (bug): raw code\PYZhy{}unit sort, wrong for most human languages}
\PY{n}{Arrays}\PY{p}{.}\PY{n+na}{sort}\PY{p}{(}\PY{n}{names}\PY{p}{)}\PY{p}{;}

\PY{c+c1}{// After (fixed): locale\PYZhy{}aware collation}
\PY{n}{Arrays}\PY{p}{.}\PY{n+na}{sort}\PY{p}{(}\PY{n}{names}\PY{p}{,}\PY{+w}{ }\PY{n}{Collator}\PY{p}{.}\PY{n+na}{getInstance}\PY{p}{(}\PY{n}{userLocale}\PY{p}{)}\PY{p}{)}\PY{p}{;}
\end{Verbatim}
\end{codebox}

\end{enumerate}


\setcounter{chapter}{9}
\chapter{File I/O}
\label{chap:10}
\label{h:10:10-file-io}
Every real program eventually has to read or write something outside its own memory: a config file, a log, a network payload, a saved object. Java gives you three overlapping toolkits for this — the old \code{java.\allowbreak{}io} streams, the newer NIO/NIO.2 (\code{java.\allowbreak{}nio}) APIs, and Java’s built-in object serialization. Each one solves a slightly different problem, and each one has sharp edges that show up constantly in code review. This chapter walks through all three, plus the modern \code{Path}/\code{Files} API that has mostly replaced \code{java.\allowbreak{}io.\allowbreak{}File}.

\textbf{Table of Contents}

\begin{itemize}
\item 
\hyperref[h:10:javaio-file-io]{java.io (File I/O)}

\item 
\hyperref[h:10:nio-and-nio2]{NIO and NIO.2}

\item 
\hyperref[h:10:paths-and-files]{Paths and Files}

\item 
\hyperref[h:10:serialization]{Serialization}

\item 
\hyperref[h:10:externalizable]{Externalizable}

\item 
\hyperref[h:10:common-code-review-interview-pitfalls]{Common Code-Review Interview Pitfalls}

\end{itemize}

\section[java.io (File I/O)]{java.io (File I/O)}
\label{h:10:javaio-file-io}
\code{java.\allowbreak{}io} is the original Java I/O library, dating back to Java 1.0. It is \emph{stream-based}: you read or write one chunk at a time, in order, from start to end. There is no random access and no built-in buffering — you add both yourself.

The most important split in \code{java.\allowbreak{}io} is \textbf{byte streams vs. character streams}.

\begin{itemize}
\item 
\textbf{Byte streams} (\code{InputStream} / \code{OutputStream}) move raw bytes (\code{0}–\code{255}). Use them for binary data: images, zip files, network sockets, serialized objects.

\item 
\textbf{Character streams} (\code{Reader} / \code{Writer}) move \code{char} values (text). Internally they still read/write bytes, but they convert bytes to characters using a \textbf{charset} (a character encoding table, e.g. UTF-8). Use them for text files.

\end{itemize}

Mixing the two up is a classic bug: reading text with a byte stream and hand-rolling your own encoding logic, or reading binary data with a character stream and corrupting bytes during “encoding conversion” that binary data was never meant to go through.

\subsection[The class hierarchy]{The class hierarchy}
\label{h:10:the-class-hierarchy}
\begin{codebox}
\begin{Verbatim}
                     byte streams                         character streams

               InputStream (abstract)                 Reader (abstract)
               /      |        \                      /      |       \
FileInputStream  ByteArrayInputStream  BufferedInputStream   FileReader  StringReader  BufferedReader
      |                                       (decorator)         |                       (decorator)
      +-- wrapped by InputStreamReader --------------------------> (bridges byte -> char)

              OutputStream (abstract)                 Writer (abstract)
               /      |        \                      /      |       \
FileOutputStream  ByteArrayOutputStream  BufferedOutputStream  FileWriter  StringWriter  BufferedWriter
      |                                       (decorator)         |                       (decorator)
      +-- wrapped by OutputStreamWriter -------------------------> (bridges char -> byte)
\end{Verbatim}
\end{codebox}

\code{InputStreamReader} and \code{OutputStreamWriter} are the bridge classes: they wrap a byte stream and a charset, and expose it as a character stream. \code{FileReader} and \code{FileWriter} are just convenience subclasses of that bridge, hardwired to a \code{File} — historically they used the \emph{platform default charset}, which is exactly why you should avoid their no-charset constructors (see below).

\subsection[InputStream / OutputStream / Reader / Writer]{InputStream / OutputStream / Reader / Writer}
\label{h:10:inputstream--outputstream--reader--writer}
These four abstract classes are the roots of everything in \code{java.\allowbreak{}io}.

\begin{codebox}
\begin{Verbatim}[commandchars=\\\{\}]
\PY{k+kn}{import}\PY{+w}{ }\PY{n+nn}{java.io.*}\PY{p}{;}
\PY{k+kn}{import}\PY{+w}{ }\PY{n+nn}{java.nio.charset.StandardCharsets}\PY{p}{;}

\PY{k+kd}{public}\PY{+w}{ }\PY{k+kd}{class} \PY{n+nc}{BasicStreams}\PY{+w}{ }\PY{p}{\PYZob{}}
\PY{+w}{    }\PY{k+kd}{public}\PY{+w}{ }\PY{k+kd}{static}\PY{+w}{ }\PY{k+kt}{void}\PY{+w}{ }\PY{n+nf}{main}\PY{p}{(}\PY{n}{String}\PY{o}{[}\PY{o}{]}\PY{+w}{ }\PY{n}{args}\PY{p}{)}\PY{+w}{ }\PY{k+kd}{throws}\PY{+w}{ }\PY{n}{IOException}\PY{+w}{ }\PY{p}{\PYZob{}}
\PY{+w}{        }\PY{c+c1}{// Byte stream: read raw bytes from a file}
\PY{+w}{        }\PY{k}{try}\PY{+w}{ }\PY{p}{(}\PY{n}{InputStream}\PY{+w}{ }\PY{n}{in}\PY{+w}{ }\PY{o}{=}\PY{+w}{ }\PY{k}{new}\PY{+w}{ }\PY{n}{FileInputStream}\PY{p}{(}\PY{l+s}{\PYZdq{}}\PY{l+s}{data.bin}\PY{l+s}{\PYZdq{}}\PY{p}{)}\PY{p}{)}\PY{+w}{ }\PY{p}{\PYZob{}}
\PY{+w}{            }\PY{k+kt}{int}\PY{+w}{ }\PY{n}{b}\PY{p}{;}
\PY{+w}{            }\PY{k}{while}\PY{+w}{ }\PY{p}{(}\PY{p}{(}\PY{n}{b}\PY{+w}{ }\PY{o}{=}\PY{+w}{ }\PY{n}{in}\PY{p}{.}\PY{n+na}{read}\PY{p}{(}\PY{p}{)}\PY{p}{)}\PY{+w}{ }\PY{o}{!}\PY{o}{=}\PY{+w}{ }\PY{o}{\PYZhy{}}\PY{l+m+mi}{1}\PY{p}{)}\PY{+w}{ }\PY{p}{\PYZob{}}\PY{+w}{   }\PY{c+c1}{// read() returns \PYZhy{}1 at end of stream}
\PY{+w}{                }\PY{c+c1}{// process one byte at a time (slow without buffering, see below)}
\PY{+w}{            }\PY{p}{\PYZcb{}}
\PY{+w}{        }\PY{p}{\PYZcb{}}

\PY{+w}{        }\PY{c+c1}{// Character stream: read text, one char at a time}
\PY{+w}{        }\PY{k}{try}\PY{+w}{ }\PY{p}{(}\PY{n}{Reader}\PY{+w}{ }\PY{n}{reader}\PY{+w}{ }\PY{o}{=}\PY{+w}{ }\PY{k}{new}\PY{+w}{ }\PY{n}{InputStreamReader}\PY{p}{(}
\PY{+w}{                }\PY{k}{new}\PY{+w}{ }\PY{n}{FileInputStream}\PY{p}{(}\PY{l+s}{\PYZdq{}}\PY{l+s}{notes.txt}\PY{l+s}{\PYZdq{}}\PY{p}{)}\PY{p}{,}\PY{+w}{ }\PY{n}{StandardCharsets}\PY{p}{.}\PY{n+na}{UTF\PYZus{}8}\PY{p}{)}\PY{p}{)}\PY{+w}{ }\PY{p}{\PYZob{}}
\PY{+w}{            }\PY{k+kt}{int}\PY{+w}{ }\PY{n}{c}\PY{p}{;}
\PY{+w}{            }\PY{k}{while}\PY{+w}{ }\PY{p}{(}\PY{p}{(}\PY{n}{c}\PY{+w}{ }\PY{o}{=}\PY{+w}{ }\PY{n}{reader}\PY{p}{.}\PY{n+na}{read}\PY{p}{(}\PY{p}{)}\PY{p}{)}\PY{+w}{ }\PY{o}{!}\PY{o}{=}\PY{+w}{ }\PY{o}{\PYZhy{}}\PY{l+m+mi}{1}\PY{p}{)}\PY{+w}{ }\PY{p}{\PYZob{}}
\PY{+w}{                }\PY{c+c1}{// process one character}
\PY{+w}{            }\PY{p}{\PYZcb{}}
\PY{+w}{        }\PY{p}{\PYZcb{}}
\PY{+w}{    }\PY{p}{\PYZcb{}}
\PY{p}{\PYZcb{}}
\end{Verbatim}
\end{codebox}

\subsection[The decorator pattern]{The decorator pattern}
\label{h:10:the-decorator-pattern}
\code{java.\allowbreak{}io} is built almost entirely on the \textbf{decorator pattern}: small classes that wrap another stream and add one piece of behavior, without changing the type. You “stack” decorators to combine behaviors.

\begin{codebox}
\begin{Verbatim}[commandchars=\\\{\}]
\PY{k+kn}{import}\PY{+w}{ }\PY{n+nn}{java.io.*}\PY{p}{;}
\PY{k+kn}{import}\PY{+w}{ }\PY{n+nn}{java.nio.charset.StandardCharsets}\PY{p}{;}

\PY{c+c1}{// FileInputStream: raw bytes from disk}
\PY{c+c1}{// \PYZhy{}\PYZgt{} wrapped by BufferedInputStream: adds an internal buffer (fewer OS calls)}
\PY{c+c1}{// \PYZhy{}\PYZgt{} wrapped by InputStreamReader: converts bytes to chars using UTF\PYZhy{}8}
\PY{c+c1}{// \PYZhy{}\PYZgt{} wrapped by BufferedReader: adds line\PYZhy{}based reading (readLine())}
\PY{k}{try}\PY{+w}{ }\PY{p}{(}\PY{n}{BufferedReader}\PY{+w}{ }\PY{n}{reader}\PY{+w}{ }\PY{o}{=}\PY{+w}{ }\PY{k}{new}\PY{+w}{ }\PY{n}{BufferedReader}\PY{p}{(}
\PY{+w}{        }\PY{k}{new}\PY{+w}{ }\PY{n}{InputStreamReader}\PY{p}{(}
\PY{+w}{                }\PY{k}{new}\PY{+w}{ }\PY{n}{BufferedInputStream}\PY{p}{(}\PY{k}{new}\PY{+w}{ }\PY{n}{FileInputStream}\PY{p}{(}\PY{l+s}{\PYZdq{}}\PY{l+s}{report.txt}\PY{l+s}{\PYZdq{}}\PY{p}{)}\PY{p}{)}\PY{p}{,}
\PY{+w}{                }\PY{n}{StandardCharsets}\PY{p}{.}\PY{n+na}{UTF\PYZus{}8}\PY{p}{)}\PY{p}{)}\PY{p}{)}\PY{+w}{ }\PY{p}{\PYZob{}}

\PY{+w}{    }\PY{n}{String}\PY{+w}{ }\PY{n}{line}\PY{p}{;}
\PY{+w}{    }\PY{k}{while}\PY{+w}{ }\PY{p}{(}\PY{p}{(}\PY{n}{line}\PY{+w}{ }\PY{o}{=}\PY{+w}{ }\PY{n}{reader}\PY{p}{.}\PY{n+na}{readLine}\PY{p}{(}\PY{p}{)}\PY{p}{)}\PY{+w}{ }\PY{o}{!}\PY{o}{=}\PY{+w}{ }\PY{k+kc}{null}\PY{p}{)}\PY{+w}{ }\PY{p}{\PYZob{}}
\PY{+w}{        }\PY{n}{System}\PY{p}{.}\PY{n+na}{out}\PY{p}{.}\PY{n+na}{println}\PY{p}{(}\PY{n}{line}\PY{p}{)}\PY{p}{;}
\PY{+w}{    }\PY{p}{\PYZcb{}}
\PY{p}{\PYZcb{}}
\end{Verbatim}
\end{codebox}

Each layer only knows about the layer directly underneath it. That is the whole point of the pattern: \code{BufferedReader} does not care whether its underlying \code{Reader} is a file, a socket, or an in-memory string.

\subsection[Why buffering matters]{Why buffering matters}
\label{h:10:why-buffering-matters}
Every unbuffered \code{read()}/\code{write()} call on a \code{FileInputStream} or \code{FileOutputStream} typically triggers a system call to the OS. System calls are expensive compared to in-memory operations — doing one per byte is enormously wasteful.

\begin{codebox}
\begin{Verbatim}[commandchars=\\\{\}]
\PY{k+kn}{import}\PY{+w}{ }\PY{n+nn}{java.io.*}\PY{p}{;}

\PY{c+c1}{// Bad: one system call per byte. On a 1 MB file, that\PYZsq{}s \PYZti{}1,000,000 syscalls.}
\PY{k}{try}\PY{+w}{ }\PY{p}{(}\PY{n}{InputStream}\PY{+w}{ }\PY{n}{in}\PY{+w}{ }\PY{o}{=}\PY{+w}{ }\PY{k}{new}\PY{+w}{ }\PY{n}{FileInputStream}\PY{p}{(}\PY{l+s}{\PYZdq{}}\PY{l+s}{big.dat}\PY{l+s}{\PYZdq{}}\PY{p}{)}\PY{p}{)}\PY{+w}{ }\PY{p}{\PYZob{}}
\PY{+w}{    }\PY{k+kt}{int}\PY{+w}{ }\PY{n}{b}\PY{p}{;}
\PY{+w}{    }\PY{k}{while}\PY{+w}{ }\PY{p}{(}\PY{p}{(}\PY{n}{b}\PY{+w}{ }\PY{o}{=}\PY{+w}{ }\PY{n}{in}\PY{p}{.}\PY{n+na}{read}\PY{p}{(}\PY{p}{)}\PY{p}{)}\PY{+w}{ }\PY{o}{!}\PY{o}{=}\PY{+w}{ }\PY{o}{\PYZhy{}}\PY{l+m+mi}{1}\PY{p}{)}\PY{+w}{ }\PY{p}{\PYZob{}}
\PY{+w}{        }\PY{c+c1}{// ...}
\PY{+w}{    }\PY{p}{\PYZcb{}}
\PY{p}{\PYZcb{}}

\PY{c+c1}{// Good: BufferedInputStream reads a large chunk (default 8 KB) into memory once,}
\PY{c+c1}{// then serves read() calls from that in\PYZhy{}memory buffer. Far fewer syscalls.}
\PY{k}{try}\PY{+w}{ }\PY{p}{(}\PY{n}{InputStream}\PY{+w}{ }\PY{n}{in}\PY{+w}{ }\PY{o}{=}\PY{+w}{ }\PY{k}{new}\PY{+w}{ }\PY{n}{BufferedInputStream}\PY{p}{(}\PY{k}{new}\PY{+w}{ }\PY{n}{FileInputStream}\PY{p}{(}\PY{l+s}{\PYZdq{}}\PY{l+s}{big.dat}\PY{l+s}{\PYZdq{}}\PY{p}{)}\PY{p}{)}\PY{p}{)}\PY{+w}{ }\PY{p}{\PYZob{}}
\PY{+w}{    }\PY{k+kt}{int}\PY{+w}{ }\PY{n}{b}\PY{p}{;}
\PY{+w}{    }\PY{k}{while}\PY{+w}{ }\PY{p}{(}\PY{p}{(}\PY{n}{b}\PY{+w}{ }\PY{o}{=}\PY{+w}{ }\PY{n}{in}\PY{p}{.}\PY{n+na}{read}\PY{p}{(}\PY{p}{)}\PY{p}{)}\PY{+w}{ }\PY{o}{!}\PY{o}{=}\PY{+w}{ }\PY{o}{\PYZhy{}}\PY{l+m+mi}{1}\PY{p}{)}\PY{+w}{ }\PY{p}{\PYZob{}}
\PY{+w}{        }\PY{c+c1}{// ...}
\PY{+w}{    }\PY{p}{\PYZcb{}}
\PY{p}{\PYZcb{}}
\end{Verbatim}
\end{codebox}

The rule of thumb in review: \textbf{any raw \code{FileInputStream}/\code{FileOutputStream}/\code{FileReader}/\code{FileWriter} that isn’t wrapped in a \code{Buffered*} decorator is a performance smell}, unless the code already reads/writes in large chunks itself (e.g., \code{readAllBytes()}).

\subsection[Always specify a charset]{Always specify a charset}
\label{h:10:always-specify-a-charset}
Text is just bytes plus an agreement on how to decode them. If you don’t specify a charset, Java uses the \textbf{platform default charset}, which depends on the OS, locale, and JVM flags. That means the exact same code can produce different (or garbled) output on different machines — a Linux CI server, a developer’s macOS laptop, and a Windows box might all disagree.

\begin{codebox}
\begin{Verbatim}[commandchars=\\\{\}]
\PY{k+kn}{import}\PY{+w}{ }\PY{n+nn}{java.io.*}\PY{p}{;}
\PY{k+kn}{import}\PY{+w}{ }\PY{n+nn}{java.nio.charset.StandardCharsets}\PY{p}{;}

\PY{c+c1}{// Bad: charset is whatever the platform happens to default to (risky, not portable)}
\PY{n}{Writer}\PY{+w}{ }\PY{n}{badWriter}\PY{+w}{ }\PY{o}{=}\PY{+w}{ }\PY{k}{new}\PY{+w}{ }\PY{n}{FileWriter}\PY{p}{(}\PY{l+s}{\PYZdq{}}\PY{l+s}{out.txt}\PY{l+s}{\PYZdq{}}\PY{p}{)}\PY{p}{;}

\PY{c+c1}{// Good: charset is explicit and reproducible everywhere}
\PY{n}{Writer}\PY{+w}{ }\PY{n}{goodWriter}\PY{+w}{ }\PY{o}{=}\PY{+w}{ }\PY{k}{new}\PY{+w}{ }\PY{n}{OutputStreamWriter}\PY{p}{(}
\PY{+w}{        }\PY{k}{new}\PY{+w}{ }\PY{n}{FileOutputStream}\PY{p}{(}\PY{l+s}{\PYZdq{}}\PY{l+s}{out.txt}\PY{l+s}{\PYZdq{}}\PY{p}{)}\PY{p}{,}\PY{+w}{ }\PY{n}{StandardCharsets}\PY{p}{.}\PY{n+na}{UTF\PYZus{}8}\PY{p}{)}\PY{p}{;}

\PY{c+c1}{// Java 11+ convenience: FileReader/FileWriter overloads that accept a Charset directly}
\PY{n}{Writer}\PY{+w}{ }\PY{n}{goodWriterJava11}\PY{+w}{ }\PY{o}{=}\PY{+w}{ }\PY{k}{new}\PY{+w}{ }\PY{n}{FileWriter}\PY{p}{(}\PY{l+s}{\PYZdq{}}\PY{l+s}{out.txt}\PY{l+s}{\PYZdq{}}\PY{p}{,}\PY{+w}{ }\PY{n}{StandardCharsets}\PY{p}{.}\PY{n+na}{UTF\PYZus{}8}\PY{p}{)}\PY{p}{;}
\end{Verbatim}
\end{codebox}

Since Java 18, the JVM’s default charset for most contexts is UTF-8 (JEP 400), which fixes a lot of historical foot-guns. But relying on “the default happens to be UTF-8 now” is still fragile — a reviewer should still flag missing explicit charsets, since the JVM can still be launched with \code{-\allowbreak{}Dfile.\allowbreak{}encoding=...} overriding it, and library code should never assume the caller’s environment.

\subsection[File and its weak error reporting]{File and its weak error reporting}
\label{h:10:file-and-its-weak-error-reporting}
\code{java.\allowbreak{}io.\allowbreak{}File} represents a path on the filesystem — but it is notorious for \textbf{silently returning \code{false} instead of throwing an exception} when something goes wrong. This makes bugs hard to diagnose because you get no information about \emph{why} an operation failed.

\begin{codebox}
\begin{Verbatim}[commandchars=\\\{\}]
\PY{k+kn}{import}\PY{+w}{ }\PY{n+nn}{java.io.File}\PY{p}{;}

\PY{n}{File}\PY{+w}{ }\PY{n}{file}\PY{+w}{ }\PY{o}{=}\PY{+w}{ }\PY{k}{new}\PY{+w}{ }\PY{n}{File}\PY{p}{(}\PY{l+s}{\PYZdq{}}\PY{l+s}{/some/protected/path/data.txt}\PY{l+s}{\PYZdq{}}\PY{p}{)}\PY{p}{;}

\PY{k+kt}{boolean}\PY{+w}{ }\PY{n}{deleted}\PY{+w}{ }\PY{o}{=}\PY{+w}{ }\PY{n}{file}\PY{p}{.}\PY{n+na}{delete}\PY{p}{(}\PY{p}{)}\PY{p}{;}
\PY{k}{if}\PY{+w}{ }\PY{p}{(}\PY{o}{!}\PY{n}{deleted}\PY{p}{)}\PY{+w}{ }\PY{p}{\PYZob{}}
\PY{+w}{    }\PY{c+c1}{// Was it \PYZdq{}file doesn\PYZsq{}t exist\PYZdq{}? \PYZdq{}Permission denied\PYZdq{}? \PYZdq{}Directory not empty\PYZdq{}?}
\PY{+w}{    }\PY{c+c1}{// File.delete() gives you NO clue. You just get `false`.}
\PY{+w}{    }\PY{n}{System}\PY{p}{.}\PY{n+na}{out}\PY{p}{.}\PY{n+na}{println}\PY{p}{(}\PY{l+s}{\PYZdq{}}\PY{l+s}{Delete failed, but we don\PYZsq{}t know why.}\PY{l+s}{\PYZdq{}}\PY{p}{)}\PY{p}{;}
\PY{p}{\PYZcb{}}

\PY{k+kt}{boolean}\PY{+w}{ }\PY{n}{created}\PY{+w}{ }\PY{o}{=}\PY{+w}{ }\PY{n}{file}\PY{p}{.}\PY{n+na}{mkdir}\PY{p}{(}\PY{p}{)}\PY{p}{;}
\PY{k}{if}\PY{+w}{ }\PY{p}{(}\PY{o}{!}\PY{n}{created}\PY{p}{)}\PY{+w}{ }\PY{p}{\PYZob{}}
\PY{+w}{    }\PY{c+c1}{// Same problem: could be \PYZdq{}parent doesn\PYZsq{}t exist\PYZdq{}, \PYZdq{}already exists\PYZdq{},}
\PY{+w}{    }\PY{c+c1}{// \PYZdq{}permission denied\PYZdq{} — all collapse into a single boolean.}
\PY{+w}{    }\PY{n}{System}\PY{p}{.}\PY{n+na}{out}\PY{p}{.}\PY{n+na}{println}\PY{p}{(}\PY{l+s}{\PYZdq{}}\PY{l+s}{mkdir failed, reason unknown.}\PY{l+s}{\PYZdq{}}\PY{p}{)}\PY{p}{;}
\PY{p}{\PYZcb{}}
\end{Verbatim}
\end{codebox}

This is one of the concrete reasons \code{java.\allowbreak{}nio.\allowbreak{}file.\allowbreak{}Files} (covered below) exists: \code{Files.\allowbreak{}delete()} and \code{Files.\allowbreak{}create\allowbreak{}Directory()} throw specific, informative exceptions (\code{NoSuchFileException}, \code{File\allowbreak{}Already\allowbreak{}Exists\allowbreak{}Exception}, \code{AccessDeniedException}, \code{Directory\allowbreak{}Not\allowbreak{}Empty\allowbreak{}Exception}) instead of a flat boolean.

\subsection[try-with-resources for streams]{try-with-resources for streams}
\label{h:10:try-with-resources-for-streams}
All \code{java.\allowbreak{}io} streams implement \code{Closeable} (which extends \code{AutoCloseable}), so they work with \code{try-\allowbreak{}with-\allowbreak{}resources}. Forgetting to close a stream leaks file handles or sockets — on a long-running server this eventually exhausts OS file-descriptor limits.

\begin{codebox}
\begin{Verbatim}[commandchars=\\\{\}]
\PY{k+kn}{import}\PY{+w}{ }\PY{n+nn}{java.io.*}\PY{p}{;}
\PY{k+kn}{import}\PY{+w}{ }\PY{n+nn}{java.nio.charset.StandardCharsets}\PY{p}{;}

\PY{c+c1}{// Bad: if readLine() throws, the stream never gets closed}
\PY{n}{BufferedReader}\PY{+w}{ }\PY{n}{reader}\PY{+w}{ }\PY{o}{=}\PY{+w}{ }\PY{k}{new}\PY{+w}{ }\PY{n}{BufferedReader}\PY{p}{(}
\PY{+w}{        }\PY{k}{new}\PY{+w}{ }\PY{n}{InputStreamReader}\PY{p}{(}\PY{k}{new}\PY{+w}{ }\PY{n}{FileInputStream}\PY{p}{(}\PY{l+s}{\PYZdq{}}\PY{l+s}{data.txt}\PY{l+s}{\PYZdq{}}\PY{p}{)}\PY{p}{,}\PY{+w}{ }\PY{n}{StandardCharsets}\PY{p}{.}\PY{n+na}{UTF\PYZus{}8}\PY{p}{)}\PY{p}{)}\PY{p}{;}
\PY{n}{String}\PY{+w}{ }\PY{n}{line}\PY{+w}{ }\PY{o}{=}\PY{+w}{ }\PY{n}{reader}\PY{p}{.}\PY{n+na}{readLine}\PY{p}{(}\PY{p}{)}\PY{p}{;}
\PY{n}{reader}\PY{p}{.}\PY{n+na}{close}\PY{p}{(}\PY{p}{)}\PY{p}{;}\PY{+w}{ }\PY{c+c1}{// never reached if an exception is thrown above}

\PY{c+c1}{// Good: try\PYZhy{}with\PYZhy{}resources guarantees close(), even on exception,}
\PY{c+c1}{// and closes multiple resources in reverse order of declaration}
\PY{k}{try}\PY{+w}{ }\PY{p}{(}\PY{n}{FileInputStream}\PY{+w}{ }\PY{n}{fis}\PY{+w}{ }\PY{o}{=}\PY{+w}{ }\PY{k}{new}\PY{+w}{ }\PY{n}{FileInputStream}\PY{p}{(}\PY{l+s}{\PYZdq{}}\PY{l+s}{data.txt}\PY{l+s}{\PYZdq{}}\PY{p}{)}\PY{p}{;}
\PY{+w}{     }\PY{n}{InputStreamReader}\PY{+w}{ }\PY{n}{isr}\PY{+w}{ }\PY{o}{=}\PY{+w}{ }\PY{k}{new}\PY{+w}{ }\PY{n}{InputStreamReader}\PY{p}{(}\PY{n}{fis}\PY{p}{,}\PY{+w}{ }\PY{n}{StandardCharsets}\PY{p}{.}\PY{n+na}{UTF\PYZus{}8}\PY{p}{)}\PY{p}{;}
\PY{+w}{     }\PY{n}{BufferedReader}\PY{+w}{ }\PY{n}{reader2}\PY{+w}{ }\PY{o}{=}\PY{+w}{ }\PY{k}{new}\PY{+w}{ }\PY{n}{BufferedReader}\PY{p}{(}\PY{n}{isr}\PY{p}{)}\PY{p}{)}\PY{+w}{ }\PY{p}{\PYZob{}}

\PY{+w}{    }\PY{n}{String}\PY{+w}{ }\PY{n}{firstLine}\PY{+w}{ }\PY{o}{=}\PY{+w}{ }\PY{n}{reader2}\PY{p}{.}\PY{n+na}{readLine}\PY{p}{(}\PY{p}{)}\PY{p}{;}
\PY{+w}{    }\PY{n}{System}\PY{p}{.}\PY{n+na}{out}\PY{p}{.}\PY{n+na}{println}\PY{p}{(}\PY{n}{firstLine}\PY{p}{)}\PY{p}{;}
\PY{p}{\PYZcb{}}\PY{+w}{ }\PY{c+c1}{// fis, isr, and reader2 are all closed automatically here}
\end{Verbatim}
\end{codebox}

\section[NIO and NIO.2]{NIO and NIO.2}
\label{h:10:nio-and-nio2}
“NIO” (New I/O, \code{java.\allowbreak{}nio}, added in Java 1.4) and “NIO.2” (the \code{java.\allowbreak{}nio.\allowbreak{}file} package, added in Java 7) are a different model from \code{java.\allowbreak{}io}. Instead of a one-directional stream of bytes, NIO works with \textbf{buffers} and \textbf{channels}, and supports non-blocking and memory-mapped I/O. NIO.2 additionally gave us the modern \code{Path}/\code{Files} API, covered in the next section.

\subsection[Buffers and channels]{Buffers and channels}
\label{h:10:buffers-and-channels}
A \textbf{channel} (\code{FileChannel}, \code{SocketChannel}, …) is a connection to an I/O source — a file, a socket. Unlike a \code{Stream}, a channel can be read from \emph{and} written to, and (for files) supports random access via \code{position()}.

A \textbf{buffer} (\code{ByteBuffer}, \code{CharBuffer}, …) is a fixed-size block of memory that channels read into and write from. You don’t read byte-by-byte with NIO; you read a whole chunk into a buffer, then process the buffer.

\begin{codebox}
\begin{Verbatim}[commandchars=\\\{\}]
\PY{k+kn}{import}\PY{+w}{ }\PY{n+nn}{java.io.IOException}\PY{p}{;}
\PY{k+kn}{import}\PY{+w}{ }\PY{n+nn}{java.nio.ByteBuffer}\PY{p}{;}
\PY{k+kn}{import}\PY{+w}{ }\PY{n+nn}{java.nio.channels.FileChannel}\PY{p}{;}
\PY{k+kn}{import}\PY{+w}{ }\PY{n+nn}{java.nio.file.Path}\PY{p}{;}
\PY{k+kn}{import}\PY{+w}{ }\PY{n+nn}{java.nio.file.StandardOpenOption}\PY{p}{;}

\PY{k+kd}{public}\PY{+w}{ }\PY{k+kd}{class} \PY{n+nc}{ChannelBufferExample}\PY{+w}{ }\PY{p}{\PYZob{}}
\PY{+w}{    }\PY{k+kd}{public}\PY{+w}{ }\PY{k+kd}{static}\PY{+w}{ }\PY{k+kt}{void}\PY{+w}{ }\PY{n+nf}{main}\PY{p}{(}\PY{n}{String}\PY{o}{[}\PY{o}{]}\PY{+w}{ }\PY{n}{args}\PY{p}{)}\PY{+w}{ }\PY{k+kd}{throws}\PY{+w}{ }\PY{n}{IOException}\PY{+w}{ }\PY{p}{\PYZob{}}
\PY{+w}{        }\PY{n}{Path}\PY{+w}{ }\PY{n}{path}\PY{+w}{ }\PY{o}{=}\PY{+w}{ }\PY{n}{Path}\PY{p}{.}\PY{n+na}{of}\PY{p}{(}\PY{l+s}{\PYZdq{}}\PY{l+s}{data.bin}\PY{l+s}{\PYZdq{}}\PY{p}{)}\PY{p}{;}

\PY{+w}{        }\PY{k}{try}\PY{+w}{ }\PY{p}{(}\PY{n}{FileChannel}\PY{+w}{ }\PY{n}{channel}\PY{+w}{ }\PY{o}{=}\PY{+w}{ }\PY{n}{FileChannel}\PY{p}{.}\PY{n+na}{open}\PY{p}{(}\PY{n}{path}\PY{p}{,}\PY{+w}{ }\PY{n}{StandardOpenOption}\PY{p}{.}\PY{n+na}{READ}\PY{p}{)}\PY{p}{)}\PY{+w}{ }\PY{p}{\PYZob{}}
\PY{+w}{            }\PY{n}{ByteBuffer}\PY{+w}{ }\PY{n}{buffer}\PY{+w}{ }\PY{o}{=}\PY{+w}{ }\PY{n}{ByteBuffer}\PY{p}{.}\PY{n+na}{allocate}\PY{p}{(}\PY{l+m+mi}{1024}\PY{p}{)}\PY{p}{;}\PY{+w}{ }\PY{c+c1}{// 1 KB heap buffer}

\PY{+w}{            }\PY{k+kt}{int}\PY{+w}{ }\PY{n}{bytesRead}\PY{+w}{ }\PY{o}{=}\PY{+w}{ }\PY{n}{channel}\PY{p}{.}\PY{n+na}{read}\PY{p}{(}\PY{n}{buffer}\PY{p}{)}\PY{p}{;}\PY{+w}{ }\PY{c+c1}{// fills the buffer, returns count (\PYZhy{}1 = EOF)}
\PY{+w}{            }\PY{n}{System}\PY{p}{.}\PY{n+na}{out}\PY{p}{.}\PY{n+na}{println}\PY{p}{(}\PY{l+s}{\PYZdq{}}\PY{l+s}{Read }\PY{l+s}{\PYZdq{}}\PY{+w}{ }\PY{o}{+}\PY{+w}{ }\PY{n}{bytesRead}\PY{+w}{ }\PY{o}{+}\PY{+w}{ }\PY{l+s}{\PYZdq{}}\PY{l+s}{ bytes}\PY{l+s}{\PYZdq{}}\PY{p}{)}\PY{p}{;}
\PY{+w}{        }\PY{p}{\PYZcb{}}
\PY{+w}{    }\PY{p}{\PYZcb{}}
\PY{p}{\PYZcb{}}
\end{Verbatim}
\end{codebox}

\subsection[ByteBuffer: position, limit, flip, clear]{ByteBuffer: position, limit, flip, clear}
\label{h:10:bytebuffer-position-limit-flip-clear}
A \code{ByteBuffer} tracks three cursors you must understand to use it correctly:

\begin{itemize}
\item 
\textbf{capacity} — total size, fixed at creation.

\item 
\textbf{position} — where the next read/write happens; advances as you go.

\item 
\textbf{limit} — how far you’re allowed to read/write; can’t go past it.

\end{itemize}

\begin{codebox}
\begin{Verbatim}[commandchars=\\\{\}]
\PY{k+kn}{import}\PY{+w}{ }\PY{n+nn}{java.nio.ByteBuffer}\PY{p}{;}
\PY{k+kn}{import}\PY{+w}{ }\PY{n+nn}{java.nio.charset.StandardCharsets}\PY{p}{;}

\PY{n}{ByteBuffer}\PY{+w}{ }\PY{n}{buffer}\PY{+w}{ }\PY{o}{=}\PY{+w}{ }\PY{n}{ByteBuffer}\PY{p}{.}\PY{n+na}{allocate}\PY{p}{(}\PY{l+m+mi}{128}\PY{p}{)}\PY{p}{;}

\PY{c+c1}{// \PYZhy{}\PYZhy{}\PYZhy{} WRITE mode \PYZhy{}\PYZhy{}\PYZhy{}}
\PY{n}{buffer}\PY{p}{.}\PY{n+na}{put}\PY{p}{(}\PY{l+s}{\PYZdq{}}\PY{l+s}{hello}\PY{l+s}{\PYZdq{}}\PY{p}{.}\PY{n+na}{getBytes}\PY{p}{(}\PY{n}{StandardCharsets}\PY{p}{.}\PY{n+na}{UTF\PYZus{}8}\PY{p}{)}\PY{p}{)}\PY{p}{;}
\PY{c+c1}{// position = 5 (just wrote 5 bytes), limit = 128 (capacity)}

\PY{c+c1}{// flip(): switch from writing to reading.}
\PY{c+c1}{// Sets limit = current position, then resets position = 0.}
\PY{n}{buffer}\PY{p}{.}\PY{n+na}{flip}\PY{p}{(}\PY{p}{)}\PY{p}{;}
\PY{c+c1}{// position = 0, limit = 5 \PYZhy{}\PYZhy{} ready to read exactly the 5 bytes we wrote}

\PY{c+c1}{// \PYZhy{}\PYZhy{}\PYZhy{} READ mode \PYZhy{}\PYZhy{}\PYZhy{}}
\PY{k+kt}{byte}\PY{o}{[}\PY{o}{]}\PY{+w}{ }\PY{n}{data}\PY{+w}{ }\PY{o}{=}\PY{+w}{ }\PY{k}{new}\PY{+w}{ }\PY{k+kt}{byte}\PY{o}{[}\PY{n}{buffer}\PY{p}{.}\PY{n+na}{remaining}\PY{p}{(}\PY{p}{)}\PY{o}{]}\PY{p}{;}\PY{+w}{ }\PY{c+c1}{// remaining() = limit \PYZhy{} position}
\PY{n}{buffer}\PY{p}{.}\PY{n+na}{get}\PY{p}{(}\PY{n}{data}\PY{p}{)}\PY{p}{;}
\PY{n}{System}\PY{p}{.}\PY{n+na}{out}\PY{p}{.}\PY{n+na}{println}\PY{p}{(}\PY{k}{new}\PY{+w}{ }\PY{n}{String}\PY{p}{(}\PY{n}{data}\PY{p}{,}\PY{+w}{ }\PY{n}{StandardCharsets}\PY{p}{.}\PY{n+na}{UTF\PYZus{}8}\PY{p}{)}\PY{p}{)}\PY{p}{;}\PY{+w}{ }\PY{c+c1}{// \PYZdq{}hello\PYZdq{}}
\PY{c+c1}{// position = 5, limit = 5}

\PY{c+c1}{// clear(): reset for writing again.}
\PY{c+c1}{// Sets position = 0, limit = capacity. Does NOT erase old data, just forgets it.}
\PY{n}{buffer}\PY{p}{.}\PY{n+na}{clear}\PY{p}{(}\PY{p}{)}\PY{p}{;}
\PY{c+c1}{// position = 0, limit = 128 \PYZhy{}\PYZhy{} ready to write fresh data}
\end{Verbatim}
\end{codebox}

Forgetting to call \code{flip()} before reading is one of the most common NIO bugs: you’ll try to read starting at \code{position} (wherever writing left it) up to \code{limit} (still the full capacity), which usually reads garbage or nothing useful.

\subsection[Direct vs. heap buffers]{Direct vs. heap buffers}
\label{h:10:direct-vs-heap-buffers}
\begin{codebox}
\begin{Verbatim}[commandchars=\\\{\}]
\PY{k+kn}{import}\PY{+w}{ }\PY{n+nn}{java.nio.ByteBuffer}\PY{p}{;}

\PY{n}{ByteBuffer}\PY{+w}{ }\PY{n}{heapBuffer}\PY{+w}{ }\PY{o}{=}\PY{+w}{ }\PY{n}{ByteBuffer}\PY{p}{.}\PY{n+na}{allocate}\PY{p}{(}\PY{l+m+mi}{1024}\PY{p}{)}\PY{p}{;}\PY{+w}{       }\PY{c+c1}{// lives in JVM heap}
\PY{n}{ByteBuffer}\PY{+w}{ }\PY{n}{directBuffer}\PY{+w}{ }\PY{o}{=}\PY{+w}{ }\PY{n}{ByteBuffer}\PY{p}{.}\PY{n+na}{allocateDirect}\PY{p}{(}\PY{l+m+mi}{1024}\PY{p}{)}\PY{p}{;}\PY{+w}{ }\PY{c+c1}{// lives outside the JVM heap}
\end{Verbatim}
\end{codebox}

{\tablefont
\begin{xltabular}{\linewidth}{@{}L{0.444}L{1.291}L{1.265}@{}}
\toprule
\rowcolor{tablehdr}
\thd{} & \thd{Heap buffer (\code{allocate})} & \thd{Direct buffer (\code{allocateDirect})} \\
\midrule
\endhead
Memory location & JVM heap & Native (OS) memory, outside GC heap \\
Allocation cost & Cheap, fast & More expensive to allocate/free \\
I/O performance & JVM may copy to a temp native buffer before an OS call & Can be passed straight to the OS, no extra copy \\
Best for & Small, short-lived, frequently allocated buffers & Large, long-lived buffers reused for repeated I/O \\
GC impact & Collected normally & Freed via cleaner mechanisms, not standard GC \\
\bottomrule
\end{xltabular}
}

For most application code, heap buffers are fine and simpler. Direct buffers pay off for high-throughput I/O (e.g., network servers, bulk file copies) where you reuse the same buffer many times.

\subsection[Memory-mapped files]{Memory-mapped files}
\label{h:10:memory-mapped-files}
\code{FileChannel.\allowbreak{}map()} maps a file’s contents directly into memory, so you can treat file bytes like an in-memory array — the OS handles paging data in and out as you access it. This is great for large files you need random access into (e.g., a large binary index) without reading the whole thing.

\begin{codebox}
\begin{Verbatim}[commandchars=\\\{\}]
\PY{k+kn}{import}\PY{+w}{ }\PY{n+nn}{java.io.IOException}\PY{p}{;}
\PY{k+kn}{import}\PY{+w}{ }\PY{n+nn}{java.nio.MappedByteBuffer}\PY{p}{;}
\PY{k+kn}{import}\PY{+w}{ }\PY{n+nn}{java.nio.channels.FileChannel}\PY{p}{;}
\PY{k+kn}{import}\PY{+w}{ }\PY{n+nn}{java.nio.file.Path}\PY{p}{;}
\PY{k+kn}{import}\PY{+w}{ }\PY{n+nn}{java.nio.file.StandardOpenOption}\PY{p}{;}

\PY{k}{try}\PY{+w}{ }\PY{p}{(}\PY{n}{FileChannel}\PY{+w}{ }\PY{n}{channel}\PY{+w}{ }\PY{o}{=}\PY{+w}{ }\PY{n}{FileChannel}\PY{p}{.}\PY{n+na}{open}\PY{p}{(}
\PY{+w}{        }\PY{n}{Path}\PY{p}{.}\PY{n+na}{of}\PY{p}{(}\PY{l+s}{\PYZdq{}}\PY{l+s}{large\PYZhy{}file.dat}\PY{l+s}{\PYZdq{}}\PY{p}{)}\PY{p}{,}\PY{+w}{ }\PY{n}{StandardOpenOption}\PY{p}{.}\PY{n+na}{READ}\PY{p}{)}\PY{p}{)}\PY{+w}{ }\PY{p}{\PYZob{}}

\PY{+w}{    }\PY{n}{MappedByteBuffer}\PY{+w}{ }\PY{n}{mapped}\PY{+w}{ }\PY{o}{=}\PY{+w}{ }\PY{n}{channel}\PY{p}{.}\PY{n+na}{map}\PY{p}{(}
\PY{+w}{            }\PY{n}{FileChannel}\PY{p}{.}\PY{n+na}{MapMode}\PY{p}{.}\PY{n+na}{READ\PYZus{}ONLY}\PY{p}{,}\PY{+w}{ }\PY{l+m+mi}{0}\PY{p}{,}\PY{+w}{ }\PY{n}{channel}\PY{p}{.}\PY{n+na}{size}\PY{p}{(}\PY{p}{)}\PY{p}{)}\PY{p}{;}

\PY{+w}{    }\PY{c+c1}{// Access any byte at any offset directly \PYZhy{}\PYZhy{} OS pages it in as needed}
\PY{+w}{    }\PY{k+kt}{byte}\PY{+w}{ }\PY{n}{firstByte}\PY{+w}{ }\PY{o}{=}\PY{+w}{ }\PY{n}{mapped}\PY{p}{.}\PY{n+na}{get}\PY{p}{(}\PY{l+m+mi}{0}\PY{p}{)}\PY{p}{;}
\PY{+w}{    }\PY{k+kt}{byte}\PY{+w}{ }\PY{n}{lastByte}\PY{+w}{ }\PY{o}{=}\PY{+w}{ }\PY{n}{mapped}\PY{p}{.}\PY{n+na}{get}\PY{p}{(}\PY{p}{(}\PY{k+kt}{int}\PY{p}{)}\PY{+w}{ }\PY{n}{channel}\PY{p}{.}\PY{n+na}{size}\PY{p}{(}\PY{p}{)}\PY{+w}{ }\PY{o}{\PYZhy{}}\PY{+w}{ }\PY{l+m+mi}{1}\PY{p}{)}\PY{p}{;}
\PY{p}{\PYZcb{}}
\end{Verbatim}
\end{codebox}

Caveats: mapped regions are capped by addressable memory (practically, \textasciitilde{}2 GB per mapping on many JVMs due to \code{int} indexing), and unmapping is not guaranteed to happen promptly — on some platforms the mapped file may stay locked until the buffer is garbage collected.

\subsection[Selector and non-blocking I/O]{Selector and non-blocking I/O}
\label{h:10:selector-and-non-blocking-io}
Regular blocking I/O ties up one thread per connection: the thread calls \code{read()} and just sits there until data arrives. That doesn’t scale to thousands of connections. NIO’s \code{Selector} lets a \textbf{single thread monitor many channels at once}, and only wakes up to handle the channels that are actually ready (readable, writable, or have a new connection) — this is the classic “one thread, many sockets” event-loop model used by things like Netty. In modern Java, this specific problem is often solved more simply with \textbf{virtual threads} (Project Loom, Java 21) instead of hand-rolled \code{Selector} loops, but understanding \code{Selector} still matters for reading legacy networking code and frameworks built on it.

\subsection[FileChannel.transferTo (zero-copy)]{FileChannel.transferTo (zero-copy)}
\label{h:10:filechanneltransferto-zero-copy}
\code{transferTo} can copy data directly between channels, and on many operating systems this happens via a \textbf{zero-copy} OS call — the data never has to pass through a Java buffer in user space at all.

\begin{codebox}
\begin{Verbatim}[commandchars=\\\{\}]
\PY{k+kn}{import}\PY{+w}{ }\PY{n+nn}{java.io.IOException}\PY{p}{;}
\PY{k+kn}{import}\PY{+w}{ }\PY{n+nn}{java.nio.channels.FileChannel}\PY{p}{;}
\PY{k+kn}{import}\PY{+w}{ }\PY{n+nn}{java.nio.file.Path}\PY{p}{;}
\PY{k+kn}{import}\PY{+w}{ }\PY{n+nn}{java.nio.file.StandardOpenOption}\PY{p}{;}

\PY{k}{try}\PY{+w}{ }\PY{p}{(}\PY{n}{FileChannel}\PY{+w}{ }\PY{n}{source}\PY{+w}{ }\PY{o}{=}\PY{+w}{ }\PY{n}{FileChannel}\PY{p}{.}\PY{n+na}{open}\PY{p}{(}\PY{n}{Path}\PY{p}{.}\PY{n+na}{of}\PY{p}{(}\PY{l+s}{\PYZdq{}}\PY{l+s}{source.dat}\PY{l+s}{\PYZdq{}}\PY{p}{)}\PY{p}{,}\PY{+w}{ }\PY{n}{StandardOpenOption}\PY{p}{.}\PY{n+na}{READ}\PY{p}{)}\PY{p}{;}
\PY{+w}{     }\PY{n}{FileChannel}\PY{+w}{ }\PY{n}{destination}\PY{+w}{ }\PY{o}{=}\PY{+w}{ }\PY{n}{FileChannel}\PY{p}{.}\PY{n+na}{open}\PY{p}{(}
\PY{+w}{             }\PY{n}{Path}\PY{p}{.}\PY{n+na}{of}\PY{p}{(}\PY{l+s}{\PYZdq{}}\PY{l+s}{destination.dat}\PY{l+s}{\PYZdq{}}\PY{p}{)}\PY{p}{,}
\PY{+w}{             }\PY{n}{StandardOpenOption}\PY{p}{.}\PY{n+na}{CREATE}\PY{p}{,}\PY{+w}{ }\PY{n}{StandardOpenOption}\PY{p}{.}\PY{n+na}{WRITE}\PY{p}{)}\PY{p}{)}\PY{+w}{ }\PY{p}{\PYZob{}}

\PY{+w}{    }\PY{k+kt}{long}\PY{+w}{ }\PY{n}{transferred}\PY{+w}{ }\PY{o}{=}\PY{+w}{ }\PY{l+m+mi}{0}\PY{p}{;}
\PY{+w}{    }\PY{k+kt}{long}\PY{+w}{ }\PY{n}{size}\PY{+w}{ }\PY{o}{=}\PY{+w}{ }\PY{n}{source}\PY{p}{.}\PY{n+na}{size}\PY{p}{(}\PY{p}{)}\PY{p}{;}
\PY{+w}{    }\PY{k}{while}\PY{+w}{ }\PY{p}{(}\PY{n}{transferred}\PY{+w}{ }\PY{o}{\PYZlt{}}\PY{+w}{ }\PY{n}{size}\PY{p}{)}\PY{+w}{ }\PY{p}{\PYZob{}}
\PY{+w}{        }\PY{n}{transferred}\PY{+w}{ }\PY{o}{+}\PY{o}{=}\PY{+w}{ }\PY{n}{source}\PY{p}{.}\PY{n+na}{transferTo}\PY{p}{(}\PY{n}{transferred}\PY{p}{,}\PY{+w}{ }\PY{n}{size}\PY{+w}{ }\PY{o}{\PYZhy{}}\PY{+w}{ }\PY{n}{transferred}\PY{p}{,}\PY{+w}{ }\PY{n}{destination}\PY{p}{)}\PY{p}{;}
\PY{+w}{    }\PY{p}{\PYZcb{}}
\PY{p}{\PYZcb{}}
\end{Verbatim}
\end{codebox}

For plain file-to-file copies, though, \code{Files.\allowbreak{}copy(\allowbreak{}Path,\allowbreak{}~\allowbreak{}Path,\allowbreak{}~\allowbreak{}...)} (NIO.2) is simpler and just as efficient in practice — reach for \code{transferTo} mainly when you’re moving data channel-to-channel (e.g., file to socket) rather than file to file.

\subsection[When NIO actually helps]{When NIO actually helps}
\label{h:10:when-nio-actually-helps}
{\tablefont
\begin{xltabular}{\linewidth}{@{}L{1.129}L{0.871}@{}}
\toprule
\rowcolor{tablehdr}
\thd{Use case} & \thd{Better fit} \\
\midrule
\endhead
Simple file read/write in a small tool or CLI & \code{java.\allowbreak{}io} / \code{Files} (NIO.2) — simpler code \\
Large file, need random access to specific offsets & NIO \code{FileChannel} + memory mapping \\
High-throughput network server, many connections & NIO \code{Selector}, or virtual threads \\
Bulk copy between files or channels & \code{FileChannel.\allowbreak{}transferTo} / \code{Files.\allowbreak{}copy} \\
Everyday text/config file processing & \code{Files.\allowbreak{}readString}, \code{Files.\allowbreak{}lines} (NIO.2) \\
\bottomrule
\end{xltabular}
}

\textbf{Reviewer’s rule of thumb}: don’t reach for raw \code{ByteBuffer}/\code{Channel}/\code{Selector} code unless there’s a real reason (performance-critical I/O, non-blocking networking, memory-mapped access). For ordinary file reading and writing, NIO.2’s \code{Files} API (next section) is simpler, less error-prone, and just as fast.

\section[Paths and Files]{Paths and Files}
\label{h:10:paths-and-files}
NIO.2 (Java 7+) introduced \code{java.\allowbreak{}nio.\allowbreak{}file.\allowbreak{}Path} and \code{java.\allowbreak{}nio.\allowbreak{}file.\allowbreak{}Files} to replace most uses of \code{java.\allowbreak{}io.\allowbreak{}File}. \code{Path} represents a filesystem location (it doesn’t have to exist); \code{Files} is a utility class full of static methods that operate on \code{Path}s.

{\tablefont
\begin{xltabular}{\linewidth}{@{}L{0.612}L{0.874}L{1.515}@{}}
\toprule
\rowcolor{tablehdr}
\thd{} & \thd{\code{java.\allowbreak{}io.\allowbreak{}File}} & \thd{NIO.2 (\code{Path} / \code{Files})} \\
\midrule
\endhead
Error reporting & Booleans (\code{false} on failure, no reason) & Specific exceptions (\code{NoSuchFileException}, etc.) \\
Symbolic link support & Poor / inconsistent & First-class (\code{Files.\allowbreak{}isSymbolicLink}, etc.) \\
File attributes & Limited (\code{lastModified}, \code{length}) & Rich (\code{BasicFileAttributes}, POSIX permissions, owner) \\
Directory walking & Manual recursion & \code{Files.\allowbreak{}walk}, \code{Files.\allowbreak{}walkFileTree} \\
Watching for changes & Not supported & \code{WatchService} \\
Reading whole file as text & Manual stream wiring & \code{Files.\allowbreak{}readString} (one line) \\
Extensibility & Tied to default filesystem & Pluggable filesystem providers (e.g., zip, in-memory) \\
\bottomrule
\end{xltabular}
}

\subsection[Creating and manipulating Paths]{Creating and manipulating Paths}
\label{h:10:creating-and-manipulating-paths}
\begin{codebox}
\begin{Verbatim}[commandchars=\\\{\}]
\PY{k+kn}{import}\PY{+w}{ }\PY{n+nn}{java.nio.file.Path}\PY{p}{;}
\PY{k+kn}{import}\PY{+w}{ }\PY{n+nn}{java.nio.file.Paths}\PY{p}{;}

\PY{c+c1}{// Two equivalent ways to build a Path. Path.of() is preferred since Java 11}
\PY{c+c1}{// (Paths.get() still works and is common in older code / examples).}
\PY{n}{Path}\PY{+w}{ }\PY{n}{p1}\PY{+w}{ }\PY{o}{=}\PY{+w}{ }\PY{n}{Paths}\PY{p}{.}\PY{n+na}{get}\PY{p}{(}\PY{l+s}{\PYZdq{}}\PY{l+s}{data}\PY{l+s}{\PYZdq{}}\PY{p}{,}\PY{+w}{ }\PY{l+s}{\PYZdq{}}\PY{l+s}{input.csv}\PY{l+s}{\PYZdq{}}\PY{p}{)}\PY{p}{;}
\PY{n}{Path}\PY{+w}{ }\PY{n}{p2}\PY{+w}{ }\PY{o}{=}\PY{+w}{ }\PY{n}{Path}\PY{p}{.}\PY{n+na}{of}\PY{p}{(}\PY{l+s}{\PYZdq{}}\PY{l+s}{data}\PY{l+s}{\PYZdq{}}\PY{p}{,}\PY{+w}{ }\PY{l+s}{\PYZdq{}}\PY{l+s}{input.csv}\PY{l+s}{\PYZdq{}}\PY{p}{)}\PY{p}{;}

\PY{n}{System}\PY{p}{.}\PY{n+na}{out}\PY{p}{.}\PY{n+na}{println}\PY{p}{(}\PY{n}{p1}\PY{p}{.}\PY{n+na}{equals}\PY{p}{(}\PY{n}{p2}\PY{p}{)}\PY{p}{)}\PY{p}{;}\PY{+w}{ }\PY{c+c1}{// true}

\PY{n}{Path}\PY{+w}{ }\PY{n}{base}\PY{+w}{ }\PY{o}{=}\PY{+w}{ }\PY{n}{Path}\PY{p}{.}\PY{n+na}{of}\PY{p}{(}\PY{l+s}{\PYZdq{}}\PY{l+s}{/home/user/project}\PY{l+s}{\PYZdq{}}\PY{p}{)}\PY{p}{;}
\PY{n}{Path}\PY{+w}{ }\PY{n}{resolved}\PY{+w}{ }\PY{o}{=}\PY{+w}{ }\PY{n}{base}\PY{p}{.}\PY{n+na}{resolve}\PY{p}{(}\PY{l+s}{\PYZdq{}}\PY{l+s}{src/Main.java}\PY{l+s}{\PYZdq{}}\PY{p}{)}\PY{p}{;}
\PY{n}{System}\PY{p}{.}\PY{n+na}{out}\PY{p}{.}\PY{n+na}{println}\PY{p}{(}\PY{n}{resolved}\PY{p}{)}\PY{p}{;}\PY{+w}{ }\PY{c+c1}{// /home/user/project/src/Main.java}

\PY{n}{Path}\PY{+w}{ }\PY{n}{relative}\PY{+w}{ }\PY{o}{=}\PY{+w}{ }\PY{n}{base}\PY{p}{.}\PY{n+na}{relativize}\PY{p}{(}\PY{n}{Path}\PY{p}{.}\PY{n+na}{of}\PY{p}{(}\PY{l+s}{\PYZdq{}}\PY{l+s}{/home/user/other/lib.jar}\PY{l+s}{\PYZdq{}}\PY{p}{)}\PY{p}{)}\PY{p}{;}
\PY{n}{System}\PY{p}{.}\PY{n+na}{out}\PY{p}{.}\PY{n+na}{println}\PY{p}{(}\PY{n}{relative}\PY{p}{)}\PY{p}{;}\PY{+w}{ }\PY{c+c1}{// ../other/lib.jar}

\PY{n}{Path}\PY{+w}{ }\PY{n}{messy}\PY{+w}{ }\PY{o}{=}\PY{+w}{ }\PY{n}{Path}\PY{p}{.}\PY{n+na}{of}\PY{p}{(}\PY{l+s}{\PYZdq{}}\PY{l+s}{/home/user/./project/../project/src}\PY{l+s}{\PYZdq{}}\PY{p}{)}\PY{p}{;}
\PY{n}{System}\PY{p}{.}\PY{n+na}{out}\PY{p}{.}\PY{n+na}{println}\PY{p}{(}\PY{n}{messy}\PY{p}{.}\PY{n+na}{normalize}\PY{p}{(}\PY{p}{)}\PY{p}{)}\PY{p}{;}\PY{+w}{ }\PY{c+c1}{// /home/user/project/src}

\PY{n}{Path}\PY{+w}{ }\PY{n}{relPath}\PY{+w}{ }\PY{o}{=}\PY{+w}{ }\PY{n}{Path}\PY{p}{.}\PY{n+na}{of}\PY{p}{(}\PY{l+s}{\PYZdq{}}\PY{l+s}{relative/path}\PY{l+s}{\PYZdq{}}\PY{p}{)}\PY{p}{;}
\PY{n}{System}\PY{p}{.}\PY{n+na}{out}\PY{p}{.}\PY{n+na}{println}\PY{p}{(}\PY{n}{relPath}\PY{p}{.}\PY{n+na}{toAbsolutePath}\PY{p}{(}\PY{p}{)}\PY{p}{)}\PY{p}{;}\PY{+w}{ }\PY{c+c1}{// resolved against the current working directory}
\end{Verbatim}
\end{codebox}

\begin{itemize}
\item 
\code{resolve} — joins a path onto another (like appending a subpath).

\item 
\code{relativize} — computes the relative path needed to get from one path to another.

\item 
\code{normalize} — removes redundant \code{.} and \code{..} segments \emph{syntactically}, without touching the filesystem.

\item 
\code{toAbsolutePath} — anchors a relative path to the current working directory (still doesn’t check if it exists).

\end{itemize}

\subsection[Reading and writing whole files]{Reading and writing whole files}
\label{h:10:reading-and-writing-whole-files}
\begin{codebox}
\begin{Verbatim}[commandchars=\\\{\}]
\PY{k+kn}{import}\PY{+w}{ }\PY{n+nn}{java.io.IOException}\PY{p}{;}
\PY{k+kn}{import}\PY{+w}{ }\PY{n+nn}{java.nio.charset.StandardCharsets}\PY{p}{;}
\PY{k+kn}{import}\PY{+w}{ }\PY{n+nn}{java.nio.file.Files}\PY{p}{;}
\PY{k+kn}{import}\PY{+w}{ }\PY{n+nn}{java.nio.file.Path}\PY{p}{;}
\PY{k+kn}{import}\PY{+w}{ }\PY{n+nn}{java.util.List}\PY{p}{;}
\PY{k+kn}{import}\PY{+w}{ }\PY{n+nn}{java.util.stream.Stream}\PY{p}{;}

\PY{n}{Path}\PY{+w}{ }\PY{n}{file}\PY{+w}{ }\PY{o}{=}\PY{+w}{ }\PY{n}{Path}\PY{p}{.}\PY{n+na}{of}\PY{p}{(}\PY{l+s}{\PYZdq{}}\PY{l+s}{notes.txt}\PY{l+s}{\PYZdq{}}\PY{p}{)}\PY{p}{;}

\PY{c+c1}{// Read the entire file as one String (Java 11+)}
\PY{n}{String}\PY{+w}{ }\PY{n}{content}\PY{+w}{ }\PY{o}{=}\PY{+w}{ }\PY{n}{Files}\PY{p}{.}\PY{n+na}{readString}\PY{p}{(}\PY{n}{file}\PY{p}{,}\PY{+w}{ }\PY{n}{StandardCharsets}\PY{p}{.}\PY{n+na}{UTF\PYZus{}8}\PY{p}{)}\PY{p}{;}

\PY{c+c1}{// Write a String to a file, creating or truncating it (Java 11+)}
\PY{n}{Files}\PY{p}{.}\PY{n+na}{writeString}\PY{p}{(}\PY{n}{file}\PY{p}{,}\PY{+w}{ }\PY{l+s}{\PYZdq{}}\PY{l+s}{Hello, NIO.2!}\PY{l+s}{\PYZbs{}}\PY{l+s}{n}\PY{l+s}{\PYZdq{}}\PY{p}{,}\PY{+w}{ }\PY{n}{StandardCharsets}\PY{p}{.}\PY{n+na}{UTF\PYZus{}8}\PY{p}{)}\PY{p}{;}

\PY{c+c1}{// Read all lines into a List\PYZlt{}String\PYZgt{} \PYZhy{}\PYZhy{} fine for small/medium files}
\PY{n}{List}\PY{o}{\PYZlt{}}\PY{n}{String}\PY{o}{\PYZgt{}}\PY{+w}{ }\PY{n}{lines}\PY{+w}{ }\PY{o}{=}\PY{+w}{ }\PY{n}{Files}\PY{p}{.}\PY{n+na}{readAllLines}\PY{p}{(}\PY{n}{file}\PY{p}{,}\PY{+w}{ }\PY{n}{StandardCharsets}\PY{p}{.}\PY{n+na}{UTF\PYZus{}8}\PY{p}{)}\PY{p}{;}

\PY{c+c1}{// Stream lines lazily \PYZhy{}\PYZhy{} for large files, so you don\PYZsq{}t load everything into memory.}
\PY{c+c1}{// IMPORTANT: Files.lines() opens a file handle that MUST be closed.}
\PY{k}{try}\PY{+w}{ }\PY{p}{(}\PY{n}{Stream}\PY{o}{\PYZlt{}}\PY{n}{String}\PY{o}{\PYZgt{}}\PY{+w}{ }\PY{n}{lineStream}\PY{+w}{ }\PY{o}{=}\PY{+w}{ }\PY{n}{Files}\PY{p}{.}\PY{n+na}{lines}\PY{p}{(}\PY{n}{file}\PY{p}{,}\PY{+w}{ }\PY{n}{StandardCharsets}\PY{p}{.}\PY{n+na}{UTF\PYZus{}8}\PY{p}{)}\PY{p}{)}\PY{+w}{ }\PY{p}{\PYZob{}}
\PY{+w}{    }\PY{k+kt}{long}\PY{+w}{ }\PY{n}{count}\PY{+w}{ }\PY{o}{=}\PY{+w}{ }\PY{n}{lineStream}\PY{p}{.}\PY{n+na}{filter}\PY{p}{(}\PY{n}{line}\PY{+w}{ }\PY{o}{\PYZhy{}}\PY{o}{\PYZgt{}}\PY{+w}{ }\PY{o}{!}\PY{n}{line}\PY{p}{.}\PY{n+na}{isBlank}\PY{p}{(}\PY{p}{)}\PY{p}{)}\PY{p}{.}\PY{n+na}{count}\PY{p}{(}\PY{p}{)}\PY{p}{;}
\PY{+w}{    }\PY{n}{System}\PY{p}{.}\PY{n+na}{out}\PY{p}{.}\PY{n+na}{println}\PY{p}{(}\PY{l+s}{\PYZdq{}}\PY{l+s}{Non\PYZhy{}blank lines: }\PY{l+s}{\PYZdq{}}\PY{+w}{ }\PY{o}{+}\PY{+w}{ }\PY{n}{count}\PY{p}{)}\PY{p}{;}
\PY{p}{\PYZcb{}}\PY{+w}{ }\PY{c+c1}{// stream (and its underlying file handle) closed here}
\end{Verbatim}
\end{codebox}

\code{Files.\allowbreak{}lines()} returns a lazily-evaluated \code{Stream\textless{}\allowbreak{}String\textgreater{}} backed by an open file descriptor. If you don’t close it (via try-with-resources, since \code{Stream} implements \code{AutoCloseable}), the file handle leaks — this is easy to miss because a \code{Stream} doesn’t \emph{look} like a resource the way an \code{InputStream} obviously does.

\begin{codebox}
\begin{Verbatim}[commandchars=\\\{\}]
\PY{k+kn}{import}\PY{+w}{ }\PY{n+nn}{java.io.BufferedReader}\PY{p}{;}
\PY{k+kn}{import}\PY{+w}{ }\PY{n+nn}{java.io.IOException}\PY{p}{;}
\PY{k+kn}{import}\PY{+w}{ }\PY{n+nn}{java.nio.charset.StandardCharsets}\PY{p}{;}
\PY{k+kn}{import}\PY{+w}{ }\PY{n+nn}{java.nio.file.Files}\PY{p}{;}
\PY{k+kn}{import}\PY{+w}{ }\PY{n+nn}{java.nio.file.Path}\PY{p}{;}

\PY{c+c1}{// Files.newBufferedReader gives you a fully\PYZhy{}buffered Reader with a chosen charset,}
\PY{c+c1}{// useful when you want line\PYZhy{}by\PYZhy{}line control instead of a Stream or a full readAllLines()}
\PY{k}{try}\PY{+w}{ }\PY{p}{(}\PY{n}{BufferedReader}\PY{+w}{ }\PY{n}{reader}\PY{+w}{ }\PY{o}{=}\PY{+w}{ }\PY{n}{Files}\PY{p}{.}\PY{n+na}{newBufferedReader}\PY{p}{(}\PY{n}{Path}\PY{p}{.}\PY{n+na}{of}\PY{p}{(}\PY{l+s}{\PYZdq{}}\PY{l+s}{big.log}\PY{l+s}{\PYZdq{}}\PY{p}{)}\PY{p}{,}\PY{+w}{ }\PY{n}{StandardCharsets}\PY{p}{.}\PY{n+na}{UTF\PYZus{}8}\PY{p}{)}\PY{p}{)}\PY{+w}{ }\PY{p}{\PYZob{}}
\PY{+w}{    }\PY{n}{String}\PY{+w}{ }\PY{n}{line}\PY{p}{;}
\PY{+w}{    }\PY{k}{while}\PY{+w}{ }\PY{p}{(}\PY{p}{(}\PY{n}{line}\PY{+w}{ }\PY{o}{=}\PY{+w}{ }\PY{n}{reader}\PY{p}{.}\PY{n+na}{readLine}\PY{p}{(}\PY{p}{)}\PY{p}{)}\PY{+w}{ }\PY{o}{!}\PY{o}{=}\PY{+w}{ }\PY{k+kc}{null}\PY{p}{)}\PY{+w}{ }\PY{p}{\PYZob{}}
\PY{+w}{        }\PY{c+c1}{// process one line at a time, low memory overhead}
\PY{+w}{    }\PY{p}{\PYZcb{}}
\PY{p}{\PYZcb{}}
\end{Verbatim}
\end{codebox}

\subsection[Copy, move, delete]{Copy, move, delete}
\label{h:10:copy-move-delete}
\begin{codebox}
\begin{Verbatim}[commandchars=\\\{\}]
\PY{k+kn}{import}\PY{+w}{ }\PY{n+nn}{java.io.IOException}\PY{p}{;}
\PY{k+kn}{import}\PY{+w}{ }\PY{n+nn}{java.nio.file.CopyOption}\PY{p}{;}
\PY{k+kn}{import}\PY{+w}{ }\PY{n+nn}{java.nio.file.Files}\PY{p}{;}
\PY{k+kn}{import}\PY{+w}{ }\PY{n+nn}{java.nio.file.Path}\PY{p}{;}
\PY{k+kn}{import}\PY{+w}{ }\PY{n+nn}{java.nio.file.StandardCopyOption}\PY{p}{;}

\PY{n}{Path}\PY{+w}{ }\PY{n}{source}\PY{+w}{ }\PY{o}{=}\PY{+w}{ }\PY{n}{Path}\PY{p}{.}\PY{n+na}{of}\PY{p}{(}\PY{l+s}{\PYZdq{}}\PY{l+s}{report.pdf}\PY{l+s}{\PYZdq{}}\PY{p}{)}\PY{p}{;}
\PY{n}{Path}\PY{+w}{ }\PY{n}{target}\PY{+w}{ }\PY{o}{=}\PY{+w}{ }\PY{n}{Path}\PY{p}{.}\PY{n+na}{of}\PY{p}{(}\PY{l+s}{\PYZdq{}}\PY{l+s}{backup/report.pdf}\PY{l+s}{\PYZdq{}}\PY{p}{)}\PY{p}{;}

\PY{c+c1}{// Copy, overwriting the target if it already exists}
\PY{n}{Files}\PY{p}{.}\PY{n+na}{copy}\PY{p}{(}\PY{n}{source}\PY{p}{,}\PY{+w}{ }\PY{n}{target}\PY{p}{,}\PY{+w}{ }\PY{n}{StandardCopyOption}\PY{p}{.}\PY{n+na}{REPLACE\PYZus{}EXISTING}\PY{p}{)}\PY{p}{;}

\PY{c+c1}{// Move (rename) a file; ATOMIC\PYZus{}MOVE guarantees no partial/half\PYZhy{}written result is ever visible}
\PY{n}{Files}\PY{p}{.}\PY{n+na}{move}\PY{p}{(}\PY{n}{source}\PY{p}{,}\PY{+w}{ }\PY{n}{target}\PY{p}{,}\PY{+w}{ }\PY{n}{StandardCopyOption}\PY{p}{.}\PY{n+na}{ATOMIC\PYZus{}MOVE}\PY{p}{)}\PY{p}{;}

\PY{c+c1}{// Delete: throws NoSuchFileException if missing (unlike File.delete()\PYZsq{}s silent `false`)}
\PY{n}{Files}\PY{p}{.}\PY{n+na}{delete}\PY{p}{(}\PY{n}{target}\PY{p}{)}\PY{p}{;}

\PY{c+c1}{// deleteIfExists: same, but does nothing (no exception) if the file is already gone}
\PY{n}{Files}\PY{p}{.}\PY{n+na}{deleteIfExists}\PY{p}{(}\PY{n}{target}\PY{p}{)}\PY{p}{;}
\end{Verbatim}
\end{codebox}

\subsection[Files.walk vs. Files.walkFileTree]{Files.walk vs. Files.walkFileTree}
\label{h:10:fileswalk-vs-fileswalkfiletree}
Both traverse a directory tree, but they differ in flexibility and resource handling.

\begin{codebox}
\begin{Verbatim}[commandchars=\\\{\}]
\PY{k+kn}{import}\PY{+w}{ }\PY{n+nn}{java.io.IOException}\PY{p}{;}
\PY{k+kn}{import}\PY{+w}{ }\PY{n+nn}{java.nio.file.*}\PY{p}{;}
\PY{k+kn}{import}\PY{+w}{ }\PY{n+nn}{java.nio.file.attribute.BasicFileAttributes}\PY{p}{;}
\PY{k+kn}{import}\PY{+w}{ }\PY{n+nn}{java.util.stream.Stream}\PY{p}{;}

\PY{c+c1}{// Files.walk: returns a lazy Stream\PYZlt{}Path\PYZgt{} of every file/dir under root.}
\PY{c+c1}{// Simple for filtering/collecting, but MUST be closed (it holds an open directory handle).}
\PY{k}{try}\PY{+w}{ }\PY{p}{(}\PY{n}{Stream}\PY{o}{\PYZlt{}}\PY{n}{Path}\PY{o}{\PYZgt{}}\PY{+w}{ }\PY{n}{paths}\PY{+w}{ }\PY{o}{=}\PY{+w}{ }\PY{n}{Files}\PY{p}{.}\PY{n+na}{walk}\PY{p}{(}\PY{n}{Path}\PY{p}{.}\PY{n+na}{of}\PY{p}{(}\PY{l+s}{\PYZdq{}}\PY{l+s}{src}\PY{l+s}{\PYZdq{}}\PY{p}{)}\PY{p}{)}\PY{p}{)}\PY{+w}{ }\PY{p}{\PYZob{}}
\PY{+w}{    }\PY{n}{paths}\PY{p}{.}\PY{n+na}{filter}\PY{p}{(}\PY{n}{p}\PY{+w}{ }\PY{o}{\PYZhy{}}\PY{o}{\PYZgt{}}\PY{+w}{ }\PY{n}{p}\PY{p}{.}\PY{n+na}{toString}\PY{p}{(}\PY{p}{)}\PY{p}{.}\PY{n+na}{endsWith}\PY{p}{(}\PY{l+s}{\PYZdq{}}\PY{l+s}{.java}\PY{l+s}{\PYZdq{}}\PY{p}{)}\PY{p}{)}
\PY{+w}{         }\PY{p}{.}\PY{n+na}{forEach}\PY{p}{(}\PY{n}{System}\PY{p}{.}\PY{n+na}{out}\PY{p}{:}\PY{p}{:}\PY{n}{println}\PY{p}{)}\PY{p}{;}
\PY{p}{\PYZcb{}}

\PY{c+c1}{// Files.walkFileTree: callback\PYZhy{}based (visitor pattern), more control,}
\PY{c+c1}{// e.g. skipping whole subtrees, handling errors per\PYZhy{}file, computing directory sizes.}
\PY{n}{Files}\PY{p}{.}\PY{n+na}{walkFileTree}\PY{p}{(}\PY{n}{Path}\PY{p}{.}\PY{n+na}{of}\PY{p}{(}\PY{l+s}{\PYZdq{}}\PY{l+s}{src}\PY{l+s}{\PYZdq{}}\PY{p}{)}\PY{p}{,}\PY{+w}{ }\PY{k}{new}\PY{+w}{ }\PY{n}{SimpleFileVisitor}\PY{o}{\PYZlt{}}\PY{n}{Path}\PY{o}{\PYZgt{}}\PY{p}{(}\PY{p}{)}\PY{+w}{ }\PY{p}{\PYZob{}}
\PY{+w}{    }\PY{n+nd}{@Override}
\PY{+w}{    }\PY{k+kd}{public}\PY{+w}{ }\PY{n}{FileVisitResult}\PY{+w}{ }\PY{n+nf}{visitFile}\PY{p}{(}\PY{n}{Path}\PY{+w}{ }\PY{n}{file}\PY{p}{,}\PY{+w}{ }\PY{n}{BasicFileAttributes}\PY{+w}{ }\PY{n}{attrs}\PY{p}{)}\PY{+w}{ }\PY{p}{\PYZob{}}
\PY{+w}{        }\PY{k}{if}\PY{+w}{ }\PY{p}{(}\PY{n}{file}\PY{p}{.}\PY{n+na}{toString}\PY{p}{(}\PY{p}{)}\PY{p}{.}\PY{n+na}{endsWith}\PY{p}{(}\PY{l+s}{\PYZdq{}}\PY{l+s}{.class}\PY{l+s}{\PYZdq{}}\PY{p}{)}\PY{p}{)}\PY{+w}{ }\PY{p}{\PYZob{}}
\PY{+w}{            }\PY{k}{return}\PY{+w}{ }\PY{n}{FileVisitResult}\PY{p}{.}\PY{n+na}{CONTINUE}\PY{p}{;}\PY{+w}{ }\PY{c+c1}{// could delete, count, etc.}
\PY{+w}{        }\PY{p}{\PYZcb{}}
\PY{+w}{        }\PY{k}{return}\PY{+w}{ }\PY{n}{FileVisitResult}\PY{p}{.}\PY{n+na}{CONTINUE}\PY{p}{;}
\PY{+w}{    }\PY{p}{\PYZcb{}}

\PY{+w}{    }\PY{n+nd}{@Override}
\PY{+w}{    }\PY{k+kd}{public}\PY{+w}{ }\PY{n}{FileVisitResult}\PY{+w}{ }\PY{n+nf}{visitFileFailed}\PY{p}{(}\PY{n}{Path}\PY{+w}{ }\PY{n}{file}\PY{p}{,}\PY{+w}{ }\PY{n}{IOException}\PY{+w}{ }\PY{n}{exc}\PY{p}{)}\PY{+w}{ }\PY{p}{\PYZob{}}
\PY{+w}{        }\PY{n}{System}\PY{p}{.}\PY{n+na}{err}\PY{p}{.}\PY{n+na}{println}\PY{p}{(}\PY{l+s}{\PYZdq{}}\PY{l+s}{Could not visit: }\PY{l+s}{\PYZdq{}}\PY{+w}{ }\PY{o}{+}\PY{+w}{ }\PY{n}{file}\PY{+w}{ }\PY{o}{+}\PY{+w}{ }\PY{l+s}{\PYZdq{}}\PY{l+s}{ (}\PY{l+s}{\PYZdq{}}\PY{+w}{ }\PY{o}{+}\PY{+w}{ }\PY{n}{exc}\PY{+w}{ }\PY{o}{+}\PY{+w}{ }\PY{l+s}{\PYZdq{}}\PY{l+s}{)}\PY{l+s}{\PYZdq{}}\PY{p}{)}\PY{p}{;}
\PY{+w}{        }\PY{k}{return}\PY{+w}{ }\PY{n}{FileVisitResult}\PY{p}{.}\PY{n+na}{CONTINUE}\PY{p}{;}\PY{+w}{ }\PY{c+c1}{// don\PYZsq{}t abort the whole walk on one bad file}
\PY{+w}{    }\PY{p}{\PYZcb{}}
\PY{p}{\PYZcb{}}\PY{p}{)}\PY{p}{;}
\end{Verbatim}
\end{codebox}

Use \code{Files.\allowbreak{}walk} for quick stream-style filtering/collecting. Use \code{Files.\allowbreak{}walkFileTree} when you need fine-grained control (skip directories, react to errors per file, avoid loading the whole path list into a stream pipeline first).

\subsection[Temp files, DirectoryStream, WatchService, attributes]{Temp files, DirectoryStream, WatchService, attributes}
\label{h:10:temp-files-directorystream-watchservice-attributes}
\begin{codebox}
\begin{Verbatim}[commandchars=\\\{\}]
\PY{k+kn}{import}\PY{+w}{ }\PY{n+nn}{java.io.IOException}\PY{p}{;}
\PY{k+kn}{import}\PY{+w}{ }\PY{n+nn}{java.nio.file.DirectoryStream}\PY{p}{;}
\PY{k+kn}{import}\PY{+w}{ }\PY{n+nn}{java.nio.file.Files}\PY{p}{;}
\PY{k+kn}{import}\PY{+w}{ }\PY{n+nn}{java.nio.file.Path}\PY{p}{;}
\PY{k+kn}{import}\PY{+w}{ }\PY{n+nn}{java.nio.file.attribute.FileAttribute}\PY{p}{;}
\PY{k+kn}{import}\PY{+w}{ }\PY{n+nn}{java.nio.file.attribute.PosixFilePermission}\PY{p}{;}
\PY{k+kn}{import}\PY{+w}{ }\PY{n+nn}{java.nio.file.attribute.PosixFilePermissions}\PY{p}{;}
\PY{k+kn}{import}\PY{+w}{ }\PY{n+nn}{java.util.Set}\PY{p}{;}

\PY{c+c1}{// Secure temp file: created with a random name, and on POSIX systems Files.createTempFile}
\PY{c+c1}{// sets permissions so only the owner can read/write it (0600) \PYZhy{}\PYZhy{} avoids TOCTOU race}
\PY{c+c1}{// conditions and world\PYZhy{}readable secrets that java.io.File.createTempFile historically allowed.}
\PY{n}{Path}\PY{+w}{ }\PY{n}{tempFile}\PY{+w}{ }\PY{o}{=}\PY{+w}{ }\PY{n}{Files}\PY{p}{.}\PY{n+na}{createTempFile}\PY{p}{(}\PY{l+s}{\PYZdq{}}\PY{l+s}{upload\PYZhy{}}\PY{l+s}{\PYZdq{}}\PY{p}{,}\PY{+w}{ }\PY{l+s}{\PYZdq{}}\PY{l+s}{.tmp}\PY{l+s}{\PYZdq{}}\PY{p}{)}\PY{p}{;}
\PY{n}{System}\PY{p}{.}\PY{n+na}{out}\PY{p}{.}\PY{n+na}{println}\PY{p}{(}\PY{n}{tempFile}\PY{p}{)}\PY{p}{;}\PY{+w}{ }\PY{c+c1}{// e.g. /tmp/upload\PYZhy{}1234567890.tmp}

\PY{c+c1}{// DirectoryStream: lazily iterate direct children of a directory (not recursive),}
\PY{c+c1}{// with an optional glob filter. Must be closed \PYZhy{}\PYZhy{} holds a directory handle.}
\PY{k}{try}\PY{+w}{ }\PY{p}{(}\PY{n}{DirectoryStream}\PY{o}{\PYZlt{}}\PY{n}{Path}\PY{o}{\PYZgt{}}\PY{+w}{ }\PY{n}{stream}\PY{+w}{ }\PY{o}{=}\PY{+w}{ }\PY{n}{Files}\PY{p}{.}\PY{n+na}{newDirectoryStream}\PY{p}{(}\PY{n}{Path}\PY{p}{.}\PY{n+na}{of}\PY{p}{(}\PY{l+s}{\PYZdq{}}\PY{l+s}{.}\PY{l+s}{\PYZdq{}}\PY{p}{)}\PY{p}{,}\PY{+w}{ }\PY{l+s}{\PYZdq{}}\PY{l+s}{*.java}\PY{l+s}{\PYZdq{}}\PY{p}{)}\PY{p}{)}\PY{+w}{ }\PY{p}{\PYZob{}}
\PY{+w}{    }\PY{k}{for}\PY{+w}{ }\PY{p}{(}\PY{n}{Path}\PY{+w}{ }\PY{n}{entry}\PY{+w}{ }\PY{p}{:}\PY{+w}{ }\PY{n}{stream}\PY{p}{)}\PY{+w}{ }\PY{p}{\PYZob{}}
\PY{+w}{        }\PY{n}{System}\PY{p}{.}\PY{n+na}{out}\PY{p}{.}\PY{n+na}{println}\PY{p}{(}\PY{n}{entry}\PY{p}{.}\PY{n+na}{getFileName}\PY{p}{(}\PY{p}{)}\PY{p}{)}\PY{p}{;}
\PY{+w}{    }\PY{p}{\PYZcb{}}
\PY{p}{\PYZcb{}}

\PY{c+c1}{// PosixFilePermissions: set exact rwx bits when creating a file, useful for}
\PY{c+c1}{// \PYZdq{}this file must only be readable by its owner\PYZdq{} security requirements}
\PY{n}{Set}\PY{o}{\PYZlt{}}\PY{n}{PosixFilePermission}\PY{o}{\PYZgt{}}\PY{+w}{ }\PY{n}{perms}\PY{+w}{ }\PY{o}{=}\PY{+w}{ }\PY{n}{PosixFilePermissions}\PY{p}{.}\PY{n+na}{fromString}\PY{p}{(}\PY{l+s}{\PYZdq{}}\PY{l+s}{rw\PYZhy{}\PYZhy{}\PYZhy{}\PYZhy{}\PYZhy{}\PYZhy{}\PYZhy{}}\PY{l+s}{\PYZdq{}}\PY{p}{)}\PY{p}{;}
\PY{n}{FileAttribute}\PY{o}{\PYZlt{}}\PY{n}{Set}\PY{o}{\PYZlt{}}\PY{n}{PosixFilePermission}\PY{o}{\PYZgt{}}\PY{o}{\PYZgt{}}\PY{+w}{ }\PY{n}{attr}\PY{+w}{ }\PY{o}{=}\PY{+w}{ }\PY{n}{PosixFilePermissions}\PY{p}{.}\PY{n+na}{asFileAttribute}\PY{p}{(}\PY{n}{perms}\PY{p}{)}\PY{p}{;}
\PY{n}{Path}\PY{+w}{ }\PY{n}{secretFile}\PY{+w}{ }\PY{o}{=}\PY{+w}{ }\PY{n}{Files}\PY{p}{.}\PY{n+na}{createFile}\PY{p}{(}\PY{n}{Path}\PY{p}{.}\PY{n+na}{of}\PY{p}{(}\PY{l+s}{\PYZdq{}}\PY{l+s}{secret.key}\PY{l+s}{\PYZdq{}}\PY{p}{)}\PY{p}{,}\PY{+w}{ }\PY{n}{attr}\PY{p}{)}\PY{p}{;}
\end{Verbatim}
\end{codebox}

\begin{codebox}
\begin{Verbatim}[commandchars=\\\{\}]
\PY{k+kn}{import}\PY{+w}{ }\PY{n+nn}{java.io.IOException}\PY{p}{;}
\PY{k+kn}{import}\PY{+w}{ }\PY{n+nn}{java.nio.file.*}\PY{p}{;}

\PY{c+c1}{// WatchService: get notified when files in a directory change, without polling.}
\PY{c+c1}{// Common use: reload config when a file on disk is edited.}
\PY{k}{try}\PY{+w}{ }\PY{p}{(}\PY{n}{WatchService}\PY{+w}{ }\PY{n}{watchService}\PY{+w}{ }\PY{o}{=}\PY{+w}{ }\PY{n}{FileSystems}\PY{p}{.}\PY{n+na}{getDefault}\PY{p}{(}\PY{p}{)}\PY{p}{.}\PY{n+na}{newWatchService}\PY{p}{(}\PY{p}{)}\PY{p}{)}\PY{+w}{ }\PY{p}{\PYZob{}}
\PY{+w}{    }\PY{n}{Path}\PY{+w}{ }\PY{n}{dir}\PY{+w}{ }\PY{o}{=}\PY{+w}{ }\PY{n}{Path}\PY{p}{.}\PY{n+na}{of}\PY{p}{(}\PY{l+s}{\PYZdq{}}\PY{l+s}{config}\PY{l+s}{\PYZdq{}}\PY{p}{)}\PY{p}{;}
\PY{+w}{    }\PY{n}{dir}\PY{p}{.}\PY{n+na}{register}\PY{p}{(}\PY{n}{watchService}\PY{p}{,}
\PY{+w}{            }\PY{n}{StandardWatchEventKinds}\PY{p}{.}\PY{n+na}{ENTRY\PYZus{}CREATE}\PY{p}{,}
\PY{+w}{            }\PY{n}{StandardWatchEventKinds}\PY{p}{.}\PY{n+na}{ENTRY\PYZus{}MODIFY}\PY{p}{,}
\PY{+w}{            }\PY{n}{StandardWatchEventKinds}\PY{p}{.}\PY{n+na}{ENTRY\PYZus{}DELETE}\PY{p}{)}\PY{p}{;}

\PY{+w}{    }\PY{n}{WatchKey}\PY{+w}{ }\PY{n}{key}\PY{+w}{ }\PY{o}{=}\PY{+w}{ }\PY{n}{watchService}\PY{p}{.}\PY{n+na}{take}\PY{p}{(}\PY{p}{)}\PY{p}{;}\PY{+w}{ }\PY{c+c1}{// blocks until an event occurs}
\PY{+w}{    }\PY{k}{for}\PY{+w}{ }\PY{p}{(}\PY{n}{WatchEvent}\PY{o}{\PYZlt{}}\PY{o}{?}\PY{o}{\PYZgt{}}\PY{+w}{ }\PY{n}{event}\PY{+w}{ }\PY{p}{:}\PY{+w}{ }\PY{n}{key}\PY{p}{.}\PY{n+na}{pollEvents}\PY{p}{(}\PY{p}{)}\PY{p}{)}\PY{+w}{ }\PY{p}{\PYZob{}}
\PY{+w}{        }\PY{n}{System}\PY{p}{.}\PY{n+na}{out}\PY{p}{.}\PY{n+na}{println}\PY{p}{(}\PY{n}{event}\PY{p}{.}\PY{n+na}{kind}\PY{p}{(}\PY{p}{)}\PY{+w}{ }\PY{o}{+}\PY{+w}{ }\PY{l+s}{\PYZdq{}}\PY{l+s}{: }\PY{l+s}{\PYZdq{}}\PY{+w}{ }\PY{o}{+}\PY{+w}{ }\PY{n}{event}\PY{p}{.}\PY{n+na}{context}\PY{p}{(}\PY{p}{)}\PY{p}{)}\PY{p}{;}
\PY{+w}{    }\PY{p}{\PYZcb{}}
\PY{+w}{    }\PY{n}{key}\PY{p}{.}\PY{n+na}{reset}\PY{p}{(}\PY{p}{)}\PY{p}{;}\PY{+w}{ }\PY{c+c1}{// must call this to keep receiving events on this key}
\PY{p}{\PYZcb{}}
\end{Verbatim}
\end{codebox}

\subsection[Atomic writes: temp file + ATOMIC\_MOVE]{Atomic writes: temp file + ATOMIC\_MOVE}
\label{h:10:atomic-writes-temp-file--atomic_move}
A common bug: writing directly to a target file. If the process crashes mid-write, readers can see a half-written, corrupted file. The safe pattern is to write to a temp file first, then atomically rename it into place — readers only ever see the old complete file or the new complete file, never a partial one.

\begin{codebox}
\begin{Verbatim}[commandchars=\\\{\}]
\PY{k+kn}{import}\PY{+w}{ }\PY{n+nn}{java.io.IOException}\PY{p}{;}
\PY{k+kn}{import}\PY{+w}{ }\PY{n+nn}{java.nio.charset.StandardCharsets}\PY{p}{;}
\PY{k+kn}{import}\PY{+w}{ }\PY{n+nn}{java.nio.file.Files}\PY{p}{;}
\PY{k+kn}{import}\PY{+w}{ }\PY{n+nn}{java.nio.file.Path}\PY{p}{;}
\PY{k+kn}{import}\PY{+w}{ }\PY{n+nn}{java.nio.file.StandardCopyOption}\PY{p}{;}

\PY{k+kd}{public}\PY{+w}{ }\PY{k+kd}{class} \PY{n+nc}{AtomicWrite}\PY{+w}{ }\PY{p}{\PYZob{}}
\PY{+w}{    }\PY{k+kd}{public}\PY{+w}{ }\PY{k+kd}{static}\PY{+w}{ }\PY{k+kt}{void}\PY{+w}{ }\PY{n+nf}{writeAtomically}\PY{p}{(}\PY{n}{Path}\PY{+w}{ }\PY{n}{target}\PY{p}{,}\PY{+w}{ }\PY{n}{String}\PY{+w}{ }\PY{n}{content}\PY{p}{)}\PY{+w}{ }\PY{k+kd}{throws}\PY{+w}{ }\PY{n}{IOException}\PY{+w}{ }\PY{p}{\PYZob{}}
\PY{+w}{        }\PY{c+c1}{// temp file in the SAME directory as target (atomic move requires same filesystem)}
\PY{+w}{        }\PY{n}{Path}\PY{+w}{ }\PY{n}{tempFile}\PY{+w}{ }\PY{o}{=}\PY{+w}{ }\PY{n}{Files}\PY{p}{.}\PY{n+na}{createTempFile}\PY{p}{(}\PY{n}{target}\PY{p}{.}\PY{n+na}{getParent}\PY{p}{(}\PY{p}{)}\PY{p}{,}\PY{+w}{ }\PY{l+s}{\PYZdq{}}\PY{l+s}{tmp\PYZhy{}}\PY{l+s}{\PYZdq{}}\PY{p}{,}\PY{+w}{ }\PY{l+s}{\PYZdq{}}\PY{l+s}{.swap}\PY{l+s}{\PYZdq{}}\PY{p}{)}\PY{p}{;}
\PY{+w}{        }\PY{k}{try}\PY{+w}{ }\PY{p}{\PYZob{}}
\PY{+w}{            }\PY{n}{Files}\PY{p}{.}\PY{n+na}{writeString}\PY{p}{(}\PY{n}{tempFile}\PY{p}{,}\PY{+w}{ }\PY{n}{content}\PY{p}{,}\PY{+w}{ }\PY{n}{StandardCharsets}\PY{p}{.}\PY{n+na}{UTF\PYZus{}8}\PY{p}{)}\PY{p}{;}
\PY{+w}{            }\PY{c+c1}{// ATOMIC\PYZus{}MOVE: the OS guarantees this rename is indivisible \PYZhy{}\PYZhy{}}
\PY{+w}{            }\PY{c+c1}{// any reader sees either the fully\PYZhy{}old or fully\PYZhy{}new file, never a mix}
\PY{+w}{            }\PY{n}{Files}\PY{p}{.}\PY{n+na}{move}\PY{p}{(}\PY{n}{tempFile}\PY{p}{,}\PY{+w}{ }\PY{n}{target}\PY{p}{,}
\PY{+w}{                    }\PY{n}{StandardCopyOption}\PY{p}{.}\PY{n+na}{ATOMIC\PYZus{}MOVE}\PY{p}{,}\PY{+w}{ }\PY{n}{StandardCopyOption}\PY{p}{.}\PY{n+na}{REPLACE\PYZus{}EXISTING}\PY{p}{)}\PY{p}{;}
\PY{+w}{        }\PY{p}{\PYZcb{}}\PY{+w}{ }\PY{k}{finally}\PY{+w}{ }\PY{p}{\PYZob{}}
\PY{+w}{            }\PY{n}{Files}\PY{p}{.}\PY{n+na}{deleteIfExists}\PY{p}{(}\PY{n}{tempFile}\PY{p}{)}\PY{p}{;}\PY{+w}{ }\PY{c+c1}{// cleanup if the move never happened}
\PY{+w}{        }\PY{p}{\PYZcb{}}
\PY{+w}{    }\PY{p}{\PYZcb{}}
\PY{p}{\PYZcb{}}
\end{Verbatim}
\end{codebox}

\subsection[Path traversal security (../)]{Path traversal security (\code{../})}
\label{h:10:path-traversal-security-}
If any part of a file path comes from user input, an attacker can inject \code{../} segments to escape the intended directory and read or write files elsewhere on the system — a classic \textbf{path traversal} vulnerability.

\begin{codebox}
\begin{Verbatim}[commandchars=\\\{\}]
\PY{k+kn}{import}\PY{+w}{ }\PY{n+nn}{java.io.IOException}\PY{p}{;}
\PY{k+kn}{import}\PY{+w}{ }\PY{n+nn}{java.nio.file.Files}\PY{p}{;}
\PY{k+kn}{import}\PY{+w}{ }\PY{n+nn}{java.nio.file.Path}\PY{p}{;}

\PY{c+c1}{// Bad: naively trusts user input, allows \PYZdq{}../../../etc/passwd\PYZdq{} style attacks}
\PY{k+kd}{public}\PY{+w}{ }\PY{k+kt}{byte}\PY{o}{[}\PY{o}{]}\PY{+w}{ }\PY{n+nf}{readUserFileUnsafe}\PY{p}{(}\PY{n}{String}\PY{+w}{ }\PY{n}{userSuppliedName}\PY{p}{)}\PY{+w}{ }\PY{k+kd}{throws}\PY{+w}{ }\PY{n}{IOException}\PY{+w}{ }\PY{p}{\PYZob{}}
\PY{+w}{    }\PY{n}{Path}\PY{+w}{ }\PY{n}{path}\PY{+w}{ }\PY{o}{=}\PY{+w}{ }\PY{n}{Path}\PY{p}{.}\PY{n+na}{of}\PY{p}{(}\PY{l+s}{\PYZdq{}}\PY{l+s}{/var/app/uploads}\PY{l+s}{\PYZdq{}}\PY{p}{)}\PY{p}{.}\PY{n+na}{resolve}\PY{p}{(}\PY{n}{userSuppliedName}\PY{p}{)}\PY{p}{;}
\PY{+w}{    }\PY{k}{return}\PY{+w}{ }\PY{n}{Files}\PY{p}{.}\PY{n+na}{readAllBytes}\PY{p}{(}\PY{n}{path}\PY{p}{)}\PY{p}{;}\PY{+w}{ }\PY{c+c1}{// could read ANY file the process can access}
\PY{p}{\PYZcb{}}

\PY{c+c1}{// Good: resolve, normalize, then verify the result is still inside the allowed root}
\PY{k+kd}{public}\PY{+w}{ }\PY{k+kt}{byte}\PY{o}{[}\PY{o}{]}\PY{+w}{ }\PY{n+nf}{readUserFileSafe}\PY{p}{(}\PY{n}{String}\PY{+w}{ }\PY{n}{userSuppliedName}\PY{p}{)}\PY{+w}{ }\PY{k+kd}{throws}\PY{+w}{ }\PY{n}{IOException}\PY{+w}{ }\PY{p}{\PYZob{}}
\PY{+w}{    }\PY{n}{Path}\PY{+w}{ }\PY{n}{root}\PY{+w}{ }\PY{o}{=}\PY{+w}{ }\PY{n}{Path}\PY{p}{.}\PY{n+na}{of}\PY{p}{(}\PY{l+s}{\PYZdq{}}\PY{l+s}{/var/app/uploads}\PY{l+s}{\PYZdq{}}\PY{p}{)}\PY{p}{.}\PY{n+na}{toRealPath}\PY{p}{(}\PY{p}{)}\PY{p}{;}\PY{+w}{ }\PY{c+c1}{// resolves symlinks too}
\PY{+w}{    }\PY{n}{Path}\PY{+w}{ }\PY{n}{requested}\PY{+w}{ }\PY{o}{=}\PY{+w}{ }\PY{n}{root}\PY{p}{.}\PY{n+na}{resolve}\PY{p}{(}\PY{n}{userSuppliedName}\PY{p}{)}\PY{p}{.}\PY{n+na}{normalize}\PY{p}{(}\PY{p}{)}\PY{p}{;}

\PY{+w}{    }\PY{k}{if}\PY{+w}{ }\PY{p}{(}\PY{o}{!}\PY{n}{requested}\PY{p}{.}\PY{n+na}{startsWith}\PY{p}{(}\PY{n}{root}\PY{p}{)}\PY{p}{)}\PY{+w}{ }\PY{p}{\PYZob{}}
\PY{+w}{        }\PY{k}{throw}\PY{+w}{ }\PY{k}{new}\PY{+w}{ }\PY{n}{SecurityException}\PY{p}{(}\PY{l+s}{\PYZdq{}}\PY{l+s}{Path traversal attempt: }\PY{l+s}{\PYZdq{}}\PY{+w}{ }\PY{o}{+}\PY{+w}{ }\PY{n}{userSuppliedName}\PY{p}{)}\PY{p}{;}
\PY{+w}{    }\PY{p}{\PYZcb{}}
\PY{+w}{    }\PY{k}{return}\PY{+w}{ }\PY{n}{Files}\PY{p}{.}\PY{n+na}{readAllBytes}\PY{p}{(}\PY{n}{requested}\PY{p}{)}\PY{p}{;}
\PY{p}{\PYZcb{}}
\end{Verbatim}
\end{codebox}

The check must happen \textbf{after} \code{resolve} + \code{normalize} (and ideally after resolving symlinks with \code{toRealPath()}), never before — normalizing first and checking after is the part reviewers most often see done wrong or skipped entirely.

\section[Serialization]{Serialization}
\label{h:10:serialization}
\textbf{Serialization} converts an object (and everything it references — its \textbf{object graph}) into a byte stream, so it can be saved to disk or sent over a network. \textbf{Deserialization} reverses the process, rebuilding the object graph from bytes.

\begin{codebox}
\begin{Verbatim}[commandchars=\\\{\}]
\PY{k+kn}{import}\PY{+w}{ }\PY{n+nn}{java.io.*}\PY{p}{;}

\PY{k+kd}{record} \PY{n+nc}{Point}\PY{p}{(}\PY{k+kt}{int}\PY{+w}{ }\PY{n}{x}\PY{p}{,}\PY{+w}{ }\PY{k+kt}{int}\PY{+w}{ }\PY{n}{y}\PY{p}{)}\PY{+w}{ }\PY{k+kd}{implements}\PY{+w}{ }\PY{n}{Serializable}\PY{+w}{ }\PY{p}{\PYZob{}}\PY{p}{\PYZcb{}}

\PY{k+kd}{public}\PY{+w}{ }\PY{k+kd}{class} \PY{n+nc}{SerializationBasics}\PY{+w}{ }\PY{p}{\PYZob{}}
\PY{+w}{    }\PY{k+kd}{public}\PY{+w}{ }\PY{k+kd}{static}\PY{+w}{ }\PY{k+kt}{void}\PY{+w}{ }\PY{n+nf}{main}\PY{p}{(}\PY{n}{String}\PY{o}{[}\PY{o}{]}\PY{+w}{ }\PY{n}{args}\PY{p}{)}\PY{+w}{ }\PY{k+kd}{throws}\PY{+w}{ }\PY{n}{IOException}\PY{p}{,}\PY{+w}{ }\PY{n}{ClassNotFoundException}\PY{+w}{ }\PY{p}{\PYZob{}}
\PY{+w}{        }\PY{n}{Point}\PY{+w}{ }\PY{n}{point}\PY{+w}{ }\PY{o}{=}\PY{+w}{ }\PY{k}{new}\PY{+w}{ }\PY{n}{Point}\PY{p}{(}\PY{l+m+mi}{3}\PY{p}{,}\PY{+w}{ }\PY{l+m+mi}{4}\PY{p}{)}\PY{p}{;}

\PY{+w}{        }\PY{c+c1}{// Serialize to a byte array (could just as easily be a FileOutputStream)}
\PY{+w}{        }\PY{k+kt}{byte}\PY{o}{[}\PY{o}{]}\PY{+w}{ }\PY{n}{bytes}\PY{p}{;}
\PY{+w}{        }\PY{k}{try}\PY{+w}{ }\PY{p}{(}\PY{n}{ByteArrayOutputStream}\PY{+w}{ }\PY{n}{baos}\PY{+w}{ }\PY{o}{=}\PY{+w}{ }\PY{k}{new}\PY{+w}{ }\PY{n}{ByteArrayOutputStream}\PY{p}{(}\PY{p}{)}\PY{p}{;}
\PY{+w}{             }\PY{n}{ObjectOutputStream}\PY{+w}{ }\PY{n}{oos}\PY{+w}{ }\PY{o}{=}\PY{+w}{ }\PY{k}{new}\PY{+w}{ }\PY{n}{ObjectOutputStream}\PY{p}{(}\PY{n}{baos}\PY{p}{)}\PY{p}{)}\PY{+w}{ }\PY{p}{\PYZob{}}
\PY{+w}{            }\PY{n}{oos}\PY{p}{.}\PY{n+na}{writeObject}\PY{p}{(}\PY{n}{point}\PY{p}{)}\PY{p}{;}
\PY{+w}{            }\PY{n}{bytes}\PY{+w}{ }\PY{o}{=}\PY{+w}{ }\PY{n}{baos}\PY{p}{.}\PY{n+na}{toByteArray}\PY{p}{(}\PY{p}{)}\PY{p}{;}
\PY{+w}{        }\PY{p}{\PYZcb{}}

\PY{+w}{        }\PY{c+c1}{// Deserialize back into an object}
\PY{+w}{        }\PY{k}{try}\PY{+w}{ }\PY{p}{(}\PY{n}{ObjectInputStream}\PY{+w}{ }\PY{n}{ois}\PY{+w}{ }\PY{o}{=}\PY{+w}{ }\PY{k}{new}\PY{+w}{ }\PY{n}{ObjectInputStream}\PY{p}{(}\PY{k}{new}\PY{+w}{ }\PY{n}{ByteArrayInputStream}\PY{p}{(}\PY{n}{bytes}\PY{p}{)}\PY{p}{)}\PY{p}{)}\PY{+w}{ }\PY{p}{\PYZob{}}
\PY{+w}{            }\PY{n}{Point}\PY{+w}{ }\PY{n}{restored}\PY{+w}{ }\PY{o}{=}\PY{+w}{ }\PY{p}{(}\PY{n}{Point}\PY{p}{)}\PY{+w}{ }\PY{n}{ois}\PY{p}{.}\PY{n+na}{readObject}\PY{p}{(}\PY{p}{)}\PY{p}{;}
\PY{+w}{            }\PY{n}{System}\PY{p}{.}\PY{n+na}{out}\PY{p}{.}\PY{n+na}{println}\PY{p}{(}\PY{n}{restored}\PY{p}{)}\PY{p}{;}\PY{+w}{ }\PY{c+c1}{// Point[x=3, y=4]}
\PY{+w}{        }\PY{p}{\PYZcb{}}
\PY{+w}{    }\PY{p}{\PYZcb{}}
\PY{p}{\PYZcb{}}
\end{Verbatim}
\end{codebox}

To be serializable, a class just implements the empty marker interface \code{Serializable} — there are no methods to implement. The JVM uses reflection to walk the object’s fields and write them out.

\subsection[serialVersionUID]{serialVersionUID}
\label{h:10:serialversionuid}
Every \code{Serializable} class has an implicit \textbf{version ID} used to check that a serialized object matches the class trying to deserialize it. If you don’t declare it explicitly, the JVM computes one automatically from the class’s structure (fields, methods, etc.) — which means \textbf{any change to the class}, even something as harmless as adding a method, can change the computed ID and break deserialization of old data with an \code{InvalidClassException}.

\begin{codebox}
\begin{Verbatim}[commandchars=\\\{\}]
\PY{k+kn}{import}\PY{+w}{ }\PY{n+nn}{java.io.Serializable}\PY{p}{;}

\PY{k+kd}{public}\PY{+w}{ }\PY{k+kd}{class} \PY{n+nc}{User}\PY{+w}{ }\PY{k+kd}{implements}\PY{+w}{ }\PY{n}{Serializable}\PY{+w}{ }\PY{p}{\PYZob{}}
\PY{+w}{    }\PY{c+c1}{// Always declare this explicitly. It lets you evolve the class}
\PY{+w}{    }\PY{c+c1}{// (add fields, add methods) without silently breaking old serialized data.}
\PY{+w}{    }\PY{k+kd}{private}\PY{+w}{ }\PY{k+kd}{static}\PY{+w}{ }\PY{k+kd}{final}\PY{+w}{ }\PY{k+kt}{long}\PY{+w}{ }\PY{n}{serialVersionUID}\PY{+w}{ }\PY{o}{=}\PY{+w}{ }\PY{l+m+mi}{1L}\PY{p}{;}

\PY{+w}{    }\PY{k+kd}{private}\PY{+w}{ }\PY{n}{String}\PY{+w}{ }\PY{n}{username}\PY{p}{;}
\PY{+w}{    }\PY{k+kd}{private}\PY{+w}{ }\PY{k+kt}{int}\PY{+w}{ }\PY{n}{age}\PY{p}{;}

\PY{+w}{    }\PY{c+c1}{// getters/setters omitted}
\PY{p}{\PYZcb{}}
\end{Verbatim}
\end{codebox}

\textbf{Rule for review}: any \code{Serializable} class without an explicit \code{serialVersionUID} is a red flag — it works today but is a landmine for future changes.

\subsection[transient]{transient}
\label{h:10:transient}
Marking a field \code{transient} excludes it from serialization. Use it for fields that shouldn’t (or can’t) be persisted: caches, derived/computed values, \code{Thread}/\code{Socket}/\code{Connection} handles, secrets that shouldn’t be written to disk.

\begin{codebox}
\begin{Verbatim}[commandchars=\\\{\}]
\PY{k+kn}{import}\PY{+w}{ }\PY{n+nn}{java.io.Serializable}\PY{p}{;}

\PY{k+kd}{public}\PY{+w}{ }\PY{k+kd}{class} \PY{n+nc}{Session}\PY{+w}{ }\PY{k+kd}{implements}\PY{+w}{ }\PY{n}{Serializable}\PY{+w}{ }\PY{p}{\PYZob{}}
\PY{+w}{    }\PY{k+kd}{private}\PY{+w}{ }\PY{k+kd}{static}\PY{+w}{ }\PY{k+kd}{final}\PY{+w}{ }\PY{k+kt}{long}\PY{+w}{ }\PY{n}{serialVersionUID}\PY{+w}{ }\PY{o}{=}\PY{+w}{ }\PY{l+m+mi}{1L}\PY{p}{;}

\PY{+w}{    }\PY{k+kd}{private}\PY{+w}{ }\PY{n}{String}\PY{+w}{ }\PY{n}{sessionId}\PY{p}{;}
\PY{+w}{    }\PY{k+kd}{private}\PY{+w}{ }\PY{k+kd}{transient}\PY{+w}{ }\PY{n}{String}\PY{+w}{ }\PY{n}{cachedDisplayName}\PY{p}{;}\PY{+w}{ }\PY{c+c1}{// recomputed on demand, not persisted}
\PY{+w}{    }\PY{k+kd}{private}\PY{+w}{ }\PY{k+kd}{transient}\PY{+w}{ }\PY{k+kt}{char}\PY{o}{[}\PY{o}{]}\PY{+w}{ }\PY{n}{temporaryToken}\PY{p}{;}\PY{+w}{    }\PY{c+c1}{// sensitive; must never hit disk}

\PY{+w}{    }\PY{c+c1}{// after deserialization, cachedDisplayName and temporaryToken are null/zero\PYZhy{}valued,}
\PY{+w}{    }\PY{c+c1}{// NOT whatever they held before serialization}
\PY{p}{\PYZcb{}}
\end{Verbatim}
\end{codebox}

After deserialization, \code{transient} fields are reset to their default value (\code{null}, \code{0}, \code{false}, etc.) — the constructor is \textbf{not} called during deserialization, so you can’t rely on constructor logic to re-populate them; use \code{readObject} (below) if you need to recompute them.

\subsection[Custom writeObject / readObject]{Custom writeObject / readObject}
\label{h:10:custom-writeobject--readobject}
You can hook into the serialization process by declaring \code{private} methods with these exact signatures. The JVM calls them via reflection if present.

\begin{codebox}
\begin{Verbatim}[commandchars=\\\{\}]
\PY{k+kn}{import}\PY{+w}{ }\PY{n+nn}{java.io.*}\PY{p}{;}

\PY{k+kd}{public}\PY{+w}{ }\PY{k+kd}{class} \PY{n+nc}{Account}\PY{+w}{ }\PY{k+kd}{implements}\PY{+w}{ }\PY{n}{Serializable}\PY{+w}{ }\PY{p}{\PYZob{}}
\PY{+w}{    }\PY{k+kd}{private}\PY{+w}{ }\PY{k+kd}{static}\PY{+w}{ }\PY{k+kd}{final}\PY{+w}{ }\PY{k+kt}{long}\PY{+w}{ }\PY{n}{serialVersionUID}\PY{+w}{ }\PY{o}{=}\PY{+w}{ }\PY{l+m+mi}{1L}\PY{p}{;}

\PY{+w}{    }\PY{k+kd}{private}\PY{+w}{ }\PY{n}{String}\PY{+w}{ }\PY{n}{owner}\PY{p}{;}
\PY{+w}{    }\PY{k+kd}{private}\PY{+w}{ }\PY{k+kd}{transient}\PY{+w}{ }\PY{k+kt}{double}\PY{+w}{ }\PY{n}{cachedBalance}\PY{p}{;}\PY{+w}{ }\PY{c+c1}{// recomputed after load, not stored directly}

\PY{+w}{    }\PY{k+kd}{private}\PY{+w}{ }\PY{k+kt}{void}\PY{+w}{ }\PY{n+nf}{writeObject}\PY{p}{(}\PY{n}{ObjectOutputStream}\PY{+w}{ }\PY{n}{out}\PY{p}{)}\PY{+w}{ }\PY{k+kd}{throws}\PY{+w}{ }\PY{n}{IOException}\PY{+w}{ }\PY{p}{\PYZob{}}
\PY{+w}{        }\PY{n}{out}\PY{p}{.}\PY{n+na}{defaultWriteObject}\PY{p}{(}\PY{p}{)}\PY{p}{;}\PY{+w}{ }\PY{c+c1}{// writes all non\PYZhy{}transient fields normally}
\PY{+w}{        }\PY{n}{out}\PY{p}{.}\PY{n+na}{writeDouble}\PY{p}{(}\PY{n}{computeAuditedBalance}\PY{p}{(}\PY{p}{)}\PY{p}{)}\PY{p}{;}\PY{+w}{ }\PY{c+c1}{// then write something extra}
\PY{+w}{    }\PY{p}{\PYZcb{}}

\PY{+w}{    }\PY{k+kd}{private}\PY{+w}{ }\PY{k+kt}{void}\PY{+w}{ }\PY{n+nf}{readObject}\PY{p}{(}\PY{n}{ObjectInputStream}\PY{+w}{ }\PY{n}{in}\PY{p}{)}\PY{+w}{ }\PY{k+kd}{throws}\PY{+w}{ }\PY{n}{IOException}\PY{p}{,}\PY{+w}{ }\PY{n}{ClassNotFoundException}\PY{+w}{ }\PY{p}{\PYZob{}}
\PY{+w}{        }\PY{n}{in}\PY{p}{.}\PY{n+na}{defaultReadObject}\PY{p}{(}\PY{p}{)}\PY{p}{;}\PY{+w}{ }\PY{c+c1}{// reads all non\PYZhy{}transient fields normally}
\PY{+w}{        }\PY{k}{this}\PY{p}{.}\PY{n+na}{cachedBalance}\PY{+w}{ }\PY{o}{=}\PY{+w}{ }\PY{n}{in}\PY{p}{.}\PY{n+na}{readDouble}\PY{p}{(}\PY{p}{)}\PY{p}{;}\PY{+w}{ }\PY{c+c1}{// then read the extra value back}
\PY{+w}{    }\PY{p}{\PYZcb{}}

\PY{+w}{    }\PY{k+kd}{private}\PY{+w}{ }\PY{k+kt}{double}\PY{+w}{ }\PY{n+nf}{computeAuditedBalance}\PY{p}{(}\PY{p}{)}\PY{+w}{ }\PY{p}{\PYZob{}}
\PY{+w}{        }\PY{k}{return}\PY{+w}{ }\PY{l+m+mf}{0.0}\PY{p}{;}\PY{+w}{ }\PY{c+c1}{// placeholder for real logic}
\PY{+w}{    }\PY{p}{\PYZcb{}}
\PY{p}{\PYZcb{}}
\end{Verbatim}
\end{codebox}

Typical uses: encrypting a sensitive field before writing it out (and decrypting on read), validating invariants after loading untrusted data, or handling fields that aren’t naturally serializable (e.g., writing a \code{Map}’s contents manually instead of the map object itself).

\subsection[writeReplace and readResolve]{writeReplace and readResolve}
\label{h:10:writereplace-and-readresolve}
These let a class substitute a \emph{different} object during serialization (\code{writeReplace}) or deserialization (\code{readResolve}) — commonly used to enforce singletons or to serialize a stable “proxy” representation instead of the real object.

\begin{codebox}
\begin{Verbatim}[commandchars=\\\{\}]
\PY{k+kn}{import}\PY{+w}{ }\PY{n+nn}{java.io.*}\PY{p}{;}

\PY{k+kd}{public}\PY{+w}{ }\PY{k+kd}{class} \PY{n+nc}{Singleton}\PY{+w}{ }\PY{k+kd}{implements}\PY{+w}{ }\PY{n}{Serializable}\PY{+w}{ }\PY{p}{\PYZob{}}
\PY{+w}{    }\PY{k+kd}{private}\PY{+w}{ }\PY{k+kd}{static}\PY{+w}{ }\PY{k+kd}{final}\PY{+w}{ }\PY{k+kt}{long}\PY{+w}{ }\PY{n}{serialVersionUID}\PY{+w}{ }\PY{o}{=}\PY{+w}{ }\PY{l+m+mi}{1L}\PY{p}{;}

\PY{+w}{    }\PY{k+kd}{public}\PY{+w}{ }\PY{k+kd}{static}\PY{+w}{ }\PY{k+kd}{final}\PY{+w}{ }\PY{n}{Singleton}\PY{+w}{ }\PY{n}{INSTANCE}\PY{+w}{ }\PY{o}{=}\PY{+w}{ }\PY{k}{new}\PY{+w}{ }\PY{n}{Singleton}\PY{p}{(}\PY{p}{)}\PY{p}{;}

\PY{+w}{    }\PY{k+kd}{private}\PY{+w}{ }\PY{n+nf}{Singleton}\PY{p}{(}\PY{p}{)}\PY{+w}{ }\PY{p}{\PYZob{}}\PY{p}{\PYZcb{}}

\PY{+w}{    }\PY{c+c1}{// Without readResolve, deserialization would create a SECOND, distinct}
\PY{+w}{    }\PY{c+c1}{// instance via reflection, breaking the singleton guarantee.}
\PY{+w}{    }\PY{k+kd}{protected}\PY{+w}{ }\PY{n}{Object}\PY{+w}{ }\PY{n+nf}{readResolve}\PY{p}{(}\PY{p}{)}\PY{+w}{ }\PY{p}{\PYZob{}}
\PY{+w}{        }\PY{k}{return}\PY{+w}{ }\PY{n}{INSTANCE}\PY{p}{;}\PY{+w}{ }\PY{c+c1}{// replace the deserialized instance with the canonical one}
\PY{+w}{    }\PY{p}{\PYZcb{}}
\PY{p}{\PYZcb{}}
\end{Verbatim}
\end{codebox}

\code{enum} types are serialization-safe singletons for free (the JVM guarantees enum constants deserialize to the same instance) — that’s one more reason the enum singleton pattern is generally preferred over a hand-rolled one.

\subsection[Why Java serialization is a security disaster]{Why Java serialization is a security disaster}
\label{h:10:why-java-serialization-is-a-security-disaster}
\code{Object\allowbreak{}Input\allowbreak{}Stream.\allowbreak{}read\allowbreak{}Object()} does something dangerous by design: it instantiates classes and calls code (\code{readObject}, constructors of superclasses in some cases, \code{readResolve}, etc.) \textbf{based purely on bytes that came from the input} — before your application logic gets any say in whether that data should be trusted.

If an attacker controls the byte stream, they can craft input referencing classes already on your classpath whose \code{readObject}/\code{finalize}/other methods have \emph{side effects} — chaining several such classes together forms a \textbf{gadget chain} that can lead to arbitrary code execution, denial of service, or data exfiltration, all before your code ever sees the resulting object. This class of vulnerability caused real, high-severity CVEs in major frameworks throughout the 2010s (e.g., in Apache Commons Collections, WebLogic, Jenkins) and is why security guidance now generally treats \textbf{deserializing untrustworthy data with native Java serialization as unsafe by default.}

\begin{codebox}
\begin{Verbatim}[commandchars=\\\{\}]
\PY{k+kn}{import}\PY{+w}{ }\PY{n+nn}{java.io.*}\PY{p}{;}

\PY{c+c1}{// DANGEROUS: never deserialize bytes from an untrusted source (network request,}
\PY{c+c1}{// user upload, unauthenticated input) with plain ObjectInputStream.}
\PY{k+kd}{public}\PY{+w}{ }\PY{n}{Object}\PY{+w}{ }\PY{n+nf}{deserializeUntrusted}\PY{p}{(}\PY{k+kt}{byte}\PY{o}{[}\PY{o}{]}\PY{+w}{ }\PY{n}{attackerControlledBytes}\PY{p}{)}\PY{+w}{ }\PY{k+kd}{throws}\PY{+w}{ }\PY{n}{Exception}\PY{+w}{ }\PY{p}{\PYZob{}}
\PY{+w}{    }\PY{k}{try}\PY{+w}{ }\PY{p}{(}\PY{n}{ObjectInputStream}\PY{+w}{ }\PY{n}{ois}\PY{+w}{ }\PY{o}{=}\PY{+w}{ }\PY{k}{new}\PY{+w}{ }\PY{n}{ObjectInputStream}\PY{p}{(}
\PY{+w}{            }\PY{k}{new}\PY{+w}{ }\PY{n}{ByteArrayInputStream}\PY{p}{(}\PY{n}{attackerControlledBytes}\PY{p}{)}\PY{p}{)}\PY{p}{)}\PY{+w}{ }\PY{p}{\PYZob{}}
\PY{+w}{        }\PY{k}{return}\PY{+w}{ }\PY{n}{ois}\PY{p}{.}\PY{n+na}{readObject}\PY{p}{(}\PY{p}{)}\PY{p}{;}\PY{+w}{ }\PY{c+c1}{// could trigger a gadget chain before you ever inspect it}
\PY{+w}{    }\PY{p}{\PYZcb{}}
\PY{p}{\PYZcb{}}
\end{Verbatim}
\end{codebox}

\subsection[Serialization filters (ObjectInputFilter, JDK 9+)]{Serialization filters (ObjectInputFilter, JDK 9+)}
\label{h:10:serialization-filters-objectinputfilter-jdk-9}
If you must use Java serialization (e.g., legacy systems), \code{ObjectInputFilter} lets you restrict \emph{which classes} are allowed to be deserialized, rejecting everything else before it’s fully constructed.

\begin{codebox}
\begin{Verbatim}[commandchars=\\\{\}]
\PY{k+kn}{import}\PY{+w}{ }\PY{n+nn}{java.io.*}\PY{p}{;}

\PY{n}{ObjectInputFilter}\PY{+w}{ }\PY{n}{filter}\PY{+w}{ }\PY{o}{=}\PY{+w}{ }\PY{n}{ObjectInputFilter}\PY{p}{.}\PY{n+na}{Config}\PY{p}{.}\PY{n+na}{createFilter}\PY{p}{(}
\PY{+w}{        }\PY{l+s}{\PYZdq{}}\PY{l+s}{com.example.model.*;java.base/*;!*}\PY{l+s}{\PYZdq{}}\PY{p}{)}\PY{p}{;}\PY{+w}{ }\PY{c+c1}{// allow our package + JDK basics, reject all else}

\PY{k}{try}\PY{+w}{ }\PY{p}{(}\PY{n}{ObjectInputStream}\PY{+w}{ }\PY{n}{ois}\PY{+w}{ }\PY{o}{=}\PY{+w}{ }\PY{k}{new}\PY{+w}{ }\PY{n}{ObjectInputStream}\PY{p}{(}\PY{k}{new}\PY{+w}{ }\PY{n}{FileInputStream}\PY{p}{(}\PY{l+s}{\PYZdq{}}\PY{l+s}{data.ser}\PY{l+s}{\PYZdq{}}\PY{p}{)}\PY{p}{)}\PY{p}{)}\PY{+w}{ }\PY{p}{\PYZob{}}
\PY{+w}{    }\PY{n}{ois}\PY{p}{.}\PY{n+na}{setObjectInputFilter}\PY{p}{(}\PY{n}{filter}\PY{p}{)}\PY{p}{;}
\PY{+w}{    }\PY{n}{Object}\PY{+w}{ }\PY{n}{result}\PY{+w}{ }\PY{o}{=}\PY{+w}{ }\PY{n}{ois}\PY{p}{.}\PY{n+na}{readObject}\PY{p}{(}\PY{p}{)}\PY{p}{;}\PY{+w}{ }\PY{c+c1}{// throws InvalidClassException if a disallowed class appears}
\PY{p}{\PYZcb{}}
\end{Verbatim}
\end{codebox}

You can also set a process-wide default filter via the \code{jdk.\allowbreak{}serialFilter} system property or security property, so every \code{ObjectInputStream} in the JVM is protected even in code you don’t control.

\subsection[The modern recommendation]{The modern recommendation}
\label{h:10:the-modern-recommendation}
Given the risks and the format’s other downsides (JVM-only, brittle across class changes, opaque binary format that’s hard to debug or audit), the strong modern recommendation is:

\textbf{Avoid Java’s built-in \code{Serializable} for anything crossing a trust boundary (network, files from outside your process, another service). Prefer JSON (Jackson, Gson) or a schema-based binary format (Protocol Buffers, Avro) instead.}

\begin{codebox}
\begin{Verbatim}[commandchars=\\\{\}]
\PY{c+c1}{// Instead of: implements Serializable + ObjectOutputStream/ObjectInputStream}
\PY{c+c1}{// Prefer something like (using a JSON library):}
\PY{c+c1}{//}
\PY{c+c1}{//   String json = objectMapper.writeValueAsString(user);}
\PY{c+c1}{//   User user = objectMapper.readValue(json, User.class);}
\PY{c+c1}{//}
\PY{c+c1}{// JSON/protobuf are language\PYZhy{}agnostic, human\PYZhy{}inspectable (JSON) or}
\PY{c+c1}{// schema\PYZhy{}validated (protobuf), and don\PYZsq{}t execute arbitrary code on parse.}
\end{Verbatim}
\end{codebox}

\code{Serializable} still shows up legitimately for things that never leave a single trusted JVM process — e.g., objects passed between nodes in a cluster you fully control, or short-lived in-memory caching frameworks — but even there, many teams now default to a safer format out of caution and long-term maintainability.

\section[Externalizable]{Externalizable}
\label{h:10:externalizable}
\code{Externalizable} is an alternative to \code{Serializable} that extends it and hands you \textbf{full manual control} over the byte format — you write every field yourself, in \code{writeExternal}, and read every field yourself, in \code{readExternal}.

\begin{codebox}
\begin{Verbatim}[commandchars=\\\{\}]
\PY{k+kn}{import}\PY{+w}{ }\PY{n+nn}{java.io.*}\PY{p}{;}

\PY{k+kd}{public}\PY{+w}{ }\PY{k+kd}{class} \PY{n+nc}{Product}\PY{+w}{ }\PY{k+kd}{implements}\PY{+w}{ }\PY{n}{Externalizable}\PY{+w}{ }\PY{p}{\PYZob{}}
\PY{+w}{    }\PY{k+kd}{private}\PY{+w}{ }\PY{n}{String}\PY{+w}{ }\PY{n}{name}\PY{p}{;}
\PY{+w}{    }\PY{k+kd}{private}\PY{+w}{ }\PY{k+kt}{double}\PY{+w}{ }\PY{n}{price}\PY{p}{;}

\PY{+w}{    }\PY{c+c1}{// Externalizable REQUIRES a public no\PYZhy{}arg constructor.}
\PY{+w}{    }\PY{c+c1}{// The JVM calls this constructor first, THEN calls readExternal()}
\PY{+w}{    }\PY{c+c1}{// to populate the fields \PYZhy{}\PYZhy{} unlike Serializable, which uses reflection}
\PY{+w}{    }\PY{c+c1}{// to bypass the constructor entirely.}
\PY{+w}{    }\PY{k+kd}{public}\PY{+w}{ }\PY{n+nf}{Product}\PY{p}{(}\PY{p}{)}\PY{+w}{ }\PY{p}{\PYZob{}}\PY{p}{\PYZcb{}}

\PY{+w}{    }\PY{k+kd}{public}\PY{+w}{ }\PY{n+nf}{Product}\PY{p}{(}\PY{n}{String}\PY{+w}{ }\PY{n}{name}\PY{p}{,}\PY{+w}{ }\PY{k+kt}{double}\PY{+w}{ }\PY{n}{price}\PY{p}{)}\PY{+w}{ }\PY{p}{\PYZob{}}
\PY{+w}{        }\PY{k}{this}\PY{p}{.}\PY{n+na}{name}\PY{+w}{ }\PY{o}{=}\PY{+w}{ }\PY{n}{name}\PY{p}{;}
\PY{+w}{        }\PY{k}{this}\PY{p}{.}\PY{n+na}{price}\PY{+w}{ }\PY{o}{=}\PY{+w}{ }\PY{n}{price}\PY{p}{;}
\PY{+w}{    }\PY{p}{\PYZcb{}}

\PY{+w}{    }\PY{n+nd}{@Override}
\PY{+w}{    }\PY{k+kd}{public}\PY{+w}{ }\PY{k+kt}{void}\PY{+w}{ }\PY{n+nf}{writeExternal}\PY{p}{(}\PY{n}{ObjectOutput}\PY{+w}{ }\PY{n}{out}\PY{p}{)}\PY{+w}{ }\PY{k+kd}{throws}\PY{+w}{ }\PY{n}{IOException}\PY{+w}{ }\PY{p}{\PYZob{}}
\PY{+w}{        }\PY{n}{out}\PY{p}{.}\PY{n+na}{writeUTF}\PY{p}{(}\PY{n}{name}\PY{p}{)}\PY{p}{;}
\PY{+w}{        }\PY{n}{out}\PY{p}{.}\PY{n+na}{writeDouble}\PY{p}{(}\PY{n}{price}\PY{p}{)}\PY{p}{;}
\PY{+w}{    }\PY{p}{\PYZcb{}}

\PY{+w}{    }\PY{n+nd}{@Override}
\PY{+w}{    }\PY{k+kd}{public}\PY{+w}{ }\PY{k+kt}{void}\PY{+w}{ }\PY{n+nf}{readExternal}\PY{p}{(}\PY{n}{ObjectInput}\PY{+w}{ }\PY{n}{in}\PY{p}{)}\PY{+w}{ }\PY{k+kd}{throws}\PY{+w}{ }\PY{n}{IOException}\PY{+w}{ }\PY{p}{\PYZob{}}
\PY{+w}{        }\PY{n}{name}\PY{+w}{ }\PY{o}{=}\PY{+w}{ }\PY{n}{in}\PY{p}{.}\PY{n+na}{readUTF}\PY{p}{(}\PY{p}{)}\PY{p}{;}
\PY{+w}{        }\PY{n}{price}\PY{+w}{ }\PY{o}{=}\PY{+w}{ }\PY{n}{in}\PY{p}{.}\PY{n+na}{readDouble}\PY{p}{(}\PY{p}{)}\PY{p}{;}
\PY{+w}{    }\PY{p}{\PYZcb{}}

\PY{+w}{    }\PY{n+nd}{@Override}
\PY{+w}{    }\PY{k+kd}{public}\PY{+w}{ }\PY{n}{String}\PY{+w}{ }\PY{n+nf}{toString}\PY{p}{(}\PY{p}{)}\PY{+w}{ }\PY{p}{\PYZob{}}
\PY{+w}{        }\PY{k}{return}\PY{+w}{ }\PY{l+s}{\PYZdq{}}\PY{l+s}{Product\PYZob{}name=\PYZsq{}}\PY{l+s}{\PYZdq{}}\PY{+w}{ }\PY{o}{+}\PY{+w}{ }\PY{n}{name}\PY{+w}{ }\PY{o}{+}\PY{+w}{ }\PY{l+s}{\PYZdq{}}\PY{l+s}{\PYZsq{}, price=}\PY{l+s}{\PYZdq{}}\PY{+w}{ }\PY{o}{+}\PY{+w}{ }\PY{n}{price}\PY{+w}{ }\PY{o}{+}\PY{+w}{ }\PY{l+s+sc}{\PYZsq{}\PYZcb{}\PYZsq{}}\PY{p}{;}
\PY{+w}{    }\PY{p}{\PYZcb{}}
\PY{p}{\PYZcb{}}
\end{Verbatim}
\end{codebox}

\begin{codebox}
\begin{Verbatim}[commandchars=\\\{\}]
\PY{k+kn}{import}\PY{+w}{ }\PY{n+nn}{java.io.*}\PY{p}{;}

\PY{k+kd}{public}\PY{+w}{ }\PY{k+kd}{class} \PY{n+nc}{ExternalizableDemo}\PY{+w}{ }\PY{p}{\PYZob{}}
\PY{+w}{    }\PY{k+kd}{public}\PY{+w}{ }\PY{k+kd}{static}\PY{+w}{ }\PY{k+kt}{void}\PY{+w}{ }\PY{n+nf}{main}\PY{p}{(}\PY{n}{String}\PY{o}{[}\PY{o}{]}\PY{+w}{ }\PY{n}{args}\PY{p}{)}\PY{+w}{ }\PY{k+kd}{throws}\PY{+w}{ }\PY{n}{IOException}\PY{p}{,}\PY{+w}{ }\PY{n}{ClassNotFoundException}\PY{+w}{ }\PY{p}{\PYZob{}}
\PY{+w}{        }\PY{n}{Product}\PY{+w}{ }\PY{n}{product}\PY{+w}{ }\PY{o}{=}\PY{+w}{ }\PY{k}{new}\PY{+w}{ }\PY{n}{Product}\PY{p}{(}\PY{l+s}{\PYZdq{}}\PY{l+s}{Widget}\PY{l+s}{\PYZdq{}}\PY{p}{,}\PY{+w}{ }\PY{l+m+mf}{9.99}\PY{p}{)}\PY{p}{;}

\PY{+w}{        }\PY{k+kt}{byte}\PY{o}{[}\PY{o}{]}\PY{+w}{ }\PY{n}{bytes}\PY{p}{;}
\PY{+w}{        }\PY{k}{try}\PY{+w}{ }\PY{p}{(}\PY{n}{ByteArrayOutputStream}\PY{+w}{ }\PY{n}{baos}\PY{+w}{ }\PY{o}{=}\PY{+w}{ }\PY{k}{new}\PY{+w}{ }\PY{n}{ByteArrayOutputStream}\PY{p}{(}\PY{p}{)}\PY{p}{;}
\PY{+w}{             }\PY{n}{ObjectOutputStream}\PY{+w}{ }\PY{n}{oos}\PY{+w}{ }\PY{o}{=}\PY{+w}{ }\PY{k}{new}\PY{+w}{ }\PY{n}{ObjectOutputStream}\PY{p}{(}\PY{n}{baos}\PY{p}{)}\PY{p}{)}\PY{+w}{ }\PY{p}{\PYZob{}}
\PY{+w}{            }\PY{n}{oos}\PY{p}{.}\PY{n+na}{writeObject}\PY{p}{(}\PY{n}{product}\PY{p}{)}\PY{p}{;}
\PY{+w}{            }\PY{n}{bytes}\PY{+w}{ }\PY{o}{=}\PY{+w}{ }\PY{n}{baos}\PY{p}{.}\PY{n+na}{toByteArray}\PY{p}{(}\PY{p}{)}\PY{p}{;}
\PY{+w}{        }\PY{p}{\PYZcb{}}

\PY{+w}{        }\PY{k}{try}\PY{+w}{ }\PY{p}{(}\PY{n}{ObjectInputStream}\PY{+w}{ }\PY{n}{ois}\PY{+w}{ }\PY{o}{=}\PY{+w}{ }\PY{k}{new}\PY{+w}{ }\PY{n}{ObjectInputStream}\PY{p}{(}\PY{k}{new}\PY{+w}{ }\PY{n}{ByteArrayInputStream}\PY{p}{(}\PY{n}{bytes}\PY{p}{)}\PY{p}{)}\PY{p}{)}\PY{+w}{ }\PY{p}{\PYZob{}}
\PY{+w}{            }\PY{n}{Product}\PY{+w}{ }\PY{n}{restored}\PY{+w}{ }\PY{o}{=}\PY{+w}{ }\PY{p}{(}\PY{n}{Product}\PY{p}{)}\PY{+w}{ }\PY{n}{ois}\PY{p}{.}\PY{n+na}{readObject}\PY{p}{(}\PY{p}{)}\PY{p}{;}\PY{+w}{ }\PY{c+c1}{// calls Product() then readExternal()}
\PY{+w}{            }\PY{n}{System}\PY{p}{.}\PY{n+na}{out}\PY{p}{.}\PY{n+na}{println}\PY{p}{(}\PY{n}{restored}\PY{p}{)}\PY{p}{;}\PY{+w}{ }\PY{c+c1}{// Product\PYZob{}name=\PYZsq{}Widget\PYZsq{}, price=9.99\PYZcb{}}
\PY{+w}{        }\PY{p}{\PYZcb{}}
\PY{+w}{    }\PY{p}{\PYZcb{}}
\PY{p}{\PYZcb{}}
\end{Verbatim}
\end{codebox}

\subsection[Serializable vs. Externalizable]{Serializable vs. Externalizable}
\label{h:10:serializable-vs-externalizable}
{\tablefontsmall
\begin{xltabular}{\linewidth}{@{}L{0.430}L{1.068}L{1.502}@{}}
\toprule
\rowcolor{tablehdr}
\thd{} & \thd{\code{Serializable}} & \thd{\code{Externalizable}} \\
\midrule
\endhead
Implementation effort & None (marker interface) — JVM uses reflection & Must implement \code{writeExternal}/\code{readExternal} yourself \\
No-arg constructor & Not required (fields set via reflection, bypassing constructors) & \textbf{Required} and must be \code{public} \\
Performance & Slower — reflection over all fields & Can be faster — you control exactly what’s written, no reflection over field metadata \\
Format control & Default format mirrors class layout; can customize via \code{writeObject}/\code{readObject} & Total control — you decide every byte \\
Versioning safety & \code{serialVersionUID} + default mechanism handles a lot automatically & You must handle format evolution entirely yourself \\
Risk of mistakes & Lower — JVM does the heavy lifting & Higher — a forgotten field, wrong read order, or missing no-arg constructor breaks everything silently or loudly \\
Typical use & Most day-to-day cases (when serialization is genuinely needed) & Performance-critical or space-constrained scenarios, or when you need a stable, hand-controlled wire format \\
\bottomrule
\end{xltabular}
}

The trade-off in one sentence: \code{Externalizable} gives you speed and control, but takes away every safety net — get the field order wrong between \code{writeExternal} and \code{readExternal}, and you’ll silently read corrupted data with no compiler or JVM warning.

\section[Common Code-Review Interview Pitfalls]{Common Code-Review Interview Pitfalls}
\label{h:10:common-code-review-interview-pitfalls}
\begin{enumerate}
\item 
\textbf{Not closing streams / using try-with-resources.}
\emph{Why it matters}: leaked file handles or sockets eventually exhaust OS limits, causing mysterious failures under load.

\begin{codebox}
\begin{Verbatim}[commandchars=\\\{\}]
\PY{c+c1}{// Before}
\PY{n}{FileInputStream}\PY{+w}{ }\PY{n}{in}\PY{+w}{ }\PY{o}{=}\PY{+w}{ }\PY{k}{new}\PY{+w}{ }\PY{n}{FileInputStream}\PY{p}{(}\PY{l+s}{\PYZdq{}}\PY{l+s}{data.txt}\PY{l+s}{\PYZdq{}}\PY{p}{)}\PY{p}{;}
\PY{c+c1}{// ... in.close() forgotten, or skipped if an exception is thrown above it}

\PY{c+c1}{// After}
\PY{k}{try}\PY{+w}{ }\PY{p}{(}\PY{n}{FileInputStream}\PY{+w}{ }\PY{n}{in}\PY{+w}{ }\PY{o}{=}\PY{+w}{ }\PY{k}{new}\PY{+w}{ }\PY{n}{FileInputStream}\PY{p}{(}\PY{l+s}{\PYZdq{}}\PY{l+s}{data.txt}\PY{l+s}{\PYZdq{}}\PY{p}{)}\PY{p}{)}\PY{+w}{ }\PY{p}{\PYZob{}}
\PY{+w}{    }\PY{c+c1}{// ...}
\PY{p}{\PYZcb{}}
\end{Verbatim}
\end{codebox}

\item 
\textbf{Using platform-default charset instead of an explicit one.}
\emph{Why it matters}: the exact same code produces different bytes/text on different machines or JVM configurations.

\begin{codebox}
\begin{Verbatim}[commandchars=\\\{\}]
\PY{c+c1}{// Before}
\PY{n}{Writer}\PY{+w}{ }\PY{n}{w}\PY{+w}{ }\PY{o}{=}\PY{+w}{ }\PY{k}{new}\PY{+w}{ }\PY{n}{FileWriter}\PY{p}{(}\PY{l+s}{\PYZdq{}}\PY{l+s}{out.txt}\PY{l+s}{\PYZdq{}}\PY{p}{)}\PY{p}{;}

\PY{c+c1}{// After}
\PY{n}{Writer}\PY{+w}{ }\PY{n}{w}\PY{+w}{ }\PY{o}{=}\PY{+w}{ }\PY{k}{new}\PY{+w}{ }\PY{n}{FileWriter}\PY{p}{(}\PY{l+s}{\PYZdq{}}\PY{l+s}{out.txt}\PY{l+s}{\PYZdq{}}\PY{p}{,}\PY{+w}{ }\PY{n}{StandardCharsets}\PY{p}{.}\PY{n+na}{UTF\PYZus{}8}\PY{p}{)}\PY{p}{;}
\end{Verbatim}
\end{codebox}

\item 
\textbf{Reading byte-by-byte or line-by-line without buffering.}
\emph{Why it matters}: each unbuffered call can trigger a system call — brutal performance hit on large files.

\begin{codebox}
\begin{Verbatim}[commandchars=\\\{\}]
\PY{c+c1}{// Before}
\PY{n}{InputStream}\PY{+w}{ }\PY{n}{in}\PY{+w}{ }\PY{o}{=}\PY{+w}{ }\PY{k}{new}\PY{+w}{ }\PY{n}{FileInputStream}\PY{p}{(}\PY{l+s}{\PYZdq{}}\PY{l+s}{big.dat}\PY{l+s}{\PYZdq{}}\PY{p}{)}\PY{p}{;}

\PY{c+c1}{// After}
\PY{n}{InputStream}\PY{+w}{ }\PY{n}{in}\PY{+w}{ }\PY{o}{=}\PY{+w}{ }\PY{k}{new}\PY{+w}{ }\PY{n}{BufferedInputStream}\PY{p}{(}\PY{k}{new}\PY{+w}{ }\PY{n}{FileInputStream}\PY{p}{(}\PY{l+s}{\PYZdq{}}\PY{l+s}{big.dat}\PY{l+s}{\PYZdq{}}\PY{p}{)}\PY{p}{)}\PY{p}{;}
\end{Verbatim}
\end{codebox}

\item 
\textbf{Trusting \code{File.\allowbreak{}delete()} / \code{File.\allowbreak{}mkdir()} booleans without understanding they hide the failure reason.}
\emph{Why it matters}: silent \code{false} gives no diagnostic information; bugs become “it just doesn’t work” reports.

\begin{codebox}
\begin{Verbatim}[commandchars=\\\{\}]
\PY{c+c1}{// Before}
\PY{k+kt}{boolean}\PY{+w}{ }\PY{n}{ok}\PY{+w}{ }\PY{o}{=}\PY{+w}{ }\PY{n}{file}\PY{p}{.}\PY{n+na}{delete}\PY{p}{(}\PY{p}{)}\PY{p}{;}

\PY{c+c1}{// After}
\PY{n}{Files}\PY{p}{.}\PY{n+na}{delete}\PY{p}{(}\PY{n}{path}\PY{p}{)}\PY{p}{;}\PY{+w}{ }\PY{c+c1}{// throws NoSuchFileException, AccessDeniedException, etc.}
\end{Verbatim}
\end{codebox}

\item 
\textbf{Not closing \code{Files.\allowbreak{}lines()} or \code{Files.\allowbreak{}walk()} streams.}
\emph{Why it matters}: these hold an open file/directory handle under the hood even though they “look like” a plain in-memory \code{Stream}.

\begin{codebox}
\begin{Verbatim}[commandchars=\\\{\}]
\PY{c+c1}{// Before}
\PY{n}{Files}\PY{p}{.}\PY{n+na}{lines}\PY{p}{(}\PY{n}{path}\PY{p}{)}\PY{p}{.}\PY{n+na}{forEach}\PY{p}{(}\PY{n}{System}\PY{p}{.}\PY{n+na}{out}\PY{p}{:}\PY{p}{:}\PY{n}{println}\PY{p}{)}\PY{p}{;}\PY{+w}{ }\PY{c+c1}{// handle never released}

\PY{c+c1}{// After}
\PY{k}{try}\PY{+w}{ }\PY{p}{(}\PY{n}{Stream}\PY{o}{\PYZlt{}}\PY{n}{String}\PY{o}{\PYZgt{}}\PY{+w}{ }\PY{n}{lines}\PY{+w}{ }\PY{o}{=}\PY{+w}{ }\PY{n}{Files}\PY{p}{.}\PY{n+na}{lines}\PY{p}{(}\PY{n}{path}\PY{p}{)}\PY{p}{)}\PY{+w}{ }\PY{p}{\PYZob{}}
\PY{+w}{    }\PY{n}{lines}\PY{p}{.}\PY{n+na}{forEach}\PY{p}{(}\PY{n}{System}\PY{p}{.}\PY{n+na}{out}\PY{p}{:}\PY{p}{:}\PY{n}{println}\PY{p}{)}\PY{p}{;}
\PY{p}{\PYZcb{}}
\end{Verbatim}
\end{codebox}

\item 
\textbf{Building file paths from user input without validating against path traversal (\code{../}).}
\emph{Why it matters}: lets attackers read/write arbitrary files outside the intended directory.

\begin{codebox}
\begin{Verbatim}[commandchars=\\\{\}]
\PY{c+c1}{// Before}
\PY{n}{Path}\PY{+w}{ }\PY{n}{p}\PY{+w}{ }\PY{o}{=}\PY{+w}{ }\PY{n}{Path}\PY{p}{.}\PY{n+na}{of}\PY{p}{(}\PY{l+s}{\PYZdq{}}\PY{l+s}{/uploads}\PY{l+s}{\PYZdq{}}\PY{p}{)}\PY{p}{.}\PY{n+na}{resolve}\PY{p}{(}\PY{n}{userInput}\PY{p}{)}\PY{p}{;}

\PY{c+c1}{// After}
\PY{n}{Path}\PY{+w}{ }\PY{n}{root}\PY{+w}{ }\PY{o}{=}\PY{+w}{ }\PY{n}{Path}\PY{p}{.}\PY{n+na}{of}\PY{p}{(}\PY{l+s}{\PYZdq{}}\PY{l+s}{/uploads}\PY{l+s}{\PYZdq{}}\PY{p}{)}\PY{p}{.}\PY{n+na}{toRealPath}\PY{p}{(}\PY{p}{)}\PY{p}{;}
\PY{n}{Path}\PY{+w}{ }\PY{n}{p}\PY{+w}{ }\PY{o}{=}\PY{+w}{ }\PY{n}{root}\PY{p}{.}\PY{n+na}{resolve}\PY{p}{(}\PY{n}{userInput}\PY{p}{)}\PY{p}{.}\PY{n+na}{normalize}\PY{p}{(}\PY{p}{)}\PY{p}{;}
\PY{k}{if}\PY{+w}{ }\PY{p}{(}\PY{o}{!}\PY{n}{p}\PY{p}{.}\PY{n+na}{startsWith}\PY{p}{(}\PY{n}{root}\PY{p}{)}\PY{p}{)}\PY{+w}{ }\PY{k}{throw}\PY{+w}{ }\PY{k}{new}\PY{+w}{ }\PY{n}{SecurityException}\PY{p}{(}\PY{l+s}{\PYZdq{}}\PY{l+s}{Traversal attempt}\PY{l+s}{\PYZdq{}}\PY{p}{)}\PY{p}{;}
\end{Verbatim}
\end{codebox}

\item 
\textbf{Writing directly to a target file instead of write-temp-then-atomic-move.}
\emph{Why it matters}: a crash mid-write leaves a corrupted, half-written file that readers may see.

\begin{codebox}
\begin{Verbatim}[commandchars=\\\{\}]
\PY{c+c1}{// Before}
\PY{n}{Files}\PY{p}{.}\PY{n+na}{writeString}\PY{p}{(}\PY{n}{target}\PY{p}{,}\PY{+w}{ }\PY{n}{content}\PY{p}{)}\PY{p}{;}

\PY{c+c1}{// After}
\PY{n}{Path}\PY{+w}{ }\PY{n}{tmp}\PY{+w}{ }\PY{o}{=}\PY{+w}{ }\PY{n}{Files}\PY{p}{.}\PY{n+na}{createTempFile}\PY{p}{(}\PY{n}{target}\PY{p}{.}\PY{n+na}{getParent}\PY{p}{(}\PY{p}{)}\PY{p}{,}\PY{+w}{ }\PY{l+s}{\PYZdq{}}\PY{l+s}{tmp\PYZhy{}}\PY{l+s}{\PYZdq{}}\PY{p}{,}\PY{+w}{ }\PY{l+s}{\PYZdq{}}\PY{l+s}{.swap}\PY{l+s}{\PYZdq{}}\PY{p}{)}\PY{p}{;}
\PY{n}{Files}\PY{p}{.}\PY{n+na}{writeString}\PY{p}{(}\PY{n}{tmp}\PY{p}{,}\PY{+w}{ }\PY{n}{content}\PY{p}{)}\PY{p}{;}
\PY{n}{Files}\PY{p}{.}\PY{n+na}{move}\PY{p}{(}\PY{n}{tmp}\PY{p}{,}\PY{+w}{ }\PY{n}{target}\PY{p}{,}\PY{+w}{ }\PY{n}{StandardCopyOption}\PY{p}{.}\PY{n+na}{ATOMIC\PYZus{}MOVE}\PY{p}{,}\PY{+w}{ }\PY{n}{StandardCopyOption}\PY{p}{.}\PY{n+na}{REPLACE\PYZus{}EXISTING}\PY{p}{)}\PY{p}{;}
\end{Verbatim}
\end{codebox}

\item 
\textbf{\code{Serializable} class with no explicit \code{serialVersionUID}.}
\emph{Why it matters}: any future change to the class can change the auto-computed ID, breaking deserialization of previously-saved data.

\begin{codebox}
\begin{Verbatim}[commandchars=\\\{\}]
\PY{c+c1}{// Before}
\PY{k+kd}{class} \PY{n+nc}{User}\PY{+w}{ }\PY{k+kd}{implements}\PY{+w}{ }\PY{n}{Serializable}\PY{+w}{ }\PY{p}{\PYZob{}}\PY{+w}{ }\PY{n}{String}\PY{+w}{ }\PY{n}{name}\PY{p}{;}\PY{+w}{ }\PY{p}{\PYZcb{}}

\PY{c+c1}{// After}
\PY{k+kd}{class} \PY{n+nc}{User}\PY{+w}{ }\PY{k+kd}{implements}\PY{+w}{ }\PY{n}{Serializable}\PY{+w}{ }\PY{p}{\PYZob{}}
\PY{+w}{    }\PY{k+kd}{private}\PY{+w}{ }\PY{k+kd}{static}\PY{+w}{ }\PY{k+kd}{final}\PY{+w}{ }\PY{k+kt}{long}\PY{+w}{ }\PY{n}{serialVersionUID}\PY{+w}{ }\PY{o}{=}\PY{+w}{ }\PY{l+m+mi}{1L}\PY{p}{;}
\PY{+w}{    }\PY{n}{String}\PY{+w}{ }\PY{n}{name}\PY{p}{;}
\PY{p}{\PYZcb{}}
\end{Verbatim}
\end{codebox}

\item 
\textbf{Deserializing untrusted input with plain \code{ObjectInputStream}.}
\emph{Why it matters}: this is the entry point for deserialization gadget-chain attacks — potential remote code execution.

\begin{codebox}
\begin{Verbatim}[commandchars=\\\{\}]
\PY{c+c1}{// Before}
\PY{n}{Object}\PY{+w}{ }\PY{n}{obj}\PY{+w}{ }\PY{o}{=}\PY{+w}{ }\PY{k}{new}\PY{+w}{ }\PY{n}{ObjectInputStream}\PY{p}{(}\PY{n}{untrustedInputStream}\PY{p}{)}\PY{p}{.}\PY{n+na}{readObject}\PY{p}{(}\PY{p}{)}\PY{p}{;}

\PY{c+c1}{// After: use a JSON/protobuf library instead, or at minimum add an ObjectInputFilter}
\PY{n}{ObjectInputStream}\PY{+w}{ }\PY{n}{ois}\PY{+w}{ }\PY{o}{=}\PY{+w}{ }\PY{k}{new}\PY{+w}{ }\PY{n}{ObjectInputStream}\PY{p}{(}\PY{n}{untrustedInputStream}\PY{p}{)}\PY{p}{;}
\PY{n}{ois}\PY{p}{.}\PY{n+na}{setObjectInputFilter}\PY{p}{(}\PY{n}{ObjectInputFilter}\PY{p}{.}\PY{n+na}{Config}\PY{p}{.}\PY{n+na}{createFilter}\PY{p}{(}\PY{l+s}{\PYZdq{}}\PY{l+s}{com.example.*;!*}\PY{l+s}{\PYZdq{}}\PY{p}{)}\PY{p}{)}\PY{p}{;}
\PY{n}{Object}\PY{+w}{ }\PY{n}{obj}\PY{+w}{ }\PY{o}{=}\PY{+w}{ }\PY{n}{ois}\PY{p}{.}\PY{n+na}{readObject}\PY{p}{(}\PY{p}{)}\PY{p}{;}
\end{Verbatim}
\end{codebox}

\item 
\textbf{Forgetting sensitive fields need \code{transient}.}
\emph{Why it matters}: secrets, tokens, or passwords silently end up serialized to disk or over the wire.

\begin{codebox}
\begin{Verbatim}[commandchars=\\\{\}]
\PY{c+c1}{// Before}
\PY{k+kd}{class} \PY{n+nc}{Session}\PY{+w}{ }\PY{k+kd}{implements}\PY{+w}{ }\PY{n}{Serializable}\PY{+w}{ }\PY{p}{\PYZob{}}\PY{+w}{ }\PY{n}{String}\PY{+w}{ }\PY{n}{authToken}\PY{p}{;}\PY{+w}{ }\PY{p}{\PYZcb{}}

\PY{c+c1}{// After}
\PY{k+kd}{class} \PY{n+nc}{Session}\PY{+w}{ }\PY{k+kd}{implements}\PY{+w}{ }\PY{n}{Serializable}\PY{+w}{ }\PY{p}{\PYZob{}}\PY{+w}{ }\PY{k+kd}{transient}\PY{+w}{ }\PY{n}{String}\PY{+w}{ }\PY{n}{authToken}\PY{p}{;}\PY{+w}{ }\PY{p}{\PYZcb{}}
\end{Verbatim}
\end{codebox}

\item 
\textbf{\code{Externalizable} class missing the required public no-arg constructor.}
\emph{Why it matters}: deserialization throws \code{InvalidClassException} at runtime — the class \emph{looks} correct until you actually try to deserialize it.

\begin{codebox}
\begin{Verbatim}[commandchars=\\\{\}]
\PY{c+c1}{// Before}
\PY{k+kd}{class} \PY{n+nc}{Product}\PY{+w}{ }\PY{k+kd}{implements}\PY{+w}{ }\PY{n}{Externalizable}\PY{+w}{ }\PY{p}{\PYZob{}}
\PY{+w}{    }\PY{n}{Product}\PY{p}{(}\PY{n}{String}\PY{+w}{ }\PY{n}{name}\PY{p}{)}\PY{+w}{ }\PY{p}{\PYZob{}}\PY{+w}{ }\PY{k}{this}\PY{p}{.}\PY{n+na}{name}\PY{+w}{ }\PY{o}{=}\PY{+w}{ }\PY{n}{name}\PY{p}{;}\PY{+w}{ }\PY{p}{\PYZcb{}}\PY{+w}{ }\PY{c+c1}{// only constructor}

\PY{c+c1}{// After}
\PY{+w}{    }\PY{k+kd}{public}\PY{+w}{ }\PY{n+nf}{Product}\PY{p}{(}\PY{p}{)}\PY{+w}{ }\PY{p}{\PYZob{}}\PY{p}{\PYZcb{}}\PY{+w}{              }\PY{c+c1}{// required}
\PY{+w}{    }\PY{n}{Product}\PY{p}{(}\PY{n}{String}\PY{+w}{ }\PY{n}{name}\PY{p}{)}\PY{+w}{ }\PY{p}{\PYZob{}}\PY{+w}{ }\PY{k}{this}\PY{p}{.}\PY{n+na}{name}\PY{+w}{ }\PY{o}{=}\PY{+w}{ }\PY{n}{name}\PY{p}{;}\PY{+w}{ }\PY{p}{\PYZcb{}}
\end{Verbatim}
\end{codebox}

\item 
\textbf{Mismatched field order between \code{writeExternal} and \code{readExternal}.}
\emph{Why it matters}: \code{Externalizable} has zero safety net — a swapped read order silently produces corrupted objects with no exception.

\begin{codebox}
\begin{Verbatim}[commandchars=\\\{\}]
\PY{c+c1}{// Before: writeExternal writes name then price, readExternal reads price then name}
\PY{n}{out}\PY{p}{.}\PY{n+na}{writeUTF}\PY{p}{(}\PY{n}{name}\PY{p}{)}\PY{p}{;}\PY{+w}{ }\PY{n}{out}\PY{p}{.}\PY{n+na}{writeDouble}\PY{p}{(}\PY{n}{price}\PY{p}{)}\PY{p}{;}
\PY{n}{price}\PY{+w}{ }\PY{o}{=}\PY{+w}{ }\PY{n}{in}\PY{p}{.}\PY{n+na}{readDouble}\PY{p}{(}\PY{p}{)}\PY{p}{;}\PY{+w}{ }\PY{n}{name}\PY{+w}{ }\PY{o}{=}\PY{+w}{ }\PY{n}{in}\PY{p}{.}\PY{n+na}{readUTF}\PY{p}{(}\PY{p}{)}\PY{p}{;}\PY{+w}{ }\PY{c+c1}{// WRONG ORDER}

\PY{c+c1}{// After: keep read/write order identical, ideally documented}
\PY{n}{out}\PY{p}{.}\PY{n+na}{writeUTF}\PY{p}{(}\PY{n}{name}\PY{p}{)}\PY{p}{;}\PY{+w}{ }\PY{n}{out}\PY{p}{.}\PY{n+na}{writeDouble}\PY{p}{(}\PY{n}{price}\PY{p}{)}\PY{p}{;}
\PY{n}{name}\PY{+w}{ }\PY{o}{=}\PY{+w}{ }\PY{n}{in}\PY{p}{.}\PY{n+na}{readUTF}\PY{p}{(}\PY{p}{)}\PY{p}{;}\PY{+w}{ }\PY{n}{price}\PY{+w}{ }\PY{o}{=}\PY{+w}{ }\PY{n}{in}\PY{p}{.}\PY{n+na}{readDouble}\PY{p}{(}\PY{p}{)}\PY{p}{;}
\end{Verbatim}
\end{codebox}

\item 
\textbf{Using \code{ByteBuffer} without calling \code{flip()} before reading.}
\emph{Why it matters}: reading resumes from wherever \code{position} was left after writing, producing empty or garbage results instead of the intended data.

\begin{codebox}
\begin{Verbatim}[commandchars=\\\{\}]
\PY{c+c1}{// Before}
\PY{n}{buffer}\PY{p}{.}\PY{n+na}{put}\PY{p}{(}\PY{n}{data}\PY{p}{)}\PY{p}{;}
\PY{n}{buffer}\PY{p}{.}\PY{n+na}{get}\PY{p}{(}\PY{n}{result}\PY{p}{)}\PY{p}{;}\PY{+w}{ }\PY{c+c1}{// reads from wrong position, likely returns nothing useful}

\PY{c+c1}{// After}
\PY{n}{buffer}\PY{p}{.}\PY{n+na}{put}\PY{p}{(}\PY{n}{data}\PY{p}{)}\PY{p}{;}
\PY{n}{buffer}\PY{p}{.}\PY{n+na}{flip}\PY{p}{(}\PY{p}{)}\PY{p}{;}\PY{+w}{ }\PY{c+c1}{// limit = position, position = 0}
\PY{n}{buffer}\PY{p}{.}\PY{n+na}{get}\PY{p}{(}\PY{n}{result}\PY{p}{)}\PY{p}{;}
\end{Verbatim}
\end{codebox}

\item 
\textbf{Using \code{java.\allowbreak{}io.\allowbreak{}File.\allowbreak{}create\allowbreak{}Temp\allowbreak{}File} for sensitive data instead of NIO.2’s secure temp file creation.}
\emph{Why it matters}: \code{File.\allowbreak{}createTempFile} predictable/world-readable permissions on some platforms invite race conditions and information leaks; \code{Files.\allowbreak{}createTempFile} sets owner-only permissions on POSIX systems by default.

\begin{codebox}
\begin{Verbatim}[commandchars=\\\{\}]
\PY{c+c1}{// Before}
\PY{n}{File}\PY{+w}{ }\PY{n}{temp}\PY{+w}{ }\PY{o}{=}\PY{+w}{ }\PY{n}{File}\PY{p}{.}\PY{n+na}{createTempFile}\PY{p}{(}\PY{l+s}{\PYZdq{}}\PY{l+s}{upload}\PY{l+s}{\PYZdq{}}\PY{p}{,}\PY{+w}{ }\PY{l+s}{\PYZdq{}}\PY{l+s}{.tmp}\PY{l+s}{\PYZdq{}}\PY{p}{)}\PY{p}{;}

\PY{c+c1}{// After}
\PY{n}{Path}\PY{+w}{ }\PY{n}{temp}\PY{+w}{ }\PY{o}{=}\PY{+w}{ }\PY{n}{Files}\PY{p}{.}\PY{n+na}{createTempFile}\PY{p}{(}\PY{l+s}{\PYZdq{}}\PY{l+s}{upload\PYZhy{}}\PY{l+s}{\PYZdq{}}\PY{p}{,}\PY{+w}{ }\PY{l+s}{\PYZdq{}}\PY{l+s}{.tmp}\PY{l+s}{\PYZdq{}}\PY{p}{)}\PY{p}{;}
\end{Verbatim}
\end{codebox}

\item 
\textbf{Choosing \code{Files.\allowbreak{}walk} (or recursive \code{File} listing) when directory traversal needs error handling or early termination per subtree.}
\emph{Why it matters}: a single unreadable subdirectory can blow up an entire \code{Stream} pipeline with no way to skip just that branch; \code{walkFileTree} lets you react per-file/per-directory and continue.

\begin{codebox}
\begin{Verbatim}[commandchars=\\\{\}]
\PY{c+c1}{// Before}
\PY{n}{Files}\PY{p}{.}\PY{n+na}{walk}\PY{p}{(}\PY{n}{root}\PY{p}{)}\PY{p}{.}\PY{n+na}{forEach}\PY{p}{(}\PY{k}{this}\PY{p}{:}\PY{p}{:}\PY{n}{process}\PY{p}{)}\PY{p}{;}\PY{+w}{ }\PY{c+c1}{// one bad file aborts everything}

\PY{c+c1}{// After}
\PY{n}{Files}\PY{p}{.}\PY{n+na}{walkFileTree}\PY{p}{(}\PY{n}{root}\PY{p}{,}\PY{+w}{ }\PY{k}{new}\PY{+w}{ }\PY{n}{SimpleFileVisitor}\PY{o}{\PYZlt{}}\PY{n}{Path}\PY{o}{\PYZgt{}}\PY{p}{(}\PY{p}{)}\PY{+w}{ }\PY{p}{\PYZob{}}
\PY{+w}{    }\PY{n+nd}{@Override}
\PY{+w}{    }\PY{k+kd}{public}\PY{+w}{ }\PY{n}{FileVisitResult}\PY{+w}{ }\PY{n+nf}{visitFileFailed}\PY{p}{(}\PY{n}{Path}\PY{+w}{ }\PY{n}{file}\PY{p}{,}\PY{+w}{ }\PY{n}{IOException}\PY{+w}{ }\PY{n}{exc}\PY{p}{)}\PY{+w}{ }\PY{p}{\PYZob{}}
\PY{+w}{        }\PY{n}{log}\PY{p}{.}\PY{n+na}{warn}\PY{p}{(}\PY{l+s}{\PYZdq{}}\PY{l+s}{Skipping unreadable file: \PYZob{}\PYZcb{}}\PY{l+s}{\PYZdq{}}\PY{p}{,}\PY{+w}{ }\PY{n}{file}\PY{p}{)}\PY{p}{;}
\PY{+w}{        }\PY{k}{return}\PY{+w}{ }\PY{n}{FileVisitResult}\PY{p}{.}\PY{n+na}{CONTINUE}\PY{p}{;}
\PY{+w}{    }\PY{p}{\PYZcb{}}
\PY{+w}{    }\PY{n+nd}{@Override}
\PY{+w}{    }\PY{k+kd}{public}\PY{+w}{ }\PY{n}{FileVisitResult}\PY{+w}{ }\PY{n+nf}{visitFile}\PY{p}{(}\PY{n}{Path}\PY{+w}{ }\PY{n}{file}\PY{p}{,}\PY{+w}{ }\PY{n}{BasicFileAttributes}\PY{+w}{ }\PY{n}{attrs}\PY{p}{)}\PY{+w}{ }\PY{p}{\PYZob{}}
\PY{+w}{        }\PY{n}{process}\PY{p}{(}\PY{n}{file}\PY{p}{)}\PY{p}{;}
\PY{+w}{        }\PY{k}{return}\PY{+w}{ }\PY{n}{FileVisitResult}\PY{p}{.}\PY{n+na}{CONTINUE}\PY{p}{;}
\PY{+w}{    }\PY{p}{\PYZcb{}}
\PY{p}{\PYZcb{}}\PY{p}{)}\PY{p}{;}
\end{Verbatim}
\end{codebox}

\end{enumerate}


\setcounter{chapter}{10}
\chapter{JVM Fundamentals}
\label{chap:11}
\label{h:11:11-jvm-fundamentals}
Every Java program you review eventually runs on the JVM (Java Virtual Machine), so understanding what happens between \code{javac} and a running thread is essential for spotting real bugs — not just style issues. This chapter walks through the JDK/JRE/JVM split, how source becomes bytecode and gets loaded and run, how the JIT (Just-In-Time compiler) makes hot code fast, how class loaders isolate and share code, and what the Java Memory Model actually guarantees about threads seeing each other’s writes. All examples target Java 21+ on HotSpot, the default JVM shipped in OpenJDK.

\section[JDK, JRE, JVM Architecture]{JDK, JRE, JVM Architecture}
\label{h:11:jdk-jre-jvm-architecture}
Three acronyms get thrown around loosely. Here is what each one actually is.

{\tablefont
\begin{xltabular}{\linewidth}{@{}L{0.403}L{0.547}L{2.050}@{}}
\toprule
\rowcolor{tablehdr}
\thd{Term} & \thd{Full name} & \thd{What it contains} \\
\midrule
\endhead
JVM & Java Virtual Machine & The engine that executes bytecode: class loader, runtime memory areas, interpreter, JIT compiler, garbage collector. \\
JRE & Java Runtime Environment & JVM + core class libraries (\code{java.\allowbreak{}lang}, \code{java.\allowbreak{}util}, …) needed to \emph{run} compiled programs. No compiler. \\
JDK & Java Development Kit & JRE + development tools: \code{javac} (compiler), \code{javap} (disassembler), \code{jar}, \code{jlink}, \code{jshell}, debuggers, etc. Needed to \emph{build} programs. \\
\bottomrule
\end{xltabular}
}

\textbf{Important nuance for Java 11+:} the JRE as a separately downloadable, separately installed artifact was removed starting with Java 11. Oracle and OpenJDK now ship only the JDK. If you need a smaller, runtime-only image (the old JRE’s job), you build one yourself with \code{jlink}, which assembles a custom minimal runtime image containing only the modules your application actually uses.

\begin{codebox}
\begin{Verbatim}[commandchars=\\\{\}]
\PY{c+c1}{\PYZsh{} Build a custom \PYZdq{}JRE\PYZhy{}like\PYZdq{} runtime image containing only java.base and java.sql}
jlink\PY{+w}{ }\PYZhy{}\PYZhy{}add\PYZhy{}modules\PY{+w}{ }java.base,java.sql\PY{+w}{ }\PY{l+s+se}{\PYZbs{}}
\PY{+w}{      }\PYZhy{}\PYZhy{}output\PY{+w}{ }my\PYZhy{}custom\PYZhy{}runtime\PY{+w}{ }\PY{l+s+se}{\PYZbs{}}
\PY{+w}{      }\PYZhy{}\PYZhy{}strip\PYZhy{}debug\PY{+w}{ }\PYZhy{}\PYZhy{}no\PYZhy{}header\PYZhy{}files\PY{+w}{ }\PYZhy{}\PYZhy{}no\PYZhy{}man\PYZhy{}pages

\PY{c+c1}{\PYZsh{} Run your app with it instead of a full JDK}
my\PYZhy{}custom\PYZhy{}runtime/bin/java\PY{+w}{ }\PYZhy{}jar\PY{+w}{ }app.jar
\end{Verbatim}
\end{codebox}

\subsection[JVM runtime data areas]{JVM runtime data areas}
\label{h:11:jvm-runtime-data-areas}
When the JVM starts, it carves up memory into distinct regions. Some are shared by all threads, some exist one-per-thread. Getting this picture straight is the foundation for understanding \code{OutOfMemoryError} variants, thread-safety, and GC behavior later.

\begin{codebox}
\begin{Verbatim}
                     JVM PROCESS
   ┌────────────────────────────────────────────────────────────┐
   │  SHARED (one per JVM instance)                             │
   │                                                            │
   │   ┌─────────────────────┐   ┌─────────────────────────┐    │
   │   │        HEAP         │   │       METASPACE         │    │
   │   │ objects, arrays,    │   │ class metadata, method  │    │
   │   │ instance fields     │   │ bytecode, constant pool │    │
   │   │ (Young + Old gen)   │   │ (native memory, not     │    │
   │   │ GC-managed          │   │  bound by -Xmx)         │    │
   │   └─────────────────────┘   └─────────────────────────┘    │
   │                                                            │
   │   ┌─────────────────────────────────────────────────┐      │
   │   │              CODE CACHE                          │      │
   │   │  JIT-compiled native machine code (C1/C2 output) │      │
   │   └─────────────────────────────────────────────────┘      │
   ├────────────────────────────────────────────────────────────┤
   │  PER-THREAD (one set per Java thread)                      │
   │                                                            │
   │  Thread A                    Thread B                      │
   │  ┌───────────────┐           ┌───────────────┐             │
   │  │  JVM STACK    │           │  JVM STACK    │             │
   │  │ frames: local │           │ frames: local │             │
   │  │ vars, operand │           │ vars, operand │             │
   │  │ stack, refs   │           │ stack, refs   │             │
   │  ├───────────────┤           ├───────────────┤             │
   │  │  PC REGISTER  │           │  PC REGISTER  │             │
   │  │ next bytecode │           │ next bytecode │             │
   │  │ instruction   │           │ instruction   │             │
   │  ├───────────────┤           ├───────────────┤             │
   │  │ NATIVE METHOD │           │ NATIVE METHOD │             │
   │  │ STACK (JNI)   │           │ STACK (JNI)   │             │
   │  └───────────────┘           └───────────────┘             │
   └────────────────────────────────────────────────────────────┘
\end{Verbatim}
\end{codebox}

{\tablefont
\begin{xltabular}{\linewidth}{@{}L{0.431}L{0.431}L{2.048}L{1.090}@{}}
\toprule
\rowcolor{tablehdr}
\thd{Region} & \thd{Shared or per-thread} & \thd{Holds} & \thd{Overflow error} \\
\midrule
\endhead
Heap & Shared & All objects and arrays, instance fields & \code{Out\allowbreak{}Of\allowbreak{}Memory\allowbreak{}Error:\allowbreak{}~\allowbreak{}Java~\allowbreak{}heap~\allowbreak{}space} \\
Metaspace & Shared & Class metadata, method bytecode, constant pools (replaced PermGen since Java 8) & \code{Out\allowbreak{}Of\allowbreak{}Memory\allowbreak{}Error:\allowbreak{}~\allowbreak{}Metaspace} \\
Code Cache & Shared & Native machine code produced by the JIT & \code{Out\allowbreak{}Of\allowbreak{}Memory\allowbreak{}Error:\allowbreak{}~\allowbreak{}Code\allowbreak{}Cache} (rare, tune \code{-\allowbreak{}XX:\allowbreak{}Reserved\allowbreak{}Code\allowbreak{}Cache\allowbreak{}Size}) \\
JVM Stack & Per-thread & Stack frames: local variables, operand stack, return addresses & \code{StackOverflowError} \\
PC Register & Per-thread & Address of the currently executing bytecode instruction & n/a \\
Native Method Stack & Per-thread & State for native (JNI) method calls & \code{StackOverflowError} (native) \\
\bottomrule
\end{xltabular}
}

A quick mental model: \textbf{heap is for objects, stacks are for method calls, metaspace is for the classes themselves.} Local primitive variables and object \emph{references} live on the stack; the objects they point to live on the heap.

\begin{codebox}
\begin{Verbatim}[commandchars=\\\{\}]
\PY{k+kt}{void}\PY{+w}{ }\PY{n+nf}{example}\PY{p}{(}\PY{p}{)}\PY{+w}{ }\PY{p}{\PYZob{}}
\PY{+w}{    }\PY{k+kt}{int}\PY{+w}{ }\PY{n}{x}\PY{+w}{ }\PY{o}{=}\PY{+w}{ }\PY{l+m+mi}{42}\PY{p}{;}\PY{+w}{                 }\PY{c+c1}{// primitive lives on this thread\PYZsq{}s stack frame}
\PY{+w}{    }\PY{n}{StringBuilder}\PY{+w}{ }\PY{n}{sb}\PY{+w}{ }\PY{o}{=}\PY{+w}{ }\PY{k}{new}\PY{+w}{ }\PY{n}{StringBuilder}\PY{p}{(}\PY{p}{)}\PY{p}{;}\PY{+w}{ }\PY{c+c1}{// \PYZdq{}sb\PYZdq{} reference on stack,}
\PY{+w}{                                             }\PY{c+c1}{// the StringBuilder object itself on the heap}
\PY{p}{\PYZcb{}}
\end{Verbatim}
\end{codebox}

\section[Compilation, Class Loading, and Execution]{Compilation, Class Loading, and Execution}
\label{h:11:compilation-class-loading-and-execution}
Getting from a \code{.\allowbreak{}java} file to running code has five distinct stages. Interview reviewers care whether you know \emph{when} things like static initializers actually fire, because that is where subtle bugs (and infinite loops, and \code{NoClassDefFoundError}s from failed static blocks) hide.

\begin{codebox}
\begin{Verbatim}
 source.java --javac--> source.class (bytecode)
       │
       ▼
 ┌───────────┐   ┌────────────────────────────┐   ┌────────────┐   ┌──────────┐
 │  LOADING  │→  │         LINKING            │→  │INITIALIZATION│→ │ EXECUTION│
 │ find      │   │  1. Verify (bytecode is    │   │ run static  │   │ run main │
 │ bytes,    │   │     safe & well-formed)    │   │ initializers│   │ or invoked│
 │ create    │   │  2. Prepare (allocate      │   │ and static  │   │ method   │
 │ Class obj │   │     static fields, defaults)│   │ blocks, top │   │           │
 │           │   │  3. Resolve (symbolic refs │   │ to bottom   │   │           │
 │           │   │     → direct references,   │   │             │   │           │
 │           │   │     can be lazy)           │   │             │   │           │
 └───────────┘   └────────────────────────────┘   └────────────┘   └──────────┘
\end{Verbatim}
\end{codebox}

\begin{enumerate}
\item 
\textbf{Compilation (\code{javac}):} Human-readable \code{.\allowbreak{}java} source is compiled into platform-independent bytecode stored in a \code{.\allowbreak{}class} file.

\item 
\textbf{Loading:} A class loader locates the bytes (from disk, a JAR, or the network) and creates a \code{java.\allowbreak{}lang.\allowbreak{}Class} object representing the type in the JVM.

\item 
\textbf{Linking}, which has three sub-steps:

\begin{itemize}
\item 
\textbf{Verify} — the bytecode verifier checks the class file is structurally valid and doesn’t violate JVM safety rules (no stack underflows, no illegal casts, etc). This is what stops hand-crafted or corrupted bytecode from crashing the JVM.

\item 
\textbf{Prepare} — static fields are allocated and set to default values (\code{0}, \code{null}, \code{false}) — \emph{not} their initializer expressions yet.

\item 
\textbf{Resolve} — symbolic references (like \code{\textquotedbl{}com/\allowbreak{}foo/\allowbreak{}Bar\textquotedbl{}}) in the constant pool are resolved to actual memory addresses/direct references. HotSpot resolves most of this lazily, on first use.

\end{itemize}

\item 
\textbf{Initialization:} Static initializer blocks and static field initializer expressions run, \textbf{in textual order, top to bottom}, exactly once per class, the first time the class is “actively used.”

\item 
\textbf{Execution:} The JVM starts interpreting (and later JIT-compiling) the bytecode of the entry point (\code{main}, or whatever invoked the class).

\end{enumerate}

\subsection[When does initialization actually trigger?]{When does initialization actually trigger?}
\label{h:11:when-does-initialization-actually-trigger}
A class is only initialized on first \textbf{active use} — not simply on loading. Active use includes:

\begin{itemize}
\item 
Creating an instance (\code{new}).

\item 
Calling a static method.

\item 
Accessing/assigning a static field (except a \code{static~\allowbreak{}final} compile-time constant).

\item 
Reflection (\code{Class.\allowbreak{}forName(\allowbreak{}name)} with default \code{initialize=\allowbreak{}true}).

\item 
Initializing a subclass (forces superclass init first).

\end{itemize}

\begin{codebox}
\begin{Verbatim}[commandchars=\\\{\}]
\PY{k+kd}{class} \PY{n+nc}{Parent}\PY{+w}{ }\PY{p}{\PYZob{}}
\PY{+w}{    }\PY{k+kd}{static}\PY{+w}{ }\PY{p}{\PYZob{}}\PY{+w}{ }\PY{n}{System}\PY{p}{.}\PY{n+na}{out}\PY{p}{.}\PY{n+na}{println}\PY{p}{(}\PY{l+s}{\PYZdq{}}\PY{l+s}{Parent init}\PY{l+s}{\PYZdq{}}\PY{p}{)}\PY{p}{;}\PY{+w}{ }\PY{p}{\PYZcb{}}
\PY{p}{\PYZcb{}}

\PY{k+kd}{class} \PY{n+nc}{Child}\PY{+w}{ }\PY{k+kd}{extends}\PY{+w}{ }\PY{n}{Parent}\PY{+w}{ }\PY{p}{\PYZob{}}
\PY{+w}{    }\PY{k+kd}{static}\PY{+w}{ }\PY{p}{\PYZob{}}\PY{+w}{ }\PY{n}{System}\PY{p}{.}\PY{n+na}{out}\PY{p}{.}\PY{n+na}{println}\PY{p}{(}\PY{l+s}{\PYZdq{}}\PY{l+s}{Child init}\PY{l+s}{\PYZdq{}}\PY{p}{)}\PY{p}{;}\PY{+w}{ }\PY{p}{\PYZcb{}}
\PY{+w}{    }\PY{k+kd}{static}\PY{+w}{ }\PY{k+kt}{int}\PY{+w}{ }\PY{n}{VALUE}\PY{+w}{ }\PY{o}{=}\PY{+w}{ }\PY{l+m+mi}{10}\PY{p}{;}
\PY{p}{\PYZcb{}}

\PY{k+kd}{public}\PY{+w}{ }\PY{k+kd}{class} \PY{n+nc}{Demo}\PY{+w}{ }\PY{p}{\PYZob{}}
\PY{+w}{    }\PY{k+kd}{public}\PY{+w}{ }\PY{k+kd}{static}\PY{+w}{ }\PY{k+kt}{void}\PY{+w}{ }\PY{n+nf}{main}\PY{p}{(}\PY{n}{String}\PY{o}{[}\PY{o}{]}\PY{+w}{ }\PY{n}{args}\PY{p}{)}\PY{+w}{ }\PY{p}{\PYZob{}}
\PY{+w}{        }\PY{n}{System}\PY{p}{.}\PY{n+na}{out}\PY{p}{.}\PY{n+na}{println}\PY{p}{(}\PY{l+s}{\PYZdq{}}\PY{l+s}{before}\PY{l+s}{\PYZdq{}}\PY{p}{)}\PY{p}{;}
\PY{+w}{        }\PY{k+kt}{int}\PY{+w}{ }\PY{n}{v}\PY{+w}{ }\PY{o}{=}\PY{+w}{ }\PY{n}{Child}\PY{p}{.}\PY{n+na}{VALUE}\PY{p}{;}\PY{+w}{   }\PY{c+c1}{// triggers Parent init, THEN Child init}
\PY{+w}{        }\PY{n}{System}\PY{p}{.}\PY{n+na}{out}\PY{p}{.}\PY{n+na}{println}\PY{p}{(}\PY{l+s}{\PYZdq{}}\PY{l+s}{after: }\PY{l+s}{\PYZdq{}}\PY{+w}{ }\PY{o}{+}\PY{+w}{ }\PY{n}{v}\PY{p}{)}\PY{p}{;}
\PY{+w}{    }\PY{p}{\PYZcb{}}
\PY{p}{\PYZcb{}}
\end{Verbatim}
\end{codebox}

Output:

\begin{codebox}
\begin{Verbatim}
before
Parent init
Child init
after: 10
\end{Verbatim}
\end{codebox}

Notice: merely referencing \code{Child} in the source doesn’t run anything. Initialization is deferred until \code{Child.\allowbreak{}VALUE} is actually touched, and superclasses always initialize before subclasses.

\begin{codebox}
\begin{Verbatim}[commandchars=\\\{\}]
\PY{c+c1}{\PYZsh{} See the class\PYZhy{}loading and initialization events for a run}
java\PY{+w}{ }\PYZhy{}Xlog:class+init\PY{o}{=}info\PY{+w}{ }Demo
\end{Verbatim}
\end{codebox}

\section[Bytecode]{Bytecode}
\label{h:11:bytecode}
Bytecode is the JVM’s instruction set — a compact, platform-independent set of opcodes that \code{javac} emits instead of native machine code. The JVM’s execution model is a \textbf{stack machine}: most instructions pop operands off an “operand stack” (part of the current stack frame) and push results back, rather than working with named CPU registers.

\subsection[A small example]{A small example}
\label{h:11:a-small-example}
\begin{codebox}
\begin{Verbatim}[commandchars=\\\{\}]
\PY{k+kd}{public}\PY{+w}{ }\PY{k+kd}{class} \PY{n+nc}{MathOps}\PY{+w}{ }\PY{p}{\PYZob{}}
\PY{+w}{    }\PY{k+kd}{public}\PY{+w}{ }\PY{k+kt}{int}\PY{+w}{ }\PY{n+nf}{add}\PY{p}{(}\PY{k+kt}{int}\PY{+w}{ }\PY{n}{a}\PY{p}{,}\PY{+w}{ }\PY{k+kt}{int}\PY{+w}{ }\PY{n}{b}\PY{p}{)}\PY{+w}{ }\PY{p}{\PYZob{}}
\PY{+w}{        }\PY{k}{return}\PY{+w}{ }\PY{n}{a}\PY{+w}{ }\PY{o}{+}\PY{+w}{ }\PY{n}{b}\PY{p}{;}
\PY{+w}{    }\PY{p}{\PYZcb{}}

\PY{+w}{    }\PY{k+kd}{public}\PY{+w}{ }\PY{k+kd}{static}\PY{+w}{ }\PY{k+kt}{int}\PY{+w}{ }\PY{n+nf}{square}\PY{p}{(}\PY{k+kt}{int}\PY{+w}{ }\PY{n}{x}\PY{p}{)}\PY{+w}{ }\PY{p}{\PYZob{}}
\PY{+w}{        }\PY{k}{return}\PY{+w}{ }\PY{n}{x}\PY{+w}{ }\PY{o}{*}\PY{+w}{ }\PY{n}{x}\PY{p}{;}
\PY{+w}{    }\PY{p}{\PYZcb{}}
\PY{p}{\PYZcb{}}
\end{Verbatim}
\end{codebox}

Compile it and disassemble with \code{javap}:

\begin{codebox}
\begin{Verbatim}[commandchars=\\\{\}]
javac\PY{+w}{ }MathOps.java
javap\PY{+w}{ }\PYZhy{}c\PY{+w}{ }MathOps.class
\end{Verbatim}
\end{codebox}

Output (trimmed to the two methods):

\begin{codebox}
\begin{Verbatim}
  public int add(int, int);
    Code:
       0: iload_1        // push local var 1 (a) onto operand stack
       1: iload_2        // push local var 2 (b) onto operand stack
       2: iadd            // pop both, push their sum
       3: ireturn          // pop and return the top of the operand stack

  public static int square(int);
    Code:
       0: iload_0        // push local var 0 (x) — note: no "this" slot for static methods
       1: iload_0
       2: imul
       3: ireturn
\end{Verbatim}
\end{codebox}

\subsection[The operand stack model]{The operand stack model}
\label{h:11:the-operand-stack-model}
Each method invocation gets its own \textbf{stack frame}, containing:

\begin{itemize}
\item 
\textbf{Local variable slots} (\code{local~\allowbreak{}var~\allowbreak{}0}, \code{1}, \code{2}, …) — parameters and local variables. For instance methods, slot 0 is always \code{this}; static methods have no such slot, which is why \code{square}’s \code{x} is slot 0 while \code{add}’s \code{a}/\code{b} are slots 1/2 (slot 0 is \code{this}).

\item 
\textbf{The operand stack} — a small scratch stack instructions push to and pop from. \code{iload\_\allowbreak{}1} pushes; \code{iadd} pops two ints and pushes their sum; \code{ireturn} pops the return value.

\end{itemize}

\subsection[Common opcodes]{Common opcodes}
\label{h:11:common-opcodes}
{\tablefontsmall
\begin{xltabular}{\linewidth}{@{}L{0.428}L{1.572}@{}}
\toprule
\rowcolor{tablehdr}
\thd{Opcode} & \thd{Meaning} \\
\midrule
\endhead
\code{iload\_\allowbreak{}n} / \code{aload\_\allowbreak{}n} & Push int / object-reference local var \code{n} onto operand stack \\
\code{istore\_\allowbreak{}n} / \code{astore\_\allowbreak{}n} & Pop top of stack into int / object-reference local var \code{n} \\
\code{iadd}, \code{imul}, \code{isub} & Pop two ints, push arithmetic result \\
\code{getfield} / \code{putfield} & Read/write an instance field \\
\code{getstatic} / \code{putstatic} & Read/write a static field \\
\code{invokevirtual} & Call an instance method using dynamic dispatch (normal polymorphic call — most common) \\
\code{invokestatic} & Call a static method (no receiver, resolved at compile time) \\
\code{invokespecial} & Call a constructor, a private method, or a superclass method via \code{super.\allowbreak{}foo()} — non-virtual, resolved statically \\
\code{invokeinterface} & Call a method through an interface reference (needs extra runtime lookup vs \code{invokevirtual}) \\
\code{invokedynamic} & Defer the call-site linking logic to a “bootstrap method” resolved at first execution — powers lambdas, method references, and string concatenation \\
\bottomrule
\end{xltabular}
}

\begin{codebox}
\begin{Verbatim}[commandchars=\\\{\}]
\PY{c+c1}{\PYZsh{} Disassemble with line numbers and the constant pool too}
javap\PY{+w}{ }\PYZhy{}c\PY{+w}{ }\PYZhy{}p\PY{+w}{ }\PYZhy{}v\PY{+w}{ }MathOps.class
\end{Verbatim}
\end{codebox}

\subsection[The constant pool]{The constant pool}
\label{h:11:the-constant-pool}
Every \code{.\allowbreak{}class} file has a \textbf{constant pool}: a table of literals, class/method/field names, and symbolic references that the bytecode indexes into instead of embedding raw strings and numbers inline. \code{javap~\allowbreak{}-\allowbreak{}v} prints it:

\begin{codebox}
\begin{Verbatim}
Constant pool:
   #1 = Methodref          #6.#15         // java/lang/Object."<init>":()V
   #2 = Class              #16            // MathOps
   #3 = Utf8               add
   ...
\end{Verbatim}
\end{codebox}

An instruction like \code{invokevirtual~\allowbreak{}\#7} doesn’t carry a class/method name directly — it carries an index into this table, which is resolved to an actual method address during linking (the “Resolve” step described earlier).

\subsection[invokedynamic: lambdas and string concatenation]{\code{invokedynamic}: lambdas and string concatenation}
\label{h:11:invokedynamic-lambdas-and-string-concatenation}
Before Java 8, every method call site was one of the four \code{invoke*} opcodes above, each resolved to a fixed target. \code{invokedynamic} (added in Java 7, exploited heavily from Java 8 onward) instead calls a \textbf{bootstrap method} the first time a call site executes; that bootstrap method decides — and caches — how the call should actually be linked. This lets \code{javac} avoid generating a synthetic inner class for every single lambda.

\begin{codebox}
\begin{Verbatim}[commandchars=\\\{\}]
\PY{n}{Runnable}\PY{+w}{ }\PY{n}{r}\PY{+w}{ }\PY{o}{=}\PY{+w}{ }\PY{p}{(}\PY{p}{)}\PY{+w}{ }\PY{o}{\PYZhy{}}\PY{o}{\PYZgt{}}\PY{+w}{ }\PY{n}{System}\PY{p}{.}\PY{n+na}{out}\PY{p}{.}\PY{n+na}{println}\PY{p}{(}\PY{l+s}{\PYZdq{}}\PY{l+s}{hi}\PY{l+s}{\PYZdq{}}\PY{p}{)}\PY{p}{;}
\end{Verbatim}
\end{codebox}

\begin{codebox}
\begin{Verbatim}[commandchars=\\\{\}]
javap\PY{+w}{ }\PYZhy{}c\PY{+w}{ }\PYZhy{}p\PY{+w}{ }Demo.class
\end{Verbatim}
\end{codebox}

\begin{codebox}
\begin{Verbatim}
  0: invokedynamic #7,  0    // InvokeDynamic #0:run:()Ljava/lang/Runnable;
  5: astore_1
\end{Verbatim}
\end{codebox}

The actual lambda body is compiled into a private synthetic method, and \code{invokedynamic} + \code{LambdaMetafactory} wire up a \code{Runnable} instance pointing at it lazily, at first call.

Since Java 9, string concatenation (\code{\textquotedbl{}a\textquotedbl{}~\allowbreak{}+\allowbreak{}~\allowbreak{}b~\allowbreak{}+\allowbreak{}~\allowbreak{}\textquotedbl{}c\textquotedbl{}}) also compiles to \code{invokedynamic} calling \code{StringConcatFactory}, instead of the old pattern of allocating a \code{StringBuilder} and chaining \code{.\allowbreak{}append()} calls:

\begin{codebox}
\begin{Verbatim}[commandchars=\\\{\}]
\PY{n}{String}\PY{+w}{ }\PY{n}{s}\PY{+w}{ }\PY{o}{=}\PY{+w}{ }\PY{l+s}{\PYZdq{}}\PY{l+s}{Count: }\PY{l+s}{\PYZdq{}}\PY{+w}{ }\PY{o}{+}\PY{+w}{ }\PY{n}{count}\PY{p}{;}
\end{Verbatim}
\end{codebox}

\begin{codebox}
\begin{Verbatim}
  invokedynamic #2,  0   // InvokeDynamic #0:makeConcatWithConstants:(I)Ljava/lang/String;
\end{Verbatim}
\end{codebox}

\begin{codebox}
\begin{Verbatim}[commandchars=\\\{\}]
\PY{c+c1}{\PYZsh{} Try it yourself}
\PY{n+nb}{echo}\PY{+w}{ }\PY{l+s+s1}{\PYZsq{}public class Cat \PYZob{} String go(int n)\PYZob{} return \PYZdq{}n=\PYZdq{} + n; \PYZcb{} \PYZcb{}\PYZsq{}}\PY{+w}{ }\PYZgt{}\PY{+w}{ }Cat.java
javac\PY{+w}{ }Cat.java\PY{+w}{ }\PY{o}{\PYZam{}\PYZam{}}\PY{+w}{ }javap\PY{+w}{ }\PYZhy{}c\PY{+w}{ }\PYZhy{}p\PY{+w}{ }Cat.java
\end{Verbatim}
\end{codebox}

\section[JIT Compilation]{JIT Compilation}
\label{h:11:jit-compilation}
Bytecode is portable but interpreting it instruction-by-instruction is slow. HotSpot (the default JVM) starts by \textbf{interpreting} bytecode, then compiles the “hot” (frequently executed) parts to native machine code on the fly — hence \emph{Just-In-Time}.

\subsection[Tiered compilation: interpreter → C1 → C2]{Tiered compilation: interpreter → C1 → C2}
\label{h:11:tiered-compilation-interpreter--c1--c2}
\begin{codebox}
\begin{Verbatim}
   bytecode
      │
      ▼
 ┌───────────┐   method called often     ┌────────────┐   still hot,
 │INTERPRETER│ ───────────────────────▶  │  C1 (client)│ ──────────────▶ ┌────────────┐
 │ (slow, no │   (thousands of calls)    │  fast compile│  long-running   │ C2 (server) │
 │ warm-up)  │                           │  light opt   │  method         │ heavy opt,  │
 └───────────┘                           └────────────┘                  │ aggressive  │
                                                                          │ inlining    │
                                                                          └────────────┘
\end{Verbatim}
\end{codebox}

\begin{itemize}
\item 
\textbf{Interpreter} — executes bytecode directly, one instruction at a time. Starts instantly, no compile delay, but slow per-iteration.

\item 
\textbf{C1 (“client compiler”)} — compiles quickly with light optimization. Good for short-lived or moderately-hot code; gets you off the interpreter fast.

\item 
\textbf{C2 (“server compiler”)} — compiles slowly but produces highly optimized native code: aggressive inlining, loop unrolling, escape analysis, etc. Reserved for genuinely hot methods.

\end{itemize}

Modern HotSpot uses \textbf{tiered compilation} by default: code moves interpreter → C1 (with several sub-tiers of profiling) → C2, escalating only if a method really is called a lot.

\subsection[Hot spots and OSR]{Hot spots and OSR}
\label{h:11:hot-spots-and-osr}
The JVM counts method invocations and \emph{loop back-edges} (loop iterations). Once a counter crosses a threshold, the method is queued for JIT compilation. If a \textbf{single long-running loop} inside a method gets hot before the method itself is recompiled, HotSpot can do \textbf{On-Stack Replacement (OSR)}: it JIT-compiles just that loop and swaps the interpreter’s execution \emph{in place}, mid-method, onto the compiled version — without waiting for the method to be called again.

\begin{codebox}
\begin{Verbatim}[commandchars=\\\{\}]
\PY{k+kt}{long}\PY{+w}{ }\PY{n+nf}{sum}\PY{p}{(}\PY{k+kt}{int}\PY{+w}{ }\PY{n}{n}\PY{p}{)}\PY{+w}{ }\PY{p}{\PYZob{}}
\PY{+w}{    }\PY{k+kt}{long}\PY{+w}{ }\PY{n}{s}\PY{+w}{ }\PY{o}{=}\PY{+w}{ }\PY{l+m+mi}{0}\PY{p}{;}
\PY{+w}{    }\PY{k}{for}\PY{+w}{ }\PY{p}{(}\PY{k+kt}{int}\PY{+w}{ }\PY{n}{i}\PY{+w}{ }\PY{o}{=}\PY{+w}{ }\PY{l+m+mi}{0}\PY{p}{;}\PY{+w}{ }\PY{n}{i}\PY{+w}{ }\PY{o}{\PYZlt{}}\PY{+w}{ }\PY{n}{n}\PY{p}{;}\PY{+w}{ }\PY{n}{i}\PY{o}{+}\PY{o}{+}\PY{p}{)}\PY{+w}{ }\PY{p}{\PYZob{}}\PY{+w}{   }\PY{c+c1}{// if n is huge, this loop alone can trigger OSR}
\PY{+w}{        }\PY{n}{s}\PY{+w}{ }\PY{o}{+}\PY{o}{=}\PY{+w}{ }\PY{n}{i}\PY{p}{;}
\PY{+w}{    }\PY{p}{\PYZcb{}}
\PY{+w}{    }\PY{k}{return}\PY{+w}{ }\PY{n}{s}\PY{p}{;}
\PY{p}{\PYZcb{}}
\end{Verbatim}
\end{codebox}

\subsection[Inlining and deoptimization]{Inlining and deoptimization}
\label{h:11:inlining-and-deoptimization}
\textbf{Inlining} replaces a call site with the callee’s body, eliminating call overhead and opening the door to further optimization (e.g., dead-code elimination once constants propagate through). C2 inlines aggressively, including through virtual calls when it can prove (via profiling) there’s effectively only one implementation in practice.

\textbf{Deoptimization} is the safety valve: if an assumption the JIT baked in turns out false at runtime (e.g., a second implementation of an interface finally shows up, invalidating a “monomorphic” inlining bet), the JVM discards the compiled code for that method, falls back to the interpreter, and re-profiles. This is invisible in correctness terms — Java’s semantics never change — but it does explain sudden latency blips in production right after a “cold” code path executes for the first time.

\begin{codebox}
\begin{Verbatim}[commandchars=\\\{\}]
\PY{k+kd}{interface} \PY{n+nc}{Shape}\PY{+w}{ }\PY{p}{\PYZob{}}\PY{+w}{ }\PY{k+kt}{double}\PY{+w}{ }\PY{n+nf}{area}\PY{p}{(}\PY{p}{)}\PY{p}{;}\PY{+w}{ }\PY{p}{\PYZcb{}}
\PY{k+kd}{class} \PY{n+nc}{Circle}\PY{+w}{ }\PY{k+kd}{implements}\PY{+w}{ }\PY{n}{Shape}\PY{+w}{ }\PY{p}{\PYZob{}}\PY{+w}{ }\PY{k+kd}{public}\PY{+w}{ }\PY{k+kt}{double}\PY{+w}{ }\PY{n+nf}{area}\PY{p}{(}\PY{p}{)}\PY{+w}{ }\PY{p}{\PYZob{}}\PY{+w}{ }\PY{k}{return}\PY{+w}{ }\PY{l+m+mf}{3.14}\PY{p}{;}\PY{+w}{ }\PY{p}{\PYZcb{}}\PY{+w}{ }\PY{p}{\PYZcb{}}
\PY{k+kd}{class} \PY{n+nc}{Square}\PY{+w}{ }\PY{k+kd}{implements}\PY{+w}{ }\PY{n}{Shape}\PY{+w}{ }\PY{p}{\PYZob{}}\PY{+w}{ }\PY{k+kd}{public}\PY{+w}{ }\PY{k+kt}{double}\PY{+w}{ }\PY{n+nf}{area}\PY{p}{(}\PY{p}{)}\PY{+w}{ }\PY{p}{\PYZob{}}\PY{+w}{ }\PY{k}{return}\PY{+w}{ }\PY{l+m+mf}{4.0}\PY{p}{;}\PY{+w}{ }\PY{p}{\PYZcb{}}\PY{+w}{ }\PY{p}{\PYZcb{}}

\PY{c+c1}{// If callers only ever pass Circle for a long time, C2 may speculate}
\PY{c+c1}{// \PYZdq{}this call site is monomorphic\PYZdq{} and inline Circle.area() directly.}
\PY{c+c1}{// The moment a Square shows up, that assumption breaks \PYZhy{}\PYZgt{} deoptimization.}
\PY{k+kt}{double}\PY{+w}{ }\PY{n+nf}{totalArea}\PY{p}{(}\PY{n}{List}\PY{o}{\PYZlt{}}\PY{n}{Shape}\PY{o}{\PYZgt{}}\PY{+w}{ }\PY{n}{shapes}\PY{p}{)}\PY{+w}{ }\PY{p}{\PYZob{}}
\PY{+w}{    }\PY{k+kt}{double}\PY{+w}{ }\PY{n}{total}\PY{+w}{ }\PY{o}{=}\PY{+w}{ }\PY{l+m+mi}{0}\PY{p}{;}
\PY{+w}{    }\PY{k}{for}\PY{+w}{ }\PY{p}{(}\PY{n}{Shape}\PY{+w}{ }\PY{n}{s}\PY{+w}{ }\PY{p}{:}\PY{+w}{ }\PY{n}{shapes}\PY{p}{)}\PY{+w}{ }\PY{n}{total}\PY{+w}{ }\PY{o}{+}\PY{o}{=}\PY{+w}{ }\PY{n}{s}\PY{p}{.}\PY{n+na}{area}\PY{p}{(}\PY{p}{)}\PY{p}{;}
\PY{+w}{    }\PY{k}{return}\PY{+w}{ }\PY{n}{total}\PY{p}{;}
\PY{p}{\PYZcb{}}
\end{Verbatim}
\end{codebox}

\subsection[Why microbenchmarks need JMH and warm-up]{Why microbenchmarks need JMH and warm-up}
\label{h:11:why-microbenchmarks-need-jmh-and-warm-up}
A naive “time it in a loop with \code{System.\allowbreak{}nanoTime()}” benchmark is almost always wrong, because:

\begin{itemize}
\item 
The \textbf{interpreter} runs the first iterations, which are much slower than the eventual steady state.

\item 
\textbf{JIT compilation happens on background threads mid-benchmark}, so timings from early iterations mix cold and hot code.

\item 
\textbf{Dead-code elimination} can let the JIT discard a whole “benchmarked” computation if its result is never used.

\end{itemize}

\href{https://github.com/openjdk/jmh}{JMH} (Java Microbenchmark Harness) solves this by running explicit \textbf{warm-up iterations} (to let JIT compilation stabilize) before measuring, and by using tricks (blackholes, forcing consumption of results) to prevent dead-code elimination from making your “fast” code fast simply because it does nothing.

\begin{codebox}
\begin{Verbatim}[commandchars=\\\{\}]
\PY{n+nd}{@BenchmarkMode}\PY{p}{(}\PY{n}{Mode}\PY{p}{.}\PY{n+na}{AverageTime}\PY{p}{)}
\PY{n+nd}{@OutputTimeUnit}\PY{p}{(}\PY{n}{TimeUnit}\PY{p}{.}\PY{n+na}{NANOSECONDS}\PY{p}{)}
\PY{n+nd}{@Warmup}\PY{p}{(}\PY{n}{iterations}\PY{+w}{ }\PY{o}{=}\PY{+w}{ }\PY{l+m+mi}{5}\PY{p}{,}\PY{+w}{ }\PY{n}{time}\PY{+w}{ }\PY{o}{=}\PY{+w}{ }\PY{l+m+mi}{1}\PY{p}{)}
\PY{n+nd}{@Measurement}\PY{p}{(}\PY{n}{iterations}\PY{+w}{ }\PY{o}{=}\PY{+w}{ }\PY{l+m+mi}{5}\PY{p}{,}\PY{+w}{ }\PY{n}{time}\PY{+w}{ }\PY{o}{=}\PY{+w}{ }\PY{l+m+mi}{1}\PY{p}{)}
\PY{n+nd}{@Fork}\PY{p}{(}\PY{l+m+mi}{1}\PY{p}{)}
\PY{k+kd}{public}\PY{+w}{ }\PY{k+kd}{class} \PY{n+nc}{StringBenchmark}\PY{+w}{ }\PY{p}{\PYZob{}}
\PY{+w}{    }\PY{n+nd}{@Benchmark}
\PY{+w}{    }\PY{k+kd}{public}\PY{+w}{ }\PY{n}{String}\PY{+w}{ }\PY{n+nf}{concat}\PY{p}{(}\PY{p}{)}\PY{+w}{ }\PY{p}{\PYZob{}}
\PY{+w}{        }\PY{k}{return}\PY{+w}{ }\PY{l+s}{\PYZdq{}}\PY{l+s}{a}\PY{l+s}{\PYZdq{}}\PY{+w}{ }\PY{o}{+}\PY{+w}{ }\PY{l+s}{\PYZdq{}}\PY{l+s}{b}\PY{l+s}{\PYZdq{}}\PY{+w}{ }\PY{o}{+}\PY{+w}{ }\PY{l+s}{\PYZdq{}}\PY{l+s}{c}\PY{l+s}{\PYZdq{}}\PY{p}{;}
\PY{+w}{    }\PY{p}{\PYZcb{}}
\PY{p}{\PYZcb{}}
\end{Verbatim}
\end{codebox}

\subsection[Observing the JIT]{Observing the JIT}
\label{h:11:observing-the-jit}
\begin{codebox}
\begin{Verbatim}[commandchars=\\\{\}]
\PY{c+c1}{\PYZsh{} Print every JIT compilation event as it happens (method, tier, time)}
java\PY{+w}{ }\PYZhy{}XX:+PrintCompilation\PY{+w}{ }MyApp

\PY{c+c1}{\PYZsh{} Typical lines:}
\PY{c+c1}{\PYZsh{}    1     3       3       java.lang.String::hashCode (60 bytes)}
\PY{c+c1}{\PYZsh{}  100    45       4       MyApp::hotLoop (120 bytes)   made not entrant   \PYZlt{}\PYZhy{} deopt!}
\end{Verbatim}
\end{codebox}

\code{made~\allowbreak{}not~\allowbreak{}entrant} in the output is the visible sign of a deoptimization — that compiled version was thrown away.

\subsection[AOT and CDS in one paragraph]{AOT and CDS in one paragraph}
\label{h:11:aot-and-cds-in-one-paragraph}
The JIT’s warm-up cost is a real problem for short-lived processes (CLIs, serverless functions, fast-restarting microservices). Two HotSpot features help: \textbf{CDS (Class Data Sharing)}, which pre-parses and serializes class metadata into a shared archive (\code{-\allowbreak{}Xshare:\allowbreak{}dump} / \code{-\allowbreak{}XX:\allowbreak{}SharedArchiveFile}) so subsequent JVM startups skip re-parsing those classes, and its evolution \textbf{AppCDS / dynamic CDS archives}, which extend this to application classes, not just JDK classes. Separately, \textbf{Project Leyden} and tools like GraalVM Native Image push further into full \textbf{AOT (Ahead-Of-Time) compilation}, compiling to native machine code before the process even starts, trading some peak-throughput JIT optimizations for near-instant startup — a trade-off worth naming in an interview when startup latency comes up.

\begin{codebox}
\begin{Verbatim}[commandchars=\\\{\}]
\PY{c+c1}{\PYZsh{} Create and use an AppCDS archive}
java\PY{+w}{ }\PYZhy{}Xshare:off\PY{+w}{ }\PYZhy{}XX:ArchiveClassesAtExit\PY{o}{=}app.jsa\PY{+w}{ }\PYZhy{}jar\PY{+w}{ }app.jar
java\PY{+w}{ }\PYZhy{}XX:SharedArchiveFile\PY{o}{=}app.jsa\PY{+w}{ }\PYZhy{}jar\PY{+w}{ }app.jar
\end{Verbatim}
\end{codebox}

\section[Class Loaders]{Class Loaders}
\label{h:11:class-loaders}
A \textbf{class loader} is the object responsible for turning a class name into a \code{Class} object — locating the bytecode and defining it into the JVM. HotSpot has a built-in hierarchy of three:

\begin{codebox}
\begin{Verbatim}
   Bootstrap ClassLoader   (native, written in C++, no Java superclass)
        loads: java.base and other core JDK modules (java.lang.*, java.util.*)
           │
           ▼ parent
   Platform ClassLoader     (was "Extension" loader before Java 9)
        loads: other JDK-provided modules not in java.base
           │
           ▼ parent
   Application ClassLoader  (a.k.a. "system" class loader)
        loads: your application's classes from the classpath / module path
\end{Verbatim}
\end{codebox}

\subsection[Parent delegation]{Parent delegation}
\label{h:11:parent-delegation}
By default, when a class loader is asked to load a class, it \textbf{first asks its parent} to try, and only attempts to load it itself if the parent fails (\code{ClassNotFoundException}). This means core JDK classes like \code{java.\allowbreak{}lang.\allowbreak{}String} are always loaded by the bootstrap loader — even if you name your own class \code{java.\allowbreak{}lang.\allowbreak{}String}, the delegation model prevents your version from shadowing the real one (and the JVM outright rejects defining classes in the \code{java.*} package from a non-bootstrap loader).

\begin{codebox}
\begin{Verbatim}[commandchars=\\\{\}]
\PY{k+kd}{public}\PY{+w}{ }\PY{k+kd}{class} \PY{n+nc}{WhichLoader}\PY{+w}{ }\PY{p}{\PYZob{}}
\PY{+w}{    }\PY{k+kd}{public}\PY{+w}{ }\PY{k+kd}{static}\PY{+w}{ }\PY{k+kt}{void}\PY{+w}{ }\PY{n+nf}{main}\PY{p}{(}\PY{n}{String}\PY{o}{[}\PY{o}{]}\PY{+w}{ }\PY{n}{args}\PY{p}{)}\PY{+w}{ }\PY{p}{\PYZob{}}
\PY{+w}{        }\PY{n}{System}\PY{p}{.}\PY{n+na}{out}\PY{p}{.}\PY{n+na}{println}\PY{p}{(}\PY{n}{String}\PY{p}{.}\PY{n+na}{class}\PY{p}{.}\PY{n+na}{getClassLoader}\PY{p}{(}\PY{p}{)}\PY{p}{)}\PY{p}{;}\PY{+w}{          }\PY{c+c1}{// null (bootstrap)}
\PY{+w}{        }\PY{n}{System}\PY{p}{.}\PY{n+na}{out}\PY{p}{.}\PY{n+na}{println}\PY{p}{(}\PY{n}{WhichLoader}\PY{p}{.}\PY{n+na}{class}\PY{p}{.}\PY{n+na}{getClassLoader}\PY{p}{(}\PY{p}{)}\PY{p}{)}\PY{p}{;}\PY{+w}{     }\PY{c+c1}{// AppClassLoader}
\PY{+w}{        }\PY{n}{System}\PY{p}{.}\PY{n+na}{out}\PY{p}{.}\PY{n+na}{println}\PY{p}{(}\PY{n}{WhichLoader}\PY{p}{.}\PY{n+na}{class}\PY{p}{.}\PY{n+na}{getClassLoader}\PY{p}{(}\PY{p}{)}\PY{p}{.}\PY{n+na}{getParent}\PY{p}{(}\PY{p}{)}\PY{p}{)}\PY{p}{;}\PY{+w}{ }\PY{c+c1}{// PlatformClassLoader}
\PY{+w}{    }\PY{p}{\PYZcb{}}
\PY{p}{\PYZcb{}}
\end{Verbatim}
\end{codebox}

\subsection[Custom class loaders]{Custom class loaders}
\label{h:11:custom-class-loaders}
Frameworks that need to load code dynamically — plugin systems, application servers deploying multiple WARs, dependency-shading/isolation tools — write custom class loaders, usually by extending \code{ClassLoader} or \code{URLClassLoader}.

\begin{codebox}
\begin{Verbatim}[commandchars=\\\{\}]
\PY{k+kd}{public}\PY{+w}{ }\PY{k+kd}{class} \PY{n+nc}{PluginClassLoader}\PY{+w}{ }\PY{k+kd}{extends}\PY{+w}{ }\PY{n}{URLClassLoader}\PY{+w}{ }\PY{p}{\PYZob{}}
\PY{+w}{    }\PY{k+kd}{public}\PY{+w}{ }\PY{n+nf}{PluginClassLoader}\PY{p}{(}\PY{n}{URL}\PY{o}{[}\PY{o}{]}\PY{+w}{ }\PY{n}{pluginJars}\PY{p}{,}\PY{+w}{ }\PY{n}{ClassLoader}\PY{+w}{ }\PY{n}{parent}\PY{p}{)}\PY{+w}{ }\PY{p}{\PYZob{}}
\PY{+w}{        }\PY{k+kd}{super}\PY{p}{(}\PY{n}{pluginJars}\PY{p}{,}\PY{+w}{ }\PY{n}{parent}\PY{p}{)}\PY{p}{;}
\PY{+w}{    }\PY{p}{\PYZcb{}}
\PY{+w}{    }\PY{c+c1}{// Override findClass()/loadClass() to change lookup order,}
\PY{+w}{    }\PY{c+c1}{// e.g. \PYZdq{}child\PYZhy{}first\PYZdq{} delegation for plugin isolation (as Tomcat\PYZsq{}s}
\PY{+w}{    }\PY{c+c1}{// webapp classloader does, breaking strict parent\PYZhy{}first order}
\PY{+w}{    }\PY{c+c1}{// deliberately, on purpose, for isolation).}
\PY{p}{\PYZcb{}}
\end{Verbatim}
\end{codebox}

\begin{codebox}
\begin{Verbatim}[commandchars=\\\{\}]
\PY{c+c1}{\PYZsh{} Load and run a class from an arbitrary jar at runtime, illustrating dynamic loading}
java\PY{+w}{ }\PYZhy{}cp\PY{+w}{ }plugin.jar\PY{+w}{ }\PYZhy{}Djava.system.class.loader\PY{o}{=}com.example.PluginClassLoader\PY{+w}{ }Main
\end{Verbatim}
\end{codebox}

\subsection[ClassNotFoundException vs NoClassDefFoundError]{\code{ClassNotFoundException} vs \code{NoClassDefFoundError}}
\label{h:11:classnotfoundexception-vs-noclassdeffounderror}
These are commonly confused in code review and interviews:

{\tablefontsmall
\begin{xltabular}{\linewidth}{@{}L{0.416}L{1.292}L{1.292}@{}}
\toprule
\rowcolor{tablehdr}
\thd{} & \thd{\code{ClassNotFoundException}} & \thd{\code{NoClassDefFoundError}} \\
\midrule
\endhead
Type & Checked \code{Exception} & \code{Error} (subclass of \code{LinkageError}) \\
When & You explicitly ask to load a class by name — \code{Class.\allowbreak{}forName(...)}, or a custom class loader’s \code{loadClass()} — and it isn’t found. & A class was \textbf{available and successfully compiled against}, but is \textbf{missing at runtime}, OR the class was found once but its \textbf{static initializer threw}, poisoning it for all future references. \\
Typical cause & Wrong string in \code{Class.\allowbreak{}forName}, missing JDBC driver, typo’d reflection lookup. & Jar removed from classpath after compiling against it; a previous load attempt failed (e.g. static init threw an exception) so the JVM marks the class permanently unusable. \\
\bottomrule
\end{xltabular}
}

\begin{codebox}
\begin{Verbatim}[commandchars=\\\{\}]
\PY{k}{try}\PY{+w}{ }\PY{p}{\PYZob{}}
\PY{+w}{    }\PY{n}{Class}\PY{p}{.}\PY{n+na}{forName}\PY{p}{(}\PY{l+s}{\PYZdq{}}\PY{l+s}{com.example.MissingDriver}\PY{l+s}{\PYZdq{}}\PY{p}{)}\PY{p}{;}\PY{+w}{ }\PY{c+c1}{// \PYZhy{}\PYZgt{} ClassNotFoundException if absent}
\PY{p}{\PYZcb{}}\PY{+w}{ }\PY{k}{catch}\PY{+w}{ }\PY{p}{(}\PY{n}{ClassNotFoundException}\PY{+w}{ }\PY{n}{e}\PY{p}{)}\PY{+w}{ }\PY{p}{\PYZob{}}
\PY{+w}{    }\PY{c+c1}{// handle: driver jar not on classpath}
\PY{p}{\PYZcb{}}
\end{Verbatim}
\end{codebox}

\begin{codebox}
\begin{Verbatim}[commandchars=\\\{\}]
\PY{k+kd}{class} \PY{n+nc}{Broken}\PY{+w}{ }\PY{p}{\PYZob{}}
\PY{+w}{    }\PY{k+kd}{static}\PY{+w}{ }\PY{k+kt}{int}\PY{+w}{ }\PY{n}{x}\PY{+w}{ }\PY{o}{=}\PY{+w}{ }\PY{l+m+mi}{1}\PY{+w}{ }\PY{o}{/}\PY{+w}{ }\PY{l+m+mi}{0}\PY{p}{;}\PY{+w}{   }\PY{c+c1}{// ArithmeticException during static init}
\PY{p}{\PYZcb{}}

\PY{k+kd}{class} \PY{n+nc}{UsesBroken}\PY{+w}{ }\PY{p}{\PYZob{}}
\PY{+w}{    }\PY{k+kt}{void}\PY{+w}{ }\PY{n+nf}{run}\PY{p}{(}\PY{p}{)}\PY{+w}{ }\PY{p}{\PYZob{}}
\PY{+w}{        }\PY{k}{new}\PY{+w}{ }\PY{n}{Broken}\PY{p}{(}\PY{p}{)}\PY{p}{;}\PY{+w}{ }\PY{c+c1}{// first call: ExceptionInInitializerError}
\PY{+w}{        }\PY{k}{new}\PY{+w}{ }\PY{n}{Broken}\PY{p}{(}\PY{p}{)}\PY{p}{;}\PY{+w}{ }\PY{c+c1}{// second call in this or another method: NoClassDefFoundError}
\PY{+w}{                       }\PY{c+c1}{// (class is now permanently \PYZdq{}erroneous\PYZdq{})}
\PY{+w}{    }\PY{p}{\PYZcb{}}
\PY{p}{\PYZcb{}}
\end{Verbatim}
\end{codebox}

\subsection[Class identity = (name, loader)]{Class identity = (name, loader)}
\label{h:11:class-identity--name-loader}
The JVM treats two classes as \textbf{the same type only if both the fully-qualified name AND the defining class loader match}. Load the “same” \code{.\allowbreak{}class} bytes through two different loaders and you get two distinct, mutually incompatible \code{Class} objects — a frequent source of baffling \code{ClassCastException}s with messages like \code{com.\allowbreak{}example.\allowbreak{}Foo~\allowbreak{}cannot~\allowbreak{}be~\allowbreak{}cast~\allowbreak{}to~\allowbreak{}com.\allowbreak{}example.\allowbreak{}Foo}.

\begin{codebox}
\begin{Verbatim}[commandchars=\\\{\}]
\PY{n}{Class}\PY{o}{\PYZlt{}}\PY{o}{?}\PY{o}{\PYZgt{}}\PY{+w}{ }\PY{n}{a}\PY{+w}{ }\PY{o}{=}\PY{+w}{ }\PY{n}{loader1}\PY{p}{.}\PY{n+na}{loadClass}\PY{p}{(}\PY{l+s}{\PYZdq{}}\PY{l+s}{com.example.Foo}\PY{l+s}{\PYZdq{}}\PY{p}{)}\PY{p}{;}
\PY{n}{Class}\PY{o}{\PYZlt{}}\PY{o}{?}\PY{o}{\PYZgt{}}\PY{+w}{ }\PY{n}{b}\PY{+w}{ }\PY{o}{=}\PY{+w}{ }\PY{n}{loader2}\PY{p}{.}\PY{n+na}{loadClass}\PY{p}{(}\PY{l+s}{\PYZdq{}}\PY{l+s}{com.example.Foo}\PY{l+s}{\PYZdq{}}\PY{p}{)}\PY{p}{;}
\PY{n}{System}\PY{p}{.}\PY{n+na}{out}\PY{p}{.}\PY{n+na}{println}\PY{p}{(}\PY{n}{a}\PY{+w}{ }\PY{o}{=}\PY{o}{=}\PY{+w}{ }\PY{n}{b}\PY{p}{)}\PY{p}{;}\PY{+w}{              }\PY{c+c1}{// false!}
\PY{n}{System}\PY{p}{.}\PY{n+na}{out}\PY{p}{.}\PY{n+na}{println}\PY{p}{(}\PY{n}{a}\PY{p}{.}\PY{n+na}{equals}\PY{p}{(}\PY{n}{b}\PY{p}{)}\PY{p}{)}\PY{p}{;}\PY{+w}{          }\PY{c+c1}{// false!}
\PY{c+c1}{// Casting an instance loaded by loader1 to the type loaded by loader2 throws}
\PY{c+c1}{// ClassCastException even though the source code is byte\PYZhy{}for\PYZhy{}byte identical.}
\end{Verbatim}
\end{codebox}

\subsection[Classloader leaks in application servers]{Classloader leaks in application servers}
\label{h:11:classloader-leaks-in-application-servers}
Servlet containers like Tomcat give \textbf{each deployed webapp its own class loader}, so redeploying a WAR without restarting the whole server means loading a fresh set of classes. If anything outside the webapp — a thread the app started, a static field in a JDK class, a \code{ThreadLocal}, JDBC driver registries — keeps a reference to a class (or an instance of a class) from the old webapp’s loader, \textbf{the entire old class loader, and every class and object it ever loaded, cannot be garbage collected}. Symptoms: \code{Metaspace} usage climbing on every redeploy, eventually \code{Out\allowbreak{}Of\allowbreak{}Memory\allowbreak{}Error:\allowbreak{}~\allowbreak{}Metaspace}. Common causes and fixes:

\begin{codebox}
\begin{Verbatim}[commandchars=\\\{\}]
\PY{c+c1}{// LEAK: a thread started by the webapp keeps running after undeploy,}
\PY{c+c1}{// its Thread object (loaded by the webapp\PYZsq{}s classloader) stays reachable}
\PY{c+c1}{// from the JVM\PYZsq{}s global thread list \PYZhy{}\PYZgt{} webapp classloader can never be freed.}
\PY{n}{Thread}\PY{+w}{ }\PY{n}{worker}\PY{+w}{ }\PY{o}{=}\PY{+w}{ }\PY{k}{new}\PY{+w}{ }\PY{n}{Thread}\PY{p}{(}\PY{k}{this}\PY{p}{:}\PY{p}{:}\PY{n}{pollForever}\PY{p}{)}\PY{p}{;}
\PY{n}{worker}\PY{p}{.}\PY{n+na}{setDaemon}\PY{p}{(}\PY{k+kc}{true}\PY{p}{)}\PY{p}{;}
\PY{n}{worker}\PY{p}{.}\PY{n+na}{start}\PY{p}{(}\PY{p}{)}\PY{p}{;}
\PY{c+c1}{// Fix: register a ServletContextListener.contextDestroyed() hook that}
\PY{c+c1}{// signals the thread to stop and joins it before the webapp is unloaded.}
\end{Verbatim}
\end{codebox}

\subsection[JPMS impact (the module system)]{JPMS impact (the module system)}
\label{h:11:jpms-impact-the-module-system}
Since Java 9, the \textbf{Java Platform Module System (JPMS)} adds a layer on top of classloaders: modules declare explicit dependencies (\code{requires}) and exposed packages (\code{exports}) in a \code{module-\allowbreak{}info.\allowbreak{}java}. This affects class loading in two practical ways reviewers should know: (1) \textbf{strong encapsulation} — internal packages that aren’t \code{exported} are inaccessible via reflection from outside the module by default, even if they were \code{public}, breaking some old reflection-heavy frameworks; (2) the platform’s own bootstrap/platform/application loaders were re-organized around modules (e.g. \code{java.\allowbreak{}base} is the module the bootstrap loader owns), and each named module effectively gets a layer that still resolves through the same parent-delegation-style rules for unnamed/automatic modules on the classpath.

\begin{codebox}
\begin{Verbatim}[commandchars=\\\{\}]
\PY{c+c1}{// module\PYZhy{}info.java}
\PY{n}{module}\PY{+w}{ }\PY{n}{com}\PY{p}{.}\PY{n+na}{example}\PY{p}{.}\PY{n+na}{app}\PY{+w}{ }\PY{p}{\PYZob{}}
\PY{+w}{    }\PY{n}{requires}\PY{+w}{ }\PY{n}{java}\PY{p}{.}\PY{n+na}{sql}\PY{p}{;}
\PY{+w}{    }\PY{n}{exports}\PY{+w}{ }\PY{n}{com}\PY{p}{.}\PY{n+na}{example}\PY{p}{.}\PY{n+na}{app}\PY{p}{.}\PY{n+na}{api}\PY{p}{;}\PY{+w}{   }\PY{c+c1}{// only this package is visible to other modules}
\PY{+w}{    }\PY{c+c1}{// com.example.app.internal is NOT exported \PYZhy{}\PYZgt{} inaccessible outside the module,}
\PY{+w}{    }\PY{c+c1}{// even via setAccessible(true) reflection, unless \PYZdq{}opens\PYZdq{} is declared.}
\PY{p}{\PYZcb{}}
\end{Verbatim}
\end{codebox}

\section[Java Memory Model (JMM)]{Java Memory Model (JMM)}
\label{h:11:java-memory-model-jmm}
The Java Memory Model is the specification (JLS Chapter 17) that answers a deceptively hard question: \textbf{when a thread writes a value, when is another thread guaranteed to see it?} Without a memory model, compilers and CPUs are free to reorder, cache, and optimize your code in ways that are safe for a single thread but disastrous once multiple threads are involved.

\subsection[Reordering by compiler and CPU]{Reordering by compiler and CPU}
\label{h:11:reordering-by-compiler-and-cpu}
Both the JIT compiler and the CPU are allowed to reorder instructions and cache values in registers, as long as a \textbf{single thread} can’t observe the difference. But another thread, watching from the outside, absolutely can.

\begin{codebox}
\begin{Verbatim}[commandchars=\\\{\}]
\PY{k+kd}{class} \PY{n+nc}{Reorder}\PY{+w}{ }\PY{p}{\PYZob{}}
\PY{+w}{    }\PY{k+kt}{int}\PY{+w}{ }\PY{n}{a}\PY{+w}{ }\PY{o}{=}\PY{+w}{ }\PY{l+m+mi}{0}\PY{p}{;}
\PY{+w}{    }\PY{k+kt}{boolean}\PY{+w}{ }\PY{n}{flag}\PY{+w}{ }\PY{o}{=}\PY{+w}{ }\PY{k+kc}{false}\PY{p}{;}

\PY{+w}{    }\PY{k+kt}{void}\PY{+w}{ }\PY{n+nf}{writer}\PY{p}{(}\PY{p}{)}\PY{+w}{ }\PY{p}{\PYZob{}}
\PY{+w}{        }\PY{n}{a}\PY{+w}{ }\PY{o}{=}\PY{+w}{ }\PY{l+m+mi}{42}\PY{p}{;}\PY{+w}{        }\PY{c+c1}{// (1)}
\PY{+w}{        }\PY{n}{flag}\PY{+w}{ }\PY{o}{=}\PY{+w}{ }\PY{k+kc}{true}\PY{p}{;}\PY{+w}{   }\PY{c+c1}{// (2) — compiler/CPU may reorder (1) and (2)!}
\PY{+w}{    }\PY{p}{\PYZcb{}}

\PY{+w}{    }\PY{k+kt}{void}\PY{+w}{ }\PY{n+nf}{reader}\PY{p}{(}\PY{p}{)}\PY{+w}{ }\PY{p}{\PYZob{}}
\PY{+w}{        }\PY{k}{if}\PY{+w}{ }\PY{p}{(}\PY{n}{flag}\PY{p}{)}\PY{+w}{ }\PY{p}{\PYZob{}}
\PY{+w}{            }\PY{n}{System}\PY{p}{.}\PY{n+na}{out}\PY{p}{.}\PY{n+na}{println}\PY{p}{(}\PY{n}{a}\PY{p}{)}\PY{p}{;}\PY{+w}{ }\PY{c+c1}{// could print 0, not 42, without synchronization}
\PY{+w}{        }\PY{p}{\PYZcb{}}
\PY{+w}{    }\PY{p}{\PYZcb{}}
\PY{p}{\PYZcb{}}
\end{Verbatim}
\end{codebox}

Each CPU core also has its own \textbf{cache}, separate from main memory. A write by one thread might sit in that core’s cache/store-buffer without being flushed to main memory (or the shared cache level other cores see) for an unbounded amount of time, from the other thread’s perspective, unless something forces synchronization.

\subsection[Happens-before]{Happens-before}
\label{h:11:happens-before}
The JMM defines correctness in terms of a \textbf{happens-before} relationship: if action X happens-before action Y, every effect of X (every write) is guaranteed visible to Y. Without an established happens-before edge between two threads’ actions, there is \textbf{no guarantee} the second thread ever sees the first thread’s writes — not “might be slow,” but genuinely unspecified/broken behavior.

Key sources of happens-before edges:

\begin{itemize}
\item 
A \code{synchronized} block’s unlock happens-before a later thread’s lock on the \emph{same monitor}.

\item 
A \code{volatile} field write happens-before every subsequent \code{volatile} read of that same field.

\item 
\code{Thread.\allowbreak{}start()} happens-before anything the started thread does.

\item 
Everything a thread does happens-before another thread observes it has finished, via \code{Thread.\allowbreak{}join()}.

\item 
Writes to \code{final} fields performed in a constructor happen-before any thread that gets a correctly-published reference to that object.

\end{itemize}

\subsection[volatile]{\code{volatile}}
\label{h:11:volatile}
\code{volatile} gives you \textbf{visibility} (every read sees the latest write, no stale caching) and prevents certain compiler/CPU reorderings around that field, but it does \textbf{not} give you atomicity or mutual exclusion for compound operations.

\begin{codebox}
\begin{Verbatim}[commandchars=\\\{\}]
\PY{k+kd}{class} \PY{n+nc}{FlagExample}\PY{+w}{ }\PY{p}{\PYZob{}}
\PY{+w}{    }\PY{k+kd}{private}\PY{+w}{ }\PY{k+kd}{volatile}\PY{+w}{ }\PY{k+kt}{boolean}\PY{+w}{ }\PY{n}{running}\PY{+w}{ }\PY{o}{=}\PY{+w}{ }\PY{k+kc}{true}\PY{p}{;}

\PY{+w}{    }\PY{k+kt}{void}\PY{+w}{ }\PY{n+nf}{stop}\PY{p}{(}\PY{p}{)}\PY{+w}{ }\PY{p}{\PYZob{}}\PY{+w}{ }\PY{n}{running}\PY{+w}{ }\PY{o}{=}\PY{+w}{ }\PY{k+kc}{false}\PY{p}{;}\PY{+w}{ }\PY{p}{\PYZcb{}}\PY{+w}{               }\PY{c+c1}{// thread A}

\PY{+w}{    }\PY{k+kt}{void}\PY{+w}{ }\PY{n+nf}{loop}\PY{p}{(}\PY{p}{)}\PY{+w}{ }\PY{p}{\PYZob{}}
\PY{+w}{        }\PY{k}{while}\PY{+w}{ }\PY{p}{(}\PY{n}{running}\PY{p}{)}\PY{+w}{ }\PY{p}{\PYZob{}}\PY{+w}{                            }\PY{c+c1}{// thread B}
\PY{+w}{            }\PY{c+c1}{// do work}
\PY{+w}{        }\PY{p}{\PYZcb{}}
\PY{+w}{    }\PY{p}{\PYZcb{}}
\PY{p}{\PYZcb{}}
\end{Verbatim}
\end{codebox}

Without \code{volatile} here, the JIT is legally allowed to hoist \code{running} into a register in \code{loop()} (since single-threaded reasoning says it never changes inside the loop body), producing an \textbf{infinite loop that never sees \code{stop()}'s write} — this is the concrete broken-visibility bug interviewers love to ask about.

\begin{codebox}
\begin{Verbatim}[commandchars=\\\{\}]
\PY{c+c1}{// BROKEN: no volatile, no synchronization \PYZhy{}\PYZgt{} may spin forever on some JVMs/JITs}
\PY{k+kd}{class} \PY{n+nc}{BrokenFlag}\PY{+w}{ }\PY{p}{\PYZob{}}
\PY{+w}{    }\PY{k+kd}{private}\PY{+w}{ }\PY{k+kt}{boolean}\PY{+w}{ }\PY{n}{running}\PY{+w}{ }\PY{o}{=}\PY{+w}{ }\PY{k+kc}{true}\PY{p}{;}
\PY{+w}{    }\PY{k+kt}{void}\PY{+w}{ }\PY{n+nf}{stop}\PY{p}{(}\PY{p}{)}\PY{+w}{ }\PY{p}{\PYZob{}}\PY{+w}{ }\PY{n}{running}\PY{+w}{ }\PY{o}{=}\PY{+w}{ }\PY{k+kc}{false}\PY{p}{;}\PY{+w}{ }\PY{p}{\PYZcb{}}
\PY{+w}{    }\PY{k+kt}{void}\PY{+w}{ }\PY{n+nf}{loop}\PY{p}{(}\PY{p}{)}\PY{+w}{ }\PY{p}{\PYZob{}}\PY{+w}{ }\PY{k}{while}\PY{+w}{ }\PY{p}{(}\PY{n}{running}\PY{p}{)}\PY{+w}{ }\PY{p}{\PYZob{}}\PY{+w}{ }\PY{c+cm}{/* spin */}\PY{+w}{ }\PY{p}{\PYZcb{}}\PY{+w}{ }\PY{p}{\PYZcb{}}
\PY{p}{\PYZcb{}}
\end{Verbatim}
\end{codebox}

\code{volatile} does \textbf{not} make increments atomic:

\begin{codebox}
\begin{Verbatim}[commandchars=\\\{\}]
\PY{k+kd}{private}\PY{+w}{ }\PY{k+kd}{volatile}\PY{+w}{ }\PY{k+kt}{int}\PY{+w}{ }\PY{n}{counter}\PY{+w}{ }\PY{o}{=}\PY{+w}{ }\PY{l+m+mi}{0}\PY{p}{;}
\PY{k+kt}{void}\PY{+w}{ }\PY{n+nf}{increment}\PY{p}{(}\PY{p}{)}\PY{+w}{ }\PY{p}{\PYZob{}}\PY{+w}{ }\PY{n}{counter}\PY{o}{+}\PY{o}{+}\PY{p}{;}\PY{+w}{ }\PY{p}{\PYZcb{}}\PY{+w}{  }\PY{c+c1}{// STILL a race: read, add, write are 3 separate steps}
\PY{c+c1}{// Fix: use AtomicInteger, or synchronize the whole read\PYZhy{}modify\PYZhy{}write.}
\end{Verbatim}
\end{codebox}

\subsection[final field freeze semantics]{\code{final} field freeze semantics}
\label{h:11:final-field-freeze-semantics}
If an object is constructed with a \code{final} field, and the reference to that object is not leaked out of the constructor before it finishes, then \textbf{any thread that later obtains a reference to that object is guaranteed to see the fully-initialized value of the \code{final} field} — without needing any additional synchronization. This is the JMM’s “safe construction via final fields” guarantee.

\begin{codebox}
\begin{Verbatim}[commandchars=\\\{\}]
\PY{k+kd}{final}\PY{+w}{ }\PY{k+kd}{class} \PY{n+nc}{Point}\PY{+w}{ }\PY{p}{\PYZob{}}
\PY{+w}{    }\PY{k+kd}{final}\PY{+w}{ }\PY{k+kt}{int}\PY{+w}{ }\PY{n}{x}\PY{p}{,}\PY{+w}{ }\PY{n}{y}\PY{p}{;}
\PY{+w}{    }\PY{n}{Point}\PY{p}{(}\PY{k+kt}{int}\PY{+w}{ }\PY{n}{x}\PY{p}{,}\PY{+w}{ }\PY{k+kt}{int}\PY{+w}{ }\PY{n}{y}\PY{p}{)}\PY{+w}{ }\PY{p}{\PYZob{}}\PY{+w}{ }\PY{k}{this}\PY{p}{.}\PY{n+na}{x}\PY{+w}{ }\PY{o}{=}\PY{+w}{ }\PY{n}{x}\PY{p}{;}\PY{+w}{ }\PY{k}{this}\PY{p}{.}\PY{n+na}{y}\PY{+w}{ }\PY{o}{=}\PY{+w}{ }\PY{n}{y}\PY{p}{;}\PY{+w}{ }\PY{p}{\PYZcb{}}
\PY{p}{\PYZcb{}}
\PY{c+c1}{// Any thread that receives a Point reference through a properly published}
\PY{c+c1}{// channel (see below) will see correct, fully\PYZhy{}initialized x and y —}
\PY{c+c1}{// guaranteed by final field \PYZdq{}freeze\PYZdq{} semantics, even with zero synchronization}
\PY{c+c1}{// on the reading side.}
\end{Verbatim}
\end{codebox}

\subsection[Safe publication]{Safe publication}
\label{h:11:safe-publication}
“Safe publication” means making an object visible to other threads in a way the JMM actually guarantees is safe — not just “it usually works.” Common safe publication mechanisms:

\begin{itemize}
\item 
Assigning the reference to a \code{volatile} or \code{static~\allowbreak{}final} field.

\item 
Storing it into a properly locked field (\code{synchronized}).

\item 
Putting it into a thread-safe collection like \code{ConcurrentHashMap} or \code{BlockingQueue}, whose own internal synchronization carries the happens-before edge.

\item 
Initializing it as a \code{static} field during class initialization (class-init has its own locking guarantees).

\end{itemize}

\begin{codebox}
\begin{Verbatim}[commandchars=\\\{\}]
\PY{c+c1}{// UNSAFE publication: another thread might see a half\PYZhy{}constructed object}
\PY{k+kd}{public}\PY{+w}{ }\PY{k+kd}{class} \PY{n+nc}{Holder}\PY{+w}{ }\PY{p}{\PYZob{}}
\PY{+w}{    }\PY{k+kd}{public}\PY{+w}{ }\PY{k+kd}{static}\PY{+w}{ }\PY{n}{Helper}\PY{+w}{ }\PY{n}{helper}\PY{p}{;}
\PY{+w}{    }\PY{k+kd}{public}\PY{+w}{ }\PY{k+kd}{static}\PY{+w}{ }\PY{k+kt}{void}\PY{+w}{ }\PY{n+nf}{init}\PY{p}{(}\PY{p}{)}\PY{+w}{ }\PY{p}{\PYZob{}}
\PY{+w}{        }\PY{n}{helper}\PY{+w}{ }\PY{o}{=}\PY{+w}{ }\PY{k}{new}\PY{+w}{ }\PY{n}{Helper}\PY{p}{(}\PY{p}{)}\PY{p}{;}\PY{+w}{ }\PY{c+c1}{// another thread reading `helper` non\PYZhy{}null}
\PY{+w}{                                 }\PY{c+c1}{// isn\PYZsq{}t guaranteed to see a fully built Helper}
\PY{+w}{                                 }\PY{c+c1}{// unless Helper\PYZsq{}s fields are all final, or}
\PY{+w}{                                 }\PY{c+c1}{// helper is volatile/static\PYZhy{}final.}
\PY{+w}{    }\PY{p}{\PYZcb{}}
\PY{p}{\PYZcb{}}
\end{Verbatim}
\end{codebox}

\begin{codebox}
\begin{Verbatim}[commandchars=\\\{\}]
\PY{c+c1}{// SAFE publication via volatile}
\PY{k+kd}{public}\PY{+w}{ }\PY{k+kd}{class} \PY{n+nc}{Holder}\PY{+w}{ }\PY{p}{\PYZob{}}
\PY{+w}{    }\PY{k+kd}{private}\PY{+w}{ }\PY{k+kd}{static}\PY{+w}{ }\PY{k+kd}{volatile}\PY{+w}{ }\PY{n}{Helper}\PY{+w}{ }\PY{n}{helper}\PY{p}{;}
\PY{+w}{    }\PY{k+kd}{public}\PY{+w}{ }\PY{k+kd}{static}\PY{+w}{ }\PY{n}{Helper}\PY{+w}{ }\PY{n+nf}{getHelper}\PY{p}{(}\PY{p}{)}\PY{+w}{ }\PY{p}{\PYZob{}}
\PY{+w}{        }\PY{k}{if}\PY{+w}{ }\PY{p}{(}\PY{n}{helper}\PY{+w}{ }\PY{o}{=}\PY{o}{=}\PY{+w}{ }\PY{k+kc}{null}\PY{p}{)}\PY{+w}{ }\PY{p}{\PYZob{}}
\PY{+w}{            }\PY{k+kd}{synchronized}\PY{+w}{ }\PY{p}{(}\PY{n}{Holder}\PY{p}{.}\PY{n+na}{class}\PY{p}{)}\PY{+w}{ }\PY{p}{\PYZob{}}
\PY{+w}{                }\PY{k}{if}\PY{+w}{ }\PY{p}{(}\PY{n}{helper}\PY{+w}{ }\PY{o}{=}\PY{o}{=}\PY{+w}{ }\PY{k+kc}{null}\PY{p}{)}\PY{+w}{ }\PY{n}{helper}\PY{+w}{ }\PY{o}{=}\PY{+w}{ }\PY{k}{new}\PY{+w}{ }\PY{n}{Helper}\PY{p}{(}\PY{p}{)}\PY{p}{;}
\PY{+w}{            }\PY{p}{\PYZcb{}}
\PY{+w}{        }\PY{p}{\PYZcb{}}
\PY{+w}{        }\PY{k}{return}\PY{+w}{ }\PY{n}{helper}\PY{p}{;}
\PY{+w}{    }\PY{p}{\PYZcb{}}
\PY{p}{\PYZcb{}}
\end{Verbatim}
\end{codebox}

\subsection[Data races vs race conditions]{Data races vs race conditions}
\label{h:11:data-races-vs-race-conditions}
These terms are related but not identical, and mixing them up in review comments is a common tell:

{\tablefontsmall
\begin{xltabular}{\linewidth}{@{}L{0.410}L{1.295}L{1.295}@{}}
\toprule
\rowcolor{tablehdr}
\thd{} & \thd{Data race} & \thd{Race condition} \\
\midrule
\endhead
Definition & Two threads access the same memory location, at least one write, with no happens-before ordering between them. & Program’s correctness depends on timing/interleaving of operations, regardless of memory visibility. \\
Scope & A specific JMM-defined term about memory access without synchronization. & A general concurrency term about outcome depending on execution order. \\
Example & Non-volatile \code{boolean~\allowbreak{}running} read/written by two threads. & Two threads both doing “check-then-act” on a bank balance with a lock each holds separately — well-synchronized memory-wise but still logically racy (TOCTOU). \\
\bottomrule
\end{xltabular}
}

A program can have a race condition without a data race (properly synchronized but logically wrong ordering), and a data race almost always implies undefined/unreliable behavior even if it “seems to work” in testing.

\subsection[synchronized: two jobs, not one]{\code{synchronized}: two jobs, not one}
\label{h:11:synchronized-two-jobs-not-one}
\code{synchronized} is often described only as “locking,” but it actually does \textbf{two separate jobs} at once, and code reviewers should check both are needed (or that neither is skippable):

\begin{enumerate}
\item 
\textbf{Mutual exclusion} — only one thread can hold a given monitor at a time, preventing concurrent execution of the guarded block.

\item 
\textbf{Visibility} — entering a \code{synchronized} block establishes a happens-before edge with the \emph{previous} thread that held (and released) the same lock, guaranteeing you see all of its writes made inside that lock, not just the fact that it finished.

\end{enumerate}

\begin{codebox}
\begin{Verbatim}[commandchars=\\\{\}]
\PY{k+kd}{class} \PY{n+nc}{Counter}\PY{+w}{ }\PY{p}{\PYZob{}}
\PY{+w}{    }\PY{k+kd}{private}\PY{+w}{ }\PY{k+kt}{int}\PY{+w}{ }\PY{n}{count}\PY{+w}{ }\PY{o}{=}\PY{+w}{ }\PY{l+m+mi}{0}\PY{p}{;}
\PY{+w}{    }\PY{k+kd}{private}\PY{+w}{ }\PY{k+kd}{final}\PY{+w}{ }\PY{n}{Object}\PY{+w}{ }\PY{n}{lock}\PY{+w}{ }\PY{o}{=}\PY{+w}{ }\PY{k}{new}\PY{+w}{ }\PY{n}{Object}\PY{p}{(}\PY{p}{)}\PY{p}{;}

\PY{+w}{    }\PY{k+kt}{void}\PY{+w}{ }\PY{n+nf}{increment}\PY{p}{(}\PY{p}{)}\PY{+w}{ }\PY{p}{\PYZob{}}
\PY{+w}{        }\PY{k+kd}{synchronized}\PY{+w}{ }\PY{p}{(}\PY{n}{lock}\PY{p}{)}\PY{+w}{ }\PY{p}{\PYZob{}}
\PY{+w}{            }\PY{n}{count}\PY{o}{+}\PY{o}{+}\PY{p}{;}\PY{+w}{              }\PY{c+c1}{// mutual exclusion: no lost updates}
\PY{+w}{        }\PY{p}{\PYZcb{}}
\PY{+w}{    }\PY{p}{\PYZcb{}}

\PY{+w}{    }\PY{k+kt}{int}\PY{+w}{ }\PY{n+nf}{get}\PY{p}{(}\PY{p}{)}\PY{+w}{ }\PY{p}{\PYZob{}}
\PY{+w}{        }\PY{k+kd}{synchronized}\PY{+w}{ }\PY{p}{(}\PY{n}{lock}\PY{p}{)}\PY{+w}{ }\PY{p}{\PYZob{}}\PY{+w}{     }\PY{c+c1}{// visibility: guaranteed to see latest count,}
\PY{+w}{            }\PY{k}{return}\PY{+w}{ }\PY{n}{count}\PY{p}{;}\PY{+w}{         }\PY{c+c1}{// even without this being volatile}
\PY{+w}{        }\PY{p}{\PYZcb{}}
\PY{+w}{    }\PY{p}{\PYZcb{}}
\PY{p}{\PYZcb{}}
\end{Verbatim}
\end{codebox}

A frequent review mistake: locking on \code{increment()} but reading \code{count} directly without synchronization elsewhere in \code{get()}. That reintroduces the visibility problem even though updates themselves are safely serialized — the lock’s visibility guarantee only applies to threads that \emph{also} go through the same lock.

\begin{codebox}
\begin{Verbatim}[commandchars=\\\{\}]
\PY{c+c1}{\PYZsh{} See lock contention / thread states live, useful when reviewing concurrency claims}
jcmd\PY{+w}{ }\PYZlt{}pid\PYZgt{}\PY{+w}{ }Thread.print\PY{+w}{ }\PY{p}{|}\PY{+w}{ }grep\PY{+w}{ }\PYZhy{}A3\PY{+w}{ }\PY{l+s+s2}{\PYZdq{}waiting to lock\PYZdq{}}
\end{Verbatim}
\end{codebox}

\section[Common Code-Review Interview Pitfalls]{Common Code-Review Interview Pitfalls}
\label{h:11:common-code-review-interview-pitfalls}
\begin{enumerate}
\item 
\textbf{Assuming the JRE still exists as a separate download.} Since Java 11, there is no standalone JRE artifact; reviewers sometimes still ask candidates to “install the JRE” — the correct modern answer is a full JDK, or a \code{jlink}-built custom runtime image.

\begin{codebox}
\begin{Verbatim}[commandchars=\\\{\}]
jlink\PY{+w}{ }\PYZhy{}\PYZhy{}add\PYZhy{}modules\PY{+w}{ }java.base\PY{+w}{ }\PYZhy{}\PYZhy{}output\PY{+w}{ }tiny\PYZhy{}runtime
\end{Verbatim}
\end{codebox}

\item 
\textbf{Confusing heap \code{OutOfMemoryError} with metaspace/native exhaustion.} \code{-\allowbreak{}Xmx} only bounds the heap; unbounded class generation (common with proxy-heavy frameworks, or dynamically generated classes never unloaded) exhausts Metaspace instead, which needs \code{-\allowbreak{}XX:\allowbreak{}MaxMetaspaceSize} to even surface as a clean error instead of eating all system memory.

\begin{codebox}
\begin{Verbatim}[commandchars=\\\{\}]
java\PY{+w}{ }\PYZhy{}Xmx512m\PY{+w}{ }\PYZhy{}XX:MaxMetaspaceSize\PY{o}{=}256m\PY{+w}{ }\PYZhy{}jar\PY{+w}{ }app.jar
\end{Verbatim}
\end{codebox}

\item 
\textbf{Believing \code{static} blocks run at class-load time.} They run at class \textbf{initialization}, which is deferred to first active use — code that assumes a static block has “already run” just because a class was referenced (e.g. in a type declaration or import) is wrong.

\begin{codebox}
\begin{Verbatim}[commandchars=\\\{\}]
\PY{k+kd}{static}\PY{+w}{ }\PY{p}{\PYZob{}}\PY{+w}{ }\PY{n}{System}\PY{p}{.}\PY{n+na}{out}\PY{p}{.}\PY{n+na}{println}\PY{p}{(}\PY{l+s}{\PYZdq{}}\PY{l+s}{I might run much later than you think}\PY{l+s}{\PYZdq{}}\PY{p}{)}\PY{p}{;}\PY{+w}{ }\PY{p}{\PYZcb{}}
\end{Verbatim}
\end{codebox}

\item 
\textbf{Treating \code{invokevirtual} and \code{invokestatic} as interchangeable in reasoning about polymorphism.} A private or static “helper” method resolved via \code{invokestatic}/\code{invokespecial} cannot be overridden — reviewers should flag any comment claiming a private method “overrides” a superclass method; it just shadows/hides it.

\begin{codebox}
\begin{Verbatim}[commandchars=\\\{\}]
\PY{k+kd}{class} \PY{n+nc}{Base}\PY{+w}{ }\PY{p}{\PYZob{}}\PY{+w}{ }\PY{k+kd}{private}\PY{+w}{ }\PY{k+kt}{void}\PY{+w}{ }\PY{n+nf}{helper}\PY{p}{(}\PY{p}{)}\PY{+w}{ }\PY{p}{\PYZob{}}\PY{p}{\PYZcb{}}\PY{+w}{ }\PY{p}{\PYZcb{}}\PY{+w}{      }\PY{c+c1}{// invokespecial, not virtual dispatch}
\PY{k+kd}{class} \PY{n+nc}{Sub}\PY{+w}{ }\PY{k+kd}{extends}\PY{+w}{ }\PY{n}{Base}\PY{+w}{ }\PY{p}{\PYZob{}}\PY{+w}{ }\PY{k+kd}{private}\PY{+w}{ }\PY{k+kt}{void}\PY{+w}{ }\PY{n+nf}{helper}\PY{p}{(}\PY{p}{)}\PY{+w}{ }\PY{p}{\PYZob{}}\PY{p}{\PYZcb{}}\PY{+w}{ }\PY{p}{\PYZcb{}}\PY{+w}{ }\PY{c+c1}{// unrelated method, NOT an override}
\end{Verbatim}
\end{codebox}

\item 
\textbf{Expecting \code{+} string concatenation to always allocate a \code{StringBuilder}.} Since Java 9, javac lowers concatenation to \code{invokedynamic} + \code{StringConcatFactory}; reviewers shouldn’t ding code for “missing an explicit \code{StringBuilder}” in simple concatenation — the compiler already optimizes it, though a raw loop with \code{+=} inside still has quadratic-cost risk and \emph{should} use an explicit \code{StringBuilder}.

\begin{codebox}
\begin{Verbatim}[commandchars=\\\{\}]
\PY{n}{String}\PY{+w}{ }\PY{n}{s}\PY{+w}{ }\PY{o}{=}\PY{+w}{ }\PY{l+s}{\PYZdq{}}\PY{l+s}{\PYZdq{}}\PY{p}{;}
\PY{k}{for}\PY{+w}{ }\PY{p}{(}\PY{n}{String}\PY{+w}{ }\PY{n}{part}\PY{+w}{ }\PY{p}{:}\PY{+w}{ }\PY{n}{parts}\PY{p}{)}\PY{+w}{ }\PY{n}{s}\PY{+w}{ }\PY{o}{+}\PY{o}{=}\PY{+w}{ }\PY{n}{part}\PY{p}{;}\PY{+w}{  }\PY{c+c1}{// still O(n\PYZca{}2) — use StringBuilder here}
\end{Verbatim}
\end{codebox}

\item 
\textbf{Trusting raw loop-based microbenchmarks in a PR’s “performance proof.”} Without JIT warm-up and dead-code-elimination guards, a \code{System.\allowbreak{}nanoTime()} loop can report numbers dominated by interpreter overhead or optimized-away code. Ask for JMH results instead.

\begin{codebox}
\begin{Verbatim}[commandchars=\\\{\}]
\PY{n+nd}{@Warmup}\PY{p}{(}\PY{n}{iterations}\PY{+w}{ }\PY{o}{=}\PY{+w}{ }\PY{l+m+mi}{5}\PY{p}{)}\PY{+w}{ }\PY{n+nd}{@Measurement}\PY{p}{(}\PY{n}{iterations}\PY{+w}{ }\PY{o}{=}\PY{+w}{ }\PY{l+m+mi}{5}\PY{p}{)}
\PY{k+kd}{public}\PY{+w}{ }\PY{k+kd}{class} \PY{n+nc}{MyBenchmark}\PY{+w}{ }\PY{p}{\PYZob{}}\PY{+w}{ }\PY{n+nd}{@Benchmark}\PY{+w}{ }\PY{k+kd}{public}\PY{+w}{ }\PY{k+kt}{int}\PY{+w}{ }\PY{n+nf}{compute}\PY{p}{(}\PY{p}{)}\PY{+w}{ }\PY{p}{\PYZob{}}\PY{+w}{ }\PY{k}{return}\PY{+w}{ }\PY{n}{heavy}\PY{p}{(}\PY{p}{)}\PY{p}{;}\PY{+w}{ }\PY{p}{\PYZcb{}}\PY{+w}{ }\PY{p}{\PYZcb{}}
\end{Verbatim}
\end{codebox}

\item 
\textbf{Assuming a hot method is always running JIT-compiled code.} A method can look identical in two profiling runs yet perform very differently if one run deoptimized it (e.g. a newly-loaded subclass broke a monomorphic call-site assumption). Look for \code{made~\allowbreak{}not~\allowbreak{}entrant} in \code{-\allowbreak{}XX:+\allowbreak{}PrintCompilation} output before blaming “the algorithm.”

\item 
\textbf{Mixing up \code{ClassNotFoundException} and \code{NoClassDefFoundError} in exception-handling review.} Catching \code{ClassNotFoundException} around \code{Class.\allowbreak{}forName} won’t help if the real problem is a static initializer that already threw once — that manifests later as \code{NoClassDefFoundError}, which is an \code{Error}, not caught by typical \code{Exception} handlers.

\begin{codebox}
\begin{Verbatim}[commandchars=\\\{\}]
\PY{k}{try}\PY{+w}{ }\PY{p}{\PYZob{}}\PY{+w}{ }\PY{n}{Class}\PY{p}{.}\PY{n+na}{forName}\PY{p}{(}\PY{l+s}{\PYZdq{}}\PY{l+s}{x.Y}\PY{l+s}{\PYZdq{}}\PY{p}{)}\PY{p}{;}\PY{+w}{ }\PY{p}{\PYZcb{}}
\PY{k}{catch}\PY{+w}{ }\PY{p}{(}\PY{n}{ClassNotFoundException}\PY{+w}{ }\PY{n}{e}\PY{p}{)}\PY{+w}{ }\PY{p}{\PYZob{}}\PY{+w}{ }\PY{c+cm}{/* won\PYZsq{}t catch a poisoned class! */}\PY{+w}{ }\PY{p}{\PYZcb{}}
\end{Verbatim}
\end{codebox}

\item 
\textbf{Comparing classes loaded by different class loaders with \code{==} or \code{instanceof} and being surprised by failures.} In app-server or plugin architectures, “the same” class loaded twice by different loaders are different types to the JVM. A \code{ClassCastException} with identical-looking type names on both sides is the signature symptom — check for multiple loaders before assuming a bug in the cast logic itself.

\item 
\textbf{Overlooking classloader leaks from long-lived threads or \code{ThreadLocal}s in app-server deployments.} A background thread or \code{ThreadLocal} value that outlives a hot-redeploy pins the entire old webapp class loader in Metaspace. Reviewers should ask: “who stops this thread / clears this ThreadLocal on shutdown?”

\end{enumerate}

\begin{codebox}
\begin{Verbatim}[commandchars=\\\{\}]
\PY{n}{contextDestroyed}\PY{p}{(}\PY{n}{ServletContextEvent}\PY{+w}{ }\PY{n}{e}\PY{p}{)}\PY{+w}{ }\PY{p}{\PYZob{}}\PY{+w}{ }\PY{n}{worker}\PY{p}{.}\PY{n+na}{interrupt}\PY{p}{(}\PY{p}{)}\PY{p}{;}\PY{+w}{ }\PY{n}{threadLocal}\PY{p}{.}\PY{n+na}{remove}\PY{p}{(}\PY{p}{)}\PY{p}{;}\PY{+w}{ }\PY{p}{\PYZcb{}}
\end{Verbatim}
\end{codebox}

\begin{enumerate}[start=11]
\item 
\textbf{Reading or writing a shared flag/field across threads with no \code{volatile}, lock, or other happens-before edge.} This is the textbook broken-visibility bug — it can compile, pass single-threaded tests, and still spin forever or read stale values in production, because nothing in the JMM obligates the writer’s update to ever become visible to the reader.

\begin{codebox}
\begin{Verbatim}[commandchars=\\\{\}]
\PY{k+kt}{boolean}\PY{+w}{ }\PY{n}{running}\PY{+w}{ }\PY{o}{=}\PY{+w}{ }\PY{k+kc}{true}\PY{p}{;}\PY{+w}{       }\PY{c+c1}{// missing volatile}
\PY{k+kt}{void}\PY{+w}{ }\PY{n+nf}{loop}\PY{p}{(}\PY{p}{)}\PY{+w}{ }\PY{p}{\PYZob{}}\PY{+w}{ }\PY{k}{while}\PY{+w}{ }\PY{p}{(}\PY{n}{running}\PY{p}{)}\PY{+w}{ }\PY{p}{\PYZob{}}\PY{p}{\PYZcb{}}\PY{+w}{ }\PY{p}{\PYZcb{}}\PY{+w}{   }\PY{c+c1}{// may never see stop()\PYZsq{}s write}
\end{Verbatim}
\end{codebox}

\item 
\textbf{Believing \code{volatile} makes compound operations atomic.} \code{volatile~\allowbreak{}int~\allowbreak{}counter;\allowbreak{}~\allowbreak{}counter++;} is still a three-step read-modify-write race between threads — visibility is guaranteed, atomicity is not. Look for \code{AtomicInteger}/\code{AtomicLong} or a lock instead.

\item 
\textbf{Publishing a partially-constructed object reference without a \code{final}/\code{volatile}/lock-based safe-publication path.} The classic double-checked-locking-without-volatile singleton bug: another thread can see a non-null but not-yet-fully-initialized instance.

\begin{codebox}
\begin{Verbatim}[commandchars=\\\{\}]
\PY{k+kd}{private}\PY{+w}{ }\PY{k+kd}{static}\PY{+w}{ }\PY{n}{Helper}\PY{+w}{ }\PY{n}{instance}\PY{p}{;}\PY{+w}{ }\PY{c+c1}{// missing volatile — double\PYZhy{}checked locking is broken here}
\end{Verbatim}
\end{codebox}

\item 
\textbf{Conflating “synchronized gives mutual exclusion” with “synchronized on this block automatically makes everything else visible too.”} Only accesses that go through the \emph{same} lock get the happens-before guarantee. A \code{synchronized} setter paired with an unsynchronized getter on the same field is a subtle, easy-to-miss review flag — it looks safe but isn’t.

\item 
\textbf{Calling something a “data race” when it’s really a logic-level “race condition” (or vice-versa) in review comments.} A well-synchronized (no data race) check-then-act sequence can still be logically racy (TOCTOU) if the check and act aren’t combined atomically — that’s a race condition, not a JMM data race, and the fix (atomic compound operations / higher-level locking) is different from “just add \code{volatile}.”

\begin{codebox}
\begin{Verbatim}[commandchars=\\\{\}]
\PY{k}{if}\PY{+w}{ }\PY{p}{(}\PY{o}{!}\PY{n}{map}\PY{p}{.}\PY{n+na}{containsKey}\PY{p}{(}\PY{n}{k}\PY{p}{)}\PY{p}{)}\PY{+w}{ }\PY{n}{map}\PY{p}{.}\PY{n+na}{put}\PY{p}{(}\PY{n}{k}\PY{p}{,}\PY{+w}{ }\PY{n}{v}\PY{p}{)}\PY{p}{;}\PY{+w}{ }\PY{c+c1}{// race condition even if map is a ConcurrentHashMap}
\PY{c+c1}{// Fix: map.putIfAbsent(k, v);}
\end{Verbatim}
\end{codebox}

\end{enumerate}


\setcounter{chapter}{11}
\chapter{Memory Management}
\label{chap:12}
\label{h:12:12-memory-management}
Memory management is where “the code compiles” and “the code survives production” part ways. A code reviewer who understands garbage collection, common leak patterns, and the diagnostic tools can catch problems that unit tests never will — a slowly growing heap, a rogue \code{ThreadLocal}, a GC pause that violates an SLA. This chapter covers how the JVM’s garbage collectors work, how to spot and debug memory leaks, what escape analysis actually does (and doesn’t) for you, and the real commands you’ll reach for when production memory goes wrong. Examples target Java 21+ HotSpot, since that is the current LTS baseline for interviews and production.

\section[Garbage Collection]{Garbage Collection}
\label{h:12:garbage-collection}
\textbf{Garbage collection (GC)} is the JVM automatically finding objects nobody can reach anymore and freeing their memory. You never call \code{free()} like in C — the collector does it for you. That’s a huge productivity win, but it means understanding \emph{how} it decides what’s garbage is part of writing (and reviewing) correct, fast Java.

\textbf{Reachability and GC roots.} An object is “alive” if it’s \emph{reachable} — if you can follow a chain of references to it starting from a \textbf{GC root}. GC roots are the starting points the collector trusts by definition: local variables and parameters on any thread’s stack, active static fields, JNI references, and a few JVM-internal references (like classes loaded by the bootstrap classloader). If an object can’t be reached by following references from any GC root, it’s garbage — even if it still has a valid memory address and even if two garbage objects point at each other (a reference cycle does \textbf{not} keep objects alive in Java, unlike naive reference counting).

\begin{codebox}
\begin{Verbatim}[commandchars=\\\{\}]
\PY{k+kd}{class} \PY{n+nc}{Node}\PY{+w}{ }\PY{p}{\PYZob{}}
\PY{+w}{    }\PY{n}{Node}\PY{+w}{ }\PY{n}{next}\PY{p}{;}
\PY{p}{\PYZcb{}}

\PY{n}{Node}\PY{+w}{ }\PY{n}{a}\PY{+w}{ }\PY{o}{=}\PY{+w}{ }\PY{k}{new}\PY{+w}{ }\PY{n}{Node}\PY{p}{(}\PY{p}{)}\PY{p}{;}\PY{+w}{      }\PY{c+c1}{// \PYZsq{}a\PYZsq{} is a GC root (local variable) \PYZhy{}\PYZgt{} object A is reachable}
\PY{n}{Node}\PY{+w}{ }\PY{n}{b}\PY{+w}{ }\PY{o}{=}\PY{+w}{ }\PY{k}{new}\PY{+w}{ }\PY{n}{Node}\PY{p}{(}\PY{p}{)}\PY{p}{;}
\PY{n}{a}\PY{p}{.}\PY{n+na}{next}\PY{+w}{ }\PY{o}{=}\PY{+w}{ }\PY{n}{b}\PY{p}{;}\PY{+w}{                }\PY{c+c1}{// B is reachable through A}
\PY{n}{a}\PY{+w}{ }\PY{o}{=}\PY{+w}{ }\PY{k+kc}{null}\PY{p}{;}\PY{+w}{                  }\PY{c+c1}{// A is no longer reachable from any root...}
\PY{c+c1}{// ...so A AND B are both garbage now, even though B still exists in memory,}
\PY{c+c1}{// because nothing reachable from a GC root points to A anymore.}
\end{Verbatim}
\end{codebox}

\textbf{Mark-sweep-compact.} This is the classic three-phase algorithm most collectors build on:

\begin{enumerate}
\item 
\textbf{Mark} — walk from GC roots and flag every reachable object.

\item 
\textbf{Sweep} — reclaim the memory used by everything unmarked.

\item 
\textbf{Compact} — slide the surviving objects together to remove gaps, so the heap has one large free region instead of fragmented holes.

\end{enumerate}

Compaction matters because fragmented free memory can make a large allocation fail even though the \emph{total} free space is enough — like having plenty of parking but no single spot big enough for a bus.

\textbf{Copying collectors.} Instead of sweeping in place, a copying collector moves all live objects from one memory region (“from-space”) to another (“to-space”) and then treats the entire from-space as free. This is very fast when most objects are garbage (which is typical for young objects) because it only touches \emph{live} objects, not dead ones — cost is proportional to what survives, not to heap size.

\textbf{The generational hypothesis.} Empirically, most objects die young — a request-scoped object, a loop temporary, a \code{String} built for logging. Few objects live a long time (caches, singletons, connection pools). Collectors exploit this by splitting the heap into generations:

\begin{itemize}
\item 
\textbf{Young generation} — where new objects are born. Collected frequently and cheaply with a copying collector. A young-gen collection is called a \textbf{minor GC}.

\item 
\textbf{Old (tenured) generation} — where objects that survive several young collections get promoted (“tenured”). Collected less often but more expensively, since it’s usually larger and mark-sweep-compact style. Collecting the old generation (or the whole heap) is a \textbf{major GC}; collecting the \emph{entire} heap including metaspace is often called a \textbf{full GC}. Full GCs are the ones most likely to cause a noticeable pause.

\end{itemize}

\begin{codebox}
\begin{Verbatim}
Young Gen                              Old Gen
+---------+---------+---------+       +---------------------+
|  Eden   | Surv. 0 | Surv. 1 |  -->  |   Tenured objects   |
+---------+---------+---------+       +---------------------+
new objects here    ping-pong          long-lived objects promoted here
\end{Verbatim}
\end{codebox}

New objects are allocated in \textbf{Eden}. When Eden fills up, a minor GC copies survivors into a \textbf{survivor space}, ping-ponging between two survivor spaces on each collection, and objects that survive enough cycles get promoted to the old generation.

\textbf{TLABs (Thread-Local Allocation Buffers).} To avoid every thread contending on a lock just to bump a heap pointer, each thread gets its own small private chunk of Eden (a TLAB) to allocate into without synchronization. Allocation becomes “bump a pointer in my own buffer” — extremely fast — until the TLAB is full, at which point the thread grabs a new one. This is why object allocation in Java is often faster than a naive C \code{malloc}.

\textbf{Stop-the-world (STW) pauses.} Most collectors need to pause all application threads at some point — for example, to safely walk thread stacks for GC roots without them changing underneath the collector. This pause is called \textbf{stop-the-world}. The whole story of modern collector design (G1, ZGC, Shenandoah) is about shrinking STW pauses from “proportional to heap size” down to a few milliseconds, regardless of how big the heap is.

\textbf{Minor, major, and full GC — a quick disambiguation.} Interviewers like this because the terms are used loosely in the wild:

{\tablefont
\begin{xltabular}{\linewidth}{@{}L{0.425}L{1.425}L{0.622}L{1.528}@{}}
\toprule
\rowcolor{tablehdr}
\thd{Term} & \thd{What it collects} & \thd{Typical cost} & \thd{Triggered by} \\
\midrule
\endhead
\textbf{Minor GC} & Young generation only (Eden + survivor spaces) & Cheap, frequent & Eden fills up \\
\textbf{Major GC} & Old generation (definitions vary slightly by collector) & Expensive, less frequent & Old gen fills up / promotion pressure \\
\textbf{Full GC} & Entire heap, and usually metaspace too & Most expensive, rarest & Old gen exhausted, explicit \code{System.\allowbreak{}gc()}, metaspace pressure \\
\bottomrule
\end{xltabular}
}

In G1, the line between “major” and “full” is blurrier than in Parallel GC, because G1 does most old-gen work incrementally and concurrently — a true G1 full GC (falling back to single-threaded mark-sweep-compact of the whole heap) is meant to be a rare emergency fallback, and seeing one repeatedly in a log is itself a red flag that the collector is under more pressure than it’s designed to handle.

\textbf{Allocation rate and promotion rate.} Two numbers worth knowing by name because tools report them directly: \textbf{allocation rate} is how many bytes per second the application creates (mostly in Eden); \textbf{promotion rate} is how many bytes per second survive long enough to be pushed into the old generation. A high allocation rate with a \emph{low} promotion rate is healthy (objects are dying young, exactly as the generational hypothesis predicts). A high promotion rate is a warning sign — it means minor GCs are working harder to fill up the old generation faster, dragging major/full GCs closer together.

\textbf{Weak, soft, and phantom references.} Normal (“strong”) references keep an object alive. Java also has three weaker kinds, all in \code{java.\allowbreak{}lang.\allowbreak{}ref}:

{\tablefontsmall
\begin{xltabular}{\linewidth}{@{}L{0.473}L{1.636}L{0.891}@{}}
\toprule
\rowcolor{tablehdr}
\thd{Reference type} & \thd{Collected when} & \thd{Typical use} \\
\midrule
\endhead
\textbf{Strong} (normal) & Never, while reachable & Everyday references \\
\textbf{Soft} (\code{SoftReference}) & GC \emph{may} clear it, typically only under memory pressure & Memory-sensitive caches \\
\textbf{Weak} (\code{WeakReference}) & GC clears it as soon as nothing else strongly reaches it & \code{WeakHashMap} keys, avoiding leaks in listener maps \\
\textbf{Phantom} (\code{PhantomReference}) & After the object is finalized/unreachable, but before memory is reclaimed; \code{get()} always returns \code{null} & Precise cleanup via \code{Cleaner} \\
\bottomrule
\end{xltabular}
}

\begin{codebox}
\begin{Verbatim}[commandchars=\\\{\}]
\PY{n}{Map}\PY{o}{\PYZlt{}}\PY{n}{Key}\PY{p}{,}\PY{+w}{ }\PY{n}{Value}\PY{o}{\PYZgt{}}\PY{+w}{ }\PY{n}{cache}\PY{+w}{ }\PY{o}{=}\PY{+w}{ }\PY{k}{new}\PY{+w}{ }\PY{n}{java}\PY{p}{.}\PY{n+na}{util}\PY{p}{.}\PY{n+na}{WeakHashMap}\PY{o}{\PYZlt{}}\PY{o}{\PYZgt{}}\PY{p}{(}\PY{p}{)}\PY{p}{;}
\PY{c+c1}{// entries disappear automatically once \PYZsq{}Key\PYZsq{} has no other strong references —}
\PY{c+c1}{// useful for caches keyed by objects you don\PYZsq{}t own the lifecycle of}
\end{Verbatim}
\end{codebox}

\textbf{Why \code{System.\allowbreak{}gc()} is wrong.} Calling \code{System.\allowbreak{}gc()} only \emph{suggests} a full GC to the JVM — it does not force one, and on some collectors it’s an expensive full GC that pauses the whole application. Application code should never call it: it’s the collector’s job to decide when collection is worthwhile, and it usually knows better than a hardcoded call buried in business logic. (A code reviewer should flag any \code{System.\allowbreak{}gc()} call outside of a benchmark harness or manual profiling session.)

\textbf{Why \code{finalize()} is wrong.} \code{Object.\allowbreak{}finalize()} was Java’s original hook to run cleanup code before an object is reclaimed. It’s a bad tool: it runs on an unpredictable finalizer thread at an unpredictable time (or never, if the JVM exits first), it can resurrect objects, it slows GC down, and a throwing finalizer silently swallows the exception. \textbf{\code{finalize()} has been deprecated for removal since Java 9} — don’t use it in new code.

\begin{codebox}
\begin{Verbatim}[commandchars=\\\{\}]
\PY{c+c1}{// Wrong — deprecated for removal, unpredictable timing}
\PY{n+nd}{@Override}
\PY{k+kd}{protected}\PY{+w}{ }\PY{k+kt}{void}\PY{+w}{ }\PY{n+nf}{finalize}\PY{p}{(}\PY{p}{)}\PY{+w}{ }\PY{p}{\PYZob{}}
\PY{+w}{    }\PY{n}{connection}\PY{p}{.}\PY{n+na}{close}\PY{p}{(}\PY{p}{)}\PY{p}{;}
\PY{p}{\PYZcb{}}

\PY{c+c1}{// Right — deterministic, explicit}
\PY{k}{try}\PY{+w}{ }\PY{p}{(}\PY{k+kd}{var}\PY{+w}{ }\PY{n}{connection}\PY{+w}{ }\PY{o}{=}\PY{+w}{ }\PY{n}{openConnection}\PY{p}{(}\PY{p}{)}\PY{p}{)}\PY{+w}{ }\PY{p}{\PYZob{}}
\PY{+w}{    }\PY{c+c1}{// use connection}
\PY{p}{\PYZcb{}}\PY{+w}{ }\PY{c+c1}{// .close() called here, guaranteed}

\PY{c+c1}{// If you truly need a GC\PYZhy{}triggered safety net (not the primary cleanup path):}
\PY{k+kn}{import}\PY{+w}{ }\PY{n+nn}{java.lang.ref.Cleaner}\PY{p}{;}

\PY{k+kd}{class} \PY{n+nc}{ResourceHolder}\PY{+w}{ }\PY{k+kd}{implements}\PY{+w}{ }\PY{n}{AutoCloseable}\PY{+w}{ }\PY{p}{\PYZob{}}
\PY{+w}{    }\PY{k+kd}{private}\PY{+w}{ }\PY{k+kd}{static}\PY{+w}{ }\PY{k+kd}{final}\PY{+w}{ }\PY{n}{Cleaner}\PY{+w}{ }\PY{n}{CLEANER}\PY{+w}{ }\PY{o}{=}\PY{+w}{ }\PY{n}{Cleaner}\PY{p}{.}\PY{n+na}{create}\PY{p}{(}\PY{p}{)}\PY{p}{;}
\PY{+w}{    }\PY{k+kd}{private}\PY{+w}{ }\PY{k+kd}{final}\PY{+w}{ }\PY{n}{Cleaner}\PY{p}{.}\PY{n+na}{Cleanable}\PY{+w}{ }\PY{n}{cleanable}\PY{p}{;}
\PY{+w}{    }\PY{k+kd}{private}\PY{+w}{ }\PY{k+kd}{final}\PY{+w}{ }\PY{n}{State}\PY{+w}{ }\PY{n}{state}\PY{p}{;}

\PY{+w}{    }\PY{k+kd}{private}\PY{+w}{ }\PY{k+kd}{static}\PY{+w}{ }\PY{k+kd}{class} \PY{n+nc}{State}\PY{+w}{ }\PY{k+kd}{implements}\PY{+w}{ }\PY{n}{Runnable}\PY{+w}{ }\PY{p}{\PYZob{}}
\PY{+w}{        }\PY{n}{Connection}\PY{+w}{ }\PY{n}{connection}\PY{p}{;}
\PY{+w}{        }\PY{n+nd}{@Override}\PY{+w}{ }\PY{k+kd}{public}\PY{+w}{ }\PY{k+kt}{void}\PY{+w}{ }\PY{n+nf}{run}\PY{p}{(}\PY{p}{)}\PY{+w}{ }\PY{p}{\PYZob{}}\PY{+w}{ }\PY{n}{connection}\PY{p}{.}\PY{n+na}{close}\PY{p}{(}\PY{p}{)}\PY{p}{;}\PY{+w}{ }\PY{p}{\PYZcb{}}\PY{+w}{ }\PY{c+c1}{// no reference back to ResourceHolder!}
\PY{+w}{    }\PY{p}{\PYZcb{}}

\PY{+w}{    }\PY{n}{ResourceHolder}\PY{p}{(}\PY{n}{Connection}\PY{+w}{ }\PY{n}{c}\PY{p}{)}\PY{+w}{ }\PY{p}{\PYZob{}}
\PY{+w}{        }\PY{n}{state}\PY{+w}{ }\PY{o}{=}\PY{+w}{ }\PY{k}{new}\PY{+w}{ }\PY{n}{State}\PY{p}{(}\PY{p}{)}\PY{p}{;}
\PY{+w}{        }\PY{n}{state}\PY{p}{.}\PY{n+na}{connection}\PY{+w}{ }\PY{o}{=}\PY{+w}{ }\PY{n}{c}\PY{p}{;}
\PY{+w}{        }\PY{n}{cleanable}\PY{+w}{ }\PY{o}{=}\PY{+w}{ }\PY{n}{CLEANER}\PY{p}{.}\PY{n+na}{register}\PY{p}{(}\PY{k}{this}\PY{p}{,}\PY{+w}{ }\PY{n}{state}\PY{p}{)}\PY{p}{;}
\PY{+w}{    }\PY{p}{\PYZcb{}}

\PY{+w}{    }\PY{n+nd}{@Override}\PY{+w}{ }\PY{k+kd}{public}\PY{+w}{ }\PY{k+kt}{void}\PY{+w}{ }\PY{n+nf}{close}\PY{p}{(}\PY{p}{)}\PY{+w}{ }\PY{p}{\PYZob{}}\PY{+w}{ }\PY{n}{cleanable}\PY{p}{.}\PY{n+na}{clean}\PY{p}{(}\PY{p}{)}\PY{p}{;}\PY{+w}{ }\PY{p}{\PYZcb{}}
\PY{p}{\PYZcb{}}
\end{Verbatim}
\end{codebox}

\code{Cleaner} is the modern replacement — it still shouldn’t be your \emph{primary} cleanup strategy (\code{try}-with-resources / explicit \code{close()} should be), but it’s a safe backstop, unlike \code{finalize()}, because its cleanup action must not hold a reference back to the object being cleaned, avoiding the resurrection problem.

\section[Serial GC]{Serial GC}
\label{h:12:serial-gc}
\textbf{How it works.} The simplest collector. One single thread does all the work — marking, sweeping, compacting — while every application thread is stopped. No parallelism, no concurrency.

\textbf{Pause characteristics.} Pauses are proportional to heap size and can be long, but there’s zero coordination overhead between GC threads, so for a \emph{small} heap the pause is small too.

\textbf{Heap sweet spot.} Small heaps — typically under a few hundred MB. Think embedded devices, small CLI tools, containers with 1-2 CPUs and constrained memory.

\textbf{Flag to enable:} \code{-\allowbreak{}XX:+\allowbreak{}UseSerialGC}

\textbf{When to pick it.} Single-core or memory-constrained environments where the simplicity and low overhead outweigh the fact that pauses aren’t parallelized. Rarely the right choice for a server with multiple cores and a heap above a few hundred MB.

Example log line — notice the pause scales with the \emph{whole} young generation size, with no “(Normal)”/“(Concurrent Start)” qualifiers because there’s only one kind of pause:

\begin{codebox}
\begin{Verbatim}
[gc] GC(5) Pause Young (Allocation Failure) 96M->12M(128M) 45.112ms
\end{Verbatim}
\end{codebox}

\section[Parallel GC]{Parallel GC}
\label{h:12:parallel-gc}
\textbf{How it works.} Same generational, mark-sweep-compact/copying design as Serial GC, but both the young and old generation collections use \textbf{multiple threads} to do the marking, copying, and compacting work. The application is still fully stopped during a collection (it’s still stop-the-world) — the parallelism just makes that pause shorter by spreading the work across cores.

\textbf{Pause characteristics.} Optimizes for \textbf{throughput} (maximum total work done by the application over time), not for pause length. Pauses can still be noticeable, especially full GCs on a large old generation, but total CPU spent on GC (versus useful work) is low.

\textbf{Heap sweet spot.} Medium to large heaps on multi-core machines, for batch jobs and throughput-sensitive workloads that don’t have a strict pause-time SLA.

\textbf{Flag to enable:} \code{-\allowbreak{}XX:+\allowbreak{}UseParallelGC}

\textbf{When to pick it.} Batch processing, offline data pipelines, scientific computing — anything where total job completion time matters more than any individual pause, and there’s no interactive user waiting on a response.

Example log line — same shape as Serial’s, but the wall-clock duration is shorter for a similarly sized young generation because multiple GC threads did the copying in parallel:

\begin{codebox}
\begin{Verbatim}
[gc] GC(5) Pause Young (Allocation Failure) 96M->10M(128M) 12.847ms
\end{Verbatim}
\end{codebox}

\section[G1 GC]{G1 GC}
\label{h:12:g1-gc}
\textbf{How it works.} \textbf{G1 (Garbage-First)} divides the heap into many equally sized \textbf{regions} (typically 1-32 MB each) rather than into two contiguous young/old blocks. Each region can play the role of Eden, survivor, or old space, and this assignment can change over time. G1 tracks how much garbage is in each region and, true to its name, collects the regions with the \emph{most garbage first} — maximizing reclaimed memory per unit of pause time. Most of its work (finding live objects, some copying) happens concurrently with the application; only short phases are stop-the-world.

\textbf{Pause characteristics.} Designed around a \textbf{pause-time goal} you set, not a fixed algorithm — G1 tries to keep pauses under a target (default 200 ms) by choosing how many regions to collect in each cycle. It trades a bit of throughput for much more predictable pauses than Parallel GC.

\textbf{Heap sweet spot.} Medium to large heaps (multi-GB, up into hundreds of GB) on multi-core machines — this is the general-purpose “good enough for almost everything” choice.

\textbf{Flag to enable:} \code{-\allowbreak{}XX:+\allowbreak{}UseG1GC} (this is also the \textbf{default collector since JDK 9}, so usually no flag is needed at all)

\textbf{When to pick it.} The default choice for most server applications: web services, microservices, anything wanting a balance of throughput and bounded pauses without special tuning. Tune the target with \code{-\allowbreak{}XX:\allowbreak{}Max\allowbreak{}GCPause\allowbreak{}Millis=\textless{}\allowbreak{}ms\textgreater{}}.

\textbf{Humongous objects.} Any object at least half the size of a single G1 region (region size defaults based on heap size, typically 1-32 MB) is classified as \textbf{humongous} and allocated directly into one or more dedicated regions, bypassing the normal young-gen path entirely. Humongous objects are never copied during young collections (too expensive to move), so they can only be reclaimed by an old-gen/full collection — a service that allocates lots of large arrays or buffers (e.g. big \code{byte[]} responses, large \code{ArrayList} backing arrays) can suffer humongous-allocation-driven full GCs even though the heap “looks” mostly empty. Symptom in the log: frequent \code{Pause~\allowbreak{}Young~\allowbreak{}(\allowbreak{}Concurrent~\allowbreak{}Start)\allowbreak{}~\allowbreak{}(\allowbreak{}G1~\allowbreak{}Humongous~\allowbreak{}Allocation)} entries. Fix by increasing \code{-\allowbreak{}XX:\allowbreak{}G1HeapRegionSize} so fewer objects cross the humongous threshold, or by reducing how often the code allocates very large single objects (e.g. streaming/chunking instead of buffering the whole payload).

Example log line — G1 pauses are qualified with \emph{why} the pause happened:

\begin{codebox}
\begin{Verbatim}
[gc] GC(41) Pause Young (Concurrent Start) (G1 Evacuation Pause) 612M->140M(2048M) 9.774ms
[gc] GC(42) Pause Remark 2.011ms
[gc] GC(43) Pause Cleanup 0.415ms
\end{Verbatim}
\end{codebox}

\code{Concurrent~\allowbreak{}Start} kicks off a background marking cycle for the old generation; \code{Remark} and \code{Cleanup} are the short stop-the-world phases that bookend that concurrent marking work.

\section[ZGC]{ZGC}
\label{h:12:zgc}
\textbf{How it works.} \textbf{ZGC (Z Garbage Collector)} is designed for extremely low pause times, largely independent of heap size. It does almost all of its work — marking, relocating objects, remapping references — concurrently with the running application, using \textbf{colored pointers} (metadata bits embedded in the reference itself) and \textbf{load barriers} (a small check the JVM inserts on reference reads) to let it move objects while the application keeps running, safely redirecting any thread that reads a stale reference. Since JDK 21, ZGC has a \textbf{generational mode} that adds young/old generations (like G1) for much better throughput, because it can collect young garbage cheaply instead of scanning the whole heap every cycle.

\textbf{Pause characteristics.} Sub-millisecond stop-the-world pauses, and — this is the headline feature — pause time does \textbf{not} scale with heap size. A 10 GB heap and a 10 TB heap have similarly tiny pauses.

\textbf{Heap sweet spot.} Large to very large heaps (multi-GB to multi-TB) where minimizing worst-case pause latency matters more than raw throughput or memory footprint (ZGC uses more memory and CPU overhead for its barriers and metadata than G1).

\textbf{Flag to enable:} \code{-\allowbreak{}XX:+\allowbreak{}UseZGC} (generational mode is the default \emph{shape} of ZGC as of JDK 21+; explicitly request it with \code{-\allowbreak{}XX:+\allowbreak{}Use\allowbreak{}ZGC~\allowbreak{}-\allowbreak{}XX:+\allowbreak{}ZGenerational} on JDK 21, and it’s simply the only mode from JDK 23 onward)

\textbf{When to pick it.} Latency-critical services with strict pause-time SLAs (trading systems, real-time bidding, anything where a 500 ms pause is a user-visible incident) and large heaps. \textbf{Non-generational ZGC is deprecated for removal starting JDK 23} — generational ZGC is the production-ready, actively developed mode going forward.

Example log line — note the tiny STW pauses; almost everything else in a ZGC cycle is logged as \code{Concurrent}, not \code{Pause}:

\begin{codebox}
\begin{Verbatim}
[gc] GC(88) Pause Mark Start 0.012ms
[gc] GC(88) Pause Mark End 0.031ms
[gc] GC(88) Pause Relocate Start 0.028ms
[gc] GC(88) Garbage Collection (Allocation Rate) 8192M(50%)->6144M(37%)
\end{Verbatim}
\end{codebox}

\section[Shenandoah]{Shenandoah}
\label{h:12:shenandoah}
\textbf{How it works.} Shenandoah (developed by Red Hat) has a similar goal to ZGC — pause times independent of heap size — but a different mechanism: it uses \textbf{Brooks pointers} (an extra forwarding pointer stored with each object) and concurrent evacuation to move objects while the application runs, along with concurrent marking and reference updating. Conceptually it’s a sibling approach to ZGC solving the same problem with different bookkeeping.

\textbf{Pause characteristics.} Very low, largely heap-size-independent pauses, comparable in spirit to ZGC, though the two have different throughput/footprint trade-offs depending on workload — benchmark both if latency is critical.

\textbf{Heap sweet spot.} Same territory as ZGC: large heaps, latency-sensitive workloads.

\textbf{Flag to enable:} \code{-\allowbreak{}XX:+\allowbreak{}UseShenandoahGC}

\textbf{When to pick it.} Available in most OpenJDK builds (Red Hat, Eclipse Temurin/Adoptium, Amazon Corretto) but notably \textbf{not shipped in Oracle’s own JDK builds}. Pick it when you’re already on a distribution that includes it and want ZGC-like pause behavior, or want to compare it against ZGC for your specific workload.

\section[Collector Comparison]{Collector Comparison}
\label{h:12:collector-comparison}
{\tablefont
\begin{xltabular}{\linewidth}{@{}L{0.446}L{1.443}L{0.918}L{1.224}L{0.787}L{1.181}@{}}
\toprule
\rowcolor{tablehdr}
\thd{Collector} & \thd{Pause goal} & \thd{Throughput} & \thd{Heap size sweet spot} & \thd{Enable flag} & \thd{Default since} \\
\midrule
\endhead
\textbf{Serial} & Long, but simple & Low (single-threaded) & Small (\textless{} few hundred MB) & \code{-\allowbreak{}XX:+\allowbreak{}UseSerialGC} & Never default on server VMs \\
\textbf{Parallel} & Long, throughput-first & Highest & Medium-large, batch jobs & \code{-\allowbreak{}XX:+\allowbreak{}UseParallelGC} & Default in JDK 8 \\
\textbf{G1} & Configurable target (default 200 ms) & High & Medium-large, general purpose & \code{-\allowbreak{}XX:+\allowbreak{}UseG1GC} & \textbf{Default since JDK 9} \\
\textbf{ZGC} & Sub-millisecond, size-independent & Good (generational mode) & Large-huge, latency-critical & \code{-\allowbreak{}XX:+\allowbreak{}UseZGC} & Not default; opt-in \\
\textbf{Shenandoah} & Sub-millisecond, size-independent & Good & Large, latency-critical & \code{-\allowbreak{}XX:+\allowbreak{}UseShenandoahGC} & Not default; opt-in, OpenJDK-only \\
\bottomrule
\end{xltabular}
}

Rule of thumb for an interview: \textbf{G1 unless you have a specific reason not to.} Reach for ZGC or Shenandoah when pause time is the dominant concern and you have the memory/CPU budget to spend on it; reach for Parallel for pure batch throughput; reach for Serial only in tiny, constrained environments.

\section[Memory Leaks]{Memory Leaks}
\label{h:12:memory-leaks}
A Java “memory leak” doesn’t mean memory vanishes — it means objects stay \textbf{reachable} long after they’re logically dead, so the GC can never reclaim them. The heap slowly grows until an \code{OutOfMemoryError}. Here are the classic causes.

\textbf{Static collections that only grow.}

\begin{codebox}
\begin{Verbatim}[commandchars=\\\{\}]
\PY{k+kd}{public}\PY{+w}{ }\PY{k+kd}{class} \PY{n+nc}{SessionRegistry}\PY{+w}{ }\PY{p}{\PYZob{}}
\PY{+w}{    }\PY{c+c1}{// Wrong: nothing ever removes entries, so this grows forever}
\PY{+w}{    }\PY{k+kd}{private}\PY{+w}{ }\PY{k+kd}{static}\PY{+w}{ }\PY{k+kd}{final}\PY{+w}{ }\PY{n}{Map}\PY{o}{\PYZlt{}}\PY{n}{String}\PY{p}{,}\PY{+w}{ }\PY{n}{Session}\PY{o}{\PYZgt{}}\PY{+w}{ }\PY{n}{SESSIONS}\PY{+w}{ }\PY{o}{=}\PY{+w}{ }\PY{k}{new}\PY{+w}{ }\PY{n}{HashMap}\PY{o}{\PYZlt{}}\PY{o}{\PYZgt{}}\PY{p}{(}\PY{p}{)}\PY{p}{;}

\PY{+w}{    }\PY{k+kd}{public}\PY{+w}{ }\PY{k+kd}{static}\PY{+w}{ }\PY{k+kt}{void}\PY{+w}{ }\PY{n+nf}{register}\PY{p}{(}\PY{n}{String}\PY{+w}{ }\PY{n}{id}\PY{p}{,}\PY{+w}{ }\PY{n}{Session}\PY{+w}{ }\PY{n}{s}\PY{p}{)}\PY{+w}{ }\PY{p}{\PYZob{}}
\PY{+w}{        }\PY{n}{SESSIONS}\PY{p}{.}\PY{n+na}{put}\PY{p}{(}\PY{n}{id}\PY{p}{,}\PY{+w}{ }\PY{n}{s}\PY{p}{)}\PY{p}{;}
\PY{+w}{    }\PY{p}{\PYZcb{}}
\PY{+w}{    }\PY{c+c1}{// no remove() is ever called}
\PY{p}{\PYZcb{}}
\end{Verbatim}
\end{codebox}

A \code{static} field lives as long as the class is loaded — effectively forever in most applications. If nothing ever calls \code{remove}, this is an unbounded leak in plain sight.

\textbf{Unbounded caches.}

\begin{codebox}
\begin{Verbatim}[commandchars=\\\{\}]
\PY{c+c1}{// Wrong: a plain HashMap used as a cache has no eviction policy}
\PY{k+kd}{private}\PY{+w}{ }\PY{k+kd}{final}\PY{+w}{ }\PY{n}{Map}\PY{o}{\PYZlt{}}\PY{n}{Long}\PY{p}{,}\PY{+w}{ }\PY{n}{Report}\PY{o}{\PYZgt{}}\PY{+w}{ }\PY{n}{cache}\PY{+w}{ }\PY{o}{=}\PY{+w}{ }\PY{k}{new}\PY{+w}{ }\PY{n}{HashMap}\PY{o}{\PYZlt{}}\PY{o}{\PYZgt{}}\PY{p}{(}\PY{p}{)}\PY{p}{;}

\PY{k+kd}{public}\PY{+w}{ }\PY{n}{Report}\PY{+w}{ }\PY{n+nf}{get}\PY{p}{(}\PY{n}{Long}\PY{+w}{ }\PY{n}{id}\PY{p}{)}\PY{+w}{ }\PY{p}{\PYZob{}}
\PY{+w}{    }\PY{k}{return}\PY{+w}{ }\PY{n}{cache}\PY{p}{.}\PY{n+na}{computeIfAbsent}\PY{p}{(}\PY{n}{id}\PY{p}{,}\PY{+w}{ }\PY{k}{this}\PY{p}{:}\PY{p}{:}\PY{n}{loadReport}\PY{p}{)}\PY{p}{;}\PY{+w}{ }\PY{c+c1}{// grows forever}
\PY{p}{\PYZcb{}}
\end{Verbatim}
\end{codebox}

\begin{codebox}
\begin{Verbatim}[commandchars=\\\{\}]
\PY{c+c1}{// Right: bound it, with size or time\PYZhy{}based eviction}
\PY{k+kd}{private}\PY{+w}{ }\PY{k+kd}{final}\PY{+w}{ }\PY{n}{Cache}\PY{o}{\PYZlt{}}\PY{n}{Long}\PY{p}{,}\PY{+w}{ }\PY{n}{Report}\PY{o}{\PYZgt{}}\PY{+w}{ }\PY{n}{cache}\PY{+w}{ }\PY{o}{=}\PY{+w}{ }\PY{n}{Caffeine}\PY{p}{.}\PY{n+na}{newBuilder}\PY{p}{(}\PY{p}{)}
\PY{+w}{        }\PY{p}{.}\PY{n+na}{maximumSize}\PY{p}{(}\PY{l+m+mi}{10\PYZus{}000}\PY{p}{)}
\PY{+w}{        }\PY{p}{.}\PY{n+na}{expireAfterAccess}\PY{p}{(}\PY{n}{Duration}\PY{p}{.}\PY{n+na}{ofMinutes}\PY{p}{(}\PY{l+m+mi}{30}\PY{p}{)}\PY{p}{)}
\PY{+w}{        }\PY{p}{.}\PY{n+na}{build}\PY{p}{(}\PY{p}{)}\PY{p}{;}
\end{Verbatim}
\end{codebox}

\textbf{Listeners/callbacks never unregistered.}

\begin{codebox}
\begin{Verbatim}[commandchars=\\\{\}]
\PY{k+kd}{public}\PY{+w}{ }\PY{k+kd}{class} \PY{n+nc}{EventBus}\PY{+w}{ }\PY{p}{\PYZob{}}
\PY{+w}{    }\PY{k+kd}{private}\PY{+w}{ }\PY{k+kd}{final}\PY{+w}{ }\PY{n}{List}\PY{o}{\PYZlt{}}\PY{n}{Listener}\PY{o}{\PYZgt{}}\PY{+w}{ }\PY{n}{listeners}\PY{+w}{ }\PY{o}{=}\PY{+w}{ }\PY{k}{new}\PY{+w}{ }\PY{n}{ArrayList}\PY{o}{\PYZlt{}}\PY{o}{\PYZgt{}}\PY{p}{(}\PY{p}{)}\PY{p}{;}
\PY{+w}{    }\PY{k+kd}{public}\PY{+w}{ }\PY{k+kt}{void}\PY{+w}{ }\PY{n+nf}{subscribe}\PY{p}{(}\PY{n}{Listener}\PY{+w}{ }\PY{n}{l}\PY{p}{)}\PY{+w}{ }\PY{p}{\PYZob{}}\PY{+w}{ }\PY{n}{listeners}\PY{p}{.}\PY{n+na}{add}\PY{p}{(}\PY{n}{l}\PY{p}{)}\PY{p}{;}\PY{+w}{ }\PY{p}{\PYZcb{}}
\PY{+w}{    }\PY{c+c1}{// if callers subscribe but never call a matching unsubscribe(l),}
\PY{+w}{    }\PY{c+c1}{// every subscriber (and everything it references) leaks forever}
\PY{p}{\PYZcb{}}
\end{Verbatim}
\end{codebox}

This is especially common in UI frameworks and long-lived event buses. Always pair \code{subscribe}/\code{addListener} with a matching \code{unsubscribe}/\code{removeListener}, ideally enforced with try-finally or \code{AutoCloseable}.

\textbf{\code{ThreadLocal} in pooled threads.}

\begin{codebox}
\begin{Verbatim}[commandchars=\\\{\}]
\PY{c+c1}{// Wrong: thread pools reuse threads, so a ThreadLocal that\PYZsq{}s set but}
\PY{c+c1}{// never removed keeps its value (and everything it references) alive}
\PY{c+c1}{// for the LIFE OF THE POOLED THREAD, across unrelated tasks}
\PY{k+kd}{private}\PY{+w}{ }\PY{k+kd}{static}\PY{+w}{ }\PY{k+kd}{final}\PY{+w}{ }\PY{n}{ThreadLocal}\PY{o}{\PYZlt{}}\PY{n}{UserContext}\PY{o}{\PYZgt{}}\PY{+w}{ }\PY{n}{CONTEXT}\PY{+w}{ }\PY{o}{=}\PY{+w}{ }\PY{k}{new}\PY{+w}{ }\PY{n}{ThreadLocal}\PY{o}{\PYZlt{}}\PY{o}{\PYZgt{}}\PY{p}{(}\PY{p}{)}\PY{p}{;}

\PY{k+kt}{void}\PY{+w}{ }\PY{n+nf}{handleRequest}\PY{p}{(}\PY{n}{Request}\PY{+w}{ }\PY{n}{r}\PY{p}{)}\PY{+w}{ }\PY{p}{\PYZob{}}
\PY{+w}{    }\PY{n}{CONTEXT}\PY{p}{.}\PY{n+na}{set}\PY{p}{(}\PY{n}{loadUserContext}\PY{p}{(}\PY{n}{r}\PY{p}{)}\PY{p}{)}\PY{p}{;}
\PY{+w}{    }\PY{n}{process}\PY{p}{(}\PY{n}{r}\PY{p}{)}\PY{p}{;}
\PY{+w}{    }\PY{c+c1}{// missing CONTEXT.remove() — the UserContext (and its whole reference}
\PY{+w}{    }\PY{c+c1}{// graph) stays attached to this pool thread until it\PYZsq{}s overwritten}
\PY{p}{\PYZcb{}}
\end{Verbatim}
\end{codebox}

\begin{codebox}
\begin{Verbatim}[commandchars=\\\{\}]
\PY{c+c1}{// Right}
\PY{k+kt}{void}\PY{+w}{ }\PY{n+nf}{handleRequest}\PY{p}{(}\PY{n}{Request}\PY{+w}{ }\PY{n}{r}\PY{p}{)}\PY{+w}{ }\PY{p}{\PYZob{}}
\PY{+w}{    }\PY{n}{CONTEXT}\PY{p}{.}\PY{n+na}{set}\PY{p}{(}\PY{n}{loadUserContext}\PY{p}{(}\PY{n}{r}\PY{p}{)}\PY{p}{)}\PY{p}{;}
\PY{+w}{    }\PY{k}{try}\PY{+w}{ }\PY{p}{\PYZob{}}
\PY{+w}{        }\PY{n}{process}\PY{p}{(}\PY{n}{r}\PY{p}{)}\PY{p}{;}
\PY{+w}{    }\PY{p}{\PYZcb{}}\PY{+w}{ }\PY{k}{finally}\PY{+w}{ }\PY{p}{\PYZob{}}
\PY{+w}{        }\PY{n}{CONTEXT}\PY{p}{.}\PY{n+na}{remove}\PY{p}{(}\PY{p}{)}\PY{p}{;}\PY{+w}{ }\PY{c+c1}{// always clean up, even on exception}
\PY{+w}{    }\PY{p}{\PYZcb{}}
\PY{p}{\PYZcb{}}
\end{Verbatim}
\end{codebox}

\textbf{ClassLoader leaks.} Common in application servers that redeploy webapps without restarting the JVM. If any object (a thread, a static reference in a shared library, a JDBC driver registered in \code{DriverManager}) keeps a reference into classes loaded by the \emph{old} deployment’s classloader, that entire classloader — and every class and static field it loaded — cannot be garbage collected, even after redeploy. Symptom: \code{Out\allowbreak{}Of\allowbreak{}Memory\allowbreak{}Error:\allowbreak{}~\allowbreak{}Metaspace} (or the older \code{PermGen~\allowbreak{}space}) that grows with every redeploy.

\textbf{\code{substring} history (old JDKs).} Before Java 7u6, \code{String.\allowbreak{}substring()} shared the backing \code{char[]} array with the original string — a small substring of a huge string kept the \emph{entire} huge array alive.

\begin{codebox}
\begin{Verbatim}[commandchars=\\\{\}]
\PY{c+c1}{// Only a real leak on JDK \PYZlt{}= 6 / early 7 (fixed since 7u6 — substring copies now)}
\PY{n}{String}\PY{+w}{ }\PY{n}{hugeLine}\PY{+w}{ }\PY{o}{=}\PY{+w}{ }\PY{n}{readEntireFileAsOneLine}\PY{p}{(}\PY{p}{)}\PY{p}{;}\PY{+w}{ }\PY{c+c1}{// 10 MB string}
\PY{n}{String}\PY{+w}{ }\PY{n}{tinyPart}\PY{+w}{ }\PY{o}{=}\PY{+w}{ }\PY{n}{hugeLine}\PY{p}{.}\PY{n+na}{substring}\PY{p}{(}\PY{l+m+mi}{0}\PY{p}{,}\PY{+w}{ }\PY{l+m+mi}{5}\PY{p}{)}\PY{p}{;}\PY{+w}{  }\PY{c+c1}{// pre\PYZhy{}fix: still references the 10 MB array!}
\end{Verbatim}
\end{codebox}

On modern JDKs (7u6+) \code{substring} always copies, so this specific leak is historical — but it’s still a favorite interview question, and the general lesson (watch what a “small” object secretly references) still applies to things like \code{ByteBuffer} slices.

\textbf{Non-static inner classes holding outer references.}

\begin{codebox}
\begin{Verbatim}[commandchars=\\\{\}]
\PY{k+kd}{public}\PY{+w}{ }\PY{k+kd}{class} \PY{n+nc}{ReportGenerator}\PY{+w}{ }\PY{p}{\PYZob{}}\PY{+w}{           }\PY{c+c1}{// large object with lots of state}
\PY{+w}{    }\PY{k+kd}{private}\PY{+w}{ }\PY{k+kt}{byte}\PY{o}{[}\PY{o}{]}\PY{+w}{ }\PY{n}{hugeBuffer}\PY{+w}{ }\PY{o}{=}\PY{+w}{ }\PY{k}{new}\PY{+w}{ }\PY{k+kt}{byte}\PY{o}{[}\PY{l+m+mi}{100\PYZus{}000\PYZus{}000}\PY{o}{]}\PY{p}{;}

\PY{+w}{    }\PY{k+kd}{class} \PY{n+nc}{ProgressListener}\PY{+w}{ }\PY{p}{\PYZob{}}\PY{+w}{              }\PY{c+c1}{// non\PYZhy{}static inner class!}
\PY{+w}{        }\PY{k+kt}{void}\PY{+w}{ }\PY{n+nf}{onProgress}\PY{p}{(}\PY{k+kt}{int}\PY{+w}{ }\PY{n}{percent}\PY{p}{)}\PY{+w}{ }\PY{p}{\PYZob{}}\PY{+w}{ }\PY{n}{System}\PY{p}{.}\PY{n+na}{out}\PY{p}{.}\PY{n+na}{println}\PY{p}{(}\PY{n}{percent}\PY{p}{)}\PY{p}{;}\PY{+w}{ }\PY{p}{\PYZcb{}}
\PY{+w}{    }\PY{p}{\PYZcb{}}

\PY{+w}{    }\PY{n}{ProgressListener}\PY{+w}{ }\PY{n+nf}{newListener}\PY{p}{(}\PY{p}{)}\PY{+w}{ }\PY{p}{\PYZob{}}
\PY{+w}{        }\PY{k}{return}\PY{+w}{ }\PY{k}{new}\PY{+w}{ }\PY{n}{ProgressListener}\PY{p}{(}\PY{p}{)}\PY{p}{;}\PY{+w}{    }\PY{c+c1}{// implicitly holds ReportGenerator.this}
\PY{+w}{    }\PY{p}{\PYZcb{}}
\PY{p}{\PYZcb{}}
\end{Verbatim}
\end{codebox}

Every instance of a non-static inner class carries a hidden reference to its enclosing instance. If that listener is registered somewhere long-lived (a static event bus, a cache), the \emph{entire} outer \code{ReportGenerator} — including its 100 MB buffer — leaks with it. Fix: make the inner class \code{static} (or a top-level class) if it doesn’t need the outer instance, and pass in only the specific data it needs.

\textbf{Unclosed resources.}

\begin{codebox}
\begin{Verbatim}[commandchars=\\\{\}]
\PY{c+c1}{// Wrong: if readAll() throws, the stream is never closed — leaks a file handle}
\PY{c+c1}{// and, since streams often hold internal buffers, JVM\PYZhy{}heap memory too}
\PY{n}{FileInputStream}\PY{+w}{ }\PY{n}{in}\PY{+w}{ }\PY{o}{=}\PY{+w}{ }\PY{k}{new}\PY{+w}{ }\PY{n}{FileInputStream}\PY{p}{(}\PY{n}{path}\PY{p}{)}\PY{p}{;}
\PY{k+kt}{byte}\PY{o}{[}\PY{o}{]}\PY{+w}{ }\PY{n}{data}\PY{+w}{ }\PY{o}{=}\PY{+w}{ }\PY{n}{in}\PY{p}{.}\PY{n+na}{readAllBytes}\PY{p}{(}\PY{p}{)}\PY{p}{;}
\PY{n}{in}\PY{p}{.}\PY{n+na}{close}\PY{p}{(}\PY{p}{)}\PY{p}{;}

\PY{c+c1}{// Right}
\PY{k}{try}\PY{+w}{ }\PY{p}{(}\PY{n}{FileInputStream}\PY{+w}{ }\PY{n}{in}\PY{+w}{ }\PY{o}{=}\PY{+w}{ }\PY{k}{new}\PY{+w}{ }\PY{n}{FileInputStream}\PY{p}{(}\PY{n}{path}\PY{p}{)}\PY{p}{)}\PY{+w}{ }\PY{p}{\PYZob{}}
\PY{+w}{    }\PY{k+kt}{byte}\PY{o}{[}\PY{o}{]}\PY{+w}{ }\PY{n}{data}\PY{+w}{ }\PY{o}{=}\PY{+w}{ }\PY{n}{in}\PY{p}{.}\PY{n+na}{readAllBytes}\PY{p}{(}\PY{p}{)}\PY{p}{;}
\PY{p}{\PYZcb{}}\PY{+w}{ }\PY{c+c1}{// closed automatically, even on exception}
\end{Verbatim}
\end{codebox}

\textbf{Direct (off-heap) buffer leaks.} \code{Byte\allowbreak{}Buffer.\allowbreak{}allocate\allowbreak{}Direct()} and memory-mapped files allocate native memory \emph{outside} the Java heap, freed only when the buffer object itself is garbage collected and its cleaner runs. Holding on to direct buffers longer than necessary (e.g., caching them the same way you’d cache a normal object) leaks native memory that \textbf{won’t show up in a heap dump or \code{-\allowbreak{}Xmx} usage at all} — it shows up as the process’s resident memory (RSS) growing past what the heap accounts for, which is a common source of confusing container OOM-kills where “the heap looks fine.”

\begin{codebox}
\begin{Verbatim}[commandchars=\\\{\}]
\PY{c+c1}{// Wrong: direct buffers cached indefinitely, never explicitly released}
\PY{k+kd}{private}\PY{+w}{ }\PY{k+kd}{static}\PY{+w}{ }\PY{k+kd}{final}\PY{+w}{ }\PY{n}{Map}\PY{o}{\PYZlt{}}\PY{n}{String}\PY{p}{,}\PY{+w}{ }\PY{n}{ByteBuffer}\PY{o}{\PYZgt{}}\PY{+w}{ }\PY{n}{BUFFER\PYZus{}CACHE}\PY{+w}{ }\PY{o}{=}\PY{+w}{ }\PY{k}{new}\PY{+w}{ }\PY{n}{HashMap}\PY{o}{\PYZlt{}}\PY{o}{\PYZgt{}}\PY{p}{(}\PY{p}{)}\PY{p}{;}
\PY{n}{ByteBuffer}\PY{+w}{ }\PY{n}{buf}\PY{+w}{ }\PY{o}{=}\PY{+w}{ }\PY{n}{ByteBuffer}\PY{p}{.}\PY{n+na}{allocateDirect}\PY{p}{(}\PY{l+m+mi}{64}\PY{+w}{ }\PY{o}{*}\PY{+w}{ }\PY{l+m+mi}{1024}\PY{+w}{ }\PY{o}{*}\PY{+w}{ }\PY{l+m+mi}{1024}\PY{p}{)}\PY{p}{;}\PY{+w}{ }\PY{c+c1}{// 64 MB native memory}
\PY{n}{BUFFER\PYZus{}CACHE}\PY{p}{.}\PY{n+na}{put}\PY{p}{(}\PY{n}{key}\PY{p}{,}\PY{+w}{ }\PY{n}{buf}\PY{p}{)}\PY{p}{;}\PY{+w}{ }\PY{c+c1}{// never removed — native memory leak, invisible on the heap}

\PY{c+c1}{// Right: bound the cache and/or size direct memory explicitly with}
\PY{c+c1}{// \PYZhy{}XX:MaxDirectMemorySize=\PYZlt{}size\PYZgt{} so a leak throws an OutOfMemoryError instead}
\PY{c+c1}{// of silently exhausting container memory}
\end{Verbatim}
\end{codebox}

\textbf{Summary of causes and fixes:}

{\tablefontsmall
\begin{xltabular}{\linewidth}{@{}L{0.648}L{1.176}L{1.176}@{}}
\toprule
\rowcolor{tablehdr}
\thd{Cause} & \thd{Why it leaks} & \thd{Fix} \\
\midrule
\endhead
Static collection that only grows & Static fields live for the program’s lifetime & Bound size, add explicit removal \\
Unbounded cache & No eviction policy & Use a cache library with \code{maximumSize}/TTL \\
Listener never unregistered & Event source outlives the “logical” subscriber & Pair every \code{subscribe} with \code{unsubscribe} \\
\code{ThreadLocal} in a pool & Pooled threads outlive individual tasks & Always call \code{remove()} in a \code{finally} \\
ClassLoader leak & A live reference reaches into an “unloaded” deployment’s classes & Deregister drivers/threads on shutdown; restart JVM on redeploy if needed \\
\code{substring} history & Pre-7u6 shared backing array & N/A on modern JDKs; still watch buffer slices \\
Non-static inner class & Hidden reference to outer instance & Make it \code{static} unless it needs the outer instance \\
Unclosed resource & Native/buffered state never released & \code{try}-with-resources \\
Direct buffer left cached & Off-heap memory tied to buffer object’s lifecycle & Bound cache; set \code{-\allowbreak{}XX:\allowbreak{}Max\allowbreak{}Direct\allowbreak{}Memory\allowbreak{}Size} \\
\bottomrule
\end{xltabular}
}

\textbf{How to find a leak: heap dump + dominator tree.}

\begin{enumerate}
\item 
Capture a heap dump while the process is exhibiting the growth (or trigger one automatically on OOM — see \hyperref[h:12:jvm-flags-and-diagnostics]{JVM Flags and Diagnostics}):

\begin{codebox}
\begin{Verbatim}[commandchars=\\\{\}]
jcmd\PY{+w}{ }\PYZlt{}pid\PYZgt{}\PY{+w}{ }GC.heap\PYZus{}dump\PY{+w}{ }/tmp/heap.hprof
\end{Verbatim}
\end{codebox}

\item 
Open it in a heap dump analyzer (Eclipse MAT, or JFR/JMC’s heap dump view, or VisualVM).

\item 
Look at the \textbf{dominator tree}, not just “biggest objects.” An object X \emph{dominates} object Y if every path from a GC root to Y passes through X — meaning if X were freed, Y would become garbage too. The dominator tree groups the \emph{retained} size (an object plus everything it exclusively keeps alive) under its dominator, so the true “root cause” object floats to the top instead of being hidden among thousands of small leaked instances.

\item 
Compare two heap dumps taken minutes apart (MAT’s “compare” feature) to see which object counts are \emph{still climbing} — that’s a much stronger leak signal than a single snapshot, since some growth is normal caching, not a leak.

\item 
Once you’ve found the dominating object (say, a \code{HashMap} field on some \code{SessionRegistry}), trace back through its “path to GC roots” to find \emph{why} it’s still reachable, then fix the reachability (add eviction, call \code{remove}, unregister the listener, etc.).

\end{enumerate}

\section[Escape Analysis]{Escape Analysis}
\label{h:12:escape-analysis}
\textbf{What it is.} Escape analysis is a JIT compiler optimization (see Chapter 11) that determines whether an object created inside a method can be observed — can “escape” — outside that method (returned, stored in a field, passed to another thread, etc.). If the JIT can \emph{prove} an object never escapes its allocating method or thread, it can apply optimizations that would be unsafe for an object visible elsewhere.

\textbf{Scalar replacement.} If an object provably never escapes, the JIT can skip allocating it as an object on the heap at all and instead break it into its individual primitive fields (“scalars”), keeping those in registers or on the stack — exactly as if you’d hand-written the method using separate local variables instead of a wrapper object.

\begin{codebox}
\begin{Verbatim}[commandchars=\\\{\}]
\PY{k+kd}{class} \PY{n+nc}{Point}\PY{+w}{ }\PY{p}{\PYZob{}}
\PY{+w}{    }\PY{k+kt}{int}\PY{+w}{ }\PY{n}{x}\PY{p}{,}\PY{+w}{ }\PY{n}{y}\PY{p}{;}
\PY{+w}{    }\PY{n}{Point}\PY{p}{(}\PY{k+kt}{int}\PY{+w}{ }\PY{n}{x}\PY{p}{,}\PY{+w}{ }\PY{k+kt}{int}\PY{+w}{ }\PY{n}{y}\PY{p}{)}\PY{+w}{ }\PY{p}{\PYZob{}}\PY{+w}{ }\PY{k}{this}\PY{p}{.}\PY{n+na}{x}\PY{+w}{ }\PY{o}{=}\PY{+w}{ }\PY{n}{x}\PY{p}{;}\PY{+w}{ }\PY{k}{this}\PY{p}{.}\PY{n+na}{y}\PY{+w}{ }\PY{o}{=}\PY{+w}{ }\PY{n}{y}\PY{p}{;}\PY{+w}{ }\PY{p}{\PYZcb{}}
\PY{p}{\PYZcb{}}

\PY{k+kt}{int}\PY{+w}{ }\PY{n+nf}{distanceSquared}\PY{p}{(}\PY{k+kt}{int}\PY{+w}{ }\PY{n}{x1}\PY{p}{,}\PY{+w}{ }\PY{k+kt}{int}\PY{+w}{ }\PY{n}{y1}\PY{p}{,}\PY{+w}{ }\PY{k+kt}{int}\PY{+w}{ }\PY{n}{x2}\PY{p}{,}\PY{+w}{ }\PY{k+kt}{int}\PY{+w}{ }\PY{n}{y2}\PY{p}{)}\PY{+w}{ }\PY{p}{\PYZob{}}
\PY{+w}{    }\PY{n}{Point}\PY{+w}{ }\PY{n}{p1}\PY{+w}{ }\PY{o}{=}\PY{+w}{ }\PY{k}{new}\PY{+w}{ }\PY{n}{Point}\PY{p}{(}\PY{n}{x1}\PY{p}{,}\PY{+w}{ }\PY{n}{y1}\PY{p}{)}\PY{p}{;}\PY{+w}{ }\PY{c+c1}{// may never actually be allocated on the heap}
\PY{+w}{    }\PY{n}{Point}\PY{+w}{ }\PY{n}{p2}\PY{+w}{ }\PY{o}{=}\PY{+w}{ }\PY{k}{new}\PY{+w}{ }\PY{n}{Point}\PY{p}{(}\PY{n}{x2}\PY{p}{,}\PY{+w}{ }\PY{n}{y2}\PY{p}{)}\PY{p}{;}\PY{+w}{ }\PY{c+c1}{// if the JIT proves neither escapes this method}
\PY{+w}{    }\PY{k+kt}{int}\PY{+w}{ }\PY{n}{dx}\PY{+w}{ }\PY{o}{=}\PY{+w}{ }\PY{n}{p1}\PY{p}{.}\PY{n+na}{x}\PY{+w}{ }\PY{o}{\PYZhy{}}\PY{+w}{ }\PY{n}{p2}\PY{p}{.}\PY{n+na}{x}\PY{p}{;}
\PY{+w}{    }\PY{k+kt}{int}\PY{+w}{ }\PY{n}{dy}\PY{+w}{ }\PY{o}{=}\PY{+w}{ }\PY{n}{p1}\PY{p}{.}\PY{n+na}{y}\PY{+w}{ }\PY{o}{\PYZhy{}}\PY{+w}{ }\PY{n}{p2}\PY{p}{.}\PY{n+na}{y}\PY{p}{;}
\PY{+w}{    }\PY{k}{return}\PY{+w}{ }\PY{n}{dx}\PY{+w}{ }\PY{o}{*}\PY{+w}{ }\PY{n}{dx}\PY{+w}{ }\PY{o}{+}\PY{+w}{ }\PY{n}{dy}\PY{+w}{ }\PY{o}{*}\PY{+w}{ }\PY{n}{dy}\PY{p}{;}
\PY{p}{\PYZcb{}}
\end{Verbatim}
\end{codebox}

\textbf{Lock elision.} Similarly, if the JIT proves an object is only ever visible to one thread (it never escapes to another thread), any \code{synchronized} block guarding that object is provably uncontended — no other thread could possibly be waiting on that monitor — so the JIT can remove the locking overhead entirely.

\begin{codebox}
\begin{Verbatim}[commandchars=\\\{\}]
\PY{k+kt}{void}\PY{+w}{ }\PY{n+nf}{process}\PY{p}{(}\PY{p}{)}\PY{+w}{ }\PY{p}{\PYZob{}}
\PY{+w}{    }\PY{n}{StringBuilder}\PY{+w}{ }\PY{n}{sb}\PY{+w}{ }\PY{o}{=}\PY{+w}{ }\PY{k}{new}\PY{+w}{ }\PY{n}{StringBuilder}\PY{p}{(}\PY{p}{)}\PY{p}{;}\PY{+w}{ }\PY{c+c1}{// local, never escapes this method}
\PY{+w}{    }\PY{k+kd}{synchronized}\PY{+w}{ }\PY{p}{(}\PY{n}{sb}\PY{p}{)}\PY{+w}{ }\PY{p}{\PYZob{}}\PY{+w}{                      }\PY{c+c1}{// lock is pointless — nothing else can see \PYZsq{}sb\PYZsq{}}
\PY{+w}{        }\PY{n}{sb}\PY{p}{.}\PY{n+na}{append}\PY{p}{(}\PY{l+s}{\PYZdq{}}\PY{l+s}{data}\PY{l+s}{\PYZdq{}}\PY{p}{)}\PY{p}{;}
\PY{+w}{    }\PY{p}{\PYZcb{}}
\PY{+w}{    }\PY{c+c1}{// JIT may elide the lock/unlock entirely}
\PY{p}{\PYZcb{}}
\end{Verbatim}
\end{codebox}

\textbf{Why “objects are always on the heap” is not strictly true.} It’s a common simplification for beginners, and it’s \emph{usually} true — but escape analysis is a real, everyday counterexample. A non-escaping object can live entirely in registers/stack, never touching the heap and never costing the GC anything. This is why microbenchmarks that appear to allocate millions of objects sometimes show suspiciously low GC activity — the JIT proved the objects never escaped and optimized the allocation away (this is also a classic JMH benchmarking pitfall — see Chapter 20).

\textbf{Why you should NOT design code around it.} Escape analysis is a best-effort, unguaranteed, JIT-tier-dependent optimization:

\begin{itemize}
\item 
It only kicks in after the JIT has profiled the method enough to compile it at a high tier (see Chapter 11 on JIT compilation) — cold code gets no benefit.

\item 
Any change that makes escape harder to prove — passing the object to an interface method the JIT can’t fully inline, storing it in a collection “just for logging,” returning it conditionally — silently disables the optimization, with no compiler warning.

\item 
It’s a JVM/version-specific implementation detail, not part of the Java Language Specification. Different JVMs (or the same JVM with different flags) may or may not apply it.

\end{itemize}

Write clear code — small objects, minimal scope, no unnecessary escapes — and let escape analysis take advantage of it \emph{if} it can. Never rely on it being active as a correctness or capacity-planning assumption.

\section[JVM Flags and Diagnostics]{JVM Flags and Diagnostics}
\label{h:12:jvm-flags-and-diagnostics}
Flags are grouped into standard (\code{-\allowbreak{}X...}), and the vast “extra” (\code{-\allowbreak{}XX:...}) category, which itself splits into product flags (safe, documented) and experimental/diagnostic flags (need \code{-\allowbreak{}XX:+\allowbreak{}Unlock\allowbreak{}Experimental\allowbreak{}VMOptions} or \code{-\allowbreak{}XX:+\allowbreak{}Unlock\allowbreak{}Diagnostic\allowbreak{}VMOptions}).

Commonly needed flags for memory work:

{\tablefontsmall
\begin{xltabular}{\linewidth}{@{}L{0.512}L{1.488}@{}}
\toprule
\rowcolor{tablehdr}
\thd{Flag} & \thd{Purpose} \\
\midrule
\endhead
\code{-\allowbreak{}Xms\textless{}\allowbreak{}size\textgreater{}} & Initial heap size (e.g. \code{-\allowbreak{}Xms512m}) \\
\code{-\allowbreak{}Xmx\textless{}\allowbreak{}size\textgreater{}} & Maximum heap size (e.g. \code{-\allowbreak{}Xmx4g}) \\
\code{-\allowbreak{}Xss\textless{}\allowbreak{}size\textgreater{}} & Thread stack size \\
\code{-\allowbreak{}XX:\allowbreak{}Metaspace\allowbreak{}Size=\textless{}\allowbreak{}size\textgreater{}} & Initial metaspace size before triggering a GC to expand it \\
\code{-\allowbreak{}XX:\allowbreak{}Max\allowbreak{}Metaspace\allowbreak{}Size=\textless{}\allowbreak{}size\textgreater{}} & Cap on metaspace (class metadata) — unbounded by default! \\
\code{-\allowbreak{}XX:+\allowbreak{}Heap\allowbreak{}Dump\allowbreak{}On\allowbreak{}Out\allowbreak{}Of\allowbreak{}Memory\allowbreak{}Error} & Automatically write a \code{.\allowbreak{}hprof} heap dump when an \code{OutOfMemoryError} is thrown \\
\code{-\allowbreak{}XX:\allowbreak{}Heap\allowbreak{}Dump\allowbreak{}Path=\textless{}\allowbreak{}path\textgreater{}} & Where to write that automatic heap dump \\
\code{-\allowbreak{}XX:+\allowbreak{}PrintFlagsFinal} & Print the effective value of every JVM flag (great for checking what a collector chose automatically) \\
\code{-\allowbreak{}Xlog:\allowbreak{}gc*:\allowbreak{}file=\allowbreak{}gc.\allowbreak{}log:\allowbreak{}time,\allowbreak{}uptime,\allowbreak{}level,\allowbreak{}tags} & Unified logging: write detailed GC logs with timestamps to a file \\
\code{-\allowbreak{}XX:\allowbreak{}Max\allowbreak{}Direct\allowbreak{}Memory\allowbreak{}Size=\textless{}\allowbreak{}size\textgreater{}} & Cap on off-heap \code{Byte\allowbreak{}Buffer.\allowbreak{}allocate\allowbreak{}Direct()} memory — unbounded by default \\
\code{-\allowbreak{}XX:+\allowbreak{}Exit\allowbreak{}On\allowbreak{}Out\allowbreak{}Of\allowbreak{}Memory\allowbreak{}Error} & Kill the JVM immediately on any \code{OutOfMemoryError}, instead of limping along in a corrupted state (useful under an orchestrator that will restart the pod/container) \\
\code{-\allowbreak{}XX:+\allowbreak{}Use\allowbreak{}String\allowbreak{}Deduplication} & (G1 only) Collapse duplicate \code{String} backing \code{char[]/\allowbreak{}byte[]} arrays during GC to save heap — helpful for apps that hold many identical strings (e.g. parsed JSON keys) \\
\code{-\allowbreak{}XX:+\allowbreak{}AlwaysPreTouch} & Commit and zero all heap pages at startup instead of lazily — trades slower startup for eliminating page-fault pauses later, useful with \code{-\allowbreak{}Xms~\allowbreak{}==\allowbreak{}~\allowbreak{}-\allowbreak{}Xmx} \\
\code{-\allowbreak{}XX:\allowbreak{}Native\allowbreak{}Memory\allowbreak{}Tracking=\allowbreak{}summary} & Turn on Native Memory Tracking (NMT) to account for \emph{non-heap} JVM memory (thread stacks, GC structures, code cache, metaspace) \\
\bottomrule
\end{xltabular}
}

\begin{codebox}
\begin{Verbatim}[commandchars=\\\{\}]
\PY{c+c1}{\PYZsh{} See exactly what heap size G1 picked by default on this machine}
java\PY{+w}{ }\PYZhy{}XX:+PrintFlagsFinal\PY{+w}{ }\PYZhy{}version\PY{+w}{ }\PY{p}{|}\PY{+w}{ }grep\PY{+w}{ }\PYZhy{}i\PY{+w}{ }heapsize

\PY{c+c1}{\PYZsh{} Turn on detailed GC logging to a rotating file}
java\PY{+w}{ }\PYZhy{}Xlog:gc*:file\PY{o}{=}gc.log:time,uptime,level,tags:filecount\PY{o}{=}\PY{l+m}{5},filesize\PY{o}{=}10M\PY{+w}{ }\PYZhy{}jar\PY{+w}{ }app.jar
\end{Verbatim}
\end{codebox}

\textbf{Native Memory Tracking (NMT).} The heap is only part of a JVM process’s memory footprint. NMT breaks down exactly where the \emph{rest} went — thread stacks, GC bookkeeping, the JIT code cache, metaspace, internal buffers — which is exactly what you need when RSS (resident memory) is much bigger than \code{-\allowbreak{}Xmx} would suggest (a classic container OOM-kill mystery).

\begin{codebox}
\begin{Verbatim}[commandchars=\\\{\}]
\PY{c+c1}{\PYZsh{} Enable at startup (small, constant overhead — safe for production)}
java\PY{+w}{ }\PYZhy{}XX:NativeMemoryTracking\PY{o}{=}summary\PY{+w}{ }\PYZhy{}jar\PY{+w}{ }app.jar

\PY{c+c1}{\PYZsh{} Query it later with jcmd}
jcmd\PY{+w}{ }\PYZlt{}pid\PYZgt{}\PY{+w}{ }VM.native\PYZus{}memory\PY{+w}{ }summary
\end{Verbatim}
\end{codebox}

\begin{codebox}
\begin{Verbatim}
Total: reserved=5324912KB, committed=2109344KB
-               Java Heap (reserved=4194304KB, committed=1048576KB)
-                   Class (reserved=1069056KB, committed=45120KB)
-                  Thread (reserved=32784KB, committed=32784KB)
-                      GC (reserved=25368KB, committed=25368KB)
-                  Symbol (reserved=15236KB, committed=15236KB)
\end{Verbatim}
\end{codebox}

If “Thread” is unexpectedly huge, that’s often thousands of leaked/unbounded threads each holding a full stack (\code{-\allowbreak{}Xss}, default \textasciitilde{}512KB-1MB each) — a different flavor of “memory leak” than a heap leak, but just as real.

Reading a GC log line (unified logging, JDK 9+):

\begin{codebox}
\begin{Verbatim}
[2026-08-07T10:15:22.123+0000][3.456s][info][gc] GC(12) Pause Young (Normal) (G1 Evacuation Pause) 512M->128M(1024M) 8.213ms
\end{Verbatim}
\end{codebox}

Read this as: at wall-clock time \code{10:\allowbreak{}15:\allowbreak{}22}, \code{3.\allowbreak{}456s} after JVM start, GC cycle \#12 was a young “Normal” pause caused by a G1 evacuation. Heap usage went from 512 MB \emph{before} the collection to 128 MB \emph{after}, out of a current heap capacity of 1024 MB, and the whole pause took 8.213 ms. Watch three things across many lines: pause \textbf{duration} (spikes = latency problem), the \textbf{before→after gap} (how much was actually reclaimed — if it’s small and shrinking over time, that’s a leak signature), and \textbf{frequency} (collections getting closer together = allocation rate is outrunning what the heap can absorb).

\section[Garbage Collection Tuning]{Garbage Collection Tuning}
\label{h:12:garbage-collection-tuning}
Tuning a collector means adjusting a handful of knobs, always \textbf{after} you’ve measured, never before:

\begin{itemize}
\item 
\textbf{Pick the pause-time target, not a lower-level knob, first.} For G1: \code{-\allowbreak{}XX:\allowbreak{}Max\allowbreak{}GCPause\allowbreak{}Millis=\allowbreak{}100}. G1 will then choose region counts and collection sets to try to hit that target — you’re telling it \emph{what you want}, not \emph{how to get there}.

\item 
\textbf{Heap sizing.} A too-small heap causes frequent GCs; a too-large heap causes rare but longer full GCs and wastes memory. Start with \code{-\allowbreak{}Xms} and \code{-\allowbreak{}Xmx} set to the \emph{same} value in server/container deployments (see \hyperref[h:12:jvm-performance-tuning]{JVM Performance Tuning}) so the heap doesn’t spend time resizing itself under load.

\item 
\textbf{New (young) generation size.} \code{-\allowbreak{}XX:\allowbreak{}NewRatio=\textless{}\allowbreak{}n\textgreater{}} (old:young ratio) or explicit \code{-\allowbreak{}XX:\allowbreak{}NewSize=\textless{}\allowbreak{}size\textgreater{}} — a larger young gen means fewer, but bigger, minor GCs; too large and minor GC pauses grow.

\item 
\textbf{Survivor space tuning.} \code{-\allowbreak{}XX:\allowbreak{}SurvivorRatio=\textless{}\allowbreak{}n\textgreater{}} — rarely needed unless promotion is happening too eagerly or too late (check with \code{-\allowbreak{}Xlog:\allowbreak{}gc+\allowbreak{}age=\allowbreak{}trace}).

\item 
\textbf{Concurrent GC threads.} \code{-\allowbreak{}XX:\allowbreak{}ConcGCThreads=\textless{}\allowbreak{}n\textgreater{}} for G1/ZGC/Shenandoah’s background work; \code{-\allowbreak{}XX:\allowbreak{}Parallel\allowbreak{}GCThreads=\textless{}\allowbreak{}n\textgreater{}} for the STW parallel phases. Defaults scale with CPU count and are usually fine.

\item 
\textbf{Watch for promotion failure / to-space exhaustion.} If the old generation doesn’t have enough free space to accept everything a minor GC needs to promote, the collector falls back to a full GC to make room — an expensive, avoidable pause. The GC log shows this explicitly:

\begin{codebox}
\begin{Verbatim}
[gc] GC(77) To-space exhausted
[gc] GC(78) Pause Full (G1 Evacuation Pause) 3800M->3750M(4096M) 1044.221ms
\end{Verbatim}
\end{codebox}

Fixes: grow the heap, grow the young generation’s headroom, or address whatever is driving the promotion rate up (see the leak-hunting steps above — a real leak often first shows up as promotion failures before it shows up as an outright OOM).

\item 
\textbf{One change at a time, measured.} Change a single flag, run a representative load test, compare GC logs before/after. Changing three flags simultaneously and seeing an improvement tells you nothing about \emph{which} flag helped (or which one is secretly hurting and being masked by the others).

\end{itemize}

\section[JVM Performance Tuning]{JVM Performance Tuning}
\label{h:12:jvm-performance-tuning}
Memory tuning is one piece of the larger performance-tuning picture:

\begin{itemize}
\item 
\textbf{\code{-\allowbreak{}Xms} == \code{-\allowbreak{}Xmx} in containers.} Setting the initial and maximum heap to the \emph{same} value avoids runtime heap-resizing pauses and — critically in Kubernetes/Docker — makes memory usage \textbf{predictable}, which matters because the container’s memory limit will hard-kill (OOMKilled) the process if the JVM heap plus off-heap usage exceeds it, regardless of what the JVM \emph{thinks} its budget is.

\item 
\textbf{\code{-\allowbreak{}XX:\allowbreak{}Max\allowbreak{}RAMPercentage=\textless{}\allowbreak{}pct\textgreater{}}} — instead of a hardcoded \code{-\allowbreak{}Xmx}, tell the JVM to use a \emph{percentage} of the container’s memory limit (default up to 25\% historically; explicitly set higher, e.g. \code{75.\allowbreak{}0}, for a heap-heavy service). This is the recommended approach for containers because it automatically adapts if the container’s memory limit changes (e.g. between environments) without editing JVM flags.

\begin{codebox}
\begin{Verbatim}[commandchars=\\\{\}]
java\PY{+w}{ }\PYZhy{}XX:MaxRAMPercentage\PY{o}{=}\PY{l+m}{75}.0\PY{+w}{ }\PYZhy{}XX:InitialRAMPercentage\PY{o}{=}\PY{l+m}{75}.0\PY{+w}{ }\PYZhy{}jar\PY{+w}{ }app.jar
\end{Verbatim}
\end{codebox}

\item 
\textbf{Container awareness.} Modern JVMs (since JDK 10, mature by JDK 17+) read cgroup limits automatically — \code{Runtime.\allowbreak{}available\allowbreak{}Processors()} and \code{MaxRAMPercentage} calculations respect the container’s \emph{actual} CPU and memory limits, not the host machine’s. Verify with \code{-\allowbreak{}XX:+\allowbreak{}PrintFlagsFinal} inside the container, not on your laptop — a value that looks right locally can be wrong inside a 512 MB container.

\item 
\textbf{CPU quota affects GC thread count too.} \code{-\allowbreak{}XX:\allowbreak{}ParallelGCThreads} and \code{-\allowbreak{}XX:\allowbreak{}ConcGCThreads} default based on the \emph{visible} CPU count. In a container with a low CPU quota (e.g. \code{limits.\allowbreak{}cpu:\allowbreak{}~\allowbreak{}\textquotedbl{}2\textquotedbl{}} on a 32-core host), an old or misconfigured JVM/runtime combination that doesn’t honor cgroup CPU quotas correctly can size GC thread pools for 32 cores, causing them to thrash for CPU time against the application’s own threads — worth checking explicitly with \code{jcmd~\allowbreak{}\textless{}\allowbreak{}pid\textgreater{}\allowbreak{}~\allowbreak{}VM.\allowbreak{}flags~\allowbreak{}\textbar{}\allowbreak{}~\allowbreak{}grep~\allowbreak{}GCThreads} rather than assuming it’s right.

\item 
\textbf{Off-heap memory needs its own budget.} Direct buffers (\code{-\allowbreak{}XX:\allowbreak{}Max\allowbreak{}Direct\allowbreak{}Memory\allowbreak{}Size}), the JIT code cache, metaspace, and thread stacks are all outside \code{-\allowbreak{}Xmx} — when sizing a container’s memory limit, budget for heap \textbf{plus} all of these, not just the heap. A common rule of thumb is heap = \textasciitilde{}70-75\% of the container limit, leaving room for everything else; verify with NMT rather than guessing.

\item 
\textbf{Metaspace sizing.} Metaspace holds class metadata (not object instances) and is off-heap by default with \textbf{no cap} unless you set one — a classloader leak (see above) can silently grow metaspace until the \emph{whole container} runs out of memory, not just “the heap.” Set \code{-\allowbreak{}XX:\allowbreak{}MaxMetaspaceSize} explicitly in production so a metaspace leak becomes a visible, boring \code{Out\allowbreak{}Of\allowbreak{}Memory\allowbreak{}Error:\allowbreak{}~\allowbreak{}Metaspace} instead of an unpredictable OS-level OOM kill.

\item 
\textbf{Choosing a pause target.} Ask “what does the SLA actually require?” before picking a collector or a \code{MaxGCPauseMillis} value. A batch job with no user waiting has no reason to pay ZGC’s throughput/memory tax; a p99-latency-sensitive API absolutely does.

\item 
\textbf{The golden rule: measure first, one change at a time.} Get a baseline (GC logs, JFR recording, actual latency/throughput numbers) \emph{before} touching any flag. Change one thing. Re-measure under the same load. Keep the change only if it measurably helps your actual metric (not a synthetic microbenchmark). Tuning by folklore (“everyone sets \code{-\allowbreak{}XX:\allowbreak{}NewRatio=\allowbreak{}2}”) without measurement on \emph{your} workload is a code smell in itself.

\end{itemize}

\section[Profiling and Monitoring]{Profiling and Monitoring}
\label{h:12:profiling-and-monitoring}
\textbf{Profiling} answers “where is time/memory actually going” with low enough overhead to run in production; \textbf{monitoring} answers “is anything wrong right now” continuously over time. For memory specifically, you want to watch: heap usage over time (sawtooth is healthy — a rising floor is a leak), GC pause frequency/duration, allocation rate (bytes/sec), and promotion rate (how much survives into old gen per minor GC). The tools below (JFR, JMC, \code{jcmd}, \code{jmap}, \code{jstack}) are the practical toolbox; a good code reviewer should be able to name which tool answers which question without reaching for a web search.

{\tablefontsmall
\begin{xltabular}{\linewidth}{@{}L{1.032}L{0.968}@{}}
\toprule
\rowcolor{tablehdr}
\thd{Question} & \thd{Reach for} \\
\midrule
\endhead
“What’s the effective value of every JVM flag right now?” & \code{jcmd~\allowbreak{}\textless{}\allowbreak{}pid\textgreater{}\allowbreak{}~\allowbreak{}VM.\allowbreak{}flags} \\
“What’s happening in this app over the last hour, in production, with low overhead?” & JFR recording, analyzed in JMC \\
“Why did we just get an OutOfMemoryError?” & Heap dump (\code{jmap}/\code{jcmd}) + Eclipse MAT or JMC \\
“Are threads stuck / is there a deadlock right now?” & \code{jstack} / \code{jcmd~\allowbreak{}Thread.\allowbreak{}print} \\
“What is allocating the most garbage?” & JFR allocation profiling, viewed in JMC \\
“Is memory usage trending up or holding steady over days?” & A metrics pipeline (Micrometer/Prometheus scraping JMX GC/heap metrics) plus dashboards, not one-off tool runs \\
“Quick, lightweight look at live heap/GC/thread stats without a full recording?” & \code{jstat} (e.g. \code{jstat~\allowbreak{}-\allowbreak{}gcutil~\allowbreak{}\textless{}\allowbreak{}pid\textgreater{}\allowbreak{}~\allowbreak{}1000}) or VisualVM’s live view \\
“Where is CPU time going, down to the line, with minimal overhead?” & \code{async-\allowbreak{}profiler} (third-party, flame graphs) or JFR’s method-profiling events \\
\bottomrule
\end{xltabular}
}

For continuous production monitoring (as opposed to point-in-time debugging), most teams expose JVM metrics through Micrometer or the JMX GC/heap/thread MBeans into Prometheus/Grafana or an APM tool, and alert on trends (heap floor rising over days, GC pause p99 crossing a threshold) rather than manually running these tools reactively. The manual tools in this chapter are for \emph{investigating} an alert or incident, not replacing that baseline monitoring.

\section[JFR (JDK Flight Recorder)]{JFR (JDK Flight Recorder)}
\label{h:12:jfr-jdk-flight-recorder}
\textbf{JFR} is a low-overhead (typically \textless{}1-2\%) profiling and event-recording engine built into the JDK itself — no external agent needed. It records hundreds of event types: object allocations (sampled), GC pauses, thread state changes, lock contention, exceptions, class loading, and more, all with timestamps, into a compact \code{.\allowbreak{}jfr} binary file you analyze after the fact (or stream live).

\begin{codebox}
\begin{Verbatim}[commandchars=\\\{\}]
\PY{c+c1}{\PYZsh{} Start a 60\PYZhy{}second recording on a running process, using the built\PYZhy{}in}
\PY{c+c1}{\PYZsh{} \PYZdq{}profile\PYZdq{} settings (more detail, still low overhead) and write it to rec.jfr}
jcmd\PY{+w}{ }\PYZlt{}pid\PYZgt{}\PY{+w}{ }JFR.start\PY{+w}{ }\PY{n+nv}{settings}\PY{o}{=}profile\PY{+w}{ }\PY{n+nv}{duration}\PY{o}{=}60s\PY{+w}{ }\PY{n+nv}{filename}\PY{o}{=}rec.jfr

\PY{c+c1}{\PYZsh{} Start recording from the moment the JVM launches (captures startup too)}
java\PY{+w}{ }\PYZhy{}XX:StartFlightRecording\PY{o}{=}\PY{n+nv}{filename}\PY{o}{=}startup.jfr,duration\PY{o}{=}60s\PY{+w}{ }\PYZhy{}jar\PY{+w}{ }app.jar

\PY{c+c1}{\PYZsh{} Check on / stop a running recording}
jcmd\PY{+w}{ }\PYZlt{}pid\PYZgt{}\PY{+w}{ }JFR.check
jcmd\PY{+w}{ }\PYZlt{}pid\PYZgt{}\PY{+w}{ }JFR.stop\PY{+w}{ }\PY{n+nv}{name}\PY{o}{=}\PY{l+m}{1}\PY{+w}{ }\PY{n+nv}{filename}\PY{o}{=}final.jfr
\end{Verbatim}
\end{codebox}

Because overhead is low enough for production, JFR is the tool of choice when you can’t reproduce an issue locally and need to capture “what actually happened” on the live system. The resulting \code{.\allowbreak{}jfr} file is best read with JMC.

Some of the event categories most relevant to memory work:

{\tablefont
\begin{xltabular}{\linewidth}{@{}L{0.571}L{1.429}@{}}
\toprule
\rowcolor{tablehdr}
\thd{Event category} & \thd{What it tells you} \\
\midrule
\endhead
\code{jdk.\allowbreak{}Object\allowbreak{}Allocation\allowbreak{}Sample} / \code{jdk.\allowbreak{}Object\allowbreak{}Allocation\allowbreak{}In\allowbreak{}New\allowbreak{}TLAB} & Which classes/call sites are allocating the most, sampled cheaply \\
\code{jdk.\allowbreak{}GarbageCollection} / \code{jdk.\allowbreak{}GCPhase*} & Every GC pause, its cause, duration, and phase breakdown \\
\code{jdk.\allowbreak{}OldObjectSample} & Objects that have survived a long time — a direct signal for leak-hunting without a full heap dump \\
\code{jdk.\allowbreak{}ThreadPark} / \code{jdk.\allowbreak{}JavaMonitorWait} & Lock contention and thread blocking, useful alongside \code{jstack} for deadlock/contention analysis \\
\code{jdk.\allowbreak{}ExecutionSample} & Periodic stack-trace sampling for CPU hot-spot / “hot method” analysis \\
\bottomrule
\end{xltabular}
}

The \code{jdk.\allowbreak{}OldObjectSample} event deserves a callout: it’s JFR’s built-in, low-overhead leak detector — it tracks a sample of objects that have lived unusually long and shows their allocation stack trace and retained size, often letting you spot a leak from a \emph{live} recording without ever needing to pull a full heap dump.

\section[JMC (JDK Mission Control)]{JMC (JDK Mission Control)}
\label{h:12:jmc-jdk-mission-control}
\textbf{JMC (JDK Mission Control)} is the GUI tool for \emph{analyzing} JFR recordings (and for live-monitoring a running JVM via JMX). Open a \code{.\allowbreak{}jfr} file in JMC and you get: a memory tab showing heap usage and GC pause timelines, an allocation tab ranking classes by bytes allocated (often the fastest way to spot “why is GC busy” — a class allocating way more than expected jumps out immediately), a “hot methods” view from sampled stack traces, and a thread/lock-contention view. JMC ships separately from the JDK (download from Adoptium/Eclipse) but reads recordings from any JDK 11+ vendor. For a code review context: if someone claims “this change reduces allocations,” JMC’s allocation-by-class view is the objective way to check.

\section[jcmd]{jcmd}
\label{h:12:jcmd}
\textbf{\code{jcmd}} is the modern, unified command-line tool for talking to a running JVM — it has mostly superseded the older single-purpose tools (\code{jstack}, older \code{jmap} use cases) because it exposes the same diagnostic commands through one consistent interface, and it works reliably even in JVMs where the legacy tools are flaky (e.g. minimal container base images). List a process’s PIDs and available commands first:

\begin{codebox}
\begin{Verbatim}[commandchars=\\\{\}]
\PY{c+c1}{\PYZsh{} List all running JVMs and their PIDs (like a Java\PYZhy{}aware \PYZsq{}ps\PYZsq{})}
jcmd\PY{+w}{ }\PYZhy{}l

\PY{c+c1}{\PYZsh{} List every diagnostic command a specific JVM supports}
jcmd\PY{+w}{ }\PYZlt{}pid\PYZgt{}\PY{+w}{ }\PY{n+nb}{help}
\end{Verbatim}
\end{codebox}

Common diagnostic commands:

\begin{codebox}
\begin{Verbatim}[commandchars=\\\{\}]
\PY{c+c1}{\PYZsh{} Heap summary: generation sizes, usage, GC algorithm in use}
jcmd\PY{+w}{ }\PYZlt{}pid\PYZgt{}\PY{+w}{ }GC.heap\PYZus{}info

\PY{c+c1}{\PYZsh{} Show every JVM flag and its current effective value}
jcmd\PY{+w}{ }\PYZlt{}pid\PYZgt{}\PY{+w}{ }VM.flags

\PY{c+c1}{\PYZsh{} Dump all thread stacks (equivalent to jstack, but built\PYZhy{}in)}
jcmd\PY{+w}{ }\PYZlt{}pid\PYZgt{}\PY{+w}{ }Thread.print

\PY{c+c1}{\PYZsh{} Trigger a heap dump without waiting for OOM}
jcmd\PY{+w}{ }\PYZlt{}pid\PYZgt{}\PY{+w}{ }GC.heap\PYZus{}dump\PY{+w}{ }/tmp/heap.hprof

\PY{c+c1}{\PYZsh{} Start a JFR profiling recording for 60 seconds}
jcmd\PY{+w}{ }\PYZlt{}pid\PYZgt{}\PY{+w}{ }JFR.start\PY{+w}{ }\PY{n+nv}{settings}\PY{o}{=}profile\PY{+w}{ }\PY{n+nv}{duration}\PY{o}{=}60s\PY{+w}{ }\PY{n+nv}{filename}\PY{o}{=}rec.jfr

\PY{c+c1}{\PYZsh{} Force a full GC right now (diagnostic use only — never in app code!)}
jcmd\PY{+w}{ }\PYZlt{}pid\PYZgt{}\PY{+w}{ }GC.run

\PY{c+c1}{\PYZsh{} Class\PYZhy{}level allocation histogram (like jmap \PYZhy{}histo:live, built into jcmd)}
jcmd\PY{+w}{ }\PYZlt{}pid\PYZgt{}\PY{+w}{ }GC.class\PYZus{}histogram

\PY{c+c1}{\PYZsh{} Native (off\PYZhy{}heap) memory breakdown, if NMT was enabled at startup}
jcmd\PY{+w}{ }\PYZlt{}pid\PYZgt{}\PY{+w}{ }VM.native\PYZus{}memory\PY{+w}{ }summary

\PY{c+c1}{\PYZsh{} Print system properties / general JVM version and uptime info}
jcmd\PY{+w}{ }\PYZlt{}pid\PYZgt{}\PY{+w}{ }VM.system\PYZus{}properties
jcmd\PY{+w}{ }\PYZlt{}pid\PYZgt{}\PY{+w}{ }VM.uptime
\end{Verbatim}
\end{codebox}

A quick reference of the commands most relevant to memory and thread debugging:

{\tablefont
\begin{xltabular}{\linewidth}{@{}L{0.686}L{1.314}@{}}
\toprule
\rowcolor{tablehdr}
\thd{Command} & \thd{What it does} \\
\midrule
\endhead
\code{GC.\allowbreak{}heap\_\allowbreak{}info} & Per-generation heap usage summary \\
\code{GC.\allowbreak{}heap\_\allowbreak{}dump~\allowbreak{}\textless{}\allowbreak{}path\textgreater{}} & Write a full heap dump (\code{.\allowbreak{}hprof}) \\
\code{GC.\allowbreak{}class\_\allowbreak{}histogram} & Live-object counts and bytes by class \\
\code{GC.\allowbreak{}run} & Request a full GC (diagnostic use only) \\
\code{VM.\allowbreak{}flags} & Effective value of every JVM flag \\
\code{VM.\allowbreak{}native\_\allowbreak{}memory~\allowbreak{}summary} & Off-heap memory breakdown (requires NMT enabled) \\
\code{Thread.\allowbreak{}print} & Full thread dump, including deadlock detection \\
\code{JFR.\allowbreak{}start} / \code{JFR.\allowbreak{}stop} / \code{JFR.\allowbreak{}check} & Control a Flight Recorder recording \\
\bottomrule
\end{xltabular}
}

\code{GC.\allowbreak{}heap\_\allowbreak{}info} output looks like this — read the “used” vs “capacity” per generation to see how full each region is and whether the heap has room to grow:

\begin{codebox}
\begin{Verbatim}
garbage-first heap   total 1048576K, used 402384K [0x...]
  region size 1024K, 1024 total, 84 available
 Metaspace       used 45231K, committed 46080K, reserved 1114112K
\end{Verbatim}
\end{codebox}

\section[jmap]{jmap}
\label{h:12:jmap}
\textbf{\code{jmap}} is the older, more narrowly-focused memory tool. Its two most useful invocations:

\begin{codebox}
\begin{Verbatim}[commandchars=\\\{\}]
\PY{c+c1}{\PYZsh{} Live\PYZhy{}object histogram: class name, instance count, total bytes — sorted by size.}
\PY{c+c1}{\PYZsh{} \PYZdq{}:live\PYZdq{} forces a GC first so only reachable objects are counted (skip it and you}
\PY{c+c1}{\PYZsh{} also count garbage waiting to be collected, which inflates the numbers).}
jmap\PY{+w}{ }\PYZhy{}histo:live\PY{+w}{ }\PYZlt{}pid\PYZgt{}

\PY{c+c1}{\PYZsh{} Full heap dump in the standard .hprof binary format, for MAT / JMC / VisualVM}
jmap\PY{+w}{ }\PYZhy{}dump:live,format\PY{o}{=}b,file\PY{o}{=}heap.hprof\PY{+w}{ }\PYZlt{}pid\PYZgt{}
\end{Verbatim}
\end{codebox}

Histogram output looks like:

\begin{codebox}
\begin{Verbatim}
 num     #instances         #bytes  class name
----------------------------------------------
   1:       842,113      67,369,040  [C                 (char[])
   2:       412,904      23,942,432  java.lang.String
   3:        88,213      21,171,120  com.example.model.OrderLine
\end{Verbatim}
\end{codebox}

Read it top-to-bottom by bytes: if \code{OrderLine} has far more instances than the number of orders you’d expect to have in memory, that’s your leak candidate — cross-reference with a heap dump’s dominator tree to find \emph{why} they’re still reachable.

\textbf{\code{jmap} is largely superseded by \code{jcmd}.} \code{jcmd~\allowbreak{}\textless{}\allowbreak{}pid\textgreater{}\allowbreak{}~\allowbreak{}GC.\allowbreak{}heap\_\allowbreak{}dump} and \code{jcmd~\allowbreak{}\textless{}\allowbreak{}pid\textgreater{}\allowbreak{}~\allowbreak{}GC.\allowbreak{}class\_\allowbreak{}histogram} do the same jobs, plus \code{jcmd} is maintained more actively and is more reliable across JVM versions and restricted environments. New tooling knowledge should default to \code{jcmd}; know \code{jmap}’s syntax mainly because it still appears in older runbooks, documentation, and — often — interview questions.

\section[jstack]{jstack}
\label{h:12:jstack}
\textbf{\code{jstack}} dumps the stack trace of every thread in a running JVM — essential for diagnosing hangs, deadlocks, and high-CPU-but-stuck-somewhere issues.

\begin{codebox}
\begin{Verbatim}[commandchars=\\\{\}]
jstack\PY{+w}{ }\PYZlt{}pid\PYZgt{}
\PY{c+c1}{\PYZsh{} or, if the process is unresponsive to normal signals:}
jstack\PY{+w}{ }\PYZhy{}F\PY{+w}{ }\PYZlt{}pid\PYZgt{}
\end{Verbatim}
\end{codebox}

Each thread entry shows its name, state, and a Java-level stack trace, e.g.:

\begin{codebox}
\begin{Verbatim}
"pool-1-thread-3" #15 prio=5 os_prio=0 tid=0x... nid=0x1a2b waiting on condition [0x...]
   java.lang.Thread.State: WAITING (parking)
        at jdk.internal.misc.Unsafe.park(Native Method)
        at java.util.concurrent.locks.AbstractQueuedSynchronizer.parkAndCheckInterrupt
        at com.example.OrderProcessor.awaitCompletion(OrderProcessor.java:88)
\end{Verbatim}
\end{codebox}

Taking two or three \code{jstack} dumps a few seconds apart and comparing them is a classic trick: a thread stuck at the \emph{exact same line} across all of them is a strong signal of a genuine hang (deadlock, infinite loop, or blocked I/O), whereas a thread that’s at a \emph{different} line each time is simply doing normal work.

\textbf{Thread states you’ll see in a dump:}

{\tablefont
\begin{xltabular}{\linewidth}{@{}L{0.400}L{1.600}@{}}
\toprule
\rowcolor{tablehdr}
\thd{State} & \thd{Meaning} \\
\midrule
\endhead
\code{NEW} & Thread object created, \code{start()} not yet called \\
\code{RUNNABLE} & Executing, or ready and waiting for CPU — includes threads doing blocking I/O in native code \\
\code{BLOCKED} & Waiting to enter a \code{synchronized} block/method held by another thread \\
\code{WAITING} & Waiting indefinitely for another thread (\code{Object.\allowbreak{}wait()}, \code{Thread.\allowbreak{}join()}, \code{LockSupport.\allowbreak{}park()} with no timeout) \\
\code{TIMED\_\allowbreak{}WAITING} & Same as \code{WAITING} but with a timeout (\code{Thread.\allowbreak{}sleep()}, \code{wait(\allowbreak{}ms)}, \code{park(\allowbreak{}nanos)}) \\
\code{TERMINATED} & Run to completion \\
\bottomrule
\end{xltabular}
}

A thread pool full of \code{BLOCKED} threads all waiting on the \emph{same} lock, with one thread holding it and itself stuck elsewhere (often in \code{WAITING} on I/O or another lock), is the classic “everything is stuck behind one slow/stuck thread” pattern — look for the single thread that’s the odd one out, not the many that are merely blocked on it.

\textbf{\code{jstack} vs \code{jcmd~\allowbreak{}Thread.\allowbreak{}print}.} They produce essentially the same information; \code{jcmd~\allowbreak{}\textless{}\allowbreak{}pid\textgreater{}\allowbreak{}~\allowbreak{}Thread.\allowbreak{}print} is the more modern, actively maintained path (same rationale as \code{jmap} vs \code{jcmd} for heap dumps), and it works over the same attach mechanism, so prefer it when starting fresh, while recognizing \code{jstack} output in existing runbooks and tickets.

\section[Worked Walkthrough: Debugging an OutOfMemoryError]{Worked Walkthrough: Debugging an OutOfMemoryError}
\label{h:12:worked-walkthrough-debugging-an-outofmemoryerror}
\textbf{Scenario:} a service occasionally crashes with \code{java.\allowbreak{}lang.\allowbreak{}Out\allowbreak{}Of\allowbreak{}Memory\allowbreak{}Error:\allowbreak{}~\allowbreak{}Java~\allowbreak{}heap~\allowbreak{}space} in production.

\textbf{Step 1 — make sure the next crash captures evidence.} Add this flag (ideally it’s \emph{always} on in production, before the incident even happens):

\begin{codebox}
\begin{Verbatim}[commandchars=\\\{\}]
java\PY{+w}{ }\PYZhy{}XX:+HeapDumpOnOutOfMemoryError\PY{+w}{ }\PYZhy{}XX:HeapDumpPath\PY{o}{=}/var/dumps/\PY{+w}{ }\PYZhy{}jar\PY{+w}{ }app.jar
\end{Verbatim}
\end{codebox}

\textbf{Step 2 — while the process is still alive but suspected of leaking, capture a live heap dump too} (don’t wait for the crash if you can catch it early):

\begin{codebox}
\begin{Verbatim}[commandchars=\\\{\}]
jcmd\PY{+w}{ }\PYZlt{}pid\PYZgt{}\PY{+w}{ }GC.heap\PYZus{}dump\PY{+w}{ }/var/dumps/pre\PYZhy{}crash.hprof
\end{Verbatim}
\end{codebox}

\textbf{Step 3 — check the GC log for the shape of the problem.} If GC logging was already enabled (\code{-\allowbreak{}Xlog:\allowbreak{}gc*:\allowbreak{}file=\allowbreak{}gc.\allowbreak{}log:...}), look at the tail of the log right before the crash:

\begin{codebox}
\begin{Verbatim}
[gc] GC(430) Pause Full (G1 Evacuation Pause) 3980M->3971M(4096M) 812.442ms
[gc] GC(431) Pause Full (G1 Evacuation Pause) 3985M->3978M(4096M) 799.113ms
[gc] GC(432) Pause Full (G1 Evacuation Pause) 3990M->3982M(4096M) 820.977ms
\end{Verbatim}
\end{codebox}

This pattern — heap usage barely drops after each \emph{full} GC, and full GCs are happening back-to-back — is the signature of a real leak (versus normal pressure, where after-GC usage would drop substantially).

\textbf{Step 4 — open the heap dump in Eclipse MAT (or JMC).} Run MAT’s built-in “Leak Suspects” report, or manually open the dominator tree, sort by retained size, and look at the top entries.

\textbf{Step 5 — read the dominator tree result}, e.g.:

\begin{codebox}
\begin{Verbatim}
Class Name                                    | Shallow Heap | Retained Heap
com.example.cache.ReportCache @ 0x7f2a...     |          48  |   3,842,199,120
  -> java.util.HashMap$Node[] entries          |     524,288 |   3,842,190,000
\end{Verbatim}
\end{codebox}

The \code{ReportCache} instance alone retains \textasciitilde{}3.8 GB — almost the entire heap. That’s the leak. Trace its “path to GC roots” in MAT to see it’s a \code{static} field with no eviction policy.

\textbf{Step 6 — fix and verify.} Add a bounded eviction policy (e.g. switch to Caffeine with \code{maximumSize}), redeploy, and confirm with a fresh GC log that full-GC after-usage now drops back down close to the old-gen baseline instead of climbing indefinitely.

\section[Worked Walkthrough: Debugging a Deadlock with jstack]{Worked Walkthrough: Debugging a Deadlock with jstack}
\label{h:12:worked-walkthrough-debugging-a-deadlock-with-jstack}
\textbf{Scenario:} a service’s health check starts timing out; CPU is idle (not spinning), suggesting threads are stuck waiting, not busy-looping.

\textbf{Step 1 — take a thread dump:}

\begin{codebox}
\begin{Verbatim}[commandchars=\\\{\}]
jstack\PY{+w}{ }\PYZlt{}pid\PYZgt{}\PY{+w}{ }\PYZgt{}\PY{+w}{ }threads.txt
\PY{c+c1}{\PYZsh{} or: jcmd \PYZlt{}pid\PYZgt{} Thread.print \PYZgt{} threads.txt}
\end{Verbatim}
\end{codebox}

\textbf{Step 2 — search for the deadlock report.} \code{jstack} detects classic lock-ordering deadlocks automatically and prints a dedicated section:

\begin{codebox}
\begin{Verbatim}
Found one Java-level deadlock:
=============================
"Thread-A":
  waiting to lock monitor 0x00007f... (object 0x000000070a, a java.lang.Object),
  which is held by "Thread-B"
"Thread-B":
  waiting to lock monitor 0x00007f... (object 0x000000070b, a java.lang.Object),
  which is held by "Thread-A"

Java stack information for the threads listed above:
===================================================
"Thread-A":
        at com.example.TransferService.debit(TransferService.java:42)
        - waiting to lock <0x000000070a> (a java.lang.Object)
        - locked <0x000000070b> (a java.lang.Object)
"Thread-B":
        at com.example.TransferService.credit(TransferService.java:58)
        - waiting to lock <0x000000070b> (a java.lang.Object)
        - locked <0x000000070a> (a java.lang.Object)
\end{Verbatim}
\end{codebox}

\textbf{Step 3 — read it as a cycle.} Thread-A holds lock B’s target and wants lock A’s target; Thread-B holds the opposite. Neither can proceed — classic circular wait, the textbook deadlock condition.

\textbf{Step 4 — trace it back to the code.} Both \code{debit()} and \code{credit()} lock two account objects in \emph{opposite order} depending on which account is the caller vs. the target:

\begin{codebox}
\begin{Verbatim}[commandchars=\\\{\}]
\PY{c+c1}{// Wrong — lock order depends on argument order, so A\PYZhy{}\PYZgt{}B and B\PYZhy{}\PYZgt{}A both happen}
\PY{k+kd}{synchronized}\PY{+w}{ }\PY{p}{(}\PY{n}{fromAccount}\PY{p}{)}\PY{+w}{ }\PY{p}{\PYZob{}}
\PY{+w}{    }\PY{k+kd}{synchronized}\PY{+w}{ }\PY{p}{(}\PY{n}{toAccount}\PY{p}{)}\PY{+w}{ }\PY{p}{\PYZob{}}
\PY{+w}{        }\PY{c+c1}{// transfer logic}
\PY{+w}{    }\PY{p}{\PYZcb{}}
\PY{p}{\PYZcb{}}
\end{Verbatim}
\end{codebox}

\textbf{Step 5 — fix with a consistent lock order} (e.g., always lock the account with the smaller ID first), or replace with a higher-level concurrency utility that avoids manual nested locking entirely (\code{java.\allowbreak{}util.\allowbreak{}concurrent} structures, or \code{ReentrantLock.\allowbreak{}tryLock} with a timeout so the app can recover instead of hanging forever). See Chapter 13 for lock-ordering strategies in depth.

\textbf{Step 6 — verify.} Under the same load test that reproduced the hang, confirm the health check stays responsive and no deadlock section appears in a follow-up \code{jstack} dump.

\section[Common Code-Review Interview Pitfalls]{Common Code-Review Interview Pitfalls}
\label{h:12:common-code-review-interview-pitfalls}
\begin{enumerate}
\item 
\textbf{Calling \code{System.\allowbreak{}gc()} in application code.} Why it matters: it’s only a hint the JVM may ignore, and on some collectors it can trigger a genuinely expensive full GC pause in production for no real benefit.

\begin{codebox}
\begin{Verbatim}[commandchars=\\\{\}]
\PY{c+c1}{// Before}
\PY{n}{cache}\PY{p}{.}\PY{n+na}{clear}\PY{p}{(}\PY{p}{)}\PY{p}{;}
\PY{n}{System}\PY{p}{.}\PY{n+na}{gc}\PY{p}{(}\PY{p}{)}\PY{p}{;}\PY{+w}{ }\PY{c+c1}{// \PYZdq{}just to be safe\PYZdq{} — actually just a pause risk}
\PY{c+c1}{// After}
\PY{n}{cache}\PY{p}{.}\PY{n+na}{clear}\PY{p}{(}\PY{p}{)}\PY{p}{;}\PY{+w}{ }\PY{c+c1}{// trust the collector; profile if there\PYZsq{}s an actual problem}
\end{Verbatim}
\end{codebox}

\item 
\textbf{Using \code{finalize()} for cleanup.} Why it matters: it’s deprecated for removal, runs at an unpredictable time (or never), and can silently swallow exceptions.

\begin{codebox}
\begin{Verbatim}[commandchars=\\\{\}]
\PY{c+c1}{// Before}
\PY{n+nd}{@Override}\PY{+w}{ }\PY{k+kd}{protected}\PY{+w}{ }\PY{k+kt}{void}\PY{+w}{ }\PY{n+nf}{finalize}\PY{p}{(}\PY{p}{)}\PY{+w}{ }\PY{p}{\PYZob{}}\PY{+w}{ }\PY{n}{socket}\PY{p}{.}\PY{n+na}{close}\PY{p}{(}\PY{p}{)}\PY{p}{;}\PY{+w}{ }\PY{p}{\PYZcb{}}
\PY{c+c1}{// After}
\PY{k}{try}\PY{+w}{ }\PY{p}{(}\PY{k+kd}{var}\PY{+w}{ }\PY{n}{socket}\PY{+w}{ }\PY{o}{=}\PY{+w}{ }\PY{n}{openSocket}\PY{p}{(}\PY{p}{)}\PY{p}{)}\PY{+w}{ }\PY{p}{\PYZob{}}\PY{+w}{ }\PY{c+cm}{/* use it */}\PY{+w}{ }\PY{p}{\PYZcb{}}\PY{+w}{ }\PY{c+c1}{// deterministic close}
\end{Verbatim}
\end{codebox}

\item 
\textbf{Not calling \code{ThreadLocal.\allowbreak{}remove()} in pooled-thread code.} Why it matters: pooled threads outlive any single task, so a forgotten value (and its whole reference graph) leaks for the pool’s lifetime and can even leak data between unrelated requests.

\begin{codebox}
\begin{Verbatim}[commandchars=\\\{\}]
\PY{c+c1}{// Before}
\PY{n}{CONTEXT}\PY{p}{.}\PY{n+na}{set}\PY{p}{(}\PY{n}{userContext}\PY{p}{)}\PY{p}{;}\PY{+w}{ }\PY{n}{process}\PY{p}{(}\PY{n}{request}\PY{p}{)}\PY{p}{;}
\PY{c+c1}{// After}
\PY{n}{CONTEXT}\PY{p}{.}\PY{n+na}{set}\PY{p}{(}\PY{n}{userContext}\PY{p}{)}\PY{p}{;}
\PY{k}{try}\PY{+w}{ }\PY{p}{\PYZob{}}\PY{+w}{ }\PY{n}{process}\PY{p}{(}\PY{n}{request}\PY{p}{)}\PY{p}{;}\PY{+w}{ }\PY{p}{\PYZcb{}}\PY{+w}{ }\PY{k}{finally}\PY{+w}{ }\PY{p}{\PYZob{}}\PY{+w}{ }\PY{n}{CONTEXT}\PY{p}{.}\PY{n+na}{remove}\PY{p}{(}\PY{p}{)}\PY{p}{;}\PY{+w}{ }\PY{p}{\PYZcb{}}
\end{Verbatim}
\end{codebox}

\item 
\textbf{Using an unbounded \code{HashMap} as a cache.} Why it matters: with no eviction policy, it grows without limit and is a slow-motion \code{OutOfMemoryError} waiting to happen.

\begin{codebox}
\begin{Verbatim}[commandchars=\\\{\}]
\PY{c+c1}{// Before}
\PY{k+kd}{private}\PY{+w}{ }\PY{k+kd}{final}\PY{+w}{ }\PY{n}{Map}\PY{o}{\PYZlt{}}\PY{n}{Long}\PY{p}{,}\PY{+w}{ }\PY{n}{Report}\PY{o}{\PYZgt{}}\PY{+w}{ }\PY{n}{cache}\PY{+w}{ }\PY{o}{=}\PY{+w}{ }\PY{k}{new}\PY{+w}{ }\PY{n}{HashMap}\PY{o}{\PYZlt{}}\PY{o}{\PYZgt{}}\PY{p}{(}\PY{p}{)}\PY{p}{;}
\PY{c+c1}{// After}
\PY{k+kd}{private}\PY{+w}{ }\PY{k+kd}{final}\PY{+w}{ }\PY{n}{Cache}\PY{o}{\PYZlt{}}\PY{n}{Long}\PY{p}{,}\PY{+w}{ }\PY{n}{Report}\PY{o}{\PYZgt{}}\PY{+w}{ }\PY{n}{cache}\PY{+w}{ }\PY{o}{=}
\PY{+w}{        }\PY{n}{Caffeine}\PY{p}{.}\PY{n+na}{newBuilder}\PY{p}{(}\PY{p}{)}\PY{p}{.}\PY{n+na}{maximumSize}\PY{p}{(}\PY{l+m+mi}{10\PYZus{}000}\PY{p}{)}\PY{p}{.}\PY{n+na}{build}\PY{p}{(}\PY{p}{)}\PY{p}{;}
\end{Verbatim}
\end{codebox}

\item 
\textbf{Registering listeners without a matching unregister.} Why it matters: the listener (and everything it closes over) stays reachable from the event source forever, even after the “subscriber” logically should be gone.

\begin{codebox}
\begin{Verbatim}[commandchars=\\\{\}]
\PY{c+c1}{// Before}
\PY{n}{eventBus}\PY{p}{.}\PY{n+na}{subscribe}\PY{p}{(}\PY{n}{listener}\PY{p}{)}\PY{p}{;}\PY{+w}{ }\PY{c+c1}{// no corresponding unsubscribe anywhere}
\PY{c+c1}{// After}
\PY{n}{eventBus}\PY{p}{.}\PY{n+na}{subscribe}\PY{p}{(}\PY{n}{listener}\PY{p}{)}\PY{p}{;}
\PY{c+c1}{// ... and on shutdown/cleanup:}
\PY{n}{eventBus}\PY{p}{.}\PY{n+na}{unsubscribe}\PY{p}{(}\PY{n}{listener}\PY{p}{)}\PY{p}{;}
\end{Verbatim}
\end{codebox}

\item 
\textbf{Making a class a non-static inner class when it doesn’t need the outer instance.} Why it matters: every instance silently holds a hidden reference to its enclosing object, which can keep a large outer object (and its state) alive far longer than intended.

\begin{codebox}
\begin{Verbatim}[commandchars=\\\{\}]
\PY{c+c1}{// Before}
\PY{k+kd}{class} \PY{n+nc}{Listener}\PY{+w}{ }\PY{p}{\PYZob{}}\PY{+w}{ }\PY{c+cm}{/* implicitly holds Outer.this, unused */}\PY{+w}{ }\PY{p}{\PYZcb{}}
\PY{c+c1}{// After}
\PY{k+kd}{static}\PY{+w}{ }\PY{k+kd}{class} \PY{n+nc}{Listener}\PY{+w}{ }\PY{p}{\PYZob{}}\PY{+w}{ }\PY{c+cm}{/* no hidden outer reference */}\PY{+w}{ }\PY{p}{\PYZcb{}}
\end{Verbatim}
\end{codebox}

\item 
\textbf{Assuming \code{-\allowbreak{}Xmx} alone controls container memory usage.} Why it matters: metaspace, thread stacks, direct buffers, and JIT code cache are all off-heap and can push total memory past the container limit even when the heap itself looks fine, causing an OOM-kill with no Java-level stack trace.

\begin{codebox}
\begin{Verbatim}[commandchars=\\\{\}]
\PY{c+c1}{\PYZsh{} Before: heap\PYZhy{}only thinking}
java\PY{+w}{ }\PYZhy{}Xmx3800m\PY{+w}{ }\PYZhy{}jar\PY{+w}{ }app.jar\PY{+w}{   }\PY{c+c1}{\PYZsh{} in a 4 GB container — no room for anything else}
\PY{c+c1}{\PYZsh{} After: leave headroom, or size as a percentage of the container limit}
java\PY{+w}{ }\PYZhy{}XX:MaxRAMPercentage\PY{o}{=}\PY{l+m}{70}.0\PY{+w}{ }\PYZhy{}jar\PY{+w}{ }app.jar
\end{Verbatim}
\end{codebox}

\item 
\textbf{Setting \code{-\allowbreak{}Xms} far below \code{-\allowbreak{}Xmx} in a server/container workload.} Why it matters: the JVM has to repeatedly resize the heap under load, causing extra pauses and unpredictable memory footprint — exactly the opposite of what you want in a container with a hard memory limit.

\begin{codebox}
\begin{Verbatim}[commandchars=\\\{\}]
\PY{c+c1}{\PYZsh{} Before}
java\PY{+w}{ }\PYZhy{}Xms256m\PY{+w}{ }\PYZhy{}Xmx4g\PY{+w}{ }\PYZhy{}jar\PY{+w}{ }app.jar
\PY{c+c1}{\PYZsh{} After}
java\PY{+w}{ }\PYZhy{}Xms4g\PY{+w}{ }\PYZhy{}Xmx4g\PY{+w}{ }\PYZhy{}jar\PY{+w}{ }app.jar
\end{Verbatim}
\end{codebox}

\item 
\textbf{Leaving metaspace uncapped in a service that dynamically loads classes (or gets redeployed).} Why it matters: a classloader leak silently grows metaspace with no upper bound until the JVM (or container) runs out of memory in a confusing way.

\begin{codebox}
\begin{Verbatim}[commandchars=\\\{\}]
\PY{c+c1}{\PYZsh{} Before: no cap at all}
java\PY{+w}{ }\PYZhy{}jar\PY{+w}{ }app.jar
\PY{c+c1}{\PYZsh{} After: explicit cap turns a silent leak into a clear, diagnosable OOM}
java\PY{+w}{ }\PYZhy{}XX:MaxMetaspaceSize\PY{o}{=}512m\PY{+w}{ }\PYZhy{}jar\PY{+w}{ }app.jar
\end{Verbatim}
\end{codebox}

\item 
\textbf{Writing code that depends on escape analysis for correctness or guaranteed performance.} Why it matters: it’s a best-effort, JIT-tier-dependent optimization with no language-level guarantee — a minor refactor can silently disable it with no warning.

\begin{codebox}
\begin{Verbatim}[commandchars=\\\{\}]
\PY{c+c1}{// Before — relying on \PYZdq{}the JIT will optimize this away\PYZdq{} as a design assumption}
\PY{n}{Point}\PY{+w}{ }\PY{n}{p}\PY{+w}{ }\PY{o}{=}\PY{+w}{ }\PY{k}{new}\PY{+w}{ }\PY{n}{Point}\PY{p}{(}\PY{n}{x}\PY{p}{,}\PY{+w}{ }\PY{n}{y}\PY{p}{)}\PY{p}{;}\PY{+w}{ }\PY{c+c1}{// assumed to never really allocate}
\PY{c+c1}{// After — write clear code; treat escape analysis as a bonus, not a plan}
\PY{k+kt}{int}\PY{+w}{ }\PY{n}{dx}\PY{+w}{ }\PY{o}{=}\PY{+w}{ }\PY{n}{x2}\PY{+w}{ }\PY{o}{\PYZhy{}}\PY{+w}{ }\PY{n}{x1}\PY{p}{,}\PY{+w}{ }\PY{n}{dy}\PY{+w}{ }\PY{o}{=}\PY{+w}{ }\PY{n}{y2}\PY{+w}{ }\PY{o}{\PYZhy{}}\PY{+w}{ }\PY{n}{y1}\PY{p}{;}\PY{+w}{ }\PY{c+c1}{// if it matters, measure, don\PYZsq{}t assume}
\end{Verbatim}
\end{codebox}

\item 
\textbf{Reaching for \code{jmap~\allowbreak{}-\allowbreak{}dump} under time pressure without knowing \code{jcmd} exists.} Why it matters: \code{jcmd} is the actively maintained, more reliable tool for the same job, and knowing it signals up-to-date operational knowledge in a review.

\begin{codebox}
\begin{Verbatim}[commandchars=\\\{\}]
\PY{c+c1}{\PYZsh{} Older habit}
jmap\PY{+w}{ }\PYZhy{}dump:live,format\PY{o}{=}b,file\PY{o}{=}heap.hprof\PY{+w}{ }\PYZlt{}pid\PYZgt{}
\PY{c+c1}{\PYZsh{} Preferred}
jcmd\PY{+w}{ }\PYZlt{}pid\PYZgt{}\PY{+w}{ }GC.heap\PYZus{}dump\PY{+w}{ }/tmp/heap.hprof
\end{Verbatim}
\end{codebox}

\item 
\textbf{Choosing ZGC/Shenandoah “because it’s the newest” without a pause-time requirement.} Why it matters: both trade extra memory and CPU overhead for lower pauses — paying that cost with no latency SLA to justify it is a wasted resource, and G1’s balance is right for most services.

\begin{codebox}
\begin{Verbatim}
# Before: reflexively latency-optimizing an offline batch job
-XX:+UseZGC
# After: throughput-first collector for a job nobody is waiting on interactively
-XX:+UseParallelGC
\end{Verbatim}
\end{codebox}

\item 
\textbf{Treating a single heap-usage snapshot as proof of (or against) a leak.} Why it matters: normal caching and generational promotion both look like “heap usage went up” in one snapshot; only a trend across multiple full-GC cycles (usage not dropping back down) distinguishes a leak from healthy behavior.

\begin{codebox}
\begin{Verbatim}
# Before: "heap usage is 80% full, must be a leak!"
# After: compare after-full-GC usage across several cycles over time;
# a leak shows a rising floor, not just a high peak.
\end{Verbatim}
\end{codebox}

\item 
\textbf{Manually locking two objects in an order that depends on caller-supplied argument order.} Why it matters: it creates a classic deadlock the moment two threads call the method with the arguments swapped — exactly the bug \code{jstack}’s deadlock detector exists to catch.

\begin{codebox}
\begin{Verbatim}[commandchars=\\\{\}]
\PY{c+c1}{// Before}
\PY{k+kd}{synchronized}\PY{+w}{ }\PY{p}{(}\PY{n}{fromAccount}\PY{p}{)}\PY{+w}{ }\PY{p}{\PYZob{}}\PY{+w}{ }\PY{k+kd}{synchronized}\PY{+w}{ }\PY{p}{(}\PY{n}{toAccount}\PY{p}{)}\PY{+w}{ }\PY{p}{\PYZob{}}\PY{+w}{ }\PY{c+cm}{/* ... */}\PY{+w}{ }\PY{p}{\PYZcb{}}\PY{+w}{ }\PY{p}{\PYZcb{}}
\PY{c+c1}{// After: always lock in a fixed, consistent order regardless of arguments}
\PY{n}{Account}\PY{+w}{ }\PY{n}{first}\PY{+w}{ }\PY{o}{=}\PY{+w}{ }\PY{n}{fromAccount}\PY{p}{.}\PY{n+na}{id}\PY{p}{(}\PY{p}{)}\PY{+w}{ }\PY{o}{\PYZlt{}}\PY{+w}{ }\PY{n}{toAccount}\PY{p}{.}\PY{n+na}{id}\PY{p}{(}\PY{p}{)}\PY{+w}{ }\PY{o}{?}\PY{+w}{ }\PY{n}{fromAccount}\PY{+w}{ }\PY{p}{:}\PY{+w}{ }\PY{n}{toAccount}\PY{p}{;}
\PY{n}{Account}\PY{+w}{ }\PY{n}{second}\PY{+w}{ }\PY{o}{=}\PY{+w}{ }\PY{n}{fromAccount}\PY{p}{.}\PY{n+na}{id}\PY{p}{(}\PY{p}{)}\PY{+w}{ }\PY{o}{\PYZlt{}}\PY{+w}{ }\PY{n}{toAccount}\PY{p}{.}\PY{n+na}{id}\PY{p}{(}\PY{p}{)}\PY{+w}{ }\PY{o}{?}\PY{+w}{ }\PY{n}{toAccount}\PY{+w}{ }\PY{p}{:}\PY{+w}{ }\PY{n}{fromAccount}\PY{p}{;}
\PY{k+kd}{synchronized}\PY{+w}{ }\PY{p}{(}\PY{n}{first}\PY{p}{)}\PY{+w}{ }\PY{p}{\PYZob{}}\PY{+w}{ }\PY{k+kd}{synchronized}\PY{+w}{ }\PY{p}{(}\PY{n}{second}\PY{p}{)}\PY{+w}{ }\PY{p}{\PYZob{}}\PY{+w}{ }\PY{c+cm}{/* ... */}\PY{+w}{ }\PY{p}{\PYZcb{}}\PY{+w}{ }\PY{p}{\PYZcb{}}
\end{Verbatim}
\end{codebox}

\item 
\textbf{Not enabling \code{-\allowbreak{}XX:+\allowbreak{}Heap\allowbreak{}Dump\allowbreak{}On\allowbreak{}Out\allowbreak{}Of\allowbreak{}Memory\allowbreak{}Error} in production.} Why it matters: without it, an \code{OutOfMemoryError} crash leaves no artifact to analyze — the exact evidence needed to diagnose the leak disappears along with the crashed process.

\begin{codebox}
\begin{Verbatim}[commandchars=\\\{\}]
\PY{c+c1}{\PYZsh{} Before}
java\PY{+w}{ }\PYZhy{}jar\PY{+w}{ }app.jar
\PY{c+c1}{\PYZsh{} After}
java\PY{+w}{ }\PYZhy{}XX:+HeapDumpOnOutOfMemoryError\PY{+w}{ }\PYZhy{}XX:HeapDumpPath\PY{o}{=}/var/dumps/\PY{+w}{ }\PYZhy{}jar\PY{+w}{ }app.jar
\end{Verbatim}
\end{codebox}

\end{enumerate}


\setcounter{chapter}{12}
\chapter{Concurrency (Core)}
\label{chap:13}
\label{h:13:13-concurrency-core}
Concurrency is what happens when more than one thread touches your program’s state at the same time. Get it wrong and bugs show up rarely, non-deterministically, and usually in production under load — which makes concurrency one of the highest-value topics in a code-review interview. This chapter covers the classic “platform threads” toolkit: threads themselves, locks, \code{volatile}, atomics, \code{ThreadLocal}, and the memory-model rules (happens-before) that explain \emph{why} all of it works. Examples target Java 21+; virtual threads and structured concurrency get their own treatment in Chapter 15.

\section[Threads]{Threads}
\label{h:13:threads}
A \textbf{thread} is an independent path of execution inside a program. A single Java process can run many threads at once, each with its own call stack, but all of them share the same heap memory — the same objects, the same static fields. That sharing is exactly what makes concurrency both powerful and dangerous.

The classic (pre-virtual-thread) way to create one is \code{Thread} backed directly by an OS thread — a “platform thread.” You can subclass \code{Thread} and override \code{run()}, or (preferred) pass it a \code{Runnable}:

\begin{codebox}
\begin{Verbatim}[commandchars=\\\{\}]
\PY{k+kd}{public}\PY{+w}{ }\PY{k+kd}{class} \PY{n+nc}{ThreadBasics}\PY{+w}{ }\PY{p}{\PYZob{}}
\PY{+w}{    }\PY{k+kd}{public}\PY{+w}{ }\PY{k+kd}{static}\PY{+w}{ }\PY{k+kt}{void}\PY{+w}{ }\PY{n+nf}{main}\PY{p}{(}\PY{n}{String}\PY{o}{[}\PY{o}{]}\PY{+w}{ }\PY{n}{args}\PY{p}{)}\PY{+w}{ }\PY{k+kd}{throws}\PY{+w}{ }\PY{n}{InterruptedException}\PY{+w}{ }\PY{p}{\PYZob{}}
\PY{+w}{        }\PY{n}{Thread}\PY{+w}{ }\PY{n}{worker}\PY{+w}{ }\PY{o}{=}\PY{+w}{ }\PY{k}{new}\PY{+w}{ }\PY{n}{Thread}\PY{p}{(}\PY{p}{(}\PY{p}{)}\PY{+w}{ }\PY{o}{\PYZhy{}}\PY{o}{\PYZgt{}}\PY{+w}{ }\PY{p}{\PYZob{}}
\PY{+w}{            }\PY{n}{System}\PY{p}{.}\PY{n+na}{out}\PY{p}{.}\PY{n+na}{println}\PY{p}{(}\PY{l+s}{\PYZdq{}}\PY{l+s}{Running on: }\PY{l+s}{\PYZdq{}}\PY{+w}{ }\PY{o}{+}\PY{+w}{ }\PY{n}{Thread}\PY{p}{.}\PY{n+na}{currentThread}\PY{p}{(}\PY{p}{)}\PY{p}{.}\PY{n+na}{getName}\PY{p}{(}\PY{p}{)}\PY{p}{)}\PY{p}{;}
\PY{+w}{        }\PY{p}{\PYZcb{}}\PY{p}{,}\PY{+w}{ }\PY{l+s}{\PYZdq{}}\PY{l+s}{worker\PYZhy{}1}\PY{l+s}{\PYZdq{}}\PY{p}{)}\PY{p}{;}\PY{+w}{ }\PY{c+c1}{// name the thread — huge help in thread dumps and logs}

\PY{+w}{        }\PY{n}{worker}\PY{p}{.}\PY{n+na}{start}\PY{p}{(}\PY{p}{)}\PY{p}{;}\PY{+w}{ }\PY{c+c1}{// schedules the thread to run; never call run() directly}
\PY{+w}{        }\PY{n}{worker}\PY{p}{.}\PY{n+na}{join}\PY{p}{(}\PY{p}{)}\PY{p}{;}\PY{+w}{  }\PY{c+c1}{// wait for it to finish before continuing}
\PY{+w}{        }\PY{n}{System}\PY{p}{.}\PY{n+na}{out}\PY{p}{.}\PY{n+na}{println}\PY{p}{(}\PY{l+s}{\PYZdq{}}\PY{l+s}{Done}\PY{l+s}{\PYZdq{}}\PY{p}{)}\PY{p}{;}
\PY{+w}{    }\PY{p}{\PYZcb{}}
\PY{p}{\PYZcb{}}
\end{Verbatim}
\end{codebox}

\textbf{\code{start()} vs \code{run()}.} Calling \code{start()} asks the JVM to create a new call stack and schedule execution. Calling \code{run()} directly just executes the code on the \emph{current} thread — no concurrency happens at all. This is a common interview trick question.

\textbf{Daemon vs user (non-daemon) threads.} A \textbf{user thread} keeps the JVM alive — the process won’t exit while any user thread is still running. A \textbf{daemon thread} is a “background helper” thread; the JVM will exit even if daemon threads are still running, and it won’t wait for them. Garbage collection and JIT compilation run on daemon threads internally.

\begin{codebox}
\begin{Verbatim}[commandchars=\\\{\}]
\PY{n}{Thread}\PY{+w}{ }\PY{n}{heartbeat}\PY{+w}{ }\PY{o}{=}\PY{+w}{ }\PY{k}{new}\PY{+w}{ }\PY{n}{Thread}\PY{p}{(}\PY{p}{(}\PY{p}{)}\PY{+w}{ }\PY{o}{\PYZhy{}}\PY{o}{\PYZgt{}}\PY{+w}{ }\PY{p}{\PYZob{}}
\PY{+w}{    }\PY{k}{while}\PY{+w}{ }\PY{p}{(}\PY{k+kc}{true}\PY{p}{)}\PY{+w}{ }\PY{p}{\PYZob{}}
\PY{+w}{        }\PY{n}{System}\PY{p}{.}\PY{n+na}{out}\PY{p}{.}\PY{n+na}{println}\PY{p}{(}\PY{l+s}{\PYZdq{}}\PY{l+s}{tick}\PY{l+s}{\PYZdq{}}\PY{p}{)}\PY{p}{;}
\PY{+w}{        }\PY{k}{try}\PY{+w}{ }\PY{p}{\PYZob{}}\PY{+w}{ }\PY{n}{Thread}\PY{p}{.}\PY{n+na}{sleep}\PY{p}{(}\PY{l+m+mi}{1000}\PY{p}{)}\PY{p}{;}\PY{+w}{ }\PY{p}{\PYZcb{}}\PY{+w}{ }\PY{k}{catch}\PY{+w}{ }\PY{p}{(}\PY{n}{InterruptedException}\PY{+w}{ }\PY{n}{e}\PY{p}{)}\PY{+w}{ }\PY{p}{\PYZob{}}\PY{+w}{ }\PY{k}{break}\PY{p}{;}\PY{+w}{ }\PY{p}{\PYZcb{}}
\PY{+w}{    }\PY{p}{\PYZcb{}}
\PY{p}{\PYZcb{}}\PY{p}{)}\PY{p}{;}
\PY{n}{heartbeat}\PY{p}{.}\PY{n+na}{setDaemon}\PY{p}{(}\PY{k+kc}{true}\PY{p}{)}\PY{p}{;}\PY{+w}{ }\PY{c+c1}{// must be called BEFORE start()}
\PY{n}{heartbeat}\PY{p}{.}\PY{n+na}{start}\PY{p}{(}\PY{p}{)}\PY{p}{;}
\PY{c+c1}{// If main() returns here, the JVM exits and heartbeat is killed mid\PYZhy{}loop —}
\PY{c+c1}{// no cleanup code in it is guaranteed to run.}
\end{Verbatim}
\end{codebox}

\textbf{Thread priorities.} Every thread has a priority from \code{Thread.\allowbreak{}MIN\_\allowbreak{}PRIORITY} (1) to \code{Thread.\allowbreak{}MAX\_\allowbreak{}PRIORITY} (10), default \code{Thread.\allowbreak{}NORM\_\allowbreak{}PRIORITY} (5). This is only a \textbf{hint} to the OS scheduler about relative importance — the JVM spec doesn’t guarantee any particular scheduling behavior, and most OSes map Java’s ten levels onto a much coarser native scale. Don’t design correctness around priorities; use them only as a mild tuning hint, if at all.

\begin{codebox}
\begin{Verbatim}[commandchars=\\\{\}]
\PY{n}{Thread}\PY{+w}{ }\PY{n}{background}\PY{+w}{ }\PY{o}{=}\PY{+w}{ }\PY{k}{new}\PY{+w}{ }\PY{n}{Thread}\PY{p}{(}\PY{p}{(}\PY{p}{)}\PY{+w}{ }\PY{o}{\PYZhy{}}\PY{o}{\PYZgt{}}\PY{+w}{ }\PY{n}{doLowPriorityWork}\PY{p}{(}\PY{p}{)}\PY{p}{)}\PY{p}{;}
\PY{n}{background}\PY{p}{.}\PY{n+na}{setPriority}\PY{p}{(}\PY{n}{Thread}\PY{p}{.}\PY{n+na}{MIN\PYZus{}PRIORITY}\PY{p}{)}\PY{p}{;}\PY{+w}{ }\PY{c+c1}{// hint: \PYZdq{}less urgent,\PYZdq{} not a guarantee}
\PY{n}{background}\PY{p}{.}\PY{n+na}{start}\PY{p}{(}\PY{p}{)}\PY{p}{;}
\end{Verbatim}
\end{codebox}

\textbf{Naming threads.} Give threads meaningful names. It costs nothing and it’s the difference between a readable thread dump and a wall of \code{Thread-\allowbreak{}0}, \code{Thread-\allowbreak{}1}, \code{Thread-\allowbreak{}2}.

\textbf{\code{Uncaught\allowbreak{}Exception\allowbreak{}Handler}.} If a thread’s \code{run()} throws an unchecked exception and nobody catches it, the thread just dies silently by default — the exception goes to \code{System.\allowbreak{}err}, but nothing else happens, and no other thread is notified. Set a handler to observe or react to that:

\begin{codebox}
\begin{Verbatim}[commandchars=\\\{\}]
\PY{n}{Thread}\PY{+w}{ }\PY{n}{risky}\PY{+w}{ }\PY{o}{=}\PY{+w}{ }\PY{k}{new}\PY{+w}{ }\PY{n}{Thread}\PY{p}{(}\PY{p}{(}\PY{p}{)}\PY{+w}{ }\PY{o}{\PYZhy{}}\PY{o}{\PYZgt{}}\PY{+w}{ }\PY{p}{\PYZob{}}
\PY{+w}{    }\PY{k}{throw}\PY{+w}{ }\PY{k}{new}\PY{+w}{ }\PY{n}{IllegalStateException}\PY{p}{(}\PY{l+s}{\PYZdq{}}\PY{l+s}{boom}\PY{l+s}{\PYZdq{}}\PY{p}{)}\PY{p}{;}
\PY{p}{\PYZcb{}}\PY{p}{)}\PY{p}{;}
\PY{n}{risky}\PY{p}{.}\PY{n+na}{setUncaughtExceptionHandler}\PY{p}{(}\PY{p}{(}\PY{n}{t}\PY{p}{,}\PY{+w}{ }\PY{n}{e}\PY{p}{)}\PY{+w}{ }\PY{o}{\PYZhy{}}\PY{o}{\PYZgt{}}
\PY{+w}{    }\PY{n}{System}\PY{p}{.}\PY{n+na}{err}\PY{p}{.}\PY{n+na}{println}\PY{p}{(}\PY{l+s}{\PYZdq{}}\PY{l+s}{Thread }\PY{l+s}{\PYZdq{}}\PY{+w}{ }\PY{o}{+}\PY{+w}{ }\PY{n}{t}\PY{p}{.}\PY{n+na}{getName}\PY{p}{(}\PY{p}{)}\PY{+w}{ }\PY{o}{+}\PY{+w}{ }\PY{l+s}{\PYZdq{}}\PY{l+s}{ died: }\PY{l+s}{\PYZdq{}}\PY{+w}{ }\PY{o}{+}\PY{+w}{ }\PY{n}{e}\PY{p}{)}\PY{p}{)}\PY{p}{;}
\PY{n}{risky}\PY{p}{.}\PY{n+na}{start}\PY{p}{(}\PY{p}{)}\PY{p}{;}
\PY{c+c1}{// Prints: Thread Thread\PYZhy{}0 died: java.lang.IllegalStateException: boom}
\PY{c+c1}{// Without the handler, this exception would be easy to miss in production logs.}
\end{Verbatim}
\end{codebox}

You can also set a JVM-wide default via \code{Thread.\allowbreak{}set\allowbreak{}Default\allowbreak{}Uncaught\allowbreak{}Exception\allowbreak{}Handler(...)}, useful for centralized alerting/metrics.

\textbf{Interruption is cooperative cancellation.} \code{Thread.\allowbreak{}interrupt()} does not forcibly stop a thread. It just sets an internal “interrupted” flag and, if the thread is blocked in something interruptible (\code{sleep}, \code{wait}, \code{join}, many blocking I/O and concurrency calls), wakes it up early with an \code{InterruptedException}. The target thread has to \emph{check} for interruption and decide to stop. This is why it’s called cooperative — both sides have to participate.

\begin{codebox}
\begin{Verbatim}[commandchars=\\\{\}]
\PY{k+kd}{public}\PY{+w}{ }\PY{k+kd}{class} \PY{n+nc}{CooperativeCancellation}\PY{+w}{ }\PY{p}{\PYZob{}}
\PY{+w}{    }\PY{k+kd}{public}\PY{+w}{ }\PY{k+kd}{static}\PY{+w}{ }\PY{k+kt}{void}\PY{+w}{ }\PY{n+nf}{main}\PY{p}{(}\PY{n}{String}\PY{o}{[}\PY{o}{]}\PY{+w}{ }\PY{n}{args}\PY{p}{)}\PY{+w}{ }\PY{k+kd}{throws}\PY{+w}{ }\PY{n}{InterruptedException}\PY{+w}{ }\PY{p}{\PYZob{}}
\PY{+w}{        }\PY{n}{Thread}\PY{+w}{ }\PY{n}{worker}\PY{+w}{ }\PY{o}{=}\PY{+w}{ }\PY{k}{new}\PY{+w}{ }\PY{n}{Thread}\PY{p}{(}\PY{p}{(}\PY{p}{)}\PY{+w}{ }\PY{o}{\PYZhy{}}\PY{o}{\PYZgt{}}\PY{+w}{ }\PY{p}{\PYZob{}}
\PY{+w}{            }\PY{k}{while}\PY{+w}{ }\PY{p}{(}\PY{o}{!}\PY{n}{Thread}\PY{p}{.}\PY{n+na}{currentThread}\PY{p}{(}\PY{p}{)}\PY{p}{.}\PY{n+na}{isInterrupted}\PY{p}{(}\PY{p}{)}\PY{p}{)}\PY{+w}{ }\PY{p}{\PYZob{}}
\PY{+w}{                }\PY{c+c1}{// do a unit of work}
\PY{+w}{            }\PY{p}{\PYZcb{}}
\PY{+w}{            }\PY{n}{System}\PY{p}{.}\PY{n+na}{out}\PY{p}{.}\PY{n+na}{println}\PY{p}{(}\PY{l+s}{\PYZdq{}}\PY{l+s}{Worker noticed the interrupt and is exiting cleanly}\PY{l+s}{\PYZdq{}}\PY{p}{)}\PY{p}{;}
\PY{+w}{        }\PY{p}{\PYZcb{}}\PY{p}{)}\PY{p}{;}
\PY{+w}{        }\PY{n}{worker}\PY{p}{.}\PY{n+na}{start}\PY{p}{(}\PY{p}{)}\PY{p}{;}
\PY{+w}{        }\PY{n}{Thread}\PY{p}{.}\PY{n+na}{sleep}\PY{p}{(}\PY{l+m+mi}{50}\PY{p}{)}\PY{p}{;}
\PY{+w}{        }\PY{n}{worker}\PY{p}{.}\PY{n+na}{interrupt}\PY{p}{(}\PY{p}{)}\PY{p}{;}\PY{+w}{ }\PY{c+c1}{// politely ask it to stop}
\PY{+w}{        }\PY{n}{worker}\PY{p}{.}\PY{n+na}{join}\PY{p}{(}\PY{p}{)}\PY{p}{;}
\PY{+w}{    }\PY{p}{\PYZcb{}}
\PY{p}{\PYZcb{}}
\end{Verbatim}
\end{codebox}

If your loop calls a blocking method, catch \code{InterruptedException} and either exit or \textbf{re-interrupt and rethrow} — never swallow it silently, because that erases the cancellation signal for anyone further up the call stack:

\begin{codebox}
\begin{Verbatim}[commandchars=\\\{\}]
\PY{k+kt}{void}\PY{+w}{ }\PY{n+nf}{pollUntilCancelled}\PY{p}{(}\PY{p}{)}\PY{+w}{ }\PY{p}{\PYZob{}}
\PY{+w}{    }\PY{k}{while}\PY{+w}{ }\PY{p}{(}\PY{k+kc}{true}\PY{p}{)}\PY{+w}{ }\PY{p}{\PYZob{}}
\PY{+w}{        }\PY{k}{try}\PY{+w}{ }\PY{p}{\PYZob{}}
\PY{+w}{            }\PY{n}{Thread}\PY{p}{.}\PY{n+na}{sleep}\PY{p}{(}\PY{l+m+mi}{500}\PY{p}{)}\PY{p}{;}
\PY{+w}{            }\PY{c+c1}{// do work}
\PY{+w}{        }\PY{p}{\PYZcb{}}\PY{+w}{ }\PY{k}{catch}\PY{+w}{ }\PY{p}{(}\PY{n}{InterruptedException}\PY{+w}{ }\PY{n}{e}\PY{p}{)}\PY{+w}{ }\PY{p}{\PYZob{}}
\PY{+w}{            }\PY{n}{Thread}\PY{p}{.}\PY{n+na}{currentThread}\PY{p}{(}\PY{p}{)}\PY{p}{.}\PY{n+na}{interrupt}\PY{p}{(}\PY{p}{)}\PY{p}{;}\PY{+w}{ }\PY{c+c1}{// restore the flag for callers}
\PY{+w}{            }\PY{k}{return}\PY{p}{;}\PY{+w}{ }\PY{c+c1}{// stop the loop}
\PY{+w}{        }\PY{p}{\PYZcb{}}
\PY{+w}{    }\PY{p}{\PYZcb{}}
\PY{p}{\PYZcb{}}
\end{Verbatim}
\end{codebox}

\textbf{Why \code{Thread.\allowbreak{}stop()} and \code{Thread.\allowbreak{}suspend()} are dead.} Both are deprecated for removal. \code{stop()} kills a thread instantly, mid-operation, which can leave shared objects in a half-updated, permanently broken state (it releases locks it holds without any guarantee the protected data is consistent). \code{suspend()} freezes a thread while it may be holding a lock, which can deadlock every other thread that needs that lock. Neither has a safe way to guarantee the target thread cleans up. \textbf{Cooperative interruption (or a \code{volatile} cancel flag) is the only sound replacement.}

\textbf{Virtual threads.} Java 21 introduced lightweight, JVM-managed threads that let you write blocking-style code cheaply at massive scale. They’re built on the same \code{Thread} API you just saw, but are covered in full in Chapter 15 (Structured Concurrency / Modern Concurrency), since they change a lot of the cost-benefit calculus for the patterns in this chapter.

\section[Runnable and Callable]{Runnable and Callable}
\label{h:13:runnable-and-callable}
\code{Runnable} and \code{Callable\textless{}\allowbreak{}V\textgreater{}} are the two “unit of work” interfaces you hand to a thread or an executor.

{\tablefont
\begin{xltabular}{\linewidth}{@{}L{0.911}L{0.536}L{1.554}@{}}
\toprule
\rowcolor{tablehdr}
\thd{} & \thd{\code{Runnable}} & \thd{\code{Callable\textless{}\allowbreak{}V\textgreater{}}} \\
\midrule
\endhead
Method & \code{void~\allowbreak{}run()} & \code{V~\allowbreak{}call()\allowbreak{}~\allowbreak{}throws~\allowbreak{}Exception} \\
Returns a result? & No & Yes, typed \code{V} \\
Can throw checked exceptions? & No & Yes \\
Introduced & Java 1.0 & Java 5 (\code{java.\allowbreak{}util.\allowbreak{}concurrent}) \\
Typical use & \code{new~\allowbreak{}Thread(...)}, fire-and-forget tasks & \code{Executor\allowbreak{}Service.\allowbreak{}submit(...)} when you need a \code{Future\textless{}\allowbreak{}V\textgreater{}} \\
\bottomrule
\end{xltabular}
}

\begin{codebox}
\begin{Verbatim}[commandchars=\\\{\}]
\PY{k+kn}{import}\PY{+w}{ }\PY{n+nn}{java.util.concurrent.*}\PY{p}{;}

\PY{k+kd}{public}\PY{+w}{ }\PY{k+kd}{class} \PY{n+nc}{RunnableVsCallable}\PY{+w}{ }\PY{p}{\PYZob{}}
\PY{+w}{    }\PY{k+kd}{public}\PY{+w}{ }\PY{k+kd}{static}\PY{+w}{ }\PY{k+kt}{void}\PY{+w}{ }\PY{n+nf}{main}\PY{p}{(}\PY{n}{String}\PY{o}{[}\PY{o}{]}\PY{+w}{ }\PY{n}{args}\PY{p}{)}\PY{+w}{ }\PY{k+kd}{throws}\PY{+w}{ }\PY{n}{Exception}\PY{+w}{ }\PY{p}{\PYZob{}}
\PY{+w}{        }\PY{n}{Runnable}\PY{+w}{ }\PY{n}{fireAndForget}\PY{+w}{ }\PY{o}{=}\PY{+w}{ }\PY{p}{(}\PY{p}{)}\PY{+w}{ }\PY{o}{\PYZhy{}}\PY{o}{\PYZgt{}}\PY{+w}{ }\PY{n}{System}\PY{p}{.}\PY{n+na}{out}\PY{p}{.}\PY{n+na}{println}\PY{p}{(}\PY{l+s}{\PYZdq{}}\PY{l+s}{logged, no result}\PY{l+s}{\PYZdq{}}\PY{p}{)}\PY{p}{;}

\PY{+w}{        }\PY{n}{Callable}\PY{o}{\PYZlt{}}\PY{n}{Integer}\PY{o}{\PYZgt{}}\PY{+w}{ }\PY{n}{withResult}\PY{+w}{ }\PY{o}{=}\PY{+w}{ }\PY{p}{(}\PY{p}{)}\PY{+w}{ }\PY{o}{\PYZhy{}}\PY{o}{\PYZgt{}}\PY{+w}{ }\PY{p}{\PYZob{}}
\PY{+w}{            }\PY{n}{Thread}\PY{p}{.}\PY{n+na}{sleep}\PY{p}{(}\PY{l+m+mi}{10}\PY{p}{)}\PY{p}{;}
\PY{+w}{            }\PY{k}{return}\PY{+w}{ }\PY{l+m+mi}{42}\PY{p}{;}
\PY{+w}{        }\PY{p}{\PYZcb{}}\PY{p}{;}

\PY{+w}{        }\PY{n}{ExecutorService}\PY{+w}{ }\PY{n}{pool}\PY{+w}{ }\PY{o}{=}\PY{+w}{ }\PY{n}{Executors}\PY{p}{.}\PY{n+na}{newFixedThreadPool}\PY{p}{(}\PY{l+m+mi}{2}\PY{p}{)}\PY{p}{;}
\PY{+w}{        }\PY{n}{pool}\PY{p}{.}\PY{n+na}{execute}\PY{p}{(}\PY{n}{fireAndForget}\PY{p}{)}\PY{p}{;}\PY{+w}{              }\PY{c+c1}{// Runnable: no Future}
\PY{+w}{        }\PY{n}{Future}\PY{o}{\PYZlt{}}\PY{n}{Integer}\PY{o}{\PYZgt{}}\PY{+w}{ }\PY{n}{future}\PY{+w}{ }\PY{o}{=}\PY{+w}{ }\PY{n}{pool}\PY{p}{.}\PY{n+na}{submit}\PY{p}{(}\PY{n}{withResult}\PY{p}{)}\PY{p}{;}\PY{+w}{ }\PY{c+c1}{// Callable: get a Future}
\PY{+w}{        }\PY{n}{System}\PY{p}{.}\PY{n+na}{out}\PY{p}{.}\PY{n+na}{println}\PY{p}{(}\PY{l+s}{\PYZdq{}}\PY{l+s}{Result: }\PY{l+s}{\PYZdq{}}\PY{+w}{ }\PY{o}{+}\PY{+w}{ }\PY{n}{future}\PY{p}{.}\PY{n+na}{get}\PY{p}{(}\PY{p}{)}\PY{p}{)}\PY{p}{;}\PY{+w}{    }\PY{c+c1}{// blocks until done: 42}
\PY{+w}{        }\PY{n}{pool}\PY{p}{.}\PY{n+na}{shutdown}\PY{p}{(}\PY{p}{)}\PY{p}{;}
\PY{+w}{    }\PY{p}{\PYZcb{}}
\PY{p}{\PYZcb{}}
\end{Verbatim}
\end{codebox}

\textbf{\code{Thread.\allowbreak{}sleep()} vs \code{Object.\allowbreak{}wait()}.} Both pause a thread, but they’re built for different purposes:

{\tablefont
\begin{xltabular}{\linewidth}{@{}L{0.564}L{0.699}L{1.737}@{}}
\toprule
\rowcolor{tablehdr}
\thd{} & \thd{\code{Thread.\allowbreak{}sleep(\allowbreak{}ms)}} & \thd{\code{Object.\allowbreak{}wait()}} \\
\midrule
\endhead
Holds locks while paused? & Yes — keeps any locks it holds & No — releases the monitor lock it’s waiting on \\
Needs a lock to call? & No & Yes — must be called while holding the object’s monitor \\
How it wakes up & Timeout elapses, or interrupted & Another thread calls \code{notify()}/\code{notifyAll()}, timeout (if given), or interrupted \\
Typical use & “Pause for a fixed duration” & “Wait until some condition, controlled by another thread, becomes true” \\
\bottomrule
\end{xltabular}
}

\begin{codebox}
\begin{Verbatim}[commandchars=\\\{\}]
\PY{k+kd}{synchronized}\PY{+w}{ }\PY{p}{(}\PY{n}{lock}\PY{p}{)}\PY{+w}{ }\PY{p}{\PYZob{}}
\PY{+w}{    }\PY{k}{while}\PY{+w}{ }\PY{p}{(}\PY{o}{!}\PY{n}{conditionIsTrue}\PY{p}{)}\PY{+w}{ }\PY{p}{\PYZob{}}
\PY{+w}{        }\PY{n}{lock}\PY{p}{.}\PY{n+na}{wait}\PY{p}{(}\PY{p}{)}\PY{p}{;}\PY{+w}{ }\PY{c+c1}{// releases \PYZsq{}lock\PYZsq{} while waiting — other threads can acquire it}
\PY{+w}{    }\PY{p}{\PYZcb{}}
\PY{p}{\PYZcb{}}
\PY{c+c1}{// vs:}
\PY{n}{Thread}\PY{p}{.}\PY{n+na}{sleep}\PY{p}{(}\PY{l+m+mi}{1000}\PY{p}{)}\PY{p}{;}\PY{+w}{ }\PY{c+c1}{// holds every lock this thread already has for the full second}
\end{Verbatim}
\end{codebox}

Sleeping while holding a lock is a common cause of accidental contention — a slow, sleeping thread blocks everyone else waiting on that same lock.

\section[Thread Lifecycle]{Thread Lifecycle}
\label{h:13:thread-lifecycle}
Every \code{Thread} is always in exactly one of six states, defined by the \code{Thread.\allowbreak{}State} enum: \code{NEW}, \code{RUNNABLE}, \code{BLOCKED}, \code{WAITING}, \code{TIMED\_\allowbreak{}WAITING}, \code{TERMINATED}.

\begin{codebox}
\begin{Verbatim}
        new Thread(...)
              |
              v
      +---------------+
      |      NEW       |   (created, start() not yet called)
      +---------------+
              | start()
              v
      +---------------+   synchronized block/method     +---------------+
      |    RUNNABLE    | -------------------------------> |    BLOCKED    |
      | (running, or   |  <------------------------------ | (waiting for  |
      |  ready to run,  |    lock becomes available        |  a monitor    |
      |  scheduler-      |                                  |  lock)        |
      |  dependent)     |                                  +---------------+
      +---------------+
        |     ^   |      ^
        |     |   |      |
 wait()/join()/  |   sleep(ms)/wait(ms)/join(ms)/
 park() (no      |   parkNanos (has a timeout)
 timeout)        |        |                  ^
        v        |        v                  |
  +---------------+   +-------------------+  |
  |   WAITING     |   |  TIMED_WAITING    |--+
  | (indefinite,   |   |  (bounded wait,   |
  |  needs another |   |   times out on    |
  |  thread to act) |   |   its own too)    |
  +---------------+   +-------------------+
        |                       |
        +-----------+-----------+
                     | notify()/notifyAll(),
                     |  interrupt, timeout,
                     |  target thread finishes
                     v
              +---------------+
              |   RUNNABLE    |  (back in the run queue)
              +---------------+
                     |
              run() method returns,
              or an uncaught exception propagates out
                     v
              +---------------+
              |  TERMINATED   |   (finished — cannot be restarted)
              +---------------+
\end{Verbatim}
\end{codebox}

Key transitions to remember for a review conversation:

\begin{itemize}
\item 
\textbf{NEW → RUNNABLE}: \code{start()}. You can only call \code{start()} once; calling it again throws \code{Illegal\allowbreak{}Thread\allowbreak{}State\allowbreak{}Exception}.

\item 
\textbf{RUNNABLE → BLOCKED}: the thread tries to enter a \code{synchronized} block/method whose lock another thread already holds.

\item 
\textbf{RUNNABLE → WAITING}: \code{Object.\allowbreak{}wait()} (no timeout), \code{Thread.\allowbreak{}join()} (no timeout), or \code{LockSupport.\allowbreak{}park()}.

\item 
\textbf{RUNNABLE → TIMED\_WAITING}: \code{Thread.\allowbreak{}sleep(\allowbreak{}ms)}, \code{Object.\allowbreak{}wait(\allowbreak{}ms)}, \code{Thread.\allowbreak{}join(\allowbreak{}ms)}, \code{Lock\allowbreak{}Support.\allowbreak{}park\allowbreak{}Nanos(...)}.

\item 
\textbf{Anything → TERMINATED}: \code{run()} returns normally, or an exception propagates out of it uncaught. There is no supported way back from \code{TERMINATED} — a finished \code{Thread} object is a dead end; make a new \code{Thread} instead.

\end{itemize}

\begin{codebox}
\begin{Verbatim}[commandchars=\\\{\}]
\PY{k+kd}{public}\PY{+w}{ }\PY{k+kd}{class} \PY{n+nc}{LifecycleDemo}\PY{+w}{ }\PY{p}{\PYZob{}}
\PY{+w}{    }\PY{k+kd}{public}\PY{+w}{ }\PY{k+kd}{static}\PY{+w}{ }\PY{k+kt}{void}\PY{+w}{ }\PY{n+nf}{main}\PY{p}{(}\PY{n}{String}\PY{o}{[}\PY{o}{]}\PY{+w}{ }\PY{n}{args}\PY{p}{)}\PY{+w}{ }\PY{k+kd}{throws}\PY{+w}{ }\PY{n}{InterruptedException}\PY{+w}{ }\PY{p}{\PYZob{}}
\PY{+w}{        }\PY{n}{Object}\PY{+w}{ }\PY{n}{lock}\PY{+w}{ }\PY{o}{=}\PY{+w}{ }\PY{k}{new}\PY{+w}{ }\PY{n}{Object}\PY{p}{(}\PY{p}{)}\PY{p}{;}
\PY{+w}{        }\PY{n}{Thread}\PY{+w}{ }\PY{n}{t}\PY{+w}{ }\PY{o}{=}\PY{+w}{ }\PY{k}{new}\PY{+w}{ }\PY{n}{Thread}\PY{p}{(}\PY{p}{(}\PY{p}{)}\PY{+w}{ }\PY{o}{\PYZhy{}}\PY{o}{\PYZgt{}}\PY{+w}{ }\PY{p}{\PYZob{}}
\PY{+w}{            }\PY{k+kd}{synchronized}\PY{+w}{ }\PY{p}{(}\PY{n}{lock}\PY{p}{)}\PY{+w}{ }\PY{p}{\PYZob{}}
\PY{+w}{                }\PY{k}{try}\PY{+w}{ }\PY{p}{\PYZob{}}\PY{+w}{ }\PY{n}{Thread}\PY{p}{.}\PY{n+na}{sleep}\PY{p}{(}\PY{l+m+mi}{200}\PY{p}{)}\PY{p}{;}\PY{+w}{ }\PY{p}{\PYZcb{}}\PY{+w}{ }\PY{k}{catch}\PY{+w}{ }\PY{p}{(}\PY{n}{InterruptedException}\PY{+w}{ }\PY{n}{ignored}\PY{p}{)}\PY{+w}{ }\PY{p}{\PYZob{}}\PY{p}{\PYZcb{}}
\PY{+w}{            }\PY{p}{\PYZcb{}}
\PY{+w}{        }\PY{p}{\PYZcb{}}\PY{p}{)}\PY{p}{;}
\PY{+w}{        }\PY{n}{System}\PY{p}{.}\PY{n+na}{out}\PY{p}{.}\PY{n+na}{println}\PY{p}{(}\PY{n}{t}\PY{p}{.}\PY{n+na}{getState}\PY{p}{(}\PY{p}{)}\PY{p}{)}\PY{p}{;}\PY{+w}{ }\PY{c+c1}{// NEW}
\PY{+w}{        }\PY{n}{t}\PY{p}{.}\PY{n+na}{start}\PY{p}{(}\PY{p}{)}\PY{p}{;}
\PY{+w}{        }\PY{n}{Thread}\PY{p}{.}\PY{n+na}{sleep}\PY{p}{(}\PY{l+m+mi}{50}\PY{p}{)}\PY{p}{;}
\PY{+w}{        }\PY{n}{System}\PY{p}{.}\PY{n+na}{out}\PY{p}{.}\PY{n+na}{println}\PY{p}{(}\PY{n}{t}\PY{p}{.}\PY{n+na}{getState}\PY{p}{(}\PY{p}{)}\PY{p}{)}\PY{p}{;}\PY{+w}{ }\PY{c+c1}{// TIMED\PYZus{}WAITING (inside Thread.sleep)}
\PY{+w}{        }\PY{n}{t}\PY{p}{.}\PY{n+na}{join}\PY{p}{(}\PY{p}{)}\PY{p}{;}
\PY{+w}{        }\PY{n}{System}\PY{p}{.}\PY{n+na}{out}\PY{p}{.}\PY{n+na}{println}\PY{p}{(}\PY{n}{t}\PY{p}{.}\PY{n+na}{getState}\PY{p}{(}\PY{p}{)}\PY{p}{)}\PY{p}{;}\PY{+w}{ }\PY{c+c1}{// TERMINATED}
\PY{+w}{    }\PY{p}{\PYZcb{}}
\PY{p}{\PYZcb{}}
\end{Verbatim}
\end{codebox}

\section[Synchronization]{Synchronization}
\label{h:13:synchronization}
\textbf{Synchronization} means controlling access to shared, mutable state so that only one thread modifies it at a time, and every thread sees the latest value. In Java, the primary built-in tool is the \code{synchronized} keyword, which uses every object’s built-in \textbf{monitor lock} (also called its intrinsic lock — see next section).

\textbf{Synchronized methods vs synchronized blocks.}

\begin{codebox}
\begin{Verbatim}[commandchars=\\\{\}]
\PY{k+kd}{public}\PY{+w}{ }\PY{k+kd}{class} \PY{n+nc}{Counter}\PY{+w}{ }\PY{p}{\PYZob{}}
\PY{+w}{    }\PY{k+kd}{private}\PY{+w}{ }\PY{k+kt}{int}\PY{+w}{ }\PY{n}{count}\PY{+w}{ }\PY{o}{=}\PY{+w}{ }\PY{l+m+mi}{0}\PY{p}{;}

\PY{+w}{    }\PY{c+c1}{// Synchronized instance method: locks \PYZsq{}this\PYZsq{}}
\PY{+w}{    }\PY{k+kd}{public}\PY{+w}{ }\PY{k+kd}{synchronized}\PY{+w}{ }\PY{k+kt}{void}\PY{+w}{ }\PY{n+nf}{incrementWhole}\PY{p}{(}\PY{p}{)}\PY{+w}{ }\PY{p}{\PYZob{}}
\PY{+w}{        }\PY{n}{count}\PY{o}{+}\PY{o}{+}\PY{p}{;}
\PY{+w}{    }\PY{p}{\PYZcb{}}

\PY{+w}{    }\PY{c+c1}{// Synchronized block: locks an explicit object, narrower scope}
\PY{+w}{    }\PY{k+kd}{private}\PY{+w}{ }\PY{k+kd}{final}\PY{+w}{ }\PY{n}{Object}\PY{+w}{ }\PY{n}{lock}\PY{+w}{ }\PY{o}{=}\PY{+w}{ }\PY{k}{new}\PY{+w}{ }\PY{n}{Object}\PY{p}{(}\PY{p}{)}\PY{p}{;}
\PY{+w}{    }\PY{k+kd}{public}\PY{+w}{ }\PY{k+kt}{void}\PY{+w}{ }\PY{n+nf}{incrementBlock}\PY{p}{(}\PY{p}{)}\PY{+w}{ }\PY{p}{\PYZob{}}
\PY{+w}{        }\PY{k+kt}{int}\PY{+w}{ }\PY{n}{expensive}\PY{+w}{ }\PY{o}{=}\PY{+w}{ }\PY{n}{computeSomethingUnrelated}\PY{p}{(}\PY{p}{)}\PY{p}{;}\PY{+w}{ }\PY{c+c1}{// not protected — doesn\PYZsq{}t need to be}
\PY{+w}{        }\PY{k+kd}{synchronized}\PY{+w}{ }\PY{p}{(}\PY{n}{lock}\PY{p}{)}\PY{+w}{ }\PY{p}{\PYZob{}}
\PY{+w}{            }\PY{n}{count}\PY{o}{+}\PY{o}{+}\PY{p}{;}\PY{+w}{ }\PY{c+c1}{// only the truly shared part is inside the lock}
\PY{+w}{        }\PY{p}{\PYZcb{}}
\PY{+w}{    }\PY{p}{\PYZcb{}}

\PY{+w}{    }\PY{k+kd}{private}\PY{+w}{ }\PY{k+kt}{int}\PY{+w}{ }\PY{n+nf}{computeSomethingUnrelated}\PY{p}{(}\PY{p}{)}\PY{+w}{ }\PY{p}{\PYZob{}}\PY{+w}{ }\PY{k}{return}\PY{+w}{ }\PY{l+m+mi}{1}\PY{p}{;}\PY{+w}{ }\PY{p}{\PYZcb{}}

\PY{+w}{    }\PY{k+kd}{public}\PY{+w}{ }\PY{k+kd}{synchronized}\PY{+w}{ }\PY{k+kt}{int}\PY{+w}{ }\PY{n+nf}{getCount}\PY{p}{(}\PY{p}{)}\PY{+w}{ }\PY{p}{\PYZob{}}
\PY{+w}{        }\PY{k}{return}\PY{+w}{ }\PY{n}{count}\PY{p}{;}\PY{+w}{ }\PY{c+c1}{// reads must be synchronized too, or you can see a stale value}
\PY{+w}{    }\PY{p}{\PYZcb{}}
\PY{p}{\PYZcb{}}
\end{Verbatim}
\end{codebox}

\textbf{What object gets locked?}

{\tablefont
\begin{xltabular}{\linewidth}{@{}L{0.786}L{1.214}@{}}
\toprule
\rowcolor{tablehdr}
\thd{Declaration} & \thd{Lock acquired} \\
\midrule
\endhead
\code{synchronized~\allowbreak{}void~\allowbreak{}method()} & the instance (\code{this}) \\
\code{static~\allowbreak{}synchronized~\allowbreak{}void~\allowbreak{}method()} & the \code{Class} object (one lock shared by \emph{all} instances) \\
\code{synchronized~\allowbreak{}(\allowbreak{}obj)\allowbreak{}~\allowbreak{}\{\allowbreak{}~\allowbreak{}...\allowbreak{}~\allowbreak{}\}} & whatever object \code{obj} refers to \\
\bottomrule
\end{xltabular}
}

Mixing an instance-level lock with a static-level lock is a classic bug: two threads think they’re synchronizing on “the same thing” but are actually holding two different locks, so no mutual exclusion happens at all.

\textbf{Reentrancy.} Intrinsic locks are \textbf{reentrant}: a thread that already holds a lock can acquire it again (e.g., calling another synchronized method on the same object from inside a synchronized method) without blocking itself. The JVM tracks a hold count per thread.

\begin{codebox}
\begin{Verbatim}[commandchars=\\\{\}]
\PY{k+kd}{public}\PY{+w}{ }\PY{k+kd}{synchronized}\PY{+w}{ }\PY{k+kt}{void}\PY{+w}{ }\PY{n+nf}{outer}\PY{p}{(}\PY{p}{)}\PY{+w}{ }\PY{p}{\PYZob{}}
\PY{+w}{    }\PY{n}{inner}\PY{p}{(}\PY{p}{)}\PY{p}{;}\PY{+w}{ }\PY{c+c1}{// same thread re\PYZhy{}acquiring the same lock — allowed, does not deadlock}
\PY{p}{\PYZcb{}}
\PY{k+kd}{public}\PY{+w}{ }\PY{k+kd}{synchronized}\PY{+w}{ }\PY{k+kt}{void}\PY{+w}{ }\PY{n+nf}{inner}\PY{p}{(}\PY{p}{)}\PY{+w}{ }\PY{p}{\PYZob{}}
\PY{+w}{    }\PY{c+c1}{// ...}
\PY{p}{\PYZcb{}}
\end{Verbatim}
\end{codebox}

\textbf{\code{wait()} / \code{notify()} / \code{notifyAll()} — and the mandatory \code{while} loop.} These three methods coordinate threads around a condition, and must be called while holding the object’s monitor. \code{wait()} releases the lock and parks the thread; \code{notify()} wakes \emph{one} waiting thread; \code{notifyAll()} wakes \emph{all} of them. Always check the condition in a \code{while} loop, never an \code{if}:

\begin{codebox}
\begin{Verbatim}[commandchars=\\\{\}]
\PY{k+kd}{public}\PY{+w}{ }\PY{k+kd}{class} \PY{n+nc}{BoundedBuffer}\PY{o}{\PYZlt{}}\PY{n}{T}\PY{o}{\PYZgt{}}\PY{+w}{ }\PY{p}{\PYZob{}}
\PY{+w}{    }\PY{k+kd}{private}\PY{+w}{ }\PY{k+kd}{final}\PY{+w}{ }\PY{n}{Queue}\PY{o}{\PYZlt{}}\PY{n}{T}\PY{o}{\PYZgt{}}\PY{+w}{ }\PY{n}{queue}\PY{+w}{ }\PY{o}{=}\PY{+w}{ }\PY{k}{new}\PY{+w}{ }\PY{n}{LinkedList}\PY{o}{\PYZlt{}}\PY{o}{\PYZgt{}}\PY{p}{(}\PY{p}{)}\PY{p}{;}
\PY{+w}{    }\PY{k+kd}{private}\PY{+w}{ }\PY{k+kd}{final}\PY{+w}{ }\PY{k+kt}{int}\PY{+w}{ }\PY{n}{capacity}\PY{p}{;}

\PY{+w}{    }\PY{k+kd}{public}\PY{+w}{ }\PY{n+nf}{BoundedBuffer}\PY{p}{(}\PY{k+kt}{int}\PY{+w}{ }\PY{n}{capacity}\PY{p}{)}\PY{+w}{ }\PY{p}{\PYZob{}}\PY{+w}{ }\PY{k}{this}\PY{p}{.}\PY{n+na}{capacity}\PY{+w}{ }\PY{o}{=}\PY{+w}{ }\PY{n}{capacity}\PY{p}{;}\PY{+w}{ }\PY{p}{\PYZcb{}}

\PY{+w}{    }\PY{k+kd}{public}\PY{+w}{ }\PY{k+kd}{synchronized}\PY{+w}{ }\PY{k+kt}{void}\PY{+w}{ }\PY{n+nf}{put}\PY{p}{(}\PY{n}{T}\PY{+w}{ }\PY{n}{item}\PY{p}{)}\PY{+w}{ }\PY{k+kd}{throws}\PY{+w}{ }\PY{n}{InterruptedException}\PY{+w}{ }\PY{p}{\PYZob{}}
\PY{+w}{        }\PY{k}{while}\PY{+w}{ }\PY{p}{(}\PY{n}{queue}\PY{p}{.}\PY{n+na}{size}\PY{p}{(}\PY{p}{)}\PY{+w}{ }\PY{o}{=}\PY{o}{=}\PY{+w}{ }\PY{n}{capacity}\PY{p}{)}\PY{+w}{ }\PY{p}{\PYZob{}}\PY{+w}{ }\PY{c+c1}{// WHILE, not if — re\PYZhy{}check after waking up}
\PY{+w}{            }\PY{n}{wait}\PY{p}{(}\PY{p}{)}\PY{p}{;}\PY{+w}{ }\PY{c+c1}{// releases the lock, parks this thread}
\PY{+w}{        }\PY{p}{\PYZcb{}}
\PY{+w}{        }\PY{n}{queue}\PY{p}{.}\PY{n+na}{add}\PY{p}{(}\PY{n}{item}\PY{p}{)}\PY{p}{;}
\PY{+w}{        }\PY{n}{notifyAll}\PY{p}{(}\PY{p}{)}\PY{p}{;}\PY{+w}{ }\PY{c+c1}{// wake up any waiting consumers}
\PY{+w}{    }\PY{p}{\PYZcb{}}

\PY{+w}{    }\PY{k+kd}{public}\PY{+w}{ }\PY{k+kd}{synchronized}\PY{+w}{ }\PY{n}{T}\PY{+w}{ }\PY{n+nf}{take}\PY{p}{(}\PY{p}{)}\PY{+w}{ }\PY{k+kd}{throws}\PY{+w}{ }\PY{n}{InterruptedException}\PY{+w}{ }\PY{p}{\PYZob{}}
\PY{+w}{        }\PY{k}{while}\PY{+w}{ }\PY{p}{(}\PY{n}{queue}\PY{p}{.}\PY{n+na}{isEmpty}\PY{p}{(}\PY{p}{)}\PY{p}{)}\PY{+w}{ }\PY{p}{\PYZob{}}\PY{+w}{ }\PY{c+c1}{// WHILE — guard condition must be re\PYZhy{}checked}
\PY{+w}{            }\PY{n}{wait}\PY{p}{(}\PY{p}{)}\PY{p}{;}
\PY{+w}{        }\PY{p}{\PYZcb{}}
\PY{+w}{        }\PY{n}{T}\PY{+w}{ }\PY{n}{item}\PY{+w}{ }\PY{o}{=}\PY{+w}{ }\PY{n}{queue}\PY{p}{.}\PY{n+na}{remove}\PY{p}{(}\PY{p}{)}\PY{p}{;}
\PY{+w}{        }\PY{n}{notifyAll}\PY{p}{(}\PY{p}{)}\PY{p}{;}\PY{+w}{ }\PY{c+c1}{// wake up any waiting producers}
\PY{+w}{        }\PY{k}{return}\PY{+w}{ }\PY{n}{item}\PY{p}{;}
\PY{+w}{    }\PY{p}{\PYZcb{}}
\PY{p}{\PYZcb{}}
\end{Verbatim}
\end{codebox}

Why \code{while}, not \code{if}:

\begin{itemize}
\item 
\textbf{Spurious wakeups.} The JVM spec permits \code{wait()} to return without any \code{notify()} ever being called. Rare, but legal — so the condition must be re-checked regardless.

\item 
\textbf{Multiple waiters.} With \code{notifyAll()}, several threads wake up but only one may get to act before the condition becomes false again (e.g., two consumers wake up but only one item was added). Each thread must re-verify before proceeding.

\end{itemize}

\textbf{Lost wakeup.} This happens when a thread calls \code{notify()}/\code{notifyAll()} \emph{before} another thread starts waiting, and the signal is lost forever because there was nobody listening yet.

\begin{codebox}
\begin{Verbatim}[commandchars=\\\{\}]
\PY{c+c1}{// BROKEN: checking the condition outside the lock creates a window for a lost wakeup}
\PY{k+kd}{public}\PY{+w}{ }\PY{k+kt}{void}\PY{+w}{ }\PY{n+nf}{putBroken}\PY{p}{(}\PY{n}{T}\PY{+w}{ }\PY{n}{item}\PY{p}{)}\PY{+w}{ }\PY{k+kd}{throws}\PY{+w}{ }\PY{n}{InterruptedException}\PY{+w}{ }\PY{p}{\PYZob{}}
\PY{+w}{    }\PY{k}{if}\PY{+w}{ }\PY{p}{(}\PY{n}{isFull}\PY{p}{(}\PY{p}{)}\PY{p}{)}\PY{+w}{ }\PY{p}{\PYZob{}}\PY{+w}{          }\PY{c+c1}{// check happens OUTSIDE the lock}
\PY{+w}{        }\PY{k+kd}{synchronized}\PY{+w}{ }\PY{p}{(}\PY{k}{this}\PY{p}{)}\PY{+w}{ }\PY{p}{\PYZob{}}
\PY{+w}{            }\PY{n}{wait}\PY{p}{(}\PY{p}{)}\PY{p}{;}\PY{+w}{           }\PY{c+c1}{// another thread could notify() in between check and wait()}
\PY{+w}{        }\PY{p}{\PYZcb{}}\PY{+w}{                     }\PY{c+c1}{// — that notify is lost, and this thread waits forever}
\PY{+w}{    }\PY{p}{\PYZcb{}}
\PY{+w}{    }\PY{c+c1}{// ...}
\PY{p}{\PYZcb{}}
\end{Verbatim}
\end{codebox}

The fix is what \code{BoundedBuffer} already does above: the check and the \code{wait()} call happen \textbf{inside the same synchronized block}, so no notification can slip through the gap.

\textbf{Lock granularity.} A single, coarse lock around an entire large object is simple and safe, but it forces every operation to queue up behind every other, even unrelated ones — that’s a \textbf{throughput} cost. Splitting into several fine-grained locks (e.g., one lock per shard/bucket) improves parallelism but is more complex and raises the risk of deadlock when a thread needs more than one lock at a time.

\textbf{Deadlock.} Two (or more) threads each hold a lock the other needs, and neither can proceed. Classic two-lock example:

\begin{codebox}
\begin{Verbatim}[commandchars=\\\{\}]
\PY{k+kd}{public}\PY{+w}{ }\PY{k+kd}{class} \PY{n+nc}{DeadlockDemo}\PY{+w}{ }\PY{p}{\PYZob{}}
\PY{+w}{    }\PY{k+kd}{private}\PY{+w}{ }\PY{k+kd}{static}\PY{+w}{ }\PY{k+kd}{final}\PY{+w}{ }\PY{n}{Object}\PY{+w}{ }\PY{n}{lockA}\PY{+w}{ }\PY{o}{=}\PY{+w}{ }\PY{k}{new}\PY{+w}{ }\PY{n}{Object}\PY{p}{(}\PY{p}{)}\PY{p}{;}
\PY{+w}{    }\PY{k+kd}{private}\PY{+w}{ }\PY{k+kd}{static}\PY{+w}{ }\PY{k+kd}{final}\PY{+w}{ }\PY{n}{Object}\PY{+w}{ }\PY{n}{lockB}\PY{+w}{ }\PY{o}{=}\PY{+w}{ }\PY{k}{new}\PY{+w}{ }\PY{n}{Object}\PY{p}{(}\PY{p}{)}\PY{p}{;}

\PY{+w}{    }\PY{k+kd}{public}\PY{+w}{ }\PY{k+kd}{static}\PY{+w}{ }\PY{k+kt}{void}\PY{+w}{ }\PY{n+nf}{main}\PY{p}{(}\PY{n}{String}\PY{o}{[}\PY{o}{]}\PY{+w}{ }\PY{n}{args}\PY{p}{)}\PY{+w}{ }\PY{p}{\PYZob{}}
\PY{+w}{        }\PY{n}{Thread}\PY{+w}{ }\PY{n}{t1}\PY{+w}{ }\PY{o}{=}\PY{+w}{ }\PY{k}{new}\PY{+w}{ }\PY{n}{Thread}\PY{p}{(}\PY{p}{(}\PY{p}{)}\PY{+w}{ }\PY{o}{\PYZhy{}}\PY{o}{\PYZgt{}}\PY{+w}{ }\PY{p}{\PYZob{}}
\PY{+w}{            }\PY{k+kd}{synchronized}\PY{+w}{ }\PY{p}{(}\PY{n}{lockA}\PY{p}{)}\PY{+w}{ }\PY{p}{\PYZob{}}
\PY{+w}{                }\PY{n}{sleepQuiet}\PY{p}{(}\PY{l+m+mi}{50}\PY{p}{)}\PY{p}{;}
\PY{+w}{                }\PY{k+kd}{synchronized}\PY{+w}{ }\PY{p}{(}\PY{n}{lockB}\PY{p}{)}\PY{+w}{ }\PY{p}{\PYZob{}}\PY{+w}{ }\PY{n}{System}\PY{p}{.}\PY{n+na}{out}\PY{p}{.}\PY{n+na}{println}\PY{p}{(}\PY{l+s}{\PYZdq{}}\PY{l+s}{t1 got both}\PY{l+s}{\PYZdq{}}\PY{p}{)}\PY{p}{;}\PY{+w}{ }\PY{p}{\PYZcb{}}
\PY{+w}{            }\PY{p}{\PYZcb{}}
\PY{+w}{        }\PY{p}{\PYZcb{}}\PY{p}{)}\PY{p}{;}
\PY{+w}{        }\PY{n}{Thread}\PY{+w}{ }\PY{n}{t2}\PY{+w}{ }\PY{o}{=}\PY{+w}{ }\PY{k}{new}\PY{+w}{ }\PY{n}{Thread}\PY{p}{(}\PY{p}{(}\PY{p}{)}\PY{+w}{ }\PY{o}{\PYZhy{}}\PY{o}{\PYZgt{}}\PY{+w}{ }\PY{p}{\PYZob{}}
\PY{+w}{            }\PY{k+kd}{synchronized}\PY{+w}{ }\PY{p}{(}\PY{n}{lockB}\PY{p}{)}\PY{+w}{ }\PY{p}{\PYZob{}}\PY{+w}{           }\PY{c+c1}{// t2 grabs B first...}
\PY{+w}{                }\PY{n}{sleepQuiet}\PY{p}{(}\PY{l+m+mi}{50}\PY{p}{)}\PY{p}{;}
\PY{+w}{                }\PY{k+kd}{synchronized}\PY{+w}{ }\PY{p}{(}\PY{n}{lockA}\PY{p}{)}\PY{+w}{ }\PY{p}{\PYZob{}}\PY{+w}{ }\PY{n}{System}\PY{p}{.}\PY{n+na}{out}\PY{p}{.}\PY{n+na}{println}\PY{p}{(}\PY{l+s}{\PYZdq{}}\PY{l+s}{t2 got both}\PY{l+s}{\PYZdq{}}\PY{p}{)}\PY{p}{;}\PY{+w}{ }\PY{p}{\PYZcb{}}\PY{+w}{ }\PY{c+c1}{// ...then wants A}
\PY{+w}{            }\PY{p}{\PYZcb{}}
\PY{+w}{        }\PY{p}{\PYZcb{}}\PY{p}{)}\PY{p}{;}
\PY{+w}{        }\PY{n}{t1}\PY{p}{.}\PY{n+na}{start}\PY{p}{(}\PY{p}{)}\PY{p}{;}
\PY{+w}{        }\PY{n}{t2}\PY{p}{.}\PY{n+na}{start}\PY{p}{(}\PY{p}{)}\PY{p}{;}
\PY{+w}{        }\PY{c+c1}{// t1 holds A, waits for B. t2 holds B, waits for A. Neither ever finishes.}
\PY{+w}{    }\PY{p}{\PYZcb{}}

\PY{+w}{    }\PY{k+kd}{private}\PY{+w}{ }\PY{k+kd}{static}\PY{+w}{ }\PY{k+kt}{void}\PY{+w}{ }\PY{n+nf}{sleepQuiet}\PY{p}{(}\PY{k+kt}{long}\PY{+w}{ }\PY{n}{ms}\PY{p}{)}\PY{+w}{ }\PY{p}{\PYZob{}}
\PY{+w}{        }\PY{k}{try}\PY{+w}{ }\PY{p}{\PYZob{}}\PY{+w}{ }\PY{n}{Thread}\PY{p}{.}\PY{n+na}{sleep}\PY{p}{(}\PY{n}{ms}\PY{p}{)}\PY{p}{;}\PY{+w}{ }\PY{p}{\PYZcb{}}\PY{+w}{ }\PY{k}{catch}\PY{+w}{ }\PY{p}{(}\PY{n}{InterruptedException}\PY{+w}{ }\PY{n}{ignored}\PY{p}{)}\PY{+w}{ }\PY{p}{\PYZob{}}\PY{p}{\PYZcb{}}
\PY{+w}{    }\PY{p}{\PYZcb{}}
\PY{p}{\PYZcb{}}
\end{Verbatim}
\end{codebox}

\textbf{The fix: consistent lock ordering.} If every thread in the system always acquires locks in the same global order, a cycle can never form:

\begin{codebox}
\begin{Verbatim}[commandchars=\\\{\}]
\PY{c+c1}{// FIXED: every thread acquires lockA before lockB, no exceptions.}
\PY{n}{Thread}\PY{+w}{ }\PY{n}{t1}\PY{+w}{ }\PY{o}{=}\PY{+w}{ }\PY{k}{new}\PY{+w}{ }\PY{n}{Thread}\PY{p}{(}\PY{p}{(}\PY{p}{)}\PY{+w}{ }\PY{o}{\PYZhy{}}\PY{o}{\PYZgt{}}\PY{+w}{ }\PY{p}{\PYZob{}}
\PY{+w}{    }\PY{k+kd}{synchronized}\PY{+w}{ }\PY{p}{(}\PY{n}{lockA}\PY{p}{)}\PY{+w}{ }\PY{p}{\PYZob{}}
\PY{+w}{        }\PY{k+kd}{synchronized}\PY{+w}{ }\PY{p}{(}\PY{n}{lockB}\PY{p}{)}\PY{+w}{ }\PY{p}{\PYZob{}}\PY{+w}{ }\PY{n}{System}\PY{p}{.}\PY{n+na}{out}\PY{p}{.}\PY{n+na}{println}\PY{p}{(}\PY{l+s}{\PYZdq{}}\PY{l+s}{t1 got both}\PY{l+s}{\PYZdq{}}\PY{p}{)}\PY{p}{;}\PY{+w}{ }\PY{p}{\PYZcb{}}
\PY{+w}{    }\PY{p}{\PYZcb{}}
\PY{p}{\PYZcb{}}\PY{p}{)}\PY{p}{;}
\PY{n}{Thread}\PY{+w}{ }\PY{n}{t2}\PY{+w}{ }\PY{o}{=}\PY{+w}{ }\PY{k}{new}\PY{+w}{ }\PY{n}{Thread}\PY{p}{(}\PY{p}{(}\PY{p}{)}\PY{+w}{ }\PY{o}{\PYZhy{}}\PY{o}{\PYZgt{}}\PY{+w}{ }\PY{p}{\PYZob{}}
\PY{+w}{    }\PY{k+kd}{synchronized}\PY{+w}{ }\PY{p}{(}\PY{n}{lockA}\PY{p}{)}\PY{+w}{ }\PY{p}{\PYZob{}}\PY{+w}{   }\PY{c+c1}{// t2 also takes A first now, matching t1\PYZsq{}s order}
\PY{+w}{        }\PY{k+kd}{synchronized}\PY{+w}{ }\PY{p}{(}\PY{n}{lockB}\PY{p}{)}\PY{+w}{ }\PY{p}{\PYZob{}}\PY{+w}{ }\PY{n}{System}\PY{p}{.}\PY{n+na}{out}\PY{p}{.}\PY{n+na}{println}\PY{p}{(}\PY{l+s}{\PYZdq{}}\PY{l+s}{t2 got both}\PY{l+s}{\PYZdq{}}\PY{p}{)}\PY{p}{;}\PY{+w}{ }\PY{p}{\PYZcb{}}
\PY{+w}{    }\PY{p}{\PYZcb{}}
\PY{p}{\PYZcb{}}\PY{p}{)}\PY{p}{;}
\end{Verbatim}
\end{codebox}

Other standard deadlock fixes: use \code{tryLock} with a timeout instead of a blocking acquire (covered under \code{ReentrantLock}), or acquire all needed locks up front and back out entirely if any single one isn’t available.

\textbf{Livelock.} Threads aren’t blocked — they’re both actively running — but they keep responding to each other in a way that stops either from making progress. A common example is two threads that each politely “back off” and retry when they detect contention, but they back off in lockstep forever and never actually get through.

\textbf{Starvation.} A thread is perpetually denied access to a resource because other threads keep getting priority — e.g., low-priority threads never running because high-priority ones keep flooding the scheduler, or a “fair” resource being repeatedly grabbed by newcomers ahead of a long-waiting thread. Using a \textbf{fair lock} (see \code{ReentrantLock} below) is one direct mitigation.

\section[Intrinsic Locks]{Intrinsic Locks}
\label{h:13:intrinsic-locks}
Every Java object has a built-in lock, historically called its \textbf{monitor} or \textbf{intrinsic lock}. It’s what \code{synchronized} uses under the hood — there’s no separate “lock object” you create; the object \emph{is} the lock. Only one thread can hold a given object’s monitor at a time.

\begin{codebox}
\begin{Verbatim}[commandchars=\\\{\}]
\PY{k+kd}{public}\PY{+w}{ }\PY{k+kd}{class} \PY{n+nc}{IntrinsicLockExample}\PY{+w}{ }\PY{p}{\PYZob{}}
\PY{+w}{    }\PY{k+kd}{private}\PY{+w}{ }\PY{k+kd}{final}\PY{+w}{ }\PY{n}{Object}\PY{+w}{ }\PY{n}{monitor}\PY{+w}{ }\PY{o}{=}\PY{+w}{ }\PY{k}{new}\PY{+w}{ }\PY{n}{Object}\PY{p}{(}\PY{p}{)}\PY{p}{;}
\PY{+w}{    }\PY{k+kd}{private}\PY{+w}{ }\PY{k+kt}{int}\PY{+w}{ }\PY{n}{balance}\PY{+w}{ }\PY{o}{=}\PY{+w}{ }\PY{l+m+mi}{0}\PY{p}{;}

\PY{+w}{    }\PY{k+kd}{public}\PY{+w}{ }\PY{k+kt}{void}\PY{+w}{ }\PY{n+nf}{deposit}\PY{p}{(}\PY{k+kt}{int}\PY{+w}{ }\PY{n}{amount}\PY{p}{)}\PY{+w}{ }\PY{p}{\PYZob{}}
\PY{+w}{        }\PY{k+kd}{synchronized}\PY{+w}{ }\PY{p}{(}\PY{n}{monitor}\PY{p}{)}\PY{+w}{ }\PY{p}{\PYZob{}}\PY{+w}{ }\PY{c+c1}{// acquire monitor\PYZsq{}s intrinsic lock}
\PY{+w}{            }\PY{n}{balance}\PY{+w}{ }\PY{o}{+}\PY{o}{=}\PY{+w}{ }\PY{n}{amount}\PY{p}{;}
\PY{+w}{        }\PY{p}{\PYZcb{}}\PY{+w}{ }\PY{c+c1}{// lock is automatically released here, even if an exception is thrown}
\PY{+w}{    }\PY{p}{\PYZcb{}}
\PY{p}{\PYZcb{}}
\end{Verbatim}
\end{codebox}

Intrinsic locks have three defining properties, all already demonstrated above:

\begin{itemize}
\item 
\textbf{Mutual exclusion} — only one thread inside the block/method at a time.

\item 
\textbf{Reentrancy} — the holding thread can re-enter without blocking itself.

\item 
\textbf{Automatic release} — released when the block exits, normally or via exception. You never write an explicit “unlock” call, which avoids one entire class of bugs (forgetting to unlock) — but it also means you can’t do things like acquire in one method and release in another, or time out on acquisition. That’s exactly the gap \code{ReentrantLock} fills.

\end{itemize}

Intrinsic locks are \textbf{not fair} by default — the JVM makes no promise about which waiting thread gets the lock next — and they offer no way to interrupt a thread that’s blocked waiting to enter a \code{synchronized} block, and no way to check “is this lock free?” without blocking. Those limitations are the whole reason \code{java.\allowbreak{}util.\allowbreak{}concurrent.\allowbreak{}locks} exists.

\section[ReentrantLock]{ReentrantLock}
\label{h:13:reentrantlock}
\code{ReentrantLock} is an explicit, class-based lock with the same reentrant mutual-exclusion guarantee as \code{synchronized}, plus extra capabilities \code{synchronized} doesn’t offer: fairness policies, timed/non-blocking acquisition, interruptible acquisition, and multiple wait-conditions per lock.

\textbf{The mandatory pattern: acquire, then \code{try}/\code{finally} to release.} Unlike \code{synchronized}, nothing releases the lock for you automatically.

\begin{codebox}
\begin{Verbatim}[commandchars=\\\{\}]
\PY{k+kn}{import}\PY{+w}{ }\PY{n+nn}{java.util.concurrent.locks.ReentrantLock}\PY{p}{;}

\PY{k+kd}{public}\PY{+w}{ }\PY{k+kd}{class} \PY{n+nc}{Account}\PY{+w}{ }\PY{p}{\PYZob{}}
\PY{+w}{    }\PY{k+kd}{private}\PY{+w}{ }\PY{k+kd}{final}\PY{+w}{ }\PY{n}{ReentrantLock}\PY{+w}{ }\PY{n}{lock}\PY{+w}{ }\PY{o}{=}\PY{+w}{ }\PY{k}{new}\PY{+w}{ }\PY{n}{ReentrantLock}\PY{p}{(}\PY{p}{)}\PY{p}{;}
\PY{+w}{    }\PY{k+kd}{private}\PY{+w}{ }\PY{k+kt}{int}\PY{+w}{ }\PY{n}{balance}\PY{+w}{ }\PY{o}{=}\PY{+w}{ }\PY{l+m+mi}{0}\PY{p}{;}

\PY{+w}{    }\PY{k+kd}{public}\PY{+w}{ }\PY{k+kt}{void}\PY{+w}{ }\PY{n+nf}{withdraw}\PY{p}{(}\PY{k+kt}{int}\PY{+w}{ }\PY{n}{amount}\PY{p}{)}\PY{+w}{ }\PY{p}{\PYZob{}}
\PY{+w}{        }\PY{n}{lock}\PY{p}{.}\PY{n+na}{lock}\PY{p}{(}\PY{p}{)}\PY{p}{;}
\PY{+w}{        }\PY{k}{try}\PY{+w}{ }\PY{p}{\PYZob{}}
\PY{+w}{            }\PY{k}{if}\PY{+w}{ }\PY{p}{(}\PY{n}{amount}\PY{+w}{ }\PY{o}{\PYZgt{}}\PY{+w}{ }\PY{n}{balance}\PY{p}{)}\PY{+w}{ }\PY{k}{throw}\PY{+w}{ }\PY{k}{new}\PY{+w}{ }\PY{n}{IllegalStateException}\PY{p}{(}\PY{l+s}{\PYZdq{}}\PY{l+s}{insufficient funds}\PY{l+s}{\PYZdq{}}\PY{p}{)}\PY{p}{;}
\PY{+w}{            }\PY{n}{balance}\PY{+w}{ }\PY{o}{\PYZhy{}}\PY{o}{=}\PY{+w}{ }\PY{n}{amount}\PY{p}{;}
\PY{+w}{        }\PY{p}{\PYZcb{}}\PY{+w}{ }\PY{k}{finally}\PY{+w}{ }\PY{p}{\PYZob{}}
\PY{+w}{            }\PY{n}{lock}\PY{p}{.}\PY{n+na}{unlock}\PY{p}{(}\PY{p}{)}\PY{p}{;}\PY{+w}{ }\PY{c+c1}{// MUST be in finally — otherwise an exception leaks the lock forever}
\PY{+w}{        }\PY{p}{\PYZcb{}}
\PY{+w}{    }\PY{p}{\PYZcb{}}
\PY{p}{\PYZcb{}}
\end{Verbatim}
\end{codebox}

\textbf{Fairness.} The no-arg constructor gives a non-fair lock (higher throughput, but a thread can theoretically be overtaken repeatedly). \code{new~\allowbreak{}Reentrant\allowbreak{}Lock(\allowbreak{}true)} gives a fair lock — the longest-waiting thread goes next, which trades some throughput for protection against starvation.

\begin{codebox}
\begin{Verbatim}[commandchars=\\\{\}]
\PY{n}{ReentrantLock}\PY{+w}{ }\PY{n}{fair}\PY{+w}{ }\PY{o}{=}\PY{+w}{ }\PY{k}{new}\PY{+w}{ }\PY{n}{ReentrantLock}\PY{p}{(}\PY{k+kc}{true}\PY{p}{)}\PY{p}{;}\PY{+w}{ }\PY{c+c1}{// FIFO\PYZhy{}ish ordering, reduces starvation risk}
\end{Verbatim}
\end{codebox}

\textbf{\code{tryLock} with a timeout.} Instead of blocking forever, ask for the lock and give up after a bound — a direct deadlock-avoidance tool:

\begin{codebox}
\begin{Verbatim}[commandchars=\\\{\}]
\PY{k+kn}{import}\PY{+w}{ }\PY{n+nn}{java.util.concurrent.TimeUnit}\PY{p}{;}

\PY{k+kd}{public}\PY{+w}{ }\PY{k+kt}{boolean}\PY{+w}{ }\PY{n+nf}{tryWithdraw}\PY{p}{(}\PY{k+kt}{int}\PY{+w}{ }\PY{n}{amount}\PY{p}{)}\PY{+w}{ }\PY{k+kd}{throws}\PY{+w}{ }\PY{n}{InterruptedException}\PY{+w}{ }\PY{p}{\PYZob{}}
\PY{+w}{    }\PY{k}{if}\PY{+w}{ }\PY{p}{(}\PY{o}{!}\PY{n}{lock}\PY{p}{.}\PY{n+na}{tryLock}\PY{p}{(}\PY{l+m+mi}{200}\PY{p}{,}\PY{+w}{ }\PY{n}{TimeUnit}\PY{p}{.}\PY{n+na}{MILLISECONDS}\PY{p}{)}\PY{p}{)}\PY{+w}{ }\PY{p}{\PYZob{}}
\PY{+w}{        }\PY{k}{return}\PY{+w}{ }\PY{k+kc}{false}\PY{p}{;}\PY{+w}{ }\PY{c+c1}{// couldn\PYZsq{}t get the lock in time — back off instead of deadlocking}
\PY{+w}{    }\PY{p}{\PYZcb{}}
\PY{+w}{    }\PY{k}{try}\PY{+w}{ }\PY{p}{\PYZob{}}
\PY{+w}{        }\PY{n}{balance}\PY{+w}{ }\PY{o}{\PYZhy{}}\PY{o}{=}\PY{+w}{ }\PY{n}{amount}\PY{p}{;}
\PY{+w}{        }\PY{k}{return}\PY{+w}{ }\PY{k+kc}{true}\PY{p}{;}
\PY{+w}{    }\PY{p}{\PYZcb{}}\PY{+w}{ }\PY{k}{finally}\PY{+w}{ }\PY{p}{\PYZob{}}
\PY{+w}{        }\PY{n}{lock}\PY{p}{.}\PY{n+na}{unlock}\PY{p}{(}\PY{p}{)}\PY{p}{;}
\PY{+w}{    }\PY{p}{\PYZcb{}}
\PY{p}{\PYZcb{}}
\end{Verbatim}
\end{codebox}

\textbf{\code{lockInterruptibly()}.} Acquire the lock, but allow this thread to be interrupted while it’s waiting — useful when a blocked-on-lock thread still needs to respond to cancellation:

\begin{codebox}
\begin{Verbatim}[commandchars=\\\{\}]
\PY{k+kd}{public}\PY{+w}{ }\PY{k+kt}{void}\PY{+w}{ }\PY{n+nf}{withdrawOrCancel}\PY{p}{(}\PY{k+kt}{int}\PY{+w}{ }\PY{n}{amount}\PY{p}{)}\PY{+w}{ }\PY{k+kd}{throws}\PY{+w}{ }\PY{n}{InterruptedException}\PY{+w}{ }\PY{p}{\PYZob{}}
\PY{+w}{    }\PY{n}{lock}\PY{p}{.}\PY{n+na}{lockInterruptibly}\PY{p}{(}\PY{p}{)}\PY{p}{;}\PY{+w}{ }\PY{c+c1}{// throws InterruptedException instead of blocking forever}
\PY{+w}{    }\PY{k}{try}\PY{+w}{ }\PY{p}{\PYZob{}}
\PY{+w}{        }\PY{n}{balance}\PY{+w}{ }\PY{o}{\PYZhy{}}\PY{o}{=}\PY{+w}{ }\PY{n}{amount}\PY{p}{;}
\PY{+w}{    }\PY{p}{\PYZcb{}}\PY{+w}{ }\PY{k}{finally}\PY{+w}{ }\PY{p}{\PYZob{}}
\PY{+w}{        }\PY{n}{lock}\PY{p}{.}\PY{n+na}{unlock}\PY{p}{(}\PY{p}{)}\PY{p}{;}
\PY{+w}{    }\PY{p}{\PYZcb{}}
\PY{p}{\PYZcb{}}
\end{Verbatim}
\end{codebox}

\textbf{\code{Condition} — like \code{wait}/\code{notify}, but per-lock and multiple per lock.} \code{synchronized} gives every object exactly one implicit wait-set. A \code{ReentrantLock} can have several independent \code{Condition}s, which is handy for producer/consumer designs where “buffer is full” and “buffer is empty” are logically distinct waits:

\begin{codebox}
\begin{Verbatim}[commandchars=\\\{\}]
\PY{k+kn}{import}\PY{+w}{ }\PY{n+nn}{java.util.concurrent.locks.*}\PY{p}{;}
\PY{k+kn}{import}\PY{+w}{ }\PY{n+nn}{java.util.*}\PY{p}{;}

\PY{k+kd}{public}\PY{+w}{ }\PY{k+kd}{class} \PY{n+nc}{BoundedBufferWithLock}\PY{o}{\PYZlt{}}\PY{n}{T}\PY{o}{\PYZgt{}}\PY{+w}{ }\PY{p}{\PYZob{}}
\PY{+w}{    }\PY{k+kd}{private}\PY{+w}{ }\PY{k+kd}{final}\PY{+w}{ }\PY{n}{ReentrantLock}\PY{+w}{ }\PY{n}{lock}\PY{+w}{ }\PY{o}{=}\PY{+w}{ }\PY{k}{new}\PY{+w}{ }\PY{n}{ReentrantLock}\PY{p}{(}\PY{p}{)}\PY{p}{;}
\PY{+w}{    }\PY{k+kd}{private}\PY{+w}{ }\PY{k+kd}{final}\PY{+w}{ }\PY{n}{Condition}\PY{+w}{ }\PY{n}{notFull}\PY{+w}{ }\PY{o}{=}\PY{+w}{ }\PY{n}{lock}\PY{p}{.}\PY{n+na}{newCondition}\PY{p}{(}\PY{p}{)}\PY{p}{;}
\PY{+w}{    }\PY{k+kd}{private}\PY{+w}{ }\PY{k+kd}{final}\PY{+w}{ }\PY{n}{Condition}\PY{+w}{ }\PY{n}{notEmpty}\PY{+w}{ }\PY{o}{=}\PY{+w}{ }\PY{n}{lock}\PY{p}{.}\PY{n+na}{newCondition}\PY{p}{(}\PY{p}{)}\PY{p}{;}
\PY{+w}{    }\PY{k+kd}{private}\PY{+w}{ }\PY{k+kd}{final}\PY{+w}{ }\PY{n}{Queue}\PY{o}{\PYZlt{}}\PY{n}{T}\PY{o}{\PYZgt{}}\PY{+w}{ }\PY{n}{queue}\PY{+w}{ }\PY{o}{=}\PY{+w}{ }\PY{k}{new}\PY{+w}{ }\PY{n}{LinkedList}\PY{o}{\PYZlt{}}\PY{o}{\PYZgt{}}\PY{p}{(}\PY{p}{)}\PY{p}{;}
\PY{+w}{    }\PY{k+kd}{private}\PY{+w}{ }\PY{k+kd}{final}\PY{+w}{ }\PY{k+kt}{int}\PY{+w}{ }\PY{n}{capacity}\PY{p}{;}

\PY{+w}{    }\PY{k+kd}{public}\PY{+w}{ }\PY{n+nf}{BoundedBufferWithLock}\PY{p}{(}\PY{k+kt}{int}\PY{+w}{ }\PY{n}{capacity}\PY{p}{)}\PY{+w}{ }\PY{p}{\PYZob{}}\PY{+w}{ }\PY{k}{this}\PY{p}{.}\PY{n+na}{capacity}\PY{+w}{ }\PY{o}{=}\PY{+w}{ }\PY{n}{capacity}\PY{p}{;}\PY{+w}{ }\PY{p}{\PYZcb{}}

\PY{+w}{    }\PY{k+kd}{public}\PY{+w}{ }\PY{k+kt}{void}\PY{+w}{ }\PY{n+nf}{put}\PY{p}{(}\PY{n}{T}\PY{+w}{ }\PY{n}{item}\PY{p}{)}\PY{+w}{ }\PY{k+kd}{throws}\PY{+w}{ }\PY{n}{InterruptedException}\PY{+w}{ }\PY{p}{\PYZob{}}
\PY{+w}{        }\PY{n}{lock}\PY{p}{.}\PY{n+na}{lock}\PY{p}{(}\PY{p}{)}\PY{p}{;}
\PY{+w}{        }\PY{k}{try}\PY{+w}{ }\PY{p}{\PYZob{}}
\PY{+w}{            }\PY{k}{while}\PY{+w}{ }\PY{p}{(}\PY{n}{queue}\PY{p}{.}\PY{n+na}{size}\PY{p}{(}\PY{p}{)}\PY{+w}{ }\PY{o}{=}\PY{o}{=}\PY{+w}{ }\PY{n}{capacity}\PY{p}{)}\PY{+w}{ }\PY{n}{notFull}\PY{p}{.}\PY{n+na}{await}\PY{p}{(}\PY{p}{)}\PY{p}{;}\PY{+w}{ }\PY{c+c1}{// like wait(), but on this Condition}
\PY{+w}{            }\PY{n}{queue}\PY{p}{.}\PY{n+na}{add}\PY{p}{(}\PY{n}{item}\PY{p}{)}\PY{p}{;}
\PY{+w}{            }\PY{n}{notEmpty}\PY{p}{.}\PY{n+na}{signal}\PY{p}{(}\PY{p}{)}\PY{p}{;}\PY{+w}{ }\PY{c+c1}{// like notify(), but only wakes threads waiting on notEmpty}
\PY{+w}{        }\PY{p}{\PYZcb{}}\PY{+w}{ }\PY{k}{finally}\PY{+w}{ }\PY{p}{\PYZob{}}
\PY{+w}{            }\PY{n}{lock}\PY{p}{.}\PY{n+na}{unlock}\PY{p}{(}\PY{p}{)}\PY{p}{;}
\PY{+w}{        }\PY{p}{\PYZcb{}}
\PY{+w}{    }\PY{p}{\PYZcb{}}

\PY{+w}{    }\PY{k+kd}{public}\PY{+w}{ }\PY{n}{T}\PY{+w}{ }\PY{n+nf}{take}\PY{p}{(}\PY{p}{)}\PY{+w}{ }\PY{k+kd}{throws}\PY{+w}{ }\PY{n}{InterruptedException}\PY{+w}{ }\PY{p}{\PYZob{}}
\PY{+w}{        }\PY{n}{lock}\PY{p}{.}\PY{n+na}{lock}\PY{p}{(}\PY{p}{)}\PY{p}{;}
\PY{+w}{        }\PY{k}{try}\PY{+w}{ }\PY{p}{\PYZob{}}
\PY{+w}{            }\PY{k}{while}\PY{+w}{ }\PY{p}{(}\PY{n}{queue}\PY{p}{.}\PY{n+na}{isEmpty}\PY{p}{(}\PY{p}{)}\PY{p}{)}\PY{+w}{ }\PY{n}{notEmpty}\PY{p}{.}\PY{n+na}{await}\PY{p}{(}\PY{p}{)}\PY{p}{;}
\PY{+w}{            }\PY{n}{T}\PY{+w}{ }\PY{n}{item}\PY{+w}{ }\PY{o}{=}\PY{+w}{ }\PY{n}{queue}\PY{p}{.}\PY{n+na}{remove}\PY{p}{(}\PY{p}{)}\PY{p}{;}
\PY{+w}{            }\PY{n}{notFull}\PY{p}{.}\PY{n+na}{signal}\PY{p}{(}\PY{p}{)}\PY{p}{;}
\PY{+w}{            }\PY{k}{return}\PY{+w}{ }\PY{n}{item}\PY{p}{;}
\PY{+w}{        }\PY{p}{\PYZcb{}}\PY{+w}{ }\PY{k}{finally}\PY{+w}{ }\PY{p}{\PYZob{}}
\PY{+w}{            }\PY{n}{lock}\PY{p}{.}\PY{n+na}{unlock}\PY{p}{(}\PY{p}{)}\PY{p}{;}
\PY{+w}{        }\PY{p}{\PYZcb{}}
\PY{+w}{    }\PY{p}{\PYZcb{}}
\PY{p}{\PYZcb{}}
\end{Verbatim}
\end{codebox}

The same \code{while}-loop rule from intrinsic locks applies here too: always re-check the condition after \code{await()} returns.

\section[ReadWriteLock]{ReadWriteLock}
\label{h:13:readwritelock}
Many workloads are read-heavy: lots of threads just want to read shared state, and only occasionally does a thread need to write it. A plain mutual-exclusion lock forces even simultaneous \emph{reads} to queue up one at a time, which wastes parallelism. \code{ReadWriteLock} splits the lock into two: any number of readers can hold the read lock at once, but a writer needs exclusive access — no readers and no other writer at the same time.

\begin{codebox}
\begin{Verbatim}[commandchars=\\\{\}]
\PY{k+kn}{import}\PY{+w}{ }\PY{n+nn}{java.util.concurrent.locks.*}\PY{p}{;}
\PY{k+kn}{import}\PY{+w}{ }\PY{n+nn}{java.util.*}\PY{p}{;}

\PY{k+kd}{public}\PY{+w}{ }\PY{k+kd}{class} \PY{n+nc}{Cache}\PY{o}{\PYZlt{}}\PY{n}{K}\PY{p}{,}\PY{+w}{ }\PY{n}{V}\PY{o}{\PYZgt{}}\PY{+w}{ }\PY{p}{\PYZob{}}
\PY{+w}{    }\PY{k+kd}{private}\PY{+w}{ }\PY{k+kd}{final}\PY{+w}{ }\PY{n}{Map}\PY{o}{\PYZlt{}}\PY{n}{K}\PY{p}{,}\PY{+w}{ }\PY{n}{V}\PY{o}{\PYZgt{}}\PY{+w}{ }\PY{n}{map}\PY{+w}{ }\PY{o}{=}\PY{+w}{ }\PY{k}{new}\PY{+w}{ }\PY{n}{HashMap}\PY{o}{\PYZlt{}}\PY{o}{\PYZgt{}}\PY{p}{(}\PY{p}{)}\PY{p}{;}
\PY{+w}{    }\PY{k+kd}{private}\PY{+w}{ }\PY{k+kd}{final}\PY{+w}{ }\PY{n}{ReadWriteLock}\PY{+w}{ }\PY{n}{rwLock}\PY{+w}{ }\PY{o}{=}\PY{+w}{ }\PY{k}{new}\PY{+w}{ }\PY{n}{ReentrantReadWriteLock}\PY{p}{(}\PY{p}{)}\PY{p}{;}
\PY{+w}{    }\PY{k+kd}{private}\PY{+w}{ }\PY{k+kd}{final}\PY{+w}{ }\PY{n}{Lock}\PY{+w}{ }\PY{n}{readLock}\PY{+w}{ }\PY{o}{=}\PY{+w}{ }\PY{n}{rwLock}\PY{p}{.}\PY{n+na}{readLock}\PY{p}{(}\PY{p}{)}\PY{p}{;}
\PY{+w}{    }\PY{k+kd}{private}\PY{+w}{ }\PY{k+kd}{final}\PY{+w}{ }\PY{n}{Lock}\PY{+w}{ }\PY{n}{writeLock}\PY{+w}{ }\PY{o}{=}\PY{+w}{ }\PY{n}{rwLock}\PY{p}{.}\PY{n+na}{writeLock}\PY{p}{(}\PY{p}{)}\PY{p}{;}

\PY{+w}{    }\PY{k+kd}{public}\PY{+w}{ }\PY{n}{V}\PY{+w}{ }\PY{n+nf}{get}\PY{p}{(}\PY{n}{K}\PY{+w}{ }\PY{n}{key}\PY{p}{)}\PY{+w}{ }\PY{p}{\PYZob{}}
\PY{+w}{        }\PY{n}{readLock}\PY{p}{.}\PY{n+na}{lock}\PY{p}{(}\PY{p}{)}\PY{p}{;}\PY{+w}{ }\PY{c+c1}{// many threads can hold this simultaneously}
\PY{+w}{        }\PY{k}{try}\PY{+w}{ }\PY{p}{\PYZob{}}
\PY{+w}{            }\PY{k}{return}\PY{+w}{ }\PY{n}{map}\PY{p}{.}\PY{n+na}{get}\PY{p}{(}\PY{n}{key}\PY{p}{)}\PY{p}{;}
\PY{+w}{        }\PY{p}{\PYZcb{}}\PY{+w}{ }\PY{k}{finally}\PY{+w}{ }\PY{p}{\PYZob{}}
\PY{+w}{            }\PY{n}{readLock}\PY{p}{.}\PY{n+na}{unlock}\PY{p}{(}\PY{p}{)}\PY{p}{;}
\PY{+w}{        }\PY{p}{\PYZcb{}}
\PY{+w}{    }\PY{p}{\PYZcb{}}

\PY{+w}{    }\PY{k+kd}{public}\PY{+w}{ }\PY{k+kt}{void}\PY{+w}{ }\PY{n+nf}{put}\PY{p}{(}\PY{n}{K}\PY{+w}{ }\PY{n}{key}\PY{p}{,}\PY{+w}{ }\PY{n}{V}\PY{+w}{ }\PY{n}{value}\PY{p}{)}\PY{+w}{ }\PY{p}{\PYZob{}}
\PY{+w}{        }\PY{n}{writeLock}\PY{p}{.}\PY{n+na}{lock}\PY{p}{(}\PY{p}{)}\PY{p}{;}\PY{+w}{ }\PY{c+c1}{// exclusive: blocks all readers and other writers}
\PY{+w}{        }\PY{k}{try}\PY{+w}{ }\PY{p}{\PYZob{}}
\PY{+w}{            }\PY{n}{map}\PY{p}{.}\PY{n+na}{put}\PY{p}{(}\PY{n}{key}\PY{p}{,}\PY{+w}{ }\PY{n}{value}\PY{p}{)}\PY{p}{;}
\PY{+w}{        }\PY{p}{\PYZcb{}}\PY{+w}{ }\PY{k}{finally}\PY{+w}{ }\PY{p}{\PYZob{}}
\PY{+w}{            }\PY{n}{writeLock}\PY{p}{.}\PY{n+na}{unlock}\PY{p}{(}\PY{p}{)}\PY{p}{;}
\PY{+w}{        }\PY{p}{\PYZcb{}}
\PY{+w}{    }\PY{p}{\PYZcb{}}
\PY{p}{\PYZcb{}}
\end{Verbatim}
\end{codebox}

\textbf{When it helps:} read-mostly data structures — caches, configuration snapshots, lookup tables — where writes are rare and reads vastly outnumber them. \textbf{When it doesn’t:} if writes are frequent, or reads are extremely short (a few nanoseconds of work), the bookkeeping overhead of tracking readers/writers can cost more than the parallelism gains — a plain \code{ReentrantLock} or even \code{synchronized} may be faster in practice. Always measure before reaching for this.

\section[StampedLock]{StampedLock}
\label{h:13:stampedlock}
\code{StampedLock} (Java 8+) is a further evolution: it supports the same read/write split as \code{ReadWriteLock}, plus a third mode — \textbf{optimistic reads} — that don’t block writers at all and cost almost nothing when there’s no contention. Instead of an object-based lock, it hands back a \code{long} \textbf{stamp} that you must validate afterward.

\begin{codebox}
\begin{Verbatim}[commandchars=\\\{\}]
\PY{k+kn}{import}\PY{+w}{ }\PY{n+nn}{java.util.concurrent.locks.StampedLock}\PY{p}{;}

\PY{k+kd}{public}\PY{+w}{ }\PY{k+kd}{class} \PY{n+nc}{Point}\PY{+w}{ }\PY{p}{\PYZob{}}
\PY{+w}{    }\PY{k+kd}{private}\PY{+w}{ }\PY{k+kd}{final}\PY{+w}{ }\PY{n}{StampedLock}\PY{+w}{ }\PY{n}{stampedLock}\PY{+w}{ }\PY{o}{=}\PY{+w}{ }\PY{k}{new}\PY{+w}{ }\PY{n}{StampedLock}\PY{p}{(}\PY{p}{)}\PY{p}{;}
\PY{+w}{    }\PY{k+kd}{private}\PY{+w}{ }\PY{k+kt}{double}\PY{+w}{ }\PY{n}{x}\PY{p}{,}\PY{+w}{ }\PY{n}{y}\PY{p}{;}

\PY{+w}{    }\PY{k+kd}{public}\PY{+w}{ }\PY{k+kt}{void}\PY{+w}{ }\PY{n+nf}{move}\PY{p}{(}\PY{k+kt}{double}\PY{+w}{ }\PY{n}{dx}\PY{p}{,}\PY{+w}{ }\PY{k+kt}{double}\PY{+w}{ }\PY{n}{dy}\PY{p}{)}\PY{+w}{ }\PY{p}{\PYZob{}}
\PY{+w}{        }\PY{k+kt}{long}\PY{+w}{ }\PY{n}{stamp}\PY{+w}{ }\PY{o}{=}\PY{+w}{ }\PY{n}{stampedLock}\PY{p}{.}\PY{n+na}{writeLock}\PY{p}{(}\PY{p}{)}\PY{p}{;}\PY{+w}{ }\PY{c+c1}{// exclusive, like a normal write lock}
\PY{+w}{        }\PY{k}{try}\PY{+w}{ }\PY{p}{\PYZob{}}
\PY{+w}{            }\PY{n}{x}\PY{+w}{ }\PY{o}{+}\PY{o}{=}\PY{+w}{ }\PY{n}{dx}\PY{p}{;}
\PY{+w}{            }\PY{n}{y}\PY{+w}{ }\PY{o}{+}\PY{o}{=}\PY{+w}{ }\PY{n}{dy}\PY{p}{;}
\PY{+w}{        }\PY{p}{\PYZcb{}}\PY{+w}{ }\PY{k}{finally}\PY{+w}{ }\PY{p}{\PYZob{}}
\PY{+w}{            }\PY{n}{stampedLock}\PY{p}{.}\PY{n+na}{unlockWrite}\PY{p}{(}\PY{n}{stamp}\PY{p}{)}\PY{p}{;}
\PY{+w}{        }\PY{p}{\PYZcb{}}
\PY{+w}{    }\PY{p}{\PYZcb{}}

\PY{+w}{    }\PY{k+kd}{public}\PY{+w}{ }\PY{k+kt}{double}\PY{+w}{ }\PY{n+nf}{distanceFromOrigin}\PY{p}{(}\PY{p}{)}\PY{+w}{ }\PY{p}{\PYZob{}}
\PY{+w}{        }\PY{k+kt}{long}\PY{+w}{ }\PY{n}{stamp}\PY{+w}{ }\PY{o}{=}\PY{+w}{ }\PY{n}{stampedLock}\PY{p}{.}\PY{n+na}{tryOptimisticRead}\PY{p}{(}\PY{p}{)}\PY{p}{;}\PY{+w}{ }\PY{c+c1}{// does NOT block, does NOT take a real lock}
\PY{+w}{        }\PY{k+kt}{double}\PY{+w}{ }\PY{n}{currentX}\PY{+w}{ }\PY{o}{=}\PY{+w}{ }\PY{n}{x}\PY{p}{,}\PY{+w}{ }\PY{n}{currentY}\PY{+w}{ }\PY{o}{=}\PY{+w}{ }\PY{n}{y}\PY{p}{;}\PY{+w}{             }\PY{c+c1}{// read the fields hopefully}
\PY{+w}{        }\PY{k}{if}\PY{+w}{ }\PY{p}{(}\PY{o}{!}\PY{n}{stampedLock}\PY{p}{.}\PY{n+na}{validate}\PY{p}{(}\PY{n}{stamp}\PY{p}{)}\PY{p}{)}\PY{+w}{ }\PY{p}{\PYZob{}}\PY{+w}{             }\PY{c+c1}{// did a writer sneak in while we read?}
\PY{+w}{            }\PY{n}{stamp}\PY{+w}{ }\PY{o}{=}\PY{+w}{ }\PY{n}{stampedLock}\PY{p}{.}\PY{n+na}{readLock}\PY{p}{(}\PY{p}{)}\PY{p}{;}\PY{+w}{             }\PY{c+c1}{// fall back to a real, blocking read lock}
\PY{+w}{            }\PY{k}{try}\PY{+w}{ }\PY{p}{\PYZob{}}
\PY{+w}{                }\PY{n}{currentX}\PY{+w}{ }\PY{o}{=}\PY{+w}{ }\PY{n}{x}\PY{p}{;}
\PY{+w}{                }\PY{n}{currentY}\PY{+w}{ }\PY{o}{=}\PY{+w}{ }\PY{n}{y}\PY{p}{;}
\PY{+w}{            }\PY{p}{\PYZcb{}}\PY{+w}{ }\PY{k}{finally}\PY{+w}{ }\PY{p}{\PYZob{}}
\PY{+w}{                }\PY{n}{stampedLock}\PY{p}{.}\PY{n+na}{unlockRead}\PY{p}{(}\PY{n}{stamp}\PY{p}{)}\PY{p}{;}
\PY{+w}{            }\PY{p}{\PYZcb{}}
\PY{+w}{        }\PY{p}{\PYZcb{}}
\PY{+w}{        }\PY{k}{return}\PY{+w}{ }\PY{n}{Math}\PY{p}{.}\PY{n+na}{sqrt}\PY{p}{(}\PY{n}{currentX}\PY{+w}{ }\PY{o}{*}\PY{+w}{ }\PY{n}{currentX}\PY{+w}{ }\PY{o}{+}\PY{+w}{ }\PY{n}{currentY}\PY{+w}{ }\PY{o}{*}\PY{+w}{ }\PY{n}{currentY}\PY{p}{)}\PY{p}{;}
\PY{+w}{    }\PY{p}{\PYZcb{}}
\PY{p}{\PYZcb{}}
\end{Verbatim}
\end{codebox}

The optimistic-read pattern is always: \textbf{grab a stamp → read the data → validate the stamp → if invalid, retry with a real lock.} Skipping the validation step is a serious bug — it silently allows reading data that a writer changed mid-read.

\textbf{Critical gotcha: \code{StampedLock} is NOT reentrant.} Unlike \code{ReentrantLock} and intrinsic locks, calling \code{writeLock()} again from the same thread that already holds it will simply deadlock the thread against itself:

\begin{codebox}
\begin{Verbatim}[commandchars=\\\{\}]
\PY{k+kt}{long}\PY{+w}{ }\PY{n}{stamp}\PY{+w}{ }\PY{o}{=}\PY{+w}{ }\PY{n}{stampedLock}\PY{p}{.}\PY{n+na}{writeLock}\PY{p}{(}\PY{p}{)}\PY{p}{;}
\PY{k}{try}\PY{+w}{ }\PY{p}{\PYZob{}}
\PY{+w}{    }\PY{c+c1}{// BUG: if this method (directly or indirectly) tries to acquire}
\PY{+w}{    }\PY{c+c1}{// stampedLock.writeLock() again on the SAME thread, it blocks forever —}
\PY{+w}{    }\PY{c+c1}{// StampedLock does not track \PYZdq{}who already owns this.\PYZdq{}}
\PY{+w}{    }\PY{n}{doSomethingThatAlsoLocks}\PY{p}{(}\PY{p}{)}\PY{p}{;}
\PY{p}{\PYZcb{}}\PY{+w}{ }\PY{k}{finally}\PY{+w}{ }\PY{p}{\PYZob{}}
\PY{+w}{    }\PY{n}{stampedLock}\PY{p}{.}\PY{n+na}{unlockWrite}\PY{p}{(}\PY{n}{stamp}\PY{p}{)}\PY{p}{;}
\PY{p}{\PYZcb{}}
\end{Verbatim}
\end{codebox}

Because of this, be very careful using \code{StampedLock} in any code path that might recurse or call back into itself. It also has no \code{Condition} support and doesn’t play with \code{synchronized}-style monitor semantics — it’s a specialized tool for exactly the read-heavy, low-contention scenario it’s built for.

\section[Volatile]{Volatile}
\label{h:13:volatile}
\code{volatile} is a field modifier that gives two guarantees, and only two: \textbf{visibility} and a limited form of \textbf{ordering}. It does \textbf{not} give atomicity.

\begin{itemize}
\item 
\textbf{Visibility}: a write to a \code{volatile} field by one thread is guaranteed to be immediately visible to any thread that subsequently reads it. Without \code{volatile}, a reading thread might keep seeing a stale, cached value indefinitely (in theory — real JITs/CPUs vary, but the JMM makes no promise otherwise).

\item 
\textbf{Ordering}: the compiler and CPU cannot reorder other instructions across a \code{volatile} read/write in ways that would break the visible sequence (this feeds directly into the happens-before rules below).

\item 
\textbf{What it does NOT give you: atomicity.} \code{count++} on a \code{volatile~\allowbreak{}int} is still read-modify-write — three separate steps — and two threads can interleave those steps and lose an update, exactly like a non-volatile field.

\end{itemize}

\begin{codebox}
\begin{Verbatim}[commandchars=\\\{\}]
\PY{k+kd}{public}\PY{+w}{ }\PY{k+kd}{class} \PY{n+nc}{VolatileIsNotAtomic}\PY{+w}{ }\PY{p}{\PYZob{}}
\PY{+w}{    }\PY{k+kd}{private}\PY{+w}{ }\PY{k+kd}{volatile}\PY{+w}{ }\PY{k+kt}{int}\PY{+w}{ }\PY{n}{count}\PY{+w}{ }\PY{o}{=}\PY{+w}{ }\PY{l+m+mi}{0}\PY{p}{;}

\PY{+w}{    }\PY{k+kd}{public}\PY{+w}{ }\PY{k+kt}{void}\PY{+w}{ }\PY{n+nf}{increment}\PY{p}{(}\PY{p}{)}\PY{+w}{ }\PY{p}{\PYZob{}}
\PY{+w}{        }\PY{n}{count}\PY{o}{+}\PY{o}{+}\PY{p}{;}\PY{+w}{ }\PY{c+c1}{// READS count, ADDS 1, WRITES count — three steps, not one}
\PY{+w}{    }\PY{p}{\PYZcb{}}
\PY{p}{\PYZcb{}}
\PY{c+c1}{// With 2 threads each calling increment() 100,000 times, the final count}
\PY{c+c1}{// is very likely LESS than 200,000 — some increments are silently lost,}
\PY{c+c1}{// even though \PYZsq{}volatile\PYZsq{} is present.}
\end{Verbatim}
\end{codebox}

\textbf{The classic correct use case: a stop flag.}

\begin{codebox}
\begin{Verbatim}[commandchars=\\\{\}]
\PY{k+kd}{public}\PY{+w}{ }\PY{k+kd}{class} \PY{n+nc}{StopFlagExample}\PY{+w}{ }\PY{p}{\PYZob{}}
\PY{+w}{    }\PY{k+kd}{private}\PY{+w}{ }\PY{k+kd}{volatile}\PY{+w}{ }\PY{k+kt}{boolean}\PY{+w}{ }\PY{n}{running}\PY{+w}{ }\PY{o}{=}\PY{+w}{ }\PY{k+kc}{true}\PY{p}{;}\PY{+w}{ }\PY{c+c1}{// visibility is ALL this needs}

\PY{+w}{    }\PY{k+kd}{public}\PY{+w}{ }\PY{k+kt}{void}\PY{+w}{ }\PY{n+nf}{doWork}\PY{p}{(}\PY{p}{)}\PY{+w}{ }\PY{p}{\PYZob{}}
\PY{+w}{        }\PY{k}{while}\PY{+w}{ }\PY{p}{(}\PY{n}{running}\PY{p}{)}\PY{+w}{ }\PY{p}{\PYZob{}}
\PY{+w}{            }\PY{c+c1}{// do a unit of work}
\PY{+w}{        }\PY{p}{\PYZcb{}}
\PY{+w}{        }\PY{n}{System}\PY{p}{.}\PY{n+na}{out}\PY{p}{.}\PY{n+na}{println}\PY{p}{(}\PY{l+s}{\PYZdq{}}\PY{l+s}{Stopped cleanly}\PY{l+s}{\PYZdq{}}\PY{p}{)}\PY{p}{;}
\PY{+w}{    }\PY{p}{\PYZcb{}}

\PY{+w}{    }\PY{k+kd}{public}\PY{+w}{ }\PY{k+kt}{void}\PY{+w}{ }\PY{n+nf}{stop}\PY{p}{(}\PY{p}{)}\PY{+w}{ }\PY{p}{\PYZob{}}
\PY{+w}{        }\PY{n}{running}\PY{+w}{ }\PY{o}{=}\PY{+w}{ }\PY{k+kc}{false}\PY{p}{;}\PY{+w}{ }\PY{c+c1}{// guaranteed to become visible to the worker thread promptly}
\PY{+w}{    }\PY{p}{\PYZcb{}}
\PY{p}{\PYZcb{}}
\end{Verbatim}
\end{codebox}

This works with plain \code{volatile} because it’s a single write and a single read of one variable — no read-modify-write, so no lost-update risk. Swap \code{volatile~\allowbreak{}boolean} for a plain \code{boolean} here and the worker thread might loop forever, never observing the change (a real, documented class of bug).

\textbf{Double-checked locking, done right.} A common (broken) attempt at lazy, thread-safe singleton initialization:

\begin{codebox}
\begin{Verbatim}[commandchars=\\\{\}]
\PY{c+c1}{// BROKEN: without \PYZsq{}volatile\PYZsq{}, another thread can see a half\PYZhy{}constructed object}
\PY{k+kd}{public}\PY{+w}{ }\PY{k+kd}{class} \PY{n+nc}{BrokenSingleton}\PY{+w}{ }\PY{p}{\PYZob{}}
\PY{+w}{    }\PY{k+kd}{private}\PY{+w}{ }\PY{k+kd}{static}\PY{+w}{ }\PY{n}{BrokenSingleton}\PY{+w}{ }\PY{n}{instance}\PY{p}{;}

\PY{+w}{    }\PY{k+kd}{public}\PY{+w}{ }\PY{k+kd}{static}\PY{+w}{ }\PY{n}{BrokenSingleton}\PY{+w}{ }\PY{n+nf}{getInstance}\PY{p}{(}\PY{p}{)}\PY{+w}{ }\PY{p}{\PYZob{}}
\PY{+w}{        }\PY{k}{if}\PY{+w}{ }\PY{p}{(}\PY{n}{instance}\PY{+w}{ }\PY{o}{=}\PY{o}{=}\PY{+w}{ }\PY{k+kc}{null}\PY{p}{)}\PY{+w}{ }\PY{p}{\PYZob{}}\PY{+w}{                 }\PY{c+c1}{// first check, no lock — fast path}
\PY{+w}{            }\PY{k+kd}{synchronized}\PY{+w}{ }\PY{p}{(}\PY{n}{BrokenSingleton}\PY{p}{.}\PY{n+na}{class}\PY{p}{)}\PY{+w}{ }\PY{p}{\PYZob{}}
\PY{+w}{                }\PY{k}{if}\PY{+w}{ }\PY{p}{(}\PY{n}{instance}\PY{+w}{ }\PY{o}{=}\PY{o}{=}\PY{+w}{ }\PY{k+kc}{null}\PY{p}{)}\PY{+w}{ }\PY{p}{\PYZob{}}\PY{+w}{          }\PY{c+c1}{// second check, with lock}
\PY{+w}{                    }\PY{n}{instance}\PY{+w}{ }\PY{o}{=}\PY{+w}{ }\PY{k}{new}\PY{+w}{ }\PY{n}{BrokenSingleton}\PY{p}{(}\PY{p}{)}\PY{p}{;}\PY{+w}{ }\PY{c+c1}{// NOT guaranteed fully\PYZhy{}constructed}
\PY{+w}{                    }\PY{c+c1}{// before this reference becomes visible to other threads without \PYZsq{}volatile\PYZsq{}}
\PY{+w}{                }\PY{p}{\PYZcb{}}
\PY{+w}{            }\PY{p}{\PYZcb{}}
\PY{+w}{        }\PY{p}{\PYZcb{}}
\PY{+w}{        }\PY{k}{return}\PY{+w}{ }\PY{n}{instance}\PY{p}{;}
\PY{+w}{    }\PY{p}{\PYZcb{}}
\PY{p}{\PYZcb{}}
\end{Verbatim}
\end{codebox}

The problem: \code{new~\allowbreak{}BrokenSingleton()} isn’t a single atomic step from the JMM’s point of view — it involves allocating memory, running the constructor, and assigning the reference. Without \code{volatile}, those steps can appear reordered to another thread, which could see a non-null \code{instance} reference that points to a not-yet-fully-initialized object.

\begin{codebox}
\begin{Verbatim}[commandchars=\\\{\}]
\PY{c+c1}{// FIXED: \PYZsq{}volatile\PYZsq{} prevents the reordering and publishes a fully\PYZhy{}built object}
\PY{k+kd}{public}\PY{+w}{ }\PY{k+kd}{class} \PY{n+nc}{CorrectSingleton}\PY{+w}{ }\PY{p}{\PYZob{}}
\PY{+w}{    }\PY{k+kd}{private}\PY{+w}{ }\PY{k+kd}{static}\PY{+w}{ }\PY{k+kd}{volatile}\PY{+w}{ }\PY{n}{CorrectSingleton}\PY{+w}{ }\PY{n}{instance}\PY{p}{;}

\PY{+w}{    }\PY{k+kd}{public}\PY{+w}{ }\PY{k+kd}{static}\PY{+w}{ }\PY{n}{CorrectSingleton}\PY{+w}{ }\PY{n+nf}{getInstance}\PY{p}{(}\PY{p}{)}\PY{+w}{ }\PY{p}{\PYZob{}}
\PY{+w}{        }\PY{k}{if}\PY{+w}{ }\PY{p}{(}\PY{n}{instance}\PY{+w}{ }\PY{o}{=}\PY{o}{=}\PY{+w}{ }\PY{k+kc}{null}\PY{p}{)}\PY{+w}{ }\PY{p}{\PYZob{}}
\PY{+w}{            }\PY{k+kd}{synchronized}\PY{+w}{ }\PY{p}{(}\PY{n}{CorrectSingleton}\PY{p}{.}\PY{n+na}{class}\PY{p}{)}\PY{+w}{ }\PY{p}{\PYZob{}}
\PY{+w}{                }\PY{k}{if}\PY{+w}{ }\PY{p}{(}\PY{n}{instance}\PY{+w}{ }\PY{o}{=}\PY{o}{=}\PY{+w}{ }\PY{k+kc}{null}\PY{p}{)}\PY{+w}{ }\PY{p}{\PYZob{}}
\PY{+w}{                    }\PY{n}{instance}\PY{+w}{ }\PY{o}{=}\PY{+w}{ }\PY{k}{new}\PY{+w}{ }\PY{n}{CorrectSingleton}\PY{p}{(}\PY{p}{)}\PY{p}{;}
\PY{+w}{                }\PY{p}{\PYZcb{}}
\PY{+w}{            }\PY{p}{\PYZcb{}}
\PY{+w}{        }\PY{p}{\PYZcb{}}
\PY{+w}{        }\PY{k}{return}\PY{+w}{ }\PY{n}{instance}\PY{p}{;}
\PY{+w}{    }\PY{p}{\PYZcb{}}
\PY{p}{\PYZcb{}}
\end{Verbatim}
\end{codebox}

(In modern Java, the simplest and safest fix for a plain singleton is usually a static holder class, which relies on class-initialization guarantees instead of double-checked locking at all — but double-checked locking remains a fair-game interview topic.)

\section[Atomic Classes]{Atomic Classes}
\label{h:13:atomic-classes}
The \code{java.\allowbreak{}util.\allowbreak{}concurrent.\allowbreak{}atomic} package gives you classes that perform single-variable read-modify-write operations \textbf{atomically}, using hardware-level compare-and-swap (CAS) instructions instead of locks. That makes them faster than a lock under light-to-moderate contention, with no risk of deadlock (there’s nothing to hold across a blocking call).

\begin{codebox}
\begin{Verbatim}[commandchars=\\\{\}]
\PY{k+kn}{import}\PY{+w}{ }\PY{n+nn}{java.util.concurrent.atomic.*}\PY{p}{;}

\PY{k+kd}{public}\PY{+w}{ }\PY{k+kd}{class} \PY{n+nc}{AtomicBasics}\PY{+w}{ }\PY{p}{\PYZob{}}
\PY{+w}{    }\PY{k+kd}{public}\PY{+w}{ }\PY{k+kd}{static}\PY{+w}{ }\PY{k+kt}{void}\PY{+w}{ }\PY{n+nf}{main}\PY{p}{(}\PY{n}{String}\PY{o}{[}\PY{o}{]}\PY{+w}{ }\PY{n}{args}\PY{p}{)}\PY{+w}{ }\PY{p}{\PYZob{}}
\PY{+w}{        }\PY{n}{AtomicInteger}\PY{+w}{ }\PY{n}{counter}\PY{+w}{ }\PY{o}{=}\PY{+w}{ }\PY{k}{new}\PY{+w}{ }\PY{n}{AtomicInteger}\PY{p}{(}\PY{l+m+mi}{0}\PY{p}{)}\PY{p}{;}
\PY{+w}{        }\PY{n}{counter}\PY{p}{.}\PY{n+na}{incrementAndGet}\PY{p}{(}\PY{p}{)}\PY{p}{;}\PY{+w}{          }\PY{c+c1}{// atomic ++counter, returns new value}
\PY{+w}{        }\PY{n}{counter}\PY{p}{.}\PY{n+na}{getAndIncrement}\PY{p}{(}\PY{p}{)}\PY{p}{;}\PY{+w}{          }\PY{c+c1}{// atomic counter++, returns old value}
\PY{+w}{        }\PY{n}{counter}\PY{p}{.}\PY{n+na}{addAndGet}\PY{p}{(}\PY{l+m+mi}{5}\PY{p}{)}\PY{p}{;}\PY{+w}{               }\PY{c+c1}{// atomic counter += 5}

\PY{+w}{        }\PY{n}{AtomicLong}\PY{+w}{ }\PY{n}{bigCounter}\PY{+w}{ }\PY{o}{=}\PY{+w}{ }\PY{k}{new}\PY{+w}{ }\PY{n}{AtomicLong}\PY{p}{(}\PY{l+m+mi}{0}\PY{n}{L}\PY{p}{)}\PY{p}{;}
\PY{+w}{        }\PY{n}{bigCounter}\PY{p}{.}\PY{n+na}{incrementAndGet}\PY{p}{(}\PY{p}{)}\PY{p}{;}

\PY{+w}{        }\PY{n}{AtomicReference}\PY{o}{\PYZlt{}}\PY{n}{String}\PY{o}{\PYZgt{}}\PY{+w}{ }\PY{n}{ref}\PY{+w}{ }\PY{o}{=}\PY{+w}{ }\PY{k}{new}\PY{+w}{ }\PY{n}{AtomicReference}\PY{o}{\PYZlt{}}\PY{o}{\PYZgt{}}\PY{p}{(}\PY{l+s}{\PYZdq{}}\PY{l+s}{initial}\PY{l+s}{\PYZdq{}}\PY{p}{)}\PY{p}{;}
\PY{+w}{        }\PY{k+kt}{boolean}\PY{+w}{ }\PY{n}{swapped}\PY{+w}{ }\PY{o}{=}\PY{+w}{ }\PY{n}{ref}\PY{p}{.}\PY{n+na}{compareAndSet}\PY{p}{(}\PY{l+s}{\PYZdq{}}\PY{l+s}{initial}\PY{l+s}{\PYZdq{}}\PY{p}{,}\PY{+w}{ }\PY{l+s}{\PYZdq{}}\PY{l+s}{updated}\PY{l+s}{\PYZdq{}}\PY{p}{)}\PY{p}{;}\PY{+w}{ }\PY{c+c1}{// CAS: only if still \PYZdq{}initial\PYZdq{}}
\PY{+w}{        }\PY{n}{System}\PY{p}{.}\PY{n+na}{out}\PY{p}{.}\PY{n+na}{println}\PY{p}{(}\PY{n}{swapped}\PY{p}{)}\PY{p}{;}\PY{+w}{        }\PY{c+c1}{// true}
\PY{+w}{        }\PY{n}{System}\PY{p}{.}\PY{n+na}{out}\PY{p}{.}\PY{n+na}{println}\PY{p}{(}\PY{n}{ref}\PY{p}{.}\PY{n+na}{get}\PY{p}{(}\PY{p}{)}\PY{p}{)}\PY{p}{;}\PY{+w}{      }\PY{c+c1}{// \PYZdq{}updated\PYZdq{}}

\PY{+w}{        }\PY{c+c1}{// updateAndGet: apply a function atomically, retrying under contention}
\PY{+w}{        }\PY{n}{AtomicInteger}\PY{+w}{ }\PY{n}{price}\PY{+w}{ }\PY{o}{=}\PY{+w}{ }\PY{k}{new}\PY{+w}{ }\PY{n}{AtomicInteger}\PY{p}{(}\PY{l+m+mi}{100}\PY{p}{)}\PY{p}{;}
\PY{+w}{        }\PY{n}{price}\PY{p}{.}\PY{n+na}{updateAndGet}\PY{p}{(}\PY{n}{current}\PY{+w}{ }\PY{o}{\PYZhy{}}\PY{o}{\PYZgt{}}\PY{+w}{ }\PY{n}{current}\PY{+w}{ }\PY{o}{+}\PY{+w}{ }\PY{l+m+mi}{10}\PY{p}{)}\PY{p}{;}\PY{+w}{ }\PY{c+c1}{// atomic \PYZdq{}price += 10\PYZdq{}, CAS\PYZhy{}retry loop internally}
\PY{+w}{    }\PY{p}{\PYZcb{}}
\PY{p}{\PYZcb{}}
\end{Verbatim}
\end{codebox}

\textbf{\code{compareAndSet} (CAS)} is the primitive underneath almost everything here: “set this value to X, but only if it’s still currently Y.” If another thread changed it in between, the CAS fails and the caller (or the atomic class internally, for methods like \code{updateAndGet}) retries.

\textbf{The ABA problem.} CAS only checks “is the value still what I last saw?” — it can’t tell if the value was changed to something else and then changed \emph{back} to the original value while you weren’t looking. If your logic assumes “unchanged” means “nothing happened in between,” ABA breaks that assumption.

\begin{codebox}
\begin{Verbatim}[commandchars=\\\{\}]
\PY{c+c1}{// Thread 1 reads reference, sees \PYZdq{}A\PYZdq{}, plans to CAS \PYZdq{}A\PYZdq{} \PYZhy{}\PYZgt{} \PYZdq{}C\PYZdq{}}
\PY{c+c1}{// Thread 2 (in between) CAS \PYZdq{}A\PYZdq{} \PYZhy{}\PYZgt{} \PYZdq{}B\PYZdq{}, then CAS \PYZdq{}B\PYZdq{} \PYZhy{}\PYZgt{} \PYZdq{}A\PYZdq{}}
\PY{c+c1}{// Thread 1\PYZsq{}s CAS(\PYZdq{}A\PYZdq{}, \PYZdq{}C\PYZdq{}) now succeeds — it thinks nothing changed,}
\PY{c+c1}{// but state actually went A \PYZhy{}\PYZgt{} B \PYZhy{}\PYZgt{} A, and any side effects tied to}
\PY{c+c1}{// that B intermediate state have already happened invisibly.}
\PY{n}{AtomicReference}\PY{o}{\PYZlt{}}\PY{n}{String}\PY{o}{\PYZgt{}}\PY{+w}{ }\PY{n}{state}\PY{+w}{ }\PY{o}{=}\PY{+w}{ }\PY{k}{new}\PY{+w}{ }\PY{n}{AtomicReference}\PY{o}{\PYZlt{}}\PY{o}{\PYZgt{}}\PY{p}{(}\PY{l+s}{\PYZdq{}}\PY{l+s}{A}\PY{l+s}{\PYZdq{}}\PY{p}{)}\PY{p}{;}
\PY{c+c1}{// state.compareAndSet(\PYZdq{}A\PYZdq{}, \PYZdq{}C\PYZdq{}) succeeds even though \PYZdq{}B\PYZdq{} happened in between.}
\end{Verbatim}
\end{codebox}

Mitigation: use a version/stamp alongside the value (\code{AtomicStampedReference}) so a CAS also checks that the “version” hasn’t moved, not just the value.

\textbf{\code{LongAdder} vs \code{AtomicLong} under contention.} \code{AtomicLong} funnels every thread’s CAS attempt at the \emph{same single memory location} — under heavy contention (many threads incrementing constantly), that creates a hot spot where CAS attempts keep failing and retrying. \code{LongAdder} (Java 8+) spreads the count across multiple internal cells and only combines them when you call \code{sum()}, dramatically reducing contention at the cost of \code{sum()} being an eventually-consistent snapshot rather than a single atomic memory location.

\begin{codebox}
\begin{Verbatim}[commandchars=\\\{\}]
\PY{k+kn}{import}\PY{+w}{ }\PY{n+nn}{java.util.concurrent.atomic.LongAdder}\PY{p}{;}

\PY{n}{LongAdder}\PY{+w}{ }\PY{n}{hits}\PY{+w}{ }\PY{o}{=}\PY{+w}{ }\PY{k}{new}\PY{+w}{ }\PY{n}{LongAdder}\PY{p}{(}\PY{p}{)}\PY{p}{;}
\PY{c+c1}{// Many threads, high contention:}
\PY{n}{hits}\PY{p}{.}\PY{n+na}{increment}\PY{p}{(}\PY{p}{)}\PY{p}{;}\PY{+w}{   }\PY{c+c1}{// internally spreads updates across cells — much less CAS contention}
\PY{k+kt}{long}\PY{+w}{ }\PY{n}{total}\PY{+w}{ }\PY{o}{=}\PY{+w}{ }\PY{n}{hits}\PY{p}{.}\PY{n+na}{sum}\PY{p}{(}\PY{p}{)}\PY{p}{;}\PY{+w}{ }\PY{c+c1}{// combines cells; fine for metrics/counters, not for tight CAS logic}
\end{Verbatim}
\end{codebox}

Rule of thumb: use \code{AtomicLong} when you need the exact current value read consistently (e.g., you’re going to \code{compareAndSet} against it). Use \code{LongAdder} for pure counters/metrics under high contention where you only need the total occasionally.

\textbf{\code{Atomic\allowbreak{}Integer\allowbreak{}Field\allowbreak{}Updater} (briefly).} Lets you get atomic CAS-style updates on a plain \code{volatile~\allowbreak{}int} field of an \emph{existing} class, without changing its declared type to \code{AtomicInteger} — useful when you can’t change a field’s type (e.g., it’s part of a serialized format or an inherited class) but still need atomic updates.

\begin{codebox}
\begin{Verbatim}[commandchars=\\\{\}]
\PY{k+kn}{import}\PY{+w}{ }\PY{n+nn}{java.util.concurrent.atomic.AtomicIntegerFieldUpdater}\PY{p}{;}

\PY{k+kd}{class} \PY{n+nc}{Node}\PY{+w}{ }\PY{p}{\PYZob{}}
\PY{+w}{    }\PY{k+kd}{volatile}\PY{+w}{ }\PY{k+kt}{int}\PY{+w}{ }\PY{n}{visits}\PY{+w}{ }\PY{o}{=}\PY{+w}{ }\PY{l+m+mi}{0}\PY{p}{;}\PY{+w}{ }\PY{c+c1}{// must be volatile, not private, and not final}
\PY{p}{\PYZcb{}}

\PY{n}{AtomicIntegerFieldUpdater}\PY{o}{\PYZlt{}}\PY{n}{Node}\PY{o}{\PYZgt{}}\PY{+w}{ }\PY{n}{updater}\PY{+w}{ }\PY{o}{=}
\PY{+w}{        }\PY{n}{AtomicIntegerFieldUpdater}\PY{p}{.}\PY{n+na}{newUpdater}\PY{p}{(}\PY{n}{Node}\PY{p}{.}\PY{n+na}{class}\PY{p}{,}\PY{+w}{ }\PY{l+s}{\PYZdq{}}\PY{l+s}{visits}\PY{l+s}{\PYZdq{}}\PY{p}{)}\PY{p}{;}
\PY{n}{Node}\PY{+w}{ }\PY{n}{node}\PY{+w}{ }\PY{o}{=}\PY{+w}{ }\PY{k}{new}\PY{+w}{ }\PY{n}{Node}\PY{p}{(}\PY{p}{)}\PY{p}{;}
\PY{n}{updater}\PY{p}{.}\PY{n+na}{incrementAndGet}\PY{p}{(}\PY{n}{node}\PY{p}{)}\PY{p}{;}\PY{+w}{ }\PY{c+c1}{// atomic increment on node.visits without an AtomicInteger field}
\end{Verbatim}
\end{codebox}

\section[ThreadLocal]{ThreadLocal}
\label{h:13:threadlocal}
\code{ThreadLocal\textless{}\allowbreak{}T\textgreater{}} gives each thread its own independent copy of a variable. Two threads reading/writing the “same” \code{ThreadLocal} never see each other’s values — it’s per-thread storage, not shared state, so it needs no locking at all.

\begin{codebox}
\begin{Verbatim}[commandchars=\\\{\}]
\PY{k+kd}{public}\PY{+w}{ }\PY{k+kd}{class} \PY{n+nc}{ThreadLocalBasics}\PY{+w}{ }\PY{p}{\PYZob{}}
\PY{+w}{    }\PY{k+kd}{private}\PY{+w}{ }\PY{k+kd}{static}\PY{+w}{ }\PY{k+kd}{final}\PY{+w}{ }\PY{n}{ThreadLocal}\PY{o}{\PYZlt{}}\PY{n}{Integer}\PY{o}{\PYZgt{}}\PY{+w}{ }\PY{n}{requestId}\PY{+w}{ }\PY{o}{=}
\PY{+w}{            }\PY{n}{ThreadLocal}\PY{p}{.}\PY{n+na}{withInitial}\PY{p}{(}\PY{p}{(}\PY{p}{)}\PY{+w}{ }\PY{o}{\PYZhy{}}\PY{o}{\PYZgt{}}\PY{+w}{ }\PY{l+m+mi}{0}\PY{p}{)}\PY{p}{;}

\PY{+w}{    }\PY{k+kd}{public}\PY{+w}{ }\PY{k+kd}{static}\PY{+w}{ }\PY{k+kt}{void}\PY{+w}{ }\PY{n+nf}{main}\PY{p}{(}\PY{n}{String}\PY{o}{[}\PY{o}{]}\PY{+w}{ }\PY{n}{args}\PY{p}{)}\PY{+w}{ }\PY{k+kd}{throws}\PY{+w}{ }\PY{n}{InterruptedException}\PY{+w}{ }\PY{p}{\PYZob{}}
\PY{+w}{        }\PY{n}{Runnable}\PY{+w}{ }\PY{n}{task}\PY{+w}{ }\PY{o}{=}\PY{+w}{ }\PY{p}{(}\PY{p}{)}\PY{+w}{ }\PY{o}{\PYZhy{}}\PY{o}{\PYZgt{}}\PY{+w}{ }\PY{p}{\PYZob{}}
\PY{+w}{            }\PY{n}{requestId}\PY{p}{.}\PY{n+na}{set}\PY{p}{(}\PY{p}{(}\PY{k+kt}{int}\PY{p}{)}\PY{+w}{ }\PY{p}{(}\PY{n}{Math}\PY{p}{.}\PY{n+na}{random}\PY{p}{(}\PY{p}{)}\PY{+w}{ }\PY{o}{*}\PY{+w}{ }\PY{l+m+mi}{1000}\PY{p}{)}\PY{p}{)}\PY{p}{;}
\PY{+w}{            }\PY{n}{System}\PY{p}{.}\PY{n+na}{out}\PY{p}{.}\PY{n+na}{println}\PY{p}{(}\PY{n}{Thread}\PY{p}{.}\PY{n+na}{currentThread}\PY{p}{(}\PY{p}{)}\PY{p}{.}\PY{n+na}{getName}\PY{p}{(}\PY{p}{)}\PY{+w}{ }\PY{o}{+}\PY{+w}{ }\PY{l+s}{\PYZdq{}}\PY{l+s}{ sees }\PY{l+s}{\PYZdq{}}\PY{+w}{ }\PY{o}{+}\PY{+w}{ }\PY{n}{requestId}\PY{p}{.}\PY{n+na}{get}\PY{p}{(}\PY{p}{)}\PY{p}{)}\PY{p}{;}
\PY{+w}{        }\PY{p}{\PYZcb{}}\PY{p}{;}
\PY{+w}{        }\PY{n}{Thread}\PY{+w}{ }\PY{n}{t1}\PY{+w}{ }\PY{o}{=}\PY{+w}{ }\PY{k}{new}\PY{+w}{ }\PY{n}{Thread}\PY{p}{(}\PY{n}{task}\PY{p}{,}\PY{+w}{ }\PY{l+s}{\PYZdq{}}\PY{l+s}{t1}\PY{l+s}{\PYZdq{}}\PY{p}{)}\PY{p}{;}
\PY{+w}{        }\PY{n}{Thread}\PY{+w}{ }\PY{n}{t2}\PY{+w}{ }\PY{o}{=}\PY{+w}{ }\PY{k}{new}\PY{+w}{ }\PY{n}{Thread}\PY{p}{(}\PY{n}{task}\PY{p}{,}\PY{+w}{ }\PY{l+s}{\PYZdq{}}\PY{l+s}{t2}\PY{l+s}{\PYZdq{}}\PY{p}{)}\PY{p}{;}
\PY{+w}{        }\PY{n}{t1}\PY{p}{.}\PY{n+na}{start}\PY{p}{(}\PY{p}{)}\PY{p}{;}\PY{+w}{ }\PY{n}{t2}\PY{p}{.}\PY{n+na}{start}\PY{p}{(}\PY{p}{)}\PY{p}{;}
\PY{+w}{        }\PY{n}{t1}\PY{p}{.}\PY{n+na}{join}\PY{p}{(}\PY{p}{)}\PY{p}{;}\PY{+w}{ }\PY{n}{t2}\PY{p}{.}\PY{n+na}{join}\PY{p}{(}\PY{p}{)}\PY{p}{;}
\PY{+w}{        }\PY{c+c1}{// t1 and t2 print different values — each has its own independent slot}
\PY{+w}{    }\PY{p}{\PYZcb{}}
\PY{p}{\PYZcb{}}
\end{Verbatim}
\end{codebox}

\textbf{Common use cases:} per-request context (user id, transaction id, locale) in web servers, \code{SimpleDateFormat}/\code{DecimalFormat} instances (historically not thread-safe, so one-per-thread avoided sharing), database connections or transaction handles tied to “whatever thread is currently servicing this request.”

\textbf{\code{remove()} in thread pools, or you leak.} A \code{ThreadLocal}’s value is stored in a map owned by the \code{Thread} object itself. In a thread pool, threads are \emph{reused} across many tasks/requests — if you never call \code{remove()}, the value from request \#1 is still sitting there when the same pooled thread later handles request \#47, which is both a memory leak (values pile up and are never GC’d while the thread lives) and a correctness bug (stale data leaking across unrelated requests).

\begin{codebox}
\begin{Verbatim}[commandchars=\\\{\}]
\PY{c+c1}{// BROKEN in a pooled environment: value from a previous task can leak into the next one}
\PY{k+kd}{private}\PY{+w}{ }\PY{k+kd}{static}\PY{+w}{ }\PY{k+kd}{final}\PY{+w}{ }\PY{n}{ThreadLocal}\PY{o}{\PYZlt{}}\PY{n}{User}\PY{o}{\PYZgt{}}\PY{+w}{ }\PY{n}{currentUser}\PY{+w}{ }\PY{o}{=}\PY{+w}{ }\PY{k}{new}\PY{+w}{ }\PY{n}{ThreadLocal}\PY{o}{\PYZlt{}}\PY{o}{\PYZgt{}}\PY{p}{(}\PY{p}{)}\PY{p}{;}

\PY{k+kt}{void}\PY{+w}{ }\PY{n+nf}{handleRequest}\PY{p}{(}\PY{n}{User}\PY{+w}{ }\PY{n}{user}\PY{p}{)}\PY{+w}{ }\PY{p}{\PYZob{}}
\PY{+w}{    }\PY{n}{currentUser}\PY{p}{.}\PY{n+na}{set}\PY{p}{(}\PY{n}{user}\PY{p}{)}\PY{p}{;}
\PY{+w}{    }\PY{n}{process}\PY{p}{(}\PY{p}{)}\PY{p}{;}
\PY{+w}{    }\PY{c+c1}{// no cleanup — the next task run on this SAME pooled thread inherits \PYZsq{}user\PYZsq{} by accident}
\PY{p}{\PYZcb{}}
\end{Verbatim}
\end{codebox}

\begin{codebox}
\begin{Verbatim}[commandchars=\\\{\}]
\PY{c+c1}{// FIXED: always clean up, in a finally block, regardless of success or failure}
\PY{k+kt}{void}\PY{+w}{ }\PY{n+nf}{handleRequest}\PY{p}{(}\PY{n}{User}\PY{+w}{ }\PY{n}{user}\PY{p}{)}\PY{+w}{ }\PY{p}{\PYZob{}}
\PY{+w}{    }\PY{n}{currentUser}\PY{p}{.}\PY{n+na}{set}\PY{p}{(}\PY{n}{user}\PY{p}{)}\PY{p}{;}
\PY{+w}{    }\PY{k}{try}\PY{+w}{ }\PY{p}{\PYZob{}}
\PY{+w}{        }\PY{n}{process}\PY{p}{(}\PY{p}{)}\PY{p}{;}
\PY{+w}{    }\PY{p}{\PYZcb{}}\PY{+w}{ }\PY{k}{finally}\PY{+w}{ }\PY{p}{\PYZob{}}
\PY{+w}{        }\PY{n}{currentUser}\PY{p}{.}\PY{n+na}{remove}\PY{p}{(}\PY{p}{)}\PY{p}{;}\PY{+w}{ }\PY{c+c1}{// returns the thread\PYZsq{}s slot to its initial, empty state}
\PY{+w}{    }\PY{p}{\PYZcb{}}
\PY{p}{\PYZcb{}}
\end{Verbatim}
\end{codebox}

\textbf{\code{InheritableThreadLocal}.} A variant that copies the parent thread’s value into any \emph{child} thread it creates at the moment the child is constructed (via an overridable \code{childValue()} hook). Useful for propagating context (like a trace id) into worker threads spawned by a request handler — but it only captures a snapshot at thread-creation time, and it doesn’t work at all with pooled worker threads that were created long before the current task existed.

\begin{codebox}
\begin{Verbatim}[commandchars=\\\{\}]
\PY{k+kd}{public}\PY{+w}{ }\PY{k+kd}{class} \PY{n+nc}{InheritableExample}\PY{+w}{ }\PY{p}{\PYZob{}}
\PY{+w}{    }\PY{k+kd}{static}\PY{+w}{ }\PY{k+kd}{final}\PY{+w}{ }\PY{n}{InheritableThreadLocal}\PY{o}{\PYZlt{}}\PY{n}{String}\PY{o}{\PYZgt{}}\PY{+w}{ }\PY{n}{traceId}\PY{+w}{ }\PY{o}{=}\PY{+w}{ }\PY{k}{new}\PY{+w}{ }\PY{n}{InheritableThreadLocal}\PY{o}{\PYZlt{}}\PY{o}{\PYZgt{}}\PY{p}{(}\PY{p}{)}\PY{p}{;}

\PY{+w}{    }\PY{k+kd}{public}\PY{+w}{ }\PY{k+kd}{static}\PY{+w}{ }\PY{k+kt}{void}\PY{+w}{ }\PY{n+nf}{main}\PY{p}{(}\PY{n}{String}\PY{o}{[}\PY{o}{]}\PY{+w}{ }\PY{n}{args}\PY{p}{)}\PY{+w}{ }\PY{k+kd}{throws}\PY{+w}{ }\PY{n}{InterruptedException}\PY{+w}{ }\PY{p}{\PYZob{}}
\PY{+w}{        }\PY{n}{traceId}\PY{p}{.}\PY{n+na}{set}\PY{p}{(}\PY{l+s}{\PYZdq{}}\PY{l+s}{trace\PYZhy{}123}\PY{l+s}{\PYZdq{}}\PY{p}{)}\PY{p}{;}
\PY{+w}{        }\PY{n}{Thread}\PY{+w}{ }\PY{n}{child}\PY{+w}{ }\PY{o}{=}\PY{+w}{ }\PY{k}{new}\PY{+w}{ }\PY{n}{Thread}\PY{p}{(}\PY{p}{(}\PY{p}{)}\PY{+w}{ }\PY{o}{\PYZhy{}}\PY{o}{\PYZgt{}}\PY{+w}{ }\PY{p}{\PYZob{}}
\PY{+w}{            }\PY{n}{System}\PY{p}{.}\PY{n+na}{out}\PY{p}{.}\PY{n+na}{println}\PY{p}{(}\PY{n}{traceId}\PY{p}{.}\PY{n+na}{get}\PY{p}{(}\PY{p}{)}\PY{p}{)}\PY{p}{;}\PY{+w}{ }\PY{c+c1}{// prints \PYZdq{}trace\PYZhy{}123\PYZdq{} — inherited at creation time}
\PY{+w}{        }\PY{p}{\PYZcb{}}\PY{p}{)}\PY{p}{;}
\PY{+w}{        }\PY{n}{child}\PY{p}{.}\PY{n+na}{start}\PY{p}{(}\PY{p}{)}\PY{p}{;}
\PY{+w}{        }\PY{n}{child}\PY{p}{.}\PY{n+na}{join}\PY{p}{(}\PY{p}{)}\PY{p}{;}
\PY{+w}{    }\PY{p}{\PYZcb{}}
\PY{p}{\PYZcb{}}
\end{Verbatim}
\end{codebox}

\textbf{Looking ahead.} Java 21 introduced (as a preview, stabilizing later) \textbf{Scoped Values} as a safer, immutable alternative to \code{ThreadLocal} for propagating context — especially well-suited to virtual threads and structured concurrency, where \code{ThreadLocal}’s mutability and lifecycle can be awkward. Scoped Values are covered in Chapter 15.

\section[Happens-Before Relationship]{Happens-Before Relationship}
\label{h:13:happens-before-relationship}
The \textbf{Java Memory Model (JMM)} defines exactly when one thread’s actions are \emph{guaranteed} to be visible, in order, to another thread. That guarantee is called \textbf{happens-before}: if action A happens-before action B, every effect of A (every write to memory) is visible to whatever executes B. Without an explicit happens-before edge between two threads’ operations, the JMM makes \textbf{no promise at all} about ordering or visibility — the reordering and staleness bugs seen earlier in this chapter are exactly what happens in that gap.

The core rules:

\textbf{1. Program order rule.} Within a single thread, each action happens-before every subsequent action in that thread’s own code, in the order the code specifies (even though the JIT/CPU may reorder instructions \emph{invisibly} to that thread — the thread itself always observes its own actions in order).

\begin{codebox}
\begin{Verbatim}[commandchars=\\\{\}]
\PY{k+kt}{int}\PY{+w}{ }\PY{n}{a}\PY{+w}{ }\PY{o}{=}\PY{+w}{ }\PY{l+m+mi}{1}\PY{p}{;}
\PY{k+kt}{int}\PY{+w}{ }\PY{n}{b}\PY{+w}{ }\PY{o}{=}\PY{+w}{ }\PY{l+m+mi}{2}\PY{p}{;}
\PY{k+kt}{int}\PY{+w}{ }\PY{n}{c}\PY{+w}{ }\PY{o}{=}\PY{+w}{ }\PY{n}{a}\PY{+w}{ }\PY{o}{+}\PY{+w}{ }\PY{n}{b}\PY{p}{;}\PY{+w}{ }\PY{c+c1}{// program order guarantees \PYZsq{}a\PYZsq{} and \PYZsq{}b\PYZsq{} are both assigned before this line runs}
\end{Verbatim}
\end{codebox}

\textbf{2. Monitor lock rule.} Unlocking a monitor happens-before every subsequent lock of that \emph{same} monitor by any thread. This is what makes \code{synchronized} work: everything a thread wrote before releasing a lock is visible to the next thread that acquires that same lock.

\begin{codebox}
\begin{Verbatim}[commandchars=\\\{\}]
\PY{k+kd}{private}\PY{+w}{ }\PY{k+kd}{final}\PY{+w}{ }\PY{n}{Object}\PY{+w}{ }\PY{n}{lock}\PY{+w}{ }\PY{o}{=}\PY{+w}{ }\PY{k}{new}\PY{+w}{ }\PY{n}{Object}\PY{p}{(}\PY{p}{)}\PY{p}{;}
\PY{k+kd}{private}\PY{+w}{ }\PY{k+kt}{int}\PY{+w}{ }\PY{n}{sharedValue}\PY{+w}{ }\PY{o}{=}\PY{+w}{ }\PY{l+m+mi}{0}\PY{p}{;}

\PY{k+kt}{void}\PY{+w}{ }\PY{n+nf}{writer}\PY{p}{(}\PY{p}{)}\PY{+w}{ }\PY{p}{\PYZob{}}
\PY{+w}{    }\PY{k+kd}{synchronized}\PY{+w}{ }\PY{p}{(}\PY{n}{lock}\PY{p}{)}\PY{+w}{ }\PY{p}{\PYZob{}}
\PY{+w}{        }\PY{n}{sharedValue}\PY{+w}{ }\PY{o}{=}\PY{+w}{ }\PY{l+m+mi}{42}\PY{p}{;}\PY{+w}{ }\PY{c+c1}{// write happens before this unlock}
\PY{+w}{    }\PY{p}{\PYZcb{}}\PY{+w}{ }\PY{c+c1}{// unlock here}
\PY{p}{\PYZcb{}}

\PY{k+kt}{void}\PY{+w}{ }\PY{n+nf}{reader}\PY{p}{(}\PY{p}{)}\PY{+w}{ }\PY{p}{\PYZob{}}
\PY{+w}{    }\PY{k+kd}{synchronized}\PY{+w}{ }\PY{p}{(}\PY{n}{lock}\PY{p}{)}\PY{+w}{ }\PY{p}{\PYZob{}}\PY{+w}{ }\PY{c+c1}{// this lock happens\PYZhy{}after the writer\PYZsq{}s unlock}
\PY{+w}{        }\PY{n}{System}\PY{p}{.}\PY{n+na}{out}\PY{p}{.}\PY{n+na}{println}\PY{p}{(}\PY{n}{sharedValue}\PY{p}{)}\PY{p}{;}\PY{+w}{ }\PY{c+c1}{// guaranteed to see 42, not 0}
\PY{+w}{    }\PY{p}{\PYZcb{}}
\PY{p}{\PYZcb{}}
\end{Verbatim}
\end{codebox}

\textbf{3. Volatile variable rule.} A write to a \code{volatile} field happens-before every subsequent read of that \emph{same} field by any thread. This is exactly why the stop-flag example above is correct.

\begin{codebox}
\begin{Verbatim}[commandchars=\\\{\}]
\PY{k+kd}{private}\PY{+w}{ }\PY{k+kd}{volatile}\PY{+w}{ }\PY{k+kt}{boolean}\PY{+w}{ }\PY{n}{ready}\PY{+w}{ }\PY{o}{=}\PY{+w}{ }\PY{k+kc}{false}\PY{p}{;}
\PY{k+kd}{private}\PY{+w}{ }\PY{k+kt}{int}\PY{+w}{ }\PY{n}{data}\PY{+w}{ }\PY{o}{=}\PY{+w}{ }\PY{l+m+mi}{0}\PY{p}{;}

\PY{k+kt}{void}\PY{+w}{ }\PY{n+nf}{writer}\PY{p}{(}\PY{p}{)}\PY{+w}{ }\PY{p}{\PYZob{}}
\PY{+w}{    }\PY{n}{data}\PY{+w}{ }\PY{o}{=}\PY{+w}{ }\PY{l+m+mi}{100}\PY{p}{;}\PY{+w}{      }\PY{c+c1}{// (a) plain write}
\PY{+w}{    }\PY{n}{ready}\PY{+w}{ }\PY{o}{=}\PY{+w}{ }\PY{k+kc}{true}\PY{p}{;}\PY{+w}{     }\PY{c+c1}{// (b) volatile write — happens\PYZhy{}before any later read of \PYZsq{}ready\PYZsq{}}
\PY{p}{\PYZcb{}}

\PY{k+kt}{void}\PY{+w}{ }\PY{n+nf}{reader}\PY{p}{(}\PY{p}{)}\PY{+w}{ }\PY{p}{\PYZob{}}
\PY{+w}{    }\PY{k}{if}\PY{+w}{ }\PY{p}{(}\PY{n}{ready}\PY{p}{)}\PY{+w}{ }\PY{p}{\PYZob{}}\PY{+w}{                 }\PY{c+c1}{// (c) volatile read}
\PY{+w}{        }\PY{n}{System}\PY{p}{.}\PY{n+na}{out}\PY{p}{.}\PY{n+na}{println}\PY{p}{(}\PY{n}{data}\PY{p}{)}\PY{p}{;}\PY{+w}{ }\PY{c+c1}{// guaranteed to see 100, NOT 0 —}
\PY{+w}{    }\PY{p}{\PYZcb{}}\PY{+w}{                              }\PY{c+c1}{// because (a) happens\PYZhy{}before (b) [program order],}
\PY{p}{\PYZcb{}}\PY{+w}{                                  }\PY{c+c1}{// and (b) happens\PYZhy{}before (c) [volatile rule],}
\PY{+w}{                                   }\PY{c+c1}{// so by transitivity (a) happens\PYZhy{}before (c)}
\end{Verbatim}
\end{codebox}

\textbf{4. Thread start/join rules.} A call to \code{Thread.\allowbreak{}start()} happens-before any action inside the started thread’s \code{run()}. And every action inside a thread happens-before another thread’s successful \code{join()} on it.

\begin{codebox}
\begin{Verbatim}[commandchars=\\\{\}]
\PY{k+kt}{int}\PY{o}{[}\PY{o}{]}\PY{+w}{ }\PY{n}{result}\PY{+w}{ }\PY{o}{=}\PY{+w}{ }\PY{k}{new}\PY{+w}{ }\PY{k+kt}{int}\PY{o}{[}\PY{l+m+mi}{1}\PY{o}{]}\PY{p}{;}
\PY{n}{Thread}\PY{+w}{ }\PY{n}{worker}\PY{+w}{ }\PY{o}{=}\PY{+w}{ }\PY{k}{new}\PY{+w}{ }\PY{n}{Thread}\PY{p}{(}\PY{p}{(}\PY{p}{)}\PY{+w}{ }\PY{o}{\PYZhy{}}\PY{o}{\PYZgt{}}\PY{+w}{ }\PY{n}{result}\PY{o}{[}\PY{l+m+mi}{0}\PY{o}{]}\PY{+w}{ }\PY{o}{=}\PY{+w}{ }\PY{l+m+mi}{99}\PY{p}{)}\PY{p}{;}\PY{+w}{ }\PY{c+c1}{// write happens inside worker\PYZsq{}s run()}
\PY{n}{worker}\PY{p}{.}\PY{n+na}{start}\PY{p}{(}\PY{p}{)}\PY{p}{;}\PY{+w}{ }\PY{c+c1}{// happens\PYZhy{}before every action inside worker.run()}
\PY{n}{worker}\PY{p}{.}\PY{n+na}{join}\PY{p}{(}\PY{p}{)}\PY{p}{;}\PY{+w}{  }\PY{c+c1}{// worker\PYZsq{}s actions happen\PYZhy{}before this returning}
\PY{n}{System}\PY{p}{.}\PY{n+na}{out}\PY{p}{.}\PY{n+na}{println}\PY{p}{(}\PY{n}{result}\PY{o}{[}\PY{l+m+mi}{0}\PY{o}{]}\PY{p}{)}\PY{p}{;}\PY{+w}{ }\PY{c+c1}{// guaranteed to see 99, safely, no lock or volatile needed}
\end{Verbatim}
\end{codebox}

\textbf{5. Final field rule.} If an object is properly constructed (its constructor doesn’t leak \code{this} before finishing), the values of its \code{final} fields are guaranteed visible to any thread that gets a reference to that object after construction — no synchronization needed for those fields specifically.

\begin{codebox}
\begin{Verbatim}[commandchars=\\\{\}]
\PY{k+kd}{public}\PY{+w}{ }\PY{k+kd}{final}\PY{+w}{ }\PY{k+kd}{class} \PY{n+nc}{ImmutablePoint}\PY{+w}{ }\PY{p}{\PYZob{}}
\PY{+w}{    }\PY{k+kd}{private}\PY{+w}{ }\PY{k+kd}{final}\PY{+w}{ }\PY{k+kt}{int}\PY{+w}{ }\PY{n}{x}\PY{p}{,}\PY{+w}{ }\PY{n}{y}\PY{p}{;}\PY{+w}{ }\PY{c+c1}{// final fields}
\PY{+w}{    }\PY{k+kd}{public}\PY{+w}{ }\PY{n+nf}{ImmutablePoint}\PY{p}{(}\PY{k+kt}{int}\PY{+w}{ }\PY{n}{x}\PY{p}{,}\PY{+w}{ }\PY{k+kt}{int}\PY{+w}{ }\PY{n}{y}\PY{p}{)}\PY{+w}{ }\PY{p}{\PYZob{}}\PY{+w}{ }\PY{k}{this}\PY{p}{.}\PY{n+na}{x}\PY{+w}{ }\PY{o}{=}\PY{+w}{ }\PY{n}{x}\PY{p}{;}\PY{+w}{ }\PY{k}{this}\PY{p}{.}\PY{n+na}{y}\PY{+w}{ }\PY{o}{=}\PY{+w}{ }\PY{n}{y}\PY{p}{;}\PY{+w}{ }\PY{p}{\PYZcb{}}
\PY{+w}{    }\PY{k+kd}{public}\PY{+w}{ }\PY{k+kt}{int}\PY{+w}{ }\PY{n+nf}{getX}\PY{p}{(}\PY{p}{)}\PY{+w}{ }\PY{p}{\PYZob{}}\PY{+w}{ }\PY{k}{return}\PY{+w}{ }\PY{n}{x}\PY{p}{;}\PY{+w}{ }\PY{p}{\PYZcb{}}
\PY{+w}{    }\PY{k+kd}{public}\PY{+w}{ }\PY{k+kt}{int}\PY{+w}{ }\PY{n+nf}{getY}\PY{p}{(}\PY{p}{)}\PY{+w}{ }\PY{p}{\PYZob{}}\PY{+w}{ }\PY{k}{return}\PY{+w}{ }\PY{n}{y}\PY{p}{;}\PY{+w}{ }\PY{p}{\PYZcb{}}
\PY{p}{\PYZcb{}}
\PY{c+c1}{// Any thread that receives a reference to a fully\PYZhy{}constructed ImmutablePoint}
\PY{c+c1}{// is guaranteed to see the correct x and y, even with zero synchronization —}
\PY{c+c1}{// as long as the reference wasn\PYZsq{}t published before the constructor finished.}
\end{Verbatim}
\end{codebox}

\textbf{6. Transitivity.} If A happens-before B, and B happens-before C, then A happens-before C — even if A and C are actions on different threads that never directly synchronize with each other. This is what let the volatile example above chain a plain write through a volatile write/read pair into a guaranteed-visible result.

\section[Visibility, Atomicity, and Ordering]{Visibility, Atomicity, and Ordering}
\label{h:13:visibility-atomicity-and-ordering}
Nearly every concurrency bug is one (or more) of exactly three failures: \textbf{visibility} (a thread doesn’t see another thread’s latest write), \textbf{atomicity} (an operation that looks like one step is actually several, and gets interrupted mid-way), and \textbf{ordering} (instructions execute or become visible in a different sequence than the source code implies). Below: three broken examples, each fixed three different ways, with trade-offs.

\subsection[Broken counter (atomicity failure)]{Broken counter (atomicity failure)}
\label{h:13:broken-counter-atomicity-failure}
\begin{codebox}
\begin{Verbatim}[commandchars=\\\{\}]
\PY{c+c1}{// BROKEN: count++ is read\PYZhy{}modify\PYZhy{}write, not atomic — updates get lost under contention}
\PY{k+kd}{public}\PY{+w}{ }\PY{k+kd}{class} \PY{n+nc}{BrokenCounter}\PY{+w}{ }\PY{p}{\PYZob{}}
\PY{+w}{    }\PY{k+kd}{private}\PY{+w}{ }\PY{k+kt}{int}\PY{+w}{ }\PY{n}{count}\PY{+w}{ }\PY{o}{=}\PY{+w}{ }\PY{l+m+mi}{0}\PY{p}{;}
\PY{+w}{    }\PY{k+kd}{public}\PY{+w}{ }\PY{k+kt}{void}\PY{+w}{ }\PY{n+nf}{increment}\PY{p}{(}\PY{p}{)}\PY{+w}{ }\PY{p}{\PYZob{}}\PY{+w}{ }\PY{n}{count}\PY{o}{+}\PY{o}{+}\PY{p}{;}\PY{+w}{ }\PY{p}{\PYZcb{}}
\PY{+w}{    }\PY{k+kd}{public}\PY{+w}{ }\PY{k+kt}{int}\PY{+w}{ }\PY{n+nf}{get}\PY{p}{(}\PY{p}{)}\PY{+w}{ }\PY{p}{\PYZob{}}\PY{+w}{ }\PY{k}{return}\PY{+w}{ }\PY{n}{count}\PY{p}{;}\PY{+w}{ }\PY{p}{\PYZcb{}}
\PY{p}{\PYZcb{}}
\PY{c+c1}{// 10 threads x 100,000 increments each: final count is reliably LESS than 1,000,000}
\end{Verbatim}
\end{codebox}

\textbf{Fix 1 — \code{synchronized}:}

\begin{codebox}
\begin{Verbatim}[commandchars=\\\{\}]
\PY{k+kd}{public}\PY{+w}{ }\PY{k+kd}{class} \PY{n+nc}{SynchronizedCounter}\PY{+w}{ }\PY{p}{\PYZob{}}
\PY{+w}{    }\PY{k+kd}{private}\PY{+w}{ }\PY{k+kt}{int}\PY{+w}{ }\PY{n}{count}\PY{+w}{ }\PY{o}{=}\PY{+w}{ }\PY{l+m+mi}{0}\PY{p}{;}
\PY{+w}{    }\PY{k+kd}{public}\PY{+w}{ }\PY{k+kd}{synchronized}\PY{+w}{ }\PY{k+kt}{void}\PY{+w}{ }\PY{n+nf}{increment}\PY{p}{(}\PY{p}{)}\PY{+w}{ }\PY{p}{\PYZob{}}\PY{+w}{ }\PY{n}{count}\PY{o}{+}\PY{o}{+}\PY{p}{;}\PY{+w}{ }\PY{p}{\PYZcb{}}
\PY{+w}{    }\PY{k+kd}{public}\PY{+w}{ }\PY{k+kd}{synchronized}\PY{+w}{ }\PY{k+kt}{int}\PY{+w}{ }\PY{n+nf}{get}\PY{p}{(}\PY{p}{)}\PY{+w}{ }\PY{p}{\PYZob{}}\PY{+w}{ }\PY{k}{return}\PY{+w}{ }\PY{n}{count}\PY{p}{;}\PY{+w}{ }\PY{p}{\PYZcb{}}
\PY{p}{\PYZcb{}}
\end{Verbatim}
\end{codebox}

\textbf{Fix 2 — \code{volatile} alone does NOT work here} (shown earlier) — you’d need \code{volatile} \emph{plus} some other mechanism to make the read-modify-write atomic, so \code{volatile} by itself is not a valid fix for a counter.

\textbf{Fix 3 — atomic class:}

\begin{codebox}
\begin{Verbatim}[commandchars=\\\{\}]
\PY{k+kn}{import}\PY{+w}{ }\PY{n+nn}{java.util.concurrent.atomic.AtomicInteger}\PY{p}{;}

\PY{k+kd}{public}\PY{+w}{ }\PY{k+kd}{class} \PY{n+nc}{AtomicCounter}\PY{+w}{ }\PY{p}{\PYZob{}}
\PY{+w}{    }\PY{k+kd}{private}\PY{+w}{ }\PY{k+kd}{final}\PY{+w}{ }\PY{n}{AtomicInteger}\PY{+w}{ }\PY{n}{count}\PY{+w}{ }\PY{o}{=}\PY{+w}{ }\PY{k}{new}\PY{+w}{ }\PY{n}{AtomicInteger}\PY{p}{(}\PY{l+m+mi}{0}\PY{p}{)}\PY{p}{;}
\PY{+w}{    }\PY{k+kd}{public}\PY{+w}{ }\PY{k+kt}{void}\PY{+w}{ }\PY{n+nf}{increment}\PY{p}{(}\PY{p}{)}\PY{+w}{ }\PY{p}{\PYZob{}}\PY{+w}{ }\PY{n}{count}\PY{p}{.}\PY{n+na}{incrementAndGet}\PY{p}{(}\PY{p}{)}\PY{p}{;}\PY{+w}{ }\PY{p}{\PYZcb{}}
\PY{+w}{    }\PY{k+kd}{public}\PY{+w}{ }\PY{k+kt}{int}\PY{+w}{ }\PY{n+nf}{get}\PY{p}{(}\PY{p}{)}\PY{+w}{ }\PY{p}{\PYZob{}}\PY{+w}{ }\PY{k}{return}\PY{+w}{ }\PY{n}{count}\PY{p}{.}\PY{n+na}{get}\PY{p}{(}\PY{p}{)}\PY{p}{;}\PY{+w}{ }\PY{p}{\PYZcb{}}
\PY{p}{\PYZcb{}}
\end{Verbatim}
\end{codebox}

\subsection[Broken flag (visibility failure)]{Broken flag (visibility failure)}
\label{h:13:broken-flag-visibility-failure}
\begin{codebox}
\begin{Verbatim}[commandchars=\\\{\}]
\PY{c+c1}{// BROKEN: no visibility guarantee — the reading thread may cache \PYZsq{}running\PYZsq{} forever}
\PY{k+kd}{public}\PY{+w}{ }\PY{k+kd}{class} \PY{n+nc}{BrokenFlag}\PY{+w}{ }\PY{p}{\PYZob{}}
\PY{+w}{    }\PY{k+kd}{private}\PY{+w}{ }\PY{k+kt}{boolean}\PY{+w}{ }\PY{n}{running}\PY{+w}{ }\PY{o}{=}\PY{+w}{ }\PY{k+kc}{true}\PY{p}{;}
\PY{+w}{    }\PY{k+kd}{public}\PY{+w}{ }\PY{k+kt}{void}\PY{+w}{ }\PY{n+nf}{stop}\PY{p}{(}\PY{p}{)}\PY{+w}{ }\PY{p}{\PYZob{}}\PY{+w}{ }\PY{n}{running}\PY{+w}{ }\PY{o}{=}\PY{+w}{ }\PY{k+kc}{false}\PY{p}{;}\PY{+w}{ }\PY{p}{\PYZcb{}}
\PY{+w}{    }\PY{k+kd}{public}\PY{+w}{ }\PY{k+kt}{void}\PY{+w}{ }\PY{n+nf}{run}\PY{p}{(}\PY{p}{)}\PY{+w}{ }\PY{p}{\PYZob{}}
\PY{+w}{        }\PY{k}{while}\PY{+w}{ }\PY{p}{(}\PY{n}{running}\PY{p}{)}\PY{+w}{ }\PY{p}{\PYZob{}}\PY{+w}{ }\PY{c+cm}{/* busy loop, may never see \PYZsq{}false\PYZsq{} */}\PY{+w}{ }\PY{p}{\PYZcb{}}
\PY{+w}{    }\PY{p}{\PYZcb{}}
\PY{p}{\PYZcb{}}
\end{Verbatim}
\end{codebox}

\textbf{Fix 1 — \code{synchronized}} (works, but heavier than needed for a single boolean):

\begin{codebox}
\begin{Verbatim}[commandchars=\\\{\}]
\PY{k+kd}{public}\PY{+w}{ }\PY{k+kd}{class} \PY{n+nc}{SynchronizedFlag}\PY{+w}{ }\PY{p}{\PYZob{}}
\PY{+w}{    }\PY{k+kd}{private}\PY{+w}{ }\PY{k+kt}{boolean}\PY{+w}{ }\PY{n}{running}\PY{+w}{ }\PY{o}{=}\PY{+w}{ }\PY{k+kc}{true}\PY{p}{;}
\PY{+w}{    }\PY{k+kd}{public}\PY{+w}{ }\PY{k+kd}{synchronized}\PY{+w}{ }\PY{k+kt}{void}\PY{+w}{ }\PY{n+nf}{stop}\PY{p}{(}\PY{p}{)}\PY{+w}{ }\PY{p}{\PYZob{}}\PY{+w}{ }\PY{n}{running}\PY{+w}{ }\PY{o}{=}\PY{+w}{ }\PY{k+kc}{false}\PY{p}{;}\PY{+w}{ }\PY{p}{\PYZcb{}}
\PY{+w}{    }\PY{k+kd}{public}\PY{+w}{ }\PY{k+kd}{synchronized}\PY{+w}{ }\PY{k+kt}{boolean}\PY{+w}{ }\PY{n+nf}{isRunning}\PY{p}{(}\PY{p}{)}\PY{+w}{ }\PY{p}{\PYZob{}}\PY{+w}{ }\PY{k}{return}\PY{+w}{ }\PY{n}{running}\PY{p}{;}\PY{+w}{ }\PY{p}{\PYZcb{}}
\PY{+w}{    }\PY{k+kd}{public}\PY{+w}{ }\PY{k+kt}{void}\PY{+w}{ }\PY{n+nf}{run}\PY{p}{(}\PY{p}{)}\PY{+w}{ }\PY{p}{\PYZob{}}
\PY{+w}{        }\PY{k}{while}\PY{+w}{ }\PY{p}{(}\PY{n}{isRunning}\PY{p}{(}\PY{p}{)}\PY{p}{)}\PY{+w}{ }\PY{p}{\PYZob{}}\PY{+w}{ }\PY{c+cm}{/* every check takes a lock */}\PY{+w}{ }\PY{p}{\PYZcb{}}
\PY{+w}{    }\PY{p}{\PYZcb{}}
\PY{p}{\PYZcb{}}
\end{Verbatim}
\end{codebox}

\textbf{Fix 2 — \code{volatile}} (the idiomatic fix — single variable, single reads/writes, no atomicity needed):

\begin{codebox}
\begin{Verbatim}[commandchars=\\\{\}]
\PY{k+kd}{public}\PY{+w}{ }\PY{k+kd}{class} \PY{n+nc}{VolatileFlag}\PY{+w}{ }\PY{p}{\PYZob{}}
\PY{+w}{    }\PY{k+kd}{private}\PY{+w}{ }\PY{k+kd}{volatile}\PY{+w}{ }\PY{k+kt}{boolean}\PY{+w}{ }\PY{n}{running}\PY{+w}{ }\PY{o}{=}\PY{+w}{ }\PY{k+kc}{true}\PY{p}{;}
\PY{+w}{    }\PY{k+kd}{public}\PY{+w}{ }\PY{k+kt}{void}\PY{+w}{ }\PY{n+nf}{stop}\PY{p}{(}\PY{p}{)}\PY{+w}{ }\PY{p}{\PYZob{}}\PY{+w}{ }\PY{n}{running}\PY{+w}{ }\PY{o}{=}\PY{+w}{ }\PY{k+kc}{false}\PY{p}{;}\PY{+w}{ }\PY{p}{\PYZcb{}}
\PY{+w}{    }\PY{k+kd}{public}\PY{+w}{ }\PY{k+kt}{void}\PY{+w}{ }\PY{n+nf}{run}\PY{p}{(}\PY{p}{)}\PY{+w}{ }\PY{p}{\PYZob{}}
\PY{+w}{        }\PY{k}{while}\PY{+w}{ }\PY{p}{(}\PY{n}{running}\PY{p}{)}\PY{+w}{ }\PY{p}{\PYZob{}}\PY{+w}{ }\PY{c+cm}{/* cheap, correct, no locking */}\PY{+w}{ }\PY{p}{\PYZcb{}}
\PY{+w}{    }\PY{p}{\PYZcb{}}
\PY{p}{\PYZcb{}}
\end{Verbatim}
\end{codebox}

\textbf{Fix 3 — atomic class} (works, but overkill — you’re not doing CAS logic on a plain flag):

\begin{codebox}
\begin{Verbatim}[commandchars=\\\{\}]
\PY{k+kn}{import}\PY{+w}{ }\PY{n+nn}{java.util.concurrent.atomic.AtomicBoolean}\PY{p}{;}

\PY{k+kd}{public}\PY{+w}{ }\PY{k+kd}{class} \PY{n+nc}{AtomicFlag}\PY{+w}{ }\PY{p}{\PYZob{}}
\PY{+w}{    }\PY{k+kd}{private}\PY{+w}{ }\PY{k+kd}{final}\PY{+w}{ }\PY{n}{AtomicBoolean}\PY{+w}{ }\PY{n}{running}\PY{+w}{ }\PY{o}{=}\PY{+w}{ }\PY{k}{new}\PY{+w}{ }\PY{n}{AtomicBoolean}\PY{p}{(}\PY{k+kc}{true}\PY{p}{)}\PY{p}{;}
\PY{+w}{    }\PY{k+kd}{public}\PY{+w}{ }\PY{k+kt}{void}\PY{+w}{ }\PY{n+nf}{stop}\PY{p}{(}\PY{p}{)}\PY{+w}{ }\PY{p}{\PYZob{}}\PY{+w}{ }\PY{n}{running}\PY{p}{.}\PY{n+na}{set}\PY{p}{(}\PY{k+kc}{false}\PY{p}{)}\PY{p}{;}\PY{+w}{ }\PY{p}{\PYZcb{}}
\PY{+w}{    }\PY{k+kd}{public}\PY{+w}{ }\PY{k+kt}{void}\PY{+w}{ }\PY{n+nf}{run}\PY{p}{(}\PY{p}{)}\PY{+w}{ }\PY{p}{\PYZob{}}
\PY{+w}{        }\PY{k}{while}\PY{+w}{ }\PY{p}{(}\PY{n}{running}\PY{p}{.}\PY{n+na}{get}\PY{p}{(}\PY{p}{)}\PY{p}{)}\PY{+w}{ }\PY{p}{\PYZob{}}\PY{+w}{ }\PY{c+cm}{/* correct, but heavier machinery than needed */}\PY{+w}{ }\PY{p}{\PYZcb{}}
\PY{+w}{    }\PY{p}{\PYZcb{}}
\PY{p}{\PYZcb{}}
\end{Verbatim}
\end{codebox}

\subsection[Broken publication (reordering failure)]{Broken publication (reordering failure)}
\label{h:13:broken-publication-reordering-failure}
\begin{codebox}
\begin{Verbatim}[commandchars=\\\{\}]
\PY{c+c1}{// BROKEN: reader thread might see \PYZsq{}configured == true\PYZsq{} before seeing the fully\PYZhy{}set config fields,}
\PY{c+c1}{// because there\PYZsq{}s no happens\PYZhy{}before edge forcing the writes to become visible in order.}
\PY{k+kd}{public}\PY{+w}{ }\PY{k+kd}{class} \PY{n+nc}{BrokenPublication}\PY{+w}{ }\PY{p}{\PYZob{}}
\PY{+w}{    }\PY{k+kd}{private}\PY{+w}{ }\PY{k+kt}{int}\PY{+w}{ }\PY{n}{timeout}\PY{p}{;}
\PY{+w}{    }\PY{k+kd}{private}\PY{+w}{ }\PY{n}{String}\PY{+w}{ }\PY{n}{host}\PY{p}{;}
\PY{+w}{    }\PY{k+kd}{private}\PY{+w}{ }\PY{k+kt}{boolean}\PY{+w}{ }\PY{n}{configured}\PY{+w}{ }\PY{o}{=}\PY{+w}{ }\PY{k+kc}{false}\PY{p}{;}

\PY{+w}{    }\PY{k+kd}{public}\PY{+w}{ }\PY{k+kt}{void}\PY{+w}{ }\PY{n+nf}{configure}\PY{p}{(}\PY{k+kt}{int}\PY{+w}{ }\PY{n}{timeout}\PY{p}{,}\PY{+w}{ }\PY{n}{String}\PY{+w}{ }\PY{n}{host}\PY{p}{)}\PY{+w}{ }\PY{p}{\PYZob{}}
\PY{+w}{        }\PY{k}{this}\PY{p}{.}\PY{n+na}{timeout}\PY{+w}{ }\PY{o}{=}\PY{+w}{ }\PY{n}{timeout}\PY{p}{;}\PY{+w}{  }\PY{c+c1}{// (1)}
\PY{+w}{        }\PY{k}{this}\PY{p}{.}\PY{n+na}{host}\PY{+w}{ }\PY{o}{=}\PY{+w}{ }\PY{n}{host}\PY{p}{;}\PY{+w}{        }\PY{c+c1}{// (2)}
\PY{+w}{        }\PY{k}{this}\PY{p}{.}\PY{n+na}{configured}\PY{+w}{ }\PY{o}{=}\PY{+w}{ }\PY{k+kc}{true}\PY{p}{;}\PY{+w}{  }\PY{c+c1}{// (3) — nothing stops (3) being seen before (1)/(2) by another thread}
\PY{+w}{    }\PY{p}{\PYZcb{}}

\PY{+w}{    }\PY{k+kd}{public}\PY{+w}{ }\PY{k+kt}{void}\PY{+w}{ }\PY{n+nf}{useConfig}\PY{p}{(}\PY{p}{)}\PY{+w}{ }\PY{p}{\PYZob{}}
\PY{+w}{        }\PY{k}{if}\PY{+w}{ }\PY{p}{(}\PY{n}{configured}\PY{p}{)}\PY{+w}{ }\PY{p}{\PYZob{}}
\PY{+w}{            }\PY{n}{System}\PY{p}{.}\PY{n+na}{out}\PY{p}{.}\PY{n+na}{println}\PY{p}{(}\PY{n}{host}\PY{+w}{ }\PY{o}{+}\PY{+w}{ }\PY{l+s}{\PYZdq{}}\PY{l+s}{:}\PY{l+s}{\PYZdq{}}\PY{+w}{ }\PY{o}{+}\PY{+w}{ }\PY{n}{timeout}\PY{p}{)}\PY{p}{;}\PY{+w}{ }\PY{c+c1}{// could print a stale/partial host or 0 timeout}
\PY{+w}{        }\PY{p}{\PYZcb{}}
\PY{+w}{    }\PY{p}{\PYZcb{}}
\PY{p}{\PYZcb{}}
\end{Verbatim}
\end{codebox}

\textbf{Fix 1 — \code{synchronized}} (both methods share a lock, establishing happens-before via the monitor rule):

\begin{codebox}
\begin{Verbatim}[commandchars=\\\{\}]
\PY{k+kd}{public}\PY{+w}{ }\PY{k+kd}{class} \PY{n+nc}{SynchronizedPublication}\PY{+w}{ }\PY{p}{\PYZob{}}
\PY{+w}{    }\PY{k+kd}{private}\PY{+w}{ }\PY{k+kt}{int}\PY{+w}{ }\PY{n}{timeout}\PY{p}{;}
\PY{+w}{    }\PY{k+kd}{private}\PY{+w}{ }\PY{n}{String}\PY{+w}{ }\PY{n}{host}\PY{p}{;}
\PY{+w}{    }\PY{k+kd}{private}\PY{+w}{ }\PY{k+kt}{boolean}\PY{+w}{ }\PY{n}{configured}\PY{+w}{ }\PY{o}{=}\PY{+w}{ }\PY{k+kc}{false}\PY{p}{;}
\PY{+w}{    }\PY{k+kd}{private}\PY{+w}{ }\PY{k+kd}{final}\PY{+w}{ }\PY{n}{Object}\PY{+w}{ }\PY{n}{lock}\PY{+w}{ }\PY{o}{=}\PY{+w}{ }\PY{k}{new}\PY{+w}{ }\PY{n}{Object}\PY{p}{(}\PY{p}{)}\PY{p}{;}

\PY{+w}{    }\PY{k+kd}{public}\PY{+w}{ }\PY{k+kt}{void}\PY{+w}{ }\PY{n+nf}{configure}\PY{p}{(}\PY{k+kt}{int}\PY{+w}{ }\PY{n}{timeout}\PY{p}{,}\PY{+w}{ }\PY{n}{String}\PY{+w}{ }\PY{n}{host}\PY{p}{)}\PY{+w}{ }\PY{p}{\PYZob{}}
\PY{+w}{        }\PY{k+kd}{synchronized}\PY{+w}{ }\PY{p}{(}\PY{n}{lock}\PY{p}{)}\PY{+w}{ }\PY{p}{\PYZob{}}
\PY{+w}{            }\PY{k}{this}\PY{p}{.}\PY{n+na}{timeout}\PY{+w}{ }\PY{o}{=}\PY{+w}{ }\PY{n}{timeout}\PY{p}{;}
\PY{+w}{            }\PY{k}{this}\PY{p}{.}\PY{n+na}{host}\PY{+w}{ }\PY{o}{=}\PY{+w}{ }\PY{n}{host}\PY{p}{;}
\PY{+w}{            }\PY{k}{this}\PY{p}{.}\PY{n+na}{configured}\PY{+w}{ }\PY{o}{=}\PY{+w}{ }\PY{k+kc}{true}\PY{p}{;}
\PY{+w}{        }\PY{p}{\PYZcb{}}
\PY{+w}{    }\PY{p}{\PYZcb{}}

\PY{+w}{    }\PY{k+kd}{public}\PY{+w}{ }\PY{k+kt}{void}\PY{+w}{ }\PY{n+nf}{useConfig}\PY{p}{(}\PY{p}{)}\PY{+w}{ }\PY{p}{\PYZob{}}
\PY{+w}{        }\PY{k+kd}{synchronized}\PY{+w}{ }\PY{p}{(}\PY{n}{lock}\PY{p}{)}\PY{+w}{ }\PY{p}{\PYZob{}}
\PY{+w}{            }\PY{k}{if}\PY{+w}{ }\PY{p}{(}\PY{n}{configured}\PY{p}{)}\PY{+w}{ }\PY{n}{System}\PY{p}{.}\PY{n+na}{out}\PY{p}{.}\PY{n+na}{println}\PY{p}{(}\PY{n}{host}\PY{+w}{ }\PY{o}{+}\PY{+w}{ }\PY{l+s}{\PYZdq{}}\PY{l+s}{:}\PY{l+s}{\PYZdq{}}\PY{+w}{ }\PY{o}{+}\PY{+w}{ }\PY{n}{timeout}\PY{p}{)}\PY{p}{;}\PY{+w}{ }\PY{c+c1}{// always consistent}
\PY{+w}{        }\PY{p}{\PYZcb{}}
\PY{+w}{    }\PY{p}{\PYZcb{}}
\PY{p}{\PYZcb{}}
\end{Verbatim}
\end{codebox}

\textbf{Fix 2 — \code{volatile}} on the guard flag (relies on the volatile happens-before rule to publish everything written before it):

\begin{codebox}
\begin{Verbatim}[commandchars=\\\{\}]
\PY{k+kd}{public}\PY{+w}{ }\PY{k+kd}{class} \PY{n+nc}{VolatilePublication}\PY{+w}{ }\PY{p}{\PYZob{}}
\PY{+w}{    }\PY{k+kd}{private}\PY{+w}{ }\PY{k+kt}{int}\PY{+w}{ }\PY{n}{timeout}\PY{p}{;}
\PY{+w}{    }\PY{k+kd}{private}\PY{+w}{ }\PY{n}{String}\PY{+w}{ }\PY{n}{host}\PY{p}{;}
\PY{+w}{    }\PY{k+kd}{private}\PY{+w}{ }\PY{k+kd}{volatile}\PY{+w}{ }\PY{k+kt}{boolean}\PY{+w}{ }\PY{n}{configured}\PY{+w}{ }\PY{o}{=}\PY{+w}{ }\PY{k+kc}{false}\PY{p}{;}

\PY{+w}{    }\PY{k+kd}{public}\PY{+w}{ }\PY{k+kt}{void}\PY{+w}{ }\PY{n+nf}{configure}\PY{p}{(}\PY{k+kt}{int}\PY{+w}{ }\PY{n}{timeout}\PY{p}{,}\PY{+w}{ }\PY{n}{String}\PY{+w}{ }\PY{n}{host}\PY{p}{)}\PY{+w}{ }\PY{p}{\PYZob{}}
\PY{+w}{        }\PY{k}{this}\PY{p}{.}\PY{n+na}{timeout}\PY{+w}{ }\PY{o}{=}\PY{+w}{ }\PY{n}{timeout}\PY{p}{;}
\PY{+w}{        }\PY{k}{this}\PY{p}{.}\PY{n+na}{host}\PY{+w}{ }\PY{o}{=}\PY{+w}{ }\PY{n}{host}\PY{p}{;}
\PY{+w}{        }\PY{k}{this}\PY{p}{.}\PY{n+na}{configured}\PY{+w}{ }\PY{o}{=}\PY{+w}{ }\PY{k+kc}{true}\PY{p}{;}\PY{+w}{ }\PY{c+c1}{// volatile write: everything above is now visible to any reader}
\PY{+w}{    }\PY{p}{\PYZcb{}}

\PY{+w}{    }\PY{k+kd}{public}\PY{+w}{ }\PY{k+kt}{void}\PY{+w}{ }\PY{n+nf}{useConfig}\PY{p}{(}\PY{p}{)}\PY{+w}{ }\PY{p}{\PYZob{}}
\PY{+w}{        }\PY{k}{if}\PY{+w}{ }\PY{p}{(}\PY{n}{configured}\PY{p}{)}\PY{+w}{ }\PY{p}{\PYZob{}}\PY{+w}{ }\PY{c+c1}{// volatile read}
\PY{+w}{            }\PY{n}{System}\PY{p}{.}\PY{n+na}{out}\PY{p}{.}\PY{n+na}{println}\PY{p}{(}\PY{n}{host}\PY{+w}{ }\PY{o}{+}\PY{+w}{ }\PY{l+s}{\PYZdq{}}\PY{l+s}{:}\PY{l+s}{\PYZdq{}}\PY{+w}{ }\PY{o}{+}\PY{+w}{ }\PY{n}{timeout}\PY{p}{)}\PY{p}{;}\PY{+w}{ }\PY{c+c1}{// guaranteed correct, by the volatile rule}
\PY{+w}{        }\PY{p}{\PYZcb{}}
\PY{+w}{    }\PY{p}{\PYZcb{}}
\PY{p}{\PYZcb{}}
\end{Verbatim}
\end{codebox}

\textbf{Fix 3 — atomic reference to an immutable snapshot} (publish one object atomically instead of three separate fields):

\begin{codebox}
\begin{Verbatim}[commandchars=\\\{\}]
\PY{k+kn}{import}\PY{+w}{ }\PY{n+nn}{java.util.concurrent.atomic.AtomicReference}\PY{p}{;}

\PY{k+kd}{public}\PY{+w}{ }\PY{k+kd}{class} \PY{n+nc}{AtomicPublication}\PY{+w}{ }\PY{p}{\PYZob{}}
\PY{+w}{    }\PY{k+kd}{private}\PY{+w}{ }\PY{k+kd}{record} \PY{n+nc}{Config}\PY{p}{(}\PY{k+kt}{int}\PY{+w}{ }\PY{n}{timeout}\PY{p}{,}\PY{+w}{ }\PY{n}{String}\PY{+w}{ }\PY{n}{host}\PY{p}{)}\PY{+w}{ }\PY{p}{\PYZob{}}\PY{p}{\PYZcb{}}
\PY{+w}{    }\PY{k+kd}{private}\PY{+w}{ }\PY{k+kd}{final}\PY{+w}{ }\PY{n}{AtomicReference}\PY{o}{\PYZlt{}}\PY{n}{Config}\PY{o}{\PYZgt{}}\PY{+w}{ }\PY{n}{config}\PY{+w}{ }\PY{o}{=}\PY{+w}{ }\PY{k}{new}\PY{+w}{ }\PY{n}{AtomicReference}\PY{o}{\PYZlt{}}\PY{o}{\PYZgt{}}\PY{p}{(}\PY{p}{)}\PY{p}{;}

\PY{+w}{    }\PY{k+kd}{public}\PY{+w}{ }\PY{k+kt}{void}\PY{+w}{ }\PY{n+nf}{configure}\PY{p}{(}\PY{k+kt}{int}\PY{+w}{ }\PY{n}{timeout}\PY{p}{,}\PY{+w}{ }\PY{n}{String}\PY{+w}{ }\PY{n}{host}\PY{p}{)}\PY{+w}{ }\PY{p}{\PYZob{}}
\PY{+w}{        }\PY{n}{config}\PY{p}{.}\PY{n+na}{set}\PY{p}{(}\PY{k}{new}\PY{+w}{ }\PY{n}{Config}\PY{p}{(}\PY{n}{timeout}\PY{p}{,}\PY{+w}{ }\PY{n}{host}\PY{p}{)}\PY{p}{)}\PY{p}{;}\PY{+w}{ }\PY{c+c1}{// single atomic reference swap, fully\PYZhy{}built object}
\PY{+w}{    }\PY{p}{\PYZcb{}}

\PY{+w}{    }\PY{k+kd}{public}\PY{+w}{ }\PY{k+kt}{void}\PY{+w}{ }\PY{n+nf}{useConfig}\PY{p}{(}\PY{p}{)}\PY{+w}{ }\PY{p}{\PYZob{}}
\PY{+w}{        }\PY{n}{Config}\PY{+w}{ }\PY{n}{c}\PY{+w}{ }\PY{o}{=}\PY{+w}{ }\PY{n}{config}\PY{p}{.}\PY{n+na}{get}\PY{p}{(}\PY{p}{)}\PY{p}{;}
\PY{+w}{        }\PY{k}{if}\PY{+w}{ }\PY{p}{(}\PY{n}{c}\PY{+w}{ }\PY{o}{!}\PY{o}{=}\PY{+w}{ }\PY{k+kc}{null}\PY{p}{)}\PY{+w}{ }\PY{n}{System}\PY{p}{.}\PY{n+na}{out}\PY{p}{.}\PY{n+na}{println}\PY{p}{(}\PY{n}{c}\PY{p}{.}\PY{n+na}{host}\PY{p}{(}\PY{p}{)}\PY{+w}{ }\PY{o}{+}\PY{+w}{ }\PY{l+s}{\PYZdq{}}\PY{l+s}{:}\PY{l+s}{\PYZdq{}}\PY{+w}{ }\PY{o}{+}\PY{+w}{ }\PY{n}{c}\PY{p}{.}\PY{n+na}{timeout}\PY{p}{(}\PY{p}{)}\PY{p}{)}\PY{p}{;}\PY{+w}{ }\PY{c+c1}{// always consistent}
\PY{+w}{    }\PY{p}{\PYZcb{}}
\PY{p}{\PYZcb{}}
\end{Verbatim}
\end{codebox}

\subsection[Trade-offs at a glance]{Trade-offs at a glance}
\label{h:13:trade-offs-at-a-glance}
{\tablefontsmall
\begin{xltabular}{\linewidth}{@{}L{1.040}L{0.532}L{0.677}L{0.484}L{1.476}L{1.790}@{}}
\toprule
\rowcolor{tablehdr}
\thd{Approach} & \thd{Visibility} & \thd{Atomicity} & \thd{Ordering} & \thd{Performance} & \thd{Best for} \\
\midrule
\endhead
\code{synchronized} & Yes (via monitor rule) & Yes (whole block) & Yes & Lock overhead; blocks other threads & Multi-field invariants, compound operations \\
\code{volatile} & Yes & No (single read/write only) & Yes (for that field) & Very cheap, no blocking & Single flags, single published references \\
Atomic classes (\code{Atomic*}, \code{LongAdder}) & Yes & Yes (single variable) & Yes & CAS retry cost under heavy contention; \code{LongAdder} reduces this & Counters, single-variable compare-and-update logic \\
\code{ReentrantLock} / \code{ReadWriteLock} / \code{StampedLock} & Yes & Yes (whole critical section) & Yes & More features than \code{synchronized}, similar or better cost & Timed/interruptible acquisition, read-heavy workloads, multiple conditions \\
\bottomrule
\end{xltabular}
}

\section[Common Code-Review Interview Pitfalls]{Common Code-Review Interview Pitfalls}
\label{h:13:common-code-review-interview-pitfalls}
\begin{enumerate}
\item 
\textbf{Calling \code{run()} instead of \code{start()}.} Why it matters: \code{run()} executes on the calling thread — no new thread, no concurrency, and the bug is easy to miss because the code still “works,” just without any parallelism.

\begin{codebox}
\begin{Verbatim}[commandchars=\\\{\}]
\PY{c+c1}{// Before}
\PY{n}{Thread}\PY{+w}{ }\PY{n}{t}\PY{+w}{ }\PY{o}{=}\PY{+w}{ }\PY{k}{new}\PY{+w}{ }\PY{n}{Thread}\PY{p}{(}\PY{n}{task}\PY{p}{)}\PY{p}{;}
\PY{n}{t}\PY{p}{.}\PY{n+na}{run}\PY{p}{(}\PY{p}{)}\PY{p}{;}\PY{+w}{ }\PY{c+c1}{// runs on the CURRENT thread}
\PY{c+c1}{// After}
\PY{n}{t}\PY{p}{.}\PY{n+na}{start}\PY{p}{(}\PY{p}{)}\PY{p}{;}\PY{+w}{ }\PY{c+c1}{// actually schedules a new thread}
\end{Verbatim}
\end{codebox}

\item 
\textbf{Swallowing \code{InterruptedException} without restoring the interrupt flag.} Why it matters: catching and ignoring it erases the cancellation signal, so any outer loop or framework relying on interruption to stop the thread never finds out.

\begin{codebox}
\begin{Verbatim}[commandchars=\\\{\}]
\PY{c+c1}{// Before}
\PY{k}{try}\PY{+w}{ }\PY{p}{\PYZob{}}\PY{+w}{ }\PY{n}{Thread}\PY{p}{.}\PY{n+na}{sleep}\PY{p}{(}\PY{l+m+mi}{100}\PY{p}{)}\PY{p}{;}\PY{+w}{ }\PY{p}{\PYZcb{}}\PY{+w}{ }\PY{k}{catch}\PY{+w}{ }\PY{p}{(}\PY{n}{InterruptedException}\PY{+w}{ }\PY{n}{e}\PY{p}{)}\PY{+w}{ }\PY{p}{\PYZob{}}\PY{+w}{ }\PY{c+cm}{/* ignored */}\PY{+w}{ }\PY{p}{\PYZcb{}}
\PY{c+c1}{// After}
\PY{k}{try}\PY{+w}{ }\PY{p}{\PYZob{}}\PY{+w}{ }\PY{n}{Thread}\PY{p}{.}\PY{n+na}{sleep}\PY{p}{(}\PY{l+m+mi}{100}\PY{p}{)}\PY{p}{;}\PY{+w}{ }\PY{p}{\PYZcb{}}\PY{+w}{ }\PY{k}{catch}\PY{+w}{ }\PY{p}{(}\PY{n}{InterruptedException}\PY{+w}{ }\PY{n}{e}\PY{p}{)}\PY{+w}{ }\PY{p}{\PYZob{}}\PY{+w}{ }\PY{n}{Thread}\PY{p}{.}\PY{n+na}{currentThread}\PY{p}{(}\PY{p}{)}\PY{p}{.}\PY{n+na}{interrupt}\PY{p}{(}\PY{p}{)}\PY{p}{;}\PY{+w}{ }\PY{k}{return}\PY{p}{;}\PY{+w}{ }\PY{p}{\PYZcb{}}
\end{Verbatim}
\end{codebox}

\item 
\textbf{Using \code{wait()}/\code{notify()} (or a \code{Condition.\allowbreak{}await()}) with \code{if} instead of \code{while}.} Why it matters: spurious wakeups and multiple competing waiters mean the condition must be re-checked after waking — an \code{if} proceeds on a stale assumption.

\begin{codebox}
\begin{Verbatim}[commandchars=\\\{\}]
\PY{c+c1}{// Before}
\PY{k}{if}\PY{+w}{ }\PY{p}{(}\PY{n}{queue}\PY{p}{.}\PY{n+na}{isEmpty}\PY{p}{(}\PY{p}{)}\PY{p}{)}\PY{+w}{ }\PY{n}{wait}\PY{p}{(}\PY{p}{)}\PY{p}{;}
\PY{c+c1}{// After}
\PY{k}{while}\PY{+w}{ }\PY{p}{(}\PY{n}{queue}\PY{p}{.}\PY{n+na}{isEmpty}\PY{p}{(}\PY{p}{)}\PY{p}{)}\PY{+w}{ }\PY{n}{wait}\PY{p}{(}\PY{p}{)}\PY{p}{;}
\end{Verbatim}
\end{codebox}

\item 
\textbf{Checking a condition outside the lock, then locking separately (lost wakeup / TOCTOU).} Why it matters: another thread can change the condition and signal in the gap between the check and the lock, and that signal is lost forever.

\begin{codebox}
\begin{Verbatim}[commandchars=\\\{\}]
\PY{c+c1}{// Before}
\PY{k}{if}\PY{+w}{ }\PY{p}{(}\PY{n}{isEmpty}\PY{p}{(}\PY{p}{)}\PY{p}{)}\PY{+w}{ }\PY{p}{\PYZob{}}\PY{+w}{ }\PY{k+kd}{synchronized}\PY{+w}{ }\PY{p}{(}\PY{k}{this}\PY{p}{)}\PY{+w}{ }\PY{p}{\PYZob{}}\PY{+w}{ }\PY{n}{wait}\PY{p}{(}\PY{p}{)}\PY{p}{;}\PY{+w}{ }\PY{p}{\PYZcb{}}\PY{+w}{ }\PY{p}{\PYZcb{}}
\PY{c+c1}{// After}
\PY{k+kd}{synchronized}\PY{+w}{ }\PY{p}{(}\PY{k}{this}\PY{p}{)}\PY{+w}{ }\PY{p}{\PYZob{}}\PY{+w}{ }\PY{k}{while}\PY{+w}{ }\PY{p}{(}\PY{n}{isEmpty}\PY{p}{(}\PY{p}{)}\PY{p}{)}\PY{+w}{ }\PY{n}{wait}\PY{p}{(}\PY{p}{)}\PY{p}{;}\PY{+w}{ }\PY{p}{\PYZcb{}}
\end{Verbatim}
\end{codebox}

\item 
\textbf{Acquiring locks in inconsistent order across different code paths (deadlock risk).} Why it matters: two threads taking the same two locks in opposite order is the textbook deadlock; it often only shows up under real production load, not in tests.

\begin{codebox}
\begin{Verbatim}[commandchars=\\\{\}]
\PY{c+c1}{// Before: thread A locks(a,b); thread B locks(b,a);}
\PY{c+c1}{// After: every thread locks in the same global order, e.g. always a before b}
\end{Verbatim}
\end{codebox}

\item 
\textbf{Using \code{ReentrantLock} without \code{try}/\code{finally}.} Why it matters: unlike \code{synchronized}, nothing releases the lock automatically — an exception between \code{lock()} and \code{unlock()} leaks the lock permanently, freezing every other thread that needs it.

\begin{codebox}
\begin{Verbatim}[commandchars=\\\{\}]
\PY{c+c1}{// Before}
\PY{n}{lock}\PY{p}{.}\PY{n+na}{lock}\PY{p}{(}\PY{p}{)}\PY{p}{;}
\PY{n}{doWork}\PY{p}{(}\PY{p}{)}\PY{p}{;}\PY{+w}{ }\PY{c+c1}{// if this throws, the lock is NEVER released}
\PY{n}{lock}\PY{p}{.}\PY{n+na}{unlock}\PY{p}{(}\PY{p}{)}\PY{p}{;}
\PY{c+c1}{// After}
\PY{n}{lock}\PY{p}{.}\PY{n+na}{lock}\PY{p}{(}\PY{p}{)}\PY{p}{;}
\PY{k}{try}\PY{+w}{ }\PY{p}{\PYZob{}}\PY{+w}{ }\PY{n}{doWork}\PY{p}{(}\PY{p}{)}\PY{p}{;}\PY{+w}{ }\PY{p}{\PYZcb{}}\PY{+w}{ }\PY{k}{finally}\PY{+w}{ }\PY{p}{\PYZob{}}\PY{+w}{ }\PY{n}{lock}\PY{p}{.}\PY{n+na}{unlock}\PY{p}{(}\PY{p}{)}\PY{p}{;}\PY{+w}{ }\PY{p}{\PYZcb{}}
\end{Verbatim}
\end{codebox}

\item 
\textbf{Assuming \code{volatile} makes compound operations like \code{count++} atomic.} Why it matters: \code{volatile} only guarantees visibility and ordering for single reads/writes; a read-modify-write is still three separate steps that can interleave across threads.

\begin{codebox}
\begin{Verbatim}[commandchars=\\\{\}]
\PY{c+c1}{// Before}
\PY{k+kd}{private}\PY{+w}{ }\PY{k+kd}{volatile}\PY{+w}{ }\PY{k+kt}{int}\PY{+w}{ }\PY{n}{count}\PY{+w}{ }\PY{o}{=}\PY{+w}{ }\PY{l+m+mi}{0}\PY{p}{;}
\PY{k+kt}{void}\PY{+w}{ }\PY{n+nf}{increment}\PY{p}{(}\PY{p}{)}\PY{+w}{ }\PY{p}{\PYZob{}}\PY{+w}{ }\PY{n}{count}\PY{o}{+}\PY{o}{+}\PY{p}{;}\PY{+w}{ }\PY{p}{\PYZcb{}}\PY{+w}{ }\PY{c+c1}{// still lossy under contention}
\PY{c+c1}{// After}
\PY{k+kd}{private}\PY{+w}{ }\PY{k+kd}{final}\PY{+w}{ }\PY{n}{AtomicInteger}\PY{+w}{ }\PY{n}{count}\PY{+w}{ }\PY{o}{=}\PY{+w}{ }\PY{k}{new}\PY{+w}{ }\PY{n}{AtomicInteger}\PY{p}{(}\PY{l+m+mi}{0}\PY{p}{)}\PY{p}{;}
\PY{k+kt}{void}\PY{+w}{ }\PY{n+nf}{increment}\PY{p}{(}\PY{p}{)}\PY{+w}{ }\PY{p}{\PYZob{}}\PY{+w}{ }\PY{n}{count}\PY{p}{.}\PY{n+na}{incrementAndGet}\PY{p}{(}\PY{p}{)}\PY{p}{;}\PY{+w}{ }\PY{p}{\PYZcb{}}
\end{Verbatim}
\end{codebox}

\item 
\textbf{Double-checked-locking singleton missing \code{volatile} on the instance field.} Why it matters: without \code{volatile}, another thread can observe a non-null reference to an object whose constructor hasn’t finished running, i.e., a partially-built object.

\begin{codebox}
\begin{Verbatim}[commandchars=\\\{\}]
\PY{c+c1}{// Before}
\PY{k+kd}{private}\PY{+w}{ }\PY{k+kd}{static}\PY{+w}{ }\PY{n}{Singleton}\PY{+w}{ }\PY{n}{instance}\PY{p}{;}
\PY{c+c1}{// After}
\PY{k+kd}{private}\PY{+w}{ }\PY{k+kd}{static}\PY{+w}{ }\PY{k+kd}{volatile}\PY{+w}{ }\PY{n}{Singleton}\PY{+w}{ }\PY{n}{instance}\PY{p}{;}
\end{Verbatim}
\end{codebox}

\item 
\textbf{Forgetting \code{remove()} on a \code{ThreadLocal} used inside a pooled executor.} Why it matters: pooled threads are reused across tasks, so stale values silently leak from one task/request into the next — both a memory leak and a data-isolation bug.

\begin{codebox}
\begin{Verbatim}[commandchars=\\\{\}]
\PY{c+c1}{// Before}
\PY{n}{currentUser}\PY{p}{.}\PY{n+na}{set}\PY{p}{(}\PY{n}{user}\PY{p}{)}\PY{p}{;}
\PY{n}{process}\PY{p}{(}\PY{p}{)}\PY{p}{;}
\PY{c+c1}{// After}
\PY{n}{currentUser}\PY{p}{.}\PY{n+na}{set}\PY{p}{(}\PY{n}{user}\PY{p}{)}\PY{p}{;}
\PY{k}{try}\PY{+w}{ }\PY{p}{\PYZob{}}\PY{+w}{ }\PY{n}{process}\PY{p}{(}\PY{p}{)}\PY{p}{;}\PY{+w}{ }\PY{p}{\PYZcb{}}\PY{+w}{ }\PY{k}{finally}\PY{+w}{ }\PY{p}{\PYZob{}}\PY{+w}{ }\PY{n}{currentUser}\PY{p}{.}\PY{n+na}{remove}\PY{p}{(}\PY{p}{)}\PY{p}{;}\PY{+w}{ }\PY{p}{\PYZcb{}}
\end{Verbatim}
\end{codebox}

\item 
\textbf{Using \code{StampedLock} recursively, assuming it behaves like \code{ReentrantLock}.} Why it matters: \code{StampedLock} is not reentrant — a thread re-acquiring its own write lock deadlocks against itself, with no exception to warn you.

\begin{codebox}
\begin{Verbatim}[commandchars=\\\{\}]
\PY{c+c1}{// Before}
\PY{k+kt}{long}\PY{+w}{ }\PY{n}{s}\PY{+w}{ }\PY{o}{=}\PY{+w}{ }\PY{n}{lock}\PY{p}{.}\PY{n+na}{writeLock}\PY{p}{(}\PY{p}{)}\PY{p}{;}
\PY{n}{recursiveCallThatAlsoLocks}\PY{p}{(}\PY{p}{)}\PY{p}{;}\PY{+w}{ }\PY{c+c1}{// deadlocks if it calls lock.writeLock() again on this thread}
\PY{c+c1}{// After: restructure so the lock is acquired exactly once per call chain, or use ReentrantLock}
\end{Verbatim}
\end{codebox}

\item 
\textbf{Skipping \code{validate()} after a \code{Stamped\allowbreak{}Lock.\allowbreak{}try\allowbreak{}Optimistic\allowbreak{}Read()}.} Why it matters: without validation, the “optimistic” read might have raced with a concurrent writer and returned torn or inconsistent data with no error at all.

\begin{codebox}
\begin{Verbatim}[commandchars=\\\{\}]
\PY{c+c1}{// Before}
\PY{k+kt}{long}\PY{+w}{ }\PY{n}{stamp}\PY{+w}{ }\PY{o}{=}\PY{+w}{ }\PY{n}{lock}\PY{p}{.}\PY{n+na}{tryOptimisticRead}\PY{p}{(}\PY{p}{)}\PY{p}{;}
\PY{k+kt}{double}\PY{+w}{ }\PY{n}{result}\PY{+w}{ }\PY{o}{=}\PY{+w}{ }\PY{n}{x}\PY{+w}{ }\PY{o}{+}\PY{+w}{ }\PY{n}{y}\PY{p}{;}\PY{+w}{ }\PY{c+c1}{// used without checking if a writer interleaved}
\PY{c+c1}{// After}
\PY{k+kt}{long}\PY{+w}{ }\PY{n}{stamp}\PY{+w}{ }\PY{o}{=}\PY{+w}{ }\PY{n}{lock}\PY{p}{.}\PY{n+na}{tryOptimisticRead}\PY{p}{(}\PY{p}{)}\PY{p}{;}
\PY{k+kt}{double}\PY{+w}{ }\PY{n}{result}\PY{+w}{ }\PY{o}{=}\PY{+w}{ }\PY{n}{x}\PY{+w}{ }\PY{o}{+}\PY{+w}{ }\PY{n}{y}\PY{p}{;}
\PY{k}{if}\PY{+w}{ }\PY{p}{(}\PY{o}{!}\PY{n}{lock}\PY{p}{.}\PY{n+na}{validate}\PY{p}{(}\PY{n}{stamp}\PY{p}{)}\PY{p}{)}\PY{+w}{ }\PY{p}{\PYZob{}}\PY{+w}{ }\PY{c+cm}{/* retry with lock.readLock() */}\PY{+w}{ }\PY{p}{\PYZcb{}}
\end{Verbatim}
\end{codebox}

\item 
\textbf{Relying on thread priority for correctness or scheduling guarantees.} Why it matters: priority is only a scheduling hint with JVM/OS-dependent behavior — no ordering or fairness guarantee exists, so “fixing” a race by tweaking priorities is not a real fix.

\begin{codebox}
\begin{Verbatim}[commandchars=\\\{\}]
\PY{c+c1}{// Before: \PYZdq{}the writer thread has higher priority, so it always runs first\PYZdq{}}
\PY{n}{writerThread}\PY{p}{.}\PY{n+na}{setPriority}\PY{p}{(}\PY{n}{Thread}\PY{p}{.}\PY{n+na}{MAX\PYZus{}PRIORITY}\PY{p}{)}\PY{p}{;}\PY{+w}{ }\PY{c+c1}{// not a synchronization mechanism}
\PY{c+c1}{// After: use an actual ordering guarantee — a lock, a latch, or happens\PYZhy{}before via volatile/join}
\end{Verbatim}
\end{codebox}

\item 
\textbf{Using \code{Thread.\allowbreak{}stop()} or \code{Thread.\allowbreak{}suspend()} to cancel work.} Why it matters: both are deprecated for removal because they can kill a thread mid-update, leaving shared state permanently corrupted, or freeze a thread while it holds a lock other threads need.

\begin{codebox}
\begin{Verbatim}[commandchars=\\\{\}]
\PY{c+c1}{// Before}
\PY{n}{worker}\PY{p}{.}\PY{n+na}{stop}\PY{p}{(}\PY{p}{)}\PY{p}{;}\PY{+w}{ }\PY{c+c1}{// can leave shared state half\PYZhy{}updated}
\PY{c+c1}{// After}
\PY{k+kd}{volatile}\PY{+w}{ }\PY{k+kt}{boolean}\PY{+w}{ }\PY{n}{cancelled}\PY{p}{;}\PY{+w}{ }\PY{c+c1}{// or Thread.interrupt() + cooperative checks}
\PY{k+kt}{void}\PY{+w}{ }\PY{n+nf}{run}\PY{p}{(}\PY{p}{)}\PY{+w}{ }\PY{p}{\PYZob{}}\PY{+w}{ }\PY{k}{while}\PY{+w}{ }\PY{p}{(}\PY{o}{!}\PY{n}{cancelled}\PY{p}{)}\PY{+w}{ }\PY{p}{\PYZob{}}\PY{+w}{ }\PY{n}{doUnitOfWork}\PY{p}{(}\PY{p}{)}\PY{p}{;}\PY{+w}{ }\PY{p}{\PYZcb{}}\PY{+w}{ }\PY{p}{\PYZcb{}}
\end{Verbatim}
\end{codebox}

\item 
\textbf{Mixing \code{synchronized(\allowbreak{}this)} in one method with \code{synchronized(\allowbreak{}Some\allowbreak{}Class.\allowbreak{}class)} in another for the same shared state.} Why it matters: these are two different locks, so “synchronizing” both methods provides zero actual mutual exclusion between them.

\begin{codebox}
\begin{Verbatim}[commandchars=\\\{\}]
\PY{c+c1}{// Before}
\PY{k+kd}{public}\PY{+w}{ }\PY{k+kd}{synchronized}\PY{+w}{ }\PY{k+kt}{void}\PY{+w}{ }\PY{n+nf}{instanceMethod}\PY{p}{(}\PY{p}{)}\PY{+w}{ }\PY{p}{\PYZob{}}\PY{+w}{ }\PY{n}{sharedStatic}\PY{o}{+}\PY{o}{+}\PY{p}{;}\PY{+w}{ }\PY{p}{\PYZcb{}}\PY{+w}{       }\PY{c+c1}{// locks \PYZsq{}this\PYZsq{}}
\PY{k+kd}{public}\PY{+w}{ }\PY{k+kd}{static}\PY{+w}{ }\PY{k+kd}{synchronized}\PY{+w}{ }\PY{k+kt}{void}\PY{+w}{ }\PY{n+nf}{staticMethod}\PY{p}{(}\PY{p}{)}\PY{+w}{ }\PY{p}{\PYZob{}}\PY{+w}{ }\PY{n}{sharedStatic}\PY{o}{+}\PY{o}{+}\PY{p}{;}\PY{+w}{ }\PY{p}{\PYZcb{}}\PY{+w}{ }\PY{c+c1}{// locks the Class object}
\PY{c+c1}{// After: pick one lock and use it consistently for all access to sharedStatic}
\end{Verbatim}
\end{codebox}

\item 
\textbf{Holding a lock while doing slow I/O, network calls, or \code{Thread.\allowbreak{}sleep()}.} Why it matters: every other thread waiting on that lock is blocked for the full duration of the slow operation, turning a local slowdown into a system-wide bottleneck.

\begin{codebox}
\begin{Verbatim}[commandchars=\\\{\}]
\PY{c+c1}{// Before}
\PY{k+kd}{synchronized}\PY{+w}{ }\PY{p}{(}\PY{n}{lock}\PY{p}{)}\PY{+w}{ }\PY{p}{\PYZob{}}\PY{+w}{ }\PY{n}{callSlowRemoteService}\PY{p}{(}\PY{p}{)}\PY{p}{;}\PY{+w}{ }\PY{n}{updateSharedState}\PY{p}{(}\PY{p}{)}\PY{p}{;}\PY{+w}{ }\PY{p}{\PYZcb{}}
\PY{c+c1}{// After}
\PY{n}{Result}\PY{+w}{ }\PY{n}{r}\PY{+w}{ }\PY{o}{=}\PY{+w}{ }\PY{n}{callSlowRemoteService}\PY{p}{(}\PY{p}{)}\PY{p}{;}\PY{+w}{ }\PY{c+c1}{// do the slow part outside the lock}
\PY{k+kd}{synchronized}\PY{+w}{ }\PY{p}{(}\PY{n}{lock}\PY{p}{)}\PY{+w}{ }\PY{p}{\PYZob{}}\PY{+w}{ }\PY{n}{updateSharedState}\PY{p}{(}\PY{n}{r}\PY{p}{)}\PY{p}{;}\PY{+w}{ }\PY{p}{\PYZcb{}}\PY{+w}{ }\PY{c+c1}{// hold the lock only for the fast part}
\end{Verbatim}
\end{codebox}

\item 
\textbf{Publishing a partially-constructed object by leaking \code{this} from inside a constructor.} Why it matters: this breaks the final-field happens-before guarantee — other threads can see the object before construction finishes, e.g., via a listener registration or a static collection add inside the constructor.

\begin{codebox}
\begin{Verbatim}[commandchars=\\\{\}]
\PY{c+c1}{// Before}
\PY{k+kd}{public}\PY{+w}{ }\PY{n+nf}{Widget}\PY{p}{(}\PY{p}{)}\PY{+w}{ }\PY{p}{\PYZob{}}\PY{+w}{ }\PY{n}{registry}\PY{p}{.}\PY{n+na}{add}\PY{p}{(}\PY{k}{this}\PY{p}{)}\PY{p}{;}\PY{+w}{ }\PY{k}{this}\PY{p}{.}\PY{n+na}{id}\PY{+w}{ }\PY{o}{=}\PY{+w}{ }\PY{n}{computeId}\PY{p}{(}\PY{p}{)}\PY{p}{;}\PY{+w}{ }\PY{p}{\PYZcb{}}\PY{+w}{ }\PY{c+c1}{// \PYZsq{}this\PYZsq{} escapes before id is set}
\PY{c+c1}{// After}
\PY{k+kd}{public}\PY{+w}{ }\PY{n+nf}{Widget}\PY{p}{(}\PY{p}{)}\PY{+w}{ }\PY{p}{\PYZob{}}\PY{+w}{ }\PY{k}{this}\PY{p}{.}\PY{n+na}{id}\PY{+w}{ }\PY{o}{=}\PY{+w}{ }\PY{n}{computeId}\PY{p}{(}\PY{p}{)}\PY{p}{;}\PY{+w}{ }\PY{p}{\PYZcb{}}
\PY{c+c1}{// then, after construction completes: registry.add(widget);}
\end{Verbatim}
\end{codebox}

\item 
\textbf{Choosing \code{AtomicLong} for a hot counter under heavy multi-threaded contention.} Why it matters: every thread’s CAS attempt targets the same memory location, so under high contention most CAS attempts fail and retry, wasting CPU; \code{LongAdder} spreads updates across cells for far better throughput.

\begin{codebox}
\begin{Verbatim}[commandchars=\\\{\}]
\PY{c+c1}{// Before}
\PY{k+kd}{private}\PY{+w}{ }\PY{k+kd}{final}\PY{+w}{ }\PY{n}{AtomicLong}\PY{+w}{ }\PY{n}{hits}\PY{+w}{ }\PY{o}{=}\PY{+w}{ }\PY{k}{new}\PY{+w}{ }\PY{n}{AtomicLong}\PY{p}{(}\PY{p}{)}\PY{p}{;}
\PY{k+kt}{void}\PY{+w}{ }\PY{n+nf}{recordHit}\PY{p}{(}\PY{p}{)}\PY{+w}{ }\PY{p}{\PYZob{}}\PY{+w}{ }\PY{n}{hits}\PY{p}{.}\PY{n+na}{incrementAndGet}\PY{p}{(}\PY{p}{)}\PY{p}{;}\PY{+w}{ }\PY{p}{\PYZcb{}}\PY{+w}{ }\PY{c+c1}{// hot spot under heavy contention}
\PY{c+c1}{// After}
\PY{k+kd}{private}\PY{+w}{ }\PY{k+kd}{final}\PY{+w}{ }\PY{n}{LongAdder}\PY{+w}{ }\PY{n}{hits}\PY{+w}{ }\PY{o}{=}\PY{+w}{ }\PY{k}{new}\PY{+w}{ }\PY{n}{LongAdder}\PY{p}{(}\PY{p}{)}\PY{p}{;}
\PY{k+kt}{void}\PY{+w}{ }\PY{n+nf}{recordHit}\PY{p}{(}\PY{p}{)}\PY{+w}{ }\PY{p}{\PYZob{}}\PY{+w}{ }\PY{n}{hits}\PY{p}{.}\PY{n+na}{increment}\PY{p}{(}\PY{p}{)}\PY{p}{;}\PY{+w}{ }\PY{p}{\PYZcb{}}\PY{+w}{ }\PY{c+c1}{// much less contention; sum() when you need the total}
\end{Verbatim}
\end{codebox}

\item 
\textbf{Treating \code{ReadWriteLock}/\code{StampedLock} as a free performance win without measuring.} Why it matters: for short critical sections or write-heavy workloads, the extra bookkeeping of tracking readers/writers can be slower than a plain \code{synchronized} block or \code{ReentrantLock} — reviewers should ask “was this actually measured?”

\begin{codebox}
\begin{Verbatim}[commandchars=\\\{\}]
\PY{c+c1}{// Before: reflexively reaching for ReentrantReadWriteLock on a tiny, write\PYZhy{}heavy map}
\PY{c+c1}{// After: benchmark synchronized/ReentrantLock vs ReadWriteLock under the REAL read/write ratio}
\PY{c+c1}{//        before committing to the more complex API}
\end{Verbatim}
\end{codebox}

\end{enumerate}


\setcounter{chapter}{13}
\chapter{Concurrency (Advanced)}
\label{chap:14}
\label{h:14:14-concurrency-advanced}
Chapter 13 covered the low-level tools: threads, locks, \code{volatile}, and atomics. Those are the building blocks, but almost nobody writes raw \code{Thread} code in production anymore. This chapter covers the higher-level tools built on top of them — thread pools that run tasks for you, composable async pipelines, queues that coordinate producers and consumers, a map built for concurrent access, and synchronizers that make threads wait for each other in specific patterns. These are the tools that actually show up in day-to-day code and in code review. Examples target Java 21+; where a virtual-thread (Chapter 15) note changes the advice, it’s called out.

\section[Executors Framework]{Executors Framework}
\label{h:14:executors-framework}
Creating a new \code{Thread} for every unit of work is expensive and unbounded — nothing stops you from creating ten thousand threads and crashing the process. A \textbf{thread pool} is a manager that keeps a fixed set of reusable worker threads and hands them tasks from a queue. The \textbf{Executors Framework} is Java’s standard toolkit for thread pools.

Three interfaces matter:

{\tablefont
\begin{xltabular}{\linewidth}{@{}L{0.444}L{1.556}@{}}
\toprule
\rowcolor{tablehdr}
\thd{Interface} & \thd{Adds} \\
\midrule
\endhead
\code{Executor} & Just one method: \code{execute(\allowbreak{}Runnable)}. Fire-and-forget. \\
\code{ExecutorService} & Adds \code{submit()} (returns a \code{Future}), lifecycle (\code{shutdown()}), and batch methods (\code{invokeAll}/\code{invokeAny}). \\
\code{Scheduled\allowbreak{}Executor\allowbreak{}Service} & Adds \code{schedule()}, \code{scheduleAtFixedRate()}, \code{schedule\allowbreak{}With\allowbreak{}Fixed\allowbreak{}Delay()} — run tasks later or repeatedly. \\
\bottomrule
\end{xltabular}
}

\begin{codebox}
\begin{Verbatim}[commandchars=\\\{\}]
\PY{k+kn}{import}\PY{+w}{ }\PY{n+nn}{java.util.concurrent.*}\PY{p}{;}

\PY{k+kd}{public}\PY{+w}{ }\PY{k+kd}{class} \PY{n+nc}{ExecutorBasics}\PY{+w}{ }\PY{p}{\PYZob{}}
\PY{+w}{    }\PY{k+kd}{public}\PY{+w}{ }\PY{k+kd}{static}\PY{+w}{ }\PY{k+kt}{void}\PY{+w}{ }\PY{n+nf}{main}\PY{p}{(}\PY{n}{String}\PY{o}{[}\PY{o}{]}\PY{+w}{ }\PY{n}{args}\PY{p}{)}\PY{+w}{ }\PY{k+kd}{throws}\PY{+w}{ }\PY{n}{InterruptedException}\PY{+w}{ }\PY{p}{\PYZob{}}
\PY{+w}{        }\PY{n}{ExecutorService}\PY{+w}{ }\PY{n}{pool}\PY{+w}{ }\PY{o}{=}\PY{+w}{ }\PY{n}{Executors}\PY{p}{.}\PY{n+na}{newFixedThreadPool}\PY{p}{(}\PY{l+m+mi}{4}\PY{p}{)}\PY{p}{;}
\PY{+w}{        }\PY{n}{pool}\PY{p}{.}\PY{n+na}{execute}\PY{p}{(}\PY{p}{(}\PY{p}{)}\PY{+w}{ }\PY{o}{\PYZhy{}}\PY{o}{\PYZgt{}}\PY{+w}{ }\PY{n}{System}\PY{p}{.}\PY{n+na}{out}\PY{p}{.}\PY{n+na}{println}\PY{p}{(}\PY{l+s}{\PYZdq{}}\PY{l+s}{running on: }\PY{l+s}{\PYZdq{}}\PY{+w}{ }\PY{o}{+}\PY{+w}{ }\PY{n}{Thread}\PY{p}{.}\PY{n+na}{currentThread}\PY{p}{(}\PY{p}{)}\PY{p}{.}\PY{n+na}{getName}\PY{p}{(}\PY{p}{)}\PY{p}{)}\PY{p}{)}\PY{p}{;}
\PY{+w}{        }\PY{n}{pool}\PY{p}{.}\PY{n+na}{shutdown}\PY{p}{(}\PY{p}{)}\PY{p}{;}
\PY{+w}{    }\PY{p}{\PYZcb{}}
\PY{p}{\PYZcb{}}
\end{Verbatim}
\end{codebox}

\subsection[Executors factory methods — and why two of them are risky]{\code{Executors} factory methods — and why two of them are risky}
\label{h:14:executors-factory-methods--and-why-two-of-them-are-risky}
{\tablefontsmall
\begin{xltabular}{\linewidth}{@{}L{0.578}L{0.671}L{0.671}L{2.081}@{}}
\toprule
\rowcolor{tablehdr}
\thd{Factory method} & \thd{Pool shape} & \thd{Queue} & \thd{Risk} \\
\midrule
\endhead
\code{newFixedThreadPool(\allowbreak{}n)} & Fixed \code{n} threads & \textbf{Unbounded} \code{LinkedBlockingQueue} & Tasks pile up forever if producers outrun consumers. No backpressure, memory grows until OOM. \\
\code{newCachedThreadPool()} & 0 core, \textbf{unbounded} max threads & \code{SynchronousQueue} (no storage) & No queueing at all — instead it just creates a new thread for every task that arrives while all existing threads are busy. Under load this can spawn thousands of threads, each with its own stack (\textasciitilde{}512KB–1MB), and the process runs out of memory or the OS runs out of native threads. \\
\code{new\allowbreak{}Single\allowbreak{}Thread\allowbreak{}Executor()} & 1 thread & Unbounded queue & Same unbounded-queue risk as fixed pool, just with one worker. \\
\code{new\allowbreak{}Scheduled\allowbreak{}Thread\allowbreak{}Pool(\allowbreak{}n)} & \code{n} threads & Unbounded delayed queue & Fine for scheduling, same queue caveat applies to backlog. \\
\code{new\allowbreak{}Virtual\allowbreak{}Thread\allowbreak{}Per\allowbreak{}Task\allowbreak{}Executor()} (Java 21) & One virtual thread per task, no pooling & N/A & Different tradeoffs entirely — see Chapter 15. \\
\bottomrule
\end{xltabular}
}

The common thread: \textbf{both risky factories hide an unbounded resource} (queue or thread count) behind a friendly one-line call. In an interview or review, treat \code{Executors.\allowbreak{}new\allowbreak{}Fixed\allowbreak{}Thread\allowbreak{}Pool} and \code{Executors.\allowbreak{}new\allowbreak{}Cached\allowbreak{}Thread\allowbreak{}Pool} as a yellow flag, not an automatic pass — ask “what happens when load exceeds capacity?”

\subsection[ScheduledExecutorService]{\code{Scheduled\allowbreak{}Executor\allowbreak{}Service}}
\label{h:14:scheduledexecutorservice}
Use this when work needs to run later, or repeatedly, instead of right now:

\begin{codebox}
\begin{Verbatim}[commandchars=\\\{\}]
\PY{k+kn}{import}\PY{+w}{ }\PY{n+nn}{java.util.concurrent.*}\PY{p}{;}

\PY{k+kd}{public}\PY{+w}{ }\PY{k+kd}{class} \PY{n+nc}{ScheduledJobs}\PY{+w}{ }\PY{p}{\PYZob{}}
\PY{+w}{    }\PY{k+kd}{public}\PY{+w}{ }\PY{k+kd}{static}\PY{+w}{ }\PY{k+kt}{void}\PY{+w}{ }\PY{n+nf}{main}\PY{p}{(}\PY{n}{String}\PY{o}{[}\PY{o}{]}\PY{+w}{ }\PY{n}{args}\PY{p}{)}\PY{+w}{ }\PY{p}{\PYZob{}}
\PY{+w}{        }\PY{n}{ScheduledExecutorService}\PY{+w}{ }\PY{n}{scheduler}\PY{+w}{ }\PY{o}{=}\PY{+w}{ }\PY{n}{Executors}\PY{p}{.}\PY{n+na}{newScheduledThreadPool}\PY{p}{(}\PY{l+m+mi}{2}\PY{p}{)}\PY{p}{;}

\PY{+w}{        }\PY{c+c1}{// run once, after a delay}
\PY{+w}{        }\PY{n}{scheduler}\PY{p}{.}\PY{n+na}{schedule}\PY{p}{(}\PY{p}{(}\PY{p}{)}\PY{+w}{ }\PY{o}{\PYZhy{}}\PY{o}{\PYZgt{}}\PY{+w}{ }\PY{n}{System}\PY{p}{.}\PY{n+na}{out}\PY{p}{.}\PY{n+na}{println}\PY{p}{(}\PY{l+s}{\PYZdq{}}\PY{l+s}{ran once after delay}\PY{l+s}{\PYZdq{}}\PY{p}{)}\PY{p}{,}\PY{+w}{ }\PY{l+m+mi}{5}\PY{p}{,}\PY{+w}{ }\PY{n}{TimeUnit}\PY{p}{.}\PY{n+na}{SECONDS}\PY{p}{)}\PY{p}{;}

\PY{+w}{        }\PY{c+c1}{// run repeatedly, every 10s, measured from the START of each run}
\PY{+w}{        }\PY{c+c1}{// (if a run takes longer than the period, the next run starts immediately after it finishes)}
\PY{+w}{        }\PY{n}{scheduler}\PY{p}{.}\PY{n+na}{scheduleAtFixedRate}\PY{p}{(}\PY{p}{(}\PY{p}{)}\PY{+w}{ }\PY{o}{\PYZhy{}}\PY{o}{\PYZgt{}}\PY{+w}{ }\PY{n}{pollForUpdates}\PY{p}{(}\PY{p}{)}\PY{p}{,}\PY{+w}{ }\PY{l+m+mi}{0}\PY{p}{,}\PY{+w}{ }\PY{l+m+mi}{10}\PY{p}{,}\PY{+w}{ }\PY{n}{TimeUnit}\PY{p}{.}\PY{n+na}{SECONDS}\PY{p}{)}\PY{p}{;}

\PY{+w}{        }\PY{c+c1}{// run repeatedly, with 10s of rest AFTER each run finishes before the next one starts}
\PY{+w}{        }\PY{n}{scheduler}\PY{p}{.}\PY{n+na}{scheduleWithFixedDelay}\PY{p}{(}\PY{p}{(}\PY{p}{)}\PY{+w}{ }\PY{o}{\PYZhy{}}\PY{o}{\PYZgt{}}\PY{+w}{ }\PY{n}{pollForUpdates}\PY{p}{(}\PY{p}{)}\PY{p}{,}\PY{+w}{ }\PY{l+m+mi}{0}\PY{p}{,}\PY{+w}{ }\PY{l+m+mi}{10}\PY{p}{,}\PY{+w}{ }\PY{n}{TimeUnit}\PY{p}{.}\PY{n+na}{SECONDS}\PY{p}{)}\PY{p}{;}
\PY{+w}{    }\PY{p}{\PYZcb{}}

\PY{+w}{    }\PY{k+kd}{private}\PY{+w}{ }\PY{k+kd}{static}\PY{+w}{ }\PY{k+kt}{void}\PY{+w}{ }\PY{n+nf}{pollForUpdates}\PY{p}{(}\PY{p}{)}\PY{+w}{ }\PY{p}{\PYZob{}}
\PY{+w}{        }\PY{n}{System}\PY{p}{.}\PY{n+na}{out}\PY{p}{.}\PY{n+na}{println}\PY{p}{(}\PY{l+s}{\PYZdq{}}\PY{l+s}{polling at }\PY{l+s}{\PYZdq{}}\PY{+w}{ }\PY{o}{+}\PY{+w}{ }\PY{n}{java}\PY{p}{.}\PY{n+na}{time}\PY{p}{.}\PY{n+na}{Instant}\PY{p}{.}\PY{n+na}{now}\PY{p}{(}\PY{p}{)}\PY{p}{)}\PY{p}{;}
\PY{+w}{    }\PY{p}{\PYZcb{}}
\PY{p}{\PYZcb{}}
\end{Verbatim}
\end{codebox}

The difference between \code{scheduleAtFixedRate} and \code{scheduleWithFixedDelay} trips people up: fixed-rate measures the interval from the \emph{start} of one run to the \emph{start} of the next (so a slow task can cause back-to-back runs with no gap, or even overlap being queued up), while fixed-delay measures the interval from the \emph{end} of one run to the \emph{start} of the next (always leaves a real gap). If a scheduled task throws an uncaught exception, that task’s future scheduled runs \textbf{stop silently} — always wrap the body in a try/catch that logs, so one bad run doesn’t quietly kill all future polling.

\subsection[Building a ThreadPoolExecutor by hand]{Building a \code{ThreadPoolExecutor} by hand}
\label{h:14:building-a-threadpoolexecutor-by-hand}
For anything that matters in production, construct a \code{ThreadPoolExecutor} explicitly so every limit is visible:

\begin{codebox}
\begin{Verbatim}[commandchars=\\\{\}]
\PY{k+kn}{import}\PY{+w}{ }\PY{n+nn}{java.util.concurrent.*}\PY{p}{;}
\PY{k+kn}{import}\PY{+w}{ }\PY{n+nn}{java.util.concurrent.atomic.AtomicInteger}\PY{p}{;}

\PY{k+kd}{public}\PY{+w}{ }\PY{k+kd}{class} \PY{n+nc}{BoundedPoolExample}\PY{+w}{ }\PY{p}{\PYZob{}}
\PY{+w}{    }\PY{k+kd}{public}\PY{+w}{ }\PY{k+kd}{static}\PY{+w}{ }\PY{k+kt}{void}\PY{+w}{ }\PY{n+nf}{main}\PY{p}{(}\PY{n}{String}\PY{o}{[}\PY{o}{]}\PY{+w}{ }\PY{n}{args}\PY{p}{)}\PY{+w}{ }\PY{p}{\PYZob{}}
\PY{+w}{        }\PY{n}{BlockingQueue}\PY{o}{\PYZlt{}}\PY{n}{Runnable}\PY{o}{\PYZgt{}}\PY{+w}{ }\PY{n}{boundedQueue}\PY{+w}{ }\PY{o}{=}\PY{+w}{ }\PY{k}{new}\PY{+w}{ }\PY{n}{ArrayBlockingQueue}\PY{o}{\PYZlt{}}\PY{o}{\PYZgt{}}\PY{p}{(}\PY{l+m+mi}{200}\PY{p}{)}\PY{p}{;}\PY{+w}{ }\PY{c+c1}{// bounded! not unbounded}

\PY{+w}{        }\PY{n}{ThreadFactory}\PY{+w}{ }\PY{n}{namedThreads}\PY{+w}{ }\PY{o}{=}\PY{+w}{ }\PY{k}{new}\PY{+w}{ }\PY{n}{ThreadFactory}\PY{p}{(}\PY{p}{)}\PY{+w}{ }\PY{p}{\PYZob{}}
\PY{+w}{            }\PY{k+kd}{private}\PY{+w}{ }\PY{k+kd}{final}\PY{+w}{ }\PY{n}{AtomicInteger}\PY{+w}{ }\PY{n}{counter}\PY{+w}{ }\PY{o}{=}\PY{+w}{ }\PY{k}{new}\PY{+w}{ }\PY{n}{AtomicInteger}\PY{p}{(}\PY{l+m+mi}{1}\PY{p}{)}\PY{p}{;}
\PY{+w}{            }\PY{n+nd}{@Override}
\PY{+w}{            }\PY{k+kd}{public}\PY{+w}{ }\PY{n}{Thread}\PY{+w}{ }\PY{n+nf}{newThread}\PY{p}{(}\PY{n}{Runnable}\PY{+w}{ }\PY{n}{r}\PY{p}{)}\PY{+w}{ }\PY{p}{\PYZob{}}
\PY{+w}{                }\PY{n}{Thread}\PY{+w}{ }\PY{n}{t}\PY{+w}{ }\PY{o}{=}\PY{+w}{ }\PY{k}{new}\PY{+w}{ }\PY{n}{Thread}\PY{p}{(}\PY{n}{r}\PY{p}{,}\PY{+w}{ }\PY{l+s}{\PYZdq{}}\PY{l+s}{order\PYZhy{}worker\PYZhy{}}\PY{l+s}{\PYZdq{}}\PY{+w}{ }\PY{o}{+}\PY{+w}{ }\PY{n}{counter}\PY{p}{.}\PY{n+na}{getAndIncrement}\PY{p}{(}\PY{p}{)}\PY{p}{)}\PY{p}{;}
\PY{+w}{                }\PY{n}{t}\PY{p}{.}\PY{n+na}{setDaemon}\PY{p}{(}\PY{k+kc}{false}\PY{p}{)}\PY{p}{;}
\PY{+w}{                }\PY{k}{return}\PY{+w}{ }\PY{n}{t}\PY{p}{;}
\PY{+w}{            }\PY{p}{\PYZcb{}}
\PY{+w}{        }\PY{p}{\PYZcb{}}\PY{p}{;}

\PY{+w}{        }\PY{n}{ThreadPoolExecutor}\PY{+w}{ }\PY{n}{executor}\PY{+w}{ }\PY{o}{=}\PY{+w}{ }\PY{k}{new}\PY{+w}{ }\PY{n}{ThreadPoolExecutor}\PY{p}{(}
\PY{+w}{                }\PY{l+m+mi}{4}\PY{p}{,}\PY{+w}{                                  }\PY{c+c1}{// core pool size}
\PY{+w}{                }\PY{l+m+mi}{8}\PY{p}{,}\PY{+w}{                                   }\PY{c+c1}{// max pool size}
\PY{+w}{                }\PY{l+m+mi}{60L}\PY{p}{,}\PY{+w}{ }\PY{n}{TimeUnit}\PY{p}{.}\PY{n+na}{SECONDS}\PY{p}{,}\PY{+w}{                }\PY{c+c1}{// keep\PYZhy{}alive for extra (non\PYZhy{}core) threads}
\PY{+w}{                }\PY{n}{boundedQueue}\PY{p}{,}
\PY{+w}{                }\PY{n}{namedThreads}\PY{p}{,}
\PY{+w}{                }\PY{k}{new}\PY{+w}{ }\PY{n}{ThreadPoolExecutor}\PY{p}{.}\PY{n+na}{CallerRunsPolicy}\PY{p}{(}\PY{p}{)}\PY{+w}{ }\PY{c+c1}{// rejection handler = backpressure}
\PY{+w}{        }\PY{p}{)}\PY{p}{;}

\PY{+w}{        }\PY{n}{executor}\PY{p}{.}\PY{n+na}{execute}\PY{p}{(}\PY{p}{(}\PY{p}{)}\PY{+w}{ }\PY{o}{\PYZhy{}}\PY{o}{\PYZgt{}}\PY{+w}{ }\PY{n}{System}\PY{p}{.}\PY{n+na}{out}\PY{p}{.}\PY{n+na}{println}\PY{p}{(}\PY{l+s}{\PYZdq{}}\PY{l+s}{handled by }\PY{l+s}{\PYZdq{}}\PY{+w}{ }\PY{o}{+}\PY{+w}{ }\PY{n}{Thread}\PY{p}{.}\PY{n+na}{currentThread}\PY{p}{(}\PY{p}{)}\PY{p}{.}\PY{n+na}{getName}\PY{p}{(}\PY{p}{)}\PY{p}{)}\PY{p}{)}\PY{p}{;}
\PY{+w}{        }\PY{n}{executor}\PY{p}{.}\PY{n+na}{shutdown}\PY{p}{(}\PY{p}{)}\PY{p}{;}
\PY{+w}{    }\PY{p}{\PYZcb{}}
\PY{p}{\PYZcb{}}
\end{Verbatim}
\end{codebox}

\subsection[How core size, max size, and the queue interact]{How core size, max size, and the queue interact}
\label{h:14:how-core-size-max-size-and-the-queue-interact}
This order surprises a lot of people, so it’s worth memorizing:

\begin{enumerate}
\item 
If fewer than \code{corePoolSize} threads exist, \textbf{a new thread is created} for the task — even if other core threads happen to be idle.

\item 
Once \code{corePoolSize} threads exist, new tasks go into the \textbf{queue} instead.

\item 
Only when the \textbf{queue is full} does the pool create additional threads, up to \code{maximumPoolSize}.

\item 
If the queue is full \textbf{and} \code{maximumPoolSize} is reached, the task is \textbf{rejected} — handled by the \code{Rejected\allowbreak{}Execution\allowbreak{}Handler}.

\end{enumerate}

This means with an \emph{unbounded} queue, step 3 never triggers — \code{maximumPoolSize} is dead code, because the queue absorbs everything first. That’s exactly why \code{newFixedThreadPool}’s extra threads (beyond the fixed count, since core == max there) never get created, and why \code{newCachedThreadPool}’s \code{SynchronousQueue} never queues — it forces step 3 immediately for every task past the core threads (core is 0), maximizing thread creation instead.

\subsection[RejectedExecutionHandler policies]{\code{Rejected\allowbreak{}Execution\allowbreak{}Handler} policies}
\label{h:14:rejectedexecutionhandler-policies}
{\tablefontsmall
\begin{xltabular}{\linewidth}{@{}L{0.434}L{1.566}@{}}
\toprule
\rowcolor{tablehdr}
\thd{Policy} & \thd{Behavior} \\
\midrule
\endhead
\code{AbortPolicy} (default) & Throws \code{Rejected\allowbreak{}Execution\allowbreak{}Exception}. \\
\code{CallerRunsPolicy} & Runs the task on the \textbf{calling thread} instead — naturally throttles the producer since it’s now busy running the task instead of submitting more. \\
\code{DiscardPolicy} & Silently drops the task. Dangerous — hides failures. \\
\code{DiscardOldestPolicy} & Drops the oldest queued task, then retries submission. \\
\bottomrule
\end{xltabular}
}

You can also implement your own \code{Rejected\allowbreak{}Execution\allowbreak{}Handler} — for example, to log a metric before falling back to \code{CallerRunsPolicy}.

\subsection[submit() vs execute() — and the exception trap]{\code{submit()} vs \code{execute()} — and the exception trap}
\label{h:14:submit-vs-execute--and-the-exception-trap}
\begin{codebox}
\begin{Verbatim}[commandchars=\\\{\}]
\PY{n}{ExecutorService}\PY{+w}{ }\PY{n}{pool}\PY{+w}{ }\PY{o}{=}\PY{+w}{ }\PY{n}{Executors}\PY{p}{.}\PY{n+na}{newFixedThreadPool}\PY{p}{(}\PY{l+m+mi}{2}\PY{p}{)}\PY{p}{;}

\PY{c+c1}{// execute(): fire\PYZhy{}and\PYZhy{}forget. An uncaught exception goes to the thread\PYZsq{}s}
\PY{c+c1}{// UncaughtExceptionHandler (usually prints to stderr) — at least it\PYZsq{}s visible.}
\PY{n}{pool}\PY{p}{.}\PY{n+na}{execute}\PY{p}{(}\PY{p}{(}\PY{p}{)}\PY{+w}{ }\PY{o}{\PYZhy{}}\PY{o}{\PYZgt{}}\PY{+w}{ }\PY{p}{\PYZob{}}\PY{+w}{ }\PY{k}{throw}\PY{+w}{ }\PY{k}{new}\PY{+w}{ }\PY{n}{RuntimeException}\PY{p}{(}\PY{l+s}{\PYZdq{}}\PY{l+s}{boom (execute)}\PY{l+s}{\PYZdq{}}\PY{p}{)}\PY{p}{;}\PY{+w}{ }\PY{p}{\PYZcb{}}\PY{p}{)}\PY{p}{;}

\PY{c+c1}{// submit(): returns a Future. The exception is CAPTURED inside the Future}
\PY{c+c1}{// and only rethrown (wrapped in ExecutionException) when you call get().}
\PY{n}{Future}\PY{o}{\PYZlt{}}\PY{o}{?}\PY{o}{\PYZgt{}}\PY{+w}{ }\PY{n}{future}\PY{+w}{ }\PY{o}{=}\PY{+w}{ }\PY{n}{pool}\PY{p}{.}\PY{n+na}{submit}\PY{p}{(}\PY{p}{(}\PY{p}{)}\PY{+w}{ }\PY{o}{\PYZhy{}}\PY{o}{\PYZgt{}}\PY{+w}{ }\PY{p}{\PYZob{}}\PY{+w}{ }\PY{k}{throw}\PY{+w}{ }\PY{k}{new}\PY{+w}{ }\PY{n}{RuntimeException}\PY{p}{(}\PY{l+s}{\PYZdq{}}\PY{l+s}{boom (submit)}\PY{l+s}{\PYZdq{}}\PY{p}{)}\PY{p}{;}\PY{+w}{ }\PY{p}{\PYZcb{}}\PY{p}{)}\PY{p}{;}
\PY{c+c1}{// If nobody ever calls future.get(), this exception disappears silently!}

\PY{k}{try}\PY{+w}{ }\PY{p}{\PYZob{}}
\PY{+w}{    }\PY{n}{future}\PY{p}{.}\PY{n+na}{get}\PY{p}{(}\PY{p}{)}\PY{p}{;}
\PY{p}{\PYZcb{}}\PY{+w}{ }\PY{k}{catch}\PY{+w}{ }\PY{p}{(}\PY{n}{ExecutionException}\PY{+w}{ }\PY{n}{e}\PY{p}{)}\PY{+w}{ }\PY{p}{\PYZob{}}
\PY{+w}{    }\PY{n}{System}\PY{p}{.}\PY{n+na}{out}\PY{p}{.}\PY{n+na}{println}\PY{p}{(}\PY{l+s}{\PYZdq{}}\PY{l+s}{caught: }\PY{l+s}{\PYZdq{}}\PY{+w}{ }\PY{o}{+}\PY{+w}{ }\PY{n}{e}\PY{p}{.}\PY{n+na}{getCause}\PY{p}{(}\PY{p}{)}\PY{p}{)}\PY{p}{;}\PY{+w}{ }\PY{c+c1}{// this is how you actually find out}
\PY{p}{\PYZcb{}}\PY{+w}{ }\PY{k}{catch}\PY{+w}{ }\PY{p}{(}\PY{n}{InterruptedException}\PY{+w}{ }\PY{n}{e}\PY{p}{)}\PY{+w}{ }\PY{p}{\PYZob{}}
\PY{+w}{    }\PY{n}{Thread}\PY{p}{.}\PY{n+na}{currentThread}\PY{p}{(}\PY{p}{)}\PY{p}{.}\PY{n+na}{interrupt}\PY{p}{(}\PY{p}{)}\PY{p}{;}
\PY{p}{\PYZcb{}}
\end{Verbatim}
\end{codebox}

\textbf{Rule of thumb:} if you use \code{submit()}, you must eventually call \code{get()} (or check \code{isDone()}/\code{isCancelled()}) or handle the exception some other way — otherwise failures vanish.

\subsection[Two-phase shutdown]{Two-phase shutdown}
\label{h:14:two-phase-shutdown}
\code{shutdown()} is graceful: it stops accepting new tasks but lets queued and running tasks finish. \code{shutdownNow()} is aggressive: it interrupts running tasks and returns the list of tasks that were still waiting in the queue, so you can decide what to do with the lost work. The standard idiom combines both:

\begin{codebox}
\begin{Verbatim}[commandchars=\\\{\}]
\PY{k+kd}{public}\PY{+w}{ }\PY{k+kt}{void}\PY{+w}{ }\PY{n+nf}{shutdownAndAwait}\PY{p}{(}\PY{n}{ExecutorService}\PY{+w}{ }\PY{n}{pool}\PY{p}{)}\PY{+w}{ }\PY{p}{\PYZob{}}
\PY{+w}{    }\PY{n}{pool}\PY{p}{.}\PY{n+na}{shutdown}\PY{p}{(}\PY{p}{)}\PY{p}{;}\PY{+w}{ }\PY{c+c1}{// phase 1: stop accepting new work, let existing work finish}
\PY{+w}{    }\PY{k}{try}\PY{+w}{ }\PY{p}{\PYZob{}}
\PY{+w}{        }\PY{k}{if}\PY{+w}{ }\PY{p}{(}\PY{o}{!}\PY{n}{pool}\PY{p}{.}\PY{n+na}{awaitTermination}\PY{p}{(}\PY{l+m+mi}{30}\PY{p}{,}\PY{+w}{ }\PY{n}{TimeUnit}\PY{p}{.}\PY{n+na}{SECONDS}\PY{p}{)}\PY{p}{)}\PY{+w}{ }\PY{p}{\PYZob{}}
\PY{+w}{            }\PY{n}{pool}\PY{p}{.}\PY{n+na}{shutdownNow}\PY{p}{(}\PY{p}{)}\PY{p}{;}\PY{+w}{ }\PY{c+c1}{// phase 2: didn\PYZsq{}t finish in time, force it}
\PY{+w}{            }\PY{k}{if}\PY{+w}{ }\PY{p}{(}\PY{o}{!}\PY{n}{pool}\PY{p}{.}\PY{n+na}{awaitTermination}\PY{p}{(}\PY{l+m+mi}{10}\PY{p}{,}\PY{+w}{ }\PY{n}{TimeUnit}\PY{p}{.}\PY{n+na}{SECONDS}\PY{p}{)}\PY{p}{)}\PY{+w}{ }\PY{p}{\PYZob{}}
\PY{+w}{                }\PY{n}{System}\PY{p}{.}\PY{n+na}{err}\PY{p}{.}\PY{n+na}{println}\PY{p}{(}\PY{l+s}{\PYZdq{}}\PY{l+s}{pool did not terminate}\PY{l+s}{\PYZdq{}}\PY{p}{)}\PY{p}{;}
\PY{+w}{            }\PY{p}{\PYZcb{}}
\PY{+w}{        }\PY{p}{\PYZcb{}}
\PY{+w}{    }\PY{p}{\PYZcb{}}\PY{+w}{ }\PY{k}{catch}\PY{+w}{ }\PY{p}{(}\PY{n}{InterruptedException}\PY{+w}{ }\PY{n}{e}\PY{p}{)}\PY{+w}{ }\PY{p}{\PYZob{}}
\PY{+w}{        }\PY{n}{pool}\PY{p}{.}\PY{n+na}{shutdownNow}\PY{p}{(}\PY{p}{)}\PY{p}{;}\PY{+w}{ }\PY{c+c1}{// this thread was interrupted while waiting; force shutdown too}
\PY{+w}{        }\PY{n}{Thread}\PY{p}{.}\PY{n+na}{currentThread}\PY{p}{(}\PY{p}{)}\PY{p}{.}\PY{n+na}{interrupt}\PY{p}{(}\PY{p}{)}\PY{p}{;}
\PY{+w}{    }\PY{p}{\PYZcb{}}
\PY{p}{\PYZcb{}}
\end{Verbatim}
\end{codebox}

\subsection[invokeAll and invokeAny]{\code{invokeAll} and \code{invokeAny}}
\label{h:14:invokeall-and-invokeany}
\begin{codebox}
\begin{Verbatim}[commandchars=\\\{\}]
\PY{n}{ExecutorService}\PY{+w}{ }\PY{n}{pool}\PY{+w}{ }\PY{o}{=}\PY{+w}{ }\PY{n}{Executors}\PY{p}{.}\PY{n+na}{newFixedThreadPool}\PY{p}{(}\PY{l+m+mi}{3}\PY{p}{)}\PY{p}{;}

\PY{n}{List}\PY{o}{\PYZlt{}}\PY{n}{Callable}\PY{o}{\PYZlt{}}\PY{n}{Integer}\PY{o}{\PYZgt{}}\PY{o}{\PYZgt{}}\PY{+w}{ }\PY{n}{tasks}\PY{+w}{ }\PY{o}{=}\PY{+w}{ }\PY{n}{List}\PY{p}{.}\PY{n+na}{of}\PY{p}{(}
\PY{+w}{        }\PY{p}{(}\PY{p}{)}\PY{+w}{ }\PY{o}{\PYZhy{}}\PY{o}{\PYZgt{}}\PY{+w}{ }\PY{p}{\PYZob{}}\PY{+w}{ }\PY{n}{Thread}\PY{p}{.}\PY{n+na}{sleep}\PY{p}{(}\PY{l+m+mi}{100}\PY{p}{)}\PY{p}{;}\PY{+w}{ }\PY{k}{return}\PY{+w}{ }\PY{l+m+mi}{1}\PY{p}{;}\PY{+w}{ }\PY{p}{\PYZcb{}}\PY{p}{,}
\PY{+w}{        }\PY{p}{(}\PY{p}{)}\PY{+w}{ }\PY{o}{\PYZhy{}}\PY{o}{\PYZgt{}}\PY{+w}{ }\PY{p}{\PYZob{}}\PY{+w}{ }\PY{n}{Thread}\PY{p}{.}\PY{n+na}{sleep}\PY{p}{(}\PY{l+m+mi}{50}\PY{p}{)}\PY{p}{;}\PY{+w}{  }\PY{k}{return}\PY{+w}{ }\PY{l+m+mi}{2}\PY{p}{;}\PY{+w}{ }\PY{p}{\PYZcb{}}\PY{p}{,}
\PY{+w}{        }\PY{p}{(}\PY{p}{)}\PY{+w}{ }\PY{o}{\PYZhy{}}\PY{o}{\PYZgt{}}\PY{+w}{ }\PY{p}{\PYZob{}}\PY{+w}{ }\PY{n}{Thread}\PY{p}{.}\PY{n+na}{sleep}\PY{p}{(}\PY{l+m+mi}{200}\PY{p}{)}\PY{p}{;}\PY{+w}{ }\PY{k}{return}\PY{+w}{ }\PY{l+m+mi}{3}\PY{p}{;}\PY{+w}{ }\PY{p}{\PYZcb{}}
\PY{p}{)}\PY{p}{;}

\PY{c+c1}{// invokeAll: blocks until ALL tasks finish (or the timeout expires), returns a Future per task}
\PY{n}{List}\PY{o}{\PYZlt{}}\PY{n}{Future}\PY{o}{\PYZlt{}}\PY{n}{Integer}\PY{o}{\PYZgt{}}\PY{o}{\PYZgt{}}\PY{+w}{ }\PY{n}{results}\PY{+w}{ }\PY{o}{=}\PY{+w}{ }\PY{n}{pool}\PY{p}{.}\PY{n+na}{invokeAll}\PY{p}{(}\PY{n}{tasks}\PY{p}{,}\PY{+w}{ }\PY{l+m+mi}{5}\PY{p}{,}\PY{+w}{ }\PY{n}{TimeUnit}\PY{p}{.}\PY{n+na}{SECONDS}\PY{p}{)}\PY{p}{;}

\PY{c+c1}{// invokeAny: blocks until the FIRST task succeeds, cancels the rest, returns that value directly}
\PY{n}{Integer}\PY{+w}{ }\PY{n}{fastest}\PY{+w}{ }\PY{o}{=}\PY{+w}{ }\PY{n}{pool}\PY{p}{.}\PY{n+na}{invokeAny}\PY{p}{(}\PY{n}{tasks}\PY{p}{)}\PY{p}{;}
\end{Verbatim}
\end{codebox}

\subsection[ExecutorService as AutoCloseable (Java 19+)]{\code{ExecutorService} as \code{AutoCloseable} (Java 19+)}
\label{h:14:executorservice-as-autocloseable-java-19}
Since JDK 19, \code{ExecutorService} extends \code{AutoCloseable}. \code{close()} calls \code{shutdown()}, waits, and escalates to \code{shutdownNow()} if needed — so try-with-resources gives you the two-phase idiom for free:

\begin{codebox}
\begin{Verbatim}[commandchars=\\\{\}]
\PY{k}{try}\PY{+w}{ }\PY{p}{(}\PY{n}{ExecutorService}\PY{+w}{ }\PY{n}{pool}\PY{+w}{ }\PY{o}{=}\PY{+w}{ }\PY{n}{Executors}\PY{p}{.}\PY{n+na}{newFixedThreadPool}\PY{p}{(}\PY{l+m+mi}{4}\PY{p}{)}\PY{p}{)}\PY{+w}{ }\PY{p}{\PYZob{}}
\PY{+w}{    }\PY{n}{pool}\PY{p}{.}\PY{n+na}{submit}\PY{p}{(}\PY{p}{(}\PY{p}{)}\PY{+w}{ }\PY{o}{\PYZhy{}}\PY{o}{\PYZgt{}}\PY{+w}{ }\PY{n}{System}\PY{p}{.}\PY{n+na}{out}\PY{p}{.}\PY{n+na}{println}\PY{p}{(}\PY{l+s}{\PYZdq{}}\PY{l+s}{work}\PY{l+s}{\PYZdq{}}\PY{p}{)}\PY{p}{)}\PY{p}{;}
\PY{p}{\PYZcb{}}\PY{+w}{ }\PY{c+c1}{// close() called automatically: shutdown() + await, escalating to shutdownNow() if it hangs}
\end{Verbatim}
\end{codebox}

\subsection[Sizing a pool: CPU-bound vs IO-bound]{Sizing a pool: CPU-bound vs IO-bound}
\label{h:14:sizing-a-pool-cpu-bound-vs-io-bound}
\begin{itemize}
\item 
\textbf{CPU-bound work} (tight loops, computation, no blocking): size the pool close to the number of CPU cores (\code{Runtime.\allowbreak{}get\allowbreak{}Runtime().\allowbreak{}available\allowbreak{}Processors()}, maybe \code{+\allowbreak{}1}). More threads than cores just adds context-switch overhead — the CPU is already saturated.

\item 
\textbf{IO-bound work} (waiting on network calls, disk, database): threads spend most of their time \emph{blocked}, not computing, so you need many more of them to keep the CPU busy while some threads wait. A common formula:

\end{itemize}

\begin{codebox}
\begin{Verbatim}
threads = cores * (1 + waitTime / computeTime)
\end{Verbatim}
\end{codebox}

If a task waits 90ms for a network call and computes for 10ms, that’s a wait/compute ratio of 9, so on an 8-core box you’d want roughly \code{8~\allowbreak{}*\allowbreak{}~\allowbreak{}(\allowbreak{}1~\allowbreak{}+\allowbreak{}~\allowbreak{}9)\allowbreak{}~\allowbreak{}=\allowbreak{}~\allowbreak{}80} threads.

Note for Chapter 15: with \textbf{virtual threads}, this whole calculation mostly disappears for IO-bound work — you can have millions of cheap virtual threads blocked on I/O without tying up an OS thread each. The pool-sizing math above is still exactly right for CPU-bound work and for platform-thread pools.

\section[ForkJoinPool]{ForkJoinPool}
\label{h:14:forkjoinpool}
A \code{ForkJoinPool} is a thread pool optimized for \textbf{divide-and-conquer} work: split a big task into smaller subtasks, run them in parallel, combine the results. It’s the engine behind parallel streams and \code{CompletableFuture}’s default async methods.

\subsection[Work stealing]{Work stealing}
\label{h:14:work-stealing}
Each worker thread has its own double-ended queue (deque) of tasks. A thread pushes/pops its \emph{own} tasks from one end (like a normal stack, cheap, no contention). When a thread runs out of work, it “steals” a task from the \emph{other end} of some other busy thread’s deque. This keeps all cores busy without a central lock that everyone contends on — the key reason \code{ForkJoinPool} scales better than a naive shared-queue pool for recursive work.

\subsection[RecursiveTask vs RecursiveAction]{\code{RecursiveTask} vs \code{RecursiveAction}}
\label{h:14:recursivetask-vs-recursiveaction}
Both extend \code{ForkJoinTask} and require implementing \code{compute()}. \code{RecursiveTask\textless{}\allowbreak{}V\textgreater{}} returns a value; \code{RecursiveAction} returns nothing (pure side effect).

\begin{codebox}
\begin{Verbatim}[commandchars=\\\{\}]
\PY{k+kn}{import}\PY{+w}{ }\PY{n+nn}{java.util.concurrent.RecursiveTask}\PY{p}{;}
\PY{k+kn}{import}\PY{+w}{ }\PY{n+nn}{java.util.concurrent.ForkJoinPool}\PY{p}{;}

\PY{k+kd}{public}\PY{+w}{ }\PY{k+kd}{class} \PY{n+nc}{ParallelSum}\PY{+w}{ }\PY{k+kd}{extends}\PY{+w}{ }\PY{n}{RecursiveTask}\PY{o}{\PYZlt{}}\PY{n}{Long}\PY{o}{\PYZgt{}}\PY{+w}{ }\PY{p}{\PYZob{}}
\PY{+w}{    }\PY{k+kd}{private}\PY{+w}{ }\PY{k+kd}{static}\PY{+w}{ }\PY{k+kd}{final}\PY{+w}{ }\PY{k+kt}{int}\PY{+w}{ }\PY{n}{THRESHOLD}\PY{+w}{ }\PY{o}{=}\PY{+w}{ }\PY{l+m+mi}{10\PYZus{}000}\PY{p}{;}\PY{+w}{ }\PY{c+c1}{// tune based on benchmarking}
\PY{+w}{    }\PY{k+kd}{private}\PY{+w}{ }\PY{k+kd}{final}\PY{+w}{ }\PY{k+kt}{long}\PY{o}{[}\PY{o}{]}\PY{+w}{ }\PY{n}{array}\PY{p}{;}
\PY{+w}{    }\PY{k+kd}{private}\PY{+w}{ }\PY{k+kd}{final}\PY{+w}{ }\PY{k+kt}{int}\PY{+w}{ }\PY{n}{start}\PY{p}{,}\PY{+w}{ }\PY{n}{end}\PY{p}{;}

\PY{+w}{    }\PY{n}{ParallelSum}\PY{p}{(}\PY{k+kt}{long}\PY{o}{[}\PY{o}{]}\PY{+w}{ }\PY{n}{array}\PY{p}{,}\PY{+w}{ }\PY{k+kt}{int}\PY{+w}{ }\PY{n}{start}\PY{p}{,}\PY{+w}{ }\PY{k+kt}{int}\PY{+w}{ }\PY{n}{end}\PY{p}{)}\PY{+w}{ }\PY{p}{\PYZob{}}
\PY{+w}{        }\PY{k}{this}\PY{p}{.}\PY{n+na}{array}\PY{+w}{ }\PY{o}{=}\PY{+w}{ }\PY{n}{array}\PY{p}{;}\PY{+w}{ }\PY{k}{this}\PY{p}{.}\PY{n+na}{start}\PY{+w}{ }\PY{o}{=}\PY{+w}{ }\PY{n}{start}\PY{p}{;}\PY{+w}{ }\PY{k}{this}\PY{p}{.}\PY{n+na}{end}\PY{+w}{ }\PY{o}{=}\PY{+w}{ }\PY{n}{end}\PY{p}{;}
\PY{+w}{    }\PY{p}{\PYZcb{}}

\PY{+w}{    }\PY{n+nd}{@Override}
\PY{+w}{    }\PY{k+kd}{protected}\PY{+w}{ }\PY{n}{Long}\PY{+w}{ }\PY{n+nf}{compute}\PY{p}{(}\PY{p}{)}\PY{+w}{ }\PY{p}{\PYZob{}}
\PY{+w}{        }\PY{k+kt}{int}\PY{+w}{ }\PY{n}{length}\PY{+w}{ }\PY{o}{=}\PY{+w}{ }\PY{n}{end}\PY{+w}{ }\PY{o}{\PYZhy{}}\PY{+w}{ }\PY{n}{start}\PY{p}{;}
\PY{+w}{        }\PY{k}{if}\PY{+w}{ }\PY{p}{(}\PY{n}{length}\PY{+w}{ }\PY{o}{\PYZlt{}}\PY{o}{=}\PY{+w}{ }\PY{n}{THRESHOLD}\PY{p}{)}\PY{+w}{ }\PY{p}{\PYZob{}}
\PY{+w}{            }\PY{k+kt}{long}\PY{+w}{ }\PY{n}{sum}\PY{+w}{ }\PY{o}{=}\PY{+w}{ }\PY{l+m+mi}{0}\PY{p}{;}
\PY{+w}{            }\PY{k}{for}\PY{+w}{ }\PY{p}{(}\PY{k+kt}{int}\PY{+w}{ }\PY{n}{i}\PY{+w}{ }\PY{o}{=}\PY{+w}{ }\PY{n}{start}\PY{p}{;}\PY{+w}{ }\PY{n}{i}\PY{+w}{ }\PY{o}{\PYZlt{}}\PY{+w}{ }\PY{n}{end}\PY{p}{;}\PY{+w}{ }\PY{n}{i}\PY{o}{+}\PY{o}{+}\PY{p}{)}\PY{+w}{ }\PY{n}{sum}\PY{+w}{ }\PY{o}{+}\PY{o}{=}\PY{+w}{ }\PY{n}{array}\PY{o}{[}\PY{n}{i}\PY{o}{]}\PY{p}{;}
\PY{+w}{            }\PY{k}{return}\PY{+w}{ }\PY{n}{sum}\PY{p}{;}\PY{+w}{ }\PY{c+c1}{// small enough: just compute directly, don\PYZsq{}t fork further}
\PY{+w}{        }\PY{p}{\PYZcb{}}
\PY{+w}{        }\PY{k+kt}{int}\PY{+w}{ }\PY{n}{mid}\PY{+w}{ }\PY{o}{=}\PY{+w}{ }\PY{n}{start}\PY{+w}{ }\PY{o}{+}\PY{+w}{ }\PY{n}{length}\PY{+w}{ }\PY{o}{/}\PY{+w}{ }\PY{l+m+mi}{2}\PY{p}{;}
\PY{+w}{        }\PY{n}{ParallelSum}\PY{+w}{ }\PY{n}{left}\PY{+w}{ }\PY{o}{=}\PY{+w}{ }\PY{k}{new}\PY{+w}{ }\PY{n}{ParallelSum}\PY{p}{(}\PY{n}{array}\PY{p}{,}\PY{+w}{ }\PY{n}{start}\PY{p}{,}\PY{+w}{ }\PY{n}{mid}\PY{p}{)}\PY{p}{;}
\PY{+w}{        }\PY{n}{ParallelSum}\PY{+w}{ }\PY{n}{right}\PY{+w}{ }\PY{o}{=}\PY{+w}{ }\PY{k}{new}\PY{+w}{ }\PY{n}{ParallelSum}\PY{p}{(}\PY{n}{array}\PY{p}{,}\PY{+w}{ }\PY{n}{mid}\PY{p}{,}\PY{+w}{ }\PY{n}{end}\PY{p}{)}\PY{p}{;}
\PY{+w}{        }\PY{n}{left}\PY{p}{.}\PY{n+na}{fork}\PY{p}{(}\PY{p}{)}\PY{p}{;}\PY{+w}{                 }\PY{c+c1}{// run left asynchronously on another worker}
\PY{+w}{        }\PY{k+kt}{long}\PY{+w}{ }\PY{n}{rightResult}\PY{+w}{ }\PY{o}{=}\PY{+w}{ }\PY{n}{right}\PY{p}{.}\PY{n+na}{compute}\PY{p}{(}\PY{p}{)}\PY{p}{;}\PY{+w}{ }\PY{c+c1}{// compute right on THIS thread (avoid an extra fork)}
\PY{+w}{        }\PY{k+kt}{long}\PY{+w}{ }\PY{n}{leftResult}\PY{+w}{ }\PY{o}{=}\PY{+w}{ }\PY{n}{left}\PY{p}{.}\PY{n+na}{join}\PY{p}{(}\PY{p}{)}\PY{p}{;}\PY{+w}{      }\PY{c+c1}{// wait for left, combine}
\PY{+w}{        }\PY{k}{return}\PY{+w}{ }\PY{n}{leftResult}\PY{+w}{ }\PY{o}{+}\PY{+w}{ }\PY{n}{rightResult}\PY{p}{;}
\PY{+w}{    }\PY{p}{\PYZcb{}}

\PY{+w}{    }\PY{k+kd}{public}\PY{+w}{ }\PY{k+kd}{static}\PY{+w}{ }\PY{k+kt}{void}\PY{+w}{ }\PY{n+nf}{main}\PY{p}{(}\PY{n}{String}\PY{o}{[}\PY{o}{]}\PY{+w}{ }\PY{n}{args}\PY{p}{)}\PY{+w}{ }\PY{p}{\PYZob{}}
\PY{+w}{        }\PY{k+kt}{long}\PY{o}{[}\PY{o}{]}\PY{+w}{ }\PY{n}{data}\PY{+w}{ }\PY{o}{=}\PY{+w}{ }\PY{k}{new}\PY{+w}{ }\PY{k+kt}{long}\PY{o}{[}\PY{l+m+mi}{1\PYZus{}000\PYZus{}000}\PY{o}{]}\PY{p}{;}
\PY{+w}{        }\PY{n}{java}\PY{p}{.}\PY{n+na}{util}\PY{p}{.}\PY{n+na}{Arrays}\PY{p}{.}\PY{n+na}{fill}\PY{p}{(}\PY{n}{data}\PY{p}{,}\PY{+w}{ }\PY{l+m+mi}{1}\PY{p}{)}\PY{p}{;}
\PY{+w}{        }\PY{n}{ForkJoinPool}\PY{+w}{ }\PY{n}{pool}\PY{+w}{ }\PY{o}{=}\PY{+w}{ }\PY{k}{new}\PY{+w}{ }\PY{n}{ForkJoinPool}\PY{p}{(}\PY{p}{)}\PY{p}{;}\PY{+w}{ }\PY{c+c1}{// defaults to availableProcessors()}
\PY{+w}{        }\PY{k+kt}{long}\PY{+w}{ }\PY{n}{total}\PY{+w}{ }\PY{o}{=}\PY{+w}{ }\PY{n}{pool}\PY{p}{.}\PY{n+na}{invoke}\PY{p}{(}\PY{k}{new}\PY{+w}{ }\PY{n}{ParallelSum}\PY{p}{(}\PY{n}{data}\PY{p}{,}\PY{+w}{ }\PY{l+m+mi}{0}\PY{p}{,}\PY{+w}{ }\PY{n}{data}\PY{p}{.}\PY{n+na}{length}\PY{p}{)}\PY{p}{)}\PY{p}{;}
\PY{+w}{        }\PY{n}{System}\PY{p}{.}\PY{n+na}{out}\PY{p}{.}\PY{n+na}{println}\PY{p}{(}\PY{n}{total}\PY{p}{)}\PY{p}{;}\PY{+w}{ }\PY{c+c1}{// 1000000}
\PY{+w}{    }\PY{p}{\PYZcb{}}
\PY{p}{\PYZcb{}}
\end{Verbatim}
\end{codebox}

\textbf{Picking the threshold matters.} Fork too aggressively (tiny tasks) and the overhead of creating/scheduling tasks dwarfs the actual work — slower than sequential. Fork too little (huge tasks) and you don’t get enough parallelism. There’s no universal number; a common starting point is “a few thousand basic operations per leaf task,” then benchmark and adjust.

\subsection[The common pool]{The common pool}
\label{h:14:the-common-pool}
\code{Fork\allowbreak{}Join\allowbreak{}Pool.\allowbreak{}common\allowbreak{}Pool()} is a shared, JVM-wide pool created lazily on first use, sized by default to \code{available\allowbreak{}Processors()\allowbreak{}~\allowbreak{}-\allowbreak{}~\allowbreak{}1}. Two big consumers use it \emph{by default}:

\begin{itemize}
\item 
\textbf{Parallel streams} (\code{.\allowbreak{}parallelStream()}, \code{.\allowbreak{}parallel()}).

\item 
\textbf{\code{CompletableFuture}}’s \code{*\allowbreak{}Async} methods, when you don’t pass an explicit \code{Executor}.

\end{itemize}

Because it’s shared, \textbf{blocking a common-pool thread hurts everyone} — an unrelated parallel stream elsewhere in the same JVM can stall because there simply aren’t enough worker threads left. This is a real production hazard: a slow blocking call buried inside a \code{.\allowbreak{}parallelStream()} or a \code{thenApplyAsync} without an executor can silently degrade throughput app-wide.

You can size the common pool with the system property \code{-\allowbreak{}Djava.\allowbreak{}util.\allowbreak{}concurrent.\allowbreak{}Fork\allowbreak{}Join\allowbreak{}Pool.\allowbreak{}common.\allowbreak{}parallelism=\allowbreak{}N} at JVM startup, or avoid the shared pool entirely by creating your own \code{new~\allowbreak{}Fork\allowbreak{}Join\allowbreak{}Pool(\allowbreak{}parallelism)} for work you want isolated from parallel streams and \code{CompletableFuture} defaults:

\begin{codebox}
\begin{Verbatim}[commandchars=\\\{\}]
\PY{n}{ForkJoinPool}\PY{+w}{ }\PY{n}{dedicatedPool}\PY{+w}{ }\PY{o}{=}\PY{+w}{ }\PY{k}{new}\PY{+w}{ }\PY{n}{ForkJoinPool}\PY{p}{(}\PY{l+m+mi}{4}\PY{p}{)}\PY{p}{;}\PY{+w}{ }\PY{c+c1}{// isolated from the common pool}
\PY{n}{dedicatedPool}\PY{p}{.}\PY{n+na}{submit}\PY{p}{(}\PY{k}{new}\PY{+w}{ }\PY{n}{ParallelSum}\PY{p}{(}\PY{n}{data}\PY{p}{,}\PY{+w}{ }\PY{l+m+mi}{0}\PY{p}{,}\PY{+w}{ }\PY{n}{data}\PY{p}{.}\PY{n+na}{length}\PY{p}{)}\PY{p}{)}\PY{p}{;}
\end{Verbatim}
\end{codebox}

\subsection[ManagedBlocker]{\code{ManagedBlocker}}
\label{h:14:managedblocker}
If a fork/join task genuinely must block (e.g., waiting on I/O), it should announce that via \code{Fork\allowbreak{}Join\allowbreak{}Pool.\allowbreak{}Managed\allowbreak{}Blocker} so the pool can compensate by spinning up a replacement worker, keeping parallelism intact instead of just losing a thread to the block:

\begin{codebox}
\begin{Verbatim}[commandchars=\\\{\}]
\PY{k+kn}{import}\PY{+w}{ }\PY{n+nn}{java.util.concurrent.ForkJoinPool}\PY{p}{;}
\PY{k+kn}{import}\PY{+w}{ }\PY{n+nn}{java.util.concurrent.locks.Lock}\PY{p}{;}

\PY{k+kd}{public}\PY{+w}{ }\PY{k+kd}{class} \PY{n+nc}{LockBlocker}\PY{+w}{ }\PY{k+kd}{implements}\PY{+w}{ }\PY{n}{ForkJoinPool}\PY{p}{.}\PY{n+na}{ManagedBlocker}\PY{+w}{ }\PY{p}{\PYZob{}}
\PY{+w}{    }\PY{k+kd}{private}\PY{+w}{ }\PY{k+kd}{final}\PY{+w}{ }\PY{n}{Lock}\PY{+w}{ }\PY{n}{lock}\PY{p}{;}
\PY{+w}{    }\PY{k+kd}{private}\PY{+w}{ }\PY{k+kt}{boolean}\PY{+w}{ }\PY{n}{acquired}\PY{+w}{ }\PY{o}{=}\PY{+w}{ }\PY{k+kc}{false}\PY{p}{;}

\PY{+w}{    }\PY{n}{LockBlocker}\PY{p}{(}\PY{n}{Lock}\PY{+w}{ }\PY{n}{lock}\PY{p}{)}\PY{+w}{ }\PY{p}{\PYZob{}}\PY{+w}{ }\PY{k}{this}\PY{p}{.}\PY{n+na}{lock}\PY{+w}{ }\PY{o}{=}\PY{+w}{ }\PY{n}{lock}\PY{p}{;}\PY{+w}{ }\PY{p}{\PYZcb{}}

\PY{+w}{    }\PY{n+nd}{@Override}
\PY{+w}{    }\PY{k+kd}{public}\PY{+w}{ }\PY{k+kt}{boolean}\PY{+w}{ }\PY{n+nf}{block}\PY{p}{(}\PY{p}{)}\PY{+w}{ }\PY{p}{\PYZob{}}
\PY{+w}{        }\PY{k}{if}\PY{+w}{ }\PY{p}{(}\PY{o}{!}\PY{n}{acquired}\PY{p}{)}\PY{+w}{ }\PY{n}{acquired}\PY{+w}{ }\PY{o}{=}\PY{+w}{ }\PY{n}{lock}\PY{p}{.}\PY{n+na}{tryLock}\PY{p}{(}\PY{p}{)}\PY{p}{;}
\PY{+w}{        }\PY{k}{return}\PY{+w}{ }\PY{n}{acquired}\PY{p}{;}
\PY{+w}{    }\PY{p}{\PYZcb{}}

\PY{+w}{    }\PY{n+nd}{@Override}
\PY{+w}{    }\PY{k+kd}{public}\PY{+w}{ }\PY{k+kt}{boolean}\PY{+w}{ }\PY{n+nf}{isReleasable}\PY{p}{(}\PY{p}{)}\PY{+w}{ }\PY{p}{\PYZob{}}
\PY{+w}{        }\PY{k}{return}\PY{+w}{ }\PY{n}{acquired}\PY{+w}{ }\PY{o}{|}\PY{o}{|}\PY{+w}{ }\PY{p}{(}\PY{n}{acquired}\PY{+w}{ }\PY{o}{=}\PY{+w}{ }\PY{n}{lock}\PY{p}{.}\PY{n+na}{tryLock}\PY{p}{(}\PY{p}{)}\PY{p}{)}\PY{p}{;}
\PY{+w}{    }\PY{p}{\PYZcb{}}

\PY{+w}{    }\PY{k+kd}{public}\PY{+w}{ }\PY{k+kd}{static}\PY{+w}{ }\PY{k+kt}{void}\PY{+w}{ }\PY{n+nf}{useLock}\PY{p}{(}\PY{n}{Lock}\PY{+w}{ }\PY{n}{lock}\PY{p}{)}\PY{+w}{ }\PY{k+kd}{throws}\PY{+w}{ }\PY{n}{InterruptedException}\PY{+w}{ }\PY{p}{\PYZob{}}
\PY{+w}{        }\PY{n}{ForkJoinPool}\PY{p}{.}\PY{n+na}{managedBlock}\PY{p}{(}\PY{k}{new}\PY{+w}{ }\PY{n}{LockBlocker}\PY{p}{(}\PY{n}{lock}\PY{p}{)}\PY{p}{)}\PY{p}{;}
\PY{+w}{        }\PY{k}{try}\PY{+w}{ }\PY{p}{\PYZob{}}
\PY{+w}{            }\PY{c+c1}{// critical section}
\PY{+w}{        }\PY{p}{\PYZcb{}}\PY{+w}{ }\PY{k}{finally}\PY{+w}{ }\PY{p}{\PYZob{}}
\PY{+w}{            }\PY{n}{lock}\PY{p}{.}\PY{n+na}{unlock}\PY{p}{(}\PY{p}{)}\PY{p}{;}
\PY{+w}{        }\PY{p}{\PYZcb{}}
\PY{+w}{    }\PY{p}{\PYZcb{}}
\PY{p}{\PYZcb{}}
\end{Verbatim}
\end{codebox}

\section[CompletableFuture]{CompletableFuture}
\label{h:14:completablefuture}
\code{CompletableFuture\textless{}\allowbreak{}T\textgreater{}} is a \code{Future} you can \textbf{compose}: chain transformations, combine multiple futures, and react to completion or failure, all without manually calling blocking \code{get()} everywhere.

\subsection[Creating one]{Creating one}
\label{h:14:creating-one}
\begin{codebox}
\begin{Verbatim}[commandchars=\\\{\}]
\PY{k+kn}{import}\PY{+w}{ }\PY{n+nn}{java.util.concurrent.*}\PY{p}{;}

\PY{n}{ExecutorService}\PY{+w}{ }\PY{n}{executor}\PY{+w}{ }\PY{o}{=}\PY{+w}{ }\PY{n}{Executors}\PY{p}{.}\PY{n+na}{newFixedThreadPool}\PY{p}{(}\PY{l+m+mi}{4}\PY{p}{)}\PY{p}{;}

\PY{n}{CompletableFuture}\PY{o}{\PYZlt{}}\PY{n}{String}\PY{o}{\PYZgt{}}\PY{+w}{ }\PY{n}{fromSupplier}\PY{+w}{ }\PY{o}{=}\PY{+w}{ }\PY{n}{CompletableFuture}\PY{p}{.}\PY{n+na}{supplyAsync}\PY{p}{(}
\PY{+w}{        }\PY{p}{(}\PY{p}{)}\PY{+w}{ }\PY{o}{\PYZhy{}}\PY{o}{\PYZgt{}}\PY{+w}{ }\PY{n}{fetchUserName}\PY{p}{(}\PY{l+m+mi}{42}\PY{p}{)}\PY{p}{,}\PY{+w}{ }\PY{n}{executor}\PY{p}{)}\PY{p}{;}\PY{+w}{ }\PY{c+c1}{// runs on the given executor, returns a value}

\PY{n}{CompletableFuture}\PY{o}{\PYZlt{}}\PY{n}{Void}\PY{o}{\PYZgt{}}\PY{+w}{ }\PY{n}{fromRunnable}\PY{+w}{ }\PY{o}{=}\PY{+w}{ }\PY{n}{CompletableFuture}\PY{p}{.}\PY{n+na}{runAsync}\PY{p}{(}
\PY{+w}{        }\PY{p}{(}\PY{p}{)}\PY{+w}{ }\PY{o}{\PYZhy{}}\PY{o}{\PYZgt{}}\PY{+w}{ }\PY{n}{System}\PY{p}{.}\PY{n+na}{out}\PY{p}{.}\PY{n+na}{println}\PY{p}{(}\PY{l+s}{\PYZdq{}}\PY{l+s}{side effect}\PY{l+s}{\PYZdq{}}\PY{p}{)}\PY{p}{,}\PY{+w}{ }\PY{n}{executor}\PY{p}{)}\PY{p}{;}\PY{+w}{ }\PY{c+c1}{// no return value}

\PY{n}{CompletableFuture}\PY{o}{\PYZlt{}}\PY{n}{String}\PY{o}{\PYZgt{}}\PY{+w}{ }\PY{n}{already}\PY{+w}{ }\PY{o}{=}\PY{+w}{ }\PY{n}{CompletableFuture}\PY{p}{.}\PY{n+na}{completedFuture}\PY{p}{(}\PY{l+s}{\PYZdq{}}\PY{l+s}{cached}\PY{l+s}{\PYZdq{}}\PY{p}{)}\PY{p}{;}\PY{+w}{ }\PY{c+c1}{// no async work at all}
\end{Verbatim}
\end{codebox}

\subsection[thenApply vs thenCompose vs thenCombine]{\code{thenApply} vs \code{thenCompose} vs \code{thenCombine}}
\label{h:14:thenapply-vs-thencompose-vs-thencombine}
\begin{codebox}
\begin{Verbatim}[commandchars=\\\{\}]
\PY{c+c1}{// thenApply: transform the result, like Optional.map / Stream.map}
\PY{n}{CompletableFuture}\PY{o}{\PYZlt{}}\PY{n}{Integer}\PY{o}{\PYZgt{}}\PY{+w}{ }\PY{n}{length}\PY{+w}{ }\PY{o}{=}\PY{+w}{ }\PY{n}{fromSupplier}\PY{p}{.}\PY{n+na}{thenApply}\PY{p}{(}\PY{n}{String}\PY{p}{:}\PY{p}{:}\PY{n}{length}\PY{p}{)}\PY{p}{;}

\PY{c+c1}{// thenCompose: chain to ANOTHER CompletableFuture\PYZhy{}returning step, like Optional.flatMap.}
\PY{c+c1}{// Using thenApply here would give you a CompletableFuture\PYZlt{}CompletableFuture\PYZlt{}Address\PYZgt{}\PYZgt{} — wrong shape.}
\PY{n}{CompletableFuture}\PY{o}{\PYZlt{}}\PY{n}{Address}\PY{o}{\PYZgt{}}\PY{+w}{ }\PY{n}{address}\PY{+w}{ }\PY{o}{=}\PY{+w}{ }\PY{n}{fromSupplier}\PY{p}{.}\PY{n+na}{thenCompose}\PY{p}{(}\PY{n}{name}\PY{+w}{ }\PY{o}{\PYZhy{}}\PY{o}{\PYZgt{}}\PY{+w}{ }\PY{n}{lookupAddressAsync}\PY{p}{(}\PY{n}{name}\PY{p}{)}\PY{p}{)}\PY{p}{;}

\PY{c+c1}{// thenCombine: combine two INDEPENDENT futures once both are done}
\PY{n}{CompletableFuture}\PY{o}{\PYZlt{}}\PY{n}{String}\PY{o}{\PYZgt{}}\PY{+w}{ }\PY{n}{userName}\PY{+w}{ }\PY{o}{=}\PY{+w}{ }\PY{n}{CompletableFuture}\PY{p}{.}\PY{n+na}{supplyAsync}\PY{p}{(}\PY{p}{(}\PY{p}{)}\PY{+w}{ }\PY{o}{\PYZhy{}}\PY{o}{\PYZgt{}}\PY{+w}{ }\PY{l+s}{\PYZdq{}}\PY{l+s}{Alice}\PY{l+s}{\PYZdq{}}\PY{p}{)}\PY{p}{;}
\PY{n}{CompletableFuture}\PY{o}{\PYZlt{}}\PY{n}{Integer}\PY{o}{\PYZgt{}}\PY{+w}{ }\PY{n}{userAge}\PY{+w}{ }\PY{o}{=}\PY{+w}{ }\PY{n}{CompletableFuture}\PY{p}{.}\PY{n+na}{supplyAsync}\PY{p}{(}\PY{p}{(}\PY{p}{)}\PY{+w}{ }\PY{o}{\PYZhy{}}\PY{o}{\PYZgt{}}\PY{+w}{ }\PY{l+m+mi}{30}\PY{p}{)}\PY{p}{;}
\PY{n}{CompletableFuture}\PY{o}{\PYZlt{}}\PY{n}{String}\PY{o}{\PYZgt{}}\PY{+w}{ }\PY{n}{combined}\PY{+w}{ }\PY{o}{=}\PY{+w}{ }\PY{n}{userName}\PY{p}{.}\PY{n+na}{thenCombine}\PY{p}{(}\PY{n}{userAge}\PY{p}{,}
\PY{+w}{        }\PY{p}{(}\PY{n}{name}\PY{p}{,}\PY{+w}{ }\PY{n}{age}\PY{p}{)}\PY{+w}{ }\PY{o}{\PYZhy{}}\PY{o}{\PYZgt{}}\PY{+w}{ }\PY{n}{name}\PY{+w}{ }\PY{o}{+}\PY{+w}{ }\PY{l+s}{\PYZdq{}}\PY{l+s}{ is }\PY{l+s}{\PYZdq{}}\PY{+w}{ }\PY{o}{+}\PY{+w}{ }\PY{n}{age}\PY{+w}{ }\PY{o}{+}\PY{+w}{ }\PY{l+s}{\PYZdq{}}\PY{l+s}{ years old}\PY{l+s}{\PYZdq{}}\PY{p}{)}\PY{p}{;}
\end{Verbatim}
\end{codebox}

\subsection[thenAccept / thenRun]{\code{thenAccept} / \code{thenRun}}
\label{h:14:thenaccept--thenrun}
\begin{codebox}
\begin{Verbatim}[commandchars=\\\{\}]
\PY{n}{fromSupplier}\PY{p}{.}\PY{n+na}{thenAccept}\PY{p}{(}\PY{n}{name}\PY{+w}{ }\PY{o}{\PYZhy{}}\PY{o}{\PYZgt{}}\PY{+w}{ }\PY{n}{System}\PY{p}{.}\PY{n+na}{out}\PY{p}{.}\PY{n+na}{println}\PY{p}{(}\PY{l+s}{\PYZdq{}}\PY{l+s}{got: }\PY{l+s}{\PYZdq{}}\PY{+w}{ }\PY{o}{+}\PY{+w}{ }\PY{n}{name}\PY{p}{)}\PY{p}{)}\PY{p}{;}\PY{+w}{ }\PY{c+c1}{// consume result, no return value}
\PY{n}{fromSupplier}\PY{p}{.}\PY{n+na}{thenRun}\PY{p}{(}\PY{p}{(}\PY{p}{)}\PY{+w}{ }\PY{o}{\PYZhy{}}\PY{o}{\PYZgt{}}\PY{+w}{ }\PY{n}{System}\PY{p}{.}\PY{n+na}{out}\PY{p}{.}\PY{n+na}{println}\PY{p}{(}\PY{l+s}{\PYZdq{}}\PY{l+s}{done}\PY{l+s}{\PYZdq{}}\PY{p}{)}\PY{p}{)}\PY{p}{;}\PY{+w}{              }\PY{c+c1}{// ignore the result entirely}
\end{Verbatim}
\end{codebox}

\subsection[The *Async variants — which thread runs what]{The \code{*\allowbreak{}Async} variants — which thread runs what}
\label{h:14:the-async-variants--which-thread-runs-what}
Every \code{then*} method has a plain version and an \code{*\allowbreak{}Async} version (\code{thenApplyAsync}, \code{thenComposeAsync}, …):

\begin{itemize}
\item 
\textbf{Plain} (\code{thenApply}): the continuation runs on whichever thread completed the \emph{previous} stage. If the previous stage was already complete when you attach the continuation, it may even run on the \textbf{calling thread}, synchronously. This is subtle and easy to get wrong in reasoning about which thread does what.

\item 
\textbf{\code{*\allowbreak{}Async} without an executor}: runs on \code{Fork\allowbreak{}Join\allowbreak{}Pool.\allowbreak{}common\allowbreak{}Pool()}.

\item 
\textbf{\code{*\allowbreak{}Async} with an executor}: runs on the executor you pass — the version you want for anything that matters.

\end{itemize}

\begin{codebox}
\begin{Verbatim}[commandchars=\\\{\}]
\PY{n}{CompletableFuture}\PY{p}{.}\PY{n+na}{supplyAsync}\PY{p}{(}\PY{p}{(}\PY{p}{)}\PY{+w}{ }\PY{o}{\PYZhy{}}\PY{o}{\PYZgt{}}\PY{+w}{ }\PY{l+s}{\PYZdq{}}\PY{l+s}{raw}\PY{l+s}{\PYZdq{}}\PY{p}{,}\PY{+w}{ }\PY{n}{executor}\PY{p}{)}
\PY{+w}{        }\PY{p}{.}\PY{n+na}{thenApplyAsync}\PY{p}{(}\PY{n}{String}\PY{p}{:}\PY{p}{:}\PY{n}{toUpperCase}\PY{p}{,}\PY{+w}{ }\PY{n}{executor}\PY{p}{)}\PY{+w}{ }\PY{c+c1}{// explicit: runs on our pool, not the common pool}
\PY{+w}{        }\PY{p}{.}\PY{n+na}{thenAccept}\PY{p}{(}\PY{n}{System}\PY{p}{.}\PY{n+na}{out}\PY{p}{:}\PY{p}{:}\PY{n}{println}\PY{p}{)}\PY{p}{;}
\end{Verbatim}
\end{codebox}

\subsection[allOf / anyOf]{\code{allOf} / \code{anyOf}}
\label{h:14:allof--anyof}
\begin{codebox}
\begin{Verbatim}[commandchars=\\\{\}]
\PY{n}{CompletableFuture}\PY{o}{\PYZlt{}}\PY{n}{String}\PY{o}{\PYZgt{}}\PY{+w}{ }\PY{n}{f1}\PY{+w}{ }\PY{o}{=}\PY{+w}{ }\PY{n}{CompletableFuture}\PY{p}{.}\PY{n+na}{supplyAsync}\PY{p}{(}\PY{p}{(}\PY{p}{)}\PY{+w}{ }\PY{o}{\PYZhy{}}\PY{o}{\PYZgt{}}\PY{+w}{ }\PY{l+s}{\PYZdq{}}\PY{l+s}{a}\PY{l+s}{\PYZdq{}}\PY{p}{,}\PY{+w}{ }\PY{n}{executor}\PY{p}{)}\PY{p}{;}
\PY{n}{CompletableFuture}\PY{o}{\PYZlt{}}\PY{n}{String}\PY{o}{\PYZgt{}}\PY{+w}{ }\PY{n}{f2}\PY{+w}{ }\PY{o}{=}\PY{+w}{ }\PY{n}{CompletableFuture}\PY{p}{.}\PY{n+na}{supplyAsync}\PY{p}{(}\PY{p}{(}\PY{p}{)}\PY{+w}{ }\PY{o}{\PYZhy{}}\PY{o}{\PYZgt{}}\PY{+w}{ }\PY{l+s}{\PYZdq{}}\PY{l+s}{b}\PY{l+s}{\PYZdq{}}\PY{p}{,}\PY{+w}{ }\PY{n}{executor}\PY{p}{)}\PY{p}{;}
\PY{n}{CompletableFuture}\PY{o}{\PYZlt{}}\PY{n}{String}\PY{o}{\PYZgt{}}\PY{+w}{ }\PY{n}{f3}\PY{+w}{ }\PY{o}{=}\PY{+w}{ }\PY{n}{CompletableFuture}\PY{p}{.}\PY{n+na}{supplyAsync}\PY{p}{(}\PY{p}{(}\PY{p}{)}\PY{+w}{ }\PY{o}{\PYZhy{}}\PY{o}{\PYZgt{}}\PY{+w}{ }\PY{l+s}{\PYZdq{}}\PY{l+s}{c}\PY{l+s}{\PYZdq{}}\PY{p}{,}\PY{+w}{ }\PY{n}{executor}\PY{p}{)}\PY{p}{;}

\PY{c+c1}{// allOf: waits for every future; itself completes with Void — you collect results manually}
\PY{n}{CompletableFuture}\PY{o}{\PYZlt{}}\PY{n}{Void}\PY{o}{\PYZgt{}}\PY{+w}{ }\PY{n}{all}\PY{+w}{ }\PY{o}{=}\PY{+w}{ }\PY{n}{CompletableFuture}\PY{p}{.}\PY{n+na}{allOf}\PY{p}{(}\PY{n}{f1}\PY{p}{,}\PY{+w}{ }\PY{n}{f2}\PY{p}{,}\PY{+w}{ }\PY{n}{f3}\PY{p}{)}\PY{p}{;}
\PY{n}{CompletableFuture}\PY{o}{\PYZlt{}}\PY{n}{List}\PY{o}{\PYZlt{}}\PY{n}{String}\PY{o}{\PYZgt{}}\PY{o}{\PYZgt{}}\PY{+w}{ }\PY{n}{collected}\PY{+w}{ }\PY{o}{=}\PY{+w}{ }\PY{n}{all}\PY{p}{.}\PY{n+na}{thenApply}\PY{p}{(}\PY{n}{v}\PY{+w}{ }\PY{o}{\PYZhy{}}\PY{o}{\PYZgt{}}\PY{+w}{ }\PY{n}{List}\PY{p}{.}\PY{n+na}{of}\PY{p}{(}\PY{n}{f1}\PY{p}{.}\PY{n+na}{join}\PY{p}{(}\PY{p}{)}\PY{p}{,}\PY{+w}{ }\PY{n}{f2}\PY{p}{.}\PY{n+na}{join}\PY{p}{(}\PY{p}{)}\PY{p}{,}\PY{+w}{ }\PY{n}{f3}\PY{p}{.}\PY{n+na}{join}\PY{p}{(}\PY{p}{)}\PY{p}{)}\PY{p}{)}\PY{p}{;}

\PY{c+c1}{// anyOf: completes as soon as ANY one completes; type is Object (loses specific type info)}
\PY{n}{CompletableFuture}\PY{o}{\PYZlt{}}\PY{n}{Object}\PY{o}{\PYZgt{}}\PY{+w}{ }\PY{n}{first}\PY{+w}{ }\PY{o}{=}\PY{+w}{ }\PY{n}{CompletableFuture}\PY{p}{.}\PY{n+na}{anyOf}\PY{p}{(}\PY{n}{f1}\PY{p}{,}\PY{+w}{ }\PY{n}{f2}\PY{p}{,}\PY{+w}{ }\PY{n}{f3}\PY{p}{)}\PY{p}{;}
\end{Verbatim}
\end{codebox}

\subsection[Error handling: exceptionally / handle / whenComplete]{Error handling: \code{exceptionally} / \code{handle} / \code{whenComplete}}
\label{h:14:error-handling-exceptionally--handle--whencomplete}
\begin{codebox}
\begin{Verbatim}[commandchars=\\\{\}]
\PY{n}{CompletableFuture}\PY{o}{\PYZlt{}}\PY{n}{Integer}\PY{o}{\PYZgt{}}\PY{+w}{ }\PY{n}{risky}\PY{+w}{ }\PY{o}{=}\PY{+w}{ }\PY{n}{CompletableFuture}\PY{p}{.}\PY{n+na}{supplyAsync}\PY{p}{(}\PY{p}{(}\PY{p}{)}\PY{+w}{ }\PY{o}{\PYZhy{}}\PY{o}{\PYZgt{}}\PY{+w}{ }\PY{p}{\PYZob{}}
\PY{+w}{    }\PY{k}{if}\PY{+w}{ }\PY{p}{(}\PY{n}{Math}\PY{p}{.}\PY{n+na}{random}\PY{p}{(}\PY{p}{)}\PY{+w}{ }\PY{o}{\PYZlt{}}\PY{+w}{ }\PY{l+m+mf}{0.5}\PY{p}{)}\PY{+w}{ }\PY{k}{throw}\PY{+w}{ }\PY{k}{new}\PY{+w}{ }\PY{n}{RuntimeException}\PY{p}{(}\PY{l+s}{\PYZdq{}}\PY{l+s}{failed}\PY{l+s}{\PYZdq{}}\PY{p}{)}\PY{p}{;}
\PY{+w}{    }\PY{k}{return}\PY{+w}{ }\PY{l+m+mi}{42}\PY{p}{;}
\PY{p}{\PYZcb{}}\PY{p}{,}\PY{+w}{ }\PY{n}{executor}\PY{p}{)}\PY{p}{;}

\PY{c+c1}{// exceptionally: only runs on failure, supplies a fallback VALUE, swallows the exception}
\PY{n}{risky}\PY{p}{.}\PY{n+na}{exceptionally}\PY{p}{(}\PY{n}{ex}\PY{+w}{ }\PY{o}{\PYZhy{}}\PY{o}{\PYZgt{}}\PY{+w}{ }\PY{p}{\PYZob{}}
\PY{+w}{    }\PY{n}{System}\PY{p}{.}\PY{n+na}{out}\PY{p}{.}\PY{n+na}{println}\PY{p}{(}\PY{l+s}{\PYZdq{}}\PY{l+s}{recovered from: }\PY{l+s}{\PYZdq{}}\PY{+w}{ }\PY{o}{+}\PY{+w}{ }\PY{n}{ex}\PY{p}{.}\PY{n+na}{getMessage}\PY{p}{(}\PY{p}{)}\PY{p}{)}\PY{p}{;}
\PY{+w}{    }\PY{k}{return}\PY{+w}{ }\PY{o}{\PYZhy{}}\PY{l+m+mi}{1}\PY{p}{;}
\PY{p}{\PYZcb{}}\PY{p}{)}\PY{p}{;}

\PY{c+c1}{// handle: runs on EITHER success or failure, gets both (result, exception) — one is always null}
\PY{n}{risky}\PY{p}{.}\PY{n+na}{handle}\PY{p}{(}\PY{p}{(}\PY{n}{result}\PY{p}{,}\PY{+w}{ }\PY{n}{ex}\PY{p}{)}\PY{+w}{ }\PY{o}{\PYZhy{}}\PY{o}{\PYZgt{}}\PY{+w}{ }\PY{n}{ex}\PY{+w}{ }\PY{o}{!}\PY{o}{=}\PY{+w}{ }\PY{k+kc}{null}\PY{+w}{ }\PY{o}{?}\PY{+w}{ }\PY{o}{\PYZhy{}}\PY{l+m+mi}{1}\PY{+w}{ }\PY{p}{:}\PY{+w}{ }\PY{n}{result}\PY{p}{)}\PY{p}{;}

\PY{c+c1}{// whenComplete: side effect only, does NOT change the result, and re\PYZhy{}throws the original exception}
\PY{n}{risky}\PY{p}{.}\PY{n+na}{whenComplete}\PY{p}{(}\PY{p}{(}\PY{n}{result}\PY{p}{,}\PY{+w}{ }\PY{n}{ex}\PY{p}{)}\PY{+w}{ }\PY{o}{\PYZhy{}}\PY{o}{\PYZgt{}}\PY{+w}{ }\PY{p}{\PYZob{}}
\PY{+w}{    }\PY{k}{if}\PY{+w}{ }\PY{p}{(}\PY{n}{ex}\PY{+w}{ }\PY{o}{!}\PY{o}{=}\PY{+w}{ }\PY{k+kc}{null}\PY{p}{)}\PY{+w}{ }\PY{n}{System}\PY{p}{.}\PY{n+na}{out}\PY{p}{.}\PY{n+na}{println}\PY{p}{(}\PY{l+s}{\PYZdq{}}\PY{l+s}{logging failure: }\PY{l+s}{\PYZdq{}}\PY{+w}{ }\PY{o}{+}\PY{+w}{ }\PY{n}{ex}\PY{p}{)}\PY{p}{;}
\PY{+w}{    }\PY{k}{else}\PY{+w}{ }\PY{n}{System}\PY{p}{.}\PY{n+na}{out}\PY{p}{.}\PY{n+na}{println}\PY{p}{(}\PY{l+s}{\PYZdq{}}\PY{l+s}{logging success: }\PY{l+s}{\PYZdq{}}\PY{+w}{ }\PY{o}{+}\PY{+w}{ }\PY{n}{result}\PY{p}{)}\PY{p}{;}
\PY{p}{\PYZcb{}}\PY{p}{)}\PY{p}{;}
\end{Verbatim}
\end{codebox}

\subsection[Timeouts (Java 9+)]{Timeouts (Java 9+)}
\label{h:14:timeouts-java-9}
\begin{codebox}
\begin{Verbatim}[commandchars=\\\{\}]
\PY{n}{CompletableFuture}\PY{o}{\PYZlt{}}\PY{n}{String}\PY{o}{\PYZgt{}}\PY{+w}{ }\PY{n}{slow}\PY{+w}{ }\PY{o}{=}\PY{+w}{ }\PY{n}{CompletableFuture}\PY{p}{.}\PY{n+na}{supplyAsync}\PY{p}{(}\PY{p}{(}\PY{p}{)}\PY{+w}{ }\PY{o}{\PYZhy{}}\PY{o}{\PYZgt{}}\PY{+w}{ }\PY{p}{\PYZob{}}
\PY{+w}{    }\PY{n}{sleepUninterruptibly}\PY{p}{(}\PY{l+m+mi}{5000}\PY{p}{)}\PY{p}{;}
\PY{+w}{    }\PY{k}{return}\PY{+w}{ }\PY{l+s}{\PYZdq{}}\PY{l+s}{late}\PY{l+s}{\PYZdq{}}\PY{p}{;}
\PY{p}{\PYZcb{}}\PY{p}{,}\PY{+w}{ }\PY{n}{executor}\PY{p}{)}\PY{p}{;}

\PY{c+c1}{// orTimeout: completes EXCEPTIONALLY with a TimeoutException if not done in time}
\PY{n}{slow}\PY{p}{.}\PY{n+na}{orTimeout}\PY{p}{(}\PY{l+m+mi}{1}\PY{p}{,}\PY{+w}{ }\PY{n}{TimeUnit}\PY{p}{.}\PY{n+na}{SECONDS}\PY{p}{)}\PY{p}{;}

\PY{c+c1}{// completeOnTimeout: completes with a FALLBACK VALUE instead of failing}
\PY{n}{slow}\PY{p}{.}\PY{n+na}{completeOnTimeout}\PY{p}{(}\PY{l+s}{\PYZdq{}}\PY{l+s}{default value}\PY{l+s}{\PYZdq{}}\PY{p}{,}\PY{+w}{ }\PY{l+m+mi}{1}\PY{p}{,}\PY{+w}{ }\PY{n}{TimeUnit}\PY{p}{.}\PY{n+na}{SECONDS}\PY{p}{)}\PY{p}{;}
\end{Verbatim}
\end{codebox}

\subsection[Always pass an explicit executor]{Always pass an explicit executor}
\label{h:14:always-pass-an-explicit-executor}
Relying on the default (common pool) mixes your application’s async work with parallel streams and any other library using the common pool. It also makes tests slower/flakier (shared global state) and hides real capacity limits. Prefer:

\begin{codebox}
\begin{Verbatim}[commandchars=\\\{\}]
\PY{n}{CompletableFuture}\PY{p}{.}\PY{n+na}{supplyAsync}\PY{p}{(}\PY{k}{this}\PY{p}{:}\PY{p}{:}\PY{n}{loadFromDatabase}\PY{p}{,}\PY{+w}{ }\PY{n}{ioExecutor}\PY{p}{)}
\PY{+w}{        }\PY{p}{.}\PY{n+na}{thenApplyAsync}\PY{p}{(}\PY{k}{this}\PY{p}{:}\PY{p}{:}\PY{n}{transform}\PY{p}{,}\PY{+w}{ }\PY{n}{cpuExecutor}\PY{p}{)}
\PY{+w}{        }\PY{p}{.}\PY{n+na}{thenAcceptAsync}\PY{p}{(}\PY{k}{this}\PY{p}{:}\PY{p}{:}\PY{n}{publish}\PY{p}{,}\PY{+w}{ }\PY{n}{ioExecutor}\PY{p}{)}\PY{p}{;}
\end{Verbatim}
\end{codebox}

\subsection[Why .join() inside a callback is dangerous]{Why \code{.\allowbreak{}join()} inside a callback is dangerous}
\label{h:14:why-join-inside-a-callback-is-dangerous}
\code{.\allowbreak{}join()} (like \code{.\allowbreak{}get()}) \textbf{blocks the current thread}. If that current thread is itself a worker thread from a \emph{bounded} pool (the common pool, or any fixed executor), and the future you’re joining needs a worker thread from that \emph{same} pool to complete, you can starve or deadlock the pool: all workers end up blocked waiting on each other, and none are left to actually run the work that would unblock them.

\begin{codebox}
\begin{Verbatim}[commandchars=\\\{\}]
\PY{c+c1}{// DANGEROUS: runs on the common pool, then blocks a common\PYZhy{}pool thread waiting on}
\PY{c+c1}{// another common\PYZhy{}pool computation. Under load, this can starve the whole pool.}
\PY{n}{CompletableFuture}\PY{p}{.}\PY{n+na}{supplyAsync}\PY{p}{(}\PY{p}{(}\PY{p}{)}\PY{+w}{ }\PY{o}{\PYZhy{}}\PY{o}{\PYZgt{}}\PY{+w}{ }\PY{p}{\PYZob{}}
\PY{+w}{    }\PY{n}{CompletableFuture}\PY{o}{\PYZlt{}}\PY{n}{String}\PY{o}{\PYZgt{}}\PY{+w}{ }\PY{n}{inner}\PY{+w}{ }\PY{o}{=}\PY{+w}{ }\PY{n}{CompletableFuture}\PY{p}{.}\PY{n+na}{supplyAsync}\PY{p}{(}\PY{p}{(}\PY{p}{)}\PY{+w}{ }\PY{o}{\PYZhy{}}\PY{o}{\PYZgt{}}\PY{+w}{ }\PY{n}{fetchRemote}\PY{p}{(}\PY{p}{)}\PY{p}{)}\PY{p}{;}
\PY{+w}{    }\PY{k}{return}\PY{+w}{ }\PY{n}{inner}\PY{p}{.}\PY{n+na}{join}\PY{p}{(}\PY{p}{)}\PY{p}{;}\PY{+w}{ }\PY{c+c1}{// blocking a pool worker while waiting for another pool worker}
\PY{p}{\PYZcb{}}\PY{p}{)}\PY{p}{;}

\PY{c+c1}{// BETTER: compose instead of blocking — no thread sits idle waiting}
\PY{n}{CompletableFuture}\PY{p}{.}\PY{n+na}{supplyAsync}\PY{p}{(}\PY{p}{(}\PY{p}{)}\PY{+w}{ }\PY{o}{\PYZhy{}}\PY{o}{\PYZgt{}}\PY{+w}{ }\PY{n}{prepareRequest}\PY{p}{(}\PY{p}{)}\PY{p}{)}
\PY{+w}{        }\PY{p}{.}\PY{n+na}{thenCompose}\PY{p}{(}\PY{n}{req}\PY{+w}{ }\PY{o}{\PYZhy{}}\PY{o}{\PYZgt{}}\PY{+w}{ }\PY{n}{fetchRemoteAsync}\PY{p}{(}\PY{n}{req}\PY{p}{)}\PY{p}{)}\PY{p}{;}\PY{+w}{ }\PY{c+c1}{// chained, nobody blocks}
\end{Verbatim}
\end{codebox}

If you must block, use a separate, appropriately-sized executor dedicated to blocking work — never the common pool.

\section[BlockingQueue]{BlockingQueue}
\label{h:14:blockingqueue}
A \code{BlockingQueue\textless{}\allowbreak{}E\textgreater{}} is a queue that can make the caller \textbf{wait}: \code{put()} waits if the queue is full, \code{take()} waits if it’s empty. It’s the standard tool for producer-consumer pipelines.

\subsection[The family]{The family}
\label{h:14:the-family}
{\tablefontsmall
\begin{xltabular}{\linewidth}{@{}L{0.516}L{0.597}L{0.446}L{2.442}@{}}
\toprule
\rowcolor{tablehdr}
\thd{Implementation} & \thd{Bounded?} & \thd{Ordering} & \thd{Notes} \\
\midrule
\endhead
\code{ArrayBlockingQueue} & Yes, fixed at creation & FIFO & Backed by an array; capacity can’t change. \\
\code{LinkedBlockingQueue} & Optional (default unbounded = \code{Integer.\allowbreak{}MAX\_\allowbreak{}VALUE}) & FIFO & Watch out — “unbounded by default” is a common trap. \\
\code{SynchronousQueue} & Zero capacity & N/A & No storage at all — a \code{put()} blocks until a \code{take()} is ready to receive it, and vice versa. Pure hand-off. \\
\code{PriorityBlockingQueue} & Unbounded & Priority order (\code{Comparable}/\code{Comparator}) & \code{put()} never blocks (unbounded); \code{take()} blocks if empty. \\
\code{DelayQueue} & Unbounded & By delay expiry & Elements implement \code{Delayed}; \code{take()} only returns an element once its delay has elapsed. \\
\code{LinkedTransferQueue} & Unbounded & FIFO & Adds \code{transfer()}: like \code{put()}, but waits until a consumer actually receives the element — combines queue and synchronous hand-off. \\
\bottomrule
\end{xltabular}
}

\subsection[put/take vs offer/poll (with timeout) vs add/remove]{\code{put}/\code{take} vs \code{offer}/\code{poll} (with timeout) vs \code{add}/\code{remove}}
\label{h:14:puttake-vs-offerpoll-with-timeout-vs-addremove}
\begin{codebox}
\begin{Verbatim}[commandchars=\\\{\}]
\PY{n}{BlockingQueue}\PY{o}{\PYZlt{}}\PY{n}{String}\PY{o}{\PYZgt{}}\PY{+w}{ }\PY{n}{queue}\PY{+w}{ }\PY{o}{=}\PY{+w}{ }\PY{k}{new}\PY{+w}{ }\PY{n}{ArrayBlockingQueue}\PY{o}{\PYZlt{}}\PY{o}{\PYZgt{}}\PY{p}{(}\PY{l+m+mi}{10}\PY{p}{)}\PY{p}{;}

\PY{n}{queue}\PY{p}{.}\PY{n+na}{put}\PY{p}{(}\PY{l+s}{\PYZdq{}}\PY{l+s}{x}\PY{l+s}{\PYZdq{}}\PY{p}{)}\PY{p}{;}\PY{+w}{              }\PY{c+c1}{// blocks forever if full}
\PY{n}{String}\PY{+w}{ }\PY{n}{taken}\PY{+w}{ }\PY{o}{=}\PY{+w}{ }\PY{n}{queue}\PY{p}{.}\PY{n+na}{take}\PY{p}{(}\PY{p}{)}\PY{p}{;}\PY{+w}{ }\PY{c+c1}{// blocks forever if empty}

\PY{k+kt}{boolean}\PY{+w}{ }\PY{n}{added}\PY{+w}{ }\PY{o}{=}\PY{+w}{ }\PY{n}{queue}\PY{p}{.}\PY{n+na}{offer}\PY{p}{(}\PY{l+s}{\PYZdq{}}\PY{l+s}{y}\PY{l+s}{\PYZdq{}}\PY{p}{,}\PY{+w}{ }\PY{l+m+mi}{500}\PY{p}{,}\PY{+w}{ }\PY{n}{TimeUnit}\PY{p}{.}\PY{n+na}{MILLISECONDS}\PY{p}{)}\PY{p}{;}\PY{+w}{ }\PY{c+c1}{// waits up to 500ms, then gives up}
\PY{n}{String}\PY{+w}{ }\PY{n}{polled}\PY{+w}{ }\PY{o}{=}\PY{+w}{ }\PY{n}{queue}\PY{p}{.}\PY{n+na}{poll}\PY{p}{(}\PY{l+m+mi}{500}\PY{p}{,}\PY{+w}{ }\PY{n}{TimeUnit}\PY{p}{.}\PY{n+na}{MILLISECONDS}\PY{p}{)}\PY{p}{;}\PY{+w}{       }\PY{c+c1}{// same, for reading}

\PY{n}{queue}\PY{p}{.}\PY{n+na}{add}\PY{p}{(}\PY{l+s}{\PYZdq{}}\PY{l+s}{z}\PY{l+s}{\PYZdq{}}\PY{p}{)}\PY{p}{;}\PY{+w}{        }\PY{c+c1}{// throws IllegalStateException immediately if full — no waiting}
\PY{n}{String}\PY{+w}{ }\PY{n}{r}\PY{+w}{ }\PY{o}{=}\PY{+w}{ }\PY{n}{queue}\PY{p}{.}\PY{n+na}{remove}\PY{p}{(}\PY{p}{)}\PY{p}{;}\PY{+w}{ }\PY{c+c1}{// throws NoSuchElementException immediately if empty — no waiting}
\end{Verbatim}
\end{codebox}

Use \code{put}/\code{take} when waiting is fine (typical producer-consumer). Use \code{offer}/\code{poll} with a timeout when you need bounded patience (e.g., don’t wait more than a second). Use \code{add}/\code{remove} only when you’ve already guaranteed capacity/non-emptiness and truly want a hard failure otherwise — rare in practice.

\subsection[Producer-consumer example]{Producer-consumer example}
\label{h:14:producer-consumer-example}
\begin{codebox}
\begin{Verbatim}[commandchars=\\\{\}]
\PY{k+kn}{import}\PY{+w}{ }\PY{n+nn}{java.util.concurrent.*}\PY{p}{;}

\PY{k+kd}{public}\PY{+w}{ }\PY{k+kd}{class} \PY{n+nc}{ProducerConsumer}\PY{+w}{ }\PY{p}{\PYZob{}}
\PY{+w}{    }\PY{k+kd}{private}\PY{+w}{ }\PY{k+kd}{static}\PY{+w}{ }\PY{k+kd}{final}\PY{+w}{ }\PY{n}{String}\PY{+w}{ }\PY{n}{POISON\PYZus{}PILL}\PY{+w}{ }\PY{o}{=}\PY{+w}{ }\PY{l+s}{\PYZdq{}}\PY{l+s}{\PYZus{}\PYZus{}STOP\PYZus{}\PYZus{}}\PY{l+s}{\PYZdq{}}\PY{p}{;}

\PY{+w}{    }\PY{k+kd}{public}\PY{+w}{ }\PY{k+kd}{static}\PY{+w}{ }\PY{k+kt}{void}\PY{+w}{ }\PY{n+nf}{main}\PY{p}{(}\PY{n}{String}\PY{o}{[}\PY{o}{]}\PY{+w}{ }\PY{n}{args}\PY{p}{)}\PY{+w}{ }\PY{k+kd}{throws}\PY{+w}{ }\PY{n}{InterruptedException}\PY{+w}{ }\PY{p}{\PYZob{}}
\PY{+w}{        }\PY{n}{BlockingQueue}\PY{o}{\PYZlt{}}\PY{n}{String}\PY{o}{\PYZgt{}}\PY{+w}{ }\PY{n}{queue}\PY{+w}{ }\PY{o}{=}\PY{+w}{ }\PY{k}{new}\PY{+w}{ }\PY{n}{ArrayBlockingQueue}\PY{o}{\PYZlt{}}\PY{o}{\PYZgt{}}\PY{p}{(}\PY{l+m+mi}{50}\PY{p}{)}\PY{p}{;}\PY{+w}{ }\PY{c+c1}{// bounded = backpressure}
\PY{+w}{        }\PY{n}{ExecutorService}\PY{+w}{ }\PY{n}{pool}\PY{+w}{ }\PY{o}{=}\PY{+w}{ }\PY{n}{Executors}\PY{p}{.}\PY{n+na}{newFixedThreadPool}\PY{p}{(}\PY{l+m+mi}{2}\PY{p}{)}\PY{p}{;}

\PY{+w}{        }\PY{n}{Runnable}\PY{+w}{ }\PY{n}{producer}\PY{+w}{ }\PY{o}{=}\PY{+w}{ }\PY{p}{(}\PY{p}{)}\PY{+w}{ }\PY{o}{\PYZhy{}}\PY{o}{\PYZgt{}}\PY{+w}{ }\PY{p}{\PYZob{}}
\PY{+w}{            }\PY{k}{try}\PY{+w}{ }\PY{p}{\PYZob{}}
\PY{+w}{                }\PY{k}{for}\PY{+w}{ }\PY{p}{(}\PY{k+kt}{int}\PY{+w}{ }\PY{n}{i}\PY{+w}{ }\PY{o}{=}\PY{+w}{ }\PY{l+m+mi}{0}\PY{p}{;}\PY{+w}{ }\PY{n}{i}\PY{+w}{ }\PY{o}{\PYZlt{}}\PY{+w}{ }\PY{l+m+mi}{100}\PY{p}{;}\PY{+w}{ }\PY{n}{i}\PY{o}{+}\PY{o}{+}\PY{p}{)}\PY{+w}{ }\PY{p}{\PYZob{}}
\PY{+w}{                    }\PY{n}{queue}\PY{p}{.}\PY{n+na}{put}\PY{p}{(}\PY{l+s}{\PYZdq{}}\PY{l+s}{job\PYZhy{}}\PY{l+s}{\PYZdq{}}\PY{+w}{ }\PY{o}{+}\PY{+w}{ }\PY{n}{i}\PY{p}{)}\PY{p}{;}\PY{+w}{ }\PY{c+c1}{// blocks if consumer falls behind — natural throttling}
\PY{+w}{                }\PY{p}{\PYZcb{}}
\PY{+w}{                }\PY{n}{queue}\PY{p}{.}\PY{n+na}{put}\PY{p}{(}\PY{n}{POISON\PYZus{}PILL}\PY{p}{)}\PY{p}{;}\PY{+w}{ }\PY{c+c1}{// tell the consumer to stop}
\PY{+w}{            }\PY{p}{\PYZcb{}}\PY{+w}{ }\PY{k}{catch}\PY{+w}{ }\PY{p}{(}\PY{n}{InterruptedException}\PY{+w}{ }\PY{n}{e}\PY{p}{)}\PY{+w}{ }\PY{p}{\PYZob{}}
\PY{+w}{                }\PY{n}{Thread}\PY{p}{.}\PY{n+na}{currentThread}\PY{p}{(}\PY{p}{)}\PY{p}{.}\PY{n+na}{interrupt}\PY{p}{(}\PY{p}{)}\PY{p}{;}
\PY{+w}{            }\PY{p}{\PYZcb{}}
\PY{+w}{        }\PY{p}{\PYZcb{}}\PY{p}{;}

\PY{+w}{        }\PY{n}{Runnable}\PY{+w}{ }\PY{n}{consumer}\PY{+w}{ }\PY{o}{=}\PY{+w}{ }\PY{p}{(}\PY{p}{)}\PY{+w}{ }\PY{o}{\PYZhy{}}\PY{o}{\PYZgt{}}\PY{+w}{ }\PY{p}{\PYZob{}}
\PY{+w}{            }\PY{k}{try}\PY{+w}{ }\PY{p}{\PYZob{}}
\PY{+w}{                }\PY{k}{while}\PY{+w}{ }\PY{p}{(}\PY{k+kc}{true}\PY{p}{)}\PY{+w}{ }\PY{p}{\PYZob{}}
\PY{+w}{                    }\PY{n}{String}\PY{+w}{ }\PY{n}{job}\PY{+w}{ }\PY{o}{=}\PY{+w}{ }\PY{n}{queue}\PY{p}{.}\PY{n+na}{take}\PY{p}{(}\PY{p}{)}\PY{p}{;}
\PY{+w}{                    }\PY{k}{if}\PY{+w}{ }\PY{p}{(}\PY{n}{job}\PY{p}{.}\PY{n+na}{equals}\PY{p}{(}\PY{n}{POISON\PYZus{}PILL}\PY{p}{)}\PY{p}{)}\PY{+w}{ }\PY{k}{break}\PY{p}{;}\PY{+w}{ }\PY{c+c1}{// sentinel: no more work coming}
\PY{+w}{                    }\PY{n}{System}\PY{p}{.}\PY{n+na}{out}\PY{p}{.}\PY{n+na}{println}\PY{p}{(}\PY{l+s}{\PYZdq{}}\PY{l+s}{processing }\PY{l+s}{\PYZdq{}}\PY{+w}{ }\PY{o}{+}\PY{+w}{ }\PY{n}{job}\PY{p}{)}\PY{p}{;}
\PY{+w}{                }\PY{p}{\PYZcb{}}
\PY{+w}{            }\PY{p}{\PYZcb{}}\PY{+w}{ }\PY{k}{catch}\PY{+w}{ }\PY{p}{(}\PY{n}{InterruptedException}\PY{+w}{ }\PY{n}{e}\PY{p}{)}\PY{+w}{ }\PY{p}{\PYZob{}}
\PY{+w}{                }\PY{n}{Thread}\PY{p}{.}\PY{n+na}{currentThread}\PY{p}{(}\PY{p}{)}\PY{p}{.}\PY{n+na}{interrupt}\PY{p}{(}\PY{p}{)}\PY{p}{;}
\PY{+w}{            }\PY{p}{\PYZcb{}}
\PY{+w}{        }\PY{p}{\PYZcb{}}\PY{p}{;}

\PY{+w}{        }\PY{n}{pool}\PY{p}{.}\PY{n+na}{submit}\PY{p}{(}\PY{n}{producer}\PY{p}{)}\PY{p}{;}
\PY{+w}{        }\PY{n}{pool}\PY{p}{.}\PY{n+na}{submit}\PY{p}{(}\PY{n}{consumer}\PY{p}{)}\PY{p}{;}
\PY{+w}{        }\PY{n}{pool}\PY{p}{.}\PY{n+na}{shutdown}\PY{p}{(}\PY{p}{)}\PY{p}{;}
\PY{+w}{        }\PY{n}{pool}\PY{p}{.}\PY{n+na}{awaitTermination}\PY{p}{(}\PY{l+m+mi}{1}\PY{p}{,}\PY{+w}{ }\PY{n}{TimeUnit}\PY{p}{.}\PY{n+na}{MINUTES}\PY{p}{)}\PY{p}{;}
\PY{+w}{    }\PY{p}{\PYZcb{}}
\PY{p}{\PYZcb{}}
\end{Verbatim}
\end{codebox}

\textbf{Poison pills} are sentinel values that mean “stop, there’s no more work.” With multiple consumers, you typically send one pill per consumer (or use a shared atomic “shutting down” flag) so every consumer thread gets the signal to exit.

\subsection[Bounded queues are backpressure]{Bounded queues are backpressure}
\label{h:14:bounded-queues-are-backpressure}
A bounded \code{BlockingQueue} is a \textbf{feedback mechanism}. When the queue fills up, \code{put()} blocks the producer — which slows it down to match the consumer’s real speed. An unbounded queue removes that feedback: the producer never feels resistance, memory quietly grows, and the failure shows up later as an \code{OutOfMemoryError} far from its actual cause. Choosing a bounded queue is a deliberate design decision to fail fast (or throttle) instead of failing silently later.

\section[ConcurrentHashMap Internals]{ConcurrentHashMap Internals}
\label{h:14:concurrenthashmap-internals}
\code{ConcurrentHashMap} is a thread-safe map designed for \textbf{high concurrency}, not just thread safety. Unlike a synchronized \code{HashMap} (one lock for everything), it splits work so many threads can operate on different parts of the map at the same time.

\subsection[No full-table lock]{No full-table lock}
\label{h:14:no-full-table-lock}
There is no single lock guarding the whole map. Reads generally don’t lock at all (volatile reads of table slots). Writes lock only what they need to.

\subsection[CAS on an empty bin]{CAS on an empty bin}
\label{h:14:cas-on-an-empty-bin}
Inserting into an \textbf{empty bin} (bucket) uses a \textbf{CAS} (compare-and-swap) — a lock-free atomic operation — instead of taking a lock. This is the fast path: most insertions into a well-sized map hit empty or nearly-empty bins.

\subsection[Per-bin synchronization]{Per-bin synchronization}
\label{h:14:per-bin-synchronization}
When a bin already has entries (a collision, or an update to an existing key), the map synchronizes on that bin’s \textbf{first node} only — not the whole table. Two threads working on different bins never contend with each other at all.

\subsection[Treeification]{Treeification}
\label{h:14:treeification}
If a single bin’s linked list grows past a threshold (8 nodes, and the table itself has at least 64 slots), that bin is converted into a small \textbf{red-black tree}, turning worst-case lookup from O(n) into O(log n) for that bin. It converts back to a linked list if the bin shrinks below 6 nodes. This is an internal defense against pathological hash collisions (e.g., poor \code{hashCode()} implementations or deliberately crafted keys) — you never interact with it directly, but it’s worth knowing it exists.

\subsection[size() / mappingCount() are estimates]{\code{size()} / \code{mappingCount()} are estimates}
\label{h:14:size--mappingcount-are-estimates}
Because there’s no global lock, there’s no way to get a perfectly exact count under concurrent modification without stopping the world. \code{size()} and \code{mappingCount()} (the \code{long}-returning, Java 8+ preferred alternative) use a scalable striped counter internally and sum it up on demand — accurate as of \emph{some} moment, but another thread could be mid-insert at the exact time you read it. Treat both as approximate under concurrent writes; don’t use them for correctness-critical logic (e.g., “if size == 0 then…”).

\subsection[Atomic compute / computeIfAbsent / merge / putIfAbsent]{Atomic \code{compute} / \code{computeIfAbsent} / \code{merge} / \code{putIfAbsent}}
\label{h:14:atomic-compute--computeifabsent--merge--putifabsent}
These perform an atomic read-modify-write for a single key, holding that key’s bin lock for the duration of the supplied function:

\begin{codebox}
\begin{Verbatim}[commandchars=\\\{\}]
\PY{n}{ConcurrentHashMap}\PY{o}{\PYZlt{}}\PY{n}{String}\PY{p}{,}\PY{+w}{ }\PY{n}{Integer}\PY{o}{\PYZgt{}}\PY{+w}{ }\PY{n}{counts}\PY{+w}{ }\PY{o}{=}\PY{+w}{ }\PY{k}{new}\PY{+w}{ }\PY{n}{ConcurrentHashMap}\PY{o}{\PYZlt{}}\PY{o}{\PYZgt{}}\PY{p}{(}\PY{p}{)}\PY{p}{;}

\PY{c+c1}{// atomic increment — no lost updates even with many threads calling this concurrently}
\PY{n}{counts}\PY{p}{.}\PY{n+na}{compute}\PY{p}{(}\PY{l+s}{\PYZdq{}}\PY{l+s}{clicks}\PY{l+s}{\PYZdq{}}\PY{p}{,}\PY{+w}{ }\PY{p}{(}\PY{n}{key}\PY{p}{,}\PY{+w}{ }\PY{n}{current}\PY{p}{)}\PY{+w}{ }\PY{o}{\PYZhy{}}\PY{o}{\PYZgt{}}\PY{+w}{ }\PY{n}{current}\PY{+w}{ }\PY{o}{=}\PY{o}{=}\PY{+w}{ }\PY{k+kc}{null}\PY{+w}{ }\PY{o}{?}\PY{+w}{ }\PY{l+m+mi}{1}\PY{+w}{ }\PY{p}{:}\PY{+w}{ }\PY{n}{current}\PY{+w}{ }\PY{o}{+}\PY{+w}{ }\PY{l+m+mi}{1}\PY{p}{)}\PY{p}{;}

\PY{c+c1}{// only computes the value if the key is absent — common for caches}
\PY{n}{counts}\PY{p}{.}\PY{n+na}{computeIfAbsent}\PY{p}{(}\PY{l+s}{\PYZdq{}}\PY{l+s}{views}\PY{l+s}{\PYZdq{}}\PY{p}{,}\PY{+w}{ }\PY{n}{key}\PY{+w}{ }\PY{o}{\PYZhy{}}\PY{o}{\PYZgt{}}\PY{+w}{ }\PY{n}{expensiveInitialValue}\PY{p}{(}\PY{p}{)}\PY{p}{)}\PY{p}{;}

\PY{c+c1}{// merge: combine an existing value with a new one, or just insert if absent}
\PY{n}{counts}\PY{p}{.}\PY{n+na}{merge}\PY{p}{(}\PY{l+s}{\PYZdq{}}\PY{l+s}{clicks}\PY{l+s}{\PYZdq{}}\PY{p}{,}\PY{+w}{ }\PY{l+m+mi}{1}\PY{p}{,}\PY{+w}{ }\PY{n}{Integer}\PY{p}{:}\PY{p}{:}\PY{n}{sum}\PY{p}{)}\PY{p}{;}

\PY{c+c1}{// putIfAbsent: insert only if key is not already present, atomically}
\PY{n}{counts}\PY{p}{.}\PY{n+na}{putIfAbsent}\PY{p}{(}\PY{l+s}{\PYZdq{}}\PY{l+s}{clicks}\PY{l+s}{\PYZdq{}}\PY{p}{,}\PY{+w}{ }\PY{l+m+mi}{0}\PY{p}{)}\PY{p}{;}
\end{Verbatim}
\end{codebox}

\subsection[The recursive-update deadlock hazard]{The recursive-update deadlock hazard}
\label{h:14:the-recursive-update-deadlock-hazard}
The function you pass to \code{computeIfAbsent}/\code{compute}/\code{merge} runs \textbf{while the bin’s lock is held}. If that function calls back into the \emph{same map} — especially for a key that could hash to the same bin, or worse, the same key — you can deadlock (the thread tries to re-acquire a lock it already holds in a way that isn’t reentrant for this purpose) or, since Java 9+, the map explicitly detects some of these cases and throws \code{Illegal\allowbreak{}State\allowbreak{}Exception:\allowbreak{}~\allowbreak{}Recursive~\allowbreak{}update}.

\begin{codebox}
\begin{Verbatim}[commandchars=\\\{\}]
\PY{n}{ConcurrentHashMap}\PY{o}{\PYZlt{}}\PY{n}{String}\PY{p}{,}\PY{+w}{ }\PY{n}{Integer}\PY{o}{\PYZgt{}}\PY{+w}{ }\PY{n}{map}\PY{+w}{ }\PY{o}{=}\PY{+w}{ }\PY{k}{new}\PY{+w}{ }\PY{n}{ConcurrentHashMap}\PY{o}{\PYZlt{}}\PY{o}{\PYZgt{}}\PY{p}{(}\PY{p}{)}\PY{p}{;}

\PY{c+c1}{// DANGEROUS: the remapping function touches the SAME map it\PYZsq{}s being called from}
\PY{n}{map}\PY{p}{.}\PY{n+na}{computeIfAbsent}\PY{p}{(}\PY{l+s}{\PYZdq{}}\PY{l+s}{a}\PY{l+s}{\PYZdq{}}\PY{p}{,}\PY{+w}{ }\PY{n}{key}\PY{+w}{ }\PY{o}{\PYZhy{}}\PY{o}{\PYZgt{}}\PY{+w}{ }\PY{p}{\PYZob{}}
\PY{+w}{    }\PY{k}{return}\PY{+w}{ }\PY{n}{map}\PY{p}{.}\PY{n+na}{computeIfAbsent}\PY{p}{(}\PY{l+s}{\PYZdq{}}\PY{l+s}{b}\PY{l+s}{\PYZdq{}}\PY{p}{,}\PY{+w}{ }\PY{n}{k2}\PY{+w}{ }\PY{o}{\PYZhy{}}\PY{o}{\PYZgt{}}\PY{+w}{ }\PY{l+m+mi}{1}\PY{p}{)}\PY{p}{;}\PY{+w}{ }\PY{c+c1}{// may deadlock or throw IllegalStateException}
\PY{p}{\PYZcb{}}\PY{p}{)}\PY{p}{;}

\PY{c+c1}{// SAFER: never call back into the same map from inside compute/computeIfAbsent/merge.}
\PY{c+c1}{// Compute the value independently first, then insert.}
\PY{n}{Integer}\PY{+w}{ }\PY{n}{bValue}\PY{+w}{ }\PY{o}{=}\PY{+w}{ }\PY{n}{map}\PY{p}{.}\PY{n+na}{computeIfAbsent}\PY{p}{(}\PY{l+s}{\PYZdq{}}\PY{l+s}{b}\PY{l+s}{\PYZdq{}}\PY{p}{,}\PY{+w}{ }\PY{n}{k2}\PY{+w}{ }\PY{o}{\PYZhy{}}\PY{o}{\PYZgt{}}\PY{+w}{ }\PY{l+m+mi}{1}\PY{p}{)}\PY{p}{;}
\PY{n}{map}\PY{p}{.}\PY{n+na}{computeIfAbsent}\PY{p}{(}\PY{l+s}{\PYZdq{}}\PY{l+s}{a}\PY{l+s}{\PYZdq{}}\PY{p}{,}\PY{+w}{ }\PY{n}{key}\PY{+w}{ }\PY{o}{\PYZhy{}}\PY{o}{\PYZgt{}}\PY{+w}{ }\PY{n}{bValue}\PY{p}{)}\PY{p}{;}
\end{Verbatim}
\end{codebox}

\subsection[Weakly-consistent iterators]{Weakly-consistent iterators}
\label{h:14:weakly-consistent-iterators}
Iterators over a \code{ConcurrentHashMap} (keys, values, entries) are \textbf{weakly consistent}: they never throw \code{Concurrent\allowbreak{}Modification\allowbreak{}Exception}, they reflect the state of the map at some point at or since the iterator was created, and they’re guaranteed not to return the same mapping twice — but they might or might not reflect an update made by another thread mid-iteration. This is different from \code{HashMap}’s fail-fast iterators, and it’s exactly what you want for a map that’s actively being written to by other threads while you iterate.

\subsection[Bulk operations: forEach / search / reduce]{Bulk operations: \code{forEach} / \code{search} / \code{reduce}}
\label{h:14:bulk-operations-foreach--search--reduce}
\begin{codebox}
\begin{Verbatim}[commandchars=\\\{\}]
\PY{n}{ConcurrentHashMap}\PY{o}{\PYZlt{}}\PY{n}{String}\PY{p}{,}\PY{+w}{ }\PY{n}{Integer}\PY{o}{\PYZgt{}}\PY{+w}{ }\PY{n}{scores}\PY{+w}{ }\PY{o}{=}\PY{+w}{ }\PY{k}{new}\PY{+w}{ }\PY{n}{ConcurrentHashMap}\PY{o}{\PYZlt{}}\PY{o}{\PYZgt{}}\PY{p}{(}\PY{p}{)}\PY{p}{;}
\PY{n}{scores}\PY{p}{.}\PY{n+na}{put}\PY{p}{(}\PY{l+s}{\PYZdq{}}\PY{l+s}{alice}\PY{l+s}{\PYZdq{}}\PY{p}{,}\PY{+w}{ }\PY{l+m+mi}{90}\PY{p}{)}\PY{p}{;}
\PY{n}{scores}\PY{p}{.}\PY{n+na}{put}\PY{p}{(}\PY{l+s}{\PYZdq{}}\PY{l+s}{bob}\PY{l+s}{\PYZdq{}}\PY{p}{,}\PY{+w}{ }\PY{l+m+mi}{75}\PY{p}{)}\PY{p}{;}
\PY{n}{scores}\PY{p}{.}\PY{n+na}{put}\PY{p}{(}\PY{l+s}{\PYZdq{}}\PY{l+s}{carol}\PY{l+s}{\PYZdq{}}\PY{p}{,}\PY{+w}{ }\PY{l+m+mi}{88}\PY{p}{)}\PY{p}{;}

\PY{c+c1}{// parallelismThreshold: below this element count, runs sequentially; above it, may run in parallel}
\PY{n}{scores}\PY{p}{.}\PY{n+na}{forEach}\PY{p}{(}\PY{l+m+mi}{1}\PY{p}{,}\PY{+w}{ }\PY{p}{(}\PY{n}{key}\PY{p}{,}\PY{+w}{ }\PY{n}{value}\PY{p}{)}\PY{+w}{ }\PY{o}{\PYZhy{}}\PY{o}{\PYZgt{}}\PY{+w}{ }\PY{n}{System}\PY{p}{.}\PY{n+na}{out}\PY{p}{.}\PY{n+na}{println}\PY{p}{(}\PY{n}{key}\PY{+w}{ }\PY{o}{+}\PY{+w}{ }\PY{l+s}{\PYZdq{}}\PY{l+s}{=}\PY{l+s}{\PYZdq{}}\PY{+w}{ }\PY{o}{+}\PY{+w}{ }\PY{n}{value}\PY{p}{)}\PY{p}{)}\PY{p}{;}

\PY{n}{String}\PY{+w}{ }\PY{n}{highScorer}\PY{+w}{ }\PY{o}{=}\PY{+w}{ }\PY{n}{scores}\PY{p}{.}\PY{n+na}{search}\PY{p}{(}\PY{l+m+mi}{1}\PY{p}{,}\PY{+w}{ }\PY{p}{(}\PY{n}{key}\PY{p}{,}\PY{+w}{ }\PY{n}{value}\PY{p}{)}\PY{+w}{ }\PY{o}{\PYZhy{}}\PY{o}{\PYZgt{}}\PY{+w}{ }\PY{n}{value}\PY{+w}{ }\PY{o}{\PYZgt{}}\PY{+w}{ }\PY{l+m+mi}{85}\PY{+w}{ }\PY{o}{?}\PY{+w}{ }\PY{n}{key}\PY{+w}{ }\PY{p}{:}\PY{+w}{ }\PY{k+kc}{null}\PY{p}{)}\PY{p}{;}

\PY{k+kt}{int}\PY{+w}{ }\PY{n}{total}\PY{+w}{ }\PY{o}{=}\PY{+w}{ }\PY{n}{scores}\PY{p}{.}\PY{n+na}{reduce}\PY{p}{(}\PY{l+m+mi}{1}\PY{p}{,}\PY{+w}{ }\PY{p}{(}\PY{n}{key}\PY{p}{,}\PY{+w}{ }\PY{n}{value}\PY{p}{)}\PY{+w}{ }\PY{o}{\PYZhy{}}\PY{o}{\PYZgt{}}\PY{+w}{ }\PY{n}{value}\PY{p}{,}\PY{+w}{ }\PY{n}{Integer}\PY{p}{:}\PY{p}{:}\PY{n}{sum}\PY{p}{)}\PY{p}{;}
\end{Verbatim}
\end{codebox}

\subsection[ConcurrentHashMap.newKeySet()]{\code{Concurrent\allowbreak{}Hash\allowbreak{}Map.\allowbreak{}new\allowbreak{}Key\allowbreak{}Set()}}
\label{h:14:concurrenthashmapnewkeyset}
A quick way to get a thread-safe \code{Set\textless{}\allowbreak{}E\textgreater{}} backed by a \code{Concurrent\allowbreak{}Hash\allowbreak{}Map\textless{}\allowbreak{}E,\allowbreak{}~\allowbreak{}Boolean\textgreater{}} internally — simpler than wrapping with \code{Collections.\allowbreak{}new\allowbreak{}Set\allowbreak{}From\allowbreak{}Map}:

\begin{codebox}
\begin{Verbatim}[commandchars=\\\{\}]
\PY{n}{Set}\PY{o}{\PYZlt{}}\PY{n}{String}\PY{o}{\PYZgt{}}\PY{+w}{ }\PY{n}{activeSessions}\PY{+w}{ }\PY{o}{=}\PY{+w}{ }\PY{n}{ConcurrentHashMap}\PY{p}{.}\PY{n+na}{newKeySet}\PY{p}{(}\PY{p}{)}\PY{p}{;}
\PY{n}{activeSessions}\PY{p}{.}\PY{n+na}{add}\PY{p}{(}\PY{l+s}{\PYZdq{}}\PY{l+s}{session\PYZhy{}123}\PY{l+s}{\PYZdq{}}\PY{p}{)}\PY{p}{;}
\PY{n}{activeSessions}\PY{p}{.}\PY{n+na}{remove}\PY{p}{(}\PY{l+s}{\PYZdq{}}\PY{l+s}{session\PYZhy{}456}\PY{l+s}{\PYZdq{}}\PY{p}{)}\PY{p}{;}
\end{Verbatim}
\end{codebox}

\section[CountDownLatch]{CountDownLatch}
\label{h:14:countdownlatch}
A \code{CountDownLatch} lets one or more threads wait until a set of events has happened. It’s initialized with a count; each event calls \code{countDown()}; waiting threads call \code{await()} and are released once the count hits zero. \textbf{It’s one-shot} — once it reaches zero, it stays there forever. It cannot be reset or reused.

Typical use: a main thread waits for several background services or worker threads to finish starting up before continuing.

\begin{codebox}
\begin{Verbatim}[commandchars=\\\{\}]
\PY{k+kn}{import}\PY{+w}{ }\PY{n+nn}{java.util.concurrent.CountDownLatch}\PY{p}{;}

\PY{k+kd}{public}\PY{+w}{ }\PY{k+kd}{class} \PY{n+nc}{ServiceStartup}\PY{+w}{ }\PY{p}{\PYZob{}}
\PY{+w}{    }\PY{k+kd}{public}\PY{+w}{ }\PY{k+kd}{static}\PY{+w}{ }\PY{k+kt}{void}\PY{+w}{ }\PY{n+nf}{main}\PY{p}{(}\PY{n}{String}\PY{o}{[}\PY{o}{]}\PY{+w}{ }\PY{n}{args}\PY{p}{)}\PY{+w}{ }\PY{k+kd}{throws}\PY{+w}{ }\PY{n}{InterruptedException}\PY{+w}{ }\PY{p}{\PYZob{}}
\PY{+w}{        }\PY{k+kt}{int}\PY{+w}{ }\PY{n}{serviceCount}\PY{+w}{ }\PY{o}{=}\PY{+w}{ }\PY{l+m+mi}{3}\PY{p}{;}
\PY{+w}{        }\PY{n}{CountDownLatch}\PY{+w}{ }\PY{n}{startupLatch}\PY{+w}{ }\PY{o}{=}\PY{+w}{ }\PY{k}{new}\PY{+w}{ }\PY{n}{CountDownLatch}\PY{p}{(}\PY{n}{serviceCount}\PY{p}{)}\PY{p}{;}

\PY{+w}{        }\PY{k}{for}\PY{+w}{ }\PY{p}{(}\PY{k+kt}{int}\PY{+w}{ }\PY{n}{i}\PY{+w}{ }\PY{o}{=}\PY{+w}{ }\PY{l+m+mi}{1}\PY{p}{;}\PY{+w}{ }\PY{n}{i}\PY{+w}{ }\PY{o}{\PYZlt{}}\PY{o}{=}\PY{+w}{ }\PY{n}{serviceCount}\PY{p}{;}\PY{+w}{ }\PY{n}{i}\PY{o}{+}\PY{o}{+}\PY{p}{)}\PY{+w}{ }\PY{p}{\PYZob{}}
\PY{+w}{            }\PY{k+kt}{int}\PY{+w}{ }\PY{n}{id}\PY{+w}{ }\PY{o}{=}\PY{+w}{ }\PY{n}{i}\PY{p}{;}
\PY{+w}{            }\PY{k}{new}\PY{+w}{ }\PY{n}{Thread}\PY{p}{(}\PY{p}{(}\PY{p}{)}\PY{+w}{ }\PY{o}{\PYZhy{}}\PY{o}{\PYZgt{}}\PY{+w}{ }\PY{p}{\PYZob{}}
\PY{+w}{                }\PY{n}{initializeService}\PY{p}{(}\PY{n}{id}\PY{p}{)}\PY{p}{;}
\PY{+w}{                }\PY{n}{System}\PY{p}{.}\PY{n+na}{out}\PY{p}{.}\PY{n+na}{println}\PY{p}{(}\PY{l+s}{\PYZdq{}}\PY{l+s}{service }\PY{l+s}{\PYZdq{}}\PY{+w}{ }\PY{o}{+}\PY{+w}{ }\PY{n}{id}\PY{+w}{ }\PY{o}{+}\PY{+w}{ }\PY{l+s}{\PYZdq{}}\PY{l+s}{ ready}\PY{l+s}{\PYZdq{}}\PY{p}{)}\PY{p}{;}
\PY{+w}{                }\PY{n}{startupLatch}\PY{p}{.}\PY{n+na}{countDown}\PY{p}{(}\PY{p}{)}\PY{p}{;}\PY{+w}{ }\PY{c+c1}{// one fewer thing to wait for}
\PY{+w}{            }\PY{p}{\PYZcb{}}\PY{p}{)}\PY{p}{.}\PY{n+na}{start}\PY{p}{(}\PY{p}{)}\PY{p}{;}
\PY{+w}{        }\PY{p}{\PYZcb{}}

\PY{+w}{        }\PY{n}{startupLatch}\PY{p}{.}\PY{n+na}{await}\PY{p}{(}\PY{p}{)}\PY{p}{;}\PY{+w}{ }\PY{c+c1}{// blocks until all 3 have called countDown()}
\PY{+w}{        }\PY{n}{System}\PY{p}{.}\PY{n+na}{out}\PY{p}{.}\PY{n+na}{println}\PY{p}{(}\PY{l+s}{\PYZdq{}}\PY{l+s}{all services ready, starting server}\PY{l+s}{\PYZdq{}}\PY{p}{)}\PY{p}{;}
\PY{+w}{    }\PY{p}{\PYZcb{}}

\PY{+w}{    }\PY{k+kd}{private}\PY{+w}{ }\PY{k+kd}{static}\PY{+w}{ }\PY{k+kt}{void}\PY{+w}{ }\PY{n+nf}{initializeService}\PY{p}{(}\PY{k+kt}{int}\PY{+w}{ }\PY{n}{id}\PY{p}{)}\PY{+w}{ }\PY{p}{\PYZob{}}
\PY{+w}{        }\PY{k}{try}\PY{+w}{ }\PY{p}{\PYZob{}}\PY{+w}{ }\PY{n}{Thread}\PY{p}{.}\PY{n+na}{sleep}\PY{p}{(}\PY{l+m+mi}{100L}\PY{+w}{ }\PY{o}{*}\PY{+w}{ }\PY{n}{id}\PY{p}{)}\PY{p}{;}\PY{+w}{ }\PY{p}{\PYZcb{}}\PY{+w}{ }\PY{k}{catch}\PY{+w}{ }\PY{p}{(}\PY{n}{InterruptedException}\PY{+w}{ }\PY{n}{e}\PY{p}{)}\PY{+w}{ }\PY{p}{\PYZob{}}\PY{+w}{ }\PY{n}{Thread}\PY{p}{.}\PY{n+na}{currentThread}\PY{p}{(}\PY{p}{)}\PY{p}{.}\PY{n+na}{interrupt}\PY{p}{(}\PY{p}{)}\PY{p}{;}\PY{+w}{ }\PY{p}{\PYZcb{}}
\PY{+w}{    }\PY{p}{\PYZcb{}}
\PY{p}{\PYZcb{}}
\end{Verbatim}
\end{codebox}

\section[CyclicBarrier]{CyclicBarrier}
\label{h:14:cyclicbarrier}
A \code{CyclicBarrier} makes a fixed number of threads (\textbf{parties}) all wait at a point until every one of them has arrived — then releases all of them at once. Unlike \code{CountDownLatch}, it’s \textbf{reusable}: once everyone passes through, it automatically resets for the next round. It also supports an optional \textbf{barrier action} — a \code{Runnable} that runs once, on one of the arriving threads, right before the barrier releases everyone.

Typical use: a simulation or algorithm that runs in rounds, where every worker thread must finish round \emph{N} before any of them can start round \emph{N+1}.

\begin{codebox}
\begin{Verbatim}[commandchars=\\\{\}]
\PY{k+kn}{import}\PY{+w}{ }\PY{n+nn}{java.util.concurrent.BrokenBarrierException}\PY{p}{;}
\PY{k+kn}{import}\PY{+w}{ }\PY{n+nn}{java.util.concurrent.CyclicBarrier}\PY{p}{;}

\PY{k+kd}{public}\PY{+w}{ }\PY{k+kd}{class} \PY{n+nc}{SimulationRounds}\PY{+w}{ }\PY{p}{\PYZob{}}
\PY{+w}{    }\PY{k+kd}{public}\PY{+w}{ }\PY{k+kd}{static}\PY{+w}{ }\PY{k+kt}{void}\PY{+w}{ }\PY{n+nf}{main}\PY{p}{(}\PY{n}{String}\PY{o}{[}\PY{o}{]}\PY{+w}{ }\PY{n}{args}\PY{p}{)}\PY{+w}{ }\PY{p}{\PYZob{}}
\PY{+w}{        }\PY{k+kt}{int}\PY{+w}{ }\PY{n}{workerCount}\PY{+w}{ }\PY{o}{=}\PY{+w}{ }\PY{l+m+mi}{4}\PY{p}{;}
\PY{+w}{        }\PY{n}{CyclicBarrier}\PY{+w}{ }\PY{n}{barrier}\PY{+w}{ }\PY{o}{=}\PY{+w}{ }\PY{k}{new}\PY{+w}{ }\PY{n}{CyclicBarrier}\PY{p}{(}\PY{n}{workerCount}\PY{p}{,}
\PY{+w}{                }\PY{p}{(}\PY{p}{)}\PY{+w}{ }\PY{o}{\PYZhy{}}\PY{o}{\PYZgt{}}\PY{+w}{ }\PY{n}{System}\PY{p}{.}\PY{n+na}{out}\PY{p}{.}\PY{n+na}{println}\PY{p}{(}\PY{l+s}{\PYZdq{}}\PY{l+s}{\PYZhy{}\PYZhy{}\PYZhy{} round complete, advancing \PYZhy{}\PYZhy{}\PYZhy{}}\PY{l+s}{\PYZdq{}}\PY{p}{)}\PY{p}{)}\PY{p}{;}\PY{+w}{ }\PY{c+c1}{// barrier action}

\PY{+w}{        }\PY{k}{for}\PY{+w}{ }\PY{p}{(}\PY{k+kt}{int}\PY{+w}{ }\PY{n}{i}\PY{+w}{ }\PY{o}{=}\PY{+w}{ }\PY{l+m+mi}{1}\PY{p}{;}\PY{+w}{ }\PY{n}{i}\PY{+w}{ }\PY{o}{\PYZlt{}}\PY{o}{=}\PY{+w}{ }\PY{n}{workerCount}\PY{p}{;}\PY{+w}{ }\PY{n}{i}\PY{o}{+}\PY{o}{+}\PY{p}{)}\PY{+w}{ }\PY{p}{\PYZob{}}
\PY{+w}{            }\PY{k+kt}{int}\PY{+w}{ }\PY{n}{id}\PY{+w}{ }\PY{o}{=}\PY{+w}{ }\PY{n}{i}\PY{p}{;}
\PY{+w}{            }\PY{k}{new}\PY{+w}{ }\PY{n}{Thread}\PY{p}{(}\PY{p}{(}\PY{p}{)}\PY{+w}{ }\PY{o}{\PYZhy{}}\PY{o}{\PYZgt{}}\PY{+w}{ }\PY{p}{\PYZob{}}
\PY{+w}{                }\PY{k}{for}\PY{+w}{ }\PY{p}{(}\PY{k+kt}{int}\PY{+w}{ }\PY{n}{round}\PY{+w}{ }\PY{o}{=}\PY{+w}{ }\PY{l+m+mi}{1}\PY{p}{;}\PY{+w}{ }\PY{n}{round}\PY{+w}{ }\PY{o}{\PYZlt{}}\PY{o}{=}\PY{+w}{ }\PY{l+m+mi}{3}\PY{p}{;}\PY{+w}{ }\PY{n}{round}\PY{o}{+}\PY{o}{+}\PY{p}{)}\PY{+w}{ }\PY{p}{\PYZob{}}
\PY{+w}{                    }\PY{n}{doRoundOfWork}\PY{p}{(}\PY{n}{id}\PY{p}{,}\PY{+w}{ }\PY{n}{round}\PY{p}{)}\PY{p}{;}
\PY{+w}{                    }\PY{k}{try}\PY{+w}{ }\PY{p}{\PYZob{}}
\PY{+w}{                        }\PY{n}{barrier}\PY{p}{.}\PY{n+na}{await}\PY{p}{(}\PY{p}{)}\PY{p}{;}\PY{+w}{ }\PY{c+c1}{// wait here until all 4 workers finish this round}
\PY{+w}{                    }\PY{p}{\PYZcb{}}\PY{+w}{ }\PY{k}{catch}\PY{+w}{ }\PY{p}{(}\PY{n}{InterruptedException}\PY{+w}{ }\PY{o}{|}\PY{+w}{ }\PY{n}{BrokenBarrierException}\PY{+w}{ }\PY{n}{e}\PY{p}{)}\PY{+w}{ }\PY{p}{\PYZob{}}
\PY{+w}{                        }\PY{n}{Thread}\PY{p}{.}\PY{n+na}{currentThread}\PY{p}{(}\PY{p}{)}\PY{p}{.}\PY{n+na}{interrupt}\PY{p}{(}\PY{p}{)}\PY{p}{;}
\PY{+w}{                        }\PY{k}{return}\PY{p}{;}
\PY{+w}{                    }\PY{p}{\PYZcb{}}
\PY{+w}{                }\PY{p}{\PYZcb{}}
\PY{+w}{            }\PY{p}{\PYZcb{}}\PY{p}{)}\PY{p}{.}\PY{n+na}{start}\PY{p}{(}\PY{p}{)}\PY{p}{;}
\PY{+w}{        }\PY{p}{\PYZcb{}}
\PY{+w}{    }\PY{p}{\PYZcb{}}

\PY{+w}{    }\PY{k+kd}{private}\PY{+w}{ }\PY{k+kd}{static}\PY{+w}{ }\PY{k+kt}{void}\PY{+w}{ }\PY{n+nf}{doRoundOfWork}\PY{p}{(}\PY{k+kt}{int}\PY{+w}{ }\PY{n}{workerId}\PY{p}{,}\PY{+w}{ }\PY{k+kt}{int}\PY{+w}{ }\PY{n}{round}\PY{p}{)}\PY{+w}{ }\PY{p}{\PYZob{}}
\PY{+w}{        }\PY{n}{System}\PY{p}{.}\PY{n+na}{out}\PY{p}{.}\PY{n+na}{println}\PY{p}{(}\PY{l+s}{\PYZdq{}}\PY{l+s}{worker }\PY{l+s}{\PYZdq{}}\PY{+w}{ }\PY{o}{+}\PY{+w}{ }\PY{n}{workerId}\PY{+w}{ }\PY{o}{+}\PY{+w}{ }\PY{l+s}{\PYZdq{}}\PY{l+s}{ doing round }\PY{l+s}{\PYZdq{}}\PY{+w}{ }\PY{o}{+}\PY{+w}{ }\PY{n}{round}\PY{p}{)}\PY{p}{;}
\PY{+w}{    }\PY{p}{\PYZcb{}}
\PY{p}{\PYZcb{}}
\end{Verbatim}
\end{codebox}

If a thread times out or is interrupted while waiting at the barrier, the barrier becomes \textbf{broken}: it throws \code{BrokenBarrierException} for every other thread currently waiting or that arrives later, and it stays broken until explicitly reset via \code{barrier.\allowbreak{}reset()}.

\section[Phaser]{Phaser}
\label{h:14:phaser}
A \code{Phaser} is a more flexible generalization of both \code{CountDownLatch} and \code{CyclicBarrier}. The key difference: the number of parties can change dynamically at runtime — threads can \code{register()} to join and \code{arriveAndDeregister()} to leave, even mid-run. Each round is called a \textbf{phase}, and phases increment automatically as parties arrive.

Typical use: anything where the set of participants isn’t fixed up front — e.g., a pool of workers that grows or shrinks between phases, or a multi-stage pipeline where different stages have different party counts.

\begin{codebox}
\begin{Verbatim}[commandchars=\\\{\}]
\PY{k+kn}{import}\PY{+w}{ }\PY{n+nn}{java.util.concurrent.Phaser}\PY{p}{;}

\PY{k+kd}{public}\PY{+w}{ }\PY{k+kd}{class} \PY{n+nc}{DynamicWorkers}\PY{+w}{ }\PY{p}{\PYZob{}}
\PY{+w}{    }\PY{k+kd}{public}\PY{+w}{ }\PY{k+kd}{static}\PY{+w}{ }\PY{k+kt}{void}\PY{+w}{ }\PY{n+nf}{main}\PY{p}{(}\PY{n}{String}\PY{o}{[}\PY{o}{]}\PY{+w}{ }\PY{n}{args}\PY{p}{)}\PY{+w}{ }\PY{p}{\PYZob{}}
\PY{+w}{        }\PY{n}{Phaser}\PY{+w}{ }\PY{n}{phaser}\PY{+w}{ }\PY{o}{=}\PY{+w}{ }\PY{k}{new}\PY{+w}{ }\PY{n}{Phaser}\PY{p}{(}\PY{l+m+mi}{1}\PY{p}{)}\PY{p}{;}\PY{+w}{ }\PY{c+c1}{// 1 party: the main thread registers itself first}

\PY{+w}{        }\PY{k}{for}\PY{+w}{ }\PY{p}{(}\PY{k+kt}{int}\PY{+w}{ }\PY{n}{i}\PY{+w}{ }\PY{o}{=}\PY{+w}{ }\PY{l+m+mi}{1}\PY{p}{;}\PY{+w}{ }\PY{n}{i}\PY{+w}{ }\PY{o}{\PYZlt{}}\PY{o}{=}\PY{+w}{ }\PY{l+m+mi}{3}\PY{p}{;}\PY{+w}{ }\PY{n}{i}\PY{o}{+}\PY{o}{+}\PY{p}{)}\PY{+w}{ }\PY{p}{\PYZob{}}
\PY{+w}{            }\PY{k+kt}{int}\PY{+w}{ }\PY{n}{workerId}\PY{+w}{ }\PY{o}{=}\PY{+w}{ }\PY{n}{i}\PY{p}{;}
\PY{+w}{            }\PY{n}{phaser}\PY{p}{.}\PY{n+na}{register}\PY{p}{(}\PY{p}{)}\PY{p}{;}\PY{+w}{ }\PY{c+c1}{// dynamically add a party for this new worker}
\PY{+w}{            }\PY{k}{new}\PY{+w}{ }\PY{n}{Thread}\PY{p}{(}\PY{p}{(}\PY{p}{)}\PY{+w}{ }\PY{o}{\PYZhy{}}\PY{o}{\PYZgt{}}\PY{+w}{ }\PY{p}{\PYZob{}}
\PY{+w}{                }\PY{k}{for}\PY{+w}{ }\PY{p}{(}\PY{k+kt}{int}\PY{+w}{ }\PY{n}{phase}\PY{+w}{ }\PY{o}{=}\PY{+w}{ }\PY{l+m+mi}{0}\PY{p}{;}\PY{+w}{ }\PY{n}{phase}\PY{+w}{ }\PY{o}{\PYZlt{}}\PY{+w}{ }\PY{l+m+mi}{2}\PY{p}{;}\PY{+w}{ }\PY{n}{phase}\PY{o}{+}\PY{o}{+}\PY{p}{)}\PY{+w}{ }\PY{p}{\PYZob{}}
\PY{+w}{                    }\PY{n}{System}\PY{p}{.}\PY{n+na}{out}\PY{p}{.}\PY{n+na}{println}\PY{p}{(}\PY{l+s}{\PYZdq{}}\PY{l+s}{worker }\PY{l+s}{\PYZdq{}}\PY{+w}{ }\PY{o}{+}\PY{+w}{ }\PY{n}{workerId}\PY{+w}{ }\PY{o}{+}\PY{+w}{ }\PY{l+s}{\PYZdq{}}\PY{l+s}{ working on phase }\PY{l+s}{\PYZdq{}}\PY{+w}{ }\PY{o}{+}\PY{+w}{ }\PY{n}{phase}\PY{p}{)}\PY{p}{;}
\PY{+w}{                    }\PY{n}{phaser}\PY{p}{.}\PY{n+na}{arriveAndAwaitAdvance}\PY{p}{(}\PY{p}{)}\PY{p}{;}\PY{+w}{ }\PY{c+c1}{// arrive, then wait for the phase to advance}
\PY{+w}{                }\PY{p}{\PYZcb{}}
\PY{+w}{                }\PY{n}{phaser}\PY{p}{.}\PY{n+na}{arriveAndDeregister}\PY{p}{(}\PY{p}{)}\PY{p}{;}\PY{+w}{ }\PY{c+c1}{// this worker is done, shrink the party count}
\PY{+w}{            }\PY{p}{\PYZcb{}}\PY{p}{)}\PY{p}{.}\PY{n+na}{start}\PY{p}{(}\PY{p}{)}\PY{p}{;}
\PY{+w}{        }\PY{p}{\PYZcb{}}

\PY{+w}{        }\PY{n}{phaser}\PY{p}{.}\PY{n+na}{arriveAndDeregister}\PY{p}{(}\PY{p}{)}\PY{p}{;}\PY{+w}{ }\PY{c+c1}{// main thread\PYZsq{}s initial registration is done too}
\PY{+w}{    }\PY{p}{\PYZcb{}}
\PY{p}{\PYZcb{}}
\end{Verbatim}
\end{codebox}

Phaser also supports \textbf{tiered} (hierarchical) structures for very large numbers of parties, and you can override \code{on\allowbreak{}Advance(\allowbreak{}phase,\allowbreak{}~\allowbreak{}registered\allowbreak{}Parties)} to control termination — returning \code{true} from it terminates the phaser after that phase.

\subsection[Latch vs Barrier vs Phaser]{Latch vs Barrier vs Phaser}
\label{h:14:latch-vs-barrier-vs-phaser}
{\tablefontsmall
\begin{xltabular}{\linewidth}{@{}L{0.471}L{1.386}L{1.046}L{1.098}@{}}
\toprule
\rowcolor{tablehdr}
\thd{} & \thd{\code{CountDownLatch}} & \thd{\code{CyclicBarrier}} & \thd{\code{Phaser}} \\
\midrule
\endhead
Reusable? & No — one-shot & Yes — resets automatically each round & Yes — advances through phases \\
Party count fixed? & Fixed at creation (the initial count) & Fixed at creation & \textbf{Dynamic} — register/deregister at runtime \\
Who signals? & Any thread can call \code{countDown()}, independent of \code{await()} & Only the parties themselves, via \code{await()} & Only the parties themselves, via \code{arrive*()} \\
Optional action on completion? & No & Yes, one barrier action & Yes, override \code{onAdvance()} \\
Typical use & Wait for N independent events (startup, N tasks done) & Wait for N \emph{same} threads to reach a sync point each round & Multi-phase work with a changing number of participants \\
\bottomrule
\end{xltabular}
}

\section[Semaphore]{Semaphore}
\label{h:14:semaphore}
A \code{Semaphore} guards access to a limited number of \textbf{permits} — think of it as a bouncer letting only N people into a room at a time. Unlike a lock, a semaphore has \textbf{no ownership}: any thread can \code{release()} a permit, not just the thread that \code{acquire()}d it. That makes it useful for resource pools where the “returning” thread isn’t necessarily the “borrowing” thread.

Typical use: limiting how many concurrent requests hit a downstream service, or guarding a fixed-size pool of expensive resources (e.g., database connections).

\begin{codebox}
\begin{Verbatim}[commandchars=\\\{\}]
\PY{k+kn}{import}\PY{+w}{ }\PY{n+nn}{java.util.concurrent.Semaphore}\PY{p}{;}

\PY{k+kd}{public}\PY{+w}{ }\PY{k+kd}{class} \PY{n+nc}{RateLimitedClient}\PY{+w}{ }\PY{p}{\PYZob{}}
\PY{+w}{    }\PY{k+kd}{private}\PY{+w}{ }\PY{k+kd}{final}\PY{+w}{ }\PY{n}{Semaphore}\PY{+w}{ }\PY{n}{permits}\PY{+w}{ }\PY{o}{=}\PY{+w}{ }\PY{k}{new}\PY{+w}{ }\PY{n}{Semaphore}\PY{p}{(}\PY{l+m+mi}{5}\PY{p}{)}\PY{p}{;}\PY{+w}{ }\PY{c+c1}{// at most 5 concurrent calls}

\PY{+w}{    }\PY{k+kd}{public}\PY{+w}{ }\PY{n}{String}\PY{+w}{ }\PY{n+nf}{callExternalApi}\PY{p}{(}\PY{n}{String}\PY{+w}{ }\PY{n}{request}\PY{p}{)}\PY{+w}{ }\PY{k+kd}{throws}\PY{+w}{ }\PY{n}{InterruptedException}\PY{+w}{ }\PY{p}{\PYZob{}}
\PY{+w}{        }\PY{n}{permits}\PY{p}{.}\PY{n+na}{acquire}\PY{p}{(}\PY{p}{)}\PY{p}{;}\PY{+w}{ }\PY{c+c1}{// blocks if 5 calls are already in flight}
\PY{+w}{        }\PY{k}{try}\PY{+w}{ }\PY{p}{\PYZob{}}
\PY{+w}{            }\PY{k}{return}\PY{+w}{ }\PY{n}{doHttpCall}\PY{p}{(}\PY{n}{request}\PY{p}{)}\PY{p}{;}
\PY{+w}{        }\PY{p}{\PYZcb{}}\PY{+w}{ }\PY{k}{finally}\PY{+w}{ }\PY{p}{\PYZob{}}
\PY{+w}{            }\PY{n}{permits}\PY{p}{.}\PY{n+na}{release}\PY{p}{(}\PY{p}{)}\PY{p}{;}\PY{+w}{ }\PY{c+c1}{// always release, even on exception — hence try/finally}
\PY{+w}{        }\PY{p}{\PYZcb{}}
\PY{+w}{    }\PY{p}{\PYZcb{}}

\PY{+w}{    }\PY{k+kd}{private}\PY{+w}{ }\PY{n}{String}\PY{+w}{ }\PY{n+nf}{doHttpCall}\PY{p}{(}\PY{n}{String}\PY{+w}{ }\PY{n}{request}\PY{p}{)}\PY{+w}{ }\PY{p}{\PYZob{}}
\PY{+w}{        }\PY{k}{return}\PY{+w}{ }\PY{l+s}{\PYZdq{}}\PY{l+s}{response for }\PY{l+s}{\PYZdq{}}\PY{+w}{ }\PY{o}{+}\PY{+w}{ }\PY{n}{request}\PY{p}{;}
\PY{+w}{    }\PY{p}{\PYZcb{}}
\PY{p}{\PYZcb{}}
\end{Verbatim}
\end{codebox}

\code{tryAcquire()} (optionally with a timeout) lets you back off instead of blocking indefinitely — useful for “fail fast if the pool is currently exhausted” behavior. A \code{Semaphore} can also be constructed as \textbf{fair} (\code{new~\allowbreak{}Semaphore(\allowbreak{}5,\allowbreak{}~\allowbreak{}true)}), which guarantees FIFO ordering among waiting threads at some throughput cost, versus the default unfair mode which allows barging for better overall throughput.

A \textbf{binary semaphore} (1 permit) looks like a mutex but isn’t one: it has no notion of “owner,” so thread A can acquire it and thread B can release it — a flexible but also easy-to-misuse property.

\section[Exchanger]{Exchanger}
\label{h:14:exchanger}
An \code{Exchanger\textless{}\allowbreak{}V\textgreater{}} is the simplest of these tools: it lets \textbf{exactly two threads} meet at a rendezvous point and atomically swap an object. Each thread calls \code{exchange(\allowbreak{}myObject)}, which blocks until the other thread also calls \code{exchange()}, and each returns the \emph{other} thread’s object.

Typical use: a two-stage pipeline where one thread fills a buffer while another drains/processes a different buffer, and they periodically swap so nobody has to wait for a full read+write cycle on a shared buffer.

\begin{codebox}
\begin{Verbatim}[commandchars=\\\{\}]
\PY{k+kn}{import}\PY{+w}{ }\PY{n+nn}{java.util.concurrent.Exchanger}\PY{p}{;}
\PY{k+kn}{import}\PY{+w}{ }\PY{n+nn}{java.util.ArrayList}\PY{p}{;}
\PY{k+kn}{import}\PY{+w}{ }\PY{n+nn}{java.util.List}\PY{p}{;}

\PY{k+kd}{public}\PY{+w}{ }\PY{k+kd}{class} \PY{n+nc}{BufferSwap}\PY{+w}{ }\PY{p}{\PYZob{}}
\PY{+w}{    }\PY{k+kd}{public}\PY{+w}{ }\PY{k+kd}{static}\PY{+w}{ }\PY{k+kt}{void}\PY{+w}{ }\PY{n+nf}{main}\PY{p}{(}\PY{n}{String}\PY{o}{[}\PY{o}{]}\PY{+w}{ }\PY{n}{args}\PY{p}{)}\PY{+w}{ }\PY{p}{\PYZob{}}
\PY{+w}{        }\PY{n}{Exchanger}\PY{o}{\PYZlt{}}\PY{n}{List}\PY{o}{\PYZlt{}}\PY{n}{String}\PY{o}{\PYZgt{}}\PY{o}{\PYZgt{}}\PY{+w}{ }\PY{n}{exchanger}\PY{+w}{ }\PY{o}{=}\PY{+w}{ }\PY{k}{new}\PY{+w}{ }\PY{n}{Exchanger}\PY{o}{\PYZlt{}}\PY{o}{\PYZgt{}}\PY{p}{(}\PY{p}{)}\PY{p}{;}

\PY{+w}{        }\PY{n}{Thread}\PY{+w}{ }\PY{n}{producer}\PY{+w}{ }\PY{o}{=}\PY{+w}{ }\PY{k}{new}\PY{+w}{ }\PY{n}{Thread}\PY{p}{(}\PY{p}{(}\PY{p}{)}\PY{+w}{ }\PY{o}{\PYZhy{}}\PY{o}{\PYZgt{}}\PY{+w}{ }\PY{p}{\PYZob{}}
\PY{+w}{            }\PY{n}{List}\PY{o}{\PYZlt{}}\PY{n}{String}\PY{o}{\PYZgt{}}\PY{+w}{ }\PY{n}{buffer}\PY{+w}{ }\PY{o}{=}\PY{+w}{ }\PY{k}{new}\PY{+w}{ }\PY{n}{ArrayList}\PY{o}{\PYZlt{}}\PY{o}{\PYZgt{}}\PY{p}{(}\PY{p}{)}\PY{p}{;}
\PY{+w}{            }\PY{k}{try}\PY{+w}{ }\PY{p}{\PYZob{}}
\PY{+w}{                }\PY{k}{while}\PY{+w}{ }\PY{p}{(}\PY{k+kc}{true}\PY{p}{)}\PY{+w}{ }\PY{p}{\PYZob{}}
\PY{+w}{                    }\PY{n}{buffer}\PY{p}{.}\PY{n+na}{add}\PY{p}{(}\PY{l+s}{\PYZdq{}}\PY{l+s}{item\PYZhy{}}\PY{l+s}{\PYZdq{}}\PY{+w}{ }\PY{o}{+}\PY{+w}{ }\PY{n}{System}\PY{p}{.}\PY{n+na}{nanoTime}\PY{p}{(}\PY{p}{)}\PY{p}{)}\PY{p}{;}
\PY{+w}{                    }\PY{k}{if}\PY{+w}{ }\PY{p}{(}\PY{n}{buffer}\PY{p}{.}\PY{n+na}{size}\PY{p}{(}\PY{p}{)}\PY{+w}{ }\PY{o}{=}\PY{o}{=}\PY{+w}{ }\PY{l+m+mi}{10}\PY{p}{)}\PY{+w}{ }\PY{p}{\PYZob{}}
\PY{+w}{                        }\PY{n}{buffer}\PY{+w}{ }\PY{o}{=}\PY{+w}{ }\PY{n}{exchanger}\PY{p}{.}\PY{n+na}{exchange}\PY{p}{(}\PY{n}{buffer}\PY{p}{)}\PY{p}{;}\PY{+w}{ }\PY{c+c1}{// hand off full buffer, get an empty one back}
\PY{+w}{                    }\PY{p}{\PYZcb{}}
\PY{+w}{                }\PY{p}{\PYZcb{}}
\PY{+w}{            }\PY{p}{\PYZcb{}}\PY{+w}{ }\PY{k}{catch}\PY{+w}{ }\PY{p}{(}\PY{n}{InterruptedException}\PY{+w}{ }\PY{n}{e}\PY{p}{)}\PY{+w}{ }\PY{p}{\PYZob{}}
\PY{+w}{                }\PY{n}{Thread}\PY{p}{.}\PY{n+na}{currentThread}\PY{p}{(}\PY{p}{)}\PY{p}{.}\PY{n+na}{interrupt}\PY{p}{(}\PY{p}{)}\PY{p}{;}
\PY{+w}{            }\PY{p}{\PYZcb{}}
\PY{+w}{        }\PY{p}{\PYZcb{}}\PY{p}{)}\PY{p}{;}

\PY{+w}{        }\PY{n}{Thread}\PY{+w}{ }\PY{n}{consumer}\PY{+w}{ }\PY{o}{=}\PY{+w}{ }\PY{k}{new}\PY{+w}{ }\PY{n}{Thread}\PY{p}{(}\PY{p}{(}\PY{p}{)}\PY{+w}{ }\PY{o}{\PYZhy{}}\PY{o}{\PYZgt{}}\PY{+w}{ }\PY{p}{\PYZob{}}
\PY{+w}{            }\PY{n}{List}\PY{o}{\PYZlt{}}\PY{n}{String}\PY{o}{\PYZgt{}}\PY{+w}{ }\PY{n}{buffer}\PY{+w}{ }\PY{o}{=}\PY{+w}{ }\PY{k}{new}\PY{+w}{ }\PY{n}{ArrayList}\PY{o}{\PYZlt{}}\PY{o}{\PYZgt{}}\PY{p}{(}\PY{p}{)}\PY{p}{;}
\PY{+w}{            }\PY{k}{try}\PY{+w}{ }\PY{p}{\PYZob{}}
\PY{+w}{                }\PY{k}{while}\PY{+w}{ }\PY{p}{(}\PY{k+kc}{true}\PY{p}{)}\PY{+w}{ }\PY{p}{\PYZob{}}
\PY{+w}{                    }\PY{n}{buffer}\PY{+w}{ }\PY{o}{=}\PY{+w}{ }\PY{n}{exchanger}\PY{p}{.}\PY{n+na}{exchange}\PY{p}{(}\PY{n}{buffer}\PY{p}{)}\PY{p}{;}\PY{+w}{ }\PY{c+c1}{// hand off empty buffer, get the full one back}
\PY{+w}{                    }\PY{n}{System}\PY{p}{.}\PY{n+na}{out}\PY{p}{.}\PY{n+na}{println}\PY{p}{(}\PY{l+s}{\PYZdq{}}\PY{l+s}{processing }\PY{l+s}{\PYZdq{}}\PY{+w}{ }\PY{o}{+}\PY{+w}{ }\PY{n}{buffer}\PY{p}{.}\PY{n+na}{size}\PY{p}{(}\PY{p}{)}\PY{+w}{ }\PY{o}{+}\PY{+w}{ }\PY{l+s}{\PYZdq{}}\PY{l+s}{ items}\PY{l+s}{\PYZdq{}}\PY{p}{)}\PY{p}{;}
\PY{+w}{                    }\PY{n}{buffer}\PY{p}{.}\PY{n+na}{clear}\PY{p}{(}\PY{p}{)}\PY{p}{;}
\PY{+w}{                }\PY{p}{\PYZcb{}}
\PY{+w}{            }\PY{p}{\PYZcb{}}\PY{+w}{ }\PY{k}{catch}\PY{+w}{ }\PY{p}{(}\PY{n}{InterruptedException}\PY{+w}{ }\PY{n}{e}\PY{p}{)}\PY{+w}{ }\PY{p}{\PYZob{}}
\PY{+w}{                }\PY{n}{Thread}\PY{p}{.}\PY{n+na}{currentThread}\PY{p}{(}\PY{p}{)}\PY{p}{.}\PY{n+na}{interrupt}\PY{p}{(}\PY{p}{)}\PY{p}{;}
\PY{+w}{            }\PY{p}{\PYZcb{}}
\PY{+w}{        }\PY{p}{\PYZcb{}}\PY{p}{)}\PY{p}{;}

\PY{+w}{        }\PY{n}{producer}\PY{p}{.}\PY{n+na}{setDaemon}\PY{p}{(}\PY{k+kc}{true}\PY{p}{)}\PY{p}{;}
\PY{+w}{        }\PY{n}{consumer}\PY{p}{.}\PY{n+na}{setDaemon}\PY{p}{(}\PY{k+kc}{true}\PY{p}{)}\PY{p}{;}
\PY{+w}{        }\PY{n}{producer}\PY{p}{.}\PY{n+na}{start}\PY{p}{(}\PY{p}{)}\PY{p}{;}
\PY{+w}{        }\PY{n}{consumer}\PY{p}{.}\PY{n+na}{start}\PY{p}{(}\PY{p}{)}\PY{p}{;}
\PY{+w}{    }\PY{p}{\PYZcb{}}
\PY{p}{\PYZcb{}}
\end{Verbatim}
\end{codebox}

\code{Exchanger} is intentionally narrow — exactly two parties, one swap at a time. For more than two participants, reach for a \code{CyclicBarrier}, \code{Phaser}, or a \code{BlockingQueue} instead.

\section[Common Code-Review Interview Pitfalls]{Common Code-Review Interview Pitfalls}
\label{h:14:common-code-review-interview-pitfalls}
\begin{enumerate}
\item 
\textbf{Reaching for \code{Executors.\allowbreak{}new\allowbreak{}Fixed\allowbreak{}Thread\allowbreak{}Pool}/\code{newCachedThreadPool} without checking their queue/thread limits.} Why it matters: a fixed pool’s unbounded queue can grow without limit under sustained load, and a cached pool can spawn unbounded threads — both are a slow-motion \code{OutOfMemoryError} waiting to happen.

\begin{codebox}
\begin{Verbatim}[commandchars=\\\{\}]
\PY{c+c1}{// Before}
\PY{n}{ExecutorService}\PY{+w}{ }\PY{n}{pool}\PY{+w}{ }\PY{o}{=}\PY{+w}{ }\PY{n}{Executors}\PY{p}{.}\PY{n+na}{newFixedThreadPool}\PY{p}{(}\PY{l+m+mi}{10}\PY{p}{)}\PY{p}{;}\PY{+w}{ }\PY{c+c1}{// unbounded queue behind the scenes}
\PY{c+c1}{// After}
\PY{n}{ExecutorService}\PY{+w}{ }\PY{n}{pool}\PY{+w}{ }\PY{o}{=}\PY{+w}{ }\PY{k}{new}\PY{+w}{ }\PY{n}{ThreadPoolExecutor}\PY{p}{(}\PY{l+m+mi}{10}\PY{p}{,}\PY{+w}{ }\PY{l+m+mi}{10}\PY{p}{,}\PY{+w}{ }\PY{l+m+mi}{0}\PY{n}{L}\PY{p}{,}\PY{+w}{ }\PY{n}{TimeUnit}\PY{p}{.}\PY{n+na}{MILLISECONDS}\PY{p}{,}
\PY{+w}{        }\PY{k}{new}\PY{+w}{ }\PY{n}{ArrayBlockingQueue}\PY{o}{\PYZlt{}}\PY{o}{\PYZgt{}}\PY{p}{(}\PY{l+m+mi}{500}\PY{p}{)}\PY{p}{,}\PY{+w}{ }\PY{k}{new}\PY{+w}{ }\PY{n}{ThreadPoolExecutor}\PY{p}{.}\PY{n+na}{CallerRunsPolicy}\PY{p}{(}\PY{p}{)}\PY{p}{)}\PY{p}{;}
\end{Verbatim}
\end{codebox}

\item 
\textbf{Calling \code{submit()} and never checking the returned \code{Future}.} Why it matters: exceptions thrown inside the task are captured in the \code{Future} and silently disappear unless you call \code{get()} (or otherwise inspect it) — a failing task looks like a succeeding one.

\begin{codebox}
\begin{Verbatim}[commandchars=\\\{\}]
\PY{c+c1}{// Before}
\PY{n}{pool}\PY{p}{.}\PY{n+na}{submit}\PY{p}{(}\PY{p}{(}\PY{p}{)}\PY{+w}{ }\PY{o}{\PYZhy{}}\PY{o}{\PYZgt{}}\PY{+w}{ }\PY{n}{riskyWork}\PY{p}{(}\PY{p}{)}\PY{p}{)}\PY{p}{;}\PY{+w}{ }\PY{c+c1}{// exception vanishes if thrown}
\PY{c+c1}{// After}
\PY{n}{Future}\PY{o}{\PYZlt{}}\PY{o}{?}\PY{o}{\PYZgt{}}\PY{+w}{ }\PY{n}{f}\PY{+w}{ }\PY{o}{=}\PY{+w}{ }\PY{n}{pool}\PY{p}{.}\PY{n+na}{submit}\PY{p}{(}\PY{p}{(}\PY{p}{)}\PY{+w}{ }\PY{o}{\PYZhy{}}\PY{o}{\PYZgt{}}\PY{+w}{ }\PY{n}{riskyWork}\PY{p}{(}\PY{p}{)}\PY{p}{)}\PY{p}{;}
\PY{n}{f}\PY{p}{.}\PY{n+na}{get}\PY{p}{(}\PY{p}{)}\PY{p}{;}\PY{+w}{ }\PY{c+c1}{// or handle it in a wrapper/logging layer}
\end{Verbatim}
\end{codebox}

\item 
\textbf{Never shutting down an \code{ExecutorService}.} Why it matters: a pool with non-daemon threads keeps the JVM alive, and a pool that’s simply forgotten leaks threads for the life of the application.

\begin{codebox}
\begin{Verbatim}[commandchars=\\\{\}]
\PY{c+c1}{// Before}
\PY{n}{ExecutorService}\PY{+w}{ }\PY{n}{pool}\PY{+w}{ }\PY{o}{=}\PY{+w}{ }\PY{n}{Executors}\PY{p}{.}\PY{n+na}{newFixedThreadPool}\PY{p}{(}\PY{l+m+mi}{4}\PY{p}{)}\PY{p}{;}
\PY{c+c1}{// ... used once, never shut down}
\PY{c+c1}{// After}
\PY{k}{try}\PY{+w}{ }\PY{p}{(}\PY{n}{ExecutorService}\PY{+w}{ }\PY{n}{pool}\PY{+w}{ }\PY{o}{=}\PY{+w}{ }\PY{n}{Executors}\PY{p}{.}\PY{n+na}{newFixedThreadPool}\PY{p}{(}\PY{l+m+mi}{4}\PY{p}{)}\PY{p}{)}\PY{+w}{ }\PY{p}{\PYZob{}}
\PY{+w}{    }\PY{n}{pool}\PY{p}{.}\PY{n+na}{submit}\PY{p}{(}\PY{k}{this}\PY{p}{:}\PY{p}{:}\PY{n}{work}\PY{p}{)}\PY{p}{;}
\PY{p}{\PYZcb{}}\PY{+w}{ }\PY{c+c1}{// auto shutdown via AutoCloseable (Java 19+)}
\end{Verbatim}
\end{codebox}

\item 
\textbf{Using \code{shutdownNow()} without handling the tasks it returns.} Why it matters: \code{shutdownNow()} cancels running tasks and hands back the still-queued ones — dropping that list silently discards work that may need to be retried or logged.

\begin{codebox}
\begin{Verbatim}[commandchars=\\\{\}]
\PY{c+c1}{// Before}
\PY{n}{pool}\PY{p}{.}\PY{n+na}{shutdownNow}\PY{p}{(}\PY{p}{)}\PY{p}{;}\PY{+w}{ }\PY{c+c1}{// returned list ignored, queued work is just gone}
\PY{c+c1}{// After}
\PY{n}{List}\PY{o}{\PYZlt{}}\PY{n}{Runnable}\PY{o}{\PYZgt{}}\PY{+w}{ }\PY{n}{abandoned}\PY{+w}{ }\PY{o}{=}\PY{+w}{ }\PY{n}{pool}\PY{p}{.}\PY{n+na}{shutdownNow}\PY{p}{(}\PY{p}{)}\PY{p}{;}
\PY{n}{abandoned}\PY{p}{.}\PY{n+na}{forEach}\PY{p}{(}\PY{n}{task}\PY{+w}{ }\PY{o}{\PYZhy{}}\PY{o}{\PYZgt{}}\PY{+w}{ }\PY{n}{log}\PY{p}{.}\PY{n+na}{warn}\PY{p}{(}\PY{l+s}{\PYZdq{}}\PY{l+s}{dropped task: \PYZob{}\PYZcb{}}\PY{l+s}{\PYZdq{}}\PY{p}{,}\PY{+w}{ }\PY{n}{task}\PY{p}{)}\PY{p}{)}\PY{p}{;}
\end{Verbatim}
\end{codebox}

\item 
\textbf{Sizing a thread pool for IO-bound work the same way as CPU-bound work.} Why it matters: a pool sized to \code{cores} for work that mostly waits on network/disk will badly under-utilize the machine — most threads are blocked, not computing.

\begin{codebox}
\begin{Verbatim}[commandchars=\\\{\}]
\PY{c+c1}{// Before: 8 threads for an HTTP\PYZhy{}call\PYZhy{}heavy workload on an 8\PYZhy{}core box}
\PY{n}{ExecutorService}\PY{+w}{ }\PY{n}{pool}\PY{+w}{ }\PY{o}{=}\PY{+w}{ }\PY{n}{Executors}\PY{p}{.}\PY{n+na}{newFixedThreadPool}\PY{p}{(}\PY{l+m+mi}{8}\PY{p}{)}\PY{p}{;}
\PY{c+c1}{// After: size based on wait/compute ratio, e.g. cores * (1 + waitTime/computeTime)}
\PY{n}{ExecutorService}\PY{+w}{ }\PY{n}{pool}\PY{+w}{ }\PY{o}{=}\PY{+w}{ }\PY{n}{Executors}\PY{p}{.}\PY{n+na}{newFixedThreadPool}\PY{p}{(}\PY{l+m+mi}{64}\PY{p}{)}\PY{p}{;}
\end{Verbatim}
\end{codebox}

\item 
\textbf{Choosing a \code{ForkJoinTask} threshold that’s too small.} Why it matters: forking tiny units of work means the bookkeeping overhead (task objects, scheduling, stealing) outweighs the actual computation, making it slower than sequential code.

\begin{codebox}
\begin{Verbatim}[commandchars=\\\{\}]
\PY{c+c1}{// Before}
\PY{k+kd}{private}\PY{+w}{ }\PY{k+kd}{static}\PY{+w}{ }\PY{k+kd}{final}\PY{+w}{ }\PY{k+kt}{int}\PY{+w}{ }\PY{n}{THRESHOLD}\PY{+w}{ }\PY{o}{=}\PY{+w}{ }\PY{l+m+mi}{2}\PY{p}{;}\PY{+w}{ }\PY{c+c1}{// forks constantly, dominated by overhead}
\PY{c+c1}{// After}
\PY{k+kd}{private}\PY{+w}{ }\PY{k+kd}{static}\PY{+w}{ }\PY{k+kd}{final}\PY{+w}{ }\PY{k+kt}{int}\PY{+w}{ }\PY{n}{THRESHOLD}\PY{+w}{ }\PY{o}{=}\PY{+w}{ }\PY{l+m+mi}{10\PYZus{}000}\PY{p}{;}\PY{+w}{ }\PY{c+c1}{// benchmark and tune from here}
\end{Verbatim}
\end{codebox}

\item 
\textbf{Blocking inside a parallel stream or common-pool task without a \code{ManagedBlocker}.} Why it matters: the common pool is shared JVM-wide; a blocking call ties up one of its few worker threads and can starve unrelated parallel streams elsewhere in the app.

\begin{codebox}
\begin{Verbatim}[commandchars=\\\{\}]
\PY{c+c1}{// Before}
\PY{n}{list}\PY{p}{.}\PY{n+na}{parallelStream}\PY{p}{(}\PY{p}{)}\PY{p}{.}\PY{n+na}{forEach}\PY{p}{(}\PY{n}{item}\PY{+w}{ }\PY{o}{\PYZhy{}}\PY{o}{\PYZgt{}}\PY{+w}{ }\PY{n}{blockingHttpCall}\PY{p}{(}\PY{n}{item}\PY{p}{)}\PY{p}{)}\PY{p}{;}\PY{+w}{ }\PY{c+c1}{// starves the common pool}
\PY{c+c1}{// After: run blocking work on a dedicated executor, not the common pool}
\PY{n}{list}\PY{p}{.}\PY{n+na}{forEach}\PY{p}{(}\PY{n}{item}\PY{+w}{ }\PY{o}{\PYZhy{}}\PY{o}{\PYZgt{}}\PY{+w}{ }\PY{n}{ioExecutor}\PY{p}{.}\PY{n+na}{submit}\PY{p}{(}\PY{p}{(}\PY{p}{)}\PY{+w}{ }\PY{o}{\PYZhy{}}\PY{o}{\PYZgt{}}\PY{+w}{ }\PY{n}{blockingHttpCall}\PY{p}{(}\PY{n}{item}\PY{p}{)}\PY{p}{)}\PY{p}{)}\PY{p}{;}
\end{Verbatim}
\end{codebox}

\item 
\textbf{Calling \code{.\allowbreak{}join()}/\code{.\allowbreak{}get()} inside a \code{CompletableFuture} callback that shares a bounded pool with its dependency.} Why it matters: it can starve or deadlock the pool — all worker threads end up blocked waiting on each other with none left to make progress.

\begin{codebox}
\begin{Verbatim}[commandchars=\\\{\}]
\PY{c+c1}{// Before}
\PY{n}{CompletableFuture}\PY{p}{.}\PY{n+na}{supplyAsync}\PY{p}{(}\PY{p}{(}\PY{p}{)}\PY{+w}{ }\PY{o}{\PYZhy{}}\PY{o}{\PYZgt{}}\PY{+w}{ }\PY{n}{inner}\PY{p}{(}\PY{p}{)}\PY{p}{.}\PY{n+na}{join}\PY{p}{(}\PY{p}{)}\PY{p}{)}\PY{p}{;}\PY{+w}{ }\PY{c+c1}{// blocks a shared\PYZhy{}pool worker thread}
\PY{c+c1}{// After}
\PY{n}{CompletableFuture}\PY{p}{.}\PY{n+na}{supplyAsync}\PY{p}{(}\PY{k}{this}\PY{p}{:}\PY{p}{:}\PY{n}{prepare}\PY{p}{)}\PY{p}{.}\PY{n+na}{thenCompose}\PY{p}{(}\PY{n}{x}\PY{+w}{ }\PY{o}{\PYZhy{}}\PY{o}{\PYZgt{}}\PY{+w}{ }\PY{n}{innerAsync}\PY{p}{(}\PY{n}{x}\PY{p}{)}\PY{p}{)}\PY{p}{;}\PY{+w}{ }\PY{c+c1}{// no blocking}
\end{Verbatim}
\end{codebox}

\item 
\textbf{Not passing an explicit \code{Executor} to \code{CompletableFuture}’s async methods.} Why it matters: it silently defaults to the shared common pool, mixing unrelated workloads together and making capacity and failure modes unpredictable.

\begin{codebox}
\begin{Verbatim}[commandchars=\\\{\}]
\PY{c+c1}{// Before}
\PY{n}{CompletableFuture}\PY{p}{.}\PY{n+na}{supplyAsync}\PY{p}{(}\PY{k}{this}\PY{p}{:}\PY{p}{:}\PY{n}{loadData}\PY{p}{)}\PY{p}{;}\PY{+w}{ }\PY{c+c1}{// runs on ForkJoinPool.commonPool()}
\PY{c+c1}{// After}
\PY{n}{CompletableFuture}\PY{p}{.}\PY{n+na}{supplyAsync}\PY{p}{(}\PY{k}{this}\PY{p}{:}\PY{p}{:}\PY{n}{loadData}\PY{p}{,}\PY{+w}{ }\PY{n}{dedicatedExecutor}\PY{p}{)}\PY{p}{;}
\end{Verbatim}
\end{codebox}

\item 
\textbf{Using an unbounded \code{LinkedBlockingQueue} for a producer-consumer pipeline.} Why it matters: \code{new~\allowbreak{}Linked\allowbreak{}Blocking\allowbreak{}Queue\textless{}\textgreater{}()} defaults to \code{Integer.\allowbreak{}MAX\_\allowbreak{}VALUE} capacity — there is no backpressure, so a slow consumer just lets memory grow until the process dies.

\begin{codebox}
\begin{Verbatim}[commandchars=\\\{\}]
\PY{c+c1}{// Before}
\PY{n}{BlockingQueue}\PY{o}{\PYZlt{}}\PY{n}{Job}\PY{o}{\PYZgt{}}\PY{+w}{ }\PY{n}{queue}\PY{+w}{ }\PY{o}{=}\PY{+w}{ }\PY{k}{new}\PY{+w}{ }\PY{n}{LinkedBlockingQueue}\PY{o}{\PYZlt{}}\PY{o}{\PYZgt{}}\PY{p}{(}\PY{p}{)}\PY{p}{;}\PY{+w}{ }\PY{c+c1}{// effectively unbounded}
\PY{c+c1}{// After}
\PY{n}{BlockingQueue}\PY{o}{\PYZlt{}}\PY{n}{Job}\PY{o}{\PYZgt{}}\PY{+w}{ }\PY{n}{queue}\PY{+w}{ }\PY{o}{=}\PY{+w}{ }\PY{k}{new}\PY{+w}{ }\PY{n}{LinkedBlockingQueue}\PY{o}{\PYZlt{}}\PY{o}{\PYZgt{}}\PY{p}{(}\PY{l+m+mi}{1000}\PY{p}{)}\PY{p}{;}\PY{+w}{ }\PY{c+c1}{// bounded = backpressure}
\end{Verbatim}
\end{codebox}

\item 
\textbf{Using \code{add()}/\code{remove()} on a queue and expecting blocking behavior.} Why it matters: \code{add()} throws \code{IllegalStateException} immediately if the queue is full (it doesn’t wait), and \code{remove()} throws \code{NoSuchElementException} if empty — surprising if you meant \code{put()}/\code{take()} or \code{offer()}/\code{poll()} with a timeout.

\begin{codebox}
\begin{Verbatim}[commandchars=\\\{\}]
\PY{c+c1}{// Before}
\PY{n}{queue}\PY{p}{.}\PY{n+na}{add}\PY{p}{(}\PY{n}{job}\PY{p}{)}\PY{p}{;}\PY{+w}{ }\PY{c+c1}{// throws immediately when full, no waiting}
\PY{c+c1}{// After}
\PY{n}{queue}\PY{p}{.}\PY{n+na}{put}\PY{p}{(}\PY{n}{job}\PY{p}{)}\PY{p}{;}\PY{+w}{ }\PY{c+c1}{// blocks until space is available}
\end{Verbatim}
\end{codebox}

\item 
\textbf{Calling back into the same \code{ConcurrentHashMap} from inside \code{compute}/\code{computeIfAbsent}/\code{merge}.} Why it matters: the remapping function runs while the bin’s lock is held; recursively touching the same map (especially the same key) can deadlock or throw \code{IllegalStateException}.

\begin{codebox}
\begin{Verbatim}[commandchars=\\\{\}]
\PY{c+c1}{// Before}
\PY{n}{map}\PY{p}{.}\PY{n+na}{computeIfAbsent}\PY{p}{(}\PY{n}{key}\PY{p}{,}\PY{+w}{ }\PY{n}{k}\PY{+w}{ }\PY{o}{\PYZhy{}}\PY{o}{\PYZgt{}}\PY{+w}{ }\PY{n}{map}\PY{p}{.}\PY{n+na}{computeIfAbsent}\PY{p}{(}\PY{n}{otherKey}\PY{p}{,}\PY{+w}{ }\PY{n}{k2}\PY{+w}{ }\PY{o}{\PYZhy{}}\PY{o}{\PYZgt{}}\PY{+w}{ }\PY{l+m+mi}{1}\PY{p}{)}\PY{p}{)}\PY{p}{;}\PY{+w}{ }\PY{c+c1}{// risky recursion}
\PY{c+c1}{// After}
\PY{k+kt}{int}\PY{+w}{ }\PY{n}{otherValue}\PY{+w}{ }\PY{o}{=}\PY{+w}{ }\PY{n}{map}\PY{p}{.}\PY{n+na}{computeIfAbsent}\PY{p}{(}\PY{n}{otherKey}\PY{p}{,}\PY{+w}{ }\PY{n}{k2}\PY{+w}{ }\PY{o}{\PYZhy{}}\PY{o}{\PYZgt{}}\PY{+w}{ }\PY{l+m+mi}{1}\PY{p}{)}\PY{p}{;}
\PY{n}{map}\PY{p}{.}\PY{n+na}{computeIfAbsent}\PY{p}{(}\PY{n}{key}\PY{p}{,}\PY{+w}{ }\PY{n}{k}\PY{+w}{ }\PY{o}{\PYZhy{}}\PY{o}{\PYZgt{}}\PY{+w}{ }\PY{n}{otherValue}\PY{p}{)}\PY{p}{;}
\end{Verbatim}
\end{codebox}

\item 
\textbf{Treating \code{Concurrent\allowbreak{}Hash\allowbreak{}Map.\allowbreak{}size()} as an exact value for correctness-critical logic.} Why it matters: under concurrent modification, \code{size()}/\code{mappingCount()} are best-effort estimates from striped counters, not a locked, precise snapshot.

\begin{codebox}
\begin{Verbatim}[commandchars=\\\{\}]
\PY{c+c1}{// Before}
\PY{k}{if}\PY{+w}{ }\PY{p}{(}\PY{n}{cache}\PY{p}{.}\PY{n+na}{size}\PY{p}{(}\PY{p}{)}\PY{+w}{ }\PY{o}{=}\PY{o}{=}\PY{+w}{ }\PY{l+m+mi}{0}\PY{p}{)}\PY{+w}{ }\PY{p}{\PYZob{}}\PY{+w}{ }\PY{c+cm}{/* assume nothing cached yet */}\PY{+w}{ }\PY{p}{\PYZcb{}}
\PY{c+c1}{// After: don\PYZsq{}t rely on size for correctness; check isEmpty()/containsKey() for the specific}
\PY{c+c1}{// guarantee you actually need, and design around the map\PYZsq{}s atomic per\PYZhy{}key operations instead.}
\end{Verbatim}
\end{codebox}

\item 
\textbf{Expecting a \code{CountDownLatch} to reset for a second round.} Why it matters: \code{CountDownLatch} is one-shot by design — once it hits zero it’s done forever; reusing the same instance for a second wave of work silently does nothing (it’s already at zero, so \code{await()} returns instantly without waiting for anyone).

\begin{codebox}
\begin{Verbatim}[commandchars=\\\{\}]
\PY{c+c1}{// Before}
\PY{n}{CountDownLatch}\PY{+w}{ }\PY{n}{latch}\PY{+w}{ }\PY{o}{=}\PY{+w}{ }\PY{k}{new}\PY{+w}{ }\PY{n}{CountDownLatch}\PY{p}{(}\PY{l+m+mi}{3}\PY{p}{)}\PY{p}{;}
\PY{c+c1}{// ... used for round 1, then reused for round 2 — await() no longer actually waits}
\PY{c+c1}{// After: use a new CountDownLatch per round, or use CyclicBarrier/Phaser which are reusable}
\PY{n}{CyclicBarrier}\PY{+w}{ }\PY{n}{barrier}\PY{+w}{ }\PY{o}{=}\PY{+w}{ }\PY{k}{new}\PY{+w}{ }\PY{n}{CyclicBarrier}\PY{p}{(}\PY{l+m+mi}{3}\PY{p}{)}\PY{p}{;}
\end{Verbatim}
\end{codebox}

\item 
\textbf{Leaving a \code{CyclicBarrier} broken after a timeout or interruption.} Why it matters: once one waiting thread times out or is interrupted, the barrier becomes broken and throws \code{BrokenBarrierException} for every other party — the barrier must be explicitly \code{reset()} before it can be used again, and skipping that hangs or fails the next round.

\begin{codebox}
\begin{Verbatim}[commandchars=\\\{\}]
\PY{c+c1}{// Before: catch BrokenBarrierException and just log it, never reset the barrier}
\PY{c+c1}{// After}
\PY{k}{catch}\PY{+w}{ }\PY{p}{(}\PY{n}{BrokenBarrierException}\PY{+w}{ }\PY{o}{|}\PY{+w}{ }\PY{n}{TimeoutException}\PY{+w}{ }\PY{n}{e}\PY{p}{)}\PY{+w}{ }\PY{p}{\PYZob{}}
\PY{+w}{    }\PY{n}{barrier}\PY{p}{.}\PY{n+na}{reset}\PY{p}{(}\PY{p}{)}\PY{p}{;}\PY{+w}{ }\PY{c+c1}{// required before the barrier can be reused}
\PY{p}{\PYZcb{}}
\end{Verbatim}
\end{codebox}

\item 
\textbf{Releasing a \code{Semaphore} permit without a \code{try}/\code{finally}.} Why it matters: if the guarded code throws, the permit is never released, permanently shrinking the effective pool size — a slow permit leak that eventually blocks everyone.

\begin{codebox}
\begin{Verbatim}[commandchars=\\\{\}]
\PY{c+c1}{// Before}
\PY{n}{permits}\PY{p}{.}\PY{n+na}{acquire}\PY{p}{(}\PY{p}{)}\PY{p}{;}
\PY{n}{doWork}\PY{p}{(}\PY{p}{)}\PY{p}{;}\PY{+w}{ }\PY{c+c1}{// if this throws, release() below never runs}
\PY{n}{permits}\PY{p}{.}\PY{n+na}{release}\PY{p}{(}\PY{p}{)}\PY{p}{;}
\PY{c+c1}{// After}
\PY{n}{permits}\PY{p}{.}\PY{n+na}{acquire}\PY{p}{(}\PY{p}{)}\PY{p}{;}
\PY{k}{try}\PY{+w}{ }\PY{p}{\PYZob{}}
\PY{+w}{    }\PY{n}{doWork}\PY{p}{(}\PY{p}{)}\PY{p}{;}
\PY{p}{\PYZcb{}}\PY{+w}{ }\PY{k}{finally}\PY{+w}{ }\PY{p}{\PYZob{}}
\PY{+w}{    }\PY{n}{permits}\PY{p}{.}\PY{n+na}{release}\PY{p}{(}\PY{p}{)}\PY{p}{;}
\PY{p}{\PYZcb{}}
\end{Verbatim}
\end{codebox}

\item 
\textbf{Ignoring \code{InterruptedException} by swallowing it instead of restoring the interrupt status.} Why it matters: swallowing the exception silently (empty catch block) breaks cooperative cancellation — code further up the call stack that checks \code{Thread.\allowbreak{}interrupted()} never finds out the thread was asked to stop.

\begin{codebox}
\begin{Verbatim}[commandchars=\\\{\}]
\PY{c+c1}{// Before}
\PY{k}{try}\PY{+w}{ }\PY{p}{\PYZob{}}\PY{+w}{ }\PY{n}{queue}\PY{p}{.}\PY{n+na}{take}\PY{p}{(}\PY{p}{)}\PY{p}{;}\PY{+w}{ }\PY{p}{\PYZcb{}}\PY{+w}{ }\PY{k}{catch}\PY{+w}{ }\PY{p}{(}\PY{n}{InterruptedException}\PY{+w}{ }\PY{n}{e}\PY{p}{)}\PY{+w}{ }\PY{p}{\PYZob{}}\PY{+w}{ }\PY{c+cm}{/* ignored */}\PY{+w}{ }\PY{p}{\PYZcb{}}
\PY{c+c1}{// After}
\PY{k}{try}\PY{+w}{ }\PY{p}{\PYZob{}}\PY{+w}{ }\PY{n}{queue}\PY{p}{.}\PY{n+na}{take}\PY{p}{(}\PY{p}{)}\PY{p}{;}\PY{+w}{ }\PY{p}{\PYZcb{}}\PY{+w}{ }\PY{k}{catch}\PY{+w}{ }\PY{p}{(}\PY{n}{InterruptedException}\PY{+w}{ }\PY{n}{e}\PY{p}{)}\PY{+w}{ }\PY{p}{\PYZob{}}\PY{+w}{ }\PY{n}{Thread}\PY{p}{.}\PY{n+na}{currentThread}\PY{p}{(}\PY{p}{)}\PY{p}{.}\PY{n+na}{interrupt}\PY{p}{(}\PY{p}{)}\PY{p}{;}\PY{+w}{ }\PY{p}{\PYZcb{}}
\end{Verbatim}
\end{codebox}

\item 
\textbf{Using \code{invokeAny} without realizing it cancels the other in-flight tasks.} Why it matters: reviewers sometimes assume all submitted tasks run to completion; \code{invokeAny} returns as soon as one succeeds and actively cancels the rest, so any side effects those tasks were relying on to complete may never finish.

\begin{codebox}
\begin{Verbatim}[commandchars=\\\{\}]
\PY{c+c1}{// Before: assuming all 3 \PYZdq{}save to replica\PYZdq{} tasks fully complete}
\PY{n}{String}\PY{+w}{ }\PY{n}{result}\PY{+w}{ }\PY{o}{=}\PY{+w}{ }\PY{n}{pool}\PY{p}{.}\PY{n+na}{invokeAny}\PY{p}{(}\PY{n}{saveToReplicaTasks}\PY{p}{)}\PY{p}{;}\PY{+w}{ }\PY{c+c1}{// the other 2 get cancelled mid\PYZhy{}flight}
\PY{c+c1}{// After: use invokeAll if every task\PYZsq{}s side effect must complete, invokeAny only for}
\PY{c+c1}{// \PYZdq{}first successful result wins, and it\PYZsq{}s fine to abandon the rest\PYZdq{}}
\PY{n}{List}\PY{o}{\PYZlt{}}\PY{n}{Future}\PY{o}{\PYZlt{}}\PY{n}{String}\PY{o}{\PYZgt{}}\PY{o}{\PYZgt{}}\PY{+w}{ }\PY{n}{results}\PY{+w}{ }\PY{o}{=}\PY{+w}{ }\PY{n}{pool}\PY{p}{.}\PY{n+na}{invokeAll}\PY{p}{(}\PY{n}{saveToReplicaTasks}\PY{p}{)}\PY{p}{;}
\end{Verbatim}
\end{codebox}

\end{enumerate}


\setcounter{chapter}{14}
\chapter{Modern Concurrency (Java 21+)}
\label{chap:15}
\label{h:15:15-modern-concurrency-java-21}
Traditional Java concurrency (Chapters 13-14) is built on \textbf{platform threads} — thin wrappers around OS threads — and unstructured coordination primitives like \code{ExecutorService}. That model works, but it has two long-standing pain points: platform threads are too expensive to create “one per request,” and fan-out/fan-in logic built from raw executors and futures is easy to get wrong (leaked tasks, missed cancellation, unclear ownership). Project Loom, delivered across JDK 19-25, attacks both problems: \textbf{virtual threads} make blocking, thread-per-request code cheap again, \textbf{structured concurrency} gives fan-out/fan-in a well-defined lifetime and error model, and \textbf{scoped values} give virtual threads a safer, cheaper alternative to \code{ThreadLocal}. This chapter targets JDK 21 (where virtual threads became final) through JDK 25 (where structured concurrency’s API was reshaped), and is explicit everywhere about what is final versus preview in which release.

\section[Virtual Threads (Project Loom)]{Virtual Threads (Project Loom)}
\label{h:15:virtual-threads-project-loom}
\subsection[The Thread-Per-Request Problem]{The Thread-Per-Request Problem}
\label{h:15:the-thread-per-request-problem}
The classic Java server model is thread-per-request: accept a connection, hand it a dedicated thread, let that thread block on I/O (database calls, HTTP calls to other services, disk reads) for as long as it needs. This model is simple to write, simple to debug (one thread = one request, so a thread dump tells you exactly what every request is doing), and works naturally with existing blocking APIs like JDBC, \code{java.\allowbreak{}io}, and synchronous HTTP clients.

The problem is that \textbf{platform threads are expensive}:

\begin{itemize}
\item 
Each platform thread reserves a large stack (commonly 512 KB-1 MB, configurable via \code{-\allowbreak{}Xss}), committed as OS memory.

\item 
Creating a platform thread means an OS-level \code{clone}/\code{pthread\_\allowbreak{}create} syscall — not free.

\item 
The OS scheduler multiplexes a limited number of platform threads across a limited number of CPU cores; context switches between them are relatively costly.

\item 
Because of the memory and scheduling cost, practical platform thread pools cap out somewhere in the low thousands, not millions.

\end{itemize}

So under load, thread-per-request architectures hit a wall: to serve more concurrent \emph{blocked} requests (e.g., waiting on a slow downstream service), you need more threads, but you can’t create enough of them without exhausting memory or thrashing the scheduler. Teams historically worked around this with either bounded thread pools (which cap throughput and can deadlock under saturation) or asynchronous/reactive programming (\code{CompletableFuture} chains, reactive streams), which scales far better but fragments the code into callbacks, breaks stack traces, and makes debugging much harder — you lose the “one thread = one request” story.

\begin{codebox}
\begin{Verbatim}[commandchars=\\\{\}]
\PY{c+c1}{// Classic thread\PYZhy{}per\PYZhy{}request server thread pool — caps around a few thousand}
\PY{c+c1}{// concurrent in\PYZhy{}flight requests no matter how much idle (blocked\PYZhy{}on\PYZhy{}I/O) time}
\PY{c+c1}{// each request spends, because platform threads are the scarce resource.}
\PY{n}{ExecutorService}\PY{+w}{ }\PY{n}{pool}\PY{+w}{ }\PY{o}{=}\PY{+w}{ }\PY{n}{Executors}\PY{p}{.}\PY{n+na}{newFixedThreadPool}\PY{p}{(}\PY{l+m+mi}{200}\PY{p}{)}\PY{p}{;}
\PY{k}{for}\PY{+w}{ }\PY{p}{(}\PY{n}{Socket}\PY{+w}{ }\PY{n}{client}\PY{+w}{ }\PY{p}{:}\PY{+w}{ }\PY{n}{incomingConnections}\PY{p}{(}\PY{p}{)}\PY{p}{)}\PY{+w}{ }\PY{p}{\PYZob{}}
\PY{+w}{    }\PY{n}{pool}\PY{p}{.}\PY{n+na}{submit}\PY{p}{(}\PY{p}{(}\PY{p}{)}\PY{+w}{ }\PY{o}{\PYZhy{}}\PY{o}{\PYZgt{}}\PY{+w}{ }\PY{n}{handleRequest}\PY{p}{(}\PY{n}{client}\PY{p}{)}\PY{p}{)}\PY{p}{;}\PY{+w}{ }\PY{c+c1}{// blocks on I/O inside handleRequest}
\PY{p}{\PYZcb{}}
\end{Verbatim}
\end{codebox}

\subsection[What a Virtual Thread Is]{What a Virtual Thread Is}
\label{h:15:what-a-virtual-thread-is}
A \textbf{virtual thread} is a \code{Thread} (same \code{java.\allowbreak{}lang.\allowbreak{}Thread} API, same \code{Runnable}/interrupt/join semantics) whose stack is \textbf{not} a fixed-size native OS stack, but a \emph{continuation} stored on the Java heap that can grow and shrink. Virtual threads are scheduled by the JVM, not the OS, using a small pool of platform threads called \textbf{carrier threads} (by default, sized to the number of CPU cores, via a \code{ForkJoinPool} in FIFO mode).

The key idea is \textbf{mounting}: when a virtual thread runs, the JVM \emph{mounts} it onto a carrier thread — the carrier thread executes the virtual thread’s code for a while. When the virtual thread does something blocking (e.g., a socket read, \code{Thread.\allowbreak{}sleep}, acquiring a \code{ReentrantLock}), the JVM recognizes this and \textbf{unmounts} the virtual thread from its carrier: it parks the continuation (saves just the relevant stack frames to the heap) and frees the carrier thread to mount and run a \emph{different} virtual thread. When the blocking operation completes, the JVM finds a free carrier thread and \textbf{remounts} the virtual thread to resume exactly where it left off.

\begin{codebox}
\begin{Verbatim}
Carrier thread (platform thread)
   │
   ├── mounts virtual-thread-A ── runs until A blocks on I/O ── unmounts A
   ├── mounts virtual-thread-B ── runs until B blocks on I/O ── unmounts B
   ├── mounts virtual-thread-C ── runs to completion ────────── unmounts C
   └── mounts virtual-thread-A again (I/O finished) ── resumes A
\end{Verbatim}
\end{codebox}

Because the “expensive” part — the stack — lives on the heap as a lightweight, resizable object rather than as a pre-allocated OS stack, and because there is no OS-level thread creation involved, you can have \textbf{millions} of virtual threads outstanding, most of them parked waiting on I/O, backed by only a handful of carrier threads actually consuming CPU.

Crucially, this is invisible to your code: you still write plain, sequential, blocking-style code. \code{Thread.\allowbreak{}sleep(...)}, blocking \code{Socket} reads, JDBC calls, and synchronous \code{HttpClient.\allowbreak{}send(...)} all “just work” and internally cooperate with the scheduler to unmount the virtual thread while waiting, instead of tying up an OS thread.

\subsection[Creating Virtual Threads]{Creating Virtual Threads}
\label{h:15:creating-virtual-threads}
\begin{codebox}
\begin{Verbatim}[commandchars=\\\{\}]
\PY{c+c1}{// 1. Thread.ofVirtual() builder — full control (name, inheritance of context, etc.)}
\PY{n}{Thread}\PY{+w}{ }\PY{n}{vt}\PY{+w}{ }\PY{o}{=}\PY{+w}{ }\PY{n}{Thread}\PY{p}{.}\PY{n+na}{ofVirtual}\PY{p}{(}\PY{p}{)}
\PY{+w}{        }\PY{p}{.}\PY{n+na}{name}\PY{p}{(}\PY{l+s}{\PYZdq{}}\PY{l+s}{order\PYZhy{}worker\PYZhy{}}\PY{l+s}{\PYZdq{}}\PY{p}{,}\PY{+w}{ }\PY{l+m+mi}{0}\PY{p}{)}
\PY{+w}{        }\PY{p}{.}\PY{n+na}{start}\PY{p}{(}\PY{p}{(}\PY{p}{)}\PY{+w}{ }\PY{o}{\PYZhy{}}\PY{o}{\PYZgt{}}\PY{+w}{ }\PY{n}{processOrder}\PY{p}{(}\PY{n}{orderId}\PY{p}{)}\PY{p}{)}\PY{p}{;}
\PY{n}{vt}\PY{p}{.}\PY{n+na}{join}\PY{p}{(}\PY{p}{)}\PY{p}{;}

\PY{c+c1}{// 2. Thread.startVirtualThread — convenience one\PYZhy{}liner, equivalent to}
\PY{c+c1}{//    Thread.ofVirtual().start(task)}
\PY{n}{Thread}\PY{+w}{ }\PY{n}{vt2}\PY{+w}{ }\PY{o}{=}\PY{+w}{ }\PY{n}{Thread}\PY{p}{.}\PY{n+na}{startVirtualThread}\PY{p}{(}\PY{p}{(}\PY{p}{)}\PY{+w}{ }\PY{o}{\PYZhy{}}\PY{o}{\PYZgt{}}\PY{+w}{ }\PY{p}{\PYZob{}}
\PY{+w}{    }\PY{n}{System}\PY{p}{.}\PY{n+na}{out}\PY{p}{.}\PY{n+na}{println}\PY{p}{(}\PY{l+s}{\PYZdq{}}\PY{l+s}{Running on: }\PY{l+s}{\PYZdq{}}\PY{+w}{ }\PY{o}{+}\PY{+w}{ }\PY{n}{Thread}\PY{p}{.}\PY{n+na}{currentThread}\PY{p}{(}\PY{p}{)}\PY{p}{)}\PY{p}{;}
\PY{p}{\PYZcb{}}\PY{p}{)}\PY{p}{;}

\PY{c+c1}{// 3. Executors.newVirtualThreadPerTaskExecutor() — an ExecutorService that}
\PY{c+c1}{//    starts a brand\PYZhy{}new virtual thread for every submitted task. This is the}
\PY{c+c1}{//    recommended way to adopt virtual threads inside existing executor\PYZhy{}based code.}
\PY{k}{try}\PY{+w}{ }\PY{p}{(}\PY{n}{ExecutorService}\PY{+w}{ }\PY{n}{executor}\PY{+w}{ }\PY{o}{=}\PY{+w}{ }\PY{n}{Executors}\PY{p}{.}\PY{n+na}{newVirtualThreadPerTaskExecutor}\PY{p}{(}\PY{p}{)}\PY{p}{)}\PY{+w}{ }\PY{p}{\PYZob{}}
\PY{+w}{    }\PY{n}{List}\PY{o}{\PYZlt{}}\PY{n}{Future}\PY{o}{\PYZlt{}}\PY{n}{String}\PY{o}{\PYZgt{}}\PY{o}{\PYZgt{}}\PY{+w}{ }\PY{n}{results}\PY{+w}{ }\PY{o}{=}\PY{+w}{ }\PY{k}{new}\PY{+w}{ }\PY{n}{ArrayList}\PY{o}{\PYZlt{}}\PY{o}{\PYZgt{}}\PY{p}{(}\PY{p}{)}\PY{p}{;}
\PY{+w}{    }\PY{k}{for}\PY{+w}{ }\PY{p}{(}\PY{k+kt}{int}\PY{+w}{ }\PY{n}{i}\PY{+w}{ }\PY{o}{=}\PY{+w}{ }\PY{l+m+mi}{0}\PY{p}{;}\PY{+w}{ }\PY{n}{i}\PY{+w}{ }\PY{o}{\PYZlt{}}\PY{+w}{ }\PY{l+m+mi}{10\PYZus{}000}\PY{p}{;}\PY{+w}{ }\PY{n}{i}\PY{o}{+}\PY{o}{+}\PY{p}{)}\PY{+w}{ }\PY{p}{\PYZob{}}
\PY{+w}{        }\PY{k+kt}{int}\PY{+w}{ }\PY{n}{id}\PY{+w}{ }\PY{o}{=}\PY{+w}{ }\PY{n}{i}\PY{p}{;}
\PY{+w}{        }\PY{n}{results}\PY{p}{.}\PY{n+na}{add}\PY{p}{(}\PY{n}{executor}\PY{p}{.}\PY{n+na}{submit}\PY{p}{(}\PY{p}{(}\PY{p}{)}\PY{+w}{ }\PY{o}{\PYZhy{}}\PY{o}{\PYZgt{}}\PY{+w}{ }\PY{n}{callDownstreamService}\PY{p}{(}\PY{n}{id}\PY{p}{)}\PY{p}{)}\PY{p}{)}\PY{p}{;}
\PY{+w}{    }\PY{p}{\PYZcb{}}
\PY{+w}{    }\PY{k}{for}\PY{+w}{ }\PY{p}{(}\PY{n}{Future}\PY{o}{\PYZlt{}}\PY{n}{String}\PY{o}{\PYZgt{}}\PY{+w}{ }\PY{n}{f}\PY{+w}{ }\PY{p}{:}\PY{+w}{ }\PY{n}{results}\PY{p}{)}\PY{+w}{ }\PY{p}{\PYZob{}}
\PY{+w}{        }\PY{n}{System}\PY{p}{.}\PY{n+na}{out}\PY{p}{.}\PY{n+na}{println}\PY{p}{(}\PY{n}{f}\PY{p}{.}\PY{n+na}{get}\PY{p}{(}\PY{p}{)}\PY{p}{)}\PY{p}{;}
\PY{+w}{    }\PY{p}{\PYZcb{}}
\PY{p}{\PYZcb{}}\PY{+w}{ }\PY{c+c1}{// executor.close() waits for in\PYZhy{}flight tasks, same as shutdown + awaitTermination}
\end{Verbatim}
\end{codebox}

\code{Thread.\allowbreak{}ofPlatform()} is the analogous builder for platform threads, useful when you need to opt back into a real OS thread deliberately (e.g., for a CPU-bound worker, discussed below).

\subsection[Final in JDK 21 (JEP 444)]{Final in JDK 21 (JEP 444)}
\label{h:15:final-in-jdk-21-jep-444}
Virtual threads shipped as a \textbf{preview feature} in JDK 19 (JEP 425) and JDK 20 (JEP 436), requiring \code{--\allowbreak{}enable-\allowbreak{}preview} to compile and run. They became a \textbf{final, standard feature in JDK 21} under \textbf{JEP 444} — no preview flag needed from JDK 21 onward. This is one of the headline reasons JDK 21 (an LTS release) is treated as the practical starting point for “modern concurrency” in production Java.

\subsection[When Virtual Threads Help — and When They Don’t]{When Virtual Threads Help — and When They Don’t}
\label{h:15:when-virtual-threads-help--and-when-they-dont}
Virtual threads help enormously for \textbf{I/O-bound, high-concurrency, blocking-style code}: web servers handling many concurrent requests that spend most of their time waiting on databases, REST calls to other services, message queues, or file/network I/O. The scheduler can run thousands of such threads on a handful of CPU cores because they’re parked (not consuming CPU) most of the time.

Virtual threads do \textbf{not} help — and can actively hurt — for \textbf{CPU-bound work}. If a virtual thread is doing tight-loop computation (image processing, cryptography, sorting large in-memory data, JSON parsing of huge payloads), it never yields to the scheduler; it simply occupies its carrier thread until done, exactly like a platform thread would. Since there are only as many carrier threads as CPU cores, spinning up a million virtual threads that are all CPU-bound just serializes them behind the same small carrier pool — you get no more parallelism than you had before, plus the (small) overhead of virtual thread bookkeeping. For CPU-bound work, the right tool is still a \textbf{bounded platform thread pool} sized to the number of cores (e.g., \code{ForkJoinPool} or a fixed \code{ExecutorService}), so you get genuine parallel CPU utilization without oversubscribing cores.

\begin{codebox}
\begin{Verbatim}[commandchars=\\\{\}]
\PY{c+c1}{// GOOD fit: thousands of concurrent virtual threads, each mostly blocked on I/O}
\PY{k}{try}\PY{+w}{ }\PY{p}{(}\PY{k+kd}{var}\PY{+w}{ }\PY{n}{executor}\PY{+w}{ }\PY{o}{=}\PY{+w}{ }\PY{n}{Executors}\PY{p}{.}\PY{n+na}{newVirtualThreadPerTaskExecutor}\PY{p}{(}\PY{p}{)}\PY{p}{)}\PY{+w}{ }\PY{p}{\PYZob{}}
\PY{+w}{    }\PY{k}{for}\PY{+w}{ }\PY{p}{(}\PY{n}{String}\PY{+w}{ }\PY{n}{url}\PY{+w}{ }\PY{p}{:}\PY{+w}{ }\PY{n}{urls}\PY{p}{)}\PY{+w}{ }\PY{p}{\PYZob{}}
\PY{+w}{        }\PY{n}{executor}\PY{p}{.}\PY{n+na}{submit}\PY{p}{(}\PY{p}{(}\PY{p}{)}\PY{+w}{ }\PY{o}{\PYZhy{}}\PY{o}{\PYZgt{}}\PY{+w}{ }\PY{n}{httpClient}\PY{p}{.}\PY{n+na}{send}\PY{p}{(}\PY{n}{request}\PY{p}{(}\PY{n}{url}\PY{p}{)}\PY{p}{,}\PY{+w}{ }\PY{n}{BodyHandlers}\PY{p}{.}\PY{n+na}{ofString}\PY{p}{(}\PY{p}{)}\PY{p}{)}\PY{p}{)}\PY{p}{;}
\PY{+w}{    }\PY{p}{\PYZcb{}}
\PY{p}{\PYZcb{}}

\PY{c+c1}{// BAD fit: CPU\PYZhy{}bound matrix multiplication on virtual threads gains nothing —}
\PY{c+c1}{// use a bounded platform\PYZhy{}thread pool (e.g., Executors.newFixedThreadPool(cores))}
\PY{c+c1}{// or ForkJoinPool.commonPool() instead.}
\PY{k}{try}\PY{+w}{ }\PY{p}{(}\PY{k+kd}{var}\PY{+w}{ }\PY{n}{executor}\PY{+w}{ }\PY{o}{=}\PY{+w}{ }\PY{n}{Executors}\PY{p}{.}\PY{n+na}{newVirtualThreadPerTaskExecutor}\PY{p}{(}\PY{p}{)}\PY{p}{)}\PY{+w}{ }\PY{p}{\PYZob{}}
\PY{+w}{    }\PY{k}{for}\PY{+w}{ }\PY{p}{(}\PY{k+kt}{double}\PY{o}{[}\PY{o}{]}\PY{o}{[}\PY{o}{]}\PY{+w}{ }\PY{n}{matrix}\PY{+w}{ }\PY{p}{:}\PY{+w}{ }\PY{n}{matrices}\PY{p}{)}\PY{+w}{ }\PY{p}{\PYZob{}}
\PY{+w}{        }\PY{n}{executor}\PY{p}{.}\PY{n+na}{submit}\PY{p}{(}\PY{p}{(}\PY{p}{)}\PY{+w}{ }\PY{o}{\PYZhy{}}\PY{o}{\PYZgt{}}\PY{+w}{ }\PY{n}{multiply}\PY{p}{(}\PY{n}{matrix}\PY{p}{)}\PY{p}{)}\PY{p}{;}\PY{+w}{ }\PY{c+c1}{// pure CPU work — no yield points}
\PY{+w}{    }\PY{p}{\PYZcb{}}
\PY{p}{\PYZcb{}}
\end{Verbatim}
\end{codebox}

\subsection[Pinning]{Pinning}
\label{h:15:pinning}
\textbf{Pinning} happens when a virtual thread blocks while it \emph{cannot} be unmounted from its carrier thread — the carrier thread is stuck waiting right along with it, defeating the whole point of virtual threads and reducing the number of usable carriers for everyone else.

In \textbf{JDK 21}, pinning is caused by:

\begin{enumerate}
\item 
\textbf{\code{synchronized} blocks or methods} — if a virtual thread blocks \emph{inside} a \code{synchronized} block (e.g., waiting on another lock, or performing blocking I/O while holding the monitor), it cannot be unmounted; the carrier thread is pinned until the virtual thread exits the \code{synchronized} region.

\item 
\textbf{Native frames} — if the virtual thread’s call stack includes a native method (JNI) between it and the blocking operation, the JVM cannot unmount it, because it cannot relocate native stack frames.

\end{enumerate}

\begin{codebox}
\begin{Verbatim}[commandchars=\\\{\}]
\PY{c+c1}{// Pinning hazard in JDK 21: blocking I/O while holding a monitor pins the}
\PY{c+c1}{// carrier thread for the duration of the network call.}
\PY{k+kd}{private}\PY{+w}{ }\PY{k+kd}{final}\PY{+w}{ }\PY{n}{Object}\PY{+w}{ }\PY{n}{lock}\PY{+w}{ }\PY{o}{=}\PY{+w}{ }\PY{k}{new}\PY{+w}{ }\PY{n}{Object}\PY{p}{(}\PY{p}{)}\PY{p}{;}

\PY{k+kt}{void}\PY{+w}{ }\PY{n+nf}{handle}\PY{p}{(}\PY{p}{)}\PY{+w}{ }\PY{p}{\PYZob{}}
\PY{+w}{    }\PY{k+kd}{synchronized}\PY{+w}{ }\PY{p}{(}\PY{n}{lock}\PY{p}{)}\PY{+w}{ }\PY{p}{\PYZob{}}
\PY{+w}{        }\PY{c+c1}{// If this blocks (slow socket, slow DB), the carrier thread is pinned —}
\PY{+w}{        }\PY{c+c1}{// it cannot run any other virtual thread until this call returns.}
\PY{+w}{        }\PY{n}{String}\PY{+w}{ }\PY{n}{response}\PY{+w}{ }\PY{o}{=}\PY{+w}{ }\PY{n}{callSlowLegacyService}\PY{p}{(}\PY{p}{)}\PY{p}{;}
\PY{+w}{        }\PY{n}{cache}\PY{p}{.}\PY{n+na}{put}\PY{p}{(}\PY{n}{key}\PY{p}{,}\PY{+w}{ }\PY{n}{response}\PY{p}{)}\PY{p}{;}
\PY{+w}{    }\PY{p}{\PYZcb{}}
\PY{p}{\PYZcb{}}
\end{Verbatim}
\end{codebox}

The fix is usually to replace \code{synchronized} with \code{java.\allowbreak{}util.\allowbreak{}concurrent.\allowbreak{}locks.\allowbreak{}Reentrant\allowbreak{}Lock}, which is pinning-safe, and to keep blocking calls out of any remaining \code{synchronized} sections:

\begin{codebox}
\begin{Verbatim}[commandchars=\\\{\}]
\PY{k+kd}{private}\PY{+w}{ }\PY{k+kd}{final}\PY{+w}{ }\PY{n}{ReentrantLock}\PY{+w}{ }\PY{n}{lock}\PY{+w}{ }\PY{o}{=}\PY{+w}{ }\PY{k}{new}\PY{+w}{ }\PY{n}{ReentrantLock}\PY{p}{(}\PY{p}{)}\PY{p}{;}

\PY{k+kt}{void}\PY{+w}{ }\PY{n+nf}{handle}\PY{p}{(}\PY{p}{)}\PY{+w}{ }\PY{p}{\PYZob{}}
\PY{+w}{    }\PY{n}{lock}\PY{p}{.}\PY{n+na}{lock}\PY{p}{(}\PY{p}{)}\PY{p}{;}
\PY{+w}{    }\PY{k}{try}\PY{+w}{ }\PY{p}{\PYZob{}}
\PY{+w}{        }\PY{n}{String}\PY{+w}{ }\PY{n}{response}\PY{+w}{ }\PY{o}{=}\PY{+w}{ }\PY{n}{callSlowLegacyService}\PY{p}{(}\PY{p}{)}\PY{p}{;}\PY{+w}{ }\PY{c+c1}{// no pinning — ReentrantLock}
\PY{+w}{        }\PY{n}{cache}\PY{p}{.}\PY{n+na}{put}\PY{p}{(}\PY{n}{key}\PY{p}{,}\PY{+w}{ }\PY{n}{response}\PY{p}{)}\PY{p}{;}\PY{+w}{                   }\PY{c+c1}{// cooperates with the scheduler}
\PY{+w}{    }\PY{p}{\PYZcb{}}\PY{+w}{ }\PY{k}{finally}\PY{+w}{ }\PY{p}{\PYZob{}}
\PY{+w}{        }\PY{n}{lock}\PY{p}{.}\PY{n+na}{unlock}\PY{p}{(}\PY{p}{)}\PY{p}{;}
\PY{+w}{    }\PY{p}{\PYZcb{}}
\PY{p}{\PYZcb{}}
\end{Verbatim}
\end{codebox}

\textbf{JDK 24 (JEP 491)} removed the \code{synchronized}-block pinning problem entirely: virtual threads can now be unmounted even while holding a monitor acquired via \code{synchronized}. This was a significant JVM-internal change (it required reworking how monitors interact with the continuation mechanism) and it means that, from JDK 24 onward, \code{synchronized} is no longer a pinning hazard — only native frames still pin. This is a common interview trap: candidates who memorized “avoid \code{synchronized} with virtual threads” as an absolute rule need to know it is version-specific (true for JDK 21-23, no longer the dominant concern from JDK 24 on), though switching to \code{ReentrantLock} remains harmless and is still good practice for other reasons (e.g., \code{tryLock}, fairness, interruptible acquisition).

\subsubsection[Detecting Pinning]{Detecting Pinning}
\label{h:15:detecting-pinning}
Two supported ways to detect pinning in production or during testing:

\begin{codebox}
\begin{Verbatim}
-Djdk.tracePinnedThreads=full
\end{Verbatim}
\end{codebox}

Run with this system property and the JVM prints a stack trace every time a virtual thread parks while pinned, showing exactly which frame (the \code{synchronized} block or native call) caused it:

\begin{codebox}
\begin{Verbatim}
Thread[#31,tid=31]
    java.base/java.lang.VirtualThread.parkOnCarrierThread(VirtualThread.java:661)
    java.base/java.lang.VirtualThread.park(VirtualThread.java:592)
    ...
    <== monitors:1
    OrderService.handle(OrderService.java:42) <-- monitor
\end{Verbatim}
\end{codebox}

Use \code{full} for complete stack traces, or \code{short} for a one-line-per-occurrence summary.

The second option is \textbf{JDK Flight Recorder (JFR)}: virtual thread pinning emits a \code{jdk.\allowbreak{}Virtual\allowbreak{}Thread\allowbreak{}Pinned} event, which you can capture with a JFR recording and inspect in JDK Mission Control (JMC) or via \code{jfr~\allowbreak{}print~\allowbreak{}--\allowbreak{}events~\allowbreak{}jdk.\allowbreak{}Virtual\allowbreak{}Thread\allowbreak{}Pinned~\allowbreak{}recording.\allowbreak{}jfr}. JFR is lower-overhead and more suitable for always-on production monitoring than the tracing flag, which is best for local debugging or CI.

\begin{codebox}
\begin{Verbatim}[commandchars=\\\{\}]
\PY{c+c1}{\PYZsh{} Capture a JFR recording and filter for pinning events after the fact}
java\PY{+w}{ }\PYZhy{}XX:StartFlightRecording\PY{o}{=}\PY{n+nv}{filename}\PY{o}{=}app.jfr\PY{+w}{ }\PYZhy{}jar\PY{+w}{ }app.jar
jfr\PY{+w}{ }print\PY{+w}{ }\PYZhy{}\PYZhy{}events\PY{+w}{ }jdk.VirtualThreadPinned\PY{+w}{ }app.jfr
\end{Verbatim}
\end{codebox}

\subsection[Why You Should NOT Pool Virtual Threads]{Why You Should NOT Pool Virtual Threads}
\label{h:15:why-you-should-not-pool-virtual-threads}
Thread pools exist to \emph{reuse} expensive platform threads and to \emph{bound} concurrency. Virtual threads are cheap to create and discard — creating a new one is closer to allocating an object than to an OS syscall — so pooling them buys you nothing on the “reuse” side, and actively hurts on the “bounding” side: pooling caps how many tasks can run concurrently by limiting the pool size, which reintroduces exactly the scalability ceiling virtual threads were meant to remove (and can also create priority inversion / deadlock risk if a pooled task blocks waiting on another task that never gets a pool slot).

The idiomatic pattern is \textbf{one virtual thread per task}, unbounded in count, with concurrency to any \emph{scarce downstream resource} controlled separately (see semaphores below) — not by limiting how many virtual threads may exist.

\begin{codebox}
\begin{Verbatim}[commandchars=\\\{\}]
\PY{c+c1}{// Anti\PYZhy{}pattern: don\PYZsq{}t do this — defeats the purpose of virtual threads}
\PY{n}{ExecutorService}\PY{+w}{ }\PY{n}{fakeVirtualPool}\PY{+w}{ }\PY{o}{=}\PY{+w}{ }\PY{k}{new}\PY{+w}{ }\PY{n}{ThreadPoolExecutor}\PY{p}{(}
\PY{+w}{        }\PY{l+m+mi}{200}\PY{p}{,}\PY{+w}{ }\PY{l+m+mi}{200}\PY{p}{,}\PY{+w}{ }\PY{l+m+mi}{0}\PY{n}{L}\PY{p}{,}\PY{+w}{ }\PY{n}{TimeUnit}\PY{p}{.}\PY{n+na}{MILLISECONDS}\PY{p}{,}\PY{+w}{ }\PY{k}{new}\PY{+w}{ }\PY{n}{LinkedBlockingQueue}\PY{o}{\PYZlt{}}\PY{o}{\PYZgt{}}\PY{p}{(}\PY{p}{)}\PY{p}{,}
\PY{+w}{        }\PY{n}{Thread}\PY{p}{.}\PY{n+na}{ofVirtual}\PY{p}{(}\PY{p}{)}\PY{p}{.}\PY{n+na}{factory}\PY{p}{(}\PY{p}{)}\PY{p}{)}\PY{p}{;}\PY{+w}{ }\PY{c+c1}{// pooling virtual threads: pointless and limiting}

\PY{c+c1}{// Idiomatic: unbounded virtual\PYZhy{}thread\PYZhy{}per\PYZhy{}task}
\PY{n}{ExecutorService}\PY{+w}{ }\PY{n}{executor}\PY{+w}{ }\PY{o}{=}\PY{+w}{ }\PY{n}{Executors}\PY{p}{.}\PY{n+na}{newVirtualThreadPerTaskExecutor}\PY{p}{(}\PY{p}{)}\PY{p}{;}
\end{Verbatim}
\end{codebox}

\subsection[ThreadLocal at Millions of Threads]{\code{ThreadLocal} at Millions of Threads}
\label{h:15:threadlocal-at-millions-of-threads}
\code{ThreadLocal} was designed under the assumption that thread counts are small (dozens to low thousands) and thread lifetimes are long (pooled and reused), so the per-thread storage cost was negligible and cheap to amortize. Virtual threads invert both assumptions: there can be millions of them, and each is typically short-lived (one per task/request).

Every \code{ThreadLocal} variable set on a virtual thread still needs its own storage slot on \emph{that} thread, and Java copies \code{ThreadLocal} values into \code{InheritableThreadLocal} children on thread creation. At scale:

\begin{itemize}
\item 
Memory cost multiplies by however many concurrent virtual threads exist — a \code{ThreadLocal\textless{}\allowbreak{}byte[]\textgreater{}} holding even a modest buffer, times a million virtual threads, is a real amount of heap.

\item 
\code{InheritableThreadLocal} inheritance means every child virtual thread created (e.g., inside a \code{StructuredTaskScope}) does a copy/snapshot of the parent’s inheritable locals, which is wasted work if most children never read them.

\item 
\code{ThreadLocal} is \textbf{mutable} and has \textbf{no defined lifetime bound} — nothing forces you to call \code{remove()}, so leaked \code{ThreadLocal} entries silently retain memory (and, in pooled-thread contexts, correctness bugs) until the thread itself is discarded. With millions of short-lived virtual threads this is somewhat self-limiting (thread death frees the map), but the per-thread map allocation and lookup overhead still adds up.

\end{itemize}

This is exactly the gap that \textbf{Scoped Values} (later in this chapter) are designed to close.

\subsection[Semaphores Instead of Thread Pools for Limiting Concurrency]{Semaphores Instead of Thread Pools for Limiting Concurrency}
\label{h:15:semaphores-instead-of-thread-pools-for-limiting-concurrency}
If the goal is not to bound \emph{how many virtual threads exist} but to bound \emph{how many concurrent requests hit a specific downstream resource} (a database connection pool, a rate-limited third-party API, a fixed-capacity in-memory cache), the correct tool is a \code{Semaphore}, acquired around just the call to that resource — not a bounded executor that throttles everything indiscriminately.

\begin{codebox}
\begin{Verbatim}[commandchars=\\\{\}]
\PY{c+c1}{// Bound concurrent calls to a downstream payment gateway to 50, while still}
\PY{c+c1}{// allowing unlimited virtual threads for everything else the request does.}
\PY{k+kd}{private}\PY{+w}{ }\PY{k+kd}{final}\PY{+w}{ }\PY{n}{Semaphore}\PY{+w}{ }\PY{n}{paymentGatewayLimiter}\PY{+w}{ }\PY{o}{=}\PY{+w}{ }\PY{k}{new}\PY{+w}{ }\PY{n}{Semaphore}\PY{p}{(}\PY{l+m+mi}{50}\PY{p}{)}\PY{p}{;}

\PY{n}{String}\PY{+w}{ }\PY{n+nf}{chargeCard}\PY{p}{(}\PY{n}{Order}\PY{+w}{ }\PY{n}{order}\PY{p}{)}\PY{+w}{ }\PY{k+kd}{throws}\PY{+w}{ }\PY{n}{InterruptedException}\PY{+w}{ }\PY{p}{\PYZob{}}
\PY{+w}{    }\PY{n}{paymentGatewayLimiter}\PY{p}{.}\PY{n+na}{acquire}\PY{p}{(}\PY{p}{)}\PY{p}{;}
\PY{+w}{    }\PY{k}{try}\PY{+w}{ }\PY{p}{\PYZob{}}
\PY{+w}{        }\PY{k}{return}\PY{+w}{ }\PY{n}{paymentGateway}\PY{p}{.}\PY{n+na}{charge}\PY{p}{(}\PY{n}{order}\PY{p}{)}\PY{p}{;}\PY{+w}{ }\PY{c+c1}{// blocking call, virtual\PYZhy{}thread friendly}
\PY{+w}{    }\PY{p}{\PYZcb{}}\PY{+w}{ }\PY{k}{finally}\PY{+w}{ }\PY{p}{\PYZob{}}
\PY{+w}{        }\PY{n}{paymentGatewayLimiter}\PY{p}{.}\PY{n+na}{release}\PY{p}{(}\PY{p}{)}\PY{p}{;}
\PY{+w}{    }\PY{p}{\PYZcb{}}
\PY{p}{\PYZcb{}}
\end{Verbatim}
\end{codebox}

Each caller still runs on its own virtual thread and can make unrelated progress (or block cheaply) while waiting for a permit; only access to the \emph{specific scarce resource} is serialized to the configured limit. This decouples “how many logical tasks can be in flight” (effectively unbounded, one virtual thread each) from “how much load a specific downstream can take” (bounded by the semaphore).

\subsection[Observability and Thread Dumps]{Observability and Thread Dumps}
\label{h:15:observability-and-thread-dumps}
Because a JVM might now host millions of virtual threads, traditional \code{jstack}-style dumps (designed for thousands of platform threads) don’t scale well as plain text. JDK 21+ added a dedicated, structured thread dump command:

\begin{codebox}
\begin{Verbatim}[commandchars=\\\{\}]
jcmd\PY{+w}{ }\PYZlt{}pid\PYZgt{}\PY{+w}{ }Thread.dump\PYZus{}to\PYZus{}file\PY{+w}{ }\PYZhy{}format\PY{o}{=}json\PY{+w}{ }threads.json
\end{Verbatim}
\end{codebox}

This produces a JSON dump that groups virtual threads by which carrier thread (if any) they’re currently mounted on, distinguishes platform threads from virtual threads, and is designed to be machine-readable so tooling can visualize or query it (rather than grepping raw text). A plain-text variant is also available:

\begin{codebox}
\begin{Verbatim}[commandchars=\\\{\}]
jcmd\PY{+w}{ }\PYZlt{}pid\PYZgt{}\PY{+w}{ }Thread.dump\PYZus{}to\PYZus{}file\PY{+w}{ }\PYZhy{}format\PY{o}{=}text\PY{+w}{ }threads.txt
\end{Verbatim}
\end{codebox}

For live, low-level thread dumps of a \emph{specific} thread, \code{Thread.\allowbreak{}dump()}-style diagnostics and JFR’s \code{jdk.\allowbreak{}ThreadDump} and \code{jdk.\allowbreak{}Virtual\allowbreak{}Thread\allowbreak{}Pinned}/\code{jdk.\allowbreak{}Virtual\allowbreak{}Thread\allowbreak{}Submit\allowbreak{}Failed} events fill in the rest of the observability story: pinning (above), submission failures (e.g., carrier pool exhaustion under a custom scheduler), and thread lifecycle events are all available for continuous monitoring via JFR rather than one-off dumps.

\subsection[Migration Advice]{Migration Advice}
\label{h:15:migration-advice}
\begin{itemize}
\item 
\textbf{Start at I/O boundaries.} The highest-value, lowest-risk place to introduce virtual threads is swapping a fixed-size platform-thread \code{ExecutorService} used for blocking I/O (handling requests, calling downstream services) for \code{Executors.\allowbreak{}new\allowbreak{}Virtual\allowbreak{}Thread\allowbreak{}Per\allowbreak{}Task\allowbreak{}Executor()}. Frameworks (Tomcat, Jetty, Spring Boot 3.2+) increasingly support flipping this via configuration rather than hand-written executor code.

\item 
\textbf{Audit \code{synchronized} usage on JDK 21-23.} Use \code{-\allowbreak{}Djdk.\allowbreak{}trace\allowbreak{}Pinned\allowbreak{}Threads=\allowbreak{}full} in a load test or staging environment before rollout; replace hot, blocking \code{synchronized} blocks with \code{ReentrantLock}. This becomes far less urgent on JDK 24+ thanks to JEP 491, but native-frame pinning (JNI) can still occur on any version.

\item 
\textbf{Remove artificial concurrency caps that were sized around thread cost.} A \code{new\allowbreak{}Fixed\allowbreak{}Thread\allowbreak{}Pool(\allowbreak{}200)} sized because “200 platform threads is about all we can afford” should usually become unbounded virtual threads plus a \code{Semaphore} around the \emph{actual} scarce resource (see above) — don’t just swap the factory and keep the artificial cap, or you gain nothing.

\item 
\textbf{Don’t rewrite CPU-bound pools.} Leave \code{ForkJoinPool}/fixed platform-thread pools alone for CPU-bound work; virtual threads are not a general replacement for all concurrency.

\item 
\textbf{Watch \code{ThreadLocal}-heavy legacy code}, especially framework code (security contexts, MDC/logging context, transaction contexts) that relies on \code{InheritableThreadLocal}. At virtual-thread scale, either confirm the library has been updated to use scoped values / lightweight inheritance, or measure the actual overhead before assuming it’s fine.

\item 
\textbf{Re-baseline connection pool sizing.} Virtual threads can generate far more concurrent \emph{logical} requests than before; connection pools (DB, HTTP) sized for a small number of platform threads may need semaphore-style gating even if they were “big enough” previously, because now thousands of virtual threads can genuinely try to use them at once.

\end{itemize}

\section[Structured Concurrency]{Structured Concurrency}
\label{h:15:structured-concurrency}
\subsection[The Problem with Unstructured Fan-Out]{The Problem with Unstructured Fan-Out}
\label{h:15:the-problem-with-unstructured-fan-out}
A common pattern is: submit several tasks to an \code{ExecutorService}, then collect their results with \code{Future.\allowbreak{}get()}. This works for the happy path, but it has real structural gaps:

\begin{codebox}
\begin{Verbatim}[commandchars=\\\{\}]
\PY{c+c1}{// Unstructured fan\PYZhy{}out — looks fine, has several latent bugs}
\PY{n}{ExecutorService}\PY{+w}{ }\PY{n}{executor}\PY{+w}{ }\PY{o}{=}\PY{+w}{ }\PY{n}{Executors}\PY{p}{.}\PY{n+na}{newVirtualThreadPerTaskExecutor}\PY{p}{(}\PY{p}{)}\PY{p}{;}
\PY{n}{Future}\PY{o}{\PYZlt{}}\PY{n}{User}\PY{o}{\PYZgt{}}\PY{+w}{ }\PY{n}{userFuture}\PY{+w}{ }\PY{o}{=}\PY{+w}{ }\PY{n}{executor}\PY{p}{.}\PY{n+na}{submit}\PY{p}{(}\PY{p}{(}\PY{p}{)}\PY{+w}{ }\PY{o}{\PYZhy{}}\PY{o}{\PYZgt{}}\PY{+w}{ }\PY{n}{fetchUser}\PY{p}{(}\PY{n}{userId}\PY{p}{)}\PY{p}{)}\PY{p}{;}
\PY{n}{Future}\PY{o}{\PYZlt{}}\PY{n}{Order}\PY{o}{\PYZgt{}}\PY{+w}{ }\PY{n}{orderFuture}\PY{+w}{ }\PY{o}{=}\PY{+w}{ }\PY{n}{executor}\PY{p}{.}\PY{n+na}{submit}\PY{p}{(}\PY{p}{(}\PY{p}{)}\PY{+w}{ }\PY{o}{\PYZhy{}}\PY{o}{\PYZgt{}}\PY{+w}{ }\PY{n}{fetchOrder}\PY{p}{(}\PY{n}{orderId}\PY{p}{)}\PY{p}{)}\PY{p}{;}

\PY{n}{User}\PY{+w}{ }\PY{n}{user}\PY{+w}{ }\PY{o}{=}\PY{+w}{ }\PY{n}{userFuture}\PY{p}{.}\PY{n+na}{get}\PY{p}{(}\PY{p}{)}\PY{p}{;}\PY{+w}{   }\PY{c+c1}{// if this throws, orderFuture keeps running —}
\PY{n}{Order}\PY{+w}{ }\PY{n}{order}\PY{+w}{ }\PY{o}{=}\PY{+w}{ }\PY{n}{orderFuture}\PY{p}{.}\PY{n+na}{get}\PY{p}{(}\PY{p}{)}\PY{p}{;}\PY{+w}{ }\PY{c+c1}{// an orphaned task, nobody waits for or cancels it}
\end{Verbatim}
\end{codebox}

Concretely:

\begin{itemize}
\item 
\textbf{Leaked/orphan tasks.} If \code{userFuture.\allowbreak{}get()} throws, execution never reaches \code{orderFuture.\allowbreak{}get()}. The order-fetching task keeps running in the background — consuming a thread, a downstream connection, CPU — with nothing watching it, joining it, or cancelling it. In a request-handling context this is a silent resource leak that repeats on every failure.

\item 
\textbf{No automatic cancellation propagation.} If one subtask fails, the \emph{others} should usually be cancelled (why keep computing a result nobody will use?), but plain futures don’t do this for you — you have to remember to call \code{cancel()} on every sibling, in every error path, including exceptions thrown from unrelated code between submissions.

\item 
\textbf{Unclear ownership and lifetime.} Nothing in the code enforces that all child tasks finish (or are cancelled) before the enclosing method returns. This violates the intuitive nesting of blocks in sequential code, where a child scope’s lifetime is visibly bounded by \code{\{\allowbreak{}~\allowbreak{}\}}. Concurrent tasks submitted to a shared executor have no such visible boundary — they can outlive the method that created them, which makes reasoning about cancellation, exception propagation, and shutdown far harder than for straight-line code.

\end{itemize}

\textbf{Structured concurrency} fixes this by treating a group of related concurrent subtasks as a single unit of work: they are forked from one \emph{scope}, and that scope cannot be exited (the \code{try}-block cannot complete) until all of its forked subtasks have completed, failed, or been cancelled. The concurrent code gets the same nesting and lifetime guarantees as sequential code — a child task’s lifetime is strictly contained within its parent’s block.

\subsection[StructuredTaskScope Basics]{\code{StructuredTaskScope} Basics}
\label{h:15:structuredtaskscope-basics}
The core type is \code{java.\allowbreak{}util.\allowbreak{}concurrent.\allowbreak{}Structured\allowbreak{}Task\allowbreak{}Scope}, used in a try-with-resources block:

\begin{codebox}
\begin{Verbatim}[commandchars=\\\{\}]
\PY{k}{try}\PY{+w}{ }\PY{p}{(}\PY{k+kd}{var}\PY{+w}{ }\PY{n}{scope}\PY{+w}{ }\PY{o}{=}\PY{+w}{ }\PY{k}{new}\PY{+w}{ }\PY{n}{StructuredTaskScope}\PY{p}{.}\PY{n+na}{ShutdownOnFailure}\PY{p}{(}\PY{p}{)}\PY{p}{)}\PY{+w}{ }\PY{p}{\PYZob{}}
\PY{+w}{    }\PY{n}{Subtask}\PY{o}{\PYZlt{}}\PY{n}{User}\PY{o}{\PYZgt{}}\PY{+w}{ }\PY{n}{userTask}\PY{+w}{ }\PY{o}{=}\PY{+w}{ }\PY{n}{scope}\PY{p}{.}\PY{n+na}{fork}\PY{p}{(}\PY{p}{(}\PY{p}{)}\PY{+w}{ }\PY{o}{\PYZhy{}}\PY{o}{\PYZgt{}}\PY{+w}{ }\PY{n}{fetchUser}\PY{p}{(}\PY{n}{userId}\PY{p}{)}\PY{p}{)}\PY{p}{;}
\PY{+w}{    }\PY{n}{Subtask}\PY{o}{\PYZlt{}}\PY{n}{Order}\PY{o}{\PYZgt{}}\PY{+w}{ }\PY{n}{orderTask}\PY{+w}{ }\PY{o}{=}\PY{+w}{ }\PY{n}{scope}\PY{p}{.}\PY{n+na}{fork}\PY{p}{(}\PY{p}{(}\PY{p}{)}\PY{+w}{ }\PY{o}{\PYZhy{}}\PY{o}{\PYZgt{}}\PY{+w}{ }\PY{n}{fetchOrder}\PY{p}{(}\PY{n}{orderId}\PY{p}{)}\PY{p}{)}\PY{p}{;}

\PY{+w}{    }\PY{n}{scope}\PY{p}{.}\PY{n+na}{join}\PY{p}{(}\PY{p}{)}\PY{p}{;}\PY{+w}{            }\PY{c+c1}{// wait for both forks to finish (or one to fail)}
\PY{+w}{    }\PY{n}{scope}\PY{p}{.}\PY{n+na}{throwIfFailed}\PY{p}{(}\PY{p}{)}\PY{p}{;}\PY{+w}{   }\PY{c+c1}{// if either failed, rethrow that failure here}

\PY{+w}{    }\PY{c+c1}{// Both succeeded — safe to read results}
\PY{+w}{    }\PY{n}{User}\PY{+w}{ }\PY{n}{user}\PY{+w}{ }\PY{o}{=}\PY{+w}{ }\PY{n}{userTask}\PY{p}{.}\PY{n+na}{get}\PY{p}{(}\PY{p}{)}\PY{p}{;}
\PY{+w}{    }\PY{n}{Order}\PY{+w}{ }\PY{n}{order}\PY{+w}{ }\PY{o}{=}\PY{+w}{ }\PY{n}{orderTask}\PY{p}{.}\PY{n+na}{get}\PY{p}{(}\PY{p}{)}\PY{p}{;}
\PY{+w}{    }\PY{k}{return}\PY{+w}{ }\PY{k}{new}\PY{+w}{ }\PY{n}{Profile}\PY{p}{(}\PY{n}{user}\PY{p}{,}\PY{+w}{ }\PY{n}{order}\PY{p}{)}\PY{p}{;}
\PY{p}{\PYZcb{}}\PY{+w}{ }\PY{c+c1}{// scope.close() guarantees no forked subtask survives past this point}
\end{Verbatim}
\end{codebox}

\begin{itemize}
\item 
\textbf{\code{fork(\allowbreak{}Callable\textless{}\allowbreak{}T\textgreater{})}} starts a subtask (on its own virtual thread by default) and returns a \code{Subtask\textless{}\allowbreak{}T\textgreater{}} handle.

\item 
\textbf{\code{join()}} blocks until all forked subtasks finish, or (for policy-aware scopes like \code{ShutdownOnFailure}) until the scope’s shutdown condition triggers.

\item 
\textbf{\code{throwIfFailed()}} (on \code{ShutdownOnFailure}) rethrows the first captured failure, wrapped, if any subtask failed — this is what makes error propagation explicit and centralized instead of scattered across each \code{Future.\allowbreak{}get()} call.

\item 
\textbf{\code{close()}}, called automatically by try-with-resources, \textbf{is the structural guarantee}: it waits for any still-running subtasks to finish and, in later API shapes, cancels/interrupts them if the scope is exiting due to an exception. You cannot leave the \code{try} block with orphaned children — that’s the whole point.

\end{itemize}

\subsection[ShutdownOnFailure — All-Must-Succeed]{\code{ShutdownOnFailure} — All-Must-Succeed}
\label{h:15:shutdownonfailure--all-must-succeed}
\code{ShutdownOnFailure} implements the common “fan out, and if \emph{any} subtask fails, cancel the rest and fail fast” policy. This is the right shape when you need every result to build a combined response — e.g., assembling a page that needs a user profile \emph{and} their recent orders \emph{and} their loyalty status; if any one lookup fails, the whole response is incomplete and there’s no point waiting for (or paying for) the others.

\begin{codebox}
\begin{Verbatim}[commandchars=\\\{\}]
\PY{k+kd}{record} \PY{n+nc}{DashboardData}\PY{p}{(}\PY{n}{User}\PY{+w}{ }\PY{n}{user}\PY{p}{,}\PY{+w}{ }\PY{n}{List}\PY{o}{\PYZlt{}}\PY{n}{Order}\PY{o}{\PYZgt{}}\PY{+w}{ }\PY{n}{orders}\PY{p}{,}\PY{+w}{ }\PY{n}{LoyaltyStatus}\PY{+w}{ }\PY{n}{loyalty}\PY{p}{)}\PY{+w}{ }\PY{p}{\PYZob{}}\PY{p}{\PYZcb{}}

\PY{n}{DashboardData}\PY{+w}{ }\PY{n+nf}{loadDashboard}\PY{p}{(}\PY{n}{String}\PY{+w}{ }\PY{n}{userId}\PY{p}{)}\PY{+w}{ }\PY{k+kd}{throws}\PY{+w}{ }\PY{n}{InterruptedException}\PY{p}{,}\PY{+w}{ }\PY{n}{ExecutionException}\PY{+w}{ }\PY{p}{\PYZob{}}
\PY{+w}{    }\PY{k}{try}\PY{+w}{ }\PY{p}{(}\PY{k+kd}{var}\PY{+w}{ }\PY{n}{scope}\PY{+w}{ }\PY{o}{=}\PY{+w}{ }\PY{k}{new}\PY{+w}{ }\PY{n}{StructuredTaskScope}\PY{p}{.}\PY{n+na}{ShutdownOnFailure}\PY{p}{(}\PY{p}{)}\PY{p}{)}\PY{+w}{ }\PY{p}{\PYZob{}}
\PY{+w}{        }\PY{n}{Subtask}\PY{o}{\PYZlt{}}\PY{n}{User}\PY{o}{\PYZgt{}}\PY{+w}{ }\PY{n}{userTask}\PY{+w}{ }\PY{o}{=}\PY{+w}{ }\PY{n}{scope}\PY{p}{.}\PY{n+na}{fork}\PY{p}{(}\PY{p}{(}\PY{p}{)}\PY{+w}{ }\PY{o}{\PYZhy{}}\PY{o}{\PYZgt{}}\PY{+w}{ }\PY{n}{userService}\PY{p}{.}\PY{n+na}{fetch}\PY{p}{(}\PY{n}{userId}\PY{p}{)}\PY{p}{)}\PY{p}{;}
\PY{+w}{        }\PY{n}{Subtask}\PY{o}{\PYZlt{}}\PY{n}{List}\PY{o}{\PYZlt{}}\PY{n}{Order}\PY{o}{\PYZgt{}}\PY{o}{\PYZgt{}}\PY{+w}{ }\PY{n}{ordersTask}\PY{+w}{ }\PY{o}{=}\PY{+w}{ }\PY{n}{scope}\PY{p}{.}\PY{n+na}{fork}\PY{p}{(}\PY{p}{(}\PY{p}{)}\PY{+w}{ }\PY{o}{\PYZhy{}}\PY{o}{\PYZgt{}}\PY{+w}{ }\PY{n}{orderService}\PY{p}{.}\PY{n+na}{recentOrders}\PY{p}{(}\PY{n}{userId}\PY{p}{)}\PY{p}{)}\PY{p}{;}
\PY{+w}{        }\PY{n}{Subtask}\PY{o}{\PYZlt{}}\PY{n}{LoyaltyStatus}\PY{o}{\PYZgt{}}\PY{+w}{ }\PY{n}{loyaltyTask}\PY{+w}{ }\PY{o}{=}\PY{+w}{ }\PY{n}{scope}\PY{p}{.}\PY{n+na}{fork}\PY{p}{(}\PY{p}{(}\PY{p}{)}\PY{+w}{ }\PY{o}{\PYZhy{}}\PY{o}{\PYZgt{}}\PY{+w}{ }\PY{n}{loyaltyService}\PY{p}{.}\PY{n+na}{status}\PY{p}{(}\PY{n}{userId}\PY{p}{)}\PY{p}{)}\PY{p}{;}

\PY{+w}{        }\PY{n}{scope}\PY{p}{.}\PY{n+na}{join}\PY{p}{(}\PY{p}{)}\PY{p}{;}
\PY{+w}{        }\PY{n}{scope}\PY{p}{.}\PY{n+na}{throwIfFailed}\PY{p}{(}\PY{n}{ExecutionException}\PY{p}{:}\PY{p}{:}\PY{k}{new}\PY{p}{)}\PY{p}{;}\PY{+w}{ }\PY{c+c1}{// custom exception wrapper}

\PY{+w}{        }\PY{k}{return}\PY{+w}{ }\PY{k}{new}\PY{+w}{ }\PY{n}{DashboardData}\PY{p}{(}\PY{n}{userTask}\PY{p}{.}\PY{n+na}{get}\PY{p}{(}\PY{p}{)}\PY{p}{,}\PY{+w}{ }\PY{n}{ordersTask}\PY{p}{.}\PY{n+na}{get}\PY{p}{(}\PY{p}{)}\PY{p}{,}\PY{+w}{ }\PY{n}{loyaltyTask}\PY{p}{.}\PY{n+na}{get}\PY{p}{(}\PY{p}{)}\PY{p}{)}\PY{p}{;}
\PY{+w}{    }\PY{p}{\PYZcb{}}
\PY{+w}{    }\PY{c+c1}{// If loyaltyService.status() throws, ShutdownOnFailure immediately signals}
\PY{+w}{    }\PY{c+c1}{// shutdown: userTask and ordersTask are cancelled/interrupted if still}
\PY{+w}{    }\PY{c+c1}{// running, join() returns promptly instead of waiting for them to finish}
\PY{+w}{    }\PY{c+c1}{// naturally, and throwIfFailed() rethrows the loyalty failure.}
\PY{p}{\PYZcb{}}
\end{Verbatim}
\end{codebox}

\subsection[ShutdownOnSuccess — First-Wins]{\code{ShutdownOnSuccess} — First-Wins}
\label{h:15:shutdownonsuccess--first-wins}
\code{ShutdownOnSuccess} implements the opposite policy: race several subtasks and take the \emph{first} one to succeed, cancelling the rest. This is the natural shape for redundant lookups — e.g., querying a primary and a fallback cache/region simultaneously and using whichever answers first, or trying multiple mirrors for the same resource.

\begin{codebox}
\begin{Verbatim}[commandchars=\\\{\}]
\PY{n}{String}\PY{+w}{ }\PY{n+nf}{fetchFromFastestMirror}\PY{p}{(}\PY{n}{String}\PY{+w}{ }\PY{n}{key}\PY{p}{)}\PY{+w}{ }\PY{k+kd}{throws}\PY{+w}{ }\PY{n}{InterruptedException}\PY{p}{,}\PY{+w}{ }\PY{n}{ExecutionException}\PY{+w}{ }\PY{p}{\PYZob{}}
\PY{+w}{    }\PY{k}{try}\PY{+w}{ }\PY{p}{(}\PY{k+kd}{var}\PY{+w}{ }\PY{n}{scope}\PY{+w}{ }\PY{o}{=}\PY{+w}{ }\PY{k}{new}\PY{+w}{ }\PY{n}{StructuredTaskScope}\PY{p}{.}\PY{n+na}{ShutdownOnSuccess}\PY{o}{\PYZlt{}}\PY{n}{String}\PY{o}{\PYZgt{}}\PY{p}{(}\PY{p}{)}\PY{p}{)}\PY{+w}{ }\PY{p}{\PYZob{}}
\PY{+w}{        }\PY{n}{scope}\PY{p}{.}\PY{n+na}{fork}\PY{p}{(}\PY{p}{(}\PY{p}{)}\PY{+w}{ }\PY{o}{\PYZhy{}}\PY{o}{\PYZgt{}}\PY{+w}{ }\PY{n}{primaryRegionCache}\PY{p}{.}\PY{n+na}{get}\PY{p}{(}\PY{n}{key}\PY{p}{)}\PY{p}{)}\PY{p}{;}
\PY{+w}{        }\PY{n}{scope}\PY{p}{.}\PY{n+na}{fork}\PY{p}{(}\PY{p}{(}\PY{p}{)}\PY{+w}{ }\PY{o}{\PYZhy{}}\PY{o}{\PYZgt{}}\PY{+w}{ }\PY{n}{secondaryRegionCache}\PY{p}{.}\PY{n+na}{get}\PY{p}{(}\PY{n}{key}\PY{p}{)}\PY{p}{)}\PY{p}{;}
\PY{+w}{        }\PY{n}{scope}\PY{p}{.}\PY{n+na}{fork}\PY{p}{(}\PY{p}{(}\PY{p}{)}\PY{+w}{ }\PY{o}{\PYZhy{}}\PY{o}{\PYZgt{}}\PY{+w}{ }\PY{n}{tertiaryRegionCache}\PY{p}{.}\PY{n+na}{get}\PY{p}{(}\PY{n}{key}\PY{p}{)}\PY{p}{)}\PY{p}{;}

\PY{+w}{        }\PY{n}{scope}\PY{p}{.}\PY{n+na}{join}\PY{p}{(}\PY{p}{)}\PY{p}{;}
\PY{+w}{        }\PY{k}{return}\PY{+w}{ }\PY{n}{scope}\PY{p}{.}\PY{n+na}{result}\PY{p}{(}\PY{p}{)}\PY{p}{;}\PY{+w}{ }\PY{c+c1}{// returns the first success; throws if all failed}
\PY{+w}{        }\PY{c+c1}{// As soon as one fork succeeds, ShutdownOnSuccess cancels the other}
\PY{+w}{        }\PY{c+c1}{// in\PYZhy{}flight forks — no wasted work, no orphaned background lookups.}
\PY{+w}{    }\PY{p}{\PYZcb{}}
\PY{p}{\PYZcb{}}
\end{Verbatim}
\end{codebox}

\subsection[Timeouts via joinUntil]{Timeouts via \code{joinUntil}}
\label{h:15:timeouts-via-joinuntil}
Instead of (or in addition to) \code{join()}, \code{join\allowbreak{}Until(\allowbreak{}Instant~\allowbreak{}deadline)} waits only until a deadline, throwing \code{TimeoutException} if the subtasks haven’t finished by then — and, combined with a shutdown-triggering policy, still ensures nothing is left running past the scope’s \code{close()}.

\begin{codebox}
\begin{Verbatim}[commandchars=\\\{\}]
\PY{n}{DashboardData}\PY{+w}{ }\PY{n+nf}{loadDashboardWithTimeout}\PY{p}{(}\PY{n}{String}\PY{+w}{ }\PY{n}{userId}\PY{p}{)}
\PY{+w}{        }\PY{k+kd}{throws}\PY{+w}{ }\PY{n}{InterruptedException}\PY{p}{,}\PY{+w}{ }\PY{n}{ExecutionException}\PY{p}{,}\PY{+w}{ }\PY{n}{TimeoutException}\PY{+w}{ }\PY{p}{\PYZob{}}
\PY{+w}{    }\PY{k}{try}\PY{+w}{ }\PY{p}{(}\PY{k+kd}{var}\PY{+w}{ }\PY{n}{scope}\PY{+w}{ }\PY{o}{=}\PY{+w}{ }\PY{k}{new}\PY{+w}{ }\PY{n}{StructuredTaskScope}\PY{p}{.}\PY{n+na}{ShutdownOnFailure}\PY{p}{(}\PY{p}{)}\PY{p}{)}\PY{+w}{ }\PY{p}{\PYZob{}}
\PY{+w}{        }\PY{n}{Subtask}\PY{o}{\PYZlt{}}\PY{n}{User}\PY{o}{\PYZgt{}}\PY{+w}{ }\PY{n}{userTask}\PY{+w}{ }\PY{o}{=}\PY{+w}{ }\PY{n}{scope}\PY{p}{.}\PY{n+na}{fork}\PY{p}{(}\PY{p}{(}\PY{p}{)}\PY{+w}{ }\PY{o}{\PYZhy{}}\PY{o}{\PYZgt{}}\PY{+w}{ }\PY{n}{userService}\PY{p}{.}\PY{n+na}{fetch}\PY{p}{(}\PY{n}{userId}\PY{p}{)}\PY{p}{)}\PY{p}{;}
\PY{+w}{        }\PY{n}{Subtask}\PY{o}{\PYZlt{}}\PY{n}{List}\PY{o}{\PYZlt{}}\PY{n}{Order}\PY{o}{\PYZgt{}}\PY{o}{\PYZgt{}}\PY{+w}{ }\PY{n}{ordersTask}\PY{+w}{ }\PY{o}{=}\PY{+w}{ }\PY{n}{scope}\PY{p}{.}\PY{n+na}{fork}\PY{p}{(}\PY{p}{(}\PY{p}{)}\PY{+w}{ }\PY{o}{\PYZhy{}}\PY{o}{\PYZgt{}}\PY{+w}{ }\PY{n}{orderService}\PY{p}{.}\PY{n+na}{recentOrders}\PY{p}{(}\PY{n}{userId}\PY{p}{)}\PY{p}{)}\PY{p}{;}

\PY{+w}{        }\PY{n}{scope}\PY{p}{.}\PY{n+na}{joinUntil}\PY{p}{(}\PY{n}{Instant}\PY{p}{.}\PY{n+na}{now}\PY{p}{(}\PY{p}{)}\PY{p}{.}\PY{n+na}{plusSeconds}\PY{p}{(}\PY{l+m+mi}{2}\PY{p}{)}\PY{p}{)}\PY{p}{;}\PY{+w}{ }\PY{c+c1}{// give downstreams 2s total}
\PY{+w}{        }\PY{n}{scope}\PY{p}{.}\PY{n+na}{throwIfFailed}\PY{p}{(}\PY{p}{)}\PY{p}{;}

\PY{+w}{        }\PY{k}{return}\PY{+w}{ }\PY{k}{new}\PY{+w}{ }\PY{n}{DashboardData}\PY{p}{(}\PY{n}{userTask}\PY{p}{.}\PY{n+na}{get}\PY{p}{(}\PY{p}{)}\PY{p}{,}\PY{+w}{ }\PY{n}{ordersTask}\PY{p}{.}\PY{n+na}{get}\PY{p}{(}\PY{p}{)}\PY{p}{,}\PY{+w}{ }\PY{k+kc}{null}\PY{p}{)}\PY{p}{;}
\PY{+w}{    }\PY{p}{\PYZcb{}}\PY{+w}{ }\PY{c+c1}{// close() cancels anything still running if joinUntil timed out}
\PY{p}{\PYZcb{}}
\end{Verbatim}
\end{codebox}

\subsection[Custom Scope Policies]{Custom Scope Policies}
\label{h:15:custom-scope-policies}
Both \code{ShutdownOnFailure} and \code{ShutdownOnSuccess} are built on the same extensible mechanism (overriding \code{handleComplete} in earlier preview shapes, or supplying a custom \code{Joiner} in the JDK 25 shape — see below). This lets you express policies neither built-in covers, such as “succeed once at least N of M subtasks succeed” (a quorum) or “collect all results, ignoring individual failures, and report which ones failed at the end” instead of failing fast.

\begin{codebox}
\begin{Verbatim}[commandchars=\\\{\}]
\PY{c+c1}{// Sketch of a quorum policy (conceptual; exact base API differs by JDK version) —}
\PY{c+c1}{// wait for at least 2 of 3 replica writes to succeed before returning.}
\PY{k}{try}\PY{+w}{ }\PY{p}{(}\PY{k+kd}{var}\PY{+w}{ }\PY{n}{scope}\PY{+w}{ }\PY{o}{=}\PY{+w}{ }\PY{k}{new}\PY{+w}{ }\PY{n}{QuorumScope}\PY{o}{\PYZlt{}}\PY{n}{Boolean}\PY{o}{\PYZgt{}}\PY{p}{(}\PY{c+cm}{/* required successes */}\PY{+w}{ }\PY{l+m+mi}{2}\PY{p}{)}\PY{p}{)}\PY{+w}{ }\PY{p}{\PYZob{}}
\PY{+w}{    }\PY{n}{scope}\PY{p}{.}\PY{n+na}{fork}\PY{p}{(}\PY{p}{(}\PY{p}{)}\PY{+w}{ }\PY{o}{\PYZhy{}}\PY{o}{\PYZgt{}}\PY{+w}{ }\PY{n}{replicaA}\PY{p}{.}\PY{n+na}{write}\PY{p}{(}\PY{n}{record}\PY{p}{)}\PY{p}{)}\PY{p}{;}
\PY{+w}{    }\PY{n}{scope}\PY{p}{.}\PY{n+na}{fork}\PY{p}{(}\PY{p}{(}\PY{p}{)}\PY{+w}{ }\PY{o}{\PYZhy{}}\PY{o}{\PYZgt{}}\PY{+w}{ }\PY{n}{replicaB}\PY{p}{.}\PY{n+na}{write}\PY{p}{(}\PY{n}{record}\PY{p}{)}\PY{p}{)}\PY{p}{;}
\PY{+w}{    }\PY{n}{scope}\PY{p}{.}\PY{n+na}{fork}\PY{p}{(}\PY{p}{(}\PY{p}{)}\PY{+w}{ }\PY{o}{\PYZhy{}}\PY{o}{\PYZgt{}}\PY{+w}{ }\PY{n}{replicaC}\PY{p}{.}\PY{n+na}{write}\PY{p}{(}\PY{n}{record}\PY{p}{)}\PY{p}{)}\PY{p}{;}

\PY{+w}{    }\PY{n}{scope}\PY{p}{.}\PY{n+na}{join}\PY{p}{(}\PY{p}{)}\PY{p}{;}
\PY{+w}{    }\PY{k}{if}\PY{+w}{ }\PY{p}{(}\PY{o}{!}\PY{n}{scope}\PY{p}{.}\PY{n+na}{quorumReached}\PY{p}{(}\PY{p}{)}\PY{p}{)}\PY{+w}{ }\PY{p}{\PYZob{}}
\PY{+w}{        }\PY{k}{throw}\PY{+w}{ }\PY{k}{new}\PY{+w}{ }\PY{n}{WriteFailedException}\PY{p}{(}\PY{l+s}{\PYZdq{}}\PY{l+s}{Could not reach write quorum for }\PY{l+s}{\PYZdq{}}\PY{+w}{ }\PY{o}{+}\PY{+w}{ }\PY{n}{record}\PY{p}{.}\PY{n+na}{id}\PY{p}{(}\PY{p}{)}\PY{p}{)}\PY{p}{;}
\PY{+w}{    }\PY{p}{\PYZcb{}}
\PY{p}{\PYZcb{}}
\end{Verbatim}
\end{codebox}

The exact base class/interface for building custom policies changed between the preview iterations and the JDK 25 API — treat the sketch above as illustrating the \emph{idea} (a policy decides, as subtasks complete, whether the scope should shut down early and what the overall outcome is), not as a copy-paste API.

\subsection[Preview Status and the JDK 25 API Reshape]{Preview Status and the JDK 25 API Reshape}
\label{h:15:preview-status-and-the-jdk-25-api-reshape}
Structured concurrency is one of the more actively evolving Loom features, and its exact shape has changed across releases — this is a real trap for anyone reciting API details from memory in an interview:

\begin{itemize}
\item 
\textbf{JDK 21}: \code{StructuredTaskScope} introduced as a \textbf{preview API} (JEP 453), requiring \code{--\allowbreak{}enable-\allowbreak{}preview} at both compile and run time. Shape: \code{new~\allowbreak{}Structured\allowbreak{}Task\allowbreak{}Scope.\allowbreak{}Shutdown\allowbreak{}On\allowbreak{}Failure()} / \code{ShutdownOnSuccess\textless{}\allowbreak{}T\textgreater{}()}, \code{fork}, \code{join}, \code{throwIfFailed}, \code{joinUntil} — as shown in the examples above.

\item 
\textbf{JDK 22-24}: Continued as a \textbf{preview API} (re-previewed, JEP 480/499), with refinements but the same overall \code{ShutdownOnFailure}/\code{ShutdownOnSuccess} subclassing shape. Still requires \code{--\allowbreak{}enable-\allowbreak{}preview}.

\item 
\textbf{JDK 25}: The API was \textbf{reshaped} around an explicit \textbf{\code{Joiner}} abstraction (finalized under JEP 505 as a preview-to-standard evolution track; check the exact JEP number against your JDK’s release notes, as this has moved during the feature’s preview lifetime). Instead of subclassing \code{ShutdownOnFailure}/\code{ShutdownOnSuccess}, you open a scope with a joiner strategy:

\end{itemize}

\begin{codebox}
\begin{Verbatim}[commandchars=\\\{\}]
\PY{c+c1}{// JDK 25\PYZhy{}style shape (illustrative) — StructuredTaskScope.open(...) with a Joiner}
\PY{k}{try}\PY{+w}{ }\PY{p}{(}\PY{k+kd}{var}\PY{+w}{ }\PY{n}{scope}\PY{+w}{ }\PY{o}{=}\PY{+w}{ }\PY{n}{StructuredTaskScope}\PY{p}{.}\PY{n+na}{open}\PY{p}{(}\PY{n}{Joiner}\PY{p}{.}\PY{o}{\PYZlt{}}\PY{n}{String}\PY{o}{\PYZgt{}}\PY{n}{awaitAllSuccessfulOrThrow}\PY{p}{(}\PY{p}{)}\PY{p}{)}\PY{p}{)}\PY{+w}{ }\PY{p}{\PYZob{}}
\PY{+w}{    }\PY{n}{scope}\PY{p}{.}\PY{n+na}{fork}\PY{p}{(}\PY{p}{(}\PY{p}{)}\PY{+w}{ }\PY{o}{\PYZhy{}}\PY{o}{\PYZgt{}}\PY{+w}{ }\PY{n}{userService}\PY{p}{.}\PY{n+na}{fetch}\PY{p}{(}\PY{n}{userId}\PY{p}{)}\PY{p}{)}\PY{p}{;}
\PY{+w}{    }\PY{n}{scope}\PY{p}{.}\PY{n+na}{fork}\PY{p}{(}\PY{p}{(}\PY{p}{)}\PY{+w}{ }\PY{o}{\PYZhy{}}\PY{o}{\PYZgt{}}\PY{+w}{ }\PY{n}{orderService}\PY{p}{.}\PY{n+na}{recentOrders}\PY{p}{(}\PY{n}{userId}\PY{p}{)}\PY{p}{)}\PY{p}{;}
\PY{+w}{    }\PY{n}{List}\PY{o}{\PYZlt{}}\PY{n}{Subtask}\PY{o}{\PYZlt{}}\PY{n}{Object}\PY{o}{\PYZgt{}}\PY{o}{\PYZgt{}}\PY{+w}{ }\PY{n}{results}\PY{+w}{ }\PY{o}{=}\PY{+w}{ }\PY{n}{scope}\PY{p}{.}\PY{n+na}{join}\PY{p}{(}\PY{p}{)}\PY{p}{;}
\PY{+w}{    }\PY{c+c1}{// ...}
\PY{p}{\PYZcb{}}

\PY{c+c1}{// First\PYZhy{}success equivalent to the old ShutdownOnSuccess}
\PY{k}{try}\PY{+w}{ }\PY{p}{(}\PY{k+kd}{var}\PY{+w}{ }\PY{n}{scope}\PY{+w}{ }\PY{o}{=}\PY{+w}{ }\PY{n}{StructuredTaskScope}\PY{p}{.}\PY{n+na}{open}\PY{p}{(}\PY{n}{Joiner}\PY{p}{.}\PY{o}{\PYZlt{}}\PY{n}{String}\PY{o}{\PYZgt{}}\PY{n}{anySuccessfulResultOrThrow}\PY{p}{(}\PY{p}{)}\PY{p}{)}\PY{p}{)}\PY{+w}{ }\PY{p}{\PYZob{}}
\PY{+w}{    }\PY{n}{scope}\PY{p}{.}\PY{n+na}{fork}\PY{p}{(}\PY{p}{(}\PY{p}{)}\PY{+w}{ }\PY{o}{\PYZhy{}}\PY{o}{\PYZgt{}}\PY{+w}{ }\PY{n}{primaryRegionCache}\PY{p}{.}\PY{n+na}{get}\PY{p}{(}\PY{n}{key}\PY{p}{)}\PY{p}{)}\PY{p}{;}
\PY{+w}{    }\PY{n}{scope}\PY{p}{.}\PY{n+na}{fork}\PY{p}{(}\PY{p}{(}\PY{p}{)}\PY{+w}{ }\PY{o}{\PYZhy{}}\PY{o}{\PYZgt{}}\PY{+w}{ }\PY{n}{secondaryRegionCache}\PY{p}{.}\PY{n+na}{get}\PY{p}{(}\PY{n}{key}\PY{p}{)}\PY{p}{)}\PY{p}{;}
\PY{+w}{    }\PY{n}{String}\PY{+w}{ }\PY{n}{result}\PY{+w}{ }\PY{o}{=}\PY{+w}{ }\PY{n}{scope}\PY{p}{.}\PY{n+na}{join}\PY{p}{(}\PY{p}{)}\PY{p}{;}
\PY{p}{\PYZcb{}}
\end{Verbatim}
\end{codebox}

The \code{Joiner} shape separates “how subtasks are combined/what triggers shutdown” (the \code{Joiner}) from “the scope mechanics” (\code{open}, \code{fork}, \code{join}), making custom policies a matter of implementing \code{Joiner} rather than subclassing \code{StructuredTaskScope}. \textbf{As of JDK 25, structured concurrency is still a preview feature} — \code{--\allowbreak{}enable-\allowbreak{}preview} is required to compile and run code using it, on every JDK version mentioned above, including 25. Do not present it as final/GA without checking the release notes of whatever JDK version is actually in play; at the time of writing it has not exited preview.

\begin{codebox}
\begin{Verbatim}[commandchars=\\\{\}]
\PY{c+c1}{\PYZsh{} Compiling and running structured concurrency code on JDK 21 or JDK 25}
javac\PY{+w}{ }\PYZhy{}\PYZhy{}release\PY{+w}{ }\PY{l+m}{21}\PY{+w}{ }\PYZhy{}\PYZhy{}enable\PYZhy{}preview\PY{+w}{ }Main.java
java\PY{+w}{ }\PYZhy{}\PYZhy{}enable\PYZhy{}preview\PY{+w}{ }Main
\end{Verbatim}
\end{codebox}

\section[Scoped Values]{Scoped Values}
\label{h:15:scoped-values}
\subsection[Why ThreadLocal Is a Poor Fit for Virtual Threads]{Why \code{ThreadLocal} Is a Poor Fit for Virtual Threads}
\label{h:15:why-threadlocal-is-a-poor-fit-for-virtual-threads}
\code{ThreadLocal} gives each thread its own independent copy of a variable, which is a fine model when threads are few and long-lived. Three specific properties make it a poor fit once threads are cheap, numerous, and short-lived (i.e., virtual threads):

\begin{enumerate}
\item 
\textbf{Mutability.} \code{ThreadLocal.\allowbreak{}set(...)} can be called at any time, by any code that has a reference to the \code{ThreadLocal}, changing the value for the rest of the thread’s lifetime (or until the next \code{set}/\code{remove}). This makes it hard to reason about what value is “in scope” at a given point — any called method could have mutated it. It also opens the door to values leaking between logical tasks if a thread is reused (a classic thread-pool bug) or forgotten to be cleared.

\item 
\textbf{Unbounded lifetime.} Nothing forces a \code{ThreadLocal} value to be cleared when the logical operation that set it is done; you must remember to call \code{remove()}, typically in a \code{finally} block, or the value (and whatever memory it retains) lives until the thread itself dies or is reused. Since there is no compiler or runtime enforcement, forgetting to clean up is a common, hard-to-spot leak.

\item 
\textbf{Inheritance cost.} \code{InheritableThreadLocal} copies parent values into every child thread at creation time. With platform thread pools this happens rarely (pools don’t spawn new threads often). With virtual threads — where creating one per task, and forking many children per \code{StructuredTaskScope}, is the \emph{normal} pattern — this copying happens constantly, at a scale (potentially millions of threads) where even a cheap copy operation adds up, and where most children may never actually read the inherited value.

\end{enumerate}

\subsection[ScopedValue — Immutable, Bounded, Structured]{\code{ScopedValue} — Immutable, Bounded, Structured}
\label{h:15:scopedvalue--immutable-bounded-structured}
\code{ScopedValue\textless{}\allowbreak{}T\textgreater{}} (introduced as a preview feature, JEP 429 in JDK 21, continuing through subsequent JDKs as a preview feature with refinements — check \code{--\allowbreak{}enable-\allowbreak{}preview} requirements per JDK) addresses all three issues by design:

\begin{itemize}
\item 
\textbf{Immutable within a scope.} Once bound via \code{where(...).\allowbreak{}run(...)}, the value cannot be changed from inside that block — there is no \code{set()} method. Rebinding is only possible by opening a \emph{new}, nested \code{where(...)} scope (see below), which is visible and structured, not a silent mutation.

\item 
\textbf{Bounded lifetime.} The binding exists only for the duration of the \code{run}/\code{call} invocation — it is automatically and unconditionally torn down when that method returns, including via exception. There’s no \code{remove()} to forget.

\item 
\textbf{Cheap, structured inheritance.} When a \code{ScopedValue}-bound thread forks children through a \code{StructuredTaskScope}, the binding is available to those children without a per-child copy — the runtime shares the immutable binding rather than duplicating it, which is both cheaper and safer at virtual-thread scale.

\end{itemize}

\begin{codebox}
\begin{Verbatim}[commandchars=\\\{\}]
\PY{k+kd}{public}\PY{+w}{ }\PY{k+kd}{class} \PY{n+nc}{RequestContext}\PY{+w}{ }\PY{p}{\PYZob{}}
\PY{+w}{    }\PY{k+kd}{private}\PY{+w}{ }\PY{k+kd}{static}\PY{+w}{ }\PY{k+kd}{final}\PY{+w}{ }\PY{n}{ScopedValue}\PY{o}{\PYZlt{}}\PY{n}{String}\PY{o}{\PYZgt{}}\PY{+w}{ }\PY{n}{REQUEST\PYZus{}ID}\PY{+w}{ }\PY{o}{=}\PY{+w}{ }\PY{n}{ScopedValue}\PY{p}{.}\PY{n+na}{newInstance}\PY{p}{(}\PY{p}{)}\PY{p}{;}

\PY{+w}{    }\PY{k+kd}{public}\PY{+w}{ }\PY{k+kd}{static}\PY{+w}{ }\PY{k+kt}{void}\PY{+w}{ }\PY{n+nf}{handleRequest}\PY{p}{(}\PY{n}{String}\PY{+w}{ }\PY{n}{requestId}\PY{p}{,}\PY{+w}{ }\PY{n}{HttpExchange}\PY{+w}{ }\PY{n}{exchange}\PY{p}{)}\PY{+w}{ }\PY{p}{\PYZob{}}
\PY{+w}{        }\PY{n}{ScopedValue}\PY{p}{.}\PY{n+na}{where}\PY{p}{(}\PY{n}{REQUEST\PYZus{}ID}\PY{p}{,}\PY{+w}{ }\PY{n}{requestId}\PY{p}{)}
\PY{+w}{                   }\PY{p}{.}\PY{n+na}{run}\PY{p}{(}\PY{p}{(}\PY{p}{)}\PY{+w}{ }\PY{o}{\PYZhy{}}\PY{o}{\PYZgt{}}\PY{+w}{ }\PY{n}{processRequest}\PY{p}{(}\PY{n}{exchange}\PY{p}{)}\PY{p}{)}\PY{p}{;}
\PY{+w}{        }\PY{c+c1}{// REQUEST\PYZus{}ID is unbound again as soon as run() returns — nothing to clean up}
\PY{+w}{    }\PY{p}{\PYZcb{}}

\PY{+w}{    }\PY{k+kd}{static}\PY{+w}{ }\PY{k+kt}{void}\PY{+w}{ }\PY{n+nf}{processRequest}\PY{p}{(}\PY{n}{HttpExchange}\PY{+w}{ }\PY{n}{exchange}\PY{p}{)}\PY{+w}{ }\PY{p}{\PYZob{}}
\PY{+w}{        }\PY{n}{log}\PY{p}{.}\PY{n+na}{info}\PY{p}{(}\PY{l+s}{\PYZdq{}}\PY{l+s}{Handling request \PYZob{}\PYZcb{}}\PY{l+s}{\PYZdq{}}\PY{p}{,}\PY{+w}{ }\PY{n}{REQUEST\PYZus{}ID}\PY{p}{.}\PY{n+na}{get}\PY{p}{(}\PY{p}{)}\PY{p}{)}\PY{p}{;}\PY{+w}{ }\PY{c+c1}{// reads the bound value}
\PY{+w}{        }\PY{n}{callDownstream}\PY{p}{(}\PY{p}{)}\PY{p}{;}\PY{+w}{ }\PY{c+c1}{// REQUEST\PYZus{}ID stays bound for the whole call tree beneath run()}
\PY{+w}{    }\PY{p}{\PYZcb{}}

\PY{+w}{    }\PY{k+kd}{static}\PY{+w}{ }\PY{k+kt}{void}\PY{+w}{ }\PY{n+nf}{callDownstream}\PY{p}{(}\PY{p}{)}\PY{+w}{ }\PY{p}{\PYZob{}}
\PY{+w}{        }\PY{n}{log}\PY{p}{.}\PY{n+na}{info}\PY{p}{(}\PY{l+s}{\PYZdq{}}\PY{l+s}{Still request \PYZob{}\PYZcb{}}\PY{l+s}{\PYZdq{}}\PY{p}{,}\PY{+w}{ }\PY{n}{REQUEST\PYZus{}ID}\PY{p}{.}\PY{n+na}{get}\PY{p}{(}\PY{p}{)}\PY{p}{)}\PY{p}{;}\PY{+w}{ }\PY{c+c1}{// visible several frames down}
\PY{+w}{    }\PY{p}{\PYZcb{}}
\PY{p}{\PYZcb{}}
\end{Verbatim}
\end{codebox}

\code{where(...).\allowbreak{}call(...)} is the equivalent for code that returns a value or throws a checked exception:

\begin{codebox}
\begin{Verbatim}[commandchars=\\\{\}]
\PY{n}{String}\PY{+w}{ }\PY{n}{result}\PY{+w}{ }\PY{o}{=}\PY{+w}{ }\PY{n}{ScopedValue}\PY{p}{.}\PY{n+na}{where}\PY{p}{(}\PY{n}{REQUEST\PYZus{}ID}\PY{p}{,}\PY{+w}{ }\PY{n}{requestId}\PY{p}{)}
\PY{+w}{                            }\PY{p}{.}\PY{n+na}{call}\PY{p}{(}\PY{p}{(}\PY{p}{)}\PY{+w}{ }\PY{o}{\PYZhy{}}\PY{o}{\PYZgt{}}\PY{+w}{ }\PY{n}{fetchAndFormat}\PY{p}{(}\PY{n}{exchange}\PY{p}{)}\PY{p}{)}\PY{p}{;}
\end{Verbatim}
\end{codebox}

\subsection[Rebinding]{Rebinding}
\label{h:15:rebinding}
Because bindings are scoped to a block, “changing” a value for a nested piece of work means opening a new \code{where(...)} inside the existing one — the outer binding is restored automatically once the inner block exits:

\begin{codebox}
\begin{Verbatim}[commandchars=\\\{\}]
\PY{n}{ScopedValue}\PY{p}{.}\PY{n+na}{where}\PY{p}{(}\PY{n}{REQUEST\PYZus{}ID}\PY{p}{,}\PY{+w}{ }\PY{l+s}{\PYZdq{}}\PY{l+s}{req\PYZhy{}1}\PY{l+s}{\PYZdq{}}\PY{p}{)}\PY{p}{.}\PY{n+na}{run}\PY{p}{(}\PY{p}{(}\PY{p}{)}\PY{+w}{ }\PY{o}{\PYZhy{}}\PY{o}{\PYZgt{}}\PY{+w}{ }\PY{p}{\PYZob{}}
\PY{+w}{    }\PY{n}{log}\PY{p}{.}\PY{n+na}{info}\PY{p}{(}\PY{n}{REQUEST\PYZus{}ID}\PY{p}{.}\PY{n+na}{get}\PY{p}{(}\PY{p}{)}\PY{p}{)}\PY{p}{;}\PY{+w}{ }\PY{c+c1}{// \PYZdq{}req\PYZhy{}1\PYZdq{}}

\PY{+w}{    }\PY{n}{ScopedValue}\PY{p}{.}\PY{n+na}{where}\PY{p}{(}\PY{n}{REQUEST\PYZus{}ID}\PY{p}{,}\PY{+w}{ }\PY{l+s}{\PYZdq{}}\PY{l+s}{req\PYZhy{}1\PYZhy{}retry}\PY{l+s}{\PYZdq{}}\PY{p}{)}\PY{p}{.}\PY{n+na}{run}\PY{p}{(}\PY{p}{(}\PY{p}{)}\PY{+w}{ }\PY{o}{\PYZhy{}}\PY{o}{\PYZgt{}}\PY{+w}{ }\PY{p}{\PYZob{}}
\PY{+w}{        }\PY{n}{log}\PY{p}{.}\PY{n+na}{info}\PY{p}{(}\PY{n}{REQUEST\PYZus{}ID}\PY{p}{.}\PY{n+na}{get}\PY{p}{(}\PY{p}{)}\PY{p}{)}\PY{p}{;}\PY{+w}{ }\PY{c+c1}{// \PYZdq{}req\PYZhy{}1\PYZhy{}retry\PYZdq{} — rebound for this inner block only}
\PY{+w}{    }\PY{p}{\PYZcb{}}\PY{p}{)}\PY{p}{;}

\PY{+w}{    }\PY{n}{log}\PY{p}{.}\PY{n+na}{info}\PY{p}{(}\PY{n}{REQUEST\PYZus{}ID}\PY{p}{.}\PY{n+na}{get}\PY{p}{(}\PY{p}{)}\PY{p}{)}\PY{p}{;}\PY{+w}{ }\PY{c+c1}{// back to \PYZdq{}req\PYZhy{}1\PYZdq{} — outer binding restored}
\PY{p}{\PYZcb{}}\PY{p}{)}\PY{p}{;}
\end{Verbatim}
\end{codebox}

\subsection[Inheritance into StructuredTaskScope Forks]{Inheritance into \code{StructuredTaskScope} Forks}
\label{h:15:inheritance-into-structuredtaskscope-forks}
\code{ScopedValue} bindings established before a \code{StructuredTaskScope} is opened are visible inside subtasks forked from that scope, without any explicit copying code — this is the natural pairing of the two features: bind request-scoped context once, fan out safely, and every forked subtask (and anything it calls) can read that context.

\begin{codebox}
\begin{Verbatim}[commandchars=\\\{\}]
\PY{n}{ScopedValue}\PY{o}{\PYZlt{}}\PY{n}{String}\PY{o}{\PYZgt{}}\PY{+w}{ }\PY{n}{REQUEST\PYZus{}ID}\PY{+w}{ }\PY{o}{=}\PY{+w}{ }\PY{n}{ScopedValue}\PY{p}{.}\PY{n+na}{newInstance}\PY{p}{(}\PY{p}{)}\PY{p}{;}

\PY{n}{DashboardData}\PY{+w}{ }\PY{n+nf}{loadDashboard}\PY{p}{(}\PY{n}{String}\PY{+w}{ }\PY{n}{userId}\PY{p}{,}\PY{+w}{ }\PY{n}{String}\PY{+w}{ }\PY{n}{requestId}\PY{p}{)}
\PY{+w}{        }\PY{k+kd}{throws}\PY{+w}{ }\PY{n}{InterruptedException}\PY{p}{,}\PY{+w}{ }\PY{n}{ExecutionException}\PY{+w}{ }\PY{p}{\PYZob{}}
\PY{+w}{    }\PY{k}{return}\PY{+w}{ }\PY{n}{ScopedValue}\PY{p}{.}\PY{n+na}{where}\PY{p}{(}\PY{n}{REQUEST\PYZus{}ID}\PY{p}{,}\PY{+w}{ }\PY{n}{requestId}\PY{p}{)}\PY{p}{.}\PY{n+na}{call}\PY{p}{(}\PY{p}{(}\PY{p}{)}\PY{+w}{ }\PY{o}{\PYZhy{}}\PY{o}{\PYZgt{}}\PY{+w}{ }\PY{p}{\PYZob{}}
\PY{+w}{        }\PY{k}{try}\PY{+w}{ }\PY{p}{(}\PY{k+kd}{var}\PY{+w}{ }\PY{n}{scope}\PY{+w}{ }\PY{o}{=}\PY{+w}{ }\PY{k}{new}\PY{+w}{ }\PY{n}{StructuredTaskScope}\PY{p}{.}\PY{n+na}{ShutdownOnFailure}\PY{p}{(}\PY{p}{)}\PY{p}{)}\PY{+w}{ }\PY{p}{\PYZob{}}
\PY{+w}{            }\PY{c+c1}{// Both forked subtasks can call REQUEST\PYZus{}ID.get() and see \PYZdq{}requestId\PYZdq{},}
\PY{+w}{            }\PY{c+c1}{// with no manual propagation and no InheritableThreadLocal copy cost.}
\PY{+w}{            }\PY{n}{Subtask}\PY{o}{\PYZlt{}}\PY{n}{User}\PY{o}{\PYZgt{}}\PY{+w}{ }\PY{n}{userTask}\PY{+w}{ }\PY{o}{=}\PY{+w}{ }\PY{n}{scope}\PY{p}{.}\PY{n+na}{fork}\PY{p}{(}\PY{p}{(}\PY{p}{)}\PY{+w}{ }\PY{o}{\PYZhy{}}\PY{o}{\PYZgt{}}\PY{+w}{ }\PY{p}{\PYZob{}}
\PY{+w}{                }\PY{n}{log}\PY{p}{.}\PY{n+na}{info}\PY{p}{(}\PY{l+s}{\PYZdq{}}\PY{l+s}{[\PYZob{}\PYZcb{}] fetching user}\PY{l+s}{\PYZdq{}}\PY{p}{,}\PY{+w}{ }\PY{n}{REQUEST\PYZus{}ID}\PY{p}{.}\PY{n+na}{get}\PY{p}{(}\PY{p}{)}\PY{p}{)}\PY{p}{;}
\PY{+w}{                }\PY{k}{return}\PY{+w}{ }\PY{n}{userService}\PY{p}{.}\PY{n+na}{fetch}\PY{p}{(}\PY{n}{userId}\PY{p}{)}\PY{p}{;}
\PY{+w}{            }\PY{p}{\PYZcb{}}\PY{p}{)}\PY{p}{;}
\PY{+w}{            }\PY{n}{Subtask}\PY{o}{\PYZlt{}}\PY{n}{List}\PY{o}{\PYZlt{}}\PY{n}{Order}\PY{o}{\PYZgt{}}\PY{o}{\PYZgt{}}\PY{+w}{ }\PY{n}{ordersTask}\PY{+w}{ }\PY{o}{=}\PY{+w}{ }\PY{n}{scope}\PY{p}{.}\PY{n+na}{fork}\PY{p}{(}\PY{p}{(}\PY{p}{)}\PY{+w}{ }\PY{o}{\PYZhy{}}\PY{o}{\PYZgt{}}\PY{+w}{ }\PY{p}{\PYZob{}}
\PY{+w}{                }\PY{n}{log}\PY{p}{.}\PY{n+na}{info}\PY{p}{(}\PY{l+s}{\PYZdq{}}\PY{l+s}{[\PYZob{}\PYZcb{}] fetching orders}\PY{l+s}{\PYZdq{}}\PY{p}{,}\PY{+w}{ }\PY{n}{REQUEST\PYZus{}ID}\PY{p}{.}\PY{n+na}{get}\PY{p}{(}\PY{p}{)}\PY{p}{)}\PY{p}{;}
\PY{+w}{                }\PY{k}{return}\PY{+w}{ }\PY{n}{orderService}\PY{p}{.}\PY{n+na}{recentOrders}\PY{p}{(}\PY{n}{userId}\PY{p}{)}\PY{p}{;}
\PY{+w}{            }\PY{p}{\PYZcb{}}\PY{p}{)}\PY{p}{;}

\PY{+w}{            }\PY{n}{scope}\PY{p}{.}\PY{n+na}{join}\PY{p}{(}\PY{p}{)}\PY{p}{;}
\PY{+w}{            }\PY{n}{scope}\PY{p}{.}\PY{n+na}{throwIfFailed}\PY{p}{(}\PY{p}{)}\PY{p}{;}
\PY{+w}{            }\PY{k}{return}\PY{+w}{ }\PY{k}{new}\PY{+w}{ }\PY{n}{DashboardData}\PY{p}{(}\PY{n}{userTask}\PY{p}{.}\PY{n+na}{get}\PY{p}{(}\PY{p}{)}\PY{p}{,}\PY{+w}{ }\PY{n}{ordersTask}\PY{p}{.}\PY{n+na}{get}\PY{p}{(}\PY{p}{)}\PY{p}{,}\PY{+w}{ }\PY{k+kc}{null}\PY{p}{)}\PY{p}{;}
\PY{+w}{        }\PY{p}{\PYZcb{}}
\PY{+w}{    }\PY{p}{\PYZcb{}}\PY{p}{)}\PY{p}{;}
\PY{p}{\PYZcb{}}
\end{Verbatim}
\end{codebox}

\subsection[Preview Status]{Preview Status}
\label{h:15:preview-status}
\code{ScopedValue} has been a \textbf{preview API since JDK 21} (JEP 429) and remains preview through the JDK versions covered by this chapter, refined along the way (JEP 481/506 in later releases — check the exact JEP number for the JDK you are targeting, since, like structured concurrency, this feature’s finalization timeline has shifted release to release). \code{--\allowbreak{}enable-\allowbreak{}preview} is required to compile and run code that uses it, on every affected JDK version. Do not assume it is final without checking your specific JDK’s release notes.

\subsection[Before/After: Migrating from ThreadLocal]{Before/After: Migrating from \code{ThreadLocal}}
\label{h:15:beforeafter-migrating-from-threadlocal}
\begin{codebox}
\begin{Verbatim}[commandchars=\\\{\}]
\PY{c+c1}{// BEFORE — ThreadLocal: mutable, must remember to clean up, costly to inherit}
\PY{c+c1}{// into children at virtual\PYZhy{}thread scale.}
\PY{k+kd}{public}\PY{+w}{ }\PY{k+kd}{class} \PY{n+nc}{TenantContext}\PY{+w}{ }\PY{p}{\PYZob{}}
\PY{+w}{    }\PY{k+kd}{private}\PY{+w}{ }\PY{k+kd}{static}\PY{+w}{ }\PY{k+kd}{final}\PY{+w}{ }\PY{n}{ThreadLocal}\PY{o}{\PYZlt{}}\PY{n}{String}\PY{o}{\PYZgt{}}\PY{+w}{ }\PY{n}{TENANT\PYZus{}ID}\PY{+w}{ }\PY{o}{=}\PY{+w}{ }\PY{k}{new}\PY{+w}{ }\PY{n}{ThreadLocal}\PY{o}{\PYZlt{}}\PY{o}{\PYZgt{}}\PY{p}{(}\PY{p}{)}\PY{p}{;}

\PY{+w}{    }\PY{k+kd}{public}\PY{+w}{ }\PY{k+kd}{static}\PY{+w}{ }\PY{k+kt}{void}\PY{+w}{ }\PY{n+nf}{handle}\PY{p}{(}\PY{n}{String}\PY{+w}{ }\PY{n}{tenantId}\PY{p}{,}\PY{+w}{ }\PY{n}{Runnable}\PY{+w}{ }\PY{n}{task}\PY{p}{)}\PY{+w}{ }\PY{p}{\PYZob{}}
\PY{+w}{        }\PY{n}{TENANT\PYZus{}ID}\PY{p}{.}\PY{n+na}{set}\PY{p}{(}\PY{n}{tenantId}\PY{p}{)}\PY{p}{;}
\PY{+w}{        }\PY{k}{try}\PY{+w}{ }\PY{p}{\PYZob{}}
\PY{+w}{            }\PY{n}{task}\PY{p}{.}\PY{n+na}{run}\PY{p}{(}\PY{p}{)}\PY{p}{;}
\PY{+w}{        }\PY{p}{\PYZcb{}}\PY{+w}{ }\PY{k}{finally}\PY{+w}{ }\PY{p}{\PYZob{}}
\PY{+w}{            }\PY{n}{TENANT\PYZus{}ID}\PY{p}{.}\PY{n+na}{remove}\PY{p}{(}\PY{p}{)}\PY{p}{;}\PY{+w}{ }\PY{c+c1}{// easy to forget; leaks the value if omitted}
\PY{+w}{        }\PY{p}{\PYZcb{}}
\PY{+w}{    }\PY{p}{\PYZcb{}}

\PY{+w}{    }\PY{k+kd}{public}\PY{+w}{ }\PY{k+kd}{static}\PY{+w}{ }\PY{n}{String}\PY{+w}{ }\PY{n+nf}{currentTenant}\PY{p}{(}\PY{p}{)}\PY{+w}{ }\PY{p}{\PYZob{}}
\PY{+w}{        }\PY{k}{return}\PY{+w}{ }\PY{n}{TENANT\PYZus{}ID}\PY{p}{.}\PY{n+na}{get}\PY{p}{(}\PY{p}{)}\PY{p}{;}
\PY{+w}{    }\PY{p}{\PYZcb{}}
\PY{p}{\PYZcb{}}

\PY{c+c1}{// AFTER — ScopedValue: immutable for the duration of run(), automatically}
\PY{c+c1}{// unbound on exit (even via exception), cheap to share with forked subtasks.}
\PY{k+kd}{public}\PY{+w}{ }\PY{k+kd}{class} \PY{n+nc}{TenantContext}\PY{+w}{ }\PY{p}{\PYZob{}}
\PY{+w}{    }\PY{k+kd}{private}\PY{+w}{ }\PY{k+kd}{static}\PY{+w}{ }\PY{k+kd}{final}\PY{+w}{ }\PY{n}{ScopedValue}\PY{o}{\PYZlt{}}\PY{n}{String}\PY{o}{\PYZgt{}}\PY{+w}{ }\PY{n}{TENANT\PYZus{}ID}\PY{+w}{ }\PY{o}{=}\PY{+w}{ }\PY{n}{ScopedValue}\PY{p}{.}\PY{n+na}{newInstance}\PY{p}{(}\PY{p}{)}\PY{p}{;}

\PY{+w}{    }\PY{k+kd}{public}\PY{+w}{ }\PY{k+kd}{static}\PY{+w}{ }\PY{k+kt}{void}\PY{+w}{ }\PY{n+nf}{handle}\PY{p}{(}\PY{n}{String}\PY{+w}{ }\PY{n}{tenantId}\PY{p}{,}\PY{+w}{ }\PY{n}{Runnable}\PY{+w}{ }\PY{n}{task}\PY{p}{)}\PY{+w}{ }\PY{p}{\PYZob{}}
\PY{+w}{        }\PY{n}{ScopedValue}\PY{p}{.}\PY{n+na}{where}\PY{p}{(}\PY{n}{TENANT\PYZus{}ID}\PY{p}{,}\PY{+w}{ }\PY{n}{tenantId}\PY{p}{)}\PY{p}{.}\PY{n+na}{run}\PY{p}{(}\PY{n}{task}\PY{p}{)}\PY{p}{;}
\PY{+w}{        }\PY{c+c1}{// no finally, no remove() — unbinding is automatic and guaranteed}
\PY{+w}{    }\PY{p}{\PYZcb{}}

\PY{+w}{    }\PY{k+kd}{public}\PY{+w}{ }\PY{k+kd}{static}\PY{+w}{ }\PY{n}{String}\PY{+w}{ }\PY{n+nf}{currentTenant}\PY{p}{(}\PY{p}{)}\PY{+w}{ }\PY{p}{\PYZob{}}
\PY{+w}{        }\PY{k}{return}\PY{+w}{ }\PY{n}{TENANT\PYZus{}ID}\PY{p}{.}\PY{n+na}{orElse}\PY{p}{(}\PY{l+s}{\PYZdq{}}\PY{l+s}{unknown}\PY{l+s}{\PYZdq{}}\PY{p}{)}\PY{p}{;}\PY{+w}{ }\PY{c+c1}{// safe default if unbound}
\PY{+w}{    }\PY{p}{\PYZcb{}}
\PY{p}{\PYZcb{}}
\end{Verbatim}
\end{codebox}

The \code{AFTER} version cannot leak a stale tenant into unrelated work on a reused thread (there’s no \code{set()} outside the scope, and no thread reuse concern for one-virtual-thread-per-task code), cannot be forgotten to be cleaned up (there’s no cleanup step at all), and shares its binding cheaply with any \code{StructuredTaskScope} forks started inside \code{task.\allowbreak{}run()}.

\section[Comparison Table: Platform Threads vs Virtual Threads]{Comparison Table: Platform Threads vs Virtual Threads}
\label{h:15:comparison-table-platform-threads-vs-virtual-threads}
{\tablefontsmall
\begin{xltabular}{\linewidth}{@{}L{0.418}L{1.165}L{1.417}@{}}
\toprule
\rowcolor{tablehdr}
\thd{Aspect} & \thd{Platform Thread} & \thd{Virtual Thread} \\
\midrule
\endhead
Creation cost & High — OS-level thread creation (syscall), noticeable latency & Very low — heap allocation, no OS thread created per virtual thread \\
Stack & Fixed-size native stack (commonly 512 KB-1 MB), reserved upfront & Grows/shrinks on the heap as a continuation; starts small \\
Scheduling & OS scheduler, preemptive, across all platform threads on the machine & JVM scheduler (carrier-thread pool, typically \textasciitilde{}one per core), cooperative unmount on blocking calls \\
Typical count & Hundreds to a few thousand before memory/scheduling cost becomes a problem & Hundreds of thousands to millions \\
Blocking behavior & Blocks the OS thread; that thread is unusable for anything else while blocked & Unmounts from its carrier on blocking I/O, freeing the carrier for other virtual threads (unless pinned) \\
Best use & CPU-bound work; long-lived worker/daemon threads; native-frame-heavy work & I/O-bound, high-concurrency, blocking-style code (thread-per-request servers, fan-out to downstreams) \\
Pooling & Yes — reuse is the point (\code{newFixedThreadPool}, etc.) & No — create one per task; never pool \\
\bottomrule
\end{xltabular}
}

\section[Worked Example: Fan-Out to Three Downstream Calls, Three Ways]{Worked Example: Fan-Out to Three Downstream Calls, Three Ways}
\label{h:15:worked-example-fan-out-to-three-downstream-calls-three-ways}
Scenario: a \code{ProfileService} needs to build a page that requires calling three independent downstream services — \code{UserService}, \code{OrderService}, and \code{RecommendationService} — and combining their results. All three calls are independent and can run concurrently; the page can only be built if all three succeed.

\begin{codebox}
\begin{Verbatim}[commandchars=\\\{\}]
\PY{k+kd}{record} \PY{n+nc}{Profile}\PY{p}{(}\PY{n}{User}\PY{+w}{ }\PY{n}{user}\PY{p}{,}\PY{+w}{ }\PY{n}{List}\PY{o}{\PYZlt{}}\PY{n}{Order}\PY{o}{\PYZgt{}}\PY{+w}{ }\PY{n}{orders}\PY{p}{,}\PY{+w}{ }\PY{n}{List}\PY{o}{\PYZlt{}}\PY{n}{Recommendation}\PY{o}{\PYZgt{}}\PY{+w}{ }\PY{n}{recommendations}\PY{p}{)}\PY{+w}{ }\PY{p}{\PYZob{}}\PY{p}{\PYZcb{}}

\PY{k+kd}{interface} \PY{n+nc}{UserService}\PY{+w}{ }\PY{p}{\PYZob{}}\PY{+w}{ }\PY{n}{User}\PY{+w}{ }\PY{n+nf}{fetch}\PY{p}{(}\PY{n}{String}\PY{+w}{ }\PY{n}{userId}\PY{p}{)}\PY{+w}{ }\PY{k+kd}{throws}\PY{+w}{ }\PY{n}{Exception}\PY{p}{;}\PY{+w}{ }\PY{p}{\PYZcb{}}
\PY{k+kd}{interface} \PY{n+nc}{OrderService}\PY{+w}{ }\PY{p}{\PYZob{}}\PY{+w}{ }\PY{n}{List}\PY{o}{\PYZlt{}}\PY{n}{Order}\PY{o}{\PYZgt{}}\PY{+w}{ }\PY{n+nf}{recentOrders}\PY{p}{(}\PY{n}{String}\PY{+w}{ }\PY{n}{userId}\PY{p}{)}\PY{+w}{ }\PY{k+kd}{throws}\PY{+w}{ }\PY{n}{Exception}\PY{p}{;}\PY{+w}{ }\PY{p}{\PYZcb{}}
\PY{k+kd}{interface} \PY{n+nc}{RecommendationService}\PY{+w}{ }\PY{p}{\PYZob{}}\PY{+w}{ }\PY{n}{List}\PY{o}{\PYZlt{}}\PY{n}{Recommendation}\PY{o}{\PYZgt{}}\PY{+w}{ }\PY{n+nf}{recommendationsFor}\PY{p}{(}\PY{n}{String}\PY{+w}{ }\PY{n}{userId}\PY{p}{)}\PY{+w}{ }\PY{k+kd}{throws}\PY{+w}{ }\PY{n}{Exception}\PY{p}{;}\PY{+w}{ }\PY{p}{\PYZcb{}}
\end{Verbatim}
\end{codebox}

\subsection[Version 1: Raw ExecutorService and Future]{Version 1: Raw \code{ExecutorService} and \code{Future}}
\label{h:15:version-1-raw-executorservice-and-future}
\begin{codebox}
\begin{Verbatim}[commandchars=\\\{\}]
\PY{k+kd}{class} \PY{n+nc}{ExecutorServiceProfileService}\PY{+w}{ }\PY{p}{\PYZob{}}
\PY{+w}{    }\PY{k+kd}{private}\PY{+w}{ }\PY{k+kd}{final}\PY{+w}{ }\PY{n}{ExecutorService}\PY{+w}{ }\PY{n}{executor}\PY{+w}{ }\PY{o}{=}\PY{+w}{ }\PY{n}{Executors}\PY{p}{.}\PY{n+na}{newVirtualThreadPerTaskExecutor}\PY{p}{(}\PY{p}{)}\PY{p}{;}
\PY{+w}{    }\PY{k+kd}{private}\PY{+w}{ }\PY{k+kd}{final}\PY{+w}{ }\PY{n}{UserService}\PY{+w}{ }\PY{n}{userService}\PY{p}{;}
\PY{+w}{    }\PY{k+kd}{private}\PY{+w}{ }\PY{k+kd}{final}\PY{+w}{ }\PY{n}{OrderService}\PY{+w}{ }\PY{n}{orderService}\PY{p}{;}
\PY{+w}{    }\PY{k+kd}{private}\PY{+w}{ }\PY{k+kd}{final}\PY{+w}{ }\PY{n}{RecommendationService}\PY{+w}{ }\PY{n}{recommendationService}\PY{p}{;}

\PY{+w}{    }\PY{n}{Profile}\PY{+w}{ }\PY{n+nf}{buildProfile}\PY{p}{(}\PY{n}{String}\PY{+w}{ }\PY{n}{userId}\PY{p}{)}\PY{+w}{ }\PY{k+kd}{throws}\PY{+w}{ }\PY{n}{InterruptedException}\PY{p}{,}\PY{+w}{ }\PY{n}{ExecutionException}\PY{+w}{ }\PY{p}{\PYZob{}}
\PY{+w}{        }\PY{n}{Future}\PY{o}{\PYZlt{}}\PY{n}{User}\PY{o}{\PYZgt{}}\PY{+w}{ }\PY{n}{userFuture}\PY{+w}{ }\PY{o}{=}\PY{+w}{ }\PY{n}{executor}\PY{p}{.}\PY{n+na}{submit}\PY{p}{(}\PY{p}{(}\PY{p}{)}\PY{+w}{ }\PY{o}{\PYZhy{}}\PY{o}{\PYZgt{}}\PY{+w}{ }\PY{n}{userService}\PY{p}{.}\PY{n+na}{fetch}\PY{p}{(}\PY{n}{userId}\PY{p}{)}\PY{p}{)}\PY{p}{;}
\PY{+w}{        }\PY{n}{Future}\PY{o}{\PYZlt{}}\PY{n}{List}\PY{o}{\PYZlt{}}\PY{n}{Order}\PY{o}{\PYZgt{}}\PY{o}{\PYZgt{}}\PY{+w}{ }\PY{n}{ordersFuture}\PY{+w}{ }\PY{o}{=}\PY{+w}{ }\PY{n}{executor}\PY{p}{.}\PY{n+na}{submit}\PY{p}{(}\PY{p}{(}\PY{p}{)}\PY{+w}{ }\PY{o}{\PYZhy{}}\PY{o}{\PYZgt{}}\PY{+w}{ }\PY{n}{orderService}\PY{p}{.}\PY{n+na}{recentOrders}\PY{p}{(}\PY{n}{userId}\PY{p}{)}\PY{p}{)}\PY{p}{;}
\PY{+w}{        }\PY{n}{Future}\PY{o}{\PYZlt{}}\PY{n}{List}\PY{o}{\PYZlt{}}\PY{n}{Recommendation}\PY{o}{\PYZgt{}}\PY{o}{\PYZgt{}}\PY{+w}{ }\PY{n}{recsFuture}\PY{+w}{ }\PY{o}{=}
\PY{+w}{                }\PY{n}{executor}\PY{p}{.}\PY{n+na}{submit}\PY{p}{(}\PY{p}{(}\PY{p}{)}\PY{+w}{ }\PY{o}{\PYZhy{}}\PY{o}{\PYZgt{}}\PY{+w}{ }\PY{n}{recommendationService}\PY{p}{.}\PY{n+na}{recommendationsFor}\PY{p}{(}\PY{n}{userId}\PY{p}{)}\PY{p}{)}\PY{p}{;}

\PY{+w}{        }\PY{k}{try}\PY{+w}{ }\PY{p}{\PYZob{}}
\PY{+w}{            }\PY{c+c1}{// If userFuture.get() throws, ordersFuture and recsFuture keep running}
\PY{+w}{            }\PY{c+c1}{// unattended — we never cancel them. That\PYZsq{}s the unstructured\PYZhy{}fan\PYZhy{}out}
\PY{+w}{            }\PY{c+c1}{// problem in practice: no automatic cancellation, no bounded lifetime.}
\PY{+w}{            }\PY{n}{User}\PY{+w}{ }\PY{n}{user}\PY{+w}{ }\PY{o}{=}\PY{+w}{ }\PY{n}{userFuture}\PY{p}{.}\PY{n+na}{get}\PY{p}{(}\PY{p}{)}\PY{p}{;}
\PY{+w}{            }\PY{n}{List}\PY{o}{\PYZlt{}}\PY{n}{Order}\PY{o}{\PYZgt{}}\PY{+w}{ }\PY{n}{orders}\PY{+w}{ }\PY{o}{=}\PY{+w}{ }\PY{n}{ordersFuture}\PY{p}{.}\PY{n+na}{get}\PY{p}{(}\PY{p}{)}\PY{p}{;}
\PY{+w}{            }\PY{n}{List}\PY{o}{\PYZlt{}}\PY{n}{Recommendation}\PY{o}{\PYZgt{}}\PY{+w}{ }\PY{n}{recs}\PY{+w}{ }\PY{o}{=}\PY{+w}{ }\PY{n}{recsFuture}\PY{p}{.}\PY{n+na}{get}\PY{p}{(}\PY{p}{)}\PY{p}{;}
\PY{+w}{            }\PY{k}{return}\PY{+w}{ }\PY{k}{new}\PY{+w}{ }\PY{n}{Profile}\PY{p}{(}\PY{n}{user}\PY{p}{,}\PY{+w}{ }\PY{n}{orders}\PY{p}{,}\PY{+w}{ }\PY{n}{recs}\PY{p}{)}\PY{p}{;}
\PY{+w}{        }\PY{p}{\PYZcb{}}\PY{+w}{ }\PY{k}{catch}\PY{+w}{ }\PY{p}{(}\PY{n}{ExecutionException}\PY{+w}{ }\PY{n}{e}\PY{p}{)}\PY{+w}{ }\PY{p}{\PYZob{}}
\PY{+w}{            }\PY{n}{userFuture}\PY{p}{.}\PY{n+na}{cancel}\PY{p}{(}\PY{k+kc}{true}\PY{p}{)}\PY{p}{;}\PY{+w}{   }\PY{c+c1}{// must remember to do this manually,}
\PY{+w}{            }\PY{n}{ordersFuture}\PY{p}{.}\PY{n+na}{cancel}\PY{p}{(}\PY{k+kc}{true}\PY{p}{)}\PY{p}{;}\PY{+w}{ }\PY{c+c1}{// in every catch block, for every sibling —}
\PY{+w}{            }\PY{n}{recsFuture}\PY{p}{.}\PY{n+na}{cancel}\PY{p}{(}\PY{k+kc}{true}\PY{p}{)}\PY{p}{;}\PY{+w}{   }\PY{c+c1}{// easy to omit, easy to get wrong under refactors}
\PY{+w}{            }\PY{k}{throw}\PY{+w}{ }\PY{n}{e}\PY{p}{;}
\PY{+w}{        }\PY{p}{\PYZcb{}}
\PY{+w}{    }\PY{p}{\PYZcb{}}
\PY{p}{\PYZcb{}}
\end{Verbatim}
\end{codebox}

\subsection[Version 2: CompletableFuture]{Version 2: \code{CompletableFuture}}
\label{h:15:version-2-completablefuture}
\begin{codebox}
\begin{Verbatim}[commandchars=\\\{\}]
\PY{k+kd}{class} \PY{n+nc}{CompletableFutureProfileService}\PY{+w}{ }\PY{p}{\PYZob{}}
\PY{+w}{    }\PY{k+kd}{private}\PY{+w}{ }\PY{k+kd}{final}\PY{+w}{ }\PY{n}{Executor}\PY{+w}{ }\PY{n}{executor}\PY{+w}{ }\PY{o}{=}\PY{+w}{ }\PY{n}{Executors}\PY{p}{.}\PY{n+na}{newVirtualThreadPerTaskExecutor}\PY{p}{(}\PY{p}{)}\PY{p}{;}
\PY{+w}{    }\PY{k+kd}{private}\PY{+w}{ }\PY{k+kd}{final}\PY{+w}{ }\PY{n}{UserService}\PY{+w}{ }\PY{n}{userService}\PY{p}{;}
\PY{+w}{    }\PY{k+kd}{private}\PY{+w}{ }\PY{k+kd}{final}\PY{+w}{ }\PY{n}{OrderService}\PY{+w}{ }\PY{n}{orderService}\PY{p}{;}
\PY{+w}{    }\PY{k+kd}{private}\PY{+w}{ }\PY{k+kd}{final}\PY{+w}{ }\PY{n}{RecommendationService}\PY{+w}{ }\PY{n}{recommendationService}\PY{p}{;}

\PY{+w}{    }\PY{n}{Profile}\PY{+w}{ }\PY{n+nf}{buildProfile}\PY{p}{(}\PY{n}{String}\PY{+w}{ }\PY{n}{userId}\PY{p}{)}\PY{+w}{ }\PY{k+kd}{throws}\PY{+w}{ }\PY{n}{ExecutionException}\PY{p}{,}\PY{+w}{ }\PY{n}{InterruptedException}\PY{+w}{ }\PY{p}{\PYZob{}}
\PY{+w}{        }\PY{n}{CompletableFuture}\PY{o}{\PYZlt{}}\PY{n}{User}\PY{o}{\PYZgt{}}\PY{+w}{ }\PY{n}{userFuture}\PY{+w}{ }\PY{o}{=}
\PY{+w}{                }\PY{n}{CompletableFuture}\PY{p}{.}\PY{n+na}{supplyAsync}\PY{p}{(}\PY{p}{(}\PY{p}{)}\PY{+w}{ }\PY{o}{\PYZhy{}}\PY{o}{\PYZgt{}}\PY{+w}{ }\PY{n}{uncheck}\PY{p}{(}\PY{p}{(}\PY{p}{)}\PY{+w}{ }\PY{o}{\PYZhy{}}\PY{o}{\PYZgt{}}\PY{+w}{ }\PY{n}{userService}\PY{p}{.}\PY{n+na}{fetch}\PY{p}{(}\PY{n}{userId}\PY{p}{)}\PY{p}{)}\PY{p}{,}\PY{+w}{ }\PY{n}{executor}\PY{p}{)}\PY{p}{;}
\PY{+w}{        }\PY{n}{CompletableFuture}\PY{o}{\PYZlt{}}\PY{n}{List}\PY{o}{\PYZlt{}}\PY{n}{Order}\PY{o}{\PYZgt{}}\PY{o}{\PYZgt{}}\PY{+w}{ }\PY{n}{ordersFuture}\PY{+w}{ }\PY{o}{=}
\PY{+w}{                }\PY{n}{CompletableFuture}\PY{p}{.}\PY{n+na}{supplyAsync}\PY{p}{(}\PY{p}{(}\PY{p}{)}\PY{+w}{ }\PY{o}{\PYZhy{}}\PY{o}{\PYZgt{}}\PY{+w}{ }\PY{n}{uncheck}\PY{p}{(}\PY{p}{(}\PY{p}{)}\PY{+w}{ }\PY{o}{\PYZhy{}}\PY{o}{\PYZgt{}}\PY{+w}{ }\PY{n}{orderService}\PY{p}{.}\PY{n+na}{recentOrders}\PY{p}{(}\PY{n}{userId}\PY{p}{)}\PY{p}{)}\PY{p}{,}\PY{+w}{ }\PY{n}{executor}\PY{p}{)}\PY{p}{;}
\PY{+w}{        }\PY{n}{CompletableFuture}\PY{o}{\PYZlt{}}\PY{n}{List}\PY{o}{\PYZlt{}}\PY{n}{Recommendation}\PY{o}{\PYZgt{}}\PY{o}{\PYZgt{}}\PY{+w}{ }\PY{n}{recsFuture}\PY{+w}{ }\PY{o}{=}
\PY{+w}{                }\PY{n}{CompletableFuture}\PY{p}{.}\PY{n+na}{supplyAsync}\PY{p}{(}
\PY{+w}{                        }\PY{p}{(}\PY{p}{)}\PY{+w}{ }\PY{o}{\PYZhy{}}\PY{o}{\PYZgt{}}\PY{+w}{ }\PY{n}{uncheck}\PY{p}{(}\PY{p}{(}\PY{p}{)}\PY{+w}{ }\PY{o}{\PYZhy{}}\PY{o}{\PYZgt{}}\PY{+w}{ }\PY{n}{recommendationService}\PY{p}{.}\PY{n+na}{recommendationsFor}\PY{p}{(}\PY{n}{userId}\PY{p}{)}\PY{p}{)}\PY{p}{,}\PY{+w}{ }\PY{n}{executor}\PY{p}{)}\PY{p}{;}

\PY{+w}{        }\PY{c+c1}{// allOf composes cancellation\PYZhy{}unaware futures; if one fails, the others}
\PY{+w}{        }\PY{c+c1}{// are NOT automatically cancelled by allOf — you\PYZsq{}d need to chain}
\PY{+w}{        }\PY{c+c1}{// .exceptionally / .whenComplete on each one yourself to cancel siblings,}
\PY{+w}{        }\PY{c+c1}{// which most CompletableFuture code in the wild does not bother doing.}
\PY{+w}{        }\PY{k}{return}\PY{+w}{ }\PY{n}{CompletableFuture}\PY{p}{.}\PY{n+na}{allOf}\PY{p}{(}\PY{n}{userFuture}\PY{p}{,}\PY{+w}{ }\PY{n}{ordersFuture}\PY{p}{,}\PY{+w}{ }\PY{n}{recsFuture}\PY{p}{)}
\PY{+w}{                }\PY{p}{.}\PY{n+na}{thenApply}\PY{p}{(}\PY{n}{v}\PY{+w}{ }\PY{o}{\PYZhy{}}\PY{o}{\PYZgt{}}\PY{+w}{ }\PY{k}{new}\PY{+w}{ }\PY{n}{Profile}\PY{p}{(}\PY{n}{userFuture}\PY{p}{.}\PY{n+na}{join}\PY{p}{(}\PY{p}{)}\PY{p}{,}\PY{+w}{ }\PY{n}{ordersFuture}\PY{p}{.}\PY{n+na}{join}\PY{p}{(}\PY{p}{)}\PY{p}{,}\PY{+w}{ }\PY{n}{recsFuture}\PY{p}{.}\PY{n+na}{join}\PY{p}{(}\PY{p}{)}\PY{p}{)}\PY{p}{)}
\PY{+w}{                }\PY{p}{.}\PY{n+na}{get}\PY{p}{(}\PY{p}{)}\PY{p}{;}
\PY{+w}{    }\PY{p}{\PYZcb{}}

\PY{+w}{    }\PY{k+kd}{private}\PY{+w}{ }\PY{k+kd}{static}\PY{+w}{ }\PY{o}{\PYZlt{}}\PY{n}{T}\PY{o}{\PYZgt{}}\PY{+w}{ }\PY{n}{T}\PY{+w}{ }\PY{n+nf}{uncheck}\PY{p}{(}\PY{n}{Callable}\PY{o}{\PYZlt{}}\PY{n}{T}\PY{o}{\PYZgt{}}\PY{+w}{ }\PY{n}{callable}\PY{p}{)}\PY{+w}{ }\PY{p}{\PYZob{}}
\PY{+w}{        }\PY{k}{try}\PY{+w}{ }\PY{p}{\PYZob{}}
\PY{+w}{            }\PY{k}{return}\PY{+w}{ }\PY{n}{callable}\PY{p}{.}\PY{n+na}{call}\PY{p}{(}\PY{p}{)}\PY{p}{;}
\PY{+w}{        }\PY{p}{\PYZcb{}}\PY{+w}{ }\PY{k}{catch}\PY{+w}{ }\PY{p}{(}\PY{n}{Exception}\PY{+w}{ }\PY{n}{e}\PY{p}{)}\PY{+w}{ }\PY{p}{\PYZob{}}
\PY{+w}{            }\PY{k}{throw}\PY{+w}{ }\PY{k}{new}\PY{+w}{ }\PY{n}{CompletionException}\PY{p}{(}\PY{n}{e}\PY{p}{)}\PY{p}{;}
\PY{+w}{        }\PY{p}{\PYZcb{}}
\PY{+w}{    }\PY{p}{\PYZcb{}}
\PY{p}{\PYZcb{}}
\end{Verbatim}
\end{codebox}

This is more composable than raw futures (chaining, combining, async pipelines) but the exception/cancellation story is still manual and easy to get subtly wrong — \code{allOf} waits for all of them regardless of failure, and nothing cancels the still-running siblings when one fails unless you add that logic yourself.

\subsection[Version 3: StructuredTaskScope]{Version 3: \code{StructuredTaskScope}}
\label{h:15:version-3-structuredtaskscope}
\begin{codebox}
\begin{Verbatim}[commandchars=\\\{\}]
\PY{k+kd}{class} \PY{n+nc}{StructuredProfileService}\PY{+w}{ }\PY{p}{\PYZob{}}
\PY{+w}{    }\PY{k+kd}{private}\PY{+w}{ }\PY{k+kd}{final}\PY{+w}{ }\PY{n}{UserService}\PY{+w}{ }\PY{n}{userService}\PY{p}{;}
\PY{+w}{    }\PY{k+kd}{private}\PY{+w}{ }\PY{k+kd}{final}\PY{+w}{ }\PY{n}{OrderService}\PY{+w}{ }\PY{n}{orderService}\PY{p}{;}
\PY{+w}{    }\PY{k+kd}{private}\PY{+w}{ }\PY{k+kd}{final}\PY{+w}{ }\PY{n}{RecommendationService}\PY{+w}{ }\PY{n}{recommendationService}\PY{p}{;}

\PY{+w}{    }\PY{n}{Profile}\PY{+w}{ }\PY{n+nf}{buildProfile}\PY{p}{(}\PY{n}{String}\PY{+w}{ }\PY{n}{userId}\PY{p}{)}\PY{+w}{ }\PY{k+kd}{throws}\PY{+w}{ }\PY{n}{InterruptedException}\PY{p}{,}\PY{+w}{ }\PY{n}{ExecutionException}\PY{+w}{ }\PY{p}{\PYZob{}}
\PY{+w}{        }\PY{k}{try}\PY{+w}{ }\PY{p}{(}\PY{k+kd}{var}\PY{+w}{ }\PY{n}{scope}\PY{+w}{ }\PY{o}{=}\PY{+w}{ }\PY{k}{new}\PY{+w}{ }\PY{n}{StructuredTaskScope}\PY{p}{.}\PY{n+na}{ShutdownOnFailure}\PY{p}{(}\PY{p}{)}\PY{p}{)}\PY{+w}{ }\PY{p}{\PYZob{}}
\PY{+w}{            }\PY{n}{Subtask}\PY{o}{\PYZlt{}}\PY{n}{User}\PY{o}{\PYZgt{}}\PY{+w}{ }\PY{n}{userTask}\PY{+w}{ }\PY{o}{=}\PY{+w}{ }\PY{n}{scope}\PY{p}{.}\PY{n+na}{fork}\PY{p}{(}\PY{p}{(}\PY{p}{)}\PY{+w}{ }\PY{o}{\PYZhy{}}\PY{o}{\PYZgt{}}\PY{+w}{ }\PY{n}{userService}\PY{p}{.}\PY{n+na}{fetch}\PY{p}{(}\PY{n}{userId}\PY{p}{)}\PY{p}{)}\PY{p}{;}
\PY{+w}{            }\PY{n}{Subtask}\PY{o}{\PYZlt{}}\PY{n}{List}\PY{o}{\PYZlt{}}\PY{n}{Order}\PY{o}{\PYZgt{}}\PY{o}{\PYZgt{}}\PY{+w}{ }\PY{n}{ordersTask}\PY{+w}{ }\PY{o}{=}\PY{+w}{ }\PY{n}{scope}\PY{p}{.}\PY{n+na}{fork}\PY{p}{(}\PY{p}{(}\PY{p}{)}\PY{+w}{ }\PY{o}{\PYZhy{}}\PY{o}{\PYZgt{}}\PY{+w}{ }\PY{n}{orderService}\PY{p}{.}\PY{n+na}{recentOrders}\PY{p}{(}\PY{n}{userId}\PY{p}{)}\PY{p}{)}\PY{p}{;}
\PY{+w}{            }\PY{n}{Subtask}\PY{o}{\PYZlt{}}\PY{n}{List}\PY{o}{\PYZlt{}}\PY{n}{Recommendation}\PY{o}{\PYZgt{}}\PY{o}{\PYZgt{}}\PY{+w}{ }\PY{n}{recsTask}\PY{+w}{ }\PY{o}{=}
\PY{+w}{                    }\PY{n}{scope}\PY{p}{.}\PY{n+na}{fork}\PY{p}{(}\PY{p}{(}\PY{p}{)}\PY{+w}{ }\PY{o}{\PYZhy{}}\PY{o}{\PYZgt{}}\PY{+w}{ }\PY{n}{recommendationService}\PY{p}{.}\PY{n+na}{recommendationsFor}\PY{p}{(}\PY{n}{userId}\PY{p}{)}\PY{p}{)}\PY{p}{;}

\PY{+w}{            }\PY{n}{scope}\PY{p}{.}\PY{n+na}{join}\PY{p}{(}\PY{p}{)}\PY{p}{;}
\PY{+w}{            }\PY{n}{scope}\PY{p}{.}\PY{n+na}{throwIfFailed}\PY{p}{(}\PY{p}{)}\PY{p}{;}\PY{+w}{ }\PY{c+c1}{// rethrows the first failure; other subtasks are}
\PY{+w}{                                    }\PY{c+c1}{// already cancelled by ShutdownOnFailure by this point}

\PY{+w}{            }\PY{k}{return}\PY{+w}{ }\PY{k}{new}\PY{+w}{ }\PY{n}{Profile}\PY{p}{(}\PY{n}{userTask}\PY{p}{.}\PY{n+na}{get}\PY{p}{(}\PY{p}{)}\PY{p}{,}\PY{+w}{ }\PY{n}{ordersTask}\PY{p}{.}\PY{n+na}{get}\PY{p}{(}\PY{p}{)}\PY{p}{,}\PY{+w}{ }\PY{n}{recsTask}\PY{p}{.}\PY{n+na}{get}\PY{p}{(}\PY{p}{)}\PY{p}{)}\PY{p}{;}
\PY{+w}{        }\PY{p}{\PYZcb{}}\PY{+w}{ }\PY{c+c1}{// close() guarantees zero orphaned subtasks, success or failure}
\PY{+w}{    }\PY{p}{\PYZcb{}}
\PY{p}{\PYZcb{}}
\end{Verbatim}
\end{codebox}

\textbf{Comparison:} all three versions run the three downstream calls concurrently and produce the same result on the happy path. They diverge sharply on the failure path: Version 1 requires you to remember manual cancellation in every catch block; Version 2’s \code{allOf} does not cancel siblings on failure unless you add explicit \code{exceptionally}/\code{whenComplete} wiring per future; Version 3 gets fail-fast cancellation of siblings and a guaranteed-no-orphans lifetime \emph{for free}, as a structural property of the scope, with the least code and the clearest error-handling story. This is why structured concurrency, despite still being preview, is the direction the language is pushing fan-out/fan-in code toward.

\section[Common Code-Review Interview Pitfalls]{Common Code-Review Interview Pitfalls}
\label{h:15:common-code-review-interview-pitfalls}
\begin{enumerate}
\item 
\textbf{Pooling virtual threads (\code{newFixedThreadPool} with a virtual thread factory).}
Why it matters: pooling exists to amortize the cost of expensive platform threads and to bound concurrency; virtual threads are cheap to create, so pooling them buys nothing and reintroduces the same concurrency ceiling virtual threads were meant to remove.

\begin{codebox}
\begin{Verbatim}[commandchars=\\\{\}]
\PY{c+c1}{// Before}
\PY{k}{new}\PY{+w}{ }\PY{n}{ThreadPoolExecutor}\PY{p}{(}\PY{l+m+mi}{200}\PY{p}{,}\PY{+w}{ }\PY{l+m+mi}{200}\PY{p}{,}\PY{+w}{ }\PY{l+m+mi}{0}\PY{p}{,}\PY{+w}{ }\PY{n}{TimeUnit}\PY{p}{.}\PY{n+na}{SECONDS}\PY{p}{,}\PY{+w}{ }\PY{n}{queue}\PY{p}{,}\PY{+w}{ }\PY{n}{Thread}\PY{p}{.}\PY{n+na}{ofVirtual}\PY{p}{(}\PY{p}{)}\PY{p}{.}\PY{n+na}{factory}\PY{p}{(}\PY{p}{)}\PY{p}{)}\PY{p}{;}
\PY{c+c1}{// After}
\PY{n}{Executors}\PY{p}{.}\PY{n+na}{newVirtualThreadPerTaskExecutor}\PY{p}{(}\PY{p}{)}\PY{p}{;}
\end{Verbatim}
\end{codebox}

\item 
\textbf{Using virtual threads for CPU-bound work and expecting a speedup.}
Why it matters: virtual threads only help when threads spend time \emph{blocked}; CPU-bound work occupies its carrier thread the whole time, so a million CPU-bound virtual threads run no faster than a small platform-thread pool sized to the core count — and adds bookkeeping overhead for no benefit.

\begin{codebox}
\begin{Verbatim}[commandchars=\\\{\}]
\PY{c+c1}{// Before}
\PY{k}{try}\PY{+w}{ }\PY{p}{(}\PY{k+kd}{var}\PY{+w}{ }\PY{n}{ex}\PY{+w}{ }\PY{o}{=}\PY{+w}{ }\PY{n}{Executors}\PY{p}{.}\PY{n+na}{newVirtualThreadPerTaskExecutor}\PY{p}{(}\PY{p}{)}\PY{p}{)}\PY{+w}{ }\PY{p}{\PYZob{}}
\PY{+w}{    }\PY{k}{for}\PY{+w}{ }\PY{p}{(}\PY{k+kd}{var}\PY{+w}{ }\PY{n}{m}\PY{+w}{ }\PY{p}{:}\PY{+w}{ }\PY{n}{matrices}\PY{p}{)}\PY{+w}{ }\PY{n}{ex}\PY{p}{.}\PY{n+na}{submit}\PY{p}{(}\PY{p}{(}\PY{p}{)}\PY{+w}{ }\PY{o}{\PYZhy{}}\PY{o}{\PYZgt{}}\PY{+w}{ }\PY{n}{multiply}\PY{p}{(}\PY{n}{m}\PY{p}{)}\PY{p}{)}\PY{p}{;}\PY{+w}{ }\PY{c+c1}{// pure CPU work}
\PY{p}{\PYZcb{}}
\PY{c+c1}{// After}
\PY{k}{try}\PY{+w}{ }\PY{p}{(}\PY{k+kd}{var}\PY{+w}{ }\PY{n}{ex}\PY{+w}{ }\PY{o}{=}\PY{+w}{ }\PY{n}{Executors}\PY{p}{.}\PY{n+na}{newFixedThreadPool}\PY{p}{(}\PY{n}{Runtime}\PY{p}{.}\PY{n+na}{getRuntime}\PY{p}{(}\PY{p}{)}\PY{p}{.}\PY{n+na}{availableProcessors}\PY{p}{(}\PY{p}{)}\PY{p}{)}\PY{p}{)}\PY{+w}{ }\PY{p}{\PYZob{}}
\PY{+w}{    }\PY{k}{for}\PY{+w}{ }\PY{p}{(}\PY{k+kd}{var}\PY{+w}{ }\PY{n}{m}\PY{+w}{ }\PY{p}{:}\PY{+w}{ }\PY{n}{matrices}\PY{p}{)}\PY{+w}{ }\PY{n}{ex}\PY{p}{.}\PY{n+na}{submit}\PY{p}{(}\PY{p}{(}\PY{p}{)}\PY{+w}{ }\PY{o}{\PYZhy{}}\PY{o}{\PYZgt{}}\PY{+w}{ }\PY{n}{multiply}\PY{p}{(}\PY{n}{m}\PY{p}{)}\PY{p}{)}\PY{p}{;}
\PY{p}{\PYZcb{}}
\end{Verbatim}
\end{codebox}

\item 
\textbf{Blocking inside a \code{synchronized} block on JDK 21-23 without realizing it pins the carrier thread.}
Why it matters: pinning ties up a carrier thread for the whole blocking call, effectively turning a cheap virtual thread back into an expensive scarce resource, and can silently reduce throughput under load until diagnosed with \code{-\allowbreak{}Djdk.\allowbreak{}trace\allowbreak{}Pinned\allowbreak{}Threads}.

\begin{codebox}
\begin{Verbatim}[commandchars=\\\{\}]
\PY{c+c1}{// Before (JDK 21\PYZhy{}23)}
\PY{k+kd}{synchronized}\PY{+w}{ }\PY{p}{(}\PY{n}{lock}\PY{p}{)}\PY{+w}{ }\PY{p}{\PYZob{}}\PY{+w}{ }\PY{n}{String}\PY{+w}{ }\PY{n}{r}\PY{+w}{ }\PY{o}{=}\PY{+w}{ }\PY{n}{callSlowService}\PY{p}{(}\PY{p}{)}\PY{p}{;}\PY{+w}{ }\PY{p}{\PYZcb{}}
\PY{c+c1}{// After}
\PY{n}{lock}\PY{p}{.}\PY{n+na}{lock}\PY{p}{(}\PY{p}{)}\PY{p}{;}
\PY{k}{try}\PY{+w}{ }\PY{p}{\PYZob{}}\PY{+w}{ }\PY{n}{String}\PY{+w}{ }\PY{n}{r}\PY{+w}{ }\PY{o}{=}\PY{+w}{ }\PY{n}{callSlowService}\PY{p}{(}\PY{p}{)}\PY{p}{;}\PY{+w}{ }\PY{p}{\PYZcb{}}\PY{+w}{ }\PY{k}{finally}\PY{+w}{ }\PY{p}{\PYZob{}}\PY{+w}{ }\PY{n}{lock}\PY{p}{.}\PY{n+na}{unlock}\PY{p}{(}\PY{p}{)}\PY{p}{;}\PY{+w}{ }\PY{p}{\PYZcb{}}
\end{Verbatim}
\end{codebox}

\item 
\textbf{Claiming \code{synchronized} always pins virtual threads, regardless of JDK version.}
Why it matters: JEP 491 (JDK 24) removed \code{synchronized}-block pinning; stating this as a universal, version-independent rule in a review or interview is factually outdated and signals stale knowledge. Native-frame pinning is still a concern on every version.

\begin{codebox}
\begin{Verbatim}
// Before: "never use synchronized with virtual threads, period"
// After: "synchronized pins carriers on JDK 21-23; JEP 491 (JDK 24+) fixed
//         this for synchronized specifically — native-frame pinning remains"
\end{Verbatim}
\end{codebox}

\item 
\textbf{Using \code{ThreadLocal}/\code{InheritableThreadLocal} for per-request context at virtual-thread scale.}
Why it matters: mutability plus unbounded lifetime plus per-child inheritance copying becomes expensive and leak-prone once there can be millions of short-lived threads; \code{ScopedValue} gives immutability, automatic unbinding, and cheap structured inheritance instead.

\begin{codebox}
\begin{Verbatim}[commandchars=\\\{\}]
\PY{c+c1}{// Before}
\PY{k+kd}{private}\PY{+w}{ }\PY{k+kd}{static}\PY{+w}{ }\PY{k+kd}{final}\PY{+w}{ }\PY{n}{ThreadLocal}\PY{o}{\PYZlt{}}\PY{n}{String}\PY{o}{\PYZgt{}}\PY{+w}{ }\PY{n}{TENANT\PYZus{}ID}\PY{+w}{ }\PY{o}{=}\PY{+w}{ }\PY{k}{new}\PY{+w}{ }\PY{n}{ThreadLocal}\PY{o}{\PYZlt{}}\PY{o}{\PYZgt{}}\PY{p}{(}\PY{p}{)}\PY{p}{;}
\PY{c+c1}{// After}
\PY{k+kd}{private}\PY{+w}{ }\PY{k+kd}{static}\PY{+w}{ }\PY{k+kd}{final}\PY{+w}{ }\PY{n}{ScopedValue}\PY{o}{\PYZlt{}}\PY{n}{String}\PY{o}{\PYZgt{}}\PY{+w}{ }\PY{n}{TENANT\PYZus{}ID}\PY{+w}{ }\PY{o}{=}\PY{+w}{ }\PY{n}{ScopedValue}\PY{p}{.}\PY{n+na}{newInstance}\PY{p}{(}\PY{p}{)}\PY{p}{;}
\end{Verbatim}
\end{codebox}

\item 
\textbf{Fan-out with raw \code{ExecutorService}/\code{Future} and no cancellation on the failure path.}
Why it matters: when one \code{Future.\allowbreak{}get()} throws, sibling tasks keep running unattended unless you explicitly cancel every one of them in every catch block — a leak that’s easy to introduce and easy to miss in review.

\begin{codebox}
\begin{Verbatim}[commandchars=\\\{\}]
\PY{c+c1}{// Before}
\PY{n}{User}\PY{+w}{ }\PY{n}{u}\PY{+w}{ }\PY{o}{=}\PY{+w}{ }\PY{n}{userFuture}\PY{p}{.}\PY{n+na}{get}\PY{p}{(}\PY{p}{)}\PY{p}{;}\PY{+w}{ }\PY{n}{Order}\PY{+w}{ }\PY{n}{o}\PY{+w}{ }\PY{o}{=}\PY{+w}{ }\PY{n}{orderFuture}\PY{p}{.}\PY{n+na}{get}\PY{p}{(}\PY{p}{)}\PY{p}{;}\PY{+w}{ }\PY{c+c1}{// orderFuture orphaned if userFuture throws}
\PY{c+c1}{// After}
\PY{k}{try}\PY{+w}{ }\PY{p}{(}\PY{k+kd}{var}\PY{+w}{ }\PY{n}{scope}\PY{+w}{ }\PY{o}{=}\PY{+w}{ }\PY{k}{new}\PY{+w}{ }\PY{n}{StructuredTaskScope}\PY{p}{.}\PY{n+na}{ShutdownOnFailure}\PY{p}{(}\PY{p}{)}\PY{p}{)}\PY{+w}{ }\PY{p}{\PYZob{}}
\PY{+w}{    }\PY{k+kd}{var}\PY{+w}{ }\PY{n}{uT}\PY{+w}{ }\PY{o}{=}\PY{+w}{ }\PY{n}{scope}\PY{p}{.}\PY{n+na}{fork}\PY{p}{(}\PY{p}{(}\PY{p}{)}\PY{+w}{ }\PY{o}{\PYZhy{}}\PY{o}{\PYZgt{}}\PY{+w}{ }\PY{n}{fetchUser}\PY{p}{(}\PY{p}{)}\PY{p}{)}\PY{p}{;}\PY{+w}{ }\PY{k+kd}{var}\PY{+w}{ }\PY{n}{oT}\PY{+w}{ }\PY{o}{=}\PY{+w}{ }\PY{n}{scope}\PY{p}{.}\PY{n+na}{fork}\PY{p}{(}\PY{p}{(}\PY{p}{)}\PY{+w}{ }\PY{o}{\PYZhy{}}\PY{o}{\PYZgt{}}\PY{+w}{ }\PY{n}{fetchOrder}\PY{p}{(}\PY{p}{)}\PY{p}{)}\PY{p}{;}
\PY{+w}{    }\PY{n}{scope}\PY{p}{.}\PY{n+na}{join}\PY{p}{(}\PY{p}{)}\PY{p}{;}\PY{+w}{ }\PY{n}{scope}\PY{p}{.}\PY{n+na}{throwIfFailed}\PY{p}{(}\PY{p}{)}\PY{p}{;}
\PY{p}{\PYZcb{}}
\end{Verbatim}
\end{codebox}

\item 
\textbf{Assuming \code{Completable\allowbreak{}Future.\allowbreak{}all\allowbreak{}Of(...)} cancels siblings when one future fails.}
Why it matters: \code{allOf} simply waits for every future regardless of outcome; it does not cancel anything on failure, so “fail fast, stop wasted work” requires extra manual wiring that’s frequently omitted.

\begin{codebox}
\begin{Verbatim}[commandchars=\\\{\}]
\PY{c+c1}{// Before: assumes cancellation happens automatically}
\PY{n}{CompletableFuture}\PY{p}{.}\PY{n+na}{allOf}\PY{p}{(}\PY{n}{f1}\PY{p}{,}\PY{+w}{ }\PY{n}{f2}\PY{p}{,}\PY{+w}{ }\PY{n}{f3}\PY{p}{)}\PY{p}{.}\PY{n+na}{get}\PY{p}{(}\PY{p}{)}\PY{p}{;}
\PY{c+c1}{// After: use StructuredTaskScope.ShutdownOnFailure for genuine fail\PYZhy{}fast cancellation}
\end{Verbatim}
\end{codebox}

\item 
\textbf{Forgetting \code{--\allowbreak{}enable-\allowbreak{}preview} (or claiming the feature is final) for structured concurrency / scoped values.}
Why it matters: both remain preview APIs through JDK 25 at the time of writing; code that compiles for a reviewer without the flag, or a claim in an interview that they are GA, is simply incorrect and easy for an interviewer to catch.

\begin{codebox}
\begin{Verbatim}[commandchars=\\\{\}]
\PY{c+c1}{\PYZsh{} Before}
javac\PY{+w}{ }Main.java\PY{+w}{   }\PY{c+c1}{\PYZsh{} fails: preview API used without \PYZhy{}\PYZhy{}enable\PYZhy{}preview}
\PY{c+c1}{\PYZsh{} After}
javac\PY{+w}{ }\PYZhy{}\PYZhy{}release\PY{+w}{ }\PY{l+m}{25}\PY{+w}{ }\PYZhy{}\PYZhy{}enable\PYZhy{}preview\PY{+w}{ }Main.java\PY{+w}{ }\PY{o}{\PYZam{}\PYZam{}}\PY{+w}{ }java\PY{+w}{ }\PYZhy{}\PYZhy{}enable\PYZhy{}preview\PY{+w}{ }Main
\end{Verbatim}
\end{codebox}

\item 
\textbf{Writing \code{StructuredTaskScope} code against the wrong API shape for the target JDK (subclassing vs \code{Joiner}).}
Why it matters: JDK 21-24 use \code{new~\allowbreak{}Structured\allowbreak{}Task\allowbreak{}Scope.\allowbreak{}Shutdown\allowbreak{}On\allowbreak{}Failure()}-style subclassing; JDK 25 reshapes this around \code{Structured\allowbreak{}Task\allowbreak{}Scope.\allowbreak{}open(\allowbreak{}Joiner...)}. Code (or an answer) that mixes the two shapes reveals unfamiliarity with the feature’s actual evolution.

\begin{codebox}
\begin{Verbatim}[commandchars=\\\{\}]
\PY{c+c1}{// JDK 21\PYZhy{}24}
\PY{k}{try}\PY{+w}{ }\PY{p}{(}\PY{k+kd}{var}\PY{+w}{ }\PY{n}{scope}\PY{+w}{ }\PY{o}{=}\PY{+w}{ }\PY{k}{new}\PY{+w}{ }\PY{n}{StructuredTaskScope}\PY{p}{.}\PY{n+na}{ShutdownOnFailure}\PY{p}{(}\PY{p}{)}\PY{p}{)}\PY{+w}{ }\PY{p}{\PYZob{}}\PY{+w}{ }\PY{p}{.}\PY{p}{.}\PY{p}{.}\PY{+w}{ }\PY{p}{\PYZcb{}}
\PY{c+c1}{// JDK 25}
\PY{k}{try}\PY{+w}{ }\PY{p}{(}\PY{k+kd}{var}\PY{+w}{ }\PY{n}{scope}\PY{+w}{ }\PY{o}{=}\PY{+w}{ }\PY{n}{StructuredTaskScope}\PY{p}{.}\PY{n+na}{open}\PY{p}{(}\PY{n}{Joiner}\PY{p}{.}\PY{n+na}{awaitAllSuccessfulOrThrow}\PY{p}{(}\PY{p}{)}\PY{p}{)}\PY{p}{)}\PY{+w}{ }\PY{p}{\PYZob{}}\PY{+w}{ }\PY{p}{.}\PY{p}{.}\PY{p}{.}\PY{+w}{ }\PY{p}{\PYZcb{}}
\end{Verbatim}
\end{codebox}

\item 
\textbf{Not bounding concurrency to a genuinely scarce downstream resource just because virtual threads themselves are unbounded.}
Why it matters: “virtual threads are free” does not mean the database, connection pool, or third-party API on the other end can handle unbounded concurrent calls; without a \code{Semaphore} (or similar) around the call site, unbounded virtual-thread fan-out can overwhelm a downstream that used to be implicitly protected by a small platform-thread pool.

\begin{codebox}
\begin{Verbatim}[commandchars=\\\{\}]
\PY{c+c1}{// Before: every virtual thread hits the DB directly, no limit}
\PY{n}{executor}\PY{p}{.}\PY{n+na}{submit}\PY{p}{(}\PY{p}{(}\PY{p}{)}\PY{+w}{ }\PY{o}{\PYZhy{}}\PY{o}{\PYZgt{}}\PY{+w}{ }\PY{n}{db}\PY{p}{.}\PY{n+na}{query}\PY{p}{(}\PY{n}{sql}\PY{p}{)}\PY{p}{)}\PY{p}{;}
\PY{c+c1}{// After}
\PY{n}{limiter}\PY{p}{.}\PY{n+na}{acquire}\PY{p}{(}\PY{p}{)}\PY{p}{;}
\PY{k}{try}\PY{+w}{ }\PY{p}{\PYZob{}}\PY{+w}{ }\PY{n}{db}\PY{p}{.}\PY{n+na}{query}\PY{p}{(}\PY{n}{sql}\PY{p}{)}\PY{p}{;}\PY{+w}{ }\PY{p}{\PYZcb{}}\PY{+w}{ }\PY{k}{finally}\PY{+w}{ }\PY{p}{\PYZob{}}\PY{+w}{ }\PY{n}{limiter}\PY{p}{.}\PY{n+na}{release}\PY{p}{(}\PY{p}{)}\PY{p}{;}\PY{+w}{ }\PY{p}{\PYZcb{}}
\end{Verbatim}
\end{codebox}

\item 
\textbf{Letting a \code{ScopedValue} binding’s lifetime outlive the logical operation by binding too broadly (e.g., at application startup).}
Why it matters: the safety benefit of \code{ScopedValue} comes from its binding being tightly scoped to one call tree; binding it once at a very high level (defeating the purpose) reintroduces the same “value visible everywhere, for too long” problem \code{ThreadLocal} had, just with extra ceremony.

\begin{codebox}
\begin{Verbatim}[commandchars=\\\{\}]
\PY{c+c1}{// Before: bound once for the whole application run — no real scoping benefit}
\PY{n}{ScopedValue}\PY{p}{.}\PY{n+na}{where}\PY{p}{(}\PY{n}{TENANT\PYZus{}ID}\PY{p}{,}\PY{+w}{ }\PY{l+s}{\PYZdq{}}\PY{l+s}{default}\PY{l+s}{\PYZdq{}}\PY{p}{)}\PY{p}{.}\PY{n+na}{run}\PY{p}{(}\PY{n}{Application}\PY{p}{:}\PY{p}{:}\PY{n}{run}\PY{p}{)}\PY{p}{;}
\PY{c+c1}{// After: bind per request, inside the request\PYZhy{}handling method}
\PY{n}{ScopedValue}\PY{p}{.}\PY{n+na}{where}\PY{p}{(}\PY{n}{TENANT\PYZus{}ID}\PY{p}{,}\PY{+w}{ }\PY{n}{tenantId}\PY{p}{)}\PY{p}{.}\PY{n+na}{run}\PY{p}{(}\PY{p}{(}\PY{p}{)}\PY{+w}{ }\PY{o}{\PYZhy{}}\PY{o}{\PYZgt{}}\PY{+w}{ }\PY{n}{handleRequest}\PY{p}{(}\PY{n}{exchange}\PY{p}{)}\PY{p}{)}\PY{p}{;}
\end{Verbatim}
\end{codebox}

\item 
\textbf{Expecting a \code{StructuredTaskScope}’s forked subtasks to survive past the enclosing try-with-resources block.}
Why it matters: \code{close()} enforces that no subtask outlives the scope — leaking a \code{Subtask} handle or its result outside the block, expecting it to keep running in the background, misunderstands the entire point of structured concurrency (bounded, visible lifetimes).

\begin{codebox}
\begin{Verbatim}[commandchars=\\\{\}]
\PY{c+c1}{// Before: expecting background work to continue after the block}
\PY{n}{Subtask}\PY{o}{\PYZlt{}}\PY{n}{String}\PY{o}{\PYZgt{}}\PY{+w}{ }\PY{n}{leaked}\PY{p}{;}
\PY{k}{try}\PY{+w}{ }\PY{p}{(}\PY{k+kd}{var}\PY{+w}{ }\PY{n}{scope}\PY{+w}{ }\PY{o}{=}\PY{+w}{ }\PY{k}{new}\PY{+w}{ }\PY{n}{StructuredTaskScope}\PY{p}{.}\PY{n+na}{ShutdownOnFailure}\PY{p}{(}\PY{p}{)}\PY{p}{)}\PY{+w}{ }\PY{p}{\PYZob{}}
\PY{+w}{    }\PY{n}{leaked}\PY{+w}{ }\PY{o}{=}\PY{+w}{ }\PY{n}{scope}\PY{p}{.}\PY{n+na}{fork}\PY{p}{(}\PY{p}{(}\PY{p}{)}\PY{+w}{ }\PY{o}{\PYZhy{}}\PY{o}{\PYZgt{}}\PY{+w}{ }\PY{n}{longRunningJob}\PY{p}{(}\PY{p}{)}\PY{p}{)}\PY{p}{;}
\PY{p}{\PYZcb{}}\PY{+w}{ }\PY{c+c1}{// scope.close() already joined/cancelled it — it did not \PYZdq{}keep running\PYZdq{}}
\PY{c+c1}{// After: if it must genuinely outlive this method, don\PYZsq{}t put it in a scope —}
\PY{c+c1}{//        use a plain virtual thread or executor with its own lifecycle}
\end{Verbatim}
\end{codebox}

\item 
\textbf{Reusing a single, wide \code{Semaphore} (or none at all) for unrelated downstream resources instead of one per resource.}
Why it matters: a shared limiter conflates unrelated concurrency budgets — a burst of calls to a fast, high-capacity service can starve permits needed by a slow, low-capacity one (or vice versa) — undermining the whole point of per-resource throttling.

\begin{codebox}
\begin{Verbatim}[commandchars=\\\{\}]
\PY{c+c1}{// Before: one limiter for every downstream call}
\PY{k+kd}{private}\PY{+w}{ }\PY{k+kd}{final}\PY{+w}{ }\PY{n}{Semaphore}\PY{+w}{ }\PY{n}{anyDownstream}\PY{+w}{ }\PY{o}{=}\PY{+w}{ }\PY{k}{new}\PY{+w}{ }\PY{n}{Semaphore}\PY{p}{(}\PY{l+m+mi}{100}\PY{p}{)}\PY{p}{;}
\PY{c+c1}{// After: one limiter per distinct scarce resource}
\PY{k+kd}{private}\PY{+w}{ }\PY{k+kd}{final}\PY{+w}{ }\PY{n}{Semaphore}\PY{+w}{ }\PY{n}{paymentGatewayLimiter}\PY{+w}{ }\PY{o}{=}\PY{+w}{ }\PY{k}{new}\PY{+w}{ }\PY{n}{Semaphore}\PY{p}{(}\PY{l+m+mi}{50}\PY{p}{)}\PY{p}{;}
\PY{k+kd}{private}\PY{+w}{ }\PY{k+kd}{final}\PY{+w}{ }\PY{n}{Semaphore}\PY{+w}{ }\PY{n}{inventoryServiceLimiter}\PY{+w}{ }\PY{o}{=}\PY{+w}{ }\PY{k}{new}\PY{+w}{ }\PY{n}{Semaphore}\PY{p}{(}\PY{l+m+mi}{200}\PY{p}{)}\PY{p}{;}
\end{Verbatim}
\end{codebox}

\item 
\textbf{Ignoring \code{InterruptedException} inside a forked subtask instead of letting cancellation propagate.}
Why it matters: \code{StructuredTaskScope}’s cancellation (on \code{ShutdownOnFailure}/\code{ShutdownOnSuccess}) works by interrupting running subtasks; a subtask that swallows \code{InterruptedException} without restoring the flag or exiting defeats the scope’s ability to actually stop wasted work when a sibling fails or wins.

\begin{codebox}
\begin{Verbatim}[commandchars=\\\{\}]
\PY{c+c1}{// Before}
\PY{n}{scope}\PY{p}{.}\PY{n+na}{fork}\PY{p}{(}\PY{p}{(}\PY{p}{)}\PY{+w}{ }\PY{o}{\PYZhy{}}\PY{o}{\PYZgt{}}\PY{+w}{ }\PY{p}{\PYZob{}}\PY{+w}{ }\PY{k}{try}\PY{+w}{ }\PY{p}{\PYZob{}}\PY{+w}{ }\PY{k}{return}\PY{+w}{ }\PY{n}{slowCall}\PY{p}{(}\PY{p}{)}\PY{p}{;}\PY{+w}{ }\PY{p}{\PYZcb{}}\PY{+w}{ }\PY{k}{catch}\PY{+w}{ }\PY{p}{(}\PY{n}{InterruptedException}\PY{+w}{ }\PY{n}{e}\PY{p}{)}\PY{+w}{ }\PY{p}{\PYZob{}}\PY{+w}{ }\PY{k}{return}\PY{+w}{ }\PY{n}{retry}\PY{p}{(}\PY{p}{)}\PY{p}{;}\PY{+w}{ }\PY{p}{\PYZcb{}}\PY{+w}{ }\PY{p}{\PYZcb{}}\PY{p}{)}\PY{p}{;}
\PY{c+c1}{// After}
\PY{n}{scope}\PY{p}{.}\PY{n+na}{fork}\PY{p}{(}\PY{p}{(}\PY{p}{)}\PY{+w}{ }\PY{o}{\PYZhy{}}\PY{o}{\PYZgt{}}\PY{+w}{ }\PY{p}{\PYZob{}}
\PY{+w}{    }\PY{k}{try}\PY{+w}{ }\PY{p}{\PYZob{}}\PY{+w}{ }\PY{k}{return}\PY{+w}{ }\PY{n}{slowCall}\PY{p}{(}\PY{p}{)}\PY{p}{;}\PY{+w}{ }\PY{p}{\PYZcb{}}
\PY{+w}{    }\PY{k}{catch}\PY{+w}{ }\PY{p}{(}\PY{n}{InterruptedException}\PY{+w}{ }\PY{n}{e}\PY{p}{)}\PY{+w}{ }\PY{p}{\PYZob{}}\PY{+w}{ }\PY{n}{Thread}\PY{p}{.}\PY{n+na}{currentThread}\PY{p}{(}\PY{p}{)}\PY{p}{.}\PY{n+na}{interrupt}\PY{p}{(}\PY{p}{)}\PY{p}{;}\PY{+w}{ }\PY{k}{throw}\PY{+w}{ }\PY{n}{e}\PY{p}{;}\PY{+w}{ }\PY{p}{\PYZcb{}}
\PY{p}{\PYZcb{}}\PY{p}{)}\PY{p}{;}
\end{Verbatim}
\end{codebox}

\item 
\textbf{Diagnosing “too many threads” performance problems on a virtual-thread-heavy service with \code{jstack}-style plain-text dumps instead of the JSON format.}
Why it matters: with hundreds of thousands of virtual threads, a flat text dump is impractical to analyze; \code{jcmd~\allowbreak{}\textless{}\allowbreak{}pid\textgreater{}\allowbreak{}~\allowbreak{}Thread.\allowbreak{}dump\_\allowbreak{}to\_\allowbreak{}file~\allowbreak{}-\allowbreak{}format=\allowbreak{}json} produces a structured dump (grouped by carrier thread) intended for tooling, and JFR’s \code{jdk.\allowbreak{}Virtual\allowbreak{}Thread\allowbreak{}Pinned} event is the supported way to find pinning at scale rather than eyeballing text.

\begin{codebox}
\begin{Verbatim}[commandchars=\\\{\}]
\PY{c+c1}{\PYZsh{} Before}
jstack\PY{+w}{ }\PYZlt{}pid\PYZgt{}\PY{+w}{ }\PYZgt{}\PY{+w}{ }dump.txt\PY{+w}{   }\PY{c+c1}{\PYZsh{} unwieldy at hundreds of thousands of virtual threads}
\PY{c+c1}{\PYZsh{} After}
jcmd\PY{+w}{ }\PYZlt{}pid\PYZgt{}\PY{+w}{ }Thread.dump\PYZus{}to\PYZus{}file\PY{+w}{ }\PYZhy{}format\PY{o}{=}json\PY{+w}{ }threads.json
\end{Verbatim}
\end{codebox}

\end{enumerate}


\setcounter{chapter}{15}
\chapter{Reflection \& Runtime Features}
\label{chap:16}
\label{h:16:16-reflection--runtime-features}
Most Java code is written to be checked and wired together at \textbf{compile time}: you call a method, the compiler verifies it exists, and the JVM runs it. \textbf{Reflection} and its related APIs flip that around — they let code inspect and call other code \emph{while the program is running}, even if the exact class or method wasn’t known when it was compiled. This is how frameworks like Spring, Hibernate, and Jackson do their “magic”: they look at your classes, fields, and annotations at runtime and build behavior around them. This chapter covers reflection itself, annotations (the metadata reflection often reads), annotation processing (a compile-time alternative), dynamic proxies, the newer and faster method handle and VarHandle APIs, what replaced the old \code{Unsafe} class, and the Service Provider Interface (SPI) mechanism for plugin-style architectures. We target Java 21+ throughout.

\section[Reflection]{Reflection}
\label{h:16:reflection}
\textbf{Reflection} is the ability of code to examine and manipulate classes, fields, methods, and constructors at runtime, using objects instead of source-code syntax. The entry point is almost always a \code{Class\textless{}?\textgreater{}} object, which represents a loaded class inside the JVM. Every object has exactly one \code{Class} object describing its type.

There are three common ways to get a \code{Class} object:

\begin{codebox}
\begin{Verbatim}[commandchars=\\\{\}]
\PY{c+c1}{// 1. From a class literal — known at compile time, checked by the compiler}
\PY{n}{Class}\PY{o}{\PYZlt{}}\PY{n}{String}\PY{o}{\PYZgt{}}\PY{+w}{ }\PY{n}{byLiteral}\PY{+w}{ }\PY{o}{=}\PY{+w}{ }\PY{n}{String}\PY{p}{.}\PY{n+na}{class}\PY{p}{;}

\PY{c+c1}{// 2. From an instance you already have}
\PY{n}{String}\PY{+w}{ }\PY{n}{s}\PY{+w}{ }\PY{o}{=}\PY{+w}{ }\PY{l+s}{\PYZdq{}}\PY{l+s}{hello}\PY{l+s}{\PYZdq{}}\PY{p}{;}
\PY{n}{Class}\PY{o}{\PYZlt{}}\PY{o}{?}\PY{o}{\PYZgt{}}\PY{+w}{ }\PY{n}{byInstance}\PY{+w}{ }\PY{o}{=}\PY{+w}{ }\PY{n}{s}\PY{p}{.}\PY{n+na}{getClass}\PY{p}{(}\PY{p}{)}\PY{p}{;}

\PY{c+c1}{// 3. By name, at runtime — this is the \PYZdq{}dynamic\PYZdq{} one, used by frameworks}
\PY{n}{Class}\PY{o}{\PYZlt{}}\PY{o}{?}\PY{o}{\PYZgt{}}\PY{+w}{ }\PY{n}{byName}\PY{+w}{ }\PY{o}{=}\PY{+w}{ }\PY{n}{Class}\PY{p}{.}\PY{n+na}{forName}\PY{p}{(}\PY{l+s}{\PYZdq{}}\PY{l+s}{java.lang.String}\PY{l+s}{\PYZdq{}}\PY{p}{)}\PY{p}{;}
\end{Verbatim}
\end{codebox}

\code{Class.\allowbreak{}forName} is special: it is the classic way to load a class whose name is only known as a \code{String} (for example, read from a config file or a JDBC driver property). It also triggers static initialization of that class.

\subsection[Inspecting Members: Fields, Methods, Constructors]{Inspecting Members: Fields, Methods, Constructors}
\label{h:16:inspecting-members-fields-methods-constructors}
Once you have a \code{Class} object, you can list its members.

\begin{codebox}
\begin{Verbatim}[commandchars=\\\{\}]
\PY{k+kn}{import}\PY{+w}{ }\PY{n+nn}{java.lang.reflect.*}\PY{p}{;}

\PY{k+kd}{public}\PY{+w}{ }\PY{k+kd}{class} \PY{n+nc}{Inspector}\PY{+w}{ }\PY{p}{\PYZob{}}
\PY{+w}{    }\PY{k+kd}{static}\PY{+w}{ }\PY{k+kd}{class} \PY{n+nc}{Person}\PY{+w}{ }\PY{p}{\PYZob{}}
\PY{+w}{        }\PY{k+kd}{private}\PY{+w}{ }\PY{n}{String}\PY{+w}{ }\PY{n}{name}\PY{p}{;}
\PY{+w}{        }\PY{k+kd}{public}\PY{+w}{ }\PY{n+nf}{Person}\PY{p}{(}\PY{n}{String}\PY{+w}{ }\PY{n}{name}\PY{p}{)}\PY{+w}{ }\PY{p}{\PYZob{}}\PY{+w}{ }\PY{k}{this}\PY{p}{.}\PY{n+na}{name}\PY{+w}{ }\PY{o}{=}\PY{+w}{ }\PY{n}{name}\PY{p}{;}\PY{+w}{ }\PY{p}{\PYZcb{}}
\PY{+w}{        }\PY{k+kd}{public}\PY{+w}{ }\PY{n}{String}\PY{+w}{ }\PY{n+nf}{greet}\PY{p}{(}\PY{p}{)}\PY{+w}{ }\PY{p}{\PYZob{}}\PY{+w}{ }\PY{k}{return}\PY{+w}{ }\PY{l+s}{\PYZdq{}}\PY{l+s}{Hi, }\PY{l+s}{\PYZdq{}}\PY{+w}{ }\PY{o}{+}\PY{+w}{ }\PY{n}{name}\PY{p}{;}\PY{+w}{ }\PY{p}{\PYZcb{}}
\PY{+w}{        }\PY{k+kd}{private}\PY{+w}{ }\PY{k+kt}{void}\PY{+w}{ }\PY{n+nf}{secret}\PY{p}{(}\PY{p}{)}\PY{+w}{ }\PY{p}{\PYZob{}}\PY{+w}{ }\PY{p}{\PYZcb{}}
\PY{+w}{    }\PY{p}{\PYZcb{}}

\PY{+w}{    }\PY{k+kd}{public}\PY{+w}{ }\PY{k+kd}{static}\PY{+w}{ }\PY{k+kt}{void}\PY{+w}{ }\PY{n+nf}{main}\PY{p}{(}\PY{n}{String}\PY{o}{[}\PY{o}{]}\PY{+w}{ }\PY{n}{args}\PY{p}{)}\PY{+w}{ }\PY{k+kd}{throws}\PY{+w}{ }\PY{n}{Exception}\PY{+w}{ }\PY{p}{\PYZob{}}
\PY{+w}{        }\PY{n}{Class}\PY{o}{\PYZlt{}}\PY{o}{?}\PY{o}{\PYZgt{}}\PY{+w}{ }\PY{n}{clazz}\PY{+w}{ }\PY{o}{=}\PY{+w}{ }\PY{n}{Person}\PY{p}{.}\PY{n+na}{class}\PY{p}{;}

\PY{+w}{        }\PY{k}{for}\PY{+w}{ }\PY{p}{(}\PY{n}{Field}\PY{+w}{ }\PY{n}{f}\PY{+w}{ }\PY{p}{:}\PY{+w}{ }\PY{n}{clazz}\PY{p}{.}\PY{n+na}{getDeclaredFields}\PY{p}{(}\PY{p}{)}\PY{p}{)}\PY{+w}{ }\PY{p}{\PYZob{}}
\PY{+w}{            }\PY{n}{System}\PY{p}{.}\PY{n+na}{out}\PY{p}{.}\PY{n+na}{println}\PY{p}{(}\PY{l+s}{\PYZdq{}}\PY{l+s}{field: }\PY{l+s}{\PYZdq{}}\PY{+w}{ }\PY{o}{+}\PY{+w}{ }\PY{n}{f}\PY{p}{.}\PY{n+na}{getName}\PY{p}{(}\PY{p}{)}\PY{+w}{ }\PY{o}{+}\PY{+w}{ }\PY{l+s}{\PYZdq{}}\PY{l+s}{ : }\PY{l+s}{\PYZdq{}}\PY{+w}{ }\PY{o}{+}\PY{+w}{ }\PY{n}{f}\PY{p}{.}\PY{n+na}{getType}\PY{p}{(}\PY{p}{)}\PY{p}{)}\PY{p}{;}
\PY{+w}{        }\PY{p}{\PYZcb{}}
\PY{+w}{        }\PY{k}{for}\PY{+w}{ }\PY{p}{(}\PY{n}{Method}\PY{+w}{ }\PY{n}{m}\PY{+w}{ }\PY{p}{:}\PY{+w}{ }\PY{n}{clazz}\PY{p}{.}\PY{n+na}{getDeclaredMethods}\PY{p}{(}\PY{p}{)}\PY{p}{)}\PY{+w}{ }\PY{p}{\PYZob{}}
\PY{+w}{            }\PY{n}{System}\PY{p}{.}\PY{n+na}{out}\PY{p}{.}\PY{n+na}{println}\PY{p}{(}\PY{l+s}{\PYZdq{}}\PY{l+s}{method: }\PY{l+s}{\PYZdq{}}\PY{+w}{ }\PY{o}{+}\PY{+w}{ }\PY{n}{m}\PY{p}{.}\PY{n+na}{getName}\PY{p}{(}\PY{p}{)}\PY{p}{)}\PY{p}{;}
\PY{+w}{        }\PY{p}{\PYZcb{}}
\PY{+w}{        }\PY{k}{for}\PY{+w}{ }\PY{p}{(}\PY{n}{Constructor}\PY{o}{\PYZlt{}}\PY{o}{?}\PY{o}{\PYZgt{}}\PY{+w}{ }\PY{n}{c}\PY{+w}{ }\PY{p}{:}\PY{+w}{ }\PY{n}{clazz}\PY{p}{.}\PY{n+na}{getDeclaredConstructors}\PY{p}{(}\PY{p}{)}\PY{p}{)}\PY{+w}{ }\PY{p}{\PYZob{}}
\PY{+w}{            }\PY{n}{System}\PY{p}{.}\PY{n+na}{out}\PY{p}{.}\PY{n+na}{println}\PY{p}{(}\PY{l+s}{\PYZdq{}}\PY{l+s}{constructor: }\PY{l+s}{\PYZdq{}}\PY{+w}{ }\PY{o}{+}\PY{+w}{ }\PY{n}{c}\PY{p}{)}\PY{p}{;}
\PY{+w}{        }\PY{p}{\PYZcb{}}
\PY{+w}{    }\PY{p}{\PYZcb{}}
\PY{p}{\PYZcb{}}
\end{Verbatim}
\end{codebox}

\subsection[getDeclaredX vs getX]{\code{getDeclaredX} vs \code{getX}}
\label{h:16:getdeclaredx-vs-getx}
This is one of the most commonly confused pairs in reflection:

{\tablefontsmall
\begin{xltabular}{\linewidth}{@{}L{1.265}L{0.958}L{0.777}@{}}
\toprule
\rowcolor{tablehdr}
\thd{Method} & \thd{Includes inherited members?} & \thd{Includes private/protected members?} \\
\midrule
\endhead
\code{getFields()} / \code{getMethods()} / \code{getConstructors()} & Yes (public members from superclasses/interfaces too) & No — public only \\
\code{getDeclaredFields()} / \code{getDeclaredMethods()} / \code{get\allowbreak{}Declared\allowbreak{}Constructors()} & No — only members declared directly in this class & Yes — all of them, regardless of visibility \\
\bottomrule
\end{xltabular}
}

Rule of thumb: use \code{getDeclaredX} when you need to see \emph{everything this exact class wrote}, including private fields. Use \code{getX} when you want the public API surface, including what it inherited.

\subsection[Instantiating and Invoking]{Instantiating and Invoking}
\label{h:16:instantiating-and-invoking}
\begin{codebox}
\begin{Verbatim}[commandchars=\\\{\}]
\PY{k+kn}{import}\PY{+w}{ }\PY{n+nn}{java.lang.reflect.*}\PY{p}{;}

\PY{k+kd}{public}\PY{+w}{ }\PY{k+kd}{class} \PY{n+nc}{ReflectiveCall}\PY{+w}{ }\PY{p}{\PYZob{}}
\PY{+w}{    }\PY{k+kd}{public}\PY{+w}{ }\PY{k+kd}{static}\PY{+w}{ }\PY{k+kt}{void}\PY{+w}{ }\PY{n+nf}{main}\PY{p}{(}\PY{n}{String}\PY{o}{[}\PY{o}{]}\PY{+w}{ }\PY{n}{args}\PY{p}{)}\PY{+w}{ }\PY{k+kd}{throws}\PY{+w}{ }\PY{n}{Exception}\PY{+w}{ }\PY{p}{\PYZob{}}
\PY{+w}{        }\PY{n}{Class}\PY{o}{\PYZlt{}}\PY{o}{?}\PY{o}{\PYZgt{}}\PY{+w}{ }\PY{n}{clazz}\PY{+w}{ }\PY{o}{=}\PY{+w}{ }\PY{n}{Class}\PY{p}{.}\PY{n+na}{forName}\PY{p}{(}\PY{l+s}{\PYZdq{}}\PY{l+s}{java.util.ArrayList}\PY{l+s}{\PYZdq{}}\PY{p}{)}\PY{p}{;}

\PY{+w}{        }\PY{c+c1}{// Instantiate via a constructor}
\PY{+w}{        }\PY{n}{Constructor}\PY{o}{\PYZlt{}}\PY{o}{?}\PY{o}{\PYZgt{}}\PY{+w}{ }\PY{n}{ctor}\PY{+w}{ }\PY{o}{=}\PY{+w}{ }\PY{n}{clazz}\PY{p}{.}\PY{n+na}{getDeclaredConstructor}\PY{p}{(}\PY{p}{)}\PY{p}{;}
\PY{+w}{        }\PY{n}{Object}\PY{+w}{ }\PY{n}{list}\PY{+w}{ }\PY{o}{=}\PY{+w}{ }\PY{n}{ctor}\PY{p}{.}\PY{n+na}{newInstance}\PY{p}{(}\PY{p}{)}\PY{p}{;}

\PY{+w}{        }\PY{c+c1}{// Invoke a method reflectively}
\PY{+w}{        }\PY{n}{Method}\PY{+w}{ }\PY{n}{add}\PY{+w}{ }\PY{o}{=}\PY{+w}{ }\PY{n}{clazz}\PY{p}{.}\PY{n+na}{getMethod}\PY{p}{(}\PY{l+s}{\PYZdq{}}\PY{l+s}{add}\PY{l+s}{\PYZdq{}}\PY{p}{,}\PY{+w}{ }\PY{n}{Object}\PY{p}{.}\PY{n+na}{class}\PY{p}{)}\PY{p}{;}
\PY{+w}{        }\PY{n}{add}\PY{p}{.}\PY{n+na}{invoke}\PY{p}{(}\PY{n}{list}\PY{p}{,}\PY{+w}{ }\PY{l+s}{\PYZdq{}}\PY{l+s}{hello}\PY{l+s}{\PYZdq{}}\PY{p}{)}\PY{p}{;}

\PY{+w}{        }\PY{n}{Method}\PY{+w}{ }\PY{n}{size}\PY{+w}{ }\PY{o}{=}\PY{+w}{ }\PY{n}{clazz}\PY{p}{.}\PY{n+na}{getMethod}\PY{p}{(}\PY{l+s}{\PYZdq{}}\PY{l+s}{size}\PY{l+s}{\PYZdq{}}\PY{p}{)}\PY{p}{;}
\PY{+w}{        }\PY{n}{System}\PY{p}{.}\PY{n+na}{out}\PY{p}{.}\PY{n+na}{println}\PY{p}{(}\PY{n}{size}\PY{p}{.}\PY{n+na}{invoke}\PY{p}{(}\PY{n}{list}\PY{p}{)}\PY{p}{)}\PY{p}{;}\PY{+w}{ }\PY{c+c1}{// 1}
\PY{+w}{    }\PY{p}{\PYZcb{}}
\PY{p}{\PYZcb{}}
\end{Verbatim}
\end{codebox}

\subsection[setAccessible(true) and Strong Encapsulation (JDK 16/17+)]{\code{setAccessible(\allowbreak{}true)} and Strong Encapsulation (JDK 16/17+)}
\label{h:16:setaccessibletrue-and-strong-encapsulation-jdk-1617}
Reflection can normally only call public members. To reach private ones, you call \code{setAccessible(\allowbreak{}true)} on the \code{Field}, \code{Method}, or \code{Constructor} object, which tells the JVM to skip its usual access checks.

\begin{codebox}
\begin{Verbatim}[commandchars=\\\{\}]
\PY{n}{Field}\PY{+w}{ }\PY{n}{nameField}\PY{+w}{ }\PY{o}{=}\PY{+w}{ }\PY{n}{Person}\PY{p}{.}\PY{n+na}{class}\PY{p}{.}\PY{n+na}{getDeclaredField}\PY{p}{(}\PY{l+s}{\PYZdq{}}\PY{l+s}{name}\PY{l+s}{\PYZdq{}}\PY{p}{)}\PY{p}{;}
\PY{n}{nameField}\PY{p}{.}\PY{n+na}{setAccessible}\PY{p}{(}\PY{k+kc}{true}\PY{p}{)}\PY{p}{;}
\PY{n}{Person}\PY{+w}{ }\PY{n}{p}\PY{+w}{ }\PY{o}{=}\PY{+w}{ }\PY{k}{new}\PY{+w}{ }\PY{n}{Person}\PY{p}{(}\PY{l+s}{\PYZdq{}}\PY{l+s}{Ada}\PY{l+s}{\PYZdq{}}\PY{p}{)}\PY{p}{;}
\PY{n}{System}\PY{p}{.}\PY{n+na}{out}\PY{p}{.}\PY{n+na}{println}\PY{p}{(}\PY{n}{nameField}\PY{p}{.}\PY{n+na}{get}\PY{p}{(}\PY{n}{p}\PY{p}{)}\PY{p}{)}\PY{p}{;}\PY{+w}{ }\PY{c+c1}{// \PYZdq{}Ada\PYZdq{} — bypassing private}
\end{Verbatim}
\end{codebox}

Since JDK 9, the Java Platform Module System (JPMS) introduced \textbf{strong encapsulation}: modules can hide their internal packages so that even reflection cannot break in, unless the module explicitly “opens” that package. Starting with \textbf{JDK 16 (illegal access denied by default) and fully enforced by JDK 17}, calling \code{setAccessible(\allowbreak{}true)} on a member of an internal JDK class (like something in \code{java.\allowbreak{}util} internals) throws an \code{Inaccessible\allowbreak{}Object\allowbreak{}Exception} instead of just printing a warning as it did in JDK 9–15.

\begin{codebox}
\begin{Verbatim}
Exception in thread "main" java.lang.reflect.InaccessibleObjectException:
Unable to make field private final byte[] java.lang.String.value accessible:
module java.base does not "opens java.lang" to unnamed module
\end{Verbatim}
\end{codebox}

To allow it anyway, you must explicitly open the package, either in \code{module-\allowbreak{}info.\allowbreak{}java}:

\begin{codebox}
\begin{Verbatim}[commandchars=\\\{\}]
\PY{n}{module}\PY{+w}{ }\PY{n}{my}\PY{p}{.}\PY{n+na}{app}\PY{+w}{ }\PY{p}{\PYZob{}}
\PY{+w}{    }\PY{c+c1}{// grants reflective access to internal fields at runtime only}
\PY{+w}{    }\PY{n}{opens}\PY{+w}{ }\PY{n}{com}\PY{p}{.}\PY{n+na}{internal}\PY{p}{.}\PY{n+na}{pkg}\PY{+w}{ }\PY{n}{to}\PY{+w}{ }\PY{n}{some}\PY{p}{.}\PY{n+na}{other}\PY{p}{.}\PY{n+na}{module}\PY{p}{;}
\PY{p}{\PYZcb{}}
\end{Verbatim}
\end{codebox}

or via a JVM flag at launch time, which is the common escape hatch for libraries (like older Mockito or Jackson versions) that still reflect into JDK internals:

\begin{codebox}
\begin{Verbatim}
java --add-opens java.base/java.lang=ALL-UNNAMED -jar app.jar
\end{Verbatim}
\end{codebox}

This matters in real projects: upgrading from Java 8/11 to 17+ is a common source of \code{Inaccessible\allowbreak{}Object\allowbreak{}Exception} failures in libraries that reflect into JDK internals.

\subsection[Performance Cost]{Performance Cost}
\label{h:16:performance-cost}
Reflective calls are slower than direct calls because the JVM must perform access checks, resolve the target dynamically, and (for \code{invoke}) box/unbox arguments through \code{Object[]}. The JIT compiler also has a much harder time inlining and optimizing through \code{Method.\allowbreak{}invoke}. In hot paths (tight loops, request handling), reflection can be 10-100x slower than a direct call. Frameworks mitigate this by caching \code{Method}/\code{Field} objects (looking them up once, not per call) and by preferring \code{MethodHandle} (covered later) which the JIT can optimize almost as well as direct calls.

\subsection[Generic Type Info at Runtime]{Generic Type Info at Runtime}
\label{h:16:generic-type-info-at-runtime}
Because of \textbf{type erasure} (see the Generics chapter), a \code{List\textless{}\allowbreak{}String\textgreater{}} and a \code{List\textless{}\allowbreak{}Integer\textgreater{}} look identical at runtime — the generic type argument is gone from the object. However, reflection can still recover generic type information that was declared on a \textbf{field, method, or superclass}, because that information is stored in the class file, not on the runtime object.

\begin{codebox}
\begin{Verbatim}[commandchars=\\\{\}]
\PY{k+kn}{import}\PY{+w}{ }\PY{n+nn}{java.lang.reflect.*}\PY{p}{;}
\PY{k+kn}{import}\PY{+w}{ }\PY{n+nn}{java.util.*}\PY{p}{;}

\PY{k+kd}{class} \PY{n+nc}{StringListHolder}\PY{+w}{ }\PY{p}{\PYZob{}}
\PY{+w}{    }\PY{k+kd}{private}\PY{+w}{ }\PY{n}{List}\PY{o}{\PYZlt{}}\PY{n}{String}\PY{o}{\PYZgt{}}\PY{+w}{ }\PY{n}{names}\PY{p}{;}
\PY{p}{\PYZcb{}}

\PY{k+kd}{public}\PY{+w}{ }\PY{k+kd}{class} \PY{n+nc}{GenericInspection}\PY{+w}{ }\PY{p}{\PYZob{}}
\PY{+w}{    }\PY{k+kd}{public}\PY{+w}{ }\PY{k+kd}{static}\PY{+w}{ }\PY{k+kt}{void}\PY{+w}{ }\PY{n+nf}{main}\PY{p}{(}\PY{n}{String}\PY{o}{[}\PY{o}{]}\PY{+w}{ }\PY{n}{args}\PY{p}{)}\PY{+w}{ }\PY{k+kd}{throws}\PY{+w}{ }\PY{n}{Exception}\PY{+w}{ }\PY{p}{\PYZob{}}
\PY{+w}{        }\PY{n}{Field}\PY{+w}{ }\PY{n}{f}\PY{+w}{ }\PY{o}{=}\PY{+w}{ }\PY{n}{StringListHolder}\PY{p}{.}\PY{n+na}{class}\PY{p}{.}\PY{n+na}{getDeclaredField}\PY{p}{(}\PY{l+s}{\PYZdq{}}\PY{l+s}{names}\PY{l+s}{\PYZdq{}}\PY{p}{)}\PY{p}{;}
\PY{+w}{        }\PY{n}{Type}\PY{+w}{ }\PY{n}{genericType}\PY{+w}{ }\PY{o}{=}\PY{+w}{ }\PY{n}{f}\PY{p}{.}\PY{n+na}{getGenericType}\PY{p}{(}\PY{p}{)}\PY{p}{;}\PY{+w}{ }\PY{c+c1}{// java.util.List\PYZlt{}java.lang.String\PYZgt{}}
\PY{+w}{        }\PY{k}{if}\PY{+w}{ }\PY{p}{(}\PY{n}{genericType}\PY{+w}{ }\PY{k}{instanceof}\PY{+w}{ }\PY{n}{ParameterizedType}\PY{+w}{ }\PY{n}{pt}\PY{p}{)}\PY{+w}{ }\PY{p}{\PYZob{}}
\PY{+w}{            }\PY{n}{System}\PY{p}{.}\PY{n+na}{out}\PY{p}{.}\PY{n+na}{println}\PY{p}{(}\PY{n}{pt}\PY{p}{.}\PY{n+na}{getActualTypeArguments}\PY{p}{(}\PY{p}{)}\PY{o}{[}\PY{l+m+mi}{0}\PY{o}{]}\PY{p}{)}\PY{p}{;}\PY{+w}{ }\PY{c+c1}{// class java.lang.String}
\PY{+w}{        }\PY{p}{\PYZcb{}}
\PY{+w}{    }\PY{p}{\PYZcb{}}
\PY{p}{\PYZcb{}}
\end{Verbatim}
\end{codebox}

This trick — capturing a generic type through a subclass so it survives erasure — is called the \textbf{TypeToken pattern} (also called “super type tokens”), and it is how libraries like Gson and Guice let you say “give me a \code{List\textless{}\allowbreak{}String\textgreater{}} at runtime”:

\begin{codebox}
\begin{Verbatim}[commandchars=\\\{\}]
\PY{k+kn}{import}\PY{+w}{ }\PY{n+nn}{java.lang.reflect.*}\PY{p}{;}

\PY{k+kd}{public}\PY{+w}{ }\PY{k+kd}{abstract}\PY{+w}{ }\PY{k+kd}{class} \PY{n+nc}{TypeToken}\PY{o}{\PYZlt{}}\PY{n}{T}\PY{o}{\PYZgt{}}\PY{+w}{ }\PY{p}{\PYZob{}}
\PY{+w}{    }\PY{k+kd}{private}\PY{+w}{ }\PY{k+kd}{final}\PY{+w}{ }\PY{n}{Type}\PY{+w}{ }\PY{n}{type}\PY{p}{;}

\PY{+w}{    }\PY{k+kd}{protected}\PY{+w}{ }\PY{n+nf}{TypeToken}\PY{p}{(}\PY{p}{)}\PY{+w}{ }\PY{p}{\PYZob{}}
\PY{+w}{        }\PY{c+c1}{// getGenericSuperclass() sees TypeToken\PYZlt{}List\PYZlt{}String\PYZgt{}\PYZgt{} because}
\PY{+w}{        }\PY{c+c1}{// we captured it in an anonymous subclass below}
\PY{+w}{        }\PY{n}{ParameterizedType}\PY{+w}{ }\PY{n}{pt}\PY{+w}{ }\PY{o}{=}\PY{+w}{ }\PY{p}{(}\PY{n}{ParameterizedType}\PY{p}{)}\PY{+w}{ }\PY{n}{getClass}\PY{p}{(}\PY{p}{)}\PY{p}{.}\PY{n+na}{getGenericSuperclass}\PY{p}{(}\PY{p}{)}\PY{p}{;}
\PY{+w}{        }\PY{k}{this}\PY{p}{.}\PY{n+na}{type}\PY{+w}{ }\PY{o}{=}\PY{+w}{ }\PY{n}{pt}\PY{p}{.}\PY{n+na}{getActualTypeArguments}\PY{p}{(}\PY{p}{)}\PY{o}{[}\PY{l+m+mi}{0}\PY{o}{]}\PY{p}{;}
\PY{+w}{    }\PY{p}{\PYZcb{}}

\PY{+w}{    }\PY{k+kd}{public}\PY{+w}{ }\PY{n}{Type}\PY{+w}{ }\PY{n+nf}{getType}\PY{p}{(}\PY{p}{)}\PY{+w}{ }\PY{p}{\PYZob{}}\PY{+w}{ }\PY{k}{return}\PY{+w}{ }\PY{n}{type}\PY{p}{;}\PY{+w}{ }\PY{p}{\PYZcb{}}
\PY{p}{\PYZcb{}}

\PY{c+c1}{// usage: an anonymous subclass \PYZdq{}bakes in\PYZdq{} the type argument into the class file}
\PY{n}{TypeToken}\PY{o}{\PYZlt{}}\PY{n}{List}\PY{o}{\PYZlt{}}\PY{n}{String}\PY{o}{\PYZgt{}}\PY{o}{\PYZgt{}}\PY{+w}{ }\PY{n}{token}\PY{+w}{ }\PY{o}{=}\PY{+w}{ }\PY{k}{new}\PY{+w}{ }\PY{n}{TypeToken}\PY{o}{\PYZlt{}}\PY{n}{List}\PY{o}{\PYZlt{}}\PY{n}{String}\PY{o}{\PYZgt{}}\PY{o}{\PYZgt{}}\PY{p}{(}\PY{p}{)}\PY{+w}{ }\PY{p}{\PYZob{}}\PY{p}{\PYZcb{}}\PY{p}{;}
\PY{n}{System}\PY{p}{.}\PY{n+na}{out}\PY{p}{.}\PY{n+na}{println}\PY{p}{(}\PY{n}{token}\PY{p}{.}\PY{n+na}{getType}\PY{p}{(}\PY{p}{)}\PY{p}{)}\PY{p}{;}\PY{+w}{ }\PY{c+c1}{// java.util.List\PYZlt{}java.lang.String\PYZgt{}}
\end{Verbatim}
\end{codebox}

\subsection[When Reflection Is Legitimate vs a Smell]{When Reflection Is Legitimate vs a Smell}
\label{h:16:when-reflection-is-legitimate-vs-a-smell}
Legitimate uses:

\begin{itemize}
\item 
\textbf{Frameworks and dependency injection} (Spring, Guice) — wiring objects together without every class knowing about every other class.

\item 
\textbf{Serialization libraries} (Jackson, Gson) — reading/writing arbitrary object fields generically.

\item 
\textbf{Testing tools} (Mockito, JUnit) — creating mocks and invoking test methods discovered by annotation.

\item 
\textbf{ORMs} (Hibernate/JPA) — mapping database columns to entity fields.

\end{itemize}

Code smell / red flag in a code review:

\begin{itemize}
\item 
Application (non-framework) business logic reaching into another class’s private fields via reflection instead of adding a getter or constructor parameter. This usually signals a design problem — two classes are too tightly coupled, or an API is missing.

\item 
Reflection used to bypass \code{final} or visibility rules “just to make the compiler happy” instead of fixing the underlying design.

\end{itemize}

\section[Annotations]{Annotations}
\label{h:16:annotations}
An \textbf{annotation} is metadata attached to code — a class, method, field, or parameter — that does nothing by itself. It only has an effect if something (the compiler, a runtime library, or an annotation processor) reads it and acts on it.

\subsection[Declaring an Annotation]{Declaring an Annotation}
\label{h:16:declaring-an-annotation}
\begin{codebox}
\begin{Verbatim}[commandchars=\\\{\}]
\PY{k+kn}{import}\PY{+w}{ }\PY{n+nn}{java.lang.annotation.*}\PY{p}{;}

\PY{n+nd}{@Retention}\PY{p}{(}\PY{n}{RetentionPolicy}\PY{p}{.}\PY{n+na}{RUNTIME}\PY{p}{)}
\PY{n+nd}{@Target}\PY{p}{(}\PY{n}{ElementType}\PY{p}{.}\PY{n+na}{METHOD}\PY{p}{)}
\PY{k+kd}{public}\PY{+w}{ }\PY{n+nd}{@interface}\PY{+w}{ }\PY{n}{Timed}\PY{+w}{ }\PY{p}{\PYZob{}}
\PY{+w}{    }\PY{n}{String}\PY{+w}{ }\PY{n+nf}{label}\PY{p}{(}\PY{p}{)}\PY{+w}{ }\PY{k}{default}\PY{+w}{ }\PY{l+s}{\PYZdq{}}\PY{l+s}{\PYZdq{}}\PY{p}{;}
\PY{p}{\PYZcb{}}
\end{Verbatim}
\end{codebox}

Usage:

\begin{codebox}
\begin{Verbatim}[commandchars=\\\{\}]
\PY{k+kd}{public}\PY{+w}{ }\PY{k+kd}{class} \PY{n+nc}{OrderService}\PY{+w}{ }\PY{p}{\PYZob{}}
\PY{+w}{    }\PY{n+nd}{@Timed}\PY{p}{(}\PY{n}{label}\PY{+w}{ }\PY{o}{=}\PY{+w}{ }\PY{l+s}{\PYZdq{}}\PY{l+s}{checkout}\PY{l+s}{\PYZdq{}}\PY{p}{)}
\PY{+w}{    }\PY{k+kd}{public}\PY{+w}{ }\PY{k+kt}{void}\PY{+w}{ }\PY{n+nf}{checkout}\PY{p}{(}\PY{p}{)}\PY{+w}{ }\PY{p}{\PYZob{}}\PY{+w}{ }\PY{c+cm}{/* ... */}\PY{+w}{ }\PY{p}{\PYZcb{}}
\PY{p}{\PYZcb{}}
\end{Verbatim}
\end{codebox}

\subsection[Meta-Annotations]{Meta-Annotations}
\label{h:16:meta-annotations}
Meta-annotations are annotations that describe \emph{other annotations}. The main ones:

\begin{itemize}
\item 
\textbf{\code{@\allowbreak{}Retention}} — how long the annotation is kept.

\item 
\textbf{\code{@\allowbreak{}Target}} — what kinds of declarations it can be placed on (methods, fields, types, parameters…).

\item 
\textbf{\code{@\allowbreak{}Inherited}} — if placed on an annotation, subclasses automatically “inherit” it from their superclass (only applies to class-level annotations, not methods or fields).

\item 
\textbf{\code{@\allowbreak{}Repeatable}} — allows the same annotation to be applied more than once to the same element.

\item 
\textbf{\code{@\allowbreak{}Documented}} — includes the annotation in generated Javadoc.

\end{itemize}

\begin{codebox}
\begin{Verbatim}[commandchars=\\\{\}]
\PY{n+nd}{@Inherited}
\PY{n+nd}{@Retention}\PY{p}{(}\PY{n}{RetentionPolicy}\PY{p}{.}\PY{n+na}{RUNTIME}\PY{p}{)}
\PY{n+nd}{@Target}\PY{p}{(}\PY{n}{ElementType}\PY{p}{.}\PY{n+na}{TYPE}\PY{p}{)}
\PY{k+kd}{public}\PY{+w}{ }\PY{n+nd}{@interface}\PY{+w}{ }\PY{n}{Auditable}\PY{+w}{ }\PY{p}{\PYZob{}}\PY{+w}{ }\PY{p}{\PYZcb{}}

\PY{n+nd}{@Auditable}
\PY{k+kd}{class} \PY{n+nc}{Base}\PY{+w}{ }\PY{p}{\PYZob{}}\PY{+w}{ }\PY{p}{\PYZcb{}}

\PY{k+kd}{class} \PY{n+nc}{Derived}\PY{+w}{ }\PY{k+kd}{extends}\PY{+w}{ }\PY{n}{Base}\PY{+w}{ }\PY{p}{\PYZob{}}\PY{+w}{ }\PY{p}{\PYZcb{}}\PY{+w}{ }\PY{c+c1}{// Derived.class.isAnnotationPresent(Auditable.class) == true}
\end{Verbatim}
\end{codebox}

\begin{codebox}
\begin{Verbatim}[commandchars=\\\{\}]
\PY{k+kn}{import}\PY{+w}{ }\PY{n+nn}{java.lang.annotation.*}\PY{p}{;}

\PY{n+nd}{@Retention}\PY{p}{(}\PY{n}{RetentionPolicy}\PY{p}{.}\PY{n+na}{RUNTIME}\PY{p}{)}
\PY{k+kd}{public}\PY{+w}{ }\PY{n+nd}{@interface}\PY{+w}{ }\PY{n}{Schedules}\PY{+w}{ }\PY{p}{\PYZob{}}
\PY{+w}{    }\PY{n}{Schedule}\PY{o}{[}\PY{o}{]}\PY{+w}{ }\PY{n+nf}{value}\PY{p}{(}\PY{p}{)}\PY{p}{;}
\PY{p}{\PYZcb{}}

\PY{n+nd}{@Retention}\PY{p}{(}\PY{n}{RetentionPolicy}\PY{p}{.}\PY{n+na}{RUNTIME}\PY{p}{)}
\PY{n+nd}{@Repeatable}\PY{p}{(}\PY{n}{Schedules}\PY{p}{.}\PY{n+na}{class}\PY{p}{)}\PY{+w}{ }\PY{c+c1}{// container annotation required for @Repeatable}
\PY{k+kd}{public}\PY{+w}{ }\PY{n+nd}{@interface}\PY{+w}{ }\PY{n}{Schedule}\PY{+w}{ }\PY{p}{\PYZob{}}
\PY{+w}{    }\PY{n}{String}\PY{+w}{ }\PY{n+nf}{cron}\PY{p}{(}\PY{p}{)}\PY{p}{;}
\PY{p}{\PYZcb{}}

\PY{k+kd}{class} \PY{n+nc}{Job}\PY{+w}{ }\PY{p}{\PYZob{}}
\PY{+w}{    }\PY{n+nd}{@Schedule}\PY{p}{(}\PY{n}{cron}\PY{+w}{ }\PY{o}{=}\PY{+w}{ }\PY{l+s}{\PYZdq{}}\PY{l+s}{0 0 * * *}\PY{l+s}{\PYZdq{}}\PY{p}{)}
\PY{+w}{    }\PY{n+nd}{@Schedule}\PY{p}{(}\PY{n}{cron}\PY{+w}{ }\PY{o}{=}\PY{+w}{ }\PY{l+s}{\PYZdq{}}\PY{l+s}{0 12 * * *}\PY{l+s}{\PYZdq{}}\PY{p}{)}
\PY{+w}{    }\PY{k+kt}{void}\PY{+w}{ }\PY{n+nf}{run}\PY{p}{(}\PY{p}{)}\PY{+w}{ }\PY{p}{\PYZob{}}\PY{+w}{ }\PY{p}{\PYZcb{}}
\PY{p}{\PYZcb{}}
\end{Verbatim}
\end{codebox}

\subsection[Reading Annotations at Runtime]{Reading Annotations at Runtime}
\label{h:16:reading-annotations-at-runtime}
Only annotations with \code{Retention\allowbreak{}Policy.\allowbreak{}RUNTIME} are visible via reflection.

\begin{codebox}
\begin{Verbatim}[commandchars=\\\{\}]
\PY{k+kn}{import}\PY{+w}{ }\PY{n+nn}{java.lang.reflect.*}\PY{p}{;}

\PY{k+kd}{public}\PY{+w}{ }\PY{k+kd}{class} \PY{n+nc}{TimedRunner}\PY{+w}{ }\PY{p}{\PYZob{}}
\PY{+w}{    }\PY{k+kd}{public}\PY{+w}{ }\PY{k+kd}{static}\PY{+w}{ }\PY{k+kt}{void}\PY{+w}{ }\PY{n+nf}{main}\PY{p}{(}\PY{n}{String}\PY{o}{[}\PY{o}{]}\PY{+w}{ }\PY{n}{args}\PY{p}{)}\PY{+w}{ }\PY{k+kd}{throws}\PY{+w}{ }\PY{n}{Exception}\PY{+w}{ }\PY{p}{\PYZob{}}
\PY{+w}{        }\PY{n}{Method}\PY{+w}{ }\PY{n}{m}\PY{+w}{ }\PY{o}{=}\PY{+w}{ }\PY{n}{OrderService}\PY{p}{.}\PY{n+na}{class}\PY{p}{.}\PY{n+na}{getMethod}\PY{p}{(}\PY{l+s}{\PYZdq{}}\PY{l+s}{checkout}\PY{l+s}{\PYZdq{}}\PY{p}{)}\PY{p}{;}
\PY{+w}{        }\PY{k}{if}\PY{+w}{ }\PY{p}{(}\PY{n}{m}\PY{p}{.}\PY{n+na}{isAnnotationPresent}\PY{p}{(}\PY{n}{Timed}\PY{p}{.}\PY{n+na}{class}\PY{p}{)}\PY{p}{)}\PY{+w}{ }\PY{p}{\PYZob{}}
\PY{+w}{            }\PY{n}{Timed}\PY{+w}{ }\PY{n}{timed}\PY{+w}{ }\PY{o}{=}\PY{+w}{ }\PY{n}{m}\PY{p}{.}\PY{n+na}{getAnnotation}\PY{p}{(}\PY{n}{Timed}\PY{p}{.}\PY{n+na}{class}\PY{p}{)}\PY{p}{;}
\PY{+w}{            }\PY{n}{System}\PY{p}{.}\PY{n+na}{out}\PY{p}{.}\PY{n+na}{println}\PY{p}{(}\PY{l+s}{\PYZdq{}}\PY{l+s}{Label: }\PY{l+s}{\PYZdq{}}\PY{+w}{ }\PY{o}{+}\PY{+w}{ }\PY{n}{timed}\PY{p}{.}\PY{n+na}{label}\PY{p}{(}\PY{p}{)}\PY{p}{)}\PY{p}{;}\PY{+w}{ }\PY{c+c1}{// \PYZdq{}checkout\PYZdq{}}
\PY{+w}{        }\PY{p}{\PYZcb{}}
\PY{+w}{    }\PY{p}{\PYZcb{}}
\PY{p}{\PYZcb{}}
\end{Verbatim}
\end{codebox}

\subsection[Retention Policies]{Retention Policies}
\label{h:16:retention-policies}
{\tablefontsmall
\begin{xltabular}{\linewidth}{@{}L{0.597}L{0.597}L{1.015}L{1.791}@{}}
\toprule
\rowcolor{tablehdr}
\thd{Policy} & \thd{Kept in \code{.\allowbreak{}class} file?} & \thd{Visible to reflection at runtime?} & \thd{Typical use} \\
\midrule
\endhead
\code{SOURCE} & No — discarded by the compiler & No & Compile-time-only checks, e.g. \code{@\allowbreak{}Override}, Lombok’s markers, annotation processors that only need source info \\
\code{CLASS} (default if unspecified) & Yes & No — JVM doesn’t load it into memory for reflection & Bytecode-level tools (e.g. some bytecode weavers); rarely used directly by application code \\
\code{RUNTIME} & Yes & Yes & Frameworks that inspect annotations via reflection at runtime, e.g. \code{@\allowbreak{}Autowired}, \code{@\allowbreak{}Test}, \code{@\allowbreak{}Entity} \\
\bottomrule
\end{xltabular}
}

\subsection[Type Annotations (JDK 8+)]{Type Annotations (JDK 8+)}
\label{h:16:type-annotations-jdk-8}
Since Java 8, annotations can also be placed on \textbf{type uses}, not just declarations — for example, on a generic type argument, a cast, or a \code{throws} clause. This is mainly used by static analysis tools (like the Checker Framework) to catch bugs such as null-pointer risks before runtime.

\begin{codebox}
\begin{Verbatim}[commandchars=\\\{\}]
\PY{k+kn}{import}\PY{+w}{ }\PY{n+nn}{java.lang.annotation.*}\PY{p}{;}

\PY{n+nd}{@Target}\PY{p}{(}\PY{n}{ElementType}\PY{p}{.}\PY{n+na}{TYPE\PYZus{}USE}\PY{p}{)}
\PY{n+nd}{@interface}\PY{+w}{ }\PY{n}{NonNull}\PY{+w}{ }\PY{p}{\PYZob{}}\PY{+w}{ }\PY{p}{\PYZcb{}}

\PY{k+kd}{public}\PY{+w}{ }\PY{k+kd}{class} \PY{n+nc}{TypeAnnotationDemo}\PY{+w}{ }\PY{p}{\PYZob{}}
\PY{+w}{    }\PY{c+c1}{// annotation on the type argument, not the field itself}
\PY{+w}{    }\PY{k+kd}{private}\PY{+w}{ }\PY{n}{List}\PY{o}{\PYZlt{}}\PY{n+nd}{@NonNull}\PY{+w}{ }\PY{n}{String}\PY{o}{\PYZgt{}}\PY{+w}{ }\PY{n}{names}\PY{p}{;}

\PY{+w}{    }\PY{c+c1}{// annotation on a cast}
\PY{+w}{    }\PY{n}{Object}\PY{+w}{ }\PY{n}{o}\PY{+w}{ }\PY{o}{=}\PY{+w}{ }\PY{l+s}{\PYZdq{}}\PY{l+s}{text}\PY{l+s}{\PYZdq{}}\PY{p}{;}
\PY{+w}{    }\PY{n}{String}\PY{+w}{ }\PY{n}{s}\PY{+w}{ }\PY{o}{=}\PY{+w}{ }\PY{p}{(}\PY{n+nd}{@NonNull}\PY{+w}{ }\PY{n}{String}\PY{p}{)}\PY{+w}{ }\PY{n}{o}\PY{p}{;}
\PY{p}{\PYZcb{}}
\end{Verbatim}
\end{codebox}

\section[Annotation Processing]{Annotation Processing}
\label{h:16:annotation-processing}
\textbf{Annotation processing} happens at \emph{compile time}, not runtime. An annotation processor is a plug-in for the compiler (\code{javac}) that reads annotations in your source code and can generate new source files (or report errors), all before the program ever runs.

\subsection[The Processor Interface and AbstractProcessor]{The \code{Processor} Interface and \code{AbstractProcessor}}
\label{h:16:the-processor-interface-and-abstractprocessor}
\begin{codebox}
\begin{Verbatim}[commandchars=\\\{\}]
\PY{k+kn}{import}\PY{+w}{ }\PY{n+nn}{javax.annotation.processing.*}\PY{p}{;}
\PY{k+kn}{import}\PY{+w}{ }\PY{n+nn}{javax.lang.model.SourceVersion}\PY{p}{;}
\PY{k+kn}{import}\PY{+w}{ }\PY{n+nn}{javax.lang.model.element.*}\PY{p}{;}
\PY{k+kn}{import}\PY{+w}{ }\PY{n+nn}{java.util.Set}\PY{p}{;}

\PY{n+nd}{@SupportedAnnotationTypes}\PY{p}{(}\PY{l+s}{\PYZdq{}}\PY{l+s}{com.example.GenerateBuilder}\PY{l+s}{\PYZdq{}}\PY{p}{)}
\PY{n+nd}{@SupportedSourceVersion}\PY{p}{(}\PY{n}{SourceVersion}\PY{p}{.}\PY{n+na}{RELEASE\PYZus{}21}\PY{p}{)}
\PY{k+kd}{public}\PY{+w}{ }\PY{k+kd}{class} \PY{n+nc}{BuilderProcessor}\PY{+w}{ }\PY{k+kd}{extends}\PY{+w}{ }\PY{n}{AbstractProcessor}\PY{+w}{ }\PY{p}{\PYZob{}}

\PY{+w}{    }\PY{n+nd}{@Override}
\PY{+w}{    }\PY{k+kd}{public}\PY{+w}{ }\PY{k+kt}{boolean}\PY{+w}{ }\PY{n+nf}{process}\PY{p}{(}\PY{n}{Set}\PY{o}{\PYZlt{}}\PY{o}{?}\PY{+w}{ }\PY{k+kd}{extends}\PY{+w}{ }\PY{n}{TypeElement}\PY{o}{\PYZgt{}}\PY{+w}{ }\PY{n}{annotations}\PY{p}{,}\PY{+w}{ }\PY{n}{RoundEnvironment}\PY{+w}{ }\PY{n}{roundEnv}\PY{p}{)}\PY{+w}{ }\PY{p}{\PYZob{}}
\PY{+w}{        }\PY{k}{for}\PY{+w}{ }\PY{p}{(}\PY{n}{Element}\PY{+w}{ }\PY{n}{element}\PY{+w}{ }\PY{p}{:}\PY{+w}{ }\PY{n}{roundEnv}\PY{p}{.}\PY{n+na}{getElementsAnnotatedWith}\PY{p}{(}\PY{n}{GenerateBuilder}\PY{p}{.}\PY{n+na}{class}\PY{p}{)}\PY{p}{)}\PY{+w}{ }\PY{p}{\PYZob{}}
\PY{+w}{            }\PY{n}{TypeElement}\PY{+w}{ }\PY{n}{classElement}\PY{+w}{ }\PY{o}{=}\PY{+w}{ }\PY{p}{(}\PY{n}{TypeElement}\PY{p}{)}\PY{+w}{ }\PY{n}{element}\PY{p}{;}
\PY{+w}{            }\PY{n}{generateBuilderClass}\PY{p}{(}\PY{n}{classElement}\PY{p}{)}\PY{p}{;}
\PY{+w}{        }\PY{p}{\PYZcb{}}
\PY{+w}{        }\PY{k}{return}\PY{+w}{ }\PY{k+kc}{true}\PY{p}{;}\PY{+w}{ }\PY{c+c1}{// true = \PYZdq{}I claimed these annotations, don\PYZsq{}t let other processors see them\PYZdq{}}
\PY{+w}{    }\PY{p}{\PYZcb{}}

\PY{+w}{    }\PY{k+kd}{private}\PY{+w}{ }\PY{k+kt}{void}\PY{+w}{ }\PY{n+nf}{generateBuilderClass}\PY{p}{(}\PY{n}{TypeElement}\PY{+w}{ }\PY{n}{classElement}\PY{p}{)}\PY{+w}{ }\PY{p}{\PYZob{}}
\PY{+w}{        }\PY{c+c1}{// Use processingEnv.getFiler() to write a new .java file (shown below)}
\PY{+w}{    }\PY{p}{\PYZcb{}}
\PY{p}{\PYZcb{}}
\end{Verbatim}
\end{codebox}

\subsection[Rounds]{Rounds}
\label{h:16:rounds}
Annotation processing runs in \textbf{rounds}. Round 1 processes the original source files. If a processor generates new source files in round 1, \code{javac} runs another round to process \emph{those} generated files too (since they might themselves have annotations). This repeats until no new files are generated. \code{Round\allowbreak{}Environment.\allowbreak{}processing\allowbreak{}Over()} tells a processor when the final round has finished, useful for cleanup or validation that should only happen once.

\subsection[The Filer: Generating Code]{The \code{Filer}: Generating Code}
\label{h:16:the-filer-generating-code}
The \code{Filer} API (\code{processing\allowbreak{}Env.\allowbreak{}get\allowbreak{}Filer()}) is how a processor writes new source or resource files during compilation.

\begin{codebox}
\begin{Verbatim}[commandchars=\\\{\}]
\PY{k+kn}{import}\PY{+w}{ }\PY{n+nn}{javax.annotation.processing.Filer}\PY{p}{;}
\PY{k+kn}{import}\PY{+w}{ }\PY{n+nn}{javax.tools.JavaFileObject}\PY{p}{;}
\PY{k+kn}{import}\PY{+w}{ }\PY{n+nn}{java.io.Writer}\PY{p}{;}

\PY{k+kt}{void}\PY{+w}{ }\PY{n+nf}{writeGeneratedClass}\PY{p}{(}\PY{n}{String}\PY{+w}{ }\PY{n}{packageName}\PY{p}{,}\PY{+w}{ }\PY{n}{String}\PY{+w}{ }\PY{n}{className}\PY{p}{,}\PY{+w}{ }\PY{n}{String}\PY{+w}{ }\PY{n}{content}\PY{p}{)}\PY{+w}{ }\PY{k+kd}{throws}\PY{+w}{ }\PY{n}{Exception}\PY{+w}{ }\PY{p}{\PYZob{}}
\PY{+w}{    }\PY{n}{Filer}\PY{+w}{ }\PY{n}{filer}\PY{+w}{ }\PY{o}{=}\PY{+w}{ }\PY{n}{processingEnv}\PY{p}{.}\PY{n+na}{getFiler}\PY{p}{(}\PY{p}{)}\PY{p}{;}
\PY{+w}{    }\PY{n}{JavaFileObject}\PY{+w}{ }\PY{n}{file}\PY{+w}{ }\PY{o}{=}\PY{+w}{ }\PY{n}{filer}\PY{p}{.}\PY{n+na}{createSourceFile}\PY{p}{(}\PY{n}{packageName}\PY{+w}{ }\PY{o}{+}\PY{+w}{ }\PY{l+s}{\PYZdq{}}\PY{l+s}{.}\PY{l+s}{\PYZdq{}}\PY{+w}{ }\PY{o}{+}\PY{+w}{ }\PY{n}{className}\PY{p}{)}\PY{p}{;}
\PY{+w}{    }\PY{k}{try}\PY{+w}{ }\PY{p}{(}\PY{n}{Writer}\PY{+w}{ }\PY{n}{writer}\PY{+w}{ }\PY{o}{=}\PY{+w}{ }\PY{n}{file}\PY{p}{.}\PY{n+na}{openWriter}\PY{p}{(}\PY{p}{)}\PY{p}{)}\PY{+w}{ }\PY{p}{\PYZob{}}
\PY{+w}{        }\PY{n}{writer}\PY{p}{.}\PY{n+na}{write}\PY{p}{(}\PY{n}{content}\PY{p}{)}\PY{p}{;}
\PY{+w}{    }\PY{p}{\PYZcb{}}
\PY{p}{\PYZcb{}}
\end{Verbatim}
\end{codebox}

\subsection[Registering a Processor via META-INF/services]{Registering a Processor via \code{META-\allowbreak{}INF/\allowbreak{}services}}
\label{h:16:registering-a-processor-via-meta-infservices}
For \code{javac} to discover your processor automatically, register it as a service (the same SPI mechanism covered later in this chapter):

\begin{codebox}
\begin{Verbatim}
src/main/resources/META-INF/services/javax.annotation.processing.Processor
\end{Verbatim}
\end{codebox}

containing one line:

\begin{codebox}
\begin{Verbatim}
com.example.BuilderProcessor
\end{Verbatim}
\end{codebox}

Modern projects usually use Google’s \code{auto-\allowbreak{}service} library to generate this file automatically via \code{@\allowbreak{}Auto\allowbreak{}Service(\allowbreak{}Processor.\allowbreak{}class)}.

\subsection[Real-World Examples]{Real-World Examples}
\label{h:16:real-world-examples}
\begin{itemize}
\item 
\textbf{Lombok} — generates getters, setters, constructors, \code{equals}/\code{hashCode}, and builders from annotations like \code{@\allowbreak{}Data} and \code{@\allowbreak{}Builder}. (Technically it hooks into javac in a slightly unofficial way, modifying the in-memory AST rather than just generating new files, but conceptually it’s the same idea.)

\item 
\textbf{MapStruct} — generates type-safe mapping code between DTOs and entities from an \code{@\allowbreak{}Mapper} interface, avoiding hand-written or reflective mapping code.

\item 
\textbf{Dagger} — generates dependency-injection wiring code at compile time instead of using runtime reflection like Spring/Guice.

\item 
\textbf{Immutables} — generates immutable value classes from an \code{@\allowbreak{}Value.\allowbreak{}Immutable}-annotated abstract class or interface.

\end{itemize}

\subsection[Why Compile-Time Beats Runtime Reflection]{Why Compile-Time Beats Runtime Reflection}
\label{h:16:why-compile-time-beats-runtime-reflection}
\begin{itemize}
\item 
\textbf{Performance}: generated code is plain Java, as fast as hand-written code — no reflective lookups at runtime.

\item 
\textbf{Fail fast}: errors are caught at compile time (\code{javac} fails the build) instead of surfacing as a \code{RuntimeException} in production.

\item 
\textbf{Works with strong encapsulation}: no \code{setAccessible} calls, so no fights with the module system or native-image tools like GraalVM, which struggle with runtime reflection.

\item 
\textbf{IDE support}: generated code is visible, debuggable, and step-through-able, unlike a black-box reflective call.

\end{itemize}

\section[Dynamic Proxies]{Dynamic Proxies}
\label{h:16:dynamic-proxies}
A \textbf{dynamic proxy} is an object, generated at runtime, that implements one or more interfaces and forwards every method call to a handler you provide. It lets you add cross-cutting behavior (logging, timing, security checks, retries) around any interface without writing a wrapper class by hand for each one.

\subsection[Proxy.newProxyInstance + InvocationHandler]{\code{Proxy.\allowbreak{}newProxyInstance} + \code{InvocationHandler}}
\label{h:16:proxynewproxyinstance--invocationhandler}
\begin{codebox}
\begin{Verbatim}[commandchars=\\\{\}]
\PY{k+kn}{import}\PY{+w}{ }\PY{n+nn}{java.lang.reflect.*}\PY{p}{;}

\PY{k+kd}{interface} \PY{n+nc}{OrderService}\PY{+w}{ }\PY{p}{\PYZob{}}
\PY{+w}{    }\PY{k+kt}{void}\PY{+w}{ }\PY{n+nf}{placeOrder}\PY{p}{(}\PY{n}{String}\PY{+w}{ }\PY{n}{item}\PY{p}{)}\PY{p}{;}
\PY{+w}{    }\PY{k+kt}{int}\PY{+w}{ }\PY{n+nf}{getOrderCount}\PY{p}{(}\PY{p}{)}\PY{p}{;}
\PY{p}{\PYZcb{}}

\PY{k+kd}{class} \PY{n+nc}{OrderServiceImpl}\PY{+w}{ }\PY{k+kd}{implements}\PY{+w}{ }\PY{n}{OrderService}\PY{+w}{ }\PY{p}{\PYZob{}}
\PY{+w}{    }\PY{k+kd}{private}\PY{+w}{ }\PY{k+kt}{int}\PY{+w}{ }\PY{n}{count}\PY{+w}{ }\PY{o}{=}\PY{+w}{ }\PY{l+m+mi}{0}\PY{p}{;}
\PY{+w}{    }\PY{k+kd}{public}\PY{+w}{ }\PY{k+kt}{void}\PY{+w}{ }\PY{n+nf}{placeOrder}\PY{p}{(}\PY{n}{String}\PY{+w}{ }\PY{n}{item}\PY{p}{)}\PY{+w}{ }\PY{p}{\PYZob{}}
\PY{+w}{        }\PY{n}{count}\PY{o}{+}\PY{o}{+}\PY{p}{;}
\PY{+w}{        }\PY{n}{System}\PY{p}{.}\PY{n+na}{out}\PY{p}{.}\PY{n+na}{println}\PY{p}{(}\PY{l+s}{\PYZdq{}}\PY{l+s}{Order placed: }\PY{l+s}{\PYZdq{}}\PY{+w}{ }\PY{o}{+}\PY{+w}{ }\PY{n}{item}\PY{p}{)}\PY{p}{;}
\PY{+w}{    }\PY{p}{\PYZcb{}}
\PY{+w}{    }\PY{k+kd}{public}\PY{+w}{ }\PY{k+kt}{int}\PY{+w}{ }\PY{n+nf}{getOrderCount}\PY{p}{(}\PY{p}{)}\PY{+w}{ }\PY{p}{\PYZob{}}\PY{+w}{ }\PY{k}{return}\PY{+w}{ }\PY{n}{count}\PY{p}{;}\PY{+w}{ }\PY{p}{\PYZcb{}}
\PY{p}{\PYZcb{}}

\PY{k+kd}{class} \PY{n+nc}{LoggingTimingHandler}\PY{+w}{ }\PY{k+kd}{implements}\PY{+w}{ }\PY{n}{InvocationHandler}\PY{+w}{ }\PY{p}{\PYZob{}}
\PY{+w}{    }\PY{k+kd}{private}\PY{+w}{ }\PY{k+kd}{final}\PY{+w}{ }\PY{n}{Object}\PY{+w}{ }\PY{n}{target}\PY{p}{;}

\PY{+w}{    }\PY{n}{LoggingTimingHandler}\PY{p}{(}\PY{n}{Object}\PY{+w}{ }\PY{n}{target}\PY{p}{)}\PY{+w}{ }\PY{p}{\PYZob{}}\PY{+w}{ }\PY{k}{this}\PY{p}{.}\PY{n+na}{target}\PY{+w}{ }\PY{o}{=}\PY{+w}{ }\PY{n}{target}\PY{p}{;}\PY{+w}{ }\PY{p}{\PYZcb{}}

\PY{+w}{    }\PY{n+nd}{@Override}
\PY{+w}{    }\PY{k+kd}{public}\PY{+w}{ }\PY{n}{Object}\PY{+w}{ }\PY{n+nf}{invoke}\PY{p}{(}\PY{n}{Object}\PY{+w}{ }\PY{n}{proxy}\PY{p}{,}\PY{+w}{ }\PY{n}{Method}\PY{+w}{ }\PY{n}{method}\PY{p}{,}\PY{+w}{ }\PY{n}{Object}\PY{o}{[}\PY{o}{]}\PY{+w}{ }\PY{n}{args}\PY{p}{)}\PY{+w}{ }\PY{k+kd}{throws}\PY{+w}{ }\PY{n}{Throwable}\PY{+w}{ }\PY{p}{\PYZob{}}
\PY{+w}{        }\PY{k+kt}{long}\PY{+w}{ }\PY{n}{start}\PY{+w}{ }\PY{o}{=}\PY{+w}{ }\PY{n}{System}\PY{p}{.}\PY{n+na}{nanoTime}\PY{p}{(}\PY{p}{)}\PY{p}{;}
\PY{+w}{        }\PY{n}{System}\PY{p}{.}\PY{n+na}{out}\PY{p}{.}\PY{n+na}{println}\PY{p}{(}\PY{l+s}{\PYZdq{}}\PY{l+s}{Calling }\PY{l+s}{\PYZdq{}}\PY{+w}{ }\PY{o}{+}\PY{+w}{ }\PY{n}{method}\PY{p}{.}\PY{n+na}{getName}\PY{p}{(}\PY{p}{)}\PY{p}{)}\PY{p}{;}
\PY{+w}{        }\PY{k}{try}\PY{+w}{ }\PY{p}{\PYZob{}}
\PY{+w}{            }\PY{k}{return}\PY{+w}{ }\PY{n}{method}\PY{p}{.}\PY{n+na}{invoke}\PY{p}{(}\PY{n}{target}\PY{p}{,}\PY{+w}{ }\PY{n}{args}\PY{p}{)}\PY{p}{;}\PY{+w}{ }\PY{c+c1}{// forward to the real object}
\PY{+w}{        }\PY{p}{\PYZcb{}}\PY{+w}{ }\PY{k}{finally}\PY{+w}{ }\PY{p}{\PYZob{}}
\PY{+w}{            }\PY{k+kt}{long}\PY{+w}{ }\PY{n}{elapsedMicros}\PY{+w}{ }\PY{o}{=}\PY{+w}{ }\PY{p}{(}\PY{n}{System}\PY{p}{.}\PY{n+na}{nanoTime}\PY{p}{(}\PY{p}{)}\PY{+w}{ }\PY{o}{\PYZhy{}}\PY{+w}{ }\PY{n}{start}\PY{p}{)}\PY{+w}{ }\PY{o}{/}\PY{+w}{ }\PY{l+m+mi}{1000}\PY{p}{;}
\PY{+w}{            }\PY{n}{System}\PY{p}{.}\PY{n+na}{out}\PY{p}{.}\PY{n+na}{println}\PY{p}{(}\PY{n}{method}\PY{p}{.}\PY{n+na}{getName}\PY{p}{(}\PY{p}{)}\PY{+w}{ }\PY{o}{+}\PY{+w}{ }\PY{l+s}{\PYZdq{}}\PY{l+s}{ took }\PY{l+s}{\PYZdq{}}\PY{+w}{ }\PY{o}{+}\PY{+w}{ }\PY{n}{elapsedMicros}\PY{+w}{ }\PY{o}{+}\PY{+w}{ }\PY{l+s}{\PYZdq{}}\PY{l+s}{ microseconds}\PY{l+s}{\PYZdq{}}\PY{p}{)}\PY{p}{;}
\PY{+w}{        }\PY{p}{\PYZcb{}}
\PY{+w}{    }\PY{p}{\PYZcb{}}
\PY{p}{\PYZcb{}}

\PY{k+kd}{public}\PY{+w}{ }\PY{k+kd}{class} \PY{n+nc}{ProxyDemo}\PY{+w}{ }\PY{p}{\PYZob{}}
\PY{+w}{    }\PY{k+kd}{public}\PY{+w}{ }\PY{k+kd}{static}\PY{+w}{ }\PY{k+kt}{void}\PY{+w}{ }\PY{n+nf}{main}\PY{p}{(}\PY{n}{String}\PY{o}{[}\PY{o}{]}\PY{+w}{ }\PY{n}{args}\PY{p}{)}\PY{+w}{ }\PY{p}{\PYZob{}}
\PY{+w}{        }\PY{n}{OrderService}\PY{+w}{ }\PY{n}{real}\PY{+w}{ }\PY{o}{=}\PY{+w}{ }\PY{k}{new}\PY{+w}{ }\PY{n}{OrderServiceImpl}\PY{p}{(}\PY{p}{)}\PY{p}{;}

\PY{+w}{        }\PY{n}{OrderService}\PY{+w}{ }\PY{n}{proxy}\PY{+w}{ }\PY{o}{=}\PY{+w}{ }\PY{p}{(}\PY{n}{OrderService}\PY{p}{)}\PY{+w}{ }\PY{n}{Proxy}\PY{p}{.}\PY{n+na}{newProxyInstance}\PY{p}{(}
\PY{+w}{                }\PY{n}{OrderService}\PY{p}{.}\PY{n+na}{class}\PY{p}{.}\PY{n+na}{getClassLoader}\PY{p}{(}\PY{p}{)}\PY{p}{,}
\PY{+w}{                }\PY{k}{new}\PY{+w}{ }\PY{n}{Class}\PY{o}{\PYZlt{}}\PY{o}{?}\PY{o}{\PYZgt{}}\PY{o}{[}\PY{o}{]}\PY{+w}{ }\PY{p}{\PYZob{}}\PY{+w}{ }\PY{n}{OrderService}\PY{p}{.}\PY{n+na}{class}\PY{+w}{ }\PY{p}{\PYZcb{}}\PY{p}{,}
\PY{+w}{                }\PY{k}{new}\PY{+w}{ }\PY{n}{LoggingTimingHandler}\PY{p}{(}\PY{n}{real}\PY{p}{)}\PY{p}{)}\PY{p}{;}

\PY{+w}{        }\PY{n}{proxy}\PY{p}{.}\PY{n+na}{placeOrder}\PY{p}{(}\PY{l+s}{\PYZdq{}}\PY{l+s}{Laptop}\PY{l+s}{\PYZdq{}}\PY{p}{)}\PY{p}{;}
\PY{+w}{        }\PY{n}{System}\PY{p}{.}\PY{n+na}{out}\PY{p}{.}\PY{n+na}{println}\PY{p}{(}\PY{l+s}{\PYZdq{}}\PY{l+s}{Count: }\PY{l+s}{\PYZdq{}}\PY{+w}{ }\PY{o}{+}\PY{+w}{ }\PY{n}{proxy}\PY{p}{.}\PY{n+na}{getOrderCount}\PY{p}{(}\PY{p}{)}\PY{p}{)}\PY{p}{;}
\PY{+w}{    }\PY{p}{\PYZcb{}}
\PY{p}{\PYZcb{}}
\end{Verbatim}
\end{codebox}

Output looks like:

\begin{codebox}
\begin{Verbatim}
Calling placeOrder
Order placed: Laptop
placeOrder took 12 microseconds
Calling getOrderCount
getOrderCount took 3 microseconds
Count: 1
\end{Verbatim}
\end{codebox}

\subsection[Interface-Only Limitation]{Interface-Only Limitation}
\label{h:16:interface-only-limitation}
\code{java.\allowbreak{}lang.\allowbreak{}reflect.\allowbreak{}Proxy} can only proxy \textbf{interfaces}, not concrete classes. This is a fundamental JDK limitation — the generated proxy class implements the given interfaces, it cannot extend an arbitrary class. If you need to proxy a class with no interface (common with plain Spring \code{@\allowbreak{}Service} classes that don’t implement one), you need a bytecode-generation library instead.

\subsection[How Frameworks Use This]{How Frameworks Use This}
\label{h:16:how-frameworks-use-this}
\begin{itemize}
\item 
\textbf{Spring AOP}: when the target bean implements at least one interface, Spring uses \textbf{JDK dynamic proxies} (exactly the mechanism above) to implement things like \code{@\allowbreak{}Transactional}, \code{@\allowbreak{}Cacheable}, and custom aspects. When the bean has no interface (or \code{proxyTargetClass=\allowbreak{}true} is set), Spring falls back to \textbf{CGLIB} (bundled, repackaged inside \code{spring-\allowbreak{}core}), which generates a real subclass at runtime that overrides methods — this works on concrete classes but can’t proxy \code{final} classes or \code{final} methods.

\item 
\textbf{ByteBuddy} is a modern, actively maintained alternative to CGLIB for generating subclasses/bytecode at runtime; it’s used internally by Mockito (for mocking classes, not just interfaces) and newer versions of Hibernate.

\end{itemize}

{\tablefont
\begin{xltabular}{\linewidth}{@{}L{0.434}L{0.681}L{1.070}L{1.815}@{}}
\toprule
\rowcolor{tablehdr}
\thd{Mechanism} & \thd{Can proxy interfaces?} & \thd{Can proxy concrete classes?} & \thd{Notes} \\
\midrule
\endhead
JDK \code{Proxy} & Yes & No & Built into the JDK, no dependency needed \\
CGLIB & Yes (via subclass) & Yes, except \code{final} classes/methods & Bundled inside Spring; effectively unmaintained upstream \\
ByteBuddy & Yes & Yes, except \code{final} classes/methods & Modern, actively maintained, used by Mockito \\
\bottomrule
\end{xltabular}
}

\section[Method Handles]{Method Handles}
\label{h:16:method-handles}
A \textbf{method handle} is a typed, directly executable reference to an underlying method, constructor, or field access, created through the \code{java.\allowbreak{}lang.\allowbreak{}invoke} package. Method handles were introduced in Java 7 specifically to be a faster, more type-safe alternative to reflection for performance-sensitive dynamic dispatch — and the JVM itself uses them internally to implement lambdas and \code{invokedynamic}.

\subsection[Lookup, MethodType, and Finding a Handle]{\code{Lookup}, \code{MethodType}, and Finding a Handle}
\label{h:16:lookup-methodtype-and-finding-a-handle}
\begin{codebox}
\begin{Verbatim}[commandchars=\\\{\}]
\PY{k+kn}{import}\PY{+w}{ }\PY{n+nn}{java.lang.invoke.*}\PY{p}{;}

\PY{k+kd}{public}\PY{+w}{ }\PY{k+kd}{class} \PY{n+nc}{MethodHandleDemo}\PY{+w}{ }\PY{p}{\PYZob{}}
\PY{+w}{    }\PY{k+kd}{public}\PY{+w}{ }\PY{k+kd}{static}\PY{+w}{ }\PY{k+kt}{void}\PY{+w}{ }\PY{n+nf}{main}\PY{p}{(}\PY{n}{String}\PY{o}{[}\PY{o}{]}\PY{+w}{ }\PY{n}{args}\PY{p}{)}\PY{+w}{ }\PY{k+kd}{throws}\PY{+w}{ }\PY{n}{Throwable}\PY{+w}{ }\PY{p}{\PYZob{}}
\PY{+w}{        }\PY{n}{MethodHandles}\PY{p}{.}\PY{n+na}{Lookup}\PY{+w}{ }\PY{n}{lookup}\PY{+w}{ }\PY{o}{=}\PY{+w}{ }\PY{n}{MethodHandles}\PY{p}{.}\PY{n+na}{lookup}\PY{p}{(}\PY{p}{)}\PY{p}{;}

\PY{+w}{        }\PY{c+c1}{// MethodType describes the signature: return type first, then parameter types}
\PY{+w}{        }\PY{n}{MethodType}\PY{+w}{ }\PY{n}{stringLength}\PY{+w}{ }\PY{o}{=}\PY{+w}{ }\PY{n}{MethodType}\PY{p}{.}\PY{n+na}{methodType}\PY{p}{(}\PY{k+kt}{int}\PY{p}{.}\PY{n+na}{class}\PY{p}{)}\PY{p}{;}
\PY{+w}{        }\PY{n}{MethodHandle}\PY{+w}{ }\PY{n}{lengthHandle}\PY{+w}{ }\PY{o}{=}\PY{+w}{ }\PY{n}{lookup}\PY{p}{.}\PY{n+na}{findVirtual}\PY{p}{(}\PY{n}{String}\PY{p}{.}\PY{n+na}{class}\PY{p}{,}\PY{+w}{ }\PY{l+s}{\PYZdq{}}\PY{l+s}{length}\PY{l+s}{\PYZdq{}}\PY{p}{,}\PY{+w}{ }\PY{n}{stringLength}\PY{p}{)}\PY{p}{;}

\PY{+w}{        }\PY{k+kt}{int}\PY{+w}{ }\PY{n}{len}\PY{+w}{ }\PY{o}{=}\PY{+w}{ }\PY{p}{(}\PY{k+kt}{int}\PY{p}{)}\PY{+w}{ }\PY{n}{lengthHandle}\PY{p}{.}\PY{n+na}{invokeExact}\PY{p}{(}\PY{l+s}{\PYZdq{}}\PY{l+s}{hello}\PY{l+s}{\PYZdq{}}\PY{p}{)}\PY{p}{;}\PY{+w}{ }\PY{c+c1}{// 5}

\PY{+w}{        }\PY{c+c1}{// findStatic for static methods}
\PY{+w}{        }\PY{n}{MethodType}\PY{+w}{ }\PY{n}{parseIntType}\PY{+w}{ }\PY{o}{=}\PY{+w}{ }\PY{n}{MethodType}\PY{p}{.}\PY{n+na}{methodType}\PY{p}{(}\PY{k+kt}{int}\PY{p}{.}\PY{n+na}{class}\PY{p}{,}\PY{+w}{ }\PY{n}{String}\PY{p}{.}\PY{n+na}{class}\PY{p}{)}\PY{p}{;}
\PY{+w}{        }\PY{n}{MethodHandle}\PY{+w}{ }\PY{n}{parseInt}\PY{+w}{ }\PY{o}{=}\PY{+w}{ }\PY{n}{lookup}\PY{p}{.}\PY{n+na}{findStatic}\PY{p}{(}\PY{n}{Integer}\PY{p}{.}\PY{n+na}{class}\PY{p}{,}\PY{+w}{ }\PY{l+s}{\PYZdq{}}\PY{l+s}{parseInt}\PY{l+s}{\PYZdq{}}\PY{p}{,}\PY{+w}{ }\PY{n}{parseIntType}\PY{p}{)}\PY{p}{;}
\PY{+w}{        }\PY{k+kt}{int}\PY{+w}{ }\PY{n}{n}\PY{+w}{ }\PY{o}{=}\PY{+w}{ }\PY{p}{(}\PY{k+kt}{int}\PY{p}{)}\PY{+w}{ }\PY{n}{parseInt}\PY{p}{.}\PY{n+na}{invokeExact}\PY{p}{(}\PY{l+s}{\PYZdq{}}\PY{l+s}{42}\PY{l+s}{\PYZdq{}}\PY{p}{)}\PY{p}{;}\PY{+w}{ }\PY{c+c1}{// 42}

\PY{+w}{        }\PY{n}{System}\PY{p}{.}\PY{n+na}{out}\PY{p}{.}\PY{n+na}{println}\PY{p}{(}\PY{n}{len}\PY{+w}{ }\PY{o}{+}\PY{+w}{ }\PY{n}{n}\PY{p}{)}\PY{p}{;}\PY{+w}{ }\PY{c+c1}{// 47}
\PY{+w}{    }\PY{p}{\PYZcb{}}
\PY{p}{\PYZcb{}}
\end{Verbatim}
\end{codebox}

\subsection[invokeExact vs invoke]{\code{invokeExact} vs \code{invoke}}
\label{h:16:invokeexact-vs-invoke}
\begin{itemize}
\item 
\textbf{\code{invokeExact}} requires the argument and return types at the call site to match the handle’s \code{MethodType} \emph{exactly} — no widening, no boxing conversions. If they don’t match, it throws \code{Wrong\allowbreak{}Method\allowbreak{}Type\allowbreak{}Exception}. This is the fast path.

\item 
\textbf{\code{invoke}} is more forgiving: it performs the same kind of adaptation reflection would (widening primitives, boxing/unboxing, casting), at a small extra cost.

\end{itemize}

\begin{codebox}
\begin{Verbatim}[commandchars=\\\{\}]
\PY{n}{MethodHandle}\PY{+w}{ }\PY{n}{mh}\PY{+w}{ }\PY{o}{=}\PY{+w}{ }\PY{n}{MethodHandles}\PY{p}{.}\PY{n+na}{lookup}\PY{p}{(}\PY{p}{)}
\PY{+w}{        }\PY{p}{.}\PY{n+na}{findStatic}\PY{p}{(}\PY{n}{Math}\PY{p}{.}\PY{n+na}{class}\PY{p}{,}\PY{+w}{ }\PY{l+s}{\PYZdq{}}\PY{l+s}{max}\PY{l+s}{\PYZdq{}}\PY{p}{,}\PY{+w}{ }\PY{n}{MethodType}\PY{p}{.}\PY{n+na}{methodType}\PY{p}{(}\PY{k+kt}{int}\PY{p}{.}\PY{n+na}{class}\PY{p}{,}\PY{+w}{ }\PY{k+kt}{int}\PY{p}{.}\PY{n+na}{class}\PY{p}{,}\PY{+w}{ }\PY{k+kt}{int}\PY{p}{.}\PY{n+na}{class}\PY{p}{)}\PY{p}{)}\PY{p}{;}

\PY{k+kt}{int}\PY{+w}{ }\PY{n}{a}\PY{+w}{ }\PY{o}{=}\PY{+w}{ }\PY{p}{(}\PY{k+kt}{int}\PY{p}{)}\PY{+w}{ }\PY{n}{mh}\PY{p}{.}\PY{n+na}{invokeExact}\PY{p}{(}\PY{l+m+mi}{3}\PY{p}{,}\PY{+w}{ }\PY{l+m+mi}{5}\PY{p}{)}\PY{p}{;}\PY{+w}{        }\PY{c+c1}{// works: types match exactly}
\PY{c+c1}{// Object r = mh.invoke((Object) 3, (Object) 5); // invoke() would adapt this automatically}
\end{Verbatim}
\end{codebox}

\subsection[Why Method Handles Are Faster Than Reflection]{Why Method Handles Are Faster Than Reflection}
\label{h:16:why-method-handles-are-faster-than-reflection}
\begin{itemize}
\item 
\code{Method.\allowbreak{}invoke} always boxes arguments into an \code{Object[]} and does a security/access check on every call (unless cached).

\item 
A \code{MethodHandle}, once resolved, behaves much more like a direct call to the JIT compiler — it can be inlined, and the JVM has special bytecode support (\code{invokedynamic}) and call-site optimization for it. Repeated calls through the same \code{MethodHandle} get progressively faster as the JIT specializes them, approaching the speed of a hand-written call.

\end{itemize}

\subsection[LambdaMetafactory: How Lambdas Are Actually Implemented]{\code{LambdaMetafactory}: How Lambdas Are Actually Implemented}
\label{h:16:lambdametafactory-how-lambdas-are-actually-implemented}
When you write a lambda expression, the compiler does \textbf{not} generate an anonymous inner class (unlike pre-Java-8 anonymous classes). Instead, it emits an \code{invokedynamic} instruction whose bootstrap method is \code{Lambda\allowbreak{}Metafactory.\allowbreak{}metafactory}. At runtime, the first time that lambda expression is executed, the JVM calls the metafactory, which uses method handles to generate a small hidden class implementing the target functional interface, and then caches it. Every subsequent execution of that same lambda reuses the cached class.

\begin{codebox}
\begin{Verbatim}[commandchars=\\\{\}]
\PY{n}{Runnable}\PY{+w}{ }\PY{n}{r}\PY{+w}{ }\PY{o}{=}\PY{+w}{ }\PY{p}{(}\PY{p}{)}\PY{+w}{ }\PY{o}{\PYZhy{}}\PY{o}{\PYZgt{}}\PY{+w}{ }\PY{n}{System}\PY{p}{.}\PY{n+na}{out}\PY{p}{.}\PY{n+na}{println}\PY{p}{(}\PY{l+s}{\PYZdq{}}\PY{l+s}{hi}\PY{l+s}{\PYZdq{}}\PY{p}{)}\PY{p}{;}
\end{Verbatim}
\end{codebox}

compiles roughly to bytecode equivalent to:

\begin{codebox}
\begin{Verbatim}[commandchars=\\\{\}]
\PY{c+c1}{// conceptual illustration — not what you write, what the JVM does internally}
\PY{n}{CallSite}\PY{+w}{ }\PY{n}{site}\PY{+w}{ }\PY{o}{=}\PY{+w}{ }\PY{n}{LambdaMetafactory}\PY{p}{.}\PY{n+na}{metafactory}\PY{p}{(}
\PY{+w}{        }\PY{n}{lookup}\PY{p}{,}
\PY{+w}{        }\PY{l+s}{\PYZdq{}}\PY{l+s}{run}\PY{l+s}{\PYZdq{}}\PY{p}{,}\PY{+w}{                                  }\PY{c+c1}{// the functional interface\PYZsq{}s method name}
\PY{+w}{        }\PY{n}{MethodType}\PY{p}{.}\PY{n+na}{methodType}\PY{p}{(}\PY{n}{Runnable}\PY{p}{.}\PY{n+na}{class}\PY{p}{)}\PY{p}{,}\PY{+w}{  }\PY{c+c1}{// factory signature}
\PY{+w}{        }\PY{n}{MethodType}\PY{p}{.}\PY{n+na}{methodType}\PY{p}{(}\PY{k+kt}{void}\PY{p}{.}\PY{n+na}{class}\PY{p}{)}\PY{p}{,}\PY{+w}{      }\PY{c+c1}{// signature of Runnable.run()}
\PY{+w}{        }\PY{n}{implementationMethodHandle}\PY{p}{,}\PY{+w}{              }\PY{c+c1}{// handle pointing at the lambda body}
\PY{+w}{        }\PY{n}{MethodType}\PY{p}{.}\PY{n+na}{methodType}\PY{p}{(}\PY{k+kt}{void}\PY{p}{.}\PY{n+na}{class}\PY{p}{)}\PY{p}{)}\PY{p}{;}
\PY{n}{Runnable}\PY{+w}{ }\PY{n}{r}\PY{+w}{ }\PY{o}{=}\PY{+w}{ }\PY{p}{(}\PY{n}{Runnable}\PY{p}{)}\PY{+w}{ }\PY{n}{site}\PY{p}{.}\PY{n+na}{getTarget}\PY{p}{(}\PY{p}{)}\PY{p}{.}\PY{n+na}{invokeExact}\PY{p}{(}\PY{p}{)}\PY{p}{;}
\end{Verbatim}
\end{codebox}

This is why lambdas are much cheaper than people expect: no class file per lambda at compile time (just one small generated class at runtime, cached), and calling a lambda is close to a direct method call once the JIT has warmed up — much faster than a reflective dispatch.

\section[VarHandles]{VarHandles}
\label{h:16:varhandles}
A \textbf{\code{VarHandle}} (Java 9+) is a typed reference to a variable — a field, an array element, or an off-heap memory location — that supports fine-grained control over how it’s read and written, including atomic and memory-ordering operations. It was introduced to replace two older mechanisms: \code{sun.\allowbreak{}misc.\allowbreak{}Unsafe}’s direct memory-access methods, and the \code{java.\allowbreak{}util.\allowbreak{}concurrent.\allowbreak{}atomic.\allowbreak{}Atomic\allowbreak{}XField\allowbreak{}Updater} classes (\code{Atomic\allowbreak{}Integer\allowbreak{}Field\allowbreak{}Updater}, \code{AtomicLongFieldUpdater}, \code{Atomic\allowbreak{}Reference\allowbreak{}Field\allowbreak{}Updater}), both of which had awkward or unsafe APIs.

\subsection[Getting a VarHandle]{Getting a VarHandle}
\label{h:16:getting-a-varhandle}
\begin{codebox}
\begin{Verbatim}[commandchars=\\\{\}]
\PY{k+kn}{import}\PY{+w}{ }\PY{n+nn}{java.lang.invoke.*}\PY{p}{;}

\PY{k+kd}{public}\PY{+w}{ }\PY{k+kd}{class} \PY{n+nc}{Counter}\PY{+w}{ }\PY{p}{\PYZob{}}
\PY{+w}{    }\PY{k+kd}{private}\PY{+w}{ }\PY{k+kd}{volatile}\PY{+w}{ }\PY{k+kt}{int}\PY{+w}{ }\PY{n}{count}\PY{p}{;}

\PY{+w}{    }\PY{k+kd}{private}\PY{+w}{ }\PY{k+kd}{static}\PY{+w}{ }\PY{k+kd}{final}\PY{+w}{ }\PY{n}{VarHandle}\PY{+w}{ }\PY{n}{COUNT\PYZus{}HANDLE}\PY{p}{;}
\PY{+w}{    }\PY{k+kd}{static}\PY{+w}{ }\PY{p}{\PYZob{}}
\PY{+w}{        }\PY{k}{try}\PY{+w}{ }\PY{p}{\PYZob{}}
\PY{+w}{            }\PY{n}{COUNT\PYZus{}HANDLE}\PY{+w}{ }\PY{o}{=}\PY{+w}{ }\PY{n}{MethodHandles}\PY{p}{.}\PY{n+na}{lookup}\PY{p}{(}\PY{p}{)}
\PY{+w}{                    }\PY{p}{.}\PY{n+na}{findVarHandle}\PY{p}{(}\PY{n}{Counter}\PY{p}{.}\PY{n+na}{class}\PY{p}{,}\PY{+w}{ }\PY{l+s}{\PYZdq{}}\PY{l+s}{count}\PY{l+s}{\PYZdq{}}\PY{p}{,}\PY{+w}{ }\PY{k+kt}{int}\PY{p}{.}\PY{n+na}{class}\PY{p}{)}\PY{p}{;}
\PY{+w}{        }\PY{p}{\PYZcb{}}\PY{+w}{ }\PY{k}{catch}\PY{+w}{ }\PY{p}{(}\PY{n}{ReflectiveOperationException}\PY{+w}{ }\PY{n}{e}\PY{p}{)}\PY{+w}{ }\PY{p}{\PYZob{}}
\PY{+w}{            }\PY{k}{throw}\PY{+w}{ }\PY{k}{new}\PY{+w}{ }\PY{n}{ExceptionInInitializerError}\PY{p}{(}\PY{n}{e}\PY{p}{)}\PY{p}{;}
\PY{+w}{        }\PY{p}{\PYZcb{}}
\PY{+w}{    }\PY{p}{\PYZcb{}}

\PY{+w}{    }\PY{k+kd}{public}\PY{+w}{ }\PY{k+kt}{boolean}\PY{+w}{ }\PY{n+nf}{incrementIfBelow}\PY{p}{(}\PY{k+kt}{int}\PY{+w}{ }\PY{n}{max}\PY{p}{)}\PY{+w}{ }\PY{p}{\PYZob{}}
\PY{+w}{        }\PY{k+kt}{int}\PY{+w}{ }\PY{n}{current}\PY{+w}{ }\PY{o}{=}\PY{+w}{ }\PY{p}{(}\PY{k+kt}{int}\PY{p}{)}\PY{+w}{ }\PY{n}{COUNT\PYZus{}HANDLE}\PY{p}{.}\PY{n+na}{getVolatile}\PY{p}{(}\PY{k}{this}\PY{p}{)}\PY{p}{;}
\PY{+w}{        }\PY{k}{while}\PY{+w}{ }\PY{p}{(}\PY{n}{current}\PY{+w}{ }\PY{o}{\PYZlt{}}\PY{+w}{ }\PY{n}{max}\PY{p}{)}\PY{+w}{ }\PY{p}{\PYZob{}}
\PY{+w}{            }\PY{c+c1}{// atomically set count = current + 1, only if it\PYZsq{}s still == current}
\PY{+w}{            }\PY{k}{if}\PY{+w}{ }\PY{p}{(}\PY{n}{COUNT\PYZus{}HANDLE}\PY{p}{.}\PY{n+na}{compareAndSet}\PY{p}{(}\PY{k}{this}\PY{p}{,}\PY{+w}{ }\PY{n}{current}\PY{p}{,}\PY{+w}{ }\PY{n}{current}\PY{+w}{ }\PY{o}{+}\PY{+w}{ }\PY{l+m+mi}{1}\PY{p}{)}\PY{p}{)}\PY{+w}{ }\PY{p}{\PYZob{}}
\PY{+w}{                }\PY{k}{return}\PY{+w}{ }\PY{k+kc}{true}\PY{p}{;}
\PY{+w}{            }\PY{p}{\PYZcb{}}
\PY{+w}{            }\PY{n}{current}\PY{+w}{ }\PY{o}{=}\PY{+w}{ }\PY{p}{(}\PY{k+kt}{int}\PY{p}{)}\PY{+w}{ }\PY{n}{COUNT\PYZus{}HANDLE}\PY{p}{.}\PY{n+na}{getVolatile}\PY{p}{(}\PY{k}{this}\PY{p}{)}\PY{p}{;}
\PY{+w}{        }\PY{p}{\PYZcb{}}
\PY{+w}{        }\PY{k}{return}\PY{+w}{ }\PY{k+kc}{false}\PY{p}{;}
\PY{+w}{    }\PY{p}{\PYZcb{}}
\PY{p}{\PYZcb{}}
\end{Verbatim}
\end{codebox}

\subsection[Access Modes]{Access Modes}
\label{h:16:access-modes}
\code{VarHandle} exposes several families of access, each with different performance/visibility trade-offs — from weakest/fastest to strongest/slowest guarantee:

{\tablefontsmall
\begin{xltabular}{\linewidth}{@{}L{0.435}L{2.025}L{0.540}@{}}
\toprule
\rowcolor{tablehdr}
\thd{Access mode} & \thd{Guarantee} & \thd{Typical method} \\
\midrule
\endhead
\textbf{Plain} & No atomicity, no ordering guarantee beyond normal Java semantics — like a plain field read/write & \code{get}, \code{set} \\
\textbf{Opaque} & Atomic for that single variable, but no ordering relative to other variables (no happens-before) & \code{getOpaque}, \code{setOpaque} \\
\textbf{Acquire/Release} & One-directional memory fence — \code{acquire} prevents later reads/writes from moving before it; \code{release} prevents earlier ones from moving after it & \code{getAcquire}, \code{setRelease} \\
\textbf{Volatile} & Full happens-before ordering in both directions, same as a \code{volatile} field & \code{getVolatile}, \code{setVolatile} \\
\bottomrule
\end{xltabular}
}

\begin{codebox}
\begin{Verbatim}[commandchars=\\\{\}]
\PY{c+c1}{// same variable, four different consistency guarantees}
\PY{k+kt}{int}\PY{+w}{ }\PY{n}{plain}\PY{+w}{    }\PY{o}{=}\PY{+w}{ }\PY{p}{(}\PY{k+kt}{int}\PY{p}{)}\PY{+w}{ }\PY{n}{COUNT\PYZus{}HANDLE}\PY{p}{.}\PY{n+na}{get}\PY{p}{(}\PY{n}{counterInstance}\PY{p}{)}\PY{p}{;}
\PY{k+kt}{int}\PY{+w}{ }\PY{n}{opaque}\PY{+w}{   }\PY{o}{=}\PY{+w}{ }\PY{p}{(}\PY{k+kt}{int}\PY{p}{)}\PY{+w}{ }\PY{n}{COUNT\PYZus{}HANDLE}\PY{p}{.}\PY{n+na}{getOpaque}\PY{p}{(}\PY{n}{counterInstance}\PY{p}{)}\PY{p}{;}
\PY{k+kt}{int}\PY{+w}{ }\PY{n}{acquire}\PY{+w}{  }\PY{o}{=}\PY{+w}{ }\PY{p}{(}\PY{k+kt}{int}\PY{p}{)}\PY{+w}{ }\PY{n}{COUNT\PYZus{}HANDLE}\PY{p}{.}\PY{n+na}{getAcquire}\PY{p}{(}\PY{n}{counterInstance}\PY{p}{)}\PY{p}{;}
\PY{k+kt}{int}\PY{+w}{ }\PY{n}{volatileRead}\PY{+w}{ }\PY{o}{=}\PY{+w}{ }\PY{p}{(}\PY{k+kt}{int}\PY{p}{)}\PY{+w}{ }\PY{n}{COUNT\PYZus{}HANDLE}\PY{p}{.}\PY{n+na}{getVolatile}\PY{p}{(}\PY{n}{counterInstance}\PY{p}{)}\PY{p}{;}
\end{Verbatim}
\end{codebox}

\subsection[compareAndSet and Array Access]{\code{compareAndSet} and Array Access}
\label{h:16:compareandset-and-array-access}
\begin{codebox}
\begin{Verbatim}[commandchars=\\\{\}]
\PY{k+kn}{import}\PY{+w}{ }\PY{n+nn}{java.lang.invoke.*}\PY{p}{;}

\PY{k+kd}{public}\PY{+w}{ }\PY{k+kd}{class} \PY{n+nc}{ArrayHandleDemo}\PY{+w}{ }\PY{p}{\PYZob{}}
\PY{+w}{    }\PY{k+kd}{public}\PY{+w}{ }\PY{k+kd}{static}\PY{+w}{ }\PY{k+kt}{void}\PY{+w}{ }\PY{n+nf}{main}\PY{p}{(}\PY{n}{String}\PY{o}{[}\PY{o}{]}\PY{+w}{ }\PY{n}{args}\PY{p}{)}\PY{+w}{ }\PY{p}{\PYZob{}}
\PY{+w}{        }\PY{k+kt}{int}\PY{o}{[}\PY{o}{]}\PY{+w}{ }\PY{n}{array}\PY{+w}{ }\PY{o}{=}\PY{+w}{ }\PY{k}{new}\PY{+w}{ }\PY{k+kt}{int}\PY{o}{[}\PY{l+m+mi}{10}\PY{o}{]}\PY{p}{;}
\PY{+w}{        }\PY{n}{VarHandle}\PY{+w}{ }\PY{n}{arrayHandle}\PY{+w}{ }\PY{o}{=}\PY{+w}{ }\PY{n}{MethodHandles}\PY{p}{.}\PY{n+na}{arrayElementVarHandle}\PY{p}{(}\PY{k+kt}{int}\PY{o}{[}\PY{o}{]}\PY{p}{.}\PY{n+na}{class}\PY{p}{)}\PY{p}{;}

\PY{+w}{        }\PY{n}{arrayHandle}\PY{p}{.}\PY{n+na}{setVolatile}\PY{p}{(}\PY{n}{array}\PY{p}{,}\PY{+w}{ }\PY{l+m+mi}{3}\PY{p}{,}\PY{+w}{ }\PY{l+m+mi}{42}\PY{p}{)}\PY{p}{;}
\PY{+w}{        }\PY{k+kt}{boolean}\PY{+w}{ }\PY{n}{updated}\PY{+w}{ }\PY{o}{=}\PY{+w}{ }\PY{n}{arrayHandle}\PY{p}{.}\PY{n+na}{compareAndSet}\PY{p}{(}\PY{n}{array}\PY{p}{,}\PY{+w}{ }\PY{l+m+mi}{3}\PY{p}{,}\PY{+w}{ }\PY{l+m+mi}{42}\PY{p}{,}\PY{+w}{ }\PY{l+m+mi}{100}\PY{p}{)}\PY{p}{;}
\PY{+w}{        }\PY{n}{System}\PY{p}{.}\PY{n+na}{out}\PY{p}{.}\PY{n+na}{println}\PY{p}{(}\PY{n}{updated}\PY{+w}{ }\PY{o}{+}\PY{+w}{ }\PY{l+s}{\PYZdq{}}\PY{l+s}{ }\PY{l+s}{\PYZdq{}}\PY{+w}{ }\PY{o}{+}\PY{+w}{ }\PY{n}{array}\PY{o}{[}\PY{l+m+mi}{3}\PY{o}{]}\PY{p}{)}\PY{p}{;}\PY{+w}{ }\PY{c+c1}{// true 100}
\PY{+w}{    }\PY{p}{\PYZcb{}}
\PY{p}{\PYZcb{}}
\end{Verbatim}
\end{codebox}

\subsection[Memory Fences]{Memory Fences}
\label{h:16:memory-fences}
For advanced lock-free algorithms, \code{VarHandle} also exposes standalone fence methods that don’t touch a specific variable but establish ordering for the surrounding code: \code{Var\allowbreak{}Handle.\allowbreak{}acquire\allowbreak{}Fence()}, \code{releaseFence()}, \code{fullFence()}, and \code{loadLoadFence()}/\code{storeStoreFence()}. These map closely to CPU-level memory barrier instructions and are rarely needed outside of writing custom concurrency primitives (most application code should just use \code{java.\allowbreak{}util.\allowbreak{}concurrent} classes instead).

\section[Unsafe Alternatives]{Unsafe Alternatives}
\label{h:16:unsafe-alternatives}
\code{sun.\allowbreak{}misc.\allowbreak{}Unsafe} was an internal, undocumented JDK class that gave direct access to raw memory operations: allocating memory outside the heap, reading/writing at arbitrary memory offsets, compare-and-swap on fields by offset, throwing exceptions without declaring them, and even instantiating objects while skipping constructors. It was never part of the official Java API — it existed only because early \code{java.\allowbreak{}util.\allowbreak{}concurrent} and other JDK-internal code needed low-level primitives that didn’t exist anywhere else yet, and it “leaked” out because it was technically public.

\subsection[Why It’s Being Removed]{Why It’s Being Removed}
\label{h:16:why-its-being-removed}
\begin{itemize}
\item 
It bypasses the type system and safety checks the JVM is supposed to guarantee, so misuse can crash the JVM instead of throwing a normal exception.

\item 
It made \textbf{strong encapsulation} (from JPMS) meaningless for anyone using it, since it can reach into any object’s memory layout directly.

\item 
\textbf{JEP 471} (targeted for a recent JDK release) deprecates and eventually removes \code{Unsafe}’s memory-access methods for removal, pushing everyone toward the supported replacements below.

\item 
\textbf{JEP 498} relates to further restricting/warning about deep reflective and unsafe memory access as part of the broader multi-release effort to close off “back door” access to internals (following on from JEP 403’s strong encapsulation of JDK internals started in JDK 17).

\end{itemize}

The practical effect for a codebase: any dependency still calling \code{Unsafe.\allowbreak{}putInt}, \code{Unsafe.\allowbreak{}allocate\allowbreak{}Instance}, etc., will eventually stop working on newer JDKs, and code reviewers should flag any \emph{new} direct use of \code{sun.\allowbreak{}misc.\allowbreak{}Unsafe} in application code as a serious red flag.

\subsection[Modern Replacements]{Modern Replacements}
\label{h:16:modern-replacements}
{\tablefontsmall
\begin{xltabular}{\linewidth}{@{}L{1.161}L{0.839}@{}}
\toprule
\rowcolor{tablehdr}
\thd{Old \code{Unsafe} use case} & \thd{Modern replacement} \\
\midrule
\endhead
Atomic field access, CAS on a field by offset & \code{VarHandle} (via \code{Method\allowbreak{}Handles.\allowbreak{}lookup().\allowbreak{}find\allowbreak{}Var\allowbreak{}Handle(...)}) \\
Fast, checked-exception-free-ish reflective invocation & \code{MethodHandle} \\
Off-heap memory allocation, pointers, native calls & \textbf{Foreign Function \& Memory API} (\code{java.\allowbreak{}lang.\allowbreak{}foreign}, finalized in Java 22, usable as a preview from Java 19+) \\
Running cleanup logic when an object becomes unreachable (replacing \code{finalize()} and ad hoc \code{Unsafe}-based tricks) & \code{java.\allowbreak{}lang.\allowbreak{}ref.\allowbreak{}Cleaner} \\
Creating a class at runtime without exposing it as a discoverable, loadable, or reflectively accessible named class (used internally for lambdas) & \textbf{Hidden classes} (\code{Method\allowbreak{}Handles.\allowbreak{}Lookup.\allowbreak{}define\allowbreak{}Hidden\allowbreak{}Class}, Java 15+) \\
\bottomrule
\end{xltabular}
}

\begin{codebox}
\begin{Verbatim}[commandchars=\\\{\}]
\PY{k+kn}{import}\PY{+w}{ }\PY{n+nn}{java.lang.foreign.*}\PY{p}{;}

\PY{k+kd}{public}\PY{+w}{ }\PY{k+kd}{class} \PY{n+nc}{ForeignMemoryDemo}\PY{+w}{ }\PY{p}{\PYZob{}}
\PY{+w}{    }\PY{k+kd}{public}\PY{+w}{ }\PY{k+kd}{static}\PY{+w}{ }\PY{k+kt}{void}\PY{+w}{ }\PY{n+nf}{main}\PY{p}{(}\PY{n}{String}\PY{o}{[}\PY{o}{]}\PY{+w}{ }\PY{n}{args}\PY{p}{)}\PY{+w}{ }\PY{p}{\PYZob{}}
\PY{+w}{        }\PY{k}{try}\PY{+w}{ }\PY{p}{(}\PY{n}{Arena}\PY{+w}{ }\PY{n}{arena}\PY{+w}{ }\PY{o}{=}\PY{+w}{ }\PY{n}{Arena}\PY{p}{.}\PY{n+na}{ofConfined}\PY{p}{(}\PY{p}{)}\PY{p}{)}\PY{+w}{ }\PY{p}{\PYZob{}}
\PY{+w}{            }\PY{n}{MemorySegment}\PY{+w}{ }\PY{n}{segment}\PY{+w}{ }\PY{o}{=}\PY{+w}{ }\PY{n}{arena}\PY{p}{.}\PY{n+na}{allocate}\PY{p}{(}\PY{l+m+mi}{4}\PY{p}{)}\PY{p}{;}\PY{+w}{ }\PY{c+c1}{// 4 bytes, off\PYZhy{}heap}
\PY{+w}{            }\PY{n}{segment}\PY{p}{.}\PY{n+na}{set}\PY{p}{(}\PY{n}{ValueLayout}\PY{p}{.}\PY{n+na}{JAVA\PYZus{}INT}\PY{p}{,}\PY{+w}{ }\PY{l+m+mi}{0}\PY{p}{,}\PY{+w}{ }\PY{l+m+mi}{42}\PY{p}{)}\PY{p}{;}
\PY{+w}{            }\PY{n}{System}\PY{p}{.}\PY{n+na}{out}\PY{p}{.}\PY{n+na}{println}\PY{p}{(}\PY{n}{segment}\PY{p}{.}\PY{n+na}{get}\PY{p}{(}\PY{n}{ValueLayout}\PY{p}{.}\PY{n+na}{JAVA\PYZus{}INT}\PY{p}{,}\PY{+w}{ }\PY{l+m+mi}{0}\PY{p}{)}\PY{p}{)}\PY{p}{;}\PY{+w}{ }\PY{c+c1}{// 42}
\PY{+w}{        }\PY{p}{\PYZcb{}}\PY{+w}{ }\PY{c+c1}{// memory is freed automatically when the arena closes}
\PY{+w}{    }\PY{p}{\PYZcb{}}
\PY{p}{\PYZcb{}}
\end{Verbatim}
\end{codebox}

\begin{codebox}
\begin{Verbatim}[commandchars=\\\{\}]
\PY{k+kn}{import}\PY{+w}{ }\PY{n+nn}{java.lang.ref.Cleaner}\PY{p}{;}

\PY{k+kd}{public}\PY{+w}{ }\PY{k+kd}{class} \PY{n+nc}{ResourceHolder}\PY{+w}{ }\PY{k+kd}{implements}\PY{+w}{ }\PY{n}{AutoCloseable}\PY{+w}{ }\PY{p}{\PYZob{}}
\PY{+w}{    }\PY{k+kd}{private}\PY{+w}{ }\PY{k+kd}{static}\PY{+w}{ }\PY{k+kd}{final}\PY{+w}{ }\PY{n}{Cleaner}\PY{+w}{ }\PY{n}{CLEANER}\PY{+w}{ }\PY{o}{=}\PY{+w}{ }\PY{n}{Cleaner}\PY{p}{.}\PY{n+na}{create}\PY{p}{(}\PY{p}{)}\PY{p}{;}

\PY{+w}{    }\PY{k+kd}{private}\PY{+w}{ }\PY{k+kd}{final}\PY{+w}{ }\PY{n}{Cleaner}\PY{p}{.}\PY{n+na}{Cleanable}\PY{+w}{ }\PY{n}{cleanable}\PY{p}{;}

\PY{+w}{    }\PY{k+kd}{public}\PY{+w}{ }\PY{n+nf}{ResourceHolder}\PY{p}{(}\PY{p}{)}\PY{+w}{ }\PY{p}{\PYZob{}}
\PY{+w}{        }\PY{c+c1}{// register cleanup logic to run if close() is never called and the object is GC\PYZsq{}d}
\PY{+w}{        }\PY{k}{this}\PY{p}{.}\PY{n+na}{cleanable}\PY{+w}{ }\PY{o}{=}\PY{+w}{ }\PY{n}{CLEANER}\PY{p}{.}\PY{n+na}{register}\PY{p}{(}\PY{k}{this}\PY{p}{,}\PY{+w}{ }\PY{p}{(}\PY{p}{)}\PY{+w}{ }\PY{o}{\PYZhy{}}\PY{o}{\PYZgt{}}\PY{+w}{ }\PY{n}{System}\PY{p}{.}\PY{n+na}{out}\PY{p}{.}\PY{n+na}{println}\PY{p}{(}\PY{l+s}{\PYZdq{}}\PY{l+s}{cleaning up native resource}\PY{l+s}{\PYZdq{}}\PY{p}{)}\PY{p}{)}\PY{p}{;}
\PY{+w}{    }\PY{p}{\PYZcb{}}

\PY{+w}{    }\PY{n+nd}{@Override}
\PY{+w}{    }\PY{k+kd}{public}\PY{+w}{ }\PY{k+kt}{void}\PY{+w}{ }\PY{n+nf}{close}\PY{p}{(}\PY{p}{)}\PY{+w}{ }\PY{p}{\PYZob{}}
\PY{+w}{        }\PY{n}{cleanable}\PY{p}{.}\PY{n+na}{clean}\PY{p}{(}\PY{p}{)}\PY{p}{;}\PY{+w}{ }\PY{c+c1}{// idempotent — safe to call once explicitly, GC won\PYZsq{}t run it twice}
\PY{+w}{    }\PY{p}{\PYZcb{}}
\PY{p}{\PYZcb{}}
\end{Verbatim}
\end{codebox}

\section[Service Provider Interface (SPI)]{Service Provider Interface (SPI)}
\label{h:16:service-provider-interface-spi}
\textbf{SPI} is a pattern built into the JDK for writing pluggable applications: you define an interface (the “service”), other jars provide implementations (the “providers”), and your application discovers those implementations at runtime \textbf{without knowing their class names in advance}. This is how JDBC drivers, \code{Charset} providers, and many logging bridges plug themselves into applications.

\subsection[ServiceLoader and META-INF/services]{\code{ServiceLoader} and \code{META-\allowbreak{}INF/\allowbreak{}services}}
\label{h:16:serviceloader-and-meta-infservices}
The classic mechanism, working since Java 6, is a text file under \code{META-\allowbreak{}INF/\allowbreak{}services/} named exactly after the fully-qualified interface name, containing one implementation class name per line.

\textbf{1. Define the service interface:}

\begin{codebox}
\begin{Verbatim}[commandchars=\\\{\}]
\PY{k+kn}{package}\PY{+w}{ }\PY{n+nn}{com.example.spi}\PY{p}{;}

\PY{k+kd}{public}\PY{+w}{ }\PY{k+kd}{interface} \PY{n+nc}{GreetingProvider}\PY{+w}{ }\PY{p}{\PYZob{}}
\PY{+w}{    }\PY{n}{String}\PY{+w}{ }\PY{n+nf}{greet}\PY{p}{(}\PY{n}{String}\PY{+w}{ }\PY{n}{name}\PY{p}{)}\PY{p}{;}
\PY{p}{\PYZcb{}}
\end{Verbatim}
\end{codebox}

\textbf{2. Write two implementations, in the same or different jars:}

\begin{codebox}
\begin{Verbatim}[commandchars=\\\{\}]
\PY{k+kn}{package}\PY{+w}{ }\PY{n+nn}{com.example.spi.impl}\PY{p}{;}
\PY{k+kn}{import}\PY{+w}{ }\PY{n+nn}{com.example.spi.GreetingProvider}\PY{p}{;}

\PY{k+kd}{public}\PY{+w}{ }\PY{k+kd}{class} \PY{n+nc}{FormalGreetingProvider}\PY{+w}{ }\PY{k+kd}{implements}\PY{+w}{ }\PY{n}{GreetingProvider}\PY{+w}{ }\PY{p}{\PYZob{}}
\PY{+w}{    }\PY{k+kd}{public}\PY{+w}{ }\PY{n}{String}\PY{+w}{ }\PY{n+nf}{greet}\PY{p}{(}\PY{n}{String}\PY{+w}{ }\PY{n}{name}\PY{p}{)}\PY{+w}{ }\PY{p}{\PYZob{}}\PY{+w}{ }\PY{k}{return}\PY{+w}{ }\PY{l+s}{\PYZdq{}}\PY{l+s}{Good day, }\PY{l+s}{\PYZdq{}}\PY{+w}{ }\PY{o}{+}\PY{+w}{ }\PY{n}{name}\PY{+w}{ }\PY{o}{+}\PY{+w}{ }\PY{l+s}{\PYZdq{}}\PY{l+s}{.}\PY{l+s}{\PYZdq{}}\PY{p}{;}\PY{+w}{ }\PY{p}{\PYZcb{}}
\PY{p}{\PYZcb{}}
\end{Verbatim}
\end{codebox}

\begin{codebox}
\begin{Verbatim}[commandchars=\\\{\}]
\PY{k+kn}{package}\PY{+w}{ }\PY{n+nn}{com.example.spi.impl}\PY{p}{;}
\PY{k+kn}{import}\PY{+w}{ }\PY{n+nn}{com.example.spi.GreetingProvider}\PY{p}{;}

\PY{k+kd}{public}\PY{+w}{ }\PY{k+kd}{class} \PY{n+nc}{CasualGreetingProvider}\PY{+w}{ }\PY{k+kd}{implements}\PY{+w}{ }\PY{n}{GreetingProvider}\PY{+w}{ }\PY{p}{\PYZob{}}
\PY{+w}{    }\PY{k+kd}{public}\PY{+w}{ }\PY{n}{String}\PY{+w}{ }\PY{n+nf}{greet}\PY{p}{(}\PY{n}{String}\PY{+w}{ }\PY{n}{name}\PY{p}{)}\PY{+w}{ }\PY{p}{\PYZob{}}\PY{+w}{ }\PY{k}{return}\PY{+w}{ }\PY{l+s}{\PYZdq{}}\PY{l+s}{Hey }\PY{l+s}{\PYZdq{}}\PY{+w}{ }\PY{o}{+}\PY{+w}{ }\PY{n}{name}\PY{+w}{ }\PY{o}{+}\PY{+w}{ }\PY{l+s}{\PYZdq{}}\PY{l+s}{!}\PY{l+s}{\PYZdq{}}\PY{p}{;}\PY{+w}{ }\PY{p}{\PYZcb{}}
\PY{p}{\PYZcb{}}
\end{Verbatim}
\end{codebox}

\textbf{3. Register both in a resource file:}

\begin{codebox}
\begin{Verbatim}
src/main/resources/META-INF/services/com.example.spi.GreetingProvider
\end{Verbatim}
\end{codebox}

containing:

\begin{codebox}
\begin{Verbatim}
com.example.spi.impl.FormalGreetingProvider
com.example.spi.impl.CasualGreetingProvider
\end{Verbatim}
\end{codebox}

\textbf{4. Load and use them:}

\begin{codebox}
\begin{Verbatim}[commandchars=\\\{\}]
\PY{k+kn}{import}\PY{+w}{ }\PY{n+nn}{java.util.ServiceLoader}\PY{p}{;}
\PY{k+kn}{import}\PY{+w}{ }\PY{n+nn}{com.example.spi.GreetingProvider}\PY{p}{;}

\PY{k+kd}{public}\PY{+w}{ }\PY{k+kd}{class} \PY{n+nc}{SpiDemo}\PY{+w}{ }\PY{p}{\PYZob{}}
\PY{+w}{    }\PY{k+kd}{public}\PY{+w}{ }\PY{k+kd}{static}\PY{+w}{ }\PY{k+kt}{void}\PY{+w}{ }\PY{n+nf}{main}\PY{p}{(}\PY{n}{String}\PY{o}{[}\PY{o}{]}\PY{+w}{ }\PY{n}{args}\PY{p}{)}\PY{+w}{ }\PY{p}{\PYZob{}}
\PY{+w}{        }\PY{n}{ServiceLoader}\PY{o}{\PYZlt{}}\PY{n}{GreetingProvider}\PY{o}{\PYZgt{}}\PY{+w}{ }\PY{n}{loader}\PY{+w}{ }\PY{o}{=}\PY{+w}{ }\PY{n}{ServiceLoader}\PY{p}{.}\PY{n+na}{load}\PY{p}{(}\PY{n}{GreetingProvider}\PY{p}{.}\PY{n+na}{class}\PY{p}{)}\PY{p}{;}
\PY{+w}{        }\PY{k}{for}\PY{+w}{ }\PY{p}{(}\PY{n}{GreetingProvider}\PY{+w}{ }\PY{n}{provider}\PY{+w}{ }\PY{p}{:}\PY{+w}{ }\PY{n}{loader}\PY{p}{)}\PY{+w}{ }\PY{p}{\PYZob{}}
\PY{+w}{            }\PY{n}{System}\PY{p}{.}\PY{n+na}{out}\PY{p}{.}\PY{n+na}{println}\PY{p}{(}\PY{n}{provider}\PY{p}{.}\PY{n+na}{greet}\PY{p}{(}\PY{l+s}{\PYZdq{}}\PY{l+s}{Ada}\PY{l+s}{\PYZdq{}}\PY{p}{)}\PY{p}{)}\PY{p}{;}
\PY{+w}{        }\PY{p}{\PYZcb{}}
\PY{+w}{    }\PY{p}{\PYZcb{}}
\PY{p}{\PYZcb{}}
\end{Verbatim}
\end{codebox}

Output:

\begin{codebox}
\begin{Verbatim}
Good day, Ada.
Hey Ada!
\end{Verbatim}
\end{codebox}

Neither the interface nor \code{SpiDemo} needed to know the implementation class names — adding a third provider jar to the classpath, with its own \code{META-\allowbreak{}INF/\allowbreak{}services} entry, would make it show up automatically, with no code changes.

\subsection[The Module System Way: provides ... with]{The Module System Way: \code{provides~\allowbreak{}...\allowbreak{}~\allowbreak{}with}}
\label{h:16:the-module-system-way-provides--with}
If you’re using JPMS (the Java Module System), the equivalent declaration goes in \code{module-\allowbreak{}info.\allowbreak{}java} instead of a text file, and is checked at compile time:

\begin{codebox}
\begin{Verbatim}[commandchars=\\\{\}]
\PY{c+c1}{// in the module that declares the service interface}
\PY{n}{module}\PY{+w}{ }\PY{n}{com}\PY{p}{.}\PY{n+na}{example}\PY{p}{.}\PY{n+na}{spi}\PY{+w}{ }\PY{p}{\PYZob{}}
\PY{+w}{    }\PY{n}{exports}\PY{+w}{ }\PY{n}{com}\PY{p}{.}\PY{n+na}{example}\PY{p}{.}\PY{n+na}{spi}\PY{p}{;}
\PY{+w}{    }\PY{n}{uses}\PY{+w}{ }\PY{n}{com}\PY{p}{.}\PY{n+na}{example}\PY{p}{.}\PY{n+na}{spi}\PY{p}{.}\PY{n+na}{GreetingProvider}\PY{p}{;}
\PY{p}{\PYZcb{}}
\end{Verbatim}
\end{codebox}

\begin{codebox}
\begin{Verbatim}[commandchars=\\\{\}]
\PY{c+c1}{// in the module that provides an implementation}
\PY{n}{module}\PY{+w}{ }\PY{n}{com}\PY{p}{.}\PY{n+na}{example}\PY{p}{.}\PY{n+na}{spi}\PY{p}{.}\PY{n+na}{impl}\PY{+w}{ }\PY{p}{\PYZob{}}
\PY{+w}{    }\PY{n}{requires}\PY{+w}{ }\PY{n}{com}\PY{p}{.}\PY{n+na}{example}\PY{p}{.}\PY{n+na}{spi}\PY{p}{;}
\PY{+w}{    }\PY{n}{provides}\PY{+w}{ }\PY{n}{com}\PY{p}{.}\PY{n+na}{example}\PY{p}{.}\PY{n+na}{spi}\PY{p}{.}\PY{n+na}{GreetingProvider}
\PY{+w}{        }\PY{n}{with}\PY{+w}{ }\PY{n}{com}\PY{p}{.}\PY{n+na}{example}\PY{p}{.}\PY{n+na}{spi}\PY{p}{.}\PY{n+na}{impl}\PY{p}{.}\PY{n+na}{FormalGreetingProvider}\PY{p}{,}\PY{+w}{ }\PY{n}{com}\PY{p}{.}\PY{n+na}{example}\PY{p}{.}\PY{n+na}{spi}\PY{p}{.}\PY{n+na}{impl}\PY{p}{.}\PY{n+na}{CasualGreetingProvider}\PY{p}{;}
\PY{p}{\PYZcb{}}
\end{Verbatim}
\end{codebox}

\code{Service\allowbreak{}Loader.\allowbreak{}load(...)} works exactly the same way at the call site — it transparently checks both \code{module-\allowbreak{}info.\allowbreak{}java} declarations and legacy \code{META-\allowbreak{}INF/\allowbreak{}services} files.

\subsection[Real JDK Uses]{Real JDK Uses}
\label{h:16:real-jdk-uses}
\begin{itemize}
\item 
\textbf{JDBC drivers}: since JDBC 4.0, drivers register themselves via \code{META-\allowbreak{}INF/\allowbreak{}services/\allowbreak{}java.\allowbreak{}sql.\allowbreak{}Driver}, which is why modern code doesn’t need \code{Class.\allowbreak{}for\allowbreak{}Name(\allowbreak{}\textquotedbl{}com.\allowbreak{}mysql.\allowbreak{}cj.\allowbreak{}jdbc.\allowbreak{}Driver\textquotedbl{})} anymore — \code{DriverManager} uses \code{ServiceLoader} internally to find it just by having the jar on the classpath.

\item 
\textbf{\code{java.\allowbreak{}nio.\allowbreak{}charset.\allowbreak{}spi.\allowbreak{}Charset\allowbreak{}Provider}}: lets libraries add custom character encodings that \code{Charset.\allowbreak{}forName(...)} can then find.

\item 
\textbf{Logging bridges}: SLF4J and the Java Logging API use SPI-like discovery to find and bind the actual logging backend (Logback, Log4j2, etc.) at startup, without the application code referencing a specific logging implementation directly.

\end{itemize}

\section[Common Code-Review Interview Pitfalls]{Common Code-Review Interview Pitfalls}
\label{h:16:common-code-review-interview-pitfalls}
\begin{enumerate}
\item 
\textbf{Using \code{Class.\allowbreak{}forName(...)} or reflection for a class that’s already known at compile time.}
Why it matters: this throws away compile-time type safety and IDE support for no benefit — reflection should be reserved for cases where the target type genuinely isn’t known until runtime.

\begin{codebox}
\begin{Verbatim}[commandchars=\\\{\}]
\PY{c+c1}{// Before}
\PY{n}{Object}\PY{+w}{ }\PY{n}{svc}\PY{+w}{ }\PY{o}{=}\PY{+w}{ }\PY{n}{Class}\PY{p}{.}\PY{n+na}{forName}\PY{p}{(}\PY{l+s}{\PYZdq{}}\PY{l+s}{com.example.OrderService}\PY{l+s}{\PYZdq{}}\PY{p}{)}\PY{p}{.}\PY{n+na}{getDeclaredConstructor}\PY{p}{(}\PY{p}{)}\PY{p}{.}\PY{n+na}{newInstance}\PY{p}{(}\PY{p}{)}\PY{p}{;}
\PY{c+c1}{// After}
\PY{n}{OrderService}\PY{+w}{ }\PY{n}{svc}\PY{+w}{ }\PY{o}{=}\PY{+w}{ }\PY{k}{new}\PY{+w}{ }\PY{n}{OrderService}\PY{p}{(}\PY{p}{)}\PY{p}{;}
\end{Verbatim}
\end{codebox}

\item 
\textbf{Calling \code{setAccessible(\allowbreak{}true)} on a field or method just to avoid adding a getter/constructor parameter.}
Why it matters: it silently breaks encapsulation, is fragile across refactors (renaming the field breaks it with no compile error, just a runtime \code{NoSuchFieldException}), and will fail outright with \code{Inaccessible\allowbreak{}Object\allowbreak{}Exception} on modularized JDK internals starting with Java 17.

\begin{codebox}
\begin{Verbatim}[commandchars=\\\{\}]
\PY{c+c1}{// Before}
\PY{n}{Field}\PY{+w}{ }\PY{n}{f}\PY{+w}{ }\PY{o}{=}\PY{+w}{ }\PY{n}{obj}\PY{p}{.}\PY{n+na}{getClass}\PY{p}{(}\PY{p}{)}\PY{p}{.}\PY{n+na}{getDeclaredField}\PY{p}{(}\PY{l+s}{\PYZdq{}}\PY{l+s}{balance}\PY{l+s}{\PYZdq{}}\PY{p}{)}\PY{p}{;}
\PY{n}{f}\PY{p}{.}\PY{n+na}{setAccessible}\PY{p}{(}\PY{k+kc}{true}\PY{p}{)}\PY{p}{;}
\PY{k+kt}{double}\PY{+w}{ }\PY{n}{balance}\PY{+w}{ }\PY{o}{=}\PY{+w}{ }\PY{p}{(}\PY{k+kt}{double}\PY{p}{)}\PY{+w}{ }\PY{n}{f}\PY{p}{.}\PY{n+na}{get}\PY{p}{(}\PY{n}{obj}\PY{p}{)}\PY{p}{;}
\PY{c+c1}{// After}
\PY{k+kt}{double}\PY{+w}{ }\PY{n}{balance}\PY{+w}{ }\PY{o}{=}\PY{+w}{ }\PY{n}{obj}\PY{p}{.}\PY{n+na}{getBalance}\PY{p}{(}\PY{p}{)}\PY{p}{;}
\end{Verbatim}
\end{codebox}

\item 
\textbf{Not caching \code{Method}/\code{Field}/\code{Constructor} lookups that happen inside a loop or a hot request path.}
Why it matters: \code{getDeclaredMethod}, \code{getField}, etc. are relatively expensive lookups; doing them on every call instead of once at startup can dominate the cost of an otherwise fast method.

\begin{codebox}
\begin{Verbatim}[commandchars=\\\{\}]
\PY{c+c1}{// Before}
\PY{k}{for}\PY{+w}{ }\PY{p}{(}\PY{n}{Order}\PY{+w}{ }\PY{n}{o}\PY{+w}{ }\PY{p}{:}\PY{+w}{ }\PY{n}{orders}\PY{p}{)}\PY{+w}{ }\PY{p}{\PYZob{}}
\PY{+w}{    }\PY{n}{Method}\PY{+w}{ }\PY{n}{m}\PY{+w}{ }\PY{o}{=}\PY{+w}{ }\PY{n}{o}\PY{p}{.}\PY{n+na}{getClass}\PY{p}{(}\PY{p}{)}\PY{p}{.}\PY{n+na}{getMethod}\PY{p}{(}\PY{l+s}{\PYZdq{}}\PY{l+s}{getTotal}\PY{l+s}{\PYZdq{}}\PY{p}{)}\PY{p}{;}
\PY{+w}{    }\PY{n}{sum}\PY{+w}{ }\PY{o}{+}\PY{o}{=}\PY{+w}{ }\PY{p}{(}\PY{k+kt}{double}\PY{p}{)}\PY{+w}{ }\PY{n}{m}\PY{p}{.}\PY{n+na}{invoke}\PY{p}{(}\PY{n}{o}\PY{p}{)}\PY{p}{;}
\PY{p}{\PYZcb{}}
\PY{c+c1}{// After}
\PY{n}{Method}\PY{+w}{ }\PY{n}{m}\PY{+w}{ }\PY{o}{=}\PY{+w}{ }\PY{n}{Order}\PY{p}{.}\PY{n+na}{class}\PY{p}{.}\PY{n+na}{getMethod}\PY{p}{(}\PY{l+s}{\PYZdq{}}\PY{l+s}{getTotal}\PY{l+s}{\PYZdq{}}\PY{p}{)}\PY{p}{;}\PY{+w}{ }\PY{c+c1}{// looked up once, outside the loop}
\PY{k}{for}\PY{+w}{ }\PY{p}{(}\PY{n}{Order}\PY{+w}{ }\PY{n}{o}\PY{+w}{ }\PY{p}{:}\PY{+w}{ }\PY{n}{orders}\PY{p}{)}\PY{+w}{ }\PY{n}{sum}\PY{+w}{ }\PY{o}{+}\PY{o}{=}\PY{+w}{ }\PY{p}{(}\PY{k+kt}{double}\PY{p}{)}\PY{+w}{ }\PY{n}{m}\PY{p}{.}\PY{n+na}{invoke}\PY{p}{(}\PY{n}{o}\PY{p}{)}\PY{p}{;}
\end{Verbatim}
\end{codebox}

\item 
\textbf{Assuming \code{RetentionPolicy.\allowbreak{}CLASS} (the default) annotations are readable via reflection.}
Why it matters: only \code{RUNTIME}-retained annotations are visible to \code{getAnnotation(...)}; forgetting \code{@\allowbreak{}Retention(\allowbreak{}Retention\allowbreak{}Policy.\allowbreak{}RUNTIME)} makes \code{isAnnotationPresent} silently return \code{false} at runtime with no compile error.

\begin{codebox}
\begin{Verbatim}[commandchars=\\\{\}]
\PY{c+c1}{// Before — no @Retention at all defaults to CLASS, invisible at runtime}
\PY{k+kd}{public}\PY{+w}{ }\PY{n+nd}{@interface}\PY{+w}{ }\PY{n}{Auditable}\PY{+w}{ }\PY{p}{\PYZob{}}\PY{+w}{ }\PY{p}{\PYZcb{}}
\PY{c+c1}{// After}
\PY{n+nd}{@Retention}\PY{p}{(}\PY{n}{RetentionPolicy}\PY{p}{.}\PY{n+na}{RUNTIME}\PY{p}{)}
\PY{k+kd}{public}\PY{+w}{ }\PY{n+nd}{@interface}\PY{+w}{ }\PY{n}{Auditable}\PY{+w}{ }\PY{p}{\PYZob{}}\PY{+w}{ }\PY{p}{\PYZcb{}}
\end{Verbatim}
\end{codebox}

\item 
\textbf{Trying to create a JDK dynamic proxy for a class instead of an interface.}
Why it matters: \code{java.\allowbreak{}lang.\allowbreak{}reflect.\allowbreak{}Proxy} only supports interfaces; attempting to pass a concrete class to \code{newProxyInstance} throws \code{Illegal\allowbreak{}Argument\allowbreak{}Exception} — the candidate needs to know to reach for CGLIB/ByteBuddy (or Spring’s \code{proxyTargetClass=\allowbreak{}true}) instead.

\begin{codebox}
\begin{Verbatim}[commandchars=\\\{\}]
\PY{c+c1}{// Before — throws IllegalArgumentException at runtime}
\PY{n}{Proxy}\PY{p}{.}\PY{n+na}{newProxyInstance}\PY{p}{(}\PY{n}{cl}\PY{p}{,}\PY{+w}{ }\PY{k}{new}\PY{+w}{ }\PY{n}{Class}\PY{o}{\PYZlt{}}\PY{o}{?}\PY{o}{\PYZgt{}}\PY{o}{[}\PY{o}{]}\PY{p}{\PYZob{}}\PY{+w}{ }\PY{n}{OrderServiceImpl}\PY{p}{.}\PY{n+na}{class}\PY{+w}{ }\PY{p}{\PYZcb{}}\PY{p}{,}\PY{+w}{ }\PY{n}{handler}\PY{p}{)}\PY{p}{;}
\PY{c+c1}{// After — proxy the interface it implements}
\PY{n}{Proxy}\PY{p}{.}\PY{n+na}{newProxyInstance}\PY{p}{(}\PY{n}{cl}\PY{p}{,}\PY{+w}{ }\PY{k}{new}\PY{+w}{ }\PY{n}{Class}\PY{o}{\PYZlt{}}\PY{o}{?}\PY{o}{\PYZgt{}}\PY{o}{[}\PY{o}{]}\PY{p}{\PYZob{}}\PY{+w}{ }\PY{n}{OrderService}\PY{p}{.}\PY{n+na}{class}\PY{+w}{ }\PY{p}{\PYZcb{}}\PY{p}{,}\PY{+w}{ }\PY{n}{handler}\PY{p}{)}\PY{p}{;}
\end{Verbatim}
\end{codebox}

\item 
\textbf{Forgetting that a Spring \code{@\allowbreak{}Transactional} (or similar AOP-advised) method call from \emph{within the same class} bypasses the proxy.}
Why it matters: AOP proxies (whether JDK or CGLIB) only intercept calls that come in through the proxy object from the \emph{outside}; an internal \code{this.\allowbreak{}otherMethod()} call skips the proxy entirely, silently disabling the transaction, caching, or security check.

\begin{codebox}
\begin{Verbatim}[commandchars=\\\{\}]
\PY{c+c1}{// Before — internal call bypasses the transactional proxy}
\PY{k+kd}{public}\PY{+w}{ }\PY{k+kt}{void}\PY{+w}{ }\PY{n+nf}{placeOrder}\PY{p}{(}\PY{p}{)}\PY{+w}{ }\PY{p}{\PYZob{}}
\PY{+w}{    }\PY{k}{this}\PY{p}{.}\PY{n+na}{chargeCard}\PY{p}{(}\PY{p}{)}\PY{p}{;}\PY{+w}{ }\PY{c+c1}{// @Transactional on chargeCard() has no effect here}
\PY{p}{\PYZcb{}}
\PY{c+c1}{// After — inject/call through the bean so the proxy intercepts it}
\PY{k+kd}{public}\PY{+w}{ }\PY{k+kt}{void}\PY{+w}{ }\PY{n+nf}{placeOrder}\PY{p}{(}\PY{p}{)}\PY{+w}{ }\PY{p}{\PYZob{}}
\PY{+w}{    }\PY{n}{orderServiceSelfProxy}\PY{p}{.}\PY{n+na}{chargeCard}\PY{p}{(}\PY{p}{)}\PY{p}{;}
\PY{p}{\PYZcb{}}
\end{Verbatim}
\end{codebox}

\item 
\textbf{Mixing up \code{invoke} and \code{invokeExact} on a \code{MethodHandle} and expecting automatic type adaptation.}
Why it matters: \code{invokeExact} requires the exact \code{MethodType}, including matching primitive vs. boxed types; a mismatch throws \code{Wrong\allowbreak{}Method\allowbreak{}Type\allowbreak{}Exception} at runtime instead of a compile error, which is easy to miss in review since the code compiles fine.

\begin{codebox}
\begin{Verbatim}[commandchars=\\\{\}]
\PY{c+c1}{// Before — mh\PYZsq{}s MethodType is (int)int, throws WrongMethodTypeException}
\PY{n}{Object}\PY{+w}{ }\PY{n}{result}\PY{+w}{ }\PY{o}{=}\PY{+w}{ }\PY{n}{mh}\PY{p}{.}\PY{n+na}{invokeExact}\PY{p}{(}\PY{p}{(}\PY{n}{Object}\PY{p}{)}\PY{+w}{ }\PY{l+m+mi}{5}\PY{p}{)}\PY{p}{;}
\PY{c+c1}{// After}
\PY{k+kt}{int}\PY{+w}{ }\PY{n}{result}\PY{+w}{ }\PY{o}{=}\PY{+w}{ }\PY{p}{(}\PY{k+kt}{int}\PY{p}{)}\PY{+w}{ }\PY{n}{mh}\PY{p}{.}\PY{n+na}{invokeExact}\PY{p}{(}\PY{l+m+mi}{5}\PY{p}{)}\PY{p}{;}
\end{Verbatim}
\end{codebox}

\item 
\textbf{Using \code{Atomic\allowbreak{}Integer\allowbreak{}Field\allowbreak{}Updater}/\code{AtomicLongFieldUpdater} in new code instead of \code{VarHandle}.}
Why it matters: the field-updater classes are older, more error-prone (the field must be exactly \code{volatile} and non-private relative to the updater’s caller, with easy-to-miss reflective setup), and \code{VarHandle} supersedes them with clearer semantics and better performance.

\begin{codebox}
\begin{Verbatim}[commandchars=\\\{\}]
\PY{c+c1}{// Before}
\PY{k+kd}{private}\PY{+w}{ }\PY{k+kd}{static}\PY{+w}{ }\PY{k+kd}{final}\PY{+w}{ }\PY{n}{AtomicIntegerFieldUpdater}\PY{o}{\PYZlt{}}\PY{n}{Counter}\PY{o}{\PYZgt{}}\PY{+w}{ }\PY{n}{UPDATER}\PY{+w}{ }\PY{o}{=}
\PY{+w}{    }\PY{n}{AtomicIntegerFieldUpdater}\PY{p}{.}\PY{n+na}{newUpdater}\PY{p}{(}\PY{n}{Counter}\PY{p}{.}\PY{n+na}{class}\PY{p}{,}\PY{+w}{ }\PY{l+s}{\PYZdq{}}\PY{l+s}{count}\PY{l+s}{\PYZdq{}}\PY{p}{)}\PY{p}{;}
\PY{c+c1}{// After}
\PY{k+kd}{private}\PY{+w}{ }\PY{k+kd}{static}\PY{+w}{ }\PY{k+kd}{final}\PY{+w}{ }\PY{n}{VarHandle}\PY{+w}{ }\PY{n}{COUNT\PYZus{}HANDLE}\PY{p}{;}
\PY{k+kd}{static}\PY{+w}{ }\PY{p}{\PYZob{}}
\PY{+w}{    }\PY{k}{try}\PY{+w}{ }\PY{p}{\PYZob{}}\PY{+w}{ }\PY{n}{COUNT\PYZus{}HANDLE}\PY{+w}{ }\PY{o}{=}\PY{+w}{ }\PY{n}{MethodHandles}\PY{p}{.}\PY{n+na}{lookup}\PY{p}{(}\PY{p}{)}\PY{p}{.}\PY{n+na}{findVarHandle}\PY{p}{(}\PY{n}{Counter}\PY{p}{.}\PY{n+na}{class}\PY{p}{,}\PY{+w}{ }\PY{l+s}{\PYZdq{}}\PY{l+s}{count}\PY{l+s}{\PYZdq{}}\PY{p}{,}\PY{+w}{ }\PY{k+kt}{int}\PY{p}{.}\PY{n+na}{class}\PY{p}{)}\PY{p}{;}\PY{+w}{ }\PY{p}{\PYZcb{}}
\PY{+w}{    }\PY{k}{catch}\PY{+w}{ }\PY{p}{(}\PY{n}{ReflectiveOperationException}\PY{+w}{ }\PY{n}{e}\PY{p}{)}\PY{+w}{ }\PY{p}{\PYZob{}}\PY{+w}{ }\PY{k}{throw}\PY{+w}{ }\PY{k}{new}\PY{+w}{ }\PY{n}{ExceptionInInitializerError}\PY{p}{(}\PY{n}{e}\PY{p}{)}\PY{p}{;}\PY{+w}{ }\PY{p}{\PYZcb{}}
\PY{p}{\PYZcb{}}
\end{Verbatim}
\end{codebox}

\item 
\textbf{Using \code{getVolatile}/\code{setVolatile} on a \code{VarHandle} when a weaker (and cheaper) access mode like \code{getOpaque} or plain access would suffice.}
Why it matters: full volatile semantics impose a memory fence on every access; if the algorithm doesn’t actually need cross-thread ordering guarantees for that particular read, using the strongest mode by default is a needless performance cost, and interviewers look for awareness that access modes are a real trade-off, not just decoration.

\begin{codebox}
\begin{Verbatim}[commandchars=\\\{\}]
\PY{c+c1}{// Before — full fence on every read, even for a value only this thread touches}
\PY{k+kt}{int}\PY{+w}{ }\PY{n}{v}\PY{+w}{ }\PY{o}{=}\PY{+w}{ }\PY{p}{(}\PY{k+kt}{int}\PY{p}{)}\PY{+w}{ }\PY{n}{handle}\PY{p}{.}\PY{n+na}{getVolatile}\PY{p}{(}\PY{n}{obj}\PY{p}{)}\PY{p}{;}
\PY{c+c1}{// After — plain access when there\PYZsq{}s no cross\PYZhy{}thread visibility requirement}
\PY{k+kt}{int}\PY{+w}{ }\PY{n}{v}\PY{+w}{ }\PY{o}{=}\PY{+w}{ }\PY{p}{(}\PY{k+kt}{int}\PY{p}{)}\PY{+w}{ }\PY{n}{handle}\PY{p}{.}\PY{n+na}{get}\PY{p}{(}\PY{n}{obj}\PY{p}{)}\PY{p}{;}
\end{Verbatim}
\end{codebox}

\item 
\textbf{Any new code calling \code{sun.\allowbreak{}misc.\allowbreak{}Unsafe} directly.}
Why it matters: \code{Unsafe} is an unsupported internal API being deprecated for removal (JEP 471/498); it bypasses the type system, breaks under strong encapsulation, and has fully supported modern replacements — flagging this in review prevents a maintenance/upgrade landmine.

\begin{codebox}
\begin{Verbatim}[commandchars=\\\{\}]
\PY{c+c1}{// Before}
\PY{n}{unsafe}\PY{p}{.}\PY{n+na}{putInt}\PY{p}{(}\PY{n}{obj}\PY{p}{,}\PY{+w}{ }\PY{n}{offset}\PY{p}{,}\PY{+w}{ }\PY{l+m+mi}{42}\PY{p}{)}\PY{p}{;}
\PY{c+c1}{// After}
\PY{n}{intHandle}\PY{p}{.}\PY{n+na}{set}\PY{p}{(}\PY{n}{obj}\PY{p}{,}\PY{+w}{ }\PY{l+m+mi}{42}\PY{p}{)}\PY{p}{;}\PY{+w}{ }\PY{c+c1}{// via VarHandle}
\end{Verbatim}
\end{codebox}

\item 
\textbf{Registering a \code{ServiceLoader} provider without a public no-arg constructor.}
Why it matters: \code{ServiceLoader} instantiates providers reflectively via a public no-arg constructor by default; a provider with only a parameterized constructor (or a private one) fails at load time with \code{Service\allowbreak{}Configuration\allowbreak{}Error}, often discovered only in production when that particular provider is finally needed.

\begin{codebox}
\begin{Verbatim}[commandchars=\\\{\}]
\PY{c+c1}{// Before — no no\PYZhy{}arg constructor, ServiceConfigurationError at load time}
\PY{k+kd}{public}\PY{+w}{ }\PY{k+kd}{class} \PY{n+nc}{DbGreetingProvider}\PY{+w}{ }\PY{k+kd}{implements}\PY{+w}{ }\PY{n}{GreetingProvider}\PY{+w}{ }\PY{p}{\PYZob{}}
\PY{+w}{    }\PY{k+kd}{public}\PY{+w}{ }\PY{n+nf}{DbGreetingProvider}\PY{p}{(}\PY{n}{DataSource}\PY{+w}{ }\PY{n}{ds}\PY{p}{)}\PY{+w}{ }\PY{p}{\PYZob{}}\PY{+w}{ }\PY{p}{.}\PY{p}{.}\PY{p}{.}\PY{+w}{ }\PY{p}{\PYZcb{}}
\PY{p}{\PYZcb{}}
\PY{c+c1}{// After — provide a public no\PYZhy{}arg constructor, or a static \PYZdq{}provider()\PYZdq{} factory method}
\PY{k+kd}{public}\PY{+w}{ }\PY{k+kd}{class} \PY{n+nc}{DbGreetingProvider}\PY{+w}{ }\PY{k+kd}{implements}\PY{+w}{ }\PY{n}{GreetingProvider}\PY{+w}{ }\PY{p}{\PYZob{}}
\PY{+w}{    }\PY{k+kd}{public}\PY{+w}{ }\PY{n+nf}{DbGreetingProvider}\PY{p}{(}\PY{p}{)}\PY{+w}{ }\PY{p}{\PYZob{}}\PY{+w}{ }\PY{k}{this}\PY{p}{.}\PY{n+na}{ds}\PY{+w}{ }\PY{o}{=}\PY{+w}{ }\PY{n}{DataSourceRegistry}\PY{p}{.}\PY{n+na}{getDefault}\PY{p}{(}\PY{p}{)}\PY{p}{;}\PY{+w}{ }\PY{p}{\PYZcb{}}
\PY{p}{\PYZcb{}}
\end{Verbatim}
\end{codebox}

\item 
\textbf{Mismatched or misspelled \code{META-\allowbreak{}INF/\allowbreak{}services} file name vs. the actual interface’s fully-qualified name.}
Why it matters: the file name must exactly match the service interface’s fully-qualified class name; a typo or a stale name after a package rename means \code{Service\allowbreak{}Loader.\allowbreak{}load(...)} silently finds zero providers instead of failing loudly, which is confusing to debug.

\begin{codebox}
\begin{Verbatim}
# Before — file named for the old package after a refactor
META-INF/services/com.example.old.GreetingProvider
# After — matches the current interface location
META-INF/services/com.example.spi.GreetingProvider
\end{Verbatim}
\end{codebox}

\item 
\textbf{Writing a custom annotation processor (or Lombok-style tool) that mutates behavior at runtime via reflection when compile-time code generation would work and be far faster.}
Why it matters: reflection-based runtime frameworks pay a lookup/invocation cost on every use and fight the module system’s strong encapsulation; compile-time generation (like MapStruct, Dagger, Immutables) produces plain, fast, debuggable Java with the same errors caught at build time instead of at runtime.

\begin{codebox}
\begin{Verbatim}[commandchars=\\\{\}]
\PY{c+c1}{// Before — reflective field copy at runtime, every call}
\PY{k}{for}\PY{+w}{ }\PY{p}{(}\PY{n}{Field}\PY{+w}{ }\PY{n}{f}\PY{+w}{ }\PY{p}{:}\PY{+w}{ }\PY{n}{Dto}\PY{p}{.}\PY{n+na}{class}\PY{p}{.}\PY{n+na}{getDeclaredFields}\PY{p}{(}\PY{p}{)}\PY{p}{)}\PY{+w}{ }\PY{p}{\PYZob{}}\PY{+w}{ }\PY{n}{f}\PY{p}{.}\PY{n+na}{setAccessible}\PY{p}{(}\PY{k+kc}{true}\PY{p}{)}\PY{p}{;}\PY{+w}{ }\PY{p}{.}\PY{p}{.}\PY{p}{.}\PY{+w}{ }\PY{p}{\PYZcb{}}
\PY{c+c1}{// After — a MapStruct\PYZhy{}generated mapper compiled ahead of time}
\PY{n}{Dto}\PY{+w}{ }\PY{n}{dto}\PY{+w}{ }\PY{o}{=}\PY{+w}{ }\PY{n}{DtoMapper}\PY{p}{.}\PY{n+na}{INSTANCE}\PY{p}{.}\PY{n+na}{toDto}\PY{p}{(}\PY{n}{entity}\PY{p}{)}\PY{p}{;}
\end{Verbatim}
\end{codebox}

\item 
\textbf{Assuming \code{@\allowbreak{}Inherited} propagates an annotation across methods or fields, not just class hierarchies.}
Why it matters: \code{@\allowbreak{}Inherited} only affects whether a \emph{class-level} annotation is visible on subclasses via \code{getAnnotation}; it has no effect on method or field annotations, so code that expects an overridden method to “inherit” an annotation from the parent’s method will get \code{null} back.

\begin{codebox}
\begin{Verbatim}[commandchars=\\\{\}]
\PY{c+c1}{// Before — expecting method\PYZhy{}level @Inherited to work}
\PY{n+nd}{@Retention}\PY{p}{(}\PY{n}{RetentionPolicy}\PY{p}{.}\PY{n+na}{RUNTIME}\PY{p}{)}
\PY{n+nd}{@Inherited}
\PY{n+nd}{@Target}\PY{p}{(}\PY{n}{ElementType}\PY{p}{.}\PY{n+na}{METHOD}\PY{p}{)}
\PY{n+nd}{@interface}\PY{+w}{ }\PY{n}{Cached}\PY{+w}{ }\PY{p}{\PYZob{}}\PY{+w}{ }\PY{p}{\PYZcb{}}
\PY{c+c1}{// Derived.class.getMethod(\PYZdq{}compute\PYZdq{}).getAnnotation(Cached.class) still returns null}
\PY{c+c1}{// After — re\PYZhy{}annotate the overriding method explicitly, or check the superclass method too}
\PY{k+kd}{class} \PY{n+nc}{Derived}\PY{+w}{ }\PY{k+kd}{extends}\PY{+w}{ }\PY{n}{Base}\PY{+w}{ }\PY{p}{\PYZob{}}
\PY{+w}{    }\PY{n+nd}{@Cached}
\PY{+w}{    }\PY{n+nd}{@Override}
\PY{+w}{    }\PY{k+kt}{void}\PY{+w}{ }\PY{n+nf}{compute}\PY{p}{(}\PY{p}{)}\PY{+w}{ }\PY{p}{\PYZob{}}\PY{+w}{ }\PY{p}{.}\PY{p}{.}\PY{p}{.}\PY{+w}{ }\PY{p}{\PYZcb{}}
\PY{p}{\PYZcb{}}
\end{Verbatim}
\end{codebox}

\item 
\textbf{Ignoring \code{Proxy.\allowbreak{}is\allowbreak{}Proxy\allowbreak{}Class(...)} / unwrapping needs when logging, serializing, or type-checking an object that might be a dynamic proxy.}
Why it matters: a dynamic proxy’s \code{getClass()} returns a synthetic proxy class, not the real implementation type; code that does \code{instanceof~\allowbreak{}Some\allowbreak{}Concrete\allowbreak{}Class} or relies on \code{get\allowbreak{}Class().\allowbreak{}get\allowbreak{}Simple\allowbreak{}Name()} for logging can behave unexpectedly when handed a proxied bean (common with Spring AOP-advised beans).

\begin{codebox}
\begin{Verbatim}[commandchars=\\\{\}]
\PY{c+c1}{// Before — logs a generated proxy class name like \PYZdq{}\PYZdl{}Proxy42\PYZdq{}, not the real service}
\PY{n}{log}\PY{p}{.}\PY{n+na}{info}\PY{p}{(}\PY{l+s}{\PYZdq{}}\PY{l+s}{Handler: \PYZob{}\PYZcb{}}\PY{l+s}{\PYZdq{}}\PY{p}{,}\PY{+w}{ }\PY{n}{handler}\PY{p}{.}\PY{n+na}{getClass}\PY{p}{(}\PY{p}{)}\PY{p}{.}\PY{n+na}{getSimpleName}\PY{p}{(}\PY{p}{)}\PY{p}{)}\PY{p}{;}
\PY{c+c1}{// After — use the interface, or AopUtils.getTargetClass(handler) in Spring}
\PY{n}{log}\PY{p}{.}\PY{n+na}{info}\PY{p}{(}\PY{l+s}{\PYZdq{}}\PY{l+s}{Handler: \PYZob{}\PYZcb{}}\PY{l+s}{\PYZdq{}}\PY{p}{,}\PY{+w}{ }\PY{n}{AopUtils}\PY{p}{.}\PY{n+na}{getTargetClass}\PY{p}{(}\PY{n}{handler}\PY{p}{)}\PY{p}{.}\PY{n+na}{getSimpleName}\PY{p}{(}\PY{p}{)}\PY{p}{)}\PY{p}{;}
\end{Verbatim}
\end{codebox}

\end{enumerate}


\setcounter{chapter}{16}
\chapter{Modules \& Native Interoperability}
\label{chap:17}
\label{h:17:17-modules--native-interoperability}
Java code does not live in a vacuum. It is organized into \textbf{packages}, packages are (optionally) organized into \textbf{modules}, and sometimes Java code needs to reach outside the JVM entirely to call \textbf{native code} or squeeze more performance out of the CPU with \textbf{vectorized} operations. This chapter walks from the oldest and most universal organizational unit (the package, present since Java 1.0) to the newest, still-incubating performance API (the Vector API). Along the way we cover the Java Platform Module System (JPMS, Java 9+) in depth, and the modern replacement for JNI: the Foreign Function \& Memory API (finalized in JDK 22). We target Java 21+ throughout, calling out version-specific behavior where it matters.

\begin{itemize}
\item 
\hyperref[h:17:packages]{Packages}

\begin{itemize}
\item 
\hyperref[h:17:naming-conventions-and-reverse-domain-rules]{Naming Conventions and Reverse-Domain Rules}

\item 
\hyperref[h:17:the-package-and-import-statements]{The \code{package} and \code{import} Statements}

\item 
\hyperref[h:17:static-imports]{Static Imports}

\item 
\hyperref[h:17:wildcard-imports]{Wildcard Imports}

\item 
\hyperref[h:17:the-default-unnamed-package]{The Default (Unnamed) Package}

\item 
\hyperref[h:17:package-private-visibility-as-a-design-tool]{Package-Private Visibility as a Design Tool}

\item 
\hyperref[h:17:package-infojava]{\code{package-\allowbreak{}info.\allowbreak{}java}}

\item 
\hyperref[h:17:split-packages]{Split Packages}

\item 
\hyperref[h:17:directory-to-package-mapping]{Directory-to-Package Mapping}

\item 
\hyperref[h:17:packages-vs-modules]{Packages vs Modules}

\end{itemize}

\item 
\hyperref[h:17:java-platform-module-system-jpms]{Java Platform Module System (JPMS)}

\begin{itemize}
\item 
\hyperref[h:17:problems-jpms-solves]{Problems JPMS Solves}

\item 
\hyperref[h:17:module-infojava-basics]{\code{module-\allowbreak{}info.\allowbreak{}java} Basics}

\item 
\hyperref[h:17:requires-requires-transitive-requires-static]{\code{requires}, \code{requires~\allowbreak{}transitive}, \code{requires~\allowbreak{}static}}

\item 
\hyperref[h:17:exports-and-exports--to]{\code{exports} and \code{exports~\allowbreak{}...\allowbreak{}~\allowbreak{}to}}

\item 
\hyperref[h:17:opens-and-opens--to]{\code{opens} and \code{opens~\allowbreak{}...\allowbreak{}~\allowbreak{}to}}

\item 
\hyperref[h:17:uses-and-provides--with]{\code{uses} and \code{provides~\allowbreak{}...\allowbreak{}~\allowbreak{}with}}

\item 
\hyperref[h:17:automatic-modules-and-the-unnamed-module]{Automatic Modules and the Unnamed Module}

\item 
\hyperref[h:17:module-path-vs-class-path]{Module Path vs Class Path}

\item 
\hyperref[h:17:readability-and-accessibility-rules]{Readability and Accessibility Rules}

\item 
\hyperref[h:17:jlink-jdeps-and-jmod]{\code{jlink}, \code{jdeps}, and \code{jmod}}

\item 
\hyperref[h:17:strong-encapsulation-and---add-exports--add-opens]{Strong Encapsulation and \code{--\allowbreak{}add-\allowbreak{}exports}/\code{--\allowbreak{}add-\allowbreak{}opens}}

\item 
\hyperref[h:17:worked-multi-module-example]{Worked Multi-Module Example}

\item 
\hyperref[h:17:migration-strategy]{Migration Strategy}

\item 
\hyperref[h:17:why-many-projects-still-dont-use-jpms]{Why Many Projects Still Don’t Use JPMS}

\end{itemize}

\item 
\hyperref[h:17:foreign-function--memory-api]{Foreign Function \& Memory API}

\begin{itemize}
\item 
\hyperref[h:17:from-jni-to-ffm-why-it-changed]{From JNI to FFM: Why It Changed}

\item 
\hyperref[h:17:arena-and-deterministic-deallocation]{\code{Arena} and Deterministic Deallocation}

\item 
\hyperref[h:17:memorysegment-memorylayout-valuelayout]{\code{MemorySegment}, \code{MemoryLayout}, \code{ValueLayout}}

\item 
\hyperref[h:17:varhandle-access-into-segments]{\code{VarHandle} Access into Segments}

\item 
\hyperref[h:17:calling-a-c-function-linker-symbollookup-functiondescriptor]{Calling a C Function: \code{Linker}, \code{SymbolLookup}, \code{FunctionDescriptor}}

\item 
\hyperref[h:17:upcalls-java-callback-into-c]{Upcalls: Java Callback into C}

\item 
\hyperref[h:17:struct-layouts]{Struct Layouts}

\item 
\hyperref[h:17:jextract]{\code{jextract}}

\item 
\hyperref[h:17:--enable-native-access-and-restricted-methods]{\code{--\allowbreak{}enable-\allowbreak{}native-\allowbreak{}access} and Restricted Methods}

\item 
\hyperref[h:17:ffm-vs-jni-comparison-table]{FFM vs JNI Comparison Table}

\end{itemize}

\item 
\hyperref[h:17:vector-api-incubator]{Vector API (Incubator)}

\begin{itemize}
\item 
\hyperref[h:17:why-its-still-incubating]{Why It’s Still Incubating}

\item 
\hyperref[h:17:vectorspecies-lanes-and-masks]{\code{VectorSpecies}, Lanes, and Masks}

\item 
\hyperref[h:17:enabling-the-incubator-module]{Enabling the Incubator Module}

\item 
\hyperref[h:17:worked-example-vectorized-dot-product-vs-scalar]{Worked Example: Vectorized Dot Product vs Scalar}

\item 
\hyperref[h:17:the-loop-plus-tail-idiom]{The Loop-Plus-Tail Idiom}

\item 
\hyperref[h:17:dependence-on-project-valhalla]{Dependence on Project Valhalla}

\item 
\hyperref[h:17:production-readiness-caveat]{Production Readiness Caveat}

\end{itemize}

\item 
\hyperref[h:17:common-code-review-interview-pitfalls]{Common Code-Review Interview Pitfalls}

\end{itemize}

\section[Packages]{Packages}
\label{h:17:packages}
A \textbf{package} is a namespace that groups related classes and interfaces together, and it has been part of Java since version 1.0. Packages solve name collisions (two libraries can each have a \code{Logger} class as long as they’re in different packages) and provide a coarse-grained visibility boundary through package-private access.

\subsection[Naming Conventions and Reverse-Domain Rules]{Naming Conventions and Reverse-Domain Rules}
\label{h:17:naming-conventions-and-reverse-domain-rules}
Package names are conventionally all-lowercase and follow the \textbf{reversed internet domain name} of the organization that owns the code, to guarantee global uniqueness without a central registry.

\begin{codebox}
\begin{Verbatim}[commandchars=\\\{\}]
\PY{c+c1}{// Domain: teampicnic.com \PYZhy{}\PYZgt{} reversed \PYZhy{}\PYZgt{} com.teampicnic}
\PY{k+kn}{package}\PY{+w}{ }\PY{n+nn}{com.teampicnic.orders.billing}\PY{p}{;}

\PY{k+kd}{public}\PY{+w}{ }\PY{k+kd}{class} \PY{n+nc}{InvoiceCalculator}\PY{+w}{ }\PY{p}{\PYZob{}}
\PY{+w}{    }\PY{c+c1}{// ...}
\PY{p}{\PYZcb{}}
\end{Verbatim}
\end{codebox}

{\tablefont
\begin{xltabular}{\linewidth}{@{}L{0.925}L{0.857}L{1.218}@{}}
\toprule
\rowcolor{tablehdr}
\thd{Rule} & \thd{Example} & \thd{Reason} \\
\midrule
\endhead
All lowercase & \code{com.\allowbreak{}teampicnic.\allowbreak{}orders}, not \code{com.\allowbreak{}TeamPicnic.\allowbreak{}Orders} & Avoids case-sensitivity clashes on some filesystems and mixed-case ambiguity \\
Reverse domain as prefix & \code{com.\allowbreak{}example.\allowbreak{}app} for \code{example.\allowbreak{}com} & Guarantees uniqueness without a central registry \\
No Java reserved words as segments & Never \code{com.\allowbreak{}example.\allowbreak{}class} & \code{class}, \code{int}, \code{package} etc. are illegal identifiers \\
Avoid underscores/digits at segment start & \code{v2.\allowbreak{}api} is illegal; use \code{apiv2} or \code{api.\allowbreak{}v2} & Java identifier rules apply to each segment \\
Segments describe increasingly specific scope & \code{com.\allowbreak{}teampicnic.\allowbreak{}orders.\allowbreak{}billing.\allowbreak{}tax} & Mirrors directory structure and narrows responsibility \\
\bottomrule
\end{xltabular}
}

Since JDK 11, \code{com.\allowbreak{}sun.\allowbreak{}tools.\allowbreak{}javac} and similar internal packages remind us why the convention matters: the JDK itself uses \code{java.*} and \code{javax.*}, reserved prefixes that application code must never use, to avoid colliding with future JDK classes.

\subsection[The package and import Statements]{The \code{package} and \code{import} Statements}
\label{h:17:the-package-and-import-statements}
The \code{package} declaration must be the first non-comment line in a source file, and a file can declare at most one package.

\begin{codebox}
\begin{Verbatim}[commandchars=\\\{\}]
\PY{k+kn}{package}\PY{+w}{ }\PY{n+nn}{com.teampicnic.orders.billing}\PY{p}{;}

\PY{k+kn}{import}\PY{+w}{ }\PY{n+nn}{java.math.BigDecimal}\PY{p}{;}
\PY{k+kn}{import}\PY{+w}{ }\PY{n+nn}{java.time.LocalDate}\PY{p}{;}
\PY{k+kn}{import}\PY{+w}{ }\PY{n+nn}{com.teampicnic.orders.catalog.Product}\PY{p}{;}

\PY{k+kd}{public}\PY{+w}{ }\PY{k+kd}{class} \PY{n+nc}{Invoice}\PY{+w}{ }\PY{p}{\PYZob{}}
\PY{+w}{    }\PY{k+kd}{private}\PY{+w}{ }\PY{k+kd}{final}\PY{+w}{ }\PY{n}{Product}\PY{+w}{ }\PY{n}{product}\PY{p}{;}
\PY{+w}{    }\PY{k+kd}{private}\PY{+w}{ }\PY{k+kd}{final}\PY{+w}{ }\PY{n}{BigDecimal}\PY{+w}{ }\PY{n}{amount}\PY{p}{;}
\PY{+w}{    }\PY{k+kd}{private}\PY{+w}{ }\PY{k+kd}{final}\PY{+w}{ }\PY{n}{LocalDate}\PY{+w}{ }\PY{n}{issuedOn}\PY{p}{;}

\PY{+w}{    }\PY{k+kd}{public}\PY{+w}{ }\PY{n+nf}{Invoice}\PY{p}{(}\PY{n}{Product}\PY{+w}{ }\PY{n}{product}\PY{p}{,}\PY{+w}{ }\PY{n}{BigDecimal}\PY{+w}{ }\PY{n}{amount}\PY{p}{,}\PY{+w}{ }\PY{n}{LocalDate}\PY{+w}{ }\PY{n}{issuedOn}\PY{p}{)}\PY{+w}{ }\PY{p}{\PYZob{}}
\PY{+w}{        }\PY{k}{this}\PY{p}{.}\PY{n+na}{product}\PY{+w}{ }\PY{o}{=}\PY{+w}{ }\PY{n}{product}\PY{p}{;}
\PY{+w}{        }\PY{k}{this}\PY{p}{.}\PY{n+na}{amount}\PY{+w}{ }\PY{o}{=}\PY{+w}{ }\PY{n}{amount}\PY{p}{;}
\PY{+w}{        }\PY{k}{this}\PY{p}{.}\PY{n+na}{issuedOn}\PY{+w}{ }\PY{o}{=}\PY{+w}{ }\PY{n}{issuedOn}\PY{p}{;}
\PY{+w}{    }\PY{p}{\PYZcb{}}
\PY{p}{\PYZcb{}}
\end{Verbatim}
\end{codebox}

\code{import} is purely a compile-time convenience: it lets you use a simple name (\code{Product}) instead of the fully qualified name (\code{com.\allowbreak{}teampicnic.\allowbreak{}orders.\allowbreak{}catalog.\allowbreak{}Product}) everywhere in the file. It has \textbf{zero runtime cost} — the compiled bytecode always references fully qualified names, and \code{import} statements produce nothing in the \code{.\allowbreak{}class} file.

Two classes are automatically visible without any \code{import}:

\begin{itemize}
\item 
Everything in \code{java.\allowbreak{}lang} (e.g., \code{String}, \code{Object}, \code{Math}) is imported implicitly.

\item 
Classes in the \emph{same package} as the current file need no import at all.

\end{itemize}

\subsection[Static Imports]{Static Imports}
\label{h:17:static-imports}
A \textbf{static import} brings static members (fields or methods) into scope so you can use them unqualified, without the enclosing class name.

\begin{codebox}
\begin{Verbatim}[commandchars=\\\{\}]
\PY{k+kn}{import static}\PY{+w}{ }\PY{n+nn}{java.util.Collections.emptyList}\PY{p}{;}
\PY{k+kn}{import static}\PY{+w}{ }\PY{n+nn}{java.lang.Math.PI}\PY{p}{;}
\PY{k+kn}{import static}\PY{+w}{ }\PY{n+nn}{java.lang.Math.pow}\PY{p}{;}

\PY{k+kd}{public}\PY{+w}{ }\PY{k+kd}{class} \PY{n+nc}{CircleMath}\PY{+w}{ }\PY{p}{\PYZob{}}
\PY{+w}{    }\PY{k+kd}{public}\PY{+w}{ }\PY{k+kd}{static}\PY{+w}{ }\PY{k+kt}{double}\PY{+w}{ }\PY{n+nf}{area}\PY{p}{(}\PY{k+kt}{double}\PY{+w}{ }\PY{n}{radius}\PY{p}{)}\PY{+w}{ }\PY{p}{\PYZob{}}
\PY{+w}{        }\PY{k}{return}\PY{+w}{ }\PY{n}{PI}\PY{+w}{ }\PY{o}{*}\PY{+w}{ }\PY{n}{pow}\PY{p}{(}\PY{n}{radius}\PY{p}{,}\PY{+w}{ }\PY{l+m+mi}{2}\PY{p}{)}\PY{p}{;}\PY{+w}{ }\PY{c+c1}{// no \PYZdq{}Math.\PYZdq{} prefix needed}
\PY{+w}{    }\PY{p}{\PYZcb{}}

\PY{+w}{    }\PY{k+kd}{public}\PY{+w}{ }\PY{k+kd}{static}\PY{+w}{ }\PY{o}{\PYZlt{}}\PY{n}{T}\PY{o}{\PYZgt{}}\PY{+w}{ }\PY{n}{java}\PY{p}{.}\PY{n+na}{util}\PY{p}{.}\PY{n+na}{List}\PY{o}{\PYZlt{}}\PY{n}{T}\PY{o}{\PYZgt{}}\PY{+w}{ }\PY{n+nf}{noResults}\PY{p}{(}\PY{p}{)}\PY{+w}{ }\PY{p}{\PYZob{}}
\PY{+w}{        }\PY{k}{return}\PY{+w}{ }\PY{n}{emptyList}\PY{p}{(}\PY{p}{)}\PY{p}{;}\PY{+w}{ }\PY{c+c1}{// no \PYZdq{}Collections.\PYZdq{} prefix needed}
\PY{+w}{    }\PY{p}{\PYZcb{}}
\PY{p}{\PYZcb{}}
\end{Verbatim}
\end{codebox}

Static imports are genuinely useful for well-known idioms — \code{assertEquals} in JUnit tests, \code{Math} constants, or \code{Collectors} factory methods in stream pipelines — where the qualifying class name adds noise without adding clarity.

\textbf{When they hurt readability:} overusing static imports removes the reader’s ability to tell, at a glance, \emph{where a symbol comes from}. A bare \code{of(\allowbreak{}1,\allowbreak{}~\allowbreak{}2,\allowbreak{}~\allowbreak{}3)} could be \code{List.\allowbreak{}of}, \code{Set.\allowbreak{}of}, \code{Stream.\allowbreak{}of}, or a custom factory — the reader has to scroll up to the imports to find out.

\begin{codebox}
\begin{Verbatim}[commandchars=\\\{\}]
\PY{c+c1}{// Avoid: which \PYZdq{}of\PYZdq{} is this? Which \PYZdq{}process\PYZdq{} is this?}
\PY{k+kn}{import static}\PY{+w}{ }\PY{n+nn}{com.teampicnic.orders.OrderUtils.*}\PY{p}{;}
\PY{k+kn}{import static}\PY{+w}{ }\PY{n+nn}{com.teampicnic.billing.BillingUtils.*}\PY{p}{;}

\PY{k+kd}{public}\PY{+w}{ }\PY{k+kd}{class} \PY{n+nc}{Reconciler}\PY{+w}{ }\PY{p}{\PYZob{}}
\PY{+w}{    }\PY{k+kt}{void}\PY{+w}{ }\PY{n+nf}{run}\PY{p}{(}\PY{p}{)}\PY{+w}{ }\PY{p}{\PYZob{}}
\PY{+w}{        }\PY{k+kd}{var}\PY{+w}{ }\PY{n}{order}\PY{+w}{ }\PY{o}{=}\PY{+w}{ }\PY{n}{of}\PY{p}{(}\PY{l+m+mi}{42}\PY{p}{)}\PY{p}{;}\PY{+w}{       }\PY{c+c1}{// OrderUtils.of or BillingUtils.of? Have to check imports.}
\PY{+w}{        }\PY{n}{process}\PY{p}{(}\PY{n}{order}\PY{p}{)}\PY{p}{;}\PY{+w}{           }\PY{c+c1}{// ambiguous at a glance}
\PY{+w}{    }\PY{p}{\PYZcb{}}
\PY{p}{\PYZcb{}}
\end{Verbatim}
\end{codebox}

Rule of thumb: static-import only widely recognized, unambiguous helpers (\code{Collectors.*}, test assertion libraries, \code{Math.*}), and avoid it for your own domain classes where the qualifying name carries meaning.

\subsection[Wildcard Imports]{Wildcard Imports}
\label{h:17:wildcard-imports}
\code{import~\allowbreak{}com.\allowbreak{}teampicnic.\allowbreak{}orders.*;} imports every top-level type directly in that package (not sub-packages).

\begin{codebox}
\begin{Verbatim}[commandchars=\\\{\}]
\PY{c+c1}{// Wildcard: convenient, but hides exactly what\PYZsq{}s being pulled in}
\PY{k+kn}{import}\PY{+w}{ }\PY{n+nn}{com.teampicnic.orders.*}\PY{p}{;}

\PY{k+kd}{public}\PY{+w}{ }\PY{k+kd}{class} \PY{n+nc}{OrderProcessor}\PY{+w}{ }\PY{p}{\PYZob{}}
\PY{+w}{    }\PY{n}{Order}\PY{+w}{ }\PY{n+nf}{create}\PY{p}{(}\PY{p}{)}\PY{+w}{ }\PY{p}{\PYZob{}}\PY{+w}{ }\PY{k}{return}\PY{+w}{ }\PY{k}{new}\PY{+w}{ }\PY{n}{Order}\PY{p}{(}\PY{p}{)}\PY{p}{;}\PY{+w}{ }\PY{p}{\PYZcb{}}
\PY{p}{\PYZcb{}}
\end{Verbatim}
\end{codebox}

\begin{codebox}
\begin{Verbatim}[commandchars=\\\{\}]
\PY{c+c1}{// Explicit: IDE\PYZhy{}generated, unambiguous, and diffs cleanly in code review}
\PY{k+kn}{import}\PY{+w}{ }\PY{n+nn}{com.teampicnic.orders.Order}\PY{p}{;}

\PY{k+kd}{public}\PY{+w}{ }\PY{k+kd}{class} \PY{n+nc}{OrderProcessor}\PY{+w}{ }\PY{p}{\PYZob{}}
\PY{+w}{    }\PY{n}{Order}\PY{+w}{ }\PY{n+nf}{create}\PY{p}{(}\PY{p}{)}\PY{+w}{ }\PY{p}{\PYZob{}}\PY{+w}{ }\PY{k}{return}\PY{+w}{ }\PY{k}{new}\PY{+w}{ }\PY{n}{Order}\PY{p}{(}\PY{p}{)}\PY{p}{;}\PY{+w}{ }\PY{p}{\PYZcb{}}
\PY{p}{\PYZcb{}}
\end{Verbatim}
\end{codebox}

{\tablefont
\begin{xltabular}{\linewidth}{@{}L{0.444}L{1.791}L{0.764}@{}}
\toprule
\rowcolor{tablehdr}
\thd{} & \thd{Wildcard \code{import~\allowbreak{}pkg.*;}} & \thd{Explicit \code{import~\allowbreak{}pkg.\allowbreak{}Type;}} \\
\midrule
\endhead
Compile-time cost & Identical — resolved once at compile time either way & Identical \\
Readability & Reader must infer which types are actually used & Every used type is listed \\
Merge conflicts & Fewer lines, but ambiguous diffs & More lines, but precise diffs \\
Name collisions & Can silently break when the package adds a new type with a name you already use & Immune — each import is explicit \\
Common tooling default & Most style guides (Google Java Style, and most IDEs by default) disallow it & Preferred by nearly every style guide \\
\bottomrule
\end{xltabular}
}

Wildcard imports are almost universally discouraged in professional style guides and in code review, precisely because they trade a small amount of typing for a real risk: if the imported package later adds a class with the same simple name as something else you’re using, your code can silently start resolving to the wrong type, or fail to compile with an ambiguity error far from where the real problem is.

\subsection[The Default (Unnamed) Package]{The Default (Unnamed) Package}
\label{h:17:the-default-unnamed-package}
A source file with no \code{package} statement lives in the \textbf{default (unnamed) package}.

\begin{codebox}
\begin{Verbatim}[commandchars=\\\{\}]
\PY{c+c1}{// No package declaration at all — lives in the default package}
\PY{k+kd}{public}\PY{+w}{ }\PY{k+kd}{class} \PY{n+nc}{QuickTest}\PY{+w}{ }\PY{p}{\PYZob{}}
\PY{+w}{    }\PY{k+kd}{public}\PY{+w}{ }\PY{k+kd}{static}\PY{+w}{ }\PY{k+kt}{void}\PY{+w}{ }\PY{n+nf}{main}\PY{p}{(}\PY{n}{String}\PY{o}{[}\PY{o}{]}\PY{+w}{ }\PY{n}{args}\PY{p}{)}\PY{+w}{ }\PY{p}{\PYZob{}}
\PY{+w}{        }\PY{n}{System}\PY{p}{.}\PY{n+na}{out}\PY{p}{.}\PY{n+na}{println}\PY{p}{(}\PY{l+s}{\PYZdq{}}\PY{l+s}{Hello from the default package}\PY{l+s}{\PYZdq{}}\PY{p}{)}\PY{p}{;}
\PY{+w}{    }\PY{p}{\PYZcb{}}
\PY{p}{\PYZcb{}}
\end{Verbatim}
\end{codebox}

\textbf{Why to never use it in real code:}

\begin{enumerate}
\item 
\textbf{It cannot be imported.} Classes in the default package are invisible to any class that belongs to a named package, since \code{import} requires a package-qualified name.

\item 
\textbf{It cannot participate in JPMS.} A module cannot export the default package, so any type there is entirely inaccessible from a modular application.

\item 
\textbf{Name collisions are almost guaranteed} as soon as the project grows — two unrelated classes both wanting to be called \code{Utils} will collide immediately with no namespace to separate them.

\item 
\textbf{Build tools and JavaDoc treat it as second class} — many tools assume every meaningful class has a package and behave oddly (or refuse to process the class at all) otherwise.

\end{enumerate}

The default package is acceptable for a one-off scratch file (\code{Test.\allowbreak{}java} you compile and delete) and nothing else. Every class that will be checked in should declare a package.

\subsection[Package-Private Visibility as a Design Tool]{Package-Private Visibility as a Design Tool}
\label{h:17:package-private-visibility-as-a-design-tool}
When no access modifier is written, a member or top-level class has \textbf{package-private} (also called \emph{default}) visibility: accessible only to code in the exact same package.

\begin{codebox}
\begin{Verbatim}[commandchars=\\\{\}]
\PY{k+kn}{package}\PY{+w}{ }\PY{n+nn}{com.teampicnic.orders.billing}\PY{p}{;}

\PY{k+kd}{public}\PY{+w}{ }\PY{k+kd}{class} \PY{n+nc}{InvoiceService}\PY{+w}{ }\PY{p}{\PYZob{}}\PY{+w}{          }\PY{c+c1}{// public API entry point}
\PY{+w}{    }\PY{k+kd}{private}\PY{+w}{ }\PY{k+kd}{final}\PY{+w}{ }\PY{n}{TaxCalculator}\PY{+w}{ }\PY{n}{tax}\PY{+w}{ }\PY{o}{=}\PY{+w}{ }\PY{k}{new}\PY{+w}{ }\PY{n}{TaxCalculator}\PY{p}{(}\PY{p}{)}\PY{p}{;}

\PY{+w}{    }\PY{k+kd}{public}\PY{+w}{ }\PY{n}{java}\PY{p}{.}\PY{n+na}{math}\PY{p}{.}\PY{n+na}{BigDecimal}\PY{+w}{ }\PY{n+nf}{totalFor}\PY{p}{(}\PY{n}{Invoice}\PY{+w}{ }\PY{n}{invoice}\PY{p}{)}\PY{+w}{ }\PY{p}{\PYZob{}}
\PY{+w}{        }\PY{k}{return}\PY{+w}{ }\PY{n}{tax}\PY{p}{.}\PY{n+na}{calculate}\PY{p}{(}\PY{n}{invoice}\PY{p}{)}\PY{p}{;}\PY{+w}{ }\PY{c+c1}{// package\PYZhy{}private class used internally}
\PY{+w}{    }\PY{p}{\PYZcb{}}
\PY{p}{\PYZcb{}}

\PY{k+kd}{class} \PY{n+nc}{TaxCalculator}\PY{+w}{ }\PY{p}{\PYZob{}}\PY{+w}{                  }\PY{c+c1}{// package\PYZhy{}private: not part of the public API}
\PY{+w}{    }\PY{n}{java}\PY{p}{.}\PY{n+na}{math}\PY{p}{.}\PY{n+na}{BigDecimal}\PY{+w}{ }\PY{n+nf}{calculate}\PY{p}{(}\PY{n}{Invoice}\PY{+w}{ }\PY{n}{invoice}\PY{p}{)}\PY{+w}{ }\PY{p}{\PYZob{}}
\PY{+w}{        }\PY{k}{return}\PY{+w}{ }\PY{n}{invoice}\PY{p}{.}\PY{n+na}{amount}\PY{p}{(}\PY{p}{)}\PY{p}{.}\PY{n+na}{multiply}\PY{p}{(}\PY{n}{java}\PY{p}{.}\PY{n+na}{math}\PY{p}{.}\PY{n+na}{BigDecimal}\PY{p}{.}\PY{n+na}{valueOf}\PY{p}{(}\PY{l+m+mf}{1.2}\PY{p}{)}\PY{p}{)}\PY{p}{;}
\PY{+w}{    }\PY{p}{\PYZcb{}}
\PY{p}{\PYZcb{}}
\end{Verbatim}
\end{codebox}

This is a deliberate design tool, not just “the case where you forgot \code{public}.” Keeping helper classes package-private lets you refactor their internals freely — rename methods, change constructors, delete the class entirely — without breaking any code outside the package, because the compiler guarantees no outside code could have depended on it. This is the same idea JPMS later formalizes at a coarser grain with \code{exports}: expose the minimum surface area needed, and hide the rest.

A common pattern: one \code{public} class or interface per package acting as the facade, and everything else package-private, wired together only within that package.

\subsection[package-info.java]{\code{package-\allowbreak{}info.\allowbreak{}java}}
\label{h:17:package-infojava}
\code{package-\allowbreak{}info.\allowbreak{}java} is a special file, one per package, that holds package-level Javadoc and package-level annotations. It contains no class body — just an optional Javadoc comment, an optional set of annotations, and the \code{package} statement.

\begin{codebox}
\begin{Verbatim}[commandchars=\\\{\}]
\PY{c+cm}{/**}
\PY{c+cm}{ * Billing calculations for customer invoices, including tax and currency}
\PY{c+cm}{ * conversion rules. All monetary values use \PYZob{}@link java.math.BigDecimal\PYZcb{}}
\PY{c+cm}{ * with \PYZob{}@code RoundingMode.HALF\PYZus{}UP\PYZcb{} unless documented otherwise.}
\PY{c+cm}{ *}
\PY{c+cm}{ * @since 2.3}
\PY{c+cm}{ */}
\PY{n+nd}{@NonNullApi}\PY{+w}{ }\PY{c+c1}{// hypothetical package\PYZhy{}wide annotation, e.g. from Spring or a custom checker}
\PY{k+kn}{package}\PY{+w}{ }\PY{n+nn}{com.teampicnic.orders.billing}\PY{p}{;}

\PY{k+kn}{import}\PY{+w}{ }\PY{n+nn}{com.teampicnic.orders.NonNullApi}\PY{p}{;}
\end{Verbatim}
\end{codebox}

Two practical uses reviewers should recognize:

\begin{itemize}
\item 
\textbf{Documentation} — Javadoc generated for the package overview page comes from here, not from any regular class’s comment.

\item 
\textbf{Package-wide annotations} — a \code{@\allowbreak{}NonNullApi}-style annotation (like Spring’s \code{org.\allowbreak{}springframework.\allowbreak{}lang.\allowbreak{}Non\allowbreak{}Null\allowbreak{}Api}) applied here means “every parameter and return type in this package is non-null unless explicitly marked \code{@\allowbreak{}Nullable},” saving you from annotating every single method individually.

\end{itemize}

\subsection[Split Packages]{Split Packages}
\label{h:17:split-packages}
A \textbf{split package} occurs when the \emph{same} package name is provided by more than one JAR (or module) on the classpath/module path. On the plain classpath this “just works” in a fragile way — classes get merged from whichever JAR the classloader finds first, in an order that’s often undocumered and build-tool-dependent.

\begin{codebox}
\begin{Verbatim}
lib-a.jar
  com/example/util/StringHelper.class

lib-b.jar
  com/example/util/StringHelper.class   <- same package AND same class name, different code!
\end{Verbatim}
\end{codebox}

On the classpath, whichever JAR appears first wins silently — no error, no warning, just whichever version happened to load, which is exactly the kind of bug that’s nearly impossible to diagnose from a stack trace alone.

\textbf{JPMS makes split packages a hard error.} Two named modules are not permitted to export the same package; the module system refuses to resolve at launch:

\begin{codebox}
\begin{Verbatim}
Error occurred during initialization of boot layer
java.lang.module.ResolutionException: Module lib.a and module lib.b export
package com.example.util to module myapp
\end{Verbatim}
\end{codebox}

This is a deliberate trade-off: JPMS trades classpath flexibility for a hard compile/launch-time guarantee that a package’s contents are unambiguous. It is one of the most common real-world blockers when modularizing an existing multi-JAR project — you often discover, for the first time, that two of your dependencies quietly ship overlapping packages.

\subsection[Directory-to-Package Mapping]{Directory-to-Package Mapping}
\label{h:17:directory-to-package-mapping}
The Java compiler and JVM require that a class’s package name \textbf{exactly mirror its directory path} relative to a source or class root.

\begin{codebox}
\begin{Verbatim}
src/
└── main/
    └── java/
        └── com/
            └── teampicnic/
                └── orders/
                    └── billing/
                        ├── Invoice.java          -> package com.teampicnic.orders.billing;
                        └── InvoiceService.java   -> package com.teampicnic.orders.billing;
\end{Verbatim}
\end{codebox}

\begin{codebox}
\begin{Verbatim}[commandchars=\\\{\}]
\PY{c+c1}{\PYZsh{} Compiling from the source root, mirroring the package structure, produces}
\PY{c+c1}{\PYZsh{} matching directories under the output root:}
javac\PY{+w}{ }\PYZhy{}d\PY{+w}{ }out\PY{+w}{ }\PY{k}{\PYZdl{}(}find\PY{+w}{ }src/main/java\PY{+w}{ }\PYZhy{}name\PY{+w}{ }\PY{l+s+s2}{\PYZdq{}*.java\PYZdq{}}\PY{k}{)}

\PY{c+c1}{\PYZsh{} out/com/teampicnic/orders/billing/Invoice.class}
\PY{c+c1}{\PYZsh{} out/com/teampicnic/orders/billing/InvoiceService.class}
\end{Verbatim}
\end{codebox}

This mapping is not a convention the compiler merely prefers — it is enforced. If \code{Invoice.\allowbreak{}java} inside a directory called \code{billing/} declared \code{package~\allowbreak{}com.\allowbreak{}teampicnic.\allowbreak{}orders.\allowbreak{}tax;}, \code{javac} would fail to compile it (or, on the classpath at runtime, the class simply couldn’t be found where the classloader expects it).

\subsection[Packages vs Modules]{Packages vs Modules}
\label{h:17:packages-vs-modules}
Packages and modules solve related but distinct problems, and mixing them up is a common interview stumbling point.

{\tablefontsmall
\begin{xltabular}{\linewidth}{@{}L{0.629}L{1.112}L{1.259}@{}}
\toprule
\rowcolor{tablehdr}
\thd{} & \thd{Package} & \thd{Module} \\
\midrule
\endhead
Introduced & Java 1.0 & Java 9 \\
Unit of & Namespace + package-private visibility & Deployment, strong encapsulation, and dependency declaration \\
Declared in & \code{package} statement at top of each \code{.\allowbreak{}java} file & Single \code{module-\allowbreak{}info.\allowbreak{}java} per module (a JAR-level concern) \\
Visibility granularity & Fine-grained (per-class, per-member) & Coarse-grained (per-package, via \code{exports}/\code{opens}) \\
Can exist without the other? & Yes — packages predate and work fully without modules & A module is \emph{composed of} one or more packages; it can’t exist without them \\
Enforces uniqueness across dependencies? & No (split packages silently merge on the classpath) & Yes (split packages are a hard error) \\
Reflection access to internals & Always allowed (subject to \code{SecurityManager}, now removed) & Blocked unless the package is \code{opens}-ed \\
\bottomrule
\end{xltabular}
}

A module \emph{contains} packages; it never replaces them. You still write \code{package~\allowbreak{}com.\allowbreak{}teampicnic.\allowbreak{}orders.\allowbreak{}billing;} at the top of every file exactly as before — JPMS adds one more file, \code{module-\allowbreak{}info.\allowbreak{}java}, that declares which of those packages the module exposes to the outside world.

\section[Java Platform Module System (JPMS)]{Java Platform Module System (JPMS)}
\label{h:17:java-platform-module-system-jpms}
JPMS (JSR 376, delivered in Java 9 via Project Jigsaw) introduced modules: named, self-describing units that declare their dependencies and their public surface explicitly, instead of relying on an implicit, all-or-nothing classpath.

\subsection[Problems JPMS Solves]{Problems JPMS Solves}
\label{h:17:problems-jpms-solves}
{\tablefontsmall
\begin{xltabular}{\linewidth}{@{}L{0.477}L{1.262}L{1.262}@{}}
\toprule
\rowcolor{tablehdr}
\thd{Problem} & \thd{Pre-JPMS reality} & \thd{JPMS fix} \\
\midrule
\endhead
\textbf{Classpath hell} & Every JAR’s classes are dumped into one flat namespace; duplicate/conflicting versions silently shadow each other; no dependency graph is declared anywhere & Modules declare explicit \code{requires}; the module system builds and validates a real dependency graph at launch, failing fast on missing or conflicting modules \\
\textbf{No strong encapsulation} & \code{public} meant public to \emph{everyone} on the classpath, forever, even for classes meant as “internal use only” (\code{sun.\allowbreak{}misc.\allowbreak{}Unsafe}, \code{com.\allowbreak{}sun.*}) & \code{exports} controls which packages are visible outside the module at compile time \emph{and} runtime; unexported packages are invisible, full stop \\
\textbf{Monolithic JDK} & The entire JDK shipped as one giant \code{rt.\allowbreak{}jar}, so even a “Hello World” needed the full runtime, and internal JDK classes were reachable via reflection by anyone & The JDK itself is split into \textasciitilde{}70 modules (\code{java.\allowbreak{}base}, \code{java.\allowbreak{}sql}, \code{java.\allowbreak{}desktop}, …); \code{jlink} builds a runtime image with only the modules you actually need \\
\textbf{Reflection reaching into internals} & Any code could call \code{setAccessible(\allowbreak{}true)} and reach into private JDK fields/methods with no declared contract & \code{opens} is required for deep reflection; without it, \code{setAccessible} throws \code{Inaccessible\allowbreak{}Object\allowbreak{}Exception} \\
\bottomrule
\end{xltabular}
}

\subsection[module-info.java Basics]{\code{module-\allowbreak{}info.\allowbreak{}java} Basics}
\label{h:17:module-infojava-basics}
Every module has exactly one \code{module-\allowbreak{}info.\allowbreak{}java} at the root of its source tree (not inside any package — it sits directly under the module’s source root, alongside the top-level package directories).

\begin{codebox}
\begin{Verbatim}[commandchars=\\\{\}]
\PY{n}{module}\PY{+w}{ }\PY{n}{com}\PY{p}{.}\PY{n+na}{teampicnic}\PY{p}{.}\PY{n+na}{billing}\PY{+w}{ }\PY{p}{\PYZob{}}
\PY{+w}{    }\PY{n}{requires}\PY{+w}{ }\PY{n}{java}\PY{p}{.}\PY{n+na}{sql}\PY{p}{;}
\PY{+w}{    }\PY{n}{exports}\PY{+w}{ }\PY{n}{com}\PY{p}{.}\PY{n+na}{teampicnic}\PY{p}{.}\PY{n+na}{billing}\PY{p}{.}\PY{n+na}{api}\PY{p}{;}
\PY{p}{\PYZcb{}}
\end{Verbatim}
\end{codebox}

\begin{itemize}
\item 
The module name (\code{com.\allowbreak{}teampicnic.\allowbreak{}billing}) conventionally follows the same reverse-domain convention as packages, but it is a \textbf{separate namespace} from package names — a module is not required to contain a package of the exact same name, though it’s a common and readable convention to have one “root” package matching the module name.

\item 
Everything inside the module’s packages is invisible outside the module \textbf{unless} explicitly \code{exports}-ed.

\end{itemize}

\subsection[requires, requires transitive, requires static]{\code{requires}, \code{requires~\allowbreak{}transitive}, \code{requires~\allowbreak{}static}}
\label{h:17:requires-requires-transitive-requires-static}
\code{requires} declares a compile-time and run-time dependency on another module.

\begin{codebox}
\begin{Verbatim}[commandchars=\\\{\}]
\PY{n}{module}\PY{+w}{ }\PY{n}{com}\PY{p}{.}\PY{n+na}{teampicnic}\PY{p}{.}\PY{n+na}{billing}\PY{+w}{ }\PY{p}{\PYZob{}}
\PY{+w}{    }\PY{n}{requires}\PY{+w}{ }\PY{n}{java}\PY{p}{.}\PY{n+na}{sql}\PY{p}{;}\PY{+w}{                     }\PY{c+c1}{// plain dependency}
\PY{+w}{    }\PY{n}{requires}\PY{+w}{ }\PY{n}{transitive}\PY{+w}{ }\PY{n}{com}\PY{p}{.}\PY{n+na}{teampicnic}\PY{p}{.}\PY{n+na}{core}\PY{p}{;}\PY{+w}{ }\PY{c+c1}{// re\PYZhy{}exported to consumers of this module}
\PY{+w}{    }\PY{n}{requires}\PY{+w}{ }\PY{k+kd}{static}\PY{+w}{ }\PY{n}{com}\PY{p}{.}\PY{n+na}{teampicnic}\PY{p}{.}\PY{n+na}{testkit}\PY{p}{;}\PY{+w}{  }\PY{c+c1}{// compile\PYZhy{}time only, optional at runtime}
\PY{p}{\PYZcb{}}
\end{Verbatim}
\end{codebox}

{\tablefontsmall
\begin{xltabular}{\linewidth}{@{}L{0.432}L{1.284}L{1.284}@{}}
\toprule
\rowcolor{tablehdr}
\thd{Form} & \thd{Meaning} & \thd{Typical use} \\
\midrule
\endhead
\code{requires~\allowbreak{}X;} & This module needs \code{X} to compile and run. & The normal case for a real dependency \\
\code{requires~\allowbreak{}transitive~\allowbreak{}X;} & This module needs \code{X}, \textbf{and} any module that \code{requires~\allowbreak{}com.\allowbreak{}teampicnic.\allowbreak{}billing} automatically also gets readability to \code{X} without declaring it itself. & Your module’s public API exposes types from \code{X} (e.g., a method returns \code{java.\allowbreak{}sql.\allowbreak{}Connection}), so callers need \code{X} too \\
\code{requires~\allowbreak{}static~\allowbreak{}X;} & This module needs \code{X} only to \textbf{compile}; it is optional at runtime — if \code{X} is absent when running, the module still loads, but any code path that actually touches \code{X}’s types fails at that point. & Compile-time-only annotations (like \code{org.\allowbreak{}jetbrains.\allowbreak{}annotations.\allowbreak{}Nullable}) that aren’t needed once compiled \\
\bottomrule
\end{xltabular}
}

\code{requires~\allowbreak{}transitive} is the module-system equivalent of a Maven “compile-scope, exposed” dependency: if \code{Invoice.\allowbreak{}get\allowbreak{}Connection()} returns a \code{java.\allowbreak{}sql.\allowbreak{}Connection}, then any caller needs \code{java.\allowbreak{}sql} on their module path too — \code{requires~\allowbreak{}transitive~\allowbreak{}java.\allowbreak{}sql;} in your module spares every consumer from having to add that \code{requires} themselves.

\subsection[exports and exports ... to]{\code{exports} and \code{exports~\allowbreak{}...\allowbreak{}~\allowbreak{}to}}
\label{h:17:exports-and-exports--to}
\code{exports} makes a package’s \code{public} types accessible (at compile time and runtime) to any module that \code{requires} this module.

\begin{codebox}
\begin{Verbatim}[commandchars=\\\{\}]
\PY{n}{module}\PY{+w}{ }\PY{n}{com}\PY{p}{.}\PY{n+na}{teampicnic}\PY{p}{.}\PY{n+na}{billing}\PY{+w}{ }\PY{p}{\PYZob{}}
\PY{+w}{    }\PY{n}{exports}\PY{+w}{ }\PY{n}{com}\PY{p}{.}\PY{n+na}{teampicnic}\PY{p}{.}\PY{n+na}{billing}\PY{p}{.}\PY{n+na}{api}\PY{p}{;}\PY{+w}{                     }\PY{c+c1}{// visible to everyone}
\PY{+w}{    }\PY{n}{exports}\PY{+w}{ }\PY{n}{com}\PY{p}{.}\PY{n+na}{teampicnic}\PY{p}{.}\PY{n+na}{billing}\PY{p}{.}\PY{n+na}{internal}\PY{+w}{ }\PY{n}{to}\PY{+w}{ }\PY{n}{com}\PY{p}{.}\PY{n+na}{teampicnic}\PY{p}{.}\PY{n+na}{billing}\PY{p}{.}\PY{n+na}{tests}\PY{p}{;}\PY{+w}{ }\PY{c+c1}{// visible only to named module(s)}
\PY{p}{\PYZcb{}}
\end{Verbatim}
\end{codebox}

\begin{itemize}
\item 
Plain \code{exports~\allowbreak{}pkg;} — anyone who \code{requires} your module can see \code{pkg}’s public types.

\item 
\textbf{Qualified export} \code{exports~\allowbreak{}pkg~\allowbreak{}to~\allowbreak{}module\allowbreak{}A,\allowbreak{}~\allowbreak{}module\allowbreak{}B;} — only the named modules can see it. This is how you can share internal-but-not-fully-public code with, say, your own test module or a sibling module, without making it part of your general public API.

\end{itemize}

Packages that are \emph{not} exported at all are invisible outside the module even though their classes may be \code{public} — this is the core of JPMS’s strong encapsulation, and it is enforced by the compiler and the runtime, not just convention.

\subsection[opens and opens ... to]{\code{opens} and \code{opens~\allowbreak{}...\allowbreak{}~\allowbreak{}to}}
\label{h:17:opens-and-opens--to}
\code{opens} is about \textbf{reflection}, not compile-time linking. A package can be exported (visible, linkable, callable normally) without being opened, and vice versa.

\begin{codebox}
\begin{Verbatim}[commandchars=\\\{\}]
\PY{n}{module}\PY{+w}{ }\PY{n}{com}\PY{p}{.}\PY{n+na}{teampicnic}\PY{p}{.}\PY{n+na}{billing}\PY{+w}{ }\PY{p}{\PYZob{}}
\PY{+w}{    }\PY{n}{exports}\PY{+w}{ }\PY{n}{com}\PY{p}{.}\PY{n+na}{teampicnic}\PY{p}{.}\PY{n+na}{billing}\PY{p}{.}\PY{n+na}{api}\PY{p}{;}
\PY{+w}{    }\PY{n}{opens}\PY{+w}{ }\PY{n}{com}\PY{p}{.}\PY{n+na}{teampicnic}\PY{p}{.}\PY{n+na}{billing}\PY{p}{.}\PY{n+na}{model}\PY{p}{;}\PY{+w}{                 }\PY{c+c1}{// deep reflection allowed for everyone}
\PY{+w}{    }\PY{n}{opens}\PY{+w}{ }\PY{n}{com}\PY{p}{.}\PY{n+na}{teampicnic}\PY{p}{.}\PY{n+na}{billing}\PY{p}{.}\PY{n+na}{internal}\PY{+w}{ }\PY{n}{to}\PY{+w}{ }\PY{n}{com}\PY{p}{.}\PY{n+na}{fasterxml}\PY{p}{.}\PY{n+na}{jackson}\PY{p}{.}\PY{n+na}{databind}\PY{p}{;}\PY{+w}{ }\PY{c+c1}{// reflection allowed only for Jackson}
\PY{p}{\PYZcb{}}
\end{Verbatim}
\end{codebox}

{\tablefontsmall
\begin{xltabular}{\linewidth}{@{}L{1.288}L{0.425}L{1.288}@{}}
\toprule
\rowcolor{tablehdr}
\thd{} & \thd{\code{exports}} & \thd{\code{opens}} \\
\midrule
\endhead
Grants normal compile-time + runtime access (calling public methods, \code{new}ing public classes) & Yes & No \\
Grants deep reflection (\code{setAccessible(\allowbreak{}true)} on private members) & No & Yes \\
Typical need & Library public API & JSON/serialization frameworks (Jackson, Gson), dependency injection (Spring, Hibernate) that reflectively construct and populate your classes, including private fields \\
\bottomrule
\end{xltabular}
}

Frameworks like Hibernate and Jackson need \code{opens} (not just \code{exports}) on your model/entity packages, because they use reflection to set private fields directly. A module can also be declared \code{open} entirely — \code{open~\allowbreak{}module~\allowbreak{}com.\allowbreak{}teampicnic.\allowbreak{}billing~\allowbreak{}\{\allowbreak{}~\allowbreak{}requires~\allowbreak{}java.\allowbreak{}sql;\allowbreak{}~\allowbreak{}exports~\allowbreak{}com.\allowbreak{}teampicnic.\allowbreak{}billing.\allowbreak{}api;\allowbreak{}~\allowbreak{}\}} — which opens \emph{every} package in the module to reflection at once, a pragmatic shortcut when a module is mostly framework-managed POJOs and enumerating each package individually isn’t worth the ceremony.

\subsection[uses and provides ... with]{\code{uses} and \code{provides~\allowbreak{}...\allowbreak{}~\allowbreak{}with}}
\label{h:17:uses-and-provides--with}
These two directives implement the \textbf{\code{ServiceLoader}} pattern at the module level — a consumer declares it \emph{uses} a service interface, and a provider module declares it \emph{provides} an implementation, without either module needing to know about the other directly.

\begin{codebox}
\begin{Verbatim}[commandchars=\\\{\}]
\PY{c+c1}{// Module: com.teampicnic.billing.api  (declares the service interface)}
\PY{n}{module}\PY{+w}{ }\PY{n}{com}\PY{p}{.}\PY{n+na}{teampicnic}\PY{p}{.}\PY{n+na}{billing}\PY{p}{.}\PY{n+na}{api}\PY{+w}{ }\PY{p}{\PYZob{}}
\PY{+w}{    }\PY{n}{exports}\PY{+w}{ }\PY{n}{com}\PY{p}{.}\PY{n+na}{teampicnic}\PY{p}{.}\PY{n+na}{billing}\PY{p}{.}\PY{n+na}{api}\PY{p}{;}
\PY{p}{\PYZcb{}}
\end{Verbatim}
\end{codebox}

\begin{codebox}
\begin{Verbatim}[commandchars=\\\{\}]
\PY{k+kn}{package}\PY{+w}{ }\PY{n+nn}{com.teampicnic.billing.api}\PY{p}{;}

\PY{k+kd}{public}\PY{+w}{ }\PY{k+kd}{interface} \PY{n+nc}{TaxRuleProvider}\PY{+w}{ }\PY{p}{\PYZob{}}
\PY{+w}{    }\PY{n}{java}\PY{p}{.}\PY{n+na}{math}\PY{p}{.}\PY{n+na}{BigDecimal}\PY{+w}{ }\PY{n+nf}{rateFor}\PY{p}{(}\PY{n}{String}\PY{+w}{ }\PY{n}{countryCode}\PY{p}{)}\PY{p}{;}
\PY{p}{\PYZcb{}}
\end{Verbatim}
\end{codebox}

\begin{codebox}
\begin{Verbatim}[commandchars=\\\{\}]
\PY{c+c1}{// Module: com.teampicnic.billing.eu  (provides an implementation)}
\PY{n}{module}\PY{+w}{ }\PY{n}{com}\PY{p}{.}\PY{n+na}{teampicnic}\PY{p}{.}\PY{n+na}{billing}\PY{p}{.}\PY{n+na}{eu}\PY{+w}{ }\PY{p}{\PYZob{}}
\PY{+w}{    }\PY{n}{requires}\PY{+w}{ }\PY{n}{com}\PY{p}{.}\PY{n+na}{teampicnic}\PY{p}{.}\PY{n+na}{billing}\PY{p}{.}\PY{n+na}{api}\PY{p}{;}
\PY{+w}{    }\PY{n}{provides}\PY{+w}{ }\PY{n}{com}\PY{p}{.}\PY{n+na}{teampicnic}\PY{p}{.}\PY{n+na}{billing}\PY{p}{.}\PY{n+na}{api}\PY{p}{.}\PY{n+na}{TaxRuleProvider}
\PY{+w}{        }\PY{n}{with}\PY{+w}{ }\PY{n}{com}\PY{p}{.}\PY{n+na}{teampicnic}\PY{p}{.}\PY{n+na}{billing}\PY{p}{.}\PY{n+na}{eu}\PY{p}{.}\PY{n+na}{EuTaxRuleProvider}\PY{p}{;}
\PY{p}{\PYZcb{}}
\end{Verbatim}
\end{codebox}

\begin{codebox}
\begin{Verbatim}[commandchars=\\\{\}]
\PY{c+c1}{// Module: com.teampicnic.billing.engine  (consumes the service)}
\PY{n}{module}\PY{+w}{ }\PY{n}{com}\PY{p}{.}\PY{n+na}{teampicnic}\PY{p}{.}\PY{n+na}{billing}\PY{p}{.}\PY{n+na}{engine}\PY{+w}{ }\PY{p}{\PYZob{}}
\PY{+w}{    }\PY{n}{requires}\PY{+w}{ }\PY{n}{com}\PY{p}{.}\PY{n+na}{teampicnic}\PY{p}{.}\PY{n+na}{billing}\PY{p}{.}\PY{n+na}{api}\PY{p}{;}
\PY{+w}{    }\PY{n}{uses}\PY{+w}{ }\PY{n}{com}\PY{p}{.}\PY{n+na}{teampicnic}\PY{p}{.}\PY{n+na}{billing}\PY{p}{.}\PY{n+na}{api}\PY{p}{.}\PY{n+na}{TaxRuleProvider}\PY{p}{;}
\PY{p}{\PYZcb{}}
\end{Verbatim}
\end{codebox}

\begin{codebox}
\begin{Verbatim}[commandchars=\\\{\}]
\PY{k+kn}{package}\PY{+w}{ }\PY{n+nn}{com.teampicnic.billing.engine}\PY{p}{;}

\PY{k+kn}{import}\PY{+w}{ }\PY{n+nn}{com.teampicnic.billing.api.TaxRuleProvider}\PY{p}{;}
\PY{k+kn}{import}\PY{+w}{ }\PY{n+nn}{java.util.ServiceLoader}\PY{p}{;}

\PY{k+kd}{public}\PY{+w}{ }\PY{k+kd}{class} \PY{n+nc}{TaxEngine}\PY{+w}{ }\PY{p}{\PYZob{}}
\PY{+w}{    }\PY{k+kd}{public}\PY{+w}{ }\PY{n}{java}\PY{p}{.}\PY{n+na}{math}\PY{p}{.}\PY{n+na}{BigDecimal}\PY{+w}{ }\PY{n+nf}{calculate}\PY{p}{(}\PY{n}{String}\PY{+w}{ }\PY{n}{countryCode}\PY{p}{)}\PY{+w}{ }\PY{p}{\PYZob{}}
\PY{+w}{        }\PY{n}{ServiceLoader}\PY{o}{\PYZlt{}}\PY{n}{TaxRuleProvider}\PY{o}{\PYZgt{}}\PY{+w}{ }\PY{n}{loader}\PY{+w}{ }\PY{o}{=}\PY{+w}{ }\PY{n}{ServiceLoader}\PY{p}{.}\PY{n+na}{load}\PY{p}{(}\PY{n}{TaxRuleProvider}\PY{p}{.}\PY{n+na}{class}\PY{p}{)}\PY{p}{;}
\PY{+w}{        }\PY{k}{for}\PY{+w}{ }\PY{p}{(}\PY{n}{TaxRuleProvider}\PY{+w}{ }\PY{n}{provider}\PY{+w}{ }\PY{p}{:}\PY{+w}{ }\PY{n}{loader}\PY{p}{)}\PY{+w}{ }\PY{p}{\PYZob{}}
\PY{+w}{            }\PY{k}{return}\PY{+w}{ }\PY{n}{provider}\PY{p}{.}\PY{n+na}{rateFor}\PY{p}{(}\PY{n}{countryCode}\PY{p}{)}\PY{p}{;}
\PY{+w}{        }\PY{p}{\PYZcb{}}
\PY{+w}{        }\PY{k}{throw}\PY{+w}{ }\PY{k}{new}\PY{+w}{ }\PY{n}{IllegalStateException}\PY{p}{(}\PY{l+s}{\PYZdq{}}\PY{l+s}{No TaxRuleProvider found on the module path}\PY{l+s}{\PYZdq{}}\PY{p}{)}\PY{p}{;}
\PY{+w}{    }\PY{p}{\PYZcb{}}
\PY{p}{\PYZcb{}}
\end{Verbatim}
\end{codebox}

\code{uses} and \code{provides} predate JPMS as an idea (\code{ServiceLoader} and \code{META-\allowbreak{}INF/\allowbreak{}services/} files existed since Java 6), but JPMS makes the wiring declarative and verifiable at compile time: the module system checks that a \code{provides} implementation actually implements the declared interface and that consumers legitimately declare \code{uses} for what they load, instead of relying on a text file in \code{META-\allowbreak{}INF/\allowbreak{}services/} that nothing validates.

\subsection[Automatic Modules and the Unnamed Module]{Automatic Modules and the Unnamed Module}
\label{h:17:automatic-modules-and-the-unnamed-module}
Not every JAR on your module path has a \code{module-\allowbreak{}info.\allowbreak{}class}. JPMS handles this with two fallback mechanisms:

{\tablefontsmall
\begin{xltabular}{\linewidth}{@{}L{0.422}L{1.623}L{0.956}@{}}
\toprule
\rowcolor{tablehdr}
\thd{} & \thd{Automatic module} & \thd{Unnamed module} \\
\midrule
\endhead
What it is & A plain JAR (no \code{module-\allowbreak{}info.\allowbreak{}class}) placed on the \textbf{module path} & Anything placed on the plain \textbf{classpath} when running a modular application \\
Module name & Derived from the JAR filename (e.g., \code{guava-\allowbreak{}33.\allowbreak{}0.\allowbreak{}0.\allowbreak{}jar} -\textgreater{} \code{guava}), or from \code{Automatic-\allowbreak{}Module-\allowbreak{}Name} in \code{META-\allowbreak{}INF/\allowbreak{}MANIFEST.\allowbreak{}MF} if present & Has no name; cannot be \code{required} by name \\
What it exports & Everything (all packages, to everyone) — no encapsulation & Everything, to anyone else in the unnamed module \\
What it reads & Reads every other module, automatically (\code{requires} any named module works without listing it) & Reads all named modules and the rest of the classpath \\
Typical use & Bridging a legacy, non-modular JAR into a modular application & Legacy applications that haven’t adopted JPMS at all \\
\bottomrule
\end{xltabular}
}

Automatic modules are the “duct tape” that made JPMS adoption practical: the vast majority of the Java ecosystem’s JARs in 2017 had no module descriptor, and automatic modules let a modular application depend on them immediately, with the tradeoff that automatic modules provide none of JPMS’s actual encapsulation guarantees.

\begin{codebox}
\begin{Verbatim}[commandchars=\\\{\}]
\PY{c+c1}{\PYZsh{} guava\PYZhy{}33.0.0.jar has no module\PYZhy{}info.class, but placing it on the module path}
\PY{c+c1}{\PYZsh{} turns it into the automatic module \PYZdq{}com.google.common\PYZdq{} (from its}
\PY{c+c1}{\PYZsh{} Automatic\PYZhy{}Module\PYZhy{}Name manifest entry), usable via a plain \PYZdq{}requires\PYZdq{}:}
java\PY{+w}{ }\PYZhy{}p\PY{+w}{ }libs:out\PY{+w}{ }\PYZhy{}m\PY{+w}{ }com.teampicnic.app/com.teampicnic.app.Main
\end{Verbatim}
\end{codebox}

\subsection[Module Path vs Class Path]{Module Path vs Class Path}
\label{h:17:module-path-vs-class-path}
{\tablefont
\begin{xltabular}{\linewidth}{@{}L{0.470}L{1.226}L{1.304}@{}}
\toprule
\rowcolor{tablehdr}
\thd{} & \thd{Class path (\code{-\allowbreak{}cp} / \code{-\allowbreak{}classpath})} & \thd{Module path (\code{-\allowbreak{}p} / \code{--\allowbreak{}module-\allowbreak{}path})} \\
\midrule
\endhead
Unit & Flat bag of \code{.\allowbreak{}class} files and JARs & Named modules, each with its own boundary \\
Duplicate packages & Silently merged/shadowed (split package risk) & Hard resolution error \\
Encapsulation & None — everything \code{public} is visible to everyone & Enforced — only \code{exports}-ed packages are visible \\
Declares dependencies? & No — implicit, whatever happens to be on the path & Yes — explicit \code{requires} graph, validated at launch \\
Can mix the two? & N/A & Yes — a “hybrid” application can use both simultaneously, with classpath JARs landing in the unnamed module \\
\bottomrule
\end{xltabular}
}

\begin{codebox}
\begin{Verbatim}[commandchars=\\\{\}]
\PY{c+c1}{\PYZsh{} Classic classpath launch — flat, no module resolution at all}
java\PY{+w}{ }\PYZhy{}cp\PY{+w}{ }\PY{l+s+s2}{\PYZdq{}libs/*:out\PYZdq{}}\PY{+w}{ }com.teampicnic.app.Main

\PY{c+c1}{\PYZsh{} Modular launch — module path, resolved and validated against module\PYZhy{}info.java}
java\PY{+w}{ }\PYZhy{}p\PY{+w}{ }\PY{l+s+s2}{\PYZdq{}libs:out\PYZdq{}}\PY{+w}{ }\PYZhy{}m\PY{+w}{ }com.teampicnic.app/com.teampicnic.app.Main
\end{Verbatim}
\end{codebox}

\subsection[Readability and Accessibility Rules]{Readability and Accessibility Rules}
\label{h:17:readability-and-accessibility-rules}
JPMS access checks layer two independent questions on top of each other:

\begin{enumerate}
\item 
\textbf{Readability} — does module A \code{requires} module B (directly, or transitively via \code{requires~\allowbreak{}transitive})? If not, A cannot see \emph{any} type in B, no matter how public.

\item 
\textbf{Accessibility} — assuming A reads B, is the specific package inside B \code{exports}-ed (for normal use) or \code{opens}-ed (for reflection) to A?

\end{enumerate}

\begin{codebox}
\begin{Verbatim}
Module A --requires--> Module B
                          |
                          +-- exports com.b.api        (A can use com.b.api's public types normally)
                          +-- (com.b.internal not exported -> invisible to A, even though B "has" it)
\end{Verbatim}
\end{codebox}

A useful mental model: readability is “can I even see this module exists,” and accessibility is “given that I can see it, which of its packages am I allowed to touch, and how (normal use vs. reflection).”

\begin{codebox}
\begin{Verbatim}[commandchars=\\\{\}]
\PY{n}{module}\PY{+w}{ }\PY{n}{com}\PY{p}{.}\PY{n+na}{teampicnic}\PY{p}{.}\PY{n+na}{app}\PY{+w}{ }\PY{p}{\PYZob{}}
\PY{+w}{    }\PY{n}{requires}\PY{+w}{ }\PY{n}{com}\PY{p}{.}\PY{n+na}{teampicnic}\PY{p}{.}\PY{n+na}{billing}\PY{p}{;}\PY{+w}{ }\PY{c+c1}{// readability: app can now see billing\PYZsq{}s exported packages}
\PY{p}{\PYZcb{}}
\end{Verbatim}
\end{codebox}

If \code{com.\allowbreak{}teampicnic.\allowbreak{}billing} only exports \code{com.\allowbreak{}teampicnic.\allowbreak{}billing.\allowbreak{}api}, then \code{com.\allowbreak{}teampicnic.\allowbreak{}app} compiles fine against \code{Invoice} (which lives in \code{.\allowbreak{}api}), but a reference to \code{com.\allowbreak{}teampicnic.\allowbreak{}billing.\allowbreak{}internal.\allowbreak{}Tax\allowbreak{}Calculator} fails to compile with \code{package~\allowbreak{}com.\allowbreak{}teampicnic.\allowbreak{}billing.\allowbreak{}internal~\allowbreak{}is~\allowbreak{}not~\allowbreak{}visible} — even though \code{app} legitimately \code{requires~\allowbreak{}billing}.

\subsection[jlink, jdeps, and jmod]{\code{jlink}, \code{jdeps}, and \code{jmod}}
\label{h:17:jlink-jdeps-and-jmod}
Three CLI tools work together around JPMS: \code{jdeps} for analysis, \code{jmod} for packaging, and \code{jlink} for producing a minimal custom runtime.

\begin{codebox}
\begin{Verbatim}[commandchars=\\\{\}]
\PY{c+c1}{\PYZsh{} jdeps: what does this JAR actually depend on, and is it already modular?}
jdeps\PY{+w}{ }\PYZhy{}\PYZhy{}print\PYZhy{}module\PYZhy{}deps\PY{+w}{ }app.jar\PY{+w}{          }\PY{c+c1}{\PYZsh{} e.g. java.base,java.sql,java.logging}
jdeps\PY{+w}{ }\PYZhy{}\PYZhy{}jdk\PYZhy{}internals\PY{+w}{ }app.jar\PY{+w}{              }\PY{c+c1}{\PYZsh{} flags usage of JDK\PYZhy{}internal (encapsulated) APIs}

\PY{c+c1}{\PYZsh{} jmod: package a module into the .jmod format \PYZhy{}\PYZhy{} mainly used for JDK modules}
\PY{c+c1}{\PYZsh{} and modules bundling native libraries; regular app modules ship as plain JARs.}
jmod\PY{+w}{ }create\PY{+w}{ }\PYZhy{}\PYZhy{}class\PYZhy{}path\PY{+w}{ }out/com.teampicnic.billing\PY{+w}{ }\PYZhy{}\PYZhy{}module\PYZhy{}version\PY{+w}{ }\PY{l+m}{1}.0\PY{+w}{ }billing.jmod

\PY{c+c1}{\PYZsh{} jlink: assemble a minimal, self\PYZhy{}contained runtime image with only the}
\PY{c+c1}{\PYZsh{} modules your app needs \PYZhy{}\PYZhy{} no full JDK required on the deployment target.}
jlink\PY{+w}{ }\PYZhy{}\PYZhy{}module\PYZhy{}path\PY{+w}{ }out:\PY{n+nv}{\PYZdl{}JAVA\PYZus{}HOME}/jmods\PY{+w}{ }\PY{l+s+se}{\PYZbs{}}
\PY{+w}{      }\PYZhy{}\PYZhy{}add\PYZhy{}modules\PY{+w}{ }com.teampicnic.app\PY{+w}{ }\PY{l+s+se}{\PYZbs{}}
\PY{+w}{      }\PYZhy{}\PYZhy{}output\PY{+w}{ }custom\PYZhy{}runtime\PY{+w}{ }\PY{l+s+se}{\PYZbs{}}
\PY{+w}{      }\PYZhy{}\PYZhy{}strip\PYZhy{}debug\PY{+w}{ }\PYZhy{}\PYZhy{}no\PYZhy{}header\PYZhy{}files\PY{+w}{ }\PYZhy{}\PYZhy{}no\PYZhy{}man\PYZhy{}pages\PY{+w}{ }\PYZhy{}\PYZhy{}compress\PY{o}{=}\PY{l+m}{2}

./custom\PYZhy{}runtime/bin/java\PY{+w}{ }\PYZhy{}m\PY{+w}{ }com.teampicnic.app/com.teampicnic.app.Main
\end{Verbatim}
\end{codebox}

\code{jlink}’s output is a real win for container images: instead of a full \textasciitilde{}300MB JDK base image, you ship only the \textasciitilde{}40-60MB of modules your app actually touches, shrinking both image size and attack surface.

\subsection[Strong Encapsulation and --add-exports/--add-opens]{Strong Encapsulation and \code{--\allowbreak{}add-\allowbreak{}exports}/\code{--\allowbreak{}add-\allowbreak{}opens}}
\label{h:17:strong-encapsulation-and---add-exports--add-opens}
Since Java 16 (JEP 396), strong encapsulation of JDK internals is \textbf{on by default} — code can no longer reflectively reach into packages like \code{jdk.\allowbreak{}internal.\allowbreak{}misc} or \code{sun.\allowbreak{}nio.\allowbreak{}ch} unless the module system is explicitly told to allow it. Libraries that relied on these internals (older versions of Lombok, Mockito, various serialization frameworks) had to add explicit flags or be updated.

\begin{codebox}
\begin{Verbatim}[commandchars=\\\{\}]
\PY{c+c1}{\PYZsh{} \PYZhy{}\PYZhy{}add\PYZhy{}exports: grants normal (non\PYZhy{}reflective) access to an otherwise\PYZhy{}unexported}
\PY{c+c1}{\PYZsh{} package, from a specific module to a specific module.}
java\PY{+w}{ }\PYZhy{}\PYZhy{}add\PYZhy{}exports\PY{+w}{ }java.base/jdk.internal.misc\PY{o}{=}com.teampicnic.app\PY{+w}{ }\PYZhy{}p\PY{+w}{ }out\PY{+w}{ }\PYZhy{}m\PY{+w}{ }com.teampicnic.app/com.teampicnic.app.Main

\PY{c+c1}{\PYZsh{} \PYZhy{}\PYZhy{}add\PYZhy{}opens: grants deep reflection access to an otherwise\PYZhy{}closed package.}
java\PY{+w}{ }\PYZhy{}\PYZhy{}add\PYZhy{}opens\PY{+w}{ }java.base/java.lang\PY{o}{=}ALL\PYZhy{}UNNAMED\PY{+w}{ }\PYZhy{}p\PY{+w}{ }out\PY{+w}{ }\PYZhy{}m\PY{+w}{ }com.teampicnic.app/com.teampicnic.app.Main
\end{Verbatim}
\end{codebox}

\code{ALL-\allowbreak{}UNNAMED} means “grant this to the unnamed module,” i.e., to whatever is on the plain classpath — this is the flag combination you’ll most often see in the wild when a classpath-based library (Mockito, older Jackson, testing frameworks) needs deep reflection into \code{java.\allowbreak{}lang} or \code{java.\allowbreak{}util} internals to do bytecode manipulation or field injection.

These are explicitly called \textbf{escape hatches}: they exist so real-world code that isn’t ready for strong encapsulation can keep running, not as a long-term architectural pattern. Starting in newer JDKs, using them (or performing “restricted” native/reflective operations without them) triggers warnings, and the trend is toward requiring them more strictly with each release — treat every \code{--\allowbreak{}add-\allowbreak{}opens} in a build file as a line item to eventually remove, not a permanent fixture.

\subsection[Worked Multi-Module Example]{Worked Multi-Module Example}
\label{h:17:worked-multi-module-example}
Two modules: a library module that exposes a \code{Greeter}, and an application module that consumes it.

\begin{codebox}
\begin{Verbatim}
project/
├── com.teampicnic.greetings/
│   ├── module-info.java
│   └── com/
│       └── teampicnic/
│           └── greetings/
│               └── Greeter.java
└── com.teampicnic.app/
    ├── module-info.java
    └── com/
        └── teampicnic/
            └── app/
                └── Main.java
\end{Verbatim}
\end{codebox}

\begin{codebox}
\begin{Verbatim}[commandchars=\\\{\}]
\PY{c+c1}{// project/com.teampicnic.greetings/module\PYZhy{}info.java}
\PY{n}{module}\PY{+w}{ }\PY{n}{com}\PY{p}{.}\PY{n+na}{teampicnic}\PY{p}{.}\PY{n+na}{greetings}\PY{+w}{ }\PY{p}{\PYZob{}}
\PY{+w}{    }\PY{n}{exports}\PY{+w}{ }\PY{n}{com}\PY{p}{.}\PY{n+na}{teampicnic}\PY{p}{.}\PY{n+na}{greetings}\PY{p}{;}
\PY{p}{\PYZcb{}}
\end{Verbatim}
\end{codebox}

\begin{codebox}
\begin{Verbatim}[commandchars=\\\{\}]
\PY{c+c1}{// project/com.teampicnic.greetings/com/teampicnic/greetings/Greeter.java}
\PY{k+kn}{package}\PY{+w}{ }\PY{n+nn}{com.teampicnic.greetings}\PY{p}{;}

\PY{k+kd}{public}\PY{+w}{ }\PY{k+kd}{class} \PY{n+nc}{Greeter}\PY{+w}{ }\PY{p}{\PYZob{}}
\PY{+w}{    }\PY{k+kd}{public}\PY{+w}{ }\PY{n}{String}\PY{+w}{ }\PY{n+nf}{greet}\PY{p}{(}\PY{n}{String}\PY{+w}{ }\PY{n}{name}\PY{p}{)}\PY{+w}{ }\PY{p}{\PYZob{}}
\PY{+w}{        }\PY{k}{return}\PY{+w}{ }\PY{l+s}{\PYZdq{}}\PY{l+s}{Hello, }\PY{l+s}{\PYZdq{}}\PY{+w}{ }\PY{o}{+}\PY{+w}{ }\PY{n}{name}\PY{+w}{ }\PY{o}{+}\PY{+w}{ }\PY{l+s}{\PYZdq{}}\PY{l+s}{!}\PY{l+s}{\PYZdq{}}\PY{p}{;}
\PY{+w}{    }\PY{p}{\PYZcb{}}
\PY{p}{\PYZcb{}}
\end{Verbatim}
\end{codebox}

\begin{codebox}
\begin{Verbatim}[commandchars=\\\{\}]
\PY{c+c1}{// project/com.teampicnic.app/module\PYZhy{}info.java}
\PY{n}{module}\PY{+w}{ }\PY{n}{com}\PY{p}{.}\PY{n+na}{teampicnic}\PY{p}{.}\PY{n+na}{app}\PY{+w}{ }\PY{p}{\PYZob{}}
\PY{+w}{    }\PY{n}{requires}\PY{+w}{ }\PY{n}{com}\PY{p}{.}\PY{n+na}{teampicnic}\PY{p}{.}\PY{n+na}{greetings}\PY{p}{;}
\PY{p}{\PYZcb{}}
\end{Verbatim}
\end{codebox}

\begin{codebox}
\begin{Verbatim}[commandchars=\\\{\}]
\PY{c+c1}{// project/com.teampicnic.app/com/teampicnic/app/Main.java}
\PY{k+kn}{package}\PY{+w}{ }\PY{n+nn}{com.teampicnic.app}\PY{p}{;}

\PY{k+kn}{import}\PY{+w}{ }\PY{n+nn}{com.teampicnic.greetings.Greeter}\PY{p}{;}

\PY{k+kd}{public}\PY{+w}{ }\PY{k+kd}{class} \PY{n+nc}{Main}\PY{+w}{ }\PY{p}{\PYZob{}}
\PY{+w}{    }\PY{k+kd}{public}\PY{+w}{ }\PY{k+kd}{static}\PY{+w}{ }\PY{k+kt}{void}\PY{+w}{ }\PY{n+nf}{main}\PY{p}{(}\PY{n}{String}\PY{o}{[}\PY{o}{]}\PY{+w}{ }\PY{n}{args}\PY{p}{)}\PY{+w}{ }\PY{p}{\PYZob{}}
\PY{+w}{        }\PY{n}{System}\PY{p}{.}\PY{n+na}{out}\PY{p}{.}\PY{n+na}{println}\PY{p}{(}\PY{k}{new}\PY{+w}{ }\PY{n}{Greeter}\PY{p}{(}\PY{p}{)}\PY{p}{.}\PY{n+na}{greet}\PY{p}{(}\PY{l+s}{\PYZdq{}}\PY{l+s}{code review}\PY{l+s}{\PYZdq{}}\PY{p}{)}\PY{p}{)}\PY{p}{;}
\PY{+w}{    }\PY{p}{\PYZcb{}}
\PY{p}{\PYZcb{}}
\end{Verbatim}
\end{codebox}

Compile both modules using \code{--\allowbreak{}module-\allowbreak{}source-\allowbreak{}path}, which lets \code{javac} resolve inter-module dependencies from source in one pass:

\begin{codebox}
\begin{Verbatim}[commandchars=\\\{\}]
mkdir\PY{+w}{ }\PYZhy{}p\PY{+w}{ }out

javac\PY{+w}{ }\PYZhy{}d\PY{+w}{ }out\PY{+w}{ }\PY{l+s+se}{\PYZbs{}}
\PY{+w}{      }\PYZhy{}\PYZhy{}module\PYZhy{}source\PYZhy{}path\PY{+w}{ }project\PY{+w}{ }\PY{l+s+se}{\PYZbs{}}
\PY{+w}{      }\PY{k}{\PYZdl{}(}find\PY{+w}{ }project\PY{+w}{ }\PYZhy{}name\PY{+w}{ }\PY{l+s+s2}{\PYZdq{}*.java\PYZdq{}}\PY{k}{)}

\PY{c+c1}{\PYZsh{} out/com.teampicnic.greetings/com/teampicnic/greetings/Greeter.class}
\PY{c+c1}{\PYZsh{} out/com.teampicnic.app/com/teampicnic/app/Main.class}
\end{Verbatim}
\end{codebox}

Run it directly off the module path:

\begin{codebox}
\begin{Verbatim}[commandchars=\\\{\}]
java\PY{+w}{ }\PYZhy{}p\PY{+w}{ }out\PY{+w}{ }\PYZhy{}m\PY{+w}{ }com.teampicnic.app/com.teampicnic.app.Main
\PY{c+c1}{\PYZsh{} Hello, code review!}
\end{Verbatim}
\end{codebox}

Build a minimal custom runtime with \code{jlink}:

\begin{codebox}
\begin{Verbatim}[commandchars=\\\{\}]
jlink\PY{+w}{ }\PYZhy{}\PYZhy{}module\PYZhy{}path\PY{+w}{ }out:\PY{n+nv}{\PYZdl{}JAVA\PYZus{}HOME}/jmods\PY{+w}{ }\PY{l+s+se}{\PYZbs{}}
\PY{+w}{      }\PYZhy{}\PYZhy{}add\PYZhy{}modules\PY{+w}{ }com.teampicnic.app\PY{+w}{ }\PY{l+s+se}{\PYZbs{}}
\PY{+w}{      }\PYZhy{}\PYZhy{}launcher\PY{+w}{ }\PY{n+nv}{greet}\PY{o}{=}com.teampicnic.app/com.teampicnic.app.Main\PY{+w}{ }\PY{l+s+se}{\PYZbs{}}
\PY{+w}{      }\PYZhy{}\PYZhy{}output\PY{+w}{ }custom\PYZhy{}runtime

./custom\PYZhy{}runtime/bin/greet
\PY{c+c1}{\PYZsh{} Hello, code review!}
\end{Verbatim}
\end{codebox}

The \code{--\allowbreak{}launcher} flag is a nice \code{jlink} feature reviewers sometimes miss: it generates a native-feeling launcher script (\code{custom-\allowbreak{}runtime/\allowbreak{}bin/\allowbreak{}greet}) so end users never need to know or type the module/main-class combination.

\subsection[Migration Strategy]{Migration Strategy}
\label{h:17:migration-strategy}
Two broad strategies exist for turning an existing multi-JAR, classpath-based application into a set of real JPMS modules.

{\tablefontsmall
\begin{xltabular}{\linewidth}{@{}L{0.417}L{1.194}L{1.194}L{1.194}@{}}
\toprule
\rowcolor{tablehdr}
\thd{Strategy} & \thd{Approach} & \thd{Pros} & \thd{Cons} \\
\midrule
\endhead
\textbf{Bottom-up} & Modularize the lowest-level, most-depended-upon libraries first (utility/core JARs), then work up toward the application & Each step is small and independently verifiable; leaf libraries have the fewest dependencies to reconcile & Slow — the application itself stays non-modular for a long time; benefits (strong encapsulation for the app) arrive last \\
\textbf{Top-down} & Add a \code{module-\allowbreak{}info.\allowbreak{}java} to the application first, treating every dependency as an automatic module, then gradually convert dependencies to real modules one at a time & Gets the app onto the module path (and gets \code{jlink} working) quickly & Automatic modules provide none of JPMS’s actual safety, so early gains are mostly tooling-level (custom runtimes), not encapsulation \\
\bottomrule
\end{xltabular}
}

In practice, most real migrations are a hybrid: start top-down to get the application itself running on the module path (with \code{jdeps} used heavily to find hidden internal-API dependencies first), then chip away bottom-up at whichever dependencies are most painful — usually anything relying on \code{sun.\allowbreak{}misc.\allowbreak{}Unsafe}, reflection into \code{java.\allowbreak{}lang}, or shading/relocation tricks that conflict with module boundaries.

\begin{codebox}
\begin{Verbatim}[commandchars=\\\{\}]
\PY{c+c1}{\PYZsh{} jdeps is the standard first step of any migration: find out what a JAR}
\PY{c+c1}{\PYZsh{} actually needs before writing a single module\PYZhy{}info.java.}
jdeps\PY{+w}{ }\PYZhy{}\PYZhy{}generate\PYZhy{}module\PYZhy{}info\PY{+w}{ }out\PYZhy{}modules/\PY{+w}{ }legacy\PYZhy{}app.jar
\end{Verbatim}
\end{codebox}

\code{jdeps~\allowbreak{}--\allowbreak{}generate-\allowbreak{}module-\allowbreak{}info} will even scaffold a best-effort \code{module-\allowbreak{}info.\allowbreak{}java} for you based on observed bytecode dependencies, which is a good starting draft — but always needs manual review, since it can’t infer intent (e.g., which packages are truly meant to be public API versus which merely happen to be used across a package boundary today).

\subsection[Why Many Projects Still Don’t Use JPMS]{Why Many Projects Still Don’t Use JPMS}
\label{h:17:why-many-projects-still-dont-use-jpms}
Despite being available since 2017, most application codebases (as opposed to the JDK itself, and a handful of major libraries) never adopt module descriptors, and that is a reasonable, defensible choice:

\begin{itemize}
\item 
\textbf{Split packages in the existing dependency tree} are extremely common (two transitive dependencies exporting the same package) and can require upgrading, forking, or shading dependencies just to get past module resolution errors — for zero new feature value.

\item 
\textbf{Spring, Hibernate, and most application frameworks} still work perfectly well on the plain classpath, and JPMS’s encapsulation benefits matter far less inside a single deployable application than for library authors shipping to unrelated consumers.

\item 
\textbf{Build tool support} has historically been fiddly, with friction around test compilation, \code{--\allowbreak{}add-\allowbreak{}opens} requirements for mocking frameworks, and IDE tooling lag.

\item 
\textbf{The main payoff for most teams is \code{jlink}} (smaller container images), which is achievable alongside simpler tools (slim base images, ArchUnit-style architecture tests) without ever writing a \code{module-\allowbreak{}info.\allowbreak{}java}.

\end{itemize}

A defensible interview answer: JPMS is essential for JDK maintainers and library authors guaranteeing encapsulation across organizational boundaries; for a typical internal application, the same goals are often reachable more cheaply through build-tool module boundaries and architecture linting — so skipping it is not automatically a red flag.

\section[Foreign Function \& Memory API]{Foreign Function \& Memory API}
\label{h:17:foreign-function--memory-api}
The Foreign Function \& Memory API (FFM API, \code{java.\allowbreak{}lang.\allowbreak{}foreign} package) was finalized in \textbf{JDK 22 via JEP 454}, after multiple rounds of incubation and preview (starting as the Foreign-Memory Access API in JDK 14). It lets Java code allocate and manipulate off-heap memory and call native © functions \textbf{without writing any native code, without JNI headers, and without \code{Unsafe}}.

\subsection[From JNI to FFM: Why It Changed]{From JNI to FFM: Why It Changed}
\label{h:17:from-jni-to-ffm-why-it-changed}
JNI (Java Native Interface, since Java 1.1) required writing actual C glue code, compiling it into a platform-specific shared library, and loading it with \code{System.\allowbreak{}loadLibrary}. \code{Unsafe} and \code{ByteBuffer} (direct buffers) let you get \emph{some} off-heap access from pure Java, but both had serious limitations.

\begin{codebox}
\begin{Verbatim}[commandchars=\\\{\}]
\PY{c+c1}{// The old ByteBuffer / Unsafe world: off\PYZhy{}heap memory, but no safe deallocation}
\PY{c+c1}{// story, no structured layout description, and no way to call arbitrary}
\PY{c+c1}{// native functions without JNI.}
\PY{n}{java}\PY{p}{.}\PY{n+na}{nio}\PY{p}{.}\PY{n+na}{ByteBuffer}\PY{+w}{ }\PY{n}{direct}\PY{+w}{ }\PY{o}{=}\PY{+w}{ }\PY{n}{java}\PY{p}{.}\PY{n+na}{nio}\PY{p}{.}\PY{n+na}{ByteBuffer}\PY{p}{.}\PY{n+na}{allocateDirect}\PY{p}{(}\PY{l+m+mi}{1024}\PY{p}{)}\PY{p}{;}
\PY{n}{direct}\PY{p}{.}\PY{n+na}{putInt}\PY{p}{(}\PY{l+m+mi}{0}\PY{p}{,}\PY{+w}{ }\PY{l+m+mi}{42}\PY{p}{)}\PY{p}{;}
\PY{k+kt}{int}\PY{+w}{ }\PY{n}{value}\PY{+w}{ }\PY{o}{=}\PY{+w}{ }\PY{n}{direct}\PY{p}{.}\PY{n+na}{getInt}\PY{p}{(}\PY{l+m+mi}{0}\PY{p}{)}\PY{p}{;}
\PY{c+c1}{// No explicit \PYZdq{}free\PYZdq{} \PYZhy{}\PYZhy{} direct buffers are reclaimed only by GC + a Cleaner,}
\PY{c+c1}{// with no deterministic timing and no way to call a C function using this memory.}
\end{Verbatim}
\end{codebox}

FFM replaces this with a small set of pure-Java classes: \code{Arena} for lifecycle-managed allocation, \code{MemorySegment} for typed, bounds-checked pointers, \code{MemoryLayout}/\code{ValueLayout} for describing native data shapes, and \code{Linker} for calling and being called by native functions — all without JNI’s C compilation step.

\subsection[Arena and Deterministic Deallocation]{\code{Arena} and Deterministic Deallocation}
\label{h:17:arena-and-deterministic-deallocation}
An \code{Arena} controls the \textbf{lifetime} of the off-heap memory it allocates. All memory allocated through an arena becomes invalid the moment the arena is closed — accessing it afterward throws \code{IllegalStateException} rather than corrupting memory or segfaulting, which is the safety guarantee JNI/\code{Unsafe} never gave you.

{\tablefont
\begin{xltabular}{\linewidth}{@{}L{0.437}L{0.573}L{1.432}L{1.559}@{}}
\toprule
\rowcolor{tablehdr}
\thd{Arena kind} & \thd{Created with} & \thd{Lifetime} & \thd{Thread confinement} \\
\midrule
\endhead
Confined & \code{Arena.\allowbreak{}ofConfined()} & Explicit \code{close()} (use try-with-resources) & Only the owning thread may access or close it \\
Shared & \code{Arena.\allowbreak{}ofShared()} & Explicit \code{close()} & Any thread may access it; any thread may close it \\
Auto & \code{Arena.\allowbreak{}ofAuto()} & Garbage-collected — no \code{close()} at all & Any thread \\
Global & \code{Arena.\allowbreak{}global()} & Never freed — lives for the whole JVM process & Any thread \\
\bottomrule
\end{xltabular}
}

\begin{codebox}
\begin{Verbatim}[commandchars=\\\{\}]
\PY{k+kn}{import}\PY{+w}{ }\PY{n+nn}{java.lang.foreign.Arena}\PY{p}{;}
\PY{k+kn}{import}\PY{+w}{ }\PY{n+nn}{java.lang.foreign.MemorySegment}\PY{p}{;}
\PY{k+kn}{import}\PY{+w}{ }\PY{n+nn}{java.lang.foreign.ValueLayout}\PY{p}{;}

\PY{k+kd}{public}\PY{+w}{ }\PY{k+kd}{class} \PY{n+nc}{ArenaDemo}\PY{+w}{ }\PY{p}{\PYZob{}}
\PY{+w}{    }\PY{k+kd}{public}\PY{+w}{ }\PY{k+kd}{static}\PY{+w}{ }\PY{k+kt}{void}\PY{+w}{ }\PY{n+nf}{main}\PY{p}{(}\PY{n}{String}\PY{o}{[}\PY{o}{]}\PY{+w}{ }\PY{n}{args}\PY{p}{)}\PY{+w}{ }\PY{p}{\PYZob{}}
\PY{+w}{        }\PY{k}{try}\PY{+w}{ }\PY{p}{(}\PY{n}{Arena}\PY{+w}{ }\PY{n}{arena}\PY{+w}{ }\PY{o}{=}\PY{+w}{ }\PY{n}{Arena}\PY{p}{.}\PY{n+na}{ofConfined}\PY{p}{(}\PY{p}{)}\PY{p}{)}\PY{+w}{ }\PY{p}{\PYZob{}}\PY{+w}{ }\PY{c+c1}{// single\PYZhy{}threaded, closed like a file handle}
\PY{+w}{            }\PY{n}{MemorySegment}\PY{+w}{ }\PY{n}{segment}\PY{+w}{ }\PY{o}{=}\PY{+w}{ }\PY{n}{arena}\PY{p}{.}\PY{n+na}{allocate}\PY{p}{(}\PY{n}{ValueLayout}\PY{p}{.}\PY{n+na}{JAVA\PYZus{}INT}\PY{p}{.}\PY{n+na}{byteSize}\PY{p}{(}\PY{p}{)}\PY{+w}{ }\PY{o}{*}\PY{+w}{ }\PY{l+m+mi}{4}\PY{p}{)}\PY{p}{;}
\PY{+w}{            }\PY{k}{for}\PY{+w}{ }\PY{p}{(}\PY{k+kt}{int}\PY{+w}{ }\PY{n}{i}\PY{+w}{ }\PY{o}{=}\PY{+w}{ }\PY{l+m+mi}{0}\PY{p}{;}\PY{+w}{ }\PY{n}{i}\PY{+w}{ }\PY{o}{\PYZlt{}}\PY{+w}{ }\PY{l+m+mi}{4}\PY{p}{;}\PY{+w}{ }\PY{n}{i}\PY{o}{+}\PY{o}{+}\PY{p}{)}\PY{+w}{ }\PY{p}{\PYZob{}}
\PY{+w}{                }\PY{n}{segment}\PY{p}{.}\PY{n+na}{setAtIndex}\PY{p}{(}\PY{n}{ValueLayout}\PY{p}{.}\PY{n+na}{JAVA\PYZus{}INT}\PY{p}{,}\PY{+w}{ }\PY{n}{i}\PY{p}{,}\PY{+w}{ }\PY{n}{i}\PY{+w}{ }\PY{o}{*}\PY{+w}{ }\PY{n}{i}\PY{p}{)}\PY{p}{;}
\PY{+w}{            }\PY{p}{\PYZcb{}}
\PY{+w}{            }\PY{k}{for}\PY{+w}{ }\PY{p}{(}\PY{k+kt}{int}\PY{+w}{ }\PY{n}{i}\PY{+w}{ }\PY{o}{=}\PY{+w}{ }\PY{l+m+mi}{0}\PY{p}{;}\PY{+w}{ }\PY{n}{i}\PY{+w}{ }\PY{o}{\PYZlt{}}\PY{+w}{ }\PY{l+m+mi}{4}\PY{p}{;}\PY{+w}{ }\PY{n}{i}\PY{o}{+}\PY{o}{+}\PY{p}{)}\PY{+w}{ }\PY{p}{\PYZob{}}
\PY{+w}{                }\PY{n}{System}\PY{p}{.}\PY{n+na}{out}\PY{p}{.}\PY{n+na}{println}\PY{p}{(}\PY{n}{segment}\PY{p}{.}\PY{n+na}{getAtIndex}\PY{p}{(}\PY{n}{ValueLayout}\PY{p}{.}\PY{n+na}{JAVA\PYZus{}INT}\PY{p}{,}\PY{+w}{ }\PY{n}{i}\PY{p}{)}\PY{p}{)}\PY{p}{;}
\PY{+w}{            }\PY{p}{\PYZcb{}}
\PY{+w}{        }\PY{p}{\PYZcb{}}\PY{+w}{ }\PY{c+c1}{// memory is deterministically freed right here}

\PY{+w}{        }\PY{c+c1}{// Using \PYZsq{}segment\PYZsq{} past this point throws IllegalStateException instead of}
\PY{+w}{        }\PY{c+c1}{// corrupting memory or segfaulting \PYZhy{}\PYZhy{} the core safety win over Unsafe.}
\PY{+w}{    }\PY{p}{\PYZcb{}}
\PY{p}{\PYZcb{}}
\end{Verbatim}
\end{codebox}

Try-with-resources is the idiomatic way to use a confined arena, mirroring exactly how you’d manage a \code{FileInputStream} — the “close what you open” instinct from Chapter 6 transfers directly.

\subsection[MemorySegment, MemoryLayout, ValueLayout]{\code{MemorySegment}, \code{MemoryLayout}, \code{ValueLayout}}
\label{h:17:memorysegment-memorylayout-valuelayout}
A \code{MemorySegment} is a typed, bounds-checked view over a contiguous region of memory — either off-heap (native) or backed by an on-heap Java array.

\begin{codebox}
\begin{Verbatim}[commandchars=\\\{\}]
\PY{k+kn}{import}\PY{+w}{ }\PY{n+nn}{java.lang.foreign.MemorySegment}\PY{p}{;}
\PY{k+kn}{import}\PY{+w}{ }\PY{n+nn}{java.lang.foreign.ValueLayout}\PY{p}{;}

\PY{k+kd}{public}\PY{+w}{ }\PY{k+kd}{class} \PY{n+nc}{SegmentDemo}\PY{+w}{ }\PY{p}{\PYZob{}}
\PY{+w}{    }\PY{k+kd}{public}\PY{+w}{ }\PY{k+kd}{static}\PY{+w}{ }\PY{k+kt}{void}\PY{+w}{ }\PY{n+nf}{main}\PY{p}{(}\PY{n}{String}\PY{o}{[}\PY{o}{]}\PY{+w}{ }\PY{n}{args}\PY{p}{)}\PY{+w}{ }\PY{p}{\PYZob{}}
\PY{+w}{        }\PY{c+c1}{// Wrap an existing Java array as a MemorySegment (heap segment, no allocation)}
\PY{+w}{        }\PY{k+kt}{int}\PY{o}{[}\PY{o}{]}\PY{+w}{ }\PY{n}{javaArray}\PY{+w}{ }\PY{o}{=}\PY{+w}{ }\PY{p}{\PYZob{}}\PY{l+m+mi}{10}\PY{p}{,}\PY{+w}{ }\PY{l+m+mi}{20}\PY{p}{,}\PY{+w}{ }\PY{l+m+mi}{30}\PY{p}{\PYZcb{}}\PY{p}{;}
\PY{+w}{        }\PY{n}{MemorySegment}\PY{+w}{ }\PY{n}{heapSegment}\PY{+w}{ }\PY{o}{=}\PY{+w}{ }\PY{n}{MemorySegment}\PY{p}{.}\PY{n+na}{ofArray}\PY{p}{(}\PY{n}{javaArray}\PY{p}{)}\PY{p}{;}
\PY{+w}{        }\PY{n}{System}\PY{p}{.}\PY{n+na}{out}\PY{p}{.}\PY{n+na}{println}\PY{p}{(}\PY{n}{heapSegment}\PY{p}{.}\PY{n+na}{getAtIndex}\PY{p}{(}\PY{n}{ValueLayout}\PY{p}{.}\PY{n+na}{JAVA\PYZus{}INT}\PY{p}{,}\PY{+w}{ }\PY{l+m+mi}{1}\PY{p}{)}\PY{p}{)}\PY{p}{;}\PY{+w}{ }\PY{c+c1}{// 20}

\PY{+w}{        }\PY{c+c1}{// Out\PYZhy{}of\PYZhy{}bounds access throws IndexOutOfBoundsException, not a crash}
\PY{+w}{        }\PY{k}{try}\PY{+w}{ }\PY{p}{\PYZob{}}
\PY{+w}{            }\PY{n}{heapSegment}\PY{p}{.}\PY{n+na}{getAtIndex}\PY{p}{(}\PY{n}{ValueLayout}\PY{p}{.}\PY{n+na}{JAVA\PYZus{}INT}\PY{p}{,}\PY{+w}{ }\PY{l+m+mi}{10}\PY{p}{)}\PY{p}{;}
\PY{+w}{        }\PY{p}{\PYZcb{}}\PY{+w}{ }\PY{k}{catch}\PY{+w}{ }\PY{p}{(}\PY{n}{IndexOutOfBoundsException}\PY{+w}{ }\PY{n}{e}\PY{p}{)}\PY{+w}{ }\PY{p}{\PYZob{}}
\PY{+w}{            }\PY{n}{System}\PY{p}{.}\PY{n+na}{out}\PY{p}{.}\PY{n+na}{println}\PY{p}{(}\PY{l+s}{\PYZdq{}}\PY{l+s}{Bounds\PYZhy{}checked: }\PY{l+s}{\PYZdq{}}\PY{+w}{ }\PY{o}{+}\PY{+w}{ }\PY{n}{e}\PY{p}{.}\PY{n+na}{getMessage}\PY{p}{(}\PY{p}{)}\PY{p}{)}\PY{p}{;}
\PY{+w}{        }\PY{p}{\PYZcb{}}
\PY{+w}{    }\PY{p}{\PYZcb{}}
\PY{p}{\PYZcb{}}
\end{Verbatim}
\end{codebox}

\code{MemoryLayout} describes the \emph{shape} of memory (its size, alignment, and — for structs — named fields), independent of any specific segment. \code{ValueLayout} is the family of primitive layouts (\code{JAVA\_\allowbreak{}INT}, \code{JAVA\_\allowbreak{}LONG}, \code{JAVA\_\allowbreak{}DOUBLE}, \code{ADDRESS}, …) used to describe how raw bytes should be interpreted.

\begin{codebox}
\begin{Verbatim}[commandchars=\\\{\}]
\PY{k+kn}{import}\PY{+w}{ }\PY{n+nn}{java.lang.foreign.ValueLayout}\PY{p}{;}

\PY{c+c1}{// ValueLayout constants describe primitive C types in platform\PYZhy{}independent Java code:}
\PY{n}{ValueLayout}\PY{p}{.}\PY{n+na}{OfInt}\PY{+w}{    }\PY{n}{javaInt}\PY{+w}{    }\PY{o}{=}\PY{+w}{ }\PY{n}{ValueLayout}\PY{p}{.}\PY{n+na}{JAVA\PYZus{}INT}\PY{p}{;}\PY{+w}{    }\PY{c+c1}{// 4\PYZhy{}byte int}
\PY{n}{ValueLayout}\PY{p}{.}\PY{n+na}{OfLong}\PY{+w}{   }\PY{n}{javaLong}\PY{+w}{   }\PY{o}{=}\PY{+w}{ }\PY{n}{ValueLayout}\PY{p}{.}\PY{n+na}{JAVA\PYZus{}LONG}\PY{p}{;}\PY{+w}{   }\PY{c+c1}{// 8\PYZhy{}byte long}
\PY{n}{ValueLayout}\PY{p}{.}\PY{n+na}{OfDouble}\PY{+w}{ }\PY{n}{javaDouble}\PY{+w}{ }\PY{o}{=}\PY{+w}{ }\PY{n}{ValueLayout}\PY{p}{.}\PY{n+na}{JAVA\PYZus{}DOUBLE}\PY{p}{;}\PY{+w}{ }\PY{c+c1}{// 8\PYZhy{}byte double}
\PY{n}{AddressLayout}\PY{+w}{        }\PY{n}{pointer}\PY{+w}{    }\PY{o}{=}\PY{+w}{ }\PY{n}{ValueLayout}\PY{p}{.}\PY{n+na}{ADDRESS}\PY{p}{;}\PY{+w}{     }\PY{c+c1}{// native pointer (4 or 8 bytes)}
\end{Verbatim}
\end{codebox}

\subsection[VarHandle Access into Segments]{\code{VarHandle} Access into Segments}
\label{h:17:varhandle-access-into-segments}
A \code{VarHandle} obtained from a layout provides typed, potentially volatile/atomic access into a \code{MemorySegment} at a computed offset — the FFM equivalent of a field accessor, but for off-heap memory.

\begin{codebox}
\begin{Verbatim}[commandchars=\\\{\}]
\PY{k+kn}{import}\PY{+w}{ }\PY{n+nn}{java.lang.foreign.Arena}\PY{p}{;}
\PY{k+kn}{import}\PY{+w}{ }\PY{n+nn}{java.lang.foreign.MemoryLayout}\PY{p}{;}
\PY{k+kn}{import}\PY{+w}{ }\PY{n+nn}{java.lang.foreign.MemorySegment}\PY{p}{;}
\PY{k+kn}{import}\PY{+w}{ }\PY{n+nn}{java.lang.foreign.ValueLayout}\PY{p}{;}
\PY{k+kn}{import}\PY{+w}{ }\PY{n+nn}{java.lang.invoke.VarHandle}\PY{p}{;}

\PY{k+kd}{public}\PY{+w}{ }\PY{k+kd}{class} \PY{n+nc}{VarHandleDemo}\PY{+w}{ }\PY{p}{\PYZob{}}
\PY{+w}{    }\PY{k+kd}{public}\PY{+w}{ }\PY{k+kd}{static}\PY{+w}{ }\PY{k+kt}{void}\PY{+w}{ }\PY{n+nf}{main}\PY{p}{(}\PY{n}{String}\PY{o}{[}\PY{o}{]}\PY{+w}{ }\PY{n}{args}\PY{p}{)}\PY{+w}{ }\PY{p}{\PYZob{}}
\PY{+w}{        }\PY{n}{MemoryLayout}\PY{+w}{ }\PY{n}{intArrayLayout}\PY{+w}{ }\PY{o}{=}\PY{+w}{ }\PY{n}{MemoryLayout}\PY{p}{.}\PY{n+na}{sequenceLayout}\PY{p}{(}\PY{l+m+mi}{4}\PY{p}{,}\PY{+w}{ }\PY{n}{ValueLayout}\PY{p}{.}\PY{n+na}{JAVA\PYZus{}INT}\PY{p}{)}\PY{p}{;}
\PY{+w}{        }\PY{n}{VarHandle}\PY{+w}{ }\PY{n}{elementHandle}\PY{+w}{ }\PY{o}{=}\PY{+w}{ }\PY{n}{intArrayLayout}\PY{p}{.}\PY{n+na}{varHandle}\PY{p}{(}
\PY{+w}{            }\PY{n}{MemoryLayout}\PY{p}{.}\PY{n+na}{PathElement}\PY{p}{.}\PY{n+na}{sequenceElement}\PY{p}{(}\PY{p}{)}\PY{p}{)}\PY{p}{;}

\PY{+w}{        }\PY{k}{try}\PY{+w}{ }\PY{p}{(}\PY{n}{Arena}\PY{+w}{ }\PY{n}{arena}\PY{+w}{ }\PY{o}{=}\PY{+w}{ }\PY{n}{Arena}\PY{p}{.}\PY{n+na}{ofConfined}\PY{p}{(}\PY{p}{)}\PY{p}{)}\PY{+w}{ }\PY{p}{\PYZob{}}
\PY{+w}{            }\PY{n}{MemorySegment}\PY{+w}{ }\PY{n}{segment}\PY{+w}{ }\PY{o}{=}\PY{+w}{ }\PY{n}{arena}\PY{p}{.}\PY{n+na}{allocate}\PY{p}{(}\PY{n}{intArrayLayout}\PY{p}{)}\PY{p}{;}
\PY{+w}{            }\PY{k}{for}\PY{+w}{ }\PY{p}{(}\PY{k+kt}{int}\PY{+w}{ }\PY{n}{i}\PY{+w}{ }\PY{o}{=}\PY{+w}{ }\PY{l+m+mi}{0}\PY{p}{;}\PY{+w}{ }\PY{n}{i}\PY{+w}{ }\PY{o}{\PYZlt{}}\PY{+w}{ }\PY{l+m+mi}{4}\PY{p}{;}\PY{+w}{ }\PY{n}{i}\PY{o}{+}\PY{o}{+}\PY{p}{)}\PY{+w}{ }\PY{p}{\PYZob{}}
\PY{+w}{                }\PY{n}{elementHandle}\PY{p}{.}\PY{n+na}{set}\PY{p}{(}\PY{n}{segment}\PY{p}{,}\PY{+w}{ }\PY{l+m+mi}{0}\PY{n}{L}\PY{p}{,}\PY{+w}{ }\PY{p}{(}\PY{k+kt}{long}\PY{p}{)}\PY{+w}{ }\PY{n}{i}\PY{p}{,}\PY{+w}{ }\PY{n}{i}\PY{+w}{ }\PY{o}{*}\PY{+w}{ }\PY{l+m+mi}{10}\PY{p}{)}\PY{p}{;}\PY{+w}{ }\PY{c+c1}{// (segment, base offset, index, value)}
\PY{+w}{            }\PY{p}{\PYZcb{}}
\PY{+w}{            }\PY{k}{for}\PY{+w}{ }\PY{p}{(}\PY{k+kt}{int}\PY{+w}{ }\PY{n}{i}\PY{+w}{ }\PY{o}{=}\PY{+w}{ }\PY{l+m+mi}{0}\PY{p}{;}\PY{+w}{ }\PY{n}{i}\PY{+w}{ }\PY{o}{\PYZlt{}}\PY{+w}{ }\PY{l+m+mi}{4}\PY{p}{;}\PY{+w}{ }\PY{n}{i}\PY{o}{+}\PY{o}{+}\PY{p}{)}\PY{+w}{ }\PY{p}{\PYZob{}}
\PY{+w}{                }\PY{n}{System}\PY{p}{.}\PY{n+na}{out}\PY{p}{.}\PY{n+na}{println}\PY{p}{(}\PY{n}{elementHandle}\PY{p}{.}\PY{n+na}{get}\PY{p}{(}\PY{n}{segment}\PY{p}{,}\PY{+w}{ }\PY{l+m+mi}{0}\PY{n}{L}\PY{p}{,}\PY{+w}{ }\PY{p}{(}\PY{k+kt}{long}\PY{p}{)}\PY{+w}{ }\PY{n}{i}\PY{p}{)}\PY{p}{)}\PY{p}{;}
\PY{+w}{            }\PY{p}{\PYZcb{}}
\PY{+w}{        }\PY{p}{\PYZcb{}}
\PY{+w}{    }\PY{p}{\PYZcb{}}
\PY{p}{\PYZcb{}}
\end{Verbatim}
\end{codebox}

\code{VarHandle} access is how the FFM API lets you do fine-grained, layout-aware reads/writes (including atomic operations like \code{compareAndSet} on off-heap memory) without ever writing native code — the layout describes the shape once, and the handle gives you type-safe, bounds-checked access forever after.

\subsection[Calling a C Function: Linker, SymbolLookup, FunctionDescriptor]{Calling a C Function: \code{Linker}, \code{SymbolLookup}, \code{FunctionDescriptor}}
\label{h:17:calling-a-c-function-linker-symbollookup-functiondescriptor}
Three pieces cooperate to call a native function: \code{SymbolLookup} finds the function’s address, \code{FunctionDescriptor} describes its signature, and \code{Linker} produces a \code{MethodHandle} (a \textbf{downcall handle}) that invokes it.

\begin{codebox}
\begin{Verbatim}[commandchars=\\\{\}]
\PY{k+kn}{import}\PY{+w}{ }\PY{n+nn}{java.lang.foreign.Arena}\PY{p}{;}
\PY{k+kn}{import}\PY{+w}{ }\PY{n+nn}{java.lang.foreign.FunctionDescriptor}\PY{p}{;}
\PY{k+kn}{import}\PY{+w}{ }\PY{n+nn}{java.lang.foreign.Linker}\PY{p}{;}
\PY{k+kn}{import}\PY{+w}{ }\PY{n+nn}{java.lang.foreign.MemorySegment}\PY{p}{;}
\PY{k+kn}{import}\PY{+w}{ }\PY{n+nn}{java.lang.foreign.SymbolLookup}\PY{p}{;}
\PY{k+kn}{import}\PY{+w}{ }\PY{n+nn}{java.lang.foreign.ValueLayout}\PY{p}{;}
\PY{k+kn}{import}\PY{+w}{ }\PY{n+nn}{java.lang.invoke.MethodHandle}\PY{p}{;}

\PY{k+kd}{public}\PY{+w}{ }\PY{k+kd}{class} \PY{n+nc}{StrlenDemo}\PY{+w}{ }\PY{p}{\PYZob{}}
\PY{+w}{    }\PY{k+kd}{public}\PY{+w}{ }\PY{k+kd}{static}\PY{+w}{ }\PY{k+kt}{void}\PY{+w}{ }\PY{n+nf}{main}\PY{p}{(}\PY{n}{String}\PY{o}{[}\PY{o}{]}\PY{+w}{ }\PY{n}{args}\PY{p}{)}\PY{+w}{ }\PY{k+kd}{throws}\PY{+w}{ }\PY{n}{Throwable}\PY{+w}{ }\PY{p}{\PYZob{}}
\PY{+w}{        }\PY{n}{Linker}\PY{+w}{ }\PY{n}{linker}\PY{+w}{ }\PY{o}{=}\PY{+w}{ }\PY{n}{Linker}\PY{p}{.}\PY{n+na}{nativeLinker}\PY{p}{(}\PY{p}{)}\PY{p}{;}
\PY{+w}{        }\PY{n}{SymbolLookup}\PY{+w}{ }\PY{n}{stdlib}\PY{+w}{ }\PY{o}{=}\PY{+w}{ }\PY{n}{linker}\PY{p}{.}\PY{n+na}{defaultLookup}\PY{p}{(}\PY{p}{)}\PY{p}{;}\PY{+w}{ }\PY{c+c1}{// the process\PYZsq{}s standard C library}

\PY{+w}{        }\PY{c+c1}{// strlen(const char *s) \PYZhy{}\PYZgt{} size\PYZus{}t   (size\PYZus{}t maps to a native long on 64\PYZhy{}bit platforms)}
\PY{+w}{        }\PY{n}{MethodHandle}\PY{+w}{ }\PY{n}{strlen}\PY{+w}{ }\PY{o}{=}\PY{+w}{ }\PY{n}{linker}\PY{p}{.}\PY{n+na}{downcallHandle}\PY{p}{(}
\PY{+w}{            }\PY{n}{stdlib}\PY{p}{.}\PY{n+na}{find}\PY{p}{(}\PY{l+s}{\PYZdq{}}\PY{l+s}{strlen}\PY{l+s}{\PYZdq{}}\PY{p}{)}\PY{p}{.}\PY{n+na}{orElseThrow}\PY{p}{(}\PY{p}{(}\PY{p}{)}\PY{+w}{ }\PY{o}{\PYZhy{}}\PY{o}{\PYZgt{}}\PY{+w}{ }\PY{k}{new}\PY{+w}{ }\PY{n}{NoSuchElementException}\PY{p}{(}\PY{l+s}{\PYZdq{}}\PY{l+s}{strlen not found}\PY{l+s}{\PYZdq{}}\PY{p}{)}\PY{p}{)}\PY{p}{,}
\PY{+w}{            }\PY{n}{FunctionDescriptor}\PY{p}{.}\PY{n+na}{of}\PY{p}{(}\PY{n}{ValueLayout}\PY{p}{.}\PY{n+na}{JAVA\PYZus{}LONG}\PY{p}{,}\PY{+w}{ }\PY{n}{ValueLayout}\PY{p}{.}\PY{n+na}{ADDRESS}\PY{p}{)}
\PY{+w}{        }\PY{p}{)}\PY{p}{;}

\PY{+w}{        }\PY{k}{try}\PY{+w}{ }\PY{p}{(}\PY{n}{Arena}\PY{+w}{ }\PY{n}{arena}\PY{+w}{ }\PY{o}{=}\PY{+w}{ }\PY{n}{Arena}\PY{p}{.}\PY{n+na}{ofConfined}\PY{p}{(}\PY{p}{)}\PY{p}{)}\PY{+w}{ }\PY{p}{\PYZob{}}
\PY{+w}{            }\PY{n}{MemorySegment}\PY{+w}{ }\PY{n}{cString}\PY{+w}{ }\PY{o}{=}\PY{+w}{ }\PY{n}{arena}\PY{p}{.}\PY{n+na}{allocateUtf8String}\PY{p}{(}\PY{l+s}{\PYZdq{}}\PY{l+s}{Hello, native world!}\PY{l+s}{\PYZdq{}}\PY{p}{)}\PY{p}{;}
\PY{+w}{            }\PY{k+kt}{long}\PY{+w}{ }\PY{n}{length}\PY{+w}{ }\PY{o}{=}\PY{+w}{ }\PY{p}{(}\PY{k+kt}{long}\PY{p}{)}\PY{+w}{ }\PY{n}{strlen}\PY{p}{.}\PY{n+na}{invoke}\PY{p}{(}\PY{n}{cString}\PY{p}{)}\PY{p}{;}
\PY{+w}{            }\PY{n}{System}\PY{p}{.}\PY{n+na}{out}\PY{p}{.}\PY{n+na}{println}\PY{p}{(}\PY{l+s}{\PYZdq{}}\PY{l+s}{Length: }\PY{l+s}{\PYZdq{}}\PY{+w}{ }\PY{o}{+}\PY{+w}{ }\PY{n}{length}\PY{p}{)}\PY{p}{;}\PY{+w}{ }\PY{c+c1}{// 20}
\PY{+w}{        }\PY{p}{\PYZcb{}}
\PY{+w}{    }\PY{p}{\PYZcb{}}
\PY{p}{\PYZcb{}}
\end{Verbatim}
\end{codebox}

A no-argument function follows the same recipe with an empty argument list in the descriptor — \code{getpid()} (returns the OS process ID) needs only \code{Function\allowbreak{}Descriptor.\allowbreak{}of(\allowbreak{}Value\allowbreak{}Layout.\allowbreak{}JAVA\_\allowbreak{}INT)} and \code{getpid.\allowbreak{}invoke()} with no arguments. Running either of these requires the native-access flag described below, and both compile/run as ordinary Java — no \code{javac~\allowbreak{}-\allowbreak{}h}, no C compiler, no shared library build step, which is the headline improvement over JNI.

\subsection[Upcalls: Java Callback into C]{Upcalls: Java Callback into C}
\label{h:17:upcalls-java-callback-into-c}
An \textbf{upcall} is the reverse direction: native code calling back into Java, such as passing a Java method as a C comparator function to \code{qsort}.

\begin{codebox}
\begin{Verbatim}[commandchars=\\\{\}]
\PY{k+kn}{import}\PY{+w}{ }\PY{n+nn}{java.lang.foreign.*}\PY{p}{;}
\PY{k+kn}{import}\PY{+w}{ }\PY{n+nn}{java.lang.invoke.MethodHandle}\PY{p}{;}
\PY{k+kn}{import}\PY{+w}{ }\PY{n+nn}{java.lang.invoke.MethodHandles}\PY{p}{;}
\PY{k+kn}{import}\PY{+w}{ }\PY{n+nn}{java.lang.invoke.MethodType}\PY{p}{;}
\PY{k+kn}{import}\PY{+w}{ }\PY{n+nn}{java.util.Arrays}\PY{p}{;}

\PY{k+kd}{public}\PY{+w}{ }\PY{k+kd}{class} \PY{n+nc}{UpcallDemo}\PY{+w}{ }\PY{p}{\PYZob{}}
\PY{+w}{    }\PY{c+c1}{// This Java method will be exposed to native code as a C function pointer.}
\PY{+w}{    }\PY{k+kd}{static}\PY{+w}{ }\PY{k+kt}{int}\PY{+w}{ }\PY{n+nf}{compareInts}\PY{p}{(}\PY{n}{MemorySegment}\PY{+w}{ }\PY{n}{a}\PY{p}{,}\PY{+w}{ }\PY{n}{MemorySegment}\PY{+w}{ }\PY{n}{b}\PY{p}{)}\PY{+w}{ }\PY{p}{\PYZob{}}
\PY{+w}{        }\PY{k+kt}{int}\PY{+w}{ }\PY{n}{x}\PY{+w}{ }\PY{o}{=}\PY{+w}{ }\PY{n}{a}\PY{p}{.}\PY{n+na}{get}\PY{p}{(}\PY{n}{ValueLayout}\PY{p}{.}\PY{n+na}{JAVA\PYZus{}INT}\PY{p}{,}\PY{+w}{ }\PY{l+m+mi}{0}\PY{p}{)}\PY{p}{;}
\PY{+w}{        }\PY{k+kt}{int}\PY{+w}{ }\PY{n}{y}\PY{+w}{ }\PY{o}{=}\PY{+w}{ }\PY{n}{b}\PY{p}{.}\PY{n+na}{get}\PY{p}{(}\PY{n}{ValueLayout}\PY{p}{.}\PY{n+na}{JAVA\PYZus{}INT}\PY{p}{,}\PY{+w}{ }\PY{l+m+mi}{0}\PY{p}{)}\PY{p}{;}
\PY{+w}{        }\PY{k}{return}\PY{+w}{ }\PY{n}{Integer}\PY{p}{.}\PY{n+na}{compare}\PY{p}{(}\PY{n}{x}\PY{p}{,}\PY{+w}{ }\PY{n}{y}\PY{p}{)}\PY{p}{;}
\PY{+w}{    }\PY{p}{\PYZcb{}}

\PY{+w}{    }\PY{k+kd}{public}\PY{+w}{ }\PY{k+kd}{static}\PY{+w}{ }\PY{k+kt}{void}\PY{+w}{ }\PY{n+nf}{main}\PY{p}{(}\PY{n}{String}\PY{o}{[}\PY{o}{]}\PY{+w}{ }\PY{n}{args}\PY{p}{)}\PY{+w}{ }\PY{k+kd}{throws}\PY{+w}{ }\PY{n}{Throwable}\PY{+w}{ }\PY{p}{\PYZob{}}
\PY{+w}{        }\PY{n}{Linker}\PY{+w}{ }\PY{n}{linker}\PY{+w}{ }\PY{o}{=}\PY{+w}{ }\PY{n}{Linker}\PY{p}{.}\PY{n+na}{nativeLinker}\PY{p}{(}\PY{p}{)}\PY{p}{;}
\PY{+w}{        }\PY{n}{MethodHandle}\PY{+w}{ }\PY{n}{comparatorHandle}\PY{+w}{ }\PY{o}{=}\PY{+w}{ }\PY{n}{MethodHandles}\PY{p}{.}\PY{n+na}{lookup}\PY{p}{(}\PY{p}{)}\PY{p}{.}\PY{n+na}{findStatic}\PY{p}{(}
\PY{+w}{            }\PY{n}{UpcallDemo}\PY{p}{.}\PY{n+na}{class}\PY{p}{,}\PY{+w}{ }\PY{l+s}{\PYZdq{}}\PY{l+s}{compareInts}\PY{l+s}{\PYZdq{}}\PY{p}{,}
\PY{+w}{            }\PY{n}{MethodType}\PY{p}{.}\PY{n+na}{methodType}\PY{p}{(}\PY{k+kt}{int}\PY{p}{.}\PY{n+na}{class}\PY{p}{,}\PY{+w}{ }\PY{n}{MemorySegment}\PY{p}{.}\PY{n+na}{class}\PY{p}{,}\PY{+w}{ }\PY{n}{MemorySegment}\PY{p}{.}\PY{n+na}{class}\PY{p}{)}\PY{p}{)}\PY{p}{;}

\PY{+w}{        }\PY{k}{try}\PY{+w}{ }\PY{p}{(}\PY{n}{Arena}\PY{+w}{ }\PY{n}{arena}\PY{+w}{ }\PY{o}{=}\PY{+w}{ }\PY{n}{Arena}\PY{p}{.}\PY{n+na}{ofConfined}\PY{p}{(}\PY{p}{)}\PY{p}{)}\PY{+w}{ }\PY{p}{\PYZob{}}
\PY{+w}{            }\PY{c+c1}{// Wrap the Java method as a native function pointer (\PYZdq{}upcall stub\PYZdq{}).}
\PY{+w}{            }\PY{n}{MemorySegment}\PY{+w}{ }\PY{n}{comparatorStub}\PY{+w}{ }\PY{o}{=}\PY{+w}{ }\PY{n}{linker}\PY{p}{.}\PY{n+na}{upcallStub}\PY{p}{(}
\PY{+w}{                }\PY{n}{comparatorHandle}\PY{p}{,}
\PY{+w}{                }\PY{n}{FunctionDescriptor}\PY{p}{.}\PY{n+na}{of}\PY{p}{(}\PY{n}{ValueLayout}\PY{p}{.}\PY{n+na}{JAVA\PYZus{}INT}\PY{p}{,}\PY{+w}{ }\PY{n}{ValueLayout}\PY{p}{.}\PY{n+na}{ADDRESS}\PY{p}{,}\PY{+w}{ }\PY{n}{ValueLayout}\PY{p}{.}\PY{n+na}{ADDRESS}\PY{p}{)}\PY{p}{,}
\PY{+w}{                }\PY{n}{arena}\PY{p}{)}\PY{p}{;}

\PY{+w}{            }\PY{n}{MethodHandle}\PY{+w}{ }\PY{n}{qsort}\PY{+w}{ }\PY{o}{=}\PY{+w}{ }\PY{n}{linker}\PY{p}{.}\PY{n+na}{downcallHandle}\PY{p}{(}
\PY{+w}{                }\PY{n}{linker}\PY{p}{.}\PY{n+na}{defaultLookup}\PY{p}{(}\PY{p}{)}\PY{p}{.}\PY{n+na}{find}\PY{p}{(}\PY{l+s}{\PYZdq{}}\PY{l+s}{qsort}\PY{l+s}{\PYZdq{}}\PY{p}{)}\PY{p}{.}\PY{n+na}{orElseThrow}\PY{p}{(}\PY{p}{)}\PY{p}{,}
\PY{+w}{                }\PY{n}{FunctionDescriptor}\PY{p}{.}\PY{n+na}{ofVoid}\PY{p}{(}\PY{n}{ValueLayout}\PY{p}{.}\PY{n+na}{ADDRESS}\PY{p}{,}\PY{+w}{ }\PY{n}{ValueLayout}\PY{p}{.}\PY{n+na}{JAVA\PYZus{}LONG}\PY{p}{,}
\PY{+w}{                                          }\PY{n}{ValueLayout}\PY{p}{.}\PY{n+na}{JAVA\PYZus{}LONG}\PY{p}{,}\PY{+w}{ }\PY{n}{ValueLayout}\PY{p}{.}\PY{n+na}{ADDRESS}\PY{p}{)}\PY{p}{)}\PY{p}{;}

\PY{+w}{            }\PY{n}{MemorySegment}\PY{+w}{ }\PY{n}{array}\PY{+w}{ }\PY{o}{=}\PY{+w}{ }\PY{n}{arena}\PY{p}{.}\PY{n+na}{allocateArray}\PY{p}{(}\PY{n}{ValueLayout}\PY{p}{.}\PY{n+na}{JAVA\PYZus{}INT}\PY{p}{,}\PY{+w}{ }\PY{l+m+mi}{5}\PY{p}{,}\PY{+w}{ }\PY{l+m+mi}{3}\PY{p}{,}\PY{+w}{ }\PY{l+m+mi}{1}\PY{p}{,}\PY{+w}{ }\PY{l+m+mi}{4}\PY{p}{,}\PY{+w}{ }\PY{l+m+mi}{2}\PY{p}{)}\PY{p}{;}
\PY{+w}{            }\PY{n}{qsort}\PY{p}{.}\PY{n+na}{invoke}\PY{p}{(}\PY{n}{array}\PY{p}{,}\PY{+w}{ }\PY{l+m+mi}{5L}\PY{p}{,}\PY{+w}{ }\PY{n}{ValueLayout}\PY{p}{.}\PY{n+na}{JAVA\PYZus{}INT}\PY{p}{.}\PY{n+na}{byteSize}\PY{p}{(}\PY{p}{)}\PY{p}{,}\PY{+w}{ }\PY{n}{comparatorStub}\PY{p}{)}\PY{p}{;}

\PY{+w}{            }\PY{k+kt}{int}\PY{o}{[}\PY{o}{]}\PY{+w}{ }\PY{n}{sorted}\PY{+w}{ }\PY{o}{=}\PY{+w}{ }\PY{n}{array}\PY{p}{.}\PY{n+na}{toArray}\PY{p}{(}\PY{n}{ValueLayout}\PY{p}{.}\PY{n+na}{JAVA\PYZus{}INT}\PY{p}{)}\PY{p}{;}
\PY{+w}{            }\PY{n}{System}\PY{p}{.}\PY{n+na}{out}\PY{p}{.}\PY{n+na}{println}\PY{p}{(}\PY{n}{Arrays}\PY{p}{.}\PY{n+na}{toString}\PY{p}{(}\PY{n}{sorted}\PY{p}{)}\PY{p}{)}\PY{p}{;}\PY{+w}{ }\PY{c+c1}{// [1, 2, 3, 4, 5]}
\PY{+w}{        }\PY{p}{\PYZcb{}}
\PY{+w}{    }\PY{p}{\PYZcb{}}
\PY{p}{\PYZcb{}}
\end{Verbatim}
\end{codebox}

The upcall stub (\code{comparatorStub}) is itself a \code{MemorySegment} — a real native function pointer that C code (\code{qsort}, in this case) can call directly, entirely unaware it’s actually calling back into the JVM.

\subsection[Struct Layouts]{Struct Layouts}
\label{h:17:struct-layouts}
\code{Memory\allowbreak{}Layout.\allowbreak{}struct\allowbreak{}Layout(...)} describes a C \code{struct} shape, with named fields accessed via \code{VarHandle}s derived from path elements — no manual offset arithmetic required.

\begin{codebox}
\begin{Verbatim}[commandchars=\\\{\}]
\PY{k+kn}{import}\PY{+w}{ }\PY{n+nn}{java.lang.foreign.*}\PY{p}{;}
\PY{k+kn}{import}\PY{+w}{ }\PY{n+nn}{java.lang.invoke.VarHandle}\PY{p}{;}

\PY{k+kd}{public}\PY{+w}{ }\PY{k+kd}{class} \PY{n+nc}{StructDemo}\PY{+w}{ }\PY{p}{\PYZob{}}
\PY{+w}{    }\PY{k+kd}{public}\PY{+w}{ }\PY{k+kd}{static}\PY{+w}{ }\PY{k+kt}{void}\PY{+w}{ }\PY{n+nf}{main}\PY{p}{(}\PY{n}{String}\PY{o}{[}\PY{o}{]}\PY{+w}{ }\PY{n}{args}\PY{p}{)}\PY{+w}{ }\PY{p}{\PYZob{}}
\PY{+w}{        }\PY{c+c1}{// Equivalent to: struct Point \PYZob{} int x; int y; \PYZcb{};}
\PY{+w}{        }\PY{n}{StructLayout}\PY{+w}{ }\PY{n}{pointLayout}\PY{+w}{ }\PY{o}{=}\PY{+w}{ }\PY{n}{MemoryLayout}\PY{p}{.}\PY{n+na}{structLayout}\PY{p}{(}
\PY{+w}{            }\PY{n}{ValueLayout}\PY{p}{.}\PY{n+na}{JAVA\PYZus{}INT}\PY{p}{.}\PY{n+na}{withName}\PY{p}{(}\PY{l+s}{\PYZdq{}}\PY{l+s}{x}\PY{l+s}{\PYZdq{}}\PY{p}{)}\PY{p}{,}
\PY{+w}{            }\PY{n}{ValueLayout}\PY{p}{.}\PY{n+na}{JAVA\PYZus{}INT}\PY{p}{.}\PY{n+na}{withName}\PY{p}{(}\PY{l+s}{\PYZdq{}}\PY{l+s}{y}\PY{l+s}{\PYZdq{}}\PY{p}{)}
\PY{+w}{        }\PY{p}{)}\PY{p}{;}

\PY{+w}{        }\PY{n}{VarHandle}\PY{+w}{ }\PY{n}{xHandle}\PY{+w}{ }\PY{o}{=}\PY{+w}{ }\PY{n}{pointLayout}\PY{p}{.}\PY{n+na}{varHandle}\PY{p}{(}\PY{n}{MemoryLayout}\PY{p}{.}\PY{n+na}{PathElement}\PY{p}{.}\PY{n+na}{groupElement}\PY{p}{(}\PY{l+s}{\PYZdq{}}\PY{l+s}{x}\PY{l+s}{\PYZdq{}}\PY{p}{)}\PY{p}{)}\PY{p}{;}
\PY{+w}{        }\PY{n}{VarHandle}\PY{+w}{ }\PY{n}{yHandle}\PY{+w}{ }\PY{o}{=}\PY{+w}{ }\PY{n}{pointLayout}\PY{p}{.}\PY{n+na}{varHandle}\PY{p}{(}\PY{n}{MemoryLayout}\PY{p}{.}\PY{n+na}{PathElement}\PY{p}{.}\PY{n+na}{groupElement}\PY{p}{(}\PY{l+s}{\PYZdq{}}\PY{l+s}{y}\PY{l+s}{\PYZdq{}}\PY{p}{)}\PY{p}{)}\PY{p}{;}

\PY{+w}{        }\PY{k}{try}\PY{+w}{ }\PY{p}{(}\PY{n}{Arena}\PY{+w}{ }\PY{n}{arena}\PY{+w}{ }\PY{o}{=}\PY{+w}{ }\PY{n}{Arena}\PY{p}{.}\PY{n+na}{ofConfined}\PY{p}{(}\PY{p}{)}\PY{p}{)}\PY{+w}{ }\PY{p}{\PYZob{}}
\PY{+w}{            }\PY{n}{MemorySegment}\PY{+w}{ }\PY{n}{point}\PY{+w}{ }\PY{o}{=}\PY{+w}{ }\PY{n}{arena}\PY{p}{.}\PY{n+na}{allocate}\PY{p}{(}\PY{n}{pointLayout}\PY{p}{)}\PY{p}{;}
\PY{+w}{            }\PY{n}{xHandle}\PY{p}{.}\PY{n+na}{set}\PY{p}{(}\PY{n}{point}\PY{p}{,}\PY{+w}{ }\PY{l+m+mi}{0}\PY{n}{L}\PY{p}{,}\PY{+w}{ }\PY{l+m+mi}{10}\PY{p}{)}\PY{p}{;}
\PY{+w}{            }\PY{n}{yHandle}\PY{p}{.}\PY{n+na}{set}\PY{p}{(}\PY{n}{point}\PY{p}{,}\PY{+w}{ }\PY{l+m+mi}{0}\PY{n}{L}\PY{p}{,}\PY{+w}{ }\PY{l+m+mi}{20}\PY{p}{)}\PY{p}{;}

\PY{+w}{            }\PY{n}{System}\PY{p}{.}\PY{n+na}{out}\PY{p}{.}\PY{n+na}{println}\PY{p}{(}\PY{l+s}{\PYZdq{}}\PY{l+s}{x=}\PY{l+s}{\PYZdq{}}\PY{+w}{ }\PY{o}{+}\PY{+w}{ }\PY{n}{xHandle}\PY{p}{.}\PY{n+na}{get}\PY{p}{(}\PY{n}{point}\PY{p}{,}\PY{+w}{ }\PY{l+m+mi}{0}\PY{n}{L}\PY{p}{)}\PY{+w}{ }\PY{o}{+}\PY{+w}{ }\PY{l+s}{\PYZdq{}}\PY{l+s}{, y=}\PY{l+s}{\PYZdq{}}\PY{+w}{ }\PY{o}{+}\PY{+w}{ }\PY{n}{yHandle}\PY{p}{.}\PY{n+na}{get}\PY{p}{(}\PY{n}{point}\PY{p}{,}\PY{+w}{ }\PY{l+m+mi}{0}\PY{n}{L}\PY{p}{)}\PY{p}{)}\PY{p}{;}
\PY{+w}{        }\PY{p}{\PYZcb{}}
\PY{+w}{    }\PY{p}{\PYZcb{}}
\PY{p}{\PYZcb{}}
\end{Verbatim}
\end{codebox}

\code{pointLayout.\allowbreak{}byteSize()} and \code{point\allowbreak{}Layout.\allowbreak{}byte\allowbreak{}Offset(\allowbreak{}Path\allowbreak{}Element.\allowbreak{}group\allowbreak{}Element(\allowbreak{}\textquotedbl{}y\textquotedbl{}))} are computed automatically from the declared fields (respecting natural alignment/padding rules), which is exactly the error-prone bookkeeping JNI forced you to do by hand in C headers.

\subsection[jextract]{\code{jextract}}
\label{h:17:jextract}
Writing every \code{FunctionDescriptor} and struct layout by hand for a large C library (say, all of \code{libcurl} or \code{sqlite3}) would be tedious and error-prone. \textbf{\code{jextract}} is a separate tool (distributed alongside OpenJDK builds, not part of the JDK proper) that parses a C header file and mechanically generates the Java FFM bindings — \code{MethodHandle}s, layouts, and constants — for every declared function and struct.

\begin{codebox}
\begin{Verbatim}[commandchars=\\\{\}]
\PY{c+c1}{\PYZsh{} jextract reads a C header and generates ready\PYZhy{}to\PYZhy{}use Java FFM bindings,}
\PY{c+c1}{\PYZsh{} so you don\PYZsq{}t hand\PYZhy{}write FunctionDescriptors for every function in sqlite3.h.}
jextract\PY{+w}{ }\PYZhy{}\PYZhy{}output\PY{+w}{ }src/generated\PY{+w}{ }\PY{l+s+se}{\PYZbs{}}
\PY{+w}{         }\PYZhy{}\PYZhy{}target\PYZhy{}package\PY{+w}{ }com.teampicnic.native\PYZus{}.sqlite\PY{+w}{ }\PY{l+s+se}{\PYZbs{}}
\PY{+w}{         }/usr/include/sqlite3.h
\end{Verbatim}
\end{codebox}

The generated code is ordinary Java calling the same \code{Linker}/\code{MemorySegment} machinery shown above — \code{jextract} is a code generator for boilerplate, not a different runtime mechanism.

\subsection[--enable-native-access and Restricted Methods]{\code{--\allowbreak{}enable-\allowbreak{}native-\allowbreak{}access} and Restricted Methods}
\label{h:17:--enable-native-access-and-restricted-methods}
Native memory access and native calls are \textbf{restricted methods}: operations the JVM cannot fully verify are memory-safe, because they cross into native code or raw memory. Since JDK 22, using them prints a runtime warning unless the module (or the unnamed module, for classpath code) is explicitly granted native access; \textbf{JDK 24 tightens this further} — by default, restricted native operations throw rather than merely warn, unless access is explicitly enabled.

\begin{codebox}
\begin{Verbatim}[commandchars=\\\{\}]
\PY{c+c1}{\PYZsh{} JDK 22/23: without this flag, restricted FFM operations print a warning to stderr:}
\PY{c+c1}{\PYZsh{}   WARNING: A restricted method in java.lang.foreign.SymbolLookup has been called}
\PY{c+c1}{\PYZsh{} JDK 24+: without this flag, restricted operations THROW at the call site instead.}
java\PY{+w}{ }\PYZhy{}\PYZhy{}enable\PYZhy{}native\PYZhy{}access\PY{o}{=}com.teampicnic.app\PY{+w}{ }\PYZhy{}p\PY{+w}{ }out\PY{+w}{ }\PYZhy{}m\PY{+w}{ }com.teampicnic.app/com.teampicnic.app.Main

\PY{c+c1}{\PYZsh{} For classpath (unnamed\PYZhy{}module) code:}
java\PY{+w}{ }\PYZhy{}\PYZhy{}enable\PYZhy{}native\PYZhy{}access\PY{o}{=}ALL\PYZhy{}UNNAMED\PY{+w}{ }\PYZhy{}cp\PY{+w}{ }out\PY{+w}{ }com.teampicnic.app.Main
\end{Verbatim}
\end{codebox}

This mirrors the philosophy of \code{--\allowbreak{}add-\allowbreak{}opens}/\code{--\allowbreak{}add-\allowbreak{}exports} for JPMS internals: native access is powerful enough to corrupt memory or crash the JVM outright (an off-by-one \code{MemorySegment} write is still possible if you compute an out-of-bounds \emph{but still confined-looking} offset, and an entirely wrong \code{FunctionDescriptor} for a real native function is undefined behavior), so the JVM requires an explicit, auditable opt-in rather than allowing it silently by default.

\subsection[FFM vs JNI Comparison Table]{FFM vs JNI Comparison Table}
\label{h:17:ffm-vs-jni-comparison-table}
{\tablefontsmall
\begin{xltabular}{\linewidth}{@{}L{0.466}L{0.983}L{1.552}@{}}
\toprule
\rowcolor{tablehdr}
\thd{} & \thd{JNI} & \thd{Foreign Function \& Memory API} \\
\midrule
\endhead
Requires writing C/C++ glue code & Yes (\code{.\allowbreak{}c}/\code{.\allowbreak{}cpp} files, \code{javah}/\code{javac~\allowbreak{}-\allowbreak{}h} header generation) & No — pure Java \\
Requires a native compiler toolchain & Yes & No \\
Memory safety & None — raw pointers, manual bookkeeping, easy to crash the JVM & Bounds-checked \code{MemorySegment}s, arena-scoped lifetimes, \code{IllegalStateException} instead of memory corruption on misuse \\
Deallocation model & Manual, via native code (\code{free}) or \code{finalize()}-era hacks & Deterministic (\code{Arena.\allowbreak{}close()}/try-with-resources), GC-tied (\code{ofAuto}), or process-lifetime (\code{global}) \\
Calling native functions & Requires a pre-built, pre-loaded shared library exposing JNI-specific function signatures & Directly calls existing native libraries (\code{libc}, \code{libsqlite3}, etc.) with no JNI-specific wrapper needed \\
Callbacks into Java & Supported, but verbose and error-prone (\code{JNIEnv*} juggling) & \code{Linker.\allowbreak{}upcallStub} — a first-class, type-checked Java construct \\
Struct/layout description & Manual offset math in C headers, kept in sync by hand & \code{MemoryLayout}/\code{StructLayout}, computed automatically \\
Tooling for large libraries & Hand-written wrappers, or third-party generators & \code{jextract} generates bindings from C headers \\
Finalized in & Java 1.1 (1997) & JDK 22 (2024), JEP 454 \\
Governing flag & None (implicit trust) & \code{--\allowbreak{}enable-\allowbreak{}native-\allowbreak{}access}, restricted-method warnings/errors \\
\bottomrule
\end{xltabular}
}

\section[Vector API (Incubator)]{Vector API (Incubator)}
\label{h:17:vector-api-incubator}
The Vector API (\code{jdk.\allowbreak{}incubator.\allowbreak{}vector}) provides a portable way to express \textbf{SIMD} (Single Instruction, Multiple Data) computations in pure Java, letting the JIT compile them down to CPU vector instructions (AVX, SVE, NEON, …) instead of relying entirely on auto-vectorization heuristics.

\subsection[Why It’s Still Incubating]{Why It’s Still Incubating}
\label{h:17:why-its-still-incubating}
The Vector API has been re-incubated in \textbf{every JDK release since JDK 16} — as of \textbf{JDK 25, it is on its 10th incubation round} — because its final shape depends on language and JVM features that are not yet finished, primarily \textbf{Project Valhalla}’s value types. An incubator module (JEP 11) is the JDK’s formal mechanism for shipping an API that is real and usable, but explicitly \emph{not yet a permanent commitment}: its package name, method signatures, and semantics can still change between releases without the usual deprecation cycle.

{\tablefont
\begin{xltabular}{\linewidth}{@{}L{0.848}L{1.152}@{}}
\toprule
\rowcolor{tablehdr}
\thd{Signal} & \thd{What it tells a reviewer} \\
\midrule
\endhead
Package is \code{jdk.\allowbreak{}incubator.\allowbreak{}vector}, not \code{java.\allowbreak{}util.\allowbreak{}vector} or similar & Not a supported, permanent API yet \\
Requires \code{--\allowbreak{}add-\allowbreak{}modules~\allowbreak{}jdk.\allowbreak{}incubator.\allowbreak{}vector} explicitly & The JDK itself is telling you “opt in, at your own risk” \\
Re-incubated \textasciitilde{}10 times across \textasciitilde{}10 releases & The API shape is still actively changing release to release \\
No compatibility guarantee across JDK versions & Code compiled against one JDK’s incubator module may not run unchanged against the next \\
\bottomrule
\end{xltabular}
}

\subsection[VectorSpecies, Lanes, and Masks]{\code{VectorSpecies}, Lanes, and Masks}
\label{h:17:vectorspecies-lanes-and-masks}
A \code{VectorSpecies\textless{}\allowbreak{}E\textgreater{}} describes a concrete vector “shape”: the element type and the number of \textbf{lanes} (elements processed per vector operation), which the JVM picks based on the actual CPU’s hardware vector width at runtime.

\begin{codebox}
\begin{Verbatim}[commandchars=\\\{\}]
\PY{k+kn}{import}\PY{+w}{ }\PY{n+nn}{jdk.incubator.vector.FloatVector}\PY{p}{;}
\PY{k+kn}{import}\PY{+w}{ }\PY{n+nn}{jdk.incubator.vector.VectorSpecies}\PY{p}{;}

\PY{k+kd}{public}\PY{+w}{ }\PY{k+kd}{class} \PY{n+nc}{SpeciesDemo}\PY{+w}{ }\PY{p}{\PYZob{}}
\PY{+w}{    }\PY{k+kd}{public}\PY{+w}{ }\PY{k+kd}{static}\PY{+w}{ }\PY{k+kt}{void}\PY{+w}{ }\PY{n+nf}{main}\PY{p}{(}\PY{n}{String}\PY{o}{[}\PY{o}{]}\PY{+w}{ }\PY{n}{args}\PY{p}{)}\PY{+w}{ }\PY{p}{\PYZob{}}
\PY{+w}{        }\PY{c+c1}{// SPECIES\PYZus{}PREFERRED picks the widest vector shape the current CPU supports}
\PY{+w}{        }\PY{c+c1}{// for float (e.g., 8 lanes on a 256\PYZhy{}bit AVX2 machine, 16 on 512\PYZhy{}bit AVX\PYZhy{}512).}
\PY{+w}{        }\PY{n}{VectorSpecies}\PY{o}{\PYZlt{}}\PY{n}{Float}\PY{o}{\PYZgt{}}\PY{+w}{ }\PY{n}{species}\PY{+w}{ }\PY{o}{=}\PY{+w}{ }\PY{n}{FloatVector}\PY{p}{.}\PY{n+na}{SPECIES\PYZus{}PREFERRED}\PY{p}{;}
\PY{+w}{        }\PY{n}{System}\PY{p}{.}\PY{n+na}{out}\PY{p}{.}\PY{n+na}{println}\PY{p}{(}\PY{l+s}{\PYZdq{}}\PY{l+s}{Lanes: }\PY{l+s}{\PYZdq{}}\PY{+w}{ }\PY{o}{+}\PY{+w}{ }\PY{n}{species}\PY{p}{.}\PY{n+na}{length}\PY{p}{(}\PY{p}{)}\PY{p}{)}\PY{p}{;}
\PY{+w}{    }\PY{p}{\PYZcb{}}
\PY{p}{\PYZcb{}}
\end{Verbatim}
\end{codebox}

A \textbf{mask} is a per-lane boolean that selectively enables or disables lanes for an operation — essential for handling array lengths that aren’t an exact multiple of the vector width, or for conditional (branchless) computation.

\begin{codebox}
\begin{Verbatim}[commandchars=\\\{\}]
\PY{k+kn}{import}\PY{+w}{ }\PY{n+nn}{jdk.incubator.vector.FloatVector}\PY{p}{;}
\PY{k+kn}{import}\PY{+w}{ }\PY{n+nn}{jdk.incubator.vector.VectorMask}\PY{p}{;}
\PY{k+kn}{import}\PY{+w}{ }\PY{n+nn}{jdk.incubator.vector.VectorSpecies}\PY{p}{;}

\PY{k+kd}{public}\PY{+w}{ }\PY{k+kd}{class} \PY{n+nc}{MaskDemo}\PY{+w}{ }\PY{p}{\PYZob{}}
\PY{+w}{    }\PY{k+kd}{static}\PY{+w}{ }\PY{k+kd}{final}\PY{+w}{ }\PY{n}{VectorSpecies}\PY{o}{\PYZlt{}}\PY{n}{Float}\PY{o}{\PYZgt{}}\PY{+w}{ }\PY{n}{SPECIES}\PY{+w}{ }\PY{o}{=}\PY{+w}{ }\PY{n}{FloatVector}\PY{p}{.}\PY{n+na}{SPECIES\PYZus{}PREFERRED}\PY{p}{;}

\PY{+w}{    }\PY{k+kd}{static}\PY{+w}{ }\PY{k+kt}{void}\PY{+w}{ }\PY{n+nf}{clampNegativesToZero}\PY{p}{(}\PY{k+kt}{float}\PY{o}{[}\PY{o}{]}\PY{+w}{ }\PY{n}{data}\PY{p}{)}\PY{+w}{ }\PY{p}{\PYZob{}}
\PY{+w}{        }\PY{k+kt}{int}\PY{+w}{ }\PY{n}{i}\PY{+w}{ }\PY{o}{=}\PY{+w}{ }\PY{l+m+mi}{0}\PY{p}{;}
\PY{+w}{        }\PY{k+kt}{int}\PY{+w}{ }\PY{n}{upperBound}\PY{+w}{ }\PY{o}{=}\PY{+w}{ }\PY{n}{SPECIES}\PY{p}{.}\PY{n+na}{loopBound}\PY{p}{(}\PY{n}{data}\PY{p}{.}\PY{n+na}{length}\PY{p}{)}\PY{p}{;}
\PY{+w}{        }\PY{k}{for}\PY{+w}{ }\PY{p}{(}\PY{p}{;}\PY{+w}{ }\PY{n}{i}\PY{+w}{ }\PY{o}{\PYZlt{}}\PY{+w}{ }\PY{n}{upperBound}\PY{p}{;}\PY{+w}{ }\PY{n}{i}\PY{+w}{ }\PY{o}{+}\PY{o}{=}\PY{+w}{ }\PY{n}{SPECIES}\PY{p}{.}\PY{n+na}{length}\PY{p}{(}\PY{p}{)}\PY{p}{)}\PY{+w}{ }\PY{p}{\PYZob{}}
\PY{+w}{            }\PY{k+kd}{var}\PY{+w}{ }\PY{n}{v}\PY{+w}{ }\PY{o}{=}\PY{+w}{ }\PY{n}{FloatVector}\PY{p}{.}\PY{n+na}{fromArray}\PY{p}{(}\PY{n}{SPECIES}\PY{p}{,}\PY{+w}{ }\PY{n}{data}\PY{p}{,}\PY{+w}{ }\PY{n}{i}\PY{p}{)}\PY{p}{;}
\PY{+w}{            }\PY{n}{VectorMask}\PY{o}{\PYZlt{}}\PY{n}{Float}\PY{o}{\PYZgt{}}\PY{+w}{ }\PY{n}{negative}\PY{+w}{ }\PY{o}{=}\PY{+w}{ }\PY{n}{v}\PY{p}{.}\PY{n+na}{compare}\PY{p}{(}\PY{n}{jdk}\PY{p}{.}\PY{n+na}{incubator}\PY{p}{.}\PY{n+na}{vector}\PY{p}{.}\PY{n+na}{VectorOperators}\PY{p}{.}\PY{n+na}{LT}\PY{p}{,}\PY{+w}{ }\PY{l+m+mf}{0f}\PY{p}{)}\PY{p}{;}
\PY{+w}{            }\PY{n}{v}\PY{p}{.}\PY{n+na}{blend}\PY{p}{(}\PY{l+m+mf}{0f}\PY{p}{,}\PY{+w}{ }\PY{n}{negative}\PY{p}{)}\PY{p}{.}\PY{n+na}{intoArray}\PY{p}{(}\PY{n}{data}\PY{p}{,}\PY{+w}{ }\PY{n}{i}\PY{p}{)}\PY{p}{;}\PY{+w}{ }\PY{c+c1}{// lanes matching the mask become 0}
\PY{+w}{        }\PY{p}{\PYZcb{}}
\PY{+w}{        }\PY{k}{for}\PY{+w}{ }\PY{p}{(}\PY{p}{;}\PY{+w}{ }\PY{n}{i}\PY{+w}{ }\PY{o}{\PYZlt{}}\PY{+w}{ }\PY{n}{data}\PY{p}{.}\PY{n+na}{length}\PY{p}{;}\PY{+w}{ }\PY{n}{i}\PY{o}{+}\PY{o}{+}\PY{p}{)}\PY{+w}{ }\PY{p}{\PYZob{}}\PY{+w}{ }\PY{c+c1}{// scalar tail}
\PY{+w}{            }\PY{k}{if}\PY{+w}{ }\PY{p}{(}\PY{n}{data}\PY{o}{[}\PY{n}{i}\PY{o}{]}\PY{+w}{ }\PY{o}{\PYZlt{}}\PY{+w}{ }\PY{l+m+mf}{0f}\PY{p}{)}\PY{+w}{ }\PY{n}{data}\PY{o}{[}\PY{n}{i}\PY{o}{]}\PY{+w}{ }\PY{o}{=}\PY{+w}{ }\PY{l+m+mf}{0f}\PY{p}{;}
\PY{+w}{        }\PY{p}{\PYZcb{}}
\PY{+w}{    }\PY{p}{\PYZcb{}}
\PY{p}{\PYZcb{}}
\end{Verbatim}
\end{codebox}

\subsection[Enabling the Incubator Module]{Enabling the Incubator Module}
\label{h:17:enabling-the-incubator-module}
Incubator modules must be explicitly requested at both compile time and run time; they are never on by default.

\begin{codebox}
\begin{Verbatim}[commandchars=\\\{\}]
javac\PY{+w}{ }\PYZhy{}\PYZhy{}add\PYZhy{}modules\PY{+w}{ }jdk.incubator.vector\PY{+w}{ }\PYZhy{}\PYZhy{}release\PY{+w}{ }\PY{l+m}{21}\PY{+w}{ }VectorDotProduct.java
java\PY{+w}{  }\PYZhy{}\PYZhy{}add\PYZhy{}modules\PY{+w}{ }jdk.incubator.vector\PY{+w}{ }VectorDotProduct
\end{Verbatim}
\end{codebox}

Omitting \code{--\allowbreak{}add-\allowbreak{}modules~\allowbreak{}jdk.\allowbreak{}incubator.\allowbreak{}vector} produces a compile error (\code{package~\allowbreak{}jdk.\allowbreak{}incubator.\allowbreak{}vector~\allowbreak{}is~\allowbreak{}not~\allowbreak{}visible}) — this is the same JPMS accessibility machinery from earlier in the chapter, deliberately used here to make incubating APIs opt-in and impossible to depend on by accident.

\subsection[Worked Example: Vectorized Dot Product vs Scalar]{Worked Example: Vectorized Dot Product vs Scalar}
\label{h:17:worked-example-vectorized-dot-product-vs-scalar}
\begin{codebox}
\begin{Verbatim}[commandchars=\\\{\}]
\PY{c+c1}{// Scalar baseline: processes one float pair per loop iteration.}
\PY{k+kd}{public}\PY{+w}{ }\PY{k+kd}{class} \PY{n+nc}{DotProductScalar}\PY{+w}{ }\PY{p}{\PYZob{}}
\PY{+w}{    }\PY{k+kd}{static}\PY{+w}{ }\PY{k+kt}{float}\PY{+w}{ }\PY{n+nf}{dotProduct}\PY{p}{(}\PY{k+kt}{float}\PY{o}{[}\PY{o}{]}\PY{+w}{ }\PY{n}{a}\PY{p}{,}\PY{+w}{ }\PY{k+kt}{float}\PY{o}{[}\PY{o}{]}\PY{+w}{ }\PY{n}{b}\PY{p}{)}\PY{+w}{ }\PY{p}{\PYZob{}}
\PY{+w}{        }\PY{k+kt}{float}\PY{+w}{ }\PY{n}{sum}\PY{+w}{ }\PY{o}{=}\PY{+w}{ }\PY{l+m+mf}{0f}\PY{p}{;}
\PY{+w}{        }\PY{k}{for}\PY{+w}{ }\PY{p}{(}\PY{k+kt}{int}\PY{+w}{ }\PY{n}{i}\PY{+w}{ }\PY{o}{=}\PY{+w}{ }\PY{l+m+mi}{0}\PY{p}{;}\PY{+w}{ }\PY{n}{i}\PY{+w}{ }\PY{o}{\PYZlt{}}\PY{+w}{ }\PY{n}{a}\PY{p}{.}\PY{n+na}{length}\PY{p}{;}\PY{+w}{ }\PY{n}{i}\PY{o}{+}\PY{o}{+}\PY{p}{)}\PY{+w}{ }\PY{p}{\PYZob{}}
\PY{+w}{            }\PY{n}{sum}\PY{+w}{ }\PY{o}{+}\PY{o}{=}\PY{+w}{ }\PY{n}{a}\PY{o}{[}\PY{n}{i}\PY{o}{]}\PY{+w}{ }\PY{o}{*}\PY{+w}{ }\PY{n}{b}\PY{o}{[}\PY{n}{i}\PY{o}{]}\PY{p}{;}
\PY{+w}{        }\PY{p}{\PYZcb{}}
\PY{+w}{        }\PY{k}{return}\PY{+w}{ }\PY{n}{sum}\PY{p}{;}
\PY{+w}{    }\PY{p}{\PYZcb{}}
\PY{p}{\PYZcb{}}
\end{Verbatim}
\end{codebox}

\begin{codebox}
\begin{Verbatim}[commandchars=\\\{\}]
\PY{c+c1}{// Vectorized version: processes SPECIES.length() float pairs per loop iteration,}
\PY{c+c1}{// falling back to a scalar tail loop for the remainder.}
\PY{k+kn}{import}\PY{+w}{ }\PY{n+nn}{jdk.incubator.vector.FloatVector}\PY{p}{;}
\PY{k+kn}{import}\PY{+w}{ }\PY{n+nn}{jdk.incubator.vector.VectorSpecies}\PY{p}{;}
\PY{k+kn}{import}\PY{+w}{ }\PY{n+nn}{jdk.incubator.vector.VectorOperators}\PY{p}{;}

\PY{k+kd}{public}\PY{+w}{ }\PY{k+kd}{class} \PY{n+nc}{DotProductVector}\PY{+w}{ }\PY{p}{\PYZob{}}
\PY{+w}{    }\PY{k+kd}{static}\PY{+w}{ }\PY{k+kd}{final}\PY{+w}{ }\PY{n}{VectorSpecies}\PY{o}{\PYZlt{}}\PY{n}{Float}\PY{o}{\PYZgt{}}\PY{+w}{ }\PY{n}{SPECIES}\PY{+w}{ }\PY{o}{=}\PY{+w}{ }\PY{n}{FloatVector}\PY{p}{.}\PY{n+na}{SPECIES\PYZus{}PREFERRED}\PY{p}{;}

\PY{+w}{    }\PY{k+kd}{static}\PY{+w}{ }\PY{k+kt}{float}\PY{+w}{ }\PY{n+nf}{dotProduct}\PY{p}{(}\PY{k+kt}{float}\PY{o}{[}\PY{o}{]}\PY{+w}{ }\PY{n}{a}\PY{p}{,}\PY{+w}{ }\PY{k+kt}{float}\PY{o}{[}\PY{o}{]}\PY{+w}{ }\PY{n}{b}\PY{p}{)}\PY{+w}{ }\PY{p}{\PYZob{}}
\PY{+w}{        }\PY{k+kt}{float}\PY{+w}{ }\PY{n}{sum}\PY{+w}{ }\PY{o}{=}\PY{+w}{ }\PY{l+m+mf}{0f}\PY{p}{;}
\PY{+w}{        }\PY{k+kt}{int}\PY{+w}{ }\PY{n}{i}\PY{+w}{ }\PY{o}{=}\PY{+w}{ }\PY{l+m+mi}{0}\PY{p}{;}
\PY{+w}{        }\PY{k+kt}{int}\PY{+w}{ }\PY{n}{upperBound}\PY{+w}{ }\PY{o}{=}\PY{+w}{ }\PY{n}{SPECIES}\PY{p}{.}\PY{n+na}{loopBound}\PY{p}{(}\PY{n}{a}\PY{p}{.}\PY{n+na}{length}\PY{p}{)}\PY{p}{;}\PY{+w}{ }\PY{c+c1}{// largest multiple of the lane count \PYZlt{}= a.length}

\PY{+w}{        }\PY{k}{for}\PY{+w}{ }\PY{p}{(}\PY{p}{;}\PY{+w}{ }\PY{n}{i}\PY{+w}{ }\PY{o}{\PYZlt{}}\PY{+w}{ }\PY{n}{upperBound}\PY{p}{;}\PY{+w}{ }\PY{n}{i}\PY{+w}{ }\PY{o}{+}\PY{o}{=}\PY{+w}{ }\PY{n}{SPECIES}\PY{p}{.}\PY{n+na}{length}\PY{p}{(}\PY{p}{)}\PY{p}{)}\PY{+w}{ }\PY{p}{\PYZob{}}
\PY{+w}{            }\PY{n}{FloatVector}\PY{+w}{ }\PY{n}{va}\PY{+w}{ }\PY{o}{=}\PY{+w}{ }\PY{n}{FloatVector}\PY{p}{.}\PY{n+na}{fromArray}\PY{p}{(}\PY{n}{SPECIES}\PY{p}{,}\PY{+w}{ }\PY{n}{a}\PY{p}{,}\PY{+w}{ }\PY{n}{i}\PY{p}{)}\PY{p}{;}
\PY{+w}{            }\PY{n}{FloatVector}\PY{+w}{ }\PY{n}{vb}\PY{+w}{ }\PY{o}{=}\PY{+w}{ }\PY{n}{FloatVector}\PY{p}{.}\PY{n+na}{fromArray}\PY{p}{(}\PY{n}{SPECIES}\PY{p}{,}\PY{+w}{ }\PY{n}{b}\PY{p}{,}\PY{+w}{ }\PY{n}{i}\PY{p}{)}\PY{p}{;}
\PY{+w}{            }\PY{n}{sum}\PY{+w}{ }\PY{o}{+}\PY{o}{=}\PY{+w}{ }\PY{n}{va}\PY{p}{.}\PY{n+na}{mul}\PY{p}{(}\PY{n}{vb}\PY{p}{)}\PY{p}{.}\PY{n+na}{reduceLanes}\PY{p}{(}\PY{n}{VectorOperators}\PY{p}{.}\PY{n+na}{ADD}\PY{p}{)}\PY{p}{;}
\PY{+w}{        }\PY{p}{\PYZcb{}}

\PY{+w}{        }\PY{k}{for}\PY{+w}{ }\PY{p}{(}\PY{p}{;}\PY{+w}{ }\PY{n}{i}\PY{+w}{ }\PY{o}{\PYZlt{}}\PY{+w}{ }\PY{n}{a}\PY{p}{.}\PY{n+na}{length}\PY{p}{;}\PY{+w}{ }\PY{n}{i}\PY{o}{+}\PY{o}{+}\PY{p}{)}\PY{+w}{ }\PY{p}{\PYZob{}}\PY{+w}{ }\PY{c+c1}{// tail: leftover elements that don\PYZsq{}t fill a full vector}
\PY{+w}{            }\PY{n}{sum}\PY{+w}{ }\PY{o}{+}\PY{o}{=}\PY{+w}{ }\PY{n}{a}\PY{o}{[}\PY{n}{i}\PY{o}{]}\PY{+w}{ }\PY{o}{*}\PY{+w}{ }\PY{n}{b}\PY{o}{[}\PY{n}{i}\PY{o}{]}\PY{p}{;}
\PY{+w}{        }\PY{p}{\PYZcb{}}
\PY{+w}{        }\PY{k}{return}\PY{+w}{ }\PY{n}{sum}\PY{p}{;}
\PY{+w}{    }\PY{p}{\PYZcb{}}
\PY{p}{\PYZcb{}}
\end{Verbatim}
\end{codebox}

The vectorized version does not change the \emph{result} — both compute the same mathematical dot product — it changes how many multiplications the CPU performs per instruction. On a machine with 256-bit AVX2 and \code{float} (4 bytes), \code{SPECIES.\allowbreak{}length()} is typically 8, so each loop iteration does the work of 8 scalar iterations in roughly the same number of cycles.

\subsection[The Loop-Plus-Tail Idiom]{The Loop-Plus-Tail Idiom}
\label{h:17:the-loop-plus-tail-idiom}
Almost every Vector API loop follows the same two-phase shape: a \textbf{vectorized main loop} that processes full lanes, followed by a \textbf{scalar tail loop} that handles whatever’s left over when the array length isn’t an exact multiple of the vector width.

\begin{codebox}
\begin{Verbatim}[commandchars=\\\{\}]
\PY{k+kd}{static}\PY{+w}{ }\PY{k+kt}{void}\PY{+w}{ }\PY{n+nf}{scaleInPlace}\PY{p}{(}\PY{k+kt}{float}\PY{o}{[}\PY{o}{]}\PY{+w}{ }\PY{n}{data}\PY{p}{,}\PY{+w}{ }\PY{k+kt}{float}\PY{+w}{ }\PY{n}{factor}\PY{p}{)}\PY{+w}{ }\PY{p}{\PYZob{}}
\PY{+w}{    }\PY{n}{VectorSpecies}\PY{o}{\PYZlt{}}\PY{n}{Float}\PY{o}{\PYZgt{}}\PY{+w}{ }\PY{n}{species}\PY{+w}{ }\PY{o}{=}\PY{+w}{ }\PY{n}{FloatVector}\PY{p}{.}\PY{n+na}{SPECIES\PYZus{}PREFERRED}\PY{p}{;}
\PY{+w}{    }\PY{k+kt}{int}\PY{+w}{ }\PY{n}{i}\PY{+w}{ }\PY{o}{=}\PY{+w}{ }\PY{l+m+mi}{0}\PY{p}{;}
\PY{+w}{    }\PY{k+kt}{int}\PY{+w}{ }\PY{n}{upperBound}\PY{+w}{ }\PY{o}{=}\PY{+w}{ }\PY{n}{species}\PY{p}{.}\PY{n+na}{loopBound}\PY{p}{(}\PY{n}{data}\PY{p}{.}\PY{n+na}{length}\PY{p}{)}\PY{p}{;}

\PY{+w}{    }\PY{c+c1}{// Main loop: full vector width each iteration}
\PY{+w}{    }\PY{k}{for}\PY{+w}{ }\PY{p}{(}\PY{p}{;}\PY{+w}{ }\PY{n}{i}\PY{+w}{ }\PY{o}{\PYZlt{}}\PY{+w}{ }\PY{n}{upperBound}\PY{p}{;}\PY{+w}{ }\PY{n}{i}\PY{+w}{ }\PY{o}{+}\PY{o}{=}\PY{+w}{ }\PY{n}{species}\PY{p}{.}\PY{n+na}{length}\PY{p}{(}\PY{p}{)}\PY{p}{)}\PY{+w}{ }\PY{p}{\PYZob{}}
\PY{+w}{        }\PY{n}{FloatVector}\PY{+w}{ }\PY{n}{v}\PY{+w}{ }\PY{o}{=}\PY{+w}{ }\PY{n}{FloatVector}\PY{p}{.}\PY{n+na}{fromArray}\PY{p}{(}\PY{n}{species}\PY{p}{,}\PY{+w}{ }\PY{n}{data}\PY{p}{,}\PY{+w}{ }\PY{n}{i}\PY{p}{)}\PY{p}{;}
\PY{+w}{        }\PY{n}{v}\PY{p}{.}\PY{n+na}{mul}\PY{p}{(}\PY{n}{factor}\PY{p}{)}\PY{p}{.}\PY{n+na}{intoArray}\PY{p}{(}\PY{n}{data}\PY{p}{,}\PY{+w}{ }\PY{n}{i}\PY{p}{)}\PY{p}{;}
\PY{+w}{    }\PY{p}{\PYZcb{}}

\PY{+w}{    }\PY{c+c1}{// Tail loop: whatever\PYZsq{}s left (0 to species.length() \PYZhy{} 1 elements)}
\PY{+w}{    }\PY{k}{for}\PY{+w}{ }\PY{p}{(}\PY{p}{;}\PY{+w}{ }\PY{n}{i}\PY{+w}{ }\PY{o}{\PYZlt{}}\PY{+w}{ }\PY{n}{data}\PY{p}{.}\PY{n+na}{length}\PY{p}{;}\PY{+w}{ }\PY{n}{i}\PY{o}{+}\PY{o}{+}\PY{p}{)}\PY{+w}{ }\PY{p}{\PYZob{}}
\PY{+w}{        }\PY{n}{data}\PY{o}{[}\PY{n}{i}\PY{o}{]}\PY{+w}{ }\PY{o}{*}\PY{o}{=}\PY{+w}{ }\PY{n}{factor}\PY{p}{;}
\PY{+w}{    }\PY{p}{\PYZcb{}}
\PY{p}{\PYZcb{}}
\end{Verbatim}
\end{codebox}

An alternative to a separate tail loop is a \textbf{masked final iteration}, using \code{SPECIES.\allowbreak{}index\allowbreak{}In\allowbreak{}Range(\allowbreak{}i,\allowbreak{}~\allowbreak{}data.\allowbreak{}length)} to build a mask that disables out-of-bounds lanes so the whole array can be processed with vector operations alone, at the cost of a mask-compare on every iteration (including the ones that don’t need it):

\begin{codebox}
\begin{Verbatim}[commandchars=\\\{\}]
\PY{k}{for}\PY{+w}{ }\PY{p}{(}\PY{k+kt}{int}\PY{+w}{ }\PY{n}{i}\PY{+w}{ }\PY{o}{=}\PY{+w}{ }\PY{l+m+mi}{0}\PY{p}{;}\PY{+w}{ }\PY{n}{i}\PY{+w}{ }\PY{o}{\PYZlt{}}\PY{+w}{ }\PY{n}{data}\PY{p}{.}\PY{n+na}{length}\PY{p}{;}\PY{+w}{ }\PY{n}{i}\PY{+w}{ }\PY{o}{+}\PY{o}{=}\PY{+w}{ }\PY{n}{species}\PY{p}{.}\PY{n+na}{length}\PY{p}{(}\PY{p}{)}\PY{p}{)}\PY{+w}{ }\PY{p}{\PYZob{}}
\PY{+w}{    }\PY{k+kd}{var}\PY{+w}{ }\PY{n}{mask}\PY{+w}{ }\PY{o}{=}\PY{+w}{ }\PY{n}{species}\PY{p}{.}\PY{n+na}{indexInRange}\PY{p}{(}\PY{n}{i}\PY{p}{,}\PY{+w}{ }\PY{n}{data}\PY{p}{.}\PY{n+na}{length}\PY{p}{)}\PY{p}{;}
\PY{+w}{    }\PY{n}{FloatVector}\PY{+w}{ }\PY{n}{v}\PY{+w}{ }\PY{o}{=}\PY{+w}{ }\PY{n}{FloatVector}\PY{p}{.}\PY{n+na}{fromArray}\PY{p}{(}\PY{n}{species}\PY{p}{,}\PY{+w}{ }\PY{n}{data}\PY{p}{,}\PY{+w}{ }\PY{n}{i}\PY{p}{,}\PY{+w}{ }\PY{n}{mask}\PY{p}{)}\PY{p}{;}
\PY{+w}{    }\PY{n}{v}\PY{p}{.}\PY{n+na}{mul}\PY{p}{(}\PY{n}{factor}\PY{p}{)}\PY{p}{.}\PY{n+na}{intoArray}\PY{p}{(}\PY{n}{data}\PY{p}{,}\PY{+w}{ }\PY{n}{i}\PY{p}{,}\PY{+w}{ }\PY{n}{mask}\PY{p}{)}\PY{p}{;}
\PY{p}{\PYZcb{}}
\end{Verbatim}
\end{codebox}

Both patterns are idiomatic; the explicit tail loop is generally easier to read and slightly faster for the common (non-final) iterations, while the masked version is more compact and avoids duplicating the operation logic.

\subsection[Dependence on Project Valhalla]{Dependence on Project Valhalla}
\label{h:17:dependence-on-project-valhalla}
The Vector API’s classes (\code{FloatVector}, \code{IntVector}, etc.) are, today, ordinary heap-allocated objects — every vector operation risks JIT-dependent allocation and indirection overhead that a hand-written intrinsic wouldn’t have. \textbf{Project Valhalla}’s value types (formerly “inline classes”) are designed to let the JVM represent a \code{FloatVector} as a flat, register-resident sequence of primitives with no object header and no heap allocation at all, the same way a CPU’s SIMD register actually works. The Vector API’s final, non-incubating form is expected to be built on top of Valhalla value types once they ship — which is precisely why the API keeps changing shape release over release: its designers are deliberately keeping the door open for a completely different underlying representation.

\subsection[Production Readiness Caveat]{Production Readiness Caveat}
\label{h:17:production-readiness-caveat}
Given the above, the Vector API should be treated, in a code review, the same way you’d treat any other incubator or preview feature:

\begin{itemize}
\item 
It requires \code{--\allowbreak{}add-\allowbreak{}modules~\allowbreak{}jdk.\allowbreak{}incubator.\allowbreak{}vector} on every build and run command, forever, until it graduates — a permanent build-configuration burden.

\item 
Its behavior and API surface are not covered by Java’s usual strict backward-compatibility guarantees; upgrading the JDK can require source changes.

\item 
Performance-critical vectorization needs (image processing, ML inference, scientific computing) are, today, more commonly delivered via a well-tested native library called through the FFM API described earlier in this chapter, or via a JIT that already does reasonably good auto-vectorization for simple loops without any special API at all.

\item 
\textbf{Recommendation for production code}: avoid the Vector API in shipped production systems until it exits incubation; it is excellent to know for interviews and to experiment with, but a reviewer should flag a \code{jdk.\allowbreak{}incubator.\allowbreak{}vector} import in application code as a real risk, not a style nit.

\end{itemize}

\section[Common Code-Review Interview Pitfalls]{Common Code-Review Interview Pitfalls}
\label{h:17:common-code-review-interview-pitfalls}
\begin{enumerate}
\item 
\textbf{Using a wildcard import (\code{import~\allowbreak{}com.\allowbreak{}example.*;}) in reviewed code.}
Why it matters: it hides exactly which types are in use, risks silent breakage if the package later adds a colliding name, and is rejected by nearly every style guide.

\begin{codebox}
\begin{Verbatim}[commandchars=\\\{\}]
\PY{c+c1}{// Before}
\PY{k+kn}{import}\PY{+w}{ }\PY{n+nn}{com.teampicnic.orders.*}\PY{p}{;}
\PY{c+c1}{// After}
\PY{k+kn}{import}\PY{+w}{ }\PY{n+nn}{com.teampicnic.orders.Order}\PY{p}{;}
\PY{k+kn}{import}\PY{+w}{ }\PY{n+nn}{com.teampicnic.orders.OrderStatus}\PY{p}{;}
\end{Verbatim}
\end{codebox}

\item 
\textbf{Leaving a scratch class in the default (unnamed) package and shipping it.}
Why it matters: it cannot be imported by named-package code, cannot participate in JPMS at all, and invites name collisions as the codebase grows.

\begin{codebox}
\begin{Verbatim}[commandchars=\\\{\}]
\PY{c+c1}{// Before}
\PY{k+kd}{public}\PY{+w}{ }\PY{k+kd}{class} \PY{n+nc}{Utils}\PY{+w}{ }\PY{p}{\PYZob{}}\PY{+w}{ }\PY{c+cm}{/* no package statement */}\PY{+w}{ }\PY{p}{\PYZcb{}}
\PY{c+c1}{// After}
\PY{k+kn}{package}\PY{+w}{ }\PY{n+nn}{com.teampicnic.common}\PY{p}{;}
\PY{k+kd}{public}\PY{+w}{ }\PY{k+kd}{class} \PY{n+nc}{Utils}\PY{+w}{ }\PY{p}{\PYZob{}}\PY{+w}{ }\PY{p}{\PYZcb{}}
\end{Verbatim}
\end{codebox}

\item 
\textbf{Confusing \code{exports} with \code{opens} when a framework needs reflection.}
Why it matters: Jackson/Hibernate/Spring need deep reflective field access; \code{exports} alone grants normal compile-time visibility but still throws \code{Inaccessible\allowbreak{}Object\allowbreak{}Exception} on \code{setAccessible(\allowbreak{}true)}.

\begin{codebox}
\begin{Verbatim}[commandchars=\\\{\}]
\PY{c+c1}{// Before (Jackson fails to deserialize private fields at runtime)}
\PY{n}{module}\PY{+w}{ }\PY{n}{com}\PY{p}{.}\PY{n+na}{teampicnic}\PY{p}{.}\PY{n+na}{billing}\PY{+w}{ }\PY{p}{\PYZob{}}\PY{+w}{ }\PY{n}{exports}\PY{+w}{ }\PY{n}{com}\PY{p}{.}\PY{n+na}{teampicnic}\PY{p}{.}\PY{n+na}{billing}\PY{p}{.}\PY{n+na}{model}\PY{p}{;}\PY{+w}{ }\PY{p}{\PYZcb{}}
\PY{c+c1}{// After}
\PY{n}{module}\PY{+w}{ }\PY{n}{com}\PY{p}{.}\PY{n+na}{teampicnic}\PY{p}{.}\PY{n+na}{billing}\PY{+w}{ }\PY{p}{\PYZob{}}\PY{+w}{ }\PY{n}{opens}\PY{+w}{ }\PY{n}{com}\PY{p}{.}\PY{n+na}{teampicnic}\PY{p}{.}\PY{n+na}{billing}\PY{p}{.}\PY{n+na}{model}\PY{p}{;}\PY{+w}{ }\PY{p}{\PYZcb{}}
\end{Verbatim}
\end{codebox}

\item 
\textbf{Forgetting \code{requires~\allowbreak{}transitive} when a public API returns a type from another module.}
Why it matters: without it, every consumer of your module must separately \code{requires} the dependency your API leaks, breaking encapsulation and causing confusing “package not visible” errors far from the real cause.

\begin{codebox}
\begin{Verbatim}[commandchars=\\\{\}]
\PY{c+c1}{// Before: billing.api returns java.sql.Connection but doesn\PYZsq{}t re\PYZhy{}export java.sql}
\PY{n}{module}\PY{+w}{ }\PY{n}{com}\PY{p}{.}\PY{n+na}{teampicnic}\PY{p}{.}\PY{n+na}{billing}\PY{+w}{ }\PY{p}{\PYZob{}}\PY{+w}{ }\PY{n}{requires}\PY{+w}{ }\PY{n}{java}\PY{p}{.}\PY{n+na}{sql}\PY{p}{;}\PY{+w}{ }\PY{p}{\PYZcb{}}
\PY{c+c1}{// After}
\PY{n}{module}\PY{+w}{ }\PY{n}{com}\PY{p}{.}\PY{n+na}{teampicnic}\PY{p}{.}\PY{n+na}{billing}\PY{+w}{ }\PY{p}{\PYZob{}}\PY{+w}{ }\PY{n}{requires}\PY{+w}{ }\PY{n}{transitive}\PY{+w}{ }\PY{n}{java}\PY{p}{.}\PY{n+na}{sql}\PY{p}{;}\PY{+w}{ }\PY{p}{\PYZcb{}}
\end{Verbatim}
\end{codebox}

\item 
\textbf{Adding \code{--\allowbreak{}add-\allowbreak{}opens}/\code{--\allowbreak{}add-\allowbreak{}exports} flags to “make the error go away” without understanding why they were needed.}
Why it matters: these are escape hatches around intentional strong encapsulation; scattering them across build scripts accumulates hidden, unreviewed access to JDK internals that can break on the next JDK upgrade.

\begin{codebox}
\begin{Verbatim}[commandchars=\\\{\}]
\PY{c+c1}{\PYZsh{} Before: blanket, undocumented flag added to silence a warning}
java\PY{+w}{ }\PYZhy{}\PYZhy{}add\PYZhy{}opens\PY{+w}{ }java.base/java.lang\PY{o}{=}ALL\PYZhy{}UNNAMED\PY{+w}{ }\PYZhy{}jar\PY{+w}{ }app.jar
\PY{c+c1}{\PYZsh{} After: scoped to the specific module that actually needs it, with a}
\PY{c+c1}{\PYZsh{} comment explaining which dependency requires it and a plan to remove it}
java\PY{+w}{ }\PYZhy{}\PYZhy{}add\PYZhy{}opens\PY{+w}{ }java.base/java.lang\PY{o}{=}com.teampicnic.legacybridge\PY{+w}{ }\PYZhy{}jar\PY{+w}{ }app.jar
\end{Verbatim}
\end{codebox}

\item 
\textbf{Assuming every project must adopt JPMS “because it’s modern.”}
Why it matters: for a typical single-deployable application (as opposed to a widely distributed library), the migration cost (split packages, \code{--\allowbreak{}add-\allowbreak{}opens} for test/mocking frameworks) often outweighs the benefit; flagging the \emph{absence} of \code{module-\allowbreak{}info.\allowbreak{}java} as a defect, on its own, is not a strong review comment.

\begin{codebox}
\begin{Verbatim}
Review comment to avoid: "This project has no module-info.java, please modularize it."
Better: ask whether the team's actual pain point (image size? internal
encapsulation?) is better solved by jlink + a slimmer base image, or by
architecture-linting tools, before demanding a full JPMS migration.
\end{Verbatim}
\end{codebox}

\item 
\textbf{Calling native functions via the FFM API without wrapping allocation in try-with-resources.}
Why it matters: memory allocated by an \code{Arena.\allowbreak{}ofConfined()}/\code{ofShared()} is not freed until \code{close()} runs; forgetting it leaks native memory exactly like forgetting to close a file handle, except it won’t show up in heap dumps.

\begin{codebox}
\begin{Verbatim}[commandchars=\\\{\}]
\PY{c+c1}{// Before}
\PY{n}{Arena}\PY{+w}{ }\PY{n}{arena}\PY{+w}{ }\PY{o}{=}\PY{+w}{ }\PY{n}{Arena}\PY{p}{.}\PY{n+na}{ofConfined}\PY{p}{(}\PY{p}{)}\PY{p}{;}
\PY{n}{MemorySegment}\PY{+w}{ }\PY{n}{s}\PY{+w}{ }\PY{o}{=}\PY{+w}{ }\PY{n}{arena}\PY{p}{.}\PY{n+na}{allocate}\PY{p}{(}\PY{l+m+mi}{1024}\PY{p}{)}\PY{p}{;}
\PY{c+c1}{// ... arena never closed, memory leaks for the arena\PYZsq{}s whole (\PYZdq{}forever\PYZdq{}) lifetime}
\PY{c+c1}{// After}
\PY{k}{try}\PY{+w}{ }\PY{p}{(}\PY{n}{Arena}\PY{+w}{ }\PY{n}{arena}\PY{+w}{ }\PY{o}{=}\PY{+w}{ }\PY{n}{Arena}\PY{p}{.}\PY{n+na}{ofConfined}\PY{p}{(}\PY{p}{)}\PY{p}{)}\PY{+w}{ }\PY{p}{\PYZob{}}
\PY{+w}{    }\PY{n}{MemorySegment}\PY{+w}{ }\PY{n}{s}\PY{+w}{ }\PY{o}{=}\PY{+w}{ }\PY{n}{arena}\PY{p}{.}\PY{n+na}{allocate}\PY{p}{(}\PY{l+m+mi}{1024}\PY{p}{)}\PY{p}{;}
\PY{p}{\PYZcb{}}
\end{Verbatim}
\end{codebox}

\item 
\textbf{Passing a \code{FunctionDescriptor} that doesn’t match the real native function signature.}
Why it matters: unlike a normal Java type mismatch, this is not caught by the compiler or even reliably at runtime — it produces undefined behavior (corrupted stack, wrong results, or a JVM crash) because the FFM API trusts the descriptor you provide.

\begin{codebox}
\begin{Verbatim}[commandchars=\\\{\}]
\PY{c+c1}{// Before: strlen actually returns size\PYZus{}t (8 bytes on 64\PYZhy{}bit), declared as 4\PYZhy{}byte int}
\PY{n}{FunctionDescriptor}\PY{p}{.}\PY{n+na}{of}\PY{p}{(}\PY{n}{ValueLayout}\PY{p}{.}\PY{n+na}{JAVA\PYZus{}INT}\PY{p}{,}\PY{+w}{ }\PY{n}{ValueLayout}\PY{p}{.}\PY{n+na}{ADDRESS}\PY{p}{)}\PY{p}{;}
\PY{c+c1}{// After}
\PY{n}{FunctionDescriptor}\PY{p}{.}\PY{n+na}{of}\PY{p}{(}\PY{n}{ValueLayout}\PY{p}{.}\PY{n+na}{JAVA\PYZus{}LONG}\PY{p}{,}\PY{+w}{ }\PY{n}{ValueLayout}\PY{p}{.}\PY{n+na}{ADDRESS}\PY{p}{)}\PY{p}{;}
\end{Verbatim}
\end{codebox}

\item 
\textbf{Running FFM downcalls/upcalls without \code{--\allowbreak{}enable-\allowbreak{}native-\allowbreak{}access} and treating the resulting warning (or, on JDK 24+, exception) as noise to suppress.}
Why it matters: the warning/error exists because restricted operations can crash the JVM or corrupt memory; suppressing it without understanding the risk removes a deliberate safety signal instead of addressing it.

\begin{codebox}
\begin{Verbatim}[commandchars=\\\{\}]
\PY{c+c1}{\PYZsh{} Before: warning ignored / logging turned down to hide it}
java\PY{+w}{ }\PYZhy{}p\PY{+w}{ }out\PY{+w}{ }\PYZhy{}m\PY{+w}{ }com.teampicnic.app/com.teampicnic.app.Main\PY{+w}{ }\PY{l+m}{2}\PYZgt{}/dev/null
\PY{c+c1}{\PYZsh{} After: explicit, scoped, and documented opt\PYZhy{}in}
java\PY{+w}{ }\PYZhy{}\PYZhy{}enable\PYZhy{}native\PYZhy{}access\PY{o}{=}com.teampicnic.app\PY{+w}{ }\PYZhy{}p\PY{+w}{ }out\PY{+w}{ }\PYZhy{}m\PY{+w}{ }com.teampicnic.app/com.teampicnic.app.Main
\end{Verbatim}
\end{codebox}

\item 
\textbf{Shipping \code{jdk.\allowbreak{}incubator.\allowbreak{}vector} code in a production code path.}
Why it matters: incubator modules have no backward-compatibility guarantee and can change or disappear between JDK releases; a production dependency on one means every JDK upgrade is a potential source-breaking event.

\begin{codebox}
\begin{Verbatim}[commandchars=\\\{\}]
\PY{c+c1}{// Before: production hot path depends on the Vector API directly}
\PY{k+kn}{import}\PY{+w}{ }\PY{n+nn}{jdk.incubator.vector.FloatVector}\PY{p}{;}
\PY{c+c1}{// After: keep vectorization behind a well\PYZhy{}tested native library called}
\PY{c+c1}{// via the finalized FFM API, or rely on JIT auto\PYZhy{}vectorization, until}
\PY{c+c1}{// the Vector API graduates out of incubation}
\end{Verbatim}
\end{codebox}

\item 
\textbf{Writing a Vector API loop with no scalar tail (or mask) for leftover elements.}
Why it matters: if the array length isn’t an exact multiple of \code{SPECIES.\allowbreak{}length()}, a naive vectorized loop either throws an out-of-bounds exception or silently skips the trailing elements, producing wrong results.

\begin{codebox}
\begin{Verbatim}[commandchars=\\\{\}]
\PY{c+c1}{// Before: silently drops elements past the last full vector}
\PY{k}{for}\PY{+w}{ }\PY{p}{(}\PY{k+kt}{int}\PY{+w}{ }\PY{n}{i}\PY{+w}{ }\PY{o}{=}\PY{+w}{ }\PY{l+m+mi}{0}\PY{p}{;}\PY{+w}{ }\PY{n}{i}\PY{+w}{ }\PY{o}{\PYZlt{}}\PY{+w}{ }\PY{n}{a}\PY{p}{.}\PY{n+na}{length}\PY{p}{;}\PY{+w}{ }\PY{n}{i}\PY{+w}{ }\PY{o}{+}\PY{o}{=}\PY{+w}{ }\PY{n}{SPECIES}\PY{p}{.}\PY{n+na}{length}\PY{p}{(}\PY{p}{)}\PY{p}{)}\PY{+w}{ }\PY{p}{\PYZob{}}
\PY{+w}{    }\PY{n}{FloatVector}\PY{p}{.}\PY{n+na}{fromArray}\PY{p}{(}\PY{n}{SPECIES}\PY{p}{,}\PY{+w}{ }\PY{n}{a}\PY{p}{,}\PY{+w}{ }\PY{n}{i}\PY{p}{)}\PY{p}{.}\PY{n+na}{mul}\PY{p}{(}\PY{l+m+mf}{2f}\PY{p}{)}\PY{p}{.}\PY{n+na}{intoArray}\PY{p}{(}\PY{n}{a}\PY{p}{,}\PY{+w}{ }\PY{n}{i}\PY{p}{)}\PY{p}{;}
\PY{p}{\PYZcb{}}
\PY{c+c1}{// After}
\PY{k+kt}{int}\PY{+w}{ }\PY{n}{upperBound}\PY{+w}{ }\PY{o}{=}\PY{+w}{ }\PY{n}{SPECIES}\PY{p}{.}\PY{n+na}{loopBound}\PY{p}{(}\PY{n}{a}\PY{p}{.}\PY{n+na}{length}\PY{p}{)}\PY{p}{;}
\PY{k+kt}{int}\PY{+w}{ }\PY{n}{i}\PY{+w}{ }\PY{o}{=}\PY{+w}{ }\PY{l+m+mi}{0}\PY{p}{;}
\PY{k}{for}\PY{+w}{ }\PY{p}{(}\PY{p}{;}\PY{+w}{ }\PY{n}{i}\PY{+w}{ }\PY{o}{\PYZlt{}}\PY{+w}{ }\PY{n}{upperBound}\PY{p}{;}\PY{+w}{ }\PY{n}{i}\PY{+w}{ }\PY{o}{+}\PY{o}{=}\PY{+w}{ }\PY{n}{SPECIES}\PY{p}{.}\PY{n+na}{length}\PY{p}{(}\PY{p}{)}\PY{p}{)}\PY{+w}{ }\PY{p}{\PYZob{}}\PY{+w}{ }\PY{c+cm}{/* vectorized */}\PY{+w}{ }\PY{p}{\PYZcb{}}
\PY{k}{for}\PY{+w}{ }\PY{p}{(}\PY{p}{;}\PY{+w}{ }\PY{n}{i}\PY{+w}{ }\PY{o}{\PYZlt{}}\PY{+w}{ }\PY{n}{a}\PY{p}{.}\PY{n+na}{length}\PY{p}{;}\PY{+w}{ }\PY{n}{i}\PY{o}{+}\PY{o}{+}\PY{p}{)}\PY{+w}{ }\PY{p}{\PYZob{}}\PY{+w}{ }\PY{n}{a}\PY{o}{[}\PY{n}{i}\PY{o}{]}\PY{+w}{ }\PY{o}{*}\PY{o}{=}\PY{+w}{ }\PY{l+m+mf}{2f}\PY{p}{;}\PY{+w}{ }\PY{p}{\PYZcb{}}\PY{+w}{ }\PY{c+c1}{// tail}
\end{Verbatim}
\end{codebox}

\item 
\textbf{Believing \code{opens}/\code{exports} provide security against a determined attacker.}
Why it matters: JPMS encapsulation is a compile-time and linkage-time boundary for well-behaved code, not a security sandbox — \code{--\allowbreak{}add-\allowbreak{}opens}, native access, or an already-loaded reflective handle can bypass it, and it must never be the sole line of defense for genuinely sensitive internals.

\begin{codebox}
\begin{Verbatim}
Review comment: "We don't export the credentials package, so it's secure."
Correction: JPMS encapsulation prevents accidental coupling between
modules, not deliberate bypass; sensitive data still needs its own
access controls (encryption at rest, proper secret management), independent
of module boundaries.
\end{Verbatim}
\end{codebox}

\item 
\textbf{Introducing a split package while adding a new dependency to a modularized project.}
Why it matters: two modules exporting the same package name fail module resolution at launch with a hard-to-diagnose \code{ResolutionException}, often only discovered in CI or production, not at compile time on the developer’s machine if only one of the two JARs was present locally.

\begin{codebox}
\begin{Verbatim}
Error: Module A and module B export package com.example.util to module C
Fix: shade/relocate one of the offending JARs, upgrade to a version that
renamed the package, or exclude the duplicate dependency.
\end{Verbatim}
\end{codebox}

\item 
\textbf{Treating an automatic module (a plain JAR on the module path) as if it had real JPMS encapsulation.}
Why it matters: automatic modules export every package to everyone and read every other module, so \code{requires} on one gives no actual safety guarantee — code review should not assume “it’s on the module path” implies “its internals are protected.”

\begin{codebox}
\begin{Verbatim}
Review comment to avoid: "guava is modularized now, so its internals are encapsulated."
Correction: guava (via Automatic-Module-Name) is an automatic module —
all of its packages remain fully open; only a real module-info.class
with explicit exports would encapsulate anything.
\end{Verbatim}
\end{codebox}

\item 
\textbf{Reviewing \code{package-\allowbreak{}private} classes as if they were dead code because “nothing outside the package uses them.”}
Why it matters: package-private visibility is a deliberate design boundary, not an accident — flagging an unused-outside-package helper class as “should be deleted” without checking whether it’s the intended internal implementation of a package’s public facade produces a wrong, noisy review comment.

\begin{codebox}
\begin{Verbatim}[commandchars=\\\{\}]
\PY{c+c1}{// Correct read: TaxCalculator below is intentionally hidden; it\PYZsq{}s InvoiceService\PYZsq{}s}
\PY{c+c1}{// private implementation detail, not \PYZdq{}orphaned\PYZdq{} code, even though no import}
\PY{c+c1}{// statement anywhere references it directly.}
\PY{k+kd}{class} \PY{n+nc}{TaxCalculator}\PY{+w}{ }\PY{p}{\PYZob{}}\PY{+w}{ }\PY{c+cm}{/* used only by InvoiceService in the same package */}\PY{+w}{ }\PY{p}{\PYZcb{}}
\end{Verbatim}
\end{codebox}

\end{enumerate}


\setcounter{chapter}{17}
\chapter{Networking \& Security}
\label{chap:18}
\label{h:18:18-networking--security}
Almost every real application talks to something over a network — a database, another service, a browser, a file on another machine. This chapter covers Java’s networking building blocks, from raw sockets to the modern HTTP Client, then moves into running external processes safely, and finishes with security fundamentals and cryptography. Security bugs are some of the most expensive bugs a reviewer can miss, so this chapter leans heavily on \textbf{anti-pattern → fix} pairs: code that looks fine but is exploitable, next to the version that isn’t. We target Java 21+ throughout.

\section[Networking (java.net)]{Networking (java.net)}
\label{h:18:networking-javanet}
The \code{java.\allowbreak{}net} package is Java’s low-level toolkit for network communication. It has been in the JDK since Java 1.0, so a lot of application code still uses it directly, even though higher-level APIs (like the HTTP Client, covered next) are usually a better choice for talking HTTP.

\subsection[InetAddress — representing a host]{InetAddress — representing a host}
\label{h:18:inetaddress--representing-a-host}
An \textbf{IP address} is a numeric address for a machine on a network (like \code{192.\allowbreak{}0.\allowbreak{}2.\allowbreak{}1}). \code{InetAddress} represents one, and can resolve a \textbf{hostname} (a human-readable name like \code{example.\allowbreak{}com}) into an IP address via \textbf{DNS} (Domain Name System — the internet’s phonebook that maps names to addresses).

\begin{codebox}
\begin{Verbatim}[commandchars=\\\{\}]
\PY{k+kn}{import}\PY{+w}{ }\PY{n+nn}{java.net.InetAddress}\PY{p}{;}
\PY{k+kn}{import}\PY{+w}{ }\PY{n+nn}{java.net.UnknownHostException}\PY{p}{;}

\PY{k+kd}{public}\PY{+w}{ }\PY{k+kd}{class} \PY{n+nc}{HostLookup}\PY{+w}{ }\PY{p}{\PYZob{}}
\PY{+w}{    }\PY{k+kd}{public}\PY{+w}{ }\PY{k+kd}{static}\PY{+w}{ }\PY{k+kt}{void}\PY{+w}{ }\PY{n+nf}{main}\PY{p}{(}\PY{n}{String}\PY{o}{[}\PY{o}{]}\PY{+w}{ }\PY{n}{args}\PY{p}{)}\PY{+w}{ }\PY{k+kd}{throws}\PY{+w}{ }\PY{n}{UnknownHostException}\PY{+w}{ }\PY{p}{\PYZob{}}
\PY{+w}{        }\PY{n}{InetAddress}\PY{+w}{ }\PY{n}{address}\PY{+w}{ }\PY{o}{=}\PY{+w}{ }\PY{n}{InetAddress}\PY{p}{.}\PY{n+na}{getByName}\PY{p}{(}\PY{l+s}{\PYZdq{}}\PY{l+s}{example.com}\PY{l+s}{\PYZdq{}}\PY{p}{)}\PY{p}{;}
\PY{+w}{        }\PY{n}{System}\PY{p}{.}\PY{n+na}{out}\PY{p}{.}\PY{n+na}{println}\PY{p}{(}\PY{n}{address}\PY{p}{.}\PY{n+na}{getHostName}\PY{p}{(}\PY{p}{)}\PY{p}{)}\PY{p}{;}\PY{+w}{    }\PY{c+c1}{// example.com}
\PY{+w}{        }\PY{n}{System}\PY{p}{.}\PY{n+na}{out}\PY{p}{.}\PY{n+na}{println}\PY{p}{(}\PY{n}{address}\PY{p}{.}\PY{n+na}{getHostAddress}\PY{p}{(}\PY{p}{)}\PY{p}{)}\PY{p}{;}\PY{+w}{  }\PY{c+c1}{// the resolved IP, e.g. 93.184.216.34}
\PY{+w}{    }\PY{p}{\PYZcb{}}
\PY{p}{\PYZcb{}}
\end{Verbatim}
\end{codebox}

DNS lookups are \textbf{blocking network calls}. Calling \code{getByName} in a hot loop, or on every request, can silently stall your application if DNS is slow — cache the result if the host rarely changes.

\subsection[Socket / ServerSocket — TCP communication]{Socket / ServerSocket — TCP communication}
\label{h:18:socket--serversocket--tcp-communication}
A \textbf{socket} is one endpoint of a two-way network connection. \textbf{TCP} (Transmission Control Protocol) is a reliable, ordered, connection-based protocol — think of it as a phone call: you dial, both sides confirm the connection, then you can talk in either direction until someone hangs up. \code{ServerSocket} listens for incoming connections; \code{Socket} represents one established connection (either the server’s side of an accepted connection, or the client’s side).

Here is a tiny \textbf{echo server} — it reads a line of text from a client and sends the same line back — plus a client that talks to it.

\begin{codebox}
\begin{Verbatim}[commandchars=\\\{\}]
\PY{c+c1}{// EchoServer.java}
\PY{k+kn}{import}\PY{+w}{ }\PY{n+nn}{java.io.*}\PY{p}{;}
\PY{k+kn}{import}\PY{+w}{ }\PY{n+nn}{java.net.ServerSocket}\PY{p}{;}
\PY{k+kn}{import}\PY{+w}{ }\PY{n+nn}{java.net.Socket}\PY{p}{;}

\PY{k+kd}{public}\PY{+w}{ }\PY{k+kd}{class} \PY{n+nc}{EchoServer}\PY{+w}{ }\PY{p}{\PYZob{}}
\PY{+w}{    }\PY{k+kd}{public}\PY{+w}{ }\PY{k+kd}{static}\PY{+w}{ }\PY{k+kt}{void}\PY{+w}{ }\PY{n+nf}{main}\PY{p}{(}\PY{n}{String}\PY{o}{[}\PY{o}{]}\PY{+w}{ }\PY{n}{args}\PY{p}{)}\PY{+w}{ }\PY{k+kd}{throws}\PY{+w}{ }\PY{n}{IOException}\PY{+w}{ }\PY{p}{\PYZob{}}
\PY{+w}{        }\PY{k}{try}\PY{+w}{ }\PY{p}{(}\PY{n}{ServerSocket}\PY{+w}{ }\PY{n}{serverSocket}\PY{+w}{ }\PY{o}{=}\PY{+w}{ }\PY{k}{new}\PY{+w}{ }\PY{n}{ServerSocket}\PY{p}{(}\PY{l+m+mi}{5000}\PY{p}{)}\PY{p}{)}\PY{+w}{ }\PY{p}{\PYZob{}}
\PY{+w}{            }\PY{n}{System}\PY{p}{.}\PY{n+na}{out}\PY{p}{.}\PY{n+na}{println}\PY{p}{(}\PY{l+s}{\PYZdq{}}\PY{l+s}{Echo server listening on port 5000}\PY{l+s}{\PYZdq{}}\PY{p}{)}\PY{p}{;}
\PY{+w}{            }\PY{k}{while}\PY{+w}{ }\PY{p}{(}\PY{k+kc}{true}\PY{p}{)}\PY{+w}{ }\PY{p}{\PYZob{}}
\PY{+w}{                }\PY{k}{try}\PY{+w}{ }\PY{p}{(}\PY{n}{Socket}\PY{+w}{ }\PY{n}{client}\PY{+w}{ }\PY{o}{=}\PY{+w}{ }\PY{n}{serverSocket}\PY{p}{.}\PY{n+na}{accept}\PY{p}{(}\PY{p}{)}\PY{p}{;}
\PY{+w}{                     }\PY{n}{BufferedReader}\PY{+w}{ }\PY{n}{in}\PY{+w}{ }\PY{o}{=}\PY{+w}{ }\PY{k}{new}\PY{+w}{ }\PY{n}{BufferedReader}\PY{p}{(}
\PY{+w}{                             }\PY{k}{new}\PY{+w}{ }\PY{n}{InputStreamReader}\PY{p}{(}\PY{n}{client}\PY{p}{.}\PY{n+na}{getInputStream}\PY{p}{(}\PY{p}{)}\PY{p}{)}\PY{p}{)}\PY{p}{;}
\PY{+w}{                     }\PY{n}{PrintWriter}\PY{+w}{ }\PY{n}{out}\PY{+w}{ }\PY{o}{=}\PY{+w}{ }\PY{k}{new}\PY{+w}{ }\PY{n}{PrintWriter}\PY{p}{(}\PY{n}{client}\PY{p}{.}\PY{n+na}{getOutputStream}\PY{p}{(}\PY{p}{)}\PY{p}{,}\PY{+w}{ }\PY{k+kc}{true}\PY{p}{)}\PY{p}{)}\PY{+w}{ }\PY{p}{\PYZob{}}

\PY{+w}{                    }\PY{n}{String}\PY{+w}{ }\PY{n}{line}\PY{+w}{ }\PY{o}{=}\PY{+w}{ }\PY{n}{in}\PY{p}{.}\PY{n+na}{readLine}\PY{p}{(}\PY{p}{)}\PY{p}{;}
\PY{+w}{                    }\PY{k}{if}\PY{+w}{ }\PY{p}{(}\PY{n}{line}\PY{+w}{ }\PY{o}{!}\PY{o}{=}\PY{+w}{ }\PY{k+kc}{null}\PY{p}{)}\PY{+w}{ }\PY{p}{\PYZob{}}
\PY{+w}{                        }\PY{n}{out}\PY{p}{.}\PY{n+na}{println}\PY{p}{(}\PY{l+s}{\PYZdq{}}\PY{l+s}{Echo: }\PY{l+s}{\PYZdq{}}\PY{+w}{ }\PY{o}{+}\PY{+w}{ }\PY{n}{line}\PY{p}{)}\PY{p}{;}
\PY{+w}{                    }\PY{p}{\PYZcb{}}
\PY{+w}{                }\PY{p}{\PYZcb{}}\PY{+w}{ }\PY{c+c1}{// client socket closed automatically, connection ends}
\PY{+w}{            }\PY{p}{\PYZcb{}}
\PY{+w}{        }\PY{p}{\PYZcb{}}
\PY{+w}{    }\PY{p}{\PYZcb{}}
\PY{p}{\PYZcb{}}
\end{Verbatim}
\end{codebox}

\begin{codebox}
\begin{Verbatim}[commandchars=\\\{\}]
\PY{c+c1}{// EchoClient.java}
\PY{k+kn}{import}\PY{+w}{ }\PY{n+nn}{java.io.*}\PY{p}{;}
\PY{k+kn}{import}\PY{+w}{ }\PY{n+nn}{java.net.Socket}\PY{p}{;}
\PY{k+kn}{import}\PY{+w}{ }\PY{n+nn}{java.time.Duration}\PY{p}{;}

\PY{k+kd}{public}\PY{+w}{ }\PY{k+kd}{class} \PY{n+nc}{EchoClient}\PY{+w}{ }\PY{p}{\PYZob{}}
\PY{+w}{    }\PY{k+kd}{public}\PY{+w}{ }\PY{k+kd}{static}\PY{+w}{ }\PY{k+kt}{void}\PY{+w}{ }\PY{n+nf}{main}\PY{p}{(}\PY{n}{String}\PY{o}{[}\PY{o}{]}\PY{+w}{ }\PY{n}{args}\PY{p}{)}\PY{+w}{ }\PY{k+kd}{throws}\PY{+w}{ }\PY{n}{IOException}\PY{+w}{ }\PY{p}{\PYZob{}}
\PY{+w}{        }\PY{k}{try}\PY{+w}{ }\PY{p}{(}\PY{n}{Socket}\PY{+w}{ }\PY{n}{socket}\PY{+w}{ }\PY{o}{=}\PY{+w}{ }\PY{k}{new}\PY{+w}{ }\PY{n}{Socket}\PY{p}{(}\PY{p}{)}\PY{p}{)}\PY{+w}{ }\PY{p}{\PYZob{}}
\PY{+w}{            }\PY{n}{socket}\PY{p}{.}\PY{n+na}{connect}\PY{p}{(}\PY{k}{new}\PY{+w}{ }\PY{n}{java}\PY{p}{.}\PY{n+na}{net}\PY{p}{.}\PY{n+na}{InetSocketAddress}\PY{p}{(}\PY{l+s}{\PYZdq{}}\PY{l+s}{localhost}\PY{l+s}{\PYZdq{}}\PY{p}{,}\PY{+w}{ }\PY{l+m+mi}{5000}\PY{p}{)}\PY{p}{,}\PY{+w}{ }\PY{l+m+mi}{3000}\PY{p}{)}\PY{p}{;}\PY{+w}{ }\PY{c+c1}{// connect timeout: 3s}
\PY{+w}{            }\PY{n}{socket}\PY{p}{.}\PY{n+na}{setSoTimeout}\PY{p}{(}\PY{l+m+mi}{3000}\PY{p}{)}\PY{p}{;}\PY{+w}{ }\PY{c+c1}{// read timeout: 3s}

\PY{+w}{            }\PY{n}{PrintWriter}\PY{+w}{ }\PY{n}{out}\PY{+w}{ }\PY{o}{=}\PY{+w}{ }\PY{k}{new}\PY{+w}{ }\PY{n}{PrintWriter}\PY{p}{(}\PY{n}{socket}\PY{p}{.}\PY{n+na}{getOutputStream}\PY{p}{(}\PY{p}{)}\PY{p}{,}\PY{+w}{ }\PY{k+kc}{true}\PY{p}{)}\PY{p}{;}
\PY{+w}{            }\PY{n}{BufferedReader}\PY{+w}{ }\PY{n}{in}\PY{+w}{ }\PY{o}{=}\PY{+w}{ }\PY{k}{new}\PY{+w}{ }\PY{n}{BufferedReader}\PY{p}{(}\PY{k}{new}\PY{+w}{ }\PY{n}{InputStreamReader}\PY{p}{(}\PY{n}{socket}\PY{p}{.}\PY{n+na}{getInputStream}\PY{p}{(}\PY{p}{)}\PY{p}{)}\PY{p}{)}\PY{p}{;}

\PY{+w}{            }\PY{n}{out}\PY{p}{.}\PY{n+na}{println}\PY{p}{(}\PY{l+s}{\PYZdq{}}\PY{l+s}{Hello, server!}\PY{l+s}{\PYZdq{}}\PY{p}{)}\PY{p}{;}
\PY{+w}{            }\PY{n}{System}\PY{p}{.}\PY{n+na}{out}\PY{p}{.}\PY{n+na}{println}\PY{p}{(}\PY{l+s}{\PYZdq{}}\PY{l+s}{Server replied: }\PY{l+s}{\PYZdq{}}\PY{+w}{ }\PY{o}{+}\PY{+w}{ }\PY{n}{in}\PY{p}{.}\PY{n+na}{readLine}\PY{p}{(}\PY{p}{)}\PY{p}{)}\PY{p}{;}
\PY{+w}{        }\PY{p}{\PYZcb{}}
\PY{+w}{    }\PY{p}{\PYZcb{}}
\PY{p}{\PYZcb{}}
\end{Verbatim}
\end{codebox}

Each connected \code{Socket} gets its own thread (or virtual thread) in a simple server like this; production servers usually use non-blocking I/O or a thread pool to handle thousands of clients without one thread per connection.

\subsection[DatagramSocket — UDP communication]{DatagramSocket — UDP communication}
\label{h:18:datagramsocket--udp-communication}
\textbf{UDP} (User Datagram Protocol) is connectionless and unreliable: you send a \textbf{packet} (called a \textbf{datagram}) and hope it arrives, with no guarantee of order or delivery. Think of it as mailing postcards instead of making a phone call — cheaper, faster, but some might get lost. It’s used where occasional loss is acceptable and low latency matters more, such as video streaming, DNS, or game state updates.

\begin{codebox}
\begin{Verbatim}[commandchars=\\\{\}]
\PY{k+kn}{import}\PY{+w}{ }\PY{n+nn}{java.net.DatagramPacket}\PY{p}{;}
\PY{k+kn}{import}\PY{+w}{ }\PY{n+nn}{java.net.DatagramSocket}\PY{p}{;}
\PY{k+kn}{import}\PY{+w}{ }\PY{n+nn}{java.net.InetSocketAddress}\PY{p}{;}

\PY{k+kd}{public}\PY{+w}{ }\PY{k+kd}{class} \PY{n+nc}{UdpSender}\PY{+w}{ }\PY{p}{\PYZob{}}
\PY{+w}{    }\PY{k+kd}{public}\PY{+w}{ }\PY{k+kd}{static}\PY{+w}{ }\PY{k+kt}{void}\PY{+w}{ }\PY{n+nf}{main}\PY{p}{(}\PY{n}{String}\PY{o}{[}\PY{o}{]}\PY{+w}{ }\PY{n}{args}\PY{p}{)}\PY{+w}{ }\PY{k+kd}{throws}\PY{+w}{ }\PY{n}{Exception}\PY{+w}{ }\PY{p}{\PYZob{}}
\PY{+w}{        }\PY{k}{try}\PY{+w}{ }\PY{p}{(}\PY{n}{DatagramSocket}\PY{+w}{ }\PY{n}{socket}\PY{+w}{ }\PY{o}{=}\PY{+w}{ }\PY{k}{new}\PY{+w}{ }\PY{n}{DatagramSocket}\PY{p}{(}\PY{p}{)}\PY{p}{)}\PY{+w}{ }\PY{p}{\PYZob{}}
\PY{+w}{            }\PY{k+kt}{byte}\PY{o}{[}\PY{o}{]}\PY{+w}{ }\PY{n}{data}\PY{+w}{ }\PY{o}{=}\PY{+w}{ }\PY{l+s}{\PYZdq{}}\PY{l+s}{ping}\PY{l+s}{\PYZdq{}}\PY{p}{.}\PY{n+na}{getBytes}\PY{p}{(}\PY{p}{)}\PY{p}{;}
\PY{+w}{            }\PY{n}{DatagramPacket}\PY{+w}{ }\PY{n}{packet}\PY{+w}{ }\PY{o}{=}\PY{+w}{ }\PY{k}{new}\PY{+w}{ }\PY{n}{DatagramPacket}\PY{p}{(}
\PY{+w}{                    }\PY{n}{data}\PY{p}{,}\PY{+w}{ }\PY{n}{data}\PY{p}{.}\PY{n+na}{length}\PY{p}{,}\PY{+w}{ }\PY{k}{new}\PY{+w}{ }\PY{n}{InetSocketAddress}\PY{p}{(}\PY{l+s}{\PYZdq{}}\PY{l+s}{localhost}\PY{l+s}{\PYZdq{}}\PY{p}{,}\PY{+w}{ }\PY{l+m+mi}{6000}\PY{p}{)}\PY{p}{)}\PY{p}{;}
\PY{+w}{            }\PY{n}{socket}\PY{p}{.}\PY{n+na}{send}\PY{p}{(}\PY{n}{packet}\PY{p}{)}\PY{p}{;}
\PY{+w}{        }\PY{p}{\PYZcb{}}
\PY{+w}{    }\PY{p}{\PYZcb{}}
\PY{p}{\PYZcb{}}

\PY{k+kd}{public}\PY{+w}{ }\PY{k+kd}{class} \PY{n+nc}{UdpReceiver}\PY{+w}{ }\PY{p}{\PYZob{}}
\PY{+w}{    }\PY{k+kd}{public}\PY{+w}{ }\PY{k+kd}{static}\PY{+w}{ }\PY{k+kt}{void}\PY{+w}{ }\PY{n+nf}{main}\PY{p}{(}\PY{n}{String}\PY{o}{[}\PY{o}{]}\PY{+w}{ }\PY{n}{args}\PY{p}{)}\PY{+w}{ }\PY{k+kd}{throws}\PY{+w}{ }\PY{n}{Exception}\PY{+w}{ }\PY{p}{\PYZob{}}
\PY{+w}{        }\PY{k}{try}\PY{+w}{ }\PY{p}{(}\PY{n}{DatagramSocket}\PY{+w}{ }\PY{n}{socket}\PY{+w}{ }\PY{o}{=}\PY{+w}{ }\PY{k}{new}\PY{+w}{ }\PY{n}{DatagramSocket}\PY{p}{(}\PY{l+m+mi}{6000}\PY{p}{)}\PY{p}{)}\PY{+w}{ }\PY{p}{\PYZob{}}
\PY{+w}{            }\PY{k+kt}{byte}\PY{o}{[}\PY{o}{]}\PY{+w}{ }\PY{n}{buffer}\PY{+w}{ }\PY{o}{=}\PY{+w}{ }\PY{k}{new}\PY{+w}{ }\PY{k+kt}{byte}\PY{o}{[}\PY{l+m+mi}{1024}\PY{o}{]}\PY{p}{;}
\PY{+w}{            }\PY{n}{DatagramPacket}\PY{+w}{ }\PY{n}{packet}\PY{+w}{ }\PY{o}{=}\PY{+w}{ }\PY{k}{new}\PY{+w}{ }\PY{n}{DatagramPacket}\PY{p}{(}\PY{n}{buffer}\PY{p}{,}\PY{+w}{ }\PY{n}{buffer}\PY{p}{.}\PY{n+na}{length}\PY{p}{)}\PY{p}{;}
\PY{+w}{            }\PY{n}{socket}\PY{p}{.}\PY{n+na}{receive}\PY{p}{(}\PY{n}{packet}\PY{p}{)}\PY{p}{;}\PY{+w}{ }\PY{c+c1}{// blocks until a packet arrives}
\PY{+w}{            }\PY{n}{String}\PY{+w}{ }\PY{n}{message}\PY{+w}{ }\PY{o}{=}\PY{+w}{ }\PY{k}{new}\PY{+w}{ }\PY{n}{String}\PY{p}{(}\PY{n}{packet}\PY{p}{.}\PY{n+na}{getData}\PY{p}{(}\PY{p}{)}\PY{p}{,}\PY{+w}{ }\PY{l+m+mi}{0}\PY{p}{,}\PY{+w}{ }\PY{n}{packet}\PY{p}{.}\PY{n+na}{getLength}\PY{p}{(}\PY{p}{)}\PY{p}{)}\PY{p}{;}
\PY{+w}{            }\PY{n}{System}\PY{p}{.}\PY{n+na}{out}\PY{p}{.}\PY{n+na}{println}\PY{p}{(}\PY{l+s}{\PYZdq{}}\PY{l+s}{Received: }\PY{l+s}{\PYZdq{}}\PY{+w}{ }\PY{o}{+}\PY{+w}{ }\PY{n}{message}\PY{p}{)}\PY{p}{;}
\PY{+w}{        }\PY{p}{\PYZcb{}}
\PY{+w}{    }\PY{p}{\PYZcb{}}
\PY{p}{\PYZcb{}}
\end{Verbatim}
\end{codebox}

\subsection[URL / URI — and why URL.equals is a trap]{URL / URI — and why \code{URL.\allowbreak{}equals} is a trap}
\label{h:18:url--uri--and-why-urlequals-is-a-trap}
\code{URI} (Uniform Resource Identifier) is just a \textbf{syntax} for identifying a resource — it does no network I/O. \code{URL} (Uniform Resource Locator) is a URI that also knows how to \emph{fetch} the resource, and that extra behavior is where the trouble starts.

\textbf{\code{URL.\allowbreak{}equals()} and \code{URL.\allowbreak{}hashCode()} perform a DNS lookup} to compare hosts, because the spec says two URLs are “equal” only if they resolve to the same IP address. This means comparing two \code{URL} objects — or worse, putting them in a \code{HashSet} or using them as \code{HashMap} keys — can trigger a real network call, block on slow or unreachable DNS, and even give different results at different times if DNS changes.

\begin{codebox}
\begin{Verbatim}[commandchars=\\\{\}]
\PY{c+c1}{// Anti\PYZhy{}pattern: URLs in a HashSet trigger silent DNS lookups on every add/contains}
\PY{k+kn}{import}\PY{+w}{ }\PY{n+nn}{java.net.URL}\PY{p}{;}
\PY{k+kn}{import}\PY{+w}{ }\PY{n+nn}{java.util.HashSet}\PY{p}{;}
\PY{k+kn}{import}\PY{+w}{ }\PY{n+nn}{java.util.Set}\PY{p}{;}

\PY{n}{Set}\PY{o}{\PYZlt{}}\PY{n}{URL}\PY{o}{\PYZgt{}}\PY{+w}{ }\PY{n}{seen}\PY{+w}{ }\PY{o}{=}\PY{+w}{ }\PY{k}{new}\PY{+w}{ }\PY{n}{HashSet}\PY{o}{\PYZlt{}}\PY{o}{\PYZgt{}}\PY{p}{(}\PY{p}{)}\PY{p}{;}
\PY{n}{seen}\PY{p}{.}\PY{n+na}{add}\PY{p}{(}\PY{k}{new}\PY{+w}{ }\PY{n}{URL}\PY{p}{(}\PY{l+s}{\PYZdq{}}\PY{l+s}{https://example.com/a}\PY{l+s}{\PYZdq{}}\PY{p}{)}\PY{p}{)}\PY{p}{;}\PY{+w}{      }\PY{c+c1}{// may block on DNS}
\PY{k+kt}{boolean}\PY{+w}{ }\PY{n}{visited}\PY{+w}{ }\PY{o}{=}\PY{+w}{ }\PY{n}{seen}\PY{p}{.}\PY{n+na}{contains}\PY{p}{(}\PY{k}{new}\PY{+w}{ }\PY{n}{URL}\PY{p}{(}\PY{l+s}{\PYZdq{}}\PY{l+s}{https://example.com/b}\PY{l+s}{\PYZdq{}}\PY{p}{)}\PY{p}{)}\PY{p}{;}\PY{+w}{ }\PY{c+c1}{// may block on DNS too}
\end{Verbatim}
\end{codebox}

\begin{codebox}
\begin{Verbatim}[commandchars=\\\{\}]
\PY{c+c1}{// Fix: use URI, which never touches the network}
\PY{k+kn}{import}\PY{+w}{ }\PY{n+nn}{java.net.URI}\PY{p}{;}
\PY{k+kn}{import}\PY{+w}{ }\PY{n+nn}{java.util.HashSet}\PY{p}{;}
\PY{k+kn}{import}\PY{+w}{ }\PY{n+nn}{java.util.Set}\PY{p}{;}

\PY{n}{Set}\PY{o}{\PYZlt{}}\PY{n}{URI}\PY{o}{\PYZgt{}}\PY{+w}{ }\PY{n}{seen}\PY{+w}{ }\PY{o}{=}\PY{+w}{ }\PY{k}{new}\PY{+w}{ }\PY{n}{HashSet}\PY{o}{\PYZlt{}}\PY{o}{\PYZgt{}}\PY{p}{(}\PY{p}{)}\PY{p}{;}
\PY{n}{seen}\PY{p}{.}\PY{n+na}{add}\PY{p}{(}\PY{n}{URI}\PY{p}{.}\PY{n+na}{create}\PY{p}{(}\PY{l+s}{\PYZdq{}}\PY{l+s}{https://example.com/a}\PY{l+s}{\PYZdq{}}\PY{p}{)}\PY{p}{)}\PY{p}{;}
\PY{k+kt}{boolean}\PY{+w}{ }\PY{n}{visited}\PY{+w}{ }\PY{o}{=}\PY{+w}{ }\PY{n}{seen}\PY{p}{.}\PY{n+na}{contains}\PY{p}{(}\PY{n}{URI}\PY{p}{.}\PY{n+na}{create}\PY{p}{(}\PY{l+s}{\PYZdq{}}\PY{l+s}{https://example.com/b}\PY{l+s}{\PYZdq{}}\PY{p}{)}\PY{p}{)}\PY{p}{;}\PY{+w}{ }\PY{c+c1}{// pure string/scheme comparison}
\end{Verbatim}
\end{codebox}

Rule of thumb: use \code{URI} for parsing, comparing, and storing. Only convert to \code{URL} (via \code{uri.\allowbreak{}toURL()}) at the last moment, right before you actually need to open a connection.

\subsection[URLConnection / HttpURLConnection — the legacy way]{URLConnection / HttpURLConnection — the legacy way}
\label{h:18:urlconnection--httpurlconnection--the-legacy-way}
Before the HTTP Client API (next section), Java code opened HTTP connections through \code{URL.\allowbreak{}openConnection()}, which returns a \code{URLConnection} (or, for HTTP, an \code{HttpURLConnection}). It’s still present in the JDK and you’ll see it in older codebases, but it’s clunky: manual stream handling, awkward redirect and header handling, and no built-in HTTP/2 or async support. Treat it as \textbf{legacy} — prefer the HTTP Client API for new code.

\begin{codebox}
\begin{Verbatim}[commandchars=\\\{\}]
\PY{k+kn}{import}\PY{+w}{ }\PY{n+nn}{java.net.HttpURLConnection}\PY{p}{;}
\PY{k+kn}{import}\PY{+w}{ }\PY{n+nn}{java.net.URL}\PY{p}{;}
\PY{k+kn}{import}\PY{+w}{ }\PY{n+nn}{java.io.*}\PY{p}{;}

\PY{k+kd}{public}\PY{+w}{ }\PY{k+kd}{class} \PY{n+nc}{LegacyGet}\PY{+w}{ }\PY{p}{\PYZob{}}
\PY{+w}{    }\PY{k+kd}{public}\PY{+w}{ }\PY{k+kd}{static}\PY{+w}{ }\PY{k+kt}{void}\PY{+w}{ }\PY{n+nf}{main}\PY{p}{(}\PY{n}{String}\PY{o}{[}\PY{o}{]}\PY{+w}{ }\PY{n}{args}\PY{p}{)}\PY{+w}{ }\PY{k+kd}{throws}\PY{+w}{ }\PY{n}{IOException}\PY{+w}{ }\PY{p}{\PYZob{}}
\PY{+w}{        }\PY{n}{URL}\PY{+w}{ }\PY{n}{url}\PY{+w}{ }\PY{o}{=}\PY{+w}{ }\PY{k}{new}\PY{+w}{ }\PY{n}{URL}\PY{p}{(}\PY{l+s}{\PYZdq{}}\PY{l+s}{https://example.com/api/status}\PY{l+s}{\PYZdq{}}\PY{p}{)}\PY{p}{;}
\PY{+w}{        }\PY{n}{HttpURLConnection}\PY{+w}{ }\PY{n}{conn}\PY{+w}{ }\PY{o}{=}\PY{+w}{ }\PY{p}{(}\PY{n}{HttpURLConnection}\PY{p}{)}\PY{+w}{ }\PY{n}{url}\PY{p}{.}\PY{n+na}{openConnection}\PY{p}{(}\PY{p}{)}\PY{p}{;}
\PY{+w}{        }\PY{n}{conn}\PY{p}{.}\PY{n+na}{setRequestMethod}\PY{p}{(}\PY{l+s}{\PYZdq{}}\PY{l+s}{GET}\PY{l+s}{\PYZdq{}}\PY{p}{)}\PY{p}{;}
\PY{+w}{        }\PY{n}{conn}\PY{p}{.}\PY{n+na}{setConnectTimeout}\PY{p}{(}\PY{l+m+mi}{3000}\PY{p}{)}\PY{p}{;}\PY{+w}{ }\PY{c+c1}{// connect timeout: 3s}
\PY{+w}{        }\PY{n}{conn}\PY{p}{.}\PY{n+na}{setReadTimeout}\PY{p}{(}\PY{l+m+mi}{3000}\PY{p}{)}\PY{p}{;}\PY{+w}{    }\PY{c+c1}{// read timeout: 3s}

\PY{+w}{        }\PY{k}{try}\PY{+w}{ }\PY{p}{(}\PY{n}{BufferedReader}\PY{+w}{ }\PY{n}{reader}\PY{+w}{ }\PY{o}{=}\PY{+w}{ }\PY{k}{new}\PY{+w}{ }\PY{n}{BufferedReader}\PY{p}{(}
\PY{+w}{                }\PY{k}{new}\PY{+w}{ }\PY{n}{InputStreamReader}\PY{p}{(}\PY{n}{conn}\PY{p}{.}\PY{n+na}{getInputStream}\PY{p}{(}\PY{p}{)}\PY{p}{)}\PY{p}{)}\PY{p}{)}\PY{+w}{ }\PY{p}{\PYZob{}}
\PY{+w}{            }\PY{n}{String}\PY{+w}{ }\PY{n}{body}\PY{+w}{ }\PY{o}{=}\PY{+w}{ }\PY{n}{reader}\PY{p}{.}\PY{n+na}{lines}\PY{p}{(}\PY{p}{)}\PY{p}{.}\PY{n+na}{reduce}\PY{p}{(}\PY{l+s}{\PYZdq{}}\PY{l+s}{\PYZdq{}}\PY{p}{,}\PY{+w}{ }\PY{p}{(}\PY{n}{a}\PY{p}{,}\PY{+w}{ }\PY{n}{b}\PY{p}{)}\PY{+w}{ }\PY{o}{\PYZhy{}}\PY{o}{\PYZgt{}}\PY{+w}{ }\PY{n}{a}\PY{+w}{ }\PY{o}{+}\PY{+w}{ }\PY{n}{b}\PY{p}{)}\PY{p}{;}
\PY{+w}{            }\PY{n}{System}\PY{p}{.}\PY{n+na}{out}\PY{p}{.}\PY{n+na}{println}\PY{p}{(}\PY{l+s}{\PYZdq{}}\PY{l+s}{Status: }\PY{l+s}{\PYZdq{}}\PY{+w}{ }\PY{o}{+}\PY{+w}{ }\PY{n}{conn}\PY{p}{.}\PY{n+na}{getResponseCode}\PY{p}{(}\PY{p}{)}\PY{+w}{ }\PY{o}{+}\PY{+w}{ }\PY{l+s}{\PYZdq{}}\PY{l+s}{ body: }\PY{l+s}{\PYZdq{}}\PY{+w}{ }\PY{o}{+}\PY{+w}{ }\PY{n}{body}\PY{p}{)}\PY{p}{;}
\PY{+w}{        }\PY{p}{\PYZcb{}}\PY{+w}{ }\PY{k}{finally}\PY{+w}{ }\PY{p}{\PYZob{}}
\PY{+w}{            }\PY{n}{conn}\PY{p}{.}\PY{n+na}{disconnect}\PY{p}{(}\PY{p}{)}\PY{p}{;}
\PY{+w}{        }\PY{p}{\PYZcb{}}
\PY{+w}{    }\PY{p}{\PYZcb{}}
\PY{p}{\PYZcb{}}
\end{Verbatim}
\end{codebox}

\subsection[Timeouts — the setting a missing default punishes you for]{Timeouts — the setting a missing default punishes you for}
\label{h:18:timeouts--the-setting-a-missing-default-punishes-you-for}
Neither \code{Socket}, \code{HttpURLConnection}, nor (as we’ll see) \code{HttpClient} has a sane default timeout baked in for every operation. Two timeouts matter:

\begin{itemize}
\item 
\textbf{Connect timeout} — how long to wait while establishing the TCP connection.

\item 
\textbf{Socket/read timeout} (\code{soTimeout}) — how long to wait for data once connected, before giving up on a read.

\end{itemize}

\begin{codebox}
\begin{Verbatim}[commandchars=\\\{\}]
\PY{c+c1}{// Anti\PYZhy{}pattern: no timeouts at all}
\PY{n}{Socket}\PY{+w}{ }\PY{n}{socket}\PY{+w}{ }\PY{o}{=}\PY{+w}{ }\PY{k}{new}\PY{+w}{ }\PY{n}{Socket}\PY{p}{(}\PY{l+s}{\PYZdq{}}\PY{l+s}{flaky\PYZhy{}host}\PY{l+s}{\PYZdq{}}\PY{p}{,}\PY{+w}{ }\PY{l+m+mi}{443}\PY{p}{)}\PY{p}{;}\PY{+w}{ }\PY{c+c1}{// may hang forever if the host never responds}

\PY{c+c1}{// Fix: always set both}
\PY{n}{Socket}\PY{+w}{ }\PY{n}{socket}\PY{+w}{ }\PY{o}{=}\PY{+w}{ }\PY{k}{new}\PY{+w}{ }\PY{n}{Socket}\PY{p}{(}\PY{p}{)}\PY{p}{;}
\PY{n}{socket}\PY{p}{.}\PY{n+na}{connect}\PY{p}{(}\PY{k}{new}\PY{+w}{ }\PY{n}{InetSocketAddress}\PY{p}{(}\PY{l+s}{\PYZdq{}}\PY{l+s}{flaky\PYZhy{}host}\PY{l+s}{\PYZdq{}}\PY{p}{,}\PY{+w}{ }\PY{l+m+mi}{443}\PY{p}{)}\PY{p}{,}\PY{+w}{ }\PY{l+m+mi}{3000}\PY{p}{)}\PY{p}{;}\PY{+w}{ }\PY{c+c1}{// connect timeout}
\PY{n}{socket}\PY{p}{.}\PY{n+na}{setSoTimeout}\PY{p}{(}\PY{l+m+mi}{5000}\PY{p}{)}\PY{p}{;}\PY{+w}{ }\PY{c+c1}{// read timeout}
\end{Verbatim}
\end{codebox}

Why this is a real production incident, not a theoretical concern: if a downstream service hangs (network partition, overloaded server, firewall silently dropping packets), a thread with no timeout can block \textbf{forever}. Multiply that by every request thread calling the same downstream service, and the whole application thread pool fills up with stuck threads — a \textbf{cascading failure} where one slow dependency takes down an entire service, even though the dependency never returned an error. Always set explicit connect and read timeouts on every network call.

\subsection[SocketChannel — non-blocking I/O]{SocketChannel — non-blocking I/O}
\label{h:18:socketchannel--non-blocking-io}
\code{SocketChannel} (from \code{java.\allowbreak{}nio.\allowbreak{}channels}) is a \textbf{non-blocking} alternative to \code{Socket}: instead of a thread parking until data arrives, you register interest in an event (readable, writable, connectable) with a \code{Selector}, and one thread can monitor many channels at once.

\begin{codebox}
\begin{Verbatim}[commandchars=\\\{\}]
\PY{k+kn}{import}\PY{+w}{ }\PY{n+nn}{java.net.InetSocketAddress}\PY{p}{;}
\PY{k+kn}{import}\PY{+w}{ }\PY{n+nn}{java.nio.ByteBuffer}\PY{p}{;}
\PY{k+kn}{import}\PY{+w}{ }\PY{n+nn}{java.nio.channels.SocketChannel}\PY{p}{;}

\PY{k+kd}{public}\PY{+w}{ }\PY{k+kd}{class} \PY{n+nc}{NonBlockingClient}\PY{+w}{ }\PY{p}{\PYZob{}}
\PY{+w}{    }\PY{k+kd}{public}\PY{+w}{ }\PY{k+kd}{static}\PY{+w}{ }\PY{k+kt}{void}\PY{+w}{ }\PY{n+nf}{main}\PY{p}{(}\PY{n}{String}\PY{o}{[}\PY{o}{]}\PY{+w}{ }\PY{n}{args}\PY{p}{)}\PY{+w}{ }\PY{k+kd}{throws}\PY{+w}{ }\PY{n}{Exception}\PY{+w}{ }\PY{p}{\PYZob{}}
\PY{+w}{        }\PY{k}{try}\PY{+w}{ }\PY{p}{(}\PY{n}{SocketChannel}\PY{+w}{ }\PY{n}{channel}\PY{+w}{ }\PY{o}{=}\PY{+w}{ }\PY{n}{SocketChannel}\PY{p}{.}\PY{n+na}{open}\PY{p}{(}\PY{p}{)}\PY{p}{)}\PY{+w}{ }\PY{p}{\PYZob{}}
\PY{+w}{            }\PY{n}{channel}\PY{p}{.}\PY{n+na}{configureBlocking}\PY{p}{(}\PY{k+kc}{false}\PY{p}{)}\PY{p}{;}
\PY{+w}{            }\PY{n}{channel}\PY{p}{.}\PY{n+na}{connect}\PY{p}{(}\PY{k}{new}\PY{+w}{ }\PY{n}{InetSocketAddress}\PY{p}{(}\PY{l+s}{\PYZdq{}}\PY{l+s}{example.com}\PY{l+s}{\PYZdq{}}\PY{p}{,}\PY{+w}{ }\PY{l+m+mi}{80}\PY{p}{)}\PY{p}{)}\PY{p}{;}

\PY{+w}{            }\PY{k}{while}\PY{+w}{ }\PY{p}{(}\PY{o}{!}\PY{n}{channel}\PY{p}{.}\PY{n+na}{finishConnect}\PY{p}{(}\PY{p}{)}\PY{p}{)}\PY{+w}{ }\PY{p}{\PYZob{}}
\PY{+w}{                }\PY{c+c1}{// could do other work here instead of spinning}
\PY{+w}{            }\PY{p}{\PYZcb{}}

\PY{+w}{            }\PY{n}{ByteBuffer}\PY{+w}{ }\PY{n}{buffer}\PY{+w}{ }\PY{o}{=}\PY{+w}{ }\PY{n}{ByteBuffer}\PY{p}{.}\PY{n+na}{wrap}\PY{p}{(}\PY{l+s}{\PYZdq{}}\PY{l+s}{GET / HTTP/1.0}\PY{l+s}{\PYZbs{}}\PY{l+s}{r}\PY{l+s}{\PYZbs{}}\PY{l+s}{n}\PY{l+s}{\PYZbs{}}\PY{l+s}{r}\PY{l+s}{\PYZbs{}}\PY{l+s}{n}\PY{l+s}{\PYZdq{}}\PY{p}{.}\PY{n+na}{getBytes}\PY{p}{(}\PY{p}{)}\PY{p}{)}\PY{p}{;}
\PY{+w}{            }\PY{n}{channel}\PY{p}{.}\PY{n+na}{write}\PY{p}{(}\PY{n}{buffer}\PY{p}{)}\PY{p}{;}
\PY{+w}{        }\PY{p}{\PYZcb{}}
\PY{+w}{    }\PY{p}{\PYZcb{}}
\PY{p}{\PYZcb{}}
\end{Verbatim}
\end{codebox}

Non-blocking I/O is what lets a server handle thousands of connections with a handful of threads, instead of one thread per connection. It’s more complex to write correctly, which is why frameworks (Netty, and Java’s own virtual threads for simpler blocking-style code) exist to hide the complexity.

\textbf{Unix domain sockets} (JDK 16+): \code{SocketChannel} and \code{ServerSocketChannel} can also bind to a Unix domain socket (a file-path-based socket for fast local inter-process communication, no network stack involved) via \code{Unix\allowbreak{}Domain\allowbreak{}Socket\allowbreak{}Address.\allowbreak{}of(\allowbreak{}path)} — useful for talking to a sidecar process on the same machine without going through TCP/IP.

\section[HTTP Client API]{HTTP Client API}
\label{h:18:http-client-api}
Since Java 11, the JDK ships a modern \code{java.\allowbreak{}net.\allowbreak{}http.\allowbreak{}Http\allowbreak{}Client} — a proper HTTP/1.1 and HTTP/2 client with synchronous and asynchronous modes, replacing the awkward \code{HttpURLConnection} for new code.

\subsection[Building a client]{Building a client}
\label{h:18:building-a-client}
\code{HttpClient} instances are \textbf{immutable and thread-safe} once built, and expensive to create (they manage connection pools). Build one and \textbf{reuse it} across your whole application instead of creating a new client per request.

\begin{codebox}
\begin{Verbatim}[commandchars=\\\{\}]
\PY{k+kn}{import}\PY{+w}{ }\PY{n+nn}{java.net.ProxySelector}\PY{p}{;}
\PY{k+kn}{import}\PY{+w}{ }\PY{n+nn}{java.net.http.HttpClient}\PY{p}{;}
\PY{k+kn}{import}\PY{+w}{ }\PY{n+nn}{java.time.Duration}\PY{p}{;}
\PY{k+kn}{import}\PY{+w}{ }\PY{n+nn}{java.util.concurrent.Executors}\PY{p}{;}

\PY{n}{HttpClient}\PY{+w}{ }\PY{n}{client}\PY{+w}{ }\PY{o}{=}\PY{+w}{ }\PY{n}{HttpClient}\PY{p}{.}\PY{n+na}{newBuilder}\PY{p}{(}\PY{p}{)}
\PY{+w}{        }\PY{p}{.}\PY{n+na}{version}\PY{p}{(}\PY{n}{HttpClient}\PY{p}{.}\PY{n+na}{Version}\PY{p}{.}\PY{n+na}{HTTP\PYZus{}2}\PY{p}{)}\PY{+w}{                 }\PY{c+c1}{// prefer HTTP/2, falls back to 1.1}
\PY{+w}{        }\PY{p}{.}\PY{n+na}{connectTimeout}\PY{p}{(}\PY{n}{Duration}\PY{p}{.}\PY{n+na}{ofSeconds}\PY{p}{(}\PY{l+m+mi}{5}\PY{p}{)}\PY{p}{)}\PY{+w}{              }\PY{c+c1}{// connection timeout}
\PY{+w}{        }\PY{p}{.}\PY{n+na}{followRedirects}\PY{p}{(}\PY{n}{HttpClient}\PY{p}{.}\PY{n+na}{Redirect}\PY{p}{.}\PY{n+na}{NORMAL}\PY{p}{)}\PY{+w}{        }\PY{c+c1}{// follow redirects, except HTTPS \PYZhy{}\PYZgt{} HTTP}
\PY{+w}{        }\PY{p}{.}\PY{n+na}{proxy}\PY{p}{(}\PY{n}{ProxySelector}\PY{p}{.}\PY{n+na}{of}\PY{p}{(}\PY{k}{new}\PY{+w}{ }\PY{n}{java}\PY{p}{.}\PY{n+na}{net}\PY{p}{.}\PY{n+na}{InetSocketAddress}\PY{p}{(}\PY{l+s}{\PYZdq{}}\PY{l+s}{proxy.internal}\PY{l+s}{\PYZdq{}}\PY{p}{,}\PY{+w}{ }\PY{l+m+mi}{8080}\PY{p}{)}\PY{p}{)}\PY{p}{)}
\PY{+w}{        }\PY{p}{.}\PY{n+na}{executor}\PY{p}{(}\PY{n}{Executors}\PY{p}{.}\PY{n+na}{newVirtualThreadPerTaskExecutor}\PY{p}{(}\PY{p}{)}\PY{p}{)}\PY{+w}{ }\PY{c+c1}{// custom executor for async callbacks}
\PY{+w}{        }\PY{p}{.}\PY{n+na}{build}\PY{p}{(}\PY{p}{)}\PY{p}{;}
\end{Verbatim}
\end{codebox}

\subsection[Building requests]{Building requests}
\label{h:18:building-requests}
\code{HttpRequest} is also immutable; build one per call with \code{Http\allowbreak{}Request.\allowbreak{}new\allowbreak{}Builder()}. \code{BodyPublishers} supplies the request body for methods like \code{POST} and \code{PUT}.

\begin{codebox}
\begin{Verbatim}[commandchars=\\\{\}]
\PY{k+kn}{import}\PY{+w}{ }\PY{n+nn}{java.net.URI}\PY{p}{;}
\PY{k+kn}{import}\PY{+w}{ }\PY{n+nn}{java.net.http.HttpRequest}\PY{p}{;}
\PY{k+kn}{import}\PY{+w}{ }\PY{n+nn}{java.net.http.HttpRequest.BodyPublishers}\PY{p}{;}
\PY{k+kn}{import}\PY{+w}{ }\PY{n+nn}{java.time.Duration}\PY{p}{;}

\PY{c+c1}{// GET}
\PY{n}{HttpRequest}\PY{+w}{ }\PY{n}{getRequest}\PY{+w}{ }\PY{o}{=}\PY{+w}{ }\PY{n}{HttpRequest}\PY{p}{.}\PY{n+na}{newBuilder}\PY{p}{(}\PY{p}{)}
\PY{+w}{        }\PY{p}{.}\PY{n+na}{uri}\PY{p}{(}\PY{n}{URI}\PY{p}{.}\PY{n+na}{create}\PY{p}{(}\PY{l+s}{\PYZdq{}}\PY{l+s}{https://api.example.com/users/42}\PY{l+s}{\PYZdq{}}\PY{p}{)}\PY{p}{)}
\PY{+w}{        }\PY{p}{.}\PY{n+na}{timeout}\PY{p}{(}\PY{n}{Duration}\PY{p}{.}\PY{n+na}{ofSeconds}\PY{p}{(}\PY{l+m+mi}{10}\PY{p}{)}\PY{p}{)}\PY{+w}{ }\PY{c+c1}{// per\PYZhy{}request timeout, overrides client default}
\PY{+w}{        }\PY{p}{.}\PY{n+na}{header}\PY{p}{(}\PY{l+s}{\PYZdq{}}\PY{l+s}{Accept}\PY{l+s}{\PYZdq{}}\PY{p}{,}\PY{+w}{ }\PY{l+s}{\PYZdq{}}\PY{l+s}{application/json}\PY{l+s}{\PYZdq{}}\PY{p}{)}
\PY{+w}{        }\PY{p}{.}\PY{n+na}{GET}\PY{p}{(}\PY{p}{)}
\PY{+w}{        }\PY{p}{.}\PY{n+na}{build}\PY{p}{(}\PY{p}{)}\PY{p}{;}

\PY{c+c1}{// POST with a JSON body}
\PY{n}{HttpRequest}\PY{+w}{ }\PY{n}{postRequest}\PY{+w}{ }\PY{o}{=}\PY{+w}{ }\PY{n}{HttpRequest}\PY{p}{.}\PY{n+na}{newBuilder}\PY{p}{(}\PY{p}{)}
\PY{+w}{        }\PY{p}{.}\PY{n+na}{uri}\PY{p}{(}\PY{n}{URI}\PY{p}{.}\PY{n+na}{create}\PY{p}{(}\PY{l+s}{\PYZdq{}}\PY{l+s}{https://api.example.com/users}\PY{l+s}{\PYZdq{}}\PY{p}{)}\PY{p}{)}
\PY{+w}{        }\PY{p}{.}\PY{n+na}{timeout}\PY{p}{(}\PY{n}{Duration}\PY{p}{.}\PY{n+na}{ofSeconds}\PY{p}{(}\PY{l+m+mi}{10}\PY{p}{)}\PY{p}{)}
\PY{+w}{        }\PY{p}{.}\PY{n+na}{header}\PY{p}{(}\PY{l+s}{\PYZdq{}}\PY{l+s}{Content\PYZhy{}Type}\PY{l+s}{\PYZdq{}}\PY{p}{,}\PY{+w}{ }\PY{l+s}{\PYZdq{}}\PY{l+s}{application/json}\PY{l+s}{\PYZdq{}}\PY{p}{)}
\PY{+w}{        }\PY{p}{.}\PY{n+na}{POST}\PY{p}{(}\PY{n}{BodyPublishers}\PY{p}{.}\PY{n+na}{ofString}\PY{p}{(}\PY{l+s}{\PYZdq{}}\PY{l+s}{\PYZob{}}\PY{l+s}{\PYZbs{}\PYZdq{}}\PY{l+s}{name}\PY{l+s}{\PYZbs{}\PYZdq{}}\PY{l+s}{:}\PY{l+s}{\PYZbs{}\PYZdq{}}\PY{l+s}{Ada}\PY{l+s}{\PYZbs{}\PYZdq{}}\PY{l+s}{\PYZcb{}}\PY{l+s}{\PYZdq{}}\PY{p}{)}\PY{p}{)}
\PY{+w}{        }\PY{p}{.}\PY{n+na}{build}\PY{p}{(}\PY{p}{)}\PY{p}{;}
\end{Verbatim}
\end{codebox}

\subsection[Sending — sync vs async]{Sending — sync vs async}
\label{h:18:sending--sync-vs-async}
\code{send} blocks the calling thread until the response arrives. \code{sendAsync} returns immediately with a \code{Completable\allowbreak{}Future\textless{}\allowbreak{}Http\allowbreak{}Response\textless{}\allowbreak{}T\textgreater{}\textgreater{}}, letting you compose follow-up work without blocking.

\begin{codebox}
\begin{Verbatim}[commandchars=\\\{\}]
\PY{k+kn}{import}\PY{+w}{ }\PY{n+nn}{java.net.http.HttpResponse}\PY{p}{;}
\PY{k+kn}{import}\PY{+w}{ }\PY{n+nn}{java.net.http.HttpResponse.BodyHandlers}\PY{p}{;}

\PY{c+c1}{// Synchronous}
\PY{n}{HttpResponse}\PY{o}{\PYZlt{}}\PY{n}{String}\PY{o}{\PYZgt{}}\PY{+w}{ }\PY{n}{response}\PY{+w}{ }\PY{o}{=}\PY{+w}{ }\PY{n}{client}\PY{p}{.}\PY{n+na}{send}\PY{p}{(}\PY{n}{getRequest}\PY{p}{,}\PY{+w}{ }\PY{n}{BodyHandlers}\PY{p}{.}\PY{n+na}{ofString}\PY{p}{(}\PY{p}{)}\PY{p}{)}\PY{p}{;}
\PY{n}{System}\PY{p}{.}\PY{n+na}{out}\PY{p}{.}\PY{n+na}{println}\PY{p}{(}\PY{n}{response}\PY{p}{.}\PY{n+na}{statusCode}\PY{p}{(}\PY{p}{)}\PY{+w}{ }\PY{o}{+}\PY{+w}{ }\PY{l+s}{\PYZdq{}}\PY{l+s}{: }\PY{l+s}{\PYZdq{}}\PY{+w}{ }\PY{o}{+}\PY{+w}{ }\PY{n}{response}\PY{p}{.}\PY{n+na}{body}\PY{p}{(}\PY{p}{)}\PY{p}{)}\PY{p}{;}
\end{Verbatim}
\end{codebox}

\begin{codebox}
\begin{Verbatim}[commandchars=\\\{\}]
\PY{c+c1}{// Asynchronous, with a CompletableFuture pipeline}
\PY{n}{client}\PY{p}{.}\PY{n+na}{sendAsync}\PY{p}{(}\PY{n}{getRequest}\PY{p}{,}\PY{+w}{ }\PY{n}{BodyHandlers}\PY{p}{.}\PY{n+na}{ofString}\PY{p}{(}\PY{p}{)}\PY{p}{)}
\PY{+w}{        }\PY{p}{.}\PY{n+na}{thenApply}\PY{p}{(}\PY{n}{HttpResponse}\PY{p}{:}\PY{p}{:}\PY{n}{body}\PY{p}{)}
\PY{+w}{        }\PY{p}{.}\PY{n+na}{thenAccept}\PY{p}{(}\PY{n}{System}\PY{p}{.}\PY{n+na}{out}\PY{p}{:}\PY{p}{:}\PY{n}{println}\PY{p}{)}
\PY{+w}{        }\PY{p}{.}\PY{n+na}{exceptionally}\PY{p}{(}\PY{n}{ex}\PY{+w}{ }\PY{o}{\PYZhy{}}\PY{o}{\PYZgt{}}\PY{+w}{ }\PY{p}{\PYZob{}}
\PY{+w}{            }\PY{n}{System}\PY{p}{.}\PY{n+na}{err}\PY{p}{.}\PY{n+na}{println}\PY{p}{(}\PY{l+s}{\PYZdq{}}\PY{l+s}{Request failed: }\PY{l+s}{\PYZdq{}}\PY{+w}{ }\PY{o}{+}\PY{+w}{ }\PY{n}{ex}\PY{p}{.}\PY{n+na}{getMessage}\PY{p}{(}\PY{p}{)}\PY{p}{)}\PY{p}{;}
\PY{+w}{            }\PY{k}{return}\PY{+w}{ }\PY{k+kc}{null}\PY{p}{;}
\PY{+w}{        }\PY{p}{\PYZcb{}}\PY{p}{)}\PY{p}{;}
\end{Verbatim}
\end{codebox}

\code{BodyHandlers} controls how the response body is consumed:

\begin{itemize}
\item 
\code{Body\allowbreak{}Handlers.\allowbreak{}of\allowbreak{}String()} — whole body as a \code{String}; fine for small JSON responses.

\item 
\code{BodyHandlers.\allowbreak{}ofLines()} — a \code{Stream\textless{}\allowbreak{}String\textgreater{}} of lines; good for line-delimited data.

\item 
\code{Body\allowbreak{}Handlers.\allowbreak{}of\allowbreak{}Input\allowbreak{}Stream()} — raw \code{InputStream}; lets you stream large responses without loading them fully into memory.

\item 
\code{Body\allowbreak{}Handlers.\allowbreak{}of\allowbreak{}File(\allowbreak{}Path)} — writes the body directly to disk; ideal for downloads.

\end{itemize}

\begin{codebox}
\begin{Verbatim}[commandchars=\\\{\}]
\PY{c+c1}{// Streaming a large download straight to disk instead of buffering it in memory}
\PY{n}{HttpRequest}\PY{+w}{ }\PY{n}{downloadRequest}\PY{+w}{ }\PY{o}{=}\PY{+w}{ }\PY{n}{HttpRequest}\PY{p}{.}\PY{n+na}{newBuilder}\PY{p}{(}\PY{p}{)}
\PY{+w}{        }\PY{p}{.}\PY{n+na}{uri}\PY{p}{(}\PY{n}{URI}\PY{p}{.}\PY{n+na}{create}\PY{p}{(}\PY{l+s}{\PYZdq{}}\PY{l+s}{https://example.com/large\PYZhy{}file.zip}\PY{l+s}{\PYZdq{}}\PY{p}{)}\PY{p}{)}
\PY{+w}{        }\PY{p}{.}\PY{n+na}{build}\PY{p}{(}\PY{p}{)}\PY{p}{;}

\PY{n}{HttpResponse}\PY{o}{\PYZlt{}}\PY{n}{java}\PY{p}{.}\PY{n+na}{nio}\PY{p}{.}\PY{n+na}{file}\PY{p}{.}\PY{n+na}{Path}\PY{o}{\PYZgt{}}\PY{+w}{ }\PY{n}{saved}\PY{+w}{ }\PY{o}{=}\PY{+w}{ }\PY{n}{client}\PY{p}{.}\PY{n+na}{send}\PY{p}{(}
\PY{+w}{        }\PY{n}{downloadRequest}\PY{p}{,}\PY{+w}{ }\PY{n}{BodyHandlers}\PY{p}{.}\PY{n+na}{ofFile}\PY{p}{(}\PY{n}{java}\PY{p}{.}\PY{n+na}{nio}\PY{p}{.}\PY{n+na}{file}\PY{p}{.}\PY{n+na}{Path}\PY{p}{.}\PY{n+na}{of}\PY{p}{(}\PY{l+s}{\PYZdq{}}\PY{l+s}{large\PYZhy{}file.zip}\PY{l+s}{\PYZdq{}}\PY{p}{)}\PY{p}{)}\PY{p}{)}\PY{p}{;}
\PY{n}{System}\PY{p}{.}\PY{n+na}{out}\PY{p}{.}\PY{n+na}{println}\PY{p}{(}\PY{l+s}{\PYZdq{}}\PY{l+s}{Saved to }\PY{l+s}{\PYZdq{}}\PY{+w}{ }\PY{o}{+}\PY{+w}{ }\PY{n}{saved}\PY{p}{.}\PY{n+na}{body}\PY{p}{(}\PY{p}{)}\PY{p}{)}\PY{p}{;}
\end{Verbatim}
\end{codebox}

\subsection[HTTP/2 multiplexing]{HTTP/2 multiplexing}
\label{h:18:http2-multiplexing}
HTTP/2 allows many logical requests to share a single TCP connection at once (\textbf{multiplexing}), instead of opening a new connection — or queuing — per request as HTTP/1.1 often does. \code{HttpClient} negotiates HTTP/2 automatically when the server supports it (via ALPN during the TLS handshake); you don’t need to change any request code. This is one reason to prefer a shared \code{HttpClient} instance: reusing it lets the connection pool actually multiplex requests instead of paying a new handshake cost every time.

\subsection[Error handling and status codes]{Error handling and status codes}
\label{h:18:error-handling-and-status-codes}
\code{HttpClient} only throws exceptions for things like connection failures, timeouts, or protocol errors. \textbf{HTTP error status codes (4xx, 5xx) are not exceptions} — they come back as a normal \code{HttpResponse} that you must check yourself.

\begin{codebox}
\begin{Verbatim}[commandchars=\\\{\}]
\PY{c+c1}{// Anti\PYZhy{}pattern: assuming a response without an exception means success}
\PY{n}{HttpResponse}\PY{o}{\PYZlt{}}\PY{n}{String}\PY{o}{\PYZgt{}}\PY{+w}{ }\PY{n}{response}\PY{+w}{ }\PY{o}{=}\PY{+w}{ }\PY{n}{client}\PY{p}{.}\PY{n+na}{send}\PY{p}{(}\PY{n}{getRequest}\PY{p}{,}\PY{+w}{ }\PY{n}{BodyHandlers}\PY{p}{.}\PY{n+na}{ofString}\PY{p}{(}\PY{p}{)}\PY{p}{)}\PY{p}{;}
\PY{n}{processUser}\PY{p}{(}\PY{n}{response}\PY{p}{.}\PY{n+na}{body}\PY{p}{(}\PY{p}{)}\PY{p}{)}\PY{p}{;}\PY{+w}{ }\PY{c+c1}{// might be a 404 or 500 error page, not user JSON!}
\end{Verbatim}
\end{codebox}

\begin{codebox}
\begin{Verbatim}[commandchars=\\\{\}]
\PY{c+c1}{// Fix: check the status code explicitly}
\PY{n}{HttpResponse}\PY{o}{\PYZlt{}}\PY{n}{String}\PY{o}{\PYZgt{}}\PY{+w}{ }\PY{n}{response}\PY{+w}{ }\PY{o}{=}\PY{+w}{ }\PY{n}{client}\PY{p}{.}\PY{n+na}{send}\PY{p}{(}\PY{n}{getRequest}\PY{p}{,}\PY{+w}{ }\PY{n}{BodyHandlers}\PY{p}{.}\PY{n+na}{ofString}\PY{p}{(}\PY{p}{)}\PY{p}{)}\PY{p}{;}
\PY{k+kt}{int}\PY{+w}{ }\PY{n}{status}\PY{+w}{ }\PY{o}{=}\PY{+w}{ }\PY{n}{response}\PY{p}{.}\PY{n+na}{statusCode}\PY{p}{(}\PY{p}{)}\PY{p}{;}
\PY{k}{if}\PY{+w}{ }\PY{p}{(}\PY{n}{status}\PY{+w}{ }\PY{o}{\PYZgt{}}\PY{o}{=}\PY{+w}{ }\PY{l+m+mi}{200}\PY{+w}{ }\PY{o}{\PYZam{}}\PY{o}{\PYZam{}}\PY{+w}{ }\PY{n}{status}\PY{+w}{ }\PY{o}{\PYZlt{}}\PY{+w}{ }\PY{l+m+mi}{300}\PY{p}{)}\PY{+w}{ }\PY{p}{\PYZob{}}
\PY{+w}{    }\PY{n}{processUser}\PY{p}{(}\PY{n}{response}\PY{p}{.}\PY{n+na}{body}\PY{p}{(}\PY{p}{)}\PY{p}{)}\PY{p}{;}
\PY{p}{\PYZcb{}}\PY{+w}{ }\PY{k}{else}\PY{+w}{ }\PY{k}{if}\PY{+w}{ }\PY{p}{(}\PY{n}{status}\PY{+w}{ }\PY{o}{=}\PY{o}{=}\PY{+w}{ }\PY{l+m+mi}{404}\PY{p}{)}\PY{+w}{ }\PY{p}{\PYZob{}}
\PY{+w}{    }\PY{k}{throw}\PY{+w}{ }\PY{k}{new}\PY{+w}{ }\PY{n}{NoSuchElementException}\PY{p}{(}\PY{l+s}{\PYZdq{}}\PY{l+s}{User not found}\PY{l+s}{\PYZdq{}}\PY{p}{)}\PY{p}{;}
\PY{p}{\PYZcb{}}\PY{+w}{ }\PY{k}{else}\PY{+w}{ }\PY{p}{\PYZob{}}
\PY{+w}{    }\PY{k}{throw}\PY{+w}{ }\PY{k}{new}\PY{+w}{ }\PY{n}{IllegalStateException}\PY{p}{(}\PY{l+s}{\PYZdq{}}\PY{l+s}{Unexpected status }\PY{l+s}{\PYZdq{}}\PY{+w}{ }\PY{o}{+}\PY{+w}{ }\PY{n}{status}\PY{+w}{ }\PY{o}{+}\PY{+w}{ }\PY{l+s}{\PYZdq{}}\PY{l+s}{: }\PY{l+s}{\PYZdq{}}\PY{+w}{ }\PY{o}{+}\PY{+w}{ }\PY{n}{response}\PY{p}{.}\PY{n+na}{body}\PY{p}{(}\PY{p}{)}\PY{p}{)}\PY{p}{;}
\PY{p}{\PYZcb{}}
\end{Verbatim}
\end{codebox}

Also handle exceptions from \code{send}/\code{sendAsync}: \code{java.\allowbreak{}net.\allowbreak{}http.\allowbreak{}Http\allowbreak{}Timeout\allowbreak{}Exception} for timeouts, \code{java.\allowbreak{}net.\allowbreak{}Connect\allowbreak{}Exception} for connection failures, and \code{java.\allowbreak{}io.\allowbreak{}IOException} more generally.

\subsection[Custom SSLContext]{Custom SSLContext}
\label{h:18:custom-sslcontext}
By default, \code{HttpClient} uses the JVM’s default \code{SSLContext}, which trusts the standard public \textbf{certificate authorities} (organizations that vouch for a server’s identity — see the Security Basics section). To connect to a server with a private/internal certificate authority, supply a custom \code{SSLContext} built from your own \textbf{trust store} (a keystore holding the certificates you trust) — never a trust-all context (shown as an anti-pattern later in this chapter).

\begin{codebox}
\begin{Verbatim}[commandchars=\\\{\}]
\PY{k+kn}{import}\PY{+w}{ }\PY{n+nn}{javax.net.ssl.SSLContext}\PY{p}{;}
\PY{k+kn}{import}\PY{+w}{ }\PY{n+nn}{javax.net.ssl.TrustManagerFactory}\PY{p}{;}
\PY{k+kn}{import}\PY{+w}{ }\PY{n+nn}{java.security.KeyStore}\PY{p}{;}

\PY{n}{KeyStore}\PY{+w}{ }\PY{n}{trustStore}\PY{+w}{ }\PY{o}{=}\PY{+w}{ }\PY{n}{KeyStore}\PY{p}{.}\PY{n+na}{getInstance}\PY{p}{(}\PY{l+s}{\PYZdq{}}\PY{l+s}{PKCS12}\PY{l+s}{\PYZdq{}}\PY{p}{)}\PY{p}{;}
\PY{k}{try}\PY{+w}{ }\PY{p}{(}\PY{k+kd}{var}\PY{+w}{ }\PY{n}{in}\PY{+w}{ }\PY{o}{=}\PY{+w}{ }\PY{n}{java}\PY{p}{.}\PY{n+na}{nio}\PY{p}{.}\PY{n+na}{file}\PY{p}{.}\PY{n+na}{Files}\PY{p}{.}\PY{n+na}{newInputStream}\PY{p}{(}\PY{n}{java}\PY{p}{.}\PY{n+na}{nio}\PY{p}{.}\PY{n+na}{file}\PY{p}{.}\PY{n+na}{Path}\PY{p}{.}\PY{n+na}{of}\PY{p}{(}\PY{l+s}{\PYZdq{}}\PY{l+s}{internal\PYZhy{}ca.p12}\PY{l+s}{\PYZdq{}}\PY{p}{)}\PY{p}{)}\PY{p}{)}\PY{+w}{ }\PY{p}{\PYZob{}}
\PY{+w}{    }\PY{n}{trustStore}\PY{p}{.}\PY{n+na}{load}\PY{p}{(}\PY{n}{in}\PY{p}{,}\PY{+w}{ }\PY{l+s}{\PYZdq{}}\PY{l+s}{changeit}\PY{l+s}{\PYZdq{}}\PY{p}{.}\PY{n+na}{toCharArray}\PY{p}{(}\PY{p}{)}\PY{p}{)}\PY{p}{;}
\PY{p}{\PYZcb{}}
\PY{n}{TrustManagerFactory}\PY{+w}{ }\PY{n}{tmf}\PY{+w}{ }\PY{o}{=}\PY{+w}{ }\PY{n}{TrustManagerFactory}\PY{p}{.}\PY{n+na}{getInstance}\PY{p}{(}\PY{n}{TrustManagerFactory}\PY{p}{.}\PY{n+na}{getDefaultAlgorithm}\PY{p}{(}\PY{p}{)}\PY{p}{)}\PY{p}{;}
\PY{n}{tmf}\PY{p}{.}\PY{n+na}{init}\PY{p}{(}\PY{n}{trustStore}\PY{p}{)}\PY{p}{;}

\PY{n}{SSLContext}\PY{+w}{ }\PY{n}{sslContext}\PY{+w}{ }\PY{o}{=}\PY{+w}{ }\PY{n}{SSLContext}\PY{p}{.}\PY{n+na}{getInstance}\PY{p}{(}\PY{l+s}{\PYZdq{}}\PY{l+s}{TLSv1.3}\PY{l+s}{\PYZdq{}}\PY{p}{)}\PY{p}{;}
\PY{n}{sslContext}\PY{p}{.}\PY{n+na}{init}\PY{p}{(}\PY{k+kc}{null}\PY{p}{,}\PY{+w}{ }\PY{n}{tmf}\PY{p}{.}\PY{n+na}{getTrustManagers}\PY{p}{(}\PY{p}{)}\PY{p}{,}\PY{+w}{ }\PY{k+kc}{null}\PY{p}{)}\PY{p}{;}

\PY{n}{HttpClient}\PY{+w}{ }\PY{n}{secureClient}\PY{+w}{ }\PY{o}{=}\PY{+w}{ }\PY{n}{HttpClient}\PY{p}{.}\PY{n+na}{newBuilder}\PY{p}{(}\PY{p}{)}
\PY{+w}{        }\PY{p}{.}\PY{n+na}{sslContext}\PY{p}{(}\PY{n}{sslContext}\PY{p}{)}
\PY{+w}{        }\PY{p}{.}\PY{n+na}{build}\PY{p}{(}\PY{p}{)}\PY{p}{;}
\end{Verbatim}
\end{codebox}

\subsection[WebSocket support]{WebSocket support}
\label{h:18:websocket-support}
\code{HttpClient} also builds \textbf{WebSocket} connections (a persistent, full-duplex connection for real-time, bidirectional messaging — unlike request/response HTTP) via \code{Http\allowbreak{}Client.\allowbreak{}new\allowbreak{}Web\allowbreak{}Socket\allowbreak{}Builder()}.

\begin{codebox}
\begin{Verbatim}[commandchars=\\\{\}]
\PY{k+kn}{import}\PY{+w}{ }\PY{n+nn}{java.net.http.WebSocket}\PY{p}{;}
\PY{k+kn}{import}\PY{+w}{ }\PY{n+nn}{java.net.URI}\PY{p}{;}
\PY{k+kn}{import}\PY{+w}{ }\PY{n+nn}{java.util.concurrent.CompletionStage}\PY{p}{;}

\PY{n}{WebSocket}\PY{p}{.}\PY{n+na}{Listener}\PY{+w}{ }\PY{n}{listener}\PY{+w}{ }\PY{o}{=}\PY{+w}{ }\PY{k}{new}\PY{+w}{ }\PY{n}{WebSocket}\PY{p}{.}\PY{n+na}{Listener}\PY{p}{(}\PY{p}{)}\PY{+w}{ }\PY{p}{\PYZob{}}
\PY{+w}{    }\PY{n+nd}{@Override}
\PY{+w}{    }\PY{k+kd}{public}\PY{+w}{ }\PY{n}{CompletionStage}\PY{o}{\PYZlt{}}\PY{o}{?}\PY{o}{\PYZgt{}}\PY{+w}{ }\PY{n+nf}{onText}\PY{p}{(}\PY{n}{WebSocket}\PY{+w}{ }\PY{n}{webSocket}\PY{p}{,}\PY{+w}{ }\PY{n}{CharSequence}\PY{+w}{ }\PY{n}{data}\PY{p}{,}\PY{+w}{ }\PY{k+kt}{boolean}\PY{+w}{ }\PY{n}{last}\PY{p}{)}\PY{+w}{ }\PY{p}{\PYZob{}}
\PY{+w}{        }\PY{n}{System}\PY{p}{.}\PY{n+na}{out}\PY{p}{.}\PY{n+na}{println}\PY{p}{(}\PY{l+s}{\PYZdq{}}\PY{l+s}{Received: }\PY{l+s}{\PYZdq{}}\PY{+w}{ }\PY{o}{+}\PY{+w}{ }\PY{n}{data}\PY{p}{)}\PY{p}{;}
\PY{+w}{        }\PY{n}{webSocket}\PY{p}{.}\PY{n+na}{request}\PY{p}{(}\PY{l+m+mi}{1}\PY{p}{)}\PY{p}{;}
\PY{+w}{        }\PY{k}{return}\PY{+w}{ }\PY{k+kc}{null}\PY{p}{;}
\PY{+w}{    }\PY{p}{\PYZcb{}}
\PY{p}{\PYZcb{}}\PY{p}{;}

\PY{n}{WebSocket}\PY{+w}{ }\PY{n}{webSocket}\PY{+w}{ }\PY{o}{=}\PY{+w}{ }\PY{n}{client}\PY{p}{.}\PY{n+na}{newWebSocketBuilder}\PY{p}{(}\PY{p}{)}
\PY{+w}{        }\PY{p}{.}\PY{n+na}{buildAsync}\PY{p}{(}\PY{n}{URI}\PY{p}{.}\PY{n+na}{create}\PY{p}{(}\PY{l+s}{\PYZdq{}}\PY{l+s}{wss://example.com/socket}\PY{l+s}{\PYZdq{}}\PY{p}{)}\PY{p}{,}\PY{+w}{ }\PY{n}{listener}\PY{p}{)}
\PY{+w}{        }\PY{p}{.}\PY{n+na}{join}\PY{p}{(}\PY{p}{)}\PY{p}{;}

\PY{n}{webSocket}\PY{p}{.}\PY{n+na}{sendText}\PY{p}{(}\PY{l+s}{\PYZdq{}}\PY{l+s}{hello}\PY{l+s}{\PYZdq{}}\PY{p}{,}\PY{+w}{ }\PY{k+kc}{true}\PY{p}{)}\PY{p}{;}
\end{Verbatim}
\end{codebox}

\subsection[AutoCloseable in JDK 21+]{AutoCloseable in JDK 21+}
\label{h:18:autocloseable-in-jdk-21}
Since Java 21, \code{HttpClient} implements \code{AutoCloseable}, so you can use it in a try-with-resources block to release its resources (connection pools, background threads) deterministically — useful for short-lived clients, though for a long-lived shared client you’d typically just let it live for the application’s lifetime and never close it.

\begin{codebox}
\begin{Verbatim}[commandchars=\\\{\}]
\PY{k}{try}\PY{+w}{ }\PY{p}{(}\PY{n}{HttpClient}\PY{+w}{ }\PY{n}{shortLivedClient}\PY{+w}{ }\PY{o}{=}\PY{+w}{ }\PY{n}{HttpClient}\PY{p}{.}\PY{n+na}{newHttpClient}\PY{p}{(}\PY{p}{)}\PY{p}{)}\PY{+w}{ }\PY{p}{\PYZob{}}
\PY{+w}{    }\PY{n}{HttpResponse}\PY{o}{\PYZlt{}}\PY{n}{String}\PY{o}{\PYZgt{}}\PY{+w}{ }\PY{n}{response}\PY{+w}{ }\PY{o}{=}\PY{+w}{ }\PY{n}{shortLivedClient}\PY{p}{.}\PY{n+na}{send}\PY{p}{(}\PY{n}{getRequest}\PY{p}{,}\PY{+w}{ }\PY{n}{BodyHandlers}\PY{p}{.}\PY{n+na}{ofString}\PY{p}{(}\PY{p}{)}\PY{p}{)}\PY{p}{;}
\PY{+w}{    }\PY{n}{System}\PY{p}{.}\PY{n+na}{out}\PY{p}{.}\PY{n+na}{println}\PY{p}{(}\PY{n}{response}\PY{p}{.}\PY{n+na}{body}\PY{p}{(}\PY{p}{)}\PY{p}{)}\PY{p}{;}
\PY{p}{\PYZcb{}}\PY{+w}{ }\PY{c+c1}{// client resources released here}
\end{Verbatim}
\end{codebox}

\section[Process API]{Process API}
\label{h:18:process-api}
Sometimes your Java application needs to run an external program — a shell script, \code{ffmpeg}, \code{git}, a native tool. The \textbf{Process API} (\code{ProcessBuilder}, \code{Process}, \code{ProcessHandle}) lets you launch and manage external processes.

\subsection[ProcessBuilder — launching a process safely]{ProcessBuilder — launching a process safely}
\label{h:18:processbuilder--launching-a-process-safely}
\code{ProcessBuilder} takes the command as a \textbf{list of strings}: the program name, then each argument as a separate list element. This is critical for security — see the anti-pattern below.

\begin{codebox}
\begin{Verbatim}[commandchars=\\\{\}]
\PY{c+c1}{// Anti\PYZhy{}pattern: building a shell command string from user input — COMMAND INJECTION}
\PY{n}{String}\PY{+w}{ }\PY{n}{filename}\PY{+w}{ }\PY{o}{=}\PY{+w}{ }\PY{n}{userInput}\PY{p}{;}\PY{+w}{ }\PY{c+c1}{// attacker\PYZhy{}controlled, e.g. \PYZdq{}; rm \PYZhy{}rf / \PYZsh{}\PYZdq{}}
\PY{n}{Process}\PY{+w}{ }\PY{n}{process}\PY{+w}{ }\PY{o}{=}\PY{+w}{ }\PY{n}{Runtime}\PY{p}{.}\PY{n+na}{getRuntime}\PY{p}{(}\PY{p}{)}\PY{p}{.}\PY{n+na}{exec}\PY{p}{(}\PY{l+s}{\PYZdq{}}\PY{l+s}{ls \PYZhy{}l }\PY{l+s}{\PYZdq{}}\PY{+w}{ }\PY{o}{+}\PY{+w}{ }\PY{n}{filename}\PY{p}{)}\PY{p}{;}
\end{Verbatim}
\end{codebox}

If \code{filename} is \code{\textquotedbl{}foo;\allowbreak{}~\allowbreak{}rm~\allowbreak{}-\allowbreak{}rf~\allowbreak{}/\allowbreak{}\textquotedbl{}}, and this string is ever passed through a shell (e.g. via \code{bash~\allowbreak{}-\allowbreak{}c}), the shell happily runs \code{rm~\allowbreak{}-\allowbreak{}rf~\allowbreak{}/} as a second command. This class of bug is called \textbf{command injection}: attacker input escapes its intended role as \emph{data} and gets interpreted as \emph{code} (here, shell syntax).

\begin{codebox}
\begin{Verbatim}[commandchars=\\\{\}]
\PY{c+c1}{// Fix: pass arguments as a list — no shell parsing, no injection}
\PY{k+kn}{import}\PY{+w}{ }\PY{n+nn}{java.util.List}\PY{p}{;}

\PY{n}{ProcessBuilder}\PY{+w}{ }\PY{n}{builder}\PY{+w}{ }\PY{o}{=}\PY{+w}{ }\PY{k}{new}\PY{+w}{ }\PY{n}{ProcessBuilder}\PY{p}{(}\PY{n}{List}\PY{p}{.}\PY{n+na}{of}\PY{p}{(}\PY{l+s}{\PYZdq{}}\PY{l+s}{ls}\PY{l+s}{\PYZdq{}}\PY{p}{,}\PY{+w}{ }\PY{l+s}{\PYZdq{}}\PY{l+s}{\PYZhy{}l}\PY{l+s}{\PYZdq{}}\PY{p}{,}\PY{+w}{ }\PY{n}{filename}\PY{p}{)}\PY{p}{)}\PY{p}{;}
\PY{n}{Process}\PY{+w}{ }\PY{n}{process}\PY{+w}{ }\PY{o}{=}\PY{+w}{ }\PY{n}{builder}\PY{p}{.}\PY{n+na}{start}\PY{p}{(}\PY{p}{)}\PY{p}{;}
\end{Verbatim}
\end{codebox}

Because each argument is a separate list element, \code{filename} is passed to the \code{ls} program as a single literal argument — \code{;}, \code{\&\&}, backticks, and other shell metacharacters have no special meaning here. \textbf{Never} build a command by concatenating a string and handing it to a shell (\code{bash~\allowbreak{}-\allowbreak{}c~\allowbreak{}\textquotedbl{}...\allowbreak{}\textquotedbl{}}, \code{sh~\allowbreak{}-\allowbreak{}c~\allowbreak{}\textquotedbl{}...\allowbreak{}\textquotedbl{}}, or string-based \code{exec}) when any part of it comes from user input.

\subsection[Redirecting IO]{Redirecting IO}
\label{h:18:redirecting-io}
By default, a child process’s stdin/stdout/stderr are separate pipes your Java code must read or write. \code{redirectOutput}, \code{redirectError}, and \code{redirectInput} connect them directly to files, or \code{inheritIO()} connects them straight to the parent JVM’s own console — handy for CLI tools where you just want to see the child’s output live.

\begin{codebox}
\begin{Verbatim}[commandchars=\\\{\}]
\PY{n}{ProcessBuilder}\PY{+w}{ }\PY{n}{builder}\PY{+w}{ }\PY{o}{=}\PY{+w}{ }\PY{k}{new}\PY{+w}{ }\PY{n}{ProcessBuilder}\PY{p}{(}\PY{l+s}{\PYZdq{}}\PY{l+s}{ping}\PY{l+s}{\PYZdq{}}\PY{p}{,}\PY{+w}{ }\PY{l+s}{\PYZdq{}}\PY{l+s}{\PYZhy{}c}\PY{l+s}{\PYZdq{}}\PY{p}{,}\PY{+w}{ }\PY{l+s}{\PYZdq{}}\PY{l+s}{3}\PY{l+s}{\PYZdq{}}\PY{p}{,}\PY{+w}{ }\PY{l+s}{\PYZdq{}}\PY{l+s}{example.com}\PY{l+s}{\PYZdq{}}\PY{p}{)}\PY{p}{;}
\PY{n}{builder}\PY{p}{.}\PY{n+na}{inheritIO}\PY{p}{(}\PY{p}{)}\PY{p}{;}\PY{+w}{ }\PY{c+c1}{// child\PYZsq{}s stdout/stderr appear directly in this program\PYZsq{}s console}
\PY{n}{Process}\PY{+w}{ }\PY{n}{process}\PY{+w}{ }\PY{o}{=}\PY{+w}{ }\PY{n}{builder}\PY{p}{.}\PY{n+na}{start}\PY{p}{(}\PY{p}{)}\PY{p}{;}
\PY{k+kt}{int}\PY{+w}{ }\PY{n}{exitCode}\PY{+w}{ }\PY{o}{=}\PY{+w}{ }\PY{n}{process}\PY{p}{.}\PY{n+na}{waitFor}\PY{p}{(}\PY{p}{)}\PY{p}{;}
\end{Verbatim}
\end{codebox}

\begin{codebox}
\begin{Verbatim}[commandchars=\\\{\}]
\PY{n}{ProcessBuilder}\PY{+w}{ }\PY{n}{builder}\PY{+w}{ }\PY{o}{=}\PY{+w}{ }\PY{k}{new}\PY{+w}{ }\PY{n}{ProcessBuilder}\PY{p}{(}\PY{l+s}{\PYZdq{}}\PY{l+s}{mytool}\PY{l+s}{\PYZdq{}}\PY{p}{,}\PY{+w}{ }\PY{l+s}{\PYZdq{}}\PY{l+s}{\PYZhy{}\PYZhy{}verbose}\PY{l+s}{\PYZdq{}}\PY{p}{)}\PY{p}{;}
\PY{n}{builder}\PY{p}{.}\PY{n+na}{redirectOutput}\PY{p}{(}\PY{k}{new}\PY{+w}{ }\PY{n}{java}\PY{p}{.}\PY{n+na}{io}\PY{p}{.}\PY{n+na}{File}\PY{p}{(}\PY{l+s}{\PYZdq{}}\PY{l+s}{output.log}\PY{l+s}{\PYZdq{}}\PY{p}{)}\PY{p}{)}\PY{p}{;}
\PY{n}{builder}\PY{p}{.}\PY{n+na}{redirectError}\PY{p}{(}\PY{n}{ProcessBuilder}\PY{p}{.}\PY{n+na}{Redirect}\PY{p}{.}\PY{n+na}{INHERIT}\PY{p}{)}\PY{p}{;}\PY{+w}{ }\PY{c+c1}{// errors still show in console}
\PY{n}{Process}\PY{+w}{ }\PY{n}{process}\PY{+w}{ }\PY{o}{=}\PY{+w}{ }\PY{n}{builder}\PY{p}{.}\PY{n+na}{start}\PY{p}{(}\PY{p}{)}\PY{p}{;}
\end{Verbatim}
\end{codebox}

\subsection[Reading stdout/stderr without deadlocking]{Reading stdout/stderr without deadlocking}
\label{h:18:reading-stdoutstderr-without-deadlocking}
If you don’t redirect a stream and don’t read it, the child process’s output buffer can fill up. Once full, the child blocks trying to write more output — and if your Java code is meanwhile blocked in \code{waitFor()} waiting for the child to exit, you get a \textbf{deadlock}: each side waiting on the other forever.

\begin{codebox}
\begin{Verbatim}[commandchars=\\\{\}]
\PY{c+c1}{// Anti\PYZhy{}pattern: reading stdout only after waitFor() — can deadlock if stderr fills up first}
\PY{n}{Process}\PY{+w}{ }\PY{n}{process}\PY{+w}{ }\PY{o}{=}\PY{+w}{ }\PY{k}{new}\PY{+w}{ }\PY{n}{ProcessBuilder}\PY{p}{(}\PY{l+s}{\PYZdq{}}\PY{l+s}{noisy\PYZhy{}tool}\PY{l+s}{\PYZdq{}}\PY{p}{)}\PY{p}{.}\PY{n+na}{start}\PY{p}{(}\PY{p}{)}\PY{p}{;}
\PY{k+kt}{int}\PY{+w}{ }\PY{n}{exitCode}\PY{+w}{ }\PY{o}{=}\PY{+w}{ }\PY{n}{process}\PY{p}{.}\PY{n+na}{waitFor}\PY{p}{(}\PY{p}{)}\PY{p}{;}\PY{+w}{ }\PY{c+c1}{// may hang forever if noisy\PYZhy{}tool blocks on a full stderr buffer}
\PY{n}{String}\PY{+w}{ }\PY{n}{output}\PY{+w}{ }\PY{o}{=}\PY{+w}{ }\PY{k}{new}\PY{+w}{ }\PY{n}{String}\PY{p}{(}\PY{n}{process}\PY{p}{.}\PY{n+na}{getInputStream}\PY{p}{(}\PY{p}{)}\PY{p}{.}\PY{n+na}{readAllBytes}\PY{p}{(}\PY{p}{)}\PY{p}{)}\PY{p}{;}
\end{Verbatim}
\end{codebox}

\begin{codebox}
\begin{Verbatim}[commandchars=\\\{\}]
\PY{c+c1}{// Fix: drain both streams concurrently, in separate threads, before waiting}
\PY{n}{Process}\PY{+w}{ }\PY{n}{process}\PY{+w}{ }\PY{o}{=}\PY{+w}{ }\PY{k}{new}\PY{+w}{ }\PY{n}{ProcessBuilder}\PY{p}{(}\PY{l+s}{\PYZdq{}}\PY{l+s}{noisy\PYZhy{}tool}\PY{l+s}{\PYZdq{}}\PY{p}{)}\PY{p}{.}\PY{n+na}{start}\PY{p}{(}\PY{p}{)}\PY{p}{;}

\PY{n}{Thread}\PY{+w}{ }\PY{n}{outputReader}\PY{+w}{ }\PY{o}{=}\PY{+w}{ }\PY{n}{Thread}\PY{p}{.}\PY{n+na}{ofVirtual}\PY{p}{(}\PY{p}{)}\PY{p}{.}\PY{n+na}{start}\PY{p}{(}\PY{p}{(}\PY{p}{)}\PY{+w}{ }\PY{o}{\PYZhy{}}\PY{o}{\PYZgt{}}\PY{+w}{ }\PY{p}{\PYZob{}}
\PY{+w}{    }\PY{k}{try}\PY{+w}{ }\PY{p}{(}\PY{k+kd}{var}\PY{+w}{ }\PY{n}{in}\PY{+w}{ }\PY{o}{=}\PY{+w}{ }\PY{n}{process}\PY{p}{.}\PY{n+na}{getInputStream}\PY{p}{(}\PY{p}{)}\PY{p}{)}\PY{+w}{ }\PY{p}{\PYZob{}}
\PY{+w}{        }\PY{n}{System}\PY{p}{.}\PY{n+na}{out}\PY{p}{.}\PY{n+na}{write}\PY{p}{(}\PY{n}{in}\PY{p}{.}\PY{n+na}{readAllBytes}\PY{p}{(}\PY{p}{)}\PY{p}{)}\PY{p}{;}
\PY{+w}{    }\PY{p}{\PYZcb{}}\PY{+w}{ }\PY{k}{catch}\PY{+w}{ }\PY{p}{(}\PY{n}{java}\PY{p}{.}\PY{n+na}{io}\PY{p}{.}\PY{n+na}{IOException}\PY{+w}{ }\PY{n}{e}\PY{p}{)}\PY{+w}{ }\PY{p}{\PYZob{}}\PY{+w}{ }\PY{c+cm}{/* log */}\PY{+w}{ }\PY{p}{\PYZcb{}}
\PY{p}{\PYZcb{}}\PY{p}{)}\PY{p}{;}
\PY{n}{Thread}\PY{+w}{ }\PY{n}{errorReader}\PY{+w}{ }\PY{o}{=}\PY{+w}{ }\PY{n}{Thread}\PY{p}{.}\PY{n+na}{ofVirtual}\PY{p}{(}\PY{p}{)}\PY{p}{.}\PY{n+na}{start}\PY{p}{(}\PY{p}{(}\PY{p}{)}\PY{+w}{ }\PY{o}{\PYZhy{}}\PY{o}{\PYZgt{}}\PY{+w}{ }\PY{p}{\PYZob{}}
\PY{+w}{    }\PY{k}{try}\PY{+w}{ }\PY{p}{(}\PY{k+kd}{var}\PY{+w}{ }\PY{n}{err}\PY{+w}{ }\PY{o}{=}\PY{+w}{ }\PY{n}{process}\PY{p}{.}\PY{n+na}{getErrorStream}\PY{p}{(}\PY{p}{)}\PY{p}{)}\PY{+w}{ }\PY{p}{\PYZob{}}
\PY{+w}{        }\PY{n}{System}\PY{p}{.}\PY{n+na}{err}\PY{p}{.}\PY{n+na}{write}\PY{p}{(}\PY{n}{err}\PY{p}{.}\PY{n+na}{readAllBytes}\PY{p}{(}\PY{p}{)}\PY{p}{)}\PY{p}{;}
\PY{+w}{    }\PY{p}{\PYZcb{}}\PY{+w}{ }\PY{k}{catch}\PY{+w}{ }\PY{p}{(}\PY{n}{java}\PY{p}{.}\PY{n+na}{io}\PY{p}{.}\PY{n+na}{IOException}\PY{+w}{ }\PY{n}{e}\PY{p}{)}\PY{+w}{ }\PY{p}{\PYZob{}}\PY{+w}{ }\PY{c+cm}{/* log */}\PY{+w}{ }\PY{p}{\PYZcb{}}
\PY{p}{\PYZcb{}}\PY{p}{)}\PY{p}{;}

\PY{k+kt}{int}\PY{+w}{ }\PY{n}{exitCode}\PY{+w}{ }\PY{o}{=}\PY{+w}{ }\PY{n}{process}\PY{p}{.}\PY{n+na}{waitFor}\PY{p}{(}\PY{p}{)}\PY{p}{;}
\PY{n}{outputReader}\PY{p}{.}\PY{n+na}{join}\PY{p}{(}\PY{p}{)}\PY{p}{;}
\PY{n}{errorReader}\PY{p}{.}\PY{n+na}{join}\PY{p}{(}\PY{p}{)}\PY{p}{;}
\end{Verbatim}
\end{codebox}

An easier fix, when you don’t need the streams programmatically, is \code{Process\allowbreak{}Builder.\allowbreak{}redirect\allowbreak{}Error\allowbreak{}Stream(\allowbreak{}true)} (merges stderr into stdout — one stream to read) or \code{inheritIO()} — both avoid the deadlock entirely by not leaving an unread pipe.

\subsection[waitFor with a timeout]{waitFor with a timeout}
\label{h:18:waitfor-with-a-timeout}
Waiting forever for a child process is as risky as a network call with no timeout — a hung external tool can hang your JVM’s thread indefinitely.

\begin{codebox}
\begin{Verbatim}[commandchars=\\\{\}]
\PY{k+kn}{import}\PY{+w}{ }\PY{n+nn}{java.util.concurrent.TimeUnit}\PY{p}{;}

\PY{n}{Process}\PY{+w}{ }\PY{n}{process}\PY{+w}{ }\PY{o}{=}\PY{+w}{ }\PY{k}{new}\PY{+w}{ }\PY{n}{ProcessBuilder}\PY{p}{(}\PY{l+s}{\PYZdq{}}\PY{l+s}{slow\PYZhy{}tool}\PY{l+s}{\PYZdq{}}\PY{p}{)}\PY{p}{.}\PY{n+na}{start}\PY{p}{(}\PY{p}{)}\PY{p}{;}
\PY{k+kt}{boolean}\PY{+w}{ }\PY{n}{finished}\PY{+w}{ }\PY{o}{=}\PY{+w}{ }\PY{n}{process}\PY{p}{.}\PY{n+na}{waitFor}\PY{p}{(}\PY{l+m+mi}{30}\PY{p}{,}\PY{+w}{ }\PY{n}{TimeUnit}\PY{p}{.}\PY{n+na}{SECONDS}\PY{p}{)}\PY{p}{;}
\PY{k}{if}\PY{+w}{ }\PY{p}{(}\PY{o}{!}\PY{n}{finished}\PY{p}{)}\PY{+w}{ }\PY{p}{\PYZob{}}
\PY{+w}{    }\PY{n}{process}\PY{p}{.}\PY{n+na}{destroyForcibly}\PY{p}{(}\PY{p}{)}\PY{p}{;}\PY{+w}{ }\PY{c+c1}{// give up and kill it}
\PY{+w}{    }\PY{k}{throw}\PY{+w}{ }\PY{k}{new}\PY{+w}{ }\PY{n}{IllegalStateException}\PY{p}{(}\PY{l+s}{\PYZdq{}}\PY{l+s}{slow\PYZhy{}tool did not finish in time}\PY{l+s}{\PYZdq{}}\PY{p}{)}\PY{p}{;}
\PY{p}{\PYZcb{}}
\PY{k+kt}{int}\PY{+w}{ }\PY{n}{exitCode}\PY{+w}{ }\PY{o}{=}\PY{+w}{ }\PY{n}{process}\PY{p}{.}\PY{n+na}{exitValue}\PY{p}{(}\PY{p}{)}\PY{p}{;}
\end{Verbatim}
\end{codebox}

\subsection[ProcessHandle (JDK 9+)]{ProcessHandle (JDK 9+)}
\label{h:18:processhandle-jdk-9}
\code{ProcessHandle} gives you information about \emph{any} process — not just ones you started — by PID (process ID), including the current JVM’s own process, without needing native code.

\begin{codebox}
\begin{Verbatim}[commandchars=\\\{\}]
\PY{n}{ProcessHandle}\PY{+w}{ }\PY{n}{current}\PY{+w}{ }\PY{o}{=}\PY{+w}{ }\PY{n}{ProcessHandle}\PY{p}{.}\PY{n+na}{current}\PY{p}{(}\PY{p}{)}\PY{p}{;}
\PY{n}{System}\PY{p}{.}\PY{n+na}{out}\PY{p}{.}\PY{n+na}{println}\PY{p}{(}\PY{l+s}{\PYZdq{}}\PY{l+s}{My PID: }\PY{l+s}{\PYZdq{}}\PY{+w}{ }\PY{o}{+}\PY{+w}{ }\PY{n}{current}\PY{p}{.}\PY{n+na}{pid}\PY{p}{(}\PY{p}{)}\PY{p}{)}\PY{p}{;}

\PY{n}{Process}\PY{+w}{ }\PY{n}{child}\PY{+w}{ }\PY{o}{=}\PY{+w}{ }\PY{k}{new}\PY{+w}{ }\PY{n}{ProcessBuilder}\PY{p}{(}\PY{l+s}{\PYZdq{}}\PY{l+s}{sleep}\PY{l+s}{\PYZdq{}}\PY{p}{,}\PY{+w}{ }\PY{l+s}{\PYZdq{}}\PY{l+s}{60}\PY{l+s}{\PYZdq{}}\PY{p}{)}\PY{p}{.}\PY{n+na}{start}\PY{p}{(}\PY{p}{)}\PY{p}{;}
\PY{n}{ProcessHandle}\PY{+w}{ }\PY{n}{handle}\PY{+w}{ }\PY{o}{=}\PY{+w}{ }\PY{n}{child}\PY{p}{.}\PY{n+na}{toHandle}\PY{p}{(}\PY{p}{)}\PY{p}{;}

\PY{n}{handle}\PY{p}{.}\PY{n+na}{info}\PY{p}{(}\PY{p}{)}\PY{p}{.}\PY{n+na}{command}\PY{p}{(}\PY{p}{)}\PY{p}{.}\PY{n+na}{ifPresent}\PY{p}{(}\PY{n}{cmd}\PY{+w}{ }\PY{o}{\PYZhy{}}\PY{o}{\PYZgt{}}\PY{+w}{ }\PY{n}{System}\PY{p}{.}\PY{n+na}{out}\PY{p}{.}\PY{n+na}{println}\PY{p}{(}\PY{l+s}{\PYZdq{}}\PY{l+s}{Command: }\PY{l+s}{\PYZdq{}}\PY{+w}{ }\PY{o}{+}\PY{+w}{ }\PY{n}{cmd}\PY{p}{)}\PY{p}{)}\PY{p}{;}
\PY{n}{handle}\PY{p}{.}\PY{n+na}{info}\PY{p}{(}\PY{p}{)}\PY{p}{.}\PY{n+na}{startInstant}\PY{p}{(}\PY{p}{)}\PY{p}{.}\PY{n+na}{ifPresent}\PY{p}{(}\PY{n}{t}\PY{+w}{ }\PY{o}{\PYZhy{}}\PY{o}{\PYZgt{}}\PY{+w}{ }\PY{n}{System}\PY{p}{.}\PY{n+na}{out}\PY{p}{.}\PY{n+na}{println}\PY{p}{(}\PY{l+s}{\PYZdq{}}\PY{l+s}{Started: }\PY{l+s}{\PYZdq{}}\PY{+w}{ }\PY{o}{+}\PY{+w}{ }\PY{n}{t}\PY{p}{)}\PY{p}{)}\PY{p}{;}

\PY{c+c1}{// List all child processes of the current JVM}
\PY{n}{ProcessHandle}\PY{p}{.}\PY{n+na}{current}\PY{p}{(}\PY{p}{)}\PY{p}{.}\PY{n+na}{children}\PY{p}{(}\PY{p}{)}\PY{p}{.}\PY{n+na}{forEach}\PY{p}{(}\PY{n}{h}\PY{+w}{ }\PY{o}{\PYZhy{}}\PY{o}{\PYZgt{}}
\PY{+w}{        }\PY{n}{System}\PY{p}{.}\PY{n+na}{out}\PY{p}{.}\PY{n+na}{println}\PY{p}{(}\PY{l+s}{\PYZdq{}}\PY{l+s}{Child PID: }\PY{l+s}{\PYZdq{}}\PY{+w}{ }\PY{o}{+}\PY{+w}{ }\PY{n}{h}\PY{p}{.}\PY{n+na}{pid}\PY{p}{(}\PY{p}{)}\PY{p}{)}\PY{p}{)}\PY{p}{;}

\PY{c+c1}{// React when the process exits, without blocking the current thread}
\PY{n}{handle}\PY{p}{.}\PY{n+na}{onExit}\PY{p}{(}\PY{p}{)}\PY{p}{.}\PY{n+na}{thenAccept}\PY{p}{(}\PY{n}{h}\PY{+w}{ }\PY{o}{\PYZhy{}}\PY{o}{\PYZgt{}}
\PY{+w}{        }\PY{n}{System}\PY{p}{.}\PY{n+na}{out}\PY{p}{.}\PY{n+na}{println}\PY{p}{(}\PY{l+s}{\PYZdq{}}\PY{l+s}{Process }\PY{l+s}{\PYZdq{}}\PY{+w}{ }\PY{o}{+}\PY{+w}{ }\PY{n}{h}\PY{p}{.}\PY{n+na}{pid}\PY{p}{(}\PY{p}{)}\PY{+w}{ }\PY{o}{+}\PY{+w}{ }\PY{l+s}{\PYZdq{}}\PY{l+s}{ exited with status }\PY{l+s}{\PYZdq{}}\PY{+w}{ }\PY{o}{+}\PY{+w}{ }\PY{n}{h}\PY{p}{.}\PY{n+na}{exitValue}\PY{p}{(}\PY{p}{)}\PY{p}{)}\PY{p}{)}\PY{p}{;}
\end{Verbatim}
\end{codebox}

\subsection[destroy vs destroyForcibly]{destroy vs destroyForcibly}
\label{h:18:destroy-vs-destroyforcibly}
\code{destroy()} requests a \textbf{graceful} termination — on most platforms this sends \code{SIGTERM} on Unix-like systems, letting the process clean up (close files, flush buffers) before exiting. \code{destroyForcibly()} requests an \textbf{immediate, forceful} kill (\code{SIGKILL}-like), which the process cannot intercept or clean up after.

\begin{codebox}
\begin{Verbatim}[commandchars=\\\{\}]
\PY{n}{Process}\PY{+w}{ }\PY{n}{process}\PY{+w}{ }\PY{o}{=}\PY{+w}{ }\PY{k}{new}\PY{+w}{ }\PY{n}{ProcessBuilder}\PY{p}{(}\PY{l+s}{\PYZdq{}}\PY{l+s}{my\PYZhy{}daemon}\PY{l+s}{\PYZdq{}}\PY{p}{)}\PY{p}{.}\PY{n+na}{start}\PY{p}{(}\PY{p}{)}\PY{p}{;}
\PY{n}{process}\PY{p}{.}\PY{n+na}{destroy}\PY{p}{(}\PY{p}{)}\PY{p}{;}\PY{+w}{ }\PY{c+c1}{// ask nicely first}
\PY{k}{if}\PY{+w}{ }\PY{p}{(}\PY{o}{!}\PY{n}{process}\PY{p}{.}\PY{n+na}{waitFor}\PY{p}{(}\PY{l+m+mi}{5}\PY{p}{,}\PY{+w}{ }\PY{n}{TimeUnit}\PY{p}{.}\PY{n+na}{SECONDS}\PY{p}{)}\PY{p}{)}\PY{+w}{ }\PY{p}{\PYZob{}}
\PY{+w}{    }\PY{n}{process}\PY{p}{.}\PY{n+na}{destroyForcibly}\PY{p}{(}\PY{p}{)}\PY{p}{;}\PY{+w}{ }\PY{c+c1}{// it ignored us — force it}
\PY{p}{\PYZcb{}}
\end{Verbatim}
\end{codebox}

Prefer \code{destroy()} first and escalate to \code{destroyForcibly()} only if the process doesn’t respond in a reasonable time, so well-behaved processes get a chance to shut down cleanly.

\section[Security Basics]{Security Basics}
\label{h:18:security-basics}
Security in Java isn’t a separate skill from writing correct code — most security bugs are just correctness bugs where the attacker controls the input. This section is deliberately defensive: every insecure pattern shown here is labelled \textbf{Anti-pattern} and immediately followed by the fix, for recognizing and rejecting these patterns in review — not for building them.

\subsection[Principle of least privilege]{Principle of least privilege}
\label{h:18:principle-of-least-privilege}
Give code, users, and processes the \textbf{minimum} access they need to do their job, nothing more. A reporting service that only reads data should use a database account with read-only permissions, not the same admin account used for migrations. A file-processing service should run as a low-privilege OS user, not \code{root}. When something does go wrong — and eventually something will — least privilege limits the damage.

\subsection[Input validation]{Input validation}
\label{h:18:input-validation}
\textbf{Never trust input from outside your program’s trust boundary} — HTTP requests, file uploads, environment variables, even data from another internal service you don’t fully control. Validate shape, type, range, and length \emph{before} using the value, ideally as close to the entry point as possible.

\begin{codebox}
\begin{Verbatim}[commandchars=\\\{\}]
\PY{c+c1}{// Anti\PYZhy{}pattern: using a user\PYZhy{}supplied value without validation}
\PY{k+kt}{int}\PY{+w}{ }\PY{n}{pageSize}\PY{+w}{ }\PY{o}{=}\PY{+w}{ }\PY{n}{Integer}\PY{p}{.}\PY{n+na}{parseInt}\PY{p}{(}\PY{n}{request}\PY{p}{.}\PY{n+na}{getParameter}\PY{p}{(}\PY{l+s}{\PYZdq{}}\PY{l+s}{pageSize}\PY{l+s}{\PYZdq{}}\PY{p}{)}\PY{p}{)}\PY{p}{;}
\PY{n}{List}\PY{o}{\PYZlt{}}\PY{n}{Item}\PY{o}{\PYZgt{}}\PY{+w}{ }\PY{n}{items}\PY{+w}{ }\PY{o}{=}\PY{+w}{ }\PY{n}{repository}\PY{p}{.}\PY{n+na}{fetch}\PY{p}{(}\PY{n}{pageSize}\PY{p}{)}\PY{p}{;}\PY{+w}{ }\PY{c+c1}{// attacker sends pageSize=2000000000}

\PY{c+c1}{// Fix: validate range before use}
\PY{k+kt}{int}\PY{+w}{ }\PY{n}{pageSize}\PY{+w}{ }\PY{o}{=}\PY{+w}{ }\PY{n}{Integer}\PY{p}{.}\PY{n+na}{parseInt}\PY{p}{(}\PY{n}{request}\PY{p}{.}\PY{n+na}{getParameter}\PY{p}{(}\PY{l+s}{\PYZdq{}}\PY{l+s}{pageSize}\PY{l+s}{\PYZdq{}}\PY{p}{)}\PY{p}{)}\PY{p}{;}
\PY{k}{if}\PY{+w}{ }\PY{p}{(}\PY{n}{pageSize}\PY{+w}{ }\PY{o}{\PYZlt{}}\PY{+w}{ }\PY{l+m+mi}{1}\PY{+w}{ }\PY{o}{|}\PY{o}{|}\PY{+w}{ }\PY{n}{pageSize}\PY{+w}{ }\PY{o}{\PYZgt{}}\PY{+w}{ }\PY{l+m+mi}{100}\PY{p}{)}\PY{+w}{ }\PY{p}{\PYZob{}}
\PY{+w}{    }\PY{k}{throw}\PY{+w}{ }\PY{k}{new}\PY{+w}{ }\PY{n}{IllegalArgumentException}\PY{p}{(}\PY{l+s}{\PYZdq{}}\PY{l+s}{pageSize must be between 1 and 100}\PY{l+s}{\PYZdq{}}\PY{p}{)}\PY{p}{;}
\PY{p}{\PYZcb{}}
\PY{n}{List}\PY{o}{\PYZlt{}}\PY{n}{Item}\PY{o}{\PYZgt{}}\PY{+w}{ }\PY{n}{items}\PY{+w}{ }\PY{o}{=}\PY{+w}{ }\PY{n}{repository}\PY{p}{.}\PY{n+na}{fetch}\PY{p}{(}\PY{n}{pageSize}\PY{p}{)}\PY{p}{;}
\end{Verbatim}
\end{codebox}

\subsection[SQL injection]{SQL injection}
\label{h:18:sql-injection}
\textbf{SQL injection} happens when user input is concatenated directly into a SQL query string, letting an attacker change the query’s \emph{meaning} instead of just supplying a value.

\begin{codebox}
\begin{Verbatim}[commandchars=\\\{\}]
\PY{c+c1}{// Anti\PYZhy{}pattern: string concatenation builds attacker\PYZhy{}controlled SQL}
\PY{n}{String}\PY{+w}{ }\PY{n}{username}\PY{+w}{ }\PY{o}{=}\PY{+w}{ }\PY{n}{request}\PY{p}{.}\PY{n+na}{getParameter}\PY{p}{(}\PY{l+s}{\PYZdq{}}\PY{l+s}{username}\PY{l+s}{\PYZdq{}}\PY{p}{)}\PY{p}{;}\PY{+w}{ }\PY{c+c1}{// e.g. \PYZdq{}\PYZsq{} OR \PYZsq{}1\PYZsq{}=\PYZsq{}1\PYZdq{}}
\PY{n}{String}\PY{+w}{ }\PY{n}{sql}\PY{+w}{ }\PY{o}{=}\PY{+w}{ }\PY{l+s}{\PYZdq{}}\PY{l+s}{SELECT * FROM users WHERE username = \PYZsq{}}\PY{l+s}{\PYZdq{}}\PY{+w}{ }\PY{o}{+}\PY{+w}{ }\PY{n}{username}\PY{+w}{ }\PY{o}{+}\PY{+w}{ }\PY{l+s}{\PYZdq{}}\PY{l+s}{\PYZsq{}}\PY{l+s}{\PYZdq{}}\PY{p}{;}
\PY{n}{ResultSet}\PY{+w}{ }\PY{n}{rs}\PY{+w}{ }\PY{o}{=}\PY{+w}{ }\PY{n}{statement}\PY{p}{.}\PY{n+na}{executeQuery}\PY{p}{(}\PY{n}{sql}\PY{p}{)}\PY{p}{;}\PY{+w}{ }\PY{c+c1}{// returns ALL users, auth bypass}
\end{Verbatim}
\end{codebox}

\begin{codebox}
\begin{Verbatim}[commandchars=\\\{\}]
\PY{c+c1}{// Fix: PreparedStatement with bound parameters — input is always treated as data}
\PY{n}{String}\PY{+w}{ }\PY{n}{sql}\PY{+w}{ }\PY{o}{=}\PY{+w}{ }\PY{l+s}{\PYZdq{}}\PY{l+s}{SELECT * FROM users WHERE username = ?}\PY{l+s}{\PYZdq{}}\PY{p}{;}
\PY{k}{try}\PY{+w}{ }\PY{p}{(}\PY{n}{PreparedStatement}\PY{+w}{ }\PY{n}{ps}\PY{+w}{ }\PY{o}{=}\PY{+w}{ }\PY{n}{connection}\PY{p}{.}\PY{n+na}{prepareStatement}\PY{p}{(}\PY{n}{sql}\PY{p}{)}\PY{p}{)}\PY{+w}{ }\PY{p}{\PYZob{}}
\PY{+w}{    }\PY{n}{ps}\PY{p}{.}\PY{n+na}{setString}\PY{p}{(}\PY{l+m+mi}{1}\PY{p}{,}\PY{+w}{ }\PY{n}{username}\PY{p}{)}\PY{p}{;}
\PY{+w}{    }\PY{n}{ResultSet}\PY{+w}{ }\PY{n}{rs}\PY{+w}{ }\PY{o}{=}\PY{+w}{ }\PY{n}{ps}\PY{p}{.}\PY{n+na}{executeQuery}\PY{p}{(}\PY{p}{)}\PY{p}{;}
\PY{p}{\PYZcb{}}
\end{Verbatim}
\end{codebox}

\code{PreparedStatement} sends the query structure and the parameter values separately to the database, so the database never re-parses attacker input as SQL syntax.

\subsection[Command injection]{Command injection}
\label{h:18:command-injection}
Covered in detail in the Process API section above: never build a shell command string from untrusted input. Always pass arguments as a list to \code{ProcessBuilder}, and avoid invoking a shell (\code{bash~\allowbreak{}-\allowbreak{}c}, \code{sh~\allowbreak{}-\allowbreak{}c}) with concatenated input at all.

\subsection[Path traversal]{Path traversal}
\label{h:18:path-traversal}
\textbf{Path traversal} (or “directory traversal”) happens when user input controls a file path and contains sequences like \code{../} that escape the intended directory.

\begin{codebox}
\begin{Verbatim}[commandchars=\\\{\}]
\PY{c+c1}{// Anti\PYZhy{}pattern: filename from the client used directly in a file path}
\PY{n}{String}\PY{+w}{ }\PY{n}{filename}\PY{+w}{ }\PY{o}{=}\PY{+w}{ }\PY{n}{request}\PY{p}{.}\PY{n+na}{getParameter}\PY{p}{(}\PY{l+s}{\PYZdq{}}\PY{l+s}{file}\PY{l+s}{\PYZdq{}}\PY{p}{)}\PY{p}{;}\PY{+w}{ }\PY{c+c1}{// e.g. \PYZdq{}../../etc/passwd\PYZdq{}}
\PY{n}{File}\PY{+w}{ }\PY{n}{file}\PY{+w}{ }\PY{o}{=}\PY{+w}{ }\PY{k}{new}\PY{+w}{ }\PY{n}{File}\PY{p}{(}\PY{l+s}{\PYZdq{}}\PY{l+s}{/var/app/uploads/}\PY{l+s}{\PYZdq{}}\PY{+w}{ }\PY{o}{+}\PY{+w}{ }\PY{n}{filename}\PY{p}{)}\PY{p}{;}
\PY{k+kt}{byte}\PY{o}{[}\PY{o}{]}\PY{+w}{ }\PY{n}{contents}\PY{+w}{ }\PY{o}{=}\PY{+w}{ }\PY{n}{Files}\PY{p}{.}\PY{n+na}{readAllBytes}\PY{p}{(}\PY{n}{file}\PY{p}{.}\PY{n+na}{toPath}\PY{p}{(}\PY{p}{)}\PY{p}{)}\PY{p}{;}\PY{+w}{ }\PY{c+c1}{// reads outside uploads/}
\end{Verbatim}
\end{codebox}

\begin{codebox}
\begin{Verbatim}[commandchars=\\\{\}]
\PY{c+c1}{// Fix: resolve, normalize, and verify the result stays within the allowed directory}
\PY{n}{Path}\PY{+w}{ }\PY{n}{uploadsDir}\PY{+w}{ }\PY{o}{=}\PY{+w}{ }\PY{n}{Path}\PY{p}{.}\PY{n+na}{of}\PY{p}{(}\PY{l+s}{\PYZdq{}}\PY{l+s}{/var/app/uploads}\PY{l+s}{\PYZdq{}}\PY{p}{)}\PY{p}{.}\PY{n+na}{toRealPath}\PY{p}{(}\PY{p}{)}\PY{p}{;}
\PY{n}{Path}\PY{+w}{ }\PY{n}{requested}\PY{+w}{ }\PY{o}{=}\PY{+w}{ }\PY{n}{uploadsDir}\PY{p}{.}\PY{n+na}{resolve}\PY{p}{(}\PY{n}{filename}\PY{p}{)}\PY{p}{.}\PY{n+na}{normalize}\PY{p}{(}\PY{p}{)}\PY{p}{;}

\PY{k}{if}\PY{+w}{ }\PY{p}{(}\PY{o}{!}\PY{n}{requested}\PY{p}{.}\PY{n+na}{startsWith}\PY{p}{(}\PY{n}{uploadsDir}\PY{p}{)}\PY{p}{)}\PY{+w}{ }\PY{p}{\PYZob{}}
\PY{+w}{    }\PY{k}{throw}\PY{+w}{ }\PY{k}{new}\PY{+w}{ }\PY{n}{SecurityException}\PY{p}{(}\PY{l+s}{\PYZdq{}}\PY{l+s}{Invalid file path: }\PY{l+s}{\PYZdq{}}\PY{+w}{ }\PY{o}{+}\PY{+w}{ }\PY{n}{filename}\PY{p}{)}\PY{p}{;}
\PY{p}{\PYZcb{}}
\PY{k+kt}{byte}\PY{o}{[}\PY{o}{]}\PY{+w}{ }\PY{n}{contents}\PY{+w}{ }\PY{o}{=}\PY{+w}{ }\PY{n}{Files}\PY{p}{.}\PY{n+na}{readAllBytes}\PY{p}{(}\PY{n}{requested}\PY{p}{)}\PY{p}{;}
\end{Verbatim}
\end{codebox}

\code{resolve} + \code{normalize} collapses \code{..} segments, and the \code{startsWith} check ensures the final path is still inside the intended base directory before touching the filesystem.

\subsection[XXE — XML External Entity injection]{XXE — XML External Entity injection}
\label{h:18:xxe--xml-external-entity-injection}
XML parsers, by default, may resolve \textbf{external entities} — references inside an XML document that pull in content from a file or URL. If an attacker controls the XML being parsed, they can use this to read local files or make the server issue outbound requests (\textbf{XXE}, XML External Entity injection).

\begin{codebox}
\begin{Verbatim}[commandchars=\\\{\}]
\PY{c+c1}{// Anti\PYZhy{}pattern: default DocumentBuilderFactory resolves external entities}
\PY{n}{String}\PY{+w}{ }\PY{n}{maliciousXml}\PY{+w}{ }\PY{o}{=}\PY{+w}{ }\PY{l+s}{\PYZdq{}\PYZdq{}\PYZdq{}}
\PY{l+s}{    \PYZlt{}?xml version=}\PY{l+s}{\PYZdq{}}\PY{l+s}{1.0}\PY{l+s}{\PYZdq{}}\PY{l+s}{?\PYZgt{}}
\PY{l+s}{    \PYZlt{}!DOCTYPE foo [ \PYZlt{}!ENTITY xxe SYSTEM }\PY{l+s}{\PYZdq{}}\PY{l+s}{file:///etc/passwd}\PY{l+s}{\PYZdq{}}\PY{l+s}{\PYZgt{} ]\PYZgt{}}
\PY{l+s}{    \PYZlt{}foo\PYZgt{}\PYZam{}xxe;\PYZlt{}/foo\PYZgt{}}
\PY{l+s}{    }\PY{l+s}{\PYZdq{}\PYZdq{}\PYZdq{}}\PY{p}{;}

\PY{n}{DocumentBuilderFactory}\PY{+w}{ }\PY{n}{factory}\PY{+w}{ }\PY{o}{=}\PY{+w}{ }\PY{n}{DocumentBuilderFactory}\PY{p}{.}\PY{n+na}{newInstance}\PY{p}{(}\PY{p}{)}\PY{p}{;}
\PY{n}{DocumentBuilder}\PY{+w}{ }\PY{n}{builder}\PY{+w}{ }\PY{o}{=}\PY{+w}{ }\PY{n}{factory}\PY{p}{.}\PY{n+na}{newDocumentBuilder}\PY{p}{(}\PY{p}{)}\PY{p}{;}
\PY{n}{Document}\PY{+w}{ }\PY{n}{doc}\PY{+w}{ }\PY{o}{=}\PY{+w}{ }\PY{n}{builder}\PY{p}{.}\PY{n+na}{parse}\PY{p}{(}\PY{k}{new}\PY{+w}{ }\PY{n}{ByteArrayInputStream}\PY{p}{(}\PY{n}{maliciousXml}\PY{p}{.}\PY{n+na}{getBytes}\PY{p}{(}\PY{p}{)}\PY{p}{)}\PY{p}{)}\PY{p}{;}\PY{+w}{ }\PY{c+c1}{// leaks /etc/passwd}
\end{Verbatim}
\end{codebox}

\begin{codebox}
\begin{Verbatim}[commandchars=\\\{\}]
\PY{c+c1}{// Fix: disable DTDs and external entities explicitly}
\PY{n}{DocumentBuilderFactory}\PY{+w}{ }\PY{n}{factory}\PY{+w}{ }\PY{o}{=}\PY{+w}{ }\PY{n}{DocumentBuilderFactory}\PY{p}{.}\PY{n+na}{newInstance}\PY{p}{(}\PY{p}{)}\PY{p}{;}
\PY{n}{factory}\PY{p}{.}\PY{n+na}{setFeature}\PY{p}{(}\PY{l+s}{\PYZdq{}}\PY{l+s}{http://apache.org/xml/features/disallow\PYZhy{}doctype\PYZhy{}decl}\PY{l+s}{\PYZdq{}}\PY{p}{,}\PY{+w}{ }\PY{k+kc}{true}\PY{p}{)}\PY{p}{;}
\PY{n}{factory}\PY{p}{.}\PY{n+na}{setFeature}\PY{p}{(}\PY{l+s}{\PYZdq{}}\PY{l+s}{http://xml.org/sax/features/external\PYZhy{}general\PYZhy{}entities}\PY{l+s}{\PYZdq{}}\PY{p}{,}\PY{+w}{ }\PY{k+kc}{false}\PY{p}{)}\PY{p}{;}
\PY{n}{factory}\PY{p}{.}\PY{n+na}{setFeature}\PY{p}{(}\PY{l+s}{\PYZdq{}}\PY{l+s}{http://xml.org/sax/features/external\PYZhy{}parameter\PYZhy{}entities}\PY{l+s}{\PYZdq{}}\PY{p}{,}\PY{+w}{ }\PY{k+kc}{false}\PY{p}{)}\PY{p}{;}
\PY{n}{factory}\PY{p}{.}\PY{n+na}{setXIncludeAware}\PY{p}{(}\PY{k+kc}{false}\PY{p}{)}\PY{p}{;}
\PY{n}{factory}\PY{p}{.}\PY{n+na}{setExpandEntityReferences}\PY{p}{(}\PY{k+kc}{false}\PY{p}{)}\PY{p}{;}

\PY{n}{DocumentBuilder}\PY{+w}{ }\PY{n}{builder}\PY{+w}{ }\PY{o}{=}\PY{+w}{ }\PY{n}{factory}\PY{p}{.}\PY{n+na}{newDocumentBuilder}\PY{p}{(}\PY{p}{)}\PY{p}{;}
\PY{n}{Document}\PY{+w}{ }\PY{n}{doc}\PY{+w}{ }\PY{o}{=}\PY{+w}{ }\PY{n}{builder}\PY{p}{.}\PY{n+na}{parse}\PY{p}{(}\PY{k}{new}\PY{+w}{ }\PY{n}{ByteArrayInputStream}\PY{p}{(}\PY{n}{maliciousXml}\PY{p}{.}\PY{n+na}{getBytes}\PY{p}{(}\PY{p}{)}\PY{p}{)}\PY{p}{)}\PY{p}{;}\PY{+w}{ }\PY{c+c1}{// throws, DTD rejected}
\end{Verbatim}
\end{codebox}

The same idea applies to \code{SAXParserFactory}, \code{XMLInputFactory} (StAX), and \code{TransformerFactory} — each has equivalent feature flags to disable DTDs and external entities. When parsing XML from an untrusted source, disable them all.

\subsection[Insecure deserialization]{Insecure deserialization}
\label{h:18:insecure-deserialization}
Java’s built-in serialization (\code{Object\allowbreak{}Input\allowbreak{}Stream.\allowbreak{}read\allowbreak{}Object()}) reconstructs arbitrary objects from a byte stream, including running code in constructors, \code{readObject} overrides, and finalizers along the way. Deserializing \textbf{untrusted} data can let an attacker construct object graphs that trigger unintended — sometimes arbitrary — code execution (a well-known class of “gadget chain” attacks).

\begin{codebox}
\begin{Verbatim}[commandchars=\\\{\}]
\PY{c+c1}{// Anti\PYZhy{}pattern: deserializing bytes received from a network client}
\PY{n}{ObjectInputStream}\PY{+w}{ }\PY{n}{in}\PY{+w}{ }\PY{o}{=}\PY{+w}{ }\PY{k}{new}\PY{+w}{ }\PY{n}{ObjectInputStream}\PY{p}{(}\PY{n}{socket}\PY{p}{.}\PY{n+na}{getInputStream}\PY{p}{(}\PY{p}{)}\PY{p}{)}\PY{p}{;}
\PY{n}{Object}\PY{+w}{ }\PY{n}{payload}\PY{+w}{ }\PY{o}{=}\PY{+w}{ }\PY{n}{in}\PY{p}{.}\PY{n+na}{readObject}\PY{p}{(}\PY{p}{)}\PY{p}{;}\PY{+w}{ }\PY{c+c1}{// attacker fully controls the class graph being built}
\end{Verbatim}
\end{codebox}

\begin{codebox}
\begin{Verbatim}[commandchars=\\\{\}]
\PY{c+c1}{// Fix: use a safe, explicit data format instead (e.g. JSON with a known schema),}
\PY{c+c1}{// or if Java serialization is unavoidable, filter allowed classes}
\PY{n}{ObjectInputStream}\PY{+w}{ }\PY{n}{in}\PY{+w}{ }\PY{o}{=}\PY{+w}{ }\PY{k}{new}\PY{+w}{ }\PY{n}{ObjectInputStream}\PY{p}{(}\PY{n}{socket}\PY{p}{.}\PY{n+na}{getInputStream}\PY{p}{(}\PY{p}{)}\PY{p}{)}\PY{p}{;}
\PY{n}{in}\PY{p}{.}\PY{n+na}{setObjectInputFilter}\PY{p}{(}\PY{n}{filterInfo}\PY{+w}{ }\PY{o}{\PYZhy{}}\PY{o}{\PYZgt{}}\PY{+w}{ }\PY{p}{\PYZob{}}
\PY{+w}{    }\PY{n}{Class}\PY{o}{\PYZlt{}}\PY{o}{?}\PY{o}{\PYZgt{}}\PY{+w}{ }\PY{n}{clazz}\PY{+w}{ }\PY{o}{=}\PY{+w}{ }\PY{n}{filterInfo}\PY{p}{.}\PY{n+na}{serialClass}\PY{p}{(}\PY{p}{)}\PY{p}{;}
\PY{+w}{    }\PY{k}{if}\PY{+w}{ }\PY{p}{(}\PY{n}{clazz}\PY{+w}{ }\PY{o}{!}\PY{o}{=}\PY{+w}{ }\PY{k+kc}{null}\PY{+w}{ }\PY{o}{\PYZam{}}\PY{o}{\PYZam{}}\PY{+w}{ }\PY{o}{!}\PY{n}{clazz}\PY{p}{.}\PY{n+na}{getName}\PY{p}{(}\PY{p}{)}\PY{p}{.}\PY{n+na}{startsWith}\PY{p}{(}\PY{l+s}{\PYZdq{}}\PY{l+s}{com.myapp.dto.}\PY{l+s}{\PYZdq{}}\PY{p}{)}\PY{p}{)}\PY{+w}{ }\PY{p}{\PYZob{}}
\PY{+w}{        }\PY{k}{return}\PY{+w}{ }\PY{n}{ObjectInputFilter}\PY{p}{.}\PY{n+na}{Status}\PY{p}{.}\PY{n+na}{REJECTED}\PY{p}{;}
\PY{+w}{    }\PY{p}{\PYZcb{}}
\PY{+w}{    }\PY{k}{return}\PY{+w}{ }\PY{n}{ObjectInputFilter}\PY{p}{.}\PY{n+na}{Status}\PY{p}{.}\PY{n+na}{UNDECIDED}\PY{p}{;}
\PY{p}{\PYZcb{}}\PY{p}{)}\PY{p}{;}
\PY{n}{Object}\PY{+w}{ }\PY{n}{payload}\PY{+w}{ }\PY{o}{=}\PY{+w}{ }\PY{n}{in}\PY{p}{.}\PY{n+na}{readObject}\PY{p}{(}\PY{p}{)}\PY{p}{;}\PY{+w}{ }\PY{c+c1}{// only com.myapp.dto.* classes are allowed through}
\end{Verbatim}
\end{codebox}

The best fix is usually to avoid Java’s native serialization for untrusted data entirely — use a safe format like JSON with a strict schema and a library that only builds known DTOs, never arbitrary classes.

\subsection[SSRF — Server-Side Request Forgery]{SSRF — Server-Side Request Forgery}
\label{h:18:ssrf--server-side-request-forgery}
\textbf{SSRF} happens when your server makes an outbound HTTP request to a URL that is (directly or indirectly) controlled by the attacker, letting them reach internal-only systems (metadata endpoints, admin panels, internal services) that would otherwise be unreachable from outside.

\begin{codebox}
\begin{Verbatim}[commandchars=\\\{\}]
\PY{c+c1}{// Anti\PYZhy{}pattern: fetching a user\PYZhy{}supplied URL with no restriction}
\PY{n}{String}\PY{+w}{ }\PY{n}{imageUrl}\PY{+w}{ }\PY{o}{=}\PY{+w}{ }\PY{n}{request}\PY{p}{.}\PY{n+na}{getParameter}\PY{p}{(}\PY{l+s}{\PYZdq{}}\PY{l+s}{avatarUrl}\PY{l+s}{\PYZdq{}}\PY{p}{)}\PY{p}{;}\PY{+w}{ }\PY{c+c1}{// e.g. \PYZdq{}http://169.254.169.254/latest/meta\PYZhy{}data/\PYZdq{}}
\PY{n}{HttpResponse}\PY{o}{\PYZlt{}}\PY{k+kt}{byte}\PY{o}{[}\PY{o}{]}\PY{o}{\PYZgt{}}\PY{+w}{ }\PY{n}{response}\PY{+w}{ }\PY{o}{=}\PY{+w}{ }\PY{n}{client}\PY{p}{.}\PY{n+na}{send}\PY{p}{(}
\PY{+w}{        }\PY{n}{HttpRequest}\PY{p}{.}\PY{n+na}{newBuilder}\PY{p}{(}\PY{n}{URI}\PY{p}{.}\PY{n+na}{create}\PY{p}{(}\PY{n}{imageUrl}\PY{p}{)}\PY{p}{)}\PY{p}{.}\PY{n+na}{build}\PY{p}{(}\PY{p}{)}\PY{p}{,}\PY{+w}{ }\PY{n}{BodyHandlers}\PY{p}{.}\PY{n+na}{ofByteArray}\PY{p}{(}\PY{p}{)}\PY{p}{)}\PY{p}{;}
\end{Verbatim}
\end{codebox}

\begin{codebox}
\begin{Verbatim}[commandchars=\\\{\}]
\PY{c+c1}{// Fix: validate against an allow\PYZhy{}list of hosts/schemes, and resolve+check the IP before connecting}
\PY{k+kd}{private}\PY{+w}{ }\PY{k+kd}{static}\PY{+w}{ }\PY{k+kd}{final}\PY{+w}{ }\PY{n}{Set}\PY{o}{\PYZlt{}}\PY{n}{String}\PY{o}{\PYZgt{}}\PY{+w}{ }\PY{n}{ALLOWED\PYZus{}HOSTS}\PY{+w}{ }\PY{o}{=}\PY{+w}{ }\PY{n}{Set}\PY{p}{.}\PY{n+na}{of}\PY{p}{(}\PY{l+s}{\PYZdq{}}\PY{l+s}{images.trusted\PYZhy{}cdn.com}\PY{l+s}{\PYZdq{}}\PY{p}{)}\PY{p}{;}

\PY{n}{URI}\PY{+w}{ }\PY{n}{uri}\PY{+w}{ }\PY{o}{=}\PY{+w}{ }\PY{n}{URI}\PY{p}{.}\PY{n+na}{create}\PY{p}{(}\PY{n}{imageUrl}\PY{p}{)}\PY{p}{;}
\PY{k}{if}\PY{+w}{ }\PY{p}{(}\PY{o}{!}\PY{l+s}{\PYZdq{}}\PY{l+s}{https}\PY{l+s}{\PYZdq{}}\PY{p}{.}\PY{n+na}{equals}\PY{p}{(}\PY{n}{uri}\PY{p}{.}\PY{n+na}{getScheme}\PY{p}{(}\PY{p}{)}\PY{p}{)}\PY{+w}{ }\PY{o}{|}\PY{o}{|}\PY{+w}{ }\PY{o}{!}\PY{n}{ALLOWED\PYZus{}HOSTS}\PY{p}{.}\PY{n+na}{contains}\PY{p}{(}\PY{n}{uri}\PY{p}{.}\PY{n+na}{getHost}\PY{p}{(}\PY{p}{)}\PY{p}{)}\PY{p}{)}\PY{+w}{ }\PY{p}{\PYZob{}}
\PY{+w}{    }\PY{k}{throw}\PY{+w}{ }\PY{k}{new}\PY{+w}{ }\PY{n}{SecurityException}\PY{p}{(}\PY{l+s}{\PYZdq{}}\PY{l+s}{URL not allowed: }\PY{l+s}{\PYZdq{}}\PY{+w}{ }\PY{o}{+}\PY{+w}{ }\PY{n}{imageUrl}\PY{p}{)}\PY{p}{;}
\PY{p}{\PYZcb{}}
\PY{n}{InetAddress}\PY{+w}{ }\PY{n}{resolved}\PY{+w}{ }\PY{o}{=}\PY{+w}{ }\PY{n}{InetAddress}\PY{p}{.}\PY{n+na}{getByName}\PY{p}{(}\PY{n}{uri}\PY{p}{.}\PY{n+na}{getHost}\PY{p}{(}\PY{p}{)}\PY{p}{)}\PY{p}{;}
\PY{k}{if}\PY{+w}{ }\PY{p}{(}\PY{n}{resolved}\PY{p}{.}\PY{n+na}{isLoopbackAddress}\PY{p}{(}\PY{p}{)}\PY{+w}{ }\PY{o}{|}\PY{o}{|}\PY{+w}{ }\PY{n}{resolved}\PY{p}{.}\PY{n+na}{isSiteLocalAddress}\PY{p}{(}\PY{p}{)}\PY{+w}{ }\PY{o}{|}\PY{o}{|}\PY{+w}{ }\PY{n}{resolved}\PY{p}{.}\PY{n+na}{isLinkLocalAddress}\PY{p}{(}\PY{p}{)}\PY{p}{)}\PY{+w}{ }\PY{p}{\PYZob{}}
\PY{+w}{    }\PY{k}{throw}\PY{+w}{ }\PY{k}{new}\PY{+w}{ }\PY{n}{SecurityException}\PY{p}{(}\PY{l+s}{\PYZdq{}}\PY{l+s}{URL resolves to a disallowed internal address}\PY{l+s}{\PYZdq{}}\PY{p}{)}\PY{p}{;}
\PY{p}{\PYZcb{}}
\PY{n}{HttpResponse}\PY{o}{\PYZlt{}}\PY{k+kt}{byte}\PY{o}{[}\PY{o}{]}\PY{o}{\PYZgt{}}\PY{+w}{ }\PY{n}{response}\PY{+w}{ }\PY{o}{=}\PY{+w}{ }\PY{n}{client}\PY{p}{.}\PY{n+na}{send}\PY{p}{(}
\PY{+w}{        }\PY{n}{HttpRequest}\PY{p}{.}\PY{n+na}{newBuilder}\PY{p}{(}\PY{n}{uri}\PY{p}{)}\PY{p}{.}\PY{n+na}{build}\PY{p}{(}\PY{p}{)}\PY{p}{,}\PY{+w}{ }\PY{n}{BodyHandlers}\PY{p}{.}\PY{n+na}{ofByteArray}\PY{p}{(}\PY{p}{)}\PY{p}{)}\PY{p}{;}
\end{Verbatim}
\end{codebox}

Note that an allow-list check on the hostname alone isn’t quite enough — DNS can be manipulated (\textbf{DNS rebinding}) between the check and the actual request, so high-value systems also pin the resolved IP or route such requests through an isolated network egress proxy.

\subsection[Hard-coded secrets]{Hard-coded secrets}
\label{h:18:hard-coded-secrets}
Credentials, API keys, and encryption keys committed to source code end up in version control history forever, visible to anyone with repo access (and often, accidentally, on a public GitHub mirror).

\begin{codebox}
\begin{Verbatim}[commandchars=\\\{\}]
\PY{c+c1}{// Anti\PYZhy{}pattern: secret baked into source code}
\PY{k+kd}{private}\PY{+w}{ }\PY{k+kd}{static}\PY{+w}{ }\PY{k+kd}{final}\PY{+w}{ }\PY{n}{String}\PY{+w}{ }\PY{n}{DB\PYZus{}PASSWORD}\PY{+w}{ }\PY{o}{=}\PY{+w}{ }\PY{l+s}{\PYZdq{}}\PY{l+s}{SuperSecret123!}\PY{l+s}{\PYZdq{}}\PY{p}{;}
\PY{k+kd}{private}\PY{+w}{ }\PY{k+kd}{static}\PY{+w}{ }\PY{k+kd}{final}\PY{+w}{ }\PY{n}{String}\PY{+w}{ }\PY{n}{API\PYZus{}KEY}\PY{+w}{ }\PY{o}{=}\PY{+w}{ }\PY{l+s}{\PYZdq{}}\PY{l+s}{sk\PYZus{}live\PYZus{}51H8x...}\PY{l+s}{\PYZdq{}}\PY{p}{;}
\end{Verbatim}
\end{codebox}

\begin{codebox}
\begin{Verbatim}[commandchars=\\\{\}]
\PY{c+c1}{// Fix: load secrets from the environment or a secrets manager at runtime}
\PY{n}{String}\PY{+w}{ }\PY{n}{dbPassword}\PY{+w}{ }\PY{o}{=}\PY{+w}{ }\PY{n}{System}\PY{p}{.}\PY{n+na}{getenv}\PY{p}{(}\PY{l+s}{\PYZdq{}}\PY{l+s}{DB\PYZus{}PASSWORD}\PY{l+s}{\PYZdq{}}\PY{p}{)}\PY{p}{;}
\PY{n}{String}\PY{+w}{ }\PY{n}{apiKey}\PY{+w}{ }\PY{o}{=}\PY{+w}{ }\PY{n}{secretsManagerClient}\PY{p}{.}\PY{n+na}{getSecret}\PY{p}{(}\PY{l+s}{\PYZdq{}}\PY{l+s}{payments/api\PYZhy{}key}\PY{l+s}{\PYZdq{}}\PY{p}{)}\PY{p}{;}
\end{Verbatim}
\end{codebox}

\subsection[Logging sensitive data]{Logging sensitive data}
\label{h:18:logging-sensitive-data}
Logs are often shipped to third-party aggregators, kept for years, and readable by many engineers — a bad place for passwords, tokens, or full card numbers to end up.

\begin{codebox}
\begin{Verbatim}[commandchars=\\\{\}]
\PY{c+c1}{// Anti\PYZhy{}pattern: logging the full request, including sensitive fields}
\PY{n}{log}\PY{p}{.}\PY{n+na}{info}\PY{p}{(}\PY{l+s}{\PYZdq{}}\PY{l+s}{Login attempt: \PYZob{}\PYZcb{}}\PY{l+s}{\PYZdq{}}\PY{p}{,}\PY{+w}{ }\PY{n}{loginRequest}\PY{p}{)}\PY{p}{;}\PY{+w}{ }\PY{c+c1}{// toString() includes the raw password}

\PY{c+c1}{// Fix: log only what\PYZsq{}s needed, and redact the rest}
\PY{n}{log}\PY{p}{.}\PY{n+na}{info}\PY{p}{(}\PY{l+s}{\PYZdq{}}\PY{l+s}{Login attempt for user: \PYZob{}\PYZcb{}}\PY{l+s}{\PYZdq{}}\PY{p}{,}\PY{+w}{ }\PY{n}{loginRequest}\PY{p}{.}\PY{n+na}{username}\PY{p}{(}\PY{p}{)}\PY{p}{)}\PY{p}{;}
\end{Verbatim}
\end{codebox}

\subsection[char[] vs String for passwords]{char[] vs String for passwords}
\label{h:18:char-vs-string-for-passwords}
\code{String} is \textbf{immutable} and Java’s string pool (or just normal GC timing) can keep a copy of a password in memory for an unpredictable, possibly long time — there’s no way to proactively erase it. A \code{char[]} can be overwritten with zeros the moment you’re done with it, shrinking the window where the password sits in memory in plaintext. This is why \code{JPassword\allowbreak{}Field.\allowbreak{}get\allowbreak{}Password()} and JDBC’s \code{Connection} APIs that accept passwords use \code{char[]}.

\begin{codebox}
\begin{Verbatim}[commandchars=\\\{\}]
\PY{c+c1}{// Anti\PYZhy{}pattern: password as a String — cannot be wiped from memory on demand}
\PY{n}{String}\PY{+w}{ }\PY{n}{password}\PY{+w}{ }\PY{o}{=}\PY{+w}{ }\PY{k}{new}\PY{+w}{ }\PY{n}{String}\PY{p}{(}\PY{n}{passwordBytes}\PY{p}{)}\PY{p}{;}
\PY{n}{authenticate}\PY{p}{(}\PY{n}{username}\PY{p}{,}\PY{+w}{ }\PY{n}{password}\PY{p}{)}\PY{p}{;}
\PY{c+c1}{// password may live in memory long after this point, out of your control}

\PY{c+c1}{// Fix: char[] can be explicitly cleared right after use}
\PY{k+kt}{char}\PY{o}{[}\PY{o}{]}\PY{+w}{ }\PY{n}{password}\PY{+w}{ }\PY{o}{=}\PY{+w}{ }\PY{n}{readPasswordFromInput}\PY{p}{(}\PY{p}{)}\PY{p}{;}
\PY{k}{try}\PY{+w}{ }\PY{p}{\PYZob{}}
\PY{+w}{    }\PY{n}{authenticate}\PY{p}{(}\PY{n}{username}\PY{p}{,}\PY{+w}{ }\PY{n}{password}\PY{p}{)}\PY{p}{;}
\PY{p}{\PYZcb{}}\PY{+w}{ }\PY{k}{finally}\PY{+w}{ }\PY{p}{\PYZob{}}
\PY{+w}{    }\PY{n}{java}\PY{p}{.}\PY{n+na}{util}\PY{p}{.}\PY{n+na}{Arrays}\PY{p}{.}\PY{n+na}{fill}\PY{p}{(}\PY{n}{password}\PY{p}{,}\PY{+w}{ }\PY{l+s+sc}{\PYZsq{}\PYZbs{}0\PYZsq{}}\PY{p}{)}\PY{p}{;}\PY{+w}{ }\PY{c+c1}{// overwrite in memory}
\PY{p}{\PYZcb{}}
\end{Verbatim}
\end{codebox}

\subsection[SecureRandom vs Random]{SecureRandom vs Random}
\label{h:18:securerandom-vs-random}
\code{java.\allowbreak{}util.\allowbreak{}Random} is a fast \textbf{pseudo-random number generator (PRNG)} seeded predictably — given the seed (or even a handful of outputs), an attacker can often predict every future value. It must never be used for anything security-sensitive: tokens, session IDs, password reset codes, cryptographic keys. \code{SecureRandom} uses a \textbf{cryptographically secure PRNG (CSPRNG)}, designed so past output reveals nothing about future output.

\begin{codebox}
\begin{Verbatim}[commandchars=\\\{\}]
\PY{c+c1}{// Anti\PYZhy{}pattern: predictable token generation}
\PY{n}{Random}\PY{+w}{ }\PY{n}{random}\PY{+w}{ }\PY{o}{=}\PY{+w}{ }\PY{k}{new}\PY{+w}{ }\PY{n}{Random}\PY{p}{(}\PY{p}{)}\PY{p}{;}
\PY{n}{String}\PY{+w}{ }\PY{n}{resetToken}\PY{+w}{ }\PY{o}{=}\PY{+w}{ }\PY{n}{Long}\PY{p}{.}\PY{n+na}{toHexString}\PY{p}{(}\PY{n}{random}\PY{p}{.}\PY{n+na}{nextLong}\PY{p}{(}\PY{p}{)}\PY{p}{)}\PY{p}{;}\PY{+w}{ }\PY{c+c1}{// guessable by an attacker}

\PY{c+c1}{// Fix: SecureRandom for anything security\PYZhy{}sensitive}
\PY{n}{SecureRandom}\PY{+w}{ }\PY{n}{secureRandom}\PY{+w}{ }\PY{o}{=}\PY{+w}{ }\PY{k}{new}\PY{+w}{ }\PY{n}{SecureRandom}\PY{p}{(}\PY{p}{)}\PY{p}{;}
\PY{k+kt}{byte}\PY{o}{[}\PY{o}{]}\PY{+w}{ }\PY{n}{tokenBytes}\PY{+w}{ }\PY{o}{=}\PY{+w}{ }\PY{k}{new}\PY{+w}{ }\PY{k+kt}{byte}\PY{o}{[}\PY{l+m+mi}{32}\PY{o}{]}\PY{p}{;}
\PY{n}{secureRandom}\PY{p}{.}\PY{n+na}{nextBytes}\PY{p}{(}\PY{n}{tokenBytes}\PY{p}{)}\PY{p}{;}
\PY{n}{String}\PY{+w}{ }\PY{n}{resetToken}\PY{+w}{ }\PY{o}{=}\PY{+w}{ }\PY{n}{java}\PY{p}{.}\PY{n+na}{util}\PY{p}{.}\PY{n+na}{Base64}\PY{p}{.}\PY{n+na}{getUrlEncoder}\PY{p}{(}\PY{p}{)}\PY{p}{.}\PY{n+na}{withoutPadding}\PY{p}{(}\PY{p}{)}\PY{p}{.}\PY{n+na}{encodeToString}\PY{p}{(}\PY{n}{tokenBytes}\PY{p}{)}\PY{p}{;}
\end{Verbatim}
\end{codebox}

\subsection[Constant-time comparison]{Constant-time comparison}
\label{h:18:constant-time-comparison}
Comparing two secrets (e.g. an HMAC signature, an API key) with \code{String.\allowbreak{}equals} or \code{Arrays.\allowbreak{}equals} stops at the \textbf{first mismatched byte} — that tiny timing difference, measured over enough repeated attempts, can leak the correct value one byte at a time (a \textbf{timing attack}).

\begin{codebox}
\begin{Verbatim}[commandchars=\\\{\}]
\PY{c+c1}{// Anti\PYZhy{}pattern: early\PYZhy{}exit comparison leaks timing information}
\PY{k}{if}\PY{+w}{ }\PY{p}{(}\PY{o}{!}\PY{n}{providedSignature}\PY{p}{.}\PY{n+na}{equals}\PY{p}{(}\PY{n}{expectedSignature}\PY{p}{)}\PY{p}{)}\PY{+w}{ }\PY{p}{\PYZob{}}
\PY{+w}{    }\PY{k}{throw}\PY{+w}{ }\PY{k}{new}\PY{+w}{ }\PY{n}{SecurityException}\PY{p}{(}\PY{l+s}{\PYZdq{}}\PY{l+s}{Invalid signature}\PY{l+s}{\PYZdq{}}\PY{p}{)}\PY{p}{;}
\PY{p}{\PYZcb{}}

\PY{c+c1}{// Fix: constant\PYZhy{}time comparison — always examines every byte}
\PY{k+kt}{byte}\PY{o}{[}\PY{o}{]}\PY{+w}{ }\PY{n}{provided}\PY{+w}{ }\PY{o}{=}\PY{+w}{ }\PY{n}{providedSignature}\PY{p}{.}\PY{n+na}{getBytes}\PY{p}{(}\PY{n}{StandardCharsets}\PY{p}{.}\PY{n+na}{UTF\PYZus{}8}\PY{p}{)}\PY{p}{;}
\PY{k+kt}{byte}\PY{o}{[}\PY{o}{]}\PY{+w}{ }\PY{n}{expected}\PY{+w}{ }\PY{o}{=}\PY{+w}{ }\PY{n}{expectedSignature}\PY{p}{.}\PY{n+na}{getBytes}\PY{p}{(}\PY{n}{StandardCharsets}\PY{p}{.}\PY{n+na}{UTF\PYZus{}8}\PY{p}{)}\PY{p}{;}
\PY{k}{if}\PY{+w}{ }\PY{p}{(}\PY{o}{!}\PY{n}{MessageDigest}\PY{p}{.}\PY{n+na}{isEqual}\PY{p}{(}\PY{n}{provided}\PY{p}{,}\PY{+w}{ }\PY{n}{expected}\PY{p}{)}\PY{p}{)}\PY{+w}{ }\PY{p}{\PYZob{}}
\PY{+w}{    }\PY{k}{throw}\PY{+w}{ }\PY{k}{new}\PY{+w}{ }\PY{n}{SecurityException}\PY{p}{(}\PY{l+s}{\PYZdq{}}\PY{l+s}{Invalid signature}\PY{l+s}{\PYZdq{}}\PY{p}{)}\PY{p}{;}
\PY{p}{\PYZcb{}}
\end{Verbatim}
\end{codebox}

\subsection[TLS defaults and certificate validation]{TLS defaults and certificate validation}
\label{h:18:tls-defaults-and-certificate-validation}
\textbf{TLS} (Transport Layer Security, the successor to SSL) provides encryption and — critically — \textbf{authentication}: the client verifies the server’s identity through its \textbf{certificate}, issued and signed by a trusted \textbf{certificate authority (CA)}. Skipping that verification defeats the entire point of TLS.

\begin{codebox}
\begin{Verbatim}[commandchars=\\\{\}]
\PY{c+c1}{// Anti\PYZhy{}pattern: a trust\PYZhy{}all TrustManager disables certificate validation entirely}
\PY{n}{TrustManager}\PY{o}{[}\PY{o}{]}\PY{+w}{ }\PY{n}{trustAllCerts}\PY{+w}{ }\PY{o}{=}\PY{+w}{ }\PY{k}{new}\PY{+w}{ }\PY{n}{TrustManager}\PY{o}{[}\PY{o}{]}\PY{+w}{ }\PY{p}{\PYZob{}}
\PY{+w}{    }\PY{k}{new}\PY{+w}{ }\PY{n}{X509TrustManager}\PY{p}{(}\PY{p}{)}\PY{+w}{ }\PY{p}{\PYZob{}}
\PY{+w}{        }\PY{k+kd}{public}\PY{+w}{ }\PY{k+kt}{void}\PY{+w}{ }\PY{n+nf}{checkClientTrusted}\PY{p}{(}\PY{n}{X509Certificate}\PY{o}{[}\PY{o}{]}\PY{+w}{ }\PY{n}{chain}\PY{p}{,}\PY{+w}{ }\PY{n}{String}\PY{+w}{ }\PY{n}{authType}\PY{p}{)}\PY{+w}{ }\PY{p}{\PYZob{}}\PY{p}{\PYZcb{}}\PY{+w}{ }\PY{c+c1}{// accepts anything}
\PY{+w}{        }\PY{k+kd}{public}\PY{+w}{ }\PY{k+kt}{void}\PY{+w}{ }\PY{n+nf}{checkServerTrusted}\PY{p}{(}\PY{n}{X509Certificate}\PY{o}{[}\PY{o}{]}\PY{+w}{ }\PY{n}{chain}\PY{p}{,}\PY{+w}{ }\PY{n}{String}\PY{+w}{ }\PY{n}{authType}\PY{p}{)}\PY{+w}{ }\PY{p}{\PYZob{}}\PY{p}{\PYZcb{}}\PY{+w}{ }\PY{c+c1}{// accepts anything}
\PY{+w}{        }\PY{k+kd}{public}\PY{+w}{ }\PY{n}{X509Certificate}\PY{o}{[}\PY{o}{]}\PY{+w}{ }\PY{n+nf}{getAcceptedIssuers}\PY{p}{(}\PY{p}{)}\PY{+w}{ }\PY{p}{\PYZob{}}\PY{+w}{ }\PY{k}{return}\PY{+w}{ }\PY{k}{new}\PY{+w}{ }\PY{n}{X509Certificate}\PY{o}{[}\PY{l+m+mi}{0}\PY{o}{]}\PY{p}{;}\PY{+w}{ }\PY{p}{\PYZcb{}}
\PY{+w}{    }\PY{p}{\PYZcb{}}
\PY{p}{\PYZcb{}}\PY{p}{;}
\PY{n}{SSLContext}\PY{+w}{ }\PY{n}{sslContext}\PY{+w}{ }\PY{o}{=}\PY{+w}{ }\PY{n}{SSLContext}\PY{p}{.}\PY{n+na}{getInstance}\PY{p}{(}\PY{l+s}{\PYZdq{}}\PY{l+s}{TLS}\PY{l+s}{\PYZdq{}}\PY{p}{)}\PY{p}{;}
\PY{n}{sslContext}\PY{p}{.}\PY{n+na}{init}\PY{p}{(}\PY{k+kc}{null}\PY{p}{,}\PY{+w}{ }\PY{n}{trustAllCerts}\PY{p}{,}\PY{+w}{ }\PY{k}{new}\PY{+w}{ }\PY{n}{java}\PY{p}{.}\PY{n+na}{security}\PY{p}{.}\PY{n+na}{SecureRandom}\PY{p}{(}\PY{p}{)}\PY{p}{)}\PY{p}{;}
\PY{c+c1}{// Any server, including an attacker performing a man\PYZhy{}in\PYZhy{}the\PYZhy{}middle attack, is now \PYZdq{}trusted\PYZdq{}}
\end{Verbatim}
\end{codebox}

\begin{codebox}
\begin{Verbatim}[commandchars=\\\{\}]
\PY{c+c1}{// Fix: use the platform default trust managers (or a specific, real trust store)}
\PY{n}{SSLContext}\PY{+w}{ }\PY{n}{sslContext}\PY{+w}{ }\PY{o}{=}\PY{+w}{ }\PY{n}{SSLContext}\PY{p}{.}\PY{n+na}{getInstance}\PY{p}{(}\PY{l+s}{\PYZdq{}}\PY{l+s}{TLSv1.3}\PY{l+s}{\PYZdq{}}\PY{p}{)}\PY{p}{;}
\PY{n}{sslContext}\PY{p}{.}\PY{n+na}{init}\PY{p}{(}\PY{k+kc}{null}\PY{p}{,}\PY{+w}{ }\PY{k+kc}{null}\PY{p}{,}\PY{+w}{ }\PY{k+kc}{null}\PY{p}{)}\PY{p}{;}\PY{+w}{ }\PY{c+c1}{// null TrustManager[] =\PYZgt{} use the JVM\PYZsq{}s default, real CA trust store}
\end{Verbatim}
\end{codebox}

If you need to trust a private/internal CA, build a real \code{TrustManagerFactory} from a keystore containing that CA’s certificate (as shown in the HTTP Client section’s \code{SSLContext} example) — never disable validation to work around a certificate problem, even “temporarily.”

\subsection[The Security Manager (JEP 411)]{The Security Manager (JEP 411)}
\label{h:18:the-security-manager-jep-411}
The \textbf{Security Manager} was Java’s mechanism for sandboxing code — restricting what a piece of code could do (file access, network access, etc.) via fine-grained permissions. As of \textbf{JEP 411} (Java 17), it is deprecated for removal, and it has been fully removed in Java 24. If you inherit code that relies on \code{SecurityManager}/\code{AccessController}, plan to migrate away from it — modern isolation is typically done at the process, container, or OS level instead (containers, seccomp profiles, dedicated low-privilege OS users), not inside the JVM.

\subsection[Dependency vulnerabilities]{Dependency vulnerabilities}
\label{h:18:dependency-vulnerabilities}
Most real-world Java vulnerabilities (Log4Shell being the most infamous recent example) come from \textbf{third-party libraries}, not code you wrote. Practices that matter in review:

\begin{itemize}
\item 
Keep an inventory of dependencies and their versions (a \textbf{Software Bill of Materials}, or SBOM).

\item 
Run a dependency vulnerability scanner (e.g. OWASP Dependency-Check, \code{mvn~\allowbreak{}versions:\allowbreak{}display-\allowbreak{}dependency-\allowbreak{}updates}, GitHub Dependabot) as part of CI.

\item 
Pin versions and update regularly — being a few major versions behind makes upgrading (and patching a critical CVE) much riskier and slower.

\item 
Don’t pull in a whole library for one function if you can avoid the dependency weight and attack surface.

\end{itemize}

\section[Cryptography APIs]{Cryptography APIs}
\label{h:18:cryptography-apis}
Java’s cryptography support is split into the \textbf{JCA} (Java Cryptography Architecture — the framework: engine classes like \code{MessageDigest}, \code{Cipher}, \code{Signature}) and the \textbf{JCE} (Java Cryptography Extension — historically export-restricted extensions, now merged into the JDK). Both are \textbf{provider-based}: you ask for an algorithm by name (e.g. \code{\textquotedbl{}AES/\allowbreak{}GCM/\allowbreak{}NoPadding\textquotedbl{}}), and a registered \textbf{Provider} (the built-in SunJCE, or a third-party one like Bouncy Castle) supplies the actual implementation. This lets you swap implementations without changing application code.

\begin{codebox}
\begin{Verbatim}[commandchars=\\\{\}]
\PY{k+kn}{import}\PY{+w}{ }\PY{n+nn}{java.security.Provider}\PY{p}{;}
\PY{k+kn}{import}\PY{+w}{ }\PY{n+nn}{java.security.Security}\PY{p}{;}

\PY{k}{for}\PY{+w}{ }\PY{p}{(}\PY{n}{Provider}\PY{+w}{ }\PY{n}{provider}\PY{+w}{ }\PY{p}{:}\PY{+w}{ }\PY{n}{Security}\PY{p}{.}\PY{n+na}{getProviders}\PY{p}{(}\PY{p}{)}\PY{p}{)}\PY{+w}{ }\PY{p}{\PYZob{}}
\PY{+w}{    }\PY{n}{System}\PY{p}{.}\PY{n+na}{out}\PY{p}{.}\PY{n+na}{println}\PY{p}{(}\PY{n}{provider}\PY{p}{.}\PY{n+na}{getName}\PY{p}{(}\PY{p}{)}\PY{p}{)}\PY{p}{;}
\PY{p}{\PYZcb{}}
\end{Verbatim}
\end{codebox}

\subsection[MessageDigest — hashing]{MessageDigest — hashing}
\label{h:18:messagedigest--hashing}
A \textbf{hash function} takes input of any size and produces a fixed-size fingerprint (\textbf{digest}). It’s one-way: you can’t recover the input from the digest. Use it for integrity checks (e.g. “did this file change?”), never for passwords (see below).

\begin{codebox}
\begin{Verbatim}[commandchars=\\\{\}]
\PY{k+kn}{import}\PY{+w}{ }\PY{n+nn}{java.security.MessageDigest}\PY{p}{;}
\PY{k+kn}{import}\PY{+w}{ }\PY{n+nn}{java.security.NoSuchAlgorithmException}\PY{p}{;}
\PY{k+kn}{import}\PY{+w}{ }\PY{n+nn}{java.nio.charset.StandardCharsets}\PY{p}{;}
\PY{k+kn}{import}\PY{+w}{ }\PY{n+nn}{java.util.HexFormat}\PY{p}{;}

\PY{k+kd}{public}\PY{+w}{ }\PY{k+kd}{class} \PY{n+nc}{HashExample}\PY{+w}{ }\PY{p}{\PYZob{}}
\PY{+w}{    }\PY{k+kd}{public}\PY{+w}{ }\PY{k+kd}{static}\PY{+w}{ }\PY{n}{String}\PY{+w}{ }\PY{n+nf}{sha256}\PY{p}{(}\PY{n}{String}\PY{+w}{ }\PY{n}{input}\PY{p}{)}\PY{+w}{ }\PY{k+kd}{throws}\PY{+w}{ }\PY{n}{NoSuchAlgorithmException}\PY{+w}{ }\PY{p}{\PYZob{}}
\PY{+w}{        }\PY{n}{MessageDigest}\PY{+w}{ }\PY{n}{digest}\PY{+w}{ }\PY{o}{=}\PY{+w}{ }\PY{n}{MessageDigest}\PY{p}{.}\PY{n+na}{getInstance}\PY{p}{(}\PY{l+s}{\PYZdq{}}\PY{l+s}{SHA\PYZhy{}256}\PY{l+s}{\PYZdq{}}\PY{p}{)}\PY{p}{;}
\PY{+w}{        }\PY{k+kt}{byte}\PY{o}{[}\PY{o}{]}\PY{+w}{ }\PY{n}{hash}\PY{+w}{ }\PY{o}{=}\PY{+w}{ }\PY{n}{digest}\PY{p}{.}\PY{n+na}{digest}\PY{p}{(}\PY{n}{input}\PY{p}{.}\PY{n+na}{getBytes}\PY{p}{(}\PY{n}{StandardCharsets}\PY{p}{.}\PY{n+na}{UTF\PYZus{}8}\PY{p}{)}\PY{p}{)}\PY{p}{;}
\PY{+w}{        }\PY{k}{return}\PY{+w}{ }\PY{n}{HexFormat}\PY{p}{.}\PY{n+na}{of}\PY{p}{(}\PY{p}{)}\PY{p}{.}\PY{n+na}{formatHex}\PY{p}{(}\PY{n}{hash}\PY{p}{)}\PY{p}{;}
\PY{+w}{    }\PY{p}{\PYZcb{}}
\PY{p}{\PYZcb{}}
\end{Verbatim}
\end{codebox}

\subsection[Password hashing — PBKDF2, and why plain hashing isn’t enough]{Password hashing — PBKDF2, and why plain hashing isn’t enough}
\label{h:18:password-hashing--pbkdf2-and-why-plain-hashing-isnt-enough}
\textbf{Never store passwords in plaintext, and never hash them with a plain, fast hash like MD5 or SHA-1/SHA-256 alone.} Fast hashes let an attacker who steals your password database try billions of guesses per second on cheap GPU hardware. Password hashing needs to be \textbf{slow and tunable} on purpose, and use a per-password \textbf{salt} (random data mixed in) so identical passwords don’t produce identical hashes.

\begin{codebox}
\begin{Verbatim}[commandchars=\\\{\}]
\PY{c+c1}{// Anti\PYZhy{}pattern: a fast, unsalted hash for passwords}
\PY{n}{MessageDigest}\PY{+w}{ }\PY{n}{digest}\PY{+w}{ }\PY{o}{=}\PY{+w}{ }\PY{n}{MessageDigest}\PY{p}{.}\PY{n+na}{getInstance}\PY{p}{(}\PY{l+s}{\PYZdq{}}\PY{l+s}{SHA\PYZhy{}256}\PY{l+s}{\PYZdq{}}\PY{p}{)}\PY{p}{;}
\PY{k+kt}{byte}\PY{o}{[}\PY{o}{]}\PY{+w}{ }\PY{n}{hash}\PY{+w}{ }\PY{o}{=}\PY{+w}{ }\PY{n}{digest}\PY{p}{.}\PY{n+na}{digest}\PY{p}{(}\PY{n}{password}\PY{p}{.}\PY{n+na}{getBytes}\PY{p}{(}\PY{n}{StandardCharsets}\PY{p}{.}\PY{n+na}{UTF\PYZus{}8}\PY{p}{)}\PY{p}{)}\PY{p}{;}
\PY{c+c1}{// Crackable at billions of attempts/second with commodity GPUs; identical passwords \PYZhy{}\PYZgt{} identical hashes}
\end{Verbatim}
\end{codebox}

\begin{codebox}
\begin{Verbatim}[commandchars=\\\{\}]
\PY{c+c1}{// Fix: PBKDF2 via SecretKeyFactory — salted, and deliberately slow (many iterations)}
\PY{k+kn}{import}\PY{+w}{ }\PY{n+nn}{javax.crypto.SecretKeyFactory}\PY{p}{;}
\PY{k+kn}{import}\PY{+w}{ }\PY{n+nn}{javax.crypto.spec.PBEKeySpec}\PY{p}{;}
\PY{k+kn}{import}\PY{+w}{ }\PY{n+nn}{java.security.SecureRandom}\PY{p}{;}

\PY{k+kd}{public}\PY{+w}{ }\PY{k+kd}{class} \PY{n+nc}{PasswordHasher}\PY{+w}{ }\PY{p}{\PYZob{}}
\PY{+w}{    }\PY{k+kd}{public}\PY{+w}{ }\PY{k+kd}{static}\PY{+w}{ }\PY{k+kt}{byte}\PY{o}{[}\PY{o}{]}\PY{+w}{ }\PY{n+nf}{hash}\PY{p}{(}\PY{k+kt}{char}\PY{o}{[}\PY{o}{]}\PY{+w}{ }\PY{n}{password}\PY{p}{,}\PY{+w}{ }\PY{k+kt}{byte}\PY{o}{[}\PY{o}{]}\PY{+w}{ }\PY{n}{salt}\PY{p}{)}\PY{+w}{ }\PY{k+kd}{throws}\PY{+w}{ }\PY{n}{Exception}\PY{+w}{ }\PY{p}{\PYZob{}}
\PY{+w}{        }\PY{n}{PBEKeySpec}\PY{+w}{ }\PY{n}{spec}\PY{+w}{ }\PY{o}{=}\PY{+w}{ }\PY{k}{new}\PY{+w}{ }\PY{n}{PBEKeySpec}\PY{p}{(}\PY{n}{password}\PY{p}{,}\PY{+w}{ }\PY{n}{salt}\PY{p}{,}\PY{+w}{ }\PY{l+m+mi}{210\PYZus{}000}\PY{p}{,}\PY{+w}{ }\PY{l+m+mi}{256}\PY{p}{)}\PY{p}{;}\PY{+w}{ }\PY{c+c1}{// iterations, key length in bits}
\PY{+w}{        }\PY{n}{SecretKeyFactory}\PY{+w}{ }\PY{n}{factory}\PY{+w}{ }\PY{o}{=}\PY{+w}{ }\PY{n}{SecretKeyFactory}\PY{p}{.}\PY{n+na}{getInstance}\PY{p}{(}\PY{l+s}{\PYZdq{}}\PY{l+s}{PBKDF2WithHmacSHA256}\PY{l+s}{\PYZdq{}}\PY{p}{)}\PY{p}{;}
\PY{+w}{        }\PY{k}{return}\PY{+w}{ }\PY{n}{factory}\PY{p}{.}\PY{n+na}{generateSecret}\PY{p}{(}\PY{n}{spec}\PY{p}{)}\PY{p}{.}\PY{n+na}{getEncoded}\PY{p}{(}\PY{p}{)}\PY{p}{;}
\PY{+w}{    }\PY{p}{\PYZcb{}}

\PY{+w}{    }\PY{k+kd}{public}\PY{+w}{ }\PY{k+kd}{static}\PY{+w}{ }\PY{k+kt}{byte}\PY{o}{[}\PY{o}{]}\PY{+w}{ }\PY{n+nf}{newSalt}\PY{p}{(}\PY{p}{)}\PY{+w}{ }\PY{p}{\PYZob{}}
\PY{+w}{        }\PY{k+kt}{byte}\PY{o}{[}\PY{o}{]}\PY{+w}{ }\PY{n}{salt}\PY{+w}{ }\PY{o}{=}\PY{+w}{ }\PY{k}{new}\PY{+w}{ }\PY{k+kt}{byte}\PY{o}{[}\PY{l+m+mi}{16}\PY{o}{]}\PY{p}{;}
\PY{+w}{        }\PY{k}{new}\PY{+w}{ }\PY{n}{SecureRandom}\PY{p}{(}\PY{p}{)}\PY{p}{.}\PY{n+na}{nextBytes}\PY{p}{(}\PY{n}{salt}\PY{p}{)}\PY{p}{;}
\PY{+w}{        }\PY{k}{return}\PY{+w}{ }\PY{n}{salt}\PY{p}{;}
\PY{+w}{    }\PY{p}{\PYZcb{}}
\PY{p}{\PYZcb{}}
\end{Verbatim}
\end{codebox}

Store the salt alongside the hash (it’s not secret) and re-derive using the same salt and iteration count to verify a login attempt. In practice, \textbf{bcrypt} or \textbf{argon2} (via a library, since the JDK doesn’t ship them natively) are usually preferred over PBKDF2 for new systems — they’re purpose-built for password hashing with better resistance to GPU/ASIC cracking. Whatever you choose: never MD5, never plain SHA-1/SHA-256, never unsalted.

\subsection[Cipher — AES-GCM done correctly]{Cipher — AES-GCM done correctly}
\label{h:18:cipher--aes-gcm-done-correctly}
\code{Cipher} performs \textbf{encryption} (reversible scrambling with a key) and decryption. \textbf{AES-GCM} (Advanced Encryption Standard, Galois/Counter Mode) is the recommended default: it’s fast, and it’s an \textbf{authenticated} mode — it detects tampering, not just confidentiality.

The single most common mistake: \textbf{ECB mode}. ECB encrypts each fixed-size block independently, so identical plaintext blocks always produce identical ciphertext blocks — patterns in the input leak straight through, famously visible if you ECB-encrypt an image.

\begin{codebox}
\begin{Verbatim}[commandchars=\\\{\}]
\PY{c+c1}{// Anti\PYZhy{}pattern: AES in ECB mode — identical plaintext blocks \PYZhy{}\PYZgt{} identical ciphertext blocks}
\PY{n}{Cipher}\PY{+w}{ }\PY{n}{cipher}\PY{+w}{ }\PY{o}{=}\PY{+w}{ }\PY{n}{Cipher}\PY{p}{.}\PY{n+na}{getInstance}\PY{p}{(}\PY{l+s}{\PYZdq{}}\PY{l+s}{AES/ECB/PKCS5Padding}\PY{l+s}{\PYZdq{}}\PY{p}{)}\PY{p}{;}
\PY{n}{cipher}\PY{p}{.}\PY{n+na}{init}\PY{p}{(}\PY{n}{Cipher}\PY{p}{.}\PY{n+na}{ENCRYPT\PYZus{}MODE}\PY{p}{,}\PY{+w}{ }\PY{n}{secretKey}\PY{p}{)}\PY{p}{;}
\PY{k+kt}{byte}\PY{o}{[}\PY{o}{]}\PY{+w}{ }\PY{n}{ciphertext}\PY{+w}{ }\PY{o}{=}\PY{+w}{ }\PY{n}{cipher}\PY{p}{.}\PY{n+na}{doFinal}\PY{p}{(}\PY{n}{plaintext}\PY{p}{)}\PY{p}{;}\PY{+w}{ }\PY{c+c1}{// patterns in plaintext are visible in ciphertext}
\end{Verbatim}
\end{codebox}

\begin{codebox}
\begin{Verbatim}[commandchars=\\\{\}]
\PY{c+c1}{// Fix: AES\PYZhy{}GCM with a fresh random IV (nonce) for every single encryption}
\PY{k+kn}{import}\PY{+w}{ }\PY{n+nn}{javax.crypto.Cipher}\PY{p}{;}
\PY{k+kn}{import}\PY{+w}{ }\PY{n+nn}{javax.crypto.SecretKey}\PY{p}{;}
\PY{k+kn}{import}\PY{+w}{ }\PY{n+nn}{javax.crypto.spec.GCMParameterSpec}\PY{p}{;}
\PY{k+kn}{import}\PY{+w}{ }\PY{n+nn}{java.security.SecureRandom}\PY{p}{;}

\PY{k+kd}{public}\PY{+w}{ }\PY{k+kd}{class} \PY{n+nc}{AesGcmExample}\PY{+w}{ }\PY{p}{\PYZob{}}
\PY{+w}{    }\PY{k+kd}{private}\PY{+w}{ }\PY{k+kd}{static}\PY{+w}{ }\PY{k+kd}{final}\PY{+w}{ }\PY{k+kt}{int}\PY{+w}{ }\PY{n}{IV\PYZus{}LENGTH\PYZus{}BYTES}\PY{+w}{ }\PY{o}{=}\PY{+w}{ }\PY{l+m+mi}{12}\PY{p}{;}
\PY{+w}{    }\PY{k+kd}{private}\PY{+w}{ }\PY{k+kd}{static}\PY{+w}{ }\PY{k+kd}{final}\PY{+w}{ }\PY{k+kt}{int}\PY{+w}{ }\PY{n}{TAG\PYZus{}LENGTH\PYZus{}BITS}\PY{+w}{ }\PY{o}{=}\PY{+w}{ }\PY{l+m+mi}{128}\PY{p}{;}

\PY{+w}{    }\PY{k+kd}{public}\PY{+w}{ }\PY{k+kd}{static}\PY{+w}{ }\PY{k+kt}{byte}\PY{o}{[}\PY{o}{]}\PY{+w}{ }\PY{n+nf}{encrypt}\PY{p}{(}\PY{k+kt}{byte}\PY{o}{[}\PY{o}{]}\PY{+w}{ }\PY{n}{plaintext}\PY{p}{,}\PY{+w}{ }\PY{n}{SecretKey}\PY{+w}{ }\PY{n}{key}\PY{p}{)}\PY{+w}{ }\PY{k+kd}{throws}\PY{+w}{ }\PY{n}{Exception}\PY{+w}{ }\PY{p}{\PYZob{}}
\PY{+w}{        }\PY{k+kt}{byte}\PY{o}{[}\PY{o}{]}\PY{+w}{ }\PY{n}{iv}\PY{+w}{ }\PY{o}{=}\PY{+w}{ }\PY{k}{new}\PY{+w}{ }\PY{k+kt}{byte}\PY{o}{[}\PY{n}{IV\PYZus{}LENGTH\PYZus{}BYTES}\PY{o}{]}\PY{p}{;}
\PY{+w}{        }\PY{k}{new}\PY{+w}{ }\PY{n}{SecureRandom}\PY{p}{(}\PY{p}{)}\PY{p}{.}\PY{n+na}{nextBytes}\PY{p}{(}\PY{n}{iv}\PY{p}{)}\PY{p}{;}\PY{+w}{ }\PY{c+c1}{// MUST be unique per encryption with the same key}

\PY{+w}{        }\PY{n}{Cipher}\PY{+w}{ }\PY{n}{cipher}\PY{+w}{ }\PY{o}{=}\PY{+w}{ }\PY{n}{Cipher}\PY{p}{.}\PY{n+na}{getInstance}\PY{p}{(}\PY{l+s}{\PYZdq{}}\PY{l+s}{AES/GCM/NoPadding}\PY{l+s}{\PYZdq{}}\PY{p}{)}\PY{p}{;}
\PY{+w}{        }\PY{n}{cipher}\PY{p}{.}\PY{n+na}{init}\PY{p}{(}\PY{n}{Cipher}\PY{p}{.}\PY{n+na}{ENCRYPT\PYZus{}MODE}\PY{p}{,}\PY{+w}{ }\PY{n}{key}\PY{p}{,}\PY{+w}{ }\PY{k}{new}\PY{+w}{ }\PY{n}{GCMParameterSpec}\PY{p}{(}\PY{n}{TAG\PYZus{}LENGTH\PYZus{}BITS}\PY{p}{,}\PY{+w}{ }\PY{n}{iv}\PY{p}{)}\PY{p}{)}\PY{p}{;}
\PY{+w}{        }\PY{k+kt}{byte}\PY{o}{[}\PY{o}{]}\PY{+w}{ }\PY{n}{ciphertext}\PY{+w}{ }\PY{o}{=}\PY{+w}{ }\PY{n}{cipher}\PY{p}{.}\PY{n+na}{doFinal}\PY{p}{(}\PY{n}{plaintext}\PY{p}{)}\PY{p}{;}

\PY{+w}{        }\PY{c+c1}{// Prepend the IV so decrypt() can read it back — the IV itself is not secret}
\PY{+w}{        }\PY{k+kt}{byte}\PY{o}{[}\PY{o}{]}\PY{+w}{ }\PY{n}{result}\PY{+w}{ }\PY{o}{=}\PY{+w}{ }\PY{k}{new}\PY{+w}{ }\PY{k+kt}{byte}\PY{o}{[}\PY{n}{iv}\PY{p}{.}\PY{n+na}{length}\PY{+w}{ }\PY{o}{+}\PY{+w}{ }\PY{n}{ciphertext}\PY{p}{.}\PY{n+na}{length}\PY{o}{]}\PY{p}{;}
\PY{+w}{        }\PY{n}{System}\PY{p}{.}\PY{n+na}{arraycopy}\PY{p}{(}\PY{n}{iv}\PY{p}{,}\PY{+w}{ }\PY{l+m+mi}{0}\PY{p}{,}\PY{+w}{ }\PY{n}{result}\PY{p}{,}\PY{+w}{ }\PY{l+m+mi}{0}\PY{p}{,}\PY{+w}{ }\PY{n}{iv}\PY{p}{.}\PY{n+na}{length}\PY{p}{)}\PY{p}{;}
\PY{+w}{        }\PY{n}{System}\PY{p}{.}\PY{n+na}{arraycopy}\PY{p}{(}\PY{n}{ciphertext}\PY{p}{,}\PY{+w}{ }\PY{l+m+mi}{0}\PY{p}{,}\PY{+w}{ }\PY{n}{result}\PY{p}{,}\PY{+w}{ }\PY{n}{iv}\PY{p}{.}\PY{n+na}{length}\PY{p}{,}\PY{+w}{ }\PY{n}{ciphertext}\PY{p}{.}\PY{n+na}{length}\PY{p}{)}\PY{p}{;}
\PY{+w}{        }\PY{k}{return}\PY{+w}{ }\PY{n}{result}\PY{p}{;}
\PY{+w}{    }\PY{p}{\PYZcb{}}

\PY{+w}{    }\PY{k+kd}{public}\PY{+w}{ }\PY{k+kd}{static}\PY{+w}{ }\PY{k+kt}{byte}\PY{o}{[}\PY{o}{]}\PY{+w}{ }\PY{n+nf}{decrypt}\PY{p}{(}\PY{k+kt}{byte}\PY{o}{[}\PY{o}{]}\PY{+w}{ }\PY{n}{ivAndCiphertext}\PY{p}{,}\PY{+w}{ }\PY{n}{SecretKey}\PY{+w}{ }\PY{n}{key}\PY{p}{)}\PY{+w}{ }\PY{k+kd}{throws}\PY{+w}{ }\PY{n}{Exception}\PY{+w}{ }\PY{p}{\PYZob{}}
\PY{+w}{        }\PY{k+kt}{byte}\PY{o}{[}\PY{o}{]}\PY{+w}{ }\PY{n}{iv}\PY{+w}{ }\PY{o}{=}\PY{+w}{ }\PY{n}{java}\PY{p}{.}\PY{n+na}{util}\PY{p}{.}\PY{n+na}{Arrays}\PY{p}{.}\PY{n+na}{copyOfRange}\PY{p}{(}\PY{n}{ivAndCiphertext}\PY{p}{,}\PY{+w}{ }\PY{l+m+mi}{0}\PY{p}{,}\PY{+w}{ }\PY{n}{IV\PYZus{}LENGTH\PYZus{}BYTES}\PY{p}{)}\PY{p}{;}
\PY{+w}{        }\PY{k+kt}{byte}\PY{o}{[}\PY{o}{]}\PY{+w}{ }\PY{n}{ciphertext}\PY{+w}{ }\PY{o}{=}\PY{+w}{ }\PY{n}{java}\PY{p}{.}\PY{n+na}{util}\PY{p}{.}\PY{n+na}{Arrays}\PY{p}{.}\PY{n+na}{copyOfRange}\PY{p}{(}\PY{n}{ivAndCiphertext}\PY{p}{,}\PY{+w}{ }\PY{n}{IV\PYZus{}LENGTH\PYZus{}BYTES}\PY{p}{,}\PY{+w}{ }\PY{n}{ivAndCiphertext}\PY{p}{.}\PY{n+na}{length}\PY{p}{)}\PY{p}{;}

\PY{+w}{        }\PY{n}{Cipher}\PY{+w}{ }\PY{n}{cipher}\PY{+w}{ }\PY{o}{=}\PY{+w}{ }\PY{n}{Cipher}\PY{p}{.}\PY{n+na}{getInstance}\PY{p}{(}\PY{l+s}{\PYZdq{}}\PY{l+s}{AES/GCM/NoPadding}\PY{l+s}{\PYZdq{}}\PY{p}{)}\PY{p}{;}
\PY{+w}{        }\PY{n}{cipher}\PY{p}{.}\PY{n+na}{init}\PY{p}{(}\PY{n}{Cipher}\PY{p}{.}\PY{n+na}{DECRYPT\PYZus{}MODE}\PY{p}{,}\PY{+w}{ }\PY{n}{key}\PY{p}{,}\PY{+w}{ }\PY{k}{new}\PY{+w}{ }\PY{n}{GCMParameterSpec}\PY{p}{(}\PY{n}{TAG\PYZus{}LENGTH\PYZus{}BITS}\PY{p}{,}\PY{+w}{ }\PY{n}{iv}\PY{p}{)}\PY{p}{)}\PY{p}{;}
\PY{+w}{        }\PY{k}{return}\PY{+w}{ }\PY{n}{cipher}\PY{p}{.}\PY{n+na}{doFinal}\PY{p}{(}\PY{n}{ciphertext}\PY{p}{)}\PY{p}{;}\PY{+w}{ }\PY{c+c1}{// throws AEADBadTagException if tampered}
\PY{+w}{    }\PY{p}{\PYZcb{}}
\PY{p}{\PYZcb{}}
\end{Verbatim}
\end{codebox}

The \textbf{IV} (Initialization Vector, also called a nonce) must be unique for every encryption performed with the same key — reusing an IV with GCM catastrophically breaks its security guarantees, potentially revealing the plaintext and letting an attacker forge messages. Generating it fresh with \code{SecureRandom} every time, as above, is the standard safe approach.

\subsection[KeyGenerator / KeyPairGenerator]{KeyGenerator / KeyPairGenerator}
\label{h:18:keygenerator--keypairgenerator}
\code{KeyGenerator} creates \textbf{symmetric} keys (the same key encrypts and decrypts, used with \code{Cipher} above). \code{KeyPairGenerator} creates \textbf{asymmetric} key pairs (a public key anyone can use to encrypt or verify, and a private key only the owner holds, to decrypt or sign).

\begin{codebox}
\begin{Verbatim}[commandchars=\\\{\}]
\PY{k+kn}{import}\PY{+w}{ }\PY{n+nn}{javax.crypto.KeyGenerator}\PY{p}{;}
\PY{k+kn}{import}\PY{+w}{ }\PY{n+nn}{javax.crypto.SecretKey}\PY{p}{;}
\PY{k+kn}{import}\PY{+w}{ }\PY{n+nn}{java.security.KeyPair}\PY{p}{;}
\PY{k+kn}{import}\PY{+w}{ }\PY{n+nn}{java.security.KeyPairGenerator}\PY{p}{;}

\PY{c+c1}{// Symmetric AES key}
\PY{n}{KeyGenerator}\PY{+w}{ }\PY{n}{keyGen}\PY{+w}{ }\PY{o}{=}\PY{+w}{ }\PY{n}{KeyGenerator}\PY{p}{.}\PY{n+na}{getInstance}\PY{p}{(}\PY{l+s}{\PYZdq{}}\PY{l+s}{AES}\PY{l+s}{\PYZdq{}}\PY{p}{)}\PY{p}{;}
\PY{n}{keyGen}\PY{p}{.}\PY{n+na}{init}\PY{p}{(}\PY{l+m+mi}{256}\PY{p}{)}\PY{p}{;}\PY{+w}{ }\PY{c+c1}{// key size in bits}
\PY{n}{SecretKey}\PY{+w}{ }\PY{n}{aesKey}\PY{+w}{ }\PY{o}{=}\PY{+w}{ }\PY{n}{keyGen}\PY{p}{.}\PY{n+na}{generateKey}\PY{p}{(}\PY{p}{)}\PY{p}{;}

\PY{c+c1}{// Asymmetric key pair}
\PY{n}{KeyPairGenerator}\PY{+w}{ }\PY{n}{keyPairGen}\PY{+w}{ }\PY{o}{=}\PY{+w}{ }\PY{n}{KeyPairGenerator}\PY{p}{.}\PY{n+na}{getInstance}\PY{p}{(}\PY{l+s}{\PYZdq{}}\PY{l+s}{RSA}\PY{l+s}{\PYZdq{}}\PY{p}{)}\PY{p}{;}
\PY{n}{keyPairGen}\PY{p}{.}\PY{n+na}{initialize}\PY{p}{(}\PY{l+m+mi}{3072}\PY{p}{)}\PY{p}{;}\PY{+w}{ }\PY{c+c1}{// modulus size in bits}
\PY{n}{KeyPair}\PY{+w}{ }\PY{n}{keyPair}\PY{+w}{ }\PY{o}{=}\PY{+w}{ }\PY{n}{keyPairGen}\PY{p}{.}\PY{n+na}{generateKeyPair}\PY{p}{(}\PY{p}{)}\PY{p}{;}
\end{Verbatim}
\end{codebox}

\subsection[Signature — proving authenticity and integrity]{Signature — proving authenticity and integrity}
\label{h:18:signature--proving-authenticity-and-integrity}
A \textbf{digital signature} lets the holder of a private key sign data so that anyone with the corresponding public key can verify it came from them and wasn’t altered — the asymmetric-key analogue of the HMAC/\code{MessageDigest.\allowbreak{}isEqual} comparison shown earlier.

\begin{codebox}
\begin{Verbatim}[commandchars=\\\{\}]
\PY{k+kn}{import}\PY{+w}{ }\PY{n+nn}{java.security.Signature}\PY{p}{;}
\PY{k+kn}{import}\PY{+w}{ }\PY{n+nn}{java.security.PrivateKey}\PY{p}{;}
\PY{k+kn}{import}\PY{+w}{ }\PY{n+nn}{java.security.PublicKey}\PY{p}{;}

\PY{c+c1}{// Signing}
\PY{n}{Signature}\PY{+w}{ }\PY{n}{signer}\PY{+w}{ }\PY{o}{=}\PY{+w}{ }\PY{n}{Signature}\PY{p}{.}\PY{n+na}{getInstance}\PY{p}{(}\PY{l+s}{\PYZdq{}}\PY{l+s}{SHA256withRSA}\PY{l+s}{\PYZdq{}}\PY{p}{)}\PY{p}{;}
\PY{n}{signer}\PY{p}{.}\PY{n+na}{initSign}\PY{p}{(}\PY{n}{privateKey}\PY{p}{)}\PY{p}{;}
\PY{n}{signer}\PY{p}{.}\PY{n+na}{update}\PY{p}{(}\PY{n}{data}\PY{p}{)}\PY{p}{;}
\PY{k+kt}{byte}\PY{o}{[}\PY{o}{]}\PY{+w}{ }\PY{n}{signatureBytes}\PY{+w}{ }\PY{o}{=}\PY{+w}{ }\PY{n}{signer}\PY{p}{.}\PY{n+na}{sign}\PY{p}{(}\PY{p}{)}\PY{p}{;}

\PY{c+c1}{// Verifying}
\PY{n}{Signature}\PY{+w}{ }\PY{n}{verifier}\PY{+w}{ }\PY{o}{=}\PY{+w}{ }\PY{n}{Signature}\PY{p}{.}\PY{n+na}{getInstance}\PY{p}{(}\PY{l+s}{\PYZdq{}}\PY{l+s}{SHA256withRSA}\PY{l+s}{\PYZdq{}}\PY{p}{)}\PY{p}{;}
\PY{n}{verifier}\PY{p}{.}\PY{n+na}{initVerify}\PY{p}{(}\PY{n}{publicKey}\PY{p}{)}\PY{p}{;}
\PY{n}{verifier}\PY{p}{.}\PY{n+na}{update}\PY{p}{(}\PY{n}{data}\PY{p}{)}\PY{p}{;}
\PY{k+kt}{boolean}\PY{+w}{ }\PY{n}{valid}\PY{+w}{ }\PY{o}{=}\PY{+w}{ }\PY{n}{verifier}\PY{p}{.}\PY{n+na}{verify}\PY{p}{(}\PY{n}{signatureBytes}\PY{p}{)}\PY{p}{;}
\end{Verbatim}
\end{codebox}

\subsection[KeyStore — storing keys and certificates]{KeyStore — storing keys and certificates}
\label{h:18:keystore--storing-keys-and-certificates}
A \code{KeyStore} is a protected container for keys and certificates, typically backed by a file (\code{PKCS12} format is the modern standard; the older proprietary \code{JKS} format still appears in legacy code).

\begin{codebox}
\begin{Verbatim}[commandchars=\\\{\}]
\PY{k+kn}{import}\PY{+w}{ }\PY{n+nn}{java.security.KeyStore}\PY{p}{;}

\PY{n}{KeyStore}\PY{+w}{ }\PY{n}{keyStore}\PY{+w}{ }\PY{o}{=}\PY{+w}{ }\PY{n}{KeyStore}\PY{p}{.}\PY{n+na}{getInstance}\PY{p}{(}\PY{l+s}{\PYZdq{}}\PY{l+s}{PKCS12}\PY{l+s}{\PYZdq{}}\PY{p}{)}\PY{p}{;}
\PY{k}{try}\PY{+w}{ }\PY{p}{(}\PY{k+kd}{var}\PY{+w}{ }\PY{n}{in}\PY{+w}{ }\PY{o}{=}\PY{+w}{ }\PY{n}{java}\PY{p}{.}\PY{n+na}{nio}\PY{p}{.}\PY{n+na}{file}\PY{p}{.}\PY{n+na}{Files}\PY{p}{.}\PY{n+na}{newInputStream}\PY{p}{(}\PY{n}{java}\PY{p}{.}\PY{n+na}{nio}\PY{p}{.}\PY{n+na}{file}\PY{p}{.}\PY{n+na}{Path}\PY{p}{.}\PY{n+na}{of}\PY{p}{(}\PY{l+s}{\PYZdq{}}\PY{l+s}{app\PYZhy{}keys.p12}\PY{l+s}{\PYZdq{}}\PY{p}{)}\PY{p}{)}\PY{p}{)}\PY{+w}{ }\PY{p}{\PYZob{}}
\PY{+w}{    }\PY{n}{keyStore}\PY{p}{.}\PY{n+na}{load}\PY{p}{(}\PY{n}{in}\PY{p}{,}\PY{+w}{ }\PY{l+s}{\PYZdq{}}\PY{l+s}{keystorePassword}\PY{l+s}{\PYZdq{}}\PY{p}{.}\PY{n+na}{toCharArray}\PY{p}{(}\PY{p}{)}\PY{p}{)}\PY{p}{;}
\PY{p}{\PYZcb{}}
\PY{n}{java}\PY{p}{.}\PY{n+na}{security}\PY{p}{.}\PY{n+na}{cert}\PY{p}{.}\PY{n+na}{Certificate}\PY{+w}{ }\PY{n}{cert}\PY{+w}{ }\PY{o}{=}\PY{+w}{ }\PY{n}{keyStore}\PY{p}{.}\PY{n+na}{getCertificate}\PY{p}{(}\PY{l+s}{\PYZdq{}}\PY{l+s}{my\PYZhy{}alias}\PY{l+s}{\PYZdq{}}\PY{p}{)}\PY{p}{;}
\end{Verbatim}
\end{codebox}

\subsection[SecureRandom seeding]{SecureRandom seeding}
\label{h:18:securerandom-seeding}
\code{SecureRandom} seeds itself automatically from the operating system’s entropy source (e.g. \code{/\allowbreak{}dev/\allowbreak{}urandom} on Linux) — you generally should \textbf{not} manually reseed it with a predictable value, which would undermine its unpredictability.

\begin{codebox}
\begin{Verbatim}[commandchars=\\\{\}]
\PY{c+c1}{// Anti\PYZhy{}pattern: reseeding with a predictable value defeats the purpose}
\PY{n}{SecureRandom}\PY{+w}{ }\PY{n}{random}\PY{+w}{ }\PY{o}{=}\PY{+w}{ }\PY{k}{new}\PY{+w}{ }\PY{n}{SecureRandom}\PY{p}{(}\PY{p}{)}\PY{p}{;}
\PY{n}{random}\PY{p}{.}\PY{n+na}{setSeed}\PY{p}{(}\PY{n}{System}\PY{p}{.}\PY{n+na}{currentTimeMillis}\PY{p}{(}\PY{p}{)}\PY{p}{)}\PY{p}{;}\PY{+w}{ }\PY{c+c1}{// now guessable from the approximate time}

\PY{c+c1}{// Fix: just use the default seeding — it\PYZsq{}s already good}
\PY{n}{SecureRandom}\PY{+w}{ }\PY{n}{random}\PY{+w}{ }\PY{o}{=}\PY{+w}{ }\PY{k}{new}\PY{+w}{ }\PY{n}{SecureRandom}\PY{p}{(}\PY{p}{)}\PY{p}{;}
\PY{k+kt}{byte}\PY{o}{[}\PY{o}{]}\PY{+w}{ }\PY{n}{output}\PY{+w}{ }\PY{o}{=}\PY{+w}{ }\PY{k}{new}\PY{+w}{ }\PY{k+kt}{byte}\PY{o}{[}\PY{l+m+mi}{32}\PY{o}{]}\PY{p}{;}
\PY{n}{random}\PY{p}{.}\PY{n+na}{nextBytes}\PY{p}{(}\PY{n}{output}\PY{p}{)}\PY{p}{;}
\end{Verbatim}
\end{codebox}

\subsection[Algorithm do / don’t]{Algorithm do / don’t}
\label{h:18:algorithm-do--dont}
{\tablefontsmall
\begin{xltabular}{\linewidth}{@{}L{0.742}L{1.060}L{0.848}L{1.351}@{}}
\toprule
\rowcolor{tablehdr}
\thd{Purpose} & \thd{Do use} & \thd{Don’t use} & \thd{Why} \\
\midrule
\endhead
Hashing (integrity) & SHA-256, SHA-3 & MD5, SHA-1 & MD5 and SHA-1 have known collision weaknesses \\
Password hashing & bcrypt, argon2, PBKDF2 (many iterations) & MD5, SHA-1, plain SHA-256, unsalted anything & fast hashes are crackable at massive scale; no salt means identical inputs collide \\
Symmetric encryption & AES-GCM, AES-256 & DES, 3DES, RC4, AES-ECB & DES/RC4 are broken; ECB leaks patterns \\
Asymmetric encryption/signing & RSA (≥3072-bit), ECDSA (P-256/P-384), Ed25519 & RSA (\textless{}2048-bit) & short keys are brute-forceable with modern hardware \\
Random values (security use) & \code{SecureRandom} & \code{java.\allowbreak{}util.\allowbreak{}Random}, \code{Math.\allowbreak{}random()} & predictable PRNGs leak future/past values \\
Transport security & TLS 1.2 / TLS 1.3 & SSL 2.0/3.0, TLS 1.0/1.1 & old SSL/TLS versions have known protocol flaws \\
Key comparison & \code{MessageDigest.\allowbreak{}isEqual} & \code{==}, \code{Arrays.\allowbreak{}equals}, \code{String.\allowbreak{}equals} & early-exit comparisons leak timing information \\
\bottomrule
\end{xltabular}
}

\section[Common Code-Review Interview Pitfalls]{Common Code-Review Interview Pitfalls}
\label{h:18:common-code-review-interview-pitfalls}
\begin{enumerate}
\item 
\textbf{Opening a socket, \code{HttpURLConnection}, or \code{HttpClient} request with no connect/read timeout.}
Why it matters: a single unresponsive dependency can hang a thread forever, and enough hung threads exhaust the whole pool — a real production outage, not a theoretical risk.

\begin{codebox}
\begin{Verbatim}[commandchars=\\\{\}]
\PY{c+c1}{// Before}
\PY{n}{Socket}\PY{+w}{ }\PY{n}{socket}\PY{+w}{ }\PY{o}{=}\PY{+w}{ }\PY{k}{new}\PY{+w}{ }\PY{n}{Socket}\PY{p}{(}\PY{l+s}{\PYZdq{}}\PY{l+s}{host}\PY{l+s}{\PYZdq{}}\PY{p}{,}\PY{+w}{ }\PY{l+m+mi}{443}\PY{p}{)}\PY{p}{;}
\PY{c+c1}{// After}
\PY{n}{Socket}\PY{+w}{ }\PY{n}{socket}\PY{+w}{ }\PY{o}{=}\PY{+w}{ }\PY{k}{new}\PY{+w}{ }\PY{n}{Socket}\PY{p}{(}\PY{p}{)}\PY{p}{;}
\PY{n}{socket}\PY{p}{.}\PY{n+na}{connect}\PY{p}{(}\PY{k}{new}\PY{+w}{ }\PY{n}{InetSocketAddress}\PY{p}{(}\PY{l+s}{\PYZdq{}}\PY{l+s}{host}\PY{l+s}{\PYZdq{}}\PY{p}{,}\PY{+w}{ }\PY{l+m+mi}{443}\PY{p}{)}\PY{p}{,}\PY{+w}{ }\PY{l+m+mi}{3000}\PY{p}{)}\PY{p}{;}
\PY{n}{socket}\PY{p}{.}\PY{n+na}{setSoTimeout}\PY{p}{(}\PY{l+m+mi}{5000}\PY{p}{)}\PY{p}{;}
\end{Verbatim}
\end{codebox}

\item 
\textbf{Storing \code{URL} objects in a \code{HashSet}/\code{HashMap} or comparing them with \code{equals}.}
Why it matters: \code{URL.\allowbreak{}equals}/\code{hashCode} perform a blocking DNS lookup, which can silently stall the application and even give inconsistent results if DNS changes.

\begin{codebox}
\begin{Verbatim}[commandchars=\\\{\}]
\PY{c+c1}{// Before}
\PY{n}{Set}\PY{o}{\PYZlt{}}\PY{n}{URL}\PY{o}{\PYZgt{}}\PY{+w}{ }\PY{n}{cache}\PY{+w}{ }\PY{o}{=}\PY{+w}{ }\PY{k}{new}\PY{+w}{ }\PY{n}{HashSet}\PY{o}{\PYZlt{}}\PY{o}{\PYZgt{}}\PY{p}{(}\PY{p}{)}\PY{p}{;}
\PY{c+c1}{// After}
\PY{n}{Set}\PY{o}{\PYZlt{}}\PY{n}{URI}\PY{o}{\PYZgt{}}\PY{+w}{ }\PY{n}{cache}\PY{+w}{ }\PY{o}{=}\PY{+w}{ }\PY{k}{new}\PY{+w}{ }\PY{n}{HashSet}\PY{o}{\PYZlt{}}\PY{o}{\PYZgt{}}\PY{p}{(}\PY{p}{)}\PY{p}{;}\PY{+w}{ }\PY{c+c1}{// URI never touches the network}
\end{Verbatim}
\end{codebox}

\item 
\textbf{Creating a brand-new \code{HttpClient} for every request instead of sharing one instance.}
Why it matters: \code{HttpClient} manages connection pools; recreating it per request throws away pooling and HTTP/2 multiplexing benefits and adds unnecessary handshake overhead.

\begin{codebox}
\begin{Verbatim}[commandchars=\\\{\}]
\PY{c+c1}{// Before}
\PY{n}{HttpResponse}\PY{o}{\PYZlt{}}\PY{n}{String}\PY{o}{\PYZgt{}}\PY{+w}{ }\PY{n}{r}\PY{+w}{ }\PY{o}{=}\PY{+w}{ }\PY{n}{HttpClient}\PY{p}{.}\PY{n+na}{newHttpClient}\PY{p}{(}\PY{p}{)}\PY{p}{.}\PY{n+na}{send}\PY{p}{(}\PY{n}{req}\PY{p}{,}\PY{+w}{ }\PY{n}{BodyHandlers}\PY{p}{.}\PY{n+na}{ofString}\PY{p}{(}\PY{p}{)}\PY{p}{)}\PY{p}{;}
\PY{c+c1}{// After}
\PY{k+kd}{private}\PY{+w}{ }\PY{k+kd}{static}\PY{+w}{ }\PY{k+kd}{final}\PY{+w}{ }\PY{n}{HttpClient}\PY{+w}{ }\PY{n}{CLIENT}\PY{+w}{ }\PY{o}{=}\PY{+w}{ }\PY{n}{HttpClient}\PY{p}{.}\PY{n+na}{newHttpClient}\PY{p}{(}\PY{p}{)}\PY{p}{;}\PY{+w}{ }\PY{c+c1}{// built once, reused everywhere}
\end{Verbatim}
\end{codebox}

\item 
\textbf{Treating any non-thrown \code{HttpClient} response as a success.}
Why it matters: 4xx/5xx responses don’t throw exceptions in the HTTP Client API — skipping the status-code check means error bodies get processed as if they were valid data.

\begin{codebox}
\begin{Verbatim}[commandchars=\\\{\}]
\PY{c+c1}{// Before}
\PY{n}{processUser}\PY{p}{(}\PY{n}{response}\PY{p}{.}\PY{n+na}{body}\PY{p}{(}\PY{p}{)}\PY{p}{)}\PY{p}{;}
\PY{c+c1}{// After}
\PY{k}{if}\PY{+w}{ }\PY{p}{(}\PY{n}{response}\PY{p}{.}\PY{n+na}{statusCode}\PY{p}{(}\PY{p}{)}\PY{+w}{ }\PY{o}{/}\PY{+w}{ }\PY{l+m+mi}{100}\PY{+w}{ }\PY{o}{!}\PY{o}{=}\PY{+w}{ }\PY{l+m+mi}{2}\PY{p}{)}\PY{+w}{ }\PY{k}{throw}\PY{+w}{ }\PY{k}{new}\PY{+w}{ }\PY{n}{IllegalStateException}\PY{p}{(}\PY{l+s}{\PYZdq{}}\PY{l+s}{HTTP }\PY{l+s}{\PYZdq{}}\PY{+w}{ }\PY{o}{+}\PY{+w}{ }\PY{n}{response}\PY{p}{.}\PY{n+na}{statusCode}\PY{p}{(}\PY{p}{)}\PY{p}{)}\PY{p}{;}
\PY{n}{processUser}\PY{p}{(}\PY{n}{response}\PY{p}{.}\PY{n+na}{body}\PY{p}{(}\PY{p}{)}\PY{p}{)}\PY{p}{;}
\end{Verbatim}
\end{codebox}

\item 
\textbf{Building a shell command by string concatenation with untrusted input.}
Why it matters: this is textbook command injection — shell metacharacters in the input let an attacker run arbitrary commands.

\begin{codebox}
\begin{Verbatim}[commandchars=\\\{\}]
\PY{c+c1}{// Before}
\PY{n}{Runtime}\PY{p}{.}\PY{n+na}{getRuntime}\PY{p}{(}\PY{p}{)}\PY{p}{.}\PY{n+na}{exec}\PY{p}{(}\PY{l+s}{\PYZdq{}}\PY{l+s}{ls }\PY{l+s}{\PYZdq{}}\PY{+w}{ }\PY{o}{+}\PY{+w}{ }\PY{n}{userInput}\PY{p}{)}\PY{p}{;}
\PY{c+c1}{// After}
\PY{k}{new}\PY{+w}{ }\PY{n}{ProcessBuilder}\PY{p}{(}\PY{l+s}{\PYZdq{}}\PY{l+s}{ls}\PY{l+s}{\PYZdq{}}\PY{p}{,}\PY{+w}{ }\PY{n}{userInput}\PY{p}{)}\PY{p}{.}\PY{n+na}{start}\PY{p}{(}\PY{p}{)}\PY{p}{;}\PY{+w}{ }\PY{c+c1}{// arguments never parsed by a shell}
\end{Verbatim}
\end{codebox}

\item 
\textbf{Calling \code{process.\allowbreak{}waitFor()} before draining stdout/stderr, risking a deadlock.}
Why it matters: if the child’s output buffer fills up while your code is blocked waiting for exit, both sides wait on each other forever.

\begin{codebox}
\begin{Verbatim}[commandchars=\\\{\}]
\PY{c+c1}{// Before}
\PY{n}{process}\PY{p}{.}\PY{n+na}{waitFor}\PY{p}{(}\PY{p}{)}\PY{p}{;}
\PY{n}{String}\PY{+w}{ }\PY{n}{out}\PY{+w}{ }\PY{o}{=}\PY{+w}{ }\PY{k}{new}\PY{+w}{ }\PY{n}{String}\PY{p}{(}\PY{n}{process}\PY{p}{.}\PY{n+na}{getInputStream}\PY{p}{(}\PY{p}{)}\PY{p}{.}\PY{n+na}{readAllBytes}\PY{p}{(}\PY{p}{)}\PY{p}{)}\PY{p}{;}
\PY{c+c1}{// After}
\PY{c+c1}{// drain stdout/stderr concurrently in separate threads, then waitFor()}
\end{Verbatim}
\end{codebox}

\item 
\textbf{Calling \code{Process.\allowbreak{}waitFor()} with no timeout.}
Why it matters: a hung external process (bad input, infinite loop, deadlock) can block the calling thread indefinitely.

\begin{codebox}
\begin{Verbatim}[commandchars=\\\{\}]
\PY{c+c1}{// Before}
\PY{n}{process}\PY{p}{.}\PY{n+na}{waitFor}\PY{p}{(}\PY{p}{)}\PY{p}{;}
\PY{c+c1}{// After}
\PY{k}{if}\PY{+w}{ }\PY{p}{(}\PY{o}{!}\PY{n}{process}\PY{p}{.}\PY{n+na}{waitFor}\PY{p}{(}\PY{l+m+mi}{30}\PY{p}{,}\PY{+w}{ }\PY{n}{TimeUnit}\PY{p}{.}\PY{n+na}{SECONDS}\PY{p}{)}\PY{p}{)}\PY{+w}{ }\PY{n}{process}\PY{p}{.}\PY{n+na}{destroyForcibly}\PY{p}{(}\PY{p}{)}\PY{p}{;}
\end{Verbatim}
\end{codebox}

\item 
\textbf{Building SQL by string concatenation instead of using \code{PreparedStatement}.}
Why it matters: this is SQL injection — one of the most common and most damaging vulnerabilities, capable of leaking or destroying an entire database.

\begin{codebox}
\begin{Verbatim}[commandchars=\\\{\}]
\PY{c+c1}{// Before}
\PY{l+s}{\PYZdq{}}\PY{l+s}{SELECT * FROM users WHERE name = \PYZsq{}}\PY{l+s}{\PYZdq{}}\PY{+w}{ }\PY{o}{+}\PY{+w}{ }\PY{n}{name}\PY{+w}{ }\PY{o}{+}\PY{+w}{ }\PY{l+s}{\PYZdq{}}\PY{l+s}{\PYZsq{}}\PY{l+s}{\PYZdq{}}
\PY{c+c1}{// After}
\PY{n}{PreparedStatement}\PY{+w}{ }\PY{n}{ps}\PY{+w}{ }\PY{o}{=}\PY{+w}{ }\PY{n}{conn}\PY{p}{.}\PY{n+na}{prepareStatement}\PY{p}{(}\PY{l+s}{\PYZdq{}}\PY{l+s}{SELECT * FROM users WHERE name = ?}\PY{l+s}{\PYZdq{}}\PY{p}{)}\PY{p}{;}
\PY{n}{ps}\PY{p}{.}\PY{n+na}{setString}\PY{p}{(}\PY{l+m+mi}{1}\PY{p}{,}\PY{+w}{ }\PY{n}{name}\PY{p}{)}\PY{p}{;}
\end{Verbatim}
\end{codebox}

\item 
\textbf{Parsing untrusted XML with a default \code{DocumentBuilderFactory}/\code{SAXParserFactory}.}
Why it matters: default XML parsers resolve external entities, opening the door to XXE attacks that read local files or trigger SSRF.

\begin{codebox}
\begin{Verbatim}[commandchars=\\\{\}]
\PY{c+c1}{// Before}
\PY{n}{DocumentBuilderFactory}\PY{p}{.}\PY{n+na}{newInstance}\PY{p}{(}\PY{p}{)}\PY{p}{.}\PY{n+na}{newDocumentBuilder}\PY{p}{(}\PY{p}{)}\PY{p}{.}\PY{n+na}{parse}\PY{p}{(}\PY{n}{untrustedInput}\PY{p}{)}\PY{p}{;}
\PY{c+c1}{// After}
\PY{n}{factory}\PY{p}{.}\PY{n+na}{setFeature}\PY{p}{(}\PY{l+s}{\PYZdq{}}\PY{l+s}{http://apache.org/xml/features/disallow\PYZhy{}doctype\PYZhy{}decl}\PY{l+s}{\PYZdq{}}\PY{p}{,}\PY{+w}{ }\PY{k+kc}{true}\PY{p}{)}\PY{p}{;}
\PY{c+c1}{// ... plus disabling external general/parameter entities, then parse}
\end{Verbatim}
\end{codebox}

\item 
\textbf{Using a trust-all \code{TrustManager} (or \code{HostnameVerifier} that always returns true) “to get past a certificate error.”}
Why it matters: this disables TLS’s core security guarantee, making the connection trivially vulnerable to a man-in-the-middle attack.

\begin{codebox}
\begin{Verbatim}[commandchars=\\\{\}]
\PY{c+c1}{// Before}
\PY{k+kd}{public}\PY{+w}{ }\PY{k+kt}{void}\PY{+w}{ }\PY{n+nf}{checkServerTrusted}\PY{p}{(}\PY{n}{X509Certificate}\PY{o}{[}\PY{o}{]}\PY{+w}{ }\PY{n}{chain}\PY{p}{,}\PY{+w}{ }\PY{n}{String}\PY{+w}{ }\PY{n}{authType}\PY{p}{)}\PY{+w}{ }\PY{p}{\PYZob{}}\PY{p}{\PYZcb{}}\PY{+w}{ }\PY{c+c1}{// trusts anything}
\PY{c+c1}{// After}
\PY{c+c1}{// use the JVM\PYZsq{}s default trust managers, or a real trust store containing your CA}
\PY{n}{sslContext}\PY{p}{.}\PY{n+na}{init}\PY{p}{(}\PY{k+kc}{null}\PY{p}{,}\PY{+w}{ }\PY{k+kc}{null}\PY{p}{,}\PY{+w}{ }\PY{k+kc}{null}\PY{p}{)}\PY{p}{;}
\end{Verbatim}
\end{codebox}

\item 
\textbf{Hashing passwords with a single fast pass of SHA-256 (or MD5/SHA-1) instead of a slow, salted algorithm.}
Why it matters: fast hashes let attackers who steal the password table brute-force billions of guesses per second on cheap hardware.

\begin{codebox}
\begin{Verbatim}[commandchars=\\\{\}]
\PY{c+c1}{// Before}
\PY{n}{MessageDigest}\PY{p}{.}\PY{n+na}{getInstance}\PY{p}{(}\PY{l+s}{\PYZdq{}}\PY{l+s}{SHA\PYZhy{}256}\PY{l+s}{\PYZdq{}}\PY{p}{)}\PY{p}{.}\PY{n+na}{digest}\PY{p}{(}\PY{n}{password}\PY{p}{.}\PY{n+na}{getBytes}\PY{p}{(}\PY{p}{)}\PY{p}{)}\PY{p}{;}
\PY{c+c1}{// After}
\PY{n}{SecretKeyFactory}\PY{p}{.}\PY{n+na}{getInstance}\PY{p}{(}\PY{l+s}{\PYZdq{}}\PY{l+s}{PBKDF2WithHmacSHA256}\PY{l+s}{\PYZdq{}}\PY{p}{)}\PY{p}{.}\PY{n+na}{generateSecret}\PY{p}{(}
\PY{+w}{        }\PY{k}{new}\PY{+w}{ }\PY{n}{PBEKeySpec}\PY{p}{(}\PY{n}{password}\PY{p}{,}\PY{+w}{ }\PY{n}{salt}\PY{p}{,}\PY{+w}{ }\PY{l+m+mi}{210\PYZus{}000}\PY{p}{,}\PY{+w}{ }\PY{l+m+mi}{256}\PY{p}{)}\PY{p}{)}\PY{p}{;}
\end{Verbatim}
\end{codebox}

\item 
\textbf{Using AES in ECB mode.}
Why it matters: ECB leaks structural patterns in the plaintext because identical blocks always encrypt to identical ciphertext — it provides much weaker confidentiality than intended.

\begin{codebox}
\begin{Verbatim}[commandchars=\\\{\}]
\PY{c+c1}{// Before}
\PY{n}{Cipher}\PY{p}{.}\PY{n+na}{getInstance}\PY{p}{(}\PY{l+s}{\PYZdq{}}\PY{l+s}{AES/ECB/PKCS5Padding}\PY{l+s}{\PYZdq{}}\PY{p}{)}\PY{p}{;}
\PY{c+c1}{// After}
\PY{n}{Cipher}\PY{p}{.}\PY{n+na}{getInstance}\PY{p}{(}\PY{l+s}{\PYZdq{}}\PY{l+s}{AES/GCM/NoPadding}\PY{l+s}{\PYZdq{}}\PY{p}{)}\PY{p}{;}\PY{+w}{ }\PY{c+c1}{// with a fresh random IV per encryption}
\end{Verbatim}
\end{codebox}

\item 
\textbf{Reusing the same IV/nonce for multiple AES-GCM encryptions with the same key.}
Why it matters: IV reuse with GCM catastrophically breaks confidentiality and integrity guarantees — it can expose plaintext and allow forged ciphertexts.

\begin{codebox}
\begin{Verbatim}[commandchars=\\\{\}]
\PY{c+c1}{// Before}
\PY{k+kd}{static}\PY{+w}{ }\PY{k+kd}{final}\PY{+w}{ }\PY{k+kt}{byte}\PY{o}{[}\PY{o}{]}\PY{+w}{ }\PY{n}{FIXED\PYZus{}IV}\PY{+w}{ }\PY{o}{=}\PY{+w}{ }\PY{k}{new}\PY{+w}{ }\PY{k+kt}{byte}\PY{o}{[}\PY{l+m+mi}{12}\PY{o}{]}\PY{p}{;}\PY{+w}{ }\PY{c+c1}{// reused every call}
\PY{c+c1}{// After}
\PY{k+kt}{byte}\PY{o}{[}\PY{o}{]}\PY{+w}{ }\PY{n}{iv}\PY{+w}{ }\PY{o}{=}\PY{+w}{ }\PY{k}{new}\PY{+w}{ }\PY{k+kt}{byte}\PY{o}{[}\PY{l+m+mi}{12}\PY{o}{]}\PY{p}{;}
\PY{k}{new}\PY{+w}{ }\PY{n}{SecureRandom}\PY{p}{(}\PY{p}{)}\PY{p}{.}\PY{n+na}{nextBytes}\PY{p}{(}\PY{n}{iv}\PY{p}{)}\PY{p}{;}\PY{+w}{ }\PY{c+c1}{// fresh IV, generated per encryption}
\end{Verbatim}
\end{codebox}

\item 
\textbf{Using \code{java.\allowbreak{}util.\allowbreak{}Random} (or \code{Math.\allowbreak{}random()}) to generate tokens, session IDs, or keys.}
Why it matters: these PRNGs are predictable; an attacker who observes some outputs (or knows the seed) can often predict others, defeating the token’s purpose.

\begin{codebox}
\begin{Verbatim}[commandchars=\\\{\}]
\PY{c+c1}{// Before}
\PY{k}{new}\PY{+w}{ }\PY{n}{Random}\PY{p}{(}\PY{p}{)}\PY{p}{.}\PY{n+na}{nextLong}\PY{p}{(}\PY{p}{)}\PY{p}{;}
\PY{c+c1}{// After}
\PY{k}{new}\PY{+w}{ }\PY{n}{SecureRandom}\PY{p}{(}\PY{p}{)}\PY{p}{.}\PY{n+na}{nextBytes}\PY{p}{(}\PY{k}{new}\PY{+w}{ }\PY{k+kt}{byte}\PY{o}{[}\PY{l+m+mi}{32}\PY{o}{]}\PY{p}{)}\PY{p}{;}
\end{Verbatim}
\end{codebox}

\item 
\textbf{Comparing secrets (signatures, tokens, MACs) with \code{equals} instead of a constant-time comparison.}
Why it matters: \code{equals}/\code{Arrays.\allowbreak{}equals} exit early on the first mismatch, leaking timing information an attacker can use to guess the secret byte by byte.

\begin{codebox}
\begin{Verbatim}[commandchars=\\\{\}]
\PY{c+c1}{// Before}
\PY{k}{if}\PY{+w}{ }\PY{p}{(}\PY{o}{!}\PY{n}{hmac}\PY{p}{.}\PY{n+na}{equals}\PY{p}{(}\PY{n}{expected}\PY{p}{)}\PY{p}{)}\PY{+w}{ }\PY{k}{throw}\PY{+w}{ }\PY{k}{new}\PY{+w}{ }\PY{n}{SecurityException}\PY{p}{(}\PY{p}{)}\PY{p}{;}
\PY{c+c1}{// After}
\PY{k}{if}\PY{+w}{ }\PY{p}{(}\PY{o}{!}\PY{n}{MessageDigest}\PY{p}{.}\PY{n+na}{isEqual}\PY{p}{(}\PY{n}{hmac}\PY{p}{.}\PY{n+na}{getBytes}\PY{p}{(}\PY{p}{)}\PY{p}{,}\PY{+w}{ }\PY{n}{expected}\PY{p}{.}\PY{n+na}{getBytes}\PY{p}{(}\PY{p}{)}\PY{p}{)}\PY{p}{)}\PY{+w}{ }\PY{k}{throw}\PY{+w}{ }\PY{k}{new}\PY{+w}{ }\PY{n}{SecurityException}\PY{p}{(}\PY{p}{)}\PY{p}{;}
\end{Verbatim}
\end{codebox}

\item 
\textbf{Storing passwords in a \code{String} instead of a \code{char[]}, or logging request/response objects that contain secrets.}
Why it matters: \code{String} immutability means the password can linger in memory unpredictably with no way to wipe it, and logging full objects routinely leaks credentials into long-lived, widely-read log stores.

\begin{codebox}
\begin{Verbatim}[commandchars=\\\{\}]
\PY{c+c1}{// Before}
\PY{n}{log}\PY{p}{.}\PY{n+na}{info}\PY{p}{(}\PY{l+s}{\PYZdq{}}\PY{l+s}{Login attempt: \PYZob{}\PYZcb{}}\PY{l+s}{\PYZdq{}}\PY{p}{,}\PY{+w}{ }\PY{n}{loginRequest}\PY{p}{)}\PY{p}{;}\PY{+w}{ }\PY{c+c1}{// toString() includes password field}
\PY{c+c1}{// After}
\PY{n}{log}\PY{p}{.}\PY{n+na}{info}\PY{p}{(}\PY{l+s}{\PYZdq{}}\PY{l+s}{Login attempt for user: \PYZob{}\PYZcb{}}\PY{l+s}{\PYZdq{}}\PY{p}{,}\PY{+w}{ }\PY{n}{loginRequest}\PY{p}{.}\PY{n+na}{username}\PY{p}{(}\PY{p}{)}\PY{p}{)}\PY{p}{;}
\end{Verbatim}
\end{codebox}

\item 
\textbf{Fetching a user-supplied URL directly, with no host allow-list (SSRF).}
Why it matters: without restriction, an attacker can direct your server to make requests to internal-only endpoints (cloud metadata services, admin panels) it would otherwise never reach.

\begin{codebox}
\begin{Verbatim}[commandchars=\\\{\}]
\PY{c+c1}{// Before}
\PY{n}{client}\PY{p}{.}\PY{n+na}{send}\PY{p}{(}\PY{n}{HttpRequest}\PY{p}{.}\PY{n+na}{newBuilder}\PY{p}{(}\PY{n}{URI}\PY{p}{.}\PY{n+na}{create}\PY{p}{(}\PY{n}{userSuppliedUrl}\PY{p}{)}\PY{p}{)}\PY{p}{.}\PY{n+na}{build}\PY{p}{(}\PY{p}{)}\PY{p}{,}\PY{+w}{ }\PY{n}{handler}\PY{p}{)}\PY{p}{;}
\PY{c+c1}{// After}
\PY{k}{if}\PY{+w}{ }\PY{p}{(}\PY{o}{!}\PY{n}{ALLOWED\PYZus{}HOSTS}\PY{p}{.}\PY{n+na}{contains}\PY{p}{(}\PY{n}{URI}\PY{p}{.}\PY{n+na}{create}\PY{p}{(}\PY{n}{userSuppliedUrl}\PY{p}{)}\PY{p}{.}\PY{n+na}{getHost}\PY{p}{(}\PY{p}{)}\PY{p}{)}\PY{p}{)}\PY{+w}{ }\PY{k}{throw}\PY{+w}{ }\PY{k}{new}\PY{+w}{ }\PY{n}{SecurityException}\PY{p}{(}\PY{p}{)}\PY{p}{;}
\end{Verbatim}
\end{codebox}

\item 
\textbf{Hard-coding API keys, passwords, or encryption keys as string literals in source code.}
Why it matters: secrets committed to version control are effectively permanent and visible to anyone with repository access, including in old commits after “removal.”

\begin{codebox}
\begin{Verbatim}[commandchars=\\\{\}]
\PY{c+c1}{// Before}
\PY{k+kd}{private}\PY{+w}{ }\PY{k+kd}{static}\PY{+w}{ }\PY{k+kd}{final}\PY{+w}{ }\PY{n}{String}\PY{+w}{ }\PY{n}{API\PYZus{}KEY}\PY{+w}{ }\PY{o}{=}\PY{+w}{ }\PY{l+s}{\PYZdq{}}\PY{l+s}{sk\PYZus{}live\PYZus{}51H8x...}\PY{l+s}{\PYZdq{}}\PY{p}{;}
\PY{c+c1}{// After}
\PY{n}{String}\PY{+w}{ }\PY{n}{apiKey}\PY{+w}{ }\PY{o}{=}\PY{+w}{ }\PY{n}{System}\PY{p}{.}\PY{n+na}{getenv}\PY{p}{(}\PY{l+s}{\PYZdq{}}\PY{l+s}{API\PYZus{}KEY}\PY{l+s}{\PYZdq{}}\PY{p}{)}\PY{p}{;}\PY{+w}{ }\PY{c+c1}{// or a secrets manager, injected at deploy time}
\end{Verbatim}
\end{codebox}

\end{enumerate}


\setcounter{chapter}{18}
\chapter{Software Engineering Best Practices}
\label{chap:19}
\label{h:19:19-software-engineering-best-practices}
Writing code that \emph{compiles} is easy. Writing code that a teammate can read, change, and trust six months from now is the actual job. This chapter is a practical tour of the habits, principles, and patterns that Java reviewers look for: clean code rules, SOLID, the highest-value tips from \emph{Effective Java}, the design patterns that keep showing up in interviews, API design guidelines, and the code smells that get flagged in almost every pull request. We target Java 21+ and use records, sealed types, and pattern matching wherever they simplify the classic advice.

\section[Clean Code in Java]{Clean Code in Java}
\label{h:19:clean-code-in-java}
“Clean code” means code that is easy to read, easy to change, and does what its name says. It is not about being clever. A reviewer should be able to understand a method in a few seconds, without running it.

\subsection[Naming]{Naming}
\label{h:19:naming}
Names are the cheapest form of documentation. A good name removes the need for a comment.

\begin{codebox}
\begin{Verbatim}[commandchars=\\\{\}]
\PY{c+c1}{// Before: what is d? what is 86400?}
\PY{k+kt}{int}\PY{+w}{ }\PY{n}{d}\PY{+w}{ }\PY{o}{=}\PY{+w}{ }\PY{l+m+mi}{86400}\PY{p}{;}
\PY{k+kt}{boolean}\PY{+w}{ }\PY{n}{f}\PY{+w}{ }\PY{o}{=}\PY{+w}{ }\PY{n}{u}\PY{p}{.}\PY{n+na}{getA}\PY{p}{(}\PY{p}{)}\PY{+w}{ }\PY{o}{\PYZgt{}}\PY{+w}{ }\PY{n}{d}\PY{p}{;}

\PY{c+c1}{// After: intention is obvious without a comment}
\PY{k+kt}{int}\PY{+w}{ }\PY{n}{secondsPerDay}\PY{+w}{ }\PY{o}{=}\PY{+w}{ }\PY{l+m+mi}{86\PYZus{}400}\PY{p}{;}
\PY{k+kt}{boolean}\PY{+w}{ }\PY{n}{accountIsOverdue}\PY{+w}{ }\PY{o}{=}\PY{+w}{ }\PY{n}{user}\PY{p}{.}\PY{n+na}{getAccountAgeInSeconds}\PY{p}{(}\PY{p}{)}\PY{+w}{ }\PY{o}{\PYZgt{}}\PY{+w}{ }\PY{n}{secondsPerDay}\PY{p}{;}
\end{Verbatim}
\end{codebox}

Rules of thumb:

\begin{itemize}
\item 
Classes and variables: nouns (\code{OrderService}, \code{pendingInvoices}).

\item 
Methods: verbs (\code{calculateTotal}, \code{isValid}).

\item 
Booleans: read like a yes/no question (\code{isEmpty}, \code{hasPermission}, \code{canRetry}).

\item 
Avoid abbreviations that only the author understands (\code{usrRegSvc} → \code{user\allowbreak{}Registration\allowbreak{}Service}).

\item 
Avoid names that lie. A \code{List\textless{}\allowbreak{}Order\textgreater{}\allowbreak{}~\allowbreak{}orderList} that later becomes a \code{Set} is now a lying name — just call it \code{orders}.

\end{itemize}

\subsection[Function Size and Single Level of Abstraction]{Function Size and Single Level of Abstraction}
\label{h:19:function-size-and-single-level-of-abstraction}
A function should do one thing, and every line inside it should be at roughly the same “zoom level.” Mixing high-level steps (“place order”) with low-level details (string parsing, raw loops) in one method forces the reader to context-switch constantly.

\begin{codebox}
\begin{Verbatim}[commandchars=\\\{\}]
\PY{c+c1}{// Before: one method, three levels of abstraction mixed together}
\PY{k+kd}{public}\PY{+w}{ }\PY{k+kt}{void}\PY{+w}{ }\PY{n+nf}{placeOrder}\PY{p}{(}\PY{n}{Order}\PY{+w}{ }\PY{n}{order}\PY{p}{)}\PY{+w}{ }\PY{p}{\PYZob{}}
\PY{+w}{    }\PY{c+c1}{// high\PYZhy{}level: validate}
\PY{+w}{    }\PY{k}{if}\PY{+w}{ }\PY{p}{(}\PY{n}{order}\PY{p}{.}\PY{n+na}{getItems}\PY{p}{(}\PY{p}{)}\PY{p}{.}\PY{n+na}{isEmpty}\PY{p}{(}\PY{p}{)}\PY{p}{)}\PY{+w}{ }\PY{p}{\PYZob{}}
\PY{+w}{        }\PY{k}{throw}\PY{+w}{ }\PY{k}{new}\PY{+w}{ }\PY{n}{IllegalArgumentException}\PY{p}{(}\PY{l+s}{\PYZdq{}}\PY{l+s}{Order has no items}\PY{l+s}{\PYZdq{}}\PY{p}{)}\PY{p}{;}
\PY{+w}{    }\PY{p}{\PYZcb{}}
\PY{+w}{    }\PY{c+c1}{// low\PYZhy{}level: compute total by hand}
\PY{+w}{    }\PY{n}{BigDecimal}\PY{+w}{ }\PY{n}{total}\PY{+w}{ }\PY{o}{=}\PY{+w}{ }\PY{n}{BigDecimal}\PY{p}{.}\PY{n+na}{ZERO}\PY{p}{;}
\PY{+w}{    }\PY{k}{for}\PY{+w}{ }\PY{p}{(}\PY{n}{OrderItem}\PY{+w}{ }\PY{n}{item}\PY{+w}{ }\PY{p}{:}\PY{+w}{ }\PY{n}{order}\PY{p}{.}\PY{n+na}{getItems}\PY{p}{(}\PY{p}{)}\PY{p}{)}\PY{+w}{ }\PY{p}{\PYZob{}}
\PY{+w}{        }\PY{n}{total}\PY{+w}{ }\PY{o}{=}\PY{+w}{ }\PY{n}{total}\PY{p}{.}\PY{n+na}{add}\PY{p}{(}\PY{n}{item}\PY{p}{.}\PY{n+na}{getPrice}\PY{p}{(}\PY{p}{)}\PY{p}{.}\PY{n+na}{multiply}\PY{p}{(}\PY{n}{BigDecimal}\PY{p}{.}\PY{n+na}{valueOf}\PY{p}{(}\PY{n}{item}\PY{p}{.}\PY{n+na}{getQuantity}\PY{p}{(}\PY{p}{)}\PY{p}{)}\PY{p}{)}\PY{p}{)}\PY{p}{;}
\PY{+w}{    }\PY{p}{\PYZcb{}}
\PY{+w}{    }\PY{c+c1}{// low\PYZhy{}level: talk to payment gateway}
\PY{+w}{    }\PY{n}{HttpURLConnection}\PY{+w}{ }\PY{n}{conn}\PY{+w}{ }\PY{o}{=}\PY{+w}{ }\PY{k+kc}{null}\PY{p}{;}
\PY{+w}{    }\PY{k}{try}\PY{+w}{ }\PY{p}{\PYZob{}}
\PY{+w}{        }\PY{n}{conn}\PY{+w}{ }\PY{o}{=}\PY{+w}{ }\PY{p}{(}\PY{n}{HttpURLConnection}\PY{p}{)}\PY{+w}{ }\PY{k}{new}\PY{+w}{ }\PY{n}{URL}\PY{p}{(}\PY{l+s}{\PYZdq{}}\PY{l+s}{https://pay.example.com}\PY{l+s}{\PYZdq{}}\PY{p}{)}\PY{p}{.}\PY{n+na}{openConnection}\PY{p}{(}\PY{p}{)}\PY{p}{;}
\PY{+w}{        }\PY{n}{conn}\PY{p}{.}\PY{n+na}{setDoOutput}\PY{p}{(}\PY{k+kc}{true}\PY{p}{)}\PY{p}{;}
\PY{+w}{        }\PY{c+c1}{// ... write request body, read response ...}
\PY{+w}{    }\PY{p}{\PYZcb{}}\PY{+w}{ }\PY{k}{catch}\PY{+w}{ }\PY{p}{(}\PY{n}{IOException}\PY{+w}{ }\PY{n}{e}\PY{p}{)}\PY{+w}{ }\PY{p}{\PYZob{}}
\PY{+w}{        }\PY{k}{throw}\PY{+w}{ }\PY{k}{new}\PY{+w}{ }\PY{n}{PaymentException}\PY{p}{(}\PY{n}{e}\PY{p}{)}\PY{p}{;}
\PY{+w}{    }\PY{p}{\PYZcb{}}
\PY{+w}{    }\PY{c+c1}{// high\PYZhy{}level: persist}
\PY{+w}{    }\PY{n}{orderRepository}\PY{p}{.}\PY{n+na}{save}\PY{p}{(}\PY{n}{order}\PY{p}{)}\PY{p}{;}
\PY{p}{\PYZcb{}}
\end{Verbatim}
\end{codebox}

\begin{codebox}
\begin{Verbatim}[commandchars=\\\{\}]
\PY{c+c1}{// After: each method reads at one level of abstraction}
\PY{k+kd}{public}\PY{+w}{ }\PY{k+kt}{void}\PY{+w}{ }\PY{n+nf}{placeOrder}\PY{p}{(}\PY{n}{Order}\PY{+w}{ }\PY{n}{order}\PY{p}{)}\PY{+w}{ }\PY{p}{\PYZob{}}
\PY{+w}{    }\PY{n}{validate}\PY{p}{(}\PY{n}{order}\PY{p}{)}\PY{p}{;}
\PY{+w}{    }\PY{n}{BigDecimal}\PY{+w}{ }\PY{n}{total}\PY{+w}{ }\PY{o}{=}\PY{+w}{ }\PY{n}{order}\PY{p}{.}\PY{n+na}{calculateTotal}\PY{p}{(}\PY{p}{)}\PY{p}{;}
\PY{+w}{    }\PY{n}{paymentGateway}\PY{p}{.}\PY{n+na}{charge}\PY{p}{(}\PY{n}{order}\PY{p}{.}\PY{n+na}{getCustomerId}\PY{p}{(}\PY{p}{)}\PY{p}{,}\PY{+w}{ }\PY{n}{total}\PY{p}{)}\PY{p}{;}
\PY{+w}{    }\PY{n}{orderRepository}\PY{p}{.}\PY{n+na}{save}\PY{p}{(}\PY{n}{order}\PY{p}{)}\PY{p}{;}
\PY{p}{\PYZcb{}}

\PY{k+kd}{private}\PY{+w}{ }\PY{k+kt}{void}\PY{+w}{ }\PY{n+nf}{validate}\PY{p}{(}\PY{n}{Order}\PY{+w}{ }\PY{n}{order}\PY{p}{)}\PY{+w}{ }\PY{p}{\PYZob{}}
\PY{+w}{    }\PY{k}{if}\PY{+w}{ }\PY{p}{(}\PY{n}{order}\PY{p}{.}\PY{n+na}{getItems}\PY{p}{(}\PY{p}{)}\PY{p}{.}\PY{n+na}{isEmpty}\PY{p}{(}\PY{p}{)}\PY{p}{)}\PY{+w}{ }\PY{p}{\PYZob{}}
\PY{+w}{        }\PY{k}{throw}\PY{+w}{ }\PY{k}{new}\PY{+w}{ }\PY{n}{IllegalArgumentException}\PY{p}{(}\PY{l+s}{\PYZdq{}}\PY{l+s}{Order has no items}\PY{l+s}{\PYZdq{}}\PY{p}{)}\PY{p}{;}
\PY{+w}{    }\PY{p}{\PYZcb{}}
\PY{p}{\PYZcb{}}
\end{Verbatim}
\end{codebox}

The payment call and the total calculation moved into their own well-named units (\code{paymentGateway.\allowbreak{}charge}, \code{order.\allowbreak{}calculateTotal}). \code{placeOrder} now reads like a short story: validate, total, charge, save.

\subsection[Argument Count]{Argument Count}
\label{h:19:argument-count}
More than three or four parameters is a warning sign. Callers can pass arguments in the wrong order, and the call site becomes unreadable.

\begin{codebox}
\begin{Verbatim}[commandchars=\\\{\}]
\PY{c+c1}{// Before: which boolean is which? easy to swap}
\PY{n}{createUser}\PY{p}{(}\PY{l+s}{\PYZdq{}}\PY{l+s}{Ana}\PY{l+s}{\PYZdq{}}\PY{p}{,}\PY{+w}{ }\PY{l+s}{\PYZdq{}}\PY{l+s}{ana@example.com}\PY{l+s}{\PYZdq{}}\PY{p}{,}\PY{+w}{ }\PY{k+kc}{true}\PY{p}{,}\PY{+w}{ }\PY{k+kc}{false}\PY{p}{,}\PY{+w}{ }\PY{k+kc}{true}\PY{p}{)}\PY{p}{;}

\PY{c+c1}{// After: a small parameter object makes the call self\PYZhy{}explanatory}
\PY{k+kd}{record} \PY{n+nc}{UserCreationRequest}\PY{p}{(}\PY{n}{String}\PY{+w}{ }\PY{n}{name}\PY{p}{,}\PY{+w}{ }\PY{n}{String}\PY{+w}{ }\PY{n}{email}\PY{p}{,}\PY{+w}{ }\PY{k+kt}{boolean}\PY{+w}{ }\PY{n}{sendWelcomeEmail}\PY{p}{,}
\PY{+w}{                            }\PY{k+kt}{boolean}\PY{+w}{ }\PY{n}{requireEmailVerification}\PY{p}{,}\PY{+w}{ }\PY{k+kt}{boolean}\PY{+w}{ }\PY{n}{isAdmin}\PY{p}{)}\PY{+w}{ }\PY{p}{\PYZob{}}\PY{p}{\PYZcb{}}

\PY{n}{createUser}\PY{p}{(}\PY{k}{new}\PY{+w}{ }\PY{n}{UserCreationRequest}\PY{p}{(}\PY{l+s}{\PYZdq{}}\PY{l+s}{Ana}\PY{l+s}{\PYZdq{}}\PY{p}{,}\PY{+w}{ }\PY{l+s}{\PYZdq{}}\PY{l+s}{ana@example.com}\PY{l+s}{\PYZdq{}}\PY{p}{,}\PY{+w}{ }\PY{k+kc}{true}\PY{p}{,}\PY{+w}{ }\PY{k+kc}{false}\PY{p}{,}\PY{+w}{ }\PY{k+kc}{true}\PY{p}{)}\PY{p}{)}\PY{p}{;}
\end{Verbatim}
\end{codebox}

Even better, use named construction so the call site itself documents what each flag means:

\begin{codebox}
\begin{Verbatim}[commandchars=\\\{\}]
\PY{n}{createUser}\PY{p}{(}\PY{n}{UserCreationRequest}\PY{p}{.}\PY{n+na}{builder}\PY{p}{(}\PY{p}{)}
\PY{+w}{        }\PY{p}{.}\PY{n+na}{name}\PY{p}{(}\PY{l+s}{\PYZdq{}}\PY{l+s}{Ana}\PY{l+s}{\PYZdq{}}\PY{p}{)}
\PY{+w}{        }\PY{p}{.}\PY{n+na}{email}\PY{p}{(}\PY{l+s}{\PYZdq{}}\PY{l+s}{ana@example.com}\PY{l+s}{\PYZdq{}}\PY{p}{)}
\PY{+w}{        }\PY{p}{.}\PY{n+na}{sendWelcomeEmail}\PY{p}{(}\PY{k+kc}{true}\PY{p}{)}
\PY{+w}{        }\PY{p}{.}\PY{n+na}{requireEmailVerification}\PY{p}{(}\PY{k+kc}{false}\PY{p}{)}
\PY{+w}{        }\PY{p}{.}\PY{n+na}{isAdmin}\PY{p}{(}\PY{k+kc}{true}\PY{p}{)}
\PY{+w}{        }\PY{p}{.}\PY{n+na}{build}\PY{p}{(}\PY{p}{)}\PY{p}{)}\PY{p}{;}
\end{Verbatim}
\end{codebox}

\subsection[Boolean Flag Parameters]{Boolean Flag Parameters}
\label{h:19:boolean-flag-parameters}
A boolean parameter often means the method secretly does two different things. That is a violation of “one thing per function” and it is easy to misread at the call site.

\begin{codebox}
\begin{Verbatim}[commandchars=\\\{\}]
\PY{c+c1}{// Before: what does \PYZdq{}true\PYZdq{} mean here? you have to open the method to find out}
\PY{k+kt}{void}\PY{+w}{ }\PY{n+nf}{printReport}\PY{p}{(}\PY{n}{Report}\PY{+w}{ }\PY{n}{report}\PY{p}{,}\PY{+w}{ }\PY{k+kt}{boolean}\PY{+w}{ }\PY{n}{detailed}\PY{p}{)}\PY{+w}{ }\PY{p}{\PYZob{}}
\PY{+w}{    }\PY{k}{if}\PY{+w}{ }\PY{p}{(}\PY{n}{detailed}\PY{p}{)}\PY{+w}{ }\PY{p}{\PYZob{}}
\PY{+w}{        }\PY{n}{printFullReport}\PY{p}{(}\PY{n}{report}\PY{p}{)}\PY{p}{;}
\PY{+w}{    }\PY{p}{\PYZcb{}}\PY{+w}{ }\PY{k}{else}\PY{+w}{ }\PY{p}{\PYZob{}}
\PY{+w}{        }\PY{n}{printSummary}\PY{p}{(}\PY{n}{report}\PY{p}{)}\PY{p}{;}
\PY{+w}{    }\PY{p}{\PYZcb{}}
\PY{p}{\PYZcb{}}
\PY{n}{printReport}\PY{p}{(}\PY{n}{report}\PY{p}{,}\PY{+w}{ }\PY{k+kc}{true}\PY{p}{)}\PY{p}{;}
\end{Verbatim}
\end{codebox}

\begin{codebox}
\begin{Verbatim}[commandchars=\\\{\}]
\PY{c+c1}{// After: split into two clearly named methods}
\PY{k+kt}{void}\PY{+w}{ }\PY{n+nf}{printDetailedReport}\PY{p}{(}\PY{n}{Report}\PY{+w}{ }\PY{n}{report}\PY{p}{)}\PY{+w}{ }\PY{p}{\PYZob{}}\PY{+w}{ }\PY{p}{.}\PY{p}{.}\PY{p}{.}\PY{+w}{ }\PY{p}{\PYZcb{}}
\PY{k+kt}{void}\PY{+w}{ }\PY{n+nf}{printSummaryReport}\PY{p}{(}\PY{n}{Report}\PY{+w}{ }\PY{n}{report}\PY{p}{)}\PY{+w}{ }\PY{p}{\PYZob{}}\PY{+w}{ }\PY{p}{.}\PY{p}{.}\PY{p}{.}\PY{+w}{ }\PY{p}{\PYZcb{}}

\PY{n}{printReport}\PY{p}{.}\PY{n+na}{printDetailedReport}\PY{p}{(}\PY{n}{report}\PY{p}{)}\PY{p}{;}
\end{Verbatim}
\end{codebox}

If the flag really represents a meaningful choice (not just “two code paths bolted together”), prefer an enum — it is self-documenting and the compiler catches typos that a \code{boolean} never would.

\begin{codebox}
\begin{Verbatim}[commandchars=\\\{\}]
\PY{k+kd}{enum}\PY{+w}{ }\PY{n}{ReportDetail}\PY{+w}{ }\PY{p}{\PYZob{}}\PY{+w}{ }\PY{n}{SUMMARY}\PY{p}{,}\PY{+w}{ }\PY{n}{DETAILED}\PY{+w}{ }\PY{p}{\PYZcb{}}

\PY{k+kt}{void}\PY{+w}{ }\PY{n+nf}{printReport}\PY{p}{(}\PY{n}{Report}\PY{+w}{ }\PY{n}{report}\PY{p}{,}\PY{+w}{ }\PY{n}{ReportDetail}\PY{+w}{ }\PY{n}{detail}\PY{p}{)}\PY{+w}{ }\PY{p}{\PYZob{}}
\PY{+w}{    }\PY{k}{switch}\PY{+w}{ }\PY{p}{(}\PY{n}{detail}\PY{p}{)}\PY{+w}{ }\PY{p}{\PYZob{}}
\PY{+w}{        }\PY{k}{case}\PY{+w}{ }\PY{n}{SUMMARY}\PY{+w}{ }\PY{o}{\PYZhy{}}\PY{o}{\PYZgt{}}\PY{+w}{ }\PY{n}{printSummary}\PY{p}{(}\PY{n}{report}\PY{p}{)}\PY{p}{;}
\PY{+w}{        }\PY{k}{case}\PY{+w}{ }\PY{n}{DETAILED}\PY{+w}{ }\PY{o}{\PYZhy{}}\PY{o}{\PYZgt{}}\PY{+w}{ }\PY{n}{printFullReport}\PY{p}{(}\PY{n}{report}\PY{p}{)}\PY{p}{;}
\PY{+w}{    }\PY{p}{\PYZcb{}}
\PY{p}{\PYZcb{}}
\end{Verbatim}
\end{codebox}

\subsection[Avoiding Deep Nesting With Guard Clauses]{Avoiding Deep Nesting With Guard Clauses}
\label{h:19:avoiding-deep-nesting-with-guard-clauses}
Deeply nested \code{if} blocks are hard to follow because the reader has to hold every condition in their head at once. A \textbf{guard clause} — an early \code{return}/\code{throw} for the exceptional case — flattens the happy path.

\begin{codebox}
\begin{Verbatim}[commandchars=\\\{\}]
\PY{c+c1}{// Before: the \PYZdq{}real\PYZdq{} logic is buried three levels deep}
\PY{k+kd}{public}\PY{+w}{ }\PY{n}{BigDecimal}\PY{+w}{ }\PY{n+nf}{calculateDiscount}\PY{p}{(}\PY{n}{Customer}\PY{+w}{ }\PY{n}{customer}\PY{p}{,}\PY{+w}{ }\PY{n}{Order}\PY{+w}{ }\PY{n}{order}\PY{p}{)}\PY{+w}{ }\PY{p}{\PYZob{}}
\PY{+w}{    }\PY{k}{if}\PY{+w}{ }\PY{p}{(}\PY{n}{customer}\PY{+w}{ }\PY{o}{!}\PY{o}{=}\PY{+w}{ }\PY{k+kc}{null}\PY{p}{)}\PY{+w}{ }\PY{p}{\PYZob{}}
\PY{+w}{        }\PY{k}{if}\PY{+w}{ }\PY{p}{(}\PY{n}{customer}\PY{p}{.}\PY{n+na}{isActive}\PY{p}{(}\PY{p}{)}\PY{p}{)}\PY{+w}{ }\PY{p}{\PYZob{}}
\PY{+w}{            }\PY{k}{if}\PY{+w}{ }\PY{p}{(}\PY{n}{order}\PY{p}{.}\PY{n+na}{getTotal}\PY{p}{(}\PY{p}{)}\PY{p}{.}\PY{n+na}{compareTo}\PY{p}{(}\PY{n}{BigDecimal}\PY{p}{.}\PY{n+na}{ZERO}\PY{p}{)}\PY{+w}{ }\PY{o}{\PYZgt{}}\PY{+w}{ }\PY{l+m+mi}{0}\PY{p}{)}\PY{+w}{ }\PY{p}{\PYZob{}}
\PY{+w}{                }\PY{k}{return}\PY{+w}{ }\PY{n}{order}\PY{p}{.}\PY{n+na}{getTotal}\PY{p}{(}\PY{p}{)}\PY{p}{.}\PY{n+na}{multiply}\PY{p}{(}\PY{n}{customer}\PY{p}{.}\PY{n+na}{getDiscountRate}\PY{p}{(}\PY{p}{)}\PY{p}{)}\PY{p}{;}
\PY{+w}{            }\PY{p}{\PYZcb{}}\PY{+w}{ }\PY{k}{else}\PY{+w}{ }\PY{p}{\PYZob{}}
\PY{+w}{                }\PY{k}{return}\PY{+w}{ }\PY{n}{BigDecimal}\PY{p}{.}\PY{n+na}{ZERO}\PY{p}{;}
\PY{+w}{            }\PY{p}{\PYZcb{}}
\PY{+w}{        }\PY{p}{\PYZcb{}}\PY{+w}{ }\PY{k}{else}\PY{+w}{ }\PY{p}{\PYZob{}}
\PY{+w}{            }\PY{k}{return}\PY{+w}{ }\PY{n}{BigDecimal}\PY{p}{.}\PY{n+na}{ZERO}\PY{p}{;}
\PY{+w}{        }\PY{p}{\PYZcb{}}
\PY{+w}{    }\PY{p}{\PYZcb{}}\PY{+w}{ }\PY{k}{else}\PY{+w}{ }\PY{p}{\PYZob{}}
\PY{+w}{        }\PY{k}{return}\PY{+w}{ }\PY{n}{BigDecimal}\PY{p}{.}\PY{n+na}{ZERO}\PY{p}{;}
\PY{+w}{    }\PY{p}{\PYZcb{}}
\PY{p}{\PYZcb{}}
\end{Verbatim}
\end{codebox}

\begin{codebox}
\begin{Verbatim}[commandchars=\\\{\}]
\PY{c+c1}{// After: guard clauses handle the edge cases up front, happy path is flat}
\PY{k+kd}{public}\PY{+w}{ }\PY{n}{BigDecimal}\PY{+w}{ }\PY{n+nf}{calculateDiscount}\PY{p}{(}\PY{n}{Customer}\PY{+w}{ }\PY{n}{customer}\PY{p}{,}\PY{+w}{ }\PY{n}{Order}\PY{+w}{ }\PY{n}{order}\PY{p}{)}\PY{+w}{ }\PY{p}{\PYZob{}}
\PY{+w}{    }\PY{k}{if}\PY{+w}{ }\PY{p}{(}\PY{n}{customer}\PY{+w}{ }\PY{o}{=}\PY{o}{=}\PY{+w}{ }\PY{k+kc}{null}\PY{+w}{ }\PY{o}{|}\PY{o}{|}\PY{+w}{ }\PY{o}{!}\PY{n}{customer}\PY{p}{.}\PY{n+na}{isActive}\PY{p}{(}\PY{p}{)}\PY{p}{)}\PY{+w}{ }\PY{p}{\PYZob{}}
\PY{+w}{        }\PY{k}{return}\PY{+w}{ }\PY{n}{BigDecimal}\PY{p}{.}\PY{n+na}{ZERO}\PY{p}{;}
\PY{+w}{    }\PY{p}{\PYZcb{}}
\PY{+w}{    }\PY{k}{if}\PY{+w}{ }\PY{p}{(}\PY{n}{order}\PY{p}{.}\PY{n+na}{getTotal}\PY{p}{(}\PY{p}{)}\PY{p}{.}\PY{n+na}{compareTo}\PY{p}{(}\PY{n}{BigDecimal}\PY{p}{.}\PY{n+na}{ZERO}\PY{p}{)}\PY{+w}{ }\PY{o}{\PYZlt{}}\PY{o}{=}\PY{+w}{ }\PY{l+m+mi}{0}\PY{p}{)}\PY{+w}{ }\PY{p}{\PYZob{}}
\PY{+w}{        }\PY{k}{return}\PY{+w}{ }\PY{n}{BigDecimal}\PY{p}{.}\PY{n+na}{ZERO}\PY{p}{;}
\PY{+w}{    }\PY{p}{\PYZcb{}}
\PY{+w}{    }\PY{k}{return}\PY{+w}{ }\PY{n}{order}\PY{p}{.}\PY{n+na}{getTotal}\PY{p}{(}\PY{p}{)}\PY{p}{.}\PY{n+na}{multiply}\PY{p}{(}\PY{n}{customer}\PY{p}{.}\PY{n+na}{getDiscountRate}\PY{p}{(}\PY{p}{)}\PY{p}{)}\PY{p}{;}
\PY{p}{\PYZcb{}}
\end{Verbatim}
\end{codebox}

\subsection[Comments That Explain Why, Not What]{Comments That Explain Why, Not What}
\label{h:19:comments-that-explain-why-not-what}
A comment that repeats what the code already says is noise. A useful comment explains a decision the code \emph{cannot} express: why an unusual approach was chosen, a business rule, or a workaround for a bug.

\begin{codebox}
\begin{Verbatim}[commandchars=\\\{\}]
\PY{c+c1}{// Before: comment adds nothing the code doesn\PYZsq{}t already say}
\PY{c+c1}{// increment i by 1}
\PY{n}{i}\PY{o}{+}\PY{o}{+}\PY{p}{;}

\PY{c+c1}{// loop through all users}
\PY{k}{for}\PY{+w}{ }\PY{p}{(}\PY{n}{User}\PY{+w}{ }\PY{n}{user}\PY{+w}{ }\PY{p}{:}\PY{+w}{ }\PY{n}{users}\PY{p}{)}\PY{+w}{ }\PY{p}{\PYZob{}}\PY{+w}{ }\PY{p}{.}\PY{p}{.}\PY{p}{.}\PY{+w}{ }\PY{p}{\PYZcb{}}
\end{Verbatim}
\end{codebox}

\begin{codebox}
\begin{Verbatim}[commandchars=\\\{\}]
\PY{c+c1}{// After: the comment carries information the code cannot}
\PY{c+c1}{// Retry with exponential backoff: the payment provider rate\PYZhy{}limits}
\PY{c+c1}{// bursts and returns 429 without a Retry\PYZhy{}After header (see INC\PYZhy{}4821).}
\PY{k}{for}\PY{+w}{ }\PY{p}{(}\PY{k+kt}{int}\PY{+w}{ }\PY{n}{attempt}\PY{+w}{ }\PY{o}{=}\PY{+w}{ }\PY{l+m+mi}{0}\PY{p}{;}\PY{+w}{ }\PY{n}{attempt}\PY{+w}{ }\PY{o}{\PYZlt{}}\PY{+w}{ }\PY{n}{MAX\PYZus{}RETRIES}\PY{p}{;}\PY{+w}{ }\PY{n}{attempt}\PY{o}{+}\PY{o}{+}\PY{p}{)}\PY{+w}{ }\PY{p}{\PYZob{}}
\PY{+w}{    }\PY{k}{try}\PY{+w}{ }\PY{p}{\PYZob{}}
\PY{+w}{        }\PY{k}{return}\PY{+w}{ }\PY{n}{paymentClient}\PY{p}{.}\PY{n+na}{charge}\PY{p}{(}\PY{n}{request}\PY{p}{)}\PY{p}{;}
\PY{+w}{    }\PY{p}{\PYZcb{}}\PY{+w}{ }\PY{k}{catch}\PY{+w}{ }\PY{p}{(}\PY{n}{RateLimitException}\PY{+w}{ }\PY{n}{e}\PY{p}{)}\PY{+w}{ }\PY{p}{\PYZob{}}
\PY{+w}{        }\PY{n}{sleep}\PY{p}{(}\PY{n}{backoffMillis}\PY{p}{(}\PY{n}{attempt}\PY{p}{)}\PY{p}{)}\PY{p}{;}
\PY{+w}{    }\PY{p}{\PYZcb{}}
\PY{p}{\PYZcb{}}
\end{Verbatim}
\end{codebox}

\subsection[Dead Code]{Dead Code}
\label{h:19:dead-code}
Dead code — commented-out blocks, unused methods, unreachable branches — costs nothing to write but confuses every future reader, who has to figure out whether it is safe to delete. Version control already remembers old code; delete it instead of commenting it out.

\begin{codebox}
\begin{Verbatim}[commandchars=\\\{\}]
\PY{c+c1}{// Before}
\PY{k+kd}{public}\PY{+w}{ }\PY{k+kt}{int}\PY{+w}{ }\PY{n+nf}{calculateShipping}\PY{p}{(}\PY{n}{Order}\PY{+w}{ }\PY{n}{order}\PY{p}{)}\PY{+w}{ }\PY{p}{\PYZob{}}
\PY{+w}{    }\PY{c+c1}{// int oldRate = order.getWeight() * 2;}
\PY{+w}{    }\PY{c+c1}{// return oldRate + BASE\PYZus{}FEE;}
\PY{+w}{    }\PY{k}{return}\PY{+w}{ }\PY{n}{order}\PY{p}{.}\PY{n+na}{getWeight}\PY{p}{(}\PY{p}{)}\PY{+w}{ }\PY{o}{*}\PY{+w}{ }\PY{n}{NEW\PYZus{}RATE}\PY{+w}{ }\PY{o}{+}\PY{+w}{ }\PY{n}{BASE\PYZus{}FEE}\PY{p}{;}
\PY{p}{\PYZcb{}}

\PY{c+c1}{// After}
\PY{k+kd}{public}\PY{+w}{ }\PY{k+kt}{int}\PY{+w}{ }\PY{n+nf}{calculateShipping}\PY{p}{(}\PY{n}{Order}\PY{+w}{ }\PY{n}{order}\PY{p}{)}\PY{+w}{ }\PY{p}{\PYZob{}}
\PY{+w}{    }\PY{k}{return}\PY{+w}{ }\PY{n}{order}\PY{p}{.}\PY{n+na}{getWeight}\PY{p}{(}\PY{p}{)}\PY{+w}{ }\PY{o}{*}\PY{+w}{ }\PY{n}{NEW\PYZus{}RATE}\PY{+w}{ }\PY{o}{+}\PY{+w}{ }\PY{n}{BASE\PYZus{}FEE}\PY{p}{;}
\PY{p}{\PYZcb{}}
\end{Verbatim}
\end{codebox}

\subsection[Magic Numbers]{Magic Numbers}
\label{h:19:magic-numbers}
A magic number is a literal whose meaning is not obvious from context. Give it a name.

\begin{codebox}
\begin{Verbatim}[commandchars=\\\{\}]
\PY{c+c1}{// Before: why 0.075? why 3?}
\PY{n}{BigDecimal}\PY{+w}{ }\PY{n}{tax}\PY{+w}{ }\PY{o}{=}\PY{+w}{ }\PY{n}{price}\PY{p}{.}\PY{n+na}{multiply}\PY{p}{(}\PY{n}{BigDecimal}\PY{p}{.}\PY{n+na}{valueOf}\PY{p}{(}\PY{l+m+mf}{0.075}\PY{p}{)}\PY{p}{)}\PY{p}{;}
\PY{k}{if}\PY{+w}{ }\PY{p}{(}\PY{n}{failedAttempts}\PY{+w}{ }\PY{o}{\PYZgt{}}\PY{o}{=}\PY{+w}{ }\PY{l+m+mi}{3}\PY{p}{)}\PY{+w}{ }\PY{n}{lockAccount}\PY{p}{(}\PY{p}{)}\PY{p}{;}

\PY{c+c1}{// After}
\PY{k+kd}{private}\PY{+w}{ }\PY{k+kd}{static}\PY{+w}{ }\PY{k+kd}{final}\PY{+w}{ }\PY{n}{BigDecimal}\PY{+w}{ }\PY{n}{SALES\PYZus{}TAX\PYZus{}RATE}\PY{+w}{ }\PY{o}{=}\PY{+w}{ }\PY{n}{BigDecimal}\PY{p}{.}\PY{n+na}{valueOf}\PY{p}{(}\PY{l+m+mf}{0.075}\PY{p}{)}\PY{p}{;}
\PY{k+kd}{private}\PY{+w}{ }\PY{k+kd}{static}\PY{+w}{ }\PY{k+kd}{final}\PY{+w}{ }\PY{k+kt}{int}\PY{+w}{ }\PY{n}{MAX\PYZus{}LOGIN\PYZus{}ATTEMPTS}\PY{+w}{ }\PY{o}{=}\PY{+w}{ }\PY{l+m+mi}{3}\PY{p}{;}

\PY{n}{BigDecimal}\PY{+w}{ }\PY{n}{tax}\PY{+w}{ }\PY{o}{=}\PY{+w}{ }\PY{n}{price}\PY{p}{.}\PY{n+na}{multiply}\PY{p}{(}\PY{n}{SALES\PYZus{}TAX\PYZus{}RATE}\PY{p}{)}\PY{p}{;}
\PY{k}{if}\PY{+w}{ }\PY{p}{(}\PY{n}{failedAttempts}\PY{+w}{ }\PY{o}{\PYZgt{}}\PY{o}{=}\PY{+w}{ }\PY{n}{MAX\PYZus{}LOGIN\PYZus{}ATTEMPTS}\PY{p}{)}\PY{+w}{ }\PY{n}{lockAccount}\PY{p}{(}\PY{p}{)}\PY{p}{;}
\end{Verbatim}
\end{codebox}

\subsection[Consistent Formatting]{Consistent Formatting}
\label{h:19:consistent-formatting}
Consistent formatting (indentation, brace style, import ordering) is not about aesthetics — it removes noise from diffs and code review. A file that mixes tabs and spaces, or reformats unrelated lines in a “fix a typo” PR, hides the real change.

\begin{codebox}
\begin{Verbatim}[commandchars=\\\{\}]
\PY{c+c1}{// Before: inconsistent brace style and spacing makes diffs noisy}
\PY{k+kd}{public}\PY{+w}{ }\PY{k+kt}{void}\PY{+w}{ }\PY{n+nf}{run}\PY{p}{(}\PY{p}{)}\PY{p}{\PYZob{}}
\PY{+w}{  }\PY{k}{if}\PY{p}{(}\PY{n}{ready}\PY{p}{)}
\PY{+w}{  }\PY{p}{\PYZob{}}
\PY{+w}{      }\PY{n}{doWork}\PY{p}{(}\PY{p}{)}\PY{p}{;}
\PY{+w}{  }\PY{p}{\PYZcb{}}
\PY{+w}{    }\PY{k}{else}\PY{+w}{ }\PY{p}{\PYZob{}}\PY{+w}{ }\PY{n}{skip}\PY{p}{(}\PY{p}{)}\PY{p}{;}\PY{+w}{ }\PY{p}{\PYZcb{}}
\PY{p}{\PYZcb{}}
\end{Verbatim}
\end{codebox}

\begin{codebox}
\begin{Verbatim}[commandchars=\\\{\}]
\PY{c+c1}{// After: one consistent style (enforced by a formatter like Spotless/google\PYZhy{}java\PYZhy{}format)}
\PY{k+kd}{public}\PY{+w}{ }\PY{k+kt}{void}\PY{+w}{ }\PY{n+nf}{run}\PY{p}{(}\PY{p}{)}\PY{+w}{ }\PY{p}{\PYZob{}}
\PY{+w}{    }\PY{k}{if}\PY{+w}{ }\PY{p}{(}\PY{n}{ready}\PY{p}{)}\PY{+w}{ }\PY{p}{\PYZob{}}
\PY{+w}{        }\PY{n}{doWork}\PY{p}{(}\PY{p}{)}\PY{p}{;}
\PY{+w}{    }\PY{p}{\PYZcb{}}\PY{+w}{ }\PY{k}{else}\PY{+w}{ }\PY{p}{\PYZob{}}
\PY{+w}{        }\PY{n}{skip}\PY{p}{(}\PY{p}{)}\PY{p}{;}
\PY{+w}{    }\PY{p}{\PYZcb{}}
\PY{p}{\PYZcb{}}
\end{Verbatim}
\end{codebox}

In real teams this is automated with a formatter and a CI check, not manual discipline — humans are inconsistent, tools are not.

\subsection[Small Classes]{Small Classes}
\label{h:19:small-classes}
A class should have a short, focused responsibility. If you cannot describe what a class does in one sentence without using “and,” it is probably doing too much.

\begin{codebox}
\begin{Verbatim}[commandchars=\\\{\}]
\PY{c+c1}{// Before: one class doing validation, persistence, email, and formatting}
\PY{k+kd}{class} \PY{n+nc}{UserManager}\PY{+w}{ }\PY{p}{\PYZob{}}
\PY{+w}{    }\PY{k+kt}{void}\PY{+w}{ }\PY{n+nf}{validate}\PY{p}{(}\PY{n}{User}\PY{+w}{ }\PY{n}{u}\PY{p}{)}\PY{+w}{ }\PY{p}{\PYZob{}}\PY{+w}{ }\PY{p}{.}\PY{p}{.}\PY{p}{.}\PY{+w}{ }\PY{p}{\PYZcb{}}
\PY{+w}{    }\PY{k+kt}{void}\PY{+w}{ }\PY{n+nf}{save}\PY{p}{(}\PY{n}{User}\PY{+w}{ }\PY{n}{u}\PY{p}{)}\PY{+w}{ }\PY{p}{\PYZob{}}\PY{+w}{ }\PY{p}{.}\PY{p}{.}\PY{p}{.}\PY{+w}{ }\PY{p}{\PYZcb{}}
\PY{+w}{    }\PY{k+kt}{void}\PY{+w}{ }\PY{n+nf}{sendWelcomeEmail}\PY{p}{(}\PY{n}{User}\PY{+w}{ }\PY{n}{u}\PY{p}{)}\PY{+w}{ }\PY{p}{\PYZob{}}\PY{+w}{ }\PY{p}{.}\PY{p}{.}\PY{p}{.}\PY{+w}{ }\PY{p}{\PYZcb{}}
\PY{+w}{    }\PY{n}{String}\PY{+w}{ }\PY{n+nf}{formatForDisplay}\PY{p}{(}\PY{n}{User}\PY{+w}{ }\PY{n}{u}\PY{p}{)}\PY{+w}{ }\PY{p}{\PYZob{}}\PY{+w}{ }\PY{p}{.}\PY{p}{.}\PY{p}{.}\PY{+w}{ }\PY{p}{\PYZcb{}}
\PY{p}{\PYZcb{}}
\end{Verbatim}
\end{codebox}

\begin{codebox}
\begin{Verbatim}[commandchars=\\\{\}]
\PY{c+c1}{// After: each responsibility gets its own small class}
\PY{k+kd}{class} \PY{n+nc}{UserValidator}\PY{+w}{ }\PY{p}{\PYZob{}}\PY{+w}{ }\PY{k+kt}{void}\PY{+w}{ }\PY{n+nf}{validate}\PY{p}{(}\PY{n}{User}\PY{+w}{ }\PY{n}{u}\PY{p}{)}\PY{+w}{ }\PY{p}{\PYZob{}}\PY{+w}{ }\PY{p}{.}\PY{p}{.}\PY{p}{.}\PY{+w}{ }\PY{p}{\PYZcb{}}\PY{+w}{ }\PY{p}{\PYZcb{}}
\PY{k+kd}{class} \PY{n+nc}{UserRepository}\PY{+w}{ }\PY{p}{\PYZob{}}\PY{+w}{ }\PY{k+kt}{void}\PY{+w}{ }\PY{n+nf}{save}\PY{p}{(}\PY{n}{User}\PY{+w}{ }\PY{n}{u}\PY{p}{)}\PY{+w}{ }\PY{p}{\PYZob{}}\PY{+w}{ }\PY{p}{.}\PY{p}{.}\PY{p}{.}\PY{+w}{ }\PY{p}{\PYZcb{}}\PY{+w}{ }\PY{p}{\PYZcb{}}
\PY{k+kd}{class} \PY{n+nc}{WelcomeEmailSender}\PY{+w}{ }\PY{p}{\PYZob{}}\PY{+w}{ }\PY{k+kt}{void}\PY{+w}{ }\PY{n+nf}{send}\PY{p}{(}\PY{n}{User}\PY{+w}{ }\PY{n}{u}\PY{p}{)}\PY{+w}{ }\PY{p}{\PYZob{}}\PY{+w}{ }\PY{p}{.}\PY{p}{.}\PY{p}{.}\PY{+w}{ }\PY{p}{\PYZcb{}}\PY{+w}{ }\PY{p}{\PYZcb{}}
\PY{k+kd}{class} \PY{n+nc}{UserFormatter}\PY{+w}{ }\PY{p}{\PYZob{}}\PY{+w}{ }\PY{n}{String}\PY{+w}{ }\PY{n+nf}{formatForDisplay}\PY{p}{(}\PY{n}{User}\PY{+w}{ }\PY{n}{u}\PY{p}{)}\PY{+w}{ }\PY{p}{\PYZob{}}\PY{+w}{ }\PY{p}{.}\PY{p}{.}\PY{p}{.}\PY{+w}{ }\PY{p}{\PYZcb{}}\PY{+w}{ }\PY{p}{\PYZcb{}}
\end{Verbatim}
\end{codebox}

This is exactly the Single Responsibility Principle, covered next.

\section[SOLID Principles]{SOLID Principles}
\label{h:19:solid-principles}
\textbf{SOLID} is an acronym for five design principles that make object-oriented code easier to maintain and extend. Each one has a classic violation and a classic fix.

\subsection[Single Responsibility Principle (SRP)]{Single Responsibility Principle (SRP)}
\label{h:19:single-responsibility-principle-srp}
A class should have only one reason to change. If a class mixes unrelated concerns, a change to one concern risks breaking the other.

\begin{codebox}
\begin{Verbatim}[commandchars=\\\{\}]
\PY{c+c1}{// Violation: Invoice mixes business data, persistence, and printing}
\PY{k+kd}{class} \PY{n+nc}{Invoice}\PY{+w}{ }\PY{p}{\PYZob{}}
\PY{+w}{    }\PY{k+kd}{private}\PY{+w}{ }\PY{n}{List}\PY{o}{\PYZlt{}}\PY{n}{LineItem}\PY{o}{\PYZgt{}}\PY{+w}{ }\PY{n}{items}\PY{p}{;}

\PY{+w}{    }\PY{n}{BigDecimal}\PY{+w}{ }\PY{n+nf}{calculateTotal}\PY{p}{(}\PY{p}{)}\PY{+w}{ }\PY{p}{\PYZob{}}
\PY{+w}{        }\PY{k}{return}\PY{+w}{ }\PY{n}{items}\PY{p}{.}\PY{n+na}{stream}\PY{p}{(}\PY{p}{)}\PY{p}{.}\PY{n+na}{map}\PY{p}{(}\PY{n}{LineItem}\PY{p}{:}\PY{p}{:}\PY{n}{price}\PY{p}{)}\PY{p}{.}\PY{n+na}{reduce}\PY{p}{(}\PY{n}{BigDecimal}\PY{p}{.}\PY{n+na}{ZERO}\PY{p}{,}\PY{+w}{ }\PY{n}{BigDecimal}\PY{p}{:}\PY{p}{:}\PY{n}{add}\PY{p}{)}\PY{p}{;}
\PY{+w}{    }\PY{p}{\PYZcb{}}

\PY{+w}{    }\PY{k+kt}{void}\PY{+w}{ }\PY{n+nf}{saveToDatabase}\PY{p}{(}\PY{p}{)}\PY{+w}{ }\PY{p}{\PYZob{}}
\PY{+w}{        }\PY{c+c1}{// JDBC code directly inside the domain class}
\PY{+w}{    }\PY{p}{\PYZcb{}}

\PY{+w}{    }\PY{k+kt}{void}\PY{+w}{ }\PY{n+nf}{printToConsole}\PY{p}{(}\PY{p}{)}\PY{+w}{ }\PY{p}{\PYZob{}}
\PY{+w}{        }\PY{n}{System}\PY{p}{.}\PY{n+na}{out}\PY{p}{.}\PY{n+na}{println}\PY{p}{(}\PY{l+s}{\PYZdq{}}\PY{l+s}{Invoice total: }\PY{l+s}{\PYZdq{}}\PY{+w}{ }\PY{o}{+}\PY{+w}{ }\PY{n}{calculateTotal}\PY{p}{(}\PY{p}{)}\PY{p}{)}\PY{p}{;}
\PY{+w}{    }\PY{p}{\PYZcb{}}
\PY{p}{\PYZcb{}}
\end{Verbatim}
\end{codebox}

\begin{codebox}
\begin{Verbatim}[commandchars=\\\{\}]
\PY{c+c1}{// Fixed: each class has exactly one reason to change}
\PY{k+kd}{class} \PY{n+nc}{Invoice}\PY{+w}{ }\PY{p}{\PYZob{}}
\PY{+w}{    }\PY{k+kd}{private}\PY{+w}{ }\PY{k+kd}{final}\PY{+w}{ }\PY{n}{List}\PY{o}{\PYZlt{}}\PY{n}{LineItem}\PY{o}{\PYZgt{}}\PY{+w}{ }\PY{n}{items}\PY{p}{;}

\PY{+w}{    }\PY{n}{BigDecimal}\PY{+w}{ }\PY{n+nf}{calculateTotal}\PY{p}{(}\PY{p}{)}\PY{+w}{ }\PY{p}{\PYZob{}}
\PY{+w}{        }\PY{k}{return}\PY{+w}{ }\PY{n}{items}\PY{p}{.}\PY{n+na}{stream}\PY{p}{(}\PY{p}{)}\PY{p}{.}\PY{n+na}{map}\PY{p}{(}\PY{n}{LineItem}\PY{p}{:}\PY{p}{:}\PY{n}{price}\PY{p}{)}\PY{p}{.}\PY{n+na}{reduce}\PY{p}{(}\PY{n}{BigDecimal}\PY{p}{.}\PY{n+na}{ZERO}\PY{p}{,}\PY{+w}{ }\PY{n}{BigDecimal}\PY{p}{:}\PY{p}{:}\PY{n}{add}\PY{p}{)}\PY{p}{;}
\PY{+w}{    }\PY{p}{\PYZcb{}}

\PY{+w}{    }\PY{n}{List}\PY{o}{\PYZlt{}}\PY{n}{LineItem}\PY{o}{\PYZgt{}}\PY{+w}{ }\PY{n+nf}{items}\PY{p}{(}\PY{p}{)}\PY{+w}{ }\PY{p}{\PYZob{}}\PY{+w}{ }\PY{k}{return}\PY{+w}{ }\PY{n}{items}\PY{p}{;}\PY{+w}{ }\PY{p}{\PYZcb{}}
\PY{p}{\PYZcb{}}

\PY{k+kd}{class} \PY{n+nc}{InvoiceRepository}\PY{+w}{ }\PY{p}{\PYZob{}}
\PY{+w}{    }\PY{k+kt}{void}\PY{+w}{ }\PY{n+nf}{save}\PY{p}{(}\PY{n}{Invoice}\PY{+w}{ }\PY{n}{invoice}\PY{p}{)}\PY{+w}{ }\PY{p}{\PYZob{}}\PY{+w}{ }\PY{c+cm}{/* JDBC code lives here */}\PY{+w}{ }\PY{p}{\PYZcb{}}
\PY{p}{\PYZcb{}}

\PY{k+kd}{class} \PY{n+nc}{InvoicePrinter}\PY{+w}{ }\PY{p}{\PYZob{}}
\PY{+w}{    }\PY{k+kt}{void}\PY{+w}{ }\PY{n+nf}{print}\PY{p}{(}\PY{n}{Invoice}\PY{+w}{ }\PY{n}{invoice}\PY{p}{)}\PY{+w}{ }\PY{p}{\PYZob{}}
\PY{+w}{        }\PY{n}{System}\PY{p}{.}\PY{n+na}{out}\PY{p}{.}\PY{n+na}{println}\PY{p}{(}\PY{l+s}{\PYZdq{}}\PY{l+s}{Invoice total: }\PY{l+s}{\PYZdq{}}\PY{+w}{ }\PY{o}{+}\PY{+w}{ }\PY{n}{invoice}\PY{p}{.}\PY{n+na}{calculateTotal}\PY{p}{(}\PY{p}{)}\PY{p}{)}\PY{p}{;}
\PY{+w}{    }\PY{p}{\PYZcb{}}
\PY{p}{\PYZcb{}}
\end{Verbatim}
\end{codebox}

Now a change to the database schema touches only \code{InvoiceRepository}; a change to print formatting touches only \code{InvoicePrinter}. \code{Invoice} itself only changes when the business rules for an invoice change.

\subsection[Open/Closed Principle (OCP)]{Open/Closed Principle (OCP)}
\label{h:19:openclosed-principle-ocp}
Classes should be open for extension but closed for modification. You should be able to add new behavior without editing existing, already-tested code.

\begin{codebox}
\begin{Verbatim}[commandchars=\\\{\}]
\PY{c+c1}{// Violation: adding a new shape means editing this method every time}
\PY{k+kd}{class} \PY{n+nc}{AreaCalculator}\PY{+w}{ }\PY{p}{\PYZob{}}
\PY{+w}{    }\PY{k+kt}{double}\PY{+w}{ }\PY{n+nf}{area}\PY{p}{(}\PY{n}{Object}\PY{+w}{ }\PY{n}{shape}\PY{p}{)}\PY{+w}{ }\PY{p}{\PYZob{}}
\PY{+w}{        }\PY{k}{if}\PY{+w}{ }\PY{p}{(}\PY{n}{shape}\PY{+w}{ }\PY{k}{instanceof}\PY{+w}{ }\PY{n}{Circle}\PY{+w}{ }\PY{n}{c}\PY{p}{)}\PY{+w}{ }\PY{p}{\PYZob{}}
\PY{+w}{            }\PY{k}{return}\PY{+w}{ }\PY{n}{Math}\PY{p}{.}\PY{n+na}{PI}\PY{+w}{ }\PY{o}{*}\PY{+w}{ }\PY{n}{c}\PY{p}{.}\PY{n+na}{radius}\PY{p}{(}\PY{p}{)}\PY{+w}{ }\PY{o}{*}\PY{+w}{ }\PY{n}{c}\PY{p}{.}\PY{n+na}{radius}\PY{p}{(}\PY{p}{)}\PY{p}{;}
\PY{+w}{        }\PY{p}{\PYZcb{}}\PY{+w}{ }\PY{k}{else}\PY{+w}{ }\PY{k}{if}\PY{+w}{ }\PY{p}{(}\PY{n}{shape}\PY{+w}{ }\PY{k}{instanceof}\PY{+w}{ }\PY{n}{Rectangle}\PY{+w}{ }\PY{n}{r}\PY{p}{)}\PY{+w}{ }\PY{p}{\PYZob{}}
\PY{+w}{            }\PY{k}{return}\PY{+w}{ }\PY{n}{r}\PY{p}{.}\PY{n+na}{width}\PY{p}{(}\PY{p}{)}\PY{+w}{ }\PY{o}{*}\PY{+w}{ }\PY{n}{r}\PY{p}{.}\PY{n+na}{height}\PY{p}{(}\PY{p}{)}\PY{p}{;}
\PY{+w}{        }\PY{p}{\PYZcb{}}
\PY{+w}{        }\PY{c+c1}{// every new shape adds another branch here}
\PY{+w}{        }\PY{k}{throw}\PY{+w}{ }\PY{k}{new}\PY{+w}{ }\PY{n}{IllegalArgumentException}\PY{p}{(}\PY{l+s}{\PYZdq{}}\PY{l+s}{Unknown shape}\PY{l+s}{\PYZdq{}}\PY{p}{)}\PY{p}{;}
\PY{+w}{    }\PY{p}{\PYZcb{}}
\PY{p}{\PYZcb{}}
\end{Verbatim}
\end{codebox}

\begin{codebox}
\begin{Verbatim}[commandchars=\\\{\}]
\PY{c+c1}{// Fixed: new shapes extend the abstraction, AreaCalculator never changes}
\PY{k+kd}{interface} \PY{n+nc}{Shape}\PY{+w}{ }\PY{p}{\PYZob{}}
\PY{+w}{    }\PY{k+kt}{double}\PY{+w}{ }\PY{n+nf}{area}\PY{p}{(}\PY{p}{)}\PY{p}{;}
\PY{p}{\PYZcb{}}

\PY{k+kd}{record} \PY{n+nc}{Circle}\PY{p}{(}\PY{k+kt}{double}\PY{+w}{ }\PY{n}{radius}\PY{p}{)}\PY{+w}{ }\PY{k+kd}{implements}\PY{+w}{ }\PY{n}{Shape}\PY{+w}{ }\PY{p}{\PYZob{}}
\PY{+w}{    }\PY{k+kd}{public}\PY{+w}{ }\PY{k+kt}{double}\PY{+w}{ }\PY{n+nf}{area}\PY{p}{(}\PY{p}{)}\PY{+w}{ }\PY{p}{\PYZob{}}\PY{+w}{ }\PY{k}{return}\PY{+w}{ }\PY{n}{Math}\PY{p}{.}\PY{n+na}{PI}\PY{+w}{ }\PY{o}{*}\PY{+w}{ }\PY{n}{radius}\PY{+w}{ }\PY{o}{*}\PY{+w}{ }\PY{n}{radius}\PY{p}{;}\PY{+w}{ }\PY{p}{\PYZcb{}}
\PY{p}{\PYZcb{}}

\PY{k+kd}{record} \PY{n+nc}{Rectangle}\PY{p}{(}\PY{k+kt}{double}\PY{+w}{ }\PY{n}{width}\PY{p}{,}\PY{+w}{ }\PY{k+kt}{double}\PY{+w}{ }\PY{n}{height}\PY{p}{)}\PY{+w}{ }\PY{k+kd}{implements}\PY{+w}{ }\PY{n}{Shape}\PY{+w}{ }\PY{p}{\PYZob{}}
\PY{+w}{    }\PY{k+kd}{public}\PY{+w}{ }\PY{k+kt}{double}\PY{+w}{ }\PY{n+nf}{area}\PY{p}{(}\PY{p}{)}\PY{+w}{ }\PY{p}{\PYZob{}}\PY{+w}{ }\PY{k}{return}\PY{+w}{ }\PY{n}{width}\PY{+w}{ }\PY{o}{*}\PY{+w}{ }\PY{n}{height}\PY{p}{;}\PY{+w}{ }\PY{p}{\PYZcb{}}
\PY{p}{\PYZcb{}}

\PY{k+kd}{class} \PY{n+nc}{AreaCalculator}\PY{+w}{ }\PY{p}{\PYZob{}}
\PY{+w}{    }\PY{k+kt}{double}\PY{+w}{ }\PY{n+nf}{area}\PY{p}{(}\PY{n}{Shape}\PY{+w}{ }\PY{n}{shape}\PY{p}{)}\PY{+w}{ }\PY{p}{\PYZob{}}
\PY{+w}{        }\PY{k}{return}\PY{+w}{ }\PY{n}{shape}\PY{p}{.}\PY{n+na}{area}\PY{p}{(}\PY{p}{)}\PY{p}{;}
\PY{+w}{    }\PY{p}{\PYZcb{}}
\PY{p}{\PYZcb{}}

\PY{c+c1}{// Adding Triangle later requires zero changes to AreaCalculator:}
\PY{k+kd}{record} \PY{n+nc}{Triangle}\PY{p}{(}\PY{k+kt}{double}\PY{+w}{ }\PY{n}{base}\PY{p}{,}\PY{+w}{ }\PY{k+kt}{double}\PY{+w}{ }\PY{n}{height}\PY{p}{)}\PY{+w}{ }\PY{k+kd}{implements}\PY{+w}{ }\PY{n}{Shape}\PY{+w}{ }\PY{p}{\PYZob{}}
\PY{+w}{    }\PY{k+kd}{public}\PY{+w}{ }\PY{k+kt}{double}\PY{+w}{ }\PY{n+nf}{area}\PY{p}{(}\PY{p}{)}\PY{+w}{ }\PY{p}{\PYZob{}}\PY{+w}{ }\PY{k}{return}\PY{+w}{ }\PY{l+m+mf}{0.5}\PY{+w}{ }\PY{o}{*}\PY{+w}{ }\PY{n}{base}\PY{+w}{ }\PY{o}{*}\PY{+w}{ }\PY{n}{height}\PY{p}{;}\PY{+w}{ }\PY{p}{\PYZcb{}}
\PY{p}{\PYZcb{}}
\end{Verbatim}
\end{codebox}

Java 21’s sealed interfaces give you the best of both worlds when you \emph{do} want the compiler to force you to handle every case (e.g. exhaustive \code{switch}), which is useful when the set of variants is meant to be closed rather than open — see the Design Patterns section for a Visitor-style example.

\subsection[Liskov Substitution Principle (LSP)]{Liskov Substitution Principle (LSP)}
\label{h:19:liskov-substitution-principle-lsp}
Subtypes must be usable anywhere their supertype is expected, without breaking the caller’s assumptions. The classic textbook violation is \code{Square~\allowbreak{}extends~\allowbreak{}Rectangle}.

\begin{codebox}
\begin{Verbatim}[commandchars=\\\{\}]
\PY{c+c1}{// Violation: Square breaks the Rectangle contract (setWidth no longer only}
\PY{c+c1}{// changes width, it silently changes height too)}
\PY{k+kd}{class} \PY{n+nc}{Rectangle}\PY{+w}{ }\PY{p}{\PYZob{}}
\PY{+w}{    }\PY{k+kd}{protected}\PY{+w}{ }\PY{k+kt}{int}\PY{+w}{ }\PY{n}{width}\PY{p}{,}\PY{+w}{ }\PY{n}{height}\PY{p}{;}
\PY{+w}{    }\PY{k+kt}{void}\PY{+w}{ }\PY{n+nf}{setWidth}\PY{p}{(}\PY{k+kt}{int}\PY{+w}{ }\PY{n}{w}\PY{p}{)}\PY{+w}{ }\PY{p}{\PYZob{}}\PY{+w}{ }\PY{k}{this}\PY{p}{.}\PY{n+na}{width}\PY{+w}{ }\PY{o}{=}\PY{+w}{ }\PY{n}{w}\PY{p}{;}\PY{+w}{ }\PY{p}{\PYZcb{}}
\PY{+w}{    }\PY{k+kt}{void}\PY{+w}{ }\PY{n+nf}{setHeight}\PY{p}{(}\PY{k+kt}{int}\PY{+w}{ }\PY{n}{h}\PY{p}{)}\PY{+w}{ }\PY{p}{\PYZob{}}\PY{+w}{ }\PY{k}{this}\PY{p}{.}\PY{n+na}{height}\PY{+w}{ }\PY{o}{=}\PY{+w}{ }\PY{n}{h}\PY{p}{;}\PY{+w}{ }\PY{p}{\PYZcb{}}
\PY{+w}{    }\PY{k+kt}{int}\PY{+w}{ }\PY{n+nf}{area}\PY{p}{(}\PY{p}{)}\PY{+w}{ }\PY{p}{\PYZob{}}\PY{+w}{ }\PY{k}{return}\PY{+w}{ }\PY{n}{width}\PY{+w}{ }\PY{o}{*}\PY{+w}{ }\PY{n}{height}\PY{p}{;}\PY{+w}{ }\PY{p}{\PYZcb{}}
\PY{p}{\PYZcb{}}

\PY{k+kd}{class} \PY{n+nc}{Square}\PY{+w}{ }\PY{k+kd}{extends}\PY{+w}{ }\PY{n}{Rectangle}\PY{+w}{ }\PY{p}{\PYZob{}}
\PY{+w}{    }\PY{n+nd}{@Override}\PY{+w}{ }\PY{k+kt}{void}\PY{+w}{ }\PY{n+nf}{setWidth}\PY{p}{(}\PY{k+kt}{int}\PY{+w}{ }\PY{n}{w}\PY{p}{)}\PY{+w}{ }\PY{p}{\PYZob{}}\PY{+w}{ }\PY{n}{width}\PY{+w}{ }\PY{o}{=}\PY{+w}{ }\PY{n}{w}\PY{p}{;}\PY{+w}{ }\PY{n}{height}\PY{+w}{ }\PY{o}{=}\PY{+w}{ }\PY{n}{w}\PY{p}{;}\PY{+w}{ }\PY{p}{\PYZcb{}}
\PY{+w}{    }\PY{n+nd}{@Override}\PY{+w}{ }\PY{k+kt}{void}\PY{+w}{ }\PY{n+nf}{setHeight}\PY{p}{(}\PY{k+kt}{int}\PY{+w}{ }\PY{n}{h}\PY{p}{)}\PY{+w}{ }\PY{p}{\PYZob{}}\PY{+w}{ }\PY{n}{width}\PY{+w}{ }\PY{o}{=}\PY{+w}{ }\PY{n}{h}\PY{p}{;}\PY{+w}{ }\PY{n}{height}\PY{+w}{ }\PY{o}{=}\PY{+w}{ }\PY{n}{h}\PY{p}{;}\PY{+w}{ }\PY{p}{\PYZcb{}}
\PY{p}{\PYZcb{}}

\PY{k+kt}{void}\PY{+w}{ }\PY{n+nf}{resize}\PY{p}{(}\PY{n}{Rectangle}\PY{+w}{ }\PY{n}{r}\PY{p}{)}\PY{+w}{ }\PY{p}{\PYZob{}}
\PY{+w}{    }\PY{n}{r}\PY{p}{.}\PY{n+na}{setWidth}\PY{p}{(}\PY{l+m+mi}{5}\PY{p}{)}\PY{p}{;}
\PY{+w}{    }\PY{n}{r}\PY{p}{.}\PY{n+na}{setHeight}\PY{p}{(}\PY{l+m+mi}{10}\PY{p}{)}\PY{p}{;}
\PY{+w}{    }\PY{k}{assert}\PY{+w}{ }\PY{n}{r}\PY{p}{.}\PY{n+na}{area}\PY{p}{(}\PY{p}{)}\PY{+w}{ }\PY{o}{=}\PY{o}{=}\PY{+w}{ }\PY{l+m+mi}{50}\PY{p}{;}\PY{+w}{ }\PY{c+c1}{// fails for Square! area() is 100}
\PY{p}{\PYZcb{}}
\end{Verbatim}
\end{codebox}

\begin{codebox}
\begin{Verbatim}[commandchars=\\\{\}]
\PY{c+c1}{// Fixed: don\PYZsq{}t force an is\PYZhy{}a relationship that doesn\PYZsq{}t hold behaviorally.}
\PY{c+c1}{// Model the shared behavior with a common interface instead of inheritance.}
\PY{k+kd}{interface} \PY{n+nc}{Shape}\PY{+w}{ }\PY{p}{\PYZob{}}
\PY{+w}{    }\PY{k+kt}{int}\PY{+w}{ }\PY{n+nf}{area}\PY{p}{(}\PY{p}{)}\PY{p}{;}
\PY{p}{\PYZcb{}}

\PY{k+kd}{final}\PY{+w}{ }\PY{k+kd}{class} \PY{n+nc}{Rectangle}\PY{+w}{ }\PY{k+kd}{implements}\PY{+w}{ }\PY{n}{Shape}\PY{+w}{ }\PY{p}{\PYZob{}}
\PY{+w}{    }\PY{k+kd}{private}\PY{+w}{ }\PY{k+kd}{final}\PY{+w}{ }\PY{k+kt}{int}\PY{+w}{ }\PY{n}{width}\PY{p}{,}\PY{+w}{ }\PY{n}{height}\PY{p}{;}
\PY{+w}{    }\PY{n}{Rectangle}\PY{p}{(}\PY{k+kt}{int}\PY{+w}{ }\PY{n}{width}\PY{p}{,}\PY{+w}{ }\PY{k+kt}{int}\PY{+w}{ }\PY{n}{height}\PY{p}{)}\PY{+w}{ }\PY{p}{\PYZob{}}\PY{+w}{ }\PY{k}{this}\PY{p}{.}\PY{n+na}{width}\PY{+w}{ }\PY{o}{=}\PY{+w}{ }\PY{n}{width}\PY{p}{;}\PY{+w}{ }\PY{k}{this}\PY{p}{.}\PY{n+na}{height}\PY{+w}{ }\PY{o}{=}\PY{+w}{ }\PY{n}{height}\PY{p}{;}\PY{+w}{ }\PY{p}{\PYZcb{}}
\PY{+w}{    }\PY{k+kd}{public}\PY{+w}{ }\PY{k+kt}{int}\PY{+w}{ }\PY{n+nf}{area}\PY{p}{(}\PY{p}{)}\PY{+w}{ }\PY{p}{\PYZob{}}\PY{+w}{ }\PY{k}{return}\PY{+w}{ }\PY{n}{width}\PY{+w}{ }\PY{o}{*}\PY{+w}{ }\PY{n}{height}\PY{p}{;}\PY{+w}{ }\PY{p}{\PYZcb{}}
\PY{p}{\PYZcb{}}

\PY{k+kd}{final}\PY{+w}{ }\PY{k+kd}{class} \PY{n+nc}{Square}\PY{+w}{ }\PY{k+kd}{implements}\PY{+w}{ }\PY{n}{Shape}\PY{+w}{ }\PY{p}{\PYZob{}}
\PY{+w}{    }\PY{k+kd}{private}\PY{+w}{ }\PY{k+kd}{final}\PY{+w}{ }\PY{k+kt}{int}\PY{+w}{ }\PY{n}{side}\PY{p}{;}
\PY{+w}{    }\PY{n}{Square}\PY{p}{(}\PY{k+kt}{int}\PY{+w}{ }\PY{n}{side}\PY{p}{)}\PY{+w}{ }\PY{p}{\PYZob{}}\PY{+w}{ }\PY{k}{this}\PY{p}{.}\PY{n+na}{side}\PY{+w}{ }\PY{o}{=}\PY{+w}{ }\PY{n}{side}\PY{p}{;}\PY{+w}{ }\PY{p}{\PYZcb{}}
\PY{+w}{    }\PY{k+kd}{public}\PY{+w}{ }\PY{k+kt}{int}\PY{+w}{ }\PY{n+nf}{area}\PY{p}{(}\PY{p}{)}\PY{+w}{ }\PY{p}{\PYZob{}}\PY{+w}{ }\PY{k}{return}\PY{+w}{ }\PY{n}{side}\PY{+w}{ }\PY{o}{*}\PY{+w}{ }\PY{n}{side}\PY{p}{;}\PY{+w}{ }\PY{p}{\PYZcb{}}
\PY{p}{\PYZcb{}}
\end{Verbatim}
\end{codebox}

\code{Square} is no longer forced to honor a mutable-rectangle contract it cannot actually satisfy. If you must keep an inheritance hierarchy, the rule is: a subclass must not weaken postconditions, strengthen preconditions, or throw new exceptions the caller of the base type does not expect.

\subsection[Interface Segregation Principle (ISP)]{Interface Segregation Principle (ISP)}
\label{h:19:interface-segregation-principle-isp}
Clients should not be forced to depend on methods they do not use. Large, “fat” interfaces force implementers to write dummy/unsupported methods.

\begin{codebox}
\begin{Verbatim}[commandchars=\\\{\}]
\PY{c+c1}{// Violation: a read\PYZhy{}only report can\PYZsq{}t sanely implement print()}
\PY{k+kd}{interface} \PY{n+nc}{Machine}\PY{+w}{ }\PY{p}{\PYZob{}}
\PY{+w}{    }\PY{k+kt}{void}\PY{+w}{ }\PY{n+nf}{print}\PY{p}{(}\PY{n}{Document}\PY{+w}{ }\PY{n}{d}\PY{p}{)}\PY{p}{;}
\PY{+w}{    }\PY{k+kt}{void}\PY{+w}{ }\PY{n+nf}{scan}\PY{p}{(}\PY{n}{Document}\PY{+w}{ }\PY{n}{d}\PY{p}{)}\PY{p}{;}
\PY{+w}{    }\PY{k+kt}{void}\PY{+w}{ }\PY{n+nf}{fax}\PY{p}{(}\PY{n}{Document}\PY{+w}{ }\PY{n}{d}\PY{p}{)}\PY{p}{;}
\PY{p}{\PYZcb{}}

\PY{k+kd}{class} \PY{n+nc}{OldPrinter}\PY{+w}{ }\PY{k+kd}{implements}\PY{+w}{ }\PY{n}{Machine}\PY{+w}{ }\PY{p}{\PYZob{}}
\PY{+w}{    }\PY{k+kd}{public}\PY{+w}{ }\PY{k+kt}{void}\PY{+w}{ }\PY{n+nf}{print}\PY{p}{(}\PY{n}{Document}\PY{+w}{ }\PY{n}{d}\PY{p}{)}\PY{+w}{ }\PY{p}{\PYZob{}}\PY{+w}{ }\PY{c+cm}{/* ok */}\PY{+w}{ }\PY{p}{\PYZcb{}}
\PY{+w}{    }\PY{k+kd}{public}\PY{+w}{ }\PY{k+kt}{void}\PY{+w}{ }\PY{n+nf}{scan}\PY{p}{(}\PY{n}{Document}\PY{+w}{ }\PY{n}{d}\PY{p}{)}\PY{+w}{ }\PY{p}{\PYZob{}}\PY{+w}{ }\PY{k}{throw}\PY{+w}{ }\PY{k}{new}\PY{+w}{ }\PY{n}{UnsupportedOperationException}\PY{p}{(}\PY{p}{)}\PY{p}{;}\PY{+w}{ }\PY{p}{\PYZcb{}}
\PY{+w}{    }\PY{k+kd}{public}\PY{+w}{ }\PY{k+kt}{void}\PY{+w}{ }\PY{n+nf}{fax}\PY{p}{(}\PY{n}{Document}\PY{+w}{ }\PY{n}{d}\PY{p}{)}\PY{+w}{ }\PY{p}{\PYZob{}}\PY{+w}{ }\PY{k}{throw}\PY{+w}{ }\PY{k}{new}\PY{+w}{ }\PY{n}{UnsupportedOperationException}\PY{p}{(}\PY{p}{)}\PY{p}{;}\PY{+w}{ }\PY{p}{\PYZcb{}}
\PY{p}{\PYZcb{}}
\end{Verbatim}
\end{codebox}

\begin{codebox}
\begin{Verbatim}[commandchars=\\\{\}]
\PY{c+c1}{// Fixed: split into small, focused interfaces}
\PY{k+kd}{interface} \PY{n+nc}{Printer}\PY{+w}{ }\PY{p}{\PYZob{}}\PY{+w}{ }\PY{k+kt}{void}\PY{+w}{ }\PY{n+nf}{print}\PY{p}{(}\PY{n}{Document}\PY{+w}{ }\PY{n}{d}\PY{p}{)}\PY{p}{;}\PY{+w}{ }\PY{p}{\PYZcb{}}
\PY{k+kd}{interface} \PY{n+nc}{Scanner}\PY{+w}{ }\PY{p}{\PYZob{}}\PY{+w}{ }\PY{k+kt}{void}\PY{+w}{ }\PY{n+nf}{scan}\PY{p}{(}\PY{n}{Document}\PY{+w}{ }\PY{n}{d}\PY{p}{)}\PY{p}{;}\PY{+w}{ }\PY{p}{\PYZcb{}}
\PY{k+kd}{interface} \PY{n+nc}{Fax}\PY{+w}{ }\PY{p}{\PYZob{}}\PY{+w}{ }\PY{k+kt}{void}\PY{+w}{ }\PY{n+nf}{fax}\PY{p}{(}\PY{n}{Document}\PY{+w}{ }\PY{n}{d}\PY{p}{)}\PY{p}{;}\PY{+w}{ }\PY{p}{\PYZcb{}}

\PY{k+kd}{class} \PY{n+nc}{OldPrinter}\PY{+w}{ }\PY{k+kd}{implements}\PY{+w}{ }\PY{n}{Printer}\PY{+w}{ }\PY{p}{\PYZob{}}
\PY{+w}{    }\PY{k+kd}{public}\PY{+w}{ }\PY{k+kt}{void}\PY{+w}{ }\PY{n+nf}{print}\PY{p}{(}\PY{n}{Document}\PY{+w}{ }\PY{n}{d}\PY{p}{)}\PY{+w}{ }\PY{p}{\PYZob{}}\PY{+w}{ }\PY{c+cm}{/* ok */}\PY{+w}{ }\PY{p}{\PYZcb{}}
\PY{p}{\PYZcb{}}

\PY{k+kd}{class} \PY{n+nc}{AllInOnePrinter}\PY{+w}{ }\PY{k+kd}{implements}\PY{+w}{ }\PY{n}{Printer}\PY{p}{,}\PY{+w}{ }\PY{n}{Scanner}\PY{p}{,}\PY{+w}{ }\PY{n}{Fax}\PY{+w}{ }\PY{p}{\PYZob{}}
\PY{+w}{    }\PY{k+kd}{public}\PY{+w}{ }\PY{k+kt}{void}\PY{+w}{ }\PY{n+nf}{print}\PY{p}{(}\PY{n}{Document}\PY{+w}{ }\PY{n}{d}\PY{p}{)}\PY{+w}{ }\PY{p}{\PYZob{}}\PY{+w}{ }\PY{p}{.}\PY{p}{.}\PY{p}{.}\PY{+w}{ }\PY{p}{\PYZcb{}}
\PY{+w}{    }\PY{k+kd}{public}\PY{+w}{ }\PY{k+kt}{void}\PY{+w}{ }\PY{n+nf}{scan}\PY{p}{(}\PY{n}{Document}\PY{+w}{ }\PY{n}{d}\PY{p}{)}\PY{+w}{ }\PY{p}{\PYZob{}}\PY{+w}{ }\PY{p}{.}\PY{p}{.}\PY{p}{.}\PY{+w}{ }\PY{p}{\PYZcb{}}
\PY{+w}{    }\PY{k+kd}{public}\PY{+w}{ }\PY{k+kt}{void}\PY{+w}{ }\PY{n+nf}{fax}\PY{p}{(}\PY{n}{Document}\PY{+w}{ }\PY{n}{d}\PY{p}{)}\PY{+w}{ }\PY{p}{\PYZob{}}\PY{+w}{ }\PY{p}{.}\PY{p}{.}\PY{p}{.}\PY{+w}{ }\PY{p}{\PYZcb{}}
\PY{p}{\PYZcb{}}
\end{Verbatim}
\end{codebox}

\code{OldPrinter} now implements exactly what it can support. No dummy methods, no \code{Unsupported\allowbreak{}Operation\allowbreak{}Exception} traps for callers.

\subsection[Dependency Inversion Principle (DIP)]{Dependency Inversion Principle (DIP)}
\label{h:19:dependency-inversion-principle-dip}
High-level code should depend on abstractions, not on concrete low-level implementations. This is what makes classes testable and swappable.

\begin{codebox}
\begin{Verbatim}[commandchars=\\\{\}]
\PY{c+c1}{// Violation: OrderService is hard\PYZhy{}wired to a concrete SMTP sender.}
\PY{c+c1}{// You cannot test it without sending real emails, and you cannot swap}
\PY{c+c1}{// providers without editing OrderService.}
\PY{k+kd}{class} \PY{n+nc}{SmtpEmailSender}\PY{+w}{ }\PY{p}{\PYZob{}}
\PY{+w}{    }\PY{k+kt}{void}\PY{+w}{ }\PY{n+nf}{send}\PY{p}{(}\PY{n}{String}\PY{+w}{ }\PY{n}{to}\PY{p}{,}\PY{+w}{ }\PY{n}{String}\PY{+w}{ }\PY{n}{body}\PY{p}{)}\PY{+w}{ }\PY{p}{\PYZob{}}\PY{+w}{ }\PY{c+cm}{/* talks to an SMTP server */}\PY{+w}{ }\PY{p}{\PYZcb{}}
\PY{p}{\PYZcb{}}

\PY{k+kd}{class} \PY{n+nc}{OrderService}\PY{+w}{ }\PY{p}{\PYZob{}}
\PY{+w}{    }\PY{k+kd}{private}\PY{+w}{ }\PY{k+kd}{final}\PY{+w}{ }\PY{n}{SmtpEmailSender}\PY{+w}{ }\PY{n}{emailSender}\PY{+w}{ }\PY{o}{=}\PY{+w}{ }\PY{k}{new}\PY{+w}{ }\PY{n}{SmtpEmailSender}\PY{p}{(}\PY{p}{)}\PY{p}{;}

\PY{+w}{    }\PY{k+kt}{void}\PY{+w}{ }\PY{n+nf}{placeOrder}\PY{p}{(}\PY{n}{Order}\PY{+w}{ }\PY{n}{order}\PY{p}{)}\PY{+w}{ }\PY{p}{\PYZob{}}
\PY{+w}{        }\PY{c+c1}{// ... business logic ...}
\PY{+w}{        }\PY{n}{emailSender}\PY{p}{.}\PY{n+na}{send}\PY{p}{(}\PY{n}{order}\PY{p}{.}\PY{n+na}{getCustomerEmail}\PY{p}{(}\PY{p}{)}\PY{p}{,}\PY{+w}{ }\PY{l+s}{\PYZdq{}}\PY{l+s}{Order confirmed}\PY{l+s}{\PYZdq{}}\PY{p}{)}\PY{p}{;}
\PY{+w}{    }\PY{p}{\PYZcb{}}
\PY{p}{\PYZcb{}}
\end{Verbatim}
\end{codebox}

\begin{codebox}
\begin{Verbatim}[commandchars=\\\{\}]
\PY{c+c1}{// Fixed: OrderService depends on an interface; the concrete}
\PY{c+c1}{// implementation is injected from outside (constructor injection).}
\PY{k+kd}{interface} \PY{n+nc}{NotificationSender}\PY{+w}{ }\PY{p}{\PYZob{}}
\PY{+w}{    }\PY{k+kt}{void}\PY{+w}{ }\PY{n+nf}{send}\PY{p}{(}\PY{n}{String}\PY{+w}{ }\PY{n}{to}\PY{p}{,}\PY{+w}{ }\PY{n}{String}\PY{+w}{ }\PY{n}{body}\PY{p}{)}\PY{p}{;}
\PY{p}{\PYZcb{}}

\PY{k+kd}{class} \PY{n+nc}{SmtpEmailSender}\PY{+w}{ }\PY{k+kd}{implements}\PY{+w}{ }\PY{n}{NotificationSender}\PY{+w}{ }\PY{p}{\PYZob{}}
\PY{+w}{    }\PY{k+kd}{public}\PY{+w}{ }\PY{k+kt}{void}\PY{+w}{ }\PY{n+nf}{send}\PY{p}{(}\PY{n}{String}\PY{+w}{ }\PY{n}{to}\PY{p}{,}\PY{+w}{ }\PY{n}{String}\PY{+w}{ }\PY{n}{body}\PY{p}{)}\PY{+w}{ }\PY{p}{\PYZob{}}\PY{+w}{ }\PY{c+cm}{/* SMTP */}\PY{+w}{ }\PY{p}{\PYZcb{}}
\PY{p}{\PYZcb{}}

\PY{k+kd}{class} \PY{n+nc}{SmsSender}\PY{+w}{ }\PY{k+kd}{implements}\PY{+w}{ }\PY{n}{NotificationSender}\PY{+w}{ }\PY{p}{\PYZob{}}
\PY{+w}{    }\PY{k+kd}{public}\PY{+w}{ }\PY{k+kt}{void}\PY{+w}{ }\PY{n+nf}{send}\PY{p}{(}\PY{n}{String}\PY{+w}{ }\PY{n}{to}\PY{p}{,}\PY{+w}{ }\PY{n}{String}\PY{+w}{ }\PY{n}{body}\PY{p}{)}\PY{+w}{ }\PY{p}{\PYZob{}}\PY{+w}{ }\PY{c+cm}{/* SMS gateway */}\PY{+w}{ }\PY{p}{\PYZcb{}}
\PY{p}{\PYZcb{}}

\PY{k+kd}{class} \PY{n+nc}{OrderService}\PY{+w}{ }\PY{p}{\PYZob{}}
\PY{+w}{    }\PY{k+kd}{private}\PY{+w}{ }\PY{k+kd}{final}\PY{+w}{ }\PY{n}{NotificationSender}\PY{+w}{ }\PY{n}{notificationSender}\PY{p}{;}

\PY{+w}{    }\PY{n}{OrderService}\PY{p}{(}\PY{n}{NotificationSender}\PY{+w}{ }\PY{n}{notificationSender}\PY{p}{)}\PY{+w}{ }\PY{p}{\PYZob{}}
\PY{+w}{        }\PY{k}{this}\PY{p}{.}\PY{n+na}{notificationSender}\PY{+w}{ }\PY{o}{=}\PY{+w}{ }\PY{n}{notificationSender}\PY{p}{;}
\PY{+w}{    }\PY{p}{\PYZcb{}}

\PY{+w}{    }\PY{k+kt}{void}\PY{+w}{ }\PY{n+nf}{placeOrder}\PY{p}{(}\PY{n}{Order}\PY{+w}{ }\PY{n}{order}\PY{p}{)}\PY{+w}{ }\PY{p}{\PYZob{}}
\PY{+w}{        }\PY{c+c1}{// ... business logic ...}
\PY{+w}{        }\PY{n}{notificationSender}\PY{p}{.}\PY{n+na}{send}\PY{p}{(}\PY{n}{order}\PY{p}{.}\PY{n+na}{getCustomerEmail}\PY{p}{(}\PY{p}{)}\PY{p}{,}\PY{+w}{ }\PY{l+s}{\PYZdq{}}\PY{l+s}{Order confirmed}\PY{l+s}{\PYZdq{}}\PY{p}{)}\PY{p}{;}
\PY{+w}{    }\PY{p}{\PYZcb{}}
\PY{p}{\PYZcb{}}

\PY{c+c1}{// In tests:}
\PY{n}{OrderService}\PY{+w}{ }\PY{n}{service}\PY{+w}{ }\PY{o}{=}\PY{+w}{ }\PY{k}{new}\PY{+w}{ }\PY{n}{OrderService}\PY{p}{(}\PY{k}{new}\PY{+w}{ }\PY{n}{NotificationSender}\PY{p}{(}\PY{p}{)}\PY{+w}{ }\PY{p}{\PYZob{}}
\PY{+w}{    }\PY{k+kd}{public}\PY{+w}{ }\PY{k+kt}{void}\PY{+w}{ }\PY{n+nf}{send}\PY{p}{(}\PY{n}{String}\PY{+w}{ }\PY{n}{to}\PY{p}{,}\PY{+w}{ }\PY{n}{String}\PY{+w}{ }\PY{n}{body}\PY{p}{)}\PY{+w}{ }\PY{p}{\PYZob{}}\PY{+w}{ }\PY{c+cm}{/* record call, assert on it */}\PY{+w}{ }\PY{p}{\PYZcb{}}
\PY{p}{\PYZcb{}}\PY{p}{)}\PY{p}{;}
\end{Verbatim}
\end{codebox}

\code{OrderService} no longer knows or cares whether emails are sent over SMTP, SMS, or a mock in a test. This is the foundation dependency-injection frameworks (Spring, Guice, CDI) are built on.

\section[Effective Java Best Practices]{Effective Java Best Practices}
\label{h:19:effective-java-best-practices}
Joshua Bloch’s \emph{Effective Java} is the most-quoted Java book in code review. Below are roughly twenty of its highest-value items, each with a one-paragraph explanation and a short example.

\subsection[1. Prefer static factory methods over constructors]{1. Prefer static factory methods over constructors}
\label{h:19:1-prefer-static-factory-methods-over-constructors}
A static factory method can have a descriptive name, is not required to create a new object every call, and can return a subtype. A constructor always has the same name as the class and always returns exactly that type.

\begin{codebox}
\begin{Verbatim}[commandchars=\\\{\}]
\PY{k+kd}{class} \PY{n+nc}{Connection}\PY{+w}{ }\PY{p}{\PYZob{}}
\PY{+w}{    }\PY{k+kd}{private}\PY{+w}{ }\PY{n+nf}{Connection}\PY{p}{(}\PY{p}{)}\PY{+w}{ }\PY{p}{\PYZob{}}\PY{p}{\PYZcb{}}

\PY{+w}{    }\PY{c+c1}{// Name explains intent; a constructor could not do this.}
\PY{+w}{    }\PY{k+kd}{public}\PY{+w}{ }\PY{k+kd}{static}\PY{+w}{ }\PY{n}{Connection}\PY{+w}{ }\PY{n+nf}{newTcpConnection}\PY{p}{(}\PY{n}{String}\PY{+w}{ }\PY{n}{host}\PY{p}{,}\PY{+w}{ }\PY{k+kt}{int}\PY{+w}{ }\PY{n}{port}\PY{p}{)}\PY{+w}{ }\PY{p}{\PYZob{}}\PY{+w}{ }\PY{p}{.}\PY{p}{.}\PY{p}{.}\PY{+w}{ }\PY{k}{return}\PY{+w}{ }\PY{k}{new}\PY{+w}{ }\PY{n}{Connection}\PY{p}{(}\PY{p}{)}\PY{p}{;}\PY{+w}{ }\PY{p}{\PYZcb{}}
\PY{+w}{    }\PY{k+kd}{public}\PY{+w}{ }\PY{k+kd}{static}\PY{+w}{ }\PY{n}{Connection}\PY{+w}{ }\PY{n+nf}{newInMemoryConnection}\PY{p}{(}\PY{p}{)}\PY{+w}{ }\PY{p}{\PYZob{}}\PY{+w}{ }\PY{p}{.}\PY{p}{.}\PY{p}{.}\PY{+w}{ }\PY{k}{return}\PY{+w}{ }\PY{k}{new}\PY{+w}{ }\PY{n}{Connection}\PY{p}{(}\PY{p}{)}\PY{p}{;}\PY{+w}{ }\PY{p}{\PYZcb{}}
\PY{p}{\PYZcb{}}
\end{Verbatim}
\end{codebox}

\subsection[2. Use a builder when a constructor would need many parameters]{2. Use a builder when a constructor would need many parameters}
\label{h:19:2-use-a-builder-when-a-constructor-would-need-many-parameters}
Telescoping constructors (many overloads with increasing parameter counts) are error-prone and hard to read at the call site. A builder makes each parameter self-documenting.

\begin{codebox}
\begin{Verbatim}[commandchars=\\\{\}]
\PY{k+kd}{class} \PY{n+nc}{Pizza}\PY{+w}{ }\PY{p}{\PYZob{}}
\PY{+w}{    }\PY{k+kd}{private}\PY{+w}{ }\PY{k+kd}{final}\PY{+w}{ }\PY{k+kt}{int}\PY{+w}{ }\PY{n}{size}\PY{p}{;}
\PY{+w}{    }\PY{k+kd}{private}\PY{+w}{ }\PY{k+kd}{final}\PY{+w}{ }\PY{k+kt}{boolean}\PY{+w}{ }\PY{n}{cheese}\PY{p}{;}
\PY{+w}{    }\PY{k+kd}{private}\PY{+w}{ }\PY{k+kd}{final}\PY{+w}{ }\PY{k+kt}{boolean}\PY{+w}{ }\PY{n}{pepperoni}\PY{p}{;}

\PY{+w}{    }\PY{k+kd}{private}\PY{+w}{ }\PY{n+nf}{Pizza}\PY{p}{(}\PY{n}{Builder}\PY{+w}{ }\PY{n}{b}\PY{p}{)}\PY{+w}{ }\PY{p}{\PYZob{}}
\PY{+w}{        }\PY{k}{this}\PY{p}{.}\PY{n+na}{size}\PY{+w}{ }\PY{o}{=}\PY{+w}{ }\PY{n}{b}\PY{p}{.}\PY{n+na}{size}\PY{p}{;}\PY{+w}{ }\PY{k}{this}\PY{p}{.}\PY{n+na}{cheese}\PY{+w}{ }\PY{o}{=}\PY{+w}{ }\PY{n}{b}\PY{p}{.}\PY{n+na}{cheese}\PY{p}{;}\PY{+w}{ }\PY{k}{this}\PY{p}{.}\PY{n+na}{pepperoni}\PY{+w}{ }\PY{o}{=}\PY{+w}{ }\PY{n}{b}\PY{p}{.}\PY{n+na}{pepperoni}\PY{p}{;}
\PY{+w}{    }\PY{p}{\PYZcb{}}

\PY{+w}{    }\PY{k+kd}{static}\PY{+w}{ }\PY{k+kd}{class} \PY{n+nc}{Builder}\PY{+w}{ }\PY{p}{\PYZob{}}
\PY{+w}{        }\PY{k+kd}{private}\PY{+w}{ }\PY{k+kd}{final}\PY{+w}{ }\PY{k+kt}{int}\PY{+w}{ }\PY{n}{size}\PY{p}{;}
\PY{+w}{        }\PY{k+kd}{private}\PY{+w}{ }\PY{k+kt}{boolean}\PY{+w}{ }\PY{n}{cheese}\PY{p}{;}
\PY{+w}{        }\PY{k+kd}{private}\PY{+w}{ }\PY{k+kt}{boolean}\PY{+w}{ }\PY{n}{pepperoni}\PY{p}{;}

\PY{+w}{        }\PY{n}{Builder}\PY{p}{(}\PY{k+kt}{int}\PY{+w}{ }\PY{n}{size}\PY{p}{)}\PY{+w}{ }\PY{p}{\PYZob{}}\PY{+w}{ }\PY{k}{this}\PY{p}{.}\PY{n+na}{size}\PY{+w}{ }\PY{o}{=}\PY{+w}{ }\PY{n}{size}\PY{p}{;}\PY{+w}{ }\PY{p}{\PYZcb{}}
\PY{+w}{        }\PY{n}{Builder}\PY{+w}{ }\PY{n+nf}{cheese}\PY{p}{(}\PY{k+kt}{boolean}\PY{+w}{ }\PY{n}{value}\PY{p}{)}\PY{+w}{ }\PY{p}{\PYZob{}}\PY{+w}{ }\PY{k}{this}\PY{p}{.}\PY{n+na}{cheese}\PY{+w}{ }\PY{o}{=}\PY{+w}{ }\PY{n}{value}\PY{p}{;}\PY{+w}{ }\PY{k}{return}\PY{+w}{ }\PY{k}{this}\PY{p}{;}\PY{+w}{ }\PY{p}{\PYZcb{}}
\PY{+w}{        }\PY{n}{Builder}\PY{+w}{ }\PY{n+nf}{pepperoni}\PY{p}{(}\PY{k+kt}{boolean}\PY{+w}{ }\PY{n}{value}\PY{p}{)}\PY{+w}{ }\PY{p}{\PYZob{}}\PY{+w}{ }\PY{k}{this}\PY{p}{.}\PY{n+na}{pepperoni}\PY{+w}{ }\PY{o}{=}\PY{+w}{ }\PY{n}{value}\PY{p}{;}\PY{+w}{ }\PY{k}{return}\PY{+w}{ }\PY{k}{this}\PY{p}{;}\PY{+w}{ }\PY{p}{\PYZcb{}}
\PY{+w}{        }\PY{n}{Pizza}\PY{+w}{ }\PY{n+nf}{build}\PY{p}{(}\PY{p}{)}\PY{+w}{ }\PY{p}{\PYZob{}}\PY{+w}{ }\PY{k}{return}\PY{+w}{ }\PY{k}{new}\PY{+w}{ }\PY{n}{Pizza}\PY{p}{(}\PY{k}{this}\PY{p}{)}\PY{p}{;}\PY{+w}{ }\PY{p}{\PYZcb{}}
\PY{+w}{    }\PY{p}{\PYZcb{}}
\PY{p}{\PYZcb{}}

\PY{n}{Pizza}\PY{+w}{ }\PY{n}{pizza}\PY{+w}{ }\PY{o}{=}\PY{+w}{ }\PY{k}{new}\PY{+w}{ }\PY{n}{Pizza}\PY{p}{.}\PY{n+na}{Builder}\PY{p}{(}\PY{l+m+mi}{12}\PY{p}{)}\PY{p}{.}\PY{n+na}{cheese}\PY{p}{(}\PY{k+kc}{true}\PY{p}{)}\PY{p}{.}\PY{n+na}{pepperoni}\PY{p}{(}\PY{k+kc}{true}\PY{p}{)}\PY{p}{.}\PY{n+na}{build}\PY{p}{(}\PY{p}{)}\PY{p}{;}
\end{Verbatim}
\end{codebox}

\subsection[3. Use a single-element enum for singletons]{3. Use a single-element enum for singletons}
\label{h:19:3-use-a-single-element-enum-for-singletons}
An enum singleton gets serialization safety and reflection-attack resistance for free, which a hand-rolled singleton does not.

\begin{codebox}
\begin{Verbatim}[commandchars=\\\{\}]
\PY{k+kd}{enum}\PY{+w}{ }\PY{n}{AppConfig}\PY{+w}{ }\PY{p}{\PYZob{}}
\PY{+w}{    }\PY{n}{INSTANCE}\PY{p}{;}
\PY{+w}{    }\PY{k+kd}{private}\PY{+w}{ }\PY{k+kd}{final}\PY{+w}{ }\PY{n}{Properties}\PY{+w}{ }\PY{n}{props}\PY{+w}{ }\PY{o}{=}\PY{+w}{ }\PY{n}{loadProperties}\PY{p}{(}\PY{p}{)}\PY{p}{;}
\PY{+w}{    }\PY{n}{Properties}\PY{+w}{ }\PY{n+nf}{properties}\PY{p}{(}\PY{p}{)}\PY{+w}{ }\PY{p}{\PYZob{}}\PY{+w}{ }\PY{k}{return}\PY{+w}{ }\PY{n}{props}\PY{p}{;}\PY{+w}{ }\PY{p}{\PYZcb{}}
\PY{p}{\PYZcb{}}

\PY{n}{AppConfig}\PY{p}{.}\PY{n+na}{INSTANCE}\PY{p}{.}\PY{n+na}{properties}\PY{p}{(}\PY{p}{)}\PY{p}{;}
\end{Verbatim}
\end{codebox}

\subsection[4. Prefer dependency injection over hardwiring dependencies]{4. Prefer dependency injection over hardwiring dependencies}
\label{h:19:4-prefer-dependency-injection-over-hardwiring-dependencies}
Hardwiring a concrete dependency (\code{new~\allowbreak{}SmtpEmailSender()}) inside a class, as shown in the DIP example above, makes that class impossible to unit test in isolation. Pass dependencies into the constructor instead.

\begin{codebox}
\begin{Verbatim}[commandchars=\\\{\}]
\PY{c+c1}{// Hardwired — cannot substitute a fake in tests}
\PY{k+kd}{class} \PY{n+nc}{Report}\PY{+w}{ }\PY{p}{\PYZob{}}\PY{+w}{ }\PY{k+kd}{private}\PY{+w}{ }\PY{k+kd}{final}\PY{+w}{ }\PY{n}{Formatter}\PY{+w}{ }\PY{n}{formatter}\PY{+w}{ }\PY{o}{=}\PY{+w}{ }\PY{k}{new}\PY{+w}{ }\PY{n}{PdfFormatter}\PY{p}{(}\PY{p}{)}\PY{p}{;}\PY{+w}{ }\PY{p}{\PYZcb{}}

\PY{c+c1}{// Injected — testable and swappable}
\PY{k+kd}{class} \PY{n+nc}{Report}\PY{+w}{ }\PY{p}{\PYZob{}}
\PY{+w}{    }\PY{k+kd}{private}\PY{+w}{ }\PY{k+kd}{final}\PY{+w}{ }\PY{n}{Formatter}\PY{+w}{ }\PY{n}{formatter}\PY{p}{;}
\PY{+w}{    }\PY{n}{Report}\PY{p}{(}\PY{n}{Formatter}\PY{+w}{ }\PY{n}{formatter}\PY{p}{)}\PY{+w}{ }\PY{p}{\PYZob{}}\PY{+w}{ }\PY{k}{this}\PY{p}{.}\PY{n+na}{formatter}\PY{+w}{ }\PY{o}{=}\PY{+w}{ }\PY{n}{formatter}\PY{p}{;}\PY{+w}{ }\PY{p}{\PYZcb{}}
\PY{p}{\PYZcb{}}
\end{Verbatim}
\end{codebox}

\subsection[5. Avoid finalizers and the old Object.finalize(); use Cleaner or try-with-resources]{5. Avoid finalizers and the old \code{Object.\allowbreak{}finalize()}; use \code{Cleaner} or try-with-resources}
\label{h:19:5-avoid-finalizers-and-the-old-objectfinalize-use-cleaner-or-try-with-resources}
Finalizers run at an unpredictable time (or never), can hurt performance, and can even resurrect objects. Try-with-resources with \code{AutoCloseable} is the correct default; \code{java.\allowbreak{}lang.\allowbreak{}ref.\allowbreak{}Cleaner} is the safety-net alternative when you need last-resort cleanup.

\begin{codebox}
\begin{Verbatim}[commandchars=\\\{\}]
\PY{k+kd}{class} \PY{n+nc}{FileHandle}\PY{+w}{ }\PY{k+kd}{implements}\PY{+w}{ }\PY{n}{AutoCloseable}\PY{+w}{ }\PY{p}{\PYZob{}}
\PY{+w}{    }\PY{k+kd}{private}\PY{+w}{ }\PY{k+kd}{final}\PY{+w}{ }\PY{n}{RandomAccessFile}\PY{+w}{ }\PY{n}{file}\PY{p}{;}
\PY{+w}{    }\PY{n}{FileHandle}\PY{p}{(}\PY{n}{String}\PY{+w}{ }\PY{n}{path}\PY{p}{)}\PY{+w}{ }\PY{k+kd}{throws}\PY{+w}{ }\PY{n}{IOException}\PY{+w}{ }\PY{p}{\PYZob{}}\PY{+w}{ }\PY{n}{file}\PY{+w}{ }\PY{o}{=}\PY{+w}{ }\PY{k}{new}\PY{+w}{ }\PY{n}{RandomAccessFile}\PY{p}{(}\PY{n}{path}\PY{p}{,}\PY{+w}{ }\PY{l+s}{\PYZdq{}}\PY{l+s}{r}\PY{l+s}{\PYZdq{}}\PY{p}{)}\PY{p}{;}\PY{+w}{ }\PY{p}{\PYZcb{}}
\PY{+w}{    }\PY{k+kd}{public}\PY{+w}{ }\PY{k+kt}{void}\PY{+w}{ }\PY{n+nf}{close}\PY{p}{(}\PY{p}{)}\PY{+w}{ }\PY{k+kd}{throws}\PY{+w}{ }\PY{n}{IOException}\PY{+w}{ }\PY{p}{\PYZob{}}\PY{+w}{ }\PY{n}{file}\PY{p}{.}\PY{n+na}{close}\PY{p}{(}\PY{p}{)}\PY{p}{;}\PY{+w}{ }\PY{p}{\PYZcb{}}
\PY{p}{\PYZcb{}}

\PY{k}{try}\PY{+w}{ }\PY{p}{(}\PY{n}{FileHandle}\PY{+w}{ }\PY{n}{handle}\PY{+w}{ }\PY{o}{=}\PY{+w}{ }\PY{k}{new}\PY{+w}{ }\PY{n}{FileHandle}\PY{p}{(}\PY{l+s}{\PYZdq{}}\PY{l+s}{data.bin}\PY{l+s}{\PYZdq{}}\PY{p}{)}\PY{p}{)}\PY{+w}{ }\PY{p}{\PYZob{}}
\PY{+w}{    }\PY{c+c1}{// use handle}
\PY{p}{\PYZcb{}}\PY{+w}{ }\PY{c+c1}{// closed deterministically, no finalizer needed}
\end{Verbatim}
\end{codebox}

\subsection[6. Obey the equals/hashCode contract]{6. Obey the \code{equals}/\code{hashCode} contract}
\label{h:19:6-obey-the-equalshashcode-contract}
If you override \code{equals}, you must override \code{hashCode} so that equal objects have equal hash codes — otherwise the object breaks in \code{HashMap}/\code{HashSet}. Since Java 16, a \code{record} generates both correctly for you.

\begin{codebox}
\begin{Verbatim}[commandchars=\\\{\}]
\PY{c+c1}{// Manual, error\PYZhy{}prone version}
\PY{k+kd}{class} \PY{n+nc}{Point}\PY{+w}{ }\PY{p}{\PYZob{}}
\PY{+w}{    }\PY{k+kd}{final}\PY{+w}{ }\PY{k+kt}{int}\PY{+w}{ }\PY{n}{x}\PY{p}{,}\PY{+w}{ }\PY{n}{y}\PY{p}{;}
\PY{+w}{    }\PY{n}{Point}\PY{p}{(}\PY{k+kt}{int}\PY{+w}{ }\PY{n}{x}\PY{p}{,}\PY{+w}{ }\PY{k+kt}{int}\PY{+w}{ }\PY{n}{y}\PY{p}{)}\PY{+w}{ }\PY{p}{\PYZob{}}\PY{+w}{ }\PY{k}{this}\PY{p}{.}\PY{n+na}{x}\PY{+w}{ }\PY{o}{=}\PY{+w}{ }\PY{n}{x}\PY{p}{;}\PY{+w}{ }\PY{k}{this}\PY{p}{.}\PY{n+na}{y}\PY{+w}{ }\PY{o}{=}\PY{+w}{ }\PY{n}{y}\PY{p}{;}\PY{+w}{ }\PY{p}{\PYZcb{}}
\PY{+w}{    }\PY{n+nd}{@Override}\PY{+w}{ }\PY{k+kd}{public}\PY{+w}{ }\PY{k+kt}{boolean}\PY{+w}{ }\PY{n+nf}{equals}\PY{p}{(}\PY{n}{Object}\PY{+w}{ }\PY{n}{o}\PY{p}{)}\PY{+w}{ }\PY{p}{\PYZob{}}
\PY{+w}{        }\PY{k}{return}\PY{+w}{ }\PY{n}{o}\PY{+w}{ }\PY{k}{instanceof}\PY{+w}{ }\PY{n}{Point}\PY{+w}{ }\PY{n}{p}\PY{+w}{ }\PY{o}{\PYZam{}}\PY{o}{\PYZam{}}\PY{+w}{ }\PY{n}{p}\PY{p}{.}\PY{n+na}{x}\PY{+w}{ }\PY{o}{=}\PY{o}{=}\PY{+w}{ }\PY{n}{x}\PY{+w}{ }\PY{o}{\PYZam{}}\PY{o}{\PYZam{}}\PY{+w}{ }\PY{n}{p}\PY{p}{.}\PY{n+na}{y}\PY{+w}{ }\PY{o}{=}\PY{o}{=}\PY{+w}{ }\PY{n}{y}\PY{p}{;}
\PY{+w}{    }\PY{p}{\PYZcb{}}
\PY{+w}{    }\PY{n+nd}{@Override}\PY{+w}{ }\PY{k+kd}{public}\PY{+w}{ }\PY{k+kt}{int}\PY{+w}{ }\PY{n+nf}{hashCode}\PY{p}{(}\PY{p}{)}\PY{+w}{ }\PY{p}{\PYZob{}}\PY{+w}{ }\PY{k}{return}\PY{+w}{ }\PY{n}{Objects}\PY{p}{.}\PY{n+na}{hash}\PY{p}{(}\PY{n}{x}\PY{p}{,}\PY{+w}{ }\PY{n}{y}\PY{p}{)}\PY{p}{;}\PY{+w}{ }\PY{p}{\PYZcb{}}
\PY{p}{\PYZcb{}}

\PY{c+c1}{// Modern version: equals/hashCode/toString generated correctly, always}
\PY{k+kd}{record} \PY{n+nc}{Point}\PY{p}{(}\PY{k+kt}{int}\PY{+w}{ }\PY{n}{x}\PY{p}{,}\PY{+w}{ }\PY{k+kt}{int}\PY{+w}{ }\PY{n}{y}\PY{p}{)}\PY{+w}{ }\PY{p}{\PYZob{}}\PY{p}{\PYZcb{}}
\end{Verbatim}
\end{codebox}

\subsection[7. Minimize mutability]{7. Minimize mutability}
\label{h:19:7-minimize-mutability}
Immutable objects are simpler to reason about, are automatically thread-safe, and can be freely shared. Make fields \code{final}, do not provide mutators, and ensure the class cannot be extended in a way that breaks invariants.

\begin{codebox}
\begin{Verbatim}[commandchars=\\\{\}]
\PY{c+c1}{// Mutable — a caller can corrupt shared state}
\PY{k+kd}{class} \PY{n+nc}{Money}\PY{+w}{ }\PY{p}{\PYZob{}}
\PY{+w}{    }\PY{k+kd}{private}\PY{+w}{ }\PY{n}{BigDecimal}\PY{+w}{ }\PY{n}{amount}\PY{p}{;}
\PY{+w}{    }\PY{k+kt}{void}\PY{+w}{ }\PY{n+nf}{setAmount}\PY{p}{(}\PY{n}{BigDecimal}\PY{+w}{ }\PY{n}{amount}\PY{p}{)}\PY{+w}{ }\PY{p}{\PYZob{}}\PY{+w}{ }\PY{k}{this}\PY{p}{.}\PY{n+na}{amount}\PY{+w}{ }\PY{o}{=}\PY{+w}{ }\PY{n}{amount}\PY{p}{;}\PY{+w}{ }\PY{p}{\PYZcb{}}
\PY{p}{\PYZcb{}}

\PY{c+c1}{// Immutable — every \PYZdq{}change\PYZdq{} returns a new instance}
\PY{k+kd}{final}\PY{+w}{ }\PY{k+kd}{class} \PY{n+nc}{Money}\PY{+w}{ }\PY{p}{\PYZob{}}
\PY{+w}{    }\PY{k+kd}{private}\PY{+w}{ }\PY{k+kd}{final}\PY{+w}{ }\PY{n}{BigDecimal}\PY{+w}{ }\PY{n}{amount}\PY{p}{;}
\PY{+w}{    }\PY{n}{Money}\PY{p}{(}\PY{n}{BigDecimal}\PY{+w}{ }\PY{n}{amount}\PY{p}{)}\PY{+w}{ }\PY{p}{\PYZob{}}\PY{+w}{ }\PY{k}{this}\PY{p}{.}\PY{n+na}{amount}\PY{+w}{ }\PY{o}{=}\PY{+w}{ }\PY{n}{amount}\PY{p}{;}\PY{+w}{ }\PY{p}{\PYZcb{}}
\PY{+w}{    }\PY{n}{Money}\PY{+w}{ }\PY{n+nf}{plus}\PY{p}{(}\PY{n}{Money}\PY{+w}{ }\PY{n}{other}\PY{p}{)}\PY{+w}{ }\PY{p}{\PYZob{}}\PY{+w}{ }\PY{k}{return}\PY{+w}{ }\PY{k}{new}\PY{+w}{ }\PY{n}{Money}\PY{p}{(}\PY{k}{this}\PY{p}{.}\PY{n+na}{amount}\PY{p}{.}\PY{n+na}{add}\PY{p}{(}\PY{n}{other}\PY{p}{.}\PY{n+na}{amount}\PY{p}{)}\PY{p}{)}\PY{p}{;}\PY{+w}{ }\PY{p}{\PYZcb{}}
\PY{p}{\PYZcb{}}
\end{Verbatim}
\end{codebox}

\subsection[8. Favor composition over inheritance]{8. Favor composition over inheritance}
\label{h:19:8-favor-composition-over-inheritance}
Inheritance couples a subclass to the internal implementation details of its superclass; a change in the superclass can silently break subclasses (“fragile base class” problem). Composition — holding a reference to another object and delegating to it — is more flexible.

\begin{codebox}
\begin{Verbatim}[commandchars=\\\{\}]
\PY{c+c1}{// Inheritance: ArrayList\PYZsq{}s internals leak through; size() may not}
\PY{c+c1}{// behave as expected if addAll() is implemented in terms of add()}
\PY{k+kd}{class} \PY{n+nc}{InstrumentedSet}\PY{o}{\PYZlt{}}\PY{n}{E}\PY{o}{\PYZgt{}}\PY{+w}{ }\PY{k+kd}{extends}\PY{+w}{ }\PY{n}{HashSet}\PY{o}{\PYZlt{}}\PY{n}{E}\PY{o}{\PYZgt{}}\PY{+w}{ }\PY{p}{\PYZob{}}
\PY{+w}{    }\PY{k+kt}{int}\PY{+w}{ }\PY{n}{addCount}\PY{+w}{ }\PY{o}{=}\PY{+w}{ }\PY{l+m+mi}{0}\PY{p}{;}
\PY{+w}{    }\PY{n+nd}{@Override}\PY{+w}{ }\PY{k+kd}{public}\PY{+w}{ }\PY{k+kt}{boolean}\PY{+w}{ }\PY{n+nf}{add}\PY{p}{(}\PY{n}{E}\PY{+w}{ }\PY{n}{e}\PY{p}{)}\PY{+w}{ }\PY{p}{\PYZob{}}\PY{+w}{ }\PY{n}{addCount}\PY{o}{+}\PY{o}{+}\PY{p}{;}\PY{+w}{ }\PY{k}{return}\PY{+w}{ }\PY{k+kd}{super}\PY{p}{.}\PY{n+na}{add}\PY{p}{(}\PY{n}{e}\PY{p}{)}\PY{p}{;}\PY{+w}{ }\PY{p}{\PYZcb{}}
\PY{+w}{    }\PY{n+nd}{@Override}\PY{+w}{ }\PY{k+kd}{public}\PY{+w}{ }\PY{k+kt}{boolean}\PY{+w}{ }\PY{n+nf}{addAll}\PY{p}{(}\PY{n}{Collection}\PY{o}{\PYZlt{}}\PY{o}{?}\PY{+w}{ }\PY{k+kd}{extends}\PY{+w}{ }\PY{n}{E}\PY{o}{\PYZgt{}}\PY{+w}{ }\PY{n}{c}\PY{p}{)}\PY{+w}{ }\PY{p}{\PYZob{}}
\PY{+w}{        }\PY{n}{addCount}\PY{+w}{ }\PY{o}{+}\PY{o}{=}\PY{+w}{ }\PY{n}{c}\PY{p}{.}\PY{n+na}{size}\PY{p}{(}\PY{p}{)}\PY{p}{;}
\PY{+w}{        }\PY{k}{return}\PY{+w}{ }\PY{k+kd}{super}\PY{p}{.}\PY{n+na}{addAll}\PY{p}{(}\PY{n}{c}\PY{p}{)}\PY{p}{;}\PY{+w}{ }\PY{c+c1}{// may double\PYZhy{}count if addAll calls add() internally}
\PY{+w}{    }\PY{p}{\PYZcb{}}
\PY{p}{\PYZcb{}}

\PY{c+c1}{// Composition: wrap a Set, forward calls, no fragile\PYZhy{}base\PYZhy{}class risk}
\PY{k+kd}{class} \PY{n+nc}{InstrumentedSet}\PY{o}{\PYZlt{}}\PY{n}{E}\PY{o}{\PYZgt{}}\PY{+w}{ }\PY{p}{\PYZob{}}
\PY{+w}{    }\PY{k+kd}{private}\PY{+w}{ }\PY{k+kd}{final}\PY{+w}{ }\PY{n}{Set}\PY{o}{\PYZlt{}}\PY{n}{E}\PY{o}{\PYZgt{}}\PY{+w}{ }\PY{n}{delegate}\PY{p}{;}
\PY{+w}{    }\PY{k+kt}{int}\PY{+w}{ }\PY{n}{addCount}\PY{+w}{ }\PY{o}{=}\PY{+w}{ }\PY{l+m+mi}{0}\PY{p}{;}
\PY{+w}{    }\PY{n}{InstrumentedSet}\PY{p}{(}\PY{n}{Set}\PY{o}{\PYZlt{}}\PY{n}{E}\PY{o}{\PYZgt{}}\PY{+w}{ }\PY{n}{delegate}\PY{p}{)}\PY{+w}{ }\PY{p}{\PYZob{}}\PY{+w}{ }\PY{k}{this}\PY{p}{.}\PY{n+na}{delegate}\PY{+w}{ }\PY{o}{=}\PY{+w}{ }\PY{n}{delegate}\PY{p}{;}\PY{+w}{ }\PY{p}{\PYZcb{}}
\PY{+w}{    }\PY{k+kt}{boolean}\PY{+w}{ }\PY{n+nf}{add}\PY{p}{(}\PY{n}{E}\PY{+w}{ }\PY{n}{e}\PY{p}{)}\PY{+w}{ }\PY{p}{\PYZob{}}\PY{+w}{ }\PY{n}{addCount}\PY{o}{+}\PY{o}{+}\PY{p}{;}\PY{+w}{ }\PY{k}{return}\PY{+w}{ }\PY{n}{delegate}\PY{p}{.}\PY{n+na}{add}\PY{p}{(}\PY{n}{e}\PY{p}{)}\PY{p}{;}\PY{+w}{ }\PY{p}{\PYZcb{}}
\PY{+w}{    }\PY{k+kt}{boolean}\PY{+w}{ }\PY{n+nf}{addAll}\PY{p}{(}\PY{n}{Collection}\PY{o}{\PYZlt{}}\PY{o}{?}\PY{+w}{ }\PY{k+kd}{extends}\PY{+w}{ }\PY{n}{E}\PY{o}{\PYZgt{}}\PY{+w}{ }\PY{n}{c}\PY{p}{)}\PY{+w}{ }\PY{p}{\PYZob{}}\PY{+w}{ }\PY{n}{addCount}\PY{+w}{ }\PY{o}{+}\PY{o}{=}\PY{+w}{ }\PY{n}{c}\PY{p}{.}\PY{n+na}{size}\PY{p}{(}\PY{p}{)}\PY{p}{;}\PY{+w}{ }\PY{k}{return}\PY{+w}{ }\PY{n}{delegate}\PY{p}{.}\PY{n+na}{addAll}\PY{p}{(}\PY{n}{c}\PY{p}{)}\PY{p}{;}\PY{+w}{ }\PY{p}{\PYZcb{}}
\PY{p}{\PYZcb{}}
\end{Verbatim}
\end{codebox}

\subsection[9. Design and document for inheritance, or prohibit it]{9. Design and document for inheritance, or prohibit it}
\label{h:19:9-design-and-document-for-inheritance-or-prohibit-it}
If a class is not explicitly designed to be extended (documented “self-use” behavior, no calls from constructors to overridable methods), mark it \code{final} or make its constructors package-private. Otherwise you invite the fragile-base-class problems from item 8.

\begin{codebox}
\begin{Verbatim}[commandchars=\\\{\}]
\PY{c+c1}{// Dangerous: superclass constructor calls an overridable method;}
\PY{c+c1}{// subclass fields are not yet initialized when overridden() runs}
\PY{k+kd}{class} \PY{n+nc}{Base}\PY{+w}{ }\PY{p}{\PYZob{}}
\PY{+w}{    }\PY{n}{Base}\PY{p}{(}\PY{p}{)}\PY{+w}{ }\PY{p}{\PYZob{}}\PY{+w}{ }\PY{n}{overridden}\PY{p}{(}\PY{p}{)}\PY{p}{;}\PY{+w}{ }\PY{p}{\PYZcb{}}
\PY{+w}{    }\PY{k+kt}{void}\PY{+w}{ }\PY{n+nf}{overridden}\PY{p}{(}\PY{p}{)}\PY{+w}{ }\PY{p}{\PYZob{}}\PY{p}{\PYZcb{}}
\PY{p}{\PYZcb{}}

\PY{c+c1}{// Safe: prevent inheritance entirely when it wasn\PYZsq{}t designed for it}
\PY{k+kd}{final}\PY{+w}{ }\PY{k+kd}{class} \PY{n+nc}{Utility}\PY{+w}{ }\PY{p}{\PYZob{}}
\PY{+w}{    }\PY{k+kd}{private}\PY{+w}{ }\PY{n+nf}{Utility}\PY{p}{(}\PY{p}{)}\PY{+w}{ }\PY{p}{\PYZob{}}\PY{p}{\PYZcb{}}
\PY{+w}{    }\PY{k+kd}{static}\PY{+w}{ }\PY{k+kt}{int}\PY{+w}{ }\PY{n+nf}{square}\PY{p}{(}\PY{k+kt}{int}\PY{+w}{ }\PY{n}{x}\PY{p}{)}\PY{+w}{ }\PY{p}{\PYZob{}}\PY{+w}{ }\PY{k}{return}\PY{+w}{ }\PY{n}{x}\PY{+w}{ }\PY{o}{*}\PY{+w}{ }\PY{n}{x}\PY{p}{;}\PY{+w}{ }\PY{p}{\PYZcb{}}
\PY{p}{\PYZcb{}}
\end{Verbatim}
\end{codebox}

\subsection[10. Prefer interfaces to abstract classes for defining types]{10. Prefer interfaces to abstract classes for defining types}
\label{h:19:10-prefer-interfaces-to-abstract-classes-for-defining-types}
A class can implement multiple interfaces but extend only one abstract class. Interfaces (with default methods) usually give you enough shared behavior without giving up multiple inheritance of type.

\begin{codebox}
\begin{Verbatim}[commandchars=\\\{\}]
\PY{k+kd}{interface} \PY{n+nc}{Flyer}\PY{+w}{ }\PY{p}{\PYZob{}}\PY{+w}{ }\PY{k}{default}\PY{+w}{ }\PY{k+kt}{void}\PY{+w}{ }\PY{n+nf}{takeOff}\PY{p}{(}\PY{p}{)}\PY{+w}{ }\PY{p}{\PYZob{}}\PY{+w}{ }\PY{n}{System}\PY{p}{.}\PY{n+na}{out}\PY{p}{.}\PY{n+na}{println}\PY{p}{(}\PY{l+s}{\PYZdq{}}\PY{l+s}{Taking off}\PY{l+s}{\PYZdq{}}\PY{p}{)}\PY{p}{;}\PY{+w}{ }\PY{p}{\PYZcb{}}\PY{+w}{ }\PY{p}{\PYZcb{}}
\PY{k+kd}{interface} \PY{n+nc}{Swimmer}\PY{+w}{ }\PY{p}{\PYZob{}}\PY{+w}{ }\PY{k}{default}\PY{+w}{ }\PY{k+kt}{void}\PY{+w}{ }\PY{n+nf}{dive}\PY{p}{(}\PY{p}{)}\PY{+w}{ }\PY{p}{\PYZob{}}\PY{+w}{ }\PY{n}{System}\PY{p}{.}\PY{n+na}{out}\PY{p}{.}\PY{n+na}{println}\PY{p}{(}\PY{l+s}{\PYZdq{}}\PY{l+s}{Diving}\PY{l+s}{\PYZdq{}}\PY{p}{)}\PY{p}{;}\PY{+w}{ }\PY{p}{\PYZcb{}}\PY{+w}{ }\PY{p}{\PYZcb{}}

\PY{c+c1}{// A Duck can be both — impossible with two abstract base classes}
\PY{k+kd}{class} \PY{n+nc}{Duck}\PY{+w}{ }\PY{k+kd}{implements}\PY{+w}{ }\PY{n}{Flyer}\PY{p}{,}\PY{+w}{ }\PY{n}{Swimmer}\PY{+w}{ }\PY{p}{\PYZob{}}\PY{p}{\PYZcb{}}
\end{Verbatim}
\end{codebox}

\subsection[11. Prefer lists to arrays]{11. Prefer lists to arrays}
\label{h:19:11-prefer-lists-to-arrays}
Arrays are covariant and not truly type-safe at runtime (\code{ArrayStoreException}), and they don’t play well with generics. Prefer \code{List\textless{}\allowbreak{}T\textgreater{}}.

\begin{codebox}
\begin{Verbatim}[commandchars=\\\{\}]
\PY{c+c1}{// Arrays: compiles, fails at runtime}
\PY{n}{Object}\PY{o}{[}\PY{o}{]}\PY{+w}{ }\PY{n}{objects}\PY{+w}{ }\PY{o}{=}\PY{+w}{ }\PY{k}{new}\PY{+w}{ }\PY{n}{String}\PY{o}{[}\PY{l+m+mi}{1}\PY{o}{]}\PY{p}{;}
\PY{n}{objects}\PY{o}{[}\PY{l+m+mi}{0}\PY{o}{]}\PY{+w}{ }\PY{o}{=}\PY{+w}{ }\PY{l+m+mi}{42}\PY{p}{;}\PY{+w}{ }\PY{c+c1}{// throws ArrayStoreException at runtime}

\PY{c+c1}{// Lists: caught at compile time instead}
\PY{n}{List}\PY{o}{\PYZlt{}}\PY{n}{String}\PY{o}{\PYZgt{}}\PY{+w}{ }\PY{n}{strings}\PY{+w}{ }\PY{o}{=}\PY{+w}{ }\PY{k}{new}\PY{+w}{ }\PY{n}{ArrayList}\PY{o}{\PYZlt{}}\PY{o}{\PYZgt{}}\PY{p}{(}\PY{p}{)}\PY{p}{;}
\PY{c+c1}{// strings.add(42); // does not compile}
\end{Verbatim}
\end{codebox}

\subsection[12. Favor generics, and use bounded wildcards to increase API flexibility]{12. Favor generics, and use bounded wildcards to increase API flexibility}
\label{h:19:12-favor-generics-and-use-bounded-wildcards-to-increase-api-flexibility}
Generics give compile-time type safety instead of runtime \code{ClassCastException}. Bounded wildcards (\code{?\allowbreak{}~\allowbreak{}extends}, \code{?\allowbreak{}~\allowbreak{}super}) let an API accept a wider range of argument types (the “PECS” rule: \textbf{P}roducer \textbf{E}xtends, \textbf{C}onsumer \textbf{S}uper).

\begin{codebox}
\begin{Verbatim}[commandchars=\\\{\}]
\PY{c+c1}{// Without wildcards: only List\PYZlt{}Number\PYZgt{} works, List\PYZlt{}Integer\PYZgt{} is rejected}
\PY{k+kt}{void}\PY{+w}{ }\PY{n+nf}{printAll}\PY{p}{(}\PY{n}{List}\PY{o}{\PYZlt{}}\PY{n}{Number}\PY{o}{\PYZgt{}}\PY{+w}{ }\PY{n}{list}\PY{p}{)}\PY{+w}{ }\PY{p}{\PYZob{}}\PY{+w}{ }\PY{n}{list}\PY{p}{.}\PY{n+na}{forEach}\PY{p}{(}\PY{n}{System}\PY{p}{.}\PY{n+na}{out}\PY{p}{:}\PY{p}{:}\PY{n}{println}\PY{p}{)}\PY{p}{;}\PY{+w}{ }\PY{p}{\PYZcb{}}

\PY{c+c1}{// With a bounded wildcard: any subtype of Number works too (producer \PYZhy{}\PYZgt{} extends)}
\PY{k+kt}{void}\PY{+w}{ }\PY{n+nf}{printAll}\PY{p}{(}\PY{n}{List}\PY{o}{\PYZlt{}}\PY{o}{?}\PY{+w}{ }\PY{k+kd}{extends}\PY{+w}{ }\PY{n}{Number}\PY{o}{\PYZgt{}}\PY{+w}{ }\PY{n}{list}\PY{p}{)}\PY{+w}{ }\PY{p}{\PYZob{}}\PY{+w}{ }\PY{n}{list}\PY{p}{.}\PY{n+na}{forEach}\PY{p}{(}\PY{n}{System}\PY{p}{.}\PY{n+na}{out}\PY{p}{:}\PY{p}{:}\PY{n}{println}\PY{p}{)}\PY{p}{;}\PY{+w}{ }\PY{p}{\PYZcb{}}

\PY{c+c1}{// Consumer \PYZhy{}\PYZgt{} super: this can accept a destination for any supertype of Integer}
\PY{k+kt}{void}\PY{+w}{ }\PY{n+nf}{addIntegers}\PY{p}{(}\PY{n}{List}\PY{o}{\PYZlt{}}\PY{o}{?}\PY{+w}{ }\PY{k+kd}{super}\PY{+w}{ }\PY{n}{Integer}\PY{o}{\PYZgt{}}\PY{+w}{ }\PY{n}{destination}\PY{p}{)}\PY{+w}{ }\PY{p}{\PYZob{}}\PY{+w}{ }\PY{n}{destination}\PY{p}{.}\PY{n+na}{add}\PY{p}{(}\PY{l+m+mi}{1}\PY{p}{)}\PY{p}{;}\PY{+w}{ }\PY{n}{destination}\PY{p}{.}\PY{n+na}{add}\PY{p}{(}\PY{l+m+mi}{2}\PY{p}{)}\PY{p}{;}\PY{+w}{ }\PY{p}{\PYZcb{}}
\end{Verbatim}
\end{codebox}

\subsection[13. Prefer enums to int constants]{13. Prefer enums to \code{int} constants}
\label{h:19:13-prefer-enums-to-int-constants}
\code{int} constants (“int enum pattern”) give no type safety and no useful \code{toString()}. Enums are type-checked, printable, and can carry behavior.

\begin{codebox}
\begin{Verbatim}[commandchars=\\\{\}]
\PY{c+c1}{// int constants — nothing stops you from passing the wrong \PYZdq{}kind\PYZdq{} of int}
\PY{k+kd}{class} \PY{n+nc}{Season}\PY{+w}{ }\PY{p}{\PYZob{}}\PY{+w}{ }\PY{k+kd}{static}\PY{+w}{ }\PY{k+kd}{final}\PY{+w}{ }\PY{k+kt}{int}\PY{+w}{ }\PY{n}{SPRING}\PY{+w}{ }\PY{o}{=}\PY{+w}{ }\PY{l+m+mi}{0}\PY{p}{,}\PY{+w}{ }\PY{n}{SUMMER}\PY{+w}{ }\PY{o}{=}\PY{+w}{ }\PY{l+m+mi}{1}\PY{p}{,}\PY{+w}{ }\PY{n}{FALL}\PY{+w}{ }\PY{o}{=}\PY{+w}{ }\PY{l+m+mi}{2}\PY{p}{,}\PY{+w}{ }\PY{n}{WINTER}\PY{+w}{ }\PY{o}{=}\PY{+w}{ }\PY{l+m+mi}{3}\PY{p}{;}\PY{+w}{ }\PY{p}{\PYZcb{}}
\PY{k+kt}{void}\PY{+w}{ }\PY{n+nf}{plant}\PY{p}{(}\PY{k+kt}{int}\PY{+w}{ }\PY{n}{season}\PY{p}{)}\PY{+w}{ }\PY{p}{\PYZob{}}\PY{+w}{ }\PY{p}{.}\PY{p}{.}\PY{p}{.}\PY{+w}{ }\PY{p}{\PYZcb{}}
\PY{n}{plant}\PY{p}{(}\PY{n}{Season}\PY{p}{.}\PY{n+na}{SUMMER}\PY{p}{)}\PY{p}{;}\PY{+w}{ }\PY{c+c1}{// fine}
\PY{n}{plant}\PY{p}{(}\PY{l+m+mi}{42}\PY{p}{)}\PY{p}{;}\PY{+w}{            }\PY{c+c1}{// compiles, wrong at runtime}

\PY{c+c1}{// enum — type\PYZhy{}checked, self\PYZhy{}documenting, switch\PYZhy{}friendly}
\PY{k+kd}{enum}\PY{+w}{ }\PY{n}{Season}\PY{+w}{ }\PY{p}{\PYZob{}}\PY{+w}{ }\PY{n}{SPRING}\PY{p}{,}\PY{+w}{ }\PY{n}{SUMMER}\PY{p}{,}\PY{+w}{ }\PY{n}{FALL}\PY{p}{,}\PY{+w}{ }\PY{n}{WINTER}\PY{+w}{ }\PY{p}{\PYZcb{}}
\PY{k+kt}{void}\PY{+w}{ }\PY{n+nf}{plant}\PY{p}{(}\PY{n}{Season}\PY{+w}{ }\PY{n}{season}\PY{p}{)}\PY{+w}{ }\PY{p}{\PYZob{}}\PY{+w}{ }\PY{p}{.}\PY{p}{.}\PY{p}{.}\PY{+w}{ }\PY{p}{\PYZcb{}}
\end{Verbatim}
\end{codebox}

\subsection[14. Return empty collections or arrays, never null]{14. Return empty collections or arrays, never \code{null}}
\label{h:19:14-return-empty-collections-or-arrays-never-null}
Returning \code{null} for “no results” forces every caller to null-check, and someone eventually forgets and gets a \code{NullPointerException}. Return an empty collection instead.

\begin{codebox}
\begin{Verbatim}[commandchars=\\\{\}]
\PY{c+c1}{// Before: every caller must remember to null\PYZhy{}check}
\PY{n}{List}\PY{o}{\PYZlt{}}\PY{n}{Order}\PY{o}{\PYZgt{}}\PY{+w}{ }\PY{n+nf}{findOrders}\PY{p}{(}\PY{n}{String}\PY{+w}{ }\PY{n}{customerId}\PY{p}{)}\PY{+w}{ }\PY{p}{\PYZob{}}
\PY{+w}{    }\PY{k}{return}\PY{+w}{ }\PY{n}{orders}\PY{p}{.}\PY{n+na}{isEmpty}\PY{p}{(}\PY{p}{)}\PY{+w}{ }\PY{o}{?}\PY{+w}{ }\PY{k+kc}{null}\PY{+w}{ }\PY{p}{:}\PY{+w}{ }\PY{n}{orders}\PY{p}{;}
\PY{p}{\PYZcb{}}

\PY{c+c1}{// After: caller can always safely iterate}
\PY{n}{List}\PY{o}{\PYZlt{}}\PY{n}{Order}\PY{o}{\PYZgt{}}\PY{+w}{ }\PY{n+nf}{findOrders}\PY{p}{(}\PY{n}{String}\PY{+w}{ }\PY{n}{customerId}\PY{p}{)}\PY{+w}{ }\PY{p}{\PYZob{}}
\PY{+w}{    }\PY{k}{return}\PY{+w}{ }\PY{n}{orders}\PY{p}{.}\PY{n+na}{isEmpty}\PY{p}{(}\PY{p}{)}\PY{+w}{ }\PY{o}{?}\PY{+w}{ }\PY{n}{Collections}\PY{p}{.}\PY{n+na}{emptyList}\PY{p}{(}\PY{p}{)}\PY{+w}{ }\PY{p}{:}\PY{+w}{ }\PY{n}{orders}\PY{p}{;}
\PY{p}{\PYZcb{}}
\end{Verbatim}
\end{codebox}

\subsection[15. Use Optional judiciously — mainly as a return type, not everywhere]{15. Use \code{Optional} judiciously — mainly as a return type, not everywhere}
\label{h:19:15-use-optional-judiciously--mainly-as-a-return-type-not-everywhere}
\code{Optional\textless{}\allowbreak{}T\textgreater{}} is great for a method’s return type when “no result” is a normal, expected outcome. It is a poor fit for fields, method parameters, or collection elements — those usages just add wrapping overhead without benefit.

\begin{codebox}
\begin{Verbatim}[commandchars=\\\{\}]
\PY{c+c1}{// Good use: return type communicates \PYZdq{}might be absent\PYZdq{}}
\PY{n}{Optional}\PY{o}{\PYZlt{}}\PY{n}{User}\PY{o}{\PYZgt{}}\PY{+w}{ }\PY{n+nf}{findByEmail}\PY{p}{(}\PY{n}{String}\PY{+w}{ }\PY{n}{email}\PY{p}{)}\PY{+w}{ }\PY{p}{\PYZob{}}\PY{+w}{ }\PY{p}{.}\PY{p}{.}\PY{p}{.}\PY{+w}{ }\PY{p}{\PYZcb{}}

\PY{n}{Optional}\PY{o}{\PYZlt{}}\PY{n}{User}\PY{o}{\PYZgt{}}\PY{+w}{ }\PY{n}{user}\PY{+w}{ }\PY{o}{=}\PY{+w}{ }\PY{n}{findByEmail}\PY{p}{(}\PY{l+s}{\PYZdq{}}\PY{l+s}{a@example.com}\PY{l+s}{\PYZdq{}}\PY{p}{)}\PY{p}{;}
\PY{n}{user}\PY{p}{.}\PY{n+na}{ifPresentOrElse}\PY{p}{(}\PY{k}{this}\PY{p}{:}\PY{p}{:}\PY{n}{greet}\PY{p}{,}\PY{+w}{ }\PY{p}{(}\PY{p}{)}\PY{+w}{ }\PY{o}{\PYZhy{}}\PY{o}{\PYZgt{}}\PY{+w}{ }\PY{n}{log}\PY{p}{.}\PY{n+na}{warn}\PY{p}{(}\PY{l+s}{\PYZdq{}}\PY{l+s}{not found}\PY{l+s}{\PYZdq{}}\PY{p}{)}\PY{p}{)}\PY{p}{;}

\PY{c+c1}{// Bad use: Optional as a field or parameter just adds ceremony}
\PY{k+kd}{class} \PY{n+nc}{Order}\PY{+w}{ }\PY{p}{\PYZob{}}
\PY{+w}{    }\PY{k+kd}{private}\PY{+w}{ }\PY{n}{Optional}\PY{o}{\PYZlt{}}\PY{n}{String}\PY{o}{\PYZgt{}}\PY{+w}{ }\PY{n}{couponCode}\PY{p}{;}\PY{+w}{ }\PY{c+c1}{// avoid — use null or a default value/empty string instead}
\PY{p}{\PYZcb{}}
\end{Verbatim}
\end{codebox}

\subsection[16. Minimize the scope of local variables]{16. Minimize the scope of local variables}
\label{h:19:16-minimize-the-scope-of-local-variables}
Declare a variable as close as possible to where it is first used, and in the narrowest block possible. This shrinks the window in which a reader has to track its value.

\begin{codebox}
\begin{Verbatim}[commandchars=\\\{\}]
\PY{c+c1}{// Before: i is declared far from where the loop actually needs it}
\PY{k+kt}{int}\PY{+w}{ }\PY{n}{i}\PY{p}{;}
\PY{c+c1}{// ... 20 lines of unrelated code ...}
\PY{k}{for}\PY{+w}{ }\PY{p}{(}\PY{n}{i}\PY{+w}{ }\PY{o}{=}\PY{+w}{ }\PY{l+m+mi}{0}\PY{p}{;}\PY{+w}{ }\PY{n}{i}\PY{+w}{ }\PY{o}{\PYZlt{}}\PY{+w}{ }\PY{n}{items}\PY{p}{.}\PY{n+na}{size}\PY{p}{(}\PY{p}{)}\PY{p}{;}\PY{+w}{ }\PY{n}{i}\PY{o}{+}\PY{o}{+}\PY{p}{)}\PY{+w}{ }\PY{p}{\PYZob{}}\PY{+w}{ }\PY{n}{process}\PY{p}{(}\PY{n}{items}\PY{p}{.}\PY{n+na}{get}\PY{p}{(}\PY{n}{i}\PY{p}{)}\PY{p}{)}\PY{p}{;}\PY{+w}{ }\PY{p}{\PYZcb{}}

\PY{c+c1}{// After: declared right where it\PYZsq{}s needed, scoped to the loop}
\PY{k}{for}\PY{+w}{ }\PY{p}{(}\PY{k+kt}{int}\PY{+w}{ }\PY{n}{i}\PY{+w}{ }\PY{o}{=}\PY{+w}{ }\PY{l+m+mi}{0}\PY{p}{;}\PY{+w}{ }\PY{n}{i}\PY{+w}{ }\PY{o}{\PYZlt{}}\PY{+w}{ }\PY{n}{items}\PY{p}{.}\PY{n+na}{size}\PY{p}{(}\PY{p}{)}\PY{p}{;}\PY{+w}{ }\PY{n}{i}\PY{o}{+}\PY{o}{+}\PY{p}{)}\PY{+w}{ }\PY{p}{\PYZob{}}\PY{+w}{ }\PY{n}{process}\PY{p}{(}\PY{n}{items}\PY{p}{.}\PY{n+na}{get}\PY{p}{(}\PY{n}{i}\PY{p}{)}\PY{p}{)}\PY{p}{;}\PY{+w}{ }\PY{p}{\PYZcb{}}
\end{Verbatim}
\end{codebox}

\subsection[17. Prefer for-each loops to traditional index-based for loops]{17. Prefer \code{for-\allowbreak{}each} loops to traditional index-based \code{for} loops}
\label{h:19:17-prefer-for-each-loops-to-traditional-index-based-for-loops}
\code{for-\allowbreak{}each} removes the off-by-one and index-mutation bugs that plague manual index loops, and it works uniformly across arrays, lists, and any \code{Iterable}.

\begin{codebox}
\begin{Verbatim}[commandchars=\\\{\}]
\PY{c+c1}{// Index\PYZhy{}based: easy to get the bound wrong}
\PY{k}{for}\PY{+w}{ }\PY{p}{(}\PY{k+kt}{int}\PY{+w}{ }\PY{n}{i}\PY{+w}{ }\PY{o}{=}\PY{+w}{ }\PY{l+m+mi}{0}\PY{p}{;}\PY{+w}{ }\PY{n}{i}\PY{+w}{ }\PY{o}{\PYZlt{}}\PY{o}{=}\PY{+w}{ }\PY{n}{items}\PY{p}{.}\PY{n+na}{size}\PY{p}{(}\PY{p}{)}\PY{p}{;}\PY{+w}{ }\PY{n}{i}\PY{o}{+}\PY{o}{+}\PY{p}{)}\PY{+w}{ }\PY{p}{\PYZob{}}\PY{+w}{ }\PY{c+c1}{// bug: should be \PYZlt{}}
\PY{+w}{    }\PY{n}{process}\PY{p}{(}\PY{n}{items}\PY{p}{.}\PY{n+na}{get}\PY{p}{(}\PY{n}{i}\PY{p}{)}\PY{p}{)}\PY{p}{;}
\PY{p}{\PYZcb{}}

\PY{c+c1}{// for\PYZhy{}each: no index to get wrong}
\PY{k}{for}\PY{+w}{ }\PY{p}{(}\PY{n}{Item}\PY{+w}{ }\PY{n}{item}\PY{+w}{ }\PY{p}{:}\PY{+w}{ }\PY{n}{items}\PY{p}{)}\PY{+w}{ }\PY{p}{\PYZob{}}
\PY{+w}{    }\PY{n}{process}\PY{p}{(}\PY{n}{item}\PY{p}{)}\PY{p}{;}
\PY{p}{\PYZcb{}}
\end{Verbatim}
\end{codebox}

\subsection[18. Beware the performance cost of the String concatenation + operator in loops]{18. Beware the performance cost of the \code{String} concatenation \code{+} operator in loops}
\label{h:19:18-beware-the-performance-cost-of-the-string-concatenation--operator-in-loops}
Each \code{+} on strings inside a loop creates a new \code{String} object, giving O(n²) behavior for n concatenations. Use \code{StringBuilder} for loops.

\begin{codebox}
\begin{Verbatim}[commandchars=\\\{\}]
\PY{c+c1}{// Before: O(n\PYZca{}2) — a new String is allocated on every iteration}
\PY{n}{String}\PY{+w}{ }\PY{n}{result}\PY{+w}{ }\PY{o}{=}\PY{+w}{ }\PY{l+s}{\PYZdq{}}\PY{l+s}{\PYZdq{}}\PY{p}{;}
\PY{k}{for}\PY{+w}{ }\PY{p}{(}\PY{n}{String}\PY{+w}{ }\PY{n}{word}\PY{+w}{ }\PY{p}{:}\PY{+w}{ }\PY{n}{words}\PY{p}{)}\PY{+w}{ }\PY{p}{\PYZob{}}
\PY{+w}{    }\PY{n}{result}\PY{+w}{ }\PY{o}{+}\PY{o}{=}\PY{+w}{ }\PY{n}{word}\PY{+w}{ }\PY{o}{+}\PY{+w}{ }\PY{l+s}{\PYZdq{}}\PY{l+s}{ }\PY{l+s}{\PYZdq{}}\PY{p}{;}
\PY{p}{\PYZcb{}}

\PY{c+c1}{// After: O(n) — one growable buffer}
\PY{n}{StringBuilder}\PY{+w}{ }\PY{n}{sb}\PY{+w}{ }\PY{o}{=}\PY{+w}{ }\PY{k}{new}\PY{+w}{ }\PY{n}{StringBuilder}\PY{p}{(}\PY{p}{)}\PY{p}{;}
\PY{k}{for}\PY{+w}{ }\PY{p}{(}\PY{n}{String}\PY{+w}{ }\PY{n}{word}\PY{+w}{ }\PY{p}{:}\PY{+w}{ }\PY{n}{words}\PY{p}{)}\PY{+w}{ }\PY{p}{\PYZob{}}
\PY{+w}{    }\PY{n}{sb}\PY{p}{.}\PY{n+na}{append}\PY{p}{(}\PY{n}{word}\PY{p}{)}\PY{p}{.}\PY{n+na}{append}\PY{p}{(}\PY{l+s+sc}{\PYZsq{} \PYZsq{}}\PY{p}{)}\PY{p}{;}
\PY{p}{\PYZcb{}}
\PY{n}{String}\PY{+w}{ }\PY{n}{result}\PY{+w}{ }\PY{o}{=}\PY{+w}{ }\PY{n}{sb}\PY{p}{.}\PY{n+na}{toString}\PY{p}{(}\PY{p}{)}\PY{p}{;}
\end{Verbatim}
\end{codebox}

\subsection[19. Refer to objects by their interface, not their implementation type]{19. Refer to objects by their interface, not their implementation type}
\label{h:19:19-refer-to-objects-by-their-interface-not-their-implementation-type}
Declaring variables and parameters using the interface type (\code{List}, \code{Map}, \code{Set}) instead of the concrete class (\code{ArrayList}, \code{HashMap}) makes it trivial to swap implementations later.

\begin{codebox}
\begin{Verbatim}[commandchars=\\\{\}]
\PY{c+c1}{// Before: tied to ArrayList everywhere}
\PY{n}{ArrayList}\PY{o}{\PYZlt{}}\PY{n}{String}\PY{o}{\PYZgt{}}\PY{+w}{ }\PY{n}{names}\PY{+w}{ }\PY{o}{=}\PY{+w}{ }\PY{k}{new}\PY{+w}{ }\PY{n}{ArrayList}\PY{o}{\PYZlt{}}\PY{o}{\PYZgt{}}\PY{p}{(}\PY{p}{)}\PY{p}{;}
\PY{k+kt}{void}\PY{+w}{ }\PY{n+nf}{printNames}\PY{p}{(}\PY{n}{ArrayList}\PY{o}{\PYZlt{}}\PY{n}{String}\PY{o}{\PYZgt{}}\PY{+w}{ }\PY{n}{names}\PY{p}{)}\PY{+w}{ }\PY{p}{\PYZob{}}\PY{+w}{ }\PY{p}{.}\PY{p}{.}\PY{p}{.}\PY{+w}{ }\PY{p}{\PYZcb{}}

\PY{c+c1}{// After: depends only on the List contract}
\PY{n}{List}\PY{o}{\PYZlt{}}\PY{n}{String}\PY{o}{\PYZgt{}}\PY{+w}{ }\PY{n}{names}\PY{+w}{ }\PY{o}{=}\PY{+w}{ }\PY{k}{new}\PY{+w}{ }\PY{n}{ArrayList}\PY{o}{\PYZlt{}}\PY{o}{\PYZgt{}}\PY{p}{(}\PY{p}{)}\PY{p}{;}
\PY{k+kt}{void}\PY{+w}{ }\PY{n+nf}{printNames}\PY{p}{(}\PY{n}{List}\PY{o}{\PYZlt{}}\PY{n}{String}\PY{o}{\PYZgt{}}\PY{+w}{ }\PY{n}{names}\PY{p}{)}\PY{+w}{ }\PY{p}{\PYZob{}}\PY{+w}{ }\PY{p}{.}\PY{p}{.}\PY{p}{.}\PY{+w}{ }\PY{p}{\PYZcb{}}
\PY{c+c1}{// Switching to LinkedList later requires changing only one line}
\end{Verbatim}
\end{codebox}

\subsection[20. Prefer primitives to boxed primitives when you don’t need null]{20. Prefer primitives to boxed primitives when you don’t need \code{null}}
\label{h:19:20-prefer-primitives-to-boxed-primitives-when-you-dont-need-null}
Boxed types (\code{Integer}, \code{Long}, \code{Double}) allow \code{null}, incur autoboxing overhead, and \code{==} compares references (not values) once outside the small cached integer range — a classic bug source. Use primitives unless you specifically need nullability or generics.

\begin{codebox}
\begin{Verbatim}[commandchars=\\\{\}]
\PY{c+c1}{// Boxed: subtle == bug and unnecessary allocation}
\PY{n}{Integer}\PY{+w}{ }\PY{n}{a}\PY{+w}{ }\PY{o}{=}\PY{+w}{ }\PY{l+m+mi}{1000}\PY{p}{;}
\PY{n}{Integer}\PY{+w}{ }\PY{n}{b}\PY{+w}{ }\PY{o}{=}\PY{+w}{ }\PY{l+m+mi}{1000}\PY{p}{;}
\PY{n}{System}\PY{p}{.}\PY{n+na}{out}\PY{p}{.}\PY{n+na}{println}\PY{p}{(}\PY{n}{a}\PY{+w}{ }\PY{o}{=}\PY{o}{=}\PY{+w}{ }\PY{n}{b}\PY{p}{)}\PY{p}{;}\PY{+w}{ }\PY{c+c1}{// false! outside the \PYZhy{}128..127 cache range}

\PY{c+c1}{// Primitive: no boxing, no reference\PYZhy{}equality trap}
\PY{k+kt}{int}\PY{+w}{ }\PY{n}{a}\PY{+w}{ }\PY{o}{=}\PY{+w}{ }\PY{l+m+mi}{1000}\PY{p}{;}
\PY{k+kt}{int}\PY{+w}{ }\PY{n}{b}\PY{+w}{ }\PY{o}{=}\PY{+w}{ }\PY{l+m+mi}{1000}\PY{p}{;}
\PY{n}{System}\PY{p}{.}\PY{n+na}{out}\PY{p}{.}\PY{n+na}{println}\PY{p}{(}\PY{n}{a}\PY{+w}{ }\PY{o}{=}\PY{o}{=}\PY{+w}{ }\PY{n}{b}\PY{p}{)}\PY{p}{;}\PY{+w}{ }\PY{c+c1}{// true, compares values}
\end{Verbatim}
\end{codebox}

\subsection[21. Document thread safety]{21. Document thread safety}
\label{h:19:21-document-thread-safety}
A class’s Javadoc should state whether it is thread-safe, conditionally thread-safe, or not thread-safe at all. Without this, callers either over-synchronize (hurting performance) or under-synchronize (causing bugs).

\begin{codebox}
\begin{Verbatim}[commandchars=\\\{\}]
\PY{c+cm}{/**}
\PY{c+cm}{ * Thread\PYZhy{}safe: all mutating methods are synchronized internally.}
\PY{c+cm}{ * Callers do not need external locking.}
\PY{c+cm}{ */}
\PY{k+kd}{class} \PY{n+nc}{Counter}\PY{+w}{ }\PY{p}{\PYZob{}}
\PY{+w}{    }\PY{k+kd}{private}\PY{+w}{ }\PY{k+kt}{int}\PY{+w}{ }\PY{n}{count}\PY{p}{;}
\PY{+w}{    }\PY{k+kd}{synchronized}\PY{+w}{ }\PY{k+kt}{void}\PY{+w}{ }\PY{n+nf}{increment}\PY{p}{(}\PY{p}{)}\PY{+w}{ }\PY{p}{\PYZob{}}\PY{+w}{ }\PY{n}{count}\PY{o}{+}\PY{o}{+}\PY{p}{;}\PY{+w}{ }\PY{p}{\PYZcb{}}
\PY{+w}{    }\PY{k+kd}{synchronized}\PY{+w}{ }\PY{k+kt}{int}\PY{+w}{ }\PY{n+nf}{get}\PY{p}{(}\PY{p}{)}\PY{+w}{ }\PY{p}{\PYZob{}}\PY{+w}{ }\PY{k}{return}\PY{+w}{ }\PY{n}{count}\PY{p}{;}\PY{+w}{ }\PY{p}{\PYZcb{}}
\PY{p}{\PYZcb{}}

\PY{c+cm}{/**}
\PY{c+cm}{ * Not thread\PYZhy{}safe. Callers must provide external synchronization}
\PY{c+cm}{ * if this instance is shared across threads.}
\PY{c+cm}{ */}
\PY{k+kd}{class} \PY{n+nc}{Buffer}\PY{+w}{ }\PY{p}{\PYZob{}}\PY{+w}{ }\PY{p}{.}\PY{p}{.}\PY{p}{.}\PY{+w}{ }\PY{p}{\PYZcb{}}
\end{Verbatim}
\end{codebox}

\subsection[22. Prefer executors and the java.util.concurrent framework to raw threads]{22. Prefer executors and the \code{java.\allowbreak{}util.\allowbreak{}concurrent} framework to raw threads}
\label{h:19:22-prefer-executors-and-the-javautilconcurrent-framework-to-raw-threads}
Creating raw \code{Thread} objects by hand gives you no lifecycle management, no bounded resource usage, and no easy way to wait for results. \code{ExecutorService} (and Java 21’s virtual threads) handle pooling, scheduling, and shutdown correctly.

\begin{codebox}
\begin{Verbatim}[commandchars=\\\{\}]
\PY{c+c1}{// Before: unmanaged raw thread, no pooling, fire\PYZhy{}and\PYZhy{}forget}
\PY{k}{new}\PY{+w}{ }\PY{n}{Thread}\PY{p}{(}\PY{p}{(}\PY{p}{)}\PY{+w}{ }\PY{o}{\PYZhy{}}\PY{o}{\PYZgt{}}\PY{+w}{ }\PY{n}{processOrder}\PY{p}{(}\PY{n}{order}\PY{p}{)}\PY{p}{)}\PY{p}{.}\PY{n+na}{start}\PY{p}{(}\PY{p}{)}\PY{p}{;}

\PY{c+c1}{// After: managed executor, bounded resources, results and errors handled}
\PY{k}{try}\PY{+w}{ }\PY{p}{(}\PY{n}{ExecutorService}\PY{+w}{ }\PY{n}{executor}\PY{+w}{ }\PY{o}{=}\PY{+w}{ }\PY{n}{Executors}\PY{p}{.}\PY{n+na}{newVirtualThreadPerTaskExecutor}\PY{p}{(}\PY{p}{)}\PY{p}{)}\PY{+w}{ }\PY{p}{\PYZob{}}
\PY{+w}{    }\PY{n}{Future}\PY{o}{\PYZlt{}}\PY{n}{Receipt}\PY{o}{\PYZgt{}}\PY{+w}{ }\PY{n}{future}\PY{+w}{ }\PY{o}{=}\PY{+w}{ }\PY{n}{executor}\PY{p}{.}\PY{n+na}{submit}\PY{p}{(}\PY{p}{(}\PY{p}{)}\PY{+w}{ }\PY{o}{\PYZhy{}}\PY{o}{\PYZgt{}}\PY{+w}{ }\PY{n}{processOrder}\PY{p}{(}\PY{n}{order}\PY{p}{)}\PY{p}{)}\PY{p}{;}
\PY{+w}{    }\PY{n}{Receipt}\PY{+w}{ }\PY{n}{receipt}\PY{+w}{ }\PY{o}{=}\PY{+w}{ }\PY{n}{future}\PY{p}{.}\PY{n+na}{get}\PY{p}{(}\PY{p}{)}\PY{p}{;}
\PY{p}{\PYZcb{}}
\end{Verbatim}
\end{codebox}

\section[Design Patterns in Java]{Design Patterns in Java}
\label{h:19:design-patterns-in-java}
A \textbf{design pattern} is a proven, reusable solution to a common design problem. Patterns are grouped into three families: \textbf{creational} (object creation), \textbf{structural} (composing classes/objects), and \textbf{behavioral} (communication between objects).

{\tablefont
\begin{xltabular}{\linewidth}{@{}L{0.449}L{0.437}L{1.412}L{1.701}@{}}
\toprule
\rowcolor{tablehdr}
\thd{Pattern} & \thd{Family} & \thd{When to use} & \thd{Modern Java alternative} \\
\midrule
\endhead
Builder & Creational & Many constructor parameters, optional fields & Record with a builder, or named factory methods \\
Factory Method & Creational & Subclass decides which concrete type to create & Static factory methods (item 1 above) \\
Abstract Factory & Creational & Create families of related objects & Dependency injection container \\
Singleton & Creational & Exactly one shared instance & Enum singleton, or DI-managed scope \\
Strategy & Behavioral & Swap an algorithm at runtime & Lambda / method reference implementing a functional interface \\
Template Method & Behavioral & Fixed skeleton, customizable steps & Default methods on an interface, or a \code{Consumer}/\code{Function} parameter \\
Observer & Behavioral & Notify many listeners of a change & \code{java.\allowbreak{}util.\allowbreak{}function} callbacks, event bus, \code{Flow}/reactive streams \\
Decorator & Structural & Add behavior without changing the class & Wrapping with lambdas, \code{Stream} pipelines \\
Adapter & Structural & Make an incompatible interface fit & A thin wrapper class or method reference \\
Facade & Structural & Simplify a complex subsystem behind one API & A service class exposing a narrow API \\
Proxy & Structural & Control access (lazy load, security, remote) & Dynamic proxies, AOP frameworks (Spring AOP) \\
Command & Behavioral & Encapsulate a request as an object & \code{Runnable}/\code{Callable} lambdas \\
Chain of Responsibility & Behavioral & Pass a request along a chain of handlers & \code{Stream} of handlers, servlet filter chains \\
Iterator & Behavioral & Traverse a collection without exposing internals & Built into \code{Iterable}, enhanced for-each \\
State & Behavioral & Behavior changes with internal state & Sealed interface + \code{switch} pattern matching \\
Visitor & Behavioral & Operate over a closed set of types & Sealed interface + exhaustive \code{switch} pattern matching \\
\bottomrule
\end{xltabular}
}

\subsection[Builder]{Builder}
\label{h:19:builder}
\begin{codebox}
\begin{Verbatim}[commandchars=\\\{\}]
\PY{k+kd}{class} \PY{n+nc}{Computer}\PY{+w}{ }\PY{p}{\PYZob{}}
\PY{+w}{    }\PY{k+kd}{private}\PY{+w}{ }\PY{k+kd}{final}\PY{+w}{ }\PY{n}{String}\PY{+w}{ }\PY{n}{cpu}\PY{p}{;}
\PY{+w}{    }\PY{k+kd}{private}\PY{+w}{ }\PY{k+kd}{final}\PY{+w}{ }\PY{k+kt}{int}\PY{+w}{ }\PY{n}{ramGb}\PY{p}{;}
\PY{+w}{    }\PY{k+kd}{private}\PY{+w}{ }\PY{k+kd}{final}\PY{+w}{ }\PY{k+kt}{boolean}\PY{+w}{ }\PY{n}{hasSsd}\PY{p}{;}

\PY{+w}{    }\PY{k+kd}{private}\PY{+w}{ }\PY{n+nf}{Computer}\PY{p}{(}\PY{n}{Builder}\PY{+w}{ }\PY{n}{b}\PY{p}{)}\PY{+w}{ }\PY{p}{\PYZob{}}
\PY{+w}{        }\PY{k}{this}\PY{p}{.}\PY{n+na}{cpu}\PY{+w}{ }\PY{o}{=}\PY{+w}{ }\PY{n}{b}\PY{p}{.}\PY{n+na}{cpu}\PY{p}{;}\PY{+w}{ }\PY{k}{this}\PY{p}{.}\PY{n+na}{ramGb}\PY{+w}{ }\PY{o}{=}\PY{+w}{ }\PY{n}{b}\PY{p}{.}\PY{n+na}{ramGb}\PY{p}{;}\PY{+w}{ }\PY{k}{this}\PY{p}{.}\PY{n+na}{hasSsd}\PY{+w}{ }\PY{o}{=}\PY{+w}{ }\PY{n}{b}\PY{p}{.}\PY{n+na}{hasSsd}\PY{p}{;}
\PY{+w}{    }\PY{p}{\PYZcb{}}

\PY{+w}{    }\PY{k+kd}{static}\PY{+w}{ }\PY{k+kd}{class} \PY{n+nc}{Builder}\PY{+w}{ }\PY{p}{\PYZob{}}
\PY{+w}{        }\PY{k+kd}{private}\PY{+w}{ }\PY{n}{String}\PY{+w}{ }\PY{n}{cpu}\PY{+w}{ }\PY{o}{=}\PY{+w}{ }\PY{l+s}{\PYZdq{}}\PY{l+s}{generic\PYZhy{}cpu}\PY{l+s}{\PYZdq{}}\PY{p}{;}
\PY{+w}{        }\PY{k+kd}{private}\PY{+w}{ }\PY{k+kt}{int}\PY{+w}{ }\PY{n}{ramGb}\PY{+w}{ }\PY{o}{=}\PY{+w}{ }\PY{l+m+mi}{8}\PY{p}{;}
\PY{+w}{        }\PY{k+kd}{private}\PY{+w}{ }\PY{k+kt}{boolean}\PY{+w}{ }\PY{n}{hasSsd}\PY{+w}{ }\PY{o}{=}\PY{+w}{ }\PY{k+kc}{true}\PY{p}{;}

\PY{+w}{        }\PY{n}{Builder}\PY{+w}{ }\PY{n+nf}{cpu}\PY{p}{(}\PY{n}{String}\PY{+w}{ }\PY{n}{cpu}\PY{p}{)}\PY{+w}{ }\PY{p}{\PYZob{}}\PY{+w}{ }\PY{k}{this}\PY{p}{.}\PY{n+na}{cpu}\PY{+w}{ }\PY{o}{=}\PY{+w}{ }\PY{n}{cpu}\PY{p}{;}\PY{+w}{ }\PY{k}{return}\PY{+w}{ }\PY{k}{this}\PY{p}{;}\PY{+w}{ }\PY{p}{\PYZcb{}}
\PY{+w}{        }\PY{n}{Builder}\PY{+w}{ }\PY{n+nf}{ramGb}\PY{p}{(}\PY{k+kt}{int}\PY{+w}{ }\PY{n}{ramGb}\PY{p}{)}\PY{+w}{ }\PY{p}{\PYZob{}}\PY{+w}{ }\PY{k}{this}\PY{p}{.}\PY{n+na}{ramGb}\PY{+w}{ }\PY{o}{=}\PY{+w}{ }\PY{n}{ramGb}\PY{p}{;}\PY{+w}{ }\PY{k}{return}\PY{+w}{ }\PY{k}{this}\PY{p}{;}\PY{+w}{ }\PY{p}{\PYZcb{}}
\PY{+w}{        }\PY{n}{Builder}\PY{+w}{ }\PY{n+nf}{hasSsd}\PY{p}{(}\PY{k+kt}{boolean}\PY{+w}{ }\PY{n}{hasSsd}\PY{p}{)}\PY{+w}{ }\PY{p}{\PYZob{}}\PY{+w}{ }\PY{k}{this}\PY{p}{.}\PY{n+na}{hasSsd}\PY{+w}{ }\PY{o}{=}\PY{+w}{ }\PY{n}{hasSsd}\PY{p}{;}\PY{+w}{ }\PY{k}{return}\PY{+w}{ }\PY{k}{this}\PY{p}{;}\PY{+w}{ }\PY{p}{\PYZcb{}}
\PY{+w}{        }\PY{n}{Computer}\PY{+w}{ }\PY{n+nf}{build}\PY{p}{(}\PY{p}{)}\PY{+w}{ }\PY{p}{\PYZob{}}\PY{+w}{ }\PY{k}{return}\PY{+w}{ }\PY{k}{new}\PY{+w}{ }\PY{n}{Computer}\PY{p}{(}\PY{k}{this}\PY{p}{)}\PY{p}{;}\PY{+w}{ }\PY{p}{\PYZcb{}}
\PY{+w}{    }\PY{p}{\PYZcb{}}
\PY{p}{\PYZcb{}}

\PY{n}{Computer}\PY{+w}{ }\PY{n}{gamingPc}\PY{+w}{ }\PY{o}{=}\PY{+w}{ }\PY{k}{new}\PY{+w}{ }\PY{n}{Computer}\PY{p}{.}\PY{n+na}{Builder}\PY{p}{(}\PY{p}{)}\PY{p}{.}\PY{n+na}{cpu}\PY{p}{(}\PY{l+s}{\PYZdq{}}\PY{l+s}{i9}\PY{l+s}{\PYZdq{}}\PY{p}{)}\PY{p}{.}\PY{n+na}{ramGb}\PY{p}{(}\PY{l+m+mi}{32}\PY{p}{)}\PY{p}{.}\PY{n+na}{hasSsd}\PY{p}{(}\PY{k+kc}{true}\PY{p}{)}\PY{p}{.}\PY{n+na}{build}\PY{p}{(}\PY{p}{)}\PY{p}{;}
\end{Verbatim}
\end{codebox}

\subsection[Factory Method]{Factory Method}
\label{h:19:factory-method}
\begin{codebox}
\begin{Verbatim}[commandchars=\\\{\}]
\PY{k+kd}{interface} \PY{n+nc}{Notification}\PY{+w}{ }\PY{p}{\PYZob{}}
\PY{+w}{    }\PY{k+kt}{void}\PY{+w}{ }\PY{n+nf}{send}\PY{p}{(}\PY{n}{String}\PY{+w}{ }\PY{n}{message}\PY{p}{)}\PY{p}{;}
\PY{p}{\PYZcb{}}

\PY{k+kd}{class} \PY{n+nc}{EmailNotification}\PY{+w}{ }\PY{k+kd}{implements}\PY{+w}{ }\PY{n}{Notification}\PY{+w}{ }\PY{p}{\PYZob{}}
\PY{+w}{    }\PY{k+kd}{public}\PY{+w}{ }\PY{k+kt}{void}\PY{+w}{ }\PY{n+nf}{send}\PY{p}{(}\PY{n}{String}\PY{+w}{ }\PY{n}{message}\PY{p}{)}\PY{+w}{ }\PY{p}{\PYZob{}}\PY{+w}{ }\PY{n}{System}\PY{p}{.}\PY{n+na}{out}\PY{p}{.}\PY{n+na}{println}\PY{p}{(}\PY{l+s}{\PYZdq{}}\PY{l+s}{Email: }\PY{l+s}{\PYZdq{}}\PY{+w}{ }\PY{o}{+}\PY{+w}{ }\PY{n}{message}\PY{p}{)}\PY{p}{;}\PY{+w}{ }\PY{p}{\PYZcb{}}
\PY{p}{\PYZcb{}}

\PY{k+kd}{class} \PY{n+nc}{SmsNotification}\PY{+w}{ }\PY{k+kd}{implements}\PY{+w}{ }\PY{n}{Notification}\PY{+w}{ }\PY{p}{\PYZob{}}
\PY{+w}{    }\PY{k+kd}{public}\PY{+w}{ }\PY{k+kt}{void}\PY{+w}{ }\PY{n+nf}{send}\PY{p}{(}\PY{n}{String}\PY{+w}{ }\PY{n}{message}\PY{p}{)}\PY{+w}{ }\PY{p}{\PYZob{}}\PY{+w}{ }\PY{n}{System}\PY{p}{.}\PY{n+na}{out}\PY{p}{.}\PY{n+na}{println}\PY{p}{(}\PY{l+s}{\PYZdq{}}\PY{l+s}{SMS: }\PY{l+s}{\PYZdq{}}\PY{+w}{ }\PY{o}{+}\PY{+w}{ }\PY{n}{message}\PY{p}{)}\PY{p}{;}\PY{+w}{ }\PY{p}{\PYZcb{}}
\PY{p}{\PYZcb{}}

\PY{k+kd}{class} \PY{n+nc}{NotificationFactory}\PY{+w}{ }\PY{p}{\PYZob{}}
\PY{+w}{    }\PY{k+kd}{static}\PY{+w}{ }\PY{n}{Notification}\PY{+w}{ }\PY{n+nf}{create}\PY{p}{(}\PY{n}{String}\PY{+w}{ }\PY{n}{type}\PY{p}{)}\PY{+w}{ }\PY{p}{\PYZob{}}
\PY{+w}{        }\PY{k}{return}\PY{+w}{ }\PY{k}{switch}\PY{+w}{ }\PY{p}{(}\PY{n}{type}\PY{p}{)}\PY{+w}{ }\PY{p}{\PYZob{}}
\PY{+w}{            }\PY{k}{case}\PY{+w}{ }\PY{l+s}{\PYZdq{}}\PY{l+s}{email}\PY{l+s}{\PYZdq{}}\PY{+w}{ }\PY{o}{\PYZhy{}}\PY{o}{\PYZgt{}}\PY{+w}{ }\PY{k}{new}\PY{+w}{ }\PY{n}{EmailNotification}\PY{p}{(}\PY{p}{)}\PY{p}{;}
\PY{+w}{            }\PY{k}{case}\PY{+w}{ }\PY{l+s}{\PYZdq{}}\PY{l+s}{sms}\PY{l+s}{\PYZdq{}}\PY{+w}{ }\PY{o}{\PYZhy{}}\PY{o}{\PYZgt{}}\PY{+w}{ }\PY{k}{new}\PY{+w}{ }\PY{n}{SmsNotification}\PY{p}{(}\PY{p}{)}\PY{p}{;}
\PY{+w}{            }\PY{k}{default}\PY{+w}{ }\PY{o}{\PYZhy{}}\PY{o}{\PYZgt{}}\PY{+w}{ }\PY{k}{throw}\PY{+w}{ }\PY{k}{new}\PY{+w}{ }\PY{n}{IllegalArgumentException}\PY{p}{(}\PY{l+s}{\PYZdq{}}\PY{l+s}{Unknown type: }\PY{l+s}{\PYZdq{}}\PY{+w}{ }\PY{o}{+}\PY{+w}{ }\PY{n}{type}\PY{p}{)}\PY{p}{;}
\PY{+w}{        }\PY{p}{\PYZcb{}}\PY{p}{;}
\PY{+w}{    }\PY{p}{\PYZcb{}}
\PY{p}{\PYZcb{}}

\PY{n}{Notification}\PY{+w}{ }\PY{n}{n}\PY{+w}{ }\PY{o}{=}\PY{+w}{ }\PY{n}{NotificationFactory}\PY{p}{.}\PY{n+na}{create}\PY{p}{(}\PY{l+s}{\PYZdq{}}\PY{l+s}{email}\PY{l+s}{\PYZdq{}}\PY{p}{)}\PY{p}{;}
\PY{n}{n}\PY{p}{.}\PY{n+na}{send}\PY{p}{(}\PY{l+s}{\PYZdq{}}\PY{l+s}{Your order has shipped}\PY{l+s}{\PYZdq{}}\PY{p}{)}\PY{p}{;}
\end{Verbatim}
\end{codebox}

\subsection[Abstract Factory]{Abstract Factory}
\label{h:19:abstract-factory}
\begin{codebox}
\begin{Verbatim}[commandchars=\\\{\}]
\PY{k+kd}{interface} \PY{n+nc}{Button}\PY{+w}{ }\PY{p}{\PYZob{}}\PY{+w}{ }\PY{k+kt}{void}\PY{+w}{ }\PY{n+nf}{render}\PY{p}{(}\PY{p}{)}\PY{p}{;}\PY{+w}{ }\PY{p}{\PYZcb{}}
\PY{k+kd}{interface} \PY{n+nc}{Checkbox}\PY{+w}{ }\PY{p}{\PYZob{}}\PY{+w}{ }\PY{k+kt}{void}\PY{+w}{ }\PY{n+nf}{render}\PY{p}{(}\PY{p}{)}\PY{p}{;}\PY{+w}{ }\PY{p}{\PYZcb{}}

\PY{k+kd}{interface} \PY{n+nc}{UiFactory}\PY{+w}{ }\PY{p}{\PYZob{}}
\PY{+w}{    }\PY{n}{Button}\PY{+w}{ }\PY{n+nf}{createButton}\PY{p}{(}\PY{p}{)}\PY{p}{;}
\PY{+w}{    }\PY{n}{Checkbox}\PY{+w}{ }\PY{n+nf}{createCheckbox}\PY{p}{(}\PY{p}{)}\PY{p}{;}
\PY{p}{\PYZcb{}}

\PY{k+kd}{class} \PY{n+nc}{DarkThemeFactory}\PY{+w}{ }\PY{k+kd}{implements}\PY{+w}{ }\PY{n}{UiFactory}\PY{+w}{ }\PY{p}{\PYZob{}}
\PY{+w}{    }\PY{k+kd}{public}\PY{+w}{ }\PY{n}{Button}\PY{+w}{ }\PY{n+nf}{createButton}\PY{p}{(}\PY{p}{)}\PY{+w}{ }\PY{p}{\PYZob{}}\PY{+w}{ }\PY{k}{return}\PY{+w}{ }\PY{p}{(}\PY{p}{)}\PY{+w}{ }\PY{o}{\PYZhy{}}\PY{o}{\PYZgt{}}\PY{+w}{ }\PY{n}{System}\PY{p}{.}\PY{n+na}{out}\PY{p}{.}\PY{n+na}{println}\PY{p}{(}\PY{l+s}{\PYZdq{}}\PY{l+s}{Dark button}\PY{l+s}{\PYZdq{}}\PY{p}{)}\PY{p}{;}\PY{+w}{ }\PY{p}{\PYZcb{}}
\PY{+w}{    }\PY{k+kd}{public}\PY{+w}{ }\PY{n}{Checkbox}\PY{+w}{ }\PY{n+nf}{createCheckbox}\PY{p}{(}\PY{p}{)}\PY{+w}{ }\PY{p}{\PYZob{}}\PY{+w}{ }\PY{k}{return}\PY{+w}{ }\PY{p}{(}\PY{p}{)}\PY{+w}{ }\PY{o}{\PYZhy{}}\PY{o}{\PYZgt{}}\PY{+w}{ }\PY{n}{System}\PY{p}{.}\PY{n+na}{out}\PY{p}{.}\PY{n+na}{println}\PY{p}{(}\PY{l+s}{\PYZdq{}}\PY{l+s}{Dark checkbox}\PY{l+s}{\PYZdq{}}\PY{p}{)}\PY{p}{;}\PY{+w}{ }\PY{p}{\PYZcb{}}
\PY{p}{\PYZcb{}}

\PY{k+kd}{class} \PY{n+nc}{LightThemeFactory}\PY{+w}{ }\PY{k+kd}{implements}\PY{+w}{ }\PY{n}{UiFactory}\PY{+w}{ }\PY{p}{\PYZob{}}
\PY{+w}{    }\PY{k+kd}{public}\PY{+w}{ }\PY{n}{Button}\PY{+w}{ }\PY{n+nf}{createButton}\PY{p}{(}\PY{p}{)}\PY{+w}{ }\PY{p}{\PYZob{}}\PY{+w}{ }\PY{k}{return}\PY{+w}{ }\PY{p}{(}\PY{p}{)}\PY{+w}{ }\PY{o}{\PYZhy{}}\PY{o}{\PYZgt{}}\PY{+w}{ }\PY{n}{System}\PY{p}{.}\PY{n+na}{out}\PY{p}{.}\PY{n+na}{println}\PY{p}{(}\PY{l+s}{\PYZdq{}}\PY{l+s}{Light button}\PY{l+s}{\PYZdq{}}\PY{p}{)}\PY{p}{;}\PY{+w}{ }\PY{p}{\PYZcb{}}
\PY{+w}{    }\PY{k+kd}{public}\PY{+w}{ }\PY{n}{Checkbox}\PY{+w}{ }\PY{n+nf}{createCheckbox}\PY{p}{(}\PY{p}{)}\PY{+w}{ }\PY{p}{\PYZob{}}\PY{+w}{ }\PY{k}{return}\PY{+w}{ }\PY{p}{(}\PY{p}{)}\PY{+w}{ }\PY{o}{\PYZhy{}}\PY{o}{\PYZgt{}}\PY{+w}{ }\PY{n}{System}\PY{p}{.}\PY{n+na}{out}\PY{p}{.}\PY{n+na}{println}\PY{p}{(}\PY{l+s}{\PYZdq{}}\PY{l+s}{Light checkbox}\PY{l+s}{\PYZdq{}}\PY{p}{)}\PY{p}{;}\PY{+w}{ }\PY{p}{\PYZcb{}}
\PY{p}{\PYZcb{}}

\PY{c+c1}{// Note: Button/Checkbox declared as functional interfaces here for brevity via lambdas}
\PY{k+kt}{void}\PY{+w}{ }\PY{n+nf}{renderUi}\PY{p}{(}\PY{n}{UiFactory}\PY{+w}{ }\PY{n}{factory}\PY{p}{)}\PY{+w}{ }\PY{p}{\PYZob{}}
\PY{+w}{    }\PY{n}{factory}\PY{p}{.}\PY{n+na}{createButton}\PY{p}{(}\PY{p}{)}\PY{p}{.}\PY{n+na}{render}\PY{p}{(}\PY{p}{)}\PY{p}{;}
\PY{+w}{    }\PY{n}{factory}\PY{p}{.}\PY{n+na}{createCheckbox}\PY{p}{(}\PY{p}{)}\PY{p}{.}\PY{n+na}{render}\PY{p}{(}\PY{p}{)}\PY{p}{;}
\PY{p}{\PYZcb{}}

\PY{n}{renderUi}\PY{p}{(}\PY{k}{new}\PY{+w}{ }\PY{n}{DarkThemeFactory}\PY{p}{(}\PY{p}{)}\PY{p}{)}\PY{p}{;}\PY{+w}{ }\PY{c+c1}{// switching the whole family is one line}
\end{Verbatim}
\end{codebox}

\subsection[Singleton]{Singleton}
\label{h:19:singleton}
\begin{codebox}
\begin{Verbatim}[commandchars=\\\{\}]
\PY{c+c1}{// Classic non\PYZhy{}enum singleton (shown for comparison — avoid in new code)}
\PY{k+kd}{class} \PY{n+nc}{LegacyConfig}\PY{+w}{ }\PY{p}{\PYZob{}}
\PY{+w}{    }\PY{k+kd}{private}\PY{+w}{ }\PY{k+kd}{static}\PY{+w}{ }\PY{n}{LegacyConfig}\PY{+w}{ }\PY{n}{instance}\PY{p}{;}
\PY{+w}{    }\PY{k+kd}{private}\PY{+w}{ }\PY{n+nf}{LegacyConfig}\PY{p}{(}\PY{p}{)}\PY{+w}{ }\PY{p}{\PYZob{}}\PY{p}{\PYZcb{}}
\PY{+w}{    }\PY{k+kd}{static}\PY{+w}{ }\PY{k+kd}{synchronized}\PY{+w}{ }\PY{n}{LegacyConfig}\PY{+w}{ }\PY{n+nf}{getInstance}\PY{p}{(}\PY{p}{)}\PY{+w}{ }\PY{p}{\PYZob{}}
\PY{+w}{        }\PY{k}{if}\PY{+w}{ }\PY{p}{(}\PY{n}{instance}\PY{+w}{ }\PY{o}{=}\PY{o}{=}\PY{+w}{ }\PY{k+kc}{null}\PY{p}{)}\PY{+w}{ }\PY{n}{instance}\PY{+w}{ }\PY{o}{=}\PY{+w}{ }\PY{k}{new}\PY{+w}{ }\PY{n}{LegacyConfig}\PY{p}{(}\PY{p}{)}\PY{p}{;}
\PY{+w}{        }\PY{k}{return}\PY{+w}{ }\PY{n}{instance}\PY{p}{;}
\PY{+w}{    }\PY{p}{\PYZcb{}}
\PY{p}{\PYZcb{}}

\PY{c+c1}{// Modern replacement: enum singleton — serialization\PYZhy{}safe, thread\PYZhy{}safe for free}
\PY{k+kd}{enum}\PY{+w}{ }\PY{n}{AppConfig}\PY{+w}{ }\PY{p}{\PYZob{}}
\PY{+w}{    }\PY{n}{INSTANCE}\PY{p}{;}
\PY{+w}{    }\PY{k+kd}{private}\PY{+w}{ }\PY{k+kd}{final}\PY{+w}{ }\PY{n}{Map}\PY{o}{\PYZlt{}}\PY{n}{String}\PY{p}{,}\PY{+w}{ }\PY{n}{String}\PY{o}{\PYZgt{}}\PY{+w}{ }\PY{n}{settings}\PY{+w}{ }\PY{o}{=}\PY{+w}{ }\PY{k}{new}\PY{+w}{ }\PY{n}{ConcurrentHashMap}\PY{o}{\PYZlt{}}\PY{o}{\PYZgt{}}\PY{p}{(}\PY{p}{)}\PY{p}{;}
\PY{+w}{    }\PY{n}{String}\PY{+w}{ }\PY{n+nf}{get}\PY{p}{(}\PY{n}{String}\PY{+w}{ }\PY{n}{key}\PY{p}{)}\PY{+w}{ }\PY{p}{\PYZob{}}\PY{+w}{ }\PY{k}{return}\PY{+w}{ }\PY{n}{settings}\PY{p}{.}\PY{n+na}{get}\PY{p}{(}\PY{n}{key}\PY{p}{)}\PY{p}{;}\PY{+w}{ }\PY{p}{\PYZcb{}}
\PY{p}{\PYZcb{}}

\PY{n}{AppConfig}\PY{p}{.}\PY{n+na}{INSTANCE}\PY{p}{.}\PY{n+na}{get}\PY{p}{(}\PY{l+s}{\PYZdq{}}\PY{l+s}{timeout}\PY{l+s}{\PYZdq{}}\PY{p}{)}\PY{p}{;}
\end{Verbatim}
\end{codebox}

\subsection[Strategy]{Strategy}
\label{h:19:strategy}
\begin{codebox}
\begin{Verbatim}[commandchars=\\\{\}]
\PY{k+kd}{interface} \PY{n+nc}{DiscountStrategy}\PY{+w}{ }\PY{p}{\PYZob{}}
\PY{+w}{    }\PY{n}{BigDecimal}\PY{+w}{ }\PY{n+nf}{apply}\PY{p}{(}\PY{n}{BigDecimal}\PY{+w}{ }\PY{n}{price}\PY{p}{)}\PY{p}{;}
\PY{p}{\PYZcb{}}

\PY{k+kd}{class} \PY{n+nc}{PercentageDiscount}\PY{+w}{ }\PY{k+kd}{implements}\PY{+w}{ }\PY{n}{DiscountStrategy}\PY{+w}{ }\PY{p}{\PYZob{}}
\PY{+w}{    }\PY{k+kd}{private}\PY{+w}{ }\PY{k+kd}{final}\PY{+w}{ }\PY{n}{BigDecimal}\PY{+w}{ }\PY{n}{percent}\PY{p}{;}
\PY{+w}{    }\PY{n}{PercentageDiscount}\PY{p}{(}\PY{n}{BigDecimal}\PY{+w}{ }\PY{n}{percent}\PY{p}{)}\PY{+w}{ }\PY{p}{\PYZob{}}\PY{+w}{ }\PY{k}{this}\PY{p}{.}\PY{n+na}{percent}\PY{+w}{ }\PY{o}{=}\PY{+w}{ }\PY{n}{percent}\PY{p}{;}\PY{+w}{ }\PY{p}{\PYZcb{}}
\PY{+w}{    }\PY{k+kd}{public}\PY{+w}{ }\PY{n}{BigDecimal}\PY{+w}{ }\PY{n+nf}{apply}\PY{p}{(}\PY{n}{BigDecimal}\PY{+w}{ }\PY{n}{price}\PY{p}{)}\PY{+w}{ }\PY{p}{\PYZob{}}\PY{+w}{ }\PY{k}{return}\PY{+w}{ }\PY{n}{price}\PY{p}{.}\PY{n+na}{subtract}\PY{p}{(}\PY{n}{price}\PY{p}{.}\PY{n+na}{multiply}\PY{p}{(}\PY{n}{percent}\PY{p}{)}\PY{p}{)}\PY{p}{;}\PY{+w}{ }\PY{p}{\PYZcb{}}
\PY{p}{\PYZcb{}}

\PY{k+kd}{class} \PY{n+nc}{Checkout}\PY{+w}{ }\PY{p}{\PYZob{}}
\PY{+w}{    }\PY{k+kd}{private}\PY{+w}{ }\PY{k+kd}{final}\PY{+w}{ }\PY{n}{DiscountStrategy}\PY{+w}{ }\PY{n}{strategy}\PY{p}{;}
\PY{+w}{    }\PY{n}{Checkout}\PY{p}{(}\PY{n}{DiscountStrategy}\PY{+w}{ }\PY{n}{strategy}\PY{p}{)}\PY{+w}{ }\PY{p}{\PYZob{}}\PY{+w}{ }\PY{k}{this}\PY{p}{.}\PY{n+na}{strategy}\PY{+w}{ }\PY{o}{=}\PY{+w}{ }\PY{n}{strategy}\PY{p}{;}\PY{+w}{ }\PY{p}{\PYZcb{}}
\PY{+w}{    }\PY{n}{BigDecimal}\PY{+w}{ }\PY{n+nf}{total}\PY{p}{(}\PY{n}{BigDecimal}\PY{+w}{ }\PY{n}{price}\PY{p}{)}\PY{+w}{ }\PY{p}{\PYZob{}}\PY{+w}{ }\PY{k}{return}\PY{+w}{ }\PY{n}{strategy}\PY{p}{.}\PY{n+na}{apply}\PY{p}{(}\PY{n}{price}\PY{p}{)}\PY{p}{;}\PY{+w}{ }\PY{p}{\PYZcb{}}
\PY{p}{\PYZcb{}}

\PY{c+c1}{// Classic: a class implementing the interface}
\PY{n}{Checkout}\PY{+w}{ }\PY{n}{classic}\PY{+w}{ }\PY{o}{=}\PY{+w}{ }\PY{k}{new}\PY{+w}{ }\PY{n}{Checkout}\PY{p}{(}\PY{k}{new}\PY{+w}{ }\PY{n}{PercentageDiscount}\PY{p}{(}\PY{n}{BigDecimal}\PY{p}{.}\PY{n+na}{valueOf}\PY{p}{(}\PY{l+m+mf}{0.10}\PY{p}{)}\PY{p}{)}\PY{p}{)}\PY{p}{;}

\PY{c+c1}{// Modern replacement: DiscountStrategy is a functional interface, so a lambda works directly}
\PY{n}{Checkout}\PY{+w}{ }\PY{n}{modern}\PY{+w}{ }\PY{o}{=}\PY{+w}{ }\PY{k}{new}\PY{+w}{ }\PY{n}{Checkout}\PY{p}{(}\PY{n}{price}\PY{+w}{ }\PY{o}{\PYZhy{}}\PY{o}{\PYZgt{}}\PY{+w}{ }\PY{n}{price}\PY{p}{.}\PY{n+na}{subtract}\PY{p}{(}\PY{n}{price}\PY{p}{.}\PY{n+na}{multiply}\PY{p}{(}\PY{n}{BigDecimal}\PY{p}{.}\PY{n+na}{valueOf}\PY{p}{(}\PY{l+m+mf}{0.10}\PY{p}{)}\PY{p}{)}\PY{p}{)}\PY{p}{)}\PY{p}{;}
\end{Verbatim}
\end{codebox}

\subsection[Template Method]{Template Method}
\label{h:19:template-method}
\begin{codebox}
\begin{Verbatim}[commandchars=\\\{\}]
\PY{k+kd}{abstract}\PY{+w}{ }\PY{k+kd}{class} \PY{n+nc}{DataExporter}\PY{+w}{ }\PY{p}{\PYZob{}}
\PY{+w}{    }\PY{c+c1}{// Template method: fixed skeleton}
\PY{+w}{    }\PY{k+kd}{final}\PY{+w}{ }\PY{k+kt}{void}\PY{+w}{ }\PY{n+nf}{export}\PY{p}{(}\PY{p}{)}\PY{+w}{ }\PY{p}{\PYZob{}}
\PY{+w}{        }\PY{n}{openConnection}\PY{p}{(}\PY{p}{)}\PY{p}{;}
\PY{+w}{        }\PY{n}{writeHeader}\PY{p}{(}\PY{p}{)}\PY{p}{;}
\PY{+w}{        }\PY{n}{writeRows}\PY{p}{(}\PY{p}{)}\PY{p}{;}
\PY{+w}{        }\PY{n}{closeConnection}\PY{p}{(}\PY{p}{)}\PY{p}{;}
\PY{+w}{    }\PY{p}{\PYZcb{}}
\PY{+w}{    }\PY{k+kt}{void}\PY{+w}{ }\PY{n+nf}{openConnection}\PY{p}{(}\PY{p}{)}\PY{+w}{ }\PY{p}{\PYZob{}}\PY{+w}{ }\PY{n}{System}\PY{p}{.}\PY{n+na}{out}\PY{p}{.}\PY{n+na}{println}\PY{p}{(}\PY{l+s}{\PYZdq{}}\PY{l+s}{Opening}\PY{l+s}{\PYZdq{}}\PY{p}{)}\PY{p}{;}\PY{+w}{ }\PY{p}{\PYZcb{}}
\PY{+w}{    }\PY{k+kt}{void}\PY{+w}{ }\PY{n+nf}{closeConnection}\PY{p}{(}\PY{p}{)}\PY{+w}{ }\PY{p}{\PYZob{}}\PY{+w}{ }\PY{n}{System}\PY{p}{.}\PY{n+na}{out}\PY{p}{.}\PY{n+na}{println}\PY{p}{(}\PY{l+s}{\PYZdq{}}\PY{l+s}{Closing}\PY{l+s}{\PYZdq{}}\PY{p}{)}\PY{p}{;}\PY{+w}{ }\PY{p}{\PYZcb{}}
\PY{+w}{    }\PY{k+kd}{abstract}\PY{+w}{ }\PY{k+kt}{void}\PY{+w}{ }\PY{n+nf}{writeHeader}\PY{p}{(}\PY{p}{)}\PY{p}{;}
\PY{+w}{    }\PY{k+kd}{abstract}\PY{+w}{ }\PY{k+kt}{void}\PY{+w}{ }\PY{n+nf}{writeRows}\PY{p}{(}\PY{p}{)}\PY{p}{;}
\PY{p}{\PYZcb{}}

\PY{k+kd}{class} \PY{n+nc}{CsvExporter}\PY{+w}{ }\PY{k+kd}{extends}\PY{+w}{ }\PY{n}{DataExporter}\PY{+w}{ }\PY{p}{\PYZob{}}
\PY{+w}{    }\PY{k+kt}{void}\PY{+w}{ }\PY{n+nf}{writeHeader}\PY{p}{(}\PY{p}{)}\PY{+w}{ }\PY{p}{\PYZob{}}\PY{+w}{ }\PY{n}{System}\PY{p}{.}\PY{n+na}{out}\PY{p}{.}\PY{n+na}{println}\PY{p}{(}\PY{l+s}{\PYZdq{}}\PY{l+s}{id,name}\PY{l+s}{\PYZdq{}}\PY{p}{)}\PY{p}{;}\PY{+w}{ }\PY{p}{\PYZcb{}}
\PY{+w}{    }\PY{k+kt}{void}\PY{+w}{ }\PY{n+nf}{writeRows}\PY{p}{(}\PY{p}{)}\PY{+w}{ }\PY{p}{\PYZob{}}\PY{+w}{ }\PY{n}{System}\PY{p}{.}\PY{n+na}{out}\PY{p}{.}\PY{n+na}{println}\PY{p}{(}\PY{l+s}{\PYZdq{}}\PY{l+s}{1,Ana}\PY{l+s}{\PYZdq{}}\PY{p}{)}\PY{p}{;}\PY{+w}{ }\PY{p}{\PYZcb{}}
\PY{p}{\PYZcb{}}

\PY{c+c1}{// Modern alternative for simple cases: pass the customizable steps as lambdas}
\PY{k+kt}{void}\PY{+w}{ }\PY{n+nf}{export}\PY{p}{(}\PY{n}{Runnable}\PY{+w}{ }\PY{n}{writeHeader}\PY{p}{,}\PY{+w}{ }\PY{n}{Runnable}\PY{+w}{ }\PY{n}{writeRows}\PY{p}{)}\PY{+w}{ }\PY{p}{\PYZob{}}
\PY{+w}{    }\PY{n}{System}\PY{p}{.}\PY{n+na}{out}\PY{p}{.}\PY{n+na}{println}\PY{p}{(}\PY{l+s}{\PYZdq{}}\PY{l+s}{Opening}\PY{l+s}{\PYZdq{}}\PY{p}{)}\PY{p}{;}
\PY{+w}{    }\PY{n}{writeHeader}\PY{p}{.}\PY{n+na}{run}\PY{p}{(}\PY{p}{)}\PY{p}{;}
\PY{+w}{    }\PY{n}{writeRows}\PY{p}{.}\PY{n+na}{run}\PY{p}{(}\PY{p}{)}\PY{p}{;}
\PY{+w}{    }\PY{n}{System}\PY{p}{.}\PY{n+na}{out}\PY{p}{.}\PY{n+na}{println}\PY{p}{(}\PY{l+s}{\PYZdq{}}\PY{l+s}{Closing}\PY{l+s}{\PYZdq{}}\PY{p}{)}\PY{p}{;}
\PY{p}{\PYZcb{}}
\PY{n}{export}\PY{p}{(}\PY{p}{(}\PY{p}{)}\PY{+w}{ }\PY{o}{\PYZhy{}}\PY{o}{\PYZgt{}}\PY{+w}{ }\PY{n}{System}\PY{p}{.}\PY{n+na}{out}\PY{p}{.}\PY{n+na}{println}\PY{p}{(}\PY{l+s}{\PYZdq{}}\PY{l+s}{id,name}\PY{l+s}{\PYZdq{}}\PY{p}{)}\PY{p}{,}\PY{+w}{ }\PY{p}{(}\PY{p}{)}\PY{+w}{ }\PY{o}{\PYZhy{}}\PY{o}{\PYZgt{}}\PY{+w}{ }\PY{n}{System}\PY{p}{.}\PY{n+na}{out}\PY{p}{.}\PY{n+na}{println}\PY{p}{(}\PY{l+s}{\PYZdq{}}\PY{l+s}{1,Ana}\PY{l+s}{\PYZdq{}}\PY{p}{)}\PY{p}{)}\PY{p}{;}
\end{Verbatim}
\end{codebox}

\subsection[Observer]{Observer}
\label{h:19:observer}
\begin{codebox}
\begin{Verbatim}[commandchars=\\\{\}]
\PY{k+kd}{interface} \PY{n+nc}{OrderListener}\PY{+w}{ }\PY{p}{\PYZob{}}
\PY{+w}{    }\PY{k+kt}{void}\PY{+w}{ }\PY{n+nf}{onOrderPlaced}\PY{p}{(}\PY{n}{Order}\PY{+w}{ }\PY{n}{order}\PY{p}{)}\PY{p}{;}
\PY{p}{\PYZcb{}}

\PY{k+kd}{class} \PY{n+nc}{OrderService}\PY{+w}{ }\PY{p}{\PYZob{}}
\PY{+w}{    }\PY{k+kd}{private}\PY{+w}{ }\PY{k+kd}{final}\PY{+w}{ }\PY{n}{List}\PY{o}{\PYZlt{}}\PY{n}{OrderListener}\PY{o}{\PYZgt{}}\PY{+w}{ }\PY{n}{listeners}\PY{+w}{ }\PY{o}{=}\PY{+w}{ }\PY{k}{new}\PY{+w}{ }\PY{n}{ArrayList}\PY{o}{\PYZlt{}}\PY{o}{\PYZgt{}}\PY{p}{(}\PY{p}{)}\PY{p}{;}
\PY{+w}{    }\PY{k+kt}{void}\PY{+w}{ }\PY{n+nf}{addListener}\PY{p}{(}\PY{n}{OrderListener}\PY{+w}{ }\PY{n}{listener}\PY{p}{)}\PY{+w}{ }\PY{p}{\PYZob{}}\PY{+w}{ }\PY{n}{listeners}\PY{p}{.}\PY{n+na}{add}\PY{p}{(}\PY{n}{listener}\PY{p}{)}\PY{p}{;}\PY{+w}{ }\PY{p}{\PYZcb{}}
\PY{+w}{    }\PY{k+kt}{void}\PY{+w}{ }\PY{n+nf}{placeOrder}\PY{p}{(}\PY{n}{Order}\PY{+w}{ }\PY{n}{order}\PY{p}{)}\PY{+w}{ }\PY{p}{\PYZob{}}
\PY{+w}{        }\PY{c+c1}{// ... business logic ...}
\PY{+w}{        }\PY{n}{listeners}\PY{p}{.}\PY{n+na}{forEach}\PY{p}{(}\PY{n}{l}\PY{+w}{ }\PY{o}{\PYZhy{}}\PY{o}{\PYZgt{}}\PY{+w}{ }\PY{n}{l}\PY{p}{.}\PY{n+na}{onOrderPlaced}\PY{p}{(}\PY{n}{order}\PY{p}{)}\PY{p}{)}\PY{p}{;}
\PY{+w}{    }\PY{p}{\PYZcb{}}
\PY{p}{\PYZcb{}}

\PY{c+c1}{// Modern replacement: OrderListener is a functional interface, use lambdas directly}
\PY{n}{OrderService}\PY{+w}{ }\PY{n}{service}\PY{+w}{ }\PY{o}{=}\PY{+w}{ }\PY{k}{new}\PY{+w}{ }\PY{n}{OrderService}\PY{p}{(}\PY{p}{)}\PY{p}{;}
\PY{n}{service}\PY{p}{.}\PY{n+na}{addListener}\PY{p}{(}\PY{n}{order}\PY{+w}{ }\PY{o}{\PYZhy{}}\PY{o}{\PYZgt{}}\PY{+w}{ }\PY{n}{System}\PY{p}{.}\PY{n+na}{out}\PY{p}{.}\PY{n+na}{println}\PY{p}{(}\PY{l+s}{\PYZdq{}}\PY{l+s}{Send confirmation email for }\PY{l+s}{\PYZdq{}}\PY{+w}{ }\PY{o}{+}\PY{+w}{ }\PY{n}{order}\PY{p}{)}\PY{p}{)}\PY{p}{;}
\PY{n}{service}\PY{p}{.}\PY{n+na}{addListener}\PY{p}{(}\PY{n}{order}\PY{+w}{ }\PY{o}{\PYZhy{}}\PY{o}{\PYZgt{}}\PY{+w}{ }\PY{n}{System}\PY{p}{.}\PY{n+na}{out}\PY{p}{.}\PY{n+na}{println}\PY{p}{(}\PY{l+s}{\PYZdq{}}\PY{l+s}{Update analytics for }\PY{l+s}{\PYZdq{}}\PY{+w}{ }\PY{o}{+}\PY{+w}{ }\PY{n}{order}\PY{p}{)}\PY{p}{)}\PY{p}{;}
\end{Verbatim}
\end{codebox}

\subsection[Decorator]{Decorator}
\label{h:19:decorator}
\begin{codebox}
\begin{Verbatim}[commandchars=\\\{\}]
\PY{k+kd}{interface} \PY{n+nc}{Coffee}\PY{+w}{ }\PY{p}{\PYZob{}}
\PY{+w}{    }\PY{n}{BigDecimal}\PY{+w}{ }\PY{n+nf}{cost}\PY{p}{(}\PY{p}{)}\PY{p}{;}
\PY{+w}{    }\PY{n}{String}\PY{+w}{ }\PY{n+nf}{description}\PY{p}{(}\PY{p}{)}\PY{p}{;}
\PY{p}{\PYZcb{}}

\PY{k+kd}{class} \PY{n+nc}{SimpleCoffee}\PY{+w}{ }\PY{k+kd}{implements}\PY{+w}{ }\PY{n}{Coffee}\PY{+w}{ }\PY{p}{\PYZob{}}
\PY{+w}{    }\PY{k+kd}{public}\PY{+w}{ }\PY{n}{BigDecimal}\PY{+w}{ }\PY{n+nf}{cost}\PY{p}{(}\PY{p}{)}\PY{+w}{ }\PY{p}{\PYZob{}}\PY{+w}{ }\PY{k}{return}\PY{+w}{ }\PY{n}{BigDecimal}\PY{p}{.}\PY{n+na}{valueOf}\PY{p}{(}\PY{l+m+mf}{2.0}\PY{p}{)}\PY{p}{;}\PY{+w}{ }\PY{p}{\PYZcb{}}
\PY{+w}{    }\PY{k+kd}{public}\PY{+w}{ }\PY{n}{String}\PY{+w}{ }\PY{n+nf}{description}\PY{p}{(}\PY{p}{)}\PY{+w}{ }\PY{p}{\PYZob{}}\PY{+w}{ }\PY{k}{return}\PY{+w}{ }\PY{l+s}{\PYZdq{}}\PY{l+s}{Coffee}\PY{l+s}{\PYZdq{}}\PY{p}{;}\PY{+w}{ }\PY{p}{\PYZcb{}}
\PY{p}{\PYZcb{}}

\PY{k+kd}{abstract}\PY{+w}{ }\PY{k+kd}{class} \PY{n+nc}{CoffeeDecorator}\PY{+w}{ }\PY{k+kd}{implements}\PY{+w}{ }\PY{n}{Coffee}\PY{+w}{ }\PY{p}{\PYZob{}}
\PY{+w}{    }\PY{k+kd}{protected}\PY{+w}{ }\PY{k+kd}{final}\PY{+w}{ }\PY{n}{Coffee}\PY{+w}{ }\PY{n}{delegate}\PY{p}{;}
\PY{+w}{    }\PY{n}{CoffeeDecorator}\PY{p}{(}\PY{n}{Coffee}\PY{+w}{ }\PY{n}{delegate}\PY{p}{)}\PY{+w}{ }\PY{p}{\PYZob{}}\PY{+w}{ }\PY{k}{this}\PY{p}{.}\PY{n+na}{delegate}\PY{+w}{ }\PY{o}{=}\PY{+w}{ }\PY{n}{delegate}\PY{p}{;}\PY{+w}{ }\PY{p}{\PYZcb{}}
\PY{p}{\PYZcb{}}

\PY{k+kd}{class} \PY{n+nc}{MilkDecorator}\PY{+w}{ }\PY{k+kd}{extends}\PY{+w}{ }\PY{n}{CoffeeDecorator}\PY{+w}{ }\PY{p}{\PYZob{}}
\PY{+w}{    }\PY{n}{MilkDecorator}\PY{p}{(}\PY{n}{Coffee}\PY{+w}{ }\PY{n}{delegate}\PY{p}{)}\PY{+w}{ }\PY{p}{\PYZob{}}\PY{+w}{ }\PY{k+kd}{super}\PY{p}{(}\PY{n}{delegate}\PY{p}{)}\PY{p}{;}\PY{+w}{ }\PY{p}{\PYZcb{}}
\PY{+w}{    }\PY{k+kd}{public}\PY{+w}{ }\PY{n}{BigDecimal}\PY{+w}{ }\PY{n+nf}{cost}\PY{p}{(}\PY{p}{)}\PY{+w}{ }\PY{p}{\PYZob{}}\PY{+w}{ }\PY{k}{return}\PY{+w}{ }\PY{n}{delegate}\PY{p}{.}\PY{n+na}{cost}\PY{p}{(}\PY{p}{)}\PY{p}{.}\PY{n+na}{add}\PY{p}{(}\PY{n}{BigDecimal}\PY{p}{.}\PY{n+na}{valueOf}\PY{p}{(}\PY{l+m+mf}{0.5}\PY{p}{)}\PY{p}{)}\PY{p}{;}\PY{+w}{ }\PY{p}{\PYZcb{}}
\PY{+w}{    }\PY{k+kd}{public}\PY{+w}{ }\PY{n}{String}\PY{+w}{ }\PY{n+nf}{description}\PY{p}{(}\PY{p}{)}\PY{+w}{ }\PY{p}{\PYZob{}}\PY{+w}{ }\PY{k}{return}\PY{+w}{ }\PY{n}{delegate}\PY{p}{.}\PY{n+na}{description}\PY{p}{(}\PY{p}{)}\PY{+w}{ }\PY{o}{+}\PY{+w}{ }\PY{l+s}{\PYZdq{}}\PY{l+s}{ + Milk}\PY{l+s}{\PYZdq{}}\PY{p}{;}\PY{+w}{ }\PY{p}{\PYZcb{}}
\PY{p}{\PYZcb{}}

\PY{n}{Coffee}\PY{+w}{ }\PY{n}{order}\PY{+w}{ }\PY{o}{=}\PY{+w}{ }\PY{k}{new}\PY{+w}{ }\PY{n}{MilkDecorator}\PY{p}{(}\PY{k}{new}\PY{+w}{ }\PY{n}{SimpleCoffee}\PY{p}{(}\PY{p}{)}\PY{p}{)}\PY{p}{;}
\PY{n}{System}\PY{p}{.}\PY{n+na}{out}\PY{p}{.}\PY{n+na}{println}\PY{p}{(}\PY{n}{order}\PY{p}{.}\PY{n+na}{description}\PY{p}{(}\PY{p}{)}\PY{+w}{ }\PY{o}{+}\PY{+w}{ }\PY{l+s}{\PYZdq{}}\PY{l+s}{: \PYZdl{}}\PY{l+s}{\PYZdq{}}\PY{+w}{ }\PY{o}{+}\PY{+w}{ }\PY{n}{order}\PY{p}{.}\PY{n+na}{cost}\PY{p}{(}\PY{p}{)}\PY{p}{)}\PY{p}{;}\PY{+w}{ }\PY{c+c1}{// Coffee + Milk: \PYZdl{}2.5}
\end{Verbatim}
\end{codebox}

\subsection[Adapter]{Adapter}
\label{h:19:adapter}
\begin{codebox}
\begin{Verbatim}[commandchars=\\\{\}]
\PY{c+c1}{// Third\PYZhy{}party class we cannot change, with an incompatible interface}
\PY{k+kd}{class} \PY{n+nc}{LegacyXmlParser}\PY{+w}{ }\PY{p}{\PYZob{}}
\PY{+w}{    }\PY{n}{String}\PY{+w}{ }\PY{n+nf}{parseXml}\PY{p}{(}\PY{n}{String}\PY{+w}{ }\PY{n}{xml}\PY{p}{)}\PY{+w}{ }\PY{p}{\PYZob{}}\PY{+w}{ }\PY{k}{return}\PY{+w}{ }\PY{l+s}{\PYZdq{}}\PY{l+s}{parsed:}\PY{l+s}{\PYZdq{}}\PY{+w}{ }\PY{o}{+}\PY{+w}{ }\PY{n}{xml}\PY{p}{;}\PY{+w}{ }\PY{p}{\PYZcb{}}
\PY{p}{\PYZcb{}}

\PY{c+c1}{// The interface our code actually wants to use}
\PY{k+kd}{interface} \PY{n+nc}{JsonParser}\PY{+w}{ }\PY{p}{\PYZob{}}
\PY{+w}{    }\PY{n}{String}\PY{+w}{ }\PY{n+nf}{parseJson}\PY{p}{(}\PY{n}{String}\PY{+w}{ }\PY{n}{json}\PY{p}{)}\PY{p}{;}
\PY{p}{\PYZcb{}}

\PY{c+c1}{// Adapter bridges the two}
\PY{k+kd}{class} \PY{n+nc}{XmlToJsonAdapter}\PY{+w}{ }\PY{k+kd}{implements}\PY{+w}{ }\PY{n}{JsonParser}\PY{+w}{ }\PY{p}{\PYZob{}}
\PY{+w}{    }\PY{k+kd}{private}\PY{+w}{ }\PY{k+kd}{final}\PY{+w}{ }\PY{n}{LegacyXmlParser}\PY{+w}{ }\PY{n}{legacy}\PY{+w}{ }\PY{o}{=}\PY{+w}{ }\PY{k}{new}\PY{+w}{ }\PY{n}{LegacyXmlParser}\PY{p}{(}\PY{p}{)}\PY{p}{;}
\PY{+w}{    }\PY{k+kd}{public}\PY{+w}{ }\PY{n}{String}\PY{+w}{ }\PY{n+nf}{parseJson}\PY{p}{(}\PY{n}{String}\PY{+w}{ }\PY{n}{json}\PY{p}{)}\PY{+w}{ }\PY{p}{\PYZob{}}
\PY{+w}{        }\PY{n}{String}\PY{+w}{ }\PY{n}{xmlEquivalent}\PY{+w}{ }\PY{o}{=}\PY{+w}{ }\PY{n}{convertJsonToXml}\PY{p}{(}\PY{n}{json}\PY{p}{)}\PY{p}{;}
\PY{+w}{        }\PY{k}{return}\PY{+w}{ }\PY{n}{legacy}\PY{p}{.}\PY{n+na}{parseXml}\PY{p}{(}\PY{n}{xmlEquivalent}\PY{p}{)}\PY{p}{;}
\PY{+w}{    }\PY{p}{\PYZcb{}}
\PY{+w}{    }\PY{k+kd}{private}\PY{+w}{ }\PY{n}{String}\PY{+w}{ }\PY{n+nf}{convertJsonToXml}\PY{p}{(}\PY{n}{String}\PY{+w}{ }\PY{n}{json}\PY{p}{)}\PY{+w}{ }\PY{p}{\PYZob{}}\PY{+w}{ }\PY{k}{return}\PY{+w}{ }\PY{l+s}{\PYZdq{}}\PY{l+s}{\PYZlt{}xml\PYZgt{}}\PY{l+s}{\PYZdq{}}\PY{+w}{ }\PY{o}{+}\PY{+w}{ }\PY{n}{json}\PY{+w}{ }\PY{o}{+}\PY{+w}{ }\PY{l+s}{\PYZdq{}}\PY{l+s}{\PYZlt{}/xml\PYZgt{}}\PY{l+s}{\PYZdq{}}\PY{p}{;}\PY{+w}{ }\PY{p}{\PYZcb{}}
\PY{p}{\PYZcb{}}
\end{Verbatim}
\end{codebox}

\subsection[Facade]{Facade}
\label{h:19:facade}
\begin{codebox}
\begin{Verbatim}[commandchars=\\\{\}]
\PY{c+c1}{// A complex subsystem with many moving parts}
\PY{k+kd}{class} \PY{n+nc}{InventoryService}\PY{+w}{ }\PY{p}{\PYZob{}}\PY{+w}{ }\PY{k+kt}{boolean}\PY{+w}{ }\PY{n+nf}{reserve}\PY{p}{(}\PY{n}{String}\PY{+w}{ }\PY{n}{sku}\PY{p}{,}\PY{+w}{ }\PY{k+kt}{int}\PY{+w}{ }\PY{n}{qty}\PY{p}{)}\PY{+w}{ }\PY{p}{\PYZob{}}\PY{+w}{ }\PY{k}{return}\PY{+w}{ }\PY{k+kc}{true}\PY{p}{;}\PY{+w}{ }\PY{p}{\PYZcb{}}\PY{+w}{ }\PY{p}{\PYZcb{}}
\PY{k+kd}{class} \PY{n+nc}{PaymentService}\PY{+w}{ }\PY{p}{\PYZob{}}\PY{+w}{ }\PY{k+kt}{void}\PY{+w}{ }\PY{n+nf}{charge}\PY{p}{(}\PY{n}{String}\PY{+w}{ }\PY{n}{customerId}\PY{p}{,}\PY{+w}{ }\PY{n}{BigDecimal}\PY{+w}{ }\PY{n}{amount}\PY{p}{)}\PY{+w}{ }\PY{p}{\PYZob{}}\PY{p}{\PYZcb{}}\PY{+w}{ }\PY{p}{\PYZcb{}}
\PY{k+kd}{class} \PY{n+nc}{ShippingService}\PY{+w}{ }\PY{p}{\PYZob{}}\PY{+w}{ }\PY{k+kt}{void}\PY{+w}{ }\PY{n+nf}{schedule}\PY{p}{(}\PY{n}{String}\PY{+w}{ }\PY{n}{orderId}\PY{p}{)}\PY{+w}{ }\PY{p}{\PYZob{}}\PY{p}{\PYZcb{}}\PY{+w}{ }\PY{p}{\PYZcb{}}

\PY{c+c1}{// Facade: one simple entry point that hides the subsystem\PYZsq{}s complexity}
\PY{k+kd}{class} \PY{n+nc}{OrderFacade}\PY{+w}{ }\PY{p}{\PYZob{}}
\PY{+w}{    }\PY{k+kd}{private}\PY{+w}{ }\PY{k+kd}{final}\PY{+w}{ }\PY{n}{InventoryService}\PY{+w}{ }\PY{n}{inventory}\PY{+w}{ }\PY{o}{=}\PY{+w}{ }\PY{k}{new}\PY{+w}{ }\PY{n}{InventoryService}\PY{p}{(}\PY{p}{)}\PY{p}{;}
\PY{+w}{    }\PY{k+kd}{private}\PY{+w}{ }\PY{k+kd}{final}\PY{+w}{ }\PY{n}{PaymentService}\PY{+w}{ }\PY{n}{payment}\PY{+w}{ }\PY{o}{=}\PY{+w}{ }\PY{k}{new}\PY{+w}{ }\PY{n}{PaymentService}\PY{p}{(}\PY{p}{)}\PY{p}{;}
\PY{+w}{    }\PY{k+kd}{private}\PY{+w}{ }\PY{k+kd}{final}\PY{+w}{ }\PY{n}{ShippingService}\PY{+w}{ }\PY{n}{shipping}\PY{+w}{ }\PY{o}{=}\PY{+w}{ }\PY{k}{new}\PY{+w}{ }\PY{n}{ShippingService}\PY{p}{(}\PY{p}{)}\PY{p}{;}

\PY{+w}{    }\PY{k+kt}{void}\PY{+w}{ }\PY{n+nf}{placeOrder}\PY{p}{(}\PY{n}{String}\PY{+w}{ }\PY{n}{orderId}\PY{p}{,}\PY{+w}{ }\PY{n}{String}\PY{+w}{ }\PY{n}{sku}\PY{p}{,}\PY{+w}{ }\PY{k+kt}{int}\PY{+w}{ }\PY{n}{qty}\PY{p}{,}\PY{+w}{ }\PY{n}{String}\PY{+w}{ }\PY{n}{customerId}\PY{p}{,}\PY{+w}{ }\PY{n}{BigDecimal}\PY{+w}{ }\PY{n}{amount}\PY{p}{)}\PY{+w}{ }\PY{p}{\PYZob{}}
\PY{+w}{        }\PY{n}{inventory}\PY{p}{.}\PY{n+na}{reserve}\PY{p}{(}\PY{n}{sku}\PY{p}{,}\PY{+w}{ }\PY{n}{qty}\PY{p}{)}\PY{p}{;}
\PY{+w}{        }\PY{n}{payment}\PY{p}{.}\PY{n+na}{charge}\PY{p}{(}\PY{n}{customerId}\PY{p}{,}\PY{+w}{ }\PY{n}{amount}\PY{p}{)}\PY{p}{;}
\PY{+w}{        }\PY{n}{shipping}\PY{p}{.}\PY{n+na}{schedule}\PY{p}{(}\PY{n}{orderId}\PY{p}{)}\PY{p}{;}
\PY{+w}{    }\PY{p}{\PYZcb{}}
\PY{p}{\PYZcb{}}

\PY{k}{new}\PY{+w}{ }\PY{n}{OrderFacade}\PY{p}{(}\PY{p}{)}\PY{p}{.}\PY{n+na}{placeOrder}\PY{p}{(}\PY{l+s}{\PYZdq{}}\PY{l+s}{O\PYZhy{}1}\PY{l+s}{\PYZdq{}}\PY{p}{,}\PY{+w}{ }\PY{l+s}{\PYZdq{}}\PY{l+s}{SKU\PYZhy{}42}\PY{l+s}{\PYZdq{}}\PY{p}{,}\PY{+w}{ }\PY{l+m+mi}{1}\PY{p}{,}\PY{+w}{ }\PY{l+s}{\PYZdq{}}\PY{l+s}{C\PYZhy{}9}\PY{l+s}{\PYZdq{}}\PY{p}{,}\PY{+w}{ }\PY{n}{BigDecimal}\PY{p}{.}\PY{n+na}{valueOf}\PY{p}{(}\PY{l+m+mf}{29.99}\PY{p}{)}\PY{p}{)}\PY{p}{;}
\end{Verbatim}
\end{codebox}

\subsection[Proxy]{Proxy}
\label{h:19:proxy}
\begin{codebox}
\begin{Verbatim}[commandchars=\\\{\}]
\PY{k+kd}{interface} \PY{n+nc}{ImageLoader}\PY{+w}{ }\PY{p}{\PYZob{}}
\PY{+w}{    }\PY{k+kt}{void}\PY{+w}{ }\PY{n+nf}{display}\PY{p}{(}\PY{p}{)}\PY{p}{;}
\PY{p}{\PYZcb{}}

\PY{k+kd}{class} \PY{n+nc}{RealImage}\PY{+w}{ }\PY{k+kd}{implements}\PY{+w}{ }\PY{n}{ImageLoader}\PY{+w}{ }\PY{p}{\PYZob{}}
\PY{+w}{    }\PY{k+kd}{private}\PY{+w}{ }\PY{k+kd}{final}\PY{+w}{ }\PY{n}{String}\PY{+w}{ }\PY{n}{path}\PY{p}{;}
\PY{+w}{    }\PY{n}{RealImage}\PY{p}{(}\PY{n}{String}\PY{+w}{ }\PY{n}{path}\PY{p}{)}\PY{+w}{ }\PY{p}{\PYZob{}}\PY{+w}{ }\PY{k}{this}\PY{p}{.}\PY{n+na}{path}\PY{+w}{ }\PY{o}{=}\PY{+w}{ }\PY{n}{path}\PY{p}{;}\PY{+w}{ }\PY{n}{loadFromDisk}\PY{p}{(}\PY{p}{)}\PY{p}{;}\PY{+w}{ }\PY{p}{\PYZcb{}}
\PY{+w}{    }\PY{k+kd}{private}\PY{+w}{ }\PY{k+kt}{void}\PY{+w}{ }\PY{n+nf}{loadFromDisk}\PY{p}{(}\PY{p}{)}\PY{+w}{ }\PY{p}{\PYZob{}}\PY{+w}{ }\PY{n}{System}\PY{p}{.}\PY{n+na}{out}\PY{p}{.}\PY{n+na}{println}\PY{p}{(}\PY{l+s}{\PYZdq{}}\PY{l+s}{Loading }\PY{l+s}{\PYZdq{}}\PY{+w}{ }\PY{o}{+}\PY{+w}{ }\PY{n}{path}\PY{p}{)}\PY{p}{;}\PY{+w}{ }\PY{p}{\PYZcb{}}
\PY{+w}{    }\PY{k+kd}{public}\PY{+w}{ }\PY{k+kt}{void}\PY{+w}{ }\PY{n+nf}{display}\PY{p}{(}\PY{p}{)}\PY{+w}{ }\PY{p}{\PYZob{}}\PY{+w}{ }\PY{n}{System}\PY{p}{.}\PY{n+na}{out}\PY{p}{.}\PY{n+na}{println}\PY{p}{(}\PY{l+s}{\PYZdq{}}\PY{l+s}{Displaying }\PY{l+s}{\PYZdq{}}\PY{+w}{ }\PY{o}{+}\PY{+w}{ }\PY{n}{path}\PY{p}{)}\PY{p}{;}\PY{+w}{ }\PY{p}{\PYZcb{}}
\PY{p}{\PYZcb{}}

\PY{c+c1}{// Proxy: defers the expensive load until display() is actually called}
\PY{k+kd}{class} \PY{n+nc}{LazyImageProxy}\PY{+w}{ }\PY{k+kd}{implements}\PY{+w}{ }\PY{n}{ImageLoader}\PY{+w}{ }\PY{p}{\PYZob{}}
\PY{+w}{    }\PY{k+kd}{private}\PY{+w}{ }\PY{k+kd}{final}\PY{+w}{ }\PY{n}{String}\PY{+w}{ }\PY{n}{path}\PY{p}{;}
\PY{+w}{    }\PY{k+kd}{private}\PY{+w}{ }\PY{n}{RealImage}\PY{+w}{ }\PY{n}{real}\PY{p}{;}
\PY{+w}{    }\PY{n}{LazyImageProxy}\PY{p}{(}\PY{n}{String}\PY{+w}{ }\PY{n}{path}\PY{p}{)}\PY{+w}{ }\PY{p}{\PYZob{}}\PY{+w}{ }\PY{k}{this}\PY{p}{.}\PY{n+na}{path}\PY{+w}{ }\PY{o}{=}\PY{+w}{ }\PY{n}{path}\PY{p}{;}\PY{+w}{ }\PY{p}{\PYZcb{}}
\PY{+w}{    }\PY{k+kd}{public}\PY{+w}{ }\PY{k+kt}{void}\PY{+w}{ }\PY{n+nf}{display}\PY{p}{(}\PY{p}{)}\PY{+w}{ }\PY{p}{\PYZob{}}
\PY{+w}{        }\PY{k}{if}\PY{+w}{ }\PY{p}{(}\PY{n}{real}\PY{+w}{ }\PY{o}{=}\PY{o}{=}\PY{+w}{ }\PY{k+kc}{null}\PY{p}{)}\PY{+w}{ }\PY{p}{\PYZob{}}
\PY{+w}{            }\PY{n}{real}\PY{+w}{ }\PY{o}{=}\PY{+w}{ }\PY{k}{new}\PY{+w}{ }\PY{n}{RealImage}\PY{p}{(}\PY{n}{path}\PY{p}{)}\PY{p}{;}\PY{+w}{ }\PY{c+c1}{// loaded only on first use}
\PY{+w}{        }\PY{p}{\PYZcb{}}
\PY{+w}{        }\PY{n}{real}\PY{p}{.}\PY{n+na}{display}\PY{p}{(}\PY{p}{)}\PY{p}{;}
\PY{+w}{    }\PY{p}{\PYZcb{}}
\PY{p}{\PYZcb{}}
\end{Verbatim}
\end{codebox}

\subsection[Command]{Command}
\label{h:19:command}
\begin{codebox}
\begin{Verbatim}[commandchars=\\\{\}]
\PY{k+kd}{interface} \PY{n+nc}{Command}\PY{+w}{ }\PY{p}{\PYZob{}}
\PY{+w}{    }\PY{k+kt}{void}\PY{+w}{ }\PY{n+nf}{execute}\PY{p}{(}\PY{p}{)}\PY{p}{;}
\PY{p}{\PYZcb{}}

\PY{k+kd}{class} \PY{n+nc}{LightOnCommand}\PY{+w}{ }\PY{k+kd}{implements}\PY{+w}{ }\PY{n}{Command}\PY{+w}{ }\PY{p}{\PYZob{}}
\PY{+w}{    }\PY{k+kd}{public}\PY{+w}{ }\PY{k+kt}{void}\PY{+w}{ }\PY{n+nf}{execute}\PY{p}{(}\PY{p}{)}\PY{+w}{ }\PY{p}{\PYZob{}}\PY{+w}{ }\PY{n}{System}\PY{p}{.}\PY{n+na}{out}\PY{p}{.}\PY{n+na}{println}\PY{p}{(}\PY{l+s}{\PYZdq{}}\PY{l+s}{Light on}\PY{l+s}{\PYZdq{}}\PY{p}{)}\PY{p}{;}\PY{+w}{ }\PY{p}{\PYZcb{}}
\PY{p}{\PYZcb{}}

\PY{k+kd}{class} \PY{n+nc}{RemoteControl}\PY{+w}{ }\PY{p}{\PYZob{}}
\PY{+w}{    }\PY{k+kd}{private}\PY{+w}{ }\PY{k+kd}{final}\PY{+w}{ }\PY{n}{List}\PY{o}{\PYZlt{}}\PY{n}{Command}\PY{o}{\PYZgt{}}\PY{+w}{ }\PY{n}{history}\PY{+w}{ }\PY{o}{=}\PY{+w}{ }\PY{k}{new}\PY{+w}{ }\PY{n}{ArrayList}\PY{o}{\PYZlt{}}\PY{o}{\PYZgt{}}\PY{p}{(}\PY{p}{)}\PY{p}{;}
\PY{+w}{    }\PY{k+kt}{void}\PY{+w}{ }\PY{n+nf}{submit}\PY{p}{(}\PY{n}{Command}\PY{+w}{ }\PY{n}{command}\PY{p}{)}\PY{+w}{ }\PY{p}{\PYZob{}}
\PY{+w}{        }\PY{n}{command}\PY{p}{.}\PY{n+na}{execute}\PY{p}{(}\PY{p}{)}\PY{p}{;}
\PY{+w}{        }\PY{n}{history}\PY{p}{.}\PY{n+na}{add}\PY{p}{(}\PY{n}{command}\PY{p}{)}\PY{p}{;}
\PY{+w}{    }\PY{p}{\PYZcb{}}
\PY{p}{\PYZcb{}}

\PY{c+c1}{// Classic: a class implementing Command}
\PY{k}{new}\PY{+w}{ }\PY{n}{RemoteControl}\PY{p}{(}\PY{p}{)}\PY{p}{.}\PY{n+na}{submit}\PY{p}{(}\PY{k}{new}\PY{+w}{ }\PY{n}{LightOnCommand}\PY{p}{(}\PY{p}{)}\PY{p}{)}\PY{p}{;}

\PY{c+c1}{// Modern replacement: Command is a functional interface, lambda works directly}
\PY{k}{new}\PY{+w}{ }\PY{n}{RemoteControl}\PY{p}{(}\PY{p}{)}\PY{p}{.}\PY{n+na}{submit}\PY{p}{(}\PY{p}{(}\PY{p}{)}\PY{+w}{ }\PY{o}{\PYZhy{}}\PY{o}{\PYZgt{}}\PY{+w}{ }\PY{n}{System}\PY{p}{.}\PY{n+na}{out}\PY{p}{.}\PY{n+na}{println}\PY{p}{(}\PY{l+s}{\PYZdq{}}\PY{l+s}{Light on}\PY{l+s}{\PYZdq{}}\PY{p}{)}\PY{p}{)}\PY{p}{;}
\end{Verbatim}
\end{codebox}

\subsection[Chain of Responsibility]{Chain of Responsibility}
\label{h:19:chain-of-responsibility}
\begin{codebox}
\begin{Verbatim}[commandchars=\\\{\}]
\PY{k+kd}{abstract}\PY{+w}{ }\PY{k+kd}{class} \PY{n+nc}{SupportHandler}\PY{+w}{ }\PY{p}{\PYZob{}}
\PY{+w}{    }\PY{k+kd}{protected}\PY{+w}{ }\PY{n}{SupportHandler}\PY{+w}{ }\PY{n}{next}\PY{p}{;}
\PY{+w}{    }\PY{n}{SupportHandler}\PY{+w}{ }\PY{n+nf}{setNext}\PY{p}{(}\PY{n}{SupportHandler}\PY{+w}{ }\PY{n}{next}\PY{p}{)}\PY{+w}{ }\PY{p}{\PYZob{}}\PY{+w}{ }\PY{k}{this}\PY{p}{.}\PY{n+na}{next}\PY{+w}{ }\PY{o}{=}\PY{+w}{ }\PY{n}{next}\PY{p}{;}\PY{+w}{ }\PY{k}{return}\PY{+w}{ }\PY{n}{next}\PY{p}{;}\PY{+w}{ }\PY{p}{\PYZcb{}}
\PY{+w}{    }\PY{k+kd}{abstract}\PY{+w}{ }\PY{k+kt}{void}\PY{+w}{ }\PY{n+nf}{handle}\PY{p}{(}\PY{n}{String}\PY{+w}{ }\PY{n}{issue}\PY{p}{)}\PY{p}{;}
\PY{p}{\PYZcb{}}

\PY{k+kd}{class} \PY{n+nc}{Tier1Support}\PY{+w}{ }\PY{k+kd}{extends}\PY{+w}{ }\PY{n}{SupportHandler}\PY{+w}{ }\PY{p}{\PYZob{}}
\PY{+w}{    }\PY{k+kt}{void}\PY{+w}{ }\PY{n+nf}{handle}\PY{p}{(}\PY{n}{String}\PY{+w}{ }\PY{n}{issue}\PY{p}{)}\PY{+w}{ }\PY{p}{\PYZob{}}
\PY{+w}{        }\PY{k}{if}\PY{+w}{ }\PY{p}{(}\PY{n}{issue}\PY{p}{.}\PY{n+na}{equals}\PY{p}{(}\PY{l+s}{\PYZdq{}}\PY{l+s}{password\PYZhy{}reset}\PY{l+s}{\PYZdq{}}\PY{p}{)}\PY{p}{)}\PY{+w}{ }\PY{p}{\PYZob{}}
\PY{+w}{            }\PY{n}{System}\PY{p}{.}\PY{n+na}{out}\PY{p}{.}\PY{n+na}{println}\PY{p}{(}\PY{l+s}{\PYZdq{}}\PY{l+s}{Tier1 handled: }\PY{l+s}{\PYZdq{}}\PY{+w}{ }\PY{o}{+}\PY{+w}{ }\PY{n}{issue}\PY{p}{)}\PY{p}{;}
\PY{+w}{        }\PY{p}{\PYZcb{}}\PY{+w}{ }\PY{k}{else}\PY{+w}{ }\PY{k}{if}\PY{+w}{ }\PY{p}{(}\PY{n}{next}\PY{+w}{ }\PY{o}{!}\PY{o}{=}\PY{+w}{ }\PY{k+kc}{null}\PY{p}{)}\PY{+w}{ }\PY{p}{\PYZob{}}
\PY{+w}{            }\PY{n}{next}\PY{p}{.}\PY{n+na}{handle}\PY{p}{(}\PY{n}{issue}\PY{p}{)}\PY{p}{;}
\PY{+w}{        }\PY{p}{\PYZcb{}}
\PY{+w}{    }\PY{p}{\PYZcb{}}
\PY{p}{\PYZcb{}}

\PY{k+kd}{class} \PY{n+nc}{Tier2Support}\PY{+w}{ }\PY{k+kd}{extends}\PY{+w}{ }\PY{n}{SupportHandler}\PY{+w}{ }\PY{p}{\PYZob{}}
\PY{+w}{    }\PY{k+kt}{void}\PY{+w}{ }\PY{n+nf}{handle}\PY{p}{(}\PY{n}{String}\PY{+w}{ }\PY{n}{issue}\PY{p}{)}\PY{+w}{ }\PY{p}{\PYZob{}}
\PY{+w}{        }\PY{n}{System}\PY{p}{.}\PY{n+na}{out}\PY{p}{.}\PY{n+na}{println}\PY{p}{(}\PY{l+s}{\PYZdq{}}\PY{l+s}{Tier2 handled: }\PY{l+s}{\PYZdq{}}\PY{+w}{ }\PY{o}{+}\PY{+w}{ }\PY{n}{issue}\PY{p}{)}\PY{p}{;}
\PY{+w}{    }\PY{p}{\PYZcb{}}
\PY{p}{\PYZcb{}}

\PY{n}{SupportHandler}\PY{+w}{ }\PY{n}{chain}\PY{+w}{ }\PY{o}{=}\PY{+w}{ }\PY{k}{new}\PY{+w}{ }\PY{n}{Tier1Support}\PY{p}{(}\PY{p}{)}\PY{p}{;}
\PY{n}{chain}\PY{p}{.}\PY{n+na}{setNext}\PY{p}{(}\PY{k}{new}\PY{+w}{ }\PY{n}{Tier2Support}\PY{p}{(}\PY{p}{)}\PY{p}{)}\PY{p}{;}
\PY{n}{chain}\PY{p}{.}\PY{n+na}{handle}\PY{p}{(}\PY{l+s}{\PYZdq{}}\PY{l+s}{database\PYZhy{}outage}\PY{l+s}{\PYZdq{}}\PY{p}{)}\PY{p}{;}\PY{+w}{ }\PY{c+c1}{// falls through to Tier2Support}
\end{Verbatim}
\end{codebox}

\subsection[Iterator]{Iterator}
\label{h:19:iterator}
\begin{codebox}
\begin{Verbatim}[commandchars=\\\{\}]
\PY{k+kd}{class} \PY{n+nc}{NameCollection}\PY{+w}{ }\PY{k+kd}{implements}\PY{+w}{ }\PY{n}{Iterable}\PY{o}{\PYZlt{}}\PY{n}{String}\PY{o}{\PYZgt{}}\PY{+w}{ }\PY{p}{\PYZob{}}
\PY{+w}{    }\PY{k+kd}{private}\PY{+w}{ }\PY{k+kd}{final}\PY{+w}{ }\PY{n}{String}\PY{o}{[}\PY{o}{]}\PY{+w}{ }\PY{n}{names}\PY{p}{;}
\PY{+w}{    }\PY{n}{NameCollection}\PY{p}{(}\PY{n}{String}\PY{p}{.}\PY{p}{.}\PY{p}{.}\PY{+w}{ }\PY{n}{names}\PY{p}{)}\PY{+w}{ }\PY{p}{\PYZob{}}\PY{+w}{ }\PY{k}{this}\PY{p}{.}\PY{n+na}{names}\PY{+w}{ }\PY{o}{=}\PY{+w}{ }\PY{n}{names}\PY{p}{;}\PY{+w}{ }\PY{p}{\PYZcb{}}

\PY{+w}{    }\PY{k+kd}{public}\PY{+w}{ }\PY{n}{Iterator}\PY{o}{\PYZlt{}}\PY{n}{String}\PY{o}{\PYZgt{}}\PY{+w}{ }\PY{n+nf}{iterator}\PY{p}{(}\PY{p}{)}\PY{+w}{ }\PY{p}{\PYZob{}}
\PY{+w}{        }\PY{k}{return}\PY{+w}{ }\PY{k}{new}\PY{+w}{ }\PY{n}{Iterator}\PY{o}{\PYZlt{}}\PY{o}{\PYZgt{}}\PY{p}{(}\PY{p}{)}\PY{+w}{ }\PY{p}{\PYZob{}}
\PY{+w}{            }\PY{k+kd}{private}\PY{+w}{ }\PY{k+kt}{int}\PY{+w}{ }\PY{n}{index}\PY{+w}{ }\PY{o}{=}\PY{+w}{ }\PY{l+m+mi}{0}\PY{p}{;}
\PY{+w}{            }\PY{k+kd}{public}\PY{+w}{ }\PY{k+kt}{boolean}\PY{+w}{ }\PY{n+nf}{hasNext}\PY{p}{(}\PY{p}{)}\PY{+w}{ }\PY{p}{\PYZob{}}\PY{+w}{ }\PY{k}{return}\PY{+w}{ }\PY{n}{index}\PY{+w}{ }\PY{o}{\PYZlt{}}\PY{+w}{ }\PY{n}{names}\PY{p}{.}\PY{n+na}{length}\PY{p}{;}\PY{+w}{ }\PY{p}{\PYZcb{}}
\PY{+w}{            }\PY{k+kd}{public}\PY{+w}{ }\PY{n}{String}\PY{+w}{ }\PY{n+nf}{next}\PY{p}{(}\PY{p}{)}\PY{+w}{ }\PY{p}{\PYZob{}}\PY{+w}{ }\PY{k}{return}\PY{+w}{ }\PY{n}{names}\PY{o}{[}\PY{n}{index}\PY{o}{+}\PY{o}{+}\PY{o}{]}\PY{p}{;}\PY{+w}{ }\PY{p}{\PYZcb{}}
\PY{+w}{        }\PY{p}{\PYZcb{}}\PY{p}{;}
\PY{+w}{    }\PY{p}{\PYZcb{}}
\PY{p}{\PYZcb{}}

\PY{k}{for}\PY{+w}{ }\PY{p}{(}\PY{n}{String}\PY{+w}{ }\PY{n}{name}\PY{+w}{ }\PY{p}{:}\PY{+w}{ }\PY{k}{new}\PY{+w}{ }\PY{n}{NameCollection}\PY{p}{(}\PY{l+s}{\PYZdq{}}\PY{l+s}{Ana}\PY{l+s}{\PYZdq{}}\PY{p}{,}\PY{+w}{ }\PY{l+s}{\PYZdq{}}\PY{l+s}{Bo}\PY{l+s}{\PYZdq{}}\PY{p}{,}\PY{+w}{ }\PY{l+s}{\PYZdq{}}\PY{l+s}{Cy}\PY{l+s}{\PYZdq{}}\PY{p}{)}\PY{p}{)}\PY{+w}{ }\PY{p}{\PYZob{}}
\PY{+w}{    }\PY{n}{System}\PY{p}{.}\PY{n+na}{out}\PY{p}{.}\PY{n+na}{println}\PY{p}{(}\PY{n}{name}\PY{p}{)}\PY{p}{;}\PY{+w}{ }\PY{c+c1}{// for\PYZhy{}each works because we implement Iterable}
\PY{p}{\PYZcb{}}
\end{Verbatim}
\end{codebox}

\subsection[State]{State}
\label{h:19:state}
\begin{codebox}
\begin{Verbatim}[commandchars=\\\{\}]
\PY{c+c1}{// Modern approach: sealed interface + exhaustive switch pattern matching}
\PY{k+kd}{sealed}\PY{+w}{ }\PY{k+kd}{interface} \PY{n+nc}{OrderState}\PY{+w}{ }\PY{n}{permits}\PY{+w}{ }\PY{n}{Placed}\PY{p}{,}\PY{+w}{ }\PY{n}{Shipped}\PY{p}{,}\PY{+w}{ }\PY{n}{Delivered}\PY{+w}{ }\PY{p}{\PYZob{}}\PY{p}{\PYZcb{}}
\PY{k+kd}{record} \PY{n+nc}{Placed}\PY{p}{(}\PY{p}{)}\PY{+w}{ }\PY{k+kd}{implements}\PY{+w}{ }\PY{n}{OrderState}\PY{+w}{ }\PY{p}{\PYZob{}}\PY{p}{\PYZcb{}}
\PY{k+kd}{record} \PY{n+nc}{Shipped}\PY{p}{(}\PY{n}{String}\PY{+w}{ }\PY{n}{trackingNumber}\PY{p}{)}\PY{+w}{ }\PY{k+kd}{implements}\PY{+w}{ }\PY{n}{OrderState}\PY{+w}{ }\PY{p}{\PYZob{}}\PY{p}{\PYZcb{}}
\PY{k+kd}{record} \PY{n+nc}{Delivered}\PY{p}{(}\PY{n}{LocalDate}\PY{+w}{ }\PY{n}{date}\PY{p}{)}\PY{+w}{ }\PY{k+kd}{implements}\PY{+w}{ }\PY{n}{OrderState}\PY{+w}{ }\PY{p}{\PYZob{}}\PY{p}{\PYZcb{}}

\PY{n}{String}\PY{+w}{ }\PY{n+nf}{describe}\PY{p}{(}\PY{n}{OrderState}\PY{+w}{ }\PY{n}{state}\PY{p}{)}\PY{+w}{ }\PY{p}{\PYZob{}}
\PY{+w}{    }\PY{k}{return}\PY{+w}{ }\PY{k}{switch}\PY{+w}{ }\PY{p}{(}\PY{n}{state}\PY{p}{)}\PY{+w}{ }\PY{p}{\PYZob{}}
\PY{+w}{        }\PY{k}{case}\PY{+w}{ }\PY{n}{Placed}\PY{+w}{ }\PY{n}{p}\PY{+w}{ }\PY{o}{\PYZhy{}}\PY{o}{\PYZgt{}}\PY{+w}{ }\PY{l+s}{\PYZdq{}}\PY{l+s}{Order placed}\PY{l+s}{\PYZdq{}}\PY{p}{;}
\PY{+w}{        }\PY{k}{case}\PY{+w}{ }\PY{n}{Shipped}\PY{+w}{ }\PY{n}{s}\PY{+w}{ }\PY{o}{\PYZhy{}}\PY{o}{\PYZgt{}}\PY{+w}{ }\PY{l+s}{\PYZdq{}}\PY{l+s}{Shipped, tracking: }\PY{l+s}{\PYZdq{}}\PY{+w}{ }\PY{o}{+}\PY{+w}{ }\PY{n}{s}\PY{p}{.}\PY{n+na}{trackingNumber}\PY{p}{(}\PY{p}{)}\PY{p}{;}
\PY{+w}{        }\PY{k}{case}\PY{+w}{ }\PY{n}{Delivered}\PY{+w}{ }\PY{n}{d}\PY{+w}{ }\PY{o}{\PYZhy{}}\PY{o}{\PYZgt{}}\PY{+w}{ }\PY{l+s}{\PYZdq{}}\PY{l+s}{Delivered on }\PY{l+s}{\PYZdq{}}\PY{+w}{ }\PY{o}{+}\PY{+w}{ }\PY{n}{d}\PY{p}{.}\PY{n+na}{date}\PY{p}{(}\PY{p}{)}\PY{p}{;}
\PY{+w}{        }\PY{c+c1}{// no default needed — the compiler knows these are the only three states}
\PY{+w}{    }\PY{p}{\PYZcb{}}\PY{p}{;}
\PY{p}{\PYZcb{}}
\end{Verbatim}
\end{codebox}

\subsection[Visitor — Replaced by Sealed Interfaces + Pattern Matching]{Visitor — Replaced by Sealed Interfaces + Pattern Matching}
\label{h:19:visitor--replaced-by-sealed-interfaces--pattern-matching}
The classic Visitor pattern exists to let you add new \emph{operations} over a fixed set of types without editing those types (double-dispatch through \code{accept}/\code{visit} methods). In modern Java, if the set of types is genuinely closed, a sealed hierarchy plus an exhaustive \code{switch} achieves the same goal with far less boilerplate — and the compiler enforces exhaustiveness for you.

\begin{codebox}
\begin{Verbatim}[commandchars=\\\{\}]
\PY{k+kd}{sealed}\PY{+w}{ }\PY{k+kd}{interface} \PY{n+nc}{Shape}\PY{+w}{ }\PY{n}{permits}\PY{+w}{ }\PY{n}{Circle}\PY{p}{,}\PY{+w}{ }\PY{n}{Square}\PY{p}{,}\PY{+w}{ }\PY{n}{Triangle}\PY{+w}{ }\PY{p}{\PYZob{}}\PY{p}{\PYZcb{}}
\PY{k+kd}{record} \PY{n+nc}{Circle}\PY{p}{(}\PY{k+kt}{double}\PY{+w}{ }\PY{n}{radius}\PY{p}{)}\PY{+w}{ }\PY{k+kd}{implements}\PY{+w}{ }\PY{n}{Shape}\PY{+w}{ }\PY{p}{\PYZob{}}\PY{p}{\PYZcb{}}
\PY{k+kd}{record} \PY{n+nc}{Square}\PY{p}{(}\PY{k+kt}{double}\PY{+w}{ }\PY{n}{side}\PY{p}{)}\PY{+w}{ }\PY{k+kd}{implements}\PY{+w}{ }\PY{n}{Shape}\PY{+w}{ }\PY{p}{\PYZob{}}\PY{p}{\PYZcb{}}
\PY{k+kd}{record} \PY{n+nc}{Triangle}\PY{p}{(}\PY{k+kt}{double}\PY{+w}{ }\PY{n}{base}\PY{p}{,}\PY{+w}{ }\PY{k+kt}{double}\PY{+w}{ }\PY{n}{height}\PY{p}{)}\PY{+w}{ }\PY{k+kd}{implements}\PY{+w}{ }\PY{n}{Shape}\PY{+w}{ }\PY{p}{\PYZob{}}\PY{p}{\PYZcb{}}

\PY{k+kt}{double}\PY{+w}{ }\PY{n+nf}{area}\PY{p}{(}\PY{n}{Shape}\PY{+w}{ }\PY{n}{shape}\PY{p}{)}\PY{+w}{ }\PY{p}{\PYZob{}}
\PY{+w}{    }\PY{k}{return}\PY{+w}{ }\PY{k}{switch}\PY{+w}{ }\PY{p}{(}\PY{n}{shape}\PY{p}{)}\PY{+w}{ }\PY{p}{\PYZob{}}
\PY{+w}{        }\PY{k}{case}\PY{+w}{ }\PY{n}{Circle}\PY{+w}{ }\PY{n}{c}\PY{+w}{ }\PY{o}{\PYZhy{}}\PY{o}{\PYZgt{}}\PY{+w}{ }\PY{n}{Math}\PY{p}{.}\PY{n+na}{PI}\PY{+w}{ }\PY{o}{*}\PY{+w}{ }\PY{n}{c}\PY{p}{.}\PY{n+na}{radius}\PY{p}{(}\PY{p}{)}\PY{+w}{ }\PY{o}{*}\PY{+w}{ }\PY{n}{c}\PY{p}{.}\PY{n+na}{radius}\PY{p}{(}\PY{p}{)}\PY{p}{;}
\PY{+w}{        }\PY{k}{case}\PY{+w}{ }\PY{n}{Square}\PY{+w}{ }\PY{n}{s}\PY{+w}{ }\PY{o}{\PYZhy{}}\PY{o}{\PYZgt{}}\PY{+w}{ }\PY{n}{s}\PY{p}{.}\PY{n+na}{side}\PY{p}{(}\PY{p}{)}\PY{+w}{ }\PY{o}{*}\PY{+w}{ }\PY{n}{s}\PY{p}{.}\PY{n+na}{side}\PY{p}{(}\PY{p}{)}\PY{p}{;}
\PY{+w}{        }\PY{k}{case}\PY{+w}{ }\PY{n}{Triangle}\PY{+w}{ }\PY{n}{t}\PY{+w}{ }\PY{o}{\PYZhy{}}\PY{o}{\PYZgt{}}\PY{+w}{ }\PY{l+m+mf}{0.5}\PY{+w}{ }\PY{o}{*}\PY{+w}{ }\PY{n}{t}\PY{p}{.}\PY{n+na}{base}\PY{p}{(}\PY{p}{)}\PY{+w}{ }\PY{o}{*}\PY{+w}{ }\PY{n}{t}\PY{p}{.}\PY{n+na}{height}\PY{p}{(}\PY{p}{)}\PY{p}{;}
\PY{+w}{    }\PY{p}{\PYZcb{}}\PY{p}{;}
\PY{p}{\PYZcb{}}

\PY{c+c1}{// Adding a new operation (e.g. perimeter) is just a new function with its own switch —}
\PY{c+c1}{// no accept()/visit() plumbing needed on Shape itself.}
\end{Verbatim}
\end{codebox}

\section[API Design Best Practices]{API Design Best Practices}
\label{h:19:api-design-best-practices}
Designing a public API — a library, a service interface, a set of public classes — is different from writing internal code: once other people depend on it, every change is a potential breaking change.

\subsection[Naming]{Naming}
\label{h:19:naming}
Follow existing JDK conventions so your API feels familiar: \code{getX}/\code{isX} for accessors, verbs for actions, plural nouns for collections.

\begin{codebox}
\begin{Verbatim}[commandchars=\\\{\}]
\PY{c+c1}{// Inconsistent, surprising names}
\PY{k+kd}{class} \PY{n+nc}{UserApi}\PY{+w}{ }\PY{p}{\PYZob{}}
\PY{+w}{    }\PY{n}{User}\PY{+w}{ }\PY{n+nf}{fetchUser}\PY{p}{(}\PY{n}{String}\PY{+w}{ }\PY{n}{id}\PY{p}{)}\PY{+w}{ }\PY{p}{\PYZob{}}\PY{+w}{ }\PY{p}{.}\PY{p}{.}\PY{p}{.}\PY{+w}{ }\PY{p}{\PYZcb{}}\PY{+w}{      }\PY{c+c1}{// getUser would match convention}
\PY{+w}{    }\PY{k+kt}{boolean}\PY{+w}{ }\PY{n+nf}{userActive}\PY{p}{(}\PY{n}{String}\PY{+w}{ }\PY{n}{id}\PY{p}{)}\PY{+w}{ }\PY{p}{\PYZob{}}\PY{+w}{ }\PY{p}{.}\PY{p}{.}\PY{p}{.}\PY{+w}{ }\PY{p}{\PYZcb{}}\PY{+w}{  }\PY{c+c1}{// isUserActive reads as a question}
\PY{+w}{    }\PY{n}{List}\PY{o}{\PYZlt{}}\PY{n}{User}\PY{o}{\PYZgt{}}\PY{+w}{ }\PY{n+nf}{userList}\PY{p}{(}\PY{p}{)}\PY{+w}{ }\PY{p}{\PYZob{}}\PY{+w}{ }\PY{p}{.}\PY{p}{.}\PY{p}{.}\PY{+w}{ }\PY{p}{\PYZcb{}}\PY{+w}{          }\PY{c+c1}{// users() is more idiomatic}
\PY{p}{\PYZcb{}}

\PY{c+c1}{// Convention\PYZhy{}following names}
\PY{k+kd}{class} \PY{n+nc}{UserApi}\PY{+w}{ }\PY{p}{\PYZob{}}
\PY{+w}{    }\PY{n}{User}\PY{+w}{ }\PY{n+nf}{getUser}\PY{p}{(}\PY{n}{String}\PY{+w}{ }\PY{n}{id}\PY{p}{)}\PY{+w}{ }\PY{p}{\PYZob{}}\PY{+w}{ }\PY{p}{.}\PY{p}{.}\PY{p}{.}\PY{+w}{ }\PY{p}{\PYZcb{}}
\PY{+w}{    }\PY{k+kt}{boolean}\PY{+w}{ }\PY{n+nf}{isUserActive}\PY{p}{(}\PY{n}{String}\PY{+w}{ }\PY{n}{id}\PY{p}{)}\PY{+w}{ }\PY{p}{\PYZob{}}\PY{+w}{ }\PY{p}{.}\PY{p}{.}\PY{p}{.}\PY{+w}{ }\PY{p}{\PYZcb{}}
\PY{+w}{    }\PY{n}{List}\PY{o}{\PYZlt{}}\PY{n}{User}\PY{o}{\PYZgt{}}\PY{+w}{ }\PY{n+nf}{users}\PY{p}{(}\PY{p}{)}\PY{+w}{ }\PY{p}{\PYZob{}}\PY{+w}{ }\PY{p}{.}\PY{p}{.}\PY{p}{.}\PY{+w}{ }\PY{p}{\PYZcb{}}
\PY{p}{\PYZcb{}}
\end{Verbatim}
\end{codebox}

\subsection[Minimal Surface Area]{Minimal Surface Area}
\label{h:19:minimal-surface-area}
Every public method, class, and field is a promise to callers. Expose only what is needed; keep everything else package-private or private. A smaller API is easier to learn, test, and evolve.

\begin{codebox}
\begin{Verbatim}[commandchars=\\\{\}]
\PY{c+c1}{// Before: implementation details leak into the public API}
\PY{k+kd}{public}\PY{+w}{ }\PY{k+kd}{class} \PY{n+nc}{OrderProcessor}\PY{+w}{ }\PY{p}{\PYZob{}}
\PY{+w}{    }\PY{k+kd}{public}\PY{+w}{ }\PY{k+kt}{void}\PY{+w}{ }\PY{n+nf}{validateOrder}\PY{p}{(}\PY{n}{Order}\PY{+w}{ }\PY{n}{order}\PY{p}{)}\PY{+w}{ }\PY{p}{\PYZob{}}\PY{+w}{ }\PY{p}{.}\PY{p}{.}\PY{p}{.}\PY{+w}{ }\PY{p}{\PYZcb{}}
\PY{+w}{    }\PY{k+kd}{public}\PY{+w}{ }\PY{k+kt}{void}\PY{+w}{ }\PY{n+nf}{calculateTax}\PY{p}{(}\PY{n}{Order}\PY{+w}{ }\PY{n}{order}\PY{p}{)}\PY{+w}{ }\PY{p}{\PYZob{}}\PY{+w}{ }\PY{p}{.}\PY{p}{.}\PY{p}{.}\PY{+w}{ }\PY{p}{\PYZcb{}}
\PY{+w}{    }\PY{k+kd}{public}\PY{+w}{ }\PY{k+kt}{void}\PY{+w}{ }\PY{n+nf}{chargeCard}\PY{p}{(}\PY{n}{Order}\PY{+w}{ }\PY{n}{order}\PY{p}{)}\PY{+w}{ }\PY{p}{\PYZob{}}\PY{+w}{ }\PY{p}{.}\PY{p}{.}\PY{p}{.}\PY{+w}{ }\PY{p}{\PYZcb{}}
\PY{+w}{    }\PY{k+kd}{public}\PY{+w}{ }\PY{k+kt}{void}\PY{+w}{ }\PY{n+nf}{process}\PY{p}{(}\PY{n}{Order}\PY{+w}{ }\PY{n}{order}\PY{p}{)}\PY{+w}{ }\PY{p}{\PYZob{}}
\PY{+w}{        }\PY{n}{validateOrder}\PY{p}{(}\PY{n}{order}\PY{p}{)}\PY{p}{;}\PY{+w}{ }\PY{n}{calculateTax}\PY{p}{(}\PY{n}{order}\PY{p}{)}\PY{p}{;}\PY{+w}{ }\PY{n}{chargeCard}\PY{p}{(}\PY{n}{order}\PY{p}{)}\PY{p}{;}
\PY{+w}{    }\PY{p}{\PYZcb{}}
\PY{p}{\PYZcb{}}

\PY{c+c1}{// After: only the one operation callers need is public}
\PY{k+kd}{public}\PY{+w}{ }\PY{k+kd}{class} \PY{n+nc}{OrderProcessor}\PY{+w}{ }\PY{p}{\PYZob{}}
\PY{+w}{    }\PY{k+kd}{public}\PY{+w}{ }\PY{k+kt}{void}\PY{+w}{ }\PY{n+nf}{process}\PY{p}{(}\PY{n}{Order}\PY{+w}{ }\PY{n}{order}\PY{p}{)}\PY{+w}{ }\PY{p}{\PYZob{}}
\PY{+w}{        }\PY{n}{validateOrder}\PY{p}{(}\PY{n}{order}\PY{p}{)}\PY{p}{;}\PY{+w}{ }\PY{n}{calculateTax}\PY{p}{(}\PY{n}{order}\PY{p}{)}\PY{p}{;}\PY{+w}{ }\PY{n}{chargeCard}\PY{p}{(}\PY{n}{order}\PY{p}{)}\PY{p}{;}
\PY{+w}{    }\PY{p}{\PYZcb{}}
\PY{+w}{    }\PY{k+kd}{private}\PY{+w}{ }\PY{k+kt}{void}\PY{+w}{ }\PY{n+nf}{validateOrder}\PY{p}{(}\PY{n}{Order}\PY{+w}{ }\PY{n}{order}\PY{p}{)}\PY{+w}{ }\PY{p}{\PYZob{}}\PY{+w}{ }\PY{p}{.}\PY{p}{.}\PY{p}{.}\PY{+w}{ }\PY{p}{\PYZcb{}}
\PY{+w}{    }\PY{k+kd}{private}\PY{+w}{ }\PY{k+kt}{void}\PY{+w}{ }\PY{n+nf}{calculateTax}\PY{p}{(}\PY{n}{Order}\PY{+w}{ }\PY{n}{order}\PY{p}{)}\PY{+w}{ }\PY{p}{\PYZob{}}\PY{+w}{ }\PY{p}{.}\PY{p}{.}\PY{p}{.}\PY{+w}{ }\PY{p}{\PYZcb{}}
\PY{+w}{    }\PY{k+kd}{private}\PY{+w}{ }\PY{k+kt}{void}\PY{+w}{ }\PY{n+nf}{chargeCard}\PY{p}{(}\PY{n}{Order}\PY{+w}{ }\PY{n}{order}\PY{p}{)}\PY{+w}{ }\PY{p}{\PYZob{}}\PY{+w}{ }\PY{p}{.}\PY{p}{.}\PY{p}{.}\PY{+w}{ }\PY{p}{\PYZcb{}}
\PY{p}{\PYZcb{}}
\end{Verbatim}
\end{codebox}

\subsection[Immutability by Default]{Immutability by Default}
\label{h:19:immutability-by-default}
Return types and parameters exposed by a public API should default to immutable unless there is a specific reason for mutability. This prevents a caller from mutating internal state through a returned reference.

\begin{codebox}
\begin{Verbatim}[commandchars=\\\{\}]
\PY{c+c1}{// Before: caller can mutate the list backing your internal state}
\PY{k+kd}{public}\PY{+w}{ }\PY{k+kd}{class} \PY{n+nc}{Team}\PY{+w}{ }\PY{p}{\PYZob{}}
\PY{+w}{    }\PY{k+kd}{private}\PY{+w}{ }\PY{k+kd}{final}\PY{+w}{ }\PY{n}{List}\PY{o}{\PYZlt{}}\PY{n}{String}\PY{o}{\PYZgt{}}\PY{+w}{ }\PY{n}{members}\PY{+w}{ }\PY{o}{=}\PY{+w}{ }\PY{k}{new}\PY{+w}{ }\PY{n}{ArrayList}\PY{o}{\PYZlt{}}\PY{o}{\PYZgt{}}\PY{p}{(}\PY{p}{)}\PY{p}{;}
\PY{+w}{    }\PY{k+kd}{public}\PY{+w}{ }\PY{n}{List}\PY{o}{\PYZlt{}}\PY{n}{String}\PY{o}{\PYZgt{}}\PY{+w}{ }\PY{n+nf}{getMembers}\PY{p}{(}\PY{p}{)}\PY{+w}{ }\PY{p}{\PYZob{}}\PY{+w}{ }\PY{k}{return}\PY{+w}{ }\PY{n}{members}\PY{p}{;}\PY{+w}{ }\PY{p}{\PYZcb{}}\PY{+w}{ }\PY{c+c1}{// leaks a mutable reference}
\PY{p}{\PYZcb{}}

\PY{c+c1}{// After: return an immutable view/copy}
\PY{k+kd}{public}\PY{+w}{ }\PY{k+kd}{class} \PY{n+nc}{Team}\PY{+w}{ }\PY{p}{\PYZob{}}
\PY{+w}{    }\PY{k+kd}{private}\PY{+w}{ }\PY{k+kd}{final}\PY{+w}{ }\PY{n}{List}\PY{o}{\PYZlt{}}\PY{n}{String}\PY{o}{\PYZgt{}}\PY{+w}{ }\PY{n}{members}\PY{+w}{ }\PY{o}{=}\PY{+w}{ }\PY{k}{new}\PY{+w}{ }\PY{n}{ArrayList}\PY{o}{\PYZlt{}}\PY{o}{\PYZgt{}}\PY{p}{(}\PY{p}{)}\PY{p}{;}
\PY{+w}{    }\PY{k+kd}{public}\PY{+w}{ }\PY{n}{List}\PY{o}{\PYZlt{}}\PY{n}{String}\PY{o}{\PYZgt{}}\PY{+w}{ }\PY{n+nf}{getMembers}\PY{p}{(}\PY{p}{)}\PY{+w}{ }\PY{p}{\PYZob{}}\PY{+w}{ }\PY{k}{return}\PY{+w}{ }\PY{n}{List}\PY{p}{.}\PY{n+na}{copyOf}\PY{p}{(}\PY{n}{members}\PY{p}{)}\PY{p}{;}\PY{+w}{ }\PY{p}{\PYZcb{}}
\PY{p}{\PYZcb{}}
\end{Verbatim}
\end{codebox}

\subsection[Nulls at Boundaries]{Nulls at Boundaries}
\label{h:19:nulls-at-boundaries}
Public API boundaries should be explicit about null: either forbid it (and fail fast) or explicitly document that it is allowed. Silent acceptance of \code{null} pushes the bug downstream to a confusing \code{NullPointerException} far from the real cause.

\begin{codebox}
\begin{Verbatim}[commandchars=\\\{\}]
\PY{k+kd}{public}\PY{+w}{ }\PY{k+kt}{void}\PY{+w}{ }\PY{n+nf}{registerUser}\PY{p}{(}\PY{n}{String}\PY{+w}{ }\PY{n}{email}\PY{p}{,}\PY{+w}{ }\PY{n}{String}\PY{+w}{ }\PY{n}{referralCode}\PY{p}{)}\PY{+w}{ }\PY{p}{\PYZob{}}
\PY{+w}{    }\PY{n}{Objects}\PY{p}{.}\PY{n+na}{requireNonNull}\PY{p}{(}\PY{n}{email}\PY{p}{,}\PY{+w}{ }\PY{l+s}{\PYZdq{}}\PY{l+s}{email must not be null}\PY{l+s}{\PYZdq{}}\PY{p}{)}\PY{p}{;}
\PY{+w}{    }\PY{c+c1}{// referralCode is allowed to be null — document it clearly}
\PY{+w}{    }\PY{c+c1}{// referralCode: nullable, meaning \PYZdq{}no referral\PYZdq{}}
\PY{+w}{    }\PY{p}{.}\PY{p}{.}\PY{p}{.}
\PY{p}{\PYZcb{}}
\end{Verbatim}
\end{codebox}

\subsection[Optional for Returns Only]{\code{Optional} for Returns Only}
\label{h:19:optional-for-returns-only}
As in the Effective Java section, restrict \code{Optional} to return types. Do not accept \code{Optional} as a method parameter in a public API — it forces callers to wrap values just to call your method, and it is not \code{Serializable}.

\begin{codebox}
\begin{Verbatim}[commandchars=\\\{\}]
\PY{c+c1}{// Avoid: forces every caller to wrap a value in Optional just to call this}
\PY{k+kd}{public}\PY{+w}{ }\PY{k+kt}{void}\PY{+w}{ }\PY{n+nf}{setDiscount}\PY{p}{(}\PY{n}{Optional}\PY{o}{\PYZlt{}}\PY{n}{BigDecimal}\PY{o}{\PYZgt{}}\PY{+w}{ }\PY{n}{discount}\PY{p}{)}\PY{+w}{ }\PY{p}{\PYZob{}}\PY{+w}{ }\PY{p}{.}\PY{p}{.}\PY{p}{.}\PY{+w}{ }\PY{p}{\PYZcb{}}

\PY{c+c1}{// Prefer: an overload, or a sentinel/nullable value with a nullability contract}
\PY{k+kd}{public}\PY{+w}{ }\PY{k+kt}{void}\PY{+w}{ }\PY{n+nf}{setDiscount}\PY{p}{(}\PY{n}{BigDecimal}\PY{+w}{ }\PY{n}{discount}\PY{p}{)}\PY{+w}{ }\PY{p}{\PYZob{}}\PY{+w}{ }\PY{p}{.}\PY{p}{.}\PY{p}{.}\PY{+w}{ }\PY{p}{\PYZcb{}}\PY{+w}{ }\PY{c+c1}{// null = no discount, documented}
\PY{k+kd}{public}\PY{+w}{ }\PY{n}{Optional}\PY{o}{\PYZlt{}}\PY{n}{BigDecimal}\PY{o}{\PYZgt{}}\PY{+w}{ }\PY{n+nf}{getDiscount}\PY{p}{(}\PY{p}{)}\PY{+w}{ }\PY{p}{\PYZob{}}\PY{+w}{ }\PY{p}{.}\PY{p}{.}\PY{p}{.}\PY{+w}{ }\PY{p}{\PYZcb{}}\PY{+w}{     }\PY{c+c1}{// Optional is fine as a return type}
\end{Verbatim}
\end{codebox}

\subsection[Parameter Validation and Fail Fast]{Parameter Validation and Fail Fast}
\label{h:19:parameter-validation-and-fail-fast}
Validate arguments at the top of a public method and throw immediately with a clear message. Failing fast at the boundary is far easier to debug than a mysterious failure deep inside the call stack.

\begin{codebox}
\begin{Verbatim}[commandchars=\\\{\}]
\PY{k+kd}{public}\PY{+w}{ }\PY{n}{BigDecimal}\PY{+w}{ }\PY{n+nf}{calculateInterest}\PY{p}{(}\PY{n}{BigDecimal}\PY{+w}{ }\PY{n}{principal}\PY{p}{,}\PY{+w}{ }\PY{k+kt}{double}\PY{+w}{ }\PY{n}{rate}\PY{p}{,}\PY{+w}{ }\PY{k+kt}{int}\PY{+w}{ }\PY{n}{years}\PY{p}{)}\PY{+w}{ }\PY{p}{\PYZob{}}
\PY{+w}{    }\PY{k}{if}\PY{+w}{ }\PY{p}{(}\PY{n}{principal}\PY{p}{.}\PY{n+na}{signum}\PY{p}{(}\PY{p}{)}\PY{+w}{ }\PY{o}{\PYZlt{}}\PY{+w}{ }\PY{l+m+mi}{0}\PY{p}{)}\PY{+w}{ }\PY{p}{\PYZob{}}
\PY{+w}{        }\PY{k}{throw}\PY{+w}{ }\PY{k}{new}\PY{+w}{ }\PY{n}{IllegalArgumentException}\PY{p}{(}\PY{l+s}{\PYZdq{}}\PY{l+s}{principal must not be negative: }\PY{l+s}{\PYZdq{}}\PY{+w}{ }\PY{o}{+}\PY{+w}{ }\PY{n}{principal}\PY{p}{)}\PY{p}{;}
\PY{+w}{    }\PY{p}{\PYZcb{}}
\PY{+w}{    }\PY{k}{if}\PY{+w}{ }\PY{p}{(}\PY{n}{years}\PY{+w}{ }\PY{o}{\PYZlt{}}\PY{+w}{ }\PY{l+m+mi}{0}\PY{p}{)}\PY{+w}{ }\PY{p}{\PYZob{}}
\PY{+w}{        }\PY{k}{throw}\PY{+w}{ }\PY{k}{new}\PY{+w}{ }\PY{n}{IllegalArgumentException}\PY{p}{(}\PY{l+s}{\PYZdq{}}\PY{l+s}{years must not be negative: }\PY{l+s}{\PYZdq{}}\PY{+w}{ }\PY{o}{+}\PY{+w}{ }\PY{n}{years}\PY{p}{)}\PY{p}{;}
\PY{+w}{    }\PY{p}{\PYZcb{}}
\PY{+w}{    }\PY{k}{return}\PY{+w}{ }\PY{n}{principal}\PY{p}{.}\PY{n+na}{multiply}\PY{p}{(}\PY{n}{BigDecimal}\PY{p}{.}\PY{n+na}{valueOf}\PY{p}{(}\PY{n}{rate}\PY{+w}{ }\PY{o}{*}\PY{+w}{ }\PY{n}{years}\PY{p}{)}\PY{p}{)}\PY{p}{;}
\PY{p}{\PYZcb{}}
\end{Verbatim}
\end{codebox}

\subsection[Overload Sparingly]{Overload Sparingly}
\label{h:19:overload-sparingly}
Too many overloads confuse callers about which one gets called, especially with autoboxing, varargs, or null arguments. Prefer distinct, clearly named methods when behavior meaningfully differs.

\begin{codebox}
\begin{Verbatim}[commandchars=\\\{\}]
\PY{c+c1}{// Confusing: which overload does log(null) call? Ambiguous / surprising.}
\PY{k+kt}{void}\PY{+w}{ }\PY{n+nf}{log}\PY{p}{(}\PY{n}{String}\PY{+w}{ }\PY{n}{message}\PY{p}{)}\PY{+w}{ }\PY{p}{\PYZob{}}\PY{+w}{ }\PY{p}{.}\PY{p}{.}\PY{p}{.}\PY{+w}{ }\PY{p}{\PYZcb{}}
\PY{k+kt}{void}\PY{+w}{ }\PY{n+nf}{log}\PY{p}{(}\PY{n}{String}\PY{+w}{ }\PY{n}{message}\PY{p}{,}\PY{+w}{ }\PY{n}{Throwable}\PY{+w}{ }\PY{n}{t}\PY{p}{)}\PY{+w}{ }\PY{p}{\PYZob{}}\PY{+w}{ }\PY{p}{.}\PY{p}{.}\PY{p}{.}\PY{+w}{ }\PY{p}{\PYZcb{}}
\PY{k+kt}{void}\PY{+w}{ }\PY{n+nf}{log}\PY{p}{(}\PY{n}{Object}\PY{+w}{ }\PY{n}{message}\PY{p}{)}\PY{+w}{ }\PY{p}{\PYZob{}}\PY{+w}{ }\PY{p}{.}\PY{p}{.}\PY{p}{.}\PY{+w}{ }\PY{p}{\PYZcb{}}

\PY{c+c1}{// Clearer: distinct names remove all ambiguity}
\PY{k+kt}{void}\PY{+w}{ }\PY{n+nf}{logMessage}\PY{p}{(}\PY{n}{String}\PY{+w}{ }\PY{n}{message}\PY{p}{)}\PY{+w}{ }\PY{p}{\PYZob{}}\PY{+w}{ }\PY{p}{.}\PY{p}{.}\PY{p}{.}\PY{+w}{ }\PY{p}{\PYZcb{}}
\PY{k+kt}{void}\PY{+w}{ }\PY{n+nf}{logError}\PY{p}{(}\PY{n}{String}\PY{+w}{ }\PY{n}{message}\PY{p}{,}\PY{+w}{ }\PY{n}{Throwable}\PY{+w}{ }\PY{n}{t}\PY{p}{)}\PY{+w}{ }\PY{p}{\PYZob{}}\PY{+w}{ }\PY{p}{.}\PY{p}{.}\PY{p}{.}\PY{+w}{ }\PY{p}{\PYZcb{}}
\end{Verbatim}
\end{codebox}

\subsection[Exception Contracts]{Exception Contracts}
\label{h:19:exception-contracts}
Document exactly which exceptions a public method can throw, and keep that contract stable — it is part of your API. Prefer specific, well-named exceptions over generic ones (see Exception Handling chapter for full details).

\begin{codebox}
\begin{Verbatim}[commandchars=\\\{\}]
\PY{c+cm}{/**}
\PY{c+cm}{ * @throws UserNotFoundException if no user exists with the given id}
\PY{c+cm}{ * @throws IllegalArgumentException if id is blank}
\PY{c+cm}{ */}
\PY{k+kd}{public}\PY{+w}{ }\PY{n}{User}\PY{+w}{ }\PY{n+nf}{getUser}\PY{p}{(}\PY{n}{String}\PY{+w}{ }\PY{n}{id}\PY{p}{)}\PY{+w}{ }\PY{k+kd}{throws}\PY{+w}{ }\PY{n}{UserNotFoundException}\PY{+w}{ }\PY{p}{\PYZob{}}\PY{+w}{ }\PY{p}{.}\PY{p}{.}\PY{p}{.}\PY{+w}{ }\PY{p}{\PYZcb{}}
\end{Verbatim}
\end{codebox}

\subsection[Backwards Compatibility: Binary vs Source Compatibility]{Backwards Compatibility: Binary vs Source Compatibility}
\label{h:19:backwards-compatibility-binary-vs-source-compatibility}
\textbf{Source compatible} means existing client code still \emph{compiles} against the new version. \textbf{Binary compatible} means existing \emph{compiled} \code{.\allowbreak{}class} files still run against the new version without recompiling. Adding a new overload is usually both; removing a public method breaks both; changing a method’s return type (even to a subtype) can break binary compatibility even though source may still compile.

\begin{codebox}
\begin{Verbatim}[commandchars=\\\{\}]
\PY{c+c1}{// Adding an overload: source and binary compatible — old callers unaffected}
\PY{k+kd}{public}\PY{+w}{ }\PY{k+kt}{void}\PY{+w}{ }\PY{n+nf}{save}\PY{p}{(}\PY{n}{Order}\PY{+w}{ }\PY{n}{order}\PY{p}{)}\PY{+w}{ }\PY{p}{\PYZob{}}\PY{+w}{ }\PY{p}{.}\PY{p}{.}\PY{p}{.}\PY{+w}{ }\PY{p}{\PYZcb{}}
\PY{k+kd}{public}\PY{+w}{ }\PY{k+kt}{void}\PY{+w}{ }\PY{n+nf}{save}\PY{p}{(}\PY{n}{Order}\PY{+w}{ }\PY{n}{order}\PY{p}{,}\PY{+w}{ }\PY{k+kt}{boolean}\PY{+w}{ }\PY{n}{async}\PY{p}{)}\PY{+w}{ }\PY{p}{\PYZob{}}\PY{+w}{ }\PY{p}{.}\PY{p}{.}\PY{p}{.}\PY{+w}{ }\PY{p}{\PYZcb{}}\PY{+w}{ }\PY{c+c1}{// new, additive}

\PY{c+c1}{// Removing/renaming a public method: breaks both source and binary compatibility}
\PY{c+c1}{// public void save(Order order) \PYZob{} ... \PYZcb{}  // deleting this breaks every existing caller}
\end{Verbatim}
\end{codebox}

\subsection[@Deprecated(since, forRemoval)]{\code{@\allowbreak{}Deprecated(\allowbreak{}since,\allowbreak{}~\allowbreak{}for\allowbreak{}Removal)}}
\label{h:19:deprecatedsince-forremoval}
Deprecate before removing. \code{since} documents when the replacement became available; \code{forRemoval~\allowbreak{}=\allowbreak{}~\allowbreak{}true} tells consumers this really is going away, not just “prefer the alternative.”

\begin{codebox}
\begin{Verbatim}[commandchars=\\\{\}]
\PY{c+cm}{/**}
\PY{c+cm}{ * @deprecated use \PYZob{}@link \PYZsh{}save(Order, boolean)\PYZcb{} instead.}
\PY{c+cm}{ */}
\PY{n+nd}{@Deprecated}\PY{p}{(}\PY{n}{since}\PY{+w}{ }\PY{o}{=}\PY{+w}{ }\PY{l+s}{\PYZdq{}}\PY{l+s}{2.3}\PY{l+s}{\PYZdq{}}\PY{p}{,}\PY{+w}{ }\PY{n}{forRemoval}\PY{+w}{ }\PY{o}{=}\PY{+w}{ }\PY{k+kc}{true}\PY{p}{)}
\PY{k+kd}{public}\PY{+w}{ }\PY{k+kt}{void}\PY{+w}{ }\PY{n+nf}{save}\PY{p}{(}\PY{n}{Order}\PY{+w}{ }\PY{n}{order}\PY{p}{)}\PY{+w}{ }\PY{p}{\PYZob{}}
\PY{+w}{    }\PY{n}{save}\PY{p}{(}\PY{n}{order}\PY{p}{,}\PY{+w}{ }\PY{k+kc}{false}\PY{p}{)}\PY{p}{;}
\PY{p}{\PYZcb{}}

\PY{k+kd}{public}\PY{+w}{ }\PY{k+kt}{void}\PY{+w}{ }\PY{n+nf}{save}\PY{p}{(}\PY{n}{Order}\PY{+w}{ }\PY{n}{order}\PY{p}{,}\PY{+w}{ }\PY{k+kt}{boolean}\PY{+w}{ }\PY{n}{async}\PY{p}{)}\PY{+w}{ }\PY{p}{\PYZob{}}\PY{+w}{ }\PY{p}{.}\PY{p}{.}\PY{p}{.}\PY{+w}{ }\PY{p}{\PYZcb{}}
\end{Verbatim}
\end{codebox}

\subsection[Versioning]{Versioning}
\label{h:19:versioning}
Follow semantic versioning (\code{MAJOR.\allowbreak{}MINOR.\allowbreak{}PATCH}): increment \code{MAJOR} for breaking changes, \code{MINOR} for backwards-compatible additions, \code{PATCH} for backwards-compatible bug fixes. This lets consumers set dependency ranges with confidence.

\begin{codebox}
\begin{Verbatim}
1.4.2 -> 1.5.0   // new feature, backwards compatible (MINOR bump)
1.5.0 -> 1.5.1   // bug fix only (PATCH bump)
1.5.1 -> 2.0.0   // removed a deprecated method (MAJOR bump, breaking)
\end{Verbatim}
\end{codebox}

\subsection[Builders for Wide Constructors]{Builders for Wide Constructors}
\label{h:19:builders-for-wide-constructors}
As covered in Clean Code and Effective Java, expose a builder (or static factory with named parameters via a record) instead of a constructor with many parameters — this is doubly important in a public API, where the constructor signature becomes a long-term commitment.

\begin{codebox}
\begin{Verbatim}[commandchars=\\\{\}]
\PY{k+kd}{public}\PY{+w}{ }\PY{k+kd}{class} \PY{n+nc}{HttpRequestConfig}\PY{+w}{ }\PY{p}{\PYZob{}}
\PY{+w}{    }\PY{c+c1}{// A public constructor with 8 parameters would be locked in forever.}
\PY{+w}{    }\PY{c+c1}{// A builder can add optional settings later without breaking callers.}
\PY{+w}{    }\PY{k+kd}{public}\PY{+w}{ }\PY{k+kd}{static}\PY{+w}{ }\PY{n}{Builder}\PY{+w}{ }\PY{n+nf}{builder}\PY{p}{(}\PY{p}{)}\PY{+w}{ }\PY{p}{\PYZob{}}\PY{+w}{ }\PY{k}{return}\PY{+w}{ }\PY{k}{new}\PY{+w}{ }\PY{n}{Builder}\PY{p}{(}\PY{p}{)}\PY{p}{;}\PY{+w}{ }\PY{p}{\PYZcb{}}

\PY{+w}{    }\PY{k+kd}{public}\PY{+w}{ }\PY{k+kd}{static}\PY{+w}{ }\PY{k+kd}{final}\PY{+w}{ }\PY{k+kd}{class} \PY{n+nc}{Builder}\PY{+w}{ }\PY{p}{\PYZob{}}
\PY{+w}{        }\PY{k+kd}{private}\PY{+w}{ }\PY{n}{Duration}\PY{+w}{ }\PY{n}{timeout}\PY{+w}{ }\PY{o}{=}\PY{+w}{ }\PY{n}{Duration}\PY{p}{.}\PY{n+na}{ofSeconds}\PY{p}{(}\PY{l+m+mi}{30}\PY{p}{)}\PY{p}{;}
\PY{+w}{        }\PY{k+kd}{private}\PY{+w}{ }\PY{k+kt}{int}\PY{+w}{ }\PY{n}{maxRetries}\PY{+w}{ }\PY{o}{=}\PY{+w}{ }\PY{l+m+mi}{3}\PY{p}{;}
\PY{+w}{        }\PY{k+kd}{public}\PY{+w}{ }\PY{n}{Builder}\PY{+w}{ }\PY{n+nf}{timeout}\PY{p}{(}\PY{n}{Duration}\PY{+w}{ }\PY{n}{timeout}\PY{p}{)}\PY{+w}{ }\PY{p}{\PYZob{}}\PY{+w}{ }\PY{k}{this}\PY{p}{.}\PY{n+na}{timeout}\PY{+w}{ }\PY{o}{=}\PY{+w}{ }\PY{n}{timeout}\PY{p}{;}\PY{+w}{ }\PY{k}{return}\PY{+w}{ }\PY{k}{this}\PY{p}{;}\PY{+w}{ }\PY{p}{\PYZcb{}}
\PY{+w}{        }\PY{k+kd}{public}\PY{+w}{ }\PY{n}{Builder}\PY{+w}{ }\PY{n+nf}{maxRetries}\PY{p}{(}\PY{k+kt}{int}\PY{+w}{ }\PY{n}{maxRetries}\PY{p}{)}\PY{+w}{ }\PY{p}{\PYZob{}}\PY{+w}{ }\PY{k}{this}\PY{p}{.}\PY{n+na}{maxRetries}\PY{+w}{ }\PY{o}{=}\PY{+w}{ }\PY{n}{maxRetries}\PY{p}{;}\PY{+w}{ }\PY{k}{return}\PY{+w}{ }\PY{k}{this}\PY{p}{;}\PY{+w}{ }\PY{p}{\PYZcb{}}
\PY{+w}{        }\PY{k+kd}{public}\PY{+w}{ }\PY{n}{HttpRequestConfig}\PY{+w}{ }\PY{n+nf}{build}\PY{p}{(}\PY{p}{)}\PY{+w}{ }\PY{p}{\PYZob{}}\PY{+w}{ }\PY{k}{return}\PY{+w}{ }\PY{k}{new}\PY{+w}{ }\PY{n}{HttpRequestConfig}\PY{p}{(}\PY{k}{this}\PY{p}{)}\PY{p}{;}\PY{+w}{ }\PY{p}{\PYZcb{}}
\PY{+w}{    }\PY{p}{\PYZcb{}}
\PY{+w}{    }\PY{c+c1}{// private fields/constructor omitted for brevity}
\PY{p}{\PYZcb{}}
\end{Verbatim}
\end{codebox}

\subsection[Documenting Thread Safety and Nullability]{Documenting Thread Safety and Nullability}
\label{h:19:documenting-thread-safety-and-nullability}
Every public class and method should document whether it is thread-safe and whether parameters/return values may be \code{null}. This is not optional polish — it directly determines how a caller is allowed to use the API.

\begin{codebox}
\begin{Verbatim}[commandchars=\\\{\}]
\PY{c+cm}{/**}
\PY{c+cm}{ * Thread\PYZhy{}safe cache. Safe to share a single instance across threads.}
\PY{c+cm}{ *}
\PY{c+cm}{ * @param key must not be null}
\PY{c+cm}{ * @return the cached value, or \PYZob{}@code null\PYZcb{} if absent (this method never throws for a cache miss)}
\PY{c+cm}{ */}
\PY{k+kd}{public}\PY{+w}{ }\PY{n}{V}\PY{+w}{ }\PY{n+nf}{get}\PY{p}{(}\PY{n}{K}\PY{+w}{ }\PY{n}{key}\PY{p}{)}\PY{+w}{ }\PY{p}{\PYZob{}}\PY{+w}{ }\PY{p}{.}\PY{p}{.}\PY{p}{.}\PY{+w}{ }\PY{p}{\PYZcb{}}
\end{Verbatim}
\end{codebox}

\section[Common Code Smells]{Common Code Smells}
\label{h:19:common-code-smells}
A \textbf{code smell} is a surface signal that deeper design trouble might be lurking — not a bug by itself, but a warning sign worth investigating.

{\tablefont
\begin{xltabular}{\linewidth}{@{}L{0.595}L{1.463}L{0.942}@{}}
\toprule
\rowcolor{tablehdr}
\thd{Smell} & \thd{Why it hurts} & \thd{Refactoring} \\
\midrule
\endhead
Long method & Hard to understand, test, and reuse pieces of & Extract method, guard clauses \\
Large class / God object & Too many responsibilities, high change risk & Split by responsibility (SRP) \\
Primitive obsession & Loses domain meaning and validation & Introduce small value types / records \\
Feature envy & A method cares more about another class’s data than its own & Move the method to the class it envies \\
Data clumps & The same group of fields/params travels together everywhere & Extract a parameter object \\
Shotgun surgery & One logical change requires edits across many classes & Consolidate the responsibility into one place \\
Anemic domain model & Objects are just data bags; logic lives elsewhere & Move behavior into the domain objects \\
Boolean parameters & Ambiguous call sites, hidden branching & Split methods, or use an enum \\
Deep inheritance & Fragile, hard to trace behavior across many levels & Favor composition \\
Utility-class dumping ground & Unrelated static methods with no cohesion & Split into focused, named classes \\
Stringly-typed code & No compile-time safety for what is really a fixed set of values & Use enums or dedicated types \\
Exception as control flow & Exceptions are slow and obscure normal logic paths & Use normal conditionals / \code{Optional} \\
Static abuse & Hidden global state, hard to test & Instance fields + dependency injection \\
Over-mocking & Tests couple to implementation, not behavior & Test through the public behavior/contract \\
\bottomrule
\end{xltabular}
}

\subsection[Long Method]{Long Method}
\label{h:19:long-method}
\begin{codebox}
\begin{Verbatim}[commandchars=\\\{\}]
\PY{c+c1}{// Before: one long method mixing many steps}
\PY{k+kt}{void}\PY{+w}{ }\PY{n+nf}{processOrder}\PY{p}{(}\PY{n}{Order}\PY{+w}{ }\PY{n}{order}\PY{p}{)}\PY{+w}{ }\PY{p}{\PYZob{}}
\PY{+w}{    }\PY{k}{if}\PY{+w}{ }\PY{p}{(}\PY{n}{order}\PY{p}{.}\PY{n+na}{getItems}\PY{p}{(}\PY{p}{)}\PY{p}{.}\PY{n+na}{isEmpty}\PY{p}{(}\PY{p}{)}\PY{p}{)}\PY{+w}{ }\PY{k}{throw}\PY{+w}{ }\PY{k}{new}\PY{+w}{ }\PY{n}{IllegalArgumentException}\PY{p}{(}\PY{l+s}{\PYZdq{}}\PY{l+s}{empty order}\PY{l+s}{\PYZdq{}}\PY{p}{)}\PY{p}{;}
\PY{+w}{    }\PY{n}{BigDecimal}\PY{+w}{ }\PY{n}{total}\PY{+w}{ }\PY{o}{=}\PY{+w}{ }\PY{n}{BigDecimal}\PY{p}{.}\PY{n+na}{ZERO}\PY{p}{;}
\PY{+w}{    }\PY{k}{for}\PY{+w}{ }\PY{p}{(}\PY{n}{OrderItem}\PY{+w}{ }\PY{n}{item}\PY{+w}{ }\PY{p}{:}\PY{+w}{ }\PY{n}{order}\PY{p}{.}\PY{n+na}{getItems}\PY{p}{(}\PY{p}{)}\PY{p}{)}\PY{+w}{ }\PY{p}{\PYZob{}}
\PY{+w}{        }\PY{n}{total}\PY{+w}{ }\PY{o}{=}\PY{+w}{ }\PY{n}{total}\PY{p}{.}\PY{n+na}{add}\PY{p}{(}\PY{n}{item}\PY{p}{.}\PY{n+na}{getPrice}\PY{p}{(}\PY{p}{)}\PY{p}{.}\PY{n+na}{multiply}\PY{p}{(}\PY{n}{BigDecimal}\PY{p}{.}\PY{n+na}{valueOf}\PY{p}{(}\PY{n}{item}\PY{p}{.}\PY{n+na}{getQuantity}\PY{p}{(}\PY{p}{)}\PY{p}{)}\PY{p}{)}\PY{p}{)}\PY{p}{;}
\PY{+w}{    }\PY{p}{\PYZcb{}}
\PY{+w}{    }\PY{k}{if}\PY{+w}{ }\PY{p}{(}\PY{n}{order}\PY{p}{.}\PY{n+na}{getCoupon}\PY{p}{(}\PY{p}{)}\PY{+w}{ }\PY{o}{!}\PY{o}{=}\PY{+w}{ }\PY{k+kc}{null}\PY{p}{)}\PY{+w}{ }\PY{p}{\PYZob{}}
\PY{+w}{        }\PY{n}{total}\PY{+w}{ }\PY{o}{=}\PY{+w}{ }\PY{n}{total}\PY{p}{.}\PY{n+na}{subtract}\PY{p}{(}\PY{n}{total}\PY{p}{.}\PY{n+na}{multiply}\PY{p}{(}\PY{n}{order}\PY{p}{.}\PY{n+na}{getCoupon}\PY{p}{(}\PY{p}{)}\PY{p}{.}\PY{n+na}{rate}\PY{p}{(}\PY{p}{)}\PY{p}{)}\PY{p}{)}\PY{p}{;}
\PY{+w}{    }\PY{p}{\PYZcb{}}
\PY{+w}{    }\PY{n}{inventory}\PY{p}{.}\PY{n+na}{reserve}\PY{p}{(}\PY{n}{order}\PY{p}{)}\PY{p}{;}
\PY{+w}{    }\PY{n}{payment}\PY{p}{.}\PY{n+na}{charge}\PY{p}{(}\PY{n}{order}\PY{p}{.}\PY{n+na}{getCustomerId}\PY{p}{(}\PY{p}{)}\PY{p}{,}\PY{+w}{ }\PY{n}{total}\PY{p}{)}\PY{p}{;}
\PY{+w}{    }\PY{n}{shipping}\PY{p}{.}\PY{n+na}{schedule}\PY{p}{(}\PY{n}{order}\PY{p}{)}\PY{p}{;}
\PY{p}{\PYZcb{}}
\end{Verbatim}
\end{codebox}

\begin{codebox}
\begin{Verbatim}[commandchars=\\\{\}]
\PY{c+c1}{// After: extracted into small, named steps}
\PY{k+kt}{void}\PY{+w}{ }\PY{n+nf}{processOrder}\PY{p}{(}\PY{n}{Order}\PY{+w}{ }\PY{n}{order}\PY{p}{)}\PY{+w}{ }\PY{p}{\PYZob{}}
\PY{+w}{    }\PY{n}{validate}\PY{p}{(}\PY{n}{order}\PY{p}{)}\PY{p}{;}
\PY{+w}{    }\PY{n}{BigDecimal}\PY{+w}{ }\PY{n}{total}\PY{+w}{ }\PY{o}{=}\PY{+w}{ }\PY{n}{calculateTotal}\PY{p}{(}\PY{n}{order}\PY{p}{)}\PY{p}{;}
\PY{+w}{    }\PY{n}{inventory}\PY{p}{.}\PY{n+na}{reserve}\PY{p}{(}\PY{n}{order}\PY{p}{)}\PY{p}{;}
\PY{+w}{    }\PY{n}{payment}\PY{p}{.}\PY{n+na}{charge}\PY{p}{(}\PY{n}{order}\PY{p}{.}\PY{n+na}{getCustomerId}\PY{p}{(}\PY{p}{)}\PY{p}{,}\PY{+w}{ }\PY{n}{total}\PY{p}{)}\PY{p}{;}
\PY{+w}{    }\PY{n}{shipping}\PY{p}{.}\PY{n+na}{schedule}\PY{p}{(}\PY{n}{order}\PY{p}{)}\PY{p}{;}
\PY{p}{\PYZcb{}}
\end{Verbatim}
\end{codebox}

\subsection[Large Class / God Object]{Large Class / God Object}
\label{h:19:large-class--god-object}
\begin{codebox}
\begin{Verbatim}[commandchars=\\\{\}]
\PY{c+c1}{// Before: one class does everything for the whole application}
\PY{k+kd}{class} \PY{n+nc}{ApplicationManager}\PY{+w}{ }\PY{p}{\PYZob{}}
\PY{+w}{    }\PY{k+kt}{void}\PY{+w}{ }\PY{n+nf}{authenticateUser}\PY{p}{(}\PY{n}{String}\PY{+w}{ }\PY{n}{u}\PY{p}{,}\PY{+w}{ }\PY{n}{String}\PY{+w}{ }\PY{n}{p}\PY{p}{)}\PY{+w}{ }\PY{p}{\PYZob{}}\PY{+w}{ }\PY{p}{.}\PY{p}{.}\PY{p}{.}\PY{+w}{ }\PY{p}{\PYZcb{}}
\PY{+w}{    }\PY{k+kt}{void}\PY{+w}{ }\PY{n+nf}{sendEmail}\PY{p}{(}\PY{n}{String}\PY{+w}{ }\PY{n}{to}\PY{p}{,}\PY{+w}{ }\PY{n}{String}\PY{+w}{ }\PY{n}{body}\PY{p}{)}\PY{+w}{ }\PY{p}{\PYZob{}}\PY{+w}{ }\PY{p}{.}\PY{p}{.}\PY{p}{.}\PY{+w}{ }\PY{p}{\PYZcb{}}
\PY{+w}{    }\PY{k+kt}{void}\PY{+w}{ }\PY{n+nf}{generateReport}\PY{p}{(}\PY{p}{)}\PY{+w}{ }\PY{p}{\PYZob{}}\PY{+w}{ }\PY{p}{.}\PY{p}{.}\PY{p}{.}\PY{+w}{ }\PY{p}{\PYZcb{}}
\PY{+w}{    }\PY{k+kt}{void}\PY{+w}{ }\PY{n+nf}{backupDatabase}\PY{p}{(}\PY{p}{)}\PY{+w}{ }\PY{p}{\PYZob{}}\PY{+w}{ }\PY{p}{.}\PY{p}{.}\PY{p}{.}\PY{+w}{ }\PY{p}{\PYZcb{}}
\PY{p}{\PYZcb{}}
\end{Verbatim}
\end{codebox}

\begin{codebox}
\begin{Verbatim}[commandchars=\\\{\}]
\PY{c+c1}{// After: split by responsibility}
\PY{k+kd}{class} \PY{n+nc}{AuthService}\PY{+w}{ }\PY{p}{\PYZob{}}\PY{+w}{ }\PY{k+kt}{void}\PY{+w}{ }\PY{n+nf}{authenticate}\PY{p}{(}\PY{n}{String}\PY{+w}{ }\PY{n}{u}\PY{p}{,}\PY{+w}{ }\PY{n}{String}\PY{+w}{ }\PY{n}{p}\PY{p}{)}\PY{+w}{ }\PY{p}{\PYZob{}}\PY{+w}{ }\PY{p}{.}\PY{p}{.}\PY{p}{.}\PY{+w}{ }\PY{p}{\PYZcb{}}\PY{+w}{ }\PY{p}{\PYZcb{}}
\PY{k+kd}{class} \PY{n+nc}{EmailService}\PY{+w}{ }\PY{p}{\PYZob{}}\PY{+w}{ }\PY{k+kt}{void}\PY{+w}{ }\PY{n+nf}{send}\PY{p}{(}\PY{n}{String}\PY{+w}{ }\PY{n}{to}\PY{p}{,}\PY{+w}{ }\PY{n}{String}\PY{+w}{ }\PY{n}{body}\PY{p}{)}\PY{+w}{ }\PY{p}{\PYZob{}}\PY{+w}{ }\PY{p}{.}\PY{p}{.}\PY{p}{.}\PY{+w}{ }\PY{p}{\PYZcb{}}\PY{+w}{ }\PY{p}{\PYZcb{}}
\PY{k+kd}{class} \PY{n+nc}{ReportService}\PY{+w}{ }\PY{p}{\PYZob{}}\PY{+w}{ }\PY{k+kt}{void}\PY{+w}{ }\PY{n+nf}{generate}\PY{p}{(}\PY{p}{)}\PY{+w}{ }\PY{p}{\PYZob{}}\PY{+w}{ }\PY{p}{.}\PY{p}{.}\PY{p}{.}\PY{+w}{ }\PY{p}{\PYZcb{}}\PY{+w}{ }\PY{p}{\PYZcb{}}
\PY{k+kd}{class} \PY{n+nc}{BackupService}\PY{+w}{ }\PY{p}{\PYZob{}}\PY{+w}{ }\PY{k+kt}{void}\PY{+w}{ }\PY{n+nf}{backupDatabase}\PY{p}{(}\PY{p}{)}\PY{+w}{ }\PY{p}{\PYZob{}}\PY{+w}{ }\PY{p}{.}\PY{p}{.}\PY{p}{.}\PY{+w}{ }\PY{p}{\PYZcb{}}\PY{+w}{ }\PY{p}{\PYZcb{}}
\end{Verbatim}
\end{codebox}

\subsection[Primitive Obsession]{Primitive Obsession}
\label{h:19:primitive-obsession}
\begin{codebox}
\begin{Verbatim}[commandchars=\\\{\}]
\PY{c+c1}{// Before: a raw String can hold any garbage, no validation, no meaning}
\PY{k+kt}{void}\PY{+w}{ }\PY{n+nf}{sendEmail}\PY{p}{(}\PY{n}{String}\PY{+w}{ }\PY{n}{email}\PY{p}{)}\PY{+w}{ }\PY{p}{\PYZob{}}
\PY{+w}{    }\PY{k}{if}\PY{+w}{ }\PY{p}{(}\PY{o}{!}\PY{n}{email}\PY{p}{.}\PY{n+na}{contains}\PY{p}{(}\PY{l+s}{\PYZdq{}}\PY{l+s}{@}\PY{l+s}{\PYZdq{}}\PY{p}{)}\PY{p}{)}\PY{+w}{ }\PY{k}{throw}\PY{+w}{ }\PY{k}{new}\PY{+w}{ }\PY{n}{IllegalArgumentException}\PY{p}{(}\PY{l+s}{\PYZdq{}}\PY{l+s}{bad email}\PY{l+s}{\PYZdq{}}\PY{p}{)}\PY{p}{;}
\PY{+w}{    }\PY{p}{.}\PY{p}{.}\PY{p}{.}
\PY{p}{\PYZcb{}}

\PY{c+c1}{// After: a small value type validates itself and carries meaning everywhere}
\PY{k+kd}{record} \PY{n+nc}{EmailAddress}\PY{p}{(}\PY{n}{String}\PY{+w}{ }\PY{n}{value}\PY{p}{)}\PY{+w}{ }\PY{p}{\PYZob{}}
\PY{+w}{    }\PY{n}{EmailAddress}\PY{+w}{ }\PY{p}{\PYZob{}}
\PY{+w}{        }\PY{k}{if}\PY{+w}{ }\PY{p}{(}\PY{o}{!}\PY{n}{value}\PY{p}{.}\PY{n+na}{contains}\PY{p}{(}\PY{l+s}{\PYZdq{}}\PY{l+s}{@}\PY{l+s}{\PYZdq{}}\PY{p}{)}\PY{p}{)}\PY{+w}{ }\PY{k}{throw}\PY{+w}{ }\PY{k}{new}\PY{+w}{ }\PY{n}{IllegalArgumentException}\PY{p}{(}\PY{l+s}{\PYZdq{}}\PY{l+s}{bad email: }\PY{l+s}{\PYZdq{}}\PY{+w}{ }\PY{o}{+}\PY{+w}{ }\PY{n}{value}\PY{p}{)}\PY{p}{;}
\PY{+w}{    }\PY{p}{\PYZcb{}}
\PY{p}{\PYZcb{}}

\PY{k+kt}{void}\PY{+w}{ }\PY{n+nf}{sendEmail}\PY{p}{(}\PY{n}{EmailAddress}\PY{+w}{ }\PY{n}{email}\PY{p}{)}\PY{+w}{ }\PY{p}{\PYZob{}}\PY{+w}{ }\PY{p}{.}\PY{p}{.}\PY{p}{.}\PY{+w}{ }\PY{p}{\PYZcb{}}\PY{+w}{ }\PY{c+c1}{// impossible to pass an invalid one}
\end{Verbatim}
\end{codebox}

\subsection[Feature Envy]{Feature Envy}
\label{h:19:feature-envy}
\begin{codebox}
\begin{Verbatim}[commandchars=\\\{\}]
\PY{c+c1}{// Before: OrderPrinter reaches deep into Order\PYZsq{}s internals to compute something}
\PY{c+c1}{// that arguably belongs on Order itself}
\PY{k+kd}{class} \PY{n+nc}{OrderPrinter}\PY{+w}{ }\PY{p}{\PYZob{}}
\PY{+w}{    }\PY{k+kt}{void}\PY{+w}{ }\PY{n+nf}{print}\PY{p}{(}\PY{n}{Order}\PY{+w}{ }\PY{n}{order}\PY{p}{)}\PY{+w}{ }\PY{p}{\PYZob{}}
\PY{+w}{        }\PY{n}{BigDecimal}\PY{+w}{ }\PY{n}{total}\PY{+w}{ }\PY{o}{=}\PY{+w}{ }\PY{n}{BigDecimal}\PY{p}{.}\PY{n+na}{ZERO}\PY{p}{;}
\PY{+w}{        }\PY{k}{for}\PY{+w}{ }\PY{p}{(}\PY{n}{OrderItem}\PY{+w}{ }\PY{n}{item}\PY{+w}{ }\PY{p}{:}\PY{+w}{ }\PY{n}{order}\PY{p}{.}\PY{n+na}{getItems}\PY{p}{(}\PY{p}{)}\PY{p}{)}\PY{+w}{ }\PY{p}{\PYZob{}}
\PY{+w}{            }\PY{n}{total}\PY{+w}{ }\PY{o}{=}\PY{+w}{ }\PY{n}{total}\PY{p}{.}\PY{n+na}{add}\PY{p}{(}\PY{n}{item}\PY{p}{.}\PY{n+na}{getPrice}\PY{p}{(}\PY{p}{)}\PY{p}{.}\PY{n+na}{multiply}\PY{p}{(}\PY{n}{BigDecimal}\PY{p}{.}\PY{n+na}{valueOf}\PY{p}{(}\PY{n}{item}\PY{p}{.}\PY{n+na}{getQuantity}\PY{p}{(}\PY{p}{)}\PY{p}{)}\PY{p}{)}\PY{p}{)}\PY{p}{;}
\PY{+w}{        }\PY{p}{\PYZcb{}}
\PY{+w}{        }\PY{n}{System}\PY{p}{.}\PY{n+na}{out}\PY{p}{.}\PY{n+na}{println}\PY{p}{(}\PY{l+s}{\PYZdq{}}\PY{l+s}{Total: }\PY{l+s}{\PYZdq{}}\PY{+w}{ }\PY{o}{+}\PY{+w}{ }\PY{n}{total}\PY{p}{)}\PY{p}{;}
\PY{+w}{    }\PY{p}{\PYZcb{}}
\PY{p}{\PYZcb{}}
\end{Verbatim}
\end{codebox}

\begin{codebox}
\begin{Verbatim}[commandchars=\\\{\}]
\PY{c+c1}{// After: the computation moves to the class that owns the data}
\PY{k+kd}{class} \PY{n+nc}{Order}\PY{+w}{ }\PY{p}{\PYZob{}}
\PY{+w}{    }\PY{n}{BigDecimal}\PY{+w}{ }\PY{n+nf}{calculateTotal}\PY{p}{(}\PY{p}{)}\PY{+w}{ }\PY{p}{\PYZob{}}
\PY{+w}{        }\PY{k}{return}\PY{+w}{ }\PY{n}{items}\PY{p}{.}\PY{n+na}{stream}\PY{p}{(}\PY{p}{)}
\PY{+w}{                }\PY{p}{.}\PY{n+na}{map}\PY{p}{(}\PY{n}{i}\PY{+w}{ }\PY{o}{\PYZhy{}}\PY{o}{\PYZgt{}}\PY{+w}{ }\PY{n}{i}\PY{p}{.}\PY{n+na}{getPrice}\PY{p}{(}\PY{p}{)}\PY{p}{.}\PY{n+na}{multiply}\PY{p}{(}\PY{n}{BigDecimal}\PY{p}{.}\PY{n+na}{valueOf}\PY{p}{(}\PY{n}{i}\PY{p}{.}\PY{n+na}{getQuantity}\PY{p}{(}\PY{p}{)}\PY{p}{)}\PY{p}{)}\PY{p}{)}
\PY{+w}{                }\PY{p}{.}\PY{n+na}{reduce}\PY{p}{(}\PY{n}{BigDecimal}\PY{p}{.}\PY{n+na}{ZERO}\PY{p}{,}\PY{+w}{ }\PY{n}{BigDecimal}\PY{p}{:}\PY{p}{:}\PY{n}{add}\PY{p}{)}\PY{p}{;}
\PY{+w}{    }\PY{p}{\PYZcb{}}
\PY{p}{\PYZcb{}}

\PY{k+kd}{class} \PY{n+nc}{OrderPrinter}\PY{+w}{ }\PY{p}{\PYZob{}}
\PY{+w}{    }\PY{k+kt}{void}\PY{+w}{ }\PY{n+nf}{print}\PY{p}{(}\PY{n}{Order}\PY{+w}{ }\PY{n}{order}\PY{p}{)}\PY{+w}{ }\PY{p}{\PYZob{}}
\PY{+w}{        }\PY{n}{System}\PY{p}{.}\PY{n+na}{out}\PY{p}{.}\PY{n+na}{println}\PY{p}{(}\PY{l+s}{\PYZdq{}}\PY{l+s}{Total: }\PY{l+s}{\PYZdq{}}\PY{+w}{ }\PY{o}{+}\PY{+w}{ }\PY{n}{order}\PY{p}{.}\PY{n+na}{calculateTotal}\PY{p}{(}\PY{p}{)}\PY{p}{)}\PY{p}{;}
\PY{+w}{    }\PY{p}{\PYZcb{}}
\PY{p}{\PYZcb{}}
\end{Verbatim}
\end{codebox}

\subsection[Data Clumps]{Data Clumps}
\label{h:19:data-clumps}
\begin{codebox}
\begin{Verbatim}[commandchars=\\\{\}]
\PY{c+c1}{// Before: street/city/zip always travel together as separate params}
\PY{k+kt}{void}\PY{+w}{ }\PY{n+nf}{shipOrder}\PY{p}{(}\PY{n}{String}\PY{+w}{ }\PY{n}{street}\PY{p}{,}\PY{+w}{ }\PY{n}{String}\PY{+w}{ }\PY{n}{city}\PY{p}{,}\PY{+w}{ }\PY{n}{String}\PY{+w}{ }\PY{n}{zip}\PY{p}{,}\PY{+w}{ }\PY{n}{Order}\PY{+w}{ }\PY{n}{order}\PY{p}{)}\PY{+w}{ }\PY{p}{\PYZob{}}\PY{+w}{ }\PY{p}{.}\PY{p}{.}\PY{p}{.}\PY{+w}{ }\PY{p}{\PYZcb{}}
\PY{k+kt}{void}\PY{+w}{ }\PY{n+nf}{billCustomer}\PY{p}{(}\PY{n}{String}\PY{+w}{ }\PY{n}{street}\PY{p}{,}\PY{+w}{ }\PY{n}{String}\PY{+w}{ }\PY{n}{city}\PY{p}{,}\PY{+w}{ }\PY{n}{String}\PY{+w}{ }\PY{n}{zip}\PY{p}{,}\PY{+w}{ }\PY{n}{Customer}\PY{+w}{ }\PY{n}{customer}\PY{p}{)}\PY{+w}{ }\PY{p}{\PYZob{}}\PY{+w}{ }\PY{p}{.}\PY{p}{.}\PY{p}{.}\PY{+w}{ }\PY{p}{\PYZcb{}}
\end{Verbatim}
\end{codebox}

\begin{codebox}
\begin{Verbatim}[commandchars=\\\{\}]
\PY{c+c1}{// After: extracted into a single cohesive type}
\PY{k+kd}{record} \PY{n+nc}{Address}\PY{p}{(}\PY{n}{String}\PY{+w}{ }\PY{n}{street}\PY{p}{,}\PY{+w}{ }\PY{n}{String}\PY{+w}{ }\PY{n}{city}\PY{p}{,}\PY{+w}{ }\PY{n}{String}\PY{+w}{ }\PY{n}{zip}\PY{p}{)}\PY{+w}{ }\PY{p}{\PYZob{}}\PY{p}{\PYZcb{}}

\PY{k+kt}{void}\PY{+w}{ }\PY{n+nf}{shipOrder}\PY{p}{(}\PY{n}{Address}\PY{+w}{ }\PY{n}{address}\PY{p}{,}\PY{+w}{ }\PY{n}{Order}\PY{+w}{ }\PY{n}{order}\PY{p}{)}\PY{+w}{ }\PY{p}{\PYZob{}}\PY{+w}{ }\PY{p}{.}\PY{p}{.}\PY{p}{.}\PY{+w}{ }\PY{p}{\PYZcb{}}
\PY{k+kt}{void}\PY{+w}{ }\PY{n+nf}{billCustomer}\PY{p}{(}\PY{n}{Address}\PY{+w}{ }\PY{n}{address}\PY{p}{,}\PY{+w}{ }\PY{n}{Customer}\PY{+w}{ }\PY{n}{customer}\PY{p}{)}\PY{+w}{ }\PY{p}{\PYZob{}}\PY{+w}{ }\PY{p}{.}\PY{p}{.}\PY{p}{.}\PY{+w}{ }\PY{p}{\PYZcb{}}
\end{Verbatim}
\end{codebox}

\subsection[Shotgun Surgery]{Shotgun Surgery}
\label{h:19:shotgun-surgery}
\begin{codebox}
\begin{Verbatim}[commandchars=\\\{\}]
\PY{c+c1}{// Before: \PYZdq{}tax rate\PYZdq{} logic is duplicated in three unrelated classes.}
\PY{c+c1}{// Changing the tax rule means hunting down and editing all three.}
\PY{k+kd}{class} \PY{n+nc}{InvoiceService}\PY{+w}{ }\PY{p}{\PYZob{}}\PY{+w}{ }\PY{n}{BigDecimal}\PY{+w}{ }\PY{n+nf}{tax}\PY{p}{(}\PY{n}{BigDecimal}\PY{+w}{ }\PY{n}{amount}\PY{p}{)}\PY{+w}{ }\PY{p}{\PYZob{}}\PY{+w}{ }\PY{k}{return}\PY{+w}{ }\PY{n}{amount}\PY{p}{.}\PY{n+na}{multiply}\PY{p}{(}\PY{n}{BigDecimal}\PY{p}{.}\PY{n+na}{valueOf}\PY{p}{(}\PY{l+m+mf}{0.08}\PY{p}{)}\PY{p}{)}\PY{p}{;}\PY{+w}{ }\PY{p}{\PYZcb{}}\PY{+w}{ }\PY{p}{\PYZcb{}}
\PY{k+kd}{class} \PY{n+nc}{ReceiptService}\PY{+w}{ }\PY{p}{\PYZob{}}\PY{+w}{ }\PY{n}{BigDecimal}\PY{+w}{ }\PY{n+nf}{tax}\PY{p}{(}\PY{n}{BigDecimal}\PY{+w}{ }\PY{n}{amount}\PY{p}{)}\PY{+w}{ }\PY{p}{\PYZob{}}\PY{+w}{ }\PY{k}{return}\PY{+w}{ }\PY{n}{amount}\PY{p}{.}\PY{n+na}{multiply}\PY{p}{(}\PY{n}{BigDecimal}\PY{p}{.}\PY{n+na}{valueOf}\PY{p}{(}\PY{l+m+mf}{0.08}\PY{p}{)}\PY{p}{)}\PY{p}{;}\PY{+w}{ }\PY{p}{\PYZcb{}}\PY{+w}{ }\PY{p}{\PYZcb{}}
\PY{k+kd}{class} \PY{n+nc}{ReportService}\PY{+w}{  }\PY{p}{\PYZob{}}\PY{+w}{ }\PY{n}{BigDecimal}\PY{+w}{ }\PY{n+nf}{tax}\PY{p}{(}\PY{n}{BigDecimal}\PY{+w}{ }\PY{n}{amount}\PY{p}{)}\PY{+w}{ }\PY{p}{\PYZob{}}\PY{+w}{ }\PY{k}{return}\PY{+w}{ }\PY{n}{amount}\PY{p}{.}\PY{n+na}{multiply}\PY{p}{(}\PY{n}{BigDecimal}\PY{p}{.}\PY{n+na}{valueOf}\PY{p}{(}\PY{l+m+mf}{0.08}\PY{p}{)}\PY{p}{)}\PY{p}{;}\PY{+w}{ }\PY{p}{\PYZcb{}}\PY{+w}{ }\PY{p}{\PYZcb{}}
\end{Verbatim}
\end{codebox}

\begin{codebox}
\begin{Verbatim}[commandchars=\\\{\}]
\PY{c+c1}{// After: one authoritative place to change}
\PY{k+kd}{class} \PY{n+nc}{TaxCalculator}\PY{+w}{ }\PY{p}{\PYZob{}}
\PY{+w}{    }\PY{k+kd}{private}\PY{+w}{ }\PY{k+kd}{static}\PY{+w}{ }\PY{k+kd}{final}\PY{+w}{ }\PY{n}{BigDecimal}\PY{+w}{ }\PY{n}{RATE}\PY{+w}{ }\PY{o}{=}\PY{+w}{ }\PY{n}{BigDecimal}\PY{p}{.}\PY{n+na}{valueOf}\PY{p}{(}\PY{l+m+mf}{0.08}\PY{p}{)}\PY{p}{;}
\PY{+w}{    }\PY{n}{BigDecimal}\PY{+w}{ }\PY{n+nf}{tax}\PY{p}{(}\PY{n}{BigDecimal}\PY{+w}{ }\PY{n}{amount}\PY{p}{)}\PY{+w}{ }\PY{p}{\PYZob{}}\PY{+w}{ }\PY{k}{return}\PY{+w}{ }\PY{n}{amount}\PY{p}{.}\PY{n+na}{multiply}\PY{p}{(}\PY{n}{RATE}\PY{p}{)}\PY{p}{;}\PY{+w}{ }\PY{p}{\PYZcb{}}
\PY{p}{\PYZcb{}}
\PY{c+c1}{// InvoiceService, ReceiptService, ReportService all delegate to TaxCalculator}
\end{Verbatim}
\end{codebox}

\subsection[Anemic Domain Model]{Anemic Domain Model}
\label{h:19:anemic-domain-model}
\begin{codebox}
\begin{Verbatim}[commandchars=\\\{\}]
\PY{c+c1}{// Before: Order is just a data bag; all logic lives in a separate \PYZdq{}service\PYZdq{}}
\PY{k+kd}{class} \PY{n+nc}{Order}\PY{+w}{ }\PY{p}{\PYZob{}}
\PY{+w}{    }\PY{n}{List}\PY{o}{\PYZlt{}}\PY{n}{OrderItem}\PY{o}{\PYZgt{}}\PY{+w}{ }\PY{n}{items}\PY{p}{;}
\PY{+w}{    }\PY{n}{List}\PY{o}{\PYZlt{}}\PY{n}{OrderItem}\PY{o}{\PYZgt{}}\PY{+w}{ }\PY{n+nf}{getItems}\PY{p}{(}\PY{p}{)}\PY{+w}{ }\PY{p}{\PYZob{}}\PY{+w}{ }\PY{k}{return}\PY{+w}{ }\PY{n}{items}\PY{p}{;}\PY{+w}{ }\PY{p}{\PYZcb{}}
\PY{+w}{    }\PY{k+kt}{void}\PY{+w}{ }\PY{n+nf}{setItems}\PY{p}{(}\PY{n}{List}\PY{o}{\PYZlt{}}\PY{n}{OrderItem}\PY{o}{\PYZgt{}}\PY{+w}{ }\PY{n}{items}\PY{p}{)}\PY{+w}{ }\PY{p}{\PYZob{}}\PY{+w}{ }\PY{k}{this}\PY{p}{.}\PY{n+na}{items}\PY{+w}{ }\PY{o}{=}\PY{+w}{ }\PY{n}{items}\PY{p}{;}\PY{+w}{ }\PY{p}{\PYZcb{}}
\PY{p}{\PYZcb{}}

\PY{k+kd}{class} \PY{n+nc}{OrderService}\PY{+w}{ }\PY{p}{\PYZob{}}
\PY{+w}{    }\PY{n}{BigDecimal}\PY{+w}{ }\PY{n+nf}{calculateTotal}\PY{p}{(}\PY{n}{Order}\PY{+w}{ }\PY{n}{order}\PY{p}{)}\PY{+w}{ }\PY{p}{\PYZob{}}
\PY{+w}{        }\PY{k}{return}\PY{+w}{ }\PY{n}{order}\PY{p}{.}\PY{n+na}{getItems}\PY{p}{(}\PY{p}{)}\PY{p}{.}\PY{n+na}{stream}\PY{p}{(}\PY{p}{)}
\PY{+w}{                }\PY{p}{.}\PY{n+na}{map}\PY{p}{(}\PY{n}{i}\PY{+w}{ }\PY{o}{\PYZhy{}}\PY{o}{\PYZgt{}}\PY{+w}{ }\PY{n}{i}\PY{p}{.}\PY{n+na}{getPrice}\PY{p}{(}\PY{p}{)}\PY{p}{.}\PY{n+na}{multiply}\PY{p}{(}\PY{n}{BigDecimal}\PY{p}{.}\PY{n+na}{valueOf}\PY{p}{(}\PY{n}{i}\PY{p}{.}\PY{n+na}{getQuantity}\PY{p}{(}\PY{p}{)}\PY{p}{)}\PY{p}{)}\PY{p}{)}
\PY{+w}{                }\PY{p}{.}\PY{n+na}{reduce}\PY{p}{(}\PY{n}{BigDecimal}\PY{p}{.}\PY{n+na}{ZERO}\PY{p}{,}\PY{+w}{ }\PY{n}{BigDecimal}\PY{p}{:}\PY{p}{:}\PY{n}{add}\PY{p}{)}\PY{p}{;}
\PY{+w}{    }\PY{p}{\PYZcb{}}
\PY{p}{\PYZcb{}}
\end{Verbatim}
\end{codebox}

\begin{codebox}
\begin{Verbatim}[commandchars=\\\{\}]
\PY{c+c1}{// After: behavior moves into the domain object it belongs to}
\PY{k+kd}{class} \PY{n+nc}{Order}\PY{+w}{ }\PY{p}{\PYZob{}}
\PY{+w}{    }\PY{k+kd}{private}\PY{+w}{ }\PY{k+kd}{final}\PY{+w}{ }\PY{n}{List}\PY{o}{\PYZlt{}}\PY{n}{OrderItem}\PY{o}{\PYZgt{}}\PY{+w}{ }\PY{n}{items}\PY{p}{;}
\PY{+w}{    }\PY{n}{Order}\PY{p}{(}\PY{n}{List}\PY{o}{\PYZlt{}}\PY{n}{OrderItem}\PY{o}{\PYZgt{}}\PY{+w}{ }\PY{n}{items}\PY{p}{)}\PY{+w}{ }\PY{p}{\PYZob{}}\PY{+w}{ }\PY{k}{this}\PY{p}{.}\PY{n+na}{items}\PY{+w}{ }\PY{o}{=}\PY{+w}{ }\PY{n}{items}\PY{p}{;}\PY{+w}{ }\PY{p}{\PYZcb{}}

\PY{+w}{    }\PY{n}{BigDecimal}\PY{+w}{ }\PY{n+nf}{calculateTotal}\PY{p}{(}\PY{p}{)}\PY{+w}{ }\PY{p}{\PYZob{}}
\PY{+w}{        }\PY{k}{return}\PY{+w}{ }\PY{n}{items}\PY{p}{.}\PY{n+na}{stream}\PY{p}{(}\PY{p}{)}
\PY{+w}{                }\PY{p}{.}\PY{n+na}{map}\PY{p}{(}\PY{n}{i}\PY{+w}{ }\PY{o}{\PYZhy{}}\PY{o}{\PYZgt{}}\PY{+w}{ }\PY{n}{i}\PY{p}{.}\PY{n+na}{getPrice}\PY{p}{(}\PY{p}{)}\PY{p}{.}\PY{n+na}{multiply}\PY{p}{(}\PY{n}{BigDecimal}\PY{p}{.}\PY{n+na}{valueOf}\PY{p}{(}\PY{n}{i}\PY{p}{.}\PY{n+na}{getQuantity}\PY{p}{(}\PY{p}{)}\PY{p}{)}\PY{p}{)}\PY{p}{)}
\PY{+w}{                }\PY{p}{.}\PY{n+na}{reduce}\PY{p}{(}\PY{n}{BigDecimal}\PY{p}{.}\PY{n+na}{ZERO}\PY{p}{,}\PY{+w}{ }\PY{n}{BigDecimal}\PY{p}{:}\PY{p}{:}\PY{n}{add}\PY{p}{)}\PY{p}{;}
\PY{+w}{    }\PY{p}{\PYZcb{}}
\PY{p}{\PYZcb{}}
\end{Verbatim}
\end{codebox}

\subsection[Boolean Parameters]{Boolean Parameters}
\label{h:19:boolean-parameters}
\begin{codebox}
\begin{Verbatim}[commandchars=\\\{\}]
\PY{c+c1}{// Before}
\PY{k+kt}{void}\PY{+w}{ }\PY{n+nf}{search}\PY{p}{(}\PY{n}{String}\PY{+w}{ }\PY{n}{query}\PY{p}{,}\PY{+w}{ }\PY{k+kt}{boolean}\PY{+w}{ }\PY{n}{caseSensitive}\PY{p}{)}\PY{+w}{ }\PY{p}{\PYZob{}}\PY{+w}{ }\PY{p}{.}\PY{p}{.}\PY{p}{.}\PY{+w}{ }\PY{p}{\PYZcb{}}
\PY{n}{search}\PY{p}{(}\PY{l+s}{\PYZdq{}}\PY{l+s}{Java}\PY{l+s}{\PYZdq{}}\PY{p}{,}\PY{+w}{ }\PY{k+kc}{true}\PY{p}{)}\PY{p}{;}\PY{+w}{ }\PY{c+c1}{// true meaning what, exactly, at a glance?}
\end{Verbatim}
\end{codebox}

\begin{codebox}
\begin{Verbatim}[commandchars=\\\{\}]
\PY{c+c1}{// After: an enum documents itself at the call site}
\PY{k+kd}{enum}\PY{+w}{ }\PY{n}{CaseMode}\PY{+w}{ }\PY{p}{\PYZob{}}\PY{+w}{ }\PY{n}{SENSITIVE}\PY{p}{,}\PY{+w}{ }\PY{n}{INSENSITIVE}\PY{+w}{ }\PY{p}{\PYZcb{}}
\PY{k+kt}{void}\PY{+w}{ }\PY{n+nf}{search}\PY{p}{(}\PY{n}{String}\PY{+w}{ }\PY{n}{query}\PY{p}{,}\PY{+w}{ }\PY{n}{CaseMode}\PY{+w}{ }\PY{n}{caseMode}\PY{p}{)}\PY{+w}{ }\PY{p}{\PYZob{}}\PY{+w}{ }\PY{p}{.}\PY{p}{.}\PY{p}{.}\PY{+w}{ }\PY{p}{\PYZcb{}}
\PY{n}{search}\PY{p}{(}\PY{l+s}{\PYZdq{}}\PY{l+s}{Java}\PY{l+s}{\PYZdq{}}\PY{p}{,}\PY{+w}{ }\PY{n}{CaseMode}\PY{p}{.}\PY{n+na}{SENSITIVE}\PY{p}{)}\PY{p}{;}
\end{Verbatim}
\end{codebox}

\subsection[Deep Inheritance]{Deep Inheritance}
\label{h:19:deep-inheritance}
\begin{codebox}
\begin{Verbatim}[commandchars=\\\{\}]
\PY{c+c1}{// Before: five levels deep, unclear which ancestor defines what}
\PY{k+kd}{class} \PY{n+nc}{Animal}\PY{+w}{ }\PY{p}{\PYZob{}}\PY{p}{\PYZcb{}}
\PY{k+kd}{class} \PY{n+nc}{Mammal}\PY{+w}{ }\PY{k+kd}{extends}\PY{+w}{ }\PY{n}{Animal}\PY{+w}{ }\PY{p}{\PYZob{}}\PY{p}{\PYZcb{}}
\PY{k+kd}{class} \PY{n+nc}{Carnivore}\PY{+w}{ }\PY{k+kd}{extends}\PY{+w}{ }\PY{n}{Mammal}\PY{+w}{ }\PY{p}{\PYZob{}}\PY{p}{\PYZcb{}}
\PY{k+kd}{class} \PY{n+nc}{Feline}\PY{+w}{ }\PY{k+kd}{extends}\PY{+w}{ }\PY{n}{Carnivore}\PY{+w}{ }\PY{p}{\PYZob{}}\PY{p}{\PYZcb{}}
\PY{k+kd}{class} \PY{n+nc}{Cat}\PY{+w}{ }\PY{k+kd}{extends}\PY{+w}{ }\PY{n}{Feline}\PY{+w}{ }\PY{p}{\PYZob{}}\PY{p}{\PYZcb{}}
\end{Verbatim}
\end{codebox}

\begin{codebox}
\begin{Verbatim}[commandchars=\\\{\}]
\PY{c+c1}{// After: flatten using composition — combine behaviors instead of stacking classes}
\PY{k+kd}{interface} \PY{n+nc}{Diet}\PY{+w}{ }\PY{p}{\PYZob{}}\PY{+w}{ }\PY{k+kt}{void}\PY{+w}{ }\PY{n+nf}{eat}\PY{p}{(}\PY{p}{)}\PY{p}{;}\PY{+w}{ }\PY{p}{\PYZcb{}}
\PY{k+kd}{interface} \PY{n+nc}{Movement}\PY{+w}{ }\PY{p}{\PYZob{}}\PY{+w}{ }\PY{k+kt}{void}\PY{+w}{ }\PY{n+nf}{move}\PY{p}{(}\PY{p}{)}\PY{p}{;}\PY{+w}{ }\PY{p}{\PYZcb{}}

\PY{k+kd}{class} \PY{n+nc}{Cat}\PY{+w}{ }\PY{p}{\PYZob{}}
\PY{+w}{    }\PY{k+kd}{private}\PY{+w}{ }\PY{k+kd}{final}\PY{+w}{ }\PY{n}{Diet}\PY{+w}{ }\PY{n}{diet}\PY{p}{;}
\PY{+w}{    }\PY{k+kd}{private}\PY{+w}{ }\PY{k+kd}{final}\PY{+w}{ }\PY{n}{Movement}\PY{+w}{ }\PY{n}{movement}\PY{p}{;}
\PY{+w}{    }\PY{n}{Cat}\PY{p}{(}\PY{n}{Diet}\PY{+w}{ }\PY{n}{diet}\PY{p}{,}\PY{+w}{ }\PY{n}{Movement}\PY{+w}{ }\PY{n}{movement}\PY{p}{)}\PY{+w}{ }\PY{p}{\PYZob{}}\PY{+w}{ }\PY{k}{this}\PY{p}{.}\PY{n+na}{diet}\PY{+w}{ }\PY{o}{=}\PY{+w}{ }\PY{n}{diet}\PY{p}{;}\PY{+w}{ }\PY{k}{this}\PY{p}{.}\PY{n+na}{movement}\PY{+w}{ }\PY{o}{=}\PY{+w}{ }\PY{n}{movement}\PY{p}{;}\PY{+w}{ }\PY{p}{\PYZcb{}}
\PY{p}{\PYZcb{}}
\end{Verbatim}
\end{codebox}

\subsection[Utility-Class Dumping Ground]{Utility-Class Dumping Ground}
\label{h:19:utility-class-dumping-ground}
\begin{codebox}
\begin{Verbatim}[commandchars=\\\{\}]
\PY{c+c1}{// Before: unrelated static methods dumped into one \PYZdq{}Utils\PYZdq{} class}
\PY{k+kd}{class} \PY{n+nc}{Utils}\PY{+w}{ }\PY{p}{\PYZob{}}
\PY{+w}{    }\PY{k+kd}{static}\PY{+w}{ }\PY{n}{String}\PY{+w}{ }\PY{n+nf}{formatDate}\PY{p}{(}\PY{n}{LocalDate}\PY{+w}{ }\PY{n}{d}\PY{p}{)}\PY{+w}{ }\PY{p}{\PYZob{}}\PY{+w}{ }\PY{p}{.}\PY{p}{.}\PY{p}{.}\PY{+w}{ }\PY{p}{\PYZcb{}}
\PY{+w}{    }\PY{k+kd}{static}\PY{+w}{ }\PY{n}{BigDecimal}\PY{+w}{ }\PY{n+nf}{calculateTax}\PY{p}{(}\PY{n}{BigDecimal}\PY{+w}{ }\PY{n}{amount}\PY{p}{)}\PY{+w}{ }\PY{p}{\PYZob{}}\PY{+w}{ }\PY{p}{.}\PY{p}{.}\PY{p}{.}\PY{+w}{ }\PY{p}{\PYZcb{}}
\PY{+w}{    }\PY{k+kd}{static}\PY{+w}{ }\PY{k+kt}{boolean}\PY{+w}{ }\PY{n+nf}{isValidEmail}\PY{p}{(}\PY{n}{String}\PY{+w}{ }\PY{n}{s}\PY{p}{)}\PY{+w}{ }\PY{p}{\PYZob{}}\PY{+w}{ }\PY{p}{.}\PY{p}{.}\PY{p}{.}\PY{+w}{ }\PY{p}{\PYZcb{}}
\PY{p}{\PYZcb{}}
\end{Verbatim}
\end{codebox}

\begin{codebox}
\begin{Verbatim}[commandchars=\\\{\}]
\PY{c+c1}{// After: each concern gets its own focused, named home}
\PY{k+kd}{class} \PY{n+nc}{DateFormatting}\PY{+w}{ }\PY{p}{\PYZob{}}\PY{+w}{ }\PY{k+kd}{static}\PY{+w}{ }\PY{n}{String}\PY{+w}{ }\PY{n+nf}{format}\PY{p}{(}\PY{n}{LocalDate}\PY{+w}{ }\PY{n}{d}\PY{p}{)}\PY{+w}{ }\PY{p}{\PYZob{}}\PY{+w}{ }\PY{p}{.}\PY{p}{.}\PY{p}{.}\PY{+w}{ }\PY{p}{\PYZcb{}}\PY{+w}{ }\PY{p}{\PYZcb{}}
\PY{k+kd}{class} \PY{n+nc}{TaxCalculator}\PY{+w}{ }\PY{p}{\PYZob{}}\PY{+w}{ }\PY{k+kd}{static}\PY{+w}{ }\PY{n}{BigDecimal}\PY{+w}{ }\PY{n+nf}{calculate}\PY{p}{(}\PY{n}{BigDecimal}\PY{+w}{ }\PY{n}{amount}\PY{p}{)}\PY{+w}{ }\PY{p}{\PYZob{}}\PY{+w}{ }\PY{p}{.}\PY{p}{.}\PY{p}{.}\PY{+w}{ }\PY{p}{\PYZcb{}}\PY{+w}{ }\PY{p}{\PYZcb{}}
\PY{k+kd}{class} \PY{n+nc}{EmailValidator}\PY{+w}{ }\PY{p}{\PYZob{}}\PY{+w}{ }\PY{k+kd}{static}\PY{+w}{ }\PY{k+kt}{boolean}\PY{+w}{ }\PY{n+nf}{isValid}\PY{p}{(}\PY{n}{String}\PY{+w}{ }\PY{n}{s}\PY{p}{)}\PY{+w}{ }\PY{p}{\PYZob{}}\PY{+w}{ }\PY{p}{.}\PY{p}{.}\PY{p}{.}\PY{+w}{ }\PY{p}{\PYZcb{}}\PY{+w}{ }\PY{p}{\PYZcb{}}
\end{Verbatim}
\end{codebox}

\subsection[Stringly-Typed Code]{Stringly-Typed Code}
\label{h:19:stringly-typed-code}
\begin{codebox}
\begin{Verbatim}[commandchars=\\\{\}]
\PY{c+c1}{// Before: a \PYZdq{}status\PYZdq{} that is really a fixed set of values, stored as a raw String}
\PY{k+kt}{void}\PY{+w}{ }\PY{n+nf}{updateStatus}\PY{p}{(}\PY{n}{String}\PY{+w}{ }\PY{n}{status}\PY{p}{)}\PY{+w}{ }\PY{p}{\PYZob{}}
\PY{+w}{    }\PY{k}{if}\PY{+w}{ }\PY{p}{(}\PY{n}{status}\PY{p}{.}\PY{n+na}{equals}\PY{p}{(}\PY{l+s}{\PYZdq{}}\PY{l+s}{PENDING}\PY{l+s}{\PYZdq{}}\PY{p}{)}\PY{p}{)}\PY{+w}{ }\PY{p}{\PYZob{}}\PY{+w}{ }\PY{p}{.}\PY{p}{.}\PY{p}{.}\PY{+w}{ }\PY{p}{\PYZcb{}}
\PY{+w}{    }\PY{k}{else}\PY{+w}{ }\PY{k}{if}\PY{+w}{ }\PY{p}{(}\PY{n}{status}\PY{p}{.}\PY{n+na}{equals}\PY{p}{(}\PY{l+s}{\PYZdq{}}\PY{l+s}{Pending}\PY{l+s}{\PYZdq{}}\PY{p}{)}\PY{p}{)}\PY{+w}{ }\PY{p}{\PYZob{}}\PY{+w}{ }\PY{p}{.}\PY{p}{.}\PY{p}{.}\PY{+w}{ }\PY{p}{\PYZcb{}}\PY{+w}{ }\PY{c+c1}{// typo/case variant silently does nothing}
\PY{p}{\PYZcb{}}
\end{Verbatim}
\end{codebox}

\begin{codebox}
\begin{Verbatim}[commandchars=\\\{\}]
\PY{c+c1}{// After: an enum makes invalid values impossible to represent}
\PY{k+kd}{enum}\PY{+w}{ }\PY{n}{OrderStatus}\PY{+w}{ }\PY{p}{\PYZob{}}\PY{+w}{ }\PY{n}{PENDING}\PY{p}{,}\PY{+w}{ }\PY{n}{SHIPPED}\PY{p}{,}\PY{+w}{ }\PY{n}{DELIVERED}\PY{p}{,}\PY{+w}{ }\PY{n}{CANCELLED}\PY{+w}{ }\PY{p}{\PYZcb{}}
\PY{k+kt}{void}\PY{+w}{ }\PY{n+nf}{updateStatus}\PY{p}{(}\PY{n}{OrderStatus}\PY{+w}{ }\PY{n}{status}\PY{p}{)}\PY{+w}{ }\PY{p}{\PYZob{}}
\PY{+w}{    }\PY{k}{switch}\PY{+w}{ }\PY{p}{(}\PY{n}{status}\PY{p}{)}\PY{+w}{ }\PY{p}{\PYZob{}}
\PY{+w}{        }\PY{k}{case}\PY{+w}{ }\PY{n}{PENDING}\PY{+w}{ }\PY{o}{\PYZhy{}}\PY{o}{\PYZgt{}}\PY{+w}{ }\PY{p}{.}\PY{p}{.}\PY{p}{.}\PY{p}{;}
\PY{+w}{        }\PY{k}{case}\PY{+w}{ }\PY{n}{SHIPPED}\PY{+w}{ }\PY{o}{\PYZhy{}}\PY{o}{\PYZgt{}}\PY{+w}{ }\PY{p}{.}\PY{p}{.}\PY{p}{.}\PY{p}{;}
\PY{+w}{        }\PY{k}{default}\PY{+w}{ }\PY{o}{\PYZhy{}}\PY{o}{\PYZgt{}}\PY{+w}{ }\PY{p}{.}\PY{p}{.}\PY{p}{.}\PY{p}{;}
\PY{+w}{    }\PY{p}{\PYZcb{}}
\PY{p}{\PYZcb{}}
\end{Verbatim}
\end{codebox}

\subsection[Exception as Control Flow]{Exception as Control Flow}
\label{h:19:exception-as-control-flow}
\begin{codebox}
\begin{Verbatim}[commandchars=\\\{\}]
\PY{c+c1}{// Before: uses an exception to detect \PYZdq{}not found\PYZdq{} during normal, expected flow}
\PY{k+kt}{int}\PY{+w}{ }\PY{n+nf}{findIndex}\PY{p}{(}\PY{n}{List}\PY{o}{\PYZlt{}}\PY{n}{String}\PY{o}{\PYZgt{}}\PY{+w}{ }\PY{n}{list}\PY{p}{,}\PY{+w}{ }\PY{n}{String}\PY{+w}{ }\PY{n}{target}\PY{p}{)}\PY{+w}{ }\PY{p}{\PYZob{}}
\PY{+w}{    }\PY{k}{try}\PY{+w}{ }\PY{p}{\PYZob{}}
\PY{+w}{        }\PY{k}{for}\PY{+w}{ }\PY{p}{(}\PY{k+kt}{int}\PY{+w}{ }\PY{n}{i}\PY{+w}{ }\PY{o}{=}\PY{+w}{ }\PY{l+m+mi}{0}\PY{p}{;}\PY{+w}{ }\PY{n}{i}\PY{+w}{ }\PY{o}{\PYZlt{}}\PY{+w}{ }\PY{n}{list}\PY{p}{.}\PY{n+na}{size}\PY{p}{(}\PY{p}{)}\PY{p}{;}\PY{+w}{ }\PY{n}{i}\PY{o}{+}\PY{o}{+}\PY{p}{)}\PY{+w}{ }\PY{p}{\PYZob{}}
\PY{+w}{            }\PY{k}{if}\PY{+w}{ }\PY{p}{(}\PY{n}{list}\PY{p}{.}\PY{n+na}{get}\PY{p}{(}\PY{n}{i}\PY{p}{)}\PY{p}{.}\PY{n+na}{equals}\PY{p}{(}\PY{n}{target}\PY{p}{)}\PY{p}{)}\PY{+w}{ }\PY{k}{throw}\PY{+w}{ }\PY{k}{new}\PY{+w}{ }\PY{n}{FoundException}\PY{p}{(}\PY{n}{i}\PY{p}{)}\PY{p}{;}
\PY{+w}{        }\PY{p}{\PYZcb{}}
\PY{+w}{    }\PY{p}{\PYZcb{}}\PY{+w}{ }\PY{k}{catch}\PY{+w}{ }\PY{p}{(}\PY{n}{FoundException}\PY{+w}{ }\PY{n}{e}\PY{p}{)}\PY{+w}{ }\PY{p}{\PYZob{}}
\PY{+w}{        }\PY{k}{return}\PY{+w}{ }\PY{n}{e}\PY{p}{.}\PY{n+na}{index}\PY{p}{;}
\PY{+w}{    }\PY{p}{\PYZcb{}}
\PY{+w}{    }\PY{k}{return}\PY{+w}{ }\PY{o}{\PYZhy{}}\PY{l+m+mi}{1}\PY{p}{;}
\PY{p}{\PYZcb{}}
\end{Verbatim}
\end{codebox}

\begin{codebox}
\begin{Verbatim}[commandchars=\\\{\}]
\PY{c+c1}{// After: a normal conditional; exceptions are reserved for actual errors}
\PY{k+kt}{int}\PY{+w}{ }\PY{n+nf}{findIndex}\PY{p}{(}\PY{n}{List}\PY{o}{\PYZlt{}}\PY{n}{String}\PY{o}{\PYZgt{}}\PY{+w}{ }\PY{n}{list}\PY{p}{,}\PY{+w}{ }\PY{n}{String}\PY{+w}{ }\PY{n}{target}\PY{p}{)}\PY{+w}{ }\PY{p}{\PYZob{}}
\PY{+w}{    }\PY{k}{for}\PY{+w}{ }\PY{p}{(}\PY{k+kt}{int}\PY{+w}{ }\PY{n}{i}\PY{+w}{ }\PY{o}{=}\PY{+w}{ }\PY{l+m+mi}{0}\PY{p}{;}\PY{+w}{ }\PY{n}{i}\PY{+w}{ }\PY{o}{\PYZlt{}}\PY{+w}{ }\PY{n}{list}\PY{p}{.}\PY{n+na}{size}\PY{p}{(}\PY{p}{)}\PY{p}{;}\PY{+w}{ }\PY{n}{i}\PY{o}{+}\PY{o}{+}\PY{p}{)}\PY{+w}{ }\PY{p}{\PYZob{}}
\PY{+w}{        }\PY{k}{if}\PY{+w}{ }\PY{p}{(}\PY{n}{list}\PY{p}{.}\PY{n+na}{get}\PY{p}{(}\PY{n}{i}\PY{p}{)}\PY{p}{.}\PY{n+na}{equals}\PY{p}{(}\PY{n}{target}\PY{p}{)}\PY{p}{)}\PY{+w}{ }\PY{k}{return}\PY{+w}{ }\PY{n}{i}\PY{p}{;}
\PY{+w}{    }\PY{p}{\PYZcb{}}
\PY{+w}{    }\PY{k}{return}\PY{+w}{ }\PY{o}{\PYZhy{}}\PY{l+m+mi}{1}\PY{p}{;}
\PY{p}{\PYZcb{}}
\end{Verbatim}
\end{codebox}

\subsection[Static Abuse]{Static Abuse}
\label{h:19:static-abuse}
\begin{codebox}
\begin{Verbatim}[commandchars=\\\{\}]
\PY{c+c1}{// Before: hidden global mutable state via static fields — hard to test in isolation}
\PY{k+kd}{class} \PY{n+nc}{OrderCounter}\PY{+w}{ }\PY{p}{\PYZob{}}
\PY{+w}{    }\PY{k+kd}{static}\PY{+w}{ }\PY{k+kt}{int}\PY{+w}{ }\PY{n}{count}\PY{+w}{ }\PY{o}{=}\PY{+w}{ }\PY{l+m+mi}{0}\PY{p}{;}
\PY{+w}{    }\PY{k+kd}{static}\PY{+w}{ }\PY{k+kt}{void}\PY{+w}{ }\PY{n+nf}{increment}\PY{p}{(}\PY{p}{)}\PY{+w}{ }\PY{p}{\PYZob{}}\PY{+w}{ }\PY{n}{count}\PY{o}{+}\PY{o}{+}\PY{p}{;}\PY{+w}{ }\PY{p}{\PYZcb{}}
\PY{p}{\PYZcb{}}
\end{Verbatim}
\end{codebox}

\begin{codebox}
\begin{Verbatim}[commandchars=\\\{\}]
\PY{c+c1}{// After: instance state, injected where needed, easy to isolate in tests}
\PY{k+kd}{class} \PY{n+nc}{OrderCounter}\PY{+w}{ }\PY{p}{\PYZob{}}
\PY{+w}{    }\PY{k+kd}{private}\PY{+w}{ }\PY{k+kt}{int}\PY{+w}{ }\PY{n}{count}\PY{+w}{ }\PY{o}{=}\PY{+w}{ }\PY{l+m+mi}{0}\PY{p}{;}
\PY{+w}{    }\PY{k+kt}{void}\PY{+w}{ }\PY{n+nf}{increment}\PY{p}{(}\PY{p}{)}\PY{+w}{ }\PY{p}{\PYZob{}}\PY{+w}{ }\PY{n}{count}\PY{o}{+}\PY{o}{+}\PY{p}{;}\PY{+w}{ }\PY{p}{\PYZcb{}}
\PY{+w}{    }\PY{k+kt}{int}\PY{+w}{ }\PY{n+nf}{get}\PY{p}{(}\PY{p}{)}\PY{+w}{ }\PY{p}{\PYZob{}}\PY{+w}{ }\PY{k}{return}\PY{+w}{ }\PY{n}{count}\PY{p}{;}\PY{+w}{ }\PY{p}{\PYZcb{}}
\PY{p}{\PYZcb{}}
\end{Verbatim}
\end{codebox}

\subsection[Over-Mocking]{Over-Mocking}
\label{h:19:over-mocking}
\begin{codebox}
\begin{Verbatim}[commandchars=\\\{\}]
\PY{c+c1}{// Before: test mocks every internal collaborator and asserts on call order,}
\PY{c+c1}{// so it breaks the moment implementation details change even if behavior doesn\PYZsq{}t}
\PY{n+nd}{@Test}
\PY{k+kt}{void}\PY{+w}{ }\PY{n+nf}{placeOrder\PYZus{}callsCollaboratorsInOrder}\PY{p}{(}\PY{p}{)}\PY{+w}{ }\PY{p}{\PYZob{}}
\PY{+w}{    }\PY{n}{OrderService}\PY{+w}{ }\PY{n}{service}\PY{+w}{ }\PY{o}{=}\PY{+w}{ }\PY{k}{new}\PY{+w}{ }\PY{n}{OrderService}\PY{p}{(}\PY{n}{mockInventory}\PY{p}{,}\PY{+w}{ }\PY{n}{mockPayment}\PY{p}{,}\PY{+w}{ }\PY{n}{mockShipping}\PY{p}{)}\PY{p}{;}
\PY{+w}{    }\PY{n}{service}\PY{p}{.}\PY{n+na}{placeOrder}\PY{p}{(}\PY{n}{order}\PY{p}{)}\PY{p}{;}
\PY{+w}{    }\PY{n}{InOrder}\PY{+w}{ }\PY{n}{inOrder}\PY{+w}{ }\PY{o}{=}\PY{+w}{ }\PY{n}{inOrder}\PY{p}{(}\PY{n}{mockInventory}\PY{p}{,}\PY{+w}{ }\PY{n}{mockPayment}\PY{p}{,}\PY{+w}{ }\PY{n}{mockShipping}\PY{p}{)}\PY{p}{;}
\PY{+w}{    }\PY{n}{inOrder}\PY{p}{.}\PY{n+na}{verify}\PY{p}{(}\PY{n}{mockInventory}\PY{p}{)}\PY{p}{.}\PY{n+na}{reserve}\PY{p}{(}\PY{n}{order}\PY{p}{)}\PY{p}{;}
\PY{+w}{    }\PY{n}{inOrder}\PY{p}{.}\PY{n+na}{verify}\PY{p}{(}\PY{n}{mockPayment}\PY{p}{)}\PY{p}{.}\PY{n+na}{charge}\PY{p}{(}\PY{n}{any}\PY{p}{(}\PY{p}{)}\PY{p}{,}\PY{+w}{ }\PY{n}{any}\PY{p}{(}\PY{p}{)}\PY{p}{)}\PY{p}{;}
\PY{+w}{    }\PY{n}{inOrder}\PY{p}{.}\PY{n+na}{verify}\PY{p}{(}\PY{n}{mockShipping}\PY{p}{)}\PY{p}{.}\PY{n+na}{schedule}\PY{p}{(}\PY{n}{order}\PY{p}{)}\PY{p}{;}
\PY{p}{\PYZcb{}}
\end{Verbatim}
\end{codebox}

\begin{codebox}
\begin{Verbatim}[commandchars=\\\{\}]
\PY{c+c1}{// After: test the observable behavior/contract, not the internal wiring}
\PY{n+nd}{@Test}
\PY{k+kt}{void}\PY{+w}{ }\PY{n+nf}{placeOrder\PYZus{}reservesInventoryAndChargesCustomer}\PY{p}{(}\PY{p}{)}\PY{+w}{ }\PY{p}{\PYZob{}}
\PY{+w}{    }\PY{n}{OrderService}\PY{+w}{ }\PY{n}{service}\PY{+w}{ }\PY{o}{=}\PY{+w}{ }\PY{k}{new}\PY{+w}{ }\PY{n}{OrderService}\PY{p}{(}\PY{n}{realInventory}\PY{p}{,}\PY{+w}{ }\PY{n}{fakePayment}\PY{p}{,}\PY{+w}{ }\PY{n}{fakeShipping}\PY{p}{)}\PY{p}{;}
\PY{+w}{    }\PY{n}{service}\PY{p}{.}\PY{n+na}{placeOrder}\PY{p}{(}\PY{n}{order}\PY{p}{)}\PY{p}{;}
\PY{+w}{    }\PY{n}{assertThat}\PY{p}{(}\PY{n}{realInventory}\PY{p}{.}\PY{n+na}{isReserved}\PY{p}{(}\PY{n}{order}\PY{p}{)}\PY{p}{)}\PY{p}{.}\PY{n+na}{isTrue}\PY{p}{(}\PY{p}{)}\PY{p}{;}
\PY{+w}{    }\PY{n}{assertThat}\PY{p}{(}\PY{n}{fakePayment}\PY{p}{.}\PY{n+na}{chargedAmountFor}\PY{p}{(}\PY{n}{order}\PY{p}{)}\PY{p}{)}\PY{p}{.}\PY{n+na}{isEqualTo}\PY{p}{(}\PY{n}{order}\PY{p}{.}\PY{n+na}{calculateTotal}\PY{p}{(}\PY{p}{)}\PY{p}{)}\PY{p}{;}
\PY{p}{\PYZcb{}}
\end{Verbatim}
\end{codebox}

\section[Common Code-Review Interview Pitfalls]{Common Code-Review Interview Pitfalls}
\label{h:19:common-code-review-interview-pitfalls}
\begin{enumerate}
\item 
\textbf{Method does more than one thing, mixing abstraction levels.}
Why it matters: it forces the reader to jump between “what” and “how,” making the method hard to skim and hard to test in isolation.

\begin{codebox}
\begin{Verbatim}[commandchars=\\\{\}]
\PY{c+c1}{// Flag: validation, calculation, and I/O all inline in one method}
\PY{k+kt}{void}\PY{+w}{ }\PY{n+nf}{placeOrder}\PY{p}{(}\PY{n}{Order}\PY{+w}{ }\PY{n}{o}\PY{p}{)}\PY{+w}{ }\PY{p}{\PYZob{}}\PY{+w}{ }\PY{c+cm}{/* validate + compute + call HTTP + save */}\PY{+w}{ }\PY{p}{\PYZcb{}}
\end{Verbatim}
\end{codebox}

\item 
\textbf{Boolean parameter with no context at the call site.}
Why it matters: \code{foo(\allowbreak{}x,\allowbreak{}~\allowbreak{}true)} tells the reader nothing about what \code{true} means without opening the method.

\begin{codebox}
\begin{Verbatim}[commandchars=\\\{\}]
\PY{n}{printReport}\PY{p}{(}\PY{n}{report}\PY{p}{,}\PY{+w}{ }\PY{k+kc}{true}\PY{p}{)}\PY{p}{;}\PY{+w}{ }\PY{c+c1}{// true = ??? — prefer an enum or split methods}
\end{Verbatim}
\end{codebox}

\item 
\textbf{\code{equals()} overridden without \code{hashCode()}.}
Why it matters: breaks \code{HashMap}/\code{HashSet} — objects that are “equal” end up in different buckets and lookups silently fail.

\begin{codebox}
\begin{Verbatim}[commandchars=\\\{\}]
\PY{k+kd}{class} \PY{n+nc}{Point}\PY{+w}{ }\PY{p}{\PYZob{}}\PY{+w}{ }\PY{k+kt}{boolean}\PY{+w}{ }\PY{n+nf}{equals}\PY{p}{(}\PY{n}{Object}\PY{+w}{ }\PY{n}{o}\PY{p}{)}\PY{+w}{ }\PY{p}{\PYZob{}}\PY{+w}{ }\PY{p}{.}\PY{p}{.}\PY{p}{.}\PY{+w}{ }\PY{p}{\PYZcb{}}\PY{+w}{ }\PY{c+cm}{/* no hashCode()! */}\PY{+w}{ }\PY{p}{\PYZcb{}}
\end{Verbatim}
\end{codebox}

\item 
\textbf{Public method returns \code{null} instead of an empty collection.}
Why it matters: pushes a mandatory null-check onto every caller; someone eventually forgets and gets an NPE.

\begin{codebox}
\begin{Verbatim}[commandchars=\\\{\}]
\PY{n}{List}\PY{o}{\PYZlt{}}\PY{n}{Order}\PY{o}{\PYZgt{}}\PY{+w}{ }\PY{n+nf}{findOrders}\PY{p}{(}\PY{n}{String}\PY{+w}{ }\PY{n}{id}\PY{p}{)}\PY{+w}{ }\PY{p}{\PYZob{}}\PY{+w}{ }\PY{k}{return}\PY{+w}{ }\PY{n}{orders}\PY{p}{.}\PY{n+na}{isEmpty}\PY{p}{(}\PY{p}{)}\PY{+w}{ }\PY{o}{?}\PY{+w}{ }\PY{k+kc}{null}\PY{+w}{ }\PY{p}{:}\PY{+w}{ }\PY{n}{orders}\PY{p}{;}\PY{+w}{ }\PY{p}{\PYZcb{}}
\end{Verbatim}
\end{codebox}

\item 
\textbf{New class extends a concrete class purely for code reuse.}
Why it matters: couples the subclass to the superclass’s internals, which is exactly the fragile-base-class problem LSP and “favor composition” warn about.

\begin{codebox}
\begin{Verbatim}[commandchars=\\\{\}]
\PY{k+kd}{class} \PY{n+nc}{InstrumentedList}\PY{+w}{ }\PY{k+kd}{extends}\PY{+w}{ }\PY{n}{ArrayList}\PY{o}{\PYZlt{}}\PY{n}{String}\PY{o}{\PYZgt{}}\PY{+w}{ }\PY{p}{\PYZob{}}\PY{+w}{ }\PY{c+cm}{/* leaks ArrayList internals */}\PY{+w}{ }\PY{p}{\PYZcb{}}
\end{Verbatim}
\end{codebox}

\item 
\textbf{Singleton implemented with a lazy-checked static field instead of an enum.}
Why it matters: hand-rolled singletons are vulnerable to reflection and serialization attacks that recreate a second instance; enum singletons are not.

\begin{codebox}
\begin{Verbatim}[commandchars=\\\{\}]
\PY{k+kd}{class} \PY{n+nc}{Config}\PY{+w}{ }\PY{p}{\PYZob{}}\PY{+w}{ }\PY{k+kd}{private}\PY{+w}{ }\PY{k+kd}{static}\PY{+w}{ }\PY{n}{Config}\PY{+w}{ }\PY{n}{instance}\PY{p}{;}\PY{+w}{ }\PY{k+kd}{static}\PY{+w}{ }\PY{n}{Config}\PY{+w}{ }\PY{n+nf}{getInstance}\PY{p}{(}\PY{p}{)}\PY{+w}{ }\PY{p}{\PYZob{}}\PY{+w}{ }\PY{p}{.}\PY{p}{.}\PY{p}{.}\PY{+w}{ }\PY{p}{\PYZcb{}}\PY{+w}{ }\PY{p}{\PYZcb{}}
\end{Verbatim}
\end{codebox}

\item 
\textbf{\code{Optional} used as a field or method parameter.}
Why it matters: \code{Optional} was designed as a return type; using it for fields/parameters adds wrapping overhead and is not serializable.

\begin{codebox}
\begin{Verbatim}[commandchars=\\\{\}]
\PY{k+kd}{class} \PY{n+nc}{Order}\PY{+w}{ }\PY{p}{\PYZob{}}\PY{+w}{ }\PY{n}{Optional}\PY{o}{\PYZlt{}}\PY{n}{String}\PY{o}{\PYZgt{}}\PY{+w}{ }\PY{n}{couponCode}\PY{p}{;}\PY{+w}{ }\PY{p}{\PYZcb{}}\PY{+w}{ }\PY{c+c1}{// flag: should be a plain nullable String}
\end{Verbatim}
\end{codebox}

\item 
\textbf{String concatenation with \code{+} inside a loop.}
Why it matters: quietly O(n²) — each iteration allocates a brand-new \code{String}; use \code{StringBuilder}.

\begin{codebox}
\begin{Verbatim}[commandchars=\\\{\}]
\PY{n}{String}\PY{+w}{ }\PY{n}{s}\PY{+w}{ }\PY{o}{=}\PY{+w}{ }\PY{l+s}{\PYZdq{}}\PY{l+s}{\PYZdq{}}\PY{p}{;}
\PY{k}{for}\PY{+w}{ }\PY{p}{(}\PY{n}{String}\PY{+w}{ }\PY{n}{w}\PY{+w}{ }\PY{p}{:}\PY{+w}{ }\PY{n}{words}\PY{p}{)}\PY{+w}{ }\PY{n}{s}\PY{+w}{ }\PY{o}{+}\PY{o}{=}\PY{+w}{ }\PY{n}{w}\PY{p}{;}\PY{+w}{ }\PY{c+c1}{// flag}
\end{Verbatim}
\end{codebox}

\item 
\textbf{Magic numbers or literals with no named meaning.}
Why it matters: the reader cannot tell if \code{3} is a retry count, a threshold, or an index without digging through history.

\begin{codebox}
\begin{Verbatim}[commandchars=\\\{\}]
\PY{k}{if}\PY{+w}{ }\PY{p}{(}\PY{n}{attempts}\PY{+w}{ }\PY{o}{\PYZgt{}}\PY{o}{=}\PY{+w}{ }\PY{l+m+mi}{3}\PY{p}{)}\PY{+w}{ }\PY{n}{lockAccount}\PY{p}{(}\PY{p}{)}\PY{p}{;}\PY{+w}{ }\PY{c+c1}{// what does 3 mean, and why 3?}
\end{Verbatim}
\end{codebox}

\item 
\textbf{Catching a broad exception and swallowing it, or catching \code{Throwable}.}
Why it matters: hides real bugs, and catching \code{Throwable} also swallows JVM-level \code{Error}s that should never be caught by application code.

\begin{codebox}
\begin{Verbatim}[commandchars=\\\{\}]
\PY{k}{try}\PY{+w}{ }\PY{p}{\PYZob{}}\PY{+w}{ }\PY{n}{risky}\PY{p}{(}\PY{p}{)}\PY{p}{;}\PY{+w}{ }\PY{p}{\PYZcb{}}\PY{+w}{ }\PY{k}{catch}\PY{+w}{ }\PY{p}{(}\PY{n}{Exception}\PY{+w}{ }\PY{n}{e}\PY{p}{)}\PY{+w}{ }\PY{p}{\PYZob{}}\PY{+w}{ }\PY{c+cm}{/* nothing */}\PY{+w}{ }\PY{p}{\PYZcb{}}\PY{+w}{ }\PY{c+c1}{// flag}
\end{Verbatim}
\end{codebox}

\item 
\textbf{New public API with a constructor that has more than four parameters.}
Why it matters: callers can silently swap two same-typed arguments (e.g. two \code{int}s) with no compiler warning; a builder or record removes the ambiguity.

\begin{codebox}
\begin{Verbatim}[commandchars=\\\{\}]
\PY{k}{new}\PY{+w}{ }\PY{n}{User}\PY{p}{(}\PY{l+s}{\PYZdq{}}\PY{l+s}{Ana}\PY{l+s}{\PYZdq{}}\PY{p}{,}\PY{+w}{ }\PY{l+s}{\PYZdq{}}\PY{l+s}{ana@x.com}\PY{l+s}{\PYZdq{}}\PY{p}{,}\PY{+w}{ }\PY{k+kc}{true}\PY{p}{,}\PY{+w}{ }\PY{k+kc}{false}\PY{p}{,}\PY{+w}{ }\PY{l+m+mi}{30}\PY{p}{,}\PY{+w}{ }\PY{l+s}{\PYZdq{}}\PY{l+s}{US}\PY{l+s}{\PYZdq{}}\PY{p}{)}\PY{p}{;}\PY{+w}{ }\PY{c+c1}{// flag: use a builder}
\end{Verbatim}
\end{codebox}

\item 
\textbf{Exposing a mutable internal collection through a public getter.}
Why it matters: any caller can mutate the object’s internal state from outside, breaking encapsulation and invariants.

\begin{codebox}
\begin{Verbatim}[commandchars=\\\{\}]
\PY{k+kd}{public}\PY{+w}{ }\PY{n}{List}\PY{o}{\PYZlt{}}\PY{n}{String}\PY{o}{\PYZgt{}}\PY{+w}{ }\PY{n+nf}{getMembers}\PY{p}{(}\PY{p}{)}\PY{+w}{ }\PY{p}{\PYZob{}}\PY{+w}{ }\PY{k}{return}\PY{+w}{ }\PY{n}{members}\PY{p}{;}\PY{+w}{ }\PY{p}{\PYZcb{}}\PY{+w}{ }\PY{c+c1}{// flag: return List.copyOf(members)}
\end{Verbatim}
\end{codebox}

\item 
\textbf{Removing or renaming a public method without deprecation.}
Why it matters: breaks both source and binary compatibility for every downstream consumer with no warning period.

\begin{codebox}
\begin{Verbatim}[commandchars=\\\{\}]
\PY{c+c1}{// Flag if this method existed in a previous public release and just disappears}
\PY{c+c1}{// public void save(Order order) \PYZob{} ... \PYZcb{}}
\end{Verbatim}
\end{codebox}

\item 
\textbf{Deep nested \code{if} blocks instead of guard clauses.}
Why it matters: the reader must hold multiple nested conditions in their head at once to find the actual logic.

\begin{codebox}
\begin{Verbatim}[commandchars=\\\{\}]
\PY{k}{if}\PY{+w}{ }\PY{p}{(}\PY{n}{a}\PY{+w}{ }\PY{o}{!}\PY{o}{=}\PY{+w}{ }\PY{k+kc}{null}\PY{p}{)}\PY{+w}{ }\PY{p}{\PYZob{}}\PY{+w}{ }\PY{k}{if}\PY{+w}{ }\PY{p}{(}\PY{n}{a}\PY{p}{.}\PY{n+na}{isValid}\PY{p}{(}\PY{p}{)}\PY{p}{)}\PY{+w}{ }\PY{p}{\PYZob{}}\PY{+w}{ }\PY{k}{if}\PY{+w}{ }\PY{p}{(}\PY{n}{b}\PY{+w}{ }\PY{o}{\PYZgt{}}\PY{+w}{ }\PY{l+m+mi}{0}\PY{p}{)}\PY{+w}{ }\PY{p}{\PYZob{}}\PY{+w}{ }\PY{n}{doWork}\PY{p}{(}\PY{p}{)}\PY{p}{;}\PY{+w}{ }\PY{p}{\PYZcb{}}\PY{+w}{ }\PY{p}{\PYZcb{}}\PY{+w}{ }\PY{p}{\PYZcb{}}\PY{+w}{ }\PY{c+c1}{// flag}
\end{Verbatim}
\end{codebox}

\item 
\textbf{A “Manager”/“Helper”/“Utils” class that keeps growing with unrelated static methods.}
Why it matters: it becomes a God object with no cohesion, and every unrelated change collides in the same file.

\begin{codebox}
\begin{Verbatim}[commandchars=\\\{\}]
\PY{k+kd}{class} \PY{n+nc}{Utils}\PY{+w}{ }\PY{p}{\PYZob{}}\PY{+w}{ }\PY{k+kd}{static}\PY{+w}{ }\PY{n}{String}\PY{+w}{ }\PY{n+nf}{formatDate}\PY{p}{(}\PY{p}{.}\PY{p}{.}\PY{p}{.}\PY{p}{)}\PY{+w}{ }\PY{p}{\PYZob{}}\PY{p}{\PYZcb{}}\PY{+w}{ }\PY{k+kd}{static}\PY{+w}{ }\PY{k+kt}{void}\PY{+w}{ }\PY{n+nf}{chargeCard}\PY{p}{(}\PY{p}{.}\PY{p}{.}\PY{p}{.}\PY{p}{)}\PY{+w}{ }\PY{p}{\PYZob{}}\PY{p}{\PYZcb{}}\PY{+w}{ }\PY{p}{\PYZcb{}}\PY{+w}{ }\PY{c+c1}{// flag}
\end{Verbatim}
\end{codebox}

\item 
\textbf{Design pattern boilerplate (Strategy/Command/Observer as full classes) where a lambda would do.}
Why it matters: in modern Java, a one-line lambda implementing a functional interface is clearer than a whole extra class file for simple cases — flag unnecessary ceremony.

\begin{codebox}
\begin{Verbatim}[commandchars=\\\{\}]
\PY{k+kd}{class} \PY{n+nc}{AlwaysApprove}\PY{+w}{ }\PY{k+kd}{implements}\PY{+w}{ }\PY{n}{DiscountStrategy}\PY{+w}{ }\PY{p}{\PYZob{}}\PY{+w}{ }\PY{k+kd}{public}\PY{+w}{ }\PY{n}{BigDecimal}\PY{+w}{ }\PY{n+nf}{apply}\PY{p}{(}\PY{n}{BigDecimal}\PY{+w}{ }\PY{n}{p}\PY{p}{)}\PY{+w}{ }\PY{p}{\PYZob{}}\PY{+w}{ }\PY{k}{return}\PY{+w}{ }\PY{n}{p}\PY{p}{;}\PY{+w}{ }\PY{p}{\PYZcb{}}\PY{+w}{ }\PY{p}{\PYZcb{}}
\PY{c+c1}{// Prefer: DiscountStrategy noDiscount = price \PYZhy{}\PYZgt{} price;}
\end{Verbatim}
\end{codebox}

\item 
\textbf{Using a raw \code{String}/\code{int} where a fixed set of values exists (stringly-typed code).}
Why it matters: no compile-time safety — typos like \code{\textquotedbl{}Pending\textquotedbl{}} vs \code{\textquotedbl{}PENDING\textquotedbl{}} compile fine and fail silently at runtime.

\begin{codebox}
\begin{Verbatim}[commandchars=\\\{\}]
\PY{k+kt}{void}\PY{+w}{ }\PY{n+nf}{setStatus}\PY{p}{(}\PY{n}{String}\PY{+w}{ }\PY{n}{status}\PY{p}{)}\PY{+w}{ }\PY{p}{\PYZob{}}\PY{+w}{ }\PY{p}{.}\PY{p}{.}\PY{p}{.}\PY{+w}{ }\PY{p}{\PYZcb{}}\PY{+w}{ }\PY{c+c1}{// flag: should be an enum OrderStatus}
\end{Verbatim}
\end{codebox}

\item 
\textbf{Test mocks every collaborator and asserts call order instead of testing behavior.}
Why it matters: the test breaks whenever internal wiring changes, even if the externally observable behavior is unchanged — a maintenance burden with no correctness benefit.

\begin{codebox}
\begin{Verbatim}[commandchars=\\\{\}]
\PY{n}{verify}\PY{p}{(}\PY{n}{mockA}\PY{p}{)}\PY{p}{.}\PY{n+na}{method}\PY{p}{(}\PY{p}{)}\PY{p}{;}\PY{+w}{ }\PY{n}{verify}\PY{p}{(}\PY{n}{mockB}\PY{p}{)}\PY{p}{.}\PY{n+na}{method}\PY{p}{(}\PY{p}{)}\PY{p}{;}\PY{+w}{ }\PY{c+c1}{// flag if this is the entire assertion}
\end{Verbatim}
\end{codebox}

\end{enumerate}


\setcounter{chapter}{19}
\chapter{Testing}
\label{chap:20}
\label{h:20:20-testing}
Code review is not just about the production code in the diff — it is at least as much about the tests that accompany it. A reviewer who cannot tell \emph{what} a test proves, or who spots a test that mocks its way around the actual bug, will (and should) block the PR. This chapter covers JUnit 5 in depth — the framework almost every Java shop uses today — and Mockito, the de-facto mocking library, with an emphasis on the patterns and anti-patterns that come up constantly in review. We target JUnit Jupiter 5.10+ and Mockito 5.x on Java 21+.

\section[Testing Fundamentals (JUnit 5)]{Testing Fundamentals (JUnit 5)}
\label{h:20:testing-fundamentals-junit-5}
\subsection[Why We Test, and the Test Pyramid]{Why We Test, and the Test Pyramid}
\label{h:20:why-we-test-and-the-test-pyramid}
Tests exist to answer one question cheaply and repeatedly: \textbf{does the system still behave correctly after this change?} Without automated tests, that question can only be answered by manual verification (slow, inconsistent, easily skipped) or by production incidents (expensive, embarrassing). A good test suite turns “did I break anything?” into a command that runs in seconds and reports the answer precisely.

The \textbf{test pyramid} is a mental model for how to allocate testing effort:

\begin{codebox}
\begin{Verbatim}
        ▲
       / \        End-to-End (E2E) tests
      /   \        - few, slow, brittle, high confidence
     /-----\
    /       \      Integration tests
   /         \      - moderate count, hit real DB/HTTP/queue boundaries
  /-----------\
 /             \    Unit tests
/---------------\    - many, fast, isolated, cheap to write and run
\end{Verbatim}
\end{codebox}

\begin{itemize}
\item 
\textbf{Unit tests} exercise a single class or a small cluster of classes in isolation (often with collaborators mocked or faked). They should be the vast majority of your suite: fast (milliseconds), deterministic, and pinpoint the exact failure.

\item 
\textbf{Integration tests} verify that your code talks correctly to a real (or realistic) database, message broker, HTTP client, etc. Fewer of these — they are slower and more fragile — but essential because unit tests alone cannot catch a wrong SQL query or a misconfigured serializer.

\item 
\textbf{End-to-end tests} drive the whole system as a black box (e.g., through its public API or UI). Very few of these: they are slow, flaky, and expensive to maintain, but they catch wiring problems nothing else can.

\end{itemize}

A reviewer’s rule of thumb: if a PR adds a new branch of logic and the diff has no matching unit test, that is a legitimate request-changes comment. If a PR tries to verify business logic exclusively through a slow end-to-end test, that is also worth challenging — push the assertion down the pyramid.

\subsection[AAA / Given-When-Then Structure]{AAA / Given-When-Then Structure}
\label{h:20:aaa--given-when-then-structure}
Every well-written test should have three clearly separated parts, whether or not you label them:

\begin{itemize}
\item 
\textbf{Arrange / Given} — set up the object under test and its inputs.

\item 
\textbf{Act / When} — perform the one action being tested.

\item 
\textbf{Assert / Then} — check the outcome.

\end{itemize}

\begin{codebox}
\begin{Verbatim}[commandchars=\\\{\}]
\PY{n+nd}{@Test}
\PY{k+kt}{void}\PY{+w}{ }\PY{n+nf}{discount\PYZus{}appliesTenPercentForPremiumCustomer}\PY{p}{(}\PY{p}{)}\PY{+w}{ }\PY{p}{\PYZob{}}
\PY{+w}{    }\PY{c+c1}{// Arrange (Given)}
\PY{+w}{    }\PY{n}{Customer}\PY{+w}{ }\PY{n}{premiumCustomer}\PY{+w}{ }\PY{o}{=}\PY{+w}{ }\PY{k}{new}\PY{+w}{ }\PY{n}{Customer}\PY{p}{(}\PY{l+s}{\PYZdq{}}\PY{l+s}{alice}\PY{l+s}{\PYZdq{}}\PY{p}{,}\PY{+w}{ }\PY{n}{CustomerTier}\PY{p}{.}\PY{n+na}{PREMIUM}\PY{p}{)}\PY{p}{;}
\PY{+w}{    }\PY{n}{PricingService}\PY{+w}{ }\PY{n}{pricingService}\PY{+w}{ }\PY{o}{=}\PY{+w}{ }\PY{k}{new}\PY{+w}{ }\PY{n}{PricingService}\PY{p}{(}\PY{p}{)}\PY{p}{;}

\PY{+w}{    }\PY{c+c1}{// Act (When)}
\PY{+w}{    }\PY{n}{BigDecimal}\PY{+w}{ }\PY{n}{finalPrice}\PY{+w}{ }\PY{o}{=}\PY{+w}{ }\PY{n}{pricingService}\PY{p}{.}\PY{n+na}{priceFor}\PY{p}{(}\PY{n}{premiumCustomer}\PY{p}{,}\PY{+w}{ }\PY{k}{new}\PY{+w}{ }\PY{n}{BigDecimal}\PY{p}{(}\PY{l+s}{\PYZdq{}}\PY{l+s}{100.00}\PY{l+s}{\PYZdq{}}\PY{p}{)}\PY{p}{)}\PY{p}{;}

\PY{+w}{    }\PY{c+c1}{// Assert (Then)}
\PY{+w}{    }\PY{n}{assertEquals}\PY{p}{(}\PY{k}{new}\PY{+w}{ }\PY{n}{BigDecimal}\PY{p}{(}\PY{l+s}{\PYZdq{}}\PY{l+s}{90.00}\PY{l+s}{\PYZdq{}}\PY{p}{)}\PY{p}{,}\PY{+w}{ }\PY{n}{finalPrice}\PY{p}{)}\PY{p}{;}
\PY{p}{\PYZcb{}}
\end{Verbatim}
\end{codebox}

Keep each test to \textbf{one logical action} and, ideally, \textbf{one reason to fail}. A test with five unrelated \code{assertEquals} calls checking five unrelated behaviors is a maintenance headache: when it fails, you have to read the whole body to figure out which behavior broke. Splitting it into five small tests (or using \code{assertAll}, covered below, when the assertions really are about a single outcome) gives you a failure message that already tells you what broke.

\subsection[Naming Tests So Failures Read Like Sentences]{Naming Tests So Failures Read Like Sentences}
\label{h:20:naming-tests-so-failures-read-like-sentences}
A test name is documentation that never goes stale — if it does, the test breaks and tells you. Favor names that describe \textbf{behavior}, not implementation, and that read naturally when a report lists a failing test.

Common conventions:

{\tablefont
\begin{xltabular}{\linewidth}{@{}L{0.869}L{1.131}@{}}
\toprule
\rowcolor{tablehdr}
\thd{Style} & \thd{Example} \\
\midrule
\endhead
\code{method\allowbreak{}Name\_\allowbreak{}condition\_\allowbreak{}expected\allowbreak{}Result} & \code{withdraw\_\allowbreak{}insufficient\allowbreak{}Funds\_\allowbreak{}throws\allowbreak{}Illegal\allowbreak{}State\allowbreak{}Exception} \\
Given-When-Then in the name & \code{given\allowbreak{}Empty\allowbreak{}Cart\_\allowbreak{}when\allowbreak{}Checkout\_\allowbreak{}then\allowbreak{}Throws\allowbreak{}Empty\allowbreak{}Cart\allowbreak{}Exception} \\
Sentence-like with backticks (\code{@\allowbreak{}DisplayName}) & \code{\textquotedbl{}withdraw()\allowbreak{}~\allowbreak{}throws~\allowbreak{}when~\allowbreak{}balance~\allowbreak{}is~\allowbreak{}insufficient\textquotedbl{}} \\
\bottomrule
\end{xltabular}
}

\begin{codebox}
\begin{Verbatim}[commandchars=\\\{\}]
\PY{c+c1}{// Weak: tells you nothing when it\PYZsq{}s in a failure list}
\PY{n+nd}{@Test}
\PY{k+kt}{void}\PY{+w}{ }\PY{n+nf}{test1}\PY{p}{(}\PY{p}{)}\PY{+w}{ }\PY{p}{\PYZob{}}\PY{+w}{ }\PY{p}{.}\PY{p}{.}\PY{p}{.}\PY{+w}{ }\PY{p}{\PYZcb{}}

\PY{c+c1}{// Better: reads like a specification}
\PY{n+nd}{@Test}
\PY{k+kt}{void}\PY{+w}{ }\PY{n+nf}{withdraw\PYZus{}insufficientFunds\PYZus{}throwsIllegalStateException}\PY{p}{(}\PY{p}{)}\PY{+w}{ }\PY{p}{\PYZob{}}\PY{+w}{ }\PY{p}{.}\PY{p}{.}\PY{p}{.}\PY{+w}{ }\PY{p}{\PYZcb{}}

\PY{c+c1}{// Best when paired with @DisplayName for human\PYZhy{}readable reports}
\PY{n+nd}{@Test}
\PY{n+nd}{@DisplayName}\PY{p}{(}\PY{l+s}{\PYZdq{}}\PY{l+s}{withdraw() throws IllegalStateException when balance is insufficient}\PY{l+s}{\PYZdq{}}\PY{p}{)}
\PY{k+kt}{void}\PY{+w}{ }\PY{n+nf}{withdraw\PYZus{}insufficientFunds\PYZus{}throwsIllegalStateException}\PY{p}{(}\PY{p}{)}\PY{+w}{ }\PY{p}{\PYZob{}}\PY{+w}{ }\PY{p}{.}\PY{p}{.}\PY{p}{.}\PY{+w}{ }\PY{p}{\PYZcb{}}
\end{Verbatim}
\end{codebox}

Avoid names like \code{testWithdraw} or \code{shouldWork} — they force the reader to open the method body to learn anything, defeating the purpose of a descriptive test suite.

\subsection[JUnit 5 Architecture: Platform, Jupiter, Vintage]{JUnit 5 Architecture: Platform, Jupiter, Vintage}
\label{h:20:junit-5-architecture-platform-jupiter-vintage}
JUnit 5 is not one library but three:

\begin{itemize}
\item 
\textbf{JUnit Platform} — the foundation. It discovers and runs tests via the \code{TestEngine} API, and is what build tools (Maven Surefire, Gradle) and IDEs actually talk to.

\item 
\textbf{JUnit Jupiter} — the new programming model and extension API (\code{@\allowbreak{}Test}, \code{@\allowbreak{}BeforeEach}, \code{assertEquals}, extensions, etc.). This is “JUnit 5” as most people mean it.

\item 
\textbf{JUnit Vintage} — a \code{TestEngine} that runs old JUnit 3/4 tests on the JUnit Platform, so you can migrate incrementally instead of rewriting everything at once.

\end{itemize}

\begin{codebox}
\begin{Verbatim}
JUnit Platform  (test discovery + execution, IDE/build tool integration)
   ├── JUnit Jupiter Engine   → runs @Test-annotated JUnit 5 tests
   └── JUnit Vintage Engine   → runs legacy JUnit 3/4 tests
\end{Verbatim}
\end{codebox}

This separation is why you can mix legacy \code{@\allowbreak{}org.\allowbreak{}junit.\allowbreak{}Test} (JUnit 4) tests and modern \code{@\allowbreak{}org.\allowbreak{}junit.\allowbreak{}jupiter.\allowbreak{}api.\allowbreak{}Test} (JUnit 5) tests in the same module during a migration — both engines plug into the same platform.

\subsubsection[JUnit 4 → JUnit 5 Annotation Mapping]{JUnit 4 → JUnit 5 Annotation Mapping}
\label{h:20:junit-4--junit-5-annotation-mapping}
{\tablefont
\begin{xltabular}{\linewidth}{@{}L{0.619}L{0.773}L{1.608}@{}}
\toprule
\rowcolor{tablehdr}
\thd{JUnit 4} & \thd{JUnit 5 (Jupiter)} & \thd{Notes} \\
\midrule
\endhead
\code{@\allowbreak{}Test} & \code{@\allowbreak{}Test} & JUnit 5’s \code{@\allowbreak{}Test} no longer supports \code{expected=}/\code{timeout=} attributes \\
\code{@\allowbreak{}Before} & \code{@\allowbreak{}BeforeEach} &  \\
\code{@\allowbreak{}After} & \code{@\allowbreak{}AfterEach} &  \\
\code{@\allowbreak{}BeforeClass} & \code{@\allowbreak{}BeforeAll} & Must be \code{static} unless \code{@\allowbreak{}Test\allowbreak{}Instance(\allowbreak{}PER\_\allowbreak{}CLASS)} \\
\code{@\allowbreak{}AfterClass} & \code{@\allowbreak{}AfterAll} & Must be \code{static} unless \code{@\allowbreak{}Test\allowbreak{}Instance(\allowbreak{}PER\_\allowbreak{}CLASS)} \\
\code{@\allowbreak{}Ignore} & \code{@\allowbreak{}Disabled} &  \\
\code{@\allowbreak{}Category(\allowbreak{}Foo.\allowbreak{}class)} & \code{@\allowbreak{}Tag(\allowbreak{}\textquotedbl{}foo\textquotedbl{})} & String-based, no marker interface needed \\
\code{@\allowbreak{}RunWith(...)} & \code{@\allowbreak{}ExtendWith(...)} & Composable — you can stack multiple extensions \\
\code{@\allowbreak{}Rule} / \code{@\allowbreak{}ClassRule} & \code{Extension} implementations & More granular lifecycle hooks (see Extensions below) \\
\code{@\allowbreak{}Test(\allowbreak{}expected~\allowbreak{}=\allowbreak{}~\allowbreak{}X.\allowbreak{}class)} & \code{assert\allowbreak{}Throws(\allowbreak{}X.\allowbreak{}class,\allowbreak{}~\allowbreak{}()\allowbreak{}~\allowbreak{}-\textgreater{}\allowbreak{}~\allowbreak{}...)} & Returns the exception so you can assert on its message too \\
\code{Assume.\allowbreak{}assumeTrue(...)} & \code{Assumptions.\allowbreak{}assume\allowbreak{}True(...)} & Same idea, new package \\
\bottomrule
\end{xltabular}
}

\subsection[Maven and Gradle Setup]{Maven and Gradle Setup}
\label{h:20:maven-and-gradle-setup}
\textbf{Maven} (\code{pom.\allowbreak{}xml}) — the BOM keeps all Jupiter artifacts on matching versions:

\begin{codebox}
\begin{Verbatim}[commandchars=\\\{\}]
\PY{n+nt}{\PYZlt{}dependencyManagement}\PY{n+nt}{\PYZgt{}}
\PY{+w}{    }\PY{n+nt}{\PYZlt{}dependencies}\PY{n+nt}{\PYZgt{}}
\PY{+w}{        }\PY{n+nt}{\PYZlt{}dependency}\PY{n+nt}{\PYZgt{}}
\PY{+w}{            }\PY{n+nt}{\PYZlt{}groupId}\PY{n+nt}{\PYZgt{}}org.junit\PY{n+nt}{\PYZlt{}/groupId\PYZgt{}}
\PY{+w}{            }\PY{n+nt}{\PYZlt{}artifactId}\PY{n+nt}{\PYZgt{}}junit\PYZhy{}bom\PY{n+nt}{\PYZlt{}/artifactId\PYZgt{}}
\PY{+w}{            }\PY{n+nt}{\PYZlt{}version}\PY{n+nt}{\PYZgt{}}5.10.2\PY{n+nt}{\PYZlt{}/version\PYZgt{}}
\PY{+w}{            }\PY{n+nt}{\PYZlt{}type}\PY{n+nt}{\PYZgt{}}pom\PY{n+nt}{\PYZlt{}/type\PYZgt{}}
\PY{+w}{            }\PY{n+nt}{\PYZlt{}scope}\PY{n+nt}{\PYZgt{}}import\PY{n+nt}{\PYZlt{}/scope\PYZgt{}}
\PY{+w}{        }\PY{n+nt}{\PYZlt{}/dependency\PYZgt{}}
\PY{+w}{    }\PY{n+nt}{\PYZlt{}/dependencies\PYZgt{}}
\PY{n+nt}{\PYZlt{}/dependencyManagement\PYZgt{}}

\PY{n+nt}{\PYZlt{}dependencies}\PY{n+nt}{\PYZgt{}}
\PY{+w}{    }\PY{n+nt}{\PYZlt{}dependency}\PY{n+nt}{\PYZgt{}}
\PY{+w}{        }\PY{n+nt}{\PYZlt{}groupId}\PY{n+nt}{\PYZgt{}}org.junit.jupiter\PY{n+nt}{\PYZlt{}/groupId\PYZgt{}}
\PY{+w}{        }\PY{n+nt}{\PYZlt{}artifactId}\PY{n+nt}{\PYZgt{}}junit\PYZhy{}jupiter\PY{n+nt}{\PYZlt{}/artifactId\PYZgt{}}
\PY{+w}{        }\PY{n+nt}{\PYZlt{}scope}\PY{n+nt}{\PYZgt{}}test\PY{n+nt}{\PYZlt{}/scope\PYZgt{}}
\PY{+w}{    }\PY{n+nt}{\PYZlt{}/dependency\PYZgt{}}
\PY{+w}{    }\PY{n+nt}{\PYZlt{}dependency}\PY{n+nt}{\PYZgt{}}
\PY{+w}{        }\PY{n+nt}{\PYZlt{}groupId}\PY{n+nt}{\PYZgt{}}org.mockito\PY{n+nt}{\PYZlt{}/groupId\PYZgt{}}
\PY{+w}{        }\PY{n+nt}{\PYZlt{}artifactId}\PY{n+nt}{\PYZgt{}}mockito\PYZhy{}junit\PYZhy{}jupiter\PY{n+nt}{\PYZlt{}/artifactId\PYZgt{}}
\PY{+w}{        }\PY{n+nt}{\PYZlt{}version}\PY{n+nt}{\PYZgt{}}5.11.0\PY{n+nt}{\PYZlt{}/version\PYZgt{}}
\PY{+w}{        }\PY{n+nt}{\PYZlt{}scope}\PY{n+nt}{\PYZgt{}}test\PY{n+nt}{\PYZlt{}/scope\PYZgt{}}
\PY{+w}{    }\PY{n+nt}{\PYZlt{}/dependency\PYZgt{}}
\PY{+w}{    }\PY{n+nt}{\PYZlt{}dependency}\PY{n+nt}{\PYZgt{}}
\PY{+w}{        }\PY{n+nt}{\PYZlt{}groupId}\PY{n+nt}{\PYZgt{}}org.assertj\PY{n+nt}{\PYZlt{}/groupId\PYZgt{}}
\PY{+w}{        }\PY{n+nt}{\PYZlt{}artifactId}\PY{n+nt}{\PYZgt{}}assertj\PYZhy{}core\PY{n+nt}{\PYZlt{}/artifactId\PYZgt{}}
\PY{+w}{        }\PY{n+nt}{\PYZlt{}version}\PY{n+nt}{\PYZgt{}}3.25.3\PY{n+nt}{\PYZlt{}/version\PYZgt{}}
\PY{+w}{        }\PY{n+nt}{\PYZlt{}scope}\PY{n+nt}{\PYZgt{}}test\PY{n+nt}{\PYZlt{}/scope\PYZgt{}}
\PY{+w}{    }\PY{n+nt}{\PYZlt{}/dependency\PYZgt{}}
\PY{n+nt}{\PYZlt{}/dependencies\PYZgt{}}

\PY{n+nt}{\PYZlt{}build}\PY{n+nt}{\PYZgt{}}
\PY{+w}{    }\PY{n+nt}{\PYZlt{}plugins}\PY{n+nt}{\PYZgt{}}
\PY{+w}{        }\PY{n+nt}{\PYZlt{}plugin}\PY{n+nt}{\PYZgt{}}
\PY{+w}{            }\PY{n+nt}{\PYZlt{}groupId}\PY{n+nt}{\PYZgt{}}org.apache.maven.plugins\PY{n+nt}{\PYZlt{}/groupId\PYZgt{}}
\PY{+w}{            }\PY{n+nt}{\PYZlt{}artifactId}\PY{n+nt}{\PYZgt{}}maven\PYZhy{}surefire\PYZhy{}plugin\PY{n+nt}{\PYZlt{}/artifactId\PYZgt{}}
\PY{+w}{            }\PY{n+nt}{\PYZlt{}version}\PY{n+nt}{\PYZgt{}}3.2.5\PY{n+nt}{\PYZlt{}/version\PYZgt{}}
\PY{+w}{        }\PY{n+nt}{\PYZlt{}/plugin\PYZgt{}}
\PY{+w}{    }\PY{n+nt}{\PYZlt{}/plugins\PYZgt{}}
\PY{n+nt}{\PYZlt{}/build\PYZgt{}}
\end{Verbatim}
\end{codebox}

\textbf{Gradle} (\code{build.\allowbreak{}gradle.\allowbreak{}kts}):

\begin{codebox}
\begin{Verbatim}[commandchars=\\\{\}]
\PY{n}{dependencies}\PY{+w}{ }\PY{p}{\PYZob{}}
\PY{+w}{    }\PY{n}{testImplementation}\PY{p}{(}\PY{n}{platform}\PY{p}{(}\PY{l+s}{\PYZdq{}}\PY{l+s}{org.junit:junit\PYZhy{}bom:5.10.2}\PY{l+s}{\PYZdq{}}\PY{p}{)}\PY{p}{)}
\PY{+w}{    }\PY{n}{testImplementation}\PY{p}{(}\PY{l+s}{\PYZdq{}}\PY{l+s}{org.junit.jupiter:junit\PYZhy{}jupiter}\PY{l+s}{\PYZdq{}}\PY{p}{)}
\PY{+w}{    }\PY{n}{testImplementation}\PY{p}{(}\PY{l+s}{\PYZdq{}}\PY{l+s}{org.mockito:mockito\PYZhy{}junit\PYZhy{}jupiter:5.11.0}\PY{l+s}{\PYZdq{}}\PY{p}{)}
\PY{+w}{    }\PY{n}{testImplementation}\PY{p}{(}\PY{l+s}{\PYZdq{}}\PY{l+s}{org.assertj:assertj\PYZhy{}core:3.25.3}\PY{l+s}{\PYZdq{}}\PY{p}{)}
\PY{p}{\PYZcb{}}

\PY{n}{tasks}\PY{p}{.}\PY{n+na}{test}\PY{+w}{ }\PY{p}{\PYZob{}}
\PY{+w}{    }\PY{n}{useJUnitPlatform}\PY{p}{(}\PY{p}{)}
\PY{p}{\PYZcb{}}
\end{Verbatim}
\end{codebox}

Forgetting \code{useJUnitPlatform()} in Gradle (or the Surefire dependency alignment in Maven) is a classic “why are my tests not running at all” bug — the build reports success with zero tests executed, which is worse than a failure because it hides silently.

\subsection[Core Annotations and Lifecycle]{Core Annotations and Lifecycle}
\label{h:20:core-annotations-and-lifecycle}
\begin{codebox}
\begin{Verbatim}[commandchars=\\\{\}]
\PY{k+kn}{import}\PY{+w}{ }\PY{n+nn}{org.junit.jupiter.api.*}\PY{p}{;}
\PY{k+kn}{import static}\PY{+w}{ }\PY{n+nn}{org.junit.jupiter.api.Assertions.*}\PY{p}{;}

\PY{k+kd}{class} \PY{n+nc}{OrderCalculatorTest}\PY{+w}{ }\PY{p}{\PYZob{}}

\PY{+w}{    }\PY{k+kd}{private}\PY{+w}{ }\PY{n}{OrderCalculator}\PY{+w}{ }\PY{n}{calculator}\PY{p}{;}

\PY{+w}{    }\PY{n+nd}{@BeforeAll}
\PY{+w}{    }\PY{k+kd}{static}\PY{+w}{ }\PY{k+kt}{void}\PY{+w}{ }\PY{n+nf}{setUpClass}\PY{p}{(}\PY{p}{)}\PY{+w}{ }\PY{p}{\PYZob{}}
\PY{+w}{        }\PY{n}{System}\PY{p}{.}\PY{n+na}{out}\PY{p}{.}\PY{n+na}{println}\PY{p}{(}\PY{l+s}{\PYZdq{}}\PY{l+s}{Runs once before any test in this class}\PY{l+s}{\PYZdq{}}\PY{p}{)}\PY{p}{;}
\PY{+w}{    }\PY{p}{\PYZcb{}}

\PY{+w}{    }\PY{n+nd}{@BeforeEach}
\PY{+w}{    }\PY{k+kt}{void}\PY{+w}{ }\PY{n+nf}{setUp}\PY{p}{(}\PY{p}{)}\PY{+w}{ }\PY{p}{\PYZob{}}
\PY{+w}{        }\PY{n}{calculator}\PY{+w}{ }\PY{o}{=}\PY{+w}{ }\PY{k}{new}\PY{+w}{ }\PY{n}{OrderCalculator}\PY{p}{(}\PY{p}{)}\PY{p}{;}
\PY{+w}{        }\PY{n}{System}\PY{p}{.}\PY{n+na}{out}\PY{p}{.}\PY{n+na}{println}\PY{p}{(}\PY{l+s}{\PYZdq{}}\PY{l+s}{Runs before every single test method}\PY{l+s}{\PYZdq{}}\PY{p}{)}\PY{p}{;}
\PY{+w}{    }\PY{p}{\PYZcb{}}

\PY{+w}{    }\PY{n+nd}{@Test}
\PY{+w}{    }\PY{n+nd}{@DisplayName}\PY{p}{(}\PY{l+s}{\PYZdq{}}\PY{l+s}{total() sums line item prices}\PY{l+s}{\PYZdq{}}\PY{p}{)}
\PY{+w}{    }\PY{k+kt}{void}\PY{+w}{ }\PY{n+nf}{total\PYZus{}sumsLineItemPrices}\PY{p}{(}\PY{p}{)}\PY{+w}{ }\PY{p}{\PYZob{}}
\PY{+w}{        }\PY{n}{calculator}\PY{p}{.}\PY{n+na}{addLine}\PY{p}{(}\PY{l+s}{\PYZdq{}}\PY{l+s}{widget}\PY{l+s}{\PYZdq{}}\PY{p}{,}\PY{+w}{ }\PY{l+m+mi}{2}\PY{p}{,}\PY{+w}{ }\PY{k}{new}\PY{+w}{ }\PY{n}{BigDecimal}\PY{p}{(}\PY{l+s}{\PYZdq{}}\PY{l+s}{5.00}\PY{l+s}{\PYZdq{}}\PY{p}{)}\PY{p}{)}\PY{p}{;}
\PY{+w}{        }\PY{n}{calculator}\PY{p}{.}\PY{n+na}{addLine}\PY{p}{(}\PY{l+s}{\PYZdq{}}\PY{l+s}{gadget}\PY{l+s}{\PYZdq{}}\PY{p}{,}\PY{+w}{ }\PY{l+m+mi}{1}\PY{p}{,}\PY{+w}{ }\PY{k}{new}\PY{+w}{ }\PY{n}{BigDecimal}\PY{p}{(}\PY{l+s}{\PYZdq{}}\PY{l+s}{12.50}\PY{l+s}{\PYZdq{}}\PY{p}{)}\PY{p}{)}\PY{p}{;}

\PY{+w}{        }\PY{n}{assertEquals}\PY{p}{(}\PY{k}{new}\PY{+w}{ }\PY{n}{BigDecimal}\PY{p}{(}\PY{l+s}{\PYZdq{}}\PY{l+s}{22.50}\PY{l+s}{\PYZdq{}}\PY{p}{)}\PY{p}{,}\PY{+w}{ }\PY{n}{calculator}\PY{p}{.}\PY{n+na}{total}\PY{p}{(}\PY{p}{)}\PY{p}{)}\PY{p}{;}
\PY{+w}{    }\PY{p}{\PYZcb{}}

\PY{+w}{    }\PY{n+nd}{@AfterEach}
\PY{+w}{    }\PY{k+kt}{void}\PY{+w}{ }\PY{n+nf}{tearDown}\PY{p}{(}\PY{p}{)}\PY{+w}{ }\PY{p}{\PYZob{}}
\PY{+w}{        }\PY{n}{System}\PY{p}{.}\PY{n+na}{out}\PY{p}{.}\PY{n+na}{println}\PY{p}{(}\PY{l+s}{\PYZdq{}}\PY{l+s}{Runs after every single test method}\PY{l+s}{\PYZdq{}}\PY{p}{)}\PY{p}{;}
\PY{+w}{    }\PY{p}{\PYZcb{}}

\PY{+w}{    }\PY{n+nd}{@AfterAll}
\PY{+w}{    }\PY{k+kd}{static}\PY{+w}{ }\PY{k+kt}{void}\PY{+w}{ }\PY{n+nf}{tearDownClass}\PY{p}{(}\PY{p}{)}\PY{+w}{ }\PY{p}{\PYZob{}}
\PY{+w}{        }\PY{n}{System}\PY{p}{.}\PY{n+na}{out}\PY{p}{.}\PY{n+na}{println}\PY{p}{(}\PY{l+s}{\PYZdq{}}\PY{l+s}{Runs once after all tests in this class}\PY{l+s}{\PYZdq{}}\PY{p}{)}\PY{p}{;}
\PY{+w}{    }\PY{p}{\PYZcb{}}
\PY{p}{\PYZcb{}}
\end{Verbatim}
\end{codebox}

By default, JUnit creates a \textbf{new instance of the test class for every test method} (\code{@\allowbreak{}Test\allowbreak{}Instance(\allowbreak{}Lifecycle.\allowbreak{}PER\_\allowbreak{}METHOD)}), which is why \code{@\allowbreak{}BeforeAll}/\code{@\allowbreak{}AfterAll} must be \code{static} — there is no single instance to hold state across tests. This also means each test starts from a completely fresh object graph, preventing one test’s mutated state from leaking into another — a common source of order-dependent flakiness in older frameworks.

You can opt into one shared instance per test class:

\begin{codebox}
\begin{Verbatim}[commandchars=\\\{\}]
\PY{n+nd}{@TestInstance}\PY{p}{(}\PY{n}{TestInstance}\PY{p}{.}\PY{n+na}{Lifecycle}\PY{p}{.}\PY{n+na}{PER\PYZus{}CLASS}\PY{p}{)}
\PY{k+kd}{class} \PY{n+nc}{ExpensiveSetupTest}\PY{+w}{ }\PY{p}{\PYZob{}}

\PY{+w}{    }\PY{k+kd}{private}\PY{+w}{ }\PY{n}{ExpensiveResource}\PY{+w}{ }\PY{n}{resource}\PY{p}{;}

\PY{+w}{    }\PY{n+nd}{@BeforeAll}
\PY{+w}{    }\PY{k+kt}{void}\PY{+w}{ }\PY{n+nf}{setUp}\PY{p}{(}\PY{p}{)}\PY{+w}{ }\PY{p}{\PYZob{}}\PY{+w}{           }\PY{c+c1}{// no longer needs to be static}
\PY{+w}{        }\PY{n}{resource}\PY{+w}{ }\PY{o}{=}\PY{+w}{ }\PY{k}{new}\PY{+w}{ }\PY{n}{ExpensiveResource}\PY{p}{(}\PY{p}{)}\PY{p}{;}
\PY{+w}{    }\PY{p}{\PYZcb{}}

\PY{+w}{    }\PY{n+nd}{@Test}
\PY{+w}{    }\PY{k+kt}{void}\PY{+w}{ }\PY{n+nf}{usesSharedResource}\PY{p}{(}\PY{p}{)}\PY{+w}{ }\PY{p}{\PYZob{}}
\PY{+w}{        }\PY{n}{assertNotNull}\PY{p}{(}\PY{n}{resource}\PY{p}{)}\PY{p}{;}
\PY{+w}{    }\PY{p}{\PYZcb{}}
\PY{p}{\PYZcb{}}
\end{Verbatim}
\end{codebox}

Use \code{PER\_\allowbreak{}CLASS} sparingly — it is convenient for genuinely expensive, read-only setup, but it reintroduces the risk of tests sharing mutable state, the exact problem \code{PER\_\allowbreak{}METHOD} was designed to avoid. If you use it, make sure the shared object is either immutable or reset in \code{@\allowbreak{}BeforeEach}.

\subsection[Assertions]{Assertions}
\label{h:20:assertions}
Jupiter’s built-in \code{Assertions} class covers the common cases:

\begin{codebox}
\begin{Verbatim}[commandchars=\\\{\}]
\PY{k+kn}{import static}\PY{+w}{ }\PY{n+nn}{org.junit.jupiter.api.Assertions.*}\PY{p}{;}

\PY{n+nd}{@Test}
\PY{k+kt}{void}\PY{+w}{ }\PY{n+nf}{assertionShowcase}\PY{p}{(}\PY{p}{)}\PY{+w}{ }\PY{p}{\PYZob{}}
\PY{+w}{    }\PY{n}{assertEquals}\PY{p}{(}\PY{l+m+mi}{4}\PY{p}{,}\PY{+w}{ }\PY{l+m+mi}{2}\PY{+w}{ }\PY{o}{+}\PY{+w}{ }\PY{l+m+mi}{2}\PY{p}{)}\PY{p}{;}
\PY{+w}{    }\PY{n}{assertTrue}\PY{p}{(}\PY{n}{List}\PY{p}{.}\PY{n+na}{of}\PY{p}{(}\PY{l+m+mi}{1}\PY{p}{,}\PY{+w}{ }\PY{l+m+mi}{2}\PY{p}{,}\PY{+w}{ }\PY{l+m+mi}{3}\PY{p}{)}\PY{p}{.}\PY{n+na}{contains}\PY{p}{(}\PY{l+m+mi}{2}\PY{p}{)}\PY{p}{)}\PY{p}{;}
\PY{+w}{    }\PY{n}{assertFalse}\PY{p}{(}\PY{n}{List}\PY{p}{.}\PY{n+na}{of}\PY{p}{(}\PY{l+m+mi}{1}\PY{p}{,}\PY{+w}{ }\PY{l+m+mi}{2}\PY{p}{,}\PY{+w}{ }\PY{l+m+mi}{3}\PY{p}{)}\PY{p}{.}\PY{n+na}{isEmpty}\PY{p}{(}\PY{p}{)}\PY{p}{)}\PY{p}{;}
\PY{+w}{    }\PY{n}{assertNull}\PY{p}{(}\PY{n}{findUser}\PY{p}{(}\PY{l+s}{\PYZdq{}}\PY{l+s}{missing}\PY{l+s}{\PYZdq{}}\PY{p}{)}\PY{p}{)}\PY{p}{;}
\PY{+w}{    }\PY{n}{assertNotNull}\PY{p}{(}\PY{n}{findUser}\PY{p}{(}\PY{l+s}{\PYZdq{}}\PY{l+s}{alice}\PY{l+s}{\PYZdq{}}\PY{p}{)}\PY{p}{)}\PY{p}{;}

\PY{+w}{    }\PY{n}{assertIterableEquals}\PY{p}{(}\PY{n}{List}\PY{p}{.}\PY{n+na}{of}\PY{p}{(}\PY{l+m+mi}{1}\PY{p}{,}\PY{+w}{ }\PY{l+m+mi}{2}\PY{p}{,}\PY{+w}{ }\PY{l+m+mi}{3}\PY{p}{)}\PY{p}{,}\PY{+w}{ }\PY{n}{computeSequence}\PY{p}{(}\PY{p}{)}\PY{p}{)}\PY{p}{;}

\PY{+w}{    }\PY{n}{IllegalArgumentException}\PY{+w}{ }\PY{n}{ex}\PY{+w}{ }\PY{o}{=}\PY{+w}{ }\PY{n}{assertThrows}\PY{p}{(}
\PY{+w}{        }\PY{n}{IllegalArgumentException}\PY{p}{.}\PY{n+na}{class}\PY{p}{,}
\PY{+w}{        }\PY{p}{(}\PY{p}{)}\PY{+w}{ }\PY{o}{\PYZhy{}}\PY{o}{\PYZgt{}}\PY{+w}{ }\PY{k}{new}\PY{+w}{ }\PY{n}{Account}\PY{p}{(}\PY{k}{new}\PY{+w}{ }\PY{n}{BigDecimal}\PY{p}{(}\PY{l+s}{\PYZdq{}}\PY{l+s}{\PYZhy{}1}\PY{l+s}{\PYZdq{}}\PY{p}{)}\PY{p}{)}
\PY{+w}{    }\PY{p}{)}\PY{p}{;}
\PY{+w}{    }\PY{n}{assertEquals}\PY{p}{(}\PY{l+s}{\PYZdq{}}\PY{l+s}{initial balance must not be negative}\PY{l+s}{\PYZdq{}}\PY{p}{,}\PY{+w}{ }\PY{n}{ex}\PY{p}{.}\PY{n+na}{getMessage}\PY{p}{(}\PY{p}{)}\PY{p}{)}\PY{p}{;}

\PY{+w}{    }\PY{n}{assertTimeout}\PY{p}{(}\PY{n}{Duration}\PY{p}{.}\PY{n+na}{ofMillis}\PY{p}{(}\PY{l+m+mi}{200}\PY{p}{)}\PY{p}{,}\PY{+w}{ }\PY{p}{(}\PY{p}{)}\PY{+w}{ }\PY{o}{\PYZhy{}}\PY{o}{\PYZgt{}}\PY{+w}{ }\PY{n}{slowButBoundedOperation}\PY{p}{(}\PY{p}{)}\PY{p}{)}\PY{p}{;}

\PY{+w}{    }\PY{c+c1}{// Groups related checks: ALL run even if one fails, and all failures are reported together}
\PY{+w}{    }\PY{n}{Order}\PY{+w}{ }\PY{n}{order}\PY{+w}{ }\PY{o}{=}\PY{+w}{ }\PY{n}{buildOrder}\PY{p}{(}\PY{p}{)}\PY{p}{;}
\PY{+w}{    }\PY{n}{assertAll}\PY{p}{(}\PY{l+s}{\PYZdq{}}\PY{l+s}{order invariants}\PY{l+s}{\PYZdq{}}\PY{p}{,}
\PY{+w}{        }\PY{p}{(}\PY{p}{)}\PY{+w}{ }\PY{o}{\PYZhy{}}\PY{o}{\PYZgt{}}\PY{+w}{ }\PY{n}{assertEquals}\PY{p}{(}\PY{l+s}{\PYZdq{}}\PY{l+s}{ORD\PYZhy{}1}\PY{l+s}{\PYZdq{}}\PY{p}{,}\PY{+w}{ }\PY{n}{order}\PY{p}{.}\PY{n+na}{id}\PY{p}{(}\PY{p}{)}\PY{p}{)}\PY{p}{,}
\PY{+w}{        }\PY{p}{(}\PY{p}{)}\PY{+w}{ }\PY{o}{\PYZhy{}}\PY{o}{\PYZgt{}}\PY{+w}{ }\PY{n}{assertEquals}\PY{p}{(}\PY{n}{OrderStatus}\PY{p}{.}\PY{n+na}{PENDING}\PY{p}{,}\PY{+w}{ }\PY{n}{order}\PY{p}{.}\PY{n+na}{status}\PY{p}{(}\PY{p}{)}\PY{p}{)}\PY{p}{,}
\PY{+w}{        }\PY{p}{(}\PY{p}{)}\PY{+w}{ }\PY{o}{\PYZhy{}}\PY{o}{\PYZgt{}}\PY{+w}{ }\PY{n}{assertFalse}\PY{p}{(}\PY{n}{order}\PY{p}{.}\PY{n+na}{lineItems}\PY{p}{(}\PY{p}{)}\PY{p}{.}\PY{n+na}{isEmpty}\PY{p}{(}\PY{p}{)}\PY{p}{)}
\PY{+w}{    }\PY{p}{)}\PY{p}{;}
\PY{p}{\PYZcb{}}
\end{Verbatim}
\end{codebox}

\code{assertAll} is the key one reviewers look for: without it, three separate \code{assertEquals} calls stop at the first failure, hiding whether the second and third would also have failed. \code{assertAll} runs every supplied executable and reports every failure together — much more useful when triaging a broken test.

\code{assertTimeout} runs the code on the \emph{same} thread and waits for it to finish (so a hang will still block the test thread); \code{assert\allowbreak{}Timeout\allowbreak{}Preemptively} runs it on a separate thread and aborts at the deadline, but can leave background work running and has subtle interactions with \code{ThreadLocal}s — prefer \code{assertTimeout} unless you specifically need preemptive cancellation.

\subsubsection[AssertJ as a Fluent Alternative]{AssertJ as a Fluent Alternative}
\label{h:20:assertj-as-a-fluent-alternative}
JUnit’s static assertions are fine, but they read backwards (\code{assert\allowbreak{}Equals(\allowbreak{}expected,\allowbreak{}~\allowbreak{}actual)}) and don’t chain well. \textbf{AssertJ} provides a fluent, chainable, and far more readable API, and is the de-facto standard paired with JUnit 5 in most professional codebases:

\begin{codebox}
\begin{Verbatim}[commandchars=\\\{\}]
\PY{k+kn}{import static}\PY{+w}{ }\PY{n+nn}{org.assertj.core.api.Assertions.assertThat}\PY{p}{;}

\PY{n+nd}{@Test}
\PY{k+kt}{void}\PY{+w}{ }\PY{n+nf}{total\PYZus{}sumsLineItemPrices\PYZus{}assertj}\PY{p}{(}\PY{p}{)}\PY{+w}{ }\PY{p}{\PYZob{}}
\PY{+w}{    }\PY{n}{calculator}\PY{p}{.}\PY{n+na}{addLine}\PY{p}{(}\PY{l+s}{\PYZdq{}}\PY{l+s}{widget}\PY{l+s}{\PYZdq{}}\PY{p}{,}\PY{+w}{ }\PY{l+m+mi}{2}\PY{p}{,}\PY{+w}{ }\PY{k}{new}\PY{+w}{ }\PY{n}{BigDecimal}\PY{p}{(}\PY{l+s}{\PYZdq{}}\PY{l+s}{5.00}\PY{l+s}{\PYZdq{}}\PY{p}{)}\PY{p}{)}\PY{p}{;}
\PY{+w}{    }\PY{n}{calculator}\PY{p}{.}\PY{n+na}{addLine}\PY{p}{(}\PY{l+s}{\PYZdq{}}\PY{l+s}{gadget}\PY{l+s}{\PYZdq{}}\PY{p}{,}\PY{+w}{ }\PY{l+m+mi}{1}\PY{p}{,}\PY{+w}{ }\PY{k}{new}\PY{+w}{ }\PY{n}{BigDecimal}\PY{p}{(}\PY{l+s}{\PYZdq{}}\PY{l+s}{12.50}\PY{l+s}{\PYZdq{}}\PY{p}{)}\PY{p}{)}\PY{p}{;}

\PY{+w}{    }\PY{n}{assertThat}\PY{p}{(}\PY{n}{calculator}\PY{p}{.}\PY{n+na}{total}\PY{p}{(}\PY{p}{)}\PY{p}{)}\PY{p}{.}\PY{n+na}{isEqualByComparingTo}\PY{p}{(}\PY{l+s}{\PYZdq{}}\PY{l+s}{22.50}\PY{l+s}{\PYZdq{}}\PY{p}{)}\PY{p}{;}

\PY{+w}{    }\PY{n}{assertThat}\PY{p}{(}\PY{n}{calculator}\PY{p}{.}\PY{n+na}{lineItems}\PY{p}{(}\PY{p}{)}\PY{p}{)}
\PY{+w}{        }\PY{p}{.}\PY{n+na}{hasSize}\PY{p}{(}\PY{l+m+mi}{2}\PY{p}{)}
\PY{+w}{        }\PY{p}{.}\PY{n+na}{extracting}\PY{p}{(}\PY{l+s}{\PYZdq{}}\PY{l+s}{name}\PY{l+s}{\PYZdq{}}\PY{p}{)}
\PY{+w}{        }\PY{p}{.}\PY{n+na}{containsExactly}\PY{p}{(}\PY{l+s}{\PYZdq{}}\PY{l+s}{widget}\PY{l+s}{\PYZdq{}}\PY{p}{,}\PY{+w}{ }\PY{l+s}{\PYZdq{}}\PY{l+s}{gadget}\PY{l+s}{\PYZdq{}}\PY{p}{)}\PY{p}{;}

\PY{+w}{    }\PY{n}{assertThatThrownBy}\PY{p}{(}\PY{p}{(}\PY{p}{)}\PY{+w}{ }\PY{o}{\PYZhy{}}\PY{o}{\PYZgt{}}\PY{+w}{ }\PY{n}{calculator}\PY{p}{.}\PY{n+na}{addLine}\PY{p}{(}\PY{l+s}{\PYZdq{}}\PY{l+s}{bad}\PY{l+s}{\PYZdq{}}\PY{p}{,}\PY{+w}{ }\PY{o}{\PYZhy{}}\PY{l+m+mi}{1}\PY{p}{,}\PY{+w}{ }\PY{n}{BigDecimal}\PY{p}{.}\PY{n+na}{TEN}\PY{p}{)}\PY{p}{)}
\PY{+w}{        }\PY{p}{.}\PY{n+na}{isInstanceOf}\PY{p}{(}\PY{n}{IllegalArgumentException}\PY{p}{.}\PY{n+na}{class}\PY{p}{)}
\PY{+w}{        }\PY{p}{.}\PY{n+na}{hasMessageContaining}\PY{p}{(}\PY{l+s}{\PYZdq{}}\PY{l+s}{quantity}\PY{l+s}{\PYZdq{}}\PY{p}{)}\PY{p}{;}
\PY{p}{\PYZcb{}}
\end{Verbatim}
\end{codebox}

AssertJ’s chained assertions produce much richer failure messages out of the box (e.g., showing exactly which elements differ in a collection), and \code{assertThatThrownBy} combined with \code{.\allowbreak{}has\allowbreak{}Message\allowbreak{}Containing(...)} is more expressive than \code{assertThrows} + a separate \code{assertEquals} on the message.

\subsection[Nested Tests]{Nested Tests}
\label{h:20:nested-tests}
\code{@\allowbreak{}Nested} groups related tests inside inner classes, which is useful for expressing “given this context, these things should be true” structures and produces much clearer test reports:

\begin{codebox}
\begin{Verbatim}[commandchars=\\\{\}]
\PY{k+kn}{import}\PY{+w}{ }\PY{n+nn}{org.junit.jupiter.api.*}\PY{p}{;}
\PY{k+kn}{import static}\PY{+w}{ }\PY{n+nn}{org.junit.jupiter.api.Assertions.*}\PY{p}{;}

\PY{k+kd}{class} \PY{n+nc}{BankAccountTest}\PY{+w}{ }\PY{p}{\PYZob{}}

\PY{+w}{    }\PY{k+kd}{private}\PY{+w}{ }\PY{n}{BankAccount}\PY{+w}{ }\PY{n}{account}\PY{p}{;}

\PY{+w}{    }\PY{n+nd}{@BeforeEach}
\PY{+w}{    }\PY{k+kt}{void}\PY{+w}{ }\PY{n+nf}{setUp}\PY{p}{(}\PY{p}{)}\PY{+w}{ }\PY{p}{\PYZob{}}
\PY{+w}{        }\PY{n}{account}\PY{+w}{ }\PY{o}{=}\PY{+w}{ }\PY{k}{new}\PY{+w}{ }\PY{n}{BankAccount}\PY{p}{(}\PY{l+s}{\PYZdq{}}\PY{l+s}{acc\PYZhy{}1}\PY{l+s}{\PYZdq{}}\PY{p}{,}\PY{+w}{ }\PY{k}{new}\PY{+w}{ }\PY{n}{BigDecimal}\PY{p}{(}\PY{l+s}{\PYZdq{}}\PY{l+s}{100.00}\PY{l+s}{\PYZdq{}}\PY{p}{)}\PY{p}{)}\PY{p}{;}
\PY{+w}{    }\PY{p}{\PYZcb{}}

\PY{+w}{    }\PY{n+nd}{@Nested}
\PY{+w}{    }\PY{n+nd}{@DisplayName}\PY{p}{(}\PY{l+s}{\PYZdq{}}\PY{l+s}{when the account has a positive balance}\PY{l+s}{\PYZdq{}}\PY{p}{)}
\PY{+w}{    }\PY{k+kd}{class} \PY{n+nc}{WithPositiveBalance}\PY{+w}{ }\PY{p}{\PYZob{}}

\PY{+w}{        }\PY{n+nd}{@Test}
\PY{+w}{        }\PY{n+nd}{@DisplayName}\PY{p}{(}\PY{l+s}{\PYZdq{}}\PY{l+s}{withdraw() succeeds for amounts within balance}\PY{l+s}{\PYZdq{}}\PY{p}{)}
\PY{+w}{        }\PY{k+kt}{void}\PY{+w}{ }\PY{n+nf}{withdraw\PYZus{}withinBalance\PYZus{}succeeds}\PY{p}{(}\PY{p}{)}\PY{+w}{ }\PY{p}{\PYZob{}}
\PY{+w}{            }\PY{n}{account}\PY{p}{.}\PY{n+na}{withdraw}\PY{p}{(}\PY{k}{new}\PY{+w}{ }\PY{n}{BigDecimal}\PY{p}{(}\PY{l+s}{\PYZdq{}}\PY{l+s}{40.00}\PY{l+s}{\PYZdq{}}\PY{p}{)}\PY{p}{)}\PY{p}{;}
\PY{+w}{            }\PY{n}{assertEquals}\PY{p}{(}\PY{k}{new}\PY{+w}{ }\PY{n}{BigDecimal}\PY{p}{(}\PY{l+s}{\PYZdq{}}\PY{l+s}{60.00}\PY{l+s}{\PYZdq{}}\PY{p}{)}\PY{p}{,}\PY{+w}{ }\PY{n}{account}\PY{p}{.}\PY{n+na}{balance}\PY{p}{(}\PY{p}{)}\PY{p}{)}\PY{p}{;}
\PY{+w}{        }\PY{p}{\PYZcb{}}

\PY{+w}{        }\PY{n+nd}{@Test}
\PY{+w}{        }\PY{n+nd}{@DisplayName}\PY{p}{(}\PY{l+s}{\PYZdq{}}\PY{l+s}{withdraw() throws for amounts above balance}\PY{l+s}{\PYZdq{}}\PY{p}{)}
\PY{+w}{        }\PY{k+kt}{void}\PY{+w}{ }\PY{n+nf}{withdraw\PYZus{}aboveBalance\PYZus{}throws}\PY{p}{(}\PY{p}{)}\PY{+w}{ }\PY{p}{\PYZob{}}
\PY{+w}{            }\PY{n}{assertThrows}\PY{p}{(}\PY{n}{InsufficientFundsException}\PY{p}{.}\PY{n+na}{class}\PY{p}{,}
\PY{+w}{                }\PY{p}{(}\PY{p}{)}\PY{+w}{ }\PY{o}{\PYZhy{}}\PY{o}{\PYZgt{}}\PY{+w}{ }\PY{n}{account}\PY{p}{.}\PY{n+na}{withdraw}\PY{p}{(}\PY{k}{new}\PY{+w}{ }\PY{n}{BigDecimal}\PY{p}{(}\PY{l+s}{\PYZdq{}}\PY{l+s}{500.00}\PY{l+s}{\PYZdq{}}\PY{p}{)}\PY{p}{)}\PY{p}{)}\PY{p}{;}
\PY{+w}{        }\PY{p}{\PYZcb{}}
\PY{+w}{    }\PY{p}{\PYZcb{}}

\PY{+w}{    }\PY{n+nd}{@Nested}
\PY{+w}{    }\PY{n+nd}{@DisplayName}\PY{p}{(}\PY{l+s}{\PYZdq{}}\PY{l+s}{when the account is frozen}\PY{l+s}{\PYZdq{}}\PY{p}{)}
\PY{+w}{    }\PY{k+kd}{class} \PY{n+nc}{WhenFrozen}\PY{+w}{ }\PY{p}{\PYZob{}}

\PY{+w}{        }\PY{n+nd}{@BeforeEach}
\PY{+w}{        }\PY{k+kt}{void}\PY{+w}{ }\PY{n+nf}{freeze}\PY{p}{(}\PY{p}{)}\PY{+w}{ }\PY{p}{\PYZob{}}
\PY{+w}{            }\PY{n}{account}\PY{p}{.}\PY{n+na}{freeze}\PY{p}{(}\PY{p}{)}\PY{p}{;}
\PY{+w}{        }\PY{p}{\PYZcb{}}

\PY{+w}{        }\PY{n+nd}{@Test}
\PY{+w}{        }\PY{n+nd}{@DisplayName}\PY{p}{(}\PY{l+s}{\PYZdq{}}\PY{l+s}{withdraw() throws regardless of balance}\PY{l+s}{\PYZdq{}}\PY{p}{)}
\PY{+w}{        }\PY{k+kt}{void}\PY{+w}{ }\PY{n+nf}{withdraw\PYZus{}whenFrozen\PYZus{}throws}\PY{p}{(}\PY{p}{)}\PY{+w}{ }\PY{p}{\PYZob{}}
\PY{+w}{            }\PY{n}{assertThrows}\PY{p}{(}\PY{n}{AccountFrozenException}\PY{p}{.}\PY{n+na}{class}\PY{p}{,}
\PY{+w}{                }\PY{p}{(}\PY{p}{)}\PY{+w}{ }\PY{o}{\PYZhy{}}\PY{o}{\PYZgt{}}\PY{+w}{ }\PY{n}{account}\PY{p}{.}\PY{n+na}{withdraw}\PY{p}{(}\PY{k}{new}\PY{+w}{ }\PY{n}{BigDecimal}\PY{p}{(}\PY{l+s}{\PYZdq{}}\PY{l+s}{1.00}\PY{l+s}{\PYZdq{}}\PY{p}{)}\PY{p}{)}\PY{p}{)}\PY{p}{;}
\PY{+w}{        }\PY{p}{\PYZcb{}}
\PY{+w}{    }\PY{p}{\PYZcb{}}
\PY{p}{\PYZcb{}}
\end{Verbatim}
\end{codebox}

Each \code{@\allowbreak{}Nested} class gets its own \code{@\allowbreak{}BeforeEach} chain (outer then inner), so \code{WhenFrozen.\allowbreak{}freeze()} runs after the outer \code{setUp()}, letting you build up scenario-specific state incrementally.

\subsection[Tags and Filtering]{Tags and Filtering}
\label{h:20:tags-and-filtering}
\code{@\allowbreak{}Tag} lets you categorize tests (e.g., \code{slow}, \code{integration}, \code{smoke}) and then include or exclude categories at build time without touching test code:

\begin{codebox}
\begin{Verbatim}[commandchars=\\\{\}]
\PY{n+nd}{@Tag}\PY{p}{(}\PY{l+s}{\PYZdq{}}\PY{l+s}{integration}\PY{l+s}{\PYZdq{}}\PY{p}{)}
\PY{k+kd}{class} \PY{n+nc}{OrderRepositoryIntegrationTest}\PY{+w}{ }\PY{p}{\PYZob{}}

\PY{+w}{    }\PY{n+nd}{@Test}
\PY{+w}{    }\PY{n+nd}{@Tag}\PY{p}{(}\PY{l+s}{\PYZdq{}}\PY{l+s}{slow}\PY{l+s}{\PYZdq{}}\PY{p}{)}
\PY{+w}{    }\PY{k+kt}{void}\PY{+w}{ }\PY{n+nf}{savesAndReloadsOrder}\PY{p}{(}\PY{p}{)}\PY{+w}{ }\PY{p}{\PYZob{}}
\PY{+w}{        }\PY{c+c1}{// hits a real (test) database}
\PY{+w}{    }\PY{p}{\PYZcb{}}
\PY{p}{\PYZcb{}}
\end{Verbatim}
\end{codebox}

Maven Surefire:

\begin{codebox}
\begin{Verbatim}[commandchars=\\\{\}]
\PY{n+nt}{\PYZlt{}plugin}\PY{n+nt}{\PYZgt{}}
\PY{+w}{    }\PY{n+nt}{\PYZlt{}groupId}\PY{n+nt}{\PYZgt{}}org.apache.maven.plugins\PY{n+nt}{\PYZlt{}/groupId\PYZgt{}}
\PY{+w}{    }\PY{n+nt}{\PYZlt{}artifactId}\PY{n+nt}{\PYZgt{}}maven\PYZhy{}surefire\PYZhy{}plugin\PY{n+nt}{\PYZlt{}/artifactId\PYZgt{}}
\PY{+w}{    }\PY{n+nt}{\PYZlt{}configuration}\PY{n+nt}{\PYZgt{}}
\PY{+w}{        }\PY{n+nt}{\PYZlt{}excludedGroups}\PY{n+nt}{\PYZgt{}}slow,\PY{+w}{ }integration\PY{n+nt}{\PYZlt{}/excludedGroups\PYZgt{}}
\PY{+w}{    }\PY{n+nt}{\PYZlt{}/configuration\PYZgt{}}
\PY{n+nt}{\PYZlt{}/plugin\PYZgt{}}
\end{Verbatim}
\end{codebox}

Gradle:

\begin{codebox}
\begin{Verbatim}[commandchars=\\\{\}]
\PY{n}{tasks}\PY{p}{.}\PY{n+na}{test}\PY{+w}{ }\PY{p}{\PYZob{}}
\PY{+w}{    }\PY{n}{useJUnitPlatform}\PY{+w}{ }\PY{p}{\PYZob{}}
\PY{+w}{        }\PY{n}{excludeTags}\PY{p}{(}\PY{l+s}{\PYZdq{}}\PY{l+s}{slow}\PY{l+s}{\PYZdq{}}\PY{p}{,}\PY{+w}{ }\PY{l+s}{\PYZdq{}}\PY{l+s}{integration}\PY{l+s}{\PYZdq{}}\PY{p}{)}
\PY{+w}{    }\PY{p}{\PYZcb{}}
\PY{p}{\PYZcb{}}
\end{Verbatim}
\end{codebox}

This is how most CI pipelines run a fast “unit only” stage on every push and reserve a slower “integration” stage for merge queues or nightly runs.

\subsection[Disabling Tests]{Disabling Tests}
\label{h:20:disabling-tests}
\begin{codebox}
\begin{Verbatim}[commandchars=\\\{\}]
\PY{n+nd}{@Test}
\PY{n+nd}{@Disabled}\PY{p}{(}\PY{l+s}{\PYZdq{}}\PY{l+s}{Flaky due to JDK\PYZhy{}XXXX clock resolution bug on CI; see TICKET\PYZhy{}1234}\PY{l+s}{\PYZdq{}}\PY{p}{)}
\PY{k+kt}{void}\PY{+w}{ }\PY{n+nf}{timeSensitiveTest}\PY{p}{(}\PY{p}{)}\PY{+w}{ }\PY{p}{\PYZob{}}\PY{+w}{ }\PY{p}{.}\PY{p}{.}\PY{p}{.}\PY{+w}{ }\PY{p}{\PYZcb{}}
\end{Verbatim}
\end{codebox}

A reviewer should always ask \textbf{why} a test is disabled and whether there is a tracked ticket — \code{@\allowbreak{}Disabled} with no reason or follow-up is a silent hole in coverage that tends to be forgotten forever.

\subsection[Assumptions]{Assumptions}
\label{h:20:assumptions}
\code{Assumptions.\allowbreak{}assume\allowbreak{}True(...)} aborts (not fails) a test when a precondition isn’t met — useful for environment-dependent tests you don’t want counted as failures:

\begin{codebox}
\begin{Verbatim}[commandchars=\\\{\}]
\PY{k+kn}{import static}\PY{+w}{ }\PY{n+nn}{org.junit.jupiter.api.Assumptions.*}\PY{p}{;}

\PY{n+nd}{@Test}
\PY{k+kt}{void}\PY{+w}{ }\PY{n+nf}{nativeLibraryTest}\PY{p}{(}\PY{p}{)}\PY{+w}{ }\PY{p}{\PYZob{}}
\PY{+w}{    }\PY{n}{assumeTrue}\PY{p}{(}\PY{n}{System}\PY{p}{.}\PY{n+na}{getProperty}\PY{p}{(}\PY{l+s}{\PYZdq{}}\PY{l+s}{os.name}\PY{l+s}{\PYZdq{}}\PY{p}{)}\PY{p}{.}\PY{n+na}{toLowerCase}\PY{p}{(}\PY{p}{)}\PY{p}{.}\PY{n+na}{contains}\PY{p}{(}\PY{l+s}{\PYZdq{}}\PY{l+s}{linux}\PY{l+s}{\PYZdq{}}\PY{p}{)}\PY{p}{,}
\PY{+w}{        }\PY{l+s}{\PYZdq{}}\PY{l+s}{This test only runs on Linux}\PY{l+s}{\PYZdq{}}\PY{p}{)}\PY{p}{;}

\PY{+w}{    }\PY{c+c1}{// ... test that depends on a Linux\PYZhy{}only native library}
\PY{p}{\PYZcb{}}
\end{Verbatim}
\end{codebox}

If the assumption fails, JUnit marks the test as \textbf{skipped}, not failed — it will not turn a CI build red, but it also will not silently pass, so it still shows up in the report.

\subsection[Parameterized Tests]{Parameterized Tests}
\label{h:20:parameterized-tests}
Parameterized tests let you run the same test logic against many inputs without copy-pasting the method.

\begin{codebox}
\begin{Verbatim}[commandchars=\\\{\}]
\PY{k+kn}{import}\PY{+w}{ }\PY{n+nn}{org.junit.jupiter.params.ParameterizedTest}\PY{p}{;}
\PY{k+kn}{import}\PY{+w}{ }\PY{n+nn}{org.junit.jupiter.params.provider.*}\PY{p}{;}
\PY{k+kn}{import static}\PY{+w}{ }\PY{n+nn}{org.junit.jupiter.api.Assertions.*}\PY{p}{;}

\PY{k+kd}{class} \PY{n+nc}{PasswordValidatorTest}\PY{+w}{ }\PY{p}{\PYZob{}}

\PY{+w}{    }\PY{k+kd}{private}\PY{+w}{ }\PY{k+kd}{final}\PY{+w}{ }\PY{n}{PasswordValidator}\PY{+w}{ }\PY{n}{validator}\PY{+w}{ }\PY{o}{=}\PY{+w}{ }\PY{k}{new}\PY{+w}{ }\PY{n}{PasswordValidator}\PY{p}{(}\PY{p}{)}\PY{p}{;}

\PY{+w}{    }\PY{n+nd}{@ParameterizedTest}
\PY{+w}{    }\PY{n+nd}{@ValueSource}\PY{p}{(}\PY{n}{strings}\PY{+w}{ }\PY{o}{=}\PY{+w}{ }\PY{p}{\PYZob{}}\PY{+w}{ }\PY{l+s}{\PYZdq{}}\PY{l+s}{abc}\PY{l+s}{\PYZdq{}}\PY{p}{,}\PY{+w}{ }\PY{l+s}{\PYZdq{}}\PY{l+s}{12345}\PY{l+s}{\PYZdq{}}\PY{p}{,}\PY{+w}{ }\PY{l+s}{\PYZdq{}}\PY{l+s}{        }\PY{l+s}{\PYZdq{}}\PY{p}{,}\PY{+w}{ }\PY{l+s}{\PYZdq{}}\PY{l+s}{\PYZdq{}}\PY{+w}{ }\PY{p}{\PYZcb{}}\PY{p}{)}
\PY{+w}{    }\PY{k+kt}{void}\PY{+w}{ }\PY{n+nf}{isValid\PYZus{}rejectsShortOrBlankPasswords}\PY{p}{(}\PY{n}{String}\PY{+w}{ }\PY{n}{candidate}\PY{p}{)}\PY{+w}{ }\PY{p}{\PYZob{}}
\PY{+w}{        }\PY{n}{assertFalse}\PY{p}{(}\PY{n}{validator}\PY{p}{.}\PY{n+na}{isValid}\PY{p}{(}\PY{n}{candidate}\PY{p}{)}\PY{p}{)}\PY{p}{;}
\PY{+w}{    }\PY{p}{\PYZcb{}}

\PY{+w}{    }\PY{n+nd}{@ParameterizedTest}
\PY{+w}{    }\PY{n+nd}{@CsvSource}\PY{p}{(}\PY{p}{\PYZob{}}
\PY{+w}{        }\PY{l+s}{\PYZdq{}}\PY{l+s}{Password1!, true}\PY{l+s}{\PYZdq{}}\PY{p}{,}
\PY{+w}{        }\PY{l+s}{\PYZdq{}}\PY{l+s}{password1!, false}\PY{l+s}{\PYZdq{}}\PY{p}{,}\PY{+w}{   }\PY{c+c1}{// no uppercase}
\PY{+w}{        }\PY{l+s}{\PYZdq{}}\PY{l+s}{PASSWORD1!, false}\PY{l+s}{\PYZdq{}}\PY{p}{,}\PY{+w}{   }\PY{c+c1}{// no lowercase}
\PY{+w}{        }\PY{l+s}{\PYZdq{}}\PY{l+s}{Password!,  false}\PY{l+s}{\PYZdq{}}\PY{p}{,}\PY{+w}{  }\PY{c+c1}{// no digit}
\PY{+w}{        }\PY{l+s}{\PYZdq{}}\PY{l+s}{Password1,  false}\PY{l+s}{\PYZdq{}}\PY{+w}{   }\PY{c+c1}{// no symbol}
\PY{+w}{    }\PY{p}{\PYZcb{}}\PY{p}{)}
\PY{+w}{    }\PY{k+kt}{void}\PY{+w}{ }\PY{n+nf}{isValid\PYZus{}enforcesComplexityRules}\PY{p}{(}\PY{n}{String}\PY{+w}{ }\PY{n}{candidate}\PY{p}{,}\PY{+w}{ }\PY{k+kt}{boolean}\PY{+w}{ }\PY{n}{expected}\PY{p}{)}\PY{+w}{ }\PY{p}{\PYZob{}}
\PY{+w}{        }\PY{n}{assertEquals}\PY{p}{(}\PY{n}{expected}\PY{p}{,}\PY{+w}{ }\PY{n}{validator}\PY{p}{.}\PY{n+na}{isValid}\PY{p}{(}\PY{n}{candidate}\PY{p}{)}\PY{p}{)}\PY{p}{;}
\PY{+w}{    }\PY{p}{\PYZcb{}}

\PY{+w}{    }\PY{n+nd}{@ParameterizedTest}
\PY{+w}{    }\PY{n+nd}{@MethodSource}\PY{p}{(}\PY{l+s}{\PYZdq{}}\PY{l+s}{complexPasswordCases}\PY{l+s}{\PYZdq{}}\PY{p}{)}
\PY{+w}{    }\PY{k+kt}{void}\PY{+w}{ }\PY{n+nf}{isValid\PYZus{}matchesMethodSourceCases}\PY{p}{(}\PY{n}{String}\PY{+w}{ }\PY{n}{candidate}\PY{p}{,}\PY{+w}{ }\PY{k+kt}{boolean}\PY{+w}{ }\PY{n}{expected}\PY{p}{)}\PY{+w}{ }\PY{p}{\PYZob{}}
\PY{+w}{        }\PY{n}{assertEquals}\PY{p}{(}\PY{n}{expected}\PY{p}{,}\PY{+w}{ }\PY{n}{validator}\PY{p}{.}\PY{n+na}{isValid}\PY{p}{(}\PY{n}{candidate}\PY{p}{)}\PY{p}{)}\PY{p}{;}
\PY{+w}{    }\PY{p}{\PYZcb{}}

\PY{+w}{    }\PY{k+kd}{static}\PY{+w}{ }\PY{n}{Stream}\PY{o}{\PYZlt{}}\PY{n}{Arguments}\PY{o}{\PYZgt{}}\PY{+w}{ }\PY{n+nf}{complexPasswordCases}\PY{p}{(}\PY{p}{)}\PY{+w}{ }\PY{p}{\PYZob{}}
\PY{+w}{        }\PY{k}{return}\PY{+w}{ }\PY{n}{Stream}\PY{p}{.}\PY{n+na}{of}\PY{p}{(}
\PY{+w}{            }\PY{n}{Arguments}\PY{p}{.}\PY{n+na}{of}\PY{p}{(}\PY{l+s}{\PYZdq{}}\PY{l+s}{Correct1!Horse}\PY{l+s}{\PYZdq{}}\PY{p}{,}\PY{+w}{ }\PY{k+kc}{true}\PY{p}{)}\PY{p}{,}
\PY{+w}{            }\PY{n}{Arguments}\PY{p}{.}\PY{n+na}{of}\PY{p}{(}\PY{l+s}{\PYZdq{}}\PY{l+s}{nope}\PY{l+s}{\PYZdq{}}\PY{p}{,}\PY{+w}{ }\PY{k+kc}{false}\PY{p}{)}\PY{p}{,}
\PY{+w}{            }\PY{n}{Arguments}\PY{p}{.}\PY{n+na}{of}\PY{p}{(}\PY{l+s}{\PYZdq{}}\PY{l+s}{Tr0ub4dor\PYZam{}3}\PY{l+s}{\PYZdq{}}\PY{p}{,}\PY{+w}{ }\PY{k+kc}{true}\PY{p}{)}
\PY{+w}{        }\PY{p}{)}\PY{p}{;}
\PY{+w}{    }\PY{p}{\PYZcb{}}

\PY{+w}{    }\PY{n+nd}{@ParameterizedTest}
\PY{+w}{    }\PY{n+nd}{@EnumSource}\PY{p}{(}\PY{n}{value}\PY{+w}{ }\PY{o}{=}\PY{+w}{ }\PY{n}{PasswordStrength}\PY{p}{.}\PY{n+na}{class}\PY{p}{,}\PY{+w}{ }\PY{n}{names}\PY{+w}{ }\PY{o}{=}\PY{+w}{ }\PY{l+s}{\PYZdq{}}\PY{l+s}{WEAK}\PY{l+s}{\PYZdq{}}\PY{p}{,}\PY{+w}{ }\PY{n}{mode}\PY{+w}{ }\PY{o}{=}\PY{+w}{ }\PY{n}{EnumSource}\PY{p}{.}\PY{n+na}{Mode}\PY{p}{.}\PY{n+na}{EXCLUDE}\PY{p}{)}
\PY{+w}{    }\PY{k+kt}{void}\PY{+w}{ }\PY{n+nf}{isValid\PYZus{}acceptsNonWeakStrengths}\PY{p}{(}\PY{n}{PasswordStrength}\PY{+w}{ }\PY{n}{strength}\PY{p}{)}\PY{+w}{ }\PY{p}{\PYZob{}}
\PY{+w}{        }\PY{n}{assertTrue}\PY{p}{(}\PY{n}{validator}\PY{p}{.}\PY{n+na}{isValid}\PY{p}{(}\PY{n}{strength}\PY{p}{.}\PY{n+na}{samplePassword}\PY{p}{(}\PY{p}{)}\PY{p}{)}\PY{p}{)}\PY{p}{;}
\PY{+w}{    }\PY{p}{\PYZcb{}}

\PY{+w}{    }\PY{n+nd}{@ParameterizedTest}
\PY{+w}{    }\PY{n+nd}{@ArgumentsSource}\PY{p}{(}\PY{n}{SqlInjectionPayloadsProvider}\PY{p}{.}\PY{n+na}{class}\PY{p}{)}
\PY{+w}{    }\PY{k+kt}{void}\PY{+w}{ }\PY{n+nf}{isValid\PYZus{}rejectsSqlInjectionLikePayloads}\PY{p}{(}\PY{n}{String}\PY{+w}{ }\PY{n}{payload}\PY{p}{)}\PY{+w}{ }\PY{p}{\PYZob{}}
\PY{+w}{        }\PY{n}{assertFalse}\PY{p}{(}\PY{n}{validator}\PY{p}{.}\PY{n+na}{isValid}\PY{p}{(}\PY{n}{payload}\PY{p}{)}\PY{p}{)}\PY{p}{;}
\PY{+w}{    }\PY{p}{\PYZcb{}}

\PY{+w}{    }\PY{k+kd}{static}\PY{+w}{ }\PY{k+kd}{class} \PY{n+nc}{SqlInjectionPayloadsProvider}\PY{+w}{ }\PY{k+kd}{implements}\PY{+w}{ }\PY{n}{ArgumentsProvider}\PY{+w}{ }\PY{p}{\PYZob{}}
\PY{+w}{        }\PY{n+nd}{@Override}
\PY{+w}{        }\PY{k+kd}{public}\PY{+w}{ }\PY{n}{Stream}\PY{o}{\PYZlt{}}\PY{o}{?}\PY{+w}{ }\PY{k+kd}{extends}\PY{+w}{ }\PY{n}{Arguments}\PY{o}{\PYZgt{}}\PY{+w}{ }\PY{n+nf}{provideArguments}\PY{p}{(}\PY{n}{ExtensionContext}\PY{+w}{ }\PY{n}{context}\PY{p}{)}\PY{+w}{ }\PY{p}{\PYZob{}}
\PY{+w}{            }\PY{k}{return}\PY{+w}{ }\PY{n}{Stream}\PY{p}{.}\PY{n+na}{of}\PY{p}{(}\PY{l+s}{\PYZdq{}}\PY{l+s}{\PYZsq{} OR \PYZsq{}1\PYZsq{}=\PYZsq{}1}\PY{l+s}{\PYZdq{}}\PY{p}{,}\PY{+w}{ }\PY{l+s}{\PYZdq{}}\PY{l+s}{\PYZsq{}; DROP TABLE users; \PYZhy{}\PYZhy{}}\PY{l+s}{\PYZdq{}}\PY{p}{)}
\PY{+w}{                          }\PY{p}{.}\PY{n+na}{map}\PY{p}{(}\PY{n}{Arguments}\PY{p}{:}\PY{p}{:}\PY{n}{of}\PY{p}{)}\PY{p}{;}
\PY{+w}{        }\PY{p}{\PYZcb{}}
\PY{+w}{    }\PY{p}{\PYZcb{}}
\PY{p}{\PYZcb{}}
\end{Verbatim}
\end{codebox}

{\tablefont
\begin{xltabular}{\linewidth}{@{}L{0.429}L{1.571}@{}}
\toprule
\rowcolor{tablehdr}
\thd{Source} & \thd{Best for} \\
\midrule
\endhead
\code{@\allowbreak{}ValueSource} & A flat list of one primitive/String argument each \\
\code{@\allowbreak{}CsvSource} & A handful of inline multi-argument rows \\
\code{@\allowbreak{}MethodSource} & Complex objects, or cases shared across multiple test methods \\
\code{@\allowbreak{}EnumSource} & Iterating over some or all values of an enum \\
\code{@\allowbreak{}ArgumentsSource} & A reusable, custom provider you want to share across test classes \\
\bottomrule
\end{xltabular}
}

Parameterized tests are one of the highest-leverage tools in the whole framework: a reviewer seeing five nearly identical \code{@\allowbreak{}Test} methods that differ only in input/output values should almost always suggest collapsing them into one \code{@\allowbreak{}ParameterizedTest}.

\subsection[Repeated Tests]{Repeated Tests}
\label{h:20:repeated-tests}
\code{@\allowbreak{}RepeatedTest} runs the same test multiple times — useful for surfacing flakiness in something that has an element of non-determinism (e.g., relies on system time, randomness, or concurrency) during investigation, though it is not a substitute for fixing the non-determinism.

\begin{codebox}
\begin{Verbatim}[commandchars=\\\{\}]
\PY{n+nd}{@RepeatedTest}\PY{p}{(}\PY{l+m+mi}{10}\PY{p}{)}
\PY{k+kt}{void}\PY{+w}{ }\PY{n+nf}{idGenerator\PYZus{}neverProducesDuplicateIds}\PY{p}{(}\PY{n}{RepetitionInfo}\PY{+w}{ }\PY{n}{info}\PY{p}{)}\PY{+w}{ }\PY{p}{\PYZob{}}
\PY{+w}{    }\PY{n}{String}\PY{+w}{ }\PY{n}{id}\PY{+w}{ }\PY{o}{=}\PY{+w}{ }\PY{n}{IdGenerator}\PY{p}{.}\PY{n+na}{next}\PY{p}{(}\PY{p}{)}\PY{p}{;}
\PY{+w}{    }\PY{n}{assertTrue}\PY{p}{(}\PY{n}{seenIds}\PY{p}{.}\PY{n+na}{add}\PY{p}{(}\PY{n}{id}\PY{p}{)}\PY{p}{,}\PY{+w}{ }\PY{l+s}{\PYZdq{}}\PY{l+s}{Duplicate on repetition }\PY{l+s}{\PYZdq{}}\PY{+w}{ }\PY{o}{+}\PY{+w}{ }\PY{n}{info}\PY{p}{.}\PY{n+na}{getCurrentRepetition}\PY{p}{(}\PY{p}{)}\PY{p}{)}\PY{p}{;}
\PY{p}{\PYZcb{}}
\end{Verbatim}
\end{codebox}

\subsection[Dynamic Tests with @TestFactory]{Dynamic Tests with \code{@\allowbreak{}TestFactory}}
\label{h:20:dynamic-tests-with-testfactory}
Regular \code{@\allowbreak{}Test} methods are fixed at compile time. \code{@\allowbreak{}TestFactory} lets you \emph{generate} tests at runtime, returning a \code{Stream}/\code{Collection}/\code{Iterable} of \code{DynamicTest}:

\begin{codebox}
\begin{Verbatim}[commandchars=\\\{\}]
\PY{k+kn}{import}\PY{+w}{ }\PY{n+nn}{org.junit.jupiter.api.DynamicTest}\PY{p}{;}
\PY{k+kn}{import}\PY{+w}{ }\PY{n+nn}{org.junit.jupiter.api.TestFactory}\PY{p}{;}
\PY{k+kn}{import static}\PY{+w}{ }\PY{n+nn}{org.junit.jupiter.api.DynamicTest.dynamicTest}\PY{p}{;}

\PY{k+kd}{class} \PY{n+nc}{FizzBuzzDynamicTest}\PY{+w}{ }\PY{p}{\PYZob{}}

\PY{+w}{    }\PY{n+nd}{@TestFactory}
\PY{+w}{    }\PY{n}{Stream}\PY{o}{\PYZlt{}}\PY{n}{DynamicTest}\PY{o}{\PYZgt{}}\PY{+w}{ }\PY{n+nf}{fizzBuzzCases}\PY{p}{(}\PY{p}{)}\PY{+w}{ }\PY{p}{\PYZob{}}
\PY{+w}{        }\PY{n}{Map}\PY{o}{\PYZlt{}}\PY{n}{Integer}\PY{p}{,}\PY{+w}{ }\PY{n}{String}\PY{o}{\PYZgt{}}\PY{+w}{ }\PY{n}{expectations}\PY{+w}{ }\PY{o}{=}\PY{+w}{ }\PY{n}{Map}\PY{p}{.}\PY{n+na}{of}\PY{p}{(}
\PY{+w}{            }\PY{l+m+mi}{1}\PY{p}{,}\PY{+w}{ }\PY{l+s}{\PYZdq{}}\PY{l+s}{1}\PY{l+s}{\PYZdq{}}\PY{p}{,}\PY{+w}{ }\PY{l+m+mi}{3}\PY{p}{,}\PY{+w}{ }\PY{l+s}{\PYZdq{}}\PY{l+s}{Fizz}\PY{l+s}{\PYZdq{}}\PY{p}{,}\PY{+w}{ }\PY{l+m+mi}{5}\PY{p}{,}\PY{+w}{ }\PY{l+s}{\PYZdq{}}\PY{l+s}{Buzz}\PY{l+s}{\PYZdq{}}\PY{p}{,}\PY{+w}{ }\PY{l+m+mi}{15}\PY{p}{,}\PY{+w}{ }\PY{l+s}{\PYZdq{}}\PY{l+s}{FizzBuzz}\PY{l+s}{\PYZdq{}}
\PY{+w}{        }\PY{p}{)}\PY{p}{;}

\PY{+w}{        }\PY{k}{return}\PY{+w}{ }\PY{n}{expectations}\PY{p}{.}\PY{n+na}{entrySet}\PY{p}{(}\PY{p}{)}\PY{p}{.}\PY{n+na}{stream}\PY{p}{(}\PY{p}{)}
\PY{+w}{            }\PY{p}{.}\PY{n+na}{map}\PY{p}{(}\PY{n}{entry}\PY{+w}{ }\PY{o}{\PYZhy{}}\PY{o}{\PYZgt{}}\PY{+w}{ }\PY{n}{dynamicTest}\PY{p}{(}
\PY{+w}{                }\PY{l+s}{\PYZdq{}}\PY{l+s}{fizzBuzz(}\PY{l+s}{\PYZdq{}}\PY{+w}{ }\PY{o}{+}\PY{+w}{ }\PY{n}{entry}\PY{p}{.}\PY{n+na}{getKey}\PY{p}{(}\PY{p}{)}\PY{+w}{ }\PY{o}{+}\PY{+w}{ }\PY{l+s}{\PYZdq{}}\PY{l+s}{) == }\PY{l+s}{\PYZdq{}}\PY{+w}{ }\PY{o}{+}\PY{+w}{ }\PY{n}{entry}\PY{p}{.}\PY{n+na}{getValue}\PY{p}{(}\PY{p}{)}\PY{p}{,}
\PY{+w}{                }\PY{p}{(}\PY{p}{)}\PY{+w}{ }\PY{o}{\PYZhy{}}\PY{o}{\PYZgt{}}\PY{+w}{ }\PY{n}{assertEquals}\PY{p}{(}\PY{n}{entry}\PY{p}{.}\PY{n+na}{getValue}\PY{p}{(}\PY{p}{)}\PY{p}{,}\PY{+w}{ }\PY{n}{FizzBuzz}\PY{p}{.}\PY{n+na}{of}\PY{p}{(}\PY{n}{entry}\PY{p}{.}\PY{n+na}{getKey}\PY{p}{(}\PY{p}{)}\PY{p}{)}\PY{p}{)}
\PY{+w}{            }\PY{p}{)}\PY{p}{)}\PY{p}{;}
\PY{+w}{    }\PY{p}{\PYZcb{}}
\PY{p}{\PYZcb{}}
\end{Verbatim}
\end{codebox}

Dynamic tests are less common than parameterized tests in day-to-day code but show up when test cases themselves need to be computed (e.g., generated from a spec file or a database of fixtures).

\subsection[Extensions]{Extensions}
\label{h:20:extensions}
Jupiter replaces JUnit 4’s \code{@\allowbreak{}Rule}/\code{@\allowbreak{}RunWith} with a single, composable \code{Extension} model, activated with \code{@\allowbreak{}ExtendWith}.

\begin{codebox}
\begin{Verbatim}[commandchars=\\\{\}]
\PY{k+kd}{class} \PY{n+nc}{LoggingTest}\PY{+w}{ }\PY{p}{\PYZob{}}
\PY{+w}{    }\PY{c+c1}{// Third\PYZhy{}party or shared extensions plug in the same way Mockito\PYZsq{}s extension does}
\PY{p}{\PYZcb{}}
\end{Verbatim}
\end{codebox}

A custom extension that times every test and logs slow ones:

\begin{codebox}
\begin{Verbatim}[commandchars=\\\{\}]
\PY{k+kn}{import}\PY{+w}{ }\PY{n+nn}{org.junit.jupiter.api.extension.*}\PY{p}{;}

\PY{k+kd}{public}\PY{+w}{ }\PY{k+kd}{class} \PY{n+nc}{TimingExtension}\PY{+w}{ }\PY{k+kd}{implements}\PY{+w}{ }\PY{n}{BeforeTestExecutionCallback}\PY{p}{,}\PY{+w}{ }\PY{n}{AfterTestExecutionCallback}\PY{+w}{ }\PY{p}{\PYZob{}}

\PY{+w}{    }\PY{k+kd}{private}\PY{+w}{ }\PY{k+kd}{static}\PY{+w}{ }\PY{k+kd}{final}\PY{+w}{ }\PY{n}{ExtensionContext}\PY{p}{.}\PY{n+na}{Namespace}\PY{+w}{ }\PY{n}{NAMESPACE}\PY{+w}{ }\PY{o}{=}
\PY{+w}{        }\PY{n}{ExtensionContext}\PY{p}{.}\PY{n+na}{Namespace}\PY{p}{.}\PY{n+na}{create}\PY{p}{(}\PY{n}{TimingExtension}\PY{p}{.}\PY{n+na}{class}\PY{p}{)}\PY{p}{;}

\PY{+w}{    }\PY{n+nd}{@Override}
\PY{+w}{    }\PY{k+kd}{public}\PY{+w}{ }\PY{k+kt}{void}\PY{+w}{ }\PY{n+nf}{beforeTestExecution}\PY{p}{(}\PY{n}{ExtensionContext}\PY{+w}{ }\PY{n}{context}\PY{p}{)}\PY{+w}{ }\PY{p}{\PYZob{}}
\PY{+w}{        }\PY{n}{getStore}\PY{p}{(}\PY{n}{context}\PY{p}{)}\PY{p}{.}\PY{n+na}{put}\PY{p}{(}\PY{l+s}{\PYZdq{}}\PY{l+s}{start}\PY{l+s}{\PYZdq{}}\PY{p}{,}\PY{+w}{ }\PY{n}{System}\PY{p}{.}\PY{n+na}{nanoTime}\PY{p}{(}\PY{p}{)}\PY{p}{)}\PY{p}{;}
\PY{+w}{    }\PY{p}{\PYZcb{}}

\PY{+w}{    }\PY{n+nd}{@Override}
\PY{+w}{    }\PY{k+kd}{public}\PY{+w}{ }\PY{k+kt}{void}\PY{+w}{ }\PY{n+nf}{afterTestExecution}\PY{p}{(}\PY{n}{ExtensionContext}\PY{+w}{ }\PY{n}{context}\PY{p}{)}\PY{+w}{ }\PY{p}{\PYZob{}}
\PY{+w}{        }\PY{k+kt}{long}\PY{+w}{ }\PY{n}{start}\PY{+w}{ }\PY{o}{=}\PY{+w}{ }\PY{n}{getStore}\PY{p}{(}\PY{n}{context}\PY{p}{)}\PY{p}{.}\PY{n+na}{remove}\PY{p}{(}\PY{l+s}{\PYZdq{}}\PY{l+s}{start}\PY{l+s}{\PYZdq{}}\PY{p}{,}\PY{+w}{ }\PY{k+kt}{long}\PY{p}{.}\PY{n+na}{class}\PY{p}{)}\PY{p}{;}
\PY{+w}{        }\PY{k+kt}{long}\PY{+w}{ }\PY{n}{durationMs}\PY{+w}{ }\PY{o}{=}\PY{+w}{ }\PY{p}{(}\PY{n}{System}\PY{p}{.}\PY{n+na}{nanoTime}\PY{p}{(}\PY{p}{)}\PY{+w}{ }\PY{o}{\PYZhy{}}\PY{+w}{ }\PY{n}{start}\PY{p}{)}\PY{+w}{ }\PY{o}{/}\PY{+w}{ }\PY{l+m+mi}{1\PYZus{}000\PYZus{}000}\PY{p}{;}
\PY{+w}{        }\PY{k}{if}\PY{+w}{ }\PY{p}{(}\PY{n}{durationMs}\PY{+w}{ }\PY{o}{\PYZgt{}}\PY{+w}{ }\PY{l+m+mi}{100}\PY{p}{)}\PY{+w}{ }\PY{p}{\PYZob{}}
\PY{+w}{            }\PY{n}{System}\PY{p}{.}\PY{n+na}{out}\PY{p}{.}\PY{n+na}{printf}\PY{p}{(}\PY{l+s}{\PYZdq{}}\PY{l+s}{SLOW TEST [\PYZpc{}s]: \PYZpc{}d ms\PYZpc{}n}\PY{l+s}{\PYZdq{}}\PY{p}{,}\PY{+w}{ }\PY{n}{context}\PY{p}{.}\PY{n+na}{getDisplayName}\PY{p}{(}\PY{p}{)}\PY{p}{,}\PY{+w}{ }\PY{n}{durationMs}\PY{p}{)}\PY{p}{;}
\PY{+w}{        }\PY{p}{\PYZcb{}}
\PY{+w}{    }\PY{p}{\PYZcb{}}

\PY{+w}{    }\PY{k+kd}{private}\PY{+w}{ }\PY{n}{ExtensionContext}\PY{p}{.}\PY{n+na}{Store}\PY{+w}{ }\PY{n+nf}{getStore}\PY{p}{(}\PY{n}{ExtensionContext}\PY{+w}{ }\PY{n}{context}\PY{p}{)}\PY{+w}{ }\PY{p}{\PYZob{}}
\PY{+w}{        }\PY{k}{return}\PY{+w}{ }\PY{n}{context}\PY{p}{.}\PY{n+na}{getStore}\PY{p}{(}\PY{n}{NAMESPACE}\PY{p}{)}\PY{p}{;}
\PY{+w}{    }\PY{p}{\PYZcb{}}
\PY{p}{\PYZcb{}}

\PY{n+nd}{@ExtendWith}\PY{p}{(}\PY{n}{TimingExtension}\PY{p}{.}\PY{n+na}{class}\PY{p}{)}
\PY{k+kd}{class} \PY{n+nc}{ReportGeneratorTest}\PY{+w}{ }\PY{p}{\PYZob{}}

\PY{+w}{    }\PY{n+nd}{@Test}
\PY{+w}{    }\PY{k+kt}{void}\PY{+w}{ }\PY{n+nf}{generatesMonthlyReport}\PY{p}{(}\PY{p}{)}\PY{+w}{ }\PY{p}{\PYZob{}}
\PY{+w}{        }\PY{c+c1}{// ...}
\PY{+w}{    }\PY{p}{\PYZcb{}}
\PY{p}{\PYZcb{}}
\end{Verbatim}
\end{codebox}

\code{@\allowbreak{}TempDir} is a built-in extension for tests that need a real filesystem directory that is automatically created before the test and deleted after:

\begin{codebox}
\begin{Verbatim}[commandchars=\\\{\}]
\PY{k+kn}{import}\PY{+w}{ }\PY{n+nn}{org.junit.jupiter.api.io.TempDir}\PY{p}{;}
\PY{k+kn}{import}\PY{+w}{ }\PY{n+nn}{java.nio.file.*}\PY{p}{;}

\PY{k+kd}{class} \PY{n+nc}{CsvExporterTest}\PY{+w}{ }\PY{p}{\PYZob{}}

\PY{+w}{    }\PY{n+nd}{@Test}
\PY{+w}{    }\PY{k+kt}{void}\PY{+w}{ }\PY{n+nf}{export\PYZus{}writesCsvFileWithHeader}\PY{p}{(}\PY{n+nd}{@TempDir}\PY{+w}{ }\PY{n}{Path}\PY{+w}{ }\PY{n}{tempDir}\PY{p}{)}\PY{+w}{ }\PY{k+kd}{throws}\PY{+w}{ }\PY{n}{IOException}\PY{+w}{ }\PY{p}{\PYZob{}}
\PY{+w}{        }\PY{n}{Path}\PY{+w}{ }\PY{n}{outputFile}\PY{+w}{ }\PY{o}{=}\PY{+w}{ }\PY{n}{tempDir}\PY{p}{.}\PY{n+na}{resolve}\PY{p}{(}\PY{l+s}{\PYZdq{}}\PY{l+s}{export.csv}\PY{l+s}{\PYZdq{}}\PY{p}{)}\PY{p}{;}

\PY{+w}{        }\PY{k}{new}\PY{+w}{ }\PY{n}{CsvExporter}\PY{p}{(}\PY{p}{)}\PY{p}{.}\PY{n+na}{export}\PY{p}{(}\PY{n}{List}\PY{p}{.}\PY{n+na}{of}\PY{p}{(}\PY{k}{new}\PY{+w}{ }\PY{n}{Row}\PY{p}{(}\PY{l+s}{\PYZdq{}}\PY{l+s}{a}\PY{l+s}{\PYZdq{}}\PY{p}{,}\PY{+w}{ }\PY{l+m+mi}{1}\PY{p}{)}\PY{p}{)}\PY{p}{,}\PY{+w}{ }\PY{n}{outputFile}\PY{p}{)}\PY{p}{;}

\PY{+w}{        }\PY{n}{List}\PY{o}{\PYZlt{}}\PY{n}{String}\PY{o}{\PYZgt{}}\PY{+w}{ }\PY{n}{lines}\PY{+w}{ }\PY{o}{=}\PY{+w}{ }\PY{n}{Files}\PY{p}{.}\PY{n+na}{readAllLines}\PY{p}{(}\PY{n}{outputFile}\PY{p}{)}\PY{p}{;}
\PY{+w}{        }\PY{n}{assertEquals}\PY{p}{(}\PY{l+s}{\PYZdq{}}\PY{l+s}{name,value}\PY{l+s}{\PYZdq{}}\PY{p}{,}\PY{+w}{ }\PY{n}{lines}\PY{p}{.}\PY{n+na}{get}\PY{p}{(}\PY{l+m+mi}{0}\PY{p}{)}\PY{p}{)}\PY{p}{;}
\PY{+w}{        }\PY{n}{assertEquals}\PY{p}{(}\PY{l+s}{\PYZdq{}}\PY{l+s}{a,1}\PY{l+s}{\PYZdq{}}\PY{p}{,}\PY{+w}{ }\PY{n}{lines}\PY{p}{.}\PY{n+na}{get}\PY{p}{(}\PY{l+m+mi}{1}\PY{p}{)}\PY{p}{)}\PY{p}{;}
\PY{+w}{    }\PY{p}{\PYZcb{}}
\PY{p}{\PYZcb{}}
\end{Verbatim}
\end{codebox}

\code{@\allowbreak{}TempDir} removes the classic “tests that write to \code{/\allowbreak{}tmp} and never clean up, eventually filling the CI disk” problem entirely.

\subsection[Testing Exceptions]{Testing Exceptions}
\label{h:20:testing-exceptions}
Always assert on both the exception \textbf{type} and enough of the \textbf{message/state} to know the right branch was hit — asserting only the type can pass even if the wrong condition threw it.

\begin{codebox}
\begin{Verbatim}[commandchars=\\\{\}]
\PY{n+nd}{@Test}
\PY{k+kt}{void}\PY{+w}{ }\PY{n+nf}{transfer\PYZus{}negativeAmount\PYZus{}throwsWithDescriptiveMessage}\PY{p}{(}\PY{p}{)}\PY{+w}{ }\PY{p}{\PYZob{}}
\PY{+w}{    }\PY{n}{Account}\PY{+w}{ }\PY{n}{source}\PY{+w}{ }\PY{o}{=}\PY{+w}{ }\PY{k}{new}\PY{+w}{ }\PY{n}{Account}\PY{p}{(}\PY{l+s}{\PYZdq{}}\PY{l+s}{A}\PY{l+s}{\PYZdq{}}\PY{p}{,}\PY{+w}{ }\PY{k}{new}\PY{+w}{ }\PY{n}{BigDecimal}\PY{p}{(}\PY{l+s}{\PYZdq{}}\PY{l+s}{100}\PY{l+s}{\PYZdq{}}\PY{p}{)}\PY{p}{)}\PY{p}{;}
\PY{+w}{    }\PY{n}{Account}\PY{+w}{ }\PY{n}{target}\PY{+w}{ }\PY{o}{=}\PY{+w}{ }\PY{k}{new}\PY{+w}{ }\PY{n}{Account}\PY{p}{(}\PY{l+s}{\PYZdq{}}\PY{l+s}{B}\PY{l+s}{\PYZdq{}}\PY{p}{,}\PY{+w}{ }\PY{k}{new}\PY{+w}{ }\PY{n}{BigDecimal}\PY{p}{(}\PY{l+s}{\PYZdq{}}\PY{l+s}{0}\PY{l+s}{\PYZdq{}}\PY{p}{)}\PY{p}{)}\PY{p}{;}

\PY{+w}{    }\PY{n}{IllegalArgumentException}\PY{+w}{ }\PY{n}{ex}\PY{+w}{ }\PY{o}{=}\PY{+w}{ }\PY{n}{assertThrows}\PY{p}{(}\PY{n}{IllegalArgumentException}\PY{p}{.}\PY{n+na}{class}\PY{p}{,}
\PY{+w}{        }\PY{p}{(}\PY{p}{)}\PY{+w}{ }\PY{o}{\PYZhy{}}\PY{o}{\PYZgt{}}\PY{+w}{ }\PY{n}{transferService}\PY{p}{.}\PY{n+na}{transfer}\PY{p}{(}\PY{n}{source}\PY{p}{,}\PY{+w}{ }\PY{n}{target}\PY{p}{,}\PY{+w}{ }\PY{k}{new}\PY{+w}{ }\PY{n}{BigDecimal}\PY{p}{(}\PY{l+s}{\PYZdq{}}\PY{l+s}{\PYZhy{}10}\PY{l+s}{\PYZdq{}}\PY{p}{)}\PY{p}{)}\PY{p}{)}\PY{p}{;}

\PY{+w}{    }\PY{n}{assertEquals}\PY{p}{(}\PY{l+s}{\PYZdq{}}\PY{l+s}{transfer amount must be positive: \PYZhy{}10}\PY{l+s}{\PYZdq{}}\PY{p}{,}\PY{+w}{ }\PY{n}{ex}\PY{p}{.}\PY{n+na}{getMessage}\PY{p}{(}\PY{p}{)}\PY{p}{)}\PY{p}{;}
\PY{p}{\PYZcb{}}
\end{Verbatim}
\end{codebox}

\subsection[Testing Time with an Injected Clock]{Testing Time with an Injected \code{Clock}}
\label{h:20:testing-time-with-an-injected-clock}
Code that calls \code{Instant.\allowbreak{}now()}, \code{LocalDate.\allowbreak{}now()}, or \code{System.\allowbreak{}current\allowbreak{}Time\allowbreak{}Millis()} directly is untestable without sleeping or mocking static methods. The fix is to depend on \code{java.\allowbreak{}time.\allowbreak{}Clock} and inject it, so tests can supply a fixed clock.

\begin{codebox}
\begin{Verbatim}[commandchars=\\\{\}]
\PY{k+kd}{public}\PY{+w}{ }\PY{k+kd}{class} \PY{n+nc}{SubscriptionService}\PY{+w}{ }\PY{p}{\PYZob{}}

\PY{+w}{    }\PY{k+kd}{private}\PY{+w}{ }\PY{k+kd}{final}\PY{+w}{ }\PY{n}{Clock}\PY{+w}{ }\PY{n}{clock}\PY{p}{;}

\PY{+w}{    }\PY{k+kd}{public}\PY{+w}{ }\PY{n+nf}{SubscriptionService}\PY{p}{(}\PY{n}{Clock}\PY{+w}{ }\PY{n}{clock}\PY{p}{)}\PY{+w}{ }\PY{p}{\PYZob{}}
\PY{+w}{        }\PY{k}{this}\PY{p}{.}\PY{n+na}{clock}\PY{+w}{ }\PY{o}{=}\PY{+w}{ }\PY{n}{clock}\PY{p}{;}
\PY{+w}{    }\PY{p}{\PYZcb{}}

\PY{+w}{    }\PY{k+kd}{public}\PY{+w}{ }\PY{k+kt}{boolean}\PY{+w}{ }\PY{n+nf}{isExpired}\PY{p}{(}\PY{n}{Subscription}\PY{+w}{ }\PY{n}{subscription}\PY{p}{)}\PY{+w}{ }\PY{p}{\PYZob{}}
\PY{+w}{        }\PY{k}{return}\PY{+w}{ }\PY{n}{subscription}\PY{p}{.}\PY{n+na}{expiresAt}\PY{p}{(}\PY{p}{)}\PY{p}{.}\PY{n+na}{isBefore}\PY{p}{(}\PY{n}{LocalDate}\PY{p}{.}\PY{n+na}{now}\PY{p}{(}\PY{n}{clock}\PY{p}{)}\PY{p}{)}\PY{p}{;}
\PY{+w}{    }\PY{p}{\PYZcb{}}
\PY{p}{\PYZcb{}}
\end{Verbatim}
\end{codebox}

\begin{codebox}
\begin{Verbatim}[commandchars=\\\{\}]
\PY{k+kd}{class} \PY{n+nc}{SubscriptionServiceTest}\PY{+w}{ }\PY{p}{\PYZob{}}

\PY{+w}{    }\PY{n+nd}{@Test}
\PY{+w}{    }\PY{k+kt}{void}\PY{+w}{ }\PY{n+nf}{isExpired\PYZus{}returnsTrue\PYZus{}whenExpiryIsInThePast}\PY{p}{(}\PY{p}{)}\PY{+w}{ }\PY{p}{\PYZob{}}
\PY{+w}{        }\PY{n}{Clock}\PY{+w}{ }\PY{n}{fixedClock}\PY{+w}{ }\PY{o}{=}\PY{+w}{ }\PY{n}{Clock}\PY{p}{.}\PY{n+na}{fixed}\PY{p}{(}
\PY{+w}{            }\PY{n}{Instant}\PY{p}{.}\PY{n+na}{parse}\PY{p}{(}\PY{l+s}{\PYZdq{}}\PY{l+s}{2026\PYZhy{}08\PYZhy{}07T00:00:00Z}\PY{l+s}{\PYZdq{}}\PY{p}{)}\PY{p}{,}\PY{+w}{ }\PY{n}{ZoneOffset}\PY{p}{.}\PY{n+na}{UTC}\PY{p}{)}\PY{p}{;}
\PY{+w}{        }\PY{n}{SubscriptionService}\PY{+w}{ }\PY{n}{service}\PY{+w}{ }\PY{o}{=}\PY{+w}{ }\PY{k}{new}\PY{+w}{ }\PY{n}{SubscriptionService}\PY{p}{(}\PY{n}{fixedClock}\PY{p}{)}\PY{p}{;}

\PY{+w}{        }\PY{n}{Subscription}\PY{+w}{ }\PY{n}{expired}\PY{+w}{ }\PY{o}{=}\PY{+w}{ }\PY{k}{new}\PY{+w}{ }\PY{n}{Subscription}\PY{p}{(}\PY{n}{LocalDate}\PY{p}{.}\PY{n+na}{parse}\PY{p}{(}\PY{l+s}{\PYZdq{}}\PY{l+s}{2026\PYZhy{}08\PYZhy{}01}\PY{l+s}{\PYZdq{}}\PY{p}{)}\PY{p}{)}\PY{p}{;}

\PY{+w}{        }\PY{n}{assertTrue}\PY{p}{(}\PY{n}{service}\PY{p}{.}\PY{n+na}{isExpired}\PY{p}{(}\PY{n}{expired}\PY{p}{)}\PY{p}{)}\PY{p}{;}
\PY{+w}{    }\PY{p}{\PYZcb{}}

\PY{+w}{    }\PY{n+nd}{@Test}
\PY{+w}{    }\PY{k+kt}{void}\PY{+w}{ }\PY{n+nf}{isExpired\PYZus{}returnsFalse\PYZus{}whenExpiryIsInTheFuture}\PY{p}{(}\PY{p}{)}\PY{+w}{ }\PY{p}{\PYZob{}}
\PY{+w}{        }\PY{n}{Clock}\PY{+w}{ }\PY{n}{fixedClock}\PY{+w}{ }\PY{o}{=}\PY{+w}{ }\PY{n}{Clock}\PY{p}{.}\PY{n+na}{fixed}\PY{p}{(}
\PY{+w}{            }\PY{n}{Instant}\PY{p}{.}\PY{n+na}{parse}\PY{p}{(}\PY{l+s}{\PYZdq{}}\PY{l+s}{2026\PYZhy{}08\PYZhy{}07T00:00:00Z}\PY{l+s}{\PYZdq{}}\PY{p}{)}\PY{p}{,}\PY{+w}{ }\PY{n}{ZoneOffset}\PY{p}{.}\PY{n+na}{UTC}\PY{p}{)}\PY{p}{;}
\PY{+w}{        }\PY{n}{SubscriptionService}\PY{+w}{ }\PY{n}{service}\PY{+w}{ }\PY{o}{=}\PY{+w}{ }\PY{k}{new}\PY{+w}{ }\PY{n}{SubscriptionService}\PY{p}{(}\PY{n}{fixedClock}\PY{p}{)}\PY{p}{;}

\PY{+w}{        }\PY{n}{Subscription}\PY{+w}{ }\PY{n}{active}\PY{+w}{ }\PY{o}{=}\PY{+w}{ }\PY{k}{new}\PY{+w}{ }\PY{n}{Subscription}\PY{p}{(}\PY{n}{LocalDate}\PY{p}{.}\PY{n+na}{parse}\PY{p}{(}\PY{l+s}{\PYZdq{}}\PY{l+s}{2026\PYZhy{}12\PYZhy{}31}\PY{l+s}{\PYZdq{}}\PY{p}{)}\PY{p}{)}\PY{p}{;}

\PY{+w}{        }\PY{n}{assertFalse}\PY{p}{(}\PY{n}{service}\PY{p}{.}\PY{n+na}{isExpired}\PY{p}{(}\PY{n}{active}\PY{p}{)}\PY{p}{)}\PY{p}{;}
\PY{+w}{    }\PY{p}{\PYZcb{}}
\PY{p}{\PYZcb{}}
\end{Verbatim}
\end{codebox}

This one pattern eliminates an entire category of flaky, date-dependent tests (the ones that mysteriously fail once a year around New Year’s Eve, or that pass locally but fail on CI in a different timezone).

\subsection[Testing Concurrency — and Why Thread.sleep Is a Smell]{Testing Concurrency — and Why \code{Thread.\allowbreak{}sleep} Is a Smell}
\label{h:20:testing-concurrency--and-why-threadsleep-is-a-smell}
A test that does this is a ticking time bomb:

\begin{codebox}
\begin{Verbatim}[commandchars=\\\{\}]
\PY{c+c1}{// Bad: racy and slow}
\PY{n+nd}{@Test}
\PY{k+kt}{void}\PY{+w}{ }\PY{n+nf}{taskRunsAsynchronously}\PY{p}{(}\PY{p}{)}\PY{+w}{ }\PY{k+kd}{throws}\PY{+w}{ }\PY{n}{InterruptedException}\PY{+w}{ }\PY{p}{\PYZob{}}
\PY{+w}{    }\PY{n}{AtomicBoolean}\PY{+w}{ }\PY{n}{ran}\PY{+w}{ }\PY{o}{=}\PY{+w}{ }\PY{k}{new}\PY{+w}{ }\PY{n}{AtomicBoolean}\PY{p}{(}\PY{k+kc}{false}\PY{p}{)}\PY{p}{;}
\PY{+w}{    }\PY{n}{executor}\PY{p}{.}\PY{n+na}{submit}\PY{p}{(}\PY{p}{(}\PY{p}{)}\PY{+w}{ }\PY{o}{\PYZhy{}}\PY{o}{\PYZgt{}}\PY{+w}{ }\PY{n}{ran}\PY{p}{.}\PY{n+na}{set}\PY{p}{(}\PY{k+kc}{true}\PY{p}{)}\PY{p}{)}\PY{p}{;}

\PY{+w}{    }\PY{n}{Thread}\PY{p}{.}\PY{n+na}{sleep}\PY{p}{(}\PY{l+m+mi}{500}\PY{p}{)}\PY{p}{;}\PY{+w}{ }\PY{c+c1}{// hope 500ms is enough... until CI is under load}

\PY{+w}{    }\PY{n}{assertTrue}\PY{p}{(}\PY{n}{ran}\PY{p}{.}\PY{n+na}{get}\PY{p}{(}\PY{p}{)}\PY{p}{)}\PY{p}{;}
\PY{p}{\PYZcb{}}
\end{Verbatim}
\end{codebox}

\code{Thread.\allowbreak{}sleep} in a concurrency test encodes a guess about timing. It is either wasteful (sleeping far longer than necessary, slowing down the whole suite) or unreliable (too short under CI load, causing intermittent failures — flakiness). The fix is to wait for the actual event, not for a fixed amount of wall-clock time.

\textbf{\code{CountDownLatch}} — signal exactly when the work is done:

\begin{codebox}
\begin{Verbatim}[commandchars=\\\{\}]
\PY{n+nd}{@Test}
\PY{k+kt}{void}\PY{+w}{ }\PY{n+nf}{taskRunsAsynchronously\PYZus{}withLatch}\PY{p}{(}\PY{p}{)}\PY{+w}{ }\PY{k+kd}{throws}\PY{+w}{ }\PY{n}{InterruptedException}\PY{+w}{ }\PY{p}{\PYZob{}}
\PY{+w}{    }\PY{n}{CountDownLatch}\PY{+w}{ }\PY{n}{latch}\PY{+w}{ }\PY{o}{=}\PY{+w}{ }\PY{k}{new}\PY{+w}{ }\PY{n}{CountDownLatch}\PY{p}{(}\PY{l+m+mi}{1}\PY{p}{)}\PY{p}{;}
\PY{+w}{    }\PY{n}{AtomicBoolean}\PY{+w}{ }\PY{n}{ran}\PY{+w}{ }\PY{o}{=}\PY{+w}{ }\PY{k}{new}\PY{+w}{ }\PY{n}{AtomicBoolean}\PY{p}{(}\PY{k+kc}{false}\PY{p}{)}\PY{p}{;}

\PY{+w}{    }\PY{n}{executor}\PY{p}{.}\PY{n+na}{submit}\PY{p}{(}\PY{p}{(}\PY{p}{)}\PY{+w}{ }\PY{o}{\PYZhy{}}\PY{o}{\PYZgt{}}\PY{+w}{ }\PY{p}{\PYZob{}}
\PY{+w}{        }\PY{n}{ran}\PY{p}{.}\PY{n+na}{set}\PY{p}{(}\PY{k+kc}{true}\PY{p}{)}\PY{p}{;}
\PY{+w}{        }\PY{n}{latch}\PY{p}{.}\PY{n+na}{countDown}\PY{p}{(}\PY{p}{)}\PY{p}{;}
\PY{+w}{    }\PY{p}{\PYZcb{}}\PY{p}{)}\PY{p}{;}

\PY{+w}{    }\PY{k+kt}{boolean}\PY{+w}{ }\PY{n}{completedInTime}\PY{+w}{ }\PY{o}{=}\PY{+w}{ }\PY{n}{latch}\PY{p}{.}\PY{n+na}{await}\PY{p}{(}\PY{l+m+mi}{2}\PY{p}{,}\PY{+w}{ }\PY{n}{TimeUnit}\PY{p}{.}\PY{n+na}{SECONDS}\PY{p}{)}\PY{p}{;}

\PY{+w}{    }\PY{n}{assertTrue}\PY{p}{(}\PY{n}{completedInTime}\PY{p}{,}\PY{+w}{ }\PY{l+s}{\PYZdq{}}\PY{l+s}{task did not complete within timeout}\PY{l+s}{\PYZdq{}}\PY{p}{)}\PY{p}{;}
\PY{+w}{    }\PY{n}{assertTrue}\PY{p}{(}\PY{n}{ran}\PY{p}{.}\PY{n+na}{get}\PY{p}{(}\PY{p}{)}\PY{p}{)}\PY{p}{;}
\PY{p}{\PYZcb{}}
\end{Verbatim}
\end{codebox}

\textbf{Awaitility} — for polling a condition that eventually becomes true (e.g., an async cache being populated), which reads far better than a hand-rolled polling loop:

\begin{codebox}
\begin{Verbatim}[commandchars=\\\{\}]
\PY{k+kn}{import static}\PY{+w}{ }\PY{n+nn}{org.awaitility.Awaitility.await}\PY{p}{;}
\PY{k+kn}{import}\PY{+w}{ }\PY{n+nn}{java.time.Duration}\PY{p}{;}

\PY{n+nd}{@Test}
\PY{k+kt}{void}\PY{+w}{ }\PY{n+nf}{cache\PYZus{}eventuallyContainsWarmedEntries}\PY{p}{(}\PY{p}{)}\PY{+w}{ }\PY{p}{\PYZob{}}
\PY{+w}{    }\PY{n}{cacheWarmer}\PY{p}{.}\PY{n+na}{warmAsync}\PY{p}{(}\PY{p}{)}\PY{p}{;}

\PY{+w}{    }\PY{n}{await}\PY{p}{(}\PY{p}{)}\PY{p}{.}\PY{n+na}{atMost}\PY{p}{(}\PY{n}{Duration}\PY{p}{.}\PY{n+na}{ofSeconds}\PY{p}{(}\PY{l+m+mi}{3}\PY{p}{)}\PY{p}{)}
\PY{+w}{           }\PY{p}{.}\PY{n+na}{pollInterval}\PY{p}{(}\PY{n}{Duration}\PY{p}{.}\PY{n+na}{ofMillis}\PY{p}{(}\PY{l+m+mi}{50}\PY{p}{)}\PY{p}{)}
\PY{+w}{           }\PY{p}{.}\PY{n+na}{until}\PY{p}{(}\PY{p}{(}\PY{p}{)}\PY{+w}{ }\PY{o}{\PYZhy{}}\PY{o}{\PYZgt{}}\PY{+w}{ }\PY{n}{cache}\PY{p}{.}\PY{n+na}{get}\PY{p}{(}\PY{l+s}{\PYZdq{}}\PY{l+s}{key}\PY{l+s}{\PYZdq{}}\PY{p}{)}\PY{+w}{ }\PY{o}{!}\PY{o}{=}\PY{+w}{ }\PY{k+kc}{null}\PY{p}{)}\PY{p}{;}

\PY{+w}{    }\PY{n}{assertEquals}\PY{p}{(}\PY{l+s}{\PYZdq{}}\PY{l+s}{value}\PY{l+s}{\PYZdq{}}\PY{p}{,}\PY{+w}{ }\PY{n}{cache}\PY{p}{.}\PY{n+na}{get}\PY{p}{(}\PY{l+s}{\PYZdq{}}\PY{l+s}{key}\PY{l+s}{\PYZdq{}}\PY{p}{)}\PY{p}{)}\PY{p}{;}
\PY{p}{\PYZcb{}}
\end{Verbatim}
\end{codebox}

Both approaches wait for the \emph{minimum} necessary time and fail fast with a clear timeout message, instead of guessing a sleep duration that is either too short (flaky) or too long (slow suite).

\subsection[Test Doubles Vocabulary]{Test Doubles Vocabulary}
\label{h:20:test-doubles-vocabulary}
“Mock” is used loosely in everyday speech, but in precise testing vocabulary there are five distinct kinds of test doubles:

{\tablefontsmall
\begin{xltabular}{\linewidth}{@{}L{0.413}L{1.868}L{0.700}L{1.019}@{}}
\toprule
\rowcolor{tablehdr}
\thd{Term} & \thd{What it does} & \thd{Verifies behavior?} & \thd{Typical use} \\
\midrule
\endhead
\textbf{Dummy} & A placeholder passed to satisfy a parameter list; never actually used & No & \code{new~\allowbreak{}Service(\allowbreak{}null,\allowbreak{}~\allowbreak{}dummy\allowbreak{}Logger)} when \code{dummyLogger} is never called \\
\textbf{Stub} & Returns canned answers to calls made during the test & No & \code{when(\allowbreak{}repo.\allowbreak{}find\allowbreak{}By\allowbreak{}Id(\allowbreak{}1)).\allowbreak{}then\allowbreak{}Return(\allowbreak{}user)} \\
\textbf{Fake} & A lightweight \emph{working} implementation, not production-grade (e.g., in-memory DB) & Indirectly, through real behavior & \code{In\allowbreak{}Memory\allowbreak{}User\allowbreak{}Repository~\allowbreak{}implements~\allowbreak{}User\allowbreak{}Repository} \\
\textbf{Spy} & A real object with some methods overridden/recorded, or a mock that records calls on a real instance & Yes (records calls, can also stub) & \code{Mockito.\allowbreak{}spy(\allowbreak{}real\allowbreak{}Service)} \\
\textbf{Mock} & A dynamically generated object programmed with expectations, used to verify interactions & Yes (that’s its purpose) & \code{verify(\allowbreak{}email\allowbreak{}Sender).\allowbreak{}send(\allowbreak{}any())} \\
\bottomrule
\end{xltabular}
}

The line between stub and mock is about \textbf{intent}: a stub just feeds data into the test (state verification — you check the \emph{result}), while a mock is used to verify that specific interactions happened (behavior verification — you check the \emph{calls}). Mockito’s \code{mock()} can be used as either, depending on whether you call \code{verify(...)} on it.

\subsection[Code Coverage — and Why 100\% Is Not the Goal]{Code Coverage — and Why 100\% Is Not the Goal}
\label{h:20:code-coverage--and-why-100-is-not-the-goal}
Coverage tools (JaCoCo, for example) report the percentage of lines/branches executed by the test suite. Coverage is a \textbf{necessary but not sufficient} signal:

\begin{itemize}
\item 
Low coverage on business-critical code is a real red flag — it means there is a good chance that logic has never been exercised by a test.

\item 
High coverage tells you code \emph{ran}, not that it was \emph{checked}. A test with no assertions can hit 100\% of a method’s lines and prove nothing.

\end{itemize}

\begin{codebox}
\begin{Verbatim}[commandchars=\\\{\}]
\PY{c+c1}{// Executes every line (100\PYZpc{} \PYZdq{}covered\PYZdq{}) but asserts nothing meaningful}
\PY{n+nd}{@Test}
\PY{k+kt}{void}\PY{+w}{ }\PY{n+nf}{processOrder\PYZus{}doesNotThrow}\PY{p}{(}\PY{p}{)}\PY{+w}{ }\PY{p}{\PYZob{}}
\PY{+w}{    }\PY{n}{orderService}\PY{p}{.}\PY{n+na}{process}\PY{p}{(}\PY{n}{order}\PY{p}{)}\PY{p}{;}\PY{+w}{ }\PY{c+c1}{// \PYZdq{}coverage\PYZdq{} achieved, but no behavior verified}
\PY{p}{\PYZcb{}}
\end{Verbatim}
\end{codebox}

Chasing 100\% coverage as a target also incentivizes testing trivial code (getters, \code{toString()}, generated code) at the expense of time that should go to edge cases and error paths. A better review question than “what’s the coverage number?” is “does this test actually assert the behavior that matters, including the failure paths?”

\subsection[Flaky Tests]{Flaky Tests}
\label{h:20:flaky-tests}
A \textbf{flaky test} passes and fails intermittently with no code change — the single most corrosive thing to trust in a test suite, because a team that has learned to click “rerun” on red builds has effectively stopped trusting CI.

Common causes:

\begin{itemize}
\item 
\textbf{Time dependence} — using \code{LocalDate.\allowbreak{}now()}/\code{Instant.\allowbreak{}now()} directly instead of an injected \code{Clock} (see above).

\item 
\textbf{Shared mutable state} — static fields, singletons, or a shared database row mutated by tests running in parallel or in a different order than expected.

\item 
\textbf{Real concurrency/timing assumptions} — \code{Thread.\allowbreak{}sleep} races, or asserting order between two async operations without synchronization.

\item 
\textbf{External dependencies} — hitting a real network service, real clock-based expiry, or an unseeded random number generator.

\item 
\textbf{Test order dependence} — a test that only passes because an earlier test happened to leave the system in the right state; JUnit does not guarantee method execution order, so this can flip at any time.

\item 
\textbf{Resource leaks across tests} — a thread pool, file handle, or connection left open by one test that starves the next one.

\end{itemize}

The fix is almost always to remove the non-determinism (inject a clock, replace \code{sleep} with a latch, isolate state per test, seed random generators) rather than to add retries — retries hide the bug, they don’t fix it.

\subsection[A Full Realistic Test Class]{A Full Realistic Test Class}
\label{h:20:a-full-realistic-test-class}
Putting several of the above together for a small service that computes shipping fees:

\begin{codebox}
\begin{Verbatim}[commandchars=\\\{\}]
\PY{c+c1}{// Production code}
\PY{k+kd}{public}\PY{+w}{ }\PY{k+kd}{class} \PY{n+nc}{ShippingFeeCalculator}\PY{+w}{ }\PY{p}{\PYZob{}}

\PY{+w}{    }\PY{k+kd}{private}\PY{+w}{ }\PY{k+kd}{static}\PY{+w}{ }\PY{k+kd}{final}\PY{+w}{ }\PY{n}{BigDecimal}\PY{+w}{ }\PY{n}{FREE\PYZus{}SHIPPING\PYZus{}THRESHOLD}\PY{+w}{ }\PY{o}{=}\PY{+w}{ }\PY{k}{new}\PY{+w}{ }\PY{n}{BigDecimal}\PY{p}{(}\PY{l+s}{\PYZdq{}}\PY{l+s}{50.00}\PY{l+s}{\PYZdq{}}\PY{p}{)}\PY{p}{;}
\PY{+w}{    }\PY{k+kd}{private}\PY{+w}{ }\PY{k+kd}{static}\PY{+w}{ }\PY{k+kd}{final}\PY{+w}{ }\PY{n}{BigDecimal}\PY{+w}{ }\PY{n}{STANDARD\PYZus{}FEE}\PY{+w}{ }\PY{o}{=}\PY{+w}{ }\PY{k}{new}\PY{+w}{ }\PY{n}{BigDecimal}\PY{p}{(}\PY{l+s}{\PYZdq{}}\PY{l+s}{4.99}\PY{l+s}{\PYZdq{}}\PY{p}{)}\PY{p}{;}

\PY{+w}{    }\PY{k+kd}{public}\PY{+w}{ }\PY{n}{BigDecimal}\PY{+w}{ }\PY{n+nf}{feeFor}\PY{p}{(}\PY{n}{BigDecimal}\PY{+w}{ }\PY{n}{orderTotal}\PY{p}{,}\PY{+w}{ }\PY{n}{ShippingSpeed}\PY{+w}{ }\PY{n}{speed}\PY{p}{)}\PY{+w}{ }\PY{p}{\PYZob{}}
\PY{+w}{        }\PY{k}{if}\PY{+w}{ }\PY{p}{(}\PY{n}{orderTotal}\PY{+w}{ }\PY{o}{=}\PY{o}{=}\PY{+w}{ }\PY{k+kc}{null}\PY{p}{)}\PY{+w}{ }\PY{p}{\PYZob{}}
\PY{+w}{            }\PY{k}{throw}\PY{+w}{ }\PY{k}{new}\PY{+w}{ }\PY{n}{IllegalArgumentException}\PY{p}{(}\PY{l+s}{\PYZdq{}}\PY{l+s}{orderTotal must not be null}\PY{l+s}{\PYZdq{}}\PY{p}{)}\PY{p}{;}
\PY{+w}{        }\PY{p}{\PYZcb{}}
\PY{+w}{        }\PY{k}{if}\PY{+w}{ }\PY{p}{(}\PY{n}{orderTotal}\PY{p}{.}\PY{n+na}{compareTo}\PY{p}{(}\PY{n}{BigDecimal}\PY{p}{.}\PY{n+na}{ZERO}\PY{p}{)}\PY{+w}{ }\PY{o}{\PYZlt{}}\PY{+w}{ }\PY{l+m+mi}{0}\PY{p}{)}\PY{+w}{ }\PY{p}{\PYZob{}}
\PY{+w}{            }\PY{k}{throw}\PY{+w}{ }\PY{k}{new}\PY{+w}{ }\PY{n}{IllegalArgumentException}\PY{p}{(}\PY{l+s}{\PYZdq{}}\PY{l+s}{orderTotal must not be negative: }\PY{l+s}{\PYZdq{}}\PY{+w}{ }\PY{o}{+}\PY{+w}{ }\PY{n}{orderTotal}\PY{p}{)}\PY{p}{;}
\PY{+w}{        }\PY{p}{\PYZcb{}}

\PY{+w}{        }\PY{k}{if}\PY{+w}{ }\PY{p}{(}\PY{n}{orderTotal}\PY{p}{.}\PY{n+na}{compareTo}\PY{p}{(}\PY{n}{FREE\PYZus{}SHIPPING\PYZus{}THRESHOLD}\PY{p}{)}\PY{+w}{ }\PY{o}{\PYZgt{}}\PY{o}{=}\PY{+w}{ }\PY{l+m+mi}{0}\PY{+w}{ }\PY{o}{\PYZam{}}\PY{o}{\PYZam{}}\PY{+w}{ }\PY{n}{speed}\PY{+w}{ }\PY{o}{=}\PY{o}{=}\PY{+w}{ }\PY{n}{ShippingSpeed}\PY{p}{.}\PY{n+na}{STANDARD}\PY{p}{)}\PY{+w}{ }\PY{p}{\PYZob{}}
\PY{+w}{            }\PY{k}{return}\PY{+w}{ }\PY{n}{BigDecimal}\PY{p}{.}\PY{n+na}{ZERO}\PY{p}{;}
\PY{+w}{        }\PY{p}{\PYZcb{}}

\PY{+w}{        }\PY{k}{return}\PY{+w}{ }\PY{k}{switch}\PY{+w}{ }\PY{p}{(}\PY{n}{speed}\PY{p}{)}\PY{+w}{ }\PY{p}{\PYZob{}}
\PY{+w}{            }\PY{k}{case}\PY{+w}{ }\PY{n}{STANDARD}\PY{+w}{ }\PY{o}{\PYZhy{}}\PY{o}{\PYZgt{}}\PY{+w}{ }\PY{n}{STANDARD\PYZus{}FEE}\PY{p}{;}
\PY{+w}{            }\PY{k}{case}\PY{+w}{ }\PY{n}{EXPRESS}\PY{+w}{ }\PY{o}{\PYZhy{}}\PY{o}{\PYZgt{}}\PY{+w}{ }\PY{n}{STANDARD\PYZus{}FEE}\PY{p}{.}\PY{n+na}{multiply}\PY{p}{(}\PY{k}{new}\PY{+w}{ }\PY{n}{BigDecimal}\PY{p}{(}\PY{l+s}{\PYZdq{}}\PY{l+s}{2}\PY{l+s}{\PYZdq{}}\PY{p}{)}\PY{p}{)}\PY{p}{;}
\PY{+w}{            }\PY{k}{case}\PY{+w}{ }\PY{n}{OVERNIGHT}\PY{+w}{ }\PY{o}{\PYZhy{}}\PY{o}{\PYZgt{}}\PY{+w}{ }\PY{n}{STANDARD\PYZus{}FEE}\PY{p}{.}\PY{n+na}{multiply}\PY{p}{(}\PY{k}{new}\PY{+w}{ }\PY{n}{BigDecimal}\PY{p}{(}\PY{l+s}{\PYZdq{}}\PY{l+s}{4}\PY{l+s}{\PYZdq{}}\PY{p}{)}\PY{p}{)}\PY{p}{;}
\PY{+w}{        }\PY{p}{\PYZcb{}}\PY{p}{;}
\PY{+w}{    }\PY{p}{\PYZcb{}}
\PY{p}{\PYZcb{}}
\end{Verbatim}
\end{codebox}

\begin{codebox}
\begin{Verbatim}[commandchars=\\\{\}]
\PY{k+kn}{import}\PY{+w}{ }\PY{n+nn}{org.junit.jupiter.api.*}\PY{p}{;}
\PY{k+kn}{import}\PY{+w}{ }\PY{n+nn}{org.junit.jupiter.params.ParameterizedTest}\PY{p}{;}
\PY{k+kn}{import}\PY{+w}{ }\PY{n+nn}{org.junit.jupiter.params.provider.CsvSource}\PY{p}{;}
\PY{k+kn}{import}\PY{+w}{ }\PY{n+nn}{java.math.BigDecimal}\PY{p}{;}

\PY{k+kn}{import static}\PY{+w}{ }\PY{n+nn}{org.junit.jupiter.api.Assertions.*}\PY{p}{;}

\PY{n+nd}{@DisplayName}\PY{p}{(}\PY{l+s}{\PYZdq{}}\PY{l+s}{ShippingFeeCalculator}\PY{l+s}{\PYZdq{}}\PY{p}{)}
\PY{k+kd}{class} \PY{n+nc}{ShippingFeeCalculatorTest}\PY{+w}{ }\PY{p}{\PYZob{}}

\PY{+w}{    }\PY{k+kd}{private}\PY{+w}{ }\PY{k+kd}{final}\PY{+w}{ }\PY{n}{ShippingFeeCalculator}\PY{+w}{ }\PY{n}{calculator}\PY{+w}{ }\PY{o}{=}\PY{+w}{ }\PY{k}{new}\PY{+w}{ }\PY{n}{ShippingFeeCalculator}\PY{p}{(}\PY{p}{)}\PY{p}{;}

\PY{+w}{    }\PY{n+nd}{@Nested}
\PY{+w}{    }\PY{n+nd}{@DisplayName}\PY{p}{(}\PY{l+s}{\PYZdq{}}\PY{l+s}{input validation}\PY{l+s}{\PYZdq{}}\PY{p}{)}
\PY{+w}{    }\PY{k+kd}{class} \PY{n+nc}{InputValidation}\PY{+w}{ }\PY{p}{\PYZob{}}

\PY{+w}{        }\PY{n+nd}{@Test}
\PY{+w}{        }\PY{n+nd}{@DisplayName}\PY{p}{(}\PY{l+s}{\PYZdq{}}\PY{l+s}{throws when orderTotal is null}\PY{l+s}{\PYZdq{}}\PY{p}{)}
\PY{+w}{        }\PY{k+kt}{void}\PY{+w}{ }\PY{n+nf}{feeFor\PYZus{}nullOrderTotal\PYZus{}throws}\PY{p}{(}\PY{p}{)}\PY{+w}{ }\PY{p}{\PYZob{}}
\PY{+w}{            }\PY{n}{IllegalArgumentException}\PY{+w}{ }\PY{n}{ex}\PY{+w}{ }\PY{o}{=}\PY{+w}{ }\PY{n}{assertThrows}\PY{p}{(}\PY{n}{IllegalArgumentException}\PY{p}{.}\PY{n+na}{class}\PY{p}{,}
\PY{+w}{                }\PY{p}{(}\PY{p}{)}\PY{+w}{ }\PY{o}{\PYZhy{}}\PY{o}{\PYZgt{}}\PY{+w}{ }\PY{n}{calculator}\PY{p}{.}\PY{n+na}{feeFor}\PY{p}{(}\PY{k+kc}{null}\PY{p}{,}\PY{+w}{ }\PY{n}{ShippingSpeed}\PY{p}{.}\PY{n+na}{STANDARD}\PY{p}{)}\PY{p}{)}\PY{p}{;}
\PY{+w}{            }\PY{n}{assertEquals}\PY{p}{(}\PY{l+s}{\PYZdq{}}\PY{l+s}{orderTotal must not be null}\PY{l+s}{\PYZdq{}}\PY{p}{,}\PY{+w}{ }\PY{n}{ex}\PY{p}{.}\PY{n+na}{getMessage}\PY{p}{(}\PY{p}{)}\PY{p}{)}\PY{p}{;}
\PY{+w}{        }\PY{p}{\PYZcb{}}

\PY{+w}{        }\PY{n+nd}{@Test}
\PY{+w}{        }\PY{n+nd}{@DisplayName}\PY{p}{(}\PY{l+s}{\PYZdq{}}\PY{l+s}{throws when orderTotal is negative}\PY{l+s}{\PYZdq{}}\PY{p}{)}
\PY{+w}{        }\PY{k+kt}{void}\PY{+w}{ }\PY{n+nf}{feeFor\PYZus{}negativeOrderTotal\PYZus{}throws}\PY{p}{(}\PY{p}{)}\PY{+w}{ }\PY{p}{\PYZob{}}
\PY{+w}{            }\PY{n}{assertThrows}\PY{p}{(}\PY{n}{IllegalArgumentException}\PY{p}{.}\PY{n+na}{class}\PY{p}{,}
\PY{+w}{                }\PY{p}{(}\PY{p}{)}\PY{+w}{ }\PY{o}{\PYZhy{}}\PY{o}{\PYZgt{}}\PY{+w}{ }\PY{n}{calculator}\PY{p}{.}\PY{n+na}{feeFor}\PY{p}{(}\PY{k}{new}\PY{+w}{ }\PY{n}{BigDecimal}\PY{p}{(}\PY{l+s}{\PYZdq{}}\PY{l+s}{\PYZhy{}1.00}\PY{l+s}{\PYZdq{}}\PY{p}{)}\PY{p}{,}\PY{+w}{ }\PY{n}{ShippingSpeed}\PY{p}{.}\PY{n+na}{STANDARD}\PY{p}{)}\PY{p}{)}\PY{p}{;}
\PY{+w}{        }\PY{p}{\PYZcb{}}
\PY{+w}{    }\PY{p}{\PYZcb{}}

\PY{+w}{    }\PY{n+nd}{@Nested}
\PY{+w}{    }\PY{n+nd}{@DisplayName}\PY{p}{(}\PY{l+s}{\PYZdq{}}\PY{l+s}{standard shipping}\PY{l+s}{\PYZdq{}}\PY{p}{)}
\PY{+w}{    }\PY{k+kd}{class} \PY{n+nc}{StandardShipping}\PY{+w}{ }\PY{p}{\PYZob{}}

\PY{+w}{        }\PY{n+nd}{@Test}
\PY{+w}{        }\PY{n+nd}{@DisplayName}\PY{p}{(}\PY{l+s}{\PYZdq{}}\PY{l+s}{is free at or above the free\PYZhy{}shipping threshold}\PY{l+s}{\PYZdq{}}\PY{p}{)}
\PY{+w}{        }\PY{k+kt}{void}\PY{+w}{ }\PY{n+nf}{feeFor\PYZus{}atThreshold\PYZus{}isFree}\PY{p}{(}\PY{p}{)}\PY{+w}{ }\PY{p}{\PYZob{}}
\PY{+w}{            }\PY{n}{BigDecimal}\PY{+w}{ }\PY{n}{fee}\PY{+w}{ }\PY{o}{=}\PY{+w}{ }\PY{n}{calculator}\PY{p}{.}\PY{n+na}{feeFor}\PY{p}{(}\PY{k}{new}\PY{+w}{ }\PY{n}{BigDecimal}\PY{p}{(}\PY{l+s}{\PYZdq{}}\PY{l+s}{50.00}\PY{l+s}{\PYZdq{}}\PY{p}{)}\PY{p}{,}\PY{+w}{ }\PY{n}{ShippingSpeed}\PY{p}{.}\PY{n+na}{STANDARD}\PY{p}{)}\PY{p}{;}
\PY{+w}{            }\PY{n}{assertEquals}\PY{p}{(}\PY{n}{BigDecimal}\PY{p}{.}\PY{n+na}{ZERO}\PY{p}{,}\PY{+w}{ }\PY{n}{fee}\PY{p}{)}\PY{p}{;}
\PY{+w}{        }\PY{p}{\PYZcb{}}

\PY{+w}{        }\PY{n+nd}{@Test}
\PY{+w}{        }\PY{n+nd}{@DisplayName}\PY{p}{(}\PY{l+s}{\PYZdq{}}\PY{l+s}{charges the standard fee below the threshold}\PY{l+s}{\PYZdq{}}\PY{p}{)}
\PY{+w}{        }\PY{k+kt}{void}\PY{+w}{ }\PY{n+nf}{feeFor\PYZus{}belowThreshold\PYZus{}chargesStandardFee}\PY{p}{(}\PY{p}{)}\PY{+w}{ }\PY{p}{\PYZob{}}
\PY{+w}{            }\PY{n}{BigDecimal}\PY{+w}{ }\PY{n}{fee}\PY{+w}{ }\PY{o}{=}\PY{+w}{ }\PY{n}{calculator}\PY{p}{.}\PY{n+na}{feeFor}\PY{p}{(}\PY{k}{new}\PY{+w}{ }\PY{n}{BigDecimal}\PY{p}{(}\PY{l+s}{\PYZdq{}}\PY{l+s}{49.99}\PY{l+s}{\PYZdq{}}\PY{p}{)}\PY{p}{,}\PY{+w}{ }\PY{n}{ShippingSpeed}\PY{p}{.}\PY{n+na}{STANDARD}\PY{p}{)}\PY{p}{;}
\PY{+w}{            }\PY{n}{assertEquals}\PY{p}{(}\PY{k}{new}\PY{+w}{ }\PY{n}{BigDecimal}\PY{p}{(}\PY{l+s}{\PYZdq{}}\PY{l+s}{4.99}\PY{l+s}{\PYZdq{}}\PY{p}{)}\PY{p}{,}\PY{+w}{ }\PY{n}{fee}\PY{p}{)}\PY{p}{;}
\PY{+w}{        }\PY{p}{\PYZcb{}}
\PY{+w}{    }\PY{p}{\PYZcb{}}

\PY{+w}{    }\PY{n+nd}{@ParameterizedTest}\PY{p}{(}\PY{n}{name}\PY{+w}{ }\PY{o}{=}\PY{+w}{ }\PY{l+s}{\PYZdq{}}\PY{l+s}{\PYZob{}0\PYZcb{} order with \PYZob{}1\PYZcb{} shipping costs \PYZob{}2\PYZcb{}}\PY{l+s}{\PYZdq{}}\PY{p}{)}
\PY{+w}{    }\PY{n+nd}{@CsvSource}\PY{p}{(}\PY{p}{\PYZob{}}
\PY{+w}{        }\PY{l+s}{\PYZdq{}}\PY{l+s}{10.00, EXPRESS,   9.98}\PY{l+s}{\PYZdq{}}\PY{p}{,}
\PY{+w}{        }\PY{l+s}{\PYZdq{}}\PY{l+s}{10.00, OVERNIGHT, 19.96}\PY{l+s}{\PYZdq{}}\PY{p}{,}
\PY{+w}{        }\PY{l+s}{\PYZdq{}}\PY{l+s}{60.00, EXPRESS,   9.98}\PY{l+s}{\PYZdq{}}\PY{p}{,}\PY{+w}{   }\PY{c+c1}{// free\PYZhy{}shipping discount only applies to STANDARD}
\PY{+w}{        }\PY{l+s}{\PYZdq{}}\PY{l+s}{60.00, OVERNIGHT, 19.96}\PY{l+s}{\PYZdq{}}
\PY{+w}{    }\PY{p}{\PYZcb{}}\PY{p}{)}
\PY{+w}{    }\PY{n+nd}{@DisplayName}\PY{p}{(}\PY{l+s}{\PYZdq{}}\PY{l+s}{express and overnight fees ignore the free\PYZhy{}shipping threshold}\PY{l+s}{\PYZdq{}}\PY{p}{)}
\PY{+w}{    }\PY{k+kt}{void}\PY{+w}{ }\PY{n+nf}{feeFor\PYZus{}expressAndOvernight\PYZus{}neverFree}\PY{p}{(}\PY{n}{BigDecimal}\PY{+w}{ }\PY{n}{total}\PY{p}{,}\PY{+w}{ }\PY{n}{ShippingSpeed}\PY{+w}{ }\PY{n}{speed}\PY{p}{,}\PY{+w}{ }\PY{n}{BigDecimal}\PY{+w}{ }\PY{n}{expectedFee}\PY{p}{)}\PY{+w}{ }\PY{p}{\PYZob{}}
\PY{+w}{        }\PY{n}{assertEquals}\PY{p}{(}\PY{n}{expectedFee}\PY{p}{,}\PY{+w}{ }\PY{n}{calculator}\PY{p}{.}\PY{n+na}{feeFor}\PY{p}{(}\PY{n}{total}\PY{p}{,}\PY{+w}{ }\PY{n}{speed}\PY{p}{)}\PY{p}{)}\PY{p}{;}
\PY{+w}{    }\PY{p}{\PYZcb{}}
\PY{p}{\PYZcb{}}
\end{Verbatim}
\end{codebox}

This class demonstrates: \code{@\allowbreak{}Nested} grouping by scenario, \code{@\allowbreak{}DisplayName} for readable reports, a parameterized test for the combinatorial cases, and assertions that check both the \emph{type} and the \emph{content} of thrown exceptions.

\section[Mocking (Mockito)]{Mocking (Mockito)}
\label{h:20:mocking-mockito}
\subsection[What to Mock and What NOT to Mock]{What to Mock and What NOT to Mock}
\label{h:20:what-to-mock-and-what-not-to-mock}
Mocking is a tool for isolating the class under test from its \textbf{collaborators with real-world side effects}: databases, HTTP clients, message queues, clocks, random number generators, email senders. It is not a tool for isolating a class from everything it touches.

\textbf{Do mock:}

\begin{itemize}
\item 
Repositories/DAOs (avoid hitting a real database in a unit test)

\item 
External service clients (payment gateway, email provider, third-party API)

\item 
Anything with side effects you don’t want during a test (file system, network)

\end{itemize}

\textbf{Do NOT mock:}

\begin{itemize}
\item 
\textbf{Value objects / DTOs} (\code{Money}, \code{Address}, records, simple POJOs). These have no behavior worth faking — just construct a real instance. Mocking a value object produces an object whose \code{equals()}/\code{getters} are unpredictable stubs instead of real values, which actively makes the test harder to trust.

\end{itemize}

\begin{codebox}
\begin{Verbatim}[commandchars=\\\{\}]
\PY{c+c1}{// Bad: mocking a plain value object}
\PY{n}{Money}\PY{+w}{ }\PY{n}{mockMoney}\PY{+w}{ }\PY{o}{=}\PY{+w}{ }\PY{n}{mock}\PY{p}{(}\PY{n}{Money}\PY{p}{.}\PY{n+na}{class}\PY{p}{)}\PY{p}{;}
\PY{n}{when}\PY{p}{(}\PY{n}{mockMoney}\PY{p}{.}\PY{n+na}{amount}\PY{p}{(}\PY{p}{)}\PY{p}{)}\PY{p}{.}\PY{n+na}{thenReturn}\PY{p}{(}\PY{k}{new}\PY{+w}{ }\PY{n}{BigDecimal}\PY{p}{(}\PY{l+s}{\PYZdq{}}\PY{l+s}{10.00}\PY{l+s}{\PYZdq{}}\PY{p}{)}\PY{p}{)}\PY{p}{;}

\PY{c+c1}{// Good: it\PYZsq{}s just data — construct it}
\PY{n}{Money}\PY{+w}{ }\PY{n}{money}\PY{+w}{ }\PY{o}{=}\PY{+w}{ }\PY{k}{new}\PY{+w}{ }\PY{n}{Money}\PY{p}{(}\PY{k}{new}\PY{+w}{ }\PY{n}{BigDecimal}\PY{p}{(}\PY{l+s}{\PYZdq{}}\PY{l+s}{10.00}\PY{l+s}{\PYZdq{}}\PY{p}{)}\PY{p}{,}\PY{+w}{ }\PY{l+s}{\PYZdq{}}\PY{l+s}{USD}\PY{l+s}{\PYZdq{}}\PY{p}{)}\PY{p}{;}
\end{Verbatim}
\end{codebox}

\begin{itemize}
\item 
\textbf{Types you don’t own} (JDK classes, third-party library classes like \code{HttpClient}, \code{Connection}, \code{File}). Mocking framework internals couples your tests to implementation details of a library that can change between versions, and often the type isn’t even designed to be mocked (\code{final} classes/methods, complex internal state). Instead, \textbf{wrap it} in a small interface you own, and mock \emph{that} interface.

\end{itemize}

\begin{codebox}
\begin{Verbatim}[commandchars=\\\{\}]
\PY{c+c1}{// Bad: mocking java.time.Clock\PYZsq{}s cousin, a third\PYZhy{}party HttpClient directly everywhere}
\PY{n}{HttpClient}\PY{+w}{ }\PY{n}{mockClient}\PY{+w}{ }\PY{o}{=}\PY{+w}{ }\PY{n}{mock}\PY{p}{(}\PY{n}{HttpClient}\PY{p}{.}\PY{n+na}{class}\PY{p}{)}\PY{p}{;}
\PY{n}{when}\PY{p}{(}\PY{n}{mockClient}\PY{p}{.}\PY{n+na}{send}\PY{p}{(}\PY{n}{any}\PY{p}{(}\PY{p}{)}\PY{p}{,}\PY{+w}{ }\PY{n}{any}\PY{p}{(}\PY{p}{)}\PY{p}{)}\PY{p}{)}\PY{p}{.}\PY{n+na}{thenReturn}\PY{p}{(}\PY{n}{fakeResponse}\PY{p}{(}\PY{p}{)}\PY{p}{)}\PY{p}{;}

\PY{c+c1}{// Good: define a narrow seam you control}
\PY{k+kd}{public}\PY{+w}{ }\PY{k+kd}{interface} \PY{n+nc}{WeatherClient}\PY{+w}{ }\PY{p}{\PYZob{}}
\PY{+w}{    }\PY{n}{WeatherReport}\PY{+w}{ }\PY{n+nf}{fetch}\PY{p}{(}\PY{n}{String}\PY{+w}{ }\PY{n}{city}\PY{p}{)}\PY{p}{;}
\PY{p}{\PYZcb{}}
\PY{c+c1}{// production implementation wraps the real HttpClient/SDK}
\PY{c+c1}{// test code mocks WeatherClient, a one\PYZhy{}method interface you own}
\PY{n}{WeatherClient}\PY{+w}{ }\PY{n}{mockClient}\PY{+w}{ }\PY{o}{=}\PY{+w}{ }\PY{n}{mock}\PY{p}{(}\PY{n}{WeatherClient}\PY{p}{.}\PY{n+na}{class}\PY{p}{)}\PY{p}{;}
\PY{n}{when}\PY{p}{(}\PY{n}{mockClient}\PY{p}{.}\PY{n+na}{fetch}\PY{p}{(}\PY{l+s}{\PYZdq{}}\PY{l+s}{Berlin}\PY{l+s}{\PYZdq{}}\PY{p}{)}\PY{p}{)}\PY{p}{.}\PY{n+na}{thenReturn}\PY{p}{(}\PY{k}{new}\PY{+w}{ }\PY{n}{WeatherReport}\PY{p}{(}\PY{l+m+mi}{18}\PY{p}{,}\PY{+w}{ }\PY{l+s}{\PYZdq{}}\PY{l+s}{cloudy}\PY{l+s}{\PYZdq{}}\PY{p}{)}\PY{p}{)}\PY{p}{;}
\end{Verbatim}
\end{codebox}

\begin{itemize}
\item 
\textbf{The class under test itself.} If you find yourself mocking half the methods on the very class you’re testing, that’s a sign the class needs to be split, not mocked around.

\end{itemize}

\subsection[Creating Mocks: mock(), @Mock, @InjectMocks, @Spy]{Creating Mocks: \code{mock()}, \code{@\allowbreak{}Mock}, \code{@\allowbreak{}InjectMocks}, \code{@\allowbreak{}Spy}}
\label{h:20:creating-mocks-mock-mock-injectmocks-spy}
The programmatic form:

\begin{codebox}
\begin{Verbatim}[commandchars=\\\{\}]
\PY{n}{UserRepository}\PY{+w}{ }\PY{n}{userRepository}\PY{+w}{ }\PY{o}{=}\PY{+w}{ }\PY{n}{mock}\PY{p}{(}\PY{n}{UserRepository}\PY{p}{.}\PY{n+na}{class}\PY{p}{)}\PY{p}{;}
\end{Verbatim}
\end{codebox}

The annotation form, wired up by \code{MockitoExtension} (Mockito’s JUnit 5 integration):

\begin{codebox}
\begin{Verbatim}[commandchars=\\\{\}]
\PY{k+kn}{import}\PY{+w}{ }\PY{n+nn}{org.junit.jupiter.api.Test}\PY{p}{;}
\PY{k+kn}{import}\PY{+w}{ }\PY{n+nn}{org.junit.jupiter.api.extension.ExtendWith}\PY{p}{;}
\PY{k+kn}{import}\PY{+w}{ }\PY{n+nn}{org.mockito.InjectMocks}\PY{p}{;}
\PY{k+kn}{import}\PY{+w}{ }\PY{n+nn}{org.mockito.Mock}\PY{p}{;}
\PY{k+kn}{import}\PY{+w}{ }\PY{n+nn}{org.mockito.Spy}\PY{p}{;}
\PY{k+kn}{import}\PY{+w}{ }\PY{n+nn}{org.mockito.junit.jupiter.MockitoExtension}\PY{p}{;}

\PY{k+kn}{import static}\PY{+w}{ }\PY{n+nn}{org.mockito.Mockito.*}\PY{p}{;}
\PY{k+kn}{import static}\PY{+w}{ }\PY{n+nn}{org.junit.jupiter.api.Assertions.*}\PY{p}{;}

\PY{n+nd}{@ExtendWith}\PY{p}{(}\PY{n}{MockitoExtension}\PY{p}{.}\PY{n+na}{class}\PY{p}{)}
\PY{k+kd}{class} \PY{n+nc}{UserRegistrationServiceTest}\PY{+w}{ }\PY{p}{\PYZob{}}

\PY{+w}{    }\PY{n+nd}{@Mock}
\PY{+w}{    }\PY{k+kd}{private}\PY{+w}{ }\PY{n}{UserRepository}\PY{+w}{ }\PY{n}{userRepository}\PY{p}{;}

\PY{+w}{    }\PY{n+nd}{@Mock}
\PY{+w}{    }\PY{k+kd}{private}\PY{+w}{ }\PY{n}{EmailSender}\PY{+w}{ }\PY{n}{emailSender}\PY{p}{;}

\PY{+w}{    }\PY{n+nd}{@Spy}
\PY{+w}{    }\PY{k+kd}{private}\PY{+w}{ }\PY{n}{AuditLog}\PY{+w}{ }\PY{n}{auditLog}\PY{+w}{ }\PY{o}{=}\PY{+w}{ }\PY{k}{new}\PY{+w}{ }\PY{n}{AuditLog}\PY{p}{(}\PY{p}{)}\PY{p}{;}\PY{+w}{ }\PY{c+c1}{// real object, selectively overridable}

\PY{+w}{    }\PY{n+nd}{@InjectMocks}
\PY{+w}{    }\PY{k+kd}{private}\PY{+w}{ }\PY{n}{UserRegistrationService}\PY{+w}{ }\PY{n}{registrationService}\PY{p}{;}\PY{+w}{ }\PY{c+c1}{// Mockito injects the mocks above via constructor}

\PY{+w}{    }\PY{n+nd}{@Test}
\PY{+w}{    }\PY{k+kt}{void}\PY{+w}{ }\PY{n+nf}{register\PYZus{}newEmail\PYZus{}savesUserAndSendsWelcomeEmail}\PY{p}{(}\PY{p}{)}\PY{+w}{ }\PY{p}{\PYZob{}}
\PY{+w}{        }\PY{n}{when}\PY{p}{(}\PY{n}{userRepository}\PY{p}{.}\PY{n+na}{existsByEmail}\PY{p}{(}\PY{l+s}{\PYZdq{}}\PY{l+s}{alice@example.com}\PY{l+s}{\PYZdq{}}\PY{p}{)}\PY{p}{)}\PY{p}{.}\PY{n+na}{thenReturn}\PY{p}{(}\PY{k+kc}{false}\PY{p}{)}\PY{p}{;}

\PY{+w}{        }\PY{n}{registrationService}\PY{p}{.}\PY{n+na}{register}\PY{p}{(}\PY{l+s}{\PYZdq{}}\PY{l+s}{alice@example.com}\PY{l+s}{\PYZdq{}}\PY{p}{,}\PY{+w}{ }\PY{l+s}{\PYZdq{}}\PY{l+s}{Alice}\PY{l+s}{\PYZdq{}}\PY{p}{)}\PY{p}{;}

\PY{+w}{        }\PY{n}{verify}\PY{p}{(}\PY{n}{userRepository}\PY{p}{)}\PY{p}{.}\PY{n+na}{save}\PY{p}{(}\PY{n}{any}\PY{p}{(}\PY{n}{User}\PY{p}{.}\PY{n+na}{class}\PY{p}{)}\PY{p}{)}\PY{p}{;}
\PY{+w}{        }\PY{n}{verify}\PY{p}{(}\PY{n}{emailSender}\PY{p}{)}\PY{p}{.}\PY{n+na}{sendWelcomeEmail}\PY{p}{(}\PY{l+s}{\PYZdq{}}\PY{l+s}{alice@example.com}\PY{l+s}{\PYZdq{}}\PY{p}{)}\PY{p}{;}
\PY{+w}{    }\PY{p}{\PYZcb{}}
\PY{p}{\PYZcb{}}
\end{Verbatim}
\end{codebox}

\code{@\allowbreak{}InjectMocks} tries constructor injection first (preferred — it fails loudly at compile time if the constructor shape changes), then setter injection, then field injection. If \code{User\allowbreak{}Registration\allowbreak{}Service} has a single constructor taking \code{UserRepository}, \code{EmailSender}, and \code{AuditLog}, Mockito matches the \code{@\allowbreak{}Mock}/\code{@\allowbreak{}Spy} fields to those parameters by type.

\code{@\allowbreak{}Spy} wraps a \textbf{real} object: calls go to the real implementation unless you explicitly stub them, and Mockito still records every call for \code{verify(...)}. Use a spy when you want most of the real behavior but need to intercept one troublesome method (e.g., a legacy class that does real I/O in one method you want to fake out for the test).

\subsection[Stubbing]{Stubbing}
\label{h:20:stubbing}
\textbf{\code{when().\allowbreak{}thenReturn(...)}} is the everyday form for stubbing non-void methods on a mock:

\begin{codebox}
\begin{Verbatim}[commandchars=\\\{\}]
\PY{n}{when}\PY{p}{(}\PY{n}{userRepository}\PY{p}{.}\PY{n+na}{findById}\PY{p}{(}\PY{l+m+mi}{1L}\PY{p}{)}\PY{p}{)}\PY{p}{.}\PY{n+na}{thenReturn}\PY{p}{(}\PY{n}{Optional}\PY{p}{.}\PY{n+na}{of}\PY{p}{(}\PY{k}{new}\PY{+w}{ }\PY{n}{User}\PY{p}{(}\PY{l+m+mi}{1L}\PY{p}{,}\PY{+w}{ }\PY{l+s}{\PYZdq{}}\PY{l+s}{Alice}\PY{l+s}{\PYZdq{}}\PY{p}{)}\PY{p}{)}\PY{p}{)}\PY{p}{;}
\end{Verbatim}
\end{codebox}

\textbf{\code{thenThrow}} stubs an exception:

\begin{codebox}
\begin{Verbatim}[commandchars=\\\{\}]
\PY{n}{when}\PY{p}{(}\PY{n}{paymentGateway}\PY{p}{.}\PY{n+na}{charge}\PY{p}{(}\PY{n}{any}\PY{p}{(}\PY{p}{)}\PY{p}{)}\PY{p}{)}\PY{p}{.}\PY{n+na}{thenThrow}\PY{p}{(}\PY{k}{new}\PY{+w}{ }\PY{n}{PaymentDeclinedException}\PY{p}{(}\PY{l+s}{\PYZdq{}}\PY{l+s}{insufficient funds}\PY{l+s}{\PYZdq{}}\PY{p}{)}\PY{p}{)}\PY{p}{;}
\end{Verbatim}
\end{codebox}

\textbf{\code{thenAnswer}} for computed/dynamic responses based on the actual arguments:

\begin{codebox}
\begin{Verbatim}[commandchars=\\\{\}]
\PY{n}{when}\PY{p}{(}\PY{n}{userRepository}\PY{p}{.}\PY{n+na}{save}\PY{p}{(}\PY{n}{any}\PY{p}{(}\PY{n}{User}\PY{p}{.}\PY{n+na}{class}\PY{p}{)}\PY{p}{)}\PY{p}{)}\PY{p}{.}\PY{n+na}{thenAnswer}\PY{p}{(}\PY{n}{invocation}\PY{+w}{ }\PY{o}{\PYZhy{}}\PY{o}{\PYZgt{}}\PY{+w}{ }\PY{p}{\PYZob{}}
\PY{+w}{    }\PY{n}{User}\PY{+w}{ }\PY{n}{user}\PY{+w}{ }\PY{o}{=}\PY{+w}{ }\PY{n}{invocation}\PY{p}{.}\PY{n+na}{getArgument}\PY{p}{(}\PY{l+m+mi}{0}\PY{p}{)}\PY{p}{;}
\PY{+w}{    }\PY{k}{return}\PY{+w}{ }\PY{n}{user}\PY{p}{.}\PY{n+na}{withId}\PY{p}{(}\PY{l+m+mi}{42L}\PY{p}{)}\PY{p}{;}\PY{+w}{ }\PY{c+c1}{// simulate the DB assigning an id}
\PY{p}{\PYZcb{}}\PY{p}{)}\PY{p}{;}
\end{Verbatim}
\end{codebox}

\textbf{Consecutive returns} — different answers on successive calls, useful for simulating retry logic:

\begin{codebox}
\begin{Verbatim}[commandchars=\\\{\}]
\PY{n}{when}\PY{p}{(}\PY{n}{flakyClient}\PY{p}{.}\PY{n+na}{fetch}\PY{p}{(}\PY{p}{)}\PY{p}{)}
\PY{+w}{    }\PY{p}{.}\PY{n+na}{thenThrow}\PY{p}{(}\PY{k}{new}\PY{+w}{ }\PY{n}{IOException}\PY{p}{(}\PY{l+s}{\PYZdq{}}\PY{l+s}{timeout}\PY{l+s}{\PYZdq{}}\PY{p}{)}\PY{p}{)}
\PY{+w}{    }\PY{p}{.}\PY{n+na}{thenReturn}\PY{p}{(}\PY{l+s}{\PYZdq{}}\PY{l+s}{success}\PY{l+s}{\PYZdq{}}\PY{p}{)}\PY{p}{;}\PY{+w}{ }\PY{c+c1}{// first call throws, second call (and beyond) returns \PYZdq{}success\PYZdq{}}
\end{Verbatim}
\end{codebox}

\subsection[doReturn/doThrow/doNothing().when(...) — and When You Must Use It]{\code{doReturn}/\code{doThrow}/\code{doNothing().\allowbreak{}when(...)} — and When You Must Use It}
\label{h:20:doreturndothrowdonothingwhen--and-when-you-must-use-it}
The \code{when(\allowbreak{}mock.\allowbreak{}method()).\allowbreak{}then\allowbreak{}Return(...)} form has a hard limitation: it actually \emph{calls} the real method on the mock first to figure out what to stub. That is invisible for a plain mock (Mockito intercepts the call before any real code runs), but it breaks down in two situations:

\begin{enumerate}
\item 
\textbf{Spies} — because a spy \emph{does} run real code, \code{when(\allowbreak{}spy.\allowbreak{}risky\allowbreak{}Method()).\allowbreak{}then\allowbreak{}Return(...)} would execute the real (possibly side-effecting or exception-throwing) method before Mockito even gets to stub it.

\item 
\textbf{Void methods} — \code{when(...)} needs a return value to chain off of; a void method call has none, so there is nothing to pass to \code{when(...)}.

\end{enumerate}

The \code{do...().\allowbreak{}when(\allowbreak{}mock)} family avoids this by stubbing without invoking the real method first:

\begin{codebox}
\begin{Verbatim}[commandchars=\\\{\}]
\PY{n+nd}{@Spy}
\PY{k+kd}{private}\PY{+w}{ }\PY{n}{OrderValidator}\PY{+w}{ }\PY{n}{validator}\PY{+w}{ }\PY{o}{=}\PY{+w}{ }\PY{k}{new}\PY{+w}{ }\PY{n}{OrderValidator}\PY{p}{(}\PY{p}{)}\PY{p}{;}\PY{+w}{ }\PY{c+c1}{// real object; validate() does real, expensive checks}

\PY{n+nd}{@Test}
\PY{k+kt}{void}\PY{+w}{ }\PY{n+nf}{skipsExpensiveValidationInThisTest}\PY{p}{(}\PY{p}{)}\PY{+w}{ }\PY{p}{\PYZob{}}
\PY{+w}{    }\PY{n}{doReturn}\PY{p}{(}\PY{k+kc}{true}\PY{p}{)}\PY{p}{.}\PY{n+na}{when}\PY{p}{(}\PY{n}{validator}\PY{p}{)}\PY{p}{.}\PY{n+na}{isPreValidated}\PY{p}{(}\PY{n}{any}\PY{p}{(}\PY{p}{)}\PY{p}{)}\PY{p}{;}\PY{+w}{ }\PY{c+c1}{// safe on a spy; when(validator...) would run the real check first}

\PY{+w}{    }\PY{c+c1}{// ...}
\PY{p}{\PYZcb{}}

\PY{n+nd}{@Mock}
\PY{k+kd}{private}\PY{+w}{ }\PY{n}{NotificationService}\PY{+w}{ }\PY{n}{notificationService}\PY{p}{;}

\PY{n+nd}{@Test}
\PY{k+kt}{void}\PY{+w}{ }\PY{n+nf}{notify\PYZus{}whenServiceUnavailable\PYZus{}doesNothingInsteadOfThrowing}\PY{p}{(}\PY{p}{)}\PY{+w}{ }\PY{p}{\PYZob{}}
\PY{+w}{    }\PY{n}{doThrow}\PY{p}{(}\PY{k}{new}\PY{+w}{ }\PY{n}{ServiceUnavailableException}\PY{p}{(}\PY{p}{)}\PY{p}{)}\PY{p}{.}\PY{n+na}{when}\PY{p}{(}\PY{n}{notificationService}\PY{p}{)}\PY{p}{.}\PY{n+na}{notifyUser}\PY{p}{(}\PY{n}{any}\PY{p}{(}\PY{p}{)}\PY{p}{)}\PY{p}{;}\PY{+w}{ }\PY{c+c1}{// void method}

\PY{+w}{    }\PY{n}{orderService}\PY{p}{.}\PY{n+na}{completeOrder}\PY{p}{(}\PY{n}{order}\PY{p}{)}\PY{p}{;}\PY{+w}{ }\PY{c+c1}{// should swallow the notification failure gracefully}

\PY{+w}{    }\PY{n}{verify}\PY{p}{(}\PY{n}{notificationService}\PY{p}{)}\PY{p}{.}\PY{n+na}{notifyUser}\PY{p}{(}\PY{n}{order}\PY{p}{.}\PY{n+na}{customerId}\PY{p}{(}\PY{p}{)}\PY{p}{)}\PY{p}{;}
\PY{p}{\PYZcb{}}

\PY{n+nd}{@Test}
\PY{k+kt}{void}\PY{+w}{ }\PY{n+nf}{suppressesAuditLoggingDuringThisScenario}\PY{p}{(}\PY{p}{)}\PY{+w}{ }\PY{p}{\PYZob{}}
\PY{+w}{    }\PY{n}{doNothing}\PY{p}{(}\PY{p}{)}\PY{p}{.}\PY{n+na}{when}\PY{p}{(}\PY{n}{auditLog}\PY{p}{)}\PY{p}{.}\PY{n+na}{record}\PY{p}{(}\PY{n}{any}\PY{p}{(}\PY{p}{)}\PY{p}{)}\PY{p}{;}\PY{+w}{ }\PY{c+c1}{// void method, explicit no\PYZhy{}op}
\PY{+w}{    }\PY{c+c1}{// ...}
\PY{p}{\PYZcb{}}
\end{Verbatim}
\end{codebox}

Rule of thumb for reviewers: \code{when(...).\allowbreak{}then\allowbreak{}Return(...)} for normal mocks and non-void methods; \code{doX().\allowbreak{}when(...)} whenever the target is a \textbf{spy} or the method is \textbf{void}.

\subsection[Argument Matchers]{Argument Matchers}
\label{h:20:argument-matchers}
Matchers let you stub/verify based on argument \emph{shape} rather than an exact instance:

\begin{codebox}
\begin{Verbatim}[commandchars=\\\{\}]
\PY{n}{when}\PY{p}{(}\PY{n}{userRepository}\PY{p}{.}\PY{n+na}{findByEmail}\PY{p}{(}\PY{n}{anyString}\PY{p}{(}\PY{p}{)}\PY{p}{)}\PY{p}{)}\PY{p}{.}\PY{n+na}{thenReturn}\PY{p}{(}\PY{n}{Optional}\PY{p}{.}\PY{n+na}{empty}\PY{p}{(}\PY{p}{)}\PY{p}{)}\PY{p}{;}
\PY{n}{when}\PY{p}{(}\PY{n}{userRepository}\PY{p}{.}\PY{n+na}{findById}\PY{p}{(}\PY{n}{eq}\PY{p}{(}\PY{l+m+mi}{1L}\PY{p}{)}\PY{p}{)}\PY{p}{)}\PY{p}{.}\PY{n+na}{thenReturn}\PY{p}{(}\PY{n}{Optional}\PY{p}{.}\PY{n+na}{of}\PY{p}{(}\PY{n}{alice}\PY{p}{)}\PY{p}{)}\PY{p}{;}
\PY{n}{when}\PY{p}{(}\PY{n}{pricingEngine}\PY{p}{.}\PY{n+na}{priceFor}\PY{p}{(}\PY{n}{argThat}\PY{p}{(}\PY{n}{order}\PY{+w}{ }\PY{o}{\PYZhy{}}\PY{o}{\PYZgt{}}\PY{+w}{ }\PY{n}{order}\PY{p}{.}\PY{n+na}{items}\PY{p}{(}\PY{p}{)}\PY{p}{.}\PY{n+na}{size}\PY{p}{(}\PY{p}{)}\PY{+w}{ }\PY{o}{\PYZgt{}}\PY{+w}{ }\PY{l+m+mi}{5}\PY{p}{)}\PY{p}{)}\PY{p}{)}
\PY{+w}{    }\PY{p}{.}\PY{n+na}{thenReturn}\PY{p}{(}\PY{k}{new}\PY{+w}{ }\PY{n}{BigDecimal}\PY{p}{(}\PY{l+s}{\PYZdq{}}\PY{l+s}{0.00}\PY{l+s}{\PYZdq{}}\PY{p}{)}\PY{p}{)}\PY{p}{;}\PY{+w}{ }\PY{c+c1}{// free for large orders, expressed as a custom predicate}
\end{Verbatim}
\end{codebox}

\textbf{The critical rule: you cannot mix raw values and matchers in the same call.} If \emph{any} argument uses a matcher, \emph{every} argument in that call must use a matcher (use \code{eq(...)} to wrap literal values).

\begin{codebox}
\begin{Verbatim}[commandchars=\\\{\}]
\PY{c+c1}{// Throws InvalidUseOfMatchersException at runtime}
\PY{n}{when}\PY{p}{(}\PY{n}{userService}\PY{p}{.}\PY{n+na}{updateEmail}\PY{p}{(}\PY{l+m+mi}{1L}\PY{p}{,}\PY{+w}{ }\PY{n}{anyString}\PY{p}{(}\PY{p}{)}\PY{p}{)}\PY{p}{)}\PY{p}{.}\PY{n+na}{thenReturn}\PY{p}{(}\PY{k+kc}{true}\PY{p}{)}\PY{p}{;}

\PY{c+c1}{// Correct: wrap the literal in eq(...) once any() is used elsewhere}
\PY{n}{when}\PY{p}{(}\PY{n}{userService}\PY{p}{.}\PY{n+na}{updateEmail}\PY{p}{(}\PY{n}{eq}\PY{p}{(}\PY{l+m+mi}{1L}\PY{p}{)}\PY{p}{,}\PY{+w}{ }\PY{n}{anyString}\PY{p}{(}\PY{p}{)}\PY{p}{)}\PY{p}{)}\PY{p}{.}\PY{n+na}{thenReturn}\PY{p}{(}\PY{k+kc}{true}\PY{p}{)}\PY{p}{;}
\end{Verbatim}
\end{codebox}

This is a very common copy-paste mistake in review — a matcher gets added to one argument while a neighboring literal is left bare, and the test blows up with a confusing runtime exception instead of a compile error.

\subsection[Verification]{Verification}
\label{h:20:verification}
Stubbing controls what a mock \emph{returns}; verification checks what was \emph{called}.

\begin{codebox}
\begin{Verbatim}[commandchars=\\\{\}]
\PY{n}{verify}\PY{p}{(}\PY{n}{emailSender}\PY{p}{)}\PY{p}{.}\PY{n+na}{sendWelcomeEmail}\PY{p}{(}\PY{l+s}{\PYZdq{}}\PY{l+s}{alice@example.com}\PY{l+s}{\PYZdq{}}\PY{p}{)}\PY{p}{;}\PY{+w}{      }\PY{c+c1}{// called exactly once, with this argument}
\PY{n}{verify}\PY{p}{(}\PY{n}{emailSender}\PY{p}{,}\PY{+w}{ }\PY{n}{times}\PY{p}{(}\PY{l+m+mi}{2}\PY{p}{)}\PY{p}{)}\PY{p}{.}\PY{n+na}{sendReminder}\PY{p}{(}\PY{n}{any}\PY{p}{(}\PY{p}{)}\PY{p}{)}\PY{p}{;}\PY{+w}{               }\PY{c+c1}{// called exactly twice}
\PY{n}{verify}\PY{p}{(}\PY{n}{emailSender}\PY{p}{,}\PY{+w}{ }\PY{n}{never}\PY{p}{(}\PY{p}{)}\PY{p}{)}\PY{p}{.}\PY{n+na}{sendWelcomeEmail}\PY{p}{(}\PY{l+s}{\PYZdq{}}\PY{l+s}{blocked@example.com}\PY{l+s}{\PYZdq{}}\PY{p}{)}\PY{p}{;}
\PY{n}{verify}\PY{p}{(}\PY{n}{emailSender}\PY{p}{,}\PY{+w}{ }\PY{n}{atLeast}\PY{p}{(}\PY{l+m+mi}{1}\PY{p}{)}\PY{p}{)}\PY{p}{.}\PY{n+na}{flush}\PY{p}{(}\PY{p}{)}\PY{p}{;}
\PY{n}{verify}\PY{p}{(}\PY{n}{emailSender}\PY{p}{,}\PY{+w}{ }\PY{n}{atMost}\PY{p}{(}\PY{l+m+mi}{3}\PY{p}{)}\PY{p}{)}\PY{p}{.}\PY{n+na}{retry}\PY{p}{(}\PY{p}{)}\PY{p}{;}

\PY{n}{InOrder}\PY{+w}{ }\PY{n}{inOrder}\PY{+w}{ }\PY{o}{=}\PY{+w}{ }\PY{n}{inOrder}\PY{p}{(}\PY{n}{userRepository}\PY{p}{,}\PY{+w}{ }\PY{n}{emailSender}\PY{p}{)}\PY{p}{;}
\PY{n}{inOrder}\PY{p}{.}\PY{n+na}{verify}\PY{p}{(}\PY{n}{userRepository}\PY{p}{)}\PY{p}{.}\PY{n+na}{save}\PY{p}{(}\PY{n}{any}\PY{p}{(}\PY{n}{User}\PY{p}{.}\PY{n+na}{class}\PY{p}{)}\PY{p}{)}\PY{p}{;}
\PY{n}{inOrder}\PY{p}{.}\PY{n+na}{verify}\PY{p}{(}\PY{n}{emailSender}\PY{p}{)}\PY{p}{.}\PY{n+na}{sendWelcomeEmail}\PY{p}{(}\PY{n}{anyString}\PY{p}{(}\PY{p}{)}\PY{p}{)}\PY{p}{;}\PY{+w}{       }\PY{c+c1}{// save() must happen before the email is sent}

\PY{n}{verify}\PY{p}{(}\PY{n}{userRepository}\PY{p}{)}\PY{p}{.}\PY{n+na}{save}\PY{p}{(}\PY{n}{any}\PY{p}{(}\PY{n}{User}\PY{p}{.}\PY{n+na}{class}\PY{p}{)}\PY{p}{)}\PY{p}{;}
\PY{n}{verify}\PY{p}{(}\PY{n}{emailSender}\PY{p}{)}\PY{p}{.}\PY{n+na}{sendWelcomeEmail}\PY{p}{(}\PY{n}{anyString}\PY{p}{(}\PY{p}{)}\PY{p}{)}\PY{p}{;}
\PY{n}{verifyNoMoreInteractions}\PY{p}{(}\PY{n}{userRepository}\PY{p}{,}\PY{+w}{ }\PY{n}{emailSender}\PY{p}{)}\PY{p}{;}\PY{+w}{            }\PY{c+c1}{// fails if any OTHER call happened on these mocks}
\end{Verbatim}
\end{codebox}

\code{verify\allowbreak{}No\allowbreak{}More\allowbreak{}Interactions} is a strong assertion — it fails the test if the code under test called anything on that mock beyond what you’ve already verified. It’s a good way to lock down that a method does \emph{exactly} what you expect and nothing extra, but overusing it on every mock in every test makes tests brittle against harmless refactors (e.g., adding a debug-log call). Use it selectively, on the interactions that genuinely matter.

\subsection[ArgumentCaptor]{ArgumentCaptor}
\label{h:20:argumentcaptor}
When you need to inspect the \emph{value} an argument had, not just that a call happened, capture it:

\begin{codebox}
\begin{Verbatim}[commandchars=\\\{\}]
\PY{n+nd}{@Captor}
\PY{k+kd}{private}\PY{+w}{ }\PY{n}{ArgumentCaptor}\PY{o}{\PYZlt{}}\PY{n}{User}\PY{o}{\PYZgt{}}\PY{+w}{ }\PY{n}{userCaptor}\PY{p}{;}

\PY{n+nd}{@Test}
\PY{k+kt}{void}\PY{+w}{ }\PY{n+nf}{register\PYZus{}savesUserWithNormalizedEmail}\PY{p}{(}\PY{p}{)}\PY{+w}{ }\PY{p}{\PYZob{}}
\PY{+w}{    }\PY{n}{registrationService}\PY{p}{.}\PY{n+na}{register}\PY{p}{(}\PY{l+s}{\PYZdq{}}\PY{l+s}{  Alice@Example.com }\PY{l+s}{\PYZdq{}}\PY{p}{,}\PY{+w}{ }\PY{l+s}{\PYZdq{}}\PY{l+s}{Alice}\PY{l+s}{\PYZdq{}}\PY{p}{)}\PY{p}{;}

\PY{+w}{    }\PY{n}{verify}\PY{p}{(}\PY{n}{userRepository}\PY{p}{)}\PY{p}{.}\PY{n+na}{save}\PY{p}{(}\PY{n}{userCaptor}\PY{p}{.}\PY{n+na}{capture}\PY{p}{(}\PY{p}{)}\PY{p}{)}\PY{p}{;}
\PY{+w}{    }\PY{n}{User}\PY{+w}{ }\PY{n}{savedUser}\PY{+w}{ }\PY{o}{=}\PY{+w}{ }\PY{n}{userCaptor}\PY{p}{.}\PY{n+na}{getValue}\PY{p}{(}\PY{p}{)}\PY{p}{;}

\PY{+w}{    }\PY{n}{assertEquals}\PY{p}{(}\PY{l+s}{\PYZdq{}}\PY{l+s}{alice@example.com}\PY{l+s}{\PYZdq{}}\PY{p}{,}\PY{+w}{ }\PY{n}{savedUser}\PY{p}{.}\PY{n+na}{email}\PY{p}{(}\PY{p}{)}\PY{p}{)}\PY{p}{;}\PY{+w}{ }\PY{c+c1}{// trimmed and lower\PYZhy{}cased}
\PY{p}{\PYZcb{}}
\end{Verbatim}
\end{codebox}

\code{ArgumentCaptor} is essential when a mock’s method returns \code{void} (or nothing you need) but you still want to assert on exactly what was passed in — verification alone tells you a call happened; the captor tells you what it looked like.

\subsection[RETURNS\_DEEP\_STUBS and Why It’s a Smell]{\code{RETURNS\_\allowbreak{}DEEP\_\allowbreak{}STUBS} and Why It’s a Smell}
\label{h:20:returns_deep_stubs-and-why-its-a-smell}
Mockito can auto-generate stubs for chained calls:

\begin{codebox}
\begin{Verbatim}[commandchars=\\\{\}]
\PY{c+c1}{// Works, but...}
\PY{n}{OrderContext}\PY{+w}{ }\PY{n}{context}\PY{+w}{ }\PY{o}{=}\PY{+w}{ }\PY{n}{mock}\PY{p}{(}\PY{n}{OrderContext}\PY{p}{.}\PY{n+na}{class}\PY{p}{,}\PY{+w}{ }\PY{n}{Mockito}\PY{p}{.}\PY{n+na}{RETURNS\PYZus{}DEEP\PYZus{}STUBS}\PY{p}{)}\PY{p}{;}
\PY{n}{when}\PY{p}{(}\PY{n}{context}\PY{p}{.}\PY{n+na}{getCustomer}\PY{p}{(}\PY{p}{)}\PY{p}{.}\PY{n+na}{getAddress}\PY{p}{(}\PY{p}{)}\PY{p}{.}\PY{n+na}{getCountry}\PY{p}{(}\PY{p}{)}\PY{p}{)}\PY{p}{.}\PY{n+na}{thenReturn}\PY{p}{(}\PY{l+s}{\PYZdq{}}\PY{l+s}{DE}\PY{l+s}{\PYZdq{}}\PY{p}{)}\PY{p}{;}
\end{Verbatim}
\end{codebox}

This “solves” a symptom while hiding the real problem: a chain like \code{context.\allowbreak{}get\allowbreak{}Customer().\allowbreak{}get\allowbreak{}Address().\allowbreak{}get\allowbreak{}Country()} violates the \textbf{Law of Demeter} — the code under test is reaching three objects deep into someone else’s internals. \code{RETURNS\_\allowbreak{}DEEP\_\allowbreak{}STUBS} makes that chain \emph{mockable}, but it does not make it \emph{good design}. The test passing this way is a strong signal that the production code should instead depend directly on a \code{Country} (passed in, or obtained through a single well-named method like \code{context.\allowbreak{}get\allowbreak{}Customer\allowbreak{}Country()}), which would need no deep stub at all.

\subsection[Strict Stubs and UnnecessaryStubbingException]{Strict Stubs and \code{Unnecessary\allowbreak{}Stubbing\allowbreak{}Exception}}
\label{h:20:strict-stubs-and-unnecessarystubbingexception}
\code{MockitoExtension} runs in \textbf{strict stubbing} mode by default: if you stub something that the test never actually calls, the test fails with \code{Unnecessary\allowbreak{}Stubbing\allowbreak{}Exception} (or, depending on version/settings, a warning that fails on \code{Mockito.\allowbreak{}validate\allowbreak{}Mockito\allowbreak{}Usage()}), rather than silently passing.

\begin{codebox}
\begin{Verbatim}[commandchars=\\\{\}]
\PY{n+nd}{@ExtendWith}\PY{p}{(}\PY{n}{MockitoExtension}\PY{p}{.}\PY{n+na}{class}\PY{p}{)}
\PY{k+kd}{class} \PY{n+nc}{DiscountServiceTest}\PY{+w}{ }\PY{p}{\PYZob{}}

\PY{+w}{    }\PY{n+nd}{@Mock}
\PY{+w}{    }\PY{k+kd}{private}\PY{+w}{ }\PY{n}{CouponRepository}\PY{+w}{ }\PY{n}{couponRepository}\PY{p}{;}

\PY{+w}{    }\PY{n+nd}{@Test}
\PY{+w}{    }\PY{k+kt}{void}\PY{+w}{ }\PY{n+nf}{appliesDiscount}\PY{p}{(}\PY{p}{)}\PY{+w}{ }\PY{p}{\PYZob{}}
\PY{+w}{        }\PY{n}{when}\PY{p}{(}\PY{n}{couponRepository}\PY{p}{.}\PY{n+na}{findByCode}\PY{p}{(}\PY{l+s}{\PYZdq{}}\PY{l+s}{SAVE10}\PY{l+s}{\PYZdq{}}\PY{p}{)}\PY{p}{)}\PY{p}{.}\PY{n+na}{thenReturn}\PY{p}{(}\PY{n}{Optional}\PY{p}{.}\PY{n+na}{of}\PY{p}{(}\PY{k}{new}\PY{+w}{ }\PY{n}{Coupon}\PY{p}{(}\PY{l+m+mi}{10}\PY{p}{)}\PY{p}{)}\PY{p}{)}\PY{p}{;}
\PY{+w}{        }\PY{c+c1}{// If discountService never actually calls couponRepository.findByCode(\PYZdq{}SAVE10\PYZdq{})}
\PY{+w}{        }\PY{c+c1}{// (e.g., because the test forgot to pass a coupon code), this line is flagged}
\PY{+w}{        }\PY{c+c1}{// as an UnnecessaryStubbingException — a leftover stub from a copy\PYZhy{}pasted test.}
\PY{+w}{        }\PY{p}{.}\PY{p}{.}\PY{p}{.}
\PY{+w}{    }\PY{p}{\PYZcb{}}
\PY{p}{\PYZcb{}}
\end{Verbatim}
\end{codebox}

This is a deliberately helpful failure: unused stubs accumulate in copy-pasted tests and quietly rot, giving a false sense of what’s actually being exercised. Treat \code{Unnecessary\allowbreak{}Stubbing\allowbreak{}Exception} as a signal to delete the stub or fix the test, never to silence it with \code{lenient()} unless there’s a genuinely good reason (e.g., a shared \code{@\allowbreak{}BeforeEach} stub that only some tests in the class need).

\begin{codebox}
\begin{Verbatim}[commandchars=\\\{\}]
\PY{n}{lenient}\PY{p}{(}\PY{p}{)}\PY{p}{.}\PY{n+na}{when}\PY{p}{(}\PY{n}{couponRepository}\PY{p}{.}\PY{n+na}{findByCode}\PY{p}{(}\PY{n}{any}\PY{p}{(}\PY{p}{)}\PY{p}{)}\PY{p}{)}\PY{p}{.}\PY{n+na}{thenReturn}\PY{p}{(}\PY{n}{Optional}\PY{p}{.}\PY{n+na}{empty}\PY{p}{(}\PY{p}{)}\PY{p}{)}\PY{p}{;}\PY{+w}{ }\PY{c+c1}{// opt out of strict checking, deliberately}
\end{Verbatim}
\end{codebox}

\subsection[Mocking Static Methods]{Mocking Static Methods}
\label{h:20:mocking-static-methods}
Mockito 3.4+ supports mocking static methods via \code{mockStatic}, using inline mocking:

\begin{codebox}
\begin{Verbatim}[commandchars=\\\{\}]
\PY{k}{try}\PY{+w}{ }\PY{p}{(}\PY{n}{MockedStatic}\PY{o}{\PYZlt{}}\PY{n}{Instant}\PY{o}{\PYZgt{}}\PY{+w}{ }\PY{n}{mockedInstant}\PY{+w}{ }\PY{o}{=}\PY{+w}{ }\PY{n}{Mockito}\PY{p}{.}\PY{n+na}{mockStatic}\PY{p}{(}\PY{n}{Instant}\PY{p}{.}\PY{n+na}{class}\PY{p}{)}\PY{p}{)}\PY{+w}{ }\PY{p}{\PYZob{}}
\PY{+w}{    }\PY{n}{mockedInstant}\PY{p}{.}\PY{n+na}{when}\PY{p}{(}\PY{n}{Instant}\PY{p}{:}\PY{p}{:}\PY{n}{now}\PY{p}{)}\PY{p}{.}\PY{n+na}{thenReturn}\PY{p}{(}\PY{n}{Instant}\PY{p}{.}\PY{n+na}{parse}\PY{p}{(}\PY{l+s}{\PYZdq{}}\PY{l+s}{2026\PYZhy{}08\PYZhy{}07T00:00:00Z}\PY{l+s}{\PYZdq{}}\PY{p}{)}\PY{p}{)}\PY{p}{;}

\PY{+w}{    }\PY{n}{Report}\PY{+w}{ }\PY{n}{report}\PY{+w}{ }\PY{o}{=}\PY{+w}{ }\PY{n}{reportGenerator}\PY{p}{.}\PY{n+na}{generateDailyReport}\PY{p}{(}\PY{p}{)}\PY{p}{;}

\PY{+w}{    }\PY{n}{assertEquals}\PY{p}{(}\PY{n}{LocalDate}\PY{p}{.}\PY{n+na}{parse}\PY{p}{(}\PY{l+s}{\PYZdq{}}\PY{l+s}{2026\PYZhy{}08\PYZhy{}07}\PY{l+s}{\PYZdq{}}\PY{p}{)}\PY{p}{,}\PY{+w}{ }\PY{n}{report}\PY{p}{.}\PY{n+na}{date}\PY{p}{(}\PY{p}{)}\PY{p}{)}\PY{p}{;}
\PY{p}{\PYZcb{}}
\end{Verbatim}
\end{codebox}

This works, and the try-with-resources scoping (the static mock is only active inside the block) is important to prevent it from leaking into other tests. But \textbf{needing this at all is usually a design smell}: it means the production code called a static method directly (\code{Instant.\allowbreak{}now()}, a static utility, a singleton accessor) instead of depending on an injectable abstraction. The preferred fix, shown earlier in this chapter, is to inject a \code{Clock} (or wrap the static call behind an interface you own) — \code{mockStatic} should be a last resort for code you cannot easily refactor (e.g., a legacy static utility from a third-party library), not a first choice for new code.

\subsection[MockedConstruction]{\code{MockedConstruction}}
\label{h:20:mockedconstruction}
Similarly, \code{mockConstruction} intercepts \code{new~\allowbreak{}SomeType(...)} calls made \emph{inside} the code under test, redirecting them to a mock — useful when legacy code directly instantiates a collaborator instead of receiving it via dependency injection:

\begin{codebox}
\begin{Verbatim}[commandchars=\\\{\}]
\PY{k}{try}\PY{+w}{ }\PY{p}{(}\PY{n}{MockedConstruction}\PY{o}{\PYZlt{}}\PY{n}{PdfRenderer}\PY{o}{\PYZgt{}}\PY{+w}{ }\PY{n}{mocked}\PY{+w}{ }\PY{o}{=}\PY{+w}{ }\PY{n}{Mockito}\PY{p}{.}\PY{n+na}{mockConstruction}\PY{p}{(}\PY{n}{PdfRenderer}\PY{p}{.}\PY{n+na}{class}\PY{p}{,}
\PY{+w}{        }\PY{p}{(}\PY{n}{mockRenderer}\PY{p}{,}\PY{+w}{ }\PY{n}{context}\PY{p}{)}\PY{+w}{ }\PY{o}{\PYZhy{}}\PY{o}{\PYZgt{}}\PY{+w}{ }\PY{n}{when}\PY{p}{(}\PY{n}{mockRenderer}\PY{p}{.}\PY{n+na}{render}\PY{p}{(}\PY{n}{any}\PY{p}{(}\PY{p}{)}\PY{p}{)}\PY{p}{)}\PY{p}{.}\PY{n+na}{thenReturn}\PY{p}{(}\PY{k}{new}\PY{+w}{ }\PY{k+kt}{byte}\PY{o}{[}\PY{o}{]}\PY{+w}{ }\PY{p}{\PYZob{}}\PY{+w}{ }\PY{l+m+mi}{1}\PY{p}{,}\PY{+w}{ }\PY{l+m+mi}{2}\PY{p}{,}\PY{+w}{ }\PY{l+m+mi}{3}\PY{+w}{ }\PY{p}{\PYZcb{}}\PY{p}{)}\PY{p}{)}\PY{p}{)}\PY{+w}{ }\PY{p}{\PYZob{}}

\PY{+w}{    }\PY{k+kt}{byte}\PY{o}{[}\PY{o}{]}\PY{+w}{ }\PY{n}{pdf}\PY{+w}{ }\PY{o}{=}\PY{+w}{ }\PY{n}{invoiceService}\PY{p}{.}\PY{n+na}{generatePdf}\PY{p}{(}\PY{n}{invoice}\PY{p}{)}\PY{p}{;}\PY{+w}{ }\PY{c+c1}{// internally does `new PdfRenderer()`}

\PY{+w}{    }\PY{n}{assertArrayEquals}\PY{p}{(}\PY{k}{new}\PY{+w}{ }\PY{k+kt}{byte}\PY{o}{[}\PY{o}{]}\PY{+w}{ }\PY{p}{\PYZob{}}\PY{+w}{ }\PY{l+m+mi}{1}\PY{p}{,}\PY{+w}{ }\PY{l+m+mi}{2}\PY{p}{,}\PY{+w}{ }\PY{l+m+mi}{3}\PY{+w}{ }\PY{p}{\PYZcb{}}\PY{p}{,}\PY{+w}{ }\PY{n}{pdf}\PY{p}{)}\PY{p}{;}
\PY{p}{\PYZcb{}}
\end{Verbatim}
\end{codebox}

Like \code{mockStatic}, treat this as a stopgap while working with un-refactorable legacy code, not a pattern to reach for in new code — the sustainable fix is to inject \code{PdfRenderer} (or a factory for it) so tests can simply pass in a mock or fake.

\subsection[BDDMockito]{BDDMockito}
\label{h:20:bddmockito}
\code{BDDMockito} provides \code{given}/\code{willReturn}/\code{then} aliases that read more naturally in Given-When-Then style tests — purely stylistic, functionally identical to the plain Mockito API:

\begin{codebox}
\begin{Verbatim}[commandchars=\\\{\}]
\PY{k+kn}{import static}\PY{+w}{ }\PY{n+nn}{org.mockito.BDDMockito.*}\PY{p}{;}

\PY{n+nd}{@Test}
\PY{k+kt}{void}\PY{+w}{ }\PY{n+nf}{register\PYZus{}newEmail\PYZus{}savesUserAndSendsWelcomeEmail\PYZus{}bdd}\PY{p}{(}\PY{p}{)}\PY{+w}{ }\PY{p}{\PYZob{}}
\PY{+w}{    }\PY{c+c1}{// Given}
\PY{+w}{    }\PY{n}{given}\PY{p}{(}\PY{n}{userRepository}\PY{p}{.}\PY{n+na}{existsByEmail}\PY{p}{(}\PY{l+s}{\PYZdq{}}\PY{l+s}{alice@example.com}\PY{l+s}{\PYZdq{}}\PY{p}{)}\PY{p}{)}\PY{p}{.}\PY{n+na}{willReturn}\PY{p}{(}\PY{k+kc}{false}\PY{p}{)}\PY{p}{;}

\PY{+w}{    }\PY{c+c1}{// When}
\PY{+w}{    }\PY{n}{registrationService}\PY{p}{.}\PY{n+na}{register}\PY{p}{(}\PY{l+s}{\PYZdq{}}\PY{l+s}{alice@example.com}\PY{l+s}{\PYZdq{}}\PY{p}{,}\PY{+w}{ }\PY{l+s}{\PYZdq{}}\PY{l+s}{Alice}\PY{l+s}{\PYZdq{}}\PY{p}{)}\PY{p}{;}

\PY{+w}{    }\PY{c+c1}{// Then}
\PY{+w}{    }\PY{n}{then}\PY{p}{(}\PY{n}{userRepository}\PY{p}{)}\PY{p}{.}\PY{n+na}{should}\PY{p}{(}\PY{p}{)}\PY{p}{.}\PY{n+na}{save}\PY{p}{(}\PY{n}{any}\PY{p}{(}\PY{n}{User}\PY{p}{.}\PY{n+na}{class}\PY{p}{)}\PY{p}{)}\PY{p}{;}
\PY{+w}{    }\PY{n}{then}\PY{p}{(}\PY{n}{emailSender}\PY{p}{)}\PY{p}{.}\PY{n+na}{should}\PY{p}{(}\PY{p}{)}\PY{p}{.}\PY{n+na}{sendWelcomeEmail}\PY{p}{(}\PY{l+s}{\PYZdq{}}\PY{l+s}{alice@example.com}\PY{l+s}{\PYZdq{}}\PY{p}{)}\PY{p}{;}
\PY{p}{\PYZcb{}}
\end{Verbatim}
\end{codebox}

Teams pick one style (plain Mockito or BDDMockito) and stay consistent — mixing \code{when(...)} and \code{given(...)} in the same codebase is a minor but real readability nit reviewers flag.

\subsection[A Full Realistic Example: Service with a Repository and a Clock]{A Full Realistic Example: Service with a Repository and a Clock}
\label{h:20:a-full-realistic-example-service-with-a-repository-and-a-clock}
\begin{codebox}
\begin{Verbatim}[commandchars=\\\{\}]
\PY{c+c1}{// Production code}
\PY{k+kd}{public}\PY{+w}{ }\PY{k+kd}{class} \PY{n+nc}{LoanEligibilityService}\PY{+w}{ }\PY{p}{\PYZob{}}

\PY{+w}{    }\PY{k+kd}{private}\PY{+w}{ }\PY{k+kd}{final}\PY{+w}{ }\PY{n}{LoanRepository}\PY{+w}{ }\PY{n}{loanRepository}\PY{p}{;}
\PY{+w}{    }\PY{k+kd}{private}\PY{+w}{ }\PY{k+kd}{final}\PY{+w}{ }\PY{n}{Clock}\PY{+w}{ }\PY{n}{clock}\PY{p}{;}

\PY{+w}{    }\PY{k+kd}{public}\PY{+w}{ }\PY{n+nf}{LoanEligibilityService}\PY{p}{(}\PY{n}{LoanRepository}\PY{+w}{ }\PY{n}{loanRepository}\PY{p}{,}\PY{+w}{ }\PY{n}{Clock}\PY{+w}{ }\PY{n}{clock}\PY{p}{)}\PY{+w}{ }\PY{p}{\PYZob{}}
\PY{+w}{        }\PY{k}{this}\PY{p}{.}\PY{n+na}{loanRepository}\PY{+w}{ }\PY{o}{=}\PY{+w}{ }\PY{n}{loanRepository}\PY{p}{;}
\PY{+w}{        }\PY{k}{this}\PY{p}{.}\PY{n+na}{clock}\PY{+w}{ }\PY{o}{=}\PY{+w}{ }\PY{n}{clock}\PY{p}{;}
\PY{+w}{    }\PY{p}{\PYZcb{}}

\PY{+w}{    }\PY{k+kd}{public}\PY{+w}{ }\PY{n}{EligibilityResult}\PY{+w}{ }\PY{n+nf}{checkEligibility}\PY{p}{(}\PY{n}{String}\PY{+w}{ }\PY{n}{customerId}\PY{p}{)}\PY{+w}{ }\PY{p}{\PYZob{}}
\PY{+w}{        }\PY{n}{List}\PY{o}{\PYZlt{}}\PY{n}{Loan}\PY{o}{\PYZgt{}}\PY{+w}{ }\PY{n}{activeLoans}\PY{+w}{ }\PY{o}{=}\PY{+w}{ }\PY{n}{loanRepository}\PY{p}{.}\PY{n+na}{findActiveLoansFor}\PY{p}{(}\PY{n}{customerId}\PY{p}{)}\PY{p}{;}

\PY{+w}{        }\PY{k+kt}{boolean}\PY{+w}{ }\PY{n}{hasOverdueLoan}\PY{+w}{ }\PY{o}{=}\PY{+w}{ }\PY{n}{activeLoans}\PY{p}{.}\PY{n+na}{stream}\PY{p}{(}\PY{p}{)}
\PY{+w}{            }\PY{p}{.}\PY{n+na}{anyMatch}\PY{p}{(}\PY{n}{loan}\PY{+w}{ }\PY{o}{\PYZhy{}}\PY{o}{\PYZgt{}}\PY{+w}{ }\PY{n}{loan}\PY{p}{.}\PY{n+na}{dueDate}\PY{p}{(}\PY{p}{)}\PY{p}{.}\PY{n+na}{isBefore}\PY{p}{(}\PY{n}{LocalDate}\PY{p}{.}\PY{n+na}{now}\PY{p}{(}\PY{n}{clock}\PY{p}{)}\PY{p}{)}\PY{p}{)}\PY{p}{;}

\PY{+w}{        }\PY{k}{if}\PY{+w}{ }\PY{p}{(}\PY{n}{hasOverdueLoan}\PY{p}{)}\PY{+w}{ }\PY{p}{\PYZob{}}
\PY{+w}{            }\PY{k}{return}\PY{+w}{ }\PY{n}{EligibilityResult}\PY{p}{.}\PY{n+na}{denied}\PY{p}{(}\PY{l+s}{\PYZdq{}}\PY{l+s}{has an overdue loan}\PY{l+s}{\PYZdq{}}\PY{p}{)}\PY{p}{;}
\PY{+w}{        }\PY{p}{\PYZcb{}}
\PY{+w}{        }\PY{k}{if}\PY{+w}{ }\PY{p}{(}\PY{n}{activeLoans}\PY{p}{.}\PY{n+na}{size}\PY{p}{(}\PY{p}{)}\PY{+w}{ }\PY{o}{\PYZgt{}}\PY{o}{=}\PY{+w}{ }\PY{l+m+mi}{3}\PY{p}{)}\PY{+w}{ }\PY{p}{\PYZob{}}
\PY{+w}{            }\PY{k}{return}\PY{+w}{ }\PY{n}{EligibilityResult}\PY{p}{.}\PY{n+na}{denied}\PY{p}{(}\PY{l+s}{\PYZdq{}}\PY{l+s}{too many active loans}\PY{l+s}{\PYZdq{}}\PY{p}{)}\PY{p}{;}
\PY{+w}{        }\PY{p}{\PYZcb{}}
\PY{+w}{        }\PY{k}{return}\PY{+w}{ }\PY{n}{EligibilityResult}\PY{p}{.}\PY{n+na}{approved}\PY{p}{(}\PY{p}{)}\PY{p}{;}
\PY{+w}{    }\PY{p}{\PYZcb{}}
\PY{p}{\PYZcb{}}
\end{Verbatim}
\end{codebox}

\begin{codebox}
\begin{Verbatim}[commandchars=\\\{\}]
\PY{k+kn}{import}\PY{+w}{ }\PY{n+nn}{org.junit.jupiter.api.Test}\PY{p}{;}
\PY{k+kn}{import}\PY{+w}{ }\PY{n+nn}{org.junit.jupiter.api.extension.ExtendWith}\PY{p}{;}
\PY{k+kn}{import}\PY{+w}{ }\PY{n+nn}{org.mockito.Mock}\PY{p}{;}
\PY{k+kn}{import}\PY{+w}{ }\PY{n+nn}{org.mockito.junit.jupiter.MockitoExtension}\PY{p}{;}

\PY{k+kn}{import}\PY{+w}{ }\PY{n+nn}{java.time.*}\PY{p}{;}
\PY{k+kn}{import}\PY{+w}{ }\PY{n+nn}{java.util.List}\PY{p}{;}

\PY{k+kn}{import static}\PY{+w}{ }\PY{n+nn}{org.mockito.Mockito.*}\PY{p}{;}
\PY{k+kn}{import static}\PY{+w}{ }\PY{n+nn}{org.junit.jupiter.api.Assertions.*}\PY{p}{;}

\PY{n+nd}{@ExtendWith}\PY{p}{(}\PY{n}{MockitoExtension}\PY{p}{.}\PY{n+na}{class}\PY{p}{)}
\PY{k+kd}{class} \PY{n+nc}{LoanEligibilityServiceTest}\PY{+w}{ }\PY{p}{\PYZob{}}

\PY{+w}{    }\PY{k+kd}{private}\PY{+w}{ }\PY{k+kd}{static}\PY{+w}{ }\PY{k+kd}{final}\PY{+w}{ }\PY{n}{Clock}\PY{+w}{ }\PY{n}{FIXED\PYZus{}CLOCK}\PY{+w}{ }\PY{o}{=}
\PY{+w}{        }\PY{n}{Clock}\PY{p}{.}\PY{n+na}{fixed}\PY{p}{(}\PY{n}{Instant}\PY{p}{.}\PY{n+na}{parse}\PY{p}{(}\PY{l+s}{\PYZdq{}}\PY{l+s}{2026\PYZhy{}08\PYZhy{}07T00:00:00Z}\PY{l+s}{\PYZdq{}}\PY{p}{)}\PY{p}{,}\PY{+w}{ }\PY{n}{ZoneOffset}\PY{p}{.}\PY{n+na}{UTC}\PY{p}{)}\PY{p}{;}

\PY{+w}{    }\PY{n+nd}{@Mock}
\PY{+w}{    }\PY{k+kd}{private}\PY{+w}{ }\PY{n}{LoanRepository}\PY{+w}{ }\PY{n}{loanRepository}\PY{p}{;}

\PY{+w}{    }\PY{k+kd}{private}\PY{+w}{ }\PY{n}{LoanEligibilityService}\PY{+w}{ }\PY{n}{service}\PY{p}{;}

\PY{+w}{    }\PY{n+nd}{@org.junit.jupiter.api.BeforeEach}
\PY{+w}{    }\PY{k+kt}{void}\PY{+w}{ }\PY{n+nf}{setUp}\PY{p}{(}\PY{p}{)}\PY{+w}{ }\PY{p}{\PYZob{}}
\PY{+w}{        }\PY{n}{service}\PY{+w}{ }\PY{o}{=}\PY{+w}{ }\PY{k}{new}\PY{+w}{ }\PY{n}{LoanEligibilityService}\PY{p}{(}\PY{n}{loanRepository}\PY{p}{,}\PY{+w}{ }\PY{n}{FIXED\PYZus{}CLOCK}\PY{p}{)}\PY{p}{;}
\PY{+w}{    }\PY{p}{\PYZcb{}}

\PY{+w}{    }\PY{n+nd}{@Test}
\PY{+w}{    }\PY{k+kt}{void}\PY{+w}{ }\PY{n+nf}{checkEligibility\PYZus{}noActiveLoans\PYZus{}isApproved}\PY{p}{(}\PY{p}{)}\PY{+w}{ }\PY{p}{\PYZob{}}
\PY{+w}{        }\PY{n}{when}\PY{p}{(}\PY{n}{loanRepository}\PY{p}{.}\PY{n+na}{findActiveLoansFor}\PY{p}{(}\PY{l+s}{\PYZdq{}}\PY{l+s}{cust\PYZhy{}1}\PY{l+s}{\PYZdq{}}\PY{p}{)}\PY{p}{)}\PY{p}{.}\PY{n+na}{thenReturn}\PY{p}{(}\PY{n}{List}\PY{p}{.}\PY{n+na}{of}\PY{p}{(}\PY{p}{)}\PY{p}{)}\PY{p}{;}

\PY{+w}{        }\PY{n}{EligibilityResult}\PY{+w}{ }\PY{n}{result}\PY{+w}{ }\PY{o}{=}\PY{+w}{ }\PY{n}{service}\PY{p}{.}\PY{n+na}{checkEligibility}\PY{p}{(}\PY{l+s}{\PYZdq{}}\PY{l+s}{cust\PYZhy{}1}\PY{l+s}{\PYZdq{}}\PY{p}{)}\PY{p}{;}

\PY{+w}{        }\PY{n}{assertTrue}\PY{p}{(}\PY{n}{result}\PY{p}{.}\PY{n+na}{approved}\PY{p}{(}\PY{p}{)}\PY{p}{)}\PY{p}{;}
\PY{+w}{    }\PY{p}{\PYZcb{}}

\PY{+w}{    }\PY{n+nd}{@Test}
\PY{+w}{    }\PY{k+kt}{void}\PY{+w}{ }\PY{n+nf}{checkEligibility\PYZus{}hasOverdueLoan\PYZus{}isDenied}\PY{p}{(}\PY{p}{)}\PY{+w}{ }\PY{p}{\PYZob{}}
\PY{+w}{        }\PY{n}{Loan}\PY{+w}{ }\PY{n}{overdue}\PY{+w}{ }\PY{o}{=}\PY{+w}{ }\PY{k}{new}\PY{+w}{ }\PY{n}{Loan}\PY{p}{(}\PY{l+s}{\PYZdq{}}\PY{l+s}{loan\PYZhy{}1}\PY{l+s}{\PYZdq{}}\PY{p}{,}\PY{+w}{ }\PY{n}{LocalDate}\PY{p}{.}\PY{n+na}{parse}\PY{p}{(}\PY{l+s}{\PYZdq{}}\PY{l+s}{2026\PYZhy{}08\PYZhy{}01}\PY{l+s}{\PYZdq{}}\PY{p}{)}\PY{p}{)}\PY{p}{;}\PY{+w}{ }\PY{c+c1}{// before fixed \PYZdq{}now\PYZdq{}}
\PY{+w}{        }\PY{n}{when}\PY{p}{(}\PY{n}{loanRepository}\PY{p}{.}\PY{n+na}{findActiveLoansFor}\PY{p}{(}\PY{l+s}{\PYZdq{}}\PY{l+s}{cust\PYZhy{}1}\PY{l+s}{\PYZdq{}}\PY{p}{)}\PY{p}{)}\PY{p}{.}\PY{n+na}{thenReturn}\PY{p}{(}\PY{n}{List}\PY{p}{.}\PY{n+na}{of}\PY{p}{(}\PY{n}{overdue}\PY{p}{)}\PY{p}{)}\PY{p}{;}

\PY{+w}{        }\PY{n}{EligibilityResult}\PY{+w}{ }\PY{n}{result}\PY{+w}{ }\PY{o}{=}\PY{+w}{ }\PY{n}{service}\PY{p}{.}\PY{n+na}{checkEligibility}\PY{p}{(}\PY{l+s}{\PYZdq{}}\PY{l+s}{cust\PYZhy{}1}\PY{l+s}{\PYZdq{}}\PY{p}{)}\PY{p}{;}

\PY{+w}{        }\PY{n}{assertFalse}\PY{p}{(}\PY{n}{result}\PY{p}{.}\PY{n+na}{approved}\PY{p}{(}\PY{p}{)}\PY{p}{)}\PY{p}{;}
\PY{+w}{        }\PY{n}{assertEquals}\PY{p}{(}\PY{l+s}{\PYZdq{}}\PY{l+s}{has an overdue loan}\PY{l+s}{\PYZdq{}}\PY{p}{,}\PY{+w}{ }\PY{n}{result}\PY{p}{.}\PY{n+na}{reason}\PY{p}{(}\PY{p}{)}\PY{p}{)}\PY{p}{;}
\PY{+w}{    }\PY{p}{\PYZcb{}}

\PY{+w}{    }\PY{n+nd}{@Test}
\PY{+w}{    }\PY{k+kt}{void}\PY{+w}{ }\PY{n+nf}{checkEligibility\PYZus{}threeActiveLoans\PYZus{}isDenied}\PY{p}{(}\PY{p}{)}\PY{+w}{ }\PY{p}{\PYZob{}}
\PY{+w}{        }\PY{n}{List}\PY{o}{\PYZlt{}}\PY{n}{Loan}\PY{o}{\PYZgt{}}\PY{+w}{ }\PY{n}{loans}\PY{+w}{ }\PY{o}{=}\PY{+w}{ }\PY{n}{List}\PY{p}{.}\PY{n+na}{of}\PY{p}{(}
\PY{+w}{            }\PY{k}{new}\PY{+w}{ }\PY{n}{Loan}\PY{p}{(}\PY{l+s}{\PYZdq{}}\PY{l+s}{loan\PYZhy{}1}\PY{l+s}{\PYZdq{}}\PY{p}{,}\PY{+w}{ }\PY{n}{LocalDate}\PY{p}{.}\PY{n+na}{parse}\PY{p}{(}\PY{l+s}{\PYZdq{}}\PY{l+s}{2026\PYZhy{}12\PYZhy{}01}\PY{l+s}{\PYZdq{}}\PY{p}{)}\PY{p}{)}\PY{p}{,}
\PY{+w}{            }\PY{k}{new}\PY{+w}{ }\PY{n}{Loan}\PY{p}{(}\PY{l+s}{\PYZdq{}}\PY{l+s}{loan\PYZhy{}2}\PY{l+s}{\PYZdq{}}\PY{p}{,}\PY{+w}{ }\PY{n}{LocalDate}\PY{p}{.}\PY{n+na}{parse}\PY{p}{(}\PY{l+s}{\PYZdq{}}\PY{l+s}{2026\PYZhy{}12\PYZhy{}01}\PY{l+s}{\PYZdq{}}\PY{p}{)}\PY{p}{)}\PY{p}{,}
\PY{+w}{            }\PY{k}{new}\PY{+w}{ }\PY{n}{Loan}\PY{p}{(}\PY{l+s}{\PYZdq{}}\PY{l+s}{loan\PYZhy{}3}\PY{l+s}{\PYZdq{}}\PY{p}{,}\PY{+w}{ }\PY{n}{LocalDate}\PY{p}{.}\PY{n+na}{parse}\PY{p}{(}\PY{l+s}{\PYZdq{}}\PY{l+s}{2026\PYZhy{}12\PYZhy{}01}\PY{l+s}{\PYZdq{}}\PY{p}{)}\PY{p}{)}
\PY{+w}{        }\PY{p}{)}\PY{p}{;}
\PY{+w}{        }\PY{n}{when}\PY{p}{(}\PY{n}{loanRepository}\PY{p}{.}\PY{n+na}{findActiveLoansFor}\PY{p}{(}\PY{l+s}{\PYZdq{}}\PY{l+s}{cust\PYZhy{}1}\PY{l+s}{\PYZdq{}}\PY{p}{)}\PY{p}{)}\PY{p}{.}\PY{n+na}{thenReturn}\PY{p}{(}\PY{n}{loans}\PY{p}{)}\PY{p}{;}

\PY{+w}{        }\PY{n}{EligibilityResult}\PY{+w}{ }\PY{n}{result}\PY{+w}{ }\PY{o}{=}\PY{+w}{ }\PY{n}{service}\PY{p}{.}\PY{n+na}{checkEligibility}\PY{p}{(}\PY{l+s}{\PYZdq{}}\PY{l+s}{cust\PYZhy{}1}\PY{l+s}{\PYZdq{}}\PY{p}{)}\PY{p}{;}

\PY{+w}{        }\PY{n}{assertFalse}\PY{p}{(}\PY{n}{result}\PY{p}{.}\PY{n+na}{approved}\PY{p}{(}\PY{p}{)}\PY{p}{)}\PY{p}{;}
\PY{+w}{        }\PY{n}{assertEquals}\PY{p}{(}\PY{l+s}{\PYZdq{}}\PY{l+s}{too many active loans}\PY{l+s}{\PYZdq{}}\PY{p}{,}\PY{+w}{ }\PY{n}{result}\PY{p}{.}\PY{n+na}{reason}\PY{p}{(}\PY{p}{)}\PY{p}{)}\PY{p}{;}
\PY{+w}{    }\PY{p}{\PYZcb{}}

\PY{+w}{    }\PY{n+nd}{@Test}
\PY{+w}{    }\PY{k+kt}{void}\PY{+w}{ }\PY{n+nf}{checkEligibility\PYZus{}queriesRepositoryExactlyOnceForTheGivenCustomer}\PY{p}{(}\PY{p}{)}\PY{+w}{ }\PY{p}{\PYZob{}}
\PY{+w}{        }\PY{n}{when}\PY{p}{(}\PY{n}{loanRepository}\PY{p}{.}\PY{n+na}{findActiveLoansFor}\PY{p}{(}\PY{n}{anyString}\PY{p}{(}\PY{p}{)}\PY{p}{)}\PY{p}{)}\PY{p}{.}\PY{n+na}{thenReturn}\PY{p}{(}\PY{n}{List}\PY{p}{.}\PY{n+na}{of}\PY{p}{(}\PY{p}{)}\PY{p}{)}\PY{p}{;}

\PY{+w}{        }\PY{n}{service}\PY{p}{.}\PY{n+na}{checkEligibility}\PY{p}{(}\PY{l+s}{\PYZdq{}}\PY{l+s}{cust\PYZhy{}42}\PY{l+s}{\PYZdq{}}\PY{p}{)}\PY{p}{;}

\PY{+w}{        }\PY{n}{verify}\PY{p}{(}\PY{n}{loanRepository}\PY{p}{,}\PY{+w}{ }\PY{n}{times}\PY{p}{(}\PY{l+m+mi}{1}\PY{p}{)}\PY{p}{)}\PY{p}{.}\PY{n+na}{findActiveLoansFor}\PY{p}{(}\PY{l+s}{\PYZdq{}}\PY{l+s}{cust\PYZhy{}42}\PY{l+s}{\PYZdq{}}\PY{p}{)}\PY{p}{;}
\PY{+w}{        }\PY{n}{verifyNoMoreInteractions}\PY{p}{(}\PY{n}{loanRepository}\PY{p}{)}\PY{p}{;}
\PY{+w}{    }\PY{p}{\PYZcb{}}
\PY{p}{\PYZcb{}}
\end{Verbatim}
\end{codebox}

Notice the fixed \code{Clock} makes the “overdue” scenario completely deterministic — no reliance on the real system date, so this test will pass identically today and five years from now.

\subsection[Over-Mocking as an Anti-Pattern: Prefer a Fake]{Over-Mocking as an Anti-Pattern: Prefer a Fake}
\label{h:20:over-mocking-as-an-anti-pattern-prefer-a-fake}
Mocking every collaborator, including simple, deterministic ones, produces tests that verify \emph{wiring} instead of \emph{behavior} — they pass even when the real logic is subtly wrong, because every answer was hand-fed by \code{when(...)}.

\textbf{Before — over-mocked, brittle, and barely tests anything real:}

\begin{codebox}
\begin{Verbatim}[commandchars=\\\{\}]
\PY{n+nd}{@ExtendWith}\PY{p}{(}\PY{n}{MockitoExtension}\PY{p}{.}\PY{n+na}{class}\PY{p}{)}
\PY{k+kd}{class} \PY{n+nc}{ShoppingCartServiceOverMockedTest}\PY{+w}{ }\PY{p}{\PYZob{}}

\PY{+w}{    }\PY{n+nd}{@Mock}
\PY{+w}{    }\PY{k+kd}{private}\PY{+w}{ }\PY{n}{InventoryRepository}\PY{+w}{ }\PY{n}{inventoryRepository}\PY{p}{;}\PY{+w}{ }\PY{c+c1}{// in\PYZhy{}memory\PYZhy{}friendly; doesn\PYZsq{}t need mocking}

\PY{+w}{    }\PY{n+nd}{@Mock}
\PY{+w}{    }\PY{k+kd}{private}\PY{+w}{ }\PY{n}{PricingRepository}\PY{+w}{ }\PY{n}{pricingRepository}\PY{p}{;}\PY{+w}{      }\PY{c+c1}{// same}

\PY{+w}{    }\PY{n+nd}{@InjectMocks}
\PY{+w}{    }\PY{k+kd}{private}\PY{+w}{ }\PY{n}{ShoppingCartService}\PY{+w}{ }\PY{n}{cartService}\PY{p}{;}

\PY{+w}{    }\PY{n+nd}{@Test}
\PY{+w}{    }\PY{k+kt}{void}\PY{+w}{ }\PY{n+nf}{addItem\PYZus{}computesTotalCorrectly}\PY{p}{(}\PY{p}{)}\PY{+w}{ }\PY{p}{\PYZob{}}
\PY{+w}{        }\PY{n}{when}\PY{p}{(}\PY{n}{inventoryRepository}\PY{p}{.}\PY{n+na}{hasStock}\PY{p}{(}\PY{l+s}{\PYZdq{}}\PY{l+s}{sku\PYZhy{}1}\PY{l+s}{\PYZdq{}}\PY{p}{,}\PY{+w}{ }\PY{l+m+mi}{2}\PY{p}{)}\PY{p}{)}\PY{p}{.}\PY{n+na}{thenReturn}\PY{p}{(}\PY{k+kc}{true}\PY{p}{)}\PY{p}{;}
\PY{+w}{        }\PY{n}{when}\PY{p}{(}\PY{n}{pricingRepository}\PY{p}{.}\PY{n+na}{priceOf}\PY{p}{(}\PY{l+s}{\PYZdq{}}\PY{l+s}{sku\PYZhy{}1}\PY{l+s}{\PYZdq{}}\PY{p}{)}\PY{p}{)}\PY{p}{.}\PY{n+na}{thenReturn}\PY{p}{(}\PY{k}{new}\PY{+w}{ }\PY{n}{BigDecimal}\PY{p}{(}\PY{l+s}{\PYZdq{}}\PY{l+s}{9.99}\PY{l+s}{\PYZdq{}}\PY{p}{)}\PY{p}{)}\PY{p}{;}

\PY{+w}{        }\PY{n}{cartService}\PY{p}{.}\PY{n+na}{addItem}\PY{p}{(}\PY{l+s}{\PYZdq{}}\PY{l+s}{sku\PYZhy{}1}\PY{l+s}{\PYZdq{}}\PY{p}{,}\PY{+w}{ }\PY{l+m+mi}{2}\PY{p}{)}\PY{p}{;}

\PY{+w}{        }\PY{c+c1}{// This assertion mostly re\PYZhy{}proves the mock\PYZsq{}s own stubbed values, not the real}
\PY{+w}{        }\PY{c+c1}{// multiplication logic inside ShoppingCartService — a bug in the multiplication}
\PY{+w}{        }\PY{c+c1}{// could still pass if the stubs happen to line up.}
\PY{+w}{        }\PY{n}{assertEquals}\PY{p}{(}\PY{k}{new}\PY{+w}{ }\PY{n}{BigDecimal}\PY{p}{(}\PY{l+s}{\PYZdq{}}\PY{l+s}{19.98}\PY{l+s}{\PYZdq{}}\PY{p}{)}\PY{p}{,}\PY{+w}{ }\PY{n}{cartService}\PY{p}{.}\PY{n+na}{total}\PY{p}{(}\PY{p}{)}\PY{p}{)}\PY{p}{;}
\PY{+w}{    }\PY{p}{\PYZcb{}}
\PY{p}{\PYZcb{}}
\end{Verbatim}
\end{codebox}

Both \code{InventoryRepository} and \code{PricingRepository} are simple, deterministic, side-effect-free lookups — exactly the kind of collaborator better served by a real (if simplified) in-memory implementation than a mock.

\textbf{After — a real in-memory fake exercises the actual logic:}

\begin{codebox}
\begin{Verbatim}[commandchars=\\\{\}]
\PY{k+kd}{class} \PY{n+nc}{InMemoryInventoryRepository}\PY{+w}{ }\PY{k+kd}{implements}\PY{+w}{ }\PY{n}{InventoryRepository}\PY{+w}{ }\PY{p}{\PYZob{}}
\PY{+w}{    }\PY{k+kd}{private}\PY{+w}{ }\PY{k+kd}{final}\PY{+w}{ }\PY{n}{Map}\PY{o}{\PYZlt{}}\PY{n}{String}\PY{p}{,}\PY{+w}{ }\PY{n}{Integer}\PY{o}{\PYZgt{}}\PY{+w}{ }\PY{n}{stock}\PY{+w}{ }\PY{o}{=}\PY{+w}{ }\PY{k}{new}\PY{+w}{ }\PY{n}{HashMap}\PY{o}{\PYZlt{}}\PY{o}{\PYZgt{}}\PY{p}{(}\PY{p}{)}\PY{p}{;}

\PY{+w}{    }\PY{k+kt}{void}\PY{+w}{ }\PY{n+nf}{setStock}\PY{p}{(}\PY{n}{String}\PY{+w}{ }\PY{n}{sku}\PY{p}{,}\PY{+w}{ }\PY{k+kt}{int}\PY{+w}{ }\PY{n}{quantity}\PY{p}{)}\PY{+w}{ }\PY{p}{\PYZob{}}\PY{+w}{ }\PY{n}{stock}\PY{p}{.}\PY{n+na}{put}\PY{p}{(}\PY{n}{sku}\PY{p}{,}\PY{+w}{ }\PY{n}{quantity}\PY{p}{)}\PY{p}{;}\PY{+w}{ }\PY{p}{\PYZcb{}}

\PY{+w}{    }\PY{n+nd}{@Override}
\PY{+w}{    }\PY{k+kd}{public}\PY{+w}{ }\PY{k+kt}{boolean}\PY{+w}{ }\PY{n+nf}{hasStock}\PY{p}{(}\PY{n}{String}\PY{+w}{ }\PY{n}{sku}\PY{p}{,}\PY{+w}{ }\PY{k+kt}{int}\PY{+w}{ }\PY{n}{quantity}\PY{p}{)}\PY{+w}{ }\PY{p}{\PYZob{}}
\PY{+w}{        }\PY{k}{return}\PY{+w}{ }\PY{n}{stock}\PY{p}{.}\PY{n+na}{getOrDefault}\PY{p}{(}\PY{n}{sku}\PY{p}{,}\PY{+w}{ }\PY{l+m+mi}{0}\PY{p}{)}\PY{+w}{ }\PY{o}{\PYZgt{}}\PY{o}{=}\PY{+w}{ }\PY{n}{quantity}\PY{p}{;}
\PY{+w}{    }\PY{p}{\PYZcb{}}
\PY{p}{\PYZcb{}}

\PY{k+kd}{class} \PY{n+nc}{InMemoryPricingRepository}\PY{+w}{ }\PY{k+kd}{implements}\PY{+w}{ }\PY{n}{PricingRepository}\PY{+w}{ }\PY{p}{\PYZob{}}
\PY{+w}{    }\PY{k+kd}{private}\PY{+w}{ }\PY{k+kd}{final}\PY{+w}{ }\PY{n}{Map}\PY{o}{\PYZlt{}}\PY{n}{String}\PY{p}{,}\PY{+w}{ }\PY{n}{BigDecimal}\PY{o}{\PYZgt{}}\PY{+w}{ }\PY{n}{prices}\PY{+w}{ }\PY{o}{=}\PY{+w}{ }\PY{k}{new}\PY{+w}{ }\PY{n}{HashMap}\PY{o}{\PYZlt{}}\PY{o}{\PYZgt{}}\PY{p}{(}\PY{p}{)}\PY{p}{;}

\PY{+w}{    }\PY{k+kt}{void}\PY{+w}{ }\PY{n+nf}{setPrice}\PY{p}{(}\PY{n}{String}\PY{+w}{ }\PY{n}{sku}\PY{p}{,}\PY{+w}{ }\PY{n}{BigDecimal}\PY{+w}{ }\PY{n}{price}\PY{p}{)}\PY{+w}{ }\PY{p}{\PYZob{}}\PY{+w}{ }\PY{n}{prices}\PY{p}{.}\PY{n+na}{put}\PY{p}{(}\PY{n}{sku}\PY{p}{,}\PY{+w}{ }\PY{n}{price}\PY{p}{)}\PY{p}{;}\PY{+w}{ }\PY{p}{\PYZcb{}}

\PY{+w}{    }\PY{n+nd}{@Override}
\PY{+w}{    }\PY{k+kd}{public}\PY{+w}{ }\PY{n}{BigDecimal}\PY{+w}{ }\PY{n+nf}{priceOf}\PY{p}{(}\PY{n}{String}\PY{+w}{ }\PY{n}{sku}\PY{p}{)}\PY{+w}{ }\PY{p}{\PYZob{}}
\PY{+w}{        }\PY{k}{return}\PY{+w}{ }\PY{n}{prices}\PY{p}{.}\PY{n+na}{getOrDefault}\PY{p}{(}\PY{n}{sku}\PY{p}{,}\PY{+w}{ }\PY{n}{BigDecimal}\PY{p}{.}\PY{n+na}{ZERO}\PY{p}{)}\PY{p}{;}
\PY{+w}{    }\PY{p}{\PYZcb{}}
\PY{p}{\PYZcb{}}

\PY{k+kd}{class} \PY{n+nc}{ShoppingCartServiceFakeTest}\PY{+w}{ }\PY{p}{\PYZob{}}

\PY{+w}{    }\PY{k+kd}{private}\PY{+w}{ }\PY{k+kd}{final}\PY{+w}{ }\PY{n}{InMemoryInventoryRepository}\PY{+w}{ }\PY{n}{inventoryRepository}\PY{+w}{ }\PY{o}{=}\PY{+w}{ }\PY{k}{new}\PY{+w}{ }\PY{n}{InMemoryInventoryRepository}\PY{p}{(}\PY{p}{)}\PY{p}{;}
\PY{+w}{    }\PY{k+kd}{private}\PY{+w}{ }\PY{k+kd}{final}\PY{+w}{ }\PY{n}{InMemoryPricingRepository}\PY{+w}{ }\PY{n}{pricingRepository}\PY{+w}{ }\PY{o}{=}\PY{+w}{ }\PY{k}{new}\PY{+w}{ }\PY{n}{InMemoryPricingRepository}\PY{p}{(}\PY{p}{)}\PY{p}{;}
\PY{+w}{    }\PY{k+kd}{private}\PY{+w}{ }\PY{k+kd}{final}\PY{+w}{ }\PY{n}{ShoppingCartService}\PY{+w}{ }\PY{n}{cartService}\PY{+w}{ }\PY{o}{=}
\PY{+w}{        }\PY{k}{new}\PY{+w}{ }\PY{n}{ShoppingCartService}\PY{p}{(}\PY{n}{inventoryRepository}\PY{p}{,}\PY{+w}{ }\PY{n}{pricingRepository}\PY{p}{)}\PY{p}{;}

\PY{+w}{    }\PY{n+nd}{@Test}
\PY{+w}{    }\PY{k+kt}{void}\PY{+w}{ }\PY{n+nf}{addItem\PYZus{}computesTotalCorrectly}\PY{p}{(}\PY{p}{)}\PY{+w}{ }\PY{p}{\PYZob{}}
\PY{+w}{        }\PY{n}{inventoryRepository}\PY{p}{.}\PY{n+na}{setStock}\PY{p}{(}\PY{l+s}{\PYZdq{}}\PY{l+s}{sku\PYZhy{}1}\PY{l+s}{\PYZdq{}}\PY{p}{,}\PY{+w}{ }\PY{l+m+mi}{5}\PY{p}{)}\PY{p}{;}
\PY{+w}{        }\PY{n}{pricingRepository}\PY{p}{.}\PY{n+na}{setPrice}\PY{p}{(}\PY{l+s}{\PYZdq{}}\PY{l+s}{sku\PYZhy{}1}\PY{l+s}{\PYZdq{}}\PY{p}{,}\PY{+w}{ }\PY{k}{new}\PY{+w}{ }\PY{n}{BigDecimal}\PY{p}{(}\PY{l+s}{\PYZdq{}}\PY{l+s}{9.99}\PY{l+s}{\PYZdq{}}\PY{p}{)}\PY{p}{)}\PY{p}{;}

\PY{+w}{        }\PY{n}{cartService}\PY{p}{.}\PY{n+na}{addItem}\PY{p}{(}\PY{l+s}{\PYZdq{}}\PY{l+s}{sku\PYZhy{}1}\PY{l+s}{\PYZdq{}}\PY{p}{,}\PY{+w}{ }\PY{l+m+mi}{2}\PY{p}{)}\PY{p}{;}

\PY{+w}{        }\PY{n}{assertEquals}\PY{p}{(}\PY{k}{new}\PY{+w}{ }\PY{n}{BigDecimal}\PY{p}{(}\PY{l+s}{\PYZdq{}}\PY{l+s}{19.98}\PY{l+s}{\PYZdq{}}\PY{p}{)}\PY{p}{,}\PY{+w}{ }\PY{n}{cartService}\PY{p}{.}\PY{n+na}{total}\PY{p}{(}\PY{p}{)}\PY{p}{)}\PY{p}{;}
\PY{+w}{    }\PY{p}{\PYZcb{}}

\PY{+w}{    }\PY{n+nd}{@Test}
\PY{+w}{    }\PY{k+kt}{void}\PY{+w}{ }\PY{n+nf}{addItem\PYZus{}insufficientStock\PYZus{}throws}\PY{p}{(}\PY{p}{)}\PY{+w}{ }\PY{p}{\PYZob{}}
\PY{+w}{        }\PY{n}{inventoryRepository}\PY{p}{.}\PY{n+na}{setStock}\PY{p}{(}\PY{l+s}{\PYZdq{}}\PY{l+s}{sku\PYZhy{}1}\PY{l+s}{\PYZdq{}}\PY{p}{,}\PY{+w}{ }\PY{l+m+mi}{1}\PY{p}{)}\PY{p}{;}

\PY{+w}{        }\PY{n}{assertThrows}\PY{p}{(}\PY{n}{InsufficientStockException}\PY{p}{.}\PY{n+na}{class}\PY{p}{,}\PY{+w}{ }\PY{p}{(}\PY{p}{)}\PY{+w}{ }\PY{o}{\PYZhy{}}\PY{o}{\PYZgt{}}\PY{+w}{ }\PY{n}{cartService}\PY{p}{.}\PY{n+na}{addItem}\PY{p}{(}\PY{l+s}{\PYZdq{}}\PY{l+s}{sku\PYZhy{}1}\PY{l+s}{\PYZdq{}}\PY{p}{,}\PY{+w}{ }\PY{l+m+mi}{2}\PY{p}{)}\PY{p}{)}\PY{p}{;}
\PY{+w}{    }\PY{p}{\PYZcb{}}
\PY{p}{\PYZcb{}}
\end{Verbatim}
\end{codebox}

The fake-based test is no more verbose, needs zero mocking annotations, reads like plain Java, and — crucially — actually exercises real lookup logic (\code{getOrDefault}, map storage) instead of blindly returning whatever was stubbed. Reserve mocks for the collaborators that genuinely have side effects (network, disk, email, payment) worth avoiding in a unit test; reach for a small fake for everything else.

{\tablefontsmall
\begin{xltabular}{\linewidth}{@{}L{0.449}L{1.390}L{0.927}L{1.226}L{1.008}@{}}
\toprule
\rowcolor{tablehdr}
\thd{} & \thd{Mock} & \thd{Stub} & \thd{Fake} & \thd{Spy} \\
\midrule
\endhead
Real logic runs? & No & No & Yes (simplified) & Yes, unless overridden \\
Primary purpose & Verify interactions & Feed canned data & Realistic lightweight substitute & Observe + selectively override \\
Created with & \code{mock()} / \code{@\allowbreak{}Mock} & \code{mock()} + \code{when()} & Hand-written class implementing the interface & \code{spy()} / \code{@\allowbreak{}Spy} \\
Review red flag & Excessive \code{verify\allowbreak{}No\allowbreak{}More\allowbreak{}Interactions} on trivial calls & Stubbing something never exercised & None — often \emph{reduces} smell & Spying on the class under test itself \\
\bottomrule
\end{xltabular}
}

\section[Common Code-Review Interview Pitfalls]{Common Code-Review Interview Pitfalls}
\label{h:20:common-code-review-interview-pitfalls}
\begin{enumerate}
\item 
\textbf{Tests with no assertions (“coverage theater”).}
Why it matters: the test executes code and reports green, but proves nothing — a regression can slip through untouched.

\begin{codebox}
\begin{Verbatim}[commandchars=\\\{\}]
\PY{c+c1}{// Before}
\PY{n+nd}{@Test}
\PY{k+kt}{void}\PY{+w}{ }\PY{n+nf}{processOrder\PYZus{}works}\PY{p}{(}\PY{p}{)}\PY{+w}{ }\PY{p}{\PYZob{}}\PY{+w}{ }\PY{n}{orderService}\PY{p}{.}\PY{n+na}{process}\PY{p}{(}\PY{n}{order}\PY{p}{)}\PY{p}{;}\PY{+w}{ }\PY{p}{\PYZcb{}}
\PY{c+c1}{// After}
\PY{n+nd}{@Test}
\PY{k+kt}{void}\PY{+w}{ }\PY{n+nf}{processOrder\PYZus{}marksOrderAsShipped}\PY{p}{(}\PY{p}{)}\PY{+w}{ }\PY{p}{\PYZob{}}
\PY{+w}{    }\PY{n}{orderService}\PY{p}{.}\PY{n+na}{process}\PY{p}{(}\PY{n}{order}\PY{p}{)}\PY{p}{;}
\PY{+w}{    }\PY{n}{assertEquals}\PY{p}{(}\PY{n}{OrderStatus}\PY{p}{.}\PY{n+na}{SHIPPED}\PY{p}{,}\PY{+w}{ }\PY{n}{order}\PY{p}{.}\PY{n+na}{status}\PY{p}{(}\PY{p}{)}\PY{p}{)}\PY{p}{;}
\PY{p}{\PYZcb{}}
\end{Verbatim}
\end{codebox}

\item 
\textbf{Meaningless test names (\code{test1}, \code{testFoo}, \code{shouldWork}).}
Why it matters: a failing test in a CI report should be diagnosable from its name alone; vague names force everyone to open the source to understand what broke.

\begin{codebox}
\begin{Verbatim}[commandchars=\\\{\}]
\PY{c+c1}{// Before}
\PY{n+nd}{@Test}\PY{+w}{ }\PY{k+kt}{void}\PY{+w}{ }\PY{n+nf}{test1}\PY{p}{(}\PY{p}{)}\PY{+w}{ }\PY{p}{\PYZob{}}\PY{+w}{ }\PY{p}{.}\PY{p}{.}\PY{p}{.}\PY{+w}{ }\PY{p}{\PYZcb{}}
\PY{c+c1}{// After}
\PY{n+nd}{@Test}\PY{+w}{ }\PY{k+kt}{void}\PY{+w}{ }\PY{n+nf}{withdraw\PYZus{}insufficientFunds\PYZus{}throwsIllegalStateException}\PY{p}{(}\PY{p}{)}\PY{+w}{ }\PY{p}{\PYZob{}}\PY{+w}{ }\PY{p}{.}\PY{p}{.}\PY{p}{.}\PY{+w}{ }\PY{p}{\PYZcb{}}
\end{Verbatim}
\end{codebox}

\item 
\textbf{Using \code{Thread.\allowbreak{}sleep} to wait for async work instead of a latch or Awaitility.}
Why it matters: fixed sleeps are either wastefully slow or flaky under load — a frequent source of “works on my machine, fails on CI.”

\begin{codebox}
\begin{Verbatim}[commandchars=\\\{\}]
\PY{c+c1}{// Before}
\PY{n}{Thread}\PY{p}{.}\PY{n+na}{sleep}\PY{p}{(}\PY{l+m+mi}{500}\PY{p}{)}\PY{p}{;}\PY{+w}{ }\PY{n}{assertTrue}\PY{p}{(}\PY{n}{flag}\PY{p}{.}\PY{n+na}{get}\PY{p}{(}\PY{p}{)}\PY{p}{)}\PY{p}{;}
\PY{c+c1}{// After}
\PY{n}{assertTrue}\PY{p}{(}\PY{n}{latch}\PY{p}{.}\PY{n+na}{await}\PY{p}{(}\PY{l+m+mi}{2}\PY{p}{,}\PY{+w}{ }\PY{n}{TimeUnit}\PY{p}{.}\PY{n+na}{SECONDS}\PY{p}{)}\PY{p}{)}\PY{p}{;}\PY{+w}{ }\PY{n}{assertTrue}\PY{p}{(}\PY{n}{flag}\PY{p}{.}\PY{n+na}{get}\PY{p}{(}\PY{p}{)}\PY{p}{)}\PY{p}{;}
\end{Verbatim}
\end{codebox}

\item 
\textbf{Calling \code{LocalDate.\allowbreak{}now()}/\code{Instant.\allowbreak{}now()} directly in production code instead of injecting a \code{Clock}.}
Why it matters: makes date/time logic untestable deterministically, leading to tests that only fail once a year or in a different timezone.

\begin{codebox}
\begin{Verbatim}[commandchars=\\\{\}]
\PY{c+c1}{// Before}
\PY{k+kt}{boolean}\PY{+w}{ }\PY{n}{expired}\PY{+w}{ }\PY{o}{=}\PY{+w}{ }\PY{n}{expiry}\PY{p}{.}\PY{n+na}{isBefore}\PY{p}{(}\PY{n}{LocalDate}\PY{p}{.}\PY{n+na}{now}\PY{p}{(}\PY{p}{)}\PY{p}{)}\PY{p}{;}
\PY{c+c1}{// After}
\PY{k+kt}{boolean}\PY{+w}{ }\PY{n}{expired}\PY{+w}{ }\PY{o}{=}\PY{+w}{ }\PY{n}{expiry}\PY{p}{.}\PY{n+na}{isBefore}\PY{p}{(}\PY{n}{LocalDate}\PY{p}{.}\PY{n+na}{now}\PY{p}{(}\PY{n}{clock}\PY{p}{)}\PY{p}{)}\PY{p}{;}\PY{+w}{ }\PY{c+c1}{// clock injected in constructor}
\end{Verbatim}
\end{codebox}

\item 
\textbf{Mixing raw argument values and Mockito matchers in the same call.}
Why it matters: throws \code{Invalid\allowbreak{}Use\allowbreak{}Of\allowbreak{}Matchers\allowbreak{}Exception} at runtime instead of failing at compile time — a confusing failure mode for anyone unfamiliar with the rule.

\begin{codebox}
\begin{Verbatim}[commandchars=\\\{\}]
\PY{c+c1}{// Before}
\PY{n}{when}\PY{p}{(}\PY{n}{userService}\PY{p}{.}\PY{n+na}{updateEmail}\PY{p}{(}\PY{l+m+mi}{1L}\PY{p}{,}\PY{+w}{ }\PY{n}{anyString}\PY{p}{(}\PY{p}{)}\PY{p}{)}\PY{p}{)}\PY{p}{.}\PY{n+na}{thenReturn}\PY{p}{(}\PY{k+kc}{true}\PY{p}{)}\PY{p}{;}
\PY{c+c1}{// After}
\PY{n}{when}\PY{p}{(}\PY{n}{userService}\PY{p}{.}\PY{n+na}{updateEmail}\PY{p}{(}\PY{n}{eq}\PY{p}{(}\PY{l+m+mi}{1L}\PY{p}{)}\PY{p}{,}\PY{+w}{ }\PY{n}{anyString}\PY{p}{(}\PY{p}{)}\PY{p}{)}\PY{p}{)}\PY{p}{.}\PY{n+na}{thenReturn}\PY{p}{(}\PY{k+kc}{true}\PY{p}{)}\PY{p}{;}
\end{Verbatim}
\end{codebox}

\item 
\textbf{Using \code{when(...).\allowbreak{}then\allowbreak{}Return(...)} on a spy for a method with real side effects.}
Why it matters: the real method runs first (before Mockito applies the stub), triggering whatever expensive/exceptional behavior you were trying to avoid.

\begin{codebox}
\begin{Verbatim}[commandchars=\\\{\}]
\PY{c+c1}{// Before}
\PY{n}{when}\PY{p}{(}\PY{n}{spyValidator}\PY{p}{.}\PY{n+na}{expensiveCheck}\PY{p}{(}\PY{n}{order}\PY{p}{)}\PY{p}{)}\PY{p}{.}\PY{n+na}{thenReturn}\PY{p}{(}\PY{k+kc}{true}\PY{p}{)}\PY{p}{;}
\PY{c+c1}{// After}
\PY{n}{doReturn}\PY{p}{(}\PY{k+kc}{true}\PY{p}{)}\PY{p}{.}\PY{n+na}{when}\PY{p}{(}\PY{n}{spyValidator}\PY{p}{)}\PY{p}{.}\PY{n+na}{expensiveCheck}\PY{p}{(}\PY{n}{order}\PY{p}{)}\PY{p}{;}
\end{Verbatim}
\end{codebox}

\item 
\textbf{Mocking value objects/DTOs instead of just constructing them.}
Why it matters: produces an object whose getters return arbitrary stubbed values instead of real data, weakening the test’s connection to reality and adding boilerplate for no benefit.

\begin{codebox}
\begin{Verbatim}[commandchars=\\\{\}]
\PY{c+c1}{// Before}
\PY{n}{Money}\PY{+w}{ }\PY{n}{mockMoney}\PY{+w}{ }\PY{o}{=}\PY{+w}{ }\PY{n}{mock}\PY{p}{(}\PY{n}{Money}\PY{p}{.}\PY{n+na}{class}\PY{p}{)}\PY{p}{;}\PY{+w}{ }\PY{n}{when}\PY{p}{(}\PY{n}{mockMoney}\PY{p}{.}\PY{n+na}{amount}\PY{p}{(}\PY{p}{)}\PY{p}{)}\PY{p}{.}\PY{n+na}{thenReturn}\PY{p}{(}\PY{n}{TEN}\PY{p}{)}\PY{p}{;}
\PY{c+c1}{// After}
\PY{n}{Money}\PY{+w}{ }\PY{n}{money}\PY{+w}{ }\PY{o}{=}\PY{+w}{ }\PY{k}{new}\PY{+w}{ }\PY{n}{Money}\PY{p}{(}\PY{n}{TEN}\PY{p}{,}\PY{+w}{ }\PY{l+s}{\PYZdq{}}\PY{l+s}{USD}\PY{l+s}{\PYZdq{}}\PY{p}{)}\PY{p}{;}
\end{Verbatim}
\end{codebox}

\item 
\textbf{Over-mocking simple, deterministic collaborators instead of using a real in-memory fake.}
Why it matters: the test ends up verifying that stubs return what they were told to return, not that the real logic (calculations, branching) is correct.

\begin{codebox}
\begin{Verbatim}[commandchars=\\\{\}]
\PY{c+c1}{// Before}
\PY{n}{when}\PY{p}{(}\PY{n}{pricingRepository}\PY{p}{.}\PY{n+na}{priceOf}\PY{p}{(}\PY{l+s}{\PYZdq{}}\PY{l+s}{sku\PYZhy{}1}\PY{l+s}{\PYZdq{}}\PY{p}{)}\PY{p}{)}\PY{p}{.}\PY{n+na}{thenReturn}\PY{p}{(}\PY{k}{new}\PY{+w}{ }\PY{n}{BigDecimal}\PY{p}{(}\PY{l+s}{\PYZdq{}}\PY{l+s}{9.99}\PY{l+s}{\PYZdq{}}\PY{p}{)}\PY{p}{)}\PY{p}{;}
\PY{c+c1}{// After}
\PY{n}{pricingRepository}\PY{p}{.}\PY{n+na}{setPrice}\PY{p}{(}\PY{l+s}{\PYZdq{}}\PY{l+s}{sku\PYZhy{}1}\PY{l+s}{\PYZdq{}}\PY{p}{,}\PY{+w}{ }\PY{k}{new}\PY{+w}{ }\PY{n}{BigDecimal}\PY{p}{(}\PY{l+s}{\PYZdq{}}\PY{l+s}{9.99}\PY{l+s}{\PYZdq{}}\PY{p}{)}\PY{p}{)}\PY{p}{;}\PY{+w}{ }\PY{c+c1}{// InMemoryPricingRepository fake}
\end{Verbatim}
\end{codebox}

\item 
\textbf{\code{@\allowbreak{}Disabled} with no reason string or tracked follow-up ticket.}
Why it matters: skipped tests are easy to forget forever, silently eroding coverage without anyone noticing.

\begin{codebox}
\begin{Verbatim}[commandchars=\\\{\}]
\PY{c+c1}{// Before}
\PY{n+nd}{@Disabled}
\PY{n+nd}{@Test}\PY{+w}{ }\PY{k+kt}{void}\PY{+w}{ }\PY{n+nf}{flakyTest}\PY{p}{(}\PY{p}{)}\PY{+w}{ }\PY{p}{\PYZob{}}\PY{+w}{ }\PY{p}{.}\PY{p}{.}\PY{p}{.}\PY{+w}{ }\PY{p}{\PYZcb{}}
\PY{c+c1}{// After}
\PY{n+nd}{@Disabled}\PY{p}{(}\PY{l+s}{\PYZdq{}}\PY{l+s}{Flaky under load; see TICKET\PYZhy{}1234}\PY{l+s}{\PYZdq{}}\PY{p}{)}
\PY{n+nd}{@Test}\PY{+w}{ }\PY{k+kt}{void}\PY{+w}{ }\PY{n+nf}{flakyTest}\PY{p}{(}\PY{p}{)}\PY{+w}{ }\PY{p}{\PYZob{}}\PY{+w}{ }\PY{p}{.}\PY{p}{.}\PY{p}{.}\PY{+w}{ }\PY{p}{\PYZcb{}}
\end{Verbatim}
\end{codebox}

\item 
\textbf{Reaching for \code{RETURNS\_\allowbreak{}DEEP\_\allowbreak{}STUBS} instead of fixing a Law-of-Demeter violation.}
Why it matters: it makes a bad dependency chain (\code{a.\allowbreak{}getB().\allowbreak{}getC().\allowbreak{}getD()}) testable without questioning why the code reaches three objects deep into someone else’s internals.

\begin{codebox}
\begin{Verbatim}[commandchars=\\\{\}]
\PY{c+c1}{// Before}
\PY{n}{when}\PY{p}{(}\PY{n}{context}\PY{p}{.}\PY{n+na}{getCustomer}\PY{p}{(}\PY{p}{)}\PY{p}{.}\PY{n+na}{getAddress}\PY{p}{(}\PY{p}{)}\PY{p}{.}\PY{n+na}{getCountry}\PY{p}{(}\PY{p}{)}\PY{p}{)}\PY{p}{.}\PY{n+na}{thenReturn}\PY{p}{(}\PY{l+s}{\PYZdq{}}\PY{l+s}{DE}\PY{l+s}{\PYZdq{}}\PY{p}{)}\PY{p}{;}
\PY{c+c1}{// After}
\PY{n}{when}\PY{p}{(}\PY{n}{context}\PY{p}{.}\PY{n+na}{getCustomerCountry}\PY{p}{(}\PY{p}{)}\PY{p}{)}\PY{p}{.}\PY{n+na}{thenReturn}\PY{p}{(}\PY{l+s}{\PYZdq{}}\PY{l+s}{DE}\PY{l+s}{\PYZdq{}}\PY{p}{)}\PY{p}{;}\PY{+w}{ }\PY{c+c1}{// narrower, single\PYZhy{}hop dependency}
\end{Verbatim}
\end{codebox}

\item 
\textbf{Asserting only the exception type, ignoring the message/state.}
Why it matters: the wrong code path can throw the same exception type for a different reason, and the test would still pass, masking a real bug.

\begin{codebox}
\begin{Verbatim}[commandchars=\\\{\}]
\PY{c+c1}{// Before}
\PY{n}{assertThrows}\PY{p}{(}\PY{n}{IllegalArgumentException}\PY{p}{.}\PY{n+na}{class}\PY{p}{,}\PY{+w}{ }\PY{p}{(}\PY{p}{)}\PY{+w}{ }\PY{o}{\PYZhy{}}\PY{o}{\PYZgt{}}\PY{+w}{ }\PY{n}{service}\PY{p}{.}\PY{n+na}{transfer}\PY{p}{(}\PY{o}{\PYZhy{}}\PY{l+m+mi}{10}\PY{p}{)}\PY{p}{)}\PY{p}{;}
\PY{c+c1}{// After}
\PY{k+kd}{var}\PY{+w}{ }\PY{n}{ex}\PY{+w}{ }\PY{o}{=}\PY{+w}{ }\PY{n}{assertThrows}\PY{p}{(}\PY{n}{IllegalArgumentException}\PY{p}{.}\PY{n+na}{class}\PY{p}{,}\PY{+w}{ }\PY{p}{(}\PY{p}{)}\PY{+w}{ }\PY{o}{\PYZhy{}}\PY{o}{\PYZgt{}}\PY{+w}{ }\PY{n}{service}\PY{p}{.}\PY{n+na}{transfer}\PY{p}{(}\PY{o}{\PYZhy{}}\PY{l+m+mi}{10}\PY{p}{)}\PY{p}{)}\PY{p}{;}
\PY{n}{assertEquals}\PY{p}{(}\PY{l+s}{\PYZdq{}}\PY{l+s}{transfer amount must be positive: \PYZhy{}10}\PY{l+s}{\PYZdq{}}\PY{p}{,}\PY{+w}{ }\PY{n}{ex}\PY{p}{.}\PY{n+na}{getMessage}\PY{p}{(}\PY{p}{)}\PY{p}{)}\PY{p}{;}
\end{Verbatim}
\end{codebox}

\item 
\textbf{Chasing 100\% line coverage instead of testing edge cases and failure paths.}
Why it matters: coverage percentage measures execution, not verification; time spent testing trivial getters is time not spent on the branch that actually breaks in production.

\begin{codebox}
\begin{Verbatim}[commandchars=\\\{\}]
\PY{c+c1}{// Before: test exists purely to hit the line}
\PY{n+nd}{@Test}\PY{+w}{ }\PY{k+kt}{void}\PY{+w}{ }\PY{n+nf}{getName\PYZus{}doesNotThrow}\PY{p}{(}\PY{p}{)}\PY{+w}{ }\PY{p}{\PYZob{}}\PY{+w}{ }\PY{n}{assertNotNull}\PY{p}{(}\PY{n}{user}\PY{p}{.}\PY{n+na}{getName}\PY{p}{(}\PY{p}{)}\PY{p}{)}\PY{p}{;}\PY{+w}{ }\PY{p}{\PYZcb{}}
\PY{c+c1}{// After: test the behavior that matters}
\PY{n+nd}{@Test}\PY{+w}{ }\PY{k+kt}{void}\PY{+w}{ }\PY{n+nf}{withdraw\PYZus{}exactBalance\PYZus{}leavesZeroBalance}\PY{p}{(}\PY{p}{)}\PY{+w}{ }\PY{p}{\PYZob{}}\PY{+w}{ }\PY{p}{.}\PY{p}{.}\PY{p}{.}\PY{+w}{ }\PY{p}{\PYZcb{}}
\end{Verbatim}
\end{codebox}

\item 
\textbf{Ignoring \code{Unnecessary\allowbreak{}Stubbing\allowbreak{}Exception} by sprinkling \code{lenient()} everywhere.}
Why it matters: strict stubbing exists to catch leftover, copy-pasted stubs that no longer reflect what the test exercises; blanket \code{lenient()} defeats that safety net.

\begin{codebox}
\begin{Verbatim}[commandchars=\\\{\}]
\PY{c+c1}{// Before}
\PY{n}{lenient}\PY{p}{(}\PY{p}{)}\PY{p}{.}\PY{n+na}{when}\PY{p}{(}\PY{n}{repo}\PY{p}{.}\PY{n+na}{findAll}\PY{p}{(}\PY{p}{)}\PY{p}{)}\PY{p}{.}\PY{n+na}{thenReturn}\PY{p}{(}\PY{n}{List}\PY{p}{.}\PY{n+na}{of}\PY{p}{(}\PY{p}{)}\PY{p}{)}\PY{p}{;}\PY{+w}{ }\PY{c+c1}{// added to silence a warning, never investigated}
\PY{c+c1}{// After}
\PY{c+c1}{// delete the unused stub, or confirm and document why it\PYZsq{}s legitimately conditional}
\end{Verbatim}
\end{codebox}

\item 
\textbf{One giant test method covering five unrelated scenarios instead of \code{@\allowbreak{}ParameterizedTest} or separate tests.}
Why it matters: a failure only says “somewhere in this 80-line test something broke,” instead of pinpointing which scenario regressed.

\begin{codebox}
\begin{Verbatim}[commandchars=\\\{\}]
\PY{c+c1}{// Before}
\PY{n+nd}{@Test}\PY{+w}{ }\PY{k+kt}{void}\PY{+w}{ }\PY{n+nf}{allDiscountScenarios}\PY{p}{(}\PY{p}{)}\PY{+w}{ }\PY{p}{\PYZob{}}\PY{+w}{ }\PY{c+cm}{/* 5 unrelated asserts */}\PY{+w}{ }\PY{p}{\PYZcb{}}
\PY{c+c1}{// After}
\PY{n+nd}{@ParameterizedTest}
\PY{n+nd}{@CsvSource}\PY{p}{(}\PY{p}{\PYZob{}}\PY{l+s}{\PYZdq{}}\PY{l+s}{REGULAR,0}\PY{l+s}{\PYZdq{}}\PY{p}{,}\PY{+w}{ }\PY{l+s}{\PYZdq{}}\PY{l+s}{PREMIUM,10}\PY{l+s}{\PYZdq{}}\PY{p}{,}\PY{+w}{ }\PY{l+s}{\PYZdq{}}\PY{l+s}{VIP,20}\PY{l+s}{\PYZdq{}}\PY{p}{\PYZcb{}}\PY{p}{)}
\PY{k+kt}{void}\PY{+w}{ }\PY{n+nf}{discount\PYZus{}variesByTier}\PY{p}{(}\PY{n}{CustomerTier}\PY{+w}{ }\PY{n}{tier}\PY{p}{,}\PY{+w}{ }\PY{k+kt}{int}\PY{+w}{ }\PY{n}{expectedPercent}\PY{p}{)}\PY{+w}{ }\PY{p}{\PYZob{}}\PY{+w}{ }\PY{p}{.}\PY{p}{.}\PY{p}{.}\PY{+w}{ }\PY{p}{\PYZcb{}}
\end{Verbatim}
\end{codebox}

\item 
\textbf{Depending on test execution order (state leaking between tests via static/shared fields).}
Why it matters: JUnit does not guarantee method order; a suite that only passes in one order is a latent flaky-test time bomb waiting for a JUnit version bump or parallel execution to expose it.

\begin{codebox}
\begin{Verbatim}[commandchars=\\\{\}]
\PY{c+c1}{// Before}
\PY{k+kd}{static}\PY{+w}{ }\PY{n}{User}\PY{+w}{ }\PY{n}{sharedUser}\PY{p}{;}\PY{+w}{ }\PY{c+c1}{// set in test A, read in test B}
\PY{c+c1}{// After}
\PY{n+nd}{@BeforeEach}\PY{+w}{ }\PY{k+kt}{void}\PY{+w}{ }\PY{n+nf}{setUp}\PY{p}{(}\PY{p}{)}\PY{+w}{ }\PY{p}{\PYZob{}}\PY{+w}{ }\PY{n}{user}\PY{+w}{ }\PY{o}{=}\PY{+w}{ }\PY{k}{new}\PY{+w}{ }\PY{n}{User}\PY{p}{(}\PY{p}{.}\PY{p}{.}\PY{p}{.}\PY{p}{)}\PY{p}{;}\PY{+w}{ }\PY{p}{\PYZcb{}}\PY{+w}{ }\PY{c+c1}{// fresh per test}
\end{Verbatim}
\end{codebox}

\item 
\textbf{Mocking a type you don’t own (JDK/third-party class) instead of wrapping it behind your own seam.}
Why it matters: couples tests to a library’s internal shape, breaks across library upgrades, and often fights \code{final} classes/methods that resist mocking cleanly.

\begin{codebox}
\begin{Verbatim}[commandchars=\\\{\}]
\PY{c+c1}{// Before}
\PY{n}{HttpClient}\PY{+w}{ }\PY{n}{mockClient}\PY{+w}{ }\PY{o}{=}\PY{+w}{ }\PY{n}{mock}\PY{p}{(}\PY{n}{HttpClient}\PY{p}{.}\PY{n+na}{class}\PY{p}{)}\PY{p}{;}\PY{+w}{ }\PY{c+c1}{// fragile, tightly coupled to the SDK}
\PY{c+c1}{// After}
\PY{n}{WeatherClient}\PY{+w}{ }\PY{n}{mockClient}\PY{+w}{ }\PY{o}{=}\PY{+w}{ }\PY{n}{mock}\PY{p}{(}\PY{n}{WeatherClient}\PY{p}{.}\PY{n+na}{class}\PY{p}{)}\PY{p}{;}\PY{+w}{ }\PY{c+c1}{// your own one\PYZhy{}method interface}
\end{Verbatim}
\end{codebox}

\item 
\textbf{\code{mockStatic}/\code{mockConstruction} used as a first choice for new code instead of dependency injection.}
Why it matters: static mocking is a workaround for code that can’t be injected; reaching for it in new code perpetuates hard-to-test designs instead of fixing them.

\begin{codebox}
\begin{Verbatim}[commandchars=\\\{\}]
\PY{c+c1}{// Before}
\PY{k}{try}\PY{+w}{ }\PY{p}{(}\PY{n}{MockedStatic}\PY{o}{\PYZlt{}}\PY{n}{Instant}\PY{o}{\PYZgt{}}\PY{+w}{ }\PY{n}{m}\PY{+w}{ }\PY{o}{=}\PY{+w}{ }\PY{n}{mockStatic}\PY{p}{(}\PY{n}{Instant}\PY{p}{.}\PY{n+na}{class}\PY{p}{)}\PY{p}{)}\PY{+w}{ }\PY{p}{\PYZob{}}\PY{+w}{ }\PY{n}{m}\PY{p}{.}\PY{n+na}{when}\PY{p}{(}\PY{n}{Instant}\PY{p}{:}\PY{p}{:}\PY{n}{now}\PY{p}{)}\PY{p}{.}\PY{n+na}{thenReturn}\PY{p}{(}\PY{n}{fixed}\PY{p}{)}\PY{p}{;}\PY{+w}{ }\PY{p}{.}\PY{p}{.}\PY{p}{.}\PY{+w}{ }\PY{p}{\PYZcb{}}
\PY{c+c1}{// After}
\PY{c+c1}{// inject java.time.Clock into the constructor and use Clock.fixed(...) in the test}
\end{Verbatim}
\end{codebox}

\item 
\textbf{Blanket \code{verify\allowbreak{}No\allowbreak{}More\allowbreak{}Interactions} on every mock in every test, making refactors brittle.}
Why it matters: adding an unrelated, harmless call (e.g., a new debug log) breaks unrelated tests, training the team to treat verification failures as noise instead of signal.

\begin{codebox}
\begin{Verbatim}[commandchars=\\\{\}]
\PY{c+c1}{// Before}
\PY{n}{verifyNoMoreInteractions}\PY{p}{(}\PY{n}{userRepository}\PY{p}{,}\PY{+w}{ }\PY{n}{emailSender}\PY{p}{,}\PY{+w}{ }\PY{n}{auditLog}\PY{p}{,}\PY{+w}{ }\PY{n}{metricsClient}\PY{p}{)}\PY{p}{;}\PY{+w}{ }\PY{c+c1}{// on every test}
\PY{c+c1}{// After}
\PY{n}{verify}\PY{p}{(}\PY{n}{emailSender}\PY{p}{)}\PY{p}{.}\PY{n+na}{sendWelcomeEmail}\PY{p}{(}\PY{n}{anyString}\PY{p}{(}\PY{p}{)}\PY{p}{)}\PY{p}{;}\PY{+w}{ }\PY{c+c1}{// verify only the interactions that matter to this test}
\end{Verbatim}
\end{codebox}

\end{enumerate}


\setcounter{chapter}{20}
\chapter{Build \& Documentation}
\label{chap:21}
\label{h:21:21-build--documentation}
A reviewer rarely gets to pick the build tool on an existing project, but they are expected to read a \code{pom.\allowbreak{}xml} or \code{build.\allowbreak{}gradle.\allowbreak{}kts} fluently, spot a dependency-scope mistake, and know when a Javadoc comment is padding rather than documentation. This chapter covers Maven and Gradle in depth — enough to read, modify, and debug real build files — and Javadoc, the API-contract documentation every reviewer is expected to enforce. We target Java 21+, Maven 3.9+, and Gradle 8.x throughout.

\section[Maven]{Maven}
\label{h:21:maven}
Maven is a \textbf{declarative, convention-over-configuration} build tool. You describe \emph{what} your project is (its coordinates, dependencies, packaging) in an XML file called \code{pom.\allowbreak{}xml} (Project Object Model), and Maven figures out \emph{how} to build it by running a fixed, well-known \textbf{lifecycle}. Compare this to Gradle, which is closer to a scripting/task-graph model — see the \hyperref[h:21:maven-vs-gradle-comparison]{Maven vs Gradle} table later in this chapter.

\subsection[Convention Over Configuration and the Standard Directory Layout]{Convention Over Configuration and the Standard Directory Layout}
\label{h:21:convention-over-configuration-and-the-standard-directory-layout}
Maven assumes a standard project layout so that a plain \code{pom.\allowbreak{}xml} with no extra configuration already knows where to find source, tests, and resources:

\begin{codebox}
\begin{Verbatim}
my-app/
├── pom.xml
├── src/
│   ├── main/
│   │   ├── java/            # application source (.java)
│   │   └── resources/       # files copied onto the classpath as-is
│   └── test/
│       ├── java/            # test source
│       └── resources/       # test-only resources
└── target/                  # generated output — never commit this
    ├── classes/
    ├── test-classes/
    └── my-app-1.0.0.jar
\end{Verbatim}
\end{codebox}

Because this layout is a convention, plugins do not need to be told where anything is. Deviating from it (e.g., putting sources in \code{source/}) means every plugin that touches source directories must be reconfigured — almost never worth it.

\subsection[A Full Annotated pom.xml]{A Full Annotated \code{pom.\allowbreak{}xml}}
\label{h:21:a-full-annotated-pomxml}
\begin{codebox}
\begin{Verbatim}[commandchars=\\\{\}]
\PY{c+cp}{\PYZlt{}?xml version=\PYZdq{}1.0\PYZdq{} encoding=\PYZdq{}UTF\PYZhy{}8\PYZdq{}?\PYZgt{}}
\PY{n+nt}{\PYZlt{}project}\PY{+w}{ }\PY{n+na}{xmlns=}\PY{l+s}{\PYZdq{}http://maven.apache.org/POM/4.0.0\PYZdq{}}
\PY{+w}{         }\PY{n+na}{xmlns:xsi=}\PY{l+s}{\PYZdq{}http://www.w3.org/2001/XMLSchema\PYZhy{}instance\PYZdq{}}
\PY{+w}{         }\PY{n+na}{xsi:schemaLocation=}\PY{l+s}{\PYZdq{}http://maven.apache.org/POM/4.0.0}
\PY{l+s}{                              http://maven.apache.org/xsd/maven\PYZhy{}4.0.0.xsd\PYZdq{}}\PY{n+nt}{\PYZgt{}}
\PY{+w}{  }\PY{n+nt}{\PYZlt{}modelVersion}\PY{n+nt}{\PYZgt{}}4.0.0\PY{n+nt}{\PYZlt{}/modelVersion\PYZgt{}}

\PY{+w}{  }\PY{c+cm}{\PYZlt{}!\PYZhy{}\PYZhy{} Coordinates: uniquely identify this artifact in any repository \PYZhy{}\PYZhy{}\PYZgt{}}
\PY{+w}{  }\PY{n+nt}{\PYZlt{}groupId}\PY{n+nt}{\PYZgt{}}com.example.orders\PY{n+nt}{\PYZlt{}/groupId\PYZgt{}}
\PY{+w}{  }\PY{n+nt}{\PYZlt{}artifactId}\PY{n+nt}{\PYZgt{}}orders\PYZhy{}service\PY{n+nt}{\PYZlt{}/artifactId\PYZgt{}}
\PY{+w}{  }\PY{n+nt}{\PYZlt{}version}\PY{n+nt}{\PYZgt{}}1.4.0\PYZhy{}SNAPSHOT\PY{n+nt}{\PYZlt{}/version\PYZgt{}}
\PY{+w}{  }\PY{n+nt}{\PYZlt{}packaging}\PY{n+nt}{\PYZgt{}}jar\PY{n+nt}{\PYZlt{}/packaging\PYZgt{}}\PY{+w}{       }\PY{c+cm}{\PYZlt{}!\PYZhy{}\PYZhy{} jar (default), war, pom, ear ... \PYZhy{}\PYZhy{}\PYZgt{}}

\PY{+w}{  }\PY{n+nt}{\PYZlt{}name}\PY{n+nt}{\PYZgt{}}Orders\PY{+w}{ }Service\PY{n+nt}{\PYZlt{}/name\PYZgt{}}
\PY{+w}{  }\PY{n+nt}{\PYZlt{}description}\PY{n+nt}{\PYZgt{}}REST\PY{+w}{ }API\PY{+w}{ }for\PY{+w}{ }managing\PY{+w}{ }customer\PY{+w}{ }orders.\PY{n+nt}{\PYZlt{}/description\PYZgt{}}

\PY{+w}{  }\PY{c+cm}{\PYZlt{}!\PYZhy{}\PYZhy{} Properties: reusable values, also read by many plugins \PYZhy{}\PYZhy{}\PYZgt{}}
\PY{+w}{  }\PY{n+nt}{\PYZlt{}properties}\PY{n+nt}{\PYZgt{}}
\PY{+w}{    }\PY{n+nt}{\PYZlt{}maven.compiler.release}\PY{n+nt}{\PYZgt{}}21\PY{n+nt}{\PYZlt{}/maven.compiler.release\PYZgt{}}
\PY{+w}{    }\PY{n+nt}{\PYZlt{}project.build.sourceEncoding}\PY{n+nt}{\PYZgt{}}UTF\PYZhy{}8\PY{n+nt}{\PYZlt{}/project.build.sourceEncoding\PYZgt{}}
\PY{+w}{    }\PY{n+nt}{\PYZlt{}junit.version}\PY{n+nt}{\PYZgt{}}5.10.2\PY{n+nt}{\PYZlt{}/junit.version\PYZgt{}}
\PY{+w}{  }\PY{n+nt}{\PYZlt{}/properties\PYZgt{}}

\PY{+w}{  }\PY{c+cm}{\PYZlt{}!\PYZhy{}\PYZhy{} dependencyManagement declares *versions/scopes* without adding the}
\PY{c+cm}{       dependency to the build. Children/modules inherit the version but}
\PY{c+cm}{       must still declare the dependency themselves. This is how BOMs work. \PYZhy{}\PYZhy{}\PYZgt{}}
\PY{+w}{  }\PY{n+nt}{\PYZlt{}dependencyManagement}\PY{n+nt}{\PYZgt{}}
\PY{+w}{    }\PY{n+nt}{\PYZlt{}dependencies}\PY{n+nt}{\PYZgt{}}
\PY{+w}{      }\PY{n+nt}{\PYZlt{}dependency}\PY{n+nt}{\PYZgt{}}
\PY{+w}{        }\PY{n+nt}{\PYZlt{}groupId}\PY{n+nt}{\PYZgt{}}org.springframework.boot\PY{n+nt}{\PYZlt{}/groupId\PYZgt{}}
\PY{+w}{        }\PY{n+nt}{\PYZlt{}artifactId}\PY{n+nt}{\PYZgt{}}spring\PYZhy{}boot\PYZhy{}dependencies\PY{n+nt}{\PYZlt{}/artifactId\PYZgt{}}
\PY{+w}{        }\PY{n+nt}{\PYZlt{}version}\PY{n+nt}{\PYZgt{}}3.3.0\PY{n+nt}{\PYZlt{}/version\PYZgt{}}
\PY{+w}{        }\PY{n+nt}{\PYZlt{}type}\PY{n+nt}{\PYZgt{}}pom\PY{n+nt}{\PYZlt{}/type\PYZgt{}}
\PY{+w}{        }\PY{n+nt}{\PYZlt{}scope}\PY{n+nt}{\PYZgt{}}import\PY{n+nt}{\PYZlt{}/scope\PYZgt{}}
\PY{+w}{      }\PY{n+nt}{\PYZlt{}/dependency\PYZgt{}}
\PY{+w}{    }\PY{n+nt}{\PYZlt{}/dependencies\PYZgt{}}
\PY{+w}{  }\PY{n+nt}{\PYZlt{}/dependencyManagement\PYZgt{}}

\PY{+w}{  }\PY{c+cm}{\PYZlt{}!\PYZhy{}\PYZhy{} Actual dependencies used by this module \PYZhy{}\PYZhy{}\PYZgt{}}
\PY{+w}{  }\PY{n+nt}{\PYZlt{}dependencies}\PY{n+nt}{\PYZgt{}}
\PY{+w}{    }\PY{n+nt}{\PYZlt{}dependency}\PY{n+nt}{\PYZgt{}}
\PY{+w}{      }\PY{n+nt}{\PYZlt{}groupId}\PY{n+nt}{\PYZgt{}}org.springframework.boot\PY{n+nt}{\PYZlt{}/groupId\PYZgt{}}
\PY{+w}{      }\PY{n+nt}{\PYZlt{}artifactId}\PY{n+nt}{\PYZgt{}}spring\PYZhy{}boot\PYZhy{}starter\PYZhy{}web\PY{n+nt}{\PYZlt{}/artifactId\PYZgt{}}
\PY{+w}{      }\PY{c+cm}{\PYZlt{}!\PYZhy{}\PYZhy{} no \PYZlt{}version\PYZgt{} needed: it comes from the imported BOM above \PYZhy{}\PYZhy{}\PYZgt{}}
\PY{+w}{    }\PY{n+nt}{\PYZlt{}/dependency\PYZgt{}}

\PY{+w}{    }\PY{n+nt}{\PYZlt{}dependency}\PY{n+nt}{\PYZgt{}}
\PY{+w}{      }\PY{n+nt}{\PYZlt{}groupId}\PY{n+nt}{\PYZgt{}}org.projectlombok\PY{n+nt}{\PYZlt{}/groupId\PYZgt{}}
\PY{+w}{      }\PY{n+nt}{\PYZlt{}artifactId}\PY{n+nt}{\PYZgt{}}lombok\PY{n+nt}{\PYZlt{}/artifactId\PYZgt{}}
\PY{+w}{      }\PY{n+nt}{\PYZlt{}version}\PY{n+nt}{\PYZgt{}}1.18.32\PY{n+nt}{\PYZlt{}/version\PYZgt{}}
\PY{+w}{      }\PY{n+nt}{\PYZlt{}scope}\PY{n+nt}{\PYZgt{}}provided\PY{n+nt}{\PYZlt{}/scope\PYZgt{}}
\PY{+w}{    }\PY{n+nt}{\PYZlt{}/dependency\PYZgt{}}

\PY{+w}{    }\PY{n+nt}{\PYZlt{}dependency}\PY{n+nt}{\PYZgt{}}
\PY{+w}{      }\PY{n+nt}{\PYZlt{}groupId}\PY{n+nt}{\PYZgt{}}org.junit.jupiter\PY{n+nt}{\PYZlt{}/groupId\PYZgt{}}
\PY{+w}{      }\PY{n+nt}{\PYZlt{}artifactId}\PY{n+nt}{\PYZgt{}}junit\PYZhy{}jupiter\PY{n+nt}{\PYZlt{}/artifactId\PYZgt{}}
\PY{+w}{      }\PY{n+nt}{\PYZlt{}version}\PY{n+nt}{\PYZgt{}}\PYZdl{}\PYZob{}junit.version\PYZcb{}\PY{n+nt}{\PYZlt{}/version\PYZgt{}}
\PY{+w}{      }\PY{n+nt}{\PYZlt{}scope}\PY{n+nt}{\PYZgt{}}test\PY{n+nt}{\PYZlt{}/scope\PYZgt{}}
\PY{+w}{    }\PY{n+nt}{\PYZlt{}/dependency\PYZgt{}}
\PY{+w}{  }\PY{n+nt}{\PYZlt{}/dependencies\PYZgt{}}

\PY{+w}{  }\PY{n+nt}{\PYZlt{}build}\PY{n+nt}{\PYZgt{}}
\PY{+w}{    }\PY{n+nt}{\PYZlt{}finalName}\PY{n+nt}{\PYZgt{}}\PYZdl{}\PYZob{}project.artifactId\PYZcb{}\PY{n+nt}{\PYZlt{}/finalName\PYZgt{}}
\PY{+w}{    }\PY{n+nt}{\PYZlt{}plugins}\PY{n+nt}{\PYZgt{}}
\PY{+w}{      }\PY{n+nt}{\PYZlt{}plugin}\PY{n+nt}{\PYZgt{}}
\PY{+w}{        }\PY{n+nt}{\PYZlt{}groupId}\PY{n+nt}{\PYZgt{}}org.apache.maven.plugins\PY{n+nt}{\PYZlt{}/groupId\PYZgt{}}
\PY{+w}{        }\PY{n+nt}{\PYZlt{}artifactId}\PY{n+nt}{\PYZgt{}}maven\PYZhy{}compiler\PYZhy{}plugin\PY{n+nt}{\PYZlt{}/artifactId\PYZgt{}}
\PY{+w}{        }\PY{n+nt}{\PYZlt{}version}\PY{n+nt}{\PYZgt{}}3.13.0\PY{n+nt}{\PYZlt{}/version\PYZgt{}}
\PY{+w}{        }\PY{n+nt}{\PYZlt{}configuration}\PY{n+nt}{\PYZgt{}}
\PY{+w}{          }\PY{n+nt}{\PYZlt{}release}\PY{n+nt}{\PYZgt{}}\PYZdl{}\PYZob{}maven.compiler.release\PYZcb{}\PY{n+nt}{\PYZlt{}/release\PYZgt{}}
\PY{+w}{        }\PY{n+nt}{\PYZlt{}/configuration\PYZgt{}}
\PY{+w}{      }\PY{n+nt}{\PYZlt{}/plugin\PYZgt{}}

\PY{+w}{      }\PY{n+nt}{\PYZlt{}plugin}\PY{n+nt}{\PYZgt{}}
\PY{+w}{        }\PY{n+nt}{\PYZlt{}groupId}\PY{n+nt}{\PYZgt{}}org.apache.maven.plugins\PY{n+nt}{\PYZlt{}/groupId\PYZgt{}}
\PY{+w}{        }\PY{n+nt}{\PYZlt{}artifactId}\PY{n+nt}{\PYZgt{}}maven\PYZhy{}surefire\PYZhy{}plugin\PY{n+nt}{\PYZlt{}/artifactId\PYZgt{}}
\PY{+w}{        }\PY{n+nt}{\PYZlt{}version}\PY{n+nt}{\PYZgt{}}3.2.5\PY{n+nt}{\PYZlt{}/version\PYZgt{}}
\PY{+w}{      }\PY{n+nt}{\PYZlt{}/plugin\PYZgt{}}
\PY{+w}{    }\PY{n+nt}{\PYZlt{}/plugins\PYZgt{}}
\PY{+w}{  }\PY{n+nt}{\PYZlt{}/build\PYZgt{}}

\PY{+w}{  }\PY{c+cm}{\PYZlt{}!\PYZhy{}\PYZhy{} Profiles: conditionally activated configuration \PYZhy{}\PYZhy{}\PYZgt{}}
\PY{+w}{  }\PY{n+nt}{\PYZlt{}profiles}\PY{n+nt}{\PYZgt{}}
\PY{+w}{    }\PY{n+nt}{\PYZlt{}profile}\PY{n+nt}{\PYZgt{}}
\PY{+w}{      }\PY{n+nt}{\PYZlt{}id}\PY{n+nt}{\PYZgt{}}ci\PY{n+nt}{\PYZlt{}/id\PYZgt{}}
\PY{+w}{      }\PY{n+nt}{\PYZlt{}activation}\PY{n+nt}{\PYZgt{}}
\PY{+w}{        }\PY{n+nt}{\PYZlt{}property}\PY{n+nt}{\PYZgt{}}\PY{n+nt}{\PYZlt{}name}\PY{n+nt}{\PYZgt{}}env.CI\PY{n+nt}{\PYZlt{}/name\PYZgt{}}\PY{n+nt}{\PYZlt{}/property\PYZgt{}}
\PY{+w}{      }\PY{n+nt}{\PYZlt{}/activation\PYZgt{}}
\PY{+w}{      }\PY{n+nt}{\PYZlt{}properties}\PY{n+nt}{\PYZgt{}}
\PY{+w}{        }\PY{n+nt}{\PYZlt{}maven.test.failure.ignore}\PY{n+nt}{\PYZgt{}}false\PY{n+nt}{\PYZlt{}/maven.test.failure.ignore\PYZgt{}}
\PY{+w}{      }\PY{n+nt}{\PYZlt{}/properties\PYZgt{}}
\PY{+w}{    }\PY{n+nt}{\PYZlt{}/profile\PYZgt{}}
\PY{+w}{  }\PY{n+nt}{\PYZlt{}/profiles\PYZgt{}}
\PY{n+nt}{\PYZlt{}/project\PYZgt{}}
\end{Verbatim}
\end{codebox}

\subsection[Coordinates: groupId / artifactId / version, SNAPSHOT vs Release]{Coordinates: groupId / artifactId / version, SNAPSHOT vs Release}
\label{h:21:coordinates-groupid--artifactid--version-snapshot-vs-release}
Every artifact in a Maven repository is uniquely addressed by three (sometimes four) coordinates:

{\tablefont
\begin{xltabular}{\linewidth}{@{}L{0.448}L{1.774}L{0.779}@{}}
\toprule
\rowcolor{tablehdr}
\thd{Coordinate} & \thd{Meaning} & \thd{Example} \\
\midrule
\endhead
\code{groupId} & Organization/namespace, usually a reversed domain & \code{com.\allowbreak{}example.\allowbreak{}orders} \\
\code{artifactId} & The specific module/artifact name & \code{orders-\allowbreak{}service} \\
\code{version} & The version of that artifact & \code{1.\allowbreak{}4.\allowbreak{}0}, \code{1.\allowbreak{}4.\allowbreak{}0-\allowbreak{}SNAPSHOT} \\
\code{packaging} & Artifact type (optional, defaults to \code{jar}) & \code{jar}, \code{war}, \code{pom} \\
\code{classifier} & Distinguishes artifacts with the same GAV & \code{sources}, \code{javadoc} \\
\bottomrule
\end{xltabular}
}

\textbf{SNAPSHOT vs release}: a version ending in \code{-\allowbreak{}SNAPSHOT} (e.g., \code{1.\allowbreak{}4.\allowbreak{}0-\allowbreak{}SNAPSHOT}) is a mutable, in-progress version — Maven will re-check the remote repository for a newer snapshot on every build (or per the repository’s update policy) because the artifact behind that version string can change. A version without \code{-\allowbreak{}SNAPSHOT} (e.g., \code{1.\allowbreak{}4.\allowbreak{}0}) is a \textbf{release}: immutable, published once, never overwritten. Deploying the same release version twice to a repository that enforces immutability will fail — that’s intentional, it protects reproducibility. Application \code{pom.\allowbreak{}xml} files should never depend on a \code{-\allowbreak{}SNAPSHOT} artifact when released to production; it means the build is not reproducible.

\subsection[The Build Lifecycle: Phases, Goals, and Plugins]{The Build Lifecycle: Phases, Goals, and Plugins}
\label{h:21:the-build-lifecycle-phases-goals-and-plugins}
Maven has three built-in \textbf{lifecycles}: \code{default} (build and deploy the artifact), \code{clean} (remove \code{target/}), and \code{site} (generate documentation site). The \code{default} lifecycle is the one developers use constantly, and it’s a fixed, ordered sequence of \textbf{phases}:

\begin{codebox}
\begin{Verbatim}
validate → initialize → generate-sources → process-sources → generate-resources
→ process-resources → compile → process-classes → generate-test-sources
→ process-test-sources → generate-test-resources → process-test-resources
→ test-compile → test → prepare-package → package → pre-integration-test
→ integration-test → post-integration-test → verify → install → deploy
\end{Verbatim}
\end{codebox}

The phases most people actually type at the command line:

{\tablefont
\begin{xltabular}{\linewidth}{@{}L{0.393}L{1.607}@{}}
\toprule
\rowcolor{tablehdr}
\thd{Phase} & \thd{What it does} \\
\midrule
\endhead
\code{validate} & Checks the project is correct and all information is available \\
\code{compile} & Compiles main source code \\
\code{test} & Runs unit tests (via Surefire) — does not package anything \\
\code{package} & Packages compiled code into a jar/war \\
\code{verify} & Runs checks on the package, including integration tests (via Failsafe) \\
\code{install} & Installs the package into the \textbf{local} repository (\code{\textasciitilde{}/.\allowbreak{}m2/\allowbreak{}repository}), for use by other local projects \\
\code{deploy} & Copies the package to a \textbf{remote} repository, for sharing with other developers/CI \\
\bottomrule
\end{xltabular}
}

Key mental model: \textbf{running a phase runs every phase before it too}. \code{mvn~\allowbreak{}install} runs validate → … → package → verify → install, in order.

A \textbf{goal} is the actual unit of work — a specific piece of functionality provided by a \textbf{plugin}, written as \code{plugin:\allowbreak{}goal} (e.g., \code{compiler:\allowbreak{}compile}, \code{surefire:\allowbreak{}test}, \code{dependency:\allowbreak{}tree}). Phases don’t \emph{do} anything by themselves; they are just named hooks that plugin goals are \textbf{bound} to.

\begin{codebox}
\begin{Verbatim}
Phase        Bound goal (default binding)          Plugin
---------    ------------------------------------   ---------------------------
compile   →  compiler:compile                    →  maven-compiler-plugin
test      →  surefire:test                       →  maven-surefire-plugin
package   →  jar:jar                             →  maven-jar-plugin
install   →  install:install                     →  maven-install-plugin
deploy    →  deploy:deploy                       →  maven-deploy-plugin
\end{Verbatim}
\end{codebox}

You can also invoke a goal directly without running a whole phase chain, useful for diagnostics:

\begin{codebox}
\begin{Verbatim}[commandchars=\\\{\}]
mvn\PY{+w}{ }dependency:tree
mvn\PY{+w}{ }compiler:compile
mvn\PY{+w}{ }versions:display\PYZhy{}dependency\PYZhy{}updates
\end{Verbatim}
\end{codebox}

\textbf{Plugin vs goal vs phase, in one sentence each:}

\begin{itemize}
\item 
A \textbf{phase} is a named step in the lifecycle (\code{compile}, \code{test}, \code{package}).

\item 
A \textbf{plugin} is a jar containing one or more executable goals (\code{maven-\allowbreak{}surefire-\allowbreak{}plugin}).

\item 
A \textbf{goal} is the actual task a plugin performs, optionally bound to a phase (\code{surefire:\allowbreak{}test} is bound to the \code{test} phase).

\end{itemize}

\subsection[Dependency Scopes]{Dependency Scopes}
\label{h:21:dependency-scopes}
{\tablefont
\begin{xltabular}{\linewidth}{@{}L{0.501}L{0.404}L{0.404}L{0.404}L{0.404}L{3.883}@{}}
\toprule
\rowcolor{tablehdr}
\thd{Scope} & \thd{On compile classpath} & \thd{On test classpath} & \thd{On runtime classpath} & \thd{Packaged in artifact} & \thd{Typical use} \\
\midrule
\endhead
\code{compile} (default) & Yes & Yes & Yes & Yes & Normal library dependency \\
\code{provided} & Yes & Yes & No & No & Servlet API, Lombok — supplied by the container/JDK at runtime \\
\code{runtime} & No & Yes & Yes & Yes & JDBC drivers — needed at runtime, not to compile against \\
\code{test} & No & Yes & No & No & JUnit, Mockito, AssertJ \\
\code{system} & Yes & Yes & No & No & Local jar via explicit \code{\textless{}\allowbreak{}systemPath\textgreater{}} — avoid, not portable \\
\code{import} & n/a & n/a & n/a & n/a & Only valid in \code{\textless{}\allowbreak{}dependencyManagement\textgreater{}}, \code{type=\allowbreak{}pom} — pulls in a BOM’s managed versions \\
\bottomrule
\end{xltabular}
}

\begin{codebox}
\begin{Verbatim}[commandchars=\\\{\}]
\PY{n+nt}{\PYZlt{}dependency}\PY{n+nt}{\PYZgt{}}
\PY{+w}{  }\PY{n+nt}{\PYZlt{}groupId}\PY{n+nt}{\PYZgt{}}org.postgresql\PY{n+nt}{\PYZlt{}/groupId\PYZgt{}}
\PY{+w}{  }\PY{n+nt}{\PYZlt{}artifactId}\PY{n+nt}{\PYZgt{}}postgresql\PY{n+nt}{\PYZlt{}/artifactId\PYZgt{}}
\PY{+w}{  }\PY{n+nt}{\PYZlt{}version}\PY{n+nt}{\PYZgt{}}42.7.3\PY{n+nt}{\PYZlt{}/version\PYZgt{}}
\PY{+w}{  }\PY{n+nt}{\PYZlt{}scope}\PY{n+nt}{\PYZgt{}}runtime\PY{n+nt}{\PYZlt{}/scope\PYZgt{}}
\PY{n+nt}{\PYZlt{}/dependency\PYZgt{}}
\end{Verbatim}
\end{codebox}

A common review flag: marking a dependency \code{compile} when it’s only needed at runtime (bloats the compile classpath and leaks transitively to consumers), or forgetting \code{provided} for a servlet container API that will already be on the classpath when deployed, causing a duplicate/incompatible class at deploy time.

\subsection[Transitive Dependencies, Conflict Resolution, and dependency:tree]{Transitive Dependencies, Conflict Resolution, and \code{dependency:\allowbreak{}tree}}
\label{h:21:transitive-dependencies-conflict-resolution-and-dependencytree}
Maven pulls in \textbf{transitive dependencies} automatically — if your dependency depends on something, you get that something too, without declaring it. This is convenient but is also the \#1 source of “it works on my machine” classpath bugs.

When two versions of the same artifact are pulled in transitively at different depths, Maven uses \textbf{nearest-wins} (also called “nearest definition”) resolution: the dependency declared closest to your project in the dependency graph wins, regardless of version number. If two dependencies are at the \emph{same} depth, the first one declared in the POM wins.

\begin{codebox}
\begin{Verbatim}
your-app
├── lib-a:1.0 → depends on commons-lang3:3.9
└── lib-b:1.0 → depends on commons-lang3:3.14
\end{Verbatim}
\end{codebox}

Here commons-lang3 is at the same depth (both direct children of \code{lib-\allowbreak{}a}/\code{lib-\allowbreak{}b}, which are depth 1, so commons-lang3 is depth 2 either way) — the version resolved is whichever is declared first in your POM’s dependency list. This is exactly why nearest-wins can silently pick an old, buggy, or incompatible version. Never trust it blindly — always inspect the tree:

\begin{codebox}
\begin{Verbatim}[commandchars=\\\{\}]
mvn\PY{+w}{ }dependency:tree
mvn\PY{+w}{ }dependency:tree\PY{+w}{ }\PYZhy{}Dincludes\PY{o}{=}commons\PYZhy{}lang3
mvn\PY{+w}{ }dependency:tree\PY{+w}{ }\PYZhy{}Dverbose\PY{+w}{   }\PY{c+c1}{\PYZsh{} shows the versions that lost, and why}
\end{Verbatim}
\end{codebox}

To force a specific version regardless of what transitive resolution would pick, either declare the dependency directly (direct declarations always win over transitive ones) or use \code{dependencyManagement} (see below), or exclude the unwanted transitive dependency:

\begin{codebox}
\begin{Verbatim}[commandchars=\\\{\}]
\PY{n+nt}{\PYZlt{}dependency}\PY{n+nt}{\PYZgt{}}
\PY{+w}{  }\PY{n+nt}{\PYZlt{}groupId}\PY{n+nt}{\PYZgt{}}com.example\PY{n+nt}{\PYZlt{}/groupId\PYZgt{}}
\PY{+w}{  }\PY{n+nt}{\PYZlt{}artifactId}\PY{n+nt}{\PYZgt{}}lib\PYZhy{}a\PY{n+nt}{\PYZlt{}/artifactId\PYZgt{}}
\PY{+w}{  }\PY{n+nt}{\PYZlt{}version}\PY{n+nt}{\PYZgt{}}1.0\PY{n+nt}{\PYZlt{}/version\PYZgt{}}
\PY{+w}{  }\PY{n+nt}{\PYZlt{}exclusions}\PY{n+nt}{\PYZgt{}}
\PY{+w}{    }\PY{n+nt}{\PYZlt{}exclusion}\PY{n+nt}{\PYZgt{}}
\PY{+w}{      }\PY{n+nt}{\PYZlt{}groupId}\PY{n+nt}{\PYZgt{}}org.apache.commons\PY{n+nt}{\PYZlt{}/groupId\PYZgt{}}
\PY{+w}{      }\PY{n+nt}{\PYZlt{}artifactId}\PY{n+nt}{\PYZgt{}}commons\PYZhy{}lang3\PY{n+nt}{\PYZlt{}/artifactId\PYZgt{}}
\PY{+w}{    }\PY{n+nt}{\PYZlt{}/exclusion\PYZgt{}}
\PY{+w}{  }\PY{n+nt}{\PYZlt{}/exclusions\PYZgt{}}
\PY{n+nt}{\PYZlt{}/dependency\PYZgt{}}
\end{Verbatim}
\end{codebox}

\subsection[dependencyManagement and BOMs]{\code{dependencyManagement} and BOMs}
\label{h:21:dependencymanagement-and-boms}
\code{dependencyManagement} centralizes version (and scope/exclusion) declarations \textbf{without adding the dependency to the build}. A child module or a plugin still needs a bare \code{\textless{}\allowbreak{}dependency\textgreater{}} (no version) to actually pull the artifact in; the version comes from the management block. This is how you avoid version drift across a multi-module project.

A \textbf{BOM} (Bill of Materials) is a \code{pom}-packaged artifact whose sole purpose is a big \code{dependencyManagement} block of mutually-tested versions (Spring Boot’s, Jackson’s, and the AWS SDK’s BOMs are common examples). Importing one gives you a consistent, tested set of versions for an entire ecosystem in one line:

\begin{codebox}
\begin{Verbatim}[commandchars=\\\{\}]
\PY{n+nt}{\PYZlt{}dependencyManagement}\PY{n+nt}{\PYZgt{}}
\PY{+w}{  }\PY{n+nt}{\PYZlt{}dependencies}\PY{n+nt}{\PYZgt{}}
\PY{+w}{    }\PY{n+nt}{\PYZlt{}dependency}\PY{n+nt}{\PYZgt{}}
\PY{+w}{      }\PY{n+nt}{\PYZlt{}groupId}\PY{n+nt}{\PYZgt{}}com.fasterxml.jackson\PY{n+nt}{\PYZlt{}/groupId\PYZgt{}}
\PY{+w}{      }\PY{n+nt}{\PYZlt{}artifactId}\PY{n+nt}{\PYZgt{}}jackson\PYZhy{}bom\PY{n+nt}{\PYZlt{}/artifactId\PYZgt{}}
\PY{+w}{      }\PY{n+nt}{\PYZlt{}version}\PY{n+nt}{\PYZgt{}}2.17.1\PY{n+nt}{\PYZlt{}/version\PYZgt{}}
\PY{+w}{      }\PY{n+nt}{\PYZlt{}type}\PY{n+nt}{\PYZgt{}}pom\PY{n+nt}{\PYZlt{}/type\PYZgt{}}
\PY{+w}{      }\PY{n+nt}{\PYZlt{}scope}\PY{n+nt}{\PYZgt{}}import\PY{n+nt}{\PYZlt{}/scope\PYZgt{}}
\PY{+w}{    }\PY{n+nt}{\PYZlt{}/dependency\PYZgt{}}
\PY{+w}{  }\PY{n+nt}{\PYZlt{}/dependencies\PYZgt{}}
\PY{n+nt}{\PYZlt{}/dependencyManagement\PYZgt{}}

\PY{n+nt}{\PYZlt{}dependencies}\PY{n+nt}{\PYZgt{}}
\PY{+w}{  }\PY{c+cm}{\PYZlt{}!\PYZhy{}\PYZhy{} version omitted on purpose: comes from the imported BOM \PYZhy{}\PYZhy{}\PYZgt{}}
\PY{+w}{  }\PY{n+nt}{\PYZlt{}dependency}\PY{n+nt}{\PYZgt{}}
\PY{+w}{    }\PY{n+nt}{\PYZlt{}groupId}\PY{n+nt}{\PYZgt{}}com.fasterxml.jackson.core\PY{n+nt}{\PYZlt{}/groupId\PYZgt{}}
\PY{+w}{    }\PY{n+nt}{\PYZlt{}artifactId}\PY{n+nt}{\PYZgt{}}jackson\PYZhy{}databind\PY{n+nt}{\PYZlt{}/artifactId\PYZgt{}}
\PY{+w}{  }\PY{n+nt}{\PYZlt{}/dependency\PYZgt{}}
\PY{n+nt}{\PYZlt{}/dependencies\PYZgt{}}
\end{Verbatim}
\end{codebox}

\subsection[Parent POMs and Multi-Module Builds]{Parent POMs and Multi-Module Builds}
\label{h:21:parent-poms-and-multi-module-builds}
A \textbf{parent POM} lets multiple modules share configuration (properties, \code{dependencyManagement}, plugin versions) via inheritance:

\begin{codebox}
\begin{Verbatim}[commandchars=\\\{\}]
\PY{c+cm}{\PYZlt{}!\PYZhy{}\PYZhy{} child module\PYZsq{}s pom.xml \PYZhy{}\PYZhy{}\PYZgt{}}
\PY{n+nt}{\PYZlt{}parent}\PY{n+nt}{\PYZgt{}}
\PY{+w}{  }\PY{n+nt}{\PYZlt{}groupId}\PY{n+nt}{\PYZgt{}}com.example\PY{n+nt}{\PYZlt{}/groupId\PYZgt{}}
\PY{+w}{  }\PY{n+nt}{\PYZlt{}artifactId}\PY{n+nt}{\PYZgt{}}orders\PYZhy{}parent\PY{n+nt}{\PYZlt{}/artifactId\PYZgt{}}
\PY{+w}{  }\PY{n+nt}{\PYZlt{}version}\PY{n+nt}{\PYZgt{}}1.0.0\PY{n+nt}{\PYZlt{}/version\PYZgt{}}
\PY{+w}{  }\PY{n+nt}{\PYZlt{}relativePath}\PY{n+nt}{\PYZgt{}}../pom.xml\PY{n+nt}{\PYZlt{}/relativePath\PYZgt{}}
\PY{n+nt}{\PYZlt{}/parent\PYZgt{}}
\PY{n+nt}{\PYZlt{}artifactId}\PY{n+nt}{\PYZgt{}}orders\PYZhy{}service\PY{n+nt}{\PYZlt{}/artifactId\PYZgt{}}
\end{Verbatim}
\end{codebox}

A \textbf{multi-module} (aggregator) build is a parent \code{pom.\allowbreak{}xml} with \code{packaging=\allowbreak{}pom} and a \code{\textless{}\allowbreak{}modules\textgreater{}} list; building the parent builds every listed module in dependency order:

\begin{codebox}
\begin{Verbatim}[commandchars=\\\{\}]
\PY{n+nt}{\PYZlt{}groupId}\PY{n+nt}{\PYZgt{}}com.example\PY{n+nt}{\PYZlt{}/groupId\PYZgt{}}
\PY{n+nt}{\PYZlt{}artifactId}\PY{n+nt}{\PYZgt{}}orders\PYZhy{}parent\PY{n+nt}{\PYZlt{}/artifactId\PYZgt{}}
\PY{n+nt}{\PYZlt{}version}\PY{n+nt}{\PYZgt{}}1.0.0\PY{n+nt}{\PYZlt{}/version\PYZgt{}}
\PY{n+nt}{\PYZlt{}packaging}\PY{n+nt}{\PYZgt{}}pom\PY{n+nt}{\PYZlt{}/packaging\PYZgt{}}

\PY{n+nt}{\PYZlt{}modules}\PY{n+nt}{\PYZgt{}}
\PY{+w}{  }\PY{n+nt}{\PYZlt{}module}\PY{n+nt}{\PYZgt{}}orders\PYZhy{}api\PY{n+nt}{\PYZlt{}/module\PYZgt{}}
\PY{+w}{  }\PY{n+nt}{\PYZlt{}module}\PY{n+nt}{\PYZgt{}}orders\PYZhy{}service\PY{n+nt}{\PYZlt{}/module\PYZgt{}}
\PY{+w}{  }\PY{n+nt}{\PYZlt{}module}\PY{n+nt}{\PYZgt{}}orders\PYZhy{}persistence\PY{n+nt}{\PYZlt{}/module\PYZgt{}}
\PY{n+nt}{\PYZlt{}/modules\PYZgt{}}
\end{Verbatim}
\end{codebox}

\begin{codebox}
\begin{Verbatim}[commandchars=\\\{\}]
mvn\PY{+w}{ }install\PY{+w}{                          }\PY{c+c1}{\PYZsh{} builds all modules, in dependency order}
mvn\PY{+w}{ }\PYZhy{}pl\PY{+w}{ }orders\PYZhy{}service\PY{+w}{ }\PYZhy{}am\PY{+w}{ }install\PY{+w}{   }\PY{c+c1}{\PYZsh{} build only orders\PYZhy{}service + its dependents}
\end{Verbatim}
\end{codebox}

A module can be both a \emph{child} (inherits from a parent) and independently list its own dependencies — parent and aggregator are separate, orthogonal concepts, even though the same \code{pom.\allowbreak{}xml} often plays both roles.

\subsection[Profiles]{Profiles}
\label{h:21:profiles}
\textbf{Profiles} activate additional configuration conditionally — by an explicit flag, an OS, a JDK version, a file’s presence, or a system/environment property. Common uses: switching database connection properties for local vs CI, enabling static analysis only on CI, or activating a native-image build only on demand.

\begin{codebox}
\begin{Verbatim}[commandchars=\\\{\}]
\PY{n+nt}{\PYZlt{}profiles}\PY{n+nt}{\PYZgt{}}
\PY{+w}{  }\PY{n+nt}{\PYZlt{}profile}\PY{n+nt}{\PYZgt{}}
\PY{+w}{    }\PY{n+nt}{\PYZlt{}id}\PY{n+nt}{\PYZgt{}}integration\PYZhy{}tests\PY{n+nt}{\PYZlt{}/id\PYZgt{}}
\PY{+w}{    }\PY{n+nt}{\PYZlt{}build}\PY{n+nt}{\PYZgt{}}
\PY{+w}{      }\PY{n+nt}{\PYZlt{}plugins}\PY{n+nt}{\PYZgt{}}
\PY{+w}{        }\PY{n+nt}{\PYZlt{}plugin}\PY{n+nt}{\PYZgt{}}
\PY{+w}{          }\PY{n+nt}{\PYZlt{}groupId}\PY{n+nt}{\PYZgt{}}org.apache.maven.plugins\PY{n+nt}{\PYZlt{}/groupId\PYZgt{}}
\PY{+w}{          }\PY{n+nt}{\PYZlt{}artifactId}\PY{n+nt}{\PYZgt{}}maven\PYZhy{}failsafe\PYZhy{}plugin\PY{n+nt}{\PYZlt{}/artifactId\PYZgt{}}
\PY{+w}{          }\PY{n+nt}{\PYZlt{}executions}\PY{n+nt}{\PYZgt{}}
\PY{+w}{            }\PY{n+nt}{\PYZlt{}execution}\PY{n+nt}{\PYZgt{}}
\PY{+w}{              }\PY{n+nt}{\PYZlt{}goals}\PY{n+nt}{\PYZgt{}}\PY{n+nt}{\PYZlt{}goal}\PY{n+nt}{\PYZgt{}}integration\PYZhy{}test\PY{n+nt}{\PYZlt{}/goal\PYZgt{}}\PY{n+nt}{\PYZlt{}goal}\PY{n+nt}{\PYZgt{}}verify\PY{n+nt}{\PYZlt{}/goal\PYZgt{}}\PY{n+nt}{\PYZlt{}/goals\PYZgt{}}
\PY{+w}{            }\PY{n+nt}{\PYZlt{}/execution\PYZgt{}}
\PY{+w}{          }\PY{n+nt}{\PYZlt{}/executions\PYZgt{}}
\PY{+w}{        }\PY{n+nt}{\PYZlt{}/plugin\PYZgt{}}
\PY{+w}{      }\PY{n+nt}{\PYZlt{}/plugins\PYZgt{}}
\PY{+w}{    }\PY{n+nt}{\PYZlt{}/build\PYZgt{}}
\PY{+w}{  }\PY{n+nt}{\PYZlt{}/profile\PYZgt{}}
\PY{n+nt}{\PYZlt{}/profiles\PYZgt{}}
\end{Verbatim}
\end{codebox}

\begin{codebox}
\begin{Verbatim}[commandchars=\\\{\}]
mvn\PY{+w}{ }verify\PY{+w}{ }\PYZhy{}Pintegration\PYZhy{}tests
\end{Verbatim}
\end{codebox}

\subsection[Properties]{Properties}
\label{h:21:properties}
Properties are \code{\$\{\allowbreak{}key\}}-substitutable values, declared in \code{\textless{}\allowbreak{}properties\textgreater{}}, usable anywhere in the POM (versions, plugin config) and readable by many plugins directly (e.g., \code{maven.\allowbreak{}compiler.\allowbreak{}release}, \code{project.\allowbreak{}build.\allowbreak{}source\allowbreak{}Encoding}). Maven also exposes implicit properties like \code{\$\{\allowbreak{}project.\allowbreak{}version\}} and \code{\$\{\allowbreak{}project.\allowbreak{}basedir\}}.

\begin{codebox}
\begin{Verbatim}[commandchars=\\\{\}]
\PY{n+nt}{\PYZlt{}properties}\PY{n+nt}{\PYZgt{}}
\PY{+w}{  }\PY{n+nt}{\PYZlt{}maven.compiler.release}\PY{n+nt}{\PYZgt{}}21\PY{n+nt}{\PYZlt{}/maven.compiler.release\PYZgt{}}
\PY{+w}{  }\PY{n+nt}{\PYZlt{}spring.version}\PY{n+nt}{\PYZgt{}}6.1.8\PY{n+nt}{\PYZlt{}/spring.version\PYZgt{}}
\PY{n+nt}{\PYZlt{}/properties\PYZgt{}}
\end{Verbatim}
\end{codebox}

\subsection[Common Plugins]{Common Plugins}
\label{h:21:common-plugins}
{\tablefontsmall
\begin{xltabular}{\linewidth}{@{}L{0.434}L{1.566}@{}}
\toprule
\rowcolor{tablehdr}
\thd{Plugin} & \thd{Purpose} \\
\midrule
\endhead
\code{maven-\allowbreak{}compiler-\allowbreak{}plugin} & Compiles Java; use \code{\textless{}\allowbreak{}release\textgreater{}\allowbreak{}21\textless{}/\allowbreak{}release\textgreater{}} (not the deprecated \code{source}/\code{target} pair) so both the language level \emph{and} the bootclasspath match Java 21 exactly \\
\code{maven-\allowbreak{}surefire-\allowbreak{}plugin} & Runs \textbf{unit} tests during the \code{test} phase \\
\code{maven-\allowbreak{}failsafe-\allowbreak{}plugin} & Runs \textbf{integration} tests during \code{integration-\allowbreak{}test}/\code{verify}; failures don’t stop the build until \code{verify}, so post-integration-test cleanup still runs \\
\code{maven-\allowbreak{}shade-\allowbreak{}plugin} & Builds an uber/fat jar by unpacking dependency classes and merging them into one jar; supports relocating packages to avoid classpath collisions \\
\code{maven-\allowbreak{}assembly-\allowbreak{}plugin} & Also builds fat jars/distribution archives, but merges by concatenation, not by understanding class content — less safe for merging \code{META-\allowbreak{}INF/\allowbreak{}services} files than shade’s transformers \\
\code{jacoco-\allowbreak{}maven-\allowbreak{}plugin} & Code coverage instrumentation and reporting \\
\code{maven-\allowbreak{}enforcer-\allowbreak{}plugin} & Fails the build if rules are violated — e.g., banned dependencies, minimum Maven/JDK version, no duplicate classes \\
\code{versions-\allowbreak{}maven-\allowbreak{}plugin} & Reports and can update outdated dependency/plugin versions (\code{versions:\allowbreak{}display-\allowbreak{}dependency-\allowbreak{}updates}) \\
\bottomrule
\end{xltabular}
}

\begin{codebox}
\begin{Verbatim}[commandchars=\\\{\}]
\PY{n+nt}{\PYZlt{}plugin}\PY{n+nt}{\PYZgt{}}
\PY{+w}{  }\PY{n+nt}{\PYZlt{}groupId}\PY{n+nt}{\PYZgt{}}org.apache.maven.plugins\PY{n+nt}{\PYZlt{}/groupId\PYZgt{}}
\PY{+w}{  }\PY{n+nt}{\PYZlt{}artifactId}\PY{n+nt}{\PYZgt{}}maven\PYZhy{}compiler\PYZhy{}plugin\PY{n+nt}{\PYZlt{}/artifactId\PYZgt{}}
\PY{+w}{  }\PY{n+nt}{\PYZlt{}configuration}\PY{n+nt}{\PYZgt{}}
\PY{+w}{    }\PY{n+nt}{\PYZlt{}release}\PY{n+nt}{\PYZgt{}}21\PY{n+nt}{\PYZlt{}/release\PYZgt{}}
\PY{+w}{    }\PY{n+nt}{\PYZlt{}parameters}\PY{n+nt}{\PYZgt{}}true\PY{n+nt}{\PYZlt{}/parameters\PYZgt{}}\PY{+w}{   }\PY{c+cm}{\PYZlt{}!\PYZhy{}\PYZhy{} keep parameter names for reflection \PYZhy{}\PYZhy{}\PYZgt{}}
\PY{+w}{  }\PY{n+nt}{\PYZlt{}/configuration\PYZgt{}}
\PY{n+nt}{\PYZlt{}/plugin\PYZgt{}}

\PY{n+nt}{\PYZlt{}plugin}\PY{n+nt}{\PYZgt{}}
\PY{+w}{  }\PY{n+nt}{\PYZlt{}groupId}\PY{n+nt}{\PYZgt{}}org.apache.maven.plugins\PY{n+nt}{\PYZlt{}/groupId\PYZgt{}}
\PY{+w}{  }\PY{n+nt}{\PYZlt{}artifactId}\PY{n+nt}{\PYZgt{}}maven\PYZhy{}shade\PYZhy{}plugin\PY{n+nt}{\PYZlt{}/artifactId\PYZgt{}}
\PY{+w}{  }\PY{n+nt}{\PYZlt{}version}\PY{n+nt}{\PYZgt{}}3.6.0\PY{n+nt}{\PYZlt{}/version\PYZgt{}}
\PY{+w}{  }\PY{n+nt}{\PYZlt{}executions}\PY{n+nt}{\PYZgt{}}
\PY{+w}{    }\PY{n+nt}{\PYZlt{}execution}\PY{n+nt}{\PYZgt{}}
\PY{+w}{      }\PY{n+nt}{\PYZlt{}phase}\PY{n+nt}{\PYZgt{}}package\PY{n+nt}{\PYZlt{}/phase\PYZgt{}}
\PY{+w}{      }\PY{n+nt}{\PYZlt{}goals}\PY{n+nt}{\PYZgt{}}\PY{n+nt}{\PYZlt{}goal}\PY{n+nt}{\PYZgt{}}shade\PY{n+nt}{\PYZlt{}/goal\PYZgt{}}\PY{n+nt}{\PYZlt{}/goals\PYZgt{}}
\PY{+w}{      }\PY{n+nt}{\PYZlt{}configuration}\PY{n+nt}{\PYZgt{}}
\PY{+w}{        }\PY{n+nt}{\PYZlt{}transformers}\PY{n+nt}{\PYZgt{}}
\PY{+w}{          }\PY{n+nt}{\PYZlt{}transformer}\PY{+w}{ }\PY{n+na}{implementation=}\PY{l+s}{\PYZdq{}org.apache.maven.plugins.shade.resource.ManifestResourceTransformer\PYZdq{}}\PY{n+nt}{\PYZgt{}}
\PY{+w}{            }\PY{n+nt}{\PYZlt{}mainClass}\PY{n+nt}{\PYZgt{}}com.example.orders.Main\PY{n+nt}{\PYZlt{}/mainClass\PYZgt{}}
\PY{+w}{          }\PY{n+nt}{\PYZlt{}/transformer\PYZgt{}}
\PY{+w}{        }\PY{n+nt}{\PYZlt{}/transformers\PYZgt{}}
\PY{+w}{      }\PY{n+nt}{\PYZlt{}/configuration\PYZgt{}}
\PY{+w}{    }\PY{n+nt}{\PYZlt{}/execution\PYZgt{}}
\PY{+w}{  }\PY{n+nt}{\PYZlt{}/executions\PYZgt{}}
\PY{n+nt}{\PYZlt{}/plugin\PYZgt{}}

\PY{n+nt}{\PYZlt{}plugin}\PY{n+nt}{\PYZgt{}}
\PY{+w}{  }\PY{n+nt}{\PYZlt{}groupId}\PY{n+nt}{\PYZgt{}}org.apache.maven.plugins\PY{n+nt}{\PYZlt{}/groupId\PYZgt{}}
\PY{+w}{  }\PY{n+nt}{\PYZlt{}artifactId}\PY{n+nt}{\PYZgt{}}maven\PYZhy{}enforcer\PYZhy{}plugin\PY{n+nt}{\PYZlt{}/artifactId\PYZgt{}}
\PY{+w}{  }\PY{n+nt}{\PYZlt{}executions}\PY{n+nt}{\PYZgt{}}
\PY{+w}{    }\PY{n+nt}{\PYZlt{}execution}\PY{n+nt}{\PYZgt{}}
\PY{+w}{      }\PY{n+nt}{\PYZlt{}goals}\PY{n+nt}{\PYZgt{}}\PY{n+nt}{\PYZlt{}goal}\PY{n+nt}{\PYZgt{}}enforce\PY{n+nt}{\PYZlt{}/goal\PYZgt{}}\PY{n+nt}{\PYZlt{}/goals\PYZgt{}}
\PY{+w}{      }\PY{n+nt}{\PYZlt{}configuration}\PY{n+nt}{\PYZgt{}}
\PY{+w}{        }\PY{n+nt}{\PYZlt{}rules}\PY{n+nt}{\PYZgt{}}
\PY{+w}{          }\PY{n+nt}{\PYZlt{}requireMavenVersion}\PY{n+nt}{\PYZgt{}}\PY{n+nt}{\PYZlt{}version}\PY{n+nt}{\PYZgt{}}[3.9,)\PY{n+nt}{\PYZlt{}/version\PYZgt{}}\PY{n+nt}{\PYZlt{}/requireMavenVersion\PYZgt{}}
\PY{+w}{          }\PY{n+nt}{\PYZlt{}requireJavaVersion}\PY{n+nt}{\PYZgt{}}\PY{n+nt}{\PYZlt{}version}\PY{n+nt}{\PYZgt{}}[21,)\PY{n+nt}{\PYZlt{}/version\PYZgt{}}\PY{n+nt}{\PYZlt{}/requireJavaVersion\PYZgt{}}
\PY{+w}{          }\PY{n+nt}{\PYZlt{}bannedDependencies}\PY{n+nt}{\PYZgt{}}
\PY{+w}{            }\PY{n+nt}{\PYZlt{}excludes}\PY{n+nt}{\PYZgt{}}\PY{n+nt}{\PYZlt{}exclude}\PY{n+nt}{\PYZgt{}}log4j:log4j\PY{n+nt}{\PYZlt{}/exclude\PYZgt{}}\PY{n+nt}{\PYZlt{}/excludes\PYZgt{}}
\PY{+w}{          }\PY{n+nt}{\PYZlt{}/bannedDependencies\PYZgt{}}
\PY{+w}{        }\PY{n+nt}{\PYZlt{}/rules\PYZgt{}}
\PY{+w}{      }\PY{n+nt}{\PYZlt{}/configuration\PYZgt{}}
\PY{+w}{    }\PY{n+nt}{\PYZlt{}/execution\PYZgt{}}
\PY{+w}{  }\PY{n+nt}{\PYZlt{}/executions\PYZgt{}}
\PY{n+nt}{\PYZlt{}/plugin\PYZgt{}}
\end{Verbatim}
\end{codebox}

\subsection[The Local Repository and Offline Mode]{The Local Repository and Offline Mode}
\label{h:21:the-local-repository-and-offline-mode}
Maven caches every artifact it downloads under \code{\textasciitilde{}/.\allowbreak{}m2/\allowbreak{}repository}, keyed by GAV coordinates. \code{mvn~\allowbreak{}install} places your own module’s artifact there too, so sibling projects can depend on it without a remote repository round-trip.

\textbf{Reproducible builds}: pin exact versions (avoid version ranges like \code{[\allowbreak{}1.\allowbreak{}0,\allowbreak{}2.\allowbreak{}0)}, which resolve differently over time), pin plugin versions explicitly (an unpinned plugin resolves to “latest” the first time it’s used and then locks to whatever was cached — inconsistent across machines/CI), and prefer \code{-\allowbreak{}SNAPSHOT}-free dependencies for anything released. Maven also supports a \code{reproducible} build mode via \code{\textless{}\allowbreak{}project.\allowbreak{}build.\allowbreak{}output\allowbreak{}Timestamp\textgreater{}} so two builds of the same source produce byte-identical jars (same file timestamps inside the archive).

\begin{codebox}
\begin{Verbatim}[commandchars=\\\{\}]
mvn\PY{+w}{ }\PYZhy{}o\PY{+w}{ }package\PY{+w}{        }\PY{c+c1}{\PYZsh{} offline: fail fast if anything is missing from \PYZti{}/.m2, don\PYZsq{}t hit the network}
mvn\PY{+w}{ }\PYZhy{}o\PY{+w}{ }dependency:tree\PY{+w}{ }\PY{c+c1}{\PYZsh{} useful to sanity\PYZhy{}check nothing is silently redownloaded}
\end{Verbatim}
\end{codebox}

\subsection[Skipping Tests: -DskipTests vs -Dmaven.test.skip]{Skipping Tests: \code{-\allowbreak{}DskipTests} vs \code{-\allowbreak{}Dmaven.\allowbreak{}test.\allowbreak{}skip}}
\label{h:21:skipping-tests--dskiptests-vs--dmaventestskip}
These are \textbf{not} the same thing, and mixing them up is a classic review flag:

{\tablefont
\begin{xltabular}{\linewidth}{@{}L{1.200}L{1.200}L{0.600}@{}}
\toprule
\rowcolor{tablehdr}
\thd{Flag} & \thd{Compiles test sources?} & \thd{Runs tests?} \\
\midrule
\endhead
\code{-\allowbreak{}DskipTests} & Yes & No \\
\code{-\allowbreak{}Dmaven.\allowbreak{}test.\allowbreak{}skip=\allowbreak{}true} & No & No \\
\bottomrule
\end{xltabular}
}

\begin{codebox}
\begin{Verbatim}[commandchars=\\\{\}]
mvn\PY{+w}{ }install\PY{+w}{ }\PYZhy{}DskipTests\PY{+w}{            }\PY{c+c1}{\PYZsh{} tests are compiled (catches compile errors) but not run}
mvn\PY{+w}{ }install\PY{+w}{ }\PYZhy{}Dmaven.test.skip\PY{o}{=}\PY{n+nb}{true}\PY{+w}{ }\PY{c+c1}{\PYZsh{} test sources aren\PYZsq{}t even compiled — fastest, but hides broken tests}
\end{Verbatim}
\end{codebox}

Using \code{-\allowbreak{}Dmaven.\allowbreak{}test.\allowbreak{}skip=\allowbreak{}true} in CI can let a broken test file merge unnoticed, since it never even compiles until someone runs tests locally.

\subsection[The Maven Wrapper]{The Maven Wrapper}
\label{h:21:the-maven-wrapper}
The \textbf{Maven Wrapper} (\code{mvnw} / \code{mvnw.\allowbreak{}cmd}) pins the exact Maven version a project needs, generated once with \code{mvn~\allowbreak{}wrapper:\allowbreak{}wrapper} and committed to the repository. Anyone cloning the repo runs \code{./\allowbreak{}mvnw} and gets a reproducible Maven version automatically downloaded — no “works with Maven 3.9 but the CI box has 3.6” surprises.

\begin{codebox}
\begin{Verbatim}[commandchars=\\\{\}]
mvn\PY{+w}{ }wrapper:wrapper\PY{+w}{ }\PYZhy{}Dmaven\PY{o}{=}\PY{l+m}{3}.9.6\PY{+w}{   }\PY{c+c1}{\PYZsh{} generate the wrapper, pin version 3.9.6}
./mvnw\PY{+w}{ }clean\PY{+w}{ }verify\PY{+w}{                  }\PY{c+c1}{\PYZsh{} anyone can build without installing Maven at all}
\end{Verbatim}
\end{codebox}

Commit \code{mvnw}, \code{mvnw.\allowbreak{}cmd}, and \code{.\allowbreak{}mvn/\allowbreak{}wrapper/\allowbreak{}maven-\allowbreak{}wrapper.\allowbreak{}properties} — never \code{.\allowbreak{}gitignore} them.

\subsection[Command Cheat-Sheet]{Command Cheat-Sheet}
\label{h:21:command-cheat-sheet}
\begin{codebox}
\begin{Verbatim}[commandchars=\\\{\}]
mvn\PY{+w}{ }clean\PY{+w}{                         }\PY{c+c1}{\PYZsh{} delete target/}
mvn\PY{+w}{ }compile\PY{+w}{                       }\PY{c+c1}{\PYZsh{} compile main sources}
mvn\PY{+w}{ }\PY{n+nb}{test}\PY{+w}{                          }\PY{c+c1}{\PYZsh{} run unit tests}
mvn\PY{+w}{ }package\PY{+w}{                       }\PY{c+c1}{\PYZsh{} produce jar/war}
mvn\PY{+w}{ }verify\PY{+w}{                        }\PY{c+c1}{\PYZsh{} run integration tests + checks}
mvn\PY{+w}{ }install\PY{+w}{                       }\PY{c+c1}{\PYZsh{} install into \PYZti{}/.m2}
mvn\PY{+w}{ }deploy\PY{+w}{                        }\PY{c+c1}{\PYZsh{} publish to remote repository}

mvn\PY{+w}{ }dependency:tree\PY{+w}{               }\PY{c+c1}{\PYZsh{} show resolved dependency graph}
mvn\PY{+w}{ }dependency:analyze\PY{+w}{            }\PY{c+c1}{\PYZsh{} find unused/undeclared direct dependencies}
mvn\PY{+w}{ }versions:display\PYZhy{}dependency\PYZhy{}updates\PY{+w}{   }\PY{c+c1}{\PYZsh{} find outdated dependency versions}
mvn\PY{+w}{ }help:effective\PYZhy{}pom\PY{+w}{            }\PY{c+c1}{\PYZsh{} print the fully resolved/merged POM}

mvn\PY{+w}{ }\PYZhy{}pl\PY{+w}{ }module\PYZhy{}a\PY{+w}{ }\PYZhy{}am\PY{+w}{ }install\PY{+w}{       }\PY{c+c1}{\PYZsh{} build module\PYZhy{}a plus everything it depends on}
mvn\PY{+w}{ }\PYZhy{}pl\PY{+w}{ }module\PYZhy{}a\PY{+w}{ }\PYZhy{}amd\PY{+w}{ }install\PY{+w}{       }\PY{c+c1}{\PYZsh{} build module\PYZhy{}a plus everything that depends on it}
mvn\PY{+w}{ }\PYZhy{}T\PY{+w}{ }\PY{l+m}{4}\PY{+w}{ }install\PY{+w}{                   }\PY{c+c1}{\PYZsh{} build with 4 threads (parallel modules)}
mvn\PY{+w}{ }\PYZhy{}X\PY{+w}{ }verify\PY{+w}{                      }\PY{c+c1}{\PYZsh{} debug/verbose logging}
mvn\PY{+w}{ }\PYZhy{}q\PY{+w}{ }verify\PY{+w}{                      }\PY{c+c1}{\PYZsh{} quiet output, errors only}
\end{Verbatim}
\end{codebox}

\section[Gradle]{Gradle}
\label{h:21:gradle}
Gradle is a general-purpose build tool built around a \textbf{task graph} rather than a fixed lifecycle: you (or a plugin) define \textbf{tasks} and their dependencies, and Gradle figures out which tasks need to run and in what order, skipping anything whose inputs haven’t changed. Build scripts are code (Groovy or Kotlin), not declarative XML, which makes Gradle more flexible but also easier to write build logic that’s hard to reason about — a common review concern.

\subsection[Groovy DSL vs Kotlin DSL]{Groovy DSL vs Kotlin DSL}
\label{h:21:groovy-dsl-vs-kotlin-dsl}
Both DSLs configure the same underlying model. Kotlin DSL (\code{build.\allowbreak{}gradle.\allowbreak{}kts}) gets full IDE type-checking and autocompletion; Groovy DSL (\code{build.\allowbreak{}gradle}) is more concise and was the historical default. Since Gradle 8.x, Kotlin DSL is the recommended default for new projects.

\begin{codebox}
\begin{Verbatim}[commandchars=\\\{\}]
\PY{c+c1}{// build.gradle (Groovy DSL)}
\PY{n}{plugins}\PY{+w}{ }\PY{o}{\PYZob{}}
\PY{+w}{    }\PY{n}{id}\PY{+w}{ }\PY{l+s+s1}{\PYZsq{}java\PYZsq{}}
\PY{+w}{    }\PY{n}{id}\PY{+w}{ }\PY{l+s+s1}{\PYZsq{}application\PYZsq{}}
\PY{o}{\PYZcb{}}

\PY{n}{repositories}\PY{+w}{ }\PY{o}{\PYZob{}}
\PY{+w}{    }\PY{n}{mavenCentral}\PY{o}{(}\PY{o}{)}
\PY{o}{\PYZcb{}}

\PY{n}{dependencies}\PY{+w}{ }\PY{o}{\PYZob{}}
\PY{+w}{    }\PY{n}{implementation}\PY{+w}{ }\PY{l+s+s1}{\PYZsq{}com.google.guava:guava:33.2.1\PYZhy{}jre\PYZsq{}}
\PY{+w}{    }\PY{n}{testImplementation}\PY{+w}{ }\PY{n}{platform}\PY{o}{(}\PY{l+s+s1}{\PYZsq{}org.junit:junit\PYZhy{}bom:5.10.2\PYZsq{}}\PY{o}{)}
\PY{+w}{    }\PY{n}{testImplementation}\PY{+w}{ }\PY{l+s+s1}{\PYZsq{}org.junit.jupiter:junit\PYZhy{}jupiter\PYZsq{}}
\PY{o}{\PYZcb{}}

\PY{n}{java}\PY{+w}{ }\PY{o}{\PYZob{}}
\PY{+w}{    }\PY{n}{toolchain}\PY{+w}{ }\PY{o}{\PYZob{}}
\PY{+w}{        }\PY{n}{languageVersion}\PY{+w}{ }\PY{o}{=}\PY{+w}{ }\PY{n}{JavaLanguageVersion}\PY{o}{.}\PY{n+na}{of}\PY{o}{(}\PY{l+m+mi}{21}\PY{o}{)}
\PY{+w}{    }\PY{o}{\PYZcb{}}
\PY{o}{\PYZcb{}}

\PY{n}{application}\PY{+w}{ }\PY{o}{\PYZob{}}
\PY{+w}{    }\PY{n}{mainClass}\PY{+w}{ }\PY{o}{=}\PY{+w}{ }\PY{l+s+s1}{\PYZsq{}com.example.App\PYZsq{}}
\PY{o}{\PYZcb{}}

\PY{n}{test}\PY{+w}{ }\PY{o}{\PYZob{}}
\PY{+w}{    }\PY{n}{useJUnitPlatform}\PY{o}{(}\PY{o}{)}
\PY{o}{\PYZcb{}}
\end{Verbatim}
\end{codebox}

\begin{codebox}
\begin{Verbatim}[commandchars=\\\{\}]
\PY{c+c1}{// build.gradle.kts (Kotlin DSL) — same build, statically typed}
\PY{n}{plugins}\PY{+w}{ }\PY{p}{\PYZob{}}
\PY{+w}{    }\PY{n}{java}
\PY{+w}{    }\PY{n}{application}
\PY{p}{\PYZcb{}}

\PY{n}{repositories}\PY{+w}{ }\PY{p}{\PYZob{}}
\PY{+w}{    }\PY{n}{mavenCentral}\PY{p}{(}\PY{p}{)}
\PY{p}{\PYZcb{}}

\PY{n}{dependencies}\PY{+w}{ }\PY{p}{\PYZob{}}
\PY{+w}{    }\PY{n}{implementation}\PY{p}{(}\PY{l+s}{\PYZdq{}}\PY{l+s}{com.google.guava:guava:33.2.1\PYZhy{}jre}\PY{l+s}{\PYZdq{}}\PY{p}{)}
\PY{+w}{    }\PY{n}{testImplementation}\PY{p}{(}\PY{n}{platform}\PY{p}{(}\PY{l+s}{\PYZdq{}}\PY{l+s}{org.junit:junit\PYZhy{}bom:5.10.2}\PY{l+s}{\PYZdq{}}\PY{p}{)}\PY{p}{)}
\PY{+w}{    }\PY{n}{testImplementation}\PY{p}{(}\PY{l+s}{\PYZdq{}}\PY{l+s}{org.junit.jupiter:junit\PYZhy{}jupiter}\PY{l+s}{\PYZdq{}}\PY{p}{)}
\PY{p}{\PYZcb{}}

\PY{n}{java}\PY{+w}{ }\PY{p}{\PYZob{}}
\PY{+w}{    }\PY{n}{toolchain}\PY{+w}{ }\PY{p}{\PYZob{}}
\PY{+w}{        }\PY{n}{languageVersion}\PY{p}{.}\PY{n+na}{set}\PY{p}{(}\PY{n}{JavaLanguageVersion}\PY{p}{.}\PY{n+na}{of}\PY{p}{(}\PY{l+m}{2}\PY{l+m}{1}\PY{p}{)}\PY{p}{)}
\PY{+w}{    }\PY{p}{\PYZcb{}}
\PY{p}{\PYZcb{}}

\PY{n}{application}\PY{+w}{ }\PY{p}{\PYZob{}}
\PY{+w}{    }\PY{n}{mainClass}\PY{p}{.}\PY{n+na}{set}\PY{p}{(}\PY{l+s}{\PYZdq{}}\PY{l+s}{com.example.App}\PY{l+s}{\PYZdq{}}\PY{p}{)}
\PY{p}{\PYZcb{}}

\PY{n}{tasks}\PY{p}{.}\PY{n+na}{test}\PY{+w}{ }\PY{p}{\PYZob{}}
\PY{+w}{    }\PY{n}{useJUnitPlatform}\PY{p}{(}\PY{p}{)}
\PY{p}{\PYZcb{}}
\end{Verbatim}
\end{codebox}

\subsection[An Annotated build.gradle.kts]{An Annotated \code{build.\allowbreak{}gradle.\allowbreak{}kts}}
\label{h:21:an-annotated-buildgradlekts}
\begin{codebox}
\begin{Verbatim}[commandchars=\\\{\}]
\PY{n}{plugins}\PY{+w}{ }\PY{p}{\PYZob{}}
\PY{+w}{    }\PY{n}{id}\PY{p}{(}\PY{l+s}{\PYZdq{}}\PY{l+s}{java\PYZhy{}library}\PY{l+s}{\PYZdq{}}\PY{p}{)}\PY{+w}{               }\PY{c+c1}{// produces api()/implementation() distinction}
\PY{+w}{    }\PY{n}{id}\PY{p}{(}\PY{l+s}{\PYZdq{}}\PY{l+s}{jacoco}\PY{l+s}{\PYZdq{}}\PY{p}{)}\PY{+w}{                     }\PY{c+c1}{// coverage}
\PY{+w}{    }\PY{n}{id}\PY{p}{(}\PY{l+s}{\PYZdq{}}\PY{l+s}{com.github.spotbugs}\PY{l+s}{\PYZdq{}}\PY{p}{)}\PY{+w}{ }\PY{n}{version}\PY{+w}{ }\PY{l+s}{\PYZdq{}}\PY{l+s}{6.0.18}\PY{l+s}{\PYZdq{}}
\PY{p}{\PYZcb{}}

\PY{n}{group}\PY{+w}{ }\PY{o}{=}\PY{+w}{ }\PY{l+s}{\PYZdq{}}\PY{l+s}{com.example.orders}\PY{l+s}{\PYZdq{}}
\PY{n}{version}\PY{+w}{ }\PY{o}{=}\PY{+w}{ }\PY{l+s}{\PYZdq{}}\PY{l+s}{1.4.0\PYZhy{}SNAPSHOT}\PY{l+s}{\PYZdq{}}

\PY{n}{repositories}\PY{+w}{ }\PY{p}{\PYZob{}}
\PY{+w}{    }\PY{n}{mavenCentral}\PY{p}{(}\PY{p}{)}
\PY{p}{\PYZcb{}}

\PY{n}{dependencies}\PY{+w}{ }\PY{p}{\PYZob{}}
\PY{+w}{    }\PY{c+c1}{// api: exposed to consumers of THIS library on their compile classpath}
\PY{+w}{    }\PY{n}{api}\PY{p}{(}\PY{l+s}{\PYZdq{}}\PY{l+s}{com.fasterxml.jackson.core:jackson\PYZhy{}databind:2.17.1}\PY{l+s}{\PYZdq{}}\PY{p}{)}

\PY{+w}{    }\PY{c+c1}{// implementation: internal detail, NOT leaked to consumers}
\PY{+w}{    }\PY{n}{implementation}\PY{p}{(}\PY{l+s}{\PYZdq{}}\PY{l+s}{org.apache.commons:commons\PYZhy{}lang3:3.14.0}\PY{l+s}{\PYZdq{}}\PY{p}{)}

\PY{+w}{    }\PY{c+c1}{// compileOnly: needed to compile, but must be supplied at runtime by the caller}
\PY{+w}{    }\PY{n}{compileOnly}\PY{p}{(}\PY{l+s}{\PYZdq{}}\PY{l+s}{org.projectlombok:lombok:1.18.32}\PY{l+s}{\PYZdq{}}\PY{p}{)}
\PY{+w}{    }\PY{n}{annotationProcessor}\PY{p}{(}\PY{l+s}{\PYZdq{}}\PY{l+s}{org.projectlombok:lombok:1.18.32}\PY{l+s}{\PYZdq{}}\PY{p}{)}

\PY{+w}{    }\PY{c+c1}{// runtimeOnly: needed at runtime, not visible at compile time}
\PY{+w}{    }\PY{n}{runtimeOnly}\PY{p}{(}\PY{l+s}{\PYZdq{}}\PY{l+s}{org.postgresql:postgresql:42.7.3}\PY{l+s}{\PYZdq{}}\PY{p}{)}

\PY{+w}{    }\PY{n}{testImplementation}\PY{p}{(}\PY{n}{platform}\PY{p}{(}\PY{l+s}{\PYZdq{}}\PY{l+s}{org.junit:junit\PYZhy{}bom:5.10.2}\PY{l+s}{\PYZdq{}}\PY{p}{)}\PY{p}{)}
\PY{+w}{    }\PY{n}{testImplementation}\PY{p}{(}\PY{l+s}{\PYZdq{}}\PY{l+s}{org.junit.jupiter:junit\PYZhy{}jupiter}\PY{l+s}{\PYZdq{}}\PY{p}{)}
\PY{+w}{    }\PY{n}{testImplementation}\PY{p}{(}\PY{l+s}{\PYZdq{}}\PY{l+s}{org.assertj:assertj\PYZhy{}core:3.26.0}\PY{l+s}{\PYZdq{}}\PY{p}{)}
\PY{p}{\PYZcb{}}

\PY{n}{java}\PY{+w}{ }\PY{p}{\PYZob{}}
\PY{+w}{    }\PY{n}{toolchain}\PY{+w}{ }\PY{p}{\PYZob{}}
\PY{+w}{        }\PY{n}{languageVersion}\PY{p}{.}\PY{n+na}{set}\PY{p}{(}\PY{n}{JavaLanguageVersion}\PY{p}{.}\PY{n+na}{of}\PY{p}{(}\PY{l+m}{2}\PY{l+m}{1}\PY{p}{)}\PY{p}{)}
\PY{+w}{    }\PY{p}{\PYZcb{}}
\PY{+w}{    }\PY{n}{withSourcesJar}\PY{p}{(}\PY{p}{)}
\PY{+w}{    }\PY{n}{withJavadocJar}\PY{p}{(}\PY{p}{)}
\PY{p}{\PYZcb{}}

\PY{n}{tasks}\PY{p}{.}\PY{n+na}{test}\PY{+w}{ }\PY{p}{\PYZob{}}
\PY{+w}{    }\PY{n}{useJUnitPlatform}\PY{p}{(}\PY{p}{)}
\PY{+w}{    }\PY{n}{maxParallelForks}\PY{+w}{ }\PY{o}{=}\PY{+w}{ }\PY{n}{Runtime}\PY{p}{.}\PY{n+na}{getRuntime}\PY{p}{(}\PY{p}{)}\PY{p}{.}\PY{n+na}{availableProcessors}\PY{p}{(}\PY{p}{)}\PY{+w}{ }\PY{o}{/}\PY{+w}{ }\PY{l+m}{2}
\PY{p}{\PYZcb{}}

\PY{n}{tasks}\PY{p}{.}\PY{n+na}{jacocoTestReport}\PY{+w}{ }\PY{p}{\PYZob{}}
\PY{+w}{    }\PY{n}{dependsOn}\PY{p}{(}\PY{n}{tasks}\PY{p}{.}\PY{n+na}{test}\PY{p}{)}
\PY{p}{\PYZcb{}}
\end{Verbatim}
\end{codebox}

\subsection[The plugins \{\} Block]{The \code{plugins~\allowbreak{}\{\}} Block}
\label{h:21:the-plugins--block}
The \code{plugins~\allowbreak{}\{\}} block applies \textbf{binary plugins} by ID, resolved from the Gradle Plugin Portal (or \code{mavenCentral()} for some). It replaces the older, discouraged \code{apply~\allowbreak{}plugin:\allowbreak{}~\allowbreak{}\textquotesingle{}java\textquotesingle{}} / buildscript-classpath style, because it lets Gradle resolve and version-check plugins before any script code runs, enabling better IDE support and the plugin version-management shown below.

\begin{codebox}
\begin{Verbatim}[commandchars=\\\{\}]
\PY{n}{plugins}\PY{+w}{ }\PY{p}{\PYZob{}}
\PY{+w}{    }\PY{n}{java}
\PY{+w}{    }\PY{n}{id}\PY{p}{(}\PY{l+s}{\PYZdq{}}\PY{l+s}{org.springframework.boot}\PY{l+s}{\PYZdq{}}\PY{p}{)}\PY{+w}{ }\PY{n}{version}\PY{+w}{ }\PY{l+s}{\PYZdq{}}\PY{l+s}{3.3.0}\PY{l+s}{\PYZdq{}}
\PY{+w}{    }\PY{n}{id}\PY{p}{(}\PY{l+s}{\PYZdq{}}\PY{l+s}{io.spring.dependency\PYZhy{}management}\PY{l+s}{\PYZdq{}}\PY{p}{)}\PY{+w}{ }\PY{n}{version}\PY{+w}{ }\PY{l+s}{\PYZdq{}}\PY{l+s}{1.1.5}\PY{l+s}{\PYZdq{}}
\PY{p}{\PYZcb{}}
\end{Verbatim}
\end{codebox}

\subsection[The Task Graph and Incremental Builds]{The Task Graph and Incremental Builds}
\label{h:21:the-task-graph-and-incremental-builds}
Every unit of work in Gradle is a \textbf{task} (\code{compileJava}, \code{test}, \code{jar}, …). Tasks declare \textbf{inputs} and \textbf{outputs}; Gradle hashes them, and if neither changed since the last run, the task is reported \code{UP-\allowbreak{}TO-\allowbreak{}DATE} and skipped entirely — this is \textbf{incremental build}. Running \code{./\allowbreak{}gradlew~\allowbreak{}build} twice in a row with no source changes does almost nothing the second time.

\begin{codebox}
\begin{Verbatim}[commandchars=\\\{\}]
\PYZdl{}\PY{+w}{ }./gradlew\PY{+w}{ }build
\PYZgt{}\PY{+w}{ }Task\PY{+w}{ }:compileJava
\PYZgt{}\PY{+w}{ }Task\PY{+w}{ }:processResources\PY{+w}{ }UP\PYZhy{}TO\PYZhy{}DATE
\PYZgt{}\PY{+w}{ }Task\PY{+w}{ }:classes
\PYZgt{}\PY{+w}{ }Task\PY{+w}{ }:test
\PYZgt{}\PY{+w}{ }Task\PY{+w}{ }:jar
\PYZgt{}\PY{+w}{ }Task\PY{+w}{ }:build

BUILD\PY{+w}{ }SUCCESSFUL
\end{Verbatim}
\end{codebox}

You can inspect why a task ran (or didn’t) and visualize the dependency graph:

\begin{codebox}
\begin{Verbatim}[commandchars=\\\{\}]
./gradlew\PY{+w}{ }build\PY{+w}{ }\PYZhy{}\PYZhy{}dry\PYZhy{}run\PY{+w}{          }\PY{c+c1}{\PYZsh{} print the task plan without executing anything}
./gradlew\PY{+w}{ }\PY{n+nb}{test}\PY{+w}{ }\PYZhy{}\PYZhy{}rerun\PY{+w}{             }\PY{c+c1}{\PYZsh{} force a task to rerun even if UP\PYZhy{}TO\PYZhy{}DATE}
\end{Verbatim}
\end{codebox}

\subsection[Configuration Cache and Build Cache]{Configuration Cache and Build Cache}
\label{h:21:configuration-cache-and-build-cache}
These are two distinct, complementary optimizations, and mixing them up is a common misunderstanding:

\begin{itemize}
\item 
\textbf{Configuration cache}: caches the \emph{result of evaluating the build scripts themselves} (the task graph), so subsequent invocations skip re-running the configuration phase. Speeds up every single build, even a clean one, because it avoids re-parsing/re-executing Groovy/Kotlin build logic. Enable with \code{org.\allowbreak{}gradle.\allowbreak{}configuration-\allowbreak{}cache=\allowbreak{}true} in \code{gradle.\allowbreak{}properties}.

\item 
\textbf{Build cache}: caches \emph{task outputs} keyed by task input hash, and can restore an already-computed output on a \emph{different} machine (with a shared remote cache) or a different branch, skipping execution entirely — not just “up to date” within one workspace, but genuinely retrieved from cache. Enable with \code{org.\allowbreak{}gradle.\allowbreak{}caching=\allowbreak{}true}.

\end{itemize}

\begin{codebox}
\begin{Verbatim}[commandchars=\\\{\}]
\PY{c+c1}{\PYZsh{} gradle.properties}
\PY{n+na}{org.gradle.configuration\PYZhy{}cache}\PY{o}{=}\PY{l+s}{true}
\PY{n+na}{org.gradle.caching}\PY{o}{=}\PY{l+s}{true}
\PY{n+na}{org.gradle.parallel}\PY{o}{=}\PY{l+s}{true}
\end{Verbatim}
\end{codebox}

\begin{codebox}
\begin{Verbatim}[commandchars=\\\{\}]
./gradlew\PY{+w}{ }build\PY{+w}{ }\PYZhy{}\PYZhy{}build\PYZhy{}cache\PY{+w}{      }\PY{c+c1}{\PYZsh{} explicitly opt in to the build cache for this invocation}
\end{Verbatim}
\end{codebox}

\subsection[Dependency Configurations]{Dependency Configurations}
\label{h:21:dependency-configurations}
{\tablefont
\begin{xltabular}{\linewidth}{@{}L{0.686}L{0.952}L{0.457}L{1.905}@{}}
\toprule
\rowcolor{tablehdr}
\thd{Configuration} & \thd{Compile classpath (consumer sees it?)} & \thd{Runtime classpath} & \thd{Typical use} \\
\midrule
\endhead
\code{api} & Yes — leaks to consumers & Yes & Types that appear in your public method signatures \\
\code{implementation} & No — internal only & Yes & Everything else; the default, safest choice \\
\code{compileOnly} & Yes (this module only) & No & Annotation-only or container-supplied deps (Lombok, servlet API) \\
\code{runtimeOnly} & No & Yes & JDBC drivers, logging backends \\
\code{testImplementation} & Test compile only & Test runtime & JUnit, Mockito \\
\code{annotationProcessor} & n/a (processor classpath) & No & Lombok, MapStruct, Dagger processors \\
\bottomrule
\end{xltabular}
}

\textbf{Why \code{implementation} beats the old \code{compile}}: the deprecated \code{compile} configuration put every dependency on the \emph{consumer’s} compile classpath transitively, whether or not the consumer’s code actually referenced it. That meant bumping an internal-only dependency’s version in a library forced every downstream module to recompile, even though nothing in the public API changed. \code{implementation} hides internal dependencies from consumers entirely, which (a) shrinks consumers’ compile classpaths, (b) means changing an internal dependency no longer forces downstream recompilation, and © makes accidental leakage of internal types through public APIs a compile error for consumers, rather than something that silently works until the dependency changes.

\begin{codebox}
\begin{Verbatim}[commandchars=\\\{\}]
\PY{n}{dependencies}\PY{+w}{ }\PY{p}{\PYZob{}}
\PY{+w}{    }\PY{c+c1}{// Consumer of this module can call methods that return ObjectMapper directly:}
\PY{+w}{    }\PY{n}{api}\PY{p}{(}\PY{l+s}{\PYZdq{}}\PY{l+s}{com.fasterxml.jackson.databind:jackson\PYZhy{}databind:2.17.1}\PY{l+s}{\PYZdq{}}\PY{p}{)}

\PY{+w}{    }\PY{c+c1}{// Used internally only, e.g., for a private helper class — invisible to consumers:}
\PY{+w}{    }\PY{n}{implementation}\PY{p}{(}\PY{l+s}{\PYZdq{}}\PY{l+s}{org.apache.commons:commons\PYZhy{}lang3:3.14.0}\PY{l+s}{\PYZdq{}}\PY{p}{)}
\PY{p}{\PYZcb{}}
\end{Verbatim}
\end{codebox}

\subsection[Version Catalogs (libs.versions.toml)]{Version Catalogs (\code{libs.\allowbreak{}versions.\allowbreak{}toml})}
\label{h:21:version-catalogs-libsversionstoml}
A \textbf{version catalog} centralizes dependency coordinates and versions in one TOML file, shared across all subprojects, with type-safe accessors generated for both DSLs:

\begin{codebox}
\begin{Verbatim}[commandchars=\\\{\}]
\PY{c+c1}{\PYZsh{} gradle/libs.versions.toml}
\PY{k}{[}\PY{k}{versions}\PY{k}{]}
\PY{n}{junit}\PY{+w}{ }\PY{o}{=}\PY{+w}{ }\PY{l+s+s2}{\PYZdq{}}\PY{l+s+s2}{5.10.2}\PY{l+s+s2}{\PYZdq{}}
\PY{n}{jackson}\PY{+w}{ }\PY{o}{=}\PY{+w}{ }\PY{l+s+s2}{\PYZdq{}}\PY{l+s+s2}{2.17.1}\PY{l+s+s2}{\PYZdq{}}

\PY{k}{[}\PY{k}{libraries}\PY{k}{]}
\PY{n}{junit\PYZhy{}jupiter}\PY{+w}{ }\PY{o}{=}\PY{+w}{ }\PY{p}{\PYZob{}}\PY{+w}{ }\PY{n}{module}\PY{+w}{ }\PY{p}{=}\PY{+w}{ }\PY{l+s+s2}{\PYZdq{}}\PY{l+s+s2}{org.junit.jupiter:junit\PYZhy{}jupiter}\PY{l+s+s2}{\PYZdq{}}\PY{p}{,}\PY{+w}{ }\PY{n}{version}\PY{p}{.}\PY{n}{ref}\PY{+w}{ }\PY{p}{=}\PY{+w}{ }\PY{l+s+s2}{\PYZdq{}}\PY{l+s+s2}{junit}\PY{l+s+s2}{\PYZdq{}}\PY{+w}{ }\PY{p}{\PYZcb{}}
\PY{n}{jackson\PYZhy{}databind}\PY{+w}{ }\PY{o}{=}\PY{+w}{ }\PY{p}{\PYZob{}}\PY{+w}{ }\PY{n}{module}\PY{+w}{ }\PY{p}{=}\PY{+w}{ }\PY{l+s+s2}{\PYZdq{}}\PY{l+s+s2}{com.fasterxml.jackson.core:jackson\PYZhy{}databind}\PY{l+s+s2}{\PYZdq{}}\PY{p}{,}\PY{+w}{ }\PY{n}{version}\PY{p}{.}\PY{n}{ref}\PY{+w}{ }\PY{p}{=}\PY{+w}{ }\PY{l+s+s2}{\PYZdq{}}\PY{l+s+s2}{jackson}\PY{l+s+s2}{\PYZdq{}}\PY{+w}{ }\PY{p}{\PYZcb{}}

\PY{k}{[}\PY{k}{bundles}\PY{k}{]}
\PY{n}{testing}\PY{+w}{ }\PY{o}{=}\PY{+w}{ }\PY{p}{[}\PY{l+s+s2}{\PYZdq{}}\PY{l+s+s2}{junit\PYZhy{}jupiter}\PY{l+s+s2}{\PYZdq{}}\PY{p}{]}

\PY{k}{[}\PY{k}{plugins}\PY{k}{]}
\PY{n}{spring\PYZhy{}boot}\PY{+w}{ }\PY{o}{=}\PY{+w}{ }\PY{p}{\PYZob{}}\PY{+w}{ }\PY{n}{id}\PY{+w}{ }\PY{p}{=}\PY{+w}{ }\PY{l+s+s2}{\PYZdq{}}\PY{l+s+s2}{org.springframework.boot}\PY{l+s+s2}{\PYZdq{}}\PY{p}{,}\PY{+w}{ }\PY{n}{version}\PY{+w}{ }\PY{p}{=}\PY{+w}{ }\PY{l+s+s2}{\PYZdq{}}\PY{l+s+s2}{3.3.0}\PY{l+s+s2}{\PYZdq{}}\PY{+w}{ }\PY{p}{\PYZcb{}}
\end{Verbatim}
\end{codebox}

\begin{codebox}
\begin{Verbatim}[commandchars=\\\{\}]
\PY{c+c1}{// build.gradle.kts — type\PYZhy{}safe accessors generated from the catalog above}
\PY{n}{plugins}\PY{+w}{ }\PY{p}{\PYZob{}}
\PY{+w}{    }\PY{n}{alias}\PY{p}{(}\PY{n}{libs}\PY{p}{.}\PY{n+na}{plugins}\PY{p}{.}\PY{n+na}{spring}\PY{p}{.}\PY{n+na}{boot}\PY{p}{)}
\PY{p}{\PYZcb{}}

\PY{n}{dependencies}\PY{+w}{ }\PY{p}{\PYZob{}}
\PY{+w}{    }\PY{n}{implementation}\PY{p}{(}\PY{n}{libs}\PY{p}{.}\PY{n+na}{jackson}\PY{p}{.}\PY{n+na}{databind}\PY{p}{)}
\PY{+w}{    }\PY{n}{testImplementation}\PY{p}{(}\PY{n}{libs}\PY{p}{.}\PY{n+na}{bundles}\PY{p}{.}\PY{n+na}{testing}\PY{p}{)}
\PY{p}{\PYZcb{}}
\end{Verbatim}
\end{codebox}

This is the Gradle equivalent of a Maven BOM plus centralized \code{\textless{}\allowbreak{}properties\textgreater{}}: one file, one place to bump a version, autocomplete in the IDE, and no version drift between subprojects.

\subsection[Multi-Project Builds and settings.gradle.kts]{Multi-Project Builds and \code{settings.\allowbreak{}gradle.\allowbreak{}kts}}
\label{h:21:multi-project-builds-and-settingsgradlekts}
\code{settings.\allowbreak{}gradle.\allowbreak{}kts} (at the repo root) declares which subprojects participate in the build — this is what makes a directory a “project” at all, before any \code{build.\allowbreak{}gradle.\allowbreak{}kts} is even read.

\begin{codebox}
\begin{Verbatim}[commandchars=\\\{\}]
\PY{c+c1}{// settings.gradle.kts}
\PY{n}{rootProject}\PY{p}{.}\PY{n+na}{name}\PY{+w}{ }\PY{o}{=}\PY{+w}{ }\PY{l+s}{\PYZdq{}}\PY{l+s}{orders\PYZhy{}platform}\PY{l+s}{\PYZdq{}}

\PY{n}{include}\PY{p}{(}\PY{l+s}{\PYZdq{}}\PY{l+s}{orders\PYZhy{}api}\PY{l+s}{\PYZdq{}}\PY{p}{,}\PY{+w}{ }\PY{l+s}{\PYZdq{}}\PY{l+s}{orders\PYZhy{}service}\PY{l+s}{\PYZdq{}}\PY{p}{,}\PY{+w}{ }\PY{l+s}{\PYZdq{}}\PY{l+s}{orders\PYZhy{}persistence}\PY{l+s}{\PYZdq{}}\PY{p}{)}

\PY{n}{dependencyResolutionManagement}\PY{+w}{ }\PY{p}{\PYZob{}}
\PY{+w}{    }\PY{n}{repositories}\PY{+w}{ }\PY{p}{\PYZob{}}
\PY{+w}{        }\PY{n}{mavenCentral}\PY{p}{(}\PY{p}{)}
\PY{+w}{    }\PY{p}{\PYZcb{}}
\PY{p}{\PYZcb{}}
\end{Verbatim}
\end{codebox}

A subproject depends on another subproject with the \code{project()} accessor, not a version string:

\begin{codebox}
\begin{Verbatim}[commandchars=\\\{\}]
\PY{c+c1}{// orders\PYZhy{}service/build.gradle.kts}
\PY{n}{dependencies}\PY{+w}{ }\PY{p}{\PYZob{}}
\PY{+w}{    }\PY{n}{implementation}\PY{p}{(}\PY{n}{project}\PY{p}{(}\PY{l+s}{\PYZdq{}}\PY{l+s}{:orders\PYZhy{}api}\PY{l+s}{\PYZdq{}}\PY{p}{)}\PY{p}{)}
\PY{p}{\PYZcb{}}
\end{Verbatim}
\end{codebox}

\begin{codebox}
\begin{Verbatim}[commandchars=\\\{\}]
./gradlew\PY{+w}{ }:orders\PYZhy{}service:build\PY{+w}{     }\PY{c+c1}{\PYZsh{} build a single subproject}
./gradlew\PY{+w}{ }build\PY{+w}{                     }\PY{c+c1}{\PYZsh{} build everything from the root}
\end{Verbatim}
\end{codebox}

\subsection[Custom Tasks]{Custom Tasks}
\label{h:21:custom-tasks}
\begin{codebox}
\begin{Verbatim}[commandchars=\\\{\}]
\PY{n}{tasks}\PY{p}{.}\PY{n+na}{register}\PY{p}{(}\PY{l+s}{\PYZdq{}}\PY{l+s}{printVersion}\PY{l+s}{\PYZdq{}}\PY{p}{)}\PY{+w}{ }\PY{p}{\PYZob{}}
\PY{+w}{    }\PY{n}{doLast}\PY{+w}{ }\PY{p}{\PYZob{}}
\PY{+w}{        }\PY{n}{println}\PY{p}{(}\PY{l+s}{\PYZdq{}}\PY{l+s}{Building version }\PY{l+s+si}{\PYZdl{}\PYZob{}}\PY{n}{project}\PY{p}{.}\PY{n+na}{version}\PY{l+s+si}{\PYZcb{}}\PY{l+s}{\PYZdq{}}\PY{p}{)}
\PY{+w}{    }\PY{p}{\PYZcb{}}
\PY{p}{\PYZcb{}}

\PY{n}{tasks}\PY{p}{.}\PY{n+na}{register}\PY{o}{\PYZlt{}}\PY{n}{Copy}\PY{o}{\PYZgt{}}\PY{p}{(}\PY{l+s}{\PYZdq{}}\PY{l+s}{copyConfig}\PY{l+s}{\PYZdq{}}\PY{p}{)}\PY{+w}{ }\PY{p}{\PYZob{}}
\PY{+w}{    }\PY{n}{from}\PY{p}{(}\PY{l+s}{\PYZdq{}}\PY{l+s}{src/main/resources/config}\PY{l+s}{\PYZdq{}}\PY{p}{)}
\PY{+w}{    }\PY{n}{into}\PY{p}{(}\PY{n}{layout}\PY{p}{.}\PY{n+na}{buildDirectory}\PY{p}{.}\PY{n+na}{dir}\PY{p}{(}\PY{l+s}{\PYZdq{}}\PY{l+s}{config}\PY{l+s}{\PYZdq{}}\PY{p}{)}\PY{p}{)}
\PY{p}{\PYZcb{}}

\PY{n}{tasks}\PY{p}{.}\PY{n+na}{named}\PY{p}{(}\PY{l+s}{\PYZdq{}}\PY{l+s}{build}\PY{l+s}{\PYZdq{}}\PY{p}{)}\PY{+w}{ }\PY{p}{\PYZob{}}
\PY{+w}{    }\PY{n}{dependsOn}\PY{p}{(}\PY{l+s}{\PYZdq{}}\PY{l+s}{printVersion}\PY{l+s}{\PYZdq{}}\PY{p}{)}
\PY{p}{\PYZcb{}}
\end{Verbatim}
\end{codebox}

Custom tasks with typed inputs/outputs (rather than \code{doLast~\allowbreak{}\{\}} closures) can participate correctly in incremental build and the build cache — a plain \code{doLast} task has no declared inputs/outputs and Gradle can never mark it \code{UP-\allowbreak{}TO-\allowbreak{}DATE}.

\subsection[The Gradle Wrapper]{The Gradle Wrapper}
\label{h:21:the-gradle-wrapper}
Exactly like Maven’s wrapper, \code{gradlew}/\code{gradlew.\allowbreak{}bat} pin the exact Gradle version a project was built and tested with, downloaded on demand.

\begin{codebox}
\begin{Verbatim}[commandchars=\\\{\}]
gradle\PY{+w}{ }wrapper\PY{+w}{ }\PYZhy{}\PYZhy{}gradle\PYZhy{}version\PY{+w}{ }\PY{l+m}{8}.8\PY{+w}{   }\PY{c+c1}{\PYZsh{} (re)generate the wrapper for version 8.8}
./gradlew\PY{+w}{ }build\PY{+w}{                        }\PY{c+c1}{\PYZsh{} anyone can build without installing Gradle}
\end{Verbatim}
\end{codebox}

\textbf{Always commit the wrapper} (\code{gradlew}, \code{gradlew.\allowbreak{}bat}, \code{gradle/\allowbreak{}wrapper/\allowbreak{}gradle-\allowbreak{}wrapper.\allowbreak{}jar}, \code{gradle/\allowbreak{}wrapper/\allowbreak{}gradle-\allowbreak{}wrapper.\allowbreak{}properties}). Without it, a teammate or CI runner using a different globally-installed Gradle version can produce a different, non-reproducible build — or fail outright on a script feature that only exists in a newer/older Gradle.

\subsection[Diagnostics: --scan, --dry-run, dependencies]{Diagnostics: \code{--\allowbreak{}scan}, \code{--\allowbreak{}dry-\allowbreak{}run}, \code{dependencies}}
\label{h:21:diagnostics---scan---dry-run-dependencies}
\begin{codebox}
\begin{Verbatim}[commandchars=\\\{\}]
./gradlew\PY{+w}{ }build\PY{+w}{ }\PYZhy{}\PYZhy{}scan\PY{+w}{          }\PY{c+c1}{\PYZsh{} publish a shareable, detailed build report (timings, deprecations, task graph)}
./gradlew\PY{+w}{ }build\PY{+w}{ }\PYZhy{}\PYZhy{}dry\PYZhy{}run\PY{+w}{       }\PY{c+c1}{\PYZsh{} show which tasks WOULD run, without running them}
./gradlew\PY{+w}{ }dependencies\PY{+w}{          }\PY{c+c1}{\PYZsh{} print the resolved dependency tree for all configurations}
./gradlew\PY{+w}{ }dependencies\PY{+w}{ }\PYZhy{}\PYZhy{}configuration\PY{+w}{ }runtimeClasspath\PY{+w}{   }\PY{c+c1}{\PYZsh{} scope it to one configuration}
./gradlew\PY{+w}{ }:orders\PYZhy{}service:dependencyInsight\PY{+w}{ }\PYZhy{}\PYZhy{}dependency\PY{+w}{ }commons\PYZhy{}lang3\PY{+w}{   }\PY{c+c1}{\PYZsh{} why is this version chosen?}
\end{Verbatim}
\end{codebox}

\subsection[Maven vs Gradle Comparison]{Maven vs Gradle Comparison}
\label{h:21:maven-vs-gradle-comparison}
{\tablefont
\begin{xltabular}{\linewidth}{@{}L{0.566}L{1.132}L{1.302}@{}}
\toprule
\rowcolor{tablehdr}
\thd{Aspect} & \thd{Maven} & \thd{Gradle} \\
\midrule
\endhead
Config style & Declarative XML & Code (Groovy or Kotlin DSL) \\
Build model & Fixed lifecycle of phases & Flexible task graph \\
Incremental builds & Limited (plugin-dependent) & First-class, built into every task \\
Build/config caching & No built-in equivalent & Configuration cache + (remote) build cache \\
Multi-module & Parent POM + \code{\textless{}\allowbreak{}modules\textgreater{}} & \code{settings.\allowbreak{}gradle.\allowbreak{}kts} + subprojects \\
Dependency management & \code{dependencyManagement} + BOM import & Version catalogs (\code{libs.\allowbreak{}versions.\allowbreak{}toml}) \\
Learning curve & Lower — one way to do things & Higher — scripting flexibility cuts both ways \\
Ecosystem & Extremely mature, huge plugin base, Java/Enterprise-centric & Dominant for Android; strong in polyglot/large monorepos \\
Performance on large builds & Slower on very large multi-module builds & Generally faster due to caching/incrementality \\
Debuggability & Very predictable, easy to reason about from the POM alone & Custom Groovy/Kotlin logic can be hard to review/debug \\
\bottomrule
\end{xltabular}
}

\textbf{When to pick which}: choose Maven for straightforward, mostly-declarative enterprise Java projects where predictability and a large, stable plugin ecosystem matter more than raw build speed, or when the team values “one obvious way to configure things.” Choose Gradle for large monorepos, Android (mandatory), polyglot builds, or when build performance (incremental/cached builds) matters enough to justify the extra scripting flexibility and the review overhead that comes with it.

\section[Javadoc]{Javadoc}
\label{h:21:javadoc}
Javadoc is the documentation format baked into the Java language itself: specially-formatted comments (\code{/**\allowbreak{}~\allowbreak{}...\allowbreak{}~\allowbreak{}*/}) immediately before a declaration, parsed by the \code{javadoc} tool (and IDEs) into structured API documentation. In code review, Javadoc quality is judged on one criterion above all others: \textbf{does it describe the contract, not the implementation?}

\subsection[The Comment Format]{The Comment Format}
\label{h:21:the-comment-format}
A Javadoc comment starts with \code{/**} (two asterisks), ends with \code{*/}, and sits directly above the declaration it documents — a class, interface, method, constructor, or field. Every other \code{*}-prefixed line is conventional but not required by the parser.

\begin{codebox}
\begin{Verbatim}[commandchars=\\\{\}]
\PY{c+cm}{/**}
\PY{c+cm}{ * Converts a monetary amount between currencies using the exchange rate}
\PY{c+cm}{ * in effect at the time this method is called.}
\PY{c+cm}{ *}
\PY{c+cm}{ * @param amount the amount to convert; must be non\PYZhy{}negative}
\PY{c+cm}{ * @param from the source currency}
\PY{c+cm}{ * @param to the target currency}
\PY{c+cm}{ * @return the converted amount, rounded to the target currency\PYZsq{}s}
\PY{c+cm}{ *         standard number of decimal places}
\PY{c+cm}{ * @throws IllegalArgumentException if \PYZob{}@code amount\PYZcb{} is negative}
\PY{c+cm}{ * @throws CurrencyUnavailableException if no exchange rate is available}
\PY{c+cm}{ *         for the given currency pair}
\PY{c+cm}{ */}
\PY{k+kd}{public}\PY{+w}{ }\PY{n}{BigDecimal}\PY{+w}{ }\PY{n+nf}{convert}\PY{p}{(}\PY{n}{BigDecimal}\PY{+w}{ }\PY{n}{amount}\PY{p}{,}\PY{+w}{ }\PY{n}{Currency}\PY{+w}{ }\PY{n}{from}\PY{p}{,}\PY{+w}{ }\PY{n}{Currency}\PY{+w}{ }\PY{n}{to}\PY{p}{)}\PY{+w}{ }\PY{p}{\PYZob{}}
\PY{+w}{    }\PY{p}{.}\PY{p}{.}\PY{p}{.}
\PY{p}{\PYZcb{}}
\end{Verbatim}
\end{codebox}

\subsection[Standard Tags]{Standard Tags}
\label{h:21:standard-tags}
{\tablefont
\begin{xltabular}{\linewidth}{@{}L{0.433}L{0.700}L{1.867}@{}}
\toprule
\rowcolor{tablehdr}
\thd{Tag} & \thd{Where used} & \thd{Purpose} \\
\midrule
\endhead
\code{@\allowbreak{}param} & Methods, constructors & Describes a parameter — name plus meaning/constraints \\
\code{@\allowbreak{}return} & Methods (non-\code{void}) & Describes the return value and its meaning \\
\code{@\allowbreak{}throws} (or \code{@\allowbreak{}exception}) & Methods, constructors & Describes when/why an exception is thrown \\
\code{@\allowbreak{}see} & Anywhere & Cross-reference to a related class/method/URL \\
\code{@\allowbreak{}since} & Classes, methods & Version this API first appeared in \\
\code{@\allowbreak{}deprecated} & Classes, methods, fields & Marks the API as obsolete; should always explain the replacement \\
\code{@\allowbreak{}author} & Classes & Author name(s) — increasingly omitted in favor of VCS history \\
\code{@\allowbreak{}serial} & Fields (of \code{Serializable} classes) & Documents the serialized form of a field for the serialization contract \\
\bottomrule
\end{xltabular}
}

\begin{codebox}
\begin{Verbatim}[commandchars=\\\{\}]
\PY{c+cm}{/**}
\PY{c+cm}{ * Legacy price calculator.}
\PY{c+cm}{ *}
\PY{c+cm}{ * @deprecated Use \PYZob{}@link PricingEngine\PYZsh{}calculate(Order)\PYZcb{} instead; this}
\PY{c+cm}{ *             class does not account for regional tax rules and will be}
\PY{c+cm}{ *             removed in 3.0.}
\PY{c+cm}{ */}
\PY{n+nd}{@Deprecated}\PY{p}{(}\PY{n}{since}\PY{+w}{ }\PY{o}{=}\PY{+w}{ }\PY{l+s}{\PYZdq{}}\PY{l+s}{2.4}\PY{l+s}{\PYZdq{}}\PY{p}{,}\PY{+w}{ }\PY{n}{forRemoval}\PY{+w}{ }\PY{o}{=}\PY{+w}{ }\PY{k+kc}{true}\PY{p}{)}
\PY{k+kd}{public}\PY{+w}{ }\PY{k+kd}{class} \PY{n+nc}{LegacyPriceCalculator}\PY{+w}{ }\PY{p}{\PYZob{}}
\PY{+w}{    }\PY{p}{.}\PY{p}{.}\PY{p}{.}
\PY{p}{\PYZcb{}}
\end{Verbatim}
\end{codebox}

Note the pairing convention: \code{@\allowbreak{}deprecated} in the Javadoc comment (human-readable explanation) should always accompany the \code{@\allowbreak{}Deprecated} annotation (machine-readable, drives compiler warnings) — one without the other is an incomplete deprecation.

\subsection[Inline Tags]{Inline Tags}
\label{h:21:inline-tags}
Inline tags are embedded inside the flowing text of a description, wrapped in \code{\{@\allowbreak{}~\allowbreak{}...\allowbreak{}~\allowbreak{}\}}.

{\tablefont
\begin{xltabular}{\linewidth}{@{}L{0.475}L{1.525}@{}}
\toprule
\rowcolor{tablehdr}
\thd{Inline tag} & \thd{Purpose} \\
\midrule
\endhead
\code{\{@\allowbreak{}link~\allowbreak{}Type\#member\}} & Clickable cross-reference, rendered as a link \\
\code{\{@\allowbreak{}linkplain~\allowbreak{}Type\#member\}} & Same as \code{\{@\allowbreak{}link\}} but rendered as plain text, not code font \\
\code{\{@\allowbreak{}code~\allowbreak{}text\}} & Renders as code font, and — critically — HTML-escapes its content, so \code{\textless{}}, \code{\textgreater{}}, \code{\&} are safe \\
\code{\{@\allowbreak{}literal~\allowbreak{}text\}} & Escapes HTML like \code{\{@\allowbreak{}code\}} but does \emph{not} apply code font \\
\code{\{@\allowbreak{}value~\allowbreak{}\#FIELD\}} & Inlines the value of a static final constant \\
\code{\{@\allowbreak{}inheritDoc\}} & Pulls the Javadoc text from the overridden/implemented method \\
\code{\{@\allowbreak{}snippet~\allowbreak{}...\}} & (JDK 18+) Embeds a validated, syntax-highlighted code example \\
\bottomrule
\end{xltabular}
}

\begin{codebox}
\begin{Verbatim}[commandchars=\\\{\}]
\PY{c+cm}{/**}
\PY{c+cm}{ * The default timeout, in milliseconds: \PYZob{}@value \PYZsh{}DEFAULT\PYZus{}TIMEOUT\PYZus{}MS\PYZcb{}.}
\PY{c+cm}{ *}
\PY{c+cm}{ * \PYZlt{}p\PYZgt{}Use \PYZob{}@code Duration.ofMillis(timeout)\PYZcb{} to convert to a \PYZob{}@link java.time.Duration\PYZcb{}}
\PY{c+cm}{ * if you need a \PYZob{}@link java.time.Duration\PYZcb{}\PYZhy{}based API instead.}
\PY{c+cm}{ *}
\PY{c+cm}{ * \PYZlt{}p\PYZgt{}Generic type bound example (needs \PYZob{}@literal @literal\PYZcb{}, not \PYZob{}@code @code\PYZcb{},}
\PY{c+cm}{ * because \PYZob{}@code List\PYZlt{}String\PYZgt{}\PYZcb{} would be parsed as an HTML tag otherwise):}
\PY{c+cm}{ * a raw \PYZob{}@literal List\PYZlt{}String\PYZgt{}\PYZcb{} is unsafe.}
\PY{c+cm}{ */}
\PY{k+kd}{public}\PY{+w}{ }\PY{k+kd}{static}\PY{+w}{ }\PY{k+kd}{final}\PY{+w}{ }\PY{k+kt}{int}\PY{+w}{ }\PY{n}{DEFAULT\PYZus{}TIMEOUT\PYZus{}MS}\PY{+w}{ }\PY{o}{=}\PY{+w}{ }\PY{l+m+mi}{5000}\PY{p}{;}
\end{Verbatim}
\end{codebox}

\begin{codebox}
\begin{Verbatim}[commandchars=\\\{\}]
\PY{c+cm}{/**}
\PY{c+cm}{ * \PYZob{}@inheritDoc\PYZcb{}}
\PY{c+cm}{ *}
\PY{c+cm}{ * \PYZlt{}p\PYZgt{}This implementation additionally logs every call at DEBUG level.}
\PY{c+cm}{ */}
\PY{n+nd}{@Override}
\PY{k+kd}{public}\PY{+w}{ }\PY{k+kt}{void}\PY{+w}{ }\PY{n+nf}{process}\PY{p}{(}\PY{n}{Order}\PY{+w}{ }\PY{n}{order}\PY{p}{)}\PY{+w}{ }\PY{p}{\PYZob{}}
\PY{+w}{    }\PY{p}{.}\PY{p}{.}\PY{p}{.}
\PY{p}{\PYZcb{}}
\end{Verbatim}
\end{codebox}

The \code{\{@\allowbreak{}snippet\}} tag (JEP 413, finalized JDK 18) replaces the old \code{\textless{}\allowbreak{}pre\textgreater{}\{@\allowbreak{}code~\allowbreak{}...\}\textless{}/\allowbreak{}pre\textgreater{}} idiom with something the \code{javadoc} tool can actually validate and highlight:

\begin{codebox}
\begin{Verbatim}[commandchars=\\\{\}]
\PY{c+cm}{/**}
\PY{c+cm}{ * Builds an immutable list.}
\PY{c+cm}{ *}
\PY{c+cm}{ * \PYZob{}@snippet lang = \PYZdq{}java\PYZdq{} :}
\PY{c+cm}{ * List\PYZlt{}String\PYZgt{} names = List.of(\PYZdq{}Ada\PYZdq{}, \PYZdq{}Grace\PYZdq{}, \PYZdq{}Linus\PYZdq{});}
\PY{c+cm}{ * names.forEach(System.out::println);}
\PY{c+cm}{ * \PYZcb{}}
\PY{c+cm}{ */}
\PY{k+kd}{public}\PY{+w}{ }\PY{k+kd}{static}\PY{+w}{ }\PY{o}{\PYZlt{}}\PY{n}{T}\PY{o}{\PYZgt{}}\PY{+w}{ }\PY{n}{List}\PY{o}{\PYZlt{}}\PY{n}{T}\PY{o}{\PYZgt{}}\PY{+w}{ }\PY{n+nf}{immutableListOf}\PY{p}{(}\PY{n}{T}\PY{p}{.}\PY{p}{.}\PY{p}{.}\PY{+w}{ }\PY{n}{items}\PY{p}{)}\PY{+w}{ }\PY{p}{\PYZob{}}\PY{+w}{ }\PY{p}{.}\PY{p}{.}\PY{p}{.}\PY{+w}{ }\PY{p}{\PYZcb{}}
\end{Verbatim}
\end{codebox}

\subsection[The First-Sentence Summary Rule]{The First-Sentence Summary Rule}
\label{h:21:the-first-sentence-summary-rule}
The \textbf{first sentence} of a Javadoc comment (up to the first period followed by whitespace or a line break, in the classic algorithm — recent \code{javadoc} also respects an explicit \code{\{@\allowbreak{}summary~\allowbreak{}...\}} tag) becomes the \textbf{summary}: what shows up in package/class overview tables, method-summary tables, and search results. This means:

\begin{itemize}
\item 
Front-load the single most important fact about the method in that first sentence — assume it’s the \emph{only} sentence some readers will ever see.

\item 
Avoid a stray period inside the first sentence (e.g., in “e.g.” or a versioned class name) — it truncates the summary at the wrong place.

\end{itemize}

\begin{codebox}
\begin{Verbatim}[commandchars=\\\{\}]
\PY{c+cm}{/**}
\PY{c+cm}{ * Returns the customer\PYZsq{}s preferred shipping address, or the billing}
\PY{c+cm}{ * address if none is set. Falls back further to \PYZob{}@code null\PYZcb{} if the}
\PY{c+cm}{ * customer record itself has no addresses at all.}
\PY{c+cm}{ */}
\PY{k+kd}{public}\PY{+w}{ }\PY{n}{Address}\PY{+w}{ }\PY{n+nf}{preferredShippingAddress}\PY{p}{(}\PY{p}{)}\PY{+w}{ }\PY{p}{\PYZob{}}\PY{+w}{ }\PY{p}{.}\PY{p}{.}\PY{p}{.}\PY{+w}{ }\PY{p}{\PYZcb{}}
\end{Verbatim}
\end{codebox}

\begin{codebox}
\begin{Verbatim}[commandchars=\\\{\}]
\PY{c+c1}{// Bad: the period after \PYZdq{}e.g\PYZdq{} truncates the rendered summary to}
\PY{c+c1}{// \PYZdq{}Formats a duration, e.\PYZdq{} — nonsensical in the summary table.}
\PY{c+cm}{/**}
\PY{c+cm}{ * Formats a duration, e.g. \PYZdq{}3h 12m\PYZdq{}, using the default locale.}
\PY{c+cm}{ */}
\end{Verbatim}
\end{codebox}

\subsection[Documenting the Contract, Not the Implementation]{Documenting the Contract, Not the Implementation}
\label{h:21:documenting-the-contract-not-the-implementation}
The most important review distinction in this chapter: Javadoc documents the \textbf{contract} — what callers can rely on — not \emph{how} the current implementation happens to work. Implementation details belong in inline \code{//} comments inside the method body, if anywhere, because they can change without breaking any caller; contract statements cannot.

\begin{codebox}
\begin{Verbatim}[commandchars=\\\{\}]
\PY{c+c1}{// Bad: describes the current implementation, which callers should never depend on}
\PY{c+cm}{/**}
\PY{c+cm}{ * Looks up the user by scanning a HashMap keyed by user ID.}
\PY{c+cm}{ */}
\PY{k+kd}{public}\PY{+w}{ }\PY{n}{User}\PY{+w}{ }\PY{n+nf}{findById}\PY{p}{(}\PY{k+kt}{long}\PY{+w}{ }\PY{n}{id}\PY{p}{)}\PY{+w}{ }\PY{p}{\PYZob{}}\PY{+w}{ }\PY{p}{.}\PY{p}{.}\PY{p}{.}\PY{+w}{ }\PY{p}{\PYZcb{}}

\PY{c+c1}{// Good: describes the contract — inputs, outputs, and failure behavior}
\PY{c+cm}{/**}
\PY{c+cm}{ * Returns the user with the given ID.}
\PY{c+cm}{ *}
\PY{c+cm}{ * @param id the user ID; must be positive}
\PY{c+cm}{ * @return the matching user}
\PY{c+cm}{ * @throws NoSuchElementException if no user with this ID exists}
\PY{c+cm}{ */}
\PY{k+kd}{public}\PY{+w}{ }\PY{n}{User}\PY{+w}{ }\PY{n+nf}{findById}\PY{p}{(}\PY{k+kt}{long}\PY{+w}{ }\PY{n}{id}\PY{p}{)}\PY{+w}{ }\PY{p}{\PYZob{}}\PY{+w}{ }\PY{p}{.}\PY{p}{.}\PY{p}{.}\PY{+w}{ }\PY{p}{\PYZcb{}}
\end{Verbatim}
\end{codebox}

If the class later switches from a \code{HashMap} to a database call, the first comment becomes a lie that nobody remembers to fix; the second comment remains true regardless of the storage mechanism.

\subsection[Documenting Nullability, Thread Safety, and Side Effects]{Documenting Nullability, Thread Safety, and Side Effects}
\label{h:21:documenting-nullability-thread-safety-and-side-effects}
These three properties are exactly the ones that cannot be inferred from a method signature alone in Java (pre-nullability-annotations, and even with them, the \emph{reasoning} still benefits from prose), and they are exactly what reviewers look for missing:

\begin{codebox}
\begin{Verbatim}[commandchars=\\\{\}]
\PY{c+cm}{/**}
\PY{c+cm}{ * Finds a cached session for the given token.}
\PY{c+cm}{ *}
\PY{c+cm}{ * \PYZlt{}p\PYZgt{}\PYZlt{}b\PYZgt{}Nullability:\PYZlt{}/b\PYZgt{} returns \PYZob{}@code null\PYZcb{} if no session is cached}
\PY{c+cm}{ * for the token; never throws for an unknown token.}
\PY{c+cm}{ *}
\PY{c+cm}{ * \PYZlt{}p\PYZgt{}\PYZlt{}b\PYZgt{}Thread safety:\PYZlt{}/b\PYZgt{} this method is safe to call concurrently}
\PY{c+cm}{ * from multiple threads without external synchronization.}
\PY{c+cm}{ *}
\PY{c+cm}{ * \PYZlt{}p\PYZgt{}\PYZlt{}b\PYZgt{}Side effects:\PYZlt{}/b\PYZgt{} refreshes the token\PYZsq{}s last\PYZhy{}access timestamp}
\PY{c+cm}{ * as a side effect of a successful lookup, which resets its expiry.}
\PY{c+cm}{ *}
\PY{c+cm}{ * @param token the session token}
\PY{c+cm}{ * @return the cached \PYZob{}@link Session\PYZcb{}, or \PYZob{}@code null\PYZcb{} if not found}
\PY{c+cm}{ */}
\PY{k+kd}{public}\PY{+w}{ }\PY{n}{Session}\PY{+w}{ }\PY{n+nf}{findSession}\PY{p}{(}\PY{n}{String}\PY{+w}{ }\PY{n}{token}\PY{p}{)}\PY{+w}{ }\PY{p}{\PYZob{}}\PY{+w}{ }\PY{p}{.}\PY{p}{.}\PY{p}{.}\PY{+w}{ }\PY{p}{\PYZcb{}}
\end{Verbatim}
\end{codebox}

Where the project uses JSR-305 / \code{org.\allowbreak{}jspecify} annotations (\code{@\allowbreak{}Nullable}, \code{@\allowbreak{}NonNull}), the annotation is machine-checkable and should be preferred over prose alone — but prose explaining \emph{why} something can be null (e.g., “null before the account is activated”) still adds value an annotation cannot.

\subsection[package-info.java and module-info.java Docs]{\code{package-\allowbreak{}info.\allowbreak{}java} and \code{module-\allowbreak{}info.\allowbreak{}java} Docs}
\label{h:21:package-infojava-and-module-infojava-docs}
Package-level and module-level Javadoc live in dedicated files, not attached to any class, because a package/module is not itself a class.

\begin{codebox}
\begin{Verbatim}[commandchars=\\\{\}]
\PY{c+c1}{// package\PYZhy{}info.java}
\PY{c+cm}{/**}
\PY{c+cm}{ * Order pricing and discount calculation.}
\PY{c+cm}{ *}
\PY{c+cm}{ * \PYZlt{}p\PYZgt{}Classes in this package are responsible for computing the final}
\PY{c+cm}{ * price of an \PYZob{}@link com.example.orders.model.Order\PYZcb{}, including tax and}
\PY{c+cm}{ * promotional discounts. Nothing in this package performs persistence}
\PY{c+cm}{ * or network I/O.}
\PY{c+cm}{ *}
\PY{c+cm}{ * @since 1.0}
\PY{c+cm}{ */}
\PY{k+kn}{package}\PY{+w}{ }\PY{n+nn}{com.example.orders.pricing}\PY{p}{;}
\end{Verbatim}
\end{codebox}

\begin{codebox}
\begin{Verbatim}[commandchars=\\\{\}]
\PY{c+c1}{// module\PYZhy{}info.java}
\PY{c+cm}{/**}
\PY{c+cm}{ * Defines the order\PYZhy{}processing service and its public API.}
\PY{c+cm}{ *}
\PY{c+cm}{ * \PYZlt{}p\PYZgt{}Only \PYZob{}@code com.example.orders.api\PYZcb{} is exported; all other packages}
\PY{c+cm}{ * are internal implementation details and not accessible to other modules.}
\PY{c+cm}{ */}
\PY{n}{module}\PY{+w}{ }\PY{n}{com}\PY{p}{.}\PY{n+na}{example}\PY{p}{.}\PY{n+na}{orders}\PY{+w}{ }\PY{p}{\PYZob{}}
\PY{+w}{    }\PY{n}{requires}\PY{+w}{ }\PY{n}{java}\PY{p}{.}\PY{n+na}{sql}\PY{p}{;}
\PY{+w}{    }\PY{n}{exports}\PY{+w}{ }\PY{n}{com}\PY{p}{.}\PY{n+na}{example}\PY{p}{.}\PY{n+na}{orders}\PY{p}{.}\PY{n+na}{api}\PY{p}{;}
\PY{p}{\PYZcb{}}
\end{Verbatim}
\end{codebox}

\subsection[HTML in Javadoc and Markdown Javadoc (JEP 467, JDK 23+)]{HTML in Javadoc and Markdown Javadoc (JEP 467, JDK 23+)}
\label{h:21:html-in-javadoc-and-markdown-javadoc-jep-467-jdk-23}
Classic Javadoc comments are parsed as HTML fragments, so structuring text requires literal HTML tags — \code{\textless{}\allowbreak{}p\textgreater{}} for paragraphs, \code{\textless{}\allowbreak{}ul\textgreater{}}/\code{\textless{}\allowbreak{}li\textgreater{}} for lists, \code{\textless{}\allowbreak{}b\textgreater{}} for bold. This is verbose and easy to get wrong (unclosed tags silently break rendering, or worse, get swallowed by \code{-\allowbreak{}Xdoclint}, discussed next).

\begin{codebox}
\begin{Verbatim}[commandchars=\\\{\}]
\PY{c+cm}{/**}
\PY{c+cm}{ * Represents an order line item.}
\PY{c+cm}{ *}
\PY{c+cm}{ * \PYZlt{}p\PYZgt{}Each line item has:}
\PY{c+cm}{ * \PYZlt{}ul\PYZgt{}}
\PY{c+cm}{ *   \PYZlt{}li\PYZgt{}a \PYZob{}@code productId\PYZcb{}\PYZlt{}/li\PYZgt{}}
\PY{c+cm}{ *   \PYZlt{}li\PYZgt{}a \PYZob{}@code quantity\PYZcb{}, always positive\PYZlt{}/li\PYZgt{}}
\PY{c+cm}{ *   \PYZlt{}li\PYZgt{}a \PYZob{}@code unitPrice\PYZcb{}\PYZlt{}/li\PYZgt{}}
\PY{c+cm}{ * \PYZlt{}/ul\PYZgt{}}
\PY{c+cm}{ */}
\PY{k+kd}{public}\PY{+w}{ }\PY{k+kd}{class} \PY{n+nc}{LineItem}\PY{+w}{ }\PY{p}{\PYZob{}}\PY{+w}{ }\PY{p}{.}\PY{p}{.}\PY{p}{.}\PY{+w}{ }\PY{p}{\PYZcb{}}
\end{Verbatim}
\end{codebox}

\textbf{JEP 467} (finalized in JDK 23) allows Javadoc comments to be written in \textbf{Markdown} instead, using \code{///} line comments rather than \code{/**\allowbreak{}~\allowbreak{}*/} blocks:

\begin{codebox}
\begin{Verbatim}[commandchars=\\\{\}]
\PY{c+c1}{/// Represents an order line item.}
\PY{c+c1}{///}
\PY{c+c1}{/// Each line item has:}
\PY{c+c1}{/// \PYZhy{} a `productId`}
\PY{c+c1}{/// \PYZhy{} a `quantity`, always positive}
\PY{c+c1}{/// \PYZhy{} a `unitPrice`}
\PY{c+c1}{///}
\PY{c+c1}{/// See [LineItem\PYZsh{}total()] for how the extended price is computed.}
\PY{k+kd}{public}\PY{+w}{ }\PY{k+kd}{class} \PY{n+nc}{LineItem}\PY{+w}{ }\PY{p}{\PYZob{}}
\PY{+w}{    }\PY{c+c1}{/// Returns the extended price: `unitPrice * quantity`.}
\PY{+w}{    }\PY{k+kd}{public}\PY{+w}{ }\PY{n}{BigDecimal}\PY{+w}{ }\PY{n+nf}{total}\PY{p}{(}\PY{p}{)}\PY{+w}{ }\PY{p}{\PYZob{}}\PY{+w}{ }\PY{p}{.}\PY{p}{.}\PY{p}{.}\PY{+w}{ }\PY{p}{\PYZcb{}}
\PY{p}{\PYZcb{}}
\end{Verbatim}
\end{codebox}

Markdown Javadoc is substantially easier to read in raw source form (in an IDE or a diff) and eliminates most of the HTML-escaping foot-guns, at the cost of requiring JDK 23+ tooling to render.

\subsection[Generating Javadoc]{Generating Javadoc}
\label{h:21:generating-javadoc}
\begin{codebox}
\begin{Verbatim}[commandchars=\\\{\}]
\PY{c+c1}{\PYZsh{} Directly with the javadoc tool}
javadoc\PY{+w}{ }\PYZhy{}d\PY{+w}{ }docs\PY{+w}{ }\PYZhy{}sourcepath\PY{+w}{ }src/main/java\PY{+w}{ }\PYZhy{}subpackages\PY{+w}{ }com.example.orders

\PY{c+c1}{\PYZsh{} Maven}
mvn\PY{+w}{ }javadoc:javadoc\PY{+w}{          }\PY{c+c1}{\PYZsh{} generates target/site/apidocs}
mvn\PY{+w}{ }javadoc:jar\PY{+w}{               }\PY{c+c1}{\PYZsh{} packages it as a javadoc.jar (needed to publish to Maven Central)}

\PY{c+c1}{\PYZsh{} Gradle}
./gradlew\PY{+w}{ }javadoc\PY{+w}{             }\PY{c+c1}{\PYZsh{} generates build/docs/javadoc}
\end{Verbatim}
\end{codebox}

\begin{codebox}
\begin{Verbatim}[commandchars=\\\{\}]
\PY{c+cm}{\PYZlt{}!\PYZhy{}\PYZhy{} pom.xml: enforce doc quality and produce a javadoc jar for publishing \PYZhy{}\PYZhy{}\PYZgt{}}
\PY{n+nt}{\PYZlt{}plugin}\PY{n+nt}{\PYZgt{}}
\PY{+w}{  }\PY{n+nt}{\PYZlt{}groupId}\PY{n+nt}{\PYZgt{}}org.apache.maven.plugins\PY{n+nt}{\PYZlt{}/groupId\PYZgt{}}
\PY{+w}{  }\PY{n+nt}{\PYZlt{}artifactId}\PY{n+nt}{\PYZgt{}}maven\PYZhy{}javadoc\PYZhy{}plugin\PY{n+nt}{\PYZlt{}/artifactId\PYZgt{}}
\PY{+w}{  }\PY{n+nt}{\PYZlt{}version}\PY{n+nt}{\PYZgt{}}3.8.0\PY{n+nt}{\PYZlt{}/version\PYZgt{}}
\PY{+w}{  }\PY{n+nt}{\PYZlt{}configuration}\PY{n+nt}{\PYZgt{}}
\PY{+w}{    }\PY{n+nt}{\PYZlt{}doclint}\PY{n+nt}{\PYZgt{}}all,\PYZhy{}missing\PY{n+nt}{\PYZlt{}/doclint\PYZgt{}}\PY{+w}{  }\PY{c+cm}{\PYZlt{}!\PYZhy{}\PYZhy{} validate HTML/tags; don\PYZsq{}t require every element documented \PYZhy{}\PYZhy{}\PYZgt{}}
\PY{+w}{    }\PY{n+nt}{\PYZlt{}failOnWarning}\PY{n+nt}{\PYZgt{}}true\PY{n+nt}{\PYZlt{}/failOnWarning\PYZgt{}}
\PY{+w}{  }\PY{n+nt}{\PYZlt{}/configuration\PYZgt{}}
\PY{+w}{  }\PY{n+nt}{\PYZlt{}executions}\PY{n+nt}{\PYZgt{}}
\PY{+w}{    }\PY{n+nt}{\PYZlt{}execution}\PY{n+nt}{\PYZgt{}}
\PY{+w}{      }\PY{n+nt}{\PYZlt{}id}\PY{n+nt}{\PYZgt{}}attach\PYZhy{}javadocs\PY{n+nt}{\PYZlt{}/id\PYZgt{}}
\PY{+w}{      }\PY{n+nt}{\PYZlt{}goals}\PY{n+nt}{\PYZgt{}}\PY{n+nt}{\PYZlt{}goal}\PY{n+nt}{\PYZgt{}}jar\PY{n+nt}{\PYZlt{}/goal\PYZgt{}}\PY{n+nt}{\PYZlt{}/goals\PYZgt{}}
\PY{+w}{    }\PY{n+nt}{\PYZlt{}/execution\PYZgt{}}
\PY{+w}{  }\PY{n+nt}{\PYZlt{}/executions\PYZgt{}}
\PY{n+nt}{\PYZlt{}/plugin\PYZgt{}}
\end{Verbatim}
\end{codebox}

\begin{codebox}
\begin{Verbatim}[commandchars=\\\{\}]
\PY{c+c1}{// build.gradle.kts}
\PY{n}{tasks}\PY{p}{.}\PY{n+na}{javadoc}\PY{+w}{ }\PY{p}{\PYZob{}}
\PY{+w}{    }\PY{p}{(}\PY{n}{options}\PY{+w}{ }\PY{k}{as}\PY{+w}{ }\PY{n}{StandardJavadocDocletOptions}\PY{p}{)}\PY{p}{.}\PY{n+na}{apply}\PY{+w}{ }\PY{p}{\PYZob{}}
\PY{+w}{        }\PY{n}{addStringOption}\PY{p}{(}\PY{l+s}{\PYZdq{}}\PY{l+s}{Xdoclint:all,\PYZhy{}missing}\PY{l+s}{\PYZdq{}}\PY{p}{,}\PY{+w}{ }\PY{l+s}{\PYZdq{}}\PY{l+s}{\PYZhy{}quiet}\PY{l+s}{\PYZdq{}}\PY{p}{)}
\PY{+w}{    }\PY{p}{\PYZcb{}}
\PY{p}{\PYZcb{}}
\end{Verbatim}
\end{codebox}

\subsection[-Xdoclint and Treating Doc Warnings as Errors]{\code{-\allowbreak{}Xdoclint} and Treating Doc Warnings as Errors}
\label{h:21:-xdoclint-and-treating-doc-warnings-as-errors}
\code{-\allowbreak{}Xdoclint} is the \code{javadoc} tool’s built-in linter: it flags malformed HTML, \code{\{@\allowbreak{}link\}} references to symbols that don’t exist, missing \code{@\allowbreak{}param}/\code{@\allowbreak{}return}/\code{@\allowbreak{}throws} tags, and other structural problems, without needing a separate static-analysis tool. Since JDK 8, \code{javac} itself also runs a subset of doclint checks and can emit warnings during \textbf{compilation}, not just during Javadoc generation.

\begin{codebox}
\begin{Verbatim}[commandchars=\\\{\}]
javadoc\PY{+w}{ }\PYZhy{}Xdoclint:all\PY{+w}{ }\PYZhy{}d\PY{+w}{ }docs\PY{+w}{ }\PYZhy{}sourcepath\PY{+w}{ }src/main/java\PY{+w}{ }com.example.orders
javac\PY{+w}{ }\PYZhy{}Xdoclint:all\PY{+w}{ }\PYZhy{}d\PY{+w}{ }out\PY{+w}{ }src/main/java/com/example/orders/*.java
\end{Verbatim}
\end{codebox}

Common \code{-\allowbreak{}Xdoclint} groups you can enable/disable independently: \code{accessibility}, \code{html}, \code{missing}, \code{reference}, \code{syntax}. A pragmatic default that reviewers commonly ask for is \code{all,-\allowbreak{}missing} — validate everything that’s present, but don’t force every single public element to have a comment (that requirement tends to produce exactly the restating-the-name comments discussed below).

Treating doc warnings as build errors (\code{failOnWarning} in Maven, or \code{-\allowbreak{}Werror}-equivalent doclint options in Gradle) prevents doc rot: a broken \code{\{@\allowbreak{}link\}} to a renamed method, or a \code{\textless{}\allowbreak{}p\textgreater{}} tag that was never closed, stays broken forever unless something makes the build fail on it.

\subsection[Publishing Javadoc]{Publishing Javadoc}
\label{h:21:publishing-javadoc}
For a library, Javadoc is typically published alongside the release artifact as a \code{-\allowbreak{}javadoc.\allowbreak{}jar} classifier (required by Maven Central), and/or as a static HTML site (e.g., via GitHub Pages) generated by CI on release. \code{javadoc.\allowbreak{}io} can render any artifact already on Maven Central without any extra publishing step, by resolving the \code{-\allowbreak{}javadoc.\allowbreak{}jar} classifier directly.

\begin{codebox}
\begin{Verbatim}[commandchars=\\\{\}]
mvn\PY{+w}{ }javadoc:jar\PY{+w}{ }deploy\PY{+w}{    }\PY{c+c1}{\PYZsh{} produces and deploys my\PYZhy{}lib\PYZhy{}1.4.0\PYZhy{}javadoc.jar alongside the main jar}
\end{Verbatim}
\end{codebox}

\subsection[When a Javadoc Comment Is Worse Than No Comment]{When a Javadoc Comment Is Worse Than No Comment}
\label{h:21:when-a-javadoc-comment-is-worse-than-no-comment}
A Javadoc comment that merely restates the method/parameter name in English adds visual noise, gives the false impression that the method is documented, and — worst of all — will drift out of sync with the real behavior over time because nobody feels the need to update “obvious” prose. It is strictly worse than no comment, because a missing comment at least invites the reader to go read the code.

\begin{codebox}
\begin{Verbatim}[commandchars=\\\{\}]
\PY{c+c1}{// Bad: pure noise, restates the signature}
\PY{c+cm}{/**}
\PY{c+cm}{ * Gets the name.}
\PY{c+cm}{ * @return the name}
\PY{c+cm}{ */}
\PY{k+kd}{public}\PY{+w}{ }\PY{n}{String}\PY{+w}{ }\PY{n+nf}{getName}\PY{p}{(}\PY{p}{)}\PY{+w}{ }\PY{p}{\PYZob{}}\PY{+w}{ }\PY{k}{return}\PY{+w}{ }\PY{n}{name}\PY{p}{;}\PY{+w}{ }\PY{p}{\PYZcb{}}

\PY{c+c1}{// Bad: restates behavior visible from the method body one line below,}
\PY{c+c1}{// adds nothing about the contract (can it return null? is it idempotent?)}
\PY{c+cm}{/**}
\PY{c+cm}{ * Sets the value.}
\PY{c+cm}{ * @param value the value}
\PY{c+cm}{ */}
\PY{k+kd}{public}\PY{+w}{ }\PY{k+kt}{void}\PY{+w}{ }\PY{n+nf}{setValue}\PY{p}{(}\PY{k+kt}{int}\PY{+w}{ }\PY{n}{value}\PY{p}{)}\PY{+w}{ }\PY{p}{\PYZob{}}\PY{+w}{ }\PY{k}{this}\PY{p}{.}\PY{n+na}{value}\PY{+w}{ }\PY{o}{=}\PY{+w}{ }\PY{n}{value}\PY{p}{;}\PY{+w}{ }\PY{p}{\PYZcb{}}
\end{Verbatim}
\end{codebox}

\begin{codebox}
\begin{Verbatim}[commandchars=\\\{\}]
\PY{c+c1}{// Good: worth documenting, because the contract isn\PYZsq{}t obvious from the signature}
\PY{c+cm}{/**}
\PY{c+cm}{ * Returns the display name shown in the UI, falling back to the account\PYZsq{}s}
\PY{c+cm}{ * email address if the user has never set one.}
\PY{c+cm}{ *}
\PY{c+cm}{ * @return the display name; never \PYZob{}@code null\PYZcb{}}
\PY{c+cm}{ */}
\PY{k+kd}{public}\PY{+w}{ }\PY{n}{String}\PY{+w}{ }\PY{n+nf}{getName}\PY{p}{(}\PY{p}{)}\PY{+w}{ }\PY{p}{\PYZob{}}\PY{+w}{ }\PY{k}{return}\PY{+w}{ }\PY{n}{name}\PY{+w}{ }\PY{o}{!}\PY{o}{=}\PY{+w}{ }\PY{k+kc}{null}\PY{+w}{ }\PY{o}{?}\PY{+w}{ }\PY{n}{name}\PY{+w}{ }\PY{p}{:}\PY{+w}{ }\PY{n}{email}\PY{p}{;}\PY{+w}{ }\PY{p}{\PYZcb{}}
\end{Verbatim}
\end{codebox}

A simple review heuristic: if you deleted the Javadoc comment and could still write the exact same sentence just by reading the method name and signature, the comment is not documenting anything and should either be deleted or expanded to cover the actual contract (nullability, exceptions, thread-safety, or a non-obvious side effect).

\section[Common Code-Review Interview Pitfalls]{Common Code-Review Interview Pitfalls}
\label{h:21:common-code-review-interview-pitfalls}
\begin{enumerate}
\item 
\textbf{Committing \code{target/} or \code{build/} output directories to version control.}
Why it matters: generated artifacts bloat the repository, go stale relative to source, and cause merge noise; they should always be build-tool-generated, never hand-edited or tracked.

\begin{codebox}
\begin{Verbatim}
# .gitignore
target/
build/
\end{Verbatim}
\end{codebox}

\item 
\textbf{Not committing the Maven/Gradle Wrapper (\code{mvnw}/\code{gradlew}).}
Why it matters: without a pinned wrapper, “works on my machine” becomes literal — a teammate or CI runner with a different globally installed build-tool version can produce a different, sometimes-broken build.

\begin{codebox}
\begin{Verbatim}[commandchars=\\\{\}]
mvn\PY{+w}{ }wrapper:wrapper\PY{+w}{ }\PYZhy{}Dmaven\PY{o}{=}\PY{l+m}{3}.9.6\PY{+w}{   }\PY{c+c1}{\PYZsh{} generate and commit mvnw, mvnw.cmd, .mvn/}
\end{Verbatim}
\end{codebox}

\item 
\textbf{Using a version range (\code{[\allowbreak{}1.\allowbreak{}0,\allowbreak{}2.\allowbreak{}0)}) instead of a pinned version for a dependency.}
Why it matters: the resolved version can silently change between builds as new releases are published, breaking reproducibility — the exact opposite of what a lockfile-free ecosystem needs to stay predictable.

\begin{codebox}
\begin{Verbatim}[commandchars=\\\{\}]
\PY{c+cm}{\PYZlt{}!\PYZhy{}\PYZhy{} Before \PYZhy{}\PYZhy{}\PYZgt{}}
\PY{n+nt}{\PYZlt{}version}\PY{n+nt}{\PYZgt{}}[1.0,2.0)\PY{n+nt}{\PYZlt{}/version\PYZgt{}}
\PY{c+cm}{\PYZlt{}!\PYZhy{}\PYZhy{} After \PYZhy{}\PYZhy{}\PYZgt{}}
\PY{n+nt}{\PYZlt{}version}\PY{n+nt}{\PYZgt{}}1.4.2\PY{n+nt}{\PYZlt{}/version\PYZgt{}}
\end{Verbatim}
\end{codebox}

\item 
\textbf{Declaring a dependency as \code{compile}/\code{api} when it’s only used internally.}
Why it matters: it leaks the dependency onto every consumer’s compile classpath, forcing consumers to recompile when an internal-only dependency changes version, and can leak internal types through the public API by accident.

\begin{codebox}
\begin{Verbatim}[commandchars=\\\{\}]
\PY{c+c1}{// Before}
\PY{n}{api}\PY{p}{(}\PY{l+s}{\PYZdq{}}\PY{l+s}{org.apache.commons:commons\PYZhy{}lang3:3.14.0}\PY{l+s}{\PYZdq{}}\PY{p}{)}
\PY{c+c1}{// After — used only inside this module\PYZsq{}s private helpers}
\PY{n}{implementation}\PY{p}{(}\PY{l+s}{\PYZdq{}}\PY{l+s}{org.apache.commons:commons\PYZhy{}lang3:3.14.0}\PY{l+s}{\PYZdq{}}\PY{p}{)}
\end{Verbatim}
\end{codebox}

\item 
\textbf{Using \code{-\allowbreak{}Dmaven.\allowbreak{}test.\allowbreak{}skip=\allowbreak{}true} in CI instead of \code{-\allowbreak{}DskipTests} (or, worse, skipping tests routinely at all).}
Why it matters: \code{maven.\allowbreak{}test.\allowbreak{}skip} doesn’t even compile the test sources, so a broken test file can merge without anyone noticing until someone runs tests locally — a much worse failure mode than merely not running passing tests.

\begin{codebox}
\begin{Verbatim}[commandchars=\\\{\}]
\PY{c+c1}{\PYZsh{} Before (hides compile errors in tests)}
mvn\PY{+w}{ }install\PY{+w}{ }\PYZhy{}Dmaven.test.skip\PY{o}{=}\PY{n+nb}{true}
\PY{c+c1}{\PYZsh{} After (at least verifies tests compile)}
mvn\PY{+w}{ }install\PY{+w}{ }\PYZhy{}DskipTests
\end{Verbatim}
\end{codebox}

\item 
\textbf{Relying on Maven’s nearest-wins transitive resolution without checking \code{dependency:\allowbreak{}tree}.}
Why it matters: nearest-wins picks a version based on graph depth, not correctness — it can silently select an old or vulnerable version of a transitive dependency, and the only way to know is to actually look.

\begin{codebox}
\begin{Verbatim}[commandchars=\\\{\}]
mvn\PY{+w}{ }dependency:tree\PY{+w}{ }\PYZhy{}Dverbose
\end{Verbatim}
\end{codebox}

\item 
\textbf{Writing a Javadoc comment that restates the method name instead of documenting the contract.}
Why it matters: it looks like documentation without providing any information beyond the signature, and gives false confidence that behavior (nullability, exceptions, thread-safety) is documented when it isn’t.

\begin{codebox}
\begin{Verbatim}[commandchars=\\\{\}]
\PY{c+c1}{// Before}
\PY{c+cm}{/** Gets the name. @return the name */}
\PY{k+kd}{public}\PY{+w}{ }\PY{n}{String}\PY{+w}{ }\PY{n+nf}{getName}\PY{p}{(}\PY{p}{)}\PY{+w}{ }\PY{p}{\PYZob{}}\PY{+w}{ }\PY{p}{.}\PY{p}{.}\PY{p}{.}\PY{+w}{ }\PY{p}{\PYZcb{}}
\PY{c+c1}{// After}
\PY{c+cm}{/** Returns the display name, or the email address if none was set. */}
\PY{k+kd}{public}\PY{+w}{ }\PY{n}{String}\PY{+w}{ }\PY{n+nf}{getName}\PY{p}{(}\PY{p}{)}\PY{+w}{ }\PY{p}{\PYZob{}}\PY{+w}{ }\PY{p}{.}\PY{p}{.}\PY{p}{.}\PY{+w}{ }\PY{p}{\PYZcb{}}
\end{Verbatim}
\end{codebox}

\item 
\textbf{Javadoc describing the current implementation instead of the caller-visible contract.}
Why it matters: implementation-describing comments become lies the moment the implementation changes (e.g., swapping a \code{HashMap} for a database call), and nobody remembers to update prose that isn’t checked by the compiler.

\begin{codebox}
\begin{Verbatim}[commandchars=\\\{\}]
\PY{c+c1}{// Before}
\PY{c+cm}{/** Looks up the user by scanning an in\PYZhy{}memory HashMap. */}
\PY{c+c1}{// After}
\PY{c+cm}{/** Returns the user with the given ID. @throws NoSuchElementException if not found */}
\end{Verbatim}
\end{codebox}

\item 
\textbf{A broken \code{\{@\allowbreak{}link\}} reference to a renamed or deleted method/class.}
Why it matters: it silently degrades documentation quality (a dead link in generated docs) and usually indicates the comment wasn’t updated alongside a refactor; \code{-\allowbreak{}Xdoclint:\allowbreak{}reference} catches this automatically if enabled and enforced in CI.

\begin{codebox}
\begin{Verbatim}[commandchars=\\\{\}]
javadoc\PY{+w}{ }\PYZhy{}Xdoclint:reference\PY{+w}{ }\PYZhy{}d\PY{+w}{ }docs\PY{+w}{ }\PYZhy{}sourcepath\PY{+w}{ }src/main/java\PY{+w}{ }com.example.orders
\end{Verbatim}
\end{codebox}

\item 
\textbf{Mixing up \code{provided}/\code{compileOnly} with \code{runtime}/\code{runtimeOnly}, causing a \code{NoClassDefFoundError} at deploy time or a bloated fat jar.}
Why it matters: \code{provided}/\code{compileOnly} dependencies (e.g., servlet API, Lombok) are intentionally excluded from the packaged artifact because the runtime environment supplies them — packaging them anyway (wrong scope) can cause classloading conflicts; the reverse mistake (using \code{provided} for something actually needed at runtime) causes a missing-class crash in production.

\begin{codebox}
\begin{Verbatim}[commandchars=\\\{\}]
\PY{c+cm}{\PYZlt{}!\PYZhy{}\PYZhy{} Before: servlet\PYZhy{}api is needed only to compile against, container supplies it \PYZhy{}\PYZhy{}\PYZgt{}}
\PY{n+nt}{\PYZlt{}scope}\PY{n+nt}{\PYZgt{}}compile\PY{n+nt}{\PYZlt{}/scope\PYZgt{}}
\PY{c+cm}{\PYZlt{}!\PYZhy{}\PYZhy{} After \PYZhy{}\PYZhy{}\PYZgt{}}
\PY{n+nt}{\PYZlt{}scope}\PY{n+nt}{\PYZgt{}}provided\PY{n+nt}{\PYZlt{}/scope\PYZgt{}}
\end{Verbatim}
\end{codebox}

\item 
\textbf{Using \code{maven-\allowbreak{}assembly-\allowbreak{}plugin} (or an unconfigured fat jar) for a Spring Boot / service-loader-heavy app instead of \code{maven-\allowbreak{}shade-\allowbreak{}plugin} with transformers.}
Why it matters: assembly’s naive concatenation can silently drop or corrupt merged \code{META-\allowbreak{}INF/\allowbreak{}services} files (used by \code{ServiceLoader}, JDBC drivers, logging backends), producing runtime failures that only show up when the fat jar is actually run, not when it’s built.

\begin{codebox}
\begin{Verbatim}[commandchars=\\\{\}]
\PY{n+nt}{\PYZlt{}transformer}\PY{+w}{ }\PY{n+na}{implementation=}\PY{l+s}{\PYZdq{}org.apache.maven.plugins.shade.resource.ServicesResourceTransformer\PYZdq{}}\PY{n+nt}{/\PYZgt{}}
\end{Verbatim}
\end{codebox}

\item 
\textbf{A custom Gradle task using \code{doLast~\allowbreak{}\{\}} with no declared inputs/outputs, breaking incremental build and the build cache.}
Why it matters: Gradle can only skip a task as \code{UP-\allowbreak{}TO-\allowbreak{}DATE} or restore it from cache if it knows what the task’s inputs and outputs are; an undeclared \code{doLast} closure always reruns, silently making every build slower than it needs to be.

\begin{codebox}
\begin{Verbatim}[commandchars=\\\{\}]
\PY{c+c1}{// Before: always reruns, never cacheable}
\PY{n}{tasks}\PY{p}{.}\PY{n+na}{register}\PY{p}{(}\PY{l+s}{\PYZdq{}}\PY{l+s}{generateDocs}\PY{l+s}{\PYZdq{}}\PY{p}{)}\PY{+w}{ }\PY{p}{\PYZob{}}\PY{+w}{ }\PY{n}{doLast}\PY{+w}{ }\PY{p}{\PYZob{}}\PY{+w}{ }\PY{c+cm}{/* ... */}\PY{+w}{ }\PY{p}{\PYZcb{}}\PY{+w}{ }\PY{p}{\PYZcb{}}
\PY{c+c1}{// After: typed task with declared @InputFiles/@OutputDirectory}
\PY{n}{tasks}\PY{p}{.}\PY{n+na}{register}\PY{o}{\PYZlt{}}\PY{n}{GenerateDocsTask}\PY{o}{\PYZgt{}}\PY{p}{(}\PY{l+s}{\PYZdq{}}\PY{l+s}{generateDocs}\PY{l+s}{\PYZdq{}}\PY{p}{)}
\end{Verbatim}
\end{codebox}

\item 
\textbf{Manually keeping dependency versions in sync across Maven modules or Gradle subprojects instead of using a BOM / version catalog.}
Why it matters: without a single source of truth for versions, modules drift apart over time, producing exactly the kind of version conflicts \code{dependency:\allowbreak{}tree}/\code{dependencyInsight} are needed to untangle later — the bug is prevented far more cheaply than it’s debugged.

\begin{codebox}
\begin{Verbatim}[commandchars=\\\{\}]
\PY{c+c1}{\PYZsh{} gradle/libs.versions.toml — one place to bump a version for every subproject}
\PY{k}{[}\PY{k}{versions}\PY{k}{]}
\PY{n}{jackson}\PY{+w}{ }\PY{o}{=}\PY{+w}{ }\PY{l+s+s2}{\PYZdq{}}\PY{l+s+s2}{2.17.1}\PY{l+s+s2}{\PYZdq{}}
\end{Verbatim}
\end{codebox}

\item 
\textbf{Declaring \code{\textless{}\allowbreak{}release\textgreater{}} inconsistently across modules (or using the old \code{source}/\code{target} pair instead of \code{release}).}
Why it matters: \code{source}/\code{target} alone can compile against a \emph{newer} JDK’s bootclasspath while claiming an older language level, letting code accidentally reference APIs unavailable on the actual target runtime; \code{\textless{}\allowbreak{}release\textgreater{}} locks both together and fails the build if that happens.

\begin{codebox}
\begin{Verbatim}[commandchars=\\\{\}]
\PY{c+cm}{\PYZlt{}!\PYZhy{}\PYZhy{} Before \PYZhy{}\PYZhy{}\PYZgt{}}
\PY{n+nt}{\PYZlt{}source}\PY{n+nt}{\PYZgt{}}17\PY{n+nt}{\PYZlt{}/source\PYZgt{}}
\PY{n+nt}{\PYZlt{}target}\PY{n+nt}{\PYZgt{}}17\PY{n+nt}{\PYZlt{}/target\PYZgt{}}
\PY{c+cm}{\PYZlt{}!\PYZhy{}\PYZhy{} After \PYZhy{}\PYZhy{}\PYZgt{}}
\PY{n+nt}{\PYZlt{}release}\PY{n+nt}{\PYZgt{}}21\PY{n+nt}{\PYZlt{}/release\PYZgt{}}
\end{Verbatim}
\end{codebox}

\item 
\textbf{Suppressing all Javadoc/doclint warnings globally instead of fixing or narrowly excusing them.}
Why it matters: a blanket \code{-\allowbreak{}Xdoclint:\allowbreak{}none} (or \code{doclint=\allowbreak{}none}) hides genuinely broken documentation (bad HTML, dead \code{\{@\allowbreak{}link\}}s, missing \code{@\allowbreak{}throws} for a checked exception) right alongside harmless style nitpicks, defeating the entire purpose of running the linter in the first place.

\begin{codebox}
\begin{Verbatim}[commandchars=\\\{\}]
\PY{c+cm}{\PYZlt{}!\PYZhy{}\PYZhy{} Before: disables everything \PYZhy{}\PYZhy{}\PYZgt{}}
\PY{n+nt}{\PYZlt{}doclint}\PY{n+nt}{\PYZgt{}}none\PY{n+nt}{\PYZlt{}/doclint\PYZgt{}}
\PY{c+cm}{\PYZlt{}!\PYZhy{}\PYZhy{} After: validates structure/links, only relaxes the \PYZdq{}every element must have a comment\PYZdq{} rule \PYZhy{}\PYZhy{}\PYZgt{}}
\PY{n+nt}{\PYZlt{}doclint}\PY{n+nt}{\PYZgt{}}all,\PYZhy{}missing\PY{n+nt}{\PYZlt{}/doclint\PYZgt{}}
\end{Verbatim}
\end{codebox}

\end{enumerate}


\setcounter{chapter}{21}
\chapter{Code-Review Interview Focused Topics}
\label{chap:22}
\label{h:22:22-code-review-interview-focused-topics}
A code-review interview hands you a snippet — sometimes a whole small class, sometimes a diff — and asks: “what’s wrong with this?” It is a different skill from solving an algorithm puzzle on a whiteboard. Nobody is grading whether you can invent clever code; they are grading whether you can \textbf{read someone else’s code critically}, recognize the same handful of mistakes that break production systems every day, and talk about them the way a helpful senior engineer would in a real pull request.

Interviewers (and good reviewers in general) judge findings against a rough priority order. Learn this order and use it to decide what to say first:

\begin{enumerate}
\item 
\textbf{Correctness} — does the code do what it claims? Off-by-one errors, wrong boolean logic, unhandled edge cases, silent data corruption.

\item 
\textbf{Safety} — will it corrupt shared state, leak a resource, deadlock, or blow up under concurrent/production load, even if the “happy path” looks fine on a single thread?

\item 
\textbf{Design} — is it maintainable and testable? Right abstraction, right exception type, right collection, sensible API shape?

\item 
\textbf{Performance} — is it needlessly slow or memory-hungry, once correctness and safety are already settled?

\item 
\textbf{Style} — naming, formatting, and other things a linter could catch.

\end{enumerate}

If you spend the first three minutes discussing a variable name (style) while a \code{HashMap} is being mutated from two threads two lines above it (safety), you have failed the interview even if every word about naming was correct. Everything in this chapter is organized around that ordering — find the highest-severity issue first, say it first, and only get to nits if there’s time left.

This chapter assumes you have already read the earlier chapters on the underlying mechanics — it does not re-derive \emph{why} \code{HashMap} isn’t thread-safe or \emph{why} checked exceptions exist. It focuses purely on the reviewer’s lens: what to look for, how to phrase it, and worked examples.

\section[The Review Checklist]{The Review Checklist}
\label{h:22:the-review-checklist}
Run through this whenever you’re handed a snippet, top to bottom, in this order. It mirrors the correctness → safety → design → performance → style priority from the intro.

\subsection[Correctness]{Correctness}
\label{h:22:correctness}
\begin{itemize}
\item 
[ ] Does every branch (including the ones with no \code{else}) do the right thing? What happens on empty input, \code{null}, zero, negative numbers, max values?

\item 
[ ] Are loop bounds right? Off-by-one on \code{\textless{}} vs \code{\textless{}=}, first/last index, empty collections.

\item 
[ ] Is \code{==} used where \code{.\allowbreak{}equals()} was needed (boxed types, \code{String}, custom objects)?

\item 
[ ] Does the method actually do what its name and signature promise?

\item 
[ ] Are there unhandled edge cases the tests (if any) don’t cover?

\end{itemize}

\subsection[Thread Safety \& Concurrency]{Thread Safety \& Concurrency}
\label{h:22:thread-safety--concurrency}
\begin{itemize}
\item 
[ ] Is any state \textbf{shared} across threads (static fields, singletons, fields on a class handed to an executor)?

\item 
[ ] If shared and mutable — is access synchronized, or is it a concurrent-safe type (\code{ConcurrentHashMap}, \code{AtomicInteger}, \code{volatile}, immutable value)?

\item 
[ ] Does construction fully finish before the object is published to another thread (no leaking \code{this})?

\item 
[ ] Are there lazy-initialized singletons, caches, or \code{SimpleDateFormat}-style non-thread-safe fields stored \code{static}?

\item 
[ ] Is every \code{ExecutorService}/thread pool shut down somewhere? Is every lock released on all paths, including exceptions?

\item 
[ ] Any risk of deadlock from inconsistent lock ordering, or of double-checked locking without \code{volatile}?

\end{itemize}

\subsection[Resource \& Exception Management]{Resource \& Exception Management}
\label{h:22:resource--exception-management}
\begin{itemize}
\item 
[ ] Is every \code{Closeable}/\code{AutoCloseable} (stream, connection, statement, lock) closed on \textbf{every} path, including exceptions? Prefer try-with-resources — see \hyperref[chap:6]{Chapter~6, \emph{Exception Handling}}.

\item 
[ ] Is any \code{catch} block empty, or does it just log-and-swallow without rethrowing or handling?

\item 
[ ] Is the \code{catch} scope too wide — wrapping code that can throw for unrelated reasons?

\item 
[ ] Is the original cause preserved when wrapping an exception (\code{new~\allowbreak{}Foo(\allowbreak{}\textquotedbl{}msg\textquotedbl{},\allowbreak{}~\allowbreak{}e)}, not \code{new~\allowbreak{}Foo(\allowbreak{}\textquotedbl{}msg\textquotedbl{})})?

\item 
[ ] Are checked vs. unchecked exceptions used appropriately at the API boundary?

\end{itemize}

\subsection[Collections \& Data Structures]{Collections \& Data Structures}
\label{h:22:collections--data-structures}
\begin{itemize}
\item 
[ ] Is the chosen collection right for the access pattern (lookup-heavy vs. order-sensitive vs. dedupe vs. FIFO)? See \hyperref[chap:7]{Chapter~7, \emph{Collections Framework}}.

\item 
[ ] Are keys used in a \code{HashMap}/\code{HashSet} immutable, with a consistent \code{equals}/\code{hashCode} pair?

\item 
[ ] Any \code{equals()} override missing a matching \code{hashCode()} override?

\item 
[ ] Any unbounded growth — cache, list, map that is only ever added to and never trimmed or evicted?

\item 
[ ] Any \code{computeIfAbsent}/\code{merge} call whose lambda has side effects or re-enters the same map?

\end{itemize}

\subsection[API Design]{API Design}
\label{h:22:api-design}
\begin{itemize}
\item 
[ ] Can any public parameter or return value legally be \code{null}? Is that documented, or should it be \code{Optional}?

\item 
[ ] Are parameters validated at the boundary (fail fast) instead of failing deep inside with a confusing error?

\item 
[ ] Is the class needlessly mutable when it could be immutable?

\item 
[ ] Do overloads behave consistently, and is the naming unambiguous about side effects (\code{getX} vs \code{computeX})?

\item 
[ ] Would this change break existing callers (removed method, changed semantics, widened exceptions)?

\end{itemize}

\subsection[Performance \& Memory]{Performance \& Memory}
\label{h:22:performance--memory}
\begin{itemize}
\item 
[ ] Any O(n²) work disguised as “just a loop inside a loop” (nested lookups, repeated linear scans)?

\item 
[ ] Any string concatenation with \code{+} inside a loop instead of \code{StringBuilder}?

\item 
[ ] Any object allocation, regex compilation, or reflection happening inside a hot loop that could be hoisted out?

\item 
[ ] Is autoboxing happening in a tight numeric loop (\code{Integer} instead of \code{int}, boxed streams)?

\item 
[ ] Any premature \code{synchronized}/parallelism added without evidence it’s needed, at the cost of readability?

\end{itemize}

\subsection[Style (only after everything above is clean)]{Style (only after everything above is clean)}
\label{h:22:style-only-after-everything-above-is-clean}
\begin{itemize}
\item 
[ ] Naming clear and consistent with the domain?

\item 
[ ] Method length / single responsibility reasonable?

\item 
[ ] Are modern idioms used well (records, \code{var}, streams, switch expressions) without overuse — see \hyperref[h:22:modern-java-feature-usage]{Modern Java Feature Usage}?

\end{itemize}

\section[How to Communicate in a Code-Review Interview]{How to Communicate in a Code-Review Interview}
\label{h:22:how-to-communicate-in-a-code-review-interview}
Finding the bug is half the exercise. How you talk about it is the other half — interviewers are evaluating whether you’d be a good teammate in an actual review, not just a bug detector.

\begin{itemize}
\item 
\textbf{Lead with the highest-severity issue.} Don’t narrate top-to-bottom in file order; scan the whole snippet first, rank what you found, and open with the worst one. “Before anything else — this singleton isn’t thread-safe” beats starting with a variable-naming nit.

\item 
\textbf{Say why, not just what.} “This will break” is a claim; “this will break \emph{because} two threads can both pass the null-check before either assigns the instance, so you can get two singletons” is a review. Explaining the mechanism is what proves you understand it rather than pattern-matching a lint rule.

\item 
\textbf{Propose a concrete fix.} Don’t stop at “this is unsafe.” Say what you’d change it to — an initialization-on-demand holder, \code{ConcurrentHashMap}, a try-with-resources block — ideally in a sentence or two of code, not just an adjective.

\item 
\textbf{Separate blocking issues from nits.} Explicitly label things: “this one’s a blocker, it will corrupt data in production” vs. “this is a nit, purely a naming preference — not worth blocking on.” Interviewers want to see you can triage, not just that you can spot fifteen things.

\item 
\textbf{Ask questions instead of accusing.} “Was this meant to run in a single-threaded context, or could multiple requests hit it concurrently?” reads as collaborative. “You clearly didn’t think about thread safety” reads as hostile — and might be wrong if there’s context you don’t have (maybe it really is confined to one thread by contract).

\item 
\textbf{Say what’s good, too.} If the error handling is solid or the naming is clear, say so briefly. It shows judgment (you’re not fault-finding for its own sake) and it’s simply what a real review looks like — real PRs aren’t 100\% criticism.

\end{itemize}

A reasonable spoken structure: \emph{“The biggest problem is X, because Y — I’d fix it by Z. There’s also a smaller issue with A. Everything else — B, C — are just nits, not blockers. One thing I like here is D.”}

\section[Common Java Interview Pitfalls]{Common Java Interview Pitfalls}
\label{h:22:common-java-interview-pitfalls}
These are the recurring, easy-to-miss patterns interviewers plant because they show up constantly in real codebases. Recognize the \emph{shape} of each one instantly — you shouldn’t need to trace through logic to spot them.

\begin{enumerate}
\item 
\textbf{\code{==} on boxed types or \code{String}.} \code{Integer~\allowbreak{}a~\allowbreak{}=\allowbreak{}~\allowbreak{}1000,\allowbreak{}~\allowbreak{}b~\allowbreak{}=\allowbreak{}~\allowbreak{}1000;\allowbreak{}~\allowbreak{}a~\allowbreak{}==\allowbreak{}~\allowbreak{}b} is \code{false} outside the \code{-\allowbreak{}128..\allowbreak{}127} cache range. Always \code{.\allowbreak{}equals()} for objects.

\item 
\textbf{Autoboxing NPE.} \code{Integer~\allowbreak{}count~\allowbreak{}=\allowbreak{}~\allowbreak{}map.\allowbreak{}get(\allowbreak{}key);\allowbreak{}~\allowbreak{}if~\allowbreak{}(\allowbreak{}count~\allowbreak{}\textgreater{}\allowbreak{}~\allowbreak{}0)} throws \code{NullPointerException} when \code{count} is \code{null} — the unboxing happens invisibly at \code{\textgreater{}}.

\item 
\textbf{Mutable \code{static} fields shared across requests/threads} — the classic \code{SimpleDateFormat}, lazily-initialized singleton, or cache-as-\code{HashMap} pattern (see Scenarios below).

\item 
\textbf{\code{equals()} without \code{hashCode()}}, or vice versa — breaks every hash-based collection silently (no compiler error, just wrong behavior).

\item 
\textbf{Using a mutable object as a \code{Map}/\code{Set} key}, then mutating it after insertion — the entry becomes unreachable by lookup.

\item 
\textbf{String concatenation with \code{+} in a loop} — quadratic behavior because each \code{+} builds a new \code{String}.

\item 
\textbf{Catching \code{Exception} (or \code{Throwable}) far wider than needed}, and/or doing nothing with it — see \hyperref[chap:6]{Chapter~6, \emph{Exception Handling}} for the full treatment.

\item 
\textbf{\code{finally} with a \code{return}/\code{throw}} that silently discards the exception that was in flight.

\item 
\textbf{Off-by-one and empty-collection bugs} — \code{for~\allowbreak{}(\allowbreak{}int~\allowbreak{}i~\allowbreak{}=\allowbreak{}~\allowbreak{}0;\allowbreak{}~\allowbreak{}i~\allowbreak{}\textless{}=\allowbreak{}~\allowbreak{}list.\allowbreak{}size();\allowbreak{}~\allowbreak{}i++)}, or forgetting a collection can be empty before calling \code{.\allowbreak{}get(\allowbreak{}0)}.

\item 
\textbf{Integer overflow from \code{int} arithmetic} — \code{int~\allowbreak{}total~\allowbreak{}=\allowbreak{}~\allowbreak{}a~\allowbreak{}*\allowbreak{}~\allowbreak{}b;} overflows silently where \code{long} wouldn’t; no exception, just a wrong answer.

\item 
\textbf{Comparing floating point with \code{==}} instead of an epsilon, or using \code{double}/\code{float} for money instead of \code{BigDecimal}.

\item 
\textbf{Modifying a collection while iterating it} with a plain \code{for-\allowbreak{}each}, throwing \code{Concurrent\allowbreak{}Modification\allowbreak{}Exception} — should use an \code{Iterator.\allowbreak{}remove()}, \code{removeIf}, or a copy.

\item 
\textbf{Ignoring the return value of an immutable operation} — \code{str.\allowbreak{}trim();} on its own line does nothing, because \code{String} is immutable and \code{trim()} returns a new value.

\item 
\textbf{Leaking \code{this} from a constructor} — starting a thread, registering a listener, or passing \code{this} to another object before the constructor finishes (see Scenario 15).

\item 
\textbf{Unbounded recursion or unbounded collection growth} with no base case / no eviction, eventually a \code{StackOverflowError} or \code{OutOfMemoryError} under load rather than at compile time.

\item 
\textbf{Checked exception swallowed and converted to a generic, cause-less \code{RuntimeException}}, losing the original stack trace.

\item 
\textbf{Using \code{Optional} as a field type, method parameter, or inside a collection} instead of only as a return type — it’s designed for one purpose, and misusing it is itself a common interview trap (see \hyperref[h:22:api-design-review]{API Design Review}).

\item 
\textbf{Relying on \code{HashMap}/\code{HashSet} iteration order}, which is unspecified and can silently change between JDK versions or after resizing.

\end{enumerate}

These patterns recur throughout the scenarios below — as you read each one, try to name which pitfall(s) from this list it is before reading the hidden answer.

\section[Code Review Scenarios]{Code Review Scenarios}
\label{h:22:code-review-scenarios}
Each scenario is realistic production-style code with a real, plantable bug (sometimes more than one). Read the snippet, form your own opinion, rank the issues by severity, then expand the review to check yourself.

\subsection[Scenario 1: The Configuration Singleton]{Scenario 1: The Configuration Singleton}
\label{h:22:scenario-1-the-configuration-singleton}
\begin{codebox}
\begin{Verbatim}[commandchars=\\\{\}]
\PY{k+kd}{public}\PY{+w}{ }\PY{k+kd}{class} \PY{n+nc}{ConfigManager}\PY{+w}{ }\PY{p}{\PYZob{}}
\PY{+w}{    }\PY{k+kd}{private}\PY{+w}{ }\PY{k+kd}{static}\PY{+w}{ }\PY{n}{ConfigManager}\PY{+w}{ }\PY{n}{instance}\PY{p}{;}
\PY{+w}{    }\PY{k+kd}{private}\PY{+w}{ }\PY{k+kd}{final}\PY{+w}{ }\PY{n}{Map}\PY{o}{\PYZlt{}}\PY{n}{String}\PY{p}{,}\PY{+w}{ }\PY{n}{String}\PY{o}{\PYZgt{}}\PY{+w}{ }\PY{n}{settings}\PY{+w}{ }\PY{o}{=}\PY{+w}{ }\PY{k}{new}\PY{+w}{ }\PY{n}{HashMap}\PY{o}{\PYZlt{}}\PY{o}{\PYZgt{}}\PY{p}{(}\PY{p}{)}\PY{p}{;}

\PY{+w}{    }\PY{k+kd}{private}\PY{+w}{ }\PY{n+nf}{ConfigManager}\PY{p}{(}\PY{p}{)}\PY{+w}{ }\PY{p}{\PYZob{}}
\PY{+w}{        }\PY{n}{settings}\PY{p}{.}\PY{n+na}{put}\PY{p}{(}\PY{l+s}{\PYZdq{}}\PY{l+s}{timeout}\PY{l+s}{\PYZdq{}}\PY{p}{,}\PY{+w}{ }\PY{l+s}{\PYZdq{}}\PY{l+s}{30}\PY{l+s}{\PYZdq{}}\PY{p}{)}\PY{p}{;}
\PY{+w}{        }\PY{n}{settings}\PY{p}{.}\PY{n+na}{put}\PY{p}{(}\PY{l+s}{\PYZdq{}}\PY{l+s}{retries}\PY{l+s}{\PYZdq{}}\PY{p}{,}\PY{+w}{ }\PY{l+s}{\PYZdq{}}\PY{l+s}{3}\PY{l+s}{\PYZdq{}}\PY{p}{)}\PY{p}{;}
\PY{+w}{    }\PY{p}{\PYZcb{}}

\PY{+w}{    }\PY{k+kd}{public}\PY{+w}{ }\PY{k+kd}{static}\PY{+w}{ }\PY{n}{ConfigManager}\PY{+w}{ }\PY{n+nf}{getInstance}\PY{p}{(}\PY{p}{)}\PY{+w}{ }\PY{p}{\PYZob{}}
\PY{+w}{        }\PY{k}{if}\PY{+w}{ }\PY{p}{(}\PY{n}{instance}\PY{+w}{ }\PY{o}{=}\PY{o}{=}\PY{+w}{ }\PY{k+kc}{null}\PY{p}{)}\PY{+w}{ }\PY{p}{\PYZob{}}
\PY{+w}{            }\PY{n}{instance}\PY{+w}{ }\PY{o}{=}\PY{+w}{ }\PY{k}{new}\PY{+w}{ }\PY{n}{ConfigManager}\PY{p}{(}\PY{p}{)}\PY{p}{;}
\PY{+w}{        }\PY{p}{\PYZcb{}}
\PY{+w}{        }\PY{k}{return}\PY{+w}{ }\PY{n}{instance}\PY{p}{;}
\PY{+w}{    }\PY{p}{\PYZcb{}}

\PY{+w}{    }\PY{k+kd}{public}\PY{+w}{ }\PY{n}{String}\PY{+w}{ }\PY{n+nf}{get}\PY{p}{(}\PY{n}{String}\PY{+w}{ }\PY{n}{key}\PY{p}{)}\PY{+w}{ }\PY{p}{\PYZob{}}
\PY{+w}{        }\PY{k}{return}\PY{+w}{ }\PY{n}{settings}\PY{p}{.}\PY{n+na}{get}\PY{p}{(}\PY{n}{key}\PY{p}{)}\PY{p}{;}
\PY{+w}{    }\PY{p}{\PYZcb{}}

\PY{+w}{    }\PY{k+kd}{public}\PY{+w}{ }\PY{k+kt}{void}\PY{+w}{ }\PY{n+nf}{set}\PY{p}{(}\PY{n}{String}\PY{+w}{ }\PY{n}{key}\PY{p}{,}\PY{+w}{ }\PY{n}{String}\PY{+w}{ }\PY{n}{value}\PY{p}{)}\PY{+w}{ }\PY{p}{\PYZob{}}
\PY{+w}{        }\PY{n}{settings}\PY{p}{.}\PY{n+na}{put}\PY{p}{(}\PY{n}{key}\PY{p}{,}\PY{+w}{ }\PY{n}{value}\PY{p}{)}\PY{p}{;}
\PY{+w}{    }\PY{p}{\PYZcb{}}
\PY{p}{\PYZcb{}}
\end{Verbatim}
\end{codebox}

\begin{sidebar}{Review}
\begin{itemize}
\item 
\textbf{Blocker — non-thread-safe lazy singleton.} \code{getInstance()} has a classic race: two threads can both see \code{instance~\allowbreak{}==\allowbreak{}~\allowbreak{}null}, and both construct a \code{ConfigManager}, handing out two different instances to different callers. This is the textbook double-checked-locking setup, minus the locking. Under low concurrency it “works” in testing and then fails intermittently in production — the worst kind of bug to debug.

\item 
\textbf{Major — \code{settings} map is a mutable \code{HashMap} with no synchronization}, and \code{set()}/\code{get()} can be called concurrently once multiple threads hold references (even to the same instance). Concurrent writes to a plain \code{HashMap} can corrupt its internal structure, not just race on values.

\item 
\textbf{Minor — no way to remove/reset settings} for tests, which tends to leak state between test cases sharing the singleton. Not a functional bug, but worth flagging as a design smell for testability.

\end{itemize}

Fixed code — initialization-on-demand holder for lazy, thread-safe singleton construction, plus a concurrent-safe backing map:

\begin{codebox}
\begin{Verbatim}[commandchars=\\\{\}]
\PY{k+kd}{public}\PY{+w}{ }\PY{k+kd}{class} \PY{n+nc}{ConfigManager}\PY{+w}{ }\PY{p}{\PYZob{}}
\PY{+w}{    }\PY{k+kd}{private}\PY{+w}{ }\PY{k+kd}{final}\PY{+w}{ }\PY{n}{Map}\PY{o}{\PYZlt{}}\PY{n}{String}\PY{p}{,}\PY{+w}{ }\PY{n}{String}\PY{o}{\PYZgt{}}\PY{+w}{ }\PY{n}{settings}\PY{+w}{ }\PY{o}{=}\PY{+w}{ }\PY{k}{new}\PY{+w}{ }\PY{n}{ConcurrentHashMap}\PY{o}{\PYZlt{}}\PY{o}{\PYZgt{}}\PY{p}{(}\PY{p}{)}\PY{p}{;}

\PY{+w}{    }\PY{k+kd}{private}\PY{+w}{ }\PY{n+nf}{ConfigManager}\PY{p}{(}\PY{p}{)}\PY{+w}{ }\PY{p}{\PYZob{}}
\PY{+w}{        }\PY{n}{settings}\PY{p}{.}\PY{n+na}{put}\PY{p}{(}\PY{l+s}{\PYZdq{}}\PY{l+s}{timeout}\PY{l+s}{\PYZdq{}}\PY{p}{,}\PY{+w}{ }\PY{l+s}{\PYZdq{}}\PY{l+s}{30}\PY{l+s}{\PYZdq{}}\PY{p}{)}\PY{p}{;}
\PY{+w}{        }\PY{n}{settings}\PY{p}{.}\PY{n+na}{put}\PY{p}{(}\PY{l+s}{\PYZdq{}}\PY{l+s}{retries}\PY{l+s}{\PYZdq{}}\PY{p}{,}\PY{+w}{ }\PY{l+s}{\PYZdq{}}\PY{l+s}{3}\PY{l+s}{\PYZdq{}}\PY{p}{)}\PY{p}{;}
\PY{+w}{    }\PY{p}{\PYZcb{}}

\PY{+w}{    }\PY{k+kd}{private}\PY{+w}{ }\PY{k+kd}{static}\PY{+w}{ }\PY{k+kd}{final}\PY{+w}{ }\PY{k+kd}{class} \PY{n+nc}{Holder}\PY{+w}{ }\PY{p}{\PYZob{}}
\PY{+w}{        }\PY{k+kd}{static}\PY{+w}{ }\PY{k+kd}{final}\PY{+w}{ }\PY{n}{ConfigManager}\PY{+w}{ }\PY{n}{INSTANCE}\PY{+w}{ }\PY{o}{=}\PY{+w}{ }\PY{k}{new}\PY{+w}{ }\PY{n}{ConfigManager}\PY{p}{(}\PY{p}{)}\PY{p}{;}
\PY{+w}{    }\PY{p}{\PYZcb{}}

\PY{+w}{    }\PY{k+kd}{public}\PY{+w}{ }\PY{k+kd}{static}\PY{+w}{ }\PY{n}{ConfigManager}\PY{+w}{ }\PY{n+nf}{getInstance}\PY{p}{(}\PY{p}{)}\PY{+w}{ }\PY{p}{\PYZob{}}
\PY{+w}{        }\PY{k}{return}\PY{+w}{ }\PY{n}{Holder}\PY{p}{.}\PY{n+na}{INSTANCE}\PY{p}{;}\PY{+w}{ }\PY{c+c1}{// JVM class\PYZhy{}loading guarantees are the \PYZdq{}locking\PYZdq{}}
\PY{+w}{    }\PY{p}{\PYZcb{}}

\PY{+w}{    }\PY{k+kd}{public}\PY{+w}{ }\PY{n}{String}\PY{+w}{ }\PY{n+nf}{get}\PY{p}{(}\PY{n}{String}\PY{+w}{ }\PY{n}{key}\PY{p}{)}\PY{+w}{ }\PY{p}{\PYZob{}}
\PY{+w}{        }\PY{k}{return}\PY{+w}{ }\PY{n}{settings}\PY{p}{.}\PY{n+na}{get}\PY{p}{(}\PY{n}{key}\PY{p}{)}\PY{p}{;}
\PY{+w}{    }\PY{p}{\PYZcb{}}

\PY{+w}{    }\PY{k+kd}{public}\PY{+w}{ }\PY{k+kt}{void}\PY{+w}{ }\PY{n+nf}{set}\PY{p}{(}\PY{n}{String}\PY{+w}{ }\PY{n}{key}\PY{p}{,}\PY{+w}{ }\PY{n}{String}\PY{+w}{ }\PY{n}{value}\PY{p}{)}\PY{+w}{ }\PY{p}{\PYZob{}}
\PY{+w}{        }\PY{n}{settings}\PY{p}{.}\PY{n+na}{put}\PY{p}{(}\PY{n}{key}\PY{p}{,}\PY{+w}{ }\PY{n}{value}\PY{p}{)}\PY{p}{;}
\PY{+w}{    }\PY{p}{\PYZcb{}}
\PY{p}{\PYZcb{}}
\end{Verbatim}
\end{codebox}

The holder idiom is thread-safe without explicit locking because the JVM guarantees a class is initialized exactly once, under a lock, the first time it’s actively used — see \hyperref[chap:13]{Chapter~13, \emph{Concurrency (Core)}} for the Java Memory Model guarantees behind this.

\end{sidebar}
\subsection[Scenario 2: The Request Counter]{Scenario 2: The Request Counter}
\label{h:22:scenario-2-the-request-counter}
\begin{codebox}
\begin{Verbatim}[commandchars=\\\{\}]
\PY{k+kd}{public}\PY{+w}{ }\PY{k+kd}{class} \PY{n+nc}{RequestCounterService}\PY{+w}{ }\PY{p}{\PYZob{}}
\PY{+w}{    }\PY{k+kd}{private}\PY{+w}{ }\PY{k+kd}{final}\PY{+w}{ }\PY{n}{Map}\PY{o}{\PYZlt{}}\PY{n}{String}\PY{p}{,}\PY{+w}{ }\PY{n}{Integer}\PY{o}{\PYZgt{}}\PY{+w}{ }\PY{n}{countsByEndpoint}\PY{+w}{ }\PY{o}{=}\PY{+w}{ }\PY{k}{new}\PY{+w}{ }\PY{n}{HashMap}\PY{o}{\PYZlt{}}\PY{o}{\PYZgt{}}\PY{p}{(}\PY{p}{)}\PY{p}{;}

\PY{+w}{    }\PY{c+c1}{// Called concurrently from many servlet request threads}
\PY{+w}{    }\PY{k+kd}{public}\PY{+w}{ }\PY{k+kt}{void}\PY{+w}{ }\PY{n+nf}{recordRequest}\PY{p}{(}\PY{n}{String}\PY{+w}{ }\PY{n}{endpoint}\PY{p}{)}\PY{+w}{ }\PY{p}{\PYZob{}}
\PY{+w}{        }\PY{n}{Integer}\PY{+w}{ }\PY{n}{current}\PY{+w}{ }\PY{o}{=}\PY{+w}{ }\PY{n}{countsByEndpoint}\PY{p}{.}\PY{n+na}{get}\PY{p}{(}\PY{n}{endpoint}\PY{p}{)}\PY{p}{;}
\PY{+w}{        }\PY{k}{if}\PY{+w}{ }\PY{p}{(}\PY{n}{current}\PY{+w}{ }\PY{o}{=}\PY{o}{=}\PY{+w}{ }\PY{k+kc}{null}\PY{p}{)}\PY{+w}{ }\PY{p}{\PYZob{}}
\PY{+w}{            }\PY{n}{countsByEndpoint}\PY{p}{.}\PY{n+na}{put}\PY{p}{(}\PY{n}{endpoint}\PY{p}{,}\PY{+w}{ }\PY{l+m+mi}{1}\PY{p}{)}\PY{p}{;}
\PY{+w}{        }\PY{p}{\PYZcb{}}\PY{+w}{ }\PY{k}{else}\PY{+w}{ }\PY{p}{\PYZob{}}
\PY{+w}{            }\PY{n}{countsByEndpoint}\PY{p}{.}\PY{n+na}{put}\PY{p}{(}\PY{n}{endpoint}\PY{p}{,}\PY{+w}{ }\PY{n}{current}\PY{+w}{ }\PY{o}{+}\PY{+w}{ }\PY{l+m+mi}{1}\PY{p}{)}\PY{p}{;}
\PY{+w}{        }\PY{p}{\PYZcb{}}
\PY{+w}{    }\PY{p}{\PYZcb{}}

\PY{+w}{    }\PY{k+kd}{public}\PY{+w}{ }\PY{k+kt}{int}\PY{+w}{ }\PY{n+nf}{getCount}\PY{p}{(}\PY{n}{String}\PY{+w}{ }\PY{n}{endpoint}\PY{p}{)}\PY{+w}{ }\PY{p}{\PYZob{}}
\PY{+w}{        }\PY{k}{return}\PY{+w}{ }\PY{n}{countsByEndpoint}\PY{p}{.}\PY{n+na}{getOrDefault}\PY{p}{(}\PY{n}{endpoint}\PY{p}{,}\PY{+w}{ }\PY{l+m+mi}{0}\PY{p}{)}\PY{p}{;}
\PY{+w}{    }\PY{p}{\PYZcb{}}

\PY{+w}{    }\PY{k+kd}{public}\PY{+w}{ }\PY{n}{Map}\PY{o}{\PYZlt{}}\PY{n}{String}\PY{p}{,}\PY{+w}{ }\PY{n}{Integer}\PY{o}{\PYZgt{}}\PY{+w}{ }\PY{n+nf}{snapshot}\PY{p}{(}\PY{p}{)}\PY{+w}{ }\PY{p}{\PYZob{}}
\PY{+w}{        }\PY{k}{return}\PY{+w}{ }\PY{n}{countsByEndpoint}\PY{p}{;}
\PY{+w}{    }\PY{p}{\PYZcb{}}
\PY{p}{\PYZcb{}}
\end{Verbatim}
\end{codebox}

\begin{sidebar}{Review}
\begin{itemize}
\item 
\textbf{Blocker — plain \code{HashMap} mutated from multiple threads.} The comment even says “called concurrently,” so this isn’t a hypothetical: concurrent \code{put()} calls on a \code{HashMap} can corrupt internal bucket structure (infinite loops during resize in older JDKs, or silently dropped entries), not just lose a counter increment.

\item 
\textbf{Major — read-then-write is not atomic} even with a concurrent map: \code{get} then \code{put(\allowbreak{}current~\allowbreak{}+\allowbreak{}~\allowbreak{}1)} is a classic lost-update race — two threads can both read the same \code{current} and both write \code{current~\allowbreak{}+\allowbreak{}~\allowbreak{}1}, losing one increment. Switching the map type alone does not fix this.

\item 
\textbf{Minor — \code{snapshot()} returns the live internal map}, not a copy, so callers can mutate internal state, and any iteration over it concurrently with a write is still unsafe even with \code{ConcurrentHashMap} (iteration is weakly consistent, which is fine for a snapshot but callers should not assume they can safely mutate what they get back).

\end{itemize}

Fixed code — \code{ConcurrentHashMap} plus an atomic update method (\code{merge}, or an \code{AtomicInteger} per key) instead of get-then-put, and a defensive copy for the snapshot:

\begin{codebox}
\begin{Verbatim}[commandchars=\\\{\}]
\PY{k+kd}{public}\PY{+w}{ }\PY{k+kd}{class} \PY{n+nc}{RequestCounterService}\PY{+w}{ }\PY{p}{\PYZob{}}
\PY{+w}{    }\PY{k+kd}{private}\PY{+w}{ }\PY{k+kd}{final}\PY{+w}{ }\PY{n}{Map}\PY{o}{\PYZlt{}}\PY{n}{String}\PY{p}{,}\PY{+w}{ }\PY{n}{AtomicInteger}\PY{o}{\PYZgt{}}\PY{+w}{ }\PY{n}{countsByEndpoint}\PY{+w}{ }\PY{o}{=}\PY{+w}{ }\PY{k}{new}\PY{+w}{ }\PY{n}{ConcurrentHashMap}\PY{o}{\PYZlt{}}\PY{o}{\PYZgt{}}\PY{p}{(}\PY{p}{)}\PY{p}{;}

\PY{+w}{    }\PY{k+kd}{public}\PY{+w}{ }\PY{k+kt}{void}\PY{+w}{ }\PY{n+nf}{recordRequest}\PY{p}{(}\PY{n}{String}\PY{+w}{ }\PY{n}{endpoint}\PY{p}{)}\PY{+w}{ }\PY{p}{\PYZob{}}
\PY{+w}{        }\PY{n}{countsByEndpoint}\PY{p}{.}\PY{n+na}{computeIfAbsent}\PY{p}{(}\PY{n}{endpoint}\PY{p}{,}\PY{+w}{ }\PY{n}{e}\PY{+w}{ }\PY{o}{\PYZhy{}}\PY{o}{\PYZgt{}}\PY{+w}{ }\PY{k}{new}\PY{+w}{ }\PY{n}{AtomicInteger}\PY{p}{(}\PY{p}{)}\PY{p}{)}
\PY{+w}{                        }\PY{p}{.}\PY{n+na}{incrementAndGet}\PY{p}{(}\PY{p}{)}\PY{p}{;}
\PY{+w}{    }\PY{p}{\PYZcb{}}

\PY{+w}{    }\PY{k+kd}{public}\PY{+w}{ }\PY{k+kt}{int}\PY{+w}{ }\PY{n+nf}{getCount}\PY{p}{(}\PY{n}{String}\PY{+w}{ }\PY{n}{endpoint}\PY{p}{)}\PY{+w}{ }\PY{p}{\PYZob{}}
\PY{+w}{        }\PY{n}{AtomicInteger}\PY{+w}{ }\PY{n}{count}\PY{+w}{ }\PY{o}{=}\PY{+w}{ }\PY{n}{countsByEndpoint}\PY{p}{.}\PY{n+na}{get}\PY{p}{(}\PY{n}{endpoint}\PY{p}{)}\PY{p}{;}
\PY{+w}{        }\PY{k}{return}\PY{+w}{ }\PY{n}{count}\PY{+w}{ }\PY{o}{=}\PY{o}{=}\PY{+w}{ }\PY{k+kc}{null}\PY{+w}{ }\PY{o}{?}\PY{+w}{ }\PY{l+m+mi}{0}\PY{+w}{ }\PY{p}{:}\PY{+w}{ }\PY{n}{count}\PY{p}{.}\PY{n+na}{get}\PY{p}{(}\PY{p}{)}\PY{p}{;}
\PY{+w}{    }\PY{p}{\PYZcb{}}

\PY{+w}{    }\PY{k+kd}{public}\PY{+w}{ }\PY{n}{Map}\PY{o}{\PYZlt{}}\PY{n}{String}\PY{p}{,}\PY{+w}{ }\PY{n}{Integer}\PY{o}{\PYZgt{}}\PY{+w}{ }\PY{n+nf}{snapshot}\PY{p}{(}\PY{p}{)}\PY{+w}{ }\PY{p}{\PYZob{}}
\PY{+w}{        }\PY{n}{Map}\PY{o}{\PYZlt{}}\PY{n}{String}\PY{p}{,}\PY{+w}{ }\PY{n}{Integer}\PY{o}{\PYZgt{}}\PY{+w}{ }\PY{n}{copy}\PY{+w}{ }\PY{o}{=}\PY{+w}{ }\PY{k}{new}\PY{+w}{ }\PY{n}{HashMap}\PY{o}{\PYZlt{}}\PY{o}{\PYZgt{}}\PY{p}{(}\PY{p}{)}\PY{p}{;}
\PY{+w}{        }\PY{n}{countsByEndpoint}\PY{p}{.}\PY{n+na}{forEach}\PY{p}{(}\PY{p}{(}\PY{n}{k}\PY{p}{,}\PY{+w}{ }\PY{n}{v}\PY{p}{)}\PY{+w}{ }\PY{o}{\PYZhy{}}\PY{o}{\PYZgt{}}\PY{+w}{ }\PY{n}{copy}\PY{p}{.}\PY{n+na}{put}\PY{p}{(}\PY{n}{k}\PY{p}{,}\PY{+w}{ }\PY{n}{v}\PY{p}{.}\PY{n+na}{get}\PY{p}{(}\PY{p}{)}\PY{p}{)}\PY{p}{)}\PY{p}{;}
\PY{+w}{        }\PY{k}{return}\PY{+w}{ }\PY{n}{copy}\PY{p}{;}
\PY{+w}{    }\PY{p}{\PYZcb{}}
\PY{p}{\PYZcb{}}
\end{Verbatim}
\end{codebox}

\end{sidebar}
\subsection[Scenario 3: The Order Deduplicator]{Scenario 3: The Order Deduplicator}
\label{h:22:scenario-3-the-order-deduplicator}
\begin{codebox}
\begin{Verbatim}[commandchars=\\\{\}]
\PY{k+kd}{public}\PY{+w}{ }\PY{k+kd}{class} \PY{n+nc}{OrderKey}\PY{+w}{ }\PY{p}{\PYZob{}}
\PY{+w}{    }\PY{k+kd}{private}\PY{+w}{ }\PY{n}{String}\PY{+w}{ }\PY{n}{customerId}\PY{p}{;}
\PY{+w}{    }\PY{k+kd}{private}\PY{+w}{ }\PY{n}{LocalDate}\PY{+w}{ }\PY{n}{orderDate}\PY{p}{;}

\PY{+w}{    }\PY{k+kd}{public}\PY{+w}{ }\PY{n+nf}{OrderKey}\PY{p}{(}\PY{n}{String}\PY{+w}{ }\PY{n}{customerId}\PY{p}{,}\PY{+w}{ }\PY{n}{LocalDate}\PY{+w}{ }\PY{n}{orderDate}\PY{p}{)}\PY{+w}{ }\PY{p}{\PYZob{}}
\PY{+w}{        }\PY{k}{this}\PY{p}{.}\PY{n+na}{customerId}\PY{+w}{ }\PY{o}{=}\PY{+w}{ }\PY{n}{customerId}\PY{p}{;}
\PY{+w}{        }\PY{k}{this}\PY{p}{.}\PY{n+na}{orderDate}\PY{+w}{ }\PY{o}{=}\PY{+w}{ }\PY{n}{orderDate}\PY{p}{;}
\PY{+w}{    }\PY{p}{\PYZcb{}}

\PY{+w}{    }\PY{k+kd}{public}\PY{+w}{ }\PY{k+kt}{void}\PY{+w}{ }\PY{n+nf}{setOrderDate}\PY{p}{(}\PY{n}{LocalDate}\PY{+w}{ }\PY{n}{orderDate}\PY{p}{)}\PY{+w}{ }\PY{p}{\PYZob{}}
\PY{+w}{        }\PY{k}{this}\PY{p}{.}\PY{n+na}{orderDate}\PY{+w}{ }\PY{o}{=}\PY{+w}{ }\PY{n}{orderDate}\PY{p}{;}
\PY{+w}{    }\PY{p}{\PYZcb{}}

\PY{+w}{    }\PY{n+nd}{@Override}
\PY{+w}{    }\PY{k+kd}{public}\PY{+w}{ }\PY{k+kt}{boolean}\PY{+w}{ }\PY{n+nf}{equals}\PY{p}{(}\PY{n}{Object}\PY{+w}{ }\PY{n}{o}\PY{p}{)}\PY{+w}{ }\PY{p}{\PYZob{}}
\PY{+w}{        }\PY{k}{if}\PY{+w}{ }\PY{p}{(}\PY{o}{!}\PY{p}{(}\PY{n}{o}\PY{+w}{ }\PY{k}{instanceof}\PY{+w}{ }\PY{n}{OrderKey}\PY{+w}{ }\PY{n}{other}\PY{p}{)}\PY{p}{)}\PY{+w}{ }\PY{k}{return}\PY{+w}{ }\PY{k+kc}{false}\PY{p}{;}
\PY{+w}{        }\PY{k}{return}\PY{+w}{ }\PY{n}{customerId}\PY{p}{.}\PY{n+na}{equals}\PY{p}{(}\PY{n}{other}\PY{p}{.}\PY{n+na}{customerId}\PY{p}{)}\PY{+w}{ }\PY{o}{\PYZam{}}\PY{o}{\PYZam{}}\PY{+w}{ }\PY{n}{orderDate}\PY{p}{.}\PY{n+na}{equals}\PY{p}{(}\PY{n}{other}\PY{p}{.}\PY{n+na}{orderDate}\PY{p}{)}\PY{p}{;}
\PY{+w}{    }\PY{p}{\PYZcb{}}

\PY{+w}{    }\PY{n+nd}{@Override}
\PY{+w}{    }\PY{k+kd}{public}\PY{+w}{ }\PY{k+kt}{int}\PY{+w}{ }\PY{n+nf}{hashCode}\PY{p}{(}\PY{p}{)}\PY{+w}{ }\PY{p}{\PYZob{}}
\PY{+w}{        }\PY{k}{return}\PY{+w}{ }\PY{n}{Objects}\PY{p}{.}\PY{n+na}{hash}\PY{p}{(}\PY{n}{customerId}\PY{p}{,}\PY{+w}{ }\PY{n}{orderDate}\PY{p}{)}\PY{p}{;}
\PY{+w}{    }\PY{p}{\PYZcb{}}
\PY{p}{\PYZcb{}}

\PY{k+kd}{public}\PY{+w}{ }\PY{k+kd}{class} \PY{n+nc}{OrderDeduplicator}\PY{+w}{ }\PY{p}{\PYZob{}}
\PY{+w}{    }\PY{k+kd}{private}\PY{+w}{ }\PY{k+kd}{final}\PY{+w}{ }\PY{n}{Map}\PY{o}{\PYZlt{}}\PY{n}{OrderKey}\PY{p}{,}\PY{+w}{ }\PY{n}{Order}\PY{o}{\PYZgt{}}\PY{+w}{ }\PY{n}{seen}\PY{+w}{ }\PY{o}{=}\PY{+w}{ }\PY{k}{new}\PY{+w}{ }\PY{n}{HashMap}\PY{o}{\PYZlt{}}\PY{o}{\PYZgt{}}\PY{p}{(}\PY{p}{)}\PY{p}{;}

\PY{+w}{    }\PY{k+kd}{public}\PY{+w}{ }\PY{k+kt}{boolean}\PY{+w}{ }\PY{n+nf}{isDuplicate}\PY{p}{(}\PY{n}{OrderKey}\PY{+w}{ }\PY{n}{key}\PY{p}{,}\PY{+w}{ }\PY{n}{Order}\PY{+w}{ }\PY{n}{order}\PY{p}{,}\PY{+w}{ }\PY{n}{LocalDate}\PY{+w}{ }\PY{n}{correctedDate}\PY{p}{)}\PY{+w}{ }\PY{p}{\PYZob{}}
\PY{+w}{        }\PY{n}{key}\PY{p}{.}\PY{n+na}{setOrderDate}\PY{p}{(}\PY{n}{correctedDate}\PY{p}{)}\PY{p}{;}\PY{+w}{ }\PY{c+c1}{// apply a late correction before checking}
\PY{+w}{        }\PY{k+kt}{boolean}\PY{+w}{ }\PY{n}{duplicate}\PY{+w}{ }\PY{o}{=}\PY{+w}{ }\PY{n}{seen}\PY{p}{.}\PY{n+na}{containsKey}\PY{p}{(}\PY{n}{key}\PY{p}{)}\PY{p}{;}
\PY{+w}{        }\PY{n}{seen}\PY{p}{.}\PY{n+na}{put}\PY{p}{(}\PY{n}{key}\PY{p}{,}\PY{+w}{ }\PY{n}{order}\PY{p}{)}\PY{p}{;}
\PY{+w}{        }\PY{k}{return}\PY{+w}{ }\PY{n}{duplicate}\PY{p}{;}
\PY{+w}{    }\PY{p}{\PYZcb{}}
\PY{p}{\PYZcb{}}
\end{Verbatim}
\end{codebox}

\begin{sidebar}{Review}
\begin{itemize}
\item 
\textbf{Blocker — mutable object used as a hash-based map key.} \code{OrderKey} has a setter (\code{setOrderDate}) and participates in \code{hashCode()}. Mutating a key after it’s already stored in a \code{HashMap} changes its hash code, so the entry ends up in the wrong bucket — it becomes permanently unreachable via \code{get}/\code{containsKey} even though it’s still sitting in the map, silently leaking memory \emph{and} breaking the dedup logic it exists to implement.

\item 
\textbf{Major — the mutation happens right before the lookup}, so even a \emph{new} key that’s mutated before first insertion could match against a stale bucket layout if the same key object is reused/cached elsewhere. Any key type used in a hash-based collection must be effectively immutable for its whole lifetime.

\item 
\textbf{Minor — \code{OrderKey} fields aren’t \code{final}}, which invites exactly this kind of mistake in the future even if today’s only caller happens to use it “safely” for a single mutation before insertion.

\end{itemize}

Fixed code — make \code{OrderKey} immutable; if a corrected date is needed, build a new key instead of mutating the old one:

\begin{codebox}
\begin{Verbatim}[commandchars=\\\{\}]
\PY{k+kd}{public}\PY{+w}{ }\PY{k+kd}{final}\PY{+w}{ }\PY{k+kd}{class} \PY{n+nc}{OrderKey}\PY{+w}{ }\PY{p}{\PYZob{}}
\PY{+w}{    }\PY{k+kd}{private}\PY{+w}{ }\PY{k+kd}{final}\PY{+w}{ }\PY{n}{String}\PY{+w}{ }\PY{n}{customerId}\PY{p}{;}
\PY{+w}{    }\PY{k+kd}{private}\PY{+w}{ }\PY{k+kd}{final}\PY{+w}{ }\PY{n}{LocalDate}\PY{+w}{ }\PY{n}{orderDate}\PY{p}{;}

\PY{+w}{    }\PY{k+kd}{public}\PY{+w}{ }\PY{n+nf}{OrderKey}\PY{p}{(}\PY{n}{String}\PY{+w}{ }\PY{n}{customerId}\PY{p}{,}\PY{+w}{ }\PY{n}{LocalDate}\PY{+w}{ }\PY{n}{orderDate}\PY{p}{)}\PY{+w}{ }\PY{p}{\PYZob{}}
\PY{+w}{        }\PY{k}{this}\PY{p}{.}\PY{n+na}{customerId}\PY{+w}{ }\PY{o}{=}\PY{+w}{ }\PY{n}{Objects}\PY{p}{.}\PY{n+na}{requireNonNull}\PY{p}{(}\PY{n}{customerId}\PY{p}{)}\PY{p}{;}
\PY{+w}{        }\PY{k}{this}\PY{p}{.}\PY{n+na}{orderDate}\PY{+w}{ }\PY{o}{=}\PY{+w}{ }\PY{n}{Objects}\PY{p}{.}\PY{n+na}{requireNonNull}\PY{p}{(}\PY{n}{orderDate}\PY{p}{)}\PY{p}{;}
\PY{+w}{    }\PY{p}{\PYZcb{}}

\PY{+w}{    }\PY{k+kd}{public}\PY{+w}{ }\PY{n}{OrderKey}\PY{+w}{ }\PY{n+nf}{withDate}\PY{p}{(}\PY{n}{LocalDate}\PY{+w}{ }\PY{n}{correctedDate}\PY{p}{)}\PY{+w}{ }\PY{p}{\PYZob{}}
\PY{+w}{        }\PY{k}{return}\PY{+w}{ }\PY{k}{new}\PY{+w}{ }\PY{n}{OrderKey}\PY{p}{(}\PY{n}{customerId}\PY{p}{,}\PY{+w}{ }\PY{n}{correctedDate}\PY{p}{)}\PY{p}{;}\PY{+w}{ }\PY{c+c1}{// new key, old one untouched}
\PY{+w}{    }\PY{p}{\PYZcb{}}

\PY{+w}{    }\PY{n+nd}{@Override}
\PY{+w}{    }\PY{k+kd}{public}\PY{+w}{ }\PY{k+kt}{boolean}\PY{+w}{ }\PY{n+nf}{equals}\PY{p}{(}\PY{n}{Object}\PY{+w}{ }\PY{n}{o}\PY{p}{)}\PY{+w}{ }\PY{p}{\PYZob{}}
\PY{+w}{        }\PY{k}{if}\PY{+w}{ }\PY{p}{(}\PY{o}{!}\PY{p}{(}\PY{n}{o}\PY{+w}{ }\PY{k}{instanceof}\PY{+w}{ }\PY{n}{OrderKey}\PY{+w}{ }\PY{n}{other}\PY{p}{)}\PY{p}{)}\PY{+w}{ }\PY{k}{return}\PY{+w}{ }\PY{k+kc}{false}\PY{p}{;}
\PY{+w}{        }\PY{k}{return}\PY{+w}{ }\PY{n}{customerId}\PY{p}{.}\PY{n+na}{equals}\PY{p}{(}\PY{n}{other}\PY{p}{.}\PY{n+na}{customerId}\PY{p}{)}\PY{+w}{ }\PY{o}{\PYZam{}}\PY{o}{\PYZam{}}\PY{+w}{ }\PY{n}{orderDate}\PY{p}{.}\PY{n+na}{equals}\PY{p}{(}\PY{n}{other}\PY{p}{.}\PY{n+na}{orderDate}\PY{p}{)}\PY{p}{;}
\PY{+w}{    }\PY{p}{\PYZcb{}}

\PY{+w}{    }\PY{n+nd}{@Override}
\PY{+w}{    }\PY{k+kd}{public}\PY{+w}{ }\PY{k+kt}{int}\PY{+w}{ }\PY{n+nf}{hashCode}\PY{p}{(}\PY{p}{)}\PY{+w}{ }\PY{p}{\PYZob{}}
\PY{+w}{        }\PY{k}{return}\PY{+w}{ }\PY{n}{Objects}\PY{p}{.}\PY{n+na}{hash}\PY{p}{(}\PY{n}{customerId}\PY{p}{,}\PY{+w}{ }\PY{n}{orderDate}\PY{p}{)}\PY{p}{;}
\PY{+w}{    }\PY{p}{\PYZcb{}}
\PY{p}{\PYZcb{}}

\PY{k+kd}{public}\PY{+w}{ }\PY{k+kd}{class} \PY{n+nc}{OrderDeduplicator}\PY{+w}{ }\PY{p}{\PYZob{}}
\PY{+w}{    }\PY{k+kd}{private}\PY{+w}{ }\PY{k+kd}{final}\PY{+w}{ }\PY{n}{Map}\PY{o}{\PYZlt{}}\PY{n}{OrderKey}\PY{p}{,}\PY{+w}{ }\PY{n}{Order}\PY{o}{\PYZgt{}}\PY{+w}{ }\PY{n}{seen}\PY{+w}{ }\PY{o}{=}\PY{+w}{ }\PY{k}{new}\PY{+w}{ }\PY{n}{HashMap}\PY{o}{\PYZlt{}}\PY{o}{\PYZgt{}}\PY{p}{(}\PY{p}{)}\PY{p}{;}

\PY{+w}{    }\PY{k+kd}{public}\PY{+w}{ }\PY{k+kt}{boolean}\PY{+w}{ }\PY{n+nf}{isDuplicate}\PY{p}{(}\PY{n}{OrderKey}\PY{+w}{ }\PY{n}{key}\PY{p}{,}\PY{+w}{ }\PY{n}{Order}\PY{+w}{ }\PY{n}{order}\PY{p}{,}\PY{+w}{ }\PY{n}{LocalDate}\PY{+w}{ }\PY{n}{correctedDate}\PY{p}{)}\PY{+w}{ }\PY{p}{\PYZob{}}
\PY{+w}{        }\PY{n}{OrderKey}\PY{+w}{ }\PY{n}{correctedKey}\PY{+w}{ }\PY{o}{=}\PY{+w}{ }\PY{n}{key}\PY{p}{.}\PY{n+na}{withDate}\PY{p}{(}\PY{n}{correctedDate}\PY{p}{)}\PY{p}{;}
\PY{+w}{        }\PY{k+kt}{boolean}\PY{+w}{ }\PY{n}{duplicate}\PY{+w}{ }\PY{o}{=}\PY{+w}{ }\PY{n}{seen}\PY{p}{.}\PY{n+na}{containsKey}\PY{p}{(}\PY{n}{correctedKey}\PY{p}{)}\PY{p}{;}
\PY{+w}{        }\PY{n}{seen}\PY{p}{.}\PY{n+na}{put}\PY{p}{(}\PY{n}{correctedKey}\PY{p}{,}\PY{+w}{ }\PY{n}{order}\PY{p}{)}\PY{p}{;}
\PY{+w}{        }\PY{k}{return}\PY{+w}{ }\PY{n}{duplicate}\PY{p}{;}
\PY{+w}{    }\PY{p}{\PYZcb{}}
\PY{p}{\PYZcb{}}
\end{Verbatim}
\end{codebox}

A record (\hyperref[chap:3]{Chapter~3, \emph{Modern Java Language Features}}) would have made this mistake impossible by construction, since record components can’t have setters.

\end{sidebar}
\subsection[Scenario 4: The Report Exporter]{Scenario 4: The Report Exporter}
\label{h:22:scenario-4-the-report-exporter}
\begin{codebox}
\begin{Verbatim}[commandchars=\\\{\}]
\PY{k+kd}{public}\PY{+w}{ }\PY{k+kd}{class} \PY{n+nc}{ReportExporter}\PY{+w}{ }\PY{p}{\PYZob{}}

\PY{+w}{    }\PY{k+kd}{public}\PY{+w}{ }\PY{n}{String}\PY{+w}{ }\PY{n+nf}{loadTemplate}\PY{p}{(}\PY{n}{String}\PY{+w}{ }\PY{n}{templatePath}\PY{p}{)}\PY{+w}{ }\PY{k+kd}{throws}\PY{+w}{ }\PY{n}{IOException}\PY{+w}{ }\PY{p}{\PYZob{}}
\PY{+w}{        }\PY{n}{FileInputStream}\PY{+w}{ }\PY{n}{fis}\PY{+w}{ }\PY{o}{=}\PY{+w}{ }\PY{k}{new}\PY{+w}{ }\PY{n}{FileInputStream}\PY{p}{(}\PY{n}{templatePath}\PY{p}{)}\PY{p}{;}
\PY{+w}{        }\PY{k+kt}{byte}\PY{o}{[}\PY{o}{]}\PY{+w}{ }\PY{n}{data}\PY{+w}{ }\PY{o}{=}\PY{+w}{ }\PY{n}{fis}\PY{p}{.}\PY{n+na}{readAllBytes}\PY{p}{(}\PY{p}{)}\PY{p}{;}
\PY{+w}{        }\PY{k}{return}\PY{+w}{ }\PY{k}{new}\PY{+w}{ }\PY{n}{String}\PY{p}{(}\PY{n}{data}\PY{p}{,}\PY{+w}{ }\PY{n}{StandardCharsets}\PY{p}{.}\PY{n+na}{UTF\PYZus{}8}\PY{p}{)}\PY{p}{;}
\PY{+w}{    }\PY{p}{\PYZcb{}}

\PY{+w}{    }\PY{k+kd}{public}\PY{+w}{ }\PY{k+kt}{void}\PY{+w}{ }\PY{n+nf}{exportToFile}\PY{p}{(}\PY{n}{String}\PY{+w}{ }\PY{n}{content}\PY{p}{,}\PY{+w}{ }\PY{n}{String}\PY{+w}{ }\PY{n}{outputPath}\PY{p}{)}\PY{+w}{ }\PY{k+kd}{throws}\PY{+w}{ }\PY{n}{IOException}\PY{+w}{ }\PY{p}{\PYZob{}}
\PY{+w}{        }\PY{n}{FileOutputStream}\PY{+w}{ }\PY{n}{fos}\PY{+w}{ }\PY{o}{=}\PY{+w}{ }\PY{k}{new}\PY{+w}{ }\PY{n}{FileOutputStream}\PY{p}{(}\PY{n}{outputPath}\PY{p}{)}\PY{p}{;}
\PY{+w}{        }\PY{n}{Connection}\PY{+w}{ }\PY{n}{conn}\PY{+w}{ }\PY{o}{=}\PY{+w}{ }\PY{n}{dataSource}\PY{p}{.}\PY{n+na}{getConnection}\PY{p}{(}\PY{p}{)}\PY{p}{;}
\PY{+w}{        }\PY{n}{PreparedStatement}\PY{+w}{ }\PY{n}{stmt}\PY{+w}{ }\PY{o}{=}\PY{+w}{ }\PY{n}{conn}\PY{p}{.}\PY{n+na}{prepareStatement}\PY{p}{(}\PY{l+s}{\PYZdq{}}\PY{l+s}{INSERT INTO exports(path) VALUES (?)}\PY{l+s}{\PYZdq{}}\PY{p}{)}\PY{p}{;}
\PY{+w}{        }\PY{n}{stmt}\PY{p}{.}\PY{n+na}{setString}\PY{p}{(}\PY{l+m+mi}{1}\PY{p}{,}\PY{+w}{ }\PY{n}{outputPath}\PY{p}{)}\PY{p}{;}
\PY{+w}{        }\PY{n}{stmt}\PY{p}{.}\PY{n+na}{executeUpdate}\PY{p}{(}\PY{p}{)}\PY{p}{;}
\PY{+w}{        }\PY{n}{fos}\PY{p}{.}\PY{n+na}{write}\PY{p}{(}\PY{n}{content}\PY{p}{.}\PY{n+na}{getBytes}\PY{p}{(}\PY{n}{StandardCharsets}\PY{p}{.}\PY{n+na}{UTF\PYZus{}8}\PY{p}{)}\PY{p}{)}\PY{p}{;}
\PY{+w}{        }\PY{n}{fos}\PY{p}{.}\PY{n+na}{close}\PY{p}{(}\PY{p}{)}\PY{p}{;}
\PY{+w}{    }\PY{p}{\PYZcb{}}

\PY{+w}{    }\PY{k+kd}{private}\PY{+w}{ }\PY{n}{DataSource}\PY{+w}{ }\PY{n}{dataSource}\PY{p}{;}
\PY{p}{\PYZcb{}}
\end{Verbatim}
\end{codebox}

\begin{sidebar}{Review}
\begin{itemize}
\item 
\textbf{Blocker — \code{loadTemplate} never closes \code{fis}.} Every call leaks a file descriptor; under load this exhausts the process’s file descriptor limit and every subsequent I/O call in the JVM starts failing, not just this method.

\item 
\textbf{Blocker — \code{exportToFile} leaks \code{conn} and \code{stmt} unconditionally}, and additionally leaks \code{fos} if \code{stmt.\allowbreak{}executeUpdate()} throws (e.g., a SQL error), because \code{fos.\allowbreak{}close()} is only reached on the successful path. Database connections are a scarcer, more expensive resource than file handles — this will exhaust the connection pool quickly.

\item 
\textbf{Major — mixing an I/O write and a DB write with no transactional/ordering thought.} If the DB insert succeeds but the file write fails (or vice versa), the export record and the exported file disappear out of sync with no rollback or compensation. Not the focus of this exercise, but worth a one-line flag.

\end{itemize}

Fixed code — try-with-resources for every \code{AutoCloseable}, guaranteeing cleanup on every path including exceptions (see \hyperref[chap:6]{Chapter~6, \emph{Exception Handling}}):

\begin{codebox}
\begin{Verbatim}[commandchars=\\\{\}]
\PY{k+kd}{public}\PY{+w}{ }\PY{k+kd}{class} \PY{n+nc}{ReportExporter}\PY{+w}{ }\PY{p}{\PYZob{}}

\PY{+w}{    }\PY{k+kd}{public}\PY{+w}{ }\PY{n}{String}\PY{+w}{ }\PY{n+nf}{loadTemplate}\PY{p}{(}\PY{n}{String}\PY{+w}{ }\PY{n}{templatePath}\PY{p}{)}\PY{+w}{ }\PY{k+kd}{throws}\PY{+w}{ }\PY{n}{IOException}\PY{+w}{ }\PY{p}{\PYZob{}}
\PY{+w}{        }\PY{k}{try}\PY{+w}{ }\PY{p}{(}\PY{n}{FileInputStream}\PY{+w}{ }\PY{n}{fis}\PY{+w}{ }\PY{o}{=}\PY{+w}{ }\PY{k}{new}\PY{+w}{ }\PY{n}{FileInputStream}\PY{p}{(}\PY{n}{templatePath}\PY{p}{)}\PY{p}{)}\PY{+w}{ }\PY{p}{\PYZob{}}
\PY{+w}{            }\PY{k}{return}\PY{+w}{ }\PY{k}{new}\PY{+w}{ }\PY{n}{String}\PY{p}{(}\PY{n}{fis}\PY{p}{.}\PY{n+na}{readAllBytes}\PY{p}{(}\PY{p}{)}\PY{p}{,}\PY{+w}{ }\PY{n}{StandardCharsets}\PY{p}{.}\PY{n+na}{UTF\PYZus{}8}\PY{p}{)}\PY{p}{;}
\PY{+w}{        }\PY{p}{\PYZcb{}}
\PY{+w}{    }\PY{p}{\PYZcb{}}

\PY{+w}{    }\PY{k+kd}{public}\PY{+w}{ }\PY{k+kt}{void}\PY{+w}{ }\PY{n+nf}{exportToFile}\PY{p}{(}\PY{n}{String}\PY{+w}{ }\PY{n}{content}\PY{p}{,}\PY{+w}{ }\PY{n}{String}\PY{+w}{ }\PY{n}{outputPath}\PY{p}{)}\PY{+w}{ }\PY{k+kd}{throws}\PY{+w}{ }\PY{n}{IOException}\PY{p}{,}\PY{+w}{ }\PY{n}{SQLException}\PY{+w}{ }\PY{p}{\PYZob{}}
\PY{+w}{        }\PY{k}{try}\PY{+w}{ }\PY{p}{(}\PY{n}{Connection}\PY{+w}{ }\PY{n}{conn}\PY{+w}{ }\PY{o}{=}\PY{+w}{ }\PY{n}{dataSource}\PY{p}{.}\PY{n+na}{getConnection}\PY{p}{(}\PY{p}{)}\PY{p}{;}
\PY{+w}{             }\PY{n}{PreparedStatement}\PY{+w}{ }\PY{n}{stmt}\PY{+w}{ }\PY{o}{=}\PY{+w}{ }\PY{n}{conn}\PY{p}{.}\PY{n+na}{prepareStatement}\PY{p}{(}\PY{l+s}{\PYZdq{}}\PY{l+s}{INSERT INTO exports(path) VALUES (?)}\PY{l+s}{\PYZdq{}}\PY{p}{)}\PY{p}{)}\PY{+w}{ }\PY{p}{\PYZob{}}
\PY{+w}{            }\PY{n}{stmt}\PY{p}{.}\PY{n+na}{setString}\PY{p}{(}\PY{l+m+mi}{1}\PY{p}{,}\PY{+w}{ }\PY{n}{outputPath}\PY{p}{)}\PY{p}{;}
\PY{+w}{            }\PY{n}{stmt}\PY{p}{.}\PY{n+na}{executeUpdate}\PY{p}{(}\PY{p}{)}\PY{p}{;}
\PY{+w}{        }\PY{p}{\PYZcb{}}
\PY{+w}{        }\PY{k}{try}\PY{+w}{ }\PY{p}{(}\PY{n}{FileOutputStream}\PY{+w}{ }\PY{n}{fos}\PY{+w}{ }\PY{o}{=}\PY{+w}{ }\PY{k}{new}\PY{+w}{ }\PY{n}{FileOutputStream}\PY{p}{(}\PY{n}{outputPath}\PY{p}{)}\PY{p}{)}\PY{+w}{ }\PY{p}{\PYZob{}}
\PY{+w}{            }\PY{n}{fos}\PY{p}{.}\PY{n+na}{write}\PY{p}{(}\PY{n}{content}\PY{p}{.}\PY{n+na}{getBytes}\PY{p}{(}\PY{n}{StandardCharsets}\PY{p}{.}\PY{n+na}{UTF\PYZus{}8}\PY{p}{)}\PY{p}{)}\PY{p}{;}
\PY{+w}{        }\PY{p}{\PYZcb{}}
\PY{+w}{    }\PY{p}{\PYZcb{}}

\PY{+w}{    }\PY{k+kd}{private}\PY{+w}{ }\PY{n}{DataSource}\PY{+w}{ }\PY{n}{dataSource}\PY{p}{;}
\PY{p}{\PYZcb{}}
\end{Verbatim}
\end{codebox}

\end{sidebar}
\subsection[Scenario 5: The Payment Processor]{Scenario 5: The Payment Processor}
\label{h:22:scenario-5-the-payment-processor}
\begin{codebox}
\begin{Verbatim}[commandchars=\\\{\}]
\PY{k+kd}{public}\PY{+w}{ }\PY{k+kd}{class} \PY{n+nc}{PaymentProcessor}\PY{+w}{ }\PY{p}{\PYZob{}}

\PY{+w}{    }\PY{k+kd}{public}\PY{+w}{ }\PY{n}{PaymentResult}\PY{+w}{ }\PY{n+nf}{process}\PY{p}{(}\PY{n}{PaymentRequest}\PY{+w}{ }\PY{n}{request}\PY{p}{)}\PY{+w}{ }\PY{p}{\PYZob{}}
\PY{+w}{        }\PY{k}{try}\PY{+w}{ }\PY{p}{\PYZob{}}
\PY{+w}{            }\PY{n}{validate}\PY{p}{(}\PY{n}{request}\PY{p}{)}\PY{p}{;}
\PY{+w}{            }\PY{n}{ChargeResult}\PY{+w}{ }\PY{n}{charge}\PY{+w}{ }\PY{o}{=}\PY{+w}{ }\PY{n}{gateway}\PY{p}{.}\PY{n+na}{charge}\PY{p}{(}\PY{n}{request}\PY{p}{.}\PY{n+na}{amount}\PY{p}{(}\PY{p}{)}\PY{p}{,}\PY{+w}{ }\PY{n}{request}\PY{p}{.}\PY{n+na}{cardToken}\PY{p}{(}\PY{p}{)}\PY{p}{)}\PY{p}{;}
\PY{+w}{            }\PY{n}{ledger}\PY{p}{.}\PY{n+na}{record}\PY{p}{(}\PY{n}{request}\PY{p}{.}\PY{n+na}{orderId}\PY{p}{(}\PY{p}{)}\PY{p}{,}\PY{+w}{ }\PY{n}{charge}\PY{p}{.}\PY{n+na}{transactionId}\PY{p}{(}\PY{p}{)}\PY{p}{)}\PY{p}{;}
\PY{+w}{            }\PY{k}{return}\PY{+w}{ }\PY{n}{PaymentResult}\PY{p}{.}\PY{n+na}{success}\PY{p}{(}\PY{n}{charge}\PY{p}{.}\PY{n+na}{transactionId}\PY{p}{(}\PY{p}{)}\PY{p}{)}\PY{p}{;}
\PY{+w}{        }\PY{p}{\PYZcb{}}\PY{+w}{ }\PY{k}{catch}\PY{+w}{ }\PY{p}{(}\PY{n}{Exception}\PY{+w}{ }\PY{n}{e}\PY{p}{)}\PY{+w}{ }\PY{p}{\PYZob{}}
\PY{+w}{            }\PY{k}{return}\PY{+w}{ }\PY{n}{PaymentResult}\PY{p}{.}\PY{n+na}{failure}\PY{p}{(}\PY{l+s}{\PYZdq{}}\PY{l+s}{Payment could not be processed}\PY{l+s}{\PYZdq{}}\PY{p}{)}\PY{p}{;}
\PY{+w}{        }\PY{p}{\PYZcb{}}
\PY{+w}{    }\PY{p}{\PYZcb{}}

\PY{+w}{    }\PY{k+kd}{private}\PY{+w}{ }\PY{k+kt}{void}\PY{+w}{ }\PY{n+nf}{validate}\PY{p}{(}\PY{n}{PaymentRequest}\PY{+w}{ }\PY{n}{request}\PY{p}{)}\PY{+w}{ }\PY{p}{\PYZob{}}
\PY{+w}{        }\PY{k}{if}\PY{+w}{ }\PY{p}{(}\PY{n}{request}\PY{p}{.}\PY{n+na}{amount}\PY{p}{(}\PY{p}{)}\PY{p}{.}\PY{n+na}{compareTo}\PY{p}{(}\PY{n}{BigDecimal}\PY{p}{.}\PY{n+na}{ZERO}\PY{p}{)}\PY{+w}{ }\PY{o}{\PYZlt{}}\PY{o}{=}\PY{+w}{ }\PY{l+m+mi}{0}\PY{p}{)}\PY{+w}{ }\PY{p}{\PYZob{}}
\PY{+w}{            }\PY{k}{throw}\PY{+w}{ }\PY{k}{new}\PY{+w}{ }\PY{n}{IllegalArgumentException}\PY{p}{(}\PY{l+s}{\PYZdq{}}\PY{l+s}{amount must be positive}\PY{l+s}{\PYZdq{}}\PY{p}{)}\PY{p}{;}
\PY{+w}{        }\PY{p}{\PYZcb{}}
\PY{+w}{    }\PY{p}{\PYZcb{}}
\PY{p}{\PYZcb{}}
\end{Verbatim}
\end{codebox}

\begin{sidebar}{Review}
\begin{itemize}
\item 
\textbf{Blocker — the exception is swallowed with no logging and no rethrow.} \code{catch~\allowbreak{}(\allowbreak{}Exception~\allowbreak{}e)\allowbreak{}~\allowbreak{}\{\allowbreak{}~\allowbreak{}return~\allowbreak{}Payment\allowbreak{}Result.\allowbreak{}failure(...);\allowbreak{}~\allowbreak{}\}} discards every detail about \emph{why} a payment failed — network timeout, declined card, a \code{NullPointerException} bug in \code{validate}, or the ledger write failing after the charge already succeeded. Support will get “payment could not be processed” tickets with zero ability to diagnose them, and real bugs will hide behind the same generic message as legitimate declines.

\item 
\textbf{Blocker — the charge can succeed but the ledger write can fail}, and both end up reported identically as a generic failure to the caller. That’s a customer charged with no record of it — a money-losing (or trust-losing) bug, not just a logging annoyance. This needs to be handled as a distinct case (compensating action, alerting, retry-safe ledger write), not lumped into the same catch.

\item 
\textbf{Major — \code{catch~\allowbreak{}(\allowbreak{}Exception~\allowbreak{}e)} is too wide.} It will also catch bugs (\code{NullPointerException}, \code{ClassCastException}) from anywhere in this call chain and report them as ordinary payment failures instead of surfacing them as defects to be fixed.

\end{itemize}

Fixed code — narrow the catch to expected failure modes, log with full context, and treat “charged but not recorded” as its own alertable case:

\begin{codebox}
\begin{Verbatim}[commandchars=\\\{\}]
\PY{k+kd}{public}\PY{+w}{ }\PY{k+kd}{class} \PY{n+nc}{PaymentProcessor}\PY{+w}{ }\PY{p}{\PYZob{}}
\PY{+w}{    }\PY{k+kd}{private}\PY{+w}{ }\PY{k+kd}{static}\PY{+w}{ }\PY{k+kd}{final}\PY{+w}{ }\PY{n}{Logger}\PY{+w}{ }\PY{n}{log}\PY{+w}{ }\PY{o}{=}\PY{+w}{ }\PY{n}{LoggerFactory}\PY{p}{.}\PY{n+na}{getLogger}\PY{p}{(}\PY{n}{PaymentProcessor}\PY{p}{.}\PY{n+na}{class}\PY{p}{)}\PY{p}{;}

\PY{+w}{    }\PY{k+kd}{public}\PY{+w}{ }\PY{n}{PaymentResult}\PY{+w}{ }\PY{n+nf}{process}\PY{p}{(}\PY{n}{PaymentRequest}\PY{+w}{ }\PY{n}{request}\PY{p}{)}\PY{+w}{ }\PY{p}{\PYZob{}}
\PY{+w}{        }\PY{n}{validate}\PY{p}{(}\PY{n}{request}\PY{p}{)}\PY{p}{;}\PY{+w}{ }\PY{c+c1}{// let IllegalArgumentException propagate — it\PYZsq{}s a caller bug, fail fast}

\PY{+w}{        }\PY{n}{ChargeResult}\PY{+w}{ }\PY{n}{charge}\PY{p}{;}
\PY{+w}{        }\PY{k}{try}\PY{+w}{ }\PY{p}{\PYZob{}}
\PY{+w}{            }\PY{n}{charge}\PY{+w}{ }\PY{o}{=}\PY{+w}{ }\PY{n}{gateway}\PY{p}{.}\PY{n+na}{charge}\PY{p}{(}\PY{n}{request}\PY{p}{.}\PY{n+na}{amount}\PY{p}{(}\PY{p}{)}\PY{p}{,}\PY{+w}{ }\PY{n}{request}\PY{p}{.}\PY{n+na}{cardToken}\PY{p}{(}\PY{p}{)}\PY{p}{)}\PY{p}{;}
\PY{+w}{        }\PY{p}{\PYZcb{}}\PY{+w}{ }\PY{k}{catch}\PY{+w}{ }\PY{p}{(}\PY{n}{GatewayException}\PY{+w}{ }\PY{n}{e}\PY{p}{)}\PY{+w}{ }\PY{p}{\PYZob{}}
\PY{+w}{            }\PY{n}{log}\PY{p}{.}\PY{n+na}{warn}\PY{p}{(}\PY{l+s}{\PYZdq{}}\PY{l+s}{Charge declined for order \PYZob{}\PYZcb{}: \PYZob{}\PYZcb{}}\PY{l+s}{\PYZdq{}}\PY{p}{,}\PY{+w}{ }\PY{n}{request}\PY{p}{.}\PY{n+na}{orderId}\PY{p}{(}\PY{p}{)}\PY{p}{,}\PY{+w}{ }\PY{n}{e}\PY{p}{.}\PY{n+na}{getMessage}\PY{p}{(}\PY{p}{)}\PY{p}{)}\PY{p}{;}
\PY{+w}{            }\PY{k}{return}\PY{+w}{ }\PY{n}{PaymentResult}\PY{p}{.}\PY{n+na}{failure}\PY{p}{(}\PY{l+s}{\PYZdq{}}\PY{l+s}{Payment declined: }\PY{l+s}{\PYZdq{}}\PY{+w}{ }\PY{o}{+}\PY{+w}{ }\PY{n}{e}\PY{p}{.}\PY{n+na}{getMessage}\PY{p}{(}\PY{p}{)}\PY{p}{)}\PY{p}{;}
\PY{+w}{        }\PY{p}{\PYZcb{}}

\PY{+w}{        }\PY{k}{try}\PY{+w}{ }\PY{p}{\PYZob{}}
\PY{+w}{            }\PY{n}{ledger}\PY{p}{.}\PY{n+na}{record}\PY{p}{(}\PY{n}{request}\PY{p}{.}\PY{n+na}{orderId}\PY{p}{(}\PY{p}{)}\PY{p}{,}\PY{+w}{ }\PY{n}{charge}\PY{p}{.}\PY{n+na}{transactionId}\PY{p}{(}\PY{p}{)}\PY{p}{)}\PY{p}{;}
\PY{+w}{        }\PY{p}{\PYZcb{}}\PY{+w}{ }\PY{k}{catch}\PY{+w}{ }\PY{p}{(}\PY{n}{LedgerException}\PY{+w}{ }\PY{n}{e}\PY{p}{)}\PY{+w}{ }\PY{p}{\PYZob{}}
\PY{+w}{            }\PY{c+c1}{// Money moved but we failed to record it — this must page someone, not just log.}
\PY{+w}{            }\PY{n}{log}\PY{p}{.}\PY{n+na}{error}\PY{p}{(}\PY{l+s}{\PYZdq{}}\PY{l+s}{CRITICAL: charge \PYZob{}\PYZcb{} succeeded but ledger record failed for order \PYZob{}\PYZcb{}}\PY{l+s}{\PYZdq{}}\PY{p}{,}
\PY{+w}{                    }\PY{n}{charge}\PY{p}{.}\PY{n+na}{transactionId}\PY{p}{(}\PY{p}{)}\PY{p}{,}\PY{+w}{ }\PY{n}{request}\PY{p}{.}\PY{n+na}{orderId}\PY{p}{(}\PY{p}{)}\PY{p}{,}\PY{+w}{ }\PY{n}{e}\PY{p}{)}\PY{p}{;}
\PY{+w}{            }\PY{n}{alerting}\PY{p}{.}\PY{n+na}{pageOnCall}\PY{p}{(}\PY{l+s}{\PYZdq{}}\PY{l+s}{Unreconciled charge }\PY{l+s}{\PYZdq{}}\PY{+w}{ }\PY{o}{+}\PY{+w}{ }\PY{n}{charge}\PY{p}{.}\PY{n+na}{transactionId}\PY{p}{(}\PY{p}{)}\PY{p}{)}\PY{p}{;}
\PY{+w}{            }\PY{k}{throw}\PY{+w}{ }\PY{k}{new}\PY{+w}{ }\PY{n}{UnreconciledChargeException}\PY{p}{(}\PY{n}{request}\PY{p}{.}\PY{n+na}{orderId}\PY{p}{(}\PY{p}{)}\PY{p}{,}\PY{+w}{ }\PY{n}{charge}\PY{p}{.}\PY{n+na}{transactionId}\PY{p}{(}\PY{p}{)}\PY{p}{,}\PY{+w}{ }\PY{n}{e}\PY{p}{)}\PY{p}{;}
\PY{+w}{        }\PY{p}{\PYZcb{}}

\PY{+w}{        }\PY{k}{return}\PY{+w}{ }\PY{n}{PaymentResult}\PY{p}{.}\PY{n+na}{success}\PY{p}{(}\PY{n}{charge}\PY{p}{.}\PY{n+na}{transactionId}\PY{p}{(}\PY{p}{)}\PY{p}{)}\PY{p}{;}
\PY{+w}{    }\PY{p}{\PYZcb{}}

\PY{+w}{    }\PY{k+kd}{private}\PY{+w}{ }\PY{k+kt}{void}\PY{+w}{ }\PY{n+nf}{validate}\PY{p}{(}\PY{n}{PaymentRequest}\PY{+w}{ }\PY{n}{request}\PY{p}{)}\PY{+w}{ }\PY{p}{\PYZob{}}
\PY{+w}{        }\PY{k}{if}\PY{+w}{ }\PY{p}{(}\PY{n}{request}\PY{p}{.}\PY{n+na}{amount}\PY{p}{(}\PY{p}{)}\PY{p}{.}\PY{n+na}{compareTo}\PY{p}{(}\PY{n}{BigDecimal}\PY{p}{.}\PY{n+na}{ZERO}\PY{p}{)}\PY{+w}{ }\PY{o}{\PYZlt{}}\PY{o}{=}\PY{+w}{ }\PY{l+m+mi}{0}\PY{p}{)}\PY{+w}{ }\PY{p}{\PYZob{}}
\PY{+w}{            }\PY{k}{throw}\PY{+w}{ }\PY{k}{new}\PY{+w}{ }\PY{n}{IllegalArgumentException}\PY{p}{(}\PY{l+s}{\PYZdq{}}\PY{l+s}{amount must be positive}\PY{l+s}{\PYZdq{}}\PY{p}{)}\PY{p}{;}
\PY{+w}{        }\PY{p}{\PYZcb{}}
\PY{+w}{    }\PY{p}{\PYZcb{}}
\PY{p}{\PYZcb{}}
\end{Verbatim}
\end{codebox}

\end{sidebar}
\subsection[Scenario 6: The Order Summary Builder]{Scenario 6: The Order Summary Builder}
\label{h:22:scenario-6-the-order-summary-builder}
\begin{codebox}
\begin{Verbatim}[commandchars=\\\{\}]
\PY{k+kd}{public}\PY{+w}{ }\PY{k+kd}{class} \PY{n+nc}{OrderSummaryBuilder}\PY{+w}{ }\PY{p}{\PYZob{}}

\PY{+w}{    }\PY{k+kd}{public}\PY{+w}{ }\PY{n}{List}\PY{o}{\PYZlt{}}\PY{n}{OrderSummary}\PY{o}{\PYZgt{}}\PY{+w}{ }\PY{n+nf}{buildSummaries}\PY{p}{(}\PY{n}{List}\PY{o}{\PYZlt{}}\PY{n}{Order}\PY{o}{\PYZgt{}}\PY{+w}{ }\PY{n}{orders}\PY{p}{)}\PY{+w}{ }\PY{p}{\PYZob{}}
\PY{+w}{        }\PY{n}{List}\PY{o}{\PYZlt{}}\PY{n}{OrderSummary}\PY{o}{\PYZgt{}}\PY{+w}{ }\PY{n}{summaries}\PY{+w}{ }\PY{o}{=}\PY{+w}{ }\PY{k}{new}\PY{+w}{ }\PY{n}{ArrayList}\PY{o}{\PYZlt{}}\PY{o}{\PYZgt{}}\PY{p}{(}\PY{p}{)}\PY{p}{;}
\PY{+w}{        }\PY{k}{for}\PY{+w}{ }\PY{p}{(}\PY{n}{Order}\PY{+w}{ }\PY{n}{order}\PY{+w}{ }\PY{p}{:}\PY{+w}{ }\PY{n}{orders}\PY{p}{)}\PY{+w}{ }\PY{p}{\PYZob{}}
\PY{+w}{            }\PY{n}{Customer}\PY{+w}{ }\PY{n}{customer}\PY{+w}{ }\PY{o}{=}\PY{+w}{ }\PY{n}{customerRepository}\PY{p}{.}\PY{n+na}{findById}\PY{p}{(}\PY{n}{order}\PY{p}{.}\PY{n+na}{customerId}\PY{p}{(}\PY{p}{)}\PY{p}{)}\PY{p}{;}
\PY{+w}{            }\PY{n}{List}\PY{o}{\PYZlt{}}\PY{n}{LineItem}\PY{o}{\PYZgt{}}\PY{+w}{ }\PY{n}{items}\PY{+w}{ }\PY{o}{=}\PY{+w}{ }\PY{n}{lineItemRepository}\PY{p}{.}\PY{n+na}{findByOrderId}\PY{p}{(}\PY{n}{order}\PY{p}{.}\PY{n+na}{id}\PY{p}{(}\PY{p}{)}\PY{p}{)}\PY{p}{;}
\PY{+w}{            }\PY{n}{BigDecimal}\PY{+w}{ }\PY{n}{total}\PY{+w}{ }\PY{o}{=}\PY{+w}{ }\PY{n}{BigDecimal}\PY{p}{.}\PY{n+na}{ZERO}\PY{p}{;}
\PY{+w}{            }\PY{k}{for}\PY{+w}{ }\PY{p}{(}\PY{n}{LineItem}\PY{+w}{ }\PY{n}{item}\PY{+w}{ }\PY{p}{:}\PY{+w}{ }\PY{n}{items}\PY{p}{)}\PY{+w}{ }\PY{p}{\PYZob{}}
\PY{+w}{                }\PY{k}{for}\PY{+w}{ }\PY{p}{(}\PY{n}{LineItem}\PY{+w}{ }\PY{n}{discountRule}\PY{+w}{ }\PY{p}{:}\PY{+w}{ }\PY{n}{discountRules}\PY{p}{)}\PY{+w}{ }\PY{p}{\PYZob{}}
\PY{+w}{                    }\PY{k}{if}\PY{+w}{ }\PY{p}{(}\PY{n}{discountRule}\PY{p}{.}\PY{n+na}{appliesTo}\PY{p}{(}\PY{n}{item}\PY{p}{)}\PY{p}{)}\PY{+w}{ }\PY{p}{\PYZob{}}
\PY{+w}{                        }\PY{n}{item}\PY{+w}{ }\PY{o}{=}\PY{+w}{ }\PY{n}{discountRule}\PY{p}{.}\PY{n+na}{apply}\PY{p}{(}\PY{n}{item}\PY{p}{)}\PY{p}{;}
\PY{+w}{                    }\PY{p}{\PYZcb{}}
\PY{+w}{                }\PY{p}{\PYZcb{}}
\PY{+w}{                }\PY{n}{total}\PY{+w}{ }\PY{o}{=}\PY{+w}{ }\PY{n}{total}\PY{p}{.}\PY{n+na}{add}\PY{p}{(}\PY{n}{item}\PY{p}{.}\PY{n+na}{price}\PY{p}{(}\PY{p}{)}\PY{p}{)}\PY{p}{;}
\PY{+w}{            }\PY{p}{\PYZcb{}}
\PY{+w}{            }\PY{n}{summaries}\PY{p}{.}\PY{n+na}{add}\PY{p}{(}\PY{k}{new}\PY{+w}{ }\PY{n}{OrderSummary}\PY{p}{(}\PY{n}{order}\PY{p}{.}\PY{n+na}{id}\PY{p}{(}\PY{p}{)}\PY{p}{,}\PY{+w}{ }\PY{n}{customer}\PY{p}{.}\PY{n+na}{name}\PY{p}{(}\PY{p}{)}\PY{p}{,}\PY{+w}{ }\PY{n}{total}\PY{p}{)}\PY{p}{)}\PY{p}{;}
\PY{+w}{        }\PY{p}{\PYZcb{}}
\PY{+w}{        }\PY{k}{return}\PY{+w}{ }\PY{n}{summaries}\PY{p}{;}
\PY{+w}{    }\PY{p}{\PYZcb{}}

\PY{+w}{    }\PY{k+kd}{private}\PY{+w}{ }\PY{n}{CustomerRepository}\PY{+w}{ }\PY{n}{customerRepository}\PY{p}{;}
\PY{+w}{    }\PY{k+kd}{private}\PY{+w}{ }\PY{n}{LineItemRepository}\PY{+w}{ }\PY{n}{lineItemRepository}\PY{p}{;}
\PY{+w}{    }\PY{k+kd}{private}\PY{+w}{ }\PY{n}{List}\PY{o}{\PYZlt{}}\PY{n}{LineItem}\PY{o}{\PYZgt{}}\PY{+w}{ }\PY{n}{discountRules}\PY{p}{;}
\PY{p}{\PYZcb{}}
\end{Verbatim}
\end{codebox}

\begin{sidebar}{Review}
\begin{itemize}
\item 
\textbf{Blocker — classic N+1 query pattern.} For every single order, \code{customer\allowbreak{}Repository.\allowbreak{}find\allowbreak{}By\allowbreak{}Id(...)} and \code{line\allowbreak{}Item\allowbreak{}Repository.\allowbreak{}find\allowbreak{}By\allowbreak{}Order\allowbreak{}Id(...)} each issue a separate round trip to the database. For 1,000 orders that’s 2,000+ queries where 2-3 would do. This is the single most common real-world performance bug found in code review, and it degrades gracefully in dev (small data) and catastrophically in production.

\item 
\textbf{Major — O(n × m) discount-rule loop nested inside the per-order loop}, itself nested inside the per-line-item loop — for \code{orders~\allowbreak{}×~\allowbreak{}items~\allowbreak{}×~\allowbreak{}rules} this is cubic in the worst case relative to the input sizes, when the rules could be pre-indexed or applied via a single pass per item without needing to be recomputed identically for every order.

\item 
\textbf{Minor — \code{discount\allowbreak{}Rule.\allowbreak{}apply(\allowbreak{}item)} reassigns the loop variable \code{item}}, which is legal but confusing: it looks like line items and discount rules are the same type (\code{LineItem}), which they likely aren’t conceptually — this is either a copy-paste type mismatch or a naming/API smell worth a question to the author.

\end{itemize}

Fixed code — batch-load customers and line items once, outside the per-order loop:

\begin{codebox}
\begin{Verbatim}[commandchars=\\\{\}]
\PY{k+kd}{public}\PY{+w}{ }\PY{k+kd}{class} \PY{n+nc}{OrderSummaryBuilder}\PY{+w}{ }\PY{p}{\PYZob{}}

\PY{+w}{    }\PY{k+kd}{public}\PY{+w}{ }\PY{n}{List}\PY{o}{\PYZlt{}}\PY{n}{OrderSummary}\PY{o}{\PYZgt{}}\PY{+w}{ }\PY{n+nf}{buildSummaries}\PY{p}{(}\PY{n}{List}\PY{o}{\PYZlt{}}\PY{n}{Order}\PY{o}{\PYZgt{}}\PY{+w}{ }\PY{n}{orders}\PY{p}{)}\PY{+w}{ }\PY{p}{\PYZob{}}
\PY{+w}{        }\PY{n}{List}\PY{o}{\PYZlt{}}\PY{n}{String}\PY{o}{\PYZgt{}}\PY{+w}{ }\PY{n}{customerIds}\PY{+w}{ }\PY{o}{=}\PY{+w}{ }\PY{n}{orders}\PY{p}{.}\PY{n+na}{stream}\PY{p}{(}\PY{p}{)}\PY{p}{.}\PY{n+na}{map}\PY{p}{(}\PY{n}{Order}\PY{p}{:}\PY{p}{:}\PY{n}{customerId}\PY{p}{)}\PY{p}{.}\PY{n+na}{distinct}\PY{p}{(}\PY{p}{)}\PY{p}{.}\PY{n+na}{toList}\PY{p}{(}\PY{p}{)}\PY{p}{;}
\PY{+w}{        }\PY{n}{List}\PY{o}{\PYZlt{}}\PY{n}{Long}\PY{o}{\PYZgt{}}\PY{+w}{ }\PY{n}{orderIds}\PY{+w}{ }\PY{o}{=}\PY{+w}{ }\PY{n}{orders}\PY{p}{.}\PY{n+na}{stream}\PY{p}{(}\PY{p}{)}\PY{p}{.}\PY{n+na}{map}\PY{p}{(}\PY{n}{Order}\PY{p}{:}\PY{p}{:}\PY{n}{id}\PY{p}{)}\PY{p}{.}\PY{n+na}{toList}\PY{p}{(}\PY{p}{)}\PY{p}{;}

\PY{+w}{        }\PY{n}{Map}\PY{o}{\PYZlt{}}\PY{n}{String}\PY{p}{,}\PY{+w}{ }\PY{n}{Customer}\PY{o}{\PYZgt{}}\PY{+w}{ }\PY{n}{customersById}\PY{+w}{ }\PY{o}{=}\PY{+w}{ }\PY{n}{customerRepository}\PY{p}{.}\PY{n+na}{findAllByIds}\PY{p}{(}\PY{n}{customerIds}\PY{p}{)}
\PY{+w}{                }\PY{p}{.}\PY{n+na}{stream}\PY{p}{(}\PY{p}{)}\PY{p}{.}\PY{n+na}{collect}\PY{p}{(}\PY{n}{Collectors}\PY{p}{.}\PY{n+na}{toMap}\PY{p}{(}\PY{n}{Customer}\PY{p}{:}\PY{p}{:}\PY{n}{id}\PY{p}{,}\PY{+w}{ }\PY{n}{c}\PY{+w}{ }\PY{o}{\PYZhy{}}\PY{o}{\PYZgt{}}\PY{+w}{ }\PY{n}{c}\PY{p}{)}\PY{p}{)}\PY{p}{;}
\PY{+w}{        }\PY{n}{Map}\PY{o}{\PYZlt{}}\PY{n}{Long}\PY{p}{,}\PY{+w}{ }\PY{n}{List}\PY{o}{\PYZlt{}}\PY{n}{LineItem}\PY{o}{\PYZgt{}}\PY{o}{\PYZgt{}}\PY{+w}{ }\PY{n}{itemsByOrderId}\PY{+w}{ }\PY{o}{=}\PY{+w}{ }\PY{n}{lineItemRepository}\PY{p}{.}\PY{n+na}{findByOrderIds}\PY{p}{(}\PY{n}{orderIds}\PY{p}{)}
\PY{+w}{                }\PY{p}{.}\PY{n+na}{stream}\PY{p}{(}\PY{p}{)}\PY{p}{.}\PY{n+na}{collect}\PY{p}{(}\PY{n}{Collectors}\PY{p}{.}\PY{n+na}{groupingBy}\PY{p}{(}\PY{n}{LineItem}\PY{p}{:}\PY{p}{:}\PY{n}{orderId}\PY{p}{)}\PY{p}{)}\PY{p}{;}

\PY{+w}{        }\PY{n}{List}\PY{o}{\PYZlt{}}\PY{n}{OrderSummary}\PY{o}{\PYZgt{}}\PY{+w}{ }\PY{n}{summaries}\PY{+w}{ }\PY{o}{=}\PY{+w}{ }\PY{k}{new}\PY{+w}{ }\PY{n}{ArrayList}\PY{o}{\PYZlt{}}\PY{o}{\PYZgt{}}\PY{p}{(}\PY{p}{)}\PY{p}{;}
\PY{+w}{        }\PY{k}{for}\PY{+w}{ }\PY{p}{(}\PY{n}{Order}\PY{+w}{ }\PY{n}{order}\PY{+w}{ }\PY{p}{:}\PY{+w}{ }\PY{n}{orders}\PY{p}{)}\PY{+w}{ }\PY{p}{\PYZob{}}
\PY{+w}{            }\PY{n}{Customer}\PY{+w}{ }\PY{n}{customer}\PY{+w}{ }\PY{o}{=}\PY{+w}{ }\PY{n}{customersById}\PY{p}{.}\PY{n+na}{get}\PY{p}{(}\PY{n}{order}\PY{p}{.}\PY{n+na}{customerId}\PY{p}{(}\PY{p}{)}\PY{p}{)}\PY{p}{;}
\PY{+w}{            }\PY{n}{List}\PY{o}{\PYZlt{}}\PY{n}{LineItem}\PY{o}{\PYZgt{}}\PY{+w}{ }\PY{n}{items}\PY{+w}{ }\PY{o}{=}\PY{+w}{ }\PY{n}{itemsByOrderId}\PY{p}{.}\PY{n+na}{getOrDefault}\PY{p}{(}\PY{n}{order}\PY{p}{.}\PY{n+na}{id}\PY{p}{(}\PY{p}{)}\PY{p}{,}\PY{+w}{ }\PY{n}{List}\PY{p}{.}\PY{n+na}{of}\PY{p}{(}\PY{p}{)}\PY{p}{)}\PY{p}{;}
\PY{+w}{            }\PY{n}{BigDecimal}\PY{+w}{ }\PY{n}{total}\PY{+w}{ }\PY{o}{=}\PY{+w}{ }\PY{n}{BigDecimal}\PY{p}{.}\PY{n+na}{ZERO}\PY{p}{;}
\PY{+w}{            }\PY{k}{for}\PY{+w}{ }\PY{p}{(}\PY{n}{LineItem}\PY{+w}{ }\PY{n}{item}\PY{+w}{ }\PY{p}{:}\PY{+w}{ }\PY{n}{items}\PY{p}{)}\PY{+w}{ }\PY{p}{\PYZob{}}
\PY{+w}{                }\PY{n}{LineItem}\PY{+w}{ }\PY{n}{discounted}\PY{+w}{ }\PY{o}{=}\PY{+w}{ }\PY{n}{applyDiscounts}\PY{p}{(}\PY{n}{item}\PY{p}{)}\PY{p}{;}
\PY{+w}{                }\PY{n}{total}\PY{+w}{ }\PY{o}{=}\PY{+w}{ }\PY{n}{total}\PY{p}{.}\PY{n+na}{add}\PY{p}{(}\PY{n}{discounted}\PY{p}{.}\PY{n+na}{price}\PY{p}{(}\PY{p}{)}\PY{p}{)}\PY{p}{;}
\PY{+w}{            }\PY{p}{\PYZcb{}}
\PY{+w}{            }\PY{n}{summaries}\PY{p}{.}\PY{n+na}{add}\PY{p}{(}\PY{k}{new}\PY{+w}{ }\PY{n}{OrderSummary}\PY{p}{(}\PY{n}{order}\PY{p}{.}\PY{n+na}{id}\PY{p}{(}\PY{p}{)}\PY{p}{,}\PY{+w}{ }\PY{n}{customer}\PY{p}{.}\PY{n+na}{name}\PY{p}{(}\PY{p}{)}\PY{p}{,}\PY{+w}{ }\PY{n}{total}\PY{p}{)}\PY{p}{)}\PY{p}{;}
\PY{+w}{        }\PY{p}{\PYZcb{}}
\PY{+w}{        }\PY{k}{return}\PY{+w}{ }\PY{n}{summaries}\PY{p}{;}
\PY{+w}{    }\PY{p}{\PYZcb{}}

\PY{+w}{    }\PY{k+kd}{private}\PY{+w}{ }\PY{n}{LineItem}\PY{+w}{ }\PY{n+nf}{applyDiscounts}\PY{p}{(}\PY{n}{LineItem}\PY{+w}{ }\PY{n}{item}\PY{p}{)}\PY{+w}{ }\PY{p}{\PYZob{}}
\PY{+w}{        }\PY{k}{for}\PY{+w}{ }\PY{p}{(}\PY{n}{DiscountRule}\PY{+w}{ }\PY{n}{rule}\PY{+w}{ }\PY{p}{:}\PY{+w}{ }\PY{n}{discountRules}\PY{p}{)}\PY{+w}{ }\PY{p}{\PYZob{}}
\PY{+w}{            }\PY{k}{if}\PY{+w}{ }\PY{p}{(}\PY{n}{rule}\PY{p}{.}\PY{n+na}{appliesTo}\PY{p}{(}\PY{n}{item}\PY{p}{)}\PY{p}{)}\PY{+w}{ }\PY{p}{\PYZob{}}
\PY{+w}{                }\PY{n}{item}\PY{+w}{ }\PY{o}{=}\PY{+w}{ }\PY{n}{rule}\PY{p}{.}\PY{n+na}{apply}\PY{p}{(}\PY{n}{item}\PY{p}{)}\PY{p}{;}
\PY{+w}{            }\PY{p}{\PYZcb{}}
\PY{+w}{        }\PY{p}{\PYZcb{}}
\PY{+w}{        }\PY{k}{return}\PY{+w}{ }\PY{n}{item}\PY{p}{;}
\PY{+w}{    }\PY{p}{\PYZcb{}}

\PY{+w}{    }\PY{k+kd}{private}\PY{+w}{ }\PY{n}{CustomerRepository}\PY{+w}{ }\PY{n}{customerRepository}\PY{p}{;}
\PY{+w}{    }\PY{k+kd}{private}\PY{+w}{ }\PY{n}{LineItemRepository}\PY{+w}{ }\PY{n}{lineItemRepository}\PY{p}{;}
\PY{+w}{    }\PY{k+kd}{private}\PY{+w}{ }\PY{n}{List}\PY{o}{\PYZlt{}}\PY{n}{DiscountRule}\PY{o}{\PYZgt{}}\PY{+w}{ }\PY{n}{discountRules}\PY{p}{;}
\PY{p}{\PYZcb{}}
\end{Verbatim}
\end{codebox}

Two round trips total, regardless of how many orders are in the list. See \hyperref[h:22:performance-optimization]{Performance Optimization} for more on spotting this pattern.

\end{sidebar}
\subsection[Scenario 7: The Audit Log Formatter]{Scenario 7: The Audit Log Formatter}
\label{h:22:scenario-7-the-audit-log-formatter}
\begin{codebox}
\begin{Verbatim}[commandchars=\\\{\}]
\PY{k+kd}{public}\PY{+w}{ }\PY{k+kd}{class} \PY{n+nc}{AuditLogFormatter}\PY{+w}{ }\PY{p}{\PYZob{}}

\PY{+w}{    }\PY{k+kd}{public}\PY{+w}{ }\PY{n}{String}\PY{+w}{ }\PY{n+nf}{formatEntries}\PY{p}{(}\PY{n}{List}\PY{o}{\PYZlt{}}\PY{n}{AuditEntry}\PY{o}{\PYZgt{}}\PY{+w}{ }\PY{n}{entries}\PY{p}{)}\PY{+w}{ }\PY{p}{\PYZob{}}
\PY{+w}{        }\PY{n}{String}\PY{+w}{ }\PY{n}{report}\PY{+w}{ }\PY{o}{=}\PY{+w}{ }\PY{l+s}{\PYZdq{}}\PY{l+s}{\PYZdq{}}\PY{p}{;}
\PY{+w}{        }\PY{n}{report}\PY{+w}{ }\PY{o}{+}\PY{o}{=}\PY{+w}{ }\PY{l+s}{\PYZdq{}}\PY{l+s}{Audit Report}\PY{l+s}{\PYZbs{}}\PY{l+s}{n}\PY{l+s}{\PYZdq{}}\PY{p}{;}
\PY{+w}{        }\PY{n}{report}\PY{+w}{ }\PY{o}{+}\PY{o}{=}\PY{+w}{ }\PY{l+s}{\PYZdq{}}\PY{l+s}{=============}\PY{l+s}{\PYZbs{}}\PY{l+s}{n}\PY{l+s}{\PYZdq{}}\PY{p}{;}
\PY{+w}{        }\PY{k}{for}\PY{+w}{ }\PY{p}{(}\PY{n}{AuditEntry}\PY{+w}{ }\PY{n}{entry}\PY{+w}{ }\PY{p}{:}\PY{+w}{ }\PY{n}{entries}\PY{p}{)}\PY{+w}{ }\PY{p}{\PYZob{}}
\PY{+w}{            }\PY{n}{report}\PY{+w}{ }\PY{o}{+}\PY{o}{=}\PY{+w}{ }\PY{n}{entry}\PY{p}{.}\PY{n+na}{timestamp}\PY{p}{(}\PY{p}{)}\PY{+w}{ }\PY{o}{+}\PY{+w}{ }\PY{l+s}{\PYZdq{}}\PY{l+s}{ | }\PY{l+s}{\PYZdq{}}\PY{+w}{ }\PY{o}{+}\PY{+w}{ }\PY{n}{entry}\PY{p}{.}\PY{n+na}{user}\PY{p}{(}\PY{p}{)}\PY{+w}{ }\PY{o}{+}\PY{+w}{ }\PY{l+s}{\PYZdq{}}\PY{l+s}{ | }\PY{l+s}{\PYZdq{}}\PY{+w}{ }\PY{o}{+}\PY{+w}{ }\PY{n}{entry}\PY{p}{.}\PY{n+na}{action}\PY{p}{(}\PY{p}{)}\PY{p}{;}
\PY{+w}{            }\PY{n}{report}\PY{+w}{ }\PY{o}{+}\PY{o}{=}\PY{+w}{ }\PY{l+s}{\PYZdq{}}\PY{l+s}{\PYZbs{}}\PY{l+s}{n}\PY{l+s}{\PYZdq{}}\PY{p}{;}
\PY{+w}{        }\PY{p}{\PYZcb{}}
\PY{+w}{        }\PY{n}{report}\PY{+w}{ }\PY{o}{+}\PY{o}{=}\PY{+w}{ }\PY{l+s}{\PYZdq{}}\PY{l+s}{Total entries: }\PY{l+s}{\PYZdq{}}\PY{+w}{ }\PY{o}{+}\PY{+w}{ }\PY{n}{entries}\PY{p}{.}\PY{n+na}{size}\PY{p}{(}\PY{p}{)}\PY{+w}{ }\PY{o}{+}\PY{+w}{ }\PY{l+s}{\PYZdq{}}\PY{l+s}{\PYZbs{}}\PY{l+s}{n}\PY{l+s}{\PYZdq{}}\PY{p}{;}
\PY{+w}{        }\PY{k}{return}\PY{+w}{ }\PY{n}{report}\PY{p}{;}
\PY{+w}{    }\PY{p}{\PYZcb{}}
\PY{p}{\PYZcb{}}
\end{Verbatim}
\end{codebox}

\begin{sidebar}{Review}
\begin{itemize}
\item 
\textbf{Major — string concatenation with \code{+=} inside a loop.} \code{String} is immutable, so every \code{+=} allocates a brand-new \code{String} and copies everything seen so far into it. For \code{n} entries, that’s roughly O(n²) total character copying. For a handful of audit entries this is invisible; for a report over tens of thousands of entries it becomes a real, measurable slowdown and a burst of short-lived garbage for the GC to clean up.

\item 
\textbf{Nit — mixing string-building style} (some lines start fresh, some append) makes the method harder to scan; consolidating onto one \code{StringBuilder} also reads more clearly as “building one thing incrementally.”

\end{itemize}

Fixed code — a single \code{StringBuilder} for the whole method:

\begin{codebox}
\begin{Verbatim}[commandchars=\\\{\}]
\PY{k+kd}{public}\PY{+w}{ }\PY{k+kd}{class} \PY{n+nc}{AuditLogFormatter}\PY{+w}{ }\PY{p}{\PYZob{}}

\PY{+w}{    }\PY{k+kd}{public}\PY{+w}{ }\PY{n}{String}\PY{+w}{ }\PY{n+nf}{formatEntries}\PY{p}{(}\PY{n}{List}\PY{o}{\PYZlt{}}\PY{n}{AuditEntry}\PY{o}{\PYZgt{}}\PY{+w}{ }\PY{n}{entries}\PY{p}{)}\PY{+w}{ }\PY{p}{\PYZob{}}
\PY{+w}{        }\PY{n}{StringBuilder}\PY{+w}{ }\PY{n}{report}\PY{+w}{ }\PY{o}{=}\PY{+w}{ }\PY{k}{new}\PY{+w}{ }\PY{n}{StringBuilder}\PY{p}{(}\PY{p}{)}\PY{p}{;}
\PY{+w}{        }\PY{n}{report}\PY{p}{.}\PY{n+na}{append}\PY{p}{(}\PY{l+s}{\PYZdq{}}\PY{l+s}{Audit Report}\PY{l+s}{\PYZbs{}}\PY{l+s}{n}\PY{l+s}{\PYZdq{}}\PY{p}{)}\PY{p}{;}
\PY{+w}{        }\PY{n}{report}\PY{p}{.}\PY{n+na}{append}\PY{p}{(}\PY{l+s}{\PYZdq{}}\PY{l+s}{=============}\PY{l+s}{\PYZbs{}}\PY{l+s}{n}\PY{l+s}{\PYZdq{}}\PY{p}{)}\PY{p}{;}
\PY{+w}{        }\PY{k}{for}\PY{+w}{ }\PY{p}{(}\PY{n}{AuditEntry}\PY{+w}{ }\PY{n}{entry}\PY{+w}{ }\PY{p}{:}\PY{+w}{ }\PY{n}{entries}\PY{p}{)}\PY{+w}{ }\PY{p}{\PYZob{}}
\PY{+w}{            }\PY{n}{report}\PY{p}{.}\PY{n+na}{append}\PY{p}{(}\PY{n}{entry}\PY{p}{.}\PY{n+na}{timestamp}\PY{p}{(}\PY{p}{)}\PY{p}{)}\PY{p}{.}\PY{n+na}{append}\PY{p}{(}\PY{l+s}{\PYZdq{}}\PY{l+s}{ | }\PY{l+s}{\PYZdq{}}\PY{p}{)}
\PY{+w}{                  }\PY{p}{.}\PY{n+na}{append}\PY{p}{(}\PY{n}{entry}\PY{p}{.}\PY{n+na}{user}\PY{p}{(}\PY{p}{)}\PY{p}{)}\PY{p}{.}\PY{n+na}{append}\PY{p}{(}\PY{l+s}{\PYZdq{}}\PY{l+s}{ | }\PY{l+s}{\PYZdq{}}\PY{p}{)}
\PY{+w}{                  }\PY{p}{.}\PY{n+na}{append}\PY{p}{(}\PY{n}{entry}\PY{p}{.}\PY{n+na}{action}\PY{p}{(}\PY{p}{)}\PY{p}{)}\PY{p}{.}\PY{n+na}{append}\PY{p}{(}\PY{l+s+sc}{\PYZsq{}\PYZbs{}n\PYZsq{}}\PY{p}{)}\PY{p}{;}
\PY{+w}{        }\PY{p}{\PYZcb{}}
\PY{+w}{        }\PY{n}{report}\PY{p}{.}\PY{n+na}{append}\PY{p}{(}\PY{l+s}{\PYZdq{}}\PY{l+s}{Total entries: }\PY{l+s}{\PYZdq{}}\PY{p}{)}\PY{p}{.}\PY{n+na}{append}\PY{p}{(}\PY{n}{entries}\PY{p}{.}\PY{n+na}{size}\PY{p}{(}\PY{p}{)}\PY{p}{)}\PY{p}{.}\PY{n+na}{append}\PY{p}{(}\PY{l+s+sc}{\PYZsq{}\PYZbs{}n\PYZsq{}}\PY{p}{)}\PY{p}{;}
\PY{+w}{        }\PY{k}{return}\PY{+w}{ }\PY{n}{report}\PY{p}{.}\PY{n+na}{toString}\PY{p}{(}\PY{p}{)}\PY{p}{;}
\PY{+w}{    }\PY{p}{\PYZcb{}}
\PY{p}{\PYZcb{}}
\end{Verbatim}
\end{codebox}

Note: a single \code{+} expression made of several pieces on one line (\code{a~\allowbreak{}+\allowbreak{}~\allowbreak{}\textquotedbl{}~\allowbreak{}\textquotedbl{}~\allowbreak{}+\allowbreak{}~\allowbreak{}b}) is fine — javac compiles that to one \code{StringBuilder} chain automatically. The bug is specifically repeated concatenation \emph{across loop iterations}, where each iteration starts from an already-large accumulated string.

\end{sidebar}
\subsection[Scenario 8: The Invoice Date Parser]{Scenario 8: The Invoice Date Parser}
\label{h:22:scenario-8-the-invoice-date-parser}
\begin{codebox}
\begin{Verbatim}[commandchars=\\\{\}]
\PY{k+kd}{public}\PY{+w}{ }\PY{k+kd}{class} \PY{n+nc}{InvoiceDateParser}\PY{+w}{ }\PY{p}{\PYZob{}}

\PY{+w}{    }\PY{k+kd}{private}\PY{+w}{ }\PY{k+kd}{static}\PY{+w}{ }\PY{k+kd}{final}\PY{+w}{ }\PY{n}{SimpleDateFormat}\PY{+w}{ }\PY{n}{FORMAT}\PY{+w}{ }\PY{o}{=}\PY{+w}{ }\PY{k}{new}\PY{+w}{ }\PY{n}{SimpleDateFormat}\PY{p}{(}\PY{l+s}{\PYZdq{}}\PY{l+s}{yyyy\PYZhy{}MM\PYZhy{}dd}\PY{l+s}{\PYZdq{}}\PY{p}{)}\PY{p}{;}

\PY{+w}{    }\PY{k+kd}{public}\PY{+w}{ }\PY{n}{Date}\PY{+w}{ }\PY{n+nf}{parse}\PY{p}{(}\PY{n}{String}\PY{+w}{ }\PY{n}{rawDate}\PY{p}{)}\PY{+w}{ }\PY{k+kd}{throws}\PY{+w}{ }\PY{n}{ParseException}\PY{+w}{ }\PY{p}{\PYZob{}}
\PY{+w}{        }\PY{k}{return}\PY{+w}{ }\PY{n}{FORMAT}\PY{p}{.}\PY{n+na}{parse}\PY{p}{(}\PY{n}{rawDate}\PY{p}{)}\PY{p}{;}
\PY{+w}{    }\PY{p}{\PYZcb{}}

\PY{+w}{    }\PY{k+kd}{public}\PY{+w}{ }\PY{n}{String}\PY{+w}{ }\PY{n+nf}{format}\PY{p}{(}\PY{n}{Date}\PY{+w}{ }\PY{n}{date}\PY{p}{)}\PY{+w}{ }\PY{p}{\PYZob{}}
\PY{+w}{        }\PY{k}{return}\PY{+w}{ }\PY{n}{FORMAT}\PY{p}{.}\PY{n+na}{format}\PY{p}{(}\PY{n}{date}\PY{p}{)}\PY{p}{;}
\PY{+w}{    }\PY{p}{\PYZcb{}}
\PY{p}{\PYZcb{}}
\end{Verbatim}
\end{codebox}

\begin{sidebar}{Review}
\begin{itemize}
\item 
\textbf{Blocker — \code{SimpleDateFormat} is not thread-safe, and here it’s a shared \code{static~\allowbreak{}final} field.} Its internal \code{Calendar} field is mutated during both \code{parse()} and \code{format()}. Two threads calling either method concurrently can interleave their mutations of that shared \code{Calendar}, producing wrong dates, garbage results, or a thrown exception — silently and intermittently, which makes it brutal to reproduce in a bug report.

\item 
\textbf{Major — even ignoring the field being \code{static}}, sharing \emph{any single} \code{SimpleDateFormat} instance across threads (static or not) has this problem; the fix is either “don’t share it” or “don’t use \code{SimpleDateFormat}.”

\end{itemize}

Fixed code — use \code{java.\allowbreak{}time}’s \code{DateTimeFormatter}, which is immutable and thread-safe by design (see \hyperref[chap:9]{Chapter~9, \emph{Date, Time \& Localization}}):

\begin{codebox}
\begin{Verbatim}[commandchars=\\\{\}]
\PY{k+kd}{public}\PY{+w}{ }\PY{k+kd}{class} \PY{n+nc}{InvoiceDateParser}\PY{+w}{ }\PY{p}{\PYZob{}}

\PY{+w}{    }\PY{k+kd}{private}\PY{+w}{ }\PY{k+kd}{static}\PY{+w}{ }\PY{k+kd}{final}\PY{+w}{ }\PY{n}{DateTimeFormatter}\PY{+w}{ }\PY{n}{FORMAT}\PY{+w}{ }\PY{o}{=}\PY{+w}{ }\PY{n}{DateTimeFormatter}\PY{p}{.}\PY{n+na}{ofPattern}\PY{p}{(}\PY{l+s}{\PYZdq{}}\PY{l+s}{yyyy\PYZhy{}MM\PYZhy{}dd}\PY{l+s}{\PYZdq{}}\PY{p}{)}\PY{p}{;}

\PY{+w}{    }\PY{k+kd}{public}\PY{+w}{ }\PY{n}{LocalDate}\PY{+w}{ }\PY{n+nf}{parse}\PY{p}{(}\PY{n}{String}\PY{+w}{ }\PY{n}{rawDate}\PY{p}{)}\PY{+w}{ }\PY{p}{\PYZob{}}
\PY{+w}{        }\PY{k}{return}\PY{+w}{ }\PY{n}{LocalDate}\PY{p}{.}\PY{n+na}{parse}\PY{p}{(}\PY{n}{rawDate}\PY{p}{,}\PY{+w}{ }\PY{n}{FORMAT}\PY{p}{)}\PY{p}{;}
\PY{+w}{    }\PY{p}{\PYZcb{}}

\PY{+w}{    }\PY{k+kd}{public}\PY{+w}{ }\PY{n}{String}\PY{+w}{ }\PY{n+nf}{format}\PY{p}{(}\PY{n}{LocalDate}\PY{+w}{ }\PY{n}{date}\PY{p}{)}\PY{+w}{ }\PY{p}{\PYZob{}}
\PY{+w}{        }\PY{k}{return}\PY{+w}{ }\PY{n}{date}\PY{p}{.}\PY{n+na}{format}\PY{p}{(}\PY{n}{FORMAT}\PY{p}{)}\PY{p}{;}
\PY{+w}{    }\PY{p}{\PYZcb{}}
\PY{p}{\PYZcb{}}
\end{Verbatim}
\end{codebox}

If \code{SimpleDateFormat} truly cannot be avoided (some legacy API demands \code{Date}), the safe fallback is a fresh instance per call, or one instance per thread via \code{Thread\allowbreak{}Local\textless{}\allowbreak{}Simple\allowbreak{}Date\allowbreak{}Format\textgreater{}} — never a shared mutable field.

\end{sidebar}
\subsection[Scenario 9: The Translation Cache]{Scenario 9: The Translation Cache}
\label{h:22:scenario-9-the-translation-cache}
\begin{codebox}
\begin{Verbatim}[commandchars=\\\{\}]
\PY{k+kd}{public}\PY{+w}{ }\PY{k+kd}{class} \PY{n+nc}{TranslationService}\PY{+w}{ }\PY{p}{\PYZob{}}

\PY{+w}{    }\PY{k+kd}{private}\PY{+w}{ }\PY{k+kd}{static}\PY{+w}{ }\PY{k+kd}{final}\PY{+w}{ }\PY{n}{Map}\PY{o}{\PYZlt{}}\PY{n}{String}\PY{p}{,}\PY{+w}{ }\PY{n}{String}\PY{o}{\PYZgt{}}\PY{+w}{ }\PY{n}{cache}\PY{+w}{ }\PY{o}{=}\PY{+w}{ }\PY{k}{new}\PY{+w}{ }\PY{n}{HashMap}\PY{o}{\PYZlt{}}\PY{o}{\PYZgt{}}\PY{p}{(}\PY{p}{)}\PY{p}{;}

\PY{+w}{    }\PY{k+kd}{public}\PY{+w}{ }\PY{n}{String}\PY{+w}{ }\PY{n+nf}{translate}\PY{p}{(}\PY{n}{String}\PY{+w}{ }\PY{n}{key}\PY{p}{,}\PY{+w}{ }\PY{n}{String}\PY{+w}{ }\PY{n}{locale}\PY{p}{)}\PY{+w}{ }\PY{p}{\PYZob{}}
\PY{+w}{        }\PY{n}{String}\PY{+w}{ }\PY{n}{cacheKey}\PY{+w}{ }\PY{o}{=}\PY{+w}{ }\PY{n}{key}\PY{+w}{ }\PY{o}{+}\PY{+w}{ }\PY{l+s}{\PYZdq{}}\PY{l+s}{:}\PY{l+s}{\PYZdq{}}\PY{+w}{ }\PY{o}{+}\PY{+w}{ }\PY{n}{locale}\PY{p}{;}
\PY{+w}{        }\PY{k}{return}\PY{+w}{ }\PY{n}{cache}\PY{p}{.}\PY{n+na}{computeIfAbsent}\PY{p}{(}\PY{n}{cacheKey}\PY{p}{,}\PY{+w}{ }\PY{n}{k}\PY{+w}{ }\PY{o}{\PYZhy{}}\PY{o}{\PYZgt{}}\PY{+w}{ }\PY{n}{callTranslationApi}\PY{p}{(}\PY{n}{key}\PY{p}{,}\PY{+w}{ }\PY{n}{locale}\PY{p}{)}\PY{p}{)}\PY{p}{;}
\PY{+w}{    }\PY{p}{\PYZcb{}}

\PY{+w}{    }\PY{k+kd}{private}\PY{+w}{ }\PY{n}{String}\PY{+w}{ }\PY{n+nf}{callTranslationApi}\PY{p}{(}\PY{n}{String}\PY{+w}{ }\PY{n}{key}\PY{p}{,}\PY{+w}{ }\PY{n}{String}\PY{+w}{ }\PY{n}{locale}\PY{p}{)}\PY{+w}{ }\PY{p}{\PYZob{}}
\PY{+w}{        }\PY{c+c1}{// expensive network call}
\PY{+w}{        }\PY{k}{return}\PY{+w}{ }\PY{n}{translationClient}\PY{p}{.}\PY{n+na}{fetch}\PY{p}{(}\PY{n}{key}\PY{p}{,}\PY{+w}{ }\PY{n}{locale}\PY{p}{)}\PY{p}{;}
\PY{+w}{    }\PY{p}{\PYZcb{}}

\PY{+w}{    }\PY{k+kd}{private}\PY{+w}{ }\PY{n}{TranslationClient}\PY{+w}{ }\PY{n}{translationClient}\PY{p}{;}
\PY{p}{\PYZcb{}}
\end{Verbatim}
\end{codebox}

\begin{sidebar}{Review}
\begin{itemize}
\item 
\textbf{Blocker — unbounded cache.} Every distinct \code{(\allowbreak{}key,\allowbreak{}~\allowbreak{}locale)} pair ever requested stays in \code{cache} for the lifetime of the process — there is no eviction, no size cap, no TTL. In a long-running service with many keys/locales (or with any attacker- or user-influenced key), this is a slow, guaranteed memory leak that eventually produces an \code{OutOfMemoryError} days or weeks after deployment, long after the code that caused it has been forgotten.

\item 
\textbf{Major — plain \code{HashMap} shared as a \code{static} field with no synchronization}, read and written from what is presumably a multi-threaded service (translation lookups happening per request). Same category of bug as Scenario 2 — needs a concurrent-safe map at minimum, independent of the eviction problem.

\item 
\textbf{Minor — string-concatenated composite key (\code{key~\allowbreak{}+\allowbreak{}~\allowbreak{}\textquotedbl{}:\allowbreak{}\textquotedbl{}~\allowbreak{}+\allowbreak{}~\allowbreak{}locale})} works but is fragile (a \code{key} containing \code{:} could collide with a different \code{locale} split) and allocates a new \code{String} per lookup; a small record \code{Translation\allowbreak{}Key(\allowbreak{}String~\allowbreak{}key,\allowbreak{}~\allowbreak{}String~\allowbreak{}locale)} is clearer and avoids the delimiter-collision risk.

\end{itemize}

Fixed code — bounded, size- and time-limited cache (e.g., Caffeine), or at minimum a bounded \code{LinkedHashMap}-based LRU, plus a proper composite key:

\begin{codebox}
\begin{Verbatim}[commandchars=\\\{\}]
\PY{k+kd}{public}\PY{+w}{ }\PY{k+kd}{class} \PY{n+nc}{TranslationService}\PY{+w}{ }\PY{p}{\PYZob{}}

\PY{+w}{    }\PY{k+kd}{private}\PY{+w}{ }\PY{k+kd}{record} \PY{n+nc}{TranslationKey}\PY{p}{(}\PY{n}{String}\PY{+w}{ }\PY{n}{key}\PY{p}{,}\PY{+w}{ }\PY{n}{String}\PY{+w}{ }\PY{n}{locale}\PY{p}{)}\PY{+w}{ }\PY{p}{\PYZob{}}\PY{p}{\PYZcb{}}

\PY{+w}{    }\PY{k+kd}{private}\PY{+w}{ }\PY{k+kd}{final}\PY{+w}{ }\PY{n}{Cache}\PY{o}{\PYZlt{}}\PY{n}{TranslationKey}\PY{p}{,}\PY{+w}{ }\PY{n}{String}\PY{o}{\PYZgt{}}\PY{+w}{ }\PY{n}{cache}\PY{+w}{ }\PY{o}{=}\PY{+w}{ }\PY{n}{Caffeine}\PY{p}{.}\PY{n+na}{newBuilder}\PY{p}{(}\PY{p}{)}
\PY{+w}{            }\PY{p}{.}\PY{n+na}{maximumSize}\PY{p}{(}\PY{l+m+mi}{10\PYZus{}000}\PY{p}{)}
\PY{+w}{            }\PY{p}{.}\PY{n+na}{expireAfterWrite}\PY{p}{(}\PY{n}{Duration}\PY{p}{.}\PY{n+na}{ofHours}\PY{p}{(}\PY{l+m+mi}{6}\PY{p}{)}\PY{p}{)}
\PY{+w}{            }\PY{p}{.}\PY{n+na}{build}\PY{p}{(}\PY{p}{)}\PY{p}{;}

\PY{+w}{    }\PY{k+kd}{public}\PY{+w}{ }\PY{n}{String}\PY{+w}{ }\PY{n+nf}{translate}\PY{p}{(}\PY{n}{String}\PY{+w}{ }\PY{n}{key}\PY{p}{,}\PY{+w}{ }\PY{n}{String}\PY{+w}{ }\PY{n}{locale}\PY{p}{)}\PY{+w}{ }\PY{p}{\PYZob{}}
\PY{+w}{        }\PY{k}{return}\PY{+w}{ }\PY{n}{cache}\PY{p}{.}\PY{n+na}{get}\PY{p}{(}\PY{k}{new}\PY{+w}{ }\PY{n}{TranslationKey}\PY{p}{(}\PY{n}{key}\PY{p}{,}\PY{+w}{ }\PY{n}{locale}\PY{p}{)}\PY{p}{,}
\PY{+w}{                }\PY{n}{k}\PY{+w}{ }\PY{o}{\PYZhy{}}\PY{o}{\PYZgt{}}\PY{+w}{ }\PY{n}{callTranslationApi}\PY{p}{(}\PY{n}{k}\PY{p}{.}\PY{n+na}{key}\PY{p}{(}\PY{p}{)}\PY{p}{,}\PY{+w}{ }\PY{n}{k}\PY{p}{.}\PY{n+na}{locale}\PY{p}{(}\PY{p}{)}\PY{p}{)}\PY{p}{)}\PY{p}{;}
\PY{+w}{    }\PY{p}{\PYZcb{}}

\PY{+w}{    }\PY{k+kd}{private}\PY{+w}{ }\PY{n}{String}\PY{+w}{ }\PY{n+nf}{callTranslationApi}\PY{p}{(}\PY{n}{String}\PY{+w}{ }\PY{n}{key}\PY{p}{,}\PY{+w}{ }\PY{n}{String}\PY{+w}{ }\PY{n}{locale}\PY{p}{)}\PY{+w}{ }\PY{p}{\PYZob{}}
\PY{+w}{        }\PY{k}{return}\PY{+w}{ }\PY{n}{translationClient}\PY{p}{.}\PY{n+na}{fetch}\PY{p}{(}\PY{n}{key}\PY{p}{,}\PY{+w}{ }\PY{n}{locale}\PY{p}{)}\PY{p}{;}
\PY{+w}{    }\PY{p}{\PYZcb{}}

\PY{+w}{    }\PY{k+kd}{private}\PY{+w}{ }\PY{n}{TranslationClient}\PY{+w}{ }\PY{n}{translationClient}\PY{p}{;}
\PY{p}{\PYZcb{}}
\end{Verbatim}
\end{codebox}

See \hyperref[h:22:memory--gc-review]{Memory \& GC Review} for the general “unbounded cache” pattern and how to spot it without a library like Caffeine available.

\end{sidebar}
\subsection[Scenario 10: The Duplicate-Free Product List]{Scenario 10: The Duplicate-Free Product List}
\label{h:22:scenario-10-the-duplicate-free-product-list}
\begin{codebox}
\begin{Verbatim}[commandchars=\\\{\}]
\PY{k+kd}{public}\PY{+w}{ }\PY{k+kd}{class} \PY{n+nc}{Product}\PY{+w}{ }\PY{p}{\PYZob{}}
\PY{+w}{    }\PY{k+kd}{private}\PY{+w}{ }\PY{k+kd}{final}\PY{+w}{ }\PY{n}{String}\PY{+w}{ }\PY{n}{sku}\PY{p}{;}
\PY{+w}{    }\PY{k+kd}{private}\PY{+w}{ }\PY{n}{String}\PY{+w}{ }\PY{n}{displayName}\PY{p}{;}

\PY{+w}{    }\PY{k+kd}{public}\PY{+w}{ }\PY{n+nf}{Product}\PY{p}{(}\PY{n}{String}\PY{+w}{ }\PY{n}{sku}\PY{p}{,}\PY{+w}{ }\PY{n}{String}\PY{+w}{ }\PY{n}{displayName}\PY{p}{)}\PY{+w}{ }\PY{p}{\PYZob{}}
\PY{+w}{        }\PY{k}{this}\PY{p}{.}\PY{n+na}{sku}\PY{+w}{ }\PY{o}{=}\PY{+w}{ }\PY{n}{sku}\PY{p}{;}
\PY{+w}{        }\PY{k}{this}\PY{p}{.}\PY{n+na}{displayName}\PY{+w}{ }\PY{o}{=}\PY{+w}{ }\PY{n}{displayName}\PY{p}{;}
\PY{+w}{    }\PY{p}{\PYZcb{}}

\PY{+w}{    }\PY{k+kd}{public}\PY{+w}{ }\PY{k+kt}{void}\PY{+w}{ }\PY{n+nf}{rename}\PY{p}{(}\PY{n}{String}\PY{+w}{ }\PY{n}{newDisplayName}\PY{p}{)}\PY{+w}{ }\PY{p}{\PYZob{}}
\PY{+w}{        }\PY{k}{this}\PY{p}{.}\PY{n+na}{displayName}\PY{+w}{ }\PY{o}{=}\PY{+w}{ }\PY{n}{newDisplayName}\PY{p}{;}
\PY{+w}{    }\PY{p}{\PYZcb{}}

\PY{+w}{    }\PY{n+nd}{@Override}
\PY{+w}{    }\PY{k+kd}{public}\PY{+w}{ }\PY{k+kt}{boolean}\PY{+w}{ }\PY{n+nf}{equals}\PY{p}{(}\PY{n}{Object}\PY{+w}{ }\PY{n}{o}\PY{p}{)}\PY{+w}{ }\PY{p}{\PYZob{}}
\PY{+w}{        }\PY{k}{if}\PY{+w}{ }\PY{p}{(}\PY{o}{!}\PY{p}{(}\PY{n}{o}\PY{+w}{ }\PY{k}{instanceof}\PY{+w}{ }\PY{n}{Product}\PY{+w}{ }\PY{n}{other}\PY{p}{)}\PY{p}{)}\PY{+w}{ }\PY{k}{return}\PY{+w}{ }\PY{k+kc}{false}\PY{p}{;}
\PY{+w}{        }\PY{k}{return}\PY{+w}{ }\PY{n}{sku}\PY{p}{.}\PY{n+na}{equals}\PY{p}{(}\PY{n}{other}\PY{p}{.}\PY{n+na}{sku}\PY{p}{)}\PY{p}{;}
\PY{+w}{    }\PY{p}{\PYZcb{}}

\PY{+w}{    }\PY{k+kd}{public}\PY{+w}{ }\PY{n}{String}\PY{+w}{ }\PY{n+nf}{getSku}\PY{p}{(}\PY{p}{)}\PY{+w}{ }\PY{p}{\PYZob{}}\PY{+w}{ }\PY{k}{return}\PY{+w}{ }\PY{n}{sku}\PY{p}{;}\PY{+w}{ }\PY{p}{\PYZcb{}}
\PY{p}{\PYZcb{}}

\PY{k+kd}{public}\PY{+w}{ }\PY{k+kd}{class} \PY{n+nc}{ProductCatalog}\PY{+w}{ }\PY{p}{\PYZob{}}
\PY{+w}{    }\PY{k+kd}{private}\PY{+w}{ }\PY{k+kd}{final}\PY{+w}{ }\PY{n}{Set}\PY{o}{\PYZlt{}}\PY{n}{Product}\PY{o}{\PYZgt{}}\PY{+w}{ }\PY{n}{products}\PY{+w}{ }\PY{o}{=}\PY{+w}{ }\PY{k}{new}\PY{+w}{ }\PY{n}{HashSet}\PY{o}{\PYZlt{}}\PY{o}{\PYZgt{}}\PY{p}{(}\PY{p}{)}\PY{p}{;}

\PY{+w}{    }\PY{k+kd}{public}\PY{+w}{ }\PY{k+kt}{void}\PY{+w}{ }\PY{n+nf}{add}\PY{p}{(}\PY{n}{Product}\PY{+w}{ }\PY{n}{product}\PY{p}{)}\PY{+w}{ }\PY{p}{\PYZob{}}
\PY{+w}{        }\PY{n}{products}\PY{p}{.}\PY{n+na}{add}\PY{p}{(}\PY{n}{product}\PY{p}{)}\PY{p}{;}
\PY{+w}{    }\PY{p}{\PYZcb{}}

\PY{+w}{    }\PY{k+kd}{public}\PY{+w}{ }\PY{k+kt}{boolean}\PY{+w}{ }\PY{n+nf}{contains}\PY{p}{(}\PY{n}{String}\PY{+w}{ }\PY{n}{sku}\PY{p}{)}\PY{+w}{ }\PY{p}{\PYZob{}}
\PY{+w}{        }\PY{k}{return}\PY{+w}{ }\PY{n}{products}\PY{p}{.}\PY{n+na}{contains}\PY{p}{(}\PY{k}{new}\PY{+w}{ }\PY{n}{Product}\PY{p}{(}\PY{n}{sku}\PY{p}{,}\PY{+w}{ }\PY{l+s}{\PYZdq{}}\PY{l+s}{\PYZdq{}}\PY{p}{)}\PY{p}{)}\PY{p}{;}
\PY{+w}{    }\PY{p}{\PYZcb{}}
\PY{p}{\PYZcb{}}
\end{Verbatim}
\end{codebox}

\begin{sidebar}{Review}
\begin{itemize}
\item 
\textbf{Blocker — \code{equals()} is overridden without overriding \code{hashCode()}.} \code{Product} still uses \code{Object}’s identity-based \code{hashCode()}, so two \code{Product} instances with the same \code{sku} are \code{equals()} but almost certainly land in different \code{HashSet} buckets. \code{contains(\allowbreak{}String~\allowbreak{}sku)} builds a throwaway \code{Product} with a different identity hash than any stored one, so it will essentially always return \code{false} — the exact opposite of what this class exists to do. This is a silent, no-compiler-warning correctness bug: the equals/hashCode contract (“equal objects must have equal hash codes”) is broken.

\item 
\textbf{Major — \code{Product} is mutable (\code{rename}) but is stored in a \code{HashSet}.} Even after fixing \code{hashCode()}, \code{sku} is the only field participating in equality, so mutating \code{displayName} is actually safe here — but this is fragile by accident, not by design; a future field added to \code{equals}/\code{hashCode} that is also mutable would reintroduce Scenario 3’s mutable-key bug.

\item 
\textbf{Nit — \code{contains(\allowbreak{}String~\allowbreak{}sku)} constructs a throwaway \code{Product(\allowbreak{}\textquotedbl{}\textquotedbl{},\allowbreak{}~\allowbreak{}sku)} just to do a lookup} — legal once \code{equals}/\code{hashCode} are fixed, but a \code{Map\textless{}\allowbreak{}String,\allowbreak{}~\allowbreak{}Product\textgreater{}} keyed directly by \code{sku} is more direct and avoids the wasted allocation per lookup.

\end{itemize}

Fixed code — add the missing \code{hashCode()}, and prefer keying directly by the identity field:

\begin{codebox}
\begin{Verbatim}[commandchars=\\\{\}]
\PY{k+kd}{public}\PY{+w}{ }\PY{k+kd}{class} \PY{n+nc}{Product}\PY{+w}{ }\PY{p}{\PYZob{}}
\PY{+w}{    }\PY{k+kd}{private}\PY{+w}{ }\PY{k+kd}{final}\PY{+w}{ }\PY{n}{String}\PY{+w}{ }\PY{n}{sku}\PY{p}{;}
\PY{+w}{    }\PY{k+kd}{private}\PY{+w}{ }\PY{n}{String}\PY{+w}{ }\PY{n}{displayName}\PY{p}{;}

\PY{+w}{    }\PY{k+kd}{public}\PY{+w}{ }\PY{n+nf}{Product}\PY{p}{(}\PY{n}{String}\PY{+w}{ }\PY{n}{sku}\PY{p}{,}\PY{+w}{ }\PY{n}{String}\PY{+w}{ }\PY{n}{displayName}\PY{p}{)}\PY{+w}{ }\PY{p}{\PYZob{}}
\PY{+w}{        }\PY{k}{this}\PY{p}{.}\PY{n+na}{sku}\PY{+w}{ }\PY{o}{=}\PY{+w}{ }\PY{n}{sku}\PY{p}{;}
\PY{+w}{        }\PY{k}{this}\PY{p}{.}\PY{n+na}{displayName}\PY{+w}{ }\PY{o}{=}\PY{+w}{ }\PY{n}{displayName}\PY{p}{;}
\PY{+w}{    }\PY{p}{\PYZcb{}}

\PY{+w}{    }\PY{k+kd}{public}\PY{+w}{ }\PY{k+kt}{void}\PY{+w}{ }\PY{n+nf}{rename}\PY{p}{(}\PY{n}{String}\PY{+w}{ }\PY{n}{newDisplayName}\PY{p}{)}\PY{+w}{ }\PY{p}{\PYZob{}}
\PY{+w}{        }\PY{k}{this}\PY{p}{.}\PY{n+na}{displayName}\PY{+w}{ }\PY{o}{=}\PY{+w}{ }\PY{n}{newDisplayName}\PY{p}{;}
\PY{+w}{    }\PY{p}{\PYZcb{}}

\PY{+w}{    }\PY{n+nd}{@Override}
\PY{+w}{    }\PY{k+kd}{public}\PY{+w}{ }\PY{k+kt}{boolean}\PY{+w}{ }\PY{n+nf}{equals}\PY{p}{(}\PY{n}{Object}\PY{+w}{ }\PY{n}{o}\PY{p}{)}\PY{+w}{ }\PY{p}{\PYZob{}}
\PY{+w}{        }\PY{k}{if}\PY{+w}{ }\PY{p}{(}\PY{o}{!}\PY{p}{(}\PY{n}{o}\PY{+w}{ }\PY{k}{instanceof}\PY{+w}{ }\PY{n}{Product}\PY{+w}{ }\PY{n}{other}\PY{p}{)}\PY{p}{)}\PY{+w}{ }\PY{k}{return}\PY{+w}{ }\PY{k+kc}{false}\PY{p}{;}
\PY{+w}{        }\PY{k}{return}\PY{+w}{ }\PY{n}{sku}\PY{p}{.}\PY{n+na}{equals}\PY{p}{(}\PY{n}{other}\PY{p}{.}\PY{n+na}{sku}\PY{p}{)}\PY{p}{;}
\PY{+w}{    }\PY{p}{\PYZcb{}}

\PY{+w}{    }\PY{n+nd}{@Override}
\PY{+w}{    }\PY{k+kd}{public}\PY{+w}{ }\PY{k+kt}{int}\PY{+w}{ }\PY{n+nf}{hashCode}\PY{p}{(}\PY{p}{)}\PY{+w}{ }\PY{p}{\PYZob{}}
\PY{+w}{        }\PY{k}{return}\PY{+w}{ }\PY{n}{sku}\PY{p}{.}\PY{n+na}{hashCode}\PY{p}{(}\PY{p}{)}\PY{p}{;}
\PY{+w}{    }\PY{p}{\PYZcb{}}

\PY{+w}{    }\PY{k+kd}{public}\PY{+w}{ }\PY{n}{String}\PY{+w}{ }\PY{n+nf}{getSku}\PY{p}{(}\PY{p}{)}\PY{+w}{ }\PY{p}{\PYZob{}}\PY{+w}{ }\PY{k}{return}\PY{+w}{ }\PY{n}{sku}\PY{p}{;}\PY{+w}{ }\PY{p}{\PYZcb{}}
\PY{p}{\PYZcb{}}

\PY{k+kd}{public}\PY{+w}{ }\PY{k+kd}{class} \PY{n+nc}{ProductCatalog}\PY{+w}{ }\PY{p}{\PYZob{}}
\PY{+w}{    }\PY{k+kd}{private}\PY{+w}{ }\PY{k+kd}{final}\PY{+w}{ }\PY{n}{Map}\PY{o}{\PYZlt{}}\PY{n}{String}\PY{p}{,}\PY{+w}{ }\PY{n}{Product}\PY{o}{\PYZgt{}}\PY{+w}{ }\PY{n}{productsBySku}\PY{+w}{ }\PY{o}{=}\PY{+w}{ }\PY{k}{new}\PY{+w}{ }\PY{n}{HashMap}\PY{o}{\PYZlt{}}\PY{o}{\PYZgt{}}\PY{p}{(}\PY{p}{)}\PY{p}{;}

\PY{+w}{    }\PY{k+kd}{public}\PY{+w}{ }\PY{k+kt}{void}\PY{+w}{ }\PY{n+nf}{add}\PY{p}{(}\PY{n}{Product}\PY{+w}{ }\PY{n}{product}\PY{p}{)}\PY{+w}{ }\PY{p}{\PYZob{}}
\PY{+w}{        }\PY{n}{productsBySku}\PY{p}{.}\PY{n+na}{put}\PY{p}{(}\PY{n}{product}\PY{p}{.}\PY{n+na}{getSku}\PY{p}{(}\PY{p}{)}\PY{p}{,}\PY{+w}{ }\PY{n}{product}\PY{p}{)}\PY{p}{;}
\PY{+w}{    }\PY{p}{\PYZcb{}}

\PY{+w}{    }\PY{k+kd}{public}\PY{+w}{ }\PY{k+kt}{boolean}\PY{+w}{ }\PY{n+nf}{contains}\PY{p}{(}\PY{n}{String}\PY{+w}{ }\PY{n}{sku}\PY{p}{)}\PY{+w}{ }\PY{p}{\PYZob{}}
\PY{+w}{        }\PY{k}{return}\PY{+w}{ }\PY{n}{productsBySku}\PY{p}{.}\PY{n+na}{containsKey}\PY{p}{(}\PY{n}{sku}\PY{p}{)}\PY{p}{;}
\PY{+w}{    }\PY{p}{\PYZcb{}}
\PY{p}{\PYZcb{}}
\end{Verbatim}
\end{codebox}

\end{sidebar}
\subsection[Scenario 11: The Per-User Session Cache]{Scenario 11: The Per-User Session Cache}
\label{h:22:scenario-11-the-per-user-session-cache}
\begin{codebox}
\begin{Verbatim}[commandchars=\\\{\}]
\PY{k+kd}{public}\PY{+w}{ }\PY{k+kd}{class} \PY{n+nc}{SessionCache}\PY{+w}{ }\PY{p}{\PYZob{}}

\PY{+w}{    }\PY{k+kd}{private}\PY{+w}{ }\PY{k+kd}{final}\PY{+w}{ }\PY{n}{Map}\PY{o}{\PYZlt{}}\PY{n}{String}\PY{p}{,}\PY{+w}{ }\PY{n}{List}\PY{o}{\PYZlt{}}\PY{n}{String}\PY{o}{\PYZgt{}}\PY{o}{\PYZgt{}}\PY{+w}{ }\PY{n}{activityByUser}\PY{+w}{ }\PY{o}{=}\PY{+w}{ }\PY{k}{new}\PY{+w}{ }\PY{n}{ConcurrentHashMap}\PY{o}{\PYZlt{}}\PY{o}{\PYZgt{}}\PY{p}{(}\PY{p}{)}\PY{p}{;}

\PY{+w}{    }\PY{k+kd}{public}\PY{+w}{ }\PY{k+kt}{void}\PY{+w}{ }\PY{n+nf}{recordActivity}\PY{p}{(}\PY{n}{String}\PY{+w}{ }\PY{n}{userId}\PY{p}{,}\PY{+w}{ }\PY{n}{String}\PY{+w}{ }\PY{n}{activity}\PY{p}{)}\PY{+w}{ }\PY{p}{\PYZob{}}
\PY{+w}{        }\PY{n}{activityByUser}\PY{p}{.}\PY{n+na}{computeIfAbsent}\PY{p}{(}\PY{n}{userId}\PY{p}{,}\PY{+w}{ }\PY{n}{id}\PY{+w}{ }\PY{o}{\PYZhy{}}\PY{o}{\PYZgt{}}\PY{+w}{ }\PY{p}{\PYZob{}}
\PY{+w}{            }\PY{n}{List}\PY{o}{\PYZlt{}}\PY{n}{String}\PY{o}{\PYZgt{}}\PY{+w}{ }\PY{n}{history}\PY{+w}{ }\PY{o}{=}\PY{+w}{ }\PY{n}{loadHistoryFromDb}\PY{p}{(}\PY{n}{id}\PY{p}{)}\PY{p}{;}\PY{+w}{ }\PY{c+c1}{// side effect: DB call inside the lambda}
\PY{+w}{            }\PY{n}{history}\PY{p}{.}\PY{n+na}{add}\PY{p}{(}\PY{n}{activity}\PY{p}{)}\PY{p}{;}
\PY{+w}{            }\PY{k}{return}\PY{+w}{ }\PY{n}{history}\PY{p}{;}
\PY{+w}{        }\PY{p}{\PYZcb{}}\PY{p}{)}\PY{p}{;}
\PY{+w}{    }\PY{p}{\PYZcb{}}

\PY{+w}{    }\PY{k+kd}{public}\PY{+w}{ }\PY{k+kt}{void}\PY{+w}{ }\PY{n+nf}{addRelatedUser}\PY{p}{(}\PY{n}{String}\PY{+w}{ }\PY{n}{userId}\PY{p}{,}\PY{+w}{ }\PY{n}{String}\PY{+w}{ }\PY{n}{relatedUserId}\PY{p}{)}\PY{+w}{ }\PY{p}{\PYZob{}}
\PY{+w}{        }\PY{n}{activityByUser}\PY{p}{.}\PY{n+na}{computeIfAbsent}\PY{p}{(}\PY{n}{userId}\PY{p}{,}\PY{+w}{ }\PY{n}{id}\PY{+w}{ }\PY{o}{\PYZhy{}}\PY{o}{\PYZgt{}}\PY{+w}{ }\PY{k}{new}\PY{+w}{ }\PY{n}{ArrayList}\PY{o}{\PYZlt{}}\PY{o}{\PYZgt{}}\PY{p}{(}\PY{p}{)}\PY{p}{)}
\PY{+w}{                }\PY{p}{.}\PY{n+na}{add}\PY{p}{(}\PY{l+s}{\PYZdq{}}\PY{l+s}{related:}\PY{l+s}{\PYZdq{}}\PY{+w}{ }\PY{o}{+}\PY{+w}{ }\PY{n}{relatedUserId}\PY{p}{)}\PY{p}{;}
\PY{+w}{        }\PY{n}{activityByUser}\PY{p}{.}\PY{n+na}{computeIfAbsent}\PY{p}{(}\PY{n}{relatedUserId}\PY{p}{,}\PY{+w}{ }\PY{n}{id}\PY{+w}{ }\PY{o}{\PYZhy{}}\PY{o}{\PYZgt{}}
\PY{+w}{                }\PY{n}{activityByUser}\PY{p}{.}\PY{n+na}{get}\PY{p}{(}\PY{n}{userId}\PY{p}{)}\PY{p}{)}\PY{p}{;}\PY{+w}{ }\PY{c+c1}{// reuse the same list reference}
\PY{+w}{    }\PY{p}{\PYZcb{}}

\PY{+w}{    }\PY{k+kd}{private}\PY{+w}{ }\PY{n}{List}\PY{o}{\PYZlt{}}\PY{n}{String}\PY{o}{\PYZgt{}}\PY{+w}{ }\PY{n+nf}{loadHistoryFromDb}\PY{p}{(}\PY{n}{String}\PY{+w}{ }\PY{n}{userId}\PY{p}{)}\PY{+w}{ }\PY{p}{\PYZob{}}\PY{+w}{ }\PY{c+cm}{/* ... */}\PY{+w}{ }\PY{k}{return}\PY{+w}{ }\PY{k}{new}\PY{+w}{ }\PY{n}{ArrayList}\PY{o}{\PYZlt{}}\PY{o}{\PYZgt{}}\PY{p}{(}\PY{p}{)}\PY{p}{;}\PY{+w}{ }\PY{p}{\PYZcb{}}
\PY{p}{\PYZcb{}}
\end{Verbatim}
\end{codebox}

\begin{sidebar}{Review}
\begin{itemize}
\item 
\textbf{Blocker — \code{recordActivity} only adds the new activity on the \emph{first} call for a user.} \code{computeIfAbsent}’s mapping function only runs when the key is absent; on every subsequent call for a user who’s already in the map, the lambda is skipped entirely and \code{activity} is silently dropped. This looks like it works in a quick manual test (first call always “works”) and then quietly loses data for every repeat visitor — the most dangerous kind of bug because the obvious test case passes.

\item 
\textbf{Blocker — \code{loadHistoryFromDb} (a blocking, potentially slow DB call) runs \emph{inside} the \code{computeIfAbsent} mapping function}, which \code{ConcurrentHashMap} executes while holding an internal lock on that key’s bin. If it’s slow, or if it happens to (directly or transitively) call back into the same map for the same key, you can get very long stalls for every other thread touching that bin, or in some JDK versions an actual deadlock/\code{IllegalStateException} (recursive update). \code{computeIfAbsent}’s mapping function must be fast and must never touch the same map.

\item 
\textbf{Major — \code{addRelatedUser} aliases two users to the \emph{same mutable \code{List} instance}.} After this call, \code{activity\allowbreak{}By\allowbreak{}User.\allowbreak{}get(\allowbreak{}user\allowbreak{}Id)} and \code{activity\allowbreak{}By\allowbreak{}User.\allowbreak{}get(\allowbreak{}related\allowbreak{}User\allowbreak{}Id)} are the identical list object; recording activity for one silently mutates the “history” of the other. This is almost certainly not intended and is very hard to spot from the call site.

\end{itemize}

Fixed code — separate the DB load from the update, keep the mapping function cheap and side-effect-free with respect to the map itself, and never alias mutable collections between keys:

\begin{codebox}
\begin{Verbatim}[commandchars=\\\{\}]
\PY{k+kd}{public}\PY{+w}{ }\PY{k+kd}{class} \PY{n+nc}{SessionCache}\PY{+w}{ }\PY{p}{\PYZob{}}

\PY{+w}{    }\PY{k+kd}{private}\PY{+w}{ }\PY{k+kd}{final}\PY{+w}{ }\PY{n}{Map}\PY{o}{\PYZlt{}}\PY{n}{String}\PY{p}{,}\PY{+w}{ }\PY{n}{List}\PY{o}{\PYZlt{}}\PY{n}{String}\PY{o}{\PYZgt{}}\PY{o}{\PYZgt{}}\PY{+w}{ }\PY{n}{activityByUser}\PY{+w}{ }\PY{o}{=}\PY{+w}{ }\PY{k}{new}\PY{+w}{ }\PY{n}{ConcurrentHashMap}\PY{o}{\PYZlt{}}\PY{o}{\PYZgt{}}\PY{p}{(}\PY{p}{)}\PY{p}{;}

\PY{+w}{    }\PY{k+kd}{public}\PY{+w}{ }\PY{k+kt}{void}\PY{+w}{ }\PY{n+nf}{recordActivity}\PY{p}{(}\PY{n}{String}\PY{+w}{ }\PY{n}{userId}\PY{p}{,}\PY{+w}{ }\PY{n}{String}\PY{+w}{ }\PY{n}{activity}\PY{p}{)}\PY{+w}{ }\PY{p}{\PYZob{}}
\PY{+w}{        }\PY{n}{List}\PY{o}{\PYZlt{}}\PY{n}{String}\PY{o}{\PYZgt{}}\PY{+w}{ }\PY{n}{history}\PY{+w}{ }\PY{o}{=}\PY{+w}{ }\PY{n}{activityByUser}\PY{p}{.}\PY{n+na}{computeIfAbsent}\PY{p}{(}\PY{n}{userId}\PY{p}{,}
\PY{+w}{                }\PY{n}{id}\PY{+w}{ }\PY{o}{\PYZhy{}}\PY{o}{\PYZgt{}}\PY{+w}{ }\PY{k}{new}\PY{+w}{ }\PY{n}{CopyOnWriteArrayList}\PY{o}{\PYZlt{}}\PY{o}{\PYZgt{}}\PY{p}{(}\PY{n}{loadHistoryFromDb}\PY{p}{(}\PY{n}{id}\PY{p}{)}\PY{p}{)}\PY{p}{)}\PY{p}{;}\PY{+w}{ }\PY{c+c1}{// fast to construct, DB call happens once, outside any recursive risk}
\PY{+w}{        }\PY{n}{history}\PY{p}{.}\PY{n+na}{add}\PY{p}{(}\PY{n}{activity}\PY{p}{)}\PY{p}{;}\PY{+w}{ }\PY{c+c1}{// always runs, regardless of whether this was the first call}
\PY{+w}{    }\PY{p}{\PYZcb{}}

\PY{+w}{    }\PY{k+kd}{public}\PY{+w}{ }\PY{k+kt}{void}\PY{+w}{ }\PY{n+nf}{addRelatedUser}\PY{p}{(}\PY{n}{String}\PY{+w}{ }\PY{n}{userId}\PY{p}{,}\PY{+w}{ }\PY{n}{String}\PY{+w}{ }\PY{n}{relatedUserId}\PY{p}{)}\PY{+w}{ }\PY{p}{\PYZob{}}
\PY{+w}{        }\PY{n}{activityByUser}\PY{p}{.}\PY{n+na}{computeIfAbsent}\PY{p}{(}\PY{n}{userId}\PY{p}{,}\PY{+w}{ }\PY{n}{id}\PY{+w}{ }\PY{o}{\PYZhy{}}\PY{o}{\PYZgt{}}\PY{+w}{ }\PY{k}{new}\PY{+w}{ }\PY{n}{CopyOnWriteArrayList}\PY{o}{\PYZlt{}}\PY{o}{\PYZgt{}}\PY{p}{(}\PY{p}{)}\PY{p}{)}
\PY{+w}{                }\PY{p}{.}\PY{n+na}{add}\PY{p}{(}\PY{l+s}{\PYZdq{}}\PY{l+s}{related:}\PY{l+s}{\PYZdq{}}\PY{+w}{ }\PY{o}{+}\PY{+w}{ }\PY{n}{relatedUserId}\PY{p}{)}\PY{p}{;}
\PY{+w}{        }\PY{n}{activityByUser}\PY{p}{.}\PY{n+na}{computeIfAbsent}\PY{p}{(}\PY{n}{relatedUserId}\PY{p}{,}\PY{+w}{ }\PY{n}{id}\PY{+w}{ }\PY{o}{\PYZhy{}}\PY{o}{\PYZgt{}}\PY{+w}{ }\PY{k}{new}\PY{+w}{ }\PY{n}{CopyOnWriteArrayList}\PY{o}{\PYZlt{}}\PY{o}{\PYZgt{}}\PY{p}{(}\PY{p}{)}\PY{p}{)}
\PY{+w}{                }\PY{p}{.}\PY{n+na}{add}\PY{p}{(}\PY{l+s}{\PYZdq{}}\PY{l+s}{related\PYZhy{}to:}\PY{l+s}{\PYZdq{}}\PY{+w}{ }\PY{o}{+}\PY{+w}{ }\PY{n}{userId}\PY{p}{)}\PY{p}{;}\PY{+w}{ }\PY{c+c1}{// own list, not a shared reference}
\PY{+w}{    }\PY{p}{\PYZcb{}}

\PY{+w}{    }\PY{k+kd}{private}\PY{+w}{ }\PY{n}{List}\PY{o}{\PYZlt{}}\PY{n}{String}\PY{o}{\PYZgt{}}\PY{+w}{ }\PY{n+nf}{loadHistoryFromDb}\PY{p}{(}\PY{n}{String}\PY{+w}{ }\PY{n}{userId}\PY{p}{)}\PY{+w}{ }\PY{p}{\PYZob{}}\PY{+w}{ }\PY{c+cm}{/* ... */}\PY{+w}{ }\PY{k}{return}\PY{+w}{ }\PY{k}{new}\PY{+w}{ }\PY{n}{ArrayList}\PY{o}{\PYZlt{}}\PY{o}{\PYZgt{}}\PY{p}{(}\PY{p}{)}\PY{p}{;}\PY{+w}{ }\PY{p}{\PYZcb{}}
\PY{p}{\PYZcb{}}
\end{Verbatim}
\end{codebox}

\end{sidebar}
\subsection[Scenario 12: The Invalid Order Filter]{Scenario 12: The Invalid Order Filter}
\label{h:22:scenario-12-the-invalid-order-filter}
\begin{codebox}
\begin{Verbatim}[commandchars=\\\{\}]
\PY{k+kd}{public}\PY{+w}{ }\PY{k+kd}{class} \PY{n+nc}{OrderValidationReport}\PY{+w}{ }\PY{p}{\PYZob{}}

\PY{+w}{    }\PY{k+kd}{public}\PY{+w}{ }\PY{n}{List}\PY{o}{\PYZlt{}}\PY{n}{String}\PY{o}{\PYZgt{}}\PY{+w}{ }\PY{n+nf}{findInvalidOrderIds}\PY{p}{(}\PY{n}{List}\PY{o}{\PYZlt{}}\PY{n}{Order}\PY{o}{\PYZgt{}}\PY{+w}{ }\PY{n}{orders}\PY{p}{)}\PY{+w}{ }\PY{p}{\PYZob{}}
\PY{+w}{        }\PY{n}{List}\PY{o}{\PYZlt{}}\PY{n}{String}\PY{o}{\PYZgt{}}\PY{+w}{ }\PY{n}{invalidIds}\PY{+w}{ }\PY{o}{=}\PY{+w}{ }\PY{k}{new}\PY{+w}{ }\PY{n}{ArrayList}\PY{o}{\PYZlt{}}\PY{o}{\PYZgt{}}\PY{p}{(}\PY{p}{)}\PY{p}{;}
\PY{+w}{        }\PY{k+kt}{int}\PY{o}{[}\PY{o}{]}\PY{+w}{ }\PY{n}{invalidCount}\PY{+w}{ }\PY{o}{=}\PY{+w}{ }\PY{p}{\PYZob{}}\PY{l+m+mi}{0}\PY{p}{\PYZcb{}}\PY{p}{;}

\PY{+w}{        }\PY{n}{orders}\PY{p}{.}\PY{n+na}{parallelStream}\PY{p}{(}\PY{p}{)}\PY{p}{.}\PY{n+na}{forEach}\PY{p}{(}\PY{n}{order}\PY{+w}{ }\PY{o}{\PYZhy{}}\PY{o}{\PYZgt{}}\PY{+w}{ }\PY{p}{\PYZob{}}
\PY{+w}{            }\PY{k}{if}\PY{+w}{ }\PY{p}{(}\PY{o}{!}\PY{n}{isValid}\PY{p}{(}\PY{n}{order}\PY{p}{)}\PY{p}{)}\PY{+w}{ }\PY{p}{\PYZob{}}
\PY{+w}{                }\PY{n}{invalidIds}\PY{p}{.}\PY{n+na}{add}\PY{p}{(}\PY{n}{order}\PY{p}{.}\PY{n+na}{id}\PY{p}{(}\PY{p}{)}\PY{p}{)}\PY{p}{;}
\PY{+w}{                }\PY{n}{invalidCount}\PY{o}{[}\PY{l+m+mi}{0}\PY{o}{]}\PY{o}{+}\PY{o}{+}\PY{p}{;}
\PY{+w}{            }\PY{p}{\PYZcb{}}
\PY{+w}{        }\PY{p}{\PYZcb{}}\PY{p}{)}\PY{p}{;}

\PY{+w}{        }\PY{n}{System}\PY{p}{.}\PY{n+na}{out}\PY{p}{.}\PY{n+na}{println}\PY{p}{(}\PY{l+s}{\PYZdq{}}\PY{l+s}{Found }\PY{l+s}{\PYZdq{}}\PY{+w}{ }\PY{o}{+}\PY{+w}{ }\PY{n}{invalidCount}\PY{o}{[}\PY{l+m+mi}{0}\PY{o}{]}\PY{+w}{ }\PY{o}{+}\PY{+w}{ }\PY{l+s}{\PYZdq{}}\PY{l+s}{ invalid orders}\PY{l+s}{\PYZdq{}}\PY{p}{)}\PY{p}{;}
\PY{+w}{        }\PY{k}{return}\PY{+w}{ }\PY{n}{invalidIds}\PY{p}{;}
\PY{+w}{    }\PY{p}{\PYZcb{}}

\PY{+w}{    }\PY{k+kd}{private}\PY{+w}{ }\PY{k+kt}{boolean}\PY{+w}{ }\PY{n+nf}{isValid}\PY{p}{(}\PY{n}{Order}\PY{+w}{ }\PY{n}{order}\PY{p}{)}\PY{+w}{ }\PY{p}{\PYZob{}}
\PY{+w}{        }\PY{k}{return}\PY{+w}{ }\PY{n}{order}\PY{p}{.}\PY{n+na}{total}\PY{p}{(}\PY{p}{)}\PY{p}{.}\PY{n+na}{compareTo}\PY{p}{(}\PY{n}{BigDecimal}\PY{p}{.}\PY{n+na}{ZERO}\PY{p}{)}\PY{+w}{ }\PY{o}{\PYZgt{}}\PY{o}{=}\PY{+w}{ }\PY{l+m+mi}{0}\PY{+w}{ }\PY{o}{\PYZam{}}\PY{o}{\PYZam{}}\PY{+w}{ }\PY{n}{order}\PY{p}{.}\PY{n+na}{customerId}\PY{p}{(}\PY{p}{)}\PY{+w}{ }\PY{o}{!}\PY{o}{=}\PY{+w}{ }\PY{k+kc}{null}\PY{p}{;}
\PY{+w}{    }\PY{p}{\PYZcb{}}
\PY{p}{\PYZcb{}}
\end{Verbatim}
\end{codebox}

\begin{sidebar}{Review}
\begin{itemize}
\item 
\textbf{Blocker — \code{invalidIds.\allowbreak{}add(...)} from a \code{parallel\allowbreak{}Stream().\allowbreak{}for\allowbreak{}Each}.} \code{ArrayList} is not thread-safe, and \code{parallelStream()} runs the lambda on multiple threads from the common \code{ForkJoinPool}. Concurrent, unsynchronized \code{add()} calls on an \code{ArrayList} can throw \code{Array\allowbreak{}Index\allowbreak{}Out\allowbreak{}Of\allowbreak{}Bounds\allowbreak{}Exception}, silently drop elements, or corrupt the list’s internal array — a data race hiding behind a one-line stream call that looks innocuous.

\item 
\textbf{Major — \code{invalidCount[\allowbreak{}0]++} is also a race}, for the same reason: the array-boxing trick to mutate a captured variable from a lambda doesn’t make the increment atomic. Multiple threads can read-increment-write the same stale value and lose counts.

\item 
\textbf{Major — reaching for \code{parallelStream()} at all here is questionable} (\hyperref[h:22:performance-optimization]{Performance Optimization} covers this pattern): \code{isValid} is a trivial, non-blocking check, so the overhead of splitting work across the shared \code{ForkJoinPool} easily outweighs any benefit, and it puts load on a JVM-wide shared pool that other unrelated parallel streams also depend on.

\item 
\textbf{Nit — mixing a raw \code{System.\allowbreak{}out.\allowbreak{}println} for what looks like log-worthy operational information} instead of a logger.

\end{itemize}

Fixed code — collect via a proper stream \code{Collector} (thread-safe by construction) and drop the unnecessary parallelism:

\begin{codebox}
\begin{Verbatim}[commandchars=\\\{\}]
\PY{k+kd}{public}\PY{+w}{ }\PY{k+kd}{class} \PY{n+nc}{OrderValidationReport}\PY{+w}{ }\PY{p}{\PYZob{}}
\PY{+w}{    }\PY{k+kd}{private}\PY{+w}{ }\PY{k+kd}{static}\PY{+w}{ }\PY{k+kd}{final}\PY{+w}{ }\PY{n}{Logger}\PY{+w}{ }\PY{n}{log}\PY{+w}{ }\PY{o}{=}\PY{+w}{ }\PY{n}{LoggerFactory}\PY{p}{.}\PY{n+na}{getLogger}\PY{p}{(}\PY{n}{OrderValidationReport}\PY{p}{.}\PY{n+na}{class}\PY{p}{)}\PY{p}{;}

\PY{+w}{    }\PY{k+kd}{public}\PY{+w}{ }\PY{n}{List}\PY{o}{\PYZlt{}}\PY{n}{String}\PY{o}{\PYZgt{}}\PY{+w}{ }\PY{n+nf}{findInvalidOrderIds}\PY{p}{(}\PY{n}{List}\PY{o}{\PYZlt{}}\PY{n}{Order}\PY{o}{\PYZgt{}}\PY{+w}{ }\PY{n}{orders}\PY{p}{)}\PY{+w}{ }\PY{p}{\PYZob{}}
\PY{+w}{        }\PY{n}{List}\PY{o}{\PYZlt{}}\PY{n}{String}\PY{o}{\PYZgt{}}\PY{+w}{ }\PY{n}{invalidIds}\PY{+w}{ }\PY{o}{=}\PY{+w}{ }\PY{n}{orders}\PY{p}{.}\PY{n+na}{stream}\PY{p}{(}\PY{p}{)}
\PY{+w}{                }\PY{p}{.}\PY{n+na}{filter}\PY{p}{(}\PY{n}{order}\PY{+w}{ }\PY{o}{\PYZhy{}}\PY{o}{\PYZgt{}}\PY{+w}{ }\PY{o}{!}\PY{n}{isValid}\PY{p}{(}\PY{n}{order}\PY{p}{)}\PY{p}{)}
\PY{+w}{                }\PY{p}{.}\PY{n+na}{map}\PY{p}{(}\PY{n}{Order}\PY{p}{:}\PY{p}{:}\PY{n}{id}\PY{p}{)}
\PY{+w}{                }\PY{p}{.}\PY{n+na}{toList}\PY{p}{(}\PY{p}{)}\PY{p}{;}

\PY{+w}{        }\PY{n}{log}\PY{p}{.}\PY{n+na}{info}\PY{p}{(}\PY{l+s}{\PYZdq{}}\PY{l+s}{Found \PYZob{}\PYZcb{} invalid orders}\PY{l+s}{\PYZdq{}}\PY{p}{,}\PY{+w}{ }\PY{n}{invalidIds}\PY{p}{.}\PY{n+na}{size}\PY{p}{(}\PY{p}{)}\PY{p}{)}\PY{p}{;}
\PY{+w}{        }\PY{k}{return}\PY{+w}{ }\PY{n}{invalidIds}\PY{p}{;}
\PY{+w}{    }\PY{p}{\PYZcb{}}

\PY{+w}{    }\PY{k+kd}{private}\PY{+w}{ }\PY{k+kt}{boolean}\PY{+w}{ }\PY{n+nf}{isValid}\PY{p}{(}\PY{n}{Order}\PY{+w}{ }\PY{n}{order}\PY{p}{)}\PY{+w}{ }\PY{p}{\PYZob{}}
\PY{+w}{        }\PY{k}{return}\PY{+w}{ }\PY{n}{order}\PY{p}{.}\PY{n+na}{total}\PY{p}{(}\PY{p}{)}\PY{p}{.}\PY{n+na}{compareTo}\PY{p}{(}\PY{n}{BigDecimal}\PY{p}{.}\PY{n+na}{ZERO}\PY{p}{)}\PY{+w}{ }\PY{o}{\PYZgt{}}\PY{o}{=}\PY{+w}{ }\PY{l+m+mi}{0}\PY{+w}{ }\PY{o}{\PYZam{}}\PY{o}{\PYZam{}}\PY{+w}{ }\PY{n}{order}\PY{p}{.}\PY{n+na}{customerId}\PY{p}{(}\PY{p}{)}\PY{+w}{ }\PY{o}{!}\PY{o}{=}\PY{+w}{ }\PY{k+kc}{null}\PY{p}{;}
\PY{+w}{    }\PY{p}{\PYZcb{}}
\PY{p}{\PYZcb{}}
\end{Verbatim}
\end{codebox}

If this really needs to run in parallel because \code{isValid} does expensive, CPU-bound, side-effect-free work over a very large list, \code{parallelStream()} is fine — but then \code{.\allowbreak{}filter(...).\allowbreak{}map(...).\allowbreak{}to\allowbreak{}List()} (a proper collector) is still the right way to gather results, never a shared mutable list touched from inside \code{forEach}.

\end{sidebar}
\subsection[Scenario 13: The Batch Email Sender]{Scenario 13: The Batch Email Sender}
\label{h:22:scenario-13-the-batch-email-sender}
\begin{codebox}
\begin{Verbatim}[commandchars=\\\{\}]
\PY{k+kd}{public}\PY{+w}{ }\PY{k+kd}{class} \PY{n+nc}{BatchEmailSender}\PY{+w}{ }\PY{p}{\PYZob{}}

\PY{+w}{    }\PY{k+kd}{public}\PY{+w}{ }\PY{k+kt}{void}\PY{+w}{ }\PY{n+nf}{sendBatch}\PY{p}{(}\PY{n}{List}\PY{o}{\PYZlt{}}\PY{n}{Email}\PY{o}{\PYZgt{}}\PY{+w}{ }\PY{n}{emails}\PY{p}{)}\PY{+w}{ }\PY{p}{\PYZob{}}
\PY{+w}{        }\PY{n}{ExecutorService}\PY{+w}{ }\PY{n}{executor}\PY{+w}{ }\PY{o}{=}\PY{+w}{ }\PY{n}{Executors}\PY{p}{.}\PY{n+na}{newFixedThreadPool}\PY{p}{(}\PY{l+m+mi}{10}\PY{p}{)}\PY{p}{;}
\PY{+w}{        }\PY{k}{for}\PY{+w}{ }\PY{p}{(}\PY{n}{Email}\PY{+w}{ }\PY{n}{email}\PY{+w}{ }\PY{p}{:}\PY{+w}{ }\PY{n}{emails}\PY{p}{)}\PY{+w}{ }\PY{p}{\PYZob{}}
\PY{+w}{            }\PY{n}{executor}\PY{p}{.}\PY{n+na}{submit}\PY{p}{(}\PY{p}{(}\PY{p}{)}\PY{+w}{ }\PY{o}{\PYZhy{}}\PY{o}{\PYZgt{}}\PY{+w}{ }\PY{n}{smtpClient}\PY{p}{.}\PY{n+na}{send}\PY{p}{(}\PY{n}{email}\PY{p}{)}\PY{p}{)}\PY{p}{;}
\PY{+w}{        }\PY{p}{\PYZcb{}}
\PY{+w}{    }\PY{p}{\PYZcb{}}

\PY{+w}{    }\PY{k+kd}{private}\PY{+w}{ }\PY{n}{SmtpClient}\PY{+w}{ }\PY{n}{smtpClient}\PY{p}{;}
\PY{p}{\PYZcb{}}
\end{Verbatim}
\end{codebox}

\begin{sidebar}{Review}
\begin{itemize}
\item 
\textbf{Blocker — the executor is never shut down.} A brand-new \code{ExecutorService} (with 10 live platform threads) is created on every call to \code{sendBatch} and simply abandoned once the method returns. \code{ExecutorService}s hold real threads that are not garbage-collected just because the reference goes out of scope — each call leaks 10 threads that live forever (or until the JVM exits), and a service that calls \code{sendBatch} repeatedly will exhaust available threads/memory and eventually fail with \code{Out\allowbreak{}Of\allowbreak{}Memory\allowbreak{}Error:\allowbreak{}~\allowbreak{}unable~\allowbreak{}to~\allowbreak{}create~\allowbreak{}new~\allowbreak{}native~\allowbreak{}thread}.

\item 
\textbf{Major — no way to know when the batch actually finished}, and no error handling for individual sends: \code{smtpClient.\allowbreak{}send(\allowbreak{}email)} inside the submitted task can throw, and since nothing calls \code{Future.\allowbreak{}get()} on the submitted tasks, that exception is silently swallowed by the executor (visible only if you dig through logs for an “uncaught exception in thread pool” trace, if that’s even configured).

\item 
\textbf{Minor — a fresh fixed pool per batch is also wasteful even if shut down correctly}; a single shared, appropriately-sized executor reused across calls (owned by the surrounding component’s lifecycle, not created ad hoc per method call) is both cheaper and easier to reason about.

\end{itemize}

Fixed code — own the executor as a long-lived field with a proper lifecycle (\code{shutdown}), and track task outcomes:

\begin{codebox}
\begin{Verbatim}[commandchars=\\\{\}]
\PY{k+kd}{public}\PY{+w}{ }\PY{k+kd}{class} \PY{n+nc}{BatchEmailSender}\PY{+w}{ }\PY{k+kd}{implements}\PY{+w}{ }\PY{n}{AutoCloseable}\PY{+w}{ }\PY{p}{\PYZob{}}
\PY{+w}{    }\PY{k+kd}{private}\PY{+w}{ }\PY{k+kd}{static}\PY{+w}{ }\PY{k+kd}{final}\PY{+w}{ }\PY{n}{Logger}\PY{+w}{ }\PY{n}{log}\PY{+w}{ }\PY{o}{=}\PY{+w}{ }\PY{n}{LoggerFactory}\PY{p}{.}\PY{n+na}{getLogger}\PY{p}{(}\PY{n}{BatchEmailSender}\PY{p}{.}\PY{n+na}{class}\PY{p}{)}\PY{p}{;}
\PY{+w}{    }\PY{k+kd}{private}\PY{+w}{ }\PY{k+kd}{final}\PY{+w}{ }\PY{n}{ExecutorService}\PY{+w}{ }\PY{n}{executor}\PY{+w}{ }\PY{o}{=}\PY{+w}{ }\PY{n}{Executors}\PY{p}{.}\PY{n+na}{newFixedThreadPool}\PY{p}{(}\PY{l+m+mi}{10}\PY{p}{)}\PY{p}{;}
\PY{+w}{    }\PY{k+kd}{private}\PY{+w}{ }\PY{k+kd}{final}\PY{+w}{ }\PY{n}{SmtpClient}\PY{+w}{ }\PY{n}{smtpClient}\PY{p}{;}

\PY{+w}{    }\PY{k+kd}{public}\PY{+w}{ }\PY{n+nf}{BatchEmailSender}\PY{p}{(}\PY{n}{SmtpClient}\PY{+w}{ }\PY{n}{smtpClient}\PY{p}{)}\PY{+w}{ }\PY{p}{\PYZob{}}
\PY{+w}{        }\PY{k}{this}\PY{p}{.}\PY{n+na}{smtpClient}\PY{+w}{ }\PY{o}{=}\PY{+w}{ }\PY{n}{smtpClient}\PY{p}{;}
\PY{+w}{    }\PY{p}{\PYZcb{}}

\PY{+w}{    }\PY{k+kd}{public}\PY{+w}{ }\PY{k+kt}{void}\PY{+w}{ }\PY{n+nf}{sendBatch}\PY{p}{(}\PY{n}{List}\PY{o}{\PYZlt{}}\PY{n}{Email}\PY{o}{\PYZgt{}}\PY{+w}{ }\PY{n}{emails}\PY{p}{)}\PY{+w}{ }\PY{p}{\PYZob{}}
\PY{+w}{        }\PY{n}{List}\PY{o}{\PYZlt{}}\PY{n}{Future}\PY{o}{\PYZlt{}}\PY{o}{?}\PY{o}{\PYZgt{}}\PY{o}{\PYZgt{}}\PY{+w}{ }\PY{n}{futures}\PY{+w}{ }\PY{o}{=}\PY{+w}{ }\PY{n}{emails}\PY{p}{.}\PY{n+na}{stream}\PY{p}{(}\PY{p}{)}
\PY{+w}{                }\PY{p}{.}\PY{n+na}{map}\PY{p}{(}\PY{n}{email}\PY{+w}{ }\PY{o}{\PYZhy{}}\PY{o}{\PYZgt{}}\PY{+w}{ }\PY{n}{executor}\PY{p}{.}\PY{n+na}{submit}\PY{p}{(}\PY{p}{(}\PY{p}{)}\PY{+w}{ }\PY{o}{\PYZhy{}}\PY{o}{\PYZgt{}}\PY{+w}{ }\PY{n}{smtpClient}\PY{p}{.}\PY{n+na}{send}\PY{p}{(}\PY{n}{email}\PY{p}{)}\PY{p}{)}\PY{p}{)}
\PY{+w}{                }\PY{p}{.}\PY{n+na}{toList}\PY{p}{(}\PY{p}{)}\PY{p}{;}

\PY{+w}{        }\PY{k}{for}\PY{+w}{ }\PY{p}{(}\PY{n}{Future}\PY{o}{\PYZlt{}}\PY{o}{?}\PY{o}{\PYZgt{}}\PY{+w}{ }\PY{n}{future}\PY{+w}{ }\PY{p}{:}\PY{+w}{ }\PY{n}{futures}\PY{p}{)}\PY{+w}{ }\PY{p}{\PYZob{}}
\PY{+w}{            }\PY{k}{try}\PY{+w}{ }\PY{p}{\PYZob{}}
\PY{+w}{                }\PY{n}{future}\PY{p}{.}\PY{n+na}{get}\PY{p}{(}\PY{p}{)}\PY{p}{;}
\PY{+w}{            }\PY{p}{\PYZcb{}}\PY{+w}{ }\PY{k}{catch}\PY{+w}{ }\PY{p}{(}\PY{n}{ExecutionException}\PY{+w}{ }\PY{n}{e}\PY{p}{)}\PY{+w}{ }\PY{p}{\PYZob{}}
\PY{+w}{                }\PY{n}{log}\PY{p}{.}\PY{n+na}{warn}\PY{p}{(}\PY{l+s}{\PYZdq{}}\PY{l+s}{Failed to send an email in batch}\PY{l+s}{\PYZdq{}}\PY{p}{,}\PY{+w}{ }\PY{n}{e}\PY{p}{.}\PY{n+na}{getCause}\PY{p}{(}\PY{p}{)}\PY{p}{)}\PY{p}{;}
\PY{+w}{            }\PY{p}{\PYZcb{}}\PY{+w}{ }\PY{k}{catch}\PY{+w}{ }\PY{p}{(}\PY{n}{InterruptedException}\PY{+w}{ }\PY{n}{e}\PY{p}{)}\PY{+w}{ }\PY{p}{\PYZob{}}
\PY{+w}{                }\PY{n}{Thread}\PY{p}{.}\PY{n+na}{currentThread}\PY{p}{(}\PY{p}{)}\PY{p}{.}\PY{n+na}{interrupt}\PY{p}{(}\PY{p}{)}\PY{p}{;}
\PY{+w}{                }\PY{k}{throw}\PY{+w}{ }\PY{k}{new}\PY{+w}{ }\PY{n}{RuntimeException}\PY{p}{(}\PY{l+s}{\PYZdq{}}\PY{l+s}{Interrupted while sending batch}\PY{l+s}{\PYZdq{}}\PY{p}{,}\PY{+w}{ }\PY{n}{e}\PY{p}{)}\PY{p}{;}
\PY{+w}{            }\PY{p}{\PYZcb{}}
\PY{+w}{        }\PY{p}{\PYZcb{}}
\PY{+w}{    }\PY{p}{\PYZcb{}}

\PY{+w}{    }\PY{n+nd}{@Override}
\PY{+w}{    }\PY{k+kd}{public}\PY{+w}{ }\PY{k+kt}{void}\PY{+w}{ }\PY{n+nf}{close}\PY{p}{(}\PY{p}{)}\PY{+w}{ }\PY{p}{\PYZob{}}
\PY{+w}{        }\PY{n}{executor}\PY{p}{.}\PY{n+na}{shutdown}\PY{p}{(}\PY{p}{)}\PY{p}{;}\PY{+w}{ }\PY{c+c1}{// let in\PYZhy{}flight sends finish; see 14\PYZhy{}concurrency\PYZhy{}advanced.md for shutdown patterns}
\PY{+w}{    }\PY{p}{\PYZcb{}}
\PY{p}{\PYZcb{}}
\end{Verbatim}
\end{codebox}

\end{sidebar}
\subsection[Scenario 14: The Import Job Runner]{Scenario 14: The Import Job Runner}
\label{h:22:scenario-14-the-import-job-runner}
\begin{codebox}
\begin{Verbatim}[commandchars=\\\{\}]
\PY{k+kd}{public}\PY{+w}{ }\PY{k+kd}{class} \PY{n+nc}{ImportJobRunner}\PY{+w}{ }\PY{p}{\PYZob{}}

\PY{+w}{    }\PY{k+kd}{public}\PY{+w}{ }\PY{n}{ImportResult}\PY{+w}{ }\PY{n+nf}{runImport}\PY{p}{(}\PY{n}{String}\PY{+w}{ }\PY{n}{filePath}\PY{p}{)}\PY{+w}{ }\PY{p}{\PYZob{}}
\PY{+w}{        }\PY{k}{try}\PY{+w}{ }\PY{p}{\PYZob{}}
\PY{+w}{            }\PY{n}{List}\PY{o}{\PYZlt{}}\PY{n}{String}\PY{o}{\PYZgt{}}\PY{+w}{ }\PY{n}{lines}\PY{+w}{ }\PY{o}{=}\PY{+w}{ }\PY{n}{Files}\PY{p}{.}\PY{n+na}{readAllLines}\PY{p}{(}\PY{n}{Path}\PY{p}{.}\PY{n+na}{of}\PY{p}{(}\PY{n}{filePath}\PY{p}{)}\PY{p}{)}\PY{p}{;}
\PY{+w}{            }\PY{n}{List}\PY{o}{\PYZlt{}}\PY{n}{Record}\PY{o}{\PYZgt{}}\PY{+w}{ }\PY{n}{records}\PY{+w}{ }\PY{o}{=}\PY{+w}{ }\PY{k}{new}\PY{+w}{ }\PY{n}{ArrayList}\PY{o}{\PYZlt{}}\PY{o}{\PYZgt{}}\PY{p}{(}\PY{p}{)}\PY{p}{;}
\PY{+w}{            }\PY{k}{for}\PY{+w}{ }\PY{p}{(}\PY{n}{String}\PY{+w}{ }\PY{n}{line}\PY{+w}{ }\PY{p}{:}\PY{+w}{ }\PY{n}{lines}\PY{p}{)}\PY{+w}{ }\PY{p}{\PYZob{}}
\PY{+w}{                }\PY{n}{Record}\PY{+w}{ }\PY{n}{record}\PY{+w}{ }\PY{o}{=}\PY{+w}{ }\PY{n}{parseLine}\PY{p}{(}\PY{n}{line}\PY{p}{)}\PY{p}{;}
\PY{+w}{                }\PY{n}{validate}\PY{p}{(}\PY{n}{record}\PY{p}{)}\PY{p}{;}
\PY{+w}{                }\PY{n}{records}\PY{p}{.}\PY{n+na}{add}\PY{p}{(}\PY{n}{record}\PY{p}{)}\PY{p}{;}
\PY{+w}{            }\PY{p}{\PYZcb{}}
\PY{+w}{            }\PY{n}{repository}\PY{p}{.}\PY{n+na}{saveAll}\PY{p}{(}\PY{n}{records}\PY{p}{)}\PY{p}{;}
\PY{+w}{            }\PY{n}{notifyStakeholders}\PY{p}{(}\PY{n}{records}\PY{p}{.}\PY{n+na}{size}\PY{p}{(}\PY{p}{)}\PY{p}{)}\PY{p}{;}
\PY{+w}{            }\PY{k}{return}\PY{+w}{ }\PY{n}{ImportResult}\PY{p}{.}\PY{n+na}{success}\PY{p}{(}\PY{n}{records}\PY{p}{.}\PY{n+na}{size}\PY{p}{(}\PY{p}{)}\PY{p}{)}\PY{p}{;}
\PY{+w}{        }\PY{p}{\PYZcb{}}\PY{+w}{ }\PY{k}{catch}\PY{+w}{ }\PY{p}{(}\PY{n}{Exception}\PY{+w}{ }\PY{n}{e}\PY{p}{)}\PY{+w}{ }\PY{p}{\PYZob{}}
\PY{+w}{            }\PY{k}{return}\PY{+w}{ }\PY{n}{ImportResult}\PY{p}{.}\PY{n+na}{failure}\PY{p}{(}\PY{l+s}{\PYZdq{}}\PY{l+s}{Import failed}\PY{l+s}{\PYZdq{}}\PY{p}{)}\PY{p}{;}
\PY{+w}{        }\PY{p}{\PYZcb{}}
\PY{+w}{    }\PY{p}{\PYZcb{}}

\PY{+w}{    }\PY{k+kd}{private}\PY{+w}{ }\PY{n}{Record}\PY{+w}{ }\PY{n+nf}{parseLine}\PY{p}{(}\PY{n}{String}\PY{+w}{ }\PY{n}{line}\PY{p}{)}\PY{+w}{ }\PY{p}{\PYZob{}}\PY{+w}{ }\PY{c+cm}{/* throws NumberFormatException on bad rows */}\PY{+w}{ }\PY{k}{return}\PY{+w}{ }\PY{k+kc}{null}\PY{p}{;}\PY{+w}{ }\PY{p}{\PYZcb{}}
\PY{+w}{    }\PY{k+kd}{private}\PY{+w}{ }\PY{k+kt}{void}\PY{+w}{ }\PY{n+nf}{validate}\PY{p}{(}\PY{n}{Record}\PY{+w}{ }\PY{n}{record}\PY{p}{)}\PY{+w}{ }\PY{p}{\PYZob{}}\PY{+w}{ }\PY{c+cm}{/* throws ValidationException */}\PY{+w}{ }\PY{p}{\PYZcb{}}
\PY{+w}{    }\PY{k+kd}{private}\PY{+w}{ }\PY{k+kt}{void}\PY{+w}{ }\PY{n+nf}{notifyStakeholders}\PY{p}{(}\PY{k+kt}{int}\PY{+w}{ }\PY{n}{count}\PY{p}{)}\PY{+w}{ }\PY{p}{\PYZob{}}\PY{+w}{ }\PY{c+cm}{/* sends an email, can throw MessagingException */}\PY{+w}{ }\PY{p}{\PYZcb{}}
\PY{p}{\PYZcb{}}
\end{Verbatim}
\end{codebox}

\begin{sidebar}{Review}
\begin{itemize}
\item 
\textbf{Major — one giant \code{try} wraps five completely different failure domains}: file I/O (\code{IOException}), per-line parsing (\code{NumberFormatException}), business validation (\code{ValidationException}), the DB save (\code{repository.\allowbreak{}saveAll} — any runtime exception it throws), and notification (\code{MessagingException}). All of them collapse into the same generic \code{\textquotedbl{}Import~\allowbreak{}failed\textquotedbl{}} message, so a caller (or a human on call) cannot tell “the file didn’t exist” apart from “row 4,502 had a wrong data type” apart from “the DB write itself failed after everything else succeeded” apart from “the import fully succeeded but the confirmation email failed to send.” Those require completely different responses.

\item 
\textbf{Major — a failure in \code{notifyStakeholders} (just sending an email) reports the entire import as failed}, even though \code{repository.\allowbreak{}save\allowbreak{}All(\allowbreak{}records)} already committed. The caller will likely re-run the import, potentially double-importing records, because of a problem in an unrelated, non-critical side effect.

\item 
\textbf{Minor — no logging at all inside the \code{catch}}, generic or otherwise — whatever the real exception was is discarded completely, not even preserved for later diagnosis.

\end{itemize}

Fixed code — scope each \code{try} to the failure domain that’s actually being handled, and don’t let a non-critical side effect (notification) mark an otherwise-successful core operation as failed:

\begin{codebox}
\begin{Verbatim}[commandchars=\\\{\}]
\PY{k+kd}{public}\PY{+w}{ }\PY{k+kd}{class} \PY{n+nc}{ImportJobRunner}\PY{+w}{ }\PY{p}{\PYZob{}}
\PY{+w}{    }\PY{k+kd}{private}\PY{+w}{ }\PY{k+kd}{static}\PY{+w}{ }\PY{k+kd}{final}\PY{+w}{ }\PY{n}{Logger}\PY{+w}{ }\PY{n}{log}\PY{+w}{ }\PY{o}{=}\PY{+w}{ }\PY{n}{LoggerFactory}\PY{p}{.}\PY{n+na}{getLogger}\PY{p}{(}\PY{n}{ImportJobRunner}\PY{p}{.}\PY{n+na}{class}\PY{p}{)}\PY{p}{;}

\PY{+w}{    }\PY{k+kd}{public}\PY{+w}{ }\PY{n}{ImportResult}\PY{+w}{ }\PY{n+nf}{runImport}\PY{p}{(}\PY{n}{String}\PY{+w}{ }\PY{n}{filePath}\PY{p}{)}\PY{+w}{ }\PY{p}{\PYZob{}}
\PY{+w}{        }\PY{n}{List}\PY{o}{\PYZlt{}}\PY{n}{String}\PY{o}{\PYZgt{}}\PY{+w}{ }\PY{n}{lines}\PY{p}{;}
\PY{+w}{        }\PY{k}{try}\PY{+w}{ }\PY{p}{\PYZob{}}
\PY{+w}{            }\PY{n}{lines}\PY{+w}{ }\PY{o}{=}\PY{+w}{ }\PY{n}{Files}\PY{p}{.}\PY{n+na}{readAllLines}\PY{p}{(}\PY{n}{Path}\PY{p}{.}\PY{n+na}{of}\PY{p}{(}\PY{n}{filePath}\PY{p}{)}\PY{p}{)}\PY{p}{;}
\PY{+w}{        }\PY{p}{\PYZcb{}}\PY{+w}{ }\PY{k}{catch}\PY{+w}{ }\PY{p}{(}\PY{n}{IOException}\PY{+w}{ }\PY{n}{e}\PY{p}{)}\PY{+w}{ }\PY{p}{\PYZob{}}
\PY{+w}{            }\PY{n}{log}\PY{p}{.}\PY{n+na}{error}\PY{p}{(}\PY{l+s}{\PYZdq{}}\PY{l+s}{Could not read import file \PYZob{}\PYZcb{}}\PY{l+s}{\PYZdq{}}\PY{p}{,}\PY{+w}{ }\PY{n}{filePath}\PY{p}{,}\PY{+w}{ }\PY{n}{e}\PY{p}{)}\PY{p}{;}
\PY{+w}{            }\PY{k}{return}\PY{+w}{ }\PY{n}{ImportResult}\PY{p}{.}\PY{n+na}{failure}\PY{p}{(}\PY{l+s}{\PYZdq{}}\PY{l+s}{Import file could not be read: }\PY{l+s}{\PYZdq{}}\PY{+w}{ }\PY{o}{+}\PY{+w}{ }\PY{n}{e}\PY{p}{.}\PY{n+na}{getMessage}\PY{p}{(}\PY{p}{)}\PY{p}{)}\PY{p}{;}
\PY{+w}{        }\PY{p}{\PYZcb{}}

\PY{+w}{        }\PY{n}{List}\PY{o}{\PYZlt{}}\PY{n}{Record}\PY{o}{\PYZgt{}}\PY{+w}{ }\PY{n}{records}\PY{+w}{ }\PY{o}{=}\PY{+w}{ }\PY{k}{new}\PY{+w}{ }\PY{n}{ArrayList}\PY{o}{\PYZlt{}}\PY{o}{\PYZgt{}}\PY{p}{(}\PY{p}{)}\PY{p}{;}
\PY{+w}{        }\PY{k}{for}\PY{+w}{ }\PY{p}{(}\PY{k+kt}{int}\PY{+w}{ }\PY{n}{i}\PY{+w}{ }\PY{o}{=}\PY{+w}{ }\PY{l+m+mi}{0}\PY{p}{;}\PY{+w}{ }\PY{n}{i}\PY{+w}{ }\PY{o}{\PYZlt{}}\PY{+w}{ }\PY{n}{lines}\PY{p}{.}\PY{n+na}{size}\PY{p}{(}\PY{p}{)}\PY{p}{;}\PY{+w}{ }\PY{n}{i}\PY{o}{+}\PY{o}{+}\PY{p}{)}\PY{+w}{ }\PY{p}{\PYZob{}}
\PY{+w}{            }\PY{k}{try}\PY{+w}{ }\PY{p}{\PYZob{}}
\PY{+w}{                }\PY{n}{Record}\PY{+w}{ }\PY{n}{record}\PY{+w}{ }\PY{o}{=}\PY{+w}{ }\PY{n}{parseLine}\PY{p}{(}\PY{n}{lines}\PY{p}{.}\PY{n+na}{get}\PY{p}{(}\PY{n}{i}\PY{p}{)}\PY{p}{)}\PY{p}{;}
\PY{+w}{                }\PY{n}{validate}\PY{p}{(}\PY{n}{record}\PY{p}{)}\PY{p}{;}
\PY{+w}{                }\PY{n}{records}\PY{p}{.}\PY{n+na}{add}\PY{p}{(}\PY{n}{record}\PY{p}{)}\PY{p}{;}
\PY{+w}{            }\PY{p}{\PYZcb{}}\PY{+w}{ }\PY{k}{catch}\PY{+w}{ }\PY{p}{(}\PY{n}{NumberFormatException}\PY{+w}{ }\PY{o}{|}\PY{+w}{ }\PY{n}{ValidationException}\PY{+w}{ }\PY{n}{e}\PY{p}{)}\PY{+w}{ }\PY{p}{\PYZob{}}
\PY{+w}{                }\PY{n}{log}\PY{p}{.}\PY{n+na}{error}\PY{p}{(}\PY{l+s}{\PYZdq{}}\PY{l+s}{Bad row \PYZob{}\PYZcb{} in \PYZob{}\PYZcb{}: \PYZob{}\PYZcb{}}\PY{l+s}{\PYZdq{}}\PY{p}{,}\PY{+w}{ }\PY{n}{i}\PY{+w}{ }\PY{o}{+}\PY{+w}{ }\PY{l+m+mi}{1}\PY{p}{,}\PY{+w}{ }\PY{n}{filePath}\PY{p}{,}\PY{+w}{ }\PY{n}{e}\PY{p}{.}\PY{n+na}{getMessage}\PY{p}{(}\PY{p}{)}\PY{p}{)}\PY{p}{;}
\PY{+w}{                }\PY{k}{return}\PY{+w}{ }\PY{n}{ImportResult}\PY{p}{.}\PY{n+na}{failure}\PY{p}{(}\PY{l+s}{\PYZdq{}}\PY{l+s}{Row }\PY{l+s}{\PYZdq{}}\PY{+w}{ }\PY{o}{+}\PY{+w}{ }\PY{p}{(}\PY{n}{i}\PY{+w}{ }\PY{o}{+}\PY{+w}{ }\PY{l+m+mi}{1}\PY{p}{)}\PY{+w}{ }\PY{o}{+}\PY{+w}{ }\PY{l+s}{\PYZdq{}}\PY{l+s}{ is invalid: }\PY{l+s}{\PYZdq{}}\PY{+w}{ }\PY{o}{+}\PY{+w}{ }\PY{n}{e}\PY{p}{.}\PY{n+na}{getMessage}\PY{p}{(}\PY{p}{)}\PY{p}{)}\PY{p}{;}
\PY{+w}{            }\PY{p}{\PYZcb{}}
\PY{+w}{        }\PY{p}{\PYZcb{}}

\PY{+w}{        }\PY{n}{repository}\PY{p}{.}\PY{n+na}{saveAll}\PY{p}{(}\PY{n}{records}\PY{p}{)}\PY{p}{;}\PY{+w}{ }\PY{c+c1}{// let a DB failure propagate as\PYZhy{}is — it\PYZsq{}s a real defect, not an expected outcome}

\PY{+w}{        }\PY{k}{try}\PY{+w}{ }\PY{p}{\PYZob{}}
\PY{+w}{            }\PY{n}{notifyStakeholders}\PY{p}{(}\PY{n}{records}\PY{p}{.}\PY{n+na}{size}\PY{p}{(}\PY{p}{)}\PY{p}{)}\PY{p}{;}
\PY{+w}{        }\PY{p}{\PYZcb{}}\PY{+w}{ }\PY{k}{catch}\PY{+w}{ }\PY{p}{(}\PY{n}{MessagingException}\PY{+w}{ }\PY{n}{e}\PY{p}{)}\PY{+w}{ }\PY{p}{\PYZob{}}
\PY{+w}{            }\PY{c+c1}{// Import already succeeded; a failed notification shouldn\PYZsq{}t undo that fact.}
\PY{+w}{            }\PY{n}{log}\PY{p}{.}\PY{n+na}{warn}\PY{p}{(}\PY{l+s}{\PYZdq{}}\PY{l+s}{Import of \PYZob{}\PYZcb{} records succeeded but notification failed}\PY{l+s}{\PYZdq{}}\PY{p}{,}\PY{+w}{ }\PY{n}{records}\PY{p}{.}\PY{n+na}{size}\PY{p}{(}\PY{p}{)}\PY{p}{,}\PY{+w}{ }\PY{n}{e}\PY{p}{)}\PY{p}{;}
\PY{+w}{        }\PY{p}{\PYZcb{}}

\PY{+w}{        }\PY{k}{return}\PY{+w}{ }\PY{n}{ImportResult}\PY{p}{.}\PY{n+na}{success}\PY{p}{(}\PY{n}{records}\PY{p}{.}\PY{n+na}{size}\PY{p}{(}\PY{p}{)}\PY{p}{)}\PY{p}{;}
\PY{+w}{    }\PY{p}{\PYZcb{}}

\PY{+w}{    }\PY{k+kd}{private}\PY{+w}{ }\PY{n}{Record}\PY{+w}{ }\PY{n+nf}{parseLine}\PY{p}{(}\PY{n}{String}\PY{+w}{ }\PY{n}{line}\PY{p}{)}\PY{+w}{ }\PY{p}{\PYZob{}}\PY{+w}{ }\PY{k}{return}\PY{+w}{ }\PY{k+kc}{null}\PY{p}{;}\PY{+w}{ }\PY{p}{\PYZcb{}}
\PY{+w}{    }\PY{k+kd}{private}\PY{+w}{ }\PY{k+kt}{void}\PY{+w}{ }\PY{n+nf}{validate}\PY{p}{(}\PY{n}{Record}\PY{+w}{ }\PY{n}{record}\PY{p}{)}\PY{+w}{ }\PY{p}{\PYZob{}}\PY{+w}{ }\PY{p}{\PYZcb{}}
\PY{+w}{    }\PY{k+kd}{private}\PY{+w}{ }\PY{k+kt}{void}\PY{+w}{ }\PY{n+nf}{notifyStakeholders}\PY{p}{(}\PY{k+kt}{int}\PY{+w}{ }\PY{n}{count}\PY{p}{)}\PY{+w}{ }\PY{k+kd}{throws}\PY{+w}{ }\PY{n}{MessagingException}\PY{+w}{ }\PY{p}{\PYZob{}}\PY{+w}{ }\PY{p}{\PYZcb{}}
\PY{p}{\PYZcb{}}
\end{Verbatim}
\end{codebox}

\end{sidebar}
\subsection[Scenario 15: The Live Price Ticker]{Scenario 15: The Live Price Ticker}
\label{h:22:scenario-15-the-live-price-ticker}
\begin{codebox}
\begin{Verbatim}[commandchars=\\\{\}]
\PY{k+kd}{public}\PY{+w}{ }\PY{k+kd}{class} \PY{n+nc}{LivePriceTicker}\PY{+w}{ }\PY{p}{\PYZob{}}

\PY{+w}{    }\PY{k+kd}{private}\PY{+w}{ }\PY{k+kd}{final}\PY{+w}{ }\PY{n}{List}\PY{o}{\PYZlt{}}\PY{n}{PriceListener}\PY{o}{\PYZgt{}}\PY{+w}{ }\PY{n}{listeners}\PY{+w}{ }\PY{o}{=}\PY{+w}{ }\PY{k}{new}\PY{+w}{ }\PY{n}{ArrayList}\PY{o}{\PYZlt{}}\PY{o}{\PYZgt{}}\PY{p}{(}\PY{p}{)}\PY{p}{;}
\PY{+w}{    }\PY{k+kd}{private}\PY{+w}{ }\PY{k+kd}{volatile}\PY{+w}{ }\PY{k+kt}{double}\PY{+w}{ }\PY{n}{lastPrice}\PY{p}{;}

\PY{+w}{    }\PY{k+kd}{public}\PY{+w}{ }\PY{n+nf}{LivePriceTicker}\PY{p}{(}\PY{n}{PriceFeed}\PY{+w}{ }\PY{n}{feed}\PY{p}{)}\PY{+w}{ }\PY{p}{\PYZob{}}
\PY{+w}{        }\PY{n}{feed}\PY{p}{.}\PY{n+na}{subscribe}\PY{p}{(}\PY{k}{this}\PY{p}{:}\PY{p}{:}\PY{n}{onPriceUpdate}\PY{p}{)}\PY{p}{;}\PY{+w}{ }\PY{c+c1}{// registers \PYZsq{}this\PYZsq{} before construction finishes}
\PY{+w}{        }\PY{k}{this}\PY{p}{.}\PY{n+na}{startBackgroundRefresh}\PY{p}{(}\PY{p}{)}\PY{p}{;}
\PY{+w}{    }\PY{p}{\PYZcb{}}

\PY{+w}{    }\PY{k+kd}{private}\PY{+w}{ }\PY{k+kt}{void}\PY{+w}{ }\PY{n+nf}{startBackgroundRefresh}\PY{p}{(}\PY{p}{)}\PY{+w}{ }\PY{p}{\PYZob{}}
\PY{+w}{        }\PY{n}{Thread}\PY{+w}{ }\PY{n}{refresher}\PY{+w}{ }\PY{o}{=}\PY{+w}{ }\PY{k}{new}\PY{+w}{ }\PY{n}{Thread}\PY{p}{(}\PY{p}{(}\PY{p}{)}\PY{+w}{ }\PY{o}{\PYZhy{}}\PY{o}{\PYZgt{}}\PY{+w}{ }\PY{p}{\PYZob{}}
\PY{+w}{            }\PY{k}{while}\PY{+w}{ }\PY{p}{(}\PY{k+kc}{true}\PY{p}{)}\PY{+w}{ }\PY{p}{\PYZob{}}
\PY{+w}{                }\PY{n}{refreshFromFeed}\PY{p}{(}\PY{p}{)}\PY{p}{;}
\PY{+w}{                }\PY{n}{sleepQuietly}\PY{p}{(}\PY{p}{)}\PY{p}{;}
\PY{+w}{            }\PY{p}{\PYZcb{}}
\PY{+w}{        }\PY{p}{\PYZcb{}}\PY{p}{)}\PY{p}{;}
\PY{+w}{        }\PY{n}{refresher}\PY{p}{.}\PY{n+na}{start}\PY{p}{(}\PY{p}{)}\PY{p}{;}\PY{+w}{ }\PY{c+c1}{// thread started from inside the constructor}
\PY{+w}{    }\PY{p}{\PYZcb{}}

\PY{+w}{    }\PY{k+kd}{public}\PY{+w}{ }\PY{k+kt}{void}\PY{+w}{ }\PY{n+nf}{addListener}\PY{p}{(}\PY{n}{PriceListener}\PY{+w}{ }\PY{n}{listener}\PY{p}{)}\PY{+w}{ }\PY{p}{\PYZob{}}
\PY{+w}{        }\PY{n}{listeners}\PY{p}{.}\PY{n+na}{add}\PY{p}{(}\PY{n}{listener}\PY{p}{)}\PY{p}{;}
\PY{+w}{    }\PY{p}{\PYZcb{}}

\PY{+w}{    }\PY{k+kd}{private}\PY{+w}{ }\PY{k+kt}{void}\PY{+w}{ }\PY{n+nf}{onPriceUpdate}\PY{p}{(}\PY{k+kt}{double}\PY{+w}{ }\PY{n}{price}\PY{p}{)}\PY{+w}{ }\PY{p}{\PYZob{}}
\PY{+w}{        }\PY{k}{this}\PY{p}{.}\PY{n+na}{lastPrice}\PY{+w}{ }\PY{o}{=}\PY{+w}{ }\PY{n}{price}\PY{p}{;}
\PY{+w}{        }\PY{k}{for}\PY{+w}{ }\PY{p}{(}\PY{n}{PriceListener}\PY{+w}{ }\PY{n}{listener}\PY{+w}{ }\PY{p}{:}\PY{+w}{ }\PY{n}{listeners}\PY{p}{)}\PY{+w}{ }\PY{p}{\PYZob{}}
\PY{+w}{            }\PY{n}{listener}\PY{p}{.}\PY{n+na}{onPrice}\PY{p}{(}\PY{n}{price}\PY{p}{)}\PY{p}{;}
\PY{+w}{        }\PY{p}{\PYZcb{}}
\PY{+w}{    }\PY{p}{\PYZcb{}}

\PY{+w}{    }\PY{k+kd}{private}\PY{+w}{ }\PY{k+kt}{void}\PY{+w}{ }\PY{n+nf}{refreshFromFeed}\PY{p}{(}\PY{p}{)}\PY{+w}{ }\PY{p}{\PYZob{}}\PY{+w}{ }\PY{c+cm}{/* ... */}\PY{+w}{ }\PY{p}{\PYZcb{}}
\PY{+w}{    }\PY{k+kd}{private}\PY{+w}{ }\PY{k+kt}{void}\PY{+w}{ }\PY{n+nf}{sleepQuietly}\PY{p}{(}\PY{p}{)}\PY{+w}{ }\PY{p}{\PYZob{}}\PY{+w}{ }\PY{k}{try}\PY{+w}{ }\PY{p}{\PYZob{}}\PY{+w}{ }\PY{n}{Thread}\PY{p}{.}\PY{n+na}{sleep}\PY{p}{(}\PY{l+m+mi}{1000}\PY{p}{)}\PY{p}{;}\PY{+w}{ }\PY{p}{\PYZcb{}}\PY{+w}{ }\PY{k}{catch}\PY{+w}{ }\PY{p}{(}\PY{n}{InterruptedException}\PY{+w}{ }\PY{n}{e}\PY{p}{)}\PY{+w}{ }\PY{p}{\PYZob{}}\PY{+w}{ }\PY{n}{Thread}\PY{p}{.}\PY{n+na}{currentThread}\PY{p}{(}\PY{p}{)}\PY{p}{.}\PY{n+na}{interrupt}\PY{p}{(}\PY{p}{)}\PY{p}{;}\PY{+w}{ }\PY{p}{\PYZcb{}}\PY{+w}{ }\PY{p}{\PYZcb{}}
\PY{p}{\PYZcb{}}
\end{Verbatim}
\end{codebox}

\begin{sidebar}{Review}
\begin{itemize}
\item 
\textbf{Blocker — \code{this} escapes the constructor twice, before construction has finished.} \code{feed.\allowbreak{}subscribe(\allowbreak{}this::\allowbreak{}on\allowbreak{}Price\allowbreak{}Update)} hands out a reference to the object (via the method reference, which captures \code{this}) to another object (\code{feed}) \emph{while the constructor is still running}. If \code{feed} calls back synchronously, or another thread reaches \code{this} through \code{feed} before the constructor returns, it can observe a half-initialized \code{LivePriceTicker} — for example, \code{listeners} might still be \code{null} if this call happened before the field initializer ran, or fields set later in the constructor simply won’t be visible yet. This is the “leaking \code{this}” pitfall, and it’s especially dangerous combined with…

\item 
\textbf{Blocker — starting a background thread from inside the constructor.} \code{refresher.\allowbreak{}start()} hands the \emph{same partially-constructed} object to a new thread of execution, with no memory barrier guaranteeing the new thread sees a fully-initialized object — the Java Memory Model does not guarantee that a background thread started during construction sees writes made later in that same constructor, unless those fields are \code{final} or otherwise safely published. This is a subtle, hard-to-reproduce race that depends on JIT/CPU reordering and timing.

\item 
\textbf{Major — the refresher thread runs \code{while~\allowbreak{}(\allowbreak{}true)} with no way to stop it}, and it’s never tracked anywhere (no field holding the \code{Thread}). There is no way to shut this down cleanly — same lifecycle problem as Scenario 13’s abandoned executor, but worse, because there isn’t even a reference to call \code{interrupt()} on later.

\item 
\textbf{Minor — \code{listeners} (an \code{ArrayList}) is read and written without synchronization}, and \code{onPriceUpdate} (which iterates it) can be invoked concurrently with \code{addListener} (which mutates it) from different threads (the feed’s callback thread vs. whatever thread calls \code{addListener}) — another \code{Concurrent\allowbreak{}Modification\allowbreak{}Exception}/corruption risk layered on top of the construction issue.

\end{itemize}

Fixed code — finish construction fully before publishing \code{this} anywhere, use a static factory to start the background work only after the object is safely built, keep a handle to the thread for shutdown, and use a concurrent-safe listener list:

\begin{codebox}
\begin{Verbatim}[commandchars=\\\{\}]
\PY{k+kd}{public}\PY{+w}{ }\PY{k+kd}{class} \PY{n+nc}{LivePriceTicker}\PY{+w}{ }\PY{p}{\PYZob{}}

\PY{+w}{    }\PY{k+kd}{private}\PY{+w}{ }\PY{k+kd}{final}\PY{+w}{ }\PY{n}{List}\PY{o}{\PYZlt{}}\PY{n}{PriceListener}\PY{o}{\PYZgt{}}\PY{+w}{ }\PY{n}{listeners}\PY{+w}{ }\PY{o}{=}\PY{+w}{ }\PY{k}{new}\PY{+w}{ }\PY{n}{CopyOnWriteArrayList}\PY{o}{\PYZlt{}}\PY{o}{\PYZgt{}}\PY{p}{(}\PY{p}{)}\PY{p}{;}
\PY{+w}{    }\PY{k+kd}{private}\PY{+w}{ }\PY{k+kd}{final}\PY{+w}{ }\PY{n}{Thread}\PY{+w}{ }\PY{n}{refresher}\PY{p}{;}
\PY{+w}{    }\PY{k+kd}{private}\PY{+w}{ }\PY{k+kd}{volatile}\PY{+w}{ }\PY{k+kt}{double}\PY{+w}{ }\PY{n}{lastPrice}\PY{p}{;}
\PY{+w}{    }\PY{k+kd}{private}\PY{+w}{ }\PY{k+kd}{volatile}\PY{+w}{ }\PY{k+kt}{boolean}\PY{+w}{ }\PY{n}{running}\PY{+w}{ }\PY{o}{=}\PY{+w}{ }\PY{k+kc}{true}\PY{p}{;}

\PY{+w}{    }\PY{k+kd}{private}\PY{+w}{ }\PY{n+nf}{LivePriceTicker}\PY{p}{(}\PY{p}{)}\PY{+w}{ }\PY{p}{\PYZob{}}
\PY{+w}{        }\PY{k}{this}\PY{p}{.}\PY{n+na}{refresher}\PY{+w}{ }\PY{o}{=}\PY{+w}{ }\PY{k}{new}\PY{+w}{ }\PY{n}{Thread}\PY{p}{(}\PY{k}{this}\PY{p}{:}\PY{p}{:}\PY{n}{refreshLoop}\PY{p}{)}\PY{p}{;}
\PY{+w}{    }\PY{p}{\PYZcb{}}

\PY{+w}{    }\PY{c+c1}{// Factory method: object is fully constructed before anything external can see it.}
\PY{+w}{    }\PY{k+kd}{public}\PY{+w}{ }\PY{k+kd}{static}\PY{+w}{ }\PY{n}{LivePriceTicker}\PY{+w}{ }\PY{n+nf}{start}\PY{p}{(}\PY{n}{PriceFeed}\PY{+w}{ }\PY{n}{feed}\PY{p}{)}\PY{+w}{ }\PY{p}{\PYZob{}}
\PY{+w}{        }\PY{n}{LivePriceTicker}\PY{+w}{ }\PY{n}{ticker}\PY{+w}{ }\PY{o}{=}\PY{+w}{ }\PY{k}{new}\PY{+w}{ }\PY{n}{LivePriceTicker}\PY{p}{(}\PY{p}{)}\PY{p}{;}
\PY{+w}{        }\PY{n}{feed}\PY{p}{.}\PY{n+na}{subscribe}\PY{p}{(}\PY{n}{ticker}\PY{p}{:}\PY{p}{:}\PY{n}{onPriceUpdate}\PY{p}{)}\PY{p}{;}\PY{+w}{ }\PY{c+c1}{// \PYZsq{}this\PYZsq{} only escapes after the constructor returned}
\PY{+w}{        }\PY{n}{ticker}\PY{p}{.}\PY{n+na}{refresher}\PY{p}{.}\PY{n+na}{start}\PY{p}{(}\PY{p}{)}\PY{p}{;}
\PY{+w}{        }\PY{k}{return}\PY{+w}{ }\PY{n}{ticker}\PY{p}{;}
\PY{+w}{    }\PY{p}{\PYZcb{}}

\PY{+w}{    }\PY{k+kd}{private}\PY{+w}{ }\PY{k+kt}{void}\PY{+w}{ }\PY{n+nf}{refreshLoop}\PY{p}{(}\PY{p}{)}\PY{+w}{ }\PY{p}{\PYZob{}}
\PY{+w}{        }\PY{k}{while}\PY{+w}{ }\PY{p}{(}\PY{n}{running}\PY{p}{)}\PY{+w}{ }\PY{p}{\PYZob{}}
\PY{+w}{            }\PY{n}{refreshFromFeed}\PY{p}{(}\PY{p}{)}\PY{p}{;}
\PY{+w}{            }\PY{n}{sleepQuietly}\PY{p}{(}\PY{p}{)}\PY{p}{;}
\PY{+w}{        }\PY{p}{\PYZcb{}}
\PY{+w}{    }\PY{p}{\PYZcb{}}

\PY{+w}{    }\PY{k+kd}{public}\PY{+w}{ }\PY{k+kt}{void}\PY{+w}{ }\PY{n+nf}{shutdown}\PY{p}{(}\PY{p}{)}\PY{+w}{ }\PY{p}{\PYZob{}}
\PY{+w}{        }\PY{n}{running}\PY{+w}{ }\PY{o}{=}\PY{+w}{ }\PY{k+kc}{false}\PY{p}{;}
\PY{+w}{        }\PY{n}{refresher}\PY{p}{.}\PY{n+na}{interrupt}\PY{p}{(}\PY{p}{)}\PY{p}{;}
\PY{+w}{    }\PY{p}{\PYZcb{}}

\PY{+w}{    }\PY{k+kd}{public}\PY{+w}{ }\PY{k+kt}{void}\PY{+w}{ }\PY{n+nf}{addListener}\PY{p}{(}\PY{n}{PriceListener}\PY{+w}{ }\PY{n}{listener}\PY{p}{)}\PY{+w}{ }\PY{p}{\PYZob{}}
\PY{+w}{        }\PY{n}{listeners}\PY{p}{.}\PY{n+na}{add}\PY{p}{(}\PY{n}{listener}\PY{p}{)}\PY{p}{;}
\PY{+w}{    }\PY{p}{\PYZcb{}}

\PY{+w}{    }\PY{k+kd}{private}\PY{+w}{ }\PY{k+kt}{void}\PY{+w}{ }\PY{n+nf}{onPriceUpdate}\PY{p}{(}\PY{k+kt}{double}\PY{+w}{ }\PY{n}{price}\PY{p}{)}\PY{+w}{ }\PY{p}{\PYZob{}}
\PY{+w}{        }\PY{k}{this}\PY{p}{.}\PY{n+na}{lastPrice}\PY{+w}{ }\PY{o}{=}\PY{+w}{ }\PY{n}{price}\PY{p}{;}
\PY{+w}{        }\PY{k}{for}\PY{+w}{ }\PY{p}{(}\PY{n}{PriceListener}\PY{+w}{ }\PY{n}{listener}\PY{+w}{ }\PY{p}{:}\PY{+w}{ }\PY{n}{listeners}\PY{p}{)}\PY{+w}{ }\PY{p}{\PYZob{}}
\PY{+w}{            }\PY{n}{listener}\PY{p}{.}\PY{n+na}{onPrice}\PY{p}{(}\PY{n}{price}\PY{p}{)}\PY{p}{;}
\PY{+w}{        }\PY{p}{\PYZcb{}}
\PY{+w}{    }\PY{p}{\PYZcb{}}

\PY{+w}{    }\PY{k+kd}{private}\PY{+w}{ }\PY{k+kt}{void}\PY{+w}{ }\PY{n+nf}{refreshFromFeed}\PY{p}{(}\PY{p}{)}\PY{+w}{ }\PY{p}{\PYZob{}}\PY{+w}{ }\PY{c+cm}{/* ... */}\PY{+w}{ }\PY{p}{\PYZcb{}}
\PY{+w}{    }\PY{k+kd}{private}\PY{+w}{ }\PY{k+kt}{void}\PY{+w}{ }\PY{n+nf}{sleepQuietly}\PY{p}{(}\PY{p}{)}\PY{+w}{ }\PY{p}{\PYZob{}}\PY{+w}{ }\PY{k}{try}\PY{+w}{ }\PY{p}{\PYZob{}}\PY{+w}{ }\PY{n}{Thread}\PY{p}{.}\PY{n+na}{sleep}\PY{p}{(}\PY{l+m+mi}{1000}\PY{p}{)}\PY{p}{;}\PY{+w}{ }\PY{p}{\PYZcb{}}\PY{+w}{ }\PY{k}{catch}\PY{+w}{ }\PY{p}{(}\PY{n}{InterruptedException}\PY{+w}{ }\PY{n}{e}\PY{p}{)}\PY{+w}{ }\PY{p}{\PYZob{}}\PY{+w}{ }\PY{n}{Thread}\PY{p}{.}\PY{n+na}{currentThread}\PY{p}{(}\PY{p}{)}\PY{p}{.}\PY{n+na}{interrupt}\PY{p}{(}\PY{p}{)}\PY{p}{;}\PY{+w}{ }\PY{p}{\PYZcb{}}\PY{+w}{ }\PY{p}{\PYZcb{}}
\PY{p}{\PYZcb{}}
\end{Verbatim}
\end{codebox}

See \hyperref[chap:13]{Chapter~13, \emph{Concurrency (Core)}} for the Java Memory Model rules behind “safe publication,” and \hyperref[chap:2]{Chapter~2, \emph{Object-Oriented Programming}} for constructor discipline in general.

\end{sidebar}
\section[Performance Optimization]{Performance Optimization}
\label{h:22:performance-optimization}
A checklist first, then the wins that show up over and over in review.

\subsection[Checklist]{Checklist}
\label{h:22:checklist}
\begin{itemize}
\item 
[ ] Is there actual evidence (a profile, a benchmark, a production metric) that this code is a bottleneck, or is this optimization speculative?

\item 
[ ] Is the algorithmic complexity (Big-O) obviously worse than it needs to be for the data sizes involved?

\item 
[ ] Is the right data structure being used for the access pattern (see \hyperref[h:22:collections-selection]{Collections Selection})?

\item 
[ ] Is boxing/unboxing happening in a numeric hot path that could use primitives?

\item 
[ ] Is a \code{Pattern} being compiled repeatedly (e.g., inside a loop or via \code{String.\allowbreak{}matches()}) instead of once?

\item 
[ ] Is I/O unbuffered, or done byte-by-byte/line-by-line in a way that causes excessive system calls?

\item 
[ ] Are database calls batched, or issued one-by-one in a loop (N+1)?

\item 
[ ] Is expensive work done eagerly at startup/construction when it could be lazy, or vice versa — lazily recomputed on every call when it could be cached once?

\item 
[ ] Is parallelism (\code{parallelStream}, manual thread pools) applied to trivial or I/O-bound work where the overhead outweighs the benefit?

\end{itemize}

\subsection[Measure First]{Measure First}
\label{h:22:measure-first}
The single biggest mistake in performance review is optimizing before measuring. “This \emph{looks} slow” is not evidence; a profiler, a JFR recording (\hyperref[chap:12]{Chapter~12, \emph{Memory Management}}), or a microbenchmark is. For method-level questions (“is \code{StringBuilder} actually faster here, and by how much”), the standard tool is \textbf{JMH} (Java Microbenchmark Harness) — it handles JIT warm-up, dead-code elimination by the optimizer, and statistical noise that a naive \code{System.\allowbreak{}nanoTime()} loop gets wrong. In an interview, you won’t run JMH, but saying “I’d confirm this with a JMH benchmark rather than guessing” is exactly the right instinct to voice.

\subsection[Big-O Awareness]{Big-O Awareness}
\label{h:22:big-o-awareness}
Recognize complexity by shape, without doing formal analysis:

\begin{itemize}
\item 
A loop calling \code{.\allowbreak{}contains()} or \code{.\allowbreak{}get()} on a \code{List} inside another loop → O(n²) (should be a \code{Set}/\code{Map} lookup, O(n)).

\item 
A loop making a DB/network call per iteration → O(n) round trips that should be O(1) (Scenario 6).

\item 
Nested loops over independent dimensions that don’t need to be recomputed per outer iteration → hoist the inner computation out or precompute it once (Scenario 6’s discount rules).

\item 
String-building with \code{+} across iterations → O(n²) character copying (Scenario 7).

\end{itemize}

\begin{codebox}
\begin{Verbatim}[commandchars=\\\{\}]
\PY{c+c1}{// Flagged: O(n * m) — .contains() on a List is a linear scan, done once per outer element}
\PY{n}{List}\PY{o}{\PYZlt{}}\PY{n}{String}\PY{o}{\PYZgt{}}\PY{+w}{ }\PY{n}{blockedIds}\PY{+w}{ }\PY{o}{=}\PY{+w}{ }\PY{n}{loadBlockedIds}\PY{p}{(}\PY{p}{)}\PY{p}{;}\PY{+w}{ }\PY{c+c1}{// a List}
\PY{k}{for}\PY{+w}{ }\PY{p}{(}\PY{n}{Order}\PY{+w}{ }\PY{n}{order}\PY{+w}{ }\PY{p}{:}\PY{+w}{ }\PY{n}{orders}\PY{p}{)}\PY{+w}{ }\PY{p}{\PYZob{}}
\PY{+w}{    }\PY{k}{if}\PY{+w}{ }\PY{p}{(}\PY{n}{blockedIds}\PY{p}{.}\PY{n+na}{contains}\PY{p}{(}\PY{n}{order}\PY{p}{.}\PY{n+na}{customerId}\PY{p}{(}\PY{p}{)}\PY{p}{)}\PY{p}{)}\PY{+w}{ }\PY{p}{\PYZob{}}\PY{+w}{ }\PY{p}{.}\PY{p}{.}\PY{p}{.}\PY{+w}{ }\PY{p}{\PYZcb{}}
\PY{p}{\PYZcb{}}
\end{Verbatim}
\end{codebox}

\begin{codebox}
\begin{Verbatim}[commandchars=\\\{\}]
\PY{c+c1}{// Fixed: O(n + m) — one linear pass to build the Set, then O(1) lookups}
\PY{n}{Set}\PY{o}{\PYZlt{}}\PY{n}{String}\PY{o}{\PYZgt{}}\PY{+w}{ }\PY{n}{blockedIds}\PY{+w}{ }\PY{o}{=}\PY{+w}{ }\PY{k}{new}\PY{+w}{ }\PY{n}{HashSet}\PY{o}{\PYZlt{}}\PY{o}{\PYZgt{}}\PY{p}{(}\PY{n}{loadBlockedIds}\PY{p}{(}\PY{p}{)}\PY{p}{)}\PY{p}{;}
\PY{k}{for}\PY{+w}{ }\PY{p}{(}\PY{n}{Order}\PY{+w}{ }\PY{n}{order}\PY{+w}{ }\PY{p}{:}\PY{+w}{ }\PY{n}{orders}\PY{p}{)}\PY{+w}{ }\PY{p}{\PYZob{}}
\PY{+w}{    }\PY{k}{if}\PY{+w}{ }\PY{p}{(}\PY{n}{blockedIds}\PY{p}{.}\PY{n+na}{contains}\PY{p}{(}\PY{n}{order}\PY{p}{.}\PY{n+na}{customerId}\PY{p}{(}\PY{p}{)}\PY{p}{)}\PY{p}{)}\PY{+w}{ }\PY{p}{\PYZob{}}\PY{+w}{ }\PY{p}{.}\PY{p}{.}\PY{p}{.}\PY{+w}{ }\PY{p}{\PYZcb{}}
\PY{p}{\PYZcb{}}
\end{Verbatim}
\end{codebox}

\subsection[Common Real Wins]{Common Real Wins}
\label{h:22:common-real-wins}
\begin{codebox}
\begin{Verbatim}[commandchars=\\\{\}]
\PY{c+c1}{// Flagged: Pattern compiled on every call — String.matches() compiles internally, every time}
\PY{k+kd}{public}\PY{+w}{ }\PY{k+kt}{boolean}\PY{+w}{ }\PY{n+nf}{isValidSku}\PY{p}{(}\PY{n}{String}\PY{+w}{ }\PY{n}{sku}\PY{p}{)}\PY{+w}{ }\PY{p}{\PYZob{}}
\PY{+w}{    }\PY{k}{return}\PY{+w}{ }\PY{n}{sku}\PY{p}{.}\PY{n+na}{matches}\PY{p}{(}\PY{l+s}{\PYZdq{}}\PY{l+s}{[A\PYZhy{}Z]\PYZob{}3\PYZcb{}\PYZhy{}}\PY{l+s}{\PYZbs{}\PYZbs{}}\PY{l+s}{d\PYZob{}4\PYZcb{}}\PY{l+s}{\PYZdq{}}\PY{p}{)}\PY{p}{;}
\PY{p}{\PYZcb{}}
\end{Verbatim}
\end{codebox}

\begin{codebox}
\begin{Verbatim}[commandchars=\\\{\}]
\PY{c+c1}{// Fixed: compile once, reuse the compiled Pattern}
\PY{k+kd}{private}\PY{+w}{ }\PY{k+kd}{static}\PY{+w}{ }\PY{k+kd}{final}\PY{+w}{ }\PY{n}{Pattern}\PY{+w}{ }\PY{n}{SKU\PYZus{}PATTERN}\PY{+w}{ }\PY{o}{=}\PY{+w}{ }\PY{n}{Pattern}\PY{p}{.}\PY{n+na}{compile}\PY{p}{(}\PY{l+s}{\PYZdq{}}\PY{l+s}{[A\PYZhy{}Z]\PYZob{}3\PYZcb{}\PYZhy{}}\PY{l+s}{\PYZbs{}\PYZbs{}}\PY{l+s}{d\PYZob{}4\PYZcb{}}\PY{l+s}{\PYZdq{}}\PY{p}{)}\PY{p}{;}

\PY{k+kd}{public}\PY{+w}{ }\PY{k+kt}{boolean}\PY{+w}{ }\PY{n+nf}{isValidSku}\PY{p}{(}\PY{n}{String}\PY{+w}{ }\PY{n}{sku}\PY{p}{)}\PY{+w}{ }\PY{p}{\PYZob{}}
\PY{+w}{    }\PY{k}{return}\PY{+w}{ }\PY{n}{SKU\PYZus{}PATTERN}\PY{p}{.}\PY{n+na}{matcher}\PY{p}{(}\PY{n}{sku}\PY{p}{)}\PY{p}{.}\PY{n+na}{matches}\PY{p}{(}\PY{p}{)}\PY{p}{;}
\PY{p}{\PYZcb{}}
\end{Verbatim}
\end{codebox}

\begin{codebox}
\begin{Verbatim}[commandchars=\\\{\}]
\PY{c+c1}{// Flagged: boxed Integer accumulator in a tight numeric loop — boxes/unboxes a million times}
\PY{n}{Integer}\PY{+w}{ }\PY{n}{total}\PY{+w}{ }\PY{o}{=}\PY{+w}{ }\PY{l+m+mi}{0}\PY{p}{;}
\PY{k}{for}\PY{+w}{ }\PY{p}{(}\PY{k+kt}{int}\PY{+w}{ }\PY{n}{i}\PY{+w}{ }\PY{o}{=}\PY{+w}{ }\PY{l+m+mi}{0}\PY{p}{;}\PY{+w}{ }\PY{n}{i}\PY{+w}{ }\PY{o}{\PYZlt{}}\PY{+w}{ }\PY{n}{prices}\PY{p}{.}\PY{n+na}{length}\PY{p}{;}\PY{+w}{ }\PY{n}{i}\PY{o}{+}\PY{o}{+}\PY{p}{)}\PY{+w}{ }\PY{p}{\PYZob{}}
\PY{+w}{    }\PY{n}{total}\PY{+w}{ }\PY{o}{+}\PY{o}{=}\PY{+w}{ }\PY{n}{prices}\PY{o}{[}\PY{n}{i}\PY{o}{]}\PY{p}{;}\PY{+w}{ }\PY{c+c1}{// unbox total, add, re\PYZhy{}box into a new Integer, every iteration}
\PY{p}{\PYZcb{}}
\end{Verbatim}
\end{codebox}

\begin{codebox}
\begin{Verbatim}[commandchars=\\\{\}]
\PY{c+c1}{// Fixed: primitive accumulator, no boxing at all}
\PY{k+kt}{int}\PY{+w}{ }\PY{n}{total}\PY{+w}{ }\PY{o}{=}\PY{+w}{ }\PY{l+m+mi}{0}\PY{p}{;}
\PY{k}{for}\PY{+w}{ }\PY{p}{(}\PY{k+kt}{int}\PY{+w}{ }\PY{n}{i}\PY{+w}{ }\PY{o}{=}\PY{+w}{ }\PY{l+m+mi}{0}\PY{p}{;}\PY{+w}{ }\PY{n}{i}\PY{+w}{ }\PY{o}{\PYZlt{}}\PY{+w}{ }\PY{n}{prices}\PY{p}{.}\PY{n+na}{length}\PY{p}{;}\PY{+w}{ }\PY{n}{i}\PY{o}{+}\PY{o}{+}\PY{p}{)}\PY{+w}{ }\PY{p}{\PYZob{}}
\PY{+w}{    }\PY{n}{total}\PY{+w}{ }\PY{o}{+}\PY{o}{=}\PY{+w}{ }\PY{n}{prices}\PY{o}{[}\PY{n}{i}\PY{o}{]}\PY{p}{;}
\PY{p}{\PYZcb{}}
\end{Verbatim}
\end{codebox}

\begin{codebox}
\begin{Verbatim}[commandchars=\\\{\}]
\PY{c+c1}{// Flagged: unbuffered, one\PYZhy{}byte\PYZhy{}at\PYZhy{}a\PYZhy{}time reads — one native call per byte}
\PY{k}{try}\PY{+w}{ }\PY{p}{(}\PY{n}{FileInputStream}\PY{+w}{ }\PY{n}{in}\PY{+w}{ }\PY{o}{=}\PY{+w}{ }\PY{k}{new}\PY{+w}{ }\PY{n}{FileInputStream}\PY{p}{(}\PY{n}{path}\PY{p}{)}\PY{p}{)}\PY{+w}{ }\PY{p}{\PYZob{}}
\PY{+w}{    }\PY{k+kt}{int}\PY{+w}{ }\PY{n}{b}\PY{p}{;}
\PY{+w}{    }\PY{k}{while}\PY{+w}{ }\PY{p}{(}\PY{p}{(}\PY{n}{b}\PY{+w}{ }\PY{o}{=}\PY{+w}{ }\PY{n}{in}\PY{p}{.}\PY{n+na}{read}\PY{p}{(}\PY{p}{)}\PY{p}{)}\PY{+w}{ }\PY{o}{!}\PY{o}{=}\PY{+w}{ }\PY{o}{\PYZhy{}}\PY{l+m+mi}{1}\PY{p}{)}\PY{+w}{ }\PY{p}{\PYZob{}}\PY{+w}{ }\PY{n}{process}\PY{p}{(}\PY{n}{b}\PY{p}{)}\PY{p}{;}\PY{+w}{ }\PY{p}{\PYZcb{}}
\PY{p}{\PYZcb{}}
\end{Verbatim}
\end{codebox}

\begin{codebox}
\begin{Verbatim}[commandchars=\\\{\}]
\PY{c+c1}{// Fixed: wrap with a BufferedInputStream (or use readAllBytes/NIO for whole\PYZhy{}file reads)}
\PY{k}{try}\PY{+w}{ }\PY{p}{(}\PY{n}{BufferedInputStream}\PY{+w}{ }\PY{n}{in}\PY{+w}{ }\PY{o}{=}\PY{+w}{ }\PY{k}{new}\PY{+w}{ }\PY{n}{BufferedInputStream}\PY{p}{(}\PY{k}{new}\PY{+w}{ }\PY{n}{FileInputStream}\PY{p}{(}\PY{n}{path}\PY{p}{)}\PY{p}{)}\PY{p}{)}\PY{+w}{ }\PY{p}{\PYZob{}}
\PY{+w}{    }\PY{k+kt}{int}\PY{+w}{ }\PY{n}{b}\PY{p}{;}
\PY{+w}{    }\PY{k}{while}\PY{+w}{ }\PY{p}{(}\PY{p}{(}\PY{n}{b}\PY{+w}{ }\PY{o}{=}\PY{+w}{ }\PY{n}{in}\PY{p}{.}\PY{n+na}{read}\PY{p}{(}\PY{p}{)}\PY{p}{)}\PY{+w}{ }\PY{o}{!}\PY{o}{=}\PY{+w}{ }\PY{o}{\PYZhy{}}\PY{l+m+mi}{1}\PY{p}{)}\PY{+w}{ }\PY{p}{\PYZob{}}\PY{+w}{ }\PY{n}{process}\PY{p}{(}\PY{n}{b}\PY{p}{)}\PY{p}{;}\PY{+w}{ }\PY{p}{\PYZcb{}}
\PY{p}{\PYZcb{}}
\end{Verbatim}
\end{codebox}

\subsection[Premature Optimization Is Also a Finding]{Premature Optimization Is Also a Finding}
\label{h:22:premature-optimization-is-also-a-finding}
Flag the opposite mistake too: adding a cache, a custom data structure, or manual parallelism for code that isn’t a bottleneck adds real cost — more moving parts, more places to get thread-safety wrong (see Scenario 9, Scenario 12), harder-to-read code — for a performance gain nobody asked for or measured. “I’d leave this as a simple loop unless we have evidence it’s hot” is a legitimate, senior-sounding review comment, not a cop-out.

\section[Thread Safety Review]{Thread Safety Review}
\label{h:22:thread-safety-review}
\subsection[Checklist]{Checklist}
\label{h:22:checklist}
\begin{itemize}
\item 
[ ] What state does this object hold, and is any of it mutable?

\item 
[ ] Can more than one thread reach this object at the same time (is it a singleton, a field on a shared service, handed to an executor, stored in a servlet-scoped or application-scoped context)?

\item 
[ ] If shared and mutable: is every access to that state properly synchronized, \code{volatile}, or backed by a concurrent-safe type?

\item 
[ ] Does the object fully finish constructing before any reference to it (or to \code{this} via a lambda/method reference) is handed to another thread?

\item 
[ ] Is the class’s thread-safety contract documented, so callers don’t have to read the implementation to find out?

\end{itemize}

\subsection[Reasoning About Shared Mutable State]{Reasoning About Shared Mutable State}
\label{h:22:reasoning-about-shared-mutable-state}
The question to ask about \emph{every} field is: \textbf{who else can see this, and can they see it while I’m changing it?} Three outcomes:

\begin{enumerate}
\item 
\textbf{Not shared} — a local variable, or an object fully confined to one thread. No synchronization needed. This is the best outcome; prefer designs that keep state confined.

\item 
\textbf{Shared but immutable} — the value never changes after construction (\code{final} fields, no setters, defensively-copied collections). Safe to share freely with no locking, because there’s nothing to race on. This is the second-best outcome — see \hyperref[chap:5]{Chapter~5, \emph{Object Class \& Common APIs}} on immutability.

\item 
\textbf{Shared and mutable} — now you need an explicit safety mechanism: a lock (\code{synchronized}, \code{ReentrantLock}), \code{volatile} for single-variable visibility, an atomic class, or a concurrent collection. This is the case that needs the most scrutiny in review, and where most real bugs (Scenarios 1, 2, 9, 15) live.

\end{enumerate}

\textbf{Prefer immutability first.} Before reaching for a lock, ask whether the mutable field even needs to be mutable — could it instead be replaced wholesale (a \code{volatile} reference to a new immutable object) rather than mutated in place? This sidesteps entire classes of bugs instead of managing them.

\textbf{Confinement} is the other escape hatch: if an object is only ever touched by one thread (e.g., a per-request object never handed off, a \code{ThreadLocal}), it doesn’t need synchronization regardless of whether the class itself is thread-safe in general — but this must be true structurally, not “true in practice today,” because it silently breaks the moment someone hands the object to an executor.

\subsection[The Escape/Publication Question]{The Escape/Publication Question}
\label{h:22:the-escapepublication-question}
Ask, for every object: \textbf{how does a reference to this object first reach another thread, and has construction fully finished by then?} This is exactly what broke in Scenario 15 — \code{this} was captured by a lambda and handed to \code{feed.\allowbreak{}subscribe(...)} before the constructor returned. Safe publication options, cheapest first:

\begin{itemize}
\item 
Return the fully-constructed object from a \code{static} factory method instead of doing any handoff inside the constructor itself.

\item 
Store the reference in a \code{volatile} or \code{final} field, or into a properly locked structure, before any other thread can see it.

\item 
Use a concurrent collection (\code{ConcurrentHashMap}, etc.) to publish it — insertion into those collections has the necessary memory-visibility guarantees built in.

\end{itemize}

\subsection[Documenting Thread Safety]{Documenting Thread Safety}
\label{h:22:documenting-thread-safety}
A class’s thread-safety contract is part of its API, not an implementation detail — callers need to know it without reading the source. State it explicitly, ideally in the class Javadoc: “this class is thread-safe; all methods may be called concurrently” or “this class is not thread-safe; instances must be confined to a single thread” or “safe for concurrent reads, writes must be externally synchronized.” See \hyperref[chap:13]{Chapter~13, \emph{Concurrency (Core)}} for the Java Memory Model vocabulary (visibility, atomicity, ordering, happens-before) to use precisely when writing or reviewing such statements.

\section[Collections Selection]{Collections Selection}
\label{h:22:collections-selection}
\subsection[Decision Table]{Decision Table}
\label{h:22:decision-table}
{\tablefontsmall
\begin{xltabular}{\linewidth}{@{}L{0.995}L{0.627}L{1.146}L{1.232}@{}}
\toprule
\rowcolor{tablehdr}
\thd{Need} & \thd{Prefer} & \thd{Avoid} & \thd{Why} \\
\midrule
\endhead
Index-based random access, mostly reads & \code{ArrayList} & \code{LinkedList} & O(1) indexed get vs. O(n) traversal \\
Frequent insert/remove at both ends & \code{ArrayDeque} & \code{LinkedList}, \code{Stack} & Better cache locality, no legacy \code{synchronized} overhead \\
Fast key lookup, order doesn’t matter & \code{HashMap} / \code{HashSet} & \code{TreeMap}/\code{TreeSet} (unless sorted order needed) & O(1) average vs. O(log n) \\
Insertion order must be preserved & \code{LinkedHashMap} / \code{LinkedHashSet} & \code{HashMap} (order unspecified) & Predictable iteration order \\
Sorted key order needed & \code{TreeMap} / \code{TreeSet} & Manually sorting a \code{HashMap}’s keys repeatedly & Maintains order incrementally, O(log n) per op \\
Multiple threads read/write a map concurrently & \code{ConcurrentHashMap} & \code{HashMap}, or \code{Collections.\allowbreak{}synchronized\allowbreak{}Map} for hot paths & Lock-striped/segment-free reads, far less contention \\
Mostly-read, occasionally-written list shared across threads & \code{CopyOnWriteArrayList} & \code{synchronized(\allowbreak{}List)} with manual locking, or plain \code{ArrayList} & Read without locking at all; write cost is acceptable when reads dominate \\
Producer/consumer handoff between threads & \code{BlockingQueue} (\code{LinkedBlockingQueue}, \code{ArrayBlockingQueue}) & A \code{List} polled in a loop (“busy-wait”) & Built-in blocking wait, no CPU-spinning \\
A fixed, never-changing collection & \code{List.\allowbreak{}of(...)} / \code{Map.\allowbreak{}of(...)} & \code{Collections.\allowbreak{}unmodifiable\allowbreak{}List(\allowbreak{}new~\allowbreak{}Array\allowbreak{}List\textless{}\textgreater{}(...))} for something that’s genuinely constant & Simpler, truly immutable, slightly more memory-efficient \\
Enum-keyed map/set & \code{EnumMap} / \code{EnumSet} & \code{HashMap\textless{}\allowbreak{}SomeEnum,\allowbreak{}~\allowbreak{}V\textgreater{}} & Backed by an array indexed by ordinal — faster and more memory-compact \\
Cache that shouldn’t pin entries in memory once nothing else references the key & \code{WeakHashMap} (or better, a real cache library) & \code{HashMap} & Entries can be GC’d once the key is otherwise unreachable \\
\bottomrule
\end{xltabular}
}

See \hyperref[chap:7]{Chapter~7, \emph{Collections Framework}} for the full breakdown of each type’s internals.

\subsection[“They Used X, Should Be Y” Cases]{“They Used X, Should Be Y” Cases}
\label{h:22:they-used-x-should-be-y-cases}
\begin{codebox}
\begin{Verbatim}[commandchars=\\\{\}]
\PY{c+c1}{// Flagged: LinkedList for pure random access by index}
\PY{n}{LinkedList}\PY{o}{\PYZlt{}}\PY{n}{Transaction}\PY{o}{\PYZgt{}}\PY{+w}{ }\PY{n}{transactions}\PY{+w}{ }\PY{o}{=}\PY{+w}{ }\PY{n}{loadTransactions}\PY{p}{(}\PY{p}{)}\PY{p}{;}
\PY{n}{Transaction}\PY{+w}{ }\PY{n}{tenth}\PY{+w}{ }\PY{o}{=}\PY{+w}{ }\PY{n}{transactions}\PY{p}{.}\PY{n+na}{get}\PY{p}{(}\PY{l+m+mi}{9}\PY{p}{)}\PY{p}{;}\PY{+w}{ }\PY{c+c1}{// O(n) traversal from the head every time}
\end{Verbatim}
\end{codebox}

Should be: \code{ArrayList\textless{}\allowbreak{}Transaction\textgreater{}} — O(1) indexed access, and \code{LinkedList} rarely wins in practice even for its supposed strengths, because of pointer-chasing and per-node allocation overhead.

\begin{codebox}
\begin{Verbatim}[commandchars=\\\{\}]
\PY{c+c1}{// Flagged: HashMap iterated in a context that assumes stable, predictable order (e.g., building a UI list or a diff)}
\PY{n}{Map}\PY{o}{\PYZlt{}}\PY{n}{String}\PY{p}{,}\PY{+w}{ }\PY{n}{BigDecimal}\PY{o}{\PYZgt{}}\PY{+w}{ }\PY{n}{balancesByAccount}\PY{+w}{ }\PY{o}{=}\PY{+w}{ }\PY{k}{new}\PY{+w}{ }\PY{n}{HashMap}\PY{o}{\PYZlt{}}\PY{o}{\PYZgt{}}\PY{p}{(}\PY{p}{)}\PY{p}{;}
\PY{c+c1}{// ... populate ...}
\PY{k}{for}\PY{+w}{ }\PY{p}{(}\PY{k+kd}{var}\PY{+w}{ }\PY{n}{entry}\PY{+w}{ }\PY{p}{:}\PY{+w}{ }\PY{n}{balancesByAccount}\PY{p}{.}\PY{n+na}{entrySet}\PY{p}{(}\PY{p}{)}\PY{p}{)}\PY{+w}{ }\PY{p}{\PYZob{}}\PY{+w}{ }\PY{n}{render}\PY{p}{(}\PY{n}{entry}\PY{p}{)}\PY{p}{;}\PY{+w}{ }\PY{p}{\PYZcb{}}\PY{+w}{ }\PY{c+c1}{// order can change silently between runs/resizes}
\end{Verbatim}
\end{codebox}

Should be: \code{LinkedHashMap} if insertion order matters for the output, or \code{TreeMap} if sorted-by-key order matters — never rely on \code{HashMap}’s iteration order, which is an implementation detail, not a contract.

\begin{codebox}
\begin{Verbatim}[commandchars=\\\{\}]
\PY{c+c1}{// Flagged: Vector or Collections.synchronizedList for a hot, contended, mostly\PYZhy{}read list}
\PY{n}{List}\PY{o}{\PYZlt{}}\PY{n}{Listener}\PY{o}{\PYZgt{}}\PY{+w}{ }\PY{n}{listeners}\PY{+w}{ }\PY{o}{=}\PY{+w}{ }\PY{k}{new}\PY{+w}{ }\PY{n}{Vector}\PY{o}{\PYZlt{}}\PY{o}{\PYZgt{}}\PY{p}{(}\PY{p}{)}\PY{p}{;}\PY{+w}{ }\PY{c+c1}{// synchronizes every single method call, even reads}
\end{Verbatim}
\end{codebox}

Should be: \code{CopyOnWriteArrayList} if reads vastly outnumber writes (typical for listener lists) — no locking at all on reads.

\begin{codebox}
\begin{Verbatim}[commandchars=\\\{\}]
\PY{c+c1}{// Flagged: a List used to check membership repeatedly}
\PY{n}{List}\PY{o}{\PYZlt{}}\PY{n}{String}\PY{o}{\PYZgt{}}\PY{+w}{ }\PY{n}{allowedRoles}\PY{+w}{ }\PY{o}{=}\PY{+w}{ }\PY{n}{List}\PY{p}{.}\PY{n+na}{of}\PY{p}{(}\PY{l+s}{\PYZdq{}}\PY{l+s}{ADMIN}\PY{l+s}{\PYZdq{}}\PY{p}{,}\PY{+w}{ }\PY{l+s}{\PYZdq{}}\PY{l+s}{OWNER}\PY{l+s}{\PYZdq{}}\PY{p}{,}\PY{+w}{ }\PY{l+s}{\PYZdq{}}\PY{l+s}{BILLING}\PY{l+s}{\PYZdq{}}\PY{p}{)}\PY{p}{;}
\PY{k}{if}\PY{+w}{ }\PY{p}{(}\PY{n}{allowedRoles}\PY{p}{.}\PY{n+na}{contains}\PY{p}{(}\PY{n}{user}\PY{p}{.}\PY{n+na}{role}\PY{p}{(}\PY{p}{)}\PY{p}{)}\PY{p}{)}\PY{+w}{ }\PY{p}{\PYZob{}}\PY{+w}{ }\PY{p}{.}\PY{p}{.}\PY{p}{.}\PY{+w}{ }\PY{p}{\PYZcb{}}\PY{+w}{ }\PY{c+c1}{// fine for 3 elements, a trap if this list grows or is called in a hot loop}
\end{Verbatim}
\end{codebox}

Should be, once the list could grow or the check is on a hot path: \code{Set.\allowbreak{}of(\allowbreak{}\textquotedbl{}ADMIN\textquotedbl{},\allowbreak{}~\allowbreak{}\textquotedbl{}OWNER\textquotedbl{},\allowbreak{}~\allowbreak{}\textquotedbl{}BILLING\textquotedbl{})} — O(1) lookup instead of O(n) linear scan, and it also documents “this is a membership set,” not an ordered sequence.

\section[Exception Design]{Exception Design}
\label{h:22:exception-design}
\subsection[Checklist]{Checklist}
\label{h:22:checklist}
\begin{itemize}
\item 
[ ] Is the exception type at this API boundary checked (caller must handle/declare) or unchecked (caller may ignore)? Is that the right choice for whether callers can realistically recover?

\item 
[ ] Does a custom exception type add real domain meaning, or is it one more type to remember with no behavioral difference from an existing one?

\item 
[ ] Is any exception ever caught and silently discarded (\code{catch~\allowbreak{}(\allowbreak{}X~\allowbreak{}e)\allowbreak{}~\allowbreak{}\{\}} or logged with no rethrow and no handling)?

\item 
[ ] Is any exception logged \textbf{and} rethrown at the same layer (double reporting up the call stack)?

\item 
[ ] When wrapping a lower-level exception, is the original passed as the \code{cause} (\code{new~\allowbreak{}Foo(\allowbreak{}msg,\allowbreak{}~\allowbreak{}e)}), never dropped?

\item 
[ ] Does validation happen at the boundary (fail fast, clear message) rather than deep inside, producing a confusing downstream \code{NullPointerException} instead?

\end{itemize}

Full mechanics — \code{finally}/\code{return} interactions, suppressed exceptions, multi-catch, \code{NullPointerException} messages — are covered in \hyperref[chap:6]{Chapter~6, \emph{Exception Handling}}; this section is the \emph{design} lens specifically for review.

\subsection[Checked vs. Unchecked at API Boundaries]{Checked vs. Unchecked at API Boundaries}
\label{h:22:checked-vs-unchecked-at-api-boundaries}
{\tablefontsmall
\begin{xltabular}{\linewidth}{@{}L{1.132}L{0.868}@{}}
\toprule
\rowcolor{tablehdr}
\thd{Situation} & \thd{Choice} \\
\midrule
\endhead
Caller can genuinely do something different on failure (retry, fallback, prompt again) & Checked, or a documented unchecked type callers are expected to catch \\
Failure represents a bug in the calling code (bad argument, wrong state) & Unchecked (\code{Illegal\allowbreak{}Argument\allowbreak{}Exception}, \code{IllegalStateException}) \\
Code composed with streams/lambdas, where checked exceptions don’t fit the functional interfaces & Unchecked, or wrap with \code{UncheckedIOException}-style adapters \\
A public library boundary where you want to force callers to consciously acknowledge a failure mode & Checked \\
\bottomrule
\end{xltabular}
}

\subsection[Never Swallow, Never Log-and-Rethrow]{Never Swallow, Never Log-and-Rethrow}
\label{h:22:never-swallow-never-log-and-rethrow}
\begin{codebox}
\begin{Verbatim}[commandchars=\\\{\}]
\PY{c+c1}{// Flagged: swallowed — the failure vanishes with no trace}
\PY{k}{try}\PY{+w}{ }\PY{p}{\PYZob{}}
\PY{+w}{    }\PY{n}{inventoryService}\PY{p}{.}\PY{n+na}{reserve}\PY{p}{(}\PY{n}{sku}\PY{p}{,}\PY{+w}{ }\PY{n}{qty}\PY{p}{)}\PY{p}{;}
\PY{p}{\PYZcb{}}\PY{+w}{ }\PY{k}{catch}\PY{+w}{ }\PY{p}{(}\PY{n}{Exception}\PY{+w}{ }\PY{n}{e}\PY{p}{)}\PY{+w}{ }\PY{p}{\PYZob{}}
\PY{+w}{    }\PY{c+c1}{// ignored}
\PY{p}{\PYZcb{}}
\end{Verbatim}
\end{codebox}

\begin{codebox}
\begin{Verbatim}[commandchars=\\\{\}]
\PY{c+c1}{// Flagged: logged AND rethrown — the same failure gets logged again by every caller up the stack}
\PY{k}{try}\PY{+w}{ }\PY{p}{\PYZob{}}
\PY{+w}{    }\PY{n}{inventoryService}\PY{p}{.}\PY{n+na}{reserve}\PY{p}{(}\PY{n}{sku}\PY{p}{,}\PY{+w}{ }\PY{n}{qty}\PY{p}{)}\PY{p}{;}
\PY{p}{\PYZcb{}}\PY{+w}{ }\PY{k}{catch}\PY{+w}{ }\PY{p}{(}\PY{n}{InventoryException}\PY{+w}{ }\PY{n}{e}\PY{p}{)}\PY{+w}{ }\PY{p}{\PYZob{}}
\PY{+w}{    }\PY{n}{log}\PY{p}{.}\PY{n+na}{error}\PY{p}{(}\PY{l+s}{\PYZdq{}}\PY{l+s}{Reservation failed}\PY{l+s}{\PYZdq{}}\PY{p}{,}\PY{+w}{ }\PY{n}{e}\PY{p}{)}\PY{p}{;}
\PY{+w}{    }\PY{k}{throw}\PY{+w}{ }\PY{n}{e}\PY{p}{;}\PY{+w}{ }\PY{c+c1}{// caller will likely log this again}
\PY{p}{\PYZcb{}}
\end{Verbatim}
\end{codebox}

\begin{codebox}
\begin{Verbatim}[commandchars=\\\{\}]
\PY{c+c1}{// Fixed: handle it here, or propagate it — pick one. If propagating, don\PYZsq{}t also log here;}
\PY{c+c1}{// log once, at the boundary that actually decides what to do about it.}
\PY{k}{try}\PY{+w}{ }\PY{p}{\PYZob{}}
\PY{+w}{    }\PY{n}{inventoryService}\PY{p}{.}\PY{n+na}{reserve}\PY{p}{(}\PY{n}{sku}\PY{p}{,}\PY{+w}{ }\PY{n}{qty}\PY{p}{)}\PY{p}{;}
\PY{p}{\PYZcb{}}\PY{+w}{ }\PY{k}{catch}\PY{+w}{ }\PY{p}{(}\PY{n}{InventoryException}\PY{+w}{ }\PY{n}{e}\PY{p}{)}\PY{+w}{ }\PY{p}{\PYZob{}}
\PY{+w}{    }\PY{k}{throw}\PY{+w}{ }\PY{k}{new}\PY{+w}{ }\PY{n}{OrderFulfillmentException}\PY{p}{(}\PY{l+s}{\PYZdq{}}\PY{l+s}{Could not reserve }\PY{l+s}{\PYZdq{}}\PY{+w}{ }\PY{o}{+}\PY{+w}{ }\PY{n}{sku}\PY{p}{,}\PY{+w}{ }\PY{n}{e}\PY{p}{)}\PY{p}{;}\PY{+w}{ }\PY{c+c1}{// wrapped, cause preserved, not logged yet}
\PY{p}{\PYZcb{}}
\end{Verbatim}
\end{codebox}

\subsection[Wrap With Cause, Fail Fast]{Wrap With Cause, Fail Fast}
\label{h:22:wrap-with-cause-fail-fast}
\begin{codebox}
\begin{Verbatim}[commandchars=\\\{\}]
\PY{c+c1}{// Flagged: cause dropped, and validation happens deep inside where the error is confusing}
\PY{k+kd}{public}\PY{+w}{ }\PY{k+kt}{void}\PY{+w}{ }\PY{n+nf}{applyDiscount}\PY{p}{(}\PY{n}{Order}\PY{+w}{ }\PY{n}{order}\PY{p}{,}\PY{+w}{ }\PY{n}{String}\PY{+w}{ }\PY{n}{code}\PY{p}{)}\PY{+w}{ }\PY{p}{\PYZob{}}
\PY{+w}{    }\PY{n}{DiscountRule}\PY{+w}{ }\PY{n}{rule}\PY{+w}{ }\PY{o}{=}\PY{+w}{ }\PY{n}{ruleRepository}\PY{p}{.}\PY{n+na}{findByCode}\PY{p}{(}\PY{n}{code}\PY{p}{)}\PY{p}{;}\PY{+w}{ }\PY{c+c1}{// returns null if not found}
\PY{+w}{    }\PY{n}{order}\PY{p}{.}\PY{n+na}{applyDiscount}\PY{p}{(}\PY{n}{rule}\PY{p}{.}\PY{n+na}{percentage}\PY{p}{(}\PY{p}{)}\PY{p}{)}\PY{p}{;}\PY{+w}{ }\PY{c+c1}{// NPE here if code was invalid — confusing stack trace}
\PY{p}{\PYZcb{}}
\end{Verbatim}
\end{codebox}

\begin{codebox}
\begin{Verbatim}[commandchars=\\\{\}]
\PY{c+c1}{// Fixed: fail fast with a clear message, and preserve any real cause when wrapping}
\PY{k+kd}{public}\PY{+w}{ }\PY{k+kt}{void}\PY{+w}{ }\PY{n+nf}{applyDiscount}\PY{p}{(}\PY{n}{Order}\PY{+w}{ }\PY{n}{order}\PY{p}{,}\PY{+w}{ }\PY{n}{String}\PY{+w}{ }\PY{n}{code}\PY{p}{)}\PY{+w}{ }\PY{p}{\PYZob{}}
\PY{+w}{    }\PY{n}{Objects}\PY{p}{.}\PY{n+na}{requireNonNull}\PY{p}{(}\PY{n}{order}\PY{p}{,}\PY{+w}{ }\PY{l+s}{\PYZdq{}}\PY{l+s}{order must not be null}\PY{l+s}{\PYZdq{}}\PY{p}{)}\PY{p}{;}
\PY{+w}{    }\PY{n}{DiscountRule}\PY{+w}{ }\PY{n}{rule}\PY{+w}{ }\PY{o}{=}\PY{+w}{ }\PY{n}{ruleRepository}\PY{p}{.}\PY{n+na}{findByCode}\PY{p}{(}\PY{n}{code}\PY{p}{)}
\PY{+w}{            }\PY{p}{.}\PY{n+na}{orElseThrow}\PY{p}{(}\PY{p}{(}\PY{p}{)}\PY{+w}{ }\PY{o}{\PYZhy{}}\PY{o}{\PYZgt{}}\PY{+w}{ }\PY{k}{new}\PY{+w}{ }\PY{n}{InvalidDiscountCodeException}\PY{p}{(}\PY{l+s}{\PYZdq{}}\PY{l+s}{No discount rule for code: }\PY{l+s}{\PYZdq{}}\PY{+w}{ }\PY{o}{+}\PY{+w}{ }\PY{n}{code}\PY{p}{)}\PY{p}{)}\PY{p}{;}
\PY{+w}{    }\PY{n}{order}\PY{p}{.}\PY{n+na}{applyDiscount}\PY{p}{(}\PY{n}{rule}\PY{p}{.}\PY{n+na}{percentage}\PY{p}{(}\PY{p}{)}\PY{p}{)}\PY{p}{;}
\PY{p}{\PYZcb{}}
\end{Verbatim}
\end{codebox}

\subsection[Custom Exception Hierarchies]{Custom Exception Hierarchies}
\label{h:22:custom-exception-hierarchies}
Keep them shallow and purposeful — one or two levels, each type adding real meaning a caller might act on differently:

\begin{codebox}
\begin{Verbatim}[commandchars=\\\{\}]
\PY{k+kd}{public}\PY{+w}{ }\PY{k+kd}{class} \PY{n+nc}{OrderException}\PY{+w}{ }\PY{k+kd}{extends}\PY{+w}{ }\PY{n}{RuntimeException}\PY{+w}{ }\PY{p}{\PYZob{}}
\PY{+w}{    }\PY{k+kd}{public}\PY{+w}{ }\PY{n+nf}{OrderException}\PY{p}{(}\PY{n}{String}\PY{+w}{ }\PY{n}{message}\PY{p}{,}\PY{+w}{ }\PY{n}{Throwable}\PY{+w}{ }\PY{n}{cause}\PY{p}{)}\PY{+w}{ }\PY{p}{\PYZob{}}\PY{+w}{ }\PY{k+kd}{super}\PY{p}{(}\PY{n}{message}\PY{p}{,}\PY{+w}{ }\PY{n}{cause}\PY{p}{)}\PY{p}{;}\PY{+w}{ }\PY{p}{\PYZcb{}}
\PY{+w}{    }\PY{k+kd}{public}\PY{+w}{ }\PY{n+nf}{OrderException}\PY{p}{(}\PY{n}{String}\PY{+w}{ }\PY{n}{message}\PY{p}{)}\PY{+w}{ }\PY{p}{\PYZob{}}\PY{+w}{ }\PY{k+kd}{super}\PY{p}{(}\PY{n}{message}\PY{p}{)}\PY{p}{;}\PY{+w}{ }\PY{p}{\PYZcb{}}
\PY{p}{\PYZcb{}}

\PY{k+kd}{public}\PY{+w}{ }\PY{k+kd}{class} \PY{n+nc}{InsufficientInventoryException}\PY{+w}{ }\PY{k+kd}{extends}\PY{+w}{ }\PY{n}{OrderException}\PY{+w}{ }\PY{p}{\PYZob{}}
\PY{+w}{    }\PY{k+kd}{public}\PY{+w}{ }\PY{n+nf}{InsufficientInventoryException}\PY{p}{(}\PY{n}{String}\PY{+w}{ }\PY{n}{sku}\PY{p}{)}\PY{+w}{ }\PY{p}{\PYZob{}}
\PY{+w}{        }\PY{k+kd}{super}\PY{p}{(}\PY{l+s}{\PYZdq{}}\PY{l+s}{Insufficient inventory for SKU: }\PY{l+s}{\PYZdq{}}\PY{+w}{ }\PY{o}{+}\PY{+w}{ }\PY{n}{sku}\PY{p}{)}\PY{p}{;}
\PY{+w}{    }\PY{p}{\PYZcb{}}
\PY{p}{\PYZcb{}}

\PY{k+kd}{public}\PY{+w}{ }\PY{k+kd}{class} \PY{n+nc}{InvalidDiscountCodeException}\PY{+w}{ }\PY{k+kd}{extends}\PY{+w}{ }\PY{n}{OrderException}\PY{+w}{ }\PY{p}{\PYZob{}}
\PY{+w}{    }\PY{k+kd}{public}\PY{+w}{ }\PY{n+nf}{InvalidDiscountCodeException}\PY{p}{(}\PY{n}{String}\PY{+w}{ }\PY{n}{message}\PY{p}{)}\PY{+w}{ }\PY{p}{\PYZob{}}\PY{+w}{ }\PY{k+kd}{super}\PY{p}{(}\PY{n}{message}\PY{p}{)}\PY{p}{;}\PY{+w}{ }\PY{p}{\PYZcb{}}
\PY{p}{\PYZcb{}}
\end{Verbatim}
\end{codebox}

A reviewer should flag a hierarchy that’s either too flat (everything is a bare \code{RuntimeException}, so callers can’t distinguish failure modes without string-matching messages — pitfall \#14 from the \hyperref[h:22:common-java-interview-pitfalls]{Common Java Interview Pitfalls} list) or too deep (five levels of subclassing with no behavioral difference, adding ceremony with no payoff).

\section[API Design Review]{API Design Review}
\label{h:22:api-design-review}
\subsection[Checklist]{Checklist}
\label{h:22:checklist}
\begin{itemize}
\item 
[ ] Can any parameter or return value be \code{null}? If so, is that documented, defended against with \code{Objects.\allowbreak{}requireNonNull}, or better, redesigned to avoid \code{null} altogether?

\item 
[ ] Is \code{Optional} used only as a return type, never as a field, a parameter, or inside a collection?

\item 
[ ] Are parameters validated at the top of the method (fail fast) with a clear message, rather than failing confusingly several calls deep?

\item 
[ ] Is the type unnecessarily mutable — could fields be \code{final} and setters removed in favor of a constructor or a “with”-style copy method?

\item 
[ ] Do method names accurately signal cost and side effects (\code{getX} should be cheap and side-effect-free; \code{computeX}/\code{fetchX} signals real work)?

\item 
[ ] Do overloads behave consistently with each other, and is it clear which one a given call site will resolve to?

\item 
[ ] Would a currently-reasonable call site break if this method’s signature, exceptions, or semantics changed — i.e., is this API is easy to evolve later?

\end{itemize}

See \hyperref[chap:19]{Chapter~19, \emph{Software Engineering Best Practices}} for the broader API-design and clean-code principles this section specializes for review purposes.

\subsection[Nulls and Optional]{Nulls and Optional}
\label{h:22:nulls-and-optional}
\begin{codebox}
\begin{Verbatim}[commandchars=\\\{\}]
\PY{c+c1}{// Flagged: null used as \PYZdq{}not found,\PYZdq{} undocumented, and Optional used as a field type}
\PY{k+kd}{public}\PY{+w}{ }\PY{k+kd}{class} \PY{n+nc}{UserLookupResult}\PY{+w}{ }\PY{p}{\PYZob{}}
\PY{+w}{    }\PY{k+kd}{private}\PY{+w}{ }\PY{n}{Optional}\PY{o}{\PYZlt{}}\PY{n}{User}\PY{o}{\PYZgt{}}\PY{+w}{ }\PY{n}{user}\PY{p}{;}\PY{+w}{ }\PY{c+c1}{// don\PYZsq{}t do this — Optional as a field}
\PY{+w}{    }\PY{k+kd}{public}\PY{+w}{ }\PY{n}{User}\PY{+w}{ }\PY{n+nf}{find}\PY{p}{(}\PY{n}{String}\PY{+w}{ }\PY{n}{id}\PY{p}{)}\PY{+w}{ }\PY{p}{\PYZob{}}
\PY{+w}{        }\PY{k}{return}\PY{+w}{ }\PY{n}{repository}\PY{p}{.}\PY{n+na}{findById}\PY{p}{(}\PY{n}{id}\PY{p}{)}\PY{p}{;}\PY{+w}{ }\PY{c+c1}{// returns null if missing, no Javadoc says so}
\PY{+w}{    }\PY{p}{\PYZcb{}}
\PY{p}{\PYZcb{}}
\end{Verbatim}
\end{codebox}

\begin{codebox}
\begin{Verbatim}[commandchars=\\\{\}]
\PY{c+c1}{// Fixed: Optional only as a return type; the field itself is a plain, possibly\PYZhy{}null\PYZhy{}free reference}
\PY{k+kd}{public}\PY{+w}{ }\PY{k+kd}{class} \PY{n+nc}{UserLookupResult}\PY{+w}{ }\PY{p}{\PYZob{}}
\PY{+w}{    }\PY{c+cm}{/**}
\PY{c+cm}{     * @return the user, or empty if no user exists with the given id.}
\PY{c+cm}{     */}
\PY{+w}{    }\PY{k+kd}{public}\PY{+w}{ }\PY{n}{Optional}\PY{o}{\PYZlt{}}\PY{n}{User}\PY{o}{\PYZgt{}}\PY{+w}{ }\PY{n+nf}{find}\PY{p}{(}\PY{n}{String}\PY{+w}{ }\PY{n}{id}\PY{p}{)}\PY{+w}{ }\PY{p}{\PYZob{}}
\PY{+w}{        }\PY{k}{return}\PY{+w}{ }\PY{n}{repository}\PY{p}{.}\PY{n+na}{findById}\PY{p}{(}\PY{n}{id}\PY{p}{)}\PY{p}{;}\PY{+w}{ }\PY{c+c1}{// repository itself now returns Optional\PYZlt{}User\PYZgt{}}
\PY{+w}{    }\PY{p}{\PYZcb{}}
\PY{p}{\PYZcb{}}
\end{Verbatim}
\end{codebox}

\subsection[Parameter Validation and Naming]{Parameter Validation and Naming}
\label{h:22:parameter-validation-and-naming}
\begin{codebox}
\begin{Verbatim}[commandchars=\\\{\}]
\PY{c+c1}{// Flagged: no validation; a bad call fails deep inside, far from the actual mistake}
\PY{k+kd}{public}\PY{+w}{ }\PY{k+kt}{void}\PY{+w}{ }\PY{n+nf}{scheduleShipment}\PY{p}{(}\PY{n}{String}\PY{+w}{ }\PY{n}{orderId}\PY{p}{,}\PY{+w}{ }\PY{k+kt}{int}\PY{+w}{ }\PY{n}{daysFromNow}\PY{p}{)}\PY{+w}{ }\PY{p}{\PYZob{}}
\PY{+w}{    }\PY{n}{LocalDate}\PY{+w}{ }\PY{n}{shipDate}\PY{+w}{ }\PY{o}{=}\PY{+w}{ }\PY{n}{LocalDate}\PY{p}{.}\PY{n+na}{now}\PY{p}{(}\PY{p}{)}\PY{p}{.}\PY{n+na}{plusDays}\PY{p}{(}\PY{n}{daysFromNow}\PY{p}{)}\PY{p}{;}
\PY{+w}{    }\PY{n}{warehouse}\PY{p}{.}\PY{n+na}{schedule}\PY{p}{(}\PY{n}{orderId}\PY{p}{,}\PY{+w}{ }\PY{n}{shipDate}\PY{p}{)}\PY{p}{;}\PY{+w}{ }\PY{c+c1}{// throws something cryptic if orderId is blank}
\PY{p}{\PYZcb{}}
\end{Verbatim}
\end{codebox}

\begin{codebox}
\begin{Verbatim}[commandchars=\\\{\}]
\PY{c+c1}{// Fixed: fail fast, at the boundary, with a message that names the actual problem}
\PY{k+kd}{public}\PY{+w}{ }\PY{k+kt}{void}\PY{+w}{ }\PY{n+nf}{scheduleShipment}\PY{p}{(}\PY{n}{String}\PY{+w}{ }\PY{n}{orderId}\PY{p}{,}\PY{+w}{ }\PY{k+kt}{int}\PY{+w}{ }\PY{n}{daysFromNow}\PY{p}{)}\PY{+w}{ }\PY{p}{\PYZob{}}
\PY{+w}{    }\PY{k}{if}\PY{+w}{ }\PY{p}{(}\PY{n}{orderId}\PY{+w}{ }\PY{o}{=}\PY{o}{=}\PY{+w}{ }\PY{k+kc}{null}\PY{+w}{ }\PY{o}{|}\PY{o}{|}\PY{+w}{ }\PY{n}{orderId}\PY{p}{.}\PY{n+na}{isBlank}\PY{p}{(}\PY{p}{)}\PY{p}{)}\PY{+w}{ }\PY{p}{\PYZob{}}
\PY{+w}{        }\PY{k}{throw}\PY{+w}{ }\PY{k}{new}\PY{+w}{ }\PY{n}{IllegalArgumentException}\PY{p}{(}\PY{l+s}{\PYZdq{}}\PY{l+s}{orderId must not be blank}\PY{l+s}{\PYZdq{}}\PY{p}{)}\PY{p}{;}
\PY{+w}{    }\PY{p}{\PYZcb{}}
\PY{+w}{    }\PY{k}{if}\PY{+w}{ }\PY{p}{(}\PY{n}{daysFromNow}\PY{+w}{ }\PY{o}{\PYZlt{}}\PY{+w}{ }\PY{l+m+mi}{0}\PY{p}{)}\PY{+w}{ }\PY{p}{\PYZob{}}
\PY{+w}{        }\PY{k}{throw}\PY{+w}{ }\PY{k}{new}\PY{+w}{ }\PY{n}{IllegalArgumentException}\PY{p}{(}\PY{l+s}{\PYZdq{}}\PY{l+s}{daysFromNow must not be negative: }\PY{l+s}{\PYZdq{}}\PY{+w}{ }\PY{o}{+}\PY{+w}{ }\PY{n}{daysFromNow}\PY{p}{)}\PY{p}{;}
\PY{+w}{    }\PY{p}{\PYZcb{}}
\PY{+w}{    }\PY{n}{warehouse}\PY{p}{.}\PY{n+na}{schedule}\PY{p}{(}\PY{n}{orderId}\PY{p}{,}\PY{+w}{ }\PY{n}{LocalDate}\PY{p}{.}\PY{n+na}{now}\PY{p}{(}\PY{p}{)}\PY{p}{.}\PY{n+na}{plusDays}\PY{p}{(}\PY{n}{daysFromNow}\PY{p}{)}\PY{p}{)}\PY{p}{;}
\PY{p}{\PYZcb{}}
\end{Verbatim}
\end{codebox}

\subsection[Mutability and Backward Compatibility]{Mutability and Backward Compatibility}
\label{h:22:mutability-and-backward-compatibility}
\begin{codebox}
\begin{Verbatim}[commandchars=\\\{\}]
\PY{c+c1}{// Flagged: needlessly mutable value type, and a public setter that could break invariants}
\PY{k+kd}{public}\PY{+w}{ }\PY{k+kd}{class} \PY{n+nc}{Money}\PY{+w}{ }\PY{p}{\PYZob{}}
\PY{+w}{    }\PY{k+kd}{private}\PY{+w}{ }\PY{n}{BigDecimal}\PY{+w}{ }\PY{n}{amount}\PY{p}{;}
\PY{+w}{    }\PY{k+kd}{private}\PY{+w}{ }\PY{n}{String}\PY{+w}{ }\PY{n}{currency}\PY{p}{;}
\PY{+w}{    }\PY{k+kd}{public}\PY{+w}{ }\PY{k+kt}{void}\PY{+w}{ }\PY{n+nf}{setAmount}\PY{p}{(}\PY{n}{BigDecimal}\PY{+w}{ }\PY{n}{amount}\PY{p}{)}\PY{+w}{ }\PY{p}{\PYZob{}}\PY{+w}{ }\PY{k}{this}\PY{p}{.}\PY{n+na}{amount}\PY{+w}{ }\PY{o}{=}\PY{+w}{ }\PY{n}{amount}\PY{p}{;}\PY{+w}{ }\PY{p}{\PYZcb{}}\PY{+w}{ }\PY{c+c1}{// no validation, no currency re\PYZhy{}check}
\PY{p}{\PYZcb{}}
\end{Verbatim}
\end{codebox}

\begin{codebox}
\begin{Verbatim}[commandchars=\\\{\}]
\PY{c+c1}{// Fixed: immutable, validated at construction, changes produce a new instance}
\PY{k+kd}{public}\PY{+w}{ }\PY{k+kd}{final}\PY{+w}{ }\PY{k+kd}{class} \PY{n+nc}{Money}\PY{+w}{ }\PY{p}{\PYZob{}}
\PY{+w}{    }\PY{k+kd}{private}\PY{+w}{ }\PY{k+kd}{final}\PY{+w}{ }\PY{n}{BigDecimal}\PY{+w}{ }\PY{n}{amount}\PY{p}{;}
\PY{+w}{    }\PY{k+kd}{private}\PY{+w}{ }\PY{k+kd}{final}\PY{+w}{ }\PY{n}{String}\PY{+w}{ }\PY{n}{currency}\PY{p}{;}

\PY{+w}{    }\PY{k+kd}{public}\PY{+w}{ }\PY{n+nf}{Money}\PY{p}{(}\PY{n}{BigDecimal}\PY{+w}{ }\PY{n}{amount}\PY{p}{,}\PY{+w}{ }\PY{n}{String}\PY{+w}{ }\PY{n}{currency}\PY{p}{)}\PY{+w}{ }\PY{p}{\PYZob{}}
\PY{+w}{        }\PY{k}{if}\PY{+w}{ }\PY{p}{(}\PY{n}{amount}\PY{p}{.}\PY{n+na}{scale}\PY{p}{(}\PY{p}{)}\PY{+w}{ }\PY{o}{\PYZgt{}}\PY{+w}{ }\PY{l+m+mi}{2}\PY{p}{)}\PY{+w}{ }\PY{k}{throw}\PY{+w}{ }\PY{k}{new}\PY{+w}{ }\PY{n}{IllegalArgumentException}\PY{p}{(}\PY{l+s}{\PYZdq{}}\PY{l+s}{amount must have at most 2 decimal places}\PY{l+s}{\PYZdq{}}\PY{p}{)}\PY{p}{;}
\PY{+w}{        }\PY{k}{this}\PY{p}{.}\PY{n+na}{amount}\PY{+w}{ }\PY{o}{=}\PY{+w}{ }\PY{n}{amount}\PY{p}{;}
\PY{+w}{        }\PY{k}{this}\PY{p}{.}\PY{n+na}{currency}\PY{+w}{ }\PY{o}{=}\PY{+w}{ }\PY{n}{Objects}\PY{p}{.}\PY{n+na}{requireNonNull}\PY{p}{(}\PY{n}{currency}\PY{p}{)}\PY{p}{;}
\PY{+w}{    }\PY{p}{\PYZcb{}}

\PY{+w}{    }\PY{k+kd}{public}\PY{+w}{ }\PY{n}{Money}\PY{+w}{ }\PY{n+nf}{withAmount}\PY{p}{(}\PY{n}{BigDecimal}\PY{+w}{ }\PY{n}{newAmount}\PY{p}{)}\PY{+w}{ }\PY{p}{\PYZob{}}
\PY{+w}{        }\PY{k}{return}\PY{+w}{ }\PY{k}{new}\PY{+w}{ }\PY{n}{Money}\PY{p}{(}\PY{n}{newAmount}\PY{p}{,}\PY{+w}{ }\PY{n}{currency}\PY{p}{)}\PY{p}{;}
\PY{+w}{    }\PY{p}{\PYZcb{}}
\PY{p}{\PYZcb{}}
\end{Verbatim}
\end{codebox}

For backward compatibility: widening what a method accepts (e.g., \code{List\textless{}?\allowbreak{}~\allowbreak{}extends~\allowbreak{}T\textgreater{}} instead of \code{List\textless{}\allowbreak{}T\textgreater{}}) or narrowing what it promises to throw is generally safe; removing a public method, changing a return type, or \emph{narrowing} accepted input is a breaking change and should be flagged as a compatibility risk even if it “looks like a small tidy-up.”

\section[Memory \& GC Review]{Memory \& GC Review}
\label{h:22:memory--gc-review}
\subsection[Checklist]{Checklist}
\label{h:22:checklist}
\begin{itemize}
\item 
[ ] Any \code{static} collection/field that only ever grows (cache, registry, listener list) with no eviction or removal path?

\item 
[ ] Any allocation happening inside a hot loop that could be hoisted outside it (a \code{new} for an object that could be reused, a boxed wrapper, a formatter)?

\item 
[ ] Any long-lived object holding a reference to something that should have a short lifetime (e.g., an outer class implicitly held by a non-static inner class registered as a long-lived listener)?

\item 
[ ] Any \code{ThreadLocal} used inside a thread-pool-backed executor, without being explicitly cleared when the task finishes?

\item 
[ ] Would a heap dump plausibly show one dominant retained-size culprit here, or is this a death-by-a-thousand-small-allocations concern instead?

\end{itemize}

Deep GC algorithm mechanics (G1, ZGC, Shenandoah, tuning flags) live in \hyperref[chap:12]{Chapter~12, \emph{Memory Management}}; this section is about spotting the \emph{code patterns} that create memory pressure or leaks in the first place, regardless of which collector is running.

\subsection[Leak-Prone Patterns]{Leak-Prone Patterns}
\label{h:22:leak-prone-patterns}
\begin{codebox}
\begin{Verbatim}[commandchars=\\\{\}]
\PY{c+c1}{// Flagged: static, ever\PYZhy{}growing collection — a permanent, unbounded leak}
\PY{k+kd}{public}\PY{+w}{ }\PY{k+kd}{class} \PY{n+nc}{SessionRegistry}\PY{+w}{ }\PY{p}{\PYZob{}}
\PY{+w}{    }\PY{k+kd}{private}\PY{+w}{ }\PY{k+kd}{static}\PY{+w}{ }\PY{k+kd}{final}\PY{+w}{ }\PY{n}{Map}\PY{o}{\PYZlt{}}\PY{n}{String}\PY{p}{,}\PY{+w}{ }\PY{n}{Session}\PY{o}{\PYZgt{}}\PY{+w}{ }\PY{n}{ACTIVE\PYZus{}SESSIONS}\PY{+w}{ }\PY{o}{=}\PY{+w}{ }\PY{k}{new}\PY{+w}{ }\PY{n}{HashMap}\PY{o}{\PYZlt{}}\PY{o}{\PYZgt{}}\PY{p}{(}\PY{p}{)}\PY{p}{;}
\PY{+w}{    }\PY{k+kd}{public}\PY{+w}{ }\PY{k+kd}{static}\PY{+w}{ }\PY{k+kt}{void}\PY{+w}{ }\PY{n+nf}{register}\PY{p}{(}\PY{n}{Session}\PY{+w}{ }\PY{n}{s}\PY{p}{)}\PY{+w}{ }\PY{p}{\PYZob{}}\PY{+w}{ }\PY{n}{ACTIVE\PYZus{}SESSIONS}\PY{p}{.}\PY{n+na}{put}\PY{p}{(}\PY{n}{s}\PY{p}{.}\PY{n+na}{id}\PY{p}{(}\PY{p}{)}\PY{p}{,}\PY{+w}{ }\PY{n}{s}\PY{p}{)}\PY{p}{;}\PY{+w}{ }\PY{p}{\PYZcb{}}\PY{+w}{ }\PY{c+c1}{// never removed}
\PY{p}{\PYZcb{}}
\end{Verbatim}
\end{codebox}

\begin{codebox}
\begin{Verbatim}[commandchars=\\\{\}]
\PY{c+c1}{// Fixed: an explicit removal path tied to the session\PYZsq{}s actual lifecycle, or a bounded/expiring cache}
\PY{k+kd}{public}\PY{+w}{ }\PY{k+kd}{class} \PY{n+nc}{SessionRegistry}\PY{+w}{ }\PY{p}{\PYZob{}}
\PY{+w}{    }\PY{k+kd}{private}\PY{+w}{ }\PY{k+kd}{static}\PY{+w}{ }\PY{k+kd}{final}\PY{+w}{ }\PY{n}{Map}\PY{o}{\PYZlt{}}\PY{n}{String}\PY{p}{,}\PY{+w}{ }\PY{n}{Session}\PY{o}{\PYZgt{}}\PY{+w}{ }\PY{n}{ACTIVE\PYZus{}SESSIONS}\PY{+w}{ }\PY{o}{=}\PY{+w}{ }\PY{k}{new}\PY{+w}{ }\PY{n}{ConcurrentHashMap}\PY{o}{\PYZlt{}}\PY{o}{\PYZgt{}}\PY{p}{(}\PY{p}{)}\PY{p}{;}
\PY{+w}{    }\PY{k+kd}{public}\PY{+w}{ }\PY{k+kd}{static}\PY{+w}{ }\PY{k+kt}{void}\PY{+w}{ }\PY{n+nf}{register}\PY{p}{(}\PY{n}{Session}\PY{+w}{ }\PY{n}{s}\PY{p}{)}\PY{+w}{ }\PY{p}{\PYZob{}}\PY{+w}{ }\PY{n}{ACTIVE\PYZus{}SESSIONS}\PY{p}{.}\PY{n+na}{put}\PY{p}{(}\PY{n}{s}\PY{p}{.}\PY{n+na}{id}\PY{p}{(}\PY{p}{)}\PY{p}{,}\PY{+w}{ }\PY{n}{s}\PY{p}{)}\PY{p}{;}\PY{+w}{ }\PY{p}{\PYZcb{}}
\PY{+w}{    }\PY{k+kd}{public}\PY{+w}{ }\PY{k+kd}{static}\PY{+w}{ }\PY{k+kt}{void}\PY{+w}{ }\PY{n+nf}{unregister}\PY{p}{(}\PY{n}{String}\PY{+w}{ }\PY{n}{sessionId}\PY{p}{)}\PY{+w}{ }\PY{p}{\PYZob{}}\PY{+w}{ }\PY{n}{ACTIVE\PYZus{}SESSIONS}\PY{p}{.}\PY{n+na}{remove}\PY{p}{(}\PY{n}{sessionId}\PY{p}{)}\PY{p}{;}\PY{+w}{ }\PY{p}{\PYZcb{}}\PY{+w}{ }\PY{c+c1}{// called on logout/expiry}
\PY{p}{\PYZcb{}}
\end{Verbatim}
\end{codebox}

\subsection[ThreadLocal in a Pool]{ThreadLocal in a Pool}
\label{h:22:threadlocal-in-a-pool}
\begin{codebox}
\begin{Verbatim}[commandchars=\\\{\}]
\PY{c+c1}{// Flagged: ThreadLocal set but never removed, on a thread that will be reused by the pool}
\PY{k+kd}{public}\PY{+w}{ }\PY{k+kd}{class} \PY{n+nc}{RequestContext}\PY{+w}{ }\PY{p}{\PYZob{}}
\PY{+w}{    }\PY{k+kd}{private}\PY{+w}{ }\PY{k+kd}{static}\PY{+w}{ }\PY{k+kd}{final}\PY{+w}{ }\PY{n}{ThreadLocal}\PY{o}{\PYZlt{}}\PY{n}{User}\PY{o}{\PYZgt{}}\PY{+w}{ }\PY{n}{CURRENT\PYZus{}USER}\PY{+w}{ }\PY{o}{=}\PY{+w}{ }\PY{k}{new}\PY{+w}{ }\PY{n}{ThreadLocal}\PY{o}{\PYZlt{}}\PY{o}{\PYZgt{}}\PY{p}{(}\PY{p}{)}\PY{p}{;}
\PY{+w}{    }\PY{k+kd}{public}\PY{+w}{ }\PY{k+kt}{void}\PY{+w}{ }\PY{n+nf}{handle}\PY{p}{(}\PY{n}{Request}\PY{+w}{ }\PY{n}{req}\PY{p}{)}\PY{+w}{ }\PY{p}{\PYZob{}}
\PY{+w}{        }\PY{n}{CURRENT\PYZus{}USER}\PY{p}{.}\PY{n+na}{set}\PY{p}{(}\PY{n}{req}\PY{p}{.}\PY{n+na}{user}\PY{p}{(}\PY{p}{)}\PY{p}{)}\PY{p}{;}
\PY{+w}{        }\PY{n}{process}\PY{p}{(}\PY{n}{req}\PY{p}{)}\PY{p}{;}
\PY{+w}{        }\PY{c+c1}{// no CURRENT\PYZus{}USER.remove() — the User object (and everything it references) stays}
\PY{+w}{        }\PY{c+c1}{// reachable from this pool thread until some *unrelated* future task overwrites it}
\PY{+w}{    }\PY{p}{\PYZcb{}}
\PY{p}{\PYZcb{}}
\end{Verbatim}
\end{codebox}

\begin{codebox}
\begin{Verbatim}[commandchars=\\\{\}]
\PY{c+c1}{// Fixed: always remove in a finally, so the pool thread doesn\PYZsq{}t retain the value between tasks}
\PY{k+kd}{public}\PY{+w}{ }\PY{k+kd}{class} \PY{n+nc}{RequestContext}\PY{+w}{ }\PY{p}{\PYZob{}}
\PY{+w}{    }\PY{k+kd}{private}\PY{+w}{ }\PY{k+kd}{static}\PY{+w}{ }\PY{k+kd}{final}\PY{+w}{ }\PY{n}{ThreadLocal}\PY{o}{\PYZlt{}}\PY{n}{User}\PY{o}{\PYZgt{}}\PY{+w}{ }\PY{n}{CURRENT\PYZus{}USER}\PY{+w}{ }\PY{o}{=}\PY{+w}{ }\PY{k}{new}\PY{+w}{ }\PY{n}{ThreadLocal}\PY{o}{\PYZlt{}}\PY{o}{\PYZgt{}}\PY{p}{(}\PY{p}{)}\PY{p}{;}
\PY{+w}{    }\PY{k+kd}{public}\PY{+w}{ }\PY{k+kt}{void}\PY{+w}{ }\PY{n+nf}{handle}\PY{p}{(}\PY{n}{Request}\PY{+w}{ }\PY{n}{req}\PY{p}{)}\PY{+w}{ }\PY{p}{\PYZob{}}
\PY{+w}{        }\PY{n}{CURRENT\PYZus{}USER}\PY{p}{.}\PY{n+na}{set}\PY{p}{(}\PY{n}{req}\PY{p}{.}\PY{n+na}{user}\PY{p}{(}\PY{p}{)}\PY{p}{)}\PY{p}{;}
\PY{+w}{        }\PY{k}{try}\PY{+w}{ }\PY{p}{\PYZob{}}
\PY{+w}{            }\PY{n}{process}\PY{p}{(}\PY{n}{req}\PY{p}{)}\PY{p}{;}
\PY{+w}{        }\PY{p}{\PYZcb{}}\PY{+w}{ }\PY{k}{finally}\PY{+w}{ }\PY{p}{\PYZob{}}
\PY{+w}{            }\PY{n}{CURRENT\PYZus{}USER}\PY{p}{.}\PY{n+na}{remove}\PY{p}{(}\PY{p}{)}\PY{p}{;}
\PY{+w}{        }\PY{p}{\PYZcb{}}
\PY{+w}{    }\PY{p}{\PYZcb{}}
\PY{p}{\PYZcb{}}
\end{Verbatim}
\end{codebox}

This matters specifically \emph{because} pool threads are long-lived and reused — a \code{ThreadLocal} left set silently extends an object’s lifetime far beyond the request that created it, and can even leak one user’s data into logic that (mistakenly) reads the \code{ThreadLocal} during a later, different request on the same pooled thread.

\subsection[Heap Dump Basics (What a Reviewer Should Know to Ask For)]{Heap Dump Basics (What a Reviewer Should Know to Ask For)}
\label{h:22:heap-dump-basics-what-a-reviewer-should-know-to-ask-for}
When a leak is suspected but not obvious from reading the code, the review-worthy question is “have we taken a heap dump (\code{jcmd~\allowbreak{}\textless{}\allowbreak{}pid\textgreater{}\allowbreak{}~\allowbreak{}GC.\allowbreak{}heap\_\allowbreak{}dump}, or \code{-\allowbreak{}XX:+\allowbreak{}Heap\allowbreak{}Dump\allowbreak{}On\allowbreak{}Out\allowbreak{}Of\allowbreak{}Memory\allowbreak{}Error}) and looked at retained size by class/dominator tree?” A handful of instances of one class holding a disproportionate amount of retained memory is the signature of exactly the patterns above — an unbounded map, a listener list nobody ever removes from, or an accidental reference chain keeping something alive. See \hyperref[chap:12]{Chapter~12, \emph{Memory Management}} for the tools (\code{jmap}, \code{jcmd}, JFR, JMC) in depth.

\section[Concurrency Review]{Concurrency Review}
\label{h:22:concurrency-review}
\subsection[Checklist]{Checklist}
\label{h:22:checklist}
\begin{itemize}
\item 
[ ] Is every \code{ExecutorService} shut down somewhere on every relevant path (normal completion \emph{and} error/shutdown paths)?

\item 
[ ] Any double-checked locking pattern — is the field \code{volatile}? Without it, the second thread can see a partially-constructed object (see \hyperref[chap:13]{Chapter~13, \emph{Concurrency (Core)}}’s Java Memory Model section).

\item 
[ ] Is \code{volatile} used where an actual atomic compound operation (\code{i++}, check-then-act) was needed instead? \code{volatile} only guarantees visibility, not atomicity of anything beyond a single read or write.

\item 
[ ] If multiple locks are acquired, is the acquisition order consistent everywhere, avoiding a lock-ordering deadlock?

\item 
[ ] Does any \code{CompletableFuture} chain swallow exceptions (no \code{.\allowbreak{}exceptionally}/\code{.\allowbreak{}handle}), or block on \code{.\allowbreak{}get()}/\code{.\allowbreak{}join()} on a thread that shouldn’t block (e.g., an event-loop thread)?

\item 
[ ] If virtual threads are used, does any code hold a lock (\code{synchronized}) or otherwise pin the platform thread across a blocking call, defeating the point of virtual threads?

\end{itemize}

\subsection[Double-Checked Locking Without volatile]{Double-Checked Locking Without \code{volatile}}
\label{h:22:double-checked-locking-without-volatile}
\begin{codebox}
\begin{Verbatim}[commandchars=\\\{\}]
\PY{c+c1}{// Flagged: classic double\PYZhy{}checked locking bug — missing volatile}
\PY{k+kd}{public}\PY{+w}{ }\PY{k+kd}{class} \PY{n+nc}{ExpensiveResource}\PY{+w}{ }\PY{p}{\PYZob{}}
\PY{+w}{    }\PY{k+kd}{private}\PY{+w}{ }\PY{k+kd}{static}\PY{+w}{ }\PY{n}{ExpensiveResource}\PY{+w}{ }\PY{n}{instance}\PY{p}{;}
\PY{+w}{    }\PY{k+kd}{public}\PY{+w}{ }\PY{k+kd}{static}\PY{+w}{ }\PY{n}{ExpensiveResource}\PY{+w}{ }\PY{n+nf}{getInstance}\PY{p}{(}\PY{p}{)}\PY{+w}{ }\PY{p}{\PYZob{}}
\PY{+w}{        }\PY{k}{if}\PY{+w}{ }\PY{p}{(}\PY{n}{instance}\PY{+w}{ }\PY{o}{=}\PY{o}{=}\PY{+w}{ }\PY{k+kc}{null}\PY{p}{)}\PY{+w}{ }\PY{p}{\PYZob{}}
\PY{+w}{            }\PY{k+kd}{synchronized}\PY{+w}{ }\PY{p}{(}\PY{n}{ExpensiveResource}\PY{p}{.}\PY{n+na}{class}\PY{p}{)}\PY{+w}{ }\PY{p}{\PYZob{}}
\PY{+w}{                }\PY{k}{if}\PY{+w}{ }\PY{p}{(}\PY{n}{instance}\PY{+w}{ }\PY{o}{=}\PY{o}{=}\PY{+w}{ }\PY{k+kc}{null}\PY{p}{)}\PY{+w}{ }\PY{p}{\PYZob{}}
\PY{+w}{                    }\PY{n}{instance}\PY{+w}{ }\PY{o}{=}\PY{+w}{ }\PY{k}{new}\PY{+w}{ }\PY{n}{ExpensiveResource}\PY{p}{(}\PY{p}{)}\PY{p}{;}\PY{+w}{ }\PY{c+c1}{// can be observed half\PYZhy{}built without volatile}
\PY{+w}{                }\PY{p}{\PYZcb{}}
\PY{+w}{            }\PY{p}{\PYZcb{}}
\PY{+w}{        }\PY{p}{\PYZcb{}}
\PY{+w}{        }\PY{k}{return}\PY{+w}{ }\PY{n}{instance}\PY{p}{;}
\PY{+w}{    }\PY{p}{\PYZcb{}}
\PY{p}{\PYZcb{}}
\end{Verbatim}
\end{codebox}

\begin{codebox}
\begin{Verbatim}[commandchars=\\\{\}]
\PY{c+c1}{// Fixed: volatile field restores the needed happens\PYZhy{}before edge, OR use the holder idiom (Scenario 1) instead}
\PY{k+kd}{public}\PY{+w}{ }\PY{k+kd}{class} \PY{n+nc}{ExpensiveResource}\PY{+w}{ }\PY{p}{\PYZob{}}
\PY{+w}{    }\PY{k+kd}{private}\PY{+w}{ }\PY{k+kd}{static}\PY{+w}{ }\PY{k+kd}{volatile}\PY{+w}{ }\PY{n}{ExpensiveResource}\PY{+w}{ }\PY{n}{instance}\PY{p}{;}
\PY{+w}{    }\PY{k+kd}{public}\PY{+w}{ }\PY{k+kd}{static}\PY{+w}{ }\PY{n}{ExpensiveResource}\PY{+w}{ }\PY{n+nf}{getInstance}\PY{p}{(}\PY{p}{)}\PY{+w}{ }\PY{p}{\PYZob{}}
\PY{+w}{        }\PY{k}{if}\PY{+w}{ }\PY{p}{(}\PY{n}{instance}\PY{+w}{ }\PY{o}{=}\PY{o}{=}\PY{+w}{ }\PY{k+kc}{null}\PY{p}{)}\PY{+w}{ }\PY{p}{\PYZob{}}
\PY{+w}{            }\PY{k+kd}{synchronized}\PY{+w}{ }\PY{p}{(}\PY{n}{ExpensiveResource}\PY{p}{.}\PY{n+na}{class}\PY{p}{)}\PY{+w}{ }\PY{p}{\PYZob{}}
\PY{+w}{                }\PY{k}{if}\PY{+w}{ }\PY{p}{(}\PY{n}{instance}\PY{+w}{ }\PY{o}{=}\PY{o}{=}\PY{+w}{ }\PY{k+kc}{null}\PY{p}{)}\PY{+w}{ }\PY{p}{\PYZob{}}
\PY{+w}{                    }\PY{n}{instance}\PY{+w}{ }\PY{o}{=}\PY{+w}{ }\PY{k}{new}\PY{+w}{ }\PY{n}{ExpensiveResource}\PY{p}{(}\PY{p}{)}\PY{p}{;}
\PY{+w}{                }\PY{p}{\PYZcb{}}
\PY{+w}{            }\PY{p}{\PYZcb{}}
\PY{+w}{        }\PY{p}{\PYZcb{}}
\PY{+w}{        }\PY{k}{return}\PY{+w}{ }\PY{n}{instance}\PY{p}{;}
\PY{+w}{    }\PY{p}{\PYZcb{}}
\PY{p}{\PYZcb{}}
\end{Verbatim}
\end{codebox}

Without \code{volatile}, another thread can see \code{instance} as non-\code{null} while still observing the \emph{partially initialized} object it points to, because the write to \code{instance} and the writes inside the constructor can be reordered from that thread’s point of view. In review, the holder idiom (Scenario 1) is almost always the simpler recommendation over fixing double-checked locking by hand.

\subsection[Deadlock Ordering]{Deadlock Ordering}
\label{h:22:deadlock-ordering}
\begin{codebox}
\begin{Verbatim}[commandchars=\\\{\}]
\PY{c+c1}{// Flagged: two methods acquire the same two locks in opposite order — classic deadlock}
\PY{k+kd}{public}\PY{+w}{ }\PY{k+kd}{class} \PY{n+nc}{AccountTransfer}\PY{+w}{ }\PY{p}{\PYZob{}}
\PY{+w}{    }\PY{k+kd}{public}\PY{+w}{ }\PY{k+kt}{void}\PY{+w}{ }\PY{n+nf}{transfer}\PY{p}{(}\PY{n}{Account}\PY{+w}{ }\PY{n}{from}\PY{p}{,}\PY{+w}{ }\PY{n}{Account}\PY{+w}{ }\PY{n}{to}\PY{p}{,}\PY{+w}{ }\PY{n}{BigDecimal}\PY{+w}{ }\PY{n}{amount}\PY{p}{)}\PY{+w}{ }\PY{p}{\PYZob{}}
\PY{+w}{        }\PY{k+kd}{synchronized}\PY{+w}{ }\PY{p}{(}\PY{n}{from}\PY{p}{)}\PY{+w}{ }\PY{p}{\PYZob{}}
\PY{+w}{            }\PY{k+kd}{synchronized}\PY{+w}{ }\PY{p}{(}\PY{n}{to}\PY{p}{)}\PY{+w}{ }\PY{p}{\PYZob{}}
\PY{+w}{                }\PY{n}{from}\PY{p}{.}\PY{n+na}{debit}\PY{p}{(}\PY{n}{amount}\PY{p}{)}\PY{p}{;}
\PY{+w}{                }\PY{n}{to}\PY{p}{.}\PY{n+na}{credit}\PY{p}{(}\PY{n}{amount}\PY{p}{)}\PY{p}{;}
\PY{+w}{            }\PY{p}{\PYZcb{}}
\PY{+w}{        }\PY{p}{\PYZcb{}}
\PY{+w}{    }\PY{p}{\PYZcb{}}
\PY{+w}{    }\PY{c+c1}{// if some other code path calls transfer(to, from, ...) concurrently, thread A can hold}
\PY{+w}{    }\PY{c+c1}{// from\PYZsq{}s lock waiting for to\PYZsq{}s, while thread B holds to\PYZsq{}s lock waiting for from\PYZsq{}s — deadlock}
\PY{p}{\PYZcb{}}
\end{Verbatim}
\end{codebox}

\begin{codebox}
\begin{Verbatim}[commandchars=\\\{\}]
\PY{c+c1}{// Fixed: always acquire locks in a consistent, global order (e.g., by account id)}
\PY{k+kd}{public}\PY{+w}{ }\PY{k+kd}{class} \PY{n+nc}{AccountTransfer}\PY{+w}{ }\PY{p}{\PYZob{}}
\PY{+w}{    }\PY{k+kd}{public}\PY{+w}{ }\PY{k+kt}{void}\PY{+w}{ }\PY{n+nf}{transfer}\PY{p}{(}\PY{n}{Account}\PY{+w}{ }\PY{n}{from}\PY{p}{,}\PY{+w}{ }\PY{n}{Account}\PY{+w}{ }\PY{n}{to}\PY{p}{,}\PY{+w}{ }\PY{n}{BigDecimal}\PY{+w}{ }\PY{n}{amount}\PY{p}{)}\PY{+w}{ }\PY{p}{\PYZob{}}
\PY{+w}{        }\PY{n}{Account}\PY{+w}{ }\PY{n}{first}\PY{+w}{ }\PY{o}{=}\PY{+w}{ }\PY{n}{from}\PY{p}{.}\PY{n+na}{id}\PY{p}{(}\PY{p}{)}\PY{p}{.}\PY{n+na}{compareTo}\PY{p}{(}\PY{n}{to}\PY{p}{.}\PY{n+na}{id}\PY{p}{(}\PY{p}{)}\PY{p}{)}\PY{+w}{ }\PY{o}{\PYZlt{}}\PY{+w}{ }\PY{l+m+mi}{0}\PY{+w}{ }\PY{o}{?}\PY{+w}{ }\PY{n}{from}\PY{+w}{ }\PY{p}{:}\PY{+w}{ }\PY{n}{to}\PY{p}{;}
\PY{+w}{        }\PY{n}{Account}\PY{+w}{ }\PY{n}{second}\PY{+w}{ }\PY{o}{=}\PY{+w}{ }\PY{n}{first}\PY{+w}{ }\PY{o}{=}\PY{o}{=}\PY{+w}{ }\PY{n}{from}\PY{+w}{ }\PY{o}{?}\PY{+w}{ }\PY{n}{to}\PY{+w}{ }\PY{p}{:}\PY{+w}{ }\PY{n}{from}\PY{p}{;}
\PY{+w}{        }\PY{k+kd}{synchronized}\PY{+w}{ }\PY{p}{(}\PY{n}{first}\PY{p}{)}\PY{+w}{ }\PY{p}{\PYZob{}}
\PY{+w}{            }\PY{k+kd}{synchronized}\PY{+w}{ }\PY{p}{(}\PY{n}{second}\PY{p}{)}\PY{+w}{ }\PY{p}{\PYZob{}}
\PY{+w}{                }\PY{n}{from}\PY{p}{.}\PY{n+na}{debit}\PY{p}{(}\PY{n}{amount}\PY{p}{)}\PY{p}{;}
\PY{+w}{                }\PY{n}{to}\PY{p}{.}\PY{n+na}{credit}\PY{p}{(}\PY{n}{amount}\PY{p}{)}\PY{p}{;}
\PY{+w}{            }\PY{p}{\PYZcb{}}
\PY{+w}{        }\PY{p}{\PYZcb{}}
\PY{+w}{    }\PY{p}{\PYZcb{}}
\PY{p}{\PYZcb{}}
\end{Verbatim}
\end{codebox}

\subsection[CompletableFuture Pitfalls]{\code{CompletableFuture} Pitfalls}
\label{h:22:completablefuture-pitfalls}
\begin{codebox}
\begin{Verbatim}[commandchars=\\\{\}]
\PY{c+c1}{// Flagged: exception in the async stage is silently lost; nothing observes the completed future}
\PY{n}{CompletableFuture}\PY{p}{.}\PY{n+na}{supplyAsync}\PY{p}{(}\PY{p}{(}\PY{p}{)}\PY{+w}{ }\PY{o}{\PYZhy{}}\PY{o}{\PYZgt{}}\PY{+w}{ }\PY{n}{paymentGateway}\PY{p}{.}\PY{n+na}{charge}\PY{p}{(}\PY{n}{request}\PY{p}{)}\PY{p}{)}
\PY{+w}{        }\PY{p}{.}\PY{n+na}{thenAccept}\PY{p}{(}\PY{n}{result}\PY{+w}{ }\PY{o}{\PYZhy{}}\PY{o}{\PYZgt{}}\PY{+w}{ }\PY{n}{ledger}\PY{p}{.}\PY{n+na}{record}\PY{p}{(}\PY{n}{result}\PY{p}{)}\PY{p}{)}\PY{p}{;}
\PY{c+c1}{// if charge() or record() throws, this failure disappears — nobody calls get()/join(),}
\PY{c+c1}{// and there\PYZsq{}s no .exceptionally/.handle to observe it}
\end{Verbatim}
\end{codebox}

\begin{codebox}
\begin{Verbatim}[commandchars=\\\{\}]
\PY{c+c1}{// Fixed: handle the failure explicitly}
\PY{n}{CompletableFuture}\PY{p}{.}\PY{n+na}{supplyAsync}\PY{p}{(}\PY{p}{(}\PY{p}{)}\PY{+w}{ }\PY{o}{\PYZhy{}}\PY{o}{\PYZgt{}}\PY{+w}{ }\PY{n}{paymentGateway}\PY{p}{.}\PY{n+na}{charge}\PY{p}{(}\PY{n}{request}\PY{p}{)}\PY{p}{)}
\PY{+w}{        }\PY{p}{.}\PY{n+na}{thenAccept}\PY{p}{(}\PY{n}{result}\PY{+w}{ }\PY{o}{\PYZhy{}}\PY{o}{\PYZgt{}}\PY{+w}{ }\PY{n}{ledger}\PY{p}{.}\PY{n+na}{record}\PY{p}{(}\PY{n}{result}\PY{p}{)}\PY{p}{)}
\PY{+w}{        }\PY{p}{.}\PY{n+na}{exceptionally}\PY{p}{(}\PY{n}{ex}\PY{+w}{ }\PY{o}{\PYZhy{}}\PY{o}{\PYZgt{}}\PY{+w}{ }\PY{p}{\PYZob{}}
\PY{+w}{            }\PY{n}{log}\PY{p}{.}\PY{n+na}{error}\PY{p}{(}\PY{l+s}{\PYZdq{}}\PY{l+s}{Async charge/record pipeline failed}\PY{l+s}{\PYZdq{}}\PY{p}{,}\PY{+w}{ }\PY{n}{ex}\PY{p}{)}\PY{p}{;}
\PY{+w}{            }\PY{n}{alerting}\PY{p}{.}\PY{n+na}{pageOnCall}\PY{p}{(}\PY{l+s}{\PYZdq{}}\PY{l+s}{Payment pipeline failure: }\PY{l+s}{\PYZdq{}}\PY{+w}{ }\PY{o}{+}\PY{+w}{ }\PY{n}{ex}\PY{p}{.}\PY{n+na}{getMessage}\PY{p}{(}\PY{p}{)}\PY{p}{)}\PY{p}{;}
\PY{+w}{            }\PY{k}{return}\PY{+w}{ }\PY{k+kc}{null}\PY{p}{;}
\PY{+w}{        }\PY{p}{\PYZcb{}}\PY{p}{)}\PY{p}{;}
\end{Verbatim}
\end{codebox}

\subsection[Virtual Thread Pinning]{Virtual Thread Pinning}
\label{h:22:virtual-thread-pinning}
\begin{codebox}
\begin{Verbatim}[commandchars=\\\{\}]
\PY{c+c1}{// Flagged: synchronized block around a blocking call, run on a virtual thread — pins the carrier platform thread}
\PY{k+kd}{public}\PY{+w}{ }\PY{k+kd}{synchronized}\PY{+w}{ }\PY{n}{String}\PY{+w}{ }\PY{n+nf}{fetchAndCache}\PY{p}{(}\PY{n}{String}\PY{+w}{ }\PY{n}{key}\PY{p}{)}\PY{+w}{ }\PY{p}{\PYZob{}}
\PY{+w}{    }\PY{n}{String}\PY{+w}{ }\PY{n}{cached}\PY{+w}{ }\PY{o}{=}\PY{+w}{ }\PY{n}{cache}\PY{p}{.}\PY{n+na}{get}\PY{p}{(}\PY{n}{key}\PY{p}{)}\PY{p}{;}
\PY{+w}{    }\PY{k}{if}\PY{+w}{ }\PY{p}{(}\PY{n}{cached}\PY{+w}{ }\PY{o}{!}\PY{o}{=}\PY{+w}{ }\PY{k+kc}{null}\PY{p}{)}\PY{+w}{ }\PY{k}{return}\PY{+w}{ }\PY{n}{cached}\PY{p}{;}
\PY{+w}{    }\PY{n}{String}\PY{+w}{ }\PY{n}{value}\PY{+w}{ }\PY{o}{=}\PY{+w}{ }\PY{n}{blockingHttpCall}\PY{p}{(}\PY{n}{key}\PY{p}{)}\PY{p}{;}\PY{+w}{ }\PY{c+c1}{// blocks while holding the monitor}
\PY{+w}{    }\PY{n}{cache}\PY{p}{.}\PY{n+na}{put}\PY{p}{(}\PY{n}{key}\PY{p}{,}\PY{+w}{ }\PY{n}{value}\PY{p}{)}\PY{p}{;}
\PY{+w}{    }\PY{k}{return}\PY{+w}{ }\PY{n}{value}\PY{p}{;}
\PY{p}{\PYZcb{}}
\end{Verbatim}
\end{codebox}

A \code{synchronized} block that’s holding its monitor while making a blocking call prevents the virtual thread from unmounting from its carrier platform thread during that block — under load, this can exhaust the small pool of carrier threads and stall unrelated virtual threads, defeating the scalability virtual threads exist to provide.

\begin{codebox}
\begin{Verbatim}[commandchars=\\\{\}]
\PY{c+c1}{// Fixed: use a java.util.concurrent.locks.ReentrantLock, which does allow unmounting, or restructure}
\PY{c+c1}{// to avoid holding any lock across the blocking call in the first place}
\PY{k+kd}{private}\PY{+w}{ }\PY{k+kd}{final}\PY{+w}{ }\PY{n}{ReentrantLock}\PY{+w}{ }\PY{n}{lock}\PY{+w}{ }\PY{o}{=}\PY{+w}{ }\PY{k}{new}\PY{+w}{ }\PY{n}{ReentrantLock}\PY{p}{(}\PY{p}{)}\PY{p}{;}

\PY{k+kd}{public}\PY{+w}{ }\PY{n}{String}\PY{+w}{ }\PY{n+nf}{fetchAndCache}\PY{p}{(}\PY{n}{String}\PY{+w}{ }\PY{n}{key}\PY{p}{)}\PY{+w}{ }\PY{p}{\PYZob{}}
\PY{+w}{    }\PY{n}{String}\PY{+w}{ }\PY{n}{cached}\PY{+w}{ }\PY{o}{=}\PY{+w}{ }\PY{n}{cache}\PY{p}{.}\PY{n+na}{get}\PY{p}{(}\PY{n}{key}\PY{p}{)}\PY{p}{;}
\PY{+w}{    }\PY{k}{if}\PY{+w}{ }\PY{p}{(}\PY{n}{cached}\PY{+w}{ }\PY{o}{!}\PY{o}{=}\PY{+w}{ }\PY{k+kc}{null}\PY{p}{)}\PY{+w}{ }\PY{k}{return}\PY{+w}{ }\PY{n}{cached}\PY{p}{;}
\PY{+w}{    }\PY{n}{String}\PY{+w}{ }\PY{n}{value}\PY{+w}{ }\PY{o}{=}\PY{+w}{ }\PY{n}{blockingHttpCall}\PY{p}{(}\PY{n}{key}\PY{p}{)}\PY{p}{;}\PY{+w}{ }\PY{c+c1}{// no lock held here at all}
\PY{+w}{    }\PY{n}{lock}\PY{p}{.}\PY{n+na}{lock}\PY{p}{(}\PY{p}{)}\PY{p}{;}
\PY{+w}{    }\PY{k}{try}\PY{+w}{ }\PY{p}{\PYZob{}}
\PY{+w}{        }\PY{n}{cache}\PY{p}{.}\PY{n+na}{putIfAbsent}\PY{p}{(}\PY{n}{key}\PY{p}{,}\PY{+w}{ }\PY{n}{value}\PY{p}{)}\PY{p}{;}
\PY{+w}{    }\PY{p}{\PYZcb{}}\PY{+w}{ }\PY{k}{finally}\PY{+w}{ }\PY{p}{\PYZob{}}
\PY{+w}{        }\PY{n}{lock}\PY{p}{.}\PY{n+na}{unlock}\PY{p}{(}\PY{p}{)}\PY{p}{;}
\PY{+w}{    }\PY{p}{\PYZcb{}}
\PY{+w}{    }\PY{k}{return}\PY{+w}{ }\PY{n}{cache}\PY{p}{.}\PY{n+na}{get}\PY{p}{(}\PY{n}{key}\PY{p}{)}\PY{p}{;}
\PY{p}{\PYZcb{}}
\end{Verbatim}
\end{codebox}

See \hyperref[chap:15]{Chapter~15, \emph{Modern Concurrency (Java 21+)}} for virtual threads and structured concurrency in depth, and \hyperref[chap:14]{Chapter~14, \emph{Concurrency (Advanced)}} for the executor/\code{CompletableFuture} framework this section leans on.

\section[Modern Java Feature Usage]{Modern Java Feature Usage}
\label{h:22:modern-java-feature-usage}
Modern syntax is not automatically good style — review both directions: missed opportunities to simplify, \emph{and} overuse that hurts readability. Background on each feature is in \hyperref[chap:3]{Chapter~3, \emph{Modern Java Language Features}}; this section is specifically about judging usage.

\subsection[Records — Good Fit]{Records — Good Fit}
\label{h:22:records--good-fit}
\begin{codebox}
\begin{Verbatim}[commandchars=\\\{\}]
\PY{c+c1}{// Flagged: a hand\PYZhy{}rolled immutable value class with all the equals/hashCode/toString boilerplate}
\PY{k+kd}{public}\PY{+w}{ }\PY{k+kd}{final}\PY{+w}{ }\PY{k+kd}{class} \PY{n+nc}{Point}\PY{+w}{ }\PY{p}{\PYZob{}}
\PY{+w}{    }\PY{k+kd}{private}\PY{+w}{ }\PY{k+kd}{final}\PY{+w}{ }\PY{k+kt}{int}\PY{+w}{ }\PY{n}{x}\PY{p}{,}\PY{+w}{ }\PY{n}{y}\PY{p}{;}
\PY{+w}{    }\PY{k+kd}{public}\PY{+w}{ }\PY{n+nf}{Point}\PY{p}{(}\PY{k+kt}{int}\PY{+w}{ }\PY{n}{x}\PY{p}{,}\PY{+w}{ }\PY{k+kt}{int}\PY{+w}{ }\PY{n}{y}\PY{p}{)}\PY{+w}{ }\PY{p}{\PYZob{}}\PY{+w}{ }\PY{k}{this}\PY{p}{.}\PY{n+na}{x}\PY{+w}{ }\PY{o}{=}\PY{+w}{ }\PY{n}{x}\PY{p}{;}\PY{+w}{ }\PY{k}{this}\PY{p}{.}\PY{n+na}{y}\PY{+w}{ }\PY{o}{=}\PY{+w}{ }\PY{n}{y}\PY{p}{;}\PY{+w}{ }\PY{p}{\PYZcb{}}
\PY{+w}{    }\PY{k+kd}{public}\PY{+w}{ }\PY{k+kt}{int}\PY{+w}{ }\PY{n+nf}{getX}\PY{p}{(}\PY{p}{)}\PY{+w}{ }\PY{p}{\PYZob{}}\PY{+w}{ }\PY{k}{return}\PY{+w}{ }\PY{n}{x}\PY{p}{;}\PY{+w}{ }\PY{p}{\PYZcb{}}
\PY{+w}{    }\PY{k+kd}{public}\PY{+w}{ }\PY{k+kt}{int}\PY{+w}{ }\PY{n+nf}{getY}\PY{p}{(}\PY{p}{)}\PY{+w}{ }\PY{p}{\PYZob{}}\PY{+w}{ }\PY{k}{return}\PY{+w}{ }\PY{n}{y}\PY{p}{;}\PY{+w}{ }\PY{p}{\PYZcb{}}
\PY{+w}{    }\PY{n+nd}{@Override}\PY{+w}{ }\PY{k+kd}{public}\PY{+w}{ }\PY{k+kt}{boolean}\PY{+w}{ }\PY{n+nf}{equals}\PY{p}{(}\PY{n}{Object}\PY{+w}{ }\PY{n}{o}\PY{p}{)}\PY{+w}{ }\PY{p}{\PYZob{}}\PY{+w}{ }\PY{c+cm}{/* ... */}\PY{+w}{ }\PY{k}{return}\PY{+w}{ }\PY{k+kc}{false}\PY{p}{;}\PY{+w}{ }\PY{p}{\PYZcb{}}
\PY{+w}{    }\PY{n+nd}{@Override}\PY{+w}{ }\PY{k+kd}{public}\PY{+w}{ }\PY{k+kt}{int}\PY{+w}{ }\PY{n+nf}{hashCode}\PY{p}{(}\PY{p}{)}\PY{+w}{ }\PY{p}{\PYZob{}}\PY{+w}{ }\PY{c+cm}{/* ... */}\PY{+w}{ }\PY{k}{return}\PY{+w}{ }\PY{l+m+mi}{0}\PY{p}{;}\PY{+w}{ }\PY{p}{\PYZcb{}}
\PY{+w}{    }\PY{n+nd}{@Override}\PY{+w}{ }\PY{k+kd}{public}\PY{+w}{ }\PY{n}{String}\PY{+w}{ }\PY{n+nf}{toString}\PY{p}{(}\PY{p}{)}\PY{+w}{ }\PY{p}{\PYZob{}}\PY{+w}{ }\PY{k}{return}\PY{+w}{ }\PY{l+s}{\PYZdq{}}\PY{l+s}{Point(}\PY{l+s}{\PYZdq{}}\PY{+w}{ }\PY{o}{+}\PY{+w}{ }\PY{n}{x}\PY{+w}{ }\PY{o}{+}\PY{+w}{ }\PY{l+s}{\PYZdq{}}\PY{l+s}{, }\PY{l+s}{\PYZdq{}}\PY{+w}{ }\PY{o}{+}\PY{+w}{ }\PY{n}{y}\PY{+w}{ }\PY{o}{+}\PY{+w}{ }\PY{l+s}{\PYZdq{}}\PY{l+s}{)}\PY{l+s}{\PYZdq{}}\PY{p}{;}\PY{+w}{ }\PY{p}{\PYZcb{}}
\PY{p}{\PYZcb{}}
\end{Verbatim}
\end{codebox}

\begin{codebox}
\begin{Verbatim}[commandchars=\\\{\}]
\PY{c+c1}{// Fixed: a record generates the constructor, accessors, equals/hashCode/toString for free}
\PY{k+kd}{public}\PY{+w}{ }\PY{k+kd}{record} \PY{n+nc}{Point}\PY{p}{(}\PY{k+kt}{int}\PY{+w}{ }\PY{n}{x}\PY{p}{,}\PY{+w}{ }\PY{k+kt}{int}\PY{+w}{ }\PY{n}{y}\PY{p}{)}\PY{+w}{ }\PY{p}{\PYZob{}}\PY{p}{\PYZcb{}}
\end{Verbatim}
\end{codebox}

\subsection[Records — Bad Fit]{Records — Bad Fit}
\label{h:22:records--bad-fit}
\begin{codebox}
\begin{Verbatim}[commandchars=\\\{\}]
\PY{c+c1}{// Flagged: a record with a mutable component defeats the entire point of using a record}
\PY{k+kd}{public}\PY{+w}{ }\PY{k+kd}{record} \PY{n+nc}{OrderBatch}\PY{p}{(}\PY{n}{List}\PY{o}{\PYZlt{}}\PY{n}{Order}\PY{o}{\PYZgt{}}\PY{+w}{ }\PY{n}{orders}\PY{p}{)}\PY{+w}{ }\PY{p}{\PYZob{}}\PY{p}{\PYZcb{}}\PY{+w}{ }\PY{c+c1}{// caller can still mutate the passed\PYZhy{}in list from outside}
\end{Verbatim}
\end{codebox}

\begin{codebox}
\begin{Verbatim}[commandchars=\\\{\}]
\PY{c+c1}{// Fixed: defensively copy in the compact constructor, and expose an unmodifiable view}
\PY{k+kd}{public}\PY{+w}{ }\PY{k+kd}{record} \PY{n+nc}{OrderBatch}\PY{p}{(}\PY{n}{List}\PY{o}{\PYZlt{}}\PY{n}{Order}\PY{o}{\PYZgt{}}\PY{+w}{ }\PY{n}{orders}\PY{p}{)}\PY{+w}{ }\PY{p}{\PYZob{}}
\PY{+w}{    }\PY{k+kd}{public}\PY{+w}{ }\PY{n+nf}{OrderBatch}\PY{p}{(}\PY{n}{List}\PY{o}{\PYZlt{}}\PY{n}{Order}\PY{o}{\PYZgt{}}\PY{+w}{ }\PY{n}{orders}\PY{p}{)}\PY{+w}{ }\PY{p}{\PYZob{}}
\PY{+w}{        }\PY{k}{this}\PY{p}{.}\PY{n+na}{orders}\PY{+w}{ }\PY{o}{=}\PY{+w}{ }\PY{n}{List}\PY{p}{.}\PY{n+na}{copyOf}\PY{p}{(}\PY{n}{orders}\PY{p}{)}\PY{p}{;}\PY{+w}{ }\PY{c+c1}{// immutable copy, breaks the aliasing hazard}
\PY{+w}{    }\PY{p}{\PYZcb{}}
\PY{p}{\PYZcb{}}
\end{Verbatim}
\end{codebox}

A record whose only mutable-looking component is an array or a mutable collection is a very common planted bug — the record’s own fields are \code{final}, but that says nothing about whether the \emph{referenced} object can still be mutated by the caller.

\subsection[Pattern Matching vs. Polymorphism]{Pattern Matching vs. Polymorphism}
\label{h:22:pattern-matching-vs-polymorphism}
\begin{codebox}
\begin{Verbatim}[commandchars=\\\{\}]
\PY{c+c1}{// Flagged: pattern\PYZhy{}matching switch reimplementing what virtual dispatch already does better}
\PY{k+kd}{public}\PY{+w}{ }\PY{k+kt}{double}\PY{+w}{ }\PY{n+nf}{area}\PY{p}{(}\PY{n}{Shape}\PY{+w}{ }\PY{n}{shape}\PY{p}{)}\PY{+w}{ }\PY{p}{\PYZob{}}
\PY{+w}{    }\PY{k}{return}\PY{+w}{ }\PY{k}{switch}\PY{+w}{ }\PY{p}{(}\PY{n}{shape}\PY{p}{)}\PY{+w}{ }\PY{p}{\PYZob{}}
\PY{+w}{        }\PY{k}{case}\PY{+w}{ }\PY{n}{Circle}\PY{+w}{ }\PY{n}{c}\PY{+w}{ }\PY{o}{\PYZhy{}}\PY{o}{\PYZgt{}}\PY{+w}{ }\PY{n}{Math}\PY{p}{.}\PY{n+na}{PI}\PY{+w}{ }\PY{o}{*}\PY{+w}{ }\PY{n}{c}\PY{p}{.}\PY{n+na}{radius}\PY{p}{(}\PY{p}{)}\PY{+w}{ }\PY{o}{*}\PY{+w}{ }\PY{n}{c}\PY{p}{.}\PY{n+na}{radius}\PY{p}{(}\PY{p}{)}\PY{p}{;}
\PY{+w}{        }\PY{k}{case}\PY{+w}{ }\PY{n}{Square}\PY{+w}{ }\PY{n}{s}\PY{+w}{ }\PY{o}{\PYZhy{}}\PY{o}{\PYZgt{}}\PY{+w}{ }\PY{n}{s}\PY{p}{.}\PY{n+na}{side}\PY{p}{(}\PY{p}{)}\PY{+w}{ }\PY{o}{*}\PY{+w}{ }\PY{n}{s}\PY{p}{.}\PY{n+na}{side}\PY{p}{(}\PY{p}{)}\PY{p}{;}
\PY{+w}{        }\PY{k}{case}\PY{+w}{ }\PY{n}{Rectangle}\PY{+w}{ }\PY{n}{r}\PY{+w}{ }\PY{o}{\PYZhy{}}\PY{o}{\PYZgt{}}\PY{+w}{ }\PY{n}{r}\PY{p}{.}\PY{n+na}{width}\PY{p}{(}\PY{p}{)}\PY{+w}{ }\PY{o}{*}\PY{+w}{ }\PY{n}{r}\PY{p}{.}\PY{n+na}{height}\PY{p}{(}\PY{p}{)}\PY{p}{;}
\PY{+w}{        }\PY{k}{default}\PY{+w}{ }\PY{o}{\PYZhy{}}\PY{o}{\PYZgt{}}\PY{+w}{ }\PY{k}{throw}\PY{+w}{ }\PY{k}{new}\PY{+w}{ }\PY{n}{IllegalStateException}\PY{p}{(}\PY{p}{)}\PY{p}{;}
\PY{+w}{    }\PY{p}{\PYZcb{}}\PY{p}{;}
\PY{p}{\PYZcb{}}
\end{Verbatim}
\end{codebox}

\begin{codebox}
\begin{Verbatim}[commandchars=\\\{\}]
\PY{c+c1}{// Fixed: if this switch would need a new case every time a new Shape type is added,}
\PY{c+c1}{// and Shape is a type you own, plain polymorphism is simpler and can\PYZsq{}t forget a case}
\PY{k+kd}{sealed}\PY{+w}{ }\PY{k+kd}{interface} \PY{n+nc}{Shape}\PY{+w}{ }\PY{p}{\PYZob{}}\PY{+w}{ }\PY{k+kt}{double}\PY{+w}{ }\PY{n+nf}{area}\PY{p}{(}\PY{p}{)}\PY{p}{;}\PY{+w}{ }\PY{p}{\PYZcb{}}
\PY{k+kd}{record} \PY{n+nc}{Circle}\PY{p}{(}\PY{k+kt}{double}\PY{+w}{ }\PY{n}{radius}\PY{p}{)}\PY{+w}{ }\PY{k+kd}{implements}\PY{+w}{ }\PY{n}{Shape}\PY{+w}{ }\PY{p}{\PYZob{}}\PY{+w}{ }\PY{k+kd}{public}\PY{+w}{ }\PY{k+kt}{double}\PY{+w}{ }\PY{n+nf}{area}\PY{p}{(}\PY{p}{)}\PY{+w}{ }\PY{p}{\PYZob{}}\PY{+w}{ }\PY{k}{return}\PY{+w}{ }\PY{n}{Math}\PY{p}{.}\PY{n+na}{PI}\PY{+w}{ }\PY{o}{*}\PY{+w}{ }\PY{n}{radius}\PY{+w}{ }\PY{o}{*}\PY{+w}{ }\PY{n}{radius}\PY{p}{;}\PY{+w}{ }\PY{p}{\PYZcb{}}\PY{+w}{ }\PY{p}{\PYZcb{}}
\PY{k+kd}{record} \PY{n+nc}{Square}\PY{p}{(}\PY{k+kt}{double}\PY{+w}{ }\PY{n}{side}\PY{p}{)}\PY{+w}{ }\PY{k+kd}{implements}\PY{+w}{ }\PY{n}{Shape}\PY{+w}{ }\PY{p}{\PYZob{}}\PY{+w}{ }\PY{k+kd}{public}\PY{+w}{ }\PY{k+kt}{double}\PY{+w}{ }\PY{n+nf}{area}\PY{p}{(}\PY{p}{)}\PY{+w}{ }\PY{p}{\PYZob{}}\PY{+w}{ }\PY{k}{return}\PY{+w}{ }\PY{n}{side}\PY{+w}{ }\PY{o}{*}\PY{+w}{ }\PY{n}{side}\PY{p}{;}\PY{+w}{ }\PY{p}{\PYZcb{}}\PY{+w}{ }\PY{p}{\PYZcb{}}
\PY{k+kd}{record} \PY{n+nc}{Rectangle}\PY{p}{(}\PY{k+kt}{double}\PY{+w}{ }\PY{n}{width}\PY{p}{,}\PY{+w}{ }\PY{k+kt}{double}\PY{+w}{ }\PY{n}{height}\PY{p}{)}\PY{+w}{ }\PY{k+kd}{implements}\PY{+w}{ }\PY{n}{Shape}\PY{+w}{ }\PY{p}{\PYZob{}}\PY{+w}{ }\PY{k+kd}{public}\PY{+w}{ }\PY{k+kt}{double}\PY{+w}{ }\PY{n+nf}{area}\PY{p}{(}\PY{p}{)}\PY{+w}{ }\PY{p}{\PYZob{}}\PY{+w}{ }\PY{k}{return}\PY{+w}{ }\PY{n}{width}\PY{+w}{ }\PY{o}{*}\PY{+w}{ }\PY{n}{height}\PY{p}{;}\PY{+w}{ }\PY{p}{\PYZcb{}}\PY{+w}{ }\PY{p}{\PYZcb{}}
\end{Verbatim}
\end{codebox}

Pattern matching earns its keep when the operation is genuinely external to the type (e.g., a serializer that needs to know about every type but the types themselves shouldn’t know about serialization), or when combined with a \code{sealed} hierarchy so the compiler enforces exhaustiveness. It’s a smell when it’s just reinventing a method that \code{area()} on the interface would express more simply — and a \code{sealed} hierarchy makes the compiler catch a missing case either way, so exhaustiveness isn’t a deciding factor between the two styles.

\subsection[var — Good and Bad]{\code{var} — Good and Bad}
\label{h:22:var--good-and-bad}
\begin{codebox}
\begin{Verbatim}[commandchars=\\\{\}]
\PY{c+c1}{// Flagged: var hides the type where the type is exactly the information the reader needs}
\PY{k+kd}{var}\PY{+w}{ }\PY{n}{result}\PY{+w}{ }\PY{o}{=}\PY{+w}{ }\PY{n}{process}\PY{p}{(}\PY{n}{data}\PY{p}{)}\PY{p}{;}\PY{+w}{ }\PY{c+c1}{// what does process() return? Not obvious from this line alone}
\end{Verbatim}
\end{codebox}

\begin{codebox}
\begin{Verbatim}[commandchars=\\\{\}]
\PY{c+c1}{// Fixed: keep the explicit type when it\PYZsq{}s not obvious from the right\PYZhy{}hand side}
\PY{n}{ValidationResult}\PY{+w}{ }\PY{n}{result}\PY{+w}{ }\PY{o}{=}\PY{+w}{ }\PY{n}{process}\PY{p}{(}\PY{n}{data}\PY{p}{)}\PY{p}{;}
\end{Verbatim}
\end{codebox}

\begin{codebox}
\begin{Verbatim}[commandchars=\\\{\}]
\PY{c+c1}{// Good use of var: the type is already obvious from the constructor call — no information lost}
\PY{k+kd}{var}\PY{+w}{ }\PY{n}{orders}\PY{+w}{ }\PY{o}{=}\PY{+w}{ }\PY{k}{new}\PY{+w}{ }\PY{n}{ArrayList}\PY{o}{\PYZlt{}}\PY{n}{Order}\PY{o}{\PYZgt{}}\PY{p}{(}\PY{p}{)}\PY{p}{;}
\PY{k+kd}{var}\PY{+w}{ }\PY{n}{connection}\PY{+w}{ }\PY{o}{=}\PY{+w}{ }\PY{n}{DriverManager}\PY{p}{.}\PY{n+na}{getConnection}\PY{p}{(}\PY{n}{url}\PY{p}{)}\PY{p}{;}
\end{Verbatim}
\end{codebox}

The rule of thumb: \code{var} is good when the type is redundant with something already visible on the same line (a constructor call, a well-named factory method); it’s a readability regression when it hides the type of something returned by a method whose return type isn’t obvious from its name.

\subsection[Streams — Good and Bad]{Streams — Good and Bad}
\label{h:22:streams--good-and-bad}
\begin{codebox}
\begin{Verbatim}[commandchars=\\\{\}]
\PY{c+c1}{// Flagged: a stream pipeline that\PYZsq{}s harder to read than the equivalent loop, with a side\PYZhy{}effecting peek}
\PY{n}{List}\PY{o}{\PYZlt{}}\PY{n}{String}\PY{o}{\PYZgt{}}\PY{+w}{ }\PY{n}{ids}\PY{+w}{ }\PY{o}{=}\PY{+w}{ }\PY{n}{orders}\PY{p}{.}\PY{n+na}{stream}\PY{p}{(}\PY{p}{)}
\PY{+w}{        }\PY{p}{.}\PY{n+na}{peek}\PY{p}{(}\PY{n}{o}\PY{+w}{ }\PY{o}{\PYZhy{}}\PY{o}{\PYZgt{}}\PY{+w}{ }\PY{n}{auditLog}\PY{p}{.}\PY{n+na}{add}\PY{p}{(}\PY{n}{o}\PY{p}{.}\PY{n+na}{id}\PY{p}{(}\PY{p}{)}\PY{p}{)}\PY{p}{)}\PY{+w}{ }\PY{c+c1}{// side effect hidden inside peek — easy to miss, and peek isn\PYZsq{}t guaranteed to run for every element in all pipelines}
\PY{+w}{        }\PY{p}{.}\PY{n+na}{filter}\PY{p}{(}\PY{n}{o}\PY{+w}{ }\PY{o}{\PYZhy{}}\PY{o}{\PYZgt{}}\PY{+w}{ }\PY{n}{o}\PY{p}{.}\PY{n+na}{total}\PY{p}{(}\PY{p}{)}\PY{p}{.}\PY{n+na}{compareTo}\PY{p}{(}\PY{n}{BigDecimal}\PY{p}{.}\PY{n+na}{TEN}\PY{p}{)}\PY{+w}{ }\PY{o}{\PYZgt{}}\PY{+w}{ }\PY{l+m+mi}{0}\PY{p}{)}
\PY{+w}{        }\PY{p}{.}\PY{n+na}{map}\PY{p}{(}\PY{n}{Order}\PY{p}{:}\PY{p}{:}\PY{n}{id}\PY{p}{)}
\PY{+w}{        }\PY{p}{.}\PY{n+na}{collect}\PY{p}{(}\PY{n}{Collectors}\PY{p}{.}\PY{n+na}{toList}\PY{p}{(}\PY{p}{)}\PY{p}{)}\PY{p}{;}
\end{Verbatim}
\end{codebox}

\begin{codebox}
\begin{Verbatim}[commandchars=\\\{\}]
\PY{c+c1}{// Fixed: keep the side effect explicit and separate from the pipeline; let the stream just transform data}
\PY{n}{List}\PY{o}{\PYZlt{}}\PY{n}{String}\PY{o}{\PYZgt{}}\PY{+w}{ }\PY{n}{ids}\PY{+w}{ }\PY{o}{=}\PY{+w}{ }\PY{n}{orders}\PY{p}{.}\PY{n+na}{stream}\PY{p}{(}\PY{p}{)}
\PY{+w}{        }\PY{p}{.}\PY{n+na}{filter}\PY{p}{(}\PY{n}{o}\PY{+w}{ }\PY{o}{\PYZhy{}}\PY{o}{\PYZgt{}}\PY{+w}{ }\PY{n}{o}\PY{p}{.}\PY{n+na}{total}\PY{p}{(}\PY{p}{)}\PY{p}{.}\PY{n+na}{compareTo}\PY{p}{(}\PY{n}{BigDecimal}\PY{p}{.}\PY{n+na}{TEN}\PY{p}{)}\PY{+w}{ }\PY{o}{\PYZgt{}}\PY{+w}{ }\PY{l+m+mi}{0}\PY{p}{)}
\PY{+w}{        }\PY{p}{.}\PY{n+na}{map}\PY{p}{(}\PY{n}{Order}\PY{p}{:}\PY{p}{:}\PY{n}{id}\PY{p}{)}
\PY{+w}{        }\PY{p}{.}\PY{n+na}{toList}\PY{p}{(}\PY{p}{)}\PY{p}{;}
\PY{n}{ids}\PY{p}{.}\PY{n+na}{forEach}\PY{p}{(}\PY{n}{auditLog}\PY{p}{:}\PY{p}{:}\PY{n}{add}\PY{p}{)}\PY{p}{;}\PY{+w}{ }\PY{c+c1}{// or better, log alongside the filtering with a plain loop if the two concerns are related}
\end{Verbatim}
\end{codebox}

Streams earn their keep for a clear filter/map/collect pipeline over a collection. They’re a smell when: the lambda has side effects (mutating something outside the stream — see Scenario 12), the pipeline is nested or nearly unreadable compared to an equivalent loop, or exceptions need awkward wrapping because the functional interfaces involved don’t declare \code{throws} (see \hyperref[chap:6]{Chapter~6, \emph{Exception Handling}} and \hyperref[chap:8]{Chapter~8, \emph{Functional Programming}}).

\subsection[Switch Expressions and Text Blocks]{Switch Expressions and Text Blocks}
\label{h:22:switch-expressions-and-text-blocks}
\begin{codebox}
\begin{Verbatim}[commandchars=\\\{\}]
\PY{c+c1}{// Flagged: old\PYZhy{}style switch statement with fall\PYZhy{}through risk and repeated assignment}
\PY{n}{String}\PY{+w}{ }\PY{n}{tier}\PY{p}{;}
\PY{k}{switch}\PY{+w}{ }\PY{p}{(}\PY{n}{score}\PY{p}{)}\PY{+w}{ }\PY{p}{\PYZob{}}
\PY{+w}{    }\PY{k}{case}\PY{+w}{ }\PY{l+m+mi}{0}\PY{p}{:}\PY{+w}{ }\PY{k}{case}\PY{+w}{ }\PY{l+m+mi}{1}\PY{p}{:}\PY{+w}{ }\PY{k}{case}\PY{+w}{ }\PY{l+m+mi}{2}\PY{p}{:}\PY{+w}{ }\PY{n}{tier}\PY{+w}{ }\PY{o}{=}\PY{+w}{ }\PY{l+s}{\PYZdq{}}\PY{l+s}{BRONZE}\PY{l+s}{\PYZdq{}}\PY{p}{;}\PY{+w}{ }\PY{k}{break}\PY{p}{;}
\PY{+w}{    }\PY{k}{case}\PY{+w}{ }\PY{l+m+mi}{3}\PY{p}{:}\PY{+w}{ }\PY{k}{case}\PY{+w}{ }\PY{l+m+mi}{4}\PY{p}{:}\PY{+w}{ }\PY{n}{tier}\PY{+w}{ }\PY{o}{=}\PY{+w}{ }\PY{l+s}{\PYZdq{}}\PY{l+s}{SILVER}\PY{l+s}{\PYZdq{}}\PY{p}{;}\PY{+w}{ }\PY{k}{break}\PY{p}{;}
\PY{+w}{    }\PY{k}{default}\PY{p}{:}\PY{+w}{ }\PY{n}{tier}\PY{+w}{ }\PY{o}{=}\PY{+w}{ }\PY{l+s}{\PYZdq{}}\PY{l+s}{GOLD}\PY{l+s}{\PYZdq{}}\PY{p}{;}
\PY{p}{\PYZcb{}}
\end{Verbatim}
\end{codebox}

\begin{codebox}
\begin{Verbatim}[commandchars=\\\{\}]
\PY{c+c1}{// Fixed: switch expression — no fall\PYZhy{}through risk, no repeated variable, exhaustiveness is clearer}
\PY{n}{String}\PY{+w}{ }\PY{n}{tier}\PY{+w}{ }\PY{o}{=}\PY{+w}{ }\PY{k}{switch}\PY{+w}{ }\PY{p}{(}\PY{n}{score}\PY{p}{)}\PY{+w}{ }\PY{p}{\PYZob{}}
\PY{+w}{    }\PY{k}{case}\PY{+w}{ }\PY{l+m+mi}{0}\PY{p}{,}\PY{+w}{ }\PY{l+m+mi}{1}\PY{p}{,}\PY{+w}{ }\PY{l+m+mi}{2}\PY{+w}{ }\PY{o}{\PYZhy{}}\PY{o}{\PYZgt{}}\PY{+w}{ }\PY{l+s}{\PYZdq{}}\PY{l+s}{BRONZE}\PY{l+s}{\PYZdq{}}\PY{p}{;}
\PY{+w}{    }\PY{k}{case}\PY{+w}{ }\PY{l+m+mi}{3}\PY{p}{,}\PY{+w}{ }\PY{l+m+mi}{4}\PY{+w}{ }\PY{o}{\PYZhy{}}\PY{o}{\PYZgt{}}\PY{+w}{ }\PY{l+s}{\PYZdq{}}\PY{l+s}{SILVER}\PY{l+s}{\PYZdq{}}\PY{p}{;}
\PY{+w}{    }\PY{k}{default}\PY{+w}{ }\PY{o}{\PYZhy{}}\PY{o}{\PYZgt{}}\PY{+w}{ }\PY{l+s}{\PYZdq{}}\PY{l+s}{GOLD}\PY{l+s}{\PYZdq{}}\PY{p}{;}
\PY{p}{\PYZcb{}}\PY{p}{;}
\end{Verbatim}
\end{codebox}

\begin{codebox}
\begin{Verbatim}[commandchars=\\\{\}]
\PY{c+c1}{// Flagged: multi\PYZhy{}line SQL/JSON built with escaped string concatenation}
\PY{n}{String}\PY{+w}{ }\PY{n}{query}\PY{+w}{ }\PY{o}{=}\PY{+w}{ }\PY{l+s}{\PYZdq{}}\PY{l+s}{SELECT id, name}\PY{l+s}{\PYZbs{}}\PY{l+s}{n}\PY{l+s}{\PYZdq{}}\PY{+w}{ }\PY{o}{+}
\PY{+w}{        }\PY{l+s}{\PYZdq{}}\PY{l+s}{FROM customers}\PY{l+s}{\PYZbs{}}\PY{l+s}{n}\PY{l+s}{\PYZdq{}}\PY{+w}{ }\PY{o}{+}
\PY{+w}{        }\PY{l+s}{\PYZdq{}}\PY{l+s}{WHERE status = \PYZsq{}ACTIVE\PYZsq{}}\PY{l+s}{\PYZbs{}}\PY{l+s}{n}\PY{l+s}{\PYZdq{}}\PY{p}{;}
\end{Verbatim}
\end{codebox}

\begin{codebox}
\begin{Verbatim}[commandchars=\\\{\}]
\PY{c+c1}{// Fixed: a text block reads exactly like the text it represents}
\PY{n}{String}\PY{+w}{ }\PY{n}{query}\PY{+w}{ }\PY{o}{=}\PY{+w}{ }\PY{l+s}{\PYZdq{}\PYZdq{}\PYZdq{}}
\PY{l+s}{        SELECT id, name}
\PY{l+s}{        FROM customers}
\PY{l+s}{        WHERE status = \PYZsq{}ACTIVE\PYZsq{}}
\PY{l+s}{        }\PY{l+s}{\PYZdq{}\PYZdq{}\PYZdq{}}\PY{p}{;}
\end{Verbatim}
\end{codebox}

These two are close to unconditional wins — flag their \emph{absence} as a missed simplification, not just misuse.

\section[Common Code-Review Interview Pitfalls]{Common Code-Review Interview Pitfalls}
\label{h:22:common-code-review-interview-pitfalls}
These are mistakes candidates make \emph{in the interview itself} — separate from the Java-language pitfalls earlier in this chapter. Interviewers see these constantly, and each one costs real points even when the candidate’s technical instincts are otherwise sound.

\begin{enumerate}
\item 
\textbf{Bikeshedding a style nit while a data race sits three lines away.} Why it matters: it signals you can’t triage severity, which is the actual skill being tested — finding bugs is only half the job.

\item 
\textbf{Not asking about the concurrency model before reviewing.} Why it matters: “is this called from multiple threads?” changes almost every finding in a snippet; guessing wrong wastes the whole review and can make you miss the one bug that mattered.

\item 
\textbf{Missing the concurrency bug entirely and only flagging surface-level issues (naming, formatting).} Why it matters: concurrency bugs are usually the actual point of the exercise; if the snippet has a shared mutable \code{HashMap} or an unsynchronized singleton, that has to be the headline finding.

\item 
\textbf{Listing every issue in file order with no prioritization.} Why it matters: a review that treats a missing \code{@\allowbreak{}Override} annotation with the same weight as a resource leak tells the interviewer you don’t understand production risk.

\item 
\textbf{Being rude or accusatory about the (fictional) author’s competence.} Why it matters: real code review is collaborative; interviewers are explicitly evaluating whether they’d want you reviewing a teammate’s PR, not just whether you can spot bugs.

\item 
\textbf{Not verifying assumptions before declaring something broken.} Why it matters: jumping to “this is definitely a bug” about, say, an overload’s resolution without checking the actual types involved can make you confidently state something wrong — worse than not spotting it, because it damages credibility on the things you got right.

\item 
\textbf{Over-engineering the fix — proposing a distributed lock, a message queue, or a rewrite for a five-line bug.} Why it matters: it shows you can’t scope a fix to the actual problem, and real reviewers are wary of suggestions that would balloon the size and risk of a PR.

\item 
\textbf{Not asking about requirements or constraints before assuming the “obvious” fix is right.} Why it matters: what looks like an unbounded cache bug might be intentional if the process is short-lived and restarted daily — the right fix depends on context you don’t have unless you ask.

\item 
\textbf{Silence — reading the code for a long time without narrating your thought process.} Why it matters: interviewers are assessing \emph{how} you think, not just your final answer; silent reading gives them nothing to evaluate and often reads as being stuck.

\item 
\textbf{Fixating on one bug and running out of time to scan the rest of the snippet.} Why it matters: most review snippets have 2-4 issues by design; spending the entire session on the first one you spot means you miss the ones planted specifically to test breadth.

\item 
\textbf{Proposing a fix without explaining the mechanism behind the bug.} Why it matters: “just make it a \code{ConcurrentHashMap}” without explaining \emph{why} the current code races sounds like memorized pattern-matching rather than understanding — interviewers probe follow-ups precisely to distinguish the two.

\item 
\textbf{Ignoring what’s actually good about the code and only producing criticism.} Why it matters: real reviews acknowledge good decisions; an interview review that’s 100\% negative, even of intentionally-buggy code, reads as a habit rather than a calibrated judgment for this specific snippet.

\item 
\textbf{Getting defensive or argumentative when the interviewer pushes back or asks “are you sure?”} Why it matters: pushback is often a genuine probe to see if you’ll re-examine your reasoning, not a challenge to win — treat it as an invitation to double-check, not an attack.

\item 
\textbf{Confusing “I would add a unit test for this” with actually explaining what the bug is.} Why it matters: proposing tests is good practice, but it doesn’t substitute for identifying and explaining the defect itself — interviewers will ask “so what’s actually wrong?” if you deflect into testing talk.

\item 
\textbf{Not distinguishing “this will definitely break” from “this could theoretically be an issue under conditions we haven’t confirmed.”} Why it matters: conflating certain bugs with speculative ones misleads the listener about urgency — say which is which.

\item 
\textbf{Trying to rewrite the entire snippet from scratch instead of making the minimal fix.} Why it matters: real reviews respect the existing design and diff size where possible; a wholesale rewrite in an interview eats your limited time and often introduces new bugs of its own.

\item 
\textbf{Forgetting to mention severity labels or any triage vocabulary at all.} Why it matters: explicit labels (blocker/major/minor/nit) are what make a review actionable for a team lead deciding whether to merge — omitting them makes your feedback harder to act on even if it’s technically correct.

\item 
\textbf{Not asking whether this code is even meant to run in production as-is, versus being a simplified teaching example.} Why it matters: some interview snippets are deliberately minimal and not every “issue” (e.g., missing logging, missing tests) is worth raising with the same weight as an actual correctness or safety bug — read the intent of the exercise.

\item 
\textbf{Spending so long on the review checklist ritual that you never get to the actual snippet.} Why it matters: the checklist in this chapter is a tool to internalize beforehand, not a script to recite out loud during the interview — using it as a mental scan, not a spoken monologue, keeps the conversation moving.

\item 
\textbf{Treating the interviewer as an adversary to be defeated rather than a colleague to collaborate with.} Why it matters: the entire point of the exercise is simulating a real code review conversation — the candidates who do best treat it exactly like one, asking questions, taking suggestions, and building toward the same goal as the interviewer: correct, safe, maintainable code.

\end{enumerate}


\end{document}